\documentclass[twoside,a5paper,12pt,fleqn,openany]{extbook}
\usepackage{polyglossia}
\usepackage{xcolor}
% ----- INCLUDE PDF AS PAGES -----
\usepackage{pdfpages}
% ----- INCLUDE PDF AS PAGES -----
% ----- GEOMETRY -----
\usepackage[left=1.5cm,right=1.5cm,top=2cm,bottom=2cm,bindingoffset=0.5cm]{geometry}
% ----- GEOMETRY -----



\usepackage[inline]{enumitem}

\usepackage[stable]{footmisc}
\usepackage{tocloft}

\renewcommand{\cftchapleader}{\cftdotfill{\cftdotsep}}
\setlength{\cftsecindent}{0pt}% Remove indent for \section
\setlength{\cftsubsecindent}{0pt}% Remove indent for \subsection



\newcommand{\ml}[3]{#2}
\newcommand{\coverfrontfilename}{cover_front.pdf}

\usepackage{tabularx}
\usepackage{multirow}
\usepackage{ltablex}

% ----- HYPHENATION -----
\usepackage{hyphenat}
% ----- HYPHENATION -----





% ----- FONTS -----
\renewcommand{\baselinestretch}{1.2}
\setmainfont{Linux Libertine}
% ----- FONTS -----

% ------ HYPERLINKS ------
\usepackage{hyperref}
\definecolor{LinkColor}{HTML}{0969DA}
\hypersetup{colorlinks=true, linkcolor=LinkColor, citecolor=LinkColor, filecolor=LinkColor, urlcolor=LinkColor}
% ------ HYPERLINKS ------

% ------ FANCY PAGE STYLE ------
\setlength{\headheight}{15pt}
\usepackage{fancyhdr}
\pagestyle{fancy}
\fancyhead[LE,RO]{\thepage}
\fancyhead[LO]{{\small{Polija XIX un XX gadsimtā}}}
\fancyhead[RE]{{\small{Vitālijs~Šalda}}}
\fancyfoot{}
    \fancypagestyle{plain}{
    \renewcommand{\headrulewidth}{0mm}
    \fancyhead{}
    \fancyfoot{}
}
% ------ FANCY PAGE STYLE ------

\setdefaultlanguage{latvian}
\setotherlanguages{english,russian,polish,german,french,lithuanian,spanish,latin,ukrainian,belarusian,greek,hebrew,romanian,italian,czech,slovene,swedish,serbian,hungarian,slovak}
\defaultfontfeatures{Ligatures=TeX,Mapping=tex-text}

% ------ ELEMENTS ------
\newcommand{\asterism}{\vspace{1em}{\centering\Large\bfseries$\ast~\ast~\ast$\par}\vspace{1em}}
\newcommand{\pltxti}[1]{\textit{\textpolish{#1}}}
\newcommand{\rutxti}[1]{\textrussian{#1}}
\newcommand{\detxti}[1]{\textit{\textgerman{#1}}}
\newcommand{\frtxti}[1]{\textit{\textfrench{#1}}}
\newcommand{\entxti}[1]{\textit{\textenglish{#1}}}
\newcommand{\lttxti}[1]{\textit{\textlithuanian{#1}}}
\newcommand{\estxti}[1]{\textit{\textspanish{#1}}}
\newcommand{\latxti}[1]{\textit{\textlatin{#1}}}
\newcommand{\betxti}[1]{\textbelarusian{#1}}
\newcommand{\uktxti}[1]{\textukrainian{#1}}
\newcommand{\rotxti}[1]{\textit{\textromanian{#1}}}
\newcommand{\ittxti}[1]{\textit{\textitalian{#1}}}
\newcommand{\cstxti}[1]{\textit{\textczech{#1}}}
\newcommand{\sltxti}[1]{\textit{\textslovene{#1}}}
\newcommand{\sktxti}[1]{\textit{\textslovak{#1}}}
\newcommand{\svtxti}[1]{\textit{\textswedish{#1}}}
\newcommand{\hutxti}[1]{\textit{\texthungarian{#1}}}
\newcommand{\srtxti}[1]{\textserbian{#1}}
\newcommand{\eltxti}[1]{\textgreek{#1}}
\newcommand{\hetxti}[1]{\texthebrew{#1}}
\newcommand{\mntxti}[1]{\textit{#1}}
\newcommand{\tttxti}[1]{\textit{#1}}
\newcommand{\zhtxti}[1]{#1}

\newcommand{\citespace}{[\dots{}]}

% ------ ELEMENTS ------

% ------ EPIGRAPH ------
\usepackage{epigraph}
\renewcommand{\epigraphsize}{\footnotesize}
\epigraphrule=0pt
\epigraphwidth=8cm
\setlength{\beforeepigraphskip}{.1\baselineskip}
\setlength{\afterepigraphskip}{.2\baselineskip}

% ------ EPIGRAPH ------

\addto\captionslatvian{\renewcommand{\contentsname}{Satura rādītājs}}

\begin{document}

% ----- FRONT COVER -----
\includepdf[pages={1}]{\coverfrontfilename}
% ----- FRONT COVER -----

% ----- EMPTY PAGE -----
\thispagestyle{plain}~
% ----- EMPTY PAGE -----

% ------ TITLE PAGE ------

\begin{titlepage}
{
\centering
{~\par}
\vspace{0.25\textheight}
{\Large{Vitālijs~Šalda}\par}
\vspace{1.3cm}
{\Huge\textbf{POLIJA}\par}
{\LARGE\textbf{XIX un XX gadsimtā}\par}
\vfill
{\includegraphics[width=6em]{publisher_logo.pdf}\par}
}
\end{titlepage}
% ------ TITLE PAGE ------

\epigraph
{Pagātne nav mirusi, tā nav pat pagājusi.}
{Viljams Folkners (\entxti{William Cuthbert Faulkner})}

\epigraph
{Meklējiet patiesību faktos.}
{Dens Sjaopins (\zhtxti{鄧小平})}

\epigraph
{Vēsturnieki nevar būt pilnīgi objektīvi.
Vēsture pēc savas būtības rada emocijas \citespace{}
Vēsture, izņemot neapšaubāmus faktus, pirmkārt balstās uz interpretāciju un zināmā mērā pieļauj daudzu patiesību esamību.}
{Žaks Le Goffs (\frtxti{Jacques Le Goff})}

\epigraph
{Vēsture nav skolotāja, bet uzraudze, magistra vitae: tā neko nemāca, bet tikai soda par mācībvielas nezināšanu.}
{Vasilijs Kļučevskis (\rutxti{Василий Осипович Ключевский})}

\epigraph
{Pagātne, saglabāta atmiņā, ir daļa tagadnes.}
{Tadeušs Kotarbiņskis (\pltxti{Tadeusz Marian Kotarbiński})}

\epigraph
{Bez īstas mīlestības pret cilvēci nav īstas dzimtenes mīlestības.}
{Anatols Franss (\frtxti{Anatole France})}



\epigraph
{Var atrast svētus cilvēkus, bet nav svētu politiķu, valdību, partiju, politisku un sabiedrisku kustību.}
{Voicehs Jaruzeļskis (\pltxti{Wojciech Witold Jaruzelski})}

\epigraph
{Nav jēgas rēķināties ar kaut kādu pasaules taisnīgumu visiem~--- viss ir atkarīgs no spēku samēra.}
{Aleksandrs Zinovjevs (\rutxti{Александр Зиновьев})}

\epigraph
{Ne viss vēsturē ir pelnījis, lai ar to lepotos un ar entuziasmu uz to atsauktos.}
{Voicehs Jaruzeļskis (\pltxti{Wojciech Witold Jaruzelski})}

\epigraph
{Katras valsts vēsturē atrodamas gan varonīgas, gan kaunpilnas lappuses~--- taču pats kaunpilnākais ir noklusēt vai lielīties ar to, no kā ir jākaunas.}
{Natans Eidelmans (\rutxti{Натан Яковлевич Эйдельман})}

\epigraph
{Tautas vēsture ir tautas raksturojums.}
{Mihails Vellers (\rutxti{Михаил Иосифович Веллер})}

\epigraph
{Vēsturiskā atmiņa nekādā gadījumā nav objektīva.}
{Ivars Austers}



\epigraph
{Patiesība ir subjektīva un vēstures izpratne tāpat.
Tā ir objektīva tikai tajā ziņā, ka vispār pastāv mūsu apziņā kā sen pagājušu notikumu pēdas.
Tāpēc veltas ir cerības uz jel kādu ``objektīvu vēsturi” vai ``vēstures tiesu”, kas nāks kā atpestīšana no mūžīgiem maldiem un noliks patiesību vispārējai apskatei.
Vēsture ir mūsu apziņa šeit šajā brīdī.
Rīt tā jau būs cita.}
{Āris Puriņš}

\epigraph
{Vēstures mūza ir lēnprātīga, zinoša un nepretencioza, bet, kad jūtas atstāta un pamesta, viņa cenšas atriebties un apžilbina tos, kas viņu atstāj novārtā.}
{Lešeks Kolakovskis (\pltxti{Leszek Kołakowski})}

\epigraph
{Vēsturnieku, kurš sagroza savu vēsturi, jāiznīcina tāpat kā naudas viltotāju tāpēc, ka viņi abi kaitē savai valstij.}
{Dons Migels de Servantess Saavedra\\(\estxti{Don Miguel de Cervantes Saavedra})}

\epigraph
{Tos, kuri saka taisnību, nemīl.
Tos, kuri melo~--- nicina.
Izvēle nav liela\ldots{}}
{Jurijs Poļakovs (\rutxti{Юрий Михайлович Поляков})}

\tableofcontents

\chapter*{Priekšvārds}
\addcontentsline{toc}{chapter}{Priekšvārds}

Strādājot Daugavpils Universitātes Vēstures katedrā, starp citiem studiju kursiem šī darba grāmatas autoram nācās lasīt arī lekciju kursu par Vidus un Austrumeiropas valstu vēsturi pēc Otrā pasaules kara.
Acīmredzot tāpēc, mācību programmām pilnveidojoties, Vēstures katedras toreizējais vadītājs asociētais profesors H.~Soms viņam uzdeva sagatavot arī Polijas XX gadsimta vēstures kursu.
Darba gaitā autoram izveidojās uzskats, ka jaunākos politiskos strīdus var saprast tikai tad, ja tiek ievērotas to vēsturiskās saknes; lai dziļāk izprastu XX~gadsimta Polijas vēstures notikumus, ir jādod plašāks ieskats par tās vēsturi XIX~gadsimtā nekā tas tiek sniegts vispārējā jauno laiku vēstures kursā.

Ievērojot to, ka bez profesoru Ē.~Jēkabsona un K.~Poča pētījumiem, kas saistīti ar Polijas vēsturi, latviešu valodā studentu rīcībā tikpat kā nav tai veltītu plašāku darbu, gadu gaitā tapa lasītāju priekšā esošais darbs.
Tas sākotnēji bija iecerēts kā mācību līdzeklis studentiem--vēsturniekiem.
Tomēr, ņemot vērā sabiedrībā pieaugošo interesi par mūsu valsts un tās kaimiņzemju vēsturi, kā arī lielo ar Polijas vēsturi saistīto materiālu apjomu, gala rezultātā ir tapis studentiem domāta mācību līdzekļa un zinātniski populāra darba apvienojums.
Autors cer, ka tādā veidolā tas atradīs lasītājus ne tikai studentu, bet arī citu vēstures interesentu vidū.

Izklāstā ir ietilpināts materiāls kā par jau lielā mērā izpētītiem, tā arī par daudziem vēl neatrisinātiem Polijas vēstures jautājumiem.
Nākotnē otro traktējums, iespējams, būs jāprecizē vai pat būtiski jāpārskata.
Tomēr autors uzskata par nepieciešamu dot priekšstatu arī par šādām problēmām.

Uzskatāmības labad autora tekstam pievienoti vairāku dokumentu tulkojumi un daudzas, galvenokārt INTERNET-ā aizgūtas ilustrācijas.
Lai lasītājiem būtu vieglāk orientēties, tekstā ar \strong{treknu druku} izcelti atsevišķi jēdzieni, problēmas, kuru iztirzājums tālāk seko.
Pieminētajām vēsturiskajām personām, arī pētniekiem un publicistiem darbam pievienotajā viņu sarakstā autors centās norādīt viņu dzīves gadus, kurus gan, diemžēl, ne vienmēr izdevās atrast.
Izklāstu ilustrējoši papildmateriāli nodrukāti citā šriftā un ar nedaudz sīkāku druku.

Bez izmantoto avotu un literatūras sarakstā uzrādītajiem izdevumiem grāmatas tapšanās gaitā samērā plaši izmantoti arī daudzi periodikā un INTERNET-ā atrodami materiāli, kuri nav norādīti bibliogrāfijas sarakstā.
Diemžēl tajos sastopamas daudzas neprecizitātes faktu izklāstā, klaji subjektīvas notikumu interpretācijas, kuras pārbaudīt pēc pirmavotiem ne vienmēr bija iespējams, jo daudzi no tiem autoram nebija pieejami.
Attēli ilustrācijām arī aizgūti galvenokārt INTERNET-ā.

Rakstot šo darbu, autoru ierobežoja arī poļu valodas zināšanu trūkums.
Tā kā poļu dokumentus un vēsturnieku darbus viņam nācās lasīt tulkojumos vācu un krievu valodā vai ar datortulkotāja palīdzību, nav izslēgtas zināmas neprecizitātes to izklāstā, varēja rasties arī kļūdas dažādu viedokļu interpretācijās.
Iespējams, ka faktu atspoguļojumu avotos, dažādu tautību un valstiskās piederības vēsturnieku, publicistu darbos atrodamo secinājumu, notikumu traktējumu, personāžu vērtējumu vidū autoram nav izdevies atrast visprecīzāko, vēstures gaitai atbilstošo atainojumu.
Viņš uzņemas visu atbildību par trūkumiem un kļūdām grāmatas saturā.

Autors arī saprot, ka daudzi viņa vērtējumi izsauks poļu radikālnacionālistu neapmierinātību, iespējams, arī prasību mīkstināt Polijas valdošo slāņu kritiku, bet uz to var viņiem ieteikt atcerēties apustuļu Pētera un Pāvila atbildi uz prasību pārtraukt sludināt Kristus mācību~--- ``\latxti{Non possumus!}” (``Mēs nevaram!”), kā arī jau XVI gadsimtā izskanējušos M.~Lutera vārdus: ``\detxti{Hier stehe ich, ich kann nicht anders}” (``Šeit es stāvu un citādi nevaru”).
Vēloties būt godīgs vēsturnieks, autors atļaujas sekot nosaukto dižo personību piemēram.

Darba pilnveidē daudz palīdzēja zinātniskā redaktore Dr.~hist., RISEBA vadošā pētniece T.~Bartele.
Tā tehniskā noformēšanā neaizstājama bija dabaszinātņu maģistra, ārsta Ņikitas Skabcova palīdzība.
Viņiem autors izsaka vissirsnīgāko pateicību.
Diemžēl objektīvu iemeslu dēļ darbam bija jāiztiek bez literārā redaktora un korektora.

\chapter*{Grāmatā minēto galveno organizāciju, iestāžu un to saīsinājumu saraksts}
\addcontentsline{toc}{chapter}{Grāmatā minēto galveno organizāciju, iestāžu un to saīsinājumu saraksts}

Abreviatūras (saīsinājumi) tekstā dotas, vadoties no partiju, organizāciju utt. nosaukumiem gan latviešu, gan poļu valodā.
Piemēram: Polijas Apvienotā strādnieku partija (\pltxti{Polska Zjednoczona Partia Robotnicza})~--- saīsinājumā saukta par PASP, nevis PZPR, jo latviešu lasītājs līdz 1990.~gadam izdotajā literatūrā un presē bieži sastapās ar pirmo, latviskoto saīsinājumu.
Taču dažādās valodās apritē jau iegājušie saīsinājumi tieši no poļu valodas atstāti arī šajā darbā.
Piemēram, \pltxti{Armia Krajowa} (burtiski~--- Dzimtenes armija) vācu valodā tiek tulkota kā ``\detxti{polnische Heimatarmee}'', bet angļu dokumentos tās nosaukums tulkots dažādi: ``\entxti{Home Army}'', ``\entxti{Underground Army}'', ``\entxti{Secret Army}'', kas īsti neatbilst nosaukumam poļu valodā.
Grāmatā lietots saīsinājums AK, nevis DA.
Polijas Sociālistiskā partija (\pltxti{Polska Partia Socjalistyczna}), kuras nosaukums literatūrā parasti tiek saīsināts kā PPS, šādi dots arī šajā grāmatā, nevis kā PSP, kā būtu vajadzējis saīsināt latviešu valodā.
Tāpat īsināti arī daļa citu nosaukumu.
Tā gan ir atteikšanās no vienotas pieejas saīsinājumu darināšanā, taču tā kā minētajiem un arī citiem saīsinājumiem vēstures literatūrā jau ir savs skanējums, to mēģināt mainīt acīmredzot būtu aplami.

Poļu vārds ``\pltxti{Polska}'' (poļu un Polijas), dažādos nosaukumos tulkots dažādi, galvenokārt atkarībā no organizācijas, kuras nosaukumā tas ietverts, izveides laika.
Ja Polija kā valsts organizācijas izveides laikā nepastāvēja, tad parasti lietots pirmais variants, ja pastāvēja, tad otrais.
Taču ir arī atkāpes no šī principa.
Piemēram, 1892.~gadā dibinātā ``\pltxti{Polska Partia Socjalistyczna}'' tiek tulkota kā Polijas (nevis Poļu) Sociālistiskā partija, arī 1942.~gada janvārī nodibinātā ``\pltxti{Polska Partia Robotnicza}'' tulkota kā Polijas Strādnieku Partija, jo šie nosaukumi jau ir iegājuši literatūrā.

Lai atvieglotu darba lasīšanu, dažādas organizācijas, iestādes, politiskās partijas grāmatas teksta blokos vispirms parasti sauktas
pilnos nosaukumus, abreviatūras izmantotas tikai pēc tam.
Retāk pieminētajām organizācijām tās nav lietotas vispār, taču lasītāju ērtības labad to nosaukums un tā tulkojums sarakstā parasti ir ietverts.

\begin{footnotesize}
\noindent
\begin{tabularx}{\linewidth}{|p{3cm}|p{3.5cm}|p{1.4cm}|p{1.6cm}|}
\hline
\strong{Organizācijas vai iestādes nosaukums latviešu valodā} & \strong{Organizācijas vai iestādes nosaukums poļu (angļu vai krievu) valodā} & \strong{Saīsinājums} & \strong{Organizācijas vai iestādes pastāvēšanas gadi} \\
\hline
Apvienotā Zemnieku partija & \pltxti{Zjednoczone Stronnictwo Ludowe} & AZP & 1949--1989 \\
\hline
Apvienoto Nāciju Organizācija & \pltxti{Organizacja Narodów Zjednoczonych} & ANO & 1945--mūsdienas \\
\hline
Apvienoto nāciju palīdzības un atjaunošanas organizācija & \entxti{United Nations Relief ana Rehabilitation Administration} & UNRRA & 1943--1947 \\
\hline
Aktīvās cīņas savienība & \pltxti{Zwiazek Walki Czynnej} & ACS & 1908--1914 \\
\hline
Brīvība un neatkarība (Pilns nosaukums: Pretošanās bez kara un diversijām „Brīvība un neatkarība”) & \pltxti{Wolność i Niezawisłość (Ruch Oporu bez Wojny i Dywersji „Wolność i Niezawisłość”)} & WiN & 1945--1952 \\
\hline
Bruņotās cīņas savienība & \pltxti{Związek Walki Zbrojnej} & BCS & 1939--1942 \\
\hline
Bruņoto spēku delegatūra & \pltxti{Delegatura Sił Zbrojnych na Kraj} & --- & 1945 \\
\hline
Centrisko spēku savienība & \pltxti{Porozumienie Centrum} & PC & 1990--1997 \\
\hline
Civīlmilicija & \pltxti{Milicja Obywatelska} & MO & 1944--1990 \\
\hline
Civīlmilicijas Brīvprātīgā rezerve & \pltxti{Ochotnicza Rezerwa Milicji Obywatelskiej} & ORMO & 1946--1989 \\
\hline
Darba partija & \pltxti{Stronnictwo Pracy} & SP & 1937--1950 \\
\hline
Demokrātiskā partija &\pltxti{Stronnictwo Demokratyczne} & SD & 1939--mūsdienas \\
\hline
Demokrātiskā savienība & \pltxti{Unia Demokratyczna} & UD & 1990--1994 \\
\hline
Demokrātisko kreiso spēku savienība & \pltxti{Sojusz Lewicy Demokratycznej} & SLD & 1991--1999 (partiju savienība), kopš 1999~--- partija \\
\hline
Eiropas Cilvēktiesību tiesa & \entxti{European Court of Human Rights} & ECT & 1998--mūsdienas \\
\hline
Galīcijas un Cešinas Silēzijas poļu sociāldemokrātiskā partija & \pltxti{Polska Partia Socjalno-Demokratyczna Galicji i Śląska Cieszyńskiego} & --- & 1897--1919 \\
\hline
Komunistiskās Internacionāles Izpildu Komiteja & \entxti{The Executive Committee of the Communist International}; \rutxti{Исполнительный комитет Коммунистического Интернационала} & KIIK & 1919--1943 \\
\hline
Krievijas Federācija & \rutxti{Российская Федерация} & KF & 1991--mūsdienas \\
\hline
Krievijas Komunistiskā (boļševiku) partija & \rutxti{Российская Коммунистическая партия (большевиков)} & KK(b)P & 1918--1925 \\
\hline
Krievijas Sociāldemokrātiskā strādnieku partija & \rutxti{Российская социал-демократическая рабочая партия} & KSDSP & 1898--1918 \\
\hline
Krievijas Sociālistiskā Federatīvā Padomju Republika & \rutxti{Российская Социалистическая Федеративная Советская Республика} & KSFPR & 1918--1922 \\
\hline
Kristīgi Nacionālā apvienība & \pltxti{Zjednoczenie Chrześcijańsko-Narodowe} & ZChN & 1989--2010 \\
\hline
Ķirzakas savienība & \pltxti{Związek Jaszczurczy} & --- & 1939--1942 \\
\hline
Liberāli demokrātiskais kongress & \pltxti{Kongres Liberalno-Demokratyczny} & KLD & 1991--1994 \\
\hline
Likums un taisnīgums & \pltxti{Prawo i Sprawiedliwość} & PiS & 2001--mūsdienas \\
\hline
Nacionālā Demokrātija & \pltxti{Narodowa Demokracja~--- Endecja} & ND & 1887--1947 \\
(Tās nosaukumi dažādos laika posmos: & ~ & ~ & ~ \\
Poļu līga & \pltxti{Liga Polska} & --- & 1887--1893 \\
Nacionālā Līga & \pltxti{Liga Narodowa} & --- & 1893--1928 \\
Nacionāli demokrātiskā partija & \pltxti{Stronnictwo Narodowo-Demokratyczne} & --- & 1897--1919 \\
Nacionāli demokrātiskā savienība & \pltxti{Związek Ludowo-Narodowy} & --- & 1919--1928 \\
Nacionālā partija & \pltxti{Stronnictwo Narodowe} & SN & 1928--1947 \\
\hline
Nacionālās atmiņas institūts & \pltxti{Instytut Pamięci Narodowej} & IPN & 2000--mūsdienas \\
\hline
Nacionālā strādnieku partija & \pltxti{Narodowa Partia Robotnicza} & NPR & 1920--1937 \\
\hline
Nacionālās vienotības padome & \pltxti{Rada Jednosci Narodowej} & --- & 1944--1945 \\
\hline
Nacionālās vienotības pagaidu valdība & \pltxti{Tymczasowy Rząd Jedności Narodowej} & NVPV & 1945--1947 \\
\hline
Nacionālie bruņotie spēki & \pltxti{Narodowe Siły Zbrojne} & NBS & 1942--1947 \\
\hline
Nacionālie bruņotie spēki~--- Ķirzakas savienība & \pltxti{Narodowe Siły Zbrojne~--- Związek Jaszczurczy} & --- & 1944--1945 \\
\hline
Padomju Krievijas / PSRS Tautas komisāru padome & \rutxti{Совет народных комиссаров} & TKP & 1917--1946 \\
\hline
Padomju Savienības Komunistiskā partija & \rutxti{Коммунистическая партия Советского Союза} & PSKP & 1952--1991 \\
\hline
Padomju Savienības Telegrāfa aģentūra & \rutxti{Телеграфное агентство Советского Союза} & TASS & 1925--1991 \\
\hline
Pilsoniskā kustība~--- demokrātiskā darbība & \pltxti{Ruch Obywatelski Akcja Demokratyczna} & --- & 1990--1991 \\
\hline
Polijas Apvienotā strādnieku partija & \pltxti{Polska Zjednoczona Partia Robotnicza} & PASP & 1948--1990 \\
\hline
Polijas armija & \pltxti{Wojsko Poliskie} & --- & 1918--1989 \\
\hline
Polijas Jaunatnes savienība & \pltxti{Związek Młodzieży Polskiej} & PJS & 1948--1957 \\
\hline
Polijas karalistes un Lietuvas Sociāldemokrātiskā partija & \pltxti{Socjaldemokracja Królestwa Polskiego i Litwy} & PKunLSD & 1900--1918 \\
\hline
Polijas karalistes Sociāldemokrātija & \pltxti{Socjaldemokracja Królestwa Polskiego} & PKSD & 1893--1900 \\
\hline
Polijas Komunistiskā strādnieku partija & \pltxti{Komunistyczna Partia Robotnicza Polski} & PKSP & 1918--1925 \\
\hline
Polijas Komunistiskā partija & \pltxti{Komunistyczna Partia Polski} & PKP & 1925--1938 \\
\hline
Polijas Kristīgo demokrātu partija & \pltxti{Polskie Stronnictwo Chrześcijańskiej Demokracji}; \pltxti{Chrześcijańska Demokracja} & --- & 1919--1937 \\
\hline
Polijas Nacionālās atbrīvošanas komiteja & \pltxti{Polski Komitet Wyzwolenia Narodowego} & PNAK & 1944 \\
\hline
Polijas Republikas Nacionālā Padome & \pltxti{Rada Narodowa Rzeczypospolitej Polskiej} & --- & 1939--1941, 1942--1945 \\
\hline
Polijas Republikas pagaidu valdība & \pltxti{Rząd Tymczasowy Rzeczypospolitej Polskiej} & PRPV & 1945 \\
\hline
Polijas Republikas pilsoniskā platforma & \pltxti{Platforma Obywatelska Rzeczypospolitej Polskiej} & PO & 2001--mūsdienas \\
\hline
Polijas Republikas Sociāldemokrātija & \pltxti{Socjaldemokracja Rzeczypospolitej Polskiej} & PRSD & 1991--1999 \\
\hline
Polijas Sociālistiskā partija & \pltxti{Polska Partia Socjalistyczna} & PPS & 1892--1948 \\
\hline
Polijas Sociālistiskā partija~--- ļevica & \pltxti{Polska Partia Socjalistyczna~--- Ļevica} & PPS-ļevica & 1906--1918, 1926--1931 \\
\hline
Polijas Sociālistiskā partija~--- revolucionārā frakcija & \pltxti{Polska Partia Socjalistyczna~--- Frakcja Rewolucyjna} & PPS-frakcija & 1906--1919 \\
\hline
Polijas Strādnieku Partija & \pltxti{Polska Partia Robotnicza} & PSP & 1942--1948 \\
\hline
Polijas Tautas armija & \pltxti{Polska Armia Ludowa} & PAL & 1943--1944 \\
\hline
Polijas Tautas republika & \pltxti{Polska Rzeczpospolita Ludowa} & PTR & 1944--1989 \\
\hline
Polijas uzvaras dienests & \pltxti{Służba Zwycięstwu Polski} & --- & 1939 \\
\hline
Polijas valdība emigrācijā & \pltxti{Rząd Rzeczypospolitej Polskiej na uchodźstwie} & PVE & 1939--1945 (formāli~1990) \\
\hline
Polijas Zemnieku partija & \pltxti{Polskie Stronnictwo Ludowe} & PSL & 1945--1949, 1990--mūsdienas \\
\hline
Polijas Zemnieku partija~--- Jaunā atbrīvošanās & \pltxti{Polskie Stronnictwo Ludowe~--- Nowe Wyzwolenie} & PSL-NW & 1946--1947 \\
\hline
Polijas Zemnieku partija~--- kreisā frakcija & \pltxti{Polskie Stronnictwo Ludowe~--- ļewica} & PSL-levica & 1914--1924 \\
\hline
Politiskā konsultatīvā komiteja & \pltxti{Polityczny Komitet Porozumiewawczy} & --- & 1940--1943 \\
\hline
Poļu Kara organizācija & \pltxti{Polska Organizacja Wojskowa} & POW & 1914--1918 \\
\hline
Poļu Komunistu Centrālais Birojs & \pltxti{Centralne Biuro Komunistów Polskich} & PKCB & 1944 \\
\hline
Poļu patriotu savienība & \pltxti{Związek Patriotów Polskich} & PPtS & 1943--1946 \\
\hline
Poļu Sociālistiskā partija~--- Brīvība, Vienlīdzība, Neatkarība & \pltxti{Polska Partia Socjalistyczna~--- Wolność, Równość, Niepodległość} & PPS-WRN & 1939--1944 \\
\hline
Poļu sociālistu strādnieku partija & \pltxti{Rabotnicza Partia Polskich Socjalistów} & PSSP & 1943--1944 \\
\hline
Poļu sociālistiskā partija Prūsijā & \pltxti{Polska Partia Socjalistyczna Zaboru Pruskiego} & --- & 1893--1919 \\
\hline
Poļu Zemnieku partija „Atbrīvošanās” & \pltxti{Polskie Stronnitctwo Ludowe~--- „Wyzwolenie”} & PSL „\pltxti{Wyzwolenie}” & 1915--1931 \\
\hline
Poļu Zemnieku partija~--- Pjasts & \pltxti{Polskie Stronnictwo Ludowe~--- Piast} & PSL-Piast & 1913--1931 \\
\hline
PSRS Iekšlietu tautas komisariāts & \rutxti{Народный комиссариат внутренних дел СССР} & IeTK & 1934--1946 \\
\hline
PSRS Valsts Aizsardzības Komiteja & \rutxti{Государственный комитет обороны СССР} & PSRS VAK & 1941--1945 \\
\hline
Rietumbaltkrievijas Komunistiskā partija & baltkr.: \betxti{Камуністычная партыя Заходняй Беларусі} & RBKP & 1923--1938 \\
\hline
Rietumukrainas Komunistiskā partija & ukr.: \uktxti{Комуністична партія Західної України} & RUKP & 1919--1938 \\
\hline
Sabiedriskās drošības ministrija & \pltxti{Ministerstwo Bezpieczeństwa Publicznego} & SDM & 1945-1954 \\
\hline
Savstarpējās Ekonomiskās Palīdzības Padome & \rutxti{Совет экономической взаимопомощи} & SEPP & 1949--1991 \\
\hline
Starptautiskā Sarkanā Krusta komiteja & \entxti{International Committee of the Red Cross} & SSKK & 1863--mūsdienas \\
\hline
Strādnieku aizsardzības komiteja & \pltxti{Komitet Obrony Robotników} & KOR & 1976--1977 \\
\hline
Strādnieku aizsardzības komitejas Sociālās pašaizsardzības komiteja & \pltxti{Komitet Samoobrony Społecznej KOR} & KOS-KOR & 1977--1981 \\
\hline
Strādnieku Zemnieku Sarkanā armija & \rutxti{Рабоче-крестьянская Красная армия} & SZSA & 1918--1946 \\
\hline
Tautas armija & \pltxti{Armia Ludowa} & AL & 1944 \\
\hline
Tautas gvarde & \pltxti{Gwardia Ludowa} & GL & 1942--1943 \\
\hline
Ukraiņu nacionālistu organizācija & ukr.: \uktxti{Організація Українських Націоналістів} & OUN & 1929--mūsdienas \\
\hline
Ukrainas Sacelšanās armija & ukr.: \uktxti{Українська повстанська армія} & UPA & 1942--1953 \\
\hline
Ukrainas Tautas Republika & ukr.: \uktxti{Українська Народна Республіка} & UTR & 1917--1920 \\
\hline
Vācijas Demokrātiskā Republika & vācu: \detxti{Deutsche Demokratische Republik} & VDR & 1949--1990 \\
\hline
Vācijas Federatīvā Republika & vācu: \detxti{Bundesrepublik Deutschland} & VFR & 1949--mūsdienas \\
\hline
Valsts Nacionālā padome & \pltxti{Krajowa Rada Narodowa} & KRN & 1944--1947 \\
\hline
Valsts Politiskā pārstāvniecība & \pltxti{Krajową Reprezentację Polityczną} & --- & 1943--1944 \\
\hline
Vēlēšanu akcija Solidaritāte & \pltxti{Akcja Wyborcza Solidarność} & AWS & 1996--2003 \\
\hline
Viskrievijas Centrālā izpildu komiteja & \rutxti{Всероссийский Центральный Исполнительный Комитет} & VCIK & 1917--1938 \\
\hline
Vispolijas arodbiedrību savienība & \pltxti{Ogólnopolskie Porozumienie Związków Zawodowych} & VAS & 1984--mūsdienas \\
\hline
Vissavienības Komunistiskā (boļševiku) partija & \rutxti{Всесоюзная коммунистическая партия большевиков} & VK(b)P & 1925--1952 \\
\hline
Zemes armija & \pltxti{Armia Krajowa} & AK & 1942--1945 \\
\hline
Zemes politiskā pārstāvniecība & \pltxti{Krajowa Reprezentacja Polityczna} & --- & 1943--1944 \\
\hline
Zemnieku bataljoni & \pltxti{Bataliony Chłopskie} & BCh & 1940--1945 \\
\hline
Zemnieku partija & \pltxti{Stronnictwo Chłopskie} & SCh & 1926--1931 \\
\hline
Zemnieku partija „Roch” & \pltxti{Stronnictwo Ludowe „Roch”} & SL Roch & 1940--1945 \\
\hline
Zemnieku partija & \pltxti{Stronnictwo Ludowe} & SL & 1931--1939, 1944--1949 \\
\hline
Zemnieku partija „Tautas griba” & \pltxti{Stronnictwo Ludowe „Wola Ludu”} & SL-WL & 1943--1944 \\
\hline
Zemnieku pašpalīdzības savienība & \pltxti{Związek Samopomocy Chłopskiej} & ZPS & 1944--1957 \\
\hline
\end{tabularx}
\end{footnotesize}

\chapter*{Ievads}
\addcontentsline{toc}{chapter}{Ievads}

\epigraph
{Īsts patriots ir tas, kurš saka patiesību pat savai zemei.}
{Žans Žoress (\frtxti{Jean Léon Jaurès})}

\epigraph
{Tēvzemes mīlestība ir ļoti laba lieta, bet ir kas vēl augstāks pār to~--- patiesības mīlestība.}
{Pēteris Čaadajevs (\rutxti{Пётр Яковлевич Чаадаев})}

\epigraph
{Ir labi būt patiesam vienmēr, arī tad, ja tas skar dzimteni. Katra pilsoņa pienākums ir, ja vajadzīgs, mirt par savu dzimteni, bet nevienu nedrīkst piespiest melot dzimtenes vārdā.}
{Šarls Luijs de Monteskjē (\frtxti{Charles-Louis de Secondat, Baron de La Brède et de Montesquieu})}

\epigraph
{Uzbrukumi tautas trūkumiem un netikumiem nav noziegums, bet nopelns, īstens patriotisms.}
{Visarions Beļinskis\\(\rutxti{Виссарион Григорьевич Белинский})}



\epigraph
{Īsts patriotisms~--- tas pirmkārt ir vēlēšanās zināt un runāt patiesību par savu Tēvzemi.}
{Vjačeslavs Kostikovs (\rutxti{Вячеслав Васильевич Костиков})}

\epigraph
{Visas nelaimes pasaulē ceļas tādēļ, ka cilvēki lieto divējādas mērauklas: vienu sev un savai dzimtai, citu~--- pārējiem.}
{Zenta Mauriņa}

\epigraph
{Es dodu priekšroku savas dzimtenes šaustīšanai, tās sarūgtināšanai, tikai lai nekrāptu to.}
{Pēteris Čaadajevs (\rutxti{Пётр Я́ковлевич Чаада́ев})}

\epigraph
{Jebkurš nacionālisms ir akls, jo meklē savu nelaimju cēloņus citās tautās, tā sākas neiecietība, netaisnība. Nacionālisms nodara ļaunumu vispirms tā paudējiem, jo nenovēršami izraisa pretreakciju no citu tautu puses.}
{Ilga Apine}

\epigraph
{Visa mūsu vēsture~--- tas ir izdomājums, kuram visi piekrīt.}
{Voltērs (\frtxti{Voltaire, Fransuā Marī Aruē})}



\epigraph
{Kas gan ir vēsture, ja ne meli, ar kuriem visi ir vienisprātis?}
{Napoleons Bonaparts (\frtxti{Napoléon Bonaparte})}

\epigraph
{Gala rezultātā mēs pētām vēsturi lai apmierinātu savas intereses un, pēc iespējas, saprastu pie tam savas paša problēmas. Taču nevienu no šiem diviem mērķiem mēs nesasniegsim, ja, atrodoties neauglīgās zinātniskās objektivitātes idejas iespaidā, neuzdrošināsimies zinātniskās problēmas izskatīt no sava skatu punkta.}
{Kārlis Poppers (\entxti{Karl Raimund Popper})}

\epigraph
{Nevar būt vēstures „pagātnei, kāda tā patiesi pastāvēja”, ir iespējamas tikai vēsturiskas interpretācijas, un ne viena no tām nav galīga.}
{Kārlis Poppers (\entxti{Karl Raimund Popper})}

\epigraph
{Vieni no lielākajiem maldiem ir domāt, ka visi jūt, redz un domā tāpat, kā mēs.}
{Pjērs Buasts (\frtxti{Pierre Boiste})}

\epigraph
{Lai kā pieklājas rakstītu vēsturi, ir jāaizmirst par savu ticību, savu tēvzemi, savu partiju.}
{Pjērs Buasts (\frtxti{Pierre Boiste})}



\epigraph
{Mūsu ienaidnieku spriedumi par mums ir tuvāk patiesībai, nekā mūsu pašu.}
{Fransuā de Larošfuko (\frtxti{François de La Rochefoucauld})}

\epigraph
{Ticība tam, ka ir tikai viena, tikai tev zināma patiesība, ir visa ļaunuma sakne.}
{Makss Borns (\entxti{Max Born})}

\epigraph
{Patiesība pasaulē nedara tik daudz laba, cik šķietamas patiesības~--- ļauna. Tieksme izskaistināt vēsturi ir instinktīva tieksme paaugstināt grupas pašnovērtējumu.}
{Mihails Vellers (\rutxti{Михаил Иосифович Веллер})}

\epigraph
{Nobriedis prāts ienīst savas valsts netikumus~--- nepilnīgs noliedz tos. Godprātīgs vēsturnieks redz pienākumu izskaust savas zemes netikumus~--- negodprātīgs grib tikai iznīcināt tās ienaidniekus.}
{Mihails Vellers (\rutxti{Михаил Иосифович Веллер})}

\epigraph
{Poļu tautas rakstura varonīgās īpašības nedrīkst mūs piespiest aizvērt aci uz tās neprātīgumu un nepateicību, kas gadsimtu garumā tai radīja neizmērojamas ciešanas.}
{Vinstons Čērčils (\entxti{Winston Leonard Spencer-Churchill})}



\epigraph
{Katru nacionālistu vajā ideja, ka pagātni ne vien var izmainīt, bet tā ir jāizmaina.}
{Džordžs Orvels (\entxti{George Orwell})}

\epigraph
{Tas, ko mēs saucam par cinismu, bieži vien ir atklāti pateikta patiesība.}
{Jezups Laganovskis}

\epigraph
{Ciniķis ir cilvēks, kurš nebaidās pateikt to, ko citi domā.}
{Aleksandrs Gordons (\rutxti{Александр Гордон})}

\epigraph
{Polijas nelaime ir tā, ka tās vēsture ir tik melīga, kā nekas cits pasaulē.}
{Ježi Gedroics (\pltxti{Jerzy Władysław Giedroyc})}

\epigraph
{Atcerēsimies, ka Polija citu tautu vidū drīzāk ir liela pele, nekā mazs zilonis.}
{Tadeušs Kotarbiņskis (\pltxti{Tadeusz Marian Kotarbiński})}

\epigraph
{Kāda Centrāleiropa? Rietumāzija!}
{Josifs Brodskis (\rutxti{Иосиф Александрович Бродский})}

Polijai vienmēr ir bijusi nozīmīga vieta Eiropas vēsturē, ko noteica gan tās atrašanās pašā Eiropas centrā, gan iespaidīgais iedzīvotāju skaits, gan arī aktīvā politika. Daudzi svarīgi notikumi ir risinājušies Polijas teritorijā, poļu tauta ir bijusi un joprojām ir iesaistīta ne tikai Eiropas, bet arī visas pasaules mēroga procesos. Poļu sabiedrisko kustību, Polijas valsts politikas iedarbība bija un ir jūtama tālu aiz tās robežām, bet īpaši tajās valstīs, kuras atrodas tai kaimiņos, tai skaitā arī Latvijā.

\strong{Poļi} (pašnosaukums \pltxti{Polacy}) ir skaitliski lielākā rietumslāvu tauta. Pēc dažādiem vērtējumiem pasaulē dzīvo no 44 līdz 60 miljonu poļu, no tiem Polijas Republikā~--- 38 miljoni. Nosaukums radies no rietumslāvu cilts (maztautas) poļānu (\pltxti{polanie}) vārda, kas jau no VIII gadsimta dzīvoja Vartas (poļu \pltxti{Warta}) upes baseinā rajonā ap mūsdienu Poznaņu (poļu \pltxti{Poznań}, vācu \detxti{Posen}). Vārda \pltxti{polanie} tulkojums ir vienkārši „lauku (tīrumu) iedzīvotāji” (\pltxti{pole}~--- lauks, tīrums). Tieši poļānu vidū konsolidējās poļu tautas kodols, šeit IX gadsimtā radās pirmā Polijas valsts Pjastu (\pltxti{Piastowie}) dinastijas (10.--14.gs.) vadībā, un šeit atradās pirmās Polijas galvaspilsētas~--- Gņezno (\pltxti{Gniezno}), Poznaņa (\pltxti{Poznań}). Konsolidējot ap sevi citas rietumslāvu ciltis Vislas baseinā, pamazām veidojās kopīga poļu identitāte.

16.~gadsimta vidū poļi kopā ar lietuviešiem un slāvu ciltīm~--- mūsdienu ukraiņu un baltkrievu priekštečiem~--- radīja Polijas-Lietuvas apvienotu valsti jeb Abu Tautu Republiku (1569--1795). To parasti sauca par Žečpospolitu (poļu \pltxti{Rzeczpospolita}, no latīņu valodas: \latxti{Res publica}~--- republika. Pilns nosaukums: poļu: \pltxti{Rzeczpospolita Obojga Narodów}, lietuv.~--- \lttxti{Žečpospolita/Abiejų Tautų Respublika}, baltkr.~--- \betxti{Рэч Паспалітая Абодвух Народаў}). Tā bija muižnieku republika ar vēlētu monarhu priekšgalā, kur poļi bija valdošā tautība.

Mūsdienu poļu historiogrāfijā poļu valstiskuma dažādiem modeļiem pieņemts dot kārtas numurus: 1569.--1795.~gadu periods tiek saukts par I Žečpospolitu, no 1918. līdz 1939.~gadam~--- par II Žečpospolitu, no 1989.~gada~--- par III Žečpospolitu. Polijas Tautas Republikas pastāvēšanas periodu (1944--1989) šajā numerācijā neievēro vai arī lieto terminu \pltxti{Rzeczpospolita Ludowa} (Tautas republika), lai uzsvērtu tās atšķirību no citām poļu republikām.

Sāka veidoties poļu nācija. Taču tās izveide noslēdzās XVIII~gadsimta otrajā pusē un XIX~gadsimtā jau smagākos apstākļos, kad poļi cīnījās par atbrīvošanos no nacionālās atkarības.

Pakļautais Polijas stāvoklis attiecībās ar citām valstīm iezīmējās jau XVIII~gadsimta sākumā. Tā otrajā pusē Žečpospolita ekonomiskā un politiskā iekšējā vājuma, šļahtas (muižniecības) ķīviņu, etnisko nesaskaņu rezultātā tika trīs reizes kaimiņvalstu sadalīta un beidza pastāvēt. XIX gadsimtā tikai atsevišķos periodos Polijā pastāvēja valstiski veidojumi ar dažādu autonomijas pakāpi (Varšavas hercogiste 1807--1814, Polijas karaliste 1815--1830, Krakovas Republika 1815--1846, autonomā Galīcija pēc 1867.~gada). Šo veidojumu esamība ļāva kaut nelielā mērā saglabāt poļu valstiskuma pārmantojamu, tomēr neapmierināja radikālāk noskaņotos poļu patriotus. Viņi nepārtraukti cīnījās par savu neatkarību, vairākkārt sacēlās pret lielvalstu kundzību, taču tas tikai pasliktināja zemes stāvokli. Atbrīvošanās kustība izplatījās visās poļu zemēs, piepildīja augošo nacionālo pašapziņu ar cīņas apofeozi, bruņotas pretošanās, kritušo un nomocīto nacionālo varoņu kultu. Valsts ar 800~gadu vēsturi sadale ienesa arī pastāvīgu nemiera elementu pārējo Eiropas valstu attiecībās.

Īpaši bagāta sarežģījumiem bija Polijas vēsture XX gadsimtā. Polijas valstiskuma atjaunošana (1918) un Versaļas miers (1919) nekļuva par galīgu Polijas jautājuma risinājumu. Jaltas vienošanās (1945) atkal pakļāva Poliju lielvalstu diktātam. Polija ilgstoši atradās t.s. „sociālistiskā” bloka sastāvā, faktiskā atkarībā no PSRS. Tikai 1989.~gadā poļu tauta atkal atguva neatkarību.

Kā raksta Polijas vēsturi daudz pētījušais britu profesors N.~Deiviss: „Polijas nacionālajai kustībai bija vissenākā vēsture, visstiprākais mandāts, vislielākā apņēmība, vissliktākā reputācija un vismazākie panākumi”. Ir saprotams, ka atzīmētās pretrunas ir skaidrojamas gan ar objektīviem, gan subjektīviem, gan starptautiskiem, gan pašu poļu darbības radītiem faktoriem. Bagātās, bet sarežģītās vēstures dažāda uztvere pašu poļu vidū arī pati par sevi ir radījusi problēmas.

Ungāru vēsturnieks I.~Bibo ir atzīmējis īpašu Austrumeiropas tautu kolektīvu psiholoģijas iezīmi~--- sevišķu jūtīgumu, pat bailes no pastāvošām vai tikai šķietamām nacionālās kopības bojā ejas briesmām. Šīm tautām šīs bailes bija saistītas sākotnēji ar turku, tad vācu, dažkārt arī poļu ekspansiju, pēc tam krievu iespiešanos reģionā. Pēdējā gadsimtā šīs bailes galvenokārt saistījās ar Vāciju, Krieviju un PSRS. Lielajām tautām šīs bailes ir mazsaprotamas. Piemēram, krievu tauta ilgstoši ir bijusi apspiesta, bet parasti apspieda viņu savas valsts vara, kuru viņi neuztvēra kā etniski svešu, ja arī imperatoru dzīslās tecēja faktiski vācu asinis. Tāpēc daudzi krievi pagātnē un bieži arī mūsdienās nevar pilnībā saprast tās tautas, kuras baidās galvenokārt no etniski svešiem apspiedējiem, pat ja šo apspiedēju varā dzīvot bija labāk nekā savu valdītāju kundzībā. Tā, nevar strikti apgalvot, ka vairākumam poļu Prūsijas/Vācijas un Austrijas/Austroungārijas kundzībā dzīvot bija sliktāk nekā savu poļu „panu”~--- savas šļahtas varā līdz Žečpospolitas dalīšanai. Taču liela poļu daļa tās pašas šļahtas vadībā XIX~gadsimtā cīnījās pret nacionālo apspiestību, sava valstiskuma atjaunošanu ne tikai pret „aziātisko” Krievijas impēriju, bet arī „eiropeiskajām” Prūsiju/Vāciju un Austriju/Austroungāriju.

\strong{Polijas attiecības ar kaimiņvalstīm} un to tautām vienmēr ir bijušas neviennozīmīgas. Vissarežģītākās attiecības Polijai gandrīz vienmēr veidojās ar kaimiņos esošajām Vāciju un Krieviju.

Vispirms par \strong{Polijas attiecībām ar Vāciju} (\pltxti{Niemcy}). Vācieši ilgstoši ir bijuši vieni no tuvākajiem un arī bīstamākajiem poļu kaimiņiem. Poļi „vācieti” (\pltxti{niemiec}) gadu simteņiem uztvēra kā ienaidnieku.

1217.~gadā pāvests Honorijs III pasludināja krusta karu pret prūšu pagāniem, ar kuriem karoja Mazovijas (poļu \pltxti{Mazowsze, Mazowiecka ziemia}, vācu ~\detxti{Masowien}) hercogs Konrāds. 1225.~gadā hercogs lūdza palīdzību Teitoņu (Vācu) ordenim (latīņu \latxti{Ordo fratrum domus Sanctae Mariae Theutonicorum Ierosolimitanorum}, \latxti{Ordo Teutonicus}, vācu \detxti{Orden der Brüder vom Deutschen Haus St. Mariens in Jerusalem}, saīsināti: \detxti{Deutscher Orden}, \detxti{Deutscher Ritterorden}~--- katoļu reliģiskais bruņinieku ordenis, dibināts XII gadsimtā Palestīnā). Teitoņi ieradās Polijā 1232.~gadā, apmetās Vislas labajā krastā, sāka sagrābt prūšu zemes, pašus tos piespiežot pieņemt kristietību. Iekarotajās zemēs ieradās vācu kolonisti. Ordenis uzurpēja tiesības kristīt pagānus austrumos, noliedzot tādas pat tiesības poļiem. Galvenais ordeņa pretinieks bija Lietuvas lielkņaziste (lietuviešu \lttxti{Lietuvos Didžioji Kunigaikštystė}, poļu \pltxti{Wielkie Księstwo Litewskie}, krievu \rutxti{Великое княжество Литовское, Русское, Жемойтске и иных}; XIII~gs.--1796.). Lai atvairītu teitoņu triecienus, lielkņaziste 1385.~gadā Krēvas pilī (baltkrievu \betxti{Крэўскі замак}, poļu \pltxti{Zamek w Krewie}, lietuviešu \lttxti{Krėvos pilis}, mūsdienu Baltkrievijas teritorijā) noslēdza Polijas un Lietuvas personālo ūniju, kad abām valstīm bija viens valdnieks no Jagelloņu (lietuv.: \lttxti{Jogailaičiai}, poļu: \pltxti{Jagiellonowie}, baltkr.: \betxti{Ягелоны)} dinastijas. Tas mainīja spēku samēru par sliktu teitoņiem. Taču arī pēc Vācu ordeņa novājināšanās, Reformācijas un ordeņa zemju sekularizācijas vācu un poļu valstu attiecības neuzlabojās.

Gan vācieši, gan poļi kopš viduslaikiem netaupīja negatīvās emocijas pret kaimiņiem.

Jau Prūsijas karalis Fridrihs II, apkopojot poļiem naidīgos uzskatus, izteicās, ka sabiedrība, „kuras uzvārdi beidzas ar --ki”, ir „visās nozīmēs nicināma nācija'' (\detxti{Die ganze „Gesellschaft mit dem Namen auf~--- ki” ist „eine in jeder Hinsicht verächtliche Nation”}.) Prūsijas vēsturnieki poļus un citus slāvus apzīmēja par „barbarisma apustuļiem” („\detxti{Apostel der Barbarei}”). Oficiālais Prūsijas valsts historiogrāfs L.~Ranke uzskatīja, ka romāņu-ģermāņu tautām vienmēr ir bijusi noteicošā loma Eiropas vēsturē, bet kas attiecās uz slāviem, to loma bija aizstāvēt „Eiropas civilizāciju” no klejotāju uzbrukumiem. Šis zinātnieks uz slāviem, tai skaitā poļiem, raudzījās no augšas, uzsverot viņu kulturālo atpalicību un Vācijas civilizējošo lomu to vidū. Arī viņa pēctecis amatā H.~Treičke ģermāņu virzību slāvu zemēs nosauca par „kultūras tautu cīņu pret barbariem”. Vācu vēstures grāmatās Polijas vēsture tika reducēta uz anarhijas, iekšējo krīžu un cīņas pret kaimiņiem attēlojumu. Pozitīvie procesi, sasniegumi saimnieciskajā un kultūras jomā nebija pienācīgi novērtēti. (Arī mūsdienās vācu vēsturiskajā literatūrā var atrast šo stereotipu atliekas.) Polijas vēstures materiālam Vācijas valdošo slāņu acīs bija jākalpo vācu imperiālisma ekspansionistisko centienu pamatojumam.

Savukārt poļu vēsturnieki jau no XIX gadsimta oponēja prūsiskajai historiogrāfijai. Šajā pretstāvē piedalījās arī poļu literatūra un māksla, negatīvās krāsās attēlojot vācu pagātni. Viedokļu pretrunīgumu labi ilustrē divi mīti par kauju pie Grīnvaldes~--- Tannenbergas (poļu \pltxti{Bitwa pod Grunwaldem}, vācu \detxti{Schlacht bei Tannenberg}, 1410).

Pēc poļu mīta Teitoņu ordenis falsificēja Romas pāvesta bullu, pēc tam veica karagājienus pret prūšiem, Lietuvu un Poliju, 1308.~gadā ar viltu sagrāba Gdaņsku un Piejūras apgabalu, mēģināja pakļaut Lietuvu, kuru izglāba 1385.~gada personālā ūnija ar Poliju. Kauja pie Grīnvaldes bija mūžīgo Polijas cīņu pret Vāciju kvintesence. Pēc uzvaras Grīnvaldes kaujā Polijā valdījušās Jagelloņu dinastijas kļūda bija atļauja no krustnešu vasaļvalsts radīt mantojamu hercogisti, kura izveidojās par Prūsiju un vēlāk piedalījās Žečpospolitas dalīšanā. Arī Krievijas impērijā ietilpinātajā Polijā 1910.~gadā tika plaši atzīmēta Grīnvaldes kaujas jubileja.

Berlīnes ieņemšana 1945.~gadā bija otra Grīnvalde, Austrumprūsijas pievienošana Polijai uz visiem laikiem pielika punktu teitoņu agresijai. Pēc 1945.~gada Polijas jauniegūtajās rietumu zemēs laikam nebija pilsētas, kurā nebūtu savas Grīnvaldes ielas. Kaut oficiālā Polijas Tautas Republikas propaganda veidoja divus vāciešu tēlus: Poliju apdraudošus, atklāti agresīvus VFR pilsoņus, un sabiedrotos no VDR, jāsaka gan, ka katoliskie poļu iedzīvotāji juta maz simpātiju pret pēdējiem, kurus parasti identificēja kā prūšus.

Arī jau mūsdienās populāro poļu politiķu brāļu Ļeha un Jaroslava Kačinsku pirmsvēlēšanu kampaņa 2005.~gadā tika uzsākta uz J.~Matejko gleznas „Grīnvaldes kauja” fona.

Pilnīgi pretējs ir vācu mīts. Teitoņu ordenis neauglīgajās Eiropas ziemeļaustrumu zemēs, pārvarot pagānu pretestību, ar krustu un zobenu iedibināja kristietību, piesaistīja sava laikmeta modernās tehnoloģijas, nodibināja 96 pilsētas, uzcēla 90 cietokšņus, radīja sava laikmeta paraugvalsti, kuras mantinieces bija Prūsija un Vācijas impērija. (Vēl nacistiskajā t.s. Trešajā reihā darbojās Teitoņu ordeņa kults.) Nedienas sākās tad, kad Polija, jūtot skaudību pret teitoņu valsts panākumiem, noniecināja kristīgo solidaritāti, vienojās ar pagāniem. 1410.~gadā tā sāka karu pret ordeni un diemžēl kaujā pie Tannenbergas, galvenokārt pateicoties lietuviešu spēkiem, sakāva to. Taču vēsturiskais taisnīgums uzvarēja~--- XVII gadsimtā Prūsijas karaļiem izdevās nomest poļu kundzību, 18.~gadsimtā~--- atgūt agrāk zaudētās teritorijas, atbrīvojot tās no poliskās nekārtības.

Vācijā no XVIII gadsimta eksistēja termins „\detxti{polnische Wirtschaft}”~--- poļu saimniecība, kas, vācu acīm raugoties, nozīmēja pilnīgu nesaimnieciskumu. (Tiesa, pēc 1934.~gadā parakstītās Vācijas un Polijas deklarācijas par spēka nepielietošanu 1936.~gadā iznākušajā V.~Noltinga grāmatā bija norādīts, ka tajā laikā minētais termins esot jānodod „grabažu noliktavā”, taču vēl ilgi daudzi vācieši poļu ekonomiskos sasniegumus augstu nevērtēja.) XIX gadsimtā Vācija Prūsijas vadībā kļuva par valstiskuma, ekonomikas un kara mākslas augstāko sasniegumu. 19.~gadsimtā otrajā pusē, kad poļi visiem spēkiem cīnījās par nacionālo pastāvēšanu, vācieši kā pirmo un galveno „senseno” ienaidnieku uzlūkoja frančus; tikai valsts austrumu teritorijās, kur dzīvoja daudz poļu, vācieši ar lielākām bažām vērās austrumu virzienā.

Pēc Pirmā pasaules kara bezdibenis starp vācu un poļu vēstures ainām vēl padziļinājās. Kā atzīmējis poļu vēsturnieks M.~Mročko, poļu priekšstatiem par „vācu virzību uz austrumiem” („\detxti{Deutschen Drang nach Osten}”) blakus pastāvēja vācu viedoklis par „slāvu pieplūdumu” („\detxti{Andrang des Slawentums}”) rietumos. Vācu vēstures zinātne Polijas vēsturi un politiku tēloja kā „spiedienu uz Rietumiem” („\detxti{Drang nach Westen”}) un kā „poļu briesmas” („\detxti{polnische Gefahr}”). Ievērojamais vācu vēsturnieks J.~Hallers 1923.~gadā iznākušajā kapitāldarbā „Vācu vēstures laikmeti'' („\detxti{Epochen der deutschen Geschichte}'') rakstīja: „Poļus mūs ar pilnām tiesībām uzskatām par vācu mūžseniem ienaidniekiem Austrumos''. („\detxti{Polen denken wir uns mit Recht als den Erbfeind der Deutschen im Osten}”.)

Arī mūsdienās, kad bundesvēra augstākie virsnieki noliek ziedus Vesterplatē (\detxti{Westerplatte}~--- pussala Baltijas jūrā, kur Otrā pasaules kara sākumā poļi varonīgi pretojās nacistiskajiem iebrucējiem), vāciešu vidū vēl nav izveidojies pamatots un daudz-maz vienots priekšstats par kaimiņtautas vēsturi. Vācu publicistikā arī XX gadsimta 80.~gados varēja lasīt dusmu pilnus vārdus par „nepateicīgajiem” poļiem, kuri gan saņēma no VFR materiālu palīdzību streikotājiem~--- protestētājiem pret „sociālistisko” pārvaldi Polijā, bet tik un tā neatzina pēckara vācu iedzīvotāju deportāciju netaisnīgumu, vietā un „nevietā” atgādināja vācu iebrucēju nodarījumus Polijas teritorijā utt. Kāda pirms Otrā pasaules kara Polijā dzīvojusi un no turienes izraidīta vāciete E.~Losere 1981.~gadā rakstīja: „Kopš kristīšanas [poļu] tauta tika pakļauta stingrai kleriķu uzraudzībai, kura traucēja personības attīstībai. Viņi nespēj izrauties no šiem spaidiem. Viņi tika tā apspiesti, ka uzkrāto agresiju arvien atkal un atkal izlādē bezprecedenta naidā pret brīvākajiem un bagātākajiem vāciešiem. \citespace{} Baznīca triumfē Polijā. Un katoļu baznīca no pašiem sākumiem bija Vācijas valsts niknākais ienaidnieks. Poļi tika un tiks izmantoti kā svira lai nolaistu vācu un Vācijas tautsaimniecības asinis”. Arī XXI~gadsimta sākumā publicistikā var atrast līdzīgus „zinātniskus'' spriedumus par citām tautām.

Kā Hamburgā izdotajā vācu-poļu žurnālā „\detxti{Dialog}” („Dialogs”) rakstījis poļu žurnālists Ā.~Kržeminskis, vāciešiem attieksme pret Poliju un poļiem arī XXI gadsimtā ir atmiņas politikas (\detxti{Gedächtnispolitik}) pārbaudes līdzeklis. Polija taču bija pirmā valsts, kura ar ieročiem cīnījās pret nacistiskās Vācijas agresiju, no pirmās dienas bija antihitleriskās koalīcijas locekle, cieta kara laikā milzīgus cilvēku un materiālos zaudējumus, un, kaut arī ne no uzvarētājām lielvalstīm, ne arī Vācijas puses netika atzīta par ar tām vienlīdzīgu, tomēr kļuva par Vācijai tik sāpīgi zaudēto teritoriju ieguvēju. Taču vāciešiem atmiņās par Otro pasaules karu Polijā ir periferiāla vieta. Kā norādījis ievērojamais vācu vēsturnieks H.U.~Vēlers, Vācijas abiturientiem un studentiem ir zināšanas par Otrā pasaules kara laikā nogalinātajiem sešiem miljoniem ebreju, taču kad viņiem saka, ka tai pat karā dzīvību zaudēja gandrīz katrs piektais polis [precīzāk gan būtu teikt~--- Polijas pavalstnieks~--- V.Š.], un jau kara sākumā no vācu okupētajiem apgabaliem tika padzīti 800~000 iedzīvotāju, nākas sastapties ar nezināšanu un izbrīnu. Sociologs un vēsturnieks Rietumu institūta direktors Poznaņā A.~Saksons raksta, ka vēl daudz ūdens ir jāaizplūst robežupē Oderā līdz vācieši mainīs savus uzskatus par poļiem. XX gadsimta beigās veiktā socioloģiskā aptauja rādīja, ka 50\% vāciešu nav priekšstata par Poliju vai arī viņiem tā neinteresē. 2000.~gadā Polija viņu acīs bija pēdējā vietā no 26 pievilcīgākajām Eiropas valstīm.

[Pamatīgāks materiāls par Polijas vēstures izpēti vācu zinātniskajā literatūrā atrodams G.~Rodes rakstā „\detxti{Der Geschichte Polens in der deutschen Geschitsschreibung}” („Polijas vēsture vācu vēstures literatūrā'') krājumā „\detxti{Nationalgeschichte als Problem der deutschen und polnischen Geschichtsschreibung”}. Braunschweig, 1983.) Par Polijai veltītās vācu historiogrāfijas vērtējumu no pašu poļu viedokļa var izlasīt H.~Olševska rakstā „\detxti{Die deutsche Historiographie über Polen aus polnischen Sicht}” („Vācu historiogrāfija par Poliju no pašu poļu viedokļa”) rakstu krājumā \detxti{Dittmar Dahlmann (Hg)} „\detxti{Hundert Jahre Osteuropäische Geschichte. Vergangenheit, Gegenwart und Zukunft}''. Stuttgart, 2005].

Negatīvi domāja arī poļi par vāciešiem. Ne velti 1939.~gada 1.~septembrī Polijas prezidenta I.~Moscicka runā izskanēja vārdi: „mūsu senais ienaidnieks ir uzsācis karadarbību pret Polijas valsti”. Poļi labi atceras, ka Varšava, faktiskā Polijas galvaspilsēta kopš 1596.~gada, pēdējos divarpus gadsimtos trīsreiz (1795., 1915., 1939.~gg.) krita vācu militāristu rokās. Aušvicas (\detxti{Auschwitz}) jeb Osvencimas, kā arī daudzas citas nacistu ierīkotās koncentrācijas nometnes atradās Polijā un arī šodien atgādina saviem apmeklētājiem par vācu iebrucēju noziegumiem kopš krustnešu laikiem līdz pat Otrajam pasaules karam, kad vārds „vāciešu” kļuva par sinonīmu vārdam „noziedzīgs”.

Protams, vēsturē atrodami daudzi fakti arī par poļu un vāciešu sadarbību. Patiesība ir tā, ka naidīguma izvirdumiem starp vāciešiem un poļiem un, vēl lielākā mērā, starp viņu politiķiem, var nostādīt pretī periodus, kad abas tautas prata sadzīvot draudzīgi. Piemēram, Boļeslavs Drosmīgais sniedza militāru palīdzību Svētās Romas impērijas ķeizaram Ottonam III, pats tika kronēts pateicoties tā atbalstam. Abi personīgi tikās t.s. Gņezno (poļu \pltxti{Gniezno}, vācu \detxti{Gnesen}) kongresā (1~000), kur tika radīta patstāvīga Gņezno arhibīskapija, kas Polijai garantēja baznīcas neatkarību no vācu baznīcas. Mūsdienās pārsvaru gūst uzskats, ka pretēji abām augstāk izklāstītajām leģendām par Grīnvaldes kauju XV~gadsimtā sadursme starp Poliju un krustnešiem nebija konflikts starp nācijām vai pat civilizācijām, bet gan starp valstīm. Pie tam abi konkurējošie bloki: gan Polija un Lietuva, gan Teitoņu ordenis, kuru atbalstīja arī Svētā Romas impērija un Čehijas karaliste, pārstāvēja daudzus etnosus.

Vācu un poļu kultūras ir savstarpēji bagātinājušas viena otru. Lielākoties vācu (un ebreju) amatnieku apdzīvotās Polijas pilsētas pārņēma vācu pilsētu Lībekas, Magdeburgas, Nirnbergas vai Halles tiesības. Neviena valoda, izņemot radniecīgo krievu valodu, nav tik stipri ietekmējusi poļu valodu kā vācu valoda. Vācu „\detxti{Rathaus}'' (rātsnams) poliski ir „\pltxti{ratusz}”, „\detxti{Bürgermeister}” (birģermeistars)~--- „\pltxti{burmistrz}” un „\detxti{Pflug}” (arkls)~--- „\pltxti{plug}”. Poļu valodnieks G.~Korbuts ir atzinis, ka no 100 poļu sarunu valodas vārdiem 16--17 ir vācu izcelsmes. Savukārt vācu valodā tādi vārdi kā „\detxti{Dolmetscher}” (tulks), „\detxti{Droschke}” (važoņa rati), „\detxti{Gurke}” (gurķis), „\detxti{Peitsche}” (pātaga), „\detxti{Quark}” (biezpiens), „\detxti{Zobel}” (sabulis) u.c. ir pārņemti no senās poļu valodas. Polijas nacionālās himnas mūzikas autora ģenerāļa J.~Dombrovska māte bija vāciete. Himnas teksta autors ģenerālis J.~Vibickis cienīja vācu kultūru un izaudzināja savus dēlus Drēzdenē vācu klasikas garā. Tādu piemēru uzskaitījumu varētu turpināt.

Taču, kaut arī šodien „mūžsenā ienaidnieka” tēls abās pusēs ir krietni pabālējis, daži tam raksturīgi priekšstati ir vēl dzīvi. Poļiem ir paruna: „\pltxti{Póki swiat swiatem, Polak Niemcowi nie bedzie bratem}”. („Tik ilgi, kamēr pasaule pastāvēs, polis nekad nebūs vācietim brālis.'') Īpaši negatīvas emocijas poļos izraisa jēdziens ''\pltxti{Prusy}''(prūši). Līdz mūsdienām Varšavā nav Berlīnes vai Prūšu ielas. Toties ir Sakšu pils, Sakšu parks, Sakšu ass (ceļš), Leipcigas iela un Drēzdenes iela. Poļu--sakšu tradicionālās attiecības ir dzīvas vēl šodien. Vārdam „saksis”, atšķirībā no jēdziena „prūsis” poļu ausīs ir pozitīva nozīme. Tomēr, domājot par Vāciju, poļi vispirms atceras Prūsiju un tās iekarojumus Polijā.

Kā atzīmējis populārais poļu publicists A.~Kržeminskis, ciktāl paši poļi skata savas vēstures tumšās lappuses un kritizē daudzus nacionālos mītus, tiktāl viss ir kārtībā, īpašu saasinājumu parasti nav. Taču, ja kritikai pievienojas kāds vācietis, pieminot, piemēram, poļu represijas pret vācu gūstekņiem un nosoda vācu deportācijas no pēckara Polijas, poļu puse nonāk līdz pat publiskiem skandāliem.

Ilgstoša saskarsme poļiem ir bijusi arī ar krieviem (\pltxti{rosjanie}) un Krieviju (\pltxti{Rosja}). Taču divas kaimiņos dzīvojošas lielas radniecīgas tautas vēstures gaita saistīja ar dažādām kristīgās ticības konfesijām, dažādiem kultūras virzieniem, kas saasināja to attīstības nevienmērību. Un tāpat kā daudzos citos gadījumos, kad starp radniecīgām tautām (Izraēlas un arābu valstu iedzīvotājiem, serbiem un horvātiem u.c.), arī starp poļiem un krieviem ilgstoši pastāv it kā nepārvaramas pretrunas. Neatkarīgās Polijas valsts vēsturē grūti atrast periodus, kad tā būtu ilgstoši sadarbojusies ar Krieviju.

Polijas un Krievijas kopīgā vēsture neļauj ne vienu no tām nosaukt tikai par upuri, ne otru tikai par agresoru. Nebūt ne vienmēr poļi ir bijuši austrumu kaimiņa agresijas objekts, kā tas šķiet daudziem Polijas iedzīvotājiem, bieži viņi paši ir tam uzbrukuši. Ne reizi vien vēstures gaitā Polija ir savā labā izmantojusi Krievijai grūtas situācijas.

No XI līdz XVII gadsimtam poļi daudzkārt iebruka Krievzemes robežās, sagrāba milzīgas teritorijas un gadsimtiem ilgi ekspluatēja tās, tai laikā kad krievu karaspēks šai laikā ne reizi nenonāca pašas Polijas teritorijā.

Tā Kijevas Krievzemes laikos poļu karalis Boļeslavs Drosmīgais pēc sava znota bijušā Kijevas kņaza Svjatopolka, kurš bija zaudējis cīņu savam brālim Jaroslavam Gudrajam, aicinājuma devās uz Volīniju, sakāva Jaroslava Gudrā karadraudzi, 1018.~gadā ieņēma Kijevu, bet tā vietā lai atdotu to valdīšanā savas meitas vīram Svjatopolkam, pats sāka tajā valdīt. Taču pilsētnieki, sašutuši par poļu karadraudzes patvaļībām, sacēlās un Boļeslavs bija spiests atstāt pilsētu. Iejaukties Kijevas Krievzemes lietās mēģināja arī Boļeslavs II Drošais, taču pēc sadursmēm ar vietējiem iedzīvotājiem bija spiests savu karadraudzi izvest.

Jau dziļāku poļu un krievu pretrunu rašanās saistāma ar Maskavas Krievzemes jeb Maskavijas ģenēzi. Vēl XIII--XIV gadsimtā Lietuvas valsts atkaroja Zelta Ordai (mongoļu \mntxti{Altan Ord}, tatāru \tttxti{Altın Urda}, mongoļu-tatāru valsts XIII--XV~gs.) lielu daļu no senkrievu zemēm. Tieši Lietuvas valsts toreiz pretendēja uz šo zemju vienotājas lomu, taču pēc jau minētās 1385.~gadā noslēgtās Krēvas ūnijas arī Lietuvā izplatījās katoļu ticība, sākās senkrievu zemju polonizācija. Tāpēc Lietuva senkrievu acīs pakāpeniski zaudēja „savas” valsts tēlu, atbalsta lomu cīņā pret svešzemju jūgu; par tādu kļuva Maskavija, kura palika uzticīga pareizticībai.

1480.~gadā Lietuvas lielkņazs un Polijas karalis Kazimirs IV mēģināja ieņemt Novgorodu un Pleskavu. Protams, arī krievu kņazi veica karagājienus pret poļiem, taču tiem faktiski bija savu zemju aizsardzības un arī atriebības raksturs.

Turpmāk no 1558.~gada (Livonijas kara sākums) līdz 1939.~gadam (oficiālais Otrā pasaules kara sākums) Polija ar Krieviju karoja 13 reizes, pie tam poļi 2 karus uzvarēja un 11 zaudēja.

Pēc 1569.~gadā noslēgtās Ļubļinas ūnijas (poļu \pltxti{Unia Lubelska}, lietuviešu \lttxti{Liublino unija}, baltkrievu \betxti{Люблінская унія}~--- līgums par Polijas karalistes un Lietuvas lielkņazistes apvienošanos konfederatīvā valstī~--- Žečpospolitā ar vēlētu karali) sākās gadsimtiem ilgusī tās un Maskavas valsts sāncensība par austrumslāvu zemju apvienošanu un pakļaušanu.

Var teikt, ka ap 200 gadus XVI --XVII gadsimtā Maskavija izjuta Žečpospolitas spiedienu. Žečpospolita uzskatīja sevi par galveno katoļu ticības izplatītāju austrumos, turpretī Maskavas valsts cīnījās ne tikai par „krievu zemju” vienību, bet arī „visu pareizticīgo” kristiešu interesēm. Maskavija un Polija saskatīja viena otrā antagonistu slāvu zemēs. Polijai kaimiņos Krievijā pastāvēja Maskavas kā „trešās Romas” koncepcija, vadoties no XV~gadsimtā izskanējušās tēzes, ka „Divas Romas [Roma un Bizantija] ir kritušas, trešā [Maskava] stāv, bet ceturtās nekad nebūs” („\rutxti{«Два Рима пали, третий стоит, а четвёртому не бывать}”). XVI gadsimtā Polijas-Lietuvas valsts robeža atradās netālu no Možaiskas. No Livonijas kara (1558--1583) laikiem, kad Maskavija un Žečpospolita cīnījās par ietekmi Baltijā un Stefans Batorijs neveiksmīgi mēģināja ieņemt Pleskavu, abu valstu savstarpējās cīņas gan apdzisa, gan uzliesmoja ar jaunu sparu.

Tās saasinājās it īpaši pēc Polijas iejaukšanās Krievijas lietās t.~s. „Juku laikos” XVII gadsimta sākumā, kad poļu karaspēks okupēja Maskavu: 1604.~gadā kopā ar Viltusdmitriju (\rutxti{Лжедимитрий}, ?--1606) un 1610.--1612.~gadā „cara” Vladislava laikā. XVII gadsimtā poļi un lietuvieši tik bieži un ilgi necīnījās ne ar vienu citu ienaidnieku kā ar Maskavas valsti. XVII gadsimta sākumā Polijā bija populāra atsaukšanās uz spāņu konkiskadoru veiktajiem iekarojumiem Amerikā. Tika spriests: spāņu taču bija ļoti nedaudz salīdzinājumā ar indiāņiem, bet viņi tos uzvarēja. Poļu ir daudz, vai nu viņi netiks galā ar „moskaļiem” [nicīgs krievu nosaukums, lietots poļu, ukraiņu u.c. vidū.~--- V.Š.]? Nabadzīgā poļu šļahta cerībā uz bagātīgu laupījumu devās palīgā Viltusdmitrijam uz Maskavu. Vairākkārt Polija ir bijusi tuvu savam mērķim nostiprināt robežu ar Krieviju pa Dņepras upi.

Nobeļa prēmijas laureāts krievu rakstnieks A.~Solžeņicins, izklāstot abu valstu savstarpējos pāridarījumus, sāka ar Krievijas puses nodarījumiem Polijai, taču pēc tam rakstīja arī par to, ka „iepriekšējos gadsimtos plaukstošā, stiprā, pašpārliecinātā Polija ne īsāku laiku un ne vājāk iekaroja un apspieda mūs [krievus~--- V.Š.]”. Īpaši viņš izcēla „juku laikus” Krievijā XVII gadsimta sākumā, kad „poļi teju neatņēma mums nacionālo neatkarību, šo briesmu dziļums bija ne mazāks kā tatāru iebrukumam, jo poļi apdraudēja arī pareizticību. Pie sevis viņi to sistemātiski apspieda, dzina ūnijā [domāta 1596.~gadā noslēgtā Brestas ūnija~--- V.Š.] \citespace{} Mūsu juku laikos Polijas ekspansiju poļu sabiedrība uztvēra kā normālu un pat pareizu politiku. Paši sevi poļi iedomājās kā Dieva izredzētu tautu, kristietības bastionu, ar uzdevumu izplatīt īstenu kristietību puspāgāniskajos pareizticīgajos, mežonīgajā Maskavijā, būt par universitāšu renesanses kultūras izplatītājiem”.

Ievērojamais krievu vēsturnieks V.~Kļučevskis apgalvoja, ka polis un tatārs bija īstena krievu cilvēka pastāvīgi ienaidnieki arī XVIII gadsimtā. Pēc vēsturnieka un publicista V.~Fiļeviča vārdiem viduslaikos „krievi sātanu iedomājās poļa izskatā” („\rutxti{русские беса представляли в виде ляха}”). Taču katru reizi poļu ekspansijai sekoja prettrieciens.

Pēc Romanovu (\rutxti{Романовы}) dinastijas nodibināšanas (vēlāk Holšteinas-Gotorpas-Romanovu dinastija~--- \rutxti{Гольштейн-Готторп-Романовская династия}), kura valdīja Krievijas caristē, pēc tam Krievijas impērijā laikā no 1613.~gada līdz 1917.~gadam, tā pakāpeniski atkaroja austrumslāvu zemes.

Dažkārt Polija un Krievija gan bija arī sabiedrotās, piemēram, XVII gadsimta beigās pret kopējo ienaidnieku~--- Turciju, Ziemeļu karā (1700--1721) pret Zviedriju. Taču tad atkal priekšplānā izvirzījās nesaskaņas. Pateicoties cara, vēlāk imperatora Pētera I īstenotajai modernizācijai, Krievija izvirzījās priekšā kaimiņvalstij. XVIII~--- XX gadsimtā jau tā diktēja poļiem politiskās spēles noteikumus. Krievijas robežas rietumos atvirzījās aiz Vislas. XVIII~gadsimtā Žečpospolita zaudēja savu valstiskumu. Taču poļu šļahta nevarēja samierināties ne ar savas neatkarības, ne ar citu tautu apdzīvoto, bet līdz tam Žečpospolitas varā esošo teritoriju zaudējumu, meklēja pret Krieviju sabiedrotos Centrālajā un Rietumu Eiropā. 1812.~gadā poļi, karojot nu jau Napoleona I (\frtxti{Napoléon I Bonaparte}, 1769--1821) karaspēka sastāvā, atkal nonāca Maskavā.

Divreiz XIX gadsimtā uzliesmoja varenas poļu sacelšanās pret Krievijas varu. Kaut poļi cīnījās par savu neatkarību pret visām to sadalījušajām valstīm, taču galvenais ienaidnieks bija Krievija, kura bija ieguvusi centrālos Polijas apgabalus. Iekšējās pretrunas gan Krievijā, gan Polijā neļāva tām miermīlīgi izšķirties pēc Pirmā pasaules kara un Oktobra apvērsuma Krievijā. Daudz sarkanarmiešu zaudēja dzīvību 1920.~gadā pie Varšavas. Poļi kā vienu no svarīgākajiem notikumiem Otrā pasaules kara gados atceras 1939.~gada 17.~septembrī notikušo Sarkanās armijas invāziju Polijas teritorijā, pēc kara ar izteiktu neapmierinātību pārcieta PSRS valdošā staļiniskā režīma ietekmes izplatību uz Poliju. Tā tūkstošgadīgajā līdzāspastāvēšanas laikā divu slāvu tautu attiecībās uzkrājās ne mazums notikumu, kurus katra puse traktēja citādi. Uz vēsturisko notikumu subjektīvas uztveres pamata formējās vispārnacionāli mīti. Krievijā~--- mīts par naidīgajiem poļiem-katoļiem, kuri jebkurā situācijā ir gatavi darboties tā, lai ieriebtu „moskaļiem”, Polijā~--- mīts par vienmēr to apdraudošo, barbarisko un impērisko Krieviju.

No dotā darba, kas veltīts Polijas XIX un XX gadsimta vēsturei, autora viedokļa ļoti būtisks ir jautājums par poļu un krievu tautu attiecībām minētajā laikā. Tieši to samežģījumi ir nesuši poļu tautai daudz ciešanu.

Par šīm attiecībām, par daudzo Polijas un Krievijas konfliktu būtību rakstījuši dažādu tautu vēsturnieki un publicisti. Dažu no viņiem skaidrojumi nevar tikt uzskatīti par zinātniskiem.

Piemēram, baltkrievu publicists A.~Tarass, atsaucoties uz nenosauktu zinātnieku ģenētiskajiem pētījumiem, raksta, ka „krievi (no vienas puses), baltkrievi un poļi (no otras puses) ir etniski atšķirīgas tautas, antropoloģiski dažādas rases” un arī „tieši iedzimtas (ģenētiskas) atšķirības izšķiroši veicināja maskaviešu (vēlāk krievu) tradicionālo noraidīšo attieksmi pret Eiropas dzīvesveidu, Eiropas kultūru, katolicismu, uniātismu un protestantismu, viņu naidu pret Lietuvas Lielkņazisti, Kungu Lielo Novgorodu, Livoniju, Polijas Karalisti.” Minētais autors arī apgalvo, ka „tieši poļu un baltkrievu etniskās kopības fakts izskaidro Žečpospolitas izveidi”. Identificējot baltkrievus ar poļiem, A.~Tarass savos darbos vēsturiskos notikumus skata no polonopfīlām un rusofobām pozīcijām. Tā kā autors nenosauc viņa uzskatus apstiprinošo „ģenētiķu” vārdus, hipotēze šai grāmatā netiek komentēta, vienīgi jāsaka, ka šāds vienkāršots sarežģītu procesu izskaidrojums ir tāls no patiesības.

Daļa vēsturnieku uzskata, ka galvenā bija tīri politiska konkurence par dominanci Austrumeiropā.

Citi aizstāv viedokli, ka visu nesaskaņu pamatā bija nesamierināma cīņa starp kristietības austrumu un rietumu atzariem. Reizē līdzās reliģiskajai konfrontācijai XVII gadsimta notikumos tiek izšķirtas arī etniskam konfliktam raksturīgas pazīmes. Tā, pazīstamais poļu izcelsmes amerikāņu vēsturnieks R.~Paips atzīmējis: „Vēsturiskajā perspektīvā poļu--krievu attiecības nekad nav bijušas labas. Starp \citespace{} [šīm] valstīm vienmēr radās sasprindzinājums. \citespace{} Krievija uzskata Poliju par slāvu tradīciju nodevēju, jo \citespace{} [tā] pieņēma katoļu ticību. Savukārt poļi nemīl krievus, jo pārāk daudz ir no viņiem cietuši”.

Nebūt neatbilst patiesībai priekšstats, ka lielas daļas poļu naidīgā izturēšanās pret visu krievisko radās kā sekas trijām Polijas dalīšanām XVIII gadsimtā, poļu sacelšanos apspiešanām XIX gadsimtā, padomiskā „sociālisma” modeļa uzspiešanai, veselai virknei notikumu, kuros cietusī puse bija poļu tauta. Kā parāda poļu vēsturnieces A.~Neviaras 2006.~gadā publicēts pētījums (Niewiara A. Moskwicin—Moskal—Rosjanin w dokumentach prywatnych. Łódź, 2006) „moskaļa-aziāta”, „moskaļa-iebrucēja,” „moskaļa-Kristus nodevēja” tēls parādās poļu rakstītajos avotos vēl sen pirms XVIII~gadsimta, kad vēl nekāda poļu apspiešana no krievu puses nebija iespējama. Līdz Ļubļinas ūnijai (1569) Polijai nebija robežas ar Krieviju, tā bija Lietuvas lielkņazistei. Pēc ūnijas noslēgšanas ar Krieviju konfliktus risināja jau apvienotā Polijas-Lietuvas valsts (Žečpospolita). Ja līdz tam Polijas galvenie ienaidnieki bija rietumos (Vācija) un dienvidos (Turcija), tad tagad tā iesaistījās karos arī austrumos (pret Krieviju). Tad arī itin dabīgi radās poļu šļahtas naids pret „moskaļiem”. Tas bija ne Polijas apspiešanas no Krievijas puses rezultāts (kaut tā to sekmēja), bet jau senāks poļu šļahtas kolektīvo interešu izpausmes veids.

Līdz „sociālisma” padomju modeļa bankrotam, PSRS sabrukumam, Polijas Tautas republikas pārveidei, historiogrāfijā poļu-padomju attiecības tika skatītas no divām pretējām, savstarpēji izslēdzošām pieejām.

Viena bija pārstāvēta PSRS un PTR vēsturnieku darbos, otra Rietumu pētnieku, poļu emigrantu vidū dzīvojošo vēsturnieku, kā arī to poļu pētnieku monogrāfijās un rakstos, kuri, bieži ar pseidonīmiem, publicējās PTR valdības cenzūrai nepakļautajā vēsturiskajā literatūrā.

Pirmā koncepcija sāka veidoties jau Otrā pasaules kara laikā un bāzējās uz padomju vadības oficiālo pozīciju, kura bija pārstāvēta plaši pieejamos politiskajos un diplomātiskajos dokumentos, padomju līderu runās. Autori, kuri pieturējās šai koncepcijai, PSRS realizēto ārpolitiku, arī pret Poliju, vērtēja kā miermīlīgu, progresīvu, internacionālistisku. Turpretī Polijas Republikas starpkaru un Otrā pasaules kara laikā realizētā politika tika skatīta kā Padomju Savienībai atklāti naidīga, nacistiskās Vācijas agresīvos centienus veicinoša. Visspilgtāk šī tendence atklājās padomju vēsturnieku darbos, veltītos Polijas vēstures apkopojumam, Otrajam pasaules karam, ārpolitikas un diplomātijas vēsturei. Analoģisku pozīciju ar nelielām niansēm ieņēma oficiālā Polijas historiogrāfija 1940.~gadu beigās~--- 80.~gados.

Otrās koncepcijas piekritēji par pamatu ņēma uzskatus, gluži pretējus augstāk izklāstītajai pieejai. Viņi viennozīmīgi saskatīja padomju ārpolitikā agresīvo, pretpolisko ievirzi, kā starpkaru, tā arī Otrā pasaules kara gados, īpaši uzsvēra represijas pret poļiem, to deportācijas, poļu virsnieku nošaušanu Katiņā u.~c. Pirmie tādi pētījumi parādījās jau drīz pēc Otrā pasaules kara. Visaptverošāk šī koncepceja izklāstīta V.~Pobog-Maļinovska darbā „Polijas jaunākā politiskā vēsture”. („\pltxti{Najnowsza historia polityczna Polski} 1864--1945”, 1--3, Paryż, Londyn, 1953.--1962.) Jāsaka, ka PSRS vadības nostāja tikai veicināja Rietumos izplatīto uzskatu izplatību arī pašā Polijā. Tā vietā, lai atslepenotu dokumentus, publicētu avotus par Polijas un Krievijas/PSRS attiecību vēsturi, atspēkotu mītus par to, tika vienkārši klusēts. Tas veda pie secinājuma~--- ja jau PSRS klusē, tad to apsūdzošie notikumu izklāsti ir patiesi, bet varbūt patiesība bija vēl briesmīgāka.

Abas koncepcijas pārstāvošajiem darbiem bija raksturīga klaja tendenciozitāte, tieksme padarīt baltu savu pusi un nomelnot pretējo, plaša sabiedriska fona iztrūkums, savstarpēja pretējās puses nacionālo interešu ignorēšana.

Vēstures mantojums arī mūsdienās bremzē labu kaimiņattiecību veidošanos starp Poliju un Krieviju. Pēc Otrā pasaules kara Polijas Tautas Republikas (1952--1989) historiogrāfijā nereti tika uzsvērts, ka neraugoties uz to, ka Krievija un Polija vēsturiski ilgus gadsimtus bija ienaidnieki, abu tautu attiecībās pārsvarā esot bijušas simpātijas un sadarbība. Taču šis apgalvojums nespēja aizsegt to, ka vēsturiskās kolīzijas, valstiskās attiecības nevarēja neietekmēt arī tautu, ierindas cilvēku savstarpējās jūtas. Tā, krievu revolucionārs P.~Lavrovs 1889.~gadā atstāstīja kā gadu pirms Parīzes komūnas (1871) „man nācās piedzīvot, ka poļu mākslas profesors atstāja telpu, dzirdot, ka Dombrovskis griežas pie manis krievu valodā.” Līdzīgi arī XX gadsimta 90.~gados šī darba autoram personīgi nācās piedzīvot, kā pret citām zemēm un tautām, arī Latviju un latviešiem, draudzīgi noskaņots inteliģents pavecāks polis, kurš prata gan krievu, gan vācu valodu, tomēr nevēlējās tajās sarunāties, ar to demonstrējot savu attieksmi pret bijušajiem „okupantiem”.

Tautu savstarpējās attiecības sevišķi traucējoši ietekmēja gan tā saucamā „sociālisma” apstākļos, gan pēc tā sabrukuma skanējušie oficiālie meli.

Kaut XX gadsimta 90.~gados sākās vēstures deideoloģizācija un depolitizācija, tā nebūt nav sekmīgi pabeigta. Kā pamatoti norādījis Daugavpils Universitātes profesors A.~Ivanovs: „\dots{}Dažādas politiskas institūcijas un politiskās elites vēsturi izmantoja un joprojām izmanto kā vienu no politikas veidošanas un īstenošanas instrumentiem, kas palīdz radīt elitei izdevīgus vēsturiskus mītus un nodrošina vienprātību sabiedrībā. \citespace{} vēstures politizēšanas pakāpi nosaka politiskais režīms: antidemokrātiskā režīma apstākļos vēsture kļūst ārkārtīgi politizēta un ideoloģizēta, tieši un atklāti kalpojot valstij un oficiālajai ideoloģijai, turpretim demokrātiskā vara ļauj vēsturniekiem relatīvi brīvi interpretēt vēstures procesu un vērtēt vēstures faktus. Kaut gan arī šī radošā brīvība nav absolūta.” Turpinot šo domu gaitu, var piebilst, ka centienu „nodrošināt vienprātību sabiedrībā” parādīšanās jebkurā valstī ir līdzvērtīga demokrātijas ierobežošanas mēģinājumiem.

Poļu vēsturnieks J.~Duračinskis ir atzīmējis, ka daudzos pēc 1989.~gada publicētajos darbos pamatoti novērsto iepriekšējo melu vietā parādās jauni meli un „gan tie vecie, gan mūsdienu meli izaug uz noteiktu (kardināli pretēju) simpātiju vai arī atsevišķu pētnieku politisko ieskatu pamata. Tas attiecas, galvenokārt uz izšķirošo notikumu un poļu politiskās skatuves galveno „aktieru” vērtējumu”. Vēlme pārskatīt uzkrāto pieredzi diemžēl bieži izpaužas vienkāršotā novecojušo uzskatu caurskatīšanā, agrāko plusu vietā liekot mīnusus un otrādi. Pēc Otrā pasaules kara norisušo sarežģīto iekšējo procesu padziļinātas interpretācijas vietā bieži tiek runāts tikai par tīri vardarbīgu padomju modeļa uzspiešanu Polijai.

[Minētā daudzu pētījumu autora J.~Duračinska historiogrāfiskie darbi var dod lasītājam diezgan pilnīgu priekšstatu par pašu poļu vēsturnieku veikumu savas dzimtenes XIX un XX~gadsimta vēstures izpētē. Daži no tiem ir pieejami arī krievu valodā. Piemēram: \rutxti{Э.~Дурачинский}. „\rutxti{О польской историографии новейшей истории}” („Par jaunāko laiku vēstures poļu historiogrāfiju”) // %http://library.by/portalus/modules/rushistory/readme.php?subaction=showfull&amp;id=1192094160&amp;archive=&amp;start_from=&amp;ucat=19&amp;category=19
% // 2.05.2010.]

Poļu vēsturnieki nereti nevēlas savas vēstures periodizācijā iekļaut arī pašu poļu izraisītās nacionālās sakāves, parasti cenšas uzsvērt, ka Polija vienmēr ir bijusi saistīta ar Rietumu civilizāciju, reizē mazinot sakaru nozīmi ar austrumu kaimiņiem. Acīmredzot to veicinājis arī tas, ka, runājot Polijas valsts darbinieka un rakstnieka S.~Kata-Mackeviča vārdiem, austrumos „mēs vienmēr bijām kungu nācija”, bet rietumos „mēs bijām tikai strādnieku un zemnieku nācija”.

Diemžēl mūsdienu poļu historiogrāfijā, bet īpaši politiskajā publicistikā krievu un poļu pretrunu un savstarpējās nepatikas jūtu cēloņi parasti tiek meklēti galvenokārt jaunāko laiku periodā~--- XVIII--XX gadsimta ietvaros, deviņos gadījumos no desmit uzmanību pievēršot tikai Polijas dalīšanai, uzsverot tās paverdzināšanu, ko veica Krievija, Prūsija/Vācija un Austrija/Austroungārija, arī nežēlīgo 1794., 1830.--1831. un 1863.--1864.~gada poļu sacelšanos apspiešanu, Otrajā pasaules karā Polijai nodarītās patiesās un iedomātās „netaisnības” un padomju modeļa uzspiešanu tai. Mazāk tiek akcentēta Polijas, tikpat cik Krievijas, agresīvā darbība Livonijas karā, poļu īstenotā Pleskavas aplenkšana (1581--1582, 1615), Krievijas un Žečpospolitas cīņa par Smoļensku, Polocku un kreisā krasta Ukrainu, poļu dalība t.s. „juku laikos” Krievijā XVII~gs. sākumā, kad norisa poļu šļahtas atbalstītā Viltusdmitrija I darbība un tam sekojošās divas poļu intervences Krievijā, kā arī krievu-poļu karš 1654.--1667.~gadā, kad Krievija atsaucās Zaporožjes (\uktxti{Запорiзька Січ}) kazaku hetmaņa B.~Hmeļnicka aicinājumam palīdzēt cīņā pret poļu apspiedējiem un pieņemt Zaporožjes kazakus Krievijas pavalstniecībā.

Ja daļa poļu vēsturnieku Polijas rīcību pret Krieviju vērstajos karagājienos piemin arī paškritiski, tad tik un tā vairākumā viņu rakstītajā jūtams zināms lepnums par poļu „varoņdarbiem” šajos iekarošanas karos. Par daudziem citiem, mazāk „spožiem” notikumiem, poļu sakāvēm, arī par poļu iekarotāju nodarītajām ciešanām citu tautu iedzīvotājiem tiek runāts klusināti. Tiesa, ir arī citi piemēri. Jau minētais poļu zinātnieks J.~Tazbirs uz žurnālista jautājumu: „Kurš pirmais sāka nemīlēt otros~--- mēs krievus vai krievi mūs?” atbildēja: „Vienlaikus. Strīdus un sadursmju objektu veidoja Lietuva. Krievzeme bija vāja un Lietuva sagrāba tās zemes, kuras agrāk bija tās (Krievzemes) īpašumā. Sāncensība turpinājās, un uz tās fona noteikti bija jānonāk līdz konfliktam.”

Diemžēl, arī šeit ir vajadzīgs precizējums. Poļu vēsturnieks citētajā frāzē it kā nevilšus visu atbildību par konfliktiem uzkrauj Lietuvai, Poliju pat nepieminot, kaut Lietuvas lielkņazistē ilgu laiku pārsvarā bija slāvu pareizticīgie iedzīvotāji, visai Rietumu Krievzemei Lietuvas lielkņaziste bija dabīgs pretestības centrs kā pret tatāriem, tā Teitoņu ordeni. Par ietekmi Austrumeiropā cīnījās gan poļu karaļi, gan krievu kņazi, gan leišu kunigaiši, taču Lietuvai pakāpeniski nonākot Polijas jaunākā partnera lomā, galveno cīņas smagumu pret Krievzemi uzņēmās tieši Polija.

Žečpospolitas un Krievijas strīds par to, kura no tām valdīs pār Zaporožjes kazakiem, beidzās ar Andrusovas pamiera noslēgšanu 1667.~gadā, ar kuru pirmā zaudēja visus apgabalus Dņepras kreisajā krastā par labu Krievijai. Tas bija ievērojams pagrieziena punkts poļu-krievu attiecībās. No šī laika pusotru gadsimtu līdz 1815.~gadam Krievija izplatīja savu ietekmi uz rietumiem, līdz liela daļa Žečpospolitas nonāca tās sastāvā. Jau 1686.~gadā karalis J.~Sobeskis noslēdza ar Krieviju t.s. Mūžīgo mieru (poļu historiogrāfijā \pltxti{pokój Grzymułtowskiego}), kurā starp citu Žečpospolita apsolīja saviem pavalstniekiem~--- pareizticīgajiem piešķirt ticības brīvību, bet Krievija apņēmās tos aizstāvēt. Tas deva Krievijai ieganstu iejaukties Žečpospolitas lietās.

Mūsdienās savstarpējos poļu~--- krievu pārmetumos parādā poļu nacionālistiski noskaņotajiem vēsturniekiem un publicistiem nepaliek arī liela daļa Krievijas vēsturnieku un publicistu. Uzstājoties vai nu joprojām no padomiskām vai vienkārši impēriskām pozīcijām, viņi nevēlas atzīt Krievijas un PSRS noziegumus pret poļu tautu, visur saskata tikai poļu darbības negatīvās sekas.

[Īsu pārskatu par krievu un padomju vēsturnieku veikumu Polijas vēstures izpētē sniedz S.~Falkovičas raksts: \rutxti{Светлана Фалькович. Польская проблематика в российской историографии} (Poļu problemātika krievu historiogrāfijā) // %http://jazon.hist.uj.edu.pl/zjazd/materialy/falkowicz.pdf
// 24.07.2011.]

Savstarpējos strīdos diemžēl bieži dzimst nevis patiesība, bet jauni aizspriedumi pret kaimiņtautām. Noteiktu liecību par poļu un krievu attiecībām sniedz arī abu valstu oficiālie svētki.

Tagadējās Krievijas Federācijas likumdevēji izvēlējās par Krievijas Tautas vienības dienu noteikt 4.~novembri, t.i.~--- dienu, kad pēc dažiem pieņēmumiem 1612.~gadā krievu zemessardze K.~Miņina un D.~Požarska vadībā atbrīvoja Maskavu no poļu iebrucējiem.

Tiesa, svētku datums tika izraudzīts bez pienācīgām konsultācijām ar speciālistiem. Mūsdienu krievu ekonomists, vēsturnieks un politiskais darbinieks V.~Šeiniss uzsvēris, ka šai dienā poļi tika padzīti tikai no vienas Maskavas pilsētas daļas (\rutxti{Китай-Город}), bet Kremlī viņi vēl palika, un viņu karogi no tā sienām tika nomesti tikai 1612.~gada 7.~novembrī. Plašāku izvērtējumu Krievijas tautas vienības dienai izvēlētajam datumam 1612.~gadā ir devis ievērojamais krievu vēsturnieks V.~Nazarovs. Šai dienā Maskavā nekas ievērojams nenotika, minētā Maskavas daļa (\rutxti{Китай-Город}) tika atbrīvota pēc jaunā stila nevis 4., bet gan 1.novembrī. Nepareizais datējums radās hronoloģiskas kļūdas rezultātā, saskaņojot Pareizticīgās baznīcas un laicīgo kalendāru. 5.novembrī Kremlī ielenktie interventi, kuru rindas bija raibas pēc nacionālā sastāva un poļi, visticamāk, bija mazākumā, parakstīja kapitulāciju, bet nākamajā dienā Kremļa garnizona padevās. Tieši Kremļa atbrīvošana kļuva par ievērojamu notikumu. Tādejādi tagad Krievijā spēkā esošā svētku diena balstās uz kļūdainu datējumu, neprecīzu vēsturisko avotu traktējumu un notikumu interpretāciju, kura neatbilst avotiem.

Lai arī kam būtu taisnība, svētku datuma izvēles politizēto raksturu nevar noliegt.

Polijas Republikā kā Polijas armijas diena (\pltxti{Świeto Wojska Polskiego}) tiek atzīmēts 15.augusts, kad Polijas armija 1920.~gadā uzvarēja Sarkano armiju kaujā pie Varšavas.

Šāda masu vēsturiskās atmiņas politizēta virzība traucē atklātības apstarotu, mūsdienīgu attiecību izveidi gan atsevišķu cilvēku, gan veselu tautu starpā. Savukārt, mūsdienās pastāvošās politiskās problēmas traucē objektīvi izvērtēt Polijas vēsturi.

Nākas ar nožēlu konstatēt, ka vēsturē diemžēl var atrast argumentus par labu gandrīz jebkuram viedoklim. Īpaši tas attiecās uz divu tautu attiecībām, kuras dzīvojušas blakus vairāk nekā tūkstoš gadu un kuras gan vieno, gan šķir tūkstošgadu savstarpējās cīņas vēsture. Šai gadījumā galīgi aplami ir jautājumi „kurš pirmais sāka?” un „kurš ir vairāk vainīgs?”. Jau attiecībā uz XX gadsimta sākuma situāciju krievu filozofs N.~Berdjajevs runāja par to, ka poļiem un krieviem ir jāatbrīvojas no atmiņām par gadsimtiem ilgo pretstāvi, grēkus, atzīstot savu vainu, ir jānožēlo visiem. Šai sakarā dziļu patiesību nes poļu vēsturnieka S.~Keņeviča doma: „Nācijai īsta patiesība, lai kāda arī tā būtu, nevar kaitēt \citespace{} Aprakstot pagājušo gadsimtu konfliktus, ir jāņem vērā ne tikai mūsu labā esošie argumenti, bet arī argumenti, ka nāk par labu mūsu bijušajiem pretiniekiem, lai neaprobežotos tikai ar sašutumu par mums nodarītajām pārestībām, bet parunātu arī par pāridarījumiem, kurus mēs esam nodarījuši citiem, lai, galu galā, nestiprinātu mūsu pārliecību par pārākumu pār citām nācijām, atceroties to, ka analoģiska citu nāciju pārliecība attiecībā par poļiem mūs aizvaino”.

Bijušais Polijas komunistu līderis un Valsts prezidents V.~Jaruzeļskis, ir stāstījis, kā viņš savas vizītes laikā Londonā tikās ar Lielbritānijas premjerministri M.~Tečeri un ieminējās tai, ka viņš ir liels Napoleona I cienītājs. M.~Tečere viņu uzveda savas mājas bēniņos, parādīja plānu ādas mapīti un teica: „Jūs zināt, kam tā piederēja? Napoleonam. Taču francūzi es nekad šurp neatvestu!” V.~Jaruzeļskis to komentēja: „Skatieties, 200~gadu ir pagājis, bet angļi nevar aizmirst [Napoleona karus]. Viņi ar frančiem dzīvo savienībā, divus pasaules karus karoja kopā. Un neraugoties uz to, katrai valstij palicis diametrāli pretējs Napoleona redzējums, sava vēstures versija. Taču tas nedrīkst traucēt sadarbību.”

Kā atzīmējis Polijas Zinātņu akadēmijas zinātniskais līdzstrādnieks, Berlīnes Brīvās universitātes goda profesors R.~Traba, poļiem ir gan vajadzīga Eiropas telpa dialogiem, taču pirmkārt viņiem ir jāatrisina „iekšējie strīdi”, kritiski jāizvērtē sava pagātne, lai „skeleti skapī” netraucētu tikt vaļā no iemantotajiem kompleksiem un politiskā balasta. Var piebilst, ka poļiem, daudz cietušiem no staļiniskajiem PSRS drošības orgāniem, nevajadzētu aizmirst, ka pirmie divi to vadītāji bija poļu šļahtiču kārtai piederīgie F.~Džeržinskis un V.~Menžinskis, tāpat no šīs kārtas nāca arī viena no staļinisma laikmeta drūmākajām figūrām~--- staļiniskais prokurors t.~s. „atklātajos procesos”~--- A.~Višinskis.

Ilgstoši apspiestas tautas nacionālais jūtīgums traucē tai un arī tās vēsturniekiem objektīvi izvērtēt vēsturi.

Izcilais mūsdienu krievu filozofs un sociologs A.~Zinovjevs ir norādījis uz interesantu likumsakarību: pēc lieliem vēsturiskiem notikumiem, tajos sakāvi cietusī puse, jau <em>post faktum</em> analizējot situāciju, apstākļus, kas noveda pie šīs sakāves, nespēj būt objektīva. Raudzīties pagātnē no zaudētāju skatu punkta ir mazproduktīvi, jo viņi taču cieta sakāvi, tātad izrādījās netālredzīgi, nesaprata vai vismaz slikti saprata, kurp ved notikumu gaita, tai skaitā to notikumu, kurus izsauca viņu pašu aktivitātes. To vispirms var teikt par revolūcijās varu zaudējušajiem, taču tāpat arī attiecināt uz ar Otro pasaules karu saistītajos notikumos varu, ietekmi zaudējušajiem Polijas, un ne tikai Polijas, bijušajiem valdošajiem slāņiem.

Poļu zinātnieki par XX gadsimta otrās puses lielāko lūzumu uzskata 1989.~gada notikumus, kas ne tikai mainīja Polijas valsts raksturu, atjaunoja Polijas Republiku (\pltxti{Rzeczpospolita Polska}), bet arī atveseļoja situāciju Eiropā kopumā. Jāuzsver, ka šis lūzums, tāpat kā PSRS sabrukums 1991.~gadā, ir pavēris arī daudz labvēlīgākas iespējas padziļināti izpētīt Polijas vēsturi, tuvināties vēsturiskajai patiesībai. Tomēr tas nenozīmē, ka šī patiesība atklāsies pati no sevis. Pēdējās desmitgadēs pēc Polijas Tautas Republikas likvidācijas, kad oficiāla cenzūra vairs neeksistē, pie varas esošie spēki ir parādījuši, ka arī ar attiecīgu finanšu politiku var panākt, lai daudzi vēsturnieki, publicisti rūpētos ne tik daudz par objektīvās patiesības atklāšanu, cik politiskās elites pasūtītas vēstures interpretācijas izplatību. Daudzos mūsdienu poļu publicistu, arī zinātnieku, tāpat kā citu Austrumeiropas tautu vēsturnieku darbos, dominē politiskās nostādnes, nevis vēsturiskā patiesība.

2004.~gadā grupa poļu vēsturnieku izvirzīja iniciatīvu izstrādāt un realizēt aktīvu vēsturisko politiku (\pltxti{polityka historyczna}). Termins (\detxti{Geschichtspolitik}) bija aizgūts no Vācijas, kur tas radās jau XX gs. 80.~gados, tiesa, ar nedaudz citu saturu. Polijā vēsturiskās politikas realizācijas piekritēji vēsturi un vēsturisko atmiņu lielākoties uzskata par politiskās cīņas arēnu pret ārējo un iekšējo ienaidnieku. Ar to viņi faktiski attaisno atkāpes no profesionālās ētikas, arī domu brīvības ierobežošanu. Apriori tiek uzskatīts, ka „ārējais pretinieks” cenšas nostiprināt savu vēstures notikumu interpretāciju, tāpēc „savu” vēsturnieku pienākums ir solidāri pretoties šādām briesmām, aizstāvēt argumentus, pretējus pretinieku lietotajiem. Pie varas esošie politiskie spēki, izmantojot valsts administratīvos un finanšu resursus, cenšas nostiprināt tiem tīkamas vēsturisko notikumu interpretācijas. Valdošie slāņi cenšas uzurpēt un monopolizēt tiesības interpretēt vēsturi, tagadni uzskatot par mērlīdzekli pagātnes izvērtēšanai. Parasti speciālas vēsturiskas politikas realizēšanas nepieciešamībai kā attaisnojums tiek minēta nepietiekamā patriotisma audzināšana vēstures stundās skolā. Nacionāli noskaņotie radikāļi, lai stiprinātu „nacionālo pašapziņu”, prasa radīt „harmonisku” vēstures ainu, panākt vēsturiskās atmiņas vienveidību.

Tā 2005.--2007.~gadā politizēti publicisti mēģināja iestāstīt Polijas sabiedrībai, ka humanitārās zinātnes pēc 1989.~gada, tai laikā radītās mācību grāmatās esot nepietiekami patriotiskas, bezierunu patriotisma vietā izplatot kritisko patriotismu. Lai īstenotu vēsturisko politiku, tiek radītas speciālas valstiskas un sabiedriskas iestādes un muzeji, iesniegti likumprojekti lai iemūžinātu „vienīgi pareizo” vēsturisko notikumu traktējumu, ieviešot pat kriminālatbildību par tā neievērošanu, kontroli pār izglītības sistēmu. Izmantoti ievērojami finansiālie resursi dažādu politiski tendētu projektu īstenošanai, komisiju dibināšanai, augsti atalgotu amatu radīšanai, ieviešot vēsturisko „varoņu” panteonu, atceres dienas utt. Godājot un nereti arī pārspīlējot savus upurus, kas nesti cīņā pret „svešajiem”, bieži tiek tīšuprāt aizmirsts, ka savi noziegumi un to rezultātā nestie citu tautu upuri nevar tikt attaisnoti pat ar savas tautas interešu aizstāvību. Bieži politisko uzskatu maiņa ved pie tā, ka tiek izcelti vieni un noniecināti citi vēstures personāži. Piemēram, Poznaņā viena no garākajām ielām tagad nes nevis Jaroslava Dombrovska, bet Jana Henrika Dombrovska vārdu, jo pirmais mūsdienu „pilsētas tēviem” bijis pārāk kreiss.

Šāda vēsturiskās politikas īstenošana ved pie diskusiju ierobežošanas savā valstī un konfliktsituāciju radīšanas ar ārpasauli. Kritiskās pieejas trūkums pret savas zemes vēsturi un savu patriotismu traucē demokrātiskas politiskās kultūras izveidei, ved pie atteikšanās no plurālisma vēsturnieku darbos. Vēsture un vēsturiskā atmiņa ir svarīga kolektīvās identitātes sastāvdaļa, taču skatot to tikai caur vēsturiskās politikas prizmu, tiek tikai sarežģīta strīdu un problēmu atrisināšana. Poļu vēsturniece A.~Volfa-Poveska norāda, ka arī mūsdienu Polijas vēsturnieku un žurnālistu vidū netrūkst cilvēku, kuri demonstrē, ka var „taisīt” politiku ar vēstures palīdzību, ka poļu vēsturiskajai politikai ir ģeopolitiskas koordinātes, tā atrodas starp divu savu lielāko kaimiņvalstu vēsturiskajām politikām.

Vērtējot vēsturiskos notikumus, jāņem vērā arī tās pārvērtības, kas notiek ar apspiestām tautām, kad tās atgūst valstiskumu un kļūst par noteicējām zemē, kur dzīvo citas mazākumtautības. Bijušajiem cietējiem ir vajadzīgs patiešām augsts kultūras un demokrātiskās attīstības līmenis, lai no apspiestajiem nepārvērstos par apspiedējiem. Ne velti neatkarīgās Polijas valsts vēsturi XX gadsimta 20--30.~gados tik atšķirīgi vērtē, no vienas puses, poļi un, no otras puses, tajā dzīvojušie ukraiņi, baltkrievi, lietuvieši, vācieši, ebreji un krievi.

Vēstures jautājumiem ir jāpaliek par speciālistu un vēstures entuziastu meklējumu un diskusiju objektu, bet politiķiem jāmeklē ceļi valstu un tautu attiecību uzlabošanai. Kā atzinis Krievijas ZA loceklis A.~Čubarjans: „Attiecības starp cilvēkiem un starp tautām kļūst stiprākas, dziļākas un godīgākas, ja tās pamatojas uz vēsturiskās patiesības atzīšanu un gatavību kritiskai analīzei, arī savas vēstures pārvērtēšanai.” Acīmredzams, ka saprašanās starp kaimiņtautām~--- poļiem un krieviem~--- būs iespējama tikai tad, ja abas puses pilnībā apzināsies tās grūtības, kuras šai ceļā nosaka vēsturiskais mantojums, kā arī ievēros un cienīs pretējās puses tradīcijas un uzskatus. Mūsdienās poļu un krievu zinātnieki no Polijas un Krievijas Zinātņu akadēmijām cenšas radīt objektīvu Polijas un Krievijas sarežģīto attiecību versiju, taču joprojām abu grupu pieejā ir daudz subjektīvisma. Krievu akadēmiķis I.~Kovaļčenko norādījis: „Nacionālā vēsture~--- tā ir autobiogrāfija, bet katra autobiogrāfija ir subjektīva, tāpēc, lai piedotu nacionālajām vēsturēm objektīvu raksturu, to uzrakstīšanai ir jāpiesaista citu valstu vēsturnieki.”

Blakus pašiem poļu pētniekiem, kuri ne vienmēr ir spējīgi būt pietiekami kritiski pret savu valsti, daudz par Polijas vēstures problēmām rakstījuši arī tās kaimiņvalstu vēsturnieki. Tas gan atvieglo, gan sarežģī tās līdzsvarotāku izvērtēšanu. Atvieglo tāpēc, ka ļauj izzināt vienas puses noklusēto, sarežģī~--- jo liek meklēt patiesību bieži diametrāli pretējās faktu interpretācijās. Bieži vien vairāk nekā pašiem poļiem un arī Polijas kaimiņvalstu vēsturniekiem var uzticēties neieinteresēto ārzemju pētnieku, piemēram, šveicieša A.~Kapellera, kanādieša M.~Dž.~Karleija u.c. vērtējumiem.

XX un XXI gs. mijā starp Poliju un Krieviju attīstījās vēsturnieku sadarbība, notika savstarpēja apmaiņa ar dokumentiem un materiāliem. Apritē nāca un joprojām nāk arvien jauni arhīvu materiāli, taču faktu uzkrāšana nevar automātiski atrisināt to interpretācijas problēmu.

Starp citu, vēsturnieku aprindās bieži skan gaušanās, ka pārāk lēnu tiek atvērti vai arī nemaz netiek atvērti Krievijas arhīvi, kas vismaz daļēji atbilst patiesībai. Taču, kā norādījis krievu diplomāts, politiķis un vēsturnieks V.~Faļins, tai pat laikā tiek aizmirsts, ka ASV uz nenoteiktu laiku ir liegusi jebkādu pieeju dokumentiem, kurus amerikāņu karavīri Otrā pasaules kara beigās ieguva nacistiskās Vācijas vadības mītnē Tīringijā. Acīmredzot tie satur pārāk daudz nepatīkama materiāla par amerikāņu politiku Eiropā, varbūt arī Polijā.

Ir vajadzīgs laiks, lai atbrīvotos no aizspriedumiem un liekām emocijām, neizbēgamām īpaši XX gadsimta vētraino notikumu vērtēšanā. Diemžēl, starptautiskā situācija ne vienmēr ir labvēlīga patiesības noskaidrošanai. Tā, Rietumvalstu un Krievijas interešu pretstāve Ukrainas teritorijā, faktiskais NATO valstu un Krievijas hibrīdkarš, īpaši pēc tiešas militāras sadursmes sākuma 2022.~gada februārī, kurš objektīvi veicināja Polijas kā Ukrainas kaimiņvalsts lomas pieaugumu, tikai sekmē arvien jaunu vairāk vai mazāk būtisku pretrunu, kā arī jau minēto aizspriedumu un emociju saasināšanos.

Tikai patiesības noskaidrošana, falsifikāciju atmaskošana spēj radīt pamatu normālām, cieņas pilnām, draudzīgām attiecībām starp tautām. Tas attiecas arī uz poļiem, vāciešiem, krieviem, ukraiņiem, baltkrieviem, lietuviešiem, ebrejiem un arī latviešiem, kuri gan ilgstoši dzīvoja kaimiņos, bet kuru attiecības veidojās dažādu iekšpolitisku un ārpolitisku apstākļu ietekmē. Lai atceramies t.s. ''poļu laikus” Latvijas teritorijā, arī to, ka XX gadsimta divdesmitajos gados par Poliju Latvijā inteliģences aprindās skanēja novērtējums „panu Polija”. Profesors J.~Tazbirs vēl XXI gadsimta sākumā norādīja, ka no vairākuma Polijā dažādos izglītības līmeņos lietojamo vēstures grāmatu grūti uzzināt, ka daudzu gadsimtu gaitā Žečpospolita nacionālā ziņā atgādināja mozaīku. Kā atzinis E.~Duračinskis, arī mūsdienās daudzi poļi bijušās Polijas iedzīvotājus: ukraiņus, baltkrievus un lietuviešus joprojām uzlūko kā „mazākumtautības”, kuras nav pelnījušas neko vairāk kā Otrās Žečpospolitas realizētās politikas pret „nacionālajiem mazākumiem” uzlabotu variantu. Tāpēc arī šajā darbā nacionālajām attiecībām starp dažādajām Polijā un tās kaimiņos dzīvojošajām tautām ierādīta nozīmīga vieta.

Bijušā Polijas prezidenta B.~Komarovska padomnieks vēstures jautājumos T.~Nalenčs intervijā presei uzsvēra, ka valstij un sabiedrībai jākoncentrē uzmanība to vēsturisko motīvu aprites veicināšanai, ar kuriem Polija var lepoties, piemēram, jāpopularizē pasaulē lielākā Malborkas (poļu \pltxti{Malbork}, vācu \detxti{Marienburg}) pils vai Ļ.~Valensas, kurš stāvēja pie „jaunās Polijas” sākumiem, personība. Pēc T.~Nalenča vārdiem, vēsturi nedrīkst izmantot kā avotu negatīvu emociju formēšanai gan valsts iekšienē, gan ārpus tās. Piekrītot tam, ka vēsture jāizmanto pozitīvu ideālu audzināšanai, autoram gan jāpiezīmē, ka ar tās palīdzību nedrīkst arī redzēt un popularizēt tikai savas un tā noniecināt citu tautu sniegumus, kā arī „neredzēt” negatīvo savas valsts vēsturē.

Lasītāja priekšā esošā darba autors centies cik iespējams objektīvi un daudzpusīgi atspoguļot Polijas vēstures XIX un XX gadsimtā sarežģītās problēmas, taču uzreiz jāsaka, ka viņš tomēr piedāvā tikai vienu~--- savu redzējumu. Jāuzsver, ka mūsdienās ir ļoti grūti pētīt, analizēt, vienkārši rakstīt par XX gadsimta notikumiem, pilnībā atsakoties no sava subjektīvā vērtējuma. Cerams, ka pēc gadiem, kad kaislības būs norimušas, objektivitāti sasniegt būs vieglāk. Taču apstākļos, kad uz politiskās skatuves darbojas spēki, kuri sevi ar lielākām vai mazākām tiesībām pozicionē kā to vai citu XX gadsimtā darbojošos partiju, strāvojumu mantiniekus, prasību raudzīties uz XX gadsimta notikumiem „pilnīgi neitrāli”, droši var nosaukt par utopisku.

Godīgi jāpasaka, ka grāmatā visi vēsturiskās izšķiršanās brīži, notikumi, to veicēji vērtēti no iespējami demokrātiskāka, vardarbību izslēdzoša attīstības ceļa sasniegšanas viedokļa. Visdažādākajos nacionālajos strīdos autora personīgās simpātijas pieder apspiestajiem nacionālajiem mazākumiem, bet ne lielvalstiskas, bieži šovinistiskas nacionālās politikas realizētājiem. Tai pašā laikā autoram nav pieņemama to pagātnē (un dažkārt arī mūsdienās) apspiesto tautu nacionālistisko spēku darbība, kuri savu mērķu sasniegšanai nereti gatavi pārkāpt citu tautu, citas tautības cilvēku tādas pat tiesības uz brīvību un vienlīdzību, neapstājoties pat noziegumu priekšā.

Stādīts mērķis veicināt no aizspriedumiem un stereotipiem brīvu un līdzsvarotu izpratni par Polijas un arī visas Eiropas vēsturi Otrā pasaules kara gados, palīdzēt lasītājiem kritiski uztvert un izvērtēt dažādus apzināti un neapzināti tendenciozus sacerējumus par Poliju Otrajā pasaules karā, kuri pēdējās desmitgadēs izplatās Austrumeiropā.

Otrā pasaules kara notikumos autora simpātijas ir antihitleriskās koalīcijas sabiedroto pusē, kuri cīnījās par vispārcilvēciskiem mērķiem. Tāpēc viņš kritiski raugās uz tādiem Austrumeiropas politiķiem, kuri, piesedzoties ar „patriotiskiem” lozungiem, kara laikā faktiski atbalstīja nacistus, bet pēc tam mēģināja uzspiest šī vispasaules nozīmības izšķiršanās punkta vērtējumam savu šauri nacionālistisko skatījumu. Arī mūsdienās darbojas viņu ideju mantinieki.

Runājot par pēckara situāciju pasaulē, autors neuzskata par iespējamu staļinisko (arī pēc J.~Staļina nāves uzlaboto) „sociālisma” modeli identificēt ar to mērķi, kuru stādīja tā daļa cilvēces attīstības ceļu meklētāju, kuri par galveno uzskatīja sociālā taisnīguma ideālu sasniegšanu. Sociālisms kā sabiedriska iekārta, kurā cilvēks neekspluatē citu cilvēku, kur pastāv sabiedrisks, bet ne valdošās elites piesavināts īpašums uz ražošanas līdzekļiem, netika uzcelts ne PSRS, ne citās t.~s. „sociālisma” valstīs. Ražošanas līdzekļu koncentrācija valsts (faktiski~--- partijiski-birokrātiskā aparāta) rokās veda pie virsmonopolizācijas, kas bremzēja ražošanas attīstību. „Sociālisma” valstīs nepastāvēja sociālais taisnīgums. Pastāvēja pat superekspluatācija, kad daļa sabiedrības (staļiniskajās nometnēs ieslodzītā) praktiski nesaņēma nekādu atlīdzību, daļa strādāja par niecīgu atalgojumu, bet sabiedrības virsslānis~--- partijiski-birokrātiskā elite saņēma savam veikumam neadekvāti augstu atalgojumu. Sociālisma princips „no katra pēc spējām, katram pēc viņa darba”, „sociālisma” valstīs netika īstenots, jo nebija jau objektīvas mērauklas, kā šo darbu novērtēt. Tirgus vietā vērtējumu deva kā valsts birokrātijas slānis kopumā, tā atsevišķi tā pārstāvji. Tāpēc, runājot par PSRS, tās satelītvalstīm, arī Polijas Tautas Republiku, par tur it kā uzcelto sabiedrisko iekārtu autors parasti raksta kā par „sociālismu”, t.i.~--- lietojot pēdiņas. Tikai runājot par sociālismu kā mērķi, kuru pieņēma un centās sasniegt daļa uz labāku dzīvi cerošo, bet lielākoties nekritiski domājošo tautas masu, tas lietots bez pēdiņām.

Pēckara notikumos autora visdziļākā cieņa pieder cīnītājiem par demokrātiju un tautas labklājību, kuri sava mērķa sasniegšanai bija gatavi lietot dažādus cīņas līdzekļus, taču atteicās izmantot vardarbību, „šķirisko”, „sociālistisko” vai „nacionālo” interešu vārdā piekopt teroru pret savas un citu tautu piederīgajiem. Ar to viņi krasi norobežojās no saviem pretiniekiem~--- totalitāro, autoritāro, etnokrātisko varas sistēmu aizstāvjiem, kuriem savukārt visi līdzekļi bija derīgi it kā „valstisko”, „nacionālo”, „šķirisko”, „sociālistisko”, arī „demokrātisko” u.tml., bet faktiski~--- savu savtīgo mērķu sasniegšanai. No šādām pozīcijām arī skatīti gan poļu, gan citzemju darbinieki, viņu pārstāvētie sabiedriskie spēki, viņu ietekmētie notikumi.

Patiesam Polijas vēstures redzējumam lasītājs tuvināsies, salīdzinot dažādus pētījumus, kritiski vērtējot kā šo, tā arī citus Polijas vēsturei veltītus darbus.

Ilustrācijas ievadam

\chapter{Polija kaimiņvalstu varā. 1795~--- 1918}

\epigraph
{Karot poļi neprot. Taču dumpoties!}
{Hugo~Kollontajs (\pltxti{Hugo Kołłątaj})}

\epigraph
{Diemžēl mēs [poļi] neprotam strādāt! Kauties, lieliski cīnīties, nomirt, uz to vienmēr esam gatavi; taču uzcītīgi strādāt, ilgstoši, bez trokšņa un uzslavām, strādāt varbūt ne priekš sevis, tas mums ir par daudz.}
{Ludvika Plātere (\lttxti{Ludwika Plater})}

\epigraph
{Polija ir unikāla valsts ar tieksmi pēc impērijas, kuras tai nekad nav bijis.}
{Dmitrijs Kuļikovs (\rutxti{Дмитрий Куликов})}



\epigraph
{Nelaimīga ir valsts, kurai nav varoņu.~--- Nē! Nelaimīga ir tā valsts, kurai ir vajadzīgi varoņi.}
{Bertolts Brehts (\detxti{Eugen Bertolt Friedrich Brecht})}

\epigraph
{Lepnums mums izmaksā dārgāk nekā bads, slāpes un aukstums.}
{Tomass Džefersons (\entxti{Thomas Jefferson})}

\epigraph
{Starptautiskajā politikā morāles nav bijis, nav un nebūs.}
{Jakovs Kedmi (\hetxti{יעקב קדמ})}

\epigraph
{Nabags nav tas, kam maz pieder, bet tas, kurš daudz grib.}
{Angļu sakāmvārds}

\epigraph
{Ja visi vainīgi, neviens nav vainīgs.}
{Pēteris Krupņikovs}

\epigraph
{Vistālāk iet tas, kurš nezina, kurp iet.}
{Olivers Kromvels (\entxti{Oliver Cromwell})}



\epigraph
{Gandrīz visi dižie līderi savai dzimtenei ir maksājuši asins jūras.}
{Mihails Vellers (\rutxti{Михаил Иосифович Веллер})}

\epigraph
{Nācijai zaudēt savu valsti, savu patstāvību un neatkarību ir liela nelaime. To var salīdzināt ar smagu slimību, kas kropļo nācijas dvēseli.}
{Nikolajs Berdjajevs (\rutxti{Николай Бердяев})}

\epigraph
{Dažas tautas traģēdijas tiek uzvestas bez starpbrīžiem.}
{Staņislavs Ježijs Lecs (\pltxti{Stanisław Jerzy Lec})}

\epigraph
{Kādreiz dedzinātāji ir pārliecināti, ka tautai priekšā nesuši lāpu.}
{Staņislavs Ježijs Lecs (\pltxti{Stanisław Jerzy Lec})}

\epigraph
{Pēc neveiksmīgām revolūcijām vienmēr seko ienīstas un atriebīgas valdības.}
{Pjērs Buasts (\frtxti{Pierre Boiste})}

\epigraph
{Visas revolūcijas beidzas ar reakciju. Tas nav novēršams. Tas ir likums. Un jo negantākas un niknākas bijušas revolūcijas, jo stiprāka bija reakcija. Revolūciju un reakcijas nomaiņā ir kāds maģisks riņķis.}
{Nikolajs Berdjajevs (\rutxti{Николай Бердяев})}



\epigraph
{Neprasme pārciest nelaimi ir liela nelaime.}
{Bions no Borisfēnas (\eltxti{Βίων Βορυσθενίτης})}

\section{Poļu zemes XVII gadsimta beigās un XIX gadsimtā}

\epigraph
{Laimīga tā tauta, kurai ir garlaicīga vēsture.}
{Šarls Luijs de Monteskjē (\frtxti{Charles-Louis de Secondat, Baron de La Brède et de Montesquieu})}

\epigraph
{Vēsture ir vislabākais skolotājs, kuram ir paši sliktākie skolnieki.}
{Indira Prijadaršinī Gandija (\entxti{Indira Priyadarshini Gandhi})}

\epigraph
{Nabadzība noved pie revolūcijas, revolūcija pie nabadzības.}
{V.~Igo (\frtxti{Victor Marie Hugo})}

\epigraph
{Tur, kur ir divi poļi, pastāv trīs viedokļi.}
{Poļu paruna}

\epigraph
{Pateicoties vienprātībai aug mazas valstis, ķildu dēļ iet bojā lielvalstis.}
{Henriks Senkevičs (\pltxti{Henryk Adam Aleksander Pius Sienkiewicz})}



\epigraph
{Piesien kādu skrandu pie spieķa, turi to augstu, un tu redzēsi, cik daudzi sekos tai kā karogam.}
{Staņislavs Ježijs Lecs (\pltxti{Stanisław Jerzy Lec})}

\epigraph
{Lielā Kartāga veda trīs karus. Pēc pirmā tā joprojām bija spēcīga. Pēc otrā tā vēl bija apdzīvota. Pēc trešās tā vairs nebija atrodama.}
{Bertolts Brehts (\detxti{Bertolt Brecht})}

\epigraph
{Pat visa zeme nav vienas veltīgi izlietas asins lāses vērta.}
{Aleksandrs Suvorovs (\rutxti{Александр Васильевич Суворов})}

\epigraph
{Kad tiek meklēti kontrrevolūcijas panākumu cēloņi, no visām pusēm var saņemt parocīgu atbildi, ka X kungs vai pilsonis Y nodeva tautu. Šī atbilde var būt pareiza un arī nē \dots{} katrā gadījumā tā nekādi nepaskaidro, kā tas notika, ka tauta ļāva sevi nodot.}
{Fridrihs Engelss (\detxti{Friedrich Engels})}

\epigraph
{Naids ir baiļu sekas, mēs vispirms baidāmies un tikai pēc tam ienīstam.}
{Sirils Konolijs (\entxti{Cyril Vernon Connolly})}



\epigraph
{Patriotisms~--- tas ir „mīlu savu”, nacionālisms~--- „nīstu svešu”.}
{Šarls de Golls (\frtxti{Charles André Joseph Marie de Gaulle})}

\epigraph
{Patriotisms ir mīlestība pret savējiem, nacionālisms~--- naids pret citiem.}
{Rihards fon Vaiczekers (\detxti{Richard Karl von Weizsäcker})}

\epigraph
{Nacionālisms daudz vairāk asociējas ar naidu pret svešu tautu nekā ar mīlestību pret savējo.}
{Nikolajs Berdjajevs (\rutxti{Николай Александрович Бердяев})}

\epigraph
{Cilvēks, kas ienīst citu tautu, nemīl arī savējo.}
{Nikolajs Dobroļubovs (\rutxti{Николай Александрович Добролюбов})}

\epigraph
{Novests līdz galējam sasprindzinājumam, nacionālisms dzen postā tautu, kas tam ļāvusies, padarot šo tautu par cilvēces ienaidnieci, jo cilvēce vienmēr izrādīsies stiprāka par atsevišķu tautu.}
{Vladimirs Solovjovs (\rutxti{Владимир Сергеевич Соловьёв})}

\subsection{Polijas dalīšanas, Varšavas hercogiste un Vīnes kongress}


Poļu etnosa attīstības un poļu nācijas veidošanās apstākļi līdz XVIII gadsimta beigām bija samērā labvēlīgi. Poļu apdzīvotā teritorija bija kompakta, nesadalīta ar dabīgiem kalnu vai ūdens šķēršļiem. Senseno poļu zemju kodols gadsimtu gaitā ietilpa vienotā valstī. Tiesa, arī aiz šīs valsts robežām atradās vairākas poļu teritorijas, bet pati valsts bija daudzetniska, taču lielākā iedzīvotāju masa tajā runāja vienā~--- poļu valodā, piederēja vienai~--- katoļu ticībai. Mūsdienu poļu vēsturnieks J.~Tazbirs gan raksta, ka pēc aptuveniem aprēķiniem tikai ap 40\% Žečpospolitas (\pltxti{Rzeczpospolita Obojga Narodów}~--- Polijas-Lietuvas apvienotā valsts jeb Abu Tautu Republika) iedzīvotāju bija poļi, pie tam tikai daļai no tiem piemita nacionālā apziņa. (Pēc šī darba autora domām attiecībā uz šo laiku pareizāk būtu runāt par etnisko, nevis nacionālo apziņu.) Sākot ar Ļubļinas ūniju 1569.~gadā, kad tika izveidota Žečpospolita un Lietuvas lielkņaziste zaudēja savu krievisko politisko komponenti, Žečpospolitā sākās Pareizticīgās baznīcas apspiešana, rietumkrievu rakstiskās valodas izraidīšana no darbvedības, visa krieviskā vajāšana. Tā īpaši pastiprinājās pēc Brestas baznīcas ūnijas (1596), kad pastiprinājās katoļu ekspansija (Brestas ūnijas rezultātā daļa Žečpospolitas pareizticīgo garīdznieku pakļāva Romas pāvestam, saglabājot daļu pareizticīgo ritu un savu hierarhiju) tajās pareizticīgo austrumslāvu zemēs, kuras bija Žečpospolita sastāvdaļa.

XVII--XXVIII gadsimtā Polijas sabiedrība dalījās trijās kārtās: muižniecībā jeb šļahtā, sīkpilsoņos un zemniekos.

Šļahta bija karojošu kungu kārta, kura sevi krasi norobežoja no pārējām~--- zemākajām kārtām, saucot to piederīgos par liellopiem (\pltxti{bydło}). Šļahtai piederēja monopols uz varu, zemes īpašumu un sabiedrisko prestižu. XV--XVI gadsimtā izveidojās, bet nākamajos gadsimtos par visas šļahtas ideoloģijas sastāvdaļu kļuva t.s. sarmatisms. Tas balstījās pieņēmumā, ka šļahtas izcelsme ir saistāma ar sarmatiem (sena nomadu tauta, kura runāja indoirāņu valodā), kas antīkajā laikmetā pakļāvuši slāvu ciltis un izveidojuši virsslāni. Sarmatisms it kā pamatoja Žečpospolitas muižniecības tiesības norobežoties no etniski ``svešajiem''~--- slāvu un lietuviešu zemniekiem. Sarmatisma ideoloģijas galvenie elementi bija~--- šļahtas neierobežota brīvība, nacionālā augstprātība, ticība Polijas vēsturiskajai izredzētībai, ksenofobija (neiecietība pret svešo). Sarmatisms kultivēja īpašas parašas. Tā, ``īsts šļahtičs'' labāk mira badā, bet neaptraipīja rokas ar tam nepiedienīgo fizisko darbu.

Starp citu, poļu publicists un politiķis A.~Vasiļevskis atzīmējis, ka poļu kultūras vēsturē gadsimtu gaitā izveidojušās divas nacionālo problēmu aprakstīšanas skolas, kur katra no tām vadījās no savas patriotisma izpratnes. Sarmatiskais virziens identificēja patriotismu ar pašslavināšanu, ar nekritisku, visa, kas bija savs, attaisnošanu. Otrs~--- demokrātiski reformistiskais virziens nebaidījās teikt tautai rūgtu patiesību, bet patriotismu identificēja ar atklātu norādīšanu uz nacionālajiem trūkumiem, kaut šādas pārdomas par saviem grēkiem bija arī sāpīgas. Šī darba autors var piebilst, ka arī mūsdienu Polijas vēstures literatūrā šie virzieni turpina pastāvēt.

Žečpospolitas spilgtākā īpatnība bija nekur citur Eiropā neredzēts šļahtas~--- šī priviliģētā slāņa daudzskaitlīgums. XVIII gadsimta vidū gandrīz ik desmitais valsts iedzīvotājs bija šļahtičs. Kaut kārtas uzvedības noteikumi paredzēja visu augstdzimušo šļahtiču pilnīgu vienlīdzību, patiesībā tā bija visai iluzora. Dižciltīgo ģerboņu īpašnieku vidū lielai daļai jau nepiederēja nedz zeme, nedz dzimtcilvēki. Šie nemantīgie šļahtiči parasti pelnīja iztiku un pajumti pie šļahtas bagātās daļas~--- t.s. magnātiem (no latīņu \latxti{magnatus}~--- liels cilvēks, \latxti{magnatis}~--- dižciltīgs cilvēks, parasti zemes lielīpašnieks. Žečpospolitā pastāvēja ap 60 magnātu dzimtu), kuru politiskais svars bija atkarīgs no t.s. ``klientu''~--- viņu atbalstītāju šļahtas rindās skaita.

Viens no Polijas neveiksmju cēloņiem bija tas, ka jebkura spēku koncentēšana centrālās varas vadībā kādu tālu mērķu sasniegšanas vārdā rada kategoriskus iebildumus Polijas nacionālajā elitē. No vienas puses, XVII gadsimtā poļu aristokrātijai bija pilnīgi nesaprotami, kāpēc saspringt, kad viss ir pieejams: gan galmi, lai tur ``spīdētu'', gan teritorijas ar līdz pusvergu stāvoklim nospiestiem iedzīvotājiem austrumos, kurus varēja pērt un no kuriem varēja vākt nodevas. No otras puses, centrālās varas nostiprināšanās, pakļaujot sev magnātus un šļahtu, nebija pēdējo interesēs.

Piemēram, atbalsta trūkuma centrālai varai dēļ Polija XVII gadsimtā palika bez kara flotes, kas bija viena no galvenajām militāro panākumu atslēgām (Šī perioda Polijas karaļi mēģināja izveidot reālu floti, taču viņu mēģinājumi atkārtoti izgāzās līdzekļu trūkuma dēļ karaļa kasē, jo Polijas šļahta neredzēja nepieciešamību pēc flotes un atteicās paaugstināt nodokļus tās celtniecībai. Vēl esošie kuģi tika pārdoti 1641.--1643.~gadā, kas nozīmēja Žečpospolitas flotes beigas.)

Nākamajā~--- XVIII gadsimtā šļahta nesaprata, kāpēc tai būtu jāatsakās no \latxti{liberum veto} (no latīņu, brīvais veto jeb brīvais aizliegums~--- Žečpospolitas parlamentārās uzbūves princips, kurš ļāva jebkuram Seima loceklim pārtraukt jautājuma apspriešanu. Deputātam vajadzēja tikai skaļi izsaukties latīniski ``\latxti{Sisto activitatem!}'' (Es pārtraucu darbības) vai poliski ``\pltxti{Nie pozwalam!}'' (Es neļauju!) un jādod vienlīdzīgas tiesības visiem valsts iedzīvotājiem. Rezultātā no 1573.~līdz 1763.~gadam, kad sanāca apmēram 150 Seimu, aptuveni trešdaļā no tiem netika pieņemts lēmums. Protams, ar vienu faktora, lai arī tik svarīga kā šļahtas nostāja, nevar visu izskaidrot.

Abas pārejās kārtas bija galvenās nodevu maksātājas, bet bez politiskām tiesībām. Pie sīkpilsoņiem piederēja pilsētu iedzīvotāji, uz kuriem attiecās pilsētu tiesības un kuri guva pastāvīgus ienākumus no tirdzniecības un amatniecības. 1790.~gadā pilsētās dzīvoja ap 16\% iedzīvotāju (bez ebrejiem gan tikai 6\%). Zemniecība sastādīja ap 75\% visu iedzīvotāju, lielākā tās daļa (ap 65\%) dzīvoja uz šļahtai piederošas zemes, pārējie~--- baznīcai un karalim piederošas. Pilsonībai (buržuāzijai)~--- lieltirgotājiem, manufaktūru īpašniekiem, baņķieriem XVIII gadsimta beigās varēja pieskaitīt tikai dažus desmitus cilvēku. Zemnieki (ap 85--90\%) pildīja klaušas, pārējie maksāja renti~--- t.s. činšu (poļu \pltxti{czynsz}, vācu \detxti{Zins}~, no latīņu \latxti{census}~--- procents~--- Polijā, Lietuvā un arī Latgalē tā sauca zemnieku nodevas jeb renti naudas vai produktu veidā). Pēc feodāļu-muižnieku gribas pāreja no činša atpakaļ uz klaušu saimniecību nesastapa grūtības un bija bieža parādība. Uz 100 zemnieku saimniecībām bija 20, kurās tika apsaimniekoti 7 līdz 8 ha zemes, 62 saimniecībās tās bija mazāk, 16 bija bezzemes saimniecības.

Ārpus kārtām stāvēja garīdzniecība, pilsētu nabagi un beztiesiskie un diskrimenētie ebreji. Augstākā un vidējā garīdzniecība izcelsmes ziņā bija tuva šļahtai, bet zemākā (īpaši no XIX gadsimta otrā ceturkšņa)~--- zemniecībai. Tā kā reliģija bija cieši saistītā ar nacionālo pašapziņu, pēc Polijas neatkarības zaudēšanas liela garīdzniecības daļa kļuva par aktīvu nacionālās atbrīvošanās kustības spēku.

Tiesa, vienprātības par garīdzniecības vērtējumu nav. Piemēram, vācu autors R.~Vingendorfs pirms Otrā pasaules kara rakstīja, ka katoļu garīdzniecībai gan esot bijusi mazāk svarīga nacionālā ideja, cik iespēja to izmantot kā cīņas līdzekli pret pareizticīgo Krieviju un protestantisko Prūsiju. Krievu vēsturnieks V.~Djakovs savukārt norādīja, ka sociālie motīvi garīdzniecībai bija mazāk nozīmīgi nekā citiem iedzīvotāju slāņiem, un kritiskos momentos tā parasti pieslējās šļahtai, tikai atsevišķi indivīdi nostājās zemniecības un citu darba cilvēku pusē.

Citi minētie slāņi, kaut piedalījās vēsturiskajos notikumos (sacelšanās u.c.), kādu patstāvīgu sabiedriski-politisku pozīciju neieņēma.

Pārdzīvojusi ``ziedu laikus'' XVII gadsimtā, Žečpospolita XVIII gadsimtā nonāca ekonomiskā un politiskā panīkumā, cīnījās par savu izdzīvošanu. Žečpospolita soli pa solim tuvojās katastrofai. Konfrontācija ar kaimiņvalstīm un pārmērīgi lielas kara izmaksas izsūca zeltu no valsts kases. Daudzie kari: ar Krieviju par ukraiņu un baltkrievu zemēm, ar Zviedriju~--- par Livoniju, ar Prūsiju~--- par tās austrumu apgabaliem, Žečpospolitu novājināja. Pirmā Žečpospolita izrādījās vāja valsts, kas nespēja aizstāvēt savu neatkarību un teritoriālo integritāti, bet tajā pašā laikā tik ķildīga, lai visos apkārtējos kaimiņos pamodinātu dedzīgu vēlmi atbrīvoties no šāda kaimiņa. Krievijai gan Žečpospolita droši piesedza tās robežas no Eiropas problēmām, un tā nekad nebūtu piekritusi šīs bufervalsts iznīcināšanu, ja tā parādītu zināmu saprātīgumu un necenstos apspiest pareizticīgos un pretoties Krievijas ietekmei.

Nebeidzamie iekšējie kari pret pareizticīgajiem iedzīvotājiem Žečpospolitas austrumu zemēs veicināja panīkumu. Karos lija poļu asinis, graujot valsts cilvēcisko potenciālu.

Valsti nomocīja dažādu iekšējo politisko spēku savstarpējās cīņas, pretrunas starp priviliģētajiem katoļiem un nevienlīdzīgajiem protestantiem, pareizticīgajiem un jūdu ticības piekritējiem, laikmetam neatbilstošā, novecojusī valsts iekārta, kad katrs pārstāvniecības iestādes~--- Seima loceklis varēja izmantot t.s. \latxti{liberum veto} tiesības, nobloķējot vairākuma pieņemta lēmuma izpildi. Poļu paruna pat ar zināmu lepnumu konstatēja, ka valsts pastāv pateicoties nekārtībām: ``\pltxti{Polska nierzadem stoi}''.

Nebija vienotības poļu un lietuviešu izcelsmes muižniecības vidū. Poļu vēsturnieks H.~Visners ir uzsvēris, ka poļu valdošie slāņi kopš Jagelloņu dinastijas (poļu \pltxti{Jagiellonowie}, lietuv. \lttxti{Jogailaičiai}~--- valdīja Lietuvā un Polijā no XIV līdz XVI gadsimtam) valdīšanas sākuma plānoja vienota Polijas politiskā organisma izveidi ar Lietuvas kunigaitiju. XVIII gadsimtā poļu programma, balstoties uz gadsimtiem ilgo savienības tradīciju, vairs neredzēja vietu suverēnai Lietuvas valstij. Arī vēlāk~--- jau XIX gadsimtā poļi nesaskatīja, ka lietuvieši rada savu etnosociālu un etnopolitisku kopību. Lietuviešu vēsturnieks un politologs A.~Kulakauskas secinājis, ka poļi uz XVIII gadsimta pārmaiņām Žečpospolitas teritrorijā, kuru viņi uzskatīja par savu īpašumu, reaģēja vēsturiski nepareizi [t.i~--- neatbilstoši savām tālākajām interesēm.~--- V.Š.]. Viņi neatzina lietuviešu tautu par pilnvērtīgu, uzskatīja Lietuvu par Polijas provinci un lietuviešus kā nacionālo mazākumu, kuri nedrīkst pat izvirzīt prasību pēc kulturālas autonomijas. Lai aizstāvētu savas tiesības pastāvēt kā patstāvīgai tautai savā valstī lietuvieši \lttxti{volens nolens} (gribot negribot) bija spiesti izvēlēties savu nacionālās attīstības ceļu, norobežojoties no poļiem.

Vēl viena problēma bija augstais ebreju iedzīvotāju īpatsvars Žečpospolitā. Ir aprēķini, ka 1800.~gadā 70\% visu pasaules ebreju dzīvoja Polijā un aptuveni 25\% bijušās Žečpospolitas iedzīvotāju bija ebreji, šeit izveidojās bagāts ebreju kultūras mantojums. Taču autohtonās etniskās grupas, kuras izsenis dzīvoja šai teritorijā, nekādi nevēlējās atzīt ebrejus par līdztiesīgiem.

Etniskās pretrunas savijās ar kārtu pretišķībām. Augstdzimusī ``poļu tauta'' (\pltxti{narod Polski}), kā viduslaikos sauca poļu šļahtu, centās turēt paklausībā vienkāršo ``kalpu tautu'' (\pltxti{narod chlopskie}).

Dažādās pretrunas Žečpospolitā labi saskatīja tālaika gaišākie Eiropas prāti. Jau Ž.~Ž.~Russo rakstīja: ``Lasot poļu valdīšanas vēsturi, ar grūtībām var saprast, kā tik dīvaini uzbūvēta valsts varēja pastāvēt tik ilgi.'' Arī lielais Apgaismības klasiķis Voltērs izteicās: ``\frtxti{Un Polonais~--- c´est un charmeur; deuz polonais~--- une begarre; trois polonais, eh bien, c´est la question polonaise.}'' (Viens polis ir apburošs cilvēks, divi poļi~--- tracis, trīs poļi, nu jā, tas jau ir Polijas jautājums.). Pagāja nedaudzi gadu desmiti, un virknes karu izpostītā Žečpospolita pat nespēja pati sevi aizstāvēt.

Tālākais Polijas liktenis, valsts sadale saistījās ar Eiropas valstu pūliņiem saglabāt spēku līdzsvaru starptautiskajā arēnā. Kompromisus varēja atrast uz tai laikā vājās Polijas rēķina. Kaimiņvalstis~--- Austrija, Prūsija un Krievija, izmantojot konfesionālās nesaskaņas, t.s. disidentu (no latīņu \latxti{dissidens (disidentis)}~--- tāds, kas nepiekrīt. Polijā, kur valdīja katolicisms, tā sauca kristiešus, kuri neatbalstīja valdošo konfesiju) cīņas pret katoļu privilēģijām, iejaucās valsts iekšējās lietās, traucēja nepieciešamo politisko reformu pieņemšanu, kas varētu stiprināt valsti.

Analizējot Žečpospolitas sadalīšanas iekšējos cēloņus, ievērojamais krievu vēsturnieks S.~Solovjovs pētījumā ''Polijas krišanas cēloņi'', (\_rutxti{Соловьёв С.М. История падения Польши, Москва, 1863}) uzsvēra, ka pirmajā vietā starp galvenajiem tās katastrofas cēloņiem liekami ne kaimiņvalstu agresīvie centieni, bet gan spēcīgā krievu (tai skaitā baltkrievu un ukraiņu) nacionālās atbrīvošanās kustība ''zem reliģiskā karoga'' pret poļu jūgu, par savām tiesībām, par vienlīdzību. Domājams, jāpiekrīt S.~Solovjova uzskatam, ka pašu poļu vaina bija tā, ka Polijas pareizticīgo vidū radās disidentu kustība, kuru galvenokārt atbalstīja Krievija, bet ne tikai. Polijas katoļu vairākums pat negribēja dzirdēt par atteikšanos no savām privilēģijām un tiesību vienlīdzību ar nekatoļiem un neuniātiem. Krievu vēsturnieks norādīja, ka 1653.~gadā, kad Maskavijas cara Alekseja Mihailoviča sūtnis kņazs B.~Repņins no Polijas valdības pieprasīja, lai pareizticīgie krievu cilvēki Žečpospolitā neciestu reliģiskos spaidus, pēdējās valdība nepiekrita šai prasībai, un sekas bija Mazkrievijas atkrišana no Žečpospolitas. Pēc simts gadiem Krievijas ķeizarienes Katrīnas~II vēstnieks, arī kņazs [N.]~Repņins, izteica tādu pašu prasību, bet saņēma atteikumu, un rezultāts bija pirmā Polijas dalīšana. S.~Solovjovs ar to gan vienpusīgi vienkāršoja vēsturisko procesu, taču būtisku tās aspektu atspoguļoja. Aizbildniecība pār t.s. disidentiem, vienlīdzīgu viņu tiesību ar katoļiem kā krievu tautā vispopulārākās lietas aizstāvība, bija īpaši svarīga Katrīnai~II,~--- rakstīja cits krievu vēsturnieks V.~Kļučevskis. Sākotnēji runa pat vairāk bija par, kā mūsdienās teiktu, cilvēktiesību aizstāvības politiku, nevis par Krievijas valsts teritoriālo paplašināšanu un liel-, maz- un baltkrievu apvienošanu Krievijas impērijas robežās.

Krievijai labvēlīgu politiku gan ieturēja 1764.~gadā ar tās atbalstu par Žečpospolitas karali ievēlētais viens no Krievijas imperatores Katrīnas~II favorītiem S.~Poņatovskis. Karalis Staņislavs~II Augusts sāka savu politisko darbību, patiesi ticot, ka tikai ar Krievijas palīdzību Polijā var īstenot viņa iecerētās reformas. Viņš atbalstīja poļu rūpniecības attīstību, tirdzniecības kompāniju dibināšanu, daudz vērības veltīja literatūrai un zinātnei. Reizē viņš mēģināja ierobežot magnātu patvaļu un nostiprināt centrālo varu. Taču drīz pierādījās, ka tāda Polijas attīstība ir nevēlama ne tikai lielai daļai poļu feodāļu, bet arī kaimiņivalstīm. Krievija pieprasīja atrisināt disidentu nevienlīdzības jautājumu. Tikai 1768.~gadā milzīga Krievijas spiediena ietekmē Seims bija spiests atzīt pareizticīgo vienlīdzību ar katoļiem Žečpospolitā. 24.~februārī tas pieņēma jaunu ``mūžīgu'' miera līgumu ar Krieviju, kurā garantēja disidentiem toleranci un vienlīdzību, bet poļu valdošajam slānim tādas ``tiesības'' kā brīvas karaļa velēšanas un \latxti{liberum veto}. Tomēr poļu virsslānis šādu vienlīdzību nepieņēma. Tam vienlīdzība tiesībās ar citticībniekiem bija līdzvērtīga visu poļu brīvību zaudēšanai. Ar līgumu neapmierinātā daļa garīdzniecīnas, magnātu un šļahtas Baras cietoksnī izveidoja konfederāciju (\pltxti{Konfederacja barska}, 1768--1772) un uzsāka sacelšanos. V.~Kļučevskis to raksturoja kā ''poļu šļahtiču ''Pugačova dumpi'' \citespace{} apspiedēju laupīšanas gājienu par tiesībām uz apspiešanu''. Izraisījās pilsoņu karš, kas savukārt izsauca kaimiņvalstu intervenci.

1772.~gadā Austrija un Prūsija, kuras baidījās, ka Krievija varētu patstāvīgi sagrābt poļu un lietuviešu zemes, iniciēja \strong{pirmo Žečpospolitas sadali}, kuras rezultātā tā šķirās no vairākām svarīgām teritorijām. Cariskā Krievija sākotnēji iebilda pret sadali, cenšoties panākt visas Žečpospolitas pakļaušanu savai ietekmei, jo pie savām rietumu robežām tai bija izdevīgāk saglabāt vāju Žečpospolitu nekā spēcīgu Prūsiju, kura jau 1648.--1721.~gadā bija pievienojusi sev daļēji poļu apdzīvotās Rietumu Pomorjes jeb Pomerānijas \pltxti{(poļu Pòmòrzé}, vācu \detxti{Pommern}) un 1740.~gadā Silēzijas (poļu \pltxti{Śląsk}, vācu \detxti{Schlesien}) teritorijas.\footnote{Par jau pirms Polijas trijām dalīšanām Prūsijas rokās nonākušajām poļu apdzīvotajām teritorijām šajā darbā ies runa tikai kopsakarā ar citām, XVIII~gs. otrajā pusē Žečpospolitā ietilpstošajām zemēm.} Taču kad starptautiskā situācija (1768.~gadā Turcija pieteica karu Krievijai, kurš ilga līdz 1774.~gadam, draudēja arī Austrijas iesaistīšanās tajā pret Krieviju) virzīja Krieviju uz savienību ar Prūsiju, tā piekrita sadalei. Katrīna II raksturoja Polijas dalīšanas cēloņus no sava skatu punkta: ``Šīs [poļu] tautas nepastāvības, tās ļaunprātības un naida pret mūsu [tautu], nemitīgas tieksmes uz izvirtību un franču negantībām rezultātā, mēs nekad tajā neatradīsim nedz mierīgu, nedz drošu kaimiņu, vienīgi kā novedot to būtiskā nespēkā un nevarenībā''.

Austrija ieguva Austrumgalīcijas (bieži literatūrā apgabalu sauc arī vienkārši par Galīciju) ar Ļvovu (ukraiņu \uktxti{Львів}, poļu \pltxti{Lwów}, vācu \detxti{Lemberg}) un t.s. Mazpoliju (\pltxti{Małopolska}~--- apgabals Polijas dienvidaustrumos Vislas vidus un augštecē, tiek saukts arī par Rietumgalīciju. Vēsturiski galvenā pilsēta~--- Krakova (Krakow), tikai bez Krakovas, kura palika Žečpospolitā. Mūsdienās Austrumgalīcija atrodas Ukrainas, bet Mazpolija jeb Rietumgalīcija Polijas teritorijā. Lai aizmaskotu Polijas sadali, Austrija atcerējas, ka kādreiz (XIV gadsimtā) Austrijas pakļautā Ungārija valdīja pār Galīciju un tāpēc esot notikusi ``atkalapvienošanās''. 1774.~gadā Austrijā pat izkala medaļu ar uzrakstu ``\latxti{Antigua jura Vindicata Galicia et Lodomeria in Fidem recepetis MDCCLXXIII}'' (Senās Galīcijas un Lodomērijas tiesības atgūtas 1873.~gadā.

Prūsija, kura sevišķi bija ieinteresēta apvienot divas līdz tam atsevišķi pastāvējušās savas valsts daļas: Brandenburgu ar Pomerāniju un Austrumprūsju, saņēma Rietumprūsiju (vācu~--- \detxti{Westpreußen}~--- apgabals Vislas upes lejteces abos krastos), gan bez Dancigas (vācu \detxti{Danzig}, poļu \pltxti{Gdansk}) un Toruņas (poļu \pltxti{Toruń}, vācu \detxti{Thorn}), Kujāvijas (poļu \pltxti{Kujawy}, vācu \detxti{Kujawien}) reģiona ziemeļdaļu (Vislas kreisajā krastā) un daļu Lielpolijas (poļu \pltxti{Wielkopolska}~--- apgabals Polijas rietumos Vartas (poļu~--- \pltxti{Warta}, vācu~--- \detxti{Warthe}) upes baseinā, kuru agrāk apdzīvoja visļanu un poļanu ciltis). 1773.~gadā Prūsijas karalis Fridrihs~II vēstulē Voltēram, tēlojot šo provinču šķietamo atpalicību, rakstīja: ``Poļu provinces nevar salīdzināt ne ar vienu no Eiropas valstīm. Augstākais, tās ir salīdzināmas ar Kanādu. Būs jāpieliek daudz pūļu un laika, lai panāktu to, kas garos gadsimtos sliktas pārvaldes dēļ ir nokavēts.'' Prūsija tagad ar Žečpospolitai atņemto piejūras apgabalu savienoja Austrumprūsiju ar pārējo valsts teritoriju. Ar to Polija zaudēja arī savus svarīgākos ārējās tirdzniecības ceļus, bet Prūsija ieguva iespēju kontrolēt vairāk nekā piektdaļu Polijas ārējās tirdzniecības, gūt no muitas nodevām, ar kurām tika apliktas pa Vislu vestās poļu preces (pie Dancigas muitas nodoklis sastādīja 12\% no preču vērtības), ienākumus, kuri pēc dažiem vērtējumiem bija lielāki par pārpalikušās Polijas ienākumiem.

Krievija savukārt ieguva Latgali ar \lttxti{Dinaburgu} (tagadējo Daugavpili) un Austrumbaltkrieviju ar Polocku, Vitebsku un Mogiļevu un t.s. Melno Krieviju (lietuviešu \lttxti{Juodoji Rusia}~--- agrākās Lietuvas lielkņazistes daļu Daugavas labajā un Berezinas upes kreisajā krastā). Prūsijas karalis Fridrihs~II atzina: ``\dots{}Krievijai ir daudz tiesību tā rīkoties ar Poliju, ko gan nevar teikt par mums ar Austriju''. Karaļa teiktais prasa komentāru. Zemes, ko sagrāba Prūsija, bija pārsvarā poļu apdzīvotas. Austrijai piešķirtajās zemēs dzīvoja galvenokārt rietumukraiņi un poļi. Vāciski runājošo tajās bija nedaudz. No šī viedokļa Krievijai, pievienojot sev slāvu apdzīvotas zemes, bija uz tām vairāk tiesību nekā tās sabiedrotajām. 1772.~gadā, ko poļi atzīmē kā Polijas pirmo dalīšanu, ievērojama Baltkrievijas daļa tika atbrīvota no poļu jūga, un Krievijai šis notikums bija ne mazāk nozīmīgs kā daļas Mazkrievijas (Ukrainas) atbrīvošanās 1654.~gadā no tāda pat jūga un apvienošanās ar Krieviju. Krievu vēsturnieki to vērtē kā faktisku visu trīs bijušo senkrievu atzaru~--- baltkrievu, lielkrievu un mazkrievu apvienošanu vienas Krievijas valsts ietvaros pēc vairāku gadsimtu šķelšanās. Par iespēju atbrīvot daļu ticības brāļu~--- baltkrievu no katoļu diskriminācijas, par viņu pievienošanu Krievijai gan nācās piešķirt brīvas rokas Prūsijai un Austrijai attiecībā uz citām poļu zemēm.

Taču jau padomju vēsturnieki uzsvēra, ka nedz Prūsija, nedz Austrija neuzdrošinājās sagrābt Žečpospolitas teritorijas, kamēr to nebija sankcionējusi cariskā Krievija. Tāpēc carismam tāpat kā tā sabiedrotajiem jānes pilna atbildība par 1772.~gada Polijas dalīšanu, kas nesa prūšu un austriešu varas nodibināšanu pār to valdījumos nonākušajām poļu un ukraiņu tautas daļām.

Jau tūlīt pēc pirmās Polijas dalīšanas neapmierinātie poļu šļahtiči uzsāka cīņu pret carisko Krieviju. Tā, daži pret Krievijas ietekmi Polijā karojošās un sakautās šļahtiču Baras konfederācijas dalībnieki, kuri tika izsūtīti uz Krievijas iekšieni, J.~Pugačova vadītā dumpja jeb Zemnieku kara (1773--1775) laikā pieslējās tam.

Austrija kopumā saņēma ap 83 tūkstošus km$^{2}$ ar 2,6 miljoniem iedzīvotāju. Prūsija kopā ieguva ap 36 tūkstošus km$^{2}$ teritorijas ar 580~000 iedzīvotāju. Krievijas guvums bija ap 92 tūkstoši km$^{2}$ ar 1,3 miljoniem iedzīvotāju. (Literatūrā sastopami arī nedaudz citādi dati. Tautas skaitīšana, kuru veica Poliju dalījušās valstis līdz 1776.~gadam, rādīja, ka Prūsijas iedzīvotāju skaits pieaudzis par 0,6~miljoniem, Austrijas~--- 2,1~miljonu, Krievijas~--- 1,3~miljonu cilvēku.)

Teritoriālās pārmaiņas apstiprināja Žečpospolitas Seims 1773.~gadā. Tikai Novogrudokas (baltkrievu \betxti{Навагрудак}, poļu \pltxti{Nowogródek}) deputāts T.~Rejtans pret to protestēja. XIX~gadsimtā poļu sabiedriskā doma viņu pacēla nacionāla varoņa augstumos. Izcilais poļu mākslinieks J.~Matejko šim notikumam veltīja vienu no savām gleznām.

Tādejādi pirmajā dalīšanā Žečpospolita šķīrās no vairāk nekā 28\% savas teritorijas un vairāk nekā trešdaļas tās iedzīvotāju. Tomēr apmēram 527~tūkstošu km$^{2}$ plašais Polijas valsts pārpalikums vēl joprojām bija tikpat liels, kā Francija vai Lielbritānija un lielāks nekā 1918.~gadā radītā Polijas valsts. Pēc pirmās dalīšanas lielākā daļa poļu, visas lietuviešu, daļa ukraiņu un baltkrievu zemju joprojām palika Žečpospolitas sastāvā. Pēc poļu vēsturnieka T.~Korzona vērtējuma pārpalikušajā Žečpospolitā dzīvoja 7,4~miljoni iedzīvotāju. (Tādejādi, pirms pirmās dalīšanas Žečpospolitā bija ap 11,4~miljoni iedzīvotāju).

Valsts sadalīšana izsauca poļu sabiedrībā šoku, patriotisma uzplūdus un vēlmi modernizēt valsti. Polijas vēstures pētnieks I.~Balabans gan rakstīja, ka pēc pirmās dalīšanas ``visi saprata, ka ir nepieciešamas reformas'', taču tad jāmin arī poļu vēsturnieka M.~Bobržinska konstatējums, ka Polijas tautu valsts dalīšana ne pārāk satrauca. Lielāku ietekmi atstāja Apgaismības kustība Eiropā, vēlāk~--- Franču revolūcijas sākums (1789). Norisa reformas izglītības un administratīvās pārvaldes jomā. Tā, Polijā pēc karaļa Staņislava Augusta priekšlikuma tika radīta laicīga Izglītības komisija (\pltxti{Komisja nad Edukacją Młodzi Szlacheckiej Dozór Mająca}, 1773--1794) pirmā Eiropā iestāde, kura pēc savām funkcijām līdzinājās visu valsti aptverošai izglītības ministrijai. Pēc Prūsijas parauga tika reorganizēta armija. Tomēr pietiekami lielas un modernas armijas radīšanu, kura varētu nodrošināt Polijas politisko neatkarību, nepieļāva valsts finansiālais vājums. (Pēc M.~Bobržinska datiem vairāk kā puse Žečpospolitas ienākumu tika tērēta armijas uzturēšanai.)

Savukārt finansiālo vājumu lielā mērā noteica valsts politiskā iekārta. Absolūtie monarhi kaimiņvalstīs daudz brīvāk varēja aplikt savus pavalstniekus ar nodokļiem nekā Polijas Seima deputāti savus vēlētājus. Kā rakstīja poļu vēsturnieks J.~Rutkovskis, galvenais traucēklis enerģiskas finansiālās politikas realizācijai bija Seimu ekonomiskā pozīcija. Seimu deputāti pārstāvēja šļahtas intereses, kura vēlējās savās rokās koncentrēt maksimāli lielu nacionālā ienākuma daļu. Agrārpolitikā šļahta aizstāvēja dzimtbūšanas saglabāšanu. Rūpniecības un tirdzniecības jomā Seimi pirmkārt rūpējās par iespēju pārdot lauksaimniecības ražojumus par iespējami augstākām un iepirkt rūpniecības preces par iespējami zemākām cenām. Tas deva šļahtai iespēju celt savu dzīves līmeni. Zemnieki nekādu labumu no šīs politikas neguva. Amatnieki par savu darbu tika atalgoti arvien mazāk. Tas traucēja amatnieciskās ražošanas pāraugšanu par kapitālistisko ražošanu.

Ap 1780.~gadu sabiedriskā aina Polijā gan bija mainījusies uz labo pusi. Neraugoties uz teritoriālajiem zaudējumiem un tirdzniecības ierobežojumiem, saimniecība Žečpospolitā attīstījās. Ja 1776.~gadā valstī ieveda preces par 48 miljoniem zlotu, bet izveda tikai par 22 miljoniem, tad 1785.~gadā eksports (pamatos lauksaimnieciskā produkcija) pirmoreiz ilgu gadu laikā pārsniedza importu. Pieauga valsts ienākumi. Lauksaimniecības ražotā galvenā prece iekšējā tirgū bija rudzi, ārējā~--- kvieši. Tomēr pakāpeniski graudkopība zaudēja savas pozīcijas. No XVIII gadsimta beigām sākās kartupeļu audzēšana. To izplatība iezīmēja progresu lauksaimniecības preču ražošanā, jo kartupeļus izmantoja pārtikā, lopbarībā un kā izejvielu pārtikas rūpniecībā. No tehniskajām kultūrām plaši tika audzēti lini un kaņepāji. Pastāvēja manufaktūras un rūpali. Paplašinājās saražoto preču sortiments. Sākās pāreja uz jaunu kurināmā veidu~--- akmeņoglēm. XVIII gadsimta 90.~gados darbojas ap 280~uzņēmumu, kuri izmantoja galvenokārt algotu darbaspēku. Dzimtcilvēki raka kanālus, kuri savienoja Baltijas un Melnās jūras baseina upju augšteces, kas ļāva daļu tirdzniecības novirzīt uz Melno jūru. Auga pilsētas kā ārējās tirdzniecības centri~--- Gdaņska (poļu \pltxti{Gdańsk}, vācu \detxti{Dancig}) un Poznaņa (poļu \pltxti{Poznań}, vācu \detxti{Posen}). Pēc jau minētā T.~Korzona aprēķiniem, kopš 1775.~gada iedzīvotāju skaits pieauga par 1,4 miljoniem un sasniedza 8,8 miljonus. Varšavas iedzīvotāju skaits 1791.~gadā sasniedza 120~tūkstošus. Tiesa, iekšējā tirdzniecība bija attīstīta vēl vāji. Izplatītākā tās organizācijas forma bija vietējie un gadatirgi. Galvenais šaurā tirgus cēlonis bija zemnieku saimniecību pusnaturālā rakstura saglabāšanās, brīvā tirgus attiecību vājā iespiešanās laukos. Zemnieki sastādīja 72,7\%, pilsētnieki-kristieši~--- 6,8\%, ebreji, kuri mita galvenokārt pilsētās un miestos,~--- 10,2\%, šļahta~--- 8\%, garīdzniecība~--- 0,6\%, pārējie~--- 1,7\%.

Uzplauka kultūra. Karalis Staņislavs II Augusts (Poņatovskis) bija Apgaismības un rokoko mākslas atbalstītājs. Viņa ciešie kontakti ar Drēzdeni veicināja poļu mūzikas dzīvi. Karalis uzstājās arī kā mākslas mecenāts un veicināja celtniecības darbus Varšavā. Viņa laikā uzbūvētā ievērojamākā celtne ir Karaļa pils jeb ``Pils uz ūdens'' (\pltxti{Pałac Na Wodzie}) Lazenkos (\pltxti{Łazienki}).

Taču valsts tautsaimniecības un kultūras attīstību traucēja dzimtbūšana, brīva darbaspēka trūkums, arī novecojusī politiskā iekārta.

Demokrātiski-patriotiskais virziens iestājās par mantojamas monarhijas ieviešanu, \latxti{liberum veto} tiesību likvidēšanu. Kad 1787.~gadā sākās jauns Krievijas~--- Turcijas karš, kas saistīja Krievijas spēkus, poļos radās cerība izmantot situāciju, lai likvidētu Polijas atkarību no tās. Krievu vēsturnieks S.~Solovjovs šai sakarā rakstīja: ``Redzēja [poļi] Krieviju apgrūtinātā stāvoklī~--- un gribēja to izmantot; nespēja to izmantot lai iedvestu jaunus spēkus paralizētajā [Polijas] valsts ķermenī, toties pilnībā izbaudīja prieku iespert lauvam, netikuši skaidrībā, ka lauva ne tikai nav tuvs nāvei, bet nav pat saslimis.'' Kopā ar Krieviju pret Turciju karoja arī Austrija. Savukārt Prūsija baidījās no Austrijas un Krievijas nostiprināšanās.

Krievijas valdniece Katrīna II cerēja iesaistīt arī Žečpospolitu savienībā pret Turciju, 1788.~gadā sanāca konfederatīvs (šeit~--- tāds, kurā lēmumus pieņēma ar vairākumu balsu) Seims (t.s. Četrgadu Seims~--- \pltxti{Sejm Czteroletni}, 1788--1792), lai apspriestu šo jautājumu. Pret Žečpospolitas savienību ar Krieviju iebilda Prūsijas karalis Fridrihs Vilhelms II, kurš baidījās no Krievijas nostiprināšanās kara rezultātā. Seims atteicās noslēgt līgumu ar Krieviju. Apstākļos, kad Eiropā brieda lieli notikumi (1789.~gada 14.~jūlijā Parīzē tika ieņemta Bastīlija), patriotiskie poļu spēki Seimā uzsāka cīņu par valsts politisku un ekonomisku reformēšanu. Viens no patriotu pārstāvjiem filozofs un rakstnieks S.~Stašics 1790.~gadā griezās pie sabiedrības ar uzsaukumu ``\pltxti{Przestroi dla Polsi}'' (``Brīdinājums Polijai''), kurā bija vārdi: ``Cik tālu Polija ir atpalikusi! Polijā ir sācies tikai XV gadsimts, kad pārējā Eiropā beidzas XVIII gadsimts!''

To, ka šī atpalicība tomēr bija visai nosacīta, rādīja \strong{1791.~gada 3.~maijā} Četrgadu Seima pieņemtā jaunā valsts \strong{Konstitūcija}, kura likvidēja Žečpopolitas konfedaratīvo raksturu, Lietuvas lielkņazistes nosacīto patstāvību, ieviesa franču filozofa un rakstnieka Š.~Monteskjē proponēto varas dalīšanas principu, nodalot likumdošanas, izpildvaru un tiesu varu, pasludināja t.s. pilsoniskās brīvības. Polijas Konstitūcija iedibināja arī citus jaunievedumus: kontrasignācijas (no latīņu \latxti{contra} pret + \latxti{signare} parakstīt) principu, kad valdības vadītājs vai atsevišķs ministrs līdzās valsts galvam parakstīja tā izdotus aktus, tā uzņemdamies politisku un juridisku atbildību; parlamentārās (politiskās) atbildības principu un neuzticības votumu (no latīņu \latxti{votum}~--- vēlēšanas, lēmums, domas, kas izteiktas balsojot) ar sekojošu ministru atsaukšanu no amata. Tika noteikta Seima kā divpalātu parlamenta uzbūve (deputātu un senatoru palāta, ar deputātu palātas izšķirošu nozīmi), Seima sēžu kārtība, likumdošanas procedūra un likumprojektu balsošanas kārtība, kvoruma (latīņu \latxti{quorum}~--- nepieciešamais sapulces dalībnieku skaits, lai tā būtu pilntiesīga) princips. Tāda parlamenta struktūra darbojas arī mūsdienās.

Šļahtas un garīdzniecības privilēģijas gan lielā mērā saglabājās, bet, likvidējot \latxti{liberum veto} un nostiprinot Seima tiesības pieņemt lēmumus ar balsu vairākumu, tika mazināta feodālā anarhija, magnātu varenība, paplašinātas augošās pilsonības (buržuāzijas) tiesības. Kā norādīja poļu vēsturnieks un valsts darbinieks H.~Jablonskis, patriotiskajām aprindām, kuras centās glābt Žečposplitu, bija skaidrs, ka, lai saglabātu jau satricināto Polijas neatkarību, pirmkārt bija jācenšas iznīcināt magnātu varu. Tāpēc 3.~Maija Konstitūcijas ieviestajai nemantīgās šļahtas, kura sastādīja magnātu galveno balstu, politisko tiesību ierobežošanai bija progresīva nozīme, kaut vecās iekārtas aizstāvji skaļi vaimanāja par ``šļahtiskās demokrātijas'' iznīcināšanu.

Polijas 1791.~gada 3.~Maija Konstitūcija bija neapšaubāms solis uz priekšu, tā bija otrā rakstītā satversme pasaulē pēc ASV 1787.~gada Konstitūcijas un pirmā Eiropā (Francijā 1791.~gada Konstitūcija tika pieņemta 3.~septembrī.) Franči Polijas satversmi izmantoja, izstrādājot savas valsts 1791., 1793. un 1795.~gada Konstitūcijas. Neraugoties uz ierobežotību, Polijas Konstitūcijai bija progresīvs raksturs un tās pieņemšanas gadadiena Polijā tiek svinīgi atzīmēta vēl joprojām. Četrgadu Seima veikums un īpaši jaunā Konstitūcija izsauca līdzjūtīgu vērtējumu ``visā Eiropā''. Tomēr, kaut arī Polijā tika ieviesta konstitucionāla monarhija, Konstitūcija bija novēlota, tā būtiski nemainīja sabiedrisko iekārtu, valstī joprojām dominēja feodālā kārtība, pilsonība palika ļoti vāja. Turpināja pastāvēt dzimtbūšana un ebreju īpašais nelīdztiesīgais stāvoklis.

Pie tam jaunās satversmes pieņemšanas ar karaļa piekrišanu diena bija izraudzīta tad, kad tās pretinieki Seimā vēl atradās brīvdienās un nevarēja piedalīties balsošanā. Pēc vācu vēsturnieka E.~Meijera datiem bija ieradušies tikai ap 100 no 500 Seima deputātiem. Ar to Konstitūcijas likumību varēja apšaubīt un tās pieņemšanu pielīdzināt valsts apvērsumam. Pie tam Konstitūcijas nostiprināšanai tika draudēts lietot arī nekonstitucionālus līdzekļus. Tā, 13.~maijā visās Varšavas ielās bija piestiprināti drukāti nezināmas izcelsmes uzsaukumi, kuri aicināja nokaut katru, kurš runāja, rakstīja, pretojās vai arī gatavojās to darīt pret 3.~maija konstitūciju. Katram tādam slepkavam tika solīta atlīdzība. Policija gan uzsaukumus noplēsa. Tomēr visi vietējie seimeļi (\pltxti{sejmik}) atzina Konstitūciju.

Taču izveidojās Žečpospolitai nelabvēlīga starptautiskā situācija. Ja poļu reformatori cerēja izmantot Austrijas un Prūsijas nesaskaņas, viņi vīlās. Pēc Austrijas ķeizara Jozefa II nāves 1790.~gada 20.~februārī jau 27.~jūlijā Austrija un Prūsija savstarpēji izlīga, bet 1791.~gada augustā noslēdza vienošanos pret Franciju. Mainījās Krievijas nostāja. Vēl 1790.~gada 12.~oktobrī Krievijas imperatore Katrīna II paziņoja ``Tā kā mēs raugāmies uz Poliju kā uz valsti, kura atrodas starp četrām spēcīgām valstīm [Prūsiju, Krieviju, Austriju, Turciju] un kalpo kā šķērslis to strīdiem, tad šo šķēršļi ir jāsaglabā un jāsargā, lai arī ko tas mums maksātu. Un mēs par to parūpēsimies. Taču tikai līdz brīdim, līdz mūsu ienaidnieku un pašas Polijas ienaidnieku ļaunie nodomi nepiespiedīs mūs mainīt mūsu politiku.'' Taču pēc krievu--turku kara (1787--1791) beigām 1792.~gada maijā Katrīna~II, kura baidījās no Polijas nostiprināšanās un tās mēģinājumiem atjaunot 1772.~gada robežas, kā arī revolucionārā ``franču mēra'' izplatīšanās, bija gatava iejaukties Polijas iekšējās lietās, aicinot poļus uz pilsonisku nepakļaušanos.

Saprotams, ka, baidoties par savām pozīcijām, magnāti un no tiem atkarīgie ``klienti'' meklēja Katrīnas~II atbalstu. Ievērojamais vācu vēsturnieks H.f.~Treičke par toreizējo situāciju Polijā rakstīja: ``Pār malām plūstošs cīņas prieks un gatavība upurēties, dedzīgas runas un brālīgi apskāvieni, vaimanājoši priesteri un augstprātīgas, skaistas sieviete, pie tā klāt punšs un mazurka, cik tik sirds vēlas, bet arī savstarpējais partiju ienaids, nepaklausība, dusmīgas apsūdzības un drošsirdīgu, sajūsminātu vīru viļņošanās vidū nevienas valstiskas galvas, neviena liela rakstura.'' Krievijas piekritēji 1792.~gada maijā organizēja t.s. Targovices konfederāciju (\pltxti{Konfederacja targowicka}, pēc miesta nosaukuma), kura pasludināja jauno Konstitūciju par spēkā neesošu un ar krievu karaspēka palīdzību vērsās pret karali un Seimu. 1792.~gada 18.~maijā Krievijas karaspēks iemaršēja Polijas teritorijā. Atsaucoties uz 1791.~gada 3.~maija poļu ``revolūciju'', to pašu darīja arī Prūsijas armija. (Interesanti, ka vācu vēsturnieks H.f.~Zitzevics šo rīcību salīdzinājis ar PSRS un tās sabiedroto, arī Polijas, karaspēka ieiešanu Prāgā 1968.~gadā.)

Poļu karaspēku komandēja karaļa brāļa dēls J.~Poņatovskis un Amerikas atbrīvošanās kara dalībnieks, ASV pilsonis ģenerālis \strong{T.~Kostjuško} u.c.

T.~Kostjuško bija dzimis mūsdienu Baltkrievijas senā, bet ne pārāk turīgā augstmaņu dzimtā. Iespējams, ka viņa pirmā valoda bija baltkrievu, un viņš tika kristīts gan pareizticībā, gan katolicismā. 1769.~gadā Kostjuško tika piešķirta karaļa stipendija mācībām Parīzē. Šeit viņš kā eksterns mācījās Parīzes kara akadēmijā. Parīzē pavadītie pieci gadi būtiski ietekmēja T.~Kostjuško vēlākos uzskatus. Piedalījās Amerikas neatkarības karā. 1776.~gada augustā ieradās ASV, dažus mēnešus vēlāk kļuva par Kontinentālās armijas galveno inženieri un vadīja Filadelfijas fortifikācijas darbus, sadraudzējās ar T.~Džefersonu. Pēc septiņiem gadiem armijas dienestā Kongress 1783.~gada 13.~oktobrī T.~Kostjuško piešķīra brigādes ģenerāļa dienesta pakāpi. Viņam tika piešķirta ASV pilsonība, īpašumi. Tomēr 1784.~gadā viņš devās atpakaļ uz Poliju, bet savu īpašumu novēlēja izmantot melno vergu izpirkšanai un izglītošanai. 1784.~gada augustā K.~Kostjuško ieradās Polijā un apmetās savā dzimtas īpašumā. Viņš būtiski uzlaboja savu dzimtcilvēku stāvokli, atvieglodams klaušas, tādejādi izpelnīdamies liberāļa slavu.

Poļi guva panākumus Zeleņces kaujā (\pltxti{Bitwa pod Zieleńcami}). Pēc kaujas karalis Staņislavs II Augusts nodibināja \latxti{Virtuti Militari} (kara nopelnu) ordeni, ar kuru apbalvoja J.~Poņatovski un T.~Kostjuško. Arī mūsdienās šis ordenis ir Polijas augstākais militārais apbalvojums. Tomēr tā kā atbilstoši stratēģiskajai situācijai poļi pārspēka priekšā turpināja atkāpties, krievu avotos kauja skaitījās kā viņu uzvarēta. Turpmāk Krievijas armija vairākās kaujās sakāva poļu un lietuviešu spēkus. 1792.~gada jūlijā arī Polijas karalis Staņislavs~II Augusts pievienojās konfederātiem, izdeva rīkojumu par savas armijas atlaišanu. Krievijas karaspēks ieņēma Varšavu.

1793.~gada janvārī Krievija un Prūsija, kuras vienoja bailes no revolucionārajiem notikumiem Francijā, to ietekmes izplatības Polijā, (Prūsija vēl arī nevēlējās pieļaut Krievijas robežu tālāku izplatību uz rietumiem) parakstīja slepenu vienošanos par jaunu, jau \strong{otro Polijas dalīšanu}. Austrija, kura tajā laikā bija cietusi vairākas sakāves karā ar Franciju, kā arī centās sagrābt Bavāriju, Polijas otrajā dalīšanā nepiedalījās. Tās noteikumi poļiem tika paziņoti 1793.~gada 27.~martā.

Krievija saņēma Rietumbaltkrieviju ar Minsku (baltkr. \betxti{Мінск)}, Labā krasta Ukrainu ar Žitomiras (ukraiņu \uktxti{Житомир}) pilsētu, Volīnijas (ukraiņu \uktxti{Волинь}, poļu \pltxti{Wołyń}) austrumu daļu un daļu Podolijas (poļu \pltxti{Podole}, ukr. \uktxti{Поділля}, \uktxti{Podilla}, krievu \rutxti{Подолье}), ar Braslavas (poļu \pltxti{Wrocław}, vācu \detxti{Breslau}, krievu \rutxti{Бреславль}) pilsētu, Prūsija~--- t.s. Lielpolijas apgabalu ar Gņezno (poļu \pltxti{Gniezno}, vācu. \detxti{Gnesen}), Poznaņas, Gdaņskas un Toruņas pilsētām. Katrīna~II, kura jau 1792.~gada decembrī Krievijas sūtnim Polijā J.~Siversam rakstīja, ka Krievija nolēmusi ``atbrīvot kādreiz Krievijai piederējušās, tās tautiešiem (\rutxti{единоплеменники}) un vienai ticībai piederīgajiem apdzīvotās zemes un pilsētas no viņus apdraudošās kārdināšanas un apspiešanas'', tagad par godu ukraiņu un baltkrievu apdzīvoto zemju ``atgriešanai'' Krievijas sastāvā lika izkalt speciālu medaļu. Tās viena pusē bija viņas pašas profils, otrajā~--- Krievijas ērglis, kurš savieno agrāko Krievijas un pievienoto zemju kartes ar uzrakstu ``\rutxti{Отторженная возвратихъ}'' (Atņemtā atgriešana).

Mūsdienās baltkrievu vēsturnieki lielākoties uzskata, ka tikai pateicoties baltkrievu zemju iekļaušanai Krievijas impērijās sastāvā tika pārtraukta to katolizācija un polonizācija, kas ļāva baltkrieviem saglabāties kā etnosam. Tieši Krievijas impērijas sastāvā noformējās patstāvīga baltkrievu nācija.

Žečpospolitas teritorija samazinājās vēl divas reizes. Tagad tā aptvēra vairs tikai 230~000 km$^{2}$ ar 4,4~miljoniem iedzīvotāju. Īpaši sāpīgs Žečpospolitai bija Gdaņskas pievienošana Prūsijai, jo ar to tika zaudēta arī pieeja jūrai. Valstiski-tiesiskās attiecības starp Krieviju un Žečpospolitu noteica traktāts, ko tās parakstīja 1793.~gada oktobrī. Pēc traktāta noteikumiem abas valstis apņēmās, ka, uzbrukuma gadījumā vienai no tām, otra palīdzēs tai ar visiem spēkiem. Apvienoto karaspēku komandēt tiesības saņemtu valsts, kura piešķirtu vairāk karaspēka. Praktiski tas nozīmēja, ka poļu karaspēks tika pakļauts krievu pavēlniecībai, jo Polijai bija tiesības turēt tikai 15~000 lielu armiju. Krievija saņēma tiesības vajadzības gadījumā ievest savu karaspēku Polijā, tiesības apstiprināt tās ārpolitiskos līgumus. Tika atcelta arī 3.~maija Konstitūcija. Faktiski pēc otrās dalīšanas Žečpospolita zaudēja savu neatkarību. Tā no iespējamā Krievijas pretinieka kļuva par tās vasaļvalsti.

Žečpospolitas sadalīšanu u.c. noteikumus apstiprināja t.s. Grodņas Seims 1793.~gada rudenī. Panākt šī lēmumu pieņemt Seima deputātus Krievijas sūtnis J.~Siverss centās ar kukuļiem, draudiem, tiešu spiedienu (Grodņas pils bija pilna krievu kareivju). Kad Seima maršals prasīja nobalsot par līgumu, deputāti vairākas stundas klusēja. Galu galā atskanēja Krakovas deputāta J.~Ankviča balss: ``Klusēšana ir piekrišana''. Tad nu Seima maršals paziņoja, ka līgums pieņemts vienbalsīgi. Krievu vēsturnieks S.~Solovjovs uzsvēra, ka šādu notikumu gaitu noteica gadsimtiem ilgā poļu tautas klusēšana, kad trokšņoja vienīgi šļahta seimos.

Taču poļu patriotu cīņas griba vēl nebija sagrauta, cerot uz revolucionārās Francijas palīdzību, viņi slepus gatavoja \strong{sacelšanos}. Par savu vadītāju viņi izvirzīja T.~Kostjuško, kurš sevi jau bija pierādījis par drosmīgu karavadoni Amerikas brīvības cīņās un arī Polijā. 1794.~gada 16.~martā Krakovas iedzīvotāji pasludināja T.~Kostjuško par republikas diktatoru un nacionālo bruņoto spēku virspavēlnieku. T.~Kostjuško komandētie poļu spēki sasniedza līdz 70~000 cilvēku (kopumā, ieskaitot mobilizētos pilsētniekus un zemniekus, tautas miliciju, poļu vienības sasniedza vairākus simtus tūkstošus cilvēku), taču bija ļoti vāji apbruņoti. 1794.~gada 4.~aprīlī T.~Kostjuško vadītajām vienībām izdevās gūt panākumus kaujā pie Raclaviciem (\pltxti{Bitwa pod Racławicami}), kur krievu karaspēks uzbruka ciešā ierindā, bet T.~Kostjuško, izmantojot pieredzi, gūtu ASV Neatkarības karā, poļu strēlniekus izvietoja izklaidus, izmantojot dabīgos aizsegus. Pats T.~Kostjuško izveda ar izkaptīm bruņoto zemnieku vienības krievu daļu aizmugurē, kur tie negaidītā straujā uzbrukumā sagrāba artilēriju un piespieda krievu kājniekus atkāpties. Taču iznīcināt tos neizdevās un kara darbība turpinājās.

Poļu panākumam vēlāk tika veltīta J.~Matejko glezna. 1894.~gadā Lembergā tika atklāta 114 metru gara, 15 metru augsta kaujas panorāma ar 38 metru diametru. 1944.~gadā kara gaitā panorāma tika daļēji bojāta un to sāka atjaunot tikai 1980.~gadā, bet 1985.~gadā atklāja Vroclavā. Nosauktie mākslas darbi joprojām kalpo romantizētam poļu vēstures izklāstam.

Pēc uzvaras pie Raclaviciem sacelšanās sākās arī citviet. Sacēlušos rokās nonāca Varšava un Viļņa. Varšavas sacelšanās gaitā liela daļa krievu garnizona gāja bojā, (viens bataljons~--- ap 500 karavīru atradās baznīcā, protams, neapbruņots. Sacēlušies ielauzās baznīcā un lielāko daļu karavīru nogalināja), otrai izdevās atstāt pilsētu. Kā rakstīja I.~Balabans, pēc krievu garnizona iznīcināšanas ``asiņainajiem skatiem'' ``franču revolūcijas piemēra iejūsminātā tauta'' izvilka no mājām arī ``Tēvzemes nodevējus'' un pakāra uz laternu stabiem. Vēsts par briesmu darbiem Varšavā sasniedza krievu armijas daļas, radot tajās spēcīgas atriebības jūtas.

Savā laikā padomju vēsturnieki V.~Djakovs un I.~Millers PSRS un Polijas sakariem veltītā rakstu krājumā rakstīja, ka ``humānā un labvēlīgā sacēlušos attieksme pret viņu gūstā nonākušajiem krievu virsniekiem un karavīriem lika tiem iestāties revolucionārās armijas rindās''. Šiem krievu karavīriem poļi nelika karot pret saviem tautiešiem, bet nosūtīja uz divīziju, kura atvairīja prūšu uzbrukumu. Nenoliedzot šādu faktu esamību, domājams, ka visiem revolucionāriem simpatizējošie padomju autori sacēlušos humānismu pārspīlēja un apžēloti kā reiz tika tikai tie, kuri izteica gatavību karot poļu pusē.

T.~Kostjuško stāvokli sarežģīja nepieciešamība rēķināties kā ar kustības rojālistisko, tā revolucionāro spārnu, kurš Varšavā jau organizēja revolucionāros tribunālus un politisko pretinieku pakāršanu. Piemēram, tika sodīts ar nāvi Viļņas bīskaps I.~Masaļskis, Lietuvas hetmanis J.~Zabello u.c. Vācu vēsturnieks H.f.~Zitzevics šo linča tiesu (angļu~--- \entxti{the Lynch law}, aizdomās par noziegumiem turēto cilvēku nogalināšanu bez tiesas un izmeklēšanas) pret īstiem vai iedomātiem sacelšanās pretiniekiem pastāvēšanu, kuras nobiedēja mēreno elementus, pat minējis kā tās sakāves galveno cēloni.

T.~Kostjuško mēģinājumi iesaistīt cīņā arī plašākus zemnieku spēkus izsauca šļahtiču neapmierinātību. 1794.~gada 7.~maijā Poļaņicas (\pltxti{Połaniec}) pilsētas apkārtnē tika pasludināts manifests (\pltxti{Uniwersał połaniecki}, poļu \pltxti{uniwersal}~--- vēstījums visiem), kurš prasīja, lai zemniekiem, kuri bija samaksājuši parādus un izpildījuši virkni citu nosacījumu, tiktu piešķirta personīgā brīvība, lai klaušu dienu skaits tiktu samazināts, lai zemes īpašnieki vai tās pārvaldnieki atbildētu tiesas priekšā par zemnieku apspiešanu kā vainīgi vēlmē pazudināt nacionālās sacelšanās lietu utt. T.~Kostjuško biedēja šļahtu, ka Maskava grib pret to sacelt poļu zemniekus, norādot uz to smago dzīvi un solot Katrīnas II labvēlību. Šļahta bija sašutusi par šādu mēģinājumu graut viņu īpašuma tiesības uz zemniekiem un universālam nebija praktisku seku, bet tikai simboliska nozīme.

Kad T.~Kostjuško nosūtīja savu pārstāvi uz Parīzi lūgt palīdzību Sabiedriskās glābšanas komitejai (\frtxti{Comité de salut publicē}, 1793--1795), tā apšaubīja poļu revolucionaritāti un atteica palīdzību. Komiteja poļu pārstāvim uzdeva jautājumus: ``Kā izskaidrot, ka jūsu Kostjuško, tautas diktators, cieš sev blakus karali, kuru pie tam, kā būtu jāzin Kostjuško, tronī iesēdināja Krievija? Kā izskaidrot, ka jūsu diktators aiz bailēm aristokrātu priekšā, kuri nevēlas atteikties no ``darba rokām'', neuzdrošinājās veikt zemnieku masu mobilizāciju? Kā izskaidrot, ka viņa proklamācijas zaudē savu revolucionāro nokrāsu, attālinoties no Krakovas? Kā izskaidrot, ka viņš nekavējoties apspieda karātavām tautas sacelšanos Varšavā? Atbildiet!'' Poļu pārstāvim nācās klusēt.

Toties pret T.~Kostjuško armiju karoja arī Austrijas un Prūsijas bruņotie spēki (visai neaktīvi, kaut daži vācu vēsturnieki apgalvo, ka tieši Prūsijas karaspēka pievienošanās Krievijas armijai noteikusi kara iznākumu).

Sacelšanās apspiešanā piedalījās arī slavenais krievu karavadonis A.~Suvorovs. 28.septembrī poļu galvenie spēki tika smagi sakauti, T.~Kostjuško ievainots un saņemts gūstā.

Pēc nostāstiem, tad viņš izteicis vārdus ``\latxti{Finis Poloniae}'' (no latīņu: ``Beigas Polijai'', literatūrā ir arī cits šo vārdu variants: ``\latxti{Finis regni Poloniae''}~--- ``Beigas Polijas karaļvalstij.''). Pats T.~Kostjuško vēlāk, 1803.~gadā gan rakstīja, ka viņš bija smagi ievainots (kājā un galvā), gūstā krita jau bez samaņas, ko atguva tikai pēc divām dienām, un nekādus pravietiskus vārdus nespēja pateikt. Viņš arī neuzskatījis sevi par pēdējo poli, ar kura nāvi Poljai pienāktu beigas. Pēc citas versijas, kad Varšavā uzzināja par notikušo, tad gan atskanējuši saucieni; ``Nav Kostjuško! Beigas tēvzemei!''.

Interesanti, ka T.~Kostjuško vadībā cīnījās arī jauns virsnieks M.~Oginskis, kurš pēc sakāves bija spiests atstāt dzimteni. Pēc vienas mūzikas zinātnieku versijas tieši 1794.~gadā arī radās viņa pasaulslavenā polonēze ``\pltxti{Pożegnanie Ojczyzny}`` (Atvadas no dzimtenes), saukta arī par ``\pltxti{Polonez Oginsky}'' (Oginska polonēzi). (Pēc citas versijas polonēze tapa pēc 1820.~gada, kad M.~Oginskis jau bija Krievijas imperatora Aleksandra I amnestēts, saņēma atpakaļ konfiscētās muižas, ieņēma senatora amatu, bet, neapmierināts ar imperatora politiku, devās uz Itāliju.)

Pēc tam A.~Suvorova vadībā krievu karaspēks devās uzbrukumā Varšavai 24.~oktobrī, ieņemot tās priekšpilsētu Prāgu (\pltxti{Praga}).

Uzbrukumā piedalījās arī tie pulki, kuri bija izcietuši poļu pēkšņo sacelšanos Varšavā. Cīņas Prāgā bija asiņainas. A.~Suvorovs ziņojumā par kauju rakstīja: ``Pārvarot visas grūtības un uzveicot pretinieka sīvo aizstāvēšanos trijos nocietinājumos, mūsu karaspēks ielauzās Prāgā. Briesmīga bija asinsizliešana, katrs solis uz ielām bija nosegts ar kritušajiem. Visi laukumu noklāti ķermeņiem, pēdējā un pati briesmīgākā iznīcināšana notika Vislas krastā, ko redzēja Varšavas tauta. Šis briesmīgais skats iedvesa viņiem šausmas…''

Poļu vēstures literatūrā uzsvērts, ka krievu karavīri zvērīgi izrīkojās ar civiliedzīvotājiem. Polijā stingri nostiprinājies viedoklis par mierīgo poļu iedzīvotāju ``masu slaktiņu''. Poļu vēsturnieki Varšavas piepilsētas civiliedzīvotāju nogalināšanu piemin bieži, neminot gan daudzus līdzīgus gadījumus tā laika karu un sacelšanos vēsturē. Šis sižets ietverts daudzās periodam veltītajās monogrāfijās un arī skolai domātās mācību grāmatās, kaut autoru vidū nav vienprātības par slaktiņa mērogiem. Tiek runāts par ``pilnīgu'' izkaušanu, ``ievērojamas daļas'' apslaktēšanu un vienkāršu ``slaktiņu''. Piemēram, ievērojamais poļu publicists, vēsturnieks, vēlāk arī diplomāts L.~Vasiļevskis rakstīja, ka Krievijas karaspēks ``apkāva un noslīcināja Vislā vairāk nekā 10 tūkstošu sieviešu un bērnu''. Viedoklis ir pārvērties par aksiomu, kuru Polijā vairs neviens neapstrīd. Jāatzīmē gan, ka mūsdienu krievu vēsturnieks A.~Širokorads atzīmējis, ka neviens no poļu vēsturniekiem nav devis atbildi uz jautājumu, kāpēc, ilgus mēnešus gatavojoties pilsētas aizsardzībai, poļu vadītāji neevakuēja Prāgas iedzīvotājus vismaz pāri Vislai un neizmitināja pārējo varšaviešu namos.

PSRS šī epizode netika pieminēta nedz vēsturnieku darbos, nedz uzziņu literatūrā. No pirmspadomju Krievijas vēsturniekiem par šo kara epizodi visplašāk rakstījis N.~Kostomarovs, kurš secināja, ka poļu stāsti par šo epizodi ``neiztur kritiku''. Atsaucoties uz krievu ``tā laika avotiem'', vēsturnieks raksta, ka pavisam gāja bojā ap 12 tūkstošu poļu (karavīru un civiliedzīvotāju), pie tam daudzi, ``glābjoties no krievu durkļiem'' noslīka Vislā, gūstā krita ap 1~tūkstotis cilvēku. Pēc krievu vēsturnieka D.~Bantiša-Kamenska datiem, Prāgas ieņemšanā tika nogalināti 13,5 tūkstoši poļu karavīru, ap 11,5~tūkstošu tika saņemti gūstā, līdz diviem tūkstošiem noslīka Vislā, cenšoties tai pārkļūt, ap tūkstotim tas izdevās. No 22~tūkstošiem krievu karavīru, kas piedalījās uzbrukumā Prāgai, 580 tika nogalināti un 960 ievainoti. (1894.~gadā Sanktpēterburgā Krievijas ģenerālštāba pulkveža N.~Orlova izdotā darbā par Prāgas ieņemšanu poļu kritušo skaits minēts no 9 līdz 10 tūkstošiem, ievainoto no 11 līdz 13 tūkstošiem, krievu zaudējumi: 300--450 nogalināto un ap 2~000 ievainoto.) Jau poļu-krievu spēku samērs vien liecina, ka, ja poļu karavīri kaut cik drosmīgi cīnījās, krievu karavīriem nebija laika nodarboties ar civiliedzīvotāju ``slaktēšanu''.

Memuāros stāstīts par briesmīgiem nežēlības gadījumiem, kādi sastopami gandrīz katra cietokšņa ieņemšanas gaitā, gan arī uzvarētāju augstsirdības gadījumiem. Ir atrodamas publicista F.~Bulgarina pierakstītas krievu ģenerāļa I.~fon Klugena atmiņas, kurš stāstīja: ``Uz mums šāva no ēku logiem un jumtiem, un mūsu karavīri, ielaužoties mājās, nogalināja visus, kas viņiem gadījās ceļā \citespace{} virsnieki vairs nespēja apturēt asins izliešanu \citespace{} Pie tilta atkal sākās slaktiņš. Mūsu karavīri šāva pūļos, neņemot vērā p~--- un spalgie sieviešu kliedzieni, bērnu brēcieni stindzināja dvēseli \citespace{} Saniknotie mūsu karavīri katrā dzīvā radībā redzēja mūsējo slepkavnieku Varšavas sacelšanās laikā. ``Nav nevienam piedošanas''~--- kliedza mūsu karavīri un nogalināja visus, neievērojot ne gadus, ne dzimumu. \citespace{} Četrās stundās notika briesmīga atriebība par mūsējo apslaktēšanu Varšavā.'' Šeit gan jāņem vērā, ka Prāgas aizstāvju vidū bija vairāki tūkstoši brīvprātīgo vietējo iedzīvotāju, kuri upuru noteikšanas gaitā varēja tikt pieskaitīti civiliedzīvotājiem, kaut karoja ar ieročiem rokās. Prāgas ieņemšanas laikā bija arī vēsturnieku minētā epizode, kad daļa tās aizstāvju izlauzās līdz Vislai un krievu karavīri tos, kuriem neizdevās upi pārpeldēt, iznīcināja kreisā krasta iedzīvotāju acu priekšā. Šie fakti varēja kalpot par pamatu leģendai par civiliedzīvotāju ``slaktiņu''. Jāatzīst, ka krievu karavīru vidū, kuriem bija izdevies sacelšanās sākumā atstāt Varšavu, dzīvas bija atmiņas par tur zvērīgi nogalinātājiem ieroču biedriem, valdīja atriebības jūtas, kuras acīmredzot atsevišķos gadījumos varēja izpausties arī pret civilistiem, taču tad nu jārunā par savstarpējiem ``slaktiņiem''.

Prāgas ieņemšana nobiedēja Varšavas iedzīvotājus, daļa metās bēgt projām no pilsētas, citi nosūtīja deputātus pie karaļa, pieprasot kapitulāciju. Viņš atbalstīja lūgumu un pie A.~Suvorova tika nosūtīta Varšavas maģistrāta delegācija. Kad daži poļu virsnieki pirms tam mēģināja ar spēku izvest no pilsētas karali un iepriekš saņemtos krievu gūstekņus, lai cīņu turpinātu, pilsētnieki to nepieļāva. Tā paši Varšavas iedzīvotāji nepieļāva sacelšanās dalībnieku mēģinājumus turpināt cīņu. Kapitulācija tika pieņemta.

Varšavas maģistrāts kopā ar sālsmaizi pasniedza pilsētas atslēgas A.~Suvorovam. 26.~oktobrī viņa armija iegāja pilsētā. Kad A.~Suvorovs novērsa iespējamo pilsētas izlaupīšanu, atbrīvoja no gūsta sākotnēji 6~000 poļu zemessargu, pēc tam vēl arī 500 sagūstīto poļu virsnieku, Varšavas maģistrāts iedzīvotāju vārdā viņam uzdāvināja zelta tabakas dozi ar briljantiem un uzrakstu ``Varšava~--- savam glābējam''. Imperatorei Katrīnai~II par Varšavas ieņemšanu A.~Suvorovs ziņoja vēstulē, sastāvošā no trim vārdiem: ``Urā! Varšava mūsu!'' Katrīnas~II atbilde bija tikpat īsa: ``Urā! Feldmaršals Suvorovs!'' Tātad par savu veikumu Polijā A.~Suvorovs saņēma feldmaršala pakāpi, kas apliecina ne tikai viņa militāro spēju novērtējumu, bet arī to nozīmi, kādu Katrīna~II piešķīra Polijas pakļaušanai.

Daļa poļu patriotu, kam izdevās izglābties no Krievijas karaspēka, gan mēģināja vēl turpināt cīņu, bet drīz tika galīgi sakauti. Var atzīmēt, ka poļu šļahtiči 1795.~gadā neveiksmīgi mēģināja izraisīt sacelšanos ukraiņu un baltkrievu zemnieku vidū ar lozungu ``Par mūsu un jūsu brīvību'' (\pltxti{Za naszą i waszą wolność).} Daļa poļu bēga uz ārzemēm un no šī laika tur pastāvīgi dzīvoja poļu emigranti. 1796.~gadā Krievijas cars Pāvils I apžēloja T.~Kostjuško un līdz ar viņu arī citus 20~000 poļu politiskos ieslodzītos, kuri bija nometināti Sibīrijā. T.~Kostjuško emigrēja uz ASV.

Kaut cietusi sakāvi, 1794.--1795.~gada sacelšanās ievadīja virkni poļu nacionālās atbrīvošanās kustības dalībnieku sacelšanās vairāku gadu desmitu garumā, nostādot Polijas jautājumu Eiropas valstu politikas dienas kārtībā. Jau minētie padomju vēsturnieki V.~Djakovs un I.~Millers uzskatīja, ka šī sacelšanās iezīmē poļu atbrīvošanās kustības sākumrobežu.

Kā vēsturisku mītu var pieminēt baumas, kas klīda par Katrīnas~II nāvi, (pēc oficiālās versijas viņa mira no asinsizplūduma smadzenēs), kas it kā bija saistīta ar pēc T.~Kostjuško sakāves un trešās Polijas dalīšanas uz Pēterburgu atvesto Pjastu (\pltxti{Piasty}~--- poļu karaļu dinastija, kura valdīja no X līdz XIV gadsimtam) dinastijas troni. Katrīna~II esot personīgi likusi savos apartamentos Ziemas pilī to apvienot ar tādu toreizējo jaunievedumu kā savu personīgo ūdens klozetu. (Poļu garīdznieks un vēsturnieks, 1830--1831~gada sacelšanās dalībnieks V.~Kaļinka raksta par S.~Poņatovska troni, uz kura sēdējuši arī S.~Batorijs, Sigizmungs Vaza un viņa dēls Vladislavs. 1796.~gada 6.(17.) novembra rītā ķeizariene pēc pamošanās devās to apmeklēt, bet pēc kāda brīža galminieki izdzirdēja krītoša ķermeņa troksni. Kādu brīdi sulaiņi šaubījās līdz tomēr iedrošinājās atvērt durvis. Ķeizariene gulēja uz grīdas bez samaņas un noasiņoja. Tūlīt izsauktie ārsti neko vairs nespēja līdzēt, pēc dažām stundām Katrīna~II nomira no vaginālās asiņošanas. Pēterburgas aristokrātiskajos salonos vēl ilgi mēļoja, ka ķeizarienes ūdens klozetā zem Pjastu troņa esot bijis noslēpies kāds poļu fanātiķis, iespējams, punduris, un no apakšas viņu ievainojis ar šķēpu vai zobenu, bet pēc tam, izmantojot apjukumu, aizbēdzis no Ziemas pils. Loti jāšaubās, vai šai mītā ir kaut grans patiesības. Iespējams, baumu izcelsmi veicināja poļu svētuma Pjastu troņa izvešana no Polijas un it kā notikusī zaimošana, kas poļu acīs brēca pēc soda.

T.~Kostjuško vadītās sacelšanās, kura ilga 238~dienas, sakāve kalpoja par pamatojumu \strong{Polijas-Lietuvas valstiskuma likvidācijai}. 1795.~gada 24.~oktobrī valstis, kuras piedalījās Polijas sadalē, noteica savas jaunās robežas. \strong{Tā bija Polijas trešā dalīšana}. Karalis Staņislavs~II Augusts Poņatovskis 1795.~gada 25.novembrī nolika savas pilnvaras. Bijušās Žečpospolitas zemju valstiski-tiesiskā situācija būtiski mainījās: nodibinājās absolūtismam raksturīga pārvalde, poļu šļahta zaudēja politisko varu. Ja līdz Žežpospolitas sadalei jau sāka veidoties kopējs Vispolijas tirgus, tās sadale šo procesu pārtrauca, sarāva daudzus tradicionāli izveidojušos ekonomiskos sakarus starp atsevišķiem Polijas apgabaliem. Nepieciešamība veidot jaunus sakarus un pielāgoties Poliju sadalījušo valstu ekonomiskajām struktūrām nostādīja poļu apgabalus nevienlīdzīgā situācijā ar šo valstu teritorijām. Poļu apgabalu ekonomika nonāca atkarībā no minēto valstu politikas, to valdību centieniem ievērot vai neievērot (mazievērot) jauno pakļauto apgabalu intereses. Taču magnātu pozīcijas tika maz skartas. Jaunajos apstākļos, lielvaru aizsardzībā viņiem vairs nebija vajadzīgs sīkās šļahtas atbalsts. Pēdējā zaudēja iespēju saņemt no saviem labvēļiem līdzekļus par tiem sniegtajiem pakalpojumiem.

Austrijas varā pārgāja Galīcijas ziemeļu daļa līdz Bugas upei, ko sāka saukt par Rietumgalīciju, (atšķirībā no tās jau agrāk iegūtās Austrumgalīcijas ar Krakovu), daļa Mazovijas (\pltxti{Mazowsze}) un daži citi novadi ar 147 tūkstošu km$^{2}$ kopējo teritoriju un 1,2~miljoniem iedzīvotāju.

Prūsija ieguva zemes uz rietumiem no Piļicas (\pltxti{Pilica}), Vislas (\pltxti{Wisła}), Bugas (\pltxti{Boh}) un Nemunas (poļu \pltxti{Niemen}, lietuviešu \lttxti{Nemunas}, vācu \detxti{Memel}) upēm ar Varšavu, no kuras tās dumpīguma dēļ atteicās citas Polijas dalītājvalstis, kā arī Rietumlietuvas (Žemaitijas, poļu \pltxti{Żmudź}) zemes, kas kopumā sastādīja 55 tūkstošus km$^{2}$ ar 1~miljonu iedzīvotāju. Krievu vēsturnieks A.~Pogodins uzsvēra, ka ``pašas vecākās, īsteni poļu zemes'' 3.~dalīšanā ieguva Prūsija.

Tātad, kad runā par XVIII gadsimtā notikušajām trijām Polijās dalīšanām, faktiski tiek runāts par Žečpospolitas~--- apvienotās Polijas un Lietuvas~--- dalīšanām.

Etnisko poļu apdzīvotās teritorijas faktiski savā starpā sadalīja Austrija ar Prūsiju.

Krievija saņēma lietuviešu (daļa lietuviešu apdzīvoto teritoriju ar Suvalkiem gan ieguva Prūsija), baltkrievu un ukraiņu zemes uz austrumiem no Bugas upes ar kopējo platību 120 tūkstošu km$^{2}$ un 1,2 miljoniem iedzīvotāju.

Tā Žečpospolitas triju dalīšanu rezultātā Krievijas rokās nonāca latviešu (Kurzeme), lietuviešu, baltkrievu (izņemot teritorijas daļu ar Belostoku (\pltxti{Białystok}), kas pārgāja Prūsijas īpašumā) un ukraiņu zemes (izņemot Austrumgalīcijas apgabalu ar Lembergu (Ļvovu), kurš palika Austrijas rokās).

Tiesa, Krievijai bija jāsamaksā smaga cena~--- jāļauj nostiprināties Prūsijai un Austrijai. Jautājums nebija par Polijas likteni. Kāpēc Katrīnai~II vajadzēja ņemt vērā Polijas intereses, ja pēdējā nevēlējās ņemt vērā Krievijas un Polijā dzīvojošo krievu, baltkrievu un ukraiņu intereses? Krievijai svarīgāks bija kas cits~--- pazuda buferis starp Krieviju un vācu valstīm. Austrija un Prūsija tagad atradās tieši pie Krievijas robežām. Bet alternatīva būtu tikai atteikšanās pievienot Krievijai radniecīgos baltkrievus un ukraiņus. Karot pret Prūsiju un Austriju, ko, iespējams, atbalstītu Anglija, lai saglabātu Poliju tās etnogrāfiskajās robežās, nu nekādi neatbilda Krievijas interesēm. Var teikt, ka XVIII gadsimtā Prūsijas un Austrijas agresīvā politika pavēra iespēju Krievijai atrisināt Rietumkrievijas jautājumu bez asiņaina kara ar Eiropas lielvalstīm.

Krievu sabiedriskā doma gan toreiz, gan vēlāk apsveica ieguvumus. Kā savās 1899.~gadā izdotajās lekcijās par Krievijas vēsturi rakstīja krievu vēsturnieks S.~Platonovs, ``attiecībā pret Poliju Krievijas uzdevums bija atbrīvot pareizticīgos krievu iedzīvotājus no katolisko poļu valdīšanas, t.i.~--- atņemt Polijai vecās krievu zemes un no šīs puses sasniegt krievu tautības etnogrāfiskās robežas.'' Kā redzam, ``krievu tautībai'' šeit tika pieskaitītas arī citas (baltkrievu un mazkrievu jeb ukraiņu) tautības.

Līdzīgs viedoklis bija arī ievērojamajam angļu vēsturniekam, kulturologam un sociologam A.~Toinbi. Viņš 1947.~gadā rakstīja: ``Rietumos valda uzskats, ka Krievija ir agresors, \citespace{} XVIII gadsimtā Polijas dalīšanas laikā Krievija sagrāba teritorijas lielāko daļu: XIX gadsimtā tā ir Polijas apspiedējs \citespace{} Novērotājs no malas, ja tāds eksistētu, teiktu, ka krievu uzvaras pār zviedriem un poļiem ir tikai pretuzbrukums, \citespace{} XIV gadsimtā labākā daļa īsteno Krievijas teritoriju~--- gandrīz visa Baltkrievija un Ukraina~--- tika atrauta un pievienota rietumu kristietībai \citespace{} Poļu iekarotās īstenās krievu teritorijas \citespace{} tika atgrieztas Krievijai tikai 1930.--1945.~gada pasaules kara pēdējā fāzē.'' Citējot šo vērtējumu kā centienu panākt objektivitāti apliecinājumu, krievu literatūrzinātnieks un publicists V.~Kožinovs norādījis, ka Rietumu uzbrukums Krievijai sākās vēl agrāk, jau ar Polijas karaļa Boļeslava Drosmīgā iebrukumu Kijevā 1018.~gadā (par to runāts šī darba ievadā). V.~Kožinovs uzsvēris, ka nevar runāt par Krievijas dalību Polijas [kā poļu apdzīvotas zemes~--- V.Š.] sadalē. Patiesībā poļu zemes dalīja Austrija un Prūsija, Krievija sev pievienoja tikai senās austrumslāvu [neredzot lietuviešus un latviešus~--- V.Š.] apdzīvotās zemes, kuras arī mūsdienās ietilpst Baltkrievijas un Ukrainas sastāvā. Tāpēc V.~Kožinovs Krievijas dalību Polijas sadalīšanā nosaucis par ``liberālu mītu.''

Arī mūsdienu krievu vēsturnieks O.~Ņemenskis 2012.~gada 21.~decembrī konferencē ''Baltkrievijas un Krievijas apvienošanās'', veltītā Žečpospolitas pirmajai sadalīšanai, uzsvēra: ''Krievija visos trijos Žečpospolitas dalīšanas posmos nesaņēma ne pēdas poļu zemes, nepārkāpa Polijas etnogrāfisko robežu. Krievijas līdzdalības ideoloģija [Žečpospolitas] dalīšanās sastāvēja tieši no iepriekš vienotās Krievzemes atkalapvienošanās''. Tādejādi pēc vēsturnieka domām apvainojumi Krievijai par dalību poļu zemju dalīšanā XVIII gadsimtā pēc būtības neatbilst patiesībai.

Taču ievērojamais krievu vēsturnieks S.~Kļučevskis savukārt atzīmējis, ka Krievija, tā vietā lai pievienotu sev Rietumkrieviju, piedalījās Polijas sadalīšanā, tādejādi ``atrisinājums neatbilda uzdevumam''. Krievija sev pievienoja ne tikai Rietumkrieviju, bet arī Lietuvu un Kurzemi, taču daļu Rietumkrievijas~--- Galīciju atdeva Austrijai. Pēc S.~Kļučevska domām Polija nebija lieks loceklis Ziemeļaustrumu Eiropas valstu ģimenē, kalpojot par vāju starpnieku starp trijiem stipriem kaimiņiem. ``Atbrīvojusies no to vājinošās Rietumkrievijas un pārveidojusi savu valsts iekārtu, kā centās tās labākie sadales laikmeta ļaudis, tā varētu sniegt lielisku pakalpojumu slāvu lietai un starptautiskajam līdzsvaram, kļūstot par balstu pret ar visiem spēkiem uz austrumiem tiecošos Prūsiju. Pēc Polijas krišanas sadursmes starp nosauktajām trijām valstīm nevājināja vairs nekāds starptautisks buferis un tām bija vissāpīgāk jāatsaucas uz Krieviju, kuras robeža pa Nemunu nekļuva drošāka tāpēc, ka tai kaimiņos atradās prūšu priekšposteņi \citespace{} Bez krievu apgabaliem, savās nacionālajās robežās, pat ar izlabotu valsts iekārtu patstāvīga Polija būtu mums nesalīdzināmi mazāk bīstama, nekā tā pati Polija austriešu un vācu provinču veidā. Visbeidzot, poļu valsts iznīcināšana neatsvabināja mūs no cīņas pret poļu tautu: nepagāja ne 70~gadu kopš trešās Polijas dalīšanas, bet Krievija jau trīs reizes karoja ar poļiem (1812., 1831. un 1863.~gg.). Žečpospolitas rēgs, ceļoties no vēsturiskā kapa, radīja dzīva tautas spēka iespaidu. Varbūt, lai izvairītos no kara pret tautu, vajadzēja saglabāt tās valsti.'' S.~Kļučevska secinājumos faktiski visdziļāk izvērtēti Krievijas politikas rezultāti attiecībā pret Poliju.

Situācija, kādā nonāca poļi, tai laikā nebūt nebija ārkārtēja. Daudzas pasaules tautas gan pirms, gan arī pēc tam nonāca daudznacionālu impēriju sastāvā. Šoreiz šis liktenis bija piemeklējis poļus. Jāpiemetina, ka ar Žečpospolitas sadalīšanu sašķelta tika ne tikai poļu tauta. Arī citu tajā dzīvojošo tautu pārstāvji (tai skaitā ap 800~000 ebreju, kam nebija savas kompaktas teritorijas), nonāca dalītājvalstu varā, kur viņu statuss un dzīves apstākļi arī bija dažādi.

Pretsparu veselas Eiropas tautas nomākšanai, tās teritorijas sadalīšanai trijās valstīs pārējā Eiropa nevēlējās dot. Laikabiedru vairākums attaisnoja notikušo ``Eiropas kartes racionalizēšanu''.

Daudzi citu tautu vēsturnieki ir pauduši viedokli, ka poļi paši bija vainīgi savā nelaimē.

Viens no marksisma pamatlicējiem F.~Engelss tieši vainoja poļu feodālo aristokrātiju~--- magnātus savienībā ar trim lielvalstīm, kuras dalīja Žečpospolitu, lai tikai izbēgtu no revolūcijas, kā rezultātā tika sadalīta valsts, kas bija tikpat liela cik Francija. Domājams, ka bailes no revolūcijas, ko it kā izjuta magnāti, F.~Engelss pārspīlēja, jo viņi taču joprojām turpināja savstarpējās ķildas. Visi poļu magnāti vēlējās panākt stipras unitāras valsts eksistenci, taču tikai tai gadījumā, ja viņi paši tajā būtu noteicēji. Tādā situācijā visas reformas bija nolemtas neveiksmei.

Arī krievu vēsturnieki parasti uzsvēra Polijas sabrukuma iekšējos cēloņus, bet apskatot ārējo faktoru, galveno atbildību par Žečpospolitas sadali piedēvēja Prūsijai, turpretī Krievijas dalību tajās dažkārt pat raksturoja kā neatbilstošu tās nacionālajām interesēm. Tā, jau minētais krievu vēsturnieks A.~Pogodins rakstīja, ka Polijas valsts gāja bojā, jo tā nepaguva laikus nostāties uz tai laikā vienīgi pareizā monarhiskās varas un reizē ar to militārisma, rūpniecības un tirdzniecības attīstības ceļa.

Kopumā sākotnēji arī poļu vēsturnieku vidū lielākā uzmanība tika pievērsta valsts neatkarības zaudēšanas iekšējiem cēloņiem. Piemēram, autoritatīvais poļu vēsturnieks un valstsvīrs (viņš 1908.--1913.~gadā bija Galīcijas un Lodomērijas vietvaldis (\detxti{Statthalter})) M.~Bobržinskis, atzīmējot, ka nav bijis valsts Eiropā, pret kuru jaunajos laikos tās kaimiņi nebūtu turējuši ļaunus nodomus, uzsvēra, ka pat salīdzinoši mazas nācijas spēja no šīs cīņas iziet kā uzvarētājas, turpretī ``poļi, kad viņus dalīja kaimiņi, nebija vājāki par katru no tiem ne pēc zemes platības, ne iedzīvotāju skaita, ne materiālās labklājības, ne garīgās attīstības, pat bija pārāki pār tiem tai vai citā jomā, taču tikai vieni poļi krita, pie tam bez cīņas, bez īstas cīņas, uz kādu tie bija spējīgi''. Atsevišķu personību gatavība nest upurus valsts labā ``nevarēja aizstāt gatavības trūkumu sabiedrībā nest upurus''. Vēsturnieks secināja ``Ne [neizdevīgās] robežas un ne kaimiņi, bet vienīgi iekšējās nesaskaņas noveda poļus pie politiskās pastāvēšanas zaudēšanas''.

Krievu vēsturnieks un filozofs N.~Karejevs 19.~gadsimta 80.~gadu beigās secināja, ka visi pētnieki ir vienoti~--- Polija gāja bojā iekšējo cēloņu rezultātā. ``Diagnoze~--- likumdevējas varas bezspēcība un pilnīgs izpildvaras sajukums.''

Tomēr poļu vēsturnieku viedokļi XIX gadsimta beigās~--- XX gadsimta sākumā sadalījās divās skolās. T.s. Krakovas skolas pārstāvji uzskatīja, ka XVIII gadsimta beigās notikušās Žečpospolitas sagrūšanas cēloņi meklējami valsts iekšējā vājumā, otras~--- Varšavas skolas uzskatu aizstāvji tos saskatīja tai apstāklī, ka pēc 1772.~gadā notikušās dalīšanas valsts nonāca intensīvas modernizācijas fāzē, bet krita par upuri savu kaimiņvalstu~--- Krievijas, Prūsijas un Austrijas negausībai. Arī ievērojamais vēsturnieks no Lembergas O.~Balcers, kurš kā pirmais no poļu zinātniekiem 1921.~gadā tika apbalvots ar Baltā ērgļa (\pltxti{Orła Białego}) ordeni, apgalvoja, ka`` \dots{} īstais, izšķirošais mūsu valstiskuma pagrimuma cēlonis, īsts šī notikuma \latxti{causa efficiens} (izraisītājs) bija apvienoto, tātad vareno, kaimiņu alkatība, kuri noslēdza savienību, lai pazudinātu Poliju.''

Neviens no virzieniem tā arī neieguva izšķirošu pārsvaru. Vairākums nonāca pie secinājuma, ka Polijas dalīšanas izsauca kā iekšējo, tā ārējo faktoru mijiedarbība. Ir izteikts diezgan pārliecinošs viedoklis, ka atsevišķi par sevi ne vieni, ne otri nenovestu pie valstiskuma zaudēšanas.

Jau pieminētais vēsturnieks no Poznaņas J.~Rutkovskis rakstīja: ``Polijas iekšējā iekārta bija pietiekami stipra, lai tā varētu realizēt savu pašas iekšējo politiku, taču bija par vāju, lai varētu aizstāvēties pret ārējo spiedienu.'' Viņš arī uzsvēra Žečpospolitas militāro vājumu, kuru lielā mērā noteica tās politiskais vājums, kad Seims nevēlējās stiprināt armiju, baidoties, ka karalis to izmantos absolūtās monarhijas radīšanai. Poļu vēsturnieks J.~Tazbirs, kurš uzsver, ka Poliju XVIII gadsimtā līdz bojā ejai noveda trūkumi valsts uzbūvē, nevis nacionāli netikumi, ir norādījis, ka poļu tauta pēc pirmās Žečpospolitas sadales ar reformām, 1791.~gada 3.~maija Konstitūcijas pieņemšanu, T.~Kostjuško vadīto sacelšanos pierādīja, ka tā ir pati spējīga labot savu valsti, cīnīties par to, līdz ar to apstrīdot tēzi, ka ``mēs kritām mūsu trūkumu rezultātā.'' Jāpiekrīt, ka norisa cīņa par pretrunu, valsts trūkumu pārvarēšanu, bet reizē jāatzīmē, ka notikušās Žečpospolitas dalīšanas taču tieši rādīja, ka panākumi šai cīņā bija nepietiekami lai saglabātu valsti. Domājams, ka pārmetumus nav pelnījusi tauta, bet gan galvenokārt tās augšslāņi. Nosacījumus Polijas sadalīšanai radīja tās bezspēcība, tās valdošo slāņu nespēja kopot tautas spēkus vispārnacionālam uzdevumam~--- suverenitātes un teritoriālās integritātes aizstāvēšanai, tā vietā nodarbojoties ar cīņu par savu privilēģiju sargāšanu un paplašināšanu. Vienkāršie zemnieki neko neieguva nedz no karaļu nomaiņām, nedz varas pārdalēm, nedz arī Žečpospolitas sadalīšanas.

Mūsdienu vēsturnieku darbos autoram nav gadījies atrast Polijas valstiskuma iznīcināšanu attaisnojošus uzskatus, bet domstarpības pastāv par tās novērtējumu XVIII gadsimta notikumu kontekstā, par citu valstu tiesībām uz tām vai citām teritorijām.

Prūsija no bijušajām poļu zemēm izveidoja trīs provinces: Rietumprūsiju, Dienvidprūsiju un Jauno Austrumprūsiju. Prūsijas poļu valdījumos sākās poļu asimilācija. Daudzas valsts un baznīcas muižas nonāca prūšu muižnieku rokās. Prūsijai piederošajā Polijas daļā jau no 1776.~gada darbojās likums, kurš šļahtiču zemes atļāva iepirkt ne tikai vācu muižniekiem, bet arī pilsētniekiem. Ar dažādām privilēģijām poļu zemēm tika piesaistīti vācu zemnieki, amatnieki un tirgotāji. Poļu ciemu vidū veidojās vācu zemnieku ciemi. No 1797.~gada visa tiesvedība un administratīvā lietvedība notika vācu valodā, poļi tika izraidīti no valsts dienesta, nodibināja vācu skolas. Poļiem tika uzspiesta visas sabiedriskās dzīves reglamentācija. Poļu vēsturnieks J.~Feldmans uzsvēra, ka Prūsija uz Polijas rēķina palielinājās vairāk nekā divas reizes un kļuva par ``modru Polijas jautājuma kapraci'' Eiropas politikā, par sīkstāko poliskuma apkarotāju agrākās Žečpospolitas teritorijā.

Zemes, kuras nonāca Austrijā, ieguva nosaukumu Galīcija un Lodomērija (vācu \detxti{Königreich Galizien und Lodomerien}, poļu \pltxti{Królestwo Galicji i Lodomerii}, kur \latxti{Lodomeria} bija latīņu valodā lietots Galīcijas un Vladimiras kņazistes nosaukums XIII--XIV gs.) un tika sadalītas 12 apvidos. 1806.~gadā Austrija no Galīcijas guva 19\% visu valsts ienākumu. Arī Galīcijā sākās poļu asimilācijas mēģinājumi, gan ne tik enerģiski kā Prūsijas zemēs. 1790.~gadā Vīnes augstmaņu slepenā memorandā sakarā ar visai biklu Konstitūcijas projektu, kuru ierosināja provinces seims un kuru noraidīja imperators Leopolds~II, bez aplinkiem bija norādīts, ka Austrijas mērķim jābūt ``galīciešu pakāpeniskai pārvēršanai par vāciešiem''.

Krievijas rokās nonākušās teritorijas tika sadalītas Kurzemes, Viļņas un Grodņas guberņās. Saglabājās kārtu pašpārvalde, poļu muižnieku kundzība pār zemniekiem nebija satricināta. Uz laiku saglabājās vietējie likumi, arī Lietuvas statuta (\lttxti{Lietuvos Statutas}~--- likumu kodekss bija radīts XVI~gs. un civillietās palika spēkā daļā no bijušās Lietuvas lielkņazistes teritorijas līdz 1840.~gadam) nozīme. Krievijas dalība Polijas sadalē uz diviem turpmākajiem gadsimtiem lielā mērā noteica Krievijas ārējo politiku, jo starptautiskā stabilitāte bija lielā mērā atkarīga no valstu~--- triju dalīšanas dalībnieču attiecībām.

Dažādi Polijas apgabali vēl līdz tās dalīšanai attīstījās nevienmērīgi, tas bija redzams arī tikai etnisko poļu apdzīvotajās zemēs. Taču šis nevienmērīgums, kā uzskatīja poļu ekonomikas vēsturnieks V.~Kula, arī stiprināja valsts vienotību, radot priekšnoteikumus darba dalīšanai starp atsevišķām zemēm. Pēc Polijas sadales ekonomiskās attīstības nevienmērīgums pieauga. Tā, t.s. Lielpolija~--- ekonomiski attīstītākais Žečpospolitas apgabals XVIII gadsimta beigās, kurš acīmredzot saglabātu šo lomu neatkarīgā valstī, pēc tās sadales kļuva par Vācijas impērijas ``aizmuguri'' (vācu \detxti{Hinterland}~--- zemi, kuru raksturoja mazāks iedzīvotāju blīvums un sliktāka ekonomiskā un infrastruktūras attīstība kā blakus esošajos attīstītākajos apgabalos), bet Austrumgalīcija kļuva par ekonomiski atpalikušāko no visām poļu provincēm. Ar sadali tika sarauti vai vismaz apgrūtināti veidojamie ekonomiskie sakari starp dažādām poļu teritorijām. Tiesa, vēlreiz jāatgādina, ka ne tikai pirms 1795., bet arī pirms 1772.~gada ne visas poļu apdzīvotās teritorijas atradās vienas valsts sastāvā. Ārpus Žečpospolitas palika tādas zemes, uz kurām tā pretendēja, kā jau pieminētās Silēzijas (poļu \pltxti{Śląsk}, vācu \detxti{Schlesien}) un Rietumu Pomorjes jeb Pomerānijas (poļu \pltxti{Pòmòrzé}, vācu \detxti{Pommern}) teritorijas. Reizē jāuzsver, ka dalīšanas, sagraujot vienus sakarus, veicināja citus. Piemēram, sakarus starp Pozenes un Silēzijas apgabaliem. Tie pat stiprināja poļu elementus Silēzijā, kas guva izpausmi jau XIX gadsimta pirmajā pusē.

Par vienu no poļu zemju attīstības nevienmērības rādītājiem var kalpot iedzīvotāju pieauguma dinamika. Tā, 1816.--1856.~gadā iedzīvotāju skaits gandrīz dubultojās Prūsijas rokās esošajā Pomerānijā un Augšsilēzijā (vācu \detxti{Oberschlesien}, poļu \pltxti{Górny Śląsk}), par 73\% pieauga Pozenes hercigistē, par 36\% Polijas karalistē un tikai nedaudz vairāk par 20\% Austrijai piederošajā Galīcijā. V.~Kula uzsvēris, ka Polijas sadale paildzināja feodālisma pastāvēšanu tās zemēs. Vājā poļu feodālā valsts, kāda bija Žečpospolita savas pastāvēšanas pēdējos gados, nespējusi feodāļiem garantēt viņu privilēģijas revolucionāras situācijas priekšā, tāpēc tie bija ieinteresēti iegūt absolūtisko kaimiņvalstu aizbildniecību, kuras tad arī nodrošināja to privilēģiju pastāvēšanu vēl vairāk nekā pusgadsimta garumā. Taču šis spriedums nevar būt viennozīmīgs. Pats V.~Kula bija spiests atzīt, ka Žečpospolitas sadale neapturēja kapitālisma elementu nobriešanu feodālisma iekšienē un ja arī to aizturēja, tad tikai uz īsu laiku, ka kapitālisma attīstība daudzās nozarēs un teritorijās pēc Žečpospolitas sadales pat paātrinājās. Viņš arī norādījis, ka reģionālās ekonomikas iezīmes lielā mērā veidojās daudzu sociāli-ekonomisku faktoru ietekmē, starp kuriem politiskajām robežām nebija galvenā loma. Tā, vēlāk Krievijas atkarībā pastāvošajā Polijas karalistē (1815--1915) esošais Kališas (poļu \pltxti{Kalisz}, vācu \detxti{Kalisch}) rajons ekonomiskā ziņā maz atšķīrās no Pozenes provinces. Dombrovas baseins karalistē faktiski veidoja vienu ekonomisku rajonu ar Augšsilēziju Prūsijā, bet Austrumgalīcijai ar Lembergu bijis mazāk kopīgā ar Rietumgalīciju ap Krakovu, kas abas atradās Austrijas sastāvā, nekā starp Rietumgalīciju un tai Polijas karalistes pusē pieguļošo Mehovas (poļu \pltxti{Miechów}) rajonu. Ar kapitālisma attīstību pieauga reģionālās atšķirības, kas nesakrita ar politiskajām robežām. V.~Kula arī norādījis, ka nevis Žečpospolitas sadale, bet kapitālisma attīstība tās sadalītajās daļās noveda pie Polijas ``A'' (ekonomiski attīstītāko rajonu) un Polijas ``B'' (ekonomiski vājāko, atpaliekošo rajonu'') izveides. Pat vienā pašā Polijas karalistē bija novērojams šāds dalījums. Karalistē abos Vislas krastos pastāvēja dažādas Polijas, kur atšķirības bija lielākas nekā starp Pozenes provinci Prūsijā un Kališas rajonu Polijas karalistē.

\asterism


Taču \strong{Žečpospolitas valstiskuma likvidācija} notika apstākļos, kad poļu nācijas konsolidācijas process jau bija izvērsies, tāpēc tā sastapa stipru tautas pretestību, \strong{izsauca poļu nacionālo kustību}. Poļu emigrācija atgādināja pasaulei par Polijas pastāvēšanu. No šī laika starptautiskajās attiecībās sāka pastāvēt ``poļu jautājums'', kuru ar Polijas dalītājām konkurējošās valstis centās izmantot savās interesēs.

Kaut Polijas valsts pārstāja pastāvēt, daudzi poļi, īpaši šļahtiči, neatmeta cerības atjaunot neatkarību, kaut dažādie valsts atjaunošanas plāni visbiežāk izrādījās nereāli. Tieši \strong{šļahtiči nostājās nacionālās atbrīvošanās kustības priekšgalā}. Polijas īpatnība bija šļahtiču lielais īpatsvars~--- ap 6 līdz pat 10\% no visiem iedzīvotājiem. Tiesa, par šļahtičiem Žežpospolitā par īpašiem nopelniem varēja kļūt arī zemāko kārtu pārstāvji, taču tas notika reti, jo jautājumu izlēma Seims. Tā, ir dati, ka 1788.--1792.~gadā, kad Polijā darbojās jau minētais Četrgadu Seims (\pltxti{Sejm Czteroletni}), par šļahtičiem kļuva ap 400 cilvēku, tikpat cik visos iepriekšējos XVIII gadsimta gados.

Tomēr bija arī ceļi, kā apiet likumu. Viens no līdzekļiem, kā pierādīt savu piederību šļahtiči kārtai, bija 12 liecinieku uzrādīšana. Bagātam plebejam vispirms bija jāaiziet pie bārddziņa un jālūdz iegriezt pāris brūču, kuras nebūtu bīstamas, bet radītu skaidri saskatāmas rētas~--- jo, kas gan par šļahtiču bez kaujās gūtām rētām! Pēc tam bija jānoalgo kāds, kas apvainotu šļahtiča statusa pretendentu, ka viņš nav šļahtičs, un tam pretī jānostāda 12 liecinieki, kuri šo apgalvojumu noliegtu. Tā ar tiesas spriedumu varēja iekļūt priviliģēto kārtā. Bez tam Lietuvas lielkņazistē līdz 1764.~gadam darbojās likums, pēc kura katrs ebrejs, kurš pārgāja kristīgajā ticībā, uzreiz kļuva par šļahtiču. Tas bija apbalvojums par ticības maiņu. Bija arī gadījumi, kad plebeji pārgāja judaismā, lai pēc dažiem gadiem atgrieztos katolicismā un tā iegūtu šļahtiča ģerboni. Tāpēc arī minētais likums galu galā tika atcelts.

Kopējā šļahtiču masā zemes īpašnieki sastādīja mazākumu. Pēc Žečpospolitas pirmās sadales tādu muižnieku bija vairāk nekā 300~000, bet t.s. sīko šļahtiču, kuru īpašumā bija tikai daļa ciema vai vispār nebija sava īpašuma~--- vairāk nekā 400~000. Pēc poļu vēsturnieka T.~Ļepkovska aprēķiniem XVIII~gadsimta beigās bezzemes šļahta sastādīja 55\%, bet 40\% šļahtiču īpašumā nebija dzimtcilvēku. Lielie zemes īpašnieki~--- magnāti, kuru īpašumā bija vairāki simti ciemu, dzīvoja galvenokārt valsts austrumos, tai skaitā zemēs, kur poļi bija etnisks mazākums. Viduspolijā par magnātu jau skaitījās šļahtičs, kura īpašumā bija vairāki desmiti ciemu. Sīko šļahtiču īpašumā parasti bija muiža ar dažiem klaušu zemniekiem. Viņi parasti algoja strādniekus, jo savu klaušinieku nepietika. Šādu šļahtiču dēli bieži nodarbojās ar tirdzniecību (1775.~gadā tika atcelts aizliegums šļahtičiem nodarboties ar amatniecību un tirdzniecību), kļuva par garīdzniekiem, dienēja pie magnātiem, pārejot bezīpašuma šļahtiču statusā. Daļa šļahtiču magnātu dienestā līdz vecumam nopelnīja pensiju, bet vairākums cieta trūkumu, atšķiroties no zemniekiem tikai ar savām iluzorajām pilsoniskajām un politiskajām tiesībām. Nabadzīgākie bezīpašuma šļahtiči rentēja zemes gabalus, tai skaitā arī no zemniekiem, kurus arī paši apstrādāja. Īpaši Mazovijas (poļu \pltxti{Mazowsze}, vēsturisks apgabals Polijas centrā) un Podlases (poļu \pltxti{Роdlаsiе}~--- no poļu ``\pltxti{Pod lasem}'' (zem meža)~--- vēsturisks apgabals tagadējās Polijas austrumos) novados varēja redzēt cilvēku ar zobenu pie sāniem, kurš ara zemi. Taču šos arājus-šļahtičus nekad neatstāja pārākuma apziņa pār saviem kaimiņiem dzimtcilvēkiem.

Jau XVIII gadsimta otrajā pusē ar magnātu karadraudžu likvidāciju magnātu aizbildniecībā esošā t.s. šļahtiču klientūra strauji mazinājās, tā vairs magnātiem nebija vajadzīga. Ar neatkarības zaudēšanu un ar kapitālistisko attiecību attīstību sīkā šļahta izrādījās vairs nevajadzīga savā agrākajā~--- karotāju veidolā. Pāreja uz činšu (renti) padarīja nevajadzīgus šļahtiču amatus magnātu saimniecībās. Arvien vairāk šļahtiču pārcēlās uz pilsētām. Ja turīgākie pilsētu šļahtiči papildināja augstāko ierēdņu, advokātu un brīvo profesiju pārstāvju rindas, tad nabadzīgākie līga darbā par apsargiem, mājkalpotājiem, reizēm kļūstot arī par klaidoņiem.

Angļu vēsturnieks N.~Deiviss pat uzskata, ka juridiskā ziņā Polijas aristokrātijai pienāca gals, kad pēc Žečpospolitas sadalīšanas tika anulēti likumi, kas noteica tās statusu. Taču N.~Deiviss nav precīzs. Protams, magnātu un šļahtas ietekme citām valstīm pievienotajās poļu teritorijās mazinājās, taču pilnībā neizzuda.

Daudzas magnātu muižas tika pārdotas, lai nomaksātu parādus. To valstu valdības, kuras sadalīja Poliju, pārņēma savā īpašumā vairākumu valsts un baznīcas zemju (tā atņemot iztiku daudziem šļahtičiem~--- karaļa zemju nomniekiem), bet par dalību 1794.~gada sacelšanās konfiscētās zemes sadalīja savas valsts augstmaņiem. Administratīvā pārvalde un tiesa nonāca nepoļu ierēdņu rokās. Prūsijā šļahta zaudēja personas neaizskaramības tiesības. Austrija arī atcēla šļahtiču personīgās brīvības un to īpašuma neaizskaramības tiesības, privilēģiju būt atbrīvotiem no nodokļiem. Kā raksta jau minētais poļu vēsturnieks J.~Tazbirs, XVIII gadsimta pirmajā pusē valdošais valsts aparāta vājums, nespēja iekasēt nodokļus nāca par labu šļahtai, kura iedzīvojās uz valsts rēķina. Kad Polija tika sadalīta, šļahtiči uzreiz izjuta, ka okupanti nodokļus prot ievākt daudz labāk nekā Žečpospolita. Lielā mērā šis apstāklis bijis pamatā vienam no šļahtas kultivētajiem mītiem~--- ka Žečpospolitas valsts iekārta bijusi gandrīz vai ideāla. Īpašu sašutumu šļahtā viesa tas, ka poļu zemes sadalījušo valstu aparāts iejaucās muižnieku un zemnieku attiecībās.

Taču šo triju lielvalstu, kuras pašas bija feodālas kārtu monarhijas, politika nevarēja būt radikāla. Tās poļu zemēs saglabāja feodālās attiecības un kārtu struktūru. Daļai poļu magnātu savu dižciltību izdevās apstiprināt Austrijā un Prūsijā. Īpaši Austrijas imperators labprāt apmaiņā pret likvidējamajiem senatoru, vojevodu, kastelānu (\pltxti{Kasztelan}, no latīņu \latxti{castellum}~--- pils, XVIII gadsimtā tie dienesta hierarhijā ieņēma otro vietu pēc vojevodām) u.c. amatiem piešķīra poļu lielmuižniekiem savulaik Žečpospolitā mazizplatītos grāfu un baronu titulus. Vairākums poļu grāfu un baronu ģimeņu savus titulus ieguva Vīnē un Berlīnē pirmajā desmitgadē pēc Žečpospolitas krišanas. Poļu šļahta zaudēja daudzas savas tiesības, taču saglabājās kā kārta.

Pēc spāņu izcelsmes angļu vēsturnieces I.~De~Madariagas datiem Krievijas muižnieku (krievu \rutxti{дворяне}) kopskaits 1795.~gadā sastādīja ap 111~600. Pēc Polijas sadales tiem pievienojās 250~974 šļahtiču, kas Krievijas impērijā sastādīja 66,22\% no visiem muižniekiem. Tātad Krievijai piederošajā Polijas daļā šļahtiču bija vairāk nekā muižnieku visā pārējā Krievijas impērijā. Līdz 1830.~gadam Polijas karalistē saglabājās bagāto zemes īpašnieku pārsvars sabiedriskajā un ierēdņu hierarhijā. Tika apstiprināti prūšu un austriešu piešķirtie tituli, tika dāvāti arī jauni. Tā kā šļahtiču īpatsvars bija salīdzinoši liels, nav brīnums, ka viņu uzskati, tikumi lielā mērā ietekmēja veidojošos poļu nāciju. Šļahtas psiholoģiskais veidols un tās tradīcijas arī turpmāk atstāja lielu iespaidu poļu kultūras attīstībā, kad no to rindām nākusī inteliģences daļa saistīja sevi ar buržuāzijas ekonomiskajām un sabiedriskajām interesēm.

Nedaudz apsteidzot notikumus, jāsaka, ka 1817.~gadā Krievijai piederošajā Polijas daļā tika izdots likums, kurš, tāpat kā to paredzēja Krievijas impērijā spēkā esošā vēl 1722.~gadā Pētera I ieviestā Rangu tabula, deva iespējas par nopelniem valsts labā iegūt muižnieka (šļahtiča) tiesības. Katrs 10 gadus nokalpojis skolotājs, katrs karavīrs, kurš uzdienēja līdz kapteiņa pakāpei, katrs, kam bija ievērojami nopelni valsts priekšā, ieguva šļahtiča tiesības. Avīze ``\pltxti{Dziennik Praw}'' (``Likumu Dienasgrāmata'') publicēja šo cilvēku sarakstus. Tas veda pie šļahtas stāvokļa pakāpeniskas nonivelēšanas, ar ko ``īstie'' šļahtiči bija neapmierināti. Vēl krasāk stāvoklis mainījās pēc 1830.~gada sacelšanās. 1836.~gadā Polijas karalistē tika izdots dekrēts par muižniecību, kura mērķis bija radīt jaunu, saistītu ar carismu muižnieku kārtu. Mantojamās muižnieku (šļahtas) tiesības bija likumīgi jāapstiprina: dokumentāli jāpierāda, ka kāda no vīriešu kārtas senčiem īpašumā līdz 1775.~gadam bija vismaz viens ciems, vai līdz 1795.~gadam viņš bija Seima deputāts, senators, valsts darbinieks, poļu ordeņa kavalieris. Poļu šļahta saglabāja zināmas privilēģijas (kara dienestā, vidējās un augstākās izglītības ieguvē), taču zaudēja monopolu uz zemes īpašumu, nodokļu atlaidēm, atbrīvošanu no kara klausības. Līdz 1861.~gadam, kad tika beigts izskatīt jautājumu par šļahtiču tiesību apstiprināšanu, kārtas privilēģijas bija zaudējuši apmēram ¼ dzimušo šļahtiču. Polijas karalistē XIX gadsimta vidū tikai ap 5~000 šļahtiču ģimeņu piederēja muižas. Saprotams, ka jau tas vien radīja pamatu cietušo neapmierinātībai.

Tiesa, arī turpmāk Krievijas muižniecība bija daudznacionāla. Pēc 1897.~gada skaitīšanas datiem tikai 53\% dzimtmuižnieku par savu dzimto nosauca krievu valodu. 28,3\% no viņiem sevi uzskatīja par poļiem. Tātad vēl XIX gadsimta beigās vairāk nekā ceturtdaļa Krievijas impērijas muižnieku bija poļi. Tomēr poļu šļahtiču īpatsvars Krievijas muižnieku vidū bija ievērojami mazinājies. Daļa bagātās šļahtas~--- zemes īpašnieki, uzkrājot kapitālu, paplašinot preču ražošanu, izmantojot algotu darbaspēku, kļuva par kapitālistiskiem uzņēmējiem, saglabājot arī sociālo prestižu un sabiedrisko ietekmi. Bezīpašuma šļahta piegādāja zemākā un vidējā līmeņa ierēdņus, tās pārstāvji ieņēma daudz virsnieku posteņu armijā, taču tikai nelielai tās daļai bija vēlēšanu tiesības vietējās pašvaldības iestādēs. Sīko un vidējo muižu īpašnieki, bet vēl jo vairāk zemes īpašumus zaudējušie šļahtiči asi izjuta briesmas, kuras viņu privilēģijām nesa kapitālisma attīstība. Liela daļa no viņiem zaudēja savu mantu, to stāvoklis tuvinājās valsts zemnieku statusam. Daļa bijušo šļahtiču veidoja deklasētu ļaužu grupu. (Piemēram, bija gadījumi, ka nabadzībā nonācis poļu šļahtičs~--- dzimtmuižnieks 19.~gadsimta beigās Rīgā strādāja par amatnieku vai strādnieku.)

Šļahta ``bija galvenā pretkrieviskā noskaņojuma barotne visa XIX gadsimta laikā'', kā to raksturojis N.~Deiviss. Šļahtas centieni panākt tādu valsts iekārtu, kuras attīstību viņi varētu ietekmēt sev labvēlīgā virzienā, galvenokārt saistījās ar Polijas valsts neatkarības atjaunošanu, pie tam atjaunošanu agrākajās robežas, kad poļu feodāļi ekspluatēja arī citu tautu zemniekus. Šļahtiči un viņu sekotāji vēl XX gadsimtā par katru cenu centās atjaunot Poliju 1772.~gada robežās ``no jūras līdz jūrai'' (\pltxti{Od morza do morza}), tātad~--- vienā valstī apvienojot gan poļu, gan lietuviešu, gan ukraiņu un baltkrievu zemes. Sociālās priekšrocības, kuras šī kārta izmantoja gadsimtu gaitā (izglītība, prestižs sabiedrībā), ļāva tai ieņemt vadošo vietu poļu nacionālajā kustībā.

Tomēr jāuzsver, ka šļahtiču politiskie uzskati bija daudzveidīgi: dažas grupas ieņēma klerikālas, feodālas pozīcijas, citas, ``maksājot nodevas'' romantismam un lielvalsts atjaunošanas ideālam, pēc uzskatiem tuvinājās buržuāziskajiem revolucionāriem, izvirzīja arī progresīvas prasības par zemnieku stāvokļa uzlabošanu, to feodālās atkarības likvidēšanu, kārtu privilēģiju likvidēšanu, vēlēšanu tiesību paplašināšanu un citu buržuāzisko politisko brīvību ieviešanu.

Par poļu šļahtičiem dažādos laikmetos ir izteikti visdažādākie viedokļi, no cildinošiem līdz paļājošiem. Piemēram, kritiski pret tiem bija noskaņots krievu publicists un beletrists F.~Bulgarins, no mātes puses polis, viņa tēvs bija karojis poļu pusē pret Krieviju un tāpēc izsūtīts uz Sibīriju. Pats viņš bija dienējis gan Krievijas armijā, gan poļu leģionos Napoleona armijā, bet pēc tam bija ierēdnis un nodarbojās ar literāru darbību Pēterburgā. F.~Bulgarins rakstīja: ``Polijā no laika gala mēļoja par brīvību un vienlīdzību, kuru īstenībā nebija nevienam, tikai bagātie pani bija pilnīgi neatkarīgi no visām varām, taču tā nebija brīvība, bet patvaļa \citespace{} Sīkā šļahta, nevaldāma un neizglītota, vienmēr atradās pilnīgā atkarībā no katra, kurš to baroja un dzirdināja, un pat stājās viszemākajos amatos pie paniem un bagātās šļahtas, un pacietīgi pacieta kāvienus,~--- ar noteikumu, lai tas notiktu ne uz kailas zemes, bet uz paklāja, taču muļķīgas lepnības dēļ nicināja nodarbošanos ar tirdzniecību un amatiem kā neatbilstošu šļahtiča nosaukumam.''

Nepakļāvīgo šļahtas attieksmi pret iekarotājiem, to monarhiem lielā mērā balstīja tās senās tradīcijas. Pat zaudējuši tiesības uz šļahtiča nosaukumu, daudzi no tiem turpināja cienīt savu ģerboni. Ilgstoši saglabājās šļahtas prestižs sabiedrībā. Pat sīkie šļahtiči, kuriem vidējās izglītības un ierēdņa vietas ieguve pilsētā nozīmēja jau lielu personīgu panākumu, centās norobežoties no nedižciltīgajiem pilsētniekiem. ``Īsts'' šļahtičs karali uzskatīja par sev līdzvērtīgu, nevis augstākstāvošu personu, aizstāvēja savas \latxti{veto} tiesības, bija gatavs uz nepakļaušanos līdz pat atklātam dumpim (\pltxti{rokosz}) savu tiesību un brīvību aizsardzības vārdā. Politiska opozīcija esošajai varai tika uzskatīta par cienījamu tikumu, pilsonisko pienākumu. Faktiski princips, ka vara valdniekam nav Dieva dota, bet viņš to saņēmis no vēlētājiem, par kuru it kā iestājās šļahta, XIX gadsimta cīņās pret nacionālajiem apspiedējiem, neliecināja par šļahtas progresivitāti, bet par tradicionālismu, jo kā pilntiesīgi valdnieka vēlētāji tika atzīti tikai paši šļahtiči. Tas arī bija viens no galvenajiem šļahtas bezspēcības cēloņiem XIX gadsimta brīvības cīņās. Kā rakstīja ievērojamais krievu vēsturnieks S.~Solovjovs: ``Valstisko un sabiedrisko aizturu iztrūkums, sava spēka apzināšanās, savas pilntiesības un neatkarības vienreizīguma apzināšanās bija par pamatu galējai personības attīstībai poļu šļahtā, tieksmei pēc neierobežotas brīvības, nemācēšanai ar savu es piekāpties vispārēja labuma priekšā.''

Pēc Polijas sadales daļa muižnieku, labi saprotot, ka ārējā vara nodrošina viņiem pastāvošās sabiedriskās iekārtas un privilēģiju saglabāšanu, piekopa samierniecisku politiku. Taču daudzi magnāti un ievērojama daļa šļahtas piedalījās atbrīvošanās kustībā. Parasti viņus vadīja ne tikai patriotisma jūtas, bet arī personīgās intereses. Tomēr, runājot par atbrīvošanās kustību Polijā, V.~Ļeņins līdz XIX gadsimta 60.~gadiem to sauca par ``šļahtas atbrīvošanās kustību'', norādot, ka tā ``ieguva milzīgu, pirmšķirīgu nozīmi ne vien no visas Krievijas, ne vien no visu slāvu, bet arī no visas Eiropas demokrātijas viedokļa.''

Katra jauna šļahtiču paaudze (vismaz daļa tās) mēģināja atbrīvoties no apspiedējiem. Vieni cerēja kooperēties ar Krieviju vai Prūsiju, lai ar to palīdzību atjaunotu Polijas vienotību. Piemēram, poļu ģenerālis J.H.~Dombrovskis, kurš jau bija karojis T.~Kostjuško vadībā, 1796.~gadā piedāvāja Prūsijas karalim Fridriham Vilhelmam II poļu karavīrus nostādīt Prūsijas pavēlniecībā. Taču karalis par to neizrādīja nekādu interesi. Citi poļu šļahtiči vai nu pievienojās to valstu, kuras sadalīja Poliju, pretiniekiem, vai paši rīkoja sacelšanās. Pēc 1794.~gada daudzi poļi devās uz Itāliju un Franciju, viņu apmešanās centri bija Venēcija un Parīze. Īpašas cerības tika saistītas ar Franciju, kura veda nacionālus revolucionārus karus pret kontrrevolucionāro monarhiju koalīciju. Izveidojoties Napoleona I vadītajai Francijas impērijai, tā pakļāva virkni Eiropas nacionālo valstu, veda imperiālistiskus iekarošanas karus, kas savukārt izraisīja pret to nacionālus atbrīvošanās karus. Taču to daudzi poļu patrioti, norūpējušies tikai par savas dzimtenes brīvību, vairs nevēlējās saskatīt.

Vīlies Prūsijā, 1796.~gada oktobrī ģenerālis J.~H.~Dombrovskis piedāvāja Francijas Direktorijai (valdībai, 1795--1799) organizēt \strong{poļu leģionu}. Nākamā gada sākumā tika izveidoti divi tādi Francijas vasaļvalsts Cisalpīnas Republikas (itāļu \latxti{Repubblica Cisalpina}) Napoleona komandētās armijas sastāvā. Tajos galvenokārt dienēja poļi (1797.~gadā~--- ap 6~000), kuri bija dienējuši Austrijas armijā un saņemti gūstā vai arī bija dezertējuši no tās. Poļu leģionāri savu cīņu jēgu saskatīja Polijas valstiskuma atjaunošanā ar Francijas palīdzību. J.~H.~Dombrovskis izteicās: ``Kā uzvarētāji mēs uzcelsim no jauna mūsu Tēvzemi''. Leģioni kļuva par poļu nacionālisma audzināšanas skolu. Poļu leģionu karavīri nēsāja poļu mundierus ar franču kokardēm. J.~Dombrovskis ierosināja Napoleonam nosūtīt leģionus caur Turcijai pakļautajām teritorijām Balkānos uz Austrijai piederošo Galīciju, kur tika cerēts izraisīt poļu šļahtas sacelšanos. Tiesa, drīz (1797.~gada aprīlī) tika noslēgts Francijas un Austrijas miera līgums, kurš leģionāriem nāca pilnīgi negaidīts. Tas izsauca dziļu vilšanos poļu emigrantu vidū Francijā. (Daži pat devās uz Krieviju un iestājās tās armijā.) Tika kaldināti arī citi plāni kā atjaunot Polijas neatkarību. Tomēr lielākā daļa leģionāru naivi ticēja, ka Francija agri vai vēlu palīdzēs īstenot sapni par Polijas brīvību. Ne bez pamata ir krievu vēsturnieku izteiktā doma, ka poļu emigrācijas vadībai leģioni bija vajadzīgi nevis kā bāze nākamo Polijas bruņoto spēku veidošanai, bet kā ``avanss'' Francijai, lai pierādītu, cik svarīgi tai būtu atjaunot Poliju kā savu sabiedroto Eiropā.

Starp citu, poļu publicists un sabiedriskais darbinieks J.~Vibickis, lai iedvesmotu leģionārus, 1797.~gadā sacerēja vārdus plaši pazīstamajai ``Dombrovska mazurkai'' (``\pltxti{Mazurek Dąbrowskiego}''), sauktai arī par Dombrovska maršu, (Pirmais nosaukums gan bija ``Poļu leģionāru Itālijā dziesma''~--- ``\pltxti{Pieśń Legionów Polskich we Włoszech''}) Dziesmas pirmais pants un piedziedājums skanēja:

\vspace{1.5em}

\noindent
\begin{minipage}{0.45\textwidth}
\pltxti{
Jescze Polska nie zginela,\\
kiedu my źyjemu!\\
Co nam obca przemosc wzęeła,\\
szabla odbierzemy.\\
Marsz, marsz, Dąbrowski,\\
z ziemi wloskiej do poliski,\\
za twoim przewodem\\
złązcym się z narodem!}
\end{minipage}
\hspace{2em}
\begin{minipage}{0.45\textwidth}
Polija nav zudusi,\\
kamēr mēs vēl dzīvi!\\
Ko mums sveša vara ņēma,\\
Zobeniem atgūsim.\\
Marš, marš, Dombrovski,\\
no Itālijas uz Poliju,\\
mēs tavā vadībā\\
ar tautu vienosimies!
\end{minipage}

\vspace{1.5em}

% TODO verse example

Dziesmu izpildīja ar tautas mūziku mazurkas ritmā. Šī dziesma visur pavadīja poļu leģionārus, to viņi dziedāja, kad ar Napoleona armiju atgriezās dzimtenē, vēlāk slepenās sazvērnieku sapulcēs, 1831.~gadā tā kļuva par poļu sacelšanās himnu. Neraugoties uz to, ka poļu nebrīves gados tā tika aizliegta, kā aicinoša uz dumpi, viņi to turpināja dziedāt. 1926.~gadā ``Dombrovska mazurka'' kļuva par oficiālo Polijas valsts himnu.

Otrā pasaules kara laikā PSRS izveidotās T.~Kostjuško vārdā nosauktās poļu divīzijas karavīri himnu dziedāja bez pantiem, kuros bija rindas ``vācietis un krievs neizturēs, kad zobenu rokā ņemsim'' un ``visi vienā balsī runā, pietiks mums šīs nebrīves''. Taču mazurkā bija arī pants par poļu karavadoni S.~Čerņecki, kurš kļuva pazīstams ne tikai ar savu varonību, bet arī ar savu nežēlību, apspiežot pret poļu varu vērstās sacelšanās Ukrainā. Viņš arī apgānīja ukraiņu nacionālā varoņa B.~Hmeļņicka kapu un mirstīgās atliekas.

Pēc Otrā pasaules kara 1950.~gadā Polijas tautas republikā (PTR) kā valsts himna tika publicēti tikai pirmie divi panti, bet 1952. un 1953.~gadā tikai 1.~pants ar piedziedājumu. Izskanēja priekšlikums sarīkot konkursu jaunai himnai, taču tas tā arī netika realizēts. Mūsdienās himna tiek dziedāta bez ``vācieša un krieva'' pieminēšanas. Toties pants par S.~Čerņecki ir palicis.

Poļu leģionāri demonstrēja izcilu varonību kaujās, taču viņu liktenis arī parādīja, cik nesavienojama ir cīņa par brīvību ar kalpošanu iekarošanas politikai. 1798.~gadā pie diviem esošajiem leģioniem tika pievienots arī trešais, kurš nesa ``Donavas'' leģiona nosaukumu, kaut tika izvietots pie Reinas. Leģionāri, ciešot ievērojamus zaudējumus, piedalījās kara darbībā pret Napoleona pretiniekiem. Poļu leģionāriem bija visai nozīmīga vieta karos Itālijā, 1798.~gadā tie iegāja Romā, piedaloties Pāvesta valsts (itāļu \latxti{Stati Pontifici}) pagaidu likvidācijā. Kad 1799.~gadā Napoleons kļuva par pirmo konsulu~--- Francijas valsts galvu, viņš leģionus pārformēja, tie arī nomināli kļuva par Francijas armijas sastāvdaļu. Daļa leģionāru palika Itālijā un 1805.--1807.~gadā piedalījās Francijas vestajos karos, daļa (ap 5~000 ``durkļu'') 1801.~gadā, neraugoties uz pretestību, tika nosūtīta uz Haiti un tika izmantota cīņā pret nēģeru brīvības cīnītājiem F.~Tusēna-Luvertīra vadībā, kaujās un no slimībām zaudējot 2/3 sastāva. Pēc dažām ziņām tikai ap trīs simtu poļu atgriezās Francijā. (Pēc citām ziņām 1814.~gadā jau kā angļu gūstekņi Eiropā atgriezās ap 500 cilvēku.) Kopumā 1797.--1807.~gadā caur poļu leģioniem izgāja ap 35~000 cilvēku. Daļa no viņiem vēlāk veidoja Varšavas hercogistes armijas kodolu, daļa (ap 8~tūkstošiem) palika Francijas dienestā. No šī laika poļu leģiona ideja līdz pat valsts atjaunošanai starptautisku sarežģījumu brīžos atkal un atkal kļuva populāra poļu patriotu aprindās. Piemēram, Krimas kara (angļu \entxti{Crimean War}, franču \frtxti{guerre de Crimée}, krievu \rutxti{Крымская война}, 1853--1856) laikā izcilais poļu dzejnieks Ā.~Mickevičs Turcijas armijā mēģināja izveidot poļu daļas karam pret Krieviju.

Tāpat kā Francija, arī tās pretinieki: Krievija, Prūsija un Austrija, diplomātisku apsvērumu vadītas, dažkārt runāja par nākotnes plāniem atjaunot Polijas valstiskumu. Tādas baumas labprāt uzklausīja zināmas poļu aristokrātijas un arī zemākās šļahtas aprindas. Krievijā šādu plānu iniciators bija kņazs Ā.~Čartorijskis. Būdams Krievijas imperatoram Aleksandram~I tuva persona, viņš 1804--1806.~gadā bija pat Krievijas ārlietu ministrs. Ā.~Čartorijskis izstrādāja Žečpospolitas atjaunošanas plānu tās agrākajās robežās Aleksandra I protektorātā. Ar dažādu kompensāciju palīdzību Austrija būtu jāpierunā atdot Krievijai Galīciju, bet Prūsijai Lielpolija (\pltxti{Wielkopolska}~--- apgabals Polijas rietumos Vartas upes baseinā) būtu jāatņem ar spēku. Aleksandrs~I šo plānu tieši nenoraidīja, bet skatījās uz to kā uz diplomātisku Prūsijas ietekmēšanas līdzekli.

Jāuzsver, ka jau Krievijas imperators Pāvils I, tāpat kā vēlāk viņa dēls Aleksandrs I, mēģināja realizēt politiku, lai Krievija kļūtu it kā par Polijas labo aizbildni uz Austrijas un Prūsijas stingrākās varas fona. Īpaši Aleksandra I valdīšanās laikā Krievijas valsts veica pasākumus, lai radītu sev piekritējus poļu valdošajā slānī. Aleksandra~I liberālie žesti (konfiscēto muižu atdošana to agrākajiem īpašniekiem vai arī zaudējumu atlīdzināšana, nodokļu atvieglinājumi), pat privilēģiju piešķiršana (neraugoties uz vispārējo labības izvešanas aizliegumu, Aleksandrs~I piešķīra poļu muižniekiem atsevišķas atļaujas eksportēt labību, izmantojot Melnās jūras ostas) rada atbalsi poļu sabiedrībā. Nemazsvarīga loma bija ekonomisko sakaru nodibināšanai, agrāko Žečpospolitas apgabalu iekļaušanai Viskrievijas tirgū. Taču tai pat laikā Krievijas valdnieki nevarēja atteikties no bijušās Žečpospolitas baltkrievu un ukraiņu apgabaliem. Polija bija tradicionāls Krievijas ienaidnieks un Krievijas imperatori nevarēja atļauties atjaunot stipru varenu Poliju, kas varētu apdraudēt Krieviju.

\asterism

1806.~gadā, kad Napoleons uzsāka karu pret Prūsiju, Polija viņa plānos ieguva ievērojamu lomu. Napoleonu tā interesēja kā militārs placdarms, izdevīgs līdzeklis diplomātiskajā spēlē. Karā iesaistījās arī Krievija. Jauno kara fāzi Napoleons nosauca par ``Pirmo Polijas karu'' un uzdeva ģenerālim J.~H.~Dombrovskim organizēt poļu spēkus. J.~H.~Dombrovskis izveidoja vairākas divīzijas (ap 30~000 vīru) un tās sekmīgi karoja pret prūšiem un krieviem. 1806.~gada novembrī J.H.~Dombrovska ietekmē Pozenē notika poļu sacelšanās, poļi paši vairākās pilsētas atbruņoja Prūsijas garnizonus. Tas, protams, atviegloja Napoleona karaspēka darbību. 28.~novembrī franči iegāja Varšavā un 19.~decembrī tās iedzīvotāji uzgavilēja Napoleonam. Taču īstenot poļu cerības un atjaunot Polijas valsti viņš nesteidzās. 1807.~gada janvārī agrākajos Polijas apgabalos, kurus tagad okupēja Francija, Napoleons nodibināja ``Valdošo komisiju'' (``\pltxti{Komisja Rządząca''}), kura darbojās franču virsvadībā.

1807.~gada jūlijā, sakāvis Prūsiju un Krieviju, Napoleons par lielu vilšanos poļiem noslēdza Tilzītes mieru, pēc kura no poļu zemēm, kuras līdz tam ietilpa Prūsijas sastāvā (izņemot Rietumprūsiju un Pomerāniju (\pltxti{Pòmòrzé}), kas palika Prūsijā, un Dancigu (Gdaņsku), kura pirmo reizi [vēlāk tas notika pēc Pirmā pasaules kara] tika atzīta par ``brīvpilsētu'', taču ar franču garnizonu), izveidoja miniatūru \strong{Varšavas lielhercogisti} (\pltxti{Księstwo Warszawskie}, 1807--1815, faktiski~--- 1813), kurā ap 104~000 km$^{2}$ lielā apgabalā dzīvoja 2,6 miljoni iedzīvotāju. Hercogistes nosaukumam nebija nekāda vēsturiska pamata. Tas radās tikai tāpēc, ka abi imperatori: Napoleons I un Aleksandrs I, parakstot Tilzītes miera līgumu, nevēlējās lai jaunās Francijas vasaļvalsts oficiālajā nosaukumā figurētu vārds ``poļi'', ``Polija'' vai to atvasinājumi. 1809.~gadā pēc Napoleona uzvaroša kara pret Austriju hercogistei pievienoja teritorijas, kuras pēc trešās Žečpospolitas dalīšanas iegāja Austrijas sastāvā. Radās ap 160~000 km$^{2}$ liela ``valsts'' ar 4~350~000 iedzīvotāju. Nacionālā ziņā iedzīvotāju sastāvs bija daudz kompaktāks nekā Žečpospolitā. Pēc 1810.~gada statistikas hercogistē dzīvoja: 79\% poļu, ap 7\% ebreju, ap 6\% vāciešu, 8\%~--- lietuviešu un baltkrievu. Varšavas hercogistē administrācijā, tiesās, skolās tika lietota poļu valoda.

Pēc zināmām Napoleona svārstībām par hercogistes nominālo valdnieku kļuva Žečpospolitas priekšpēdējā karaļa Fridriha Augusta mazdēls, Saksijas karalis Fridrihs Augusts, kuru par likumīgu Polijas troņa mantinieku atzina vēl 1791.~gada Konstitūcija. Polijas valsts gan līdz tam jau nepastāvēja, taču poļu apdzīvotās teritorijas ar to nepilnu četrdesmit gadu laikā tika kārtējo reizi sadalītas un, kaut gan vēstures literatūrā to nav pieņemts darīt (Piemēram, Vācijas~--- Polijas attiecību speciālists vācu vēsturnieks V.~Jakobmeijers 1939.~gadā notikušo Polijas dalīšanu uzskata par piekto, pirms tam pieminot 1772., 1793., 1795. un 1815.~gada dalīšanas, bet aizmirstot par 1807.~gadā notikušo Napoleona īstenoto Varšavas hercogistes izveidi, kārtējo (ceturto) reizi pārdalot agrākās Žečpospolitas teritoriju), tomēr pēc autora viedokļa \strong{Napoleona veikto pārdali var kvalificēt par ceturto Polijas dalīšanu}. Pie tam Varšavas hercogiste pastāvēja vairākus gadus, šīs ceturtās dalīšanas rezultātā izveidotās robežas saglabājās ilgāk nekā pēc otrās dalīšanas radītās. Var jau uzskatīt Varšavas hercogistes radīšanu par Polijas atjaunošanas mēģinājumu, taču aiz šī valstiskā veidojuma palika daudzas etniskās poļu zemes.

Faktiski Varšavas hercogiste bija Francijas protektorāts, kurš piegādāja pēdējai pārtiku un ``lielgabalu gaļu''. Napoleons 1806.~gadā paziņoja: ``Polija ir pārāk sarežģīts jautājums: [poļi] pieļāva sadalīšanu, pārstāja būt par [vienotu] tautu, zaudēja sabiedrības garu, šļahtai tur ir pārāk liela loma, tautai pārāk maza. Tas ir līķis, kurā vispirms ir jāiedveš dzīvība, pirms es sākšu domāt, ko ar to darīt \citespace{} Es iegūšu no viņiem karavīrus, virsniekus, bet pēc tam paskatīšos [ko darīt tālāk].'' Tomēr pilnībā no Francijas atkarīgās hercogistes radīšanu daudzi poļi uzskatīja par Polijas atdzimšanas sākumu. 1807.~gada 22.~jūlijā Drēzdenē Napoleons izdeva hercogistes ``Konstitūciju'' (``\frtxti{Statut konstitutionennel du Duché de Varsovie''}). Napoleons Polijas pārvaldē balstījās uz tās muižniecību, taču visos oficiālos aktos izvairījās no vārdu ``polis'' vai ``poļu'' lietošanas, lai neradītu Austrijā un Krievijā bailes par to piesavināto poļu zemju drošību un pārliecinātu tās, ka Varšavas hercogistes radīšana nebūt vēl nenozīmē Polijas valsts atjaunošanu tās vēsturiskajās robežās. Dāvājot Konstitūciju ``Varšavas un Lielpolijas'' iedzīvotājiem, Napoleons uzņēmās saistības ``saskaņot to ar kaimiņvalstu mieru''. Konstitūcijā nebija norāžu uz vārda, preses, sapulču, biedrošanās brīvībām, personas un mantas neaizskaramību. Tiesa, atšķirībā no 1791.~gada 3.~maija Konstitūcijas visi iedzīvotāji bez kārtu atšķirībām tika pasludināti vienlīdzīgi likuma priekšā, taču šī vienlīdzība netika attiecināta uz ebrejiem. Seims tika vēlēts, otrā palāta~--- Senāts~--- monarha iecelts. Parlamenta kompetence bija ierobežota, tas varēja apspriest tikai budžeta, civilās un kriminālās likumdošanas jautājumus. 1807.~gadā tika atcelta dzimtbūšana, ieviesta zemnieku personīgā brīvība, taču zeme palika šļahtas rokās kā privātīpašums. Atstājot savu muižnieku, zemniekam bija pienākums ``atdot muižniekam pēdējā zemes īpašumu, kas sastāv no ēkām, inventāra un sējumiem''. Šļahtiči varēja zemniekus padzīt no viņu saimniecībām. Reformas īstenošana veicināja zemes koncentrāciju šļahtas rokās. 1810.~gadā hercogistē 54\% zemes atradās folvarkos, 2\%~--- baznīcas un 44\%~--- zemnieku rokās. Ap 20\% valsts muižu Napoleons izdāvāja saviem maršaliem un ģenerāļiem.

Interesantu Napoleona un poļu šļahtiču attiecību aspektu atklājis jau pieminētais poļu ekonomikas vēsturnieks V.~Kula. Tā Polijas daļa, kura pēc trim dalīšanām bija Prūsijas varā, bija nonākusi arī tās kredītiestāžu darbībā sfērā. Šīs iestādes bija radītās, lai atvieglotu prūšu muižniekiem (junkuriem) pāreju uz intensīvāku saimniekošanu. Iespēju saņemt valsts kredītus ieguva arī poļu muižnieki. Viņi to arī pilnā mērā izmantoja. Taču pēc Napoleona uzvaras pār Prūsiju viņš ar uzvarētāja tiesībām pasludināja sevi par šo kredītu īpašnieku un pieprasīja tos atmaksāt viņam. Izrādījās, ka atšķirībā no prūšu muižniecības lielākās daļas poļu muižnieki pārsvarā bija izmantojuši kredītus nevis savu muižu attīstībai, bet vienkārši izšķērdējuši. Poļu muižnieciskā historiogrāfija pat radīja tēzi, ka kredīti poļu šļahtičiem tikuši piešķirti ar nolūku tos izputināt, novest līdz bankrotam un tādā veidā iznīdēt arī ``poļu gara cietoksni''~--- šļahtiču muižas. Kad noskaidrojās, ka poļu šļahta nespēj samaksāt savus parādus Napoleonam, sākās sarunas. Varšavas hercogistes valdības delegācija sekoja Napoleonam pa visu Eiropu, lai vestu ar viņu cinisku kaulēšanos. Napoleons parādus pamazām atlaida~--- kā maksu par katru jaunu poļu pulku, kurš devās karot uz Spāniju. Tikai krietni vēlāk, XIX gadsimta vidū poļu muižnieki iemācījās izmantot naudu lai, atbilstoši tirgus principam, tā nestu jaunu naudu.

Arī hercogistei kopumā nācās ciest no Napoleona pastāvīgajām naudas prasībām. Tās saimniecisko attīstību traucēja viņa ievestā kontinentālā blokāde (1806--1814).

1808.~gadā tika ieviests Napoleona kodekss (\frtxti{Code Napoléon, arī Code Civil des Français}) kā hercogistes civillikums, kurš privātīpašuma principu sabiedriskajā dzīvē padarīja par izšķirošo, radot privātīpašniekiem sabiedriskās augšupejas iespējas neatkarīgi no kārtu piederības. Tas deva pirmo pamatīgo triecienu kārtu iekārtai, ievērojami sekmēja kapitālistisko attiecību attīstību feodālo vietā. Tomēr nedrīkst noklusēt, ka Napoleona kodeksa ieviešanu Varšavas hercogistes Seims pat īsti neapsprieda, uzstājās trīs deputāti, kuri ierosināja to pieņemt, uzreiz notika balsošana un ar 105 pret 2 balsīm kodekss tika pieņemts.

Tradicionālās pašvaldības tika atceltas un hercogisti pārvaldīja pēc Francijas departamentu sistēmas parauga. Hercogiste sākotnēji dalījās sešos departamentos un tie katrs desmit apriņķos. Pēc Napoleona uzvaras 1809.~gadā pār Austriju klāt nāca vēl četri departamenti no bijušajām Austrijas teritorijām. Bez tam Napoleons arī radīja atsevišķu 30~000 cilvēku lielu poļu armiju, no kuras daļa devās Napoleonam palīgā karot Spānijā. 1810.~gadā tās karavīru skaitu palielināja līdz 60~000 un 1812.~gadā līdz 75~000. Kopumā Varšavas hercogistes armijā ilgāku vai īsāku laiku dienēja ap 200~000 poļu. Komandēja armiju hercogistes kara ministrs J.~Poņatovskis, pēdējā Žečpospolitas karaļa Staņislava II Augusta Poņatovska brāļa dēls.

Kritiskāk noskaņotie poļi ironiski runāja par jauno valsti: ``Hercogiste Varšavas, monētas prūšu, armija poļu, karalis sakšu, bet kodekss franču'', taču apzinoties savu atkarību no Napoleona labvēlības, neuzdrošinājās atklāti kurnēt.

Par Varšavas hercogistes pastāvēšanas laiku eksistē anekdote.

\begin{quote}
Kāds ārzemnieks, nonācis Varšavas hercogistē, vietējiem iedzīvotājiem apjautājās:

``Kas tā ir par zemi?''

``Varšavas hercogiste,''~--- skanēja atbilde.

``Un kas ir tās valdnieks?''

``Saksijas karalis.''

``Un kāda viņam ir armija?''

``Poļu.''

``Bet kādas tiesības šeit pastāv?''

``Franču.''

``Kāda nauda?''

``Prūšu.''

``Tad gan jūs dzīvojat kā pie Bābeles torņa,''~--- nošūpoja galvu ārzemnieks.
\end{quote}


Poļu dižciltīgie baidījās par savu stāvokli, taču Napoleons, paredzot tos izmantot pret Krieviju, centās īstenot viņu interesēm atbilstošu politiku.

Lai arī Varšavas hercogiste nebija poļu sapņu piepildījums, tās radīšanai bija pozitīva nozīme. Hercogistē attīstījās nacionālā kultūra, zinātne, izglītība. Jau pašam faktam, ka pastāv valstisks veidojums, kurš turpina poļu valstiskuma tradīcijas, bija liela nozīme laikabiedru acīs. Polijas jautājums atkal tika izvirzīts aktuālās politikas ierindā, ar to nācās rēķināties Vīnes kongresam 1815.~gadā. Daudzi franču ieviestie likumi un institūcijas saglabāja savu iespaidu arī pēc hercogistes likvidācijas līdz pat XX gadsimtam.

Daļa poļu leģionāru arī pēc Varšavas hercogistes izveides palika Francijas armijā un palīdzēja Napoleonam okupēt Spāniju. Spāņi šo karu sauca~--- ``karš par Spānijas neatkarību'' (\estxti{Guerra de la Independencia Española}). Poļi cerēja, ka tādā ceļā viņi tuvina savas valsts atjaunošanu, kaut mūsdienās ir grūti izprast, kā, apspiežot citu tautu, var cīnīties par savas tautas neatkarību.

1810.--1811.~gadā Napoleons, cenšoties piesaistīt savā pusē poļu šļahtu, solīja tai Polijas atjaunošanu 1772.~gada robežās. Tiesa, ja solījums arī tiktu īstenots, šai valstij bija nolemts palikt Francijas atkarībā. 1810.~gada Francijas un Krievijas konvenciju, kurā Francija solīja nepaplašināt Varšavas hercogistes teritoriju, Napoleons atteicās ratificēt, taču šis solis vairāk vērtējams kā Aleksandra I šantažēšanas līdzeklis, nekā liecība par īstenajiem Napoleona nodomiem. Atbildot uz to, Aleksandrs I inspirēja baumas par autonomas Lietuvas lielkņazistes radīšanu no Krievijas rietumu guberņām, kura varētu kļūt par nākamās poļu valsts kodolu. Taču sarežģīto diplomātisko spēli pārtrauca 1812.~gadā sācies karš.

Neilgi pirms tā sākuma~--- 1812.~gada 10.~februārī poļu ģenerālis M.~Sokolņickis iesniedza Napoleonam detalizētu kara vešanas plānu pret Krieviju ar nosaukumu ``Par līdzekļiem kā Eiropai atbrīvoties no Krievijas ietekmes, bet pateicoties tam, arī no Anglijas ietekmes''. Tajā bija izstrādāts ne tikai Napoleona armijas virzīšanās maršruts, bet arī Krievijas sadalīšanas un kā neatkarīgas valsts likvidēšanas scenārijs. Francijas imperators varēja būt apmierināts, ka tā iekarošanas plānu realizēšanai atradās tāds iegansts kā Polija. 1812.~gada 22.jūnijā Napoleons parakstīja savu uzsaukumu Francijas ``Lielajai armijai'' (\frtxti{Grande Armée}): ``Karavīri, otrais poļu karš ir sācies. Pirmais beidzās Fridlandē [pilsēta~--- tagad Pravdinska, pie kuras notika uzvaroša Napoleona kauja ar Krievijas armiju 1807.~gada 14.~jūnijā] un Tilzītē. \citespace{} Tātad dosimies uz priekšu, pāriesim Nemunu, ienesīsim karu tās [Krievijas] teritorijā. Otrais poļu karš nesīs tādu pat slavu franču ieročiem kā pirmais. Taču miers, kuru mēs noslēgsim, tiks nodrošināts un nesīs beigas tai postošajai ietekmei, kurus Krievija nu jau 50 gadu kā atstāj uz Eiropu''. Sarunās ar poļu darbiniekiem Napoleons tieši norādīja, ka gaida poļu šļahtas uzstāšanās pret Krieviju lietuviešu, baltkrievu un ukraiņu zemēs. Taču 1812.~gada martā noslēgtajā līgumā ar Austriju Napoleons tai garantēja Galīcijas paturēšanu. Kāda varētu kļūt poļu valsts Napoleona uzvaras gadījumā, tā arī nebija skaidrs.

Naktī uz 24.~jūniju sākās pārcelšanās pār Nemunu un 300 poļu huzāru kā pirmie to šķērsoja. Poļu sajūsma bija liela. Ā.~Mickevicš to vēlāk (1834) tēloja poļu nacionālajā eposā ``\pltxti{Pan Tadeusz}'' (Tadeuša kungs). 26.~jūnijā Varšavas hercogistes Seims pasludināja, ka priekšā stāv Polijas ``atkalapvienošanās''. Poļu karavīri no Spānijas atgriezās Varšavas hercogistē. Napoleona pusē karojošais poļu korpuss bija visuzticamākais no ``Lielās armijas'' cittautiešu daļām. Polijas atbrīvošana kļuva par vienu no lozungiem, kas ļāva Napoleona armijā ievilināt daudz poļu. 1812.~gadā pirms karagājienu pret Krieviju Francijas imperatora rīcībā bija ap 85~000 (pēc citiem datiem~--- ne mazāk par 120~000) poļu karavīru. Kara gājienā uz Maskavu Napoleona \frtxti{Grande Armee} (Lielā armija) sastāvā bija ap 100~000 poļu karavīru (Citi dati: pret Krieviju Napoleona armijas sastāvā karoja 70 tūkstošu poļu ar 105 lielgabaliem.) Tādejādi katrs piektais Napoleona karavīrs bija polis. Viņi sastādīja gandrīz vai pusi no Napoleona kavalērijas. Taču visus poļus apvienot vienkopus Napoleons nevēlējās. Daļa no tiem tika izkaisīti starp citām franču armijas vienībām. Tikai ap 30 tūkstošu (citi dati~--- 37 tūkstošu) vīru lielu poļu kontingentu komandēja J.~Poņatovskis. Kopā ar viņu cīnījās arī poļu ģenerāļi J.H.~Dombrovskis, J.~Zaijončeks u.c.

T.~Kostjuško, kurš, atgriezies Eiropā, 1798.~gadā apmetās Parīzes apkaimē un piedalījās poļu politiskajā dzīvē, arī atbalstīja Poļu leģionu izveidi, 1799.~gada nogalē T.~Kostjuško tikās ar Napoleonu, taču, nesaņēmis no viņa solījumu par konstitucionālas Žečpospolitas atjaunošanu bijušajās robežās, uzskatīja par sev pazemojošu akli kalpot Francijas imperatoram. Viņš nepārcēlās arī uz Varšavas hercogisti, kaut 1807.~gadā sarunā ar Francijas policijas ministru J.~Fušē, piesolīja Napoleonam savu palīdzību, ja tas dotu rakstisku solījumu, publicētu avīzēs, ka Polijā tiks izveidota valsts iekārta līdzīga Anglijas iekārtai, ka zemnieki tiks atbrīvoti ar zemi un Polijas robežas sniegsies no Dancigas līdz Ungārijai, iekļaujot arī Galīciju. Atbildei uz to Napoleons rakstīja J.~Fušē: ``Es nepiešķiru nekādu nozīmi Kostjuško. Viņam savā zemē nav tās ietekmes, kurai viņš pats tic. Vispār, visa viņa uzvedība vieš pārliecību, ka viņš ir vienkārši muļķis. Vajag ļaut viņam darīt, ko vēlas, nepiegriežot viņam nekādu uzmanību''. Par Napoleonu T.~Kostjuško izteicās: ``Viņš domā tikai par sevi, nicina ikvienu lielu tautību un vēl vairāk neatkarības garu. Viņš ir tirāns.''

Poļu daļas pirmās iegāja krievu karaspēka pamestajā Viļņā. Taču poļi drīz bija vīlušies. Atkarojis Krievijai Lietuvu, Napoleons to nepievienoja Varšavas hercogistei, bet radīja patstāvīgu hercogisti, kura gan ilgi nepastāvēja.

Poļi piedalījās arī Borodinas u.c. kaujās pret krievu karaspēku, ciešot lielus zaudējumus. Pēc tam, kad jau bija pierādījusies krievu tautas pretestība iebrucējiem, Napoleona franču ģenerāļi ieteica viņam pasludināt dzimtbūšanas atcelšanu Krievijā, ko imperators tomēr nedarīja. Vācu vēsturnieks E.~Veiss uzskata, ka tas, iespējams, notika tāpēc, ka viņš negribēja kaitēt savai sabiedrotajai~--- poļu šļahtai. Tomēr jāsaka, ka pēc tam, kad krievu zemnieki bija izjutuši Napoleona karavīru pārestības, kuri atņēma tiem pēdējos mājlopus un pēdējo gabalu maizes, cerēt uz vietējo iedzīvotāju atbalstu iebrucēji diezin vai varēja, pat ja paziņotu par dzimtbūšanas jūga atcelšanu.

1812.~gada beigās, kad Varšavā nonāca ziņas par Napoleona armijas grūto stāvokli Krievijā, tika mēģināts savākt jaunus poļu brīvprātīgos, bet cerēto 30~tūkstošu vietā pieteicās vien ap 400~cilvēku. Atgriezušās Polijā, izretinātās poļu daļas ar J.~Poņatovski priekšgalā gan vēlējās aizstāvēt Varšavu pret Krievijas armiju, taču spēku nepietika un poļu karavīri devās līdzi Napoleona armijai, pametot Poliju.

Napoleons, atkāpjoties no Krievijas, izstrādāja divus variantus nākotnei. Pēc pirmā viņš bija gatavs kompromisa mieram ar Krieviju, solot tai nepieļaut Polijas valsts atjaunošanu. Viņš teica: ``Ja poļi neizpildīs savu pienākumu [domāta brīvprātīga masu mobilizācija karam pret Krieviju], tad Francijai un visai pasaulei miera jautājums vienkāršosies, jo tad ar Krieviju noslēgt mieru būs viegli.'' Otrais variants paredzēja piekāpšanos Austrijai, piešķirot tai teritorijas Balkānos, bet tai piederošo Galīciju atdodot Varšavas hercogistei un pasludinot Polijas atjaunošanu. Par Polijas karali tad vajadzētu kļūt kādam no Napoleona tuvākajiem līdzgaitniekiem. Izmantojot poļu sajūsmu par valstiskuma atgūšanu, Napoleons vēl cerēja vest uzvarošu karu pret Krieviju. Taču neviens no variantiem neīstenojās, jo Eiropas monarhi nevēlējās slēgt nekādas vienošanās ar Napoleonu.

1813.~gadā Napoleona pusē vēl karoja ap 40~000 poļu. J.~Poņatovskis kļuva par vienīgo ārzemnieku, kurš Leipcigas t.s. ``Tautu kaujas'' laikā (1813.~gada 16.--19.~oktobrī) no Napoleona rokām saņēma Francijas maršala zizli, taču pēc tam~--- arī pavēli piesegt franču daļu atkāpšanos. Franči priekšlaikus uzspridzināja tiltu un viņš, nevēloties padoties, ievainots mēģināja pārpeldēt Elsteres upi un noslīka.

Karojot Napoleona vadībā, 1812.~gadā poļi atstāja smagas atmiņas Krievijas civiliedzīvotājiem. Galvenie vardarbību un laupīšanu veicēji Krievijā bija nevis Napoleona disciplinētie francūži, bet vācieši un poļi. Par to rakstījuši daudzi krievu autori, sākot no A.~Puškina līdz J.~Tarlem. Kā rakstīja baltkrievu izcelsmes vēsturnieks profesors M.~Kojalovičs: ``Maskavā tauta neieredzēja un baidījās no tiem francūžiem, kuri saprata krieviski un runāja krieviski. Tie bija poļi.'' Krievijas apmija, sekojot Napoleona armijai un nonākot Polijā, bija gatava to pārvērst par tuksnesi, taču poļus no tā paglāba imperators Aleksandrs I, kurš Rietumeiropā vēlējās ieiet kā atbrīvotājs, nevis atriebējs. Krievu muižnieku aprindās gan klīda baumas, ka imperatora labvēlīgā attieksme pret poļiem ir saistīta ar viņa faktisko otro, neoficiālo sievu, dzimušo polieti A.~Nariškinu, ar kuru viņam bija vairāki bērni. Poļu šļahta kā savu līderi pie Aleksandra I nosūtīja Ā.~Čartorijski, taču viņš tā arī nesaņēma no Krievijas imperatora skaidru atbildi par Polijas nākotni. Līdz jautājuma galīgai starptautiskai izlemšanai Aleksandrs I kā augstāko pārvaldes iestādi lika radīt pagaidu Augstāko padomi (\pltxti{Rada Tumsczasowa Najwyźsza}), kuras sastāvā bija divi krievi, divi poļi un viens prūsis.

1814.~gadā Napoleona pusē vēl karoja ap 4~000 poļu. Pat pēc Napoleona atteikšanās no troņa viņam uzticīgie poļu brīvprātīgie izveidoja vieglās kavalērijas eskadronu, kurš pavadīja bijušo imperatoru uz Elbas salu. Šis eskadrons piedalījās arī Napoleona ``100 dienu'' epopejā un pilnībā gāja bojā Vaterlo (\entxti{Waterloo}, 1815.~gada 18.~jūnijā) kaujā. Vēl visu XIX gadsimtu poļu sabiedrībā pastāvēja kas līdzīgs Napoleona kultam.

Vācu vēsturniece A.~Šmidte-Roslere uzskata, ka to nosacīja Napoleona mīlas sakars ar poļu muižnieci M.~Vaļevsku un viņu kopējais dēls A.~Vaļevskis, kurš kļuva par Francijas pavalstnieku, bet tomēr piedalījās 1830--1831.~gada poļu sacelšanās norisē, tika apbalvots, atgriezās Francijā, kur Napoleona III laikā kļuva par tās ārlietu ministru. Vēsturniece atsaucas uz B.~Prusa romānu ``Lelle'' (``\pltxti{Lalka}''). Domājams, ka tomēr lielāka nozīme bija poļu patriotu utopiskajām cerībām, ka Napoleons uzvaras gadījumā pār Krieviju atjaunos Žežpospolitu, tāpēc simpātijas pret viņu neizdzisa vēl ilgi.

Šai sakarā gan krievu dzejnieks kņazs P.~Vjazemskis, kurš savu karjeru sāka Varšavā, jaunībā bija liberāli noskaņots un 1830.~gadā uzstājās pret krievu karaspēka ievešanu Polijā, savā dienasgrāmatā rakstīja: ``Lai cik arī poļi būtu truli, taču nevar iedomāties, ka vesela tauta brīvprātīgi ietu nāvē, pretī neglābjamai bojā ejai \dots{} Napoleons viņus sagūstīja ar divām, trijām frāzēm \citespace{} Ko viņš ir izdarījis Polijas labā? Griezies pie tās ar vairākiem madrigāliem savās proklamācijās, izdalījis tai dažus Goda Leģiona krustus, nopirktus ar poļu asins straumēm. Tas arī viss.'' Taču poļu patrioti kā slīcējs pie salmiņa pieķērās jebkurai cerībai, kas tiem solīja valstiskuma atjaunošanu.

Jau 1812.~gada decembrī, ieradies Viļņā, Krievijas imperators Aleksandrs I pasludināja vispārēju amnestiju Krievijas pavalstniekiem~--- poļiem, kuri bija karojuši Napoleona pusē, vēlāk atļāva izdzīvojušajiem poļu karavīriem atgriezties dzimtenē. Tur viņi tika iekļauti jaunveidojamās Polijas karalistes (\pltxti{Królestwo Polskie}) armijā. Ģenerālis J.~H.~Dombrovskis saņēma no Aleksandra~I kavalērijas ģenerāļa pakāpi, piedalījās Krievijas pakļautībā esošas Polijas karalistes jaunveidojamās armijas radīšanā un kļuva par Polijas senatoru. Ģenerālis J.~Zaijončeks kļuva par Aleksandra I vietvaldi (\rutxti{наместник}) Polijas karalistē, saņēma kņaza titulu.

Arī T.~Kostjuško 1814.~gadā griezās pie Aleksandra I ar vēstuli, kura saturēja padomus, kā labiekārtot Poliju. T.~Kostjuško lūdza Aleksandru I pasludināt sevi par Polijas karali, dot Polijai Konstitūciju, līdzīgu tai, kāda pastāvēja Anglijā utml. Sākotnēji Krievijas imperators pret poļu ģenerāli izturējās labvēlīgi, pat piedāvāja viņam vadīt Polijas karalistes administrāciju, taču pēc būtības atbildēja izvairīgi, ka cerot ``paveikt varonīgās tautas atdzimšanu''. T.~Kostjuško, uzzinot, ka Polija netiks atjaunota 1772.~gada robežās, atteicās sadarboties ar uzvarētājiem karā un devās uz Šveici, kur arī mira. T.~Kostjuško ķermenis tika iebalzamēts un 1818.~gadā pārvests uz Krakovu, apglabāts Vāvelas (\pltxti{Wawel}) pils katedrālē, kur atrodas Polijas karaļu kapenes.

\asterism

% page 69


1815.~gadā pēc Napoleona galīgas sakāves un uzvarētājvalstu sasauktā \strong{Vīnes kongresa} Napoleona radītā Varšavas hercogiste pazuda no kartēm. Padzenot Napoleona armiju, tās teritoriju (kas līdz hercogistes radīšanai piederēja Prūsijai un daļēji Austrijai) bija ieņēmis Krievijas karaspēks. Līdz Vīnes kongresam Krievijas rokās faktiski atradās 9/10 bijušās Žečpospolitas. Vīnes kongresā Aleksandrs I vēlējās panākt Varšavas hercogistes pievienošanu Krievijai, tikai vissliktākajā gadījumā piekāpjoties Prūsijai un atdodot tai daļu Pozenes novada. Taču Rietumvalstis, īpaši Austrija, Francija un Anglija, baidoties no Krievijas nostiprināšanās, nevēlējās visas Varšavas hercogistes zemes atstāt tai. Tā netika ievērotas arī poļu tautas objektīvās intereses radīt pēc iespējas vienotu teritoriju, kuras iedzīvotājiem būtu vieglāk cīnīties arī par savu valstiskumu.

Pēc Vīnes kongresa lēmuma Varšavas hercogistes rietumu daļa tika nodota Prūsijai un sāka saukties par Pozenes lielhercogisti (vācu \detxti{Großherzogtum Posen}; poļu \pltxti{Wielkie Księstwo Poznańskie}, 1818--1848). Tās teritorija gandrīz sasniedza 29~000 km$^{2}$, tur dzīvoja ap 0,9 miljoniem cilvēku. Varšavas hercogistes lielāko, austrumu daļu ar pašu Varšavas pilsētu pasludināja par Polijas karalisti (krievu~--- \rutxti{Царство Польское}, poļu~--- \pltxti{Królestwo Polskie}) un nodeva Krievijai (No 161~600 km$^{2}$ lielās Varšavas hercogistes 128~500 km$^{2}$ ar 2,7~miljoniem iedzīvotāju nonāca Krievijas rokās). Tā Krievija starptautiska lēmuma rezultātā kļuva arī par poļu etnisko zemju īpašnieci.

Krievu militārās tehnikas speciālists, kurš publicē arī darbus par vēsturi, A.~Širokorads raksta, ka ne Krievijas cars Aleksejs Mihailovičs, ne Katrīna II, ne J.~Staļins vēstures gaitā nepaņēma ne pēdu īstenas poļu zemes, bet atgrieza tikai Krievijai pirms tam piederošās, bet zaudētās teritorijas. Minētais publicists arī norādījis, ka zaudētu teritoriju atgūšanas piemēru Eiropas vēsturē ir daudz. Un nevienam neienāk prātā saukt Spānijas un Portugāles monarhus par agresoriem par to, ka rekonkististas (\estxti{Reconquista}~--- atkarošana) vairāku gadsimtu (722--1492) laikā viņi atkaroja Ibērijas pussalu mauriem, vai Francijas karaļus, kuri Simtgadu (1337--1453) kara laikā atguva britu sagrābtās zemes, kaut gan spāņi, gan francūži veica arī daudz zvērību, tostarp pret gūstekņiem un civiliedzīvotājiem.

Tad nu nākas atgādināt, ka Krievijā bija arī valdnieki, sākot ar Aleksandru I, kuri tomēr faktiski pievienoja poļu zemes Krievijas impērijai. Tiesa, šai gadījumā Eiropas valdnieki faktiski ar Polijas karalisti tikai atlīdzināja Krievijai par tās ieguldījumā karā pret Napoleonu, kura uzticami sabiedrotie bija poļi. Pēc šīs loģikas poļu zemju nonākšanas krievu rokās vaininieki bija paši poļi. Viņu zemju ar Varšavu nodošanu Krievijai varēja uzskatīt kā ``sodu'' poļiem par aktīvo Napoleona atbalstīšanu. Protams, viņiem bija cita loģika~--- poļu tautai bija tādas pat tiesības uz valstiskumu kā citām to ieguvušām tautām. Taču sabiedrotie tika meklēti pēc principa: ``Mana ienaidnieka ienaidnieks~--- mans draugs''. Tika aizmirsts, ka šim ``ienaidnieka ienaidniekam'' pirmajā vietā vienmēr bija (un būs) savas, nevis poļu intereses.

Polijas zemju pievienošana Krievijas impērijai no tās puses bija starptautiskās sabiedrības atbalstīta uzurpācija, neapstrīdama tik tiešām etnisko poļu zemju sagrābšana. Protams, imperators Aleksandrs I domāja pirmkārt par savām un Krievijas impērijas, nevis poļu interesēm. Turklāt Aleksandrs I 1815.~gadā formāli atjaunoja Polijas valstiskumu, tā saukto Polijas karalisti, un tā tika apvienota ar Krievijas impēriju uz personiskās ūnijas noteikumiem. Pašu poļu uztverē tas gan bija tikai jauns Krievijas solis Polijas atsevišķu daļu aneksijā.

PTR pastāvēšanas laikā oficiālie poļu autori rakstīja, ka poļu varonība Napoleona kara laikā esot nesusi savus augļus un tāpēc Vīnes kongresa esot nolemts saglabāt Polijas valstiskuma iedīgli, radot Polijas karalisti. Apgalvojums gan netika argumentēts. Faktiski Polijas karaliste radās dažādu valstu interešu sadursmju rezultātā, dalot poļu apdzīvotās teritorijas atbilstoši Eiropā pastāvošajam spēku samēram.

Mūsdienu baltkrievu autors, pedagoģijas un informātikas speciālists A.~Tarass, acīmredzot balstoties uz Aleksandra I Polijas karalistei piešķirto valstisko atribūtiku, uzskata, ka poļi ne velti karojuši zem Napoleona karogiem, jo 25 gadus no 1807. līdz 1831.~gadam viņu zeme bijusi suverēna valsts, lai arī atkarīga sākotnēji no Francijas, pēc tam no Krievijas. Var jau teikt, ka nav nevienas pilnībā suverēnas valsts, visas ir zināmā mērā atkarīgas no citām, taču minētais autors izliekas neredzot, ka suverēnas valstis vismaz piedalās savu robežu noteikšanā un tieši suverenitātes nepilnība noveda poļus pie 1830.~--- 1831.~gada sacelšanās.

Austrijai tika jau agrāk tai piederējusī Galīcija un vēl Veļičkas (\pltxti{Wieliczka}) apriņķis ar sāls raktuvēm. Austrija zaudēja teritorijas, ko tā bija guvusi trešajā un daļu~--- pirmajā Žečpospolitas dalīšanā. Tās nonāca Polijas karalistē. Austrijas Galīcijas apgabals bija 77~000 km$^{2}$ liels, tajā dzīvoja līdz 3,5~miljoniem iedzīvotāju. Krakova tika pasludināta par brīvpilsētu visu triju valstu aizbildniecībā. Tās apgabala (ap 1~000 km$^{2}$ ar 96~000 iedzīvotāju, no kuriem vairāk kā 20 tūkstošu dzīvoja pašā pilsētā, bet pārējie~--- tās apkārtnes ciemos) ``brīvība'' gan bija iluzora, jo faktiski valdīja triju uzvarētāju valstu rezidenti. Pēc Lielbritānijas pieprasījuma dalītājvalstis gan apņēmās ievērot poļu tautības un poļu valodas tiesības, taču stingri noteiktas tās nebija.

Vīnes kongress 1815.~gadā radīja tādu kārtību Eiropā, kad patstāvīga Polija nebija pastāvošās sistēmas nepieciešama sastāvdaļa. Polija palika sadalīta, Polijas jautājums~--- neatrisināts. Tā atbilstoši Vīnes kongresa lēmumiem 1815.~gadā faktiski atkal notika nu jau kārtējā \strong{Polijas teritorijas dalīšana}. (Kaut vācu vēsturnieki: A.~Šmidte-Roslere, L.~Dralle, poļu vēsturnieks G.~Labuda u.c. to sauc par ceturto Polijas dalīšanu, šī darba autors tomēr uzskata, ka pareizāk būtu runāt par trijām Polijas-Lietuvas valsts (Žečpospolitas) dalīšanām, kas notika XVIII gadsimtā, taču, ja šo uzskaiti turpina, par kārtējām bijušās valsts teritorijas dalīšanām jāuzskata gan Varšavas hercogistes radīšanu no agrāk Prūsijai un Austrijai piederošajiem pārsvarā poļu apdzīvotajiem apgabaliem, gan tās likvidēšanu. Tā XIX gadsimtā notika divas~--- ceturtā un piektā~--- šīs teritorijas dalīšanas, bet XX gadsimtā jau nākamās~--- sestā (1939.~gadā) un septītā (1945). Tiesa, jāuzsver vēl viena aspekts. XVIII gadsimtā trīs reizes tika dalīta Polijas-Lietuvas valsts (Žečpospolitas) daudzu tautu apdzīvota teritorija, XIX gadsimtā tika dalīta nu jau vienīgi poļu teritorijas, bet XX gadsimtā (1939. un 1945.~gadā) jau Polija zaudēja 1918.~--- 1920.~gadu laikā ar spēku sagrābtās teritorijas, kuras apdzīvoja kā poļi, tā arī vācieši, lietuvieši, ukraiņi, baltkrievi.

G.~Labuda uzsvēris, ka Vīnes kongresa veiktā dalīšana daudz būtiskāk kā iepriekšējās salauza saimniecisko, politisko un etnisko poļu zemju vienotību.

Marksisma teorētiķi Polijas dalīšanu kvalificēja kā ``vislielāko noziegumu''. Līdzdalība Polijas dalīšanā vairākus gadu desmitus saistīja valstis~--- tās dalībnieces ar kopējām interesēm ne tikai poļu tautas nacionālās atbrīvošanās kustības, bet zināmā mērā arī citu tautu nacionālās kustības apspiešanā. Tieši ar Vīnes kongresu radās triju monarhiju: Austrijas, Krievijas un Prūsijas t.s. Svētā savienība (franču \frtxti{La Sainte-Alliance}, vācu \detxti{Heilige Allianz}, krievu \rutxti{Священный союз}), izteikti konservatīvs spēks Eiropā. Dalība Polijas dalīšanā līdz pat 1914.~gadam cementēja triju monarhiju sadarbību un sargāja Austrumeiropu pret kariem un revolucionāriem satricinājumiem.

Kopumā Prūsija saņēma 8\%, Austrija 10\% un Krievija 82\% agrākās Žečpospolitas teritorijas. Padomju vēsturnieks A.~Manusevičs ne bez pamata apgalvoja, ka Polijas karalistes radīšana nesa zināmu atvieglojumu poļu tautas liktenī salīdzinājumā ar iepriekšējo stāvokli un radīja labvēlīgākus apstākļus attīstībai. Taču tas nevar mīkstināt secinājumu, ka faktiski kārtējā Polijas teritorijas pārdale, ko veica Vīnes kongresā, bija tāds pats noziegums pret poļu tautu kā Žečpospolitas dalīšanas XVIII gadsimtā. Pie tam ar centrālā Polijas apgabala nonākšanu Krievijas impērijas rokās poļu tautai radās vēl viens, jauns apspiedējs~--- Krievijas carisms, pret kuru poļiem nācās cīnīties par savu brīvību. Padomju vēsturnieki gan uzsvēra, ka līdz ar to poļu tauta ieguva arī sabiedroto cīņā pret carismu~--- krievu tautu (pareizāk gan būtu teikt~--- Krieviju apdzīvojošās tautas), taču autoram jāatzīmē, ka šis sabiedrotais tolaik bija vājāks par jauniegūto apspiedēju, tāpēc no poļu tautas viedokļa notikumu attīstību nekādi nevar vērtēt kā pozitīvu.

Žečpospolitas sadales, Napoleona kari postoši ietekmēja gandrīz visus banķierus un lieluzņēmējus. Rezultātā pilsonībai (buržuāzijai), kas veidojās 19.~gadsimtā, nebija pamatu iepriekšējā laikmetā, nebija savu tradīciju. Otra jaunveidojamās pilsonības īpatnība bija tā, ka lielā mērā tā veidojās no nepoļu izcelsmes elementiem. Tāda pilsonība politiski ilgi vēl nevarēja konkurēt ar šļahtu. Līdz XIX gadsimta 60.~gadiem buržuāzijā valdošie bija valdībai labvēlīgi uzskati, ko lielā mērā noteica tās atbalsts rūpniecības un tirdzniecības attīstībai un dažviet (piemēram, Lodzā) nepoļu pilsonības pārsvars.

Pēc Polijas sadales svešzemnieku vara veda pie tā, ka poļu etniskās teritorijas nomalēs norisa pretrunīgi procesi. No vienas puses, lauku iedzīvotāji sāka zaudēt nacionālo (etnisko) raksturu. Denacionalizējās daļa sīkās šļahtas un inteliģences, kura saistīja savu pastāvēšanu ar kalpošanu svešajām varām. Aristokrātijai, lielajiem magnātiem vispār bija raksturīgas kosmopolītiskas tendences. Taču no otras puses, tai pat laikā nacionālā apspiešanas politika, kuru piekopa Žečpospolitu sadalījušās valstis, izsauca pretestību un stimulēja poļu nacionālās pašapziņas izaugsmi. Tādejādi, ja nacionāla valstiskuma, nacionālas valodas un kultūras pastāvēšanai poļu nācijas, tās pašapziņas veidošanā bija nenoliedzami pozitīva loma, tad Polijas sadalei un nacionālu spaidu ieviešanai bija divejāda: gan bremzējoša, gan stimulējoša ietekme.

\asterism

\strong{Poļi XIX gadsimtā dzīvoja arī ārpus nosauktajām teritorijām}. Viņu kopējais skaits sniedzās līdz 9 miljoniem. Salīdzinoši daudz viņu bija \strong{Silēzijā} (vācu \detxti{Schlesien}, poļu \pltxti{Śląsk}), kura jau iepriekš ietilpa Prūsijas sastāvā un veidoja īpašu Silēzijas provinci (\detxti{Provinz Schlesien}).

Lejassilēzijā (vācu \detxti{Niederschlesien}, poļu Dolny Śląsk) sakarā ar tekstilrūpniecības, bet kopš XIX gadsimta sākuma arī metalurģijas strauju attīstību un valdības pasākumiem vācu kolonizācijas veicināšanai ieplūda daudz vācu iedzīvotāju un plaši izplatījās vācu valoda. Šeit pastāvēja Prūsijā lielākie muižnieku zemes īpašumi. XIX gadsimta beigās vairāk nekā 90\% Lejassilēzijas iedzīvotāju par dzimto uzskatīja vācu valodu.

Augšsilēzijā (Vācu \detxti{Oberschlesien}, poļu \pltxti{Górny Śląsk}) kur jau XVIII gadsimta otrajā pusē tika atklātas pirmās akmeņogļu ieguves šahtas (pati pirmā~--- 1751.~gadā), agrāk kā citur (jau XIX gadsimta 40.~gados) notika rūpniecības apvērsums, šahtās sāka pielietot tvaika mašīnas ūdens atpumpēšanai, attīstījās melnā un krāsainā metalurģija. XIX gadsimta pirmajā pusē akmeņogļu ieguves šahtu skaits pieauga trīs reizes, dzelzs ražošana~--- sešas reizes, cinka ražošana~--- trīs reizes un sastādīja ap 75\% no Prūsijas un 40\% no visas pasaules ieguves. 1862.~gadā Augšsilēzijā parādījās Prūsijā pirmās martena krāsnis. Raktuves, metalurģiskie uzņēmumi u.c. atradās Prūsijas valdības un prūšu aristokrātijas rokās, kas ekspluatēja poļu darbaļaudis. Augšsilēzijā jau sākotnēji šķiriskais dalījums sakrita ar nacionālo. Poļu valoda gan ilgstoši kalpoja kā mājas sarunu līdzeklis, bet arī šeit sākās poļu ģermanizācijas process. Tāpēc līdz Polijas valsts atjaunošanai 1918.~gadā sabiedriskie procesi Silēzijā tikai atbalsoja nacionālās cīņas citos poļu apdzīvotajos rajonos.

\strong{Poļi dzīvoja arī} Polijas karalistē neietilpstošajās \strong{Krievijas Ziemeļrietumu novada} Viļņas, Kauņas, Grodņas, Minskas, Mogiļevas un Vitebskas (XIX gadsimtā dalījums nedaudz mainījās) \strong{guberņās}, tai skaitā arī Vitebskas guberņā ietilpstošajā Latgalē, \strong{un Dienvidrietumu novada} (jeb labā krasta Ukrainas) Podolijas, Volīnijas un Kijevas \strong{guberņās}. Gandrīz pusi, ap 40\% no visiem šo apgabalu iedzīvotāju sastādīja katoļi un uniāti, ap 10\% ebreji, 5\% protestanti. Izglītības sistēmā valdošā bija poļu valoda. Viļņas, Grodņas, Minskas, Volīnijas un Podolijas guberņās sākotnēji poļi civilajā un militārajā pārvaldē ne tikai saglabāja, bet arī nostiprināja savu lomu. Šeit poļu muižnieki baudīja privilēģijas uz pārējo iedzīvotāju rēķina, taču sakarā ar poļu nacionālās atbrīvošanās cīņu pakāpeniski viņu stāvoklis sarežģījās. Citāds stāvoklis bija Vitebskas, Mogiļevas un Kijevas guberņās, kur gandrīz visu ierēdniecību veidoja krievi vai pārkrievojušies ārzemnieki.

Atsevišķo poļu apdzīvoto apgabalu būtiskās atšķirības ekonomiskās attīstības līmenī noteica ne tik daudz darba dalīšana starp vienas zemes novadiem, cik dažādie politiski-ekonomiskie apstākļi, kādos tie nonāca. Līdzās vāciešiem un itāļiem poļi bija pēc skaita trešā lielākā Eiropas tauta, kura XIX~gadsimtā cīnījās par nacionālo apvienošanos un vienota valstiskuma izveidi. Atšķirībā no pirmajām divām, poļiem tas neizdevās. Pārējās valstis, panākušas savu neatkarību, to nedeva Polijai, jo tas apdraudētu stabilitāti Eiropā.

Jāatzīmē, ka ārpus Polijas karalistes teritorijas palikušās agrākās Polijas-Lietuvas zemes, kuras apdzīvoja ukraiņi, baltkrievi, lietuvieši un latvieši (latgalieši) un kuras tagad ietilpa tieši Krievijas impērijas, bet nevis Polijas karalistes sastāvā, paši poļi toreiz sauca par sagrābtajām, t.i.~--- viņiem atņemtajām zemēm. Turpretī, piemēram, no baltkrievu nacionālā viedokļa šo novadu pievienošana Krievijai pārsvarā jāvērtē pozitīvi, jo notika radniecīgu, vēsturiski, kultūras un valodas ziņā tuvu slāvu tautu apvienošanās vienā spēcīgā valstī. Šeit stabilizējās politiskā situācija, pārtrūka piespiedu katolizācija un polonizācija (citu, nepoļu tautu asimilācija, pārpoļošana). Tas gan notika ne uzreiz. Šajās bijušajās Polijas-Lietuvas valsts austrumu zemēs poļi bija ievērojamā mazākumā, tomēr ierēdņu vidū viņu bija daudz, arī augstākajos amatos, jo pārējie iedzīvotāji izglītības līmeņa ziņā no poļiem atpalika. Poļu iecelšanu ierēdņu amatos veicināja arī lielkņazs Konstantīns, kuram Aleksandrs I piešķīra militāro un civilo varu arī Krievijas rietumu guberņās. Tikai Nikolaja I valdīšanas laikā (1825--1855) poļu ietekmi šeit sāka ierobežot.

Protams, pievienošana Krievijai arī ukraiņiem un baltkrieviem nesa ēnas puses, carisms savas politikas realizācijai izmantoja poļu valdošos slāņus, ukraiņu un baltkrievu tautai netika piešķirta valstiskums kaut vai autonomijas formā, to apdzīvotie apgabali tika sadalīti dažādās guberņās, saglabājās feodālā apspiestība. Poļu muižnieki, ja viņi neatteicās dot uzticības zvērestu Krievijas imperatoriem, saglabāja savu varu pār vietējiem zemniekiem.

Arī Latgalē stāvoklis bija līdzīgs kā lietuviešu, ukraiņu un baltkrievu apdzīvotajos apgabalos. Kā to parāda latviešu vēsturnieks Ē.~Jēkabsons, poļu muižnieku attieksme pret vietējiem zemniekiem neatšķīrās no vācu muižniecības attieksmes pret viņiem citās Latvijas daļās. Katram, kurš kaut ko gribēja dzīvē sasniegt, bija jārunā poliski un jākļūst par poli. Pieņemot darbiniekus muižās, priekšroka tika dota tiem, kuri runāja poliski. Atbildīgu amatu pildītājus, piemēram, muižnieku bērnu skolotājus, bieži aicināja no Polijas.

\subsection{Galīcija Austrijas / Austroungārijas valdījumā}

Austrija (no 1866.~gada Austroungārija) izrādījās visvājākā no trim lielvalstīm, tas ietekmēja arī poļu stāvokli tās pārvaldītajās teritorijās. Jau imperators Jozefs II 1781.~gadā uzsāka zemnieku brīvlaišanu (tika mazināta zemnieku personīgā atkarība no muižniekiem, zemnieku bērni ieguva iespēju doties uz pilsētām u.tml., taču klaušas vēl saglabājās), veicinot zemnieku zemes izpirkšanu, taču šis process ievilkās, jaunu impulsu saņemot tikai 1848.~gada revolūcijas laikā. Galīcija jeb pēc oficiālā nosaukuma~--- Galīcijas un Lodomērijas karaliste (vācu \detxti{Königreich Galizien und Lodomerien,} poļu \pltxti{Królestwo Galicji i Lodomerii),} kur karalistei bija pievienots arī Bukovinas novads (ukraiņu \uktxti{Буковина}, rumāņu \rotxti{Bucovina}, 1775.~gadā Krievija to atdeva Austrijai, tā to 1786.~gadā pievienoja Galīcijai, bet 1849.~gadā pārvērta par atsevišķu apgabalu; mūsdienās daļa Bukovinas ietilpst Ukrainā, otra daļa Rumānijā) pēc Vīnes kongresa tika sadalīta 18 apgabalos (\detxti{Kreis}), atjaunota centralizēta pārvalde, birokrāti bija galvenokārt etniskie vācieši (austrieši) vai pārvācoti čehi. Valsts valoda bija vācu, tajā darbojās pārvaldes iestādes. Provincei nebija vietējās autonomijas. Tika gan atjaunots Seims (\pltxti{Sejm Stanowy Królewstwa Galicji i Lodomerii}), kurā ievēlēja augstāko slāņu~--- garīdzniecības, muižniecības~--- magnātu un šļahtas~--- un arī Lembergas (Ļvovas) pilsētas~--- Galīcijas administratīvā un politiskā centra~--- pārstāvjus, taču tā pilnvaras bija niecīgas. Tāpēc vācu vēsturnieks E.~Meijers uzskatīja, ka pēc 1815.~gada poļiem Galīcijā bija mazāk brīvību nekā pārējās Polijas daļās.

Ekonomika šai novadā attīstījās vājāk kā Prūsijai un Krievijai pakļautajos poļu apgabalos. Galīcija bija arī viena no Austrijas visvairāk atpalikušajām provincēm. Vācieši to dēvēja par Habsburgu monarhijas nabagmāju (\detxti{Armenhaus}). Poļu valodā to bieži sauca par \pltxti{Golicija i Głodomorija} [vārdu spēle: kails \pltxti{(goły)} un izbadējies (\pltxti{głodny})]. Atpalicība no citiem poļu novadiem sāka mazināties tikai pēc 1830.--1831.~gada sacelšanās Polijas karalistē, kad daudzi poļu emigranti bēga uz Galīciju. Tiesa, liela daļa no viņiem Austrijas valdības spiediena rezultātā bija spiesta aizbraukt uz Franciju, taču daļa tomēr ieguva pilsonību. Šeit visilgāk valstī saglabājās zemnieku dzimtbūtnieciskā atkarība (personīgo brīvību zemnieki Galīcijā saņēma tikai 1848.~gadā), novecojušas saimniekošanas formas. Dzelzceļa celtniecība šeit sākās vēlāk nekā pārējās poļu zemēs. Krakova tikai 1855.~gadā tika savienota ar Vīni, 1861.~gadā~--- ar Lembergu, nākamajos gados dzelzceļa stigu turpināja līdz Krievijas un Rumānijas robežai. Pēc tam pāri Karpatu kalniem tika uzbūvētas četras dzelzceļa līnijas, kuras savienoja Galīciju ar Ungāriju, taču tām bija galvenokārt militāri-stratēģiska nozīme.

Kā uzskata krievu autori, tieši iepriekšējā gadsimtiem ilgā poļu kundzība bija novedusi novadu pie pilnīga panīkuma. Galīcijā zemniekus nospieda ļoti augsts nodevu slogs, kas neļāva modernizēt saimniecību. Ebreju izcelsmes poļu vēsturnieks Š.~Aškenazi rakstīja, ka Galīcijas sīkpilsonis visa veida nodokļos maksāja 16 reizes vairāk nekā sīkpilsonis Polijas karalistē. Pie tam, ja pēdējā valsts kases ieņēmumi tika izlietoti tikai vietējām vajadzībām, tad no Galīcijā ievāktajiem līdzekļiem 25\% ``aizpeldēja'' uz Vīni. 80\% Galīcijas iedzīvotāju bija zemnieki, 27\% no viņiem piederēja saimniecības mazākas par 1 ha, 42\% no tām bija starp 1 un 5 ha. Tā kā nenotika lauksaimnieciskās ražošanas modernizācija un mehanizācija, produkcijas ražošana nepieauga. XX gs sākumā Galīcijā vajadzēja pat ievest labību no Ungārijas. Ja 1869.~gadā Galīcijā bija 5,45~miljoni iedzīvotāju, līdz 1910.~gadam viņu skaits pieauga līdz 8 miljoniem. No 1880. līdz 1910.~gadam 600~000 cilvēku devās uz ārzemēm, lielākoties pāri okeānam. Nabadzība un higiēnas trūkums veda pie augstākās bērnu mirstības Habsburgu monarhijā. Arī izglītības ziņā Galīcija, gan kopā ar Bukovinu un Dalmāciju, atradās pēdējā vietā Austroungārijas impērijā. 1890.~gadā 64,8\% vīriešu un 71,6\% sieviešu Galīcijā neprata lasīt un rakstīt. Muižnieki pat iedomāties nevarēja, ka viņi brīvprātīgi varētu likvidēt klaušu sistēmu, viņi arī neko nedarīja, lai uzlabotu savu zemnieku stāvokli.

Rūpniecība praktiski neattīstījās, jo bija vājš iekšējais tirgus. Turklāt vairāk attīstītajās austriešu un čehu apdzīvotajās zemēs pastāvēja uzņēmumi, kuru konkurenci Galīcijas uzņēmumi nevarēja izturēt. 1822.~gadā Galīcijā gan pastāvēja ap 40 nelielas metālizstrādājumu fabrikas, kas pārstrādāja zemas kvalitātes dzelzsrūdu, taču to skaits pastāvīgi samazinājās. Pēc 1841.~gada datiem Galīcijā un Bukovinā, kur dzīvoja 28,7\% visu impērijas iedzīvotāju, atradās tikai 3,4\% visu Austrijas (bez Ungārijas) uzņēmumu. 22\% Galīcijas rūpnieciskās produkcijas veidoja šņabis. XIX gadsimta otrajā pusē, dzelzceļa satiksmes ar Austriju izveides un ar to saistītās metālizstrādājumu ievešanas rezultātā Galīcijā iezīmējās krasa metalurģijas krīze.

Tomēr atzīmējams, ka kopš XIX gadsimta otrās puses Galīcijā Drohobičā (\pltxti{Drohobycz}) darbojās lielākās naftas ieguves Eiropā, atklātas 1846.~gadā. 1866.~gadā šeit tika uzcelta pirmā naftas pārstrādes rūpnīca Centrāleiropā, 1910.~gadā~--- otrā. Līdz 1900.~gadam iegūstamais naftas daudzums sasniedza 1.~miljonu tonnu, bet 1912.~gadā Austroungārijā, galvenokārt Galīcijā, jau ieguva 2,9~miljonus tonnu. Ar to Austroungārija bija trešā lielākā naftu iegūstošā valsts pasaulē pēc ASV un Krievijas.

Stāvokli sarežģīja \strong{raibais nacionālais sastāvs}. Muižnieki~--- zemes īpašnieki vairumā bija poļi jeb ungāri, bet zemnieki, īpaši provinces austrumos~--- rusīni jeb \strong{rutēņi} (Habsburgu impērijā no 1848.~gada oficiālajos dokumentos vācu valodā tika lietots nosaukums \detxti{Ruthenen}, tā izdalot atsevišķu no poļiem un krieviem nacionālo grupu. Tāpēc, runājot par XIX gadsimtu, autors lieto terminu \lttxti{rutēņi}, kaut arī krievu vēsturnieki dod priekšroku terminam ``rusīni'', kas uzsver to kopību ar krievu tautu). Vietējā muižniecība, arī tirgotāji un augstākā garīdzniecība bija pārpoļota, tikai mājās vēl lietojot tautas valodu. Savai tautībai uzticīgi bija palikuši tikai sīkie amatnieki, zemnieki un zemākie garīdznieki. Pēc dažiem datiem Galīcijā dzīvoja 47\% poļu, 45\% rutēņu un 6\% ebreju (Pēc citām ziņām poļu bija 45,9\%, rutēņu~--- 42\%, ebreju~--- 10\%. Krievu vēsturniece A.~Bahturina raksta, ka XIX gadsimta vidū novada iedzīvotāju vairākumu sastādījuši rusīni~--- 43,7\% un ebreji~--- 11,8\%. Jāsaka, ka sarežģītā Austrumgalīcijas attīstība ļāva krieviem tur saskatīt krievu, poļiem~--- poļu vairākumu.). Rutēņu inteliģences Galīcijā gandrīz nebija, tāpat kā nebija augstāko un vidējo mācību iestāžu. Skolās, kuru sistēma šai laikā Austrijā vēl nebija sevišķi attīstīta, notika pārvācošana (ģermanizācija), taču sākumskolās, īpaši lauku, mācības norisa dzimtajā valodā. Habsburgu monarhija īpaši necentās pārvācot vietējos iedzīvotājus. Tomēr pakāpeniskā politiskās sistēmas liberalizācija veda arī pie Galīcijas rutēņu nacionālās kustības pacēluma. Tās priekšgalā nostājās uniātu baznīca.

No otras puses, poļu nacionālajai kustībai šeit bija senas tradīcijas. 1817.~gadā atjaunotajā Lembergas (Ļvovas) universitātē (dibināta 1784.~gadā) mācības gan norisa latīņu un vācu valodā, bet pēc diviem gadiem tur nodibināja poļu valodas katedru. Poļu muižniecības bērnu vajadzībām bija atvērtas vairākas ģimnāzijas. No Galīcijas aktīvākie poļu patrioti devās uz Polijas karalisti, lai piedalītos 1830.~gada sacelšanās. Tieši no šīs kustības izplatības baidījās Austrijas valdošie slāņi, jo Galīcija iedzīvotāju skaita ziņā sastādīja gandrīz 1/7, bet teritorijas ziņā~--- 1/8 daļu no toreizējās Habsburgu impērijas. Austrieši ievēroja rutēņu masu noskaņojumu un veicināja to pretpoļu un proaustriešu aktivitātes.

Galīcijā bija trīs katoļu eparhijas, kurās darbojās ap 700~draudžu un 900~baznīcu, ap 50~vīriešu un 15~sieviešu klosteru ar ap 500~mūkiem un mūķenēm. Bez tam novadā pastāvēja arī divas uniātu eparhijas ar vairāk nekā 200~draudžu un ap 300~baznīcu. Uniātu baznīcas rokās atradās gandrīz visas rutēņu nacionālās kustības aktivitātes, vērstas uz atbrīvošanos no poļu un krievu ietekmes.

Jaatzīmē, ka Austrijas sastāvā līdz ar Polijai atsavinātajām teritorijām nonāca arī ap 200~000 ebreju. Austrijas imperators Jozefs II jau 1781.~gadā izdeva dekrētu par reliģisko toleranci, ar ko uzlaboja arī ebreju stāvokli. Taču mērķis bija ebreju asimilācija. Tikai pēc tās viņi varētu iegūt līdztiesību.

Blakus esošajā Krakovas brīvpilsētā (\pltxti{Rzeczpospolita Krakowska}), kura 1815.~gadā saņēma savu satversmi, saglabājās demokrātiskā kustība un Polijas valsts atjaunošanas ideja. 1364.~gadā dibinātajā Krakovas universitātē veidojās nacionāli noskaņoto pulciņi. Autonomā Krakova kalpoja par vispolijas tirdzniecības centru. Kad no kaimiņos esošās Polijas karalistes pēc 1830./1831.~gada sacelšanās šurp bēga tās dalībnieki, Krievijas karaspēks uz laiku okupēja brīvpilsētu. Tās autonomijas tiesības tika ierobežotas. Taču drīz atkal radās konspiratīvas poļu patriotu grupas. Tad pilsētu līdz 1841.~gadam okupēja Austrijas karaspēks. Tomēr Krakova turpināja eksistēt kā brīvdomības un poļu patriotisma centrs.

Stāvoklis Galīcijā saasinājās jau 1846.~gadā. Poļu šļahta mēģināja sacelties pret Austriju. Iniciatīvu parādīja poļu revolucionārie demokrāti, kuru vidū liela loma bija filozofam, literātam un sabiedriskam darbiniekam E.~Dembovskim. Šī grupa cīņu par Polijas nacionālo atbrīvošanos centās apvienot ar zemnieku antifeodālo prasību apmierināšanu. Sacelšanos atbalstīja arī Galīcijas liberāļi. Revolucionārās cīņas sākās Krakovā (\pltxti{Rewolucja Krakowska}). 22.~februārī sacēlušies izveidoja t.s. Polijas Republikas Nacionālo valdību (\pltxti{Rząd Narodowy Rzeczypospolitej Polskiej}), kura manifestā ``Visiem poļiem, kuri prot lasīt!'' aicināja tautu celties cīņai pret trijām Polijas sadalītājvalstīm par nacionālo neatkarību, pirmo reizi Polijas vēsturē pasludināja radikālas sociālās reformas: feodālo klausību atcelšanu, nekavējošos zemnieku apstrādātās zemes nodošanu viņu īpašumā, zemes piešķiršanu bezzemniekiem, algas paaugstināšanu strādniekiem un amatniekiem, sabiedrisko darbu ieviešanu tiem. Uzsaukumā ``Brāļiem izraēlītiem'' tika pasludināta ebreju emancipācija. Tā Krakovā 1846.~gadā, divus gadus agrāk nekā pārējā Eiropā, izvērsās revolucionāras cīņas ar vispārnacionālu nozīmi. Taču t.s. Nacionālā valdība tā arī neko būtisku neizdarīja lai pārnestu revolucionāro kustību ārpus Krakovas, izplatītu sacelšanos zemniecībā. Neliela E.~Dembovska vadīta vienība, kura devās iesaistīt cīņā apkārtējos zemniekus, 27.~februārī cieta sakāvi sadursmē ar austriešu karaspēku, pats viņš gāja bojā. 2.~martā ap 1~500 poļu~--- Polijas Republikas Nacionālās valdības piekritēju bēga uz Pozeni, 4.~martā austriešu armija ieņēma Krakovu.

Blakus Krakovai Galīcijas laukos austriešu administrācijai izdevās pārliecināt vietējo zemniecību nostāties pret poļu ``paniem''~--- saviem muižniekiem, kuri šos zemniekus nežēlīgi apspieda. Jau pirms sacelšanās Krakovā bija sākušies zemnieku nemieri (\pltxti{rabacja}). Tagad austriešu varasvīru uzkūdītie zemnieki pavasarī izvērsa tā saucamo ``Galīcijas slaktiņu'' (\pltxti{Rzeź galicyjska). Tā bija visasākā zemnieku antifeodālā uzstāšanās XIX gadsimtā,} virzīta galvenokārt pret poļu muižniekiem, daļēji arī inteliģenci: tika izlaupītas muižas, nogalināti muižnieki, nereti arī vienkārši nepazīstami iebraucēji, kurus pūlis uzskatīja par sacēlušos aģitatoriem. Pat poļu tautības zemnieki Galīcijas rietumdaļā atbalstīja nevis savus panus~--- muižniekus, bet austriešu administrāciju. Zemnieki izdemolēja vairāk nekā 400~muižu un nogalināja vairāk nekā tūkstoti (pēc dažiem datiem līdz diviem tūkstošiem) poļu zemes īpašnieku un viņu ģimenes locekļu. Šis zemnieku trieciens poļu neatkarības aizstāvjiem bija ļoti sāpīgs kā politiski, tā morāli. Aprīlī sacelšanos apspieda austriešu karaspēks. Sacēlušos zemnieku vadītājs J.~Scela tika internēts un izsūtīts uz Bukovinu. Poļu XIX gadsimta daiļliteratūrā viņš iegājis kā antivaronis. Piemēram, S.~Vispjaņska populārajā drāmā ``\pltxti{Vesele}'' (``Kāzas'') viņš tēlots kā asinīm aptraipīts gars. Turpretī XX gadsimtā Polijas Tautas republikas laikā viņš tika vērtēts kā zemnieku sacelšanās pret apspiedējiem simbols.

Pēc sacelšanās apspiešanas Krievija, Prūsija un Austrija 1846.~gada 6.~novembrī parakstīja līgumu par Krakovas kā brīvpilsētas likvidāciju. Tā pēdējā formāli neatkarīgā poļu teritorija tika pievienota Austrijai. Austrijas Habsburgu dinastijas pārstāvji turpmāk saviem tituliem pievienoja vēl vienu~--- ``Krakovas hercogs'' (``\detxti{Herzog von Krakau}''). Pēc šiem notikumiem Krakovas poļu šļahtiču un inteliģences attieksme pret austriešu administrāciju vēl vairāk pasliktinājās. Nacionālās un sociālās pretrunas nebija likvidētas.

Krakovas sacelšanās neveiksme ietekmēja Galīcijas poļu izturēšanos turpmākajos notikumos. Bija stipri novājināts nacionālās kustības revolucionārais spārns. Galīcijas iedzīvotāji salīdzinoši vāji piedalījās 1848.~gada revolūcijā, jeb ``tautu pavasarī'', kur pārējās Eiropas tautas blakus sociālajiem izvirzīja arī nacionālos mērķus. Demokrāti, kuri iestājās par zemnieku atbrīvošanu, politiskajām brīvībām, Galīcijā bija ļoti vāji. Tomēr ziņas par revolucionāriem notikumiem Vīnē 1848.~gada martā izsauca demonstrācijas Krakovā un Lembergā. Tā, 1848.~gada 19.~martā Lembergā notika masu manifestācija un 12~000 cilvēku (aptuveni 1/6 no visiem pilsētas iedzīvotājiem) parakstīja petīciju imperatoram Ferdinandam I, ar prasību piešķirt pilsoniskās un nacionālās tiesības, ieviest mācības poļu valodās, kā arī piešķirt tai oficiālu statusu, pilnībā atcelt klaušas zemniekiem. Taču imperators atteicās pat apspriest prasības par pārmaiņām Galīcijā. 5.~aprīlī Krakovā tika radīta Nacionālā komiteja (\pltxti{Komitet Narodowy}), bet Lembergā~--- Nacionālā padome (\pltxti{Rada Narodova}). Uz vietām radās komitejas, kas pakļāvās Nacionālajai padomei.

Galīcijas poļu liberāļi joprojām kā galveno izvirzīja neatkarīgas poļu valsts atjaunošanas lozungu, bet demokrātiski pārveidojumi viņiem palika otrajā vietā. Tomēr daļa poļu liberāļu, tai skaitā viens no poļu autoritatīvākajiem politiskajiem darbiniekiem emigrācijā, bijušais Krievijas ārlietu ministrs un imperatora Aleksandra I personīgais draugs, kurš aktīvi piedalījās 1830.~gada sacelšanās, tās laikā bija Pagaidu valdības (\pltxti{Rząd Tymczasowy), pēc tam Nacionālās valdības (Rząd Narodowy) priekšsēdētājs} kņazs Ā.~Čartorijskis griezās pie muižniecības ar aicinājumu atteikties no savām vecajām tiesībām uz zemi. Taču Galīcijas šļahtiči kopumā baidījās tā pazaudēt arī tiesības uz atlīdzību par to. Tai pat laikā sākās zemnieku nemieri, tie masveidā atteicās pildīt klaušas. To izmantoja austriešu ģenerālgubernators grāfs F.~Stadions, kurš, saņēmis imperatora piekrišanu, 22.~aprīlī paziņoja par klaušu atcelšanu, zemnieku apstrādātās zemes nodošanu to īpašumā ar vēlāku valsts atlīdzību agrākajiem tās īpašniekiem. Varas iestādes veica zemes sadali zemniekiem lai atņemtu poļu patriotiem popularitāti, ko viņi daļēji bija guvuši pateicoties sociāli-ekonomiskajām prasībām. Galīcijā no augšas ievestā agrārā reforma bija liberālāka nekā blakus esošajās Krievijai un Prūsijai piederošajās ukraiņu un poļu zemēs. Tas vēl vairāk nošķīra vietējos zemniekus no poļu nacionālās kustības, novājināja to. Galīcijā no 1848. līdz 1857.~gadam tika īstenota agrārā reforma, kura pārvērta zemniekus par viņu apstrādājamās zemes īpašniekiem, taču to rokās bija tikai nelieli zemes gabali, vairāk nekā 2/3 no tiem piederēja mazāk par 10 morgiem (5,7 ha). Zemnieki asi izjuta aramzemes, bet īpaši pļavu un mežu trūkumu. Tas tikai pieauga, dalot zemi mantojumā daudzbērnu ģimenēs.

1848.~gada 25.~aprīlī Krakovā sākās nekārtības, iedzīvotāju sadursmes ar austriešu karavīriem. Nākamajā dienā austrieši apšaudīja pilsētu no Vavelas (\pltxti{Wawel}~--- kalns Krakovā pie Vislas ar monumentālu celtņu ansambli, tai skaitā karaļa pili, celtu XIII--XIV gadsimtā, un XIV gadsimta gotisku katedrāli, kur apglabāti karaļi un vēlāk arī citas izcilas Polijas personības) un aplenca pilsētu ar karaspēku. Sacēlušies kapitulēja un Nacionālā komiteja pārtrauca savu darbību. Tas bija pirmais kontrrevolūcijas panākums Eiropā.

Taču Lembergā poļu nacionālā kustība turpinājās. Tur tika radīta 20~000 liela poļu Nacionālā gvarde. Apstākļos, kad visā Habsburgu monarhijā notika revolūcija, tas deva cerības uz panākumiem. Taču Austrijas armija vispirms apspieda sacelšanos Prāgā, tad itāļu provincēs, oktobrī pašā Austrijā un varēja pievērsties arī Galīcijai. 2.~novembrī Nacionālā Padome Lembergā tika padzīta, pilsēta apšaudīta ar artilēriju, visā Galīcijā ievests kara stāvoklis, visas nacionālās organizācijas atlaistas. Revolūcija Galīcijā bija apspiesta. Vēl gan palika cerības uz revolūciju Ungārijā, kur ģenerāļa J.~Bema vadībā cīnījās arī tūkstošiem poļu, bet 1849.~gada maijā tajā pēc jaunā Austrijas imperatora Franča Jozefa aicinājuma iegāja Krievijas karaspēks ģenerālfeldmaršala I.~Paskeviča vadībā un triju mēnešu laikā sacelšanos apspieda. Ar to galīgi sabruka poļu cerības uz Eiropas palīdzību. Vairs nebija izredžu tuvākajā laikā izraisīt visas tautas bruņotu brīvības cīņu. Kara stāvoklis Galīcijā gan turpinājās vēl līdz 1854.~gadam.

Nākamo triecienu poļu nacionālajai kustībai Galīcijā deva \strong{rutēņu nacionālās kustības rašanās}. Kā raksta daži poļu vēsturnieki, Austrijā varas iestāžu mākslīgi atbalstītais konflikts starp šļahtu un zemniekiem paralizēja poļu nacionālo aktivitāti. Tiesa, var strīdēties par to, vai šis konflikts pie tā paša rezultāta nevestu arī bez varas iestāžu atbalsta. Pie tam jāuzsver, ka zemnieku kustībai bija ne tikai sociāla, bet arī nacionāla~--- rutēņu nokrāsa.

Austriešu ģenerālgubernators grāfs F.~Stadions, baidoties no tā, ka Galīcija varētu pārvērsties par ``poļu Pjemontu'' (itāļu \ittxti{Piemonte}, Pjemontas apgabals Itālijā XIX gadsimtā ieguva nozīmīgu lomu itāļu tautas nacionālās atbrīvošanās kustībā, kur ap Sardīnijas karalisti (itāļu \ittxti{Regno di Sardegna}), faktiski~--- Pjemontu 1859.--1860.~gadā notika Itālijas apvienošanās),~--- placdarmu, no kura varētu sākties neatkarīgas poļu valsts atjaunošana (Tik tiešām, vēlāk, 1918.~gadā, Galīcija, pateicoties saviem kultūras centriem, pēc Polijas neatkarības atjaunošanas kļuva par sava veida ``Polijas Pjemontu'', kura spēja sagādāt nākamajai Polijas valstij daudz politiķu, pārvaldes darbinieku, virsnieku, zinātnieku), meklēja pretsvaru poļiem un atrada tādu Galīcijas rutēņos.

Rutēņu intereses šai laikā vismaz daļēji sakrita ar Austrijas imperatora un valdības interesēm. 1847.~gadā Galīcijā iznāca 32 rutēņu izdevumi, bet 1848.~gadā~--- jau 156 (tiesa, tas bija rekords un netika pārspēts nākamo 30~gadu laikā). Ar viņa atbalstu 1848.~gadā tika radīta pirmā rutēņu politiskā organizācija~--- Galvenā Tautas padome (ukraiņu \uktxti{Головна Руска Рада}, krievu \rutxti{Галицийская рада}), kuru vadīja bīskaps G.~Jahimovičs. Tā laika rutēņu nacionālās kustības īpatnība bija naidīgums pret poļu liberāļiem, kuri izteica poļu muižnieku intereses, un uzsvērta lojalitāte Austrijas imperatoram. Pēc Galīcijas padomes iniciatīvas veidojās arī vietējās rutēņu komitejas no garīdzniecības un vēl tikai topošās inteliģences pārstāvjiem. Tika izvirzīta prasības paplašināt rutēņu tiesības Galīcijā un sadalīt provinci divās daļās: rietumu~--- poļu, un austrumu~--- rutēņu daļā. 1848.~gada jūnijā šo jautājumu apsprieda Slāvu tautu kongresā Prāgā, tur arī tika atzīta visu Galīcijas tautību līdztiesība.

Savu ietekmi atstāja arī Krievijas armijas dalība sacelšanās apspiešanā Galīcijai blakus esošajā Ungārijā, izmantojot arī Galīcijas teritoriju. Kā atzīmē daži mūsdienu krievu vēsturnieki, krievu armijas, kura runāja rutēņiem saprotamā valodā, parādīšanos Galīcijā vietējie zemnieki uztvēra ar sajūsmu un tā kļuva par impulsu krievu kultūras atdzimšanas sākumam. Tas esot izsaucis arī vēlmi tuvināties Krievijai, kas varēja apdraudēt Austrijas, bet vēlāk (no 1867.~gada) Austroungārijas vienotību. Tiesa, sajūsma par krievu armijas parādīšanos neatturēja pazīstamo rutēņu aktīvistu, juristu, rakstnieku Ā.~Dobrjaņski, kurš 1848.~gadā tika ievēlēts Ungārijas parlamentā, norobežoties no panslāvisma. Viņš paziņoja: ``Ungāru brīvība mums ir tuvāka nekā krievu patvaldība, tāpat kā mīkstais Ungārijas klimats ir tuvāks nekā Sibīrijas ziema.'' Vēlāk gan uz etnisko mazākumu asimilāciju vērstā Ungārijas varas iestāžu politika atstūma no Budapeštas rutēņu kustības vadītājus un pat pārvērtā viņus par rusofīliem. Arī minētais Ā.~Dobrjaņskis emigrēja uz Krieviju.

Uz rutēņu prasībām pēc lielākām tiesībām tika dota Galīcijas ģenerālgubernatora F.~Stadiona atbilde: ``Jūs varat rēķināties ar valdības atbalstu tikai tajā gadījumā, ja gribēsiet būt patstāvīga tauta un atteiksieties no apvienošanā ar tautu ārpus valsts, proti~--- Krievijā, ja gribēsiet būt rutēņi, bet ne krievi. Jums nekaitētu pieņemt jaunu nosaukumu, lai atšķirtos no krieviem, kuri dzīvo ārpus Austrijas''. Pēc dažu krievu vēsturnieku domām, austriešu varas iestādes, dažādi veicinot rutēņu etnisko atšķirību attīstību, atšķirīgu no krievu valodas rakstību, palīdzējušas rasties ukraiņu valodai, ukraiņu vēsturei, lai tikai atšķeltu vietējos iedzīvotājus no krievu tautas. Domājams, ka pretkrievisko nodomu vadītās Austrijas nacionālās politikas efektivitāte un arī rutēņu tendence apvienoties ar krieviem šeit tiek ievērojami pārspīlētas, taču ir nenoliedzams, ka Austrijas varas iestādes starpnacionālo attiecību laukā sekmīgi īstenoja saukli ``skaldi un valdi''. Tāpat nenoliedzams, ka 1848--1849.~gada revolucionārais pacēlums sekmēja Galīcijas rutēņu kultūras un izglītības darbības organizatorisku noformēšanos un pārtapšanu nacionāli--politiskā kustībā.

1861.~gadā tika izveidots vietējais Galīcijas Seims (poļu \pltxti{Sejm Krajowy}, vācu \detxti{Galizischer Landtag}) vietējo jautājumu risināšanai. Baidoties no 1863.~gadā Polijas karalistē sākušās sacelšanās izplatības arī Galīcijā, daži vietējie poļu muižnieki prasīja ieviest šeit aplenkuma stāvokli. 1864.~gadā Galīcijā arī tika ieviests kara stāvoklis, kurš pastāvēja līdz 1865.~gada maijam. Austrija pastiprināja robežas apsardzību ar Polijas karalisti, pār kuru virzījās palīdzība sacelšanās dalībniekiem.

No otras puses, vietējie poļu politiķi baidījās, ka Krievijas un Vācijas antipoliskā politika (Pēc 1863.--1864.~gada poļu sacelšanās sakāves cara valdība uzsāka represijas ne tikai pret poļu nacionālo kustību, bet arī pret ukraiņu valodas un kultūras propagandu) var tikt īstenota arī Galīcijā, ka varas iestādes var pret viņiem izspēlēt ``ukraiņu kārti''~--- sākušos ukraiņu nacionālās neatkarības kustību. Acīmredzot nemazsvarīga loma bija arī tam, ka Austrija, atšķirībā no Krievijas un Vācijas, bija katoliska valsts. Tāpēc viņi pakāpeniski atteicās nodalīt Galīciju no pārējās Austrijas/Austroungārijas.

Pēc 1863--1864.~gada sacelšanās sakāves Polijas karalistē, Austrijas sakāves karā ar Prūsiju (1866) un Austroungārijas izveides (1867) tās varas iestādes arī mainīja savu politiku, atteicās no ģermanizācijas. Austroungārijā praksē tika realizēta \strong{trīs ``kungu tautu''}~--- vāciešu (austriešu), ungāru (maģāru) un Galīcijas poļu~--- uzkundzēšanās pārējām tautām. Administratīvās struktūras bija izveidotas tā, ka vienā tās daļā~--- Cisleitānijā~--- vadošā tauta bija austrieši, bet otrā~--- Transleitānijā~--- ungāri. Galīcija ietilpa Cisleitānijā. Nacionālajā ziņā austrieši Bohēmijā (čehu \cstxti{Čechy}, vācu \detxti{Böhmen}~--- vēsturisks reģions mūsdienu Čehijas rietumdaļā) un Morāvijā (čehu \cstxti{Morava}, vācu \detxti{Mären}, vēsturisks novads Čehijas austrumdaļā) varēja turēt pakļautībā čehus, Istrijā (horvātu un slovēņu \sltxti{Istra}, itāļu \ittxti{Istria}~--- pussala Adrijas jūrā, Horvātijas ziemeļos)~--- slovēņus un horvātus, maģāri Ungārijā~--- slovākus, rumāņus un horvātus, bet poļi Galīcijā~--- rutēņus. Maģāru tīkojumus palielināt savu ietekmi novadā austrieši atvairīja, sadarbojoties ar konservatīvajiem poļu aristokrātiem.

Galīcijā vairāk kā citur Austroungārijā tika meklēti kompromisi ar cittautiešiem, tajā netika realizēta denacionalizācijas politika, kuru varētu salīdzināt ar ģermanizāciju Vācijai un rusifikāciju Krievijai piederošajās poļu teritorijās.

Arī Galīcijas poļu vidū auga vēlme pārskatīt iespējamos attīstības ceļus liberālisma gaismā. Buržuāziskie liberāļi, apzinoties poļu ekonomiskā potenciāla vājumu salīdzinājumā ar austriešu, prūšu un krievu, pirmkārt meklēja ceļus kā paaugstināt poļu kapitāla konkurences spēju. Pasludinot lozungu ``Ekonomika ir svarīgāka par politiku'', viņi ne tikai nosodīja romantisko ideju par poļu nacionālo vienreizīgumu, kura bija šļahtas revolucionaritātes pamatā, bet aktuālo cīņas problēmu pret nacionālo apspiešanu centās piepildīt ar citu saturu.

Jau 1864.~gada beigās Krakovas žurnālists L.~Povidajs rakstā ``\pltxti{Polacy i Indianie}'' (``Poļi un indiāņi''), kurš izsauca diskusijas visās poļu zemēs, brīdināja poļus neatkārtot Ziemeļamerikas indiāņu likteni, kuri esot noniecinājuši materiālo labklājību, lepojušies ar savu ``folkloras pārākumu'' un rezultātā gājuši bojā kolonizatoru spiediena rezultātā. (Protams, indiāņu nolemtības cēloņu attēlojums bija vairāk nekā apšaubāms.~--- V.Š.) Izmantojot prūšu politiķu izteikumus par poļiem kā ``Eiropas indiāņiem'', nolemtiem iznīcībai, L.~Povidajs minēja, pēc viņa viedokļa, briesmīgās prūšu realizētās ģermanizācijas sekas poļu zemēs, īpaši Silēzijā, kur rūpniecība un tirdzniecība bija vācu rokās. ``Svešu''~--- nepoļu kapitālu ieplūšanu rūpniecībā un tirdzniecībā viņš uzskatīja par nacionālā jautājuma asuma saglabāšanās cēloni visās poļu zemēs. Jautājuma atrisināšanas ceļu viņš saskatīja ``savu''~--- poļu uzņēmēju darbības aktivizēšanā. Taču šī aktivizācija nevarēja notikt bez politiskiem priekšnoteikumiem.

Viens no Galīcijas poļu liberālajiem līderiem F.~Smolka 1868.~gadā Galīcijas seimā iesniedza rezolūciju, kura aicināja uz cīņu par Galīcijas autonomiju un Austroungārijas federāciju, kur Galīcija būtu viena no sastāvdaļām. Tuvā franču-vācu kara, Austroungārijas dalības tajā un novājināšanās perspektīvas Galīcijas gan konservatīvajos, gan liberālajos poļu politiķos stiprināja cerības par sava mērķa panākšanu. Taču Francijas sakāve 1870.~gada karā un Austroungārijas neitralitāte tajā apraka visus federālistiskos plānus. Cīņa turpinājās par autonomiju.

Galīcija pakāpeniski ieguva vairāk politisku tiesību nekā citas Austroungārijas provinces. 60.--70.~gados Vīne ar daudziem likumdošanas aktiem faktiski ieviesa t.s. \strong{Galīcijas autonomiju}. 1873.~gadā, kad tika ieviestas tiešas (bez vietējā Seima starpniecības) vēlēšanas reihsrātā, Galīcijas poļu mantīgo slāņu prasību vairākums tika realizēts. Austrijas parlaments sastāvēja no divām palātām. Augšpalātā, kurā ietilpa imperatora ģimenes locekļi, aristokrātisko dzimtu pārstāvji un speciāli imperatora nozīmētas personas, Galīciju pārstāvēja trīs arhibīskapi un vēl 11 aristokrāti. Apakšpalātā, kuru veidoja vietējo seimu delegācijas, Galīcija ieguva 38~vietas.

Līdz pat XX gadsimta sākumam Galīcijas poļus politiski pārstāvēja lielie zemes īpašnieki, kuri bija sociāli tuvi austriešu aristokrātijai. Viņi bija galvenokārt konservatīvi orientēti. Pastāvēja vietējā pašvaldība (\pltxti{samorząd}). Poļi provincē ieguva saimnieciski un politiski vadošo lomu. Lembergā (Ļvovā) darbojās Galīcijas Seims. Tas tika vēlēts pēc kūriju principa, t.~i.~--- pēc profesionālā un mantas statusa, un tikai nelielai iedzīvotāju daļai bija vēlēšanu tiesības, priviliģētākie bija šļahtiči un bagātie pilsētnieki. No 150 Seima deputātiem 141 tika vēlēti. 9 vietas ``pēc amata'' saņēma bīskapi un universitāšu rektori. Seimā bija poļu vairākums. Ukraiņiem, kuri sastādīja pusi no Galīcijas iedzīvotājiem, bija tikai līdz ¼ vietu. Tikai XX gs. sākumā Austroungārijā realizēto reformu rezultātā Galīcijas Seimā paplašinājās arī trūcīgo iedzīvotāju pārstāvniecība. Seima kompetencē bija Galīcijas lauksaimniecības un mežsaimniecības jautājumi, novada budžets, labdarības, skolu un baznīcas lietas, Galīcijā izvietoto armijas daļu apgāde.

1866.~gadā Seims pieņēma lēmumu par poļu valodas ieviešanu Galīcijas skolās. 1867.~gadā Galīcijā tika izveidota novada skolu padome (\pltxti{Rada Szkolna Krajowa}), kura organizēja izglītības iestāžu darbu. 1869.~gadā poļu valoda Galīcijā ieguva oficiālu vietējās jeb zemes valodas (\detxti{Landessprache}) statusu. Poļu valodu atzina arī vietējās administratīvajās un tiesu iestādēs, tā tika iekļauta skolu programmās. Valodu, kādā notika mācības skolās, varēja izvēlēties. 1870.~gadā Krakovas (Jagelloņu) universitātē no vācu pārgāja uz poļu mācību valodu, 1871.~gadā tika paplašināta pasniegšana poļu valodā Lembergas universitātē, tehniskajās mācību iestādēs. Poļu mācību iestādes Austroungārijā kļuva populāras arī Pozenes un Polijas karalistes poļu vidū. Tā, 1909./1910.~mācību gadā starp 3~250 Krakovas universitātes studentiem 639 bija Krievijas pavalstnieki.

Lemberga kā Galīcijas administratīvais un politiskais centrs auga visai strauji. 1772.~gadā, kad pilsēta nonāca Austrijas sastāvā, tajā bija 20~000 iedzīvotāju. 1900.~gadā~--- jau 159~000, bet 1910.~gadā pat 200~000. 1861.~gadā dzelzceļš Lembergu saistīja ar Vīni. XIX gadsimta otrajā pusē Lemberga bija jau salīdzinoši moderna tirdzniecības pilsēta, ko kādreiz sauca arī par ``Austrumu Vīni''. Tajā darbojās Vācijas, Francijas, Lielbritānijas un Dānijas konsulāti. Ielas bija bruģētas, ierīkota kanalizācija, pacēlās modernas ēkas~--- Landtāga nams (1881), teātris (1900), dzelzceļa stacija (1904). 1894.~gadā tajā kā ceturtajā pilsētā Eiropā ierīkoja elektrisko tramvaju.

Uzplauka arī otra lielākā Galīcijas pilsēta~--- Krakova. 1869.~gadā Krakovā Vavelas (\pltxti{Wawel}) pils katedrālē tika restaurētas poļu karaļu kapenes, 1890.~gadā šeit svinīgi no Francijas tika pārapbedīts dzejnieks Ā.~Mickēvičs. 1905.~gadā, kad austriešu militāristi atstāja Vavelas pili, sākās tās atjaunošanas darbi. Krakovā tika sarīkoti nacionālo varoņu T.~Kostjuško un J.~Poņatovska atceres svētki, 1910.~gadā uzvaras Grīnvaldes kaujā 500~gadu jubilejai par godu par pianista un politiķa J.~Paderevska līdzekļiem tika uzstādīts karaļa Vladislava II Jagello piemineklis (1939.~gadā to iznīcināja vācu nacisti, atjaunots tika 1975.~gadā). Krakovā 1871.~gadā tika nodibināta Zinību Akadēmija (\pltxti{Akademia Umiejętności}), 1873.~gadā~--- Mākslas akadēmija. Tas viss pārvērta Krakovu par ``svētceļojumu'' vietu, īpaši poļiem no Pozenes provinces un Polijas karalistes.

Poļu ietekme Galīcijā bija nesalīdzināmi spēcīgāka nekā rutēņu. Poļu muižniecība ar grāfu A.~Goluhovski priekšgalā enerģiski cīnījās par to, lai visās dzīves jomās nodrošinātu poļu prioritāti. Poļu muižniecība un pilsonība atteicās atzīt ukraiņu valodas vienlīdzību ar poļu valodu Seimā. A.~Goluhovskis sākotnēji bija Galīcijas vietvalža padomnieks, pēc tam pats vairākkārt ieņēma šo posteni. 1871.~gadā Austroungārijas valdībā tika nodibināts īpašs Galīcijas lietu ministra amats, kurš līdz Austroungārijas sabrukumam parasti bija polis. Poļi varēja iegūt arī citus ministru posteņus. Ministru un Galīcijas vietvaldi imperators nozīmēja pēc saskaņošanas ar poļu deputātiem no t.s. ``poļu kolo'' (``\pltxti{Koło Polskie}''~--- poļu loks) Austrijas parlamentā. XIX gadsimta pēdējā trešdaļā Galīcijas poļu politiķi ieguva samērā lielu ietekmi Vīnē, kur 12 ministri un 2 ministru prezidenti nāca no Galīcijas. Var teikt, ka poļi Galīcijā baudīja priviliģētu stāvokli pretstatā ģermanizācijai Vācijā un carisma jūgam Krievijā. Lai aizsargātos no šiem ļaunumiem poļu virsslāņi atbalstīja Habsburgu monarhiju. Taču, kā norādījis angļu vēsturnieks A.~Dž.~P.~Teilors, tai pašā laikā viņi uzlūkoja Galīciju kā poļu valsts paraugu kādā attālā nākotnē un vēlējās saglabāt Habsburgu monarhiju tādā formā, lai Galīciju varētu atšķelt no tās jebkurā brīdī.

Galīcijas autonomija radīja poļiem vislabvēlīgākos apstākļus kopš tiem laikiem, kad eksistēja konstitucionālā Polijas-Lietuvas valsts. Austroungārijas valdošie slāņi faktiski tā atlīdzināja poļu muižniecībai par atbalstu Habsburgu monarhijai un uzstāšanos pret citu nacionālo grupu autonomijas prasībām. Poļi līdz pat impērijas sabrukumam palika pats uzticamākais un nelokāmākais tās elements. Šādas politikas būtību saskatīja arī padomju diktators un nacionālo lietu speciālists J.~Staļins. Viņš rakstīja: ``Austrieši pietuvināja sev poļus, deva tiem privilēģijas, lai poļi palīdzētu austriešiem stiprināt savas pozīcijas Polijā, un par to deva poļiem iespēju žņaugt Galīciju. Tā ir īpaša, tīri austriešu sistēma~--- izcelt dažas tautības un dot tām privilēģijas, lai pēc tam tiktu galā ar pārējām''.

Tādejādi Galīcija kļuva par vienīgo bijušās Polijas daļu, kas ieguva daļēju politisku un kulturālu brīvību, kā rezultātā Lemberga un Krakova kļuva par galvenajiem poļu politiskajiem un kultūras centriem.

Toties Galīcijas poļu zemnieku vidū nacionālās pašapziņas veidošanās norisa īpaši grūti. Autonomajā Galīcijā vara atradās poļu muižnieku rokās, kuri gan gribētu iegūt zemnieku atbalstu nacionālajā jautājumā, taču tikai ne ar politisku vai ekonomisku piekāpšanos. Konservatīvo muižnieku un buržuāzijas elementu nostiprināšanās pie varas noveda pie tā, ka Galīcijas lauki palika ``sprādziena bīstami'' kā sociālajā, tā nacionālajā aspektā. Ekonomiskā stagnācija Galīcijā veicināja sociālo attiecību arhaisku formu saglabāšanos, kas savukārt vājināja visu vietējo poļu tautības iedzīvotāju kopību. Tiesa, daļa proletarizēto vai pusproletarizēto zemnieku meklēja darbu Polijas karalistē, Pozenes lielhercogistē, arī Silēzijā~--- visur, kur attīstījās rūpniecības uzņēmumi, tā veidojot vispolijas darbaspēka tirgu, un līdz ar to veicinot dažādo poļu apgabalu iedzīvotāju nacionālo kopību.

Lembergā attīstījās arī ukraiņu kultūra, daudzās jomās pilsēta bija paraugs pārējām ukraiņu apdzīvotajām teritorijām. Rutēņu kustībā konkurēja divi strāvojumi. To piekritējus varētu saukt par ``ukrainofīliem'' un ``maskavafīliem''. Pirmais bija lojāls Habsburgu monarhijai, otrs orientējās uz Krieviju. 70.--80.~gados rutēņu aktīvisti, tā saucamie ``tautībnieki'' (\rutxti{народовцы}) arvien vairāk orientējās uz kopīgas ukraiņu nācijas attīstību. Noraidot kompromisu ar Galīcijas poļiem, viņi arī nevarēja pieņemt Krievijas orientācijas piekritēju uzskatus, kuri rutēņus uzskatīja par krievu tautas sastāvdaļu. Kā raksta angļu vēsturnieks N.~Deiviss, XIX~gs. otrajā pusē rutēņi sāka sevi saukt par ukraiņiem, lai izvairītos no maldinošā un aizskarošā apzīmējuma ``mazkrievs''. ``Ukrainis'' bija vienkārši politiski apzinīgs rutēnis. Ja nomaina vārdu ``politiski'' ar ``nacionāli'', tad konstatējumam var piekrist.

XX gadsimta sākumā vēsturnieks S.~Tomašivskis uzsvēra, ka Galīcijā pēc 1866.~gada poļi, palikdami nacionāls mazākums Austroungārijas valstī, bet iegūdami gandrīz neierobežotu varu pār apgabala vēsturiskajiem pamatiedzīvotājiem~--- ukraiņiem, kļuva par šķērsli ukraiņu nacionālajai kustībai, kas noveda pie pastāvīgām cīņām abu tautu starpā. Krievijas propaganda, vērsta uz vietējo rutēņu piesaistīšanu Krievijas ietekmei, lika poļiem, izvēloties starp rusofīlo un ukrainisko virzienu rutēņu attīstībā, atbalstīt otro. Taču ārpus valodas jautājuma poļi nevēlējās atbalstīt ukrainiskos elementus. Pēc nosauktā vēsturnieka domām, ukraiņu vairākums Galīcijā pirms Pirmā pasaules kara ilgstošas cīņas rezultātā vietējā pārvaldē bija ieguvuši ietekmi, ko salīdzinājumā ar poļu ietekmi izteica proporcija 1:4, taču, piemēram, ukraiņu vidusskolu bija tikai 1/10 no poļu vidusskolu skaita. Ukraiņu valodas tiesības sabiedriskās dzīves jomā joprojām bija ierobežotas, tā tika vienīgi pieciesta. Austroungārijas varas iestādes pretojās ukraiņu inteliģences centieniem vecā vietējo iedzīvotāju nosaukuma ``rutēņi'' nomaiņai ar ``ukraiņi''. Veco terminu mākslīgi uzturēja arī skolas. Savukārt saprotams, ka Krievija Galīcijas polonizāciju uzņēma kā visai netīkamu, vēl jo vairāk tāpēc, ka tā sekoja īsi pēc Polijas karalistes autonomijas likvidēšanas.

Neatkarīgas ukraiņu valsts radīšanas ideja apdraudēja kā Romanovu, tā Habsburgu impēriju vienotību, taču otrajai tā bija mazāk bīstama, jo tās sastāvā bija daudz mazāk ukraiņu apdzīvotu zemju nekā Krievijā. Galīcijā konflikts starp rutēņiem un poļiem te saasinājās, te atkal pieklusa. Atbalstot te vienu, te otru pusi, Austroungārijas valdība radīja provincē zināmu līdzsvaru. Rutēņu bērni varēja sākumskolās mācīties dzimtajā valodā. Šeit tika izdota literatūra ukraiņu valodā (Krievijā tā bija ierobežota). Lembergā strādāja vairāki ievērojami ukraiņu sabiedriskie darbinieki. Piemēram, šurp uz laiku pārvācās vēsturnieks, folklorists, publicists M.~Dragomanovs. Lembergas universitātē ilgstoši strādāja viens no ukraiņu nacionālās kustības līderiem vēsturnieks, profesors M.~Gruševskis. Viņš apgalvoja, ka Galīcija ir ``ukraiņu tautas avangards, kurš jau sen ir apdzinis nabago Krievijas Ukrainu'', ka ``līdz šim Galīcija gāja, bet Ukraina stāvēja vai sekoja Galīcijai''.

Vairāk nekā 10\% iedzīvotāju Galīcijā sastādīja \strong{ebreji}. Viņu tirgotāji un amatnieki veidoja vidusslāni starp lielajiem zemes īpašniekiem (galvenokārt~--- poļiem) un nabagajiem zemniekiem (galvenokārt rutēņiem). Ebreji šeit, tāpat kā citās Polijas daļās, pelnīja iztiku galvenokārt divējādi. Pirmkārt, viņu darbojās kā starpnieki starp muižniekiem~--- zemes īpašniekiem un zemniekiem, pildīja muižu pārvaldnieku, krodzinieku uzdevumus, kas krīžu periodos izraisīja pret viņiem dusmas no abām pusēm. Otrkārt, mazpilsētiņās viņi centās nodrošināt savu eksistenci ar sīktirdzniecību un amatniecību. Ebreju vairākums Galīcijā dzīvoja pilsētās vai arī savos miestiņos (jidišā \pltxti{Štetl}). 1880.~gadā ar vairum-, mazum- un ārējo tirdzniecību nodarbināto tirgotāju vidū ebreji sastādīja 84,5\%.

No poļu nacionālisma viedokļa ebreji bija jāizstumj no savām saimnieciskajām un sabiedriskajām pozīcijām. Saimniecībai bija jākļūst ``poliskai'' un jāpanāk attīstītākās valsts daļas industrializācijas gaitā, ko Galīcija kā Austrijas un Austroungārijas \latxti{de facto} lauksaimnieciska kolonija līdz šim bija ``nogulējusi''. Poļu muižniecība pati dibināja saimnieciskas sabiedrības (\pltxti{kólka rolnicze}) un, lai izstumtu no tautsaimniecības ebreju tirgotājus, manufaktūru īpašniekus un amatniekus, atbalstīja citu kārtu poļus, kuri vēlējās nodibināt savus uzņēmumus. Varas iestādes sistemātiski darbojās pret ebreju interesēm, atbalstīja antisemītisku aģitāciju. Katoļu baznīca pieļāva, ka atdzima vecas leģendas par it kā ebreju veiktām rituālajām slepkavībām. Jau no 1871.~gada skanēja aicinājumi boikotēt ebrejus, bet 90.~gados jau vairojās vardarbīgi uzbrukumi ebrejiem. Tas viss radīja grautiņiem labvēlīgu noskaņojumu sabiedrībā un virzīja ebrejus uz domām par emigrēšanas nepieciešamību. Pēc Polijas sadalīšanas ebrejiem pavērās iespēja migrēt uz līdz tam viņiem neaizsniedzamiem apgabaliem. Galīcijas ebreji, kas kļuva par Austrijas pavalstniekiem, sāka pārcelties uz citām Habsburgu monarhijas provincēm, bet vēlāk arī tās galvaspilsētu Vīni. Tomēr uz vietas palikušie Galīcijas ebreji saglabāja savu nacionālo identitāti.

\strong{Politisko partiju veidošanos Galīcijā} traucēja tās nošķirtība no pārējiem poļu apdzīvotajiem apgabaliem, tomēr līdz XX gs. sākumam arī šeit radās poļu politiskās partijas: sociālistiskā un zemnieku.

Jau 80.~gados šeit pastāvēja sociāldemokrātiskas grupas. 1890.~gadā no dažiem strādnieku pulciņiem izveidojās Galīcijas strādnieku partija, (``\pltxti{Galicyjska Partia Robotnicza}''). 1892.~gadā tā pārveidojās par Galīcijas sociāldemokrātiskā partiju (\pltxti{Galicyjska Partia Socjaldemokratyczna}, arī \pltxti{Socjaldemokratyczną Partią Galicji}). Partija formāli darbojās kā vienotās Austrijas sociāldemokrātiskās partijas sastāvdaļa.

1896.~gadā strādnieku un zemnieku kustības pacēluma apstākļos Austroungārijas valdība nolēma ieviest valstī jaunu vēlēšanu sistēmu. Jau esošajām četrām pēc mantas cenza sadalītajām vēlēšanu kūrijām tika pievienota piektā, uz vispārējo vēlēšanu tiesību pamata dibināta. Tajā ievēlēja 72~deputātus, 15~vietas piešķirot Galīcijai. Taču tas nenozīmēja tālejošu demokratizāciju, jo 1~deputātu pirmajā (muižnieku) kūrijā ievēlēja 63~vēlētāji, bet piektajā~--- ap 70~tūkstošu vēlētāju. Tomēr rezultātā poļu sociāldemokrātu pārstāvji I.~Dašinskis un J.~Kozakevičs pirmo reizi tika ievēlēti Austrijas parlamentā.

1897.~gadā Austrijas sociāldemokrātiskās partijas kongress pārveidoja partiju par sešu nacionālo partiju federāciju. Katra no tām tika veidota nevis pēc teritoriālā, bet nacionālā principa. Daudzi marksisti uzskatīja to par strādnieku šķelšanu. 1897.~gadā Galīcijas strādnieku partija arī pieņēma Galīcijas un Cešinas Silēzijas poļu sociāldemokrātiskās partijas (\pltxti{Polska Partia Socjalno-Demokratyczna Galicji i Śląska Cieszyńskiego}) nosaukumu. Turpmāk partija panāca vairāku savu pārstāvju ievēlēšanu Galīcijas Seimā un Austrijas parlamentā. Ievērojamākais partijas līderis bija minētais I.~Dašinskis.

Attīstījās \strong{zemnieku kustība}. Kapitālistiskās ražošanas attīstības lauksaimniecībā apstākļos sīko un vidējo zemes īpašnieku masas īpaši sāpīgi izjuta tradicionālās dzīves straujās pārvērtības, sava sociālā stāvokļa nestabilitāti, savas nākotnes nedrošību. Tāpat kā feodālisma laikmetā zemniecība palika pats lielākais iedzīvotāju slānis, kurš cieta no zemes trūkuma, smagas ekspluatācijas, bieži pirmskapitālistiskās formās. Krasā sociālā nevienlīdzība ar citiem slāņiem noveda pie tā, ka neraugoties uz citu partiju centieniem piesaistīt zemniecību sev sabiedrības politiskās organizācijas sistēmā radās īpašs virziens. Tas uzstājās pret feodālisma paliekām, par labvēlīgu nosacījumu radīšanu lauksaimniecības attīstībai, sociālo problēmu atrisināšanu par labu lauku sīkražotājiem, par zemnieku pielīdzināšanu tiesībās citiem slāņiem. Pati dzīve virzīja zemniecību uz apvienošanos organizācijās lai aizstāvētu savas ekonomiskās un politiskās intereses. Ar zemniecību saistītā inteliģence savukārt uzsāka kustības ideoloģisko pamatu izstrādi.

1886.~gadā Pēterburgā inženiera izglītību ieguvušais poļu šļahtičs B.~Vislouhs publicēja rakstu sēriju, kurā viņš izteica viedokli, ka politiskai organizācijai, kas vēlas panākt Polijas neatkarību, jābalstās uz darba cilvēkiem, pirmkārt~--- zemniekiem, kuri sastāda nācijas vairākumu. Nākotnes neatkarīgo Poliju viņš saskatīja etniskajās robežās, atzīstot ukraiņu, baltkrievu un lietuviešu tautu tiesības uz valstiskumu, izteicās par poļu un ukraiņu zemniecības sadarbību cīņā pret muižniekiem un veidojošos buržuāziju. B.~Vislouha programma bija virzīta uz feodālo palieku likvidāciju un par mērķi stādīja zemniecības politisku aktivizēšanu.

1888.~gadā Galīcijā izveidojās zemnieku grupa, kura darbojās ``kristīgā sociālisma'' garā. Jau radikālākās pozīcijās stāvēja no 1889.~gada Lembergā B.~Vislouha izdotā avīze ``\pltxti{Przyjaciel ludu}'', (``Tautas draugs''). Ap avīzi grupējās cilvēki, kuri centās piešķirt zemnieku kustībai visas tautas kustības raksturu. Tā ilgstoši varēja izplatīties tikai Galīcijā. (Īpaši Polijas karalistē zemnieku organizācijām praktiski bija liegtas iespējas darboties legāli.) Galīcijas Seima vēlēšanās 1889.~gada jūlijā 18~povjatos (poļu \pltxti{powiat}~--- vidēja lieluma administratīvi teritoriāla vienība, līdzīga apriņķim) radās zemnieku vēlēšanu komitejas, kuras ar inteliģences demokrātisko pārstāvju atbalstu panāca četru zemnieku deputātu ievēlēšanu, kuru darbībai tajā gan nebija reālas nozīmes.

1895.~gadā Rešuvā (poļu \pltxti{Rzeszow}, krievu \rutxti{Жешув}, ukraiņu \uktxti{Ряшiв}, vācu \detxti{Reichshof}) Galīcijā tika izveidota zemnieku partija ``\pltxti{Stronnictwo Ludowe}'' (No poļu \pltxti{lud}~--- tauta, zemniecība,~--- Zemnieku jeb Tautas partija). Tās piekritējus tautā sauca par \strong{``ļudoviešiem''}. (Tiesa, lietojot vārdu ``partija'', jāsaprot zināma īpatnība. Programmas dokumentos, zinātniskajā literatūrā, publicistikā šis vārds tiek lietots. Taču Polijā izplatītais nosaukums ``\pltxti{stronnictwo}'' burtiski nozīmē ``piekritēju savienība''. Tieši šis apzīmējums precīzāk izsaka zemnieku tā laika organizētības pakāpi, ja to, piemēram, salīdzina ar revolucionārajām strādnieku partijām.) Par tās priekšsēdētāju kļuva advokāts K.~Ļevakovskis, kurš gan pēc diviem gadiem atteicās no priekšsēdētāja amata. Partijas vadībā darbojās arī B.~Vislouhs, H.~Revakovičs, J.~Stapiņskis u.c.

Ļudovieši pārstāvēja nacionāli orientētos apzinīgākos lauku iedzīvotāju, galvenokārt sīkzemnieku, slāņus. Partijas galvenās prasības bija neatkarīgas Polijas valsts atjaunošana; Galīcijas autonomijas panākšana; vēlēšanu likumu maiņa parlamentā Vīnē un Galīcijas Seimā, lai radītu demokrātiskāku parlamentāro sistēmu; vienlīdzīgas tiesības zemniecībai, vispārējo vēlēšanu tiesību ieviešana un sabiedriskās dzīves demokratizācija; plaša zemniecības pārstāvniecība likumdevējās iestādēs; agrārās reformas īstenošana un zemnieku saimniecību skaita palielināšana. Ļudovieši centās likvidēt muižnieku privilēģijas ekonomiskajā, politiskajā un sabiedriskajā dzīvē. Seima vēlēšanās apgabalos ar jauktu~--- poļu un rutēņu iedzīvotāju sastāvu partija uzstājās pret šļahtas kandidātiem. 1895.~gada Seima vēlēšanās tajā ievēlēja 9~ļudoviešus, no tiem 7~zemniekus. Tā kā konservatoriem, muižniecībai bija cieši sakari ar baznīcu, ļudovieši drīz nonāca konfliktā ar to. Baznīcai piederēja ievērojamas zemes platības. Tā ļudoviešu parādīšanos zemnieku vidē uzskatīja par apdraudējumu savām monopoltiesībām uz ticīgo poļu zemnieku prātu kontroli, kā ksendza tiesību ierobežošanu. No kanceles pār ļudoviešiem sāka līt apvainojumi, taču viņi, aizstāvot tiesības uz apziņas brīvību, pasvītroja savu dziļo cieņu pret baznīcu, pilnīgu atbalstu kristīgajai morālei, taču noraidīja baznīcas iejaukšanos valsts politiskajā dzīvē. Kustības kreisais spārns prasīja baznīcas atdalīšanu no valsts. Praktiskajā darbībā ļudovieši centās nejaukt politiku ar reliģiju.

Ļudoviešu idejiskās pozīcijas pamatā bija Eiropas pilsoniskās un sīkpilsoniskajā kustībā XIX~gadsimtā visai plaši izplatītā \strong{agrārisma teorija}. Tā radās kā atbildes meklējumi uz pārdzīvojamajām tirgus saimniecības grūtībām. Teorijas pamatā atradās zemniecības garīgo, morālo, kulturālo un sociālo īpatnību un tās darbības atzīšana par unikālu un nepārejošu vērtību visas cilvēces attīstībā. Rūpniecības preču pārprodukcija, rūpnīcu slēgšana, masu bezdarbs daļā ekonomistu radīja uzskatu, ka lielrūpniecības attīstības iespējas ir izsmeltas, toties lauksaimniecība, ar valsts protekcionisma (no latīņu \latxti{protectio}~--- aizbildniecība; iekšējā tirgus aizsardzības sistēma pret ārzemju konkurenci ar augstu muitas nodevu u.c. līdzekļu palīdzību) politikas atbalstu var cīnīties ar ekonomiskajiem satricinājumiem.

Nometinot ``liekos'' strādniekus laukos, pārvēršot viņu par sīkzemniekiem, kuri apstrādā savu lauciņu, pavērās cerība ``nolīdzināt'' krīžu gaitu, novērst sociālo satricinājumu cēloņus. Agrārisma teorētiķi lika priekšā veicināt lauksaimniecības attīstību ar lētu kredītu piešķiršanu, nodokļu samazināšanu, ``cenu šķēru'' uz rūpniecības un lauksaimniecības preču likvidēšanu, ``kooperatīvu asociāciju'' izveidi. Šos agrārisma postulātus piemērojot dažādu valstu apstākļiem, zemnieku kustībā tika izstrādātas dažādas pastāvošās iekārtas uzlabošanas koncepcijas sīkīpašnieku un sīkražotāju labā. Rekonstruējot pastāvošo agrāro iekārtu un dažādiem līdzekļiem attīstot lauksaimniecību, zemnieku partijām parlamentārā ceļa iekarojot varu, tika plānota pakāpeniska, bez satricinājumiem pāreja no pastāvošā kapitālisma uz ``zemnieku kooperatīvo republiku''.

Vēlreiz jāuzsver, ka ka ārpus Austroungārijas citās poļu zemēs zemnieku kustība XIX gadsimtā nespēja iegūt stingrus organizatoriskus pamatus. Polijas karalistes un Pozenes novada teritorijā ļudoviešu kustība plašāk izvērsās tikai pēc Polijas neatkarības atgūšanas 1918.~gadā.

\subsection{Poznaņa Prūsijas / Vācijas sastāvā}

Pēc Vīnes kongresa radītā Pozenes province (vācu \detxti{Provinz Posen}, poļu \pltxti{Prowincja Poznańska),} aptvēra tās Polijas daļas, ko Prūsija bija jau saņēmusi 2. un 3. Žečpospolitas dalīšanas reizē. Šīs agrākās Polijas zemes kādreiz veidoja t.s. ``Lielās Polijas'' (\pltxti{Wielkopolska}) kodolu, tur X gadsimtā Gņezno (poļu~--- \pltxti{Gniezno}, vācu~--- \detxti{Gnesen}) bija poļu galvaspilsēta, līdz XIV gadsimtam tur kronēja karaļus, no XI~gadsimta tur atradās arhibīskapa sēdeklis. (Zemes, ko Prūsija bija sev pievienojusi jau agrāk~--- Rietumprūsija, Pomerānija u.c., bija tieši iekļautas Prūsijas sastāvā.)

Pozenes zeme, kas kā tilts savienoja Austrumprūsiju ar Pomerāniju un Silēziju, faktiski tika iekļauta Prūsijā kā autonoma province un līdz 1866.~gadam neietilpa Vācu savienībā. Tā sākotnēji gan tika dēvēta par Pozenes lielhercogisti (vācu \detxti{Großherzogtum Posen,} poļu \pltxti{Wielkie Księstwo Poznańskie}), taču faktiski province atradās Prūsijas karalistes sastāvā. Prūsijas karalim Pozenes lielhercoga tituls bija vajadzīgs vienīgi tāpēc, lai to pretstatītu Polijas karaļa titulam, kuru pieņēma Krievijas imperators. Ar 1848.~gada revolūciju tika atmests arī lielhercogistes nosaukums. Tāpat kā Rietumprūsijā, arī Pozenes hercogistē, īpaši tās rietumu daļā, vēl no Teitoņu ordeņa laikiem saglabājās liels skaits vācu iedzīvotāju. Parasti Rietumprūsiju un Pozenes zemi, neraugoties uz to atšķirīgo statusu, minēja reizē, tāpēc XIX gadsimta otrajā pusē tās sāka saukt par ``\detxti{Ostmark}'' (Austrumu apgabals).

Pozenes lielhercogistē dzīvoja ap 0,9 miljoniem cilvēku, pēc citiem datiem~--- ap 521~000 katolisku poļu, ap 218~000 lielākoties protestantisku vāciešu (pēc konfesionālās piederības Pozenes vāciešu vidū dominēja luterāņi, kaut pastāvēja arī maza grupiņa t.s. reformātu-kalvinistu), un ap 50~000 ebreju. Tajos novados, kas robežojās ar Krievijai piederošo Polijas karalisti, poļu daļa iedzīvotāju vidū sniedzās līdz 85\%, bet rietumu novados bija vācu vairākums. Kopumā provincē varēja runāt par attiecību 1:2 par labu poļiem.

XIX~gadsimta sākumā nevar runāt par kādu nacionālu naidu starp vietējiem vāciešiem un poļiem. Abas tautības nebija skaidri nošķirtas, eksistēja daudz polonizētu vāciešu un ģermanizētu poļu. (Ebreji kopumā labprātāk sliecās uz vācu nekā poļu ietekmes pieņemšanu.) Par savstarpējiem kontaktiem liecināja jauktās laulības. XIX~gadsimta pirmajā pusē tās sasniedza 20\% no visām. Jāatzīmē gan, ka poļi biežākus kontaktus uzturēja ar vācu katoļiem nekā ar vācu protestantiem un otrādi~--- vācieši labprātāk kontaktējās ar poļu protestantiem nekā katoļiem. Savstarpējā tuvināšanās veda pie asimilācijas, sākotnēji vairāk polonizējās vācieši, pat prūšu ierēdņi.

Sākotnēji Pozenē nacionālo attiecību laukā valdīja ``klusie gadi'', jo Prūsija 1815.~gadā bija ieguvusi arī Vestfāli un Reinzemi, kur arī bija jāorganizē jauna pārvalde, tāpēc nevarēja visu uzmanību veltīt poļu apgabaliem. Bez tam Prūsijai bija jārēķinās arī ar pārējo Poliju sadalījušo valstu realizēto politiku, īpaši Krievijas politiku Pozenes novadam blakus esošajā Polijas karalistē. Tur Aleksandrs I bija atļāvis nosacīti liberālas Konstitūcijas pieņemšanu (par to sīkāk stāstīts nākamajā apakšnodaļā), karalistes pavalstnieki baudīja zināmu autonomiju, tikai ārlietas bija stingri Krievijas imperatora noteiktas. Prūsija neuzdrošinājās Pozenes poļu iedzīvotājiem radīt sliktākus apstākļus nekā Polijas karalistē, jo tad vajadzētu rēķināties ar viņu simpātijām Krievijai.

Prūsijas toreizējais karalis Fridrihs Vilhelms III solīja jaunajiem pavalstniekiem sargāt ne tikai viņu reliģiju, bet arī valodu, nodrošināt vietējiem iedzīvotājiem pieejamību amatiem pārvaldē. Viņa uzsaukums poļiem skanēja: ``Jūs tiksiet iekļauti manā monarhijā tā, lai nevajadzētu noliegt savu tautību''. Šādu attieksmi noteica arī tas, ka trūka prūšu ierēdņu, kuri zinātu arī poļu valodu. Pozenes lielhercogistē darbojās pašvaldība: landtāgs. Tajā 24 (vēlāk~--- 26) vietas piederēja muižniekiem, 16~--- pilsētām un 8~--- lauku kopienām. Landtāgā bija poļu vairākums un viņi to izmantoja, lai aizstāvētu savas tiesības pret vēlākajiem valdnieku un ierēdņu tīkojumiem veikt ģermanizāciju. Landtāga sēdēs tika izmantota gan vācu, gan poļu valoda. Tās abas izmantoja arī pārvaldē, tiesu un izglītības iestādēs. Prūsija pat atļāva izveidot Pozenes korpusu Prūsijas armijas sastāvā ar poļu virsniekiem un poļu komandu valodu.

1815.~gadā par Prūsijas karaļa vietvaldi Pozenē kļuva poļu un prūšu politiķis, zemes lielīpašnieks un komponists firsts (kņazs) A.~Radzivills. Par savu uzdevumu viņš uzskatīja samierināt poļus ar Prūsiju, kaut Pozenes poļi pirmkārt centās iegūt pašpārvaldi. Viņa karjera gan beidzās ar sacelšanos Polijas karalistē (1830), kur viens no vadītājiem bija viņa brālis Mihals. Prūsijas valdībā kā izglītības ministrs darbojās liberāls un tālredzīgs politiķis K.~Altenšteins, kurš 1822.~gada decembrī, atbildot uz Pozenes ierēdņu jautājumu par vācu valodas ieviešanu, sacīja: ``Protams, valdība ir ieinteresēta, lai poļu iedzīvotāji zinātu vācu valodu un varētu izmantot valsts iestāžu pakalpojumus; taču tas nebūt nenozīmē, ka poļus ir jāpārvāco \citespace{} Jo kā cilvēka, tā arī tautas audzināšanu var nostiprināt tikai dzimtā valoda: atņemt to \citespace{} un censties mākslīgi iepotēt svešu valodu, tāds ceļš ir nepareizs pat attiecībā pret atsevišķu cilvēku, nemaz nerunājot par veselu tautu, pat ja tai nebūtu tik bagātas, tik izstrādāts un nobeigtas valodas, kāda ir poļu [valoda] \citespace{} Vara, kura atzīst, ciena un aizstāv tautas reliģiju un tās valodu, divus lielākos tās svētumus, var būt pārliecināta, ka piesaistīs sev pavalstnieku sirdis. Taču ja vara ir vienaldzīga pret tiem, vai vēl vairāk~--- uzliek tiem savu roku, tā apvaino un uzbudina tautu un reizē ar to gatavo sev neuzticamus un saniknotus pavalstniekus''. Diemžēl šo dziļo patiesību dažādu tautību šovinisti nav sapratuši līdz pat mūsdienām.

Neatbilst patiesībai J.~Staļina valdīšanas periodā tapušajā padomju vēstures literatūrā propagandētais apgalvojums, ka visvairāk cieta tie Polijas iedzīvotāji, kuri pēc valsts sadalīšanas nonāca Prūsijas varā. Prūsijai piederošajos apgabalos ekonomika attīstījās sekmīgi, īpaši jau pirms Žečpospolitas sadales Prūsijas rokās nonākušajā Silēzijā, lauksaimniecība arī Pozenē attīstījās raitāk nekā Krievijai piederošajās poļu zemēs. Prūsijā poļu iespējas iegūt izglītību savā valodā XIX gadsimta pirmajā pusē bija krietni labākas kā Krievijas vai Austrijas teritorijās. Ja 1815.~gadā Pozenes provincē bija 164 (citi dati~--- 540) skolas, tad 1845.~gadā to jau bija 1~142. Sākotnēji tikai dažas vācu politiskās grupas iebilda pret samierināšanās kursu un vēlējās realizēt poļu asimilāciju pēc Rietumprūsijas parauga. Tās virsroku sāka gūt tikai pēc poļu sacelšanās 1830.~gadā, kura ietekmēja situāciju arī Pozenē.

Tikai turpmāk arvien uzskatāmāk parādījās vācu un poļu interešu pretstati. Sākotnējo diezgan plašo Pozenes iekšējo autonomiju Prūsija pakāpeniski sašaurināja. Jau no XIX gadsimta 20.~gadu sākuma hercogistes autonomās iestādes sāka darboties kā Prūsijas administrācijas turpinājums, amatos tajās nozīmēja tikai etniskos vāciešus. Dažus gadus pēc Vīnes kongresa tika precizēta provinces pārvaldes sistēma. 1819.~gadā tika radīts Pozenes pārvaldes iecirknis ar 17 apgabaliem (\detxti{Kreis}) un Brombergas (\pltxti{Bydgoszcz}) iecirknis ar 9 apgabaliem. Tāds iedalījums pastāvēja līdz pat 1887.~gadam, kad Prūsijas Ministru padomes priekšsēdētājs un Vācijas impērijas kanclers O.~f.~Bismarks reorganizēja provinces. Līdz Pirmajam pasaules karam Pozenes province bija sadalīta Pozenes iecirknī ar 28 apgabaliem un Brombergas iecirknī ar 14 apgabaliem. Gnezenes arhibīskapija tika pārveidota par Gnezenes-Pozenes arhibīskapiju, kur pats arhibīskaps rezidēja Pozenē. Luterāņu baznīcas priekšgalā atradās Konsistorija ar superintendantu.

Prūsijā jau 1807.~gadā bija oficiāli atcelta zemnieku personīgā dzimtbūtnieciskā atkarība, 1811.~gadā zemniekiem tika atļauts izpirkt savas feodālās saistības, samaksājot 25-kārtīgu to gada apjomu vai atsakoties no 1/3 līdz 1/2 savas apstrādājamās zemes. Taču vēl ilgi zemnieki neieguva līdztiesību, gadu desmitiem saglabājās izpirkuma maksājumi. 1823.~gadā šo kārtību attiecināja arī uz Pozenes lielhercogisti. Arī tur notika agrārā reforma, turīgie zemnieki varēja iegūt privātīpašumā daļu apstrādājamās zemes, ar atteikšanos no otras daļas izpērkot savas feodālās saistības~--- klaušas un nodevas. Izpirkšanas maksājumus noteica valdības ieceltas speciālas komisijas. Zemes cenas Pozenē bija zemākas nekā pašā Prūsijā, kur varas iestādes vairāk rēķinājās ar muižnieku interesēm. Tādejādi valdība starp zemniekiem un muižniekiem šeit it kā izveidoja zemniekus sargājošu sienu. Kaut reformas notika ļoti lēni, zemes izpirkšana ievilkās līdz 1865.~gadam, zemnieki zaudēja apmēram 1/3 viņu agrāk apstrādājamās zemes, tomēr radās zemnieku~--- zemes īpašnieku slānis, kurš jau XIX gadsimta beigās kļuva par ietekmīgu politisku faktoru. Jāpasvītro, ka poļu zemniecība Pozenē turpmāk ļoti auksti izturējās pret šļahtas sacelšanās mēģinājumiem. Labklājības pieaugums šai zemniecībai vēlāk, kad tā nonāca pie nacionālās pašapziņas, ļāva sekmīgi cīnīties pret ģermanizāciju, sevišķi ekonomiskajā un kultūras jomā. Līdz pat Pirmajam pasaules karam lauksaimniecības attīstībā Pozene bija priekšā citām poļu apdzīvotajām teritorijām Krievijas un Austrijas sastāvā.

Taču reformas bija sāpīgas nabadzīgajiem zemniekiem, kuri nespēja izpirkt savu zemi un pakāpeniski tika no tās padzīti. Gandrīz pilnībā izzuda zemnieku sīksaimniecības. Tiesa, šo padzīšanas procesu nevajadzētu vienkāršot kā izraidīšanu uz četrām debess pusēm. Bieži, zemniekam atstājot savu agrāko saimniecību, viņiem tika ierādīts cita dzīves vieta~--- savs ``stūris'' (\pltxti{kąt}), iedalīts zemes gabals dārzeņu un kartupeļu audzēšanai un radītas citas atkarības formas no muižnieka. Lai nepaliktu bez jumta virs galvas, šādi zemnieki pēc pirmā aicinājuma ieradās darbā muižā. Ja šie cilvēki mēģināja nodarboties ar amatniecību un tirdzniecību, viņi sastapās ar profesionālu amatnieku un tirgotāju konkurenci, kuru vairākums bija vācieši.

Feodālisma apstākļos gadsimtu gaitā zemnieki, neizturējuši ekspluatācijas slogu, bēga no sava muižnieka. Muižnieki turpretī visiem līdzekļiem centās viņus noturēt, atrast bēgušos. Valsts aparāts palīdzēja muižniekiem notvert un atgūt savus dzimtbūšanas atkarībā esošos zemniekus. XIX gadsimta pirmajā pusē situācijas radikāli mainījās: muižnieki padzina daļu nabadzīgāko zemnieku no zemes, bet pēdējie savukārt pretojās, centās saglabāt savu saimniecību. Ir skaidrs, ka muižnieki padzīt zemniekus un pievienot viņu saimniecības saviem laukiem varēja tikai tad, ja bez citiem nosacījumiem (pārpalikušo zemnieku darba intensitātes pieauguma, tehnikas attīstības) bija pārliecināti, ka nepieciešamības gadījumā atradīs pietiekami daudz brīva darbaspēka, kuru varēs algot. Tātad, neraugoties uz feodālisma pastāvēšanu, jau eksistēja, lai arī ierobežots, brīva algota darbaspēka tirgus. Tas bija radies pateicoties zemnieku bēgšanai un klīšanai darba meklējumos. Kad šis darbaspēka tirgus jau bija pietiekami liels, ``liekos'' zemniekus varēja arī padzīt no viņu apstrādājamās zemes, tā savukārt atkal paplašinot algotā darbaspēka rezerves. Kartupeļu ražošanas izplatība deva iespēju iegūt produktu, ar kuru varēja pabarot zemnieku ģimenes no ievērojami mazāka zemes gabala. Tā muižnieki, pat neatņemot zemniekam viņa saimniecību, bet tikai samazinot tās platību, varēja pievienot savai muižai jaunas zemes platības.

No agrārās reformas īstenošanas laika sākās poļu zemnieku proletarizēto grupu izceļošana darba meklējumos uz Rūras ogļu ieguves uzņēmumiem, Reinas-Vestfāles metalurģiskajām un Berlīnes mašīnbūves rūpnīcām. Daudzi devās arī uz strauju attīstības ceļu uzsākušās Augšsilēzijas rūpniecības uzņēmumiem. Tā nežēlīgiem līdzekļiem tika veidota racionāla saimniekošanas sistēma.

Rūpniecības attīstība pašā Pozenē sastapās ar lielām grūtībām, jo bija jākonkurē ar daudz attīstītāko vācu apgabalu rūpniecību. Pozenes province toreiz piederēja pie visvājāk attīstītajiem Prūsijas apgabaliem. Tomēr Prūsijas varā esošajās poļu zemēs dzelzceļu celtniecība norisa straujāk un to garums bija lielāks nekā pārējās Polijas daļās. Taču satiksme tika nodrošināta pirmkārt ar Vācijas teritorijām. Pozenei nebija tiešas dzelzceļa stigas ar Varšavu. 1820.~gados Pozenes vadmala vēl tika izvesta uz Polijas karalisti, taču 30.~gados eksports gandrīz pilnībā apsīka. Konkurence, muitas robeža ar Krieviju noveda pie tā, ka tika slēgtas vairākas manufaktūras. Prūsijas un pēc tam arī Vācijas varas iestādes nebija ieinteresētas rūpniecības attīstībā poļu apdzīvotajos apgabalos, nepretojās rūpnieku, tai skaitā arī vācu, pārbraukšanai uz Polijas karalisti. Ar to Pozenes ekonomika saņēma smagu triecienu, ieguva vienpusīgāku agrāru raksturu. Šūšanas izstrādājumu ražošanā mašīnas sāka izmantot tikai 1859., apavu ražošana~--- 1861.~gadā. Puse no amatniekiem nelietoja algotu darbu. Tomēr nav apstrīdams, ka Prūsijā, bet vēlāk Vācijas impērijas (1871--1918) Prūsijas pavalstī Pozenes novads, neraugoties uz minētajām grūtībām, iedzīvotāju nacionālajām pretrunām, piedzīvoja ievērojamu saimniecisko augšupeju, tajā attīstījās moderns satiksmes tīkls, tika veikta plaša sabiedriskā un privātā celtniecība, ar industrializāciju tika vairākkārt kāpināta novada saimnieciskā ražība.

Kā smags slogs uz poļu iedzīvotājiem gūlās nepieciešamība dot rekrūšus Prūsijas armijai. Prūsijā militāro dienestu pildīja zemnieki un tie pilsētnieki, kuri nodarbojās ar lauksaimniecību. Šī kārtība tika attiecināta arī uz Prūsijā nonākušajām poļu zemēm, kuras agrāk tādu nepazina. Poļu zemnieki nevēlējās dienēt viņiem svešas valsts armijā, kura pie tam bija slavena ar savu nūjas disciplīnu. Pēc pirmās Žečpospolitas Polijas dalīšanas Prūsijas sastāvā nonākušie bēga atpakaļ uz pārpalikušo poļu-lietuviešu valsti. Pēc tās galīgas sadales šīs iespējas zuda.

\asterism

Sociālā ziņā Pozenes vāciešu vidū bija maza muižnieku grupa, jau pamatīgs pilsonības, arī amatnieku un zemnieku, slānis. Visstraujāk auga vācu inteliģences slānis. Īpaši palielinājās vācu ierēdņu skaits. Poļu ierēdņu jau bija maz. 1848.~gadā Pozenē no 700 ierēdņiem tikai 30 bija poļu tautības. Augstskolās izglītoto vāciešu īpatsvars pieauga arī citās nozarēs. Kopā ar ebrejiem vācieši sastādīja advokātu vairākumu. Puse no ārstiem bija vācieši. Auga vācu skolotāju skaits. Protestantu garīdznieki (1870.~gadā to bija 210) gandrīz bez izņēmuma bija vācieši. Katoļu garīdznieku vidū vāciešu bija 10--15\%. Tiesa, līdz 1870.~gadam vācu arhitektu, mākslinieku, literātu un aktieru skaits bija neliels~--- daži desmiti, kas bija izskaidrojams ar niecīgo pieprasījumu pēc mākslas.

1832.~gadā 185 muižas (gandrīz 22\% no visām) atradās vācu rokās. Vācu muižnieku vairākuma īpašumā galvenokārt bija no 200 līdz 1~000 ha lielas muižas. Parasti 1 cilvēkam piederēja 1 muiža, kas atviegloja tās pārvaldi. Vācu muižnieki Pozenes provincē reti piederēja kādai vecai aristokrātiskai dzimtai, kaut bija arī tādas. Piemēram, kopš 1819.~gada 25~641 ha piederēja firstiem \detxti{Thurn und Taxis}. Vairākums vācu muižu īpašnieku nebija dižciltīgie, bet gan agrāk bijuši augsti ierēdņi provinces pārvaldē ar zināšanām lauksaimniecībā. Tāpēc viņu apsaimniekošanas rezultāti parasti bija labāki nekā viņu poļu kaimiņiem. Vācu muižnieki ierīkoja solīdas saimniecības ēkas, centās paaugstināt ražas ar intensīvu saimniekošanu. Viņi pielietoja modernus mēslošanas līdzekļus, ieviesa jaunas kultūras, piemēram, kartupeļus, veica mājlopu selekciju, investēja produkcijas pārstrādē. Protams, bija arī izņēmumi, kad muižas īpašnieki interesējās tikai par ātrāku ienākumu saņemšanu, vai, neraugoties uz valsts palīdzību, nespēja pielāgoties konjunktūrai. Tika veidotas vācu lauksaimniecības biedrības. Pirmā radās 1826.~gadā. 1852.~gadā tās jau apvienojās biedrību savienībā. Līdz 1848.~gadam pastāvēja gan arī kopējas poļu un vācu organizācijas, arī lauksaimniecības biedrības.

Vāciešu vidū Pozenes provincē pats lielākais bija zemnieku slānis. Daļa no viņiem bija šai apgabalā apmetusies vēl pirms Žečpospolitas pirmās dalīšanas~--- XVIII gs. pirmajā pusē. Jaunākā grupa ieceļoja pēc 1815.~gada. Pozenes provincē bija salīdzinoši lēti nopērkama zeme, tāpēc šurp labprāt devās vācu zemnieki. Valdība viņus atbalstīja finansiāli. Brombergas pārvaldes iecirknī vēl pirms 1830.~gada tika ierīkotas 13 vācu kolonijas. Tomēr visu Pozenes zemnieku vidū vāciešu skaits nebija visai liels, sasniedza apmēram 1/3. Viņu ekonomiskais stāvoklis bija krietni labāks nekā poļu zemniekiem. Vācu zemniekiem piederēja pārsvarā vidējas un lielas saimniecības. Tam par cēloni bija kā labvēlīgie pārcelšanās noteikumi, tā pašu vācu pārceļotāju augstākais saimniekošanas līmenis. Ar valsts palīdzību tika veidotas izglītības biedrības, krājkases.

Pozenes provinces rietumu un ziemeļu pierobežas pilsētās dzīvoja ne tikai skaitliski, bet arī ekonomiski spēcīgs vācu pilsonības slānis. Pašā Pozenē vairāk nekā 2/3 pilsētas robežās esošas zemes atradās vāciešu rokās. Pēc īpašuma cenza pasīvās vēlēšanu tiesības bija tikai 115 poļiem un 315 vāciešiem. Taču vācu pilsonības spēks bija tikai relatīvs. Vācu pilsētu iedzīvotāju vidū lielāko daļu veidoja sīkpilsonība. Proletariāta īpatsvars pilsētnieku vidū bija neliels. Pat Pozenes pilsētā XIX gs. pirmajā pusē tas sasniedza augstākais 20\% no visiem iedzīvotājiem.

Poļu attieksme pret vāciešiem bija dažāda, tomēr vairākums juta līdzi mērķim atjaunot neatkarīgu Polijas valsti. Kad 1830.~gadā (un arī vēlāk 1863.~gadā) Krievijas piederošajā Polijas daļā sākās sacelšanās, daudzi poļu tautības Prūsijas pavalstnieki tajā piedalījās. Ap 12 tūkstošu Prūsijas poļu iestājās sacēlušos rindās Polijas karalistē. (Ir arī citi dati~--- ap 2~500 Prūsijas poļu piedalījušies cīņā pret Krievijas armiju.) Tie pārsvarā bija jauni cilvēki~--- amatnieku zeļļi, ģimnāzisti, dienas strādnieki. Sakauti viņi atkāpās uz Prūsijas teritoriju, caur Pozeni bēga uz Rietumiem. Taču jau 1830.~gada septembrī Pozenes ielās varēja atrast skrejlapas un plakātus ar draudiem vāciešiem. Tika atklāta gatavošanās uz poļu sacelšanos, kas, protams, nevarēja neizraisīt varas iestāžu pretsoļus.

Kaut Pozenes vāciešu lielākā daļa pret 1830.~gada poļu sacelšanos izturējās vienaldzīgi, prūšu ierēdņu acīs poļi vairs nebija ``uzticami pavalstnieki'' kā tas tika pausts agrāk, jo bija pierādījies, ka vietējo poļu mērķis nav uz mūžīgiem laikiem iekļauties Prūsijas sastāvā, bet viņi vēlas atjaunot Polijas valsti. Tas lika arī Prūsijai pievilkt grožus stingrāk. Savukārt Pozenes vāciešu pārstāvji pat devās kā brīvprātīgie uz Polijas karalisti, lai piedalītos poļu sacelšanās apspiešanā. 1830.~gada sacelšanās, tāpat kā Krievijā, arī Pozenē ievadīja izmaiņas t.s. ``poļu politikā'', pastiprinājās ģermanizācija. Salīdzinājumā ar Polijas karalisti, kur rusifikācija notika pakāpeniski, Prūsijā poļu administratīvās un izglītības iestādes tika likvidētas visai ātri. Kā norādījis vācu vēsturnieks M.~Aleksanders, iekļaušana vācu pārvaldes un tautsaimniecības sistēmā nesa vietējiem poļiem ekonomiskas priekšrocības, taču tas maz ietekmēja viņu apziņu, jo tiem nācās neatlaidīgi cīnīties par savas valodas un reliģijas saglabāšanu. Laikabiedriem nemanot Prūsija izaudzināja sev ienaidniekus, kuru nākamās paaudzes jau pēc Polijas valstiskuma atjaunošanas tos pašus ``nacionālos'' ieročus vērsa pret vāciešiem.

Par Pozenes virsprezidentu tika nozīmēts E.~Flotvels. Viņa amata laiku: 1830.--1841.~gadu Pozenes novadam veltītajā literatūrā sauc par \strong{``Flotvela ēru''}. Viņa oficiālais uzdevums bija lauzt poļu muižniecības pārmērīgo ietekmi lielhercogistē, bet faktiskais~--- pilnībā to sapludināt ar Prūsijas karalisti. Lielhercogiste tika pārvērsta par parastu Prūsijas provinci. Lai to panāktu, E.~Flotvels lika priekšā, pirmkārt, uzpirkt poļu muižniecībai piederošo zemi. To izdarīt nevajadzēja būt grūti, jo poļu šļahta bija iestigusi parādos. Otrkārt, bija jālikvidē poļu valodas pasniegšana vidējās un augstākajās mācību iestādēs. Treškārt, bija jāpakļauj valsts kontrolei katoļu garīdznieku izglītošana. Ceturtkārt, viņš ieteica poļu rekrūšus norīkot dienestā vācu pulkos. E.~Flotvels arī mēģināja nosaukto programmu īstenot. 1832.~gadā par iestāžu darba valodu oficiāli tika izsludināta vācu valoda, tikai zemākie garīdznieki, lauku un pilsētu ierēdņi drīkstēja sarakstīties poliski. 1833.~gadā apgabalu (\detxti{Kreis}) pārstāvju sapulcēm tika atņemtas tiesības izvirzīt kandidātus uz augstāko ierēdņu~--- landrātu amatiem. Ja pirms tam tie lielākoties bija poļi, tagad to vietas ieņēma vācieši. Šai pašā gadā tika nodibināts fonds, kurš iepirka par parādiem pārdodamos poļu muižnieku zemes īpašumus un atkal pārdeva tos nu jau vāciešiem. 1834.~gadā tiesas Pozenē tika pielīdzinātas pārējām Prūsijas tiesām. 1836.~gadā norisa pašvaldību reforma, poļu muižniecība un garīdzniecība zaudēja vairākas priekšrocības. Varas iestādes nonāca konfliktā ar katoļu baznīcu. 1838.~gadā Gnēzenes-Pozenes arhibīskaps M.~Duņins izdeva rīkojumu viņam pakļautajiem garīdzniekiem slēgt laulības starp dažādu konfesiju piederīgajiem tikai tad, ja tie deva solījumu bērnus audzināt katoļu garā. Tas varētu vest pie katoļu un protestantu laulībās augušo bērnu polonizācijas, tāpēc pret viņu tika uzsākts kriminālprocess. Arhibīskaps nevēlējās piekāpties, tika arestēts, notiesāts un ieslodzīts uz pusgadu cietumā. No šī laika Pozenē gan uzturējās un darbojās poļu intelektuāļi, bet galvenās nacionālās cīņas izvērtās blakus~--- Polijas karalistē. Pretpoļu rīcībai no Prūsijas puses vajadzēja stiprināt tās draudzību ar Krieviju.

``Flotvela ēras'' laikā veiktos pasākumus nevarēja nosaukt par gudru politiku. Tās mērķis~--- uz pretestību noskaņotās šļahtas vietā piesaistīt poļu zemniekus un pilsētniekus Prūsijas valstij palika nesasniegts. Kad 1840.~gada jūnijā par Prūsijas karali kļuva Fridrichs Vilhelms IV, lai gūtu Pozenes poļu muižniecības un garīdzniecības atbalstu, viņš vēlējās panākt izlīgumu ar šiem slāņiem. Jau 1840.~gadā arhibīskaps M.~Duņins ar triumfu atgriezās amatā, neatceļot savu rīkojumu par jauktajām laulībām. Tai pat gadā karalis paziņoja amnestiju 1830.~gada sacelšanās dalībniekiem. Drīzumā~--- 1841.~gada martā no amata tika atsaukts arī E.~Flotvels. 1842.~gadā Prūsijas Kultūras ministrijā tika radīta Katoļu nodaļa, kura galvenokārt nodarbojās ar poļu skolu jautājumu.

Relatīvās samierināšanās periodā poļu patrioti mēģināja savu slepeno organizāciju centrus, kas gatavoja sacelšanos, pārcelt uz Pozeni. Tomēr saspringuma atslābums poļu un vācu attiecībās bija tikai īslaicīgs. Prūsijas policija saņēma ziņās par poļu pagrīdes darbību, kura gan Pozenā nebija diez cik sekmīga.

Kad 1846.~gadā Pozenē ieradušies poļu revolucionāru pārstāvji no Polijas karalistes un emigranti no Parīzes, lai aģitētu par nacionālu sacelšanos, tie nesaņēma atbalstu no vietējiem poļu zemniekiem, kuri atcerējušies poļu šļahtas nodarītās pārestības. Vācu literatūrā aprakstīts gadījums, kad kāds poļu dižciltīgais ieradies ciema krogā, licis ieliet tur sanākušajiem zemniekiem šņabi un uzaicinājis tos tvert ieročus cīņai par Polijas brīvību. Uz to piecēlies sirms zemnieks un ar vārdiem ``\pltxti{Panie, dziękuję za waszą wolność}'' (``Kungs, pateicos par jūsu brīvību'') atpogājis kreklu un parādījis pletnes pēdas, kuras bija palikušas uz viņa muguras vēl no poļu ``brīvības'' laikiem. Saprotams, ka šādu ``brīvību'' zemnieki nevēlējās.

1846.~gada sākumā Pozenē tika arestēts poļu pagrīdes organizācijas vadītājs 1830.~gada sacelšanās dalībnieks L.~Mieroslavskis un ap 500 sazvērestības dalībnieku. No 1847.~gada augusta līdz decembrim 254 sazvērnieki tika tiesāti atklātā procesā Berlīnē. Apsūdzētie ar to ieguva atklātu tribīni, kur varēja izklāstīt savus mērķus. Viņi veikli saistīja Polijas un Vācijas jautājumus, jo abu tautu priekšā stāvēja jautājums, kā tās nonāks pie savas nacionālas valsts? Šādas runas vācu sabiedrībā guva pozitīvu atbalsi. Kā norādījis vācu vēsturnieks E.~Meijers, vācu liberāļi simpatizēja poļiem tāpēc, ka baidījās no tā, ka Krievija necietīs Vācijā liberālas reformas un poļi krievu iebrukuma gadījumā varētu būt vācu sabiedrotie. Arī tiesneši nepalika kurli pret izskanējušajiem argumentiem. Vairāk nekā puse apsūdzēto tika attaisnota. Pārpalikušie saņēma vairāk vai mazāk stingrus cietumsodus. Tikai astoņiem, tostarp arī L.~Mieroslavskim, piesprieda augstāko~--- nāves sodu. Tiesa, neviena dzīvība netika atņemta. 1848.~gada revolūcijas laikā 20.~martā berlīnieši viņus atbrīvoja no cietuma un tie ar triumfu varēja atgriezties Pozenē. Uz Berlīni tika nosūtīts pateicības vēstījums, kurā bija izteikta cerība, ka nāk laiks, kad ``brīvās Vācijas drošībai tiks uzcelta neatkarīgā Polija kā priekšnocietinājums pret aziātu spiedienu''.

Kad 1848.~gadā revolūcijas vilnis vēlās pār Eiropu, tās notikumos piedalījās arī poļi. Ā.~Mickevičs Romā nodibināja poļu leģionu, lai atbalstītu itāļu brīvības cīnītājus. Bādenē sacēlušos pusē 1849.~gadā cīnījās vairāk nekā 300 poļu virsnieku un karavīru. Otrs 1830.~gada sacelšanās dalībnieks ģenerālis J.~Bems vadīja Vīnes sacēlušos pret Austrijas ķeizara armiju. Pēc Vīnes sacelšanās sakāves viņš pievienojās revolucionāriem Ungārijā, kur L.~Košuta vadībā jau karoja ap 3~000 poļu leģionāru. Pēc sacelšanās sakāves J.~Bems bēga uz Krievijas ienaidnieces Turcijas teritoriju, pieņēma tur islāmu, vadīja turku karaspēku pret sacēlušamies arābiem, mira kā turku ģenerālis. J.~Bema piemiņa dzīvo Ungārijā viņa bijušā adjutanta, dzejnieka Š.~Petefī dzejā. Budapeštā ir uzcelts J.~Bema piemineklis. Interesanti, ka \strong{skanēja arī šādu poļu aktivitāti nosodošas balsis}. Piemēram, poļu publicists J.~Janovskis pauda poļu egocentrismu, 1848.~gadā uzstājoties par poļu asiņu ``taupīšanu'' cīniņā par ``svešām'' lietām, kalpošanu Polijai turpretī uzskatot par ``kalpošanu cilvēcei''.

Daļa Eiropas revolucionāru cerēja, ka tiešu poļu piemērs iejūsminās cīņai pārējas tautas. Vācu revolucionārais dzejnieks G.~Hervegs uzrakstīja dzejoli, kurā poļu patriotu vārdā griezās pie Eiropas tautām:

\vspace{1.5em}

\noindent
\begin{minipage}{0.6\textwidth}
\detxti{
``An dich, Europa, richten wir die Frage:\\
Veläβt du uns zum zweitenmal?''}
\end{minipage}
\hspace{1em}
\begin{minipage}{0.45\textwidth}
``Tev, Eiropa, mēs jautājam:\\
Vai atstāsi mūs otru reiz?''
\end{minipage}

\vspace{1.5em}

% page 92


Pēc Francijas karaļa Luija Filipa I atkāpšanās arī Vācijā aktivizējās revolucionārie spēki, izvirzījās iniciatīva izveidot Visvācijas parlamentu. Tā sagatavošanas gaitā 1848.~gada 31.~martā tika sasaukts t.s. ``Visvācijas priekšparlaments'' (\detxti{Vorparlament}) Frankfurtē. Tā dalībnieki aicināja noturēt brīvas parlamenta vēlēšanas, izteicās arī par Polijas apvienošanos. Priekšparlaments pasludināja Polijas dalīšanas par ``kaunpilnu netaisnību'' un ``vācu tautas svētu pienākumu veicināt Polijas atjaunošanu''.

Tāpat kā citās valstīs, arī Vācijā revolucionāri prasīja Polijas neatkarības atjaunošanu, lai tā sagrautu ``Svēto savienību''. Sociālists un viens no t.s. ``zinātniskā'' komunisma pamatlicējiem K.~Markss atzina Polijas nacionālās atbrīvošanas nepieciešamību, jo, pēc viņa domām, demokrātija Polijā nebija iespējama bez feodālo tiesību iznīcināšanas, bez agrāras kustības, kura pārvērstu feodāli atkarīgos zemniekus par brīviem īpašniekiem, bet agrārā revolūcija nebija iespējama bez vienlaicīgas nacionālās patstāvības izcīnīšanas. K.~Markss Polijas atbrīvošanu uzskatīja par pašas Vācijas atbrīvošanas priekšnoteikumu no patriarhāli feodālā absolūtisma.

Taču krievu vēsturnieks N.~Uļjanovs, kurš darbojās emigrācijā Rietumos, parādījis arī citu t.s. ``marksisma klasiķu'' nostājas interpretāciju. Viņi atbalstīja poļu revolucionāro kustību tāpēc, ka tā bija galvenokārt vērsta pret Krieviju, kura savukārt šai laikā Eiropā uzstājās kā Vīnes kongresā radītās starptautisko attiecību sistēmas garants. Pret Krieviju bija vai nu jāorganizē Eiropas valstu karš vai no tās jānorobežojas ar neatkarīgu Polijas valsti.

K.~Markss un F.~Engelss apzinājās, ka revolucionāru lomā Polijā pārsvarā uzstājas muižnieki-feodāļi, kuri tikai vārdos atzina sociālo atbrīvošanos. Taču šai gadījumā viņus maz interesēja, kas virza nacionālās atbrīvošanās kustību~--- demokrāti vai aristokrāti-muižnieki. Par 1830.~gada poļu sacelšanos K.~Markss un F.~Engelss rakstīja: ``Sauklis ``Lai dzīvo Polija!'' pats par sevi nozīmēja: nāvi Svētajai savienībai, nāvi Krievijas, Prūsijas un Austrijas militārajam despotismam, nāvi mongoļu kundzībai pār mūsdienu sabiedrību!'' Pie tam viņi atbalstīja poļu nacionālās atbrīvošanās kustību galvenokārt Polijas karalistē, noklusējot par Pozeni, kurai bija jāpaliek vāciskai.

Taču Polijas valstiskuma atjaunošana nevarēja būt nekas vairāk kā revolucionāru nodoms. Polijas sadali taču veica tās pašas trīs valstis: Prūsija, Austrija, Krievija, no kurām bija atkarīgs, vai tiks panākta Vācijas vienotība, bet Polijas neatkarības atjaunošanas plāns neglābjami pārvērta tās par iecerētās nacionālās apvienošanās pretiniecēm. Tā, Vācijas impērijas radīšana bija iespējama tikai Krievijas neitralitātes apstākļos, bet to varēja nodrošināt tikai atsakoties no Polijas patstāvības atjaunošanas. Tāpēc drīz ``marksisma klasiķi'' bija vīlušies savās cerībās. 1851.~gada 23.~maijā, sakarā ar izveidojušos starptautisko situāciju, kur Francija, Itālija un Polija bija ieinteresētas Vācijas sadrumstalotībā, un tai, kā uzskatīja F.~Engelss, izņemot Ungāriju, varēja būt tikai viens sabiedrotais~--- Krievija, kurā būtu uzvarējusi zemnieku revolūcija, viņš uzrakstīja garu vēstuli K.~Marksam, kurā izklāstīja savas pārdomas par poļiem, kā ``demoralizētu nāciju'' (\frtxti{nation foutue}): ``Viņus nākas lietot tikai kā līdzekli, un tikai līdz tam laikam, kamēr pati Krievija nepiedzīvos agrāru revolūciju. No šī brīža Polija zaudē jebkādas tiesības uz pastāvēšanu''. Attieksme pret poļiem F.~Engelsam bija radusies ļoti kritiska: ``Poļi vēsturē nekad nav darījuši neko citu, kā tikai tēlojuši drosmīgu un nebēdnīgu muļķību''. ``Poļos ir nemirstīga tieksme uz ķildām bez mazākā iemesla.'' Un visbeidzot: ``nevar atrast nevienu brīdi, kad Polija, kaut pret Krieviju, sekmīgi pārstāvētu progresu vai vispār paveiktu ko tādu, kam būtu vēsturiska nozīme. Pretēji tai Krievija tiešām iemieso progresu attiecībā pret Austrumiem.'' F.~Engelss Krievijā atrada daudz vairāk civilizācijas, izglītības un rūpniecības attīstības elementu nekā ``šļahtiski-miegainajā Polijā''. ``Nekad Polija neprata asimilēt nacionālā ziņā svešus elementus. Vācieši poļu pilsētās ir un paliek vācieši. Bet katrs Krievijas vācietis jau otrā paaudzē ir dzīvs piemērs tam, kā Krievija prot rusificēt vāciešus un ebrejus. Pat ebrejiem tur izaug slāvu vaigu kauli.'' `` Ceturtdaļa Polijas runā lietuviski, ceturtdaļa~--- rutēniski, neliela daļa~--- pa pusei krievu valodā, bet pati poļu daļa par aptuveni trešdaļu ir ģermanizēta''. F.~Engelss izteica pārliecību, ka ``Krievijā agrārā revolūcija notiks agrāk nekā Polijā, kā krievu nacionālā rakstura, tā arī vairāk attīstīto buržuāzisko elementu rezultātā. Ko gan nozīmē Varšava un Krakova pret Pēterburgu, Maskavu, Odesu utt.!''. Par F.~Engelsa ``pragmatismu'' liecināja konstatējums: ``Nācijai, kura var uzstādīt maksimāli 20~000~--- 30~000 cilvēku [domāts, karavīru~--- V.Š.] nevar būt balss tiesību. Bet Polija, acīmredzot, daudz vairāk nespēs uzstādīt.'' Tikpat revolucionāri-``pragmātisks'' bija secinājums: ``ar aizsardzības ieganstu paņemt no poļiem rietumos visu ko vien var, ieņemt viņu cietokšņus, īpaši Pozeni, ar vāciešiem, ļaut viņiem saimniekot, sūtīt viņus ugunī, aprīt viņu produktus, barot viņus ar Rīgas un Odesas solījumiem, bet gadījumā, ja izdotos iesaistīt kustībā krievus, apvienoties ar tiem un piespiest poļus piekāpties.'' Acīmredzot, nav īpaši jāsaka, ka F.~Engelss šeit pirmkārt rūpējās par Vācijas interesēm, kura viņam bija galvenā progresa un revolūcijas nesēja Eiropā. K.~Markss tāpat uzskatīja, ka Prūsijas Poliju, ar pilsētām, kuras apdzīvo vācieši, nav jāatdod ``tautai, kura līdz šim vēl nav pierādījusi savu spēju izkļūt no pusfeodālā dzīves veida, kurš balstās uz lauku iedzīvotāju nebrīvi.''

1848.~gada revolūcijas laikā Vācijā, nemieri (\pltxti{powstanie wiełopolskie}) notika arī Pozenē. Tur kā poļu pārstāvniecība radās Nacionālā komiteja (\pltxti{Komitet Narodowy}). Uz vietām poļu komitejas atņēma varu prūšu administrācijai. Prūsijas karalis Fridrichs Vilhelms IV 1848.~gada 24.~martā apsolīja ``nacionālu reorganizāciju'' vēl ar poļu vairākumu esošajā Pozenes lielhercogistē.

28.~martā Pozenē ieradās L.~Mieroslavskis un sāka organizēt brīvprātīgas poļu vienības (tajās sapulcējās ap 9~000 galvenokārt izkaptīm bruņotu vīru), kuru dalībnieki cerēja uz kopīgu vācu un poļu karu pret Krieviju. Pēc dažu poļu revolucionāru plāniem poļu korpusam no Pozenes bija jāiebrūk Polijas karalistē, lai ar to izraisītu Prūsijas un Krievijas karu. Prūsijas varas iestādes piesardzīgi 3.~aprīlī izsludināja aplenkuma stāvokli, sākotnēji uzsāka ar L.~Mieroslavski pārrunas, taču kad viņš pieprasīja nākamajai Polijas valstij arī Pozeni, stāvoklis saasinājās. Pozenes vācieši, kuri provinces robežapgabalos veidoja vairākumu, nevēlējās ietilpt poliskā valstiskā veidojumā, pieprasīja sadalīt provinci pēc nacionālās piederības principa.

14.~un 26.~aprīlī tika izdoti divi Prūsijas valdības rīkojumi. Atbilstoši tiem Pozenes lielhercogiste bija jāreorganizē. Vienā tās daļā tika solīts ierīkot vairāk vai mazāk autonomu poļu pašpārvaldi, tajā būtu garantēta izglītība, lietvedība un tiesas poļu valodā, ierēdņi varētu būt poļi. Toties lielākā, lielhercogistes rietumu daļa ar Pozenes pilsētu tiktu iekļauta Vīnes konferences radītajā Vācu savienībā (1815--1866). Poļi tam pretojās. Tā kā Pozenes lielhercogiste līdz tam gan bija piederīga Prūsijai, bet ne Vācu savienībai, šo rīkojumu saņemšana izsauca bruņotas sadursmes starp abām pusēm. 29.~aprīlī Prūsijas armija vērsās pret saformētajām poļu vienībām. Notika vairākas sīkas sadursmes. Nacionālā komiteja pašlikvidējās. Lai gan divās sadursmēs L.~Mieroslavskis uzvarēja, pret labi apmācīto prūšu armiju poļiem nebija tikpat kā nekādu izredžu. Nedēļas laikā poļu sacēlušos sakāva. 6.~maijā (citi dati~--- 9.maijā) viņi bija spiesti kapitulēt, taču ne tikai ierindas poļu kaujinieku, bet arī to vadoņi varēja mierīgi atgriezties dzīves vietās. L.~Mieroslavskis gan tika saņemts gūstā, bet 1848.~gada jūnijā pēc franču diplomātu iejaukšanās viņu atbrīvoja un jau nākamajā gadā viņš komandēja revolucionāru vienības Sicīlijā un Bādenē (Vācijā). Kaut poļu aktivitātes tika apspiestas ar bruņotu spēku, vācu vēsturnieki uzsver, ka par sacelšanās ātro sakāvi jāpateicas ne tik daudz prūšu varai, kura šai brīdī bija vāja, bet gan tam, ka poļu zemnieki nevēlējās ar savu ādu maksāt par poļu šļahtiču vēlmju piepildīšanu. Sacēlušos sagrāve 1848.~gada maijā gan nenozīmēja, ka uzreiz sākās strauja poļu ģermanizācija. Taču arī valdības aprīlī izdotie rīkojumi netika atcelti.

Frankfurtes parlaments (\detxti{Frankfurter Nationalversammlung}~--- pirmais visas Vācijas parlaments, kurš darbojās no 1848.~gada 18.~maija līdz 1849.~gada 31.~maijam) tika iesaistīts nacionālajos konfliktos, tai skaitā Polijas jautājuma izskatīšanā. Debates par poļu jautājumu notika 1848.~gada jūlijā. Centrālais bija jautājums, vai jaunajā Vācijas impērijā, kas tika projektēta, jāietilpst visai Pozenes provincei vai tikai tai tās daļai, kuru pārsvarā apdzīvoja vācieši. Tikai neliela deputātu daļa izteicās par atteikšanos no Pozenes poļu daļas. Tā sadursmē ar dzīves realitāti sāka mainīties vācu nacionālās idejas kā brīvības, neatkarības un vienotības idejas raksturs. Ar lielu balsu vairākumu Nacionālā sapulce pieņēma lēmumu par Pozenes rietumdaļas iekļaušanu Vācijas impērijā, bet no idejas atjaunot neatkarīgu Poliju atteicās.

Izklāstītie notikumi saasināja attiecības starp poļiem un vāciešiem Pozenē, parādīja, ka vācu vidū sajūsma par poļu cīņu par brīvību un pret Krieviju bija noplakusi. Vietējie vācu iedzīvotāji vērsās pret reformu plāniem.

Kopš 1848.~gada Berlīnē oficiāli runāja ne vairs par Pozenes lielhercogisti, bet Pozenes provinci, kaut arī vēl pēc Vācijas apvienošanās (1871) tās imperators lietoja arī Pozenes lielhercoga titulu. Prūsijas iekšējā organizācijā arī vairs nelietoja jēdzienu ``Prūsijas poļi'' ,bet gan ``poļu prūši''. Vienai novada daļai solītā autonomija tā arī nekad netika sagaidīta. Turpmāk vietējās saimnieciskās organizācijas tika veidotas jau pēc nacionālā principa: vāciešiem~--- savas un poļiem~--- savas. Tāpat kā Galīcijā, arī Pozenē vairs nebija cerību izraisīt visas tautas bruņotu brīvības cīņu. Izsīka arī Pozenes poļu emigrantu enerģija.

\strong{Ģermanizācijas politika} Pozenē tika izvērsta arvien noteiktāk. ``Izlīgums ar poļiem,'' jau 1848.~gada 20.~aprīlī O.~v.~Bismarks rakstīja avīzē ``\detxti{Magdeburger Zeitung}'' (``Magdeburgas Avīze'') - ``ir neiespējams; tos pilnībā iznīcināt ir nehumāni un turklāt neiespējami, ja vien paaudzes par to nenodarbosies''. Tika slēgtas poļu avīzes, aizliegtas poļu organizācijas. Kopš 1852.~gada dažādi oficiāli rīkojumi un norādījumi tika izplatīti tikai vācu valodā. Tomēr vēl XIX~gadsimta vidū pastāvēja laba vācu un poļu sadarbība izglītības jomā. Kopīgiem vācu un poļu spēkiem 1853.~gadā izdevās panākt reālskolas atvēršanu Pozenē. Abu tautu pārstāvji kopīgi apmeklēja sporta biedrības, piemēram, kopš 1838.~gada eksistējošo šaha klubu, 1860.~gadā dibināto vingrošanas biedrību. Taču 50.~gados Pozenes provincē palika vairs tikai trīs poļu ģimnāzijas, 1856.~gadā poļu ģimnāzijām tika pavēlēts latīņu klasiķus jau no jaunākajām klasēm tulkot tikai vācu valodā. 1858.~gadā visās pilsētas skolās visi mācību priekšmetu, izņemot dziedāšanu un poļu valodu, bija jāpasniedz vāciski. 1867.~gadā visās poļu sākumskolās ieviesa obligātu vācu valodas apmācību. Poļu valoda pamazām ieņēma arvien nenozīmīgāku lomu. Kā atzīmē vācu vēsturnieki, vācu liberālā opozīcija gan juta līdzi poļu patriotiem, bet rezultātu tam nebija. Turpretī auga nepatika pret poļu kaimiņiem. Vācu rakstnieka G.~Freitāga 1854.~gadā izdotajā romānā ``\detxti{Soll und Haben}'' (Vajadzēt un būt), kurš līdz gadsimta beigām piedzīvoja 40~izdevumu, kā tipiskas poļu īpašības tika nosauktas nesaimnieciskums, nolaidība, viltība.

1857.~gadā no pavisam 1~440 muižām Pozenes provincē 536 (jau 37\%) atradās vācu rokās. Par šo straujo pieaugumu bija jāpateicas valsts politikai. Valsts iestādes gan Berlīnē, gan Pozenē uzskatīja par savu uzdevumu palielināt vācu īpašumus provincē. Tomēr krasi noskaņotie vācu nacionālisti arī tad nebija mērā ar it kā nenoteikto valdības politiku. 1861.~gadā kādam parlamenta deputātam domātā vēstulē minētais G.~Freitāgs rakstīja: ``Mani tagad nodarbina Pozenes apsaimniekošanas un kolonizācijas projekts un es vēlos Berlīnē par to aģitēt, lai ar poļiem ātri tiktu galā. Šī ilgstošā vājuma apkaunojums kļūst pārāk liels. Ja valdībai nav tam spēka, mums ir jāmēģina to sagādāt privāti.'' Jāatzīmē, ka daudzi poļu muižnieki netika vaļā no agrākajiem parādiem, nespēja atteikties no agrāk ierastā dzīves stila, pielāgoties jaunajiem modernajiem saimniekošanas paņēmieniem un bija spiesti muižas pārdot pat bez īpaša spiediena.

Nākamais poļu un vācu attiecību saasinājuma periods sākās ar poļu nemieriem blakus esošajā Polijas karalistē 19.~gadsimta 60.~gados. Prūsijas premjerministrs O.~f.~Bismarks baidījās, ka sacelšanās Polijas karalistē varētu izplatīties arī uz Pozeni un tāpēc Prūsija 1863.~gada 27.~janvārī (8.~februārī) noslēdza ar Krieviju tā saucamo Anvenslēbena konvenciju (\detxti{Die Alvenslebensche Konvention}), nosauktu pēc Prūsijas pārstāvja sarunās ģenerāļa G.~f.~Alvenslēbena vārda, kura atļāva abām valstīm vajāt sacēlušos arī tad, ja tie bēga uz otras valsts teritoriju. Vienošanās nāca par labu Krievijas karaspēkam, kurš varēja sekot poļiem arī Prūsijas teritorijā. Taču literatūrā ir ziņas, ka 1863.--1864.~gadā vairāk poļu sacelšanās dalībnieku Prūsijā gāja bojā no prūšu karavīru, nekā no poļiem sekojošo krievu karavīru lodēm. Arī Prūsijā poļu patrioti tika tiesāti par palīdzību sacēlušamies Polijas karalistē. Kā uzskatīja poļu vēsturnieks A.~Žeļagovskis, šī konvencija bija O.~f.~Bismarka ``meistarstiķis'', kurš uz gadu desmitiem noteica apstākļus Eiropas austrumos, jo padarīja Aleksandram II neiespējamu jebkādu samierināšanās politiku ar poļiem un nodrošināja Prūsijai rīcības brīvību pret Austriju un Franciju.

Joprojām gan daļa Pozenes vāciešu ne tikai ar nepatiku skatījās uz Krievijas un arī Prūsijas varas iestāžu antipolisko darbību, bet pat atbalstīja sacēlušos, šujot tiem uniformas, palīdzot transportēt ieročus. Prūsijā tika runāts par Polijas valsts atjaunošanas nepieciešamību. Taču poļu kustībā, īpaši zemākajos slāņos, parādījās arī antivāciskas izpausmes, tāpēc pēc 1863.~gada poļu sacelšanās sakāves vāciešos auga antipoliskie noskaņojumi. Tomēr ikdienā dažādu etnisko grupu kopējā dzīve lielākoties norisa mierīgi, pat draudzīgi. Poļu neapmierinātība pārsvarā bija vērsta pret vācu iestādēm, bet ne vācu tautības līdziedzīvotājiem. Protams, bija arī izņēmumi.

Ģermanizācijas politika vēl pastiprinājās, kad XIX gadsimta otrajā pusē mainījās Prūsijas un citu valstu spēku samērs. Antipoliskie noskaņojumi bija arī rezultāts Prūsijas militārajiem panākumiem karos pret Dāniju un Austriju un centieniem pēc Vācijas vienotības pieauguma. Īpaši karos pret Austriju (1866) un Franciju (1870/1871) tika radīti priekšnoteikumi Vācijas impērijas radīšanai ar Prūsiju priekšgalā. Kaut poļu deputāti Prūsijas landtāgā un Vācijas reihstāgā protestēja pret viņu zemju iekļaušanu Ziemeļvācijas valstu savienībā (1866) un Vācijas impērijas sastāvā (1871), Pozenes province kopā ar Prūsiju kļuva par šo valstisko veidojumu sastāvdaļu. Vācijas iedzīvotāju vairākums bija sajūsmā par vienotas Vācijas impērijas nodibināšanu, ar neizpratni un aizvainojumu vērsās pret vietējiem poļiem, kuri ne tikai neizrādīja šādu sajūsmu, bet prasīja autonomiju un sapņoja par savu valsti. Vācieši tagad arī redzēja, ka Vācijai nav ko baidīties no Krievijas, arī tāpēc ātri pieauga nepatika pret dumpīgajiem poļiem. Arī 1871.~gada Vācijas Konstitūcija runāja tikai par ``vācu tautu'', neparedzot nekādas mazākumtautību tiesības.

Vācu uzvara karā pār Franciju, Vācijas impērijas izveide nostiprināja Prūsiju, ļāva tai aktivizēt ģermanizācijas politiku Pozenē. Tiesa, Vācijas impērijas nodibināšana nesa arī noteiktu tiesību paplašināšanos Pozenes poļiem. Atšķirībā no Prūsijas landtāga triju kārtu vēlēšanām, Vācijas reihstāga vēlēšanās varēja piedalīties visi pieaugušie vīrieši. Poļu deputātu skaits tajā svārstījās no 13 līdz 20. Gandrīz visos jautājumos, izņemot balsošanu 1893.~gadā par armijas palielināšanu, viņi kopā ar pārējo Vācijas mazākumtautību pārstāvjiem (dāņiem, elzasiešiem) balsoja kopā ar opozīciju. Panākt savu priekšlikumu ievērošanu viņi, protams, nevarēja, taču uzskatīja par vajadzīgu demonstrēt savu naidīgumu valdībai.

Vācu vēsturnieks H.~f.~Zizevics atzīmējis, ka periodā no 1862.~gada, kad O.~f.~Bismarks kļuva par Prūsijas Ministru padomes priekšsēdētāju, līdz pat 1916.~gadam Prūsijas un Vācijas politikā pret poļiem pastāvēja divi strāvojumi, kuri nepārtraukti mēģināja viens otru pārvarēt un iegūt vairākumu parlamentā. H.~f.~Zitzevics tos nosaucis par administratīvi-likumīgo virzienu un politisko virzienu.

Pirmais prasīja nekādā gadījumā nepārkāpt valsts likumus. Poļiem bija jākļūst par labiem Vācijas pilsoņiem, bet likumiem un citiem rīkojumiem bija jāpalīdz viņiem tādiem kļūt.

Otrais bāzējās uz sociāldarviniskiem priekšstatiem. Vācietība reprezentēja pilnvērtīgo, valdošo rasi, slāvi~--- nepilnvērtīgo rasi, pār kuru bija vajadzīga aizbildniecība un vadība. Attiecīgi aizbildņi varēja noteikt aizbildināmo tiesības un pienākumus.

Autoram gan jāsaka, ka otrais virziens bija redzamāks, darbojās aktīvāk. Augstākās personas Vācijas impērijā bija noskaņotas antipoliski.

Tā kanclers O.~Bismarks poļus konsekventi sauca par ``impērijas ienaidniekiem''. Pēc viņa vārdiem ``\dots{}tikai tas, kurš nepavisam nepazīst poļus, var šaubīties par to, ka viņi paliks zvērināti mūsu ienaidnieki\dots{}''

Stiprākais ģermanizācijas līdzeklis Pozenes provincē bija kolonizācija. Pēc vērtējumiem Pozenes provincē pēc tās pievienošanas Prūsijai bija no 200~000 līdz 300~000 jeb 25 līdz 35\% vāciešu. Līdz XIX gadsimta 60.~gadu sākumam \strong{vācu iedzīvotāju skaits} arvien \strong{pieauga}. Pieaugums sastādīja ap 300~000 cilvēku un vāciešu daļa visu provinces iedzīvotāju vidū palielinājās par 5 līdz 8\%. Tā, sakarā ar Pozenes cietokšņa izbūvi no 1828. līdz 1834.~gadam ieceļoja ap 1~000 vācu strādnieku. 1846--1848.~gada notikumi gan piebremzēja vāciešu ieceļošanu, taču 50.~gados tā atkal turpinājās. 1850.~gadā Prūsijā noslēdzās zemes piešķiršana zemniekiem. Zemnieki no dažādiem Vācijas novadiem Pozenē varēja iegādāties zemi ar vislabākajiem nosacījumiem. Līdz 1870.~gadam ievērojami pieauga vācu zemes īpašnieku slānis.

Taču jau līdz 1870.~gadam varēja konstatēt arī vāciešu izceļošanu no provinces. No 1816. līdz 1825.~gadam auga lētu tekstilizstrādājumu ievedums Pozenē no Silēzijas un Saksijas, kas audēju amatu noveda smagā krīzē, jo ar blakus esošo Krieviju pastāvēja muitas robeža. Daudzi vācu audēji pameta Pozeni un pārvācās uz poļu apdzīvotajiem apgabaliem Krievijā. 30. un 40.~gados sakarā ar varas iestāžu neiecietību no Pozenes uz Austrāliju izceļoja grupa t.s. ``vecluterāņu'' (šādi sauca luterāņus, kas uzstājās pret 1827.~gadā notikušo jaunas baznīcas liturģijas ieviešanu pēc Prūsijas parauga), bet lielākais izceļotāju skaits devās uz ASV. Kopš XIX gadsimta vidus vācu izceļotāji no Pozenes pārsvarā devās uz Vācijas rietumu rūpnieciskajiem apgabaliem. Līdz 60.~gadiem vāciešu pieplūdums Pozenes provincē bija lielāks nekā atplūdums, no 1824. līdz 1870.~gadam vācu ieceļotāju skaits par 31~000 pārsniedza izceļotāju skaitu. Pozenes pārvaldes iecirknī vācieši sastādīja 30 līdz 35\%, bet Brombergas (\pltxti{Bydgoszcz}) iecirknī gandrīz 50\%.

Tiesa, pret prūšu ierēdņu sniegtajiem datiem nākas izturēties ar neuzticību. 1860.~gadā viņi paziņoja, ka 44,4\% no 1,4~miljoniem Pozenes provinces iedzīvotāju ir vācieši. Taču sabiedrībā šis rezultāts tika pamatoti apšaubīts. Pēc konfesiju statistikas 1864.~gadā, pat ja visus ebrejus pieskaita vāciešiem, arī tad to īpatsvars sasniedza knapi 43\%. Kad 1871.~gadā vēlreiz notika uzskaite, to daļa bija samazinājusies līdz 41\%. Pēc citiem datiem, vāciešu un tiem pieskaitīto ebreju Pozenes provincē 1830.~gadā bija 40\%, 1860.~gadā~--- 45\%, bet 1910.~gadā~--- tikai 38\%, turpretī poļu~--- attiecīgi 60, 55 un 62\%.

Pēc 1861.~gada datiem gandrīz 40\% vāciešu poļu apgabalos dzīvoja pilsētās. Īpaši Pozenes pilsētā strauji auga vācu militāristu, ierēdņu un strādnieku skaits, vācieši tur kļuva par stiprāko tautību un palika tādi līdz 1870.~gadam. Pēc tam pārsvaru guva tādi faktori kā augstāka dzimstība poļu kā katoļu vidū un vācu tautības iedzīvotāju lielākā nekā poļu bēgšana no austrumiem (vācu \detxti{Ostflucht}, poļu \pltxti{ucieczka ze wschodu}), izceļojot uz Vācijas industriāli attīstītākajiem apgabaliem. Berlīnē atskanēja uztrauktas balsis: kā gan Pozenes provinci var ģermanizēt, ja šeit mītošo vāciešu skaits samazinās? Kā var pamatot, ka visur jārunā vāciski, ja to skaits, kuriem tā bija dzimtā valoda, samazinās? Gala secinājums, kuru izdarīja vācu birokrāti-nacionālisti bija: ir jāpaātrina poļu ģermanizācija.

Prūsijas konservatīvie, kuri neko negribēja dzirdēt par tautību izlīgumu, sauca to par [mazākuma] tautību reiboni (\detxti{Nationalitätenschwindel}), uzskatīja, ka ``poļu problēmai'' pirmkārt ir agrārpolitisks raksturs. Izturīgie un pieticīgie poļu sezonas strādnieki nodrošināja Vācijas centrālo un austrumu novadu muižu rentabilitāti. Pie tam, līdz Pirmajam pasaules karam augot Vācijas bruņotajiem spēkiem, neraugoties uz polisko izcelsmi un katolisko ticību, divvalodīgie Prūsijas poļu pavalstnieki lieti noderēja armijas papildināšanai. Vietējo vāciešu vairākums vairs nevarēja iedomāties Pozenes provinces atdalīšanos no Prūsijas un Vācijas. Tikai vācu liberāļi vēl piekrita provinces autonomijai Prūsijas sastāvā. Turpretī poļus autonomija vairs neapmierināja pat kā minimālā programma. Ja sākotnēji konflikts aprobežojās ar abu tautību elitēm, tālāk tas vērsās plašumā. Zemākajos slāņos izplatījās negatīvi priekšstati par otras tautas pārstāvjiem. Vācieši smējās par ``poļu saimniecību'' (\pltxti{polnische Wirtschaft}, t.i.~--- nesaimnieciskumu), poļi vāciešus lamāja par švābiem (vācu \detxti{Schwaben}~--- vācieši, runājoši īpašā dienvidvācu dialektā, kuri dzīvo Švābijā~--- \detxti{Schwabenland}) un ``kartupeļu rijējiem''. Kontakti sāka aprobežoties tikai ar oficiālo jomu.

Vācijas kancleram un Prūsijas premjerministram O.~f.~Bismarkam poļu nepakļāvība ļoti nepatika. Savās runās viņš atkal un atkal norādīja, ka Pozenes un Rietumprūsijas provincēm ir nepieciešama sasaiste ar Austrumprūsiju, Pomerāniju un Silēziju ne tikai stratēģisku, bet arī ekonomisku u.~c. apsvērumu dēļ. Pēc viņa uzskatiem poļiem bija jākļūst par poļu valodā runājošiem prūšiem, taču tam bija vajadzīgs, lai viņi mācītos arī vācu valodu, jo tikai tās zināšanas varētu viņiem dot visas Prūsijas pavalstnieku priekšrocības. O.~f.~Bismarks bija norūpējies par poļu nacionālismu, rakstīja, ka Polijas valsts atjaunošana nozīmētu automātisku sabiedrotā radīšanu jebkuram Prūsijas ienaidniekam. Kā cilvēkam viņam esot simpātijas pret poļiem, bet no vācu nacionālā viedokļa neesot citas izejas, kā tikai tos iznīdēt. Īpaši kanclers bija neapmierināts ar katoļu baznīcas priesteru aģitāciju pret Prūsiju. Viņš uzskatīja, ka katoļu baznīca var radīt Vācijai grūtības ārpolitikā, jo Austroungārijā katoļu reliģija bija valdošā un no šā viedokļa Prūsijas poļi varētu gribēt pievienoties tai. Blakus valodu politikai izvirzījās arī otrs ģermanizācijas līdzeklis: cīņa pret katolicismu.

Izvēršoties t.s. kultūrkampfam (vācu \detxti{Kulturkampf}~--- O.~f.~Bismarka un viņam pieslējušos vācu nacionālliberāļu cīņa pret katolisko Centra partijas un katoļu baznīcas ietekmi Vācijā. Latviešu valodā par to sīkāk var izlasīt autora grāmatā ``Vācijas impērija. 1871--1918''), jācieš bija arī poļu garīdzniecībai un draudzēm. O.~f.~Bismarks savā darbā ``Domas un atmiņas'' (``\detxti{Gedanken und Erringerungen}'', 1890) kultūrkampfam veltīto nodaļu sāka ar vārdiem: ``Kultūrkampfa sākumu man pārsvarā noteica tā poliskā puse''. 1872.--1874.~gadā visās ģimnāzijās, kur līdz tam mācības norisa poļu valodā, par vienīgo kļuva vācu valoda, poļu valodu pasniedza tikai kā mācību priekšmetu, bet visās zemākajās skolās poļu valodu varēja izmantot tikai reliģijas mācības stundās. Ja līdz šim poļi prūšu skolas Pozenē un Rietumprūsijā uzskatīja par kultūras veicināšanas līdzekli, tad tagad tās tika vērtētas kā apdraudējums savai tautībai. Atbildes reakcija bija poļu privātskolu skaita straujš pieaugums.

Kā vēlāk atzina pats O.~f.~Bismarks, kultūrkampfs bija ``vairāk vērsts pret polonismu, nekā pret katolicismu''. 1872.~gadā tika izdots Prūsijas likums par skolu uzraudzību, kas poļu garīdzniecībai atņēma skolu uzraudzības tiesības. Pret to protestēja Romas pāvests. Gnēzenes-Pozenes arhibīskaps M.~Ļedohovskis aicināja ticīgos nepakļauties valsts rīkojumiem. Viņam par vairākiem likumu pārkāpumiem pietiesāja naudas sodu 80~000 marku apmērā. Arhibīskaps nemaksāja, tad viņam kā pirmajam no Prūsijas bīskapiem piesprieda 2 gadus cietumsoda. Kad viņš arī tad nenolika savas pilnvaras, valsts reliģisko lietu tiesa paziņoja, ka viņš ir no amata atstādināts. Kad 1876.~gadā arhibīskapu atbrīvoja no aresta, pāvests viņu jau bija iecēlis par kardinālu un viņš devās uz Romu. Apcietināti tika arī apmēram 100 citu garīdznieki, daudzu vietas draudzēs palika neaizpildītas. Taču kultūrkampfam negatīva ietekme uz poļu pašapziņu bija tikai Austrumprūsijas Mazūru (poļu \pltxti{Mazurzy}, vācu \detxti{Masuren}) rajonā, kur vietējie luterticīgie iedzīvotāji gan runāja vienā no poļu valodas dialektiem, bet bez poļu pašapziņas padevās arvien lielākai vāciskuma ietekmei, nošķīrās no poļu tautas. Mazūri bieži sevi neuzskatīja nedz par vāciešiem, nedz poļiem, bet par poļu prūšiem. Pozenes novādā šī vācu ietekme bija minimāla. Katrs varas iestāžu konflikts ar katoļu baznīcu stiprināja saikni ``polis~--- tas ir katolis''. Poļu vēsturnieki atzīmē, ka kultūrkampfs pat veicināja poļu kultūras un ekonomiskās darbības pieņemšanos spēkā. Ticības, valodas un kultūras aizsardzība pret to veidoja poļu nacionālo pašapziņu. Tiesa, poļu publicists R.~Šimaņskis 1870.~gadā ar bažām atzīmēja, ka gan Pozenē, gan Silēzijā mantīgie slāņi un inteliģence, kam bija galvenā loma sabiedriskajās attiecībās, pārsvarā sastāv no ``nepoļu elementa''. Tā cēloņus R.~Šimaņskis saskatīja poļu muižnieku šķiriskajā egoismā, kuri nerūpējās par poļu buržuāzijas interešu aizsardzību. Tomēr 1873.~gadā, kad attiecības starp poļiem un vāciešiem kļuva arvien sliktākas, poļu lielmuižnieks M.~Jackovskis, uzskatot, ka poļu stāvokli, viņu nacionālo pašapziņu var nostiprināt ar sociāliem un kultūrizglītības pasākumiem, nodibināja pirmo poļu zemnieku biedrību. Līdz 1880.~gadam tādu bija jau ap 120. Biedrības sniedza saviem biedriem tehnisku palīdzību, deva kredītus, stiprināja poļu tradīcijas un parašas, pozitīvi ietekmēja arī tikumus. Bet, galvenais, tās ļāva cīnīties pret Prūsijas valdības antipoliskajiem pasākumiem, poļiem piederošās zemes izpārdošanu un nonākšanu vācu rokās. Tika dibinātas arī dažādas kultūras un izglītības biedrības.

1876.~gadā tika izdots likums, ar kuru valsts un pašvaldību iestādēs, tiesās vienīgā atļautā kļuva vācu valoda. Kopš 1815.~gada praktizētā divvalodība ar to tika pilnībā atcelta. Tāpēc poļi Vācijas impērijā jutās kā otrās šķiras cilvēki.

O.~f.~Bismarka cīņa pret katoļu baznīcu ilga līdz 1880.~gadam, kad bija redzams, ka tā ir beigusies neveiksmīgi, dienas kārtībā ir izvirzījušās citas nopietnas problēmas. Taču Vācijas kanclera vienošanās ar Romas pāvestu nenozīmēja, ka valdības pretpoliskā politika vājinātos. Izmainījās tikai tās taktika: ``poļu jautājums'' tika nodalīts no ``katolicisma jautājuma''. Taču, kā tagad atzinuši vācu vēsturnieki, sekmīgais asimilācijas process, kurš palēnām norisa kopš 1815.~gada, XIX gadsimta otrajā pusē tika negaidīti pārtraukts. Kā konstatējis vācu vēsturnieks M.~Broskats, ``beznodoma ģermanizācija'' pēc Vīnes kongresa bija saimnieciskā un sociālā progresa rezultāts. Tā vietā XIX gadsimta otrajā pusē sākās valsts realizēta poļu ierobežošanas politika. Tas izsauca poļu kopības sajūtu (\pltxti{społeczeństwo}), kas bija izplatīta visos slāņos un radīja ``valsti valstī''.

Pēc valodu statistikas datiem 1880.~gadā Pozenes provincē dzīvoja 1,7 miljoni iedzīvotāju, no tiem ap 705~000 labi runāja vāciski, nepilns miljons runāja poliski. 1885.~gadā kopējais iedzīvotāju skaits bija pieaudzis par gandrīz 13~000, taču vāciešu skaits bija samazinājies par ap 5~000, toties poļu skaits pieaudzis par 18~000. Tas bija noticis tāpēc, ka no Vācijas austrumu apgabaliem, ne tikai no Pozenes, arvien vairāk vāciešu pārcēlās uz Berlīni, bet īpaši Vācijas rietumiem: rūpniecības apgabaliem pie Reinas un Rūrā. Turpretī daudzi poļi no Polijas Austroungārijas un Krievijas daļām labprāt pārcēlās uz Pozeni vai Rietumprūsiju, jo atrada šeit labvēlīgus apstākļus peļņai.

Atbildot tam, O.~f.~Bismarks izvērsa cīņu pret cilvēkiem, kuriem nebija Vācijas pavalstniecības. Pēc 1880.~gada tautas skaitīšanas Pozenes provincē vien dzīvoja pāri par 10~000 nevāciešu, kuri nebija Prūsijas un Vācijas pavalstnieki. 1885.~gadā O.~f.~Bismarks nolēma izsūtīt visus nenaturalizējušos iebraucējus. Tā kā to vidū bija arī ap 1/3 ebreju, šai politikai ne tikai antipoliska, bet arī antiebrejiska (antisemītiska) virzība. Kopā no Vācijas īsā laikā bija jādeportē 30~000 cilvēku, no tiem 5~200 no Pozenes provinces. Taču plāni sastapa plašu atbalsi un pretestību kā iekšzemē, tā ārzemēs. O.~f.~Bismarks tomēr aizstāvēja savu kursu, lai veicinātu vācietību Pozenes un Rietumprūsijas provincēs. Kopš 1885.~gada ar Prūsijas valdības lēmumu no Vācijas izraidīja 32~tūkstošus cilvēku (galvenokārt poļus, bet 1/3 bija ebreji). Lielākoties tie bija cilvēki, kuri jau gadu desmitus bija dzīvojuši Prūsijā, atraduši tur savu otro dzimteni. Ar 1885.~gadu jauniem iebraucējiem no Krievijas, Austroungārijas ceļš uz Pozeni bija liegts. Tiesa, pēc vairākiem gadiem, ar 1890.~gadu, kad t.s. poļu aizsprosts \detxti{(Polensperre}) tika atcelts, jo muižniekiem vajadzēja darba spēku, Prūsijā atkal ieplūda tūkstoši poļu sezonas strādnieku. To skaits pārsniedza 80.~gados izraidīto skaitu. Viņu stāvoklis gan palika beztiesisks. Pēc mazākā vietējā ierēdņa, bieži arī vācu muižas īpašnieka mājiena viņus varēja izraidīt pāri robežai uz pastāvīgo dzīves vietu. Kā norādīja vācu sociologs M.~Vēbers, strādnieku ieceļošana un izceļošana bija cīņas līdzekļi latentā cīņā starp darbu un īpašumu. Izceļošana bija pielīdzināma latentam streikam, bet poļu ievešana~--- atbilstošs cīņas līdzeklis pret to.

1886.~gada aprīlī parlaments pieņēma ``Likumu par vācu pārceļošanas veicināšanu uz Rietumprūsiju un Pozeni'' (``\detxti{Gesetz betreffend die Beförderung der deutschen Ansiedlung in den Provinzen Westpreuβen und Pozen}''), ar kuru piešķīra finanšu līdzekļus poļu muižnieku zemju uzpirkšanai, lai piešķirtu tās vācu zemnieku un strādnieku pārceļotājiem uz minētajām provincēm. Tika nodibināta karaliskā Prūsijas pārceļošanas komisija (\detxti{Königlich-preuβischen Ansiedlungskomission}) ar sēdekli Pozenē. Līdz 1914.~gadam ap 0,5~miljonu ha zemju iepirkšanai tika iztērēti ap 500~miljonu marku. Rezultātā vācu iedzīvotāju skaits pieauga par 22~000 ģimenēm (citi dati~--- ap 20~000) jeb ap 150~000 cilvēku, taču ap ¼ no viņiem jau bija Pozenes vai Rietumprūsijas iedzīvotāji. Gala rezultātā vācieši no nosauktajiem apgabaliem vairāk izceļoja uz Vācijas rietumu rajoniem, nekā tajos ieceļoja, bet pārceļošanas komisija vairāk iepirka vācu, nevis poļu muižniekiem piederošo zemi.

Ja citās Vācijas provincēs vietējās kopienas varēja līdzrunāt skolotāju izvēlē, Pozenē tā bija pilnīgi valsts rokās. 1887.~gadā no skolām pilnībā padzina poļu valodu. Uzstājoties Prūsijas parlamentā, O.~f.~Bismarks 1887.~gada 28.~janvārī no vienas puses paziņoja, ka pasākumus pret poļiem izsaukusi pašu Pozenes iedzīvotāju~--- poļu izturēšanās, ka 45~gados nav izdevies panākt poļu dižciltīgo labvēlību Prūsijas valsts idejai utt., bet turpat arī apgalvoja: ``Mēs garantējam poļu zemniekiem, uzticamiem Prūsijas pavalstniekiem, ka mēs nevēlamies pret viņu valodu izturēties naidīgi, mēs gribam vienīgi dot viņiem iespēju apgūt vācu valodu''. Lozunga ``Skaldi un valdi'' realizācijai noderēja arī administratīvā reforma, izdalot atsevišķus apgabalus no tiem, kuros dzīvoja pārsvarā poliski runājoši iedzīvotāji. No 1887.~gada līdz 1918.~gadam Pozenes provincē bija 42 apgabali (\detxti{Kreis})~--- par 15 vairāk nekā 1815.~gadā.

Pēc O.~f.~Bismarka atstādināšanas (var piebilst, ka jau minētais poļu vēsturnieks J.~Feldmans 1924.~gadā rakstīja, ka droši vien nav cita vārda, kurš poļu ausīs skanētu tik drūmi, kā Vācijas ``dzelzs kanclera'' O.~f.~Bismarka vārds) kanclera L.~f.~Kaprivi laikā antipoliskais kurss tika mīkstināts, jo jaunajam valdības vadītājam vajadzēja poļu deputātu atbalstu parlamentā. 1891.~gadā skolās atkal atļāva mācīt poļu valodu. Vietējos cilvēkus iecēla par Pozenes virsprezidentiem (vietējās pārvaldes vadītājiem), poļu tautības prelāts kļuva par Gnēzenes-Pozenes arhibīskapu. Novadā tika atļauts iebraukt poļiem no citām valstīm~--- jo lauksaimniecībai trūka strādnieku. Poļu deputāti tāpēc atbalstīja likumus, kuri paredzēja Vācijas bruņoto spēku palielināšanu, ar to rādot, ka viņi nav valsts ienaidnieki. Taču antipoļu politikas mīkstināšana bija īslaicīga.

Poļu deputāti bija pieļāvuši kļūdu, cerot uz vācu ``pīrāgu'', drīz atkal viņi ieraudzīja ``pātagu''. Sakarā ar savu uzņemšanu Tornā (Toruņā) ķeizars Vilhelms II 1894.~gada 22.~septembrī paziņoja: ``Man ir kļuvis zināms, ka poļu līdzpilsoņi šeit neuzvedas tā, kā no viņiem to sagaida un var vēlēties. Viņi var sev atgādināt, ka tie tāpat kā vācieši pilnā mērā varēs paļauties uz manu gādību un līdzjūtību tikai tādā gadījumā, ja bez ierunām jutīsies kā Prūsijas pavalstnieki\dots{}''

Atkal sākās vācu valodas uzspiešana, poļu valodas izskaušana no skolām. 1894.~gadā nodibinājās ``Savienība vācietības (\detxti{Deutschtums}) veicināšanai Ostmarkā''. Palielinājās fonds vācu pārceļotāju atbalstam. Pozenes pilsētā 1900.~gadā poļi sastādīja vairs tikai 55\% iedzīvotāju. No 1900.~gada arī reliģijas mācību visās skolās bija jāmāca tikai vācu valodā.

Arī liberāli, demokrātiski noskaņotie vācieši atbalstīja centienus ģermanizēt poļus. Tika uzskatīts, ka Pozene un Rietumprūsija ar Prūsiju un Vāciju jau tā saaugušas kopā, ka šī saikne kļuvusi nesaraujama uz mūžīgiem laikiem. Pat vācu sociāldemokrāti, XIX gadsimta beigās uzskatīja Polijas neatkarību par neiespējamu. Partijas valdes loceklis I.~Auers 1898.~gadā kolēģei no Polijas R.~Luksemburgai paziņoja: ``\detxti{Wir alle fünf im Parteivorstand betrachten die Unabhängigkeit Polens als einen Unsinn, eine Phantasie''; ``Man kann ``den polnischen Arbeitern keinen größeren Gefallen tun, als sie zu germanisieren, aber man darf es den Leuten nicht sagen.''} (``Mēs visi pieci partijas valdes locekļi Polijas neatkarību uzskatām par neprātu, fantāziju''; ``Nevar izdarīt poļu strādniekiem lielāku pakalpojumu kā viņus ģermanizēt, taču to nedrīkst šiem ļaudīm sacīt.'')

Tomēr ģermnizatoru cerētie rezultāti izpalika. Kaut līdz 1914.~gadam uz Vācijas austrumu provincēm bija pārcēlušies ap 120~000 vāciešu, viņu īpatsvars turējās pastāvīgi pie 38\% robežas. Toties poļu iedzīvotāju īpatsvars, neraugoties uz to nemitīgo izceļšanu, no 62,3\% 1887.~gadā pieauga uz 67,8\% 1905.~gadā. Pozenes provincē izauga moderna, pašapzinīga poļu paaudze, kura bija pārdzīvojusi nacionālu konkurenci un mācījusies no tās atteikšanos no kompromisiem. Valsts pasākumi vācu īpatsvara paaugstināšanai, īpaši poļiem piederošās zemes izpirkšana un piešķiršana vāciešiem, nevarēja būtiski ietekmēt iedzīvotāju nacionālo sastāvu, tikai saasināja nacionālpolitisko konfliktu. Piemēram, kāda poļu avīze 1901.~gadā sludināja: ``Neviens polis nedrīkst precēties ar vācieti vai citu svešu meiteni, tas ir nāves grēks \citespace{} Labāk, ja mūsu meitenes paliek neprecētas līdz nāvei, nekā viņas apprecas ar vācieti.''

Mūsdienās vācu vēstures literatūrā ir parasts Prūsijas-Vācijas politiku ķeizaru Vilhelma I un Vilhelma II un īpaši O.F.~Bismarka laikā uzskatīt kā neizdevušos un vērtēt to negatīvi. Nav noliedzams, ka tās mērķi netika sasniegti. Taču H.f.~Zitzevics norādījis, ka šai politikai bija arī ievērojami panākumi, īpaši kulturālajā un saimnieciskajā jomā. Vispārējais izglītības līmenis (1815.~gadā tikai 20\% skolas vecuma bērnu apmeklēja skolas, XIX gadsimta beigās tas, ka visi bērni apmeklēja skolas, bija pats par sevi saprotams), kā arī pilsētnieku un zemnieku, kā arī visu vidusslāņu saimnieciskais stāvoklis pārdzīvoja tādu attīstību, kādu nesasniedza Krievijai un Austroungārijai piederošajos poļu apgabalos. Piemēram skat. \ref{tab:table1}. tabulu.

\noindent
\begin{table}[h!]
\caption{Lauksaimniecības kultūru ražība hektolitros Pozenes provincē un Polijas karalistē pirms Pirmā pasaules kara} \label{tab:table1}
\begin{tabularx}{\linewidth}{|p{5cm}|p{1.5cm}|p{3.5cm}|}
\hline
\strong{Lauksaimniecības kultūra} & \strong{Pozene} & \strong{Polijas karaliste} \\
\hline
Kvieši & 22,0 & 11 \\
\hline
Rudzi & 18,3 & 10 \\
\hline
Mieži & 43,4 & 11 \\
\hline
Auzas & 22,6 & 9 \\
\hline
Kartupeļi & 156,0 & 62 \\
\hline
Cukurbietes & 305,0 & 198 \\
\hline
\end{tabularx}
\end{table}

% page 101


Taču šie panākumi nestiprināja vācu pozīcijas provincē. Vācu iedzīvotāju skaits samazinājās, poļu vidū auga nacionālistiski noskaņojumi. No poļu valodā runājošiem Prūsijas pavalstniekiem tapa poļi ar nacionālo pašapziņu.

Lai nu kā, Pozenē intensīvais poļu zemnieku-preču ražotāju, kuri iekļāvās tirgus ekonomikas mehānismā, veidošanās process, kā likās, radīja priekšnoteikumus zemnieku pilsoniski-demokrātiskas kustības izveidei. Taču ģermanizācijas politika un poļu zemju kolonizācija atbīdīja sociālās pretrunas otrajā plānā, tuvināja dažādus poļu slāņus, vienoja tos nacionālā cīņā pret ģermanizācijas draudiem. Šādos apstākļos Pozenes poļu zemnieki nonāca aktīvākās poļu sabiedrības daļas~--- mantīgo slāņu, vispirms jau nacionāldemokrātu partijas (par to skat. Polijas karalistei veltītajā apakšnodaļā), ietekmē. Līdz pat 1918.~gadam šeit nepastāvēja atsevišķa zemniecības partija.

Tomēr salīdzinājumā ar kaimiņos esošo Polijas karalisti poļu mazākuma stāvoklis Pozenes provincē bija krietni labāks. Lai citējam tikai tur iznākošā laikraksta ``\pltxti{Oredownik}'' (``Aizstāvis'') 1900.~gada 26.~janvāra numuru: ``Ko mūsu tautieši karalistē (Krievijas Polijā) dara slepus, mēs zem prūšu sceptera darām atklāti, izmantojot Prūsijas Konstitūciju, likumus un valsts iestādes. Visur mēs tiecamies un strādājam pie tā, lai uzturētu spēkā mūsu īpašo nacionālo stāvokli un mūsu raksturu saskaņotu ar mūsu pienākumu pret valsti, kurai mēs piederam. \citespace{} Šeit Prūsijā mums ir tiesības uz īpašu nacionālo stāvokli. \citespace{} Mēs zem prūšu sceptera veicam iekšējo darbu, atklāti un legāli stiprinot un attīstot mūsu tautiskumu, tikai pateicoties tam, ka Prūsijas Konstitūcija mums dod tiesisku pamatu šim iekšējam nacionālajam darbam.''

Pozenes \strong{poļu politiskās dzīves} spektrā ilgstoši vadošo vietu saglabāja konservatīvie elementi, kuri arī pārstāvēja vietējos poļus Vācijas reihstāgā, taču ar laiku to lomu sāka apstrīdēt nacionālistiskās grupas, kas darbojās nacionāldemokrātu (par tiem skat. apakšnodaļu par Polijas karalisti) ietekmē. Par šo slāņu preses orgānu kļuva avīze ``\pltxti{Oredownik}'' (``Aizstāvis'').

80.~gados radās pirmie strādnieku pulciņi, 90.~gados Polijas ``prūšu'' daļa sociālistiskajā kustībā darbojās ap 400 cilvēku. 1893.~gadā arī Pozenē (pēc Polijas karalistes) tika nodibināta Poļu sociālistiskā partija Prūsijā (poļu \pltxti{Polska Partia Socjalistyczna Zaboru Pruskiego}, vācu \detxti{Polnische Sozialistische Partei in Preußen}). Vācijas Sociāldemokrātija (\detxti{Sozialdemokratische Partei Deutschlands}) to uzskatīja par savu sastāvdaļu. Turpretī poļu sociālisti uzskatīja savu partiju par pilnībā autonomu. Attīstījās šīs partijas sakari ar Polijas Sociālistisko partiju (PPS) Polijas karalistē.

80.~gadu otrajā pusē un 90.~gados pastāvošajai zemnieku kustībai bija pārsvarā ekonomisks raksturs. Salīdzinājumā ar citās valstīs esošajiem poļu apgabaliem Pozenē politisko partiju ietekme bija neliela.

Īpaša tēma ir jautājums par \strong{ebreju stāvokli} Pozenes provincē. No vienas puses, ebreji, kļuvuši par Prūsijas pavalstniekiem, tagad varēja bez ierobežojumiem pārvietoties uz Berlīni, Breslavu (vācu \detxti{Breslau}, poļu \pltxti{Wrocław}) un citām Prūsijas, vēlāk Vācijas pilsētām. Tas pavēra viņiem jaunas iespējas. No otras puses, apkārtējie iedzīvotāji labprāt redzētu viņu asimilāciju. Taču tā bija saistīta ar viņu emancipāciju, vienlīdzīgu tiesību iegūšanu, ko citu tautu pārstāvji visbiežāk nelabprāt atzina. Pozenes lielhercogistē 1825.~gadā dzīvoja ap 65~000 ebreju (No 1~032~000 iedzīvotāju tas sastādīja 6,3\%). 1843.~gadā viņu skaits palielinājās līdz 80~000. Lielākoties viņi mita lielākajās pilsētās (Pozenē viņu bija 22\% iedzīvotāju), tad mazākajās, pavisam maz laukos. 63,3\% ebreju nodarbojās ar tirdzniecību, naudas maiņu, uzņēmējdarbību, 34\%~--- ar amatniecību.

1833.~gadā visus Pozenes ebrejus ar likumu sadalīja divās daļās: ``naturalizētajos'' un ``nenaturalizētajos''~--- provincē tikai pieciešamajos. Lai naturalizētos, ebrejam bija jāpierāda, ka viņam ir nevainojama reputācija, jāuzņemas saistības visās sabiedriskās vietās lietot vācu valodu, jāpieņem uzvārds, viņam bija jābūt pastāvīgai dzīves vietai Pozenē kopš 1815.~gada, pietiekošiem eksistences līdzekļiem: profesijai vai zemes īpašumam 2~000 tāleru vērtībā, vai arī kapitālam 5~000 tāleru vērtībā. Naturalizāciju varēja sasniegt arī ar ``patriotisku rīcību''. Pieciešamajiem bija jāievēro daudzi ierobežojumi: viņi drīkstēja precēties tikai pēc 24 gadu sasniegšanas, tikai ārpus pilsētām viņiem privātīpašumā vai nomā drīkstēja būt sava saimniecība, viņiem bija aizliegts būt par pauniniekiem, krogus atļauts turēt tikai pilsētās utt. Abas kategorijas tika piespiedu kārtā apvienotas t.s. ``izraēļu korporācijā'',~--- ebreju kopienā. Līdz 1848.~gadam tikai niecīgs skaits~--- 5,5\% ebreju kļuva par ``naturalizētajiem''.

Ebreju asimilācijas nodoms nebija īstenojams bez viņu izglītošanas. Prūsijas varas iestādes ilgstoši šaubījās~--- ierīkot viņiem atsevišķas mācību iestādes, vai ļaut apmeklēt vispārējās. Galu galā ebrejiem atļāva gan apmeklēt kristiešu skolas, gan uzturēt privātskolas ar valsts ierēdņu ``pārbaudītiem'' skolotājiem. 1845.~gadā uz visiem ebrejiem attiecināja kara klausību. Ebreju kopienai lika nodarboties tikai ar savām garīgajām, reliģiskajām lietām, taču viņiem palika slēgta pieeja nodarbei ar izglītību un vispār kultūru. 1850.~gadā paziņoja, ka ebreji nedrīkst kļūt par tiesnešiem, jo ``viņi nedrīkst pieņemt kristiešu zvērestu''. Liela daļa ebreju bija sašutuši. Viens no viņu rabīniem 1848.~gadā paziņoja: ``Mēs esam un gribam būt tikai vācieši! Mēs nevēlamies nekādu citu tēvzemi kā vācu. Tikai pēc ticības mēs esam izraelīti, visā citā mēs piederam valstij, kurā dzīvojam''. Taču ilgstoši liegt formālu līdztiesību vairs nevarēja. 1869.~gadā Ziemeļvācijas savienības parlamentā pieņēma likumu par ebreju vienlīdzību tiesībās ar pārējiem līdziedzīvotājiem. 1871.~gadā ar Vācijas impērijas nodibināšanos likums tika attiecināts uz visu valsti. Taču pēc tam sākās izceļošanas vilnis, kurš būtiski samazināja ebreju skaitu Pozenes provincē. 1861.~gadā viņu Pozenes provincē bija 74~000~ (5\%), 1880.~gadā~--- 57~000 (3,3\%), 1890.~gadā 44~000 (2,5\%), 1910.~gadā~--- 26~500 (1,3\%). Izceļojot ebreji meklēja lielākas iespējas savai saimnieciskajai darbībai. Pēc 1870.~gada proletārisko ebreju izceļotāju vietā arvien vairāk bija pilsonisko. Attīstoties poļu uzņēmēju darbībai, poļu akciju sabiedrībām, poļu uzņēmēji, konkurējot ar vāciešiem, arvien vairāk izspieda ebrejus. Pie tam darbojās arī citi faktori. No 70.~gadiem daudzi ebreji izceļoja tāpēc, ka Pozenē nebija universitātes, citu augstskolu. Bet ebreju vidū tieksme pēc izglītības bija lielāka nekā kristīgo iedzīvotāju vidū.

Arī ebreju pensionāri (pareizāk~--- rantjē) labprātāk savu mūža nogali pavadīja Rietumu pilsētās, arī Berlīnē, kur bija mazāk antisemītisma izpausmju. Pozenē valdošais kastu gars daudziem saindēja dzīvi (Vācu virsnieki, ierēdņi, uzņēmēji un skolotāji dzīvoja nošķirtu dzīvi no citiem vāciešu slāņiem). Tas nenozīmē, ka starp atsevišķiem ebrejiem un vāciešiem vai poļiem personīgā līmeni nepastāvēja draudzīgas attiecības. Taču poļu un vācu attieksme pret ``savu'' ebreju, piemēram, maiznieku vai kurpnieku, varēja ievērojami atšķirties no attieksmes pret svešu šīs tautības piederīgo.

Daļas iedzīvotāju nepatiku izraisīja ebreju nevēlēšanās nostāties kādas nacionālās vai politiskās grupas atbalstītāju rindās. XIX gadsimta beigās un XX gadsimta sākumā īpašu gan vācu, gan poļu kritiku izsauca cionistu uzskatu sludināšana par ebreju valsts dibināšanas Palestīnā nepieciešamību.

\subsection{Polijas karaliste (Vislas novads) Krievijas varā}

Vīnes kongresā par Polijas karalisti (\pltxti{Królestwo Polskie}, \rutxti{Царство Польское}) pasludināto Polijas austrumu daļu sarunu valodā, ņemot vērā Vīnes kongresa lomu tās izveidē, dēvēja arī par ``kongresa Poliju'' (\pltxti{Królestwo Kongresowe}) un ``kongresovku'' (\pltxti{Kongresówka}). Tajā ietilpa apmēram 1/7 agrākās Žečpospolitas teritorijas ar 1/5 iedzīvotāju. Karaliste dalījās 8~vojevodistēs: Augustovas (\pltxti{Augustowska}), Kališas (\pltxti{Kaliska}), Krakovas (\pltxti{Krakowskie}, pati Krakovas pilsēta tajā neietilpa, bet vojevodistes nosaukums deva mājienu, ka ar laiku arī pilsēta varētu pievienoties tai), Ļubļinas (\pltxti{Lubelskie}), Mazovijas (\pltxti{Mazowiecki}), Plockas (\pltxti{Płockie}), Podļesjes (\pltxti{Podlaskie}), Sandomiras (\pltxti{Sandomierskie}).

Runājot par bijušās Polijas apgabaliem Krievijas varā, noteikti jāievēro, ka Polija XIX gadsimtā ietilpa Krievijas impērijas rietumu nomalē. Iedzīvotāju blīvums un skaits šeit bija krietni lielāks nekā Iekškrievijā. Viņu sociāli ekonomiskās un kultūras attīstības līmenis redzami pārsniedza ne tikai citu~--- austrumu, dienvidu, ziemeļu~--- valsts nomaļu, bet arī Krievijas centrālo apgabalu iedzīvotāju attīstības līmeni. 1816.~gadā karalistē mita ap 3,3~miljoni iedzīvotāju. No tiem pēc valodas (ko šai gadījumā var uzskatīt par diezgan precīzu nacionālā sastāva atspoguļotāju) 75\% bija poļi, 10\% ebreji, 7,5\% vācieši, 5\% lietuvieši un 2,5\% ukraiņi. Pēc ticības 83,5\% bija katoļu, 2,5\% uniātu,10\% jūdaistu un 3,75\% protestantu. Pēc citiem datiem ap 73\% iedzīvotāju bija zemnieki, 12\% sīkpilsoņi, 7,5\% šļahtič, 7,5\% beztiesīgie ebreji. (Dati par ebrejiem ir atšķirīgi, jo viņi izvairījās no tautas skaitīšanas).

Vēl pirms Polijas otrās dalīšanas 1791.~gadā ar imperatores Katrīnas II dekrētu Krievijas teritorijā tika izsludināta nometinājuma josla (\rutxti{Черта оседлости}), kurā drīkstēja dzīvot ebreji. (Citur drīkstēja apmesties tikai kristītie ebreji, kā arī jūdaisma ticību saglabājušie 1.~ģildes tirgoņi, personas ar augstāko izglītību, reģistrētās prostitūtas, bijušās militārpersonas u.c.). Pēc otrās Polijas dalīšanas tajā ietilpa arī Latgale. Pēc trešās Žečpospolitas dalīšanas nometinājuma joslā ietilpa arī Viļņas un Grodņas guberņas, kur bija daudz ebreju. Nometinājuma josla aptvēra speciāli noteiktas pilsētas un miestiņus, jo ebrejiem laukos arī bija aizliegts dzīvot. Galīgu juridisku noformējumu nometinājuma josla ieguva ar 1804.~gadā izdotajiem ``Noteikumiem ebreju iekārtošanai'' (``\rutxti{Положение об устройстве евреев}''), kas uzskaitīja tās guberņas un teritorijas, kurās ebreji drīkstēja apmesties un tirgot. Pēc Vīnes kongresa ebreju nometinājuma joslā nonāca arī visas Polijas karalistes guberņas.

Mūsdienu krievu sociālās vēstures pētnieks B.~Mironovs ir izdalījis vairākus Krievijas politikas principus tās inkorporētajās nacionālajās nomalēs.

\begin{enumerate}
\item Krievija ilgstoši tajās saglabāja to administratīvo kārtību, vietējos likumus un iestādes, zemes īpašuma attiecības, ticību, valodu un kultūru, kas pastāvēja pirms to pievienošanas Krievijai. Pagāja vairākas desmitgades, līdz pakāpeniski Viskrievijas kārtība tika attiecināta arī uz nomalēm. Cara valdība plaši sadarbojās ar vietējām feodālajām elitēm, kuras saņēma tādas pat tiesības kā krievu muižnieki. Tas centrālajai varai atviegloja jauno teritoriju pārvaldīšanu.

\item Nekrieviem tika radītas dažas priekšrocības salīdzinājumā ar krieviem. Tautas, kuras nepazina dzimtbūšanas attiecības, ar pievienošanu Krievijai tā arī neuzzināja, kas tā tāda. Ar nodokļu sistēmas palīdzību, kur nekrievu tautas maksāja mazākus nodokļus, tika radīts stāvoklis, ka nacionālajās nomalēs dzīvojošo nekrievu materiālais stāvoklis bija relatīvi labāks nekā krievu stāvoklis. (To atspoguļoja tāds apkopojošs labklājības rādītājs kā vidējais augums. Tā 1874.--1883.~gadā armijā iesaucamo vidējais augums no 28 lielkrievu guberņām bija 1~620~mm, bet no 22 pārējām guberņām~--- 1~627~mm). Teiktais nenozīmē, ka nekrievi neizjuta carisma spaidus, taču krieviem šai ziņā bija dažas ``priekšrocības''.

\item Etniskie kritēriji nebija noteicošie virzībā pa sociālajām kāpnēm. Krievijas politiskā, militārā, kultūras un zinātnes elite bija daudznacionāla. Uzticība tronim, dižciltība un profesionālisms tika vērtēti augstāk kā etniskā vai konfesionālā piederība. Ierēdņu vidū 1850.~gadā nekrievu bija 16\%, augstākās birokrātijas sastāvā 1853.~gadā nekrievu bija 32,7\%, 1917.~gadā~--- 11,8\%. Armijas virsnieku vidū 1912.~gadā nepareizticīgo, tātad cittautiešu, bija 11\%.
\end{enumerate}

Kopumā šo principu darbību varēja vērot arī Polijas karalistē, taču ar vienu būtisku papildinājumu, uz kuru norādījis arī pats B.~Mironovs. Ja Krievijai pievienotās tautas demonstrēja lojalitāti centrālajai varai, tad to zināmā autonomija palielinājās (kā Somijai XIX gadsimtā), ja tās izrādīja naidīgumu vai separātismu~--- tā tika atņemta. Polijas karaliste bija visspilgtākais piemērs, kā notikumi attīstījās pēc otrā varianta.

Īpaši jāuzsver, ka politisko procesu intensitāte Polijas karalistē bija neizmērojami augstāka nekā citās nomalēs. Ekonomiski atrašanās milzīgās impērijas sastāvā bija izdevīga Polijas karalistes rūpniecības attīstībai. Karalistes tautsaimniecība attīstījās krietni straujāk nekā visas Krievijas saimniecība kopumā. Taču pēc diviem nesekmīgiem sacelšanās mēģinājumiem poļiem tika atņemtas daudzas politiskās tiesības, ieviesti daudzi ierobežojumi. Tas, tāpat kā atmiņas par Žečpospolitas slaveno pagātni un arī konfesionālās atšķirības, poļu tautā izraisīja pastāvīgu neapmierinātību par atrašanos Krievijas impērijas sastāvā.

Poļu sabiedrība dažādi vērtēja situāciju un reaģēja uz laikmeta izaicinājumiem. Jau XVIII gadsimta beigās poļu dzejnieks K.~Kozmjans par poļu dzīvi Krievijas rietumu guberņās rakstīja: ``No zināma skatu punkta mēs šeit dzīvojam labāk nekā republikas [domāta Žečpospolita~--- V.Š.] laikā; mēs lielā mērā saglabājam to, ko mums deva dzimtene. Mums tagad nenākas baidīties no Umaņas slaktiņa (\pltxti{Rzeź humańska,}~--- 1768.~gadā sacēlušos ukraiņu zemnieku veikta Umaņas pilsētas un tās apkārtnes iedzīvotāju masu slepkavošana, kurā pēc dažādiem datiem gāja bojā no 12 līdz 20 tūkstošu ebreju, poļu un ukraiņu-uniātu) atkārtošanās; kaut Polijas vairs nav, mēs dzīvojam Polijā, un mēs~--- poļi''. Taču vairākums poļu šļahtiču sapņoja par nedalītas savas varas atjaunošanu neatkarīgā Polijā.

Sākotnēji Krievijas poļu zemes pārvaldīja Pagaidu augstākā padome, radīta 1813.~gada 1.~martā, Polijā ienākot Krievijas karaspēkam. Taču Pagaidu padome nebija vienīgā pārvaldes iestāde. 1814. gada beigās par poļu armijas virspavēlnieku tika iecelts Krievijas cara brālis lielkņazs Konstantīns, kurš faktiski pildīja vietvalža funkcijas. Pēdējās Vīnes kongresa dienās~--- 1815.~gada 22.~maijā tika parakstīti ``Polijas karalistes Konstitūcijas pamati''. Dokumentā bija uzsvērta Polijas karalistes saistība personālā ūnijā ar Krieviju, kā arī tas, ka tās Konstitūcijai jābalstās uz vēl pirms trešās Žečpospolitas dalīšanas 1791.~gadā 3.~maijā pieņemtās Konstitūcijas principiem.

Jaunās \strong{Polijas karalistes Konstitūcijas} sagatavošanā aktīvi piedalījās Krievijas imperatoram Aleksandram I tuvu stāvošais bijušais Krievijas ārlietu ministrs (1804--1806) kņazs Ā.~Čartorijskis, grāfs L.~Plāters un citi poļu magnāti. Taču tā paredzēja pārāk lielas tiesības nacionālajam Seimam, un Aleksandrs I lika tās sašaurināt, vienlaikus paplašinot monarha tiesības. Ā.~Čartorijskim Aleksandrs I gan to skaidroja ar negatīvo noskaņojumu pret poļiem Krievijas sabiedrībā: ``Poļu armijas uzturēšanās veids pie mums, laupīšanas Smoļenskā un Maskavā, visas valsts nopostīšana ir atdzīvinājušas agrāko naidu''. Tomēr jaunā Polijas karalistes Konstitūcijas noteiktā iekārta bija liberālāka nekā savulaik Varšavas hercogistē pastāvējusī.

Pēc labojumiem 1815.~gada 17.novembrī Aleksandrs I Konstitūciju parakstīja. Tā atzina Polijas karalistes mūžīgu saistību ar Krievijas impēriju, abu ārējo politiku realizēja Ārlietu ministrija Pēterburgā. Katoļu reliģija bija pielīdzināta pārējām, taču atradās īpašā valdības aizbildniecībā.

Pēc Konstitūcijas izpildvaru īstenoja karalis (vienlaikus Krievijas imperators), likumdošanas varu dalot ar Seimu. Karalim piederēja likumdošanas iniciatīva un \latxti{veto} tiesības. Viņš varēja pasludināt karu, slēgt līgumus, iecelt visus ierēdņus.

Pasīvās vēlēšanu tiesības tika piešķirtas visiem (izņemot ebrejus, kuriem politisko tiesību nebija), kas maksāja ne mazāk kā 100~zlotu tiešo nodokļu un bija sasnieguši 30 gadu vecumu. Aktīvās vēlēšanu tiesības baudīja šļahtiči~--- zemes īpašnieki no 21~gada vecuma, kā arī garīdznieki, skolotāji, amatnieku darbnīcu īpašnieki, nomnieki un tirgotāji, kuru īpašumā bija preces 10~tūkstošu zlotu vērtībā. Kopā vēlēšanu tiesības ieguva ap 3\% no iedzīvotājiem. (Salīdzinājumam, Francijā 1820.~gadā vēlēšanu tiesības bija tikai 80~000, bet Polijas karalistē~--- 100~000 iedzīvotāju. Tas gan bija saistīts nevis ar imperatora demokrātismu, bet ar muižnieku lielāku īpatsvaru Polijā nekā Francijā.) Seimā pastāvēja divas palātas. Augšpalātu (\pltxti{Izba Senatorska}) nozīmēja monarhs, apakšpalātā (\pltxti{Izba Posolska}) bija 128 pakāpeniskās vēlēšanās ievēlēti deputāti. 77~--- no dižciltīgajiem, 51~--- no pilsētu un lauku kopienām.

Pārvaldes priekšā atradās imperatora iecelts vietvaldis. Pirmais vietvaldis bija jau minētais ģenerālis J.~Zaijončeks~--- bijušais jakobīnis, kurš bija dienējis Napoleona armijā, taču pēc tam uzticīgi kalpoja Krievijai. Vietvaldim palīgā tika nozīmēta Valsts padome (\pltxti{Rada Stanu}), un Administratīvā (jeb Pārvaldes) padome (\pltxti{Rada Administracyjna}). Valsts padome, sastāvoša no ministriem, padomes locekļiem, un monarha īpaši uzaicinātām personām, sastādīja likumprojektus, apsprieda atsevišķu ministriju ziņojumus, risināja juridiskus strīdus u.tml. Administratīvajā padomē, kuras rokās atradās izpildvara, ietilpa ministri un atkal monarha īpaši uzaicinātas personas. Tika nodibinātas piecas komisijas (ministrijas): Kara, Tieslietu, Iekšlietu, Finanšu, Ticības lietu un Tautas izglītības. Ministri parasti nāca no poļu aristokrātijas vidus. Karalistes teritorija tika sadalīta astoņās vojevodistēs: Augustovas, Kališas, Krakovas, Ļubļinas, Mazovijas, Plockas, Radomas un Sandomežas (\pltxti{Województwo augustowska, kaliskie krakowskie, lubelskie, mazowieckie, płockie, radomske, sandomierskie}).

Krievijas armijas sastāvā atradās Poļu korpuss, kurš saglabāja savu formu. Miera laikā pastāvēja 30--35~000 vīru armija, to drīkstēja izmantot dzimtenes aizsardzībai tikai karalistes robežās un tās lielums bija atkarīgs no karalistes budžeta. Izdevumi armijas uzturēšanai sastādīja lielu budžeta daļu. Polijas armijas apbruņojumu pilnībā ražoja Krievijas fabrikās. Saglabājās arī poļu ordeņi: Baltā ērgļa (\pltxti{Order Orła Białego}), Svētā Staņislava (\pltxti{Order Świętego Stanisława}) un Kara nopelnu (\latxti{Virtuti Militari}). Par kara ministru un poļu armijas virspavēlnieku tika nozīmēts jau minētais cara brālis lielkņazs Konstantīns, ne visai psihiski nosvērts cilvēks, kurš iejaucās arī citu ministru kompetencēs, jo ģenerālis J.~Zaijončeks bija vāja rakstura cilvēks, nevēlējās ķildoties ar lielkņazu. Lielkņazs pat apguva poļu valodu, kaut karalistē šļahtas ikdienas sarunu valoda bija franču. Kad Aleksandrs I 1819.~gadā Konstantīnam pakļāva arī ārpus Polijas karalistes esošo Krievijas rietumu guberņu karaspēku, poļos pat radās utopiskas cerības, ka varētu notikt ``atkalapvienošanās'' ar šiem kādreiz Žečpospolitas sastāvā esošajiem apgabaliem.

Konstitūcijā tika deklarēta personas neaizskaramība, ticības un preses brīvība, tiesu neatkarība, poļu valoda tika atzīta par oficiālu karalistes valodu, kuru lietoja arī administrācija, tiesas un armija. Visi valsts amati tika piešķirti poļiem. Katram nākamajam Krievijas imperatoram bija atsevišķi jākronējas ar Polijas karalistes kroni, jāzvēr uzticība Polijas Konstitūcijai. Neraugoties uz Aleksandra I ieviestajiem ierobežojumiem, kas pavēra plašas iespējas Konstitūcijas garantiju vājināšanai, Polijas Konstitūcija tobrīd bija viena no pašām liberālākajām Eiropā. 1818.~gadā imperators Aleksandrs~I, personīgi atklājot Seima sēdi, paziņoja: ``Iepriekšējā zemes organizācija man ļāva ieviest to, kuru es jums dāvāju, iedarbinot liberālas iestādes. Šīs pēdējās vienmēr būs manu rūpju objekts, un es ar Dieva palīgu ceru izplatīt to ietekmi uz visām zemēm, kuras ar Dieva gribu dotas man pārvaldē.'' Seims pieņēma visus likumprojektus, kurus piedāvāja imperators, izņemot likumprojektu par civillaulību. Aleksandrs~I bija apmierināts ar tā darbību.

Pirmajos Polijas karalistes pastāvēšanas gados poļu varas iestādes baudīja iespēju noteikt karalistes sociāli ekonomiskās politikas virzību, reāli izmantoja administratīvo varu un naudas resursus. Taču samērā liberālās Konstitūcijas ieviešana drīz nonāca pretrunā ar Krievijas imperatora patvaldniecisko varu, viņa valdības un galma tieksmi tādu realizēt arī Polijas karalistē.

Pēc 1826.~gada īpašs vietvaldis karalistē vairs netika iecelts, tā pienākumus pildīja lielkņazs Konstantīns. Viņš neievēroja Karalistes Konstitūciju, iedzīvotāju tiesības. Tas izsauca lielu vietējo iedzīvotāju neapmierinātību. Tomēr lielkņazs Konstantīns vēlāk kļuva iecietīgāks, to uzskatīja par viņa ``polonozācijas'' sekām, jo viņu ietekmēja morganiskā sieva~--- poliete I.~Grudzinska, kura ieguva kņazienes Lovičas titulu. Precību ar viņu dēļ lielkņazs bija spiests atteikties no tiesībām uz Krievijas troņa mantošanu.

Polijas karalistē izmaiņas notika arī baznīcas organizācijā. 1818.~gadā pāvests noteica 8 katoļu bīskapiju un 1 uniātu bīskapijas robežas. Varšavas bīskaps kļuva par Polijas karalistes primasu (latīņu~--- \latxti{primas}~--- pirmajā vietā esošs, Romas katoļu baznīcā valsts galvenā bīskapa goda apzīmējums). Uniātu baznīcai piederēja 287 draudžu baznīcas ar 220~000 ticīgajiem, 4 klosteri ar 25 mūkiem un neliela seminārija.

Kopumā politiskā iekārta Polijas karalistē sākotnēji bija liberālāka nekā pārējā Krievijas impērijas daļā. Taču drīzumā cariskā administrācija sāka būtiski ierobežot konstitucionālās tiesības un brīvības. 1819.~gadā tika ieviesta cenzūra, 1821.~gadā bija spiests atstāt dienestu Polijas karalistes izglītības ministrs S.~Potockis, kura laikā tika atvērta Varšavas universitāte un ievērojami pieauga skolu skaits. 1825.~gadā, gatavojoties Seima sasaukšanai, tika pieņemts Konstitūcijas ``papildus pants'', kurš atcēla Seima sēžu atklātumu.

Žečpospolitas otrā un trešā dalīšana pārtrauca XVIII gadsimta beigu augšupejošo poļu \strong{ekonomikas} attīstības procesu, kara darbība poļu zemēs XVIII gadsimta sākumā atnesa sev līdzi lielus materiālus zaudējumus Polijai. Taču ekonomiskās attīstības rezultāts~--- feodāli dzimtbūtniecisko attiecību pagrimums noteica arī turpmāko Polijas karalistes ekonomisko attīstību. Sākotnēji ekonomikas jomā Polijas karaliste bija gandrīz neatkarīga no Krievijas. Saglabājās naudas vienība zlots (\pltxti{złoty}). Taču finanšu stāvoklis pirmajos gados bija smags. Karalistei bija jāmaksā agrāk izveidojušies parādi Austrijai un Prūsijai. Pirmajos gados karaspēka uzturēšanai Krievija piešķīra karalistei avansu, taču no 1817.~gada armijas uzturēšana pilnībā gūlās uz karalistes budžetu, bet jau izmaksātais avanss tika pievienots parādiem. Armijas pavēlnieks lielkņazs Konstantīns pieprasīja arvien lielākus izdevumus tās labā. Karalistes iestādes, ievācot nodevu parādus par iepriekšējiem gadiem, tos ieskaitīja tekošajā budžetā. Nolīdzināt iepriekšējo budžetu deficītu nekādi neizdevās. 1821.~gadā Aleksandrs I pat piedraudēja poļiem atņemt finansiāli-ekonomisko autonomiju, ja tie nespēs tā organizēt valsts kasi, lai uzturētu savu armiju.

Ekonomiskie procesi pirmajos Karalistes pastāvēšanas piecpadsmit gados ietvēra sevī divas strāvas. No vienas puses, tie balstījās uz Apgaismības laikmeta, Četrgadu Seima un Varšavas hercogistes laikā veikto reformu pamatiem. No otras puses daudzas norises ekonomikā iedīglī jau nesa jaunas kvalitatīvas izmaiņas, kuras sagatavoja pāreju uz jaunu kapitālistisko attīstības fāzi. Šim periodam bija raksturīga jaunu ražošanas formu~--- fabriku (no latīņu ~\latxti{fabrica}~--- darbnīca, rūpniecības uzņēmumu ar lielāku nekā darbnīcā strādnieku skaitu, parasti ražojošs preces nevis pēc pasūtījuma, bet tirgum), jaunu ražošanas attiecību~--- algota darba, jaunu sociālo grupu~--- algoto strādnieku, buržuāzijas veidošanās, veco ražošanas attiecību transformāciju jaunajās, kapitālistiskajās.

Ja XVI--XVIII~gadsimtā Žečpospolita Rietumeiropas valstīm piegādāja galvenokārt labību un kokmateriālus, XIX gadsimta sākuma kari izpostīja pastāvējušo darba dalīšanas sistēmu. Rūpniecības attīstība Rietumeiropā veda pie labības tirgus paplašināšanās un muižu izaugsmes uz zemnieku zemes rēķina Austrumeiropā. 1815.~gadā Polijas karaliste bija vēl agrāra zeme, 80\% iedzīvotāju dzīvoja no lauksaimniecībā gūtajiem ienākumiem, 70\% no nacionālā ienākuma deva zemes īpašums. Pilsētas bija attīstītas vāji. Lielākā bija Varšava ar vairāk nekā 100~000 iedzīvotāju. Pavisam Polijas karalistē bija ap 480 pilsētu, taču tikai 59 no tām iedzīvotāju bija no 3 līdz 10 tūkstošiem. 1816.~gadā pilsētās dzīvoja 19,4\% visu iedzīvotāju.

Varšavas hercogistē 1808.~gadā ieviestais Napoleona kodekss saglabāja savu spēku arī Polijas karalistē. Zeme atradās muižnieku, valsts un zemnieku īpašumā. Darbojās 1807.~gada Varšavas hercogistes Konstitūcijas pants par zemnieku atbrīvošanu no dzimtbūšanas.

Parasti muižniekiem nepiederēja mājlopi un zemes apstrādei nepieciešamais inventārs, viņu zemi klaušu kārtībā apstrādāja zemnieki ar savu inventāru. Muižnieki maz lietoja algotu darbu. No muižniekiem tikai mazai daļai (ap 4~000) piederēja lielas muižu saimniecības~--- latifundijas (latīņu \latxti{lātus} ``plašs'' + \latxti{fundus} ``ferma, nekustams īpašums''). No laukos dzīvojošajiem vairāk nekā 40~000 šļahtiču ap 4~100 nomāja muižas, vairāk nekā 32~000 sīko šļahtiču vai nu nebija savas zemes vai tie bija nelieli zemes gabali. Šie sīkie šļahtiči pēc ekonomiskā stāvokļa bija tuvi zemniekiem, tikai nepildīja feodālās saistības. Poļu vēsturnieks V.~Smoļenskis, kurš pētījis sīko šļahtiču dzīvi, rakstīja, ka viņu īpašumā parasti bija vairāki sīki skrejgabaliņi, un tāpēc tie bija atkarīgi no kaimiņiem, nespēja patstāvīgi saimniekot uz savas zemes.

Valstij piederēja t.s. valsts jeb nacionālās muižas, kuras bija izveidotas no bijušajām karaļu, klosteru, kā arī franču maršaliem un ģenerāļiem atsavinātajām zemēm.

Zemnieku zeme parasti vēl neatradās to pilnīgā privātīpašumā, bet gan bija aplikta ar dažādām feodālajām saistībām. (Klaušas sasniedza ne mazāk, bet atsevišķos gadījumos arī vairāk nekā 3 līdz 4 dienas nedēļā.) 1815.~gada Konstitūcija gan apstiprināja jau ar 1807.~gada Konstitūciju un Napoleona kodeksu pasludinātās zemnieku tiesības uz brīvību, bet muižnieku tiesības uz zemi, taču punkts par zemnieku tiesībām lielākoties bija fiktīvs, jo no 1818.~gada par vietējo tiesu iestāžu, kur nonāca zemnieku sūdzības, vadītājiem kļuva muižnieki. Pret zemniekiem, kuri atteicās pildīt feodālās saistības, vajadzības gadījumā tika sūtīts karaspēks. Vēl XIX gadsimta 20.~gados Polijas karalistē klaušu zemnieki sastādīja 51\% no visiem zemniekiem. Zemnieku pusnaturālo saimniecību pārvēršanās par preču ražotājām norisa visai gausi.

Sākotnējā kapitāla uzkrāšanās Polijas karalistē bija raksturīga ar lielāku nekā citās Eiropas valstīs muižnieku lomu šai procesā, ievērojamu valsts lomu, kas veicināja kapitāla uzkrāšanos lielo zemes īpašnieku un dažādo topošās buržuāzijas grupu rokās, ārējo avotu izpalikšanā. Sākotnēji izšķiroša loma kapitāla uzkrāšanā bija muižniekiem un valstij, mazāka~--- tirgotājiem un augļotājiem. Kā uzsvērusi krievu vēsturniece L.~Obušenkova XIX gadsimta pirmajā trešdaļā Karalistē sākotnējā kapitāla uzkrāšana rūpniecībā atpalika no tā uzkrāšanas tirdzniecībā, bet pēdējā~--- no tā uzkrāšanas muižnieku rokās lauksaimniecībā. Folverki (no vācu \detxti{Vorwerk}, poļu \pltxti{folwark}) jeb lopu muižas, pusmuižas kļuva par preču lielražošanas centriem, pārvērtās par puskapitālistiskiem, pusfeodāliem uzņēmumiem. Kapitāla uzkrāšana lauksaimniecībā nozīmēja galvenokārt tā uzkrāšanu to muižnieku rokās, kuru saimniecības sāka evolucionēt pa kapitālisma ceļu. Notika tas galvenokārt divos veidos~--- padzenot zemniekus no viņu apstrādātās zemes, lai to pievienotu muižai, un izpārdodot valsts muižas. Līdz 1830.~gadam bija pārdotas ap 150 valsts muižas, kas veda pie muižnieku zemes īpašuma pieauguma. Tiesa, bija vairāki gadījumi, kad valsts muižas pārdeva arī zemniekiem. Tad viņi ieguva šo zemi pilnā privātīpašumā. Muižnieki kopumā zemi zemniekiem nevēlējās pārdot. Tiesa, var pieminēt gadījumu, kad jau agrāk minētais poļu filozofs, rakstnieks un mecenāts S.~Stašičs pirms nāves novēlēja savas zemes sadalīt zemniekiem. Šis eksperiments kļuva par savdabīgu pieminekli XVIII gadsimta otrajā pusē tā arī neīstenotajiem sociālajiem ideāliem. Taču tas bija izņēmuma gadījums.

1827.~gadā ap 30\% no visiem lauku iedzīvotājiem bija bezzemnieki. Ja tiem pieskaita vēl sīku zemes gabalu turētājus, no kuriem to ģimenes nevarēja pārtikt, šis skaitlis pieauga līdz 50\%. Bezzemnieki tika nodarbināti muižās kā algota darba strādnieki.

Poļu historiogrāfijā ir izteikti divi galvenie viedokļi par kapitālistisko attiecību iespiešanos Polijas karalistes lauksaimniecībā. Profesors A.~Grodeks uzskatīja, ka XVIII gadsimta beigās un XIX gadsimta pirmajā pusē zemnieku padzīšanai no zemes nebija nekā kopīga ar kapitāla pirmatnējo uzkrāšanu, ka zemnieku noslāņošanās bija ar feodālu, nevis kapitālistisku raksturu. Tomēr vairākums poļu vēsturnieku, kuri pētījuši agrārās problēmas Polijas karalistē, tam nepiekrīt. Viņi muižnieku realizēto zemnieku padzīšanu no zemes, viņu zemes pievienošanu muižām saista ar kapitālistisko attiecību veidošanos lauksaimniecībā.

Norisa izmaiņas lauksaimnieciskās ražošanas struktūrā. Pieauga mājlopu skaits, palielinājās sugas lopu audzēšanas gadījumi. Tā, ja 1822.~gadā karalistē bija 1~527~tūkstošu smalkvilnas aitu, tad 1827.~gadā~--- jau 2~477~tūkstoši. Kopumā tomēr lopkopība atradās zemā attīstības līmenī. 1827.~gadā uz vienu lauku iedzīvotāju bija 0,129~zirgi un 0,463~ragulopi. Vasarās mājlopus izveda ganībās, bet ziemās baroja ar salmiem. Sienu zirgiem un vēršiem deva tikai smagāko lauku darbu periodos.

Viens no svarīgākajiem lauksaimniecības attīstības faktoriem bija kartupeļu audzēšana, kam grūdienu deva 1816.--1817.~gadu neražas. Kartupeļu un lopbarības kultūru iekļaušanu sējas apritē veicināja pāreju no trīslauku uz daudzlauku sistēmu, jaunas lauksaimniecības tehnikas izmantošanu. Lielus ienākumus muižniekiem deva degvīna dedzināšana. 1829.~gadā tika uzcelta pirmā cukurfabrika.

Tomēr līdz 1830.~gadam kapitālisma attīstība lauksaimniecībā tikai vēl sākās. Galvenā feodālās rentes forma bija klaušas. Straujāk kapitālisma virzienā evolucionēja muižu saimniecības Polijas karalistes rietumos.

Viens no kapitāla uzkrāšanas avotiem bija tirgotāju un augļotāju kapitāla uzkrāšana un pārvēršanās par rūpniecisko kapitālu. Pirmajos Polijas karalistes pastāvēšanas gados sākotnējā kapitāla uzkrāšanās procesu negatīvi ietekmēja pilsētu un vietējā tirgotāju kapitāla vājums. Lieltirgotāji koncentrējās galvenokārt Varšavā, citās pilsētās pārsvarā bija sīktirgotāji, augļotāji un pārpircēji.

Vīnes kongresā bija pieņemts lēmums, ka starp poļu zemēm, kuras pēc dalīšanām bija nonākušas dažādu valstu sastāvā, jāpastāv brīvai tirdzniecībai, taču šis lēmums tika ievērots tikai pirmajos gados. Pēc Polijas karalistes nodibināšanas tās ekonomiskajai attīstībai pavērās divi ceļi: 1) frītrēderisma (angļu \entxti{free trade}~—brīvā tirdzniecība, politika, kas atbalsta brīvu preču un pakalpojumu kustību pāri administratīvajām robežām un valsts neiejaukšanos privātražošanas sfērā) un 2) protekcionisma (no latīņu \latxti{protectio}~--- aizbildniecība, aizsardzība; iekšējā tirgus aizsardzības politika pret konkurenci ar zināmu ierobežojumu: importa un eksporta nodevu, subsīdiju u.c. līdzekļiem). Attīstības ceļa izvēlē izšķiroša bija valsts loma.

Frītrēderisma politiku aizstāvēja Polijas karalistes iekšlietu ministrs grāfs T.~Mostovskis. Viņš cerēja izmantot poļu zemju ģeogrāfisko stāvokli un atjaunot Polijas kā starpnieka lomu tirdzniecībā starp Rietumeiropu un Tuvējiem austrumiem. Tāpēc viņš uzstājās kā Varšavas gadatirgu iniciators. Par to noturēšanu divreiz gadā~--- pavasarī un rudenī~--- vietvaldis paziņoja 1817.~gada februārī. Gadatirgi tika uzskatīti par līdzekli kā likvidēt budžeta deficītu. 1818.~gadā uz gadatirgiem ieradās vairāk nekā 200 ārzemju tirgotāju. Brīvās tirdzniecības politikas realizācijā bija ieinteresēti arī muižu īpašnieki kā labības pārdevēji un kā ar muitas nodevu neapliktu ārzemju preču pircēji. 1816.--1819.~gadā Polijas karalistē pastāvēja brīvās tirdzniecības sistēma. Muitas tarifs vadījās tikai no valsts kases fiskālajām interesēm.

Taču nacionālās tautsaimniecības attīstībai tolaik vairāk atbilstoša bija protekcionisma politika. Enerģiski pret brīvās tirdzniecības politiku iestājās Polijas tirgotāji. Īpaši viņi iebilda pret privilēģiju piešķiršanu ārzemju tirgotājiem Varšavas gadatirgos.

1820.~gadā Polijas karaliste tika iekļauta Krievijas impērijas muitas robežās. Faktiski cara valdība neiebilda, ka arī poļi pāriet uz šādu sistēmu. Tādejādi kā Krievijā, tā Polijas karalistē tika ieviesta vienota muitas sistēma. Muitas robežas atcelšana starp Polijas karalisti un Krieviju radīja noieta tirgu poļu precēm, Tās bija kvalitatīvākas nekā dzimtbūtnieciskajā Krievijā izgatavotās, sāka strauji iekarot Krievijas (daļēji arī Ķīnas) tirgu. Taču, no otras puses, Krievijas varas iestādes pārliecinājās, ka Polijas karalistes iekļaušana vienotā muitas sistēmā un vienotas muitas robežas noteikšana Krievijai nav izdevīga. Krievijas rūpniecībai nebija izdevīga konkurence ar rūpniecības ražojumiem, kurus varēja ievest bez muitas no Austrijas un vācu valstiņām vispirms Polijas karalistē, bet no turienes arī Krievijā. Turpretī rūpniecības attīstība pašā Polijas karalistē Krievijai zināmā mērā bija izdevīga, jo tā sākotnēji netraucēja Krievijas rūpniecības attīstību. Krievijas tirgus pat bija ieinteresēts saņemt dažus poļu tekstilrūpniecības izstrādājumus. Muitas nodevas lielumu noteica krievu ierēdņi, pirmkārt ievērojot Krievijas rūpniecības intereses, kuras tika noskaidrotas pastāvīgos strīdos.

1820.--1821.~gadā pirmo vietu Polijas karalistes tirdzniecībā ieņēma Prūsija, otro un trešo~--- Krievija un Austrija. Karalistei ārējā tirdzniecībā bija negatīvs balanss, vairāk preču tika ievests nekā izvests. Tā bija pārpludināta ar ārzemēs ražotajām tekstila un metālapstrādes precēm. 1820.~gadā karalistē tika pārdots ārzemēs ražoto audumu par 24~miljoniem zlotu, bet vietējo, poļu ražoto~--- par 1 miljonu zlotu, ārzemēs ražotais metāls un tā izstrādājumi tika pārdots par 3~miljoniem zlotu, bet poļu metāls un ražojumi~--- par 707~tūkstošiem zlotu. Tādejādi zuda apstākļi tekstil un metālrūpniecības attīstībai pašā Polijas karalistē.

Polijas karalistes valdības politika strauji mainījās, kad par tās finanšu ministru kļuva F.~Druckis-Ļubeckis, kurš bija konsekvents protekcionisma piekritējs, un ar jauna muitas tarifa ieviešanu Krievijā 1822.~gada martā, kurš Polijas karalistei atjaunoja muitas autonomiju. Karalistē tika ieviests savs protekcionistisks muitas tarifs, kurš noteica augstas muitas nodevas kā no ārzemēm ievedāmajai lauksaimniecības produkcijai, tā rūpniecības precēm un greznumpriekšmetiem. Toties tirdzniecībā starp karalisti un Krieviju muitas nodevas bija krietni mazākas. Izejvielu un neapstrādātu izstrādājumu tirdzniecībā starp tām nodevas vispār tika atceltas.

Šis tarifs aizsargāja veidojošos vietējo rūpniecību no lēto austriešu un vācu preču konkurences, bet muitas nodevas kļuva par vienu no galvenajiem budžeta ienākumu avotiem. F.~Druckim-Ļubeckim izdevās faktiski aizliegt vilnas rūpniecības izstrādājumu, īpaši vadmalas audumu, ievešanu karalistē un izmantot lielās austrumu tirgus iespējas poļu vilnas tekstilizstrādājumu pārdošanai. Poļu tirgotāji un rūpnieki guva ievērojamu peļņu no muitas tarifu samazināšanas ar Krieviju. Tā, 1828.--1830.~gadā par vadmalas ievešanu Krievijā pēc atvieglotā tarifa viņi samaksāja 1,4~miljonus rbļ. muitas nodevu. Ja tiktu lietots parastais ārējās tirdzniecības, bet nevis atvieglotais, Polijas karalistei piemērojamais muitas tarifs, viņiem nāktos samaksāt nodevās 74,54 miljonu rbļ.

Pāreja uz protekcionismu gan izraisīja t.s. muitas karus ar Prūsiju (1823--1825) un Austriju (1824--1828). Šo valstu (īpaši Austrijas) daļa Polijas karalistes tirdzniecībā samazinājās, toties auga tirdzniecība ar Krieviju. Tiesa, Polijas karalistes tirdzniecības apjoms ar rūpnieciski attīstīto Prūsiju joprojām bija lielāks nekā ar Krieviju. 1826.~gadā pirmo reizi Polijas karalistes tirdzniecības bilance ar Krieviju kļuva aktīva. Muitas sistēma, kad uz robežas ar Prūsiju un Austriju muitas nodevas bija augstas, bet ar Krieviju~--- zemas, Polijas uzņēmējiem bija izdevīga, deva papildus priekšrocības. Kā raksta jau minētā krievu vēsturniece L.~Obušenkova, abu valstu: Krievijas un Polijas karalistes muitas tarifi veicināja lauksaimniecības un tekstilrūpniecības attīstību karalistē. Liela nozīme bija arī brīvam poļu preču tranzītam uz Krievijas ostām pie Baltijas jūras, arī ārzemju preču ievešanai Polijas karalistē caur Liepāju, Rīgu un Ventspili. Mūsdienu poļu vēsturnieki atzīst, ka no Krievijas puses tai laikā nebija Polijas ekonomiskas ekspluatācijas.

F.~Druckis-Ļubeckis pilnveidoja arī valsts monopoltiesību sistēmu. Vislielākos ienākumus deva tabakas, pēc tam sāls monopols. Valsts savas monopoltiesības pārdeva atsevišķiem tirgotājiem, kuri par tām maksāja lielu naudu. Tā 1816.--1821.~gadā valsts ienākumi no sāls monopoltirdzniecības sastādīja 5 miljonus zlotu, bet, kad F.~Druckis-Ļubeckis ķērās pie tās reorganizācijas,~--- 14~miljonus zlotu.

Tiesa, šļahtiči bija nemierā ar 20.~gados realizēto nodokļu politiku. Absolūtos skaitļos tie gan nelikās lieli. Salīdzinājumam, ja Polijas karalistē 1830.~gadā vienam iedzīvotājam gadā vidēji nodokļos bija jānomaksā 26,25~franču franki, tad Francijā~--- 33, bet Anglijā~--- 65~franki. Te gan jāņem vērā, ka ekonomiskās attīstības rādītājos karaliste ievērojami atpalika no nosauktajām Rietumu valstīm.

Ja no XIX gadsimta beigām par rūpniecības attīstību poļu vēsturnieku vidū ilgstoši norisa strīdi, vai tā nav mākslīgi ienesta Polijas karalistē, tad mūsdienās viennozīmīgi valda priekšstats, ka tai bija savi iekšējie priekšnoteikumi. Poļu zemēs tradicionālās rūpniecības nozares bija tekstilrūpniecība, kalnrūpniecība un metalurģija.

Pirmās, vēl pie muižām dibināmās manufaktūras Polijā parādījās jau XVIII gadsimta sākumā. Gadsimta beigās~--- 80. un 90.~gados parādījās jau karaļa un pilsētnieku īpašumā esošas manufaktūras. Tomēr rūpniecības attīstības manufaktūru periods kā tāds Polijas karalistē sākās XIX gadsimta 20.~gados, bet noslēdzās 40. vai, pēc citiem vērtējumiem, 50.~gados. Tad Polijas karalistē rūpniecības attīstība varēja balstīties uz jau pastāvošajām tradīcijām.

Taču vietējo iedzīvotāju vidū bija daudz mazāk preču patērētāju nekā Rietumeiropas valstīs. Pēc pētnieku domām, diezgan ievērojama pirktspēja bija tikai 10\% Polijas karalistes iedzīvotāju: lielajiem zemes īpašniekiem, muižu nomniekiem, augstākās garīdzniecības, ierēdniecības un virsniecības pārstāvjiem, nelielai lielāko tirgotāju grupai. Bez tam no tirgus bija atkarīgi ierēdņi un kalpotāji, brīvo profesiju pārstāvji, algotā darba strādnieki, pilsētnieku zemākie slāņi. Zemnieki pirka lauksaimniecības inventāru, apģērbu, sadzīves priekšmetus. Liela loma preču apgrozībā bija valsts iepirkumiem gan armijas, gan citām vajadzībām (ceļu un kanālu, sabiedrisko ēku celtniecībai iepērkamajiem būvmateriāliem, instrumentiem u.c.).

Valsts veicināja rūpniecības attīstību ar kvalificētu speciālistu piesaisti. Katrs ārzemju (galvenokārt vācu) amatnieks vai zemnieks, kurš pārcēlās uz Polijas karalisti, tika uz sešiem gadiem atbrīvots no visiem nodokļiem, karaklausības, nemaksāja muitu par savas kustamās mantas pārvešanu. Dati par iebraucēju skaitu no 1815 līdz 1830.~gadam gan ir atšķirīgi~--- pētnieki to skaitu vērtē robežās no 20 līdz pat 300 tūkstošiem. Atbraucēju vidū lielākoties bija dažādu specialitāšu tekstilamatnieki, to vidū bija daudz sīkuzņēmēju, kuri izjuta finansiālas grūtības dzimtenē, arī zeļļu bez darba rīkiem.

Liela nozīme rūpniecības attīstībā bija satiksmes ceļu attīstībai~--- šoseju un ūdens ceļu celtniecībai. Ar transporta sistēmas reorganizāciju nodarbojās vietējās poļu varas iestādes, tādēļ pirmajā vietā tika stādītas vietējās tautsaimniecības un tirdzniecības attīstības intereses. 1824.~gadā tika aprēķinātas izmaksas ap 10~t (precīzi~--- 98~280 kg) akmeņogļu piegādei no Dombrovas uz Varšavu. Šī daudzuma ogļu ieguve izmaksāja ap 1 tūkstoti zlotu, bet transportēšana~--- 3 tūkstošus, t.i.~--- bija trīs reizes dārgāka. Tāpēc saprotams, ka transporta problēma uztrauca kā karalistes vietējās, tā centrālās varas iestādes. Pirmkārt ceļu attīstībā bija ieinteresēta valdošā šķira. Taču to būvniecība prasīja lielus kapitālieguldījumus, kuri vai nu vispār neatmaksājās vai atmaksājās ļoti ilgā periodā. Tāpēc izdevumus uzņēmās valsts. Lai atvieglotu ekonomiskos sakarus starp augošajiem rūpniecības centriem tika būvēti ceļi un kanāli. 1829.~gadā Polijas banka piešķīra kredītus bruģētu ceļu~--- t.s. traktu izbūvei starp lielākajām pilsētām. Tika atjaunoti arī vecie ceļi. Augošā poļu buržuāzija pamatīgi nopelnīja no ceļu būvniecības. Pelnīja arī muižniecība. Blakus nospraustajām ceļu stigām esošajās muižās tika pirkti būvmateriāli, arī klaušinieku darbaspēks, kuri nebūt necentās panākt labāko darba kvalitāti. Tomēr valsts nodrošināja savas stratēģiskās intereses, muižnieki ieguva labus ceļus savas produkcijas nogādei tirgos, uzņēmēji guva peļņu, akumulēja kapitālu. Līdz 1830.~gadam tika uzbūvēti 1~030 km jaunu bruģētu ceļu, nerēķinot vecos, pastāvīgi atjaunojamos traktus. Tā rezultātā tika izveidota regulāra satiksme starp rūpniecības centriem, lauku novadiem un upju ostām.

Ja XVIII gadsimta beigās magnāti ceļoja karietēs, turpretī pilsētnieki savās tirdzniecības, bet ierēdņi dienesta darīšanās, šļahtiču atvases uz mācību iestādēm devās speciāli salīgta zemnieka vai ebreju sīktirgotāja ratos, tad ar ceļu izbūvi saistījās jauna transporta veida~--- diližansa rašanās. Tas pastāvīgi kursēja starp lielākajām pilsētām, bija salīdzinoši lēts, pārvadāja pasažierus, pastu, nelielus dārgu preču daudzumus.

Kuģojamās upes tika atbrīvotas no aizgruvumiem, aizsprostiem, dzirnavām. Tika būvēti kanāli. No tiem ievērojamākais bija Augustovas kanāls, kura būvei bija ne tikai ekonomisks, bet arī politisks aspekts. Muižnieki tā centās atbrīvoties no Prūsijas atkarības, kuras valdījumā atradās Vislas grīva. Kanāls savienoja Vislu ar Nemunu un caur Mēmeli (Klaipēdu) deva poļu precēm izeju uz jūru. Tā atklāšana notika 1839.~gadā. Tā garums sniedzās pāri 100 kilometriem. Bija izbūvētas 18 slūžas, daudz aizsprostu un tiltu.

1815.--1830.~gadā galvenā rūpniecības nozare bija tekstilrūpniecība. 1815.--1830.~gados jau bija vērojama pāreja no preču sīkražošanas pie manufaktūrām un pirmajām fabrikām. Visas tās nozares attīstījās uz algotā darbaspēka izmantošanas pamata. Savukārt tekstilnozares iekšienē pirmā vieta piederēja vilnas audumu ražošanai, otrā~--- linaudekla, bet kokvilnas audumu ražošana aizsākās tikai XIX gadsimta 20.~gados. Vājā poļu mašīnbūvniecība nespēja piegādāt nepieciešamās iekārtas. Polijā darbojošies tekstilrūpnieki mašīnas parasti ieveda no Anglijas. Kad tajā darbojās aizliegums uz to pārdošanu, imports notika kontrabandas ceļā caur Beļģiju. Tā mašīnu imports veicināja ražošanas spēku attīstību telstilrūpniecībā un celtniecībā, toties bremzēja mašīnbūvniecībā.

Vilnas audumu, galvenokārt vadmalas, produkcija rada noietu kā iekšējā, tā arī ārējā tirgū. Liels atspaids tekstilražošanai bija vadmalas valsts iepirkumi armijas mundieriem. Arī zemnieki nēsāja rupjvilnas vadmalas apģērbu. Poļu vadmala tika arī izvesta uz Krievijas guberņām un Ķīnu. Sakarā ar to arī attīstījās jau minētā aitkopība, tikai XIX gadsimta otrajā pusē aitu vilna no Austrālijas izspieda Eiropā ražoto.

Valsts iestādes palīdzēja ražotājiem sagādāt izejvielas, ierīkoja to noliktavas, subsidēja jaunu ražošanas iekārtu iegādi ārzemēs (ja 1812.~gadā darbojās tikai 100 aužamās stelles, tad 1830.~gadā~--- jau 5 tūkstoši), regulēja attiecības dažādu tekstilrūpniecības apakšnozaru vidū. Manufaktūru izveide radīja pamatu tālākai pārejai uz fabriku dibināšanu, pirmās no tām radās jau 20.~gadu beigās.

Attīstīt kaut cik plaši linaudekla ražošanu neizdevās. Armijas apgādē ar linaudeklu priekšroka tika dota augstākas kvalitātes audumam no Saksijas. No XIX gadsimta 20.~gadiem samazinājās linu audzēšana un linaudekla darināšana, to izspieda kokvilnas audumi. Kokvilnas auduma ražošana bija pati jaunākā tekstilnozare poļu zemēs. Pirmā kokvilnas audumus ražojošā fabrika tika uzbūvēta 1820.~gadā pie Varšavas. Pāreja uz kokvilnas kā izejvielas izmantošanu bija raksturīga kapitālistiskajai ražošanai, šī produkcija varēja būt sliktākas kvalitātes, bet tās priekšrocība bija lētums. Daudzie no savas saimniecības padzītie un ``stūrī'' iedzītie zemnieki bija slikti pircēji, jo tiem trūka naudas, taču tādu bija jau simti tūkstošu un pirka viņi lētāko kokvilnas audumu. Tieši kokvilnas audumi bija domāti iekšējam tirgum un pieprasījums pēc tiem izrādījās stabils. Kaut cara valdība, aizsargājot krievu ražotāju intereses, aizliedza ievest krievu guberņās poļu ražotos kokvilnas audumus un to izstrādājumus, iekšējā tirgus pieprasījums pēc lētiem audumiem bija tik liels, ka to ražošana attīstījās ļoti strauji, no 1825. līdz 1830.~gadam pieauga gandrīz 5 reizes. Toreizējā Mazovijas vojevodistē (\pltxti{Województwo Mazowieckie}) radās vairāki kokvilnas ražošanas centri. Ievērojamākais no tiem bija Lodza (\pltxti{Łódź}). 1820.~gadā ar lielkņaza Konstantīna rīkojumu tā tika pasludināta par fabriku pilsētu. Tai laikā Lodzā bija 112 namu ar 799 iedzīvotājiem. Taču jau 1829.~gadā Lodza bija kļuvusi par īstu pilsētu ar 4~273 iedzīvotājiem.

Poļu vēsturnieki raksta, ka Lodzas kā jauna tekstilrūpniecības centra uzplaukums sākās jau no XIX gadsimta 20.~gadiem, kad šeit attīstījās vadmalas un kokvilnas audumu ražotnes. Lodzas uzņēmēju ražojumus lielā daudzumā uzpirka Krievijas armijas apgādes dienests. Karalistes valdība parasti uzticējās lieluzņēmējiem un atbalstīja tos arī finansiāli. Sīkajiem uzņēmējiem bija grūti bez ārējas palīdzības paplašināt savu biznesu. Vācu vēsturnieki gan uzskata, ka Lodzas apgabala jeb, kā tas kādreiz tika saukts, ``Polijas Mančesteras'' stabila augšupeja sākās tikai XIX gadsimta vidū. Vācu zinātniskajā literatūrā uzsvērts, ka Lodzas rūpniecības īpatnība bija tā, ka šeit tekstilnozarē strādāja galvenokārt vācieši, kurus šurp bija aicinājuši poļu muižnieki. Vairākas pilsētiņas ar Lodzu priekšgalā XIX gadsimta pirmajā pusē apdzīvoja galvenokārt vācieši, kuri līdz 90\% bija luterāņi. Tikai kopš 1848.~gada fabriku pilsētās tika atļauts apmesties ebrejiem un turpmāk viņu skaits līdz XIX gs. beigām auga tāpat kā vāciešu skaits. Turpretī poļu tautsaimniecības vēsturnieks V.~Kula uzsvēris, ka Lodzas tekstilrūpniecība mantoja Lielpolijas un Silēzijas vadmalas ražošanas tradīcijas. Lodzas apgabalā ienāca emigranti no Lielpolijas, Ļubušas zemes (poļu \pltxti{Lubusz, ziemia lubuska}, vācu \detxti{Land Lebus}~--- vēsturiska novads Polijā un Vācijā abos Oderas upes krastos) un Lejas Silēzijas, lielā mērā ar poliskām etniskajām saknēm. Tas lielā mērā noteica ātro vācu strādnieku polonizāciju Lodzā.

Arī kalnrūpniecība un metalurģiskā rūpniecība nebija jaunas nozares. Poļu zemēs jau izsenis ieguva derīgos izrakteņus. Tie bija koncentrēti galvenokārt divu~--- Sandomiras un Krakovas vojevodistu teritorijā. Īpaši Dombrovas (\pltxti{Dąbrowa Górnicza}) apkārtnē atradās bagātas rūdu un ogļu iegulas. Dzelzs ieguvei bija liela loma jau senatnē, taču īpaši ievērojami pieprasījums pieauga līdz ar dažādu dzelzs mašīnu ražošanas sākumu. Jau XVIII gadsimta beigās Žečpospolitā darbojās 165 kalnrūpniecības uzņēmumi, taču XIX~gadsimta sākumā, neraugoties uz izjūtamo dzelzs trūkumu, dzelzs manufaktūras panīka sakarā ar nepietiekamo pieprasījumu pēc to produkcijas. Feodālisma laikmetā manufaktūras ražoja galvenokārt pirmās nepieciešamības preces vai pusfabrikātus (piemēram, dzelzs loksnes), no kuriem tālāk sīkākas darbnīcas izgatavoja gala produktu. Tālāk manufaktūras, pārtopot par rūpnīcām, no pusfabrikātu ražošanas pārgāja uz ražošanas līdzekļu ražošanu topošajiem spirta brūžiem, alus darītavām, cukurfabrikām, metalurģiskajām un tekstilrūpnīcām. Polijas karalistes valdība aizliedza jebkādas rūdas pārdošanu ārzemēm, jo tas atbalstītu konkurējošo (īpaši Saksijas) rūpniecību. 1816.~gadā tika panākta vienošanās par poļu metalurģijas ražojumu ievešanu Krievijā bez muitas. Tika atjauninātas vietējo uzņēmumu iekārtas. 1824.~gadā Polijas karalistē darbojās 37 dzelzsrūdas raktuves, 9 domnas, 6 liešanas fabrikas, 32 puddlinga krāsnis, kurās no čuguna ieguva dzelzi, 2 lokšņu dzelzs izgatavošanas cehi un vairāk nekā 100 sīkas darbnīcas.

Atsevišķi jāpasaka par akmeņogļu ieguvi. XVIII gadsimta beigās tās galvenokārt ieguva Austrijai piederošajos Javoznas (\pltxti{Jaworzno}) un Ščakovas (\pltxti{Szczakowa}~--- tur 1767.~gadā tika izveidota pirmā ogļu šahta Polijā) rajonos. Taču to ieguve, kaut ogļu slāņi atradās tuvu zemes virsmai un tika izmantots klaušu darbs, izrādījās neizdevīga lielo transporta izdevumu dēļ. Apstākļos, kad mežu vēl bija daudz un metalurģiskā tehnika vēl iztika ar kokoglēm, šahtas drīz tika pamestas. Taču, kad metalurģijas attīstība noteica akmeņogļu nepieciešamību lielākas temperatūras radīšanai domnās lai novērstu nevēlamus piejaukumus metālā, akmeņogles uzvarēja kokogles. Vispirms tās sāka plaši lietot metālapstrādē, cinka iegūšanā un kuģniecībā, vēlāk arī melnajā metalurģijā, uz dzelzceļa un ēku apsildīšanā. Vācu kapitālisti arī pirka oglēm bagātos zemes gabalus, taču nevis lai tos izstrādātu, bet gan lai neļautu tos izstrādāt citiem un tā mazinātu konkurenci Vācijā iegūtajām oglēm.

Par ievērojamu rūpnieciskās ražošanas centru attīstījās Varšava. 1825.~gadā tajā tika uzskaitītas 5~808 fabrikas un amatnieku darbnīcas. Polijas karalistē 1815.--1830.~gadā gūtie ekonomiskie panākumi apliecināja poļu nācijas dzīvotspēju, kaut arī tās etniskā teritorija bija saskaldīta.

Uz tautsaimniecības attīstības fona mainījās \strong{noskaņojums sabiedrībā}.

Sākotnēji ziņa par Polijas karalistes izveidi Vislas krastos tika saņemta ar eiforiju. Imperatoru Aleksandru I poļu intelektuāļi gandrīz uzreiz pasludināja par ``nācijas glābēju'' (\pltxti{zbawcą narodu}). Pateicība viņam kā jaunajam Polijas monarham tika izteikta dažādos veidos: par viņu tika sacerētas dzejas un dziedājumi, aizlūgts Dievs, viņa vārdā nosauca jaundzimušos bērnus, pie sienām tika izkārti litografēti imperatora~--- labdara portreti. Patriotiskā aizgrābtība aculieciniekos atsauca atmiņā līdzīgas ainas 1806.~gada decembrī, kad Napoleons iegāja no prūšiem atkarotajā Varšavā. Tiesa, francūži toreiz smagi piekrāpa radušās cerības, radot tikai necilo Varšavas hercogisti. Taču tagad poļi bija guvuši ko vairāk~--- Polijas vārds atgriezās politiskajās kartēs, troni ieņēma nevis kāds parasts hercogs, bet karalis!

Protams, jau no paša sākuma netrūka arī neapmierināto. Vienus tik ļoti bija savaldzinājusi Napoleona personība, ka viņi vēl ilgi nespēja atteikties no šī valdzinājuma. Citi, jau lielākā skaitā esošie, sūkstījās, ka Polijas karaliste teritorijas ziņā ir mazāka nekā bija Varšavas hercogiste, un to nekādi nevarēja salīdzināt ar Žečpospolitas plašajām teritorijām. Trešie ar aizdomām uzlūkoja dažus Aleksandra I praktiskos soļus.

Taču šī neapmierinātība, šaubas bija gandrīz nemanāmas izglītotās poļu sabiedrības laimes sajūtas priekšā: atkal pastāv Polija, atkal tā ir karaliste ar savu Konstitūciju, kura atsaucas uz 1791.~gada pamatlikumu, viss pārējais~--- sīkumi, kam nav vērts piešķirt īpašu nozīmi. Palika nepamanīts, ka Aleksandra I patvaldnieciskās reformas lielā mērā tikai izrietēja no cara varas tradīcijām. Neraugoties uz to, ka imperators pats deva poļiem Konstitūciju, viņš to neuzskatīja par sev saistošu dokumentu. Viņš arī drīz saprata, ka divu varas sistēmu: patvaldības un konstitucionālas monarhijas samierināšana nav viegla.

Zināmā autonomija, liberālā politiskā iekārta, progresīvāka nekā pašu krievu apdzīvotajā Krievijas daļā tomēr neapmierināja ievērojamu poļu tautas daļu. Tūkstošgadu valstiskuma tradīcija, daudziem sasniegumiem bagātā vēsture, katolicisms, zināma pārākuma apziņa pār iekarotājiem neļāva poļiem samierināties ar suverenitātes zaudēšanu. Reliģija sniedza poļiem drošu atbalstu nacionālajās brīvības cīņās. Viņi akcentēja savu konfesionālo piederību, tā lieku reizi uzsverot savu savdabību, pretstatot to apspiedējiem. (Tiesa, šis nacionālās atbrīvošanās kustības katoliskais raksturs neļāva tai gūt plašāku atbalstu pareizticīgajā krievu sabiedrībā.)

No otras puses, kā uzskata krievu vēsturnieks B.~Mironovs, pēc 1815.~gada Krievijas forsētā valstiskuma izbūve Polijā radīja stāvokli, kad šis valstiskuma radītājs~--- Krievija~--- poļiem vairs nebija vajadzīgs un viņi centās no tā tikt vaļā. Tiesa, viņu neapmierinātība bija virzīta ne tik daudz pret Aleksandru I personīgi, kā pret viņa ieceltajiem ierēdņiem: vietvaldi, kara ministru un poļu armijas virspavēlnieku, (vēlāk, pēc J.~Zaijončeka nāves arī vietvaldi) lielkņazu Konstantīnu, citiem ministriem.

Tālredzīgākie Krievijas valsts darbinieki paredzēja, ka poļu neapmierinātība var kārtējo reizi beigties ar sacelšanos.

Kad 1818.~gadā Varšavā notika svinības sakarā ar Seima atklāšanu divu krievu ģenerāļu starpā notika zīmīga saruna. I.~Paskevičs pajautāja grāfam A.~Ostermanam, kādas tam būs sekas. Pēdējais nedomājot atbildēja: ``Pēc desmit gadiem Tu ar savu divīziju dosies pret viņiem triecienā.'' Paredzētājs kļūdījās tikai par 3 gadiem.

Neapmierinātība auga arī Polijas karalistes armijā. Tā bija neliela, bet no Napoleona karu laika palikušo virsnieku skaits~--- salīdzinoši liels. Tāpēc izredzes sekmīgi kāpt pa karjeras kāpnēm tiem bija visai nelielas. Tas bija viens no cēloņiem kāpēc auga revolucionāri noskaņojumi virsnieku vidē. Lielu neapmierinātību poļu virsniekos izsauca lielkņaza Konstantīna nievājošā izturēšanās pret viņiem. Pirmajos četros gados gandrīz 50 virsnieku beidza dzīvi pašnāvībā, nevarot paciest aizvainojumus no Konstantīna puses. Daudzi virsnieki, kam bija citi iztikas avoti, atstāja dienestu. Virsnieku sašutumu izsauca arī nelikumības un Konstitūcijas neievērošanas gadījumi.

Par poļu aristokrātiju šai laikā poļu vēsturnieki V.~Rostockis un V.~Zajevskis raksta, ka tā pēc saviem politiskajiem nodomiem dalījās divās daļās. Veco aristokrātu dzimtu, kuras carisms bija atstūmis no valsts pārvaldes, pārstāvji bija noskaņoti opozicionāri pret carismu. Šīs daļas jūtu un interešu redzamākais izteicējs bija Ā.~Čartorijskis. Otra daļa poļu aristokrātijas, saistīta ar elementiem, kuri ar saviem kapitāliem piedalījās rūpniecības attīstībā, augstākā birokrātija (kā F.~Druckis-Ļubeckis) un armijā atstātie poļu ģenerāļi bija visstingrākie sacelšanās pret imperatoru pretinieki, perspektīvas tālākai attīstībai saskatīja vienīgi savienībā ar carismu. Abas daļas rīkojās atbilstoši saviem uzskatiem.

Nacionālās neatkarības kustības dalībnieki sākotnēji izmantoja tādu organizācijas formu kā masonu jeb brīvmūrnieku ložas. Jau Varšavas hercogistē (1807--1813) darbojās vairākas organizācijas, radušās uz masonu kustības pamata (masonu kustība ir saskaņā ar brālības principiem veidota organizācija, noteiktā mērā arī vispasaules savienība, kuras dalībniekus vieno kopīgi morāles, ētikas, filosofijas, metafizikas principi un ideāli). Tikai Polijas karalistes teritorijā 1815.~gadā pastāvēja 13 ložas, 1819.~gadā Polijas karalistē radās nacionālā masonu kustība (\pltxti{Wolnomularstwo Narodowe}), kuras priekšgalā stāvēja Polijas karalistes armijas virsnieks V.~Lukasiņskis. Līdz 1821.~gadam ložu skaits pieauga līdz 21. (Poļu politiķis un vēstures pētnieks Č.~Vicehs pat apgalvoja, ka kopā Krievijai pakļautajā Polijas daļā darbojās 47 masonu ložas ar ap 5~000 locekļu.) 1821.~gada 1.~maijā masonu pārstāvji sapulcējās mežā Varšavas pievārtē un nodibināja Patriotisko biedrību (\pltxti{Towarzystwo Patriotyczne}). Tā bija ievērojamākā no XIX gadsimta 20.~gadu sākumā galvenokārt no šļahtičiem veidojamajām slepenajām biedrībām, kuras stādīja uzdevumu izcīnīt Polijas neatkarību. Jau 1822.~gadā notika pirmie aresti. Arī V.~Lukasiņskis tika apcietināts un, pavadījis cietumā 46~gadus, tur arī nomira, bet imperators Aleksandrs~I pavēlēja visas masonu ložas slēgt.

Patriotiskajai biedrībai 1822.~gadā izveidojās kontakti ar krievu dekabristu organizācijām. Iniciatīvu parādīja krievu dekabristu ``Dienvidu biedrība'' (``\rutxti{Южное общество}''), kura piedāvāja kopīgu sacelšanos pret carismu, republikas nodibināšanu. ``Dienvidu biedrības'' vadītājs P.~Pestelis uzrakstīja programmatisku dokumentu ``Krievu tiesa'' (``\rutxti{Русская правда}''), kuras galvenos atzinumus 1823.~gadā atbalstīja ``Dienvidu biedrība''. Tajā bija teikts: ``Polijas patstāvības atjaunošanai jānotiek uz tādiem pamatiem un tādos apstākļos, kuri pilnā mērā nodrošinātu Krieviju nākamībā pret visdažādāko darbību, kura varētu būt pretēja tās drošībai vai pilnīgam tās mieram''.

1825.~gada sākumā notika Patriotiskās biedrības un krievu dekabristu ``Dienvidu biedrības'' (``\rutxti{Южное общество}'') pārstāvju sarunas. Jau sarunu sākumā poļu pārstāvji paziņoja savu vēlmi par Polijas atjaunošanu agrākajās robežās, ieskaitot Volīniju, kā arī baltkrievu un lietuviešu zemes. Tāda jautājuma nostādne nesolīja panākumus sarunām. P.~Pestelis izteica domu, ka vistaisnīgākā būtu strīdīgo teritoriju iedzīvotāju aptauja: ``lai paši izvēlas, kādai tautai vēlas piederēt''. Sarunu noslēgumā tomēr tika panākta vienošanās par vienlaicīgu kopēju bruņotu uzstāšanos. Tikšanās notika arī turpmāk. Taču negaidītā Aleksandra I nāve pārsteidza dekabristus nesagatavotus un viņu rīkotā sacelšanās neizdevās. Arī kopīga uzstāšanas ar poļu sazvērniekiem nenotika, jo viņi situāciju Eiropā vērtēja kā bezcerīgu, gribēja nogaidīt sacelšanās Krievijā rezultātus un tikai tad rīkoties. Vēlāk viens no 1830.~gada sacelšanās ievērojamiem darbiniekiem M.~Mohnackis nosodīja Patriotiskās biedrības neizlēmību: ``Mēs palaidām garām, gandrīz noraidījām labvēlīgu apstākli, kurš ne tik drīz atkal radīsies''.

1825.~gadā Krievijas un Polijas karalistes troni ieņēma Nikolajs I. Viņš solīja saglabāt Polijas karalistes konstitucionālās iestādes, kaut sarakstē ar brāli Konstantīnu neslēpa savu riebumu pret konstitucionālo iekārtu. Nikolajam I jau no bērnības bija ieaudzināta antipātija pret poļiem un viņš to nekad neslēpa. 1826.~gadā pēc vietvalža J.~Zaijončeka nāves viņa vietā jauns netika nozīmēts, bet tā pilnvaras nododot Administratīvajai padomei. Nikolajs I arī kategoriski nepiekrita poļu šļahtiču plāniem panākt lietuviešu, baltkrievu un ukraiņu apdzīvoto guberņu iekļaušanu Polijas karalistē. Lieta tā, ka 1817. un 1819.~gadā ar imperatora Aleksandra I dekrētiem lielkņaza Konstantīna kā Polijas karalistes armijas pavēlnieka vara tika izplatīta arī uz Krievijas Ziemeļrietumu novada guberņām, kur tika izveidots Lietuvas atsevišķais korpuss (\rutxti{Литовский отдельный корпус}), daļa no kura bija pastāvīgi izvietota Polijas karalistes teritorijā. Pat administratīvajās un civillietās vietējās iestādes sāka griezties pie Konstantīna Varšavā. Izglītības jomā izvērsās polonizācija. 1825.~gadā, neilgi pirms savas nāves apmeklējot Poliju, Aleksandrs I sarunās ar lielkņazu Konstantīnu apstiprinājis viņam savu nodomu pievienot Lietuvu Polijas karalistei. Turpretī Nikolaja I nodomi bija pavisam pretēji, viņš vēlējās Krievijas rietumu guberņas administratīvā, reliģiskā un nacionālā ziņā unificēt ar pārējām impērijas guberņām. Arī lielkņaza Konstantīna vara jaunajam imperatoram šķita pārlieku liela, uz tās rēķina viņš vēlējās nostiprināt monarha autoritāti. Nikolajs I, nākot tronī, uzsāka Lietuvas civilās administrācijas, kā arī Lietuvas atsevišķā korpusa rusifikāciju. (1827.~gadā korpusu sāka veidot no Grodņas, Minskas, Volīnijas rekrūšiem, bet iesaucamos no Viļņas un Podolijas guberņām nosūtīja uz impērijas iekšieni, nomainot Lietuvas korpusā ar krievu Pleskavas un Tveras guberņas rekrūšiem. 1831.~gadā Lietuvas korpuss tika pārdēvēts par 6.~korpusu.) Sākās varas iestāžu spiediens pret uniātu baznīcu ar mērķi pievienot to pareizticīgajai.

Iespējams, ka lielkņazs Konstantīns, nākot tronī Nikolajam I, savu varu Varšavā, tiesības uz mūžu pārvaldīt Polijas karalisti un daļēji arī Krievijas Ziemeļrietumu novadu uzskatīja kā kompensāciju par Krievijas troņa atdošanu jaunākajam brālim. Lielkņazs Konstantīns kļuva par sava veida starpnieku starp Krievijas imperatoru un Polijas karalistes sabiedrību. Ar šādu starpniecību nemierā bija abas puses.

Dekabristu lietas izmeklēšanas gaitā 1826.~gadā tika arestēti arī atsevišķi poļu Patriotiskās biedrības dalībnieki. Pret 8 no tiem tika celta apsūdzība. Lietu nodeva Seima tiesas izskatīšanai. Pretēji imperatora vēlmei 1828.~gada pavasarī tiesa noraidīja apvainojumus valsts nodevībā un atstāja spēkā vienīgi apsūdzības par piederību slepenām biedrībām, pieņemot salīdzinoši mīkstu spriedumu. Visbargākais sods~--- 3 gadi ieslodzījumā 1828.~gadā tika piespriests tikai vienam cilvēkam. Tomēr imperators spriedumu apstiprināja. Taču visi tiesājamie tika nogādāti Pēterburgā, it kā lai konfrontētu ar krievu sazvērniekiem.

Krievu~--- turku kara laikā (1828--1829) Nikolajs I vēlējās izmantot arī Polijas karalistes armiju, taču pret to enerģiski iebilda un savu panāca lielkņazs Konstantīns. Tā kā sākotnēji kara iznākums bija neskaidrs, Nikolajs I nevēlējas uztraukt poļu sabiedrību un sarežģīt situāciju Polijā. (Pēc Varnas ieņemšanas kā savas labvēlības apliecinājumu viņš ziedoja Varšavai divpadsmit lielgabalus, saņemtus kā trofejas, nedaudz vēlāk Varšavas katedrālei vairākus turkiem atņemtus karogus.)

Varšavā Nikolajs I kronējās gan tikai 1829.~gadā. Tiesa, viņš bija atvedis līdzi 11 gadus veco troņmantnieku Aleksandru (nākamo Aleksandru II), kurš tērpās poļu gvardes strēlnieka mundierī un tekoši runāja poliski. Arī pats Nikolajs I centās iekarot poļu augstmaņu, bet sevišķi virsnieku, simpātijas, bez konvoja pastaigājās pa pilsētu. Tomēr pirms aizbraukšanas Nikolajs I parakstīja dekrētu par jaunu senatoru iecelšanu, kas bija pretrunā ar 1815.~gada Konstitūciju, kas ierobežoja to skaitu. Tas izsauca poļu nepatiku. Valdīja neapmierinātība ar Krievijas kundzību, tika gatavotas sazvērestības. Tagad poļu sabiedriskie darbinieki centās uzsvērt Aleksandra I nopelnus, tā vēloties ietekmēt Nikolaju I turpināt viņa brāļa realizēto kursu attiecībās ar poļiem. Katru reizi, kad Nikolajam I tika atgādināts par Polijas karalistes Konstitūcijas pārkāpumiem, tika pieminēta arī Aleksandra I labvēlīgā attieksme pret poļiem. Kā rakstīja poļu vēsturnieks S.~Smoļenskis, tauta nejuta pret Nikolaju I pateicību kā pret Aleksandru I, tajā auga nepatika un naids pret Polijas Konstitūcijas pārkāpēju.

Mūsdienu poļu vēsturnieki M.~Timovskis, J.~Kenevičs un J.~Holcers ir atzinuši: ``No mūsdienu pozīcijām, pēc pusotra gadsimta, var konstatēt, ka lielos vilcienos piecpadsmitgade ``pēc Vīnes kongresa'' bija pati labvēlīgākā no poļu nacionālo interešu viedokļa visā periodā no 1795. līdz 1918.~gadam (jo Varšavas hercogiste bija tikai epizode, kura maksāja milzīgus cilvēku un materiālos upurus, saistītus ar pastāvīgiem kariem)''. Diemžēl šie vēsturnieki izvairās izsvērt, kam un cik liela atbildība bija jānes par nosauktā posma pārtraukšanu, un vai tai (pārtraukšanai) bija citas un kādas alternatīvas. Autors no savas puses var tikai teikt, ka bez skaidri redzamās Krievijas carisma vainas, noteikti jārunā arī par poļu radikāļu atbildību savas tautas priekšā par visiem tālākajiem notikumiem, poļu tautas ciešanām.

\strong{1830.--1831.~gadā} notika \strong{poļu sacelšanās}, pazīstama arī kā Novembra sacelšanās (\pltxti{Powstanie listopadowe}). (Baltkrievu autors A.~Tarass to sauc par 1830.--1831.~gada poļu-krievu karu). Šā gada augustā Varšavā nonāca ziņa par sacelšanos Parīzē un Burbonu (franču \frtxti{Bourbon}~--- dinastija Eiropā) dinastijas gāšanu. Sekoja baumas, ka Krievijas imperators Nikolajs I revolūcijas apspiešanai Francijā vēlas sūtīt krievu un poļu spēkus. Patiesībā Nikolajs I gan tikai paziņoja Konstantīnam, lai viņš veic pasākumus krievu karaspēka izmetināšanai, kurš caur Poliju dosies apspiest revolūciju Beļģijā. Poļu sazvērnieki, baidoties, ka arī karalistes armija var tikt izvesta no tās teritorijas, forsēja sacelšanās sākumu, nolēma carismam pretstatīt tautas atbrīvošanās centienus.

Polijas karalistē sazvērnieki varēja rēķināties ar ap 10~000 poļu karavīru pret 7~000 krievu, pie kam daudzi no pēdējiem bija dzimuši un auguši bijušajās Polijas teritorijās. 1830.~gada sacelšanās faktiski bija vienīgā no visām poļu bruņotajām akcijām, kurai bija kaut niecīgas izredzes uz panākumiem. Polijai bija lai maza, tomēr sava armija, kura guva virkni uzvaru cīņā pret Krievijas karaspēku.

Sacelšanās dalībniekiem nebija kādas izstrādātas rīcības un mērķu programmas. Šļahtiču apziņā sacelšanās ideja pati par sevi bija tik diža, ka aizēnoja visas citas sociāli-politiskās problēmas, kuru uzstādīšana, skarot šļahtiču kārtas interese, varēja būtiski mazināt viņu plašo aizrautību ar nacionālās atbrīvošanās ideju.

1830.~gada 17.(29.)~novembrī pusotra desmita bruņotu sazvērnieku, kursanti no Varšavas apakšvirsnieku skolas P.~Visocka vadībā iebruka lielkņaza Konstantīna rezidencē Belvedēras pilī (\pltxti{pałac Belwederski}) Varšavā. Taču kamerdīners (vācu \detxti{Kammerdiener}~--- istabas kalpotājs) paguva lielkņazu brīdināt par briesmām, viņš spēja noslēpties un tā izglābās, sacēlušies vienīgi nogalināja kādu krievu ģenerāli, ievainoja dažus galminiekus un apkalpotājus. Pēc tam pilsētas ielās viņi vēl nogalināja arī vairākus Krievijai uzticīgus poļu ģenerāļus, tai skaitā kara ministru M.~Hauki, kurš bija karojis gan vēl T.~Kostjuško, gan Napoleona vadībā, bet tagad atteicās izpildīt revolucionāru prasību un nostāties sacēlušos pusē, griezās pie tiem ar runu, kurā viņu darbošanos nosauca par neprātīgu un aicināja doties mājup. Vēlāk izrādījās, ka ministra krūtis viņa sievas un bērnu klātbūtnē bija cauršautas 19 vietās.

Imperators Nikolajs I vēlāk lika nogalinātajiem poļu ģenerāļiem, kuri labāk krita no sacēlušos rokas, nekā tiem pievienojās, uzcelt lielu obelisku, izrādīja īpašas rūpes par viņu ģimenēm. Tā, M.~Haukes meita Jūlija kļuva par Krievijas ķeizarienes galma dāmu, apprecēja Hessenes princi Aleksandru, viņu pēcteči saradojās ar daudziem Eiropas valdnieku namiem.

Turpretī mūsdienās Polijā 29.~novembris tiek atzīmēts kā ievērojams Polijas vēstures notikums. Šai dienā karaskolu beidzējiem svinīgi tiek piešķirtas virsnieku pakāpes.

Krievu publicists un vēsturnieks N.~Bergs rakstīja, ka, ja sacelšanās nebūtu krievus pārsteigusi pēkšņi, ``uguni varētu apdzēst tai pat mirklī, un [1830.~gada] 30.~novembra rītā daudzi iedzīvotāji pat nezinātu, ka tika gatavots kaut kāds nopietns sprādziens''. Arī poļu vēsturnieks V.~Smoļenskis uzskatīja, ka lielkņazam Konstantīnam pietika karaspēka, lai apspiestu sacelšanos pašā tās iedīglī. ``Sacelšanās izdzīvoja tikai pateicoties Konstantīna neaktivitātei, kurš, pastaigājoties krievu karaspēka ierindas priekšā Ujazadovska alejā (\pltxti{Aleje Ujazdowskie}), nepārtraukti atkārtoja, ka pašiem poļiem ir jānomāc viņu tautiešu radītās nekārtības''. Iespējams, ka tas bija pārāk optimistisks stāvokļa vērtējums. Poļu armijas daļas nostājās sacēlušos pusē un kopā ar pilsētniekiem ieņēma ieroču arsenāla un krievu ulānu (viegli bruņotas kavalērijas paveids) pulka kazarmu ēkas. Drīz visa Varšava atradās jau sacēlušos rokās. Poļu armiju atbalstīja iedzīvotāji. Sākusies kā militāra apvērsuma mēģinājums, kustība izvērsās par tautas sacelšanos. Krievu garnizons Varšavā neizrādīja pienācīgu pretošanos un pilsētu atstāja. Lielkņazs Konstantīns paziņoja, ka ``katra izlieta asins lāse tikai sabojās lietu'' un ar krievu daļām uzsāka atkāpšanos uz Krievijas robežu.

Varšavā izveidojās Pagaidu valdība, kura, kustības aktīvistu spiesta, savukārt par diktatoru nozīmēja Napoleona karu dalībnieku ģenerāli J.~Hlopicki. Sabiedrības radikālā daļa diktatūru sagaidīja ar prieku. Tā uzskatīja, ka J.~Hlopickis izdarīs to, ko nedarīja Pagaidu valdība: pasludinās sakaru saraušanu ar Nikolaju~I, aicinās Lietuvas un Krievijas provinces uz sacelšanos un sāks cīņu par neatkarīgu Poliju. Taču ģenerālis, kurš neticēja sacelšanās uzvarai, centās atrisināt konfliktu sarunu ceļā, nosūtot sūtni pie Nikolaja I. Pagaidu valdība kā noteikumus, lai Polijas karaliste paliktu imperatora pakļautībā, izvirzīja prasības, lai Nikolajs I ievēro 1815.~gada Konstitūciju, to izplata arī uz lietuviešu, baltkrievu un ukraiņu zemēm, izsludina amnestiju, apņemas nesūtīt krievu armiju Polijā.

Populārā izklāstā šo noteikumu izpilde nozīmētu virknes pamatos nepoļu apdzīvotu teritoriju pievienošanu Polijas karalistei. Jākonstatē, ka poļu nacionālās atbrīvošanās cīņa norisa reizē ar viņu cīņu par svešu tautu pakļaušanu. Poļi kādreiz šais zemēs bija valdījuši, taču nonākot pašiem citas valsts pakļautībā, nebija iemācījušies cienīt arī citu tautu brīvību. Tas, ka poļu nacionālās apziņas formēšanās bija gājusi tālāk kā ukraiņiem un baltkrieviem, nedeva pirmajiem tiesības ignorēt gadsimtu gaitā daudzskaitlīgās sacelšanās skaidri parādīto šo tautu nevēlēšanos pakļauties poļu paniem.

Nikolajs I apsolīja vienīgi amnestiju sacelšanās pārtraukšanas gadījumā. Krievijā šai laikā nebija organizēta spēka, kurš varētu sniegt poļu sacelšanās dalībniekiem reālu palīdzību.

Decembrī sanāca Seims, kurš nosodīja carisma noziegumus, bet sacelšanos pasludināja par visas poļu tautas lietu. Taču imperators pavēlēja atjaunot agrākās varas struktūras, bet poļu armiju koncentrēt Pinskas rajonā karalistes austrumos. Saasinājās uz samierināšanos orientētā diktatora attiecības ar citiem līderiem un 1831.~gada 18.~janvārī J.~Hlopickis atteicās no diktatora pilnvarām. Seims par diktatoru iecēla vēl vienu T.~Kostjuško un Napoleona karu dalībnieku M.~Radzivillu, kurš gan arī drīz no amata atteicās.

25.~janvārī Seimam kļuva zināms imperatora brīdinājums, ka pirmais lielgabala šāviens no poļu puses nozīmēs Polijas karalistes patstāvības beigas. Atbildei uz to Seims bez balsošanas pieņēma lēmumu gāzt Nikolaju I no Polijas karalistes troņa, atzīstot to par vakantu. Tas bija līdzvērtīgi kara pieteikumam.

Dažas dienas vēlāk~--- 28.~janvārī Seims izveidoja Nacionālo valdību (\pltxti{Rzad narodowy}) ar kņazu Ā.~Čartorijski priekšgalā. Jau minētais V.~Zajevskis uzskata, ka Ā.~Čartorijskis un viņa piekritēji pēc sacelšanās sākuma nolēma nostāties tās priekšgalā, lai piešķirtu tai mērenu, konservatīvu raksturu. Visus spēkus viņi veltīja, lai izmantotu sacelšanos izdevīgu sarunu vešanai ar imperatoru, taču nevēlējās piekrist bezierunu kapitulācijai. Nacionālajā valdībā ietilpa arī demokrātisko aprindu pārstāvji, kā Viļņas universitātes docētājs vēsturnieks J.~Lelevels. Viņš nedaudz izmainīja jau XVIII gadsimta beigās (1795), kad poļu šļahta bez panākumiem mēģināja iesaistīt savā cīņā ukraiņu un baltkrievu dzimtcilvēkus, radušos lozungu ``\pltxti{Za naszą i waszą wolność}'' (Par mūsu un jūsu brīvību). Uz sacēlušos karoga tagad bija rakstīts ``\pltxti{W imię Boga za Naszą i Waszą Wolność}'' (Dieva vārdā par mūsu un jūsu brīvību). Šādā variantā tas arī kļuva starptautiski populārs.

3.~februārī Seims apstiprināja zvērestu, kuru pats arī pirmais nodeva, ka sacelšanās mērķis ir ``tās neatkarības, tās pakāpes starp tautām, kurai Radītājs iecerējis poļu tautu, iegūšana''. Ar īpašu manifestu Seims mudināja tautu nodot šo zvērestu.

Uz Parīzi, Vīni un citām Eiropas galvaspilsētām tika nosūtīti aģenti, kuru uzdevums bija rosināt šo valstu valdības iejaukties Polijas notikumos. Kņazs Ā.~Čartorijskis ar brāļa K.~Čartorijska starpniecību saistījās ar Austrijas galmu, piedāvājot Polijas kroni vienam no austriešu erchercogiem, ja vien Austrija atbalstīs sacelšanās mērķus. Sākotnēji Vīnes galms ar apmierinājumu uzņēma ziņas par sacelšanos Polijā, taču, kad tā radikalizējās un par vienu no līderim kļuva J.~Lelevels, Austrija ne tikai neatbalstīja sacelšanos, bet arī noslēdza savu robežu, traucējot poļu sakarus ar rietumu valstīm. Kaut, pēc poļu vēsturnieka Š.~Aškenazi vērtējuma, Prūsija bija ieinteresēta Polijas karalistes sabrukumā, kas radītu iespēju iegūt jaunas teritorijas (tiesa, šo secinājumu minētais zinātnieks nepamatoja un šādas vēlmes noteikti bija pretrunā ar Prūsijas drošības un stabilitātes interesēm), Prūsija tāpat aizvēra savu robežu, noslēdza ar Krieviju traktātu, apņemoties uz poļu robežas nostādīt 60~tūkstošu vīru lielu korpusu, gatavu ieņemt Varšavu. Francijā sabiedrība sacelšanos Polijā uztvēra ar saucieniem ``Palīgā poļiem!''. 1831.~gada 13.~februārī pūlis Parīzē apmētāja akmeņiem Krievijas sūtņa namu, izsitot tam logus. Kopš 1831.~gada februāra par Francijas bruņotu iesaistīšanos cīņa dumpīgo poļu pusē pret Nikolaju I dedzīgi iestājās un vāca iemaksas sacelšanās vajadzībām Polijas sacelšanās atbalsta komitejas priekšsēdētājs marķīzs M.~De~Lafajets. Taču šo sajūsmu neatbalstīja Jūlija revolūcijas rezultātā tronī kāpušais karalis Ludvigs-Filips. Tikpat nesekmīgi bija poļu sūtņu pūliņi Anglijā. Nekādu būtisku palīdzību no ārzemēm poļu sacēlušies nesaņēma.

Polijas sadalei, poļu, baltkrievu, ukraiņu, baltiešu zemju iekļaušanai Krievijas sastāvā acīmredzot vajadzēja pēdējo nostiprināt. Bet jebkura Krievijas nostiprināšanās Rietumos parasti tika uzskatīta par tiešu izaicinājumu savai drošībai. Tāpēc arī pārliecība, ka ''Eiropa mūs aizsargās'', stiprināja poļu virsslāņa apņēmību cīnīties par Žečpospolitas atjaunošanu un iespēju apspiest nepoļu iedzīvotājus. Poļu sacelšanos vadoņi neņēma vērā tikai vienu: lielvalstīm nav mūžīgu ienaidnieku, tām ir tikai mūžīgas intereses. Ja tās nesakrita ar poļu interesēm, pēdējiem uz kādu palīdzību bija velti cerēt.

Revolucionārāk noskaņotie poļu nacionālās kustības elementi no ``Patriotu biedrības'' (\pltxti{Towarzystwo Patriotyczne}) centās sacelšanās īstenošanā iesaistīt zemniekus. Tikai zemnieku atbrīvošana no klaušām varētu nostādīt tos sacelšanās pusē. Taču likumprojektu par klaušu nomaiņu ar nodevām Seims noraidīja, jo tas izsauktu neapmierinātību Prūsijā un Austrijā, apdraudētu arī pašu poļu šļahtiču, kuri atbalstīja sacelšanos, vairākuma intereses. Tāpēc šai virzienā arī praktiski nekas netika darīts. Ar to sacelšanās jau sākotnēji bija nolemta neveiksmei. (Tiesa, krievu vēsturnieks A.~Pogodins uzskatīja, ka, ja Seims 1830.~gadā pat pasludinātu zemnieku brīvlaišanu, zemnieki diezin vai sekotu tam.) Plašas zemnieku masas, īpaši ukraiņu un baltkrievu zemnieki, sacelšanos neatbalstīja.

1831.~gada pavasarī poļu armijā bija ap 79~tūkstoši cilvēku ar 158~lielgabaliem. Tajā karoja arī sievietes, piemēram, Līksnā pie Daugavpils dzimusī grāfiene E.~Plātere. Par viņu stāsta leģendas, viņa esot līdzinājusies Francijas nacionālajai varonei Žannai d’Arkai, viņas izveidotā un vadītā vienība esot iedvesusi šausmas ienaidniekam. Viņai veltīts arī Ā.~Mickeviča dzejolis ``Pulkveža nāve'' \pltxti{(``Śmierć pułkownika'').} Tiesa, sacelšanās liecinieks prelāts, bīskaps, Varšavas garīgās akadēmijas rektors P.~Butkevičs, būdams jau Parīzē, uzrakstīja memuārus, kuros rādīja, kā E.~Plātere, pieķerta intīmos sakaros ar pakļautajiem virsniekiem, kļuva par izzobošanas objektu, pēc sakāvēm pameta savu vienību likteņa varā, bēguļoja un pirms 1831.~gada 23.~decembrī iestājušās nāves nožēloja savu iepriekšējo rīcību. Iespējams gan, ka minētais autors vadījās no citiem, nevis patiesības noskaidrošanas motīviem.

Šai laikā radās viena no vispopulārākajām poļu patriotiskajām dziesmām, 1830.~gada Novembra sacelšanās muzikālais simbols ``1831.~gada (lai atšķirtu no 1905.~gada tāda pat nosaukuma dziesmas) Varšavjanka'' (``\pltxti{Warszawianka 1831 roku}''). Parīzes dzejnieks Ž.~Delavins 1831.~gada februārī sacelšanās Varšavā iespaidā franču valodā pēc ``Marseljēzes'' (``\frtxti{La Marseillais''}~--- Francijas himna) parauga uzrakstīja dzejoli ``\frtxti{La Varsovienne''}. Pēc mēneša tas nonāca Varšavā, tika pārtulkots, publicēts, Varšavas operas teātra diriģents un komponists K.~Kurpinskis tam sacerēja melodiju. 1831.~gada 5.~aprīlī dziesma tika pirmo reizi publiski nodziedāta operas izrādes laikā un publika to uzņēma ar entuziasmu. Turpmākajos gadu desmitos dziesma kļuva par vienu no poļu nacionālās atbrīvošanās kustības simboliem. Līdz 1926.~gadam, kad ``Dombrovska mazurka'' kļuva par oficiālo Polijas himnu, ``1831.~gada Varšavjanka'' bija viena no piecām kandidātēm uz himnas vietu. Tā bija arī viena no populārākajām 1944.~gada Varšavas sacelšanās dziesmām.

1831.~gada februārī feldmaršala I.~Dibiča, pēc tam feldmaršala I.~Paskeviča vadītā Krievijas armija, kurā bija ap 115~000 karavīru, neraugoties uz holēras epidēmiju (tā aiznesa vairāk nekā 27 tūkstošu krievu karavīru dzīvību, par tās upuriem kļuva arī feldmaršals I.~Dibičs un lielkņazs Konstantīns), sakāva poļus, kuru spēki sniedzās līdz 55 tūkstošiem. (Tāpat kā 1794.~gada sacelšanās laikā arī 1831.~gadā ap 180~--- 200 dažādu tautību krievu karavīru un virsnieku pārgāja poļu sacēlušos pusē). 26.~septembrī, tai pat dienā, kad pirms gadiem notika Borodinas kauja, tika ieņemta Varšava. Sakarā ar to A.~Puškins uzrakstīja dzejoli ``Borodinas gadadienā'', kurā bija rindas:

\vspace{1.5em}

\noindent
\begin{minipage}{0.55\textwidth}
\rutxti{
Сбылось~--- и в день Бородина\\
Вновь наши вторглись знамена\\
В проломы падшей вновь Варшавы;\\
И Польша, как бегущий полк,\\
Во прах бросает стяг кровавый~---\\
И бунт раздавленный умолк.}
\end{minipage}
\hspace{1em}
\begin{minipage}{0.5\textwidth}
Ir noticis~--- Borodinas dienā\\
Jau atkal slejas mūsu karogi\\
Jau atkal kritušajā Varšavā;\\
Un Polija, kā bēgošs pulks,\\
Met pīšļos asiņaino standartu~---\\
Un dumpis nomākts klusē.
\end{minipage}

\vspace{1.5em}

% page 121


Taču pirms Varšavas ieņemšanas, kā liecina kņazienes N.~Goļicinas atmiņas, sacēlušies izrēķinājās ar sagūstītajiem krievu karavīriem un vairākiem galminiekiem. Kņaziene to nosauca ``par bojā ejošās nācijas pēdējo ārprāta lēkmi''. Nākas tikai konstatēt, ka visdažādākajos laikos patriotisms un cīņas griba ne vienmēr tiek apliecināta atklātā kaujas laukā, bet dažkārt arī vairāk nekā apšaubāmā veidā. Nebija izņēmums arī 1830.--1831.~gada poļu sacelšanās. Divas dienas ilgajā Varšavas ieņemšanā krievu karaspēks vēl zaudēja ap 10 tūkstošu cilvēku, poļu~--- līdz 11 tūkstošu. Trīs tūkstoši tika saņemti gūstā. Tika iegūti arī 132 lielgabali.

Pēc tam sacelšanās turpinājās mazāk nekā mēnesi. Daļa poļu karaspēka: ap 15 tūkstošu cilvēku ar 42 lielgabaliem pārgāja Austrijas un ap 20 tūkstošu ar 96 lielgabaliem~--- Prūsijas robežu, kur tika internēta. Vēl gan notika mēģinājums atjaunot bruņotu cīņu. 1833.~gadā Krievijas impērijas robežu šķērsoja vairākas sīkas poļu revolucionāru-emigrantu vienības J.~Zaļivska vadībā. Izraisīt sacelšanos neizdevās, ekspedīcija tika sakauta, daļu tās dalībnieku sagūstīja, pakāra, nošāva vai nosūtīja katorgā uz Sibīriju. Pats J.~Zaļivskis bēga uz Austriju, tika tur arestēts un notiesāts uz 20~gadiem cietumā. Pēc amnestijas 1848.~gadā viņš devās uz Franciju, kur arī mira.

I.~Paskevičs par sacelšanās apspiešanu saņēma Varšavas kņaza (\rutxti{князь Варшавский}) titulu, kļuva arī par Polijas karalistes vietvaldi (1832--1856). Kā raksta Polijas vēstures pētnieks krievu vēsturnieks G.~Matvejevs, ar šo laiku Krievija sāka uzskatīt, ka tā ir iekarojusi Polijas karalisti, nevis ieguvusi to uz Vīnes kongresa (1815) lēmumu pamata. Tiesa, pēc sacelšanās sakāves 1831.~gada beigās iznāca imperatora rīkojums, ar ko sacelšanās dalībnieki tika amnestēti, izņemot notikumu dalībniekus Varšavā 1830.~gada novembrī, Seima deputātus, kuri bija balsojuši par Polijas karalistes troņa atņemšanu Nikolajam I, un poļu virsniekus, aizbēgušus uz ārzemēm.

Nikolajs I 1830.--1831.~gada sacelšanos uzskatīja par poļu nodevību, poļu nacionālās jūtas~--- par nepiepildāmu fantāzijas augli un savus ārpolitiskajos plānos vadījās no Krievijas valdošās Romanovu dinastijas interesēm. 1831.~gada vasarā viņš pat nopietni apsvēra plānu apmainīt ar Prūsiju un Austriju to valdījumos nonākušās poļu zemes. (Par trim Polijas karalistes rietumu vojevodistēm viņš no Prūsijas vēlējās saņemt Toruņu, Klaipēdu un Nemunas grīvu, bet par Krakovas vojevodisti no Austrijas Tarnopoles apriņķi.) Tomēr acīmredzot galu galā imperators saprata, ka šāda apmaiņa bija pārāk sarežģīta un tāpēc nerealizējama un no tās atteicās.

Reizē varēja runāt par Aleksandra I realizētās ``mīkstās'', samērā liberālās politikas neveiksmi Polijas karalistē. Krievu vēsturnieks B.~Mironovs par pēc Napoleona kariem XIX gadsimta sākumā izveidojošos situāciju Polijā Krievijas impērijas sastāvā ir norādījis: ``Polija sniedz piemēru, kā liela autonomija, liberāla politiskā iekārta, kas bija daudz progresīvāka salīdzinājumā ar iekārtu pašā Krievijā, neradīja apmierinājumu iekarotās tautas prātos. Kā to var izskaidrot? No vienas puses, tūkstoš gadu vecā valstiskuma tradīcija, ar daudziem sasniegumiem iezīmēta vēsture, katolicisms, pārākuma sajūta pār uzvarētājiem neļāva poļu tautai samierināties ar suverenitātes zaudēšanu. No otras puses, kā man šķiet, ka Krievijas forsētā valstiskuma celtniecība Polijā 1815.~gadā radīja tādu situāciju, ka Krievija~--- Polijas valstiskuma radītāja~--- poļiem kļuva praktiski nevajadzīga pēc tam, kad jaunā Polijas valsts bija izveidota. Poļi steidzās atbrīvoties no nevajadzīgas aizbildniecības.''

Kā pamatoti atzinis vācu vēsturnieks E.~Meijers, sacelšanās dalībnieki, to uzsākot, faktiski paši lika uz spēles visus tos sasniegumus, kuri bija panākti pusotrā gadu desmitā pēc Vīnes kongresa un gandrīz visu zaudēja.

Eiropā sabiedrība līdzjūtīgi pieņēma poļu neatkarības cīnītājus. Francijas karaļa Luija Filipa valdībā neliedza viņiem patvērumu. Morālu atbalstu viņiem sniedza Polijas atbalsta komiteja, kuru vadīja ģenerālis M.~Lafajets. Tiesa, drīzumā pēc imperatora Nikolaja I lūguma Luisa Filipa valdība aizliedza darbību, vērstu pret Krieviju. Tāpēc, piemēram, J.~Lelevels bija spiests pamest Franciju un apmesties Briselē.

Arī Polijai blakus esošajā Vācijā, sabiedriskā doma juta līdzi poļu brīvības centieniem. Radās ap 100~poļu draugu biedrību (\detxti{Polenfreundvereine}), kuras gan vēlāk to politizācijas dēļ tika aizliegtas. Poļu emigrantus sagaidīja silta uzņemšana. R.~Vāgners 1836.~gadā uzrakstīja uvertīru ``\detxti{Polonia''} (``Polija''), tapa vairāku citu vācu komponistu ``poļu dziesmas'', virkne vācu dzejnieku veltīja dzejoļus poļu brīvības cīņām. Tikai daži no vācu dzejniekiem, kā H.~Heine un E.~Arndts, spēja būt arī kritiski. 40.~gados gan Vācijā ``sajūsma par poļiem'' noplaka. Acīmredzot to ietekmēja tas, ka no ap 50~000 pēc sacelšanās no dzimtenes bēgušajiem tās dalībniekiem pēc cara amnestijas pasludināšanas ap 40~000 atgriezās Polijas karalistē, tā netieši atzīstot, ka nemaz tik ``nepanesami'' apstākļi tajā nebija.

Ebreju izcelsmes poļu vēsturnieks Š.~Aškenazi īpaši uzsvēris, ka tieši poļu cerību sabrukums uz Lietuvas [faktiski daļas Krievijas Ziemeļrietumu un Dienvidrietumu novadu~--- V.Š.] pievienošanu Polijas karalistei Nikolaja I valdīšanas laikā bija galvenais cēlonis, kāpēc sākās 1830.~gada poļu sacelšanās. Viņš rakstīja, ka ``Novembra revolūcija bija pirmkārt Polijas karš ar Krieviju Lietuvas dēļ. Otrs cēlonis, kaut tas bija skaidrāk saredzams, jo atradās sabiedriskās apziņas virspusē, bija ar samērā mazāku nozīmi: tas bija konstitucionālais jautājums, pastāvīga [Polijas karalistes] konstitucionālā statusa pārkāpšana.'' Vēsturnieks arī atzīmējis, ka sava nozīme bija arī ortšķirīgiem cēloņiem, kā pretrunām starp valdošā nama brāļiem Nikolaju un Konstantīnu, krievu--turku karam, kurš solīja labas izredzes uz sacelšanās izdošanos, revolūcijai Rietumos, kura deva signālu sprādzienam.

Nav pilnīgu datu par sacelšanās dalībnieku sociālo stāvokli. Taču ziņu par 16~tūkstošu Lietuvas, Baltkrievijas un Ukrainas teritorijās cara valdības represētajiem sacelšanās dalībniekiem analīze rāda, ka to vidū ap 33\% piederēja pie nodevu maksātāju, pirmkārt zemnieku, slāņa, gandrīz 50\%~--- šļahtiču kārtai, ap 5\%~--- garīdzniecībai. Kaut nav datu par pašu poļu zemēm, tomēr esošās ziņas ļauj izdarīt zināmus secinājumus par sacelšanos kopumā. Statistiskais materiāls rāda, ka sacelšanās kustībā bija divi pamatelementi: pirmkārt, pārsvarā šļahtiču vestā nacionālās atbrīvošanās kustība; otrkārt, zemnieku antifeodālā cīņa. Abas attīstījās paralēli, taču ne sinhroni. Abas tās stiprināja viena otru, kaut pilnīgas saskaņas starp tām gandrīz nekur nebija. Krievijas rietumu guberņu zemnieki bija lietuvieši, baltkrievi, ukraiņi. Daļa no viņiem sacēlušos formējumu sastāvā atradās piespiedu kārtā, pēc savu muižnieku pavēles. Tā daļa, kura apzināti stājās sacēlušos rindās, sekoja saviem mērķiem, kuri bieži atšķīrās no sacēlušos valdības pasludinātajiem. Daudzi zemnieki, īpaši Lietuvā, uzstājās ne tik daudz pret cara valdību, cik pret klaušu sistēmu, saviem, pārsvarā poļu, muižniekiem.

Pēc sacelšanās poļu emigrācijas lielākā daļa, sastāvoša no vidējiem un zemākajiem šļahtiskās un buržuāziskās inteliģences slāņiem, nāca pie secinājuma, ka kopumā Polijas krišanas galvenais cēlonis bija šļahtas anarhistiskā kundzība, bet 1830.--1831.~gada sacelšanās sakāves cēlonis~--- sacelšanās vadošo slāņu konservatīvisms, bailes iesaistīt cīņā iedzīvotāju pamatmasas~--- zemniecību. Tā, viens no bijušajiem Patriotiskās biedrības kreisā spārna vadītājiem T.~Krepovickis, vēlāk emigrācijā analizējot 1830.--1831.~gada sacelšanās mācības, šļahtičus~--- muižu īpašniekus sauca par ``Tēvzemes un cilvēces nodevējiem'', apgalvojot, ka tās pārstāvji ``nebija nedz dižas idejas, nedz revolucionāras domas iedvesmoti''.

1848.~gadā, pieminot 1830./1831.~gada notikumus Polijā, F.~Engelss konstatēja, ka Polijas valdošā šķira izrādījās tikpat egoistiska, aprobežota un bailīga likumdošanas sapulcē, cik pilna entuziasma un vīrišķības tā bija kaujas laukā. Pēc F.~Engelsa vērtējuma ``1830.~gada sacelšanās nebija ne nacionāla revolūcija (tā atstāja aiz borta trīs ceturtdaļas Polijas), ne sociāla vai politiska revolūcija, tā neko neizmainīja tautas iekšējā stāvoklī, tā bija konservatīva revolūcija''. Arī poļu vēsturnieki V.~Bortnovskis, J.~Dutkevičs, V.~Zajevskis u.c. XX gadsimta 50.--60.~gados uzsvēra, ka sacēlušos kreisais spārns, kurš sapratis, ka sacelšanās var uzvarēt, tikai ejot pa radikālu politisku un sociālu reformu ceļu, izrādījās pārāk vājš, tāpēc vara nonāca aristokrātisko grupējumu rokās, kuri nevēlējās nekādas reformas un karadarbības jēgu saskatīja tikai cīņā ar carisko Krieviju par ukraiņu, baltkrievu un lietuviešu zemēm, kādreiz piederējušām Žečpospolitai. Pēc autora domām, nevajadzētu noliegt sacelšanās nacionālās atbrīvošanās un arī politisko raksturu, taču citādi ar nelielām iebildēm F.~Engelsa bargajam spriedumam var piekrist. Sacelšanās cieta sakāvi, jo tās dalībnieki nespēja izvirzīt skaidru un tālejošu nacionālu un sociālu programmu.

Padomju literatūrā tika izplatīti spriedumi par minētās sacelšanās starptautisko nozīmi, jo tā veicināja revolucionārās kustības attīstību daudzās Eiropas valstīs, to sveica dekabristi un krievu jaunākās paaudzes revolucionāri. Šeit gan jāatgādina pašu marksistu atzinums, ka katrā revolūcijā galvenie ir iekšējie cēloņi. Polijas pēc F.~Engelsa vārdiem ``konservatīvā revolūcija'' kalpot par paraugu attīstītākajām Eiropas zemēm varēja gandrīz vienīgi ar savu faktu, nevis programmu un taktiku.

Mūsdienu krievu publicists M.~Muravjovs uzsvēris: ``Vēlme sagrābt ``atņemto novadu'' (tā poļi sauca Krievijas impērijas rietumu guberņas, kuras ietilpa Žečpospolitas sastāvā līdz tās sadalei) kļuva par uzmācīgu poļu ``nacionālo ideju''. Loģika šajos centienos bija vienkārša: Polija savās etnogrāfiskajās robežās, praktiski sakrītošās ar ``kongresovku'', [Publicists šeit nav korekts, poļu etnogrāfiskās robežas ievērojami pārsniedza t.s. ``kongresa Polijas'' ietvarus, par ko jau augstāk runāts~--- V.Š.] iegūstot neatkarību, būs maza un neietekmīga valsts. Toties, paplašinot robežas uz austrumiem, Polija atkal kļūs par lielvalsti. Tāpēc visas poļu sacelšanās pret Krieviju bija ar nacionālas sagrābšanas raksturu. To apspiešana no Krievijas puses bija nepieciešama pašaizsardzība.'' Publicists gan nepamatoti nostājies cariskās Krievijas režīma, kurš vairāk rūpējās par valsts rietumos dzīvojošo tautu ekspluatāciju, nevis aizsardzību pret poļu kundzību, apoloģēta lomā, taču viņa norāde uz faktiski Polijas šļahtas interesēs esošā mērķa~--- citu etnosu apdzīvoto teritoriju iekļaušanas Polijas valstī pacelšanu ``nacionālās idejas'' augstumos ir pelnījusi ievērību. Kādreizējā Žečpospolitas ``vēsturiskā diženuma'' atjaunošanas ideja vēl mūsdienās darbojas kontrproduktīvi Polijas attīstības ceļu meklējumos.

Pēc 1830.~gada sacelšanās sakāves \strong{Krievija izvērsa cīņu pret} valsts rietumos dzīvojošo tautu \strong{polonizāciju}. Poļu literatūrā jēdziens ``polonizācija'' tiek traktēts kā ``pašpolonizācija'', kā cilvēku brīvprātīga izvēle. Jāievēro, ka Žečpospolita nebija nacionāla valsts mūsdienu izpratnē, nāciju veidošanās laikmetā Polijas valsts nepastāvēja un poļu rokās nebija valstisku līdzekļu asimilācijas realizēšanai. Taču XIX gadsimta sākumā nedrīkst arī pārāk zemu novērtēt tādu reliģiozu, kultūras un sociālu institūtu kā katoļu baznīcas, oficiālās skolas, vietējās administrācijas (līdz 1863--1864.~gada sacelšanās laikam), muižas ietekmi uz nepoļu iedzīvotājiem poliskuma izplatībā pat apstākļos, kad nepastāvēja Polijas valsts. Protams, jāuzsver arī, ka Krievijas iestāžu vēršanās pret polonizāciju reizē nozīmēja virzību uz rusifikāciju.

Ievērojamākajiem sacelšanās dalībniekiem~--- emigrantiem, kā Ā.~Čartorijskim un J.~Lelevelam, aizmuguriski tika piespriests nāves sods. Šļahtičiem~--- sacelšanās dalībniekiem karalistē tika atsavinātas ap 3~000 muižu (apmēram 1/10 no lielsaimniecībām).

1833.~gada septembrī Krievija un Austrija Minhengrecā (vācu \detxti{Münchengrätz}, čehu \cstxti{Mnichovo Hradiště)} noslēdzot līgumu par kopēju politiku pret Turciju, saskaņoja arī konvenciju par \latxti{status quo} ievērošanu poļu zemēs un politisko noziedznieku savstarpēju izdošanu. Pēc mēneša tai pievienojās arī Prūsija. Konvencija faktiski deva trīspusējas garantijas tai kārtībai, kura tika radīta Krievijai piederošajā Polijas daļā pēc 1830.--1831.~gada sacelšanās.

Pat vēl vairāk kā Polijas karalisti carisma sodu politika skāra bijušās Lietuvas lielkņazistes teritorijas. Poliskās izglītības iestādes tika rusificētas, Viļņas universitāte 1832.~gadā slēgta. Daļa poļu šļahtiču Lietuvā un vēlāk arī tagadējās Ukrainas un Baltkrievijas teritorijā zaudēja savas dižciltīgo tiesības, arī muižas, tika pazemināti līdz viensētniekiem (\rutxti{однодворцы}), kuri, kaut bija personīgi brīvi, bet maksāja valstij nodevas. Sākusies jau XIX gadsimta 20.~gados dokumentu, kas apliecinātu šļahtiču dižciltīgo izcelsmi, pārbaude stipri ieilga un faktiski nebija pa spēkam smagnējai Krievijas birokrātiskajai mašīnai. Tomēr tā veicināja muižnieku norobežošanos no saviem nabagajiem kārtas brāļiem. Tiesa, lielākā daļa muižu palika to agrāko īpašnieku rokās, carisms turpināja aizstāvēt muižniecības intereses. 1836.~gadā Nikolajs I poļu muižniecības tiesības un privilēģijas pielīdzināja krievu muižniecības kā kārtas tiesībām.

Jautājumu par poļu ierēdņu izmantošanu Krievijas impērijas dienestā noteica 1852.~gada 14.~novembrī imperatora apstiprināts nolikums. Tajā bija uzskaitītas 28~guberņas, kuru iestādēs varēja strādāt poļu jaunieši. Šai uzskaitījumā nebija Pēterburgas guberņas, Baltijas, Besarābijas, toreizējās Mazkrievijas (Ukrainas), Kaukāza, Sibīrijas apgabalu. No tām guberņām, kuras robežojās ar Rietumu novadu, sarakstā bija tikai divas~--- Pleskavas un Smoļenskas guberņas.

Daļu poļu aktīvistu izsūtīja no Polijas karalistes un nometināja citās Krievijas impērijas guberņās, arī Kaukāzā un Sibīrijā. (Pēc vācu vēsturnieka H.v.~Zitzevica datiem tika izsūtīts ap 45~000 poļu šļahtiču.) Taču nevajadzētu pārspīlēt viņu ciešanas. Kā raksta mūsdienu poļu vēsturniece V.~Sļivovska, vairumā gadījumu pat uz katorgas darbiem notiesātie poļi, izņemot pirmās dienas, faktiski neizcieta piespriesto sodu. Katordznieku jūgu patiesībā nesa tikai kriminālnoziedznieki. Lieta tā, ka visā Sibīrijā turīgajiem tirgoņiem, vietējās administrācijas pārstāvjiem bija vajadzīgi ārsti, skolotāji, kuri spētu sagatavot viņu bērnus iestājai ģimnāzijā vai augstskolā, audzinātāji, kuri varētu mācīt bērniem mūziku, dejas, glezniecību. To visu varēja veikt izsūtītie poļi. Galvaspilsētas bija tālu, uz vietas visu lēma ģenerālgubernatori un tiem pakļautie ierēdņi. Poļus labprāt ņēma darbā dažādās kancelejās ``papīra darbiem''. Ar laiku daži kļuva pat par zelta raktuvju administratoriem. Oficiālajās atskaitēs viņi visi nesa tiem uzlikto sodu, ir zināms, ka Demidovu (\rutxti{Демидовы}~--- bagāta krievu uzņēmēju dzimta) rūpnīcā viņiem bija izgatavotas arī speciālas, viegli noņemamas važas~--- speciāli revizoru atbraukšanas gadījumiem. Pēc 1856.~gada amnestijas avīze ``\rutxti{Сибирь}'' (``Sibīrija'') ar nožēlu rakstīja par to, ka gandrīz visi šeit mitušie poļi aizbraukuši.

1832.~gada 14.~februārī Polijas 1815.~gada Konstitūcijas vietā tika ievests Organiskais statuts (poļu \pltxti{Statut Organiczny Królestwa Polskiego}, krievu \rutxti{Органический статут Царства Польского}). Tiesa, krievu vēsturnieks A.~Korņilovs uzsvēra, ka tas pilnībā tā arī nekad netika ieviests, jo poļu revolucionāri atkal gatavojās sacelties. Tomēr faktiski Polijas karaliste pārvērtās par vienkāršu Krievijas provinci, neatņemamu Krievijas impērijas sastāvdaļu. Īpaša kronēšana ar tās kroni vairs nebija paredzēta. Līdz 1856.~gadam Krievijas Polijā darbojās kara stāvoklis. Seims tika likvidēts. Saglabājās tikai dažas autonomas iestādes. Pārvaldes padomes locekļi vairs nebija poļi. Padomes kompetence tika aprobežota ar reliģijas, kultūras un tirdzniecības lietām. Tikai līdz 1841.~gadam saglabājās Valsts padome kā konsultatīva iestāde pie Krievijas Valsts padomes (1861.~gadā, sākoties revolucionāram pacēlumam, to gan atkal atjaunoja). 1832.~gadā poļu valūta~--- zlots tika nomainīta ar Krievijas rubli. Uz monētām, kuras tika kaltas 1832--1842 gadā, bija uzraksts «\rutxti{15 коп.}~--- \pltxti{1 zloty}» (15 kapeikas~--- 1 zlots).

Poļu armija tika likvidēta, tās vietā ieviesta rekrūšu ņemšana Krievijas armijā. Polijā tika izvietota stipras Krievijas armijas daļas. 1832.--1834.~gadā augstienē pie Varšavas pēc imperatora Nikolaja I pavēles 1832.--1834.~gadā tika uzbūvēts cietoksnis, kurš izmaksāja tolaik milzīgu summu~--- 11 miljonus rubļu. To kā sodu par sacelšanos bija jāsedz Varšavas pilsētai un tagadējās Polijas Nacionālās bankas priekštecei. Miera laikā šai Varšavas citadelē atradās 5~000 karavīru. Cietokšņa lielgabali bija notēmēti uz pilsētu. Imperators draudēja ar tiem nolīdzināt Varšavu no zemes virsas mazākās nepakļaušanās gadījumā. Cietoksnis tika izmantots arī kā ieslodzījuma vieta. Tā cellēs varēja turēt līdz 3~000 ieslodzīto. Citadelē tika izpildīti arī nāves sodi. Apmeklējot 1835.~gadā pilsētu, Nikolajs I atteicās pieņemt pilsētnieku delegāciju ar padevības apliecinājumiem, paziņojot, ka tā viņš atbrīvo tos no nepieciešamības melot.

Tiesa, Nikolajs I tā arī nespēja realizēt savu nodomu radīt Polijas karalistē spēcīgu krievu ierēdniecības slāni. Neraugoties uz aicinājumiem kalpot tēvijai Polijā, citiem imperatora gribas izpaudumiem, krievi tā īsti nespēja iedzīvoties poļu sabiedrībā, bija pārstāvēti galvenokārt Krievijas armijas sastāvā.

Turpretī Pēterburgā gan pastāvēja daudzskaitlīga poļu kolonija, kura kļuva par poļu-krievu sabiedrisko sakaru uzturēšanas centru.

Ir jāuzver, ka nebūt ne visi poļu izcelsmes aristokrāti, šļahtiči bija nelokāmi Krievijas valdīšanas Polijā pretinieki. Daudzi poļu muižniecības pārstāvji uzticīgi kalpoja Krievijas imperatoram un Polijas karalim. XIX~gadsimta 50.~gados Pēterburgas ierēdniecībā poļu īpatsvars sastādīja 6\%. To, ka daudzi poļi ieguva izglītību un izcilu stāvokli sabiedrībā, pateicoties dzīvei Krievijā, patriotiski noskaņotie poļi ilgstoši nevēlējās redzēt. Viņi uzskatīja, ka, piemēram, Pēterburgā, pozitīvi vērtējams polis varēja atrasties vai nu kā ieslodzītais, kā T.~Kostjuško, vai kā troni zaudējis monarhs, kā S.~Poņatovskis. Patiesībā Pēterburgā pastāvīgi dzīvoja poļi~--- Krievijas Valsts padomes locekļi (vēlāk, jau 20.~gadsimtā, arī Krievijas Valsts domes deputāti), poļu aristokrātijas pārstāvji, galma titulu nesēji, gvardes pulku virsnieki un Kadetu korpusa audzēkņi, tātad cilvēki, kuri jau bija dzīvē ko sasnieguši vai arī gatavojās kāpt pa karjeras kāpnēm.

No izcilā poļu dzejnieka Ā.~Mickeviča biogrāfijas Polijā parasti atceras viņa ieslodzījumu Viļņā, izsūtīšanu uz Krieviju un gadus, pavadītus emigrācijā. Taču vairāk nekā gadu viņš nodzīvoja Pēterburgā, apmeklējot vietējos salonus, uzturot kontaktus ar krievu sabiedriskajiem un kultūras darbiniekiem. Kaut Varšava bija visai dinamisks kultūras centrs, kurš saistīja arī daudzus krievus, tomēr salīdzinājumā ar Pēterburgu visu XIX gadsimtu tā bija kultūras province. Poļu inteliģenti, kuri visai sarkastiski atsaucās par visu krievisko, ar pasaules literatūras sasniegumiem bieži iepazinās, pateicoties labiem un lētiem to tulkojumiem krievu valodā, kuri tika iespiesti impērijas galvaspilsētā. Kā noskaidrojis poļu zinātnieks L.~Baziļevs, no XVIII gadsimta beigām līdz Pirmā pasaules kara sākumam Pēterburgā vien īsāku vai garāku laiku (no dažām dienām līdz visai savai dzīvei) pabija apmēram 1/4 miljons poļu. Tiek uzskatīts, ka XIX gadsimta 40.~gados katrs trešais no Pēterburgas universitātes 750 studentiem varēja būt polis no Lietuvas vai Polijas karalistes. Viņi uzskatīja sevi par Krievijas monarha pavalstniekiem, realizēja viņa impērijas politiku, darbojās apzinīgi un uzcītīgi, saņēma par to apbalvojumus un paaugstinājumus amatos, tai pat laikā nezaudējot savu poļu apziņu un ieinteresētību poļu stāvokļa uzlabošanā, bieži palīdzot poļu iestādēm gan finansiāli, gan ar savu ietekmi. XIX~gadsimta vidū 3\% no visiem Krievijas ierēdņiem bija poļi.

Samērā daudz šļahtiču dienēja Krievijas armijā. 1862.~gadā Krievijas armijas virsnieku vidū pareizticīgo bija 69,37\%, katoļu~--- 20,06\% un protestantu~--- 9,33\%. Statistika pēc tautībām netika vesta, taču katoļi gandrīz visi bija poļi, bet protestanti~--- vācieši. Tātad katrs piektais imperatora armijas virsnieks bija polis, kaut Krievijas iedzīvotāju vidū poļi sastādīja mazāk nekā 6\%. Kā raksta poļu vēsturnieki, nevar noliegt, ka daļa poļu patriotu dienestu valstu~--- Polijas pakļāvēju~--- armijās uzskatīja par iespēju apgūt karavīra prasmes, kuras pēc tam varētu izmantot nacionālās lietas labā. Virkne poļu izcelsmes virsnieku~--- Z.~Serakovskis, Z.~Padļevskis, R.~Trauguts u.c.~--- savas militārās zināšanas un prasmes vēlāk lika lietā pret Krieviju 1863.~gada sacelšanās laikā, par ko pēc sacelšanās sakāves arī tika pakārti vai nošauti.

Kā poļu revolucionāra dzīves gājuma piemēru var uzskatīt Z.~Serakovska biogrāfiju. Viņš bija dzimis nabadzīga šļahtiča ģimenē Volīnijā, mācoties Pēterburgas universitātē, iesaistījās revolucionāra studentu pulciņa darbībā, tika arestēts un nosūtīts karadienestā. Pēc astoņiem dienesta gadiem viņš saņēma virsnieka pakāpi, atgriezās Pēterburgā, beidza Ģenerālštāba akadēmiju, dienēja Kara ministrijā. Tās uzdevumā pabija komandējumos ārzemēs, kur iepazinās ar Dž.~Garibaldi, poļu emigrantiem, krievu revolucionāro demokrātu A.~Hercenu. Krievijā viņš uzturēja sakarus ar revolucionārajiem demokrātiem N.~Černiševski un N.~Dobroļubovu. 1859.--1860.~gadā Z.~Serakovska vadībā tika izveidots Krievijas Ģenerālštābā dienošo poļu virsnieku pulciņš, kura darbībā piedalījās arī poļu civilierēdņi un studenti. Pēc sacelšanās Polijā sākuma viņš palūdza atvaļinājumu ārzemju apmeklējumam, bet pats devās uz Lietuvu, kur kļuva par vienu no sacelšanās vadītājiem.

Taču bija arī pretēji, uzticīga dienesta Krievijai piemēri.

Lai nosaucam dažus no tiem. S.~Grabovskis XVIII~gadsimta beigās un XIX~gadsimta sākumā piedalījās vairākos poļu karos pret Krieviju, pabija arī izsūtījumā Sibīrijā, pēc apžēlošanas Napoleona armijas sastāvā atkal iesaistījās karagājienā pret Krieviju, līdz 1813.~gadā kaujā pie Leipcigas krita krievu gūstā. Kopš 1815.~gada viņš dienēja Polijas karalistē, bija tuvs lielkņazam Konstantīnam. 1825.~gadā dekabristu sacelšanās laikā, būdams Pēterburgā, spēja nodemonstrēt savu uzticību Nikolajam I, kļuva par Polijas karalistes valsts sekretāru. 1830.~gada sacelšanās laikā S.~Grabovskis to neatbalstīja, bet kļuva par Krievijas ieceltās Polijas Pagaidu valdības un Krievijas Valsts padomes locekli.

F.~Druckis-Lubeckis piedzima Pēterburgā un, neraugoties uz poļu izcelsmi, palika vienmēr uzticīgs Krievijas monarham un iestājās par Polijas un Krievijas vienošanos. Viņš piedalījās A.~Suvorova vadītajos karagājienos uz Itāliju un Šveici. Napoleona karagājiena uz Krieviju laikā F.~Druckis-Lubeckis dienēja Krievijā un bija atbildīgs par tās armijas apgādi. 1813.--1815.~gadā bija Krievijas okupētās Varšavas hercogistes Pagaidu augstākās padomes loceklis, 1821.~gadā kļuva par Polijas karalistes finanšu ministru, centās paplašināt karalistes autonomiju. 1830.~gada sacelšanās laikā viņš veda Nikolajam I sacēlušos prasības, taču pēc atbildes nogādāšanas Varšavā atgriezās Pēterburgā, ar to saglabājot imperatora uzticību. 1832.~gadā F.~Druckis-Lubeckis tika nozīmēts par Krievijas Valsts padomes locekli.

Ļoti daudzos XIX gadsimtā Krievijas vestajos karos drosmīgi cīnījās ģenerālis Ā.~Ževuzskis. Kad sākās 1830.~gada poļu sacelšanās, viņš, neraugoties uz savu polisko izcelsmi, saglabāja uzticību Krievijas imperatoram un aktīvi piedalījās sacelšanās apspiešanā, par ko tika apbalvots ar zelta zobenu, uz kura bija uzraksts ``Par drosmi'', un vairākiem ordeņiem. Arī 1863.--1864.~gadā Ā.~Ževuzskis, būdams Kijevas kara apgabala pavēlnieks, palika atstatus no sacelšanās Polijā.

Tādejādi nevajadzētu uzskatīt, ka jebkurš polis patriotisma vārdā bija gatavs cīnīties pret Poliju pakļāvušajām valstīm, aizmirstot par visu citu. Daļa no viņiem spēja saskatīt, ka Polijas karaliste sacelšanās rezultātā ne tikai zaudēja to uzplaukumu, kuru tā sasniedza ap F.~Drucki-Ļubecki saliedēto darbinieku ekonomiskās politikas rezultātā, bet arī savu Konstitūciju un vismaz nominālo patstāvību, kļuva par Krievijas provinci, kaut arī zināma autonomija vēl saglabājās. Kā rakstīja poļu vēsturnieks O.~Haleckis, Polijas īpašais valstiskais statuss tika nomainīts ar ``simulētu autonomiju''.

1837.~gadā poļu vojevodistes pielīdzināja Krievijas guberņām, kuru priekšgalā atradās iecelti gubernatori. Skolu lietas tika nodotas tieši Tautas Izglītības ministrijas Pēterburgā pārziņā. Tika slēgta Varšavas un Viļņas universitātes. (Interesanti, ka Nikolajs I, apzinoties katoļu reliģijas nepieciešamību kārtības uzturēšanai, šo augstskolu teoloģijas fakultātes saglabāja kā atsevišķas mācību iestādes.) Īpaši otrās slēgšana negatīvi ietekmēja arī lietuviešu un baltkrievu sabiedriski politisko un nacionāli kulturālo dzīvi. Samazinājās ģimnāziju skaits. Tika ieviesta stingra cenzūra. Bija aizliegts ne tikai drukāt poļu dzejnieku un citu kultūras darbinieku darbus, bet arī pieminēt viņu vārdus.

Vēloties mazināt poļu ietekmi uz citām slāvu tautām, Krievijas valdība centās ierobežot katoļu bīskapiju skaitu. Saskaņā ar draudžu skaita krišanos samazinājās arī katoļu baznīcu skaits. Katoļu draudzi varēja nodibināt, ja tā aptvēra 100 ēku un 400 ticīgo. Pēc 1830.--1831.~gada sacelšanās valdība spēra soļus, lai vēl vairāk ierobežotu katoļu garīdznieku, kuru vairākums netieši bija atbalstījuši sacelšanos, ietekmi. Viņiem aizliedza atstāt savas draudzes bez speciālas atļaujas. Baznīcu zemes īpašumus konfiscēja. Pieauga kā pareizticīgās, tā katoļu baznīcas atkarība no valsts. 1841.--1843.~gadā tās nonāca valsts apgādībā. 1842.~gadā garīdzniekiem noteica algas. Ar laiku arvien redzamāka kļuva valsts labvēlība tieši pareizticīgajai baznīcai. 1839.~gadā Polockā (krievu \rutxti{Полоцк}, baltkrievu \rutxti{Полацк}) notika uniātu baznīcas koncils, kas nolēma pievienot uniātu draudzes pareizticīgo baznīcai. Sākās aktīva uniātu draudžu pievienošana pareizticīgajām.

Arī Polijas karalistei blakus esošajās guberņās (Grodņas, Kauņas, Minskas, Viļņas, Vitebskas) tika ievadīti rusifikācijas pasākumi. Pēc Krievijas piemēra Prūsija un Austrija, kuras atbalstīja Krievijas pasākumus pēc sacelšanās apspiešanas, tagad varēja arī pašas daudz asāk vērsties pret poļu nacionālajām prasībām. Tā Prūsijas varas iestādes izdeva Krievijai lielu daļu poļu karavīru, kuri pēc sacelšanās bija nonākuši tās teritorijā. Poļi~--- Prūsijas pavalstnieki, kuri bija karojuši sacēlušos rindās, tika neklātienē notiesāti.

Ar poļu 1830.--1831.~gada sacelšanās apspiešanu notikušais pavērsiens Krievijas impērijas politikā Polijā izpaudās arī Krievijas sabiedriskajā dzīvē. 1831.~gadā tapa A.~Puškina dzejolis ``\rutxti{Клеветникам России}'' (``Krievijas apmelotājiem''). Lūk, fragments no tā:

\vspace{1.5em}

\noindent
\begin{minipage}{0.55\textwidth}
\pltxti{Уже давно между собою\\
Враждуют эти племена;\\
Не раз клонилась под грозою,\\
То их, то наша сторона.\\
Кто устоит в неравном споре:\\
Кичливый лях иль верный росс?\\
Славянские ль ручьи\\
сольются в русском море?\\
Оно ль иссякнет? Вот вопрос.}
\end{minipage}
\hspace{1em}
\begin{minipage}{0.6\textwidth}
Sen sens ir naids\\
Starp ciltīm šīm;\\
Ne reizi vien ir nācies vētrā liekties,\\
Gan viņējai, gan mūsu pusei.\\
Kurš uzvarēs šai strīdā nevienlīdzīgā:\\
Vai lepnais poļu pans vai uzticamais krievs?\\
Vai krievu jūrā\\
slāvu strauti saplūdīs?\\
Vai izsīks tā? Lūk, jautājums.
\end{minipage}

\vspace{1.5em}

% page 128




Dzejolī atklājās dzejnieka domu gaita: gadsimtiem ilgās cīņās Krievija bija pieveikusi Poliju. Tāpēc Krievijas pavalstniekiem arī turpmāk bija jāpaliek ``uzticamiem krieviem''~--- sava cara kalpiem, bet ``lepnie poļu pani'', kas priekšroku deva savām feodālajām privilēģijām, faktiski bija zaudējuši savu tautu, atdodot to svešzemnieku (prūšu un austriešu) kundzībā, no kā bija tikai viens glābiņš~--- ``ieplūst krievu jūrā''.

Nebija nejaušība arī tā, ka 1836.~gadā pirmuzvedumu piedzīvoja krievu komponista M.~Gļinkas opera ``Dzīvību par caru'' (\rutxti{Жизнь за царя}), vēlākais nosaukums~--- ``Ivans Susaņins'' (\rutxti{Иван Сусанин}). Operas darbība notiek 1613.~gadā, kad ``juku laikos'' ieslīgušajā Krievijā tronī nāca Romanovu dinastijas dibinātājs Mihails. Tā vēstīja par krietno krievu zemnieku I.~Susaņinu, kurš, samaksājot ar savu dzīvību, izglāba caru no poļu iebrucējiem. Jau 1862.~gadā krievu-ukraiņu vēsturnieks N.~Kostomarovs šo faktu apstrīdēja. Arī mūsdienu literatūrā norādīts, ka leģendu neapstiprina zinātniski dati, tomēr Krievijā tā joprojām tiek izmantota patriotisma audzināšanā. M.~Gļinkas estētiskās koncepcijas pamatā bija varonīgo krievu un nelietīgo poļu pretnostatījums, ko atspoguļo gan divi kori~--- poļu un krievu, gan divi kontrastējoši scenogrāfijas un mūzikas stili. Poļus raksturo dižmanība, viņi dzied tikai korī un dejo ceremoniālas masu dejas polonēzes, mazurkas un krakovjaka ritmā. Turpretī krievi dzied vai nu jaukas tautasdziesmas, vai romantiskas ārijas. Pēc I.~Susaņina nogalināšanas epiloga kulminācija ir aina ar Mihaila Romanova triumfālo ierašanos Maskavā Sarkanajā laukumā. Ar M.~Gļinkas darbu cariskajā Krievijā atklāja ikvienu operas sezonu Maskavā un Sankt-Pēterburgā. To 1866.~gadā iestudēja Prāgā čehu valodā, 1878.~gadā~--- Rīgā latviešu valodā un 1899.~gadā~--- Pozenes Vācu teātrī Prūsijā. Taču nekad to neuzveda Varšavā vai Krakovā.

Sakarā ar sacelšanās apspiešanu Polijā Krievijas sabiedrībā iezīmējās negatīvas tendences. Ne bez pamata Nikolajam I tuvi cilvēki sūdzējās, ka ``poļu revolūcija aizturēja Krievijas attīstību par 30~gadiem''. Krievijas slavofīli uz Poliju raudzījās visai aizdomīgi, cīnītājus par Polijas neatkarību vērtēja kā Krievijas un visu slāvu tautu vienotības ienaidniekus. Arī vēsturnieks A.~Pogodins uzskatīja, ka Krievija, kura jau bija izkļuvusi no Austrijas un tās ārlietu ministra, viena no Svētās savienības arhitektiem K.~Meterniha ietekmes un tuvinājās Anglijai, atkal tika iesviesta Svētās savienības apkampienos. Pēc poļu sacelšanās Nikolajs I pulcēja ap sevi konservatīvi noskaņotus cilvēkus. Šajā laikā, kā rakstīja vēsturnieks, ``Krievijā gāja bojā Puškins,~Ļermontovs, Gogolis, apklusa jebkura brīvdomība, zēla korupcija''.

Pašā Polijas karalistē pēc 1830--1831.~gada sacelšanās nacionālās atbrīvošanās kustības veidols ievērojamā mērā demokratizējās. No 1835.līdz 1845.~gadam kustībā ievērojami auga tā slāņa īpatsvars, kuru poļu vēsturnieki visbiežāk sauc par inteliģenci: sīkie ierēdņi, ārsti, skolotāji, pārvaldnieki, t.s. brīvo profesiju pārstāvji, studenti. Tas bija tipisks pārejas perioda slānis. Tā atvērtība un augošā loma padarīja garīgo darbu pievilcīgu kā nabadzībā grimstošajai šļahtiču daļai, kura centās pretoties savai noslīkšanai sabiedrības padibenēs, kā arī pilsētniekiem, kuriem izglītības iegūšana solīja vilinošu perspektīvu uzkāpt pa sabiedriskās hierarhijas kāpnēm. Runājot par poļu inteliģenci kopumā, jāuzsver, ka materiāli nodrošināto slānis tās rindās bija ļoti plāns, pārsvarā bija grupas, kuras knapi varēja iztikt, nebija pārliecinātas par savu rītdienu, saistīja savu nākotni ar politiskām pārmaiņām.

Pēc dažiem datiem no 1835. līdz 1841.~gadam šī inteliģences slāņa īpatsvars sastādīja 43\%, bet no 1842.~līdz 1845.~gadam~--- 61\% no visiem atbrīvošanas kustības dalībniekiem. Kustība bija ļoti raiba pēc sava sastāva, arī tās ideoloģija bija pretrunīga. Taču nacionālo prasību jomā visas politiskās grupas un partijas prasīja atjaunot Poliju 1772.~gada robežās. Tomēr arī šeit vienotība bija pārsvarā ārēja. Vieni gribēja atjaunot veco Žečpospolitu, citi tās teritorijā cerēja nodibināt demokrātisku republiku, kurai būtu jāveic radikāli sociāli pārveidojumi, bet pēc tam jādod ukraiņiem, baltkrieviem un lietuviešiem pašiem lemt savu likteni.

Poļu zemēs darbojās virkne nelegālu organizāciju.

40.~gadu sākumā Varšavā darbojās konspiratīva organizācija, saistīta ar publicistu E.~Dembovski, kurš minēts jau sakarā ar notikumiem Galīcijā. Lietas izmeklēšanas gaitā tā ieguva nosaukumu ``Demokrātiskā biedrība''. Nosaukums nebija nejaušs. Biedrības programmatiskie dokumenti nav saglabājušies, taču tās dalībnieku izteikumi ļauj spriest, ka par savu mērķi tā stādīja Polijas valstiskuma atjaunošanu, sociāli-ekonomisko attiecību pārveidi.

Vienu no konspiratīvajām poļu organizācijām, kura pastāvēja 1840.--1844.~gadā, vadīja zemnieku izcelsmes katoļu garīdznieks P.~Scegennijs. No tiem šīs organizācijas biedriem, kuri tika saukti pie atbildības par pretvalstisku darbību, apmēram puse bija zemnieki, starp pārējiem bija divi muižnieku dēli, daži muižu rentnieki un amatnieki, pārējie~--- sīki ierēdņi, kā arī skolotāji, privātkalpotāji, ģimnāzisti, studenti.

Organizācijas politiskā programma saturēja prasības pēc tautvaldības un visu līmeņu pārvades iestāžu ievēlējamības. Nākotnes valstiskumu P.~Scegennijs un viņa līdzgaitnieki iedomājās kā daudznacionālu republikānisku federāciju, kurā vadošā vieta piederētu poļu tautai. Federācijas robežas bija paredzēts noteikt pa Oderas, Nisas, Daugavas un Dņepras upēm. Tika uzskatīts, ka federācijai pievienosies arī Austrijas impērijā dzīvojošie slāvi.

P.~Scegennija un viņa līdzgaitnieku sociālā programma paredzēja feodāli-dzimtbūtnieciskās kārtības, kārtu, nacionālās un reliģiskās nelīdztiesības likvidēšanu, zemes pasludināšanu par visas tautas īpašumu, izdalot lietošanā ar mantošanas tiesībām saimniecības tiem zemniekiem, kuri tās apstrādāja. Organizācijas dalībnieki bija pārliecināti, ka viņu prasību īstenošana radīs sabiedrību, kurā visi būs laimīgi, un uz mūžīgiem laikiem valdīs brīvība, vienlīdzība un brālība. Faktiski šī programma zem utopiskā sociālisma formas izteica darbaļaužu, pirmkārt zemnieku, intereses buržuāziskās revolūcijas gaitā. Tā atspoguļoja tās izmaiņas, kuras notika poļu atbrīvošanās kustībā. Pat cara vietvaldis I.~Paskēvičs par 40.~gadu konspiratoriem rakstīja, ka viņi ``stāda mērķi ne tikai gāzt krievu valdību Polijā, bet reizē iznīcināt arī poļu šļahtičus-muižniekus.''

Poļu muižniecības augšslāņi bija vislabāk piemērojušies valdošajai kārtībai. Krievu vēsturnieks N.~Bergs rakstīja, ka vietvaldis I.~Paškevičs, kurš 1856.~gada sākumā šķīrās no dzīves, ``atstāja Poliju pilnīgi mierīgu \citespace{} Aristokrāti apspieda visu, bija sava veida vietvalža līdzvaldnieki. Ko tikai toreiz nevarēja poļu tiesās ietekmīgs poļu aristokrāts, īpaši ievērojot poļu ierēdņu neparasto korumpētību! Ar aristokrātu palīdzību Paškēvičs darbināja visas savas atsperes \citespace{} viņš uzskatīja šādu pārvaldes metodi ne tikai par labāko, bet arī par neizbēgamu, ievērojot vispārējo lietu stāvokli impērijā un toreizējos vispārējos uzskatus.'' Pēc tā paša N.~Berga domām nākamais vietvaldis kņazs M.~Gorčakovs aristokrātiem deva pārāk lielu vaļu un neizrādīja pietiekamu stingrību pret poļiem.

Taču daudzi šļahtiči turpināja pretestību carismam, devās emigrācijā. Pēc poļu publicista un vēsturnieka L.~Vasiļevska datiem pēc 1831.~gada emigrācijā atradās ap 70~tūkstošu sacelšanās dalībnieku un atbalstītāju. Pat neraugoties uz cara izsludināto amnestiju, daļa 1830.~gada sacelšanās dalībnieku, turpināja dzīvot \strong{emigrācijā} (galvenokārt Šveicē un Francijā), ko poļu vēsturiskajā literatūrā sauc par ``lielo''. Emigrācijā dzīvoja dzejnieks Ā.~Mickevičs, komponists F.~Šopēns u.c. Pēc citiem, vācu vēsturnieku datiem Eiropā darbojās līdz 10~000 poļu emigrantu. 1847.~gadā tikai Francijā vien bija 4~790, bet Anglijā ap 800 iebraukušo no Polijas. Poļu historiogrāfija pirmkārt norāda uz šīs emigrācijas garīgo nozīmīgumu, bet ne tās skaitu. Emigranti bija galvenokārt virsnieki un inteliģences pārstāvji, cēlušies no šļahtičiem, daļēji arī no pilsonības. Poļu nacionālās kustības smaguma centrs pārvietojās uz ārzemēm, īpaši Franciju. Lielā emigrācija padziļināja poļu patriotisma antikrievisko, šovinistisko raksturu. Eiropas demokrātiskās sabiedrības atbalsts, liberāļu simpātijas daudzos emigrantos stiprināja apziņu, ka poļiem lemts kā savulaik Jēzum Kristum kļūt par mesiju apspiesto tautu atbrīvošanās kustībā.

Emigranti, kaut arī dalījās aristokrātiskajā un demokrātiskajā spārnā, lielākoties bija gatavi piedalīties katrā Eiropas valstīs notiekošā revolucionārā pasākumā. Taču drīzumā viņi arī organizatoriski nodalījās divās grupās: t.s. ``baltajos'' (mērenajos) un ``sarkanajos'' (demokrātos). Abas grupas vēlējās panākt Polijas valstisku neatkarību, taču tās sasniegšanas ceļus redzēja savādāk.

Vislabāk materiāli nodrošināts un vispopulārākais ārvalstu valstsvīru vidū bija \strong{``balto''} līderis, jau vairākkārt minētais kņazs Ā.~Čartorijskis, kurš vadīja emigrācijas konservatīvo spārnu. Tas pulcējās ap viņa rezidenci \frtxti{Hôtel Lambert} (grezna savrupmāja, celta XVII~gadsimtā, kuru nopirka Ā.~Čartorijska sieva). 1843.~gada martā Ā.~Čartorijska vadībā konstituējās ``\pltxti{Towarzystwo Monarchicne Fundatorow i Przyjaciόl Trzeciego Maja''} (Monarhistiskā 3.~maija [domāta 1791.~gada 3.~maija Konstitūcija] dibinātāju un draugu biedrība)~--- konservatīvo vadošais politiskais kodols. Pēc viņu domām Polijas jautājumu varēja risināt tikai visas Eiropas krīzes gadījumā. Šīs politikas būtību izteica Ā.~Mickeviča vārdi: ``\pltxti{O wojnę powszechną za wolność ludόw / prosimy się Panie}'' (``Vispārēju karu par tautu brīvību / lūdzam Tev, Kungs''). Poļu sacelšanās būtu dienas kārtībā, tikai pastāvot labvēlīgiem apstākļiem. Līdz tam par galveno uzdevumu tika uzskatīts Eiropas lielvalstu valdošajās aprindās uzturēt poļu lietai labvēlīgu noskaņojumu. Padomju vēsturnieki rakstīja, ka ``baltie'' un \frtxti{Hôtel Lambert} aprindas vēlējās, lai sacelšanās Polijā ``gruzdētu'', kas rādītu Rietumeiropas valstu valdībām, ka poļu tauta vēl ir dzīva un aizstāv savas nacionālās un reliģiskās tiesības. Ā.~Čartorijska vadībā faktiski tika uzturētas neoficiālas diplomātiskas attiecības ar Lielbritānijas, Francijas un Turcijas valdībām. Reizē tika mēģināts turpināt poļu leģionu tradīciju. Taču mēģinājumi to nodibināt Beļģijā, Portugālē, Ēģiptē un Turcijā beidzās nesekmīgi. 1834.~gadā ap 1~000 poļu iestājās franču Ārzemnieku leģionā, taču jau nākamajā gadā pārstrukturēšanas rezultātā tā vienības zaudēja ``nacionālo'' raksturu.

No otras puses, ``baltie'' emigranti centās nodalīt jautājumu par nacionālo atbrīvošanos no jautājuma par sociālo problēmu revolucionāru atrisināšanu. Ā.~Čartorijska piekritēji izvirzīja lozungu: ``Vispirms būt, pēc tam~--- kā būt''. Visi vēlākie notikumi rādīja, ka šāda pieeja tikai traucēja atrisināt jautājumu ``būt vai nebūt''.

\strong{``Sarkanie''} jeb nosacīti ``demokrāti'' skaita ziņā bija pārsvarā, taču baudīja mazāku starptautisko atbalstu. Viņu pirmā organizācija bija 1831.~gadā Parīzē nodibinātā ``\pltxti{Komitet Narodowy Patriotyczne}'' (``Patriotiskā nacionālā komiteja'') jau minētā vēsturnieka un politiķa J.~Lelevela vadībā. Viņš norādīja, ka pastāv divas Eiropas; tautas masu Eiropa, kura cīnās pret feodālismu un despotismu, un reakcionāro valdību Eiropa, un aicināja visus poļus sadarboties ar pirmo.

No poļu emigrantu organizācijām skaitliski pati lielākā bija ``\pltxti{Towarzystwo Demokratyczne Polskie}'' (``Poļu demokrātiskā biedrība'', 1832--1862). Tai bija savi pārstāvji Polijā. 1836.~gadā tā izdeva savu pamatprogrammu~--- ``\pltxti{Wielki Manifest}'' (``Dižais manifests''), kurš pasludināja ``tautas revolūcijas'' lozungu, ko varētu īstenot galvenokārt zemniecības spēkiem, bet saistībā ar visas Eiropas revolūciju, feodālo pienākumu likvidēšanu un zemniekiem iedalīto zemes platību nodošanu bez atlīdzības to īpašumā. Tā sludināja šķiru mieru poļu nācijas iekšienē. Tiktu likvidētas visas privilēģijas, ieviesta brīvība un vienlīdzība, atgūtas Žečpospolitas 1772.~gada robežas. Taču reizē deklarācijā bija izcelta Polijas īpašā loma Eiropā. ``Polija pagātnē aizsargāja Rietumus pret barbariskajiem tatāru, turku un ``moskaļu'' iebrukumiem. Polija gāja bojā tāpēc, ka tad, kad Rietumos cilvēku brīvības doma pasludināja karu vecajai kārtībai, par kuras aizstāvībai iestājās krievu despotisms, Polija, pildot savu vēsturisko misiju, iesaistījās cīņā pret šo spēku, taču tika uzvarēta. Eiropas glābšana tika atlikta.'' No sacītā tika secināts, ka Polijas glābšana nav tikai pašas Polijas, bet visas cilvēces lieta. Deklarācija rādīja, ka tās autori Polijas misiju uzskatīja būt par sargu Eiropas Austrumos. Tatāru un turku briesmas nebija vairs aktuālas, tātad Eiropu bija jāglābj no ``moskaļiem''.

Manifesta vājā vieta bija cerības uz brīvprātīgu šļahtas atteikšanos no savām privilēģijām. Tas aicināja griezties pie šļahtas ``ne ar ercenģeļa zobenu, bet ar tēvzemes vēstures grāmatu rokās''. Protams, aicinājums no tā izskanēšanas brīža bija nolemts neveiksmei. Jau 1846.~gada sacelšanās Krakovā parādīja, ka nav iespējams samierināt šļahtiču--muižnieku un zemnieku intereses. Taču ``Poļu demokrātiskā biedrība'' neatteicās no savām ilūzijām, vienīgi aktivizēja kārtu solidaritātes idejas paušanu, ``nācijas vienības'' vārdā. Tai pat laikā arī organizācijas iekšienē pastāvēja domstarpības, daļa biedru bija liberāļi vai mēreni demokrāti, citi~--- sociālisti-utopisti.

Emigrācijas revolucionārais spārns vienojās organizācijā ``\pltxti{Lud Polski}'' (``Polijas tauta'', 1835--1846, darbojās galvenokārt Lielbritānijas teritorijā), kura uz laiku pat izstājās no Demokrātiskās biedrības un iestājās par feodālisma, kārtu nevienlīdzības, zemes privātīpašuma likvidēšanu zemnieku sacelšanās ceļā. Polijas valstiskuma atjaunošana tika saistīta ar tādas iekārtas, kur tikai priviliģēto saujiņa izmantoja ``visu kopīgo zemi un tās augļus'' likvidēšanu, sludināta cīņa par demokrātiju, kura padarīs visus poļus ``bez vārda, dzimtas un reliģijas atšķirības'' brīvus un laimīgus. Poļu sociālisti-utopisti, kaut pieslējās t.s. ``kristīgā sociālisma'' ievirzei, atšķirībā no šo ideju atbalstītājiem Rietumos aicināja iznīcināt feodālismu vardarbības ceļā. Tomēr 1846.~gada Krakovas revolūcijas mācība jau 1848.~gadā lika demokrātiem darboties ar liberāļiem vienotās organizācijās.

Šajā laikā poļu vidū joprojām bija populārs lozungs ``\pltxti{Za naszą i waszą wolność}'' (Par mūsu un jūsu brīvību), kurš aicināja arī krievu revolucionārus uz kopīgu cīnu pret carismu. Toties Rietumu revolucionāri gandrīz nespēja iedomāties kopīgu krievu un poļu cīņu pret carismu, kas tai pavērtu jaunas perspektīvas. Tā, F.~Engelss 1851.~gada 23.~maijā rakstīja: ``Ja krievus var pamudināt uz aktīvu rīcību, tad jāveido savienība ar viņiem un jāpiespiež poļus piekāpties''.

Kā jau teikts, 1846--1848.~gadā Polijā atkal notika sacelšanās mēģinājumi, poļi piedalījās sacelšanās Austrijā, Ungārijā, taču Polijas karalistē nemieru gandrīz nebija.

Pēc Krievijas imperatora Nikolaja I nāves un Aleksandra II nākšanas tronī 1856.~gadā sākās zināma liberalizācijas fāze carisma politikā Polijā. Taču, kad imperators Aleksandrs II, ar kuru visas Eiropas liberāļi saistīja lielas cerības, 1856.~gadā ieradās Varšavā un poļu sabiedrības pārstāvji iesniedza viņam petīciju ar lūgumu atjaunot Polijas autonomiju un pārtraukt ierēdņu patvaļu, viņa atbilde bija kā auksta ūdens šalts: ``Jūs esat tuvi manai sirdij tāpat kā somi un citi Krievijas pavalstnieki, taču es nevēlos, lai tā kārtība, ko noteica mans tēvs, tiktu mainīta. Tāpēc, nekādu ilūziju, kungi! Es pratīšu apstādināt to centienus, kuri iedomāsies pakļauties sapņiem. Es pratīšu rīkoties tā, ka šie sapņi nepārkļūs pār sapņotāju iztēles robežu. Polijas laime slēpjas tās pilnīgā saplūšanā ar manas Impērijas tautām. Viss, ko paveica mans tēvs, ir paveikts labi.'' Sabiedrībā īpaši tika uzsvērti viņa vārdi: ``\frtxti{Point de rệveries, Messiers}'' (franču: Nekādu ilūziju, kungi!) Tiesa, jau nākamajā dienā pēc citētās runas Aleksandrs II parakstīja manifestu par politieslodzīto poļu amnestiju un atļauju poļu emigrantiem atgriezties dzimtenē. Tiesības atgriezties no ārzemēm ieguva gan tikai ``grēkus'' nožēlojušie, kuri saņēma attiecīgu apliecinājumu no Krievijas konsulārajām iestādēm. Četru gadu laikā Polijā atgriezās ap deviņi tūkstoši izsūtīto un emigrantu. Pēc N.~Berga vērtējuma tā ``sarkano'' pulciņi Polijas karalistē ieguva vadītājus, kuri ieradās no Sibīrijas un Eiropas, bet ``sarkano'' ietekmē aktivizējas arī ``balto'' pulciņi. Tika atcelts kara stāvoklis, Varšavā atvērta Medicīniski-ķirurģiskā akadēmija, dažādas svētdienas skolas un arodskolas. Tika mīkstināta cenzūra. Varēja izdot agrāk aizliegtu rakstnieku darbus. Piemēram, Varšavā izdeva dažus 1855.~gadā mirušā Ā.~Mickeviča darbus. Iznāca jaunas avīzes un žurnāli. Atļāva kultūrizglītības un ekonomisko biedrību darbību.

Starptautiskajā līmenī Polijas jautājums neoficiāli tika skarts miera konference pēc Krievijai neveiksmīgā Krimas kara. Anglija un Francija samierinājās ar Aleksandra II īstenoto agrāk sodīto poļu amnestiju, citas prasības neizvirzīja.

\strong{Saimniecība Polijas karalistē XIX gs. vidū}, neraugoties uz nacionālo apspiešanu, attīstījās ļoti sekmīgi, straujiem tempiem. Krievijas tirgus pavēra Polijas tautsaimniecībai lielas iespējas. Polijas karaliste piesaistīja amatniekus ar izdevīgajiem apmešanās noteikumiem un atbrīvošanu no nodokļiem.

Jau 1847.~gadā tika atklāta tvaikoņu kustība pa Vislu. Kopš XIX gadsimta 40.~gadiem Polijas karalistē tika būvēts dzelzceļš. Tieši Polijā pēc dzelzceļa līnijas Pēterburga~--- Carskoe selo būves 23~km garumā kā otrā Krievijā tika būvēta līnija no Varšavas līdz Austrijas robežai 308~km garumā. To uzbūvēja 1845.--1848.~gadā. Šis ceļš savienoja Varšavu ar Dombrovas ogļu ieguves baseinu, kas veicināja strauju ogļraktuvju attīstību. 1848.~gadā Varšavu savienoja ar Vīni. Kaut, vadoties no stratēģiskiem apsvērumiem, tālākā dzelzceļa celtniecība notika diezgan lēni un tā tīkls, salīdzinot ar citām Eiropas zemēm, bija samērā rets, tomēr Krievija tirgus deva lieliskas attīstības iespējas tekstilrūpniecībai Lodzas apgabalā un smagajai rūpniecībai Augšsilēzijā. 1862.~gadā tika pabeigta Pēterburgas--Varšavas un Skernevices (\pltxti{Skierniowice})~--- Brombergas (poļu \pltxti{Bydgoszcz}, vācu \detxti{Bromberg)} dzelzceļa līnija. 1871.~gadā Varšavu dzelzceļš savienoja ar Maskavu.

Uzņēmēji, kuri apņēmās uzbūvēt attiecīgos dzelzceļa posmus, muižnieki, kuri iedzīvojās no sliktas kvalitātes pārtikas pārdošanas strādnieku brigādēm, kuras virzījās pa jaunbūvējamo stigu, pat akcionāri, kuri, cerot uz peļņu, riskēja nopirkt dažas dzelzceļa sabiedrības akcijas, nevarēja pilnā mērā apjaust to labumu, ko jaunais transporta veids nesīs visai zemei. Parasti visi vēlējās iedzīvoties bagātībā uzreiz, vēl celtniecības laikā. Tāpēc notika vairāk vai mazāk fiktīvi celtniecības kompāniju bankroti, bija vērojama zema darbu kvalitāte, kura parasti atklājās tikai tad, kad pa jaunajiem uzbērumiem un sliežu ceļiem sāka kursēt vilcieni. Sliedes tika vestas no Anglijas, vagoni~--- no Vācijas. Tomēr pāris desmitgažu laikā notika milzu pārmaiņas. Dzelzceļš, kuram, kā likās, ienākumi bija jāsaņem galvenokārt no pasažieru transportēšanas, negaidīti sāka nest nesalīdzināmi lielāku peļņu no preču pārvadāšanas. Dzelzceļš tuvināja dažādus Polijas novadus, pat esošus dažādu valstu atkarībā.

Dzelzceļu celtniecībai Polijas karalistē gan bija ievērojams trūkums~--- tā tomēr nespēja pilnībā apmierināt rūpniecības vajadzības. Ar Prūsijai piederošo Polijas daļu Varšava bija savienota tikai četros punktos, Austrijai piederošo~--- tikai vienā. Vīnes, Brombergas un Lodzas līnijām bija Rietumeiropai parastais (šaurais) sliežu platums, pārējām~--- pēc Krievijas parauga~--- platais. Vienota standarta trūkums apgrūtināja pārvadājumus pat pašas Polijas karalistes robežās.

\strong{Buržuāzijas} jeb pilsonības veidošanās gaitā Polijas karalistē bija arī savas specifiskas īpatnības, ko noteica kapitāla sākotnējās uzkrāšanas process, kam lielā mērā bija raksturīga kapitāla gūšana no avotiem, kuri nebija saistīti ar pievienotās vērtības ražošanu. Jauno, kapitālistisko iekārtu simbolizējošā lielburžuāzija poļu zemēs veidoja mazskaitlīgu slāni un lielā mērā vēl izmantoja kapitāla sākotnējās (nekapitālistiskās) uzkrāšanas metodes. Tā Polijas karalistē XIX gadsimta pirmajā pusē vislielākie kapitāli piederēja nelielai komersantu grupai, kura atpirka no valsts dažādus tai piederošus monopolus~--- sāls, tabakas u.c. pārdošanas tiesības. Netiešo nodokļu iegūšanas tiesību atpirkšana, valsts pasūtījumi, valdības dotācijas, valsts uzņēmumu iegūšana uz izdevīgiem noteikumiem nomā~--- tādi bija visdrošākie līdzekļi straujai tikšanai pie bagātības. To zaudēšana draudēja ar gandrīz neizbēgamu bankrotu. Tas viss spieda buržuāziju censties iegūt cara valdības, īpaši finanšu iestāžu korumpēto ierēdņu labvēlību. Tas noteica arī šīs buržuāzijas daļas mazo interesi par nacionālās atbrīvošanās lietu. Desmitgadēs pēc Polijas dalīšanas poļu nacionālo varoņu rindās neizvirzījās neviens komersants vai rūpnieks, kuriem~--- bieži pamatoti~--- bija augļotāju un spekulantu reputācija. Tā kā ārpus agrārā sektora bija izjūtams kapitālu trūkums, komersantu un uzņēmēju vidū izvirzījās liela grupa muižnieku, kuri gan nereti rūpes par saviem uzņēmumiem uzticēja pārvaldniekiem, nākušiem no tirdzniecības un banku sektora.

Īpašu noti buržuāzijas un arī pilsētu iedzīvotāju sastāva izveidē ienesa ebreji, kuri, kā jau minēts, pēc senas tradīcijas bija juridiski nepilntiesīgi. Ebrejiem tiesības nepiešķīra pat samērā liberālās Varšavas hercogistes un Polijas karalistes Konstitūcijas. Gluži pretēji, XIX gadsimta 20.~gados tika pastiprināti aizliegumi viņiem dzīvot lauku apvidos, tāpēc 60.~gadu vidū ebreji jau sastādīja gandrīz pusi karalistes pilsētu iedzīvotāju. Tikai 1862.~gadā pēc A.~Veļepoļska, (par viņu tālāk tiks runāts plašāk), kurš vēlējās veicināt ``ebreju izcelsmes poļu'' veidošanos, iniciatīvas tika atcelti ierobežojumi ebrejiem, saistīti ar dzīves vietas izvēli, dažādu amatu ieņemšanu, papildus nodokļi tiem. Tomēr, neraugoties uz nodalīšanu no poļiem gan juridiskiem, gan ar kultūras un sadzīves šķēršļiem, ebreji ļoti aktīvi piedalījās Polijas pilsētu un miestu saimnieciskajā dzīvē. Amatniecībā, tirdzniecībā, kredīta operācijās viņu loma bija visai ievērojama, dažkārt pat noteicošā.

\strong{Strādniecība} jeb rūpnieciskais proletariāts galvenokārt veidojās no poļu un citu tautību amatniekiem, pilsētu plebsa, zemniekiem, daļas deklasētās šļahtas. Strādnieku darba un dzīves apstākļi bija neapmierinoši. Jau 1823.~gadā Varšavā notika pirmais streiks. 1848.~gadā neapmierinātība parādījās jau plašās strādnieku aprindās. Par strādnieku apziņu var spriest pēc tā, ka vēl 60.~gadu sākumā Lodzas audēji, tāpat kā angļu luditi, (angļu \entxti{luddites}~--- no viņu vadoņa N.~Luda angļu strādnieku grupa, kura XIX gadsimta sākumā, protestējot pret izmaiņām savā stāvoklī, kuras nesa rūpniecības apvērsums, bieži iznīcināja mašīnas un citas iekārtas) uz bezdarbu atbildēja ar mašīnu iznīcināšanu.

Attīstību bremzēja valsts muižu zemnieku klaušu darba izmantošana. XIX~gadsimta pirmajā pusē lai nodibinātu fabriku kapitālistam, kaut retāk nekā 18.~gadsimtā, tomēr bieži vien vajadzēja iegādāties muižu ar tās izejvielu un, galvenais, darbaspēka rezervēm. (Jāuzsver, ka ne Prūsijai, ne Austrijai piederošajās poļu zemēs rūpniecībā nepastāvēja valsts uzņēmumi, kuri izmantotu klaušu darbu. Galīcijā XIX gadsimta pirmajā pusē klaušu darbs tika izmantots tikai muižniekiem piederošajās manufaktūrās.) Klaušu darba izmantošana uzņēmumos strauji samazinājās 50.~gados, kad lielākā daļa muižu pārgāja no klaušām uz činšu un auga algoto strādnieku skaits. Zemnieku noslāņošanās nodrošināja pastāvīgu darbaspēka pieplūdumu rūpniecībai. Beidzot kalnrūpniecībā un metalurģijā tika atļauts iesaistīties arī privātkapitālam.

1845.~gadā karalistē bija 46~tūkstoši strādnieku. No Varšavas un Lodzas fabrikām tikai dažās strādnieku skaits sasniedza vairākus simtus. Pārējās strādnieku bija krietni vien mazāk. Rūpniecība saražoja produkciju 9,8~miljonu rubļu vērtībā. Turpretī amatniecība, ar kuru 1842.--1843.~gadā nodarbojās 65,5~tūkstoši cilvēku, deva produkciju 6,2~miljonu rubļu vērtībā. Turpmākajos gados notika lielas izmaiņas. 1860.~gadā strādnieku skaits rūpniecībā pieauga līdz 75~tūkstošiem, viņu saražotās produkcijas vērtība līdz 32~miljoniem rubļu. 1864.~gadā 78~tūkstoši strādnieku ražoja produkciju 50~miljonu rubļu vērtībā. Tātad 20.~gados rūpnieciskās produkcijas vērtība bija pieaugusi 20~reizes. Ar amatniecību 1860.~gadā nodarbojās 91~tūkstotis cilvēku, kas ražoja produkciju par vairāk nekā 18~miljoniem rubļu.

Tālāk attīstījās tekstilrūpniecība. Vācu vēsturnieki gan uzsver, ka vadmalas aušana pēc uzplaukuma 30.~gados pēkšņi piedzīvoja strauju kritumu. Cēloņi tam bija meklējami daudzu aušanas uzņēmumu izvietojumā Viduspolijā, kas radīja lielu konkurenci; bez tam tika ievests daudz ārzemēs ražota kokvilnas auduma. Mehanizācija kokvilnas audumu ražošanā noveda īpaši pie vilnas audumu sīkražošanas panīkuma. Visvairāk tomēr traucējusi 1831.~gadā ierīkotā muitas robeža starp Poliju un Krieviju. Padomju vēsturnieks A.~Manusevičs gan norādīja, ka vienkāršākais muitas robežas pārvarēšanas līdzeklis bija uz Krievijas tirgu strādājošās tekstilražošanas pārnešana uz kaimiņos esošo Belostokas (\pltxti{Białystok}) apgabalu, kurš pēc trešās Žečpospolitas dalīšanas atradās Prūsijas, bet pēc Tilzītes miera (1807)~--- Krievijas impērijas sastāvā, un kurp bez grūtībām varēja piegādāt izejvielas no Polijas. Otrs muitas traucēkļu pārvarēšanas līdzeklis bija mašīnu ieviešana un darba intensitātes palielināšana, tāpēc 1831.~gada muitas tarifs gala rezultātā tikai veicināja rūpniecības attīstību Polijas karalistē. 1847.--1849.~gadā tika izstrādāts un 1850.~gadā apstiprināts vienots muitas tarifs Krievijai un Polijas karalistei, iekšējā muitas robeža ar to 1851.~gadā tika likvidēta (tāpēc arī vācu vēsturnieki XIX gadsimta vidu uzsver kā robežšķirtni). Polijas karalistes loma Krievijas tirdzniecībā īpaši pieauga Krimas kara laikā, kad Krievijas tirgotāji, iepirkuši preces Rietumos, Krievijas jūras blokādes dēļ bija spiesti tās ievest pa sauszemes ceļu caur Poliju.

50.--60.~gados Polijā, īpaši tekstilrūpniecībā, norisa rūpniecības apvērsums. Kapitālistiskā kokvilnas auduma ražošanas manufaktūra ap gadsimta vidu pārvērtās par mehanizētu fabriku. Mehanizēts tika viss tehnoloģiskais process: sākumā vērpšana (1840), pēc tam arī aušana (1864). 1860.~gadā tekstilrūpniecībā darbojās 360~uzņēmumu ar 36~tūkstošiem strādnieku. Uzplaukumu veicināja kā iekšējā tirgus augšana, tā arī poļu preču ieveduma pieaugums Krievijā, īpaši pēc muitas robežas likvidēšanas starp Krieviju un Polijas karalisti. Stāvokli vilnas audumu ražošanā tā gan neglāba, toties veicināja ārzemju investīciju (galvenokārt vācu un ebreju kapitāla) pieplūdumu Lodzā. Auga arī pati pilsēta. XIX un XX gs. mijā tajā mita jau 400~000 iedzīvotāju. Blakus kokvilnas audumu ražošanai Lodzā ražoja arī zīdu, attīstījās arī lenšu, trikotāžas izstrādājumu ražošana.

Svarīga nozare Polijas karalistē bija kalnrūpniecība. To savās rokās turēja valsts, privātnoma un ārzemju kapitāla izmantošana kalnrūpniecībā bija aizliegti. Plaši sāka pielietot tvaika mašīnas, mehanizēt atsevišķus izejvielu ieguves un ražošanas procesus. Jau XIX gadsimta 20.~gados radās pirmās mašīnbūves fabrikas, kurās ražoja iekārtas tekstilrūpniecības, kalnrūpniecības, metalurģiskajām un metalapstrādes rūpnīcām. Progress bija tik ievērojams, ka jau 20.~gadu beigās poļu mašīnbūvniecības uzņēmumi varēja izturēt konkurenci arī ārējos tirgos. 1833.~gadā kalnrūpniecība un metalurģija pārgāja Polijas bankas, bet 1843.~gadā~--- valsts kases pārziņā. Pastāvēja uzņēmumi gan ar vecu tehnoloģiju, kuri lietoja kokogles un ūdens enerģiju, gan celti jauni (Dombrovas baseinā), kuri kā kurināmo izmantoja akmeņogles.

% TODO: Нужен заголовок таблицы!

\noindent
\begin{table}[h!]
\caption{Polijas karalistē saražotās rūpnieciskās produkcijas aptuvenā kopīgā vērtība (miljonos rubļu)} \label{tab:table2}
\begin{tabularx}{\linewidth}{|p{4cm}|p{4cm}|}
\hline
1845.~g. & 9,8 \\
\hline
1857.~g. & 21,3 \\
\hline
1864.~g. & 50,0 \\
\hline
\end{tabularx}
\end{table}

% page 136


Kaut tautsaimniecība attīstījās, tomēr karaliste vēl joprojām palika agrāra zeme ar muižniecisko saimniecību pārsvaru. No 40.~gadu beigām līdz 1863.~gadam Polijas karaliste bija vienīgā no poļu zemēm, kur vēl saglabājās zemniecības feodālā atkarība. Zemnieku apstrādājamā zeme skaitījās muižnieku īpašums un atradās tikai zemnieku lietošanā. Galvenie ieguvēji bija muižnieki, kuri centās atsavināt zemnieku saimniecību zemi, nemazinot to pienākumus. Ja XIX gadsimta otrajā gadu desmitā zemnieku zeme vēl sastādīja 58\% no visas lauksaimniecībā izmantojamās zemes, tad 50.~gados tās bija palicis vairs tikai ap 32\%. (Pēc citiem datiem, 50.~gadu beigās 60\% lauksaimniecībā izmantojamo zemju bija muižu un t.s. folverku rīcībā.) Jāpiebilst, ka muižu rīcībā bija ne tikai lielākā, bet arī labākā, auglīgākā, izdevīgāk izvietotā zemes daļa, jau agrāk vai šai laikā grupēta lielos masīvos. Šais saimniecībās bija vērojams arī lielāks ražības pieaugums. Ja 1840.~gadā salīdzinājumā ar iesēto sēklu graudaugu raža bija 3,18~reizes lielāka, tad 1861.~gadā~--- jau 4,9~reizes. Tā kā pieauga arī sējumu platība, graudu kopprodukcija minētajos gados pieauga par 177\%. Palielinājās arī citu lauksaimniecības nozare produktivitāte. Ja 1822.~gadā Polijas karalistē bija 1,5~miljoni aitu, tad 1861.~gadā~--- jau 3,7~miljoni. 1839.~gadā bija 1,17~miljonu liellopu, bet 1861.~gadā~--- 2,07~miljoni. 1839.~gadā karalistē bija 442~tūkstoši zirgu, bet 1861.~gadā~--- 605~tūkstoši.

XIX gadsimta pirmajā pusē Polijas karalistes lauksaimniecībā arvien vairāk ekstensīvo saimniecības metožu vietā tika pielietotas intensīvās. Tas pirmkārt nozīmēja pāreju no trīslauku sistēmas uz daudzlauku, ievērojot racionālu augu seku. Pāreja uz daudzlauku sistēmu nebija iespējama bez racionālas lopkopības. Tā prasīja lopbarības bāzes paplašināšanu, kas nebija iedomājama bez pārejas uz zaļbarību, kuras sagāde savukārt prasīja pāreju uz daudzlauku sistēmu. Tas viss kopā prasīja kapitālieguldījumu palielināšanu: dzīvā un lietiskā inventāra iegādi, kura līdz tam muižām, apstrādājamām ar zemnieku arkliem un vēršiem, nebija. Muižniekiem piederoša dzīvā un lietiskā inventāra aprūpei bija vajadzīgs pastāvīgs personāls, kvalificētāks un atbildīgāks nekā zemnieki-klaušinieki. Tehnisko kultūru ieviešana augu sekā radīja lielu darbaspēka pieprasījumu. Šeit strādnieki varēja būt nekvalificēti un fiziski mazspējīgi, taču viņu bija vajadzīgs daudz un viņu darbam bija jābūt lētam. Lopkopības attīstība un tehnisko kultūru audzēšana pavēra agrāk nebijušas iespējas folverkos izmantot sieviešu un bērnu darbu. Tehnisko kultūru un lauksaimniecības produktu pārstrādes rūpnīcas (piemēram, tikai sezonas laikā strādājošās cukurfabrikas) izmainīja sezonas darbaspēka pieprasījumu.

Pakāpeniski muižas vērtību noteica ne vairs tās platības lielums, bet tajā investēto kapitālu daudzums, kurš, saprotams, arī nodrošināja tās ražību. Šļahtiči-sīkmuižnieki centās izvairīties no ``liekiem'' kapitālieguldījumiem, toties lielmuižnieki jau meklēja jaunas investīciju iespējas. Ne katram zemes īpašniekam pāreja uz kapitālistiskām ražošanas metodēm bija pa spēkam. Ekonomiski spēcīgākajiem tā bija vieglāka. Vājākie folverki, nespēdami izturēt konkurenci, izmisīgi cīnījās par pastāvēšanu, izmantojot tradicionālās feodālās metodes. Turpretī lielākās, spēcīgākās muižnieku saimniecības enerģiski cīnījās par ārējiem tirgiem. Muižnieku-kapitālistu akciju sabiedrības centās panākt produktu izvešanu vismaz pēc pirmās to pārstrādes (izvest miltus, nevis graudus), tātad~--- par attiecīgi augstāku cenu. Tiesa, vēl joprojām visai izdevīga bija neapstrādātu baļķu pārdošana. Uz Angliju, Vāciju tika vesti poļu muižnieku mežos augušie koki, kuru nociršana un pludināšana izmaksāja lēti, bet kuru cenas ārzemju tirgos sakarā ar strauji augošo kuģu būvniecību un celtniecību cēlās galvu reibinošos tempos. Tāds stāvoklis XIX gadsimtā saglabājās vēl ilgi, nesot muižniekiem ienākumus: vieniem greznībai, citiem~--- ražošanas attīstībai.

Auga pieprasījums pēc lauksaimniecības mašīnām. 1833.~gadā nodibinājās pirmie nelielie lauksaimniecības mašīnu ražošanas uzņēmumi. 1840.~gados to skaits pieauga. 30.~gados sāka audzēt cukurbietes un radās pirmās cukurfabrikas. 1849.--1850.~gadā pastāvēja jau 33 tādi uzņēmumi. 50.~gadu sākumā eksistēja arī ap 3~000 spirta brūžu. Zemes apstrādes intensifikācija noteica ražīguma pieaugumu. No XIX gadsimta 30. līdz 50.~gadiem tas palielinājās apmēram pusotras reizes. Pieauga arī lauksaimniecības produkcijas ražošana uz vienu iedzīvotāju.

Pateicoties kapitālistisko saimniecības metožu pielietošanai, ievērojami pieauga muižnieku saimniecībās ražotā preču produkcijas daļa. Graudus un vilnu tās audzēja pārdošanai, kartrupeļus~--- šņabja, bietes~--- cukura ražošanai. Taču, no vienas puses, folverki ražoja masu produkciju tirgum, bet no otras~--- tās ražošanu nodrošināja ar klaušu darbu. Ļoti lēni norisa klaušu nomaiņa ar algotu darbu. Pēc 1846.~gada datiem renti jeb činšu maksāja 1/3~visu zemnieku saimniecību, pārējās vēl joprojām pildīja klaušas. No 3,3~miljoniem Polijas karalistes zemnieku 1859.~gadā ap 1,34~miljonus sastādīja bezzemnieki, kas veidoja 40,5\% no visiem zemniekiem.

Zemniecības pamatmasu stāvokļa pasliktināšanās rezultātā radās stagnācija iedzīvotāju pieaugumā. 1846.~gadā Polijas karalistē bija 4~867~tūkstoši iedzīvotāju, bet 1859.~gadā to skaits saruka līdz 4~764~tūkstošiem. Ebreju izcelsmes poļu ekonomists un vēsturnieks S.~Kempners norādīja: ``Tas bija iedzīvotāju zemnieciskās daļas, kura gandrīz vai mira badā, pauperizācijas rezultāts.''

Taču ekonomiskie panākumi nespēja nolīdzināt politiskās un sociālās pretrunas. Baidoties, ka zemnieku sacelšanās no Galīcijas var pārsviesties arī uz Polijas karalisti, 1846.~gadā Nikolajs I izdeva dekrētu, ar kuru aizliedza muižniekiem zemniekus izdzīt no viņu lietojumā esošajām zemēm vai samazināt to apjomu, solīja uzraudzīt činša līgumu slēgšanu. Poļu vēsturnieks V.~Grebeņskis komentēja tālāko: ``Iekšlietu komisiju, realizējot šo dekrētu, vadīja zemnieku naida pret šļahtu pamodināšanas tendence un tā sēja zemnieku iedzīvotājos sabiedriskās cīņas sēklas. Pārņemot savās rokās tiesības izšķirt strīdus starp muižniekiem un zemniekiem, komisija apzināti radīja dažādus sarežģījumus, kuri paplašināja bezdibeni starp zemnieka būdu un muižas centru.'' Šeit nu jāatzīst, ka cariskā administrācija nekādi nevarēja vadīties no šķiru cīņas saasināšanas motīva. Cita lieta, ka izpildot dekrētu, tika radītas neērtības šļahtičiem zemnieku ekspluatācijā. Svarīgākais vēsturnieka~--- šļahtas interešu izteicēja rakstītajā ir atzinums par ``bezdibeni starp zemnieka būdu un muižas centru'', un tas bija radies jau sen pirms Nikolaja I dekrēta.

1857.~gadā cars atļāva radīt Zemkopības biedrību (\pltxti{Towarzystwo Rolnicze}), ko vadīja grāfs A.~Zamoiskis. Biedrība atspoguļoja mantīgās šļahtas intereses, kura vēlējās mīkstināt politisko režīmu Polijas karalistē, atjaunot tās autonomiju, pievienot tai lietuviešu, ukraiņu un baltkrievu apdzīvotās zemes, zemniecības jautājumā centās panākt klaušu atcelšanu un nomaiņu ar činšu. Šādas prasības atbalstīja arī pilsētnieki, jo to izpilde veicinātu tautsaimniecības attīstību. Grāfs A.~Zamoiskis vispār uzskatīja, ka Polijai ir izdevīgāk būt vienā valstī ar Krieviju, jo Polijai atbrīvojoties, Krievija atkal mēģinātu to pakļaut, kas liktu Polijai tērēt milzīgus līdzekļus aizsardzībai. Dažu gadu laikā biedrība izauga, tajā iestājās vairāki tūkstoši cilvēku, strādāja pastāvīgas komisijas. Biedrība ieguva ne tikai ekonomisku, bet arī lielu politisku nozīmi, daļēji kļuva par parlamenta aizvietotāju. Poļu cerību uzplaukumu veicināja arī sabiedriskās aktivitātes Krievijā, jaunu reformu gatavošana. Poļu vidū radās kustība, vērsta uz sadarbību ar Krieviju, īpaši sociālekonomiskos jautājumos.

Taču 50.~gadu beigās Krievijā izveidojās revolucionāra situācija. Poļu sabiedrībā uzskatīja, ka briestošie satricinājumi Krievijā rada labvēlīgu augsni poļu mērķiem. Itālijā sākās karš par zemes atbrīvošanu no austriešu kundzības. Itāļiem palīdzēja Francija, kura bija ieinteresēta Austrijas novājināšanā. Francijas imperators Napoleons III pasludināja t.s. ``nacionālo principu''~--- nacionālo neatkarību. Poļu vidū Napoleons III ieguva milzu popularitāti.

Arī Krievijas impērijas sastāvā ieejošajās poļu zemēs 1859.~gadā pastiprinājās pret carismu vērsti noskaņojumi: tika nēsāti tradicionālie poļu apģērbi, lietotas dažādas rotas lietas ar nacionālajām emblēmām, demonstratīvi svinētas vēsturiskas gadadienas tādiem notikumiem, kā, piemēram, Ļubļinas ūnija (1569), 3.~Maija Konstitūcijas (1791) pieņemšana, Novembra (1830) sacelšanās. Tika dziedātas patriotiskas dziesmas, sarīkoti ``kaķu koncerti'' zem carismam uzticīgo poļu sabiedrības pārstāvju logiem.

1860.~gadā Polijā sākās pirmās masu manifestācijas. Oktobrī triju Poliju sadalījušo valstu monarhu tikšanās laikā Varšavā patrioti izvērsa aģitāciju par tikšanās boikotēšanu. Aleksandram II iebraucot pilsētā, tās ielas kļuva gluži tukšas. 20.~oktobrī operas teātrī pirms izrādes sākuma cara ložu aplēja ar sērskābi, no galerijas tika nomesti flakoni ar smakojošu šķidrumu, izplatījās tāds smārds, ka sapulcējusies publika bija spiesta atstāt skatītāju zāli.

1861.~gada sākumā Polija gaidīja divus svarīgus notikumus. 21.~februārī bija jāsākas ``Zemkopības biedrības'' kārtējam kongresam, kuram vajadzēja pieņemt svarīgus lēmumus zemnieku jautājumā, un 25.~februārī apritēja 30.~gadadiena kopš kaujas pie Grohovas (\pltxti{Bitwa pod Olszynką Grochowską}) 1831.~gada poļu sacelšanās laikā. (Tajā pēc poļu datiem krita 7~000 poļu un 10~000 krievu, bet pēc krievu datiem 12~000 poļu un 9~400 krievu.) Viens no poļu emigrācijas vadoņiem L.~Mieroslavskis, kurš 60.~gadu sākumā mēreno poļu demokrātu, Ā.~Čartorijska pretinieku vidū bija ieguvis lielu ietekmi, no Parīzes pat atsūtīja direktīvu Grohovas kaujas gadadienā sākt sacelšanos. Tiesa, izrādījās, ka viņa aprēķini par iespējamo Francijas palīdzību ir nereāli un viņš pats atteicās no sava plāna.

1861.~gada 24.~februārī ``Zemkopības biedrības'' pieņemtajā rezolūcijā tika izteikta tikai vēlme visur nomainīt klaušas ar činšu (atgādinājumam: poļu \pltxti{czynsz}, vācu \detxti{Zins}~--- procents, no latīņu. \latxti{Census}, Žežpospolitā tā sauca nodevas graudā jeb naudā, kuras zemnieki maksāja par zemes lietošanu muižniekam), bet pēc tam gadu desmitu laikā ļaut zemniekiem izpirkties brīvībā un kļūt par zemes īpašniekiem. (Līdz 1859.~gadam uz činšu bija pārgājušas 41\% zemnieku saimniecību, taču daudzās no tām zemnieki nodevas maksāja, reizē pildot arī klaušas. Vairākumā gadījumu pāreju uz činšu pavadīja daļas zemnieku saimniecību zemes ``nogriešana'', apmaiņā pret sliktākas kvalitātes zemi.)

25.~februārī poļu revolucionāru pulciņi organizēja vairākus tūkstošus dalībnieku lielu demonstrāciju sakarā ar kaujas pie Grohovas gadadienu, kurai vajadzēja ne tikai demonstrēt sabiedrības noskaņojumu, bet arī ietekmēt ``Zemkopības biedrību'' radikālākas darbības virzienā. Speciāli izsūtīts žandarmu puseskadrons demonstrāciju izklīdināja nagaikām un zobeniem, arestēja daudzus tās dalībniekus. Kārtība tika atjaunota, taču ne uz ilgu laiku. Pēc divām dienām 27.~februārī notika nākamā demonstrācija. Demonstranti, sastapuši ceļā karaspēka aizsprostu, to apmētāja akmeņiem. Karaspēks atbildēja ar zalvi, rezultātā pieci cilvēki tika nogalināti, bija arī ievainotie. Demonstranti nogalināto ķermeņus vadāja pa pilsētu un rādīja iedzīvotājiem. No poļu buržuāzijas un inteliģences pārstāvjiem sastādīta komiteja ieradās pie Krievijas vietvalža kņaza M.~Gorčakova un lūdza izmeklēt incidentu, ļaut organizēt upuru apbedīšanu. Lūgumi tika apmierināti. Apbedīšana norisa 2.~martā un izvērtās ap 100~000 cilvēku demonstrācijā. Kristīgajā ceremonijā piedalījās arī ebreju garīdzniecība savos nacionālajos tērpos. Uz ielām nebija nedz kareivju, nedz policijas. ``Zemkopības biedrības'' vadība nosūtīja caram vēstījumu, saskaņotu ar buržuāziski-inteliģentisko komiteju, kurā paziņoja, ka zeme nevarēs attīstīties, kamēr poļu tautas nacionālie principi netiks īstenoti visās sabiedriskās dzīves jomās. Imperators Aleksandrs II vēstījumu pieņēma, savā vēstulē vietvaldim apsolot īstenot dažas reformas.

Stāvokli Polijā īpaši saasināja dzimtbūšanas atcelšana Krievijā (parakstīta 1861.~gada 19.~februārī, izsludināta 5.~martā). Agrārā reforma Krievijā neattiecās uz Polijas karalisti, bet poļu zemnieki cerēja, ka arī viņi varēs saņemt zemi. Gandrīz vairākas nedēļas pēc februāra demonstrācijām Varšavā valdīja diezgan liela sapulču un vārda brīvība. Demokrātiskie pulciņi gaidīja norādījumus no L.~Mieroslavska, taču, kā rāda poļu vēsturnieks S.~Keņēvičs, viņš šai laikā tikai ieteica paplašināt kustību uz austrumiem~--- uz ukraiņu, baltkrievu un lietuviešu zemēm, prasot lai vietējie šļahtiči piešķirtu zemniekiem zemi uz labvēlīgākiem nosacījumiem, nekā to solīja valdība, bet pašā Polijā nesteigties, iet kopā ar šļahtu un pretoties sociālās revolūcijas tendencēm.

Sabiedrības neapmierinātības pieaugums piespieda cara valdību uzsākt \strong{reformas} arī \strong{Polijā}. 1861.~gada 27.~martā tika publicēti rīkojumi par Polijas karalistes Valsts padomes atjaunošanu, tika pavērtas plašākas iespējas poļu valodas un kultūras apguvei skolās. Tika radīta ticības un izglītības lietu komisija ar lielmuižnieku marķīzu A.~Veļepoļski priekšgalā. Viņa mērķis bija Polijas pārveidošana sadarbībā ar Krieviju un pakāpeniska Polijas autonomijas atjaunošana Organiskā statuta robežās. A.~Veļepoļskis bija pārliecināts, ka ar vardarbību poļi nespēs atbrīvoties no Krievijas kundzības un, ka vispār nevis Krievija un krievi, bet gan vācieši ir galvenie poļu ienaidnieki. Viņš nepretendēja uz vecās Žečpospolitas bijušajiem austrumu apgabaliem, uzskatot, ka tas būtu pilnīgi nereāli. Šais apgabalos tobrīd dzīvoja ap 10 miljoniem iedzīvotāju, no kuriem tikai nedaudz vairāk par miljonu bija poļi. A.~Veļepoļska reformu programma paredzēja pārvaldes repolonizāciju, izglītības reformu, agrārā un ebreju jautājumu atrisināšanu. Vismaz divos pirmajos no nosauktajiem virzieniem iezīmējās panākumi. Satiksmes un pasta jautājumos karaliste kļuva neatkarīga no Krievijas. 1861.~gadā tika nodibinātas ap 300~poļu skolas. Tomēr poļu sabiedrībā A.~Veļepoļskim bija maz piekritēju. Pat tie poļu aristokrātijas pārstāvji, kuri bija apmierināti ar cariskās Krievijas sociālo politiku un saistījuši savas intereses ar Krieviju, centās panākt visu agrāk Polijai piederējušo ukraiņu, baltkrievu un lietuviešu apdzīvoto apgabalu apvienošanu ar Polijas karalisti vienā autonomā valstī ar īpašu Konstitūciju. Arī šo apgabalu poļu iedzīvotāji juta līdzi tautiešiem rietumos un pauda vēlmi apvienot Lietuvu, Baltkrieviju un Labā krasta Ukrainu ar Polijas karalisti.

Uzsāktie pārkārtojumi radīja nepamatotas cerības poļu sabiedrībā. Īpaši jaunatne, pārvērtējot savus spēkus, cerēja, ka varēs uzspiest Krievijai arvien jaunus piekāpšanās soļus. Daļa poļu virsslāņu pārstāvju, kuri pēc saviem sociāli-ekonomiskajiem uzskatiem ļoti maz atšķīrās no krievu liberāļiem un pat konservatoriem, atšķirībā no krievu virsslāņiem, bija gatavi uz bruņotu cīņu pret patvaldību, lai radītu muižnieciski-buržuāzisku Poliju. Reizē jāuzsver, ka viņi nebija ar nepārvaramu sienu atdalīti arī no pilsoniskās nacionālās kustības, no plebejiskajiem revolucionāriem. Atsevišķi publicisti ir izteikuši viedokli, ka labvēlīgas notikumu attīstības gadījumā nacionālā kustība varēja savīties ar sociālo.

1861.~gada 6.~aprīlī varas iestādes atlaida ``Zemkopības biedrību'', tā vājinot to poļu slāņu pozīcijas, kuri gribēja izvairīties no sacelšanās. Sākās poļu pilsoņu manifestācijas. 8.~aprīlī atkal radikāļu organizēto demonstrāciju Varšavā ar šāvieniem izklīdināja karaspēks, nogalinot vairāk nekā 100 cilvēku. Atskanēja kara trauksmes signāls un seši lielgabala šāvieni pavēstīja varšaviešiem, ka pilsētā tiek ievests kara stāvoklis. Karaspēks ieņēma svarīgākās vietas. Taču nekāda sacelšanās vēl nesākās, jo nebija kam nostāties patriotiskās kustības priekšgalā. Virkne slepenu organizāciju tika sagrautas.

Laukos tai pat laikā zemnieki ``Zemkopības biedrības'' lēmumu ar ieteikumu nomainīt klaušas ar čiņšu uztvēra kā reālu varas iestāžu rīkojumu. Kad laukos nonāca ziņas par reformu Krievijā, daudzu ciemu zemnieki nolēma atteikties no klaušu pildīšanas. Aprīļa vidū kustība aptvēra 176~ciemus, bet mēneša beigās jau 377~muižas, 954~ciemus. Pavisam līdz maija vidum pretošanās kustībā muižniekiem iesaistījās vairāk nekā 22~tūkstošu zemnieku sētu ar vairāk nekā 180~tūkstošiem cilvēku, kas sastādīja 15 līdz 20\% visu klaušinieku.

1861.~gada 2.~oktobrī Krievijas vietvaldis Polijas karalistē ģenerālis K.~Lamberts, kurš nomainīja maijā mirušo M.~Gorčakovu, centās rīkoties bez vardarbības, tomēr, atbildot uz virkni nemieru Polijas pilsētās, bija spiests izsludināt visā karalistē kara stāvokli. Pēc dienas~--- 3.~oktobrī, lielas ļaužu masas Varšavā piepildīja baznīcas, tur notika demonstratīvs aizlūgums par nacionālo varoni T.~Kostjuško, tika dziedātas patriotiskas himnas. Ar nolūku arestēt visus no baznīcām iznākušos vīriešus, tās aplenca un bloķēja krievu karaspēks. Taču pēc dievkalpojumu beigām neviens baznīcas neatstāja. Pagāja diena, iestājās nakts. Tad tika dots rīkojums karaspēkam ielauzties baznīcās. Vienā no tām karavīri neatrada nevienu cilvēku, visi bija paguvuši pazust pa sānejām. Pārējās baznīcās pēc dažādām ziņām tika arestēti no 1~500 līdz 3~000 cilvēku, kuri izrādīja pretošanos. Tiesa, pēc K.~Lamberta pavēles lielāko daļu arestēto drīz atbrīvoja. (Tāda nekonsekvence izsauca Varšavas ģenerālgubernatora A.~Hercencveiga protestu. Notika asa izskaidrošanās un, lai izbēgtu sodam par dalību klasiskā divkaujā, tika pieņemts tās ``amerikāņu'' variants~--- izlozēts, kam dzīve jābeidz pašnāvībā. Nepaveicās ģenerālgubernatoram, viņš nošāvās.) Protestējot pret iebrukumu baznīcās, katoļu garīdzniecība uz laiku slēdza visus Varšavas dievnamus. Tie nedarbojās līdz pat 1862.~gada sākumam. K.~Lamberts tika atsaukts no vietvalža amata. 2.~novembrī no sava amata atteicās un uz Pēterburgu devās A.~Veļepoļskis. Tur ilgstošās pārrunās viņam izdevās pārliecināt Aleksandru II turpināt ``kontrolētas reformas''.

Tomēr ziemā Varšavā uz ielām tika izvietotas karaspēka daļas, notika pilsētnieku atbruņošana. Konfiscēja vairāk nekā septiņus tūkstošus medību bišu, daudz citu ieroču. Taču pagrīdē turpinājās poļu patriotu organizēšanās, veidojās dažādas komitejas.

Svarīgi, ka manifestācijās gandrīz nepiedalījās viens no poļu sabiedrības pamatslāņiem~--- zemniecība. Poļu konservatīvais vēsturnieks un notikumu aculiecinieks V.~Pšiborovskis par lauku baznīcās rīkotajām patriotiskajām manifestācijām rakstīja: ``Ļoti bieži gadījās, ka zemnieki, padzirduši aizliegtu dziesmu, devās laukā no baznīcas un bariem metās bēgt, it kā viņus biedētu tās šļahtiskās Polijas rēgs, par kuru ``kungi'' lūdza dievu.'' Minētais vēsturnieks atstāstīja faktu par Sandomiras vojevodistē 1862.~gada septembra vidū organizētu lielu patriotisku pasākumu ar dievlūgšanu. Kad ksendzs pēc uzbudinoša sprediķa pieprasīja lai sapulcējušies paceļ roku kā zīmi zvērestam, ka viņi cīnīsies par tēvzemi līdz pēdējai asins lāsei, iestājās klusums, kuru pārtrauca sauciens: ``Lai neviens neceļ roku!''~--- un zemnieki pūlī devās projām no baznīcas.

Arī Polijas karalistes sabiedrībā pēc emigrācijas parauga 1861.~gada beigās noformējās divas galvenās politiskās nometnes~--- \strong{``baltie''} (konservatīvie, arī daļa liberāļu) un \strong{``sarkanie''} (radikālie demokrāti).

``Baltie'' daļēji nāca no ``Zemkopības biedrības'', izteica zemes īpašnieku~--- šļahtas un lielpilsonības intereses, atbalstīja pasīvas opozīcijas taktiku. Sociālajā sfērā ``baltie'' iestājās par feodālo attiecību likvidēšanas pēc Prūsijas parauga~--- pēc iespējas mazāk zemi atstājot zemniekiem. A.~Veļepoļskim neizdevās piesaistīt ``baltos'' savam kursam, jo pēc to domām viņš pārāk maz uzmanības veltīja agrāk Žečpospolitai piederošo apgabalu pievienošanai Polijas karalistei. ``Baltie'' 1862.~gada decembrī izveidoja savu ``direktorātu'', kurā galvenokārt darbojās bijušās ``Zemkopības biedrības'' locekļi.

``Sarkanie'' galvenokārt veidojās no sīkburžuāziskajiem~--- ierēdņu, pilsētnieku, studentu u.c.~--- slāņiem, iestājās par bruņotas sacelšanās sagatavošanu, lai tādā ceļā atjaunotu Polijas neatkarību. Sarkano paraugs bija Itālijas apvienošanās Dž.~Garibaldi un K.~Kavūra vadībā. ``Sarkano'' lozungs skanēja: ``\ittxti{Polonija farà da se}'' (itāļu ``Polijai pašai sevi jāatbrīvo''). ``Sarkano'' labējais spārns, tai skaitā L.~Mieroslavskis, pirmajā vietā stādot Polijas atjaunošanu 1772.~gada robežās, pietiekami nenovērtēja sociālās problēmas. Kreisais, revolucionāri-demokrātiskais ``sarkano'' spārns: Z.~Serakovskis, J.~Dombrovskis, K.~Kaļinovskis, A.~Mackevičs atzina lietuviešu, baltkrievu, ukraiņu tiesības uz pašnoteikšanos, sociālajā sfērā iestājās par radikālākām reformām. Vienoja abas nometnes vēlme atjaunot Žečpospolitu agrākajās robežās.

Situācija bija nestabila, ``rūgšana'' tautā turpinājās. Baznīcu darbības atjaunošana 1862.~gada pavasarī neuzlaboja stāvokli. Varšavā faktiski turpināja eksistēt kara stāvoklis. Šādos apstākļos Aleksandrs II atkal mēģināja arī savos poļu valdījumos atgriezties pie piekāpšanās politikas. Aprīlī tika apžēloti daudzi politiskie ieslodzītie. Polijas karaliste ieguva diezgan plašu autonomiju ar savu pārvaldi. Maijā par vietvaldi nozīmēja ar saviem liberālajiem uzskatiem pazīstamo cara brāli lielkņazu Konstantīnu Nikolajeviču (nejaukt ar Aleksandra I brāli Konstantīnu Pavloviču). Tiesa, savā instrukcijā brālim Aleksandrs II rakstīja, ka Polijas karalistei ``vienmēr ir jāpaliek Krievijas īpašumā'' un viņš ``nekādā gadījumā'' nepieļaus lai tā iegūtu Konstitūciju un savu armiju. Augstākais, ko tai pakāpeniski varētu piešķirt~--- zināmā mērā atsevišķu pārvaldi.

Būdams dedzīgs Aleksandra II realizējamo reformu piekritējs, lielkņazs Konstantīns Polijas karalistē nesekmīgi centās realizēt samierniecisku politiku. Lielkņazs plaši izmantoja viņam piešķirtās apžēlošanas tiesības, līdz 1862.~gada septembrim no 499~notiesātajiem viņš piedeva 289. Par vietvalža palīgu civillietās kļuva A.~Veļepoļskis. Civillietu nodošanai poļa rokās vajadzēja simbolizēt imperatora labvēlību poļiem. A.~Veļepoļskis tūlīt publicēja jau iepriekš sagatavotus rīkojumus par reformām.

Pirmkārt, 1862.~gada 16.~maijā tika izdots rīkojums par klaušu nomaiņu no 1.~oktobra ar čiņšu. Nedz muižnieki, nedz zemnieki nevarēja no tā atteikties. Tas tika noteikts zemāks kā agrāk, taču zemnieki zaudēja servitūtus~--- tiesības izmantot agrāk kopīpašumā, bet tagad muižnieku īpašumā pārgājušās pļavas, mežus, ūdeņus. Rīkojums neskāra citus zemnieku pienākumus muižnieku un valsts priekšā, nabagākos zemniekus muižnieki tāpat kā iepriekš varēja padzīt no zemes. Tā rīkojums kopumā atbilda muižnieku interesēm, kuri pat pieaugošās zemnieku kustības apstākļos nevēlējās šķirties no savām feodālajām tiesībām, bet neapmierināja zemniekus. Antifeodāli nemieri aptvēra simtiem ciemu. Zemnieki vēlējās panākt, lai klaušu atcelšana notiktu bez atlīdzības un apstrādātā zeme nonāktu viņu īpašumā.

Otrkārt, jūnijā tika izdoti rīkojumi par Polijas karalistes Valsts padomes atjaunošanu un vietējās pašvaldības iestāžu~--- vēlētu pilsētu, apriņķu un guberņu padomju izveidi. Kaut radikāļi aicināja boikotēt šo padomju vēlēšanas, tās notika, pilsētās tika ievēlēti buržuāzijas, apriņķos~--- muižnieku pārstāvji. Manifestāciju kustība no Polijas karalistes pārsviedās arī uz Lietuvu, Baltkrieviju, Labā krasta Ukrainu. Lietuvā jau 1861.~gada augustā varas iestādes bija spiestas izsludināt kara stāvokli.

Treškārt, tika atcelti ierobežojumi ebrejiem dzīves vietas izvēlē, privātīpašuma iegādē, sabiedrisko pienākumu izpildē. Tika atcelti īpašie nodokļi ebrejiem. Saglabājās gan arī daži ierobežojumi~--- iegādāties zemi, kur tika pildītas klaušas, ieņemt ciemu vecāko (\pltxti{wójt}) amatus. Atceļot daudzos ierobežojumus ebrejiem, rīkojums bija solis uz priekšu nacionālā jautājuma risināšanā. Tiesa, netika piešķirtas kādas politiskās brīvības.

Ceturtkārt, īpaši svarīgs bija rīkojums, veltīts izglītības sfērai. Pilsētās, miestos un lielos ciemos tika dibinātas valsts sākumskolas ar bezmaksas apmācību poļu valodā. Tika atļauts dibināt jaunas poļu ģimnāzijas. Mācību maksa tika pazemināta. Uz Kara medicīniskās akadēmijas bāzes kā \pltxti{Szkoła Główna Warszawska} (Varšavas Galvenā skola) faktiski tika atjaunota 1832.~gadā slēgtā Varšavas universitāte. Galvenajā skolā darbojās 4 nodaļas: vēstures un filoloģijas, fizikas un matemātikas, juridiskā un medicīniskā. Salīdzinājumā ar iepriekšējo izglītības sistēmu, tas bija ievērojams progress.

Bez jau norādītā, tika veikti vēl daži soļi režīma daļējai mīkstināšanai. Tā, tika likvidēts ģenerālgubernatora amats, atlaisti daži sabiedrības visneieredzētākie ierēdņi.

Taču poļu neatkarības centieni nenorima. Jebkuras Krievijas piekāpšanās poļu patstāvības centieniem nāca par vēlu un šķita tiem pārāk niecīgas, tās izsauca poļos tikai vēl lielāku neapmierinātību ar pastāvošo stāvokli. Radikāļi, cenšoties traucēt reformu īstenošanu, kas varētu mazināt atbalstu revolucionāru atbalstu tautā, izvērsa teroru. Poļu revolucionāri darbojās, īstenojot lozungu ``Jo sliktāk [tautai]~--- jo labāk [revolūcijai]'', ar teroru centās izprovocēt bargus atbildes soļus, lai tos izvirzītu kā apvainojumu saviem politiskajiem pretiniekiem. Kā raksta vācu vēsturnieks E.~Meijers, ``sarkanie'' rīkoja masu pasākumus nolūkā izsaukt sadursmes ar Krievijas armijas daļām un provocēt asins izliešanu. Vēsturnieks gan nedod pierādījumus, ka tiešu tādi bija revolucionārpatriotu nodomi, bet notikumi attīstījās tieši tādā virzienā un saprātīgi cilvēki to varēja paredzēt. Diemžēl, radikāļi savu mērķu sasniegšanai ir gatavi upurēt ne tikai savas, bet arī daudzu citu cilvēku dzīvības.

Pēc vietvalža Konstantīna ierašanās Varšavā pret viņu notika atentāts, taču lielkņazs tika tikai viegli ievainots. Arī pret A.~Veļepoļski tika veikti divi neveiksmīgi atentāti. Radikālpatrioti tā centās nepieļaut viņu plāniem neatbilstošas reformas, kuras varētu nomierināt tautu. Visus notvertos teroristus 1862.~gada augustā pakāra Varšavas citadelē. Taču atentātu politiskais mērķis bija sasniegts, revolucionāru provokācijas izraisīja sabiedrībā sašutumu nevis pret slepkavības mēģinājumu realizētājiem, bet to sodītājiem. Izvērsās dažādu aizlūgumu, dievkalpojumu ar politisku ievirzi noturēšana, pieauga sacelšanos gatavojošo ietekme.

Gan Rietumu sabiedriskā doma, gan poļu, gan arī padomju vēsturnieki tolaik Polijā notiekošos procesus traktēja kā plašu demokrātisku kustību pret krievu cara tirāniju ar diviem nesaraujami saistītiem uzdevumiem: sociālo~--- zemnieku atbrīvošanu un nacionālo~--- visas tautas atbrīvošanu. Turpretī mūsdienu krievu vēsturnieks A.~Širokorads pievērsis uzmanību tam, ka 60.~gadu sākums bija Aleksandra II reformu izvēršanās posms Krievijā, norisa zemnieku brīvlaišana, tika gatavota zemstu ieviešana, tiesu u.c. reformas. Cita lieta, ka krievu revolucionāri prasīja vēl radikālākus pasākumus. Turpretī poļu nemiernieki pēc A.~Širokarada domām nestādīja mērķi veikt kādas demokrātiskas vai ekonomiskas reformas, bet, gluži pretēji, vairākums t.s. ``izglītoto poļu'' ar sašutumu uzņēma Aleksandra II 1861.~gada manifestu par zemnieku atbrīvošanu no dzimtbūšanas. Poļu sabiedrības galvenais mērķis palika Polijas atjaunošana ``no jūras līdz jūrai''. 1863.~gada sacelšanās viņa vērtējumā bija ``vienīgi panu un ksendzu inspirēta no augšas''. Abiem viedokļiem var atrast argumentus.

Var ilgi diskutēt par to, kā reformas, Aleksandra II un jaunā vietvalža lielkņaza Konstantīna piekāpšanās poļu prasību priekšā ietekmēja stāvokli, vai tā veicināja vai kavēja jaunas poļu sacelšanās sākumu. Piemēram, konservatīvais krievu publicists kņazs V~Meščerskis poļu sacelšanās izraisīšanā vainoja ne tikai revolucionāro demokrātu A.~Hercenu, bet arī liberālo izdevēju A.~Krajevski, Krievijas iekšlietu ministru P.~Valujevu, ietekmīgo galmā Sanktpēterburgas ģenerālgubernatoru kņazu A.~Suvorovu, Polijas karalistes vietvaldi lielkņazu Konstantīnu, kuri, koķetējot ar poļiem, parādījuši nepieļaujamu vājumu. Būtībā visa krievu konservatīvā prese lielkņaza Konstantīna politikā saskatīja pārlieku piekāpību poļu priekšā un nedaudz vēlāk~--- 1863.~gada rudenī viņš nolika vietvalža pilnvaras. Arī mūsdienu krievu vēstures literatūrā ir sastopams viedoklis, ka Krievijas varas iestādes pārāk bieži lietoja ``pīrāga'' metodi~--- periodiski notika politisko noziedznieku amnestijas un vispār šīm iestādēm trūka nepieciešamās stingrības, bet sazvērnieki to uztvēra kā vājuma pazīmi.

Atbildot uz Aleksandra II gatavības uz sadarbību apliecinājumiem, uz Konstantīna iniciatīvām, ap trīssimt sabiedrības pārstāvju, ``balto'' piekritēju, 1862.~gada jūlijā griezās ar lūgumu pie A.~Zamoiska paziņot lielkņazam poļu prasības: pieņemt liberālu, Polijas autonomiju garantējošu satversmi, radīt pašu poļu armiju un pievienot tai atdalītos apgabalus. ``Kā poļi mēs varam atbalstīt valdību tikai tad, kad tā kļūs par poļu valdību, un kad visi apgabali, kas sastāda mūsu dzimteni, tiks apvienoti un varēs lietot Konstitūciju un brīvas iestādes''~--- tā tika izteikti ``balto'' uzskati. Tikai uz šādiem noteikumiem viņi vēlējās sadarboties ar Krieviju. Citiem vārdiem, tika prasīta ne tikai demokrātijas ieviešana, bet arī citu tautību iedzīvotāju apdzīvoto apgabalu pievienošana Polijas karalistei. A.~Zamoiska iesniegumā bija arī vārdi: ``\pltxti{Žądać nie mame prawę, a prosić o coś nie chemy}'' (``Mums nav tiesību ko prasīt, bet lūgt mēs nevēlamies''), kurus varēja arī tulkot vismaz kā brīdinājumu, ja ne draudus. Sakarā ar sabiedrības radikālisma pieaugumu Aleksandrs II izsauca A.~Zamoiski uz Pēterburgu. Kad A.~Zamoiskis Aleksandram II apstiprināja, ka ``tikai Lietuvas un atņemto provinču pievienošana karalistei apmierinās poļus'', viņu ``palūdza'' uz laiku doties uz ārzemēm.

No 1862.~gada vasaras sacelšanās gatavošanu vadīja Centrālā Nacionālā Komiteja (\pltxti{Centralny Komitet Narodowy}), kuru ``sarkanie'' bija izveidojuši 1861.~gada oktobrī (gan zem cita nosaukuma, tie mainījās). Viens no tās locekļiem bija Krievijas ģenerālštāba poļu virsnieku pulciņa dalībnieks J.~Dombrovskis. 1862.~gada februārī viņš ieradās Varšavā, uzsāka nelegālu darbību, izvirzījās par vienu no nacionālās kustības vadītājiem.

Jūlijā Centrālā Nacionālā Komiteja publicēja ``instrukciju'', kura faktiski saturēja sacelšanās programmu. Par organizācijas mērķi tika paziņota Polijas neatkarības atjaunošana 1772.~gada robežās, kas dotu visiem pilsoņiem ``pilnīgu brīvību un vienlīdzību likuma priekšā'', bet tautām, dzīvojošām valstī, to ``tiesību cienīšana''. Tika sludināta ``brālība starp šķirām'' un zemes piešķiršana zemniekiem ar valdības atlīdzību muižniekiem.

Ar to 1863.~gada sacelšanās, no vienas puses, atšķirībā no sacelšanās 1830.--1831.~gadā nacionālās atbrīvošanās lietu saistīja ar pārmaiņām īpašuma un sociālajā sfērā. No otras puses, no šļahtiču un pilsonības aprindām nākušie sacelšanās dalībnieki tomēr nespēja piedāvāt pietiekami dziļas izmaiņas agrārajās attiecībās, likvidējot muižnieku zemes īpašumu, kas varētu piesaistīt viņiem plašas zemnieku masas.

Mūsdienu krievu dzejnieks un publicists, žurnāla ``\rutxti{Наш современник}'' (``Mūsu laikabiedrs'') galvenais redaktors S.~Kuņajevs, kurš no krievu nacionālista pozīcijām polemizē ar poļu publicistiem, poļu šļahtas neapmierinātības cēloni saskata tās lielajā īpatsvarā kopējā iedzīvotāju skaitā, kad Krievijā tikai katrs divsimtais tās Eiropas daļas iedzīvotājs bijis muižnieks, bet Polijā katrs desmitais, un Krievijas imperatora veiktajā zemnieku atbrīvošanā no dzimtbūšanas. Patiesībā Polijas karalistē dzimtbūšana bija atcelta jau XIX gs. sākumā. Runa varētu būt tikai par to, ka daļa poļu muižnieku zaudēja varu pār zemniekiem ārpus Polijas karalistes esošajās Krievijas rietumu guberņās, bet pašā Polijas karalistē muižnieki varēja būt neapmierināti tikai ar klaušu aizstāšanu ar renti. Tikai sacelšanās rezultātā Polijas karalistē Krievijas imperators izdeva dekrētu par zemes piešķiršanu zemniekiem. Šis piemērs rāda, kā strīdos tiek aizmirsta vēsturiskā precizitāte.

Interesantus spriedumus ir izteicis poļu ekonomikas vēsturnieks V.~Kula. Parasti par vienu no topošās buržuāzijas antifeodālo noskaņojumu cēloņiem uzskata tās neapmierinātību ar pastāvošo kārtību, kad feodāļi turēja zemniekus dzimtbūtnieciskā atkarībā un kā vienīgie izmantoja to darbaspēku. Analizējot situāciju, V.~Kula norādījis, ka XIX gadsimta pirmajā pusē Polijas karalistē darba roku trūkumu izjuta nevis jaunā buržuāzija, bet gan muižniecība. Tāpēc arī buržuāzijai nebija pamata būt neapmierinātai ar pastāvošo šauro, bet tās vajadzības tomēr apmierinošo darbaspēka tirgu. Abas šķiras~--- gan muižniecība, gan buržuāzija izmantoja tām garantētās darbaspēka ekspluatācijas sfēras. Tāpat buržuāzija jau varēja ieguldīt savus kapitālus zemes iegādē, saimniekojot uz tās ar kapitālistiskām metodēm. Tāpēc arī poļu buržuāzija nebija ieinteresēta izvirzīt konsekventus antifeodālus lozungus, iesaistīties cīņā pret feodālismu. Arī šļahtas noslāņošanās process vēl nebija beidzies. Daļa no tās, zaudējusi savu īpašumu, bija pārgājusi citu slāņu rindās, taču daļa vēl joprojām visiem spēkiem centās saimniekot savās muižas ar vecajām feodālajām metodēm, jo tai nebija izvēles. Savu līdzekļu saimniecības kapitālistiskai pārveidošanai šai šļahtas daļai nebija, kredītu saņemt sakarā ar lielajiem parādiem tā nevarēja, pārdot daļu īpašumu, lai iegūtu līdzekļus,~--- arī nē, jo šie īpašumu jau tā bija nelieli. Ceļš uz kapitālismu šai šļahtas daļai draudēja ar izputēšanu. Tāpēc tā, neatkarīgi no saviem politiskajiem uzskatiem, sociālajos jautājumos, īpaši agrārajā, bija viens no konservatīvākajiem spēkiem.

Padomju vēsturnieks V.~Zaicevs savukārt izdalījis divas zemniecības daļas: 1) uz valsts un 2) uz privātīpašnieku zemes saimniekojošas. Valsts īpašumā esošajās zemēs kā tās pārvaldītājs uzstājās vietējā administrācija. Tāpēc valsts zemnieku sociālie centieni (pirmkārt pēc zemes iegūšanas) sakrita ar nacionālo atbrīvošanās kustību, bija vērsti pret kopējo ienaidnieku~--- Krievijas impēriju. Citāda situācija bija privātajā īpašumā esošajās saimniecībās. Šeit zemnieku sociālās intereses sadūrās ar zemes īpašnieku interesēm. Pēdējie, jau attīstošos preču-naudas attiecību apstākļos centās iegūt no sava īpašuma maksimālu peļņu un savās muižās paaugstināja zemnieku ekspluatāciju. Tāpēc privātajos īpašumos pastāvēja stipras sociālās pretrunas, kuras, no vienas puses, kavēja zemnieku nacionālās pašapziņas rašanos, bet, no otras, stiprināja viņos monarhistiskas ilūzijas.

Kā pārliecinoši rāda poļu vēsturnieks S.~Keņevičs, lai piesaistītu nacionālajai kustībai zemniecību, bija jārealizē radikāla agrārā reforma. Muižniekiem zemnieku labā bija jāziedo daļa savas zemes, arī baznīcas institūcijām jāatsakās no saviem zemes valdījumiem. Tikai tādos apstākļos bija iespējama vienota nacionāla fronte. Zemes piešķiršana zemniekiem radītu tajos nacionālas jūtas un veidotu sacelšanās kustībai plašu masu pamatu. Taču kā rādīja gan iepriekšējo, gan šīs sacelšanās pieredze, muižnieki nebija spējīgi pat uz nelieliem upuriem savu tautiešu--zemnieku labā. Sacēlušos agrārā programma stādīja uzdevumu veidot ``brālību starp šķirām''. Praksē tas nozīmēja to, ka nacionālā sacelšanās kustība pakļāva zemnieku intereses muižnieku interesēm. Dabīgi, ka tas neapmierināja zemniekus.

Arī nostādnes nacionālajā jautājumā neveicināja plašu cittautiešu masu atbalstu sacelšanās kustībai. Poļu sabiedrisko darbinieku demokrātisma pārbaudei sava veida lakmusa papīrs bija attieksme pret lietuviešu, ukraiņu un baltkrievu tautām. Jā, Centrālās Nacionālās Komitejas programmā gan tika runāts par ar Poliju apvienoto tautu ``tiesību cienīšanu''. Taču tāpat tajā gāja runa par to, ka sacēlušos mērķis ir atjaunot Poliju 1772.~gada robežās, tātad~--- iekļaujot tajā arī lietuviešus, ukraiņus, baltkrievus un daļēji arī latviešus. Šīs prasības izpilde nozīmētu ukraiņu un latviešu apdzīvoto teritoriju sadali divās daļās, bet lietuviešu un baltkrievu zemju pilnīgu iekļaušanu Polijā. Tātad apvienojot vienu (poļu) tautu, tiktu dalītas un pakļautas poļiem citas tautas.

Par agrākās Žečpospolitas robežu atjaunošanu, šo tautu apdzīvoto teritoriju pievienošanu Polijai vairāk vai mazāk apzināti iestājās vairākums poļu patriotu. Ā.~Čartorijska piekritēji t.s. austrumu nomaļu (\pltxti{Kresy Wschodnie}) iedzīvotājus uzskatīja par ``vienotas poļu tautas'' daļu. L.~Mieroslavskis un viņa piekritēji tāpat noklusēja kārtu un nacionālo apspiešanu ``kresos'' un uzskatīja lietuviešus, ukraiņus un baltkrievus tikai par poļu nācijas ``atzarojumiem'', uzstājās pret kādu ``īpašu'' tiesību piešķiršanu tiem. ``Vienotas tautas teoriju'' atbalstīja Ā.~Mickēvičs un J.~Lelevels (kurš nākotnē redzēja tikai demokrātiskās Polijas pilsoņus, kuri runās poliski, ukrainiski, lietuviski), kaut arī viņi 1772.~gada robežas prasīja jaunai, pārveidotai Polijas valstij, centās apvienot visas tautības pret apspiedējiem. Emigrantu organizācija ``\pltxti{Lud Polski}'' (``Polijas tauta'') 1772.~gada robežās vēlējās radīt līdztiesīgu tautu valsti, uz šī pamata jau bija jāizveidojas vienotai nācijai no poļiem, lietuviešiem, ukraiņiem un baltkrieviem. Atsevišķas poļu emigrantu grupas runāja par piekāpšanos lietuviešiem un ukraiņiem valodas un kultūras jomā, citas~--- par zināmas autonomijas piešķiršanu. Tādejādi radikālākās poļu aprindas 1772.~gada robežu jautājumu saistīja ar sabiedrisko attiecību demokratizāciju šajās teritorijās. Taču poļu revolucionāru vidū bija arī tādi: P.~Scegennijs, B.~Švarce, J.~Dombrovskis, kuri pret ``vienotas tautas teoriju'' izturējās visai skeptiski.

Loģiski, ka poļu patrioti domāja galvenokārt par to, kāda situācija pēc nacionālās atbrīvošanās kustības uzvaras izveidosies pašā Polijā, bet ne aiz tās robežām. To viņiem nevar pārmest. Taču reizē jākonstatē, ka cerība ar pašu spēkiem un ar Rietumeiropas palīdzību izcīnītu uzvaru veda gan pie pretinieku~--- triju Poliju sadalījušo lielvalstu spēku nenovērtēšanas, gan lietuviešu, ukraiņu un baltkrievu tautu nacionālās un sociālās atbrīvošanās centienu neievērošanas. Tikai poļu sacelšanās vadītājs Lietuvā Z.~Serakovskis izvirzīja ideju par Polijas, Ukrainas, Baltkrievijas un Lietuvas autonomiju pārveidotas demokrātiskas Krievijas sastāvā.

Poļu patriotu kreisais spārns gan panāca, ka Centrālā Nacionālā Komiteja (\pltxti{Centralny Komitet Narodowy}) 1862.~gada septembrī vēstulē avīzes ``\rutxti{Колокол}'' (``Zvans'', 1857--1867) redakcijai iekļāva vārdus: ``Mēs vēlamies atjaunot Poliju tās agrākajās robežās, atstājot tautām tajās dzīvojošām: lietuviešiem un rutēņiem brīvību palikt savienībā ar Poliju vai rīkoties ar savu likteni pēc savas gribas''. Centrālā Nacionālā Komiteja solīja veltīt visus spēkus lai piesaistītu ``\pltxti{kresus}'' kopīgai cīņai ar poļiem. Taču skaidras nostādnes, kā to izdarīt, nebija.

Sakarā ar samierniecisko pozīciju sociālajos jautājumos Centrālajā Nacionālajā Komitejā nebija arī vienotas nostājas, kāda valsts iekārta nākotnē tiks pasludināta. Augustā tika arestēts un Varšavas citadelē ieslodzīts J.~Dombrovskis. (Viņu notiesāja uz 15~gadiem katorgas darbos, taču 1864.~gada beigās viņam izdevās izbēgt no pārsūtīšanas cietuma Maskavā un doties uz ārzemēm.) Neraugoties uz komitejas locekļu pašu republikānisko pārliecību, tā savos dokumentos vairījās lietot vārdu ``republika'', jo baidījās zaudēt sabiedrības virsslāņu atbalstu. Tiesa, savukārt baidoties no tautas masu neapmierinātības, nekas netika teikts arī par monarhijas pasludināšanu.

Var secināt, ka neraugoties uz visiem poļu revolucionāru demokrātiskajiem, nacionālās atbrīvošanās centieniem, lielā mērā taisnība bija Krievijas ārzemju ticību garīgā departamenta (\rutxti{Департамент духовных дел иностранных исповеданий}) direktoram V.~Skripicinam, kurš vēstulē Krievijas kara ministram D.~Miļutinam rakstīja: ``Toreiz [Žečpospolitas laikā] poļu muižniecība bija kolektīva valdošā dinastija, bet tagad tā ir kolektīvs pretendents [uz varu], kurš tāpat kā visi pretendenti, nekad neatteiksies no savām zaudētajām tiesībām un patiesi nepakļausies nekādai augstākai varai, kura nenāktu no viņa paša.''

Septembrī Centrālā Nacionālā Komiteja publicēja paziņojumu, ka tā turpmāk darbosies kā ``nacionāla valdība'' un sāka veidot savu varas aparātu. Agrāko piecu guberņu vietā tika veidotas astoņas vojevodistes, uz vietām tika nozīmēti komisāri, kuriem bija jārūpējas par Komitejas rīkojumu izpildi. Varšavā tika organizēta slepena nacionālā policija, kura apsargāja Centrālo komiteju. Šis policijas rindās iestājās arī daudz oficiālās policijas ierēdņu, kuri par pēdējās darbību ziņoja sazvērniekiem.

Bez tam Centrālā Nacionālā Komiteja, lai traucētu poļu iesaukšanu Krievijas armijā, publicēja dekrētus, ar kuriem lika visiem poļiem~--- pilsētu, apriņķu un guberņu padomju locekļiem atlaist šīs iestādes, noteica nodokli, kurš bija jāmaksā visiem poļiem. Tas sastādīja 0,5\% no nekustāmā īpašuma vērtības un 5\% no ienākumiem. Ja pirmajam dekrētam vairums turīgo padomju locekļu nepakļāvās, vienīgi izsakot savu protestu pret nelikumīgo iesaukšanu, bet dažas padomes atteicās sūtīt savus pārstāvjus iesaukšanas komisijās, tad otrais dekrēts tika pildīts visai cerīgi. Nodokļa iekasētāji apstaigāja mājas un, saņēmuši vajadzīgo summu, atstāja numurētas kvītis. To numurus publicēja avīze, tā pierādot, ka naudu saņēma nacionālā valdība. Kā liecinājis laikabiedrs V.~Pšiborovskis, nodokli visi maksāja bez pretošanās un negriezās pie oficiālās varas pēc aizsardzības. Maksāja ne tikai sarkanie, bet arī baltie un pat daži cara iestāžu ierēdņi. Pašā Centrālajā Nacionālajā Komitejā, kuras sastāvs pastāvīgi mainījās, norisa strīdi par sacelšanās sākšanas laiku.

1862.~gada beigās konspiratīvā organizācija aptvēra ap 20~tūkstošus cilvēku. Visur tika dibinātas šūniņas triju cilvēku sastāvā, katrs tās dalībnieks zināja tikai savus līdzbiedrus un ``desmitnieku'', kas apgrūtināja sazvērestības atklāšanu. Tika izstrādāti vairāki sacelšanās plāna projekti, par tiem norisa strīdi organizācijas vadība, kuras locekļu sastāvs savstarpējās idejiskās cīņās bieži mainījās. ``Sarkano'' provinciālās komitejas darbojās arī Lietuvā un Labā krasta Ukrainā.

Centrālās Nacionālās Komitejas pārstāvji veda pārrunas ar krievu revolucionāriem. 1862.~gada septembrī Londonā notika sarunas ar carismam opozicionārās avīzes ``\rutxti{Колокол}'' izdevējiem A.~Hercenu un N.~Ogarevu. Viens no viņiem~--- A.~Hercens~--- jau 50.~gadu vidū izvērsa poļu neatkarības un krievu un poļu revolucionāru savienības nepieciešamības propagandu. Būdams krievu patriots, bet nosodīdams Krievijas valdības rīcību Polijā, viņš rakstīja: ``Mēs vēlamies Polijas neatkarību, jo mēs vēlamies Krievijas brīvību'', ``Polijas atbrīvošana, tai pieguļošo apgabalu un Krievijas atbrīvošana~--- nav šķiramas.'' Avīze ``\rutxti{Колокол}'' sludināja: ``Polijai pieder \emph{neatņemamas} pilnas tiesības uz valstisku patstāvību, \emph{neatkarīgu} no Krievijas.'' Reizē A.~Hercens noteikti noraidīja nacionālistiskos apgalvojumus par to, ka Polijas tiesības uz neatkarību ir saistītas ar ``vēsturiskajām'' tiesībām uz Žečpospolitas veco robežu atjaunošanu, ieskaitot ukraiņu, baltkrievu un lietuviešu zemes. Viņš rakstīja: ``Nevienu nevajag ne rusificēt, ne polonizēt \citespace{} nevienam nevajag traucēt runāt un domāt, mācīties un rakstīt kā viņš vēlas''. Tai laikā tā bija revolucionāra nacionālā jautājuma izpratne. A.~Hercens izteica ticību, ka poļu patrioti un krievu demokrātija varēs pilnībā samierināties, pauda cerību par poļu tuvināšanos slāvu pasaulei. Viņš prasīja, lai krievu karavīri zvērētu poļu sacelšanās gadījumā necelt ieročus pret Poliju, bet uzskatīja, ka Polijai jāizpērk ``sava atsvešinātība, savs rietumnieciskais aristokrātisms, uzticība pāvestam''. Tiesa, A.~Hercena izvirzītais slāvu tautu federācijas lozungs bija utopisks, tāpat kā viņa ideja par zemnieku kopienas sociālisma iespējamību. Taču lozungam bija demokrātisks saturs, jau tāpēc vien tas augstu vērtējams.

Pēterburgā 1862.~gada novembrī~--- decembrī norisa poļu revolucionāru sarunas ar slepenās revolucionārās biedrības ``Zeme un Brīvība'' (``\rutxti{Земля и воля}'', 1861--1864) Centrālo komiteju. Ar poļu revolucionāriem bija saistīti ``Krievu virsnieku komitejas'' (''\rutxti{Комитет русских офицеров}'', 1861--1863) Polijā pārstāvji.

Kā ievērības vērta jāuzsver poļu vēsturnieka S.~Keņēviča doma, ka šajā laikā poļu un krievu revolucionāru solidaritātes idejai nācās lauzt gadsimtu gaitā starp tautām uzcelto aizspriedumu, nelabvēlības, asu savstarpēju pretenziju sienu. Pret revolucionāru sadarbības piekritējiem tika izvirzīti pārmetumi par ``pārdošanos'' mūžsenajam ienaidniekam.

Rīkojās arī oficiālās varas iestādes. Lai izolētu gatavojamās sacelšanās potenciālos dalībniekus, tās pēc A.~Veļepoļska iniciatīvas iecerēja veikt 20--30~gadu vecu vīriešu rekrūšu iesaukumu armijā. Atšķirībā no nesen izsludinātajiem jaunajiem iesaukšanas noteikumiem, tam bija jānotiek nevis lozējot, bet pēc jau iepriekš sagatavotiem sarakstiem. Tajos bija ietverti, pēc dažādām ziņām no 10~000 līdz 12~000 poļu jauniešu, kuri bija zināmi varas iestādēm kā demonstrāciju un ielu nekārtību dalībnieki. Visiem bija skaidrs, kāds ir šādas iesaukšanas kārtības mērķis. Tai bija vai nu jāsagrauj revolucionāru organizācija, vai jāizsauc priekšlaicīga sacelšanās, kuru varētu viegli apspiest. Pat daudzi augstākie ierēdņi un militārpersonas, ieskaitot pašu vietvaldi lielkņazu Konstantīnu iebilda pret šādu provokāciju. Taču A.~Veļepoļskis panāca sava projekta apstiprināšanu Pēterburgā. ``Augonis ir nobriedis un to vajag uzšķērst,~--- paziņoja A.~Veļepoļskis saviem oponentiem.~--- Sacelšanos es apspiedīšu nedēļas laikā un tad varēšu pārvaldīt [karalisti].''

Paziņojums par priekšā stāvošo iesaukumu, kura datums netika atklāts, izsauca lielu satraukumu sacelšanās gatavotāju nometnē. Organizācija nebija gatava sākt sacelšanos, bet tās biedri, kuriem draudēja iesaukums, prasīja rīkoties. Lai nezaudētu lielu daļu atbalstītāju un autoritāti tautā, Centrālā Nacionālā komiteja paziņoja, ka tā aizstāvēs jauniešus pret iesaukšanu un ``pirms naidnieks pagūs veikt barbarisko rīcību, jau sitīs atmodas stunda''. Tādejādi faktiski tika atzīts, ka rekrūšu iesaukšana nozīmēs sacelšanās sākumu. Uz Parīzi tika nosūtīta cilvēku grupa ieroču iepirkšanai, virsnieki emigranti tika aicināti atgriezties dzimtenē.

Krievu revolucionāri vienprātīgi iestājās pret nekavējošos sacelšanās sākumu Polijas karalistē, jo Krievijā viņi vēl nebija gatavi sākt bruņotu cīņu pret carismu, atsevišķas poļu sacelšanās panākumiem neticēja, bet gan uzskatīja, ka tās sakāve tikai nostiprinās carismu un uz ilgu laiku aizturēs revolūciju Krievijā. Avīzes ``\rutxti{Колокол}'' izdevēji, griežoties pie krievu virsniekiem Polijā, brīdināja: ``Priekšlaicīgs sprādziens Poliju neatbrīvos, jūs pazudinās un noteikti apstādinās mūsu krievu lietu''. Krievu revolucionāri ieteica poļu biedriem sākt sacelšanos ne agrāk kā 1863.~gada pavasarī, kad Krievijā pēc dzimtbūšanas atcelšanas bija jābeidzas zemnieku ``pagaidu atkarības stāvoklim'' (\rutxti{временнообязанное состояние}) no muižniekiem un varēja gaidīt, ka zemnieki pieprasīs zemi. Tā tika cerēts apvienot revolucionāro kustību Krievijā un Polijā. Taču A.~Veļepoļskim izdevās uzspiest Centrālajai Nacionālajai komitejai kauju valdības izraudzītajā brīdī.

Faktiski 1862.~gadā, tāpat kā 1830.~gadā, lēmumu par sacelšanās sākšanu pieņēma neliela cilvēku grupa, kurai nebija konkrētas programmas un darbības plānu nākotnei. Taču poļu vēsturnieki uzsver, ka sacelšanās vēriens, tās ilgums liekot atmest pieņēmumus par sacēlušos vieglprātību un politiskā brieduma trūkumu. Nacionālās emocijas bijušas tik saspringtas, ka varējušas izlauzties jebkurā brīdī. Tas esot bijis vairāk goda, nekā politikas jautājums. Pēc autora domām, jautājumos, kur iet runa par tūkstošiem cilvēku dzīvībām un labklājību, revolucionāri, vai tie būtu jakobīņi, boļševiki vai poļu nacionālisti, diemžēl parasti izturas ar pārlieku pašpārliecinātību, padodas sava mesiānisma sajūtai.

Tostarp 23.~decembrī cara policija atklāja Centrālās Nacionālās komitejas galveno tipogrāfiju, arestēja tās locekli B.~Švarci. Pastāvēja aizdomas, ka ar viņu policijas rokās nonākuši arī Komitejas dokumenti. 24.~decembrī tika saņemta ziņa no Parīzes, ka arestēti visi turp izbraukušie iepirkt ieročus. Baidoties no jauniem arestiem, steidzami tika izveidots atkal jauns Centrālās Nacionālās Komitejas sastāvs, kurš apstiprināja iepriekšējo lēmumu par sacelšanās sākumu līdz ar rekrūšu iesaukšanu. Arestētais J.~Dombrovskis no cietuma atsūtīja plānu, kurš paredzēja 4~000 sacelšanās dalībnieku uzbrukumu Novogeorgijevskas cietoksnim (krievu \rutxti{Новогеоргиевская крепость}, poļu \pltxti{Twierdza Modlin}), tā arsenāla sagrābšanu un reizē sacelšanās sākšanos visā karalistē. Centrālā Komiteja plānu pieņēma un uzsāka pēdējos sacelšanās sagatavošanas pasākumus.

Centrālā Nacionālā komiteja gatavojās sākt sacelšanos 1863.~gada 14.~(26.)~janvārī. Taču varas iestādes nolēma pagrīdniekus apsteigt. Naktī uz 3.~(15.)~janvāri sākās rekrūšu iesaukums armijā, par kuru gan sacelšanās gatavotāji savlaicīgi uzzināja un brīdināja organizācijas biedrus. Apmēram tūkstotis jauniešu bēga no Varšavas, pulcējās mežos, organizēja pirmās sacēlušos vienības, apbruņotas ar izkaptīm, nažiem, daļēji medību bisēm. Tomēr simti pilsētā vēl palikušo, kas neslēpās mežā un kam nebija nodoma pretoties, tika sapulcināti citadelē. Šķita, ka iesaukšana ir izdevusies un A.~Veļepoļskis bija apmierināts.

Bija skaidrs, ka tuvākajās dienās piespiedu iesaukšana turpināsies. Centrālā Nacionālā komiteja nolēma sākt sacelšanos ātrāk nekā bija plānots~--- naktī no 10.~(22.). uz 11.~(23.)~janvāri.

Pirms paša sacelšanās sākuma 10.~janvārī Centrālā Nacionālā komiteja pasludināja sevi par \strong{Nacionālo Pagaidu valdību} (\pltxti{Tymczasowy Rząd Narodowy}), publicēja savus programmatiskos dokumentus: manifestu un agrāros dekrētus.

Manifestā tika pasludināts, ka Polija ``negrib un nevar'' bez pretestības pakļauties tai vardarbībai, ko pār to realizē krievu carisms~--- nelikumīgajai rekrūšu iesaukšanai. Visas kārtu un reliģiskās privilēģijas tika pasludinātas par atceltām, proklamēta ebreju vienlīdzība. Centrālā Nacionālā komiteja kā tagad vienīgā likumīgā poļu nacionālā valdība aicināja Polijas, Lietuvas un Krievijas tautas uz cīņu par atbrīvošanos. Noslēgumā dokumenta autori griezās pie krievu tautas. Centrālā Nacionālā komiteja paziņoja, ka tā nevaino krievu tautu par noziegumiem pret Poliju, jo tā pati cieš no carisma jūga, izteica cerību, ka tā neatbalstīs tirānu un brīdināja, ka pretējā gadījumā būs neizbēgams karš starp divām tautām.

Tika publicēti arī dekrēti par zemnieku apstrādājamās zemes platības nodošanu viņu rokās ar sekojošu kompensāciju zemes īpašniekiem no valsts kases, kā arī par neliela zemes gabala (3 morgu jeb ap 1,5 ha) piešķiršanu sacelšanās dalībniekiem~--- bezzemniekiem no nacionālajiem fondiem.

Vērtējot manifestu un agrāros dekrētus, padomju vēsturnieks M.~Misko uzsvēris, ka agrārajā jautājumā tie bija liels solis uz priekšu salīdzinājumā ar 1831.~gadu. Pirmkārt, tie pilnībā atcēla visus zemnieku feodālos pienākumus~--- gan klaušas, gan činšu. Otrkārt, zemniekiem tika piešķirtas visas to rīcībā esošās zemes (agrāk līdz 3 morgiem lielos zemes gabalus muižnieks varēja zemniekam atsavināt). Treškārt, reforma attiecās ne tikai uz muižnieku, bet arī uz valsts un baznīcas zemniekiem. Beidzot, ceturtkārt, tā klusu ciešot atzina servitūtu palikšanu zemnieku rokās. Kopumā sacēlušos izsludinātā reforma zemniekiem deva daudz vairāk nekā cara valdības un A.~Veļepoļska dekrēti. Iniciatīva bija sacēlušos rokās. Tomēr dekrētos nekas nebija teikts par zemes piešķiršanu bezzemniekiem, izņemot tos, kuri piedalīsies sacelšanās. Un arī tad zeme bija solīta tikai pēc sacelšanās beigām un tikai trīs morgu apmērā, kas bija nepietiekoši ģimenes uzturēšanai. Bez tam, dekrētos nekas nebija teikts par pēdējos gadu desmitos zemniekiem atsavināto zemju likteni, tātad tās paliktu muižnieku rokās. Beidzot, reforma paredzēja atlīdzību muižniekiem par zaudētajām zemēm, kas būtībā nozīmēja, ka zemniekiem (gan ar citu iedzīvotāju slāņu palīdzību) būtībā būs jāmaksā izpirkšanas maksa par savu zemi. Reformas īstenošanas princips nebija pietiekami skaidrs. To vajadzēja realizēt vietējām varas iestādēm un sacēlušos pavēlniecībai. Taču sacēlušos vienības sākumā izveidojās tikai dažviet un tā reformas īstenošana tika bremzēta. Reformas realizēšanā netika aicināti iesaistīties paši zemnieki. Arī priekšlikums uzdot reformu īstenot muižniekiem tika noraidīts, jo sacēlušies baidījās izsaukt to pretestību un vispārēju sajukumu valstī. Vēlāk sacēlušos militārajiem priekšniekiem tika uzdots nepieļaut sociālu apvērsumu. Zemnieku uzbrukumu gadījumā muižniekiem vai to mantai vainīgie bija jāsoda, viņu mājas jānodedzina. Tādejādi progresīvie dekrēti tikai daļēji atbilda sākušās sacelšanās izvēršanas veicināšanai.

Tas pats jāsaka par nacionālā jautājuma risināšanu. Centrālās Nacionālās Komitejas manifests neko neteica par ukraiņu, baltkrievu, lietuviešu pašnoteikšanās tiesībām, kuras taču Komiteja bija it kā atzinusi vēstulē žurnāla ``\rutxti{Колокол}'' redakcijai 1862.~gada septembrī. Tādejādi ``vienotās poļu nācijas'' teorija, mēģinājumi praksē notušēt nacionālās pretrunas, faktiski vienoja poļu demokrātus ar konservatīvajiem.

\strong{Sacelšanās} sākās nelabvēlīgos apstākļos. Atšķirībā no 1830.~gadā poļu rīcībā nebija labi apmācītas armijas. Sacelšanās dalībniekiem nebija nedz ieroču, nedz finanšu. Valstī bija savākts tikai ap 600~medību bišu. Organizācijas kasē bija ap 7,5 tūkstoši rubļu. Sacelšanās dalībnieki nebija militāri apmācīti. Zemnieki nebija gatavi sākt sacelšanos. Krievu revolucionāri savu sacelšanos plānoja vēlā pavasarī. Poļu revolucionāri uzsāka cīņu ziemā, kura neapmācītām vienībām radīja lielākus škēršļus kā labāk sagatavotajai krievu armijai. Tai Polijas karalistē bija ap 100~000 vīru ar 176 lielgabaliem.

Sacelšanās militāro vadību 18.~janvārī Centrālā Nacionālā komiteja uzticēja ģenerālim L.~Mieroslavskim, kurš tika pasludināts par diktatoru. Virkne vēsturnieku vērtē šo lēmumu kā kļūdainu, jo viņš savas personīgās intereses stādījis augstāk par nacionālajām, pēc sava rakstura neesot bijis spējīgs apvienot sacēlušos vadošo kodolu. 1861.~gada oktobrī viņš ar Itālijas valdības atļauju bija organizējis Dženovā (\pltxti{Genova}) virsnieku skolu (\pltxti{Polska Szkoła Wojskowa w Genui}) sacelšanās militārās vadības sagatavošanai. Tajā iestājās daudz emigrācijā esošu poļu jauniešu. Tiesa, drīz skolā sākās nesaskaņas sakarā ar L.~Mieroslavska dižmanību un viņš skolu atstāja, skaļi paužot savu neapmierinātību. Taču 1862.~gada jūnijā Krievija atzina jauno nacionālās atbrīvošanās kustības rezultātā izveidoto Itālijas valsti, viens no līguma nosacījumiem ar to bija poļu militārās skolas slēgšana. Itālija to arī lika poļiem izdarīt.

Naktī uz 11.~(23.)~janvāri cīņās iesaistījās ap 6~000 sacelšanās dalībnieku, sapulcinātu vairāk nekā 30~vienībās. Sajukums organizatoriskajā sacelšanās sagatavošanā noveda pie tā, ka iecerētā poļu un krievu revolucionāri noskaņoto virsnieku kopīgā uzstāšanās nenotika, tā vieta norisa 18 uzbrukumi Krievijas armijas garnizoniem. (Pavisam tie atradās 180~vietās.) Taču uzbrukumi lielākoties bija neveiksmīgi. To militārais efekts bija niecīgs, toties politiskās sekas~--- katastrofālas. Tā, mēģinājums ieņemt Plockas (\pltxti{Płock}) pilsētu, kur sacēlušies gribēja ierīkot savu centru, beidzās ar neveiksmi, ap 150 sacēlušos tika sagūstīti. Neviena guberņas pilsēta netika ieņemta. (Uz Varšavu, kur atradās liels krievu armijas garnizons, sākotnēji sacelšanās neizplatījās. Tas paglāba to no apšaudes ar lielgabaliem no citadeles. Cietoksnī atradās 555 lielgabali. Tāpēc nav brīnums, ka sacēlušies cietoksni ieņemt nevarēja. Revolucionāru organizācijas autoritāte tika satricināta. Jau pēc sacelšanās sākuma Varšavas citadeles garnizons tika pastiprināts līdz 16~000 vīriem.) Poļu sacēlušos uzbrukumos cara armija cieta niecīgus zaudējumus, taču tās personālsastāvā dzima naids pret uzbrucējiem. Kaut daļa krievu revolucionāri noskaņoto virsnieku pārgāja sacēlušos pusē, daži pat neveiksmīgi mēģināja aizvest sev līdzi uz sacēlušos nometnēm veselas karaspēka daļas, jau pirmās sacelšanās dienas parādīja, ka plašākas poļu un krievu sadarbības nodomi ir nereāli. Visi plāni par krievu revolucionāru vienību (družīnu) izveidi, kas dotos palīgā poļiem, izrādījās neizpildāmi, kaut kā atsevišķi cīnītāji sacēlušos pusē darbojās vairāki simti krievu, ukraiņu, baltkrievu un citu Krievijas tautu pārstāvju. (Šveices armijas virsnieks H.f.~Erlahs, kurš sacelšanās laikā kā novērotājs atradās sacēlušos rindās, pēc tās publicētā grāmatā atzīmēja, ka ``Starp ik dienas vienībā ienākušajiem brīvprātīgajiem vienmēr noteikts procents ir krievu bēgļu''. Vēlāk cariskās varas iestādes īpaši bargi sodīja poļu pusē pārgājušos Krievijas armijas karavīrus. No 183 Krievijas gūstā kritušajiem pārbēdzējiem 89 tika izpildīts nāves sods.)

Viens no krievu revolucionāru vadītājiem M.~Bakuņins, kurš jau no 40.~gadiem bija saistīts ar poļu revolucionāriem, ierosināja izveidot krievu leģionu, kurš cīnītos sacēlušos pusē. Ar 1863.~gada 9.~(21.)~februāri datētā vēstulē poļu revolucionāriem viņš rakstīja: ``Nav nekādu šaubu, ka zem lieliskā ``Zemes un brīvības'' karoga pastāvošs krievu nacionālais leģions radītu lielu iespaidu krievu armijā, kas virzīta ne tikai pret jums, bet pret visu Krieviju. Pats tā pastāvēšanas fakts būtu līdzvērtīgs vairākām uzvarētām kaujām. Par nelaimi, tagad šī projekta īstenošana ir kļuvusi daudz grūtāka, nekā, piemēram, pirms mēneša''. Viens no krievu virsniekiem--revolucionāriem Polijā A.~Potebņa gan vēl mēģināja gūt viena no populāriem sacēlušos vienību vadītājiem M.~Langeviča atbalstu krievu leģiona idejai, taču pēdējais vēsi izturējās gan pret poļu un krievu revolucionāru sadarbību, gan konkrēto ideju. A.~Potebņa drīz krita kaujā pret Krievijas karaspēku un ideja tā arī palika nerealizēta, krievu karavīru masveida pāreja sacēlušos pusē nenotika. Faktiski tā arī nevarēja notikt, jo sacēlušies poļi turpināja uz krieviem raudzīties kā ``barbariem''. ``Prom uz Āziju, Čingizhana [mongoļu karavadonis un valdnieks] pēcteči''~--- šie vārdi no kādas poļu nacionālās atbrīvošanās kustības laiku dziesmām ataino noturīgo poļu masu apziņas līmeņa stereotipu par Maskavas ``barbarijas'' civilizācijas līmeni.

L.~Mieroslavskis, šķērsojis pie Pozenes Polijas robežu, savāca ap 400--500 cilvēku lielu vienību, taču divās sadursmēs krievu karaspēks to izklīdināja, iegūstot trofejas, arī L.~Mieroslavska saraksti. Pats viņš bēga atpakaļ uz Parīzi, taču turpināja sevi uzskatīt par diktatoru.

``Balto'' vadība neticēja sacelšanās panākumiem, nejuta līdzi sacēlušamies un sākotnēji pat plānoja viņus publiski nosodīt. Viņi baidījās, ka L.~Mieroslavskis varētu piešķirt kustībai sociālas revolūcijas raksturu, virzīt zemnieku jautājumu uz sociālistisku risinājumu. Nevarot novērst sacelšanos, ``baltie'' nolēma tai pievienoties, lai pārņemtu tās vadību.

Izmantojot L.~Mieroslavska neveiksmes, ``baltie'' panāca, ka februāra beigās par diktatoru izsludināja M.~Langeviču, bet jau pēc nedēļas krievu karaspēka spiediena rezultātā viņš pārgāja Galīcijas robežu un Austrijas varas iestādes viņu internēja. No sacelšanās dalībnieku viedokļa ļaunākais bija tas, ka M.~Langevičs to izdarīja tikai sava štāba pavadībā un slepus no vienkāršajiem cīnītājiem. Visi bija sašutuši, un, saucot viņu par ``nodevēju'', arī paši devās pāri Austrijas robežai. Zemnieki sagūstīja un varas iestādēm nodeva ap 200~sacelšanās dalībnieku. Lielkņazs Konstantīns 1863.~gada 17.~februārī pat telegrafēja Aleksandram II: ``Zemnieki labākajā noskaņojumā, priecājas, ka viss ir beidzies''. Līdz šim nenorimst strīdi par M.~Langeviča rīcības motīviem. Parasti poļu vēsturnieki atzīst, ka viņš vēlējies Galīcijā rast papildspēkus un ieročus, bet daži autori viņam pārmet ``negodīgumu'', savukārt konservatīvais vēsturnieks V.~Pšiborovskis ir pievedis liecību, ka viņš vēlējies pārtraukt bezcerīgo sacelšanos.

Krievu karaspēks tika koncentrēts, tā izvietošanas vietu skaits samazinājās četras reizes. Tas deva iespēju sacelšanās dalībniekiem ieņemt plašas teritorijas, ieskaitot daudzas apriņķu pilsētas. Tomēr februārī Krievijas karaspēks uzsāka uzbrukumus lieliem spēkiem un sacēlušos stāvoklis pasliktinājās. Janvārī notika 58~karaspēka sadursmes ar sacelšanās dalībniekiem, februārī~--- 76. Visā Polijas karalistē atkal ieveda kara stāvokli, izsludināja, ka tie dumpinieki, kas kritīs gūstā ar ieročiem rokās, tiks tiesāti lauku kara tiesā pēc saīsinātas tiesvedības un sodi tiks izpildīti nekavējoties, visu sacelšanās dalībnieku īpašumiem uzlika sekvestru (rīcības ierobežojumus).

Tomēr 1863.~gada 31.~martā cars izdeva manifestu, kurā piedāvāja amnestiju visiem, kuri līdz 1.~maijam noliks ieročus. Manifestā Aleksandrs II arī solīja piešķirt karalistei vietējās pašpārvaldes tiesības. Carisms vēl neatteicās no piekāpšanās politikas. Taču tad pat aprīlī ``Centrālajā Nacionālajā komitejā'' vairākumu ieguvušie un maijā par Nacionālo valdību (\pltxti{Rząd narodowy}) to pārdēvējušie ``baltie'' piedāvājumu noraidīja. Tika cerēts, ka Rietumvalstis veiks intervenci pret Krieviju, lai palīdzētu poļiem. Tā kā Ā.~Čartorijskis bija jau miris un viņa vietu labējās emigrācijas vadībā ieņēma viņa dēls V.~Čartorijskis, viņš arī tika iecelts par sacēlušos galveno diplomātisko pārstāvi ārzemēs. Plaši izplatījās baumas par drīzu Francijas un Anglijas iejaukšanos. Tāpēc sacelšanās it kā varēja nevis degt pilnās liesmās, bet tikai pamazām gruzdēt, lai tikai saglabātu pamatu ārvalstu intervencei.

Aprīlī un maijā karalistē notika 127~bruņotas sadursmes. Tikmēr Krievijas karaspēks koncentrējās vairākos desmitos Polijas karalistes pilsētu. Vasaras mēnešos Polijas karalistē norisa 238 bruņotas sadursmes. Sacēlušos kara darbība pārvērtās par partizānu cīņām. Drīz kļuva redzams, ka poļi ar militārām metodēm uzvarēt nespēs. 1863.~gada maijā Polijas karalistē par vietvalža pienākumu izpildītāju, bet oktobrī~--- par vietvaldi tika iecelts ģenerālis F.~Bergs, vecs karotājs, tai skaitā arī 1831~.~gada poļu kampaņas dalībnieks. A.~Veļepoļskis jūlijā palūdza atvaļinājumu un devās uz ārzemēm, bet augustā tika atlaists no visiem amatiem. Tas arī nozīmēja, ka viņa politika bija piedzīvojusi sakāvi. Poļu sabiedrībā dažādi vērtēta viņa personība un politiskā pozīcija. Revolucionāri viņu nosodīja kā nacionālo interešu nodevēju, kurš stādījis uzdevumu poļu mantīgos slāņus samierināt ar carismu, saglabājot Krievijas kundzību poļu zemēs. Liberāļi A.~Veļepoļska realizētajā politikā saskatīja kreiso uzstāšanās rezultātā neizmantotu iespēju uzlabot Polijas stāvokli.

Ar jauna vietvalža nozīmēšanu iezīmējās Krievijas pāreja uz terora politiku Polijas karalistē. F.~Bergs ieveda stingru policijas un vietējās administrācijas uzraudzību. Sākās sacelšanās dalībnieku publiska sodīšana. Ja agrāk viņu nošaušana notika citadelē, tagad to rīkoja Varšavas laukumos. Tika pastiprināti policijas spēki. Varšavā agrāko tūkstoš policistu vietā darbojās 3 tūkstoši. Policijā iesauca regulārās armijas karavīrus un virsniekus. Tika aizliegta pārvietošanās naktīs, izbraukšana no pilsētas bez speciālas atļaujas. Naktīs pilsētu apsargāja daudzas kājnieku un jātnieku patruļas. Sākās poļu tautības ierēdņu nomaiņa ar speciāli atlasītiem un uz Poliju nosūtītiem krievu ierēdņiem. Tika ieviesta iedzīvotāju kolektīvā atbildība. Ja kādā ēkā atrada ieročus vai slapstošos sacelšanās dalībnieku, tās īpašnieku un iedzīvotājus tiesāja kara tiesa. Atentātu liecinieki, kuri nepalīdzēja izmeklēšanai notvert vainīgos, tika uzskatīti par līdzdalībniekiem un atbildēja pēc kara stāvokļa likumiem. Svarīgs iedzīvotāju ietekmēšanas līdzeklis bija to aplikšana ar kontribūcijām, naudas sodi un konfiskācijas. Plaši tika praktizētas sacelšanās dalībnieku muižu un citu īpašumu konfiskācijas.

Taču ``baltie'' turpināja izvairīties no radikālas agrārā jautājuma risināšanas, cerot uz Rietumvalstu palīdzību pret Krieviju, nevis enerģiski kopojot tautas spēkus. Arī materiālie resursi netika mobilizēti pietiekamā mērā. Mantīgo šķiru pārstāvji gan maksāja nacionālās nodevas un pirka nacionālā aizņēmuma obligācijas. Diezgan plaši līdzekļi ieplūda no ārzemēs dzīvojošajiem emigrantiem, kuri grupējās ap \pltxti{Hotel Lambert} un atbalstīja ``baltos''. Taču daudzi bagātnieki arī izvairījās no finansiālajām saistībām. Ar to mantīgās šķiras kopumā sacelšanās labā ziedoja daudz mazāku daļu savu līdzekļu, nekā vēlāk zaudēja no cariskās valdības uzliktajiem sodiem un kontribūcijām. Ja sākotnēji sacēlušies savāca 400 tūkstošus zlotu (atcerēsimies, 1 zlots bija līdzvērtīgs 15 kapeikām), tad 1863.~gada jūnijā Varšavā no Polijas karalistes galvenās kases tika nolaupīti vairāk nekā trīs miljoni rubļu, citur~--- vēl viens miljons. (Tiesa, no galvenās kases nolaupītajā summā tikai 0,5 miljoni bija skaidra nauda, citi līdzekļi bija vērtspapīros, kurus Nacionālā valdība tā arī nespēja realizēt). Bez tam sacēlušies pret ``kvītīm'' rekvizēja zirgus, pajūgus, apģērbu un pārtiku. Naudu vāca arī, pieprasot nomaksāt nodevas divus gadus uz priekšu, pat aplaupot iedzīvotājus.

Sagādātos līdzekļus izmantoja ieroču iepirkšanai ārzemēs un pašā Polijā. Pēc dažām ziņām sacelšanās laikā tika iepirkts ap 100~000 karabīņu. Taču nelaime bija, tā, ka daļu ieroču sacelšanās vadītāji turēja noliktavās, daļu konfiscēja prūšu un austriešu varas iestādes, pirms tie tika nogādāti uz Polijas karalisti, daļa nonāca cara varas iestāžu rokās, tā kā sacēlušies ieguva tikai mazu daļu. 1863.~gada jūlijā ar karabīnēm bija apbruņoti tikai ap 13~000 sacēlušos. Puse no sacelšanās dalībniekiem karoja, bruņoti ar izkaptīm un pīķiem.

Sacelšanās dalībnieku vienības sastāvēja no kājniekiem, nelielām kavalērijas daļām, tehniskajiem vadiem un visai nelielas dažās vienībās esošas artilērijas, kuru daļēji veidoja koka lielgabali. Pēc sociālā sastāva visdažādākā bija kavalērija, tomēr vairākumu tajā veidoja šļahtiči. Kājnieki bija strēlnieki un ``izkaptnieki'' (\pltxti{kosynierzy}), kuri galvenokārt nāca no zemnieku vidus. Gandrīz katrā vienībā bija ārsts un ksendzs.

Kaut 22.~janvāra agrārie dekrēti tika īstenoti, no 1863.~gada 1.~aprīļa tika pārtraukta činša iekasēšana, nacionālās varas iestādes brīdināja muižniekus nepieprasīt atcelto saistību izpildi un stingri sodīja likuma neievērotājus, tomēr, kaut daudzi poļu zemnieki piedalījās sacelšanās jau no pirmajām tās dienām, zemnieku pamatmasa palika pasīva. Tai pat laikā daudzi zemnieki uzstājās arī pret sacelšanās dalībniekiem, uzlūkojot tos par muižnieku interešu aizstāvjiem. Sacēlušies pirms tās sākuma nebija spējuši pierādīt zemniekiem savu mērķu tautisko, atbrīvošanās raksturu. Tās bija arī sekas šļahtiču privilēģijām un sacelšanās dalībnieku neieinteresētībai panākt vispārēju vienlīdzību. Agrārajai reformai netika veltīta vajadzīgā uzmanība. Nabadzīgākie zemnieki un bezzemes algādži turpināja palikt ārpus reformas, muižnieki prasīja no tiem agrāko saistību izpildi. Arī par zemes piešķiršanu sacelšanās dalībniekiem~--- solītajiem 3 morgiem~--- tika runāts reti. Šis solījums faktiski netika izpildīts. No otras puses, cara varas iestādes visādi centās nomelnot sacelšanās kustību kā muižniecisku un pret zemniekiem vērstu. Padomju vēsturnieks M.~Misko uzsvēris, ka zemnieku uzstāšanās pret sacelšanos tās pirmajās nedēļās jāuzskata nevis kā antinacionāla, bet kā sociāla, pret muižniekiem virzīta kustība.

Atbildot uz zemnieku uzstāšanos pret sacelšanās dalībniekiem--muižniekiem, sacēlušies pielietoja represijas. Daudzi zemnieki tika apvainoti spiegošanā, aģitācijā pret ``nacionālo'' valdību un nošauti vai pakārti. Ir grūti noteikt precīzu sacelšanās upuru skaitu, bet kā uzskatīja poļu profesors S.~Kenevičs, gada laikā Poļesjes (\pltxti{Polésie}) rajonā vien no sacēlušos rokām gāja bojā 200--300 zemnieku. Krievu vēsturnieks emigrācijā A.~Kersnovskis rakstīja, ka kopumā sacēlušies nogalinājuši līdz 5~000 mierīgo iedzīvotāju, galvenokārt etnisko poļu. Toties ir zināms tikai viens gadījums, kad pēc viena no sacēlušos vadītājiem Z.~Padļevska pavēles arī kāds muižnieks tika nošauts par atteikšanos sniegt palīdzību sacelšanās dalībniekiem un griešanos pie cara varas iestādēm ar lūgumu sniegt viņam militāru apsardzi.

Neraugoties uz vairākām neveiksmēm, 1863.~gada vasarā sacelšanās kustība aptvēra gandrīz visu Polijas karalisti, kā arī daļu Lietuvas, Baltkrievijas un Labā krasta Ukrainas teritorijas. Kopā kustībā piedalījās ap 200~000 cilvēku.

Jau 1862.~gada vasaras beigās izveidojās lietuviešu un baltkrievu demokrātu organizācija, kas uzturēja sakarus ar Z.~Serakovska pulciņu Pēterburgā. Redzamākie baltkrievu un lietuviešu demokrāti: publicists, dzejnieks K.~Kaļinovskis un garīdznieks (ksendzs) A.~Mackēvičs atbalstīja poļu atbrīvošanās kustību. Tā K.~Kaļinovskis uzskatīja, ka baltkrievu revolucionāriem ir jāorientējas uz Poliju, ka poļu atbrīvošanās kustībā varēs garantēt pašvaldību arī baltkrievu un lietuviešu apdzīvotajām teritorijām. Viņa devīze bija: ``Poļu lieta~--- tā ir mūsu lieta, tā ir brīvības lieta''. Arī A.~Mackēvičs cerības uz lietuviešu tautas atbrīvošanos saistīja ar cīņu pret carismu kopā ar poļu, krievu un baltkrievu tautām.

Pēc poļu sacelšanās sākuma 1863.~gada martā~--- aprīlī tā plaši izplatījās arī Lietuvā un Rietumbaltkrievijā. Lietuvā sacelšanās ieguva plašas zemnieku uzstāšanās raksturu, kura attīstījās agrāras revolūcijas virzienā. Martā beigās šeit militāro vadību uzņēmās Z.~Serakovskis. Baltkrievijā sacelšanās attīstījās vājāk. Vairākums sadursmju ar carisko armiju norisa Grodņas guberņā. Krievija veica pasākumus, lai sacelšanās neizplatītos tās rietumu guberņās. Šo guberņu zemnieki jau 1863.~gada pavasarī tika atbrīvoti no ``pagaidu atkarības'' stāvokļa un kļuva par pilntiesīgiem savas zemes saimniekiem, un viņu izpirkšanas maksājumi samazināti. Tas būtiski traucēja sacelšanās izplatību. Tomēr galvenais sacelšanos bremzējošais faktors bija tas, ka lietuviešu, baltkrievu un ukraiņu zemnieki neatbalstīja prasību atjaunot Žečpospolitu 1772.~gada robežās.

Var piebilst, ka arī jaunlatviešu izdotās ``Pēterburgas Avīzes'' nesimpatizēja poļiem, jo poļu kustības darbinieki bija muižnieki un garīdznieki, un paši poļu, bet īpaši leišu zemnieki tiem dažviet izrādīja pretestību. ``Pēterburgas Avīzes'' nosodīja poļu sacelšanos.

Piemēram, kad baltkrievu apdzīvotajos novados kara darbību uzsāka no poļu apgabaliem ieradušās vienības, tās šeit lielākoties ātri tika sakautas. Baltkrievu zemnieki dažkārt pat kopā ar cara armiju cīnījās pret K.~Kaļinovska vienībām. (Neraugoties uz to, viņš līdz pat savam arestam 1864.~gada janvārī turpināja atbalstīt poļu brīvības cīnītājus.) Arī Vladimiras--Volinskas (poļu \pltxti{Włodzimierz Wołyński}, krievu \rutxti{Владимир-Волынскuй}) rajonā ap pusotra tūkstoša zemnieku ar izkaptīm pievienojās krievu karaspēkam lai ``iztīrītu'' apkaimi no sacelšanās dalībniekiem. Kaut daudzi šļahtiču izcelsmes poļi arī juta līdzi apspiestajām ukraiņu, baltkrievu un lietuviešu tautām, pat piedalījās to nacionālās atbrīvošanās kustībā, poļu muižas šais zemēs, tāpat kā carisms, bija nacionāls un sociāls apspiedējs. Ukraiņu, baltkrievu un lietuviešu nāciju veidošanās norisinājās cīņā gan pret carismu, gan poļu muižu.

Līdzīgi notikumi norisa \strong{Latgalē}. Tur jau pirmajās sacelšanās dienās ieradās vairāki emisāri no Polijas un Lietuvas, kuri centās piesaistīt kustībai arī vietējos zemniekus, tāpat rīkojās vietējie muižnieki un mācītāji, aicinot uz cīņu par Polijas brīvību, aizmirstot naidu pret kungiem un ar tiem kopā dodoties cīņā pret ''moskaļiem'' (krieviem). Kā konstatējis poļu tautas vēsturi daudz pētījušais vēsturnieks Ē.~Jēkabsons, notika savdabīga sacensība par Latgales zemnieku simpātijām. Krievijas valdība 1863.~gada martā atļāva zemniekiem iepirkt zemi bez muižnieku piekrišanas un izsludināja likumu par sacelšanās dalībnieku īpašumu konfiskāciju. Savukārt Polijas Pagaidu valdība aizliedza muižniekiem prasīt, lai zemnieki pildītu klaušas. Taču muižnieki nevēlējās aizliegumam pakļauties. Rezultātā Polijas atjaunošanas ideju zemnieki saistīja ar panu (muižnieku) varas pastiprināšanos. 1863.~gada aprīlī Latgales poļu muižnieku vienība grāfa L.~Plātera vadībā uzbruka vāji apsargātam krievu armijas ieroču transportam ceļā no Daugavpils (vācu \detxti{Dünaburg}, poļu \pltxti{Dyneburg}, krievu \rutxti{Динабург}, \rutxti{Двинск}) uz Krāslavu (poļu \pltxti{Krasław}, krievu \rutxti{Креславка}). Krievijas varas iestādes savukārt aicināja zemniekus uzbrukt L.~Plātera vienībai, kurā darbojās arī vairāki īpaši viņu neieredzēti muižnieki. Zemnieki sakāva un sagūstīja poļu nodaļu. L.~Plāters tika ieslodzīts Dinaburgas cietoksnī, kur viņam arī tika piespriests un izpildīts nāves sods. Savukārt aktīvākajiem zemniekiem, kas bija piedalījušies sacelšanās dalībnieku vajāšanā, žandarmērija izmaksāja balvās 450~sudraba rubļus.

Samērā maz sacēlušos vidū bija ebreju. Ebreji varētu to atbalstīt, ja tās uzvara nestu viņu stāvokļa uzlabošanos. Taču nopietns šķērslis ebreju dalībai sacelšanās kustībā bija poļu šovinisms. Poļu turīgo slāņu vidū bija daudz cilvēku, kuri naidīgi izturējās pret ebrejiem un bija pret vienlīdzīgu tiesību piešķiršanu viņiem, jo baidījās no to konkurences. Tikai sākotnējā, manifestāciju periodā ebreji plaši piedalījās kustībā. Iziet uz ielas vai organizēt aizlūgumu sinagogā viņi varēja brīvi, daudz grūtāk bija iestāties sacēlušos vienībās, kuras komandēji nacionālistiski noskaņotie šļahtiči. Pēc dažām ziņām sacelšanās kustībā piedalījās tikai 100 līdz 200 ebreju. Vairāk viņu bija sacēlušos civilajās varas iestādēs. Pārējās ebreju masas parasti neatšķīrās no citiem līdzpilsoņiem: akurāti maksāja nacionālos nodokļus, sniedza sacelšanās organizācijām dažādus pakalpojumus.

XIX gadsimta 60.~gados atkal, tiesa, pēdējo reizi šī gadsimta vēsturē, \strong{Polijas jautājums nonāca Eiropas diplomātijas uzmanības centrā}. Romas pāvests lika baznīcās aizlūgt par cietēju Poliju, taču pirmkārt viņš rūpējās par katoļu baznīcas pozīciju stiprināšanu tajā. K.~Markss un F.~Engelss, vērtējot poļu sacelšanos pirmkārt no starptautiskā viedokļa, kļūdaini uzskatīja, ka ``Eiropā atkal ir plaši atklāta revolūciju ēra''. F.~Engelss rakstīja K.~Marksam: ``Poļi ir malači. Un ja viņi noturēsies līdz 15.~martam, tad visa Krievija sāks kustēties''. Un atkal tā bija kļūda. Masu kustības atplūdi Krievijā, sākušies 1861.~gadā, šajā attīstības etapā izrādījās vēsturiski neatgriezeniski.

Sarežģītāka bija pasaules \strong{lielvalstu valdību pozīcija}. Visvairāk ieinteresēta atbalstīt poļu tautas cīņu bija Francija. Neatkarīga Polija tai būtu izdevīga sabiedrotā ne tikai pret Prūsiju un Austriju, bet arī Krieviju, kura bija viens no galvenajiem Francijas sāncenšiem kontinentā. Napoleons III labprāt vēlētos pārdalīt Eiropas karti, apvienot savā ietekmē Itāliju, pievienot Francijai Reinas kreiso krastu. To nevarēja panākt bez Krievijas atbalsta. Meklējot šādu atbalstu, Napoleons solīja Krievijai atteikties no oficiālas palīdzības poļu kustībai. Tiesa, ievērojot savas valsts sabiedrisko domu, Napoleons III ieteica Krievijai veikt piekāpšanās poļu priekšā. Cerības uz Krievijas atbalstu Francijai nepiepildījās un Francijas imperators Napoleons~III, kurš bija velti cerējis uz Francijas tuvināšanos Krievijai, vērstu pret Prūsiju, tad nu iestājās poļu pusē. Franciju vismaz vārdos atbalstīja arī Lielbritānija un Austrija. Taču Krievijas imperators Aleksandrs~II Francijas sūtnim Krievijā L.~de Montebello paziņoja, ka viņš principā piekrīt Francijas, Lielbritānijas un Austrijas priekšlikumiem, bet nespēj tos realizēt praksē: ``Ja varētu atgriezt Polijas karalistei tās neatkarību, tad mans tēvs to izdarītu un es arī rīkotos tāpat. Taču [tās] neatkarība ir praktiski neiespējama: Polija nespēj dzīvot savās robežās, un paši poļi to neslēpj. Viņi vēlas iegūt savas 1772.~gada robežas, citiem vārdiem, vēlas sadalīt Krieviju''.

Lielbritānijai faktiski nerūpēja poļu intereses. Polijas valsts atjaunošana stiprinātu britu galveno sāncensi kontinentā~--- Franciju un vājinātu sabiedroto~--- Prūsiju. Gan Lielbritānijas valdība, gan parlaments, gan prese izteica līdzjūtību poļu tautai, taču tai pat laikā pauda viedokli, ka nekādu palīdzību poļu kustībai viņu valsts nesniegs. Lielbritānija tikai diplomātiski atbalstīja sacēlušos. Taču, kā norādījuši poļu vēsturnieki J.~Klačko un H.~Vereščuckis, par līdzsvara saglabāšanu Eiropā norūpējušās Lielbritānijas acīs neatkarīga Polija būtu pastāvīgs nemiera izraisītāja Eiropas centrā, jo Prūsija un Krievija censtos atgūt savus agrākos valdījumus, nestabilitāte traucētu brīvo tirdzniecību, kurā Lielbritānija bija ieinteresēta. Bez tam Polijas atrašanās Krievijas sastāvā Lielbritānijai bija izdevīga, jo tā pastāvīgi izraisīja Krievijas rūpes un kalpoja tās kā Lielbritānijas sāncenses vājināšanai. Tāpēc Lielbritānijas ārlietu ministrs Dž.~Rossels piedāvājumu atjaunot Poliju, iekļaujot tajā arī Pozeni un Galīciju, nosauca par ``neprātu''.

Kad Polijas karalistē 1863.~gadā sākās atklāta sacelšanās, A.~Hercens savā pirmajā tai veltītajā rakstā uzdeva jautājumu: ``Ko darīs Eiropa? Vai viņa pieļaus [poļu sacelšanās apspiešanu]? Vai viņa vēlreiz pieļaus?'' un pats atbildēja: ``Pieļaus''. Tā tas arī notika. Tikai atsevišķi slāvu un Rietumeiropas zemju demokrātisko aprindu pārstāvji kā brīvprātīgie devās cīnīties sacēlušos rindās. Īpaši spēcīgu atbalsi poļu sacelšanās sastapa Itālijā un Ungārijā. Ungāru revolucionāri vēlējās atdalīties no Austrijas. Itāļu patrioti cerēja atbrīvot no Austrijas kundzības Venēciju. Abu tautu revolucionāriem galvenais ienaidnieks bija Austrija. Pret to viņi vēlējās noslēgt savienību ar poļiem, lai kopīgiem spēkiem izraisītu sacelšanās pret Austriju Galīcijā, Ungārijā, Čehijā un Balkānos. Poļu vadītāji turpretī nevēlējās izvērst cīņu pret Austriju, cerot, ka tā netraucēs sacēlušos vienību organizācijai Galīcijā.

Austrija un Prūsija, kā valstis, kuras kopīgi ar Krieviju bija sadalījušas Poliju, bija pret piekāpšanos poļiem Polijas karalistē, kas varētu stimulēt poļu nacionālo kustību arī kaimiņvalstīs. Prūsija, kā jau minēts, 1863.~gada 8.~februārī (pēc vecā stila~--- 27.~janvārī) pat noslēdza ar Krieviju t.s. Alvenslēbena konvenciju, kura atļāva Krievijas karaspēkam vajāt poļu sacēlušos arī Prūsijas teritorijā. Austrijas valdība gan nepievienojās Prūsijas un Krievijas 8.~februāra konvencijai, bet pavēlēja atbruņot poļu sacēlušos vienības, kas šķērsoja valsts robežu, paziņoja aizliegumu Galīcijas teritorijā vervēt brīvprātīgos sacēlušos vienībās un izvirzīja papildus sava karaspēka daļas pie Krievijas robežas. Austrijas salīdzinoši ``mīkstā'' politika pret poļu sacelšanos, tās dalībnieku vidū un arī Francijā radīja zināmas cerības par Austrijas gatavību palīdzēt poļu lietai, taču tās nepiepildījās.

Tomēr poļu sacēlušos vadītāji utopiski cerēja uz Rietumu atbalstu. 1863.~gada 30.~jūnijā britu avīze ``\entxti{Morning Standart}'' rakstīja: ``Poļu dumpis beigtos pats no sevis, ja tā vadītāji necerētu uz Rietumvalstu militāru iejaukšanos''. Tika domāts par poļu spēku taupīšanu, lai izmantotu tos vēlāk, kad palīgā Polijai nāks ``visa Eiropa''. Kad dažviet vietējie militārie priekšnieki veica pasākumus, lai mobilizētu vīriešus vecumā no 18 līdz 30~gadiem, sacēlušos valdība izdeva dekrētu, kas atļāva pieņemt tikai tādu skaitu brīvprātīgo, kādam pietika ieroču. Atsevišķi sacēlušos vienību komandieri vairījās iesaistīties kaujā pret cara armijas daļām, tā vietā veicot nogurdinošus pārgājienus, kuri veda pie to dalībnieku spēku izsīkuma, bada, dezertēšanas. Atbilde uz pārmetumiem par tādu taktiku bija: ``Kaut ar mazu vienību, bet sagaidīt [Rietumu] intervenci.'' Sadursmes notika galvenokārt tad, kad cara karaspēks panāca kādu sacēlušos vienību un tai nācās vest kauju bez iepriekšējas sagatavošanās.

Tāpat tika uzskatīts, ka nevajag steigties ar ārzemju revolucionāru palīdzības pieņemšanu, jo tā varēja pasliktināt Rietumvalstu valdību attieksmi pret Poliju. Maijā notā, adresētā Francijas valdībai, Nacionālā valdība rakstīja, ka uzskata Rietumvalstu karu pret Krieviju par ``neizbēgamu'', tā kā tikai ar spēku iespējams piespiest Krieviju piekāpties; ka Polija ārvalstu armiju aizsardzībā ir spējīga organizēt 150~000 vīru lielu savu armiju un ir gatava samaksāt visus kara izdevumus. Vienlaikus izskanēja arī brīdinājums, ka palīdzības aizkavēšanās ``piespiedīs meklēt palīdzību revolucionāro partiju rindās''.

Tiesa, Francijas un Lielbritānijas valdības poļu sacelšanās sakarā sūtīja notas uz Pēterburgu, prasot ievērot Vīnes kongresa (1815) lēmumus. (Austrija tās nevarēja atbalstīt, jo pati 1846.~gadā bija tos pārkāpusi, likvidējot Krakovas republiku). Ar grūtībām Lielbritānijas, Francijas un Austrijas valdības saskaņoja un jūnijā nosūtīja Krievijai savas prasības: 1)~pasludināt pilnīgu amnestiju sacelšanās dalībniekiem, 2)~radīt nacionālu poļu valdību, 3)~pieļaut plašu poļu dalību administratīvajā pārvaldē, 4)~ieviest apziņas un katoliskās reliģijas ritu izpildes brīvību, 5)~poļu valodu atzīt par oficiālu visās iestādēs, 6)~noteikt likumīgu kārtību iesaukšanai armijā. Bez tam Krievijai tika piedāvāts noslēgt pamieru ar sacelšanās dalībniekiem un sasaukt lielvalstu, kuras bija parakstījušas Vīnes līgumu, konferenci, lai noteiktu kārtību, kā norādītās prasības jāīsteno dzīvē. Austrijas valdība, piedaloties šo prasību izstrādē, centās panākt, lai tās dotu karalistes iedzīvotājiem ne vairāk tiesību kā poļiem to bija Galīcijā. Prasības neatbilda sacēlušos vēlmēm iegūt pilnīgu savas tautas neatkarību. Sacēlušos valdība bija vīlusies, sagaidīja priekšlikumus ar ``sāpēm un izbrīnu'', vērtēja tos kā ``vislielāko pakalpojumu'' cara režīmam. Tā bija sarūgtināta, ka pamieram bija jāattiecas tikai uz Polijas karalisti, neskarot Krievijas rietumu guberņas (kaut bija zināms, ka Rietumu lielvalstis nekad pat nedomāja atzīt šo zemju piederību poļiem).

Cara valdība, pārliecinājusies, ka vismaz Austrija un Lielbritānija negatavojas iesaistīties militārā konfliktā, triju valstu prasības noraidīja, vienlaikus piekrītot sarīkot triju Poliju sadalījušo valstu konferencei, lai apspriestu stāvokli šo valstu poļu valdījumos un saskaņotu to atbilstoši ``laikmeta prasībām un apstākļiem''. Rietumu sabiedrībā, īpaši Francijā, Krievijas atbilde izsauca sašutumu, taču viena pati karot Francija nespēja, bet tās partnervalstis karu nevēlējās. Lielbritānija ar Polijas sacelšanos savu mērķi jau bija panākusi~--- nepieļāvusi Krievijas un Francijas tuvināšanos un vājinājusi Krievijas pozīcijas pasaulē. Tagad tā varēja arī piekāpties.

Jūlijā Nacionālā valdība publicēja aicinājumu ``Eiropas tautām un valdībām'', kurā sauca tās pārtraukt jebkādas attiecības ar carisko Krieviju un atzīt, ka tai nav nekādu tiesību valdīt Polijā.

Taču ``Eiropas valdības'' atteicās atzīt Poliju par karojošo pusi, kas varētu atvērt tai iespēju legāli iepirkt ārzemēs ieročus un vervēt brīvprātīgos. Sacēlušies tā arī palika dumpinieku statusā.

ASV Ziemeļu štati tai laikā karoja pret vergturu intereses atbalstošo Dienvidu štatu Konfederāciju. Lielbritānija vēlējās izmantot konfliktu savas ietekmes stiprināšanai Amerikas kontinentā. Krievija, kura 1861.~gadā bija atcēlusi dzimtbūšanu, ieņēma draudzīgu pozīciju pret ziemeļniekiem un pat nosūtīja divas eskadras uz Amerikas krastiem. Kā raksta krievu kara vēstures speciālists A.~Širokorads, Vidusjūrā uz Lielbritānijas jūras komunikācijām izgāja fregate ``\rutxti{Олег}'' un korvete ``\rutxti{Сокол}'' (``Vanags''). Vēl pirms tam Orenburgas ģenerālgubernators A.~Bezaks uzsāka ekspedīcijas korpusa formēšanu karagājienam uz Afganistānu un Indiju. Tas gan notika slepus, bet sekoja informācijas noplūde britu presei. Rietumu biržās sākās panika, kuģniecības kompānijas paaugstināja pārvadājumu cenas, apdrošinātāji mainīja apdrošināšanas noteikumus. Krievijas militāri-diplomātiskie soļi izrādījās pareizi. Galu galā Lielbritānija atteicās no sava nodoma iejaukties Amerikas Pilsoņu karā konfederātu pusē, ziemeļnieki turpināja dienvidnieku piekrastes blokādi, Dienvidiem katastrofāli sāka trūkt munīcijas un viņi karu zaudēja. ASV atbalstīja Krieviju Polijas jautājumā.

Galu galā iespējamo Krievijas pretinieku intereses bija pārāk atšķirīgas, lai tās nopietni iejauktos notikumos. Tā neviena valdība neatsaucās aicinājumam. Neviena valsts nevēlējās uzsākt karu ar Krieviju lai panāktu poļu brīvību. Rietumu lielvalstu politika vienīgi radīja veltas ilūzijas sacēlušos vidū, bet to palīdzības izpalikšana izsauca poļos lielu vilšanos.

1863.~gadā starptautiskā situācija pakāpeniski kļuva poļiem nelabvēlīga. Prūsija un Austrija tuvinājās kopīgai cīņai pret Dāniju. Drīz sacēlušies zaudēja savu atbalsta bāzi Galīcijā. 1864.~gada februārī Austrija tur ieveda kara stāvokli, arestēja vairākus tūkstošus cilvēku, likvidēja organizācijas, kuras palīdzēja cīnītājiem Polijas karalistē. Labvēlīgas izturēšanās pret poļu sacelšanos maska ar to tika nomesta. Poļu patriotu cerības uz Austrijas palīdzību galīgi izplēnēja.

Gan iekšējo, gan starptautisko notikumu pavērsiena rezultātā ``baltie'' no sacelšanās vadības tika izstumti un no 1863.~gada septembra vadošo vietu atkal ieguva ``sarkanie'', taču vienotības viņu rindās nebija. Oktobrī par diktatoru kļuva R.~Trauguts. Viņš kā virsnieks bija dienējis Krievijas armijā, 1862.~gadā atstājis dienestu, no 1863.~gada aprīļa vadījis sacēlušos vienību Grodņas un Minskas guberņās. R.~Trauguts sevi nepieskaitīja nevienai revolucionārajai partijai, pieturējās mēreniem uzskatiem, bet noteikti atbalstīja 22.~janvāra manifesta un agrāro dekrētu principus. Viņš mēģināja piesaistīt sacelšanās kustībai zemniekus. Taču neejot tālāk par 22.~janvāra dekrētiem, nesolot zemniekiem piešķirt vismaz daļu muižniekiem atsavinātas zemes, nedodot zemi bezzemniekiem, uz ko R.~Trauguts un viņa palīgi nebija spējīgi, zemnieku ietekmēšanai nebija citu reālu līdzekļu. Mēģinājumi aktivizēt sacelšanos izrādījās neveiksmīgi. Atsevišķi darbojošās sacēlušos vienības cieta sakāvi pēc sakāves. Decembrī tika reorganizēti sacēlušos bruņotie spēki. No atsevišķu vojevodistu vienībām tika veidota regulāra armija, sastāvoša no korpusiem, divīzijām, pulkiem, bataljoniem. Ar to tika stiprināta disciplīna, savstarpējā sadarbība. Taču tad pat~--- rudens un ziemas mēnešos skaidri iezīmējās poļu mantīgo slāņu, tai skaitā arī Pozenē un Galīcijā, atiešana no sacelšanās. Daudzi muižnieku un buržuāzijas pārstāvji atteicās maksāt nacionālo nodokli, devās uz ārzemēm.

Sacelšanās dalībnieku darbību traucēja arī bargie laika apstākļi. Tūkstoši cilvēku no sacēlušos vienībām tās pameta un vai nu bēga uz ārzemēm vai ieradās pie varas iestādēm nožēlot nodarīto un lūgt žēlastību. Ziemā aktīvo sacelšanās dalībnieku skaits sasniedza vien dažus tūkstošus. Tie formāli dalījās četros korpusos, no kuriem divi tā arī faktiski netika izveidoti. No 1863.~gada septembra līdz 1864.~gada aprīlim Polijas karalistē notika 392~bruņotas sadursmes, lielākoties gan visai nelielas.

Krievu sabiedrības lielākā daļa poļu sacelšanos uztvēra klaji negatīvi. Carisms, izmantojot revolucionārās kustības kritumu Krievijā pēc dzimtbūšanas atcelšanas, varēja koncentrēt spēkus cīņai pret Polijas atbrīvošanās kustību. Biedrība ``Zeme un brīvība'' gan mēģināja organizēt sacelšanos Pievolgā un Urālu apkaimē, bet neveiksmīgi. Arī poļu revolucionārs H.~Keņevičs ar vairāku Kazaņas garnizona virsnieku palīdzību, kuru bruņojumā bija daži revolveri, mēģināja organizēt zemnieku sacelšanos, izplatot Tambovas, Penzas un Saratovas guberņās viltotus it kā cara izdotus manifestus. Visi grupas dalībnieki tika arestēti un sodīti ar nāvi 1864.~gada jūnijā. Ir gan zināmi gadījumi, kad sacēlušos pusē pārgāja krievu virsnieki un karavīri, bet tā savu gatavību aizstāvēt citas tautas brīvību izpauda tikai atsevišķi cilvēki. Šādos apstākļos Krievijā veikt revolucionāro darbu kļuva ļoti grūti. Pakāpeniski novājinājās ``Zemes un brīvības'' organizācija. Drīz Krievijā vairs nebija revolucionāra centra. ``Acīmredzot,~--- rakstīja F.~Engelss K.~Marksam 1863.~gada jūnijā,~--- poļu sacelšanās šai ziņā radīja nelabvēlīgu ietekmi''. Secinājums atbilda patiesībai.

Liela daļa izglītotās krievu sabiedrības ieguva pārliecību, ka poļu sacelšanos izraisīja vienīgi Krievijas valdības pārliekā piekāpība uz nacionāliem sapņiem tendēto poļu priekšā. Muižnieku sapulces, pilsētu domes, universitātes, konservatori un liberāļi, baznīcas kalpotāji sūtīja Krievijas imperatoram savus vēstījumus, kuros atbalstīja soda ekspedīciju darbību Polijā. Kā atzina A.~Hercens, ``muižniecība, liberāļi, literāti, zinātnieki un pat skolnieki ir visi inficēti; viņu asinīs un audos ir iesūcies patriotiskais sifiliss''. Jau pēc dažiem mēnešiem pēc tam, kā A.~Hercens bija atklāti nostājies sacēlušos poļu pusē un paziņojis, ka ``neklusēs veselas tautas noslepkavošanas priekšā'', žurnāla ``\rutxti{Колокол}'', kurš pirms tam ar carisko ierēdņi patvaļas atmaskošanu Krievijā bija iekarojis milzu popularitāti, tirāža kritās piecas reizes. Poļu sacelšanos atbalstošie A.~Hercens un M.~Bakuņins strauji zaudēja popularitāti. Konservatīvā krievu prese A.~Hercenu un viņa piekritējus nesauca citādi kā par nodevējiem un poļu kustības aģentiem. (Pēc dažiem gadiem~--- 1867.~gada 29.~aprīlī A.~Hercens atzina, ka žurnāla ``\rutxti{Колокол}'' atbalsts poļiem 1860.--1863.~gadā bija ``kolosāla kļūda''. Tagad viņš bija pārliecināts ``par pilnīgu šīs drosmīgās un dumjās nācijas nespēju, bīstamību un trulumu'', kuru dēļ viņš un N.~Ogarevs pazudināja savu reputāciju.) Iespējams, nesekmīgā poļu sacelšanās aizstāvība ietekmēja A.~Hercena veselību, viņš šķīrās no dzīves, būdams salauzts, draugu pusaizmirsts un ienaidnieku apmelots. Galu galā arī pats viņš savos memuāros ``\rutxti{Былое и думы}'' (``Pagājība un apceres'') ar rūgtumu rakstīja: ``Poļos katolicisms ir attīstījis to mistisko ekzaltāciju, kura pastāvīgi viņus notur rēgu pasaulē \citespace{} Mesiānisms ir sajaucis galvu simtiem poļu un pašam Mickevičam''.

Smagu triecienu sacēlušies saņēma, kad 1864.~gada 19.~februārī (2.~martā) Krievijas imperators Polijas karalistē izdeva dekrētu par zemes piešķiršanu zemniekiem, un visas pēc 1846.~gada viņiem muižnieku atņemtās zemes atdošanu tiem atpakaļ. Ja 1863.~gadā cara valdība nežēlīgi vajāja tos, kuri sludināja un īstenoja Centrālās Nacionālās Komitejas 22.~janvāra agrāros dekrētus, tagad tā zemnieku reformu Polijas karalistē ņēma savās rokās.

Faktiski 1863.--1864.~gada sacelšanās bija pagrieziena punkts carisma agrārajā politikā poļu zemēs, kura rezultātā poļu muižniekiem nācās šķirties no cerības saglabāt visu zemi savās rokās. Jauno Krievijas kursu noteica vairāki faktori. Pirmkārt, valdība Pēterburgā nevarēja nerēķināties ar sacēlušos valdības 1863.~gada 22.~janvāra dekrētiem, kuri paredzēja zemnieku apsaimniekotās zemes nodošanu viņu rokās pret izpirkšanas maksu, carismam bija jākonkurē ar sacēlušos šļahtu un jānotur poļu zemnieki savā pusē. Otrkārt, neraugoties uz zemnieku kustības karalistē visumā mierīgo raksturu, 60.~gadu sākumā tā tomēr signalizēja par zemnieku neapmierinātību, lika carismam ar sevi rēķināties. Treškārt, sacelšanās sākums pārsvītroja carisma mēģinājumus izlīgt ar poļu muižniekiem, tāpēc tagad ar to interesēm varēja nerēķināties.

Īstenojot reformu, zemnieki tika ``izspēlēti'' pret šļahtu. Viņi tagad saņēma visu savu agrāk apstrādāto zemi, kā arī ēkas, mājlopus, darba rīkus un sēklu privātīpašumā, par ko maksāja zemes nodokli 2/3 apmērā no agrākā činša. Muižnieki saņēma kompensāciju par zaudējumiem no valsts fonda, kuru veidoja zemnieku maksājamais zemes nodoklis. Arī bezzemnieki ieguva nelielus zemes gabalus no valsts zemēm. Zemnieku zemes platība pieauga par 28\%, rezultātā zemnieku saimniecību skaits jau 1864.~gadā sasniedza 514~tūkstošus, kad 1859.~gadā to bija 424~tūkstoši. Pēc citiem datiem ap 200~000 bezzemnieku ģimeņu ieguva nelielus zemes gabaliņus. 40.~gados pēc reformas muižu platība karalistē samazinājās par 14\%, taču vēl joprojām 61,4\% visas zemes atradās muižnieku rokās. Kaut zemnieku ieguvums bija neliels, vairs nebija cerību zemi saņēmušos iesaistīt cīņā pret lielvaru.

Svarīgi, ka Polijas karalistē zemnieki zemi ieguva privātīpašumā, bet nevis kopienas īpašumā, kā tas notika Krievijā. Pie tam poļu zemnieki ieguva visu savu agrāk lietoto zemi, bez kādiem ``nogriezumiem'' kā Krievijā, viņi kļuva par tās īpašniekiem tūlīt, bez ``pagaidu atkarības stāvokļa''. Zemnieki tika atbrīvoti arī no muižnieku aizbildniecības, to pašvaldības bija brīvas no šļahtas ietekmes, bet bija atkarīgas no apriņķa varas iestādēm. Protams, tāds Krievijas varas ``radikālisms'' bija izskaidrojams ar centieniem apturēt zemnieku dalību sacelšanās kustībā. Jāuzsver, ka tieši poļu nacionālās atbrīvošanās kustība imperatoru, tā valdību virzīja uz konsekventāku, demokrātiskāku zemnieku jautājuma risināšanu, nekā uz to bija spējīga šļahta. Zemniekos tika stiprinātas monarhiskās ilūzijas, cerības uz valdības pasākumiem zemnieku aizsardzībai pret muižnieku patvaļu. Cara reformai poļu zemnieku acīs bija arī tā priekšrocību, ka to realizēja ``likumīga'' vara, kuras rokās bija visi līdzekļi tās realizēšanai, nevis ``dumpinieki'', kuru pašu vidū agrārajā jautājumā nebija vienotības. Zemnieku sacelšanos vairs nebija. Zemnieki, norūpējušies par pēc iespējas izdevīgāku sev cara valdības pasludinātās agrārās reformas īstenošanu, atteicās atbalstīt sacēlušos, dažreiz pat izdeva tos cara varas iestādēm. Tiesa, vairākums no jaunajām saimniecībām bija līdz 1,5~ha lielas (pareizāk~--- mazas). Norisa intensīvs zemnieku saimniecību platību sarukšanas process, galvenokārt mantošanas rezultātā. No 1864. līdz 1889.~gadam sīko saimniecību (no 1,5 līdz 7,5~ha) skaits ātri pieauga. Šo pusproletārisko saimniecību īpašnieki bija spiesti meklēt darbu ārpus saimniecības, muižnieku laukos. Tā tika radīts lēts darbaspēks, kas bija svarīgs nosacījums muižnieku saimniecību pārkārtošanai uz kapitālistiskiem pamatiem.

Reforma neatrisināja servitūtu jautājumu par zemnieku tiesībām lietot ganības un mežu, kur tie varēja iegūt būvmateriālus un kurināmo. Pēc reformas servitūti kļuva gandrīz vai par galveno gminas (\pltxti{gmina}~--- no vācu \detxti{Gemeinde}, kopiena, arī administratīvi-teritoriāla vienība, līdzīga pagastam) un muižnieku pretrunu punktu.

Arī blakus Polijai esošajā Lietuvā, Baltkrievijā un Labā krasta Ukrainā 60.~gados agrārā reforma tika pasteidzināta uz zemniekiem izdevīgākiem nosacījumiem.

Sacēlušos vadība neatrada pārliecinošu atbildi uz cara valdības reformu. Agrāro reformu ``visa poļu sabiedrība'' vērtēja kā demagoģisku pasākumu, un tiešu uzbrukumu svētajam šļahtas privātīpašumam. R.~Trauguts izdeva uzsaukumu, kurā pārmeta cara valdībai liekuļošanu, paziņojot, ka zemnieku pašvaldības neesot iespējams realizēt viņu analfabētisma dēļ, un, visbeidzot, draudēja, ka visus ievēlētos kopienu vecākos Nacionālā valdība uzlūkos kā ``brāļu slepkavas'' un vajās. Protams, tāds uzsaukums nesastapa labvēlīgu attieksmi zemnieku vidē.

Polijas karalistes dienvidos un Lietuvā sacēlušies vēl turpināja karot arī 1863./1864.~gada ziemā, bet domstarpības to vidū palīdzēja carismam uzvarēt. 1864.~gada aprīlī R.~Trauguts tika arestēts, zemākie komandieri bēga uz ārzemēm. Faktiski 1864.~gada maijā sacelšanās jau bija apspiesta. Aprīlī vēl notika 34~sadursmes, maijā~--- 14, jūnijā~--- 4, jūlijā~--- 1. Vasaras beigās likvidēti tika arī pēdējie pretošanās centri Polijas karalistes dienvidos. R.~Traugutam un nelegālās valdības locekļiem tiesa piesprieda nāves sodu un augustā viņi tika pakārti. Pēdējā nemiernieku vienība ksendza S.~Bžuzkas vadībā bēguļoja pa mežiem līdz 1865.~gada aprīlim, kad viņš tika saņemts gūstā un pakārts.

Sacelšanās ilga gadu un četrus mēnešus. Bruņotā cīņā no poļu puses piedalījās ap 50~000 cilvēku. Pavisam notika 1~229 bruņotas sadursmes, Polija karalistē 956, Lietuvā un Baltkrievijā 237, Ukrainā 35. Sacēlušos zaudējumi sasniedza 20~000 (ir pat dati, ka ap 30~000) cilvēku. Brīvprātīgi padevās ap 15~000, tikpat tika ieslodzīti un izvesti uz Sibīriju. Pēc sacelšanās sagrāves uz Rietumiem bēga ap 7~tūkstoši tās dalībnieku. Krievu armijas zaudējumi bija 3~343 cilvēki, no tiem 2~169 ievainoto. Īpaši skarbi cara varas iestādes izrēķinājās ar sacelšanās dalībniekiem~--- bijušajiem Krievijas armijas karavīriem. No 183 to rokās nonākušajiem krieviem 89 tika sodīti ar nāvi. Pēdējais nāves sods tika izpildīts 1866.~gada 22.~novembrī (4.~decembrī). Krievija, apspiežot sacelšanos rīkojās nežēlīgi, taču tikpat nežēlīgi, ja ne vēl briesmīgāk, tai pat laikā rīkojās rietumnieciskā Lielbritānija, apspiežot sipaju sacelšanos Indijā (\entxti{India's First War of Independence, the Great Rebellion}, 1857--1859).

1863.--1864.~gada sacelšanās, tāpat kā iepriekšējās poļu bruņotās uzstāšanās, neaptvēra visas poļu apdzīvotās teritorijas, taču visām tām bija vispārnacionāls raksturs tai ziņā, ka tās izsauca poļu nacionālās jūtas arī citos apgabalos, saņēma no turienes palīdzību, veicināja nacionālo ideju izplatību arī ārpus sacelšanās aptvertās teritorijas.

Daži vēsturnieki 1863.--1864.~gada sacelšanos vērtē galvenokārt kā poļu muižniecības (dažkārt literatūrā poļus arī šai periodā vēl sauc par ``šļahtiču nāciju'') kārtējo mēģinājumu sasniegt Žečpospolitas atjaunošanu ar ārvalstu palīdzību, kurš kārtējo reizi noveda zemi vēl lielākā nelaimē.

Turpretī revolucionāri, īpaši marksisti, augstu vērtēja poļu tautas pūliņus sasniegt brīvību un demokrātiju, jo tās panākumi būtu visas Eiropas demokrātijas interesēs, varētu vājināt Krievijas carismu. Sacelšanos sakāve izsauca dziļu F.~Engelsa nožēlu. Viņš rakstīja: ``Paies daudz laika, līdz Polija atkal varēs pacelties, pat ar citu palīdzību, bet Polija mums ir pilnīgi nepieciešama.'' Acīmredzot ar ``mums'' bija domātas visu Eiropas revolucionāru intereses.

1874.~gadā F.~Engelss uzsvēra, ka Polijas neatkarība un revolūcija Krievijā bija savstarpēji saistītas. Abu uzvara nozīmētu to, ka Vācijas valdība un buržuāzija turpmāk varētu paļauties tikai uz saviem spēkiem, ar kuriem laika gaitā vācu strādnieki paši tiktu galā. Jāatzīst gan, ka marksisma klasiķis pārāk romantiski vērtēja revolūcijas perspektīvas gan Krievijā, gan Vācijā.

Padomju vēsturnieki, kuri bija spiesti visur pārspīlēt šķiru cīņas nozīmi vēsturē, tomēr dažādi vērtēja tās raksturu. Vieni uzskatīja, ka tai piemita nacionālās atbrīvošanās cīņas un antifeodāls raksturs. Sacelšanās pēc būtības bijusi buržuāziska revolūcija. Citi uzsvēra, ka 1863.--1864.~gada sacelšanās nekļuva par buržuāziski--demokrātisku revolūciju, bet Polijas vēsturē iegāja kā neizdevusies buržuāziska revolūcija. Pēc šī darba autora domām sacelšanās spilgtākā komponente bija nacionālā, kura sevī saturēja arī ievērojamu demokrātisma lādiņu, sacelšanās uzvara varētu pavērt ceļu arī demokrātiskiem procesiem (ja vien poļu nacionālisti nesāktu konsekventi apspiest citas tautas). Taču tā nenotika ne tikai Krievijas militārā pārspēka, Rietumeiropas valstu neieinteresētības, bet arī poļu šļahtas nacionālā un sociālā egoisma rezultātā.

Pēc sacelšanās izskanēja interesants secinājums. To izteica poļu marķīzs A.~Veļepoļskis, vietvalža lielkņaza Konstantīna palīgs: ``poļu labā var vēl ko izdarīt, taču kopā ar poļiem~--- nekad.'' Tiesa, vēlāk A.~Veļepoļskis precizēja, ka tie nebija viņa vārdi, viņš tikai pievērsis tiem tautiešu uzmanību. Iespējams arī, ka dusmu brīdī tos bija pateicis pats vietvaldis. Taču neatkarīgi no autora personības zināmu patiesības graudu tie saturēja.

Pastāv pamatots viedoklis, ka 1863.~gada sacelšanās varēja uzvarēt tikai tad, ja tā būtu nacionāla revolūcija ar radikālas buržuāziskas agrārās reformas programmu, ja to atbalstītu visaptveroša sociāla zemnieku revolūcija.

Baltkrievu publicists A.~Tarass izdalījis četrus galvenos \strong{sacelšanās sakāves cēloņus}:

\begin{enumerate}

\item Reālas palīdzības trūkums no Rietumeiropās valstu puses. Bet sacelšanās tika gatavota tieši cerot uz Rietumu iejaukšanos!

\item Sacelšanās vadītājiem trūka pārdomāta kaujas darbību plāna.

\item Carismam izdevās piesaistīt sev poļu, ukraiņu, baltkrievu un lietuviešu zemniecību. (Autoram jāpiebilst, ka carismam tas neizdotos, ja sacēlušies būtu mērķtiecīgi aizstāvējuši zemniecības intereses.)

\item Nežēlīgās M.~Muravjova un F.~Berga represijas. (Atkal jāpiebilst, ka pret pretiniekiem nežēlīgas represijas pielietoja arī sacēlušies. Ir kontrproduktīvi skaidrot, kura puse pirmā lietoja represijās, kura bija nežēlīgāka. Kara gaitā tādus rēķinus neviens nestādīja. Sacēlušos vadītāji, ja viņi nebija garīgi atpalikuši, bet viņi tādi nebija, nevarēja ar šādām sekām nerēķināties, tomēr uzņēmās atbildību par kustības izraisīšanu).

\end{enumerate}

% page 161


Sacelšanās Polijā un 1864.~gada agrārā reforma nosacīti bija arī robežšķirtne starp diviem laikmetiem Polijā, kas nodalīja feodālismu un kapitālismu. Jau minētie padomju vēsturnieki V.~Djakovs un I.~Millers gan precizēja, ka ne visas poļu zemes šo šķirtni pārvarēja vienlaikus. Teritorijās, kura bija pakļautas Prūsijai un Austrijai, tā tika sasniegta jau 1848.~gadā, bet Krievijai pakļautajās~--- 1864.~gadā. Vairākums vēsturnieku gan par tādu nosacīti pieņem otro gadskaitli. Tomēr, kā norāda poļu autors V.~Kula, nenākas fetišizēt vienu datumu. Robeža starp feodālismu un kapitālismu visas Polijas mērogā veidojās XIX gadsimta vidū, kad Pozenes un Galīcijas apgabalos tiesiski zemniecība jau bija izcīnījusi atbrīvošanos no ārpusekonomiskajiem spaidiem. T.s. ``kongresa Polijā'' agrārā reforma norisa kā pēdējā no trijām poļu zemēm.

Pati 1863.--1864.~gada sacelšanās bija vismasveidīgākais un ilgstošākais poļu tautas atbrīvošanās mēģinājums no svešinieku kundzības. Taču ar tās sakāvi uz gadu desmitiem cīņa par ``brīvību un neatkarību'' (\pltxti{wolność i niepodległość}) tika atvirzīta citu, ikdienišķāku uzdevumu priekšā.

Krievijas impērijā pēc 1863.--1864.~gada tika forsēts kurss uz nacionālo nomāļu integrāciju. Viens no cēloņiem tam bija valdības vilšanās cerībās uz mierīgas un liberālas nacionālās politikas realizācijas sekmību. Pārsvaru guva doma, ka tikai stingra nacionālā politika palīdzēs saglabāt un nostiprināt impērijas vienotību.

1863.--1864.~gada sacelšanās, neraugoties uz neveiksmi, veicināja feodālo attiecību likvidāciju Polijas karalistē, piespiežot carismu īstenot tajā agrāro reformu uz zemniekiem labvēlīgākiem nosacījumiem nekā pašā Krievijā. Pēc 1864.~gada zemnieku reformai sekoja citas. Tai pat demokrātiskajā garā bija sastādīts 1864.~gada likums par gminu pašpārvaldi. Likuma uzdevums bija atbrīvot zemniekus no vietējo muižnieku ietekmes. (Pēc krievu vēsturnieka A.~Pogodina vērtējuma tas gan veda pie pretējā~--- nodibinājās kopējas vietējo iedzīvotāju materiālās un garīgās (nacionālās) intereses. Vietējie muižnieki kļuva tuvāki gminām un saskarsmē ar zemniecību ietekmēja zemnieku pasaules uzskata veidošanos. Zemnieki tāpat kā šļahtiči un garīdznieki sajuta, ka nacionālajai pastāvēšanai draud briesmas un mērķtiecīgi uzsāka tās aizstāvēšanu. Tomēr autoram šķiet, ka izklāstīto vērtējumu būtu grūti pamatot, jo vienas pretrunas~--- feodālās tika nomainītas ar citām~--- kapitālistiskām).

Poļu revolucionārs B.~Švarce, kurš gan tika arestēts un ieslodzīts neilgi pirms sacelšanās sākuma, bet pēc atbrīvošanas pievērsās sacelšanās izpētei un izvērtēšanai, savos darbos enerģiski apgalvoja, ka visas pozitīvās pārmaiņas, kuras notika Polijas karalistē pēc 1863.~gada, bet īpaši agrārā reforma, bija sacelšanās, tās carismam iedvesto baiļu rezultāts. Poļu vēsturnieks, sociologs, politiķis (sociālists) B.~Ļimanovskis uzskatīja, ka carisma realizētā agrārā reforma bija gandrīz vai norakstīta no poļu Centrālās Nacionālā komitejas agrārajiem dekrētiem. Pēc šī darba autora domām, starp carisma 1864.~gada agrāro reformu Polijas karalistē un sacelšanos bija vistiešākais sakars, tā kalpoja sacelšanās apspiešanai. A.~Tarass raksta, ka likteņa ironijas rezultātā poļu šļahta [sacelšanās gaitā] bija nolikusi savas galvas par to, lai uzlabotu stāvokli zemniekiem, kurus tās vairākums dziļi nicināja.

Daudzas citas pozitīvas pārmaiņas lauza sev ceļu, pārvarot carisma pretestību.

Pēc 1863.~gada poļu sacelšanās Krievijas valdība zaudēja cerības mierīgā liberālā ceļā panākt cittautiešu asimilāciju. Krievijas valdošie slāņi secināja, ka piekāpšanās, mēģinājumi pārvaldīt novadu ar liberālām metodēm poļu vidū izsauc tikai jaunas prasības, ved pie sacelšanās. Krievijas valdošajās aprindās radās sava veida ``poļu sindroms''~--- pārliecība par to, ka liberāls kurss ved pie dumpjiem, bet stingrs~--- pie nomierināšanās. Šis secinājums kļuva par pārliecinošu argumentu stingrāka kursa attaisnošanai un pārejai uz aktīvāku Polijas integrāciju impērijas sastāvā. Poļu nacionālisma vilnis izsauca krievu nacionālisma atbildes vilni. Tika cerēts, ka stingra nacionālā politika ļaus nostiprināt valsts vienotību.

Izplatījās panslāvisma ideoloģija, kuras pamatā bija ideja par visu slāvu tautu politisku apvienošanos uz kopīgiem etniskajiem, kultūras un valodas pamatiem, dzimusi XVIII gadsimta beigās~--- XIX gadsimta pirmajā pusē (terminu 1826.~gadā pirmo reizi lietoja slovāku jurists un rakstnieks J.~Herkels, XIX gadsimta 40.~gados to plaši izmantoja vācu un ungāru publicisti lai apzīmētu domājamos Krievijas draudus) pašā Krievijā sākotnēji nespēlēja nekādu lielo lomu, jo krieviem atšķirībā no citām slāvu tautām jau bija sava valsts. Situācija mainījās tikai pēc zaudētā Krimas kara, kad slāvu tautas tika uzlūkotas kā vēlams sabiedrotais cīņā par Krievijas ietekmi. Tomēr Krievijas impērijas ārpolitikā panslāvisma loma nebija necik nozīmīga, kaut tās politiskie pretinieki apgalvoja pretējo. Toties iekšpolitikā krievu augšslāņi panslāvismu centās realizēt, to lielākoties saprotot kā pankrievismu jeb vienkārši krievu nacionālismu, pārejošu šovinismā. Krievu panslāvisti pat mēģināja noliegt poļu piederību slāviem, apgalvot, ka tie piederot Rietumu sabiedrībai, jo nodevuši slāvu ideālus. Ar to klusējot tika atzīts, ka lielkrievu ideāli ir arī visu slāvu ideāli. No tā varēja secināt, ka poļiem ir būtiski jāmainās, lai tos atzītu par piederīgiem slāviem.

Krievijas administrācijas kurss uz nacionālo nomaļu integrāciju kļuva vispārējs un tika forsēts. Tam pievienojās unifikācija valodas jomā rusifikācijas formā. Jaunais kurss iezīmēja mērķi valodas un kultūras ziņā unificēt ne tikai Krievijas Rietumu zemes, bet arī Baltijas guberņas, Somiju, Aizkauāzu. Integrācijas politika izpaudās kā skolu denacionalizācija, preses, grāmatu dzimtajā valodā izdošanas, iestāšanās universitātēs un ģimnāzijās (sevišķi ebrejiem) ierobežojumi. Carisms poļu sacelšanos izmantoja kā ieganstu, lai pastiprinātu nacionālo spiedienu pret citām Krievijas mazākumtautām.

\strong{Cariskās Krievijas rīcība pēc} 1863.--1864.~gada \strong{sacelšanās sakāves} bija vēl krasāka nekā pēc 1830.--1831.~gada sacelšanās. Poļu pretošanās apspiešanas metodes Aleksandra II~--- cara-reformatora laikā izrādījās krietni skarbākas nekā viņa konservatīvā tēva Nikolaja I laikā. Nākamajās desmitgadēs Krievijas politika Polijā stādīja mērķi galīgi atrisināt tās jautājumu ar represiju un vardarbīgas integrācijas palīdzību. Taču tas nenozīmēja, kā to bez pamatojuma apgalvo baltkrievu publicists A.~Tarass, ka ``krievi kā mērķi stādīja poļu kultūras, poļu valodas un paša poļu dumpīgā gara likvidēšanu''. Polijas 1863.~gada sacelšanās apspiešana, poļu bargā sodīšana par nacionālās brīvības centieniem, parādīja pārējām slāvu tautām, ko tās var sagaidīt no Krievijas impērijas pakļaušanās un ko pretošanās gadījumā.

Ievērojamais krievu filozofs V.~Solovjovs pat izteica domu: ``Ja Vīnes kongresā toreizējais patvaldnieciskais imperators Aleksandrs I vairāk domātu par krieviem, nekā par poļu interesēm'', tad ``pamatiedzīvotāju apdzīvoto Poliju būtu atgriezis Prūsijai'', kuras sastāvā tā bija pēc Žečpospolitas dalīšanām. ``Tādā gadījumā,~--- sprieda krievu domātājs,~--- mums, iespējams, nevajadzētu daudz runāt par Poliju'', poļu jautājumu atrisinātu ``neizbēgama ģermanizācija'', no kuras ``Krievija 1814.~gadā izglāba Poliju''. Bijušajā Pozenes lielhercogistē, kuru arī Vīnes kongress izveidoja bijušajās Žečpospolitas zemēs, bet pievienoja Prūsijai, visi runāja vāciski, jo tur ``poļu elements nespēj pretoties vāciešiem un arvien vairāk un vairāk tiek to absorbēts.'' V.~Solovjovam nebija šaubu par to, ``kas būtu noticis, ja Prūsijas vācieši būtu saimnieki Polijas galvenajā daļā''. Tomēr, neraugoties uz savu atzinumu, ka ``Polijas ķermeni saglabāja un uzaudzināja Krievija'', krievu filozofs bija spiests konstatēt, ka ``poļu patrioti drīzāk piekristu noslīkt vācu jūrā, nekā patiesi samierināties ar Krieviju''. No tā izrietēja secinājums: ``Naidam ir dziļāks, garīgs iemesls.'' Tas ir ``mūžīgā strīda starp Rietumiem un Austrumiem izpausme, un Polijas jautājums ir tikai lielā Austrumu jautājuma fāze''. Un, ja tas tā ir, tad ``saprasties ar poļiem nav iespējams ne uz sociāliem, ne valstiskiem pamatiem.'' V.~Solovjova spriedumus XX gadsimtā faktiski apstiprināja poļu vēsturnieks un publicists A.~Bočeņskis, kurš nebūt nebija rusofīls, sapņoja par PSRS sabrukumu. Tūlīt pēc Otrā pasaules kara publicētajā grāmatā ``\pltxti{Dzieje głupoty w Polsce. Pamflety dziejopisarskie}'' (``Muļķības vēsture Polijā. Vēsturiski pamfleti. 1--4, Varšava, 1947.) viņš secināja, ka attiecībā uz poļiem, sākot ar Katrīnu Lielo, ``visi cari sekoja dinastiskās, nevis valstiskās un jo īpaši nacionālās absorbcijas līnijai''. Tomēr šī līnija ``vairakkārt tika pārtraukta, un gandrīz vienmēr to izdarīja Polijas puse'', ko ``vienmēr pavadīja neracionāls naids pret Maskavu, kas tika cītīgi kultivēts un uzpūsts līdz augstākajām robežām''.

Turpmāk Polijas karalistes sabiedriski--politiskajā iekārtā salīdzinājumā ar Krievijas citiem novadiem tika ieviests mazāk pārkārtojumu, kas ļautu piemēroties tirgus ražošanai. Tas nespēja apturēt ātro kapitālistiskās saimniekošanas veida ieviešanu Polijas karalistē, tomēr bremzēja to.

Kamēr Polija nebija pilnībā kļuvusi par Krievijas impērijas sastāvdaļu, vienmēr varēja gaidīt Eiropas valstu iejaukšanos Krievijas attiecībās ar to. Tāpēc likās nepieciešams likvidēt Polijas jautājumu kā starptautisku, kā pamatojošos uz 1815.~gada Vīnes traktātu. Visizdevīgāk to bija izdarīt 1864.--1865.~gadā, kad kļuva skaidrs, ka Polijas dēļ neviens negatavojas karot ar Krieviju, kad Prūsija un Austrija bija gatavas poļu sacelšanās laikā aizmirst Vīnes traktātu. Polijas karaliste tika pilnībā inkorporēta Krievijas sastāvā. Pasta un satiksmes ceļu pārvalde tika pakļauta tieši Pēterburgas ministrijām. Lai pastiprinātu administratīvo varu, bijušo piecu guberņu vietā izveidoja desmit. Pēc F.~Berga nāves 1874.~gadā tika likvidēts vietvalža institūts, kas pēc būtības jau bija zaudējis savu nozīmi. Viņa varas pārmantotājs P.~Kocebu tika nozīmēts par Varšavas ģenerālgubernatoru. Faktiski tika likvidēta Polijas ekonomiskā autonomija. Polijas karaliste zaudēja savu budžetu un patstāvīgu finanšu pārvaldes sistēmu. Valsts padomi un Pārvaldes padomi atlaida. Krievu administrācija Polijas karalistes nosaukumu nomainīja ar Vislas novada (krievu \rutxti{Привислинский край}, poļu \pltxti{Kraj Przywiślański, Kraj Nadwiślański}) vai Vislas guberņu (\rutxti{Привисленские губернии}) nosaukumu. Tas tika lietots jau 70.~gados (ir dati, ka no 1867.~gada), bet 1888.~gadā parādījās arī likumdošanas aktos, saprotot ar to 10 bijušās Polijas karalistes guberņas. Vislas novads bija arī viena no Krievijas ģenerālgubernatūrām. Kaut Polijas karalistes vārds gan turpināja vēl eksistēt, šeit netika ieviestas zemstes, kā tas notika Krievijā. Kaut carisms centās unificēt poļu zemes ar pārējām impērijas teritorijām, šeit netika ievestas Krievijā realizētās liberālās reformas (tiesu, cenzūras u.c.)

Kara tiesu represēto statistiskie aprēķini aptver vairāk nekā 20~tūkstošus sacelšanās dalībniekus un ar tiem sadarbojušos personu. No sodītajiem sacelšanās dalībniekiem 30,85\% bija zemnieki, 17,31\%~--- sīkpilsoņi, 36,42\%~--- tie šļahtiči, kuri bija pierādījuši savu ``augstdzimušo'' statusu, 10,21\%~--- tie šļahtiči, kuri šo statusu bija zaudējuši. Tādejādi gandrīz puse sodīto oficiāli vai neoficiāli piederēja šļahtiču kārtai. ``Sadarbojušos'' vidū šļahtiču bija vairākums. (Padomju vēsturnieks V.~Zaicevs, izanalizējot nepilnīgos Krievijas Kara ministrijas Auditoru departamenta datus par ap 8 tūkstošiem represētajiem sacelšanās dalībniekiem, gan devis nedaudz atšķirīgus skaitļus~--- šļahtiču šo represēto vidū bija līdz 30\%. Tomēr ievērojot, ka Polijas karalistes iedzīvotāju vidū pēc Krievijas senatora, valsts sekretāra Polijas karalistes lietās N.~Miļutina vadībā 1864.~gadā veiktā pētījuma, viņi kopā ar garīdzniekiem sastādīja tikai 5,27\%, arī tad šļahtiču īpatsvars represēto vidū jāvērtē kā visai augsts.) Var secināt, ka tieši viņi bija aktīvākie sacelšanās dalībnieki. Tikai pēc tam aktivitātes ziņā sekoja sīkpilsoņi un zemnieki. Pēc V.~Zaiceva datiem, jo tālāk uz Polijas karalistes austrumiem un dienvid-austrumiem, jo samazinājās nodokļus maksājošo un palielinājās priviliģēto slāņu represēto pārstāvju skaits.

Par dalību 1863.--1864.~gada sacelšanās ar nāvi sodīja ap 400~cilvēkus, ap 18~000 nosūtīja katorgā un izsūtījumā uz Sibīriju. Taču pēc poļu vēsturnieces V.~Sļivovskas savāktajām ziņām viņu dzīve bija tāla no mīta par ``pie ķerrām piekaltajiem katordzniekiem''. Par Usoļjes (\rutxti{Усолье-Сибирскоe}, pilsēta netālu no Irkutskas) katorgu kāds poļu memuārists rakstījis, ka nekad savā dzīvē viņš nav sastapis tik daudz interesantu cilvēku un dzirdējis tādas lekcijas un diskusijas kā šajā vietā. Protams, tā nebija visur. Dažu, īpaši garīdznieku liktenis bija visai smags.

Pati traģiskākā lappuse bija t.s. Aizbaikala sacelšanās (\rutxti{Забайкальское восстание}) 1866.~gadā. To organizēja 1863.~gada sacelšanās dalībnieki, kuri bija izsūtīti uz Sibīriju, kur tai laikā atradās ap 18~tūkstošu notiesāto un izsūtīto poļu revolucionāru. Krasnojarskā un Kanskā viņi kopā ar krievu revolucionāriem gatavoja sacelšanos. Poļu revolucionāri bija nozīmēti trakta (ceļa) celtniecībā no Irkutskas līdz Verhņeudinskai (\rutxti{Верхнеудинск}, tagad Ulanude (\rutxti{Улан-Удэ})). Ap 700~poļu 50 līdz 150 cilvēku grupās strādāja dažādos ceļa celtniecības posmos 200~km garumā. Aresti 1866.~gada sākumā izjauca sākotnējos plānus, sacelšanās ar mērķi izrauties uz Ķīnu sākās pienācīgi nesagatavota. Steidzīgi atsūtītā karaspēka priekšā sacēlušies atkāpās taigā, cenšoties sasniegt Krievijas robežu. Slikti bruņotos sacelšanās dalībniekus sakāva kazaku vienības, pēc divām nedēļām viņi bija spiesti padoties. Sacelšanās vadītājus nošāva Irkutskā 1866.~gada novembrī, simtus tās dalībnieku notiesāja mūža katorgā.

Pamatīgu triecienu saņēma šļahtiču kārta. Tika veikta tās ``tīrīšana'', ap 200~000 šļahtiču zaudēja savu statusu. Polijā muižniecība necieta vēl vairāk tikai tāpēc, ka to nevarēja aizvietot ar krievu vai lietuviešu inteliģenci. Turpretī ārpus Polijas karalistes poļu muižniecību šie represīvie pasākumi gandrīz neskāra, jo tur arī nebija plaša mēroga sacelšanās. Tūkstošus šļahtiču ģimeņu pārcēla uz Krieviju, ar to faktiski iznīka daudzi sīkie poļu ciemi. Pēc dažiem datiem uz Krieviju deportēja ap 70~000 cilvēku, pēc citiem~--- pat vairāk nekā 92~000. Sodītie poļi tika izmantoti mazapdzīvotu teritoriju kolonizācijai. Šim mērķim viņi saņēma aizdevumus un labus zemes gabalus. Izsūtīto vidū bija labi pelnoši ārsti, skolotāji, kantoru darbinieki. Tiesa, poļu pieņemšanu darbā dažādu nozaru iestādēs regulēja nevis likums, bet daudzas slepenas, nekad nepublicētas instrukcijas, cirkulāri, noteikumi, dažkārt arī mutiski izteikti aizliegumi, par kuru saturu var spriest tikai pēc faktiski pastāvošās situācijas.

Izsūtīto poļu vidū bija arī ievērojamais 1863.~gada sacelšanās dalībnieks, tās valdības komisārs Baltkrievijā un Lietuvā ģeogrāfs, zoologs, sabiedriskais darbinieks B.~Dibovskis. Viņam tika piespriests nāves sods pakarot, taču pateicoties vācu zoologu aizstāvībai un O.~f.~Bismarka vidutājībai, sods tika mīkstināts un viņš tika izsūtīts uz 12~gadiem Sibīrijā. Tur viņš pētīja Baikāla ezeru, Amūras upi, bagātināja zinātni ar virkni atklājumu zooloģijā, 1877.~gadā atgriezās dzimtenē, turpināja zinātnisko darbību. Dzīvi B.~Dibovskis beidza jau neatkarīgajā Polijā, būdams pasaulē atzīts zinātnieks, arī PSRS Zinātņu akadēmijas korespondētājloceklis.

Sibīrijas izpētē daudz devuši arī divi citi nedaudz vēlāk izsūtīti poļi: etnogrāfs, literāts, politiķis V.~Seroševskis un etnogrāfs, lingvists E.~Pekarskis. Pirmais par dalību strādnieku kustībā un pretošanos policijai tika izsūtīts uz Jakutiju, kur vāca etnogrāfiskus materiālus. 1896. un 1900.~gadā krievu valodā iznāca viņa darbs par jakutiem. Pēc trimdas viņš turpināja ražīgu zinātnisko, literāro un sabiedrisko darbību, 1933.--1939.~gadā bija Polijas Literatūras akadēmijas prezidents. Otrais 1880.~gadā tika arestēts par piederību sociālistu revolucionāru partijai un nelegālas literatūras glabāšanu notiesāts uz 15~gadiem katorgā. Šeit viņš nodarbojās ar dārzeņu audzēšanu, medībām, zveju. Apguvis jakutu valodu, viņš aizstāvēja vietējo iedzīvotāju intereses presē, rakstīja lūgumrakstus, uzsāka jakutu~--- krievu valodas vārdnīcas sastādīšanu. 1891.~gadā ieguvis tiesības atgriezties Krievijas Eiropas daļa, viņš tās neizmantoja. Tikai XX gadsimta sākumā E.~Pekarskis pārcēlās uz Pēterburgu, sākot tur strādāt Zinātņu akadēmijas Antropoliģijas un etnogrāfijas muzejā. 1907.~gadā iznāca viņa sastādītās jakutu valodas vārdnīcas pirmais laidiens. Līdz 1917.~gada oktobrim tika publicēti pieci vārdnīcas laidieni. 1930.~gadā iznāca tā pēdējais~--- 13.~laidiens. Bez vārdnīcas E.~Pekarskis uzrakstīja arī vairākus etnogrāfiskus darbus. 1927.~gadā viņš kļuva par PSRS Zinātņu akadēmijas korespondētājlocekli, 1931.~gadā~--- tās goda locekli.

Tā kā Baltkrievijā un Lietuvā katoļi nevarēja ieņemt valsts ierēdņu, arī skolotāju, amatus, daudzi no to pretendentiem bija spiesti paši meklēt darbu Iekškrievijas guberņās. Šādu izsūtīto un pārceļotāju pēcteči bija izcilais komponists D.~Šostakovičs, par kura vecvectēva poļu vai baltkrievu saknēm gan var strīdēties, un rakstnieks A.~Grīns. Poļu inteliģence deva lielu ieguldījumu Krievijas nomaļu attīstībā. Kad 1897.~gadā avīzē ``\rutxti{Новоe время}'' (``Jaunais Laiks'') tika publicēti dati par poļu speciālistu koncentrāciju dzelzceļa maģistrāļu celtniecībā impērijas austrumos, par to tika ziņots arī imperatoram. Nikolajs II tikai noplātīja rokas, krievu kadru bija pārāk maz, lai tie varētu nodrošināt milzīgās valsts industriālo attīstību. Piemēram, 1898.~gadā žurnāls ``\rutxti{Русский вестник}'' (``Krievu ziņotājs'') rakstīja par to, ka Permas-Kotlasas dzelzceļa maģistrāles celtniecība pilnībā atrodas poļu rokās, un tajā nodarbinātie krievi sūdzas par spaidiem no poļu puses. Tiesa, Krievijas valdības un sabiedrības zināmas aprindas arī pozitīvi vērtēja poļu darbu impērijas nomalēs. Tā sabiedriskais darbinieks un publicists F.~Umaņecs rakstīja: ``Krievijas pavalstnieki no ārzemniekiem Austrumos ir tādi pat eiropeiski-kristīgās idejas nesēji kā mēs \citespace{}, viņi šai stāvoklī vienmēr meklēs atbalstu tais pat valstiskuma pamatos, uz kuriem balstāmies arī mēs \citespace{} Rezultātā ir jāizdara praktisks secinājums, ka vācieši, zviedri, poļi, norvēģi utt. Usūrijā, pie Amūras, Turkestānā utt. ir tāds pat vēlams kolonizācijas elements, kā dzimušie Tveras, Maskavas vai Kurskas guberņās.''

Represijām pakļāva Polijas katoļu garīdzniecību. Šādam pavērsienam tā bija devusi iemeslu sacelšanās laikā, kad daudzi garīdznieki to atbalstīja. Laikabiedrs prūšu majors E.~Knorrs aprakstīja baznīcas kalpotāju izturēšanos sacelšanās laikā: ``Gandrīz visur klosteri un baznīcas kalpo kā ieroču un pārtikas noliktavas, kā slēptuves atsevišķiem insurgentiem [sacēlušamies] un veselām bandām''. Tāpēc ar 1864.~gada 8.~novembra dekrētu tika noteikta baznīcas īpašumu konfiskācija un klosteru, kā arī Varšavas garīgās akadēmijas slēgšana. No 197 klosteriem ar 2~187 garīgo ordeņu brāļiem turpināja darboties 35 klosteri ar 470~ordeņbrāļiem. Tiem bija aizliegts uzņemt jaunus brāļus. 1867.~gadā poļu katoļu baznīca tika pakļauta Romas katoļu kolēģijas Pēterburgā uzraudzībai. Virkne bīskapu, tostarp arī Varšavas, tika izsūtīti uz Krieviju. Romas pāvests ar to visu nevarēja samierināties, ar Krieviju tika pārtrauktas attiecības un beidza darboties 1847.~gada konkordāts. 1870.~gadā no 15 diocēzēm 12 nebija bīskapu. 1875.~gadā Krievijas pareizticīgās baznīcas Sinode (\rutxti{Святейший Правительствующий Синод}) nolēma likvidēt ūniju. 1879.~gadā tika oficiāli likvidēta uniātu baznīca, piespiedu kārtā apvienojot to ar pareizticīgo. Taču baznīca turpināja darboties pagrīdē. Tikai 80.~gadu sākumā notika daļējs Vatikāna un Krievijas izlīgums. Iedzīvotāji pasākumus pret katoļu garīdzniecību uztvēra kā vērstus pret katoļu reliģiju.

Lietvedība administratīvajās un tiesu iestādēs ``krievu'' Polijā kopš 1868.~gada oficiāli norisa krievu valodā. No 1875.~gada tā kļuva par obligātu tiesās. Varšavas Galveno skolu 1869.~gadā pārvērta par universitāti, kur mācības notika krievu valodā un darbojās Viskrievijas universitāšu statūti, kas ievērojami ierobežoja tās autonomiju. No 1872.~gada mācības visās vidējās izglītības iestādēs notika krievu valodā, nākamajā gadā sākās sākumskolu pāreja uz krievu valodu. (Varšavas apgabala un 9 Krievijas rietumu guberņu augstākajās un vidējās mācību iestādēs gan poļu valodu varēja pasniegt kā mācību priekšmetu.) Ar 1885.~gadu sākumskolas tika pakļautas pilnīgai rusifikācijai. Tikai ticības mācību katoļiem varēja pasniegt poļu valodā. Poļu zemnieki izrādīja pasīvu pretošanos skolu rusifikācijai, centās izvairīties no nodevām sākumskolu uzturēšanai. Arī poļu virsslāņi bija neapmierināti ar rusifikāciju, vienmēr uzsvēra carisma atbildību par kultūras attīstības bremzēšanu poļu zemēs, taču pēc oficiālajiem datiem paši maz ko darīja tautas izglītības līmeņa celšanai. Polijas karalistes lauku skolu budžetā privātziedojumu bija tikai 2,8~procenti, tas bija vismazākais īpatsvars starp visām Krievijas impērijas guberņām un apgabaliem. Diemžēl autora rīcībā nav datu, cik lielu atbalstu zemnieki un muižnieki sniedza slepenai poļu valodas mācīšanai, kura bija izplatīta ciematos.

Rusifikācijas politika augstāko punktu sasniedza 1883.--1894.~gados, kad Varšavas ģenerālgubernatora pienākumus pildīja ģenerālis J.~Gurko. Rusifikācijas izglītības jomā faktiski bija stiprākais līdzeklis kā modināt poļu zemnieku masu nacionālo pašapziņu, kad lielai daļai inteliģences nebija vairs spēka protestēt un bērnus tā sūtīja krievu skolās. Rusifikācijas rezultātā bija vērojama skolu sniegtā izglītības līmeņa pazemināšanās (skolas beidzēji nebija ieguvuši izglītību dzimtajā valodā un nebija apguvuši arī krievu valodu) un skolu tīkla sašaurināšanās. Ik gadu pieauga skolas neapmeklējušo bērnu skaits, kā rezultātā auga arī lasīt un rakstīt nepratēju skaits. Piemēram, Varšavā 1882.~gadā tādu bija 48,6\%, bet 1897.~gadā~--- 46,5\% no visiem pilsētas iedzīvotājiem.

Krievijas valdītājiem bija skaidrs, ka savu kundzību, novada pārvaldi nevar bāzēt uz poļu ierēdniecību, liela daļa no kuras sacelšanās laikā izrādījās nelojāla. To bija jānomaina ar krievu ierēdniecību. Krievijas pareizticīgās baznīcas Sinodes virsprokurors K.~Pobedonoscevs 1983.~gadā vēstulē imperatoram Aleksandram III rakstīja; ``nekādā gadījumā nedrīkst ar uzticību paļauties uz katoļiem-poļiem\dots{}'' Tā kā pat daudzi cariskās administrācijas ierēdņi krievi nevēlējās piedalīties patvaldības politikas realizācijā Vislas novadā, valdība radīja virkni privilēģiju tiem ierēdņiem, kuri dienēja Polijā. Tāpēc šurp devās dažādi karjeristi, pašlabuma meklētāji. Ierēdņu un vietējo iedzīvotāju attiecības pasliktinājās. Deklasētajiem poļu šļahtičiem bija slēgts ceļš uz valsts, sabiedriskā, kultūras darba virsotnēm, viņi varēja vienīgi iestāties darbā valsts iestādēs zemākajos amatos bieži pārlieku pašpārliecinātu un aprobežotu krievu ierēdņu uzraudzībā. Šādi šļahtiči saglabāja romantisku jūsmu par Polijas pagātni un naidu pret carisko patvaldību. Diemžēl naidu pret carismu viņi bieži pārnesa arī uz visu krievu tautu.

Šļahtiču nesamierināmību ar carismu uzskatāmi demonstrēja nabadzīgs šļahtičs A.~Berezovskis, mūzikas skolotāja dēls, kurš pretēji tēva gribai bija piedalījies 1863.~gada sacelšanās, pēc tās apspiešanas emigrējis un dzīvoja Parīzē, strādājot atslēdznieka darbnīcā. 1867.~gadā Parīzē notika Vispasaules izstāde, kuru apmeklēja arī imperators Aleksandrs II. 6.~jūnijā Buloņas mežā A.~Berezovskis šāva uz Alekasndru II, kurš brauca ekipāžā ar diviem saviem dēliem un Francijas imperatoru Napoleonu III. Pistoles lādiņš bija pārāk stiprs, tā uzsprāga, ievainojot rokā pašu A.~Berezovski, bet lode trāpīja kāda imperatora pavadoņa zirgu. Pūlis tūlīt sagrāba šāvēju. Viņam piesprieda mūža katorgu, bet 1906.~gadā amnestēja.

Taču, kā norādījis latviešu vēsturnieks Aldis Miņins, impērijas attieksme pret poļiem nebija totāli represīva. Pat pēc 1863.~gada poļu sacelšanās apspiešanas 6\% impērijas augstākajā ierēdniecībā veidoja poļu muižnieki, bet Polijas karalistes pārvaldē etniskie poļi vienmēr veidoja vairākumu, attiecīgi XIX gadsimta 60.~gados~--- 80\%, bet 90.~gados~--- nedaudz virs 50\% no visiem ierēdņiem. Tiesa, poļu ierēdniecības procentuālā īpatsvara kritums Polijas karalistē pēc sacelšanās apspiešanas bija iespaidīgs~--- gandrīz 30\%.

Sakāves sajūta un represijas, kuras Krievijas vara vērsa kā pret sacelšanās dalībniekiem, tā arī visu poļu sabiedrību, tikai padziļināja un pastiprināja poļos pretkrievisko noskaņojumu. Vairākuma prātos nostiprinājās pārliecība, ka Krievija ir bijusi un būs mūžīgi naidīga Polijai. Tautas vairākumam pēc 1863.~gada sacelšanās sakāves spēku gan pietika tikai pasīvai pretošanās demonstrēšanai.

Polijas karalistē un blakus esošajās Krievijas rietumu guberņās izplatījās lapiņas ar aicinājumiem cīnīties pret nacionālo apspiestību, draudu vēstules cara administratoriem. Bieži bija ``viņa majestātes'' (imperatora) apvainošanas gadījumi, aizlūgumi baznīcās atbrīvošanās cīņu dalībnieku piemiņai, to portretu izplatīšana, demonstratīva sēru drānu un atceres zīmju nēsāšana datumos, saistītos ar 1863.~gada sacelšanos. Manifestācijās iesaistījās diezgan daudz patriotisko spēku pārstāvju. Pilsētu un lauku apakšslāņu pārstāvji centās izvairīties no dienesta armijā. Tikai 1867.--1870.~gadā Varšavas apgabaltiesa piesprieda sodus 1504~poļu rekrūšiem par atteikšanos nodot militāro zvērestu vai bēgšanu uz ārzemēm.

Veidojās slepeni pulciņi. Tajos tika izstrādāti jauni bruņotas cīņas plāni. Nacionālās cīņas aktīvisti īpaši aktivizējās saasinoties starptautiskajai situācijai un kaut gan Francijas sakāve 1870.~gada karā pret Vāciju (padomju literatūrā to parasti sauca par franču-prūšu karu, kaut tajā piedalījās arī citas vācu valstis) nobīdīja poļu jautājumu no starptautiskās dienas kārtības, poļi tomēr loloja cerības uz tā atkārtotu izvirzīšanos.

60.~gados arī Krievijas pilsētās pastāvēja daudz dažādu virzienu nelegālu poļu pulciņu. Tie uzturēja sakarus ar poļu izsūtītajiem, emigrantiem ārzemēs. Vairāki poļu revolucionāri aktīvi piedalījās krievu narodņiku teroristiskajā darbībā. Tā tieši poļa I.~Griņevicka mestā bumba ņēma dzīvību gan imperatoram Aleksandram II, gan viņam pašam. L.~Mirskis 1879.~gadā neveiksmīgi šāva uz Krievijas žandarmērijas šefu A.~Drentelnu. L.~Kobiļanskis piedalījās nāvējošā atentātā pret Harkovas gubernatoru kņazu D.~Kropotkinu. Vēl daudzi citi poļi tika tiesāti par dalību terorismā, kā arī revolucionārā propagandā.

Toties poļiem juta līdzi ``visa'' ``progresīvā'' Eiropas sabiedrība. Tās acīs poļu sacelšanos apspiešana veidoja Krievijas kā despotiskas monarhijas tēlu, kura apspiež visas brīvības izpausmes, bremzē progresu, simbolizē reakciju un atpalicību. Kā norādījis poļu--krievu attiecību pētnieks E.~Gorizontovs, ``Poļu rusofobija būtiski ietekmēja Krievijas tēlu Rietumu acīs un atstāja iespaidu arī uz ukraiņu un baltkrievu kustību ideoloģiju.'' Turpretī par poļu tautu Eiropā nostiprinājās priekštats, ka tai ir īpaša loma cīņā ``par mūsu un jūsu brīvību''. 1867.~gadā K.~Markss rakstīja: ``\dots{}Eiropai ir atlicis izvēlēties vienu no diviem. Vai nu aziātiskais barbarisms ar moskaļiem priekšgalā kā lavīna nobruks pār to, vai tai nāksies atjaunot Poliju, ar to norobežojoties no Āzijas ar divdesmit miljoniem varoņu un izmantojot brīvās minūtes jaunu sabiedrisku pārveidojumu veikšanai''. Tiesa, Eiropas sabiedrisko un nacionālo atbrīvošanās kustību atbalsts nekādi nepalīdzēja Polijai, taču atstāja savu iespaidu tautas politiskajā apziņā.

Dažus gadus vēlāk būtiski mainījās situācija Eiropā, radās spēcīga, apvienota Vācijas impērija. Jau 70.~gados ``triju imperatoru savienība'' (tā dēvēja Krievijas, Vācijas un Austroungārijas imperatoru savienību, nostiprinātu ar 1873., 1881., un 1884.~gada līgumiem) apstiprināja Eiropas politiskās kartes nemainību un Polijas jautājums pārstāja būt par starptautisku un kļuva par Krievijas, Vācijas un Austroungārijas savstarpējo attiecību jautājumu. Poļiem zuda cerības bez nopietna starptautiska konflikta sekmīgi cīnīties pret apspiedējvalstīm Krieviju, Vāciju un Austroungāriju.

Pēc 1863.~gada sacelšanās sakāves veco emigrāciju papildināja jauni poļu pulki. Atšķirībā no vecās, galvenokārt šļahtiču veidotās, emigrācijas, šī jaunā emigrācija pārsvarā sastāvēja no inteliģences un amatniekiem. Taču tā vairs nebija vecā, ``lielā'' emigrācija. Ja pēc 1831.~gada sacelšanās emigrantu vidū pārsvarā valdīja optimisms un cīņas spars, ar viņiem sakaros stājās Rietumu valdības, tad pēc 1864.~gada, kā atzinis L.~Vasiļevskis, emigrantu vidū valdīja bezcerība, starp viņiem nebija Rietumos atzītu izcilu personību. Rietumvalstu valdības pārtrauca sakarus ar viņiem. Jaunie emigranti pēc sociālā sastāva un uzskatiem bija ievērojami demokrātiskāki nekā iepriekšējā viņu paaudze. Tas izsauca arvien dziļākas pretrunas starp poļu emigrantu grupām jautājumā par Polijas neatkarības atjaunošanas ceļiem.

Eiropas valstīs strādnieku vidū pastāvēja solidaritātes kustība ar 1863--1864.~gada poļu sacelšanos. 1863.~gada 23.~jūlijā pēc solidaritātes mītiņa Londonā tikās angļu, franču un poļu strādnieku delegācijas. Tika nolemts sagatavot uzsaukumu franču strādniecībai ar aicinājumiem radīt starptautisku strādnieku organizāciju un ``uzstāties kopīgi Polijas atbrīvošanas labā''. Kopīgajai poļu tautas aizstāvībai bija jākļūst par pirmo starptautiskās strādnieku savienības rīcības aktu. Franču strādnieku pozitīvā atbilde tika nolasīta 1864.~gada 28.~septembrī Londonā lielā mītiņā, kurā arī tika pasludināta Starptautiskās Strādnieku asociācijas (vēlāk sauktas par I Internacionāli, 1864--1872) dibināšana. Tāpēc arī K.~Markss rakstīja, ka 1863.~gada poļu sacelšanās, kura izsauca angļu un franču strādnieku protestu pret savu valdību starptautiskajām ļaundarībām, kalpoja par Internacionāles sākumpunktu.

Vēlāk I Internacionāles vadītāji pastāvīgi uzsvēra, ka Polijas atbrīvošanās jautājums ir viens no svarīgākajiem straptautiskajam proletariātam, ka Polijas atbrīvošana nozīmētu visu valstu, bet pirmkārt jau Vācijas un Krievijas, strādnieku atbrīvošanu, ka pats Polijas stāvoklis revolucionarizē to. Jautājums par ``Maskavas draudiem Eiropai un neatkarīgas un vienotas Polijas atjaunošanu'' tika iekļauts Starptautiskās Strādnieku asociācijas pirmās konferences darba kārtībā. Kaut K.~Markss un F.~Engelss vienmēr uzsvēra, ka tieši poļu valdošo slāņu egoisms ir pazudinājis Poliju, 1865.~gadā K.~Markss tomēr izteica cerību, ka pēdējās sacelšanās mācības šļahtai nav bijušas gluži veltas. (Jāsaka gan, ka no revolucionārās strādniecības viedokļa šāds tās politisko vadītāju revolucionārs romantisms kopumā neattaisnojās.)

1865.~gadā Londonā tika izveidota I Internacionāles Poļu sekcija, kuras sastāvā ietilpa vairāki desmiti emigrantu. Sekcijā iestājās gan Londonas poļu radikālo organizāciju, kam bija sociālistiska ievirze, dalībnieki, gan ``sarkano'' kreisā spārna pārstāvji~--- bijušie 1863.--1864.~gada sacelšanās dalībnieki, nonākušie emigrācijā. Vairāki poļi tika ievēlēti Internacionāles Ģenerālajā padomē. Poļu revolucionāri saņēma arī materiālu palīdzību no Internacionāles. Policijas ziņojumos bija pausta trauksme par Internacionāles ideju iespiešanos Polijā. Taču teiktais nenozīmēja, ka visi poļu emigranti pieņēma marksisma ideoloģiju. Daudzus piesaistīja nevis sociālistiskie ideāli, bet Internacionāles pozīcija Polijas jautājumā. Daļa poļu demokrātu, arī ``sarkano'' mērenā spārna pārstāvji pret Internacionāli izturējās noraidoši, pat naidīgi. Kopumā visas poļu emigrācijas rindās I Internacionāles piekritēju grupa veidoja nenozīmīgu mazākumu.

Daudzi emigrācijā nonākušie poļu revolucionāri atbalstīja nacionālās un sociālās atbrīvošanās cīņu citās Eiropas valstīs. Viņi piedalījās itāļu tautas cīņā pret austriešu apspiedējiem. Vācijas un Francijas kara (1870.--1871.) otrajā posmā simti poļu revolucionāru un tai skaitā arī J.~Dombrovskis cīnījās Francijas Republikas pusē, cerot, ka tās uzvara atvieglos Polijas ``atjaunošanu''. Vēlāk ap pieci simti poļu sociālistu un demokrātu darbojās Parīzes komūnā. Starp augstākajiem Komūnas militārajiem komandieriem apmēram 20 bija poļi. Kā atzina F.~Engelss, poļi Parīzes komūnas aizstāvju rindās bija vienīgie uzticamie un spējīgie karavadoņi. J.~Dombrovskis tika ievainots un no ievainojuma nomira. Versaļas valdības izveidotās kara tiesas gūstā kritušos poļus parasti lika nošaut. Piemēram, nāves spriedums neklātienē tika piespriests vienam no Komūnas spēku vadītājiem V.~Vrubļevskim. Viņš nonāca Londonā, izveidoja tur draudzīgas attiecības ar K.~Marksu un F.~Engelsu, pēc viņu priekšlikuma kļuva par I Internacionāles Ģenerālpadomes locekli, līdz mūža beigām turpināja simpatizēt sociālismam. 1875.~gada decembrī starptautiskā strādnieku mītiņā Londonā V.~Vrubļevskis paziņoja, ka īstais Polijas karogs ir sociāldemokrātiskais karogs, zem kura jāvienojas visiem poļu tautas progresīvajiem spēkiem cīņā par dzimtenes neatkarību un sociālo progresu. Mērķa sasniegšanas iespēju viņš saskatīja vienīgi nodibinot ciešu poļu tautas savienību ar rietumvalstu un īpaši Krievijas darbaļaudīm. ``Mums ir jāgatavojas uz sacelšanos,~--- teica V.~Vrubļevskis,~--- kopā ar krievu sociālistiem, tāpēc poļiem ir jāpalīdz krieviem Krievijā, tāpat kā krieviem jāpalīdz poļiem Polijā.'' K.~Markss augstu vērtēja poļu starptautisko darbību: ``Polija lēja savas asinis karā par Amerikas neatkarību, viņas leģioni cīnījās zem pirmās Francijas Republikas karogiem; 1830.~gadā tā ar savu revolūciju novērsa iebrukumu Francijā, kuru toreiz nolēma veikt Polijas dalīšanas dalībnieces: 1846.~gadā Krakovā Polija pirmā Eiropā pacēla sociālās revolūcijas karogu; 1848.~gadā tās dēli ieņēma izcilu vietu revolucionārajās cīņās Ungārijā, Vācijā un Itālijā; beidzot, 1871.~gadā tā deva Parīzes komūnai labākos ģenerāļus un pašus varonīgākos kareivjus''. Ievērojot Polijas kā nemiera un nestabilitātes perēkļa lomu Eiropā, K.~Marksa revolucionāro jūsmu var saprast. Taču, kā atzīmējusi krievu vēsturniece S.~Falkoviča, piedaloties notikumos Eiropā, cīnoties par citu tautu interesēm, poļi gandrīz vienmēr uzstājās arī par Poliju, meklēja kādu labumu poļu nacionālajai lietai. Tā, krievu-turku karā (1877--1878) pret Krieviju karoja arī poļi, kopā 145 cilvēki. Tomēr neizdevās turku mēģinājums, solot naudu, izraisīt poļu sacelšanos.

Pēc 1863.--1864.~gada sacelšanās Polijas karalistē \strong{īpaši cieta apkārtējie apgabali}. Jau 1863.~gada maijā par Krievijas Ziemeļrietumu novada (Lietuvas un Baltkrievijas) ģenerālgubernatoru tika nozīmēts ģenerālis M.~Muravjovs. (Revolucionāru vidū viņš tika iesaukts par ``Muravjovu--kārēju'', jo, kad 1831.~gadā M.~Muravjovu iecēla par Grodņas gubernatoru, viņam pajautāja, vai viņš nav pakārtā dekabrista S.~Muravjova-Apostola radinieks. M.~Muravjovs uz to atbildēja, ka ir cēlies nevis no tiem Muravjoviem, kurus kar, bet no tiem Muravjoviem, kuri paši kar.) Krievijas valdošajās aprindās pastāvēja gan birokrātu grupa naidīga Ziemeļrietumu novada ģenerālgubernatoram M.~Muravjovam. Tā, iekšlietu ministrs grāfs P.~Valujevs uzskatīja, ka ģenerālgubernators M.~Muravjovs vēlas pilnībā atstumt centrālās varas iestādes no Ziemeļrietumu novada pārvaldes. Tomēr galu galā arī viņš atzina vietvalža darbības rezultātus: ``Viņš rīkojas kā diktators, bet rīkojas!'' M.~Muravjovs uzskatīja, ka 1863.~gada sacelšanās gaitai bija jāpārliecina valdību, ka līdz tam realizētā piekāpšanās politika, jebkuras reformas poļu muižniecība uztver kā vājuma apliecinājumu un tikai pastiprina tās dumpīguma garu. Viņš rakstīja, ka poļu muižnieki ``\dots{}ar ``tēvzemes atjaunošanas'' ieganstu vēlas tikai \citespace{} iegūt agrākās patvarīgās tiesības pār zemniekiem un pārējiem nodokļu maksātājiem, kādas viņi neierobežoti izmantoja Žečpospolitas laikā''. M.~Muravjova valdīšanas laikā Krievijas Ziemeļrietumu guberņās represijām tika pakļauti ap 9~000 cilvēku, tai skaitā ar nāvi sodīti 128, katorgā nosūtīti ap tūkstoti, pārējie nodoti karadienestā vai arestantu rotās.

Krievijas kara ministrijas augsts ierēdnis I.~Jākobsons 1863.~gada 12.~jūnijā savā dienagrāmatā atstāja ierakstu par M.~Muravjova rīcību; ``Viņš pakāra Viļņā kādu ksendzu un insurgentu [sacelšanās dalībnieku] un atstāja viņus uz karātavām apgānīšanai veselu dienu. Kad viņu līķus apraka zemē, saradās daudz cilvēku, lai paklanītos viņu kapavietām, tika turēti aizlūgumi, tika ņemta zeme no kapiem par piemiņu, un tie tika dažādi godināti. Policijmeistars par to ziņoja Muravjovam un lūdza norīkot karaspēku, lai izdzenātu cilvēkus. Muravjovs tam nepiekrita, bet lika apliet kapus ar ekskrementiem un šai nolūkā atvest vairākus vezumus no slimnīcu atejām. Pavēle tika izpildīta, un tik tiešām vairāk neviens netuvojās kapiem. Taču kāds iespaids tika radīts fanātiskajiem katoļiem?''

Tautā plaši bija izplatīts M.~Muravjova teiciens: ``Ko nepabeidza krievu durklis, to pabeigs krievu skola''. Atbilstoši tam notika skolu rusifikācija. Skolās aizliedza poļu valodu, 1865.~gadā krievu valodu ieviesa kā vienīgo pārvaldes valodu. Par augstākajiem ierēdņiem nozīmēja tikai krievus. No ``nomierināšanas'' viedokļa M.~Muravjovs darbojās tik efektīvi, ka pēc amata atstāšanas saņēma grāfa titulu ar piebildi ``Viļņas'' (``\rutxti{Виленский''}) un Krievijas impērijas augstāko~--- Aleksandra Pirmdzimtā (\rutxti{Александра Первозванного}) ordeni. Aleksandrs II vēlāk izteicās par M.~Muravjovu, ka viņš esot bijis vienīgais ierēdnis, kurš pratis poļus ``turēt rokās''. Arī mūsdienu Krievijas historiogrāfijā tiek uzsvērts, ka M.~Muravjovs lielas pūles veltīja baltkrievu zemnieku emancipācijai un viņu atkarības no poļu muižniecības vājināšanai un viņa darbība pierādīja, ka krievu skola, krievu garīdznieki un ierēdņi izrādījās efektīvāki nekā krievu durkļi. Tiek pausts arī viedoklis, ka tieši viņš kā Krievijas impērijas augstākā amatpersona ielika pamatus mūsdienu baltkrievu nācijai.

Poļu muižnieku kārta Krievijas Ziemeļrietumu novadā pēc 1863.~gada sacelšanās sakāves piedzīvoja bēdu dienas. Kā norādījis poļu vēsturnieks V.~Smoļenskis, viena no galvenajām M.~Muravjova aktivitātēm bija sava veida kontribūcijas uzlikšana 10\% apmērā no ienākumiem visiem poļu zemes īpašumiem. Kaut arī sacelšanās beidzās, bet sods par to poļiem turpinājās. Tiesa, 1869.~gadā dekrēts samazināja M.~Muravjova noteiktās iemaksas, bet aizstāja to ar pastāvīgu nodokli ar 5\% nodevas nosaukumu no poļu izcelsmes personām. Šis nodoklis, kas lietuviešu apdzīvotajās guberņās ik gadu ienesa valsts kasē apmēram 2 miljonus rubļu, pakāpeniski izputināja poļu zemes īpašniekus.

Ja Polijas karalistē poļiem tika konfiscētas 1~660 muižas, tad Krievijas rietumu rajonos pat vairāk~--- 1~760. Dzimtbūšana, muižnieku privilēģijas bija atceltas. Daudzi muižnieki bija aizbēguši vai izsūtīti, saimniekot vajadzēja viņu sievām, kuras to neprata, vēl jo vairāk tāpēc, ka tagad vajadzēja izmantot algotu darbaspēku. Situāciju izmantoja bagātākie ebreji, kuri piedāvājās muižu saimniecēm, kā arī palikušajiem, bet saimniekot nespējīgajiem poļu šļahtičiem iznomāt viņu īpašumu. Īpašnieki saņēma pieklājīgu nomas maksu un tiem nebija jārūpējas par saimniecību. Protams, arī muižu apsaimniekotāji ebreji nepalika zaudētājos.

XIX gadsimta otrajā pusē blakus Polijas karalistei esošajā Grodņas guberņas Kameņecas (\rutxti{Каменец-Литовск}) pilsētiņā dzīvojušais ebreju tirgotājs E.~Kotiks savās atmiņās krāsaini aprakstījis vietējo iedzīvotāju, īpaši ebreju, attiecības ar poļu muižniekiem, kuras līdzīgas bija arī Polijas karalistē.

Pēc E.~Kotika liecības ap 65\% Grodņas guberņas muižu apmetās saimniekot ebreji. Muižnieki ar viņu saimniekošanu bija apmierināti. Atgriežoties mājup pēc izsūtījuma, muižnieki nepazinuši savas labiekārtotās saimniecības. Daži mēģinājuši turpināt saimniekot paši, bez ebreju nomniekiem, tomēr pēc diviem-trim gadiem muižas atkal panīkušas un muižnieki atkal tās nodevuši nomā ebrejiem. Taču, ja sākumā muižnieki savas muižas iznomāja par zemu cenu, tad turpmāk viņi to pamatīgi pacēla un pēc gadiem desmit muižu apsaimniekošana vairs ebrejiem nebijusi ``nemaz tik izdevīga''. (Tiesa, atmiņu autors nestāsta, kā ebreju saimniekošanas veidu vērtēja vietējais darbaspēks~--- zemnieki-baltkrievi. Acīmredzot viņos auga nepatika ne tikai pret poļu muižniekiem, bet arī nomniekiem--ebrejiem, kuri saimniekoja jau ar kapitālistiskām metodēm.) Ap pilsētu dzīvoja pāris simtu muižnieku, kam katram piederēja ap divi simti vai vairāk dzimtcilvēku. Katram muižniekam miestā bija viens vai divi uzticami ebreji-veikalnieki, ar kuriem viņš kārtoja savus darījumus. Ja ap muižnieku ``grozījās'' divi ebreji, tad viens no tiem parasti bija cienījams tirgotājs, bet otrs ``sīkāks'' kā tirdzniecības, tā tikumu ziņā. Pirmais tika izmantots galvenokārt kā padomdevējs, otrs~--- kā īpašu uzdevumu veicējs. Tomēr abi vienādi baidījās no muižnieka, kaut bieži dzīvoja uz tā rēķina. Pēc pirmās iegribas muižnieks varēja arī nopērt ``savu'' ebreju, pie tam vēl piebilstot: ``Turēsi muti~--- paliksi man kalpot, ja nē~--- paņemšu citu ebreju, bet tu man tik un tā neko nepadarīsi, jo gan policijas iecirkņa uzraugs, gan apriņķa policijas priekšnieks ir manās rokās''. Kopumā tā arī bija. Ebreji klusēja, domājot: ``Lai arī muižnieks mani var nopērt, taču es pie viņa nopelnīšu savu maizes gabalu. Arī kad nomiršu, mans dēls no viņa gūs peļņu''. Ja ebrejs nomira, muižnieks parasti pieņēma darbā sev bijušā izpalīga dēlu vai citu radinieku. Bez šī ``tirdzniecības lietu kārtotāja'' katram muižniekam miestā bija arī savi amatnieki: kurpnieks, drēbnieks, skārdnieks u.c. Tikai šim savam amatniekam tika piešķirti muižnieka pasūtījumi. Amatniekiem nopelnīt sev iztiku bija daudz grūtāk nekā tirgotājiem. Muižnieka rīcībā varēja būt arī savs ebrejs, saukts par faktoru, kurš par komisijas naudu kārtoja viņa darījumus, dzīvojot nevis pilsētā, bet muižā, kā arī ebrejs-nomnieks, kurš pārvaldīja viņam iznomāto to vai citu muižu. Ja muižniekam piederēja vairākas muižas, tad visās tajās nereti dzīvoja šādi faktori un nomnieki. Protams, ebreji augsti vērtēja savu pietuvinātību muižniekam, nevēlējās to zaudēt. Visu, kas muižniekam bija nepieciešams, viņš sagādāja ar sava ebreja starpniecību, uzskatot to par gudru, viltīgu, tomēr godīgu radījumu, jo ``savs, paša'' ebrejs taču neuzdrošināsies apkrāpt savu labvēli! Visus pārējos ebrejus muižnieki uzskatīja par negodīgiem. Ja pilsētas ebrejs atbrauca pie muižnieka, viņš vispirms gaidīja, kad kāds no vietējiem zemniekiem aizvedīs viņu pie faktora, jo pieeju muižas ēkai sargāja suņu bars. Tikai pēc tam kāds no mājas kalpotājiem atbraukušo pavadīja pie muižnieka. Ja muižnieks bija labvēlīgs, atpakaļ līdz vārtiem ebreju pavadīja sulainis, ja nē~--- viņam nācās iet vienam, mēģinot atkauties no suņiem. Kā rakstīja E.~Kotiks, to mocību apraksts, ko šādā gadījumā cieta ebreji no muižnieku suņiem, varētu saturēt daudz lappušu.

Līdz ar poļiem nācās ciest arī to sacelšanos neatbalstījušajiem citu tautu pārstāvjiem, arī latgaliešiem Vitebskas guberņā. Pret polonizāciju vērstais latīņu drukas aizliegums tika attiecināts arī uz Latgali, skarot arī latviešu valodā izdodamos darbus. 1864.~gadā izsludinātais un līdz 1904.~gadam spēkā esošais latīņu iespieddarbu aizliegums liedza iespēju Latgalē drukāt un izplatīt latviešu grāmatas. Latgaliešiem, tāpat kā poļiem, bija aizliegts iegādāties nekustāmo īpašumu bez gubernatora atļaujas utt.

Tiesa, attīstoties lietuviešu, baltkrievu un ukraiņu nacionālajām kustībām, kas pastiprināja tradicionālās reģionālās atšķirības, daļa Krievijas rietumu guberņu poļu izcelsmes šļahtas zaudēja savu poļu identitāti. Literatūrā var atrast piemērus, kur tuvi radinieki, dzimuši Krievijas Ziemeļrietumu novadā, bet kā personības veidojušies dažādos apstākļos: Viļņā, Varšavā un Pēterburgā, apzinājās sevi gan kā lietuvieši, gan poļi, gan baltkrievi.

Ukrainā jeb Krievijas Dienvidrietumu novadā (Kijevas, Podolijas un Volīnijas guberņās), kur arī dzīvoja poļi, stāvoklis tiem bija ievērojami labvēlīgāks, jo šeit kā sacelšanās kustībai bija mazāks vēriens, tā arī kā ģenerālgubernators saimniekoja daudz mērenākais ģenerālis N.~Anņenkovs.

\strong{XIX gs. pēdējā trešdaļā saimniecība Polijas karalistē} attīstījās samērā sekmīgi. 1807.~gadā pasludinātā dzimtcilvēku personīgā atbrīvošana, plašu mazzemnieku un bezzemnieku masu pastāvēšana, daļas muižnieku--šļahtiču tendence jau gadsimta pirmajā pusē pielietot kapitālistiskas saimniekošanas formas, muižu pastāvēšana, kurās jau tika pārsvarā pielietots algots darbs, radīja labvēlīgus priekšnoteikumus 1864.~gada agrārās reformas norisei.

Pirmie gadu desmiti pēc reformas sakrita ar samērā augstu labības cenu periodu pasaules tirgū. Auga arī labības pieprasījums iekšējā tirgū. Šie abi apstākļi noteica muižnieku saimniecību attīstības profilu. Tās turpināja galvenokārt ražot graudus. Muižnieku saimniecības uz kapitālistiskās saimniekošanas ceļu pārgāja pakāpeniski, bez straujām pārmaiņām ražošanas organizācijā un agrotehnikā.

Poļu slānis, kurš guva vislielāko labumu no Krievijas impērijas politikas pēc 1863.~gada sacelšanās apspiešanas, bija zemnieki. Pavisam ap 2 miljoni zemnieku (ar ģimenes locekļiem) ieguva zemes no muižnieku īpašumiem, vēl 110~000~--- no valsts un sekularizētajām zemēm. Zemnieku saimniecību skaits no ap 0,5~miljoniem 1864.~gadā pieauga līdz 0,6~miljoniem 1873.~gadā un līdz 0,78~miljoniem 1904.~gadā. Šai ziņā Polijā situācija bija daudz labāka nekā pārējā Krievijā pēc 1861.~gada reformas. Zemnieku zeme 1872.~gadā sasniedza 46\% no lauksaimniecībā izmantojamās zemes, tās īpatsvars auga arī turpmāk. Ar 1888.~gadu Zemnieku banka (\rutxti{Крестьянский поземельный банк}) izplatīja savu darbību arī Polijas karalistē, izsniedzot zemniekiem hipotekārus kredītus. Pēc L.~Vasiļevska datiem 42\% no visas zemnieku iepirktās zemes ieguva bezzemnieki. Tiesa, vidējā nopirktā zemes gabala lielums bija tikai ap 5 desetīnām (1 desetīna~--- 1,0925~ha). Trešdaļā zemnieku sētu nebija zirga, tomēr radās arī spēcīgs vidējās zemniecības slānis. Ciemos, lai mazinātu muižniecības ietekmi, tika organizētas gminas (kopienas). Ar to, lai cik paradoksāli arī tas skanētu, faktiski Krievijas apspiešanas politika palīdzēja veidoties modernajai poļu nācijai. 1899.~gadā Mehovas (\pltxti{Miechów}) apriņķī tika nodibināts pirmais lauksaimniecības artelis. Tas kļuva par paraugu citām lauksaimniecības biedrībām Polijas karalistē.

Taču pastāvēja arī virkne negatīvu parādību. Lēni auga lauksaimnieciskās tehnikas pielietojums. Piemēram, Varšavas apgabalā 1895.~gadā tikai 4\% saimniecību bija kuļmašīnas, kad Pozenes provincē Prūsijā ar tām strādāja jau 52\% saimniecību. Tikai pasaules lauksaimniecības krīze 80.--90.~gados, saistīta ar lētu graudu ievedumu no aizokeāna valstīm, paātrināja pārmaiņas, lielsaimniecību pārvēršanos par kapitālistiskiem uzņēmumiem. Taču tieši šai laikā muižnieki arī samazināja graudaugu sējumu platību un daļu muižu, kuras nedeva regulārus ienākumus, pārdeva, tai skaitā arī zemniekiem.

Lauku iedzīvotāju skaits auga tik strauji, ka pat esošo lielsaimniecību parcelācija (sadalīšana sīkākās) nevarēja novērst bezzemnieku skaita pieaugumu. 1870.~gadā bija 220~000~bezzemnieku, bet 1891.~gadā~--- jau 849~000. XX gs sākumā viņu bija ap 1,2 miljoni.

Kaut carisma ekonomiskā politika pret Polijas karalisti bija daļa no tā vispārējā kursa uz ``centra'' atbalstīšanu par ļaunu ``nomalēm'', \strong{Polijas karalistes rūpniecības attīstībā} vislielākā loma bija tam, ka tā varēja realizēt saražoto produkciju Viskrievijas tirgū. Pēc 1861.~gada dzimtbūšanas atcelšanas Krievijā strauji veidojās iekšējais tirgus, kas labvēlīgi ietekmēja arī Polijas karalisti, tās rūpniecības rajonus.

Liela loma Polijas karalistes attīstībā joprojām bija dzelzceļu celtniecībai. T.s. Lodzas fabriku dzelzceļa līnijas \rutxti{(Лодзинская фабричная железная дорога}) celtniecība (1867) zināmā mērā kompensēja to, ka agrāk uzceltā Varšavas--Vīnes līnija atstāja malā Lodzas rūpniecības rajonu. Pēc 70.~gados notikušās Terespoles (\pltxti{Terespol}) dzelzceļa savienošanas ar Maskavas-Brestas un Kijevas-Brestas dzelzceļiem Polijas karaliste ieguva sakarus ar Krievijas impērijas centrālajiem un dienvidu rajoniem. 80.~gados dzelzceļu tīkls Polijas karalistē bija visattīstītākais Krievijā. Ja 1870.~gadā Polijas karalistē visu dzelzceļu garums sastādīja 960~km, tad 1910.~gadā~--- jau 3~810~km. Sevišķi svarīga bija novada sasaiste ar Vīni (1845--1848), Pēterburgu (1862), Berlīni (1862) un Maskavu (1870). Varšavas dzelzceļa mezgls atpalika tikai no Maskavas mezgla.

XIX gs. 60--70.~gados Polijas karalistē norisa \strong{rūpniecības apvērsums}. (Metalurģijā un ogļrūpniecībā tas noslēdzās 80.~gados.) Ja 1850.~gadā Lodzā vēl strādāja tikai četras tvaika mašīnas ar 150 zirgspēku jaudu, tad jau 60.~gadu sākumā tur pastāvēja 10~fabrikas, kurās darbojās tvaika un mehāniskie dzinēji. Pēc tam, kad 60.~gados aušana tika mehanizēta, tekstilrūpniecības peļņa pieauga no 60~miljoniem rubļu 1880.~gadā uz 341~miljonu rubļu 1910.~gadā. 80.~gados ap 3/4 kokvilnas rūpniecības ražojumu izveda uz Krieviju vai caur to uz austrumu tirgiem. 1883.~gadā tekstilrūpniecību gan skāra krīze, bet tā nebija ilgstoša. 1885.~gadā tekstilrūpniecība un citas vieglās rūpniecības nozares deva vairāk nekā 50\% no visas karalistē saražotās rūpniecības produkcijas vērtības.

XIX gadsimta pēdējā desmitgadē karalistes rūpniecība mehanizācijas ziņā bija līdere Viskrievijas mērogā, bet 1867.~gadā no Varšavas, Kališas un Kelces guberņu daļām izveidotā Petrakovas guberņa (\pltxti{Gubernia Piotrkowska}) 1892.~gadā šai rādītājā ieņēma pirmo vietu, apsteidzot Pēterburgas, Jekaterinoslavas un Maskavas guberņas.

Auga rūpnieciskās ražošanas koncentrācija lielos uzņēmumos. XIX gadsimta beigās lielās rūpnīcas (kuras nodarbināja vairāk nekā 50~strādniekiem) sastādīja 31\% no visu rūpnīcu skaita, taču ražoja 88\% visas rūpnieciskās produkcijas, tajās bija koncentrēti 86\% visu strādnieku.

Kapitālistiskās saimniecības sistēmas attīstība prasīja modernu kredīta sistēmu, radās jauna tipa bankas~--- akciju sabiedrības. 1870.~gadā Varšavā tika dibināta Tirdzniecības banka, kurā tika iesaistīti lielo zemes īpašnieku kapitāli. Banka veica starpniecību graudu un cukura tirdzniecībā, piedalījās investīcijās rūpniecībā. 70.~gados radās banka Lodzā, kura apkalpoja tekstilrūpniecību un tās ražojumu eksportu uz Krieviju, arī citas.

Lai veiktu jaunas investīcijas rūpniecībā, bija nepieciešami lieli kapitāli. Tāpēc auga akciju sabiedrību veidoti uzņēmumi. XIX/XX gadsimta mijā akciju sabiedrību īpašumā esoši uzņēmumi sastādīja tikai 9\% no kopējā uzņēmumu skaita, taču ražoja vairāk nekā 50\% visas rūpniecības produkcijas. To īpatsvars saražotās produkcijas kopējā vērtībā kalnrūpniecības un metalurģiskajā rūpniecībā sastādīja 92,1\%, metālapstrādes rūpniecībā 67,8\%, cukurrūpniecībā 53,1\%.

Sakarā ar Krievijas realizēto protekcionisma politiku, raksturīgu ar augstām muitas nodevām uz no ārzemēm ievedamajām precēm, kuras citādi varētu izkonkurēt iekšzemes ražojumus, Polijas karaliste, tāpat kā Krievija kopumā, kļuva arī par ārzemju kapitāla iespiešanās objektu. XX gadsimta sākumā franču, beļģu un vācu kapitālam piederēja apmēram 1/4 no visiem karalistes metalurģiskās, ķīmiskās un tekstilrūpniecības uzņēmumiem, taču tas nodarbināja vairāk nekā 69\% no kopējā strādnieku skaita un radīja vairāk nekā 60\% rūpnieciskās produkcijas. Tā 1901./1902.~gadā ārzemju kapitālam Polijas karalistē piederēja 669~fabrikas, tas nodarbināja 37,6~tūkstošus strādnieku, radīja produkciju 272~miljonu rubļu vērtībā. Nacionālajam kapitālam piederēja 1971~fabrika, tas nodarbināja 84~tūkstošus strādnieku un ražoja produkciju 181~miljona rubļu vērtībā. Minētie dati ļauj konstatēt, ka viens ārzemju kapitālam piederošs uzņēmums vidēji ražoja produkciju 406~tūkstošu rubļu vērtībā, nodarbinot vidēji 56~strādniekus, kad vietējam kapitālam piederošs uzņēmums, nodarbinot vidēji 43~strādniekus, ražoja produkciju tikai 92~tūkstošu vērtībā. Pēc dažiem datiem ārzemju kapitāls sastādīja 39\% no karalistes rūpnieciskā kapitāla.

Tirdzniecība ar Krieviju un ārzemju, galvenokārt franču, beļģu un vācu kapitāli, investēti Polijas karalistes rūpniecībā, bija tās tautsaimniecības pamats. Tiesa, kā norādījis poļu ekonomikas vēsturnieks V.~Kula, ne visus ar vāciskiem uzvārdiem saistītos kapitālieguldījumus Polijā varēja uzskatīt par ārzemju kapitāliem. Tā, Lodzas tekstilrūpnīcu īpašnieki vācieši L.~Geiers un K.~Šeiblers savus līdzekļus bija ieguldījuši Polijā, paši tos pārvaldīja uz vietas, peļņu izmantoja ražošanas paplašināšanai, nevis izveda uz ārzemēm.

Šai sakarā jāpiemin, ka XIX gadsimta beigās~--- XX gadsimta sākumā zināmu izplatību ieguva t.s. kapitālisma Polijā attīstības ``koloniālā teorija'', kura kapitālismu uzskatīja nevis par iekšēju sociāli-ekonomisko pārmaiņu rezultātā radušos, bet svešu, ārpus Polijas dzimušu un tās augsnē ``mākslīgi'' ``pārstādītu'' parādību. To raksturojot ārzemju kapitāli, noieta tirgi ārzemēs, ārzemju izcelsmes buržuāzija un arī lielā daudzumā no ārzemēm iebraukuši strādnieki. Kā uzskatīja V.~Kula, minētajā posmā un īpaši pēc Polijas valstiskās neatkarības atjaunošanas 1918.~gadā šai teorijai bija uzdevums pierādīt, ka poļu rūpniecība ir radusies tikai pateicoties ``austrumu tirgiem'' un poļiem ir nacionāli vienoti jācīnās gan pret no ārzemēm nākušo lielburžuāziju, gan reizē par šo tirgu saglabāšanu. Pats V.~Kula marksistiski vērtēja kapitālisma veidošanos kā nesaraujami vienotu vēsturisku procesu, ko izsauc ražošanas spēku attīstība un progress sabiedriskā darba dalīšanā, iedzīvotāju noslāņošanās, jaunu kapitālismam raksturīgu šķiru~--- buržuāzijas un proletariāta rašanās, bagātības akumulācija vienā sabiedrības polā un nabadzības otrā, kapitāla uzkrāšana vienu rokās un otru pāreja algotu darbinieku kategorijā. Tomēr, atzīstot, ka komplekss tirgus ekonomikas izveides skaidrojums ir tuvāks īstenībai, iekšējie procesi bija noteicošie, nedrīkst arī vienkāršot situāciju un noliegt, ka gan poļu zemju kopumā, gan Polijas karalistes tautsaimniecība atsevišķi, nebija izolēta no apkārtējo valstu, īpaši Poliju sadalījušo, ekonomiskās attīstības, ka starptautiskās norises, arī kapitāla imports, ietekmēja poļu zemju ekonomisko attīstību.

No otras puses, Polijas karalistes rūpniecības arvien lielāka piemērošanās tirgiem austrumos padarīja to pārlieku atkarīgu no tiem. Dārgo izejvielu imports no Rietumiem palielināja poļu produkcijas pašizmaksu un neļāva tai sekmīgi konkurēt rietumu tirgos. Karalistes rūpniecība tāpēc kļuva īpaši jūtīga pret konjunktūras pasliktināšanos vai politisko attiecību izmaiņām.

Joprojām pirmo vietu rūpniecības attīstības ziņā Polijas karalistē ieņēma Lodzas (\pltxti{Łόdz}) rajons. Tajā ražoja 80\% visu kokvilnas audumu. Interesanti, ka jau 60.~gados Lodzā bija vērojams plašs katolisko poļu strādnieku pieplūdums. (XIX gadsimta beigās tajā bija nodarbināti 95 tūkstoši strādnieku.) Turpretī uzņēmēju vidū poļu bija maz. Rūpnieki galvenokārt bija vācieši. Poļi uzņēmēju Lodzā vēl pirms Pirmā pasaules kara sastādīja tikai 6,5\% un arī tie pārsvarā bija sīkīpašnieki. Vācieši šeit, balstoties uz saimniecisko spēku, varēja droši aizsargāt savu nacionālo identitāti. Kontakti dažādu tautību uzņēmēju starpā bija drīzāk izņēmums, nevis parasta parādība. Salīdzinājumā, piemēram, ar Varšavu vācu uzņēmēji Lodzā dzīvoja izolācijā. No viņiem tikai dažas ģimenes bija polonizētas. XIX gs. beigās starp Lodzas tekstiluzņēmējiem vācieši bija 49\%, ebreji~--- 45,7\%. XX gs. sākumā proporcija pat mainījās ebreju uzņēmējiem par labu. (Lodzā, tāpat kā Varšavā, radās asimilētas ebreju pilsonības slānis.) 1913.~gadā Lodzā ebreju uzņēmēju vidū bija jau 47,1\%, vāciešu 44\%. Tomēr lieluzņēmēju vidū vāciešu bija vairāk un viņi saglabāja tekstilrūpniecībā vadošo vietu. Tikai neliels skaits vāciešu ieradās Polijā jau ar vērā ņemamu kapitālu. Sekmīgai viņu darbības izvēršanai ļoti liela nozīme bija no valsts piešķirtajiem kredītiem. Krīžu laikā, kad kredīti aptrūka, tas veda pie lielām grūtībām. Tomēr XIX gs. beigās daļa uzņēmēju darbojās jau ar savu kapitālu, kļuva neatkarīgi no valsts kredītiem. Liela daļa vācu uzņēmēju tehniskās zināšanas bija jau apguvuši dzimtenē. Viņi arī Polijā uzturēja sakarus ar Vāciju, ar to Rietumu sasniegumus ienesa Polijā. Kā rakstīja poļu literatūras klasiķis V.~Reimonts, radās t.s. ``Lodzas cilvēks''~--- ``auksts, gudrs, lietišķs cilvēks, kurš bija gatavs uz visu''.

Tekstilrūpniecības attīstība, protams, bija saistīta ar metālizstrādājumu ražošanu~--- metalurģiju un mašīnrūpniecību, jo bija vajadzīgi darbgaldi, stelles u.c. iekārtas. Lodzā līdz ar tekstilrūpniecības attīstību, pilsētas iedzīvotāju skaita pieaugumu attīstījās arī pārtikas rūpniecība, papīra ražošana, tipogrāfiju darbība. Pirms Pirmā pasaules kara Lodzā bija jau ap 20~multimiljonāru ģimeņu. Vācu uzņēmēji Polijas karalistē izturējās ļoti ekspansīvi, ieguldīja savus kapitālus ne tikai savu pamatinterešu, bet arī citās saimniecības nozarēs tuvākajos Krievijas rietumu apgabalos. Taču rajonos, kur stingri dominēja krievu kapitāls~--- Aizkaukāzā, Centrālajā Āzijā, vācieši darboties negribēja, jo jau tā Lodzas uzņēmējiem nācās sīvi konkurēt ar Maskavas tekstilrūpniekiem.

Otrais pēc nozīmes Polijas karalistē bija Varšavas rūpniecības rajons, kur attīstījās dažādas rūpniecības nozares un pastāvēja gan lieli uzņēmumi, gan liels skaits sīku amatnieku darbnīcu. Tā 1895.~gadā šeit bija 44~000 rūpniecības strādnieku un 55~000 amatnieku.

Pakāpeniski arvien lielāku lomu ieguva Dombrovas ogļu baseins (mūsdienās ietilpst Silēzijas vojevodistē). XIX gadsimtā tas, atšķirībā no Augšsilēzijas, kura atradās Prūsijas sastāvā, ietilpa Polijas karalistē. Līdz pat Pirmajam pasaules karam akmeņogles tur ieguva galvenokārt daudzās sīkās, slikti aprīkotās šahtās. Tomēr, ja 1870.~gadā šeit ieguva 300~000 tonnu ogļu, tad 1890.~gadā~--- jau 3~000~000 t. XIX gs. beigās baseinā ieguva ap 4 miljonus tonnu ogļu gadā. Pirms Pirmā pasaules kara Dombrovas ogļu baseins deva jau 19\% no visas Krievijas impērijas akmeņoglēm. Kopumā Polijas karalistes ogļu ieguves rūpniecībā 19 šahtas piederēja 13 firmām, tai skaitā piecām lielām. Starp tām četras, kā arī divas mazākas pārstāvēja ārzemju kapitālu. Poļu kapitālam piederēja tikai divas firmas. 1897.~gadā sešas ogļu ieguves sabiedrības Dombrovas baseinā noslēdza sadarbības līgumu. Tām piederošajās šahtās ieguva 91\% no visām akmeņoglēm.

Neapturami samazinājās amatniecības pienesums rūpnieciskās produkcijas ražošanā. No 1855. līdz 1866.~gadam ap 25~000, no 1866. līdz 1889.~gadam vēl ap 22~000~amatnieku kļuva par rūpniecības strādniekiem. Ja 50.~gadu beigās amatnieki saražoja preces gandrīz par tādu pat summu kā manufaktūras un fabrikas, tad jau 1877.~gadā rūpnīcu un fabriku ražojumu vērtība pārsniedza amatniecības izstrādājumu vērtību vairāk nekā četras reizes.

Kaut arī agrākā feodālā cunfšu reglamentācija visur bija likvidēta, amatniecībā gan vēl saglabājās cunfšu organizāciju struktūra: māceklis~--- zellis~--- meistars. Zeļļi centās kļūt par meistariem, iekārtot savu darbnīcu, taču veiksmes gadījumi bija rets izņēmums. No 1871. līdz 1880.~gadam tikai kokvilnas audumu ražošanā iznīka ap 7~000~sīko darbnīcu. Dažas amatnieku specialitātes~--- audēju, vērpēju, kariešu izgatavotāju u.c., neizturot konkurenci ar fabrikām, vispār izzuda. Ja 1870.~gadā karalistē bija mazāk nekā 70~tūkstošu, tad 1880.~gadā Polijas karalistē strādāja jau 150~000 rūpniecības strādnieku (33\% no tiem~--- tekstilrūpniecībā). Salīdzinājumam~--- Vācijai un Austroungārijai piederošajos poļu apdzīvotajos apgabalos bija tikai 30~000 strādnieku. 1895.~gadā Polijas karalistē bija jau ap 250~tūkstošu strādnieku. Lielākā daļa viņu nāca no bezzemnieku un trūcīgo zemnieku vidus. 1897.--1901.~gadā rūpniecības strādnieku vidū gandrīz 1/4 sastādīja sievietes. Tekstilrūpniecībā viņu bija vairāk par pusi.

1890.~gadu sākumā Polijas karalistē saražoja 19,5\% visas Krievijas produkcijas (ir dati, ka pat 1/4), kaut iedzīvotāju tajā bija tikai 7\% no visiem impērijā dzīvojošajiem. Uz vienu cilvēku ražojamās produkcijas ziņā karaliste atpalika tikai no Pēterburgas un Maskavas. Ogļu un tērauda ražošanā Polija bija otrajā vietā aiz Doņeckas baseina. Polijā no visiem Krievijā ražotajiem izstrādājumiem tika izgatavots 42\% kokvilnas auduma, 26,9\% vilnas auduma, 77\% dzijas, 76\% adījumu. Lielākā daļa ražojumu (ap 70\%) tika pārdoti Krievijā. Laikā no 1880. līdz 1910.~gadam karalistes tirdzniecība ar Rietumu valstīm palielinājās divreiz, bet ar Krieviju~--- 11~reizes. Galvenokārt turp tika izvestas tekstilpreces, dzija, mašīnas, ogles. Šāds stāvoklis izraisīja lielu neapmierinātību krievu rūpniekos. Taču Polijas karaliste ne tikai izveda preces uz Krieviju, bet arī pati bija Krievijas noieta tirgus. No Krievijas karaliste importēja graudus, miltus, putraimus, 60\% nepieciešamās kokvilnas, 50\% vilnas, lielu daudzumu zīda, linu, dzijas, dzelzs rūdu no Krivoirogas (krievu \rutxti{Кривой Рог}, ukr. \uktxti{Кривий Ріг}).

Īpaši asa konkurence bija tekstilrūpniecībā un ogļrūpniecībā. Var teikt, ka risinājās Maskavas cīņa ar Lodzu. Maskaviešiem lētāk izmaksāja izejvielu piegāde, toties dārgāk kurināmais. Poļu ražotājiem bija pieejami tikai dārgāki kredīti kā krievu fabrikantiem, toties viņiem bija jāmaksā gandrīz divreiz mazāki nodokļi. Strādnieku algas Polijas fabrikās bija ievērojami augstākas nekā Krievijas centrā, kas ļāva maskaviešiem ietaupīt ievērojamus līdzekļus. Turklāt Maskavas fabrikās darba dienas ilgums bija 14 stundas, bet Polijā 12--13~stundas. Tiesa, Polijā gadā fabrikas strādāja 7~dienas vairāk. Kopumā ražošanas apstākļi bija aptuveni vienādi. No otras puses, pēc dažiem vērtējumiem, Maskavas fabrikas ievērojami atpalika darba organizācijā. Nav pamata apgalvot, ka Polijas rūpniecība varētu nest ievērojamus zaudējumus krievu ražotājiem, Polijas tirdzniecības balanss ar Krieviju pastāvīgi bija pasīvs, taču krievu publicisti un rūpnieki bieži pārspīlēja Polijas ekonomiskās izaugsmes rādītājus.

Jāuzsver, ka starp Polijas karalistes rūpniekiem, finansistiem un Rietumueiropas, kā arī Krievijas kolēģiem pastāvēja plaši sakari. Polijas karalistes kapitāls centās veikt investīcijas arī Krievijas tirgos. Piemēram, tas finansēja dzelzs ražošanu Krivoirogā un Kramatorskā (ukraiņu \uktxti{Краматорськ}~--- pilsēta Doņeckas apgabalā Ukrainā). Polijas karalistes teritorijā darbojās lielo Krievijas banku filiāles.

Lai nu kā, Polijas karaliste 19.~gadsimta beigās stingri ieņēma trešo vietu ražošanas apjoma ziņā starp Krievijas rūpniecības rajoniem aiz Maskavas un Baltijas rajona. Asā konkurences cīņa ar Krievijas rūpniekiem atstāja arī jūtamu iespaidu uz poļu buržuāzijas politiskās platformas izveidi.

Salīdzinājumā ar poļu zemēm dienvidos un arī rietumos (izņemot Silēziju) Polijas karaliste bija vairāk rūpnieciski attīstīta, taču salīdzinājumā ar Rietumeiropas pirmrindas valstīm tā ievērojami atpalika.

Sekmīgā tautsaimniecības attīstība nodrošināja salīdzinoši labvēlīgu demogrāfisko situāciju. Līdz XIX~gadsimta beigām Polijas karalistē strauji pieauga iedzīvotāju skaits. Ja 1863.~gadā tajā bija 5 miljoni iedzīvotāju, tad 1897~--- 9,5 miljoni, sasniedzot 75 cilvēkus uz km$^{2}$. Ap 3/4 iedzīvotāju bija poļi, 14\% ebreji, 4,7\% krievi un ukraiņi, (Pēc 1897.~gada skaitīšanas datiem no Polijas iedzīvotājiem tikai 2,8\% bija krievi. No citiem Krievijas rajoniem viņu vēl mazāk bija tikai Somijā~--- 0,23\%) 4,4\%~--- vāciešu, 3,3\%~--- lietuviešu. Pēc konfesionālās piederības katoļu bija 75\%, judaistu~--- 14\%, protestantu~--- 6\% un pareizticīgo 5\%. Visā karalistē līdz 2/3 iedzīvotāju nodarbojās ar zemkopību, taču jau 27\% pelnīja iztiku ar tirdzniecību, amatniecību un rūpniecību. Šļahtiči sastādīja vairs tikai 1,5\% iedzīvotāju).

Tā kā 1870.--1890.~gados rūpniecība Polijā intensīvi attīstījās, strauji pieauga pilsētu iedzīvotāju skaits. Ja XIX gadsimta vidū Polijas karalistē to bija ap 16\%, tad gadsimta beigās jau 35\%. Varšavā 1858.~gadā bija 160~000, bet 1897.~gadā~--- jau 638~000 iedzīvotāju, taču Lodzā pieauguma tempi bija 2~reizes straujāki: no 40~000 līdz 315~000, kas izskaidrojams ar tās rūpniecisko attīstību.

1897.~gadā Polijas karalistē bija ap 250~000 rūpniecības strādnieku, vairāk nekā 67~000 dzelzceļnieku un ceļu būvētāju, apkalpojošā sfērā bija nodarbināti vairāk nekā 270~000 cilvēku. Varšavas metālistu un Dombrovas kalnraču vidū augsts bija kvalificēto strādnieku īpatsvars. Var teikt, ka tas viss liecināja par salīdzinoši viduvēju ekonomiskās attīstības līmeni. Strādnieku pamatmasa Polijas karalistē bija poļu tautības. Pārējo daļu veidoja ebreju un nelielā skaitā arī vācu strādnieki. Ebreju strādnieku liela daļa koncentrējās Lodzā un Belostokā, vāciešu~--- Lodzā un Dombrovas-Sosnoveckas rajonā. Toties augsts bija vāciešu procents visas karalistes rūpniecības vidējā tehniskā personāla vidū.

Turpinājās tālāka \strong{poļu nacionālās apziņas veidošanās}. Gan nepārtrauktās sakāves, ko piedzīvoja poļu nacionālisti, gan pārmaiņas Polijas karalistes sociāli-ekonomiskajos apstākļos veicināja jaunu idejisku procesu attīstību. Jāatzīmē, ka Polijas karalistes tautsaimniecības attīstībai labvēlīgā cariskās Krievijas ekonomiskā politika bija ne tikai priekšnoteikums tās ekonomiskajam uzplaukumam, bet arī poļu nacionālās pašapziņas pieaugumam, sabiedriskās aktivitātes kāpumam. Šļahtas loma mazinājās. Liela inteliģences daļa gan vēl nāca no šļahtas, bet, kā vēlāk atzina viens no poļu sociālistu (PPS) ideologiem, L.~Vasiļevskis, izteica jau rūpnieciskās buržuāzijas intereses.

Tieši ar kapitālisma attīstību poļu zemēs bija saistīta liberālā virziena attīstība. Liberāļi prasīja likvidēt feodālisma paliekas un šļahtiskās tradīcijas, aicināja visus slāņus apvienoties cīņā par nacionālajām interesēm. Taču Polijā liberālisms nespēja attīstīties tādā varenā strāvojumā, kāds tas bija valstīs, kuras izvirzījās priekšgalā kapitālisma attīstībā. Salīdzinoši vēlākā poļu zemju nostāšanās uz kapitālistiskās attīstības ceļa noteica dažādu sociālo un politisko procesu pārbīdi un sakrišanu laikā. Polijā salīdzinoši ātri attīstījās t.s. ``trešās šķiras'' (sīkpilsonības) diferenciācija, kas laupīja liberālismam to sociālo pamatu, ko šī ``trešā šķira'' varēja radīt, uzstājoties pret feodālismu kā vienots vesels. Salīdzinoši ātri norisa strādniecības un zemniecības politiskā pašnoteikšanās.

Pēc 1863.--1864.~gadu sacelšanās sakāves jaunākā poļu nacionālistu paaudze nereti deva priekšroku praktiskam darbam savas tautas labā, nevis cīņai par nacionālas valsts izveidi, kurai, kā toreiz likās, nav perspektīvu. Poļu aristokrāti Vislas novadā augstu vērtēja savas krievu aristokrātijai pielīdzinātās privilēģijas, arī poļu lielburžuāzija parasti bija konservatīva, jo ar carismu to saistīja milzīgā Krievijas tirgus izmantošanas intereses.

1876.--1881.~gadu periodu var uzskatīt par zināmu ``atkusni'' carisma antipoliskajā politikā. Krievijas varas iestādes šai laikā vairāk bija nodarbinātas ar ārpolitiskām akcijām un cīņu pret revolucionāro kustību pašā Krievijā. Poļu sabiedrībā auga cerības uz pašvaldības tiesību piešķiršanu karalistei, plaši tika apspriesta iespēja panākt samierināšanos ar cara valdību.

XIX gadsimta 70.--80.~gados mēreno poļu politiķu vidū pilsoniskās jeb buržuāziskās kustības ietvaros radās ``pozitīvisma'' (šeit~--- politiskās un tiesiskās domas virziens XIX gs. otrajā pusē, kurš par vienīgajām atzina valsts izveidotās tiesību normas) jeb tā saucamās ``organiskā darba'' (\pltxti{praca organiczna}) skolas socioloģijā piekritēji. (Šī skola bija radusies XIX gs. 40.~gados Vācijā un pēc tās uzskatiem sabiedrība ir organisms un attīstās pēc bioloģiskām likumsakarībām.) Viņu atbalstītā ``darba pie pamatiem'' (\pltxti{praca u podstaw}) un ``organiskā darba'' teorija virzīja uz tautas ekonomisko un kultūras sasniegumu tālāku attīstību, atstājot politiskās prasības tikai pašvaldību līmenī, neprasot valstiskuma atjaunošanu. Šie mērenie politiķi uzsvēra, ka konspiratīvā darbība un bruņota cīņa par nacionālo neatkarību ir kaitīga. ``Organiskā darba'' skolas programma paredzēja koncentrēšanos ekonomisko, sociālo un kultūras jautājumu risināšanai, propagandēja rūpniecības un tirdzniecības, zinātnes, tehnikas attīstību, noraidot ticību poļu mesiānismam un cerībām ko panākt ar ieročiem.

Šai sakarā var pieminēt vēsturisku anekdoti. Kad Francijas valstsvīram A.~Tjēram kāds poļu sabiedriskais darbinieks pajautājis, kas poļiem jādara, lai atkal tiktu uz augšu, viņš atbildējis: ``\frtxti{Enrichissez vous'}!'' (Kļūstiet bagāti!)

Par ļoti svarīgu tika uzskatīta apzināta darbība sabiedrības attīstībā. Liberāļi-pozitīvisti pasludināja, ka ``lai arī kādā stāvoklī būtu sabiedrība, vienmēr atradīsies sfēra, kurā tā var strādāt savā labā''. Liela loma tika piešķirta izglītībai.

Krievijā šo virzienu šai laikā propagandēja viens no narodņiku līderiem N.~Mihailovskis ar savu ``mazo darbu'' teoriju. XIX gs. 80.~gadu sākumā valstu~--- Polijas sadalītāju savstarpējās nesaskaņas, arī zināma Krievijas politikas liberalizācija veicināja domu par kompromisa ar Krieviju atrašanas nepieciešamību.

Polijā viens no pirmajiem ``organiskā darba'' sludināšanu uzsāka žurnāls ``\pltxti{Przegląd Tygodniowy}'' (''Iknedēļas Apskats''). Tas daudz vērības veltīja zemes ekonomiskajai attīstībai, propagandēja rūpniecības un tirdzniecības, zinātnes un tehnikas attīstību, uzsvēra, ka sabiedrības uzdevums ir materiālo vērtību uzkrāšana. ``Organiskā darba'' teorija noraidīja uz vardarbību virzītu ideoloģiju, atzina mierīgu, reformistisku darbību. Politikā pozitivisma ietekmē arī radās t.s. samierināšanās (\pltxti{ugoda}) programma, no šejienes šī grupējuma pārstāvjus arī sāka saukt par ugodoviešiem. Viņi centās pārliecināt cara valdību par poļu uzticamību un nepieciešamību tāpēc ieviest Polijas karalistē zemstes, pilsētu pašvaldības un zvērināto piesēdētāju tiesas, veikt citas reformas. Tiesa, nacionāli-patriotisko tradīciju aizstāvji nereti ugodoviešus uzskatīja par nodevējiem, renegātiem.

Kā atzīmējis E.~Meijers, pūliņiem veicināt visu slāņu, arī zemnieku, labklājību bija panākumi, bet tie veicināja arī lietuviešu un ukraiņu nacionālo apziņu, to, ka šīs tautas par savu ienaidnieku uzskatīja ne tik daudz Krievijas impēriju, bet gan poļu augšslāni, iepriekšējās kopības starp poļiem, lietuviešiem un ukraiņiem izjukšanu. Pēc autora domām par šādas kopības pastāvēšanu līdz XIX gadsimta beigām var runāt tikai nosacīti. Tā gan darbojās kopēju cīņu laikā pret ārējiem iebrucējiem, bet sociālajām sadursmēm starp apspiestajiem un apspiedējiem jau iepriekšējos gadsimtos bija arī nacionāla nokrāsa, jo apspiedēji lielākoties bija poļi. Toties jāpiekrīt minētā vēsturnieka konstatējumam, ka jēdziens ``Polija'' šai laikā sašaurinājās līdz poļu apdzīvotajām teritorijām. Krievijas rietumguberņu iedzīvotāji uz savām teritorijām to vairs neattiecināja.

Uz Vislas novadu lielu iespaidu atstāja izmaiņas Krievijā valdošajā dinastijā. 1881.~gada 1.~(13.)~martā, krievu revolucionārās organizācijas ``\pltxti{Нaродная воля}'' (Tautas griba) poļu-baltkrievu izcelsmes biedra I.~Griņevicka mestās bumbas sprādziena rezultātā gāja bojā Aleksandrs II. (Baltkrievu publicists A.~Tarass, padevies asinsatriebības garam un ignorējot to, ka atentātā gāja bojā 3 nevainīgi cilvēki un vēl 17 tika ievainoti, šai sakarā rakstīja, ka Aleksandrs II ``beidzot, ar savu dzīvību samaksāja par Kaļinovska, Serakovska, Trauguta un tūkstošu citu poļu un baltkrievu revolucionāru dzīvībām''.) Tronī nonāca Aleksandrs III, vara spēja pārdzīvot narodņiku izraisītos satricinājumus, tā atkal pārgāja uzbrukumā pret tiem, ko uzskatīja par valstij naidīgiem spēkiem. Krievijas valdības attieksme pret poļiem atkal pasliktinājās. 80.~gadu vidū kļuva skaidrs, ka ``organiskais darbs'' nav tik radikāls līdzeklis kā panākt uzlabojumus, kā agrāk tas šķita pozitīvistiem.

Aleksandrs III gan valdīja samērā neilgi, tomēr no revolucionāru-marksistu viedokļa raugoties, Krievijas stāvoklis, un arī viņu attieksme pret to būtiski mainījās. Ja vēl 1893.~gadā K.~Markss rakstīja, ka ``Krievijā var būt tikai tāds vai šāds dumpis, pie tam cietīs vācu tērpos ģērbtie, bet revolūcijas nekādas un nekad nebūs'', tad 1881.~gadā viņš paziņoja, ka ``Krievija ir revolucionārās kustības priekšpulks Eiropā''. Šāds vērtējums balstījās uz Aleksandra II nogalināšanas faktu, krievu teroristu skaļo darbošanos. Tika cerēts, ka cara patvaldība, kura viņaprāt bremzēja revolūcijas attīstību Eiropā, tiks iznīcināta. Mūsdienās skaidri redzams, ka secinājums, pamatots nevis ar strādniecības augšanu, tās ``idejisko briedumu'', kura svarīgumu tik ļoti uzsvēra marksisti, bet gan ar teroristu organizācijas īslaicīgu veiksmi, bija utopisks. Taču to uztvēra K.~Marksa sekotāji Krievijā, kuri sāka runāt par XIX gadsimta beigās notikušo ``pasaules revolucionārās kustības centra pārvietošanos uz Krieviju''. Līdz ar revolucionāro pretrunu nobriešanu Krievijā Polija nevarēja nezaudēt savu izcilo vietu Eiropas revolucionārajā kustībā.

Imperatora Nikolaja II nākšana Krievijas un Polijas tronī (1894) izsauca poļu cerību atdzimšanu uz valdošā režīma izmaiņām. Šai pat gadā notikusī vecā Varšavas ģenerālgubernatora J.~Gurko nomaiņa arī pastiprināja samierinātāju cerības. Daudziem šķita, ka rusifikācijas smagums jau ir aiz muguras. Tika organizēta virkne deputāciju uz Pēterburgu, lai panāktu pārmaiņas Polijas karalistē. Viena no lielākajām savu pavalstniecisko jūtu demonstrēšanas akcijām notika 1897.~gadā, kad Nikolajs II apmeklēja Varšavu. Rusofīlie elementi organizēja krāšņu viņa sagaidīšanu. Nikolajs II ziedoja 1 miljonu rubļu Varšavas tehniskās augstskolas dibināšanai, atļāva uzcelt pieminekli nacionālajam dzejniekam Ā.~Mickēvičam. Krievija uzņēma 30~gadus atpakaļ pārtrauktos kontaktus ar Vatikānu. Taču poļu ugodoviešu cerības neattaisnojās, jau 1898.~gadā atkal tika atjaunota stingra izturēšanās pret ``polonismu''.

Tā poļu historiogrāfija, kas balstās tikai uz aktīvo cīņas formu pret apspiedējiem slavināšanu, līdz šai dienai pozitīvisma teorijas valdīšanas periodu sabiedrībā raksturo kā ``sociālo marasmu''. Arī padomju vēsturnieki uzsvēra, ka Polijā, kur kapitālisma attīstība bija sākusies vēlāk nekā Rietumos, toties tā attīstība notika straujāk, tāpēc arī šķiru cīņa stimulēja buržuāzijas ātrāku nobriešanu un tā, nepaguvusi līdz galam izveidoties, neieguva to revolucionāro potenciālu, kāds bija pilsonībai Rietumos buržuāzisko revolūciju periodā. Tāpēc arī poļu pilsonība nevarējusi vadīt buržuāziski-demokrātisko kustību Polijā. Tā, latviešu vēsturnieks Indulis Ronis rakstīja, ka pozitīvisma programma ``atspoguļoja valdošo šķiru centienus pilnā mērā izmantot tās ekonomiskās izmaiņas un iespējas, ko pavēra kapitālisma attīstība, un ar reformu palīdzību buržuāziski pārveidot sabiedrību. Tādēļ, no vienas puses, pozitīvisti kritizēja agrāko šļahtas ideoloģiju, tās feodālās privilēģijas, bet, no otras~--- bija klaji naidīgi politiskajai, bet it īpaši revolucionārajai cīņai''. Mūsdienās nākas konstatēt, ka šī pieeja bija dogmatiska, mākslīgi nodalot revolucionāro cīņu no cīņas par reformām, cīņu par pilsonības ekonomiskajām interesēm no cīņas par demokrātiju. Atkarībā no vēsturiskajiem apstākļiem arī poļu buržuāzijas interešu priekšplānā izvirzījās dažādi uzdevumi. Kā rādīja vēlākie XX~gadsimta notikumi, tā spēja uzstāties arī kā demokrātisks spēks.

XIX gadsimta pēdējā trešdaļā pieauga \strong{sasprindzinājums attiecībās starp Vislas novada iedzīvotāju vairākumu~--- poļiem un nacionālajiem mazākumiem}. Poļu zemēs dzīvojošajās nepoļu tautu pārstāvji sāka paust savu nacionālismu

Polijas karalistē dzīvojošie \strong{vācu} zemnieki jau no XIX gadsimta 30.~gadiem izrādīja vēlmi izceļot no Polijas uz blakus esošo Krievijas Volīniju (ukraiņu \uktxti{Волинь}, poļu \pltxti{Wołyń}, mūsdienās novads aptver vairākus apgabalus Ukrainā un Ļubļinas vojevodistē Polijā), jo apkārtējie poļi izturējās pret tiem naidīgi. Kad Vācijā 1871.~gadā nodibinājās vienota impērija un tās kanclers O~.f.~Bismarks tajā realizēja pret poļiem naidīgu ģermanizācijas politiku Vācijas austrumu provincēs, pasliktinājās poļu attieksme pret vietējiem vāciešiem arī t.s. ``kongresa Polijā''. Vācu ieceļotāji tika uzlūkoti kā zemei naidīgi. Tika runāts par vācu tieksmi iespiesties austrumos, realizēt ``\detxti{Drang nach Osten}'' (dziņu uz austrumiem). 1880.~gada kāds poļu žurnālists rakstīja: ``Viss jau būtu labi, ja vācieši dzīvotu poļu ciemos kā vieninieki un drīz, kā tas notiek pilsētās, kļūtu par poļiem. Taču viņi apmetas noslēgtos savos ciematos, tāpēc ir no tiem jābrīdina''. Poļi vāciešus sauca par ``kartupeļu rijējiem'', švābiem, prūšiem. Galu galā vācu vidū radās sakāmvārds ``\detxti{Solange die Welt eine Welt wird sein, wird der Deutsche nicht des Polen Bruder sein}'' (Kamēr pasaule pastāvēs, vācietis polim nekad nebūs brālis). Vācieši, acīmredzot augstākas izglītības, darba tikuma dēļ, bija veiksmīgāki. Poļi tāpēc viņus uztvēra kā visu nelaimju cēloni. Tas atspoguļojās arī literatūrā, piemēram, B.~Prusa garstāstā ``Priekšpostenis'' (``\pltxti{Placowka}'', 1885.~gads), kurā rakstnieks stāstīja par poļu sadursmēm ar vāciešiem. Tomēr daļa poļu inteliģences spēja pacelties pāri nacionālajiem aizspriedumiem. 1901.~gadā B.~Pruss avīzē ``\pltxti{Gazeta Polska}'' (``Poļu Avīze'') rakstīja: ``Ar vācu tautu mums vienmēr bijušas labākās attiecības. No viņiem pārņēmām mēs gotisko stilu būvniecībā, kokgriešanu, lielu daudzumu dažādu iekārtu, trauku, amatnieku instrumentu, lielu daudzumu zinātnes atziņu, amatu, tirdzniecību, daudzas parašas, organizācijas formas \citespace{} Nekautrēsimies no patiesības: šai cēlajai tautai mums jāpateicas par lielāko mūsu civilizācijas daļu. Kā pretdevumu tam mūsu vidū apmetušies vācu iedzīvotāji izmanto īpašas privilēģijas un laucinieki lielus atvieglojumus. Kā viņi jūtas mūsu vidū, rāda tas, ka simti tūkstošu vāciešu brīvprātīgi bez kādas piespiešanas pieņēmuši mūsu tautību~--- un pateiksim to skaļi~--- dāvinājuši mums labākos strādniekus un cienījamākos pilsoņus. Mūsu zeme ir kļuvusi viņiem laba māte un viņi~--- tai labi dēli.'' Tiesa gan, jāatzīst, ka rakstnieks vairāk rādīja ideālu, kādām būtu jābūt savstarpējām poļu un vācu attiecībām, bet dzīvē tās parasti nebūt tādas nebija.

Polijas karalistē poļu vairākuma vidū dzīvojošie ukraiņi un baltkrievi izrādīja tendenci pārorientēties no poļu vērtību sistēmas, ko viņi bija uztvēruši Žečpospolitas laikā, uz krievu kultūru, taču pilnībā to izdarīt nespēja. Uz poļu nacionālajiem uzdevumiem orientētās politiskās organizācijas un darbinieki neuzskatīja mazākumtautas par potenciāliem sabiedrotajiem, ignorēja tās, un uz lietuviešu, ukraiņu un baltkrievu nacionālo kustību skatījās kā uz ``krievu'' vai ``vācu'' iecerētu pret poļiem vērstu intrigu. Lietuviešu, ukraiņu, baltkrievu un ebreju nacionālismam bija katram sava pozīcija, kas kopīgu cīņu kā pret Krievijas impērijas, tā poļu muižniecības apspiešanu padarīja neiespējamu.

Ar Polijas dalīšanu Krievijas sastāvā nonāca daudz \strong{ebreju}. Krievijā ebreji ilgi palika izolēta sabiedriska grupa un, kā uzskatīja pazīstamais ebreju vēstures pētnieks S.~Dubnovs, par Krievijas ebrejiem var runāt tikai no XIX gs. 70.~gadiem. Līdz tam tos pareizāk bija saukt par ``Polijas ebrejiem'', jo vēsturiskās saknes viņiem saistījās ar Poliju.

Tā kā jūdaisma normas noliedza dzimstības regulēšanu un salīdzinājumā ar citiem apkārtējiem iedzīvotājiem ebrejiem bija augstākas higiēnas normas, viņu skaits strauji pieauga. 1914.~gadā Krievijas impērijā dzīvoja apmēram 5,250~miljonu ebreju (puse no visiem ebrejiem pasaulē).

Agrākajā Polijas--Lietuvas valstī ebreji bija kapitālistisko tendenču nesēji, tieši tāpēc, lai viņi veicinātu saimniecības attīstību, viņus uz Poliju aicināja karaļi. Ebreji parasti dzīvoja savos \emph{geto} (no itāļu \ittxti{ghetto}~--- pilsētas daļa, kurā viduslaikos nometināja pie svešas reliģijas piederīgus iedzīvotājus, galvenokārt ebrejus), runāja apkārtējiem nesaprotamā valodā~--- \emph{jidišā} (no vācu \detxti{jüdisch}~--- ģermāņu valodu grupas valoda, ko stipri ietekmējušas citas valodas un kurā runāja ebreji Centrālajā un Austrumu Eiropā), valkāja savu tradicionālo apģērbu, kurš faktiski bija ap XVI gadsimtu lietotais poļu apģērbs, tikai poļi to nomainīja, ebreji~--- nē. Ebreju vīrieši nēsāja bārdas un vaigu bārdas, sievietes turpretī skuva galvas un valkāja parūkas. Loti aktīvi ebreji bija tirdzniecībā~--- gan vairumā, gan mazumā, uzstājās kā starpnieki. Viņi bieži darbojās kā muižu pārvaldnieki, pilnībā aizvietojot šļahtičus, nomāja lauksaimniecības objektus, vietējos uzņēmumus, īpaši krogus, faktiski kontrolēja sīko un vidējo kredītu izsniegšanas jomu~--- t.i.~--- nodarbojās ar augļošanu. Nav brīnums, ka šāda viņu darbība izsauca apkārtējo kristīto iedzīvotāju sašutumu. Taču līdz Polijas dalīšanas laikam ebreju ietekme saimniecībā bija jūtami mazinājusies. Īpaši finanšu jomā ebreji nespēja izturēt bagāto muižnieku un klosteru konkurenci. Tomēr priekšstati par ebrejiem kā apkārtējo kristīto, sevišķi zemnieku, ekspluatatoriem izrādījās visai noturīgi.

Krievijas sastāvā ebreju stāvoklis ievērojami mainījās. Katrīna II aizliedza kā aizskarošu lietot vārdu ``žīds'', izdeva likumus, kas pielīdzināja ebrejus tiesībās citiem pilsētu iedzīvotājiem, taču realizēt tās ebrejiem neizdevās, jo tam pretojās vietējie iedzīvotāji, īpaši poļu muižnieki. Nereti viņus vienkārši nepielaida pašvaldību vēlēšanās. Varas iestādes nespēja panākt likuma izpildi. Tā pati Katrīna II, krievu tirgotāju, kuri nevēlējās ciest ebreju konkurenci, spiediena rezultātā 1791.~gadā noteica Krievijā ``ebreju apmešanās zonu''. Ebrejiem, kas kļuva par Krievijas impērijas pavalstniekiem, bija atļauts apmesties tikai viņiem atvēlētajā nometinājuma zonā gar valsts rietumu robežu. Vēlākajos gados tā tika gan paplašināta, gan sašaurināta. XIX gs. beigās tajā ietilpa 15 guberņas: Besarābijas, Viļņas, Vitebskas, Volīnijas, Grodņas, Jekaterinoslavas, Kauņas, Minskas, Mogiļevas, Podoļskas, Poltavas, Taurijas, Hersonas, Čerņigovas un Kijevas. Taču likumi bieži tika pārkāpti un lielākajās Krievijas pilsētās sāka veidoties ebreju kopienas, it sevišķi Sanktpēterburgā, Maskavā, Kijevā un Odesā. XIX otrajā pusē ebreju migrācijas apjoms ievērojami pieauga. Zināmā mērā to var izskaidrot ar pieaugošo demogrāfisko spiedienu, kā arī pastāvīgajiem modernizācijas un urbanizācijas procesiem. Uz Krievijas iekšējiem rajoniem devās arī ebreji no Polijas.

Krievijas varas iestāžu realizētajā politikā apkārtējo iedzīvotāju ``aizsardzība'' no ebreju ekspluatācijas kļuva par vienu no stūrakmeņiem. Ebreji bija ``grēkāzis'', uz ko novelt vainu par grūto dzīvi. Bieži tika izdoti rīkojumi ar aizliegumiem viņiem apmesties pilsētās un miestos, par to izsūtīšanu no laukiem, ierobežojumiem nodarboties ar tādu vai citu saimniecisko darbību. Ebreji atbilstoši Krievijas likumiem nevarēja iegādāties zemi ārpus pilsētas robežām. Tā kā ebreju skaits strauji auga, valdības politikas rezultātā drīz viņu pamatmasa nonāca nabadzībā. XX gadsimta sākumā ebreju amatnieka ienākumi vidēji bija 1,5 līdz 2~reizes zemāki kā zemnieka ienākumi (attiecīgi 150--300 un 400--500 rbļ.).

Bez tam Krievijas varas iestādes par savu uzdevumu uzskatīja cīnīties pret ``ebreju fanātismu'', kurš izpaužoties sevi uzskatot par Dieva izredzētu tautu, nicinot apkārt dzīvojošos kristiešus, jo ebreji dzīvojot noslēgti no tiem, neparādot lojalitāti pret pastāvošo valsts varu, savu reliģisko normu izpildi uzskatot par svarīgāku nekā valsts likumu ievērošanu! (Tiešām, draudzības saites ebrejiem ar kristiešiem veidojās reti. Piemēram, ebrejs nevarēja doties viesos pie kristieša, jo tā mājā ēdienu gatavoja, neievērojot ebreju reliģiskās prasības.) Savus uzplūdus cīņa pret ``ebreju fanātismu'' sasniedza Nikolaja I valdīšanas laikā. Krievijā no 1827.~gadā uz ebrejiem tika attiecināts arī rekrūšu došanas pienākums. Pie tam viņiem iesaukšanas normas bija noteiktas trīs reizes lielākas nekā no citām iedzīvotāju grupām. Varas iestādes karaklausības ieviešanu ebrejiem un viņu dienestu armijā izmantoja, lai piespiestu ebrejus kristīties. Tikai 1874.~gadā Krievijā tika ieviesta vispārējā kara klausība.

Krievu liberāļi uzskatīja, ka ebrejiem vispirms ir jādod tādas pat tiesības kā visiem citiem, tad viņi ``labosies'', turpretī konservatori domāja, ka vispirms ir jāpanāk ebreju ``labošanās'' un tikai tad varēs viņiem dot vienlīdzīgas tiesības. Asimilācijas tendences XIX gadsimtā aptvēra arī ebrejus, kā rezultātā radās ebreju izcelsmes, bet kristīti ārsti un advokāti, žurnālisti, augstskolu profesori, aktieri un režisori. Tas savukārt saasināja citu tautību nepatiku pret viņiem. Ja agrāk krieviem, poļiem u.c. nācās konkurēt ar ebreju veikalniekiem, krodzniekiem vai rentniekiem, tagad jutās apdraudēti arī brīvo profesiju pārstāvji. Tāpēc tika aizmirsts, ka juridiski tie vairs nav ebreji, bet kristieši.

Ebrejiem Krievijā bija piešķirtas zināmas pašpārvaldes tiesības, kopiena jeb \emph{kagals} (\rutxti{кахал}), tās vēlētā pārstāvniecība ievāca kopīgus nodokļus, t.s. kastes ievākumu (\rutxti{коробочный сбор}), jo savāktos līdzekļus glabāja kastē, ko uzraudzīja kopienas pilnvarots cilvēks), regulēja saimniecisko darbību, lai ierobežotu savstarpējo konkurenci, centās ierobežot patēriņu, greznību.

Varas iestādes centās ierobežot ebreju izglītošanās iespējas, taču par spīti visiem šķēršļiem, viņi pēc tās tiecās arvien vairāk. Aleksandra II laikā (60.--70.~gados) tika izdoti vairāki likumi, ar kuriem ārpus apmešanās zonas atļāva dzīvot jūdu ticības iedzīvotājiem ar augstāko izglītību, kā arī dantistiem, farmaceitiem, feldšeriem, pēc tam amatniekiem, bijušajiem karavīriem. Ebreji ieguva tiesības iestāties valsts dienestā, piedalīties pilsētu un zemstu pašvaldību darbā. Tomēr nelīdztiesība, smagie dzīves apstākļi noveda pie tā, ka jau 70.~gados arvien vairāk ebreju iesaistījās revolucionārajā kustībā. Ja 1871.--1873.~par politiskām lietām izmeklēšanā nonākošo ebreju skaits sastādīja 4--5\%, kas atbilda viņu skaitam iedzīvotāju vidū, tad XIX gs. 80.~gadu beigās viņi jau sastādīja 35--40\% visu revolucionāru.

Lūzuma moments iestājās 1881.~gadā, kad sakarā ar cara Aleksandra II nogalināšanu (1881.~gada 1.~martā) Krievijā izveidojās nestabila politiskā situācija. Izplatījās baumas, ka cara slepkavas esot ebreji un varas iestādes devušas netiešus norādījumus visur sarīkot ebreju grautiņus. Ebreju iedzīvotāju nelīdztiesība, antisemītiski noskaņojumi un daudzu iedzīvotāju neapmierinātība ar ebreju konkurenci veicināja šādu baumu plašu izplatību. Vietējās administrācijas, policijas un karaspēka pasivitāte it kā apliecināja baumu pamatotību. Pirmais tāds grautiņš notika Jeļizavetgradā (tagad Kropivņicka (\uktxti{Кропивницький})) Ukrainas dienvidos 15-16.~aprīlī, pēc tam turpinājās citur Krievijas impērijas rietumu un Dienvidrietumu apgabalos. 13.~decembrī notika ebreju grautiņš Varšavā, ko bija inspirējušas Krievijas varas iestādes. Grautiņā piedalījās arī poļu sabiedrības apakšslāņi. Inteliģence to nosodīja, katoļu garīdznieki devās ielās, cenšoties pierunāt grautiņu dalībniekus izklīst. Tikai trešajā dienā karaspēks apspieda nekārtības. Bija sagrauts ap pusotra tūkstoša ebreju dzīvokļu, citu telpu, ievainoti 24 cilvēki. Atsevišķas vietās ebreju grautiņi turpinājās līdz pat 1884.~gadam. Valdība gan mēģināja tos novērst, gan arī reizē centās atbildību par tiem uzkraut pašiem ebrejiem.

1887.~gadā tika ieviesta t.s. ``procentu norma'' viņu iestājai augstākajās un vidējās mācību iestādēs~--- 10\% ebreju apmešanās zonā, 5\%~--- ārpus tās, 3\%~--- galvaspilsētās. Atsevišķās vietās attieksmi pret ebrejiem noteica augstākās varas personas. Piemēram, kad par Maskavas ģenerālgubernatoru 1891.~gadā kļuva lielkņazs Sergejs Aleksandrovčs, no pilsētas un guberņas izsūtīja ap 20~tūkstošu ebreju amatnieku un pat Nikolaja I laika karavīru.

Kā raksta poļu vēsturnieks J.~Tazbirs, poļu--ebreju attiecības saasinājās XIX gadsimta beigās, jo poļu apdzīvotajās teritorijās ieceļoja daudz t.s. ļitvaku (no jidiša \hetxti{ליטװאַק}, lietuviešu~--- \lttxti{litvakai}~--- ebreji, kuri dzīvoja Baltkrievijas, Lietuvas, Latvijas, dažos Krievijas un Polijas rajonos). Viņi tikuši izsūtīti no Krievijas iekšienes, slikti zinājuši poļu valodu un atradušies krievu kultūras lokā. Poļi tos uztvēruši kā rusifikācijas īstenotājus, kuri bremzējuši asimilācijas procesus, kas norisinājušies Polijas ebreju vidū. Tas novedis pie tā, ka ne tikai poļu tirgotāju un amatnieku, bet arī inteliģentu vidū atbalsi guvuši antisemītiski lozungi. Domājams, ka pret šo vērtējumu jāizturas kritiski. Vajadzēja pacensties, lai ebrejus, cietušus no antisemītisma izpausmēm Krievijā, attēlotu kā rusifikācijas īstenotājus. Acīmredzot darbojās tā pati vien konkurence starp poļu un ebreju amatniekiem, tirgotājiem, brīvo profesiju pārstāvjiem. Protams, poļu antisemīti to nevarēja atzīt.

Nabadzības, nevienlīdzības un arī grautiņu rezultātā sākās plašs ebreju emigrācijas vilnis, galvenokārt uz ASV. Līdz Pirmajam pasaules karam no Krievijas impērijas emigrēja vairāk nekā 2~miljoni ebreju. Grautiņu radītais šoks, ierobežojumi izsauca ideju par ebreju atgriešanos Palestīnā. Gan visā Krievijas impērijā, gan Polijas karalistē atbalstu guva cionisma kustība, kas bija vērsta uz ebreju valsts Izraēlas izveidi un nostiprināšanu un ebreju nācijas attīstību.

Daži teikumi jāpasaka arī par \strong{ārpus etniskās Polijas} palikušajiem bijušajiem Polijas-Lietuvas apgabaliem.

\strong{Lietuviešu} apdzīvotās teritorijas pēc Polijas trešās dalīšanas nonāca Krievijas impērijas sastāvā. Vārdu ``Lietuva'' tagadējās Lietuvas valsts teritorijas apzīmēšanai sāka lietot tikai XIX gadsimta otrajā pusē. Pēc nošķiršanas no Žečpospolitas lietuvieši daudz mazākā mērā tika pakļauti polonizācijai, viņos sāka veidoties sava nacionālā apziņa. Arī \strong{ukraiņu} nacionālajā inteliģencē Krievijā XIX gadsimta vidū veidojās sava nacionālā pašapziņa. Vēsturnieks N.~Kostomarovs, kurš daudz nodarbojās ar Ukrainas vēsturi, publicēja rakstu ``Divas krievu tautības'' (``\rutxti{Две русские народности}''), kura nosaukums jau izsaka tā saturu. 1905.~gadā Krievijas Zinātņu akadēmija oficiāli atzina ukraiņu valodas esamību. Trešās ``krievu tautības''~--- \strong{baltkrievu} nacionālās pašapziņas veidošanās un viņu kā atsevišķas ``tautības'' atzīšana norisa krietni lēnāk. Baltkrievu iedzīvotāju vairākums neapzinājās savu nacionālo identitāti, uz jautājumiem par to deva sekojošas atbildes: ``tautība~--- ``\pltxti{tutejšij}'', ticība~--- ``\pltxti{viera polska}'', valoda~--- ``\pltxti{jazik prostij}''. Vietējās baltkrievu apdzīvoto apgabalu inteliģences prātos cīkstējās poļu un t.s. ``rietum-krieviskuma'' (``\rutxti{западно-руссизм}'') pašapziņa. Saskaņā ar otro baltkrievi bija tikai viena no krievu tautas etnogrāfiskajām grupām. Baltkrievu tautas patstāvības ideja ieguva skanējumu tikai XIX gadsimta 80.~gados Pēterburgā, kur bija studenti no baltkrievu novadiem. Taču vēl vairākus gadu desmitus šīs idejas aizstāvjiem nācās pierādīt tās tiesības uz eksistenci. Un tomēr jēdziens ``baltkrievi'', lai arī vēl tīri etnogrāfiskā nozīmē, pakāpeniski iespiedās masu apziņā. 1897.~gada Krievijas iedzīvotāju skaitīšanā 74\% no tās Ziemeļ-Rietumu apgabala iedzīvotājiem par savu dzimto atzina baltkrievu valodu.

Tādejādi tajos bijušajos Žečpospolitas, bet vēlāk Krievijas apgabalos, kuri neietilpa Polijas karalistē, kaut dažādiem tempiem, bet tomēr attīstījās agrāk poļiem pakļauto tautu pašapziņa, kas arvien vairāk apdraudēja poļu cerības apvienot šīs teritorijas atdzimušā Žežpospolitā. Ja XVIII gadsimta beigās t.s. [Dņepras] Labā krasta Ukrainā poļi sastādīja valdošā slāņa vairākumu, tad XIX gadsimtā, īpaši pēc abu poļu sacelšanos sakāvēm, viņi bija zaudējuši savas pozīcijas, ar nostaļģiju atceroties Žečpospolitas ``labos laikus''. Taču kādu nopietnu cerību tos atgriezt viņiem nebija. Ārpus Polijas karalistes esošajos baltkrievu apgabalos dzīvojošo poļu vidū pastāvēja divas tendences. Vieni pirmajā vietā izvirzīja visa novada un tajā dzīvojošo tautu intereses, etniskās nācijas vietā izvirzot nācijas kā politiskas vienības kategoriju, uzsverot lietuviešu un baltkrievu apdzīvoto zemju atšķirību no etniskās Polijas. Otrie uzsvēra poļu nacionālo uzdevumu prioritāti.

XIX gadsimta beigās Polijā sāka veidoties \strong{politiskās partijas} mūsdienu izpratnē (ar vairāk vai mazāk izteiktām partijas pazīmēm~--- savu programmu, statūtiem, vadošajām iestādēm un vietējām organizācijām, savu propagandas aparātu un piekritēju kustību).

Līdz ar kapitālisma attīstību Polijā iezīmējās galvenie politiskās attīstības virzieni, atspoguļojoši ne tikai šķirisko diferenciāciju, raksturīgu tirgus saimniecībai, bet arī sabiedrības uzbūves sarežģītību, daudzslāņainību, kurā saglabājās vēl pirmskapitālistiskās struktūras. Neizdzēšamu iespaidu uz Polijā radušos politisko partiju un sabiedrisko strāvojumu veidolu atstāja nacionālā jautājuma specifika, kā arī zemi sadalījušo monarhiju iekšējā kārtība, kura pilnībā vai daļēji liedza poļu partijām likumdošanas un citas parlamentārās un pat vispār legālās darbības iespējas.

Pārmaiņas ekonomiskajā un sabiedriskajā dzīvē, galvenokārt no Rietumiem nākošu ideju izplatība, tai skaitā augstāk izklāstītās ``pozitīvisma'' un ``organiskā darba'' idejas, ietekmēja Polijā divu politisku strāvojumu veidošanos, kam, tāpat kā citās Eiropas valstīs, bija visievērojamākā loma mūsdienu vēsturē. Tie bija nacionālisms un sociālisms. Abi aptvēra visus trīs poļu reģionus dažādās valstīs, saglabājot savstarpējos kā intelektuālos, tā zināmā mērā arī organizatoriskos sakarus.

60.~gadu beigās~--- 70.~gadu sākumā notika pirmās poļu \strong{strādnieku šķiras} uzstāšanās. Rūpniecības centri Lodza, Varšava kļuva par strādnieku streiku arēnu. Polijas agrīnās strādnieku kustības (kā plebejiski-nacionālistiskās, tā arī revolucionāri-sociālistiskās, jeb t.s. ``proletāriskās'') kodolu veidoja bijušie amatnieki un no izputējušajiem šļahtičiem nākušie strādnieki. Polijā taču dzimtbūšana bija atcelta ātrāk nekā Krievijā un kapitālisms attīstījās straujāk. Jau XIX gs. vidū eksistēja ievērojams bezzemnieku slānis, kuri, dodoties uz pilsētu, papildināja fabriku strādnieku skaitu. No šļahtičiem nākušie strādnieki saglabāja kādreiz līdz smieklīgumam nonākošu šļahtičiem raksturīgu savas pašapziņas sajūtu, taču tā bieži arī veda viņus pie varonīgas rīcības. Šie strādnieki nevēlējās ciest pazemojumus no ekspluatatoriem, bija gatavi dumpoties. Poļu strādniekos deklasētie šļahtiči ienesa iepriekšējā posma atbrīvošanās cīņas tradīcijas, it īpaši nacionālajā jautājumā.

\strong{Poļu sociālistiskais virziens} veidojās ciešā saistībā ar krievu narodņikiem, iegūstot piekritējus poļu studentu vidū Krievijas augstskolās. Tai pat laikā Varšavas studentu rokās nonāca sociālistiskās brošūras no Rietumiem. Poļu sociālistiskā kustība izauga uz Rietumu un pirmkārt Vācijas sociālistiskās kustības pieredzes pamata. Šis sociālisms jau bija orientēts uz strādniekiem. Pirmo poļu sociālistu darbība izvērsās augsnē, ko jau bija sagatavojuši dažādi poļu sabiedriskās kustības virzieni 1830.--1831.~gada un 1863.--1864.~gada nacionālo sacelšanos laikmetā. Kā uzsver poļu vēsturnieki, poļu sociālisti smēla iedvesmu arī ``pozitīvismā''.

Strādnieku kustība Polijā attīstījās jau no XIX gs. 70.~gadiem, kad plaši streiki notika Varšavā, Lodzā u.c. pilsētās. Pirmais poļu studentu sociālistiskais pulciņš izveidojās ārpus Polijas~--- Pēterburgā 1874.~gadā. 1876.~gadā arī Varšavā tika organizēts pirmais ``tīri sociālistiskais'' pulciņš. 1877.~gadā tika izdota pirmā sociālistiskā brošūra poļu valodā, paredzēta izplatīšanai strādnieku vidū. 1878.~gadā notika pirmās strādnieku sapulces, kas lika pamatus vēlāk dibinātajām strādnieku organizācijām, sākās organizācijas statūtu un programmas projekta izstrāde. Krievijas narodņiku ietekmē pirmajos poļu sociālistiskajos pulciņos bija izplatīti teroristiski noskaņojumi. Tika gatavots atentāts pret ģenerālgubernatoru grāfu P.~Kocebu, organizācijas vadītājiem ar grūtībām izdevās atturēt aktīvistus no šī soļa. 1880.~gadā strādnieki nogalināja organizācijā iekļuvušu provokatoru. Pielietoja arī ``papīra'' (ekonomisko) teroru, kad rūpniecības uzņēmumu īpašniekiem nosūtīja prasības paaugstināt darba algu, uzlabot izturēšanos pret strādniekiem, pretējā gadījumā draudot iznīcināt kā uzņēmumus, tā pašus tā īpašniekus. Tomēr izskanēja arī marksistiski secinājumi. Tā, 1880.~gada 29.~novembrī starptautiskā mītiņā Ženēvā L.~Variņskis kā pirmais poļu sociālists atzīmēja, ka ``Krievija vairs nav reakcijas balsts'', ka sociālistu priekšā stāv nevis slāvu federācijas ideja vai jautājums par poļu valsts robežām, bet cīņa par ``visas pasaules proletariātam kopīgu ekspluatējamo solidaritāti cīņā pret apspiedējiem.'' 1881.~gada oktobrī starptautiskajā sociālistu kongresā Kurā (\svtxti{Kur}, Šveice) L.~Variņskis paziņoja, ka poļu proletariāta uzdevums ir nevis nacionālas sacelšanās gatavošana, bet sociāla apvērsuma sagatavošanā, kura gaitā tiks likvidēta arī nacionālā apspiestība. Tiesa, L.~Variņskis uzskatīja, ka sociālismam vispār ``nav nekā kopīga ar nacionālo jautājumu\dots{}''

1881.~gadā, protestējot pret bezdarbu, Varšavas strādnieki izgāja pirmajā manifestācijā. 1882.~gadā L.~Variņskis nodibināja pirmo Starptautisko sociāli revolucionāro partiju \strong{``Proletariāts''} (\pltxti{Międzynarodowa Socjalno-Rewolucyjna Partia ``Proletariat''}). Partijai izveidojās piekritēju pulciņi arī Pozenes un Krakovas novados, tai izdevās izdot pirmo nelegālo avīze, kas tā arī saucās~--- ``\pltxti{Proletariat}'', arī virkni brošūru un uzsaukumu. Izveidojās sakari starp partiju un krievu revolucionāru~--- narodovoļcu (``\rutxti{Народная воля}''~--- ``Tautas griba'') organizāciju. Partijas programmā tika atzīmēta strādnieku un zemnieku interešu sakritība, uzsvērta nepieciešamība solidarizēties ar visiem apspiestajiem, neatkarīgi no viņu tautības. Nacionālais jautājums programmā bija skatīts visai šauri~--- tikai kā visu tautību un konfesiju piederīgo tiesību nolīdzināšana.

1883--1885.~gadā notika plaši aresti, četri revolucionāri 1886.~gadā tika pakārti, L.~Variņski ieslodzīja Šliselburgas cietoksnī, kur viņš nomira. Partija zaudēja nozīmību. 1888.~gadā to mēģināja atjaunot. Jaunā partija iegāja vēsturē ar nosaukumu \strong{``Proletariāts II''} (\pltxti{Socjalno-Rewolucyjna Partia ``Proletariat'' II)}. Tam pievienojās arī poļu sociālistu grupa Rīgā. Pilsētā nelegālā tipogrāfijā tika drukātas sociālistiska satura brošūras, kuras izmatoja propagandai ne tikai Polijas karalistē, bet arī citur Krievijā. Taču atšķirībā no L.~Variņska paaudzes sociālistiem ``Proletariāts II'' darbinieku vidū valdošais bija uzskats, ka reiz jau patvaldība Krievijā ir tikusi galā ar revolucionāro spēku spiedienu un pārvarējusi pagaidu krīzi, poļu sociālistu uzdevums ir nevis tieša sociālistiskās revolūcijas sagatavošana, bet cīņa par Konstitūciju Krievijā un Polijas karalistes autonomiju Krievijas valsts sastāvā. Par visiedarbīgāko politiskās cīņas līdzekli ``Proletariāts II'' uzskatīja teroru, jo neredzēja citu iespēju kā izsaukt masu kustību valdošā despotisma apstākļos. Terora koncepcija bija pārņemta no krievu narodovoļciem. Terora lomas izcelšana neļāva pareizi novērtēt ekonomiskās cīņas nepieciešamību, lai iesaistītu strādniekus masu kustībā. Kaut kopumā šai laikā sociālistu idejiski-politiskā ietekme uz Varšavas sabiedrisko dzīvi bija neliela, ``Proletariāts II'' tomēr nodrošināja poļu sociālistiskās kustības nepārtrauktību gados pēc ``Proletariāta'' sagrāves.

1889.~gadā krāsotājs J.~Marhļevskis, žurnālists Ā.~Varskis, virkne Varšavas studentu un strādnieku nodibināja Poļu strādnieku savienību (\pltxti{Związku Robotników Polskich}). Par savas darbības paraugu tā uzskatīja Vācijas sociāldemokrātiju. No ``Proletariāta II'' savienība atšķīrās ar uzsvērti negatīvu attieksmi pret terora taktiku visās tā izpausmēs. 1890.~gadā Varšavas strādnieki pirmie Krievijas impērijā atzīmēja darbaļaužu starptautiskās solidaritātes dienu 1.~maiju, 1892.~gadā Lodzā notika pirmais Polijā vispārējais politiskais streiks. Tomēr tikai 1892.--1893.~gadā radās sociālistiskas politiskas organizācijas šā vārda pilnā nozīmē.

Sociāli politisko procesu specifika 90.~gados, liberālā un demokrātiskā virziena vājums poļu pilsoniskajās organizācijās, nacionālā jautājuma asums ietekmēja sociālistiskās kustības attīstību. Tajā izveidojās divas strāvas. Viena pirmajā vietā izvirzīja šķiriskos, otrā~--- nacionālos uzdevumus.

1893.~gadā, apvienojoties ``Proletariātam II'' un Poļu strādnieku savienībai, nodibinājās Polijas karalistes sociāldemokrātija (\pltxti{Socjaldemokracja Królestwa Polskiego}), kura 1900.~gadā tika pārdēvēta par \strong{Polijas karalistes un Lietuvas sociāldemokrātiju} (PKunLSD, \pltxti{Socialdemokracja Krolestwa Polskiego i Litwy}). Tās vadītāji bija jau minētais J.~Marhļevskis un R.~Luksemburga. J.~Marhļevskis 1896.~gadā beidza Cīrihes universitāti, aizstāvēja disertāciju ``Fiziokrātisms vecajā Polijā'', kurā pirmo reizi poļu historiogrāfijā demonstrēja marksistisku pieeju tā laika poļu sabiedrības problēmām. R.~Luksemburga arī studēja Cīrihē, aizstāvēja disertāciju ``Rūpniecības attīstība Polijā''. Vēl partijā darbojās tādi pazīstami revolucionāri kā J.~Tiška, F.~Dzeržinskis u.c.

Minētie sociālisti stadijā, kad tiem vēl nebija plašas vispārdemokrātiskas platformas, nespēja izstrādāt arī pozitīvu programmu nacionālajā jautājumā. Partija pirmajā vietā stādīja sociālos uzdevumus, bet nacionālas valsts izveidi uzskatīja par šķērsli ceļā uz sociālismu. Partijas rindās pārsvaru guva uzskati, ka atsevišķas Polijas daļas jau ir tik cieši saistītās ar Poliju sadalījušajām valstīm, ka vienotas Polijas valsts izveide nav vairs iespējama. Idejiski partija lielā mērā atradās R.~Luksemburgas ietekmē, kura 90.~gados noformulēja teoriju par Polijas ekonomisko attīstību kā poļu sociāldemokrātu politiskās programmas pamatu. Viņa aizstāvēja uzskatu, ka rūpniecība Polijas karalistē pastāv un attīstās tikai pateicoties politiskajai piederībai Krievijai, kapitālisma attīstības tendence ir arvien ciešāka dažādu teritoriju (šai gadījumā Polijas karalistes un Krievijas) ekonomiska saplūšana, bet nacionālu valstu (piemēram, Polijas) izveide ir pretrunā ar kapitālistiskās sabiedrības attīstības likumiem. ``Polijas neatkarība,~--- rakstīja R.~Luksemburga,~--- bija līdzkapitālistiskā, šļahtiskā, naturālās saimniecības perioda ideāls''. Poļu strādnieki, pēc viņas domām, nedrīkstēja iekļaut savās prasībās Polijas neatkarību, jo tad viņi kļūtu par ekonomiskā progresa pretiniekiem. Tāpat kā kapitālisti, arī strādnieki varēja apvienoties ārpus valstu un nāciju robežām, strādnieki bija tendēti uz šķiru, nevis nacionālo atbrīvošanas kustību. Nacionālo problēmu loku viņa ierobežoja tikai ar kultūras sfēru. Tā šaura, nedialektiska pieeja ekonomiskās attīstības likumu traktēšanai noveda šo marksisma teorētiķi pie nacionālā jautājuma nepamatotas ignorēšanas. R.~Luksemburga un viņas sekotāji, vedot principiālu cīņu pret nacionālismu, pieļāva kļūdu, uzstājoties pret neatkarības lozungu, kurš Polijas politiskās dzīves konkrētajos apstākļos XIX/XX gadsimta mijā viņai likās par poļu-krievu revolucionāras savienības, tautu brīvības un sociālisma pretstatu. Kļūdaina bija prasība, lai krievu sociāldemokrāti atteiktos no prasības par nāciju pašnoteikšanās tiesībām, arī poļu tiesībām uz valstisku neatkarību. Aizraušanās ar cīņu pret nacionālismu noveda pie strādnieku šķiras cīņas internacionālā un nacionālā aspektu pareiza samēra zaudēšanas. R.~Luksemburgas principiāla kļūda bija arī Polijas karalistes rūpniecības pastāvēšanas ārpus Krievijas iespēju noliegšana.

Kā pirmais R.~Luksemburgas kļūdas vēsturiskajos uzskatos kritizēja jau bijušais partijas ``Proletariāts'' darbinieks lietuviešu revolucionārs L.~Janovičs. Citi marksisti, tai skaitā V.~Ļeņins, šai laikā atzina, ka labākos apstākļus kapitālisma attīstībai rada tieši nacionāla valsts. Carisma ekonomiskā politika virknē gadījumu traucēja Polijas rūpniecības attīstību, arī savu laiku pārdzīvojusī politiskā iekārta bremzēja poļu tautas progresu. Taču PKuLSD atteicās no nacionālās neatkarības lozunga izvirzīšanas, uzskatot, ka poļu proletariātam ir jāapvieno savi spēki ar citu zemju proletariātu un šī kopīgi realizētā sociālā revolūcija atrisinās visas strādnieku šķiras problēmas. Ar to tika noliegts nāciju pašnoteikšanās demokrātiskais saturs. Galvenokārt jautājumā par nāciju pašnoteikšanās tiesībām radās domstarpības ar krievu sociāldemokrātiem, tāpēc arī PKunLSD atteicās apvienoties ar Krievijas Sociāldemokrātisko strādnieku partiju. 1893.--1895.~gadā Polijas karalistes sociāldemokrātija cieta smagus zaudējumus arestos.

Kā pretspēks internacionāli orientētām sociālistu organizācijām no dažādām sociālistiskām grupām 1892.~gadā Parīzē nodibinājās Poļu sociālistu ārzemju savienība (\pltxti{Związek Zagraniczny Socjalistów Polskich}), kura pasludināja uzticību 1863.~gada sacelšanās tradīcijām, stādīja uzdevumu organizēt sacelšanos pret Krieviju. Attiecībā pret Vāciju un Austroungāriju tāds uzdevums netika stādīts. No savienības 1893.~gadā jau noformējās \strong{Polijas Sociālistiskā partija} (\pltxti{Polska Partia Socjalistyczna}) jeb \strong{PPS}. Tā lielu uzmanību veltīja nacionālajam jautājumam. Rietumu sociālistu izvirzītie uzdevumi partijas programmā bija apvienoti ar tieksmi pēc nacionālās neatkarības. Partijas radītāji ticēja, ka kapitālisma progress radīs spēcīgu strādnieku šķiru, kura varēs cīnīties kā par privātīpašuma likvidāciju, tā par nacionālo neatkarību. PPS nepieciešamo sociālo reformu īstenošanu saistīja ar neatkarīgas demokrātiskas Polijas Republikas izveidi. Pēc PPS ideologu uzskatiem par neatkarīgu Polijas Republiku bija jācīnās vēl līdz sociālistiskajai revolūcijai. Neatkarīgajai valstij bija jārada poļu strādnieku šķirai labvēlīgi apstākļi nākotnes cīņai par sociālistiskas iekārtas nodibināšanu. Polijas neatkarības idejas aizstāvji balstījās uz pesimistisku Krievijas revolucionārās kustības perspektīvu novērtējumu. Viņu žurnāls ``\pltxti{Przedswit}'' (``Rītausma'') 1894.~gadā rakstīja, ka ``nopietnai revolucionārai kustībai Krievijā izredžu nav''. Ja krievu sociāldemokrāti dažādu tautu nacionālajā kustībā saskatīja palīgspēku savu sociālistisko mērķu labā, šie poļu sociālisti Krievijas revolucionārajā kustībā redzēja līdzekli, izmantojamu savu nacionālo mērķu sasniegšanai.

1896.~gadā PPS uzstāja, lai II Internacionāles IV kongress Londonā iekļautu Polijas neatkarības prasību savā programmā. PPS atsaucās uz marksisma pamatlicēju atbalstu Polijas neatkarībai. Tomēr kongress šai prasībai nepakļāvās, reizē atzīstot visu nāciju pašnoteikšanās tiesības, bet arī aicinot visu nāciju proletariātu uz internacionālu vienotību. Taču PPS turpināja aizstāvēt Polijas izņēmuma stāvokli, 1897.~gadā savā IV kongresā izstrādāja līguma noteikumus ar Krievijas sociālistisko partiju (kuras nodibināšana bija paredzama nākotnē), kuru būtība bija tā, ka Krievijas sociālistiem bija jāizvirza Polijas neatkarības prasība un aktīvi jāaģitē krievu vidū par šādas prasības izpildi. Faktiski partijas pozīcijā bija guvuši izpausmi jau novecojušie uzskati par poļu jautājuma izcilo, izņēmuma lomu starptautiskajā revolucionārajā kustībā. Kad nākamajā~--- 1898.~gadā notika Krievijas Sociāldemokrātiskās Strādnieku partijas dibināšanas kongress, PSS pārstāvji uz to netika lūgti, jo krievu sociāldemokrātu grupas, kas organizēja kongresu, PPS neuzskatīja par sociāldemokrātisku.

Der sīkāk pateikt par vienu no partijas vadītājiem \strong{J.~Pilsudski}. J.~Pilsudskis bija dzimis senā, bet nabadzīgā šļahtiču (viņa mātes pūrā gan bija vairāk nekā 10~000 ha zemes, bet tēva novatoriskie saimnieciskie pasākumi praksē visi beidzās neveiksmīgi un par iztikas avotu kļuva banku kredīti) ģimenē Zalavā (toreiz Krievijā, tagad Lietuvā), netālu no Viļņas. Ģimenē valdīja poļu sacelšanos varoņu kults. Jau tajos gados jaunā šļahtiča apziņā notika carisma pielīdzināšana krievu tautai, kurai tika piestādīts rēķins par patvaldības pāridarījumiem poļu tautai. Vēlāk viņš rakstīja: ``Tai laikā visi mani sapņi koncentrējās ap sacelšanos un bruņotu cīņu pret moskaļiem, kurus es nīdu ar visu dvēseli, uzskatot katru no viņiem par nelieti un zagli.'' Ar laiku antikrieviskums kļuva par vienu no viņa politiskās programmas galvenajiem postulātiem. 1885.~gadā Juzefs beidza Viļņas ģimnāziju, uzsāka medicīnas studijas Harkovas universitātē, piedalījās kādā studentu politiskā demonstrācijā, par ko dažas dienas pavadīja arestā. 1887.~gadā, 20 gadu vecumā, saistībā ar cara Aleksandra III slepkavības mēģinājumu (Juzefs sava vecākā brāļa Broņislava \emph{uzdevumā pārnesa gatavojamās bumbas sastāvdaļas, pats nezinot par to}) jaunieti arestēja un izsūtīja uz 5 gadiem uz Austrumsibīriju. (Starp citu krievu politologs S.~Čerņahovskis ir pievērsis uzmanību tam, ka gan V.~Ļeņina brālis, gan J.~Pilsudska brālis piedalījās atentātā pret Aleksandru III, bet neviens nepievēršot uzmanīgu tam, ka V.~Ļeņina brālis Aleksandrs tika sodīts ar nāvi, bet J.~Pilsudska brālis tikai izsūtīts. Papildinot teikto, gan jāpasaka, ka Juzefa Pilsudska brālim Broņislavam, tāpat kā A.~Uļjanovam un virknei citu, tika piespriests nāves sods, bet imperators, apstiprinot spriedumu, nomainīja Broņislavam nāves sodu ar 15 gadiem katorgas darbos Sahalīnā.)

Tuvāka iepazīšanās ar krievu tautas pārstāvjiem neizsauca J.~Pilsudskī cieņas un labvēlības jūtas pret tiem. Lūk, viņa vārdi: ``Visi viņi ir vairāk vai mazāk slēpti imperiālisti. Viņu vidū ir daudz anarhistu, taču, dīvaina lieta, republikāņus viņu vidū es nemaz nesastapu.'' Spriežot pēc šī izteiciena, J.~Pilsudska sociāli politiskā orientācija bija visai miglaina. Grūti saprast kā var būt par anarhistu, kurš neatzīst kādu personību, grupu vai šķiru varu pār tautu, tai pat laikā esot imperiālistam. Pēc atgriešanās Viļņā 1892.~gadā viņš loloja plānus patstāvīgi nokārtot universitātes gala pārbaudījumus juridiskajās zinātnēs (Tādi precedenti bija. Piemēram, 1891.~gadā pašmācības ceļā jurista licenci ieguva V.~Uļjanovs-Ļeņins, taču J.~Pilsudskim pietrūka uzņēmības savus plānus īstenot.) Šajā laikā J.~Pilsudskis iepazinās ar savu nākamo dzīves biedri M.~Juškeviču, par kuras labvēlību viņam nācās cīnīties arī pret R.~Dmovski, vēlāko nacionāldemokrātiskās partijas līderi. Sāncensība, kura vēlāk apauga leģendām, iespējams, ietekmēja arī abu vēlākās politiskās attiecības.

Veidojoties PPS, tajā iekļāvās arī J.~Pilsudskis. Publicistikā ir izteikts viedoklis, ka notika tas nevis viņa pasaules uzskata attīstības, bet gan aizraušanās ar savu nākamo dzīves biedri rezultātā. (Var atzīmēt, ka ļoti daudzi revolucionāri par tādiem ir kļuvuši sabiedrisku un personīgu cēloņu mijiedarbības rezultātā.) 1899.~gadā viņš salaulājās ar M.~Juškeviču. Tā kā viņa bija šķirtene, laulības nevarēja notikt katoļu baznīcā. Abi nolēma pieņemt protestantismu. Tikai 1916.~gadā J.~Pilsudskis atgriezās katoļu baznīcā. (Kad viņš jau neatkarīgajā Polijā ieņēma augstākos valsts amatus, par šo atteikšanās faktu no katoļu reliģijas bija nolaists klusēšanas plīvurs, jo katoliskajā poļu sabiedrībā tas nevarēja stiprināt politiķa popularitāti. Tikai 1934.~gadā fakts kļuva zināms atklātībai.) 1900.~gadā J.~Pilsudskis atkal tika arestēts, ieslodzīts Varšavas Citadelē. Viņš simulēja garīgi slimo, tika pārvests uz psihiatrisko slimnīcu Pēterburgā, no kuras izbēga ar ārsta, PPS biedra V.~Mazurkeviča palīdzību.

Jau no darbības partijā sākotnes J.~Pilsudskis aktīvi propagandēja savus uzskatus. Tā 1893.~gadā žurnālā ``\pltxti{Przedswit}'' viņš pārmeta krievu sociālistiem to, ka viņi neuzskatīja Lietuvu, Baltkrieviju un Ukrainu par Polijas daļu, apgalvoja, ka krievu sociālisti nevēlas vairāk attīstītās Polijas atdalīšanos no Krievijas, jo ``tas nobremzētu sociālistiskās kustības attīstību Krievijā. Tika apgalvots, ka tur, kur darbojas PPS grupas (bet par savu darbības sfēru PPS uzskatīja visu XVIII gadsimta Žečpospolitas teritoriju), krievu revolucionāri drīkst darboties tikai ar poļu grupu piekrišanu, un ``to kontrolē''.

J.~Pilsudskis aktīvi propagandēja cīņu pret carisko Krieviju, daudz spēku veltot nelegāla laikraksta ``\pltxti{Robotnik}'' (``Strādnieks'') izdošanai. 1895.~gadā viņš izteica pārliecību, ka ``poļu strādnieku šķira, kura jau tagad sagādā carismam ne mazums raižu, vedīs aiz sevis cīņā citu paverdzināto tautu strādājošo masas un, gūstot atbalstu no pašas Krievijas revolucionārās kustības, vedīs tās uz uzvaru, kas nodrošinās brīvību un atbrīvošanos ne tikai poļiem, bet visiem pārējiem, kuri cietuši no carisma.'' Turpmāk gan viņa uzskatu sistēmā uzsvari pārvietojās. Poļu strādnieku šķiras intereses tika atbīdītas nacionālo interešu vārdā, proletariāts viņam bija tikai līdzeklis poļu nācijas atbrīvošanai. Lai sasniegtu šo mērķi, gūtu tam sabiedrības atbalstu, J.~Pilsudskis un viņa sabiedrotie bija spiesti izvirzīt sociālistiskus mērķus, kaut paši nopietni centās panākt tikai nacionālas buržuāziski demokrātiskas republikas izveidi. Krievijas revolucionārā kustība vairs netika vērtēta kā uzticams ``citu paverdzināto tautu strādājošo masu'' sabiedrotais. J.~Pilsudska rakstos PPS presē ne vienreiz vien atkārtojās tēze, ka, kamēr Polija ietilpst Krievijas impērijā, tā nevar normāli attīstīties, jo krievu tautas stāv uz ievērojami zemākas kulturālās attīstības pakāpes nekā poļi. Krievijas dzīves apstākļi vienmēr traucēs gan pašas krievu tautas, gan it īpaši poļu virzību uz priekšu. Kā primārais uzdevums tika izvirzīta poļu nacionālās atbrīvošanās cīņa. Sabiedriskiem pārkārtojumiem bija jānoris vēlāk, jau neatkarīgajā Polijā. Ar strādnieku šķiru J.~Pilsudski saistīja nevis sociālo tiesību izcīnīšanas programma vai vēlme mainīt sabiedrisko iekārtu, bet cerība, ka proletariāta pieaugošo spēku izdosies izmantot cīņā par Polijas neatkarību. Saviem draugiem sociālistiem J.~Pilsudskis esot pateicis: ``Mēs kādu laiku varam braukt vienā vilcienā. Taču pieturvietā ``Polijas neatkarība'' es izkāpšu''. Kā K.~Markss strādnieku šķirā saskatīja līdzekli, ar ko panākt sociālisma ideālu īstenošanu, tā J.~Pilsudskis tajā redzēja līdzekli poļu nacionālo interešu sasniegšanai. Kā atzīmējis padomju autors A.~Manusevičs, 90.~gadu strādnieku streiku izvēršanās, marksisma plašas izplatīšanās apstākļos nevarēja pat domāt par strādnieku šķiras iesaistīšanu cīņā par Polijas neatkarību, nesavienojot nacionālās neatkarības ideju ar cīņu par demokrātiju un sociālismu. Proletariāta tieksmi izcīnīt sociālismu J.~Pilsudskis u.c. PPS līderi vēlējās izmantot carisma jūga gāšanai, agrākajai Žečpospolitai līdzīga valstiskuma radīšanai ``brīvas federācijas'' veidā. Šis mērķis jau sākotnēji bija pretrunīgs, tā kā saistīja taisnīgo poļu tautas cīņu par neatkarību ar lietuviešu, baltkrievu, ukraiņu apdzīvoto zemju iekļaušanu nākamās Polijas sastāvā kaut vai federācijas veidā. Ar to Polijas brīvība tika pretstatīta Lietuvas, Baltkrievijas, Ukrainas tautu tiesībām uz pašnoteikšanos, sava valstiskuma iedibināšanu.

Viens no J.~Pilsudska līdzgaitniekiem V.~Jodko-Narkevičs žurnālā ``\pltxti{Przedswit}'' rakstā ``Etapi'' centās pierādīt, ka, ja agrāk~--- pirmajā poļu strādnieku kustības etapā partija ``Proletariāts'' centās panākt vienīgi revolūcijas uzvaru pār carismu, tad tagad poļu sociālisti iet tālāk, cenšoties apvienot visas poļu zemes. Taču runa negāja par poļu apdzīvoto un Krievijas, Vācijas un Austroungārijas pakļauto teritoriju apvienošanu, bet gan par Krievijā ietilpstošo, kādreiz Žečpospolitas sastāvā esošo lietuviešu, baltkrievu, ukraiņu u.c. tautu apdzīvoto zemju iekļaušanu Polijā. Mērķa sasniegšanai tika ieteikts censties vienoties ar sociālistiskajām organizācijām tais tautās, ``kuras sastāda ar poļiem vienu zemi''. Poļu sociālistu uzdevums bija panākt dažādo Krievijas tautu opozicionāro grupu vidū ``separātisku tendenču'' attīstību. Žurnāls ``\pltxti{Przedswit}'' rakstīja, ka lietuvieši un latvieši, tāpat kā ukraiņi ``tikai savienībā ar poļu sociālistiem un zem viņu programmas karoga varētu kļūt par ievērojamu spēku''.

XX~gadsimta sākumā J.~Pilsudskis jau mēģināja arī latviešu revolucionārus pārliecināt par nepieciešamību cīnīties ne tikai par sociālismu un autonomiju Krievijas sastāvā, bet arī par pilnīgu neatkarību. Piemēram, Miķelis Valters atzina poļu sociālistu un paša J.~Pilsudska ietekmi uz savu uzskatu veidošanos.

XIX gadsimta beigās Vācijas sociāldemokrātu teorētiskajā žurnālā ``\detxti{Die Neue Zeit}'' (``Jaunais laiks'') iedegās diskusija par Polijas jautājuma vietu starptautiskā proletariāta cīņā. Pārstāvēti bija trīs uzskati. PPS aizstāvēja viedokli, ka Polijas jautājums XIX gadsimta otrajā pusē nav mainījies salīdzinājumā ar iepriekšējo laikmetu un starptautiskās strādnieku kustības gandrīz vai pats svarīgākais uzdevums ir cīņa par Polijas neatkarību. Polijas karalistes sociāldemokrātijas pārstāvji krita citā galējībā un, sekojot R.~Luksemburgai, apgalvoja, ka Polijas neatkarības prasība ir utopiska, neizpildāma, ir pretrunā ar proletariāta sociālistiskajām interesēm un poļu zemju ekonomiskās attīstības interesēm. Visbeidzot, savu viedokli pauda arī Vācijas sociāldemokrātu atzītais teorētiķis K.~Kautskis. Viņš uzskatīja, ka sakarā ar spēcīgas revolucionārās kustības attīstību Krievijā Polijas jautājums neapšaubāmi ir zaudējis savu izcilo starptautisko nozīmi, taču tas nenozīmē, ka poļu sociāldemokrāti nedrīkst izvirzīt Polijas neatkarības prasību. Šim viedoklim piekrita arī toreiz viens no jaunākajiem krievu sociāldemokrātu līderiem V.~Leņins.

Poļu sociālisti cerēja arī uz starptautiskiem satricinājumiem savu nacionālo mērķu īstenošanai. PPS vadība ar J.~Pilsudski priekšgalā uzskatīja, ka Eiropu aptveroša konflikta rezultātā var rasties labvēlīgi apstākļi kārtējās poļu sacelšanās realizācijai.

Zīmīga bija PPS pozīcija jautājumā par Krētas iedzīvotāju sacelšanos pret Turcijas jūgu (1896--1897), kurai juta līdzi lielākā daļa eiropiešu. Arī Polijas karalistes sociāldemokrātija izteica atbalstu sacelšanās dalībniekiem. Turpretī PPS uzskatīja, ka Turcijas novājināšanās var būt izdevīga Krievijai, uzstājās par atbalstu Turcijai cīņā pret Krētas tautas nacionālās atbrīvošanās cīnītājiem. Žurnāls ``\pltxti{Przedswit}'' centās pārliecināt lasītājus, ka ziņas par turku zvērībām Krētā, Armēnijā, slāvu zemēs esot pārspīlētas, ka musulmaņus raksturojot liels iecietīgums, bet sultāna valsts kristieši ``esot zemākā attīstības līmenī'' nekā turku iedzīvotāji, ka jebkāda palīdzība krētiešu sacelšanās kustībai esot palīdzība Krievijai. ``Turcija ne ar ko neapdraud demokrātiju Eiropā, bet Krievija ir tās ienaidnieks un sliktākais no visiem ienaidniekiem'',~--- apgalvoja ``\pltxti{Przedswit}''.

PPS ārpolitisko raksturoja tās attieksme pret angļu-būru karu (1899—1902). Vārdos deklarējot, ka PPS piekrīt vispārējai līdzjūtībai pret būriem, žurnāls ``\pltxti{Przedswit}'' tālāk paziņoja, ka angļu strādnieku dzīves līmenis tik tālu pārsniedz Transvālas darbaļaužu dzīves līmeni, ka to interesēs būtu savas zemes pārvēršana par angļu koloniju. Taču poļiem nav svarīgi tas, cik slikti vai labi dzīvo cilvēki Dienvidāfrikā, jo ``mūsu galvenais un nesamierināmais ienaidnieks ir krievu carisms \citespace{} viss kas to novājina, atvieglo mūsu cīņu \citespace{} Kaitējumam Krievijai ir jābūt par revolucionārā proletariāta starptautiskās politikas izejas punktu''. Pēc žurnāla domām PPS nevar ``vēlēties tāda nemazsvarīga Krievijas pretinieka, kā Anglija, novājināšanu \citespace{} un mēs nekādā gadījumā nevaram pievienoties skarbajam Anglijas ienaidnieku un būru draugu korim''. Acīmredzot šādu pozīciju jāvērtē kā aklu nacionālu egoismu.

PPS mazākums, neatsakoties no Polijas neatkarības idejas, vairāk cerēja uz Eiropas mēroga sociālo revolūciju, kā arī uz liberalizāciju Poliju sadalījušajās valstīs. Kā norāda krievu pētnieki I.~Jažborovska un N.~Buharins, jau no 1900.~gada PPS rindās kristalizējās kreisas grupas, kuras gan arvien tika no partijas izspiestas.

Kā jau minēts, Polijas karalistes poļu sociālistu organizatoriskajam piemēram sekoja arī Prūsijas un Galīcijas sociālistiski orientētie poļi, nodibinot Galīcijas un Cešinas Silēzijas poļu sociāldemokrātisko partiju (\pltxti{Polska Partia Socjalno-Demokratyczna Galicji i Śląska Cieszyńskiego}) un Poļu sociālistisko partiju Prūsijā (poļu \pltxti{Polska Partia Socjalistyczna Zaboru Pruskiego}, vācu \detxti{Polnische Sozialistische Partei in Preußen}). Tās abas izvirzīja lozungus par strādnieku stāvokļa uzlabošanu, demokrātiskām reformām, apvienotu un neatkarīgu Poliju, taču abas partijas bija diezgan vājas. Abas bija daudz tuvākas PPS nekā Polijas karalistes sociāldemokrātijai.

Īpaši jāatzīmē, ka arī Polijas ebreju vidū izveidojās sava sociālistiska partija. Ebreju inteliģenci, sīkpilsonību, strādniecību blakus cionismam saistīja arī sociālima mācība. 1897.~gadā izveidojās Vispārējā Ebreju strādnieku savienība Lietuvā, Polijā un Krievijā (jidišā \hetxti{אַלגעמיין ייִדיש לייבער יוניאַן פון ליטע, פוילן און רוסלאַנד}, poļu \pltxti{Powszechny Żydowski Związek Robotniczy na Litwie, w Polsce i Rosji}, krievu \rutxti{Всеобщий еврейский рабочий союз в Литве, Польше и России}), ko saīsināti sauca par Bundu (savienību). Partija iestājās par demokrātiju un ražošanas līdzekļu sabiedriskošanu. Marksismu Bunda biedri traktēja, savienojot to ar tradicionālajiem priekšstatiem par ebreju tautas īpašo misiju, kaismīgi aizstāvot ideju par ebreju proletāriešu, kā vairāk par citiem apspiestiem, īpašajām vajadzībām. Tāpēc arī tam bija jādibina sava nacionāla partija, nevis jāapvienojas kopā ar citu tautību pārstāvjiem. Nacionālajā jautājumā Bunds aizstāvēja nacionāli kulturālo autonomiju Austrumeiropas ebrejiem, laicīgu izglītību jidiša valodā. Bunda biedri ticēja, ka tādā veidā ebreji varēs saglabāt savu identitāti, nevis asimilēsies. Bundisti arī bija antireliģiski noskaņoti un noliedza nepieciešamību ebrejiem doties uz Palestīnu, ko aizstāvēja cionisti. 1898.~gadā Bunds piedalījās Krievijas Sociāldemokrātiskās strādnieku partijas I kongresa sagatavošanā un iegāja šai partijā kā organizācija autonoma jautājumos, kas skar ebreju proletariātu.

90.~gados blakus iepriekšējai poļu emigrācijai Eiropā, parādījās jauna~--- sociāldemokrātiska.

1889.~gadā uz ārzemēm bēga R.~Luksemburga, nākamajā gadā J.~Tiška. 1892.~gadā poļu sociāldemokrātu grupā, kas grupējās ap R.~Luksemburgu, iekļāvās A.~Varskis. 1893.~gadā Cīrihē ieradās J.~Marhļevskis. Grupa sāka uzstāties ar programmatiskiem paziņojumiem, polemizēja ar PPS ideologiem. Reizē katorgā un izsūtījumā Sibīrijā dažādos gados atradās tūkstoši poļu revolucionāru. Tā, no 1897. līdz 1912.~gadam F.~Dzeržinskis tika arestēts sešas reizes. Vienpadsmit gadu viņš atradās cietumā un katorgā, trīs reizes bija izsūtīts, tai skaitā uz Sibīriju. B.~Veselovskis no 46 savas dzīves gadiem ap 20 pavadīja cietumos un izsūtījumā. No 1900. līdz 1917.~gadam septiņas reizes arestēts bija J.~Unšlihts. Četras reizes arestēts, bet 1906.~gadā izsūtīts uz Turuhanskas novadu bija E.~Pruhņaks. Ap 20~gadu Sibīrijā sabija turp izsūtītais F.~Kons. Jāsaka gan, ka mūsu priekšstati par dzīvi izsūtījumā bieži neatbilst patiesībai. Piemēram, atrodoties izsūtījumā Jeņisejas guberņas Šušenskas ciemā, nākamais krievu boļševiku vadonis V.~Ļeņins ar bisi rokās dienām ilgi pastaigājās pa mežiem un laukiem kopā ar poļu revolucionāru J.~Prominski.

Strādnieku un sociālistiskās kustības panākumus tieši Polijas karalistē salīdzinājumā ar poļu zemēm Vācijā un Austroungārijā padomju vēsturnieki saistīja ne tikai ar karalistes straujāku ekonomisko attīstību, bet arī ar jau minēto marksistisko ideju par XIX gadsimta beigās notikušo ``pasaules revolucionārās kustības centra pārvietošanos uz Krieviju''. Līdz ar to poļu strādnieku kustība ieguvusi varenu sabiedroto~--- krievu strādnieku kustību, kura radījusi varenu marksistisku partiju ar revolucionāru teoriju. Nenoliedzami, ka savijoties Krievijas politiskajai atpalicībai un straujajai ekonomiskajai attīstībai XIX~gs. beigās, pretrunas tajā izpaudās asāk nekā Vācijā un Austroungārijā. Tas ietekmēja situāciju arī Polijas karalistē, strādnieku kustību tajā. Taču nevajadzētu arī aizmirst to, ka tieši Polijas karalistē, salīdzinājumā ar Vācijas un Austroungārijas valdījumā esošajām poļu zemēm, bija vismazāk demokrātijas, arī nacionālā apspiestība sasniedza visaugstāko līmeni. Tāpēc cīņa par nacionālajām interesēm Polijas karalistē norisa ar vislielāko sasprindzinājumu, prasīja vislielākos upurus. 1909.~gada rudenī V.~Ļeņins kategoriski rakstīja, ka ``Polijas brīvība nav iespējama bez Krievijas brīvības''. Ja raugās tikai uz Polijas karalistes XIX gadsimta vēsturi, secinājums atbilda patiesībai, taču plašākā skatījumā visas Polijas brīvība nebija iespējama bez visu triju to sadalījušo monarhiju atbrīvošanās no feodālisma paliekām un politiskas demokratizācijas.

\strong{Nacionālistiskais virziens} mūsdienu politiskas kustības izpratnē poļu zemēs sāka formēties vēlāk nekā sociālistiskais. Acīmredzot tas izskaidrojams ar nacionālo interešu aizstāvju labāko materiālo stāvokli, kas kavēja to politiskās aktivitātes izvēršanos, ar draudiem turīgajiem cīņas gaitā par saviem mērķiem zaudēt ne tikai brīvību (vai pat dzīvību), bet arī mantu, kuras parasti nebija cīnītājiem par sociālo taisnīgumu.

Nacionālisms kļuva par poļu buržuāzijas opozīcijas formu svešzemju valsts varai un galveno tās idejiski politisko virzienu. Tas guva atbalsi arī citos sabiedrības slāņos kā reakcija uz nacionālajiem spaidiem.

Daži poļu vēsturnieki uzsvēruši, ka arī pēc poļu sacelšanās apspiešanas, neraugoties uz grūtībām, ko izsauca īpašumu konfiskācijas un sāktās reformas, šļahtiči nebija salauzti nedz ekonomiski, nedz arī politiski. Viņi saglabāja savu ietekmi uz zemniekiem. Turklāt tas, ka zemnieki saņēma zemi, sekmēja ekonomiskos un sociālos procesus, pastiprinājās lauku iedzīvotāju pieplūdums pilsētās, paātrinājās dzelzceļu celtniecība, kas vienoja lauksaimnieciskos novadus. Zemnieku lietas tagad tika skatītas krieviskajās iestādēs un tiesās, kur bez krievu valodas nevarēja iztikt. Zemnieki sāka saprast, ka viņu nacionālā piederība viņus nostāda otrās šķirās Krievijas pavalstnieku stāvoklī. Tātad, pateicoties privātīpašuma ieguvei, pastiprinājās poļu zemnieku kopējās intereses ar šļahtu.

Pilsētnieku vairākums piederēja mazizglītotajiem un mazturīgajiem iedzīvotāju slāņiem. Strādnieku slānis no agrāko gadsimtu pilsētu plebsa mantoja nacionālās jūtas. Pilsētu apdzīvotības blīvums un to iedzīvotāju kulturālā aktivitāte veicināja strādnieku nacionālās pašapziņas veidošanos, kaut arī, kā raksta poļu vēsturnieki, intensīvais papildinājuma pieplūdums no laukiem pazemināja šīs nacionālās pašapziņas līmeni. (Pēc grāmatas autora domām, šis jautājums vēl ir pētāms, jo bijušie nesenie zemnieki, kuri vēl nebija iesaistījušies šķiru cīņā, nesa sev līdzi šļahtiču ietekmi un varēja stiprināt pilsētnieku nacionālo apziņu pretstatā šķiriskajai.)

Blakus ugodoviešu (samierināšanās ar carismu piekritēju) grupējumam izveidojās radikālāks strāvojums, kurš uzstājās par Žečpospolitas atjaunošanu uz nacionāli-liberālu ideju pamata. Ievērojamākās personības tajā bija J.~Poplavskis un S.~Balickis. J.~Poplavskis sāka izstrādāt virziena idejisko platformu, kura paredzēja tautas iesaistīšanu ``nacionālās politikas ietvaros'', ar to liekot idejiskos pamatus nacionāl-demokrātiskajai kustībai.

1887./1888.~gadā Šveicē nodibinājās nelegāla liberāla virziena grupa~--- \pltxti{Liga Polska} (``Poļu līga''), kura aicināja pārtraukt lojalitātes politiku pret carismu, cīnīties par Polijas neatkarības atjaunošanu uz federālisma pamata. Organizācijas vadītājs sākotnēji bija Z.~Milkovskis. Viņš aizstāvēja konsekventas cīņas pret Polijas sadalītājvalstīm nepieciešamību, uzstājās par Polijas atjaunošanu tās 1772.~gada robežās. Līga sludināja gatavību ``uz aktīvu nacionālu uzstāšanos''. Pēc vairākiem darbības gadiem tajā nostiprinājās tradicionālās pretkrievu pozīcijas, balstoties uz Austrijas un Vācijas savienību. Par lietderīgu tika atzīts viss, kas varēja vājināt Romanovu impēriju. Ar laiku Līgas priekšgalā izvirzījās J.~Poplavskis, T.~Valigorskis, \strong{R.~Dmovskis}.

Pēdējais bija Līgas biedrs no 1889.~gada. R.~Dmovskis arī piedalījās nelegālas studentu organizācijas ``Polijas jaunatnes savienība ``\pltxti{Zet}'' (``\pltxti{Związek Młodzieży Polskiej ``Zet''}) darbībā. 1890.~gadā, kad Varšavas studenti apsprieda jautājumu vai viņiem atbalstīt Krievijā sākušās studentu uzstāšanās pret valdību, viņš neslēpa savu antikrievisko nostāju: ``Polijas revolucionārās kustības uzdevums ir izkarot Polijas neatkarību no Krievijas, neatkarīgi no tā, kas to pārvalda un pārvaldīs, patvaldniecisks cars vai parlaments, kaut pats liberālākais. Sabiedrībā ir jāstiprina neatkarības gars, visur jāuzsver mūsu mērķu atšķirība no krievu mērķiem un centieniem''. 1791.~gada 3.~maija Polijas Konstitūcijas pieņemšanas simtajā gadadienā R.~Dmovskis organizēja studentu demonstrāciju, tāpēc tika arestēts un izsūtīts uz Mītavu (tagad Jelgava). 1893.~gadā R.~Dmovskis izvirzījās par līgas idejisko vadītāju, bet vēlāk kļuva redzamākais vēlākā Polijas valsts vadītāja J.~Pilsudska konkurents XX gadsimta Polijas sabiedrības vēsturē.

Kopš 1893.~gada grupa veica pagriezienu izteikta nacionālisma virzienā, tika pārdēvēta par ``Nacionālo līgu'' (\pltxti{Liga Narodowa}). Kopš 1895.~gada Lembergā (tagad~--- Ļvova), pēc tam Krakovā iznāca Līgas preses orgāns ``\pltxti{Przegląd Wszechpolski}'' (``Vispolijas Apskats''). ``Nacionālās līgas'' šūniņas pastāvēja arī Prūsijā un Austrijā. No 1895.~gada uz Galīciju devās tie Līgas vadītāji, kuriem draudēja arests Krievijā. Prūsijā Līgas darbība attīstījās sarežģītāk, šeit pirmos biedrus tā ieguva tikai 1899.~gadā. Līga pirmkārt centās ietekmēt poļu zemniekus, no 1896.~gada izdeva tiem domātu žurnālu ``\pltxti{Polak}'', (``Polis'') taču nevairījās uzrunāt arī strādniekus. Nacionālā līga izstrādāja cīņas programmu gan pret krieviem, gan vāciešiem.

S.~Balickis un R.~Dmovskis arī centās izlabot ``poļu raksturu'', viņu nacionālo psihi. Viņi vadījās no uzskata, ka nācijai kā ``organizācijai'' sabiedriskai ir jātiek vadītai un ir jāpastāv īpašiem vadības instrumentiem. ``Normālos'' apstākļos tādi instrumenti ir tiesību normas, morāle, tradīcijas. Sadalīšanas un iekļaušanas citos valstiskos organismos rezultātā poļu nācija ieguvusi ``svešas'' tiesību normas, ``sliktas'' tradīcijas un ``sabojātu'' morāli. Tāpēc nāciju vajadzēja ``garīgi atveseļot''. No nācijas bioloģiskās teorijas, attieksmes pret to kā dzīvu organismu izauga pārliecība par nācijas ``labošanas'' iespēju, piemēram, neveselīgo elementu amputācijas ceļā. Pēc S.~Balicka poļu nācijas ``nenormālo'' stāvokli noteica tas, ka tā bija spiesta dzīvot blakus ar ``naidīgām nacionalitātēm'', tai skaitā tādām, kuras agrāk bija ietilpušas Žežpospolitas sastāvā. No šejienes viņš izdarīja secinājumu par poļu kā nācijas, kura atradās ``ārkārtējā stāvoklī'' īpašām tiesībām un uzdevumiem. Tā kā poļu dzīves apstākļi prasīja dzelzs raksturu, stingru gribu, karavīra sīkstumu un spēku, pilsoņu tiesības šādos apstākļos varēja saņemt tikai karotājs. Tāda karotāja-pilsoņa pienākums bija aizstāvēt etniskās, valodas, reliģiskās un psiholoģiskās atšķirības. ``Nacionālās līgas'' idejisko pamatu tad arī veidoja ideja par organizācijas biedru īpašo misiju, viņu tiesībām uz morālu un politisku diktatūru.

1897.~gadā no Nacionālās līgas attīstījās \strong{Nacionāldemokrātiskā partija} (\pltxti{Stronnictwo Narodowa Democratyczne} jeb \pltxti{Narodowa Demokracja}, saīsināti \pltxti{Endecja}), sākotnēji kā slepena organizācija, bet pēc 1905.~gada revolūcijas (Austroungārijai piederošajā Galīcijā~--- jau kopš 1903.~gada) tā varēja darboties arī legāli. Taču turpināja pastāvēt arī nelegāla Nacionālā līga, gan sastāva ziņā, gan idejiski cieši saistīta ar Nacionāldemokrātisko partiju. 1905.~gadā Nacionālajā līgā bija 585~biedri. Tikai 1928.~gadā Nacionālā līga tika likvidēta un iekļauta Nacionāldemokrātiskās partijas struktūrā. (Visu šo organizāciju biedrus tautā saīsināti sauca par endekiem.)

Diemžēl kā ikviens šāda veida virziens, arī poļu nacionālisms dzīvo nacionālo pašapziņu pārvērta par abstraktu principu, kurš savas nācijas intereses izvirzīja kā noteiktu pretējību citām tautām. No tā izrietēja naidīgums ne tikai pret tiešajiem apspiedējiem, bet arī pret veselām tautām, kā arī etniskām grupām, kuras dzīvoja poļu zemēs, atvirzot otrajā plānā kopējās sociālās problēmas un šķiru pretišķības. Nacionāldemokrāti pretendēja uz visas poļu sabiedrības politisko vadību, neraugoties uz tās sociālo neviendabību. R.~Dmovskis atklāti aicināja polonizēt ukraiņus, baltkrievus un lietuviešus. Par galveno nacionālistu uzbrukuma objektu kļuva ebreji, tika prasīta viņu izraidīšana no Polijas. Ja agrāk poļiem svešā ebreju kultūra, īpaši ticība vienkārši izsauca negatīvas emocijas, tad ar tirdzniecības un rūpniecības attīstību, kur ebreji spēlēja aktīvu lomu, radās arvien vairāk pretrunu, kuras savā propagandā centās izmantot poļu nacionālisti. 1903.~gadā iznāca R.~Dmovska darbs ``\pltxti{Mysli nowoczesnnego Polaka}'' (``Mūsdienu poļa pārdomas''), kurā viņš, gan atzīmējot nacionālisma nehumānās izpausmes, tomēr nāciju cīņu par savu izdzīvošanu uzskatīja par sabiedrības attīstības likumu, velkot paralēles ar Č.~Darvina teoriju par dabisko izlasi. Pakāpeniski notika nacionāldemokrātu mērķu korekcija. Ja 1897.~gadā parijas kongress par tās galveno uzdevumu stādīja ``nacionālu sacelšanos un dzimtenes atbrīvošanu ar spēku'', bet 1899.~gadā avīze ``\pltxti{Przegląd Wszechpolski}'' paziņoja, ka atbrīvošanos no krievu jūga nāksies panākt ar ``uguni un dzelzi'', tad drīzumā pozīcija mainījās, tuvinoties ar ugodoviešiem. Gadsimtu mijā, atbildot uz nacionālistiskās propagandas pieaugumu, uzbrukumiem slāvu tautām Vācijā, arī poļu nacionālisti bija spiesti pārorientēties, strauji pāriet no pārsvarā pretkrieviskas propagandas uz pārsvarā pretvācisku. Glābiņu no ``\detxti{Drang nach Osten}'' R.~Dmovskis piedāvāja meklēt poļu vienotībā ar visiem slāviem uz jauniem, antivāciskiem pamatiem.

Runājot par poļu politisko darbību tomēr noteikti jāievēro, ka varas iestāžu noliedzošā attieksme pret pavalstnieku politisko darbību un pārstāvniecības iestāžu ierobežotā loma, to izveidošanas nedemokrātiskais veids (ne tikai Polijas karalistē, bet arī citās Polijas daļās), atstāja iedzīvotāju lielāko vairumu ārpus politisko partiju ietekmes, kas bija raksturīgi autoritārai sabiedriskai kārtībai. Tautas masas lielākoties saglabāja konservatīvu politisku vienaldzību.

\section{Polija XX~gadsimta sākumā}

\epigraph
{Vēloties mainīt valsts iekārtu, ir jāņem vērā pavalstnieku gatavība tam.}
{Tits Līvijs (\latxti{Titus Livius})}

\epigraph
{Naids pret buržuāziju~--- tikumības sākums.}
{Gistavs Flobērs (\frtxti{Gustave Flaubert})}

\epigraph
{Tā tauta, kas domā, ka cilvēkus godīgus padara ticība, bet nevis labi likumi, man šķiet visai atpalikusi.}
{Denī Didro (\frtxti{Denis Diderot})}

\epigraph
{Revolūcijas periodos politisko karjeru visbiežāk taisa vismazāk tikumīgie un sociāli atbildīgie darbinieki, kuri vienmēr ir gatavi sacīt: ``Ir tāda partija!''}
{Igors Kons (\rutxti{Игорь Семёнович Кон})}

\epigraph
{Galvenā cilvēku kļūda ir tā, ka no šodienas nelaimēm viņi baidās vairāk nekā no rītdienas nelaimēm.}
{Kārlis fon Klauzevics (\detxti{Carl Philipp Gottlieb von Clausewitz})}

\epigraph
{Lielvalstis vienmēr ir uzvedušās kā bandīti, bet mazās valstis kā prostitūtas.}
{Stenlijs ~Kubriks (\entxti{Stanley Kubrick})}

\epigraph
{Pirmais pasaules karš bija divdesmitā gadsimta galvenā katastrofa, jo visas pārējās vēlākās katastrofas ir tikai tā sekas.}
{Pēteris Krupņikovs}

\subsection{1905.~gada revolūcija poļu teritorijās}

% page 195


Poļu zemes XX gadsimta sākumā bija ievērojami pavirzījušās uz priekšu savā ekonomiskajā attīstībā. Ja XIX gadsimta sākumā tajās ar lauksaimniecību nodarbināto iedzīvotāju īpatsvars sasniedza 85\%, tad XX gadsimta sākumā tas jau bija samazinājies līdz 56\%. Poļu zemes no pārsvarā lauksaimnieciskām pārvērtās par agrāri-rūpnieciskām. Braucieni uz ārzemēm darba meklējumos un lauku iedzīvotāju migrācija uz pilsētām aptvēra jau ap 9 miljoniem cilvēku. Norisa industrializācija un urbanizācija, kas bija saistītas ar iedzīvotāju proletarizāciju. Dažādu valstu valdījumos šie procesi gan risinājās ar dažādu intensitāti. Gadsimtu mijā izveidojās trīs sociāli-ekonomiskās struktūras tipi. Galīcijā tas bija agrārs, kur līdz 80\% iedzīvotāju iztikas avots bija lauksaimniecība, Polijas karalistē un t.s. Lielpolijā ap Vartas upi, Piejūrā agrāri-rūpniecisks (ap 60\% iedzīvotāju), Silēzijā~--- rūpniecisks (mazāk par 40\%). Visos apgabalos lauksaimniecībā valdošais stāvoklis joprojām palika lielīpašnieku rokās. Polijas karalistē tiem piederēja ap 56,6\%, Prūsijas poļu apgabalos~--- ap 50\%, Galīcijā~--- ap 42\% zemes. Taču politiski poļu sabiedrība savā attīstībā ievērojami atpalika no rietumos esošajiem paraugiem. Sava valstiskuma trūkums, nacionālā apspiestība, poļu buržuāzijas vājās ekonomiskās pozīcijas bremzēja tās politisko un organizatorisko konsolidāciju.

XX gadsimta sākumā \strong{Polijas karaliste} bija kļuvusi par nozīmīgu Krievijas ekonomisku rajonu, kurā bija koncentrēti ap 15\% visu impērijas rūpniecības strādnieku. Taču 1900.--1903.~gadā Polijā, tāpat kā visā pasaulē, izpaudās ekonomiskā krīze. Dažādas saimniecības nozares tā skāra dažādā laikā un nevienādā mērā. 1904.~gada janvārī ar Japānas jūras kara spēku uzbrukumu Krievijas flotei Portarturas bāzē sākās krievu--japāņu karš. Karš traucēja saimniecības attīstību ne tikai pašā Krievijā, bet arī Polijas karalistē. Kopumā XX gadsimta sakumā Polijas karalistē rūpniecības produkcijas ražošana saruka par 35\%. 1904.~gadā Polijā bija zema raža. Samazinājās rūpniecības uzņēmumu skaits. Kritās algas rūpniecībā, auga bezdarbnieku skaits. Kritās strādājošo dzīves līmenis. Pilnīgs vai daļējs bezdarbs Varšavā aptvēra 60\%, bet Lodzā 70\% strādnieku, darba alga samazinājās par 20--25\%, ievērojami pieaugot pirmās nepieciešamības preču cenām.

Tas viss kopā radīja neapmierinātības noskaņojumu tautā. Sabiedrības uzmanību piesaistīja poļu sociāldemokrātiskās kustības veterāna M.~Kaspšaka rīcība 1904.~gada 27.~aprīlī Varšavā. Viņa ierīkoto nelegālās literatūras spiestuvi atklāja policija. M.~Kaspšaks nevēlējās padoties, atšaudoties nogalināja četrus un ievainoja vēl vienu policistu. (Pēc dažiem mēnešiem viņam piesprieda nāves sodu, ko nomainīja ar mūža ieslodzījumu, taču 1905.~gada septembrī viņš Varšavas citadelē pakārās.) Jau 1904.~gada otrajā pusē politiskās aktivitātes Polijas karalistē ievērojami pieauga. Vietējie politiskie spēki uz to reaģēja dažādi.

Konservatīvie poļu politiķi (ugodovieši) centās demonstrēt cara impērijai atbalstu grūtā stundā, lai pēc kara saņemtu piekāpšanos nacionālajā sfērā. Viņi organizēja katoļu hospitāli frontei, piedalījās pieminekļa atklāšanā Katrīnei II Viļņā. 1905.~gadā šī kustība noformējās par \strong{Reālās politikas partiju} (\pltxti{Stronnictwo Polityki Realne}j), kuras biedrus tautas valodā sauca arī par ``reālistiem'' jeb ``pozitīvistiem''. Partiju vadīja grāfs Z.~Veļepoļskis un publicists E.~Piļcs. Tās programma ietvēra poļu vienlīdzības prasību, kas praksē nozīmēja municipālo pašvaldību~--- zemstu ieviešanu, zvērināto piesēdētāju tiesas, prasības pēc poļu pielaides administratīvajiem amatiem, vienādiem nodokļiem, reliģiskās iecietības un poļu valodas pielietošanas izglītībā, kas būtu pirmais solis uz Polijas karalistes autonomiju. Faktiski programma negāja tālāk par vēl 1815.~gadā Vīnē pieņemto Polijas sadales noteikumu atjaunošanu. Partija neguva plašu atbalstu Polijas karalistes iedzīvotājos, tai bija maz biedru. Taču karalistē tai bija diezgan liela ietekme, jo tās vadoņiem bija labi sakari ar valdošajām aprindām Pēterburgā.

Salīdzinājumā ar XIX gadsimta 90.~gadiem būtiski no pretkrieviskās orientācijas uz pretvācisku bija mainījusies \strong{Nacionāldemokrātiskās partijas} (\pltxti{Stronnictwo Narodowo-Demokratyczne}) nostāja. \strong{R.~Dmovskis}, izanalizējis Vācijas un Austroungārijas vadošo politiķu attiecības pret poļu mazākumiem savās valstīs, nonāca pie konstatējuma, ka Vācija veic spiedienu uz Habsburgu impēriju, virzītu uz tās pakļaušanu Vācijas ekspansionistiskajiem plāniem. No tā izrietēja tālāks secinājums: poļu orientācijai uz Austroungāriju nav perspektīvu, ir jāmeklē citas politiskās alternatīvas.

Sadarbības ar Krieviju izdevīgums kļuva tik acīm redzams, ka Nacionāldemokrātiskās partijas kongress 1903.~gadā savā rezolūcijā deklarēja: ``Neatsakoties no sava Polijas atjaunošanas ideāla, dotajā vēsturiskajā situācijā mēs uzskatām par nepieciešamu censties izcīnīt tai pēc iespējas lielāku patstāvību Krievijas ietvaros''. R.~Dmovskis neatbalstīja poļu sacelšanās plānus, pamatoti uzskatot, ka tai nebūs izredžu uz panākumiem, taču izteicās par nacionālo apziņu stiprinošām manifestācijām. 1904.~gadā Nacionāldemokrātiskā partija pilnībā atteicās no nacionālās neatkarības prasības un aicināja cīnīties par visu poļu zemju (ietilpstošu Krievijā, Vācijā un Austroungārijā) apvienošanu vienotā autonomā Polijā Krievijas impērijas sastāvā. Politiskās un kulturālās autonomijas ieguve bija paredzēta sadarbības un sarunu ceļā ar Krieviju. Par poļu nācijas galvenajiem ienaidniekiem nacionāldemokrāti uzskatīja vāciešus un ebrejus. Partijas programma pamatojās politiskajā pragmatismā, paredzēja ``organisku darbu''~--- tautsaimniecības attīstību, tautas izglītošanu.

Tomēr savās prasībās pret Krieviju Nacionāldemokrātiskā partija ar R.~Dmovski priekšgalā gāja tālāk kā konservatīvie poļu politiķi. Attiecībā pret Krieviju partijas programma izvirzīja prasību pēc kulturālas un politiskas autonomijas, ko vajadzētu sasniegt sadarbības un sarunu ceļā. Nacionāldemokrātija prasīja poļu valodas ieviešanu ne tikai skolās, bet arī administrācijā. Kopumā R.~Dmovskis bija pārliecināts, ka Krievijas impērija pārvarēs krīzi un tāpēc uzstājās pret jebkādām radikālām prasībām. Nacionāldemokrāti jeb tautas valodā endeki tikai sākotnēji atbalstīja streikus un demonstrācijas, bet turpmāk bija gatavi pret tiem cīnīties, ja Polijai tiktu piešķirta autonomija. Vēlāk R.~Dmovskis rakstīja: ``Nacionāldemokrātiskā nometne jau no paša sākuma noteikti uzstājās pret visiem sacelšanās vai revolucionāras kustības [izvēršanas] mēģinājumiem.'' R.~Dmovskim 1905.~gada sākumā ierodoties Varšavā, tika veikti pasākumi partijas organizatoriskai nostiprināšanai, hierarhiskas struktūras izveidei. 1905.~gada jūnijā pie Nacionāldemokrātiskās partijas izveidojās arī Nacionālā strādnieku savienība (\pltxti{Narodowy Zwiazek Robotniczy}). Nacionāldemokrātiskajai partijai nebija savas kaujinieku organizācijas, taču tāda darbojās Nacionālajā strādnieku savienībā. Īslaicīgi darbojās arī dažas citas endeku ietekmē esošas kaujas grupas.

``Reālisti'' un ``endeki'' veda sarunas ar krievu liberāļiem (nākamajiem kadetiem), kuri principā piekrita Polijas autonomijai, bet nekādas konkrētas norunas negribēja slēgt. Nacionāldemokrātu un ``reālistu'' prasība pēc autonomijas, ar kuras izpildi, iespējams, varētu sekmīgi apstādināt revolucionāro kustību, tā arī palika varas iestāžu nesadzirdēta.

Principiāli citu pozīciju uzturēja \strong{Polijas Sociālistiskā partija} (\pltxti{Polska Partia Socijalistyczna}) \strong{(PPS)}. Tajā t.s. ``jaunie'' sākotnēji ticēja, ka drīzumā var sākties revolūcija Vācijā, sacelšanās Polijā un progresīvās Eiropas karš pret reakcijas citadeli~--- carisko Krieviju. Taču PPS rindās uzvarēja t.s. ``vecie'' ar \strong{J.~Pilsudski} priekšgalā, kuru līnija veda uz kara un revolūcijas Krievijā izmantošanu, lai atbrīvotu Polijas karalisti no tās. J.~Pilsudskis uzskatīja, ka labvēlīgs brīdis kārtējai bruņotai sacelšanās par Polijas brīvību jau ir pienācis vai arī tūlīt pienāks. ``Vecie'' paļāvās ne uz Krievijas revolucionāro kustību, bet uz citām tautām, kuras apspieda carisms. Šī grupa arvien vairāk attālinājās no marksistiskās ideoloģijas. PPS vadītāji nolēma sadarboties ar Japānas valdību. J.~Pilsudskis 1904.~gada jūlijā ieradās Tokijā, lai organizētu poļu leģionus no gūstā saņemtajiem Krievijas armijas poļiem cīņai par Polijas neatkarību. Japānas imperatora ierēdņiem viņš arī piedāvāja finansēt PPS pretkrieviskās akcijas, apgādāt to ar ieročiem, solot traucēt armijas mobilizācijai Krievijā, veikt izlūkošanu un sabotāžu un pie labvēlīgiem nosacījumiem sarīkot sacelšanos kopā ar citām Krievijas impērijas apspiestajām tautām.

J.~Pilsudska uzskati bija koncentrēti izklāstīti memorandā, kuru viņš 1904.~gada 13.~jūlijā iesniedza Ārlietu ministrijas ierēdnim Tokijā. Pēc poļu politiķa domām Krievijas impērijas galvenais vājuma avots, tās ``Ahileja papēdis'', kurā visiem tās ienaidniekiem ir jācenšas trāpīt, ir tās vienotības trūkums. Ne visām nācijām impērijā esot vienāda vērtība. J.~Pilsudskis līdzīgi vācu nacionālistiem visas nācijas dalīja ``vēsturiskajās'' un ``nevēsturiskajās''. Pie pirmo kategorijas viņš pieskaitīja poļus, somus, gruzīnus, zināmā mērā'' armēņus, pie otrajām~--- lietuviešus, baltkrievus un latviešus. Viņš gandrīz pilnībā noliedza iespēju, ka patstāvīgi politiski varētu uzstāties rusīni (rutēņi), tātad ukraiņi. Poļiem viņa memorandā bija veltīta īpaša uzmanība. Ja pārējie nacionāļi patstāvīgi varēja tikai radīt politisku opozīciju valdībai un sabotēt tās darbību, tad ``poļi ir spējīgi novest notikumus līdz atklātai cīņai un piesaistīt sev pārējās nācijas uz bruņotu cīņu par atbrīvošanos no krievu paverdzinātāja''. Revolūcija pašā Krievijā tika vērtēta kā maz iespējama.

Taču Japānā J.~Pilsudskis sastapās ar R.~Dmovski, kurš bija atbraucis pārliecināt japāņus, ka nevajag provocēt poļu sacelšanos. R.~Dmovski pieņēma Japānas Ģenerālštāba priekšnieks, pēc kura lūguma viņa sastādīja divus plašus ziņojumus: par stāvokli Krievijā un par Polijas jautājumu. R.~Dmovskis centās pierādīt, ka par nacionāla mēroga sacelšanos Polijas nevar būt runas, bet lokālas uzstāšanās nesīs ļaunumu gan Polijai, gan Japānai, jo Krievija, baidoties no nemieriem, pie Vislas tur veselu armiju (tā mazinot savu kaujas potenciālu Tālajos Austrumos), kura noslīcinās asins jūrā jebkuru dumpi. Ja Polijā sāktos sacelšanās, Vācija labprāt palīdzētu Krievijai to apspiest. Jebkura ārēja palīdzību poļu separātistiem tikai veicinātu Krievijas sabiedrības saliedēšanos.

Japānas valdošās aprindas ar neuzticību raudzījās uz J.~Pilsudska plāniem, radīt leģionu netaisījās, taču pētīja iespējas izmantot poļu nacionālistus sev derīgas informācijas ieguvei un pēc dažām ziņām piešķīra šim nolūkam 20 tūkstošus (daži poļu vēsturnieki uzskata, ka pat ap 30 tūkstošu) sterliņu mārciņu. Lauvas tiesa no šīs naudas tika izmantota ieroču iegādei un PPS darbības finansēšanai.

Atgriezies Polijā, J.~Pilsudskis nodibināja kaujas organizāciju (\pltxti{bojòwka}). Jau 1904.~gada aprīlī PPS kaujas družīnas sāka t.s. ``spiegu'' (faktiski~--- savu politisko pretinieku, turētu aizdomās par sadarbību ar varas iestādēm) ``sodīšanu''~--- slepkavošanu. Kaujinieki arī nodarbojās ar ekspropriācijām~--- banku un pasta vilcienu aplaupīšanām partijas finansiālām vajadzībām. Tiesa, ilgu laiku J.~Pilsudskis pats personīgi šajos bīstamajos pasākumos nepiedalījās, aprobežojoties ar dziļi konspirēta vadītāja lomu, kas, kā rakstīja žurnāliste un sabiedriska darbiniece I.~Pannenkova, kura 1922.~gadā ar pseidonīmu \emph{Jan Lipecki} publicēja darbu ``\pltxti{Legenda Piłsudskiego}'' (``Pilsudska leģenda''), viņa nelabvēļu vidū izsauca runas par bailīgumu.

1904.~gada 13.~novembrī PPS organizēja Varšavā grandiozu protesta demonstrāciju pret poļu karavīru iesaistīšanu krievu-japāņu karā. Pret policiju, kas centās traucēt demonstrācijas norisi, darbojās poļu kaujinieku grupa. Policistu vidū bija viens nogalinātais un vairāki ievainotie. Sākās sadursmes ar policiju arī citās Polijas pilsētās. PPS popularitāte, sevišķi strādnieku vidū strauji auga, 1905.~gadā tajā bija ap 50~000 biedru.

J.~Pilsudskis un viņa atbalstītāji PPS uzskatīja teroru par līdzekli kā uzkurināt revolūciju Polijas karalistē. Kaujas organizācijas tika uzskatītas par kaujas armijas iedīgli, kurai bija jāizraisa karš pret carismu. T.s. ``jaunie'' PPS sastāvā, kuri noraidīja bruņotu cīņu, jau 1905.~gada martā panāca iepriekšējās partijas vadības atlaišanu, taču bruņošanās tolaik tikai pastiprinājās.

Turpretī \strong{Polijas karalistes} un Lietuvas sociāldemokrātija uzskatīja, ka revolūcijā, lai gāztu carisko režīmu, ir jāsadarbojas ar Krievijas revolucionāriem. PKunLSD bija cieši saistīta ar Krievijas sociāldemokrātisko strādnieku partiju (KSDSP). Taču, kā jau iepriekš teikts, poļu sociāldemokrātiem bija savi patstāvīgi uzskati. 1903.~gadā PKunLSD savā IV kongresā nolēma, ievērojot noteiktas prasības, apvienoties ar KSDSP. Saskaņā ar to KSDSP II kongress 1903.~gada jūlijā uzaicināja tajā ar padomdevēju balsīm piedalīties arī PKunLSD pārstāvjus. Šie pārstāvi: J.~Gaņeckis un A.~Varskis uzstājās pret KSDSP programmas projektā ietverto prasību pēc nāciju pašnoteikšanās tiesībām, domājot, ka tas nāks par labu poļu nacionālistiem no PPS. Kad projekts ar minēto prasību tika pieņemts, viņi atstāja kongresu. Abu partiju apvienošanās notika jau revolūcijas gaitā KSDSP IV kongresā 1906.~gada pavasarī, kad poļu sociāldemokrāti atteicās no savas nostājas nacionālajā jautājumā un vairs neprasīja lai krievu sociāldemokrāti pārskata savu nostāju par nāciju pašnoteikšanās tiesību atzīšanu. Šai kongresā PKunLSD pārstāvji jau sadarbojās ar krievu boļševikiem. A.~Varskis un F.~Dzeržinskis kļuva par KSDSP CK locekļiem.

1904.~gada beigās PKunLSD cara patvaldības gāšanas, politisko brīvību izcīnīšanas lozungus papildināja ar demokrātiskas republikas, satversmes sapulces sasaukšanas un poļu zemju autonomijas lozungiem. Partija, kura līdz 1905.~gadam noliedza teroru, sākoties revolūcijai, pieņēma to kā ``taktisku'' cīņas līdzekli. Sociāldemokrātu kaujas organizācijas uzdevums bija aizstāvēt strādniekus pret cara varas iestāžu patvaļu un vēlāk arī pret endeku kaujiniekiem.

1905.~gada 9.~(22.)~janvāris~--- ``asiņainā svētdiena'' arī Polijā izsauca plašu rezonansi. Praktiski uzreiz kā Poliju sasniedza ziņas par notikumiem Pēterburgā, visas politiskās partijas griezās ar paziņojumiem pie tautas. Dažas no tām, kā PPS, uzsvēra nacionālo jautājumu. Tagad grūti pateikt vai tautas vairākums ticēja iespējai panākt nacionālo autonomiju, vai vienkārši cerēja uz liberālām pārmaiņām valsts iekārtā. Faktiski gan vienu, gan otru mērķi varēja panākt tikai cīņā pret cara monarhiju, ko pirmajā vietā izvirzīja sociāldemokrāti. Atsaucoties uz ``asiņaino svētdienu'' Pēterburgā, daudzās Polijas karalistes pilsētās uzliesmoja streiki, kuri ilga trīs nedēļas, streikoja ap 93\% visu strādnieku. Viņu ekonomiskās un daļa politisko prasību (par strādnieku komiteju atzīšanu) tika apmierinātas, taču Varšavā, Lodzā u.c. norisa sadursmes ar policiju, lija asinis.

Līdz februāra vidum oficiāli jau bija uzskaitīti 19 vardarbīgas nāves gadījumi, to faktiskais skaits droši vien bija lielāks. 1905.~gada jūnijā, kad, tika apšauta tekstilrūpniecības strādnieku demonstrācija Lodzā (bojā gāja 25~cilvēki), pilsētu pārklāja barikādes. To skaits bija ap 100. Nākamajās četrās dienās bojā gāja jau ap 200~cilvēku. Faktiski tā bija kā nacionālās, tā sociālās cīņas izpausme, stihiska bruņota uzstāšanās, pirmā pēc 1863.~gada sacelšanās.

Notikumi Polijā bija Krievijas revolūcijas sastāvdaļa. To vēriens, iznākums bija atkarīgs no revolucionārās cīņas vēriena un iznākuma visā Krievijā. Polijas karaliste kļuva par vienu no svarīgākajiem revolucionārās cīņas centriem Krievijā. Notikumiem Polijas karalistē bija pat lielāka nozīme nekā citās, impērijas krievvalodīgajās provincēs. Pēc dažu pētnieku datiem Polijas karalistē šai laikā dzīvoja 13,5\% no visiem Krievijas fabriku un rūpnīcu strādnieku, taču streikotāju skaits šeit sasniedza 29\% no visiem streiku dalībniekiem Krievijā. Pēc citiem datiem starp visiem Krievijas impērijas streikotājiem Polijas karalistes strādnieki 1905.~gadā sastādīja 34\%, 1906.~gadā~--- 46\%, bet 1907.~gadā~--- 14\%. Ievērojami pieauga revolucionāro partiju biedru skaits. Sociāldemokrātu skaits no 2~000 izauga gandrīz līdz 40~000, PPS rindās skaitījās 55~000~--- 60~000~cilvēku.

1905.~gada pavasarī aktivizējās muižu laukstrādnieku kustība, kuri prasīja darba apmaksas paaugstināšanu. Zemnieki necentās sadarboties ar laukstrādniekiem, taču savukārt ap 50~apriņķos zemnieki izvērsa cīņu par servitūtiem. Drīzumā savās sapulcēs viņi prasīja poļu valodas atjaunošanu zemākajās iestādēs gminās, tiesās un skolās. Tika panākts, ka kopā ar krievu valodu gminās varēja lietot arī poļu valodu. 1905.~gadā Polijas karalistē uzliesmoja ``skolu revolūcija''. Streikoja ne tikai Varšavas universitātes studenti, bet ari vidējo un pat sākumskolu skolnieki. Šo streiku 1905.~gada 6.~februārī (jaunais stils) atbalstīja arī vecāku mītiņi. 1905.~gada pirmajā pusē imperators Nikolajs II pēc Ministru komitejas ierosinājuma Polijas karalistes sākumskolās atļāva virkni mācību priekšmetu mācīt dzimtajā valodā. Tika atcelti ierobežojumi poļu ekonomiskajai darbībai Krievijas Rietumu guberņās. Tai skaitā viņiem atļāva šeit iepirkt nelielus zemes gabalus izmantošanai rūpnieciskām vajadzībām. Tomēr 1905~.~gada 1.~maijā imperatora apstiprinātajā Ministru komitejas žurnālā tika ietverts slepens punkts, ar kuru tika aizliegts deviņās rietumu guberņās valsts amatus ieņemt ``poļu izcelsmes personām''. Tai skaitā bija arī krievu valodas, vēstures, ģeogrāfijas un pedagoģijas skolotāju amati. Ar 1905.~gada 1.~oktobra cara dekrētu tika atļauts dibināt privātskolas ar poļu mācību valodu. Taču šo skolu absolventiem nebija tiesību iestāties Krievijas impērijas augstākajās mācību iestādēs. Praksē tas noveda pie tā, ka viņi turpināja izglītību ārzemēs.

Pēc cara Nikolaja II 1905.~gada 17.~oktobra manifesta, kurš solīja pavalstniekiem dot vienkāršākās pilsoniskās tiesības, legalizējot politiskās partijas, ieviešot ko līdzīgu vispārējām vēlēšanu tiesībām un sasaucot Valsts Domi kā centrālo likumdevēju iestādi, Polijā sākās ``brīvības dienas'' (\pltxti{dni Wolności}), notika plaši mītiņi un demonstrācijas. Tautas masas nesamierinājās ar carisma piekāpšanās apjomu. Varšavas ģenerālgubernators G.~Skalons 26.~oktobrī Iekšlietu ministrijai Pēterburgā ziņoja: ``Revolucionārā kustība man uzticētajā novadā redzami un ātri pastiprinās, aptverot arvien lielākus ļaužu lokus un nokļūstot pat zemnieku masās. Situācija ir visai nopietna, un es redzu tikai vienu izeju~--- tūlītēju kara stāvokļa ievešanu visā Polijas karalistē. Jebkuru vilcināšanos es uzskatu par bīstamu''. 28.~oktobrī Nikolajs II parakstīja dekrētu par kara stāvokļa ievešanu Polijas karalistē. Uzreiz sākās masu aresti, tika slēgtas revolucionāro avīžu redakcijas, tika izdota pavēle šaut uz demonstrantiem. Taču daudzās Krievijas pilsētās izvērsās solidaritātes kustība. Uz notikumiem Polijā reaģēja ne tikai revolucionāri, bet arī Krievijas liberālā kustība. Konstitucionālo demokrātu (kadetu) partijas I kongress 1905.~gada 12.--18.~oktobrī jau bija pieņēmis programmu, kura atzina, ka Polijas karalistei ir jādod autonomijas tiesības, saglabājot valsts vienotību. Tajā bija paredzēts, ka robeža starp Polijas karalisti un kaimiņu guberņām var tikt labota saskaņā ar tautību sastāvu un vietējo iedzīvotāju vēlmēm, pie kam Polijas karalistē ir jādarbojas visas valsts garantētām pilsoniskajām brīvībām un tautību tiesībām uz kulturālo autonomiju, jābūt nodrošinātām mazākumu tiesībām. Pēc kara stāvokļa ievešanas Polijas karalistē liberāļi organizēja lielus mītiņus, kuru rezolūcijas prasīja atcelt kara stāvokli Polijas karalistē un piešķirt tai autonomiju. Prasību atbalstīja arī Krievijas zemstu darbinieku kongress 6.--13.~novembrī. Pēc krievu liberāļu uzskatiem autonomijas piešķiršana likvidētu ``karstu punktu'' Krievijas valstī, ko nespēja varas iestādes ar represiju politiku. Ievērojamais krievu sabiedriskais un politiskai darbinieks V.~Maklakovs rakstīja: ``Poļu autonomijas ideja bija tam laikam visvalstiskākā no idejām, kuras bija mūsu sabiedrības rīcībā.'' Tiesa, krievu oktobristi (``oktobristi'' jeb ``17.~oktobra savienība'' 1905--1917.~gadā bija labēja politiska partija Krievijā) uzstājās pret autonomijas piešķiršanu Polijas karalistei, norādot, ka vēsturiski tā pati ir vainīga tās zaudēšanā.

Drīzumā Krievijas Ministru padomes priekšsēdētājs S.~Vitte griezās pie ģenerālgubernatora G.~Skalona ar jautājumu, vai nevajag atcelt kara stāvokli. Pēdējais iesniedza šādu priekšlikumu. Valdība paziņoja, ka kara stāvokļa ieviešana ``atskurbinājusi'' Poliju un tas pēc ģenerālgubernatora lūguma tiek atcelts. Tomēr Polijas karalistē tika izvietoti ap 300~000 krievu karavīru. Kad 1905.~gada nogalē centrālajā Krievijā notikumi sasniedza kulmināciju, Polijā streiku skaits jau samazinājās. Tomēr līdz ar vispārēja politiska streika, kas pārauga sacelšanās, sākšanos Maskavā, baidoties no līdzīgas notikumu attīstības Polijā, ģenerālgubernators 21.~decembrī novadā atkal ieviesa kara stāvokli. Tomēr apstākļi, lai sekmīgi izvērstu sacelšanos, karaspēka appludinātajā Polijā nebija un streiki neieguva tādu apjomu kā 1905.~gada janvārī--februārī un arī oktobrī--novembrī. Tomēr 1905.~gadā strādnieku dalības streikos koeficients Polijas karalistē bija 2,1~reizes augstāks nekā visā Krievijas impērijā kopā.

Saasinājās arī pretrunas pašā poļu sabiedrībā. 1905.~gada oktobrī Lodzā jau poļi šāva uz poļiem, jo endeki centās neitralizēt PPS, kura pēc viņu domām veda poļus uz veltu asins izliešanu. PPS iestājās par konfrontāciju, nacionāldemokrāti~--- par kooperāciju ar varas iestādēm. Starp nacionālistiem un sociālistiem arī vēl 1906./1907.~gada ziemā Lodzā notika asiņainas cīņas. Tas radīja savstarpēju rūgtu naidu, kurš pastāvēja arī vēlāk~--- jau patstāvīgajā Polijas valstī.

Krievu vēsturnieks N.~Postņikovs, kurš pētījis poļu teroristu darbību 1905.~gada revolūcijā, aprakstījis poļu kaujinieku darbību. Visu partiju kaujas organizāciju struktūra bija apmēram līdzīga: Partijas vadībai pakļāvās kaujas nodaļa, kura vadīja guberņu kaujas organizācijas, kas savukārt dalījās ``desmitos'' vai ``pieciniekos''. Šo organizāciju darbība bija stingri konspirēta. Kaujas grupās pastāvēja stingra disciplīna, notika militārā apmācība pēc armijas parauga. J.~Pilsudska vadībā tika gatavoti vadītāji bruņotām cīņām. Par kaujiniekiem kļuva dažādu slāņu pārstāvji, taču parasti viņi bija paši aktīvākie savu partiju biedri. Skaitliski lielākā un kaujas spējīgākā bija PPS Kaujas organizācija, kurā 1905.--1907.~gadā sastāvēja 7~631 cilvēks. Endekiem pakļautajās, īpaši Nacionālas strādnieku savienības kaujas organizācijā, bija ap 1~000 cilvēku. PKunLSD kaujinieku skaits svārstījās no 1~000 līdz 1~500 cilvēkiem.

Kaujinieku realizētie terora akti bija dažādi. Tos parasti veica viena ``piecnieka'' biedri. Piemēram, uz ielas viņi piegāja pie izraudzītā upura (kāda policista, virsnieka vai ierēdņa, kuri visi nēsāja mundierus), tiešā tēmējumā viņu nošāva un nozuda no nozieguma vietas. Cits plaši izplatīts teroristisku uzbrukumu veids bija bumbu sviešana valsts iestādēs (policijas iecirkņos, karaspēka daļu sardzes telpās). Bumbu meta viens kaujinieks, pārējie viņu piesedza.

Pavisam 1905.--1907.~gadā Polijas karalistes teritorijā tika veikti 3~166 kaujinieku uzbrukumi (1905.~gadā~--- 1~441, 1906.~gadā~--- 1~090, 1907.~gadā~--- 635) 1~108 apdzīvotās vietās. 1905.~gadā un četros pirmajos 1906.~gada mēnešos bija nogalināti 127 policisti, žandarmi, militāristi, ievainoto skaits sasniedza 207. Kopā 1905.--1906.~gadā karalistē tika nogalināti 790~varas pārstāvji, vairāk nekā 860~ievainoti. No 1906.~gada terora akti tika vērsti arī pret fabrikantiem, streiklaužiem, visiem, kurus uzskatīja par ``nodevējiem''. Par streiklaužiem uzskatīja veikalu, kafejnīcu, restorānu īpašniekus, kuri streiku laikā nepārtrauca savu uzņēmumu darbu, lauciniekus, kuri pieveda pilsētām produktus u.tml. Bez tam karalistes teritorijā tika veiktas 443~ekspropriācijas, sagrābjot 600~143 rubļus.

Īpaši PPS Kaujas organizācijas akcijas poļu atmiņā atdzīvināja nacionālo sacelšanos romantiskās tradīcijas. 1905.--1907.~gadu laikā tā veica 70,2\% visu Polijas karalistē īstenoto kaujas akciju. Tikai 1906.~gadā uzskaitīti J.~Pilsudska vadīto kaujinieku veikti 678~atentāti, kuriem par upuri krita 336~krievu militāristi vai ierēdņi. 1906.~gada augustā šī darbība sasniedz augstāko punktu, notika vairāki uzbrukumi pasta vilcieniem. ``Asiņainajā trešdienā'' (\pltxti{Krwawa środa})~--- 1906.~gada 15.~augustā vienlaikus notika sekmīgi uzbrukumi vairākiem policijas iecirkņiem, kuros tika nogalināti 72 policisti un žandarmi, bet 18.~augustā~--- 19~gadus vecā V.~Krahaļska veica nesekmīgu atentātu pret ģenerālgubernatoru G.~Skalonu. Endeku kaujinieki pavisam veica 214 akciju (desmit reižu mazāk nekā PPS kaujinieki), apmēram tikpat kaujas akciju bija uz visu pārējo Polijas karalistē darbojošos partiju kaujinieku rēķina.

1906.~gada beigās sāka mainīties kaujinieku sastāvs. Ja agrāk par tiem cilvēki kļuva politisku motīvu vadīti, tad tagad viņu vidū iesaistījās arvien vairāk puskriminālu elementu, kuri izmantoja ieročus krimināliem mērķiem. Nacionālās strādnieku savienības kaujinieku sadursmes ar sociālistiem jau atgādināja atklātu pilsoņu karu un ar to izsauca neapmierinātību sabiedriskajā domā. ``Kaujinieku'' uzbrukumi parastiem iedzīvotājiem izsauca dziļi negatīvu attieksmi.

Valsts reakcija pret teroristu darbību bija nekavējoša. Polijas karalistes cietumi bija pārpildīti. Tika piespriesti un izpildīti nāves sodi. 280 PPS Kaujas organizācijas dalībniekiem tika piespriests nāves sods, 747~--- katorgas ieslodzījums. Pavisam 1904.--1907.~gadā tika apcietināti 2~739~kaujinieki. Taču pēc Polijas valstiskuma atjaunošanas virkne PPS Kaujas organizācijas dalībnieku (arī pēc nāves) par saviem nopelniem tika apbalvoti ar ordeņiem.

1906.~gada vasarā atkal saasinājās cīņa laukos. Muižu algādži sāka iznīcināt muižniekiem piederošos lopus, zemnieki sagrāba gan muižniekiem, gan valstij piederošās zemes platības, dedzināja muižu ēkas, iesaistījās sadursmēs ar viņus ``nomierināt'' atsūtīto karaspēku.

1906.~gada septembrī J.~Pilsudskis mēģināja nodibināt sadarbību ar Austroungārijas izlūkdienestu un tā pārstāvim stādījās priekšā kā 70~000 cilvēku lielas bruņotas partijas priekšstāvis. Taču abu lielvalstu: Austroungārijas un Krievijas militārs konflikts bija vēl tālu priekšā un austrieši nesteidzās ar atbildi.

Šī J.~Pilsudska darbība izsauca asas nesaskaņas starp PPS spārniem. Tās saasināja PPS Kaujas organizācijas 48~kaujinieku uzbrukums karavīru apsargātam pasta vagonam pie Rogovas (\pltxti{Rogów}, Lodzas tuvumā), kad tika nogalināts 1 policists, ievainoti trīs pasta darbinieki, divi slepenpolicisti un 13~karavīri. Uzbrucēju laupījums bija ap 30~000 rubļu. PPS vadība uzskatīja, ka šādas akcijas nedrīkst novest pie upuriem karavīru vidū, jo tas varēja traucēt revolucionārajai propagandai armijā un nolēma apturēt kaujas organizācijas darbību. Uz to atbildot šīs organizācijas vadība noturēja savu konferenci, kura pieprasīja, lai partijas tuvākajā kongresā tiktu pieņemts lēmums par gatavošanos uz bruņotu sacelšanos un bruņotas milicijas radīšanu kaujas organizācijas vadībā.

PPS IX kongresā 1906.~gada novembrī kaujas organizācijas darbības apspriešana kļuva par galveno jautājumu. Delegātu vairākums pieņēma rezolūciju, ka kaujas organizācija ``realizēja politiku un taktiku, kura atradās klajā pretrunā ar partijas līniju'', kas veda ``pie revolucionārās kustības Polijā nošķiršanās no šīs kustības visā Krievijas valstī''. Pēc tam 7 delegāti atstāja kongresu. Viņi noturēja savu konferenci, kurā nolēma norobežoties atsevišķā grupā. Tā radās PPS~--- ``Revolucionārā frakcija'' (\pltxti{Polska Partia Socjalistyczna}~--- \pltxti{Frakcja Rewolucyjna}) jeb ``fraki'', kā to parasti sauca, ar J.~Pilsudski priekšgalā. Tajā no 46~tūkstošiem PPS biedru pārgāja ap 14~tūkstošu. Pēdējos atbalstīja arī Galīcijas sociāldemokrāti. ``Revolucionārā frakcija'' aizstāvēja bruņotas cīņas un nacionālās sacelšanās organizācijas idejas. PPS-frakcijas īpatnība bija tā, ka sociāli-politiskās programmas ierobežotais raksturs tajā apvienojās ar dumpīgumu, radikālismu, kas bija raksturīgi iepriekšējo sacelšanos tradīcijām, šļahtičiem raksturīgu ``varoņu'' kultu. Pēc būtības ``revolucionārajā frakcija'' kļuva par sīkburžuāzisku militāri-politisku organizāciju ar uz nacionālās atbrīvošanās cīņu orientētu virzību. Vadoties no tās darbības, V.~Ļeņins rakstīja: ``Fraki'' ir nevis proletāriska, sociālistiska, bet sīkburžuāziska, nacionālistiska partija''. Jāatzīmē, ka no ``fraku'' kaujas vienību dalībnieku vidus nāca daudzi vēlākie augstākie Polijas valsts darbinieki. Tomēr ``fraku'' ietekme masās ievērojami kritās. Ja 1907.~gadā to rindās bija ap 14~tūkstošu cilvēku, tad 1909.~gadā~--- vairs tikai ap 3~tūkstoši. 1911.~gadā ``fraku'' vidū izveidojās opozīcija, tomēr, sākoties karam 1914.~gadā, tā atgriezās ``revolucionārās frakcijas'' rindās.

PPS IX kongresa vairākuma pozīcijas stāvošie pieņēma nosaukumu Polijas Sociālistiskā partija~--- ļevica (\pltxti{Polska Partia Socjalistyczna~--- Ļevica}), kura stāvēja tuvu Polijas karalistes un Lietuvas sociāldemokrātijai (daļa no PPS-ļevicas biedriem 1909.~gadā arī iestājās PKunLSD). No PPS uz PPS-ļevicu pārgāja vairākums biedru. PPS-ļevica nevēlējās atbalstīt bruņotu teroristisku cīņu, turpināja darboties patstāvīgi galvenokārt arodbiedrībās, dažādās kultūrizglītības biedrībās un tamlīdzīgās organizācijās, vēlāk~--- jau Pirmā pasaules kara laikā ieturēja internacionālistisku pozīciju, 1918.~gada decembrī, apvienojoties ar PKuLSD, izveidoja Polijas Komunistisko strādnieku partiju (\pltxti{Komunistyczna Partia Robotnicza Polski}), pie tam nacionālā jautājuma atrisināšanu atvirzot otrajā plānā.

Strādnieku kustība zem sociālistiskiem lozungiem, plašais terora izvirdums nobiedēja konservatorus. ``Reālisti'' atklāti vērsās pie Krievijas varas iestādēm, aicinot apspiest revolucionāro kustību. Līdz tam viņi nebija tik tālu gājuši vēlmē sadarboties ar impērijas varu. Viņus atbalstīja arī katoļu baznīca, kura baidījās no sociālistu vidū izplatītās ateistiskās domāšanas.

R.~Dmovskis griezās pie Krievijas valdības vadītāja S.~Vittes ar priekšlikumu piešķirt Polijas karalistei autonomiju, brīdinot, ka citādi neizdosies pārvarēt anarhiju. Taču Krievijas varas iestādes labi atcerējās XIX gadsimta 60.~gadu sākumu, kad zināmas piekāpšanās noveda pie nacionālistiskās kustības eskalācijas.

No 1906.~gada vidus par strādnieku šķiru cīņa arēnu atkal kļuva Lodza. Šeit notika virkne ekonomisku un politisku streiku, kā arī no 1906.~gada decembra līdz 1907.~gada martam norisa masu lokauts (uzņēmumu slēgšana un strādnieku masveida atlaišana, lai uzņēmēji varētu uzspiest tiem savus noteikumus.) Fabrikanti izvēlējās ziemu, kad strādnieki visvairāk cieta no sala un bada, lai piespiestu tos padoties. Ieganstu viņi atrada kādas strādnieku grupas konfliktā ar inženieri vienā no pilsētas rūpnīcām. Tās direkcija slēdza rūpnīcu un darba atjaunošanai uzstādīja strādniekiem jaunus darba noteikumus un paziņoja par 98~aktīvistu atlaišanu. Kad strādnieki atteicās piekrist, 17.~decembrī direkcija pasludināja lokautu, bet 29.~decembrī to pieteica visas lielākās Lodzas firmas. Bez darba izrādījās ap 30~000 strādnieku. Viņu biedri visā Polijas karalistē vāca līdzekļus to atbalstam. Tomēr pēc triju mēnešu cīņas, strādnieku spēku izsīkumu rezultātā lokauts beidzās ar fabrikantu uzvaru, strādnieki pieņēma viņu noteikumus. Tomēr kopumā poļu strādnieki streiku cīņu turpināja.

Krievijas varas iestādes no 1905.~gada sākuma gatavoja \strong{Krievijas Valsts domes} sasaukšanu. Izstrādājot tās vēlēšanu noteikumus, etniskajā jautājumā sadūrās divas nostājas. Viena~--- vai nu pielietot vienotas (vai gandrīz vienotas) normas visos pēc etno-konfesionālajām un kultūr-ekonomiskajām pazīmēm dažādajos reģionos vai, otrādi, katrā reģionā ieviest īpašus noteikumus. Notikumiem attīstīties varas iestādēs nostiprinājās viedoklis, ka dažādās nekrievu tautas, to pārstāvniecības Domē būs opozīcijas pusē. Pats Nikolajs II prasīja, lai, īstenojot reformu, ``tiktu nodrošināts, lai katrs krievs visā impērijā varētu justies savās mājās'' un uzsvēra nepieciešamību izdalīt ``krievu nacionālo kodolu'' (t.~i.~--- dot tam privilēģijas). Sākotnēji bija paredzēts, ka par Domes deputātiem nevarēs kļūt cilvēki, kuri ``nezin krievu valodu'' un ``neprot lasīt un rakstīt krieviski'', taču, ņemot vērā Iekškrievijas guberņu zemnieku analfabētismu, no otrās prasības nācās atteikties. Īpaši tika skatīts jautājums par vēlēšanu tiesību piešķiršanu ebrejiem. Bija cilvēki, kuri prasīja ``kaitīgās ebreju nācijas nepielaišanu Domē'', tomēr pēc S.~Vittes priekšlikuma ebrejiem piešķīra gan pasīvās, gan aktīvās vēlēšanu tiesības, jo tika cerēts, ka Domē iekļūs tikai pāris ebreju, kas nevarēs ietekmēt pārējo deputātu masu. (Patiesībā I Valsts domē bija 13 ebreju deputātu, bet II Domē~--- 6, tikai pēc vēlēšanu likuma izmaiņas III un IV Domē bija pa trijiem ebreju deputātiem.) Poļu apdzīvotajām guberņām tika nolemts piešķirt tādu pat pārstāvniecību kā Eiropas Krievijas guberņām, jo tajās bija augsts izglītības un saimnieciskās attīstības līmenis. Polijas karalistei tika piešķirtas 36 deputātu vietas. Eiropas Krievijā viens deputāts bija vidēji uz 227~800, Polijas karalistē~--- uz 254~100 iedzīvotājiem. Pēc 17.~oktobra manifesta izdošanas Polijas karaliste ieguva vēl 2 deputātu vietas. (Pavisam Domē bija 524 vietas). Nacionāldemokrātiskā partija cara 17.~oktobra manifestu un Domes vēlēšanas uztvēra kā izšķirošu panākumu, nosodīja sociālistus par streiku turpināšanu. Daudzi poļu pilsonības pārstāvji uzskatīja, ka Krievijas un Polijas attiecībās sākas jauns posms. Arī krievu vēsturnieks A.~Pogodins 1915.~gadā rakstīja, ka pati dzīve pierādījusi, ka pēc 1863.--1864.~gada poļu sacelšanās Krievijas valdības izdarītie secinājumi bijuši ``nemērķtiecīgi un politiski nepareizi'', bet 1905.~gada 17.~oktobra manifests iezīmējis jaunu robežlīniju, ``lai arī cik nemanāma būtu pāreja''.

1906.~gada martā notika vēlēšanas uz \strong{I Valsts domi}. Sociālisti (PPS, KSDSP, PKunLSD u.c.) tās boikotēja (boikotu atbalstīja tikai lielāko rūpniecības centru proletārieši un tas nespēja traucēt vēlēšanas). I~Dome izrādījās zemnieciska (vairāk kā 50\% deputātu bija tuvi zemniecībai) un daudznacionāla (vairāk nekā 40\% bija nekrievi). Nacionāldemokrāti ar līderiem R.~Dmovski un V.~Grabski priekšgalā izcīnīja 34 no iespējamajiem 38 mandātiem. Vēl 19 poļu tautības deputātus ievēlēja no Lietuvas, Baltkrievijas un Ukrainas.

Par poļu frakcijas (\pltxti{kolo}) priekšsēdētāju Domē kļuva R.~Dmovskis. Frakcija aizstāvēja Polijas autonomiju Krievijas impērijas sastāvā, kas paredzēja poļu valodu kā oficiālo, ticības brīvību un katoļu baznīcas tiesības, no Polijas iesaukto tiesības miera laikā dienēt savā zemē, patstāvīgu karalistes budžetu. Endeki piedāvāja bijušās Polijas-Lietuvas valsts teritorijā dzīvojušajiem ``leišiem un rusīniem'' atbalstīt endeku programmu, citādi solot izvērst pret tiem politisku cīņu. Viņi draudēja arī izspiest ebrejus no sabiedriskās dzīves. Kā uzsver krievu vēsturnieki, ksenofobijas un antisemītisma elementi endeku programmā ieņēma ievērojamu vietu. R.~Dmovskis un viņa piekritēji gribēja radīt tādu Polijas valsti, kurā etniskie poļi būtu pārsvarā, bet nacionālie mazākumi būtu pietiekami vāji, lai tos varētu asimilēt. Tāpēc, pēc endeku viedokļa, poļiem nevajadzēja cīnīties par Polijas atjaunošanu 1772.~gada robežās, jo tad šie nacionālie mazākumi būtu pārāk lieli. R.~Dmovskis 1906.~gadā pasludināja ``bezierunu'' atbalstu ``neoslāvismam'', cenšoties nostiprināt slāvu ietekmi Austroungārijā un tuvināt to ar uz konstitucionālisma ceļa nostājušos Krieviju. R.~Dmovskis mēģināja nodibināt savienību ar krievu kadetiem, bet reizē viņš arī nevēlējās, lai Krievijas varas iestādes poļu nacionāldemokrātus identificētu ar krievu liberālo opozīciju.

Poļu deputāti I Valsts domē varēja jau izmantot pieredzi, kuru bija guvušas poļu frakcijas Vācijas Reihstāgā un Austrijas Reihsrātā. Poļu pārstāvji Domē veidoja Autonomistu savienības kodolu. Savienībā bija 63~deputāti, bet kopā ar tādiem, kuri vienlaikus ietilpa citās frakcijās un arī autonomistu grupā, viņu bija ap 120~cilvēku. Savienības galvenais mērķis bija panākt Krievijas nacionālo reģionu pārvaldes decentralizāciju. No nacionālajiem jautājumiem kā pirmais Domē tika izvirzīts tieši poļu jautājums. 1906.~gada 30.~aprīlī 27~poļu deputāti iesniedza Domei dokumentu, kurā, atsaucoties uz 1815.~gada Vīnes kongresa aktu un Valsts Pamatlikumiem, argumentēja nepieciešamību atjaunot Polijas karalistes tiesības un autonomiju.

Jāuzsver, runa gāja tikai par autonomiju, bet nevis neatkarību. Kaut Krievijā pastāvošais administratīvais nacionālisms nekrievu tautās izsauca negatīvu reakciju, ātrā rūpniecības, tirdzniecības, satiksmes ceļu attīstība valstī tās attīstītāko nomaļu~--- Polijas, Baltijas, Somijas atdalīšanos no Krievijas padarīja ekonomiski gauži neizdevīgu. Poļu \pltxti{kolo} Domē iesniegtais likumprojekts pasludināja: ``Polijas karaliste kā neatņemama Krievijas valsts sastāvdaļa savā iekšējā dzīvē tiek pārvaldīta ar īpašiem nolikumiem uz īpašas likumdošanas pamata''. Prasot radīt Polijā Seimu ar likumdošanas tiesībām, īpašu valsts kasi un tiesu iekārtu, likumprojekts paredzēja visas impērijas kompetencē atstāt ``pareizticīgās baznīcas, ārlietu, armijas, flotes, monetāros jautājumus; likumdošanu muitas, akcīzes, pasta, dzelzceļu satiksmes, preču zīmju un privilēģiju, literārā un mākslas īpašuma, valsts aizņēmumu un saistību, dumpju pret centrālo varu un valsts nodevību lietās''. Šajā laikā Polijas nacionālie līderi bija gatavi aizmirst divu nacionālo sacelšanos rūgtumu, lai piedalītos sekmīgajā daudznacionālās Krievijas saimnieciskajā un pēc 1905.~gada revolūcijas arī cerīgajā politiskajā attīstībā.

Nacionālajiem reģioniem aktuāli bija arī agrārais jautājums, kadetu un nacionālo partiju ierosinājums atcelt nacionālos un reliģiskos ierobežojumus. Likumprojektu prezentējušais kadets F.~Kokoškins nosodīja daudzos ``poļu izcelsmes'' personu ierobežojumus, norādot, ka likums pat nepasaka to, kā var pierādīt, ka konkrētā persona ir tieši ``poļu izcelsmes''. Tādejādi I Dome sāka plaši iztirzāt cariskās valdības nacionālās neiecietības un ksenofobijas politiku. Taču jau 1906.~gada jūnijā I Valsts dome tika atlaista.

Jāatzīmē, ka I Domē Poļu \pltxti{kolo} neatbalstīja kadetu agrāro likumprojektu, pēc tās atlaišanas neparakstīja t.s. Viborgas uzsaukumu, kurš aicināja sabiedrību uz pilsonisko nepakļaušanos (\rutxti{гражданское неповиновение}) varai~--- nemaksāt nodokļus, nepakļauties iesaukumam armijā u.t.t., ar to iegūstot zināmā mērā labvēlīgu attieksmi valdības sfērās.

\strong{II Domes} (tā darbojās no 1907.~gada marta līdz jūlijam) vēlēšanās no poļu partijām nepiedalījās vienīgi PPS. PKunLSD uzskatīja par nepieciešamu izmantot vēlēšanu kampaņu aģitācijas un propagandas nolūkiem. Ievēlētajā 34~cilvēku lielajā II Domes poļu deputātu grupā dominēja nacionāldemokrāti. Viņi ieguva 29~mandātus. Kopā ar Lietuvas un Baltkrievijas Teritoriālo frakciju viņi kopā izveidoja 46~deputātu lielu grupu. Poļu kolo prezīdiju vadīja nacionāldemokrāts R.~Dmovskis. Jaunā Krievijas premjera P.~Stolipina valdība mēģināja piedāvāt II Domei līdzdarboties, īstenojot valdības sociāli ekonomisko kursu. Reizē P.~Stolipins savā deklarācijā pacentās izvairīties no nacionālā jautājuma atzīšanas Krievijā, ignorēja nacionālās pašvaldības un autonomijas jautājumus. Taču, uzsverot pareizticīgās baznīcas lomu, viņš norādīja, ka tai ir jāsaņem no valsts puses īpaša aizsardzība. Faktiski tas nozīmēja citu reliģiju, konfesiju noniecināšanu. Valdības vadītājs solīja Krievijas Rietumu apgabalā un Polijā ieviest pašpārvaldi ``ar dažām vietējām īpatnībām'', mainot dažas administratīvās robežas, izdalot īpašā vienībā tīri krieviskos iedzīvotājus, ``kuriem ir savas specifiskās intereses''. Šo programmu P.~Stolipins mēģināja realizēt kā II, tā arī III Valsts domē.

II Domē poļu deputātiem nācās karsti diskutēt ar pareizticīgās garīdzniecības pārstāvjiem. Tā apspriežot likumprojektu par sodu atcelšanu par ``slepenu apmācību'' (t.i.~--- nelegālu skolu uzturēšanu) bīskaps Jevlogijs apvainoja poļu skolas ``tīri krievisko iedzīvotāju polonizācijā''. R.~Dmovskis uz to norādīja, ka visas tautas vēlas, lai to bērni varētu mācīties dzimtajā valodā, bet Polijā šis jautājums atšķiras no citiem apgabaliem ar to, ka valdības uzraudzītās poļu sākuma, vidējās un augstākās skolas bija iznīcinātas. Tiesa, kā reakcija uz 1905.~gadā notikušajiem ``skolu streikiem'', bija atļauts dibināt privātskolas ar poļu mācību valodu. Nodibinājās Skolu savienība (\pltxti{Macierz Szkolna}), kurā bija ap 100~000 biedru un kura atvēra ap 800~sākumskolu. Asas debates izvērsās, apspriežot Izglītības ministrijas budžetu. Izrādījās, ka, kaut Krievijas Ziemeļrietumu novadā bija atļauta poļu un lietuviešu valodas mācīšana, taču par to bija jāmaksā skolēnu vecākiem. Kad P.~Stolipins mēģināja atcelt virkni ierobežojumu ebrejiem, Nikolajs II viņam to neatļāva darīt, norādot: ``Iekšējā balss man arvien noteiktāk saka, lai es neuzņemtos pieņemt šo lēmumu''. Daudzus likumprojektus, kas paredzēja likvidēt nacionālo nevienlīdzību, iesniedza opozīcijas partijas. Arī agrārās reformas nepieciešamības apspriešanas gaitā tika skarts Polijas jautājums. Kāds zemnieku deputāts [A.~Krakovskis] no Rietumu apgabala R.~Dmovska uzstāšanos ar prasību piešķirt pašvaldības tiesības Polijas karalistei iztēloja kā ``poļu muižnieku galveno trumpi, lai pakļautu savai ietekmei pareizticīgos zemniekus''. Viņš arī apgalvoja, ka ieguvuši kontroli pār vietējām skolām, ``poļi to izmantos, lai ``pavedinātu pareizticīgos katolicismā''. Tā II Domē skaidri parādījās dažādas nostājas Polijas jautājumos. Valdība vēlējās nacionālo jautājumu novirzīt reliģiski-konfesionālā gultnē, novēršot virkni nacionāli reliģisko problēmu. Domes opozīcija cīnījās par nacionālo un reliģisko vienlīdzību, paplašināt reģionālu vietējo pašpārvaldi.

Kopumā II Domē poļu deputāti izraisīja valdības neapmierinātību, jo viņu balsīm dažkārt izrādījās izšķiroša nozīme. Pēdējās II Domes pastāvēšanas dienās poļu deputāti iesniedza likumprojektu par Polijas autonomiju. Valdība kļuva pret viņiem atklāti naidīga.

Poļi bija pārstāvēti arī Krievijas Valsts padomē, kura sastāvēja kā no vēlētiem, tā cara nozīmētiem locekļiem. Vēlētos Valsts padomes locekļus ievēlēja uz 9 gadiem, ik pēc 3 gadiem notika viņu rotācija. Ar 1906.~gada 20.~februāra dekrētu Polijas muižnieki ieguva tiesības ievēlēt 6~Padomes locekļus, kas netika atļauts, piemēram Aizkaukāza muižniecībai. Kā vērtēja padomju vēsturnieks A.~Stepanskis, tas liecināja par carisma vēlmi izlīgt ar poļu muižniecību. Arī no deviņām Lietuvas, Baltkrievijas un Labā krasta Ukrainas guberņām, kurās pārsvarā bija poļu muižnieku zemes īpašums, Valsts padomē bija pārstāvēti pēdējie. Pavisam Padomē bija 18~poļu, gandrīz 10\% no visiem tās locekļiem, dažkārt viņu balsis varēja izšķirt balsojumu iznākumu. Ja I un II Valsts domē poļu deputātu vidū vairākums bija nacionāldemokrātiem, tad Valsts padomes poļu locekļi kopumā pārstāvēja t.s. ugodoviešu jeb ``reālistu'' virzienu. Kopumā šī grupa Valsts padomē uzstājās kā vienots vesels, sākumā parasti pieslējās oktobristiski-melnsimtnieciskajam (melnsimtnieki~--- Krievijas galēji labējo organizāciju pārstāvju kolektīvais nosaukums 1905.-1917.~gadā, kas uzstājās ar patvaldības un pareizticības stiprināšanas saukļiem) Padomes vairākumam. Tā poļu grupas līderis I.~Korvins-Miļevskis aizstāvēja oktobristu un melnsimtnieku tēzi par to, ka agrārā jautājuma risinājums jāmeklē ne muižnieku zemes atsavināšanā, bet zemnieku saimniecību intensifikācijā. Taču labējo polonofobija drīz piespieda poļus norobežoties no tiem. Poļi Valsts padomē atrada kopēju valodu ar t.~s. ``centru'', kaut arī tas noraidīja pat Polijas autonomijas lozungu. Kaut kādus demokrātiskus risinājumus no saviem Valsts padomes locekļiem poļu tautai nenācās gaidīt ne revolūcijas gados, ne arī pēc tās.

Var teikt, ka 1904.--1907.~gada notikumi Polijas karalistē kļuva par XX gadsimta sākuma poļu politiskās vēstures kvintesenci. Notika dabīga politisko ideju un organizāciju atlase, kuras izteica atsevišķu sociālo un nacionālo grupu intereses un centienus. Vēsturiskajā arēnā iznāca šķiras un partijas, kuras pusgadsimta garumā turpmāk noteica Polijas vēstures gaitu. Revolūcija atklāja un saasināja ideju konfliktus, atšķirības starp galvenajiem politiskajiem spēkiem.

Notikumu Polijas karalistē un visā Krievijas impērijā ēnā palika procesi, kas norisa šai laikā Pozenē un Galīcijā, kaut arī šeit XX gadsimta sākumā dzīvoja tikai nedaudz mazāk iedzīvotāju nekā Polijas karalistē. (1900.~gadā karalistē dzīvoja ap 10~miljonu, bet Austroungārijai pievienotajās poļu zemēs~--- 7,6~miljoni un Vācijas valdījumos~--- 7,3~miljoni.) Tomēr 1905.~gada revolūcija Krievijā guva atskaņas arī citās poļu zemēs.

Vācu iestādes \strong{Pozenes} provincē XX gs. sākumā turpināja agrāk aizsākto poļu ierobežošanas politiku. Kanclers B.~Bīlovs runāja par ``sociālistiskajām un poļu briesmām Prūsijai un impērijai''. Sarunā ar žurnālistiem B.~Bīlovs salīdzināja poļus ar trušiem, kuru ``vairošanos ir jāapstādina''. Pretpoļu nostāju atbalstīja arī ķeizars Vilhelms II.

Spilgtākais tās izpaudums bija 1901.~gadā notikušais t.s. ``Vrešenas (vācu \detxti{Wreschen}, poļu \pltxti{Wrzesnia}, apdzīvota vieta, kurā no 5~000~iedzīvotājiem 70\% bija poļi) incidents''(vācu \detxti{Wreschener Schulstreik}, poļu \pltxti{Strajk dzieci wrzesińskich}). Tika izdots rīkojums, ka poļu bērniem ticības (reliģijas) mācībā jāizmanto tikai vācu valodā iespiestas grāmatas un jāatbild vāciski. 118~skolnieki un skolnieces priestera J.~Laskovska vadībā atteicās pieņemt vācu valodā drukāto katķismu un atbildēt vāciski. Izglītības inspektors pavēlēja bērnus nopērt. Vecāki pieprasīja atcelt sodus. Nonāca līdz nemieriem. 18~Vrešenas pilsoņu--poļu par miera traucēšanu tika notiesāti uz mēnešiem ilgu apcietinājumu. Šis Vrešenas ``skolu streiks'' izsauca plašu sabiedrības rezonansi. Būtiski sabojājās attiecības vietējo iedzīvotāju~--- poļu un vāciešu starpā. Varas pārstāvji nevēlējās padoties. Prūsijas iekšlietu ministrs H.~f.~Hammeršteins 1904.~gadā paziņoja: ``Poļu jautājums mūs pārstās interesēt tikai tad, kad pie mums nebūs poļu, bet būs tikai prūši, kas runā poliski. Manuprāt, mūsu pacietība ir pārāk liela. Mums nav darīšana ar līdzīgu pretinieku, mūsu lieta~--- pavēlēt, viņu~--- paklausīt''.

1905.~gada pavasarī un vasarā notika virkne strādnieku streiku Pozenes pilsētās. Rudenī, kad Krievijas revolūcijā bija vērojams augstākais pacēlums, Vācijas varas iestādes poļu apdzīvotajās zemēs koncentrēja karaspēku. Situācija atkal saasinājās 1906.~gada pavasarī. Lejassilēzijas pilsētas Breslavas (vācu \detxti{Breslau}, poļu \pltxti{Wrocłow}) metalurģisko rūpnīcu īpašnieki pasludināja lokautu, atlaižot ap 10~000 strādnieku. Viņi organizēja plašu demonstrāciju, pret kuru tika raidītas policijas vienības. Simti strādnieku tika ievainoti, no tiem vairāki desmiti~--- smagi.

Kad 1906.~gadā vēl 200~tautskolas pārgāja uz reliģijas mācības pasniegšanu vācu valodā, 1906./1907.~gadā Pozenes provincē notika jauns liels ``skolu streiks''. Streiks šoreiz vairs nebija lokāli ierobežots, izplatījās ne tikai visā Pozenes provincē, kurā toreiz mācījās 370~000 skolēnu, (No tiem 118~000 atzina par dzimto vācu valodu, bet 240~000~--- poļu valodu, bet ap 20~000 nevarēja izšķirties par vienu no tām, jo bija no jauktajām ģimenēm), bet arī Rietumprūsijā un Augšsilēzijā, kur tika lietotas abas~--- poļu un vācu valodas. 755~Pozenes provinces skolās no ap 90~000 skolēnu streikoja gandrīz 47~000. No 1906.~gada rudens līdz 1907.~gada beigām Prūsijas tiesas poļu rietumu zemēs bija pārslogotas ar lietām, ierosinātām pret poļu nacionālās kustības dalībniekiem, kuri atbalstīja skolu streiku. Pasaules sabiedriskā doma stāvēja poļu pusē, jo visi saprata, ka sākumskolā ir jāmācās dzimtajā valodā. Poļu dzejniece M.~Konopņicka, kura bija aktīvi piedalījusies jau Vrešenas streika izraisītajos notikumos, 1908.~gadā uzrakstīja dzejoli ``\pltxti{Rota}'' (``Zvērests''), kura pirmās rindas skanēja:

\vspace{1.5em}

\noindent
\begin{minipage}{0.5\textwidth}
\pltxti{Nie rzucim ziemi skąd nasz ród,\\
Nie damy pogrześć mowy,\\
Polski naród, polski lud,\\
Królewski szczep piastowy.\\
Nie damy by nas zniemczył wróg!\\
Tak nam dopomóż Bóg!\\
Tak nam dopomóż Bóg!}
\end{minipage}
\hspace{1em}
\begin{minipage}{0.5\textwidth}
Neatdosim zemi, kurā dzimām,\\
Nemirs mūsu valoda,\\
Polijas mēs bērni, poļu tauta,\\
Karaliskā Pjastu celma.\\
Ienaidnieks mūs nepārvācos!\\
Dievs mums palīdzēs!\\
Dievs mums palīdzēs!
\end{minipage}

\vspace{1.5em}

% page 207


Ar komponista F.~Novovijska mūziku tas kļuva par iemīļotu dziesmu, kura līdz 1918.~gadam kalpoja kā poļu skautu kustības himna un tika arī apspriesta kā iespējamā Polijas valsts himna.

Otrs sāpīgs punkts XX gadsimta sākumā vācu un poļu attiecībās radās ar t.s. ``ugunsdrošības'' likuma pieņemšanu. 1904.~gadā Prūsijas landtāgs, bet pēc tam arī Vācijas reihstāgs pieņēma likumu, kas deva tiesības vietējai administrācijai nepieļaut ēku celtniecību jau iepirktajā zemē. Likums bija vērsts uz to, lai ierobežotu jaunu poļu tautības zemnieku saimniecību ierīkošanu. Tas izsauca gan traģiskas, gan komiskas sekas.

Augšsilēzijā poļu kalnracis, darba invalīds F.~Hroščs, neraugoties uz vairākiem lūgumiem atļaut viņam ierīkot mitekli, saņēma atteikumu. Tad viņš uz sev piederoša maza zemes gabaliņa uzcēla šķūni, zem tā izraka dziļu pagrabu, kurā arī apmetās. Kad vietējās varas iestādes lika viņam no pagraba izvākties, viņš to atteicās darīt. Kad ieradās žandarmi un sāka graut pagrabu, F.~Hroščs vienu no tiem nogalināja, bet pēc tam, pārējo aplenkts, nošāvās arī pats.

Pozenes provincē poļu zemnieks M.~Džimala, kuram tāpat tika aizliegts uz savas 1904.~gadā iepirktās zemes ierīkot dzīvojamo māju ar apkuri, lai apietu likumu, sāka dzīvot šķūni, bet pēc tam nopirka pārvietojamu cirka furgonu, kuram bija izcirsti logi un ierīkota dzelzs krāsniņa. Taču ierēdņi noteica, ka furgons, ja tas atrodas vienā vietā, nav pārvietošanās līdzeklis, bet gan māja un M.~Džimalam tā jāatstāj. Tad nu viņš ik pēc dienas pārbīdīja savu mājokli un turpināja tajā dzīvot un saimniekot uz savas zemes. Laikam ejot, auga viņa popularitāte, cilvēki no visām Polijas malām sāka viņu apmeklēt, lai paskatītos uz īpatnējo dzīvojamo ēku. Drīz viņš par sabiedrības saziedotiem līdzekļiem varēja nopirkt jaunu, labāku furgonu un smīnēdams skatījās, kā viņu pašu un viņa dzīves vietu fotografēja un zīmēja karikatūrās dažādu laikrakstu korespondenti. Sakarā ar šo gadījumu par prūšu varas iestādēm zobojās ne tikai poļi visās trijās kaimiņvalstīs, kurās tie dzīvoja, bet arī ārzemnieki. V.~Džimalas piemēram sekoja arī citi zemnieki. To skaits kļuva tik ievērojams, ka kāds Pozenes fabrikants uzskatīja par izdevīgu uzsākt speciāli tādu furgonu ražošanu un pārdošanu poļu zemniekiem. Varas iestādes nelikās mierā, lika aizvākt no furgona tajā iebūvēto krāsni, jo tās atrašanās pārvietošanās līdzeklī neatbilda priekšrakstiem. Pēc tam viņam pieprasīja iznest no furgona mēbeles un guļam piederumus, jo ``guļamistabas'' ierīkošana transporta līdzeklī arī bija ``nelikumīga''. ``Tiesas darbi'' ilga vairākus gadus, atnesa poļu zemniekam atpazīstamību arī ārzemēs. Par viņu iestājās arī tādi slaveni rakstnieki kā H.~Senkevičs, Ļ.~Tolstojs. Tomēr pēc vairākus gadus ilgiem tiesas darbiem un arī aresta M.~Džimalam nācās kapitulēt un 1909.~gadā savu zemi pārdot. Tiesa, par saziedotiem līdzekļiem viņš kaimiņu ciemā varēja nopirkt citu zemes gabalu ar jau esošu dzīvojamo ēku, pret kuru iestādēm nevarēja būt pretenziju. Vēlāk viņš pameta Pozeni un pārcēlās uz Galīciju. Viņa otrais furgons tika pārvests uz Krakovu un izstādīts muzejā kā liecība poļu cīņai pret prūšu patvaļu. Pēc Pirmā pasaules kara jaunā Polijas valsts M.~Džimalam kā nacionālam varonim Poznaņas novadā uzdāvināja saimniecību un piešķīra goda renti. Pēc nāves viņš tika apbalvots ar Polijas atdzimšanas 5.~pakāpes ordeni, viņa vārdā par \pltxti{Drzymałowo} tika nosaukts Podgradovices ciems. Otrā pasaules kara laikā vācu okupanti iznīcināja viņa kapa vietu, mūsdienās tā ir atjaunota.

Jau agrāk pieminētais Vācijas valstiskais atbalsts poļu muižnieku zemju uzpirkšanai un piešķiršanai vācu pārceļotājiem izsauca poļu pretestību, tie dibināja bankas, kas arī nodarbojās ar šādu zemju uzpirkšanu un tālāku pārdošanu poļiem. Karaliskā Prūsijas pārceļošanas komisija reāli varēja iepirkt zemi tikai jau no vācu tautības īpašniekiem. Berlīnē radās ideja dot komisijai tiesības īpašos gadījumos poļu lielīpašnieku zemes ne tikai pirkt, bet pat atsavināt. Strīdi par šāda pastāvošā likuma papildinājuma lietderīgumu ilga trīs gadus, līdz beidzot 1908.~gadā tas arī tika pieņemts. Vācu-poļu attiecības vēl vairāk saasinājās. Vācijas impērija viennozīmīgi pārkāpa savu pilsoņu vienlīdzības principu. Tikai 1912.~gadā likuma papildinājums tika arī piemērots, kopā gan tikai četros gadījumos. Pret atlīdzību tika atsavināti 1~700 ha poļu īpašnieku zemes. Pēc vācu pētnieka L.~Dralles viedokļa precedents nodarīja lielu ļaunumu vācu un poļu attiecībām pirms Pirmā pasaules kara sākuma. Ne tikai Pozenē, bet arī Rietumprūsijā un Silēzijā vācieši un poļi dzīvoja savstarpējā naidā.

Diemžēl, ne vienmēr apspiestie poļi spēja saglabāt objektivitāti savos pārmetumos vāciešiem. Vācu nacionālisti tāpēc publicēja 400 lappušu plašu t.s. ``Poļu spoguli'', ar apakšvirsrakstu ``Poļu musināšana viņu pašu presē'' (``\detxti{Polenspiegel: Die Umtriebe Der Polen Nach Ihrer Eigenen Presse}''), kur bija pārpublicēti poļu preses izteikumi par vāciešiem. Piemēram, tur ar atsaucēm uz konkrētiem laikrakstiem bija atrodami brīdinājumi: ``Neviens polis nedrīkst precēt vācu vai kādu citu svešu meiteni, tas ir nāves grēks'', ``Neviens poļu katolis nedrīkst apmeklēt vācu sprediķi, jo ar to viņš atbalstītu savus niknākos ienaidniekus'', ``No negaisa, bada nāves, sērgām, ugunsgrēka, no ātras, ļaunas nāves bez grēksūdzes un no ģermanizācijas pasargi mūs, Dievs'' utt. Vāciešiem šī publikācija rādīja, ka ``tādi ir poļi''.

Vācietis no Toruņas R.~Vingendorfs 1939.~gadā Vācijā izdotā grāmatā gan apgalvoja, ka Prūsijas poļi nacionālajā jautājumā aprobežojušies vienīgi ar reliģijas brīvības, poļu avīžu izdošanas un poļu elementu iekļaušanas Prūsijas pārvaldē prasībām, pat nesapņojuši par Polijas valsti, taču tūlīt bija spiests pieminēt, ka katoļu garīdzniecība vairāk nekā gadsimtu uzturējusi cerības uz Polijas kā katolicisma bastiona atjaunošanu. Pārējiem Prūsijas poļiem šī sapņa īstenošana pēc Vācijas zaudējuma karā 1918.~gadā nākusi pilnīgi negaidīti. Turpretī R.~Dmovskis jau 1908.~gadā rakstīja, ka Prūsijai piederošajos [poļu] apgabalos poļu tautas masu garīgais līmenis ir augstāks nekā kur citur, un nacionālā enerģija ir sasniegusi augstāko pakāpi. Lai nu kā, ar vai bez tālākiem nākotnes plāniem, tomēr poļu nacionālā kustība turpinājās arī Vācijas poļu vidū.

\strong{Galīcija} XX gadsimta sākumā rūpnieciskās attīstības ziņā piederēja pie atpalikušākajām Austroungārijas provincēm. Visas Galīcijā saražotās rūpnieciskās produkcijas vērtība 1900.~gadā sastādīja tikai 2,14\% no visas Austroungārijas impērijas ražotās rūpnieciskās produkcijas vērtības. Galīcijas rūpniecībā, ieskaitot mājražošanu, XIX gadsimta beigās bija iesaistīta tikai 1/10 novada iedzīvotāju. Arī rūpniecības apvērsums šeit bija sācies tikai XIX gadsimta beigās, un noritēja tas lēnām. Galīcijas rūpniecībā darbojās pārsvarā sīki uzņēmumi. Provinces teritorija sastādīja 28\% un tajā arī dzīvoja 28\% Cisleitānijas (Austroungārijas austriešu daļas) iedzīvotāju, bet tajā bija tikai 15,8\% visu Cisleitānijas rūpniecības uzņēmumu, kuri nodarbināja 10,1\% tās strādnieku. 1902.~gadā Galīcijas rūpniecībā un amatniecībā kopā bija nodarbināti 336~000 cilvēku, taču fabriku strādnieki bija tikai 33~000--50~000.

Poļu kapitāls šeit uzstājās kā austriešu un ārzemju kapitāla jaunākais partneris. Lielu vietu Galīcijas ekonomikā ieņēma iegūstošā rūpniecība. Šeit ieguva sāli, ogles, dzelzs, svina un cinka rūdu. Galīcijas īpatsvars naftas ieguvē XX gadsimta sākumā gan samazinājās. 1900.~gadā Galīcijā ieguva 1,18\% naftas pasaulē, taču naftas rūpniecībā no visa kapitāla poļiem piederēja tikai 0,5\%. Apmēram 1/3 Galīcijā iegūtās naftas šeit pat arī tika pārstrādāta, pārējo izveda galvenokārt uz Vāciju. Parafīna ražošanā Galīcija ieņēma monopolstāvokli.

XX gadsimta sakumā Galīcijā norisa virkne laukstrādnieku streiku. Ļudoviešu darbinieki 1903.~gada februārī partijas Galvenās padomes sēdē Žešovā (poļu \pltxti{Rzeszów}) pieņēma jaunu B.~Vislouha sagatavotu partijas programmu. Par kustības galveno principu tika pasludināta tiekšanās pēc tautas nacionāla, politiska, ekonomiska un kulturāla pacēluma. Programma saturēja demokrātiskas vēlēšanu reformas, Galīcijas autonomijas paplašināšanas, tās tiesību pielīdzināšanas Ungārijas karalistei prasības. Atzīstot rūpniecības attīstības nepieciešamību, uzsvars tomēr tika likts uz lauksaimniecību. Nacionālajā jautājumā tika atzīta poļu, lietuviešu, ukraiņu līdztiesība. Programmā bija runāts ne tikai par Galīciju, bet arī par Polijas karalisti un poļu zemēm Vācijas sastāvā, ar ko tika uzsvērta poļu tautas vienotība un izteikta nepieciešamība par neatkarīgas Polijas valsts atjaunošanu. Tika mainīts arī partijas nosaukums~--- uz ``\pltxti{Polskie Stronnictwo Ludowe''} (``\pltxti{Poļu Zemnieku partija}''). Ar to bija jāsaprot, ka partija darbosies visā Polijā.

Arī Galīcijā 1905.~gada revolūcija Krievijā izsauca sabiedriskās dzīves aktivizēšanos. Lembergā un citās pilsētās notika streiki, kuros ekonomiskās prasības sabalsojās ar politiskajām. To centrā atradās vispārējo vēlēšanu tiesību prasība. Sociālisti un ļudovieši centās panākt vēlēšanu demokratizāciju. 1905.~gada septembrī~--- oktobrī Krakovas ielās izgāja no 10 līdz 20~tūkstošiem strādnieku, bet 2.~novembrī demonstrācija izvērtās sadursmē ar policiju. 28.~novembrī~--- dienā, kad visā Austrijā norisa vispārējs streiks ar prasību realizēt vēlēšanu reformu, Galīcijā streikoja desmiti tūkstošu cilvēku, norisa plašas demonstrācijas, Lembergā vien ielās izgāja vairāk nekā 40~tūkstoši. Rezultātā 1907.~gada sākumā notika vēlēšanu reforma, tika likvidētas kūrijas, ieviestas vispārējas vēlēšanas uz Vīnes Reihsrātu, taču pēc patvaļīgi noteiktām nacionālajām kvotām. Tā, rutēņi, kuri pēc oficiālās statistikas sastādīja 13,2\% austriešu zemju iedzīvotāju, saņēma tikai 6,4\% deputātu vietu, no 115~Galīcijas deputātiem rutēņi bija tikai 33. Vēlēšanu reforma saasināja nacionālās pretrunas visā Austrijā, arī Galīcijā. Tomēr tai bija svarīga nozīme, tā deva triecienu konservatīvo kundzībai, to skaits parlamentā samazinājās. Taču vietējā Seimā, kura vēlēšanas reforma neskāra, konservatori joprojām saglabāja vairākumu.

\subsection{Poļu sabiedrība Pirmā pasaules kara priekšvakarā}

% page 210


1907.~gada 3.~jūnijā Krievijā notika valsts apvērsums, valdība padzina II Domi un mainīja vēlēšanu likumu, ko bez Domes piekrišanas tai nebija tiesību darīt. Ar šiem notikumiem var uzskatīt 1905.~gada revolūciju par noslēgušos. Revolūcijas viļņi norima arī Polijā. Ja 1907.~gadā Polijas karalistē notika 869~streiki ar 98~063~dalībniekiem, tad 1908.~gadā streiku skaits bija sarucis līdz 244 ar 31~709 dalībniekiem. 1909.~gadā notika tikai 67~streiki un tajos piedalījās 10~033~strādnieki. Vienlaikus pieauga represijas pret demokrātiskajiem un revolucionārajiem spēkiem. Tikai vienas nedēļas laikā no 1907.~gada 30.~jūnija līdz 6.~jūlijam Polijas karalistē tika arestēti 144~cilvēki no PKunLSD, PPS, Bunda u.c. revolucionāro organizāciju biedri.

1905.--1907.~gada revolūcijas rezultāti Polijā bija pretrunīgi. Norisa pozitīvas pārmaiņas. Kaut arī neizdevās gāzt monarhiju un iegūt Polijai autonomiju, kaut revolūcijā zaudēja dzīvību tūkstošiem cilvēku, tomēr tās gaitā salīdzinājumā ar pirmsrevolūcijas laiku tika izcīnītas carisma piekāpšanās kā ekonomiskajā, tā politiskajā sfērā, panākta darba apstākļu un tā apmaksas uzlabojumi rūpniecībā un lauksaimniecībā, tika mīkstināta valdības politika valodas un ticības jomā. Krievijas neveiksmes Krievu-japāņu karā un 1905.~gada revolūcija veda pie jūtamas rusifikācijas politikas vājināšanās Polijas karalistē.

Masu uzstāšanās aktivizēja lielas Polijas karalistes iedzīvotāju daļas, kura agrāk reti vai nemaz nepiedalījās politiskajā dzīvē, sabiedriski politisko apziņu. Aktivizējās cilvēki laukos. Strādnieku kustība demonstrēja savu masveida raksturu. Politika pārstāja būt elites lieta. Politiskās partijas savus uzsaukumus drukāja daudzu tūkstošu eksemplāru lielās tirāžās, pirmoreiz tik plaši izmantojot šo masu informācijas līdzekli. Jāatzīmē, ka revolūcijas dalībnieku rindas bija daudznacionālas. Plecu pie pleca cīnījās poļi un ebreji, krievi un vācieši. Nacionālajā jomā atkal atdzima poļu cerības uz nacionālo atbrīvošanos. Kaut daļa revolūcijas dalībnieku Polijas neatkarības ieguves jautājumu vēl atbīdīja uz nākotni, tomēr nacionālā jautājuma klātbūtne revolucionārajos notikumos ļauj runāt par tās nacionālās atbrīvošanās kustības raksturu.

1905.~gada revolūcija deva iespējas straujāk attīstīties poļu kultūras dzīvei. Tika mīkstināta cenzūra, atļauts dibināt poļu privātskolas. Nacionāldemokrātiem tuvu stāvošā Poļu skolu savienība (\pltxti{Polska Macierz Szkolna)} līdz savai atlaišanai 1907.~gada decembrī paguva nodibināt 800~pamatskolu un daudz bērnu dārzu, skolu bibliotēku. Radās daudz biedrību. (Lielāko vingrošanas biedrību ``\pltxti{Sokol}'' (``Vanags'') gan 1906.~gada septembrī aizliedza.) Sāka darboties daudz zinātnisko biedrību. Līdz Pirmajam pasaules karam poļu sabiedrība bija sasniegusi jau ievērojamu izglītības līmeni.

Taču, no otras puses, pēc 1905.~gada revolūcijas situācija Polijā zināmā mērā pasliktinājās. Saasinājās attiecības starp poļiem un cittautiešiem.

Pieauga konkurence starp Lodzas flaneļa (parķa) un Maskavas katūna ražotājiem, arī konkurence pašu Polijas ražotāju vidū, poļi naidīgi uztvēra veiksmīgākos vāciešus un ebrejus. Tāpat konkurēja vācu un poļu zemnieki. Saasinājās arī poļu inteliģences attieksme pret vāciešiem. Piemēram, 1907.~gadā tika izdota karte, kurā tika atzīmētas vācu kolonijas un vācu muižnieku īpašumi, lai visi redzētu no kurienes poļiem draud briesmas. Vācieši, kuri t.s. kongresa Polijā sastādīja tikai 4,2\% iedzīvotāju, tika tēloti kā ``grēku plūdi'', kā ``stratēģiskas briesmas'' Krievijas dzelzceļiem, kuru tuvumā parasti vācu kolonijas atradās. Pēc Krievijas statistikas datiem Polijas karalistē dzīvoja 552~000 vāciešu. 116~tūkstošu mita pilsētās, 417~tūkstošu~--- uz laukiem, pārējie~--- miestos. Taču 308~tūkstoši, tātad vairāk nekā puse no visiem, dzīvoja 101~kopienā, kur tie sastādīja vairāk nekā 20\% no iedzīvotājiem, tāpēc varēja saglabāt savu nacionālo identitāti. Tas radīja neapmierinātību poļos.

Nebija gludas arī \strong{poļu un ebreju attiecības}. R.~Dmovskis, pirmkārt balstoties uz to, ka ebreji runāja vācu valodai tuvajā jidiš valodā, uzskatīja ebrejus par vācu ``\detxti{Drang nach Osten}'' avangardu, kas gan bija pilnīgā pretrunā ar vēstures attīstības gaitu, jo neievēroja, ka tieši vāciešu izceļošana no savas zemes sagatavoja ceļu ebreju, t.s. aškenazi, izceļošanai no Vācijas XV un XVI gadsimtā. Taču R.~Dmovska uzskats, ka arī Polijas karalistē ebreji ir ``ģermanizējošs elements'', guva poļos plašu izplatību. Tiesa, ebreji parasti prata vācu valodu, jo tā bija vajadzīga viņu darījumos. Bet tāpat viņi prata arī poļu un krievu valodu. Tā kā ebreji tiecās pēc izglītības, viņu bērni bieži gāja vācu skolās, vācu mācītāja un ebreju rabīna dēls varēja sēdēt blakus vienā solā un viņus mācīt varēja ne tikai vācu, bet arī ebreju skolotājs. Poļu skolās labas attiecības starp poļu katoļu un ebreju bērniem nevarēja pastāvēt, tāpēc arī ebreji tās neapmeklēja.

Ebreju ekonomists un filozofs J.~Bukšpans 1913.~gadā uzsvēra, ka poļu sabiedrība tradicionāli ir piesātināta ar antisemītismu, īpaši sadzīves līmenī, un antisemītisms kļūst par poļu kultūras sastāvdaļu.

\strong{Antisemītisma rašanos} poļu vidū vācu vēsturnieks E.~Meijers saista ar Jasna Goras (\pltxti{Jasna Gora}) klostera Čenstohovā (\pltxti{Częstochowa}) aizstāvēšanas panākumu pret zviedru daudzkārtēju pārspēku 1656.~gadā. Tas izsauca entuziasma uzplūdus Polijā un zviedru padzīšanu, ko Polijā daudzi vērtēja kā klosterī glabājamās Jaunavas Marijas svētbildes (\pltxti{Częstochowska Ikona Matki Bożej (Bogurodzicy})) paveiktu brīnumu. Karalis Jans II Kazimirs pasludināja Jaunavu Mariju par Polijas karalieni~--- aizbildni (\pltxti{Polską Królową}) un uzvaras gadījumā pār zviedriem apsolīja padzīt no valsts visus ebrejus. No tā laika daudziem poļiem poliskums, katolicisms un antisemītisms cieši savijās. XIX gadsimtā Polijas karalistes intelektuālā un politiskā elite pret ebrejiem izturējās ar nevērību, ebrejus uztverot gan kā konkurentus ekonomiskajā jomā, taču ne kā aktīvus politiskā procesa dalībniekus. Tāpēc tad ``ebreju jautājums'' bija atbīdīts otrajā plānā uz poļu cīņas pret rusifikāciju fona. Taču XX gadsimta sākumā, kad ebreji sastādīja Varšavas pilsētas vēlētāju vairākumu un atbalstīja sociālistus pret poļu nacionālajām partijām, situācija radikāli mainījās. Tas izsauca poļu nacionālistu un pat liberāļu sašutuma vētru. Nacionāldemokrāti un citu pilsonisko partiju pārstāvji aicināja boikotēt ebreju veikalus un citas iestādes, poļu uzņēmējus aicināja pārtraukt sakarus ar ebreju kapitālu, atbrīvot no darba visus ebrejus, sākas uzbrukumi ebreju jauniešiem, pat māksliniekiem. Tiesa, daļa nacionāldemokrātu vadītāju sāka saprast šādas plašas antisemītiskas kampaņas iespējamās postošas sekas un aicināja to ierobežot vienīgi ar ekonomisko sfēru. Antisemītisms bija sava veida indikators, ar kuru varēja noteikt politisko institūtu un iedzīvotāju politiskās kultūras briedumu.

J.~Bukšpans atzīmēja antisemītisma specifiku Polijā. Pirmkārt, ebreji Polijas karalistē nebija paklausīgs nenozīmīgs iedzīvotāju mazākums, viņi veidoja līdz 15\% no reģiona iemītniekiem. Otrkārt, ebreji un poļi bija saistīti daudzām ekonomiskām, kultūras un ikdienas saiknēm. Treškārt, poļiem nebija savas valsts, tāpēc poļu antisemīti bija spiesti vienmēr atskatīties uz Pēterburgas nostāju. Pie tam ievērojama Polijas ebreju daļa ekonomiski bija maz atkarīga no poļiem, jo orientējās uz tirdzniecības attīstību ar Krieviju. Tāpēc viņu labklājība bija atkarīga nevis no iekšējiem Polijas notikumiem, bet gan ražas vai neražas Krievijā, rubļa stabilitātes un Krievijas impērijas politikas kopumā. Savas relatīvās neatkarības no poļu lietām sajūta deva ebrejiem arī zināmu brīvību dialogā ar poļu eliti. Ebreju jaunatne diezgan ātri apguva krievu valodu, un daļa pie tās slikti vai nemaz nezināja poļu valodu. Polijas ebreji viegli dibināja kontaktus ar tautiešiem citās Krievijas impērijas daļās, pārņemot to ``krieviskumu''. Tas parādījās arī ikdienas dzīvē un politiskajā orientācijā.

Pēc J.~Bukšpana domām Polijas ebreju kopiena izrādījās visai mobila, atšķirībā no poļiem diezgan ātri integrējās industriālajā sabiedrībā, gūstot no tā arī labumu materiālajā, sociālajā un kultūras jomā. Ebreji atmeta agrākos nacionālos un vēsturiskos aizvainojumus, piemērojoties XX gadsimta realitātei. Turpretī poļi demonstrēja savu tradicionālismu, ar grūtībām pieņemot XX gadsimta izaicinājumus. Īpaši tas attiecās uz konservatīvo poļu zemniecību. Ne velti poļu masu apziņā ebrejiskā pilsēta, slīgstoša netikumos, tika pretstatīta it kā tikumīgajiem poļu laukiem.

No otras puses, ebreju pamatmasa nepieņēma polonizāciju sava zemā izglītības un kultūras līmeņa dēļ, kuram polonizātori nepievērsa vajadzīgo uzmanību. Varšavas un citu pilsētu ebreju zemākie slāņi nebija galīgi sarāvuši saiknes ar mazo ebreju miestiņu pasaules uztveri. Polonizātoru aicinājumi asimilēties poļu sabiedrībā tos maz skāra. Tikai neievērojama ebreju inteliģences daļa pilnībā polonizējās, papildinot poļu zinātnes un kultūras pārstāvju rindas XIX--XX gadsimtu mijā.

Tikai ``reālās politikas partija'' un sociālisti nosodīja antisemītismu Polijā. ``Reālisti'' jebkurā vardarbībā un arī ebreju boikotā saskatīja masu terora priekšvēstnešus un pašu poļu sabiedrības degradācijas sākumu. Sociālisti stāvēja internacionālisma pozīcijās.

\strong{Saimnieciskās attīstības jomā} arī pēc 1905.~gada revolūcijas, tāpat kā iepriekš, Krievijas Finanšu ministrija centās dažādi atbalstīt Polijas tautsaimniecības pakļaušanu visas Krievijas interesēm.

Jau XIX--XX gs. mijā visai svarīga bija problēma par Vislas novada jeb Polijas karalistes lomu visas Krievijas impērijas budžeta formēšanā. Par diskusijas izejas punktu kļuva ekonomista N.~Jasnopoļska XIX gs. 90.~gadu sākumā publicētais sacerējums ``\rutxti{О географическом распределении государственных доходов и расходов в России}'' (``Par Krievijas valsts ienākumu un izdevumu ģeogrāfisko sadali''), kurš bija balstīts uz Valsts kontroles datiem par 1860--1880.~gadu. Pēc ziņojuma, poļu guberņas impērijā ieņēma trešo vietu ienesīguma ziņā, dodot ap 8\% no visiem valsts ienākumiem. Tā mazauglīgās Polijas karalistes zemes bija visai svarīgas impērijas budžetam. Salīdzinājumā ar Sibīriju, Aizkaukāzu un Turkestānu, kuras deva tikai 1/20~daļu valsts ienākumu, Polija no ekonomiskā viedokļa izskatījās par visai ``izdevīgu'' perifēriju. Taču pēc izdevumiem (6,6\% no visiem) Polija arī atpalika tikai no galvaspilsētu guberņām. Neraugoties uz to, ka Polijas dotie ienākumi pārsniedza izdevumus (1868.--1890.~gadā ik gadus par 6.~miljoniem rbļ.) Krievijas ``centra'' pārstāvji bija sašutuši par lielajiem izdevumiem. 1899.~gadā Krievijas Senāta Finanšu komitejas loceklis P.~Saburovs publicēja ``Materiālus par krievu finanšu vēsturi 1866--1897.~gadā'' (''\rutxti{Материалы для истории русских финансов 1866--1897 гг.}''). Aprēķinot ienākumus, viņš no to summas izslēdza ienākumus no dzelzceļiem, no vīna [un degvīna] pārdošanas un muitas. Pēc šādiem aprēķiniem iznāca, ka uz vienu pašas Krievijas Eiropas daļas guberņās dzīvojošo iznāca 1,3~reizes vairāk tiešo maksājumu nekā Polijā un nomaļu iedzīvotāji, tai skaitā Polijā un Baltijā, ik gadu it kā ``iedzīvojās'' par 12 līdz 22~rbļ. uz vienu vīriešu kārtas cilvēku. To izskaidroja ar valsts lielajiem izdevumiem armijas uzturēšanai pierobežas apgabalos. Pēc šīs loģikas militārās administrācijas tērētās summas tiešā ceļā ieplūda pierobežas iedzīvotāju kabatās. Iekškrievijas guberņās izdevumi tur izvietoto armijas daļu uzturēšanai ievērojami atpalika no Polijas karalistes rādītāja. Tātad P.~Saburova materiāli netieši apvainoja Polijas guberņu iedzīvotājus, ka viņi ir slogs Krievijas citu guberņu iedzīvotājiem. Tāda tendence valsts ierēdņu nostājā bija vērojama arī turpmāk. (Mūsdienu krievu vēsturnieks A.~Širokorads gan norādījis, ka Vislas novada guberņu ekonomiku ievērojami palīdzēja attīstīt Krievijas ģenerāļu korupcija. Piemēram, par smago lielgabalu lafešu, šaursliežu dzelzceļa iekārtu un cita aprīkojuma cietokšņiem pagatavošanu Urālu rūpnīcām tika maksāts vismaz divreiz mazāk nekā Varšavas un citu Polijas pilsētu rūpnīcām. Taču pasūtījumus saņēma šie Vislas novada uzņēmumi, kuru īpašnieki vairākumā bija vācieši un ebreji.) Viedoklis, ka Polija ir vainojama valsts centra ekonomiskā vājumā, Krievijā kļuva visai izplatīts.

Poļu izcelsmes autori: jurists un publicists V.~Spasovičs un žurnālists, politiskais darbinieks E.~Piļcs apstrīdēja šo uzskatu, rādot, ka liela daļa no izdevumiem poļu guberņās aiziet tiešajiem militārajiem mērķiem un tikai maza daļa no tiem tiek tērēta vietējo iedzīvotāju darba atlīdzināšanai. Pēc viņu ziņām Polijas iedzīvotāji salīdzinājumā ar Krievijas centrālajām guberņām nesa smagāku tiešo nodokļu nastu. Centrālajās guberņās viens cilvēks gadā maksāja 6,32~rbļ. nodokļu, bet Polijā~--- 14,26~rbļ. Valsts lielākie izdevumi poļu guberņās rādīja nevis to priviliģēto, bet pretēji~--- apspiesto stāvokli, jo kalpoja svešas armijas uzturēšanai.

Mūsdienu Krievijas sociālās vēstures pētnieks B.~Mironovs ir devis valsts ienākumu (visu tiešo un netiešo ienākumu, izpirkšanas maksājumu, muitas u.c.) un izdevumu (pārvaldei, armijas uzturēšanai u.c.) ainu XIX gadsimta beigu dažos gados. Tajā var redzēt arī 10~Vislas novada un 30~lielkrievu apdzīvoto guberņu (Bez Pēterburgas kā galvaspilsētas guberņas) stāvokli. Skat. \ref{tab:table3}. tabulu.

\noindent
\begin{table}[h!]
\caption{Vidējie valsts ienākumi un izdevumi 1892.--1895.~gadā Vislas novada guberņās salīdzinājumā ar 30 lielkrievu guberņām} \label{tab:table3}
\begin{tabularx}{\linewidth}{|p{1.5cm}|p{1.5cm}|p{1.5cm}|p{1.5cm}|p{1.5cm}|p{1.5cm}|p{1.5cm}|}
\hline
\strong{Guberņu grupa} & \strong{Iedzīvotāju skaits (tūkstošos)} & \strong{Ienākumi (tūkstošos rbļ.)} & \strong{Izdevumi (tūkstošos rbļ.)} & \strong{Ienākumi uz 1 iedzīvotāju (rubļos)} & \strong{Izdevumi uz 1 iedzīvotāju (rubļos)} & \strong{Starpība starp 1 iedzīvotāja ienākumiem un izdevumiem (rubļos)} \\
\hline
10 Vislas novada guberņas & 9~443 & 99~425 & 48~069 & 10,53 & 5,09 & 5,44 \\
\hline
30 lielkrievu apdzīvotās guberņas & 52~578 & 388~706 & 195~306 & 7,39 & 3,71 & 3,68 \\
\hline
\end{tabularx}
\end{table}

% page 214


Tātad valsts ienākumi minētajos gados uz 1 iedzīvotāju Vislas novada guberņās pārsniedza ienākumus no 1 iedzīvotāja 30 lielkrievu guberņās par 3,14~rbļ., bet starpība starp ienākumiem un izdevumiem valsts labā no 1~Vislas novada guberņu iedzīvotāja sastādīja 5,44~rbļ. salīdzinājumā ar 3,68~rbļ. starpību uz 1~iedzīvotāju no 30~lielkrievu apdzīvotajām guberņām. Jāsaka, ka B.~Mironovs devis arī datus par tiešajiem nodokļiem no 1~iedzīvotāja 1886.--1895.~gadā. Vislas novada guberņās tas sastādīja 1,35~rbļ., 30~lielkrievu guberņās~--- 1,72~rbļ. Tātad, kaut periodi nesakrīt, var pieņemt, ka tiešos nodokļus lielkrievu guberņu iedzīvotāji arī tabulā rādītajos 1892.--1895.~gados maksāja lielākus nekā Vislas novada guberņu iedzīvotāji. Taču šī starpība (nosacīti 0,37~rbļ.) bija krietni mazāka nekā kopējā starpība starp valsts ienākumiem un izdevumiem uz 1~iedzīvotāju, kurā taču tiešie nodokļi jau bija ieskaitīti. Pats B.~Mironovs uzsver to, ka valsts ienākumi uz 1~iedzīvotāju lielkrievu guberņās bija lielāki nekā nekrievu guberņās, taču tas, kā redzējām, neattiecās uz Vislas novada guberņām. To rāda viņa paša minētie skaitļi.

Kad sāka darboties Valsts dome, diskusijas tika pārceltas uz to. Poļu deputāti prasīja izlīdzināt nodokļu slogu dažādās guberņas. Labējie un krievu nacionālisti, atkārtojot tēzi par Polijas priviliģēto stāvokli un krievu tautas ekspluatāciju, kategoriski iebilda pret nodokļu nolīdzināšanu. ``Nabagā'' un ``nekulturālā'' Polija esot ``apbērta'' ar krievu zeltu. Šīs debates bija tieši saistītas ar krievu un poļu rūpnieku konkurences cīņām. Savstarpējo cīņu noteica krievu rūpnieku pretenzijas uz kundzību iekšējā tirgū un Polijas baņķieru un rūpnieku aktivitāte, kuru viņi izplatīja ne tikai uz valsts nomalēm, bet arī centru. Pretrunas izsauca arī tas, ka Polijas karalistes tautsaimniecība, kura orientējās uz Krievijas tirgiem, auga uz vācu kapitālu rēķina un orientējās uz Eiropas tehnoloģiju un darba organizāciju. Uz 1914.~gada sākumu Vislas novada guberņās dzīvoja nepilni 7\% visas Krievijas impērijas iedzīvotāju, bet ražots tika ap 12\% visas rūpnieciskās produkcijas. 1914.~gadā Polijas karalistē saražotās rūpnieciskās produkcijas vērtība sasniedza 1~200 miljonus rubļu.

Tā, kaut Polija veiksmīgi iegāja Krievijas tirgos, tas radīja jaunas problēmas un pretrunas. Reizē jāsecina, ka Krievijas valdības realizētie vienveidības un administratīvās centralizācijas principi izrādījās mazefektīvi, jo dažādu reģionu ekonomiskās sistēmas noteica arī to dziļās sociālās un kultūras atšķirības.

Saimniecības attīstība Polijas karalistē veicināja iedzīvotāju pieaugumu. Iedzīvotāju skaits auga no 9,5~miljoniem 1897.~gadā līdz 13~miljoniem 1913.~gadā. Tas saasināja sociālās problēmas. Sekas bija izceļošana. Pavisam līdz 1914.~gadam 2,6~miljoni poļu no karalistes pameta savu dzimteni, lielākoties dodoties uz ASV.

1907.~gadā cara valdības atteikšanās no iepriekšējā pārstāvniecības modeļa un tā nomaiņa ar 3.~jūnija sistēmu nozīmēja, ka varai uz laiku ir izdevies radīt ilūziju, ka impērija ir atkal nostiprinājusies. Taču faktiski turpinājās stagnācija, īpaši etnoreģionālajā aspektā. Nacionālo nomaļu pārstāvju skaita strauja samazināšana III Domē nozīmēja, ka varas iestādes atsakās no to integrācijas valsts politiskajā sistēmā. Tāpēc, ka politiskās pārmaiņas tika novilcinātas gadu desmitiem un zināmas nepietiekamas izmaiņas notika tikai revolūcijas rezultātā, politiskā sistēma pārdzīvoja lielu pārslodzi. Dome nedarbojās kā tilts starp monarhiju, valdību, kura bija noskaņota pārāk konservatīvi, un sabiedrību, kura bija noskaņota pārāk radikāli. Domes darbība vienīgi atspoguļoja bezdibeni starp tām.

Deputātu skaits \strong{III Domē} bija samazināts par 82~vietām (uz 442). 3.~jūnija cara Manifestā par II~Domes atlaišanu bija paziņots, ka ``Domei pēc gara ir jābūt krieviskai. Citām tautām, kuras ietilpst mūsu valstī , jābūt savu vajadzību pārstāvjiem, taču ne tādā skaitā, kas dotu tiem iespēju lemt tīri krieviskus jautājumus.'' Poļu guberņu deputātu skaits tika samazināts no 37 uz 12, pie tam divi mandāti no tiem tika piešķirti pareizticīgajiem. Ja Krievijas Eiropas daļas iedzīvotāji vidēji ievēlēja 1~deputātu no 284,2~tūkstošiem cilvēku, tad Polijas karalistē pārstāvniecības norma bija 1 Domes deputāts uz 671,6~tūkstošiem iedzīvotāju. Vēlēšanu likums nošķīra daudz iedzīvotāju no vēlēšanām, sevišķi ebrejus un poļus no mazām pilsētiņām. Pavisam III Domē tika ievēlēti 22~poļi, kuri darbojās poļu \pltxti{kolo} un poļu-lietuviešu-baltkrievu grupā.

III Dome, tāpat kā iepriekšējās, kļuva par impērijas varas un nekrievu tautu sadursmju arēnu. Krievu lielburžuāzijas pārstāvju (oktobristu un progresistu) pozīciju formulēja grāfs A.~Uvarovs no Saratovas guberņas: ``Mēs, oktobristi, ar lielu prieku dosim Polijai visu to, ko paši saņemsim centrā, mēs ar prieku dosim Polijai zemstu pašpārvaldi, plašu pilsētu pārvaldi, taču tikai, kungi, ar to nosacījumu, ka nomales neprasīs to, kā nav Krievijas centrālajai daļai.'' Pret šādu nostāju uzreiz, apspiežot Domes atbildi uz imperatora vēstījumu tai un valdības deklarāciju, pēc būtības uzstājās R.~Dmovskis, kurš ierosināja pēc vēlējuma ``celt vispārējo labklājību'' pievienot vēlējumu ``apmierināt to tautību taisnīgos centienus, kuras ieiet valsts sastāvā''. 6~miljonu ``nabago un pilnībā beztiesīgo'' ebreju vārdā labojumu atbalstīja arī deputāts no Kurzemes Lazars~Nisselovičs. Neraugoties uz kadetu un kreiso atbalstu, R.~Dmovska labojums tika noraidīts. Atbildot uz to, poļu \pltxti{kolo} atteicās balsot par atbildi kopumā. R.~Dmovskis paziņoja, ka ``poļu tauta nekad nesamierināsies ar otrās šķiras pilsoņu stāvokli valstī un nevarēs samierināties ar valsti, kurā viņai piešķirts tāds stāvoklis''. Uz to tūlīt atskanēja premjera P.~Stolipina skarba atbilde, ka ``tam pašvaldības spēkam, uz kuru balstīsies valdība, vienmēr ir jābūt nacionālam'', piebilstot, ka Polijā ir ne tikai maz skolu, bet nav arī augstākās mācību iestādes, ``tāpēc, ka tie pilsoņi, kuri sevi tikko nosauca par otrās šķiras pilsoņiem'', negrib augstskolā lietot visas valsts valodu~--- krievu valodu. P.~Stolipins ieteica poļiem vispirms iemācīties cienīt Krievijas pilsoņa nosaukumu. Viņš piedāvāja: ``Lūk, saliedējaties vispārnacionālā cementā un tad, kungi, prasiet no mums decentralizāciju''. Kā redzam, nacionālisti, esot pie varas, visās valstīs un laikos rīkojas vienādi: vispirms prasa citiem atteikties no savas tautības un tikai tad sola tiem kādas tiesības. Viens no liberālās krievu pilsonības [kadetu] līderiem F.~Rodičevs, atbalstot nacionālo frakciju uzstāšanās, atzīmēja, ka ``katram ir jābūt tiesībām uzskatīt Krieviju par savu tēvzemi''. Viņš vērsa uzmanību uz to, ka pasaulslavenā zinātniece, radioaktivitātes atklājēja M.~Sklodovska-Kirī nevar savā valodā lasīt lekcijas Krievijā, bet varas iestādes uz to saka: ``lai mūs mīl par to''. Liberālais ekonomists P.~Strūve, norādot, ka Krievijas varenības pieaugums ir iespējams Tuvo Austrumu un Klusā okeāna virzienā, uzstāja, ka tā sasniegšanai ir nepieciešams atzīt divu ``nomaļu''~--- Polijas un Somijas asimilācijas neiespējamību. Viņš uzskatīja, ka, balstoties uz Polijas ekonomisko saistību ar Krieviju, ir jāpanāk, lai poļu iedzīvotāji ``būtu apmierināti ar savu likteni''.

Nacionāldemokrāti bija vienīgā poļu politiskā partija, kura konsekventi vēlējās apvienot visus poļu slāņus. Pēc II Domes atlaišanas un vēlēšanu likuma izmaiņas R.~Dmovskis neatbalstīja protestu pret 3.jūnija sistēmu Krievijā, kaut partijas iekšienē bija arī grupas, kuras vēlējās aktīvāku cīņu par neatkarības atjaunošanu. Nacionāldemokrātiem vienīgajiem bija arī izstrādāta programma Polijas kā nacionālas valsts atjaunošanai. Tai bija jābalstās uz draudzību ar Krieviju, kas pasargātu to no galvenajām prūšu-vācu briesmām. 1908.~gadā R.~Dmovskis publicēja grāmatu ``\pltxti{Niemzy, Rosja i kwestia Polska}'' (Vācija, Krievija un poļu jautājums), kurā izklāstīja savus ģeopolitiskos uzskatus. Pār galveno poļu ienaidnieku viņš uzskatīja Vāciju kā spēcīgāko asimilatorisko spēku, brīdināja no vācu ekspansionisma briesmām, kurām poļi varot pretoties, vienīgi saņemot Krievijas atbalstu. Iespējamajā karā Eiropā poļiem bija jānostājas Krievijas pusē pret Vāciju un jācīnās par visu poļu zemju apvienošanu Pēterburgas varā. Ar savu pretvācisko virzību partijai izdevās uzrunāt plašus poļu tautas slāņus. R.~Dmovska sludinātais ``neoslāvisms'' ar piesliešanos Krievijai poļu tautā gan rada mazu atsaucību, tomēr tas ļāva šai pilsoniski-nacionālistiskajai un arī antisemītiski noskaņotajai partijai samērā sekmīgi darboties. Taču R.~Dmovska tēzi, ka Polijai bīstamāka ir Vācija, bija arvien grūtāk pierādīt, jo cara varas iestādes centās atņemt poļiem visus revolūcijas laikā panāktos iekarojumus. Arī 1908.~gada oktobrī sakarā ar Austroungārijas veikto Bosnijas un Hercogovinas aneksiju izcēlusies starptautiskā krīze deva triecienu R.~Dmovska koncepcijai. Tā no vienas puses parādīja Krievijas vājumu, no otras puses~--- relatīvo Austroungārijas un Vācijas bloka vienotību. Krievijas patvaldība bija spiesta atgriezties pie krievu-vācu attiecību uzlabošanas kursa, kurš zināmā mērā arī noteica Krievijas iekšējo ``poļu'' politiku. Šāda perspektīva būtiski apdraudēja krievu un poļu tuvināšanos uz pretvāciska pamata.

Oficiālo Pēterburgas pozīciju Polijas jautājumā ietekmēja arī Vācija, Krievijas un Vācijas imperatoru radniecīgās attiecības. Ķeizars Vilhelms II brīdināja Nikolaju II, ka autonomijas piešķiršana Polijai būs nedraudzīgs akts pret viņa valsti. Krievijas ārlietu ministrs S.~Sazonovs vēlāk savās atmiņās rakstīja, ka Krievijas politiku ievērojamā mērā noteica Berlīnes ietekme, kura izpaudās kā draudzīgi padomi un brīdinājumi katru reizi, kad Vācijas valdība atklāja Pēterburgā vismazāko nodomu izlīgt ar Poliju.

Krievijas valdošie slāņi neuzticējās R.~Dmovskim, vēl jo vairāk tāpēc, ka viņš pats draugu pulkā runāja par savas pozīcijas taktisko, pagaidu raksturu. Par kādu savu kritiķi R.~Dmovskis teica: ``Viņš grib karot ar krieviem jau šodien, bet es dodu priekšroku pagaidīt līdz rītam''. Starp citu, krievu politiķu un arī sabiedrību aizvainoja R.~Dmovska demonstrētā šļahtičiem raksturīgā augstprātība. Tā, viņš izteicās, ka patiesu vienošanos starp krievu un poļu tautām traucē tas, ka ``mēs poļi esam eiropieši, bet krievi~--- aziāti''. Pēc laika biedru liecības P.~Stolipins šai sakarā piezīmēja: ``Ja tas tiešam ir tā, es Dmovska kungam drīz parādīšu, kur šī Āzija sākas!'' Krievijas valdība ne bez pamata baidījās, ka autonomijas piešķiršana Polijai kļūs par pirmo soli tās ceļā uz atdalīšanos no Krievijas.

R.~Dmovski savukārt biedēja Pēterburgas un Berlīnes savienības iespēja, tāpēc viņš visādi centās atvieglot Francijas--Anglijas un Krievijas tuvināšanos. Kad Krievijas valdība neizrādīja vēlmi veidot dialogu poļu jautājumā, 1909.~gada janvārī R.~Dmovskis, vīlies Krievijas Valsts Domes darbībā, atteicās ne tikai no frakcijas vadītāja amata Domē, bet arī vispār nolika savu Domes deputāta mandātu, taču palika politiski aktīvs. Piemēram, 1911.~gadā viņš guva bēdīgu slavu, organizējot ebreju uzņēmumu boikotu.

XX gadsimta sākumā Polijas karalistē valdīja sarežģīta situācija izglītības laukā. Pēc analfabētu skaitā karaliste bija viens no atpalikušākajiem Eiropas reģioniem, sāka pat atpalikt no Krievijas iekšzemes guberņām, kur, pateicoties zemstu darbībai un arī valdības iestādēm, strauji auga skolnieku skaits. Kaut valdība deva atļauju dibināt privātskolas ar poļu mācību valodu, to beidzēji nevarēja iestāties valsts augstākajās mācību iestādēs, jo tiem nebija oficiālu valsts izsniegtu gatavības apliecību. Atbilstoši 1908.~gadā pieņemtajam nolikumam poļu izcelsmes personas nedrīkstēja mācīt vēsturi un ģeogrāfiju.

1910.~gada oktobrī Domē tika skatīts jautājums par vispārējās sākumizglītības ieviešanu Krievijā. Likumprojekts paredzēja, ka mācībām jānotiek krievu valodā, taču pirmos divus mācību gadus var atļaut mācības dzimtajā valodā, ar to nosacījumu, ka jau pirmajos trijos mēnešos sāksies arī krievu valodas stundas. Izraisījās asas debates. Nacionālās frakcijas dzimtās valodas statusu uzskatīja par pilnīgi neapmierinošu. Rezultātā likums tā arī netika pieņemts.

Atbilstoši likumam par izglītību, ko cars parakstīja 1912.~gadā, par skolotājiem varēja strādāt tikai pareizticīgie, kad Polija karalistē šai laikā 85\% skolotāju bija katoļi.

Valdība centās mobilizēt krievu nacionālismu pret nacionālo nomaļu iedzīvotājiem. Atsevišķos gadījumos izvērsās veseli ksenofobijas uzplūdi. Viens no tādiem gadījumiem bija sakarā ar it kā notikušo rituālo slepkavību 1911.--1913.~gadā izvērstā ``Beilisa lieta'', taču tā notika ārpus Polijas Kijevas guberņā. Poliju skāra otrs tikpat spilgts gadījums, kad no t.s. Vislas novada 1912.~gadā mākslīgi tika izdalīta Holmas (poļu \pltxti{Chełm}, ukraiņu \uktxti{Холм}) guberņa un pārcelta Krievijas Rietumu novada sastāvā, bet Vislas novadā palika 9~guberņas. Jaunā Holmas guberņa sastāvēja no Sedlecas guberņas (\rutxti{Седлецкая губерния}, poļu \pltxti{Gubernia siedlecka}) austrumu daļas un Vislas novada agrākās Ļubļinas guberņas. Saknes šim projektam bija meklējamas vēl XIX~gadsimta vidū un bija saistīts ar uniātu baznīcas likvidēšanu šai teritorijā. Projekts par Holmas guberņas radīšanu Krievijas Iekšlietu ministrijā tika apspriests jau no 1836.~gada un vienmēr ministrija to noraidīja, uzskatot, ka ``strauji pārkrievošanas un gadsimtiem pastāvējušo Aizbugas iedzīvotāju sakaru ar Polijas karalistes civiliekārtu laušanas līdzekļi izsauks visādā ziņā nevēlamu poļu iedzīvotāju prātu sajukumu'' (''\pltxti{крутые меры обрусения и ломки вековых связей Забужья с гражданским укладом Царства Польского вызовут во всех отношениях нежелательное смятение умов польского населения}''.) Gadsimta pēdējā ceturksnī pieauga tā saukto ``stūrgalvīgo'' (\rutxti{упорствующие}) skaits, t.i.~--- to bijušo uniātu skaits, kuri nevis pārgāja pareizticībā, bet pieslējās katoļu ticībai. Birokrātijā spēcīga bija tendence kā vainīgos uzskatīt ``kūdītājus'', šai gadījumā katoļu garīdzniekus. 1902.~gadā tika izveidota pat ``Sevišķā padome'' Holmas guberņas izveides jautājuma izskatīšanai, kura ieteica to neveidot. Imperators toreiz piekrita šim ieteikumam.

Arī 1905.~gada 17.~aprīļa dekrēts ``Par tolerances nostiprināšanu ticības lietās'' (''\pltxti{Об укреплении начал веротерпимости}'') nebija radikāls pavērsiens varas iestāžu reliģiskajā nostājā. Tas neatcēla pareizticības priviliģēto stāvokli. Galvenais šī dekrēta saturs bija atļauja personām, vecākām par 14~gadiem, pēc savas izvēles mainīt konfesionālo piederību, atļaujot arī atteikties no pareizticības. Tiesa, ticības maiņas procedūra bija visai sarežģīta un varēja vilkties ilgi. Jau dekrēta sagatavošanas gaitā tika prognozēts, ka tā pieņemšanas rezultātā ap 100 līdz 150~tūkstošu cilvēku no pareizticības pāries katoļticībā, taču revolūcijas apstākļos dekrēts tomēr tika izsludināts. Prognoze attaisnojās un agrākajos uniātu apgabalos pareizticīgo baznīca zaudēja ap 1/3~savu ticīgo, nācās slēgt veselas draudzes. Ar varu uzspiestā pareizticība bijušajiem uniātiem nelikās pievilcīga. Pie tam daudzu zemnieku acīs katolicisms bija ``kungu'' (panu) reliģija un pāreja tajā saistījās ar sociālā prestiža pieaugumu. Citādi to tulkoja poļu nacionālisti. Pāreju katolicismā viņi saistīja ar poliskās nacionālās apziņas pieaugumu. Poļu muižnieki aktīvi ietekmēja no viņiem atkarīgos zemniekus pāriet katoļticībā. Daudzviet algādžiem paziņoja, ka tikai katoļi varēs saglabāt darbu, pārējie tiks atlaisti. Kā vēlāk rakstīja jau minētais bīskaps Jevlogijs: ``Bija draudi, bija solījumi: grāfiene Zamoiska (\pltxti{Zamoyski}~--- poļu dižciltīgo dzimta) solīja govi katrai ģimenei, kura pāries katoļticībā''.

Pēc 1905.~gada polemika par Vislas novada dienviddaļas ap 300~000 lielās iedzīvotāju grupas ``īsto'' nacionālo piederību ieguva asu raksturu. Katra puse operēja ar sev izdevīgiem statistiskajiem datiem. Atzīmējams, ka gan krievi, gan poļi pilnībā ignorēja Sedļecas un Ļubļinas guberņās dzīvojošos daudzos ebrejus.

Savdabīgais katolicisma revanšs nemazināja krievu nacionālistu centienus konsolidēt krievu nacionālo teritoriju. Protams, atkal skanēja atsauces uz katoļu garīdznieku ``vainu'', kuri ``atvilinot'' ticīgos no pareizticīgo draudzēm. Holmas pareizticīgo bīskaps Jevlogijs ļoti enerģiski uzstājās par īpašas Holmas guberņas izveidi, kuras iedzīvotāju vairākums būtu pareizticīgie, tātad krievi. 1909.~gadā III Valsts dome ķērās pie šīs problēmas apspriešanas. 1912.~gadā ziņojumā Domē tika atzīts, ka katoļu šai teritorijā ir vairāk. Taču tika atzīts, ka daļa no pareizticības atkritušo, lai ar grūtībām, ne vienmēr apzināti, bet tomēr saglabājot savu dzimto valodu un ierašas. Kopā ar pareizticīgajiem iedzīvotājiem viņiem esot nepieciešama savas atsevišķas guberņas izveide, kurā tie varētu droši aizstāvēt savas interese. Iekšlietu ministrija 1912.~gada 12.~maijā likumprojektu iesniedza Valsts domei. Pēc tā Holmas guberņu bija jāpievieno Krievijas Dienvidrietumu novadam, bet Sedļecas guberņa jālikvidē. Poļu deputāti uzsvēra, ka jaunās guberņas izveide ne tikai mainīs likumdošanas ceļā noteiktās Polijas karalistes robežas, bet nozīmēs arī to, ka tajā tiks ieviesti visi noteikumi, kuri vērsti pret poļiem un katoļu baznīcu Krievijas rietumu guberņas. Kad Domē norisa debates par Holmas guberņas izveidi, Lembergā (Ļvovā) pie Ā.~Mickēviča pieminekļa notika poļu demonstrācija, kuras dalībnieki dziedāja poļu patriotiskas dziesmas un sadedzināja Nikolaja II portretu. Pret jaunas guberņas izveidi Domē uzstājās kadeti un kreisās frakcijas. Sociāldemokrātiskā frakcija, uzstājoties Domē pret Holmas atdalīšanu, to raksturoja kā ``ceturto Polijas dalīšanu'', kam nevar piekrist, jo tā bija tikai Polijas karalistes dalīšana. Boļševiku avīze ``\rutxti{Правда}'' (``Taisnība'') paziņoja, ka krievu strādniekam nav nekā kopīga ar carisma rusifikātorisko politiku, ka poļu proletariāts krievu proletariātā atradīs uzticamu sabiedroto pret nacionālo apspiestību. Pret īpašas Holmas guberņas izveidi nostājās arī daži galēji labējie krievu politiķi. Tā N.~Markovs apgalvoja, ka likumprojekts ir kaitīgs, jo ``nostiprina melīgo mācību, it kā pastāvētu kāda reāla Polijas karaliste\ldots{}'' Projektu apsprieda 17~sēdēs un pēc asas cīņas Dome to tomēr pieņēma, bet Valsts padome apstiprināja. 1913.~gada septembrī oficiāli tika ``atklāta'' Holmas guberņa. Jaunajā guberņā pareizticīgie veidoja 40\%, katoļi 39\%, bet jūdu ticības piederīgie~--- 16\% iedzīvotāju.

Soļi pret polonizāciju tajā gandrīz pilnībā atkārtoja tos, kuri bija īstenoti rietumu guberņās XIX~gs. 60.~gados. Īpaši vērīgi sekoja poļu valodas izskaušanai no sabiedriskās dzīves, par slepenu poļu valodas lietošanu skolās bija noteikti stingri sodi, pastā nepieņēma korespondenci ar adresi poļu valodā u.t.t. Saprātīgākie krievu tautības ierēdņi atzina, ka neprātīgā krievu nacionālisma kampaņa, atlaižot no valsts dienesta poļus un ebrejus, kuri bija valstij ne mazāk uzticami, kā krievi, nedod valstij kādu labumu, bet rada tikai neapmierinātību cietušajos.

Pareizticīgo Krievijas varas iestāžu cīņa pret polonizāciju, par katolicisma ietekmes samazināšanu poļu apdzīvotajos apgabalos atstāja ēnā attiecības starp katoļiem un protestantiem. 1914.~gadā kongresa Polijā dzīvoja arī ap 450~000 vācu un 35~000 poļu protestantu. Vācu iedzīvotāji lielpilsētās, īpaši Varšavā, bija pakļauti polonizācijai, uz tiem lielu iespaidu atstāja poļu šļahtiču kultūra. Piemēram, vācu izglītības iestāžu skaits Varšavā samazinājās. Piesliešanās poļiem deva sociālās izaugsmes iespējas. Process, neraugoties uz carisma cīņu pret polonismu, turpinājās līdz pat 1914.~gadam.

Paralēli Holmas lietai valdībā un likumdošanas iestādes skatīja arī jautājumu par \strong{zemstu ieviešanu} Krievijas rietumu guberņās: Lietuvā, Baltkrievijā un Labā krasta Ukrainā. Tās neietilpa Polijas karalistē, bet tajās dzīvoja arī poļi un kur, galvenais, bija poļu muižniecība. (Piemēram, pēc 1909.~gada oficiālajiem tautas skaitīšanas datiem Labā krasta Ukrainā dzīvoja 10~288~342 cilvēki. Poļu no tiem bija 646~562, taču zemes lielākā daļa piederēja poļu muižniekiem: Volīnijā~--- 47,9\%, Podolijā~--- 55,4\%, Kijevas guberņā~--- 44,1\%.) Zemstu ieviešana pielīdzinātu šo guberņu iedzīvotāju tiesības citu Krievijas guberņu iemītnieku tiesībām, reizē varētu nodrošināt krievu muižniecības un pareizticīgās garīdzniecības pārstāvjiem vairākumu, samazinot viņu mantas cenzu salīdzinājumā ar poļu muižniecību.

Jau 1906.~gadā tika strādāts pie likumprojekta par proporcionālas nacionālas pārstāvniecības ievešanu zemstēs, garantējot krieviem kā minimumu vairākumu amatu tajās. Debates Domē ievilkās četru ar pusi gadu garumā. Viens no varas iestāžu neapmierinātības iemesliem bija tas, ka rietumu guberņās parlamenta augšpalātā~--- Valsts padomē savus pārstāvjus vēlēja tikai [lielo] zemes īpašnieku kūrijai piederīgie un tāpēc šīs guberņas Padomē pārstāvēja vienīgi poļu muižnieki, kuri tur veidoja poļu \pltxti{kolo}. Bija domāts 6 rietumu guberņās (izņemot Viļņas, Kauņas un Grodņas guberņas, kur arī ar proporcionālas nacionālas pārstāvniecības ieviešanu nevarēja atņemt poļu muižniekiem vairākumu) ieviest poļu un krievu kūrijas, piešķirot otrajai vairākumu deputātu vietu. Apspriešanas gaitā izpaudās divas pretējas pozīcijas. Poļu \pltxti{kolo} uzskatīja par nepieņemamu poļu tiesību ierobežošanu, krievu nacionālisti šo ierobežošanu uzskatīja par nepietiekamu. P.~Stolipins atgādināja poļu deputātiem savu teikto vēl pirms II Domes atlaišanas, ka ``politikā nav atriebības, bet ir sekas'', uzsvēra, ka ``Rietumu novads būs krievisks uz visiem laikiem''. 1910.~gada 1.~jūnijā Dome projektu atbalstīja. Taču krievu labējo politiķi cīņā pret P.~Stolipinu izvirzīja jaunus iebildumus. Monarhists V.~Puriškēvičs apgalvoja, ka, ja tagad no Volīnijas, Podoļskas, Mogiļevas un Minskas guberņās uz Valsts domi tiek sūtīti cara kalpi, tad ar zemstu ieviešanu tajās parādīsies sociāldemokrāti, sociālrevolucionāri un separātisti. Valsts padomes locekļu vairākums ar S.~Vitti priekšgalā uzsvēra, ka nacionālās kūrijas nav savienojamas ar daudznacionālas impērijas uzbūves principiem. 1911.~gada martā Valsts padome noraidīja pantu par nacionālo kūriju ievešanu. P.~Stolipins griezās pēc atbalsta pie Nikolaja II, lai atlaistu gan Valsts domi, gan Valsts padomi un saskaņā ar Pamatlikumu 87.~pantu to prombūtnes laikā likumprojektu pieņemtu ar valdnieka dekrētu. Nikolajs II necieta tādu spiedienu uz sevi un tas vēl vairāk saasināja pretrunas starp caru un valdības vadītāju. Gala rezultātā 1911.~gada martā ar īpašu cara dekrētu (\rutxti{указ}) sešās guberņās (Vitebskas, Minskas, Mogiļevas, Kijevas, Volīnijas un Podolijas) zemstes gan tika ieviestas un vietējā poļu šļahta tika dziļi aizvainota, taču pēc šīs sadursmes jautājums par premjerministra nomaiņu faktiski bija jau izlemts. P.~Stolipina slepkavība 1911.~gada 1.~septembrī tikai ``atrisināja'' jautājumu automātiski.

Kā atzīmējusi mūsdienu krievu vēsturniece A.~Bahturina, Krievijas realizētājā ''poļu'' politikā par dominējošo pakāpeniski kļuva stingri noteiktas nacionāli-kulturālās robežas noteikšana starp Polijas karalisti un Krievijas impērijas ziemeļrietumu un dienvidrietumu guberņām. Galvenās briesmas krievu valdošo slāņu politiķi saskatīja nevis poļu kultūras attīstībā pašā Polijas karalistē, bet gan tās izplatībā baltkrievu, ukraiņu un lietuviešu apdzīvotajās teritorijās, kuras kādreiz ietilpa Žečpospolitas sastāvā. Krievijas valdība, pieļaujot dažas piekāpšanās poļu prasībām Polijas karalistē, rūpīgi raudzījās lai tās nepārkāptu tām nospraustās robežas.

Paralēli lielkrievu ietekmes stiprināšanai Polijas karalistē un Krievijas rietumu guberņās \strong{Krievijas valdība} bija norūpējusies par iespējamo ukraiņu separātisma pieaugumu un tāpēc centās vairāk vai mazāk tieši \strong{iejaukties Austroungārijas iekšējās dzīves jautājumos}, stiprinot tur prokrieviskās pozīcijas. Tiešām, Galīcijā rutēņiem draudēja polonizācijas briesmas. Ja pēc 1846.~gada tautas skaitīšanas datiem Galīcijā dzīvoja divarpus miljoni rutēņu un mazāk par diviem miljoniem poļu, tad 1910.~gadā iedzīvotāju skaitīšana konstatēja 4,8 miljonus poļu un nedaudz vairāk par 3 miljoniem rutēņu. Tiesa, tās bija ne tikai polonizācijas, bet arī tautas skaitīšanas īpatnību sekas. Tautas skaitīšanu 1846.~gadā veica Austrijas impērijas ierēdņi, pārsvarā vācieši, un par poļiem tika ierakstīti tikai tādi, kas sevi par tādiem nosauca. 1910.~gadā skaitīšanu veica jau ierēdņi~--- poļi, kuri par poļiem ierakstīja visus, kuri sevi neuzskatīja par nevienas citas tautības pārstāvjiem. Arī ebreji tika ieskaitīti kā poļi. Kaut Galīcijā ukraiņi sastādīja apmēram 1/2 visu iedzīvotāju, no 38~pastāvošajām ģimnāzijām tikai piecās mācības notika ukraiņu valodā.

Taču kad 1905.~gadā Krievijā tika atcelti 1876.~gadā ievestie ierobežojumi ukraiņu valodai, sākās straujš ukraiņu kultūras un izglītības uzplaukums. 1868.~gadā Galīcijas ``tautībnieku'' (\rutxti{народовцы}) nodibinātā biedrība ``\uktxti{Просвита}'' (``Izglītība'') tagad uzsāka savu darbību arī Krievijai piederošajā Ukrainā, plaši izvēršot grāmatu, žurnālu, avīžu izdošanu, dibinot ukraiņu valodas un vēstures mācību kursus. Tika izpirkts 100~000 eksemplāru metienā izdotais tautas kalendārs, 200~000 eksemplāru ukraiņu tautas dzejnieka T.~Ševčenko darbu kopojums ``\uktxti{Кобзар}'' (Tautas dziesminieks). Tāds ukraiņu kultūras uzplaukums Galīcijas ukraiņu (rutēņu) vidē radīja tieksmi vienoties ar Krievijā dzīvojošajiem tautiešiem. Galīcijas ukraiņos kļuva diezgan populārs vairs ne lozungs par autonomu Ukrainu Austroungārijas impērijas sastāvā, bet sauklis par Galīciju kopā ar pārējo Ukrainu Krievijas federācijas sastāvā. T.s. ``mazepieši'' (nosaukums atvasināts no hetmaņa I.~Mazepas, kurš cīnījās par Ukrainas nodalīšanos no Krievijas Ziemeļu kara laikā (1700--1721), nodeva Krieviju un pārgāja Zviedrijas pusē) sāka zaudēt cīņā par ietekmi ``bogdanoviešiem'' (nosaukums no hetmaņa B.~Hmeļņicka, kurš iestājās par Ukrainas apvienošanos ar Krieviju). Krievija veicināja ``bogdanoviešu'' darbību. Jau 1902.~gada 15.~decembrī Pēterburgā bija nodibināta Galīcijas krievu labdarības biedrība (\rutxti{Галицко-русское благотворительное общество}). Krievijas Iekšlietu ministrijas apstiprinātajos statūtos biedrība stādīja mērķi ``sniegt dažāda veida morālu un materiālu palīdzību krievu galīciešiem un to ģimenēm, kuras uz laiku vai pastāvīgi dzīvo Pēterburgā. Bez labdarības biedrības mērķis bija arī iepazīstināt krievu sabiedrību ar ``Piekarpatu Krievijas'' dzīvi, tās pagātni un tagadni. 1914.~gadā biedrībā darbojās ap 700 cilvēku.

Tomēr tai pašā laikā Krievijas valdošajās aprindās pastāvēja bailes no ukraiņu nacionālās kustības un Austroungārijas ietekmes tajā. 1908.~gadā Krievijas Ārlietu ministrija sagatavotajā ziņojumā bija norādīts, ka Austroungārijas varas iestāžu politikas rezultātā Galīcijā ``ukraiņi apvienosies vienā patstāvīgā tautā un tad cīņa ar separātismu kļūs neiespējama. Kamēr Galīcijā vēl ir dzīvs krievu gars, Krievijai ukraiņu nacionālā kustība (\rutxti{украинство}) vēl nav tik bīstama, bet tiklīdz austriešu-poļu valdībai izdosies īstenot savu sapni, iznīcinot visu krievisko Galīcijā un likt uz gadsimtiem aizmirst par kādreizējās Sarkanās (ukr. \uktxti{Червона Русь}, poļu \pltxti{Ruś Czerwona}~--- rutēņu apdzīvotās Galīcijas vēsturiskais nosaukums) pareizticīgās Krievijas pastāvēšanu, būs par vēlu un Krievija ar naidnieku vairs netiks galā.'' Bailes no ukraiņu separātisma iekļūšanas Krievijā lika Krievijas Iekšlietu un finanšu ministrijām pieņemt lēmumu par regulāru līdzekļu piešķiršanu ``palīdzībai piekarpatu rusīniem''. 1911.~gadā P.~Stolipins vienreizēji piešķīra 15~000 rubļu ``Maskavas orientācijas'' organizācijām Austrijas parlamenta vēlēšanās. Ik gadus 25~tūkstošus rubļu Finanšu ministrija un 60~tūkstošus Iekšlietu ministrija piešķīra jaunslāvu kustības līderim Valsts domes deputātam grāfam V.~Bobrinskim krievu un ``piekarpatu slāvu'' kultūras un izglītības iestāžu uzturēšanai un attīstības veicināšanai. Bez valsts dotācijām vēl 10--12 tūkstošus rubļu ienesa privātie ziedojumi. Centrālā vieta kultūras darbā tika veltīta krievu valodas izplatīšanai Galīcijā.

Protams, Austroungārijas valdība nevēlējās Krievijas ietekmes pieaugumu savā provincē. 1908.~gadā kāds rutēņu nacionālists nogalināja austriešu vietvaldi Galīcijā. 1910.~gadā izraisījās vardarbīgas sadursmes Lembergas Universitātē. Kaut pilsētas iedzīvotāju sastāvā ap 20\% bija rutēņi, viņi gandrīz nebija pārstāvēti pārvaldē. Vēl mazāk viņu bija iestādēs laukos. Tāpēc šai pašā gadā Austrijas (Cisleitānijas) parlamenta (Reihsrāta) Ukraiņu deputātu klubs, norādot, ka rutēņi gan pēc to pārstāvniecības ierēdņu, skolotāju un tiesnešu vidū, gan pēc to skolu skaita, gan vispārējās attieksmes pret tiem, ir ne tikai nostādīti sliktāk nekā citas tautas valstī, bet to stāvokli pat nevar salīdzināt ar pārējo tautu stāvokli (kas gan bija pārspīlējums, ar Austrijas varā esošo Balkānu tautu stāvokli to salīdzināt varēja~--- V.Š), griezās pie Reihsrātā pārstāvētajām partijām ar aicinājumu palīdzēt rutēņiem iegūt tiem pienākošās tiesības. Protams, panākumu nebija. Gluži pretēji, 1910.~gadā austriešu varas iestādes slēdza visas ``rusofīlās'' organizācijas Bukovinā (ukraiņu \uktxti{Буковина}, poļu \pltxti{Bukowina}, rumāņu \rotxti{Bucovina}) apgabals Karpatos, mūsdienās atrodas Ukrainas un Rumānijas teritorijā) un Černovcos (ukraiņu \uktxti{Чернівці}, rumāņu \rotxti{Cernăuţi}, vācu. \detxti{Czernowitz}, poļu. \pltxti{Czerniowce}, pilsēta Ukrainā), apvainojot tās spiegošanā un valsts nodevībā, kas gan nebija pierādāms.

Krievijā bija arī reāli domājoši politiķi, kuri iebilda pret krievu nacionālistu avantūristisko politiku. Spilgtākais no viņiem bija Krievijas iekšlietu ministrs P.~Durnovo. 1914.~gadā, kad pirmskara apstākļos krievu nacionālisti sapņoja par visu ``krievu zemju'' apvienošanu, viņš rakstīja par Galīciju: ``Mums nav izdevīgi nacionālās sentimentalitātes vārdā mūsu tēvzemei pievienot apgabalu, kurš ir zaudējis ar mums ikkatru dzīvu saikni. Uz niecīgu saujiņu pēc gara krieviskiem galīciešiem, cik mēs saņemsim poļu, ebreju, ukrainizētu uniātu? Tā saucamā mazepiešu kustība tagad mums nav bīstama, taču nedrīkst dot tai nostiprināties, palielinot nemierīgo ukraiņu elementu skaitu, jo šajā kustībā ir ļoti bīstamā ukraiņu separātisma dīglis, kurš labvēlīgos apstākļos var sasniegt gluži negaidītus apmērus.''

\strong{IV Valsts dome} (1912--1917) nacionālajā ziņā bija vēl monolītāka nekā III Dome. Tajā tika ievēlēti tikai 17 poļu (tas gan bija vislielākais nekrievu deputātu skaits), 9 no kuriem veidoja poļu \pltxti{kolo}, 5~--- baltkrievu-lietuviešu-poļu grupu. Valsts domes poļu deputātu skaits, attieksme pret tiem uzskatāmi parādīja, ka esošā vara nacionālajā jautājumā nevēlējās ieklausīties pat mērenās prasībās, ja tās neatbilda tradicionālajai Krievijas valdošo slāņu politikai. Kaut ārlietu ministrs S.~Sazonovs 1914.~gada janvārī rakstīja imperatoram Nikolajam II, ka tieši ``Krievijas lielvalstisko uzdevumu'' labad ``ir nepieciešams kaut kādā mērā apmierināt poļu sabiedrības saprātīgās vajadzības pašpārvaldes, valodas, skolas un baznīcas jomās\dots{}'', viņa balss netika uzklausīta. Valdība nezināja citas receptes nacionālajā politikā kā ``turēt un nelaist'' (\rutxti{держать и не пущать}). Daudziem mazākumtautību deputātiem bija zudušas cerības atrisināt rietumu apgabalu etnokonfesionālās problēmas ar darbu Domē. Faktiski Domi nekrievu deputāti varēja izmantot tikai kā tribīni un legālu centru, kurā koordinēt ārpusdomes politisko darbību.

Poļu iedzīvotāju dzīve \strong{Austroungārijas un Vācijas valdījumā esošajās provincēs} pirms Pirmā pasaules kara gan nepiedzīvoja tik krasas pārvērtības, kādas notika blakus esošajā Polijas karalistē, taču problēmu netrūka arī tur.

\strong{Galīcija} arī XX gadsimta sākumā palika agrāra province, kur ar lauksaimniecību nodarbojās 73\% poļu un 94\% ukraiņu iedzīvotāju. Zemes īpašums bija sadalīts ļoti nevienmērīgi. 992~000 zemnieku saimniecību, kuru lielums nepārsniedza 20~ha, īpašumā bija 44,2\% zemes, bet pārējie gandrīz 56\% atradās tikai 17~000~saimniecību rokās, kuras bija lielākas par 20~ha. Starp pēdējām bija arī 21~latifundija, kurām piederēja 8,5\% visas zemes. Tā kā muižnieki bija poļi vai pārpoļoti elementi, agrārais jautājums cieši savijās ar nacionālo.

XX gs. sākumā Galīcijā lauksaimniecības mašīnas tika lietotas 26,4\% visu saimniecību. Audzējamo kultūru vidū pārsvarā bija graudaugi, kuri aizņēma 63\% visu sējumu platību, un kartupeļi (13\%). Graudaugu raža no 1900. līdz 1914.~gadam pieauga par 23\%, kartupeļu~--- par 6\%. Taču pašu ražotās lauksaimniecības produkcijas Galīcijas iedzīvotāju patēriņam nepietika. Labība lielā apjomā tika ievesta no blakus esošās Ungārijas.

Pirms kara Galīcijā auga ogļu ieguve. Arī XX gadsimta sākuma depresija tās izaugsmi nepārtrauca (1913.~gadā tajā ražoja ap 2~miljoniem t ogļu). Auga arī svina rūdas, sāls ieguve, toties samazinājās iegūtās dzelzs un cinka rūdas daudzums. Ekonomikas lejupslīde, saistīta ar Bosnijas krīzi (1908) un Balkānu kariem (1912--1913) ietekmēja naftas ieguvi~--- tā 1913.~gadā samazinājās līdz 1907.~gada līmenim (1,1~miljons t).

Pēc ilgas cīņas Galīcijas politiskajiem spēkiem izdevās panākt vietējā Seima reformu. Tā vēlēšanu likumā nacionālo kūriju sistēma gan saglabājās, tomēr vēlēšanām bija jākļūst tiešām un aizklātām, vēlēšanu tiesības pienācās visiem, kuri piedalījās Austrijas Reihsrāta vēlēšanās, 27\% vietu tika piešķirtas rutēņiem. Diemžēl, sākoties Pasaules karam, likums tā arī nestājās spēkā. Ar cīņas par Seima vēlēšanu reformu noslēgšanos mainījās spēku samērs starp ietekmīgākajiem politiskajiem grupējumiem. Konservatīvo spēki gāja mazumā, toties pieauga ļudoviešu ietekme. 1907.~gadā Reihsrāta vēlēšanās, pateicoties demokrātiskas sociālekonomiskas un politiskas programmas izvirzīšanai, viņi guva ievērojamus panākumus. Cerības panākumu vēlēšanās pārvērst politiskā guvumā virzīja autoritatīvāko ļudoviešu līderi Galīcijā J.~Stapiņski uz savienības noslēgšanu ar konservatīvajiem. 1908.~gada tā arī tika noslēgta un līdz 1913.~gadam partija slēdza kompromisus lai saņemtu no valdības sīkas politiskas un ekonomiskas piekāpšanās. 1908.~gadā ļudovieši ieguva 19~vietas vietējā Seimā, 1911.~gadā~--- 24~mandātus Reihsrātā. Tiesa, vietējo ļudoviešu rindās nebija vienotības, darbojās dažādas grupas. No 1913.~gada Galīcijā zemnieku kustības labējā spārna pārstāvji nodibināja savu partiju ``Pjasts'' (\pltxti{Polskie Stronictwo Ludowe Piast}, pēdējais vārds partijas nosaukumā nāca no partijas preses orgāna nosaukuma), kas propagandēja dažādu sociālo grupu sadarbības ideju, kura it kā esot valdījusi pirmās Polijas valdnieku Pjastu dinastijas laikā. Partija izvirzīja lauksaimniecības kooperācijas programmu, propagandēja sadarbību ar valdību, buržuāzijas partijām. Partijas pārstāvjiem bija raksturīga noraidoša attieksme pret rutēņiem un ebrejiem. Uzstājoties par brīvu Poliju nākotnē, bet jau tuvākajā laikā plašu Galīcijas autonomiju, viņi noraidīja novada sadalīšanu poļu un ukraiņu daļās.

Galīcijas poļi, dzīvojot labvēlīgākos apstākļos nacionālajai attīstībai nekā parējās poļu zemēs, solidarizējās ar to iedzīvotājiem cīņā pret nacionālo apspiešanu, protestēja pret pretpoļu likumiem Prūsijā, 1912.~gadā visu Galīciju aptvēra protesta kampaņa pret Holmas guberņas izveidi.

Divdesmitā gadsimta sākumā daļa Austroungārijas politiķu ar erchercogu Franci-Ferdinandu priekšgalā stādīja uzdevumus Austroungārijas imperatora varā apvienot tās rokās esošo Galīciju ar Krievijai atņemamo Polijas karalisti. Šie plāni nepalika apslēpti arī Krievijas diplomātiem. 1909.~gadā Krievijas Ārlietu ministrijas ierēdņi atzīmēja, ka valdība Vīnē izstrādā autonomas poļu valsts izveides plānus un ved slepenu karu pret Krievijas ietekmi Polijas karalistē. Apstākļos, kad pasaulē saasinājās pretrunas starp diviem blokiem: Antanti jeb Trejsabiedrību Lielbritānijas, Francijas un Krievijas sastāvā un Trejsavienību Vācijas, Austroungārijas un formāli arī Itālijas sastāvā, kā arī starp divām poļu zemes anektējušajām valstīm: Austroungāriju un Krieviju, Galīcijas poļu virsslāņu vidū iezīmējās divas orientācijas.

Daļa poļu politiķu Austroungārijā par labāko ``poļu jautājuma'' risinājumu uzskatīja Austroungārijas pārvēršanu no duālistiskas par triālistisku monarhiju, kur Galīcija ieņemtu vienādu statusu ar Austriju un Ungāriju, paplašinoties uz Polijas karalistes un ukraiņu zemju rēķina. Taču otra poļu politiķu daļa orientējās uz Krieviju, jo daļa vietējās poļu buržuāzijas, nespējot konkurēt ar austriešu kapitālu, cerēja, ka Krievijā viņu intereses tiks vairāk ievērotas. Baidoties no ukraiņu nacionālās atbrīvošanās kustības, šādu orientāciju atbalstīja arī daļa Austrumgalīcijas poļu muižniecības. Šo slāņu intereses galvenokārt izteica nacionāldemokrāti, kuru propaganda bija vērsta pret ukraiņu un ebreju tiesību paplašināšanu un sludināja paļāvību uz Krieviju.

Galīcijas sociāldemokrāti un šeit apmetušies t.s. ``fraki'' (PPS Revolucionārās frakcijas biedri) savukārt apkaroja prokrievisko orientāciju. Galīcijas sociāldemokrāti ar I.~Dašinski priekšgalā atradās līdzīgās ar J.~Pilsudski teorētiskās pozīcijas. Abas partijas orientējās uz karu, kurā poļiem būtu jāatbalsta Austroungārija.

I.~Dašinskis 1912.~gada oktobrī, aicinot proletariātu cīnīties pret kara briesmām, reizē apstiprināja poļu gatavību cīnīties pret Krievijas carismu austriešu-krievu militāra konflikta gadījumā. Viņš vēlāk konstatēja: ``Internacionālē (Domāta 1889.~gadā dibinātā starptautiskā sociālistisko un strādnieku partiju apvienība, ko parasti sauca par II Internacionāli~--- V.Š.) nebija daudz sociālistisko partiju, kuras tā sapņotu par karu kā poļu sociālisti, bet precīzāk divas poļu partijas: PPS un Galīcijas un Silēzijas PSD''(\pltxti{Polska Partia Socjalno-Demokratyczna Galicji i Śląska Cieszyńskiego}).

J.~Pilsudskis un viņa piekritēji orientējās uz vispārēju sacelšanos Polijas karalistē, kuras laikā varētu radīt pastāvīgu poļu armiju, kura karotu Vācijas un Austroungārijas pusē un spētu uzvaras gadījumā panākt šo valstu piekrišanu patstāvīgas Polijas valsts radīšanai no karalistes un citām Polijai atsavinātajām zemēm. Šādai pozīcijai principiālu vērtējumu deva V.~Ļeņins: ``Būt par visas Eiropas karu tikai vienas Polijas atjaunošanas dēļ~--- tas nozīmē būt sliktākās markas marksistam, nostādīt nelielā poļu skaita intereses augstāk par simtiem miljonu cilvēku interesēm, kuri cieš no kara.''

Izmantojot vēsturisko 50.~gadadienu kopš 1863.~gada sacelšanās, 1913.~gadā Austroungārijā J.~Pilsudska piekritēji sarīkoja pretkrievisku kampaņu, kura atrada piekrišanu dažādos poļu sabiedrības slāņos. Mītiņi vai pārrunas par 1863.~gada sacelšanos nereti beidzās ar rezolūciju pieņemšanu, kurās bija uzsvērta nepieciešamība trenēties šaušanā. Tiesa, PPS-ļevica un PKunLSD apkaroja Galīcijas un Silēzijas PSD, ``fraku'' un dažādo militarizēto organizāciju biedru uzskatus.

Jāuzsver, ka Austroungārijas varas iestādes labvēlīgi uztvēra jebkādu pret Krieviju vērstu darbību, tāpēc arī tās teritorijā varēja izvērsties ne tikai poļu nacionālisti, bet arī krievu revolucionāri. Tā, no 1912.~gada jūnija līdz 1914.~gada augustam Krakovā un citur Austroungārijas rokās esošajās poļu zemēs uzturējās Krievijas boļševiku vadonis V.~Ļeņins. Viņš šeit varēja nodoties mierīgam teorētiskam darbam, uzrakstīt vairāk nekā 300~rakstu. Galīcija deva patvērumu arī Polijas karalistes nelegālo revolucionāro partiju (PPS ļevicas, PKunLSD) darbībai. Tās šeit izvērsa plašu izdevniecību darbu.

Vācu iestādes \strong{Pozenes provincē} gados pirms Pirmā pasaules kara turpināja agrāk aizsākto poļu ierobežošanas politiku. Lielus rūpestus tām sagādāja 1908.~gadā Prūsijas landtāga pieņemtais zemes atsavināšanas likums, kas paredzēja poļu muižniekiem piederošo muižu atsavināšanu un piešķiršanu vācu zemniekiem~--- ieceļotājiem poļu apgabalos. (Tiesa, likums Vācijas Reihstāgā izskanējušo protestu rezultātā tā gandrīz arī netika piemērots. Tikai 1912.~gadā bija 4 tā pielietošanas gadījumi.) Nobela prēmijas literatūrā laureāts H.~Senkevičs griezās ar aicinājumu pie visas pasaules literatūras un mākslas pārstāvjiem pievienoties protestiem pret nacionālo apspiešanu. Protestos aktīvi iesaistījās desmiti zinātnieku, rakstnieku un citu kultūras darbinieku. Vācijas oficiālie dokumenti rādīja, ka gadu desmitos pirms Pirmā pasaules kara kolonizācijas vajadzībām iegūtā zeme galvenokārt (72\%) bija iepirkta nevis no poļu, bet gan no vācu zemes īpašniekiem. Poļu rokās esošo zemju daudzums pat pieauga. Vācu šovinistu plāni izstumt poļu iedzīvotājus no viņu apmešanās vietām beidzās ar neveiksmi. Faktiski salīdzinājumā ar Krievijas impēriju, kur Polijas karalistē poļiem atsavināja daudz vairāk zemes, Vācijā tas tika darīts krietni mazākos apmēros, taču jāievēro, ka patvaldnieciskā Krievija nemaz nepretendēja uz tiesiskas valsts statusu Eiropas sabiedrības acīs, bet Vācija to darīja.

Rūpniecības attīstībā Pozenes novads joprojām atpalika no Polijas karalistes. 1907.~gadā rūpniecības strādnieku skaits šeit sasniedza 162~tūkstošus. Poļu sīkburžuāziskie elementi ar poļu akciju sabiedrību atbalstu pirms Pirmā pasaules kara nostiprināja savas pozīcijas. Visu poļu rūpniecisko uzņēmumu, ieskaitot arī pašus mazākos, īpatsvars no 1882. līdz 1907.~gadam pieauga no 50,1\% līdz 58\%, tirdzniecības uzņēmumu~--- no 23,7\% līdz 41,6\%. XX gadsimta sākumā Pozenes buržuāzija finansiāli bija atkarīga no muižnieku kapitāliem. Visai izplatītajās šeit akciju sabiedrībās vairāk nekā puse akciju piederēja muižniekiem.

Joprojām turpinājās dažādas nacionālās sadursmes. 1908.~gadā Prūsijā tika pieņemts likums, kurš cita starpā pieprasīja poļu biedrībām savas sapulces noturēt vācu valodā. Tiesa, tas paredzēja arī izņēmumus apdzīvotām vietām, kur cittautiešu īpatsvars pārsniedza 60\% no visiem iedzīvotājiem. Tajās šie nevācu iedzīvotāji vēl 20~gadus varēja sapulcēs lietot savu dzimto valodu. Tomēr 45\% Prūsijas poļu sanāksmēs poliski vairs runāt nedrīkstēja.

Poļu mantīgo slāņu pārstāvji spēja saglabāt ietekmi masās tikai ar nacionālā radikālisma ideju palīdzību. Konservatīvo lojalitāte varas iestādēm padarīja tos nepopulārus. Tāpēc nacionāldemokrāti Pozenes novadā, darbojoties dažādās sabiedriskās organizācijās, lietoja radikālas frāzes, uzstājās pret lojālismu, sludināja cīņu pret visu vācisko. Rezultātā Pozenes provincē ievērojami pieauga Nacionāldemokrātiskās partijas biedru skaits. Vēlēšanās pieauga par poļu nacionālistu kandidātiem balsojušo skaits. Tā, 1898.~gada reihstāga vēlēšanās par tiem bija nodots 244~000 balsu, 1903.~gadā~--- 348~000, 1907~--- 454~000, bet 1912.~gadā~--- 442~000.

Nevienmērīgi attīstījās Vācijas Sociāldemokrātijas un Poļu sociālistiskās partijas Prūsijā (\pltxti{Polska Partia Socjalistyczna Zaboru Pruskiego) attiecības. Ja 1906.~gadā to vienošanās rezultātā otrā pēc būtības kļuva par} Vācijas Sociāldemokrātijas autonomu daļu, tad 1913.~gadā šī vienošanās tika anulēta un notika galīga šo partiju attiecību saraušana. Šis fakts atspoguļoja Poļu sociālistiskās partijas Prūsijā tuvināšanos PPS-frakcijai un poļu nacionālistiem, kuri iestājās par bruņotu cīņu pret Krieviju Austroungārijas un Vācijas pusē.

Neraugoties uz poļu nacionālisma pieaugumu, 1912.--1914.~gadā Prūsijā tika pieņemta virkne jaunu pretpoļu likumu. Liberāli-konservatīvais vācu vēsturnieks un politiķis H.~Delbruks šai sakarā norādīja, ka Vācijas austrumu nomaļu politika (``\detxti{Ostmarkenpolitik}'') ne tikai bija kļūdaina un neko nedeva vācietībai, ne tikai dižciltīgajos un garīdzniecībā, bet visā poļu sabiedrībā radīja nepārvaramu neuzticēšanos visam vāciskajam. Vācijas reihstāga vairākums, kuru veidoja sociāldemokrāti, poļu frakcija, katoliskās Centra partijas un Elzasas deputāti, pat izteica Vācijas kancleram T.~Betmanam-Holvegam neuzticības votumu, bet tas neko nemainīja.

Tomēr jāatzīst, ka Pozenes provincē dzīvojošo poļu ģermanizācija risinājās sekmīgāk nekā rusifikācija Vislas novadā, jo Pozenes provincē poļu īpatsvars bija mazāks nekā Krievijas valdījumos.

Vērtējot kopumā situāciju visās poļu teritorijās, jāsaka, ka 1914.~gadā nedz nacionālās, nedz sociālās pretrunas, neraugoties uz to sāpīgumu, pašas par sevi vēl nevarēja izsaukt Polijā revolūciju. Tā varēja notikt tikai tad, kad pasaules karš tās bija daudzkārt saasinājis. Pastāvošās pretrunas mīkstināja emigrācija, it sevišķi no ekonomiski atpalikušās Galīcijas. Sezonas migrācija aptvēra ap pusmiljonu cilvēku gadā, bet patstāvīgā emigrācijā no līdz 1914.~gadam bija devušies no 3,5 līdz 4~miljoniem cilvēku. Visvairāk emigrantu no visiem poļu apdzīvotajiem apgabaliem apmetās ASV, sastādot tur vairākus miljonus lielu poļu kopienu. Polijas sabiedrībā arī nebija izveidojušies vienoti uzskati par savām interesēm iespējamajā karā.

Kad sākās Pirmais pasaules karš, poļu iedzīvotāju vairākums ļāva sevi iesaukt visu triju valstu armijās bez kādas lielākas aizrautības no R.~Dmovska sludinātā ``neoslāvisma'' vai J.~Pilsudska ``aktīvisma''. Ir dati, ka pirms Pirmā pasaules kara līdz 10\% no Krievijas armijas virsniekiem bija poļi. Arī kara gados līdz Februāra revolūcijai Krievijā no poļu iedzīvotāju puses nebija vērojama masveida atteikšanās no dienesta karojošo valstu armijās vai kādas masveida demonstrācijas pret karu.

\subsection{Polijas teritorijas Pirmajā pasaules karā. Valsts proklamēšana}

Jaunas iespējas cīņai par nacionālo neatkarību Polijai atnesa Pirmais pasaules karš, kurš sākās 1914.~gada augustā.

Polijas karalistes iedzīvotāju vidū, protams, bija tādi, kuri vēlēja Krievijai sakāvi karā, taču uz kopējā fonā šie noskaņojumi bija vāji pamanāmi. Poļu sabiedrība atradās vai nu tolaik pret Krieviju labvēlīgi noskaņoto nacionāldemokrātu ietekmē vai bija neitrāla. Līdzīgi lojāli pret savām valdībām izturējās Vācijas un Austroungārijas poļu tautības pavalstnieki.

Tajā laikā pasaulē atradās ļoti maz politisko spēku, kuri iedrošinājās principiāli un atklāti nostāties pret karu kā imperiālistisko valstu, to valdošo slāņu izvērstu citu valstu aplaupīšanas akciju. Poļu deputāti savu valstu parlamentos balsoja par kara kredītu piešķiršanu valdībām un vēlēja savām valdībām un imperatoriem uzvaru karā.

Arī poļu deputāti Krievijas IV Valsts domē pret karu neprotestēja, bet vienīgi domāja, kā to izmantot poļu labā. Tā daļa Domes poļu \pltxti{kolo} deputātu 1914.~gada oktobrī-decembrī piedalījās kopīgās sanāksmēs ar krievu kadetiem Pēterburgā, Maskavā, Kijevā un Varšavā par poļu jautājumu. Šajās sanāksmēs krievu puse kā galveno izvirzīja jautājumu par poļu attieksmi pret nacionālajiem mazākumiem nākamajā apvienotajā Polijā. Krievu kadeti uzskatīja, ka poļiem jāgarantē nacionālajiem mazākumiem politisku un valodas brīvību, kā arī citas priekšrocības sociālajā likumdošanā un ekonomikā. Poļu dalībnieki sajutās aizskarti un paziņoja, ka attiecību nodibināšana ar nacionālajiem mazākumiem ir nākamās autonomās Polijas iekšēja lieta. Tā pirmā domu apmaiņa starp kadetiem un poļu politiskajām aprindām noveda pie jūtamas attiecību pasliktināšanās.

Pat sociāldemokrātu vairākums atbalstīja savu valstu valdības, aicināja savus piekritējus cīnīties par nacionālajām interesēm. Tikai dažās zemēs daļa sociāldemokrātu palika konsekventi internacionālās pozīcijās. Galēji kreisās poļu politiskās partijas kara jautājumā nostājās revolucionārās pozīcijās.

Starptautiskā Sociālistiskā biroja pēdējā sēdē 1914.~gada 29--30.~jūlijā (pēc jaunā stila) PKunLSD un PPS-ļevicas pārstāvji R.~Luksemburga un M.~Horvics kopīgi paziņoja, ka ``kara gadījumā mēs uzstāsimies pret carismu, bet roku rokā ar Krievijas proletariātu'', nosodot PPS biedrus, ``kuri solīja austriešiem, ja sāksies karš, uzsākt Polijā provācisku sacelšanos''. Sākoties karam, nosauktās partijas saglabāja šo principiālo nostāju. Jau 1914.~gada 1.~augustā, kad visur notika iesaucamo mobilizācija, sociāldemokrāti Varšavā izkāra plakātus ``Nost karu!''. 2.~augustā Varšavā pieņemtais PKunLSD, PPS-ļevicas un Bunda manifests karu raksturoja kā imperiālistisku, iekarošanas, prettautisku (PKunLSD rindās šai kara kvalifikācijai piekrita gan ``žondovieši'', gan ``razlomovieši'', kuri 1916.~gadā apvienojās) un aicināja proletariātu iet ``tikai revolucionārās cīņas ceļu''. Manifestā tika tieši paziņots, ka proletariātam, ``lai atrisinātu savus uzdevumus'' ``ir jāiekaro vara, tā jāsagrābj savās rokās'', ka: ``proletariāts piesaka cīņu savām valdībām, saviem apspiedējiem~--- un bargās strādnieku revolūcijas elpa tuvojas \citespace{} Mūsu zemes proletariāts savas revolucionārās uzstāšanās saskaņos ar tādām pat Krievijā un Eiropā, cik tās izrietēs no revolucionāriem centieniem gāzt tagadējās valdības un nodibināt tautas valdības''. 1915.~gadā PKunLSD uzskati kara jautājumā ievērojami tuvinājās V.~Ļeņina pozīcijai. PKunLSD valdes sagatavotajā 1.~maija uzsaukumā bija formulēts uzdevums ``buržuāzisko karu pārvērst par sociāli-revolucionāru karu'', kas faktiski atkārtoja V.~Ļeņina izvirzīto lozungu par imperiālistiskā kara pārvēršanu pilsoņu karā. Analoģiski lozungi bija ietverti arī citos partijas dokumentos.

Jau vēlāk, 1915.~gada vasarā sakarā ar krievu karaspēka atkāpšanos, PKunLSD paziņoja: ``Proletariāts ir pats sev sabiedrotais. Tā ienaidnieks ir buržuāzija, tā ienaidnieki vienādā mērā ir gan carisma, gan vācu un austriešu valdības. Proletariāts nepazīst nekādu ``orientāciju''. Tam ir tikai viena ``orientācija''~--- sociāldemokrātiska.'' Taču poļu sociāldemokrāti arī kara laikā joprojām noraidīja nāciju pašnoteikšanās lozungu. Tiem nebija skaidras nostājas Polijas jautājumā laikā, kad nacionālistiskās noskaņas, tāpat kā visur, arī poļu tautā kara laikā bija pārsvarā pār internacionālām, šķiriskām un pat vispārcilvēciskām interesēm.

Dažādo poļu apdzīvoto apgabalu iedzīvotāju vidū valstiskās un nacionālās piederības sajūtas veidoja ļoti sarežģītu ainu.

Piemēram, kādreizējais Tornas (Toruņas) iedzīvotājs vācietis R.~Vingebdorfs vēlāk, jau starpkaru periodā, stāstīja kā 1915.~gadā, nonācis krievu gūstā, viņš kopā ar citiem gūstekņiem tika vests caur toreiz vēl Krievijas armijas rokās esošo Varšavu. Vietējie poļu iedzīvotāji lamāja un apspļaudīja gūstekņus. Taču neviens no poļiem~--- Vācijas pavalstniekiem pat domās neesot pieļāvis iespēju sacīt: ``Es neesmu vācietis, es esmu polis''. Kad vācu gūstekņi vēlāk tika šķiroti pēc tautībām un slāvu izcelsmes, Centrālo valstu pavalstnieki tika atstāti Krievijas Eiropas daļā, bet vācieši sūtīti uz Sibīriju, Prūsijas poļi labāk devušies uz Sibīriju, un Krievijā dibinātajos poļu leģionos nebijis gandrīz neviena Prūsijas poļa. Gūstā nonākušie Austroungārijas poļi, ar kuriem autors sarunājies ceļā uz Krievijas austrumiem, sapņojuši gan par Polijas valsts atjaunošanu, tomēr nepretendējuši uz ap Toruņu esošā, kādreiz Vācu ordenim piederējušā apgabala pievienošanu tai. (Te nu gan jānorāda, ka, iespējams, minētie Austroungārijas poļi vienkārši nevēlējās izraisīt nevajadzīgu strīdu ar vāciskas izcelsmes gūsta biedriem par toreiz neparedzamās nākotnes jautājumiem.)

Savukārt latviešu vēsturnieks un diplomāts Alfrēds Bīlmanis 20.~gadu sakumā rakstīja, ka, karam sākoties, Krievija sagaidījusi arī tūlītēju poļu vispārēju sacelšanos ``Pievislas apgabalā'', bet vācieši bijuši poļu tik ienīsti, ka ``viņi atrada pat krievu slogu par vairāk pieņemamu un izvēlējās no diviem ļaunumiem mazāko'' un ``tamdēļ arī negaidīti gludi noritēja Polijā mobilizācija''. Ar to gan pretrunā ir paša A.~Bīlmaņa atzinums, ka ``tikpat gludi [tā notika] arī Vācijā un Austrijā''. Jāsecina, ka arī poļu apdzīvotajos apgabalos, tāpat kā citur, vietējie iedzīvotāji kara sākumā lielā mērā bija dezorientēti un nespēja saskatīt ceļu uz nacionālo un sociālo atbrīvošanos.

Daudzu poļu iedzīvotāju, domājams, ka arī visas poļu tautas vairākuma, apziņā nacionālās atbrīvošanās uzdevums tika stādīts augstāk par sociālās apspiestības likvidācijas uzdevumu, kā arī pacifistiskiem lozungiem. Tiesa gan, šīs nacionālās atbrīvošanās cerības saistījās ar dažādām karojošajām pusēm. Kā atzīmējuši poļu vēsturnieki J.~Holzers un J.~Molenda, kuri pētījuši poļu noskaņojumu Pirmā pasaules kara laikā, kara sākumā Galīcijas un arī daļas Polijas karalistes iedzīvotāju simpātijas bija Austroungārijas pusē, kur poļu nacionālā apspiestība bija relatīvi vājāka un daudzi no viņiem paši uzstājās nacionālo apspiedēju lomā. Taču drīz noskaidrojās, ka Austroungārijas pavēlniecība nedomā piešķirt J.~Pilsudskim kaut cik plašas pilnvaras un negatavojas apvienot Polijas valsti, kaut vai Austroungārijas ietvaros. Poļu zemēs, kuras ietilpa Vācijas sastāvā, diezgan izteiktas bija prokrieviskās simpātijas. Dažādo kara notikumu gaitā noskaidrojās, ka vislabāk pret vietējiem iedzīvotājiem izturējās krievu karavīri, vietējo iedzīvotāju palīdzība Krievijas armijai tika labi apmaksāta. Tāpēc radās, ``ja ne labvēlīga, tad, mazākais, neitrāla attieksme pret Krieviju''.

Dažādu tautu nacionālisma ``sprādziens'' bija sekas gan vispārējām totālā kara grūtībām, gan karojošo valstu jaunajai politikai cīņā pret saviem pretiniekiem. To pašu var teikt par poļu nacionālismu. Tas jau pirms kara bija valdošā ideoloģija Polijā. Pēc kara sākuma tā ietekme vēl pieauga. Taču situācija, kad Poliju sadalījušās lielvaras karoja savā starpā, ne tikai deva cerības uz paverdzinātāju sakāvi, uz Polijas patstāvības atjaunošanu, bet poļu tautai nacionālās brīvības centienu īstenošanā radīja arī papildus grūtības.

Pirmkārt, poļiem nācās karot dažādu valstu armijās, poļi karoja pret poļiem. Visas Poliju sadalījušās valstis vēlējās izspēlēt ``poļu kārti'' pret saviem pretiniekiem, iesaistīt vietējos iedzīvotājus informācijas vākšanai par notiekošo frontes otrā pusē, izmantot poļu dzīvo spēku savā labā. 1914--1915.~gados abās frontes pusēs karoja apmēram 2~miljoni, pēc citiem datiem~--- 2,5~miljoni poļu un pat 3~miljoni cilvēku ``no poļu zemēm'' (tātad~--- arī nepoļi). Līdz kara beigām ārvalstu dienestā nonāca ap 3,4~miljonu Polijas iedzīvotāju, zaudēja dzīvību ap 400~000, ievainoti bija ap 800--900~tūkstošu.

Otrkārt, poļu apdzīvotās teritorijas jau kara sākumā kļuva par kaujas lauku. Vācijas un Austroungārijas impēriju austrumu un Krievijas impērijas rietumu provinces kara laikā kļuva par kara darbības zonu, tur notika iedzīvotāju migrācija, gan brīvprātīga, gan piespiedu, mainījās varas.

Visaktīvākā kara darbība poļu zemēs norisa 1914.~gada augustā~--- novembrī, kad krievu karaspēks ieņēma Austroungārijai piederošo Austrumgalīciju, bet atdeva Vācijai Polijas karalistes rietumu rajonus. Pēc tam situācija stabilizējās līdz 1915.~gada maijam, kad austriešu un vācu armijas pārrāva fronti Karpatos, atkaroja Galīciju, piespieda Krievijas armiju atstāt Polijas karalisti. 5.~augustā Vācijas armija ieņēma Varšavu, septembra vidū Viļņu, austrieši okupēja Polijas karalistes dienvidaustrumu daļu.

Protams, poļu teritorijas tika izpostītas. Izņēmums bija vienīgi daļa Prūsijas varā esošo poļu apdzīvoto teritoriju un Rietumgalīcija.

Interesanti, ka Krievijas varas iestādes poļu rūpniecībai nodarīja mazākus zaudējumus (tie sastādīja 18\%), kā citas karojošās valstis. Krieviem Polijas karaliste bija viena no valsts provincēm, pie tam šai laikā pilnīgi lojāla, tāpēc līdz tās atstāšanai viņi izturējās korekti. Centrālo valstu politiķi un militāristi uz šīm teritorijām raudzījās citādi. Viņiem tās bija uz laiku okupētas, to pilnīga inkorporācija pēc kara netika plānota. Tāpēc tās arī netika saudzētas. Vācijas un Austroungārijas acīs bijušās Krievijas teritorijas kļuva par negausīgas ekspluatācijas objektu. Attiecīgi šo valstu poļu rūpniecībai nodarītie zaudējumi pēc poļu vēsturnieka S.~Keņeviča datiem, sastādīja 52\% un 22\%. Pie tam tikai 4\% zaudējumu bija nodarīti kara darbības rezultātā. Versaļas miera konferencē 1919.~gadā poļu delegācija materiālos zaudējumus no kara darbībās un abu pušu izlaupīšanas politikas noteica 73~miljardu franču franku apmērā.

Drīz visas Polijas, bet galvenokārt Vācijas un Austroungārijas rokās nonākušo bijušo Krievijas apgabalu, poļu pilsoniskie un sīkpilsoniskie iedzīvotāji nosacīti sadalījās divās grupās: t.s. ``pasīvistos'' un ``aktīvistos.

``Pasīvisti'', neizrādot kādu pretestību okupantiem, atteicās ar tiem politiski sadarboties, uzskatot Vācijas impēriju par galveno poļu ienaidnieku, savas cerības, tāpat kā pirms kara, saistīja ar poļu apgabalu apvienošanu Krievijas sastāvā, kas ļautu augt to patstāvībai, ja arī nedotu uzreiz pilnīgu neatkarību. Arī tad, kad Krievija cieta sakāves karā, ''pasīvisti'' gaidīja Antantes uzvaru. ``Pasīvistu'' kodolu veidoja nacionāldemokrāti.

``Aktīvisti'' piekrita plaši sadarboties ar vāciešiem un austriešiem, savas cerības saistīja ar Vācijas un Austroungārijas uzvaru karā. ``Aktīvistu'' kodolu veidoja J.~Pilsudska piekritēji, kuri sociālā ziņā nāca galvenokārt no izputējušo šļahtiču un pilsētu sīkburžuāzijas aprindām. Tiesa, viņu cerības ātri iegūt plašu tautas slāņu un pirmkārt zemniecības atbalstu nepiepildījās. ``Aktīvisti'' okupētajos apgabalos baudīja priviliģētu stāvokli. Darbojās viņu partija~--- PPS-Revolucionārā frakcija, pastāvēja viņu vadītie militārie formējumi~--- poļu leģioni. ``Aktīvistu priekšgalā nostājās bijušais ``sociālistu'' līderis, (Kā norādījis krievu vēsturnieks G.~Matvejevs, šai laikā J.~Pilsudskis uz visiem laikiem šķīrās no sociālistu partijas, kas norāda, ka viņš nekad nebija bijis sociālists šā vārda precīzā nozīmē. PPS viņam kalpoja tikai kā instruments galvenā dzīves mērķa~--- Polijas karalistes atbrīvošanai no Krievijas.) militarizēto poļu vienību ``komandants'' Galīcijā J.~Pilsudskis, kuru angļu vēsturnieks A.Dž.P.~Teilors šai sakarā nosauca par ``militāru avantūristu''. J.~Pilsudskis cerēja pēc Krievijas sakāves panākt Polijas karalistes apvienošanu ar Galīciju (Vācijai piederošās zemes gan paliktu ārpus jaunā poļu zemju veidojuma) un Austroungārijas pārveidi no duālistiskas par triālistisku valsti, kurā līdzās Austrijai un Ungārijai kā līdzvērtīga sastāvdaļa ietilptu arī Polija. 1915.~gada vasarā vēstulē vienam no Galīcijas politiķiem V.~Javorskim viņš rakstīja: ``\dots{}Kara politiskais mērķis, kuru es no paša sākuma sev stādīju, bija un līdz šim paliek Galīcijas un [Polijas] karalistes saplūšana Austroungārijas monarhijas sastāvā. Es neuzskatīju un neuzskatu, ka šajā karā var izcīnīt labākus nosacījumus Polijai''. Taču, kā saka, ``apetīte nāca ēdot'', J.~Pilsudska plāni kara gaitā mainījās. Ierobežoto mērķi kara gaitā nomainīja cits~--- panākt Polijas neatkarību.

Pēc krievu politiķa V.~Černova atmiņām, J.~Pilsudskis jau 1914.~gada sākumā Parīzē Ģeogrāfijas biedrības zālē nolasītā lekcijā paredzējis, ka sākoties Eiropā karam, tajā agri vai vēlu Rietumu valstu pusē iesaistīsies arī ASV un karš beigsies ar Centrālo valstu sakāvi. Taču J.~Pilsudska biogrāfijas autori Darja un Tomašs Nalenči, izanalizējuši dokumentus, kuri attiecas uz doto lekciju, ir nonākuši pie secinājuma, ka, iespējams, J.~Pilsudskis tādu versiju arī izteicis, taču tas nevarēja notikt minētajos apstākļos. Pirmkārt, publiski izklāstīt savus plānus būtu politiski naivi, savienības ar Austroungāriju pagaidu rakstura atzīšana varēja vest pie runātāja izsūtīšanas no Galīcijas, ko viņš nekādā ziņā nevarēja vēlēties. Otrkārt, no uz lekciju sanākušajiem apmēram 500~poļu diasporas Francijā pārstāvjiem neviens nav apstiprinājis it kā V.~Černova dzirdēto. Treškārt, lekcijā piedalījās cariskās \emph{ohrankas} (Apsardzes daļas) aģents, kurš uz karstām pēdām par to uzrakstīja desmit mašīnraksta lappušu garu ziņojumu. Tajā ne vārda nav par J.~Pilsudska izteikto kara prognozi.

Tomēr literatūrā tiek citēts J.~Pilsudska it kā 1914.~gadā teiktais: ``Polija būs brīva, ja Vācija sakaus Krieviju un Francija sakaus Vāciju. Mums ir tam jāpalīdz''. Arī pēc vairākuma mūsdienu poļu vēsturnieku uzskata, viņš rēķinājies ar to, ka Vācija un Austroungārija sakaus Krieviju, pirms tās pašas savukārt zaudēs karu Rietumvalstīm. Pēc tam uz triju lielvaru drupām varētu rasties neatkarīga Polijas valsts. Tāpēc poļiem, lai tie kara beigās tiktu uzklausīti, ir jārada savi bruņotie spēki Centrālo valstu pusē un vispirms jāpalīdz tām salauzt Krievijas vara. Ja tas izdotos, tad vajadzības gadījumā varētu vērsties arī pret agrākajiem sabiedrotajiem.

Šādu skaidru notikumu prognozēšanu jau no paša sākuma gan var apšaubīt, jo pirms kara J.~Pilsudskis Galīcijā gatavoja poļu jaunatni cīņām tikai pret Krieviju, līdz 1917.~gadam faktiski sadarbojās tikai ar Austroungārijas un Vācijas, bet ne Antantes valstu valdībām.

Arī paša J.~Pilsudska izteikumi ir pretrunā ar versiju, ka viņš ģeniāli paredzējis notikumu attīstību. Jau 1913.~gada sākumā viņš izteicās: ``\dots{}nākotne~--- tas ir vienādojums ar daudziem nezināmajiem, kuru, attiecīgi, nav iespējams atrisināt.'' Šo pašu domu J.~Pilsudskis attīstīja kādā intervijā 1924.~gada februārī: ``Es no paša sākuma pretēji visiem citiem apgalvoju, ka karš turpināsies ievērojami ilgāk, nekā daudzi paredzēja. Rezultātā abas puses, uzvarētājs un uzvarētais, būs izsmēluši spēkus un novājināti. \citespace{}

Tad radīsies iespēja ietekmēt Polijas likteni, ja mūsu rīcībā būs tādi spēki. Es neizdarīju nekādus pieņēmumus, kas uzvarēs~--- vieni vai otri. Es biju pārliecināts, ka gan vieni, gan otri kļūs vājāki un mēs spēsim to izmantot \dots{} Minēšu vairākus savus jokus un aforismus par šo tēmu. Es ne reizi vien atkārtoju, taču kā joku, ka Krieviju sākumā sakaus centrālās valstis, bet pēc tam tās pašas sakaus Antante. \citespace{} Es salīdzināju mūsu uzdevumu ar jāšanas sacensībām. Zirgi joņo uz finišu, bet mēs kā muša sēžam uz viena no tiem auss. Finiša brīdim tuvojoties, muša aizlido no nogurušā zirga un sasniedz mērķi pirmā''. Faktiski pretēji tam, ko no 1918.~gada sludināja J.~Pilsudska piekritēji, viņam kara laikā ilgi nebija priekšstata par radušās situācijas attīstības galīgo rezultātu. (Tāda nebija un arī nevarēja būt arī citiem Eiropas politiķiem.) Par to liecina arī tas, ka poļu ``muša'' bija nosēdusies uz tāda ``zirga'' kā Austroungārija auss, kura poļiem bija devusi zināmu autonomiju, ļāva savā teritorijā gatavot poļu militāros kadrus, bet kurai, kā vēlāk izrādījās, bija niecīgas izredzes nonākt pie finiša kā uzvarētājai. Domājams, ka sākotnējo ``zirga'' izvēli noteica vienīgi Austroungārijas salīdzinoši labvēlīgākā attieksme pret poļiem.

Pats J.~Pilsudskis savu savienību ar Austroungāriju vēlāk tēloja kā tīri taktisku, pirmkārt vērtējot iespēju radīt nākamās poļu armijas militāros kadrus. Pēc viņa vārdiem, ja kādā citā valstī ``lai tā būtu mūsu austrumu kaimiņš vai pat Vācija'' tas būtu vieglāk izdarāms, viņš ``bez svārstībām dotos turp''. Jau minētajā lekcijā Parīzē 1914.~gada sākumā, pēc cariskās ohrankas aģenta ziņojuma, J.~Pilsudskis, atvairot pārmetumus par kalpošanu Austrijas interesēm, paziņoja: ``Jā, Austrija ir valsts, sastāvoša no dažādām daļām, un viena no tām ir Galīcija, nevis pakļauta tai, bet veidojoša daļu no šīs savienības, un šajā gadījumā kalpot Austrijai~--- nozīmē kalpot Polijas lietai. Pievienot Polijas karalisti Austrijai~--- nozīmē cīnīties par Polijas neatkarību''.

Visus savus pūliņus J.~Pilsudskis veltīja poļu armijas izveidei, redzot tajā galveno ``trumpi'' starptautiskajā kara ``spēlē''. Viņa proponētais kurss turpināja Napoleona laika poļu leģionu un XIX gadsimta lielo sacelšanos tradīcijas, tam piemita romantiska ievirze, tas vienus poļu slāņus piesaistīja, bet citus atbaidīja un lika tiem atbalstīt ``saprātīgāko'' ``pasīvistu'' ``reālismu''.

Jāuzsver, ka, kaut abas koncepcijas kara sākumā gan bija it kā diametrāli pretējas, vēlāk, veidojoties Polijas valstij, dažu mēnešu laikā to piekritēji, neraugoties uz daudzām pretrunām, spēja sastrādāties.

Sākotnēji vēstures priekšplānā izvirzījās ``aktīvistu'' darbība. Pirms izvirzīt kādas prasības, J.~Pilsudskis vēlējās nodemonstrēt savu spēku, parādīt visiem, ka ``Polija vēl ir dzīva'', tāpēc 1914.~gada 6.~augustā viena viņa ietekmē esošo poļu militarizēto organizāciju izveidota rota pārgāja Austroungārijas~--- Krievijas robežu Krakovas tuvumā, tā faktiski atklājot kara darbību Austrumu frontē. (Faktiski atkārtojās 1812.~gada situācija, kad poļi sveša karaspēka sastāvā iebruka Krievijā.)

Krievu publicistikā ironiski atzīmēts, ka cerībā aiz robežas atrast kavalērijai derīgus zirgus, poļu strēlnieki bija pat paņēmuši līdzi seglus. Tiesa, kad parādījušies krievu kazaki, poļi bijuši tik saprātīgi, ka atkāpušies uz Austroungārijas teritoriju. Tik tiešām, vācu autora V.~Noltinga grāmatā izdevās atrast datus, ka pirmajā no astoņiem vīriem sastāvošajā izlūkpatruļā trīs vīri nesuši līdzi seglus. Taču autora rīcībā nav ziņu par kādām poļu strēlnieku sadursmēm ar kazakiem.

Krievu karaspēks no pierobežas bija atvilkts, nekādas sadursmes nenotika. Nākamajā dienā rotai pievienojās vēl ap 450~cilvēku, kopā poļu kaujinieku skaits sasniedza ap 600~cilvēku. 10.~augustā tiem sekoja vēl ap 2~500 poļu strēlnieku. 12.~augustā J.~Pilsudskis savas vienības priekšgalā iesoļoja guberņas pilsētā Kelcē (\pltxti{Kielce)}, un apmetās gubernatora pilī. Tas tad arī bija Polijas karalistes statusa maiņas plāna realizācijas kulminācijas punkts.

J.~Pilsudskis uzskatīja, ka krievu armija aizstāvēs tikai Vislas līniju, tās kreisajā krastā atstājot vienīgi nelielas aizsega vienības, ar kurām poļu strēlnieki, kurus atbalstīs vietējie iedzīvotāji, tiks galā. Netika izslēgta arī kārtējās poļu sacelšanās pret Krievijas impēriju iespēja. (Uz to cerēja arī Krievijas pretinieču~--- Vācijas un Austroungārijas karaspēku pavēlniecības. Polijas karalistes teritorijā vēl pirms kara tika izplatītas skrejlapas ar aicinājumiem kara uzsākšanas brīdī uzsākt sacelšanos pret Krieviju.) Taču Krievijas armija nedomāja bez kaujām atdot daļu Polijas karalistes teritorijas, gatavoja uzbrukumu Krakovai un Lembergai, tāpēc nevarēja atstāt Varšavu, kas uzreiz pasliktinātu stratēģisko stāvokli. Izvēršoties krievu armijas uzbrukumam Galīcijā, poļu un austriešu daļām nācās atkāpties. Austrieši 1914.~gada septembra pirmajā dekādē atvilka spēkus uz agrākajām pozīcijām, kopā ar tiem atkāpās arī poļi, kuri šai laikā tā arī nebija piedalījušies nevienā lielākā kaujā. Tomēr 13.~septembrī viņi saņēma pateicību no austriešu Virspavēlniecības par priekšzīmīgo atkāpšanos. Nepiepildījās J.~Pilsudska godkārīgais plāns pirms austriešiem un vāciešiem ieiet Varšavā, radīt poļu armiju.

Taču galvenais bija tas, ka nepiepildījās cerības Krievijai piederošajās poļu apdzīvotajās teritorijās izraisīt sacelšanos. J.~Pilsudskis pat paziņoja par Nacionālās valdības izveidi Varšavā. Tas gan bija blefs, kaut kādas sacelšanās Polijas karalistē nenotika. Visi J.~Pilsudska piekritēju mēģinājumi to izsaukt beidzās neveiksmīgi. ``Atbrīvotajos'' ciemos poļu iedzīvotāji nekādu entuziasmu un vēlmi cīnīties pret carismu neizrādīja. Slikti bruņotu cilvēku ar poļu zīmotnēm uz cepurēm parādīšanās, kuri pie tam uzreiz uzsāka pārtikas un zirgu rekvizīcijas, neizsauca patriotisku jūtu uzplūdus. 9.~augustā to bija spiests atzīt arī pats afēras organizators. J.~Pilsudskis bija dziļi vīlies, ļoti kritiski izteicās par poļu tautu. Vienā no savām runām 1914.~gada beigās viņš pateica rūgtus vārdus, ka poļi izrādījušies ``ar sūnām apaudzis akmens, bet nevis dinamīta lādiņš.'' 1915.~gada februārī atbildot uz jautājumu, kas notiks ar Poliju pēc kara, viņš izteicās vēl skarbāk: ``Bet kas tad var notikt ar Poliju, ar šo tautu, kura nav spējīga ne uz kādu darbību, kura ir Eiropas apkaunojums un trīsdesmit miljonu liels traips uz cilvēces ķermeņa''.

J.~Pilsudska un viņa piekritēju darbību vēlāk asi kritizēja viņa politiskais pretinieks nacionāldemokrātu līderis R.~Dmovskis.

1925.~gadā viņš sarkastiski atzīmēja, ka J.~Pilsudska komandētā poļu leģiona~--- it kā Nacionālās valdības Varšavā spēku~--- karagājiens no Krakovas Polijas karalistes iekšienē vēl nebija nekāda traģēdija, bet tā būtu traģēdija, ja Polijas karalistē tik tiešām sāktos plaša sacelšanos, kuru izmantojot, Vācija varētu atbrīvoties no karošanas Austrumu frontē, pārsviest savus spēkus uz Rietumiem un īsā laikā sakaut Franciju. Tad poļi būtu ``atbrīvojuši Eiropu no ilgstoša kara, bet sevi no rūpēm, saistītām ar savas valsts atdzimšanu''. Tātad pēc R.~Dmovska domām poļu sacelšanās nozīmētu Vācijas uzvaru karā. Acīmredzot arī R.~Dmovskis pārvērtēja Polijas faktora nozīmi karā, bet nenovērtēja Krievijas pretošanās spējas. Taču var piekrist viņam, ka poļu sacelšanās, ar kuru J.~Pilsudskis saistīja savas cerības, apgrūtinātu Antantes valstu stāvokli. Modelējot tālāk situāciju, var droši teikt, ka tā neļautu Krievijai 1914.~gada augustā izvērst uzbrukumu Vācijai Austrumprūsijā, līdz ar to nevar pateikt, kā beigtos kauja pie Marnas [Tā notika septembra sākumā pie Marnas upes starp Vācijas un Francijas un Lielbritānijas spēkiem un beidzās ar vācu atkāpšanos, kā rezultātā netika īstenots Vācijas plāns izcīnīt ātru uzvaru Rietumu frontē un piespiest kapitulēt Franciju~--- V.Š.]. Nevar pat izslēgt iespēju, ka Centrālās valstis karā uzvarētu, savstarpēji sadalītu poļu apdzīvotās teritorijas un nostiprinātu tur savu varu. Ticamāk gan, ka Centrālās valstis tā arī nespētu uzvarēt karā, taču Antantei uzvara nāktu ar lielākām pūlēm un nevar pateikt, kādas tad būtu Polijas izredzes uz neatkarību, pat ja J.~Pilsudskis ar poļu leģioniem kara beigu posmā pārietu Antantes pusē.

Pēc aprakstītajiem notikumiem Austroungārijas valdība J~Pilsudskim uzstādīja ultimātu: viņam savas vienības bija vai nu jāatlaiž, vai jāiekļauj Austroungārijas armijā. Komandants izvēlējās otro iespēju.

Lai nepieļautu izformēt poļu militarizētās vienības un ietekmētu Galīcijas nacionāldemokrātus (endekus), Krakovas konservatīvie un tiem tuvie elementi nolēma organizēt visu politisko partiju un kustību pārstāvniecību un ar Austroungārijas valdības atļauju 1914.~gada 16.~augustā (pēc jaunā stila) Krakovā tika radīta poļu \strong{``Augstākā nacionālā komiteja''} (\pltxti{Naczelny Komitet Narodowy}), kurā bija pārstāvētas gandrīz visas partijas un kurai Berlīnes un Vīnes aizbildniecībā bija jāvada pret Krieviju vērsta nacionāla kustība. J.~Pilsudskis sakarā ar to piekrita atlaist fiktīvo ``valdību'' Varšavā. 1914.~gada 16.~augustā tika arī izsludināta uzņemšana Austroungārijas armijas \strong{poļu leģionos}. Bija paredzēts izveidot divus~--- pa vienam Rietumu un Austrumu Galīcijā. Iniciatīvu parādīja tā pati proaustriešu orientācijas poļu politiķu grupa. J.~Pilsudskis bija spiests piekrist. Leģionos iestājās galvenokārt viņa piekritēji. Leģionāriem un pašam J.~Pilsudskim nācās zvērēt uzticību Habsburgu monarhijai. J.~Pilsudska uzticības zvērests imperatoram Francim Jozefam tikai apstiprināja viņam piemītošo pilnīgo politisko principu trūkumu.

1915.--1916.~gadā leģionos dienēja ap 15~tūkstošiem karavīru. Leģionāri nēsāja poļu atšķirības zīmes, pavēles tika dotas poļu valodā. Leģionus pakļāva poļu izcelsmes Austroungārijas ģenerāļiem. Līdz ar to leģionu dalībnieki ieguva kombatantu statusu (no franču \frtxti{combattant}~--- karavīrs, cīnītājs. Jēdzienu noformulēja otrā Hāgas konference Konvencijā par sauszemes karu likumiem 1907.~gadā. Par kombatantiem sauca brīvprātīgo vienību dalībniekus, kuru priekšgalā atradās persona, kas nesa atbildību par savu padoto darbību, kuri atklāti nēsāja ieročus, kā arī noteiktas un no tālienes labi saskatāmas atšķirības zīmes un ievēroja kara vešanas likumus.) Tas bija svarīgi, jo krievu karaspēks, rīkojoties pilnā atbilstībā starptautiski atzītajiem kara likumiem, 13.~augustā nošāva piecus Kelces apkārtnē ar ieročiem rokās sagūstītus poļu strēlniekus kā vienkāršus laupītājus.

Pats J.~Pilsudskis tika iecelts par Rietumu (Krakovas) poļu leģiona 1.~pulka, vēlāk 1.~brigādes komandieri jeb, kā toreiz teica, brigadieri (kas toreiz atbilda jaunākā ģenerāļa statusam), taču ne par visu poļu spēku virspavēlnieku, kā viņš to vēlējās. Tā strēlnieku pārformēšana leģionos nesa J.~Pilsudskim militāru pazeminājumu no visu poļu spēku komandanta posteņa līdz vienas vienības komandiera amatam. Iecerētais Austrumu (Lembergas) poļu leģions drīz pēc izveidošanas tika atlaists, jo tā karavīri pēc Krievijas uzvaras Galīcijas kaujas rezultātā (1914.~gada augustā~--- septembrī) atteicās karot pret Krieviju. Rietumu leģionā toties izveidoja trīs brigādes. Savā vienībā J.~Pilsudskis saglabāja strēlnieku naģenes~--- t.s. macejevkas (\pltxti{maciejówka}), kaut pārējās nēsāja atšķirīgas cepures~--- t.s. konfederatkas \pltxti{(rogatywka}~--- cepure ar četriem stūriem), kuras t.s. Bāras konfederācijas (1768--1772) laikā bija iecienītākās konfederātu galvassegas.

Pats J.~Pilsudskis no 1914.~gada augusta līdz 1916.~gada rudenim nepārtraukti intriģēja pret leģiona pavēlniecību, Augstākās nacionālās komitejas kara departamentu, kuru vadīja V.~Sikorskis. J.~Pilsudskis centās panākt arī pēc iespējas aktīvāku poļu dalību kaujās pret Krieviju, lai atdzimtu leģenda par bruņoto cīņu par Polijas brīvību. 1.~brigāde karoja 1914.~gada decembrī pie Lovičas (poļu \pltxti{Łowicz}, vācu \detxti{Lowitsch}), 1915.~gadā pie Konariem (\pltxti{Konary}). Tika radīta sava brigādes himna~--- dziesma ``\pltxti{My, piersza brigada}'' (``Mēs, pirmā brigāde''), pazīstama arī kā \pltxti{``Marsz Pierwszej Brygady}'' (``Pirmās brigādes maršs''), kura mūsdienās ir Polijas bruņoto spēku oficiāla himna. Samērā sekmīgi cīnījās 2.~brigāde, piemēram, veiksmīgs bija poļu ulānu uzbrukums pie Rokitnas (\pltxti{Rokitna}) 1915.~gada jūnijā, tomēr lielu panākumu nebija. Kā atzīmējušas krievu vēsturnieces I.~Jažborovska un V.~Parsadanova, ``visas J.~Pilsudska piekritēju leģendām apvītās kaujas beidzās ar sakāvēm, kas sāpīgi dzēla ``komandanta'' patmīlību''. Ar pamatīgu, lieliem zaudējumiem pavadītu austroungāru karaspēka, kura pusē cīnījās arī poļu leģionāri, sakāvi noslēdzās arī viņu ievērojamākā kauja pie Kostjuhnovkas (\rutxti{Костюхновка}) Volīnijā 1916.~gada jūnijā.

Tomēr līdz pat 1917.~gadam poļu leģioni dūšīgi karoja Centrālo valstu pusē. Paša J.~Pilsudska nopelnus Habsburgu monarhija novērtēja, 1915.~gadā apbalvojot viņu ar Dzelzs kroņa ordeni \detxti{(Orden der Eisernen Krone}). Vēloties karot par Polijas neatkarību, leģionāri objektīvi aizstāvēja Austroungārijas intereses. Leģionāri gan nevarēja īsti saprast, kāpēc viņiem jāzvēr uzticība Austroungārijas imperatoram, kā tas nāks par labu Polijai. Tomēr viņi vadījās no devīzes: ``Komandants zina, ko dara'', kura it kā atrisināja visas problēmas. Bezgalīga uzticība J.~Pilsudskim lielā mērā kompensēja skaidras politiskas platformas trūkumu. Komandanta personība personificēja cīņu par Polijas neatkarību. Veidojās noturīga J.~Pilsudska piekritēju grupa. (No 1918.~gada tā ieguva pilsudčiku nosaukumu.)

Poļu zinātnieks M.~Gertihs, izanalizējis karojošo pušu armiju skaitlisko sastāvu un stāvokli, ir atzinis, ka J.~Pilsudska organizētajiem leģioniem nebija nozīmes karā. Taču atzīmējama arī poļu vēsturnieka S.~Keņeviča doma, ka poļu leģionu potenciālam ``bez šaubām, nebija nekādas lomas pasaules kara liktenī, taču leģionu leģenda sacēla kājās veselu paaudzi''. Arī Otrā pasaules kara laikā tā kalpoja kā patriotisma piemērs poļu karavīriem.

Mērķtiecīgi tika veidota paša J.~Pilsudska un viņa komandētās 1.~brigādes leģenda. To veicināja gan paša brigadiera kā sociālistiska aģitatora un propagandista pieredze, gan tas, ka viņa vadībā dienēja daudz izglītotu, radošu cilvēku, kuri labi saprata drukātā vārda iedarbību sabiedriskās domas veidošanā un plaši to pielietoja. Gan kara laikā, gan pēc tā viņu daiļradē J.~Pilsudskis tika tēlots kā 1.~brigādes veidotājs un vadītājs, kurš gandrīz vai vienatnē izlauza savai tautai ceļu uz neatkarību. 2. un 3.~brigādes kauju ceļš faktiski neatšķīrās no 1.~brigādes veikuma, taču ne tās, ne to komandieri nekad neguva tik plašu slavu. Var piebilst, ka J.~Pilsudskis pastiprināti gatavoja sevi nākamās Polijas militāro spēku vadītāja amatam, studēja Napoleona karu pieredzi, vācu militāristu K.~Klauzevica un H.~Moltkes darbus, angļu-būru kara, poļu nacionālās sacelšanos 1863.--1864.~gadā pieredzi. Protams, tas nevarēja aizstāt profesionālu militāro sagatavotību, taču deva zināmās zināšanas, lai viņš nejustos pilnīgs profāns kara lietās.

1916.~gadā leģiona rindās bija ap 25~000 karavīru. (Pēc citām ziņām, leģionā iestājušies 15--20 tūkstošu, vēl pēc citām~--- caur leģioniem izgājuši ap 30~ tūkstoši cilvēku. Tiesa, publicistikā atrodami arī citi, acīmredzot pazemināti skaitļi, ka J.~Pilsudska vadītajā leģionā Austroungārijas pusē pret Krieviju reāli karojuši tikai ap 3~000 karavīru.) Leģionu sastāvā cīnījās arī trīs nākamie Polijas maršali~--- pats J.~Pilsudskis, E.~Ridz-Smiglijs un M.~Žimerskis (pseidonīms \pltxti{Rola}).

Var arī atzīmēt, ka Austroungārijas armijā pastāvēja arī t.s. Ukrainas Sečas strēlnieku (ukraiņu \uktxti{Українські Січові Стрільці})~--- austriešu orientācijas ukraiņu-galīciešu vienības, kuras tāpat kā poļu leģioni karoja pret Krieviju, cerot uz Ukrainas statusa izmaiņu. Vēlāk~--- 1918./1919.~gadā poļu leģionāri un ukraiņu Sečas strēlnieki karoja savā starpā.

1914.~gada 4.~novembrī Austroungārijas imperators Francis Jozefs I parādīja labdarības žestu poļiem: pasludināja Galīcijas autonomiju un noteica, ka tā ``var kārtot savas lietas piederības ietvaros kopīgajai valstij''. Taču Galīcijas nacionāldemokrāti centās nepieļaut cīņu pret Krieviju zem poļu karogiem un drīz izstājās no ``Augstākās nacionālās komitejas''. Arī J.~Pilsudskis nebija apmierināts, ka austrieši ierobežo viņa patstāvību, taču viņa mēģinājumi tuvināties Vācijai (atrast citu ``zirgu'' poļu ``mušai'') cieta neveiksmi. J.~Pilsudskim nebija tiešas ietekmes arī uz ``Augstāko nacionālo komiteju'', taču viņa pūliņi veicināja poļu nacionālās politiskās darbības attīstību vācu un austriešu ieņemtajās Polijas karalistes teritorijās. Tur J.~Pilsudska vadībā tika veidota nelegāla Poļu kara organizācija (\pltxti{Polska Organizacja Wojskowa}), kuras locekļi piedalījās arī izlūkošanas un diversiju operācijās pret Krieviju. Viņš cerēja, ka šīs poļu iniciatīvas karā kļūs par trumpi nākamajās sarunās ar Vācijas un Austroungārijas valdībām.

Kamēr ``aktīvistu'' galvenā uzmanība bija pievērsta Austroungārijai un notikumiem Galīcijā, Polijas karalistē joprojām stingras palika ``pasīvistu'' pozīcijas.

Kara sākumā \strong{Krievijas iestādes}, kā nu prata, \strong{centās nostiprināt savu ietekmi Polijas karalistē}. Tās guberņās uzreiz tika izsludināts kara stāvoklis, ierobežotas jau tā šaurās pilsoniskās brīvības. Varšavas ģenerālgubernators un apgabalā izvietotā Krievijas karaspēka pavēlnieks ģenerālis J.~Žilinskis aizliedza noturēt sapulces, gājienus, manifestācijas, pārdot uz ielām avīzes, žurnālus un grāmatas bez speciālām atļaujām, tika ieviesta administratīva to tieslietu, sods par kurām nepārsniedza 3 mēnešu ieslodzījumu cietumā vai naudas sodu līdz 3~000 rbļ., izskatīšanas kārtība. Kritās papīra naudas vērtība. Tirgotāji atteicās pieņemt kā maksāšanas līdzekli Krievijas papīrnaudu, pieprasot zelta, sudraba naudu un sīkmonētas. Varšavas kara gubernators izdeva rīkojumu vainīgos par atteikšanos pieņemt papīra naudu ne pēc tās vērtības arestēt un sodīt, ``nāves sodu ieskaitot''.

Visu karojošo valstu valdību priekšā stāvēja problēma, ko darīt ar tām naidīgo valstu pilsoņiem un to mantu. Visur šo ārvalstu pilsoņu tiesības tika ierobežotas, viņi tika deportēti vai internēti. Taču nevienā valstī šī vēršanās pret citu valstu pavalstniekiem nebija tik asa kā Krievijā, īpaši tās rietumu nomalēs.

Kā norādījis piespiedu migrāciju pētnieks PSRS teritorijā P.~Poļans, pasaules vēsture pazīst milzumu ``deportāciju'' vai ``piespiedu migrāciju'' piemēru. Lai atceramies tikai Veco derību, kurā daudz epizožu veltīti ebreju ``specizsūtīto'' liktenim Ēģiptē, Babilonijā u.c. Aiz visām tām stāvēja tāds vai citāds politisko un pragmatisko faktoru apvienojums. Politiskie motīvi~--- novērst iespējamās sacelšanās, vājināt vai apspiest protestus, padarīt viendabīgākus izsūtījuma vai ieceļošanas rajonus parasti bija prioritāri, taču arī ekonomiskā faktora loma ir nozīmīga, ar laiku tā (lēta, gandrīz bezmaksas darbaspēka, pārsviesta pēc deportāciju veicēja gribas uz tam vajadzīgo vietu, iegūšana) varēja pat aizēnot politisko impulsu. Tā, no XVI līdz pat XIX gadsimta vidum uz Ameriku tika deportēti ap 11 līdz 15 miljonu afrikāņu, kuri, nonākot jaunajā mītnes zemē, kļuva par vergiem. Sevišķi bagāta piespiedu migrācijām ir ebreju tautas vēsture.

Arī Krievijas vēsturē vēl līdz Oktobra revolūcijai ir atrodami daudzi piespiedu izsūtīšanas piemēri. Tā vēl XIX~gadsimta sākumā krievu-turku karu laikā no Melnās jūras piekrastes uz Krimu tika izsūtīti visi musulmaņi (tatāri, tjurki, nogajieši), pēc Krievijas uzvaras Kaukāza karā (1817--1864) notika ap 400~tūkstošu Rietumkaukāza kalniešu vardarbīga pārvietošana uz Turciju. 1891.--1892.~gadā no Maskavas tika izsūtīti līdz 38~tūkstošiem ebreju. Taču īpaši plašas deportācijas Krievijas impērijā saistījās ar Pirmo pasaules karu.

Tieši Krievijā ārzemju kapitāls, pirmkārt~--- vācu, spēlēja ļoti lielu lomu tautsaimniecībā, sevišķi Krievijas rietumu apgabalos. Turklāt Krievijā dzīvoja daudz vācu kolonistu, bet Vācijas varas iestādes XX~gs. sākumā atklāti demonstrēja, ka uzskata vācu diasporu par vācu nācijas daļu, bez izšķirības, kādas valsts pavalstnieki vācieši bija. Daudz šādu cilvēku dzīvoja arī Polijā. Protams, Krievijas varas iestādes neuzticējās vietējiem vāciešiem. No 1914.~gada septembra no poļu guberņām sākās aizdomās turēto par spiegošanu izsūtīšana uz Krievijas iekšējām guberņām, pat tad, ja aizdomām nekādu pierādījumu nebija. Izsūtīšana tika oficiāli apstiprināta ar Varšavas ģenerālgubernatora 1914.~gada 7.~septembra rīkojumu. 1915.~gada 2.~februārī tika izdots likums par vācu koloniju likvidēšanu Polijas karalistē un Volīnijā. Šo provinču rietumu daļā gan drīz ielauzās Vācijas armija un vācu iedzīvotāji tika paglābti, bet ap 120~000 vāciešu no Polijas karalistes austrumiem un Volīnijas tomēr tika izsūtīti uz Krievijas iekšieni un Sibīriju. (Tā gan nebija pirmā masu izsūtīšana (deportācija) XX gadsimtā, jo pirms tam pēc Otrā Balkānu kara 1913.~gada jūnijā-jūlijā notika nacionālo mazākumu pārvietošana (apmaiņa) vienošanas rezultātā starp Bulgāriju un Turciju.) Kopumā Krievijā izsūtīšanai Pirmā pasaules kara laikā bija pakļauti ap 200~--- 330 tūkstošu vāciešu un austriešu. Izsūtīti tika arī ebreji. 1914.--1916.~gadā no Polijas, Lietuvas, Baltkrievijas un Kurzemes guberņas it kā to nelojalitātes dēļ pret Krieviju tika izsūtīti ap 250--350~tūkstošu ebreju. Tas notika neraugoties uz to, ka gandrīz katrā ebreju ģimenē kāds bija iesaukts armijā un ebreju jauniešus tajā turpināja iesaukt arī no izsūtīto rindām. (Ir arī dati, ka deportācijas Krievijas rietumos skāra ap 1~miljonu cilvēku, puse no kuriem bija vācieši, trešā daļa~--- ebreji.)

Krievijas armijas virspavēlnieks lielkņazs Nikolajs Nikolajevičs tūlīt kara sākumā~--- 1.~(14.)~augustā~--- izdeva manifestu, kuru laikabiedri pierakstīja viņa personīgajai iniciatīvai.

Taču, kā raksta vispusīgi manifesta tapšanas vēsturi pētījusī krievu vēsturniece A.~Bahturina, aiz tā faktiski stāvēja pats Nikolajs~II un tā tapšanā bija piedalījušies Krievijas ārlietu ministrs S.~Sazonovs, kara ministrs V.~Suhomļinovs, Ģenerālštāba priekšnieks ģenerālis N.~Januškevičs. Manifests tika apspriests valdībā un pieņemts ar balsu vairākumu, diviem ministriem atturoties. Manifesta publicēšanas ideju atbalstīja arī imperators. Kāpēc viņš pats arī neparakstīja dokumentu, savās atmiņās stāstīja Francijas sūtnis Krievijā M.~Paleologs. Ārlietu ministrs S.~Sazonovs viņam teicis, ka konservatīvie Krievijas valdības locekļi pārliecinājuši Nikolaju~II par to, ka, kamēr Galīcijas un Pozenes poļi vēl atrodas austriešu un prūšu pakļautībā, Krievijas imperatora uzruna tiem esot priekšlaicīga. Toties lielkņazs Nikolajs Nikolajevičs to varot darīt kā tās armijas virspavēlnieks, kura ieiet poļu apdzīvotajās teritorijās. Lielkņazam Nikolajam Nikolajevičam palūdza to parakstīt nevis ar vārdu ``Nikolajs'', bet poļiem pierastākajā formā ``Mikolajs'', ko viņš arī izdarīja. Pēc dažām ziņām, līdzdalību manifesta apspriešanā bija ņēmis arī P.~Dmovskis, kurš rādījis Krievijas ārlietu ministram S.~Sazonovam vēlamo Polijas robežu rietumos (to P.~Dmovskis vēlāk aizstāvēja arī Parīzes miera konferencē). Tas, ka manifests nebija nedz paša cara, nedz tā valdības parakstīts, varēja atvieglot Krievijas manevrus nākotnē. Iespējams, oficiālās Krievijas aprindas vēlāk varēja to arī neatzīt kā sev saistošu.

Dokumentā gāja runa par poļu zemju apvienošanu, ``kuras būtu brīvas savā ticībā, valodā un pašvaldībā''. Tātad manifests solīja faktiski vienotai Polijai autonomiju Krievijas sastāvā.

Jāuzsver, ka idejai par visu poļu zemju apvienošanu pirms kara bija daudz pretinieku gan augstāko krievu birokrātijas, gan militāristu vidū. Viņi ar šaubām izturējās pret ideju, ka Austrumgalīcijas iedzīvotāji ir vardarbīgi atrauti Krievijai un Galīciju nepieciešams pievienot Krievijai. Vēl 1900.~gadā toreizējais Krievijas kara ministrs ģenerālis A.~Kuropatkins ziņojumā Nikolajam~II izteicās pret Krievijas rietumu robežu ``uzlabošanu'' un Austrumprūsijas un Galīcijas pievienošanu Krievijai, jo tās varētu kļūt par sava veida ``Austrumeiropas Elzasu-Lotringu'' (vācu \detxti{Elsaß-Lothringen}, franču \frtxti{Alsace-Lorraine}~--- pēc franču-vācu kara 1871.~gadā abas provinces tika pievienotas Vācijai, taču Francijā tāpēc radās iespaidīga revanšistiska kustība. Taču kara sākumā arī šie skeptiķi padevās nacionālisma uzplūdiem. Arī A.~Kuropatkins ziņojumā ``Krievijas robežas 1914.--1915.~gada kara rezultātā'' atzina savu 1900.~gadā izteikto viedokli par nepareizu un ieteica uzstādīt sabiedrotajiem tādus noteikumus, kuri ``var nodrošināt Krievijas intereses XX gadsimtā'', paredzot pievienot Krievijai gan Austrumprūsiju, gan citas Vācijas un Austroungārijas teritorijas ar poļu iedzīvotājiem.

Gan pats dokuments, gan vēlākā Krievijas Ārlietu ministrijas darbība liecināja, ka savu jauno statusu Polijai bija jāsaņem no Krievijas rokām, nevis starptautiskas vienošanās ceļā. Krievijas sabiedrotās valstis to tā arī saprata. Nākamajai miera konferencei vajadzētu tikai apstiprināt jau noteiktās Polijas robežas. Šis manifests tajā brīdī piedāvāja Polijai visvairāk salīdzinājumā ar pretinieku valstu solījumiem un tās to nevarēja pārspēt, jo nedz Vācija, nedz Austroungārija neuzdrošinājās solīt Polijas apvienošanu, tā aizskarot sava sabiedrotā intereses. Aprēķini, ka karš drīz uzvaroši beigsies, lika Krievijas valdībai rīkoties tieši tā: pēc iespējas ātrāk noteikt vēlamos teritoriālos ieguvumus. Bez tam Krievijas valstsvīri centās apsteigt Vāciju un Austroungāriju ar savu deklarāciju ar attieksmi pret slāvu tautām. Krievijas valdošā slāņa priekšā stāvēja arī uzdevums saglabāt poļu iedzīvotāju lojalitāti un stabilu politisko situāciju kara darbības joslā. Šis mērķis tika sasniegts. Nikolaja Nikolajeviča manifesti poļiem un citām Austroungārijas slāvu tautām formulēja Krievijas kara ideoloģiju tā sākuma posmā, nosakot kā mērķi cīņu par ``mazo tautu'' atbrīvošanu un slāvu ``cilšu'' apvienošanu.

Vēl pirms dokumenta publicēšanas Krievija ar to iepazīstināja sabiedrotās valstis. Kaut pēc būtības sabiedrotie bija nostādīti fakta priekšā, ka Krievija pretendē uz Vācijai un Austroungārijai piederošajām poļu zemēm, Lielbritānija un Francija uztvēra to labvēlīgi, kara sākumposmā Polijas jautājumu uzskatot par Krievijas iekšēju lietu. Pēc manifesta publicēšanas, to pārpublicēja arī ārzemju avīzes. Ārvalstu valdības un sabiedrība to uztvēra ar lielu interesi. Uz daudziem pieprasījumiem tuvāk izskaidrot manifesta tekstu ārlietu ministrs S.~Sazonovs atbildēja, ka viss Virspavēlnieka manifestā poļiem izklāstītais ``izsaka īstenus un nopietnus valdības nodomus'', taču ``par paredzamo brīvību robežām un apjomu nevar pateikt neko noteiktu un pats jautājums prasa likumdevēju izskatīšanu''.

1914.~gada septembrī, kad Galīcijas austrumu daļa jau bija Krievijas karaspēka rokās, Krievijas ārlietu ministrs S.~Sazonovs strādāja pie plāniem, pēc kuriem Rietumgalīcijai, Silēzijai un Poznaņai būtu jāapvienojas ar Polijas karalisti zem Romanovu sceptera. S.~Sazonovs Francijas un Lielbritānijas sūtņiem Krievijā M.~Paleologam un D.~Bjukenenam izklāstīja kara rezultātā panākamā miera galvenās iezīmes. Nākotnē sasniedzamo teritoriālo izmaiņu pamatā bija jāliek nacionālais princips. Nemunas (lietuviešu \lttxti{Nemunas}, vācu \detxti{Memel}, poļu \pltxti{Niemen}) upes lejtece un Austrumgalīcija bija jāsaņem Krievijai; Poznaņa, Silēzija un Rietumgalīcija bija jāiekļauj Polijas sastāvā. Ņemot vērā S.~Sazonova vienošanās ar sabiedroto diplomātiskajiem pārstāvjiem, Krievijas valdība plānoja izlemt, pēc kādas pazīmes tiks veidotas robežas: vai tās noteiks tikai etnogrāfiskā pazīme vai arī tiks ņemtas vērā arī vēsturiskās Polijas valstiskās tradīcijas. 1914.~gada novembrī audiencē Francijas sūtnim M.~Paleologam arī imperators Nikolajs~II runāja par kara mērķiem. Viņš paskaidroja, ka Poznaņa un daļa Silēzijas būs vajadzīga ``Polijas atjaunošanai'' bet Galīcijas un Bukovinas (ukraiņu \uktxti{Буковина}, rumāņu \rotxti{Bucovina}) ziemeļu daļas (ukraiņu \uktxti{Північна Буковина}, rumāņu \rotxti{Bucovina de Nord}, mūsdienās daļa no Černivcu apgabala~--- \uktxti{Чернівецька область}~--- Ukrainā) iegūšana ļaus Krievijai sasniegt savas ``dabīgās robežas''.

Ārlietu ministram S.~Sazonovam izdevās saņemt sabiedroto valstu piekrišanu, ka līdz karam Austroungārijai un Vācijai piederējušie pārsvarā poļu apdzīvotie rajoni~--- Rietumgalīcija, Silēzija, daļa Prūsijas un Lielpolija tiks apvienoti ar Polijas karalisti vienotā Polijas valstī, kura atrastos personālas ūnijas attiecībās ar Krievijas impēriju.

Poļu sabiedrībā lielkņaza Nikolaja Nikolajeviča manifests tika uztverts ar entuziasmu. Tā, Lodzā notika manifestācija, kurā piedalījās ap 20~000 cilvēku. Formāli poļu sabiedriski-politisko organizāciju vairākums izteica savu lojalitāti cara režīmam. Poļu politisko partiju vārdā lielkņazam Nikolajam Nikolajevičam tika nosūtīts uzticību sološs vēstījums, kurā bija teikts, ka ``poļi karsti gaida krievu ieroču uzvaru \citespace{} un tās pilnīgu triumfu kaujas laukā \citespace{} Ticam, ka Polijas dēlu asinis, izlietas kopā ar Krievijas dēliem cīņā pret kopīgo ienaidnieku kļūs par labāko ķīlu jaunai dzīvei pasaule un divu slāvu tautu draudzībai.'' Vēstījumā gan nekas nebija teikts par Polijas vienotības pasludināšanu zem Krievijas patvaldnieka sceptera. Taču arī tādā veidā Krievijas varas iestādes pacentās to maksimāli izmantot. Tā, Varšavas ģenerālgubernatora kanceleja saņēma 50~000 vēstījuma eksemplāru izdalīšanai iedzīvotājiem. Tie tika izsūtīti uz guberņām. Polijas karalistes iedzīvotāji kopumā saglabāja lojalitāti Krievijai līdz pat brīdim, kad 1915.~gadā Vācijas karaspēks ieņēma Varšavu.

Tomēr Krievijas neskaidrie solījumi atrisināt Polijas jautājumu veda pie poļu nacionālā separātisma pieauguma un stiprināja viņu nacionālo pašapziņu. Lielkņaza Nikolaja Nikolajeviča manifests atbilda R.~Dmovska vadīto ``endeku'' nostājai. Uz Nikolaja Nikolajeviča manifestu viņi atbildēja ar lojalitātes apliecinājumiem Krievijas Valsts domē. Poļu pārstāvji Domē demonstrēja sevi kā ``atbrīvošanas'' kara piekritēji, atbalstīja ``slāvu pasauli'' ar Krieviju priekšgalā pret ``ģermāņu pasauli''.

Poļu pārstāvis Krievijas Valsts padomē A.~Meištovičs 1914. gada 8.~augustā uzsvēra, ka poļi vienmēr ir cīnījušies pret ``\detxti{Drang nach Osten}'' (``Spiediens uz austrumiem''), kā apliecinājumu minot Grīnvaldes kauju (poļu \pltxti{Bitwa pod Grunwaldem}, vācu \detxti{Schlacht bei Tannenberg}, 1410). Viņš atzina, ka ir neizbēgami lieli upuri, bet ``poļu asinis tiks izlietas krietnas lietas labā''. Pie tam viņš ne vārda neteica par poļu jautājuma atrisināšanu. Šī poļu politiķu grupa joprojām uzskatīja Vāciju par galveno Polijas ienaidnieku, vēlēja Antantei uzvaru un darbojās tās pusē. Viņu mērķis bija poļu zemju apvienošana Krievijas kontrolē un autonomijas statusa iegūšana. Viņi bija gatavi vervēt brīvprātīgos, kuri karotu Krievijas pusē. (Pirmajam Pasaules karam sākoties Krievijas armijā dienēja vairāki simti tūkstošu poļu, 1917.~gada februārī~--- jau 700~000, tai skaitā 119~ģenerāļu un ap 20~tūkstošu virsnieku. Pēc dažiem vērtējumiem 1914.--1917.~gadā dažādās karaspēka daļās Krievijā dienējis ap miljonu poļu.)

Tika cerēts Krievijas pusē radīt 200 līdz 300~tūkstošu lielu poļu armiju. R.~Dmovskis domāja, ka pēc kara Krievija nespēs apvienoto Poliju ``sagremot'' un agri vai vēlu piešķirs tai autonomiju. Kara sākumposmā lielākā Polijas karalistes pilsonības daļa ieņēma ``rusofīlu'' pozīciju, mazākā bija Vācijas un Austroungārijas pusē. Ne vieni, ne otri nevēlējās saskatīt, ka vienas vai otras puses bezierunu uzvara Pirmajā pasaules karā nestu Polijai jaunu pārdali, bet ne nacionālo neatkarību.

``Reālisti'' un endeki ar varas iestāžu atbalstu sāka veidot pilsoņu komitejas, kurās saskatīja nākamās pašvaldības dīgļus. 1914.~gada novembrī Varšavā izveidojās \strong{Poļu nacionālā komiteja}. Tajā ietilpa gan ``reālists'' kņazs Z.~Ļubomirskis, endeki R.~Dmovskis, V.~Grabskis, gan bijušais PPS loceklis S.~Voicehovskis u.c. Komiteja publicēja uzsaukumu poļu tautai, kurā bija teikts: ``Tautas priekšā ir nostājies tikai viens uzdevums: sagraut Vācijas draudīgo varenību un apvienot Poliju zem krievu monarha sceptera''. Poļu emigrantu vidū Rietumeiropā Lozannā tika radīts ``Preses birojs'', kurš uzturēja sakarus ar Varšavā, pēc tam Petrogradā pastāvošo Poļu nacionālo komiteju. Taču Krievijas varas iestādes pret komiteju, tās plāniem veidot poļu armijas daļas izturējās ar aizdomām, domājot, ka tās mērķis ir izvirzīt arvien jaunas prasības poļu jautājumā.

Realitātē cara valdība ne tikai nesāka realizēt solītās pašvaldības tiesības, bet nesteidzās arī likvidēt poļu tautas tiesiskos ierobežojumus. Jau manifesta sagatavošana un pieņemšana parādīja, ka Krievijas sabiedrībā pret to ir dažādi noskaņojumi. Liberāļi cītīgi atgādināja, ka Polijas sadalīšanas iniciatore XVIII~gadsimtā bija nevis Krievija, bet Prūsija, uzsvēra, ka Krievija tagad uzstājas par demokrātiskas Polijas atjaunošanu tās etniskajās robežās~--- bez Austrumgalīcijas, Baltkrievijas un Lietuvas. Krievijas diplomāti un militāristi centās nostādīt nākamo miera konferenci jau gatava jautājuma risinājuma priekšā un panākt poļu apdzīvotajos apgabalos mieru un drošību. Turpretī birokrātijas vidē skaidri bija redzama vēlme nezaudēt ``gardo kumosu''. Varšavas ģenerālgubernatora J.~Žilinska palīgi administratīvajās un policijas lietās A.~Essens un Ļ.~Utgofs Virspavēlnieka manifestu uztvēra ar nepatiku, kā ``pārpratumu''. Varšavā tika aizliegts tā izdošanas sakarā izkārt nacionālos balt-sarkanos karogus un dziedāt himnu~--- Dombrovska mazurku. Krievijas iekšlietu ministrs N.~Maklakovs sāka apgalvot, ka manifestā iet runa tikai par ``austriešu'' un ``vācu'' Poliju, bet ``krievu'' daļa tas neskar. Tā Krievijā vieni, saprotot, ka atstāt visu pa vecam nav iespējams, centās kādi risināt problēmu, citi nikni pretojās pat no augšas nākušajai iniciatīvai, saskatot manifestā ``valsts pamatu graušanu'' un tā parādot milzu inertumu.

Krievijas varas iestādes kara laikā vairākkārt apsprieda Polijas pašnoteikšanās jautājumu, taču konservatīvie, kuri nevēlējās piekāpties poļiem, panāca jautājuma izlemšanas atlikšanu līdz kara beigām, kuras, kā zināms, Krievijas impērija tā arī nesagaidīja.

Tā, 1914.~gada augusta beigās (citi dati~--- 2.~oktobrī) Krievijas Virspavēlniecības štāba priekšnieks ģenerālis N.~Januškevičs griezās ar vēstuli pie Ministru padomes priekšsēdētāja I.~Goremikina, kurā sakarā ar krievu karaspēka ieiešanu Galīcijā uzstāja, lai Ministru padome noteiktu savu viedokli Polijas un Ukrainas jautājumos, konkretizētu 1.~(14.)~augusta manifesta solījumus. Piecās sēdēs oktobrī un novembrī Ministru padome debatēja jautājumu par Poliju, taču netika tālāk par ierobežotas pašvaldības ar Seimu priekšgalā programmu. Pirmajā sēdē 20.~oktobrī vairākums tās dalībnieku nonāca pie secinājuma, ka iespējamo Polijas karalistes pārvaldes reformu gaitā ir jānodrošina tās nesaraujama saistība ar Krievijas impēriju. Virkne augstāko ierēdņu brīdināja imperatoru, ka pat ierobežotu reformu programmas pieņemšana nozīmētu valsts sairšanas sākumu. Konservatīvie krievu politiķi, tādi kā jau minētais iekšlietu ministrs N.~Maklakovs u.c. par galveno uzskatīja Krievijas impērijas nostiprināšanu, pievienojot tai jaunas teritorijas: Austroungārijai piederošo Austrumgalīciju, Aizkarpatu Ukrainu un Ziemeļbukovīnu. Turpretī par Pozeni un citām Vācijai piederošajām vācu zemēm neesot vērts karot, bet labāk atstāt slāvu zemes Vācijas sastāvā, lai to novājinātu.

5.--12.~novembrī Ministru padome sprieda par līdzekļiem, kā nodrošināt Krievijas impērijas un Polijas karalistes vienotību, vietējo pašvaldību Polijā un tās robežām. Sēdes dalībnieki nonāca pie secinājuma, ka robežām jātiek noteiktām pēc etnogrāfiskās pazīmes. Tika uzsvērts, ka Krievijas uzdevums karā ir nevis Polijas atjaunošana tās vēsturiskajās robežās, bet gan to poļu atbrīvošana, kuri dzīvo Vācijas un Austroungārijas varā. Attiecīgi, pēc Krievijas Ministru padomes uzskatiem, nākamajā Polijā bija jāietilpst Polijas karalistes guberņām, daļai Galīcijas, kur pārsvarā dzīvoja poļu iedzīvotāji (Galīcijas rietumu daļai) un Poznaņas apgabalam. Polijas pārvaldi bija jābūvē uz ticības brīvības pamatiem. Tika atzīta iespēja paplašināt valodas autonomiju, atzīstot nepieciešamību administratīvajos un tiesu amatos ierēdņiem zināt poļu valodu, reizē saglabājot iespēju vietējiem krievu iedzīvotājiem griezties pie novada iestādēm krievu valodā. Visās mācību iestādēs bija jāievieš poļu mācību valoda. Novada pārvaldei bija paredzēts iecelt imperatora vietvaldi, pie kura izveidotajā Padomē būtu jādarbojas gan valdības ieceltiem cilvēkiem, gan vietējo vēlēto iestāžu pārstāvjiem. Polijas karalistē bija paredzēts ieviest zvērināto tiesu, pilsētu un zemstu pašvaldības pēc impērijas parauga. Tā 1914.~gada nogalē Ministru padome nosprauda ierobežotu pārveidojumu plānu Polijas karalistei, kurš bija virzīts galvenokārt uz zināmas kulturāli-nacionālās autonomijas piešķiršanu, maksimāli saglabājot centrālās varas prerogatīvas tiešā novada pārvaldē. Faktiski netika pieņemti principiāli lēmumi nedz ``krievu Polijas'', nedz tās Polijas, kurai bija jārodas kara rezultātā, pārvaldes jautājumos. To noteica domstarpības pašā Ministru padomē starp liberālā un konservatīvā kursa piekritējiem. Visi cara valdības solījumi bija nosacīti, atkarīgi no Krievijas uzvaras karā, vāciešiem un austriešiem piederošo poļu zemju pievienošanas tai. 1914./1915.~gada ziemā Krievijas valdošās aprindas galvenokārt centās ``poļu jautājuma'' apspriešanu novilcināt līdz kara beigām.

% page 238


Konservatīvie pretojās gandrīz jebkurām izmaiņām kā valdības nostādnēs, tā Polijas karalistes pārvaldē. Lielkņazam Nikolajam Nikolajevičam tikai izdevās 1914.~gada 2.~decembrī atbrīvot no amatiem abus jau minētos Varšavas ģenerālgubernatora palīgus A.~Essenu un Ļ.~Utgofu. Taču poļu sabiedrība to bija gaidījusi tūlīt pēc lielkņaza manifesta publicēšanas, tāpēc šādu ierēdņu ilgstoša palikšana amatos ļāva viņiem darboties pilnīgi pretēji solītajai politikai un tas sarūgtināja uz reformām cerošos poļus. Pēc tam tika atcelti vairāki ministri, tai skaitā arī N.~Maklakovs, taču Krievijas politika mainījās gaužām lēni. Tikai 1915.~gada 3.~martā Ministru padome atļāva poļu guberņu izglītības sistēmā noturēt nodarbības poļu valodā. 30.~martā uz Polijas karalisti tika attiecināts Krievijas likums par pilsētu pašvaldībām~--- ar padomdevēju funkcijām administrācijai, likvidējot līdz tam pastāvošos pilsētu maģistrātus. Taču Krievijas valdība bija jau sākusi atkāpšanos frontē un reāli Krievijas varas iestāžu kontrolē bija jau tikai 3 no 10~poļu guberņām. Poļu sabiedriskās aprindas prasīja nopietnākas reformas un vēlējās pašas piedalīties to sagatavošanā. Poļu politiķi mēģināja nostiprināt savu tautiešu stāvokli arī ārpus poļu guberņām. Tā poļu deputāti Krievijas Valsts domē un Valsts padomē iesniedza memorandu, norādot, ka Lietuvā un daļā Baltkrievijas un Ukrainas poļu zemes īpašums sastāda vairāk nekā pusi no visas zemes platības, tāpēc sekoja lūgums~--- atcelt poļiem uzliktos ierobežojumus. Tas gan netika izpildīts.

Smagu triecienu poļu~--- Krievijas piekritēju cerībām deva Krievijas armijas sakāves Prūsijā un Polijā. Sākotnēji Krievijas armija, lai palīdzētu Francijai izturēt pirmo Vācijas triecienu, iebruka Austrumprūsijā. Vācija bija spiesta noņemt no savas Rietumu frontes tur tik nepieciešamos spēkus un pārsviest uz austrumiem. Divas Krievijas armijas tika sagrautas, taču Vācijas plāni Francijā neīstenojās un karš ievilkās. Polijas karalistes teritorijā vācu armijas nonāca gandrīz līdz Varšavai. Tiesa, 1914.~gada oktobrī un novembrī vācieši arī šeit bija spiesti atkāpties, taču Polijas karalistes teritorija bija kļuvusi par kara darbības zonu.

1915.~gadā Vācijas karaspēks atkal enerģiski centās ielenkt un iznīcināt krievu spēkus Polijā. Tas gan neizdevās, bet Krievijas armija 1915.~gada jūlijā sāka atkāpšanos plašā frontē, atstājot Varšavu, Viļņu, Kauņu, Brest-Ļitovsku. Fronti izdevās stabilizēt tikai līnijā Rīga~--- Dvinska (tagad~--- Daugavpils)~--- Pinska (baltkrievu \rutxti{Пінск})~--- Rovno (ukraiņu \uktxti{Рівне}, poļu \pltxti{Równe})~--- Ternopole (krievu \rutxti{Тернополь}, ukraiņu \uktxti{Тернопіль}, poļu \pltxti{Tarnopol}). 1915.~gada augusta sākumā Vācijas karaspēks ieņēma Varšavu. Jau līdz tam 1915.~gada vasarā Krievijas varas iestādes uzsāka daļas Polijas karalistes rūpniecības uzņēmumu un poļu strādnieku masveida evakuāciju uz Urāliem, Pievolgu un citiem Iekškrievijas rajoniem. Faktiski Krievijas varas pielietoja ``izdedzinātās zemes'' taktiku, izveda arhīvus, virkni mācību iestāžu, izejvielu un pārtikas rezerves, mājlopus. Tika evakuēti valsts un pašvaldību iestāžu kalpotāji, arī zemnieki, tai skaitā vācu kolonisti. Mūsdienu poļu vēsturnieks K.~Zeļinskis aprēķinājis, ka pirmajā kara gadā no Pievislas guberņām tika evakuēti 1~miljons 200~tūkstošu krievu, poļu u.c. tautību dzelzceļnieku, ierēdņu, garīdznieku, skolotāju, kara rūpnīcu strādnieku. Daļa iedzīvotāju, kara darbības ietekmēti, paši devās bēgļu gaitās. Pēc citiem aprēķiniem no Polijas karalistes kara gados uz Krievijas iekšējām guberņām evakuējās apmēram 2~miljoni cilvēku, tai skaitā vairāk nekā viens miljons poļu. Viņu repatriācija uz dzimteni ievilkās līdz 1924.~gadam.

Tomēr Vācijas cerības piespiest Krieviju izstāties no kara nepiepildījās, tā ienesa ievērojamu ieguldījumu nākamajā Antantes valstu uzvarā, taču kā krievu tautā, tā arī tās poļu atbalstītāju vidū noskaņojums pasliktinājās.

Poļu politiķi, kuri orientējās uz sadarbību ar Krieviju, atstāja Poliju reizē ar krievu karaspēku un Poļu nacionālā komiteja sāka darboties Petrogradā. Pēc tam, kad Polijas karalisti okupēja Vācijas un Austroungārijas karaspēks, krievu un poļu pilsonisko partiju preses izdevumi Krievijā sāka izvirzīt cīņas par ``neatkarīgu Poliju'' lozungu. To atbalstīja arī krievu meņševiki, kuri tāpat aizstāvēja lozungu par ``mieru bez aneksijām''. Krievu boļševiki savukārt norādīja savu ceļu uz pilnīgu Polijas neatkarību un tās sasniegšanas ceļu. Norādot, ka abi minētie lozungi sastop arī cara valdības atbalstu, boļševiku vadonis V.~Ļeņins rakstīja, ka lozungs ``mieru bez aneksijām'' ``\dots{}\emph{praksē} nozīmē lozungu~--- uzvara pār Vāciju, kura atņēmusi Krievijai Poliju'', ka ``Krievijas sociāldemokrātiem jāatmasko tas, ka carisms \emph{apkrāpj} tautu, kad \emph{tagad} Krievijā tiek izvirzīti lozungi par ``mieru bez aneksijām'' un ``Polijas neatkarību'', jo abi šie lozungi tagadējā stāvoklī nozīmē tieksmi turpināt karu un attaisno tādu tieksmi.''

Drīz vadošais impērijā poļu politiķis R.~Dmovskis, sapratis, ka Krievijas valdība, nebūdama pārliecināta par kara iznākumu, neuzņemsies risināt Polijas jautājumu, uz visiem laikiem novērsās no Krievijas, 1915.~gadā devās uz Eiropu, Šveici, lai tur pie Antantes valstu pārstāvniecībām lobētu Polijas intereses. Rietumu lielvalstu aizbildnībā viņš sāka cīnīties par Polijas neatkarību.

Krievijas armijas neveiksmes 1915.~gada pavasarī lika arī Nikolajam~II un daļai cariskās birokrātijas vismaz uz laiku mainīt savu konservatīvo nostāju, meklēt sabiedrotos ``sabiedrības'' aprindās, spert jaunus soļus ``poļu jautājumā''. 1915.~gada maijā--jūnijā ar cara piekrišanu norisinājās ``krievu un poļu darbinieku apspriede''. Apspriedes laikā izvērsusies diskusija tika traktēta kā Virspavēlnieka manifestā pasludināto mērķu realizācijas kārtības pirmējā apspriešana. Bez augstākajiem Krievijas ierēdņiem apspriedē piedalījās arī Polijas karalistes guberņu pārstāvji: R.~Dmovskis, J.~Dobeckis, J.~Garusevičs, S.~Grabskis, S.~Veļepoļskis un I.~Šebeko. Visi apspriedes dalībnieki piekrita rusifikācijas politikas pārtraukšanai poļu zemēs. Taču tālāk domas dalījās. Norisa cīņa par Polijas jautājuma risināšanas virzību. Poļu dalībnieki centās Krievijas kara mērķu un nostādnes uz slāvu savienību nākotnē iezīmētajās robežās izvirzīt Polijas patstāvības ideju. Poļu teritorijas tika nosauktas par ``visas slāvu pasaules aizsardzības galveno ieroci pret vācu mierīgo iekarošanu. Krievi savukārt uzsvēra, ka Polijas lomas pieaugumam slāvu pasaulē nav jānonāk pretrunā ar Krievijas valstiskumu un nepieciešams saglabāt ciešu Polijas saikni ar impērijas centru. Pie tam Ministru padomes priekšsēdētājs I.~Goremikins uzsvēra, ka solījums atjaunot brīvu Poliju tiks realizēts tikai tad, ja Krievijas karaspēks ieņems Poznaņas apgabalu: ``ir Poznaņa utt., ir arī autonomija, nav Poznaņas,~--- nav arī autonomijas.'' Nosakot visas impērijas un Polijas kopīgas likumdošanas priekšmetu, poļu dalībnieki vēlējās aprobežoties ar kopīgu ārpolitiku, aizsardzību, naudas sistēmu, Krievu Pareizticīgās baznīcas un imperatora galma lietām. Krievu dalībnieki vēlējās bez nosauktajām jomām kopīgajā sfērā iekļaut arī finanses, tiesu iekārtu, satiksmes ceļus, pastu un telegrāfu. Poļu un krievu apspriedes dalībniekiem tā arī neizdevās vienoties. Apspriedes poļu dalībnieki kategoriski iebilda pret zemstu tipa apgabalu pārvaldes ieviešanu Polijā, uzskatot, ka jārunā par ``Polijas kā īpaša novada, īpašas tautas, ne pēc skaita un spēka, bet pēc savas kultūras un garīgās dzīves līdztiesīgas krievu tautai, nacionālu pašpārvaldi''. Diskusijas gaitā atklājās interesanta parādība, vismaz daļa krievu dalībnieku sāka sliekties uz to, ka Krievijas interesēs ir labāk radīt neatkarīgu, nevis autonomu Poliju. Kņazs D.~Svjatopolks-Mirskis pat paziņoja: ``Krievijai Polija nav vajadzīga''. Autonomijas ieviešana Polijā tika uzskatīta par Krievijas impērijas pārvaldes principu pārkāpšanu, kas varēja izraisīt ķēdes reakciju arī citos novados. (XX~gadsimta sākuma juridiskajā literatūrā termins ``autonomija'' tika traktēts kā tiesības uz patstāvīgu likumdošanu. Kādas teritorijas autonomija nozīmēja ``valsts daļu, kurai ir vietējais parlaments vai arī kāda cita vietējās likumdošanas varas organizācija''. Krievijā līdz Pirmajam pasaules karam valsts iestādes par autonomiju nerunāja, bet gan lietoja terminu ``vietējā pašvaldība''.)

Diskusija tika turpināta Ministru padomē. Pēdējo poļiem veltīto žestu šai periodā jau pēc Varšavas krišanas IV~Valsts domes kārtējās sesijas atklāšanas laikā 19.~jūlijā izdarīja Ministru padomes priekšsēdētājs I.~Goremikins. Viņš, atsaucoties uz imperatora norādījumiem, oficiāli paziņoja, ka pēc kara poļi iegūs tiesības ``brīvi noteikt savu nacionālo, kultūras un saimniecisko dzīvi uz autonomijas pamatiem, saglabājot vienotu valstiskumu [ar Krieviju]''. Ārlietu ministrs S.~Sazonovs vēlāk savās atmiņās rakstīja, ka pirms tam Ministru padomes sēdē 16.~jūlijā viņš protestējis pret šādu paziņojumu, jo ``jautājums par poļu autonomiju prasa nekavējošu atrisinājumu ar Viņa Augstības manifestu'', pat negaidot Valsts domes sēžu sākumu. Šāda manifesta uzdevums bija piesaistīt Krievijai poļu simpātijas. Justīcijas ministrs A.~Hvostovs turpretī uzskatīja, ka manifests par poļu autonomiju būs veltīgs, jo ``poļi sapņo par neatkarību un ar autonomiju viņus nevarēs apmierināt''.

Daļa carisko ierēdņu, galvenokārt no Ārlietu un Kara ministrijas, savlaicīgi no augšas veiktu reformu piekritēji, vēlējās pēc iespējas ātrāk piešķirt Polijai ko līdzīgu kulturāli-nacionālai autonomijai. 1915.~gada augustā, jau pēc vācu uzvarām Polijas teritorijā izveidojies Valsts domē Progresīvais bloks (kadeti, progresisti, kreisie oktobristi, Centra grupa, ``progresīvie krievu nacionālisti'') savā programmā izvirzīja prasību nekavējoties piešķirt Polijai autonomiju, likvidēt poļu tiesību ierobežojumus. Taču Krievijas valdība, ievērojot vācu un austriešu atturīgo attieksmi pret poļu sabiedrības mēģinājumiem uzzināt okupantu nodomus par Polijas nākotni, arī nolēma atturēties no tālākas poļu jautājuma apspriešanas.

Tiesa, kā pretsvars Austroungārijā J.~Pilsudska iniciētajam poļu leģionam Krievijas rokās esošajā Polijas pilsētiņā Pulavā (\pltxti{Puławy}) arī tika veidots poļu brīvprātīgo leģions (\pltxti{Legion Puławski}) izvērsta bataljona (ap tūkstoti cilvēku) apmērā. 1915.~gada pavasarī tas nonāca frontē, piedalījās kaujās. Formēja arī otru~--- Ļubļinas leģionu (\pltxti{Legion Lubelski}), kurš kaujās nepiedalījās, bet tika izmantots Pulavas leģiona papildināšanai. 1915.~gada oktobrī abus leģionus pārveidoja par Poļu strēlnieku brigādi. Taču Poļu nacionālā komiteja, kas sākotnēji piedalījās leģionu formēšanā, vīlusies vēlāk pārtrauca to darīt, jo, tāpat kā agrāk, vēl vācu neieņemtajā Polijas karalistē turpinājās rusifikācijas politika, kad Galīcijā poļiem bija daudz vairāk nacionālās brīvības.

Tikai 1916.~gada pavasarī sakarā ar situācijas maiņu Vācijā pēc ārlietu ministra S.~Sazonova iniciatīvas poļu jautājuma apspriešana Krievijas valdošajās aprindās atsākās. Krievijas ārlietu ministrs gribēja pārtvert Vācijas iniciatīvu radīt bufervalsti starp Krieviju un Vāciju. Viņš norādīja, ka, lai novērstu Polijas jautājuma starptautisku risinājumu, Krievijai tas ir pašai pēc iespējas ātrāk jāatrisina. Savās vēlāk tapušajās atmiņās viņš rakstīja, ka ``krievu un poļu starpā bija izliets pārāk daudz brāļu asiņu, lai to izlīgums notiktu citādi kā uz augstākā taisnīguma un savstarpējo vēsturisko tiesību pilnīgas atzīšanas pamata.'' Pēc S.~Sazonova priekšlikuma tika izstrādāts arī manifesta projekts, kurš guva imperatora atbalstu un tika iesniegts Ministru padomei. Projekta ievaddaļā bija teikts, ka tiek atzīts par nepieciešamu piešķirt Polijai visas brīvības, kuras ir savienojamas ar ``kopēju abām tautām [poļu un krievu~--- V.Š.] valstiskumu''. Taču Ministru padomes priekšsēdētājs B.~Štjurmers devās pie Nikolaja II un pēc šīs vizītes S.~Sazonovs, acīmredzot lielā mērā par to, ka viņš bija autonomijas piešķiršanas Polijai piekritējs, 1916.~gada 7.~jūlijā tika atstādināts no amata, bet ārlietu ministra pienākumus uzņēmās B.~Štjurmers. Manifesta projekts 13., 16. un 18.~jūlijā Ministru padomē tika apspriests klaji konservatīvā garā. Ministri baidījās, ka ``poļu zemju tiesiska norobežošana'' var radīt Krievijā virkni iekšpolitisku sarežģījumu. Esot jāievēro, ka Krievijā dzīvo arī citas tautas, kuras ``izsenis pārņemtas sapņiem par cilšu pašnoteikšanos''. Ministru padomes vairākums vienojās, ka līdz kara beigām ``Polijas jaunās iekārtas izlemšana varētu radīt pārpratumus''. Manifesta izdošanas vietā tika ierosināts pabeigt izstrādāt likumprojektu par Poliju un iesniegt to izskatīšanai Valsts domei un Valsts padomei. Tā, neraugoties uz to, ka visa Polijas karaliste sen bija Vācijas un Austroungārijas armijas okupēta un Krievijas karaspēks zaudējis kontroli pār visām teritorijām, kur pārsvarā bija poļu iedzīvotāji, Polijas autonomiju Krievija tā arī neizsludināja. (Vēlāk izrādījās, ka ar to no Krievijas puses tika pieļauta ievērojama kļūda, jo dažus mēnešus vēlāk~--- 1916.~gada 23.~oktobrī (5.~novembrī) Vācija un Austroungārija paziņoja par neatkarīgas Polijas monarhijas radīšanu.) Var apgalvot, ka tuvredzīgā carisko iestāžu politika tieši veicināja starpnacionālo pretrunu, no kurām tās tā baidījās, pieaugumu. Pat uzdevums mobilizēt visus valsts spēkus cīņai pret ārējo ienaidnieku nepiespieda Krievijas varas iestādes pietiekamā mērā ievērot nekrievu tautu vēlmes.

Krievijas Rietumu sabiedrotie bija tendēti pēc iespējas ātrāk izlemt Polijas jautājumu, lai Krievija piekāptos poļiem un Antantes pusē karotu poļu karavīri. Taču bez Krievijas to nekādi nevarēja panākt. Rietumu lielvalstis līdz pat 1917.~gadam uzskatīja poļu jautājumu, visu, kas attiecās uz Vācijas un Austroungārijas sastāvā esošajām poļu zemēm, kā pilnībā ietilpstošu vienīgi Krievijas kompetencē. Lielbritānija neizrādīja interesi par Austrumeiropu. Arī Francija bija spiesta ievērot sava sabiedrotā~--- Krievijas pozīciju.

Citādi kā Polijas karalistē notikumi risinājās Austroungārijai piederošajā Galīcijā. 1914.~gada rudenī Galīcijas kaujas rezultātā (5.~(18.)~augustā~--- 13.~(26.)~septembrī) \strong{Krievijas karaspēks ieņēma praktiski visu Austrumgalīciju} un daļēji arī Rietumgalīciju. No 1914.~gada rudens Krievijas politika Galīcijā dalījās divās daļās, paredzot Rietumgalīcijas iekļaušanos nākamajā Polijā un Austrumgalīcijas pievienošanu Krievijas impērijai. Kara apstākļos tika mēģināts uzreiz realizēt Austrumgalīcijas sapludināšanu ar Krievijas impēriju. Krievu politiķu plāni tika pamatoti ar šo teritoriju vēsturisko vienotību, vēsturisko Galīcijas saistību ar Senkrievu valsti un Ukrainu (Mazkrieviju), kura atradās Krievijas sastāvā.

Ļvovā (agrākajā Lembergā) tika radīta Galīcijas ģenerālgubernatūra no vairākām guberņām (sākotnēji Ļvovas un Ternopoles (ukraiņu \uktxti{Тернопіль}, poļu \pltxti{Tarnopol)}, pēc tālākas Krievijas armijas pavirzīšanās arī Černivcu (ukraiņu \uktxti{Чернівці́}, poļu \pltxti{Czerniowcez}) un Peremišļas (poļu \pltxti{Przemyśl}, ukraiņu \uktxti{Перемишль}) ar grāfu G.~Bobrinski priekšgalā. Tā darbojās līdz 1915.~gada jūlijam. Vietējie iedzīvotāji gan ieguva tikai tulku un sīko ierēdņu amatus, jo lielāko daļu krieviskās orientācijas inteliģentu austriešu varas iestādes jau pašā kara sākumā bija represējušas, bet pārpalikušajiem krievu administrācija neuzticējās. Rietumgalīcijā, kur dzīvoja samērā daudz poļu, par ierēdņiem iecēla Krievijas izcelsmes poļus. Kara apstākļos tas nenācās viegli. Arī krievu ierēdņi nelabprāt pārcēlās dienēt uz piefrontes guberņām, tāpēc to kvalitatīvais sastāvs bija samērā zems. Zēla kukuļņemšana un dienesta stāvokļa ļaunprātīga izmantošana.

Jau no 1914.~gada septembra Galīcijas vidējās un sākumskolās, ``lai garīgi tuvinātu'' galīciešus ar krievu tautu, tika uzsākta krievu valodas mācīšana. Austrumgalīcijā tika slēgtas visas skolas ar vācu, poļu un ukraiņu mācību valodu, pret ko protestēja poļu pārstāvji Krievijas Valsts domē un Valsts padomē. Tikai ar 1915.~gada janvāri ģenerālgubernators G.~Bobrinskis Ļvovā un Galīcijā atļāva atvērt vairākas privātas krievu un poļu skolas. Pēdējās krievu valoda bija jāmāca vismaz 5~stundas nedēļā, jālieto Krievijas Izglītības ministrijas atzītas vēstures, ģeogrāfijas un poļu valodas un literatūras mācību grāmatas, skolotāju sastāvam bija jābūt ģenerālgubernatora apstiprinātam.

Taču saglabājās Austrijas tiesu sistēma un likumi. Krievijas justīcijas ministrs I.~Ščeglovitovs uzskatīja, ka pirms Galīcijas pievienošanas Krievijai tiesu sistēmas pārveidojumi ir priekšlaicīgi, jo var izsaukt vietējo iedzīvotāju neapmierinātību. Pārkārtojumi tiesu sistēmā norisa tikai tajā lietojamo valodu jomā, jo Austrijā sarakste ar augstākajām tiesu instancēm notika vāciski, vietējā lietvedība un tiesu procesi notika poļu valodā. 1915.~gada sākumā tika atļauta krievu un ukraiņu valodas un ``pagaidām'' arī poļu valodas lietošana tiesās. Taču ieņemtajās teritorijās tūlīt tika pielietota Krievijas impērijas antisemītiskā likumdošana: ebrejiem aizliedza būt zemes īpašniekiem (Austroungārijas sastāvā tas bija atļauts), sākās viņu zemes ekspropriācija, kuru pārtrauca tikai novada zaudēšana 1915.~gada vasarā. Faktiski uz visiem Galīcijas ebrejiem raudzījās kā uz potenciāliem spiegiem un nodevējiem. Simti ebreju tika deportēti uz frontes aizmuguri.

Kaut Krievijas imperators Nikolajs~II un armijas Virspavēlnieks lielkņazs Nikolajs Nikolajevičs deva norādījumu okupācijas varas iestādēm nelikt nekādus šķēršļus vietējo uniātu darbībai, reliģisko jautājumu risināt ar maksimālu piesardzību un ``nelikt vienlīdzības zīmi starp politisku neuzticamību un reliģisku atšķirtību'' (``\rutxti{не отождествлять политическую неблагонадежность с религиозной разъединенностью}''), pareizticīgo arhibīskaps Jevlogijs, kurš bija nozīmēts vadīt baznīcas lietas okupētajos novados, ar saviem piekritējiem nemaz nevēlējās rīkoties uzmanīgi un, kad 1914.~gada rudenī kļuva skaidrs, ka nav perspektīva uniātu garīdzniecības piesaistīšana pareizticībai, realizēja vardarbīgu uniātu pievienošanas pareizticībai politiku. Jau tikai Jevlogija ierašanos vien Galīcijā vietējie poļu iedzīvotāji uztvēra kā antipoliska un antikatoliska kursa realizācijas sākumu. No novada izsūtīja daudzus uniātu baznīcas garīdzniekus, arī metropolītu A.~Šepticki, kurš kara sākumā vēstījumā draudzei bija paziņojis, ka ``Maskavas cars nevarēja paciest to, ka mums Austrijas valstī pastāv ticības un tautības brīvība'', pretojās uniātu pakļaušanai pareizticīgajai baznīcai arī turpmāk. Metropolītu izsūtīja uz Kijevu, tad Ņižņijnovgorodu, Kursku, Jaroslavļu, beidzot uz vienu no Suzdaļas pareizticīgo klosteriem. Brīvību viņš atguva tikai 1917.~gada martā. Uniātu garīdznieku vietā arhibīskaps Jevlogijs ar ģenerālgubernatora atļauju iecēla pareizticīgos garīdzniekus. Pavisam pēc ģenerālgubernatora kancelejas datiem tika iecelti 86~pareizticīgie garīdznieki. Taču nereti draudžu oficiāla pāreja no uniātisma pareizticībā bija atkarīga no tā, kura puse~--- vietējie uniāti vai pareizticīgie~--- spēja piekukuļot apriņķa priekšnieku vai policijas ierēdņus. Pēc oficiālajiem pareizticīgo baznīcas datiem divos mēnešos pēc Galīcijas okupācijas pareizticībā bija pārgājuši 30~000 uniātu. Taču par tādiem uzskatīja visus, kuri bija izsūdzējuši grēkus pareizticīgo baznīcā vai noklausījušies pareizticīgā garīdznieka sprediķi, kad vietējie iedzīvotāji pareizticīgās baznīcas apmeklējumu nebūt nesaistīja ar atteikšanos no savas ticības. Austrumgalīcijā tika radītas 50--100~pareizticīgo draudzes. Ievērojot, ka pirms kara Galīcijā dzīvoja ap 3,5~miljoniem uniātu, 1914.~gadā pareizticībā no tiem pārgāja 0,85\%. No ap 2~000 Galīcijas uniātu garīdzniekiem tikai 29~pārgāja pareizticībā. Tomēr arhibīskapa Jevlogija darbība, izsaucot vietējo iedzīvotāju neapmierinātību, arī atbalsi ārzemēs, darīja bažīgas militārās un civilās amatpersonas. Rezultātā 1915.~gada aprīlī arhibīskaps ar godu tika atsaukts no Galīcijas.

1914.~gada rudenī tika mēģināts organizēt nodokļu ievākšanu no Galīcijas iedzīvotājiem, taču jau oktobrī kļuva pilnīgi skaidrs, ka viņi ir kara izpostīti. Izrādījās, ka no viņiem ne tikai nevarēja gūt kādus ienākumus valsts kases papildināšanai, bet bija jādomā vēl par atbalstu tiem.

Kopumā Krievijas varas iestāžu politika Galīcijā bija nekonsekventa. Poļu ekonomists un politiķis S.~Grabskis atzīmēja, ka dažādi apriņķu priekšnieki realizēja atšķirīgu politiku. Vieni īstenoja atbrīvošanas kara ideju, otri iekarotāju politiku. Šī politika nereti izsauca arī vietējo rusofīlu, kuri bija konsekventāki prokrieviskās un antipoliskās politikas piekritēji nekā Krievijas administrācijas ierēdņi, neapmierinātību. Novadā saglabājās ievērojama poļu muižnieku, poļu kultūras, katoļu un uniātu baznīcas ietekme.

1915.~gada aprīlī Vācijas un Austroungārijas karaspēks Galīcijā pārgāja uzbrukumā. Tā pusē bija liels materiālais pārsvars, īpaši ložmetēju, smagās artilērijas un munīcijas ziņā. Krievijas karaspēks atstāja Peremišļu, Ļvovu un lielāko daļu Galīcijas. Lai neļautu pretiniekam pilnvērtīgi izmantot Galīcijas saimniecisko potenciālu, pirms tam tika izpostīti ap 200~naftas urbumu un sadedzinātas 350~000~t jēlnaftas. Līdz 1915.~gada augustam Krievijas aizmugurē bija nonākuši arī ap 100~tūkstošu bēgļu no Galīcijas. Bez tam 1915.~gadā vien Krievijas iestādes saņēma, pēc dažādiem datiem, no 5,5 līdz 7~simtiem ķīlnieku. Parasti tie bija uzņēmēji, banku direktori u.tml., starp tiem daudz ebreju. Tika paskaidrots, ka ebreji tiek ņemti par ķīlniekiem tāpēc, ka pēc viņu ticības brāļu norādēm austriešu iestādes jau no Krievijas armijas atkarotajos rajonos esot vajājušas ``rusīnus'' (ukraiņus) par sadarbību ar krievu okupācijas varas iestādēm. Kopumā 1914.--1915.~gados cara varas iestādes administratīvā kārtā no Galīcijas izsūtīja gandrīz divus tūkstošus cilvēku.

Pēc Galīcijas atkarošanas Krievijai šeit tiešām izvērsās austriešu terors pret vietējiem iedzīvotājiem, kuri tika turēti aizdomās par ``nodevību''~--- simpātijām Krievijai. Palikušos vietējos rusofīlus, pēc dažām ziņām līdz 20~tūkstošiem, austriešu iestādes nosūtīja uz koncentrācijas nometnēm Terezinā (\detxti{Theresienstadt}) un Talerhofā (\detxti{Talerhof}).

1916.~gada vasarā Galīcijas austrumdaļu atkal ieņēma Krievijas karaspēks, taču tagad tā ieiešana šai teritorijā vairs nebija tik iespaidīga kā 1914.~gadā, vairs neviens nerunāja arī par gaidāmo Galīcijas apvienošanos ar Krieviju. 1916.~gada jūlijā atkal tika izveidota ģenerālgubernatūra, šoreiz jau ar ģenerāli F.~Trepovu vadībā. Viņš, balstoties uz militārās pavēlniecības nostādnēm, centās izvairīties no ``politiskas tendences novada pārvadīšanā'', atcēla ierobežojumus poļu un ukraiņu valodām skolās, nesteidzās radīt Galīcijā krievu skolas. Tas izsauca toreizējā ārlietu ministra B.~Štjurmera iebildumus, kurš uzskatīja, ka ``krievu valodas provinciālie dialekti'' jeb ukraiņu valoda var gan tikt lietota pašvaldībās, bet skolās mācībām jānotiek ``vienīgajā nacionālajā valodā'', proti~--- krievu valodā. Tomēr uzvarēja militāro iestāžu viedoklis, kuras vairāk rūpējās par miera un kārtības uzturēšanu piefrontes joslā, nekā par nacionālo audzināšanu. Drīz gan sekoja Februāra revolūcija Krievijā, tās armijas sabrukums, kas veda arī pie Galīcijas atstāšanas.

\strong{Kara laikā Vācijas un Austroungārijas okupētās poļu zemes} faktiski dalījās divās daļās: vācu un austriešu. Robeža starp tām gāja pa Varšavas~--- Vīnes dzelzceļa līniju. Sadales līnija aptuveni sakrita ar to robežu, kura starp Prūsiju un Austriju izveidojās pēc Žečpospolitas sadales 1795.~gadā.

Teritorijā, kuru kontrolēja Vācija, krievu valoda administratīvajās iestādēs, kuru struktūra pamatos palika nemainīga, tika nomainīta ar vācu valodu. Poļu valoda tika saglabāta iepriekšējā apjomā. Vācija objektīvi nebija ieinteresēta patstāvīgas Polijas valsts izveidē, pat tad, ja tajā tiktu iekļauta tikai Polijas karalistes teritorija, jau tāpēc vien, ka jaunā valsts noteikti pieprasītu sev vācu valdījumā esošo Pozenes provinci. Taču kara laikā jautājums par Polijas, pareizāk par tās Polijas daļas, kura līdz tam piederēja Krievijai, likteni kļuva par visai dzīvas diskusijas objektu Vācijā, tās presē, dažādās sapulcēs u. tml.

``Krievijas tautu atbrīvošana'' kā līdzeklis sašķelt Krieviju kļuva par vienu no galvenajiem Vācijas uzdevumiem Austrumu frontē. Vācijas propagandisti aktīvi darbojās Krievijas mazākumtautību vidū. Piemēram, vācu varas iestādes radīja īpašas karagūstekņu nometnes ukraiņiem, kur ar viņiem strādāja propagandisti no Galīcijas. Kad atklājās, ka ``krievu Ukrainas'' iedzīvotāji atšķiras no galīciešiem, šie propagandisti tika nomainīti ar citiem, jau no pašu ``krievu ukraiņu'' vidus.

Polijas iedzīvotājus Vācija un Austroungārija vēlējās izmantot bruņotai cīņai pret Krieviju. Jau 1914.~gadā, Vācijas un Austroungārijas armijām ieejot t.s. ``kongresa Polijas'' (Polijas karalistes) teritorijā, tika paziņots: ``Mēs nesam jums brīvību un neatkarību \citespace{} lai austrumu barbarisms atkāpjas rietumu civilizācijas priekšā, kura ir kopīga gan jums, gan mums. Saceļaties, pievienojaties mūsu karaspēkam, un kopīgiem spēkiem padzīsim no Polijas robežām aziātu ordas.''

Taču vācu karavīri pret poļu iedzīvotājiem izturējās visai nežēlīgi. Preses konferencē Pozenē ģenerālis T.~f.~Zaubercveigs paziņoja: ``Nedrīkst būt mīkstam attiecībā pret poļiem, jārīkojas cik iespējams nežēlīgi. Vācijas armijas un Vācijas tēvijas īpašums ir augstāks par visu. Virspavēlniekam vairāk nozīmē viena vācieša dzīvība kara laukā un jebkura vācieša uzturēšana Vācijā nekā visi okupētas zemes iedzīvotāji.'' Šāds uzstādījums tika arī pildīts. Jau 1914.~gada 3.~augustā, vācu armijai ieņemot Polijas karalistē Kališas (poļu \pltxti{Kalisz}, vācu \detxti{Kalisch)} pilsētu, kuru krievu karaspēks jau bija atstājis, nenoskaidrotos apstākļos pēkšņi sākās nekārtīga šaudīšanās, un vācu artilērijas apšaudē 20\% no pilsētas tika pārvērsti krāsmatās. Bija daudz upuru civiliedzīvotāju vidū, kas poļos izsauca sašutumu par prūšu barbarismu. Vēlāk kara darbībā tika izpostīta rūpnieciskā Žirardovas (\pltxti{Żyrardów}) pilsēta, stipri cieta Lodza, daudzas citas apdzīvotas vietas. Par to, ka Vācijas armija ne sevišķi augstu vērtēja civilizācijas sasniegumus poļu zemēs un, iespējams, arī tāpēc necentās saudzēt poļu materiālās un kultūras vērtības, liecina toreiz vēl ģenerāļa P.~f.~Hindenburga atmiņas: ``Mēs sasniedzām Krievijas Poliju un tūlīt varējām pilnībā iepazīties ar to, ko kāds franču ģenerālis savā Napoleona 1806.~gada kara gājiena aprakstā apzīmēja kā vietējā kara vešanas elementu, proti~--- ar dubļiem. Dubļiem katrā iedomājamā formā, ne tikai brīvā dabā, bet arī tā saucamajos cilvēku mitekļos un pie pašiem cilvēkiem. Ar mūsu robežas šķērsošanu mēs nonācām pavisam citā pasaulē. Nevilšus radās jautājums: kā tas ir iespējams, ka uz Eiropas zemes robežakmeņi starp Pozeni un Poliju iezīmē tik krasu norobežošanās līniju starp tā paša celma tautas divām kultūras pakāpēm?''

Jautājumu par Polijas likteni Vācija ilgstoši atlika. Piemēram, poļu avīzēm bija aizliegts savās slejās skart Polijas likteņa jautājumu, pat lietot tādus vārdus kā ``neatkarība'' u.tml.

Varšavas nonākšana Vācijas karaspēka rokās 1915.~gada augustā pastiprināja vācu valdošo aprindu interesi par Poliju. Vācijas ķeizars Vilhelms~II, tās kanclers T.~f.~Betmans Holvegs Polijas karalistes pievienošanu Vācijai uzskatīja par nelietderīgu. Tikai atsevišķas tās teritorijas, pēc viņu domām, Vācijai bija nepieciešamas tās austrumu robežu ``stratēģiskai uzlabošanai''. Tāpēc no daļas Polijas karalistes un baltkrievu zemēm bija paredzēts radīt poļu valsti, kura atrastos militārā un ekonomiskā savienībā ar Vāciju (faktiski~--- atkarībā no tās). Savukārt vācu militāristi ar ģenerālfeldmaršalu P.~Hindenburgu priekšgalā uzskatīja, ka Polijas karalisti ir jāsadala starp Vāciju un Austroungāriju okupācijas zonās.

Reizē abās valstīs: Vācijā un Austroungārijā bija arī politiķi, kuri cerēja noslēgt ar Krieviju separātu mieru. Taču 1915.~gadā, neraugoties uz Krievijas militārajām sakāvēm, tas neizdevās.

Austroungārija gan bija ieinteresēta Polijas jautājuma ``austroungāriski-poliskā'' risinājumā, taču tā bija vājākā no karojošajām Eiropas lielvalstīm. Tās iekšējās problēmas arī neļāva izvirzīt skaidrus kara mērķus. Austrieši Polijas karalistes teritorijā centās uzstāties kā poļu atbrīvotāji no ilgstošas citas valsts kundzības. Austroungārija tās kontrolētajās poļu teritorijās ieviesa pārvaldi, kāda pastāvēja Austrijas Galīcijā. Krievu valodas vietā aktīvi tika ieviesta poļu valoda. Faktiski tie bija mēģinājumi viņu ieņemtajās Polijas karalistes guberņās radīt autonomiju, kuru Krievija solīja tikai pēc kara beigām. Tomēr tos poļus, kuri orientējās uz Austroungāriju, gaidīja daudzi sarūgtinājumi. 1915.~gada septembrī tapa zināms, ka Austroungārijai nav nodoma aizstāvēt Polijas karalistes nedalāmību un tās pievienošanu Galīcijai, jo Austroungārija bija vairāk ieinteresēta paplašināt savu teritoriju uz Ukrainas rēķina. Tiesa, 1916.~gada ziemā austriešu un liela daļa ungāru politiķu tomēr atzina Polijas karalistes pievienošanas Austroungārijai vēlamību. 1916.~gada sākumā Austrijas valdība veica virkni pasākumu, lai nostiprinātu savu aizmuguri Polijā. Tika paziņots par Holmas guberņas likvidēšanu un iekļaušanu Ļubļinas guberņā, poļu sabiedrisko organizāciju iesaistīšanu tautas izglītības darbā un iedzīvotāju apgādē ar pārtiku. Faktiski tā bija reakcija uz Krievijas armijas uzbrukuma sākumu Dienvidrietumu frontē. Toties ukraiņu presē tika rakstīts, ka Austrija atbalsta poļu ``imperiālismu''. Sevišķi ukraiņus uztrauca Holmas liktenis un baumas par austriešu mēģinājumiem latinizēt ukraiņu alfabētu, ieviešot to ukraiņu skolās Galīcijā. Tiesa, krievu armijas uzbrukuma gaitā Austroungārijas politiķi solīja ukraiņiem atvērt Lembergā (Ļvovā) ukraiņu universitāti.

Vācijas un Austroungārijas valdības apsprieda projektu par Polijas karalistes, pat nedaudz to paplašinot ziemeļrietumos, pievienošanu Austroungārijai. Austriešu valstsvīri izstrādāja savu risinājumu~--- ``kongresa Polija'' un Galīcija būtu jāapvieno, radot trešo partneri Habsburgu monarhijā, kas atbilda J.~Pilsudska un viņa piekritēju cerībām. Tam gan pretojās ungāru politiķi, kuri protestēja pret trīs partneru principu neatkarīgi no tā vai šis trešais būtu dienvidslāvi vai poļi, un kuri gribēja saglabāt Galīciju kā unitāras Austrijas daļu, lai poļi tajā kaut daļēji neitralizētu vāciešu ietekmi. Par Polijas palikšanu katoliskās Austroungārijas sastāvā iestājās Vatikāns, nevēloties neko dzirdēt par neatkarīgu Poliju.

Vācija bija gatava atdot Polijas karalisti Habsburgu monarhijai gadījumā, ja tā noslēgtu ilgstošu militāru un arī ekonomisku savienību ar Vāciju, kurā Vācijai, protams, būtu vadošā vieta. Tā būtu tā saucamās Viduseiropas koncepcijas realizācija, kas paredzēja nodibināt Centrālajā Eiropā Vācijas hegemoniju. Taču daļa vācu politiķu baidījās no izmaiņām Polijā, jo neatkarīgs vai Habsburgu valstī iekļauts poļu valstiskums varētu ietekmēt poļu zemes Vācijas sastāvā (kā jau teikts, daļēji apvienotā Polija neizbēgami pieprasītu sev arī Pozenes apgabalu) un ar to iznīcinātu visas iepriekšējās pūles tās ģermanizējot. Vācijas sabiedrībā izskanēja dažādi viedokļi. Pazīstamais vācu sociologs M.~Vēbers, baidoties, ka ``vācu nacionālā politika'' var padarīt poļus par nāvīgiem Vācijas ienaidniekiem un Krievijas piekritējiem, uzskatīja par lietderīgu radīt autonomas poļu, ukraiņu u.c. valstis, kuru teritorijā cietokšņi un stratēģiskie ceļi atrastos pilnīgi centrālo lielvalstu rokās. Vēsturnieks un politologs H.~Delbruks ieteica Poliju saistīt ar Saksiju, kā tas kādreiz jau bija Polijas vēsturē. Tad, formāli neatrodoties Vācijas sastāvā, Polija tomēr būtu ar to politiski, ekonomiski un dinastiski cieši saistīta. Vācu militāristi savukārt gatavojās anektēt lielu Polijas daļu, faktiski veikt sesto (ja trijām dalīšanām XVIII gadsimtā pieskaita arī divas XIX~gadsimtā notikušās) Polijas teritorijas dalīšanu. Polijas karalistes teritorijā tad varētu izveidot nelielu, no Vācijas atkarīgu valstiņu. Kanclers T.~f.~Betmans Holvegs, kurš nojauta, ka ``mums [Vācijai] labvēlīgs un drošs Polijas problēmas risinājums vispār neeksistē'', gan pret to iebilda.

Vietējie Polijas karalistes vācieši parasti palika visai pasīvi. Vācijas realizētā karalistes okupācija viņiem nāca negaidīti, kaut kāda viņu pozīcija nebija izstrādāta. Pirms tam viņi kā Krievijas pavalstnieki bija devušies dienēt Krievijas armijā, arī karojuši pret Vāciju. Tagad tikai pakāpeniski orientācija mainījās uz Vācijai labvēlīgu. Piemēram, Lodzas tekstilrūpniekiem provinces atraušana no Krievijas nāca tikai par sliktu. Viņi bija prokrieviski noskaņoti. Vācijas okupācijas varas iestādes gan poļu, gan vācu fabrikantu preces un izejvielas konfiscēja, mašīnas demontēja un aizveda uz Vāciju. Sākotnēji pastāvēja tikai neliela grupiņa vietējo vācu t.s. ``aktīvistu'', kuri centās uzkurināt vietējo vāciešu pašapziņu un stiprināt novada saistību ar Vāciju. Vācu okupācijas varas iestādes gan pret šiem vācu ``aktīvistiem'' izturējās piesardzīgi, jo negribēja saasināt attiecības ar poļiem. Tikai 1916.~gada 5.~martā vācu ``aktīvistu'' grupa nodibināja Vācu savienību (\detxti{Deutscher Verein}), kurā bija ap 230~vietējo grupu un vairāk nekā 30~000 biedru. Savienība nodarbojās ar nacionālo kultūras darbu, taču tam bija arī politiska nokrāsa.

Drīz Vācija sāka atteikties no saviem solījumiem Austroungārijai. Strīdi par Polijas likteni gan Vācijas valdībā un sabiedrībā, gan starp Vāciju un Austroungāriju turpinājās. Protams, vācu un austriešu pretenzijas uz poļu apdzīvotajām Krievijas teritorijām poļu iedzīvotāju vidū nekādu sajūsmu neizsauca.

1915.~gada vasarā Vācijas ieņemtajā poļu teritorijā tika nodibināta Varšavas ģenerālgubernatūra ar ģenerāli H.~f.~Bezeleru priekšgalā un Austroungārijas ieņemtajā poļu teritorijā~--- Ļublinas militārā gubernatūra ar ģenerāli K.~f.~Kuku priekšgalā. (Jau minētais latviešu vēsturnieks A.~Bīlmanis, neievērojot Polijas teritorijas dalīšanu XIX gadsimtā, to nosaucis par ``it kā ceturto Polijas sadalīšanu starp vāciem un austriešiem''.) Varšavas ģenerālgubernators ģenerālis H.~f.~Bezelers uzrakstīja memorandu, kurā apgalvoja, ka ``patstāvīgas'' Polijas radīšana ļaus uzstādīt vismaz trīs divīzijas ar poļu brīvprātīgajiem, kuri cīnīsies vācu pusē. Vācijas ieņemtajā Polijas daļā tika publicēts aicinājums poļiem veidot brīvprātīgo daļas cīņai pret Krieviju, jo pēc Krievijas impērijas īstenotās mobilizācijas un vietējo iedzīvotāju evakuācijas uz Iekškrieviju Polijas karalistē joprojām bija aptuveni viens miljons militārajam dienestam derīgu vīriešu.

Kad 1915.~gada rudenī visu Polijas karalisti okupēja Centrālo valstu karaspēks, lai arī visa augstākā vara palika militāro iestāžu rokās, tika veiktas arī zināmas pārmaiņas civilajā pārvaldē. 1915.~gada novembrī Varšavā tika atvērtas poļu augstākās mācību iestādes~--- Universitāte un Politehniskais institūts. Lai mazinātu poļu pretestību, okupantiem vācu militārās iestādes bija spiestas piekrist poļu pašpārvaldes iestāžu ievešanai, kuras gan vairāk nodarbojās ar kultūras jautājumiem. Pēc gadiem ilgās rusifikācijas skolās ieviesa poļu mācību valodu. Poļu valodu izmantoja gminās un tiesvedībā. Pārmaiņu rezultātā radās jaukta militāri~--- civila pārvaldes sistēma. 1916.~gadā ar ģenerālfeldmaršala P.~f.~Hindenburga pavēli baltkrievu zemēs tika atļauts izmantot vietējo iedzīvotāju~--- poļu, lietuviešu, baltkrievu valodas, bet aizliegts izglītībā, presē un administrācijā lietot krievu valodu.

Būtībā kara pirmajā posmā, līdz 1916.~gadam, poļu jautājums saglabāja savu iepriekšējo, Poliju sadalījušo impēriju iekšējas lietas raksturu, bet nacionāldemokrātu un J.~Pilsudska piekritēju mēģinājumi to izkustināt no sastinguma punkta savu pirmskara koncepciju ietvaros nedeva kādus reālus rezultātus. Triju impēriju varas iestādes labvēlīgi pieņēma tām uzticību demonstrējošas deklarācijas, taču nespēra nekādus praktiskus soļus poļu nacionālo vēlmju apmierināšanai.

Kara laikā ievērojami pasliktinājās \strong{poļu iedzīvotāju ekonomiskais stāvoklis}. Poļu apdzīvotās teritorijas tika nesaudzīgi ekspluatētas. Īpaši Vācijas kontrolēto apgabalu iedzīvotājos radās pret okupantiem vērsti noskaņojumi. Austrieši uzvedās daudz-maz mēreni, turpretī vācieši rīkojās bez ceremonijām, dažkārt īsti nežēlīgi. Jau pirmajā okupācijas pusgadā piespiedu darbā uz Vāciju tika nosūtīti 80~000 strādnieku. Kopumā vairāk nekā 200~tūkstošu poļu strādnieku no bijušās Polijas karalistes teritorijas tika vardarbīgi sūtīti uz Vāciju darbam dzelzceļu un zemes ceļu celtniecībā, militārajā, metalurģiskajā, meža un kalnrūpniecībā. Vācijā palika arī ap 300~000 poļu sezonas strādnieku, kurus karš tur pārsteidza.

Poļu zemēs auga bezdarbs. 1916.~gadā visā vācu okupācijas zonā bija nodarbināti tikai 22\% no kopīgā strādnieku skaita. Bezdarbs spieda strādniekus meklēt darbu uz laukiem, pārdot savas darbarokas par iztiku. Rezultātā lielākās rūpniecības pilsētas zaudēja 1/4 un pat vairāk savu iedzīvotāju. Okupācijas režīma apstākļos strādnieki zaudēja visus savus pirmskara sociālos iekarojumus. Darba diena tika pagarināta līdz 10--12 stundām, pārtraukumi darba laikā bija samazināti vai atcelti, soda naudas paaugstinātas utt.

Pēc oficiāliem datiem kara laikā (1914--1918) no Polijas uz Vāciju tika izvesti 60~tūkstoši vagonu rūpniecības iekārtu, 380~tūkstošu vagonu lopu, ap 600~tūkstošu vagonu izejvielu un dažādu preču. Liela daļa lauksaimniecības produkcijas (mājlopi un labība) tika konfiscēta. Armijas vajadzībām rekvizēja vairāk par miljonu zirgu, kas bija galvenais vilcējspēks lauksaimniecībā. 1917.~gadā Vācijas okupācijas zonā lauksaimniecības pamatkultūru ievākums salīdzinājumā ar 1913.~gadu bija samazinājies: kviešiem par 29,9\%, rudziem~--- 58,7\%, auzām~--- 41,7\%, miežiem~--- 36,1\%, kartupeļiem~--- 53,9\%. Tika ieviestas pārtikas produktu kartītes, turklāt uz tām saņemamo produktu normas, ko noteica okupācijas varas iestādes, bija ievērojami zemākas nekā Vācijā. Tās pakāpeniski vēl samazinājās, bet cenas nepārtraukti cēlās. 1917.~gadā pārtikas deva dienai bija 131~gr maizes, 205~gr kartupeļu un 7~gr putraimu. 1917.~gadā tikai Varšavā no bada nomira 2~250~cilvēku. Strauji pieauga saslimstība ar tuberkulozi. Mirstība no tās 1917.~gadā bija piecas līdz sešas reizes augstāka nekā pirms kara. Austriešu Ļubļinas ģenerālgubernatūras informācijas nodaļa 1917.~gada jūlija sākumā ziņoja ārlietu ministram O.~Čerņinam, ka Polijā pilsētās pūļi uzbrūk un izlaupa pārtikas veikalus, bet laukos badā esošie izrok un apēd tikko iestādītos kartupeļus. Zemnieki sūdzējās par varas iestāžu rekvizīcijām, par to, ka viņiem tika atņemtas rokas dzirnaviņas graudu samalšanai, lai viņi tos nemaltu slepus savās mājās, un tāpēc tos nākas sasmalcināt piestās.

Vācieši poļu apgabalos pat organizēja ražošanas iekārtu, izgatavotu no metāla, pārkausēšanu, kas deva iespēju ne tikai paplašināt ieroču ražošanu, bet novērst vācu rūpniecībai nevēlamo konkurenci. Poļu strādnieku reālā darba alga 1916.~gadā (rēķinot pa pusgadiem) sastādīja 23 un 20\%, bet 1917.~gadā~--- 17 un 12\% salīdzinājumā ar priekškara laiku.

Poļu vidū brieda apziņa, ka abas karojošās puses nebūt nedomā par poļu labklājību. Kara gadu pieredze palīdzēja veidoties nacionālajai pašapziņai, vēlmei atjaunot nacionālo neatkarību.

Karam turpinoties, Vācijas valdošajās aprindās pārsvaru sāka gūt koncepcija par nepieciešamību starp Vāciju un Krieviju radīt bufervalsti~--- Poliju. Faktiski stratēģiski-ekonomiskā realitāte diktēja savas vajadzības. Vācijā kara gaitā arvien vairāk bija izjūtams kā karavīru, tā darbaroku un resursu trūkums. Kā dabīga izeja no situācijas šķita poļu zemju resursu izmantošana.

Tomēr tikai 1916.~gadā smagie vācu un austriešu zaudējumi Rietumu frontē pie Verdenas un Sommas, un t.s. Brusilova pārrāvuma (krievu armijas uzbrukuma Galīcijā un Volīnijā) laikā izsmēla Vācijas un Austroungārijas spēkus, ienesa redzamas izmaiņas Vācijas un Austroungārijas ``poļu politikā''. Abām valstīm bija nepieciešams jauns dzīvā spēka avots. Tāds varēja būt okupēto apgabalu iedzīvotāji. Vācijas armijas pavēlniecība īpaši cerēja uz pret Krieviju noskaņotu poļu karavīru izmantošanu. Vēlāk viens no Vācijas armijas vadītājiem E.~Ludendorfs atzinās: ``Manās acīs vienīgā nozīme bija jautājumam par poļu armiju, kas dotu spēku pieaugumu turpmākajai kara vešanai''. Tika lēsts, ka poļi var dot no 900~000 līdz 1 miljonam karotāju. Iesaukt tos savās armijā mobilizācijas kārtībā nebija iespējams, jo starptautiskās tiesības aizliedza okupēto valstu iedzīvotāju mobilizāciju iekarotāju armijās, bija arī gaidāma mobilizācijas sabotāža.

Palika viena izeja~--- izveidot īpašu ``poļu valsti''. Tas, kā cerēja Vācijas varas iestādes, ļautu Vācijas politiķiem reklamēt sevi par poļu valstiskuma atjaunotājiem un atvieglotu poļu vervēšanu armijā. 1916.~gada 4.--5.~aprīlī vizītē Berlīnē Austroungārijas ārlietu ministrs I.~Burians no Vācijas kanclera T.~f.~Betmana Holvega uzzināja par Vācijas nodomu radīt Polijas karalistē marionešu bufervalsti.

Sarunās poļu jautājumā 1916.~gadā 11.--12.~augustā Vācija un Austroungārija vienojās paziņot par ``patstāvīgas Polijas karalistes'' izveidi, taču, tāpat kā iepriekš, nenosakot precīzi tās atkarības pakāpi. Taču bija skaidrs, ka šī karaliste atradīsies kopīgā Vācijas un Austroungārijas kontrolē, politiskā, ekonomiskā un militārā ziņā būs pilnībā no tām atkarīga.

Vācijas kanclers T.~f.~Betmans Holvegs un Austroungārijas ārlietu ministrs I.~Burians norunāja nodalīt no okupētajām zemēm papildus apgabalus nolūkā ``izlabot savas robežas ar nākotnes Poliju''. Bija paredzēts daļu vācu-poļu pierobežas zemju, kuras līdz tam atradās Polijas karalistē, pievienot Vācijai, bet daļu baltkrievu un ukraiņu zemju ārpus Polijas karalistes~--- pievienot nākamajai Polijas valstij. Tika arī nolemts, ka ``nākotnes Polijai'' nebūs savas ārpolitikas, bet tās bruņotie spēki atradīsies Vācijas kontrolē. Austroungārijai tas nozīmēja atteikšanos no plāniem par Polijas karalistes pakļaušanu savai ietekmei, taču kā vājākajai no sabiedrotajām valstīm tai nebija citas izvēles, kā piekrist Vācijas pozīcijai. Austroungārija principā bija ar mieru pakļaut poļu armiju Vācijai, taču nevēlējās pilnībā atdot iniciatīvu Vācijas rokās, izstrādāja arī pati savus plānus, centās iesaistīt poļus pēc iespējas lielākā skaitā savos bruņotajos spēkos. Par Galīcijas tālāko likteni, tās attiecībām ar Polijas karalisti, kā arī politisko kontroli pār jauno valsti Vācijai un Austroungārijai bija jāvienojas tālākās sarunās. Precīzs jaunās poļu valsts proklamēšanas datums netika noteikts.

Pēc I.~Buriana apgalvojuma Vāciju uz tik krasu soli poļu jautājumā virzīja divi apstākļi: 1) izplatījušās baumas, ka Krievija gatavojas solīt poļiem visu to zemju apvienošanu un neatkarību; 2) jau minētā vēlme saformēt poļu karaspēku, ko nevarēja izdarīt līdz Polijas karaliste nebija ieguvusi atsevišķas valsts statusu.

T.~f.~Betmans Holvegs bez vācu militāristu vēlmes iegūt savā rīcībā ``sabiedroto'' poļu armiju, kā vācu pozīcijas mainās cēloņus nosauca arī cerību uz separātu mieru ar Krieviju izzušanu; Vācijas nevēlēšanos visu Polijas karalisti atdot Austroungārijai, tai skaitā baidoties, ka poļi tajā var iegūt pārlieku, vāciešiem nevēlamu svaru; arī sapratni, ka pēc kara nebūs iespējams saglabāt \latxti{statys quo}, tāpēc nepieciešams jau laicīgi kārtējo reizi sadalīt poļu tautu un tās teritorijas.

Vēsturnieki kā vēl vienu no pieņemtā lēmuma cēloņiem nosauc arī vācu centienus īstenot Vācijas kontrolētas t.~s. ``Viduseiropas'' (``\detxti{Mitteleuropa}'') izveides tieksmes. Tā bija jāveido uz blakus valstu teritoriju rēķina. Jaunajai ``Polijas valstij'', pēc šīs ģeopolitiskās koncepcijas īstenotāju viedokļa, bija jākļūst par pirmo ``patstāvīgo'' valsti Krievijas nomalē un jāveido ``buferis'' starp Krieviju un Vāciju. Pēc tam līdzīgas valstis varētu radīt Baltijā, Somijā un Ukrainā. Tas atbrīvotu Vāciju no rūpēm par līdzās esošo vareno kaimiņu~--- Krieviju un veicinātu vācu koloniālās politikas īstenošanu. Tāpēc arī Vācija bija gatava veicināt Krievijā apspiesto tautu separātismu.

Vācijas un Austroungārijas vadošo politiķu plāni neatbilda J.~Pilsudska un viņa piekritēju centieniem. Pēc Polijas karalistes nonākšanas Centrālo valstu rokās J.~Pilsudskis uzskatīja, ka poļiem nav jāpalīdz militāri Vācijai un Austroungārijai, ja tās nenosaka karalistes nākamo statusu. Viņš, pretēji ``Augstākās nacionālās komitejas'' lēmumam un arī viena no saviem tuvākajiem līdzbiedriem V.~Sikorska domām, izteicās pret tālāku poļu vervēšanu leģionā. J.~Pilsudskis centās izveidot arī jaunu, neatkarīgu no ``Augstākās nacionālās komitejas'' pārstāvniecību, kura varētu vest sarunas ar okupācijas varu. Jau 1915.~gada decembrī viņš radīja \strong{Centrālo nacionālo komiteju}, kuras sastāvā sākotnēji gan bija tikai PPS un dažu sīku politisku grupu pārstāvji. Tiesa, pēc divām nedēļām tai pievienojās no dažādām zemnieku--ļudoviešu grupām jaunizveidotā partija \pltxti{Wyzwolenje} (\pltxti{Polskie Stronnitctwo Ludowe ``Wyzwolenje''}~--- Poļu Zemnieku partija ``Atbrīvošanās'', pēdējais vārds partijas nosaukumā nāca no partijas preses orgāna nosaukuma), kurā sastapa sapratni J.~Pilsudska aicinājums uz cīņu pret Krievijas impēriju. Šis ļudoviešu spārns aktīvi piedalījās dažādu nelegālu politisku un pusmilitāru organizāciju darbībā, kuras faktiski vadīja J.~Pilsudskis un viņam tuvi stāvoši cilvēki. PSL ``\pltxti{Wyzwolenie}'' kļuva par vienu no lielākajām partijām bijušās Krievijas impērijas poļu teritorijās. Šī partija bija opozīcijā okupantiem, prasīja muižnieku zemes sadali (ar atlīdzību) zemnieku labā. Lai stiprinātu savas pozīcijas, J.~Pilsudskis meklēja kontaktus arī ar endekiem, taču bez rezultātiem.

Tai pat laikā J.~Pilsudskis vēl palika Poļu leģiona 1.~brigādes komandiera amatā. 1916. gada jūlijā viņš, vienojies ar leģiona pulku komandieriem, nosūtīja Austroungārijas Virspavēlniecībai prasību sniegt skaidrus paskaidrojumus Polijas jautājumā un pārformēt leģionu vienotā korpusā, kurš kļūtu par poļu armijas pamatu. Lai izdarītu spiedienu uz austriešiem viņš pēc divām nedēļām iesniedza atlūgumu. Septembrī tas tika apmierināts. J.~Pilsudska piekritēji vēlāk ir dažādi interpretējuši šo notikumu. Vieni to uzskatīja par pierādījumu, ka viņš bija nolēmis saraut sakarus ar Centrālajām valstīm, citi to uzskatīja par spiediena līdzekli, kurš pēc dažām nedēļām noveda pie 1916.~gada 5.~novembra Vācijas un Austroungārijas ķeizaru deklarācijas par ``neatkarīgas Polijas'' proklamēšanu (par to tiks runāts turpmāk). Taču pēc poļu vēsturnieku D. un T.~Nalenču viedokļa šiem apgalvojumiem ir maz kopīga ar patieso ainu. Kā jau teikts, J.~Pilsudska lūgums atbrīvot viņu no leģionu 1.brigādes komandiera pienākumiem bija iecerēts kā spiediena līdzeklis uz austriešiem. Tā apmierināšana viņam pašam bija pārsteigums. Neiztika bez Vācijas militāristu iejaukšanās, tiem bija savi poļu izmantošanas plāni, kuru īstenošanai, kā viņiem likās, J.~Pilsudskis tikai traucēja. Tā 1916.~gada septembrī neveiksmīgi beidzās J.~Pilsudska mēģinājumi panākt neatkarīgas Polijas valdības un arī armijas izveidi. 1916.~gada 19.~septembrī poļu leģioni tika pārdēvēti par poļu palīgkorpusu (vācu \detxti{Polnisches Hilfskorps}, poļu \pltxti{Polski Korpus Posiłkowy}). Komandēja to pulkvedis S.~Šeptickis un pakļāvās tas Austroungārijas pavēlniecībai. Pēc dažām ziņām 1916.~gada novembrī korpusā dienēja ap 20~000 vīru.

Pēc atvaļināšanas no poļu leģiona 1.~brigādes komandiera amata J.~Pilsudskis, izmantojot leģiona radīto slavu, varēja izvērst politisku darbību Varšavā. Pēc viņa norādījumiem leģionāru vidū sākās kampaņa, lūdzot tos atbrīvot no dienesta. Vienlaikus savu konspiratīvo darbību izvērsa J.~Pilsudska 1914.~gadā izveidotā un kontrolētā Poļu kara organizācijā (\pltxti{Polska Organizacja Wojskowa}). Tajā 1916.~gada beigās sastāvēja jau 11~tūkstošu cilvēku un tā turpināja augt. Visā bijušās Polijas karalistes teritorijā tika rīkoti pusmilitārās formās tērptu tās locekļu gājieni, mācības, pietuvinātas kara apstākļiem, dažkārt vācu virsnieku~--- novērotāju klātbūtnē. Organizācijai bija savi virsnieku kursi, kuru beidzējiem piešķīra militāras dienesta pakāpes. Tas viss piesaistīja organizācijai poļu, īpaši jauniešu~--- studentu un vecāko klašu skolnieku, uzmanību.

Galīgi Vācijas un Austroungārijas pozīcija jautājumā par Polijas karalistes nākotni tika saskaņota augstāko politisko un militāro vadītāju sanāksmē 1916.~gada 18.~oktobrī. Tika nolemts pēc kara radīt poļu valsti, bet līdz tam saglabāt ģenerālgubernatūras, kaut arī vācu militāristi uzstāja uz to apvienošanu, lai atvieglotu vervēšanu Polijas armijā.

1916.~gada 5.~novembrī abu okupācijas varu ģenerālgubernatori savu ķeizaru vārdā publicēja deklarāciju -- t.s. ``divu imperatoru aktu'' par \strong{``neatkarīgas'' Polijas kā mantojamas monarhijas} (poļu \pltxti{Regencyjne Królestwo Polskie}, vācu \detxti{Regentschaftskönigreich Polen}) \strong{proklamēšanu}. Tas paredzēja ``poļu valsts'' un poļu armijas izveidi no ``krievu kundzībai atrautajiem apgabaliem''. Bija teikts, ka šī valsts tiks veidota kā Centrālo valstu sabiedrotā ``ar mantojamu monarhiju un konstitucionālu iekārtu'', bet poļu armijas ``organizācija, apmācība un vadība tiks noregulēta, savstarpēji vienojoties.''. Arī valsts robežas netika paziņotas, to noteikšana atlikta uz nākotni.

Lēmums par Polijas karaļvalsts radīšanu bija rupjš starptautisko tiesību, kuras atļāva veikt teritoriālas pārmaiņas tikai miera konferencēs ar visu ieinteresēto pušu piedalīšanos, pārkāpums. Tomēr vācu historiogrāfijā Polijas valstiskums tiek skaitīts no 1916.~gada 5.~novembra, nevis 1918.~gada 11.~novembra, kā pēc 1926.~gada apvērsuma tika pieņemts poļu historiogrāfijā.

Pēc vācu-austriešu slepenas vienošanās ``neatkarība'' tika attiecināta tikai uz agrāk Krievijas sastāvā ietilpstošo Polijas daļu, jaunās ``valsts'' sastāvā bija jāietilpst tikai bijušās Polijas karalistes teritorijām. Austroungārijas imperators Francis Jozefs~I dienu pirms ``divu imperatoru akta'' pasludināšanas izdeva reskriptu par Galīcijas autonomijas paplašināšanu Austroungārijas sastāvā. Vācijai un Austroungārijai bija pat jāsaņem jaunas poļu zemes rietumos, bet jaunproklamētajai Polijai~--- kompensācija uz ukraiņu un baltkrievu apdzīvoto teritoriju rēķina. Par poļiem, kuri tā cerēja nākotnē paplašināt savu valsti, angļu vēsturnieks A.~Dž.~P.~Teilors rakstīja: ``\dots{}poļi līdz pašām beigām prasīja brīvību sev un pakļaušanos~--- no citiem; viņi gribēja atgūt monopola stāvokli Galīcijā, kas bija mazinājusies kara un [Krievijas] okupācijas apstākļos, tomēr lūdza Habsburgu armijas palīdzību, lai atņemtu Ukrainu nevarīgajai Krievijai''. Faktiski ``poļu valsti'' okupanti paredzēja veidot bez Galīcijas, Silēzijas, Lielpolijas, Pomoržes (poļu \pltxti{Pomorze Wschodnie}; vācu \detxti{Pommerellen}) un Mazūru apgabaliem, kuri paliktu Austroungārijas un Vācijas sastāvā. Šis projekts deva arī triecienu Polijas karalistes un Galīcijas apvienošanas vienā katoliskā Habsburgu monarhijā piekritējiem.

Bijušais Vācijas kanclers B.~f.~Bīlovs gan uzskatīja, ka Polijas neatkarības pasludināšana bija vislielākā Vācijas kļūda, jo traucēja sagraut Krievijas savienību ar Rietumiem. Par niecīgu skaitu poļu karavīru Vācija samaksājusi ar tās Krievijas, kura varēja būt Vācijas potenciālā sabiedrotā, atsvešinātību. Šodien ir skaidrs, ka ekskanclers pārvērtēja iespējas ķeizariskajai Vācijai izlīgt ar carisko Krieviju, germanofīlu lomu Krievijā, taču pareizi vērtēja to skaitu poļu vidū.

Arī daži padomju vēsturnieki rakstīja, ka paši Vācijas politiķi, pasludinot 1916.~gada 5.~novembrī ``neatkarīgas Polijas'' nodibināšanu, izsaukuši Krievijas vadības neapmierinātību un izjaukuši 1916.~gada jūlijā Stokholmā Vācijas valdības pārstāvju baņķiera M.~Varburga un rūpnieka H.~Stinesa ievadītās sarunas ar Krievijas Valsts domes vicepriekšsēdētāju A.~Protopopovu par separāta miera noslēgšanu. Domājams, ka šīs diplomātiskās zondāžas loma tika speciāli pārvērtēta, lai pēc iespējas vairāk nomelnotu Krievijas impērijas valdošos slāņus.

Sākotnēji J.~Pilsudskis un viņa piekritēji ar sajūsmu uzņēma 5.~novembra vācu-austriešu deklarāciju. J.~Pilsudskis to pat vērtēja kā ``savu ideju triumfu''. Nākamajai ``poļu valstij'' taču nebija jābūt pat saistītai ar Austroungāriju, tai tika solīta brīvība. Sakarā ar tās proklamēšanu 12.~novembrī sapulcējušies visu brigāžu leģionāru un Galīcijas poļu sabiedrības pārstāvji nosūtīja J.~Pilsudskim apsveikumu, kurā godināja ``leģionu dibinātāju'' un izteica cerību, ka ``tas, kurš ar savu dižo garu radīja leģionus, organizēs Poļu armiju, kas novedīs pie brīvas poļu tautas atdzimšanas''. J.~Pilsudskis un viņa piekritēji cerēja panākt poļu valdības izveidi un poļu karavīru nodošanu tās pārziņā.

Taču cerības nepiepildījās. Poļu vidū psiholoģiski ļoti nelabvēlīgu iespaidu radīja tas, ka jau dažas dienas pēc Polijas ``neatkarības'' proklamēšanas~--- 9.~novembrī Centrālās valstis izplatīja poļiem domātu uzsaukumu: ``Vēl cīņa pret Krieviju nav beigusies. Tāpēc iestājaties brīvprātīgi mūsu armijā, lai palīdzētu pabeigt mūsu uzvaru pār jūsu apspiedēju.'' Decembrī tika izveidota t.s. ``poļu armija'' (poļu \pltxti{Polska Siła Zbrojna}, vācu \detxti{Polnische Wehrmacht}). Tai bija jānēsā vācu uniforma, jābūt Vācijas pavēlniecībā un jādalās galīciešu un ``kongresa Polijas'' vienībās. Tiesa, veidojās ``poļu armija'' gaužām lēni. Literatūrā sastopami dati, ka līdz 1917.~gada janvāra beigām vācu kontrolētajā ``poļu'' karaspēkā kā brīvprātīgie bija pieteikušies tikai ap divi tūkstoši cilvēku, bet kopumā tajā bijis 5 līdz 6 tūkstoši vīru.

Visai lēni ar okupācijas varu rīkojumiem tika dibinātas jaunās ``valsts'' iestādes. Atbildot uz poļu prasību ``bez Polijas valdības~--- nekādas poļu armijas'', 1916.~gada 6.~decembrī tika pasludināts par 25-vietīgas ``Pagaidu valsts padomes'' (\pltxti{Tymczasowa Rada Stanu}) ar padomdevēja funkcijām dibināšanu, kuru nozīmēja vāciešu un austroungāru ģenerālgubernatori. (15~pārstāvji bija no Vācijas okupācijas zonas un 10~no Austroungārijas zonas.) Starp citu, Padomē ietilpa arī Daugavpilī dzimušais provāciski noskaņotais poļu politiķis un publicists V.~Studņickis.

Šādos apstākļos J.~Pilsudskis piedāvāja vācu Varšavas ģenerālgubernatoram ģenerālim H.~f.~Bezeleram savus pakalpojumus. 5.~decembrī pēdējais saņēma J.~Pilsudska sūtījumu, kurā viņš izklāstīja savu redzējumu par brīvprātīgo vervēšanu ``poļu armijā''. ``Komandants'' deva saprast, ka bez viņa okupantiem neizdosies radīt poļu armiju. Viņš rakstīja: ``Poļi nav organizēta tauta, kuru var vadīt racionālā ceļā, \latxti{via Vernuft} [caur saprātu]. Tā ir plūstošu smilšu kaudze, kura kustas tikai tad, kad izdodas prasmīgi organizēt stipru vēju vajadzīgajā virzienā \citespace{} Tagad viss ir atkarīgs no tā vai izdosies, izmantojot 5.~novembra faktu, prasmīgi organizēt stipru vēju, kurš \citespace{} padarīs armiju un karadienestu populāru prātos un sirdīs''. Iepazinies ar vēstījumu, vācu ģenerālgubernators 8.~decembrī ielūdza J.~Pilsudski ierasties Varšavā, taču vēl nepiedāvāja nekādu amatu Polijas varas institūtos.

12.~decembrī daudzie J.~Pilsudska piekritēji ar sajūsmu sagaidīja viņa ierašanos Varšavā. Kad viņš, iznākot no stacijas, sēdās ekipāžā, jaunieši izjūdza zirgus un, ieņēmusi to vietu, aizveda savu elku uz viesnīcu. J.~Pilsudskis tikās ar ievērojamākajiem poļu politiskajiem darbiniekiem, viņu pieņēma arī H.~f.~Bezelers. Taču divas stundas ilgās sarunas nedeva konkrētus rezultātus. Domstarpības pastāvēja galvenajā jautājumā. J.~Pilsudskis vēlējās, lai armijas formēšana tiktu uzticēta Pagaidu valsts padomei, bet H.~f.~Bezelers uzskatīja, ka tā ir okupācijas varas iestāžu, bet nevis poļu kompetences joma. J.~Pilsudskis vācu ģenerālgubernatoram bija vajadzīgs, lai ar savu autoritāti pievilinātu armijā pēc iespējas vairāk brīvprātīgos. Rezultātā J.~Pilsudskis Pagaidu valsts padomē ieņēma kara komisijas priekšsēdētāja vietu, taču padomei nebija nekādu pilnvaru. Tās sastāvā iegājušie J.~Pilsudska atbalstītāji bija mazākumā. Visi poļu pūliņi organizēt viņu pakļautībā esošu armiju, iegūt savās rokās dažas valsts pārvaldes funkcijas, atvieglot poļu iedzīvotāju sakarus starp dažādām okupācijas zonām, beidzot, ievēlēt Polijas karalistes reģentu (no latīņu \latxti{regens}~--- valdošais, pagaidu valsts galvas funkciju izpildītājs, kad monarha vai nu nav vai tas ir nepilngadīgs vai rīcībnespējīgs) un izveidot valdību nenesa panākumus.

Kaut Vācija bija ieinteresēta pēc iespējas lielāka poļu karavīru skaita izmantošanā, vācu militāristi nevēlējās ne dzirdēt par patiesu Polijas autonomiju. T.s. panģermāņi (Vācijā 1891.--1939.~gados pastāvošā Panģermāņu savienība~--- \detxti{Alldeutscher Verband}~--- bija šovinistiska organizācija, līdz 1918.~gadam ļoti aktīva, iestājās par Vācijas pasaules kundzību, virknes apgabalu atņemšanu Krievijai) loloja plānus par poļu izsūtīšanu no Vācijas iekarotajām teritorijām. Vietējie vācieši plānus par poļu valsts izveidi Vācijas aizbildniecībā akceptēja, vienīgi t.s. vācu aktīvisti bija vīlušies, jo vēlējās panākt novada ciešāku sasaisti ar Vāciju. Vācijas armijas Ģenerālštāba priekšnieks P.~f.~Hindenburgs, bet īpaši viņa vietnieks Ē.~Lūdendorfs vēlējās lai poļu daļas zvēr uzticību Vācijas ķeizaram, nevis Polijas valstij. Vācijas kanclers T.~f.~Betmans Holvegs gan nebija sajūsmā par šādu soli, taču neiedrošinājās iebilst abiem toreiz populārajiem militāristiem.

Poļiem drīz kļuva skaidrs, ka īstu neatkarību Vācija Polijai nedos. Lozungs ``Bez Polijas valdības~--- nekādas poļu armijas'' ietekmēja ne tikai poļu ``pasīvistus''. Vācijā pirms kara realizētā poļu zemes īpašumu atsavināšanas politika iespaidoja poļu attieksmi ne tikai pret Vāciju, bet arī pret Centrālo valstu bloku kopumā. Izrādījās, ka Vācijas antipoliskā pirmskara politika arī kara laikā sarežģīja tās attiecības ar poļiem.

1917.~gada aprīlī austrieši nodeva poļu palīgkorpusu Pagaidu valsts padomes, bet faktiski vācu ģenerāļa H.~f.~Bezelera pakļautībā. Taču vācu imperiālistu cerības iegūt ievērojamu skaitu poļu karavīru, panākt poļu brīvprātīgo plašu dalību cīņā pret Krieviju nepiepildījās. Pretēji Berlīnes un Vīnes cerībām Polijas iedzīvotāji nesteidzās brīvprātīgi karot par Centrālo valstu interesēm.

Vāciešu formējamās poļu daļas veidojās ļoti lēni, un kad beidzot tās varēja sūtīt uz fronti, izrādījās, ka tām nav ar ko karot, jo Krievijas armija sabruka. Doma par to, ka tās varētu nosūtīt uz Rietumu fronti pret francūžiem, bija utopiska.

Jau 1917.~gada pavasarī Polijā parādījās pirmās atklātās neapmierinātības pazīmes. No Polijas valstiskās patstāvības pasludināšanas brīža poļu aprindās gan pašā Polijā, gan Krievijā pārsvaru sāka gūt viedoklis, ka autonomija Poliju neapmierinās un kā mērķis ir jāizvirza neatkarīgas poļu valsts radīšana. Saprotams, ka ``Polijas neatkarību'' atzina tikai Vācijas un Austroungārijas impērijas. Pirmā uz notikumiem Polijā jutās spiesta reaģēt Krievijas valdība. Tās paziņojumā 1916.~gada 2.~(15).novembrī bija norādīts, ka centrālo valstu gubernatoru akts neievēro starptautiskās tiesības, tā mērķis ir nevis neatkarīgas Polijas valsts radīšana, bet gan rekrūšu iesaukšana armijā. Tika vēlreiz deklarēts, ka Krievija vēlas radīt ``vienotu Poliju no visām poļu zemēm, piešķirot tai pēc kara beigām tiesības brīvi noteikt savu nacionālo, kultūras un saimniecisko dzīvi uz autonomijas pamatiem zem krievu valdnieku sceptera un saglabājot vienotu valstiskumu''.

Gatavojoties 1917.~gada uzbrukumam Polijas teritorijā, gada nogalē Krievija vēlreiz atgriezās pie Polijas statusa. 1916.~gada 12.~(pēc jaunā stila~--- 25.)~decembrī, pēc B.~Štjurmera atstādināšanas no valdības vadītāja amata Nikolajs~II izdeva armijai un flotei pavēli N 870, kurā starp citiem Krievijas uzdevumiem karā nosauca arī ``brīvas Polijas radīšanu no visiem tās līdz šim atšķirtajiem apgabaliem'', tātad~--- Polijas karalistes apvienošanu ar Vācijas un Austroungārijas pakļautībā esošajām teritorijām. Tika arī solīts, ka Polija iegūs savu valstisku pārvaldi ar savām likumdošanas iestādēm un savu armiju. Tādā kārtā imperators apstiprināja lielkņaza Nikolaja Nikolajeviča 1914.~gada augusta manifestā izteikto solījumu. Var arī teikt, ka beidzot Krievijas augstākajās pārvaldes sfērās bija uzvarējusi pozīcija par nepieciešamību piešķirt Polijas karalistei autonomiju un pievienot tai citas poļu zemes, kuru aizstāvēja R.~Dmovskis ar saviem piekritējiem. Uz īsu laiku par Ministru padomes priekšsēdētāju kļuvušais konservatīvais A.~Trepovs paziņoja par nodomu izveidot brīvu Poliju personālā ūnijā ar Krieviju. Tas jau gāja pat tālāk par agrākajiem autonomijas solījumiem. Vairāki citi krievu konservatīvie politiķi uzstāja, ka, ja nāksies piešķirt poļiem nacionālo un pilsonisko brīvību, tad Poliju vajag pilnīgi atdalīt no Krievijas, citādi Polijas nestabilitāte apdraudēs arī Krieviju. Kā norādījusi jau minētā krievu vēsturniece A.~Bahturina, autonomas Polijas Krievijas sastāvā radīšanas perspektīva nozīmētu ne tikai Polijas statusa maiņu, bet arī centra un nomaļu attiecību izmaiņu. Tas varēja kļūt par sākumu Krievijas nomaļu autonomizācijai.

Pēdējais Krievijas impērijas valdības vadītājs kņazs N.~Goļicins 1917.~gada 8., 9.~un 12.~februārī noturēja vairākas apspriedes Polijas nākamās iekārtas pamatu izstrādāšanai, kurās tika izskatīti divi varianti: Polijas valsts saistīta ar Krieviju personālā ūnijā, bet neatkarīga iekšējās lietās (kā Ungārija Austroungārijas vai Somija Krievijas sastāvā), vai no Krievijas atdalījusies pilnīgi neatkarīga valsts. Militāristi~--- ģenerāļi M.~Aļeksejevs, A.~Brusilovs u.c.~--- centās pierādīt, ka autonomija poļu zemēs vairs nevienu neapmierinās. Atbildei skanēja galmam tuvo birokrātu viedoklis, ka piekāpšanās Polijai vedīs arī pie Baltijas un Kaukāza atsvešināšanās no Krievijas. Tomēr 17.~gada 12.~februārī Ministru Padome nolēma, ka pilnīga Polijas neatkarība no Krievijas būtu vistaisnīgākais un vislietderīgākais poļu jautājuma risinājums. Vairāki ministri balsoja par šādu risinājumu tādēļ, ka, ja Polija saņemtu autonomiju, tas būtu pretrunā ar Krievijas valstiskuma pamatprincipu. Nikolajs~II vai nu nevēlējās vai nepaguva apstiprināt šo lēmumu. Lai kā arī nebūtu, tikai pēdējā Krievijas impērijas valdība ārkārtējos kara apstākļos piekrita neatkarīgas Polijas valsts atjaunošanai, un nevar zināt, vai šis lēmums tiktu realizēts, ja Krievijas impērija turpinātu pastāvēt un būtu kara uzvarētāju vidū. Ievērojot carisma piedzīvotās militārās neveiksmes, ir pamats apšaubīt Ministru padomes apņēmību īstenot lēmumu. Tam bija propagandas raksturs, īpaši ņemot vērā gatavojamo Krievijas armijas uzbrukumu 1917.~gadā. Tomēr jākonstatē, ka pretstatā kara sākumposmā Krievijas valdībā valdošajam viedoklim par nepieciešamību stiprināt nomaļu, arī Polijas karalistes, saistību ar lielkrievisko centru, kara gaitā arvien vairāk dažādos krievu sabiedriskajos slāņos ieskanējās doma par Krievijas konsolidāciju uz tīri nacionāla pamata. Kņaza N.~Goļicina vadīto apspriežu lēmums liecināja, ka arī valdībā tā bija guvusi atbalstu. Neatkarīgas Polijas noformēšana likvidētu poļu jautājumu kā tādu un neradītu autonoma veidojuma pastāvēšanas precedentu Krievijas impērijā. Pie tam cara valdība plānoja neatkarīgās Polijas ciešu saistību ar Krieviju. Pamatuzsvars tika likts uz ekonomiskās vienotības saglabāšanu, kura varēja daudz stiprāk nekā tas vai cits tiesiskais statuss nodrošināt Krievijas un Polijas saistību. Taču tam visam vairs nebija tiešas ietekmes uz Vācijas un Austroungārijas okupēto Poliju.

Kā personīgā sarakstē S.~Sazonovam rakstīja Krievijas sūtnis Francijā A.~Izvoļskis, ``Polijas jautājums sāk izslīdēt no Krievijas kontroles, iegūst plašāku un patstāvīgāku starptautisku jautājumu, to var atrisināt tikai ar sabiedroto valstu līdzdalību''. Tā arī notika.

Ar Vācijas un Austroungārijas ``Polijas neatkarības'' atzīšanu Polijas jautājums atkal tika izvirzīts kā starptautiskas uzmanības vērta tēma. Vācijas un Austroungārijas paziņojums par valstiskuma piešķiršanu Polijai, saistīts ar poļu armijas uzstādīšanas mēģinājumu pret Antanti, izsauca 1916.~gada 5.~(18.)~novembrī Francijas, Itālijas un Lielbritānijas pārstāvju apspriedē izstrādātu protestu pret šo Centrālo valstu aktu. Protestā bija uzsvērts, ka ``militāras okupācijas fakts \citespace{} nevar izsaukt suverēnu tiesību uz okupēto teritoriju iegūšanu, un attiecīgi dot kādam tiesības piešķirt kādam šīs teritorijas.'' Taču dokumentā nekas nebija teikts par Polijas valstiskumu.

ASV prezidents V.~Vilsons, cenšoties uzstāties starpnieka lomā starp Centrālajām valstīm un Antanti, lūdza abām pusēm paskaidrojumus par to mērķiem karā. Antantes atbilde V.~Vilsonam kļuva par pirmo oficiālo sabiedroto dokumentu, kurš rādīja, ka Polijas jautājuma noregulēšana pēc kara kļūs par svarīgu problēmu. Pēc Krievijas prasības tajā bija pateikts, ka Krievijas imperatora nodomi tikuši ``skaidri izklāstīti pavēlē, kuru viņš izdeva savām armijām'' [1916.~gada 12.~(25.)~decembrī]. Taču 1917.~gada 22.~janvārī izcilais poļu pianists un komponists I.~Paderevskis, atrazdamies ASV, izklāstīja ASV prezidentam V.~Vilsonam, kāpēc kara rezultātā būtu jārodas ``vienotai, neatkarīgai un autonomai Polijai''. Arī V.~Vilsona atbildes vēstījumā Antantes valstīm 1917.~gada janvārī bija teikts, ka ir jātop ``vienotai, neatkarīgai, autonomai Polijai''.

Tomēr neilgi pirms Februāra revolūcijas Krievijā 1917.~gada 14.~februārī Krievija un Francija noslēdza slepenu vienošanos, kurā bija paredzēts, ka Krievija sekmēs Francijas mērķus savas pēckara austrumu robežas noteikšanā, bet savukārt Francija palīdzēs Krievijai panākt sev labvēlīgu rietumu robežu. Faktiski ar to Francija, kura ne vienu reizi vien bija izteikusi savas simpātijas poļu nacionālās atbrīvošanās centieniem, pilnībā piekrita Krievijas patvaldības iekarošanas plāniem. Taču tai pat laikā Francijas valdība nevairījās dāvāt zināmas cerības arī poļu aprindām Krievijā un Rietumeiropā.

Kad \strong{Krievijā} notika \strong{Februāra revolūcija}, situācija mainījās arī Polijā. Krievijas Pagaidu valdība bija orientēta atlikt jautājumu par nacionālo autonomiju valstī dzīvojošajām mazākumtautām līdz Satversmes sapulces sasaukšanai. Taču 14./27.~martā Petrogradas strādnieku un kareivju deputātu padome griezās pie poļu tautas ar uzsaukumu, kurā bija teikts: ``Cariskais režīms, kurš pusotra gadsimta garumā līdz ar krieviem apspieda arī poļu tautu, apvienotiem proletariāta un armijas spēkiem ir gāzts. Ziņojot poļiem par brīvības uzvaru pār viskrievijas žandarmu, Petrogradas strādnieku un kareivju deputātu padome paziņo, ka demokrātiskā Krievija stāv uz tautu nacionāli-politiskās pašnoteikšanās pamata, un proklamē, ka Polijai ir tiesības būt pilnībā neatkarīgai valsts-starptautiskajās attiecībās.'' Tiesa, tuvākus paskaidrojumus, kā jātiek veidotai šai valstij, paziņojums nesniedza, bet no tikko sanākušās iestādes to arī objektīvi nevarēja gaidīt. Interesanti, ka PPS Centrālā strādnieku komiteja (\pltxti{Centralny Komitet Robotniczy PPS}) sakarā ar Petrogradas strādnieku un kareivju deputātu padomes uzsaukumu izvirzīja prasību lai Krievija atsakās ne tikai no poļu, bet arī lietuviešu zemēm.

Bija spiesta reaģēt arī Krievijas Pagaidu valdība. Pēc divām dienām tā kā vienu no pirmajiem saviem aktiem 1917.~gada 17.~martā nāca klajā ar ``Uzsaukumu poļiem'' (``\rutxti{Воззвание к полякам}''). Tajā tika pasludināts, ka Polijai jākļūst par ``brīvu un neatkarīgu'' valsti, pilntiesīgu starptautisko attiecību subjektu. Tiesa gan, tam bija jānotiek nevis pēc pašas poļu tautas, bet pēc Viskrievijas Satversmes sapulces lēmuma. Krievijas Pagaidu valdība aicināja poļus ``vienotas Polijas atbrīvošanas un neatkarības vārdā'' plecu pie pleca ar krievu armiju cīnīties pret kopējo ienaidnieku~--- vācu iekarotājiem. Taču tālāk uzsaukumā bija teikts, ka Polijai jāpaliek ``brīvā mūžīgā vienotā militārā savienībā ar Krieviju'' (``\rutxti{в свободном вечном военном союзе с Россией}''). Tātad Polija būtu militārā ziņā atkarīga no Krievijas. Šis Polijas ``neatkarības'' variants pēc būtības līdzinājās pirms dažiem mēnešiem~--- 1916.~gada 5.~novembrī~--- publicētajai vācu un austroungāru deklarācijai par ``neatkarīgās'' Polijas monarhijas proklamēšanu. Krievijas valdošie slāņi nevēlējās pieļaut iespēju, ka pie Krievijas rietumu robežām atrastos tai naidīga poļu valsts. ``Brīvo mūžīgo vienoto militāro savienību'' bija paredzēts noslēgt līdz ar ierobežotās ``neatkarības'' piešķiršanu. Tā, Lielbritānijas sūtnis Krievijā Dž.~Bjukinens savos memuāros liecināja, ka Pagaidu valdības ārlietu ministrs P.~Miļukovs šajā laikā, pretēji viņa apgalvojumiem pēc Oktobra apvērsuma, sliecās Polijai piešķirt nevis neatkarību, bet gan tikai autonomiju. Liela daļa Rietumu un krievu vēsturnieku ir vienisprātis, ka Pagaidu valdība piekrita piešķirt Polijai neatkarību tāpēc, ka Vācija un Austroungāriju bija ieņēmušas poļu apdzīvotās teritorijas un bija nepieciešams pretoties to vēlmei nostiprināt savu ietekmi tajās. Pagaidu valdība darbojās taktisku mērķu vadīta, politiskās lietderības vārdā, cerot, ka Polijas neatkarības atzīšana stimulēs poļus cīņai pret vācu un austroungāru iebrucējiem.

Jautājumā par Polijas robežām Pagaidu valdība turpināja Krievijas impērijas politisko kursu uz vienotas Polijas radīšanu, kurā ietilptu līdz tam gan Krievijai, gan Vācijai un Austroungārijai piederējušās poļu teritorijas. Minētajā ``Uzsaukumā'' bija teikts, ka Polijas robežu ar Vāciju un Austroungāriju vai valstīm~--- tās mantiniecēm jānosaka starptautiskai miera konferencei, bet robeža starp Poliju un Krieviju tika saukta par ``etnisku'', tās guberņas un apriņķi, kurās bija poļu iedzīvotāju vairākums, nonāktu Polijā, bet kurās poļi būtu mazākumā~--- paliktu Krievijā. Šo otro robežu bija jānosaka Krievijas un Polijas Satversmes sapulcēm. Daudznacionālo vietējo iedzīvotāju vēlmju noskaidrošana nebija paredzēta. Tika atrunātas vienīgi nacionālo mazākumu tiesības. Pagaidu valdības ``Uzsaukumā poļiem'' bija teikts: ``Krievija tic, ka gadsimtiem ilgā kopīgā dzīvē saistītās ar Poliju tautas saņems \citespace{} savas pilsoniskās un nacionālās pastāvēšanas stingru nodrošinājumu''. To gan varēja tulkot arī kā iejaukšanos ``neatkarīgas'' valsts iekšējās lietās. Poļu vēsturnieks Š.~Aškenazi, salīdzinot Krievijas Pagaidu valdības un Petrogradas strādnieku un kareivju deputātu padomes dokumentus par iespējamo Polijas valstiskuma atjaunošanu, secināja, ka pirmā to saistīja ar atkarības no Krievijas saglabāšanu, turpretī otrā nekādus noteikumus neizvirzīja.

Pirmo reizi pēc 1863.~gada poļu sacelšanās Polijas jautājums atkal kļuva starptautisks. Krievu vēsturnieks G.~Matvejevs uzsvēris, ka Februāra revolūcija mainīja arī Polijas jautājuma starptautisko skanējumu. Polijas karalistes nākamais statuss principā bija skaidrs. Antantes valstu uzvaras gadījumā šeit rastos iekšējo jautājumu risināšanā (autonoma) valsts, tāpat kā tas notiktu, ja kara laime uzsmaidītu Centrālajām valstīm. Taču vācu un austriešu rokās esošo Polijas daļu nākotne kļuva par atklātu jautājumu, jo Krievija atteicās no pretenzijām uz tām, kaut sabiedrotie pēdējo reizi bija apstiprinājuši to atzīšanu vēl 1917.~gada februārī.

Savukārt J.~Pilsudska biogrāfi poļu vēsturnieki D. un T.~Nalenči ir atzīmējuši, ka pēc Februāra revolūcijas viņa iepriekšējā politiskā līnija zaudēja jēgu. Tās galvenais saturs taču bija Polijas neatkarības atjaunošana cīņā pret Krieviju. Šim mērķim bija pakļauta armijas formēšanas programma. Pretkrieviska nostāja bija arī galvenais priekšnoteikums savienībai ar Austroungāriju un Vāciju. Carisma gāšana grāva šīs koncepcijas pamatus, jo revolucionārā Krievija atzina poļu tiesības uz savu valsti, solīja tiem vairāk nekā imperatoru valdības Berlīnē un Vīnē.

Šai sakarā jāatzīmē, ka vācu un austriešu okupācijas varu nodibinātā ``Pagaidu valsts padome'' (\pltxti{Tymczasowa Rada Stanu}), reaģējot uz Krievijas pagaidu valdības uzsaukumu, speciālā deklarācijā 1917.~gada 16.~(29.)~martā paziņoja, ka Polijas neatkarība, ko tikai tagad atzinusi Pagaidu valdība, esot jau iegūta 1916.~gada 5.~novembrī ar abu okupētājvalstu gubernatoru aktu, kaut arī šis akts ``vēl nenosprauda'' valsts teritoriālās robežas. Krievijas Pagaidu valdībai tika pārmests, ka tā ``ierobežo nācijas'' gribu, uzspiežot ``vienotu militāru savienību'', un nepareizi traktē teritoriālo jautājumu, piešķirot poļiem teritorijas, ``uz kurām neizplatās tās vara''. Ja pirmajam pārmetumam var piekrist, tad par otro gan jāsaka, ka pati ``Pagaidu valsts padome'' neievēroja pirms kara jau atzītās starptautiskās tiesības, kas noteica, ka daļēja vai pilnīga valsts okupācija nenozīmē tās suverenitātes un valstiskuma zaudēšanu, kā arī okupētās teritorijas valstiskās piederības zaudēšanu. Tā kā ``Pagaidu valsts padome'' savu jurisdikciju izplatīja tikai uz Krievijas impērijas poļu zemēm, kuras nepiederēja ne Vācijai, ne Austroungārijai, tad par šo (bijušās Krievijas impērijas) zemju piederību no starptautisko tiesību viedokļa varēja lemt Krievijas Satversmes sapulce. Tiesa, ja jautājums tiktu izšķirts demokrātiski, pirmkārt bija jāievēro vietējo iedzīvotāju vēlmes.

Boļševiku pozīcija Polijas jautājumā šai laikā bija vēl noteiktāka, sasaistot kopā nacionālās atbrīvošanās un strādnieku kustības uzdevumus. KSDS(b)P VII konferences rezolūcijā ``Sakarā ar Borgbjerga [Frederiks Borgbjergs~--- dāņu sociālistu līderis.~--- V.Š.] priekšlikumu'' bija teikts: ``Kas attiecas uz Polijas neatkarības atjaunošanu, tad tā ir krāpniecība kā no vācu-austriešu kapitālistu, tā arī krievu Pagaidu valdības puses, kura runā par it kā ``brīvu'' Polijas un Krievijas militāru savienību. Jo, lai tiešām noteiktu tautu gribu, visos anektētajos apgabalos nepieciešams izvest karaspēku un veikt brīvu iedzīvotāju aptauju. Tikai šāda līdzekļa piemērošana visai Polijai (t.i. ne tikai krievu, bet arī vācu un austriešu sagrābtajai) \citespace{} būtu solis valdības doto solījumu īstenošanā''. (Uzreiz gan jāpasaka, ka 1920.~gadā, kad liela daļa Polijas bija viņu varā, boļševiki tur iedzīvotāju aptaujas bez karaspēka klātbūtnes netaisījās rīkot.)

Pēc Februāra revolūcijas uzvaras Krievijā tika sperti arī praktiski soļi Polijas jautājuma risināšanā. Petrogradā no krievu un poļu sabiedriskajiem darbiniekiem tika radīta \strong{Likvidācijas komisija} (\pltxti{Komisja Likwidacyjna do Spraw Królestwa Polskiego}), kurai bija jānodarbojas ar jau evakuēto Krievijas valsts iestāžu Polijas karalistē likvidāciju, kā arī ar uz Krieviju izvesto Polijas materiālo un kultūras vērtību atdošanu poļu tautas īpašumā. Jautājums bija svarīgs ne tikai Vācijas un Austroungārijas okupētājos apgabalos palikušajiem, bet arī Krievijas teritorijās izkaisītajiem poļiem, jo 1917.~gada sākumā Krievijas pretinieku neokupētajās teritorijās pēc dažādiem aprēķiniem atradās no 2,5 līdz 4~miljoniem poļu. Tai skaitā bija ap 1,5~miljoni evakuēto no poļu zemēm, no 300 līdz 700~tūkstošu mobilizēto armijā, līdz 100 tūkstošu karagūstekņu no Vācijas un Austroungārijas pavalstnieku~--- poļu skaita. (Pēc krievu vēsturnieka N.~Rajska datiem 1917.~gada 1.~janvārī Krievijas guberņās bija 3~112~180 bēgļu no Polijas karalistes).

Poļu politiskie darbinieki Krievijā kā vieni no pirmajiem nacionālo mazākumu pārstāvjiem stājās sakaros ar Krievijas jaunajiem vadītājiem, stimulējot demokrātisko reformu attīstību.

Likvidācijas komisija strādāja Ziemas pilī Petrogradā. Tās atklāšanā ar runām uzstājās Pagaidu valdības vadītājs kņazs G.~Ļvovs un ārlietu ministrs P.~Miļukovs. Komisija aktīvi darbojās līdz Oktobra apvērsumam, kaut nomināli pastāvēja vēl līdz 1920.~gada augustam (lai gan 1919.~gada janvārī tika izdots KSFPR Tautas komisāru padomes dekrēts par komisijas likvidēšanu, nododot tās lietas un mantu KSFPR Ārlietu tautas komisariātam). Tās dalībnieki no Krievijas puses panāca poļu pārstāvju piekrišanu Holmas guberņas un Lietuvas ietilpšanai Krievijā, Polijas proporcionālas līdzatbildības par Krievijas impērijas parādiem atzīšanu. Lielu uzmanību sarunām veltīja Austroungārijā dzīvojošo citu slāvu tautu politiķi, kuri cerēja panākt savu tautu neatkarību uz tādiem pat principiem. Savukārt Vācijas un Austroungārijas valstsvīri bija norūpējušies par iespējamā Polijas un Krievijas noregulējuma atskaņām savās pārvaldītajās teritorijās.

Krievijā tika risināti arī citi praktiski jautājumi. Pagaidu valdības kara un jūras lietu ministrs A.~Gučkovs nodarbojās ar Polijas ``brīvās mūžīgās vienotās militārās savienības ar Krieviju'' praktisko realizāciju. Viņš uzskatīja, ka savlaicīga Krievijas valstisko struktūru dalība Polijas armijas radīšanā nodrošinās nākotnē Krievijai attiecīgu ietekmi gan uz šo armiju, gan Polijas valsti. Vēl 1914.~gada oktobrī Pulavā radīto poļu leģionu (brigādi) 1917.~gada janvārī pārformēja par strēlnieku divīziju un poļu ulānu pulku, bet ar no 1917.~gada jūlija par Krievijas armijas virspavēlnieku ieceltā ģenerāļa L.~Korņilova pavēli tai pat jūlijā par 1.~poļu korpusu ar ģenerāli J.~Dovboru-Musņicki priekšgalā. Daļa krievu politiķu bija pret tādu pārformēšanu, uzskatot, ka tā var novest Krievijas armiju pie sabrukuma. Vēlākais kara un jūras lietu ministrs A.~Kerenskis mēģināja kavēt poļu korpusa veidošanu, tomēr 1917.~gada beigās tajā bija jau 12 strēlnieku pulki, 3~ulānu pulki, artilērija un pat lidotāju vienība ar 24~lidmašīnām. Korpusa formēšanai atbalstu sniedza ne tikai ģenerālis L.~Korņilovs, bet arī Rietumu frontes pavēlnieks A.~Deņikins. (Vēlāk J.~Dovbors-Musņickis savukārt atlīdzināja viņiem, palīdzot bēgt no ieslodzījuma Bihovā pēc t.s. korņiloviādes neveiksmes.) Poļu karavīri arī revolucionāro notikumu laikā saglabāja disciplīnu (Iestājoties korpusā, bija jāparaksta solījums pakļauties visām komandieru pavēlēm). J.~Dovbora-Musņicka komandētajā 1.~poļu korpusā dienēja gan poļi, gan baltkrievi~--- katoļi. Nacionāldemokrāti ar sajūsmu apsveica poļu karaspēka daļu veidošanu, kreisie, kuri uzstājās pret atsevišķu nacionālu armijas vienību radīšanu, skatījās uz tām ar bažām, ja ne pilnībā noraidoši.

Antantes lielvalstis, kuras arī pēc Februāra revolūcijas bija ieinteresētas saglabāt Krieviju kā sabiedroto karā, oficiāli atbalstīja Pagaidu valdības politiku Polijas jautājumā, līdz ar to arī plānojamo tās atkarību no Krievijas. Taču tagad Krievijas valdības pozīcija Polijas jautājumā deva Rietumu sabiedrotajām lielākas nekā agrāk manevra iespējas. Tās vairs nesaistīja solījumi cara valdībai uzskatīt poļu jautājumu par Krievijas iekšēju lietu. Radās apstākļi starptautiskai šī jautājuma apspriešanai. Krievijas sabiedrotās Lielbritānija, Francija un Itālija 30.~aprīlī kopējā notā izteica tai savu solidaritāti un gatavību atzīt Polijas neatkarību, ja to atzīst Krievijas valdība. Antantes valstu valdības bija norūpējušās par situāciju Krievijā, tāpēc sāka aktīvāk izvirzīt arī savas iniciatīvas poļu jautājumā.

Pēc Februāra revolūcijas Krievijā lielāku aktivitāti parādīja Francija, vēloties izmantot poļu spēkus karā savā labā. Sakarā ar minētās Likvidācijas komisijas radīšanu Krievijā 1917.~gada 4.~jūnijā tika pieņemts Krievijas puses priekšlikums par poļu armijas radīšanu Francijā un tās prezidents R.~Puankare parakstīja dekrētu par ``poļu autonomās armijas'' formēšanu no gūstekņiem, vēlāk arī no poļu diasporas pārstāvjiem pasaulē (ASV vien dzīvoja ap 3~miljoni poļu) un bēgļiem no Krievijas.

Krievijas Pagaidu valdības 1917.~gada martā pasludinātā Polijas neatkarība radīja bažas Vācijā un Austroungārijā. Poļu deputāti šo valstu parlamentos tāpat kā Poļu Nacionālā komiteja Parīzē izvirzīja vienotas, neatkarīgas un neitrālas Polijas prasību.

1917.~gada martā Austroungārija ar slepenu protokolu atsacījās par labu Vācijai no savām pretenzijām uz Polijas karalisti, apmaiņai saņemot lielu daļu Rumānijas. Maijā Vācijas kanclers T.~Betmans-Holvegs un Austroungārijas ārlietu ministrs O.~Čerņins parakstīja vienošanos par ietekmes sfēru sadali. Vācijai pēc tās bija tiesības pievienot sev Lietuvu un Kurzemi un pilnībā pakļaut ``poļu valsti''. Austroungārija atteicās no ietekmes ``poļu valstī'' un solīja koncentrēt savus pūliņus dienvidu virzienā (Taču vēl 1917.~gada rudenī Austroungārijas valdošie slāņi tomēr cerēja realizēt austroungāriski--polisko risinājumu).

Poļu attieksme pret Februāra revolūciju bija dažāda. Nostāja pret okupantiem~--- vāciešiem un austriešiem bija negatīva. Zemniecību putināja rekvizīcijas, pilsētu iedzīvotāji cieta badu. Auga neapmierinātība leģionos. Taču skaidrības, kā situāciju var mainīt revolūcija Krievijā, nebija. Dažādi to vērtēja poļu politiskās partijas.

PSL ``\pltxti{Wyzwolenje}'' biedru daļa saistīja ar Februāra revolūciju cerības uz nacionālo atbrīvošanos, otra daļa, īpaši partijas vadītāji uzskatīja, ka ``mēs nevaram cerēt uz krievu revolūciju, kā mēs necerējām uz krievu caru''. Partijas vadība meklēja kontaktus ar Antantes vadītājiem, sadarbojoties ar pilsoniskajām partijām Polijā.

Šajā laikā skaidri parādījās arī strādnieku kustības attīstības nevienmērība. XIX gadsimta beigās un XX gadsimta sākumā Polijas strādnieku kustība bija visai spēcīga un veidoja būtisku visas Krievijas strādnieku kustības sastāvdaļu, taču Pirmā pasaules kara laikā, kad poļu zemes kļuva par kaujas arēnu, tika okupētas, attīstījās bēgļu kustība, daļa strādniecības tika piespiedu kārtā izvesta no dzimtenes, vietējā rūpniecībā valdīja panīkums un bezdarbs, arī strādnieku kustībā iestājās panīkums. Apstākļos, kad nebija skaidrs kara iznākums, strādnieku kustību šķēla arī dažādo kreiso partiju konkurence.

Pēc Februāra revolūcijas Krievijā PKunLSD un PPS-ļevica paziņoja: ``poļu strādnieku šķiras uzdevums mazāk par visu ir atbalstīt karojošo valstu imperiālistisko politiku ar ``militāru savienību''\dots{}''. Griežoties pie Krievijas proletariāta, abas partijas aicināja: ``\dots{}Ne karš līdz uzvarai, bet revolūcija līdz uzvarai! Lai dzīvo Krievijas, lai dzīvo Eiropas revolūcija!'' Abas poļu proletāriskās partijas uzskatīja, ka carisma gāšana nozīmē nacionālās apspiešanas likvidāciju ne tikai Krievijā, bet arī bijušajā Polijas karalistē, un poļu proletariāta priekšā stāv kara izbeigšanas un sociālistiskās revolūcijas uzdevums.

1917.~gada jūnijā Varšavā sakarā ar puses strādnieku atlaišanu no trijām militārā aprīkojuma rūpnīcām izcēlās ievērojami nemieri. Uzliesmoja streiks un strādnieki devās demonstrācijā. Pret to tika nosūtīts karaspēks, kurš arestēja ap pusotru tūkstoti strādnieku un veda tos uz citadeli. Ceļā konvoju sagaidīja barikādes, no kurām strādnieki karavīrus apmētāja akmeņiem. Tā bija pirmā strādnieku masveida sadursme ar karaspēku kopš kara sākuma. Izmantojot iestājušos apjukumu, ap 900~arestēto aizbēga. Sadursme turpinājās trīs stundas. Abās pusēs bija kritušie un ievainotie. No pārpalikušajiem un uz citadeli aizvestajiem ap 600 cilvēkiem trīs tika nošauti, pārējie nosūtīti uz koncentrācijas nometnēm un spaidu darbos. Kaut cik ievērojama sociālistu ietekme šais notikumos gan nebija manāma.

PKunLSD organizācijas darbojās arī Krievijas teritorijā. Viena no pirmajām PKunLSD grupām radās Petrogradā. To vadīja ievērojami poļu sociāldemokrātijas dibinātāji un darbinieki B.~Veselovskis, J.~Unšlihts, J.~Leščinskis, Z.~Faberkevičs u.c. Jau pirmajā savā sēdē tā nolēma iestāties KSDSP sastāvā. Pie KSDSP Petrogradas organizācijas rajoni komitejām tika izveidotas poļu sekcijas. 1917.~gada 1.~martā revolucionārais karaspēks Maskavā ieņēma Butirku cietumu un atbrīvoja ieslodzītos. To vidū bija arī F.~Dzeržinskis un viņa cīņu biedrs E.~Pruhņaks. Viņi uzreiz sāka pulcēt poļu revolucionārus pie Maskavas KSDSP komitejas. Arī Maskavā veidojās PKunLSD sekcijas. Poļu sociāldemokrātu grupas radās arī citās Krievijas lielākajās pilsētās, kurp bija evakuējušies lielākie rūpniecības uzņēmumi no Polijas. 1917.~gada jūnijā visu PKunLSD grupu Krievijā darbības apvienošanai un vadībai tika izveidota Izpildu komiteja. Tajā ietilpa jau nosauktie F.~Dzeržinskis, J.~Unšlihts, J.~Leščinskis, E.~Pruhņaks u.c.

Poļu revolucionāri-internacionālisti 1917.~gadā aktīvi piedalījās revolucionārajos notikumos Krievijā, Oktobra bruņotā apvērsuma sagatavošanā. Atgriezušies tikko no katorgas un izsūtījuma poļu revolucionāri F.~Dzeržinskis, J.~Unšlihts u.c. piedalījās KSDS(b)P VII (Aprīļa) konferencē, kā apspiestas nācijas pārstāvji iestājas par visu Krievijas un Polijas revolucionāro spēku apvienošanu. Taču viņi vēl joprojām nepietiekami novērtēja nacionālās atbrīvošanās cīņas potenciālu, apspiedējas un apspiestās nācijas pozīciju atšķirību. Tā, F.~Dzeržinskis vēl joprojām uzskatīja, ka Polijas neatkarības lozunga pasludināšana no sociāldemokrātu puses nāks par labu poļu nacionālistiem. Viņš konferencē apgalvoja, ka, tā kā poļu un krievu strādnieku šķiriskās pozīcijas ir identiskas, ``mūsu pozīcijām jābūt vienādām''. Tomēr PKunLSD Maskavas un Odesas organizācijas, kā arī virkne partijas darbinieku izvirzīja neatkarības un poļu zemju apvienošanas uzdevumus, balstoties uz revolucionāro Krieviju. Kopumā poļu marksistiskā teorētiskā doma, kaut arī izstrādāja vairākas pareizas pieejas demokrātiskās, nacionālās atbrīvošanās un strādnieku kustības problēmām, tomēr līdz Polijas neatkarības atgūšanai nepaguva tās noformēt vienotā zinātniskā koncepcijā.

Astoņi poļu sociāldemokrāti piedalījās KSDSP VI kongresā un visos jautājumos atbalstīja kongresa kursu uz bruņotas sacelšanās gatavošanu. Kongresa ievēlētajā boļševiku partijas CK sastāvā ietilpa arī F.~Dzeržinskis. Viņš un vēl viens poļu sociāldemokrāts J.~Doleckis 16.~(29.)~oktobrī piedalījās arī KSDSP CK paplašinātajā apspriedē, kurā tika pieņemts lēmums par bruņoto sacelšanos.

Jau no 1917.~gada marta savas sekcijas sāka veidot PPS-ļevica. Tā negatīvi vērtēja Krievijas pagaidu valdības darbību, nosaucot solījumus atrisināt poļu jautājumu par ``demagoģiju'' un ``tukšu plātīšanos''. PPS-ļevicas biedri apsveica Februāra revolūciju, taču neredzēja tai iespēju pāriet uz sociālistisku revolūciju, pirms tāda nebūs sākusies citās valstīs. Boļševiku ideju par šādu pāreju viņi uzskatīja par utopisku.

PPS-frakcija Februāra revolūciju sagaidīja ar gandarījumu, taču uzskatīja to Polijai kā ``ārēju, ārzemju'' lietu, par savu mērķi pasludinot neatkarīgas poļu valsts celtniecību. PPS prasīja nodot pilnībā varu Polijā tikai ar padomdevēja funkcijām apveltītajai Pagaidu valsts padomei. Paziņojot par atbalstu Polijas Republikas lozungam, tā vienlaikus piekrita atzīt arī monarhiju ``kā pagaidu formu''. Kā uzsver krievu vēsturnieces I.~Jažborovska un V.~Parsadanova, 1917.~gada pavasarī PPS gan pārgāja opozīcijā Pagaidu valsts padomei, kritizēja okupācijas režīmu, taču palika integrēta pastāvošajā politiskajā sistēmā.

Tomēr 1917.~gada jūnijā notikušais PPS XIII kongress zināmā mērā atspoguļoja Februāra revolūcijas ietekmi uz Poliju. Tā lēmumos gāja runa par poļu zemju apvienošanas iespēju neatkarīgā valstī, pieaugušo demokrātijas un sociālisma ietekmi Eiropā. Reizē gan bija izteikta cerība uz ``brīvprātīgu Polijas valstisku savienību ar Lietuvu un Baltkrieviju''. Tas liecināja, ka PPS stāveja uz poļu ``vēsturisko'' teritoriju, kuras atradās ārpus poļu etniskajā zemēm, sagrābšanas pozīcijām. Poļu sociālistiskajai partijai bija un palika nacionālistiska, lielvalstiska ievirze.

Boļševiki, pēc PPS-frakcijas biedru ieskatiem, bija atrauti no dzīves realitātes, kritizēja to nevēlēšanos sadarboties ar citām demokrātiskajām partijām, kuras noraidīja sociālistisko revolūciju.

PPS-frakcija vēlējās panākt, lai poļu sociālisti un sociāldemokrāti Krievijā, kuri līdz tam bija piederējuši dažādām frakcijām, tagad apvienojas vienotās poļu sociālistu organizācijās, ``lietas labā atmetot agrāko partijisko cīņu''. Atsevišķās Krievijas pilsētās~--- Harkovā, Saratovā, Taganrogā u.c. tādas arī uz laiku izveidojās, taču domstarpības biedru vidū bija tik lielas, ka vairākums to līdz 1917.~gada rudenim saira. Virknē Krievijas pilsētu, kur šādas apvienotas poļu sociālistu organizācijas neizdevās radīt, PPS-frakcijas biedri veidoja savas grupas. Tādas pastāvēja Petrogradā, Maskavā, Kijevā, Irkutskā u.c. To mērķis bija cīņa par neatkarīgas Polijas radīšanu. Kā šīs cīņas varonis tika atzīts J.~Pilsudskis. PPS-frakcija Krievijā centās atjaunot kara laikā pārtrūkušo sadarbību ar citām Krievijas sociāl-nacionālistiskajām organizācijām~--- armēņu dašnakiem (Armēņu revolucionāro federāciju ``\emph{Dašnakcjutun}''), Bundu, gruzīnu sociālistiem-federālistiem, latviešu sociālistiem-revolucionāriem u.c. Taču tā kā šo organizāciju vairākums stāvēja uz valstiska federālisma un nacionāli-kulturālās autonomijas pozīcijām, bet PPS-frakcija uzskatīja, ka Krievija jāsašķeļ sīkās valstiņās un daļa no tām tādā vai citādā veidā jāvieno ap Poliju, nekāda cieša savienība starp šīm grupām nebija iespējama.

Toties tuvinājās poļu pilsonisko partiju kursi. Tās visas cerības saistīja ar Antanti. Jau 1917.~gada martā nacionāldemokrātu līderis R.~Dmovskis iesniedza Lielbritānijas ārlietu ministram A.~Balfūram notu, kurā izklāstīja poļu nacionāldemokrātu partijas uzskatus nākamās Polijas valsts robežu jautājumā. Nacionāldemokrātu prasības pēc būtības maz atšķīrās no Poļu konservatīvajās un labējās sociālistiskajās aprindās radītajiem uzmetumiem. (Tā, vēl 1915.~gadā sociālistu autori rakstīja, ka nākamajai Polijai jāiekļauj Polijas karalisti, Galīciju, Cešinas Silēziju (mūsdienās Čehijas teritorija), Opoļjes Silēziju (Augšsilēzijas daļa), Poznaņu un Rietumprūsiju, tātad apgabalus, kur poļi tai laikā sastādīja tikai 62,6\% no visiem iedzīvotājiem.)

1917.~gada 15.~augustā R.~Dmovskis Lozannā uz endeku un ``reālistu'' vienošanās pamata no jauna nodibināja \strong{Poļu Nacionālo komiteju} (nosaukta tā tika tāpat kā jau 1914.~gadā Petrogradā nodibinātā komiteja, arī darbībā bija vērojama zināma pēctecība), kura 1918.~gada sākumā pārcēlās uz Parīzi. 20.~septembrī Francija, bet nākamajās nedēļās arī citas sabiedrotās valstis (Lielbritānija 15.~oktobrī, Itālija 30.~oktobrī, ASV 10.~novembrī) to atzina par ``oficiālu poļu iestādi'' (pārstāvniecību). Tās sastāvā darbojās arī I.~Paderevskis. Komiteja tūlīt uzsāka cīņu par Polijai labvēlīgām robežām pēc Vācijas sakāves, uzskatot, ka tajā jāietilpst ne tikai etnogrāfiskajiem poļu rajoniem, bet arī teritorijām ar spēcīgu poļu elementa klātbūtni. Šajā laikā poļi arvien vairāk sāka saistīt savas nacionālās cerības ar Antantes valstīm, to uzvaru karā. Francijā jau 1917.~gadā pēc prezidenta R.~Puankarē un Poļu nacionālās komitejas pārstāvja I.~Paderevska vienošanās tika veidota jauna poļu armija. Tās kodolu veidoja poļi, kuri karoja jau Francijas armijā vai Krievijas ekspedīcijas korpusā Francijā, poļu tautības karagūstekņi, līdz tam dienējušie Vācijas un Austroungārijas armijās (pavisam ap 35~000), kā arī poļi no ASV (ap 23~tūkstoši)., citām pasaules valstīm (ap 300~cilvēku atbrauca pat no Brazīlijas).

Jaunajos apstākļos orientācija uz Centrālajām valstīm zaudēja jebkādu perspektīvu. Jau maija beigās--jūnijā sāka mainīties J.~Pilsudska nostāja. Viņu vairs neapmierināja Centrālo valstu uzticama sabiedrotā loma. Redzot, ka gan Austroungārija, gan Vācija nedomā piešķirt Polijai neatkarību, ka poļu daļas tiek pakļautas ģenerālim H.~f.~Bezeleram, bet ne Pagaidu valsts padomei, J.~Pilsudskis uzsāka sarunas ar endekiem par vispārēju neitralitāti pret iekarotājiem (vienošanās gan netika panākta). Šajā laikā J.~Pilsudskis jau vairāk uzmanības veltīja nelegālās Polijas kara organizācijas (\pltxti{Polska Organizacja Wojskowa}) izbūvei, nekā legālajai armijai. Viņš bija gatavs pat uzstāties pret Vāciju, ja tas notiktu ar neatkarības lozungu. Arī Pagaidu valsts padome maija vidū uz diviem mēnešiem pārtrauca darbību. Kā norāda krievu vēsturnieks G.~Matvejevs savā pētījumā par J.~Pilsudski, viņš gan veda sarunas par atgriešanos Austroungārijas dienestā (taču tika noraidīts), gan nopietni apsvēra iespēju pāriet frontes līniju Krievijas pusē (viņa tuvākie līdzgaitnieki izstrādāja variantus to izdarīt ar aeroplānu vai kopā ar poļu ulānu pulku), bet galu galā no idejas atteicās. Viņa līdzbiedri apliecina, ka šajā laikā J.~Pilsudskim jau nebija šaubu par Centrālo valstu gaidāmo sakāvi karā. Viņš nonāca pie secinājuma, ka poļu karavīriem nav citas izejas kā atteikties nodot zvērestu Vācijai. Tā teksta izstrāde, kuru pirms tam bija atzinis arī pats J.~Pilsudskis, kā reiz beidzās un 3.~jūlijā to apstiprināja Pagaidu valsts padome. 9.~jūlijā notika poļu karavīru uzticības zvēresta nākamajam Polijas karalim un ``ieroču brālībai ar Centrālajām valstīm'' nodošanas ceremonija, kura beidzās ar poļu palīgkorpusa krīzi.

Galu galā J.~Pilsudskis ieteica karavīriem atteikties pieņemt zvērestu. Ap 1~100 karavīru, galvenokārt no poļu leģiona 1.~un 3.~brigādēm, kā arī artilērijas pulka, kuri bija dzimuši apgabalos, kas piederēja Austroungārijai vai Vācijai, zvērestu nodeva. Ap 3~300 cilvēku~--- bijušo Krievijas pavalstnieku, kuri atteicās zvērēt, tika internēti un ieslodzīti koncentrācijas nometnēs.

Pēc dažiem datiem, kopumā zvērēt uzticību Centrālajām valstīm atteicās ap 2/3 poļu virsnieku un karavīru. Nereti poļu vēstures literatūrā atteikšanās nodot zvērestu ``Centrālajām valstīm'' tiek traktēta kā dzīli patriotisks solis. Taču jāatceras, ka tie paši poļu leģionāri, to skaitā arī pats J.~Pilsudskis, savā laikā jau bija zvērējuši uzticību Austroungārijas monarhijai, un tas netika uzskatīts par nepatriotisku soli. Ir skaidrs, ka tāds simbolisks žests, ka atteikšanās no zvēresta nodošanas J.~Pilsudskim bija vajadzīgs, lai nodemonstrētu H.~f.~Bezeleram un Berlīnei savu ietekmi un tā radītu platformu tālākai politiskai darbībai.

Daļa poļu, tai skaitā daudzi no bijušā poļu leģiona 2.~brigādes, kuru komandēja J.~Hallers, gan nesekoja J.~Pilsudska ieteikumam. (Vēlāk viņi, nonākot Austrumu frontē, pārgāja Krievijas pusē.) Poļus~--- Austroungārijas pavalstniekus iekļāva tās armijā, daļu (ap 3~000) nosūtīja uz Itālijas fronti, ap 7~500 palika Austrijā. Daļu no viņiem, īpaši galīciešus, sadalīja pa austriešu armijas daļām, citus iekļāva poļu palīgkorpusā.

Vācijai pakļautajās poļu daļās palika visai nedaudz karavīru, to pieplūdums gandrīz pilnībā apstājās. 1918.~gada aprīlī tajās bija ap 2~700 karavīru. (Tomēr jāatzīmē, ka vēlāk~--- 1918.~gadā, kad Centrālās valstis piedzīvoja sakāvi karā, Polijā bijušie leģionāri veidoja militāri politisku spēku, kurš spēja nopietni ietekmēt valsts veidošanos.)

Šādā situācijā nedrīkstēja nereaģēt arī leģionu radīšanas idejas autors. ``Tad es,~--- kā vēlāk rakstīja J.~Pilsudskis,~--- uzrakstīju vēstuli Bezeleram un paziņoju, ka gribu dalīties liktenī ar maniem internētajiem karavīriem. Pēc šīs vēstules mani arestēja''. (G.~Matvejevs gan precizē, ka nosūtīt minēto vēstuli J.~Pilsudskis nepaguva, bet aresta brīdī nodeva to vācu policijas priekšniekam Varšavā. Arests notika 1917.~gada 22.~jūlijā.)

Kaut J.~Pilsudska iecerētā savas ietekmes demonstrācija izdevās tikai daļēji, vāciešu sabiedrotās Austroungārijas Ļubļinas ģenerālgubernatūra ziņoja savai Ārlietu ministrijai, ka J.~Pilsudska arests ``plašās poļu aprindās izsauca krasu niknumu pret vāciešiem'', ka Varšavā un citās pilsētās šai sakarā notika demonstrācijas, demonstranti izsituši stiklus ``Pagaidu valsts padomes'' sēžu zālei un vairākām laikrakstu redakcijām. Arī Austroungārijas okupācijas zonā notika demonstrācijas un vienas dienas streiks. Poļu karavīros J.~Pilsudska arests izsaucis ``papildus rūgšanu'', viņa autoritāte sabiedrībā atkal pieaugusi.

Vācieši J.~Pilsudski un viņa štāba priekšnieku K.~Sosnkovski no 1917.~gada vasaras līdz 1918.~gada novembrim turēja apcietinājumā Magdeburgas cietoksnī. J.~Pilsudski ieslodzīja virsnieku virssardzē, viņa rīcībā bija divstāvu koka ēkas (pēc kara to pārveda uz Varšavu un uzstādīja Belvedēra pils, kura kļuva par valsts vadītāja rezidenci, parkā kā muzeja eksponātu) trīs istabas. J.~Pilsudski uzskatīja par augsta ranga sagūstītu virsnieku, viņam piešķīra speciālas privilēģijas, personīgām vajadzībām maksāja 250 marku mēnesī, viņa saimniecību veda feldfēbelis, kurš dzīvoja tai pat ēkā. Represēta tika arī virkne Polijas kara organizācijas biedru~--- J.~Pilsudska līdzgaitnieku. Viņus ieslodzīja Varšavas citadelē. 1918.~gada sākumā, atrodoties ieslodzījumā, J.~Pilsudskis par Polijas kara organizācijas vadītāju nozīmēja savu līdzgaitnieku E.~Ridz-Smigliju.

Austroungārijas ārlietu ministrs grāfs O.~Čerņins ir atzīmējis, ka J.~Pilsudska arests padarīja viņu par ``mocekli kā reiz tad, kad viņa zvaigzne sāka bālēt'' un kļuva par viņa ``lielas popularitātes'' avotu. Arī poļu sociāldemokrāts J.~Marhļevskis rakstīja, ka arests J.~Pilsudskim izrādījās ``necerēta laime'', domājot ar to, ka it kā nejaušais arests ļāva viņam izvairīties no pārmetumiem lētticībā, sadarbībā ar svešzemju iebrucējiem, saglabāt vienkāršo savu piekritēju uzticību kā cīnītājam par Polijas neatkarību, cietējam no okupācijas varas, pat kļūt par cīņas pret okupantiem simbolu. Autors gan uzskata, ka patiesībā arests nebija ``necerēta laime'', J.~Pilsudskis, redzot savas iepriekšējās politikas bezperspektivitāti, spēja laikus pārorientēties, faktiski pats izprovocēja savu arestu, kurš viņam ļāva saglabāt tālredzīga un gudra tautas vadoņa lomu. Viņa piekritēji arī sava līdera aresta laikā neļāva tautai par viņu aizmirst, 1918.~gada februārī pat uzsāka īpašu Magdeburgas ieslodzītā apsveikuma akciju sakarā ar viņa vārda dienu 19.~martā. Speciāla ``19.~marta komiteja'' aicināja visu triju Polijas daļu iedzīvotājus piedalīties tajā. Lai nu kā, J.~Pilsudskis saglabāja autoritāti arī daudzu tālu no aktīvas politikas stāvošu tautiešu vidū.

Turpretī Pagaidu valsts padome, kura gan prasīja savas kompetences robežu paplašināšanu, bet nekādi nepretojās okupantiem, pēc krīzes sakarā ar leģionāru atteikšanos no zvēresta nodošanas zaudēja jebkādu autoritāti. Poļu demokrātiski-liberālajās aprindās izplatījās apziņa, ka Polija nedrīkst karot pret atjaunotnes ceļu ejošo Krieviju un tās ieņēma neitrālu pozīciju. PPS un ļudoviešu pārstāvji izstājās no Pagaidu Valsts padomes, bet augustā beigās tā bija spiesta pilnā sastāvā atkāpties. Vācieši nepaklausīgajiem poļiem atriebās, slēdzot Varšavas Universitāti, bet tas tikai kāpināja neapmierinātību.

Iespējams, baidoties no poļu sacelšanās, 1917~.~gada 12.~septembrī \strong{Centrālās valstis paziņoja par poļu valsts varas izveidi}. Divi imperatori~--- Vācijas un Austroungārijas~--- apņēmās radīt augstāko (līdz varas nodošanai karalim) iestādi~--- \strong{Reģentu padomi} (\pltxti{Rada Regencyjna Krόlestwa Polskiego}) no 3 cilvēkiem, kura nozīmētu valdību. Ļoti nosacīti var teikt, ka 1917.~gada 12.~septembrī bija pieņemta Polijas konstitūcija, kura paredzēja monarhiju, divpalātu parlamentu, nepolitiskus ministrus. Taču saglabājās visas abu ģenerālgubernatoru tiesības, valdības vadītāju Reģentu padome varēja iecelt tikai ar viņu piekrišanu. Pastāvēja dažādi ierobežojumi starptautisko sakaru, likumdošanas, administratīvās un tiesu varas jomā.

Sadarboties ar okupantiem bija gatavi tikai paši konservatīvākie poļu politiķi, no kuru rindām arī nāca padomes locekļi. Reģentu padome, sastāvoša no Varšavas arhibīskapa A.~Kakovska un diviem aristokrātiem~--- Z.~Ļubomirska un J.~Ostrovska, kam piešķīra tiesības pārvaldīt daļu ``kongresa Polijas'', bet ne visu Centrālo valstu okupēto teritoriju, tika izveidota 27.~oktobrī.

Pēc tam 7.~decembrī (pēc jaunā stila) tika izveidota valdība ar publicistu un vēsturnieku J.~Kuharževski priekšgalā.

J.~Kuharževskis vēlāk, laikā no 1923. līdz 1935.~gadam uzrakstīja septiņus sējumus aptverošu fundamentālu darbu ``No baltā līdz sarkanajam caram'' (``\pltxti{Od białego do czerwonego caratu'')}, kurā pirmo reizi nosauca PSRS diktatoru J.~Staļinu par caru.

Vācieši poļu pārziņā nodeva izglītības un tiesvedības jautājumus. Valdības sastāvs līdz 1918.~gada novembrim vairākkārt mainījās. Atzīmējams, ka tajā iegāja arī līdz tam esošie ``pasīvisti'' ar nacionāldemokrātu, prominento tautsaimniecības profesoru V.~Grabski priekšgalā. Taču Poļu Nacionālā komiteja Parīzē Reģentu padomē vēl joprojām saskatīja tikai vācu iebrucēju ielikteņus, vēlējās pati saņemt varu Polijā. Arī J.~Pilsudskis, atrodoties arestā, nevarēja piedalīties tekošajās ``politiskajās spēlēs'', palika nesaistīts ar Reģentu padomi, tās dibināto valdību. Šīm institūcijām arvien vairāk kompromitējoties, auga viņa politiskais kapitāls. Laiks darbojās viņa labā.

Tā faktiski 1917.~gada rudenī eksistēja divas pretēju politisko nometņu veidotas Polijas valdības, viena~--- pie Reģentu padomes Varšavā, kura sadarbojās ar vācu okupantiem, otra~--- Poļu Nacionālā komiteja Parīzē, kuru atbalstīja Antante.

\section{Poļu kultūra XIX~gadsimtā un XX~gadsimta sākumā}

\epigraph
{Kultūra ir tas viss, kas paliek pāri, kad viss cits izzūd.}
{Eduards Erio (\frtxti{Edouard Herriot})}

\epigraph
{Vienīgi kultūra padara pūli par tautu.}
{Helēna Jampoļska (\rutxti{Елена Александровна Ямпольская})}

\epigraph
{Atšķirība starp patriotu un nacionālistu ir tā, ka patriots mīl savu zemi, bet nacionālists grib to visiem izrādīt.}
{K.~Novaks (\entxti{Christoph Novak})}

\epigraph
{Vienīgais pierādījums Dieva pastāvēšanai ir māksla.}
{Arturs Šopenhauers (\detxti{Arthur Schopenhauer})}

\epigraph
{Amatu māte ir vajadzība, mākslas māte~--- pārpalikums.}
{Arturs Šopenhauers (\detxti{Arthur Schopenhauer})}

\epigraph
{Mākslas mērķis ir atspoguļot nevis lietu ārējo izskatu, bet to iekšējo nozīmi.}
{Aristotelis (\eltxti{Ἀριστοτέλης})}

\epigraph
{Nācija ir ļaužu apvienība, ko vieno ilūzija par kopīgiem senčiem un kopīgs naids pret kaimiņiem.}
{Viljams Ralfs Indžs (\entxti{William Ralph Inge})}

\epigraph
{Nacionālisms ir nevis nācijas spēka, bet gan vājuma izpausme. Ar nacionālismu sirgst galvenokārt vājas tautas.}
{Dmitrijs Ļihačovs (\rutxti{Дмитрий Сергеевич Лихачёв})}

\epigraph
{Nacionālisms~--- bērnu slimība, cilvēces masaliņas.}
{Alberts Einšteins (\detxti{Albert Einstein})}

\epigraph
{Poļi ir dzejnieki politikā un politiķi dzejā.}
{Oto fon Bismarks (\detxti{Otto Eduard Leopold von Bismarck-Schönhausen})}

\epigraph
{Tautai, kurai atņemta sabiedriskā brīvība, literatūra ir vienīgā tribīne, no kura augstuma tā piespiež dzirdēt sava sašutuma un savas sirdsapziņas kliedzienu.}
{Aleksandrs Hercens (\rutxti{Александр Иванович Герцен})}

\epigraph
{Romantisms~--- tā ir savdabīga reakcija uz Franču revolūciju.}
{Kārlis Markss (\detxti{Karl Heinrich Marx})}

\epigraph
{Domas netiek muitotas?~--- Jā, ja tās nešķērso robežas.}
{Staņislavs Ježijs Lecs (\pltxti{Stanisław Jerzy Lec})}

% page 265


Poļu nacionālā pašapziņa sāka veidoties gadsimtu gaitā uz kopīgas teritorijas, valodas, materiālās un garīgās kultūras elementu pamata. Sabiedrībai attīstoties, tā bagātinājās ar idejām par piederību vienai valstij, uzticību karalim, dinastijai, reliģijai. ``Tēvzemes'' jēdziens sākotnēji radās valdošās kārtas pārstāvju vidū. Atšķirībā no Rietumeiropas Polijā stipras centrālās varas vietā, kas balstījās ne tikai uz muižniecību, bet arī ``trešo'' kārtu, nostiprinājās vienas kārtas~--- šļahtas diktatūra. XVIII~gadsimtā vēl pastāvēja ``šļahtas nācijas'' pašapziņa, kurā nebija vietas ne zemniecībai, ne pilsētniekiem. Tikai gadsimta beigās progresīvā publicistika sāka graut kārtu barjeras poļu sabiedrības apziņā, atteicās no ``nācijas'' pielīdzināšanas šļahtai. Veidojās nācijas kā brīvu un mantīgu sabiedrības slāņu kopības, bet radikāļu vidū~--- nācijas kā tautas, kura ietver visas kārtas, izpratne. Arī izglītoto šļahtiču aprindās blakus dižciltībai arvien augstāk tika vērtēta bagātība, nopelni sabiedrības priekšā, izglītība kā vērtības, kuras deva tiesības un svaru sabiedrībā. Tomēr Polijai bija raksturīga liela starpība starp mantīgo un sabiedrības zemāko slāņu nacionālo apziņu. Atsvešinātā no valsts zemniecība nevarēja kļūt par nacionālās atbrīvošanās idejas nesēju. Zemnieku pašapziņa feodālo attiecību kundzības periodā bija vairāk saistīta ar kārtas un reliģiskās piederības izjūtām. Arī buržuāzija Polijā izveidojušās īpatnējās sociālās struktūras rezultātā bija vāja un nevarēja kļūt par aktīvu spēku nacionālās pašapziņas izveidē. Tāpēc poļu inteliģencei šai procesā bija īpaši svarīga loma.

Kā sabiedrisks slānis poļu \strong{inteliģence} sāka veidoties no XVIII gadsimta beigām. Tās pirmie pārstāvji (zinātnieki, skolotāji, juristi, ārsti, mērnieki u.c.) nāca no šļahtas un sīkpilsonības vidus. Tad pat radās arī t.s. radošā inteliģence jeb brīvo profesiju pārstāvji (rakstnieki, publicisti, žurnālisti, aktieri, gleznotāji, arhitekti). Šī slāņa izveides gaitā liela nozīme bija karaļa Staņislava Augusta, Četrgadu Seima darbinieku, atsevišķu mecenātu mērķtiecīgai darbībai. Piemēram, Polijas karalistē pēc dažiem aprēķiniem visu garīgā darba strādnieku skaits 1830.~gadā sasniedza ap 7,5~tūkstošus, nedaudz vairāk nekā 0,5\% visu profesionāli aktīvo iedzīvotāju. Politiski inteliģence pārsvarā atradās šļahtas ietekmē.

Apskatāmajā periodā ievērojamo \strong{pieauga pilsētu} kā vietu, kur attīstījās rūpnieciskā ražošana, pieauga dažādu pakalpojumu un izklaižu piedāvājums, kā politisko un kultūras centru \strong{loma}. Pilsētu iedzīvotāji, tai skaitā nesen atstājušie laukus, pakāpeniski zaudēja sakarus ar tradicionālo sadzīves kultūru, pieņēma jaunu dzīves stilu, kuru radīja inteliģence vai kurš tika ienests no ārzemēm. Pilnveidojās pilsētu saimniecība, mainījās administratīvo un dzīvojamo ēku veidols, pie tam ne tikai to centrālajos vai aristokrātiskajos rajonos, bet arī nomalēs, kur tika būvētas samērā lētas dzīvojamās mājas ar savam laikam atbilstošām ērtībām. Koka apbūvi pakāpeniski nomainīja mūra celtnes, paplašinājās ūdensvada un kanalizācijas tīkls, tomēr XX gadsimta sākumā tikai dažās pilsētās bija ierīkoti ūdensvads un kanalizācija, kādēļ bieži sākās epidēmiskās slimības. Paaugstinājās ielu labiekārtošanas līmenis, eļļas lampu vietā XIX gadsimta vidū parādījās gāzes lukturi, bet pēc tam arī elektriskais apgaismojums. Taču saglabājās arī graustu rajoni, kur trūka pašu dzīvei nepieciešamāko nosacījumu. Arī lauku ciemos sāka parādīties mūra ēkas. 1914.~gadā Prūsijai piederošajā daļa tādās dzīvoja jau 71\% zemnieku, citās zemēs to īpatsvars gan vēl bija neliels un svārstījās no 1,4\% līdz 11\%.

Mainījās apģērbi un apavi, īpaši ātri pilsētas šļahtiču, bet arī citu pilsētnieku slāņi vidū. Tērpi tika ``eiropizēti'' kā siluetu, tā arī materiālu, no kuriem tie bija izgatavoti, ziņā. Izplatījās smēķēšana, sākotnēji modē bija pīpes, bet vēlāk lētākie un ``demokrātiskākie'' papirosi. Pilsētas dzīves veids iespiedās arī lauku vidē, īpaši uzņēmīga pret jaunajām vēsmām bija lauku jaunatne. 1901.~gadā tika palaists pirmais elektriskais tramvajs Krakovā, bet 1908.~gadā~--- Varšavā.

Jau no XVIII gadsimta pēdējām desmitgadēm Polijā iznāca ap 80 periodisko izdevumu. Dažu no tiem tirāža sasniedza 2~000 eksemplāru. Nacionālās atbrīvošanās kustības pacēluma periodos šo izdevumu skaits un to tirāžas palielinājās. 1848.~gadā iznāca 105, 1861.--1864.~gadā ap 100 pārsvarā nelegālu izdevumu, kuru tirāžas sasniedza 5 līdz 8~tūkstošus eksemplāru. Galvenokārt tiem bija informatīvi-politisks, kultūras, zinātnieki-tehnisks raksturs. XIX gadsimta vidū parādījās bagātīgi ilustrēti žurnāli, kā arī specializēti periodiski izdevumi bērniem un ``vienkāršajai tautai''.

XVIII gadsimta vidū 93--95\% poļu bija analfabēti, pat šļahtiču vidū katrs otrais neprata ne lasīt, ne rakstīt. 1870.~gadu sakumā Polijas karalistē un Galīcijā lasīt un rakstīt pratēju bija ap 20\%, bet Pozenes hercogistē un Silēzijā~--- 60--70\% iedzīvotāju. Analfabētu skaits pakāpeniski samazinājās. XX~gadsimta sākumā tādu vairs nebija Prūsijai piederošajās poļu zemēs, kaut Krievijas un Austroungārijas valdījumos tie vēl sastādīja lielāko iedzīvotāju daļu. Pēc vācu autora V.~Noltinga datiem XX gadsimta sākumā Polijas karalistē bija tikai 1~skola uz 2~840 iedzīvotājiem, kad jau gadsimtu pirms tam 1~skola bija uz 2~600 iedzīvotājiem. Līdz 1914.~gadam stāvoklis uzlabojās, taču sākumskolu vēl joprojām trūka.

Salīdzinoši plaši sākumskolu tīkls bija izplatīts Prūsijas poļu zemēs, kuras atradās katoļu baznīcas gādībā. XIX~gadsimta beigās visi skolas vecuma bērni tās apmeklēja. Tiesa, mācību līmenis nebija visai augsts, pie tam apmēram puse no sākumskolām bija ģermanizētas. Augstāko un vidējo mācību iestāžu ar poļu mācību valodu vispār nebija.

1815.~gadā Polijas karalistē pastāvēja 845~skolas, kurās strādāja ap 1~100 skolotāju un mācījās ap 29~000 skolnieku. 1817.~gadā skolnieku skaits pieauga līdz 36~000. Taču jau pēc 1830.~gada sacelšanās, kad skolu lietas tika nodotas tieši Tautas Izglītības ministrijas Pēterburgā pārziņā, mācības poļu valodā tika ierobežotas. Tautas izglītības sistēma attīstījās ārkārtīgi lēni un gadījās, ka skolu skaits laukos pat pazeminājās. Jauniešu došanās uz ārzemēm lai iegūtu izglītību bija aizliegta. No XIX~gadsimta 70.~gadiem arī sākumskola tika pakļauta rusifikācijai. Izglītības līmenis bija zems, skolotāju, īpaši laukos, trūka. Daļēji arī tāpēc 1901.~gadā karalistē skolas apmeklēja tikai 82\% skolas vecuma bērnu. Trūkumus tautas izglītības jomā centās mazināt dažādas politiskas nokrāsas izglītības un kultūras biedrības. Pašizglītību, vispārējo un profesionālo zināšanu izplatību, kā arī vienlaikus patriotisko audzināšanu sekmēja dažādu pulciņu, bibliotēku, lasītavu tīkls. Pastāvēja izdevniecības, kas izdeva populāru literatūru.

Turpretī Austrijai piederošajā Galīcijā XIX gadsimta sākumā sākumskolās mācības norisa dzimtajā valodā. Vairākumā skolu tā bija rutēņu (ukraiņu) valoda. No 1815.~gada to varēja ieviest arī ģimnāzijās. Taču, kad 60.--70.~gados tika ieviesta t.s. Galīcijas autonomija, tās Seims pieņēma lēmumu par poļu valodas ieviešanu Galīcijas skolās. Poļu rokās esošā novada skolu padome (\pltxti{Rada Szkolna Krajowa}), organizēja izglītības iestāžu darbu.

Polijas karalistē un Galīcijā atsevišķos periodos varēja iegūt arī vidējo un augstāko izglītību poļu valodā. Austriešu un prūšu teritorijās pārsvarā bija klasiskās ģimnāzijas, bet Polijas karalistē, īpaši pēc 1905.~gada revolūcijas ievērojams bija reālskolu un komercskolu īpatsvars.

1816.~gadā tika nodibināta Varšavas universitāte (\pltxti{Uniwersytet Warszawski}), kurā drīz studēja ap 800~cilvēku. Universitāti stingri kontrolēja varas iestādes. Poļi studēja arī 1570.~gadā atklātajā Viļņas akadēmijā (\latxti{Almae Academia et Universitas Vilnensis Societatis Jesu}), kuru 1802.~gadā pārvērta par laicīgu universitāti. Pēc 1830.~gada sacelšanās Varšavas un Viļņas universitātes slēdza. Taču tā kā bija nepieciešami izglītoti kadri, 1833.~gadā Polijas karalistes lielākajās pilsētās tika dibinātas 8~klašu ģimnāzijas, kuras gatavoja jauniešus ne tikai darbam civilajā un militārajā dienestā, kā arī rūpniecībā, bet arī iestājai universitātēs Krievijā.

60.~gadu sākuma liberalizācijas periodā Polijas karalistē atkal tika nodibināta augstākā mācību iestāde~--- 1862.~gadā notika Varšavas Galvenās skolas (\pltxti{Szkoła Główna Warszawska}) atklāšana. Varšavas Galveno skolu 1869.~gadā pārvērta par vienu no Krievijas universitātēm (\rutxti{Императорский Варшавский университет}), kur mācības notika krievu valodā un darbojās Viskrievijas universitāšu statūti, kas ievērojami ierobežoja tās autonomiju.

Bez universitātes Varšavā tika izveidotas kara skola (\rutxti{Варшавское пехотное юнкерское училище}, 1860--1887) un Politehniskais institūts (\pltxti{Politechnika Warszawska}, 1898), 1818.~gadā~--- Meža skola (\pltxti{Szkoła Leśna}), 1820.~gadā~--- Agronomiskais institūts (\pltxti{Instytut Agronomiczny}). 1816.~gadā Kelcē tika nodibināta Polijā pirmā tehniskā universitāte (\pltxti{Szkoła Akademiczno-Górnicza}). Karalistē pastāvēja arī kadetu korpuss Kališā (\pltxti{Korpus Kadetów w Kaliszu}, 1793--1831), skolotāju institūts Lovičā (\pltxti{Instytut Nauczycielski w Łowiczu}, no 1854), luterāņu garīgā akadēmija un ebreju rabīnu skola.

Labvēlīgāki apstākļi augstākās izglītības attīstībai bija Austrijas poļu zemēs. 1364.~gadā dibinātajā Krakovas Akadēmijā (\pltxti{Akademia Krakowska}, vēlāk \pltxti{Szkoła Główna Koronna}, no XIX gadsimta 70.~gadiem \pltxti{Uniwersytet Jagielloński}) ar poļu zemju iekļaušanu Habsburgu impērijā par studiju valodu kļuva vācu valoda, taču 1809.~gadā poļu valoda atgriezās mācību procesā. No 1870.~gada universitāte pilnībā no vācu pārgāja uz poļu mācību valodu. Krakovas universitātes autonomija bija stipri ierobežota, vēlāk to pilnīgi likvidēja.

1871.~gadā tika paplašināta pasniegšana poļu valodā Lembergas universitātē (oficiāli dibināta 1661.~gadā, līdz 1824.~gadam mācības norisa latīņu, pēc tam vācu valodā), arī tehniskajās mācību iestādēs (1844.~gadā Lembergā tika radīta Tehniskā Akadēmija (\pltxti{Technische Akademie}), vēlāk Politehniskā skola (\pltxti{Polytechnische Schule}), 1856.~gadā Dubļanu ciemā pie Lembergas nodibināta Lauksaimniecības akadēmija (\pltxti{Wyższą Szkołę Rolniczą w Dublanach}).) Poļu mācību iestādes Austro-Ungārijā kļuva populāras arī Pozenes un Polijas karalistes poļu vidū.

Liela loma poļu inteliģences izveidē bija augstākās izglītības ieguves iespējām Rietumu un Krievijas augstskolās Vīnē, Berlīnē, Breslavā (Vroclavā), Pēterburgā, Maskavā, Tērbatā u.c. Mācības svešumā veicināja poļu studentu patriotiskas vides veidošanos.

Pirms Pirmā pasaules kara Pēterburgā vien mācījās jau ap 1~400 poļu studentu. Kopējais poļu skaits pilsētā sasniedza 70~tūkstošu, un bija pats lielākais nacionālais mazākums Krievijas galvaspilsētā. Nacionālās piederības sajūtu nodrošināja katoļu baznīca. Pirms 1917.~gada Pēterburgā, Carskoje Selo un Pētergofā darbojās 16~katoļu dievnami un 17~kapelas, kā arī Romas katoļu Garīgās akadēmija un Garīgais seminārs. Poļu kopiena Pēterburgā sniedza plašu materiālo atbalstu poļu izglītības, kultūras un labdarības iestādēm kā pašā Pēterburgā, tā arī Polijā. Vēlāk Otrās Žečpospolitas laikā tās valsts un sabiedrisko darbinieku, rakstnieku un gleznotāju, zinātnieku un militāristu vidū bija vairāki simti Pēterburgas ģimnāziju, virsnieku skolu, Pēterburgas Universitātes un Romas katoļu Garīgās akadēmijas absolventu. Tālāku poļu diasporas uzplaukumu apstādināja 1917.~gada revolūcija. 1919.--1924.~gadā uz dzimteni no Petrogradas repatriējās ap 60~tūkstošu poļu. Līdz 1921.~gadam Polijas armijā veiktajā virsnieku pāratestācijā noskaidrojās, ka no 18~tūkstošiem gandrīz 1/3 bija bijušās cara armijas virsnieki, bieži pat visai augsta ranga.

Māksla un literatūra, kā vienmēr, arī XIX~gadsimtā atspoguļoja vēsturisko esamību. Bieži valstīs, kurās tiek ierobežotas politiskās un citas brīvības, aktīvākie sabiedrības locekļi pastiprināti pievēršas darbībai kultūras jomā, tā savdabīgi kompensējot ierobežojumus citās sfērās. Līdzīgi XIX~gadsimtā attīstījās kultūra Polijā. Valstiskās neatkarības zaudēšana nebūt neveda pie kultūras dzīves izdzišanas. Vairāku poļu sacelšanos sakāve tikai kāpināja kultūras nozīmīgumu. Poļu apdzīvoto teritoriju sašķeltība palielināja kultūras nozīmi vienotas nācijas izveidē, sabiedrības saliedēšanā. Kā reiz kultūra kalpoja kā nacionālās vienotības pierādījums politiski sadalītajā Polijā. Protams, sašķeltība traucēja arī kultūras attīstībai, tai nebija tāda dabīgā pamata kā nacionālā skola, tās nostiprināšanos bremzēja dažādi ierobežojumi. Taču māksla, un sevišķi literatūra, tāpat kā baznīca kļuva par efektīvu cīņas līdzekli kā uzturēt un attīstīt nacionālo apziņu.

Padomju gados (1983) LZA loceklis Indulis Ronis, uzskaitījis faktus par poļu kultūras attīstību Galīcijā, rakstīja par to kā ``par vienīgo Rečas Pospoļitas [Žečpspolitas~--- V.Š.] apgabalu, kur XIX~gs. otrajā pusē zināmā mērā varēja attīstīties poļu nacionālā kultūra\dots{}'', kas nepavisam neatbilst vēsturiskajai patiesībai. Galīcijā apstākļi poļu kultūras attīstībai minētajā laikā bija labvēlīgāki kā Polijas karalistē un Vācijas Pozenes provincē, taču neraugoties uz nacionālo jūgu, arī tajās poļu kultūra piedzīvoja nemitīgu attīstību, Galīcija nebūt nebija ``vienīgais'' apgabals, kur varēja attīstīties poļu nacionālā kultūra.

Jau XIX~gadsimta pirmā trešdaļa bija poļu kultūras ievērojamas augšupejas laiks. Šai laikā labvēlīgāki nekā citās Polijas daļās apstākļi kultūras attīstībai bija Varšavas hercogistē (1807--1812), pēc tam Polijas karalistē (līdz 1830.~gadam).

Svarīgākie kultūras centri bija Varšava, kā arī aiz poļu zemju robežām esošā Viļņas universitāte (1803--1832), Kremeņecas licejs (\pltxti{Liceum Krzemienieckie}, arī \pltxti{Liceum Wołyńskie}, 1805--1833)~--- Kremeņecas (poļu \pltxti{Krzemieniec}, ukraiņu \uktxti{Кременець}, krievu \rutxti{Кременец}) pilsētā Volīnijā pastāvošā poļu mācību iestāde, vēlāk pārcelta uz Kijevu un pārveidota Kijevas sv. Vladimira universitāti (\rutxti{Киевский Императорский университет святого Владимира}). Abās mācību iestādēs tika veikti pētījumi dabas zinātnēs un medicīnā.

Lembergā 1817.~gadā ar bibliogrāfa un vēsturnieka grāfa J.~Osoļinska gādību tika nodibināta grāmatu krātuve ``\latxti{Ossolineum}''. Par zinātnes un kultūras centriem kļuva arī magnātu dzimtu Čartorijsku (\pltxti{Czartoryscy}) bibliotēka Krakovā, Krasiņsku (\pltxti{Krasińscy}) un Zamoisku (\pltxti{Zamoyscy}) bibliotēkas Varšavā.

1800.~gadā Varšavā tika nodibināta pirmā nozīmīgā ``Zinātnes draugu biedrība'' (\pltxti{Towarzystwo Warszawskie Przyjaciół Nauk}, 1800--1832). Tās priekšsēdētājs 1808--1826.~gadā bija ievērojamais sabiedriskais darbinieks un zinātnieks S.~Stašics. Pēc šīs biedrības parauga līdzīgas dibinājās arī citās Polijas pilsētās. Tā, 1815.~gadā tika izveidota Krakovas zinātniskā biedrība (\pltxti{Towarzystwo Naukowe Krakowskie}), kura no 1872.~gada pieņēma Zinību akadēmijas (\pltxti{Akademia Umiejętności}) nosaukumu. Arvien aktīvāk sāka izpausties pirmās poļu inteliģences paaudzes darbība.

Pēc 1830.--1831.~gada sacelšanās sakāves poļu kultūras dzīves centrs pārvietojās uz Austrijas impērijas poļu zemēm. Vēlāk radās zinātniski centri Prūsijai pakļautajās poļu zemēs. 1857.~gadā Pozenē tika nodibināta Zinātnes draugu biedrība (\pltxti{Poznańskie Towarzystwo Przyjaciół Nauk)}. Pēc 1830.~gada poļu zinātniski-kultūras centri izveidojās arī ārzemēs: Parīzē, Rapersvillā (\frtxti{Rapperswill}) Šveicē u.c.

Kaut pēc Polijas sadales mākslinieciskais process dažādās tās daļās norisa dažādi, pastāvēja kopēju tēmu un jēdzienu loks, vispārnacionāla problemātika, vispārnacionāli ideāli, kurus pauda kultūras darbinieki.

No XVIII gadsimta vidus Polijā valdošā bija Rietumos izplatītās \strong{Apgaismības} kā idejiskas kustības \strong{ietekme}. Rietumos Apgaismība iezīmēja pāreju no feodālisma pie brīvā tirgus saimniecības, atspoguļoja dzimstošās buržuāzijas un tautas masu cīņu pret feodālismu. Polijā, lai gan atpalika no Rietumiem kapitālisma attīstības ziņā, tomēr atradās feodāli dzimtbūtnieciskās iekārtas sairuma stadijā. Antifeodālo pārkārtojumu objektīvā nepieciešamība tajā kļuva par Apgaismības ideju izplatības pamatu. Apgaismības ideoloģija atrada piekritējus pirmkārt progresīvāk noskaņotajā šļahtas daļā, kura iestājās par centrālās varas nostiprināšanu un magnātu patvaļas ierobežošanu, modificējot pastāvošās iekārtas sociāli-ekonomiskos pamatus. Bez tam Polijā Apgaismība bija cieši saistīta ar nacionālās atbrīvošanās kustību. Polijas dalīšanas ievērojami ietekmēja nacionālās kultūras likteni, taču nemainīja tās attīstību apgaismības ideju gultnē. Poļu autori bija pilni vēlmes saglabāt nacionālo kultūru arī svešzemju kundzības apstākļos. Apgaismības kultūra ieguva racionālismam raksturīgas iezīmes, to veicināja klasicisma stils mākslā, kurš jaunā līmenī atgriezās pie antīkā un Renesanses laikmetu mantojuma un kuram piemita cildena formu vienkāršība un skaidrība.

Apgaismības ideju nesēji Polijā bija šļahtiču izcelsmes inteliģenti. Poļu apgaismotāju priekšgalā atradās izcili zinātnieki H.~Kollontajs, S.~Stašics un brāļi Jans un Andžejs Sņedecki~--- kārtu iekārtas un dzimtbūšanas pretinieki, nāciju vienlīdzības, nacionālās suverenitātes neaizskaramības, dabīgo tiesību un sabiedriskā līguma ideju aizstāvji. H.~Kollontajs daudz paveica izglītības sistēmas reformā Polijā, piedalījās sacelšanās T.~Kostjuško vadībā. S.~Stašics savos pamatdarbos ``Brīdinājums Polijai'' (``\pltxti{Przestrogi dla Polski}'', 1790) un ``Cilvēku dzimta'' (''\pltxti{Ród Ludzki}'', 1820) kritizēja feodālismu, magnātu visvarenību, iestājās par tirdzniecības un amatniecības attīstību, muižnieku aplikšanu ar nodokļiem. Izglītības jomā viņš uzstājās par patriotiskās audzināšanas nostiprināšanu, matemātikas, dabaszinātņu, ģeogrāfijas un vēstures ieviešanu skolu programmās.

J.~Sņedeckis bija izcils astronoms, kurš sarakstījis arī darbus matemātikā un filozofijā, arī literāri-kritiskus rakstus. A.~Sņedeckis bija ķīmijas, bioķīmijas un fizioloģijas zinātņu nozaru pamatlicējs Polijā.

Polijas valstiskuma likvidēšana iezīmēja pēdējā etapa (turpinājās līdz XIX gadsimta 20.~gadiem) sākumu poļu Apgaismībā. Ārvalstu varas drukāto vārdu šai laikā uzskatīja par līdzekli, ar kuru var tikt izplatītās idejas, pret kurām cīnījās valdošie nami Eiropā. Cenzūra darbojās visās sadalītās Polijas daļās, kas ievērojami traucēja kultūras, sevišķi literatūras, attīstībai. \strong{Literatūra} sadalījās emigrantu izdotajā, kura iznāca ārzemēs bez cenzūras un kontrabandas ceļā nonāca poļu zemēs, un uz vietas drukātajā. Ja pirmajā pilnā mērā izpaudās visai Eiropai raksturīgie idejiskie meklējumi, tad otrā bija stipra ar novērojumu dziļumu un notiekošo sociālo procesu atspoguļojumu Polijā.

Šai laikā priekšplānā kā ``oficiāla'' tendence, kuru nostiprināja arī Varšavas ``Zinātnes draugu biedrība'', izvirzījās \strong{klasicisms}, kurš gan XIX gadsimta sākumā citur jau atradās norieta stadijā. Pēc Polijas trešās dalīšanas poļu literāti par savu pienākumu uzskatīja sludināt patriotismu. Viņi pēc franču rakstnieku parauga, tēlojot savas tautas vēsturi, centās radīt jaunu žanru~--- nacionālo traģēdiju. Arī historiogrāfijā par valdošo pakāpeniski kļuva varonības momentu idealizācija pagātnē, lai iejūsminātu šļahtu uz nacionālās atbrīvošanās cīņu.

Kā ievērojamākais šī laika dzejnieks, dramaturgs, literatūras teorētiķis jānosauc ``Zinātnes draugu biedrības'' loceklis valodnieks, dzejnieks \strong{A.~Felinskis}. Lielus panākumus guva viņa 1811.~gadā pabeigtā, 1817.~gadā uz skatuves uzvestā un 1820.~gadā izdotā traģēdija ``\pltxti{Barbara Radziwiłłόwna}''. Tā kļuva par poļu pseidoklasiskās traģēdijas paraugu. Laikabiedri traģēdijā saskatīja Polijas agrākās varenības apofeozi.

Četrgadu Seima deputāts, 3.~maija Konstitūcijas piekritējs, Varšavas hercogistes ministrs-valsts sekretārs, Aleksandra I laikā tautas izglītības un ticības lietu ministrs grāfs \strong{S.~Potockis} bija pazīstams arī kā rakstnieks. Savas teorijas viņš attīstīja nelielās brošūrās ``\pltxti{Swistki krytyczne}'' (``Kritiskie svilpieni''). 1820.~gadā viņš izdeva poētisku fantāziju ``\pltxti{Podroz do Ciemnogrodu}'' (``Ceļojums uz Pilsētu Tumšo''), kurā izsmēja māņticību, ticēšanu leģendām un teikām.

Senators K.~Kozmians literatūrā plašāk kļuva pazīstams, pateicoties savām poēmām, kurās atdarināja senās Romas dzejnieka Vergīlija darbus.

Vēlīnie klasicisma poļu epigoņi idealizēja Polijas un tās šļahtas pagātni, slavēja tās agrāko varenību, bet nespēja norādīt ceļu atdzimšanai. Uz šī fona izcēlās poļu leģionāru dzeja, kurā viņi aicināja turpināt cīņu par atbrīvošanos, pauda ticību uzvarai. Jau minēts J.~Vibickis un viņa radītā himna ``\pltxti{Jeszcze Polska nie zginęła}''. Pazīstams dzejnieks bija leģionārs C.~Godebskis. Viņš apdziedāja leģionāru varonību, ar rūgtumu rakstīja par poļu patriotiem, Napoleona nosūtītiem uz Haiti.

Kaut XIX~gadsimtā Polijā literatūras attīstību traucēja poļu tautas dzīves apstākļi: zemais lasītprasmes līmenis, pilsētu iedzīvotāju mazais īpatsvars, kas noteica lasītāju loka šauro apjomu, tomēr notika arī literatūras komercializācija. Rakstnieku sacerējumus sāka publicēt ne tikai žurnāli, bet arī avīzes. Līdz XIX~gadsimta beigām izveidojās literāta--profesionāļa tips, kurš atlīdzību par savu darba saņēma nevis no mecenāta, bet no lasītāja. (Tiesa, vēl gadsimta vidū tāds ražīgs rakstnieks kā J.~Kraševskis, kurš sacerēja ap divus simtus tematikas ziņā dažādu~--- gan vēsturisku, gan laikmetīgu romānu, nevarot apgādāt ģimeni ar literāro darbību, bija spiests nomāt muižu.)

No Napoleona karu beigām poļu sabiedrības garīgo dzīvi galvenokārt noteica \strong{romantisma} (\pltxti{romantyzm}) strāvojums. J.~G.~Herdera ideju ietekmē liela uzmanība tika pievērsta tautiskajai kultūrai, nacionālajai savdabībai. Viņš taču ieteica griezties pie agrīnajiem savas tautas vēstures periodiem, smelties tajos spēku dzīves un mākslas atjaunotnei. Romantisms, kurš sludināja ticību garīgumam, indivīda kultu, labi atbilda poļu nacionālās atbrīvošanās kustības specifikai. Ticība cilvēkam un tautai caurvija visu poļu romantisma vēsturi. Dažādos veidos tika pausta doma par individuālo ziedošanos, par vienu no centrālajiem kļuva motīvs par egoistiski lepna, patvaļīga, tikai saviem mērķiem dzīvojoša cilvēka pārvēršanos par idejisku cīnītāju, vēsturiska uzdevuma īstenotāju. Romantisma gars atspoguļoja poļu autoru noskaņojumus, ietekmēja viņu izturēšanos pret Poliju sadalījušajām varām, radīja garīgu pamatu XIX~gadsimta poļu sacelšanās kustībām. Romantisma laikmetā bija vērojams ievērojams poļu kultūras uzplaukums. Vadošā literārā centra lomu XIX~gadsimta sākumā ieguva Varšava. 20.~gados tā kopā ar citiem centriem, kuri atradās ārpus Polijas karalistes, kļuva par poļu romantisma izcelsmes vietu.

Romantisma laikmetā faktiski sakrita galvenie nacionālās atbrīvošanas kustības un literatūras attīstības etapi. Līdz 20.~gadu beigām bija vērojams kā slepeno biedrību, tā romantisma veidošanās process. Privātie mecenāti, bet nereti arī valstiskās iestādes labprāt finansēja beletristiskus un vēsturiskus sacerējumus ar patriotisku saturu, motivējot to ar rūpēm par tautas izglītību, nepieciešamību visu nāciju iesaistīt cīņa par neatkarību. Reizē šāds mecenātisms ieguva nacionālā mesiānisma propagandas raksturu.

Poļu romantiskās literatūras sākumu speciālisti datē ar 1822.~gadu, kad iznāca Ā.~Mickēviča darbu pirmais sējums. Romantiskā gaisotnē, bet no 40.~gadiem arī romantisma un reālisma robežjoslā, ciešā saistībā ar nacionālo atbrīvošanās kustību attīstījās daudzu tolaik pazīstamu poļu dzejnieku daiļrade. Saikne bija ne tikai emocionāla, bet arī personiska. Praktiski nebija neviena poļu pagrīdnieka, pie kura aresta laikā neatrastu pazīstamu romantiķu tekstus vai izrakstus no tiem, bet bieži arī paša mēģinājumus dzejā. Lielākā dzejnieku~--- romantiķu daļa vai nu paši piedalījās pagrīdes organizāciju darbībā, vai uzturēja ar tām vai poļu revolucionārajiem emigrantiem ārzemēs pastāvīgus sakarus.

Jau līdz 1830.~gada sacelšanās sākumam poļu literatūrā tika izmantoti visi tie žanri (balādes, poētiski garstāsti, fantastiskas drāmas), kuri pastāvēja Eiropas valstīs. 1830.--1831.~gads bija ne tikai sacelšanās gads, bet iezīmēja arī romantiskās literatūras uzplaukumu emigrācijā, kad romantisms jau izspieda citus virzienus, kļuva par galveno virzienu literatūrā. Pēc poļu sacelšanās sakavēs 1831.~gadā romantisms saistījās ar cerībām panākt revolucionāru pavērsienu Polijas liktenī.

Romantisma laikmetā par poļu tautas centienu izteicējiem un ``tautas varoņiem''. kļuva mākslinieki, īpaši literāti. Rakstnieki uzņēmās garīgo, bet dažkārt arī politisko līderu, morālo autoritāšu, tiesību aizstāvju, vadoņu lomu. Literatūra kļuva gandrīz vai par vienīgo nacionālās un kulturālās pašapziņas saglabāšanas formu. XIX~gadsimta poļu literatūrā rakstnieku, bet īpaši dzejnieku vārdi nereti ieguva augstākā taisnīguma, tiesību, gandrīz vai atklāsmes vērtību. Dzejniekus stādīja uz pjedestāla kā nacionālus praviešus, literatūru uzskatīja par ``cēlu uzdevumu, misiju''. Tiesa, tādām prasībām varēja atbilst tikai izcilākie XIX~gs. dzejnieki: Ā.~Mickēvičs, J.~Slovackis, Z.~Krasiņskis, K.~Norvids.

Romantisma idejas plaši izplatījās poļu emigrantu aprindās, kuras garīgi saistīja sadalītos poļu apgabalus. Politiskās misijas apziņa emigrācijā saistījās ar reliģisku misticismu, izpaudās t.s. \strong{``mesiānismā''} (pretenzijās glābt tautu vai pat visu cilvēci). Īpaši skaidri tas izpaudās A.~Toviaņska filozofiskajā mācībā, kurai bija liela ietekme poļu emigrantu aprindās Parīzē. Pēc šīs mācības ar atsevišķu cilvēku starpniecību darbojas nemirstīgi gari, Dieva sūtņi, kuri realizē viņa gribu (pie tādiem A.~Toviaņskis pieskaitīja Napoleonu un arī sevi). Pēc A.~Toviaņska mācības Polija bija vienīgā īsteni kristīgā valsts~--- Mesija. Tās sūtība bija cilvēces grēku izpirkšana, kuras gaitā tā kā Mesija tika sista krustā. Kā Vecajā derībā ebreji, Dieva izredzētā tauta, tika izkaisīta pa pasauli, tā par to, ka tā neprata aizstāvēt dzimteni, tika izkaisīta arī Jaunās Derības izredzētā tauta~--- poļi. Taču nāks laiks, poļi pulcēsies savā zemē un Polija atdzims, pasaule pārdzims jaunā veidolā. Mesiānisma idejas atrada izpausmi gan literatūrā, gan mākslā, bet politiskajā plaknē tās pārtapa par poļu hegemoniju citu, pirmkārt~--- slāvu, tautu vidū.

\strong{Ādams~Mickēvičs}~--- nabadzīga šļahtiča dēls, kura māte bija cēlusies no kristītu Lietuvas ebreju dzimtas, kurš pats par sevi rakstīja ``Mana dzimtene~--- Lietuva'', bija poļu romantisma iedvesmotājs un līdz mūsdienām ievērojamākais poļu dzejnieks. Viņš pirmais pacēla poļu dzeju kā līdztiesīgu Eiropas tautu literatūrā. Jau pirmais viņa dzejas krājums ``\pltxti{Poezje}'' (``Dzejas'', 1822), kurā bija iekļautas balādes un romances (``\pltxti{Ballady i Romanse}''), kā arī priekšvārds ``\pltxti{O poezji romantycznej}'' (Par romantisko dzeju), kļuva par romantisma virziena manifestu poļu literatūrā.

Pēc gada iznāca ``Dzeju'' otrais krājums ar romantisko liriski-episko poēmu ``\pltxti{Grażyna}'' un dramatiskās poēmas ``\pltxti{Dziady}'' (``Vectēvi'') 2.~un 4.~daļu. Vēlāk turpinātajā poēmā dota jauna poļu nācijas sūtības versija, izklāstīta ``poļu mesiānisma'' ideja, atbilstoši kurai Polijas ciešanas ir saistītas ar īpašu tautas-mocekles jeb ``Kristus tautas'' sūtību. Polija šeit tēlota kā Jēzus Kristus, bet Francija pilda Jūdejas pārvaldnieka Poncija Pilāta lomu. Mesiānisms dziļi ietekmēja poļu pašapziņu līdz pat 1863.~gada sacelšanās sakāvei.

Šī pati ideja ir attīstīta mākslinieciski--publicistiskajā sacerējumā ``\pltxti{Księga narodu polskiego i pielgrzymstwa polskiego}'' (``Poļu tautas un poļu piligrimu grāmata''). Krievijā iznāca viņa darbi ``\rutxti{Сонеты}'' (``Sonētas''), poēma ``\rutxti{Конрад Валленрод}'' (``Konrads Vallenrods''). Pēdējā vēsta par lietuviešu cīņu pret krustnešiem. Nosaukumā minētais Konrads Vallenrods tēlots kā traģisks varonis, cīnītājs~--- vienpatnis. Lietuvietis pēc izcelsmes, spiests atteikties no savas dzimtenes un kļuvis par Teitoņu ordeņa maģistru, viņš veda to pretī sakāvei. [Pretēji Ā.~Mickēviča interpretācijai patiesībā K.~Vallenrode~--- Teitoņu ordeņa 24.~lielmestrs (1391--1393) bija cēlies no senas franku bruņinieku dzimtas, cīņā par vadību ordenī sastapās ar stipru opozīciju, kaujās ar ordeņa ārējiem ienaidniekiem cieta arī dažas neveiksmes, bet viņu nevar vainot apzinātā ordeņa vešanā pretī sakāvei.~--- V.Š.] Laikabiedri poēmā saskatīja poļu cīņu pret apspiedējiem, morāles un politikas samierināšanu, kad ārēji drīkst samierināties ar ienaidnieku, iet tā kalpībā, vienlaikus darbojoties pret to. No šī laika vēsturiskajā literatūrā pat ir ieviesies termins ``vallenrodisms''.

Apjomīgākais dzejnieka sacerējums bija tai laikā Eiropas literatūrā vienīgais varoņeposs ``\pltxti{Pan Tadeusz}'' (darba pilns nosaukums~--- ``\pltxti{Pan Tadeusz, czyli Ostatni zajazd na Litwie}''~--- ``Tadeuša kungs, vai pēdējais brauciens uz Lietuvu''), kurā radīts krāsains, humora un nostalģijas pilns vēsturiskai nolemtībai pakļauto šļahtiču tikumu attēlojums. Darbs tiek uzskatīts par poļu nacionālo epopeju un vārda mākslas šedevru. Poļu emigrācijas kreisā spārna ideologs S.~Vorcels, to nosauca par ``kapakmeni, ģēnija roka uzliktu vecajai Polijai''. 1999.~gadā režisors A.~Vajda eposu ekranizēja.

Savas dzīves laikā Ā.~Mickēvičs iemantoja nesatricināmu nacionālajā jūgā nospiestās poļu tautas vadoņa lomu. Viņa slava bija krietni lielāka nekā pārējiem nacionālajiem romantiķiem. 1855.~gadā Ā.~Mickēvičs, vēloties organizēt jaunu poļu leģionu cīņai pret Krieviju (risinājās Krimas karš), devās uz Konstantinopoli, kur saslima ar holēru un mira. 1890.~gadā viņa pīšļi tika pārvesti uz Krakovu un apbedīti Vavelas katedrālē.

Izcils romantisma skolas pārstāvis bija dzejnieks un dramaturgs \strong{Jūliušs Slovackis}. 1830.~gada sacelšanās laikā viņš kā Nacionālās valdības diplomātiskais kurjers nonāca ārzemēs, kur arī radīja savus nozīmīgākos darbus. Virtuozi pārvaldot poētisko valodu, augstākās virsotnes viņš sasniedza nepabeigtajā poēmā~--- romantiskā subjektīvisma un literārās polemikas šedevrā ``\pltxti{Beniowski}'' (1840--1841) un epopejā ``\pltxti{Król-Duch}'' (``Karalis-gars''), kurā tēloja poļu tautas dzīvi no senajiem līdz viduslaikiem. Dzejnieka jaunradē izpaudās daudzveidīgi žanru meklējumi, bagātīgas ieceres, daudzas no kurām tā arī netika līdz galam īstenotas. Īpaši atzīmējams J.~Slovacka ieguldījums romantiskās drāmas attīstībā. Drāmas ``\pltxti{Kordian}'' (1834), ``\pltxti{Balladyna}'' (1834), ``\pltxti{Lilla Weneda}'' (1839) pieder poļu dramaturģijas šedevriem. Individuālā traģisma motīvi tajās savijas ar nacionālajām filozofiski--vēsturiskajām problēmām. Abās pēdējās drāmās izteikta tēze, ka feodālisms un katolicisms ir uzspiests no ārienes, pārkāpjot īsteni slāvisko tautvaldības principu.

J.~Slovacka lirikā dominēja patriotiskas grūtsirdības, uzticības revolucionāra sūtībai motīvi. Daudzi darbi bija veltīti 1830.~gada sacelšanās sakāvei un poļu emigrantu un izsūtīto liktenim. Dzejnieks uzstājās no poļu šļahtas viedokļa, taču skeptiskā attieksme pret varoni--individuālistu bija ekvivalenta valdošās kārtas kritikai. Viņš atmaskoja Rietumeiropas valstu valdību un arī Vatikāna vienaldzību pret Polijas likteni. J.~Slovacka spalvai pieder arī vairāki sacerējumi vēsturiskās prozas žanrā. J.~Slovacka daiļrade atstāja milzu iespaidu uz poļu literatūras attīstību. 1927.~gadā viņa pīšļi, tāpat kā 1890.~gadā A.~Mickēviča mirstīgās atliekas, tika pārapbedītas Krakovā Vavelas katedrāles kapenēs.

Arī dzejnieks un dramaturgs \strong{Zigmunts Krasiņskis} ir pieskaitāms ievērojamākajiem romantikas pārstāvjiem Eiropā. Pirmie viņa sacerējumi bija V.~Skota darbu atdarinājumi, taču nākamie jau veltīti Polijas sociālajiem jautājumiem, mēģinot tos pacelt vispārcilvēcisko problēmu līmenī. Nozīmīgākie viņa darbi radīti XVIII gadsimta 30.~gados. Ļoti populāra kļuva anonīmi publicētā drāma ``\pltxti{Nie-Boska Komedia}'' (``Nedievišķā komēdija''). Tajā autors ar pārsteidzošu tālredzību traktēja tajā laikā vēl neapzinātās šķiru cīņas un sociālās revolūcijas tēmas. Taču, neredzot iespēju piepildīt savus sapņus un piedzīvot Polijas atbrīvošanu no ``moskaļiem'', viņš meklēja glābiņu misticismā.

\strong{Kiprians Norvids} piederēja tai jauno rakstnieku paaudzei, kura 40.~gados gribēja saraut saites ar tradicionālo šļahtas pagātni. Viņš nepiekrita romantiķu vidū izplatītajam uzskatam par poļu mesiānisko lomu, radīja oriģinālu humānā universālisma koncepciju un no tās viedokļa centās analizēt svarīgākās vēstures, politikas un kultūras problēmas. Neapšaubāmu talantu apliecinošie darbi reizē ir apveltīti mistiskām noskaņām. Taču dzīves laikā tika publicēts maz viņa kā dzejas, tā prozas darbu. Tikai XX~gadsimta sākumā tie tika no jauna ``atklāti'' (daudzi no tiem diemžēl jau bija zuduši) un K.~Norvids ieņēma savu godpilno vietu blakus pirms tam nosauktajiem poļu literātiem. XX gadsimta 70.~gados Varšavā iznāca viņa kopotie raksti vienpadsmit sējumos.

Var minēt arī citus romantiķus: pirmās poļu romantiskās poēmas ``\pltxti{Maria. Powieść ukraińska}''; (``Marija, ukraiņu stāsts'', 1825), kura ieguva milzīgu popularitāti, autoru A.~Malčevski, poļu zemnieku dzīvi idealizējošo, tautiskās tradīcijās rakstītās, ar skaistu formu apveltītās dzejas radītāju T.~Lenartoviču u.c.

Ievērojamu vietu XIX gadsimta pirmās trešdaļas literatūrā ieņem J.~Niemcevičs, kurš gan nav tieši pieskaitāms romantiķu~--- izcilnieku grupai. J.~Niemcovičs radīja nacionālās balādes žanru. Lielu popularitāti ieguva arī viņa vēsturisko dziesmu krājums.

Sasniedzis XIX gadsimta vidū savu augstāko attīstību poļu romantisms iegāja norieta fāzē, kaut vēl līdz 1863.~gadam bija valdošais virziens literatūrā. Romantisms atklāja poļiem to vēsturisko pagātni, virzīja viņi domas par vēstures būtību un darbības mehānismiem, sagatavojot vēsturiskā romāna žanra uzplaukumu XIX gadsimta beigās. Šajā laikā sasniegumi bija vērojami arī prozā, taču dzeja vēl joprojām sabiedrībā tika vērtēta augstāk nekā stāsti un romāni, jo tieši tā bija tradicionāli saistīta ar romantiskā protesta patosu.

Par romantisma ietekmi poļu tautā nedrīkst spriest tikai pēc datiem par izdevumu skaitu, atsauksmēm u.c. formāliem rādītājiem. Ziņas par cilvēkiem, kuri tika vajāti par revolucionāru darbību, rāda, ka šajā vidē romantiķu darbi tika plaši lasīti, pašu pārrakstīti gāja no rokas rokā. Tā perioda revolucionāru audzināšanā dzejai, jādomā, bija pat lielāka loma nekā teorētiskiem darbiem. Poļu literatūrā romantisms nostiprināja sabiedriskā aktivitātes tendenci, šajā periodā pamatos izveidojās poļu nacionālā pašapziņa, kura turpmākajā cīņa par nacionālo atbrīvošanos nostiprinājās un pilnveidojās.

Taču jau pirms romantisma laikmeta beigām poļu literatūrā izpaudās reālisma nepieciešamība. Sabiedrībā notiekošo procesu, tās tikumu, dažādu slāņu attiecību atspoguļošanas uzdevums līdz ar kapitālisma attīstību kļuva arvien aktuālāks.

1846.--1848.~gadā auga ne tikai revolucionārā kustība, bet arī demokrātiskie noskaņojumi literatūrā, pēc šo gadu revolucionārās kustības neveiksmēm mazinājās arī romantisma loma literatūrā. 1863.--1864.~gada poļu sacelšanās sakāve nesa galu romantiskajam pasaules uzskatam, noteica jauna, var teikt, antiromantiska virziena rašanos Polijas literatūrā. Pēc 1863.--1864.~gada sacelšanās sakāves aizsākās t.s. pozitīvisma periods politikā, bet romantismu literatūrā nomainīja reālisms, attīstījās ``organiskā darba'' (\pltxti{praca organiczna}) koncepcija, kura paredzēja ar praktisku saimniecisku un kulturālu darbību, sagatavot poļus neatkarībai. Valstiskā neatkarība kļuva par tālu nākotnes mērķi un atkāpās lojalitātes pret Poliju sadalījušajām valstīm priekšā.

Tomēr bija arī cilvēki, kuri nezaudēja revolucionāros ideālus. 1882.~gadā Žēnēvā tika izdots poļu revolucionāru Varšavas un Pozenes cietumos radītā revolucionārā dzeja. Tai skaitā bija arī B.~Červiņska dzejolis ``\pltxti{Czerwony sztandar}'' (Sarkanais karogs), kurš vēlāk kļuva par populāru strādnieku dziesmu.

Valdošais virziens literatūrā kļuva \strong{reālisms}. Dažkārt šo laiku raksturo kā ``pozitīvisma periodu'', jo pozitīvisma sociāli-politiskās idejas ietekmēja daudzu rakstnieku darbību. Publicistikā un arī daiļdarbos tika propagandētas ``organiskā darba'', zemes ekonomiskās attīstības veicināšanas, feodālo šļahtiču aizspriedumu likvidācijas idejas. Saimnieciskais, sociālais un kulturālais progress tika vērtēts augstāk nekā politiskā brīvība. Literatūra ieguva didaktiskas funkcijas, norādot lasītājam uz viņa vietu un uzdevumiem sabiedriskajā dzīvē, nepieciešamību atteikties no romantiskas sapņošanas un sekot pozitīvisma principam ``zināšanas ir spēks''. Uzplaukumu piedzīvoja proza. Romānos, garstāstos un stāstos tika atainota plaša poļu sabiedrības, ieskaitot arī zemniecību un pilsētu zemākos slāņus, panorāma, parādot sabiedriskās attīstības pretrunas, jaunus sociāli-psiholoģiskos tipus.

Par reālistiskā virziena pamatlicējiem poļu literatūrā kļuva jau minētais J.~Kraševskis un J.~Korženovskis.

Rakstnieka, dzejnieka un dramaturga \strong{Juzefa~Kraševska} spalvai pieder vairāk kā 500~sējumu sacerējumu. Viņš bija viens no ražīgākajiem literātiem pasaulē, reizē arī literatūrkritiķis, izdevējs un sabiedrisks darbinieks. Milzīgajā viņa literārajā mantojumā vērtīgākie ir 144~sadzīves un 88~vēsturiski (par Polijas un Lietuvas pagātni) garstāsti un romāni. Jau 40.--50.~gados radītajiem zemnieku dzīves tēmai veltītajiem J.~Kraševska darbiem bija svarīga vieta poļu reālisma attīstībā. Viņa prozā sadzīvo gan romantiski, gan reālistiski, gan arī naturālistiski momenti.

\strong{Juzefs Korženovskis} bija pedagogs, kurš paveica lielu darbu poļu mācību iestāžu organizēšanā. Literārās darbības sākumā viņš pārstāvēja vēl romantisko virzienu, taču 40.~gados viņš kļuva par vienu no poļu kritiskā reālisma aizsācējiem. Viņa darbiem raksturīgi izvērsti šļahtiču sadzīves tēlojumi, satīra par tiem. Tiesa, sadzīves detaļu patiesums tajos sadzīvo ar klaju didaktismu.

Pozitīvisma idejas savos darbos ļoti aktīvi pauda žurnālists, rakstnieks un dramaturgs A.~Svetohovskis. Viņš bija karsts progresa, izglītības un kultūras veicinātājs, cīnījās par sieviešu, ebreju tiesību vienlīdzību, pret konservatīvismu, atpalicību un klerikālismu.

Līdz ar reālisma atnākšanu par pašu populārāko literatūras žanru kļuva romāns. Izcils poļu reālisma pārstāvis bija \strong{Boļeslavs Pruss}. Viņš gadu desmitiem nodarbojās ar žurnālistiku, bija sava laika sabiedriskās dzīves centrā. B.~Prusa literāro darbu varoņi bija visdažādāko slāņu pārstāvji, arī strādnieki un zemnieki. Sācis kā jautru humoristisku pastāstu autors, viņš drīzumā pievērsās sociālajām novelēm, parādot vienkāršo darba cilvēku~--- zemnieku, amatnieku, strādnieku, sīku ierēdņu dzīvi. B.~Prusa varoņi ir atsaucīgi, gatavi pašuzupurēties tādu pat dzīves pabērnu, kādi ir viņi paši, labā.

Rakstnieks atainoja arī naudas varas graujošo ietekmi uz mākslu. Plašu popularitāti ieguva garstāsts ``\pltxti{Placówka}'' (Priekšpostenis), jauna tipa sociāli-psiholoģiskais romāns ``\pltxti{Lalka}'' (Lelle)~--- plaša Varšavas dzīves panorāma un reizē trīs poļu ideālistu paaudžu vēsture. (Pēc romāna uzņemta poļu režisora V.~Hasa mākslas filma (1968) un televīzijas filma (1977).

Apzināta novēršanās no vēsturiskā romantisma parādījās B.~Prusa romānā ``\pltxti{Faraon}'' (Faraons, 1895--1896) par Senās Ēģiptes dzīvi, kurā skartas daudzas svarīgas problēmas: tautas loma valsts dzīvē, tās smagais stāvoklis, politiskās varas funkcionēšana valstī u.c. Romānā B.~Pruss uz vēsturiska materiāla pamata mēģinājis ne tikai iemiesot savu valsts pārvaldes ideālu, bet parādīt arī savu sabiedriskās krīzes izjūtu. Romāna sižetu sastāda izdomātas vēsturiskas personas~--- faraona Ramzesa XII cīņa ar visspēcīgo priesteru kārtu. Faraons piesaista savā pusē tautas masas, solot atvieglot tās dzīvi. Ramzess XII iet bojā, bet pie varas nonākušie priesteri ir spiesti nedaudz piekāpties tautai. Pēc romāna 1966.~gadā uzņemta režisora J.~Kavaleroviča tāda pat nosaukuma filma.

Plaša Varšavas un provinces dzīves aina parādīta romānā ``\pltxti{Emancypantki}'' (''Emancipētās sievietes'', 1890--1893). Tas liecināja par sabiedrisko attiecību sarežģītības pieaugumu, cilvēku arvien pieaugošajām rūpēm par rītdienu, kad darbs un tikums netiek pienācīgi novērtēti, godīgiem cilvēkiem arvien grūtāk kļūst sagādāt iztiku. Nožēlojams ir ``emancipēto'' sieviešu liktenis, kurām ekonomiskā nepieciešamība liek kļūt par darba armijas ierindniecēm. Pie tādiem neiepriecinošiem secinājumiem nonāca rakstnieks, kurš kādreiz optimistiski skatījās uz buržuāziskās civilizācijas perspektīvām.

Tuva pēc gara B.~Prusam bija rakstniece \strong{Elīza Ožeškova}. Par vienu no galvenajām viņas daiļradē kļuva zemniecības tēma. 60.--70.~gados tapušo viņas romānu, kuros plaši radīta šļahtas un buržuāzijas dzīve, māksliniecisko vērtību mazina trūkumi kompozīcijā, pārlieka sentimentalitāte. Toties 80.~gados viņas talants sasniedza radošo pilnbriedumu. Labākajā savā romānā ``\pltxti{Nad Niemnem}'' (``Pie Nemunas''), kurš pamatoti tiek atzīts par vienu no izcilākajiem poļu literatūrā, rakstniece rādīja šļahtas dzīvi bijušās Lietuvas lielkņazistes teritorijā. Tā darbība norisa pēc 1864.~gada agrārās reformas. Ar lielu sirsnību, bet visai skopi, ievērojot cenzūras noteikumus, rakstniece rakstīja arī par 1863.~gada sacelšanos. E.~Ožeškova izteica domu par šļahtas vienības ar tautu nepieciešamību. Kopumā bija radīts pilnasinīgs reālistisks darbs.

Milzīgu popularitāti Polijā ieguva ``pozitīvists'' \strong{Henriks Senkevičs}. Viņš no 1883. līdz 1888.~gadam publicēja tautā iemīļoto triloģiju ``\pltxti{Ogniem i mieczem}'', ``\pltxti{Potop}'', ``\pltxti{Pan Wołodyjowski}'' (''Ar uguni un zobenu'', ''Plūdi'', ''Pans Volodijevskis''). 90.~gados viņa darbi arī rādīja, ka arī XIX~gadsimta beigu krīzes izjūtas literatūrā var vest pie jauniem mākslinieciskiem risinājumiem. H.~Senkeviča romāns ``\pltxti{Bez dogmatu}'' (``Bez dogmāta'', 1889--1890) īsā laikā ieguva visas Eiropas ievērību, kļuva par karstu strīdu izraisītāju. Pirmoreiz poļu literatūrā parādījās varonis, kurš iemiesoja sevī pozitīvo ideālu krīzi. Romāna atmosfēra liecināja, ka ``pozitīviskais'' etaps ir izsmēlis sevi, ka radušies daudzi jautājumi, uz kuriem inteliģence vēl nezināja atbildes. Toties nākamais viņa romāns ``\pltxti{Rodzina Połanieckich''} (''Poļeņecku ģimene'', 1893--1894) daļā sabiedrības izsauca neapmierinātību, jo autors sniedza pārburžuāziskotās šļahtas apoloģiju. Pēc krievu rakstnieka A.~Čehova vērtējuma ``romāna mērķis ir ieaijāt buržuāziju tās zelta sapņos'', nomierināt to ``ar domu, ka ir iespējams gan kapitālā iedzīvoties, gan nevainību saglabāt, reizē rīkoties zvēriski un būt laimīgam''. 90.~gados iznāca arī divi kārtējie H.~Senkeviča vēsturiskie romāni ``\latxti{Quo Vadis?}'' (''Kurp ej?'', 1896) un ``\pltxti{Krzyżacy}'' (''Krustneši'', 1900), kuri gan salīdzinājumā ar rakstnieka agrākajiem darbiem šajā jomā neienesa neko jaunu, taču izpelnījās lielu ievērību. Vācu kritiķi pat rakstīja, ka daļai H.~Senkeviča darbu nav nekā kopīga ar nopietnu literatūru, jo tie saturot naida un rīdīšanas plūdus pret vāciešiem. H.~Senkevičs ``par izciliem nopelniem episka rakstura darbu radīšanā'' 1905.~gadā kļuva par Nobela prēmijas laureātu literatūrā (Nevis, kā bieži kļūdaini tiek apgalvots, par romānu ``\latxti{Quo vadis?}'', kurš bija veltīts kristietības pirmsākumiem.) Mūsdienās daudzi H.~Senkeviča darbi ir ekranizēti.

Reālisma valdošais stāvoklis literatūrā neveicināja dzejas attīstību, tomēr var izcelt vismaz divus ievērojamus XIX gadsimta otrās puses dzejniekus: Ā.~Asniku un M.~Konopņicku.

\strong{Ādams~Asniks} bija 1863.~gada sacelšanās dalībnieks, debitēja kā romantisma tradīciju turpinātājs, taču viņu ietekmēja arī pozitīvisms. Ā.~Asnika dzeja ir bagāta sabiedriskiem motīviem, apliecina uzticību patriotiskajiem ideāliem. Īpaši tas atklājas sonetu ciklā ``\pltxti{Nad głębiami}'' (''Virs dzelmes'', 1883--1894). Ā.~Asniks bija arī formas meistars. H.~Senkevičs par viņa dzejoļiem teica, ka tie ``atgādina Benvenuto Čellini (izcils Renesanses laikmeta itāļu skulptors un juvelieris~--- V.Š.) darbus''. Ā.~Asniks rakstīja arī noveles un drāmas, taču tās mūsdienās ir pilnīgi aizmirstas.

Propagandējot kalpošanu tautai, lauku un pilsētu trūcīgo slāņu dzīvi savās novelēs tēloja \strong{Marija~Konopņicka}. Poļu emigrantu dzīves grūtībām viņa piešķīra visas tautas ciešanu ceļa simbolisko nozīmi. Visbiežāk viņas daiļdarbu varoņi bija bērni, kas kāpināja to emocionalitāti. M.~Konopņickas sacerētie teksti vēl viņas dzīves laikā kļuva par tautas dziesmām. Pilsoniski motīvi parādījās viņas dzejā, veltītā 1905.~gada revolūcijai. Poēmā ``\pltxti{Pan Balcer w Brazylii}'' (Pans Balcers Brazīlijā, 1892--1910) kulminācijas punkts ir strādnieku demonstrācijas attēlojums ostā, kurā piedalās arī bijušie poļu zemnieki-emigranti.

Pozitīvisma krīze XIX gadsimta 80.~gadu beigās un naturālisma izcelsme iezīmēja arī jaunu laikmetu poļu literatūrā. \strong{XIX gadsimta beigās} Polijā, tāpat kā Rietumeiropā, bija \strong{vērojama atiešana no pozitīvisma un reālisma}. Gadsimta pēdējā desmitgadē un XX gadsimta pirmajās divās desmitgadēs poļu kultūrā izveidojās jauns, neoromantisks virziens. Gan tas, gan vesels posms poļu literatūras attīstībā, ieguva nosaukumu ``\pltxti{Młoda Polska}'' (\strong{Jaunā Polija}). Kā atzīmē poļu mākslas un literatūras vēsturniece D.~Ļeščinska, izšķiroša nozīme jaunas literātu apziņas veidošanā un tās atiešanā no pozitīvisma bija diviem faktoriem: sociālo konfliktu nobriešanai un kaujinieciska patriotisma atdzimšanai. Jaunais virziens attīstījās ciešā mijiedarbībā ar tādiem mākslinieciskiem strāvojumiem kā naturālisms, impresionisms, simbolisms, ekspresionisms un neoklasicisms. Jau nosaukto virzienu daudzums runā par to, ka ``Jaunās Polijas'' estētiskā programma bija pretrunīga, tā atspoguļoja kā daļas radošās inteliģences elitāros noskaņojumus, tā individuālistiskus protestus pret buržuāzisko dzīves veidu un tās kultūru. Atšķirībā no romantisma skolas ``Jaunās Polijas'' pārstāvjiem raksturīga arī literatūras nacionālo uzdevumu noliegšana un to vietā sekošana principam ``\frtxti{L’art pour l’art}'' (franču Māksla mākslai). Aiz šīs devīzes kā burvju formulas tolaik savos manifestos slēpās ļoti daudzi mākslinieki. Apskatot sava laika poļu literatūru, XIX un XX gadsimtu mijā strādājušais poļu literatūras kritiķis J.~Lorentovičs kādā no savām esejām rakstīja: ``Jauns~--- tas ir tas, kurš norobežojas no visāda veida rutīnas, no dažādiem novītušiem literāriem aizspriedumiem, kurš nemitīgi atkaro savu individualitāti''. Šī savas individualitātes ``atkarošana'' likumsakarīgi noveda pie tā, ka daudzu ``jaunpoļu'' radošo stilu nevarēja pieskaitīt kādam vienam noteiktam virzienam. Kopumā viņus vienoja dzīves disharmonijas sajūta, radošās brīvības un mākslinieka neatkarības apliecināšana, saistīta ar jaunu mākslinieciskās izteiksmes līdzekļu meklējumiem, sadumpošanās pret mietpilsonisko sabiedrību un tās morāli. ``Jaunpoļu'' radošais virziens parādījās gan literatūrā (īpaši dzejā), mūzikā, teātra mākslā.

Centrālā vieta ``Jaunās Polijas'' prozā pieder \strong{Stefanam Žeromskim}, jauna tipa romāna, kurā reālistisks, dažkārt pat naturālistisks skatījums kontrastēja ar liriski sakāpinātu skanējumu, autoram. Cēlies no šļahtiču aprindām, bet agri zaudējis vecākus, spiests atstāt universitāti, S.~Žeromskis pelnīja iztiku kā mājskolotājs, badojās, iepazina darba cilvēku dzīvi dažādās nomalēs, klaiņoja pa ārvalstīm. Nebūdams pēc politiskajiem uzskatiem ne sociālists, ne arī radikālis, viņš nepiedalījās revolucionārā darbībā, taču šļahtiča patriotisms noteica viņa daiļrades kaujinieciskumu, radikālismu. S.~Žeromskis radošo darbību uzsāka apstākļos, kad poļu valdošajos slāņos valdīja samiernieciski noskaņojumi un Polijas valstiskuma atjaunošana likās neaizsniedzama utopija. Tādos aiz robežas izdotos darbos kā ``\pltxti{Rozdziobią nas kruki, wrony}'' (Noknābs mūs kraukļi, vārnas, 1895) vai ``\pltxti{Syzyfowe prace}'' (``Sizifa darbs'', 1898), viņš spilgti raksturoja tā laika mietpilsoniskās poļu inteliģences morālo pagrimumu. Raksturīgi, ka meklējot iespējas atkal iededzināt patriotisko degsmi, rakstnieks savus literāros varoņus ved ne tikai uz vēsturiskiem arhīviem, bet arī uz kapsētu. Veci, nekopti agrāko sacelšanās dalībnieku kapi daudzos viņa darbos kalpo kā tādas nacionālas jūsmas avots, kādu vairs nevarēja atrast reālajā dzīvē. S.~Žeromskis juta dziļu sašutumu pret šļahtu, kura, vadoties no sava šķiriskā egoisma, atstūma zemniekus un ar to pazudināja Polijas neatkarību. Pie šī motīva rakstnieks savos vēsturiskajos sacerējumos atgriezās vēl un vēl. Vēsturiskajā romānā ``\pltxti{Popioły}'' (Pelni, 1902--1903) S.~Žeromskis spilgti un patiesi atainoja Napoleona I laika poļu leģionāru traģēdiju. Romāns rāda, kā poļu leģionāri no cīnītājiem par savas zemes brīvību kļuva par Napoleona iekarošanas karu dalībniekiem un bija spiesti cīnīties pret itāļu un spāņu brīvības cīnītājiem. Romāns ir ekranizēts A.~Vajdas tāda pat nosaukuma filmā 1965.~gadā.

Arī citos darbos: ``\pltxti{Duma o hetmanie}'' (Doma par hetmani, 1905) u.c. S.~Žeromskis šaustīja šļahtas šķirisko tuvredzību. Pie tam viņš nestādīja mērķi revolucionizēt masas, viņš tikai vēlējās izlabot to cilvēku tikumus, kuru rokās bija vara un nauda. Taču viņa atmaskojumi faktiski rādīja, ka nav pamata cerībām izlabot apspiedējus. Rakstnieks pirmais parādīja poļu sabiedrības dziļo noslāņošanos. Jau pati dziļo sociālo kontrastu iezīmēšana piešķīra viņa darbiem revolucionāru skanējumu, pēc 1905.~gada revolūcijas sakāves tapušajos sacerējumos šī tendence izpaudās vēl asāk. S.~Žeromska darbos priekšplānā gandrīz vienmēr stāv varonis-vienpatnis, kurš cenšas radīt apstākļus Polijas atjaunošanai, neraugoties uz apkārtējo buržuāzijas savtīgo brīvības interešu nodevību un proletariāta internacionālismu.

Cits ievērojams rakstnieks \strong{Vladislavs Raimonts} savos agrīnajos darbos augošajai pilsētu civilizācijai pretstatīja lauku tīrību un morālo veselību, taču labākie viņa stāsti rāda, ka rakstnieks redz kapitālisma agrīnās stadijas nežēlīgo likumu iespiešanos arī laukos. Par viņa daiļrades virsotni tiek uzskatīts romāns ``\pltxti{Chłopi}'' (Zemnieki), kurā parādīta lauku dzīves sociālās pretrunas un psiholoģiskie konflikti. Par šo romānu autors 1924.~gadā saņēma Nobela prēmiju literatūrā. Vācu kultūras vēsturnieki uzsvēruši, ka V.~Raimonta darbi vācu tulkojumos guva plašāku piekrišanu pasaulē nekā dzimtenē. Arī Nobela komitejai izskatīšanai bija iesniegts viņa darba vāciskais tulkojums.

``Jaunās Polijas'' virzienam bija raksturīga arī simboliskās dzejas un dramaturģijas, no vienas puses, un naturālistiskās prozas un dramaturģijas, no otras puses, līdzāspastāvēšana. Jaunais virziens visspilgtāk izpaudās Prūsijai piederošajā Toruņā (poļu \pltxti{Toruń}, vācu \detxti{Thorn}) dzīvojušā romānista un dramaturga S.~Pšebiševska, kurš sāka daiļradi vācu, bet vēlāk turpināja arī poļu valodā, un Krakovā dzimušā dzejnieka, rakstnieka, dramaturga K.~Tetmeijera darbos. Abi viņi pauda dekadences un garīgās krīzes idejas.

\strong{Staņislava~Pšebiševska} raksts ``\latxti{Confiteor}'' (latīņu: Es atzīstu, 1899) kļuva par ``Jaunās Polijas'' manifestu. Autors tajā paziņoja: ``Mākslai nav nekāda mērķa, tā ir mērķis pats par sevi, absolūts, jo tā atspoguļo absolūtu dvēseli'', ``tendencioza māksla, didaktiski-tikumiska māksla un māksla--izklaide, patriotiskā māksla, māksla, kam ir kaut kāds morāls vai sabiedrisks mērķis, pārstāj būt māksla''. S.~Pšebiševska romāns ``\latxti{Homo sapiens}'' (1901) un citi darbi ir naturālisma, simbolisma un reālisma, vēlāk~--- XX gadsimta pirmajā gadu desmitā~--- arī ekspresionisma sakausējums.

\strong{Kazimira~Tetmeijera} labākais veikums prozā ir pēc poļu kalniešu leģendām tapušais stāstu cikls ``\pltxti{Na skalnym Podhalu}'' (``Klinšainajā Podhalā'', 1903--1910), kuram pieslejas vēsturiskais romāns ``\pltxti{Legenda Tatr}'' (Tatru leģenda, 1910--1911) par XVII gadsimta vidus kariem un zemnieku sacelšanos. Viņš bija arī vairāku drāmu autors. Spilgti ``jaunās poēzijas'' idejiskās un mākslinieciskās tendences pauda K.~Tetmeijera lirika. To caurauž individuālistiski noskaņojumi, mākslinieka-radītāja pretnostatījums mietpilsoniskajam pūlim, tajā arvien no jauna atkārtojas bezgalīgas un nedzēšamas grūtsirdības, dvēseles noguruma motīvi.

No ``Jaunās Polijas'' \strong{dzejniekiem} var izdalīt J.~Kasproviču, T.~Micinski un B.~Lesmjanu. XIX un XX gadsimtu mijā dzejā parādījās arī novirzieni, kas vismaz daļēji jau izgāja arī ārpus ``Jaunās Polijas'' strāvojuma.

\strong{Jans Kasprovičs} ilgstoši strādāja Lembergā (Ļvovā), sākotnēji kā žurnālists, vēlāk kā mācību spēks universitātē, bija Lembergas Universitātes rektors. Līdz 90.~gadiem viņš reālistiski tēloja zemnieku dzīvi, poēma ``\pltxti{Chrystus}'' (Kristus, 1890), meklēja sociālo pretrunu atrisinājumu cilvēces ētiskā atdzimšanā. Tad Lembergas mākslinieciskās bohēmas ietekmē viņa daiļradē notika modernistisks lūzums un nostiprinājās simboliskā poētika, idejas par cilvēka ciešanu nepieciešamību. Viņa dzejas ciklā ``\latxti{Salve Regina}'' (''Sveika karaliene'') (1902), poētiskās prozas grāmatā ``\pltxti{O bohaterskim koniu i walącym się domie}'' (''Par varonīgo zirgu un brūkošo māju'', 1906), krājumā ``\pltxti{Księga ubogich}'' (''Nabagu grāmata'', 1916) u.c. valda reliģiski-mistiskas noskaņas.

Reālistiskās poētikas elementi stipri jūtami dzejnieka, prozaiķa un dramaturga \strong{Tadeuša Micinska}, romānos, īpaši romānā ``\pltxti{Xiądz Faust}'' (''Ksendzs Fausts'', 1913). Katoļu garīdznieka (ksendza) Fausta fantastikas pilnajos dzīves stāstos iepītas Eiropas un Polijas dzīves realitātes. T.~Micinska dzeja ir pilna simboliskām un ekspresionistiskām tendencēm, demonstrē dzejnieka aktīvi-dumpīgu pozīciju. Īstenības deformācija, uz sapņu un realitātes robežas esoši tēli, kam jārada haosa un ciešanu apjausma, tuvina dzejnieku vēlākajiem sirreālistiem un ekspresionistiem. Revolūcijas laikā T.~Mičinskis uzrakstīja simboliski-ekspresionistisku drāmu ``\pltxti{Kniaź Patiomkin}'' (''Kņazs Potjomkins'', 1906), veltītu matrožu dumpim uz bruņukuģa ``Potjomkins''. Drāmā ir arī ainas, kurās atbilstoši īstenībai tiek rādīta sacelšanās gaita, tomēr šķiru cīņas idejai autors pretstata ticību kā pamatu tautas morālai atdzimšanai.

Ievērojamākais ebreju izcelsmes poļu simbolisma dzejnieks \strong{Boļeslavs Lesmjans}, kurš rakstīja gan poļu, gan krievu valodā, debitēja jau ``Jaunās Polijas'' norieta laikmetā, publicējoties kā poļu, tā krievu simbolistu žurnālos. Pilnā apjomā viņā daiļrade gan uzplauka XX gadsimta 20.--30.~gados, taču galvenie viņa poētikas principi tapa jau gadsimta sākumā. Dzejnieks atainoja pasauli kā pastāvīgas atjaunotnes procesu. Sācis jaunradi kā poļu simbolisma pārstāvis, B.~Lesmjans kļuva par vienu no oriģinālākajiem XX gadsimta poļu liriķiem. Viņa darbos fantastiski pasaku tēli mijas ar reālās dzīves un poētiskām dabas ainām. Ironijas un patētikas, naivi-jūsmīgas pasaules uztveres un dziļa filozofiskuma savīšanās dzejnieka darbos bieži realizējas ar misticismu iekrāsotos slāvu folkloras tēlos vai austrumu leģendu motīvos. Dabas tīrība un bagātība kalpo dzejniekam kā antitēze cilvēku kaislībām, viņa likteņa traģismam un nāves neizbēgamībai. B.~Lesmjanu uzskata par saikni starp simbolistiem un ekspresionistiem, arī sirreālistiem.

% page 278



\strong{XIX un XX gadsimtu mijā} atkal \strong{atdzima romantisms}, īpaši to dzejnieku vidū, kuri grupējās toreiz uzplaukušajā kurortcentrā Zakopanē. Īpaši Krakovas bohēma mīlēja uzturēties šai Tatru kalnu pilsētiņā un baudīt dabu. Tā rezultātā radās daudzi liriski darbi, veltīti šīs apkārtnes skaistajai dabai. Gadsimtu mijā hronoloģiski pēdējais dzejā izpaudās neoklasicisms. Izcilākais dzejnieks saistīts ar šo novirzienu bija \strong{Leopolds Staffs}. Savā pirmajā dzejas krājumā ``\pltxti{Sny o potędze}'' (``Sapņi par varenību'', 1901) viņš vēl sekoja ``Jaunās Polijas'' simboliskajiem meklējumiem, bet reizē jau uzstājās pret dekadenci (šeit~--- mākslas strāva, kas atsakās no reālisma, sociālajām un nacionālajām idejām, pievēršas anormālajam, mistiskajam u.tml.), slavinot pozitīvu attieksmi pret dzīvi un aktīvu darbību. L.~Staffa poētiskā filozofija, atšķirībā no simbolisma, ir racionāla, tajā ir stiprs optimisma kodols. Turpmākajos krājumos ``\pltxti{Dzień duszy}'' (``Dvēseles diena'', 1903), ``\pltxti{Ptakom niebieskim}'' (''Debesu putniem'', 1905), ``\pltxti{Gałąź kwitnąca}'' (``Ziedošais zars'', 1908), ``\pltxti{Uśmiechy godzin}'' (``Mirkļu smaids'', 1919) u.c. bija jūtama tendence iemiesot viņa dzejas ideālu, kas līdzinātos pēc antīkajiem paraugiem, apdziedātu dzīves skaistumu un harmoniju. L.~Staffa dzeja demonstrēja orientāciju uz klasiskajām tradīcijām, interesi par antīko un renesanses laikmeta kultūru, mūžīgajām ētiskajām un estētiskajām vērtībām.

Strādnieku kustības sākums veicināja t.s. \strong{``proletāriskās dzejas''} pirmo paraugu rašanos. 1881.~gadā poļu dzejnieks-sociālists B.~Červinskis pēc 1871.~gada Parīzes komunāru dziesmas ``\frtxti{Le drapeau rouge}'' (``Sarkanais karogs'') parauga uzrakstīja dzejoli ar nosaukumu ``\pltxti{Krew naszą długo leją katy}'' (``Bendes ilgi lēja mūsu asinis''), kurš gan vēlāk ieguva citu nosaukumu~--- ``\pltxti{Czerwony Sztandar}'' (''Sarkanais karogs''). Mūziku tam radīja J.~Kozakevičs. Dziesma tika dziedāta streiku, demonstrāciju laikā, kļuva par Polijas Sociālistiskās partijas himnu, ieguva starptautisku popularitāti, tulkota 12~valodās. Krievu valodā to pārtulkoja G.~Kržižanovskis, vācu~--- R.~Luksemburga.

Otras populārākās dziesmas tekstu jau 1879.~gada uzrakstīja Varšavas citadelē ieslodzītais poļu sociālists V.~Svecickis. 1883.~gadā tas tika publicēts žurnālā ``\pltxti{Proletariat}''. Kā melodija tika izmantots jau 1863.~gada poļu sacelšanās laikā popularitāti ieguvušais ``Zuāvu maršs'' (1863.--1864.~gadā Polijā pēc franču koloniālā karaspēka parauga par zuāviem mēdza saukt brīvprātīgo-nāvinieku daļas), tad maršu dziedāja ar V.~Voļska tekstu, bet par īstu strādnieku cīņas dziesmu tā kļuva 1905.~gada revolūcijas laikā. Tad to arī nosauca par ``1905.~gada Varšavjanku'' (''\pltxti{Warszawianka 1905 roku}''), lai atšķirtu no 1831.~gada Varšavjankas. Dziesma kļuva ļoti iecienīta arī krievu revolucionāru vidū. (Tāpat kā ``\pltxti{Czerwony Sztandar}'', krievu valodā to Butirkas cietumā (\rutxti{Бутырская тюрьма, ``Бутырка''}) pārtulkoja G.~Kržižanovskis.) Spāniski tulkota tā kļuva par spāņu anarhosindikālistisko arodbiedrību (Confederación Nacional del Trabajo) himnu, Grieķijā zem nosaukuma ``\eltxti{Άνεμοι Θύελλες}'' (''Vēji un vētras'') komunisti to dziedāja itāļu un vācu okupācijas un pilsoņu kara (1944--1949) laikā. Latviešu valodā 1905.~gada revolūcijas laikā to tulkoja Jānis Akuraters.

Dziesmas 1.~pants un piedziedājums:

\vspace{1.5em}

\noindent
\begin{minipage}{0.65\textwidth}
poļu valodā:\\
~\\
\pltxti{Śmiało podnieśmy sztandar nasz w góre,\\
Choć burza wrogich żywiołów wyje,\\
Choć nas dziś gnębią siły ponure,\\
Chociaż niepewne jutro niczyje.\\
O, bo to sztandar całej ludzkości,\\
To hasło święte, pieśń zmartwychwstania,\\
To tryumf pracy, sprawiedliwości,\\
To zorza wszystkich ludów zbratania!\\
~\\
Naprzód Warszawo!\\
Na walkę krwawą, (bis)\\
Świętą a prawą!\\
Marsz, marsz, Warszawo!
}
\end{minipage}
\hspace{1em}
\begin{minipage}{0.6\textwidth}
latviešu valodā:\\
~\\
Ienaida negaisi pāri mums krājas,\\
Tumsības spēki mūs spaida un māc.\\
Kaujā pret naidnieku biedri jau stājas~---\\
Cīņai vēl gala vis neparedz skats;\\
Tomēr mēs pacelsim lepni un brīvi\\
Karogu cīņā par strādnieku dzīvi,\\
Zem viņa tautas lai salasās barā\\
Labākas pasaules brīvības karā!\\
~\\
Uz asins kauju,\\
Uz taisnības kauju,\\
Uz priekšu trauc\\
Darbaļaužu daudz.
\end{minipage}

\vspace{1.5em}

% page 280


Ar sabiedriski-politiskās situācijas maiņu, īpaši 1905.--1907.~gadā, notika daudzu ``Jaunās Polijas'' programmatisko nostādņu pārvērtēšana. Var pat teikt, ka ``Jaunās Polijas'' periods dalās divos posmos~--- pirms un pēc 1905.~gada revolūcijas. Nebija neviena rakstnieka, kura daiļradē revolūcijas ietekmē nenotiktu izmaiņas literāro darbu tematikā un formā. Tās atspoguļoja literātu cerības, ilūzijas un vilšanās.

Grūti noteikt ``Jaunās Polijas'' noslēguma laiku. Kaut daudzi vēsturnieki uzskata, ka par tādu var uzskatīt 1918.~gadu, tomēr šis viedoklis nav neapšaubāms. Daudzi no izcilākajiem ``jaunpoļiem'': S.~Žeromskis, V.~Raimonts, J.~Kasprovičs u.c. aktīvi darbojās arī starpkaru periodā.

Vēl piebilstams, ka šai periodā par angļu literatūras klasiķi kļuva poļu izcelsmes rakstnieks \strong{Džozefs Konrāds}). Angļu un franču piedzīvojumu romānu iespaidā poļu zēns Juzefs izvēlējās kļūt par jūrnieku, devās uz Marseļu, kur uzsāka darbu jūrā. 1894.~gadā viņš apmetās Londonā, nolēmis beigt jūrnieka gaitas. Jau 1895.~gadā viņš debitēja literatūras jomā ar pirmo romānu. Sekoja citi, rakstnieka popularitāte auga. 1924.~gadā D.~Konrāds atteicās no britu bruņinieka kārtas piešķiršanas, jo viņš jau bija poļu šļahtičs ar savu dzimtas ģerboni.

Daudzsējumu \strong{poļu literatūras vēsturi}, kā arī lielu monogrāfiju par Ā.~Mickēviču radīja Lembergas universitātes profesors P.~Hmeļevskis.

XVIII~gadsimtā (Apgaismības laikmetā) svarīgākie \strong{mākslas} centri izveidojās pie karaļa galma, lielo magnātu muižās. Polijas sadalīšana, kas izpaudās visas poļu sabiedriskās un kultūras dzīves jomas, mākslas attīstību īpaši ietekmēja ar to, ka tā zaudēja svarīgu attīstības stimulu~--- karaļu kā mecenātu atbalstu. Izzuda krāšņais galms Varšavā, daudzi mākslinieki bija spiesti savu darbību saistīt ar citiem kultūras centriem. Zināmā mērā karaļu lomu mākslas atbalstīšanā pārņēma magnāti. Piemēram, Čartorijski apzināti centās pārvērst savu rezidenci Pulavos (\pltxti{Puławy}, tagad Lubļinas vojevodistē) par vispolijas nozīmes patriotisku centru. Pagātnes kults, romantiskie ideāli spilgti izpaudās Radzivillu rezidencē ``Arkādija'' netālu no Varšavas. Arī citas magnātu dzimtas ievērojami atviegloja dzīvi poļu māksliniekiem. Tomēr mākslas atbalstam ziedoto magnātu līdzekļu nepietika lai radītu apstākļus intensīvai arhitektūras attīstībai, sabiedrisku ēku celtniecībai. Mākslas institūtu, plaša muzeju tīkla, gleznu galeriju, izstāžu zāļu trūkums bremzēja mākslas attīstību. (Tāpēc par izcilu notikumu kļuva mākslas darbu izstāžu ēkas ``\pltxti{Zachęta}'' (Stimuls, pilns nosaukums~--- \pltxti{Zachęta Narodowa Galeria Sztuki}) atklāšana Varšavā 1898.~gadā (arhitekts S.~Šillers). Savus darbus šeit izstādīja pazīstamākie poļu mākslinieki.)

Vēl nozīmīgāks nekā magnātu-mecenātu atbalsts mākslas attīstībai kļuva plašākas sabiedrības brīvprātīga aizbildniecība pār mākslu un māksliniekiem. Mākslā, tāpat kā literatūrā, sabiedriskais mecenātisms atbalstīja nacionālās, patriotiskās kultūras tendences, darbojās ``kalpošanas nācijai imperatīvs''. Patriotiskais strāvojums mākslā īpaši pieauga XIX gadsimta pirmajā trešdaļā. Par to, piemēram, liecināja gleznotāju V.~Vaņkoviča, M.~Stahoviča daiļrade. 1832.~gadā par sabiedrības savāktiem līdzekļiem Varšavā tika uzcelts piemineklis J.~Poņatovskim. Tā autors bija izcilais dāņu skulptors B.~Torvaldsens. Viņa darinājums bija arī par sabiedrības saziedotajiem līdzekļiem 1829.~gadā turpat Varšavā uzceltais piemineklis astronomam N.~Kopernikam. Brīžos, kad Polijā eksistēja valstiski veidojumi (Varšavas hercogiste, Polijas karaliste autonomijas periodā) mecenāta lomu uzņēmās valsts. Tā, piemēram, Polijas karalistes laikā tika uzcelta slavenā Belvedēra pils (\pltxti{Pałac Belweder}). Ēku 1824.~gadā klasicisma stilā projektēja arhitekts J.~Kubickis par Krievijas valdības līdzekļiem lielkņaza Konstantīna vajadzībām. (Mūsdienās pilī atrodas viena no Polijas prezidenta rezidencēm.)

Poļu gleznotāji un skulptori strādāja ne tikai Varšavā, kur 1816.--1831.~gadā pie vietējās Universitātes darbojās speciāla daiļmākslu nodaļa (\pltxti{Oddział Sztuk Pięknych}), ne tikai Krakovā, kur 1818.~gadā tika atvērta Mākslas akadēmija (\pltxti{Akademia Sztuk Pięknych w Krakowie}), bet arī Viļņā un Lembergā, kur gadsimtu mijā izveidojās patstāvīgi mākslas centri.

Blakus mecenātu darbībai mākslinieki arī saistījās ar mākslas patērētājiem, izmantojot tirgus attiecības. Tika rīkotas izstādes, izpārdošanas, darbojās mākslas saloni. Taču tēlotājmākslā nebija nekā līdzīga plašajiem strīdiem literatūrā. Formas, mākslinieciskās kvalitātes jautājumi maz kam bija svarīgi, mākslu tā laika sabiedrība pirmkārt uzlūkoja no pragmatiskā pielietošanas viedokļa. Tai vēl bija jāpierāda nepieciešamība sabiedrībai, gatavība kalpot cīņai par nacionālo atbrīvošanos. Tā bija savdabīga maksa par sabiedrības gatavību atbalstīt mākslu.

Viens no poļu literātiem S.~Gosčinskis 1841.~gadā enerģiski iestājās par nacionālo tematiku (nacionālo ainavu, nacionālo tipu) atveidošanu, reizē uzsverot poļu mākslas un atbrīvošanās kustības saikni, apgalvojot, ka mākslinieki ir tautas sastāvdaļa, kas veido tās likteni. Cits poļu rakstnieks un kritiķis M.~Grabovskis daiļrades tautiskumu saskatīja saistībā ar reliģiozitāti, sludināja, ka poļu mesiānisma idejai jāizpaužas arī mākslā. Īpašu asumu jautājums par poļu glezniecības lomu izraisījās sakarā ar 1857.~gadā publicēto ebreju ģimenē dzimušā un Rietumeiropā ilgstoši dzīvojušā, Parīzē mītnes zemi atradušajai poļu emigrācijai tuvā publicista J.~Klačko rakstu, kurā autors noraidīja pat poļu skolas glezniecībā pastāvēšanas iespējamību, norādot, ka poļu tauta visā savas pastāvēšanas vēsturē nav demonstrējusi spējas tēlotājmākslas jomā. Viņš paziņoja: ``Slāvi~--- mēs varam būt tikai vārda meistari.'' Vājie, salīdzinājumā ar literatūru, poļu tēlotājmākslas sasniegumi, kā šķita, deva pamatu tik pesimistiskam vērtējumam. Raksts izsauca polemiku, kurā tika asi kritizēta J.~Klačko teorija. Diskusija parādīja, ka poļu sabiedrība savas nacionālās tēlotājmākslas pastāvēšanu vērtē kā vienu no nacionālas kultūras saglabāšanas un attīstības priekšnoteikumiem.

Var teikt, ka XIX gadsimta 40.--60.~gados Polijā kristalizējās vienota nacionālā skola \strong{glezniecībā}. Tai raksturīga žanru daudzveidība, izteiksmes līdzekļu bagātība. No 1864.~gada Varšavā darbojās V.~Gersona zīmēšanas skola, Krakovā pastāvēja daiļo mākslu skola.

XIX gadsimta pirmajā trešdaļā glezniecībā vadošais virziens vēl bija akadēmiskais \strong{klasicisms}, bet to pakāpeniski nomainīja poļu \strong{romantisms}, kurš atspoguļoja šļahtas revolucionāro virzību. Augstākie sasniegumi tika gūti portretu glezniecībā un (nedaudz vēlāk) vēsturiskajā glezniecībā. Romantismam poļu tēlotājmākslā tāpat kā literatūrā saknes sniedzās Apgaismības laikmetā.

Starp šī laika māksliniekiem portretu glezniecībā izcēlās \strong{Antonijs Brodovskis}. Viņa darbi ir vienkārši mākslinieciskās izteiksmes līdzekļos un reālistiski, bet reizē apdvesti ar romantismu, bagāti saturā.

Batālā žanra glezniecībā ar panākumiem darbojās \strong{Aleksandrs Orlovskis}. Viņš 1802.~gadā pārcēlās uz Pēterburgu, baudīja tur Ā.~Čartorijska un lielkņaza Konstantīna Pavloviča labvēlību, 1809.~gadā kļuva par Krievijas Mākslas akadēmijas locekli, radīja 1812.~gada Krievijas Tēvijas kara varoņu portretu galeriju, taču pēc Konstantīna Pavloviča nāves nonāca materiālās grūtībās un mira trūkumā. A.~Orlovskis tiek vienlaikus uzskatīts gan par poļu, gan krievu gleznotāju.

Ievērojama figūra poļu glezniecībā bija \strong{Pjotrs~Mihalovskis}, ļoti spilgts gleznotājs-romantiķis, franču mākslinieku T.~Žeriko un E.~Delakrua sekotājs. Viņa otai pieder daudz veiksmīgu kauju ainu, visbiežāk veltītu Napoleona karu vēsturei. Tāpat liela mākslinieciskā vērtība ir viņa sadzīviskajām gleznām, arī zemnieku portretiem. Gleznotājs lieliski rīkojās ar akvareļa krāsām. Sadzīves ainiņas, radītas šai tehnikā, raksturojās ar virtuozu izpildījumu.

XIX gadsimta 30.--40.~gadi poļu tēlotājmākslā raksturojas jau ar vispārēju \strong{reālisma} uzplaukumu, tā nostiprināšanos kā portretā, tā vēsturiskajās gleznās, tā arī dzimstošajā sadzīves žanrā.

40.--60.~gados turpinājās radošā procesa un tēlotājmākslas uztveres demokratizācija. Daudzi gleznotāji centās iedziļināties vienkāršā cilvēka iekšējā pasaulē. Tādas iezīmes bija raksturīgas J.~Pivarska, J.~Glovacka darbiem.

Jau šajos 40.--60.~gados Polijā radās ievērojami gleznotāji, kuri ieguva atpazīstamību visā Eiropā. Viens no tādiem bija \strong{H.~Rodakovskis}~--- lielisks portretists, kurš pievērsās arī vēsturiskajai glezniecībai, sadzīves tematikai, dekoratīvajai glezniecībai. Gadsimta vidū viņa portreti ļoti augsti tika vērtēti Parīzes mākslinieciskajās aprindās.

Pasaules slavu ieguva poļu izcelsmes krievu gleznotājs \strong{Henrihs Semiradskis}. Jau mācoties Pēterburgas Mākslas akadēmijā, viņš ``savāca'' visus iespējamos apbalvojumus glezniecības jomā, bet 1870.~gadā par gleznu ``Maķedonijas Aleksandrs un viņa ārsts Filips'' ieguva zelta medaļu un ārzemju stipendiju sešiem gadiem. Turpmāk viņš lielākoties dzīvoja un strādāja ārzemēs, taču bieži atgriezās Pēterburgā, kur baudīja cara galma labvēlību.

Viens no labākajiem zemnieku tēmas gleznotājiem bija \strong{Aleksandrs Kotsis}. Tautas ikdienas dzīve bija viņa daiļrades pamatā. Mākslinieka vēlme likt skatītājiem domāt par vecākus zaudējušu bērnu vai līdz galējai nabadzības robežai nonākušas ģimenes likteni ļāva viņam radīt satraucošas cilvēciskas drāmas. Viņš lieliski prata izmantot gaismas un krāsu risinājumus gleznas emocionālās iedarbības pastiprināšanai.

Sarežģītie vēsturiskie apstākļi noteica to, ka tēlotājmākslā viena no galvenajām atveidotajām tēmām kļuva nacionālā atbrīvošanas kustība. Vēsturiskā glezniecība ieguva arvien lielāku nozīmi, tās uzplaukums pēc 1863.~gada sacelšanās sakāves saistāms ar jau minētā A.~Grotgera un J.~Matejko vārdu.

Kaut trīs Polijas daļas attīstījās dažādi, notikumi, īpaši sacelšanās, vienā no tām, ietekmēja norisēs mākslā arī pārējās teritorijās. Spilgts piemērs tam bija \strong{Artura Grotgera} daiļrade. Viņš dzīvoja Galīcijā, nepiedalījās 1863.~gada sacelšanās, bet bez citiem darbiem (portretiem, žanra ainām) radīja arī veselus ciklus uz kartona (ogle, zīmulis, krītiņi), veltītus tās notikumiem. Šie viņa darbi kļuva par vienu no XIX gadsimta poļu glezniecības virsotnēm. Reālisms un patriotisks heroisms viņa darbos sadzīvoja ar akadēmisma un mistiskas simbolikas elementiem. Atbrīvošanās cīņas tēma, kā tā tēlota A.~Grotgera darbos, kļuva par visu to mākslinieku orientieri, kuri vēlāk pievērsās šai tematikai.

\strong{Jans Matejko} par savu aicinājumu uzskatīja reliģisko jaunradi, bet 1863.~gada sacelšanā neveiksme pamudināja viņu pievērsties vēsturiskajai glezniecībai un turpmāk viņš daudz gleznoja kaujas un citus vēsturiskus notikumus. 60.--80.~gados viņa darbos vēsturiskā glezniecība sasniedza virsotni. Nacionāli-pilsoniskais patoss, dramatiskā izteiksmība, reālistiski pārliecinošie tēli J.~Matejko daudzfigurālajās kompozīcijās apvienoti ar spilgtu krāsu gammu. Gleznotājs kļuva par gleznu autoru, kuras atainoja svarīgākos notikumus Polijas vēsturē. Ne velti viņu kādreiz sauc par Polijas vēstures ilustratoru. Mūsdienās viņa gleznas glabājas Polijas nacionālās mākslas dārgumu krātuvēs. Ja 1857.~gadā J.~Klačko noraidīja poļu mākslas tiesības uz pastāvēšanu, tad 1878.~gadā Krakovas sabiedrība pasludināja J.~Matejko par poļu gleznotāju ``karali'', nostādot to vienā rindā ar vārda māksliniekiem Ā.~Mickēciču, J.~Slovacki un Z.~Krasiņski.

Tiesa, J.~Mateiko darbos sastopamas arī virkne vēsturisku neprecizitāšu. Piemēram, gleznā ``Stefans Batorijs pie Pleskavas'' tiek rādīta pilsētas padošanās, kad patiesībā šis Polijas karalis tā arī nespēja to ieņemt.

Demokrātiskā reālisma tendences glezniecībā savu augstāko pakāpi sasniedza 70.--80.~gados gleznotāju brāļu Maksimiliāna un \strong{Aleksandra Gerimsku un J.~Heļmonska} daiļradē.

\strong{Maksimiliāna Gerimska} jaunrades galvenā tēma bija ainava. Tās atveide balstījās uz dabas novērojumiem dažādos gada un diennakts laikos. Ir apbrīnojama viņa dabas un krāsu toņu izjūta. Bez tam M.~Gerimska otai pieder arī gleznas no pilsētas nabadzīgāko slāņu dzīves, 1863.--1864.~gada sacelšanās ainas. Viņš galvenokārt dzīvoja un strādāja Minhenē, ap viņu grupējās citi antiakadēmiskā virziena gleznotāji. Par medību skatu attēlojumu viņš tika ievēlēts par Vācijas Mākslas akadēmijas goda locekli. Viņa darbi tika izstādīti Londonā un Parīzē, mākslinieks saņēma daudz pasūtījumu, bet dzimtenē ilgi nevēlējās viņu atzīt par nopietnu gleznotāju. M.~Gerimskis mira jauns un atzinību Polijā viņa darbi guva tikai pēc viņa nāves.

\strong{Aleksandrs Gerimskis} daudz ceļoja, bieži dzīvoja trūkumā, tāpat kā brālis dzīves laikā Polijā bija maz pazīstams. Arī viņš gleznoja ainavas, pilsētas dzīves skatus, bija lielisks gaismas attēlošanas meistars.

\strong{Juzefs Heļmonskis} sākotnēji dzīvoja un strādāja Minhenē un Parīzē. Pēdējā darbojās kā avīzes ``\frtxti{Le Monde}'' (``Pasaule'') ilustrators. 1887.~gadā atgriezās Polijā, gleznoja galvenokārt Polijas un Ukrainas ainavas, arī žanra un medību skatus.

Batālajā glezniecībā daudz darbu radījis \strong{Julians Kossaks}. Viņš ilgstoši dzīvoja Parīzē, Minhenē, studēja izcilāko ārzemju meistaru darbus, tāpēc daži vācu autori rakstījuši, ka viņš bijis ``pasaules pilsonis'' ne tikai dzīvē, bet arī savā mākslā. Taču vēlāk J.~Kossaks vadīja Varšavas laikraksta ``\pltxti{Tygodnik Ilustrowany}'' (``Ilustrētais Iknedēļas Izdevums'') māksliniecisko izdošanu, mūža nogali pavadīja Krakovā. Viņa māja (``\pltxti{Kossakówka}'') kļuva par Krakovas radošās inteliģences~--- jau minēto Ā.~Asnika, H.~Senkeviča, J.~Heļmonska u.c. satikšanās vietu. Viņa gleznu galvenās tēmas bija kauju skati no Polijas agrīnās vēstures, viņš gleznoja arī romantiskus žanra skatus~--- ūsainus panus kuntušos (\pltxti{Kontusz}~--- garš mētelis ar izgrieztām atverēm rokām) un meičas lauku puķu vainagos uz ļoti ``poliskas'' dabas fona. Patika viņam gleznot arī medību ainas un dažādas krāsas zirgus.

Ļoti populārs bija akadēmiskā virziena gleznotājs V.~Čahorskis. Viņa darbi bija tik pieprasīti, ka pasūtītājiem nācās gaidīt savu kārtu līdz diviem gadiem.

XIX~gadsimta vidū ievērojamu uzplaukumu piedzīvoja poļu \strong{grafika}. Virkne talantīgu meistaru, tādu kā jau minētie A.~Grotgers un J.~Kossaks, arī \strong{F.~Kostševskis} daudz strādāja grafikas jomā, radot augsti mākslinieciskus darbus. F.~Kostševskis gleznoja galvenokārt žanra ainas un ainavas. Viņa ilustrācijas un karikatūras publicēja Varšavas periodika. Bez tam plaši attīstījās ilustrāciju grafika, grafikas tehnikā (galvenokārt attīstījās ksilogrāfija un litogrāfija) tika atveidoti pazīstamu gleznotāju darbi.

XIX un XX~gadsimtu mijā, kad literatūrā valdīja ``Jaunās Polijas'' virziens, poļu glezniecībā izplatījās \strong{impresionisms}.

Reālistiskas un reizē impresionistiskas ainavas gleznoja \strong{Julians Falats}. Studējis Krakovas un Minhenes Mākslas akadēmijās, viņš daudz arī ceļoja pa Eiropu un Āziju un savus iespaidus realizēja romantisma gara pilnos darbos. Medībās kņaza A.~Radzivilla pilī viņš iepazinās ar Prūsijas princi Vilhelmu (nākamo Vilhelmu II). Pēdējam kļūstot par Vācijas ķeizaru, J.~Falats kļuva par galma gleznotāju Berlīnē. Viņa otai piederēja daudz vēsturisku gleznu, dzīvnieku atveidojumu un ziemas ainavu. 1895.~gadā viņš atgriezās Krakovā, kļūstot par mākslas profesoru, reformēja Krakovas mākslas skolu, ieviešot jaunas mācību metodes, balstītas uz impresionistisko pieredzi.

XIX gadsimta demokrātisko mākslu ar jaunā laikmeta mākslinieciskajiem principiem veiksmīgi savienoja \strong{Leons Vičulkovskis}. Sākotnēji viņš gleznoja reālistiskā manierē, taču XIX gadsimta beigās pēc brauciena uz Parīzi viņa daiļradē auga impresionisma ietekme un mākslinieks radīja ainavas un žanra gleznas ar dažādiem gaismas efektiem. Strādāja viņš dažādās tehnikās: viņu saistīja eļļas glezniecības mākslinieciskās iespējas, akvareļu artistiskums, pasteļkrāsu smalkums, tušas zīmējumu efekti, litogrāfiju krāsu saskaņas sarežģītība. Savas dzīves laikā viņš radīja lielu skaitu portretu, daudzos no kuriem iemūžināti poļu kultūras pārstāvji. XIX un XX~gadsimtu mijā L.~Vičulkovskis radīja arī daudz grafikas darbu, veltītu Polijas vecajām pilsētām~--- Varšavai, Krakovai, Ļubļinai un to arhitektūras pieminekļiem. Viņa litogrāfijām ar mīkstu gaismas un ēnu gradāciju bija liela loma poļu grafikas attīstībā jau XX~gadsimtā, kad grafika piedzīvoja jaunu uzplaukumu.

Ievērojams notikums tēlotājmākslas attīstībā bija Raclavicas panorāmas (\pltxti{Panorama Racławicka}) izveide, kura attēloja kauju pie Raclavicas (\pltxti{bitwa pod Racławicami}) 1794.~gada 4.~aprīlī starp poļu sacelšanās dalībniekiem T.~Kostjuško vadībā un ģenerāļa A.~Tormasova komandēto Krievijas armiju, kad poļi guva vienu no nedaudzajām savām uzvarām. Panorāma tika radīta 1893.--1894.~gadā uz kaujas simtgadi pēc Ļvovas pilsētas, kura tad atradās Austroungārijas sastāvā, pasūtījuma. Tās galvenie autori bija gleznotāji J.~Stuka un V.~Kossaks. Panorāmu atklāja 1894.~gada 5.~jūnijā Ļvovā notiekošās vispārējās poļu izstādes (\pltxti{Powszechna Wystawa Krajowa}) gaitā speciāli šim nolūkam uzbūvētā ēkā.

1944.~gadā glezna daļēji tika bojāta bombardēšanas rezultātā Ļvovas atbrīvošanas no nacistiem laikā. 1946.~gadā PSRS to nodeva Polijai un to pārveda uz Vroclavu, kur uzglabāja muzejā saritinātā veidā. Politiski cēloņi (Polijas varas iestāžu nevēlēšanās riskēt ar padomju vadības neapmierinātību par uzvaras pār Krievijas armiju slavināšanu) bremzēja lēmuma pieņemšanu par jaunas ēkas celtniecību panorāmas vajadzībām, tā sākās tikai 1980.~gadā. Tika restaurēta arī pati glezna. Panorāma tika atklāta 1985.~gada 14.~jūnijā Vroclavā. Kopš tā laiku to ir apmeklējuši daudzi miljoni skatītāju.

Izcils oforta, litogrāfijas meistars bija viens no ``Jaunās Polijas'' izcilākajiem darbiniekiem \strong{Juzefs Mehofers}, kurš radīja ļoti daudz t.s. ``sadzīves grafikas'' darbu: grāmatu vākus un ilustrācijas, vinjetes, grāmatzīmes, plakātus. Viņš veidoja arī daudz, tai skaitā radošo personību, portretu. Mākslinieks ir pazīstams arī kā logu vitrāžu autors Krakovas dievnamos.

Poļu \strong{tēlniecība} XIX~gadsimta pirmajos gadu desmitos bija attīstīta vājāk nekā glezniecība, sasniegumi mazāk redzami. Lielā mēra vainojama bija Polijas sadalīšana, situācija nozīmīgām pilsoniskām un patriotiskām tēmām veltītas monumentālās skulptūras attīstībai bija nelabvēlīga. Trūka arī bagātu aizbildņu ziedojumu lielus izdevumus prasošajiem mākslinieku darbiem. Tomēr arī šai mākslas jomā tika radīts ne mazums interesantu un vērtīgu darbu.

Arī skulptūrā nostiprinājās nacionāli-patriotiskā virzība, kas izpaudās galvenokārt tēlnieku daiļrades sižetiski-tematiskajā ievirzē. Viņu meklēja iedvesmas avotus vēsturē un literatūrā, izmantoja alegorijas. Galvenokārt tēlniecība attīstījās Polijas karalistē, īpaši Varšavā. Tās bija memoriālās skulptūras (nereti~--- kapu pieminekļi) un plastika, saistīta ar jaunu arhitektūras kompleksu noformēšanu: reljefi, skulpturāli izrotājumi. Šajos reljefos valdīja klasicisma formas. Par modes lietu kļuva arī izcilu poļu darbinieku portretu un krūšutēlu darināšana.

Diskusija par to, kādai jābūt nacionālajai tēlniecībai, Polijā izraisījās pirms diskusijas par glezniecību. Lieta tā, ka 1817.~gadā dāņu skulptoram B.~Torvaldsenam tika pasūtīts piemineklis J.~Poņatovskim. 1829.~gadā Varšavā izgatavoja tā modeli. Piemineklis bija traktēts romiešu heroikas stilā. B.~Torvaldsens poļu kņazu bija attēlojis līdzīgi Romas imperatora Marka Aurēlija piemineklim Kapitolijā, ietērpjot to antīkā tērpā. Tas izsauca protestu poļu sabiedrības dažādos slāņos. Atzīstot autora talantu un tā darinājuma augstās mākslinieciskās vērtības, J.~Poņatovska cienītāji sašuta: ``Kāda mums daļa par šo plikni krekliņā''; saskatīja puskailajā antīkajā varonī ``sastingumu, dzīvības un izteiksmes trūkumu, sirds un atmiņu trūkumu''. B.~Torvaldsenam nācās izdarīt pieprasītos labojumus. Reljefus uz J.~Poņatovska pieminekļa pjedestāla gatavoja poļu skulptors J.~Tatarkevičs. Tajos jau parādījās reālistiskas konkrētības tendence notikumu atveidē.

Bronzā pieminekli izlēja 1832.~gadā, taču sakarā ar represijām pēc Novembra sacelšanās to nevarēja uzstādīt tam paredzētajā vietā. Monuments tika konfiscēts un aizvests uz Krieviju. 1840-jos gados Nikolajs~I to uzdāvināja sacelšanās apspiedējam I.~Paskevičam, kurš to novietoja pie savas pils Gomeļā. Piemineklis tika atgriezts Polijai tikai 1922.~gadā sakarā ar Rīgas miera noslēgšanu (1921). 1923.~gadā monumentu uzstādīja Sakšu pils (\pltxti{Pałac Saski}) un Nezināmā kareivja kapavietas (\pltxti{Grób Nieznanego Żołnierza w Warszawie}) priekšā. Pēc Varšavas sacelšanās (1944) apspiešanas vācu okupanti līdz ar Sakšu pili to uzspridzināja. 1948.--1951.~gadā Dānijas valdība pēc modeļa, kas bija saglabājies D.~Torvaldsena muzejā Kopenhāgenā, lika izgatavot jaunu pieminekli un uzdāvināja to Polijai. Pieminekli uzstādīja Lazenku parkā, bet kopš 1965.~gada tas atrodas bijušās Vietvalža pils (\pltxti{Pałac Namiestnikowski}), bet tagadējās Polijas prezidenta pils (\pltxti{Pałac Prezydencki w Warszawie}) priekšā.

J.~Tatarkevičs veidoja (1834--1836) S.~Potocka kapa pieminekli. S.~Potocka lielie nopelni poļu kultūras labā stādīja viņu ``dižo poļu'' rindā, kuru kults un arī pieminekļi kalpoja nacionālās pašapziņas celšanai. Šai pašā tradīcijā J.~Tatarkevičs radīja arī citu poļu darbinieku pieminekļus un krūšutēlus. Pieminekļus un krūšutēlus veidoja arī citi, mazāk pazīstami XIX gadsimta pirmās puses poļu tēlnieki.

70--90.~gadu skulptūrās stiprāk kā iepriekš izpaudās reālisms. Tā, gan akadēmisma, gan reālisma iezīmes var atrast K.~Godebska darbos. Viņš darinājis poļu zinātnes un kultūras darbinieku: astronoma N.~Kopernika, gleznotāja J.~Matejko, dramaturga A.~Fredro, dzejnieka Ā.~Mickēciča, muzikanta S.~Moņuško pieminekļus.

XIX~gadsimta beigās un XX~gadsimta sākumā ``Jaunās Polijas'' laikmets veicināja tēlniecības nozīmības pieaugumu. Reālisms sadzīvoja ar modernismu un impresionismu. V.~Šimanovskis pārstāvēja neoromantisma un simbolisma virzienus. Viņš ir izcilā poļu mūziķa F.~Šopēna pieminekļa autors Varšavā Lazenku rajonā, kurš raksturīgs ar emocionālu pacēlumu un kompozīcijas dinamismu, stilizētu formu elastību.

Virkni monumentālu skulptūru radīja K.~Duņikovskis. Akadēmiskajā skulptūrā izcēlās P.~Velonskis, kurš bija mācījies Pēterburgas Mākslas akadēmijā, strādājis Itālijā, Francijā un Vācijā. Viņa darbi bronzā raksturojās ar rūpīgu materiāla izstrādi, nobeigtību.

Kaut tēlotājmākslā, tāpat kā literatūrā, radošie gari guva iedvesmu konkrētajā sociālajā pieredzē, tomēr arī viņu darbos izpaudās nacionālās jūtas, lai arī darbojās viņi dažādos centros kā sadalītajā Polijā, tā arī ārzemēs.

\strong{Arhitektūrā} XIX~gs. pirmajos gadu desmitos vēl turpināja attīstību no XVIII~gadsimta beigām mantotās klasicisma formas, kuras ieguva izteiktāku monumentalitāti un stingru vienkāršību. 20.~gados sākās jauns celtniecības vilnis. Turpinājās piļu celtniecība. 20.--30.~gados noformējās Varšavas centrs. Arhitekta A.~Koracci vadībā notika Banku laukuma izbūve ar Valsts ienākumu un kases komisijas (\pltxti{Pałac Komisji Rządowej Przychodów i Skarbu}, 1825) vareno ēku, kura veido vienotu kompleksu, taču tajā nav ne miņas no monotonuma. Mierīgs kolonu ritms atvieglo fasādi, iekļauj gaismēnas efektus ēkas veidolā. A.~Koracci būvēja arī Varšavas Lielo teātri (\pltxti{Teatr Wielki}, 1825--1833) un daudzas citas ievērojamas celtnes, kā Stašica pili (\pltxti{Pałac Staszica}, 1820--1823), pieminekli lojālajiem virsniekiem, nogalinātiem Novembra sacelšanās laikā (\pltxti{Pomnik oficerów-lojalistów poległych w Noc Listopadową w Warszawie}, 1841) u.c. A.~Koracci darbi bija spilgtākā poļu klasicisma izpausme.

Neoklasicisma stilā Pēterburgā daudzas ēkas uzbūvēja poļu arhitekts M.~Lalevicš. Viņš ieguva arhitektūras akadēmiķa grādu, bija Pēterburgas arhitektu biedrības priekšsēdētājs. Pēc 1918.~gada dzīvoja Polijā, strādāja Varšavas Politehniskajā institūtā (\pltxti{Politechnika Warszawska}), vadīja daudzu ēku celtniecību.

Kopš 1830.~gada attīstījās eklektiskais historisma stils, kurš vienlaikus pavēra plašas iespējas jaunradei, veicinot ģeogrāfiski un etniski visdažādāko kultūras elementu sakausēšanai, nodzēšot robežas starp pagātni, tagadni un nākotni. Jau XIX~gadsimta 40.--60.~gados arhitektūra piedzīvoja pārmaiņas. Norisa iepriekšējām desmitgadēm raksturīgās ar lielu ansambļu celtniecību saistītās stila vienotības noārdīšana. Arvien plašāk parādījās eklektisms, raksturīgs industriālā laikmeta attīstības sākumposmam. XIX~gadsimta 50.~gados populārākā bija t.s. ``neorenesanse''. Daudz tai raksturīgu ēku radīja poļu arhitekts H.~Markoni.

Rūpniecībai, tirdzniecībai nepieciešamo ēku celtniecība veda pie jaunu konstruktīvo risinājumu pielietošanas, jaunu ēku tipu izstrādes, tādu materiālu kā čuguns un dzelzs pielietošanas celtniecībā. Taču pilnībā šīs iezīmes izpaudās nākamajos gadu desmitos. XIX~gs. nogalē populārs kļuva modernisma stils, ko Polijā parasti sauca par ``secesiju''. Tam raksturīga bija plaša betona izmantošana celtniecībā. Sāka veidoties jaunas pieejas pilsētu un atsevišķu dzīvojamo rajonu plānošanai.

XIX gadsimta pirmajās desmitgadēs \strong{mūzikā} darbojās vācu izcelsmes poļu komponists J.~Elsners, J.~Kozlovskis, M.~Oginskis u.c. Šajā periodā pieauga interese par poļu mūzikas tautiskajām tradīcijām. Tapa daudz polonēžu, mazurku, krakovjaku un citu instrumentālu skaņdarbu, kā arī dziesmas un romances ar poļu tekstiem.

J.~Elsners bija dzimis vācu katoļu ģimenē, 1799.~gadā pārcēlās uz Varšavu, kur līdz 1824.~gadam bija Nacionālā teātra direktors. Viņš sacerējis daudz operu, tai skaitā ``\pltxti{Leszek Biały}'', (Baltais Ļešeks) un ``\pltxti{Król Łokietek}'' (Karalis Loketeks), astoņas simfonijas, polonēzes, maršus, kamermūziku.

Pēterburgā karjeru veidoja J.~Kozlovskis, kurš 1786.~gadā iestājās Krievijas armijā, bet jau 1791.~gadā guva ievērību Pēterburgā kā komponists, 90.~gadu vidū sacerēja romanses ``\rutxti{Российские песни}'' (Krievijas dziesmas). Populāras bija viņa polonēzes. Jau iepriekš ir pieminēta viena no M.~Oginska polonēzēm. Sākotnēji polonēze bija poļu tautas deja, bet pārvērtās par mūziku, ko izpildīja koncertos.

M.~Oginskis aktīvi nodarbojās ar politisku darbību, līdz trešajai Žečpospolitas sadalīšanai bija tās sūtnis Holandē, Anglijā, kā virsnieks cīnījās par to. Pēc neatkarības zaudēšanas emigrēja, Francijā veda sarunas ar Direktoriju par Polijas valstiskuma atjaunošanu, tad 1810.~gadā kļuva par Krievijas imperatora Aleksandra I uzticības personu, sagatavoja Lietuvas lielkņazistes atjaunošanas plānu, kurš tika noraidīts. Tai pat laikā M.~Oginskis nodarbojās ar komponēšanu, tai skaitā uzrakstīja ap 20 dažādu polonēžu.

J.~Elsnera skolniece bija virtuozā piāniste un komponiste M.~Šimanovska. 1823.--1827.~gados viņa koncertēja visā Eiropā. Iespējams, bija pirmā piāniste, kura savus sacerējumu izpildīja bez notīm.

Visas šīs muzikālās aktivitātes sagatavoja augsni \strong{Frederika Šopēna} daiļradei, kura piesaista visas pasaules mūzikas speciālistu un vienkārši cienītāju uzmanību. Polietes un franču emigranta, kurš piedalījās T.~Kostjuško vadītajā sacelšanās pret Krieviju, dēls, viņš jau bērnībā demonstrēja muzikāla brīnumbērna dotības. Viņa skolotāji bija čehu vijolnieks V.~Živnijs un J.~Elsners. Tāpat kā literātiem Ā.~Mickevičam un J.~Slovackim arī F.~Šopēnam~--- visu laiku ievērojamākajam poļu pianistam un komponistam, nācās dzīvot un nodarboties ar jaunradi emigrācijas apstākļos. 1830.~gada sacelšanās sakāve liedza viņam atgriešanos dzimtenē. Kad L.~v.~Bēthovena klasicismu kā galveno virzienu Eiropas mūzikā aizstāja romantisms, par vienu no tā galvenajiem pārstāvjiem kļuva F.~Šopēns. Lielākā viņa kompozīciju daļa ir solo darbi klavierēm. Visiecienītākās F.~Šopēna formas bija etīdes, mazurkas, noktirnes un valši. No lielajām formām jāatzīmē divi klavierkoncerti. Daudzi tā laika mākslinieki un literāti apjūsmoja F.~Šopēnu. Mūsdienās viņa darbi atrodami daudzu piānistu repertuārā. Kopš 1927.~gada Varšavā notiek Starptautiskais Šopēna vārdā nosauktais pianistu konkurss. Turpat 2010.~gadā sakarā ar komponista 200 dzimšanas dienu pēc rekonstrukcijas un modernizācijas tika atklāts viņa muzejs.

No 1810 līdz 1842.~gadam Varšavas operas teātra diriģents bija K.~Karpinskis.Viņš radīja virkni operu ar poļu vēstures sižetiem, aizstāvēja nacionālo kolorītu mūzikā, paziņojot: ``Dziedāsim skaisti, taču visai poļu tautai''.

Ievērojamāko vietu poļu muzikālās kultūras attīstībā ieņēma \strong{Staņislavs Moņuško}. Dzimis Baltkrievijā, savu radošo darbību viņš saistīja ar poļu kultūru. No 1843.~gada sacerēja dziesmas (kopā ap 400), kuras lika poļu vokālās mūzikas pamatus. ``Es neradu neko jaunu! Klaiņojot pa poļu zemēm, es piepildos ar tautas dziesmu garu. No tām, ārpus manas gribas, iedvesma ieplūst visos manos sacerējumos'',~--- tā viņš raksturoja savu daiļradi mūža beigās. Vokālā mūzika kļuva par komponista radošās darbības virsotni, tajā parādījās viņa melodiskā apdāvinātība, novatorisms. Taču viņu vilināja poļu nacionālas operas radīšanas ideja. Pirmā no tām ``\pltxti{Halka}'', koncertizpildījumā izskanēja 1848.~gadā, bet uzvesta tika 1854.~gadā Viļņā. Otrā redakcijā izskanēja 1858.~gadā Varšavā. S.~Maņuško radīja vairāk nekā 15 operu, arī operetes, baletus, mūziku dramatiskajām izrādēm.

Par S.~Maņuško darba turpinātājiem kļuva vairāki ievērojami muzikanti.

V.~Želenska operas ``\pltxti{Konrad Wallenrod}'', ``\pltxti{Goplana}'' ``\pltxti{Janek}'', ``\pltxti{Stara Baśń}'' (Veca pasaka) izceļas ar nacionālo skanējumu, kaut tajās jūtami arī Rietumeiropas mūzikas aizguvumi. Izcils diriģents un muzikālais pedagogs bija Z.~Noskovskis.

Organisks poļu tradīciju savienojums ar rietumu neoromantismu izpaudās L.~Ružicka darbos. 1909.--1912.~gadā viņš strādāja kā diriģents un konservatorijas profesors Lembergā, pēc tam devās turnejā pa Eiropu. 1912.--1919.~gadā pārsvarā dzīvoja Berlīnē, kur viņu pārsteidza karš. 1919.~gadā viņš atgriezās dzimtenē jau kā komponists ar pasaules vārdu.

Ar lielu emocionalitāti, lirismu bija pārpildīta M.~Karloviča daiļrade, īpaši romances. Viņš izmantoja ne tikai poļu, bet arī baltkrievu un lietuviešu tautas dziesmu motīvus.

Komponistu daiļrades attīstība norisa reizē ar visas muzikālās dzīves uzplaukumu. Poļu rietumu zemēs tika dibinātas koru biedrības, darbojās konservatorijas un mūzikas skolas, no 1858.~gadā Muzikālais institūts Varšavā (\pltxti{Warszawski Instytut Muzyczny}).

Izcilu vietu poļu muzikālās kultūras attīstībā ieņēma \strong{brāļi} Henrihs un Juzefs \strong{Veņavski}. Pirmais bija vijolnieks, XIX~gadsimta 50.~gados kā virtuozs kļuva slavens Rietumeiropā un Krievijā. No Pēterburgas konservatorijas dibināšanas 1862.~gadā līdz 1865.~gadam bija tās profesors, 1866.--1869.~gadā bija Maskavas, bet no 1878.~gada~--- līdz mūža beigām~--- Briseles konservatorijas profesors. Kā vijolnieks triumfu piedzīvoja Vispasaules izstādē Parīzē 1878.~gadā. Kā komponists H.~Veņavskis radījis daudzus skaņdarbus vijolei, kuri raksturojas ar lielu gaumi un domas oriģinalitāti. Kopš 1935.~gada notiek starptautiskais Veņavska vārdā nosauktais vijolnieku konkurss. Juzefs Veņavskis bija pianists, daudz koncertēja, kā komponists popularizēja poļu muzikālo kultūru aiz robežas, apvienojot sekošanu F.~Šopēna iedibinātajām tradīcijām ar sava laika harmoniskām novitātēm.

Lielu popularitāti ieguva muzikantu Kontsku (\pltxti{Kątski}) ģimene, īpaši pianists, komponists, pedagogs Antons Konskis kā arī vijolnieks, komponists un pedagogs Apolinārijs Konskis. Otrais 1861.~gadā Varšavā nodibināja Muzikālo institūtu (\pltxti{Instytut Muzyczny}), kurā arī strādāja.

XIX~gadsimta beigās pasaules slavu iekaroja pianists \strong{Jans Paderevskis}. Īpaši atzītas bija viņa F.~Šopēna darbu interpretācijas. J.~Paderevskis kļuva pazīstams arī kā komponists. Viņa opera ``\pltxti{Manru}'' (1893--1901) tika ar panākumiem uzvesta gan Eiropā, gan ASV. Tā ir vienīgā opera poļu valodā, kura uzvesta Metropoliten (\pltxti{Metropolitan Opera}) teātrī Ņujorkā.

Pasaules slavu ieguva poļu muzikālā teātra aktieri~--- operdziedātājas M.~Zemriha-Kohaņska (koloratūrsoprāns), J.~Koroļevič-Vaidova (soprāns), operdziedātājs V.~Merčinskis (tenors). Viņi dziedāja arī S.~Maņuško operu uzvedumos, kuri faktiski radīja muzikālo teātri Polijā. Maskavā, Pēterburgā, Rietumeiropas centros uzstājās A.~Boļska (liriskais soprāns). Krievu baleta vēsturē ievērojama vieta pieder poļu izcelsmes māksliniekiem M.~Kšešinskai un V.~Nežinskim. Pirmā kļuva pazīstama ne tikai kā Pēterburgas Marijas teātra soliste, vēlāk baleta pedagoģe, bet arī vairāku Romanovu dinastijas pārstāvju mīļākā. V.~Nežinskis spoži sāka savu dejotāja un horeogrāfa karjeru, taču tā nebija ilga, jo viņš ilgstoši slimoja.

1901.~gadā tika nodibināta ``Varšavas filharmonija''~--- akciju sabiedrība, kura par mērķi stādīja simfoniskās un kamermūzikas koncertu organizēšanu. Šim nolūkam tika radīts liels simfoniskais orķestris (iepriekš Varšavā simfoniskos koncertus sniedza operas orķestris) un uzcelta koncertzāle ar vietām 2~tūkstošiem klausītāju. Diemžēl filharmonijai pelņas nolūkos nācās orientēties uz ārzemju, nevis vietējiem izpildītājiem.

Poļu mūzikas vēstures attīstību veicināja mūzikas zinātņu katedru izveidošana Krakovas (1911) un Lembergas (1912) universitātēs.

Muzikā šīs mākslas jomas specifikas rezultātā visspilgtāk izpaudās nacionālās mākslas nedalāmības princips, tās nepakļaušanās valstu robežām.

Poļu \strong{teātra mākslas} attīstība bija cieši saistīta ar mūziku. Lugas, kuras uz Varšavas un citu pilsētu skatuvēm uzveda kā profesionālu aktieru, tā amatieru trupas, bieži iekļāva dziesmas, dejas, atsevišķas ainas pavadīja instrumentāla mūzika.

Jau XVIII~gadsimtā, Apgaismības laikmetā bija dibināts Polijas `` Nacionālais teātris''. Tad pat radās arī pirmie poļu tautības dramaturgi. Var atzīmēt, ka pirmā sociālpolitiskā komēdija ``\pltxti{Powrót posła}'' (``Sūtņa atgriešanas'', 1791), kuras autors bija J.~Ņemcovičs, izsauca Seimā sašutumu, jo tā atmaskoja lielas šļahtas daļas garīgo atpalicību.

Pakāpeniski poļu teātru repertuārā līdz ar V.~Šekspīra, Ž.-B.~Moljēra, F.~Šillera, G.~Kalderona un F.~Lope de Vegas lugām nostiprinājās arī vietējo dramaturgu darbi.

Romantisma laikmetā poļu teātrī pastāvēja divi pretēji virzieni. Vienu pārstāvēja Aleksandrs Fredro~--- dižciltīgs muižnieks (grāfs), kurš bija karojis Napoleona karaspēkā, apbalvots ar poļu ordeņa ``\latxti{Virtuti Militari}'' Zelta krustu un franču ``Goda leģiona'' ordeni (\frtxti{Ordre national de la Légion d'honneur}), un atgriezies Polijā. 1818.~gadā viņš uzrakstīja savu pirmo lugu ``\pltxti{Pan Geldhab}''. Uzveda to 1821.~gadā. Pēc tam viņš rakstīja komēdijas par šļahtas sadzīvi, dzeju, poēmas, aforismus. Kaut A.~Fredro dzīvoja romantisma valdīšanas laikmetā, viņš balstījās uz Apgaismības tradīcijām un dzīvo teātra praksi, ienesa komēdijas žanrā jaunu sociāli-sadzīvisku saturu, tuvināja skatuves valodu sarunu valodai. Viņa komēdijām bija liela loma poļu reālistiskā teātra izveidē. Viņa paraugi bija franču aktieris, režisors un dramaturgs Ž.~B.~Moljērs un venēciešu dramaturgs K.~Goldoni, bet situācijas un personāži viņa lugās bija poliski. A.~Fredro komēdijām bija ievērojama loma poļu reālistiskā teātra tapšanā, tās pieskaitāmas poļu literatūras zelta fondam un arī mūsdienās tiek uzvestas uz daudzu pasaules teātru skatuvēm.

Pretēji A.~Fredro \strong{Juzefs Korženovskis}, kurš strādāja kā pedagogs, mēģināja pārvarēt franču teātra ietekmi, viņa drāmās atspoguļojās personīgie sociālo apstākļu vērojumi un to izmaiņas uz saimniecisko pārmaiņu pamata. Pirmajos literārā darba gados (1825--1846) J.~Korženovskis rakstīja drāmas, vēlāk arvien tālāk atejot no romantisma pie reālisma, pievērsās prozai.

XIX~gs. otrajā pusē poļu teātris pārdzīvoja t.s. ``zvaigžņu laikmetu'', kas saistās ar izcilu aktrišu kā \strong{Helēna Modzejevska} un aktieru kā \strong{Aloizs Žolkovskis} uzstāšanos. Ne tikai Varšavā, bet arī Lembergā, Viļņā, Pozenē nostiprinājās profesionālais teātris. Īpaši uzplauka teātris Krakovā režisora T.~Pavļikovska vadībā. Šajā laikā poļu teātri arī sāka uzvest lielo romantiķu Ā.~Mickēviča, J.~Slovacka darbus.

Ievērojamu ieguldījumu poļu un visas pasaules kultūrā ienesa XIX~gadsimta beigu--XX~gadsimta sākuma poļu dramaturģija. Izcilākā figūra tajā bija \strong{Staņislavs Vispjaņskis}~--- rakstnieks, gleznotājs, teātra darbinieks, mūsdienu modernā poļu teātra iedibinātājs, pazīstams arī ar savām dekorācijām. Būdams dzejnieks--dramaturgs, laikabiedru atzīts par dzejnieku~--- vadoni, dzejnieku~--- pravieti, kādi savā laikā bija Ā.~Mickevičs un J.~Slovackis, viņš nepieturējās kādai noteiktai politiskai orientācijai un viņam nebija savas kaut cik izstrādātas sociālās programmas. Turpinot savu priekšgājēju ``pravietiskās'' tradīcijas, S.~Vispjaņskis izvērsa arī polemiku pret vēstures romantisku izpratni, pret politisko romantismu. Viņa drāmas un traģēdijas, kurās mākslinieks radīja metaforiskus, monumentālus vēstures un tagadnes vispārinājumus, raksturojas ar sarežģītu simbolisku, fantastisku un reālistisku līniju savijumu. Tajās parādīta Polijas varonīgā pagātne~--- ``\pltxti{Bolesław Śmiały}'' (Boļeslavs Drosmīgais, 1903), 1830.~gada sacelšanās~--- ``\pltxti{Warszawianka}'' (Varšavjanka, 1898); ``\pltxti{Noc Listopadowa}'' (Novembra nakts, 1904); sava laikmeta dzīves un poļu tautas atbrīvošanās cīņas problēmas~--- ``\pltxti{Wesele}'' (Kāzas, 1901; ``\pltxti{Akropolis}'' (Akropole, 1904); antīki sižeti~--- ``\pltxti{Meleager}'' (Meleagrs (grieķu mitoloģijas varonis), 1899; ``\pltxti{Powrót Odysa}'' (Odiseja atgriešanās, 1907) u.c. Par S.~Vispjaņska dramaturģijas šedevru tiek atzīta drāma-pamflets un vienlaikus drāma-pasaka ``\pltxti{Vesele}'' (Kāzas). Tajā izmantoti poļu tautas leļļu teātra Ziemassvētku gadatirgus uzveduma paņēmieni, kad uzstājas tradicionāli tipāži: kalps, šļahtičs, ksendzs, ebrejs. Kopā ar viņiem darbojas arī fantastiski tēli. Sižets vijas ap Krakovas šļahtiča-inteliģenta precībām ar zemnieku meiteni. Tā pamata bija konkrēts vēsturisks fakts. Kāzu ainas mijas ar epizodēm no Polijas pagātnes, vienotas ar pastāvošā stāvokļa un muižniecības pasivitātes kritiku. Dažādu sabiedrisko un ideoloģisko pozīciju, kuras pārstāv kāzu dalībnieki, sadursmē parādījās zemnieku un šļahtas pasaules nesamierināmība. Ironizējot par šķiru solidaritātes mītu, nosodot inteliģenci par nespēju ieņemt tautas garīgā vadoņa vietu, S.~Vispjaņskis neredzēja spēku, kurš varētu izraut poļu sabiedrību no ``guļas'', no nevarības un bezdarbības, pesimistiski vērtējot poļu atbrīvošanās kustības perspektīvas. S.~Vispjaņska daiļrade ir tik savdabīga, ka nenoveda pie kādas īpašas skolas vai virziena rašanās, kas turpinātu realizēt viņa radošos principus. Sava laikmeta dramaturģijā viņš bija spilgta, bet vientuļa personība. Mūsdienās Krakovā ēkā, kur kādreiz dzīvoja S.~Vispjaņskis, ierīkots viņa muzejs. Rakstnieka lugas arī mūsdienās pastāvīgi atrodas poļu teātru repertuārā.

Rezonansi Eiropā ieguva S.~Pšebiševska drāmas ``\pltxti{Dla szczęścia}'' (Laimes vārdā, 1900), ``\pltxti{Goście}'' (Viesi, 1901), ``\pltxti{Złote runo}'' (Zelta aunāda, 1901), ``\pltxti{Śnieg}'' (Sniegs, 1903) u.c., saistītas ar vācu, skandināvu un arī krievu modernismu. Tomēr S.~Pšebiševska mēģinājumi radīt naturālistiski-simbolisku dramaturģisku sintēzi bija visai mazauglīgi. Visās viņa drāmās variējas viens un tas pats konflikts starp fatālu erotisku tieksmi un pienākuma apziņu, kas noved pie traģiskām sekām.

XIX~gadsimta beigās un XX~gadsimta sākumā reālistiskās dramaturģijas virsotnes sasniedza \strong{Gabriela Zapoļska}. Viņas spalvai pieder gan romāni un stāsti, gan arī drāmas. Plaši pazīstama kļuva komēdija ``\pltxti{Moralność pani Dulskiej}'' (Panes Duļskas morāle, 1908), kurā autore parādījusi, ka pilsoniskajā sabiedrībā aiz ``augstiem'' morāles principiem slēpjas liekulība, skopums, gara nabadzība un iztapība bagātības priekšā. Arī drāmās ``\pltxti{Panna Maliczewska}'', (Pana Maličevska, 1912) u.c. parādās rakstnieces novērošanas spējas, lieliskā mietpilsoņu vides tikumu pārzināšana. Tas ļāvis viņai spīdoši atmaskot šo aprindu liekulīgo un melīgo morāli. G.~Zapoļskas dramatiskos darbus raksturo prasmīga skatuves darbībās attīstība, raits dialogs, lieliski atainojošs atsevišķu personāžu individuālās īpašības. Arī viņas lugas ir ieņēmušas stingru vietu poļu teātru repertuārā. ''Panes Duļskas morāle'' tika uzvesta vairākos padomju teātros, 1956.~gadā arī J.~Raiņa Latvijas PSR Valsts Akadēmiskajā Dailes teātrī.

1913.~gadā darbību uzsāka \strong{\pltxti{Teatr Polski w Warszawie}} (Poļu teātris Varšavā), pirmais teātris Polijā ar grozāmo skatuvi. Ar to saistās tādu ievērojamu režisoru kā \strong{Arnolds Sufmans}, \strong{Leons Šillers} darbība. L.~Šillera nopelns bija arī režisoru augstākās izglītības iedibināšana Polijā (1933). Daudzi viņa skolnieki pēc Otrā pasaules kara ieņēma vadošas pozīcijas teātra mākslā.

Jau 1894.~gadā K.~Prošinskis konstruēja aparātu, ar kuru uzņēma īsas dokumentālas \strong{kinolentas}. 1908.~gadā Varšavā tika uzņemta pirmā poļu mēmā komēdija ``\pltxti{Antoś po raz pierwszy w Warszawie}'' (Atons pirmo reizi Varšavā). Pirmās poļu filmas bija darinātas pēc franču filmu parauga. Priekšroka tika dota lielu literāru darbu inscenēšanai, lomās tika aicināti pazīstami teātru aktieri. 1911.~gadā režisors un producents A.~Hercs nodibināja pirmo poļu kinostudiju ``\pltxti{Sfinks}'' (Sfinksa).

Kinoteātri Polijas teritorijā parādījās ap 1899.~gadu, bet pirmās kinofirmas un studijas~--- pēc 1902.~gada. Kinomāksla sāka ieņemt arvien redzamāku vietu kultūras dzīvē. Līdz Pirmajam pasaules karam poļu zemēs darbojās ap 400~kinoteātru, kuros tika demonstrētas ne tikai ārzemju, bet no 1908.~gada arī pašražotās filmas.

Žečpospolitas valstiskās neatkarības zaudēšana 1795.~gadā veda pie apstākļu pasliktināšanās \strong{zinātniskajai darbībai}. Daļa poļu tautības zinātnieku tika vajāta par dalību nacionālajā kustībā, bija spiesta emigrēt. Polijas ekonomiskā atpalicība izpaudās arī dabas un tehnisko zinātņu materiālās bāzes vājumā, tāpēc daudzi poļu zinātnieki nevarēja atrast darbu augstākajās mācību iestādēs un dzīvoja un strādāja \strong{ārpus Polijas}, jo tur apstākļi zinātniskai darbībai bija labvēlīgāki. Aiz robežas strādāja matemātiķis un filozofs J.~Vronskis; ceļotājs un ģeologs P.~Stržeļeckis pētīja kalnus Austrālijā; tagadējās Baltkrievijas teritorijā dzimušais ģeologs un ģeogrāfs I.~Domeiko, kurš bija mācījies Viļņas universitātē, strādāja Čīlē, 16 gadu bija Čīles universitātes (\estxti{Universidad de Chile}) rektors.

Jāpiezīmē, ka zinātnisku karjeru, tāpat kā agrāk ierēdņa karjeru, varēja taisīt Pēterburgā un citos Krievijas centros, pat nepārejot pareizticībā, kas Polijas karalistē pēc 1863.~gada parasti nebija iespējams.

Krievijas universitātēs Pēterburgā, Maskavā, Kijevā un Odesā strādāja pazīstamie poļu zinātnieki: lingvists J.~Boduēns de Kurtenē, juristi A.~Ļedņickis un V.~Spasovičs, tiesībzinātnieks, sociologs, filozofs L.~Petražickis, antīkās literatūras un reliģijas speciālists T.~Zeļinskis, matemātiķis J.~Sohockis, ķīmiķis A.~Verigo u.c.

Neraugoties uz nacionālo apspiestību un ziņatnisko iestāžu trūkumu, Apgaismības laikmeta tradīciju ietekmē arī pašā Polijā zinātne turpināja attīstīties. Tā, Varšavā zinātniskā darbība norisa sabiedrības dibinātās organizācijās. 1875.~gadā pilsētā tika nodibināts rūpniecības un lauksaimniecības muzejs (\pltxti{Muzeum Przemysłu i Rolnictwa}), pie kura darbojās arī labaratorijas, tika lasītas lekcijas. 1881.~gadā par godu Varšavas Galvenās skolas rektoram J.~Mianovskim bijušie tās audzēkņi nodibināja palīdzības kasi (``\pltxti{Kasę pomocy dla osób pracujących na polu naukowem im. dr. med. Józefa Mianowskiego}''), kuras uzdevums bija ar stipendiju piešķiršanu, zinātnisku publikāciju un konferenču, zinātnieku tikšanos finansēšanu veicināt zinātnes attīstību Polijā. Līdz 1907.~gadam kase bija viens no galvenajiem zinātnes atbalsta centriem Polijā. 1907.~gadā Varšavā tika radīta \strong{Varšavas Zinātnes biedrība} (\pltxti{Towarzystwo Naukowe Warszawskie}), kura turpināja ``Zinātnes draugu biedrības'' darbību, organizēja virkni zinātniski pētniecisko institūtu. Polijā tika izdoti arī zinātniski žurnāli ``\pltxti{Prace matematyczno-fizyczne}'' (Darbi matemātikā un fizikā, no 1888), ``\pltxti{Nowosci matematycznе}'' (Matemātikas ziņas, no 1897) u.c.

Sevišķi ievērojami bija panākumi ķīmijas zinātnē, kuru uz priekšu virzīja akmeņogļu pārstrādes, pārtikas un tekstilrūpniecības vajadzības. Jānosauc vairāki izcili tā laika poļu zinātnieki.

Poļu \strong{ķīmiķi} K.~Olševskis un Z.~Vrubļevskis, kuri strādāja Krakovā, pirmoreiz zinātnes vēsturē 1883.~gadā ieguva šķidru skābekli, pēc tam šķidru slāpekli.

Armēņu izcelsmes poļu farmaceits un ķīmiķis-tehnologs J.~Lukaševičs 1852.~gadā radīja metodi, kā no jēlnaftas destilācijas ceļā iegūstama petroleja. Nākamajā~--- 1853.~gadā viņš izgudroja petrolejas lampu. Reizē J.~Lukaševičs bija poļu naftas ieguves rūpniecības radītājs, jau 1854.~gadā viņš veica pirmo naftas urbumu un 1856.~gadā nodibināja pirmo naftas pārstrādes rūpnīcu Polijā.

Pasaulslavens eksperts naftas atradņu jautājumos bija ģeologs R.~Zubers, kurš, meklējot naftu, strādāja daudzas valstīs.

\strong{Medicīnas} darbinieku vidū godpilna vieta pienāks T.~Halubiņskim, kura nopelni sevišķi lieli bija kurortoloģijā. Viņš lika pamatus kūrorta attīstībai poļu Karpatos~--- Zakopanē.

60.--70.~gados pieauga poļu inteliģences interese par \strong{dabaszinātnēm}. Poļu valodā tika pārtulkoti Č.~Darvina un E.~Hekkeļa darbi. Daudzus jauniešus tie piesaistīja dabas izpētei.

Statistiskajā fizikā pasaules nozīmes darbus veica M.~Smoluhovskis; augu morfoloģijas u anatomijas jomā sekmīgi strādāja poļu-vācu botāniķis E.~Strasburgers; adrenalīna darbības fizioloģisko darbību atklāja fiziologs N.~Cibuļskis; botāniķis E.~Godļevskis bija viens no augu anaerobās elpošanas procesa atklājējiem. Viņš noskaidroja, ka augi spēj uzsūkt vairākus simtus reižu vairāk ogļskābās gāzes nekā tā parasti ir gaisā.

Poļu \strong{bioķīmiķis} K.~Funks, kurš strādāja gan Polijā, gan aiz tās robežām, publicējot darbus par slimībām, saistītām ar vitamīnu trūkumu, 1912.~gadā ieviesa pašu jēdzienu ``vitamīni''. Jau vēlāk~--- 1936.~gadā viņam izdevās atšifrēt vitamīna B1 struktūru un sintezēt to. Viņš arī pirmais izolēja nikotīnskābi (vitamīnu B3).

Ļoti sekmīgi strādāja ievērojamākais poļu \strong{inženierbūvju} speciālists Z.~Jasinskis. Viņš projektēja un uzcēla daudzus dzelzceļa līnijas un stacijas Galīcijā.

Neapšaubāmi izcilākā poļu zinātniece bija divkārtējā Nobeļa prēmijas laureāte \strong{Marija Sklodovska-Kirī}. Savu atklājumu ceļu viņa uzsāka XIX~gadsimta beigās. Nespējot iegūt izglītību Polijā, viņa to 1891.~gadā atstāja un uzsāka studijas Sorbonas universitātē Parīzē. Kopā ar vīru Pjēru Kirī pētīja urāna radioaktivitāti. No 1898. līdz 1902.~gadam viņi pārstrādāja 8 tonnas urāna rūdas un ieguva 1/100 gr jaunas vielas~--- radija. Vēlāk tika iegūts vēl viens jauns elements~--- polonijs, nosaukts tā par godu M.~Sklodovskas-Kirī dzimtenei. Ar to abi zinātnieki atspēkoja tēzi par to, ka ķīmiskie elementi ir pastāvīgi un nemainīgi. Tika izdarīts pirmais solis pretī atoma sairšanas izpētei. 1903.~gadā viņa bija pirmā sieviete pasaulē, kura tika apbalvota ar Nobela prēmiju (fizikā). Vēlāk kļuva par pirmo sievieti~--- Sorbonas universitātes profesori. 1911.~gadā vēlreiz saņēma Nobela prēmiju (ķīmijā). M.~Sklodovskas-Kirī patriotisms izpaudās palīdzībā radioloģijas attīstībai Polijā. Viņas vizīte Varšavā 1913.~gadā izsauca entuziasmu plašu poļu tautas slāņu vidū. Pirmā pasaules kara laikā viņa organizēja apakšvienību ar pārvietojamiem rentgena aparātiem, kuri ļāva ievainotajiem izdarīt rentgena uzņēmumus turpat frontes līnijā. Sekas ilgajam darbam ar radioaktīvajiem materiāliem bija zinātnieces saslimšana ar leikēmiju, ko kuras viņa arī mira.

Daudzi poļu inteliģenti pievērsās \strong{humanitārajām zinātnēm}.

Filozofiskajā domā poļu inteliģence galvenokārt sekoja Rietumos notiekošajai attīstībai, neradot oriģinālus darbus. Toties Apgaismības ietekme spilgti izpaudās nacionālās \strong{valodas zinātnes} attīstībā. XVIII~gadsimta beigās literārās valodas funkcijas blakus poļu pildīja arī latīņu valoda, kaut arī tās lietošanas sfēra jau ievērojami sašaurinājās. Taču arī vairākums poļu valodā tapušo sacerējumu saturēja daudz vārdu un iestarpinājumu, aizgūtu no latīņu valodas. Tāpēc pakāpeniski izvērsās cīņa par poļu valodas tīrību. Garīgās dzīves sfērā latīņu valodas vietu ieņēma poļu (daļēji~--- arī franču) valoda. Poļu valodnieks O.~Kopčinskis jau XVIII~gadsimta beigās uzrakstīja ``\pltxti{Grammatika języka polskiego i łacińskiego dla szkól narodowych}'' (Poļu un latīņu valodas gramatika tautskolām), kura tika izdota vairāk nekā septiņdesmit reižu, un kuru vēlāk autors izdeva pārstrādātu kā ``\pltxti{Grammatika języka polskiego}'' (Poļu valodas gramatika). XIX~gadsimta sākumā izcils ieguldījums zinātnē bija S.~Lindes sagatavotā poļu valodas vārdnīca sešos sējumos. Ievērojamus panākumus poļu valodniecība guva XIX un XX~gadsimtu mijā. Krakovas akadēmija 1901.--1911.~gadā izdeva J.~Karloviča sastādīto ``Poļu dialektu vārdnīcu'' (\pltxti{Słownik gwar polskich}) sešos sējumos. Vienlaikus 1900.~gadā Varšavā parādījās ``Poļu valodas vārdnīcas'' (\pltxti{Słownik języka polskiego}), kuru J.~Karlovičs sastādīja kopā ar A.~Kriņski un V.~Nedzvecki, pirmais sējums. (Pēdējais~--- astotais~--- sējums iznāca 1927.~gadā.) Krakova kļuva par lingvistikas centru.

Liela loma poļu nacionālās pašapziņas un patriotisma jūtu audzināšanā bija \strong{vēstures zinātnei}. No nodibināja XVIII~gadsimta pēdējās trešdaļas sabiedriski-politisko domu nodarbināja problēma par Polijas nespējas cēloņiem pretoties Austrijas, Krievijas un Prūsijas teritoriālajām pretenzijām. Tas veda pie intereses par pagātni pieauguma. Publicisti plaši operēja ar vēstures faktiem. Taču Napoleona kari, kuros ierautas bija arī Polijas teritorijas, to pārdale traucēja vēsturiskās patiesības meklējumus. Sadalītājvalstis nebija ieinteresētas poļu vēsturnieku pētījumu izvēršanā.

XIX~gadsimta sākumā bijušās slavenās poļu universitātes Krakovā un Lembergā bija pakļautas ģermanizācijai, tāpēc poļu studenti vēsturi studēt devās uz ārzemēm, lielākoties Vāciju. Virkne poļu vēsturnieku bija mācījušies pie vācu vēsturniekiem. Tikai pēdējos XIX~gadsimta gadu desmitos, kad poļi atkal atguva ietekmi Krakovas un Lembergas universitātēs, poļu vēsturnieku darba apstākļi uzlabojās. Taču abās universitātēs bija štatu ziņā mazas vēstures katedras un daudzi apdāvināti pētnieki bija spiesti meklēt darbu citur.

Poļu historiogrāfija XIX~gadsimtā jau sākotnēji dalījās divos virzienos: demokrātiskajā un mērenajā. Pirmo pārstāvēja jau agrāk minētie Apgaismības ideju nesēji H.~Kollontajs un S.~Stašics. Abi viņi Polijas panīkuma cēloņus saskatīja tajā valdošajā politiskajā anarhijā, šļahtas kārtu privilēģiju saglabāšanā, magnātu egoismā.

H.~Kollontajs pēc T.~Kostjuško vadītās sacelšanās sakāves vairākus gadus pavadīja Austrijas cietumos. Pēc tam nodarbojās ar zinātnisku darbu vēstures, filozofijas, ētikas jomās. Vēsturi viņš atzina par zinātni, kam ir liela politiska un morāla nozīme, uzskatīja, ka ir kompleksi jāpēta visas cilvēka darbības sfēras: ražošana, tikumi, likumdošana, zinātne utt. Lielu nozīmi viņš piešķīra avotu vākšanai un to kritiskai analīzei. Galveno atbildību par Polijas sadalīšanu viņš iedalīja kaimiņvalstīm, kuras izmantoja Polijas panīkumu savtīgos mērķos.

S.~Stašics sabiedrības attīstību apskatīja kā objektīvu cīņas un sociālo konfliktu procesu, uzstājās pret šļahtas demokrātijas idealizāciju.

Mēreno virzienu pārstāvēja pirmkārt garīdznieka, vēlāk bīskapa, pirmā modernā poļu vēsturnieka A.~Naruševiča darbi. Viņš aizstāvēja stipra monarha varu, kura nodrošinātu līdzsvaru starp kārtām, garantētu kārtību valstī. Atbildīgus par Polijas novājināšanos viņš uzskatīja magnātus. Žečpospolitas novājināšanās cēloņus viņš saskatīja arī pārejā no poļu uz vācu tiesībām, kuras izrādījās izdevīgākas lielajiem zemes īpašniekiem un veda pie karaļa varas ierobežošanas. A.~Naruševičs nespēja izvairīties no šļahtiskā ``sarmatisma'' ideoloģijas ietekmes, slavināja šļahtu kā kārtu, kura pārstāv visu tautu, taču jēdzienā ``poļu tauta'' bez šļahtas iekļāva arī pilsētniekus.

Viņš apšaubīja vēsturisko leģendu patiesumu, aicināja balstīties uz laikabiedru liecībām, vākt dokumentus, kas nosacīja valsts uzbūvi un attīstību. Lai sagatavotu pamatīgu darbu par Polijas vēsturi, A.~Naruševičs radīja arhīvistu un kopētāju štatus, kuri vāca, apstrādāja, kopēja dokumentus. Tas ļāva savākt lielu materiālu daudzumu~--- 230~sējumu, kurus pēc tam izmantoja daudzas poļu vēsturnieku paaudzes.

A.~Naruševičam bija daudz sekotāju, kuri centās izplatīt Apgaismības idejas sabiedrībā. Taču pakāpeniski apdzisa Apgaismības tradīciju kritiskā virzība. Pārsvaru guva tendence idealizēt varonības momentus pagātnē, kam bija jāmodina patriotisms šļahtā. Ievērojamā loma šādā vēsturisko uzskatu attīstībā bija Varšavas ``Zinātnes draugu biedrībai''.

Tās biedrs J.~Ņemcevičs tāpat kā A.~Naruševičs savos darbos slavināja Polijas pagātni, tās valdniekus. Viņš arī attīstīja ideju par īpašu poļu raksturu, kam ir savas pozitīvās un negatīvās puses. 1816.~gadā viņš publicēja 33 tautas dziedājumu krājumu ``\pltxti{Śpiewy historyczne}'' (Vēsturiskās dziesmas) ar komentāriem. Turpmākajos 100 gados tas tika atkārtoti izdots 20~reizes, kaut autors sevišķi nerūpējoties par atbilstību patiesībai, centās izcelt poļu varonības lappuses pagātnē.

``Zinātnes draugu biedrības'' loceklis bija arī T.~Čackis, kurš galvenokārt nodarbojās ar tiesību vēstures jautājumiem, bet neatstāja bez uzmanības arī ekonomiskās un kultūras vēstures jautājumus. Tendence reabilitēt, izskaistināt pagātni, absolutizējot nacionālos momentus, tuvināja viņu romantiķiem.

Apgaismības idejas bija pārsvarā valdošās vēsturnieku darbos līdz XIX~gadsimta 20.~gadiem, tad to vietā nāca romantisms. Tas gan centās sintezēt Apgaismības un romantisma ideālus, taču priekšplānā izvirzīja nācijas, nacionālas valsts un kultūras savdabības idejas. Tās saturēja vēsturiskās evolūcijas kā procesa, sastāvoša no vairākiem etapiem, atzīšanu. Romantismam bija raksturīgs arī subjektīvisms, pagātnes apoloģija, taču nacionālās atbrīvošanas kustības apstākļos romantiskā historiogrāhija veicināja nacionālās pašapziņas izaugsmi, patriotisku noskaņojumu veidošanos. Historiogrāfija izrādījās politiskās cīņas centrā, vēsturnieku darbi kalpoja kā ieroči polemikā starp dažādiem sabiedriskajiem spēkiem.

Romantiskā virziena iedibinātājs un galvenais pārstāvis poļu historiogrāfijā bija \strong{Joahims Lelevels}, vairāku augstskolu profesors, aktīvs sabiedriska darbinieks. Pēc 1830.~gada sacelšanās sakāves dzīvoja Rietumos, bija viens no poļu emigrācijas pilsoniski demokrātiskā spārna līderiem, uzturēja draudzīgas attiecības ar K.~Marksu, kontaktējās ar krievu progresīvajiem sabiedriskajiem darbiniekiem. Šajā periodā vēsturnieks~--- romantiķis arī radīja savus vērtīgākos darbus. Poļu vēsturnieks M.~Friedbergs salīdzinājis J.~Lelevela darbības apstākļus ar tai pat laikā dzīvojošā tāda pat mēroga vācu vēsturnieka L.~f.~Rankes dzīves un darba apstākļiem. Ja vācu vēsturnieks varēja mierīgi strādāt Prūsijas un Bavārijas karaļu aizbildniecībā, viņa poļu kolēģis kopš 1831.~gada vadīja savu dzīvi emigrācijā, cīnoties ar materiālām grūtībām.

J.~Lelevela uzmanības centrā bija Polijas vēsture, tai bija veltīts vairums viņa pētījumu, taču virkni darbu viņš uzrakstīja arī par vispārējo vēsturi, vēstures zinātnes palīgdisciplīnām.

J.~Lelevels, tāpat kā citi romantiķi, idealizēja slāvu kopienas iekārtu un uzskatīja, ka to iztraucēja ``rietumu civilizācijas'' iespiešanās. Ceļu uz Polijas atdzimšanu viņš redzēja kā atgriešanos pie ``nacionālajām saknēm'', demokrātijas un taisnīguma ieviešanu. Par poļu ``mūžīgo ideālu'' nesējiem viņš uzskatīja tautu, to saprotot kā sīko šļahtu un zemniecību. Žečpospolitas novājināšanās cēloņus viņš meklēja tās sociāli ekonomiskajās un politiskajās struktūrās, zemnieku jautājuma neatrisinātībā un kaimiņvalstu iekarošanas politikā. 1794. un 1830.--1831.~gada sacelšanos sakāves viņš izskaidroja ar to vadītāju bailēm no tautas masām un centieniem tās atbīdīt no aktīvas dalības cīņā. Pats viņš aktīvi iestājās par zemnieku atbrīvošanu no dzimtbūšanas ar zemi, kā arī politisku tiesību piešķiršanu tiem. Tikai tas, kā viņš uzskatīja, var nodrošināt nacionālās atbrīvošanās kustības panākumus. Kā tās sabiedrotos viņš redzēja Rietumeiropas un Krievijas tautu atbrīvošanās kustības. Šos secinājumus vēsturnieks un politiķis izmantoja arī politiskajā cīņā, atmaskojot poļu magnātus un to realizēto ``poļu tautai naidīgo'' monarhistisko politiku. Viņš rakstīja: ``\dots{}ir mirusi šļahtas Polija, nebrīves un privilēģiju Polija, Polija, kuru pārstāvēja tikai viena tās iedzīvotāju šķira. Bet tā Polija, kura atdzims, būs tautas Polija.'' Augstu J.~Lelevela darbus vērtēja K.~Markss, kurš rakstīja: ``Ar rūpīgo ekonomisko apstākļu izpēti, kas pārvērta brīvo poļu zemnieku no brīva cilvēka par dzimtcilvēku vecais Lelevels ir izdarījis daudz vairāk savas dzimtenes paverdzināšanas cēloņu noskaidrošanā, nekā vesels lērums rakstnieku, kuru visa bagāža sastāvēja no Krievijai adresētiem lamuvārdiem.'' XIX~gadsimta 20.--50.~gados J.~Lelevela koncepcija kalpoja cīņai par demokrātiju, agrārā jautājuma atrisināšanu, nacionālās pašapziņas atraisīšanu, Polijas buržuāziski demokrātisku pārkārtošanu. Viņa darbi atstāja lielu ietekmi uz daudziem poļu revolucionāriem.

J.~Lelevela uzskatus atbalstīja daudzi XIX~gadsimta pirmās puses poļu vēsturnieki. Taču vairākums no viņiem senslāvu kopienas idealizāciju attiecināja arī uz vēlākiem laikiem, tā pūloties reabilitēt šļahtu, centās modināt poļu nacionālo pašapziņu, tai pat laikā apejot sociālos jautājumus. Redzamākie no tiem bija J.~Moračevskis, K.~Šainoha un V.~Macijevskis.

Taču XIX~gadsimta 40.~gados daži J.~Lelevela ietekmēti vēsturnieki mēģināja attīstīt arī viņa koncepcijas sociālo pusi. Viens no tiem bija Krakovas sacelšanās vadītājs 1846.~gadā E.~Dembovskis. Pagātnes notikumos viņš saskatīja šķiru cīņu, kura, pēc tā domām, neizbēgami ved pie revolūcijas, kā rezultātā jānodibinās iekārtai, brīvai no ekspluatācijas. Progresu viņš traktēja ideālistiski, kā brīvības un taisnīguma ideju attīstību. Viss vēsturiskais process tika skatīts kā cilvēciskā saprāta pilnveidošanās. Polijas pagātnē viņš saskatīja zemnieku cīņas pret muižniekiem vēsturi. E.~Dembovskis arī cīnījās pret poļu demokrātiskajos slāņos izplatītajām ilūzijām, ka šļahtas patriotisms var to piespiest atteikties no šķiriskajām interesēm. Pēc viņa vērtējuma nacionālās atbrīvošanās kustības pamats bija tautas masas.

Pret romantiķiem uzstājās muižnieciski monarhistiski noskaņotie vēsturnieki, kuri galvenokārt sekoja A.~Naruševiča monarhistiskajiem ideāliem. Ievērojamākais no tiem bija K.~Hofmans. Viņš uzstājās pret pagātnes idealizāciju, romantiķiem raksturīgajiem ``nacionālā gara'' meklējumiem tajā.

Rietumos 1848.--1849.~gada revolūciju rezultātā radusies pozitīvisma filozofija ietekmēja arī vēstures zinātni. Notika atteikšanās no secinājumiem, kuri balstījās vienīgi uz prāta spekulācijām, kas neatbilda vēsturiskajai īstenībai. Tie turpmāk bija jābalsta uz stingru faktu analīzi. Poļu zemēs pozitīvisma izplatība norisa reizē ar romantisma vēstures zinātnē kritiku kā no liberālās pilsonības, tā arī šļahtas konservatīvo aprindu puses. Taču nacionālo spaidu pieaugums Vislas novadā pēc 1863.--1864.~gada sacelšanās bremzēja kultūras un zinātnes attīstību. Labākos apstākļos atradās 60.~gados autonomiju saņēmusī Galīcija. Pie Krakovas Zinību akadēmijas (\pltxti{Akademia Umiejętności}) nodibināta komisija nodarbojās ar vēstures avotu vākšanu un izdošanu. Tā publicēja trīs sērijas dokumentu izdevumu (51 sējums), izdeva krājumus, žurnālus, monogrāfijas. 1880.~gadā pēc tās iniciatīvas Krakovā notika poļu vēsturnieku I~kongress. 1886.~gadā nu jau Lembergā tika izveidota Poļu Vēsturiskā biedrība (\pltxti{Polskie Towarzystwo Historyczne}). Pozitīva nozīme bija vēstures žurnāla ``\pltxti{Kwartalnik historycny}'' (Vēstures ikceturkšņa žurnāls) izdošanai no 1886.~gada Lembergā un ``\pltxti{Przegląd historycny}'' (Vēsturiskais apskats) iznākšanai no 1905.~gada Varšavā.

XIX~gadsimta otrajā pusē poļu historiogrāfijā vadošās bija divas skolas: Krakovas, cieši saistīta ar Galīcijas muižnieciski-pilsoniskajiem slāņiem, un Varšavas, kura orientējās uz Polijas karalistes liberāli-pilsoniskajām aprindām.

\strong{Krakovas vēstures skola} lielā mērā atbalstīja šļahtas ideoloģiju, katoļu pasaules uzskatu. Krakovas skolas pārstāvji atradās opozīcijā J.~Lelevela romantiskajai, demokrātiskajai vēstures skolai. Viņi veltīja ne mazums pūļu lai pamatotu politisko kursu, kurš veda uz izlīgumu ar Habsburgu monarhiju.

Reizē Krakovas skolai bija raksturīga padziļināta avotu kritika, detalizēta līdz tam maz pētītu Polijas vēstures periodu izstrāde. Tās pārstāvji V.~Kaļinka, J.~Šujskis, M.~Bobžinskis, S.~Smolka u.c. stāvēja nacionālistiskās pozīcijās, iestājās par ukraiņu un baltkrievu pārpoļošanu. Atzīts Krakovas skolas dibinātājs bija V.~Kaļinka. Kaut jaunībā piedalījies nacionālās atbrīvošanas kustībā, pēc 1863.--1864.~gadu sacelšanās viņš tuvinājās konservatīvajiem, idealizēja patriarhālās attiecības poļu laukos, pat čiņša maksājumu ieviešanu zemniecībai XVIII~gadsimta beigās uzskatīja par priekšlaicīgu, asi kritizēja Apgaismības laikmeta darbiniekus, reformatorus. Viņa darbos pozitīvi vērtēta katoļu baznīcas loma Polijas vēsturē, toties Apgaismības idejām esot bijusi graujoša ietekme uz poļu tautu.

Par Krakovas skolas vadītāju tika uzskatīts J.~Šuiskis. Arī viņš no atbrīvošanās kustības dalībnieka evolucionēja pie labējiem uzskatiem. J.~Šuiskis aicināja atteikties no revolucionāras darbības, sadarboties ar Austrijas valdību. Ļoti naidīga bija viņa attieksme pret Krieviju un tās tautu. Nozīmīgāko vietu vēstures procesā J.~Šuiskis piešķīra providencei, Dieva gribai. Pretēji J.~Lelevelam viņš apgalvoja, ka Polijas jau sākotnē nevaldījusi ``demokrātija'', bet ``kņazu absolūtisms''. Vēlāk gan tādu poļiem raksturīgu īpašību kā ``tieksme uz anarhiju'' rezultātā monarhistiskā vara novājinājusies. Atbildību par poļu valstiskuma vājumu viņš uzkrāva poļu tautai, uzskatīja, ka atrodoties tālu no Romas un citiem Eiropas centriem, Polija savā attīstībā visu laiku atpalikusi no Rietumiem. J.~Šuiska darbus caurvij stipras monarhistiskas varas ideja, kuras īstenošanas paraugu viņš saskatīja Austroungārijā.

Arī M.~Bobržinskis pretstatīja monarhiju tautai. Viņš lielu uzmanību veltīja valsts tiesību jautājumiem, stiprā valsts birokrātiskajā aparātā saskatīja sabiedriskās kārtības garantiju. Parlamentārisms un šļahtas brīvības, kuras bija pretrunā ar ``kārtību un likumu'', pēc viņa domām bija vienīgie Polijas krišanas cēloņi.

S.~Smolka, būdams atzīts Polijas viduslaiku pētnieks, lielu vērību veltīja sociālajām problēmām. Analoģijas Polijas vēsturei viņš meklēja citu slāvu tautu pagātnē. Reizē S.~Smolka idealizēja katoļu baznīcas lomu, pierakstot tai milzu lomu nacionālās pašapziņas atmodināšanā, feodālās sašķeltības likvidēšanā un ``civilizācijas'' izplatībā uz Austrumiem.

Konservatīvie poļu ideologi Polijas vēstures kritiku uzskatīja par nepatriotisku. Pret šādiem uzskatiem vērsās poļu nacionālās atbrīvošanās kustības, 1830.~gada sacelšanās dalībnieks, publicists J.~Janovskis, kurš par patriotisku rīcību uzskatīja nevis savas dzimtenes kaišu slēpšanu, bet to nožēlošanu, nevis pagātnes idealizāciju, bet tikai labāko nacionālo tradīciju pārņemšanu.

Krakovas skolai sevi pretnostatīja \strong{Varšavas skola}: T.~Korzons, V.~Smoļenskis u.c. Ja Krakovas vēsturnieki pārsvarā nodarbojās ar politiskās vēstures problēmām, tad varšavieši galvenokārt nodarbojās ar Polijas XVIII~gs. vēstures, kad Žečpospolitā nobrieda kapitālistiskās iekārtas elementi, pētniecību, īpašu uzmanību veltījot ekonomikas un kultūras attīstībai, kritizēja kārtu sistēmu, klerikālismu. Atbrīvošanās cīņas lozungiem viņi pretstatīja kapitālistiskās attīstības nepieciešamību, buržuāziskā progresa vārdā bija gatavi izlīgt ar carismu. Kā liberālās buržuāzijas ideologi Varšavas skolas pārstāvji centās uzsvērt pilsētu iedzīvotāju nozīmību Polijas vēsturē. Virknē jautājumu (par Žečpospolitas krišanas cēloņiem, par aristokrātijas, katoļu baznīcas lomu) viņu uzskati bija tuvi Apgaismības un romantisma tradīcijām. Tomēr Varšavas skola noraidīja sacelšanos nepieciešamību neatkarības izcīnīšanai, par vienīgo pieņemamo uzskatīja reformu ceļu, kuras veic paši valdošie slāņi.

T.~Korzons bija pazīstams kā vairāku monogrāfiju, mācību grāmatu vispārīgajā vēsturē, kā arī darbu par Žečpospolitas vēsturi XVIII~gadsimta otrajā pusē autors. Viņš izmantoja milzīgu daudzumu publicēto un arhīvu dokumentu, pētniecībā plaši pielietoja statistisko un salīdzinošo metodes. Uz šī pamata viņš pierādīja, ka XVIII~gadsimta otrajā pusē Žečpospolitā buržuāzisko attiecību veidošanā bija jau ievērojami pavirzījusies uz priekšu. Vecie spēki (aristokrātija, augstākie baznīckungi), pēc viņa domām, veda valsti pretī iznīcībai, bet jaunie (izglītotākie pilsētnieki un patriotiskā šļahta)~--- uz atdzimšanu un glābiņu. Magnāti un to sabiedrotie, nevēloties nekādas pārmaiņas, atdevuši valsti citu kaimiņvalstu varā. T.~Korzons gan bija sašutis par šļahtas konservatīvismu un egoismu, taču reizē neatzina agrārā jautājuma atrisināšanas revolucionāru ceļu.

V.~Smoļenskis savos darbos galveno vērību veltīja poļu kultūras un sabiedriski politiskās domas attīstībai XVIII~gadsimta otrajā pusē. Savos agrīnajos sacerējumos viņš atmaskoja šļahtas negatīvo lomi vēsturē, uzsvēra pilsētnieku lomu tā laika reformu kustībā. XX~gadsimta sākumā viņš evolucionēja labējo virzienā, tuvinājās nacionālajiem demokrātiem.

Starp abām vēstures skolām norisa polemika. Piemēram, T.~Korzons asi kritizēja Krakovas vēsturniekus par vēstures zinātnes pakļaušanu Galīcijas valdošo slānu interesēm. V.~Smoļenskis uzsvēra, ka Krakovas skolas vēsturnieku doktrīnai ir raksturīgs konservatīvisma, katolicisma un aristokrātijas apoloģētika. Taču abu skolu polemika nenesa pārsvaru kādai no pusēm. Diskusija pakāpeniski izsmēla sevi. Marksistiskie vēsturnieki apgalvoja, ka tas noticis sakarā ar strādnieku lomas pieaugumu sabiedrībā, kas lika abu skolu pārstāvjiem, vadoties no savām šķiriskajām interesēm, meklēt kopīgu valodu. Taču jāņem vērā, ka XIX~gadsimta beigās un XX~gadsimta sākumā historiogrāfijā parādījās arī jauni virzieni, kuri priekšplānā izvirzīja jau citus jautājumus.

Vēl viena vēstures zinātnes skola veidojās Austrijai piederošajā Lembergā. Tur gan vēsturnieku bija mazāk nekā Krakovā un Varšavā. Toties viņi arī bija mazāk saistīti ar valdošo slāni nekā Krakovā. Skolas dibinātājs bija K.~Liske, pazīstams kā daudzu kritisku rakstu, apskatu, recenziju autors. Pēc viņa iniciatīvas 1868.~gadā sāka iznākt dokumentu krājums ``\pltxti{Akta Grodzkie i Ziemskie}'' (Pilsētu un zemju akti). Ievērojams medievistikas speciālists bija T.~Vojcehovskis. Vēsturisko attīstību viņš uzskatīja par atkarīgu galvenokārt no teritoriāli etniskā faktora. Iebilstot vācu vēsturniekiem, viņš piestādīja pārliecinošus pierādījumus tam, ka slāvi ir poļu zemju autohtonie (kas no laika galā dzīvo šai teritorijā) iedzīvotāji.

XIX~gadsimta 60.--90.~gados pieauga interese par vēsturi poļu sabiedrībā. Plašas iedzīvotāju masas (tiesa gan, galvenokārt Galīcijā) piedalījās tādu vēsturisku jubileju kā Jagelloņu universitātes (\pltxti{Uniwersytet Jagielloński}, 1364) piecsimtgades, Lubļinas ūnijas (\pltxti{Unia lubelska}, 1569) trīssimtgades, J.~Sobeska uzvaras pār osmaņiem pie Vīnes (1683) divsimtgades, 3.~maija Konstitūcijas (1791) un T.~Kostjuško sacelšanās (1794) simtgades svinībās.

XIX un XX~gadsimtu mijā, kad daļa poļu pilsonības un inteliģences bija vīlusies pozitīvisma ietekmē esošo vēsturnieku piedāvātājā vēstures kā taisnvirziena evolucionāra procesa attēlojumā un izvirzīja Polijas autonomijas vai pat neatkarības izcīnīšanas uzdevumu, pilsoniskā historiogrāfija nespēja pārliecinoši izskaidrot pagātni. Pēc III poļu vēsturnieku kongress 1900.~gadā radās jauni virzieni vēstures zinātnē: vēsturiski-juridiskais, ekonomiskais un neoromantiskais. Vēsturnieki aktīvāk iesaistījās politiskajās cīņās.

Vēsturiski-juridiskā virziena dibinātājs bija Lembergas universitātes profesors O.~Balcers. Viņš pakāpeniski evolucionēja no ``pesimistiska'' uz ``optimistisku'' Polijas vēstures skatījumu, centās reabilitēt Žečpospolitas tiesiskos un politiskos institūtus. Salīdzinot Žečpospolitas iekārtu ar dažu Rietumeiropas valstu tiesiskajiem institūtiem, viņš secināja, ka tā atbilda laika garam. Analizējot vienīgi tiesību normas un institūtus un neievērojot sabiedriski politiskos apstākļus, kādos tie darbojās, viņš atbildību par valsts sabrukumu noņēma no iekšējiem spēkiem un pārcēla uz Polijas kaimiņiem. To pašu virzienu pārstāvēja S.~Kutšeba, kurš gan centās jau ievērot sociāli ekonomisko faktoru ietekmi uz tiesisko institūtu veidošanos. Viņš uzskatīja, ka valsts iekārtas tiesiskās formas poļi pārņēma no Rietumiem un tikai centrālās varas novājināšanās rezultātā no XVI~gadsimta Polija sāka atpalikt attīstībā no Rietumeiropas.

Nacionālās atbrīvošanās kustības aktivizēšanās XX gadsimta sākumā atjaunoja nacionāli romantiskos ideālus. Nacionālo problēmu atrisināšanas labā tika aizmirstas, atbīdītas otrajā plānā šķiriskās pretrunas. Viens no neoromantiskā virziena iedibinātājiem bija S.~Zakševskis. Viņš uzstājās pret Krakovas vēstures skolu, uzskatot, ka tā noniecinot poļu tautas sasniegumus. S.~Zakševskis arī atmeta pozitīvisma metodoloģiju, uzsverot, ka pagātnes izzināšana nevar balstīties vienīgi uz faktiem, bet vēstures izpratnei jāņem vērā arī intuīcija. Viņš aizstāvēja uzskatu, ka vēsturi veido spēcīgas personības, centās visas vēsturiskās parādības skatīt cīņas par Polijas neatkarību gaismā. Vēsturnieka un viņa sekotāju darbos nacionālā ideja ieguva lielvalstisku raksturu. Ne velti neoromantiskā koncepcija kļuva par oficiālo ``sanācijas'' režīma laikā (1926--1939). Arī V.~Sobeskis absolutizēja nacionālo ideju, saistīja to ar nacionālo psiholoģiju, kura it kā nosakot zemes likteni. Neoromantisko ideju ietekmē atradās arī S.~Aškenazi. Atšķirībā no citiem neoromantiķiem viņa darbiem raksturīga godprātīga attieksme pret vēstures avotiem. Taču arī viņam nebija sveša ``varonības'' izcelšana poļu vēsturē, pastiprināta uzmanība nacionālajam jautājumam. Nacionālās atbrīvošanās kustības vēsturē viņš izcēla tās antikrievisko virzību, taču viņa simpātijas piederēja mēreni liberālajiem darbiniekiem, kuriem tas pūlējās piešķirt varonības oreolu. Darbā ``\pltxti{Napoleon a Polska}'' (Napoleons un Polija, 1918--1919) vēsturnieks centās izcelt Polijas jautājuma vietu Eiropas dzīvē. Viņš bija arī pirmais, kurš Polijas sadalīšanas laiku uzskatīja par izšķirošu poļu nācijas izveides gaitā.

Kapitālistiskās iekārtas attīstība, ekonomiskā materiāla izmantošana agrāko paaudžu vēsturnieku darbos, socioloģijas attīstība Rietumos veicināja intereses pieaugumu par sociāli ekonomisko problemātiku poļu vēsturniekos. Viņi pievērsās jauna veida avotiem~--- statistikas krājumiem, dažādiem saimnieciskiem rēķiniem, inventāra sarakstiem utt. Rezultātā noformējās ekonomiskais virziens poļu vēsturnieku darbos.

Zināma ietekme virziena izveidē bija ekonomista un vēsturnieka, vēlāk divkārtēja Polijas premjerministra V.~Grabska darbiem. Viņš pētīja agrārās attiecības un poļu zemnieku stāvokli XIX~gadsimtā, nonāca pie secinājuma, ka kapitālisms, laužot sev ceļu ekonomikā, prasīja veikt pārkārtojumus, taču tai pat laikā iestājās par muižnieku un turīgo zemnieku pozīciju saglabāšanu.

Viens no tautsaimniecības vēstures kā zinātniskas disciplīnas aizsācējiem bija F.~Bujaks, kurš pētīja Galīcijas agrāro vēsturi, piesaistot arī etnogrāfijas, statistikas, ģeogrāfijas materiālu. Tiesa, ekonomiskās problēmas tika skatītas atrautībā no politiskajām.

Uz J.~Rutkovska kā vēsturnieka attīstību ievērojamu ietekmi atstāja franču ``Annāļu'' skola (\frtxti{École des Annales}). Pētot agrārās attiecības Polijā XVII--XVIII~gadsimtā, viņš parādīja zemnieku stāvokļa pasliktināšanos sakarā ar klaušu-folvarku sistēmas attīstību.

Kopumā pilsoniskie vēsturnieki radīja darbus, kuros bagāts faktu materiāls, skrupuloza tā analīze aizsedza dziļu secinājumu trūkumu par Polijas vēsturiskās attīstības galvenajām problēmām. Poļu profesionālie vēsturnieki gandrīz pilnībā apgāja nesenās pagātnes notikumus, veltot savus darbus galvenokārt viduslaiku posmam. Polijas pagātne vairākumā vēsturnieku darbu bija tēlota idealizētas savdabības toņos. Poļu vēstures savdabības idealizācija uz ilgu laiku nostiprinājās poļu vēsturnieku uzskatos.

Blakus minētajiem pilsoniskajiem virzieniem vēstures zinātne attīstījās arī marksistiskais virziens.

Kad 70.--80.~gados poļu zemēs sāka izplatīties marksisma mācība, pirmie poļu marksisti aicināja visus notikumus vēsturē skatīt kā šķiru cīņas izpausmes. Poļu valodā tika pārtulkoti vairāki K.~Marksa un F.~Engelsa darbi. Aktīvi pret historiogrāfiskiem darbiem, kuri centās izplatīt uzskatu, ka poļu tauta pagātne esot bijusi vienota, uzstājās L.~Varinskis. Antropologs, sociologs un ekonomists un L.~Kšivickis no marksistiskām pozīcijām kritizēja pozitīvismu.

Turpmāk marksistiskajā virzienā pastāvēja divas strāvas: revolucionārā un reformistiskā.

Vispirms par revolucionāro. XIX~gadsimta beigās iznāca poļu marksistu J.~Marhļevska un R.~Luksemburgas darbi, veltīti Polijas vēstures problēmām. Pirmais 1897.~gadā izdeva grāmatu ``\pltxti{Fizjokratyzm w dawnej Polsce}'' (Fiziokrātisms vecajā Polijā), kurā attēloja, kā fiziokrāti XVIII~gadsimta otrajā pusē centās ierobežot šļahtas un garīdzniecības patvaļu, bet tas tiem neizdevās un viņi nespēja izvest valsti no pagrimuma. 1903.~gadā darbā ``\pltxti{Stosunki społeczno-ekonomiczne pod panowaniem pruskim}'' (Sociāli ekonomiskās attiecības Prūsijas varā esošajās poļu zemēs), J.~Marhļevskis rādīja, ka poļu rietumu zemju pievienošana Prūsijai veda šeit pie rūpniecības panīkuma un feodālās kārtības iekonservēšanas laukos.

Jau pirmajā R.~Luksemburgas darbā, kurš iznāca vācu valodā, ``\detxti{Die industrielle Entwickelung Polens}'' (Rūpniecības attīstība Polijā) uz statistisko datu pamata viņa pierādīja, ka XIX~gadsimta pēdējā trešdaļā poļu zemēs ir uzvarējušas kapitālistiskās attiecības. Diemžēl, viņa arī kļūdaini uzskatīja, ka šādos apstākļos kustība par Polijas neatkarību ir utopija un poļu strādnieku kustībai cīņa par nācijas pašnoteikšanos ir neaktuāla. J.~Marhļevskis un R.~Luksemburga arī uzsāka poļu strādnieku kustības vēsturi izpēti.

Blakus revolucionāri orientētiem marksistiskiem darbiem tapa arī reformistiski. To autori~--- labējo sociālistu pārstāvji bija noskaņoti nacionālistiski. Viens no šī virziena un arī PPS dibinātājiem bija Latgalē dzimušais žurnālists, vēsturnieks un sociologs B.~Ļimanovskis. Viņa pirmā publikācija bija ``\pltxti{Listy z Inflant}'' (Vēstules no Livonijas, 1860), kurā bija tēlota latviešu nacionālās kustības vēsture. B.~Ļimanovskis rakstīja arī par Polijas 1830.~gada sacelšanās dalībnieci E.~Plāteri. Galvenie viņa darbi bija veltīti poļu tautas cīņai par neatkarību, poļu demokrātijas vēsturei pēc Žečpospolitas sadales. B.~Ļimanovska uzskati laika gaitā gan mainījās. Tā, piemēram, no pozitīva uz negatīvu mainījās viņa vērtējums par poļu un krievu revolucionāru sadarbību. Revolucionāri viņam pārmeta šķiru cīņas un revolucionāro spēku lomas nenovērtēšanu poļu atbrīvošanās kustībā.

Publicists un politiskais darbinieks L.~Kuļčickis daudz uzmanības veltīja sociālistiskās kustības vēsturei. Uz plašas faktiskā materiāla pamata par dažādu strādnieku organizāciju darbību viņš parādīja, ka poļu strādnieku kustībā pastāv divi: internacionālistiskais un nacionālistiskais strāvojumi.

Viens no PPS vadītājiem F.~Perls grāmatā ``\pltxti{Dzieje ruchu socjalistycznego w zaborze rosyjskim}'' (Sociālistiskās kustības vēsture Polijas karalistē, 1910) poļu revolucionāro kustību pretstatīja krievu revolucionārajai kustībai, kurai esot raksturīgas ``teroristiskas'' cīņas formas.

Dotajā nodaļā apskatītajā laika posmā, īpaši XIX~gadsimta pirmajā pusē, poļu nacionālā kultūra attīstījās strauji un vispusīgi. Nostiprinājās tās sakari ar Rietumeiropas, kā arī Krievijas un citu slāvu tautu kultūru. Taču tai pat laikā poļu kultūra ne tikai nezaudēja savu nacionālo veidolu, bet par spīti nacionālajai sadalītībai triju valstu sastāvā, vēl vairāk nostiprināja sev raksturīgās nacionālās iezīmes. Sākotnēji feodālais kultūras veidols pakāpeniski ieguva tik daudz pilsonisku jeb buržuāzisku iezīmju, ka tās jau XIX~gadsimta 60.~gados kļuva par noteicošām. Tas izpaudās kā izglītības un zinātnes pacēlumā, garīgo vērtību radīšanas un patērēšanas procesa demokratizācijā, tā arī radošās inteliģences statusa izmaiņās. Vēsturiskās tradīcijas caurvija visas kultūras jomas, tā stiprinot tās vienotību, ko apdraudēja sadalīšana dažādās valstīs. Tas lielā mērā arī izskaidroja romantisma dzīvotspēju. Tas sadzīvoja ar reālismu arī pozitīvisma nostiprināšanās apstākļos, bet pēc tam piedzīvoja atdzimšanu.

\section{Dokumenti 1.~nodaļai}

\subsection[Izraksts no N.~Karamzina vēstules Aleksandram~I sakarā ar Polijas atjaunošanas projektu]{Izraksts no N.~Karamzina\footnote{N.~Karamzins~--- izcils vēsturnieks, ``Krievijas valsts vēstures'' (1--12 sēj., 1803.--1826.) autors, sentimentālisma pārstāvis krievu literatūrā.} vēstules Aleksandram~I sakarā ar Polijas atjaunošanas projektu}

\dots{}Jūs domājat par \emph{senās} Polijas karalistes atjaunošanu, bet vai šāda atjaunošana saskan ar Krievijas valstiskās labklājības likumu? vai saskan ar Jūsu svētajiem pienākumiem, Jūsu mīlestību pret Krieviju un pret taisnīgumu? Pirmkārt, (nerunājot par Prūsiju) jautāju: vai Austrija labprātīgi atdos Galīciju? Vai Jūs, \emph{Svētās Savienības}\footnote{Svētā Savienība~--- (angļu: \entxti{Holy Alliance}, \entxti{Grand Alliance}, franču: \frtxti{La Sainte-Alliance}, krievu: \rutxti{Священный союз}, vācu: \detxti{Heilige Allianz}) Eiropas valstu pakts par starptautisko attiecību pamatprincipiem alianses ietvaros, ko Vīnes kongresa laikā 1815.~gada janvārī noslēdza trīs Napoleonu uzvarējušo kontinentālo lielvalstu Austrijas, Krievijas un Prūsijas valdnieki. 1815.~gada 20.~novembrī Parīzē aliansei pievienojās arī Lielbritānijas karaliste, bet 1818.~gadā arī Francijas karaliste.} radītājs, varat pieteikt tai karu, nesaderīgu ne tikai ar kristīgo ticību, bet arī valstisko taisnīgumu? Jo Jūs pats atzināt Galīciju par likumīgu Austrijas valdījumu. Otrkārt, vai varat ar mierīgu sirdsapziņu atņemt mums Baltkrieviju, Lietuvu, Volīniju, Podoliju, vēl pirms Jūsu valdīšanas apstiprinātu Krievijas īpašumu? Šīs zemes jau bija \emph{Krievija}, kad metropolīts Platons\footnote{Metropolīts Platons 1801.~gada 15.~septembrī kronēja Aleksandru I.} Jums pasniedza Monomaha, Pētera un Katrīnas, kuru Jūs pats nosaucāt par \emph{Lielo}, kroni. Vai kāds teiks, ka viņa nelikumīgi sadalīja Poliju? Taču Jūs rīkotos vēl nelikumīgāk, ja iedomātos nolīdzināt viņas netaisnību ar pašas Krievijas sadali. Mēs ieguvām Poliju ar zobenu: lūk mūsu tiesības, kurām visas valstis ir parādā par savu esamību, jo visas tās sastāv no iekarojumiem. Katrīna ir atbildīga Dievam, ir atbildīga Vēsturei par savu veikumu, bet tas ir paveikts, un Jums jau ir svēts: Jums Polija ir likumīgs Krievijas valdījums. Politikā nav \emph{vecu cietokšņu}, citādi mums būtu jāatjauno gan Kazaņas, gan Astrahaņas caristes, gan Novgorodas republika, Rjazaņas lielkņaziste un daudzas citas. Turklāt arī \emph{vecie cietokšņi} Baltkrievija, Volīnija, Podolija? kopā ar Galīciju kādreiz bija iesakņots Krievijas valdījums. Ja Jūs atdosiet tos, tad Jums pieprasīs gan Kijevu, gan Čerņigovu, gan Smoļensku, jo tās tāpat ilgi piederēja naidīgajai Lietuvai. Vai nu visu vai neko. Līdz šim mūsu valsts likums bija: \emph{ne pēdas, ne naidniekam, ne draugam!} Napoleons varēja iekarot Krieviju: bet Jūs, kaut arī patvaldnieks, nevarējāt atteikties viņa labā kaut no vienas krievu būdas. Tāds ir mūsu raksturs un valsts gars. Vai Jūs, mīlot likumīgo pilsonisko brīvību, pielīdzināsiet Krieviju bezdvēseliskam, bezvārda īpašumam? Vai patvaļīgi dalīsiet to daļās un dāvāsiet tās, kam par labu uzskatīsit? Krievija, Valdniek, ir mēma Jūsu priekšā, taču ja atjaunotos senā Polija (lai Dievs nedod!) un kādreiz laistu pasaulē vēsturnieku godājamu, patiesu, objektīvu, tad viņš, Valdniek, nosodītu Jūsu augstsirdību, kā kaitīgu Jūsu īstenajai tēvijai, krietnajai, stiprajai Krievijai. Šis vēsturnieks pateiktu pavisam ne to, ko tagad saka Jums poļi, kuriem mēs piedodam, bet Jums mēs, krievi, nepiedotu, ja Jūs viņu aplausu dēļ iegrūstu mūs izmisumā. Valdniek, tagad slavenais, dižais, laipnais! atbildu Jums ar savu galvu par šādu nenovēršamu pilnīgas Polijas atjaunošanas iedarbību. Es dzirdu krievus un pazīstu viņus: mēs zaudētu ne tikai brīnišķīgus apgabalus, bet arī mīlestību pret caru; mūsu dvēseles atsaltu pret tēvzemi, redzot to kā patvaldīgas patvaļas spēļu laukumu; vājinātos ne tikai ar valsts samazināšanos, bet arī garā; pazemotos citu un savā priekšā. Pils, protams nekļūtu tukša: Jums arī tad būtu ministri, ģenerāļi: bet viņi kalpotu ne tēvzemei, bet vienīgi savam personīgajam izdevīgumam, kā algotņi, kā īsti vergi \dots{} Bet Jūs, Valdniek taču nicināt verdzību un gribat dot mums brīvību!

Vārdu sakot \dots{} un Dievs siržu pazinējs lai slēdz man lūpas šai pašā brīdī, ja saku Jums nepatiesību \dots{} vārdu sakot, Polijas atjaunošana būs Krievijas krišana, vai mūsu dēli slacīs poļu zemu savām asinīm un atkal triecienā ieņems Prāgu!

Nē, Valdniek, nekad poļi nebūs mums ne īsti brāļi, ne uzticami sabiedrotie. Tagad viņi ir vāji un niecīgi: vājie nemīl stipros; kad nostiprināsiet viņus, tie gribēs neatkarību, un pirmais mēģinājums būs atkāpšanās no Krievijas, protams, ne Jūsu valdīšanas laikā, taču Jūs, Valdniek, skatāties tālāk par savu mūžu, un ja neesat nemirstīgs ar ķermeni, tad esat nemirstīgs ar slavu! Valsts lietās jūtas un pateicība klusē, bet neatkarība ir pilsonisku sabiedrību galvenais likums. Lietuva, Volīnija vēlas Polijas karalisti, bet mēs vēlamies vienotu Krievijas impēriju. Kādu balsi vairāk jādzird Jūsu sirdij? Viņi kara, starp citu nemaz ne iespējama (jo kas tagad var sacelties pret Krieviju?), gadījumā var nodot mūs: tad sodīsim nodevību ar spēku un tiesībām: tiesībām vienmēr ir savs spēks, bet dumpis, kā nelikumība, atņem to. Poļi ar likumu apstiprināti īpašas, valstiskas tautas cieņā, mums ir bīstamāki kā poļi~--- Krievijas pavalstnieki.

Valdniek! Dievs Jums ir devis tādu spēku un tādu valsti, ka Jūs bez nepateicības, bez kristietības grēka un bez godkārības. ko nosoda pati cilvēciskākā politika, nedrīkstat vēlēties neko vairāk kā miera nostiprināšanu Eiropā un Krievijas labiekārtošanu: pirmo ar nesavtīgu, augstsirdīgu starpniecību; otro ar labiem likumiem un vēl labāku pārvaldi. Jūs jau esat ieguvis Lielā nosaukumu\footnote{Oficiālajā Krievijas impērijas historiogrāfijā Aleksandrs~I tika dēvēts par ``Svētīto'' (\rutxti{Благословенный}).}, iegūstiet vēl mūsu Tēva nosaukumu. Lai pastāv un pat zeļ Polijas karaliste, tāda kā tā tagad pastāv: taču lai pastāv, lai zeļ arī Krievija, kāda tā ir un kādu to Jums atstāja Katrīna! Katrīna mīlēja Jūs maigi, mīlēja arī mūsu tēvzemi: viņas ēna ir šeit klāt \dots{} apklustu.

\emph{Carskoje Selo. 1819.~gada 17~oktobrī.}


Avots: Н.~М.~Карамзин. О древней и новой России. Избранная проза и публицистика. М., ``Жизнь и мысль'', 2002, С. 437--438

Tulkots no: Н.~М.~Карамзин:~Мнение русского гражданина [Письмо Александру~I по поводу проекта восстановления Польши] // ~Русский обозреватель 13.10.2008~ // % http://www.rus-obr.ru/library/915 - 19.11.2009

% page 301


\subsection{Izraksts no Ezekela Kotika atmiņām par 1863.--1864.~gada poļu sacelšanos}

Un te, 1863.~gada sākumā Polijā un Lietuvā\footnote{Poļu sacelšanās~--- 1863.~gada janvāris~--- 1864.~gada maijs~--- aptvēra ne tikai Polijas Karalisti, bet arī Lietuvu, daļu Baltkrievijas un Labā krasta Ukrainas.} pēkšņi uzliesmoja sacelšanās. Tā gan bija nelaime ebrejiem! Muižnieki un šļahta bariem devās mežos ar ieročiem, sāka sagrābt mazpilsētas, kur nebija spēris kāju krievu karavīrs. Ienākot miestos viņi uzreiz norāva krievu ērgli no pagastvaldēm\footnote{Pagastvalde~--- pagasta pašpārvaldes iestāde XIX~gadsimtā. Tajā ietilpa pagasta vecākais un pagasta valdes locekļi, kurus ievēlēja pagasta pilnsapulce. Pagasts~--- vairākus ciemus aptverošs novads.}, izkarināja poļu ģerboni\footnote{Poļu ģerbonis~--- balta viengalvaina ērgļa attēls.} un pie tam kliedza ``\pltxti{Vzenta Rosia}''\footnote{\pltxti{Vzenta Rosia}~--- ieņemta Krievija.}. Pie tam muižnieki-revolucionāri apgājās ar ebrejiem ļoti ``savdabīgi''.

Pilsētās ebreji dzīvoja bailēs, viņi baidījās iziet uz ielas. Ebrejam, kurš gāja pa ielu un ieraudzīja Polijas karavīru, nemaz nerunājot par virsnieku (tie bija no muižniekiem), bija jānorauj cepure un jāklanās goddevīgi trijos līkumos~--- vēl jo vairāk, ja tas bija virsnieks. Te nu vajadzēja krist četrrāpus, liekties un klanīties.

Un, ja muižnieks bija neapmierināts ar to, kā klanās ebrejs, viņu nekavējoties satvēra un vilka aiz bārdās pie pulkveža uz tiesu. Te nu pār ebreju ņirgājās pilnā sparā: par viņa cices\footnote{Cices~--- ebreju apģērba sastāvdaļa, vilnas diegu bārkstis, kas nokarājās no žaketes. To uzdevums bija atgādināt ticīgajiem ebrejiem viņu reliģiskos pienākumus.}, par viņa bārdu un vaigubārdu, viņu raustīja, grūstīja, kratīja un ebrejam to visu vajadzēja paciest, kamēr pulkvedis nedeva pavēli atlaist ``žīdu'', kuram turpat uz vietas bija jāapzvēr sava uzticība poļu valdībai\footnote{Poļu valdība (poļu \pltxti{Rząd Narodowy})~--- galvenā sacēlušos varas iestāde, komiteja, kas vadīja sacelšanos.}.

Ilgi palikt pilsētiņā nemiernieki nevarēja, pēc dienām divām vai trijām ieradās krievu karavīri, kazaki ar artilēriju, un muižnieki bēga. Daudzi slēpās pie ebrejiem bēniņos un pagrabos, krāsns apakšā, kur audzē cāļus\footnote{Telpa zem krievu krāsns tika izmantota kā inkubators.}, un aizkrāsnē.

Bet slēpjoties, muižnieki baidījās, ka ebreji~--- ja karavīri viņus piespiedīs~--- var tos nodot. Un muižnieki ņēma ebreju ķīlnieku sev līdzi aizkrāsnē vai kādā midzenī bēniņos. Stāsta~--- un tā nav leģenda~--- ka muižnieks, paslēpies tādā midzenī kopā ar kā ķīlnieku paņemto ebreju, pieprasīja lai viņš noņemtu cepuri un slēptos ar kailu galvu: kā citādi, viņš taču slēpās kopā ar muižnieku! Pat ebreju jermolku\footnote{Jermolka vai kipa~--- tradicionāla ebreju vīriešu galvas sega~--- maza adīta, vai šūta cepurīte, kas nosedz pakausi.} ebrejs nedrīkstēja paturēt galvā.

Krievu karaspēks medīja muižniekus, pakarot un postot tos, bet piebeigt viņus uzreiz bija grūti. Tiklīdz krievu karaspēks atstāja pilsētu, atkal bija klāt jaunu muižnieku masas. Veselu gadu norisa cīņas starp krieviem un poļiem.

Zemniekus mocīja bailes un šausmas, ka tikai poļi nesagrābj Krieviju. Viņi baidījās, ka tad, ja tas notiktu, atkal tiktu ieviesta dzimtbūšana.\footnote{Atmiņu autors 1863.~gada Poļu sacelšanos uzskata nevis par nacionālās atbrīvošanās un demokrātisku kustību, bet gan vairāk par poļu muižnieku reakciju uz dzimtbūšanas atcelšanu Krievijā. Pie kam, pēc viņa domām, vietējie baltkrievu zemnieki pareizticīgajā cara administrācijā redzēja savu tiesību aizstāvi pret nīstajiem muižniekiem.} Tāpat viņi baidījās iziet ārā no ciema, lai tos nevarētu piespiest ar spēku piedalīties sacelšanās.

Bet poļi neņēma zemniekus [savās vienībās], jo neuzticējās tiem. Un, ja atsevišķus zemniekus aizturēja uz ceļa vai mežā, tad iesaistīja kādā palīgdienestā, bet nelika nest ieročus.

Bet ebrejiem klājās tūkstoš reižu sliktāk. Kad sagrāba ebreju, kurš apstaigāja ciemus [tirgojoties] vai devās kaut kur peļņā , tad viņš izjuta īstas nāves bailes. Pirmkārt, viņam teica, ka viņš tiks pakārts, pēc tam kārtīgi izņirgājās pār ``žīdu'', un kad nogura ņirgāties, ņēma virvi, lēnām aplika ebrejam ap kaklu un lika sūdzēt grēkus, tajā pašā laikā krītot pie zemes no smiekliem.

Ebrejs skaitīja pēdējo grēksūdzi\footnote{Atbilstoši jūdaisma reliģisko, ētisko un tiesisko likumu krājumam Talmūdam cilvēkam pirms nāves, piemēram, uz nāvi notiesātam vai smagi slimam, ir jāizsūdz grēki.}, ko pavadīja [sacelšanās dalībnieku] dzīvnieciski smiekli.

Tas vilkās briesmīgi ilgi, inkvizitoriski, nežēlīgi. Pēc divām stundām, kad ebrejs, izlējis asaru upes, pabeidza savu pirmsnāves grēksūdzi, bez spēka vairs raudāt, viņam noņēma virvi un teica:~--- Tu taču zini, žīd, ka mēs esam labi cilvēki. Tu domāji, ka mēs tevi pakārsim? To spēj krievi, bet nevis mēs, poļi. Zvēri, ka būsi mums uzticīgs. Un, ja satiksi ceļā krievu, nesaki, kur mēs atrodamies,~--- un vācies!

Mājup ebrejs (vai vairāki ebreji) ieradās ar nāvīgi bālu seju, tik--tikko dzīvs, līdz nāvei nobiedējot ar savu izskatu sievu un bērnus. Daži pat drīzumā mira pēc šāda ``labo'' poļu ``joka''.

Cīņa bija šāda: ja poļu bija trīs līdz četras reizes vairāk nekā krievu, tad notika smaga kauja, kuras rezultātā krievi, esot regulāra armija, sakāva poļus vai saņēma tos gūstā. Bet, ja pārākums poļiem bija mazāks, cīņa bija īsa. Poļi vai uzreiz padevās vai bēga kur kurais. Taču krievi bieži aptvēra visu mežu, aplenca to un pakāpeniski saņēma visus poļus gūstā.

Kādreiz pie Černjavčicas miesta\footnote{Baltkrievu \betxti{Чарнаўчыцы}~--- miests Grodņas guberņas Brestas apriņķī, mūsdienās Baltkrievijas Brestas apgabalā.}~--- tas ir netālu no mums~--- poļi kā parasti uzlika kādam ebrejam cilpu kaklā. Tai brīdī ieradās krievu karavīru bataljons\footnote{Bataljons~--- 19.~gadsimtā Krievijas armijas vienība, sastāvoša no četrām rotām, ar apmēram 900~kareivjiem.} ar kazaku eskadronu\footnote{Eskadrons~--- kavalērijas vienība. Krievijas armijā 19.~gadsimtā tajā bija ap 120--150~jātnieku.}. Poļu bija tikai daži simti, starp tiem~--- visvairāk ievērojamie muižnieki. Krievu tos aplenca un., tuvojoties ieraudzīja, ka poļi gatavojas pakārt ebreju. Ebreju no virvēm nekavējoties atbrīvoja un krievu pulkvedis, kas bija pazīstams ar poļu pulkvedi, viņam jautāja:

--- Kāpēc tu vēlējies pakārt ebrejus?

--- Mēs vēlējāmies pajokot,~--- atbildēja poļu pulkvedis. Ir skaidrs, ka krievu pulkvedis vēlējās tagad pajokot par poļu pulkvedi \citespace{}

Sākotnēji, kā jau tika teikts, zemnieki baidījās no poļiem: baidījās atstāt ciemus, lai poļi vardarbīgi nepiespiestu viņus karot. Bet beigās, kad karojošo poļu kļuva arvien mazāk, un viņi jau zaudēja drosmi, bet atgriezties mājās nevarēja, jo bija iekļauti [krievu] militārā priekšnieka sarakstā, kā tādi, kas atradās mežā, lūk tad zemnieki sāka atriebties muižniekiem.

Ieradās pie sava muižnieka naktī, izvilka viņu no gultas, sasēja, nopēra, bet pēc tam aizveda pie militārā priekšnieka un teica, ka viņi muižnieku atraduši mežā. Ko pēc tam darīja ar muižnieku~--- ir saprotams.

Tā zemnieki atriebās saviem saimniekiem. Daudzi muižnieku krita savu zemnieku, bijušo dzimtcilvēku rokās, un tie varētu darīt ar viņiem ko vēlējās. Tika iznīcināts milzīgs skaits muižnieku: pakārti, nošauti, nodurti, sakropļoti, izsūtīti uz Sibīriju un Dievs zina, kur.\footnote{Sacelšanās apspiešanas rezultātā tika sodīti ar nāvi ap 400~cilvēku, 2,5~tūkstošiem piesprieda katorgu, līdz 20~000 izsūtīja uz Krievijas iekšējām guberņām un Sibīriju.} Daudzi īpašumi tika konfiscēti, un krievu virsnieki pirka lielākos īpašumu par simtnieku. Par muižu, kura pēc tam ienesa trīs tūkstošus rubļu gadu nomas, krievi maksā nieka naudu vai vispār neko.

Poļu sakāve~--- viens no visbriesmīgākajiem pēdējo gadsimtu notikumiem.

Tulkots no avota: Ехезкел Котик. Мои воспоминания. Санкт-Петербург~--- Москва~--- Иерусалим, 2009, С. 218--222

\subsection{[Krievijas] Augstākā Virspavēlnieka uzsaukums poļiem}

Poļi, ir situsi stunda, kad jūsu tēvu un vectēvu senais sapnis var piepildīties.

Pirms pusotra gadsimta dzīvais Polijas ķermenis tika saraustīts gabalos, taču nenomira tās dvēsele. Tā dzīvoja ar cerību, ka nāks poļu tautas atdzimšanas diena, kad notiks tās brālīga samierināšanās ar dižo Krieviju. Krievu karaspēks nes jums šīs samierināšanās labo ziņu. Lai tiek nolīdzinātas robežas, kas sašķeļ daļās poļu tautu. Lai tā apvienojas zem Krievu Cara sceptera. Zem šī sceptera apvienosies Polija brīva savā ticībā, valodā, pašvaldībā.

Tikai vienu gaida no jums Krievija: tādu pat cieņu pret visu to tautību tiesībām, ar kurām jūs saistījusi vēsture. Ar atklātu sirdi, ar brālīgi pastieptu roku nāk pie jums dižā Krievija. Tā tic, ka nav ierūsējis zobens, kurš satrieca naidnieku pie Grīnvaldes. No Klusā okeāna līdz Ziemeļu jūrām virzās krievu armijas.

Jums paveras jaunas dzīves ausma. Lai atmirdz šai ausmā krusta zīme, tautu ciešanu un atdzimšanas simbols.

\emph{Augstākais Virspavēlnieks ģenerāl-adjutants Nikolajs.}

\emph{1914.~gada 1.~(14.)~augustā}

Avots: История России. ХХ век: 1894~--- 1939.~--- Москва: Астрель: АСТ, 2009, С.~351--352

Tulkots no: % http://www.hrono.ru/dokum/k_polyakam.html - 17.11.2009.

\subsection{Vācijas imperatora Vilhelma~II un Austroungārijas imperatora Franča Jozefa~I t.~s. ``Divu ķeizaru uzsaukums'' poļiem ar Polijas karalistes pasludināšanu 1916.~gada 5.~novembrī}

''Viņa Majestāte Vācijas ķeizars un viņa Apustuliskā Majestāte Austrijas ķeizars un Ungārijas karalis, vadoties no vēlmes vest Krievijas kundzībai atrautos novadus pretī laimīgākai nākotnei, ir vienojušies no šiem apgabaliem izveidot patstāvīgu valsti ar mantojamu monarhiju un konstitucionālu iekārtu. Precīzāka karalistes robežu noteikšana notiks nākotnē. Tā [karaliste] radīs atbalstu abās sabiedrotajās valstīs, kuras veicinās brīvu tās spēku attīstību. Mūsdienu lielajā karā tās armijā ir jāatdzimst pagājušo laiku poļu armijas slavenajām tradīcijām. Tās organizācija, apmācība un vadība tiks noregulēta, savstarpēji vienojoties.''

Tulkots no: Schmidt-Rösler Andrea. Polen. Vom Mittelalter bis zur Gegenwart. München, 1996. S.121.

\subsection{Petrogradas strādnieku un kareivju deputātu padomes vētījums poļu tautai 1917.~gada 14.~(27.)~martā}

Poļu tautai. Carisko režīmu, kas pusotru gadsimtu apspieda poļu tautu vienlaikus ar krieviem, gāza proletariāta un karaspēka apvienotie spēki.

Informējot poļu tautu par šo brīvības uzvaru pār Viskrievijas žandarmu, Petrogradas Strādnieku un kareivju deputātu padome paziņo, ka Krievijas demokrātija stāv uz tautu nacionāli-politiskās pašnoteikšanās atzīšanas pamata un pasludina, ka Polijai ir tiesības būt pilnīgi neatkarīgai valstiski-starptautiskajās attiecībās.

Sūtam savus sveicienus brālīgajai poļu tautai un vēlam tai panākumus gaidāmajā cīņā par demokrātiskas un republikāniskas sistēmas izveidošanu neatkarīgajā Polijā.

\emph{Petrogradas strādnieku un kareivju deputātu padome}

Tulkots no: Документы и материалы по истории советско-польских отношений. М., 1963. Т. I. С. 26.

\section{Ilustrācijas 1.~nodaļai}

% TODO Illustrations.
% TODO Автору: приложения в конец книги, все вместе
% TODO Автору: Иллюстрации в текст, на наиболее подходящие места

\chapter{Polijas Republika. 1918--1939}

\epigraph
{Cilvēki labprāt tic tam, kam vēlas ticēt.}
{Gajs Jūlijs Cēzars (\latxti{Gaius Iulius Caesar})}

\epigraph
{Patriotisms~--- pārliecība, ka tava dzimtene ir labāka par citām tāpēc, ka tieši tu tajā esi piedzimis.}
{Džordžs Bernards Šovs (\entxti{George Bernard Shaw})}

\epigraph
{Samaitāta tauta, iegūstot brīvību, var to saglabāt tikai ar milzu pūlēm.}
{Tits Līvijs (\latxti{Titus Livius})}

\epigraph
{Demokrātija nevar nostāties augstākā līmenī par to cilvēcisko materiālu, no kura veidoti tās vēlētāji.}
{Džordžs Bernards Šovs (\entxti{George Bernard Shaw})}

\epigraph
{Slikta demokrātija tomēr ir labāka par labu diktatūru.}
{Arveds Bergs}

\epigraph
{Liliputu valstī skatīties uz valsts galvu ir atļauts tikai caur palielināmo stiklu.}
{Staņislavs Ježijs Lecs (\pltxti{Stanisław Jerzy Lec)}}

\epigraph
{Vienādi briesmīgs ir kā valstiskais, tā reliģiozais fanātisms.}
{Vitālijs Ginzburgs (\rutxti{Виталий Лазаревич Гинзбург})}

\epigraph
{Agresīvais nacionālisms ir tikai cits vārds aprobežotībai.}
{Zenta Mauriņa}

\epigraph
{No visiem ismiem, kurus ir izgudrojusi cilvēce, manās acīs pats negodīgākais ir nacionālisms. Kā mazas, tā lielas tautas.}
{Vladlens Dozorcevs (\rutxti{Владлен Леонидович Дозорцев})}

\epigraph
{Patriotismu nosaka kauna pakāpe, ko cilvēks izjūt par noziegumiem, kuri veikti viņa tautas vārdā.}
{Ādams Mihņiks (\pltxti{Adam Michnik})}

\epigraph
{Neradi elku no savas valsts, jo visām nācijām, ieskaitot arī tavējo, pastāv vispārcilvēciskas vērtības.}
{Mihaels Martess (\detxti{Michael Martes})}

\epigraph
{Vāja paaudžu atmiņa stiprina leģendas.}
{Staņislavs Ježijs Lecs (\pltxti{Stanisław Jerzy Lec})}

\epigraph
{Pēc vēstures attīrīšanas no meliem ne vienmēr pāri paliek patiesība, reizēm~--- tikpat kā nekas.}
{Staņislavs Ježijs Lecs (\pltxti{Stanisław Jerzy Lec})}

\section{Valsts izveide. Cīņa par teritorijām. 1918--1921}

\epigraph
{Impērijas sabrukums nebūt nenozīmē, ka tagad būs labāk. Asīrija, Persija, Roma, Lielbritānija~--- milzums piemēru.}
{Mihails Vellers. (\rutxti{Михаил Веллер})}

\epigraph
{Tēvzemes mīlestība nepazīst svešas robežas.}
{Staņislavs Ježijs Lecs (\pltxti{Stanisław Jerzy Lec})}

\epigraph
{Mazu tautu nacionālisms ir izolācijas un pašapmierinātības izpausme. Lielu tautu nacionālisms ir imperiālistiska ekspansija.}
{Nikolajs Berdjajevs (\rutxti{Николай Бердяев})}

\epigraph
{Politika pasludina dižus principus, taču atzīst tikai stiprā tiesības.}
{Pjērs Buasts (\frtxti{Pierre Boiste})}

\epigraph
{Gandrīz katrs ļaunums notiek ar viltus ieganstu darīt labu.}
{Pjērs Buasts (\frtxti{Pierre Boiste})}

\epigraph
{Meistars ir tas, kurš spēj sevi ierobežot.}
{Johans Volfgangs fon Gēte (\detxti{Johann Wolfgang von Goethe})}

\epigraph
{Revolūcijas ir tādas slimības, kuru gaitā tūkstoši veiklu šarlatānu prot iegūt sev ievērojamu labumu.}
{Pjērs Buasts (\frtxti{Pierre Boiste})}

\epigraph
{Revolūciju ikdienā cilvēki, kuri knapi derīgi airēšanai, nokļūst pie stūres.}
{Pjērs Buasts (\frtxti{Pierre Boiste})}

\epigraph
{To, kas ir ar vardarbību iegūts, var tikai ar vardarbību paturēt.}
{Mahatma Gandhi (\entxti{Mohandas Karamchand Gandhi})}

% page 337


Milzu ietekmi uz Polijas likteni atstāja 1917.~gada oktobra notikumi Krievijā. Pēc Oktobra apvērsuma Viskrievijas Otrais strādnieku un kareivju deputātu padomju kongresā piedalījās arī PKunLSD un PPS-ļevicas grupu pārstāvji un Viskrievijas Centrālajā Izpildu komitejā tika ievēlēti arī vairāki poļi. Kongress pieņēma Dekrētu par mieru, kurš karojošās tautas aicināja to noslēgt bez aneksijām un kontribūcijām. Uzstājoties pret svešu zemju aneksiju, ar šo dekrētu arī boļševiki faktiski starptautiskajās tiesībās ieveda tautu pašnoteikšanās tiesības. (Tikai pēc diviem mēnešiem tās plaši proponēja ASV prezidents V.~Vilsons.) ``Ar aneksiju jeb svešu zemju sagrābšanu'' tika saprasta jebkuras tautības pievienošana stiprai valstij ``neatkarīgi no tā, kad šī pievienošana notikusi, neatkarīgi arī no tā, cik attīstīta vai atpalikusi ir nācija, ko varmācīgi pievieno vai varmācīgi tur attiecīgas valsts robežās''. Teiktais tika attiecināts arī uz poļu tautu.

Padomju Krievijas valdība bez kādiem iebildumiem atzina Polijas neatkarību. Tā uzskatīja, ka šī jautājuma noregulēšanai nav vajadzīgas nekādas formalitātes. Tam pamatā bija boļševiku cerības uz drīzu pasaules revolūciju, kuras ievadījums, pēc viņu domām, bija Oktobra apvērsums Krievijā. Poļu revolucionārs F.~Dzeržinskis II Padomju kongresā, norādot, ka Polija vienmēr atradīs kopēju valodu ar revolucionāro Krieviju, paziņoja: ``Mēs neuzstādām sevis atdalīšanas no revolucionārās Krievijas lozungu. Ar to mēs vienmēr sarunāsim. Mums būs viena brālīga tautu saime, bez strīdiem un ķildām''. Polijas karalistes un Lietuvas Sociāldemokrātijas Petrogradā izdotā avīze ``\pltxti{Trybuna}'' (``Tribīne'') 1917.~gada 28.~oktobrī (pēc jaunā stila~--- 10.~novembrī) rakstīja: ``Sociālisma triumfs novedīs pie neatkarīgas Polijas~--- brīvas starp brīvām un vienlīdzīgas starp vienlīdzīgām brāļu tautu saimē, kuras savstarpēji palīdz viena otrai.''

Mērenās Galīcijas un Cešinas Silēzijas poļu sociāldemokrātiskās partijas (\pltxti{Polska Partia Socjalno-Demokratyczna Galicji i Śląska Cieszyńskiego}) centrālais izdevums ``\pltxti{Naprzód}'' (``Uz priekšu'') 1917.~novembrī rakstīja: ``No visām tagadējām Krievijas politiskajām partijām tikai boļševiku partija, bez ierunām un patiesi stāv \citespace{} Polijas neatkarības pozīcijās. Ļeņins grib neatkarību Polijai bez kādiem ``noteikumiem'' no Krievijas puses''. Tā gan bija nepilnīga boļševiku pozīcijas interpretācija, jo boļševiki cerēja, ka Polijā tāpat kā Krievijā pie varas nāks revolucionāri-internacionālisti. Arī PPS pirmās atsauksmes uz Oktobra apvērsumu bija labvēlīgas, dažkārt pat pilnas entuziasma.

Kopā ar PKunLSD biedriem Padomju varas nodibināšanā un nostiprināšanā Krievijā piedalījās arī PPS~--- ļevicas biedri. Tiesa, tās vadītāji daudzkārt arī kritizēja krievu boļševiku partijas un padomju valdības realizēto politiku. Vārdus ``proletariāta diktatūra'', ``strādnieku kontrole'', ``sociālistiskās revolūcijas uzvara'' u.c. partijas izdevumi lika pēdiņās, jo uzskatīja tos par utopiskiem kapitālisma zemās attīstības un ekonomiskā sabrukuma Krievijā apstākļos. PPS--ļevicas konference Harkovā 1917.~gada decembrī raksturoja Oktobra apvērsumu kā zemnieciski-buržuāzisku revolūciju. Tika nolemts nevairīties no sadarbības ar boļševikiem, taču, tā kā pēc šo poļu sociālistu domām ``boļševiki novirzās no marksisma ceļa uz blankismu'' (fr. \frtxti{Le blanquisme}~--- Luija Ogista Blankī izveidota teorētiska koncepcija, kas postulēja sabiedrības sociālās evolūcijas uz komunismu paātrināšanu, izmantojot organizētas bruņotas sacelšanās iespēju), viņiem arī turpmāk bija jāuzsver sava patstāvība. Tomēr revolucionārie notikumi virzīja kreisos poļu sociālistus uz arvien ciešāku sadarbību ar boļševikiem.

PPS-frakcijas Krievijā nostāja bija pretrunīga. Tā atzina Krievijas revolūcijas šķirisko raksturu, to, ka tā norāda ceļu uz tautu nacionālo atbrīvošanos, politiskās iekārtas pārbūvi, uzskatīja, ka padomju varas sabiedriski-politiskie soļi ir virzīti strādnieku šķiras labā, ka poļu revolucionāriem ir izdevīgi, lai revolūcija Krievijā uzvarētu, aicināja savus biedrus saliedēties ap Viskrievijas padomju kongresu, poļu karavīrus atstāt kontrrevolucionāros formējumus un iestāties Sarkanajā armijā. Taču tāpat PPS-frakcija arī pēc Oktobra apvērsuma turpināja uzskatīt, ka Krievijas ``tumšās, mazkulturālās'' iedzīvotāju masas nesaprot, kas notiek valstī un boļševiki nespēs realizēt savus plānus. Politiskajos pārveidojumos Krievijā tā saskatīja vienas boļševiku partijas diktatūras nodibināšanos.

Gandrīz 1/5 daļa Polijas iedzīvotāju (vairāk nekā 3,5~miljoni) tolaik atradās revolucionārās Krievijas teritorijā. Viņu vidū diezgan liels bija poļu strādnieku īpatsvars, izbraukušu no Polijas pēc 1905.~gada revolūcijas un evakuētu no Polijas karalistes sakarā ar Vācijas uzbrukumu 1915.~gadā (ap 1,5 miljoniem). Daudz bija arī Krievijas armijā iesaukto karavīru (500--600~tūkstošu) un karagūstekņu-poļu no Vācijas un Austroungārijas armijām (ap 500~tūkstošu).

Daudzi poļi deva lielu ieguldījumu Oktobra apvērsumā, Padomju Krievijas, Sarkanās armijas, Viskrievijās Ārkārtējas komisijas (Čekas) u.tml. iestāžu izveidē un vadībā.

Pavisam ap 50 tūkstošu poļu bija Oktobra notikumu dalībnieki. Pēc 1967.~gadā Polijā izdotās ``Poļu~--- Oktobra revolūcijas dalībnieku grāmatas'' (\pltxti{Księga Polaków uczestników Rewolucji Październikowej 1917--1920. Biografie}) datiem katrs desmitais no viņiem bija fabriku un rūpnīcu komiteju, pulku, rotu u.c. kareivju komiteju loceklis; katrs trešais tieši piedalījās Oktobra bruņotajā apvērsumā Petrogradā, Maskavā u.c., sastāvēja Sarkanajā gvardē vai tautas milicijā; katrs ceturtais strādāja Padomēs, revolucionārajās komitejās; gandrīz katrs otrais sastāvēja Sarkanajā Armijā, citās bruņotajās vienībās. Tikai Petrogradas Kara-revolucionārās komitejas, kura tieši organizēja Oktobra apvērsumu, sastāvā ietilpa poļu revolucionāri F.~Dzeržinskis, J.~Unšlihts, F.~Gželščaks.

No apmēram 300~000 dažādu tautību internacionālistu, kuri vēlāk~--- Pilsoņu kara laikā karoja Sarkanās armijas un padomju varu atbalstošo partizānu vienību sastāvā, 1/3~jeb ap 100~000 bija poļi. (Poļu literatūrā figurēja arī apgalvojums, ka cīņā ``par Oktobra uzvaru'' piedalījušies ap 200~000 poļu.) To Sarkanās armijas komandieru vidū, kuri ieņēma ``ģenerāļu'' amatus, arī bija vairāki desmiti poļu.

Daudzi poļu revolucionāri ieņēma vadošus amatus boļševiku radītajās varas iestādēs. S.~Budzinskis bija Maskavas un Maskavas apgabala sociālo lietu tautas komisārs, V.~Figelskis~--- Turkestānas Autonomās sociālistiskās padomju republikas tautas komisāru padomes priekšsēdētājs, A.~Sluckis~--- Taurijas Republikas (\rutxti{Советская Социалистическая Республика Тавриды}) Tautas komisāru padomes priekšsēdētājs Krimā utt.

Pie KSFPR Tautību lietu komisariāta 1917.~gada novembrī tika radīts Poļu nacionālo lietu komisariāts ar uzdevumu izskatīt visas lietas~--- bēgļu, militārās u.c., saistītas ar poļu iedzīvotāju stāvokli Krievijā un Padomju valdības nacionālās politikas realizāciju pret poļu tautu.

Viens no komisariāta uzdevumiem bija poļu bēgļu un citu ``darbaļaužu'' atbrīvošana no ``kaitīgās poļu nacionālistu ietekmes''. Komisariātu vadīja Polijas karalistes un Lietuvas Sociāldemokrātijas un boļševiku partijas biedrs J.~Ļeščinskis--Ļenskis. Tika atklāts analfabētisma likvidēšanas skolu un dažādu kursu tīkls pieaugušajiem, daudzās Krievijas pilsētās poļu bērni skolās varēja mācīties dzimtajā valodā, komisariāts izdeva avīzi, sekmēja grāmatu iznākšanu poļu valodā. Pie Komisariāta tika izveidota Polijas revolucionāri demokrātisko organizāciju padome, kurā darbojās pārstāvji no PKunLSD, PPS-ļevicas un PPS-frakcijas.

Komisariāts uzsāka arī poļu revolucionāro karaspēka daļu formēšanu. Pirmais un plašāko slavu ieguvušais bija Belgorodas (\rutxti{Белгородский}) pulks, kura sastāvā vienu brīdi bija pat ap 17~000 karavīru. Šis pulks revolūcijas pusē pārgāja jau 1917.~gadā korņiloviādes laikā, pēc tam aizstāvēja jauno Padomju Krievijas valsti pret vācu uzbrukumu 1918.~gada sākumā. Martā tas iegāja Sarkanās armijas sastāvā kā 1.~Revolucionārās Varšavas sarkanais pulks (\rutxti{1 Красный полк революционной Варшавы}). 1918.~gada vasarā tas kļuva par internacionālās Rietumu divīzijas kodolu. Pie divīzijas tika izveidota pirmā poļu virsnieku skola. Poļu revolucionārās daļas 1918.~gada aprīlī Maskavā piedalījās anarhistu atbruņošanā, turpat t.s. ``kreiso eseru dumpja'' apspiešanā, cīņā pret labējiem eseriem Jaroslavļā 1918.~gada vasarā, vēlākajās Sarkanās armijas cīņās pret ``baltajiem''.

Padomju Krievijas Tautas Komisāru Padome 1918.~gada 19.~janvārī publicēja dekrētu par poļu tautai piederošā un kara laikā uz Iekškrieviju evakuētā vēstures un mākslas mantojuma: bibliotēku, arhīvu, muzeju eksponātu nodošanu valsts aizsardzībā Poļu nacionālo lietu komisariāta personā. (Pēc Pilsoņu un poļu-padomju kara Padomju Krievija šīs vērtības nodeva Polijai. To vidū bija izcili mākslas darbi.) Taču pēc tam, kad Padomju Krievija bija atzinusi Polijas neatkarību, nekādas valstiskas robežas netika nospraustas. Boļševiki Krievijā, kuru vidū bija arī ievērojami poļu izcelsmes revolucionāri~--- jau minētie F.~Dzeržinskis, J.~Marhļevskis u.c., ar ietekmi Polijas kreiso aprindās, loloja cerības uz pasaules revolūciju un arī poļu teritoriju revolucionizēšanu. Padomju Krievijas vadība uzskatīja, ka nav lietderīgi vest strīdus par dažādām teritorijām un to robežām pirms sociālistiskas Eiropas izveides. Tika cerēts, ka Eiropas un pasaules padomju valstu federācija atcels teritoriālās problēmas, kuras sociālisma apstākļos nedrīkstēja kļūt par tautu pretnostatīšanas faktoru. Taču vēlāk, pēc 1918.~gada Novembra revolūcijas Vācijā lielākā poļu daļa centās atgriezties dzimtenē. Daļa no repatriantiem bija iemantojuši simpātijas pret revolucionāro Krieviju, taču pārējie uz tiem skatījās kā uz poļu lietas nodevējiem, kuri gatavi sadarboties ar agrākajiem apspiedējiem.

Sākotnēji arī nacionālistiski noskaņotie poļi cerēja izmantot Oktobra apvērsumu savās interesēs. 1917./1918.~gada mijā uz Maskavu nosūtītais J.~Pilsudska pārstāvis T.~Holovko panāca pieņemšanu pie V.~Ļeņina un Ļ.~Trocka, piedāvāja tiem radīt KSFPR poļu karaspēka daļas, ar kuru palīdzību varētu atjaunot Polijas neatkarību. Taču šis priekšlikums boļševiku vadoņu vidū neatrada dzirdīgas ausis, jo Brestā norisa Padomju Krievijas sarunas ar Vāciju, kurās gāja runa par kara darbības pārtraukšanu. T.~Holovko atrada atbalstu tikai kreiso eseru aprindās, kuri vēlējās turpināt karu pret imperiālistisko Vāciju.

Arī jautājums par Polijas statusu pagaidām palika neatrisināts. Tas īpaši netika skatīts arī Krievijas un Centrālo valstu miera sarunās Brestļitovskā. Kad 1917.~gada oktobrī vara Petrogradā nonāca boļševiku rokās, viņi piedāvāja visām karojošajām valstīm noslēgt mieru bez aneksijām un kontribūcijām. 2.~(pēc jaunā stila~--- 15.)~decembrī tika noslēgts Padomju Krievijas pamiers ar Vāciju. Pēc dažām dienām Brestļitovskā, jeb īsāk~--- Brestā, sākās miera sarunas. 5.~(18.)~janvārī Vācijas delegācija pieprasīja atsavināt Krievijai ap 150~km$^{2}$ teritorijas, tai skaitā Poliju, Lietuvu, daļu Latvijas, plašus ukraiņu un baltkrievu apdzīvotus apgabalus. 1918.~gada 27.~janvarī (9.~februārī) Vācija un Austroungārija noslēdza līgumu ar Ukrainas Centrālo Radu, ar ko atzina Ukrainas suverenitāti. Lai nostiprinātu Centrālās Radas pozīcijas, līgums paredzēja pēc tā ratifikācijas Holmas (\pltxti{Chełm}) guberņas, kuras rietumdaļu apdzīvoja poļi, un dažu citu poļu zemju nodošanu Radai, ar to atkal realizējot daļēju Polijas dalīšanu, bet noraidīja Ukrainas pretenzijas uz Austroungārijai piederošajām teritorijām~--- Bukovinu un Austrumgalīciju. Pēdējai ar Ļvovu bija jākļūst par īpašu Austrijas ukraiņu provinci, kaut līdz tam tās dzīvi faktiski kontrolēja poļi.

No 10.~februāra ukraiņu pārstāvji piedalījās Brestas miera sarunās, kad Polija tajās nebija pārstāvēta. Polijas marionešu valdība 1918.~gada sākumā publicēja notu par Brestas sarunām, kurā paziņoja, ka poļu tauta neatzīs līgumu, ja tas tiks noslēgts bez tās pārstāvju dalības. Arī boļševiki sākotnēji pieprasīja tiesības piedalīties Polijas likteņa izlemšanā, taču Centrālās valstis prasību neņēma vērā. Vienīgi PKunLSD pārstāvji S.~Bobiņskis un K.~Radeks, kuri bija kopējās Viskrievijas delegācijas locekļu vidū, guva iespēju 25.~janvārī (7.~februārī) sarunās uzstāties ar deklarāciju, kurā norādīja, ka poļu tautas masas tajās nav pārstāvētas un to vārdā izvirzīja prasības par nacionālo spaidu likvidēšanu, policijisko šķēršļu atcelšanu starp poļu zemju trijām daļām, brīvu savas valsts iekārtas veidošanu, poļu zemju atbrīvošanu no okupantiem u.c. Kā liecina deklarācijas teksts, PKunLSD pārstāvji gāja pat tālāk par savas partijas uzskatiem nacionālajā un poļu jautājumā, faktiski nostājās poļu tautas pašnoteikšanās atzīšanas pozīcijās. Vācijas un Austroungārijas delegācijas atteicās apspriest poļu sociāldemokrātu prasības. Tomēr vēlāk, jau daudzus gadus pēc Brestas miera sarunām, 1933.~gadā toreiz Polijā valdošā t.s. ``sanācijas'' režīma apstākļos armijas Galvenā štāba vēstures biroja sagatavotajā nepublicējamā pārskatā par ``poļu jautājuma'' attīstību Brestas miera sarunu laikā kā ``vēsturiski svarīgs fakts'' bija uzsvērts tas, ka ``pirmoreiz pēc ilgiem gadiem polis piedalījās starptautiskās sarunās, uzstājoties poļu jautājumā. [Padomju Krievijas] Nacionālo lietu komisariāta loceklis Staņislavs Bobiņskis asā un noteiktā formā uzstājās pret okupantiem\dots{}''

Sarunās Ukrainas Centrālās Radas pārstāvis paziņoja, ka Padomju Krievijas vara neattiecas uz Ukrainas teritoriju un Ukraina patstāvīgi aizstāvēs savas intereses. Nākamajā dienā Vācija ultimatīvi pieprasīja, lai Krievija paraksta līgumu, pēc kura tā zaudētu Poliju, Baltiju, Ukrainu, daļu baltkrievu apdzīvoto teritoriju. Uz to Padomju Krievijas ārlietu tautas komisārs Ļ.~Trockis Padomju valdības vārdā deklarēja, ka tā atsakās parakstīt aneksionistisko miera līgumu un reizē paziņo par kara stāvokļa izbeigšanu un armijas demobilizāciju. Taču 16.~februārī Vācija pasludināja kara darbības atsākšanu un 18.~februārī uzsāka uzbrukumu frontē, strauji virzoties uz priekšu un vajājot kaujas nespējīgās Krievijas armijas daļas. Šādā situācijā Padomju Krievija 3.~martā Brestā parakstīja mieru ar vēl smagākiem noteikumiem, kā tas bija sākotnēji prasīts, zaudējot Ukrainu, Baltiju, poļu un arī pārsvarā baltkrievu iedzīvotāju apdzīvotās guberņas. Kaut līgumā ne vārda nebija teikts par Poliju, faktiski ar tā noslēgšanu KSFPR atteicās no tiesībām uz Poliju. Polijas karalistes un Lietuvas sociāldemokrātijas un PPS-ļevicas pārstāvji Krievijā protestēja pret boļševiku piekāpšanos un Vācijas noteikumu pieņemšanu. Tomēr, kad 23.~februārī boļševiku partijas CK izskatīja Vācijas izvirzītos miera noteikumus, F.~Dzeržinskis partijas vienības saglabāšanas interesēs nebalsoja pret to pieņemšanu. Arī KK(b)P VII kongresā 1918.~gada marta sākumā neviens no tajā pārstāvētajiem poļu sociāldemokrātiem pret Brestas mieru neuzstājās. Toties PPS-ļevica pēc Brestas miera noslēgšanas atsauca savu pārstāvi no Poļu nacionālo lietu komisariāta, drīz pēc tam tā paša iemesla dēļ no komisariāta aizgāja PPS-frakcijas pārstāvis. Poļu sociālisti boļševiku valdības noslēgto vienošanos ar imperiālistiskajām valdībām, neizdevīgo miera līguma noteikumu pieņemšanu uzlūkoja kā būtisku revolūcijas sakāvi. Līgums tika uztverts kā Padomju Krievijas tiesiska sankcija Vācijas imperiālistiem okupēt Polijas, Lietuvas, Baltkrievijas un Ukrainas teritorijas. Tikai 1918.~gada vidū, kad boļševiki apspieda pirmos bruņotos mēģinājumus gāzt viņu varu, poļu sociālistu kreisā spārna kritiskā attieksme pret Brestas mieru pakāpeniski sāka izzust. 1918.~gada novembrī PPS-ļevicas vadība Krievijā griezās pie KK(b)P CK ar iesniegumu, ka tās biedri padomju republiku teritorijā iestājas KK(b)P. Kad pašā Polijā 1918.~gada decembrī PKunLSD un PPS-ļevica apvienojās Polijas Komunistiskajā strādnieku partijā (\pltxti{Komunistyczna Partia Robotnicza Polski}) Krievijas teritorijā no PKunLSD un PPS-ļevicas sekcijām izveidojās apvienotas PKSP grupas. Interesanti, ka arī PPS-frakcijas sekciju konference 1918.~gada oktobrī Krievijā aicināja savus biedrus atbalstīt boļševikus, arī praktiskos jautājumos PPS-frakcijas biedru nostāja bija daudz revolucionārāka nekā PPS pozīcija Polijā.

Vācu un austroungāru okupētajā Baltkrievijā un Ukrainā atradās daudz poļu bēgļu. Viņu lielākā daļa centās atgriezties mājup. Uz dzimteni devās arī daudzi no izjukušās Krievijas armijas demobilizētie poļu karavīri. Taču vācu militārās iestādes piekrita bēgļus un citus repatriantus pieņemt tikai tad, ja tie sev līdzi vestu pārtiku, kas pietiktu to iztikai līdz jaunai ražai, faktiski pusgadam. No kara izmocītās Krievijas izvest tik daudz pārtikas praktiski nebija iespējams. Okupācijas varas iestādes nebija ieinteresētas gādāt transportu, pārtiku repatriantiem. Viņi rīkojās uz savu roku, atgriešanās notika gaužām lēni. Daudzi bija spiesti mitināties pie demarkācijas līnijas, ciešot milzu trūkumu.

Polijas sabiedrībā Centrālo valstu vienošanās ar Ukrainas Centrālo Radu un teritorijas sadale izsauca sašutumu. Reģentu Padome izdeva manifestu, kurā okupācijas varu līgumu ar Radu kvalificēja kā ``ceturto Polijas dalīšanu'' un atteicās to atzīt. Protestējot pret Holmas (\pltxti{Chełm)} atdošanu Ukrainai, atkāpās J.~Kuharževska vadītā valdība, protestus iesniedza poļu frakcijas Vācijas Reihstāgā, Prūsijas landtāgā un Austroungārijas Reihsrātā. Notika masu demonstrācijas Polijas lielākajās pilsētās. Pret okupantu rīcību protestēja kā pilsoniskās partijas, tā arī PPS. Arī Polijas karalistes un Lietuvas Sociāldemokrātijas Galvenā valde Polijā 12.~februārī izdeva skrejlapu, kurā nosodīja Holmas guberņas nodošanu Ukrainai, uzskatot to par jaunu Polijas dalīšanu. Skrejlapā bija teikts: ``Jauna Polijas dalīšana! Lūk, kāpēc tika lietas [poļu] leģionu asinis, lūk, kādēļ bija vajadzīgs Bezeleram pakļauts ``poļu karaspēks!''

Kad Vācija 1918.~gada 23.~martā atzina arī ``Lietuvas karalistes'' neatkarību, poļu nacionālistu nepatika pret Vāciju vēl pieauga. Jāpiezīmē, ka Rietumbaltkrievija, Rietumukraina vēsturiski bija bijušās Lietuvas lielkņazistes teritorijas, kuras gan kādreiz iegāja Žežpospolitas sastāvā, bet Lietuvai paziņojot par savu neatkarību~--- tas notika 1918.~gada 16.~februārī~--- un atsakoties no savienības ar Poliju, Polijas pretenzijām zuda arī šāds it kā ``vēsturisks'' pamats. Faktiski poļi vēlējās šīs teritorijas piesavināties tikai tāpēc, ka tur dzīvoja poļu mazākums, kuram gan piederēja ievērojami īpašumi. Pēc aculiecinieka (R.~Vingendorfa) stāstītā 1918.~gadā poļu naids pret vāciešiem Varšavā salīdzinājumā ar 1915.~gadu, kad tā atradās vēl Krievijas varā, nebija mazinājies, tikai kļuvis klusāks; vācu valodā runāt atteicās pat tie poļi, kuri to lieliski pārvaldīja.

Turpretī Vācijā auga sašutums par ``nekaunīgajiem'' poļiem ``kongresa Polijā'', kuri pieprasīja Ukrainas pievienošanu, par poļiem Vīnes Reihsrātā kuri prasīja visu Galīciju, par poļiem Vācijas Reihstāgā, kuri ar sava pārstāvja A.~Stuhela starpniecību pieprasīja Pozenes un visas Rietumprūsijas pievienošanu topošajai Polijas valstij.

Sarežģījās Austroungārijas stāvoklis. Tā kā tās delegāti Brestas sarunās bija apsveikuši Ukrainas Republikas izveidi un Holmas guberņas nodošanu tai, Austroungārijas izveidotās okupēto zemju ģenerālgubernatūras centrā Ļubļinā poļu iedzīvotāji publiski sadedzināja Vācijas un Austroungārijas imperatoru portretus un Austroungārijas karogu. Poļu tautības augstākie Austroungārijas ierēdņi (slepenpadomnieki) atteicās no savām pilnvarām, daudzi poļi atteicās no Austroungārijas piešķirtajiem ordeņiem. Poļi, kā konstatē angļu vēsturnieks A.~Dž.~P.~Teilors, ``beidzot novērsās no Habsburgiem. Tā bija katastrofa Austrijas valdībai, kas bez poļu atbalsta nespēja uzturēt pat kontroles šķitumu Reihsrātā. Tāpēc notika pasākumi, lai anulētu Holmas darījumu~--- Ukrainai tika izteikti jauni piedāvājumi, Polijai~--- jaunas piekāpšanās; taču it viss palika bez risinājuma: poļi bija atsvešināti, bet ukraiņi nebija kļuvuši par sabiedrotajiem''. Austroungārijas ierēdņi stādīja dažādus plānus, kā panākt Habsburgu impērijas interešu apmierināšanu, piemēram, paredzot, ka Polijas karalistei vajadzētu nekavējoties pasludināt savu neatkarību, poļu karaspēkam organizēt puču Minskā, proklamējot novada patstāvību. (Šajās norisēs, pēc austriešu viedokļa, īpaši noderīgs varētu būt J.~Pilsudskis.) Jauno valstisko veidojumu pievienošanu Habsburgu valstij tad varētu atlikt uz vēlāku laiku. Saprotams, ka šādi plāni tika būvēti uz smiltīm.

Jau 1918.~gada janvārī poļu muižnieki, dzīvojošie Ukrainā, lūdza palīdzību Reģentu padomei pret Ukrainas Radas pasludināto privātīpašuma atcelšanu uz saimniecībām, kuras netika apstrādātas ar savu personīgo darbaspēku. Tika uzsvērtas briesmas, kuras draudot poļiem t.s. labā krasta Ukrainā. Tās tika pamatotas ar apgalvojumu, ka ``atšķirībā no poļu sentimentalitātes, piekāpības un altruisma, ukraiņos atrodams briesmīgs egoisms, austrumnieciska viltība un aizdomīgums, kas [viņiem] nes reālu labumu arī politikā''.

Reģentu padome ar Vācijas starpniecību piedāvāja Padomju Krievijai nodibināt diplomātiskas attiecības, bet tā 16.~jūnijā no priekšlikuma atteicās, neuzskatot Reģentu padomi par poļu tautas, bet gan tikai par okupantu gribas izteicēju. 1918.~gada 29.~augustā saskaņā ar Brestas līgumu Padomju valdība izdeva dekrētu, ar kuru paziņoja: ``Visi līgumi un akti, ko noslēgusi bijušās Krievijas impērijas valdība ar Prūsijas Karalisti un Austroungārijas impēriju, kas skar Polijas sadali, tiek neatgriežami atcelti kā pretrunā esoši tautu pašnoteikšanās principam un krievu tautas revolucionārajai tiesību apziņai. Polijas tautai tiek atzītas neatņemamas tiesības uz patstāvību un vienību'' (pasvītrojumi mani~--- V.Š.). Tas bija veikls propagandisks trieciens Vācijas plāniem Polijā. Vācija un Austroungārija ir nedomāja labprātīgi atteikties no līgumiem, kuri tām varēja vēl noderēt strīdos par teritorijām (abas valstis to izdarīja tikai pēc sakāves, 1919.~gadā), toties Padomju Krievija demonstrēja poļu pašnoteikšanās tiesību atzīšanu, ar to noskaņojot poļus pret okupantiem. Ar šo Padomju Krievijas aktu tika radīts starptautiski tiesisks pamats Polijas valstiskuma atjaunošanai.

Tuvojoties Pirmā pasaules kara beigām, nu jau Padomju Krievijas valdība, vēloties iegūt operatīvāku informāciju par notikumiem Polijā, piedāvāja Reģentu padomei akreditēt poļu sociālistiskās kustības dalībnieku, vienu no poļu komunistu līderiem, toreiz Krievijas Ārlietu tautas komisariāta darbinieku, vēlāk Padomju Krievijas valstsvīru, Kominternes darbinieku J.~Marhļevski kā KSFPR diplomātisko pārstāvi Varšavā. Šoreiz neatbildēja Reģentu padome, kura baidījās no boļševiku ietekmes pieauguma Polijā.

Poļu sabiedrība šajā laikā arvien vairāk savas cerības saistīja ar Antanti. Zināmu iespaidu kā uz visu starptautisko sabiedrību, tā arī uz Rietumvalstu politiķiem bija atstājusi poļu sabiedrisko un kultūras darbinieku ilgstošā propaganda, ka Polijas vēsturiskā loma ir kalpot par barjeru starp rietumu civilizāciju no aziātiski-bizantiskajiem austrumiem, ka poļu kultūra ir devusi redzamu ieguldījumu pasaules kultūrā, ka ir izveidojies specifisks ``poļu raksturs'', kurš var kalpot par garantiju sekmīgai nacionālas demokrātiskas valsts attīstībai. Vēl Parīzes miera konferences priekšvakarā Bernē poļu un Parīzē franču valodā tika izdota enciklopēdija ``Polija'' (\frtxti{Pologne}. Paris, 1918) kurā bija akumulētas nosauktās tēzes.

Rietumeiropas valstīs un īpaši ASV XX~gadsimta sākumā plaši izplatījās nacionālu valstu izveides ideja. Pie šīs idejas piekritējiem piederēja ASV prezidents T.~Rūzvelts, kurš atzina, ka Eiropai jātiek uzbūvētai pēc nacionālā principa (\entxti{the principal of nationality}). Šo ideju jau I pasaules kara laikā atbalstīja arī toreizējais ASV prezidents V.~Vilsons, kurš paziņoja, ka tautai ir tiesības uz savas nacionālas valsts izveidi.

Taču lielāka ietekme uz Antantes valstu nostāju bija Oktobra apvērsumam Krievijā un tās separātajam mieram ar Vāciju. Šie notikumi galīgi atbrīvoja Rietumu sabiedrotos no saistībām Krievijas priekšā. Arī Polijas jautājumā tie tagad varēja, neatskatoties uz agrākajiem solījumiem Krievijai, atklāti aizstāvēt savas intereses.

1918.~gada 5.~janvārī par Polijas neatkarību publiski izteicās Lielbritānijas premjerministrs D.~Loids Džordžs, paziņojot, ka neatkarīga Polija, ``kura iekļautu visus poļu elementus, kuri vēlēsies kļūt par tās daļu'' ir nepieciešams Austrumeiropas stabilizācijas nosacījums. Pēc trim dienām~--- 8.~janvārī ASV prezidents V.~Vilsons proklamēja savu miera programmu~--- t.s. ``14 punktus'', kuros bija paredzēts: ``Jātiek radītai neatkarīgai Polijai, kurai jāietver visas teritorijas ar neapstrīdami poļu iedzīvotājiem, un kurai jānodrošina brīva un droša izeja pie jūras, bet politiskā un ekonomiskā neatkarība, kā arī teritoriālā integritāte jāgarantē starptautiskiem līgumiem''. (``\entxti{An independent Polish state should be erected which should include the territories inhabited by indisputably Polish populations, which should be assured a free and secure access to the sea, and whose political and economic independence and territorial integrity should be guaranteed by international covenant}''.) Tiesa, V.~Vilsona programmas 10.~punkts paredzēja plašu autonomiju Austroungārijas tautām (to varēja tulkot kā Austroungārijas integritātes saglabāšanu, atstājot tajā arī poļu apdzīvotos apgabalus). Pie tam šo tautu skaitā bija arī ukraiņi (rutēņi), kuri nebūt nevēlējās atrasties Polijas sastāvā, bet varētu prasīt vai nu neatkarību vai vismaz autonomiju Austroungārijā. Pēc būtības šie paziņojumi nesolīja pievienot nākamajai Polijas valstij Pozenes novadu, Dancigu un Silēziju. Tomēr, kad minētie dokumenti un īpaši V.~Vilsona ``14 punkti'' tika publicēti, tie, neraugoties uz savu pretrunīgumu, atviegloja arī poļu pārstāvniecības Varšavā cīņu par Polijas valsts radīšanu.

Bez tam 18.~februārī Lielbritānijas valdība nosūtīja Poļu Nacionālajai komitejai Parīzē notu ar paziņojumu, ka tā devusi norādījumu savam pārstāvim Kijevā uzstāties ar deklarāciju par Vācijas un Austroungārijas līguma ar Centrālo radu neatzīšanu. Līdzīgu paziņojumu sniedza arī Francijas valdība.

Antantes valstu uzstāšanās Polijas jautājumā pēc Brestas miera rādīja to vēlmi izmantot poļus pret Centrālajām valstīm, izolēt tos no Padomju Krievijas, taču, paredzot iespēju atjaunot Polijas valstiskumu, deva poļiem jaunas cerības. Tiesa, kā parādījis amerikāņu vēsturnieks Č.~Seimūrs, nedz Antantes valstis, nedz ASV līdz Pirmā pasaules kara beigām negatavojās sagraut Habsburgu monarhiju, baidoties no anarhijas izcelšanās tās teritorijās. Tomēr, kad Vācija, pēc Brestas miera noslēgšanas lai nostiprinātu savas pozīcijas austrumos, uzsāka aktīvu Vācijas kontrolē esošas ``Centrāleiropas'' (\detxti{Mitteleuropa}) izveidi, cenšoties arī pakļaut arvien straujāk spēkus zaudējošo Austroungāriju, bet apspiesto nāciju, tai skaitā arī poļu, politiķi, kas uzstājās par nacionālo neatkarību, intensificēja pūliņus nacionālu bruņotu spēku un savu varas iestāžu radīšanai, Antantes valstis izšķīrās par šo centienu atbalstīšanu.

1918.~gada 3.~jūnijā Lielbritānijas, Francijas un Itālijas valdību vadītāji Lielbritānijas premjera un ASV prezidenta proponēto neatkarīgas Polijas radīšanas principu atzina kopīgā Versaļas deklarācijā, paziņojot, ka ``Vienotas un neatkarīgas poļu valsts izveide ar brīvu izeju pie jūras ir viens no stingra un taisnīga miera un tiesiskās kārtības nosacījumiem Eiropā''. Tiesa, R.~Dmovskis bija neapmierināts un norādīja, ka Rietumvalstu Polijai paredzētā ``brīvā izeja pie jūras'' nenozīmē vis Baltijas jūras piekrastes daļas nodošanu Polijai, bet gan tikai ``brīvu izeju pie jūras svešā teritorijā''.

Gan Austrumos, gan Rietumos darbojās \strong{poļu militārie formējumi}.

Sākotnēji pēc Padomju valdības nodibināšanās Petrogradā ģenerāļa J.~Dovbora-Musņicka vadītais 1.~poļu korpuss formāli saglabāja neitralitāti (tādu nostāju 10.~(23.)~decembrī deklarēja arī vairāku fronšu poļu karavīru kongresu delegācijas), taču pēc būtības bija noskaņots pret boļševikiem. 1917.~gada 27.~oktobrī (9.~novembrī) J.~Dovbors-Musņickis ziņoja Krievijas armijas augstākā virspavēlnieka štābam, ka viņa vadītās daļas pretosies jebkuram mēģinājumam no personu puses, kuras nav ``tautas vai Pagaidu valdības pilnvarotas'', arestēt vai padzīt no amatiem pavēlniecības locekļus, jo tas ``draud ar Krievijas un reizē ar to arī Polijas bojā eju''. Drīzumā izpaudās pretrunas starp poļu korpusa vadību un vietējām Padomju iestādēm. Mogiļevas guberņas Zemnieku deputātu padomes Izpildu komiteja pieprasīja izvest korpusu no guberņas teritorijas. Padomju varas nozīmētais armijas virspavēlnieks N.~Kriļenko pavēlēja korpusā noturēt komisāru vēlēšanas, ko poļu virsnieki ir nedomāja pildīt. Padomju Krievijas Rietumu frontes pavēlnieks A.~Mjasņikovs savukārt, atsaucoties uz Minskas guberņas zemnieku sūdzībām par poļu karavīru iejaukšanos zemnieku veicamajā muižnieku zemju dalīšanā un patvaļīgajām pārtikas rekvizīcijām, pavēlēja tam mainīt dislokācijas rajonu. 1918.~gada 12.~(25.)~janvārī ģenerālis J.~Dovbors-Musņickis, atsaucoties uz neitralitāti Krievijas iekšējās lietās, atteicās pakļauties Padomju valdības rīkojumiem. Viņa komandētās daļas ieņēma Rogačevas pilsētu Baltkrievijā. Tad 21.~janvārī (3.~februārī) Padomju Krievijas Rietumu frontes pavēlnieks A.~Mjasņikovs izdeva pavēli par korpusa izformēšanu un tā karavīru demobilizāciju. J.~Dovbors-Musņickis atteicās pakļauties pavēlei. 31.~janvārī (12.~februārī) Jukuma Vācieša vadībā latviešu strēlnieku 1.~Daugavgrīvas un 4.~Vidzemes pulks kopā ar Baltijas jūras matrožu vienību ieņēma Rogačevu.

Šādos apstākļos februāra vidū notika radiogrammu apmaiņa starp J.~Dovboru-Musņicki un Vācijas virspavēlniecības štāba priekšnieku Austrumu frontē ģenerāli M.~Hofmanu. J.~Dovbors-Musņickis paziņoja, ka viņa korpuss par vienīgo augstāko varu atzīst Reģentu padomi Varšavā un lūdza palīdzību saistīties ar to. Sakari tika nodibināti un Reģentu padome faktiski sekmēja korpusa pakļaušanos Vācijas pavēlniecībai. Tika panākta vienošanās par poļu un vācu spēku sadarbību pret ``kopējo ienaidnieku~--- boļševismu''. Poļu korpuss darbojās kopā ar Vācijas karaspēku, kurš ar pilsonisko spēku radītās Baltkrievijas Radas (\betxti{Беларуская Народная Рaда}~---  baltkrievu nacionālo spēku izveidota pārstāvniecības iestāde) atbalstu ieņēma Minsku. Kad pēc Brestas miera līguma noslēgšanas (1918.~gada 3.~martā) visa Baltkrievijas teritorija nonāca Vācijas rokās, 1.~poļu korpuss ar ģenerāli J.~Dovboru-Musņicki priekšgalā pilnībā pakļāvās Vācijas pavēlniecībai, izpildot okupācijas karaspēka funkcijas Baltkrievijā.

Jau februāra beigās J.~Dovbors-Musņickis lūdza M.~Hofmanam atļauju sarīkot soda ekspedīciju pret nepaklausīgajiem baltkrievu zemniekiem. Vācu varas iestādes nodeva poļu militāristu pārvaldē piecus baltkrievu apriņķus. Tajos nekavējoties tika atjaunots muižnieku zemes īpašums, kas, kā teica poļu virsnieki, ``protams, nenotika bez upuriem''. Poļu daļas nevēlējās pāriet vācu armijas apgādē un ``pārtika uz šo piecu apriņķu rēķina''. Lai paturētu karavīrus paklausībā, tiem tika solīta drīza atgriešanās dzimtenē. Baltkrievu vēsturnieki mūsdienās runā par J.~Dovbors-Musņicka komandētā karaspēka ``zvērisku izrēķināšanos'' ar Bobruiskas, Igumenskas un Minskas apriņķu baltkrievu zemniekiem.

Maijā-jūlijā Vācijas pavēlniecība, baidoties no poļu nacionālās kustības pieauguma, veica poļu daļu atbruņošanu un izformēšanu, jau 20.~maijā pieprasīja 1.~poļu korpusa ``pilnīgu izformēšanu''. Reģentu padome pret to neprotestēja. Nākamajā dienā tika parakstīts arī atbilstošs poļu korpusa pavēlniecības līgums ar Vācijas armijas Austrumos pavēlniecību. J.~Dovbora-Musņicka komandētā 1.~poļu korpusa atliekas un pats ģenerālis devās uz Poliju. (1918.~gada beigās, 1919.~gadā viņš vēl piedalījās jaunās Polijas armijas veidošanā, taču vēlāk zaudēja konkurences cīņā par ietekmi armijā J.~Pilsudskim un bija spiests aiziet no dienesta.)

Viegli neklājās arī poļu karavīriem Austroungārijas armijā. Pēc 1917.~gada marta nemieriem neskarta palikusī poļu leģionu 2.~brigāde (Austroungārijas Poļu palīgkorpusa daļa) pulkveža J.~Hallera vadībā naktī no 1918.~gada 15.~uz 16.~februāri, pirms izvērsta Vācijas uzbrukuma Padomju Krievijai ar kauju pret austroungāru daļām pārgāja pāri frontes līnijai uz Padomju Ukrainas teritoriju, lai apvienotos ar poļu daļām Krievijas armijā. Pēc paša J.~Hallera atmiņām padomju iestādes draudzīgi sagaidīja negaidītos frontes pārgājējus, uzskatot tos par sabiedrotajiem. V.~Ļeņins, kuram tika paziņots par notikumu, ierosināja nodibināt ar tiem labas attiecības un sniegt palīdzību. Militāro delegāciju, ko J.~Hallers nosūtīja uz Maskavu, pieņēma VCIK priekšsēdētājs J.~Sverdlovs. J.~Hallera vadībā izveidojās 2.~poļu korpuss. Ukrainā un Baltkrievijā izvietoto poļu karavīru kopējais skaits sasniedza ap 30~tūkstošu vīru. Taču padomju varas iestāžu un poļu nacionālistu sadarbība nevarēja būt cieša. Kā atmiņās rakstīja pats J.~Hallers, viņa komandētie karavīri, aizstāvot kādas muižnieces īpašumu, sarīkoja soda ekspedīciju, nošāva zemnieku vadoņus un nodedzināja sādžu. Pēc 18.~februārī sāktā vācu uzbrukuma arī teritorijas, kur atradās poļu karavīri, nonāca Vācijas armijas rokās. Kad vācieši mēģināja poļu daļas atbruņot, lielāko pretošanos izrādīja tieši 2.~korpuss. Pēc sīvām cīņām korpuss tomēr kapitulēja. Tikai nelielai poļu karavīru daļai izdevās izlauzties caur vācu spēkiem uz padomju teritoriju un tā izbēgt gūstam. Līdz ar to 1918.~gada jūnijā beidza pastāvēt visi Centrālo valstu pusē 1914.~gadā radītie poļu formējumi. Pats J.~Hallers bēga uz Maskavu, šeit uzņēma sakarus ar krievu baltgvardiem, iesaistījās sazvērestībā pret boļševiku varu, pēc tās izgāšanās ar franču sanitāro vilcienu tika izvests uz Murmansku, no tās devās uz Franciju, kur ieradās 1918.~gada jūlijā.

Vācijas uzbrukums Padomju Krievijai paātrināja poļu daļu veidošanos Sarkanās armijas sastāvā. Jau minēta Belgorodas poļu rezerves pulks tika pārveide par Sarkanās Varšavas revolucionāro pulku. 2.~augustā tā izvadīšanā uz fronti piedalījās arī V.~Ļeņins. Veidojās arī citas poļu daļas.

Padomju Krievijas aizmugurē (Čerņigovā, Harkovā, Kijevā, Maskavā, Odesā, Petrogradā u.c.) darbojās J.~Pilsudska radītā Poļu kara organizācija (\pltxti{Polska Organizacja Wojskowa}), kura formēja poļu vienības cīņai par Polijas interesēm, orientējās uz sadarbību ar Antanti un pretboļševistiskajiem spēkiem, sūtīja poļu brīvprātīgos uz Murmansku un Arhangeļsku, kur martā un augustā izcēlās britu desants. Antantes pārstāvji šos brīvprātīgos izmantoja cīņai pret boļševikiem. Daļa Krievijā esošo poļu karavīru iesaistījās krievu ``balto'' karā pret Padomju varu, taču daļa atteicās iejaukties Krievijas iekšējās lietās.

Kā iepriekšējā nodaļā minēts, Francijā jau 1917.~gadā pēc prezidenta R.~Puankarē un Poļu nacionālās komitejas pārstāvja I.~Paderevska vienošanās tika veidota jauna poļu armija. 1918.~gada sākumā armijas politiskā vadība pārgāja Poļu nacionālās komitejas rokās. Pēc komitejas lūguma Francijas valdība 1918.~gada 28.~septembrī to atzina par ``sabiedroto un karu vedošu armiju''. Tajā iestājās arī bēgļi no Krievijas. Pēdējo vidū bija arī nu jau ģenerālis J.~Hallers, kurš 1918.~gada jūlijā nonāca Parīzē un uzņēmās armijas komandēšanu. Tā paguva piedalīties arī cīņās pret Vāciju. Pirmā pasaules kara beigās šai t.s. ``Hallera'' vai ``zilajā armijā'' (\pltxti{Błękitna Armia}, tā nosaukta pēc formas tērpu krāsas) bija jau 6 divīzijas. (Līdz 1919.~gada maijam armijā sapulcējās ap 70~000 poļu.) Ģenerālis J.~Hallers 1918.~gada 6.~oktobrī pavēlē sakarā ar savu stāšanos visu poļu spēku, kuri karoja pret Centrālo valstu bloku, virspavēlnieka amatā kā uzdevumu tiem izvirzīja ne tikai Varšavas un Krakovas, bet arī Viļņas un Ļvovas ``atbrīvošanu''.

Vācijas līgums ar Ukrainu, Brestas miers, kā arī Austroungārijas paziņojums par patstāvīgas ukraiņu provinces izveidi Austrumgalīcijā, kur dzīvoja ap 3,2~miljoniem ukraiņu un ap 2~miljoniem poļu, veda pie straujas neapmierinātības augšanas poļu vidū. No okupācijas varām atkarīgā Reģentu padome vēl joprojām meklēja Polijas valstij izdevīgāko variantu sadarbībā ar Vāciju un Austroungāriju. 1918.~gada 19.~augustā no Reģentu padomes un Ministru padomes sastāvošā Troņa padome sprieda, ka Polijai pieņemamāks būtu t.s. austro-poliskais risinājums, jo tas dotu perspektīvu pārējo poļu apdzīvotās teritorijās apvienot ar Galīciju, austrieši arī labāk izturējās pret poļiem. Turpretī Vācijas vadītāji apmaiņā pret poļu atteikšanos no šī varianta solīja tiem savu atsacīšanos no vācu-poļu robežas caurskatīšanas, Holmas novada atgriešanu Polijai, kā arī ekspansiju uz austrumiem~--- Belostokas un Grodņas apgabalu pievienošanu Polijai, Dancigas ostas ``atvēršanu'' Polijas tirdzniecībai. Konkrētu izšķiršanos Troņa padome tā arī neizdarīja, nolemjot turpināt mēģinājumus izkaulēt no abām okupācijas varām Polijai labvēlīgākas robežas. Arī Berlīnē un Vīnē notika drudžainas sanāksmes Polijas jautājumā, taču tās beidzās bez rezultātiem.

1918.~gada 23.~septembrī Bulgārijas un Antantes valstu noslēgtā pamiera rezultātā sabruka Centrālo valstu savienība. Drastiski pasliktinājās Vācijas un Austroungārijas stāvoklis. Naktī uz 4.~oktobri Vācijas valdība nosūtīja ASV prezidentam V.~Vilsonam notu, kurā bija izteikta gatavība vest miera sarunas uz viņa piedāvāto “14 punktu'' pamata. Austroungārijas ķeizars Kārlis I 16.~oktobra manifestā proklamēja dubultmonarhijas pārvēršanu par ``brīvu tautu savienību'', faktiski~--- tās federalizāciju. Šādu plānu realizācija skartu arī poļu apdzīvotās teritorijas. Polijas ``austriešu'' daļai Austroungārijas federatīvas valsts izveide nesolīja apvienošanos ar pārējo Poliju.

Taču Reģentu padome Varšavā sakarā ar paredzamo Centrālo valstu sakāvi jau 7.~oktobrī paziņoja par Troņa padomes atlaišanu, pasludināja plānu radīt ``vienotu neatkarīgu Poliju'' atbilstoši 3.~jūnija Antantes valstu Versaļas deklarācijai, solīja paplašināt valdības sastāvu ar plašu tautas slāņu pārstāvjiem, noturēt Likumdošanas Seima vēlēšanas, bet 12.~oktobrī aicināja poļu karaspēka daļas aizstāvēt ``tēvzemi, poļu valsti un Reģentu padomi''. Ģenerālis H.~f.~Bezelers bija spiests uzsākt sarunas par varas nodošanu. 23.~oktobrī Reģentu padome paziņoja par jaunas daudzpartiju valdības izveidi. Par jauno valdības vadītāju Padome iecēla nacionāldemokrātu J.~Svežinski, kurš 3.novembrī izveidoja pirmo Polijas nacionālo valdību. Tajā kā kara ministram bija jādarbojas arī J.~Pilsudskim, kurš vēl atradās ieslodzījumā. Kā ārlietu ministrs valdībā iegāja ievērojams nacionāldemokrāts S.~Globinskis, saistīts ar Poļu nacionālo komiteju Parīzē, kam vajadzēja simbolizēt bijušo ``aktīvistu'' un ``pasīvistu'' spēku apvienošanos valsts izveides gaitā. Taču vienotību neizdevās panākt un jau 4.~novembrī Reģentu padome atlaida valdību.

27.~oktobrī Reģentu padome pasludināja Polijas armijas izveidi, iekļaujot tajā visus poļu militāros formējumus. Tikko izveidotās armijas Ģenerālštāba priekšnieks (\pltxti{Szef Sztabu Generalnego Wojska Polskiego}) ģenerālis T.~Rozvodovskis izdeva pavēli visiem Austroungārijas armijā dienošajiem poļu karavīriem no tās izstāties un iestāties poļu vienībās. Savukārt Austrijas Reihsrāta poļu deputāti 9.~oktobrī pieņēma lēmumu par bijušo Žečpospolitas zemju, ieskaitot Galīciju, apvienošanu Polijas valstī. Mēneša beigās sākās Austroungārijas sabrukums. Oktobra beigās darbu arī uzsāka poļu komisija austriešu varas likvidēšanai poļu zemēs. To vadīja ļudovietis autodidakts ar darba pieredzi Austrijas parlamentā V.~Vitoss. Komisijas pienākums bija veikt pasākumus, kas nodrošinātu Galīcijas apvienošanu ar citām poļu zemēm neatkarīgas Polijas valsts ietvaros. Poļu rokās nonāca vara Austrumgalīcijā, Cešinas Silēzijā, Ļubļinas ģenerālgubernatūrā. Kā konstatējis britu vēsturnieks A.~Dž.~P.~Teilors, ``visu karu stāvot ar vienu kāju katrā no pretinieku nometnēm'', tagad poļi ``kļuva par Antantes piekritējiem, paturēdami priekšrocības, ko deva sadarbība ar Habsburgiem \citespace{}. Būtībā poļi bija Brestļitovskas līguma mantinieki, vienlaikus iegūstot Pozenes apgabalu un ``Polijas koridoru'' kā Vācijas sakāves rezultātus''.

Neapšaubāmi, ka Polijas valstiskums varēja atdzimt pateicoties ilgstošai tautas cīņai par neatkarību, taču jāuzsver, ka panākumu nepieciešams nosacījums bija tāda spēku samēra izveidošanās starptautiskajā arēnā, kad dažādu valstu imperiālisti nevarēja traucēt apvienotas Polijas valsts izveidei. Tam nepietika viena no imperiālistisko valstu grupējumu sakāves, jo Poliju sadalījušās valstis taču atradās kara stāvoklī. Ja Pirmajā pasaules karā būtu uzvarējušas Centrālās valstis, poļiem labākajā gadījumā varētu rasties neliels, no agrākās Polijas karalistes teritorijas sastāvošs, no Vācijas un, iespējams, arī Austroungārijas, atkarīgs valstisks veidojums. Viennozīmīga Antantes valstu uzvara Polijas jautājumu saistīja ar Krievijas pozīciju, tāpēc nestu poļiem uz Vācijas un Austroungārijas zemju rēķina paplašinātas Polijas plašāku vai šaurāku autonomiju Krievijas sastāvā. Februāra revolūcijas uzvara radīja iespēju radīt no Krievijas atkarīgu Polijas valsti. Oktobra apvērsums Krievijā, Antantes uzvara karā bez Krievijas pavēra iespēju neatkarīgas Polijas valsts atjaunošanai.

Notikumu gaitas rezultātā, taču arī pateicoties R.~Dmovska un viņa piekritēju pūliņiem, Rietumu valstis, no kuru nostājas tagad bija atkarīgs Eiropas liktenis, arī poļu jautājuma risinājums, bija piekritušas vienotas, neatkarīgas Polijas valsts izveidei. Taču nacionāldemokrātu līderis jau vairākus gadus dzīvoja ārpus Polijas, atradās Francijā, viņa sakari ar dzimteni bija ierobežoti. Nacionāldemokrātu partija nepiedalījās okupantu radīto iestāžu darbībā, tai nebija ietekmes poļu militārajos formējumos dzimtenē. Sabrūkot Centrālo valstu militārajai varai, tā nespēja ātri radīt Polijā valsts institūcijas, kas pārņemtu varu. R.~Dmovska mēģinājumi panākt J.~Hallera komandētās armijas ātru pārdislocēšanu uz Poliju nedeva panākumus.

No partijām Polijā pirmo vijoli spēlēja PPS (\pltxti{Polska Partia Socijalistyczna}) un PSL (\pltxti{Polskie Stronnitctwo Ludowe ``Wyzwolenie''}). PPS popularitāte sevišķi auga strādnieku vidū, ļudoviešu~--- zemnieku vidū. Kamēr nacionāldemokrāti jeb endeki cerēja, ka kādā brīdī Reģentu padome, ievērojot Rietumvalstu atbalstu Poļu Nacionālajai komitejai Parīzē, pati tiem piedāvās ņemt varu savās rokās un nodos arī tolaik vēl no 5~000 līdz 9~000 vīru lielos poļu bruņotos spēkus, šīs partijas izrādījās visvairāk sagatavojušās stāvoklim, kad poļu teritorijās vairs nebija spēku, kas varētu traucēt varas pāreju to rokās.

Pilsoniskās partijas Galīcijā pēc īslaicīga apjukuma, ko izsauca Vācijas un Austroungārijas sakāve karā, samērā ātri pārorientējās uz Antanti un autonomijas prasības vietā izvirzīja prasību atjaunot Polijas valstiskumu. Krakovā un citās pilsētās kā sēnes pēc lietus izauga dažādas pilsoniskās komitejas un komisijas, kuras centās iesaistīt tautas masas cīņā par Polijas valstiskuma atjaunošanu.

Stāvokli sarežģīja kreiso un labējo politisko spēku savstarpējās cīņas. Novembra sākumā Polijas pilsētās (Ļubļinā, Varšavā, Lodzā u.c.) tika dibinātas Strādnieku deputātu padomes (to vēlēšanās piedalījās vairāk nekā pusmiljons cilvēku), bet konservatīvajai Reģentu padomei un arī Strādnieku padomēm opozīcijā stāvošie sociālisti un ļudovieši (no \pltxti{Polskie Stronictwo Ludowe ``Piast''} un \pltxti{Polskie Stronnitctwo Ludowe ``Wyzwolenje''}) 7.~novembrī Ļubļinā izveidoja savu Tautas valdību ar Austroungārijas poļu sociālistu līderi, V.~Vitosa vadītās likvidācijas komisijas locekli I.~Dašinski priekšgalā. I.~Dašinskis ilgus gadus bija sadarbojies ar J.~Pilsudski. Valdībā ietilpa arī J.~Pilsudska piekritēji E.~Moračevskis, T.~Arciševskis, par kara ministru šai valdībā kļuva Poļu kara organizācijas vadītājs E.~Ridzs-Smiglijs.

Valdība Ļubļinā pasludināja Reģentu padomes varu par gāztu un diezgan radikālu reformu programmu: Polijas Tautas republikas, pašvaldību un civilmilicijas izveidi (pastāvēja līdz 1919.~gada vidum), pilsoniskās brīvības, 8~stundu darba dienu, solīja iesniegt nākamajam Satversmes Seimam izskatīšanai likumprojektus par agrārās reformas īstenošanu, dažu rūpniecības nozaru nacionalizāciju un strādnieku dalību privātuzņēmumu pārvaldē. Šī programma nobiedēja mērenās un labējās partijas. Radās reālas briesmas, ka, vēl nepaguvusi konstituēties, valsts jau varētu tikt ierauta dziļā konfliktā, iespējams, pat pilsoņu karā. Šīs valdības radikālisms varētu negatīvi ietekmēt karā uzvarējušo Rietumvalstu attieksmi pret Poliju. Taču Ļubļinas valdība pastāvēja tikai pusotru nedēļu.

Tā 1918.~gada rudenī Parīzē, Varšavā un Ļubļinā vienlaikus pastāvēja trīs dažādi valsts varas aizmetņi ar dažādu partijiski politisko sastāvu. Šādā situācijā Reģentu padome centās panākt, lai Vācija atbrīvotu J.~Pilsudski, saskatot viņa personā visautoritatīvāko spēku, kurš varētu uzņemties atbildību par poļu tālāko likteni.

Tuvojoties neizbēgamajai kapitulācijai, par vienu no Vācijas valdošo aprindu uzdevumiem kļuva izdarīt visu iespējamo lai Polijas valsts atjaunošanas gaitā vismaz daļēji saglabātu savu ietekmi poļu zemēs, panāktu tādas varas izveidošanos tajās, kas neapdraudētu Vācijas austrumu robežas. Galvenais tuvākais Vācijas valdības uzdevums bija nodrošināt sava karaspēka pēc iespējas netraucētu evakuāciju no austrumos okupētajām, revolucionāro cīņu aptvertajām teritorijām. Lai to izdarītu vajadzēja panākt, lai netiktu radīti šķēršļi karaspēka evakuācijai caur poļu apdzīvotajām teritorijām. (Nav izslēgts, ka vācieši baidījās, ka poļu zemēs vara varētu nonākt nacionāldemokrātu rokās, kuri neslēpa savu naidīgumu pret ``ģermāņiem''.) Piemērotākā kandidatūra šim mērķim vācu acīs bija J.~Pilsudskis.

Vācijas valdība 31.~oktobrī nosūtīja savu pārstāvi grāfu H.~Kesleru uz Magdeburgu vest sarunas ar J.~Pilsudski. Viņš apgalvoja, ka poļi nekaros pret Vāciju par Pozeni, ka boļševiki ir kā Vācijas, tā Polijas ienaidnieki, apsolīja pēc dažiem mēnešiem radīt armiju, kura spēs apspiest draudošo revolūciju. Pēc atbrīvošanas no Magdeburgas cietokšņa 6.~novembrī H.~Keslers nogādāja J.~Pilsudski un K.~Sosnkovski Berlīnē, kur viņi veda sarunas ar Vācijas augstākajām amatpersonām, bet 10.~novembrī ar speciāli piešķirtu vilcienu Varšavā. Vācieši arī parūpējās, lai vilciena ierašanās laiks nebūtu plaši zināms, tāpēc J.~Pilsudski sagaidīja tikai Reģentu padomes loceklis kņazs Z.~Ļubomirskis un vēl daži cilvēki.

J.~Pilsudskis jau Varšavas stacijā pavēlēja Poļu Kara organizācijas biedriem atbruņot Varšavā esošo vācu garnizonu. Vācieši virknē gadījumu izrādīja pretestību, notika bruņotas sadursmes ar cilvēku upuriem. Taču Vācijā notiekošas revolūcijas ietvaros arī Varšavas vācu garnizonā tika radītas karavīru padomes, kuras nosūtīja pie J.~Pilsudska savus pārstāvjus ar piedāvājumu nodot ieročus un apsargājamos objektus pret netraucētu viņu izbraukšanu uz dzimteni. Priekšlikums tika pieņemts. Tā vara bez īpašiem traucējumiem nonāca poļu rokās.

Sakarā ar \strong{J.~Pilsudska} kā Polijas neatkarības cīnītāja, bijušā sociālista, poļu leģionu vadītāja lielo autoritāti 1918.~gada 11.~novembrī Polijas \strong{Reģentu padome iecēla viņu par militāro pavēlnieku}, bet jau 14.~novembrī pašlikvidējoties \strong{nodeva viņam arī politisko varu}. Tā \strong{123~gadus pēc valstiskuma zaudēšanas Polija to atkal atguva}. Interesanti, ka par Polijas (Otrās Žečpospolitas) dibināšanas (neatkarības atjaunošanas) dienu 1926.~gadā tika pasludināts 11.~novembris, kad J.~Pilsudskis kļuva par armijas virspavēlnieku, bet nevis 14.~novembris~--- datums, kad viņš kļuva par oficiālu varas turētāju.

Tiesa, kā raksta mūsdienu poļu vēsturnieks K.~Pilavskis, Otrajā Polijas Republikā 11.~novembra svinības saistījās ar Pirmā pasaules kara beigām, bet J.~Pilsudskis tikai izdarīja meistarīgu soli, pieņemot kara noslēguma datumu arī par Polijas neatkarības atjaunošanas dienu. Toties vācu autors R.~Vingendorfs šai sakarā rakstīja, ka tas, ka Polijas valsts 11.~novembri, kad Varšavā ``cilvēki, kuri vēl nesa Vācijas kareivja goda nosaukumu, nometa vai bez pretestības atdeva ieročus no pagrīdes iznākušajai Poļu kara organizācijai'', svin kā dibināšanas dienu, paliks kā rūgts negoda traips Vācijas vēsturē.

Bijušais Reģentu padomes loceklis kņazs Z.~Ļubomirskis atmiņās rakstīja, ka padomei bija jāatdod vara J.~Pilsudskim, jo leģionāru vadītājs simbolizēja neatkarību. ``Toreiz visas partijas, no galēji labējām līdz kreisajām, pieprasīja, lai mēs nodotu varu Pilsudskim.'' Sabiedrības priekšstati par ``stipras rokas'' nepieciešamību izejai no krīzes darbojās viņa labā. Viņa piekritēji kara laikā bija nemitīgi centušies celt sava līdera popularitāti un ietekmi sabiedrībā. Viņu atbalstošā Poļu kara organizācija 1918.~gada rudenī apvienoja jau 15~tūkstošus (citi dati~--- 25~tūkstošus) aktīvu, entuziasma pilnu biedru. Tautiešu acīs J.~Pilsudskim izdevās iegūt demokrātijas un likumības piekritēja, ne labēja, ne kreisa, bet visas tautas vārdā runājoša politiķa veidolu, kurš cenšas apvienot visu nāciju valsts izveides vārdā, sasaukt Satversmes Seimu, kas izteiktu tautas gribu. Jāuzsver arī, ka valsts dibināšanu ievērojami atviegloja tas, ka J.~Pilsudska rīcībā tika nodots jau gatavs līdz tam Reģentu padomei padots pārvaldes aparāts.

Tiesa, krievu vēsturniece M.~Bobrika pamatoti atzīmējusi: ``Ja neatsakās no alternatīvas pieejas vēstures izpētei, tad var pilnīgi noteikti teikt, ka, ja nebūtu ne Pilsudska, ne viņa leģionāru, arī tad sapnis par neatkarīgu Poliju kļūtu par realitāti, ko apstiprina pati notikumu gaita poļu zemēs, kuru atbrīvošana no vācu okupantiem sākās rudenī. Tai laikā pats Pilsudskis vēl bija ``Magdeburgas ieslodzītais'', bet viņa leģionāru vairākums bija internēti. Taču pilnībā noliegt viņu ietekmi uz poļu jautājuma risinājumu mēs nevaram''. Tomēr jāuzsver, ka J.~Pilsudska un viņa piekritēju darbība Polijai izšķirošajā posmā daudzkārt ieguva noteicošo lomu. Arī bez viņa Polija iegūtu neatkarību, taču tā būtu cita Polija.

J.~Pilsudskis nācijai svarīgajā brīdī uzskatīja par nepieciešamu realizēt kursu uz tās politiskās elites konsolidāciju. Gan valstij veidojoties, gan līdz pat tās demokrātiskās attīstības beigām viņš pat necentās radīt savu politisku partiju. Acīmredzot tas tika darīts, lai neapgrūtinātu dialogu ar dažādiem politiskiem spēkiem visas valsts, bet nevis kādas atsevišķas politiskas vai sociālas grupas, interesēs.

Lai izlīdzinātu pretrunas ar valdību Ļubļinā, J.~Pilsudskis sākotnēji par ministru prezidentu iecēla I.~Dašinski (ar to valdība Ļubļinā beidza pastāvēt), bet tai nebija labējo atbalsta (I.~Dašinska sarunas ar nacionāldemokrātiem par koalīcijas valdības izveidi beidzās bez rezultātiem) un jau 18.~novembrī, sastādot Varšavā nākamo Pagaidu valdību, viņš nomainīja I.~Dašinski ar mērenāku sociālistu un bijušo leģionāru E.~Moračevski, kuram bija zināmi kontakti ar nacionāldemokrātiem. Tomēr šī valdība, tāpat kā iepriekšējā, sastāvēja tikai no labējo sociālistu un ļudoviešu partiju pārstāvjiem, par kuru lojalitāti sev J.~Pilsudskis varēja būt drošs. Pats viņš tajā sākotnēji izpildīja kara ministra pienākumus.

Vienā no 22.~novembrī izlaistiem uzsaukumiem J.~Pilsudskis sevi nosauca par \strong{Valsts pagaidu priekšnieku} (\pltxti{Tumczasowy Naczelnik Panstwa}), kas atbilda T.~Kostjuško laika tradīcijām. Šo nosaukumu viņš nesa līdz parlamenta (Seima) sasaukšanai. Tai pašā 22.~novembrī ar kopēju valsts priekšnieka un valdības vadītāja dekrētu tika noteiktas savstarpējās kompetences. Visa augstākā likumdošanas un izpildvara atradās J.~Pilsudska kā valsts priekšnieka rokās. Tikai viņš varēja iecelt un atlaist valdību, kura arī bija atbildīga viņa priekšā. Viņš apstiprināja likumprojektus un budžetu, kurus pieņēma valdība. Viņš apstiprināja valdības ieteiktos augstākos ierēdņus. Pēc būtības viņa rokās koncentrējās diktatora pilnvaras. Tādas līdz tam Polijā vēl nevienam nebija piederējušas, vismaz pēc Jagelloņu valdīšanas laikiem. Tiesa, valsts pagaida priekšnieka parakstītajiem aktiem bija jātiek kontrasignētiem (\emph{kontrasignācija}~--- no latīņu \latxti{contra}~--- pret un \latxti{signare}~--- parakstīt; blakus valsts galvas parakstam bija jābūt arī valdības vadītāja vai atsevišķa ministra parakstam). Taču apstākļos, kad ministrus un valdības vadītāju nozīmēja valsts galva un tie bija atbildīgi viņa priekšā, tas ļoti nenozīmīgi ierobežoja viņa tiesības.

Kaut jau tad, kad J.~Pilsudskis iecēla I.~Dašinski par ministru prezidentu, visi līdz tam izdotie t.s. Ļubļinas valdības dekrēti zaudēja spēka, E.~Moračevska kabinets vēl turpināja Lubļinas valdības līniju, izdodot virkni dekrētu sociālajā jomā, kuri bija virzīti uz to, lai Polijā ieviestu standartus, kuri bija pieņemti attīstītajās Rietumu valstīs: astoņu stundu darba dienu, saīsinātu darba laiku sestdienās, sociālo apdrošināšanu slimības gadījumā, noteica minimālo algu valsts uzņēmumos, dzīvokļu īrētāju tiesību aizsardzību u.c. Pagaidām reformām, vai pareizāk~--- to solījumiem, sevišķi nepretojās arī labējie politiskie spēki. Šai sakarā var minēt nacionāldemokrātu līdera R.~Dmovska vārdus, teiktus 1919.~gadā un skaidrojošus valdošo slāņu attieksmi pret sociālajām reformām: ``Mēs esam no visām pusēm ielenkti no revolūcijām: krievu, ungāru, vācu. Mums nav armijas. Tāpēc nepieciešamības spiesti mēs izdarījām to, ko dabā dara vāji zvēri bez zobiem un nagiem: mēs tērpāmies aizsargkrāsā''.

Poļu vēsturnieks B.~Radļaks atzīmējis, ka I.~Dašinska un E.~Moračevska valdību veidošanās gaitā uzvarēja ``demokrātijas konsolidācijas'', bet ne ``nācijas konsolidācijas'' koncepcija, kam jāpateicas nevis poļu sociālistu vadītājiem, bet nepieciešamībai rēķināties ar plašu masu noskaņojumu. Tāpēc arī abas valdības vairījās atklāti pasludināt savu attieksmi pret revolucionāro Krieviju. Tomēr jau valdības deklarācijā 20.~novembrī, bet vēl vairāk E.~Moračevska paziņojumā 25.~novembrī bija izteikta tās politiskā orientācija. ``Valdības sociālajai politikai,~--- deklarēja E.~Moračevskis,~--- ir jāiet pa Rietumu reformu līniju, bet ne Austrumu, no kurienes draud haoss''. Acīmredzot lielā mērā ir taisnība krievu vēsturniekam G.~Matvejevam, ka I.~Dašinska un E.~Moračevska valdību sociālajā jomā darbaļaužu interesēs spertie soļi ne tikai nostādīja Poliju šai ziņā blakus attīstītajām valstīm, bet ļāva arī izbēgt no sabiedrības tālākas radikalizēšanās un sašķelšanās nesamierināmās nometnēs. Tiesa, visu nozīmīgāko reformu realizācija tika atlikta līdz Seima sasaukšanai.

Strīdīgs ir jautājums par J.~Pilsudska attieksmi šajā laikā pret sociālistiem. Kad 1918.~gada nogalē pie viņa ieradušies PPS pārstāvji un uzrunājuši viņu ar ``biedri Pilsudski'', viņš šo uzrunu noraidījis un lūdzis sevi saukt par ``panu''. Tomēr labējie sociālisti apgalvoja, ka viņš, tāpat kā pirms kara, ir to piekritējs. Toties radikāli kreisi noskaņotie J.~Pilsudska pretinieki pauda viedokli, ka viņš ir tikai ierocis buržuāzijas rokās.

Vēsturnieks P.~Olšanskis konstatējis, ka šajā laikā noformējās J.~Pilsudska piekritēju grupējums, ko mēdza saukt par ``Belvedēras nometni'' pēc valsts priekšnieka uzturēšanās vietas~--- Belvedēras pils Varšavā. Tā bija raiba gan pēc uzskatiem, gan pēc šķiriskās piederības, dažādu kreisu politisku grupu un organizāciju, kuras agrāk bija zināmas kā neatkarības aizstāves, apvienība. Tajā ietilpa bijušie leģionāri, Poļu kara organizācijas biedri, sociālisti, ļudoviešu partijas ``Atbrīvošanās'' un citu grupu locekļi. Šie cilvēki, atrodoties dažādās politiskās grupās, kurām bija savas programmas, tomēr kopumā darbojās vienā virzienā, ietekmējot daudzas sabiedriski-politiskas, kultūras un izglītības u.c. organizācijas. T.s. ``Belvedēras nometnei'' bija arī sava veida aģenti pat tajās organizācijās, kura sevi pretstatīja valsts pagaidu priekšnieka politikai. Tai pieslējās arī nebūt ne kreisi noskaņotu muižnieku grupas no ukraiņu, baltkrievu un lietuviešu apdzīvotajām zemēm, kuri cerēja, ka J.~Pilsudska vadībā viņiem izdosies iegūt tās Polijas sastāvā. Apvienības pretinieki, īpaši no kreisajām grupām, sāka tās piekritējus saukt arī par pilsudčikiem (\pltxti{pilsudczycy}). Ja agrāk tā galvenokārt dēvēja J.~Pilsudska vadītos leģionārus, tad tagad nosaukums ieguva politisku nokrāsu.

Pats valsts pagaidu priekšnieks bija uztraukts par to, ka liela sabiedrības daļa Polijā un ārzemju politiķi vēl joprojām viņu uzskata par sociālistu. Viens no viņa līdzgaitniekiem L.~Vasiļevskis atmiņās rakstīja, ka kādā šajā laikā notikušā sarunā J.~Pilsudskis izteicies: ``Ak, cik labi būtu, ja boļševiki organizētu pret mani atentātu, mestu bumbu vai izdarītu vēl ko līdzīgu \dots{} Dabīgi, atentāts neizdotos, bet kādu efektu tas izsauktu ārzemēs! Viņi uzreiz pārliecinātos, ka viss, ko runā par Moračevska valdību, ir muļķības.''

Sapnis par iespējamo atentātu drīz piepildījās. Tikai to veica nevis kreisie komunisti, bet labējie nacionāldemokrāti--endeki. Naktī no 1919.~gada 4.~uz 5.~janvāri sazvērnieku grupa mēģināja saņemt varu Varšavā. Tika arestēts valdības vadītājs un vairāki ministri. Viena no grupām vēlējās arestēt arī J.~Pilsudski. Taču sazvērestība bija slikti sagatavota un tika drīz apspiesta. J.~Pilsudskim uzticamās karaspēka daļas pierādīja savu spēku. Valsts pagaidu priekšnieks izsauca sazvērniekus uz Belvedēru, kur atradās viņa rezidence, rupjiem vārdiem nostrostēja un \dots{} atlaida. Daži no sazvērniekiem vēlāk pat taisīja veiksmīgu karjeru valsts dienestā. Sazvērestības neveiksme stiprināja J.~Pilsudska stāvokli. Klīda pat baumas, ka viņš pats to inspirējis. Taču jāatzīst, ka valsts pagaidu priekšnieka ietekme tai laikā nebija tik liela, lai viņš varētu kontrolēt arī savu politisko pretinieku darbību. Mīkstā attieksme pret sazvērniekiem ir skaidrojama ar viņa nevēlēšanos saasināt domstarpības ar labējiem. Taču vismaz daļēji izvairīties no savas kreisās pagātnes mantojuma J.~Pilsudskis varēja tikai nomainot kreiso E.~Moračevska valdību, kas arī drīz notika.

Laikā, kad gan Polijā (Varšavā, Ļubļinā, Poznaņā), gan aiz tās robežām (Parīzē) eksistēja dažādi politiskie spēki, kuri pretendēja uz veidojamās valsts vadību, E.~Moračevska valdība atradās ļoti sarežģītos apstākļos, tai no visām pusēm tika izteikti pārmetumi. Revolucionāri viņam kā sociālistam pārmeta oportunismu, strādnieku šķiras interešu nodevību, turpretī labējiem un centristiem viņa valdība likās pārāk radikāla. E.~Moračevska valdība dekretēja dažus agrākās I.~Dašinska valdības pieņemtos soļus sociālajā jomā: 8 stundu darbadienu, obligātu apdrošināšanu slimības gadījumiem, darba inspekcijas ieviešanu, darba algas minimumu. Taču E.~Moračevska valdības programmā nebija iekļauti vai bija atlikti solījumi par izmaiņām materiālo attiecību sfērā (par agrārās reformas īstenošanu, dažu rūpniecības nozaru nacionalizāciju u.c.). To varēja izskaidrot ar draudošo nacionālo spēku, visas poļu tautas šķelšanās. Taču, kā atzīmējis poļu jurists un tiesību vēsturnieks K.~Gržibovskis, atšķirībā no vairākuma citu valstu, kuras konstituējās sakarā ar agrākās tiesību sistēmas gāšanu, Polija uzsāka jauno tiesību radīšanu bez tās jauno pamatprincipu deklarēšanas. Tas sekmēja pēc būtības vecās sociālekonomiskās iekārtas saglabāšanu.

Tā kā pagaidām Poliju kā suverēnu valsti neatzina arī Antantes valstis, politiskās sadursmes pārliecināja J.~Pilsudski par to, ka nepieciešams panākt kompromisu ar nacionāldemokrātiem un R.~Dmovski, kompromisu ar nacionālajām poļu pārstāvniecības iestādēm gan Pozenē, gan Parīzē. 1918.~gada decembrī R.~Dmovska pārstāvji, pirmkārt jau bijušais sociālists, kurš bija pārgājis nacionāldemokrātu pozīcijās, S.~Grabskis uzsāka sarunas ar J.~Pilsudski. Tika panākta vienošanās par to, ka kreisas valdības atrašanās pie varas var izsaukt Rietumu valstu neuzticību Polijai, kurā notikumi varētu attīstīties līdzīgi kā Krievijā, tāpēc šī valdība ir jāreorganizē. J.~Pilsudskis atteicās no partiju cīņas saasināšanas, atlaida E.~Moračevska valdību un 17.~janvārī nomainīja to ar jaunu. To vadīja un reizē ārlietu ministra pienākumus pildīja I.~Paderevskis, kurš bija steidzami ar angļu kreiseri atgādāts no ASV.

Pa ceļam pēc Pozenes Augstākās Tautas padomes (\pltxti{Naczelna Rada Ludowa} vai \pltxti{Centralny Komitet Obywatelski}), kura kā pāri partijām stāvoša Prūsijas valdījumā esošo teritoriju poliska komiteja bija nodibināta jau 1916.~gadā (lielākā ietekme tajā bija nacionāldemokrātu partijai) un kuru 1918.~gada 5.~decembrī Pozenē sapulcējušies poļu pārstāvji atzina par likumīgu poļu pārstāvniecību, aicinājuma I.~Paderevskis apstājās Pozenē, izsaucot nacionāli sajūsminātu poļu manifestācijas. Ja sākotnēji Augstākā Tautas padome runāja par mierīgu varas pārņemšanu Pozenes poļu apdzīvotajos apgabalos, tad tagad Poļu kara organizācija un citi J.~Pilsudska piekritēji izmantoja iespēju lai apspēlētu Tautas padomi, vardarbīgi padzītu vācu militārās un civilās varas iestādes un saņemtu provincē varu poļu rokās. Jau otrajā dienā pēc I.~Paderevska ierašanās Pozenē~--- 27.~decembrī atskanēja pirmie šāvieni. 28.~decembrī Pozenes pilsēta atradās jau sacēlušos rokās. Sacelšanās drīz aptvēra visu Pozenes provinci un 1919.~gada janvāra vidū tā jau gandrīz visa bija poļu rokās. Februārī Augstākā Tautas padome pasludināja mobilizāciju un poļu vienību dalībnieku skaits izauga līdz 70~000. Kad 1919.~gada 16.~februārī Trīrā (\detxti{Trier}) Vācija bija spiesta parakstīt Kompjenas pamiera pagarinājumu, tā apņēmās atteikties no naidīgas darbības pret poļiem. Atsevišķas sadursmes gan turpinājās. 1919.~gada 19.~augustā Augstākā Tautas padome nodeva savus pilnvaras valdībai Varšavā.

Par I.~Paderevski kā valstsvīru kāds notikumu aculiecinieks (ārzemnieks) rakstīja: ``Poļi paši apgalvo, ka viņš visu laiku, kamēr ieņēma premjera postenī, palika vairāk pianists nekā valsts darbinieks''. Kaut oficiāli jaunā I.~Paderevska vadītā valdībā skaitījās bezpartijiska, tajā nebija vairs cilvēku, pret kuriem iebilda endeki. PPS un zemnieku partiju pārstāvji tajā vairs neietilpa. Atbildot uz dažu savu padomdevēju kritiku par kompromisu ar labējiem, J.~Pilsudskis sadusmots atbildēja: ``Man priekšā stāv cīņa ar Krieviju, bet ne ar Dmovski''. I.~Paderevska valdības izveide nozīmēja kompromisu starp Petrogradas-Parīzes Poļu nacionālās komitejas aprindām un J.~Pilsudska piekritējiem.

PPS vadības nostāja atspoguļojās partijas XV kongresa lēmumos 1918.~gada decembra vidū. Galvenajā ziņojumā, ko kongresā nolasīja F.~Perls, tika izvirzīta tēze par apstākļu trūkumu sociālistiskai revolūcijai tuvākajā nākotnē. Ziņotājs aicināja īstenot Polijā sociālās reformas, kuras pretstatīja ``asiņainai diktatūrai''.

Vilšanās PPS vadības politikā un tās atbalstītajā E.~Moračevska valdībā, labējo apvērsuma mēģinājums 1919.~gada janvāra sākumā, pēc tam E.~Moračevska valdības atlaišana un I.~Paderevska valdības izveide~--- tas viss noveda pie kreiso grupas ar T.~Žarski un A.~Landu priekšgalā izveidi PPS sastāvā. Šī grupa atzina proletariāta diktatūras nodibināšanas nepieciešamību (neatsakoties no Polijas neatkarības), uzskatīja, ka sociālistiem jācenšas revolucionārā ceļā izveidot sociālistisku valdību. Drīz grupa tālākas evolūcijas rezultātā sarāva sakarus ar PPS un 1920.~gada augustā iestājās Polijas Komunistiskajā strādnieku partijā.

Valdībai klajā opozīcijā atradās kreisie no PKunLSD un PPS-ļevicas, kuri organizēja strādnieku deputātu padomes. 1918.~gada 16.~decembrī abas minētās sociālistiskās partijas apvienojās un izveidoja Polijas Komunistisko strādnieku partiju (\pltxti{Komunistyczna Partia Robotnicza Polski}, no 1926.~gada~--- Polijas Komunistiskā partija~--- \pltxti{Komunistyczna Partia Polski}). Tā no PKunLSD mantoja negatīvo attieksmi pret nāciju pašnoteikšanās lozungu, nacionālās neatkarības iegūšanu un nacionālas valsts izveidi it kā derīgu tikai pirmsmonopolistiskā kapitālisma laikmetam un konkrētajos apstākļos nesavienojamu ar proletārisko internacionālismu. Komunisti gan pasludināja par savu galveno uzdevumu proletāriskās revolūcijas sagatavošanu un īstenošanu un sociālisma uzcelšanu Polijā, taču nespēja kavēt varas pāreju pilsonības rokās. Poļu komunisti jau ar pirmajām jaunās valsts pastāvēšanas dienām visā savā propagandā un aģitācijā uzsvēra, ka tā nav patiesi neatkarīga, ka tauta šai valstī nevar pati izlemt savu likteni, tās vietā lemj pašmāju un ārzemju buržuāzija. Pelnīti un nepelnīti pārmetumi pašmāju uzņēmējiem un muižniekiem mijās ar pārspīlētu ārzemju kapitālistu lomas izcēlumu, kuri it kā pilnībā nosakot jaunās valsts rīcību. Īpaši asi izskanēja to poļu politiķu nosodījums, kuri Pasaules kara laikā ``izkalpojās Franča Jozefa un Vilhelma priekšā un pārdeva Austrijai un Vācijai leģionos esošo poļu asinis''. Kreisie spēki nepietiekami novērtēja ar nacionālo atbrīvošanos saistīto eiforiju, kura valdīja tautā. Pie tam kreisajiem nācās konkurēt ne tikai ar labējiem buržuāziskajiem spēkiem, bet arī ar sociālistisko PPS, kura kopumā atbalstīja valsts priekšnieka un tā ieceltās valdības pasākumus. Pie tam PKSP nodibināšanu pavadīja tik plaša mēroga tās piekritēju aresti, ka partija praktiski nevarēja darboties legāli.

Strādnieku deputātu padomju veidošana balstījās uz zināmu revolucionārā viļņa pacēlumu Polijā. Taču to veidošanā piedalījās revolucionāri noskaņoti strādnieki ar dažādiem uzskatiem. Krievu vēsturnieces I.~Jažborovska un V.~Parsadanova, izpētījušas 19~lielāko strādnieku padomju, kuras pārstāvēja 221 tūkstoti strādnieku, materiālus, konstatēja, ka par PPS balsoja ap 70 tūkstošu cilvēku, par komunistiem~--- ap 63 tūkstoši. Attiecīgi arī padomju sastāvā sociālistu bija vairāk nekā komunistu (izņemot Dombrovskas baseinu, kur padomēs komunisti bija pārsvarā). Pavisam bijušās Polijas karalistes teritorijā pilsētās tika radītas ap 120~strādnieku padomes. Bez tām laukos veidojās zemnieku un laukstrādnieku padomes. Pret visām tām jau decembrī ar bruņotu spēku vērsās nacionālistiski noskaņotās poļu karavīru vienības. Strādnieku vidū 1918.~gada beigās populāra kļuva rūpniecības uzņēmumu nacionalizācijas ideja. 12.~decembrī streikojošie Varšavas strādnieki ielu manifestācijas laikā izvirzīja lozungu ``Mēs pieprasām rūpnīcu sabiedriskošanu un to darba atsākšanu''. Daudzos uzņēmumos tika nodibināta strādnieku kontrole, strādnieki griezās pie E.~Moračevska valdības pēc finansiālas palīdzības rūpnīcu darba atjaunošanai. ``Valdība,~--- vēlāk atzinās E.~Moračevskis,~--- sistemātiski atteica strādnieku delegācijām izsniegt naudu, jo tās, pirmkārt, nebija, bet galvenokārt tāpēc, ka ar to lika strādniekiem atteikties no vilinājuma patstāvīgi sakārtot ražošanu''. Strādnieku deputātu padomes atbalstīja nacionalizācijas prasību. Nespējot streikot, dīkā stāvošo Lodzas uzņēmumu bez darba esošie strādnieki ieņēma rūpnīcas, konfiscēja izejvielas un gatavās produkcijas krājumus, aplenca saimnieku dzīves vietas un kantorus, pieprasot atjaunot darbu uzņēmumos, sniegt materiālu palīdzību tajos strādājošajiem. Taču streiki, demonstrācijas, rūpnīcu ieņemšana u.c. gandrīz nekur nevarēja apstādināt dzīves līmeņa krišanos. (Vienīgi Dombrovskas baseinā trīs nedēļu ilgs ogļraču streiks nodrošināja 8 stundu darba dienas ieviešanu, algas paaugstināšanu.) Sociālistu ietekme pieauga. Pie tam galvenajos strādnieku centros drīz tika ievests ārkārtējais stāvoklis: janvārī Varšavā, martā Ļubļinas vojevodistē un Dombrovskas baseinā, pēc tam arī citur. Pret strādniekiem tika nosūtīts karaspēks. Dombrovskas baseinā strādnieku izveidotā Sarkanā gvarde ap pusotru mēnesi pretojās karaspēkam.

PKunLSD un PPS-ļevica, bet pēc tam PKSP, vērtēja strādnieku deputātu padomes kā nākamās proletariāta diktatūras iestādes. Turpretī PPS vadība centās to lomu novirzīt vienīgi ekonomiskās cīņas laukā, ierobežot ar saimnieciskiem un kultūras pasākumiem. Padomju politiskā loma tika aprobežota ar I.~Dašinska un E.~Moračevska valdību atbalstīšanu. Partiju pretējās politiskās pozīcijas bija viens no strādnieku deputātu padomju vājuma cēloņiem. Līdz 1919.~gada vidum tās tika likvidētas. Kopumā poļu strādnieku kustība bija Pasaules kara novājināta, paši strādnieki marginalizējušies un nespēja aizstāvēt savas šķiriskās intereses. (Tā 1918.~gada beigās bijušās Polijas karalistes teritorijā bija palikuši tikai ap 47~000 strādnieku, kas sastādīja 14\% no pirmskara skaita.) Komunistiem sekoja tikai daļa no tiem.

Šai laikā boļševiki Krievijā cerēja uz to, ka Polijā pie varas nāks revolucionārie spēki. KSFPR VCIK 1918.~gada 13.~novembra dekrētā bija uzsvērts, ka Polijas darbaļaudis ``tagad ir aicināti paši lemt savu likteni'', bet abu valstu attiecību pamatā ``var būt tikai tie principi, kuri atbilst visu zemju un nāciju darbaļaužu brālīgajām attiecībām un kurus ir pasludinājusi Oktobra revolūcija''. Arī padomju vēsturnieki savā laikā apgalvoja, ka 1918.~gadā Polijā esot pastāvējuši priekšnoteikumi sociālistiskai revolūcijai. Faktiskie notikumi, poļu tautas vairākuma atbalsts savas nacionālas valsts izveidei, nevis sociālā taisnīguma ideālu īstenošanai gan parādīja, ka apgalvojums neatbilda patiesībai.

Tā pirms 1914.~gada gandrīz neiedomājamais triju virsvaru: Krievijas, Vācijas un Austroungārijas sabrukums deva iespēju veidoties jaunajai Polijas valstij. Šādas iespējas rašanos ievērojamā mērā veicināja boļševiku darbība, graujot Pirmā pasaules laikā Krievijas militāro varenību un īstenojot Oktobra apvērsumu. Krievijai paliekot līdz kara beigām Antantes rindās, saglabājot militāro potenciālu, tā visticamāk nepieļautu savas provinces Polijas atdalīšanos. Augstākais uz ko šādā gadījumā varētu cerēt poļi~--- visu viņu apdzīvoto apgabalu apvienošana un autonomija Krievijas sastāvā.

Latviešu vēsturnieks Aldis Miņins pamatoti pievērsis uzmanību kanādiešu vēsturnieces M.~Makmillanas secinājumam, ka Polijas labās izredzes kļūt par ievērojamu starptautisku spēku noteica trīs lietas: pirmkārt, visi trīs Polijas sadalītāji~--- Austroungārija (sadalīšanas brīdī Austrija), Vācija (sadalīšanas brīdī Prūsija) un Krievija~--- bija Pirmā pasaules kara zaudētājvalstis. Otrkārt, atšķirībā no citām Krievijas postimpērijas Rietumu perifērijas teritoriālām vienībām, tikai Polijai bija nepārprotams atbalsts Rietumu sabiedroto vidū patstāvīgam valstiskumam. Viena no V.~Vilsona prasībām bija brīva Polija. Treškārt, poļiem bija ``pilna nācija'' ar pabeigtu sociālo struktūru. Turklāt nacionālā ideja bija svarīga visām poļu sabiedrības grupām~--- šļahtičiem, inteliģencei, darba devējiem un darba ņēmējiem.

Tomēr jaunā valsts tapa smagos apstākļos. Daudzas pilsētas un ciemi bija ievērojami cietuši karā, saimnieciskie sakari tā rezultātā bija sarauti, bija jāveido vienotas valsts tautsaimniecība no bijušām triju valstu nomalēm. Nebija savas naudas sistēmas (kursēja piecas valūtas), savu pārbaudītu, tautas uzticību baudošu pārvaldes iestāžu, beidzot, nebija nospraustas valsts robežas.

Parasti poļu literatūrā tiek norādīts, ka KSFPR ar Brestas miera līgumu atsacījās no tiesībām uz Lietuvu, Rietumu Baltkrieviju un Rietumu Volīniju, bet 1918.~gada augustā, pildot šo līgumu, KSFPR Tautas Komisāru Padomes atcēla arī visus Krievijas impērijas līgumus ar citām valstīm par Polijas sadalīšanu. Taču tiek ignorēts fakts, ka ar Antantes valstu un Vācijas 11.~novembrī parakstīto Kompjenas pamieru un 13.~novembrī izdoto Padomju valdības dekrētu tika anulēts Brestas miera līgums un līdz ar to arī KSFPR TKP augusta dekrēts. Ar to Padomju Krievija tiesiski varēja pretendēt uz visām bijušajām Krievijas impērijas teritorijām. (Starp citu, pēc Kompjenas pamiera noteikumiem Vācija apņēmās izvest savu karaspēku tikai no tām zemēm, kuras tā bija sagrābusi Pirmā pasaules kara laikā, bet ne no agrāk Prūsijai piederējušajām teritorijām.) Vācijas un Austroungārijas karaspēks 1918.~novembrī atstāja bijušās Polijas karalistes un Rietumgalīcijas teritoriju. 1918.~gada 12.~decembrī Sarkanās armijas daļas saņēma pavēli ieņemt teritoriju līnijā: Paņeveža, Viļņa, Lida, Baranoviči. Pavēles izpilde īpašas grūtības neradīja, jo Sarkanās armijas daļas vienkārši sekoja atejošajam Vācijas karaspēkam. Saskaņā ar starptautiskajiem tiesību aktiem to nevarēja uzskatīt par Padomju Krievijas agresiju. (Arī ģenerāļa A.~Deņikina komandētā ``balto'' karaspēka uzbrukumu Ukrainā 1919.~gadā, ne arī Polijā Rietumu lielvalstis nekvalificēja kā agresiju.) Ar Poliju padomju puse saprata tikai bijušās Polijas karalistes teritorijas, bet rajonus, kuri iegāja Krievijas Rietumu novada sastāvā par poļiem piederošiem neuzskatīja. Tos Padomju Krievija atzina par piederīgiem Padomju Lietuvai un Padomju Baltkrievijai.

Topošās Polijas valsts vadība šādus starptautisko tiesību ``smalkumus'' neņēma vērā. 16.~novembrī J.~Pilsudskis paziņoja visu valstu, izņemot Padomju Krievijas, valdībām par neatkarīgas Polijas valsts izveidi. (Der atzīmēt, ka boļševiku valdība Krievijā, kura bija padzinusi Satversmes sapulci, atšķirībā no Vācijas sociāldemokrātu sastādītās Tautas pilnvaroto padomes (\detxti{Rat der Volksbeauftragten}), kura organizēja demokrātiskas Nacionālās sapulces (\detxti{Nationalversammlung}) sasaukšanu, vēl ilgi palika starptautiski neatzīta.)

18.~novembrī Poļu Nacionālā komiteja Parīzē apsprieda ziņu, ka Sabiedroto karaspēka virspavēlnieks maršals F.~Fošs saņēmis telegrammu no J.~Pilsudska ar lūgumu Polijas aizsardzībai pret boļševikiem okupēt to ar sabiedroto un Francijā esošo poļu karaspēku. Komiteja šai sakarā nolēma vest sarunas ar Francijas Ārlietu ministriju, lai panāktu J.~Hallera armijas pārcelšanu uz Poliju. Savukārt Francijas Ārlietu ministrija, pirms nosūtīt savus pārstāvjus uz Varšavu, konsultējās ar Poļu Nacionālo komiteju, kura bija gatava sadarboties ar J.~Pilsudski cīņā pret boļševismu, bet uzsvēra, ka īsteno tautas gribu jāpauž Satversmes sapulcei.

Stājoties valsts priekšgalā, J.~Pilsudskis noslēdza vienošanos ar Vācijas Centrālo karavīru padomi, ka apmaiņā pret vācu karavīru atgriešanās garantijām Vācijā tā apņemas atstāt poļiem pārtikas krājumus un ieročus. Pēdējo bija daudz un pietika pirmējai Polijas armijas apbruņošanai. Arī par pirmo valsti, kura 1918.~gada 18.~novembrī atzina Polijas valsti, kļuva Vācija, kuras toreizējā valdība~--- Tautas pilnvaroto padome kā sūtni uz Varšavu nosūtīja jau agrāk ar J.~Pilsudski kontaktus uzņēmušo grāfu H.~Kesleru, liekot viņam arī paziņot, ka Vācija piekrīt Pozenes provinces daļas nodošanai Polijai. Grāfs H.~Keslers sākotnēji tika draudzīgi sagaidīts, taču iespaidīgs Nacionāldemokrātu partijas un Parīzē pastāvošās Poļu nacionālās komitejas spiediens lika J.~Pilsudskim decembra vidū paziņot par attiecību saraušanu ar Vāciju. To saglabāšana nostādītu Poliju neitrālas valsts stāvoklī un neļautu tai piedalīties Pirmā pasaules kara uzvarētājvalstu sasaucamajā Parīzes miera konferencē, kur poļi cerēja ar Rietumvalstu atbalstu iegūt teritorijas, kuras agrāk atradās Vācijas sastāvā. H.~Keslers atstāja Varšavu. Daudzos vāciešos poļu ``nepateicība'' tikai pastiprināja rūgtumu. Feldmaršals P.~f.~Hindenburgs memuāros rakstīja: ``Man nekad agrāk nebija personīgu antipātiju pret poļu tautu, no otras puses \citespace{} es nešaubījos, ka mēs no poļiem nekad nesaņemsim ne miņas no pateicības, ka mēs viņus ar mūsu zobenu un asinīm atbrīvojām no krievu pātagas, tik maz mēs bijām saņēmuši atzīšanas par to saimniecisko un garīgo pacēlumu, ko piedzīvoja poļu tautas daļa Prūsijā.''

Kaut Antantes valstis sākotnēji vairāk uzticējās Poļu nacionālajai komitejai Parīzē, kura jau sen orientējās uz Rietumu sabiedrotajiem, nevis J.~Pilsudskim, kurš Pirmā pasaules kara pirmajā pusē aktīvi sadarbojās ar Centrālajām valstīm, tās bija gatavas sadarboties arī ar ietekmīgāko poļu politiķi. Pats J.~Pilsudskis centās sabiedroto acīs attīrīties no pārmetumiem par iepriekšējo sadarbību ar Centrālajām valstīm pret Antanti, ``nomazgāt'' savu pagātni. Intervijās Rietumu avīzēm J.~Pilsudskis pieslīpēja argumentāciju, kāpēc viņš bija karojis Centrālo valstu pusē. Lūk, viens variants, izskanējis 1919.~gada sākumā: ``Es vienmēr biju Antantes draugs, bet man pirmkārt nācās rūpēties par manas tēvzemes labklājību. Tā diktēja manu cīņu pret carismu, kas nenozīmēja vēlmi cīnīties pret Antanti. Tāpēc pēc carisma sabrukuma es pārtraucu cīņu pret Krieviju un atlaidu manus leģionus. Atgriezies Polijā pēc Centrālo valstu sakāves, es sapratu, ka mana personība var traucēt tuvu attiecību izveidi starp manu tēvzemi un Antanti. Es uzdevu sev jautājumu, vai man jāpaliek vai jānoiet malā. Es pārdzīvoju briesmīgu iekšēju cīņu. Un šoreiz mīlestība pret tēvzemi lika man palikt, jo es sapratu, ka mana klātbūtne paglābs Poliju no vislielākajām briesmām~--- postošas revolūcijas, bet pretēji, ļaus tai nostāties uz mierīga evolūcijas ceļa. Tāpēc es paliku. Un tagad, kad starp mani un Antanti vairs nestāv carisms, es varu atklāti runāt par manām draudzīgajām jūtām pret Antanti, ar kuru mums arī turpmāk jāstiprina draudzīgas saiknes''. Tiesa, šai retorikai daudzi tā arī nenoticēja. Lielbritānijas premjerministrs D.~Loids Džordžs 1921.~gada 13.~maijā skaidri pateica: ``Ja mēs rēķinātos ar poļu karaspēku, kurš cīnījās Centrālo valstu pusē, tad Polija šodien būtu vācu vai austriešu novads.''

Rietumu valstis 1919.~gada 18.~janvārī Poliju atzina kā asociētu (ar tām saistītu) varu un ļāva piedalīties Parīzes miera konferencē. Faktiski ar to tās jau atzina Poliju kā suverēnu valsti. Oficiāli pirmā Rietumu lielvalsts, kas 1919.~gada 30.~janvārī atzina Polijas valsts neatkarību \latxti{de iure} bija ASV. Iespējams, tas notika tāpēc, ka ASV mita miljoni poļu izcelsmes vēlētāju, poļu emigranti organizēja vērienīgas masu kampaņas ar mērķi ietekmēt prezidentu un citus varas nesējus, lai viņi ieņemtu lojālu pozīciju pret atjaunojamo Polijas valsti. Februārī par Polijas \latxti{de iure} atzīšanu paziņoja arī Francija, Lielbritānija un Itālija.

Uz agrāko poļu militāro formējumu bāzes veidojās jaunās valsts armija. Tā Polijai šai laikā bija objektīvi nepieciešama. Tikai balstoties uz to varēja radīt stipru valdību, nosargāt Poliju no nemieriem, kuri bija aptvēruši bijušo Krievijas impēriju un kuru izplatību uz rietumiem pagaidām vēl novērsa Centrālo valstu armijas, bet bija redzams, ka ilgi tās to vairs nespēs izdarīt. Protams, Polijas pilsonībai armija bija vajadzīga arī cīņai pret pašu poļu radikālajiem spēkiem~--- PKunLSD un PPS-ļevicu. Bez armijas nebija arī iedomājama Polijas politiķu iecerētā neatkarīgajai valstij labvēlīgu robežu iegūšana.

Sākotnēji par būtisku spēku Polijā kļuva Poļu kara organizācija (POW), kuru vadīja E.~Ridzs-Smiglijs. Tās vienības apsargāja valdību, 11.~novembrī sāka atbruņot Vācijas armijas karavīrus, savākt to ieročus. Poļu militārais vēsturnieks Z.~Zaluskis sniedzis datus, ka vācu atbruņošanā Varšavā un citur piedalījās nedaudz vairāk nekā desmit tūkstošu poļu. Armija sākumā formējās uz brīvprātīguma pamata no POW biedriem, bijušajiem leģionāriem, t.s. poļu vērmahta, citiem Austroungārijas, Vācijas armiju poļu tautības karavīriem, Krievijā pastāvējušo poļu brīvprātīgo korpusu dalībniekiem, ģimnāziju vecāko klašu audzēkņiem un studentiem, arī jauniešiem no strādnieku un zemnieku aprindām. Taču teritoriālās pretenzijas, kuras drīzumā izvirzīja valdība, prasīja daudz lielāku armiju. Brīvprātīguma princips neļāva tādu radīt. Tāpēc drīz tika ieviesta visu Polijas pilsoņu vispārējā karaklausība. Visas armijas daļas, kuru karavīri bija dienējuši dažādu valstu armijās, vajadzēja unificēt. Tika meklēta sasaiste ar Polijas varonīgo pagātni. Tā, ārzemnieki brīnījās, redzot poļu karavīrus sasveicinoties paceltai rokai izstiepjot tikai rādītāju un vidējo pirkstu, bet pārējos saliecot. Viņiem paskaidroja, ka tā ir veca tradīcija: ``Tā sveicinājās jau karaļu laikos. Divi izstieptie pirksti nozīmēja uzticību Dievam un karalim. Tagad viņu nozīme: uzticību Dievam un valstij''.

J.~Pilsudskis vispirms centās radīt militāristu grupu, kas būtu viņam pilnīgi uzticīgi, atrastos pilnīgā viņa rīcībā. Atbildīgākajos amatos viņš parasti nozīmēja personīgi uzticamus cilvēkus, pat ja tie ne visai labi tika galā ar saviem pienākumiem. Taču šādu militāru kadru nebija daudz. Pats J.~Pilsudskis, kuram bija visaugstākā pakāpe bijušo strēlnieku un leģionāru vidū, bija tikai brigadieris (brigādes ģenerālis)~--- tātad ar viszemāko ģenerāļa pakāpi. Leģionāru vidū bija ap pusotra desmita pulkvežu, dažas reizes vairāk majoru. Vairākumā tie bija virsnieki--praktiķi, kuriem nebija pat vidējās militārās izglītības. Taču armijai bija vajadzīgi kara ministrs, Ģenerālštāba priekšnieks, fronšu, operatīvo grupu utt. komandieri. Tāpēc nācās iesaistīt armijā arī virsniekus ar augstāko izglītību, kuri bija dienējuši Vācijas, Austroungārijas un Krievijas armijās. Poļu vēsturnieks S.~Jačinskis saskaitījis, ka Polijas armijā 1920.~gada pavasarī bija 93 ģenerāļi, kuri bija guvuši kaujas pieredzi trīs impēriju (Austroungārijas, Krievijas un Vācijas) armijās virsnieku statusā. Viņiem, atšķirībā no leģionāriem, trūka saliedētības, uzticības savam vadonim. Tāpēc vēlāk leģionāri sekmīgāk virzījās pa karjeras kāpnēm. Kā piemērus var nosaukt K.~Sosnkovski, E.~Ridzs-Smigliju, V.~Sikorski. Ar J.~Pilsudska atbalstītājiem konkurējošie pilsoniskie spēki kā pretsvaru militāri ``neizglītotajiem'' leģionāru komandieriem parasti pretstatīja ``īstus'' virsniekus, kuri jau pirms Pirmā pasaules kara bija dienējuši Poliju sadalījušo valstu armijās. Tā, viens no šo spēku kandidātiem uz poļu armijas Ģenerālštāba priekšnieka amatu bija ģenerāļa J.~Hallera brālēns S.~Hallers, bijušais Austroungārijas armijas pulkvedis, beidzis tās Kara akadēmiju, pēc tam Polijas armijas ģenerālis.

J.~Pilsudskis lika jauno bruņoto spēku ietvaros atjaunot bijušās leģionāru daļas, paplašinot to sastāvu līdz trijām divīzijām. Virspavēlnieks demonstrēja agrāko aizvainojumu aizmiršanu. Jaunajos formējumos tika uzņemti arī tie, kuri 1917.~gada jūlijā nodeva zvērestu Vācijai un ar to nepaklausīja viņa ieteikumam. Līdz pat Otrajam pasaules karam uz leģionu pamata veidotās divīzijas Polijā pildīja sava veida gvardes lomu.

1919.--1920.~gadā armija formējās galvenokārt tikai no poļiem un daļēji ebrejiem. Tāds gandrīz mononacionāls sastāvs labvēlīgi ietekmēja tās kaujas garu. Tiesa, pieņemtais teritoriālais armijas daļu veidošanas princips kavēja dažādu novadu iedzīvotāju savstarpējo integrāciju, taču ievērojot to, ka vairāk kā simt gadu šo novadu iedzīvotāji bija dzīvojuši šķirti, tas acīmredzot bija pareizākais ceļš kā izvairīties no nesaprašanās un ķildām. Pirmajos valsts pastāvēšanas gados kā visdisciplinētākās, ar vislielākajām kaujas spējām, tiesa, arī kā visvairāk antisemītiski noskaņotas, sevi apliecināja daļas, veidotas Poznaņas (agrākās Pozenes) novadā, kur pats ietekmīgākais politiskais spēks bija nacionāldemokrāti.

Pēc dažām ziņām 1919.~gada janvārī Polijas armijā bija jau ap 100~000~karavīru. 1919.~gada pirmā ceturkšņa beigās armijas skaitliskais sastāvs tuvojās 200~000. Aprīlī no Francijas uz Poliju tika pārdislocēta arī ap 70~000 vīru lielā ģenerāļa J.~Hallera armija. Līdz 1919.~gada augustam Polijas armijas skaitliskais sastāvs pieauga aptuveni no 100~000 līdz 500~000 karavīru,

Reizē J.~Pilsudska vadībā jau 1919.~gada 26.~janvārī teritorijās, kuras tolaik kontrolēja Polija (bijušajā Polijas karalistē un Rietumgalīcijā), tika organizētas \strong{Satversmes seima} vēlēšanas (citur vēlēšanas notika vēlāk).

Tuvākais iekšpolitiskais uzdevums bija panākt sabiedrības konsolidāciju jaunās valsts atbalstam, tās armijas radīšanai. Radīt Polijā tādu iekārtu, kas būtu pievilcīga arī citām tautām, kuras J.~Pilsudskis tajā vēlējās iekļaut, kā arī saņemt Rietumu demokrātisko valstu atbalstu, bez kura nebija iedomājama jaunās valsts pastāvēšana, varēja tikai demonstrējot uzticību demokrātijai.

J.~Pilsudskis 1918.~gada beigās iestājās par parlamentāras, demokrātiskas valsts izveidi. Viņš labi saprata, ka šai laikā karā uzvarējušajām rietumvalstīm tikai brīvās vēlēšanās ievēlēts parlaments varēja liecināt par jaunveidojamo valstu leģitīmo raksturu. Parlamentāram režīmam, pēc valsts priekšnieka domām, bija jānobremzē radikālo politisko elementu ietekmes augšana masās, kuras bija neapmierinātas ar smagajiem dzīves apstākļiem. Viņš arī cerēja, ka Seims atbrīvos to no nepieciešamības stingri uzraudzīt valdības darbību un ļaus tam koncentrēt uzmanību armijas radīšanai un valsts robežu problēmai. 1918.~gada 8.~decembrī J.~Pilsudskis tai laikā pie varas esošo partiju~--- sociālistiskās (PPS) un ļudoviešu (\pltxti{Polskie Stronnitctwo Ludowe ``Wyzwolenie''})~--- pārstāvjiem paziņoja: ``Mani galvenie centieni tagad veltīti Seima sasaukšanai. Polijā visi kliedz, ka tiem ir vairākums. Taču tikai Seims noskaidros un noteiks, kur ir šis vairākums un ko viņš grib''. Ar to valsts pagaidu priekšnieks brīvprātīgi atteicās no savām diktatora pilnvarām, lai sasniegtu savus iekšpolitiskos un ārpolitiskos mērķus. Saviem līdzstrādniekiem J.~Pilsudskis paaugstinātos toņos skaidroja: ``Iekšējās lietas izlems Seims, kuru tādēļ es arī sasaucu. Kāds tas būs~--- kreiss vai labējs~--- paskatīsimies. Visiem maniem pūliņiem ir jātiek vērstiem uz armiju. \citespace{} Kad man būs armija, viss būs manās rokās.''

Kā norādījuši poļu vēsturnieki D.~un T.~Nalenči, J.~Pilsudska rūpes par Seima sasaukšanu nenozīmēja, ka viņš parlamentu vērtētu kā pašu pilnīgāko varas realizācijas institūtu. Viņš uzskatīja, ka pats labāk saprot un realizē tautas intereses; nevis viņam ir jākļūst par Seima gribas izpildītāju valsts celtniecībā, bet gluži pretēji~--- parlamentam ir jāpalīdz viņam šī uzdevuma izpildē. Turpretī Seima deputāti uzskatīja sevi par īstajiem Polijas Republikas saimniekiem. Šīs savas lomas dažādajā izpratnē arī slēpās vēlāko konfliktu galvenais cēlonis.

Kā galvenie spēki Seima vēlēšanās uzstājās politiskās partijas. Tajā laikā katra no tām un pastāvošā partiju sistēma kopumā vēl neveidoja stingru politisku organismu. Pat vienāda idejiski politiskā un sociālā rakstura partiju saplūšana bieži sastapa šķēršļus, ko izsauca sadarbības tradīciju trūkums, atšķirīga iepriekšējā politiskā pieredze vai iepriekšējo idejiski politisko programmu nianses. Bez lielām grūtībām notika tikai dažādu sociālistu grupu, kuras orientējās uz sociālo reformu īstenošanu, apvienošanās vienotā PPS, taču arī tā noslēdzās 1919.~gada aprīlī, jau pēc Satversmes seima vēlēšanām. To laikā un Satversmes seima darbības sākumposmā nacionāldemokrātiem arī izdevās ap sevi apvienot dažādas labējas buržuāziskas un muižnieciskas grupas. Citas partijas turpretī vēl joprojām uzstājās kā atsevišķu Polijas daļu interešu izteicējas. Kopumā vienotai valstij raksturīgas partiju sistēmas izveide vēl bija priekšā.

Sociālisti (PPS) uz vēlēšanām gāja, solot parlamentārā ceļā īstenot sociālistiskus pārkārtojumus. Ļudovieši (\pltxti{Polskie Stronnitctwo Ludowe ``Wyzwolenie''}, \pltxti{Polskie Stronnictwo Ludowe ``Piast''}, \pltxti{Polskie Stronnictwo Ludowego~--- Lewica} u.c.) kā galveno izvirzīja agrārās reformas lozungu. Šīs dažādās zemnieku partijas atšķīrās gan pēc ceļu izvēles cīņā par nacionālo brīvību (\pltxti{Piast} orientējās uz sadarbību ar nacionālajiem demokrātiem, \pltxti{Wyzwolenie} un \pltxti{Lewica}~--- ar J.~Pilsudska piekritējiem), gan pēc savu sociāli ekonomisko programmu radikālisma. Kā jau rāda nosaukums, kreisajā flangā atradās \pltxti{Lewiva}, uz centru tiecās \pltxti{Wyzwolenie}, bet labējākas pozīcijas tolaik ieņēma \pltxti{Piast}.

Labējās partijas (nacionāldemokrāti, kristīgie demokrāti u.c.) solīja nodibināt valstī ``kārtību'', ``likumību'', ``sociālu līdzsvaru'', reformu īstenošanu uz ``kristīgā taisnīguma'' pamata, rūpniecības, tirdzniecības un kultūras polonizāciju. To aģitācijai bija nacionālistisks, nereti arī antisemītisks saturs.

Liela loma vēlēšanās bija katoļu baznīcai. Sabiedrības demokratizācijas perspektīva, jaunu masu iesaistīšana politiskajā dzīvē ar antiklerikāliem lozungiem draudēja mazināt reliģijas tradicionālo ietekmi, apdraudēt tās mantisko stāvokli. Tāpēc daudzi garīdznieki tieši iesaistījās politiskajā darbībā, balotējās vēlēšanās labējo un centrisko partiju sarakstos. Sakarā ar vēlēšanām izplatītajā bīskapu vēstījumā draudzēm tika nosodīta sociālo reformu ideja un ticīgie aicināti balsot par labējām partijām.

J.~Pilsudskis un viņa tuvākie līdzgaitnieki neveidoja kādu politisku partiju, neizvirzīja savas kandidatūras vēlēšanās, bet baudot nacionālistiski noskaņoto masu uzticību, varēja ietekmēt situāciju ar savu autoritāti un militāristu atbalstu. J.~Pilsudska un tā līdzgaitnieku E.~Ridzs-Smiglija, K.~Sosnkovska u.c. piekritējus, kuri oficiāli bija bezpartijiski, šodien arī literatūrā pieņemts saukt par pilsudčikiem.

Vēlēšanās dažādos apgabalos piedalījās no 60 līdz 94\% balsotāju. Seima darbības sākumā tajā bija 335~deputāti, bet darbības beigās 1922.~gadā~--- 432 deputāti.

Seima vēlēšanas parādīja, ka nacionāldemokrāti ir pati ietekmīgākā partija Polijā. (1919.~gadā tā pārdēvējās par Nacionāldemokrātisko savienību (\pltxti{Związek Ludowo-Narodowy}), bet 1928.~gadā par Nacionālo partiju (\pltxti{Stronnictwo Narodowe}).) Tai pieslējās arī citas labējās partijas. Nacionāldemokrāti un viņu sabiedrotie guva relatīvi lielāko balsu (ap 37\%) un arī deputātu mandātu skaitu. Izrādījās, ka, neraugoties uz sabiedrības radikalizācijas šķietamību, vēlētāju vairākums stāv konservatīvās vai mērenās pozīcijās. Vēlēšanas arī parādīja, ka poļu sabiedrība nebūt nenosodīja nacionāldemokrātus par viņu sadarbību ar Krieviju un neizjuta pastiprinātas simpātijas pret partijām, kuras aktīvi cīnījās par Polijas karalistes atbrīvošanu no Krievijas impērijas jūga. Ar to arī taktiski pareiza izrādījās J.~Pilsudska atteikšanās no balstīšanās uz viņam agrāk tuvajām kreisajām partijām.

Centrs, kam sākotnēji bija pieskaitāma gan zemnieku partija ``\pltxti{Piast}'', kura, lai gan radusies 1913.~gadā Galīcijā, no 1918.~gada savu darbību izplatīja visā Polijā, gan citas sīkākas partijas, saņēma 15,4\% balsu. Var atzīmēt, ka Satversmes Seimā bija ievēlēti 127~zemnieku deputāti, taču viņi piederēja dažādām partijām. Apmēram puse zemniecības bija nobalsojusi par partijām, kuras nebija saistītas ar zemnieku kustību.

Kreisie (galvenokārt sociālisti un ļudovieši no partijas ``\pltxti{Wyzwolenje}'', arī citi nekomunisti, kuri bija atbalstījuši t.s. ``tautas'' valdības ar I.~Dašinski un E.~Moračevski priekšgalā) vēlēšanās ieguva 34\% balsu. PPS teorētiski pasludināja parlamentu par savu programmatisko mērķu sasniegšanas instrumentu, bet praktiski tai nebija iespaida, lai panāktu kaut mazākos pārveidojumus. Partijā arī nebija vienotības.

11,5\% balsu saņēma ebreju un vācu nacionālo mazākumu organizācijas. Tā, ebreju partijas panāca 11~savu pārstāvju ievēlēšanu. (Literatūrā ir sastopami arī citi dati par vēlēšanu rezultātiem, jo autori dažādi grupējuši politiskās partijas un tām pieslējušās grupas.)

Vēl atzīmējams, ka J.~Pilsudska piekritēji bija gandrīz visu politisko partiju ievēlēto deputātu vidū, tāpēc partiju iezīmētās politiskās vadlīnijas bieži tika pārkāptas.

Komunisti, uzskatot, ka valsts atrodas sociālistiskās revolūcijas priekšvakarā, parlamenta vēlēšanas boikotēja. Tikai 1921.~gadā, diviem deputātiem no citām partijām pārejot komunistu rindās, Seimā tika izveidota arī komunistu frakcija, kura par savu mērķi uzskatīja ``ekspluatatoru iekārtas'' atmaskošanu no parlamenta tribīnes.

Seima deputātu sastāvā tika iekļauti arī Vācijas Reihstāga un Austrijas Reihsrāta poļu deputāti. Laika gaitā, kad Polijai tika pievienotas jaunas teritorijas, deputātu skaits papildinājās papildus vēlēšanās.

Satversmes Seima uzdevums bija izstrādāt valsts konstitūciju, bet tā pilnvaru termiņš nebija noteikts. Pirms konstitūcijas izstrādāšanas Satversmes Seims 1919.~gada 20.~februārī pieņēma pagaidu pamatlikumu, t.s. \strong{``Mazo konstitūciju''}. Pēc tās visa varas pilnība atradās Satversmes Seima rokās. Tam bija pakļauta izpildvara~--- valsts priekšnieks un valdība. Valsts priekšnieks, kuru ievēlēja parlaments, un kura priekšā viņš bija atbildīgs, pārstāvēja valsti ārlietās, bija Seima lēmumu augstākais izpildītājs. Viņš nozīmēja valdību, saskaņojot to ar Seimu. Arī tā bija atbildīga Seima priekšā. Valsts priekšnieka izdotie dokumenti ieguva likumīgu spēku tikai pēc attiecīga ministra kontrasignācijas.

Nacionāldemokrātu un ļudoviešu ar V.~Vitosu priekšgalā, ar kuriem J.~Pilsudskim bija ne pašas labākās attiecības, stiprās pozīcijas Seimā nesolīja viņam vieglu sadarbību ar to. Valsts pagaidu priekšnieka piekritēji pat uztraucās, vai Seims apstiprinās viņa pilnvaras. Iespējams, tāpēc J.~Pilsudskis savas pilnvaras nolika nevis parlamenta pirmajā plenārajā sēdē 10.~februārī (kā tam bija jābūt pēc Reģentu padomes un E.~Moračevska valdības dekrētiem), bet tikai 10~dienas vēlāk. Šai laikā J.~Pilsudska piekritējiem izdevās pārliecināt citus deputātus saglabāt viņa rokās valsts priekšnieka posteni līdz nākamajām Seima vēlēšanām. 1919.~gada 20.~februārī Satversmes Seims arī apstiprināja J.~Pilsudski par valsts vadītāju (\pltxti{Naczelnik Państwa}), gan ar ierobežotu varu. Neizdevās nacionāldemokrātu mēģinājums vispār likvidēt valsts priekšnieka posteni, tā pilnvaras nododot Seima priekšsēdētājam, kas sastapa stipru pretestību no pilsudčiku un centra partiju puses (kopā 188~deputāti pret 109~nacionāldemokrātu pārstāvjiem). Tomēr J.~Pilsudskis, kurš formāli saglabāja savu amatu, no diktatora pārvērtās par parlamentāras republikas galvu. Vienlaikus viņš pildīja arī armijas virspavēlnieka pienākumus. Abi amati kopā deva J.~Pilsudskim ļoti plašas pilnvaras, lielu ietekmi, darbojoties ārpus parlamenta. Viņš pat nemēģināja parlamentā radīt sev uzticamu deputātu grupu, tradicionāli izmantojot kreiso partiju atbalstu, bet neidentificējot sevi ar to programmām un darbību, jo idejiski bija jau sarāvis ar tām sakarus. Pēc kāda literāta (\emph{Ā.~Novačinska}) vārdiem viņš ``izkāpa no partijas vilciena stacijā, kura saucas Neatkarība''.

Kaut labējās partijas bija vairākumā, droša absolūtā vairākuma tām trūka, tāpēc tās apstiprināja I.~Paderevska valdību. Nacionāldemokrātu līderis R.~Dmovskis tika iecelts par Polijas delegācijas vadītāju Parīzes miera konferencē. Savukārt Poļu nacionālās komitejas un Polijas delegācijas miera konferencē sastāvā tika iekļauti arī J.~Pilsudska piekritēju pārstāvji. Jaunās valsts izveides darbs spieda abus politiskos sāncenšus~--- J.~Pilsudski un R.~Dmovski~--- sadarboties.

1919.~gada vidū nacionāldemokrātu un labējo partiju pozīcijas Seimā nostiprinājās, jo guva panākumus Seima vēlēšanas tajos Polijas apgabalos, kur tās nevarēja notikt janvārī. Satversmes Seims apstādināja vēl I.~Dašinska valdības pasludinātās reformas, arī agrāro, spēkā palika vienīgi tās ieviestā 8~stundu darba diena rūpnīcās. Tā kā lielākā daļa valsts līdzekļu tika izlietota militārajiem mērķiem (valsts budžetā ārkārtējie, pirmkārt militārie, izdevumi sastādīja 73\%), rūpniecība nesaņēma nepieciešamā daudzumā investīcijas, iekārtas, kurināmo. 1919.~gada 7.~februārī Varšavā notika vispārējs streiks pret I.~Paderevska vadītās valdības politiku. Solidarizējoties ar to streikoja Dombrovskas baseina, Lodzas, Ļubļinas u.~c. pilsētu strādnieki. Parasti ekonomiskās prasības savijās ar politiskajām~--- pārtraukt strādnieku organizāciju vajāšanu, to vadītāju arestus, atbrīvot politieslodzītos u.c. Šķiru cīņa sasniedza tādu pakāpi, ka valdība ar 1919.~gada 1.~aprīli pasludināja ārkārtējo stāvokli visā bijušās Polijas karalistes teritorijā. Ar to karaspēku likumīgi varēja izmantot pret streikojošajiem strādniekiem vai zemnieku nemieru dalībniekiem.

Īpaši saasinājās stāvoklis 1919.~gada rudenī, kad Polijas pārtikas ministrija pati konstatēja, ka pilsētu apgāde ar maizi kļuvusi par īstu katastrofu. Valdība mēģināja ieviest valsts monopolu uz tirdzniecību ar labību, aizliegt tās izmantošanu spirta tecināšanai, taču muižnieki un turīgie zemnieki atteicās nodot labību par noteiktajām cenām, bet pārdeva to melnajā tirgū vai tecināja spirtu. Valdība bija spiesta padoties, atteikties no labības monopola un atļaut brīvo tirdzniecību. (Stāvokli būtiski atviegloja ASV. Pēc oficiāliem amerikāņu datiem no 1918.~gada 1.~decembra līdz 1919.~gada 31.~augustam no ASV uz Poliju tika nosūtītas dažādas kravas vairāk nekā 122~miljonu dolāru vērtībā.)

Protams, trūcīgie pilsētu iedzīvotāji badojās. Sakarā ar pārtikas trūkumu strādnieku un policijas sadursmes, kurās bija cilvēku upuri, septembrī notika Lodzā, oktobrī Krakovā un tās apkaimes pilsētās. Streikoja ogļrači Dombrovskas baseinā, dzelzceļnieki, kara rūpniecībā nodarbinātie, laukstrādnieki u.c. Daudzi no streikiem beidzās ar sadursmēm pret žandarmēriju un policiju.

Lai nomierinātu zemnieku kustību, valdība iesniedza Seimam agrārās reformas projektu. Tās nepieciešamību neapšaubīja neviena partija, taču par reformas principiem izraisījās karstas debates. Ļudovieši vēlējās paātrināt reformas īstenošanu, zemes īpašuma ierobežošanu ar 50--150~ha, baznīcu un klosteru zemju atsavināšanu, mazu atlīdzību par atsavināmo zemi, mežu nacionalizāciju. Pašu ļudoviešu vidū arī pastāvēja domstarpības. To kreisas spārns vēlējās panākt radikālāku reformu. Labējais spārns~--- ``\pltxti{Piast}'' un citi grupējumi~--- uzstādīja mērenākas prasības. Tiesa, gan ``\pltxti{Wyzwolenie}'', gan ``\pltxti{Piast}'' vienmēr izrādīja gatavību saglabāt lielo zemes īpašnieku rokās daļu no to valdījumiem.

Nacionāldemokrāti (endeki) centās panākt, lai zemes maksimums būtu ne mazāks par 500--750~ha, augstu atlīdzību par atsavināmo zemi, baznīcas un klosteru zemes īpašuma saglabāšanu. Sociālisti iestājās par lielo un vidējo zemes īpašumu atsavināšanu un pārvēršanu par sabiedrisku īpašumu. Garīdzniecība enerģiski uzstājās pret tīkojumiem atsavināt baznīcas un klosteru zemes (pāri par 200~000~ha). Galu galā 1919.~gada jūlijā Satversmes Seims tikai ar vienas balss vairākumu (ar 183 pret 182 balsīm) pieņēma likumu par zemes īpašumu izmēru ierobežošanu. Rūpnieciskajos rajonos zemes maksimums tika noteikts 60~ha, bet lauksaimnieciskās teritorijās~--- 180~ha. Pie tam rietumu apgabalos, kuri agrāk ietilpa Vācijas sastāvā, un austrumu teritorijās, kuras apdzīvoja lietuvieši, baltkrievi un ukraiņi, zemes maksimālais daudzums tika palielināts līdz 400~ha. Par zemes atsavināšanu tika noteikta atlīdzība~--- 1/2~no tās vērtības. Jauno zemes īpašumu izpirkšanai tika paredzēts kredīts 75\% apmērā no to vērtības. Ik gadu tika paredzēts atsavināt un sadalīt ne mazāk kā 100~000~ha. Taču likums neparedzēja, ne kā praktiski varētu tikt atsavinātas liekās zemes platības, ne kā tās tiks turpmāk izlietotas.

No poļu politiķu viedokļa visaktuālākais bija \strong{jautājums par jaunās valsts robežām}. Starptautiskā situācija bija labvēlīga poļu radikālo nacionālistu mērķiem. Ja Krievijā pēc Oktobra apvērsuma nesāktos Pilsoņu karš, Polijai gandrīz nebūtu izredžu iegūt zemes uz austrumiem no bijušās Polijas karalistes. Lai cik Rietumvalstīm nepatika boļševiku vara Krievijā, tūlīt pēc Pirmā pasaules kara tās nebija spējīgas tik aktīvi iesaistīties Krievijas notikumos, lai Polija, to izmantojot, varētu sagrābt ukraiņu un baltkrievu apdzīvotās teritorijas. Šķelšanās un Pilsoņu karš Krievijā deva Polijai izdevību pretendēt uz kādreiz Žečpospolitas īpašumā bijušajām zemēm. ``Sarkanie'' un ``baltie'' cīnījās galvenokārt par kontroli valsts centros~--- Petrogradā un Maskavā, perifērijai veltījot mazāku uzmanību. Ja perifēriju boļševiki gaidīja palīdzību no krievu boļševikiem, tad vietējie nacionālisti, nobijušies no boļševikiem, meklēja pret tiem palīdzību kur vien varēja: gan pie Antantes valstīm un ASV, gan jau sakautās Vācijas, gan arī pie Polijas, kurā salīdzinoši ātri veidojās jauna valsts, kas ieguva starptautisku atzīšanu. Līdz 1918.~gadam uzkrātā vēsturiskā pieredze liecināja, ka mazās valstis nevar realizēt patstāvīgu ārpolitiku, bet savu drošību var panākt vienīgi pateicoties lielvalstu aizbildniecībai, maksājot par to ar daļu savas suverenitātes. To labi saprata poļu politiķi. Taču viņi vēlējās panākt arī ko vairāk nekā poļu etnisko teritoriju neatkarību.

Vēl nepaguvusi piedzimt kā suverēna valsts, vēl neieguvusi savā kontrolē poļu apdzīvotās teritorijas, Polija jau uzstādīja \strong{prasības par nepoļu zemju pievienošanu}. Nākas konstatēt, ka poļu vairākums neatzina to pašsaprotamo faktu, ka XX~gadsimta sākumā ukraiņi, baltkrievi, lietuvieši, latvieši jau apzinājās savu nacionālo oriģinalitāti un vēlējās radīt savas valstis. Žečpospolita viņiem asociējās ar šļahtas zemes īpašumu un pilnīgu poļu vienvaldību.

Poļu nacionālā komiteja Parīzē 1918.~gada 12.~oktobrī, vēl pirms Vācijas kapitulācijas, iesniedza Sabiedroto valstu valdībām memorandu ar prasību Francijā dislocētās poļu t.s. ``Hallera armijas'' spēkiem okupēt Kameņecas-Podoļskas, Brestļitovskas un Kauņas apgabalus. ``Šī okupācija~--- skanēja memorands,~--- garantētu Polijai drošību austrumos un varētu kalpot par nākotnes bāzi sabiedroto kara operācijām Krievijā''. Poļu politiķi netaisījās ievērot to, ka ukraiņu, lietuviešu, bet ar Pirmā pasaules kara sākumu arī baltkrievu politiskie grupējumi cīnījās par savu nacionālo valstiskumu un nebūt negatavojās saistīt savu likteni ar poļiem. Poļu nacionālā komiteja Parīzē arī bija tā, kas prasīja Polijas valstij pievienot poļu apdzīvotos Vācijas apgabalus, kad ne to iedzīvotāji, ne arī J.~Pilsudskis vēl neizrādīja tādu iniciatīvu.

1918.~gada 2.~novembrī sabiedroto sanāksmē Francijas ārlietu ministrs S.~Pišons izklāstīja Francijas vēlmes nākamo Polijas robežu jautājumā. ``Es gribētu stingri iestāties lai zem [vācu] evakuējamajām teritorijām tiktu saprastas visas tās teritorijas, kuras sastādīja Polijas karalisti līdz pirmajai sadalei 1772.~gadā''. Uz to Lielbritānijas ārlietu ministrs A.~Balfūrs paziņoja: ``Es izdzirdu šo priekšlikumu ar bažām. Jūs sakāt, ka 1918.~gada Polijai ir jābūt 1772.~gada Polijai. Tas nav tas, pēc kā mēs tiecāmies un ko mēs apsolījām. Mēs solījām atjaunot Poliju, kuru apdzīvotu poļi. 1772.~gada Polija neatbilst šim mērķim, tā nesastāvēja tikai no poļiem. Tajā bija iekļautas arī nepoļu teritorijas, tai laikā kad poļu teritorijas sastādīja tikai tās daļu. Tādejādi šī formula cieš gan no nepilnības, gan no pārspīlējuma. Precīza robežu delimitācija [starptautiska vienošanās par valstu savstarpējām robežām un to iezīmēšana] ir tik sarežģīta lieta, ka es jūs lūdzu neienest to pamiera noteikumos [ar Vāciju].'' ASV prezidenta pārstāvis pulkvedis E.~Hauzs prezidenta V.~Vilsona vārdā paziņoja, ka viņš pievienojas A.~Balfūra priekšlikumam. S.~Pišonam nācās piekāpties. Tēze par 1772.~gada Polijas robežām neiekļuva Kompjenas pamiera noteikumos.

Antantes valstu uzskatus par Polijas vietu pasaulē galvenajos vilcienos izteica maršals F.~Fošs 1919.~gada 15.~septembrī, norādot, ka Polijas būtu jābūt ``no vienas puses, barjerai pret boļševismu, bet, no otrs puses, — pret Vāciju''. Šos noskaņojumus Antantes valstīs Polijas diplomāti centās maksimāli izmantot, izceļot poļu nopelnus un noniecinot citu kaimiņu tautu sasniegumus ekonomikā un kultūrā.

\strong{Kādas teritorijas tad bija jāaptver jaunajai Polijas valstij?} Cīņa jautājumā par nacionālo valsti un tās robežām atspoguļoja poļu politiķu uzskatus par Polijas ģeopolitisko vietu Eiropā.

Vai ``poļu zemes'', par kurām tik daudz runāja politiķi, bija jāsniedzas līdz Dņeprai un pat Melnajai jūrai un jāietver arī Silēzija un Baltijas jūras piekraste ar Dancigu? Poļu sabiedrībā nepastāvēja stingru kritēriju jautājumā par to, ko saukt par ``poļu'' zemēm. 1919.~gadā 1.~janvārī poļu karavīru avīze rakstīja: ``Mūsu tēvzeme, Polija, ir plaša, skaista un bagāta zeme. No Tatru kalnu stāvajām virsotnēm dienvidos līdz Baltijas jūras zilganajiem ūdeņiem ziemeļos, no Odras augšienes līdz Daugavai un Dņeprai sniedzas apgabali, ko mēs saucam par Poliju. Ne visus tos apdzīvo vienīgi poļi \citespace{} taču pat neatkarīgi no ciešajām kopīgajām tradīciju saiknēm visas provinces mums jāuzskata kā ar mums cieši saistītām, kurās, mazākais, daļa iesakņoto vietējo iedzīvotāju savu piederību poļu nācijai pierāda vārdos un darbos. Šādus iedzīvotājus mēs atrodam kā Pomerānijā [Piejūrā~--- \pltxti{Pomorze}] tā arī Baltkrievijā un Ukrainā. Un tāpēc mēs uzskatām šos apgabalus par poliskiem.''

Teorētiski jautājumā par Polijas robežām pastāvēja trīs galvenie varianti:

\begin{enumerate}

\item poļu nacionāla valsts aptvertu tikai etniskās poļu apdzīvotās teritorijas (kaut, ņemot vērā ebreju diasporas eksistenci, arī šāda valsts nevarētu būt pilnīgi etniski viendabīga);

\item valsts iekļautu sevī plašākas teritorijas nekā tikai etnisko poļu apdzīvotās, bet tajās būtu jābūt ievērojamam poļu īpatsvaram, kas nodrošinātu priekšrocības poļiem un radītu apstākļus to sekmīgai nacionālai attīstībai;

\item daudznacionāla valsts līdzinātos Žežpospolitai līdz tās sadalei, kur vienota valstiskuma princips būtu pārsvarā pār nacionālu valstu izveides principu.

\end{enumerate}


% page 363


Reāli Polijā politiķu un arī vienkāršo iedzīvotāju vidū jau līdz Polijas neatkarības proklamēšanai pastāvēja un savstarpēji konkurēja tikai divas koncepcijas~--- divi kā pēdējie minētie priekšstati par nākamo valsti, kuri cieši saistījās ar poļu darbinieku jau realizējamo politiku. Pirmo no tiem pārstāvēja nacionāldemokrātu līderis R.~Dmovskis, otro~--- J.~Pilsudskis. Tos vienoja viens~--- abi \strong{Polijas valsti saskatīja arī ārpus poļu vairākuma apdzīvotajiem apgabaliem}. Abi poļu politiķi uzskatīja, ka Polijas sastāvā noteikti jāietilpst Viļņas apgabalam (liet. \lttxti{Vilniaus kraštas}, poļu \pltxti{Wileńszczyzna}, baltkr. \betxti{Віленшчына}, krievu \rutxti{Виленский край, Вильнюсский край}), Rietumbaltkrievijai, Galīcijai un Rietumu Volīnijai. Atšķirīga pieeja bija tikai par uz austrumiem no nosauktajiem esošo apgabalu likteni.

Poļu uzskatos šo apgabalu pievienošana Polijai būtu atlīdzība par viņiem nodarīto ``netaisnību'', trijām Žečpospolitas dalīšanām utt. Šajā aprēķinā cieši savijās dotā reģiona nacionālā sastāva vērtējums ar vēsturiskās atriebības ideju Krievijai~--- uzsver krievu vēsturnieks V.~Pihalovs~--- Polijas saprašanos ar Padomju Krievijas taču apgrūtināja ne tikai atšķirīgās valdošās ideoloģijas, bet arī visa agrākās Krievijas impērijas un Polijas attiecību vēsture, kā arī Polijā, īpaši tās inteliģencē, ilgstoši izplatītā ideja par ``Krievijas vēsturisko vainu'' Polijas priekšā. Šī ideja bija dziļi iesakņojusies nacionāli noskaņotajā poļu inteliģencē. To pēc 1915.~gada, kad Krievijas karaspēks tika padzīts no poļu zemēm, atbalstīja arī vācu un austroungāru propaganda. Poļu pilsoniskajiem politiķiem šī ideja kalpoja viņu pretenziju attaisnošanai pret saviem kaimiņiem. Abu uzskatu pārstāvji it kā ``aizmirsa'', ka strīdīgajos apgabalos gan lietuviešiem, ukraiņiem un baltkrieviem arī bija tiesības uz pašnoteikšanos.

J.~Pilsudskis pārstāvēja t.s. \strong{``jagellonisko'' priekšstatu}, kas kā uz pamatu jaunās Polijas robežām balstījās uz Jagelloņu dinastijas pārvaldītajām Polijas-Lietuvas valsts teritorijām pirms tās dalīšanas. J.~Pilsudskis un viņa piekritēji vadījās no fakta, ka teritorijas starp Poliju un Krieviju apdzīvoja nekrievu tautas (baltkrievi, ukraiņi, lietuvieši, ebreji u.c.), bet starp tām bija poļu ielāsmojumi, kas viņiem kalpoja par pamatojumu uz pretenzijām pēc 1772.~gada robežām. Kā norādījis mūsdienu baltkrievu vēsturnieks I.~Meļņikovs, jagelloņu ideja lietuviešu, ukraiņu, baltkrievu acīs bija tradicionālā poļu imperiālisma forma. (Pēc būtības Polijas ``vēsturiskās tiesības'', vēlmi iekļaut tās sastāvā arī 1772.~gada robežās esošās ukraiņu, baltkrievu un lietuviešu zemes varēja tikt balstītas ne uz iedzīvotāju nacionālo sastāvu, bet gan vienīgi uz zemes īpašuma tiesībām. 1916.~gadā poļu muižnieku zemes īpašums Ukrainā, Baltkrievijā un Lietuvā sastādīja 15~miljonus morgu jeb 8,4~miljonus ha, kas veidoja no 50~līdz 70\% visa zemes īpašuma dažādos apriņķos.)

Galvenās briesmas Polijai J.~Pilsudskis saskatīja Krievijā. Tāpēc tā bija jānovājina, atraujot no tās nacionālās nomales, radot virkni nacionāli-valstisku veidojumu, kuri atrastos Polijas ietekmē. J.~Pilsudskis saprata, ka XVIII gadsimta beigu Žečpospolitas iekārtas atjaunošana nebija iespējama, tāpēc tika izstrādāts plāns poļu, lietuviešu, ukraiņu un baltkrievu apdzīvotos apgabalus apvienot federatīvā valstī, protams, poļu vadībā. Ja tas neizdotos, būtu jāveido valstu federācija, protams, Polijas vadībā. Jo ilgāk Krievijā turpinātos sajukums, jo lielākas teritorijas varētu savā kontrolē iegūt Polija. Tas ļautu Polijai kļūt par lielvalsti, kura Austrumeiropā ieņemtu agrāko Krievijas vietu. Pēc Pirmā pasaules kara izveidojies nosacītais varas vakuums Austrumeiropā, Vācijas un Krievijas novājināšanās, arī poļu elites Polijas izdevīgā stāvokļa Centrāleiropā apzināšanās, tās pārākuma apziņa pār citām reģiona valstīm deva bagātīgu augsni federālistiskām un reizē hegemoniālām idejām. ``Jagellonisko'' ideju īstenošana, kā norādījis poļu vēsturnieks J.~Kirhmeijers, ``gala rezultātā nozīmēja noteiktu Krievijas atbīdīšanu uz austrumiem''. ``Jagellonisko'' priekšstatu atbalstīja konservatīvie poļu politiķi, PPS, vadošie militāristi un arī lielākā daļa nacionāli noskaņotās poļu inteliģences.

R.~Dmovskis un viņa piekritēji nacionāldemokrāti par galveno draudu Polijai uzskatīja Vāciju, aizstāvēja t.s. \strong{``pjastu'' priekšstatu}, kas paredzēja atjaunot Pjastu dinastijas valsti, kas austrumos aprobežotos tikai ar pārsvarā poļu apdzīvotajiem apgabaliem (sākotnēji tika runāts, ka austrumos varētu palikt 1815.~gadā noteiktā robeža), bet rietumos uz Vācijas rēķina Polijai būtu jāpievieno zemes līdz Oderai, kur agrāk X gadsimtā, kad veidojās poļu valsts, dzīvoja slāvi. Daudzi poļu vēsturnieki, arheologi, ģeogrāfi, valodnieki, publicisti un rakstnieki jau pirms Pirmā pasaules kara bija pielikuši daudz pūļu, lai pierādītu, ka Polijas austrumos esošie Prūsijas rietumu apgabali: Pozene (vācu \detxti{Posen}, poļu \pltxti{Poznań}), Ermlande (vācu \detxti{Ermland}, poļu \pltxti{Warmia)}, Rietumprūsija (vācu \detxti{Westpreussen}, poļu \pltxti{Prusy Zachodnie}), Silēzija, (vācu: \detxti{Schlesien}, poļu \pltxti{Śląsk)}, Pomerānija (vācu \detxti{Pommern}, poļu \pltxti{Pomorze}), Mazūrija (poļu \pltxti{Mazury}; vācu \detxti{Masuren}), daļēji arī Brandenburga (vācu \detxti{Brandenburg}) ir senas poļu zemes, kuras agrākajos gadsimtos galvenokārt Prūsijas pūliņiem daļēji ģermanizētas un ietilpinātas vācu kultūras lokā. Ģermanizācija tika tēlota vienīgi kā piespiedu kārtā veikta, tātad kā netaisnība. Taču no poļu viedokļa tā nespējusi mainīt minēto zemju polisko raksturu. (Protams, vācu zinātnieki, kas apkalpoja vācu toreizējās nacionālās intereses, neatzina poļu puses argumentus.)

R.~Dmovskis jau 1917.~gada martā iesniedza Lielbritānijas Ārlietu ministrijas valsts sekretāram A.~Balfūram notu, kurā bija norādīts, ka Polijas robežas nedrīkst būt nedz vēsturiskas (tās būtu pārāk plašas), nedz etniskas (pārāk šauras). Pēc R.~Dmovska viedokļa Polijas teritorijai bija jāiekļauj 1)~Austrijas Poliju~--- Galīciju un daļu Austrijas Silēzijas, 2)~Krievijas Poliju~--- Polijas karalisti un Kauņas, Viļņas, Grodņas guberņas, daļu Minskas guberņas un Volīniju, 3)~Vācijas Poliju~--- Pozeni (Poznaņu), Rietumprūsiju ar Dancigu (Gdaņsku), Augšsilēziju un Austrumprūsijas dienviddaļu.

Nacionāldemokrātu ideāls bija nacionāli viendabīga valsts, taču sekošana šim ideālam konkrētajos apstākļos nozīmētu piekrist modificētai Napoleona laiku Varšavas hercogistes versijai, kam viņi 1919.~gadā nebija gatavi. Izeja tika atrasta valsts daudznacionālā rakstura pagaidu saglabāšanā. Tālāk tika spriests sekojoši: ja no etniskās daudzveidības nevar izvairīties, tad vismaz vajag to ierobežot, forsēti polonizējot nacionālos mazākumus un izolējot nepakļāvīgos elementus. Pēc R.~Dmovska uzskatiem, poļu valstij nebija jābūt etniski homogēnai, tajā līdzās etnogrāfiski poliskām teritorijām, jāiekļauj arī rajoni, kur poļu elements bija redzami pārstāvēts. 1919.~gadā Poļu nacionālās komitejas sēdē R.~Dmovskis atzina, ka visvājākais ``poļu teritorijas punkts ir zemes, kuras atrodas uz austrumiem no Brestas. Šī teritorija, tāpat kā Grodņas guberņa, ir puse, kur kompakti dzīvo nepoļu iedzīvotāji, taču neraugoties uz to, ievērojot ģeogrāfisko faktoru, šo zemi nevar Polijai nogriezt''. Viņš piekrita arī tam, ka Austrumgalīcijā poļu iedzīvotāji ir mazākumā, taču šo rajonu pievienošanu Polijai attaisnoja ne tikai ar stratēģiskiem, bet arī ideoloģiskiem argumentiem~--- šeit dzīvojošo tautu it kā nepietiekamo gatavību izveidot savu valstiskumu. R.~Dmovskis vēsturē saskatīja tautu nevienmērīgu nobriešanu, ko poļi dotajā situācijā varētu izmantot savā labā. R.~Dmovskis nenoliedza ukraiņu un baltkrievu nacionālos centienu pastāvēšanu, taču, pēc viņa domām, šīm tautām bija jātiek polonizētām, jo tās atpalika no poļiem nacionālajā un kultūras briedumā. Kopumā endeku doktrīna bija pretrunā kā ar liberāli-demokrātiskajām un sociālistiskajām, tā arī konservatīvo tradīcijām uzskatos par nacionālu valsti.

Pēc ``jagelloņu'' varianta realizēšanas Polijai vajadzēja rēķināties galvenokārt ar Krievijas naidīgumu, pēc ``pjastu’~--- ar Vācijas.

Abu~--- gan ``pjastu'', gan ``jagelloņu'' uzskatu aizstāvji bija atbildīgi par to, ka jaunā poļu valsts nekavējoties sabojāja attiecības ar gandrīz visām citām Austrumeiropas nacionālajām kustībām. Jau pirmie jaunās Polijas soļi rādīja, ka kā plānojamā ``jagelloniskā'' federācija, tā pārsvarā ``pjastu'' tradīcijas turpinošā poliskā valsts var balstīties tikai uz poļu pārākumu un spēku. Polija praktiski pārvērtās par agresīvu valsti, kura militārām metodēm uzsāka savas teritorijas paplašināšanu uz sabrukušo impēriju drupām. Par Polijas ārpolitiku ironizējis krievu žurnālists O.~Jalovenko: ``Acīmredzot, kad runa sākas par Polijas robežu paplašināšanu jebkurā virzienā, Varšavā sāk darboties ``hotentotu morāle'': ``Kas ir labi? Tas ir tad, kad mēs atņemam govis zulusiem. Bet kas ir slikti? Tas ir tad, kad zulusi atņem govis mums.''

Atšķirībā no nosauktajiem nacionālās teorētiskās domas un praktiskās politikas virzieniem, poļu sociālisti un augošās zemnieku kustības līderi (ļudovieši) valstisko sakaru attīstībā saskatīja pamatfaktoru, kuram bija jāvieno visus Polijas pilsoņus, neatkarīgi no reliģiskajiem uzskatiem un nacionālās piederības. Tieši Polijas pilsonībai, pēc viņu domām, bija jākļūst par integrējošu spēku, kas vienotu dažādās iedzīvotāju etniskās grupas. Perspektīvā tam bija jānoved pie vienas ``valsts nācijas'' (Mūsdienās to parasti sauc par ``politisko nāciju''~--- V.Š.) izveides. Taču, neraugoties uz savu pievilcību, pēdējai koncepcijai bija utopisks raksturs, tā neievēroja poļu nacionālisma uzplūdus, politiķu vēlmi tos izmantot savā labā, citu nāciju pretdarbību, spēku samēru starp tautām, kuras dzīvoja agrākajās Žečpospolitas robežās.

J.~Pilsudskis kā valsts vadītājs un I.~Paderevskis kā valdības vadītājs ap sevi vienoja federatīvas valsts izveides piekritējus, kurā Polija ar sev pievienotām poļu apdzīvotajām agrāk Vācijai piederējušajām zemēm kļūtu par hegemonu valstu blokam, kurā bez Lietuvas un Baltkrievijas varētu ietilpt arī pārējās Baltijas valstis (Latvija un Igaunija) un Ukraina. Tikai šo valstu piekrišana nebija prasīta.

Pašā Polijā šī programma sastapa asus endeku iebildumus, kuri neticēja šādas starptautiskas federācijas iespējai. R.~Dmovskis uzskatīja, ka varbūtējie federācijas partneri austrumos ir pārāk vāji, lai sekmīgi aizstāvētos pret poļu hegemoniju, savukārt Polija ir pārāk vāja, lai nostātos Austrumeiropas valstu federācijas priekšgalā. Taču nacionāldemokrātu alternatīvā programma bija vēl vairāk nacionālistiski-šovinistiska, paredzēja Viļņas un Ļvovas pievienošanu Polijai, bez kaut kādas piekāpšanās lietuviešiem un ukraiņiem. Citiem vārdiem~--- agrāk apspiestajiem poļiem pašiem bija jākļūst par apspiedējiem. Kā norādījis krievu vēsturnieks~--- polonists G.~Matvejevs, Polija centās savā valstī iekļaut nozīmīgas austrumu nomales teritorijas, kas tai piederēja pirms 1772.~gada. ``Polijas mērķis bija atdalīt no Krievijas daļu bijušo poļu teritoriju un tādējādi vājināt šo valsti'', bet ``jautājums kā to izdarīt: saskaņā ar Dmovska inkorporācijas vai Pilsudska federālas iekļaušanas politiku, bija sekundārs''.

No 1919.~gada 18.~janvāra līdz 1920.~gada 21.~janvārim Parīzē norisinājās miera konference, kurā piedalījās Francija, Lielbritānija, ASV, Itālija, kā arī citas valstis, kuras bija karojušas Antantes pusē vai vismaz sarāvušas diplomātiskās attiecības ar Vāciju un tās sabiedrotajām valstīm. Tikai Padomju Krievija nebija uzaicināta.

Sarunās piedalījās arī pilntiesīga Polijas delegācija. Par pamatu Polijas uzaicināšanai uz konferenci kalpoja fakts, ka poļu daļas taču no 1917.~gada bija cīnījušās Francijā, karojušas arī Krievijas armijas sastāvā. To, ka daudzi poļi paši pēc savas gribas poļu leģionos bija sparīgi karojuši arī Austroungārijas un Vācijas pusē, R.~Puankarē, atklājot konferenci, vienkārši noklusēja. Polijas delegāciju vadīja Rietumu sabiedroto uzticību baudošais R.~Dmovskis, kurš propagandēja jaunās valsts izveidi maksimāli plašās robežās. Tiesa, viņš rēķinājās ar tām izmaiņām, kuras bija notikušas 150~gados pēc trešās Žečpospolitas dalīšanas. Tāpēc Polija atteicās no 1772.~gadā tajā ietilpstošas Kijevas guberņas, gandrīz visas Podolijas, lielākās daļas Volīnijas, visas Mogiļevas, trešdaļas Minskas guberņas, diviem Viļņas guberņas povjatiem un visas Vitebskas guberņas~--- kopumā no 311~007 km$^{2}$ (ar 10,6~miljoniem iedzīvotāju) no tiem 758~144 km$^{2}$, kas ietilpa Žečpospoļitā 1772.gadā.

Toties atbilstoši ``pjastu'' variantam Polijas robežās bija jāiekļauj Lietuvu un agrāk Vācijai piederošās zemes: Rietumprūsiju ar Dancigu, Cešinas (vācu \detxti{Teschen}, poļu \pltxti{Cieszyn}, čehu \cstxti{Těšín}) apgabalu un Augšsilēziju. Austrumprūsija bija jāatdala no Vācijas. Kopā no Vācijas bija jāatsavina 84~198 km$^{2}$ teritorijas. Pēc šī projekta Polijas teritorija būtu 481~386 km$^{2}$ ar 36,8 miljoniem (tai brīdī) iedzīvotāju.

Polijas delegācijas plānus enerģiski atbalstīja Francija. (Savos memuāros R.~Dmovskis nepārprotami atzina Francijas lomu Polijas valsts izveides sākumperiodā: ``Par visu, ko mēs sasniedzām, \citespace{} mums pirmkārt jāpateicas Francijai.'') Francija vēlējās radīt spēcīgu Polijas valsti, kura Eiropas austrumos varētu pildīt Francijas sabiedrotā lomu kā pret Vāciju, tā arī Padomju Krieviju. Tas, kā šķita Francijas politiķiem, kalpotu kā Francijas hegemonijas balsts Eiropā. Bet tieši tas izsauca Lielbritānijas, kura nevēlējās pārlieku Francijas nostiprināšanos, sīkstu pretošanos. Lielbritānija, kura tradicionāli centās uzturēt Eiropā spēku balansu, baidījās arī no jaunu konflikta apgabalu kā Elzasa un Lotringa rašanās, ja tiktu apmierinātas visas poļu vēlmes. Lielbritānijas premjers D.~Loids Džordžs cīnījās par Polijas spēku ierobežošanu un tautas nobalsošanu Austrumprūsijā un Augšsilēzijā. Viņš arī pamatoti norādīja: ``Es neredzu tos lielos nopelnus, kurus Polija ieguldījusi Vācijas sakāvē, tai pašā laikā šis valsts prasības ir kā galvenajiem uzvarētājiem''.

Zināmas cerības poļu delegācija saistīja arī ar ASV, kaut to prezidents V.~Vilsons iestājās par jaunu valstu radīšanu Eiropā uz etniskajiem pamatiem un nomaļu apgabalu (bez Polijas) atstāšanu Krievijas sastāvā.

Tā Parīzes konferencē starp Franciju, no vienas puses, un Lielbritāniju, no otras puses, norisa cīņa par Polijas vietu Eiropā, par tās rietumu un austrumu robežām, Silēzijas un Dancigas piederību. Neraugoties uz Francijas premjera Ž.~Klemanso pūliņiem, Francijai neizdevās īstenot Polijas 1772.~gada robežās atjaunošanas programmu.

Kā raksta V.~Kovaļskis u.c. poļu vēsturnieki, bijusi arī psiholoģiska nianse, kura apgrūtinājusi R.~Dmovskim viņa plānu īstenošanu. Lieta tā, ka R.~Dmovskis vadīja partiju, kurai bija raksturīgs izteikts antisemītisms. Toties Lielbritānijas premjers bijis cieši saistīts ar savas valsts ebreju kapitālu. R.~Dmovskis nespējis attaisnoties par ebreju grautiņiem Polijā. Autoram gan šķiet, ka poļu vēsturnieki pārvērtē šīs ``nianses'' lomu. D.~Loids Džordžs, protams, ņēma vērā gan sabiedrības negatīvo attieksmi pret ebreju grautiņiem, gan dažādu kapitāla grupu viedokļus, taču galvenokārt vadījās no savas valsts interesēm. Lielbritānija parasti realizēja līdzsvara politiku Eiropā un tai nebija vēlama ar Polijas atbalstu sasniegta Francijas dominance kontinentā. Ne velti 1919.~gada 2.~martā Polijas Nacionālās komitejas sēdē Parīzē, kur tika apspriestas Polijai vēlamās robežas un to iegūšanas ceļi, tuvais J.~Pilsudskim profesors A.~Sujkovskis izteica domu, ka, ievērojot Anglijas vēlmi ``ierobežot Polijas apetīti'', Polijas prasība pievienot baltkrievu apdzīvotās teritorijas var izsaukt apvainojumus imperiālismā. Šis pats apvainojums draudot arī no amerikāņu puses, kuri etnogrāfiskos apstākļus uzskatot par izšķirošajiem. Lai pierādītu miera konferencei, ka poļi neesot imperiālisti, esot jāpieņem federācijas princips. Tā tikai ``Lietuvas pasludināšana par neatkarīgu, lai arī tikai skata pēc neatkarīgu, noņemtu no Polijas apvainojumus imperiālismā'' un ``ļautu mums [poļiem] iespiesties aiz Dņepras un Daugavas ar mērķi radīt milzīgu stratēģisku Poliju, kura, kā tagad tiek runāts, aizstāvētu Eiropu no boļševikiem. Citu jautājuma nostādni Eiropa, dotajā gadījumā Antante, var nepieņemt.''. Atspoguļojot ``jagelloniskā'' un ``pjastu'' valstiskuma idejas aizstāvju cīņas, A.~Sujkovskim iebilda R.~Dmovskis, kurš uzsvēra, ka federācijas izveidei ir nepieciešama spēja piekrist kompromisam, bet Lietuva un pārējās tautas ``austrumos'' no Eiropas tautām esot vismazāk spējīgas kompromisam, tāpēc vēlme nodibināt ar tām federāciju tikai radītu nekārtības un anarhiju. Valsts robežām gan esot jābūt ``cik iespējams plašām'', taču tās nedrīkstot pārsniegt robežas, kur Polijas Seims pārstātu būt nacionāls, citādi arī valsts pārstātu būt nacionāla un pārtrauktu realizēt ``poļu ideju un poļu politiku''. R.~Dmovskis lūdza atmest federācijas ideju, cīnīties tikai par tādu apgabalu pievienošanu Polijai, kas ļautu poļiem saglabāt nacionālo pārsvaru. Viņš paziņoja: ``Protams, labi būtu saņemt visu [vēlamo], taču tas nav iespējams. Neatkārtosim to kļūdu, kura noveda pie Krievijas valsts bojā ejas. Krievijai bija raksturīgs, ka apetīte tai bija lielāka, nekā kuņģis: tā aiz negausības apēda vairāk nekā spēja sagremot.'' Neapsteidzot notikumus, šeit tikai jāpasaka, ka nākotne parādīja R.~Dmovska brīdinājuma Polijai pamatotību.

Parīzes miera konferencē Polijas delegācija pēc būtības centās atrisināt tikai divus jautājumus: panākt Polijas kā suverēnas starptautiskās politikas subjekta atzīšanu un noteikt tās robežas ar Vāciju. Visus citus robežjautājumus Varšava cerēja atrisināt ar pašu spēkiem. Jaunā Polijas valsts paļāvās uz savu armiju, kura lēni, bet nepārtraukti virzījās uz austrumiem, tālu izejot ārpus poļu etniskajām robežām. Parīzes miera konference, pasaules sabiedrība tika nostādīta fakta priekšā.

Parīzes miera konference izstrādāja miera līgumus ar Vāciju (Versaļas miera līgums, parakstīts 1919.~gada 28.~jūnijā), kā arī ar Austriju (Senžermēnas miera līgums, parakstīts 1919.~gada 10.~septembrī), Ungāriju (Trianonas miera līgums, parakstīts 1920.~gada 4.~junijā), Bulgāriju (Neijī līgums, parakstīts 1919.~gada 27.~novembrī) un Turciju (Sevras līgums, parakstīts 1920.~gada 10.~augustā), kuri lika pamatus pēckara t.s. Versaļas--Vašingtonas sistēmai pasaulē.

Parīzes miera konference arī nodibināja Tautu Savienību (angļu \entxti{League of Nations}, franču \frtxti{Société des Nations}, krievu \rutxti{Лига наций}. Savienība oficiāli darbību uzsāka 1920.~gada 10.~janvārī). Savienības mērķis bija iespējamo karu novēršana, radot kolektīvo drošību un veicot vispārēju atbruņošanos, konfliktu atrisināšana sarunu ceļā un globāla dzīves kvalitātes uzlabošana. Versaļas miera līgumu parakstīja arī Polijas pārstāvji R.~Dmovskis un I.~Paderevskis.

Atbilstoši Versaļas līguma 116.~pantam Vācija atzina ``visu teritoriju, kuras 1914.~gada 1.~augustā ietilpa Krievijas impērijas sastāvā'', neatkarību. Līguma 117.~pants apšaubīja padomju režīma leģitimitāti Krievijā un uzlika Vācijai par pienākumu ``atzīt visus līgumus ar valstīm, kuras bija radušās vai rodas visā Krievijas impērijas teritorijā vai tās daļās.'' Versaļas līgumu Vācija uzskatīja par tās sakāves simbolu un pazemojumu, paziņojot, ka tas paredz ``nedzirdētu netaisnību'', tiek tai uzspiests ar spēku. Vācijas valdība to parakstīja ļoti negribīgi. Versaļas līguma noslēgšanas rezultātā Polija ieguva sev ievērojamu ienaidnieku rietumos~--- Vāciju. Polija ar to rēķinājās. Lūk, kā stāvokli vērtēja poļu vēsturnieks J.~Kozenskis: ``Versaļas līgums ne tikai neapmierināja nevienu no pusēm, bet radīja bīstamu un neizdevīgu situāciju Polijai. Pastāvīgie draudi Gdaņskai [Dancigai], Polijas rietumu robežām lika Polijai pastāvīgi būt modrai līdz pat Otrā pasaules kara sākumam.''

Austrija un Ungārija bez īpašas pretestības piekrita Polijas u. c.~valstu robežām, kuras noteiks sabiedrotie, atsakoties no XVIII un XIX gadsimtā sagrābtajām teritorijām. Polija savukārt uzņēmās izmaksāt daļu Krievijas parādu un citu finanšu saistību, proporcionāli tās rokās nonākušajai teritorijai.

Kodolīgi par Versaļas līgumu ir izteicies angļu vēsturnieks A.~Dž.~P.~Teilors: ``Centrāleiropas politiķi iedomājās, ka ir paveikuši brīnumdarbu un izglābušies no Vācijas hegemonijas, nekļūstot atkarīgi no Krievijas. Šī brīnumdarba pamats bija alianse ar Franciju \citespace{} Valstis-mantinieces paļāvās uz Francijas stiprumu; savukārt Francija cerēja saņemt no tām spēku, kas pašai trūka.''

Vienlaikus ar Versaļas līgumu Polijai, tāpat kā citām jaunizveidotajām valstīm, atbilstoši Versaļas līguma 93.~pantam, bija jāparaksta īpašs līgums ar Sabiedrotajām valstīm par nacionālo, valodas un reliģiozo mazākumu aizsardzību, uzņemoties sevišķas saistības minoritāšu tiesību ievērošanas jomā.

Nacionālu valstu radīšanas principu pēckara Eiropā nevarēja konsekventi realizēt, tas tika realizēts tikai daļēji, jo noteikt nacionāli taisnīgas robežas praktiski nebija iespējams, radās nacionālo mazākumu, kuriem bija jākļūst par jauno valstu~--- Polijas, Čehoslovākijas, Ungārijas u.c. iedzīvotājiem, problēma. Pastāvēja divi tās atrisināšanas varianti.

Pirmais~--- iedzīvotāju pārvietošana. To jau 1916.~gadā bija piedāvājis franču antropologs un etnologs Ž.~Montandons. Pēc Pirmā pasaules kara tikai Balkānos tika realizēta nacionālo mazākumu iedzīvotāju pārvietošana. Tai par pamatu kļuva 1913.~gadā Turcijas un Bulgārijas noslēgtais līgums par iespējamo iedzīvotāju apmaiņu pierobežas rajonos un 1919.~gadā Neijī (\frtxti{Neuilly}, Francijā) parakstītā Grieķijas un Bulgārijas konvencija un līdzīga 1923.~gadā Lozannā parakstīta Grieķijas un Turcijas konvencija. To realizācija radīja lielus ekonomiskus un sociālus sarežģījumus, Balkānu pieredze tika vērtēta negatīvi un neviens no toreizējiem vadošajiem Eiropas politiķiem neizvirzīja un neatbalstīja projektus par nacionālo mazākumu pārvietošanu.

Parīzes konferencē priekšroka tika dota citai~--- nacionālo mazākumu tiesību aizsardzības idejai. Sākotnējo iniciatīvu šāda līguma radīšanai deva ebreju organizācijas ASV un Francijā, kuras bija norūpējušās par antisemītiskajiem noskaņojumiem un pat ebreju grautiņiem Austrumeiropā un jo sevišķi Polijā 1918.--1919.~gadā.

Par noskaņojumu pret ebrejiem poļu sīkpilsoņu aprindās ir liecinājusi arī kādu laiku Latvijā dzīvojušā ievērojamā ebreju vēsturnieka S.~Dubnova meita S.~Dubnova-Ērliha. 1919.~gadā viņa dzīvoja Varšavā un veda savus bērnus pastaigās uz parku pie Sakšu pils (\pltxti{Pałac Saski}). Ar bērniem viņa sarunājās krievu valodā. Tas nereti izsauca dažādas aizvainojošas replikas no parka apmeklētāju puses un padomus iemācīties poļu valodu. Kad S.~Dubnova-Ērliha mierīgi atbildēja, ka poļu valodā viņa runā ne sliktāk par nelūgtajiem padomdevējiem, tas izsauca jaunu uzbudinājuma vilni. Viņai tika ieteikts doties vai nu pie ``sava'' Trocka vai arī atbilstoši tradīcijai uz Palestīnu. Gandrīz visos poļu preses izdevumos, izņemot sociālistiskos, varēja sastapt antisemītiskus izlēcienus. Kaut virkne ievērojamu poļu rakstnieku, kā E.~Ožeškova, M.~Konopņicka u.c. vienmēr centās iedvest lasītājiem, ka ebreji ir nesaraujami saistīti pagātnes un tagadnes saitēm ar vietējiem iedzīvotājiem un ir šis zemes pilsoņi, ``tagad aicinājumi izskaust nacionālo naidu slīka šovinistisko noskaņojumu trakošanā, ko radīja jaunā valstiskuma rašanās atmosfēra''.

Polijai neizdevīgu noteikumu iekļaušanu nacionālo mazākumu aizsardzības līgumā objektīvi novērtēja laikraksts ``\pltxti{Gazeta Polska}'' (``Poļu Avīze'') 1934.~gadā: ``Pamatā pienākumu uzkraušanai Polijai pret tās nacionālajiem mazākumiem kalpoja zooloģiskā nacionālisma vilnis, kurš vēlās pār poļu zemēm pēc poļu valstiskuma atjaunošanas.''

Jaunveidojamajām valstīm bija pienākums garantēt nacionālajiem mazākumiem tiesības uz dzīvību un brīvību, tiesības ierīkot savas sociālās un reliģiskās iestādes, apgabalos, kur tām bija ievērojams īpatsvars, dzimtās valodas lietošanu skolās, arī privātu skolu uzturēšanu. Polija bija pirmā valsts, kura 1919.~gada 28.~jūnijā parakstīja šādu traktātu ar Tautu Savienību \pltxti{(Traktat Mniejszošciowy}). Viņai sekoja citas jaunās valstis. Ar to tika radīta nacionālo mazākumu tiesību aizsardzības sistēma, kura deva tiem daudz vairāk savu interešu aizstāvības iespēju, nekā tas bija iepriekšējā periodā. Taču sistēmai bija arī būtiskas nepilnības.

Tautu Savienības un Polijas līguma 12.~pants noteica, ka līguma nosacījumi nevar tikt mainīti bez Tautu Savienības Padomes vairākuma piekrišanas, katrs Padomes loceklis guva tiesības griezties pie tās, norādot uz līguma pārkāpšanu vai tās draudiem no Polijas puses. Savukārt Padome varēja lemt, kā rīkoties, nepieciešamības gadījumā pat iesniedzot jautājumu izlemšanai starptautiskai tiesai, kuru paredzēja Tautu Savienības statūtu 14.~punkts. Jebkura valsts~--- Tautu Savienības locekle varēja pievērst uzmanību faktam, ka Polijā tiek pārkāpti tās pienākumi pret nacionālajiem mazākumiem. Tādā gadījumā Tautu Savienības Padome ar balsu vairākumu varēja pieņemt lēmumu, kurš bija obligāts Polijai. (Turpmāk ļoti bieži šo iespēju izmantoja Vācija). Tiesa, līgumā nebija paredzēts nacionālo mazākumu stāvokļa kontroles mehānisms, katra puse (šai gadījumā Polija) varēja pati noteikt savu pienākumu apjomu, kas būtiski mazināja nacionālo mazākumu tiesību starptautisko garantiju lomu. Nepastāvēja līguma pārkāpēju ietekmes mehānisms, ja kāda valsts nepildīja uzņemtās saistības. Nebija arī runāts par nacionālo mazākumu kolektīvajām tiesībām, starptautiski bija paredzēts aizstāvēt tikai ar nacionālo izcelsmi vienotu indivīdu tiesības. Jaunizveidotās valstis, kuras atbilstoši Versaļas līgumam uzņēmās saistības nacionālo mazākumu aizsardzības jomā, arī norādīja, ka sistēma neattiecas vienādi uz visām Tautu Savienības loceklēm.

Arī Polijai 1919.~gadā Versaļā tika uzspiestas tādas saistības mazākumtautību tiesību jomā, kādas netika prasītas ne no Rietumu lielvalstīm, ne no Vācijas, kur bija liela poļu minoritāte. Tādēļ jau kopš 30.~gadu sākuma Polija Tautu Savienībā un starptautiskajās konferencēs arvien uzstājīgāk aicināja šādas saistības vienādā mērā attiecināt uz visām valstīm. (Polija valdība apgalvoja, ka līgums Polijai ticis uzspiests ar varu un nav taisnīgs, bet 1934.~gadā to uzteica.)

Pēc šī līguma parakstīšanas Polija paziņoja, ka tā ``automātiski un bez kādām formalitātēm'' atzīst par Polijas pilsoņiem tos bijušos Vācijas, Austrijas, Ungārijas un Krievijas pavalstniekus, kuri līguma stāšanās brīdī pastāvīgi dzīvo Polijai piešķirtajās teritorijās, un šīs personas, kuras ir vecākas par 18~gadiem, varēs veikt izvēli par labu jebkurai citai viņiem pieejamai pilsonībai. Tāpat atjaunotā valsts uzņēmās saistības, ka ``netiks izdoti nekādi ierobežojumi Polijas pilsoņiem brīvi lietot jebkuru valodu, vai tas būtu privātajā, tirdznieciskajā, reliģiskajā sfērā, vai presē, dažāda veida publikāciju lietās, vai publiskās sapulcēs''. Neraugoties uz Polijas valdības ieviesto poļu valodu valsts iestādes, citas valodas lietojošajiem bija jābūt atbilstošiem atvieglojumiem tiesās u.c. Piemēram, apgabalos ar ievērojamu nacionālo mazākumu pārstāvju skaitu to bērniem skolās bija jānodrošina tiesības mācīties dzimtajā valodā, ko poļu nacionāldemokrāti uzskatīja par neattaisnojamu ``privilēģiju''. Tai pat laikā mazākumtautību skolās poļu valodas apmācība bija jāievieš kā obligāta.

Jāuzsver, ka neraugoties uz teritoriālajiem ieguvumiem, Versaļas līgumu, kopā ar t.~s. ``mazo Versaļas līgumu'' par nacionālo mazākumu tiesībām arī poļu sabiedrībā vērtēja kā daļēju diplomātisku sakāvi un pāridarījumu Polijai. 1919.~gada 31.~jūlijā Seimā līgums tika ratificēts tikai ar 265 (citur literatūrā~--- 268) no 416~balsīm. (Pēc ļudoviešu kustības līdera V.~Vitosa liecības J.~Pilsudskis, izmantojot kā ieganstu to, ka I.~Paderevskis nespējis aizstāvēt Polijai vajadzīgās austrumu robežas, atlaidis viņu no premjerministra amata tikko bija izkaltusi tinte uz viņa parakstītā līguma. Iespējams, V.~Vitoss arī neprecīzi traktē notikumu gaitu. Literatūrā ir izteikts cits viedoklis~---  sakarā ar pieaugošajām J.~Pilsudska atbalstītāju un nacionāldemokrātu nesaskaņām I.~Paderevska valdība 1919.~gada novembrī bija spiesta atkāpties un decembrī par premjerministru kļuva L.~Skuļskis, kurš atbalstīja J.~Pilsudska politiku.)

Jau 1919.~gadā iznāca ievērojamā poļu zinātnieka (jurista un vēsturnieka) un politiķa S.~Kutšebas grāmata ``\pltxti{Kongres, traktat i Polska}'' (Kongress, līgums un Polija), kurā viņš vērtēja Versaļas līgumu kā Polijas suverenitāti pārkāpjošu un Poliju nostādošu atkarībā no lielvalstu politikas. Tiesa, autors aizvēra acis uz Polijas delegācijas pārmērīgajām prasībām robežu jautājumā, kuras atradās pretrunā ar Eiropas sabiedriskajā domā dominējošo pārliecību, ka robežu nospraušanas pamatā ir jāliek etniskais princips. No tā laika līdz pat mūsu dienām, izņemot vienīgi t.s. Polijas Tautas republikas eksistences posmu, tāda nostāja dominē Polijas nacionālajā historiogrāfijā.

\strong{Polijas cīņa}, arī bruņota, \strong{par izdevīgām robežām} konfliktu rajonos sākās vēl pirms Parīzes miera konferences, turpinājās tās laikā, nenorima pēc līgumu noslēgšanas.

J.~Pilsudska pretinieki viņam gan pārmeta poļu apdzīvoto zemju rietumos nenovērtēšanu, aizraušanos ar labvēlīgu austrumu robežu ieguvi. Viņš tik tiešām cīņai par izdevīgām rietumu robežām veltīja mazāk vērības kā par austrumu robežām. Taču jāpiekrīt poļu vēsturniekiem D. un T.~Nalenčiem, ka J.~Pilsudskis tā rīkojās nevis tāpēc, ka viņa sirdij bija tuvāks nevis Pjastu, bet Jagelloņu Polijas variants, bet gan tāpēc, ka viņš skaidri apzinājās, ka Polijas rietumu robežas, īpaši jau ar Vāciju, galu galā noteiks uzvarētājas karā~--- Antantes valstis un Polija šo procesu var maz ietekmēt. 1919.~gada 31.~maijā vēstulē I.~Paderevskim viņš izklāstīja savu uzskatu, ka rietumu robežu jautājumā Polija par 9/10 ir atkarīga no Antantes, tāpēc, kamēr šis jautājums nebūs atrisināts, jānovilcina visas citas lietas, kurās Polija varēja nonākt konfliktā ar Antanti. Tikai pēc tam Polija austrumos kļūs par spēku, ar kuru Antantei būs jārēķinās. ``Tad būs viegli, izmantojot ieganstus, kuri vienmēr atradīsies, risināt šīs lietas savā labā, izmantojot to situāciju, ka mēs būsim citiem vairāk vajadzīgi, nekā viņi mums.''

Cīņa par Vācijas sastāvā ietilpstošo poļu apdzīvoto apgabalu pievienošanu Polijas valstij norisa asās formās.

Līdz Pirmā pasaules kara beigām \strong{Pozenes province}, kurā pēc vācu datiem dzīvoja ap 60\% poļu un 38\% vāciešu, palika Vācijas rokās. Taču 1918.~gada 10.~novembrī arī tajā izplatījās revolūcija. Lielākajās novada pilsētās veidojās strādnieku un karavīru padomes, 20.~novembrī šo padomju kongresā bija pārstāvētas 54~padomes, taču to sastāvā bija maz strādnieku un tās ātri zaudēja ietekmi, ``izšķīda'' citās vietējās varas iestādēs, bet 1919.~gada 8.~janvārī tika atlaistas. 15.~novembrī dažu Prūsijas landtāga poļu deputātu vadībā Pozenē/Poznaņā tika radīta Augstākā tautas padome (\pltxti{Naczelna Rada Ludowa}), kura sastāvēja galvenokārt no nacionāldemokrātiem un sākotnēji pat neatzina ``sociālistisko'' J.~Pilsudska valdību. 27.~decembrī Polijas apdzīvotajos Prūsijas rajonos sākās t.s. Lielpolijas sacelšanās (\pltxti{Powstanie wielkopolskie}), kuras galvenie notikumi risinājās Poznaņā, Bidgoščā (poļu \pltxti{Bydgoszcz}, vācu \detxti{Bromberg}) un Toruņā (poļu \pltxti{Toruń}, vācu \detxti{Thorn}). Radās t.s. Lielpolijas fronte. Par poļu daļu pavēlnieku Lielpolijā J.~Pilsudskis nozīmēja J.~Dovbor-Musņicki. Norisa asiņainas sadursmes starp no frontes mājup pārnākušajiem vācu karavīriem, un poļu sacelšanās dalībniekiem, kuri veidoja poļu milicijas vienības. Janvārī un februāra sākumā sacēlušies atbrīvoja no vācu daļām lielu daļu Lielpolijas un Austrumpomorjes (poļu \pltxti{Pomorze Wschodnie}, vācu \detxti{Pommerellen}). Tikai tajos provinces apvidos, kur bija vācu iedzīvotāju vairākums, viņiem izdevās apturēt sacēlušos. Abas puses izmantoja sadursmes, lai aplaupītu un izvērstu teroru pret citas tautības vietējiem iedzīvotājiem. Šie notikumi vēl vairākus gadu desmitus uz priekšu apgrūtināja abu tautu attiecības. Sadursmes beidzās ar poļu uzvaru. Vācu ierēdņi vairs nevarēja turpināt pārvaldi. 1919.~gada 16.~februārī tika noslēgts pamiers uz \latxti{status quo} (uz to brīdi pastāvoša stāvokļa) pamata. Bija beigusies simtgadīgā vācu kundzība Pozenes/Poznaņas provincē. Lielpolijas sacelšanās veiksme nostiprināja J.~Pilsudskī pārliecību, ka savu politisko mērķu sasniegšana ir jābalstās uz ``notikušu faktu'' politiku, kaut tā arī nenovēršami bojāja attiecības ar tautām, pret kurām tika pielietota.

Kad Versaļas miera līgums 1920.~gada 10.~janvārī stājās spēkā, Polija saņēma no Vācijas 42~927 km$^{2}$ teritorijas (literatūrā norādītais iegūtās teritorijas lielums nedaudz atšķiras), apmēram pusi, ko tā bija sākotnēji pieprasījusi. Pēc 1910.~gada Prūsijas tautas skaitīšanas datiem šais apgabalos dzīvoja pavisam 2~962~000 cilvēku. No viņiem 1~714~000 runāja poliski, 1~080~000 vāciski, 105~000 kašubu un 9~000 mazūru valodā. No Vācijas tika nodalīta lielākā daļa no Pozenes provinces, kur, kā iepriekš minēts, vara jau bija nonākusi poļu rokās, un Rietumprūsijas; tāpat daļa no Lejassilēzijas un divi Austrumprūsijas apgabali. No šiem apgabaliem tika izveidots diezgan šaurs (vidējais tā platums nepārsniedza 200~km, bet šaurākā vietā tas sasniedza tikai 30~km) t.~s. Poļu koridors (vācu \detxti{Polnischer Korridor}; poļu \pltxti{Korytarz poliski}, arī \pltxti{Korytarz Gdański}; angļu \entxti{Polish Corridor gdański}), kurš veda caur plašu Vācijai piederošu zemju masīvu un nodrošināja Polijai pieeju jūrai, bet Vāciju ``sagrieza'' divās daļās. Polijas rokās nonāca 71~km Baltijas jūras piekrastes. J.~Pilsudskis šādu risinājumu uzņēma ļoti piesardzīgi: ``Tā ir liktenīga līnija, tas ir ļoti slikti''. Viņaprāt Polijai pietiktu ar Dancigas (Gdaņskas) iegūšanu. Turpretī ``poļu koridora'' izveide radīja nākamo vācu-poļu pretrunu avotu.

Arī daudzi ieinteresēti un neieinteresēti vērotāji brīdināja par sekām, kādas nākotnē var būt šādam teritoriālajam risinājumam. Franču ģenerālis M.~Veigans par koridoru teica, ka tas ``ir nevajadzīgs miera laikā un nav aizstāvams kara laikā''. Neilgi pirms nāves 1929.~gadā Francijas un Polijas maršals F.~Fošs, skatoties kartē, uzlika roku ``Poļu koridora'' teritorijai un pateica ``Tur atrodas nākamā kara saknes''. Angļu liberālis J.~Kenvorts savā grāmatā, veltītā ``poļu koridora'' problēmai, rakstīja: ``Vienu Eiropā var pateikt pilnīgi noteikti: tagadējās Polijas robežas [ilgi] nepastāvēs''.

Tomēr 1920.~gadā ģenerālis J.~Hallers vadīja Pomerānijas ieņemšanas operāciju un pie Puckas (poļu \pltxti{Puck}, vācu \detxti{Putzig}) pilsētiņas Polijas iekšlietu ministra (vēlākā Polijas prezidenta) S.~Vojcehovska, daudzu ģenerāļu un sabiedrisko darbinieku klātbūtnē noturēja ar zināmu misticismu caurstrāvotu simbolisku Polijas saderināšanās ar jūru ceremoniju. Pirms ceremonijas Gdaņskas poļu kopienas iedzīvotāji ģenerālim J.~Halleram uzdāvināja divus platīna gredzenus. Ģenerālis zirgā piejāja pie jūras un vienu no gredzeniem ar vārdiem ``Gaišības Polijas Žečpospolitas vārdā es, ģenerālis Hallers, pieņemu īpašumā šo mūsu seno slāvu Baltiju'' iemeta jūrā. Otru gredzenu viņš uzvilka sev pirkstā. Tam sekoja 21~artilērijas zalve. Pēc dievkalpojuma tika iesvētīts un pacelts Polijas jūras kara spēku karogs, kā zīme tam, ka tie uzņemas sardzi Baltijas jūrā. Vietējais kašubu kapteinis nodeva poļu jūrniekiem simbolisku sardzi. Tika parakstīts speciāls akts, kurš apliecināja ``saderināšanos'', kurš mūsdienās glabājas Nacionālajā muzejā Krakovā. ``Saderināšanās'' vietā tika uzstādīts stabs ar piemiņas uzrakstu.

Polija pretendēja arī uz Dancigu (Gdaņsku). Pilsēta bija piedzīvojusi dažādus laikus, vairākus gadsimtus piederējusi gan Vācu ordenim, gan Prūsijai, gan bijusi brīvpilsēta Polijas karalistē, apdzīvoja to pārsvarā vācieši. Taču pēc Pirmā pasaules kara pilsēta vēl nenonāca Polijas rokās. Danciga ar apgabalu bez iedzīvotāju nobalsošanas tika pasludināta par brīvpilsētu Tautu Savienības protektorātā, kurš gan vēlāk izrādījās visai mazefektīvs. Formāli Danciga bija suverēna valsts ar savu karogu un savu naudu (Dancigas guldeni). Tikai Dancigas ārpolitikai bija jāpaliek Polijas kontrolē. Pretrunas starp vācu pamatiedzīvotājiem un poļu nacionālistiem noveda pie tā, ka Polija nevarēja efektīvi izmantot Dancigas ostu. Tā padomju-poļu kara laikā Dancigas ostas strādnieku streika dēļ Polijai pārtrūka piegādes caur šo ostu. Britu pētnieks M.~Folliks 1929.~gadā par ``poļu koridora'' un Dancigas brīvpilsētas radīšanu rakstīja: ``Poļu koridora radīšana ir tūkstošreiz smagāks noziegums kā, ja Vācija savas uzvaras karā rezultātā radītu koridoru, pieņemsim, caur tagadējo Kaledonijas kanālu (\entxti{Caledonian Canal}~--- atrodas Skotijā, šķērso visu Lielbritāniju no austrumu līdz rietumu piekrastei) un nodotu šo 10~jūdžu plato joslu Holandei ar vienīgo mērķi novājināt Britāniju. Apmēram tā rīkojās Francija, piešķirot Polijai koridoru, kas šķērsoja vienu no auglīgākajiem Vācijas apgabaliem. Piekrītot šim noziedzīgajam aktam, Francijas sabiedrotie paveica vienu no pašiem smagākajiem zināmajiem \citespace{} noziegumiem pret civilizāciju \citespace{} Lai iedotu Polijai jūras ostu tika veikts cits noziegums pret Vāciju: tai atņēma Dancigu un pasludināja to par brīvpilsētu. Bet no visa vāciskā Vācijā Danciga ir pati vāciskākā.'' M.~Folliks uzskatīja, ka Polijai ir jāizdara vienkārša izvēle: vai nu jāatdod ``koridors'' un Danciga Vācijai~--- vai jāgatavojas karam ar savu kaimiņvalsti, tikko tā būs “piecēlusies kājās'' pēc sakāves karā. Šai gadījumā, pēc pētnieka viedokļa, Polijai bija jābūt gatavai ``pašam postošākajam karam ar Vāciju, anarhijai, un, iespējams, atkrišanai verdzībā, no kuras tā tikai nesen atbrīvojās''. Ne visam no citētā var piekrist, taču jāatzīst, ka visas valstis Parīzes konferencē pirmkārt rūpējās par savām interesēm, bet Francijas un Polijas vēlmju īstenošana, izcilā britu militārā vēsturnieka D.~Fullera vārdiem sakot, veda pie ``pūķa zobu iesēšanas, no kuriem bija jāizaug vēl postošākam konfliktam, nekā ar vardarbīgo [Versaļas] mieru apdzēstajam.''

Polijas pieprasītajā Austrumprūsijas dienvidu daļā un Augšsilēzijā sabiedroto uzraudzībā vajadzēja notikt tautas nobalsošanai par piederību Vācijai vai Polijai. Francijai nācās piekrist plebiscīta rīkošanai.

Nobalsošanā 1920.~gada 11.~jūlijā Austrumprūsijas dienvidu daļā [Varmijas (vācu \detxti{Ermland}, poļu \pltxti{Warmia}), Mazūru (vācu \detxti{Masuren}, poļu \pltxti{Mazurzy}) un Pievislas (poļu \pltxti{Powisle}) apgabalos] pāri par 90\% iedzīvotāju izteicās par labu piederībai Vācijai. Par pievienošanos Polijai nobalsoja nedaudz vairāk nekā 15~tūkstoši, pret to~--- 447~tūkstoši.

Jāatzīmē, ka šādi nobalsošanas rezultāti rādīja, ka sadzīvē lietotā valoda un nacionālās piederības apziņa daudziem bija atšķirīga, jo gandrīz visi mazūriski un ļoti daudzi poļu valodā runājošie nobalsoja par labu palikšanai Vācijas sastāvā. Iespējams, daudzi balsotāji ņēma vērā arī abu valstu ekonomiskās attīstības līmeni, to, ko tās varēja nodrošināt saviem iedzīvotājiem. Poļi zaudējumu gan mēģināja izskaidrot ar nelabvēlīgo starptautisko situāciju, kad šķita, ka Sarkanā armija var ieņemt Varšavu un Polija pārdzīvo savas pēdējās dienas, neatzīstot iedzīvotāju vairākuma sevis identifikāciju par vāciešiem.

Notikumu laikā nacionāldemokrāti, bet pēc Otrā pasaules kara komunisti izteica J.~Pilsudskim un viņa piekritējiem pārmetumu, ka viņi atteikušies no cīņas par Polijai piederīgajiem rietumu rajoniem. Mūsdienās daži poļu vēsturnieki un arī krievu vēsturnieks G.~Matvejevs izsaka viedokli, ka tas noticis tāpēc, ka Polijas valdība, aizņemta ar karu austrumos, nesniegusi nosauktajos apgabalos dzīvojušajiem poļiem nepieciešamo palīdzību. Poļu vēsturnieks J.~Holcers, atzīstot, ka lielāka uzmanība bija veltīta austrumu novadu iegūšanai, tomēr apgalvo, ka tieši J.~Pilsudska piekritēji devuši ierosmi poļu sacelšanās kustībai Augšsilēzijā, jo viņi mazāk kā nacionāldemokrāti turējušies pie diplomātiskiem risinājumiem, bet ticējuši notikušu faktu spēkam. Dotās grāmatas autoram visi šie apgalvojumi nešķiet pilnībā pārliecinoši. Diezin vai Polijai jebkuros apstākļos bija izdevīgi izraisīt tiešu diplomātisku konfliktu ar Antantes valstīm un reālu militāru sadursmi ar Vāciju. Tāpēc rietumos tā rīkojās netieši, ar savu ietekmes aģentu starpniecību.

Asiņaini notikumi risinājās Silēzijā. Šeit ļoti aktīvi kā poļu nacionālistu vadonis darbojās nacionāldemokrāts, Vācijas Reihstāga un Prūsijas landtāga deputāts V.~Korfantu. Poļi mēģināja pārņemt šos apgabalus savā kontrolē pirms vēl bija izskanējis kāds starptautisks lēmums par teritoriju piederību, vairākkārt mēģināja sacelties, savukārt vietējie vācu iedzīvotāji dibināja pašaizsardzības vienības, aicināja palīgā bijušos frontiniekus, kuri veidoja brīvprātīgo korpusus (\detxti{Freikorps)} un steidzās izpalīdzēt tautiešiem. Kopumā triju gadu laikā notika trīs poļu sacelšanās.

Faktiski ar Polijas valsts netiešu atbalstu Silēzijas poļu nacionālisti pārkāpa Versaļas miera līgumu jau tūlīt pēc tā noslēgšanas. Kad kļuva redzams, ka lielvalstis nevēlas bez ierunām atdot Polijai Augšsilēziju (otro pēc Rūras nozīmīgāko Vācijas rūpniecības rajonu), poļu puse 1919.~gada augustā uzsāka jau iepriekš gatavotu mēģinājumu bruņotā ceļā iegūt šo provinci. Taču spēku poļu nacionālistiem izrādījās pārāk maz un vācieši samērā viegli sacelšanos apspieda.

Kopš 1920.~gada janvāra kārtības uzturēšanai Augšsilēzijā atradās angļu, franču un itāļu karaspēka daļas. No februāra darbojās šo valstu pārstāvju veidota sabiedroto komisija. Taču vācu un poļu sadursmes turpinājās.

Gadījās arī dažādi kuriozi. Piemēram, Silēzijas rietumos esošais vācu Šventes (\detxti{Schwenten}) ciems, tuvojoties poļu daļām un baidoties no izpostīšanas, pasludināja sevi par ``brīvvalsti'' (\detxti{Freistaat}) un tik tiešām ar to izvairījās no ieņemšanas un postījumiem. Pēc pamiera nodibināšanas no garīdznieka, mežziņa un viesnīcnieka pastāvošā ``brīvvalsts'' valdība sarīkoja plebiscītu, kurā pēc iedzīvotāju gribas pasludināja ``pievienošanos'' Vācijai. (Var piebilst, ka arī pēc galīga starptautiska noregulējuma ``Šventes Republikas'' novads palika Vācijas sastāvā.) Neveiksmīga bija arī otrā sacelšanā 1920.~gada augustā, kas sakrita ar Sarkanās armijas pietuvošanos Varšavai.

Augšsilēzijas strīdīgajā daļā, kur dzīvoja divvalodīgi iedzīvotāji, 1921.~gada martā notika plebiscīts. Tajā ap 60\% Augšsilēzijas iedzīvotāju nobalsoja par palikšanu Vācijā, bet ap 40\%~--- par pievienošanos Polijai. (Vācijas sastāvā vēlējās palikt 707~605 cilvēki, par pievienošanos Polijai izteicās 479~359 plebiscīta dalībnieki, 3~882 biļeteni tika atzīti par nederīgiem). 597~ kopienās bija poļu, bet 664~--- vācu vairākums. (Pēc citiem datiem poļi uzvarēja 678~gminās, vācieši~--- 844). Lielbritānija un Itālija, bija par plebiscīta rezultātu atzīšanu un attiecīgu teritorijas sadali. Taču iebilda Francija, vadoties no vēlmes maksimāli novājināt Vāciju.

3.~maijā vietējie poļi sacēlās kārtējo (trešo) reizi. Sacēlušos pusē karoja arī Polijas regulārās armijas karavīri. Vācu vēsturnieki pārmet arī franču okupācijas karaspēkam, ka tas nepretojies sacelšanās dalībniekiem. Šoreiz poļi spēja ieņemt lielāko daļu strīdīgo teritoriju. Taču pēc mēnesi ilgušām kaujām starp karojošajām pusēm nostājās Rietumvalstu spēki, poļi nevarēja panākt visas izrakteņiem bagātās Augšsilēzijas atdalīšanu no Vācijas.

Nevēloties arī turpmāk Eiropas centrā saglabāt bruņota saspringuma centru, Tautu Savienības Padome 1921.~gada 12.~oktobrī pretrunā ar plebiscīta rezultātiem piedāvāja nodot Polijai lielāko apgabala daļu pa uguns pārtraukšanas līniju. Augšsilēzijas apgabals pēc ``Sabiedroto valdību plebiscīta komisijas'' (vācu \detxti{Interallierten Regierungs und Plebiscitkommission}) lēmuma tika sadalīts starp Poliju un Vāciju. 20.~oktobrī Antantes valstis apstiprināja šo priekšlikumu. Rezultātā ap 71\% Augšsilēzijas teritorijas saņēma Vācija, bet ap 29\% teritorijas ar 46\% iedzīvotāju (tas sastādīja 3~213 km$^{2}$ teritorijas ar gandrīz miljonu iedzīvotāju), toties gandrīz visu rūpnieciski nozīmīgo daļu ar Katovices (vācu \detxti{Kattowitz}, poļu \pltxti{Katowice}) un Hožovas (vācu \detxti{Königshütt}e, poļu \pltxti{Chorzów}) pilsētām, 97\% dzelzs rūdas ieguves vietām, 76\% ogļu šahtu, 82\% cinka un 71\% alvas raktuvēm, visiem cinka un alvas ražošanas uzņēmumiem, 50\% koksa ķīmiskajiem uzņēmumiem, 50\% domnas krāsnīm~--- Polija. Vācu vēsturnieks L.~Dralle norādījis, ka Polijai pienākošajās daļā atradās 75\% visu rūpniecisko iekārtu. Krievu vēsturnieki uzsver, ka Polijas sastāvā nonākušajos rajonos bija koncentrēti lielākie rūpniecības uzņēmumi, kuri palika vācu tautības īpašnieku rokās, bet Vācija paturēja to Augšsilēzijas daļu, kurā atradās izejvielu avoti. Pēc oficiālas Augšsilēzijas daļas nodošanas Polijai 1922.~gada jūnijā strauji aktivizējās iedzīvotāju migrācijas procesi. Tā, no poļu rokās nonākušās Augšsilēzijas daļas emigrēja simti tehniķu un inženieru. Tomēr abās pusēs robežai palika simti tūkstoši cilvēku, kuri piederēja nacionālajiem mazākumiem.

Faktiski Polijas puse šai gadījumā šeit rīkojās no spēka pozīcijām, pārkāpjot agrāk noslēgtās vienošanās, piespiežot no kara nogurušās lielvalstis sankcionēt soli, kurš rupji pārkāpa vietējo iedzīvotāju vairākuma intereses. Kā rakstīja D.~Loids Džordžs: ``Es biju pārliecināts, ka Augšsilēzijas nodošana Polijai ir liela netaisnība, ka tā apdraudēs mieru visā Eiropā, jo izveidos Centrāleiropā jaunu Elzasu-Lotringu.'' Taču ir tad vēl Polijā īpaši ``pjastu'' varianta aizstāvji~--- nacionāldemokrāti palika neapmierināti.

Protams, vācu autori lielākoties klaji negatīvi vērtēja jaunās robežas, kuras viņuprāt bija novilktas patvaļīgi, līnijā, kur norisinājās nacionālās cīņas. Tās pat tika sauktas par ``degošajām robežām'' (``\detxti{Brennende Grenzen}''), naida robežām, kuras ar ``poļu koridoru'' šķīra vācu tautas daļas. Ne velti toreizējais Vācijas bruņoto spēku (reihsvēra) pavēlnieks ģenerālis H.v.~Zēkts 1920.~gada 20.~februārī paziņoja: ``Nevienam vācietim nav jāpakustina ne pirksts lai glābtu no boļševisma Poliju, šo nāvīgo Vācijas ienaidnieku, Francijas radīto un sabiedroto, vācu kultūras grāvēju, un ja velns parautu Poliju, mums vajadzētu tam palīdzēt.''

Protams, arī vācu nacionālisti Vācijas valstij un vācu tautai nodarītās netaisnības bija tendēti pārspīlēt. Kā apgalvoja laikabiedrs, vācu nacionālists R.~Vingendorfs, vienīgais ar patiešām poļu iedzīvotājiem apdzīvotais apgabals bija bijusi ``kongresa Polija'' (Polijas karaliste), kura jaunajā valstī sastādīja nedaudz vairāk nekā 1/3 teritorijas, bet kurā dzīvoja 2/3 visu valstī dzīvojošo poļu. Izmantojot saviem ekspansionistiskajiem plāniem labvēlīgo konjunktūru, poļi tik tālu paplašinājuši savas valsts robežas, ka tās atgādinot savam īpašniekam daudz par lielu mēteli, kuru viņš nespēj ar savu ķermeni aizpildīt. Daudznacionālajā valstī (\detxti{Vielvölkerstaat}) poļi mēģinot realizēt savu kundzību, gluži tāpat kā savulaik šļahta jagelloņu valstī. Cīņa pret nacionālajiem mazākumiem poļiem aizvietojot ``pozitīvu valsts ideju''. Pār Polijas robežu uz Vāciju nākot naida plūsma, kura radusies no poļu vainas apziņas par nelikumīgi iegūto īpašumu. Vācijā jaunradītā Vācijas~--- Polijas robeža tikusi vienmēr uzlūkota kā mūžīgs Vācijas apdraudējums. Šis pats autors arī uzsvēris, ka cīņās par robežu ar Vāciju krasi atšķirīga bija J.~Pilsudska piekritēju, kuri galveno vērību veltīja Polijas austrumu robežām, un nacionāldemokrātu nostāja, kuri galveno uzmanību pievērsa valsts rietumu robežai. Šo pēdējo vērtējumu atzinuši arī daudzi mūsdienu vēsturnieki.

Var arī atzīmēt, ka pēc Otrā pasaules kara padomju vēsturnieki pārmeta jau bijušajiem Polijas valdošajiem slāņiem kopumā, ka tie 20.~gadu sākumā nepietiekoši atbalstījuši sacēlušos, kuru vidū strādnieki cīnījušies pret saviem sociālajiem apspiedējiem~--- vācu kapitālistiem. Padomju vēsturnieks A.~Manusevičs 1952.~gadā izdotā grāmatā, tendenciozi attaisnojot Otrā pasaules kara rezultātā PSRS panākto Polijas ``pārbīdīšanu'' no austrumiem uz rietumiem, kritizējot šai sakarā Rietumvalstis, īpaši uzsvēra, ka, ``nespējot pēc Lielās Oktobra sociālistiskās revolūcijas aizkavēt neatkarīgas poļu valsts rašanos, Antantes valstu vadītāji jau iepriekš atteicās iekļaut Polijas sastāvā sensenās poļu rietumu teritorijas, kuru iedzīvotāju vidū ilgstošas ģermanizācijas rezultātā poļi pārstāja sastādīt vairākumu. Ar to vēl karā gaitā ar Vāciju ASV un Anglijas imperiālisti apsargāja vācu imperiālistu intereses, paredzot tiem ievērojamu lomu cīņā pret Padomju Krieviju''. Tam var iebilst, ka ``ASV un Anglijas imperiālistiem'' šai laikā ``cīņa pret Padomju Krieviju'' bija tikai viens no risināmajiem uzdevumiem. Plānojot pēckara pasaules uzbūvi, viņiem prātā nevarēja ienākt, ka boļševiki ilgstoši noturēsies pie varas, nerunājot jau par to, ka tie varētu prognozēt tādu PSRS un Vācijas sadursmi, kāda notika pēc gandrīz gadsimta ceturkšņa Otrā pasaules kara laikā. Šajā gadījumā paradoksālā veidā faktiski sakrita poļu radikāļu nacionālistiskā un padomju vēsturnieka ``marksistiskā'' pozīcija.

Tāpat Polijai bija nesaskaņas ar Čehoslovākijas Republiku par Olsas upes apgabalu ar Cešinas pilsētu. Tur 1919.~gada janvāra beigās izvērsās īsts karš. Uzbrukumu 23.~janvārī sāka čehi, kuri saprata, ka, ja 26.~janvārī šeit notiks poļu gatavojamās Polijas Satversmes Seima vēlēšanas, tas viņiem nozīmēs teritorijas zaudējumu. Nedēļas laikā čehi spēja ieņemt lielāko daļu teritorijas, kura kopš 5.~novembra atradās poļu kontrolē. Galvenie poļu spēki šajā laikā veda cīņas ar ukraiņiem Rietumgalīcijā un nespēja iejaukties Cešinas notikumos. Poļu papildspēki gan apstādināja čehu uzbrukumu, bet lai panāktu ko vairāk rezervju nebija. Konfliktā iejaucās lielvalstis, to spiediena rezultātā kara darbība tika apstādināta. 3.~februārī R.~Dmovskis un Čehoslovākijas ārlietu ministrs E.~Benešs noslēdza vienošanos par pagaidu apgabala sadali pa Olsas upi. Apmēram 70~000 poļu palika dzīvot Čehoslovākijas pusē. Pretinieki vienojās strīdu atrisināt ar plebiscīta palīdzību, taču tas tā arī nenotika. 1920.~gada martā konflikts uzliesmoja ar jaunu sparu. Poļu masu uzstāšanās Cešinas apgabalā piespieda Čehoslovākijas valdību ievest šeit kara stāvokli. Čehoslovākijas prezidents T.~Masariks paziņoja, ja konflikts netiks nokārtots par labu viņa valstij, tā iesaistīsies poļu~--- padomju karā Padomju Krievijas pusē. Polija, baidoties no kara divās frontēs, bija spiesta piekāpties. 1920.~gada 28.~jūlijā, kad Sarkanā armija strauji virzījās uz Varšavu, Rietumu sabiedrotie konferencē Beļģijā galīgi sadalīja apgabalu, lielāko daļu no Cešinas Silēzijas, Spišas un Oravas apgabaliem, kur bija daudz izrakteņu, attīstītāka rūpniecība, piešķirot Čehoslovākijai (tā ieguva 1~200 km$^{2}$ no 2~222 km$^{2}$ strīdīgās teritorijas un 293~tūkstošiem iedzīvotāju no pavisam 435~tūkstošiem. Polija saņēma nabagāko austrumu Cešinas daļu. Abas puses palika neapmierinātas. Strīdi par robežu Augstajos Tatros (Poļu \pltxti{Tatry Wysokie}, slovāku \sktxti{Vysoké Tatry}, vācu \detxti{Hohe Tatra}) turpinājās līdz 1924.~gadam. Tikai 1925.~gadā vienošanās ceļā tika noregulēti visi jautājumi, saistīti ar Cešinas sadali.

Noslēdzot jautājumu par Polijas rietumu robežu noteikšanu, jāmin mūsdienu ukraiņu publicista S.~Lozuņko viedoklis. Viņš pamatoti kritizējis Parīzes miera konferences dalībniekus par žonglēšanu ar principiem, lai panāktu sev tīkama miera līguma noslēgšanu, taču reizē arī norādījis, ka ``Polija ar tās nesamērīgo apetīti un ambīcijām, bez šaubām, ir galvenā vaininiece par to, ka Vilsona principi, un pats galvenais no tiem~--- tautu tiesības uz pašnoteikšanos, neuzvarēja. Protams, tas notika ar sabiedroto piekāpību un pat tiešu palīdzību, taču, ja poļi ieņemtu atbildīgāku un saprātīgāku pozīciju, ja viņi paskatītos tālāk par savu šļahtisko degunu, arī miers varēja būt daudz taisnīgāks.'' Par secinājumu gan var strīdēties, taču tas satur lielu patiesības devu.

\strong{Polijas austrumu robežas} nefiksēja nedz Versaļas, nedz Sen-Žermenas līgumi, kurus sabiedrotie noslēdza ar Vāciju un Austriju. Vēl daudz asiņaināk kā rietumos robežu jautājuma risināšana sākās austrumos, jo Krievija un tās teritorijā radušies pagaidu valstiskie veidojumi nebija aicināti uz Parīzes miera konferenci un tās dalībnieki savukārt nevarēja būtiski ietekmēt jautājuma nokārtošanu.

Valstis~--- uzvarētājas kopumā atbalstīja 1917.~gada martā izskanējušo Krievijas Pagaidu valdības viedokli, kas pamatojās uz etniskā principa. Pie tam tās neatzina ukraiņu un baltkrievu tautas par pastāvīgiem etnosiem, bet uzskatīja par vienotas krievu nācijas sastāvdaļām. Tāpat kā XIX gadsimtā uz šādu uzskatu par vienotām valstīm izveidojās Vācija un Itālija, tā pēc Pirmā pasaules kara tika veidotas Čehoslovākija (čehu un slovāku: \sktxti{Československo}) un Serbu, horvatu un slovēņu karaliste (serbu: \srtxti{Краљевина Срба, Хрвата и Словенаца}). Tādai pat pieejai bija jābūt attiecībā pret Poliju.

1919.~gada 12.~februārī Parīzes miera konferencē tika izveidota speciāla komisija Ž.~Kambona vadībā poļu lietu izskatīšanai, kurai tika arī uzdota priekšlikumu sagatavošana par Polijas austrumu robežām. Komisija noraidīja Polijas notas šai jautājumā, kuras paredzēja iekļaut Polijas sastāvā ievērojamas teritorijas ar ukraiņu, baltkrievu un lietuviešu iedzīvotājiem. Tāpat tika noraidīts Polijas oficiālais priekšlikums radīt federāciju. Jau 22.~aprīlī komisija piedāvāja pagaidu Polijas robežu, noteikt pēc etniskā principa (apmēram pa Bugas upes līniju). Parīzē dominējošais viedoklis par Polijas atjaunošanu etniskajās robežās gandrīz pilnībā sakrita ar bijušas Polijas karalistes robežu austrumos. Dokuments noslēdzās ar vārdiem: ``Tiesības, ar kurām Polija varētu pretendēt uz teritorijām uz austrumiem no norādītās līnijas, ir stingri ierobežotas''. Komisijas viedokli speciālā deklarācijā apstiprināja un publicēja sabiedroto Augstākā padome 1919.~gada 8.~decembrī. Šo lēmumu J.~Pilsudskis vienkārši ignorēja, risinot teritoriālos jautājumus ar militāra spēka palīdzību.

Ja, nosakot ziemeļu, rietumu un dienvidu robežas, Polija bija spiesta kaut zināmā mērā rēķināties ar Versaļā pasludināto etnisko principu, tad austrumu robežu nospraušanu tā risināja militārā ceļā, neievērojot arī V.~Vilsona programmas 13.~punktu, kurš pasludināja, ka neatkarīgajā Polijai jāietilpst ``teritorijām ar neapšaubāmi poļu iedzīvotājiem''. Ja nacionālās sadursmes starp poļiem un vāciešiem Poznaņā un Silēzijā jau 1919.~gada sākumā pārliecināja J.~Pilsudski par to, ka viņam Polijas iekšienē nepieciešams panākt kompromisu ar nacionāldemokrātiem un R.~Dmovski, ar nacionālo pārstāvniecību Parīzē, tad kompromisi ar citām nācijām, ja neskaita neskaidrus aicinājumus uz sadarbību, nedz rietumos, nedz vēl jo vairāk austrumos, faktiski netika meklēti. Netika ievērots, ka 1919.~gada 28.~jūnijā Versaļā parakstītajā Sabiedroto valstu un Polijas līgumā par nacionālo, valodas un reliģiozo mazākumu aizsardzību bija noteikts, ka Polija īsteno ``suverenitāti pār bijušās Krievijas impērijas daļu, kuru vairākumā apdzīvo poļi''. Tādejādi bija jānosaka robeža, kur poļi veido vairākumu, un kur nē. Taču Polijas valdītāji vēlējās noteikt robežu ar spēku, nevis ievērojot nacionālo sastāvu. J.~Pilsudskis 1919.~gada 7.~februārī diplomātam V.~Baranovskim, kurš devās uz Parīzes konferenci, norādīja, ka rietumu robežu noteikšana ir pilnīgā Antantes varā, bet austrumos ``ir durvis, kuras atveras un aizveras un viss ir atkarīgs no tā, kurš un cik plaši tās atvērs ar spēku''. Šo domu viņš daudzkārt atkārtoja arī vēlāk. Ja nacionāldemokrāti savu mērķu sasniegšanā robežu jautājumā vairāk rēķinājās ar starptautiskajām tiesībām, uzvarētājvalstu atbalstu, tad J.~Pilsudska piekritēji vadījās no ``jau notikušu faktu principa''. Gala rezultātā valsts politikā uzvarēja spēka pielietošanas variants~--- inkorporēt visu, ko vien izdosies iegūt. Jāsaka, ka Rietumvalstis, cenšoties novājināt uzvarēto Vāciju un boļševistisko Krieviju, arī nerīkojās pietiekami piesardzīgi, nedeva pretsparu poļu vadītāju imperiālistiskajiem centieniem un galu galā atļāva tiem radīt tādu valsts konfigurāciju, kura jau no sākuma nesa sevī nākamo starpvalstu pretrunu iedīgli.

Jau 1918.~gada novembrī un decembrī Padomju Krievijas valdība trīsreiz piedāvāja Polijai nodibināt diplomātiskas attiecības, taču tā izvairījās no priekšlikuma pieņemšanas. 26.~novembrī, nepieminot Padomju valdības uzaicinājumu uzņemt Polijā kā Padomju Krievijas pārstāvi J.~Marhļevski, Polijas ārlietu ministrs L.~Vasiļevskis, pieprasīja atbrīvot Maskavā arestētos Polijas Reģentu padomes pārstāvjus. 28.~novembrī KSFPR ārlietu tautas komisārs G.~Čičerins nosūtīja L.~Vasiļevskim notu ar paskaidrojumu, ka Reģentu padomes pārstāvniecības locekļi nav arestēti, bet gan atrodas brīvībā, ka Padomju valdība neuzskatīja Reģentu padomes pārstāvniecība par suverēnas Polijas iestādi, jo tāda okupācijas laikā neeksistēja, bet likvidējusi to esot nevis Padomju valdība, bet [boļševiku organizētā] Poļu bēgļu komiteja. Reizē L.~Vasiļevskis tika aicināts atbildēt uz jau iepriekš izteikto Padomju valdības jautājumu, vai Polijas Ārlietu ministrija ir ar mieru uzņemt J.~Marhļevski. Uz to Polijas valdība atteicās ``atbildēt uz jautājumu par savām attiecībām ar Krievijas Padomju republiku'' pirms nebūs likvidēts incidents ar ``mūsu misiju Maskavā''.

Pēc autora domām, jaunās Polijas diplomātija pieļāva politisku kļūdu, atzīstot Reģentu padomes pārstāvniecības locekļus Maskavā par ``mūsu misiju''. Enerģiskā vēršanās pret varmācību, kas skāra Reģentu padomes pārstāvniecības darbiniekus, visnotaļ jāvērtē atzinīgi, taču to varēja darīt, aizstāvot viņus kā jaundibinātās valsts pilsoņus, nevis faktiski atzīstot pēctecību ar okupantu radītu iestādi. Var apstrīdēt arī izvēlētā cīņas līdzekļa~--- atteikšanās no diplomātisko attiecību nodibināšanas, pareizību. Acīmredzot to, kāpēc Polijas valdība nevēlējās nekavējoties nodibināt diplomātiskas attiecības ar Padomju Krieviju, rāda oficiālajā Polijas valdības avīzē ``\pltxti{Monitor Polski}'' (``Polijas Monitors'') publicētais paziņojums: ``\dots{}Padomju valdība ir griezusies pie poļu valdības ar priekšlikumu pieņemt doktoru Marhļevski kā oficiālu krievu diplomātisko pārstāvi Polijā. Jautājuma atrisināšana ir jāatliek sakarā ar nepieciešamību aizstāvēt poļu intereses Krievijā, kā arī sakarā ar izveidojušos starptautisko stāvokli.'' Jaunveidojamās Polijas valdība baidījās stāties sakaros ar starptautiski neatzīto boļševiku valdību. Tas diemžēl laupīja Krievijas teritorijā nonākušajiem simtiem tūkstošu poļu bēgļu cerības drīzumā atgriezties dzimtenē. Krievu publicistikā pastāv arī viedoklis, ka J.~Pilsudskis un viņu atbalstošie poļu politiķi tīšuprāt nevēlējušies poļu bēgļu atgriešanos dzimtenē, lai izmantotu viņu atrašanos Polijas kaimiņapgabalos kā pamatojumu to militārai sagrābšanai. Viedoklis gan nav pierādīts ar faktiem.

Polijas valdības nostāju kompromitēja fakts, ka tā publicēja tikai savas Padomju Krievijas valdībai sūtāmās notas, bet slēpa no sabiedrības saņemtās Krievijas atbildes, padomju priekšlikumus nodibināt diplomātiskas attiecības. Arī pēc L.~Vasiļevska nomaiņas ar I.~Paderevski ārlietu ministra amatā padomju notas netika publiskotas. Kad 1919.~gada 10.~februārī G.~Čičerins nu jau kārtējo reizi piedāvāja Polijai nodibināt ``normālas attiecības'' un viņa notu Varšavā publicēja komunistisks nedēļas izdevums, to konfiscēja un izdevniecību slēdza. Polijas Komunistiskā strādnieku partija šo faktu izmantoja savos uzsaukumos, rādot, ka valdība slēpj dokumentus, lai varētu ``kliegt'' par draudošo ``boļševistisko iebrukumu'' un tā ievilkt tautu karā pret revolucionāro Krieviju, Lietuvu un Baltkrieviju.

1918.~gada 18.~decembrī Varšavā ieradās Krievijas Sarkanā Krusta delegācija piecu cilvēku sastāvā, kuru vidū bija arī divas sievietes, ar vienu no bijušajiem Polijas un Lietuvas sociāldemokrātijas vadītājiem, Krievijas Centrālās Izpildu komitejas sekretāru B.~Veselovski priekšgalā. (Vēlāk Polijas valdība mēģināja apgalvot, ka delegācija Polijā ieradusies nelegāli, taču tiesā noskaidrojās, ka patiesībā jautājums bija saskaņots un tai arī izdots attiecīgs dokuments ar Polijas ārlietu ministra L.~Vasiļevska parakstu. Acīmredzot bija notikusi jauno poļu birokrātu kļūda.) Delegācijas oficiālais uzdevums bija sniegt palīdzību krievu karagūstekņiem, kuri no Vācijas caur Poliju atgriezās dzimtenē. Tā apmeklēja Polijas iestādes~--- Darba ministriju un Galveno aizbildniecības padomi, tomēr drīzumā delegācijas locekļus arestēja, tad pēc viņu pieprasījuma nosūtīja uz austrumu robežu, taču netālu no tās 1919.~gada 2.~janvārī nošāva. Četri delegācijas locekļi, to vidū arī B.~Veselovskis, gāja bojā, piektais~--- L.~Alters bija tikai ievainots, bet noturēts par mirušu un atstāts mežā. Viņš izglābās un visu izstāstīja. Uzreiz pēc notikuma Polijas Ārlietu ministrija publiski apgalvoja, ka delegācija tikusi aizvadīta līdz Grodņas guberņai, bet tās nogalināšana tikai liecinot, ka tā nelegāli atgriezusies Polijas teritorijā.

Tā nelaimīgi sākās padomju-poļu attiecības un vaina par to bija jāuzņemas poļu pusei. Taču tā nevēlējās to darīt. Poļu presē parādījās spriedumi, ka šī slepkavība bijusi ``tipisks kolektīvas paštiesāšanas akts'', ``sabiedrības pašaizsardzības parādība''. Toties protestēja Dombrovskas baseina strādnieku deputātu padome. Tās 1919.~gada 2.~janvārī rezolūcija skanēja: ``Baseina apvienotās Padomes par noziegumu uzskata atbildīgu ``tautas'' valdību, kura ar savu kūdīšanu un provokatīvo izturēšanos pret delegāciju izsauca izrēķināšanos pār krievu biedriem.''

Notikušais izraisīja enerģisku KSFPR pretdarbību. Jau 1918.~gada 31.~decembrī KSFPR ārlietu ministrs G.~Čičērins nosūtīja protesta notu par Krievijas Sarkanā Krusta delegācijas ieslodzīšanu. Pēc viņu noslepkavošanas tika arestēti Polijas Reģentu padomes misijas locekļi Krievijā, kuriem šais dienās bija jādodas mājup. Viņi tika paturēti kā ķīlnieki. Par to arī tika paziņots Polijas pusei. Izraisītais skandāls lika Polijas iestādēm enerģiski rīkoties, lai noskaidrotu atbildīgos par Sarkanā Krusta pārstāvju nogalināšanu. Tika veikta aptveroša izmeklēšana. 1919.~gada martā Polijas Ārlietu ministrijas pārstāvis paziņoja G.~Čičerinam, ka vainīgie neizbēgs sodam, lūdza izteikt krievu Sarkanā Krusta biedrībai līdzjūtību sakarā ar tā darbinieku noslepkavošanu, ``vēl jo vairāk tāpēc, ka šīs drausmīgais noziegums noticis ne bez mūsu žandarmērijas aģentu vainas''. Atbildot uz vainas atzīšanu, padomju puse atbrīvoja Polijas Reģentu padomes darbiniekus.

Tikai 1920.~gada martā vainīgie nonāca tiesas priekšā, kura tiem piesprieda no gada līdz diviem ieslodzījumā. Tiesa, neviens no viņiem cietumā tā arī nenokļuva. Acīmredzot tāpēc, ka tai laikā sākās jauns poļu-padomju kara uzliesmojums un Polijas valdītāju politiskie apsvērumi izrādījās spēcīgāki par tiesiskajiem principiem.

Par starptautisko tiesību ievērošanu nerūpējās arī padomju puse. KSFPR valdība arī vēlāk turpināja prettiesisko ķīlnieku ņemšanas praksi. Tā, 1919.~gada 6.~maijā tā paziņoja, ka sakarā ar poļu armijas nežēlīgo rīcību tās okupētajās lietuviešu un baltkrievu teritorijās Padomju Krievijā arestēti kā ķīlnieki virkne poļu buržuāzijas pārstāvju. Padomju valdība piedāvāja izdot ķīlniekus apmaiņai pret Polijā arestētajiem padomju darbiniekiem un poļu komunistiem. Arī turpmāk, minot daudzskaitlīgus faktus par poļu zvērībām pret padomju karagūstekņiem un civilpersonām, ebreju grautiņiem, padomju puse piedāvāja veikt ķīlnieku apmaiņu pret Polijas varas iestāžu arestētajiem. Uzdevumu vest sarunas par šo jautājumu saņēma J.~Marhļevskis. Jūnijā, ceļā no Vācijas uz Maskavu būdams Varšavā, viņš tikās ar atsevišķiem Polijas valdības locekļiem un piedāvāja savus pakalpojumus sarunu vešanā. Jūlijā Belovežas (\pltxti{Białowieža}) dzelzceļa stacijā sākās viņa sarunas ar Polijas pārstāvjiem. Rezultātā, lai arī neregulāri, ar dažādiem starpgadījumiem, tomēr notika apmaiņa ar arestētajiem un ķīlniekiem, tika ievadītas arī politiskas sarunas.

Tikmēr 1919.~gada sākumā gan Padomju Krievijas, gan Polijas valdība centās izplatīt savu varu pār maksimāli lielāku teritoriju. Jautājums bija~--- kurai izrādīsies vairāk spēka.

Ja Polijas puse vēlējās atjaunot Žečpospolitas 1772.~gada robežas, iegūt kontroli pār Ukrainu (ar Donbasu), Baltkrieviju, Lietuvu, kā arī iegūt ģeopolitisko dominanci Austrumeiropā, tad boļševiki kā minimālo stādīja uzdevumu iegūt kontroli pār agrākajām Krievijas impērijas teritorijām un sovjetizēt tās, bet kā maksimālo~--- paplašināt ietekmi uz Poliju un Vāciju, radīt labvēlīgus apstākļus pasaules revolūcijai. Pēckara Vācijā pastāvošā nestabilitāte veicināja šādu uzskatu nostiprināšanos (Līdz 1920.~gadam jaunā Vācijas Republika pārdzīvoja trīs puča mēģinājumus un piecus valdības vadītājus).

\strong{Baltkrievija un Lietuva} kļuva par asu politiskās un nacionālās cīņas lauku. Padomju Krievijas valdība uzskatīja par savu internacionālo pienākumu palīdzēt Krievijas nomaļu tautām izveidot un nostiprināt padomju varu. Krievijā veidojamā Sarkanā armija kopā ar bijušās Krievijas impērijas rietumu nomaļu tautu nacionālajām daļām virzījās Igaunijas, Latvijas, Lietuvas un Baltkrievijas teritorijās. Padomju Krievijas Rietumu armijas uzdevums bija nodibināt kontroli pār Baltkrieviju un Lietuvu. 1918.~gada 10.~decembrī tā, sekojot Vācijas karaspēkam, kurš atkāpās no okupētajiem apgabaliem, ieņēma Minsku, pēc tam gandrīz visu Baltkrieviju. Sarkanās armijas virzība uz rietumiem izsauca Polijas protestu un uztrauca Antantes valsts.

Polija pretendēja uz daļu no Lietuvas un Baltkrievijas. Ar šīm teritorijām tai bija senas vēsturiskas saites. Tās savulaik bija ietilpušas Žečpospolitā. Poļus un lietuviešus it kā vienoja katoļu ticība. Taču arī šo faktoru savā labā vēlējās izmantot stiprākie~--- poļi. Jau pirms Pirmā pasaules kara poļu garīdzniecībai un politiķiem izdevās gandrīz pilnībā pārliecināt Vatikānu, ka tikai Polijas valsts spēj izplatīt katolicismu uz austrumiem. Turpretī lietuviešu garīdzniecības pārstāvji praktiski tika izolēti. Arī Pasaules kara laikā, vācu okupācijas varai netraucējot, notika lietuviešu teritoriju polonizācija.

Lietuvas un Baltkrievijas poļu iedzīvotāji organizēja savas ``pašaizsardzības'' vienības, saformētas no poļu korpusu bijušajiem karavīriem un citiem brīvprātīgajiem, kuras pretojās Sarkanās armijas daļām, kad tās ieņēma vācu okupantu atstātās teritorijas, izveidoja Austrumu nomaļu aizsardzības komiteju (``\pltxti{Komitetu Obrony Kresów Wschodnich}''), kas griezās pēc palīdzības pie Polijas valdības. Šo darbību J.~Pilsudskis raksturoja kā preventīvu Polijas aizsardzību pret briesmām no ārienes un pašaizsardzības vienības 7.~decembrī pasludināja par Polijas armijas sastāvdaļu. Vēlāk viņš rakstīja, ka ``Polijai kara sākums ar Padomēm bija jau 1918.~gadā.'' Tālāk viņš uzsvēra, ka jau 1918.~gadā sev stādījis skaidru mērķi karā pret Padomēm, proti, izlēmis sasprindzināt visus spēkus, lai neļautu vēlreiz Polijai uzspiest svešu varu.

Jau 1918.~gada novembrī J.~Pilsudskis pavēlēja radīt t.s. lietuviešu-baltkrievu pulkus, bet uz to pamata poļu vadītu lietuviešu-baltkrievu divīziju. Viņu nebūt neatturēja tas, ka Viļņā jau eksistēja nacionāla Lietuvas valdība. Jau 12.~decembrī Polijas valdības sēdē tika spriests par karaspēka nosūtīšanu uz Viļņu. Valdības vadītājs E.~Moračevskis paziņoja savu viedokli, ka būtu vēlams Lietuvā ar bruņotu spēku radīt \latxti{fait accompli} (latīņu~--- noticis fakts) situāciju. (Var piebilst, ka tādu pat viedokli par Polijas valdības un armijas uzdevumiem robežu paplašināšanā pauda daudzi Polijas politiķi un militāristi).

Vietējo poļu novadnieciskās intereses tika pakļautas vispārnacionālajām interesēm. Īpaši poļu muižnieki, baidoties zaudēt savus īpašumus, cerēja uz Polijas valsts atbalstu, bija gatavi nest arī materiālus upurus. Piemēram, Baltkrievijas Moziras apriņķī poļu muižnieki jau 1918.~gada janvārī griezās pie Reģentu padomes ar lūgumu viņu zemes pievienot Polijas valstij. Pēc gada viņi bija gatavi bez atlīdzības nodot tai 1/3 savu zemju piešķiršanai poļu karavīriem, kuri aizstāvētu ``austrumu nomales pret krievu agresiju'', sniegt arī finansiālu palīdzību Polijas valdībai, ieķīlājot savas pārējās zemes.

Var konstatēt, ka poļu-padomju militārais konflikts norisa starptautiski nevienlīdzīgos apstākļos. Polija, kaut vēl ar neskaidrām robežām, bija starptautiski atzīta valsts ar ietekmīgiem sabiedrotajiem Eiropas lielvalstu vidū, bet KSFPR ne tikai nebija starptautiski atzīta, bet bija spiesta izcīnīt Pilsoņu karu, kurš nevarētu tik plašā mērā izvērsties bez Rietumvalstu palīdzības ``baltajiem''. Šāda situācija nāca par labu poļu elites ekspansijai pret kaimiņtautām, īpaši dzīvojošām austrumos no Polijas.

Krievu vēsturnieks K.~Novikovs ir atzīmējis, ka padomju-poļu attiecības pilnā mērā atspoguļoja visu bijušās impērijas zemju teritoriālās norobežošanās problēmas smagumu. Poļu inteliģences vidū bija visai populāra ap simt gadu garumā propagandētā ideja par Krievijas vainu Polijas priekšā. Tā aktīvi tika izplatīta arī plašās poļu masās. J. Pilsudskim un viņa sekotājiem šī ideja kalpoja poļu nacionālisma un pretenziju uz savu austrumu kaimiņu~--- lietuviešu, ukraiņu un baltkrievu zemēm~--- attaisnošanai.

Revolūcija un Pilsoņu kara notikumi Krievijā noveda pie tā, ka ar nacionālajām pretrunām cieši savijās arī sociāli ekonomiskās pretrunas. Polijas vadītāji varēja sevi tēlot Eiropas un pasaules acīs ne tikai par ``civilizēto Rietumu'' aizstāvjiem pret ``barbariskajiem'' Austrumiem, bet arī kā šķērsli komunisma izplatīšanai visā Eiropā.

Teorētiski Polija varēja rīkoties dažādi. Tā varēja kopīgi ar Rietumu sabiedrotajiem palīdzēt ``baltajiem'' Krievijā un censties par šo palīdzību iegūt sev maksimāli labvēlīgākus pastāvēšanas nosacījumus. Taču pret šādu ceļu runāja visa iepriekšējā poļu-krievu attiecību vēsture. ``Baltie'' poļu valstij labākajā gadījumā piekristu piešķirt etnogrāfisku robežu (ko līdzīgu vēlākajai Kērzona līnijai, par kuru tiks runāts tālākajā tekstā). Otrs variants bija: palīdzot Lietuvā, Ukrainā, Baltkrievijā un arī Latvijā esošajiem pretboļševistiskajiem spēkiem, vienlaikus ar to pacensties nodrošināt sev vēlamos teritoriālos ieguvumus. Taču, ja Polija tiktu izveidota atbilstoši nacionālās pašnoteikšanās principam, tās teritorija aptvertu ap 150~tūkstošu km$^{2}$ teritorijas ar līdz 20~miljoniem iedzīvotāju. Kā uzskatīja poļu politiķi un militāristi, tas neļautu Polijai būt drošībā par savu neatkarību, blakus tai esot tādiem spēcīgiem kaimiņiem kā Vācija un Krievija. Protams, šādi apsvērumi neņēma vērā iespēju, ka Austrumeiropā blakus Polijai varētu izveidoties (un jau veidojās) virkne citu neatkarīgu nacionālu valstu, sadarbība ar kurām varētu sniegt drošību gan Polijai, gan tām. Tiesa, tad būtu jāatsakās no vienpusīgā poļu nacionālisma, jāatzīst agrāk poļiem pakļauto tautu vienlīdzība, kas bija pretrunā ar agrākajām šļahtas tradīcijām. Polijas vadība uzskatīja par labāku rīkoties pēc trešā varianta: ar spēku sagrābt tik daudz teritoriju, cik vien būs iespējams.

1918.~gada 22.~decembrī Padomju Krievijas Tautas komisāru Padome deva direktīvu Sarkanās armijas virspavēlniekam J.~Vācietim apturēt Sarkanās armijas virzību, sasniedzot agrākās Polijas karalistes robežu. Tai pat dienā KSFPR ārlietu tautas komisārs G.~Čičerins paziņoja Polijas valdībai, ka no Krievijas puses tai briesmas nedraud un Sarkanā armija ne tikai atrodas tālu no Polijas, bet ir nošķirta no tās ar Lietuvas, Baltkrievijas un Ukrainas teritoriju. Jāsaka, ka Padomju Krievijas miera piedāvājumi Polijai ir dokumentāli pierādāmi un nebija tikai vienkāršs propagandas paņēmiens lai parādītu savu internacionālismu. Tie atbilda gan pasludinātajai padomju ārpolitikas koncepcijai, gan arī Padomju valsts kritiskajam militārajam un ekonomiskajam stāvoklim. Taču Polijas puse gan 1918.~gadā, gan vēlāk visus izskanējušos Padomju Krievijas miera nodibināšanas aicinājumus acīmredzot vērtēja kā vājuma pazīmi, ar jebkuru ieganstu tos noraidīja, par tiem neziņoja presē, apzināti slēpjot no poļu un starptautiskās sabiedrības.

30.~decembrī Polijas valdība nosūtīja KSFPR Ārlietu tautas komisariātam protesta notu sakarā ar ``Padomju valdības, kuras karaspēks pārpludina Lietuvu un Baltkrieviju, imperiālistisko politiku'', apgalvojot, ka daļa šo teritoriju ir ``neapstrīdami piederošas poļiem''. Pēc nedēļas KSFPR atbildēja ar notu, norādot, ka Sarkanā armija nekur nav iegājusi teritorijā, ``kura varētu tikt uzskatīta par piederošu Polijas Republikai.'' Faktiski abas puses vēlējās ieņemt pēc iespējas lielākas pārsvarā lietuviešu, baltkrievu un ukraiņu apdzīvotās strīdīgās teritorijas, lai otru pusi nostādītu fakta priekšā. Notika virkne dažādu vienību sadursmju par apdzīvotām vietām, kur šo vienību piederību Sarkanajai armijai vai Polijas bruņotajiem spēkiem bija grūti klasificēt. Sarkanās armijas virspavēlnieks J.~Vācietis ziņoja V.~Ļeņinam, ka Padomju Krievijas tieksme sasniegt bijušās impērijas robežas var radīt ``\latxti{casus belli}'' (ieganstu karam) ar Poliju, jo dažviet Sarkanā armija jau atradās bijušās Polijas karalistes teritorijā. 1919.~janvāra vidū uz rietumiem no Lidas (\betxti{Лида}~--- pilsēta Baltkrievijas rietumdaļā) padomju un poļu karaspēka daļas sastapās. (Pēc dažiem datiem, pirmā kauja starp Padomju Krievijas un Polijas bruņotajiem spēkiem notika 1919.~gada 28.~janvārī pie Volkoviskas (baltkrievu \betxti{Ваўкавыск}, poļu \pltxti{Wołkowysk}, pilsēta Baltkrievijā), ko februārī ieņēma poļi.) 8.~februārī J.~Pilsudskis ar Francijas palīdzību noslēdza vienošanos ar Vācijas spēku pavēlniecību Baltkrievijā par poļu daļu izlaišanu caur vācu rindām, lai tās varētu nonākt saskarsmē ar Sarkanās armijas daļām.

Uzrunājot Satversmes Seima deputātus tā atklāšanas sēdē 1919.~gada 10.~februārī, valsts priekšnieks J.~Pilsudskis paziņoja: ``Neatdosim nevienu pēdu Polijas zemes un neļausim, lai mūsu robežas, uz kurām mums ir tiesības, tiktu samazinātas. Mūsu kaimiņu tieksmju rezultātā mums ar visiem viņiem pašlaik ir atklāts kara stāvoklis vai vismaz ļoti saspringtas attiecības''. Protams, tikpat pamatoti var teikt, ka tieši Polijas valstsvīru ``tieksmju'' rezultātā Polijai ar kaimiņiem izveidojās ``atklāts kara stāvoklis vai vismaz ļoti saspringtas attiecības.''

12.~februārī padomju pusē tika radīta vienota Rietumu fronte. 13.~februārī Baranoviču apkārtnē notika kārtējā poļu un padomju daļu sadursme, kurā poļu gūstā nokļuva pirmie 60 sarkanarmieši. Var uzskatīt, ka šai dienā sākās nepasludinātais poļu~--- padomju karš, kurš ilga 20~mēnešu un norisa galvenokārt zemēs, kuras nevarēja uzskatīt ne poļiem, ne krieviem piederošām, bet kuras abas puses uzskatīja par savām.

Februāra vidū Sarkanās armijas virzība apstājās apmēram 100~km uz austrumiem no tagadējās (XXI gadsimta sākuma) Polijas austrumu robežas, bet sadursmes ar poļu vienībām turpinājās. Tomēr Sarkanās armijas pavēlniecības direktīvā 22.~februārī kā uzbrukuma mērķis nebija norādīts neviens apdzīvota punkts bijušās Polijas karalistes teritorijā. Ar to padomju puse vēlreiz demonstrēja savu piekrišanu etnisku robežu nospraušanai.

Īsi jāpasaka par situāciju Polijas kaimiņapgabalos.

\strong{Lietuvas valsts} patstāvības atjaunošanu Lietuvas Padome (\lttxti{Lietuvos Taryba}) ar Vācijas valdības piekrišanu bija proklamējusi jau 1918.~gada 16.~februārī. Taču līdz 1918.~gada beigām Lietuva atradās vācu karaspēka kontrolē. 1918.~gada beigās, atkāpjoties Vācijas armijai, izveidojās vairāki valstiskuma centri: Lietuvas Padomes, kas vēlāk pārdēvējās par Lietuvas Valsts padomi (\lttxti{Lietuvos Valstybės Taryba)}, izveidotā valdība, ko finansēja Vācija; lietuviešu komunistu grupējumi un poļu apvienības. Visi šie spēki cīnījās par varu.

Ar Vācijas materiālu palīdzību Lietuvas nacionālā valdība ar A.~Voldemarasu priekšgalā sāka formēt savus bruņotos spēkus. Mūsdienu lietuviešu vēsturnieks Č.~Laurinavičus apraksta Lietuvas nostāju: ``Lietuvieši par katru cenu alka pēc Viļņas, lai iegūtu kopēju robežu ar Krieviju un tādā veidā izskaustu poliskuma saknes Lietuvā''.

Lietuviešu komunisti Maskavā 1918.~gada 8.~decembrī radīja Lietuvas strādnieku un zemnieku pagaidu valdību (\lttxti{Lietuvos laikinoji revoliucinė darbininkų ir valstiečių vyriausybė}) ar V.~Mickēviču--Kapsuku priekšgalā. 22.~decembrī KSFPR Tautas komisāru padome atzina Padomju Lietuvas neatkarību. 1918.~un 1919.~gada mijā Padomju Lietuvas valdība atradās Sarkanās armijas ieņemtajā Daugavpilī. Viļņā faktiski bija četras konkurējošas varas~--- vācu militāristu, lietuviešu nacionālistu, lietuviešu boļševiku un poļu nacionālistu.

Polija piedāvāja Lietuvas nacionālajai valdībai ievest Viļņā savu karaspēku. A.~Voldemaras piedāvājumu uztvēra labvēlīgi, par ko arī 19.~decembrī paziņoja Polijas valdībai. Arī Lietuvas valdība kopumā bija gatava atbildēt ar savu piekrišanu, tikai kā tās nosacījumu izvirzīja prasību, lai Polija ``atzītu patstāvīgu Lietuvas valdību ar galvaspilsētu Viļņā''. Tāda netika saņemta, jo Polija pati pretendēja uz Viļņu.

Senajā Lietuvas galvaspilsētā gadsimtu gaitā bija izveidojies poļu iedzīvotāju vairākums, ko Polijas vadība varēja izmantot gan savu pretenziju pamatošanai, gan savu militāro spēku atbalstam. (Tā, poļu vēsturnieks P.~Madajčiks norāda, ka Viļņas pilsētā 1916.~gadā lietuviešu bija mazāk par 3\%, vēlāk to īpatsvars kritās zem 1\%, taču nesaka, ka šāds stāvoklis bija izveidojies piespiedu polonizācijas rezultātā, kura īpaši saasinājās pēc Viļņas iekļaušanas Žečpospolitas sastāvā.) Lietuvas iestādes gan apgalvoja, ka vairākums šeit dzīvojošo poļu ir polonizēti lietuvieši un Lietuvai ir tādas pat tiesības radīt savu daudznacionālu valsti kā Polijai.

Polijas valdība A.~Voldemarasa 19.~decembra vēstuli iztulkoja kā piekrišanu poļiem aizstāvēt Viļņu. Polijas armijas pavēlniecība 29.~decembrī organizēja Lietuvas un Baltkrievijas kara apgabalu. Vietējās militarizētās poļu vienības tika pakļautas Polijas armijas pavēlniecībai.

Kad 1918./1919.~gadu mijā (datumi dažādās valstīs izdotā literatūrā ir atšķirīgi) vācu karaspēks kopā ar Lietuvas valdību evakuējās uz Kauņu, Viļņas pilsētas Strādnieku deputātu padome paziņoja par varas pārņemšanu Viļņā savās rokās, taču vietējās poļu militārās vienības gāza padomju varu un ieņēma galvenos objektus. Kāda poļu ielenkta vietējo komunistu grupa tika iznīcināta, pārējie bēga no pilsētas vai paslēpās. Tika radīta poļu administrācija, taču pastāvēja tā tikai dažas dienas, jo 1919.~gada 5.~janvārī Viļņu ieņēma Sarkanās armijas daļas un turp pārcēlās V.~Mickēviča--Kapsuka vadītā valdība. 1919.~gada sākumā jau ``sarkano'' daļas ieņēma lielu daļu Lietuvas. Polija šajā brīdī nespēja sniegt būtisku palīdzību ``Austrumu nomaļu aizsardzības komitejai'', jo daļa veidojamo poļu armijas vienību bija iesaistītas cīņās Galīcijā, robežu konfliktā ar Čehoslovākiju un gatavojās iespējamajam konfliktam ar Vāciju par Silēziju. Pie tam Polijas rietumu novados vēl atradās vācu armijas daļas. Tikai pēc Antantes iejaukšanās 5.~februārī tika parakstīts līgums, ka vācieši izlaidīs poļus uz austrumiem. Formāli tas nozīmēja, ka Polija reģionā uzņemas boļševisma apturēšanas funkcijas, bet faktiski tā ieguva iespēju realizēt šeit ekspansionistisku politiku.

\strong{Baltkrievijā} pēc Vācijas impērijas sabrukuma gan poļi, gan baltkrievi izvirzīja pretenzijas uz Minsku, bet līdz lielām sadursmēm lieta nenonāca. Baltkrievu nacionālisti veda ar poļiem sarunas par federācijas izveidi, kur Baltkrievija būtu tās patstāvīgs loceklis, taču poļi drīz atbruņoja dažas vācu radītās baltkrievu daļas, sākās baltkrievu aktīvistu aresti. Baltkrievijas patstāvību poļu nacionālisti negribēja atzīt. Baltkrievijā sāka augt atbalsts boļševikiem.

1919.~gada 1.~janvārī boļševiki proklamēja Baltkrievijas Sociālistisko padomju republiku (\betxti{Беларуская Сацыялістычная Савецкая Рэспубліка) kā KSFPR sastāvdaļu. 7.~janvārī tās valdība ieradās Minskā. 1919.~gada 31.~janvārī formāli Baltkrievijas SPR izstājās no KSFPR sastāva, 27.~februārī apvienojās ar Lietuvas Padomju republiku un tās abas ieguva vienu: Lietuvas-Baltkrievijas Padomju Sociālistiskās republikas }(lietuviešu \lttxti{Lietuvos~--- Baltarusijos Tarybų Socialistinė Respublika (Litbelas}); baltkrievu \betxti{Літоўска-Беларуская Савецкая Сацыялістычная Рэспубліка (Літбел)}, arī \betxti{Савецкая Сацыялістычная Рэспубліка Літвы і Беларусі)} nosaukumu. Republikas galvaspilsētai bija jābūt Viļņai. Poļu komunisti aktīvi piedalījās šīs republikas izveidē. K.~Cihovskis kļuva par tās Centrālās Izpildu komitejas priekšsēdētāju, J.~Unšlihts~--- par kara lietu, J.~Leščinskis-Ļenskis~--- izglītības, M.~Kozlovskis~--- justīcijas, S.~Bersons~--- valsts kontroles komisariātu vadītājiem. K.~Cihovskis un J.~Unšlihts bija republikas Aizsardzības padomes locekļi. Arī vairāki citi poļu komunisti ieņēma jaunizveidotajā republikā svarīgus amatus.

Polija neatzina ne Lietuvas Padomju valdību, ne Baltkrievijas Sociālistiskās padomju republikas, bet vēlāk Lietuvas un Baltkrievijas Padomju sociālistiskās republikas nodibināšanu. Kad 1919.~gada 16.~februārī Lietuvas un Baltkrievijas Padomju sociālistiskās republikas nosūtīja Polijai priekšlikumu nekavējoties izveidot komisiju valstu robežu noteikšanai, atbilde netika saņemta. Tautai tas tika noklusēts. Kad Polijas Komunistiskā partija savā avīzē pieprasīja sniegt paskaidrojumus, avīze tika nekavējoties slēgta. Tomēr komunistiem izdevās, kaut mazā tirāža, padomju priekšlikumus publicēt. Paši poļu komunisti šai laikā propagandēja miera un labu kaimiņattiecību nodibināšanu. Rezolūcijas pret karu ar padomju republikām pieņēma vēl pastāvošās strādnieku padomes. Pretkara rezolūcijas pieņēma arī PPS organizācijas, taču vienlaikus tās iestājās par Polijai labvēlīgām robežām.

1919.~gada februārī Polija, izmantojot par ieganstu cīņu pret boļševikiem, uzsāka kara darbību pret lietuviešu un baltkrievu padomju valstiskajiem veidojumiem, forsējot Nemunas upi. Saskaņā ar Vācijas armijas pavēlniecību, poļu daļām bija nodrošināta iziešana caur vācu ieņemtajām teritorijām, Vācu armija nodeva poļu daļām virkni stratēģisku punktu. Poļu karaspēka uzbrukums bija veiksmīgs arī tāpēc, ka padomju aizmugurē notika nemitīgas sacelšanās pret boļševiku varu. 9.~februārī poļi ieņēma Brestu. Viņiem pretī virzījās Sarkanās armijas daļas.

Februāra vidū virknē Baltkrievijas rajonu (kā arī Ukrainā~--- Volīnijā) notika pirmās poļu un padomju armijas daļu sadursmes. Ir ziņas, ka pirmā sadursme Baltkrievijā notika 14.~februārī pulksten sešos no rīta, kad poļu kavalērijas patruļa pārsteidza boļševiku nometni brokastojam. Kā raksta angļu vēsturnieks N.~Deiviss, Polijas un Padomju Krievijas spēku regulāras sadursmes ``gandrīz vai nejauši'' aizsākās 1919.~gada februārī un turpinājās 20~mēnešus.

Kaut 1919.~gada februārī Polijas bruņotajos spēkos nomināli jau skaitījās ap 150~000 vīru, to spēki Baltkrievijā bija visai nelieli~--- ap 8~tūkstoši cilvēku. Padomju Rietumu armijā skaitījās ap 45~000 cilvēku, taču šo daļu kaujas spējas bija ļoti zemas. (Ziņas par abu pušu bruņotajiem spēkiem visa 1919.~gada garumā, to kaujas spējām literatūrā ir ļoti pretrunīgas, balstītas uz atšķirīgiem avotiem.)

16.~aprīlī poļu kavalērijas spēki pārrāva fronti un uzsāka virzību uz Viļņu. Poļus atbalstīja liela daļa Viļņas apgabala iedzīvotāju, īpaši Poļu kara organizācijas (\pltxti{Polska Organizacja Wojskowa}) biedri. 21.~aprīlī pēc smagām trīs dienas ilgušām kaujām Polijas armija ieņēma Viļņu. (Boļševiku varas krišanu Viļņas apgabalā lielā mērā nosacīja viņu 1919.~gada sākumā realizētā neveiksmīgā agrārā un pārtikas politika. Vietējie iedzīvotāji to uzskatīja par leģitimētu laupīšanu un tāpēc neatbalstīja padomju varu.) Poļu gūstā nonāca un bez tiesas un izmeklēšanas tika nošauts arī Lietuvas-Baltkrievijas Padomju Sociālistiskās republikas valsts kontroles komisariāta vadītājs polis S.~Bersons. KSFPR vadība protestēja pret Viļņas ieņemšanu, tur notikušajām zvērībām pret ebreju (dažkārt arī kristiešu) iedzīvotājiem, taču, ievērojot Krievijā notiekošo Pilsoņu karu, turpināja neatlaidīgus mēģinājumus atrisināt jautājumus par poļu bēgļu repatriāciju no Krievijas uz Poliju, sadursmju novēršanu, normālu diplomātisku attiecību nodibināšanu.

Redzot, ka poļi negrasās atzīt Lietuvas patstāvību, arī Lietuvas nacionālā valdība ar A.~Voldemarasu priekšgalā pasludināja, ka poļu ieiešana Viļņā ir tās okupācija. Viļņas ieņemšana ar spēku slēdza Polijas valdībai iespēju turpmāk sadarboties ar lietuviešu un baltkrievu nacionālajām politiskajām elitēm, vairāk nekā problemātiska kļuva Polijas-Lietuvas-Baltkrievijas federatīvas valsts izveide. Polijas rīcība virzīja lietuviešu politiķus uz sadarbību ar Padomju Krieviju pret Poliju.

Tā Lietuvas un Baltkrievijas teritorijā \strong{izveidojās poļu~--- padomju fronte}. Tiesa, sākotnēji kara darbība norisa ne visai aktīvi, taču drīzumā, izplatoties arī uz Ukrainas teritorijām, Polijai tā kļuva par galveno. Faktiski ar to bija sācies poļu-padomju karš, kurš ilga divus gadus, un bija 1918.--1920.~gada lielākā militārā sadursme pēc tajā dažādās pusēs iesaistītā karaspēka skaita.

Šis bruņotais konflikts historiogrāfijā pazīstams ar vairākiem nosaukumiem. Starpkaru periodā poļu militāristi to sauca vienkārši par poļu-krievu karu, kas, pēc krievu vēsturnieka G.~Matvejeva domām, visprecīzāk atspoguļo konflikta saturu, jo karš starp Krieviju un Poliju norisa par teritorijām, bet nevis kādu ideoloģisku apsvērumu dēļ. (Kā vēlāk redzēsim, ``ideoloģiskajiem apsvērumiem'' bija svarīga loma karā, taču tā pamatā tik tiešām bija cīņa par teritorijām.) Tā konflikta būtību vērtēja arī J.~Pilsudskis. Pēc M.~Kosakovska atmiņām valsts priekšnieks 1919.~gada vasarā teica, ka tagad Polijai ir unikāla iespēja panākt austrumos lielu sasniegumu~--- ``ieņemt Krievijas vietu'', kļūt par hegemonu reģionā un beidzot nodalīt Poliju no Krievijas.

Padomju historiogrāfijā nostiprinājās cits nosaukums: padomju-poļu karš. Pēc Otrā pasaules kara tā to sāka saukt arī Polijas Tautas republikas vēsturnieki. Padomju Savienībā oficiāli šis karš tika skatīts ``šķiriski''~--- kā ārvalstu intervences jeb tā saucamā ``trešā Antantes kara gājiena'' pret Padomju Krieviju sastāvdaļa Krievijas Pilsoņu kara laikā, kur karoja buržuāziskie ``baltie'' pret boļševistiskajiem ``sarkanajiem''. Polijas mēģinājumi izcīnīt sev labvēlīgas robežas tika tēloti kā atbalsts ``baltajiem'' pret ``sarkanajiem''. Mūsdienu krievu historiogrāfijā parasti tiek lietots nosaukums ``Padomju--Polijas 1919.--1921.~gada karš'' (\rutxti{Советско-польская война 1919--1921 годов}).

Mūsdienu poļu vēsturnieku darbos lielākoties to sauc par ``poļu-boļševiku karu'' (\pltxti{wojna polsko-bolszewicka}) vai arī ``karš ar boļševikiem'' (\pltxti{wojna bolszewicka}). Eksistē arī apzīmējums ``1920.~gada karš'' (\pltxti{wojna 1920 roku}). Faktiski gan vairāk vai mazāk aktīva kara darbība starp Poliju un Padomju Krieviju norisa arī jau gandrīz visu 1919.~gadu, kaut formāli tām abām nebija kopējas robežas un sadursmes norisa ukraiņu, baltkrievu un lietuviešu teritorijās.

Polijas Zinātņu akadēmijas Vēstures institūta un Krievijas Zinātņu akadēmijas Krievijas vēstures institūta 1998.~gadā uzsāktā savstarpējā projekta ``PSRS un Polija XX gadsimtā'' ietvaros veiktā darba rezultātā Polijas zinātnieki uz rūpīgas dokumentu analīzes pamata secināja, ka 1920.~gada karu ir jāvērtē kā poļu-padomju konfliktu, jo tobrīd Padomju Krievija nevēlējās karot un kara darbības izvēršanos noteica J.~Pilsudska vēlme radīt Poliju ``no jūras līdz jūrai''. Kļūdainais priekšstats par šā lozunga realizācijas iespējamību noveda pie tūkstošiem poļu nāves. (Protams, tikpat kļūdaina bija Padomju vadības ideja 1920.~gadā nodibināt Polijā padomju varu. Šīs kļūdas cena bija tūkstošu sarkanarmiešu nāve kauju laukos Polijā un vēlāk gūstekņu nometnēs.)

Vēlreiz jāuzsver, ka Polijas vadība izvērsa bruņotu cīņu ne tikai pret kaimiņu tautu boļševikiem, bet arī pret to nacionālistiem. Parīzes miera konferencē R.~Dmovskis prasīja Lietuvu ar Viļņu pievienot Polijai. Viņš izteicās: ``Lietuvieši arī ir poļi. Vajag tikai nošaut simt lietuviešu fanātiķu un tad būs miers.'' Tā kā J.~Pilsudskis cerēja izveidot visu ar Krieviju robežojošos republiku federāciju no Somijas līdz Gruzijai, Polijas nemitīgās nesaskaņas ar Lietuvu šim plānam bija kā sprungulis riteņos.

Pretpadomju lozungi kalpoja, lai attaisnotu vēršanos, tai skaitā arī bruņotu, pret kaimiņu tautām. Tā, 1919.~gada martā Polijas delegācija savā notā Parīzes konferencei lietuviešu ``separātiskās tendences, kam bija naids pret poļu nāciju'' (Pēc Viļņas ieņemšanas Kauņā nonākušie lietuviešu nacionāli noskaņotie politiķi atteicās piedalīties Lietuvas lielkņazistes atjaunošanā, kura atrastos Polijas aizbildniecībā) mēģināja izskaidrot ar cariskās Krievijas izplatīto pretpoļu propagandu, nemēģinot saskatīt arī poļu muižnieku vainu šo pretpoļu noskaņojumu izraisīšanā lietuviešos. Izrādot it kā pretimnākšanu lietuviešiem, delegācija paziņoja: ``sakarā ar to, ka nacionālā lietuviešu kustība, kaut arī ļoti jauna, tomēr ir guvusi ievērojamus panākumus, Polijas valdība vienlaikus uzskata, ka lietuviešu valodas teritorijai~--- Kauņas guberņai, daļai no Viļņas guberņas ziemeļrietumos no līnijas Traķi~--- Svenčone~--- Zarasi, daļai no Suvalku guberņas (Kongresa karaliste) ziemeļos no Sejniem [poļu: \pltxti{Sejny}, lietuv. \lttxti{Seinai}~--- pilsēta Polijas ziemeļaustrumos~--- V.Š.] galu galā daļai no Austrumprūsijas, kas aptver Nemunas lejteci un grīvu, ir jābūt organizētai atsevišķā zemē Polijas valsts robežās un tai ir jāsaņem speciāla iekārta, kas balstīta uz lietuviešu tautības tiesībām''. Tādejādi lietuviešu nācijai, kuras valstiskums tradīciju ziņā varēja droši konkurēt ar poļu valstiskumu, Polija piesolīja tikai autonomiju ``Polijas valsts robežās''. Tādos apstākļos gluži dabīgi jaunās Lietuvas valsts politiķi Poliju uzskatīja par savu galveno ienaidnieku un pat dzirdēt neko negribēja par vienošanos ar to.

Latvijas iedzīvotājiem der zināt, ka Polijas delegācija Parīzes konferencē ``ģeogrāfisku un ekonomisku iemeslu'' dēļ paplašināja savas pretenzijas arī līdz teritorijai ziemeļos no Liepājas, un apgalvoja, ka šos iemeslus ``vēl pastiprina fakts, ka puse Liepājas iedzīvotāju ir poļi un lietuvieši'', kad patiesībā pirms Pirmā pasaules kara Liepājā bija tikai 17\% katoļticībai piederīgu~--- galvenokārt poļu un lietuviešu~--- iedzīvotāju.

Nedaudz vēlāk Polija, (starp citu~--- tāpat kā Padomju Krievija), izmantoja kaimiņvalstu šķelšanas taktiku. Piemēram, 1919.~gada novembrī, J.~Pilsudskis piedāvāja Somijas diplomātiskajam pārstāvim Varšavā vispirms savstarpēji vienoties, bet pēc tam runāt ar pārējām Baltijas valstīm. Saskaņā ar viņa priekšlikumu, Polija atzītu par Somijas ietekmes sfēru Igauniju, bet Somija par Polijas ietekmes sfēru~--- Latviju un Lietuvu. Pēc mēneša Somijas armijas virspavēlniekam K.~Mannerheimam tika piedāvāts sadalīt atbildību par minēto valstu drošību.

Jau ar 1919.~gada sākumu poļu ekspansija lietuviešu un citās Polijas kaimiņzemēs zaudēja brālīgas palīdzības pret boļševiku despotismu auru un kļuva par to, kas tā bija jau no paša sākuma~--- cīņu par pēc iespējas lielāku teritoriju iegūšanu. Kā raksta krievu vēsturniece I.~Jažborovska, nepārliecinoši ir mēģinājumi likt vienlīdzības zīmi starp J.~Pilsudska mērķiem un ukraiņu, baltkrievu, lietuviešu un citu tautu interesēm. XX gadsimta otrā gadu desmita beigās nedrīkstēja šo ilgstošu sociālās un nacionālās attīstības ceļu izgājušos tautu likteņus risināt pēc XVIII gadsimta parauga.

Sabiedroto lielvalstu ``etnogrāfiskā'' nostāja Parīzes konferencē neatbilda Polijas vēlmēm. Tā 1919.~gada jūnijā Poļu nacionālā komiteja atzina, ka ``no Polijas interešu viedokļa būtu kaitīga patstāvīgu, ar to nesaistītu mazu valstu, tādu kā Baltkrievija vai Ukraina, pastāvēšana''. Komitejas neapmierinātību izsauca ziņa, ka 29.~martā ASV, Francijas, Itālijas, Japānas un Lielbritānijas pārstāvji vienbalsīgi pieņēma principus, kuriem bija jābūt Parīzes konferences komisiju, kas nodarbojās ar Polijas robežu noteikšanu, darba pamatā. Šie principi noteica, ka:

Polijas austrumu robežu noteikšanā jāvadās no Krievijas Pagaidu valdības 1917.~gada 17.~marta deklarācijas, kurā bija atzīts, ka:

\begin{itemize}
\item Polijai jākļūst par valsti, kuru apdzīvo galvenokārt poļi;

\item pēc etniskā sastāva un iedzīvotāju tieksmēm ``šaubīgas'' teritorijas nevar tikt pievienotas Polijai;

\item ka atsevišķos apgabalos strīdīgajās teritorijās būtu vēlams veikt iedzīvotāju aptauju;

\item galīgi jautājums par Polijas austrumu robežām tiks pieņemts pēc [neboļševistiskas] Krievijas valdības izveides, ar kuru lielvalstis šo jautājumu varētu apspriest.
\end{itemize}

Taču Antantes valstis pamatoti skeptiski vērtēja lietuviešu un baltkrievu nacionālos bruņotos formējumus, to spējas cīnīties pret boļševikiem. ASV militārās misijas loceklis Baltijā apakšpulkvedis E.~Dž.~Douls, iepazinies ar stāvokli, 1919.~gada maijā rakstīja: ``\dots{}tagad vācieši pilnībā kontrolē Lietuvu, it kā tā būtu Vācijas daļa. Lietuviešu vara ir tikai nomināla. Ja vācu karaspēks tiks izvests, radīsies briesmas, ka boļševiki pilnībā sagrābs varu.''

Bija gan arī citi vērtējumi. Lielbritānijas premjers D.~Loids Džordžs 1919.~gada 19.~februāra vēstulē britu delegācijas pārstāvim Parīzē rakstīja: ``Kas attiecas uz Poliju, Somiju un citām valstīm, kurām ir jābūt atgrieztām no Krievijas un nodotām Tautu Savienības aizbildniecībā, tad te, es domāju, mums vajadzēs sniegt ne tikai morālu, bet arī materiālu un, ja būs nepieciešams, pilnīgu militāru atbalstu šīm jaunizveidotajām valstīm, lai tās aizsargātu no boļševistiskā iebrukuma no padomju Krievijas. Tomēr pašlaik es neredzu nekādas pazīmes, ka Padomju valdībai ir kaut kādi nodomi vai vēlēšanās iebrukt šajās teritorijās.'' 1919.~gada 11.~aprīlī arī ASV pārstāvis Antantes valstu misijā Polijā ģenerālis Dž.~Kernans ziņoja prezidentam V.~Vilsonam: ``\dots{}kaut Polijā visās ziņās un sarunās tiek runāts par boļševiku agresiju, es nevarēju pamanīt it neko tamlīdzīgu. Gluži pretēji, es ar apmierinājumu novēroju, ka pat sīkas sadursmes uz Polijas austrumu robežas liecināja drīzāk par poļu agresīvajām darbībām un vēlmi pēc iespējas ātrāk ieņemt krievu zemes un pavirzīties pēc iespējas tālāk. Vieglums, ar kuru tiem tas izdevās, pārliecinoši pierāda, ka poļiem pretī nestāvēja organizēti padomju bruņotie spēki.'' ASV ģenerālis arī atzina, ka ``Neviens patlaban neuzbrūk Polijai. Gluži pretēji, skumji noraudzīties, ka valstī, kurā valda tāds trūkums, kur visiem valdības pūliņiem un visiem līdzekļiem vajadzētu tikt virzītiem iedzīvotāju stāvokļa un valsts pārvaldes uzlabošanai, visur valda militārisks gars. Šis militāriskais gars Polijas nākotnei rada lielākas briesmas nekā boļševisms''. Taču ģenerāļa-reālista balss netika ievērota.

Pārspīlētie boļševisma izplatības draudi aizsedza Sabiedroto principiālo nostāju nospraust robežas, ievērojot dažādu tautu nacionālās intereses. Lielvalstu valdību acīs pirmajā vietā nostājās Krievijas kaimiņtautu spējas apkarot boļševismu. Polijai šai jomā piederēja neapšaubāma prioritāte.

Poļi visai veikli izmantoja Rietumvalstu bailes no Vācijas un boļševiku varbūtējās nostiprināšanās Baltijā un Baltkrievijā. 1919.~gada 8.~aprīlī J.~Pilsudskis savam līdzgaitniekam Polijas Austrumu politikas realizācijā, pirmajam Polijas ārlietu ministram (no 1918.~gada 17.~novembra līdz 1919.~gada 16.~janvārim) L.~Vasiļevskim rakstīja: ``.. tuvākajā laikā es varēšu nedaudz atvērt politikas, saistītas ar Lietuvu un Baltkrieviju, durvis. Tu zini manus uzskatus par to. Tas ir ietverts viedoklī, ka mūsu Dieva pasaulē, liekas, sāk uzvarēt muldēšana par cilvēku un tautu brālību, un amerikāņu doktrīneles [acīmredzot domāti V.~Vilsona ``14 punkti''~--- V.Š.], tāpēc es nevēlos būt ne imperiālists, ne federālists līdz tam, kamēr man nav iespējas runāt par šīm lietām tā vai citādi, nopietni, bet ar revolveri kabatā.'' Šie vārdi skaidri rādīja, ka J.~Pilsudskim ``federālisms'' nebija koncepcija, balstīta uz nāciju pašnoteikšanās tiesībām, bet gan aneksiju aizmaskošanas līdzeklis. Federālisma princips bija ērts, lai izmantotu vietējo politiķu demokrātiskās idejas, politisko romantismu un ilūzijas Polijas valdības iekārotajās teritorijās, arī Rietumu lielvalstu vēlmi pievarēt boļševismu. Šis ``federālisms'' vēlāk viegli pārvērtās ceļā uz aneksiju un inkorporāciju. Tiesa gan, lai īstenotu J.~Pilsudska plānu, pastāvēja trīs galvenie šķēršļi: neskaidra lielvalstu pozīcija; vietējo nacionālo politisko elišu pretestība, kas gribēja veidot savas valstis neatkarīgas ne tikai no Krievijas, bet arī no Polijas (J.~Pilsudskis cerēja pārvarēt šo opozīciju vai nu ar militāru spēku~--- Austrumu Galīcijā, Rietumvolīnijā vai diplomātiskā ceļā~--- Viļņas jautājumā); visbeidzot, gan ``sarkanās'', gan ``baltās'' Krievijas nevēlēšanās atzīt Polijas tiesības uz minētajām teritorijām.

Kad 1919.~gada 21.~aprīlī Polijas karaspēks ieņēma Viļņu, Valsts priekšnieks un armijas virspavēlnieks J.~Pilsudskis griezās ar aicinājumu pie ``bijušās Lietuvas lielkņazistes iedzīvotājiem'', solot nodibināt brīvību ``šai novadā'' un aicinot atgriezties pie Žečpospolitas laiku ūnijas. Uzsaukumā bija arī teikts: ``Gribu dot Jums iespēju risināt iekšējās, nacionālās un ticības problēmas tādā veidā, kā paši to vēlaties, bez nekādas vardarbības un spiediena no Polijas puses'', taču vēlāk solījumi tika aizmirsti. Polijas Seima pieņemtajā t.s. ``Viļņas deklarācijā'' tika pasludināts: ``Kostjuško, Mickēviča un Trauguta tēvzeme pieder Polijai kā nedalāma sastāvdaļa''. Visas šīs nosauktās vēsturiskās personas bija dzimušas tagadējās Baltkrievijas teritorijā, bet darbojās poļu kultūras lokā.

Kāpēc Polijas armija bija iegājusi Lietuvas zemēs, vajadzēja izskaidrot arī vienkāršajiem poļu karavīriem. Tāpēc tika izdotas propagandiska rakstura brošūras. Vienas no tām autors, talantīgs poļu publicists M.~Vankovičs, rakstīja: ``Mēs atnācām šurp uz lietuviešu~--- baltkrievu zemi kā poļu armijas karavīri un Polijas valsts pilsoņi. Zemnieki, ebreji un pilsētnieki, kuri mūs šeit redz, pirmo reizi skata poļu karaspēku, bet par poļu valsti gandrīz neko nezin. Viņi ir tumši un neizglītoti cilvēki. Pie mums Polijā cilvēki gan avīzes vairāk lasa, gan ģērbjas un dzīvo pieklājīgāk, saprot vairāk lietu, kā cilvēki šeit~--- vienu vārdu sakot, pie mums Polijā cilvēki ir vairāk izglītoti. Tāpēc mums, izmitinātiem dzīvokļos, ir nepieciešams vietējiem cilvēkiem izskaidrot, kāpēc mēs esam šeit ieradušies.'' Pēc tam autors pamatoja Polijas nākotnes ģeopolitisko koncepciju. Tās būtību viņš skaidroja ar Polijas centieniem atjaunot ``ceļu no varjagiem uz grieķiem'', jo tas dotu iespēju izveidot valsti no jūras līdz jūrai un tās abas saistītu upes, ļautu izvērst ekspansiju uz austrumiem, kur ``ir daudz zemes, vairāk nekā pie mums Polijā'', bet cilvēku šais zemēs ir mazāk, tos tur varētu izvietot divas, trīs reizes vairāk nekā patlaban. Pie tam zeme Lietuvā ir četras, piecas reizes lētāka, poļu karavīri to varētu iepirkt. Pēc tam autors pamācīja kā karavīriem uzvesties: ``Parastais cilvēks daudz ko nesaprot. Viņš spriež pēc tā, ko redz deguna priekšā. Ja karavīrs viņam paņems pēdējo teļu no kūts, apkaus vistas, izsmels auzas un sienu, pēc tam vēl pašu piekaus, viņš zemniekam kļūs par ienaidnieku.'' Noslēgumā sekoja solījums pēc kādas-nekādas kārtības nodibināšanas lietuviešu un baltkrievu zemēs sarīkot Seima velēšanas, kurās varēs piedalīties visi iedzīvotāji.

Lietuvas un Baltkrievijas SPR valdība pārcēlās uz Minsku, taču 1919.~gada 8./9.~augustā poļu daļas ieņēma arī to. Šādu notikumu attīstību veicināja arī līdz 70~tūkstošu vīru lielās poļu armijas ierašanās Polijā no Francijas. 1919.~gada septembra beigās Polijas armijā, kuru pastāvīgi papildināja brīvprātīgie, bija jau ap 540~000 karavīru, no kuriem ap 230~000 bija izvietoti austrumu, poļu-padomju frontē. Viņu militāro apmācību vadīja angļu un franču militāra misija, kurā pēc padomju vēsturnieka A.~Manuseviča datiem ietilpa 9~ģenerāļi, 28~pulkveži, 668~citi virsnieki un 2~500 ierindnieku un apakšvirsnieku. Turpretī Sarkanā armija šai laikā bija spiesta rezerves virzīt pret A.~Deņikina vadīto krievu ``balto'' armiju, kura jūlijā uzsāka uzbrukumu Maskavas virzienā.

Taču tagad jau poļu karaspēka aizmugurē izvērsās partizānu kustība. Kaut daļa baltkrievu nacionālistu sadarbībā ar Poliju cerēja īstenot neatkarīgas Baltkrievijas ideju, runāja par to, ka ``Baltkrievija, Lietuva un Polija vienmēr ir bijušas viens politisks veidojums'', ka nepieciešams ``uz federācijas pamatiem apvienot baltkrievu un poļu tautas likteni'', taču drīz kļuva skaidri saskatāms, ka aiz poļu nacionālistu aicinājumiem cīnīties ``par mūsu un jūsu brīvību'', vēlējumiem redzēt ``šo novadu'' brīvu, slēpjas tikai viens mērķis~--- nostiprināt poļu valdošo ietekmi baltkrievu zemēs. Poļu ieņemtajās teritorijās tika slēgtas baltkrievu skolas, kultūrizglītības iestādes, avīzes. Pēc J.~Pilsudska pavēles Grodņā tika atbruņotas nacionālajai Baltkrievu tautas republikas padomei (Radai) pakļautās nedaudzās karaspēka daļas. Polija baltkrieviem solīja nevis federāciju, bet tikai autonomiju Polijas valsts ietvaros. Daļa baltkrievu nacionālo politiķu tika arestēta. Gala rezultātā Baltkrievu tautas republika tā arī netika nodibināta. Tiesa, galvenais cēlonis tam gan bija nevis poļu karaspēka pretdarbība, bet baltkrievu tautas masu atbalsta trūkums.

1919.~gada jūlijā poļu spēki ieņēma Baltkrievijas pilsētu Molodečno (baltkrievu \betxti{Маладзечна}), 8.~augustā~--- Minsku, bet 29.~augustā~--- Bobruisku (poļu \pltxti{Bobrujsk}, baltkrievu \betxti{Бабруйск}). Pēdējā epizodē nozīmīga loma bija ap 120~franču tankiem, kurus J.~Hallera armija atveda sev līdzi no Francijas, kur jau 1919.~gada sākumā bija izveidots poļu tanku pulks. Kā raksta mūsdienu baltkrievu vēsturnieki, poļu okupanti terorizēja vietējos iedzīvotājus, veica vardarbības un slepkavības. Sagrautajā Minskā viņi nošāva simtus cilvēku, vairāk nekā tūkstoti arestēja un ieslodzīja cietumā, tikai Igumenskas apriņķī nodedzināja ap 2~180~zemnieku sētu. No 1919.~gada jūnija līdz novembrim poļu okupanti izveda no Baltkrievijas 6~000 vagonu salaupīto materiālo vērtību, ievērojamu daļu no kurām sastādīja pārtikas produkti. Poļu karaspēka nodarītie materiālie zaudējumi tiek vērtēti uz 52~miljoniem rbļ. zeltā.

Interesanti, ka 1919.~gada 15.~septembrī Parīzes konferences Augstākajā padomē tika apspriests Polijas valdības vadītāja I.~Paderevska priekšlikums pēc Sabiedroto vēlmes nosūtīt pusmiljonu lielu poļu karaspēku ``uz Maskavu'', saņemot par to 600~000 sterliņu mārciņu dienā. Viņš solīja nepieciešamības gadījumā dažu mēnešu laikā vēl savākt papildus 480~000~cilvēku. Ja Sabiedrotie nevēlētos, lai Polija atbilstoši rīkotos (un, tātad, atteiktos apmaksāt kara izdevumus), I.~Paderevskis bija gatavs noņemt savu priekšlikumu, jo jau minētajās J.~Marhļevska un poļu pārstāvju sarunās Belovežā (\pltxti{Białowieža}) padomju puse piedāvāja Polijai ļoti izdevīgus miera noteikumus, atsakoties no daudzām teritorijām. Rietumvalstu pārstāvju attieksme pret I.~Paderevska priekšlikumu bija dažāda. Lielbritānijas valdības vadītājs D.~Loids Džordžs lika priekšā pateikt I.~Paderevskim, ka Sabiedrotie nevēlas, lai viņš sāktu darboties, jo tad viņš pieprasītu finansiālu atbalstu. Francijas premjerministrs Ž.~Klemanso norādīja, ka censties uzvarēt Krieviju ar poļu palīdzību ir sliktākais ceļš. Angļu un franču spēku iesaistīšana cīņās Krievijā nozīmētu, ka tie darbojas visas Eiropas interesēs, bet poļu spēku izmantošana izsauktu visas Krievijas pretsparu. ASV pārstāvis A.~Polks paziņoja, ka ASV ir gatava zināmā mērā palīdzēt poļiem, kuri atrodas grūtā ekonomiskā stāvoklī, taču tām nav vēlēšanās dot naudu, lai poļi gūtu iespēju vest karu. Izkristalizējās viedoklis, ka Sabiedrotie ir galvenokārt ieinteresēti vēl joprojām ārpus Vācijas izvietoto tās militāro spēku atvilkšanā uz dzimteni, to pretestības gadījumā paredzot izmantot pret tiem Polijas armiju. I.~Paderevskis tam piekrita, gan norādot, ka Lietuvas valdība atrodoties pilnīgā Vācijas ietekmē un lūdza viņu nevainot, ja poļu armijas ieņemtās teritorijas pasludinās savu piederību Polijai. Tādejādi Sabiedrotie atteicās finansēt Polijas iecerēto karagājienu uz Maskavu.

Poļu-padomju frontē pēc neveiksmīga Sarkanās armijas pretuzbrukuma oktobrī kara darbība pieklusa (Ir tad, kā raksta krievu vēsturniece I.~Mihutina, 1919.~gadā Polijas gūstā vien bija krituši jau 30~tūkstoši sarkanarmiešu).

Sākās ilgstošas Polijas un padomju Krievijas pārstāvju sarunas. To sākumu noteica Antantes valstu un Krievijas dienvidos esošo ``balto'' spēku pavēlnieka ģenerāļa A.~Deņikina nevēlēšanās atbalstīt Polijas tālākas ekspansijas plānus.

Šeit jāpasaka par \strong{Polijas un krievu ``balto'' attiecībām}. Jau Parīzes miera konferencē pilnā apjomā nostājās problēma, kā tās dalībvalstīm izturēties pret Krieviju, valsti, kura kaut Oktobra apvērsuma rezultātā bija noslēgusi separātu mieru ar Vāciju un tās sabiedrotajiem, tomēr bija ienesusi ievērojamu ieguldījumu uzvarā pār Centrālajām valstīm. Antantes valstis vēlējās pēc iespējas ātrāku padomju varas krišanu un likumīgas varas atjaunošanu Krievijā, taču tām nebija efektīvu līdzekļu, lai to panāktu. Radās projekts savākt visas pret lieliniekiem karojošas puses Marmora jūras Prinču salās, pārtraucot šai laikā Pilsoņu karu, ar lielvalstu starpniecību panākt visu pušu vienošanos un kopējas Krievijas delegācijas nosūtīšanu uz Parīzi. Padomju Krievija uzaicinājumu uz sarunām nesaņēma, taču paziņoja par savu gatavību tajās piedalīties. ``Baltie'' protestēja pret boļševiku~--- ``šīs slepkavu un laupītāju bandas'' pielaišanu sarunās. Galu galā sarunas nenotika, atbildību par to uzveļot boļševikiem, kuri nepiekrita to laikā pārtraukt kara darbību.

1919.~gada maijā Antantes Augstākā padome nolēma atzīt Krievijas Augstākā pārvaldītāja (\rutxti{Верховный правитель России}~--- tādu titulu viņam piešķīra Omskā esošās apvienotās antiboļševistiskās t.s. Direktorijas Ministru padome) admirāļa A.~Kolčaka, valdību kā vienīgi legālo Krievijas interešu pārstāvi, reizē izvirzot tai prasību garantēt, ka tās mērķis ``būs tāds pat kā sabiedroto valstu mērķi''. A.~Kolčaks savā atbildē piekrita visiem Antantes noteikumiem. Jau pieminētā Parīzes miera konferencē izveidotā Ž.~Kambona vadītā komisija, kurai bija uzdots noteikt Polijas robežas, nosūtīja A.~Kolčakam notu, kurā tika uzdots jautājums par Polijas neatkarības atzīšanu un tās robežām. Atbildes notā A.~Kolčaks paziņoja, ka Krievijas valdība uzskata par savu pienākumu apstiprināt Pagaidu valdības lēmumu par Polijas neatkarību, bet jautājumu par robežām var izlemt vienīgi Krievijas Satversmes sapulce. Krievijas sastāvā bija jāietilpst visām lietuviešu, baltkrievu un ukraiņu zemēm, arī Holmas guberņai un Austrumgalīcijai. Tādejādi A.~Kolčaks pieņēma Polijas izveides principu tās etniskajās robežās uz jau Krievijas Pagaidu valdības 1917.~gadā pasludinātajiem pamatiem, taču neko neminēja par Pagaidu valdības prasību par Polijas obligātu savienību ar Krieviju, kas nozīmēja atteikšanos no tās. Polijas suverenitāte tika atzīta bez kādiem papildus noteikumiem. Taču Polijas pusi šāda pieeja neapmierināja, tā vēlējās par savu atbalstu ``balto'' kustībai cīņā pret boļševikiem nekavējoties, negaidot Krievijas Satversmes sapulces sasaukšanu, garantijas par Polijas teritoriju paplašināšanu uz austrumiem uz Krievijas rēķina. Tiesa, poļu pusei bija dibināts pamats baidīties, ka pēc boļševiku sakāves sanākušajā Krievijas Satversmes sapulcē pārsvaru var ņemt šovinistiski noskaņojumi un Polijas neatkarība un tai vēlamās robežas tā arī netiks atzītas. (Šai sakarā var piezīmēt, ka Polijas komunisti savā izdevumā ``\pltxti{Gromada}'' 1919.~gada jūnijā pārmeta Polijas valdībai savienību ar A.~Kolčaku, kurš cīnoties par vecās kārtības atjaunošanu no Vislas līdz Vladivostokai, no Arhangēļskas līdz Persijas un Turcijas robežai. Apgalvojums neatbilda patiesībai ne attiecībā uz A.~Kolčaka mērķiem, ne Polijas attieksmi pret viņu pašu.)

Arī ``Krievijas dienvidu bruņoto spēku'' (\rutxti{Вооруженные Силы Юга России}) virspavēlnieks ģenerālis A.~Deņikins, kurš pats bija dzimis Polijas karalistē un kura māte bija poliete, būdams pozitīvi noskaņots pret Polijas valstiskuma atjaunošanu, neuzskatīja par iespējamu līdz Satversmes sapulces sasaukšanai noteikt Krievijas valsts iekārtu un robežas. 1919.~gada 18.~augustā Polijas armijas pārstāvis A.~Deņikina štābā pēc sarunas ar ģenerāli savai vadībai ziņoja, ka viņš uzskata par nepieciešamu robežu noteikšanā ņemt vērā etnogrāfiskos, vēsturiskos faktorus, kā arī abu kaimiņvalstu divpusējās intereses garantēt pastāvīgas draudzīgas attiecības.

Nosauktie Krievijas ``balto'' vadītāji centās apvienot pēc iespējas plašākus spēkus ap pretboļševiku cīņas, ``varenas, vienotas un nedalāmas Krievijas'' lozungiem. Otrais gan neapmierināja nedz Donas un Kubaņas kazakus, kuri gribēja iegūt autonomiju nākotnes Krievijā, nedz arī Ukrainas, Baltijas un Aizkaukāza nacionālistus, kuri redzēja reālu iespēju iegūt savu tautu valstisku patstāvību, nedz Polijas vadītājus, kuri jau bija iemantojuši Antantes labvēlību. J.~Pilsudskim sadarbība ar ``pustautieti'' A.~Deņikinu, kurš atzina Polijas neatkarību, bet reizē stāvēja ``vienotas un nedalāmas'' Krievijas aizstāvja pozīcijās, nebija pieņemama. Taču atklāti uzstāties pret to viņš neuzdrošinājās, ja tad Antantes valstis varētu pārtraukt Polijas armijas apgādi ar ieročiem, munīciju, pārtiku utt. Piemēram, kad Polijas armija savas virzības gaitā bija ieguvusi Krievijas pārtikas krājumus ``vairāku miljonu'' vērtībā un, pēc A.~Deņikina ziņām, gatavojās tos izvest uz Varšavu, ``baltie'' 1919.~gada decembrī protestēja ar Francijas Ārlietu ministrijas starpniecību. Francijas valdība vēlējās, lai pretboļševiku cīņas nolūkos saglabātos ``sirsnīgas'' attiecības starp Polijas valdību un ``baltajiem'' un no savas puses ieteica Polijas pusei apmierināt ``balto'' protestu. (Tai pat laikā Londonā Francijas un Lielbritānijas premjerministru sarunās ar ASV un Itālijas sūtņiem tika saskaņots viedoklis, ka ap Padomju Krieviju ir jāizveido ``dzeloņstiepļu žogs'', aiz kura Antantei nav jāiejaucas notikumos Krievijā, jāpārtrauc finansiāla palīdzība ``kārtības atjaunošanai'' tajā, bet Polija ir jāpārvērš ``barjerā pret Krieviju un šķērsli Vācijai'').

Cenšoties nekaitēt Polijas interesēm, valsts priekšnieks uzsāka zināmu tuvināšanos ar A.~Deņikinu. Polija stājās tiešās attiecībās ar ``baltajiem''. Pie ``Krievijas dienvidu bruņoto spēku'' štāba no 1919.~gada jūlija līdz 1920.~gada vasarai darbojās Polijas militārā misija ar J.~Pilsudska pārstāvi ģenerāli A.~Karņicki priekšgalā. 1919.~gada septembrī Taganrogā notika ģenerāļu A.~Deņikina un A.~Karņicka sarunas. Polijas kontrolētajās teritorijās tika atļauta brīvprātīgo vervēšana ``balto'' armijai. (Vēlāk, 1920.~gada janvārī-februārī, kad Sarkanā armija sekmīgi uzbruka Ukrainas teritorijā, J.~Pilsudskis atļāva patverties Polijas teritorijā no ``balto'' pamatspēkiem nogrieztajai ģenerāļa N.~Bredova vienībai.) Taču A.~Deņikins paziņoja, ka pēc uzvaras pār Padomju Krieviju gatavs atzīt neatkarīgās Polijas pastāvēšanu tikai bijušās Polijas Karalistes teritorijā, kas gan atbilda Parīzes miera konferences ieteiktajai Polijas robežai, bet Polijas pārstāvjiem, kuri gribēja pievienot savai valstij citu tautu apdzīvotās teritorijas, tas nebija pieņemami. Sarunas beidzās bez rezultātiem. (Tikai 1920.~gada martā A.~Deņikinu ``Krievijas dienvidu bruņoto spēku'' virspavēlnieka amatā nomainījušais barons P.~Vrangels demonstrēja elastīgāku nostāju, piemēram, pasludinot, ka nākotnē Krievijai būs federatīva uzbūve. Taču bija par vēlu, krievu ``balto'' un Polijas kopēju cīņu pret boļševikiem tā arī neizdevās noorganizēt.)

Polija šajā situācijā nebija ieinteresēta ``balto'' kustības uzvarā, kura varētu Krievijā ieviest likumību un kārtību. V.~Leņins, boļševiki ar nāciju pašnoteikšanās lozunga atzīšanu bija J.~Pilsudskim pieņemamāki. Tiesa, viņš bieži, īpaši Rietumu pārstāvju klātbūtnē, uzsvēra savu noliedzošo attieksmi pret komunismu un boļševiku realizētajām sabiedriskajām pārmaiņām. Taču valsts priekšnieks saprata, ka Antante var aizmirst visus solījumus poļiem, ja Krievijā nodibinātos stabila ``balto'' vara. Polijas atkarība no Antantes faktiski izslēdza iespēju tai vest karu par teritorijām pret ``balto'' Krieviju, bet teritoriju atņemšanu ``sarkanajai'' Krievijai Sabiedrotie atbalstītu. Tāpēc J.~Pilsudskim vairāk bija tīkama boļševiku uzvara pār A.~Deņikinu. Viņš bija arī gatavs atzīt boļševiku valdību, jo domāja, ka tā iegrūdīs Krieviju haosā, kas novedīs pie tās lielvalsts lomas zaudēšanas un nāks par labu Polijai. Starp citu, kā savās atmiņās liecinājis krievu un vēlāk padomju diplomāts J.~Solovjovs, arī ``Polijas iedzīvotāji, neraugoties uz sabiedroto pastiprinātu propagandu, vairāk nekā vienaldzīgi izturējās pret krievu baltgvardu kustību dienvidos.''

Ja poļi 1919.~gada sākumā aktīvi karoja pret boļševikiem, tad rudenī, kad ``baltie'' A.~Deņikina vadībā tuvojās Maskavai, viņi kara darbību pret ``sarkanajiem'' pārtrauca. Baltkrievijas Mikaševiču miestā (krievu \rutxti{Микашевичи}, baltkrievu \betxti{Мікашэвічы}) 1919.~gada 10.~oktobrī~--- 13.~decembrī slepus no A.~Deņikina un Antantes valstīm notika sarunas starp boļševiku pārstāvi J.~Marhļevski un Polijas delegāciju, kura uzstājās it kā Polijas Sarkanā Krusta organizācijas vārdā un kuru vadīja ievērojams poļu sabiedrisks darbinieks grāfs M.~Kosakovskis. J.~Marhļevskis piedāvāja slepenu vienošanos, poļu puse par to izrādīja lielu interesi.

Sarunu gaitā boļševiki vēlējās iegūt laiku, bet poļi pieprasīja atdot Latvijai tolaik boļševiku rokās esošo Daugavpili, pārtraukt karu pret Ukrainas Tautas Republiku. 23.~oktobrī J.~Marhļevskis paziņoja par padomju valdības gatavību atteikties no ķīlnieku ņemšanas, norādot, ka tā radās Polijas varas iestāžu nevērības rezultātā pret Krievijas Sarkanā krusta misijas nogalināšanu 1919.~gada janvārī, kā arī padomju darbinieku nošaušanām poļu ieņemtajās teritorijās. Šis paziņojums kļuva par pamatu vienošanās noslēgšanai 2.~novembrī, kurā padomju puse apņēmās atbrīvot un nosūtīt uz dzimteni Padomju Krievijā par ķīlniekiem saņemtos poļus. Abas puses parakstīja svinīgu solījumu ``uz savstarpējiem pamatiem turpmāk nepielietot ķīlnieku sistēmu''. Tika sastādīts 1~574 poļu saraksts, kuri atradās ieslodzījumā Padomju Krievijā, un 307~komunistu saraksts Polijas cietumos, kuri bija savstarpēji jāapmaina.

1919.~gada 3.~novembrī J.~Pilsudska uzticības persona kapteinis I.~Berners iesniedza J.~Marhļevskim pagaidu pamiera noteikumus, kuros solīja, ka poļu karaspēks nevirzīsies tālāk par frontes līniju, kura gāja no Ptičas (baltkrievu~--- \betxti{Пціч}, \betxti{Птыч}) upes līdz Daugavai. 21.~novembrī J.~Marhļevskis atveda Padomju valdības piekrišanu noteikumiem un piedāvāja parakstīt par to vienošanos. Kaut kāds dokuments par pagaidu pamieru gan netika parakstīts. Vienošanās bija mutiska. Polija negribēja sev saistīt ar to rokas, paturēja rīcības brīvību. J.~Pilsudskis sauca boļševiku valdību un A.~Deņikina atbalstītājus par politiskām grupām, vēlējās, lai tās arī turpmāk sekmīgi apkaro viena otru. Tomēr Polijas pārstāvis J.~Marhļevskim paziņoja: ``Mums ir svarīgi, lai jūs uzvarētu Deņikinu. Ņemat savus pulkus, sūtiet tos vai pret Deņikinu vai Judeniču [``balto'' ģenerālis, kurš vadīja cīņu pret boļševikiem pie Petrogradas~--- V.~Š.]. Mēs jūs neaiztiksim''. Tā abas puses faktiski noslēdza nosacītu pamieru. Līdz šī ``pus-pamiera'' darbības beigām 1920.~gada 25.~aprīlī, Polijas spēki to gan divreiz pārkāpa~--- 1920.~gada janvārī Latgalē, kad, atbalstot Latvijas Republikas uzbrukumu, ieņēma Daugavpili, un tā paša gada martā Baltkrievijas Poļesjē (baltkrievu \betxti{Палессе}, poļu \pltxti{Polesie}~---  kultūrvēsturisks un ģeogrāfisks novads tagadējās Baltkrievijas Republikas Gomeļas un Brestas apgabalos).

Sīkākas poļu-padomju bruņotas sadursmes Baltkrievijā gan turpinājās nepārtraukti. Jau 1919.~gada septembrī Polijas premjerministrs I.~Paderevskis, runājot par Polijas armijas panākumiem, uzsvēra, ka tā ir saņēmusi gūstā no 30~līdz 35~tūkstošiem sarkanarmiešu. Tomēr, kad ar Antantes pārstāvju starpniecību A.~Deņikins pieprasīja, lai poļi uzsāktu uzbrukumu pret ``sarkanajiem'' un pirmkārt dotu triecienu Moziras (baltkrievu \betxti{Мазыр}) virzienā, J.~Pilsudskis izvairījās no militāras palīdzības ``baltajiem''. Kā uzsver mūsdienu krievu vēsturnieki, J.~Pilsudskis no ``baltajiem'', kuros viņš saskatīja vecās, viņa nīstās Krievijas turpinātājus~--- ``imperiālistus un reakcionārus'', baidījās vairāk nekā no šīs vecās Krievijas iznīcinātājiem~--- boļševikiem. Domājams, ka galvenās tomēr bija nacionālās intereses, kā tās saprata valsts priekšnieks, nevis agrāko gadu antipātijas.

Lai nu kā, kad 1919.~gada vēlā rudenī A.~Deņikina armija atradās tikai 100~jūdžu no Maskavas, poļi nekustējās no vietas. Kā raksta mūsdienu poļu publicists B.~Skaradzinskis, ``Pilsudska Polija ar rīcību nobalsoja par labu boļševikiem un komunismam''. To izmantojot, ``sarkanie'' pārsvieda ap 40~000~karavīru pret A.~Deņikinu un sakāva to. Īpaši atzīmējams, ka uz Dienvidus fronti tika pārdislocēta Sarkanajā armijā elitārā Latviešu strēlnieku divīzija, kura veidoja trieciengrupas kodolu, kas uzbruka A.~Deņikina komandēto ``balto'' flangā. Novembrī dramatiskā vēstulē A.~Deņikins lūdza J.~Pilsudskim militāru palīdzību. Velti. Vēlāk J.~Pilsudskis ne reizi vien dižojās, ka viņš izšķīris Pilsoņu kara likteni Krievijā boļševiku labā. Tāpat uzskatīja arī virkne laikabiedru gan no ``balto'', gan ``sarkano'' nometnes. Savos memuāros vainu par boļševiku galīgo uzvaru A.~Deņikins uzvēla J.~Pilsudskim, kaut faktiski ``balto'' neveiksmi veicināja gan J.~Pilsudska bailes no ``vienotās un nedalāmās Krievijas'' piekritēju uzvaras, gan arī paša A.~Deņikina nevēlēšanās saprast un akceptēt Polijas nacionālos centienus. Poļu ģenerālis un militārais vēsturnieks M.~Kukels par A.~Deņikina neveiksmes faktoriem uzskatīja kā Sarkanās armijas rīcībā esošās lielās teritorijas, kuras tā uz laiku varēja atdot pretiniekam, tā arī ``balto'' armijas, kuru neatbalstīja ne zemnieki, ne strādnieki, ne arī nekrievu tautas, šauro sociālo bāzi.

Arī poļu vēsturnieks A.~Novaks noraida gan A.~Deņikina pārmetumus J.~Pilsudskim, gan atspēko pēdējā dižošanos. Vēsturnieks norāda, ka 40~tūkstošu karavīru pārvietošana no Polijas frontes pret Denikinu nevarēja noteikt A.~Deņikins zaudējumu un sarkano uzvaru. ``Vienkārši boļševikiem bija daudz vairāk resursu un spēka. Varbūt, ja ne pamiers starp poļiem un sarkanajiem, Deņikins būtu pavirzījies vēl 20--50 kilometrus tālāk, bet katrā ziņā boļševikiem tad vienkārši bija nesalīdzināmi vairāk resursu.''

Pēc A.~Deņikina uzbrukuma Maskavai neveiksmes poļu-padomju sarunas tika pārtrauktas. J.~Pilsudskis tagad gatavojās turpināt karu pret Padomju Krieviju patstāvīgi. 1920.~gada sākumā Polijas armijas Volīnijas frontes pavēlniekam ģenerālim A.~Ļistovskim viņš paziņoja: ``Mēs nevaram ticēt Antantes solījumiem un fabulām par Krievijas aplenkšanu ar dzeloņstiepli! Ir jāpiespiež boļševikus izdarīt mums izdevīgas piekāpšanās.'' Tiklīdz A.~Kolčaka un A.~Deņikina sakāve kļuva acīm redzama, Polijas valsts galva J.~Pilsudskis deva pavēli sākt militārās operācijas, un Polijas armija drīz bez kāda brīdinājuma atjaunoja pilna apjoma kara darbību. Ar to Polijas puse pagaidu ``pus-pamieram'' ar padomju Krieviju tā arī neļāva pārtapt īstā pamierā un mierā, bet izraisīja nu jau īstu liela mēroga karu.

\strong{Ukrainā} pēc Vācijas karaspēka atkāpšanās izraisījās iekšējās cīņas starp vietējiem nacionālistiem un boļševikiem. Polija iejaucās arī šajos notikumos.

Pēc Sen-Žermenas līguma, kuru Rietumu sabiedrotie noslēdza ar Austriju, pēdējā Čehoslovākijai atdeva Aizkarpatu Ukrainu, Polijai~--- Rietumgalīciju ar Ļvovu un Krakovu, bet no tiesībām uz Austrumgalīciju tikai atteicās. Ukrainas valdības~--- gan nacionālistiskā, gan padomju~--- pretendēja uz visu Galīciju.

Polija kategoriski iebilda pret Galīcijas, vai pat tikai tās daļas pievienošanu Ukrainai. 1919.~gada 23.~maijā Polijas Seims atzina ukraiņu centienu radīt savu valstiskumu taisnīgumu, taču tikai Krievijas impērijas teritorijā. Kā Rietumgalīcijai, tā Austrumgalīcijai bija jānokļūst Polijas sastāvā. J.~Pilsudska vadībā poļu nacionālisti centās īstenot vecveco polonizatoru lozungu ``\latxti{Natione polonus, gente Ruthenus}'' (latīņu ``Nācija poļu, tautība~--- rutēņu''.) Polijas politika 1919.~gadā grāva ukraiņu cerības nostiprināt sava valstiskuma aizmetņus.

Savukārt vietējo rutēņu (ukraiņu) nacionālā kustība, kura jau ilgstoši bija cīnījusies par pašnoteikšanos atsevišķi kā no Austroungārijas un Krievijas, tā arī no Polijas, pieprasīja atzīt Galīcijas kā ukraiņu valsts neatkarību. Austroungārijas sabrukuma rezultātā 1918.~gada nogalē Rietumgalīcijas teritorijā radās vairāki ukraiņu (rutēņu) anklāvi, piemēram, Tarnobžegas (\pltxti{Republika Tarnobrzeska}), Lemku republika (\pltxti{Rusińska Ludowa Republika Łemków)}, Komančas republika (\pltxti{Republika Komańczańska}), kuri pasludināja savu neatkarību, bet poļu spēki tās drīz likvidēja. Piemēram, pēc dažām ziņām, apspiežot Tarnobžegas zemnieku kustību, 4~tūkstoši no tiem tika nopērti, 350~arestēti un tiesāti.

Lembergā (Ļvovā) tika izveidota Ukraiņu nacionālā padome, kura pretendēja uz agrākās Žečpospolitas ukraiņu apdzīvoto zemju valdības lomu, prasīja Austrumgalīcijas kā ukraiņu valsts atzīšanu. Šo padomi atbalstīja t.s. Sečas strēlnieki (\uktxti{Українські Січові Стрільці})~--- ukraiņu nacionālās daļas bijušajā Austroungārijas armijā. 9.~novembrī Padome pasludināja Austrumgalīcijas un Bukovinas teritorijā \strong{Rietumukrainas Tautas Republiku} (\pltxti{Sachidna Ukrainska Narodna Respublika)}. Vietējo poļu iedzīvotāju mazākums, kurš uzskatīja visu Galīciju par poļu zemi, cerēja uz tās pievienošanu Polijai. Kaut lauku rajonos pārsvarā bija ukraiņi, Lembergā iedzīvotāju vairākums bija poļi un viņi ne bez pamata uzskatīja pilsētu par vienu no svarīgākajiem Polijas vēsturiskajiem un kultūras centriem. Pēc A.~Tarassa datiem potenciālajā Rietumukrainas Tautas Republikas teritorijā ukraiņi un poļi veidoja attiecīgi 60\% un 25\% īpatsvaru. Etniskais sastāvs noteica konflikta neizbēgamību Rietumukrainas Tautas Republikas un Polijas starpā.

Jau 1918.~gada 1.~novembrī Ļvovā sākās poļu un ukraiņu sadursmes. 4.~novembrī tajās iesaistījās poļu armijas daļas. Krievu vēsturnieces I.~Jažborovska un V.~Parsadanova šo militāro kampaņu nosaukušas par ``pirmo jaunveidojamās Polijas valsts ārpolitisko aktu.'' Tam var piekrist tikai ar ierunām. Polijā par valsts dibināšanas dienu taču tiek uzskatīts 11.novembris. 4.~novembrī vēl darbojās okupantu ieceltā Reģentu padome. Pareizāk būtu teikt, ka Polijas valsts šīs sadursmes saņēma mantojumā no iepriekšējās varas, taču pārtraukt tās nedomāja.

Izraisījās Polijas karš pret Rietumukrainas Tautas Republiku, kura saņēma palīdzību no Ukrainas Tautas republikas ar Direktoriju priekšgalā. Jau novembrī poļu armija okupēja Galīciju, bet Rumānija Bukovīnu.

Diemžēl uzvaru poļu daļas ``nosvinēja'' ar ebreju grautiņu Ļvovā. Rietumukrainas Tautas Republikas Nacionālās radas programmas deklarācijā bija teikts, ka ``Ukrainas valstī visiem pilsoņiem bez valodas, ticības un izcelsmes atšķirībām būs vienādas tiesības Likuma priekšā''. Tas ebrejos radīja vēlmi atbalstīt ukraiņu valsts izveidi. To veicināja arī Piekarpatu pilsētās notikušie antisemītiskie ekscesi, kuros piedalījās arī poļu karavīri. Taču, lai nenoskaņotu poļus pret ebrejiem Ebreju nacionālās padomes Ļvovā, Staņislavovā (ukr. \uktxti{Станислав} vai \uktxti{Станіславів}; poļu \pltxti{Stanisławów}; vācu \detxti{Stanislau}; mūsdienās Ivanofrankovska~--- ukraiņu \uktxti{Івано-Франківськ}, poļu \pltxti{Iwano-Frankiwsk}) un citās Galīcijas pilsētās paziņoja, ka ebreji ukraiņu un poļu kara gaitā par varu Galīcijā saglabās neitralitāti. Ukraiņu puse ar izpratni pieņēma ebreju taktiku, toties poļu puse to uzņēma kā pret Poliju vērstu izlēcienu.

Ebreju enciklopēdijā lasām: ``Poļi daudzkārt apsūdzēja ebrejus viņu pretinieku atbalstīšanā, ``antipoliskā'', ``proboļševistiskā'' un pat ``proukrainiskā'' noskaņojumā.'' (Jāatzīst, ka vispār ebreji cieta visvairāk un no itin visiem bruņotajiem spēkiem, kas šai laikā darbojās Krievijas rietumu nomalēs.)

Protams, cieta arī citi iedzīvotāji. Bijušais leģionārs un nākamais Polijas ārlietu ministrs J.~Beks vēlāk stāstīja savam tēvam, kā viņš un viņa līdzgaitnieki 1918.~gada beigās darbojās ``boļševizētajā'' Ukrainā: ``Ciemos mēs nogalinājām visus pēc kārtas un pie mazākajām aizdomām par nevaļsirdību visu nodedzinājām. Es personīgi strādāju ar [šautenes] laidni.''

Rietumukrainas Tautas republikas bruņotie spēki atkāpās uz austrumiem un apvienojās ar Ukrainas Tautas republikas armiju. Rietumukrainas valdība, atrodoties Staņislavas pilsētā, 1919.~gada 22.~janvārī parakstīja līgumu (``\uktxti{Акт злуки}'') ar Ukrainas Direktoriju, atzina tās virsvaru un paziņoja par apvienošanos ar Ukrainas Tautas Republiku vienotā ukraiņu valstī. Taču 1919.~gada martā ukraiņu spēki tika atspiesti no Ļvovas.

Cenšoties samierināt karojošās puses, ASV Prezidents V.~Vilsons, Francijas un Itālijas premjerministri Ž.~Klemanso un V.~Orlando piedāvāja Polijai un Ukrainai pārtraukt karu Galīcijas frontē un noslēgt pamieru, par pamatu ņemot esošo frontes līniju. Polijai tika paziņots, ka poļu-ukraiņu robežu noteiks Antante, bet kara darbībās laikā ieņemtās teritorijas tās lēmumu neietekmēs. Ukrainas puse Antantes valstu priekšlikumam piekrita, bet Polijas~--- noraidīja to. Polijas valdība nolēma apspēlēt ukraiņu nacionālistus un ``pāri'' Direktorijai noslēdza pamieru ar Padomju Ukrainas valdību, saskaņā ar kuru demarkācijas līnijai bija jāiet pa Zbručas (poļu \pltxti{Zbrucz}, ukraiņu \uktxti{Збруч}, vācu \detxti{Sbrutsch}) upi (Dņestras pieteku, kura vairākus gadsimtus kalpoja par Turcijas un Polijas, bet no XVIII gadsimta beigām par Austrijas un Krievijas robežu.) Ar to ukraiņu boļševiki piekrita atdot Polijai arī ukraiņu apdzīvotās teritorijas. Rietumukrainas Tautas Republikas armija tomēr turpināja karot ar Polijas armiju līdz 1919.~gada maijam, kad t.s. ``Hallera armija'' no Francijas tika pārdislocēta uz Galīciju un izspieda ukraiņu spēkus no tās teritorijas aiz Zbručas upes.

Bruņotā cīņa par Austrumgalīciju maksāja ap 15~000 dzīvību ukraiņiem un ap 10~000 poļiem. Cieta arī neitrālie ebreji. Kā norāda ebreju vēsturnieks Š.~Etingers, Hallera armijas karavīri atkal ielauzās ebreju mājās, slepkavoja un laupīja, izvaroja sievietes, apgraizīja vaigu bārdas un pilnbārdas vīriešiem, dažādi ņirgājās par upuriem. To skaits sasniedza vairākus tūkstošus. Ziņas par šīm zvērībām lika ASV prezidentam V.~Vilsonam nosūtīt uz Poliju poļu un ebreju attiecību izpētei komisiju, kuras priekšgalā atradās ievērojamais ebreju tautības amerikāņu politiskais darbinieks H.~Morgentau. Tādas pat komisijas tika nosūtītas no Lielbritānijas un Francijas. Tās savā darbībā sastapās ar poļu iedzīvotāju pretsparu, poļu šovinistiskā prese ``sacēla troksni'' par starptautisko ebreju aprindu jaukšanos Polijas lietās, tika publicēti ķengājoši raksti par Polijas ebrejiem. Komisiju darbība nespēja mainīt poļu un ebreju attiecības.

Taču 1919.~gada pavasarī Sarkanā armija sakāva Ukrainas Direktorijas karaspēku, 1919.~gada martā nodibinājās Ungārijas Padomju republika (\hutxti{Magyarországi Tanácsköztársaság}) un tikai Polija stāvēja ceļā starp revolucionāro Krieviju un Ungāriju, tāpēc 1919.~gada 25.jūnijā Lielbritānijas, Francijas, ASV, Itālijas ārlietu ministru padome Versaļas konferencē Eiropas likteņus lemjošo Rietumu lielvalstu vārdā atļāva Polijai okupēt Austrumgalīciju, kaut jautājums par tās piederību vēl nebija galīgi izlemts. Jūlija vidū Polijas karaspēks sasniedza Zbručas upi, Rietumukrainas Tautas republikas bruņotie spēki atkāpās un pievienojās A.~Deņikna komandētajai ``balto'' armijai, kura, neatzīstot Ukrainas neatkarību, karoja pret Direktorijas spēkiem. 21.~novembrī Antantes Augstākā padome nolēma, ka Polija var saņemt mandātu uz 25~gadus ilgu Austrumgalīcijas pārvaldi, bet ukraiņu apdzīvotajiem apgabaliem jāsaņem autonomijas tiesības. Pēc 25~gadiem Tautu Savienībai atkal būtu jālemj par apgabala likteni. Vēlāk padomju-poļu kara gaitā 1920.~gadā uz īsu laiku (jūlijā-septembrī) tika pasludināta Galīcijas Sociālistiskā padomju republika, bet pēc Rīgas miera noslēgšanas 1921.~gadā (par ko skat. zemāk šai pašā apakšnodaļā) šī teritorija ietilpa Polijas sastāvā.

(Aizsteidzoties notikumiem priekšā, šeit tomēr der pateikt, ka jautājums par Austrumgalīcijas piederību tika cilāts arī turpmāk. Lai it kā ievērotu Tautu Savienības nostāju, 1922.~gada 26.~septembrī Polijas Seims atzina Austrumgalīcijas ukraiņiem autonomijas tiesības. To ievērojot, 1923.~gada 14./15.~martā sūtņu konference Parīzē Francijas spiediena rezultātā pieņēma lēmumu par Polijas austrumu robežām, tai skaitā par Austrumgalīcijas nodošanu tai. Taču autonomija Austrumgalīcijai tā arī nekad netika piešķirta. Padomju Ukrainas un Krievijas vadība šai sakarā demonstrēja nostāju, ka tā nav samierinājusies ar ukraiņu un baltkrievu teritoriju nonākšanu Polijas sastāvā. Tā Ukrainas Sociālistiskās Padomju Republikas ārlietu tautas komisārs H.~Rakovskis 1923.~gada 12.~marta notā paziņoja lielvalstu valdībām, ka Polijas veiktā Austrumgalīcijas okupācija ir vardarbīgs akts un PSRS valdība uzskata par nelikumīgu jebkāda režīma izveidi tajā bez savas iepriekšējas piekrišanas un pašu iedzīvotāju aptaujas. 13.~martā analoģisku paziņojumu lielvalstīm iesniedza KSFPR ārlietu komisārs G.~Čičerins. Viņš arī brīdināja, ja Austrumgalīcijas liktenis tiks lemts bez padomju republiku līdzdalības, tad tā rezultātā ``nākotnē radīsies jaunu sadursmju perēkļi''. 1923.~gada 5.~maijā PSRS Tautas Komisāru padome paziņoja, ka PSRS neatzīst Antantes valstu tiesības lemt citu tautu un teritoriju likteņus, sūtņu konferences Parīzē lēmumus par Viļņas apgabala, Rietumukrainas un Rietumbaltkrievijas apstiprināšanu Polijas sastāvā. 1924.~gada 12.~augustā Lielbritānijas un PSRS sarunās Londonā padomju delegācija atkal apšaubīja Austrumgalīcijas atrašanos Polijas sastāvā, uzsvēra, ka tā uzskata par nepieciešamu piešķirt Austrumgalīcijas iedzīvotājiem pašnoteikšanās tiesības. Polijas valdība, atsaucoties uz 1921.~gada Polijas un Padomju Krievijas miera līgumu, padomju puses pretenzijas noraidīja kā iejaukšanos Polijas iekšējās lietās. Tik tiešām PSRS valdības nostāja bija pretrunīga: no vienas puses tā bija noslēgusi ar Poliju miera līgumu, kurā bija noteiktas arī abu valstu robežas, bet no otras puses neatzina Austrumgalīcijas piederību Polijai. Austrumgalīcijas jautājumam PSRS valdība pievērsās arī citos savos paziņojumos. Notikumu attīstība rādīja, ka ar spēku citam atņemtais ieguvums, mainoties spēku samēram, arvien ir apdraudēts.

1919.~gada 8.~decembrī Antantes Augstākā padome rekomendēja par Polijas austrumu robežu uzskatīt līniju, kas sakrita ar Bugas upi un ar poļu etnogrāfiskās izplatības austrumu robežu, un faktiski bija arī Polijas-Lietuvas valsts trešās dalīšanas robežlīnija. (Lielbritānijas ārlietu ministrs lords Dž.~N.~Kērzons 1920.~gadā piedāvāja poļiem un to pretiniekiem noslēgt pamieru uz šīs līnijas un tāpēc literatūrā tā iegājusi kā \strong{Kērzona līnija}.) Uz rietumiem no tās paliktu zemes ar pārsvarā poļu tautības, bet uz austrumiem ar pārsvarā ukraiņu, baltkrievu un lietuviešu tautības iedzīvotājiem.

Arī Padomju puse, spaidīgo Pilsoņu kara apstākļu spiesta, oficiāli piedāvāja Polijai uzsākt miera sarunas. 1919.~gada 22.~decembrī tās ierosināja KSFPR ārlietu tautas komisārs G.~Čičerins. 1920.~gada 28.~janvārī arī KSFPR Tautas komisāru padome griezās pie Polijas tautas un valdības ar radiogrammu, kurā apstiprināja bezierunu Polijas neatkarības atzīšanu un piedāvāja uzsākt miera sarunas. Šis piedāvājums bija jau piektais kopš poļu-padomju kara sākšanās. Padomju valdības piedāvātais robežas variants atradās 30 līdz 150 km tālāk uz austrumiem, nekā 14 mēnešus vēlāk Polijas ar Rīgas miera līgumu noteiktā robeža. (Šai sakarā var piebilst, ka boļševiku vadītājus šai laikā ne sevišķi uztrauca ukraiņu un baltkrievu tautu sadalīšana. Acīmredzot tas tika uzskatīts par pagaidu risinājumu, kuru ar laiku varētu mainīt revolūcijas jauni uzplūdi pasaulē.) Vienīgais nosacījums, ko padomju puse izvirzīja Polijai bija Ukrainas Sociālistiskās Padomju Republikas atzīšana. 2.~februārī KSFPR ``parlaments''~--- Viskrievijas Centrālā Izpildu Komiteja no savas puses griezās ar miera piedāvājumu pie Polijas tautas un apgalvoja, ka ``Polijas brīvība ir nepieciešams nosacījums Krievijas brīvai attīstībai'', ka krievu komunisti ``nedomā un pat nevar domāt par komunisma vardarbīgu uzspiešanu svešām zemēm'', ka nav neviena teritoriāla, ekonomiska vai cita jautājuma, kuru nevarētu atrisināt mierīgi, sarunu, savstarpējas piekāpšanās un vienošanās ceļā. 22.~februārī noslēgt mieru Polijas valdībai piedāvāja Ukrainas sociālistiskā padomju republika. Poļu vēsturnieks P.~Lossovskis raksta, ka Krievija no decembra beigām līdz sarunu izjukšanai īstenoja ``miera uzbrukuma'' taktiku. To ar saviem paziņojumiem atbalstīja arī poļu komunisti gan pašā Polijā, gan padomju Krievijā. Arī PPS vadība nosūtīja savām vietējām organizācijām aicinājumu pieprasīt miera sarunu ar Padomju Krieviju uzsākšanu. Taču Polijas vadība nevēlējās slēgt mieru, it kā baidoties no boļševiku divkosības. Vairums poļu vēsturnieku nevēlas kritizēt Polijas pusi par šādu politiku. Tikai daži mūsdienu vēsturnieki to iedrošinās. Tā, piemēram, M.~Pruščinskis raksta: ``Šodien no pagājušo gadu perspektīvas šķiet neapstrīdami, ka padomju miera priekšlikumi 1919.--1920.~gadu mijā bija neliekuļoti''.

4.~februārī Polijas ārlietu ministrs S.~Pateks atbildēja uz G.~Čičerina notu, ka Polijas valdība izskatīs padomju puses priekšlikumus. Tai pat laikā S.~Pateks konsultējās Parīzē ar lielvalstu pārstāvjiem, uzsverot, ka viņš ir pret mieru ar Padomju Krieviju, bet Polijas pozīcija būs atkarīga no lielvalstu nostājas. D.~Loids Džordžs uzskatīja, ka lēmums jāpieņem pašai Polijai, taču neieteica sākt uzbrukumu. Viņš 1920.~gada 26.~janvārī personīgā sarunā ar S.~Pateku pamatoja savu uzskatu, ka ``ārpus savas teritorijas boļševiki nerada nopietnus militārus draudus''. Galvenās grūtības miera noslēgšanai starp Poliju un Padomju Krieviju britu premjers saskatīja apstāklī, ka ``poļu armija ir pavirzījusies tālu pāri etnogrāfiskajām robežām un ieņēmusi ievērojamas teritorijas, kuru iedzīvotāji ir pārsvarā krievi.'' Ja poļi mēģinātu noslēgt taisnīgu mieru, bet boļševiki atteiktos to darīt, vai arī pēc līguma noslēgšanas to pārkāptu, Lielbritānija bija gatava sniegt Polijai visu nepieciešamo palīdzību. Toties Francijas un Itālijas pārstāvji ieteica turpināt karu, kas vairāk atbilda J.~Pilsudska vēlmei. Francijas premjerministrs A.~Miljerāns īpaši bija norūpējies par to, ka Padomju Krievijas un Polijas miera noslēgšana vēl vairāk pasliktinātu jau tā neapskaužamo krievu ``balto'' stāvokli. Tomēr Francija nevēlējās rīkoties bez Lielbritānijas atbalsta. 24.~februārī Antantes Augstākā padome nolēma pagaidām nedibināt diplomātiskas attiecības ar Padomju Krieviju. Tādejādi Polijai palika zināma rīcības brīvība, taču J.~Pilsudskis to noteikti virzīja uz jaunu teritoriju iekarošanu un ietekmes palielināšanu. Poļu politiskais darbinieks St.~Kozickis liecinājis, ka jau februārī S.~Pateks viņam pastāstīja, ka valsts priekšnieks gatavojas karagājienam pret Krieviju, jo ir pārliecināts, ka ``viņa karaspēks virzīsies uz austrumiem tik tālu, cik viņš uzskatīs par vajadzīgu''.

Neguvis atbalstu no Sabiedrotajiem, cīņā pret Padomju Krieviju J.~Pilsudskis bija gatavs izmantot arī krievu antiboļševistisko spēku, ``vienotās un nedalāmās Krievijas'' piekritēju palīdzību. 1920.~gada janvārī J.~Pilsudskis veda sarunas ar Varšavā iebraukušajiem krievu boļševiku pretiniekiem N.~Čaikovski un bijušo Krievijas eseru kaujas organizācijas līderi, Krievijas kara ministra biedru (vietnieku) Krievijas Pagaidu valdībā 1917.~gadā, vēlāk aktīvu cīnītāju pret boļševikiem B.~Savinkovu par valstu savienības izveidi, kas būtu vērsta pret Padomju Krieviju. 18.~janvārī viņš parakstīja vienošanos ar nosauktajiem ``balto'' pārstāvjiem. Pēc vienošanās Polija apņēmās atbalstīt ``balto'' cīņu Krievijā un atteicās no vēlmes atjaunot Žečpospolitu 1772.~gada robežās. Jau janvārī J.~Pilsudskis uzsāka B.~Savinkova vadītās Krievu politiskās komitejas (\rutxti{Русский политический комитет}) organizācijas finansēšanu. Savukārt šie krievu politiķi atteicās no nodoma atjaunot Krieviju pirmskara robežās, piekrita atzīt Igaunijas, Latvijas, Somijas un Ukrainas neatkarību, kā arī plebiscīta par valstisko piederību sarīkošanu Baltkrievijā un Lietuvā. Februārī J.~Pilsudska pārstāvis turpināja sarunas par sadarbību ar krievu emigrantiem Parīzē.

1920.~gada jūlijā Krievu politiskā komiteja B.~Savinkova vadībā sāka Polijā organizēt cīņai pret Sarkano armiju t.s. 3.~krievu armiju (faktiski~--- korpusu), kuras sastāvā bija jāietilpst divām divīzijām~--- kājnieku un kavalērijas. Sākotnēji ar tās formēšanu nodarbojās ģenerālis P.~Glazenaps, vēlāk ģenerālis B.~Peremikins. Tiesa, armija bija neliela. Jau izveidotā tajā bija 4~tūkstoši cilvēku, no tiem 700~virsnieku. Tās vadība novilcināja došanos uz fronti, prasot krievu daļas nosūtīt pie P.~Vrangeļa. 1920.~gada oktobra sākumā armija sāka uzbrukumu Dienvidukrainā ar uzdevumu savienoties ar P.~Vrangeļa karaspēku, taču pēc divām dienām 1920.~gada 12.oktobrī Rīgā tika noslēgts Padomju Krievijas pamiers ar Poliju. Polijas armija kara darbību pārtrauca. 3.~krievu armija zem Sarkanās armijas triecieniem atkāpās uz Polijas teritoriju, kur tika atbruņota, tās sastāvs internēts. 1921.~gadā Francijas spiediena rezultātā Polija krievu karavīrus atbrīvoja no internēto nometnēm. Lielākā daļa boļševikiem opozīcijā esošo krievu politiķu nosodīja B.~Savinkova un viņa līdzgaitnieku vienošanos ar poļiem.

Reizē Polijas vadība sagatavoja savu miera sarunu sākšanas nosacījumu projektu. Poļu militārais vēsturnieks Z.~Zaluskis ir uzsvēris, ka šajā dokumentā gandrīz katrs punkts sākās ar vārdiem ``padomju puse apņemas'' un neviens ar vārdiem ``poļu puse apņemas''. Padomju pusei bija jāatzīst ``bijušās Krievijas impērijas teritorijā radušos neatkarīgo nacionālo valstu, kurām dotajā momentā bija savas valdības [neatkarīgi no tā, vai tām bija tautas atbalsts~--- V.Š.] neatkarība'', jāsamaksā Polijai kompensācija zeltā un preču piegāžu veidā par Polijas zaudējumiem tās dalīšanu rezultātā, jāļauj Krievijas teritorijā ar eksteritorialitātes tiesībām darboties daudzām Polijas iestādēm, jānodrošina Polijas pilsoņiem iespējas Krievijā nodarboties ar tirdznieciski-rūpniecisku darbību, neievērojot padomju likumdošanu. Svarīgs bija paragrāfs, kurš prasīja nekavējoties izvest padomju karaspēku ārpus Polijas 1772.~gada robežām. Uz šī projekta var izlasīt J.~Pilsudska palīga austrumu jautājumos L.~Vasiļevska piezīmi: ``Prasības ievērojamā mērā formulētas tā, it kā tās būtu adresētas valstij, kuru Polija pilnībā sakāvusi. Tā ir būtiska iejaukšanās Krievijas iekšējās lietās, kas ierobežo tās suverenitāti\dots{}'' Izskanēja arī Padomju Krievijai, kura vēl karoja ar ``baltajiem'', iesniedzama prasība nekavējoties demobilizēt pusi savas armijas, nākamā gadā laikā vēl vairāk samazināt tās sastāvu. (1920.~gada 17.~aprīļa vidū Polijas armijas informācijas un pretizlūkošanas dienesta darbinieku sanāksmē tika paziņots, ka padomju armijas demobilizācijas prasība esot bijusi ``nepieciešama, jo tikai demobilizācija ir vienīgā reālā miera noteikumu izpildes garantija'', ka miera līgums ar Padomju Krieviju pastāvēs vienīgi uz papīra, kamēr pastāvēs Krievijas armija, bet to likvidēt varot tikai divos ceļos: vai demobilizācijas vai sakaujot kaujas laukā.)

Vēl nesaņēmis Polijas nosacījumus, 1920.~gada 27.~februārī V.~Ļeņins rakstīja Krievijas Republikas Revolucionārajai kara padomei: ``\dots{}visas pazīmes liecina, ka Polija mums uzstādīs absolūti neizpildāmus, pat nekaunīgus noteikumus'' un pieprasīja visu izmanību veltīt Rietumu frontes nostiprināšanai.

Tā, neraugoties uz padomju miera priekšlikumiem, 1919./1920.~gada ziemā abas puses gatavojās jaunām cīņām. Vienlaikus abu valstu~--- gan Padomju Krievijas, gan Polijas~--- līderi V.~Ļeņins un J.~Pilsudskis sniedza intervijas ārzemju žurnālistiem, cenšoties ietekmēt sabiedrisko domu sev par labu. Neviens no viņiem negribēja uzņemties atbildību par karu. 6.~martā G.~Čičerins nosūtīja Polijas ārlietu ministram notu, kurā izteica nožēlu, ka Polijas nav devusi savu atbildi uz padomju priekšlikumiem, bet 5.~martā poļu spēki atsākuši uzbrukumu. (Formālu atbildi uz padomju priekšlikumiem ar norādi, ka ``tiek izstrādāti miera ar Padomju Krieviju noslēgšanas noteikumi'', Polija deva tikai 27.~martā. Var piebilst, ka poļu notas nosūtīšanas dienā Francijas militārās misijas Varšavā vadītājs ģenerālis P.~Anri kā reiz beidza izstrādāt Polijas armijas operāciju plānu austrumos.) V.~Ļeņins bija nobažījies par kara iespējamību. 1920.~gada 11.~martā viņš šifrētā telegrammā paziņoja Ļ.~Trockim: ``maz iespēju izbēgt no kara'', ``ir jābūt gataviem uz visļaunāko'', ka arī padomju diplomāts M.~Ļitvinovs brīdina, ka Polija karos, jo padomju miera iniciatīvas uzskata par vājumu.

Polijā un tās Satversmes Seimā bija izplatīti dažādi, arī negatīvi noskaņojumi pret J.~Pilsudska plāniem Austrumu frontē. A.~Širokorads, raksta, ka 1919.~gada beigās un 1920.~gada sākumā Varšavā klīda baumas par J.~Pilsudska vēlēšanos kļūt pat Polijas karali. Iespējams, ka aizdomas par to stiprināja tas, ka Polijas kara ministrs J.~Lesņevskis februārī~--- martā izveidoja speciālu komisiju Polijas karaļu regāliju meklēšanai.

Arī daļa Seima deputātu skeptiski vērtēja J.~Pilsudska plānus. Padomju puses miera priekšlikumi taču bija izdevīgi Polijai, tās rokās paliktu gandrīz visas 1919.~gadā ieņemtās teritorijas, tai skaitā gandrīz visa Baltkrievija. Pragmātiski noskaņotie nacionāldemokrāti aicināja pieņemt padomju priekšlikumus. 1920.~gada 11.~februārī pie J.~Pilsudska ieradās PPS, bet 20.~februārī ļudoviešu partijas (\pltxti{Stronnitctwo Ludowe ``Wyzwolenie''}) delegācijas, kuras prasīja noslēgt mieru. Polijas pretizlūkošanas ziņojumā par PPS nostāju pret karu austrumos bija minēti fakti par to, kā redzami sociālistu darbinieki I.~Dašinskis, M.~Nedzjalkovskis izteicās pret nepārdomātu kara vešanu austrumos. Otrais no minētajiem paziņoja, ka Polijai piešķirot pašnoteikšanās tiesības visām austrumu kaimiņtautām, būtu iespējams izveidot ``Austrumeiropas tautu bloku, kura dvēsele būtu Polija''. Viņš arī uzsvēra, ka ``kaimiņtautu demokrāti tieksies tuvināties Polijai tikai tad, ja Varšavā pie varas būs īsteni demokrātiska valdība''. Taču pilsonisko partiju un arī PPS nostāja poļu-padomju kara jautājumā bija nekonsekventa, atkarīga no tā, kura puse dotajā brīdī guva panākumus.

J.~Pilsudskis bija spiests Seimam sniegt paskaidrojumus par valsts politikas mērķiem austrumos. Cita starpā viņš paziņoja: ``Kara mērķis~--- Polijas aizsardzība un mūsu kaimiņu atbrīvošana no boļševikiem, iespējas piešķiršana šiem kaimiņiem pašiem lemt savu likteni. Mieru patlaban noslēgt nevar, jo nav garantijas, ka viņi ievēros mieru, bet nedrošības stāvoklis ir vienāds ar kara stāvokli.'' Tādejādi, gan boļševiki, gan J.~Pilsudskis solīja ``kaimiņiem'' iespēju ``pašiem lemt savu likteni''. Nebija pārliecinoša J.~Pilsudska norāde uz boļševiku, kuri bija iestiguši Pilsoņu karā, iespējamo miera neievērošanu, arī ``nedrošības stāvokļa'' pielīdzināšana kara stāvoklim. Taču J.~Pilsudskis panāca, ka Seims 24.~februārī pieņēma rezolūciju ar prasību izvest padomju karaspēku ārpus 1772.~gada Žečpospolitas robežām un referenduma par pašnoteikšanos sarīkošanu ``neitrālajā joslā''~--- faktiski Ukrainā. J.~Pilsudskis cerēja panākt nomināli patstāvīgas, bet faktiski no Polijas atkarīgas Ukrainas izveidi, kura nošķirtu Poliju no Krievijas. Krievu vēsturnieks A.~Vorobjovs savā disertācijā ir uzsvēris, ka Polijas austrumu politikas pamatakmens bija neatkarīgas Ukrainas atbalstīšana, bet nevis boļševisma sagrāve, pret kuru J.~Pilsudskis izturējies daudz iecietīgāk nekā pret ``vienotās un nedalāmās'' Krievijas piekritējiem. Tāpēc arī poļu karaspēks ticis 1920.~gadā virzīts uz Kijevu. Tomēr secinājums ir jāprecizē. J.~Pilsudskis vēlējās, lai ``neatkarīgā'' Ukraina būtu atkarīga no Polijas.

1920.~gada sākumā Polijas armijā bija jau ap 700~tūkstošu vīru. Rietumu sabiedrotie materiāli apgādāja Polijas armiju. Francija piešķīra kredītu 1~miljarda franku apjomā. 1919.~gada gaitā tikai no Francijas uz Poliju nogādāja 1~106 lielgabalus, 3~401 ložmetēju, 321~000 šauteņu, 350 lidmašīnu, 800 kravas auto, 4~500 ratu (literatūrā ir arī citi, nedaudz atšķirīgi dati), lielu daudzumu munīcijas, mundieru, apavu, citus kara materiālus. 1919.~gadā līdz 1920.~gada vasarai Francija piešķīra pusmiljardu franku lielu kredītu. ASV piešķīra Polijai kredītu \$159,6 apmērā un 1920.~gada pirmajā pusē piegādāja vairāk nekā 200~bruņumašīnu, vairāk nekā 300~lidmašīnu, 3~miljonus mundieru, 4~miljonus pāru apavu. Piegādes nāca arī no Lielbritānijas. Interesants fakts~--- Polijas apbruņošanai lieti noderēja arī sakautās un atbruņotās Vācijas ieroči. No tās arsenāliem poļi saņēma 1~200 ložmetēju, ap pusotru tūkstoti lielgabalu u.c. ieroču.

Tas viss deva pamatu agresīvi noskaņoto poļu politiķu un militāristu optimistiskām cerībām. 1920.~gada februārī franču žurnālista jautāts par Polijas plāniem attiecībā pret Ukrainu, Baltkrieviju un Lietuvu, J.~Pilsudskis atbildēja: ``Uz durkļiem nesam šīm nelaimīgajām zemēm brīvību''.

Sarkanās armijas pavēlniecība rēķinājās ar iespējamo Polijas armijas aktivizēšanos austrumos pēc tam, kad tās daļas atbrīvojās Poznaņas apgabalā un Galīcijā, taču nepareizi noteica poļu uzbrukuma galveno virzienu, uzskatot, ka galvenās cīņas tiks izcīnītas Rietumu frontē Baltkrievijas virzienā un turp arī nosūtīja rezerves. Polijas Ģenerālštābs noskaidroja, ka no 1919.~gada decembra beigām līdz 1920.~gada vidum uz Rietumu fronti no A.~Deņikina frontes (dienvidos) bija pārdislocēti ap 32~000 durkļu, 2 tūkstoši zobenu un 60 lielgabalu. No tā poļu militāristi izdarīja secinājumu, ka martā-maijā boļševiki gatavojas uzbrukumam.

Šai sakarā var pieminēt poļu vēsturnieku B.~Musialu, kurš 1985.~gadā emigrēja uz VFR, ieguva tur izglītību, pēc tam atgriezās Polijā. 2008.~gadā viņš publicēja monogrāfiju ``Kaujas lauks Vācija: Staļina kari pret Rietumiem'' (\detxti{Kampfplatz Deutschland. Stalins Kriegspläne gegen den Westen, Berlin: Propyläen, 2008}), kurā apgalvo, ka kopš 1920.~gada janvāra padomju vadība gatavoja karu pret Poliju. Tiesa, to viņš ``pierāda'' visus avotos esošos norādījumus uz padomju puses pasākumiem, virzītiem uz aizsardzības stiprināšanu, iztulkojot kā padomju īsteno nolūku aizmaskošanu. Pēc B.~Musiala teiktā, J.~Pilsudskis, pateicoties izlūkošanai, savlaicīgi uzzinājis par Sarkanās armijas gaidāmo uzbrukumu Polijai, un tāpēc pats pavēlējis uzsākt uzbrukumu Ukrainā. Situāciju 1920.~gadā šis vēsturnieks pielīdzināja 1941.~gada notikumiem, kad J.~Staļins gatavojis uzbrukumu Vācijai, bet Ā.~Hitlers to apsteidzis. Taču tā kā Vācija 1941.~gadā neko nezinājusi par PSRS paredzamo uzbrukumu, no Ā.~Hitlera puses tas nebijis preventīvs karš, turpretī 1920.~gadā Polijas trieciens bijis apsteidzošs (preventīvs).

Kā pamatoti norādījis viņu kritizējušais vācu vēsturnieks J.~Gancenmillers, B.~Musiala konstatējums, ka Polijas uzbrukums Padomju Krievijai, par kuru, kā redzams jau viņa paša citētajos avotos, padomju vadība tika brīdināta, bijis tai negaidīts, nonāk klajā pretrunā ar viņa iepriekšminēto avotu interpretāciju. Toties poļu vēsturnieks neko nesaka par Polijas gatavošanos karam, tā vietā tikai apgalvojot, ka J.~Pilsudskim neesot bijis agresīvu plānu. Visus, kuri nepiekrīt B.~Musiala viedoklim, viņš uzskata par padomju propagandas sekotājiem, kuri vedot mērķtiecīgu kampaņu pret nesatricināmiem antikomunistiem.

5.~un 6.~martā notika sava veida izlūkošana ar kauju. Baltkrievijas dienvidos Polijas armija pārrāva padomju fronti un ieņēma svarīgos Moziras (baltkrievu \betxti{Мазыр}) un Kaļinkoviču (baltkrievu \betxti{Калінкавічы}) dzelzceļa mezglus. Sarkanā armija šai laikā mēģināja uzbrukt Ukrainā, taču neveiksmīgi, un bija spiesta pāriet aizsardzībā. Padomju vadība šīs sadursmes vēl neuztvēra kā pilna apjoma karu, kaut 7.~martā Padomju Krievijas un Ukrainas ārlietu tautas komisāri iesniedza kopīgu notu Antantes valstīm sakarā ar poļu uzbrukumu.

8.~martā Polijas valdība apstiprināja miera sarunu pamatnostādnes. Tās paredzēja cīņu par 1772.~gada robežām, reizē arī ņemot vērā iespēju tās gan nesasniegt, bet tomēr pārbīdīt robežu Polijai par labu. Taču arī tad vēl šīs pamatnostādnes bija jāprecizē speciālai Ārlietu ministrijas komisijai. Beidzot 27.~martā, lai demonstrētu poļu tautai it kā gatavību miera noslēgšanai ar Padomju Krieviju, Polijas valdība sankcionēja sarunu sākšanos uz frontes līnijas esošajā Borisovā (baltkrievu \betxti{Барысаў}, poļu \pltxti{Borysów}~--- pilsēta Baltkrievijā, Minskas apgabalā), kara darbību pārtraucot tikai šajā šaurajā iecirknī. Padomju priekšlikums vest sarunas Parīzē, Londonā (kas nozīmētu Padomju Krievijas neoficiālu atzīšanu no Rietumvalstu puses), Petrogradā, Maskavā vai Varšavā, bet pamieru izsludināt visā frontē, tika noraidīts. Padomju un krievu vēsturnieki to skaidro ar Polijas vadības vēlmi nepieļaut Sarkanās armijas darbību Baltkrievijā, kur bija koncentrēts ap 85\% visu rietumos esošo padomju spēku, kad paši poļi izvērstu uzbrukumu Ukrainā.

Īslaicīgās sarunās Borisovā poļu puse uzstāja uz Žečpospolitas 1772.~gada robežu atjaunošanu, kam piekrišanu nesaņēma. Tajā brīdī poļu armija atradās krietni tālāk uz austrumiem no t.s. Kērzona līnijas, tāpēc Polijas vadība ar argumentu, ka, pieņemot padomju priekšlikumus, ārpus Polijas robežām paliktu daudz poļu iedzīvotāju, netaisījās atstāt kaut daļu no ar spēku ieņemtajām teritorijām un ignorēja Antantes un Padomju Krievijas priekšlikumus. Par to, ka Polijā varētu nonākt daudz citu tautību iedzīvotāju, Polijas valdība neuztraucās. Pēc būtības Polija tikai centās iegūt laiku, lai labāk sagatavotos uzbrukumam KSFPR. 1920.~gada pavasarī Polijas tauta arvien vairāk kļuva par sava līdera J.~Pilsudska, viņu atbalstošo politiķu un ģenerāļu ambīciju ķīlnieci. J.~Pilsudskim izdevās, pēc R.~Dmovska vārdiem, ``ievilkt valsti militārā avantūrā''.

8.~aprīlī Padomju valdība paziņoja, ka sarunas izjukušas, pusēm nespējot vienoties par sarunu vietu, kam nebijis ``precedenta starptautisko attiecību vēsturē''. Padomju Krievijas pretinieks Nr.~1 šai laikā bija nu jau ģenerāļa P.~Vrangeļa vadītā ``balto'' armija Krievijas dienvidos. 1920.~gada 12.~aprīlī Sarkanās armijas virspavēlnieks S.~Kameņevs telegrammā Dienvid-Rietumu frontes pavēlniekam norādīja, ka Krimas ieņemšanas operācijai ir galvenā loma, tai jāveltī visi frontes spēki, pat nevairoties no pagaidu frontes ``poļu iecirkņa novājināšanas''. Vēl aprīļa vidū, iespējams, cerot uz miera sarunām, S.~Kameņevs lika visus brīvos spēkus pārsviest uz dienvidiem, pret P.~Vrangeļa armiju Krimā.

Šai laikā, atšķirībā no \strong{Polijas} attiecībām ar citām kaimiņvalstīm, \strong{ar Latviju} tai izveidojās \strong{sadarbība}, ko var izskaidrot kā ar īslaicīgu abu jaunveidojamo valstu interešu sakritību, tā arī Latvijas ģeogrāfisko atrašanos attālāk no Polijas un starp tai nedraudzīgajām valstīm. Mūsdienu Polijas vēsturnieks T.~Paļušinskis norādījis, ka, lai gan Polijai pašai pēc Versaļas miera līguma tika piešķirta pieeja pie jūras, tā 1919.~un 1920.~gada mijā perspektīvajos plānos Baltijas valstu, arī Latvijas, ostas redzēja kā kara bāzes savai jūras kara flotei. Taču atšķirībā no lietuviešiem poļi tai brīdī neizvirzīja teritoriālas pretenzijas uz Latgali. Tas ļāva Latvijas valdībai meklēt iespēju sadarbībai ar Poliju.

Latvijas un Polijas diplomāti kontaktējās jau Parīzes miera konferencē, attiecības turpinājās arī pēc tās. 1919.~gada septembrī par Polijas militāro pārstāvi Latvijā tika norīkots poļu kapteinis A.~Meškovskis, kurš 1919.~gada rudenī uzsāka Polijas militārā izlūkošanas tīkla izveidi Latvijā. Šai laikā norisa arī Lietuvas un Polijas armiju, kuras šeit sadarbojās cīņā pret boļševikiem, sadursmes ar līdz Latgalei sarukušās Padomju Latvijas teritorijā dislocētās KSFPR Rietumu frontes sastāvā esošās 15.~armijas daļām. Poļi un lietuvieši uzbruka Daugavpils virzienā. Tās aizsardzību vadīja latviešu komunists Ansis Daumanis, Sarkanajā armijā cīnījās arī latviešu daļas. Daugavpils tika pamatīgi apšaudīta ar artilēriju, pēc poļu gūstekņu liecības katram poļu karavīram par Daugavpils ieņemšanu tika solīta pudele šņabja un 500~cigarešu, taču poļi un lietuvieši ieņemt to nespēja. Poļu sapieri vien uzspridzināja abus Daugavpils tiltus.

Jaunu posmu Latvijas un Polijas attiecībās ievadīja ar vācu atbalstu radītās P.~Bermonta-Avalova vadītās Rietumkrievijas Brīvprātīgo armijas uzbrukums Rīgai. Vēl pirms tā sākšanās (1919.~gada 8.~oktobrī) Latvijas armijas virspavēlnieks ģenerālis Dāvids Sīmansons 6.~oktobrī pilnvaroja Latvijas militāro pārstāvi Polijā Mārtiņu Hartmani lūgt Polijas valdību atklāt kara darbību pret P.~Bermonta daļām. 1919.~gada beigās Polijas Ģenerālštābs sāka apsvērt iespēju par karadarbības izvēršanu pret P.~Bermonta-Avalova armiju, saskatot iespējas savas ietekmes palielināšanai Lietuvas un Latvijas teritorijās, kur koncentrējās bermontieši. Taču Lietuvas valdība apzinājās Polijas ekspansionistiskos nolūkus un atteicās izlaist Polijas daļas caur Lietuvas teritoriju ``kamēr netiks noregulētas politiskās lietas'' (Viļņas piederības jautājums, ko bija ieņēmusi Polija). Tā Polijas~--- Lietuvas konflikta dēļ iespējamā Polijas palīdzība Latvijai izpalika.

Latvijas attiecības ar Poliju vislabākās bija 1919.~gada beigās, 1920.~gada sākumā, kad poļu karaspēks iesaistījās Latvijas Brīvības cīņās, palīdzot tai atbrīvot Dienvidlatgali no Sarkanās armijas. 1919.~gada 30.~decembrī tika parakstīts Latvijas un Polijas armiju virspavēlniecību līgums par kopīgu uzbrukumu sarkanajai armijai Latgalē un tās atbrīvošanu. Kā atzīmējis latviešu vēsturnieks Ē.~Jēkabsons: ``Polijas militāri politiskās intereses bija diezgan tālejošas, jo tā~--- gan bez īsta pamata~--- cerēja piesaistīt Latviju kā sabiedroto ne tikai pret Padomju Krieviju, bet arī Lietuvu.'' 1920.~gada janvāra sākumā E.~Ridzs-Smiglija vadībā poļu karaspēks ieņēma Daugavpili, ko pēc tam nodeva Latvijas varas iestādēm. Par nopelniem Latgales kaujās ar Latvijas piešķirto Lāčplēša kara ordeni tika apbalvoti 63 Polijas armijas karavīri. Arī J.~Pilsudskis saņēma ordeni, pie tam~--- pirmās šķiras, kad cīņās piedalījušies ģenerāļi~--- tikai otrās šķiras. 1920.~gada janvāra beigās J.~Pilsudskis ieradās Daugavpilī, šeit notika viņa tikšanās ar Latvijas armijas virspavēlnieku Jāni~Balodi. Polijas valsts priekšnieks, būdams Daugavpilī, neslēpa Polijas interesi pievienot vairākus Latvijas apriņķus savai valstij.

Paralēli notikumiem Latgalē 1920.~gada 5.--21.~janvārī Helsinkos notika Baltijas valstu konference, kurā \strong{Polija centās savā vadībā apvienot Somiju, Igauniju, Latviju un Lietuvu}. Tika pat pieņemts militāras sadarbības projekts, kurš gan neguva tālāku attīstību.

Francija atteicās oficiāli piedalīties Helsinku konferencē, jo uzskatīja, ka konferences sasaukšana neatbilst sabiedroto kopējam uzskatam ``Krievijas jautājumā''. Tā baidījās, ka iespējamā atklātas konvencijas noslēgšana pret Padomju Krieviju varētu nozīmēt tās \latxti{de facto} atzīšanu. Tomēr, atturoties oficiāli piedalīties konferencē, Francijas diplomāti ar Polijas starpniecību centās panākt Francijai vadītājas pozīcijas plānojamā Baltijas savienībā un iespēju veidot to tālaika apstākļos par līdzekli iespējamai intervencei pret Padomju Krieviju. J.~Pilsudskis mēģināja realizēt ``Starpjūru valstu savienības'' (poļu \pltxti{Międzymorze}, latīņu \latxti{Intermarium}, krievu \rutxti{Междуморье}, vācu \detxti{Zwischenmeerland}) projektu, kurš paredzēja sākotnēji Polijas, Lietuvas, Ukrainas, Baltkrievijas, bet vēlāk arī Latvijas, Igaunijas, (iespējams~--- arī Somijas), Rumānijas, Ungārijas, Dienvidslāvijas un Čehoslovākijas konfederācijas izveidi, kura varētu starptautiskajā jomā nostāties kā pret Padomju Krieviju, tā pret Vāciju. Nosaukums bija radies no idejas apvienot valstis starp Baltijas un Melno jūru. Šādas konfederācijas izveidei pēc J.~Pilsudska domām bija jāturpina bijušās daudznacionālās Žečpospolitas tradīcijas, kas ļautu izbēgt no Vācijas vai Krievijas dominēšanas Centrālajā Eiropā.

Bez rezultātiem beidzās arī 1920.~gada martā Francijas rosinātā Polijas vadībā sasauktā konference Varšavā. Bez Polijas tajā piedalījās Somija, Latvija un Rumānija. Oficiāli tika paziņots, ka konference sanāk, lai izstrādātu vienotus noteikumus miera sarunās ar Padomju Krieviju. Taču dalībvalstis ļoti piesardzīgi vērtēja Polijas iniciatīvas. Latvijai izveidot tuvākus kontaktus ar Poliju un Franciju konkrētajos apstākļos nozīmēja zaudēt Lielbritānijas atbalstu. J.~Pilsudska plānu neatbalstīja arī lielvalstis, izņemot jau minēto Franciju. Tas sastapa pretestību arī pašā Polijā. Ukraiņu, baltkrievu un lietuviešu nacionālisti baidījās, ka apvienotajā valstī nepoļi, bet īpaši nekatoļi, izrādīsies otrās šķiras pilsoņu statusā. Arī pašu poļu vidū priekšroka tika dota tīri nacionālas Polijas valsts izveidei.

Pēc Latgales atbrīvošanas poļu sapieri gan atjaunoja Daugavpils dzelzceļa tiltu un uzcēla pontontiltu, taču jau pēc ``sarkano'' padzīšanas poļi civiliedzīvotājiem nodarīja jaunus materiālos zaudējumus (1920.~gada maijā bija ienācis 7~000~(!) vietējo iedzīvotāju sūdzību par poļu nodarījumiem). Daudzi poļu karavīri uzskatīja Latgali par iekarotu teritoriju, pieļāva disciplīnas pārkāpumus, nodarbojās ar laupīšanu, arī slepkavošanu. Pat radikālā poļu nacionāldemokrātu prese atzina, ka poļi Latgalē nav pienācīgi uzvedušies, zinot, ka šī ``poliskā'' zeme būs jāatdod Latvijai. Poļu virsnieki ar nepatiku uztvēra Latvijas iestāžu centienus nostiprināt savu varu Latgalē. Attiecības poļu karavīru un latviešu vidū pasliktinājās.

Pēc 1920.~gada 30.~janvārī notikušā Latvijas un KSFPR pamiera noslēgšanas un miera sarunu sākšanās Polija vilcināja sava karaspēka izvešanu no Latvijas.

Jau no 1920.~gada 8.~marta notika Polijas un Latvijas sarunas. Tās norisa divās komisijās: militārajā un politiskajā. Latvijas puse piekrita pārdomāt Polijas piedāvātās militāras konvencijas priekšlikumu. Pēc J.~Pilsudska norādījuma tā bija izstrādāta divos variantos, kur pirmais paredzēja armiju kopdarbību karā pret Padomju Krieviju, otrais~--- kopīgu uzbrukumu Lietuvai. Politiskajā komisijā galvenais bija jautājums par robežu, ko Polijas delegācija mēģināja izmantot kā kaulēšanās objektu, piedāvājot robežu izveidot pēc 1772.~gada parauga, kas nozīmēja Latgales iekļaušanu Polijā. Latvijas puse tam kategoriski pretojās. Polijas delegācija gan ātri piekāpās šai jautājumā, bet pieprasīja atstāt tās kontrolē pirmstilta nocietinājumus pie Daugavpils un izīrēt tai joslas Latvijas ostās, pirmkārt Liepājā. Šādas prasības no Polijas puses tikai pasliktināja tās attiecības ar Latviju. Sarunas beidzās ar fiasko.

Polijas valdība, sarūgtināta par Latvijas nevēlēšanos iesaistīties militāri politiskā savienībā, sāka izvirzīt dažādas prasības un atteicās atdot Latvijai Ilūkstes apriņķa sešus pagastus un Grīvas pilsētu. Šo teritoriju Latvijas spēki pārņēma tikai 1920.~gada jūlijā, kad sakarā ar Sarkanās armijas ofensīvu Polijā Daugavpilī un Ilūkstes apriņķī izvietotajām poļu vienībām draudēja atšķelšana no galvenajiem spēkiem un tās steigšus atkāpās. Tā pateicoties notikumu attīstībai poļu~--- padomju frontē, Sarkanās armijas uzbrukumam, seši Ilūkstes apriņķa pagasti, kuri līdz tam faktiski jau bija iekļauti Polijā, tika pievienoti Latvijai. Diplomātijas jomā teritoriālais strīds starp abām valstīm gan turpinājās līdz 1929.~gadam.

Latvijai neatkarības sākuma posmā Polijas atbalsts bija ļoti būtisks un izdevīgs, taču Latvijas valdības attieksme pret kaimiņzemi bija samērā piesardzīga. Tāpat Latvija centās ievērot formālu neitralitāti Polijas un Lietuvas strīdā par Viļņu. Latvijas pārstāvis Polijā Atis Ķeniņš 1920.~gada 3.~janvārī Ārlietu ministrijai ziņoja: ``Jākonstatē, ka šovinistiski imperiālistiskā tendence~--- nodibināt lielu un stipru Poliju, pārejot tās etnogrāfiskās robežas, piemīt, vairāk vai mazāk, visām Seima partijām un dominē sabiedrībā; tamdēļ arī [ministru] kabinets nevar būt citāds.'' 1920.~gada 22.~februārī A.~Ķeniņš ārlietu ministram Zigfrīdam Anna Meierovicam ziņoja par sarunu ar J.~Pilsudski, kurā tas mēģināja Latvijas pārstāvi pārliecināt par nepieciešamību darboties savienībā ar Poliju. A.~Ķeniņš rakstīja: ``Runājot par jaunām valstīm un viņu atzīšanu \latxti{de~iure}, Pilsudskis vairākkārt uzsvēra, ka Latvijas ģeogrāfiskais stāvoklis esot šajā ziņā daudz labvēlīgāks, nekā Igaunijas. Polijai pārceļot savu robežas līniju uz austrumiem, Latvijai radīšoties liela kopēja (?) robežu līnija ar Poliju. Latvija un Polija tad sastādīšot zināmu kompleksu, kuru Lielkrievija nevarēšot tik viegli apdraudēt kā Igauniju. Tamdēļ arī Poliju protežējošās Vakareiropas valstis jutīšoties aicinātas aizstāvēt reizē ar Poliju arī Latviju.'' A.~Ķeniņš, komentējot J.~Pilsudska vārdus, tālāk rakstīja: ``Šie Pilsudska pēdējie izteicieni palika mums mīklaini. Skaidrs ir tikai, ka Polija pie miera slēgšanas nekādā ziņā neatteiksies no okupētiem Baltkrievijas apgabaliem un ir ļoti varbūtīgi, ka dosies varbūt pat vēl tālāk uz austrumiem. Bet kur meklējama tā lielā kopējā robežu līnija starp Latviju un Poliju?'' Latvijas pārstāvis iztirzāja variantus, kādi varētu izveidoties, ja Latvija vai Polija okupētu Krievijas Vitebskas un Pleskavas apgabalus, bet nepieminēja vēl vienu iespēja, ja Polija pievirzītos Latvijai, pievienojot sev Lietuvu. Acīmredzot A.~Ķeniņš pat diplomātiskajā pastā nevēlējās runāt par Polijas draudiem Lietuvai. Taču visticamāk, ka ``liela kopēja robežu līnija ar Poliju'' Latvijai varēja izveidoties tieši tad, ja Polijai sev pievienotu Lietuvu. Tad varētu tikt tieši apdraudēta arī Latvija.

Kā norādīja A.~Ķeniņš, Nacionāldemokrātu partijas centrālais orgāns ``\pltxti{Gazeta Warszawska}'' (``Varšavas Avīze'') mēģināja ``aizstāvēt poļu kulturālās, vēsturiskās un pat etnogrāfiskās tiesības uz Latgali''. Polijas Ārlietu ministrija gan skaidroja, ka valdība atzīst Latvijas suverēnās tiesības uz Latgali, arī J.~Pilsudskis ``deva sajust, ka viņš nesimpatizē minētām N[acionālo] D[emokrātu] pretenzijām uz Latgali, tomēr noteiktas un gaišas atbildes uz šo jautājumu nedeva''. Iespējams, Latvijas tiesību uz Latgali atzīšanu no Polijas puses tās ``priekšnieks'' vēlējās izmantot kā ietekmes līdzekli uz Latviju, lai tā kopā ar Poliju turpinātu karu pret Padomju Krieviju.

1920.~gada 19.~martā, J.~Pilsudska vārda dienā notika ārēji spožs notikums viņa militārajā karjerā. Līdz tam viņš joprojām nesa brigadiera pakāpi, nēsāja pelēku strēlnieka frenci bez kādām zīmotnēm. 1919.~gada vasarā kara ministrs ģenerālis J.~Ļesņevskis armijas ģenerāļu vārdā griezās pie Seima ar lūgumu, lai tas ``pašu cienījamāko Polijas Republikas kareivi'' ieceļ par Polijas pirmo maršalu. Taču ``valsts priekšnieka'' un virspavēlnieka politiskie pretinieki nebūt nevēlējās paši savām rokām paaugstināt konkurentu. Seima vecāko konvents novērtēja priekšlikumu kā priekšlaicīgu. Tad nu armijnieki nolēma paši godināt savu virspavēlnieku. Kopējā atestācijas komisija, kura nodarbojās ar poļu virsnieku dažādu valstu armijās iegūto militāro dienesta pakāpju apstiprināšanu, griezās J.~Pilsudska vārda dienā pie viņa ar lūgumu pieņemt Pirmā Polijas maršala (\pltxti{Pierwszy Marszałek Polski}) nosaukumu. Nespēdams pretoties godkārības vilinājumam un vēloties uz Polijas armijas panākumu viļņa iegūt personīgu pagodinājumu, J.~Pilsudskis, neraugoties uz Seima pretestību, (tajā nacionāldemokrātu deputāts, kara kapelāns ksendzs K.~Lutoslavskis tūlīt pārmeta viņam pilnīgu militāru nekompetenci, nelikumīgu maršala nosaukuma un ordeņa ``\latxti{Virtuti Militari}'' piešķiršanu pašam sev), izdeva pavēli ``Polijas pirmā maršala pakāpi pieņemu un apstiprinu'', kas J.~Pilsudski uzreiz pacēla virs ģenerāļiem, pārnākušiem Polijas dienestā no Austroungārijas un Krievijas armijām. Maršala biogrāfijas autore I.~Pannenkova vēlāk atgādināja, ka Krievijas valdnieks Nikolajs~II, kurš, kļūstot par imperatoru, bija tikai pulkvedis, šo militāro pakāpi saglabāja arī tad, kad Pirmā pasaules kara laikā kļuva par Krievijas armijas virspavēlnieku un Lielbritānijas armijas feldmaršalu.

Kad 1920.~gada 14.~novembra svētdienā Polijā pirmo reizi atzīmēja Neatkarības dienu (tikai 1926.~gadā to saistīja ar 11.~novembri), par svinību centrālo notikumu kļuva pēc poļu mākslinieka M.~Kotarbiņska projekta no sudraba izgatavota maršala zižļa pasniegšana J.~Pilsudskim. Taču ceremonijā demonstratīvi nepiedalījās nedz Seima maršals (priekšsēdētājs), nedz valdības vadītājs. J.~Pilsudskis saņēma maršala zizli no ordeņa ``\latxti{Virtuti Militari}'' kavaliera kareivja J.~Venžika rokām. No šī brīža J.~Pilsudska līdzgaitnieki, griežoties pie viņa, sāka lietot dažādas uzrunas. Radinieki un vistuvākie draugi vēl no PPS agrīnā darbības perioda lietoja viņa bērnības iesauku ``\pltxti{Zjuk}'' vai pseidonīmu ``\pltxti{Dziadek''} (``vectēvs''). Tie, kas vēl pirms Pirmā pasaules kara viņa pakļautībā darbojās strēlnieku vienībās, joprojām sauca viņu par komandantu, tie, kuri bija viņa vadībā dienējuši leģionā~--- par brigadieri. Visi pārējie~--- par maršalu.

Pēc poļu un boļševiku miera sarunu neveiksmes kara darbība atjaunojās. Atbilstoši J.~Pilsudska ``Starpjūru valstu savienības'' projektam nozīmīga vieta plānotajā federācijā tika piešķirta \strong{Ukrainai}. Ar ukraiņu nacionālistiem J.~Pilsudskis uz laiku panāca izlīgumu.

Poļu historiogrāfija Polijas karagājienu uz Kijevu pūlas izskaidrot ar to, ka J.~Pilsudskis vēlējies nesavtīgi palīdzēt ukraiņu tautai iegūt neatkarību. Taču tas ir kārtējais ar Polijas vēsturi saistītais mīts. Ukraiņu tautā nebija lielas vēlēšanās aizstāvēt Ukrainas Tautas Republikas Direktoriju (valdību), tādēļ tās vadītāji bija spiesti vienoties ar Poliju par sadarbību. Ukrainas nacionālistu pārstāvji būtiski piekāpās Polijai teritoriālos un saimnieciskos jautājumos, faktiski atsakoties no ukraiņu nacionālajām interesēm, pretī saņemot politisku un militāru atbalstu cīņā pret boļševikiem. 1920.~gada 21.~aprīlī viens no ukraiņu nacionālistu līderiem, no 1918.~gada novembra Ukrainas Tautas Republikas (UTR) karaspēka virspavēlnieks, no 1919.~gada februāra līdz 1920.~gada novembrim UTR Direktorijas, par kuru toreiz tautā smējās, ka ``vagonā ir Direktorija, zem vagona~--- teritorija'' (\rutxti{в вагоне Директория, под вагоном территория)}, priekšsēdētājs S.~Petļura UTR vārdā vienojās ar Poliju. Šajā laikā S.~Petļuras vadītā Ukrainas Tautas Republika jau bija zaudējusi patstāvīga politiska un militāra spēka lomu, tās militārie formējumi bija atkāpušies poļu kontrolētajās teritorijās. Tikai ar Polijas palīdzību S.~Petļura varēja cerēt izveidot nacionālu ukraiņu valsti.

1920.~gada 21.~aprīlī Varšavā noslēgtais slepenais līgums, kurš stājās spēkā tūlīt pēc tā parakstīšanas bez ratifikācijas tautas pārstāvniecības iestādēs (acīmredzot tāpēc, lai neizsauktu sašutumu ukraiņu tautā), noteica abu valstu robežu pa Zbručas upi, ar to atzīstot Galīcijas iekļaušanu Polijas sastāvā. Tādejādi S.~Petļura ``labprātīgi'' atteicās Polijas labā no ievērojamas Rietumukrainas daļas. Pārējās Ukrainas robežas palika nenoteiktas. S.~Petļura arī apsolīja savā nākamajā valdībā iekļaut divus ministrus~--- poļus. Līgumam tika pievienotas arī militāra konvencija un tirdznieciski-ekonomiska vienošanās. Ukrainai kopš kopīgās karadarbības pret boļševikiem sākuma bija jāapgādā Polijas armija ar pārtiku un jāuztur tā savā teritorijā, tāpēc Polija apgādātu ar ieročiem ukraiņu armiju. Bija paredzētas poļu puses tiesības rekvizēt pārtiku. Uzvaras pār boļševikiem gadījumā Ukraina tiktu nostādīta gandrīz puskoloniālā atkarībā no Polijas. Ukrainai bija jāiznomā Polijai Krivoirogas dzelzs rūdas ieguves vietas (\uktxti{Криворізький залізорудний басейн}), dzelzceļa līnijas, piestātnes savās jūras ostās uz 99~gadiem. Vietējie iedzīvotāji uzstājās pret šādiem ``atbrīvošanas'' plāniem, nevēlējās nedz tādu ``atbrīvošanu'', nedz tādus ``atbrīvotājus''. Saasinājās S.~Petļuras un Rietumukrainas nacionālistu vadoņa J.~Petruševiča piekritēju attiecības, un S.~Petļura zaudēja šo nacionālo spēku atbalstu. Šos faktus mūsdienās sāk atzīt arī atsevišķi poļu vēsturnieki.

Var piebilst, ka Polijas uzbrukumu cerēja izmantot arī baltkrievu nacionālisti. Viņi nosūtīja uz Maskavu delegāciju ar Baltkrievijas Tautas republikas valdības vadītāju V.~Lastovski priekšgalā, lai noslēgtu mieru ar Padomju Krieviju. Baltkrievijas Tautas republikas teritorijā, pēc delegācijas domām, bija jāietilpst arī Viļņai. Tiesa, plāni izmantot Padomju Krievijas grūto stāvokli nepiepildījās, boļševiki ar baltkrievu nacionālistiem varēja nerēķināties to vājuma dēļ.

Poļu puses pārstāvji 1920.~gadā atsaucās uz Padomju Krievijas, boļševiku partijas un Komunistiskās internacionāles vadītāju paziņojumiem par pasaules sociālistiskās revolūcijas neizbēgamību, kuras ceļš no Krievijas uz Rietumeiropu ietu caur Poliju. Arī vēl mūsdienās daudzi poļu vēsturnieki lieto šo argumentāciju. Piemēram, Jagelloņu universitātes profesors Krakovā A.~Novaks kādā rakstā apgalvo, ka tas nebijis vienkāršs poļu-krievu karš, Sarkanā armija karojusi nevis par Krievijas, bet ``komunisma uzvaru'', gribējusi iet ne tikai pāri ``baltās Polijas līķim'', bet arī ``baltās'' (nekomunistiskās) Krievijas un ``vecās Eiropas līķim''. Tiesa, savā rakstā šis autors pat nepiemin, ka tieši Polija iebruka Ukrainā, ka Sarkanajai armijai, kura vēl nebija noslēgusi Pilsoņu karu, 1920.~gada pirmajā pusē nācās aizstāvēties pret ``balto'' Poliju.

Notikumu būtība daudzu poļu vēsturnieku darbos tiek klāstīta sekojoši. Ukrainas valdība~--- Direktorija Sarkanās armijas draudu priekšā griezās pēc palīdzības pie Polijas, solot par to atteikties no suverenitātes pār Rietumukrainu. Boļševiki, vēloties realizēt savus pasaules revolūcijas plānus, gatavojās uzbrukt Eiropai, lai tur ievestu komunismu, pirmā valsts viņu armijas ceļā bija Polija. Tāpēc tā atradās nāves briesmās, varēja nonākt kārtējā Krievijas jūgā. Lai ``aizsteigtos priekšā'' boļševiku triecienam, poļu armija pirmā deva triecienu austrumu virzienā. Sākumā poļiem veicās, tie pat ieņēma Kijevu, taču drīz sākās pretuzbrukums un viņiem nācās atkāpties.

Cik lielā mērā šī versija atbilst patiesībai?

Pēc autora domam, boļševiku vadoņos samērā ilgi valdīja cerības izmantot pēc Pirmā pasaules kara Rietumos pastāvošo darbaļaužu trūkumu un neapmierinātību ar savu stāvokli revolūcijas uzkurināšanai Eiropā. Boļševiki Oktobra apvērsuma uzvaru Krievijā sākotnēji uzskatīja par pasaules revolūcijas sākumu.

Pasaules revolūcijas koncepcijas autors bija V.~Ļeņins. Pārējie viņa līdzgaitnieki labākajā gadījumā bija tikai talantīgi tās propagandētāji. Taču jau 1918.~gada sākumā V.~Ļeņins saprata, ka šī koncepcija īstenosies ne tik drīz un ne tik viegli kā tas šķita agrāk. Taču visa Padomju Krievijas pirmo gadu ārpolitika balstījās uz šo koncepciju. Oficiāli atteikties no tās nozīmēja atzīt Oktobra bruņotās sacelšanās un Padomju varas nodibināšanas priekšlaicīgumu. Tomēr bija jāmeklē atkāpšanās ceļš pārāk optimistiskajiem revolūcijas perspektīvu vērtējumiem. Tāpēc 1919.--1920.~gadā V.~Ļeņins sāka pamatot pasaules revolūcijas viļņa ``atplūdus'', Padomju Krievijas un citu padomju republiku bruņotas sadursmes ar visu kapitālistisko pasauli vietā sludināt ekonomisku un politisku sāncensību. Tas no Padomju Krievijas viedokļa bija vienīgais pareizais ceļš, bet citi boļševiku vadītāji to kaut cik jūtami neatbalstīja. V.~Ļeņina viedokli neviens atklāti neapstrīdēja, bet padomju valsts iekšējās un ārējās politikas aparāts turpināja darboties saskaņā ar 1917.~gada oktobra nostādnēm. No otras puses, padomju republiku darbinieki, neraugoties uz dažu savu vadītāju kareivīgajiem noskaņojumiem, sākotnēji ilgi ticēja Polijas ``sovjetizācijai'' no iekšienes. Vēl 1919.~gada martā V.~Ļeņins, uzstājoties KK(b)P VIII~kongresā, runāja par to, ka Polijā ``pēc mūsu revolucionārā kalendāra vairs nav tālu Oktobris''. Pie tam Sarkanās armijas spēki Pilsoņu kara brāzmainajos notikumos bija izsmelti jau līdz 1920.~gadam, tai vairs nebija kaujas spara. 1920.~gadā padomju vadītājiem vairs nebija lielu cerību uz drīzu pasaules revolūciju, veidojās vismaz daļēja sapratne, ka tā jāatliek uz nenoteiktu laiku.

Laikā, kamēr Krievijā vēl nebija sakauti ``baltie'', Eiropā valdīja pagurums pēc Pirmā pasaules kara un reizē iedzīvotāju nevēlēšanās nonākt vēl briesmīgākās ekonomiskās grūtībās, pasaules revolūcijas plāni bija nereāli. Augstākais ko spēja krievu boļševiki~--- palīdzēt nostiprināt padomju varu vairākās bijušās Krievijas impērijas nacionālajās nomalēs. Pie tam ne tikai Sarkanā armija no Krievijas nāca palīgā vietējiem boļševikiem Ukrainā, Baltkrievijā u.c., bet arī Polija jau 1919.~gadā iebruka kaimiņtautu apdzīvotajās teritorijās, cerot izmantot juku laikus savu lielvalstisko ambīciju apmierināšanai.

Paziņojumi par drīz gaidāmu pasaules revolūciju padomju propagandā gan izskanēja arī 1920.~gadā un vēlāk. Tomēr jāņem vērā, ka starp šiem jūsmīgi ``rrrevolucionārajiem'' paziņojumiem un KSFPR tai laikā realizēto ārpolitiku bija ievērojama atšķirība. 1920.~gada 2.~februārī Padomju Krievija jau bija noslēgusi miera līgumu ar Igauniju, pat piekāpjoties strīdīgo teritoriju, kurās dzīvoja jaukts~--- krievu un igauņu iedzīvotāju sastāvs, jautājumā. Tika vestas sarunas par tādu līgumu noslēgšanu ar Somiju, Lietuvu un Latviju, tika izvirzīti priekšlikumi Rumānijai. Polija nebija izņēmums. Pilsoņu karam Krievijā turpinoties, Padomju Krievijas vadība, neraugoties uz visiem pasaules revolūcijas lozungiem, bija ieinteresēta nekavējoties noslēgt ar to mieru, arī piekāpjoties teritoriāli. Taču Polijas vadītāji ar J.~Pilsudski priekšgalā loloja citus plānus.

Jau 17.~aprīlī J.~Pilsudskis deva pavēli Polijas armijai 25.~aprīlī sākt uzbrukumu. Pirms tam, jau 19.~aprīlī poļu aviācija nometa uz Kijevu bumbas. Uzlidojumā gāja bojā 10 un ievainoti tika 14~cilvēki. 1920.~gada 25.~aprīlī tieši Polijas armija ar S.~Petļuras vadīto vienību atbalstu uzsāka uzbrukumu Kijevai. Sākās Polijas ``\pltxti{оfensywa na Wschod}'' (uzbrukums austrumos). Kā norādījis ukraiņu vēsturnieks S.~Lozuņko, faktiski Polija bija pirmā no Tautu Savienības (Oficiāli dibināta 1920.~gada 10.~janvārī) dalībniecēm, kura uzsāka agresīvu karu pret kaimiņvalsti.

Boļševiki 1920.~gada pavasarī gatavojās aktīvi karot pret Poliju, kura nepārtraukti centās paplašināt savu teritoriju, taču, tā kā notikumi ievirzījās citādi, neviens nevar apgalvot, ka padomju uzbrukums Polijai tiešām sāktos. Tikai Polijas uzbrukums izsauca jau atlikto pasaules revolūcijas plānu izvilkšanu no ``tālās atvilktnes''. Krievu vēsturniece I.~Mihutina uz arhīvu materiālu pamata secinājusi, ka pasaules revolūcijas ideja poļu-padomju kara pirmajā posmā nebūt neprevalēja padomju politikā un propagandā, pirmajā plānā tā izvirzījās tikai pēc pirmajiem Sarkanās armijas panākumiem. Tieši J.~Pilsudska gājiens uz Kijevu izprovocēja Sarkanās armijas atbildes karagājienu uz Varšavu, pasaules revolūcijas idejas atdzimšanu boļševikos un Polijas valstiskuma dziļu krīzi.

Literatūrā un publicistikā ir izteikti vairāki pieņēmumi par Polijas vadības, J.~Pilsudska rīcības motīviem un izvēlēto rīcības brīdi.

Pirmkārt, jāņem vērā, ka Polijas un Padomju Krievijas attiecības nevarēja neietekmēt 1917.--1919.~gada revolūcijas Austrumeiropā un Pilsoņu kara notikumi Krievijā. Tajos nacionālās cīņas cieši savijās ar sociāli-politiskajām. Tas tikai uzkurināja abu valstu valdošo slāņu savstarpējo nepatiku. Taču pagaidām (līdz Sarkanās armijas ielaušanās brīdim Polijā) galvenais strīdus jautājums, kurš slēpās aiz idejiski politiskiem lozungiem, bija Padomju Krievijas un Polijas teritoriālā norobežošanās, kuru katra puse gribēja realizēt pēc saviem ieskatiem.

Otrkārt, 1920.~gada pavasarī jau kļuva skaidrs, ka Krievijā boļševiki uzvarēs t.s. kontrrevolucionārus. J.~Pilsudskis par savu galveno ienaidnieku tagad uzskatīja jau nevis ``baltos'', bet ``sarkanos'', nolēma izvērst cīņu pret tiem.

Treškārt, Antantes valstis, kuras uzstājās kā poļu etnogrāfiskās un valsts robežas sakrišanas piekritējas, šai laikā vairāk uzmanības pievērsa robežu noregulēšanai Eiropas centrā, nevis tās austrumos.

Tāpēc Polija šo momentu vēlējās izmantot, lai ieņemtu pēc iespējas lielākas teritorijas, lai ne tikai nostiprinātu Polijas valstiskumu, radītu Ukrainas nacionālistu marionešu valsti S.~Petļuras vadībā, atkarīgu no Polijas, bet reizē arī nodarītu pēc iespējas lielāku kaitējumu savam konkurentam austrumos~--- Krievijai, kas savukārt, pēc godkārīgo poļu politiķu viedokļa, ļautu Polijai iegūt lielvalsts statusu. Pie tam starptautisko statusu Polijai bija it kā ``jāmanto'' no sagrautās Krievijas impērijas. Pēdējais pieņēmums gan nav pienācīgi apstiprināts ar pētījumiem, bet loģikas tam netrūkst. Taču J.~Pilsudskis nepietiekami novērtēja boļševiku spēkus.

Jāsaka, ka Polijas politiķu vidū bija arī cilvēki, kuri protestēja pret tās valdības un J.~Pilsudska avantūristisko politiku. Pret J.~Pilsudska plāniem iebilda ne tikai poļu komunisti. Zemniecību pārstāvošie ļudovieši un arī sociālisti vēlējās panākt miera sarunu sākumu. Nacionāldemokrāti J.~Pilsudska plānus uzskatīja ne tikai par nereāliem, bet arī kaitīgiem poļu nacionālajai valstij. Īpaši asas domstarpības bija par vienošanos ar S.~Petļuru. Endeku prese un deputāti Seimā par viņu runāja nicinoši kā par klaidoni, kuru var pieciest tikai aiz žēlsirdības. Endeki labāk vēlētos noslēgt savienību ar ģenerāli P.~Vrangeli. Kad 1920.~gada aprīļa beigās pēc Ukrainas Tautas Republikas un Polijas līguma parakstīšanas J.~Pilsudskis deva savām armijām pavēli maršēt uz Kijevu, viens no nacionāldemokrātu līderiem S.~Grabskis, uzskatot, ka valsts pārdzīvo ``brutālu demagoģijas laikmetu'', atstāja Seima Ārlietu komisijas priekšsēdētāja vietu, lai protestētu ar visiem līdzekļiem pret ``Ukrainas avantūru''. Arī citi nacionāldemokrātu politiķi kritiski vērtēja J.~Pilsudska ārpolitiku. Kad 1920.~gada jūnijā Polijas armija bija spiesta uzsākt strauju atkāpšanos, nacionāldemokrāti R.~Dmovska vadībā Poznaņā jau veica pasākumus, lai konsolidētu poļu spēkus pēc iespējamās Varšavas krišanas.

Kā jau teikts, 1920.~gada 25.~aprīlī, izmantojot to, ka galvenie ``sarkano'' spēki rietumos bija koncentrēti Baltkrievijā, Polija ar S.~Petļuras vadīto ukraiņu nacionālistu atbalstu uzsāka uzbrukumu boļševiku spēkiem Ukrainā visā frontes līnijā no Černobiļas līdz Rumānijas robežai. Pēc padomju militārā vēsturnieka N.~Kakurina datiem 1920.~gada pavasarī Polijas armijas sastāvā bija 738 ~tūkstošu cilvēku. Pēc mūsdienu krievu vēsturnieču I.~Jažborovskas un V.~Parsadanovas vērtējuma poļu-padomju spēku attiecība Ukrainā bija 3:1, bet Baltkrievijā 0,8:1. (Padomju vēsturnieki poļu pārsvaru vērtēja krietni augstāk.) Poļu armijas uzbrukums norisa pēc Francijas militārās misijas Varšavā vadītāja ģenerāļa P.~Anri izstrādātā plāna. Tas paredzēja aplenkt un iznīcināt ``sarkano'' spēkus Ukrainā, ieņemt Kijevu un Odesu, pēc tam pārgrupēt spēkus un ieņemt visu Baltkrieviju. Sākotnēji notikumi risinājās poļu labā. Viņi saņēma gūstā vairāk nekā 25 tūkstošus sarkanarmiešu, ieguva lielas trofejas (to rokās nonāca 2 bruņu vilcieni, 120 lielgabalu un 418 ložmetēju.) Viņu pusē 27.~aprīlī pārgāja no Rietumukrainas Tautas republikas armijas bijušajiem karavīriem saformētās trīs Sarkanās armijas brigādes (ap 12~000 karavīru). 7.~maijā poļi ieņēma Kijevu (tā bija jau piecpadsmitā varas maiņa pilsētā pēdējo triju gadu laikā), izveidoja placdarmu arī Dņepras kreisajā krastā. Starp citu, poļu izlūki iebrauca pilsētas centrā ar tramvaju, jo Sarkanā armija, lai pilsētu neizpostītu, bija to jau atstājusi.

Poļiem panākumu gūšanu lielā mērā atviegloja tas, ka Ukrainas teritorijā darbojās dažādas politiskās nokrāsas partizānu vienības un vienkārši puskriminālas bandas, kuras tikai izmantoja politiskus lozungus, bet būtībā nodarbojās ar laupīšanu.

J.~Pilsudskis jau 26.~aprīlī Žitomiras (ukraiņu \uktxti{Житомир}) pilsētā izdeva uzsaukumu, ar kuru griezās pie ukraiņu tautas, skaidrojot, ka poļi padzenot no ukraiņu apdzīvotās teritorijas iebrucējus, pret kuriem šī tauta saceļoties ar ieročiem rokās, solīja tai neatkarību un patstāvību valsts izbūvē. Poļu karaspēkam Ukrainā bija jāpaliek līdz laikam, ``kamēr ukraiņu tautas armija spēs aizstāvēt savu zemi pret jauniem iebrukumiem''. (Poļu vēsturnieks A.~Novaks apgalvo, ka J.~Pilsudskis'' ``uzbruka Kijevai, lai Petļura, kurš kopā ar poļiem soļoja kā palīgformējums, varētu tur nodibināt Ukrainas valdību un nostiprināties. Plāns bija tāds, ka apmēram trīs mēnešus Polijas armija aizstāvēs Petļuru''.) Faktiski tā bija ukraiņu teritoriju okupācija.

S.~Petļura savukārt savā aicinājumā tautai uzsvēra poļu un ukraiņu savienības nesatricināmību. Taču cerētā ukraiņu sacelšanās pret padomju varu nenotika. Gluži pretēji, 8.~maijā poļu militārās administrācijas Ukrainā vadītājs E.~Ridzs--Smiglijs bija spiests izdot pavēli par lauku kara tiesu ievešanu, kurām bija tiesības par ``dumpi'' u.c. pārkāpumiem piespriest nāves sodu. Krievu vēsturnieks M.~Meļķuhovs aprakstījis, kā par atteikšanos dot okupantiem produktus virknes ukraiņu ciemu iedzīvotāji tika apšauti ar ložmetējiem.

Ukraiņu nacionālistu cerības ar poļu rokām radīt savu valsti bija gauži nereālas. Jau kopējā uzbrukuma laikā S.~Petļura sūdzējās J.~Pilsudskim, ka tiek ignorētas ``Ukrainas valsts elementārās neatkarības, valodas, valdības institūtu tiesības, netiek ievēroti tās juridiskie akti, notiek nepamatota šīs valsts militāro un civilo aizstāvju tiesību atņemšana, bezkaunīga valsts īpašuma sagrābšana un piesavināšanās miljardu rubļu vērtībā.'' Vēlāk, jau 1921.~gadā S.~Petļura žēlojās, ka Polijas palīdzība Ukrainai bija daudz mazāka nekā tās vērtības, ko poļu karaspēks izveda no Ukrainas.

Mūsdienu ukraiņu publicists O.~Pidluckis stāsta, kā notika poļu muižnieku īpašuma tiesību atjaunošana Ukrainā. Tā viens no lietuviešu izcelsmes kņazu Sanguško (\pltxti{Sanguszkowie}) dinastijas pārstāvjiem, atgriezies savā muižā Podolijā (ukr. \uktxti{Поділля}, poļu \pltxti{Podole}), ``piespieda visus ciema iedzīvotājus astoņas stundas viņa priekšā stāvēt uz ceļiem, bet katru, kurš bez pana atļaujas piecēlās, ``sabiedroto'' karavīri bez žēlastības nogalināja.'' Acīmredzamā S.~Petļuras pakļaušanās Polijas diktātam, poļu karavīru, kuriem šis karš bija tikai Pirmajā pasaules karā norisušās cīņas pret carisko Krieviju turpinājums, bezceremoniālā uzvedība pret vietējiem iedzīvotājiem, marodierisms, pārtikas un furāžas konfiskācijas izsauca ukraiņu neapmierinātību, kura virknē gadījumu pārauga partizānu karā. Daudziem pat pārliecinātiem Ukrainas patstāvības piekritējiem (\uktxti{самостiйнiк}) tagad likās, ka šādos apstākļos ukraiņiem mazākais ļaunums būs boļševistiskā Krievija, nevis ``panu'' Polija. S.~Petļuras vadītā ukraiņu militārā un civilā administrācija tā ir nepārbrauca no Viņņicas (\uktxti{Вінниця}) uz Kijevu. Divus mēnešus vēlāk, Polijai ciešot sakāves, beidza pastāvēt arī Ukrainas Tautas republika.

Tā kā Polija bija uzbrūkošā puse, tas radīja morālu pārsvaru Padomju Krievijai. (Šai laikā formāli neatkarīgās padomju republikas, tai skaitā arī Padomju Ukraina, bija noslēgušas savstarpēju militārās savienības līgumu. Reāli gan galvenie lēmumi tika pieņemti Maskavā.)

Padomju vadība karā pret Poliju lielu lomu piešķīra ideoloģiskajam propagandas darbam pretinieka vidū. 1920.~gada 24.~martā un turpmāk vēl pirms kara sākšanās propagandai pretinieka vidū Rietumu frontei tika iedalīti 5~miljoni rubļu, 32~tonnas papīra un 1,6~tonnas poļu šrifta, tika atlasīti un speciālu sagatavošanos izgāja 83~poļi, baltkrievi un lietuvieši, kuru uzdevums bija izplatīt aģitācijas un propagandas rakstura literatūru ienaidnieka vidū. Jau 1920.~gada aprīlī Rietumu frontē bija izdoti 32~nosaukumu drukātu materiālu ar pusotra miljona tirāžu. Arī Lietuvas un Baltkrievijas Komunistiskās partijas CK aprīlī-maijā izplatīšanai Rietumu frontē izdeva 103~nosaukumu materiālus vairāk nekā miljonu eksemplāru tirāžā poļu, lietuviešu, baltkrievu un jidiša valodās. VCIK un TKP uzsaukumā ``Visiem Krievijas strādniekiem, zemniekiem un godīgiem pilsoņiem'' kā jau par izlemtu lietu tika paziņots: ``pēc Pilsudska bandu sakāves mums Polijas neatkarība paliks neaizskarama. Ar poļu proletariātu un poļu zemniecību, kuri kļūs par pilnīgiem savas zemes saimniekiem, mēs bez grūtībām nodibināsim brālīgu savienību. Tikai mūsu kopīgais šļahtiski-buržuāziskais ienaidnieks mūs šķir. Tad uz cīņu pret šo ienaidnieku.'' Rietumu frontes karavīriem tika solīts: ``Par jums būs poļu strādnieki, kuriem jūs palīdzēsiet atbrīvoties no buržuāzijas jūga!''. Jāuzsver, ka ar šādu propagandu uz ilgu laiku gandrīz pilnīgi tika aizmirstas agrāk bieži pieminētās poļu kā katoļu un krievu kā pareizticīgo pretrunas. Reliģiskais aspekts karā, ar to saistītajā ideoloģiskajā pretstāvē nebija dominējošs, to atbīdīja padomju šķiriskuma un poļu nacionālisma uzsvērums.

Poļu karaspēks nespēja izpildīt savu uzdevumu aplenkt un iznīcināt ``sarkano'' armijas. Tās spēja atkāpties un pārgrupēties. Padomju Ukrainas vadošās iestādes, vadoties no tā, ka ``sākusies izšķiroša cīņa starp strādnieku un zemnieku republikām Krieviju un Ukrainu un poļu šļahtu'', aicināja triecienam atbildēt ar triecienu.

Rietumos auga skepse pret ``Kijevas dēku'' un noraidoša attieksme pret Poliju. Darba cilvēki, kuri tagad pārliecinājās, ka uzbrucēji ir nevis boļševiki, bet Rietumu atbalstītā Polija, izvērsa kustību ar lozungu ``Rokas nost no Krievijas''.

Arī pašā Polijā karš austrumos nebija populārs. Daudzi poļi nevēlējās tikt nogalināti nevis karā par savas dzimtenes neatkarību, bet gan svešu tautu teritoriju iekarošanas gaitā. Tāpēc gribētāju karot Polijas civiliedzīvotāju vidū bija gaužām maz. Kara ministra vietnieks ģenerālis K.~Sosnkovskis kara apgabalu pavēlniekiem rakstīja, ka 1900.--1901.~gadā dzimušo iesaukšana rādījusi, ka daudzi no iesaukšanas izvairījās. 1920.~gada 30.~aprīlī K.~Sosnkovskis izdeva slepenu pavēli, kurā citastarp bija teikts: ``Kara ministrija 12.~un 15.~maijā nozīmē visā zemē no karadienesta izvairījušos personu, it sevišķi 1900. un 1901.~gadā dzimušo iesaucamo, vispārēju tvarstīšanu.'' Jūlijā Lietuvas un Baltkrievijas frontes pavēlnieks ģenerālis S.~Šeptickis ziņoja Ģenerālštābam, ka ``ievērojama daļa to karavīru, kuri atgriežas no atvaļinājumiem un no jaunā papildinājuma ir aģitatori, kuri iestājas pret tālāku karu tā bezmērķības dēļ un tāpēc, ka ``Pilsudskis ir pārdevies kungiem un sviež karavīrus lielgabalu ugunīs''.

Arī daļa poļu inteliģences nosodīja karu pret padomju republikām. 14.~maijā virkne poļu rakstnieku publicēja protestu pret poļu militāristu iekarošanas plāniem. Tiesa, viņu motīvi bija dažādi. Piemēram, S.~Žeromskis 1920.~gadā rakstīja: ``Ir smieklīgi, ka neiekarojuši vēl rietumu un ziemeļu robežas \citespace{} mēs ar mūsu bruņinieku varonīgajām asinīm slakām priekštilta nocietinājumus pie Dņepras un iekarojam Sožas [Dņepras pieteka, tek pa Krievijas, Baltkrievijas teritoriju un daļēji pa Ukrainas robežu~--- V.Š.] upi! Poliju nicina visa rietumu strādnieku pasaule. Mēs iesaistījāmies karā ar Maskavu, atmetot rietumus un jūru\dots{}'' 1.~jūlijā Seimā zemnieku deputāts ksendzs E.~Okons izvirzīja pārmetumus Polijas valdībai, Seimam un militārajai sistēmai: ``Pārāk daudz ir izdarīts kļūdu, pārāk lielu, lai neiesēstos peļķē, pie tam bez mazākās vajadzības. Šīs kļūdas ir notikušas mūsu administrācijas neprasmes dēļ, kura bieži dara pāri lauku iedzīvotājiem, utopiskās austrumu politikas ar postošo Ukrainas mirāžu dēļ; izdarītas stratēģiskas kļūdas nevajadzīga un kaitīga karagājiena uz Kijevu dēļ; izdarīja kļūdas gan ģeneralitāte, gan militāristi, gan mūsu intendantūra, kura barojas un aptaukojas uz tēvzemes rēķina, kad tā ir briesmās; un pieļauta kļūda~--- palaista garām iespēja noslēgt mieru, kad Krievija sniedza mums roku un ar Čičerina muti teica: ``Poļi, pietiek karot! Mēs gribam mieru!'' Kaut PPS vairākums veicināja partijas biedru un tās ietekmē esošo strādnieku un zemnieku iestāšanos armijā, opozīcija sociālistu rindās nosodīja uzbrukumu padomju republikām, aicināja boikotēt mobilizāciju. 1920.~gada jūlijā notika Varšavas pilsētas domes ārkārtas sēde. Tajā ebreju sociālistu (Bunda) pārstāvis partijas vārdā deklarēja, ka valstij nepieciešami nevis teritoriāli ieguvumu, bet sociāli ekonomiskas reformas, iespējamas tikai apstākļos, kad valda miers ar kaimiņiem. Uz to no varas iestāžu puses pret Bundu tika izvērstas represijas, tika slēgti partijas laikraksti, provincē notika aresti.

Vēl negatīvāk Polijas karu vērtēja tās ieņemto ukraiņu, baltkrievu un lietuviešu apgabalu iedzīvotāji. Tā, 1920.~gada 20.~jūnijā Polijas 1.~armijas, kura darbojās tagadējās Baltkrievijas un Lietuvas teritorijā, izlūkošanas nodaļa ziņoja Ģenerālštābam, ka vietējie iedzīvotāji labprāt slēpj poļu dezertierus. ``Ja dezertieris saka, ka bēdzis no padomju gūsta, iedzīvotāji atsakās viņam sniegt jebkādu palīdzību. Ja viņš saka, ka bēg no frontes, viņš saņem palīdzību natūrā un norādījumus tālākajam ceļam.'' Jūlija vidū dažos Ļubļinas kara apgabala apriņķos iesaukšanas punktos ieradās ne vairāk kā 30\% iesaucamo, no kuriem 90\% prasīja iesaukšanas atlikšanu. No ukraiņu ciemiem, izņemot Helmas (krievu Holmas) novadu, iedzīvotāji vispār neatsaucās mobilizācijai. Dezertieri izturējās arvien drošāk, bija gadījumi, kad notika to apšaudes ar policistiem un žandarmiem. 24.~jūlijā Polijas kara ministrija izdeva pavēli par lauku kara tiesu ieviešanu cīņai pret dezertēšanu. To kompetence tika paplašināta arī uz civilpersonām, kuras varēja tiesāt par kūdīšanu nepakļauties iesaukumam un citiem noziegumiem. Tika pavēlēts šo tiesu spriedumus plaši izplatīt iedzīvotāju vidū.

Turpretī Krievijā Polijas iebrukums Ukrainā un Baltkrievijā kļuva par patriotiskas kustības rašanās cēloni. Pats Kijevas~--- ``visu krievu pilsētu mātes'' (\rutxti{Киев~--- мать городов русских}) sagrābšanas fakts daudzos Krievijas pilsoņos, neatkarīgi no viņu politiskās pārliecības, pamodināja dziļas patriotisma jūtas. Zināmā mērā var pat salīdzināt situāciju ar 1812.~gadu, kad iekšpolitiskās problēmas atvirzījās otrajā un trešajā plānā nepieciešamības dot pretsparu ārējam ienaidniekam priekšā, atbrīvot Krieviju no iebrucējiem. Slavenais krievu ģenerālis A.~Brusilovs 1.~maijā vēstulē vienam no Sarkanās armijas darbiniekiem (tā nonāca arī V.~Ļeņina rokās) rakstīja: ``Brīvā Krievija \citespace{}, atbrīvojusi poļus un devusi viņiem iespēju pašnoteikties un nokārtot dzīvi pēc saviem ieskatiem, ir tiesīga prasīt to pašu arī no viņiem, un poļu iebrukums zemēs, kuras izsenis piederējušas pareizticīgajai tautai, ir nepieciešams atvairīt ar spēku.'' Ģenerālis griezās pie padomju varas iestādēm ar aicinājumu spert soļus ``tautas patriotisma veicināšanai, bez kura nebūs stipras kaujas spējīgas armijas.'' Atbildot uz to, KSFPR Revolucionārā kara padome pie virspavēlnieka radīja ``Īpašo apspriedi spēku kāpināšanas jautājumos cīņai pret poļu kontrrevolūcijas uzbrukumu''. Tās priekšgalā tika nozīmēts A.~Brusilovs. 30.~maijā ``Īpašā apspriede'' publicēja ``Aicinājumu visiem bijušajiem virsniekiem, lai kur arī viņi atrastos''. Tajā visi bijušie Krievijas virsnieki dzimtenes mīlestības vārdā tika aicināti aizmirst pārestības un doties dienēt Sarkanajā armijā lai ``par katru cenu aizstāvētu mūsu dārgo Krieviju''. Bija norādīts, ka ``īpaši tas ir nepieciešams tagadējā bargajā laikā, kad brālīgā un mums dārgā poļu tauta, kura pati izbaudījusi smagu ārvalstu jūgu, tagad pēkšņi grib atņemt mums zemes ar senseniem krievu iedzīvotājiem un atkal pakļaut tos poļu apspiedējiem''. Tika publicēti arī nesen ne bez netiešas Polijas palīdzības sagrautās A.~Deņikina armijas bijušo virsnieku paziņojumi par gatavību doties cīņā pret poļiem par padomju Krieviju. Arī Ļ.~Trockis šim aicinājumam piebiedroja savu balsi: ``Visi, kas ir godīgi inteliģencē! Krievu virsniecība~--- tā, kura sapratusi, ka Sarkanā armija izglābs krievu tautas brīvību un neatkarību! Jūs visus aicina Rietumu fronte!'' Aicinājuma ietekmēti daudzi virsnieki sāka pieteikties Sarkanajā armijā. Tā poļu cīņa par savām nacionālajām interesēm vismaz daļā Krievijas sabiedrības radīja rūpes par savas valsts nacionālajām interesēm. Ja savstarpējās ``balto'' un ``sarkano'' cīņās tās tika aizmirstas, ārvalsts~--- Polijas uzbrukums tās pamodināja. Viens no augstākajiem poļu militāristiem K.~Sosnkovskis atzina, ka ``Brusilova aicinājums savu ir paveicis, vēlamo iespaidu sasniedzis, visi bijušās krievu armijas virsnieki ieplūda boļševiku armijas rindās, un karš ieguva kara par Krieviju raksturu.''

Protams, boļševiki krievu patriotismu izmantoja nepieciešamības spiesti, taktiskos nolūkos, nekad neaizmirstot savus stratēģiskos šķiriskos mērķus, taču līdz poļu padzīšanai no citu tautu teritorijām internacionālistu-boļševiku un krievu patriotu uzdevumi sakrita. Tiesa, pēc Sarkanās armijas ieiešanas Polijas teritorijā boļševiki jau viennozīmīgi krievu patriotismu izmantoja ļaunprātīgi.

``Kērzona līnijas'' pāriešana no Polijas puses sarežģīja Polijas attiecības ar ģenerāļa P.~Vrangeļa vadītajiem krievu ``baltajiem''. P.~Vrangelis un viņa pārstāvji izteica gatavību kopā ar poļiem karot pret boļševikiem, taču nevēlējās tiem atdot zemes uz austrumiem no ``Kērzona līnijas''. Ģenerāli atbalstīja Lielbritānija un ASV. 1920.~gada 1.~augustā Lielbritānijas valdība uzsvēra, ka tā neatbalsta Polijas tieksmi paplašināt savas robežas uz austrumiem no ``Kērzona līnijas''. 21.~augustā ASV Valsts departaments aicināja Polijas vadību ``atteikties no jebkādas Krievijas teritoriālās vienotības pārkāpšanas''.

Jāpasaka arī par \strong{poļu komunistu nostāju}. 1920.~gada 2.~februāra Krievijas VCIK sēdē poļu komunistu vārdā uzstājās J.~Marhļevskis. Viņš norādīja, ka ``kādas armijas no ārienes'' ielaušanās Polijā ``dotu iespēju poļu buržuāzijai uzturēt tautas masās, kuras pusotra gadsimta bija cietušas no nacionālajiem spaidiem, nacionālisma garu ar visām no tā izrietošajām sekām.'' Martā poļu komunisti arī protestēja pret Komunistiskās Internacionāles uzsaukumā visu tautu strādniekiem sakarā ar Polijas jautājumu lietoto terminu ``neatkarīgā Polija'' bez pēdiņām, tā demonstrējot Polijas Republikas suverenitātes atzīšanu. Tiesa, tai pat laikā viņi izteica savu gatavību, ja Polijas valdība turpinās uzbrukumu, aizstāvēt proletārisko Krieviju. Polijas komunisti propagandēja uzskatu, ka J.~Pilsudskis, tāpat kā 1918.~gadā Vācijas ģenerālis Ē.~Lūdendorfs grib nožņaugt Ukrainas Padomju republiku, iesēdināt tās hetmaņa tronī savu ielikteni S.~Petļuru, tāpat kā Ē.~Lūdendorfs gribēja to izdarīt ar P.~Skoropadski. (Starp citu, arī mūsdienu ukraiņu vēsturnieks S.~Lozuņko atzīst, ka J.~Pilsudskis pēc būtības vēlējās kopēt Vācijas rīcību, kura 1918.~gadā atzina Ukrainas Centrālo Radu, bet pēc tam to padzina, varu Kijevā formāli nododot Vācijai paklausīgajam hetmanim P.~Skoropadskim. Visi lieliski sapratuši, ka J.~Pilsudskim ``karš par Ukrainas neatkarību ir vienīgi maskarāde''.)

PKSP I konference 1920.~gada 3.~maijā, tātad jau pēc Polijas armijas iebrukuma Ukrainā sākuma, tā formulēja partijas uzdevumus: ``Šim periodam ir jābeidzas ar valsts organizācijas sabrukumu un varas saņemšanu poļu proletariāta rokās. \citespace{} Cīņai par Ukrainu ir jābeidzas ar cīņu par Polijas strādnieku deputātu Padomju republikas nodibināšanu''. Šo poļu komunistu aktīvo pozīciju padomju vadītāji uztvēra kā poļu darbaļaužu gatavības apliecinājumu iesaistīties pasaules revolūcijā, kaut patiesībā arodbiedrības u.c. masu organizācijas komunistus neatbalstīja.

Īpaša nostāja bija Polijas cietumos pabijušajiem un apmaiņas kārtībā Padomju Krievijā nonākušajiem poļu komunistiem. Nonākuši Maskavā, viņi apgalvoja, ka Polijas armiju var viegli sakaut, sarīkojot sacelšanos tās aizmugurē.

Savs redzējums bija Krievijā un citās padomju republikās jau ilgstoši dzīvojošajiem poļu komunistiem. XX~gadsimta sākumā, kad kapitālisma kvalitatīvo izmaiņu un tās pavadošo grūtību rezultātā krasi pasliktinājās darba cilvēku dzīves apstākļi, internacionālisma ideja bija ļoti populāra. Tā kļuva par ļoti saistošu un pievilcīgu, pasaulē to atbalstīja miljoni cilvēku. Tieši šī ideja daudzu zemju daudziem iedzīvotājiem kļuva par galveno viņu dzīvē, nostumjot visu pārējo, arī nacionāli-patriotiskās jūtas otrajā plānā. Viens no spilgtākajiem piemēriem ir poļu darbaļaužu, kuri dažādu iemeslu dēļ 1920.~gadā bija nonākuši padomju republikās, dalība karā pret savu dzimteni~--- Poliju. Sarkanajā armijā bija daudz poļu un Polijas ebreju (tāpat kā daudz tādu nāciju pārstāvju, kuri nekad nebija pastāvīgi dzīvojuši Krievijas impērijas robežās). Daļu poļu sarkanarmiešu sastādīja pārliecināti revolucionāri. Pēc krievu vēsturnieku vērtējumiem 1919.~gadā Padomju zemē darbojās ap 18~000 poļu komunistu, kad pašā Polijā to bija tikai ap 10~000. Daudzi no viņiem brīvprātīgi iestājās Sarkanajā armijā. Tomēr daļa no sarkanarmiešiem bija poļu bēgļi, kuri iestājās armijā, nevis tādēļ, ka bija guvuši padomisku pārliecību, bet lai gūtu iztiku un jumtu virs galvas. Daudz bija arī tādu, kuri Sarkanajā armijā bija nonākuši parastajā mobilizācijas kārtībā, jo netika atzīta viņu Polijas pilsonība. Jau Pirmajā pasaules karā Krievijas armijas gūstā nonāca daudz poļu. Grūtajos gūsta apstākļos, lai izbēgtu no pusbada dzīves, pazemojumiem, gūstekņi aktīvi reaģēja uz notiekošo Krievijā. Taču Februāra revolūcija maz ko mainīja viņu dzīvē. Turpretī boļševiku Oktobra apvērsums pavēra arī ar Krieviju karojošo valstu karavīriem, ja vien viņi bija gatavi iesaistīties notikumos, necerētas iespējas. Vislielākais vilinājums bija pieslieties boļševikiem, kuram sekoja arī daudzi poļi. Arī 1918.--1920.~gadā Sarkanās armijas gūstā saņemtie poļi nereti iestājās Sarkanajā armijā, jo dzīve gūstā nevienam neradīja prieku, bet sarkanarmiešu rindās pavērās iespēja atgriezties dzimtenē. Protams, kā ikvienā armijā bija arī karjeristi, vienkārši avantūristi.

1920.~gada 3.--5.~maijā Maskavā notika Otrā Viskrievijas poļu komunistu konference, kura nolēma, ka visiem tās dalībniekiem jādodas uz fronti cīņā pret poļu paniem. Konferences rezolūcijā, nosūtītā visām partijas guberņu komitejām, bija norādīts, ka sakarā ar ``plēsonīgo buržuāziskās Polijas valdības uzbrukumu Padomju Krievijai'' poļu komunisti proletāriskās republikas aizstāvībai ``kārtojas rindās'', ka uz fronti tiks nosūtīts maksimāls skaits poļu. 1920.~gada 19.~maijā poļu komunistu sapulcē Maskavā viens no viņu līderiem J.~Ļeščinskis norādīja: ``Poļu komunistu Krievijā uzdevums patreizējā brīdī nosūtīt cik iespējams vairāk [partijas] biedru darbam poļu karaspēka vidū piefrontes joslā, kur ievērojamu iedzīvotāju daļu sastāda poļi.''

Daudzi poļu komunisti Krievijā patiesi vēlējās ātrāk nonākt frontē, lai aizstāvētu revolūciju Krievijā un izraisītu tādu Polijā. Pēc padomju vēsturnieka P.~Goluba datiem cīņās pret ``baltpoļiem'' piedalījās 5~700 poļu komunistu, kas sastādīja 25--30\% no visu poļu komunistu skaita Krievijā, bet pavisam padomju varu ar ieročiem aizstāvēja 18~tūkstošu poļu. Gandrīz visiem viņiem bija pagrīdes, Pilsoņu kara bruņoto cīņu pieredze, viņi bija ietekmīgs spēks cīņā par poļu iedzīvotāju prātiem. Aģitācijas darbs poļu vidū tika veikts ar izdomu. Plaši tika izmantots lozungs ``Par mūsu un jūsu brīvību''. Zinot poļu reliģiozitāti, KK(b)P CK Poļu birojs (pastāvēja 1919--1931) iemanījās Krievijas pilsētās vietējo poļu vidū izdalīt revolucionāro literatūru katoļu baznīcu priekšā.

Tomēr KK(b)P CK Poļu birojs Padomju Krievijas vadītājiem izteica arī Polijas Komunistiskās Strādnieku partijas vadības šaubas un prognozes par notikumu iespējamo attīstību, ja Sarkanā armija ieietu Polijas teritorijā. PKSP vadītāji uzskatīja, ka ir jāievēro poļu iedzīvotāju noskaņojums, kurš ``prasa visupirms turpināt mūsu miera priekšlikumu politiku'', ka Sarkanās armijas tuvošanās poļu teritorijām var izraisīt ``nacionālu karu'', kurā uzvarētu ideja ``aizstāvēt tēvzemi pret ``moskaļiem''''. Viņi uzsvēra, ka ir jānosaka skaidri Padomju Krievijas mērķi gadījumam, ja, neraugoties uz miera sarunām, tai nāksies karot ar Poliju.

Tādejādi poļu komunistu pozīcija bija neviennozīmīga, tajā cerības uz Sarkanās armijas durkļiem atnestu padomju varu Polijai mijās ar reālistiskām bažām par ``nacionālā kara'' izraisīšanās iespēju.

Sarkanās armijas Rietumu fronti Baltkrievijā 29.~aprīlī uzsāka komandēt viens no panākumiem bagātākajiem Sarkanās armijas komandieriem, revolucionārā kara teorētiķis, toreiz 27~gadus vecais M.~Tuhačevskis, kurš bija cēlies no Smoļenskas muižnieku (sastopami arī apgalvojumi, ka no poļu šļahtiču) ģimenes, dienējis cara armijā, bijis vācu gūstā, Pilsoņu karā jau sevi labi atestējis kaujās pret admirāļa A.~Kolčaka un ģenerāļa A.~Deņikina vadītajām armijām. Dienvidrietumu frontes Ukrainā pavēlnieks bija ievērojamais ``sarkano'' komandieris A.~Jegorovs.

M.~Tuhačevska komandētā Rietumu fronte 14.~maijā sāka pretuzbrukumu, drīz tas apstājās, taču atviegloja Sarkanās armijas stāvoli Dienvidrietumu frontē. Poļu karavīri, atkāpjoties laupīja, veica vardarbības aktus, neievērojot 1919.~gada beigās Mikaševičos noslēgto vienošanos, ņēma ķīlniekus (šī vienošanās pēc viņu traktējuma neesot bijusi attiecināta uz Ukrainas teritoriju), nesaudzīgi iznīcināja apdzīvotās vietas. Tā, pēc tam, kad Borisovas pilsētu 25.~maijā bija ieņēmusi Sarkanā armijā, poļu artilērija to divas reizes apšaudīja tik nežēlīgi, ka faktiski visa pilsēta tika iznīcināta. Simti iedzīvotāju gāja bojā. Par šo barbarisma aktu KSFPR ārlietu tautas komisārs G.~Čičerins iesniedza protestus Poliju atbalstošajām Lielbritānijas, Francijas, ASV un Itālijas valdībām.

26.~maijā sākās Sarkanās armijas pretuzbrukums Kijevas rajonā. Turp no Kaukāza bija pārdislocēta S.~Budjonnija komandētā Pirmā Jātnieku armija. Sākotnējie mēģinājumi triecienā forsēt Dņepru neizdevās, tikai 1.~jūnijā Sarkanās armijas daļām ar lielām grūtībām izdevās ieņemt vairākus placdarmus uz ziemeļiem no Kijevas. Izšķirošu ieguldījumu deva Pirmā Jātnieku armija, kura 5.--7.jūnijā pārrāva poļu aizsardzību un virzījās tās aizmugurē. Sekmīgais uzbrukums radīja daudzos boļševikos rožainas cerības. Tā 1920.~gada 2.~jūnijā boļševiku avīze ``\rutxti{Правда}'' (``Taisnība'') publicēja Petrogradas boļševiku vadītāja G.~Zinovjeva apgalvojumu, ka lozungu ``Nāvi poļu paniem un Petļuram!'' drīz sadzirdēs arī ``Sēnas un Temzas krastos''. Tā Sarkanās armijas panākumi iedvesmoja revolūcijas eksporta piekritējus.

10.~jūnijā poļi, baidoties no aplenkuma, atstāja Kijevu, taču pirms tam izvedot no ierindas elektrostaciju, pilsētas kanalizāciju, dzelzceļa pasažieru un preču stacijas. Acīmredzot poļu militāristi bija sapratuši, ka šeit viņi vairs neatgriezīsies.

Polijas armija retumis centās pārtvert iniciatīvu, taču tas tai neizdevās. Uz 17.~jūniju tā atradās apmēram tais pašās pozīcijās, kā pirms sava uzbrukuma aprīlī. Sarkanās armijas Dienvidrietumu fronte gan nespēja aplenkt un iznīcināt Ukrainā iebrukušo poļu karaspēku. Taču Kijeva bija atbrīvota, Polijas armijas pavēlniecība spiesta pārsviest daudz karaspēka no Baltkrievijas uz Ukrainu, kā rezultātā poļu kreisais flangs tika atkailināts. Tas ļāva padomju pavēlniecībai uzsākt uzbrukumu uz Varšavu caur Baltkrieviju pa taisnāko ceļu.

Kad poļu armija uzsāka atkāpšanos, sabiedrībā pret J.~Pilsudski vēlās skarbas kritikas vilnis. Viņu pat sauca par boļševiku aģentu, kurš uztur sakarus ar V.~Ļeņinu pa slepenu telefona līniju. L.~Skuļska valdība Polijā bija spiesta atkāpties, viņa vietā 23.~jūnijā par premjeru kļuva nacionāldemokrāts V.~Grabskis, taču J.~Pilsudskis saglabāja varu. 1920.~gada 1.~jūlijā Seims radīja Valsts aizsardzības padomi (\pltxti{Rada Obrony Państwa}), kuras uzdevums bija risināt kara un miera jautājumus. Tajā ietilpa J.~Pilsudski (priekšsēdētājs), premjerministrs (sākotnēji V.~Grabskis, pēc tam V.~Vitoss), galvenie kabineta locekļi, parlamenta frakciju, militāristu pārstāvji.

Jūlija sākumā padomju Rietumu fronte atkal pārgāja uzbrukumā. M.~Tuhačevskis izdeva pavēli, kurā bija vārdi ``Rietumos tiek lemts jautājums par pasaules revolūcijas likteni, pāri baltās Polijas līķim iet ceļš uz vispasaules ugunsgrēku. Uz Viļņu, Minsku, Varšavu, marš!'' (Polijā šī pavēle tiek ļoti bieži citēta, taču nereti aizmirstot epitetu ``baltās''.) Sarkanās armijas oficiālais ienaidnieks bija ``poļu pans''.

1920.~gada aprīlī V.~Ļeņins nosūtīja KK(b)P CK Sekretariātam zīmīti, kurā ierosināja dot padomju avīžu redaktoriem norādījumu nekrist šovinismā, vienmēr nodalīt poļu panus un kapitālistus no strādniekiem un zemniekiem. Šo nostāju atspoguļojot, padomju dzejnieks V.~Majakovskis 1920.~gadā rakstīja:

\vspace{1.5em}

\noindent
\begin{minipage}{0.45\textwidth}
\rutxti{Мы воюем с панским родом,\\
а не с польским народом.}
\end{minipage}
\hspace{2em}
\begin{minipage}{0.45\textwidth}
Mēs pret poļu panu dzimtu,\\
Bet ne tautu karojam.
\end{minipage}

\vspace{1.5em}

% page 413


Taču poļu strādnieku un zemnieku vairākums šādās ``niansēs'' neiedziļinājās, valdošais bija uzskats, ka krievu boļševiki uzbrūk viņu tēvzemei~--- Polijai. Kā norādījis poļu militārais vēsturnieks \emph{Z.~Zaluskis}, termins ``baltpoļi'' sākotnēji kalpoja lai atšķirtu īstos poļus no poļu kontrrevolucionāriem, taču tuvojoties Vislai, šo ``baltpoļu'' kļuva arvien vairāk, bet pie Varšavas radās iespaids, ka citu poļu nemaz nav.

1920.~gada 5.~maijā viens no boļševiku līderiem Ļ.~Trockis, uzstājoties VCIK, Maskavas Strādnieku deputātu padomes un Arodbiedrību padomes apvienotajā sēdē, skaidroja boļševiku cerības: ``Poļu strādnieku šķira negribēja un negrib karu. Poļu zemniecība no jaunā režīma, no Pilsudska režīma ir saņēmusi tikai drumslas, pareizāk sakot, ir saņēmusi tikai solījumus iedot drumslas. Tā nevar ilgi balstīt šo režīmu, tā nevar tiekties uz karu, kas atnesīs tai valsts nodokļu padziļināšanos, bet ja karš ieilgs, tad arī pilnīgu grimšanu nabadzībā, pilnīgu spēku izsīkumu. Tas viss ir neapstrīdami, taču šodien to vēl visa poļu zemniecība nesaprot. Tajā vēl ir stipri nacionālie aizspriedumi. Polija pēc ilga jūga vēl pārāk nesen ir ieguvusi neatkarīgas republikas stāvokli. Nacionālais jūtīgums vēl pārlieku svaigs, vēl nav izdzīvoti valstiskās neatkarības medus mēneši, un uz šiem noskaņojumiem cenšas būvēt savu politiku ``poļu valsts priekšnieks'' Pilsudskis. Vēl nav izzudusi visatpalikušāko zemniecības masu neuzticība un naids pret Krieviju un krievisko, jo tās atmiņā Krievija un krieviskais ir cars un cariskais.'' Taču blakus šiem diezgan reālistiskajiem vērtējumiem Ļ.~Trockis arī paziņoja: ``Poļu šļahtiči saka, ka krievus, šos barbarus un skitus, vajag atvirzīt tālāk uz austrumiem. Bet mēs tiecamies uz rietumiem, pretī Eiropas proletariātam, kurš zin, ka mēs varam sastapties ne citādi, kā tikai pār baltgvardu Polijas līķi, brīvā un neatkarīgā strādnieku un zemnieku Polijā.''

Reālistiski situāciju Polijas karaspēka aizmugurē avīzē ``\rutxti{Правда}'' 25.~un 26.~maijā vērtēja Dienvidrietumu frontes Kara padomes loceklis J.~Staļins: ``Atšķirībā no Kolčaka un Deņikina aizmugures poļu karaspēka aizmugure ir viendabīga un nacionāli vienota. No šejienes arī vienotība un izturīgums. Tās pārsvarā esošais noskaņojums~--- ``tēvzemes sajūta'' \citespace{} Protams, Polijas aizmugure nav viendabīga (un nevar būt viendabīga!) šķiriskā ziņā, taču šķiriskie konflikti vēl nav sasnieguši tādu spēku, lai pārrautu nacionālās vienotības sajūtu un inficētu ar pretrunām dažādu šķiriskā ziņā fronti. Ja poļu karaspēks darbotos pašas Polijas rajonā, ar to, neapšaubāmi, būtu grūti cīnīties.''

Taču Padomju vadītāji karu pret Poliju uzskatīja arī kā cīņu pret visu tai laikā jau pamatos radīto Versaļas sistēmu, pasaules revolūcijas procesa sastāvdaļu. Tāpēc viņi ļoti vēlējās pēc iespējas ātrāk redzēt Ļ.~Trocka pieminēto ``baltgvardu Polijas līķi''. Tas aizmigloja skatu arī tiem, kuri tikko vēl bija runājuši par poļu nacionālo vienotību.

Poļu~--- padomju kara gaitā radušos situāciju varēja izmantot \strong{Lietuva}. No 1919.~gada aprīļa, kad, padzenot Sarkano armiju, Viļņu ieņēma Polijas karaspēks un faktiski iekļāva to Polijas sastāvā, visi Polijas valdības mēģinājumi noregulēt Polijas un Lietuvas attiecības atdūrās pret Lietuvas prasību atdot tai senās lietuviešu zemes ar Viļņu, ko savukārt pirmā kategoriski atteicās darīt. Situācija mainījās 1920.~gada vasarā, sākoties Sarkanās armijas pretuzbrukumam Polijai. Jūlijā Viļņu ieņēma Sarkanā armija, 12.~jūlijā tika parakstīts Lietuvas un Padomju Krievijas miera līgums, saskaņā ar kuru Lietuva ieguva tiesības uz Viļņu ar apkārtējo apgabalu. Padomju Krievijas vadība acīmredzot cerēja, ka Polijai zaudējot karu, arī Lietuvā varēs nodibināt Padomju varu, tāpēc nepretojās Viļņas nodošanai Lietuvas rokās. Reāli šīs zemes Padomju Krievija nodeva Lietuvai tikai pēc pusotra mēneša, kad kļuva skaidrs, ka Sarkanā armija pie Varšavas ir cietusi sagrāvi. Tā, pateicoties padomju-poļu karam, Lietuva panāca savas neatkarības atzīšanu, ieņēma Viļņu un Suvalku rajonu, uz kuru pretendēja arī Polija. Pie tam Lietuvas karaspēks, izmantojot labvēlīgo situāciju, iegāja poļu aizmugurē un ieņēma Trāķus (lietuviešu \lttxti{Trakai}, poļu \pltxti{Troki) un Landvarovas} (lietuviešu \lttxti{Lentvaris}, poļu \pltxti{Landwarowo, Landwarów}, krievu \rutxti{Ландварово}) dzelzceļa mezglu. Boļševiku vadīto Lietuvas un Baltkrievijas padomju republiku apvienošana nebija devusi cerētos rezultātus, Lietuvā pie varas bija nacionālā valdība, tāpēc 1920.~gada 31.~jūlijā atkal tika radīta ar Lietuvu nesaistīta Baltkrievijas Sociālistiskā Padomju Republika, kuru pēc PSRS izveides 1922.~gadā pārdēvēja par Baltkrievijas Padomju Sociālistisko Republiku.

Ātrums, ar kādu Sarkanā armija 1920.~gada jūlijā virzījās uz rietumiem, bija pārsteidzošs. Īsā laikā tā pavirzījās uz priekšu vairāk nekā 600 km un 11.~jūlijā ieņēma Minsku, 19.~jūlijā~--- Grodņu. 26.~jūlijā Belostokas (poļu \pltxti{Białystok}, lietuviešu \lttxti{Balstogė}, baltkrievu \betxti{Бялисток}) rajonā tā iegāja pašas Polijas teritorijā. 1.~augustā tās rokās krita Bresta. Sarkanā armija sasniedza Kērzona līniju. Ukrainas Direktorija ar S.~Petļuras armiju evakuējās uz Poliju. Kopā emigrēja ap 40~tūkstošu cilvēku.

Poļu historiogrāfijā vēl līdz šai dienai ilgst strīdi, kādi bija galvenie cēloņi, kāpēc Polijas armija cieta tik pamatīgu sakāvi 1920.~gada jūnijā~--- jūlijā. Vieni vaino J.~Pilsudski, kurš uzsāka šo avantūristisko karagājienu uz Kijevu, citi~--- ģenerāli S.~Šepticki, kurš tai laikā komandēja Lietuvas-Baltkrievijas fronti, vēl citi raksta par spēku pārsvaru M.~Tuhačevska vadīto Sarkanās armijas daļu pusē. Acīmredzot atbilde, kura apmierinātu visus strīdniekus, tā arī nekad netiks atrasta. Kā uzsvēris krievu vēsturnieks G.~Matvejevs, tajā laikā ar pastāvošajiem sakaru līdzekļiem precīzu un drošu karaspēka vadību, kad vispārēja uzbrukuma vai atkāpšanās apstākļos virspavēlniecību pavēles pienāca ar nokavēšanos vai vispār nesasniedza adresātu, neviens nevarēja garantēt, un dažādos līmeņos tika pieļauts milzums kļūdu. Lai nu kā, bet poļu pamatspēki, tāpat kā 1920.~gada aprīlī--maijā Sarkanās armijas daļas, izglābās. Tomēr daudziem likās, ka Polijas valsts liktenis tobrīd karājas mata galā.

Antantes Augstākās padomes konference Spā (Beļģija) 1920.~gada jūlijā vēlreiz apstiprināja jau minēto Parīzes miera konferences speciālā komisijā izstrādāto un Augstākās padomes 1919.~gada 8.~decembrī speciālā deklarācijā apstiprināto Polijas pagaidu robežlīniju, noteiktu pēc etniskā principa (apmēram pa Bugas upi). Polijas premjerministrs V.~Grabskis 9.~jūlijā konferencē Spā paziņoja: ``Polija saprot, ka tā pati ir vainīga, nonākot tādā stāvoklī, ka tai jāmaina sava politika kā pret kaimiņiem, tā arī sabiedrotajām lielvalstīm. Tā atzīst nepieciešamību nodot jautājumu par savām dzīves interesēm, pat jautājumu par savām robežām, sabiedrotajām lielvalstīm. Līdz šim brīdim, kaut tai tika doti padomi šādi rīkoties, tā paļāvās uz savu pašas militāro spēku. Stipras personības ar plašiem plāniem novirzīja Poliju no pareizā ceļa, taču šie plāni neatbilda ne veselajam saprātam, ne milzīgas tautas daļas patriotisma jūtām. Ne Kijeva, ne 1772.~gada robeža nav nacionāli mērķi. Tas bija pie varas esošo sapnis, bet ne iedzīvotāju vairākuma vēlme. Patreizējā brīdī stāvoklis ir ļoti bīstams, tāpēc tagad nav īstais brīdis nodarboties ar savstarpējiem uzbrukumiem un apvainojumiem. Valsts galva Pilsudska kungs atzīst savu kļūdu, un Polijas vadība izsaka savu piekrišanu neuzsākt strīdus, bet iet pa jaunu ceļu''. V.~Grabskis tāpat piekrita Lielbritānijas premjerministram D.~Loidam Džordžam, kurš atzina, ka Polijas politika attiecībā pret Austrumgalīciju ``bija vēl viena kļūda''.

Jāpasaka, ka D.~Loida Džordža nostāja vismaz vārdos bija principiāla. Viņš atzina, ka ``ir bezmērķīgi apvainot vienam otru, taču tomēr ir jāatgādina pagātne, pirms solīt palīdzību nākotnē. Par nelaimi, Polija nebija apmierināta ar savām valstiskajām robežām. Tā saķildojās ar Lietuvu un sagrāba lietuviešu teritoriju. Tā iebruka Krievijā un okupēja teritoriju, kuru apdzīvo 20 līdz 30~miljonu krievu [Britu premjers šeit par krieviem nosaucis arī ukraiņus un baltkrievus~--- V.Š.]. Tā uzstāja 3~500 tūkstošu ukraiņu [Domāti Austrumgalīcijas ukraiņi~--- V.Š.] pievienošanai, neraugoties uz viņu aso protestu. Tā nespēja vienoties ar čehiem. Tā radīja grūtības Dancigas jautājumā. Polija faktiski ir saķildojusies ar visiem saviem kaimiņiem, ieskaitot Vāciju. Par to visu ir nepieciešams pateikt nevis tāpēc, lai pārmestu Polijai, bet tāpēc, ka to prasa lieta.'' D.~Loids Džordžs arī paziņoja, ka ``nevar griezties pie Anglijas parlamenta pēc palīdzības Polijai, pirms [viņš] spēs pateikt, ka Polija pilnībā ir atteikusies no aneksijām, un lūdz vienīgi nodrošināt neatkarību tās etnogrāfiskajās robežās, un pirms spēs pateikt, ka tā ir gatava noslēgt taisnīgu vienošanos ar Lietuvu, likt miera ar Krieviju pamatā etnogrāfiskās robežas, ievērot Austrumgalīcijas tautas gribu šīs zemes nākamā likteņa noteikšanai, vienoties ar čehiem un necensties turpmāk panākt poļu kundzību Dancigā''. D.~Loids Džordžs nepiekrita Rietumu valstu militārai intervencei Polijā, bet apsolīja tikai daļēji piekāpties tai politiski Lietuvas, Dancigas, poļu un čehu attiecību jautājumos. Lielbritānijas kolēģi atbalstīja arī Francijas premjerministrs A.~Miljerāns. V.~Grabskis atbildēja, ka pirms izbraukšanas no Polijas viņš tika pilnvarots pieņemt šādus noteikumus, Polija tagad vēlas vienīgi drošību savas nacionālās pastāvēšanai un ir gatava rēķināties ar kaimiņtautu interesēm. Tādejādi Polijas premjerministrs, ievērojot šai brīdī nelabvēlīgo kara gaitu, apmaiņā pret Lielbritānijas starpniecību sarunās ar Padomju Krieviju par karadarbības pārtraukšanu, un lai saņemtu no Rietumu lielvalstīm kara materiālus, bija spiests piekrist to spiedienam. Viņš 1920.~gada 10.~jūlijā sabiedroto konferencē Spā parakstīja vienošanos ar Lielbritānijas un Francijas valdībām, piekrītot pamieram ar Padomju Krieviju atbilstoši līnijai, Antantes Augstākās padomes noteiktai jau 1919.~gada 8.~decembrī. Šai pašas Polijas izraisītajā grūtajā situācijā 10.~jūlijā S.~Grabskis arī piekrita, ka Viļņas likteni izlems sabiedroto Augstākā padome un pagaidām pilsētu, lai tajā neieietu Sarkanā armija, pārņems Lietuva. Tā Polijas premjers oficiāli atteicās no pretenzijām uz ukraiņu, baltkrievu un lietuviešu zemēm. (Pēc situācijas maiņas Polijas valsts priekšnieks un valdība šo attiekšanos ignorēja.)

Polijas Valsts aizsardzības padome 13.~jūlijā apstiprināja konferencē Spā panāktās vienošanās, neraugoties uz dažu tās locekļu iebildumiem.

Uzskatot J.~Pilsudski par personīgi atbildīgu par militārajām un politiskajām sakāvēm, nacionāldemokrātu līderi R.~Dmovskis un S.~Glabinskis 19.~jūlijā Valsts aizsardzības padomes sēdē izvirzīja jautājumu par J.~Pilsudska atstādināšanu no virspavēlnieka amatā. Taču J.~Pilsudskis faktiski uzstādīja Padomei ultimātu. Norādot, ka kritika gan Aizsardzības padomē, gan presē grauj viņa autoritāti, viņš lika priekšā Padomei vai nu atbrīvot viņu no valsts priekšnieka amata, vai nodrošināt vienprātīgu atbalstu. Pēc tam viņš un padomes locekļi militāristi atstāja sēžu zāli. Aizsardzības padomes locekļi~--- civilisti neuzdrošinājās šādā valstij grūtajā brīdī sašķelt spēkus un izteica ``pilnīgu uzticību valsts priekšniekam un virspavēlniekam''. R.~Dmovskis un S.~Glabinskis izstājās no Valsts aizsardzības padomes sastāva. (Sekojošais ``brīnums pie Vislas'', par kuru tiks tālāk runāts, atkal cēla J.~Pilsudska autoritāti un notikumi pavērsās citā virzienā.)

Boļševiku partijas vadībā domas par Antantes priekšlikumu dalījās. Revolucionārās kara padomes priekšsēdētājs un KSFPR kara un jūras lietu tautas komisārs Ļ.~Trockis ierosināja kā pirmo formulēt tēzi, ka Polija ir neatkarīga valsts, kuras neaizskaramību boļševiki nekad nav nolieguši. Viņš piedāvāja pieņemt Lielbritānijas vidutājību, garantējot Kērzona līnijas neaizskaramību, taču nenosakot galīgi robežu, jo tas ir tautu pašnoteikšanās jautājums. Ārlietu tautas komisārs G.~Čičerins uzsvēra: ``Polijas sovjetizācija ar moskaļu durkļiem būtu avantūra''. Turpretī KK(b)P CK Politbiroja loceklis Ļ.~Kameņevs uzskatīja, ka vienīgā garantija pret kara turpināšanos var būt Polijas sovjetizācija (Lai atceramies J.~Pilsudska uzskatus par miera garantijām).

Kopumā Krievijā, Ukrainā un Baltkrievijā valdīja maldīgs priekšstats par spēku samēru Polijā, jo informācija no tās pienāca pretrunīga, bieži nepareiza. Tā, piemēram, 1919.~gada janvārī KK(b)P CK locekļa kandidāts, Padomju Krievijas VCIK loceklis A.~Joffe ziņoja, ka nacionālisma dvinga tuvākajā laikā aptvers ne tikai buržuāziju, sīkpilsonību, zemniecību, bet arī proletariātu. Taču padomju līderi šādam vērtējumam nevēlējās ticēt. Citu, arī augstāk situētu padomju darbinieku viedoklis bija atšķirīgs. Tā F.~Dzeržinskis 6.~augustā V.~Ļeņinam ziņoja, ka lietas Polijā virzās lieliski: buržuāzija ir panikā, brīvprātīgo, kuri vēlas karot pret Sarkano armiju, tik maz, ka to formācijās uzņem pat divpadsmitgadīgus bērnus. Pēc 11 dienām, kad stāvoklis frontē strauji saasinājās, F.~Dzeržinskis jau mainīja optimistisko toni. Viņš ziņoja, ka Varšavas strādnieki, kaut arī gaida Sarkanās armijas atnākšanu, taču paši nekādu aktivitāti neizrāda, norādīja, ka Polijas valdība arestē poļu komunistus un ebreju bundistus, bet PPS līderi, kam ir liela ietekme masās, aicina aizstāvēt Polijas galvaspilsētu pret ``sarkanajiem''.

Savulaik KSDS(b)P VII konferencē 1917.~gada aprīlī V.~Ļeņins norādīja: ``Neviens tā nav apspiedis poļus kā krievu tauta. Krievu tauta cara rokās bija Polijas brīvības bende. Nav tautas, ko tā būtu pārņēmis naids pret Krieviju, nav tautas, kas tik briesmīgi neieredzētu Krieviju kā poļi, un no tā izriet dīvaina parādība. Poļu buržuāzijas dēļ Polija ir traucēklis sociālistiskajai kustībai. Lai kaut visa pasaule deg ugunīs,~--- ka tikai Polija būtu brīva.'' Tālāk, runājot par sociāldemokrātu attieksmi pret apspiesto nāciju tiesībām atdalīties no apspiedējām valstīm, viņš gan uzsvēra: ``Mums, krieviem, jāuzsver atdalīšanās brīvība, bet Polijā~--- savienošanās brīvība''. KK(b)P VIII~kongresā 1919.~gada 19.~martā, referējot par partijas programmu, V.~Ļeņins saprātīgi teica: ``Neatzīt to, kas ir, nedrīkst: tas pats piespiedīs sevi atzīt. Dažādās zemēs proletariāta un buržuāzijas norobežošanās norit savdabīgā ceļā. Šajā ceļā mums jārīkojas ārkārtīgi uzmanīgi. Sevišķi uzmanīgiem jābūt pret dažādām nācijām, jo nav nekā ļaunāka par nācijas neuzticēšanos. Pie poļiem notiek proletariāta pašnoteikšanās.\citespace{} Bet \citespace{} pašlaik stāvoklis ir tāds, ka vairākums poļu strādnieku, kas ir vairāk attīstīti, kulturālāki nekā mūsējie, turas pie sociālaizsardzībniecības, sociālpatriotisma viedokļa. Vajag nogaidīt \citespace{} nav ne mazāko šaubu, ka pašlaik nedrīkst neatzīt poļu nācijas pašnoteikšanos \citespace{} strādniekus tur biedē ar to, ka moskaļi, lielkrievi, kas vienmēr poļus žņauguši, Polijā gribot ienest savu lielkrievisko šovinismu, kas maskēts ar komunisma nosaukumu.''

1920.~gada 15.~jūlijā V.~Ļeņins nosūtīja Rietumu frontes Revolucionārajai kara padomei telegrammu uz tās locekļa J.~Unšlihta vārda. Tajā bija teikts: ``Paziņojiet kā jūs un citi biedri vērtē tādu taktiku.

\begin{enumerate}
\item Mēs ļoti svinīgi pasludinām, ka nodrošinām poļu strādniekiem un zemniekiem robežu uz austrumiem no tās, kuru dod Kērzons un Antante.

\item Mēs sasprindzinām visus spēkus lai sakautu Pilsudski.

\item Mēs ieejam pašā Polijā tikai uz īsu laiku, lai apbruņotu strādniekus un tūlīt no tās aizejam.

\item Vai jūs uzskatāt par iespējamu un cik ātri padomju apvērsumu Polijā.
\end{enumerate}

Ļeņins.''

% page 417


Telegramma rāda, ka boļševiku vadonis svārstījās jautājumā par stāvokļa vērtējumu Polijā un iespējamo krievu boļševiku taktiku tajā. Zīmīgi, ka šis dokuments Polijā tika publicēts tikai 1964.~gadā, bet PSRS~--- vēl pēc 26~gadiem.

Tomēr galu galā Padomju Krievijas valdības vadītājs V.~Ļeņins iestājās par uzbrukuma paātrināšanu Polijai. Pārsvaru padomju vadībā guva V.~Ļeņina un to boļševiku līderu viedoklis, kuri uzstāja par kara turpināšanu nu jau Polijas teritorijā. Pēc neilgām konsultācijām un domu apmaiņas padomju vadībā galīgi nobrieda uzskats par nepieciešamību turpināt karu. KK(b)P CK plēnums 16.~jūlijā Dž.~Kērzona priekšlikumu ar balsu vairākumu noraidīja. 17.~jūlijā padomju valdība deva norādījumu nepārtraukt militārās operācijas Rietumu un Dienvidrietumu frontē arī tad, ja ``mēs būtu spiesti uz laiku pāriet Antantes noteikto robežu.'' Jau 1920.~gada jūlijā politiskajā darbā Sarkanajā armijā pamatakcents tika likts vairs ne uz padomju teritorijas atbrīvošanu, kas jau bija acīmredzama, bet uz kara turpināšanu sociālisma uzvaras labā visā pasaulē.

Jāatzīst, ka 1920.~gadā V.~Ļeņinam pašam pietrūka tik nepieciešamās cieņas pret poļu tautas ``atdalīšanās brīvību'', pacietības rīkoties ``uzmanīgi'', viņš pārāk optimistiski cerēja uz poļu ``proletariāta pašnoteikšanos'', tā iesaistīšanos cīņā nevis par nacionālo neatkarību, bet par proletariāta vispasaules uzvaru. Lēmumu pāriet robežu ar Poliju noteica padomju politiķu pārliecība par to, ka šo soli atbalstīs Polijas sabiedriskā doma, poļu darbaļaužu praktiskie darbi. Pirmkārt, boļševiki vēlējās redzēt tieši tādu notikumu attīstību, pieņēma vēlamo par esošo. Otrkārt, informācija no Polijas par tās iedzīvotāju noskaņojumu pienāca ļoti pretrunīga. Boļševiki taču apgalvoja, ka darbaļaudīm kapitālisma apstākļos nav tēvzemes, patriotisms var attīstīties tikai sociālisma apstākļos. Arī poļu komunisti kļūdaini ticēja, ka poļu tautas masās internacionālisma idejas ir iesakņojušās dziļāk nekā patriotisms, uz šīs ticības pamata sūtīja faktiski dezorientējošas ziņas par masu noskaņojumu.

Vēl vairāk nekā boļševiku vadītājiem ``uzmanības'' pret poļu nacionālajām jūtām pietrūka praktiskajiem ``proletariāta pašnoteikšanās'' realizētājiem~--- Sarkanās armijas komandieriem un politdarbiniekiem. 5.~augustā Rietumu frontes pavēlnieks M.~Tuhačevskis izdeva pavēli visus vietējos iedzīvotājus skarošos rīkojumus un sludinājumus izdot poļu valodā, mītiņos ļaut runāt tikai tādiem krieviem, kuri pietiekoši orientējas ``poļu lietās'', neaizskart poļu iedzīvotāju reliģiskās jūtas, nemainīt poļu izkārtnes pret krievu, ieteikt visiem militārajiem orķestriem bez ``Internacionāles'' atskaņot arī ``Varšavjanku'' un citas poļu revolucionārās dziesmas, aģitācijas plakātos krieviskos uzrakstus nomainīt pret poļu. Arī poļu komunisti uzsvēra nepieciešamību Sarkanajai armijai ievērot savu tautiešu nacionālās jūtas. Taču, kā trāpīgi rakstīja poļu vēsturnieks Z.~Zaluskis: ``Velti Juliāns Marhļevskis aicināja rūpēties par ``poliskuma'' saglabāšanu. Pievolgas un Donas divīziju revolucionārās komitejas padomju varu proklamēja krieviski un žargonā.'' (Ar ``žargonu'' domāta jidiša valoda, kuru pārsvarā lietoja Krievijā un Polijā dzīvojušie ebreji. Pavisam XX gadsimta sākumā šai valodā runāja ap 11 miljonu ebreju visā pasaulē~--- V.Š.)''.

Krievu vēsturnieks S.~Poltoraks vērtē, ka poļu-padomju karš bija avantūra ``kvadrātā'', divpusēja avantūra. J.~Pilsudska un viņa sekotāju ambīcijas ar karu nostiprināt Polijas valsts un valdības stāvokli sadurās ar padomju vadītāju atkal uzplaukušajām cerībām panākt Polijas sovjetizāciju. Padomju līderi sprieda, ja jau karš ir sācies, tas jānoved līdz galam, lai Padomju Krievijai būtu robeža ar strādnieku-zemnieku Poliju, nevis ar ``balto panu'' Poliju, kā arī uzpūstu pasaules revolūcijas liesmu.

Kad Sarkanā armija pārgāja agrāk ar Polijas karalisti noteikto robežu, radās kvalitatīvi jauna situācija. Tas bija solis, kura nozīmi, kā uzskata S.~Poltoraks, nesaprata neviens no padomju līderiem, arī V.~Ļeņins. Pēc padomju republiku vadītāju domām tas bija karagājiena par pasaules proletariāta atbrīvošanu sākums. KK(b)P CK plēnums 16.~jūlijā kā tiešo karadarbības uzdevumu formulēja Polijas sovjetizāciju. Tiesa, tā kā šai laikā bija sākušās Padomju Krievijas tirdzniecības sarunas ar dažām Rietumu valstīm un, pirmkārt, Lielbritāniju, boļševiku vadība uzskatīja par labāku sovjetizāciju veikt ne tikai ar Sarkanās armijas spēkiem, bet uzspiežot Polijai mieru noteikumus, kas liktu tai faktiski atbruņoties~--- samazināt armiju līdz 50~tūkstošiem vīru, radīt strādnieku miliciju u.c. Pēc tam bija jāizveido Polijā revolucionāra varas iestāde (23.~jūlijā Smoļenskā arī tika nodibināta Polijas Pagaidu Revolucionārā komiteja~--- \pltxti{Tymczasowy Komitet Rewolucyjny Polski, Polrewkom}), kura organizētu savu Sarkano armiju, kas kopā ar strādnieku miliciju nodibinātu Polijā padomju varu.

Ar to karš no iekarošanas kara, kāds tas līdz tam bija no Polijas puses, un atbrīvošanās kara no padomju puses, pārvērtās par Polijas aizstāvēšanās karu un Padomju republiku vestu iekarošanas karu. Abu pušu līdzšinējās pozīcijas it kā apmainījās vietām.

Tagad poļu iedzīvotāji varēja skatīties uz padomju karavīriem kā uz okupantiem, aizmirstot, ka nupat-nupat okupanti bija paši poļi. Turpretī daudziem Sarkanās armijas karavīriem un virsniekiem, neraugoties uz boļševiku izvērsto pasaules revolūcijas un internacionālā pienākuma propagandu, zuda stimuls turpināt karu. Tie Krievijas, Ukrainas, Baltkrievijas un Lietuvas patrioti, kuri līdz tam bija cīnījušies par savu padomju dzimteni, tagad uzskatīja mērķi par sasniegtu, nevēlējās karot svešā zemē.

Boļševiku vadoņi cerēja ar Sarkanās armijas ieiešanu Polijā, bet pēc tam arī Vācijā, kur pēc viņu aprēķiniem sāktos ar Versaļas līgumu neapmierināto vācu strādnieku sacelšanās, izraisīt revolucionārus satricinājumus un veicināt pasaules revolūcijas sākumu. G.~Čičerina notā Lielbritānijas valdībai bija norādīts, ka laikā, kad Polija gatavoja savu uzbrukumu Padomju Krievijai, Lielbritānija neko nebija darījusi lai to novērstu, bet, kad Polija bija cietusi militāru sakāvi, piedāvā vidutājību, kas neļauj to uzskatīt par neieinteresētu pusi. KSFPR Tautas Komisāru Padome paziņoja Lielbritānijas valdībai, ka tās vidutājības piedāvājums tiek noraidīts un deklarēja savu gatavību vest sarunas tieši ar pašu Poliju, vienlaikus solot, ka tiks noteikta ``poļu tautai izdevīgāka teritoriālā robeža'' nekā Kērzona līnija. Savukārt 20.~jūlijā Dž.~Kērzons notā padomju valdībai norādīja, ka, ja arī pēc Polijas valdības pamiera lūguma padomju karaspēks turpinās uzbrukumu, Lielbritānija sniegs Polijai tās lūgto palīdzību.

V.~Grabska vadītā Polijas valdība 22.~jūlijā griezās pie Padomju Krievijas ar priekšlikumu uzsākt tiešas sarunas. Taču sarunu uzsākšanu vilcināja gan padomju, gan poļu puse. 23.~jūlijā padomju pārstāvji Polijai paziņoja, ka Sarkanās armijas pavēlniecība saņēmusi norādījumus nekavējoties sākt sarunas ar Polijas armijas pavēlniecību par pamiera noslēgšanu un miera līguma sagatavošanu, taču vienlaikus Sarkanajai armijai tika dota pavēle doties uz Varšavu ``lai galīgi sakautu pretinieku''. Padomju puse norādīja, ka visos līdz šim notikušajos militārajos konfliktos pirms nebija noslēgts pamiers, abas puses varēja izvērst kara darbību. Polijas un Padomju Krievijas delegācijas gan sastapās 1.~augustā, bet poļu delegācijai bija dotas instrukcijas vest sarunas tikai par pamieru. Padomju Krievija, izmantojot sev labvēlīgo situāciju, pasludināja poļu delegācijas pilnvaras par nepietiekamām. Polijas vadība, cerot uz Antantes diplomātisku un militāru atbalstu, nesteidzās piešķirt saviem delegātiem mandātu sākt sarunas arī par mieru. Sarunu sākšanās kavējās, par ko abas puses centās uzvelt vainu viena otrai.

Polijas valdība, 11.~augustā apspriežot miera noteikumus, uzdeva savai delegācijai miera sarunās aizstāvēt viedokli, ka Ukrainai pašai jānosaka sava pārvaldes forma, tā joprojām atsakoties redzēt S.~Petļuras vadīto ukraiņu nacionālistu zaudējumu cīņā par ietekmi Ukrainā; atzīt Baltkrievijas neatkarību, bet sarunas par tās robežām nevest; protestēt pret Viļņas nodošanu Lietuvai. Ar to Polijas valdība faktiski pakāpeniski atgāja no solījumiem Rietumu lielvalstīm piekrist savas valsts etnogrāfiskām robežām.

Augstākajā Sarkanās armijas panākumu punktā 12.~augustā tika publicēti Padomju Krievijas valdības priekšlikumi Polijas valdībai, kuri bija izstrādāti KK(b)P CK plēnumā 16.~jūlijā. Par robežu gan tika piedāvāts atzīt Kērzona līniju, bet reizē Padomju valdība izvirzīja Polijai nepieņemamas prasības: samazināt armiju līdz 50~000 vīru (Padomju Krievija tad pie Polijas robežas turētu ``tikai'' 200~000 karavīru), iekšējās kārtības uzturēšanai radīt strādnieku miliciju, piešķirt bez maksas zemi visu karā kritušo, ievainoto un darba spējas zaudējušo poļu ģimenēm, realizēt pilnīgu politisku un militāru amnestiju utt. Šīs prasības krievu-ukraiņu-poļu miera sarunās, reizē deklarējot Polijas neatkarības un suverenitātes atzīšanu, izklāstīja Padomju Krievijas un Padomju Ukrainas delegācijas vadītājs latviešu lielinieks Jūlijs Kārlis Daniševskis. Tiesa, 14.~augustā, dienu pirms izšķirošās kaujas, V.~Ļeņins lika viņam paziņot, ka savu runu tam jāsāk ar svinīgu paziņojumu par Polijas neatkarības un suverenitātes atzīšanu, robežu piešķiršanu, kuras būtu labākas par Kērzona līniju, padomju puses atteikšanos no kontribūcijām. Taču šīs principiālās deklarācijas nevarēja aizmaskot citu padomju noteikumu Polijai nepieņemamo raksturu. Iespējams arī, ka padomju priekšlikumi bija izvirzīti, lai tos pretējā puse noraidītu un uzvarošs Sarkanās armijas karagājiens varētu turpināties. Kā konstatējis krievu vēsturnieks V.~Dainess: ``Mieru varēja sasniegt arī krietni agrāk un tādēļ nebija nepieciešamības pāriet Polijas etnogrāfisko robežu''. Bijušais Polijas prezidents V.~Jaruzeļskis norādījis, ka ``šī kara apstākļi, kā poļu ulānu veiktā Kijevas sagrābšana, Sarkanās armijas atbildes gājiens uz Varšavu noveda pie postošām sekām, krasi negatīvi ietekmēja tālākās poļu-padomju attiecības. Tika iesaldēts pagātnes palieku pārvarēšanas un mūsu tautu tuvināšanās process.''

Dziļākos cēloņus, kāpēc boļševiku vadoņi tā rīkojās, parāda tikai pēc PSRS sabrukuma publiskotais V.~Ļeņina nolasītais KK(b)P CK atskaites ziņojums partijas IX~konferencē 1920.~gada 22.~septembrī jau pēc Sarkanās armijas sakāves pie Varšavas. ``Es lūdzu mazāk pierakstīt; tas nedrīkst nonākt presē'', brīdināja V.~Ļeņins konferences delegātus. Konferencē viņš runāja atklāti, diplomātiskus izteicienus nemeklēja, partijas vadības atbildību par militāro sakāvi nenoliedza. Viņš atzina, ka Antantes piedāvātā Kērzona līnija ``mums bija ļoti izdevīga''. Taču karu ar Poliju V.~Ļeņins kvalificēja kā lūzuma momentu proletariāta un imperiālisma spēku cīņā, kad revolucionārajiem spēkiem bija radusies iespēja no aizsardzības pozīcijas pāriet uzbrukumā lai palīdzētu sovjetizēt Poliju. ``Bet tagad mēs teicām: mēs tagad esam kļuvuši stiprāki un uz katru jūsu uzbrukuma mēģinājumu mēs atbildēsim ar pretuzbrukumu \citespace{} fakti, ka Polijā labi attīstīts proletariāts un labāk audzināts lauku proletariāts, šie fakti saka mums: jums ir jāpalīdz tiem sovjetizēties.'' V.~Ļeņins atzinās: ``\dots{}savā starpā mēs runājām, ka mums ar durkli jāpārbauda vai Polijā nav nobriedusi proletariāta sociālā revolūcija''. Tātad boļševiku vadoņi subjektīvi nedomāja nedz par nāciju pašnoteikšanās tiesību principa realizēšanu, nedz par Krievijas nacionālajām interesēm, bet gan par utopiskajiem pasaules revolūcijas realizācijas plāniem. Cita lieta, cik lielā mērā šie plāni objektīvi bija īstenojami.

To realizēšanā tika cerēts uz poļu ``darbaļaužu'' atbalstu. Taču poļu komunisti atradās sarežģītā situācijā. PKSP CK līdz 1920.~gada jūnija beigām nepieņēma nevienu oficiālu lēmumu jautājumā par to, vai ir vērts tiekties pēc Polijas etnogrāfiskās robežas pāriešanu, Varšavas ieņemšanas un zemes sovjetizācijas. Sākotnēji, kad vēl Polijas armija atradās ukraiņu un baltkrievu teritorijās, poļu komunisti atradās sakāvnieciskās pozīcijās (t.i.~--- vēlēja savai valstij sakāvi). Viņu pašu nespēja izraisīt revolucionāru pacēlumu daļā no tiem radīja vēlmi, lai revolūcijai Polijā palīdzētu Padomju Krievijas karaspēks. Taču Sarkanajai armijai virzoties jau Polijas teritorijā vietējie komunisti arvien vairāk izjuta rusofobos (antikrieviskos) noskaņojumus. Poļu komunisti sašķēlās. Daļa enerģiski iebilda pret padomju armiju ieiešanu Polijas etniskajās robežās. V.~Ļeņins tai laikā tādus uzskatīja par ``nacionālistiem'' un ``pacifistiem''. Daļa palika ``internacionālistiskās'' pozīcijās, uzskatot, ka poļu revolucionāriem ir tiesības saņemt palīdzību no ārienes, taču nespēja dot skaidru atbildi uz jautājumu par Polijas valsts nākotni. Viens no PKSP darbiniekiem P.~Levinsons-Lapiņskis 1920.~gada jūlija beigās ziņoja Komunistiskās Internacionāles Izpildu komitejai, ka pēdējā laikā Polijas valdības veiktie komunistu aresti ir noveduši pie partijas darbības pilnīgas apstāšanās. 1.~augustā PKSP CK izdeva uzsaukumu, kurā propagandēja uzskatu, ka tekošais ``karš ir šķirisks, nevis nacionāls'', taču nekas netika teikts par to, ka Padomju Krievija atzīst poļu tautas tiesības uz pašnoteikšanos. No uzsaukuma izrietēja, ka komunisti cerēja, tuvojoties Sarkanajai armijai, sarīkot streikus, sasaukt mītiņus, radīt bruņotas grupas, izraisīt masu uzstāšanās un ``augsta masu revolucionārā pacēluma'' apstākļos noorganizēt strādnieku deputātu padomju vēlēšanas, kā arī formēt poļu Sarkano armiju, kura darbotos kopā ar Krievijas Sarkano armiju un gāztu ``kapitālistisko varu visā Polijā''. Sarkanajai armijai tuvojoties Varšavai, PKSP CK aicināja strādniekus un kareivjus uz bruņotu sacelšanos. Taču faktiski kaut kāda revolucionārās cīņas izvēršanās Polijā tolaik nebija iespējama, ne tikai represiju, bet galvenokārt masu atbalsta trūkuma dēļ. Poļu vairākums pirmkārt domāja par neatkarīgas valsts izcīnīšanu un aizstāvēšanu, jautājumu~--- kāda tā būs, atstājot uz vēlāku laiku.

Kā jau augstāk minēts, 23.~jūlijā Smoļenskā tika izveidota \strong{Polijas Pagaidu revolucionārā komiteja} (\pltxti{Tymczasowy Komitet Rewolucyjny Polski, Polrewkom}) jeb PPRK, kurai vajadzēja uzņemties varas funkcijas pēc Varšavas ieņemšanas un J.~Pilsudska gāšanas. Par to boļševiki paziņoja 1.~augustā Belostokā~--- pirmajā pilsētā aiz Kērzona līnijas, kuras tuvumā arī izvietojās PPRK. Komitejas sastāvā ietilpa KK(b)P biedri J.~Marhļevskis, F.~Dzeržinskis, J.~Unšlihts, F.~Kons, E.~Pruhņaks u.c. PPRK pēc būtības bija Padomju Krievijas radīta marionešu valdība. Visi poļu revolucionāri~--- komitejas locekļi ieņēma ievērojamus amatus padomju republiku partijas un valsts iestādēs. 30.~jūlijā PPRK priekšsēdētājs J.~Marhļevskis nosūtīja telegrammu V.~Ļeņinam, kurā ne tikai ziņoja par komitejas izveidi, bet arī izteica cerību uz drīzu uzvaru karā. Pie tam īpašas cerības tika saistītas ar ``visas pasaules strādnieku šķiras atbrīvotāju~--- varonīgo Sarkano armiju''. Tai pat dienā V.~Ļeņins parakstīja TKP lēmumu par nosūtīšanu F.~Dzeržinskim 1~miljardu rubļu izdalīšanai vietējām revolucionārajām komitejām Sarkanās armijas okupētajā Polijas teritorijā. PPRK izdeva F.~Dzeržinska uzrakstīto ``Aicinājumu visiem pilsētu un lauku darbaļaudīm'', kurā tika pasludināta Polijas Padomju Republikas dibināšana, zemes nacionalizācija, baznīcas atdalīšana no valsts, aicinājums darba cilvēkiem padzīt kapitālistus un muižniekus, ieņemt fabrikas un rūpnīcas, nepakļauties Polijas armijas pavēlēm, veidot revolucionārās kara komitejas kā varas iestādes (tika izveidotas 65~komitejas). Tika izdota avīze, izveidota Ārkārtējā komisija (Čeka) un Revolucionārie tribunāli. Taču tāpat kā XIX~gadsimtā poļu šļahtiči, kuri vēlējās atbrīvoties no Krievijas kundzības, bet ne par ko nevēlējās šķirties no saviem īpašumiem un privilēģijām, arī PPRK nevēlējās muižniekiem atsavināto zemi piešķirt pēc tās alkstošajiem poļu zemniekiem. PPRK visu muižniekiem, valstij, baznīcai un klosteriem piederošo zemi pasludināja par konfiscētu, taču to nevis sadalīja bezzemniekiem un mazzemniekiem, bet nodeva laukstrādnieku komiteju pārvaldē valsts saimniecību organizācijai. Vēl vairāk, PPRK arī zemnieku saimniecības, lielākas par 20~ha, pasludināja par tautas īpašumu. Līdz ar to PPRK neievēroja nedz tajā laikā jau zināmo Ungārijas Padomju republikas, nedz Latvijas Sociālistiskās padomju republikas 1919.~gada negatīvo pieredzi, kas veda pie padomju varas zaudēšanas, nedz Krievijas ar 1917.~gada ``Dekrētu par zemi'' saistīto pozitīvo pieredzi zemnieku piesaistīšanā padomju varai.

15.~augustā tika parakstīta Rietumu frontes pavēlnieka pavēle par 1.~Poļu Sarkanās armijas formēšanu, kurai vajadzētu ietilpt frontes sastāvā. Par šīs armijas pavēlnieku tika iecelts etniskais polis R.~Longva, bijušais cara armijas štābkapteinis, Oktobra apvērsuma dalībnieks Petrogradā, kurš poļu-padomju kara laikā komandēja divīziju, par varonību kaujās pret poļiem Baltkrievijā bija apbalvots ar Sarkanā Karoga ordeni. 17.~augustā, kad jau bija sācies poļu pretuzbrukums pie Varšavas, J.~Unšlihts telegrammā V.~Ļeņinam lūdza izdot pavēli par visu poļu komunistu, kuri dienēja dažādos Sarkanās armijas štābos, daļās un iestādēs, nosūtīšanu 1.~Poļu Sarkanās armijas rīcībā. Bija paredzēts vervēt karavīrus Krievijas iekšienē esošajās gūstekņu nometnēs. Septembrī pat tika izveidoti poļu sarkano komandieru kursi ar vairāk nekā tūkstoti kursantiem. Taču īstu kaujas spējīgu armiju izveidot neizdevās. Tajā nepieteicās pat pāris simtu brīvprātīgo. Sarkanās armijas sakāvju rezultātā poļu armijas formēšana drīz tika apstādināta, pēc tam pilnībā pārtraukta. Tomēr poļu komunistu šūniņas Krievijā vēl septembrī, pat decembrī turpināja mobilizēt savus biedrus nosūtīšanai uz Sarkano armiju.

Polijai, lai izturētu Sarkanās armijas iebrukumu, bija jāveic ārkārtēji pasākumi.

Lai iegūtu poļu zemniecības atbalstu, kura nacionālistiem tik ļoti bija pietrūcis XIX gadsimtā notikušo poļu sacelšanos laikā, 15.~jūlijā uz zemnieku partijas (ļudoviešu) līdera V.~Vitosa priekšlikumu Polijas Seims steidzamības kārtā izskatīja un pieņēma jaunu agrāro likumu, papildinošu 1919.~gada 10.~jūlijā pieņemtos agrārās reformas principus. Atbilstoši jaunajam likumam lielajiem zemes īpašniekiem atsavināmās zemes daudzums palielinājās, vidējas zemnieku saimniecības lielumam vajadzēja sasniegt 15--20~ha. Pirmkārt tiesības uz zemi tika piešķirtas Polijas armijas invalīdiem un karavīriem, bezzemniekiem, mazzemniekiem un laukstrādniekiem. Visiem viņiem valdība solīja piešķirt ilglaicīgu kredītu.

1920.~gada 24.~jūlijā V.~Grabska ``labējā'' valdība iesniedza atlūgumu, tai pat dienā ar J.~Pilsudska rīkojumu sāka darboties jauna daudzpartiju nacionālās aizsardzības valdība ar premjeru V.~Vitosu priekšgalā. Jau šīs valdības laikā 6.~augustā vairākos miljonos eksemplāru tika izplatīts uzsaukums ``Dzimtene briesmās'' ar aicinājumu aizstāvēt dzimteni. Sarkanās armijas darbība tika raksturota kā ``mēģinājums paverdzināt'' Poliju. Tie poļu zemnieki, kuri bija devušies armijā brīvprātīgi vai ilgi karojuši, varēja rēķināties, ka saņems zemi pirmie, īpaši izcēlušies~--- bez maksas, kritušo ģimenēm tika solīta valsts apgāde. (Tik tiešām desmiti tūkstošu kara dalībnieku pēc tā saņēma zemi Rietumukrainā un Rietumbaltkrievijā.) Pavisam brīvprātīgi armijā iestājās ap 106~000 poļu. Brīvprātīgo vidū bija arī daudz ebreju.

Taču nevajag arī pārspīlēt poļu patriotisma uzplūdus. 14.~augustā Polijas armijas ģenerālštāba priekšnieks T.~Rozvadovskis pavēlēja Ziemeļu frontē aiz frontes ierīkot aizsprosta līniju dezertieriem. Tika izvērstas represijas pret PKSP, lielākā daļa tās vadītāju tika ieslodzīta cietumos. (Starp citu, jau 1920.~gada 24.~jūlijā padomju puse lūdza Latvijas ārlietu ministram Z.~A.~Meierovicam paziņot Polijas valdībai, ka Krievija uzskata sagūstītos poļu virsniekus kā ķīlniekus par Polijā arestētajiem komunistiem un piedāvā to savstarpēju apmaiņu.)

Katoļu baznīca, izmantojot poļu reliģiozitāti, pasludināja, ka viņu dzimtenei ir uzbrucis Antikrists Sarkanās armijas veidolā. Varšavā Romas pāvesta legāts, vēlākais pāvests Pijs XI paziņoja, ka ir gatavs pats personiski stāties pretī Antikrista ordām.

Poļu zemnieki kopumā naidīgi sagaidīja Sarkano armiju, pat par naudu nevēlējās dot sarkanarmiešiem pārtiku, slēpa mežos mājlopus, pārtiku, dažkārt pat šāva uz tiem no slēptuvēm. Lai kaut kādi uzturētu kārtību, Sarkanās armijas daļas ņēma ķīlniekus. Tas viss tikai pastiprināja savstarpējo naidu.

Citāda situācija bija agrākajās Vācijai piederējušajās teritorijās, kuras tikai pēc Pirmā pasaules kara bija piešķirtas Polijai. Šeit vietējie vācu iedzīvotāji Sarkanajā armijā saskatīja atbrīvotājus un izturējās pret to ļoti labvēlīgi. Kad Sarkanās armijas daļas ienāca vācu apdzīvotās pilsētās un vietās, vietējie iedzīvotāji tās sagaidīja ar mūziku, karogiem, smaidiem un rokas spiedieniem, cienāja sarkanarmiešus ar mājas ēdienu un alu.

Augusta sākumā \strong{Polijas stāvoklis kļuva kritisks}. Smaga bija valsts ekonomiskā situācija. 1920.~gada pavasarī valstī čuguna ražošana sastādīja tikai 10,2\%, tērauda~--- 11,6\% salīdzinājumā ar 1913.~gada rādītājiem. Kara vešana bija novedusi pie finanšu kraha. Polijas ārējais parāds sasniedza 709,1 miljardus marku (Markas bija ieviesusi Vācija Pirmā pasaules kara laikā pēc Polijas okupācijas). Markas kurss arvien pazeminājās. Bezdarbnieku skaits 1920.~gada aprīlī pārsniedza 350~000 cilvēku. Vācija un Čehoslovākija, ar kurām, tāpat kā ar Padomju Krieviju, Polija bija iesaistījusies konfliktā par robežām, atteicās atļaut caur savu teritoriju tai domāto ieroču un munīcijas tranzītu. Vācu vēsturnieks, publicists un politiķis O.~Hoetšs, kurš Vāckonservatīvo partijā (\detxti{Deutschkonservative Partei}) reihstāgā darbojās kā eksperts Austrumeiropas jautājumos, 27.~jūlijā no reihstāga tribīnes paziņoja: ``Es personīgi mierīgi saku, ka uz šo krievu kara mērķi raugos ar pilnām simpātijām \citespace{} Polija var veidot tiltu tikai tad, kad tiks atjaunotas arī izvaroto vāciešu pašnoteikšanās tiesības Rietumprūsijā un Pozenē \citespace{} līdz tam laikam starp mums un Poliju pastāv nāvīgs ienaids.'' Poļu vēsturnieks J.~Holcers pat norāda, ka izšķirošajā poļu-padomju kara brīdī Vācija devusi ``diskrētu, bet sadzirdamu mājienu par savu gatavību piedalīties jaunā Polijas sadalīšanā''. Tikai Ungārija, kur nesen bija likvidēta padomju vara, piedāvāja Polijai militāru palīdzību. Taču Čehoslovākijas robežas slēgšana neļāva to nosūtīt. Vienīgais preču piegādes ceļš palika Danciga, taču arī tur neiztika bez incidentiem. Ierodoties angļu kuģim ``\entxti{Triton}'' ar kara materiāliem, ostas strādnieki to atteicās izkraut. Viņu vietā darbos tika iesaistīti angļu karavīri. Taču daži desmiti no tiem arī atteicās izkraut kuģi un tika arestēti. Kad par to uzzināja pilsētā, dokeri metās apcietinātos atbrīvot. Strādnieki aizturēja dzelzceļa stacijā poļu komandantu, neļāva doties pa dzelzceļu uz Poliju poļu virsniekiem un karavīriem. Angļu varas iestādes ar karaspēku apspieda nemierus Dancigā, taču vismaz daļa pilsētas iedzīvotāju drīzāk juta līdzi Padomju Krievijai nekā Polijai. Tuvojoties Sarkanajai armijai, no Varšavas steidzīgi sāka evakuēties ārvalstu diplomātiskās misijas.

Jau 1920.~gada 23.~jūlijā sākās Sarkanās armijas Varšavas operācija, kura ilga 33~diennaktis, beidzoties ar sakāvi. Tās sākumā padomju karaspēks atradās apmēram simt, vietām pat astoņdesmit km attālumā no Varšavas. Padomju militāristi bija pārliecināti, ka Varšava tūlīt, tūlīt kritīs. Kad V.~Ļeņins apjautājās partijas CK locekļiem, kuri atradās frontē, un Sarkanās armijas pavēlniecības pārstāvjiem, kādi ir viņu uzskati par kara darbības perspektīvām, Republikas revolucionārā kara padome izteica drošu pārliecību, ka 16.~augustā Sarkanā armija ieies Varšavā. Pamats šādai pārliecībai bija. Pusotra mēneša kaujās, no 1.~jūlija līdz 16.~augustam, poļu armija zaudēja 80~tūkstošu karavīru~--- 56\% no personālā sastāva, kurš atradās frontē. Kara laikā bija ievērojami pieaugušas Polijas armijas aizmugures struktūras, no gandrīz 1~miljona armijas rindās esošo 1920.~gada augustā frontē bija tikai 165~tūkstoši cilvēku. (Pēc Polijas Kara ministrijas datiem pavisam uz fronti tika nosūtīti: sešu mēnešu laikā no 1919.~gada 1.~decembra līdz 1920.~gada 1.~jūnijam 4~080 virsnieku un 260~000~karavīru, septiņās nedēļās no 24.~jūlija līdz 15.~septembrim 1~986~virsnieki un 163~889~karavīri. Vislielākajā spēku sasprindzinājuma brīdī Polijas armijā bija mobilizēti 16~gadu gājumi un armijas kopējais sastāvs sasniedza ap 1~200~000~vīru.) No 6. līdz 15.~augustam, kad pret Varšavu faktiski izvērsās uzbrukums, karojošajās poļu armijas daļās brīvprātīgi iestājās ap 60~tūkstošu cilvēku. Padomju vadībā to neviens negaidīja. Kā 1920.~gada 4.~augustā G.~Ordžonikidzem rakstīja K.~Vorošilovs: ``\dots{} no poļu strādniekiem mēs gaidījām sacelšanās un revolūcijas, bet saņēmām šovinismu un trulu naidu pret ``krieviem''.

Polijas pusē dažādu iemeslu dēļ karoja arī daudzi krievi, ukraiņi, baltkrievi, nemaz nerunājot par citu valstu pārstāvjiem. Piemēram, Polijas armijā 1920.~gada jūnijā bija 117~lidmašīnas, no tām kaujas spējīgas~--- 75.~Pilotu vidū bija 54~francūži, 17~amerikāņi un 3~angļi. Tā T.~Kostjuško vārdā nosauktā eskadriļa, kura karoja pret 1.~Jātnieku armiju, sastāvēja no amerikāņu lidotājiem. (Salīdzinājumam var minēt, ka Sarkanās armijas Rietumu frontes rīcībā tai laikā bija 220~lidmašīnu, tai skaitā 84~kaujas nespējīgas).

Augusta vidū ``sarkano'' kavalērija sasniedza Vislas līkumu pie Toruņas, kas atradās tikai piecu dienu pārgājiena attālumā no Berlīnes. Boļševiki cerēja nonākt līdz Vācijas galvaspilsētai, kur pēc viņu aprēķiniem varētu sākties ar Versaļas līgumu neapmierināto vācu strādnieku sacelšanās. M.~Tuhačevska plāns paredzēja forsēt Vislu un apiet Varšavu no rietumiem. Pēc dažiem pieņēmumiem Sarkanajai armijai bija jānovirzās uz ziemeļiem, lai dziļi ieietu Poļu koridorā, ātrāk sasniegtu Vācijas robežu un paātrinātu revolūcijas sākumu tajā. Taču ar to tika novājināts uzbrūkošo kreisais flangs. Padomju armijas bija tālu atrāvušās no savām apgādes daļām, tām trūka munīcijas un pārtikas. Jau sākot uzbrukumu bija skaidrs, ka fronte var tikt apgādāta ne vairāk kā divus kauju mēnešus. Jau augustā sarkanarmieši badojās, tiem izsniedza pusi, bet dažkārt arī tikai trešdaļu no paredzētās pārtikas devas. Karavīri ēda visdažādākos zaļumus, kas noveda pie dizentērijas epidēmijas uzliesmošanas. Morālais līmenis bija kritis. Nonāca pat līdz dumpjiem.

Izvēršoties Sarkanās armijas uzbrukumam, Polijas Sociālistiskā Partija, kura 1918.~gadā ne tikai piedalījās neatkarīgās Polijas atjaunošanā, bet vēlāk nebija iebildusi arī pret Polijas karaspēka iebrukumu lietuviešu, ukraiņu un baltkrievu apgabalos, iesaistījās V.~Vitosa vadītajā koalīcijas valdībā. Tas gan noveda pie partijas šķelšanās. No tās nodalījās pret karu noskaņotā biedru daļa ar S.~Lancucki un J.~Sohacki priekšgalā. Opozīcijas vadītāji radīja jaunu PPS-ļevicu (\pltxti{Polska Partia Socjalistyczna~--- ļevica}), kura 1921.~gadā iestājās PKSP (\pltxti{Komunistyczna Partia Robotnicza Polski}).

Polijas Sociālistiskās partijas nostādnes daļēji noteica principiāli, daļēji~--- konjunktūras rakstura apsvērumi. PPS programmas projektā, kuru partijas Programmas komisija piedāvāja tās XVII~kongresam bija teikts: ``Represiju un pat terora līdzekļu pārvēršana par veselu sistēmu [Tā bija nepārprotama norāde uz Padomju Krievijā nodibināto boļševiku varas sistēmu.~--- V.~Š.] \citespace{} nav savienojama ar sociālisma raksturu un nevar novest pie strādnieku šķiras atbrīvošanas, tāpēc PPS noraida tā saucamo un pielietojamo ``proletariāta diktatūru'', tās vietā izvirzot sociālistisku valdību, kuras balstītos uz pilsētu un lauku darba masām, atbilstu sabiedrības vairākuma gribai, būtu visu pilsoņu kontrolētas, prasību.'' Taču Sarkanajai armijai tuvojoties Varšavai, PPS Centrālā Izpildkomiteja griezās pie tautas ar uzsaukumu, kurā pēkšņi ``atcerējās'': ``Bēdas uzvarētajiem! [no latīņu ``\latxti{Vae victis}!''~--- ``Bēdas uzvarētajiem!''] ir lozungs, kas jau kopš pasaules sākuma ierakstīts visu tautu visu uzvarētāju valstu karogos''. Saturot tēzēs par to, ka ``šobrīd pie varas Polijā ir reakcija''; ka ``karš~--- tā ir slepkavošana, tvans, uguns un bads!'', uzsaukums tomēr ar it kā ``sociālistiskiem'' argumentiem, kuri faktiski viegli pārtapa nacionālistiskos, aicināja tautu cīņā par neatkarību. Lūk, daži citāti: ``Šķiru cīņu ar poļu buržuāziju līdz galam jāveic poļu strādnieku šķirai tikai ar pašas spēkiem. Tā uzvarēs, ja būs stiprāka par buržuāziju, zaudēs, ja būs par to vājāka. Svešais Sarkanās armijas iebrukums ir nebeidzama badā novārdzinātu cilvēku straume, kas grib pieēsties, tā ir jauna [Polijas] sadalīšana. Neatkarība ir apdraudēta. To vajag aizstāvēt. Vajag visu strādnieku šķiras spēku sasprindzinājumu, lai iebrukumu atsistu, lai izvairītos no sadalīšanas.'' \citespace{} ``Vispārējās mobilizācijas ceļā veidots karaspēks, nevar būt reakcionārs. Tie taču ir zemnieku un strādnieku dēli, kauls no poļu tautas kaula, asinis no tautas asinīm.'' \citespace{} ``Mūsu buržuāzija lieliski izprot situāciju, skaidri to saredz pašreizējā valdība, un gribot ilgāk palikt pie varas, vēlas noslēgt mieru tā vietā, lai balstītos uz tautas spēku. Vēl jo vairāk tāpēc, ka miers ar Krieviju ir neizmērojams tirgus ar visiem Austrumiem, liels milzīgu, pastāvīgu ienākumu avots poļu kapitālismam.''

Tātad poļu sociālisti poļu zemniekus un strādniekus it kā aicināja cīņā gan pret savu buržuāziju, gan pret savas zemes izlaupīšanu no ``nebeidzamas badā novārdzinātu cilvēku straumes, kas grib pieēsties'', gan pret Polijas sadalīšanu. Grūti spriest, cik lielā mērā tāds pretrunīgs aicinājums ietekmēja sociālisma piekritējus valstī, taču jākonstatē, ka kaut kādu ievērojamu šķiru cīņas uzliesmojumu Polijas aizmugurē nebija, padomju komunistu cerības piesaistīt sev poļu strādnieku un trūcīgo zemnieku simpātijas neattaisnojās.

Vēl kāda piebilde. PPS uzsaukumā bija teikts: ``Tātad miers! Taču miers taisnīgs, bez dalīšanām, bez vardarbības, kurā slēpjas nākamo karu aizmetņi. Ne Lietuva, ne Baltkrievija, ne Ukraina nav krievu zemes. Šo zemju tautām ir tiesības lemt savu likteni bez krievu vai poļu spiediena.'' Diemžēl šo grūtā brīdī apjausto patiesību pēc uzvaras pār Sarkano armiju aizmirsa ne tikai pilsonības, bet arī PPS vadītāji un lietuviešu, baltkrievu, ukraiņu tautu liktenis tika izlemts ar poļu militāru spiedienu.

Ievērojamu militāro palīdzību Polijai sniedza Francija. Franču militāro misiju, kuras uzdevums bija palīdzēt strādāt Polijas Ģenerālštābam, vadīja ģenerālis M.~Veigans. Francijas palīdzības loma Sarkanās armijas sakāvē pie Varšavas vēlāk padomju literatūrā tika pārspīlēta, taču tās misijas sastāvā ietilpa 9~ģenerāļi, 29~pulkveži, 196~kapteiņi, 425~leitnanti, vairāki simti karavīru un instruktoru. Franču virsnieku vidū bija arī vēlākais Francijas nacionālais varonis un prezidents Š.~de~Golls. Kā raksta N.~Deiviss, Š.~de~Gollu 1920.~gada kampaņa ar plašiem karaspēka manevriem iedvesmoja formulēt savas idejas par moderno kara darbību.

Padomju un arī daļa mūsdienu krievu vēsturnieku uzsver, ka padomju-poļu karš norisa vienlaikus ar Antantes valstu intervenci Krievijā un tāpēc lielvalstis to uzskatīja par cīņas pret boļševistisko režīmu sastāvdaļu. Antantes valstis sniedza Polijai ekonomisku, militāru un diplomātisku atbalstu. Padomju vēsturnieki apgalvoja, ka bez Rietumu atbalsta Polija nespētu sekmīgi karot. Tomēr nākas atzīt, ka Rietumvalstu viedoklis par nepieciešamību sniegt Polijai palīdzību karā bija visai atšķirīgs un šo palīdzību, kaut nozīmīgu, neklājas pārvērtēt. (Domājams, bez šīs palīdzības Polijas un Padomju Krievijas/PSRS robežas būtu citādas, bet Polijas valstiskums saglabātos).

ASV un Francija uzskatīja par nepieciešamu visiem spēkiem atbalstīt J.~Pilsudska kursu, aktīvi piedalījās Polijas armijas apgādē. Tikai 1919.~gadā no ASV noliktavām Eiropā, kuras pēc Pirmā pasaules kara bija izrādījušās nevajadzīgas, Polija saņēma militāro aprīkojumu par 60~miljoniem, bet 1920.~gadā~--- jau par 100~miljoniem dolāru.

Turpretī Lielbritānija iestājās par kompromisu, tiecās ierobežot palīdzību Polijai, bet pēc tam pat sliecās uz neitralitāti tās konfliktā ar Padomju Krieviju, faktiski pārtrauca sniegt Polijai militāru un ekonomisku palīdzību. 1920.~gada augustā leiboristu un arodbiedrību konference nobalsoja par vispārējā streika sākšanu gadījumā, ja valdība nolemtu arī turpmāk atbalstīt Poliju. Munīcijas iekraušana Polijai turpmāk tika sabotēta. Domstarpības Antantes valstu vidū par palīdzību Polijai 1920.~gadā veicināja šo valstu atbalsta pārtraukšanu ``balto'' kustībai Krievijā un vēlāko Padomju Savienības starptautisko atzīšanu.

Malā no notikumiem nevarēja palikt arī Latvija. Kara darbības atjaunošana poļu-padomju frontē iespaidoja Latvijas un Padomju Krievijas miera sarunas. Vēl pirms poļu uzbrukuma uzsākšanas 1920.~gada 6.~aprīlī laikraksta ``Republikas Sargs'' darbiniekiem Polijas sūtnis Latvijā paziņoja, ka, ja Latvija noslēgs mieru ar Krieviju, bet Polija turpinās karu, Polija savu karaspēku no Latgales neizvedīs, jo Daugavpils būs nepieciešama Polijai par atbalsta punktu tās armijas kreisajam spārnam. Tādejādi Polija nekautrējās pielietot draudus, mēģinot atturēt kaimiņvalsti no miera līguma slēgšanas ar Padomju Krieviju.

Turpretī Lielbritānijas valdība ar D.~Loidu Džordžu priekšgalā atbalstīja Baltijas valstu centienus noslēgt mieru ar Padomju Krieviju un arī pati no 1920.~gada aprīļa Kopenhāgenā veda ar to sarunas ekonomiskajos jautājumos. ``\entxti{Times}'' 1920.~gada 3.~jūlijā rakstīja, ka britu valdība paziņojusi Polijas valdībai savu vēlēšanos, lai Polija atzītu Lietuvas, ieskaitot Viļņas, neatkarību, kā arī lai šī valsts iesaistītos sarunās ar Lietuvu. Mērķis~--- nodibināt draudzīgas attiecības starp abām valstīm. Teritoriālajos jautājumos Lielbritānija atbalstīja Latviju un Lietuvu. Latvijas laikrakstos sāka atklāti parādīties priekšlikumi noslēgt Baltijas valstu savienību bez Polijas. Sociāldemokrāts F.~Menders Satversmes sapulcē 1920.~gada 17.~jūnijā izteica spriedumu, ka Latvijas savienība ar Poliju ``var dot vislielāko politisko atkarību no Polijas, kura mums saimnieciski neko nevar dot vai var dot ļoti maz''. Neraugoties uz to, Latvijas valdība, rūpējoties par savām interesēm un neņemot vērā bijušās sabiedrotās Polijas intereses, kura 1920.~gada janvārī bija devusi galveno ieguldījumu Latgales atbrīvošanā no Sarkanās armijas, 11.~augustā, laikā, kad Sarkanā armija atradās pie Varšavas, noslēdza mieru ar Padomju Krieviju.

Poļu pretuzbrukuma plāns paredzēja spēcīgu triecienu no dienvidaustrumiem Rietumu frontes aizmugurē. Uzbrukuma plānošanā un realizācijā ļoti palīdzēja poļu radioizlūkošanas un atšifrēšanas dienesta darbinieku sasniegumi. Jau 1919.~gada augustā~--- septembrī Polijas radioizlūkošanas šifrētāju nodaļa poručika (virsleitnanta) J.~Kovaļevska vadībā ``uzlauza'' gandrīz visus kā ``balto'' t.s. Brīvprātīgo armijas, tā arī Sarkanās armijas šifrus, bet 1920.~gada janvārī arī Vācijas armijas šifrus. Īpaši rezultatīvi nodaļa strādāja 1920.~gada vasarā, kad tika dešifrētas ap 500~radiotelegrammas mēnesī, kuras bija parakstījuši padomju karavadoņi. Tieši kaujas par Varšavu laikā bija atšifrēti 420~padomju komandieru ziņojumi un pavēles, tai skaitā arī M.~Tuhačevska direktīvas karaspēka daļām, kuras viņš izdeva Minskā, simtiem km attālumā no frontes līnijas. Uz to pamata J.~Pilsudskis un viņa ģenerāļi varēja pieņemt pareizus lēmumus, kas ievērojami atviegloja ceļu uz uzvaru.

Starp citu, arī ``sarkano'' rokās 13.~augustā nonāca poļu pavēle par trieciena sagatavošanu, taču padomju pavēlniecība to uzskatīja par dezinformāciju ar nolūku aizkavēt uzbrukumu Varšavai. Drīz gadījuma pēc izdevās pārtvert arī otru pavēli, kura attīstīja pirmo. Un atkal attiecīgi secinājumi netika izdarīti. Vēlāk bijušais Dienvidrietumu frontes pavēlnieks A.~Jegorovs pamatīgi kritizēja Rietumu frontes pavēlnieku M.~Tuhačevski par neprasmi pārorientēties atbilstoši kaujas apstākļiem.

Kā uzskata krievu vēsturnieks M.~Meļķuhovs, poļus izglāba Sarkanās armijas virspavēlnieka S.~Kameņeva kļūda, kurš pārāk optimistiski novērtējis uzvaras iespējas un tā vietā lai dotu vienu triecienu Varšavas virzienā pieļāvis Dienvidrietumu frontes virzību uz Ļvovu, tādejādi vēršot divus triecienus Varšavas un Ļvovas virzienā. Tiesa, ievērojot kritisko situāciju, S.~Kameņevs pavēlēja pārsviest no Dienvidrietumu frontes divas armijas (tai skaitā arī Pirmo Jātnieku armiju) uz Rietumu fronti, taču Dienvidrietumu frontes pavēlniecība, pie kuras piederēja arī J.~Staļins, vēloties ieņemt Ļvovu, to savlaicīgi neizpildīja. Pēc vairāku vēsturnieku domām tas lielā mērā ietekmēja notikumu attīstību par ļaunu Sarkanajai armijai. Par to, kam no militārajiem Sarkanās armijas vadītājiem bija jāuzņemas atbildība par zaudējumu, vēsturnieku strīdi jau ir ilgstoši un diezin vai kad tiks noslēgti.

Vēl var atzīmēt, ka Padomju Krievijai Pilsoņu kara laikā arvien no jauna atradās kāda ``šaurā vieta'', kas nelaikā ``plīsa'' un bija ``jālāpa''. Kaujas pie Varšavas bija saistītas ar lielām grūtībām Sarkanajai armijai. Jau augusta sākumā V.~Ļeņins telegrafēja J.~Staļinam, ka milzu apmērā pieaugušas Vrangeļa briesmas, ka KK(b)P CK aug noskaņojums tūlīt noslēgt mieru ar Poliju. KK(b)P CK Politbirojs 2.~augustā nolēma Dienvidrietumu frontes daļu pie Krimas izdalīt atsevišķā Dienvidu frontē cīņai pret P.~Vrangeļa komandētajiem ``baltajiem''. To noteica ne tikai pieaugušās briesmas no ģenerāļa P.~Vrangeļa uzbrukuma, bet arī kazaku sacelšanās Kubaņā. Sekoja pavēles par cilvēku, kara materiālu, citu resursu nosūtīšanu jaunajai frontei, atsakot Rietumu frontei. Vēl tagad krievu literatūrā var atrast viedokli, ka kauju pie Varšavas uzvarēja ģenerālis P.~Vrangelis, bet ne ``baltpoļi''.

Savas bažas par stāvokli poļu-padomju frontē izrādīja arī Rietumi. Francijas maršals F.~Fošs, kurš par militārajiem nopelniem Pirmajā pasaules karā saņēma arī Lielbritānijas feldmaršala (1919) un Polijas maršala (1923) pakāpes, noteikti ieteica J.~Pilsudskim atstāt savās rokās tikai valsts priekšnieka amatu, bet virspavēlniecību nodot kādam poļu ģenerālim, kurš tad rīkotos saskaņā ar Antantes padomiem. Līdzīgu prasību izvirzīja franču-britu konference 1920.~gada 8.--10.~augustā.

J.~Pilsudskis 12.~augustā (citi dati~--- jau 6.~augustā) nodeva premjerministram V.~Vitosam iesniegumu, kurā viņš rakstīja par savu atbildību par katastrofālo militāro stāvokli un savu atkāpšanos gan no virspavēlnieka, gan Valsts priekšnieka amata. Kad darīt to zināmu sabiedrībai, bija jāizlemj premjeram. Viņš savukārt uzskatīja par saprātīgāku Polijai tik kritiskā laikā neinformēt sabiedrību par J.~Pilsudska iesniegumu. Kad karā notika pavērsiens, tika gūta uzvara, jautājums par valsts priekšnieka un virspavēlnieka demisiju atkrita automātiski. J.~Pilsudska lūgums tā arī nekļuva zināms tautai.

Kaut 13.~augustā KK(b)P CK Politbirojs pieņēma lēmumu no blakus Rietumu frontei pret Poliju virzītās Dienvidrietumu frontes sastāva izvest kaujas spējīgākās daļas, lai nostiprinātu pret ģenerāļa P.~Vrangeļa komandētajiem ``baltajiem'' veidojamo Dienvidu fronti, tai pat dienā sākās Sarkanās armijas Rietumu frontes uzbrukums pie Radzminas (\pltxti{Radzymin}) pilsētas 23 km no Varšavas. Poļu daļas atkāpās. Ģenerālštāba priekšnieks T.~Rozvadovskis pavēlēja atbruņot 11.~divīzijas daļas, kuras bija atstājušas Radzminu un nosūtīt uz soda rotām to karavīrus. Taču 15.~augustā sākās jau poļu spēku uzbrukums pie Radzminas, tā drīz tika ieņemta. J.~Pilsudskis šai laikā notikumu vadībā nepiedalījās, savā auto atradās tūlīt aiz uzbrūkošajām daļām. Kaut vēl joprojām Sarkanajai armijai bija dzīvā spēka pārsvars proporcijā 1,2:1, poļi izveidoja ievērojamu pārspēku trieciena virzienos, sekmīgi uzbruka un radīja visu Varšavas rajonā esošo Sarkanās armijas spēku aplenkuma draudus. Tie kaujā pie Vislas 13.--15.~augustā cieta smagu sakāvi, uzsāka haotisku atkāpšanos. 18.~augustā J.~Pilsudskis kā Valsts priekšnieks griezās pie sabiedrības ar aicinājumu neļaut no poļu zemes aiziet nevienam dzīvam sarkanarmietim. Pusotras nedēļas laikā etniskās poļu teritorijas bija atbrīvotas no pretinieka.

Tikai 20.~augustā, kad Rietumu frontes karaspēks jau nekārtībā atkāpās uz austrumiem, Pirmā Jātnieku armija sāka virzīties uz ziemeļiem. Šai laikā 19.~augustā poļi ieņēma Brestu, 23.~augustā Belostoku. Tai pat laikā pēc dažādiem datiem no 40 līdz pat 100~tūkstošiem sarkanarmiešu pārgāja Vācijas robežu un tika internēti. (Pēc mūsdienās veiktiem arhīvu pētījumiem Sarkanās armijas Mobilizācijas pārvalde par internētiem uzskatīja 43~tūkstošus sarkanarmiešu un komandieru.) Pēc vērtējumiem Sarkanā armija pie Varšavas zaudēja ap 25~tūkstošiem kritušo, 16~tūkstošu gūstā nonākušo, vairāki tūkstoši karavīru pazuda bez vēsts. (Literatūrā sastopami arī dati, ka gūstā krituši no 60 līdz 95~tūkstošiem sarkanarmiešu.) Tika zaudēts arī daudz artilērijas un citas tehnikas. Poļu zaudējumi sastādīja ap 15~tūkstošu nogalināto un bez vēsts pazudušo un 22 tūkstošus ievainoto.

Krievu vēsturnieks un filozofs A.~Burovskis, rakstot par Padomju Krievijas mēģinājumiem XX~gadsimta 20.~gados nest ``sociālismu'' citām valstīm, atzinis, ka ``nekāda Padomju Krievijas vai PSRS agresija nevarēja nodibināt padomju varu bez sagrābto valstu iedzīvotāju palīdzības. Polija nekļuva padomiska ne tāpēc, ka tā [bija] stipra un liela, bet tāpēc, ka pret padomju agresiju cēlās pat poļu komunisti un anarhisti'' [Pēdējie vārdi par poļu komunistu cīņu pret padomju agresiju gan ir pārspīlējums~--- V.Š.] Galvenais Sarkanās armijas sakāves cēlonis pie Varšavas bija klaja poļu nācijas apņēmības nosargāt savu neatkarību, tās kaujas gatavības nenovērtēšana laikā, kad Sarkanās armijas daļas turpināja uzbrukt tikai inerces pēc, zaudējušas tam stimulu. Protams, bija kļūdas Sarkanās armijas stratēģiskajā vadībā, taču tās tikai pastiprināja galvenā cēloņa iedarbību.

Jau XX gadsimta 20.~gados sākās poļu vēsturnieku strīdi par to, kas bija uzvaras plāna autors. Pēc vienas versijas tas bija virspavēlnieks J.~Pilsudskis. (Viņš diezgan sīki aprakstīja tā tapšanu savā darbā ``\pltxti{1920.~rok}'' (``1920.~gads'').) Oficiālā propaganda ar entuziasmu uzvarētāja harizmu piešķīra J.~Pilsudskim. Turpretī nacionāldemokrāti izplatīja versiju, ka idejas autors bija franču ģenerālis M.~Veigans, kurš no 25.~jūlija vadīja Antantes militāri-diplomātisko misiju Polijā un bija nozīmēts par tās Ģenerālštāba padomnieku. Trešie par plāna autoru uzskatīja Ģenerālštāba priekšnieku (\pltxti{Szef Sztabu Generalnego Wojska Polskiego}) ģenerāli T.~Rozvadovski. (1920.~gada 27.augustā Valsts aizsardzības padomē J.~Pilsudskis paziņoja, ka M.~Veigana plāns paredzējis poļu atkāpšanos līdz Sanas upei un tāpēc viņš to noraidījis, bet T.~Rozvodovska plāns paredzējis spēku koncentrāciju pie Varšavas un tāpēc ar valsts priekšnieka un virspavēlnieka labojumiem ticis pieņemts.) Katras versijas piekritēji to balstīja ar saviem argumentiem. Jau mūsdienās J.~Pilsudska biogrāfijas un citu darbu Polijas vēsturē autors G.~Matvejevs uzsvēris, ka strīds bija bez īsta pamata, jo virspavēlniekam jebkurā gadījumā piederēja izšķirošais vārds, pārējiem bija tikai padomdevēja balss. Taču praktiski uzbrukuma operāciju sākumā Polijas armiju vadīja ģenerālis T.~Rozvadovskis. Daudzi poļu vēsturnieki tieši šim profesionālim piešķir nozīmīgāko vietu uzvaras izplānošanā un īstenošanā. Lielu ieguldījumu panākumu kaldināšanā deva ģenerāļa V.~Sikorska komandētā armija, kura devās triecienā divu ``sarkano'' armiju saskarsmes vietā. Pēc M.~Veigana domām bez V.~Sikorska, tā komandētās armijas ``\dots{}kauja netiktu uzvarēta''. Tā laika valdības vadītājs V.~Vitoss vēlāk ar rūgtumu konstatēja, ka tādi uzvaras autori kā T.~Rozvadovskis, V.~Sikorskis, J.~Hallers tika aizmirsti.

Atšķirībā no pirms Otrā pasaules kara tapušajiem poļu vēsturnieku darbiem mūsdienu poļu historiogrāfijā ir nostiprinājies uzskats par J.~Pilsudska izšķirošo lomu uzvaras plāna sagatavošanā. Uzreiz pēc uzvaras pie Varšavas J.~Pilsudskis demonstrēja pret ģenerāli M.~Veiganu tik augstprātīgu izturēšanos, ka viņš apvainojās un atstāja Poliju jau 24.~augustā.

Kauja pie Varšavas tika nosaukta par \strong{``brīnumu pie Vislas''} (\pltxti{Cud nad Wisla}). Interesanti, ka nosaukuma autors bija J.~Pilsudska nesamierināmais oponents S.~Stronskis, kurš 1920.~gada 14.~augustā avīzē ``\pltxti{Rzeczpospolita}'' ievietoja rakstu ar nosaukumu ``Ak, brīnums pie Vislas'', kurā izteica domu, ka ar tādu virspavēlnieku tikai brīnums spēs glābt Poliju. J.~Pilsudska apkārtne centās šo dūrienu neievērot. Taču ``brīnums'' nebija vajadzīgs, Polijas armija uzvarēja arī bez tā. Tomēr pēc tam nosaukums tika izmantots propagandas nolūkos, nu jau ar pretēju ievirzi virspavēlnieka vērtējumam. Kaujas sākuma gadadienā~--- 15.~augustā Polijā sāka svinēt Polijas armijas svētkus. (Kopš 1990.~gada 15.~augusts atkal ir Polijas nacionālie svētki.) Poļu skolu mācību grāmatās tā tika pielīdzināta Grīnvaldes (1410.~gadā apvienotie poļu-leišu-krievu spēki šeit sakāva Teitoņu ordeni un apturēja tā virzību uz austrumiem) un Vīnes kaujai (1683.~gadā poļu--austriešu--vācu spēki sakāva turkus, kauja kļuva par lūzuma momentu trīs gadsimtus ilgušajā Centrālās Eiropas valstu cīņā pret Osmaņu impēriju). Pēc aculiecinieku stāstiem 2010.~gadā ``brīnuma pie Vislas'' gadadiena Polijā tika svinēta ar lielāku vērienu nekā uzvaras par nacistisko Vāciju 65.~gadadiena.

Lielbritānijas speciālais sūtnis E.~D’Aberons 1920.~gadā Varšavā izdeva grāmatu ``Varšavas kauja~--- astoņpadsmitā nozīmīgākā pasaules vēsturē''(``\pltxti{Osiemnasta decydująca bitwa ш dziejach świata под Warszawa}''). Poļu mācību grāmatās uzsvērts, ka kaujā pie Varšavas esot izglābta ne tikai Polija, bet arī visa Eiropa no boļševiku iebrukuma. Šādi apgalvojumi sastopami ne tikai presē (populārā poļu avīze ``\pltxti{Gazeta Wyborcza}''. (``Vēlēšanu avīze'') 2009.~gada 22.~augustā rakstīja ``Pie Varšavas tika izglābta Polija, pus-Eiropas, iespējams, arī pasaule.''), bet arī nopietnu vēsturnieku darbos. Piemēram, Varšavas universitātes profesors T.~Nalenčs, tēlojot situāciju, kāda varētu rasties Sarkanās armijas uzvaras rezultātā, zīmē apokaliptiskas ainas, kad pāri ``Polijas līķim'' Sarkanā armija ielauztos Vācijā, saņemot ievērojamu palīdzību no vācu komunistiem, ar krāpšanu un teroru tur atrastu ``kopēju valodu'' ar miljoniem vāciešu, pievarētu Pasaules karā izmocīto Franciju, kurā daudzi franči bijuši kreisi radikāli noskaņoti, un, iespējams, tikai Lamanšs paglābtu Lielbritāniju no iebrukuma. Pēc nosauktā vēsturnieka domām Otro pasaules karu tad izraisītu nevis nacisms, bet ``agresīvais komunisms'', un Mao Dzeduns nevis 1949.~gadā, bet krietni agrāk pasludinātu Ķīnas Tautas Republiku. Polija ar uzvaru pie Vislas ietekmējusi ``visas pasaules likteņus, kura neapjaušot stāvokļa bīstamību, to gandrīz nepamanīja''. Minētais autors šeit, nerūpējoties par paša uzburtās iespējamības atbilstību vēsturiskajai situācijai, iedomāto uztver kā neapšaubāmu īstenību.

Toties ievērojamais poļu vēsturnieks A.~Čubinskis rakstīja, ka ``krievu revolucionāriem 1920.~gadā nebija Eiropas iekarošanas plāna. Viņi bija pārāk vāji. Viņi varēja palīdzēt revolucionāriem centieniem blakus esošajās teritorijās, taču nevarēja pacelt revolucionāru vilni Eiropas mērogā.'' Arī šī darba autors uzskata, ka, ja Sarkanā armija pat ieņemtu Varšavu, Polijā boļševiku varu sagaidītu tāds pats liktenis, kādu tā 1919.~gadā piedzīvoja Latvijā, kad latviešu lielinieki ar utopisko saimniecisko politiku un represijām nostādīja pret sevi lielāko tautas daļu. Pasaules revolūcijas plāni nekādi nevarētu piepildīties. Boļševikiem pat vislabvēlīgākajā situācijā neizdotos nopietni apdraudēt Vāciju, kurā valdīja briesmīgs nogurums no kara, un kura jau bija pārdzīvojusi savu komunistu mēģinājumu saņemt varu. Arī Rietumvalstis bija pietiekami stipras, lai Krievijas Pilsoņu karā novārgušo Sarkano armiju nelaistu tālāk uz rietumiem. Kad jau pēc poļu-padomju kara un arī Pilsoņu kara beigām Krievijā toreizējie boļševiku vadoņi ar Kominternes un daļēji arī PSRS militārā izlūkdienesta spēkiem mēģināja ``pagrūst'' revolūciju Bulgārijā (1923), Vācijā (1923), Igaunijā (1924), Ķīnā (1927), tie visi cieta neveiksmi ne tikai nosaukto valstu iekšējo apstākļu, bet arī PSRS vājuma rezultātā. Jautājums par ``boļševisma briesmu'' realitāti tai vai citai valstij, tai vai citā brīdī ir diskutabls.

Runājot par Polijas kā ``Eiropas glābēja'' no boļševisma lomu, jāmin arī mūsdienu ukraiņu vēsturnieka S.~Lozuņko viedoklis. Viņš uzskata, ka Polijas naidīgo attieksmi pret Padomju Krieviju noteica nevis ``antisovjetisms'' un ``antiboļševisms'', bet gan tajā izplatītie ``vēsturiskie kompleksi'', gadsimtos izveidojies naidīgums pret Krieviju un tālejošie lielvalstiskie Polijas ``no jūras līdz jūrai'' atjaunošanas plāni 1772.~gada robežās. Politiskajam režīmam, kāds Krievijā pastāvēja, Polijas valdītāju acīs bija otršķirīga loma. Ja Krievijā tai laikā būtu cita vara, Polija tik un tā uzsāktu ar to konfliktu, pasludinot sevi par ``bastionu cīņā pret krievu barbarismu''. ``Cīņas pret boļševismu'' aizsegs Polijai bija vajadzīgs, lai attaisnotu savus agresijas aktus un nāciju pašnoteikšanās tiesību pārkāpumus, kā arī lai saņemtu no Antantes valstīm militāru palīdzību. Šī darba autoram jāsaka, ka fakti: J.~Pilsudska atteikšanās palīdzēt ``baltajiem'' cīņā pret ``sarkanajiem'' 1919.~gadā, boļševiku un ``krievu barbaru'' tēlu sapludināšana poļu propagandā, cīņa ne tikai pret boļševikiem, bet arī kaimiņnāciju nacionālistiem (ja tie nepadevās Polijas vadībai) liecina par to, ka ukraiņu vēsturniekam lielā mērā ir taisnība. To netieši apstiprina arī tas, ka jau mūsdienās, gadu desmitus pēc padomju varas sabrukuma, Polijas republikas attiecības ar Krievijas Federāciju nav kļuvušas daudz draudzīgākas.

Taču, runājot par Varšavas kauju, noteikti jāatzīst, ja ne poļu cīņas spars, aizstāvot dzimteni, karš Polijai nodarītu vēl daudz ļauna. Diemžēl, J.~Pilsudskis pēc uzvaras tajā nerūpējās par kara ātrāku pārtraukšanu, bet gan par lielāku teritoriju iekarošanu, kas prasīja jaunus upurus abās karojošajās pusēs.

Augusta beigās Pirmā Jātnieku armija devās triecienā Ļubļinas virzienā. Poļi pēc radio ziņu atšifrēšanas bija jau tai pretī izvirzījuši rezerves, cerot to ielenkt un iznīcināt. 31.~augustā notika lielākā kopš 1813.~gada kavalērijas kauja, kurā cīņās par Ļvovu un garajā pārgājienā nogurusī 1.~Jātnieku armija gan tika atsviesta, taču ielenkt un iznīcināt poļiem to neizdevās. Līdzi tai arī pārējās Dienvidrietumu frontes armijas bija spiestas atstāt Ļvovu un pāriet aizsardzībā. Sarkanajai armijai zaudējot poļu apdzīvotās teritorijas, nācās pārtraukt savu darbību Polijas pagaidu revolucionārajai komitejai.

No Sarkanās armijas atbrīvotajās teritorijās Polijas administrācija krasi pastiprināja represīvo darbību. Cietums vai pat nāves sods pēc ārkārtējās kara tiesas sprieduma draudēja visiem, kuri kaut kādi bija palīdzējuši Sarkanajai armijai un poļu revolucionārajām komitejām. Tur, kur pēdējās bija darbojušās, izplatību guva arī nekontrolētas linča tiesas. Plaši bija izplatīti miesas sodi. Pērti tika galvenokārt maz- un bezzemnieki, kuri bija sveikuši Sarkanās armijas atnākšanu, piedalījušies laukstrādnieku revolucionāro komiteju darbībā. Tā 1920.~gada vasara un rudens daudziem kļuva par asiņainu ``patriotisma'' mācību stundu. Plockā (\pltxti{Płock}) bija gadījums, kad ostraksismam (izraidīšanai no sabiedrības) tika pakļauti vecāko klašu skolnieki, kuri nebija pieteikušies kā brīvprātīgie Polijas armijā. Jau augusta beigās poļu Muižnieku savienības galvenā valde nosūtīja valdībai memorandu, kurā izteica neapmierinātību ar tikko pirms dažām nedēļām pieņemto agrāro reformu un uzstādīja virkni prasību. 13.--15.~oktobrī šīs savienības visas valsts sapulce oficiāli nosodīja reformu un pieprasīja atņemt zemniekiem jau piešķirto zemi. Savukārt Laukstrādnieku arodbiedrības Galvenā valde septembra sākumā nosūtīja vēstuli Polijas Ministru padomei ar lūgumu novērst muižnieku atriebību laukstrādniekiem teritorijās, kuras poļu karaspēks bija atkarojis Sarkanajai armijai.

Tai pat laikā no 17.~augusta Minskā turpinājās pēc Polijas puses priekšlikuma uzsāktās miera sarunas. J.~K.~Daniševskis tajās vēl uzstādīja prasības samazināt Polijas bruņotos spēkus līdz 50~000 vīru, izveidot strādnieku miliciju, atgriezts poļu iekarotāju atņemto īpašumu baltkrievu un ukraiņu iedzīvotājiem u.c., kuras bija izvirzītas Sarkanajai armijai uzbrūkot. Arī padomju delegācijas deklarācija saturēja šīs prasības, kuras bija aktuālas pirms poļu uzbrukuma atjaunošanās. 23.~augustā Polijas delegācija paziņoja, ko ``neviena tauta nevarētu piekrist tādai savas suverenitātes ierobežošanai, kādu KSFPR mēģina uzspiest Polijai''. Tai pat laikā arī Polijas delegācijas paziņojums bija pretrunīgs, saturēja jau poļu pretenzijas. Tajā tika norādīts, ka Kērzona līnija neesot ne ar ko pamatota, kaut atzina tautu, dzīvojošu starp Poliju un Krieviju, ``demokrātiskas pašnoteikšanās tiesības''. 25.~augustā J.~K.~Daniševskis jau paziņoja, ka padomju delegācijas priekšlikumi nav ultimatīvi, bet tādi ir 23.~augustā izvirzītie Polijas delegācijas priekšlikumi. Padomju delegācijas vadītājs pārmeta Polijas delegācijai, ka tā, noraidot Kērzona līniju, neizvirza nekādus citus pamatotākus priekšlikumus. 3.~septembrī notika pēdējās sarunas Minskā, pēc trijām nedēļām tās atsākās Rīgā.

% page 430


21.~augustā ASV Valsts departaments nosūtīja Polijas valdībai notu ar prasību nevirzīt savu armiju tālāk pāri Kērzona līnijai uz austrumiem, taču J.~Pilsudskis šajā brīdī jau varēja atļauties ignorēt šādus brīdinājumus. J.~Pilsudskis Polijas Aizsardzības padomē paziņoja, ka Kērzona līnija ir ``bezjēdzība'' un nepieciešamas turpināt uzbrukumu arī pēc tās sasniegšanas. 27.~augustā Aizsardzības padome pieņēma slepenus lēmumus censties ieņemt teritorijas aiz Kērzona līnijas; izspiest, nepieciešamības gadījumā~--- ar spēku, lietuviešus no to kontrolētajām teritorijām bijušajā Suvalku guberņā, kur bija lietuviešu iedzīvotāju vairākums; ieņemt Padomju Krievijas Lietuvai atdotās teritorijas; miera sarunās prasīt teritorijas uz austrumiem no Kērzona līnijas.

18.~septembrī ASV vēstnesis Varšavā H.~Gibsons ziņoja ASV valsts sekretāram, ka, neraugoties uz visu poļu slāņu nogurumu no kara, cerībām uz mieru, J.~Pilsudski ielenkusī militāristu grupa, izmantojot Padomju Krievijas cīņas pret ģenerāļa P.~Vrangeļa vadītajiem ``baltajiem'', nevēlas pieļaut miera noslēgšanu un gatavo Poliju jaunām militārām akcijām.

Arī boļševiku vadība vēl cerēja uz sev labvēlīgu pavērsienu karā un vilcināja laiku. Dažas dienas pēc sakāves pie Vislas sākās Sarkanās armijas uzbrukums Krimai, kur bija nocietinājušies ģenerāļa P.~Vrangeļa vadītie krievu ``baltie''. Padomju Krievijas vadība, atrodoties izvēles priekšā~--- pret ko koncentrēt spēkus~--- Poliju jeb P.~Vrangeļa ``balto'' armiju, izvēlējās otro. Septembrī Varšavā ieradās P.~Vrangeļa pārstāvis ģenerālis P.~Mahrovs. Tikmēr kara darbība poļu-padomju frontē turpinājās. Teorētiski varēja izveidoties abu pretboļševistisko spēku savienība, taču krievu ``baltie'' bija jau parādījuši savu vājumu, Polijas vadība uzskatīja par drošāku noslēgt mieru ar Padomju Krieviju.

Pēc atkāpšanās no Polijas Sarkanā armija nostiprinājās Baltkrievijā, saņēma papildinājumus. Rietumu frontē atgriezās arī 30~tūkstošu pirms tam Austrumprūsijā internēto sarkanarmiešu. 12.~septembrī M.~Tuhačevskis izdeva pavēli gatavot jaunu uzbrukumu. Taču poļi to uzsāka pirmie un ieņēma Koveļas pilsētu. Sarkanā armija bija spiesta atkal atkāpties. 20.~septembrī sākās asiņainas kaujas par Grodņas pilsētu. Tad pat poļu spēki iebruka Lietuvā un, sagrābuši tiltu pār Nemunu, izgāja Sarkanās armijas aizmugurē. Tā steigā atkāpās, zaudējot daudz gūstekņu un ieroču. 12.~oktobrī poļi ieņēma Minsku. J.~Pilsudskis bija iecerējis aplenkt un sagraut Padomju Krievijas Rietumu frontes karaspēku, taču nogurušie poļu spēki nespēja izpildīt pavēli. Nogurušas bija abas karojošās puses. Sarkanās armijas spēku izsīkumu apstiprināja arī aģitācijas un propagandas literatūras tirāžas samazināšanās. Tās nosaukumu skaits no septembra līdz novembrim samazinājās četras reizes, tikpat reižu samazinājās tirāža. Acīmredzot boļševiki šai laikā bija zaudējuši ticību arī drukātā vārda spēkam.

Kara gaitā gandrīz miljons vīru katrā pusē sešu mēnešu laikā veica turpat tūkstoš jūdžu turp un atpakaļ, kas noteica milzu nogurumu abās frontes pusēs. Arī Antantes palīdzības ierobežojumi pēc uzvaras pie Varšavas lielā mērā noteica to, ka poļi nevarēja līdz galam sakaut Sarkanās armijas Rietumu frontes daļas, veicināja miera līguma noslēgšanu. Sarkanās armijas daļas paguva pārgrupēties un ieņemt aizsardzību.

Šajā laikā arī notika jau pieminētā KK(b)P IX konference. V.~Ļeņins, analizējot Sarkanās armijas neveiksmes cēloņus Polijas karagājiena laikā, atzina, ka Polija nav gatava sociālajai revolūcijai. Sarkanajai armijai arī neizdevās sasniegt rūpnieciski attīstītos rajonus ar proletariātu, uz kura atbalstu tā cerēja boļševiki. Boļševiku līderis atzina pirmā ``uzbrūkošās'' politikas realizācijas mēģinājuma neveiksmi, taču izteica apņemšanos: ``Mēs vēlreiz un vēlreiz pāriesim no aizsardzības politikas pie uzbrukuma politikas, kamēr mēs visus [mūsu pretiniekus] nesakausim līdz galam.'' Boļševiki nenožēloja izdarīto, principiāli atzina ``uzbrūkošas darbības'' likumību.

Taču pagaidām Padomju Krievijai nebija spēku lai tuvākajā laikā mēģinātu vēlreiz realizēt pasaules revolūcijas plānus. 1920.~gada 23.~septembrī boļševiku vadoņi, redzot savas Rietumu frontes nespēju pievarēt Polijas armiju, bija spiesti atteikties no Polijas sovjetizācijas un idejas izraisīt revolūciju Vācijā. Padomju Krievija nevarēja atļauties ziemas karu, bija gatava teritoriāli ievērojami piekāpties Polijai, kā arī daļēji atzīt tās prasības par ``Polijas dalības Krievijas [impērijas] attīstībā'' atlīdzināšanas prasību. Polija pieprasīja 300~miljonus rubļu zeltā, padomju puse piedāvāja 10~reizes mazāku summu, kura gala rezultātā arī tika pieņemta. 12.~oktobrī tika parakstīts ``Līgums par pamieru un preliminārajiem (iepriekšējiem) miera noteikumiem starp Krievijas SFPR un Ukrainas SPR no vienas puses un Poliju~--- no otras''. To drīzumā ratificēja visas puses, tas stājās spēkā 18.~oktobrī. Poļu armijas daļas atstāja Minsku un atkāpās uz noteikto līniju, tika nosprausta arī 30~km neitrālā josla. V.~Ļeņins pēc preliminārā miera noslēgšanas rakstīja, ka ``miers tika parakstīts uz noteikumiem, kas Polijai mazāk izdevīgi nekā toreizējie [1920.~gada sākuma] noteikumi'' Polija atteicās turpmāk atbalstīt P.~Vrangeļa ``baltos'' un S.~Petļuras ukraiņu nacionālistus. Ģenerālis P.~Vrangelis savās atmiņās par to rakstīja: ``Poļi savā divkosībā palika sev uzticīgi.''

Sarkanā armija guva iespēju pārsviest karaspēku uz dienvidiem un sakaut P.~Vrangeļa vadīto armiju Krimā.

Kā atzīmējis vēsturnieks Ē.~Jēkabsons, 1920.~gada septembrī un oktobra sākumā šķita, ka beidzot iespējama vismaz daļēja attiecību noregulēšana starp Poliju un Padomju Krieviju, kā arī Baltijas valstīm. Par to liecināja gan poļu~--- padomju miera sarunas Rīgā, gan poļu~--- lietuviešu sarunas Suvalkos, kuras sākās 29.~septembrī pēc Tautu Savienības iniciatīvas. Taču jau 1920.~gada septembrī Polijas armija atkal sāka ieņemt teritorijas, ko Padomju Krievija bija atdevusi Lietuvai, jo Polija neatzina KSFPR un Lietuvas miera līgumu. Lielbritānijas un Francijas spiediena rezultātā Polija gan pagaidām bija spiesta atzīt Lietuvas tiesības uz Viļņu un 1920.~gada 7.~oktobrī abas valstis Suvalkos parakstīja līgumu par pamieru. Tas fiksēja demarkācijas līniju starp abām armijām, atstājot lietuviešu kontrolē arī galveno līdzšinējo nesaskaņu cēloni~--- Viļņu. Īpaši Lielbritānija izdarīja stipru spiedienu uz Poliju, lai tā piekristu šādai Polijas~--- Lietuvas robežai. Arī Polijas pārstāvis Tautu Savienībā 20.~septembrī piekrita lietuviešu ieņemto apgabalu neitralizācijai. Līgumam bija jāstājas spēkā 10.~oktobrī. Galīgi Viļņas liktenis bija jāizlemj Tautu Savienības Kontroles komisijai, kurai vajadzēja ierasties pilsētā ap 9.~oktobri. (Ja mūsdienās (2011) Vilņā dzīvo 63,2\% lietuviešu un tikai 16,5\% poļu, tad pēc 1897.~gada tautas skaitīšanas datiem lietuviešu tajā bija tikai 2,1\%, baltkrievu 4,25\%, toties poļu~--- 30,9\%, ebreju~--- 40\%, krievu~--- 20,1\%. Tiesa, lauku apgabals ap pilsētu bija pārsvarā lietuvisks.) Šķita, ka Polijas~--- Lietuvas konflikts tuvojas atrisinājumam.

Taču J.~Pilsudskis, negribot samierināties ar Viļņas zaudējumu, centās izmantot labvēlīgo situāciju karā pret Padomju Krieviju, lai veiktu daļas \strong{Lietuvas teritoriju okupāciju un aneksiju}. Faktiski Polija atteicās no ar Lietuvu parakstīto līgumu izpildes. Polijas vadību neinteresēja Tautu Savienības ideālu realizēšana, tā bija orientēta uz maksimālu savu teritoriju palielināšanu. Šī mērķa dēļ tika ignorētas Tautu Savienības prasības.

J.~Pilsudskis jau augustā bija sācis gatavot agresīvu akciju. Lai vēlāk būtu iespējams attaisnot dažu karaspēka daļu ieplānoto ``nepaklausību'', tika ieviests sadalījums regulārajās un neregulārajās karaspēka daļās, pēdējo kategorijā ieskaitot arī tā dēvēto Pirmo Lietuvas--Baltkrievijas divīziju (\pltxti{1~Dywizja Litewsko-Białoruska}), kuru komandēja ģenerālis L.~Žeļigovskis un kurā galvenokārt dienēja poļi no austrumu rajoniem.

Pats L.~Žeļigovskis bija dzimis Lietuvā, Viļņas guberņā. Pirmā pasaules kara laikā komandēja pulku Krievijas armijā. Pēc Februāra revolūcijas viņš kļuva par vienu no poļu nacionālo daļu organizētājiem Krievijā, komandēja brigādi 1.~Poļu korpusā, 1918.~gadā izveidoja 4.~poļu strēlnieku divīziju Kubaņā, kura kopā ar ģenerāļa A.~Deņikina komandētajiem baltajiem karoja pret boļševikiem. 1919.~gada aprīlī caur Odesu un Besarābiju divīzija atgriezās Polijā un iekļāvās tās armijā.

1920.~gada 20.~septembrī Polijas armijas Virspavēlniecība izsauca L.~Žeļigovski uz tās mītni Belostokā. J.~Pilsudskis pēc gadiem rakstīja: ``Es izvēlējos ģenerāli Žeļigovski, jo es pats, valsts galva un Polijas augstākais virspavēlnieks, nevarēju pārkāpt saistības. Es izvēlējos ģenerāli, par kuru biju visvairāk pārliecināts, ka pēc sava rakstura viņš spēs noturēties attiecīgā līmenī un militārajā darbībā pretoties valdības pavēlēm un prasībām, kā arī manām pavēlēm un prasībām''. L.~Žeļigovskis 29.~septembrī ieradās Belostokā. J.~Pilsudskis lūdza L.~Žeļigovskim it kā sadumpoties un ieņemt Viļņu, raksturojot viņam situāciju: ``Var iestāties brīdis, kad pret Jums pretī nostāsies ne tikai pasaules, bet arī Polijas viedoklis. Var iestāties brīdis, kad pat es būšu spiests iet pret Jums. Vajadzēs to visu uzņemties uz sevis. To es nevaru pavēlēt. Tādas lietas nepavēl. Taču es vēršos pie Jūsu labās gribas un jums kā viļnieša''. Ģenerālis pēc zināmā pārdomām piekrita izpildīt valsts galvas un virspavēlnieka neoficiālo, bet tāpēc ne mazāk saistošo uzdevumu. Savās atmiņās viņš rakstīja: ``Es gribēju noskaidrot, vai pastāv vēl kāda iespēja atņemt [Lietuvai] Viļņu. Atbildēts tika, ka tādas iespējas nav. Ja neieņemsim, tā mums zudīs uz visiem laikiem''.

Tika sagatavoti plāni valdības surogāta ar Pagaidu Valdības komisijas vārdu radīšanai. Ieņemtajai teritorijai bija jāsaucas par Vidus Lietuvu (poļu \pltxti{Litwa Środkowa}), kas liecināja par to, ka ir iespējama federācijas radīšana. (J.~Pilsudskis cerēja radīt sabiedrotu ar Poliju Lietuvas federāciju, sastāvošu no Ziemeļu Lietuvas (ar Kauņu kā centru), vidus Lietuvas (ar Viļņu) un Dienvidu Lietuvas (ar Novogrudokas (poļu \pltxti{Nowogródek}, baltkrievu \betxti{Навагрудак, Наваградак}) pilsētu, kas tagad atrodas Baltkrievijā.) Tiesa, daļa Lietuvas--Baltkrievijas divīzijas virsnieku izteica protestu pret plānojamo operāciju. Tikai pēc J.~Pilsudska prasības 3.~armijas, kurā divīzija ietilpa, pavēlnieka ģenerāļa V.~Sikorska veiktā iejaukšanās nomierināja prātus. Toties absolūtais vairākums no Polijas austrumu provincēm nākušo karavīriem ar entuziasmu apsveica gājienu uz Viļņu. Par šiem ``austrumu kresu’’ iedzīvotājiem L.~Žeļigovskis liecināja: ``Visi saprata, par ko ir runa. Viņos es atradu morālu atbalstu''.

L.~Žeļigovskis izpildīja J.~Pilsudska neoficiālo rīkojumu un 1920.~gada 8.~oktobrī nosūtīja Polijas armijas pavēlniecībai paziņojumu, kurā bija norādīts: ``Secinot, ka ar Kauņas valdību noslēgtajā līgumā paredzētās pamiera līnijas jau iepriekš, mums Viļņas, Grodņas un Lidas zemes iedzīvotājiem~--- par sliktu kopā ar polisko Viļņu atdod lietuviešu rokās mūsu zemes, izlēmu ar ieročiem rokās aizstāvēt [mūsu] dārgās tēvzemes iedzīvotāju pašnoteikšanās tiesības un, nespējot rīkoties pretēji savai sirdsapziņai un paklausot pilsoņa pienākuma izjūtai, uzņēmos vadību pār karavīriem, kuri cēlušies no šīs zemes. Ar nožēlu pasludinu, ka atstāju savus pienākumus un aizeju no grupas vadības. Audzināti disciplīnā, uzticīgi tēvzemes atbrīvošanas idejai, man pakļautie komandieri un karaspēks klausa manām pavēlēm\dots{}''. Tai pat dienā poļu daļas devās uz nedaudz vairāk nekā 50~km attālumā esošo Viļņu. Faktiski J.~Pilsudskis ar L.~Žeļigovska rokām pārkāpa tikko noslēgto Polijas-Lietuvas līgumu par pamieru, ignorējot arī Tautu Savienības pozīciju, nostājās uz atklātas agresijas ceļa. Karavīri saņēma pavēli~--- gūstā saņemtos lietuviešus atbruņot un atlaist, izskaidrojot, ka poļi negrasās karot ar lietuviešiem, ``bet vienkārši pēc kara atgriežas mājup''. Lietuviešiem poļu karaspēka darbības zonā bija tikai viens pulks un vēl divi Viļņā. Tie nespēja efektīvi aizstāvēties. Ziņojumi par poļu uzbrukumu izsauca paniku, vakarā tika izsludināta pilsētas evakuācija un Lietuvas valdības pilnvarotais nodeva varu Tautu Savienības delegātam, franču pulkvedim K.~Rebulam. Viņš pasludināja sevi par Viļņas pagaidu gubernatoru un izsludināja pilsētā aplenkuma stāvokli.

9.~oktobra otrajā pusē L.~Žeļigovska vienības ieņēma Viļņu, kur tās sajūsmināti sveica poļu tautības iedzīvotāji. Ģenerāli pieņēma gubernators, piedāvājot, lai Viļņa kļūtu par brīvpilsētu. L.~Žeļigovskis šo priekšlikumu noraidīja, pieprasīja lai līdz 12.~oktobrim Viļņu atstāj visi sabiedroto valstu pārstāvji. Tāpat viņš izrīkojās ar 10.~oktobrī Viļņā ieradušos Tautu Savienības delegāciju un 12.~oktobrī ar notām Polijas, Lietuvas un Antantes valstu valdībām pasludināja ``neatkarīgu'' Vidus Lietuvas valsti (poļu \pltxti{Republika Litwy Środkowej}, lietuviešu \lttxti{Vidurio Lietuvos Respublika}).

Lietuvieši asi protestēja, griežoties pie lielvalstīm un Tautu Savienības. Francija juta līdzi Polijai, ASV valdība paziņoja, ka to šī problēma neinteresē. Asi reaģēja tikai Lielbritānija. Polijas valdība centās pierādīt, ka tai nav nekā kopēja ar L.~Žeļigovska izrīcību. 14.~oktobrī Seimā premjerministrs V.~Vitoss nosodīja ģenerāļa ``patvaļīgo'' operāciju, taču pēc tam paziņoja, ka Polijas valdība saprot motīvus, no kuriem viņš vadījies un deva mājienu, ka ir iespējams nokārtot attiecības ar Vidus Lietuvu. Tautu Savienības Padome 25.~oktobrī pieņēma lēmumu, kurš paredzēja sarīkot Viļņas zemē plebiscītu Savienības uzraudzībā.

Polija tam piekrita, gan uzstājot, ka teritorijai, kur jānotiek plebiscītam, jāsniedzas līdz Kauņas pievārtei, kur arī dzīvoja poļi. Lietuvieši arī piekrita plebiscītam dažās teritorijās, taču izņemot Viļņu un tās apgabalu. 18.~oktobrī lietuvieši uzsāka pretuzbrukumu, taču tika smagi sakauti. L.~Žeļigovskis pat sāka plānot uzbrukumu Kauņas virzienā, taču Polijas valdības norādījumu rezultātā no tā atteicās. 29.~novembrī atkal tika noslēgts pamiers. Turpmākās Polijas, Lietuvas un Tautu Savienības sarunas nedeva rezultātus.

Okupētā tā dēvētās Vidus Lietuvas teritorija bija aptuveni 15~000~km$^{2}$ liela, tajā dzīvoja ap 500~000 iedzīvotāju. No tiem, pēc poļu datiem 68\% bija poļi, 18\%~--- lietuvieši. Lietuviešu dati atšķīrās, taču nemainīja tendenci, tāpēc Lietuva nevēlējās apgabala piederības jautājumu izlemt ar plebiscīta palīdzību.

1922.~gada janvārī notika Viļņas Zemes (Vidus Lietuvas) vēlēšanas, kuram bija jānosaka Viļņas piederība. Jau pirms tām Latvijas pārstāvniecības Polijā pagaidu vadītājs Pēteris~Oliņš reāli vērtēja situāciju: ``Nevar, zināms, šaubīties, kāda būs ``tautas griba'', kuru izteiks zem Žeļigovska iestāžu spiediena izvēlētais Seims''. Tik tiešām vēlēšanās, kuras lietuvieši un daļa baltkrievu un ebreju boikotēja (piedalījās 64\% balsstiesīgo iedzīvotāju), ievēlētais Seims 1922.~gada februārī gandrīz vienbalsīgi (6 deputātiem atturoties) pasludināja zemes piederību Polijai. Polijas Seims Varšavā 24.~martā ratificēja Viļņas Seima lēmumu. Lietuva to neatzina. Arī Tautu Savienība atzina Viļņas Seima vēlēšanas par nelikumīgām, taču atturējās spert kādus konkrētus soļus lai sodītu agresoru. Vēl vairāk~--- piekāpās Polijas realizētās \latxti{fait accompli} (latīņu~--- noticis fakts) politikas priekšā~--- 1923.~gada 3.~februārī ar Tautu Savienības lēmumu Viļņas novads tika atzīts par Polijas sastāvdaļu. Faktiski Polija anektēja Viļņu un tās apgabalu. (Der atzīmēt, ka 1922.~gada februārī lietuvieši ``nokopēja'' L.~Žeļigovska akciju un ar safabricēta ``apvērsuma'' palīdzību ieņēma Klaipēdu. Jau 1921.~gadā Padomju Krievija tāpat izrīkojās Gruzijā.) 1923.~gadā J.~Pilsudskis, uzstājoties Viļņā, pats atklāja, ka L.~Žeļigovska ``dumpis'' un iebrukums Viļņā bija inspirēti pēc viņa tieša norādījuma.

Lietuva, neraugoties uz Tautu savienības mēģinājumiem noregulēt konfliktu, neatzina Viļņas apgabala piederību Polijai, atteicās uzturēt ar Poliju diplomātiskās attiecības. 1927.~gadā Lietuvā demonstratīvi tika mainīta konstitūcija. To papildināja ar punktu: ``Lietuvas galvaspilsēta ir Viļņa. Ar atbilstošu likumu tā uz laiku var tikt pārcelta citur''. Atbilde bija lietuviešu vajāšana Viļņas apgabalā. Poļu nacionāldemokrāti savā propagandā apgalvoja, ka nacionālajās nesaskaņās esot vainīgas Vācija un Krievija (Padomju Savienība), kuras negatīvi ietekmējot ``kļūdaini vadīto lietuviešu brāļu tautu''. Starp Polijas un Lietuvas valstīm pastāvēja ``ne kara, ne miera'' stāvoklis. Pat pasta sakari norisa caur Kēnigsbergu. Līdz pat 1938.~gadam Polijas attiecības ar Lietuvu nebija noregulētas. Nebija šaubu, ka Polija jebkurā brīdī varēja iekarot vājāko Lietuvu. No tā atturēja vienīgi bailes par kaimiņu, pirmkārt jau Padomju Savienības reakciju. Diemžēl jākonstatē, ka arī mūsdienu Polijas valdošajās aprindās nav nožēlas par pret Lietuvu pieļauto vardarbību. Tā, 2009.~gada maijā, uzstājoties televīzijā un runājot par vāciešu vidū izplatīto uzskatu, ka pēc Otrā pasaules kara Eiropas karte tika par jaunu piegriezta, bet viņu tautieši nelikumīgi izdzīti no tagad Polijai piederējušajām teritorijām, Polijas ārlietu ministrs R.~Sikorskis pieminēja arī lietuviešus: ``Lietuva uzskata, ka starpkaru periodā mēs okupējām Viļņu, taču mēs tā neuzskatām.''

Sakarā ar Viļņas un tās apgabala sagrābšanu arī Latvijas sabiedrībā radās bažas par Polijas iespējamo tālāko rīcību, nodrošinot sev maksimāli izdevīgas robežas uz Baltijas valstu rēķina. Latvijas valdībai jau sākotnēji bija skaidrs, ka L.~Žeļigovskis rīkojās saskaņā ar J.~Pilsudski. Turklāt Latvijas stāvokli sarežģīja tas, ka daļu Ilūkstes apriņķa šai laikā ieņēma Lietuvas armija, tādēļ radās draudi, ka Polija varētu okupēt arī šo teritoriju. Tādēļ 1920.~gada 12.~oktobrī Latvijas armijas virspavēlnieks J.~Balodis paziņoja Lietuvas armijas virspavēlniekam, ka atbilstoši ziņām par L.~Žeļigovska pret Lietuvu uzsākto darbību ``Latvijas armija ir spiesta darīt visu, lai nodrošinātu savu teritoriju'', ieņemot bijušo Kauņas un Kurzemes guberņu robežu. Tā arī notika, kad lietuviešu daļas bez pretestības atkāpās. 1920.~gada otrajā pusē Latvijai bīstamāks bija ilglaicīgs robežkonflikts ar Poliju par sešiem Ilūkstes apriņķa pagastiem. Tiesa, lai nesaasinātu attiecības vēl arī ar Latviju, poļu daļas saņēma pavēli nepāriet bijušās Kurzemes guberņas robežu un uzskatīt to par demarkācijas līniju līdz laikam, kad abu valstu valdības vienosies par robežu. Taču vēlākajos gados turpinājās nesaskaņas šai jautājumā un Polija neatzina Latvijas robežu par noteiktu. Tikai 1937.~gadā tā tika abpusēji noteikta un iezīmēta dabā.

Pēc Latvijas izlūkdienesta ziņām 1920.~gada nogalē Ilūkstes apriņķī uzradās vairāki bijušie tā iedzīvotāji, par kuriem bija zināms, ka viņi dienēja Polijas armijas izlūkošanas un pretizlūkošanas struktūrās. Grīvā divi tādi arī tika arestēti un pie viņiem atrada dažādas Polijas, Lietuvas, Krievijas un Latvijas pases, dažādus zīmogus u.c. Kā atzinis Latvijas-Polijas attiecības pētījušais vēsturnieks Ē.~Jēkabsons, Polijai un arī ģenerālim L.~Žeļigovskim reāli nebija agresīvu nodomu pret Latviju, Polijas valdība un Ārlietu ministrija darīja visu, lai nerastos nesaprašanās. Un nebūtu arī loģiski Polijai, tikko novērsušai smagāku atmaksu par uzbrukumu Padomju Krievijai, iestigušai konfliktā ar Lietuvu, meklēt vēl kādu ienaidnieku. Tomēr jau ar draudiem Lietuvas pastāvēšanai vien tika apdraudēts arī Latvijas stāvoklis, jo bija skaidrs, ka Varšavai pakļauta Polijas--Lietuvas valsts noteikti tiektos paplašināt savas robežas arī uz Latvijas rēķina. Latvijā sabiedriskā doma pauda nemieru ar Polijas agresīvo rīcību. Dzejnieks Eduards Virza presē rakstīja, ka nedrīkst samazināt izdevumus armijas vajadzībām, kamēr ``Žeļigovskis posta leišu brīvību un viņa divīzijas vērstas arī pret Daugavpili''. 20.~novembrī Latvijas ārlietu ministrs Z.~A.~Meierovics iesniedza Polijas sūtnim V.~Kameņeckim konfidenciālu notu, kurā deklarēja Latvijas neieinteresētību Viļņas valstiskās piederības jautājumā, taču brīdināja ``ja ģenerāļa Žeļigovska karaspēks sāktu gājienu uz Kauņu, Latvija būtu spiesta nostāties Lietuvas pusē.'' Ģenerāļa L.~Želigovska vadītais iebrukums 1920.~gada oktobrī Viļņas apgabalā un Polijas centieni ar koridora izveidi starp Lietuvu un Padomju Baltkrieviju sasniegt pieeju Daugavai vēl vairāk samežģīja un apgrūtināja piecu Baltijas valstu savienības noslēgšanu.

Viens no Viļņas apgabala sagrābšanas rezultātiem bija tiešas 104~km garas Polijas~--- Latvijas robežas izveidošanās, tādejādi uz diviem gadu desmitiem izolējot Lietuvu no Padomju Krievijas un neļaujot tai arī vēlāk saņemt no Padomju Krievijas / Padomju Savienības lielāku atbalstu pret Poliju.

Arī turpmākajos gados Latvija formāli ieturēja neitralitāti Viļņas jautājumā, taču faktiski ar to tika atzīta pilsētas atrašanās Polijā. Viļņas konflikts bija viens no galvenajiem iemesliem, kas trīsdesmitajos gados liedza vienoties Baltijas valstīm, traucēja izveidot drošības sistēmu, kura, iespējams, būtu varējusi novērst vai vismaz būtiski izmainīt 1939.--1940.~gada notikumus reģionā, bet varbūt arī visā Eiropā.

Latviju \latxti{de~iure} Polija atzina tikai 1921.~gada 27.~janvārī, pēc tam, kad to bija atzinušas visas ievērojamākās Eiropas valstis. Kā atzīmējis mūsdienu poļu vēsturnieks T.~Paļišinskis, Polijas kavēšanos noteica tās politiskā atkarība no Francijas diplomātijas, no brīža, kad protektors no Sēnas krastiem atzītu Latvijas Republiku. Bez tam, Polijas Ārlietu ministrija esot bijusi pārliecināta, ka Francija pirms Baltijas republiku atzīšanas konsultēsies ar Polijas Republikas valdību. Tomēr tas nenotika un tāpēc Polija nokavēja.

Latviešu vēsturnieks Ē.~Jēkabsons ir aprakstījis arī Polijas pašas diplomātisko spēli Latvijas atzīšanas jautājumā. Tā 1920.~gada oktobra beigās no Polijas puses Latvijai tika doti mājieni par iespējamu Latvijas atzīšanu, ja tā apņemtos ievērot pilnīgu neitralitāti Polijas~--- Lietuvas konfliktā, dot Polijai ostu nomas tiesības uz 99~gadiem, neīstenot Latgalē zemes reformu un noslēgt ar Poliju militāru konvenciju. Decembrī Polija izteica vēlmi vispirms noslēgt ar Latviju saimniecisku vienošanos, kas Polijai ļautu izmantot Latvijas ostas. 15.~decembrī Ženēvā balsojumā par Baltijas valstu uzņemšanu Tautu Savienībā Polija pārstāvis I.~Paderevskis atturējās, atsaucoties uz to, ka to nevēloties Polijas ``labvēļi''~--- Rietumu lielvalstis. 1920./1921.~gada mijā Lielbritānija un Francija bija zaudējušas cerības uz boļševistiskā režīma drīzu gāšanu, pirmā arī nostājusies uz tirdzniecības sakaru nodibināšanas ceļa ar Padomju Krieviju, tāpēc 1921.~gada 26.janvārī Antantes Augstākā padome atzina Latviju ``\latxti{de~iure}''. Tad nu arī Polijas diplomātiskie pārstāvji paziņoja par Latvijas atzīšanu ar atpakaļejošu datumu~--- 1920.~gada 31.~decembri. Latvijas oficiālie pārstāvji neiebilda, taču žurnālisti, pēc tam arī visa sabiedrība, uzzināja patiesību un tas atstāja negatīvu iespaidu uz Polijas prestižu Latvijā. Jāpiekrīt Ē.~Jēkabsonam, ka ``Polijas galvenie politiskie grupējumi 1919.--1920.~gadā pret Baltijas valstīm izturējās augstprātīgi, protekcionistiski un ar pārākuma apziņu.''

1920.~un 1921.~gada mijā Polijas militārā vadība nolēma Latvijā izveidot poļu izlūkošanas centru, kurā tiktu koordinēta visu Krieviju aptverošā izlūkdienesta darbība. No poļu izlūkdienesta vajadzību viedokļa Latvijai bija izdevīga ģeopolitiskā atrašanās vieta. Nemazsvarīgu lomu spēlēja arī līdzšinējā militārā sadarbība. Poļu izlūkdienesta struktūras Latvijā darbojās ar Latvijas izlūkdienesta ziņu, sadarbībā ar to, galvenokārt pret Padomju Krieviju, bieži savstarpēji apmainoties ar saņemto informāciju.

Polijas Kara ministrija arī turpināja veicināt krievu ``balto'', ukraiņu nacionālistu un arī tāda Pilsoņu kara laikam raksturīga avantūrista kā S.~Bulak-Balahoviča vienību formēšanu Polijas armijas sastāvā. Viņš no 1915.~līdz 1920.~gadam bija kāpis pa karjeras kāpnēm no praporščika (jaunākā leitnanta) līdz ģenerālmajoram, kalpodams gan cara, gan ``sarkano'', gan ``balto'' armijās, gan Igaunijas un pilsoniskās Baltkrievijas Tautas Republikas, visbeidzot no 1920.~gada pavasara arī Polijas armijā. J.~Pilsudskis B.~Savinkovam par S.~Bulak-Balahoviču izteicās: ``Jā, bandīts, bet ne tikai. Šodien~--- viņš ir krievs, rīt~--- polis, parīt~--- baltkrievs, nākamajā dienā~--- nēģeris'', taču pēc tam piebilda: ``Lai arī Balahovičs ir bandīts, taču tā kā nav izvēles, labāk, tomēr, lietas kārtot ar Balahoviču, nekā ar ``zeltuzplečainajiem'' ģenerāļiem''. (Domāti krievu ``balto'' ģenerāļi, kuri nevēlējās aizmirst Krievijas intereses~--- V.Š.) S.~Bulak-Balahovičs no jūnija piedalījās arī poļu-padomju karā, veicot daudz nelikumību.

Kara gaitā abas: gan padomju, gan Polijas puses paveica daudz \strong{vardarbību pret civiliedzīvotājiem}. Atbaidošā veidā izpaudās Pirmā pasaules kara laikā radušie morāles kropļojumi, kuri saplūda ar šķirisko un nacionālo naidu pret ekspluatatoriem un apspiedējiem~--- kā ``poļu paniem'', tā ``moskaļiem'', kuri pie tam vēl bija boļševiku varas nesēji. Pie tam īpaši cieta ebreji. Kā Sarkanās armijas, tā arī Polijas armijas pavēlniecība gan veica šādu gadījumu izmeklēšanu, lai sodītu vainīgos, centās tādus ekscesus nepieļaut, taču panākumi bija pieticīgi.

Polijas un Baltkrievijas arhīvos glabājas dokumenti par 1919.--1921.~gadā notikušo ap 30 Rietumbaltkrievijas ciematu iznīcināšanu ar visiem iedzīvotājiem. 1919.~gada martā, kad Pinskas (baltkrievu \betxti{Пінск}, krievu \rutxti{Пинск}) pilsētu ieņēma poļu vienības, tās komandants lika nošaut 40~ebrejus, kuri bija sapulcējušies uz dievlūgšanu, noturot šo dievkalpojumu par boļševiku sapulci. Aprīlī, ieņemot Viļņu, notika izrēķināšanās ar gūstā kritušajiem sarkanarmiešiem, padomju varas līdzjutējiem. Polijas armijas uzbrukumu Ukrainā pavadīja ebreju grautiņi un masu apšaušanas. Pēc Ebreju enciklopēdijas ziņām padomju-poļu kara laikā daudzi ebreji ``tika nogalināti, jo viņus apvainoja spiegošanā Sarkanās armijas labā''. Pēc krievu vēsturnieku datiem tika šauti cilvēki Rovnas (ukraiņu \uktxti{Рівне}, poļu \pltxti{Równe}) pilsētā, nodedzināti vairāki ciemati. Poļu vēsturnieki gan šos datus apstrīd, norādot, ka grautiņus un citas vardarbības rīkojušas ukraiņu nacionālistu vienības. Taču ebreju sabiedriskais darbinieks un publicists D.~Pasmaniks izsvērti norādījis, ka ebreju grautiņi ``sākās ``atjaunotajā'' Polijā uzreiz nākamajā dienā pēc tam, kad tā bija atkal kļuvusi brīva un neatkarīga valsts. Savukārt pašā Krievijā grautiņus sāka demokrāta Petļuras un sociālista Viņņičenko ukraiņu karaspēks \citespace{} Ukraiņi grautiņus pārvērtā ikdienišķā parādībā.'' Mūsdienu krievu vēsturnieks G.~Kostirčenko, apkopojot dažādus ebreju avotus, uzskaitījis līdz pat 900~masveida grautiņu, no kuriem 40\% bija petļuroviešu, Ukrainas Direktorijas piekritēju rīkoti, 25\%~--- Ukrainas ``\emph{batjku}'' bandu rīkoti, 17\%~--- deņikiniešu un 8,5\% S.~Budjonnija 1.~jātnieku armijas un citu ``sarkano'' rīkoti. Var piebilst, ka, atriebjot ebreju slepkavības, 1926.~gadā ebreju publicists, anarhists Š.~Švarcbards nogalināja S.~Petļuru, bet Francijas zvērināto tiesa viņu attaisnoja.

Tomēr īpašu vietu terorā pret civiliedzīvotājiem ieņem tieši sākotnēji Polijas dienestā esošā, vēlāk it kā ``atsevišķi'' karojošā ģenerāļa S.~Bulaka-Balahoviča tā saucamo ``baltkrievu'' daļu darbība.

Jau 23.~septembrī pēc Ģenerālštāba priekšnieka ģenerāļa T.~Rozvadovska pieprasījuma tika radīta izmeklēšanas komisija, kura izdarot vietējo iedzīvotāju aptauju viņa vienību izvietojuma vietās, noskaidroja, ka tur masveidā notika vardarbības, slepkavības un laupīšanas, bet poļu varas iestādes tam nepievērsa uzmanību. Kad 1920.~gada 12.~oktobrī tika noslēgts pamiers, Polijas Seims pieprasīja lai militārā pavēlniecība visas ``sabiedroto'' (krievu, ukraiņu) daļas vai nu atbruņo vai līdz 2.~novembrim izraida no Polijas. Tas attiecās arī uz S.~Bulak-Balahoviča vadītajām daļām. Tad arī B.~Savinkova un S.~Bulak-Balahoviča saformētā t.~s. Krievu tautas brīvprātīgā armija (\rutxti{Русская народная добровольческая армия}~---  nejaukt ar A.~Deņikina komandēto krievu ``balto'' Brīvprātīgo armiju), kura pēc dažām ziņām sasniedza no 10~līdz 20~tūkstošu vīru skaitu, neraugoties uz noslēgto Polijas un Padomju Krievijas vienošanos par kara darbības izbeigšanu, iebruka Baltkrievijas teritorijā ar uzdevumu radīt Baltkrievijas Republiku un veiksmes gadījumā uzbrukt arī Maskavas virzienā. Polijas valdība oficiāli izlikās, ka tai nav nekāda sakara ar šo ``pašdarbību''.

S.~Bulak-Balahoviča ``armija'' rīkoja masveida grautiņus, laupīšanas, slepkavoja civiliedzīvotājus. 1920./1921.~gada ziemā visā Baltkrievijā dažādi politiski krāsoti bandīti veica 40~grautiņus, tai skaitā Moziras apriņķī, kur darbojās bulahovieši~--- 21.

Poļu kara prokurors, kurš izmeklēja sūdzības par to darbību, ziņoja: ``\dots{}Balahoviča armija ir laupītāju banda, kura pārsūta salaupīto zeltu. Lai ieņemtu kādu pilsētu, tiek sūtīta armija, kuras karavīri laupa un slepkavo. Un tikai pēc daudziem grautiņiem, divas dienas vēlāk, ierodas Balahovičs ar savu štābu. Pēc laupīšanas sākas pļēgurošana \citespace{} Kas attiecas uz Balahoviču, viņš ļauj laupīt, citādi viņi atteiktos doties uz priekšu \citespace{} katrs virsnieks, kurš iestājas Balahoviča armijā, aplej sevi ar mēsliem, kurus ne ar ko nevar nomazgāt.'' Izmeklēšanā noskaidrojās, ka tikai Turovas (baltkrievu \betxti{Тураў}, krievu \rutxti{Туров}) pilsētiņā balahovieši izvaroja 70~ebreju meiteņu vecumā no 12~līdz 15~gadiem.

Pēc Izraēlas institūta ``\emph{Yad Vashem}'' datiem, ceturtā daļa ebreju, kas gāja bojā grautiņos Pilsoņu kara laikā, ir uz S.~Bulak-Balahoviča ``armijas'' rēķina.

1920.~gada 27.~oktobrī Berlīnē iznākošā krievu emigrantu avīze ``\rutxti{Голос России}'' (``Krievijas Balss'') publicēja ziņas par šīs ``armijas'' noziegumiem: laupīšanām, iedzīvotāju spīdzināšanu, sieviešu, tai skaitā 12~gadīgu meitenīšu, izvarošanu utt.

Poļu vēsturnieks M.~Kabanovskis 1993.~gadā Varšavā izdotā grāmatā “\pltxti{General Stanislaw Bulak-Balachowicz}''(``Ģenerālis Staņislavs Bulak-Balahovičs'') publicējis divu liecinieku stāstīto: ``\dots{}pa ceļam sastapām kapteini~--- balahovieti. Viņš pajautāja: ``Ko vedat~--- Ebrejus \dots{}~--- Nošaut viņus. Ar mums bija vēl viens ebrejs~--- Maršalkovičs. Konvojētāji pavēlēja mums nolaist apakšbikses un laizīt vienam otra dibenu. Pēc tam tāpat piespieda urinēt vienam otra mutē un darīt citas derdzības \dots{} Bet zemniekus sapulcināja apkārt un pavēlēja tajā visā noskatīties. Lika mums veikt dzimumaktu ar telīti. Izvaroja mūs un apķēzīja sejas \citespace{} Blūmenkranks nevarēja paciest ņirgāšanos un lūdza viņu nošaut.''

Acīmredzot pēc kārtējo ziņu saņemšanas par šīm zvērībām, V.~Ļeņins Tautas komisāru padomes sēdes laikā dusmās uzrakstīja zīmīti revolucionārās kara padomes priekšsēdētāja vietniekam J.~Skļanskim ar priekšlikumu sodīt tās valstis (V.~Ļeņins kļūdaini domāja, ka tās ir Latvija un Igaunija, iespējams tāpēc, ka S.~Bulak-Balahovičs pirms tam bija uzturējies to teritorijā) ievedot Sarkanās armijas daļas to teritorijā un pakarot ``100--200 to ierēdņus un bagātniekus''. Drīz pēc tam V.~Ļeņins uzrakstīja vēl vienu zīmīti, kurā ierosināja ``zaļo'' izskatā iebrukt 10-20~verstis pretinieka teritorijā un tur pakārt ``kulakus, popus, muižniekus.'' Oficiāli gan KSPFR 1920.~gada 28.~oktobrī nosūtīja notu Lielbritānijas valdībai, kurā norādīja, ka pēc pamiera ar Poliju un miera līgumu ar Igauniju, Latviju un Lietuvu noslēgšanas ``kara stāvoklis faktiski turpinās'', Baltkrievijā un Rietumukrainā bruņotas bandas, kuras nepakļaujas nevienai valdībai un kuras komandē ``Balahovičs un Petļura'', turpina darboties pret abu republiku mierīgajiem iedzīvotājiem. Tika izteikts brīdinājums, ka Krievija un Ukraina ir nolēmušas darīt galu šo bandu nelikumīgajai darbībai. Kad 10.~novembrī S.~Bulak-Balahoviča ``armija'' ieņēma Moziras (baltkrievu~--- \betxti{Мазыр}) pilsētu, viņš pat pasludināja sevi par ``Baltkrievijas valsts priekšnieku'' un nozīmēja par jaunās valdības premjeru baltkrievu nacionālo darbinieku P.~Aļeksjuku. Taču drīz Sarkanajai armijai izdevās sakaut S.~Bulak-Balahoviča ``armiju'', pēc desmit dienām ieņemt Moziru. S.~Bulak-Bulahoviča daļas ar grūtībām izlauzās atpakaļ uz Poliju un tur tika internētas Padomju Krievijas pārstāvja klātbūtnē. Padomju valdība pieprasīja Polijai izdot S.~Bulak-Balahoviču, taču J.~Pilsudskis atteicās to darīt.

Daļa bijušās S.~Bulak-Bulahoviča ``armijas'' karavīru turpināja teroristiskus uzbrukumus padomju Baltkrievijai līdz pat 1922.~gadam. (Pēc 1921.~gada 18.~marta Rīgas miera noslēgšanas uz Franciju bija spiests izbraukt B~Savinkovs, bet vīzu atpakaļceļam uz Poliju nesaņēma. Pretēji J.~Pilsudska gribai tad no Polijas tika izsūtīti arī S.~Bulak--Bulahovičs un S.~Petļura. 1940.~gadā jau vācu okupētājā Polijā S.~Bulak--Bulahovičs gāja bojā nenoskaidrotos apstākļos, pēc dažām ziņām~--- iereibis saķildojies ar vācu patruļu, tāpēc nošauts uz vietas.) Tā neveiksmīgi beidzās kārtējā avantūra, kura, iespējams, no tās poļu atbalstītāju puses bija iecerēta kā līdzīga Viļņas pievienošanai Polijai.

Padomju pusē galvenās grautiņu veicējas slavu ieguva 1.~Jātnieku armija. Pēc mūsdienu Krievijas avīzē ``\rutxti{Совершенно секретно}'' (``Pilnīgi Slepeni'') publicētajām armijas vadības ziņojumiem un pavēlēm redzams, ka tikai no 1920.~gada 18.~līdz 22.~septembrim šīs armijas 6.~kavalērijas divīzijas karavīri veica vairāk nekā 30~grautiņus, 29.~septembrī grautiņā nogalināja 60~cilvēku, naktī uz 3.~oktobri nogalināja 21, ievainoja 12~cilvēku ``un izvaroja daudz sieviešu''. Kad divīzijas komandieris mēģināja apturēt grautiņu, viņš pats tika nogalināts. Par noziegumiem armijas pavēlniecība divīziju izformēja, divi brigāžu komandieri un vairāki simti ierindnieku tika tiesāti, 157~nošauti. Pēc karagājiena dalībnieka, vēlākā rakstnieka I.~Bābela teiktā, pēc Ļ.~Trocka pavēles par ebreju miestiņu laupīšanu tika nošauti ap 300 ``sarkano'' kazaku. Taču notikušajā lielā mērā bija vainīga visa padomju sistēma, kura nespēja nodrošināt armijas daļu normālu apgādi un reizē mācīja meklēt ``šķiras ienaidniekus'', par kādiem viegli varēja kļūt gan turīgākie zemnieki, gan mūžīgie ``grēkāži''~--- ebreji. Nebija brīnums, ka dažādu tautību vietējie iedzīvotāji bieži Sarkanās armijas daļas sagaidīja ar cerībām, pavadīja~--- ar lāstiem.

Taču uzreiz jāuzsver, ka šajā laikā notika jau Sarkanās armijas atkāpšanās un šie ekscesi norisa nevis poļu, bet ukraiņu un ebreju apdzīvotajās teritorijās (kas tāpēc tos tāpat nekādi neattaisno). Acīmredzot Sarkanās armijas daļām, kad tās iegāja etniskās Polijas teritorijā, bija doti stingri norādījumi nekompromitēt revolūcijas ideju un tās cīnītājus ar laupīšanām. Kā šie norādījumi darbojās liecina Varšavas kara apgabala izlūkošanas un pretizlūkošanas nodaļas 1920.~gada 19.~septembra ziņojums: ``Padomju karaspēka izturēšanās visā okupācijas perioda laikā bija nevainojama, pierādīts, ka līdz atkāpšanās brīdim tas nepieļāva nekādas nevajadzīgas laupīšanas un vardarbības. Rekvizīcijas tas centās īstenot formāli un maksāja prasītās cenas naudā, kaut arī zaudējušā vērtību. Nevainojamā padomju karaspēka uzvedība salīdzinājumā ar vardarbību un nevajadzīgām laupīšanām, ko atkāpjoties veica mūsu karaspēks, būtiski grāva uzticību Polijas varas iestādēm''. Diezin vai, sākot strauju atkāpšanos, šī stingrā kārtība Sarkanajā armijā saglabājās. Šādos apstākļos apgāde taču gandrīz nefunkcionēja, un karavīriem vajadzēja pašiem gādāt sev iztiku.

Karā pret Poliju 1920.~gadā Sarkanā armija pēc aptuveniem aprēķiniem zaudēja ap 130~000 kritušo, ar ievainotajiem kopējie zaudējumi sasniedza ap 232~tūkstošus. Bez tam abās karojošajās armijās plosījās epidēmijas: tīfs, holēra, bakas, dizentērija. No 1920.~gada aprīļa līdz oktobrim Polijas armija bija zaudējusi 184~246~cilvēkus~--- gandrīz 20\% no sava maksimālā sastāva. Kopumā kara rezultātā abās pusēs dzīvību un veselību zaudēja vairāk nekā pusmiljons karavīru un civiliedzīvotāju.

Neveiksmīgā karagājiena uz Varšavu laikā un neorganizētā Sarkanās armijas atkāpšanās veda pie lieliem zaudējumiem, arī liela gūstekņu skaita. Ir zināms, ka militārie ziņojumi nekad nav bijuši precīzi attiecībā par gūstā saņemtajiem pretinieka karavīriem, kā arī par saviem karavīriem, nonākušiem pretinieka gūstā. Pilnā mērā tas sakāms arī par poļu-padomju karu.

Precīzu ziņu par kopējo \strong{gūstā kritušo sarkanarmiešu} un internēto padomju Krievijas civiliedzīvotāju skaitu Polijā diemžēl nav.

Krievu vēsturnieks G.~Matvejevs norādījis, ka pieejamie poļu dokumenti nesatur kaut-cik sistemātiskus datus par gūstā nonākušajiem sarkanarmiešiem. Polijā, neraugoties uz atkārtotām augstāko instanču pavēlēm un instrukcijām, nebija sakārtota stingra gūstekņu un ar tiem notiekošo izmaiņu uzskaites sistēma, kā to prasīja starptautiskās tiesības. Politiskajos memorandos un arī vēstures literatūrā lietojamie dati mainās arī atkarībā no politiķu un vēsturnieku nacionālās piederības.

Kā raksta krievu vēsturnieks M.~Meļķuhovs, pēc ``sociālistiskās'' iekārtas nodibināšanas Polijā Padomju Savienībā nemēdza daudz runāt par padomju-poļu karu, lai nebojātu abu valstu attiecības. Netika runāts ne par poļu iebrucēju ļaundarībām Baltkrievijā un Ukrainā, ne padomju karagūstekņu ``nogalināšanu ar aukstumu un badu'' poļu nometnēs. Padomju filmas rādīja notikumus visās citās Pilsoņu kara frontēs, bet divu gadu kara darbība padomju-poļu kara frontē tika apieta ar uzmanību. Dažkārt Poliju apmeklējošos padomju tūristus arī brīdināja, lai tie neiedomājas dziedāt Padomju Savienībā populāro D.~Pokrasa un A.~Surkova ``Jātnieku armijas dziesmu'', kurā gāja runa par ``poļu paniem'' un spējajiem ``mūsu jātnieku triecieniem''. Taču tautu atmiņā fakti par nodarītajām pārestībām palika, to noklusēšana tikai veda pie mītu rašanās.

Runājot par gūstā nonākušajiem sarkanarmiešiem, nevar neminēt arī nošaušanu bez tiesas un izmeklēšanas. Poļu-padomju kara laikā netika ievēroti starpvalstu karos pieņemtie noteikumi, tai bija raksturīga šķiru naida asas izpausmes no abām pusēm. Gūstā saņemtos poļu virsniekus sarkanarmieši nošāva uz vietas, tāpat darīja poļi ar gūstā kritušajiem boļševiku komisāriem. Tādi fakti bija vērojami abās frontes pusēs. Taču par to mērogiem gandrīz neko nevar teikt, jo dokumentu par to praktiski nav, ir tikai atsevišķu aculiecinieku liecības. Ka padomju pusē bija šādi gadījumi, liecina rakstnieka Ī.~Bābela tapusī ``Jātnieku armijas dienasgrāmata'', uz kuras pamata viņš arī vēlāk uzrakstīja stāstu krājumu ``Jātnieku armija'' (\rutxti{Конармия}). Ka tas pats notika arī poļu pusē atzina J.~Pilsudska personīgais sekretārs 1918.--1922.~gadā K.~Svitaļskis: ``Boļševiku armijas dezorganizācija ar dezertēšanu ir apgrūtināta ar mūsu karavīru īstenoto nikno un nežēlīgo karagūstekņu iznīcināšanu''.

Pēc kāda sarkanarmieša (V.~Valujeva), kurš gūstā nonāca 18.~augustā pie Novominskas (\pltxti{Mińsk Mazowiecki}, apriņķa pilsēta netālu no Varšavas) sagūstītāji ``no visa sastāva izvēlējās komunistus, komandējošo sastāvu, komisārus un ebrejus, pie tam turpat visu sarkanarmiešu acu priekšā vienu komisāru ebreju piekāva un pēc tam nošāva.'' Viņš arī liecināja, ka visiem sagūstītajiem atņēma formas tērpus, bet kuri uzreiz neizpildīja pavēli to atdot, poļu leģionāri piekāva līdz nāvei. Tāpēc nenākas brīnīties, ka poļu avīzes bieži rakstīja par skrandās tērptajiem padomju gūstekņiem.

Krievu publicisti J.~Prudņikova un I.~Čigirins raksta: ``Ir jau saprotams~--- valsts tikko izveidojusies, tās teritorijā notiek karš, ir grūtības ar pārtiku, pie tam vēl fenomenālais poļu bardaks. Bet Antante taču ziedoja Polijai 3~miljonus karavīru formas tērpu! Tas četras reizes pārsniedza poļu karaspēka vajadzības. Neviens nesaka, ka gūstekņus jāietērpj jaunos mundieros, taču ko~--- otrā nēsāšanas termiņa apģērbu viņiem nevarēja piešķirt? Vēl jo vairāk tāpēc, ka izģērba gūstekņus poļu karavīri, par kuru rīcību valdībai taču vajadzētu nest kādu atbildību? Vai arī Antantes mundierus gaidīja tai pašā tirgū?''

Necilvēcīgi apstākļi valdīja arī Polijas ierīkotajās gūstekņu nometnēs. Poļu puse gan jau sākotnēji slēpa statistiku par nometnēs bojā gājušajiem sarkanarmiešiem, taču krievu vēstures literatūrā minēti daudzi jo daudzi fakti par gūstekņu piekaušanu, locekļu salaušanu, sagūstīto sieviešu izvarošanu, apstākļu radīšanu, kuri tieši veda pie gūstekņu nāves.

Jau kara laikā un tūlīt pēc tā Padomju Krievija pārmeta Polijai nežēlību pret gūstekņiem, to apzinātu iznīcināšanu. Par poļu attieksmi pret sarkanarmiešiem toreiz kā avīžu ``\rutxti{Правда}'' (Taisnība) un ``\rutxti{Известия}'' (``Ziņas'') korespondents padomju-poļu frontē strādājušais rakstnieks A.~Serafimovičs rakstīja: ``Drausmīgās spīdzināšanas un ņirgāšanās pār mūsu gūstekņiem izsauc šausmas poļu armijas ierindnieku vidū. Taču tās virsnieki gandrīz vienbalsīgi pavēl ``iznīcināt visiem līdzekļiem ``sarkanos suņus'', krievu ``okupantus''. Padomju pavēlniecības priekšlikumus par humānu apiešanos ar gūstekņiem un vietējiem iedzīvotājiem Varšava ignorēja. Mūsu puses mēģinājumi iesaistīt šai lietā kā starpniekus Tautu Savienību vai Polijas kaimiņus izrādījās bez panākumiem~--- baltpoļu obstrukcijas dēļ\dots{}'' Tiesa, izsauc šaubas avīzes korespondenta informētība par notiekošo otrā frontes pusē.

Taču arī Polijas teritorijā esošā prese apstiprināja augstāk zīmēto ainu. 20.~gada beigās fakti par necilvēcīgas izturēšanas un ņirgāšanās gadījumiem par padomju karagūstekņiem kļuva zināmi poļu presei un daļā no tās izsauca protestu. 1920.~gada 22.~decembrī Ļvovas avīzē ``\rutxti{Вперед}'' (``Uz Priekšu'') tika ziņots, ka 9.~decembrī Tuholas (\detxti{Tucheln}, \pltxti{Tuchola}, \rutxti{Тухоля}, \betxti{Тухола}, \uktxti{Тухоль})~--- nometnē vienā dienā nomira 45~krievu gūstekņi. Cēlonis bija tas, ka viņus salā un vējā ``pusapģērbtus un basām kājām'' veda uz pirti ar betona grīdu, bet pēc tam izmitināja zemnīcās bez grīdas. Publikācijas rezultātā poļu varas iestādes bija spiestas veikt izmeklēšanu, kura, itin ``dabīgi'', atspēkoja avīzes ziņojumu. Komisija noskaidroja, ka tai dienā no izsitumu tīfa nomiruši 10 cilvēku, pirts bijusi sasildīta un gūstekņi pēc mazgāšanās izvietoti pirms tam dezinficētās barakās. Avīze ``\rutxti{Вперед}'' pēc izmeklēšanas rezultātiem tika uz nenoteiktu laiku slēgta par ``pārspīlētu un tendenciozu ziņu publicēšanu.'' Zinot gūstekņu dzīves vispārējos apstākļus, autors atļaujas neuzticēties oficiālās komisijas ziņojumam. Runāt par kādu starptautisko normu ievērošanu attiecībā pret padomju karagūstekņiem nevar. Kad padomju delegācija miera sarunās Rīgā mēģināja pievērst Polijas delegācijas uzmanību gūstekņu turēšanas neatbilstībai 1907.~gada Hāgas konvencijai ``Par sauszemes kara likumiem un ierašām'', tā neguva panākumus. KSFPR ārlietu tautas komisāra G.~Čičerina 1921.~gada 9.~septembra memorandā Polijas sūtniecībai Maskavā bija teikts, ka ``divu gadu laikā no 130~000 krievu gūstekņiem Polijā nomiruši 60~000''. Krievijas un Ukrainas delegācija ar Ukrainas ārlietu tautas komisāru D.~Manuiļski priekšgalā, kura 1921.--1923.~gadā mēģināja noskaidrot internēto likteni, savā ziņojumā rakstīja, ka viņus ``poļu iestādes uzskatīja par beztiesīgiem vergiem \citespace{} dažās nometnēs viņiem lika zirgu vietā pārvadāt pašu fekālijas. Sodi izceļas ar barbarisku nežēlību. Sarunas šajos jautājumos sašķīda pret poļu delegācijas sabotāžu.''

Padomju karagūstekņus turēja vairākās nometnēs. Daļa krievu vēsturnieku tās nereti sauc par koncentrācijas nometnēm. No poļu un krievu ``demokrātu'' puses skan iebildumi, ka gūstekņu nometnes nav korekti saukt par koncentrācijas nometnēm, kaut smagos uzturēšanās apstākļus tajās viņiem nākas atzīt. Dzīves apstākļi nometnēs bija necilvēcīgi, par gūstekņiem ņirgājās, jūdza ratos, arklos, ecēšās. 1993.~gadā Krievijas periodikā (\rutxti{Военно-исторический журнал}) bija publicēti dati, ka padomju gūstekņus poļu šautuvēs izmantoja par dzīviem mērķiem. Mirstība nometnēs sasniedza līdz 20\%. Galvenais nāves cēlonis bija epidēmijas, kas vairāk nekā pieticīgās ēdināšanas, cilvēku sablīvētības un medicīniskās palīdzības trūkuma rezultātā veda pie augsta letāla iznākuma līmeņa. Piemēram, 1920.~gada 2.~decembrī Polijas Augstākais ārkārtējais komisārs cīņā pret epidēmijām E.~Godļevskis vēstulē kara ministram K.~Sosnkovskim ziņoja par situācijas pārbaudi karagūstekņu sadales stacijā Pulavos (\pltxti{Puławy~--- pilsēta Polijas austrumos}). Tajā 3.~novembrī no Koveļas (ukr. \uktxti{Ковель}~--- pilsēta Ukrainas ziemeļrietumos) pa dzelzceļu ieradās 700~karagūstekņu, izbadējušies tiktāl, ka liela daļa no viņiem nespēja patstāvīgi izkāpt no vagoniem. 15~cilvēku pirmajā dienā pēc atbraukšanas nomira. Noskaidrojās, ka vilciens ceļa pavadīja 4~diennaktis, kuru laikā gūstekņus vispār nebaroja. Pie tam ar to pašu vilcienu tika vesta gaļa, paredzēta gūstekņu ēdināšanai (Gaļas tālākais liktenis palika neskaidrs). Pulavos gūstekņus baroja tā, ka pēc E.~Godļevska vārdiem, viņiem bada dēļ ``bija jāizmirst 111~dienu laikā.'' Sadales stacijas vadītājs sūdzējās Augstākajam ārkārtējam komisāram, ka ``nepanesamie gūstekņi'' pastāvīgi apēdot kartupeļu mizas ``no mēslu kaudzes, kura atrodas nometnē, un to nāksies apjozt ar dzeloņstieplēm''. Augstākais ārkārtējais komisārs stāvokli nometnēs raksturoja kā ``vienkārši necilvēcīgu un pretrunā esošu ne tikai higiēnas, bet vispār kultūras prasībām''. Var piebilst, ka sešos 1920.~gada rudens~--- 1921.~gada pavasara mēnešos no 1~100 gūstekņiem Pulavos nomira 900.

Liela daļa gūstekņu strādāja valsts uzņēmumos vai pie privātpersonām, galvenokārt lauksaimniecībā. Miera sarunās Rīgā Polijas puse plānoja padomju pusei piestādīt rēķinu par sagūstīto sarkanarmiešu uzturēšanu. Taču padomju diplomāti pēc ap 3~000 jau atgriezušos gūstekņu aptaujas izrēķināja, ka gūstekņu paveiktā darba vērtība Polijā ievērojami pārsniedza viņu uzturēšanai izmantotos līdzekļus, tāpēc poļu delegācija šo jautājumu vairs necilāja.

Pēc krievu un baltkrievu emigrantu datiem \strong{padomju gūstekņu skaits Polijā} bija ne mazāks par 73~000.

Poļu pētnieki kopējo padomju karagūstekņu skaitu vērtēja kā 80~000~līdz 110~000 lielu (šādus datus attiecībā uz 1920.~gada 18.~oktobri minēja jau J.~Pilsudskis grāmatā ``1920.~gads'', kad M.~Tuhačevskis savukārt ziņoja par 95 tūkstošiem pazudušajiem bez vēsts un nokļuvušajiem gūstā). Iespējams, ka atšķirtības skaitļos ietekmējis arī tas, ka 1920.~gada augustā-oktobrī Polijas karaspēks mēdza kaujās laukā atstāt bez palīdzības ievainotos sarkanarmiešus, tai pat laikā ierakstot tos atskaitēs kā sagūstītus, ko rāda baltkrievu vēsturnieku publicētie dokumenti. Arī mūsdienu poļu vēsturnieks Z.~Karpuss kopējo krievu-ukraiņu karagūstekņu skaitu vērtē kā 110~000 lielu.

Padomju un daži mūsdienu krievu vēsturnieki min datus par 157--165 tūkstošiem (dažkārt publicisti runā pat par 200 tūkstošiem) gūstekņu un no 40~līdz 80 tūkstošiem bojā gājušo.

Piemēram, krievu vēsturniece I.~Mihutina uzskata, ka 1919.--1921.~gadā poļu gūstā nomira 83,5 tūkstošu krievu karavīru. 1998.~gada Krievijas Federācijas Ģenerālprokuratūras vēstulē Polijas Prokuratūrai tika norādīts, ka poļu karagūstekņu nometnēs no nežēlīgas apiešanās un necilvēcīgiem dzīves apstākļiem bojā gāja vairāk nekā 80~000 gūstā kritušo sarkanarmiešu. Ja dati par ap 200~000 gūstekņiem un 80~000 bojā gājušajiem ir patiesi, tad paliek neskaidrs vēl ap 50~tūkstošu liktenis. Polijas valdība gan noraidīja šādus datus.

Jau pieminētais krievu vēsturnieks M.~Meļķuhovs 2001.~gadā izdotajā grāmatā ``Padomju--poļu kari. Militāri politiskais pretstatījums 1918--1939 gg.'' (\rutxti{Советско-польские войны. Военно-политическое противостояние 1918--1939 гг.}) rakstīja, ka 1919--1920.~gadā ``kara laikā poļu karaspēks saņēma gūstā vairāk nekā 146~tūkstošus cilvēku, kuru uzturēšana Polijā bija ļoti tālu no kādiem humāniem standartiem. Īpašas ņirgāšanās bija jāizcieš komunistiem vai aizdomās turamajiem par piederību viņiem, bet gūstā saņemtie sarkanarmieši-vācieši vispār tika nošauti uz vietas. Taču arī parastie gūstekņi bieži kļuva par poļu militāro iestāžu patvaļas upuriem. Plaši bija izplatīta gūstekņu aplaupīšana, ņirgāšanās par sagūstītajām sievietēm. Acīmredzot tāda attieksme pret padomju karagūstekņiem ievērojamā mērā bija sekas ilgus gadus vestajai propagandai par Krievijas ``vainu'' Polijas priekšā. Tas viss noveda pie tā, ka ap 60~tūkstošu padomju karagūstekņu nomira poļu nometnēs. Uz 1921.~gada 21.~novembri no Polijas atgriezās 75~699 bijušie gūstekņi (932 cilvēki atteicās atgriezties).'' Taču minētie skaitļi nav vispāratzīti. Vēsturnieks G.~Matvejevs uz Polijas armijas Ģenerālštāba dokumentu (Polijas armijas Virspavēlniecības Operatīvās nodaļas atskaišu) izpētes pamata noskaidrojis, ka 20~kara mēnešos (no 1919.~gada 13.~februāra līdz 1920.~gada 18.~oktobrim) gūstā nokļuva 206~877 sarkanarmieši, no kuriem kara ministrijas rīcībā tika nodoti un reāli izjuta nebrīves smagumu 157~tūkstoši. Šī vēsturnieka un viņa meitas V.~Matvejevas 2011.~gadā izdotā grāmatā norādīts, ka no gūsta mājās atgriezās 75~699 cilvēki (kas sakrīt ar M.~Meļķuhova minētajiem datiem), apmēram 25-29 tūkstoši mira no bada un aukstuma, citi pārgāja dienēt dažādos pretpadomju formējumos.

Ir pieejami arī citi dati. Pēc dažiem jaunākajiem poļu, vācu, lietuviešu un krievu vēsturnieku vērtējumiem no 120~000~līdz 130~000 internētajiem padomju civiliedzīvotājiem un 1919.--1920.~gada kara kampaņās gūstā kritušajiem sarkanarmiešiem (pēc citām ziņām Polijā bija 66~000 karagūstekņu un 50~000 internēto civiliedzīvotāju) pēc vienošanās par kara gūstekņu un internēto apmaiņu parakstīšanas 1921.~gada 24.~februārī 65~000--70~000 cilvēku atgriezās KSFPR. Pie ``baltajiem'', tai skaitā Polijā, Vācijā, Rumānijā, Baltijā, pārgāja ne mazāk kā 6~tūkstoši. Uz pastāvīgu dzīvi Polijā un tās kaimiņvalstīs (izņemot KSFPR-PSRS) palika apmēram 1,5~tūkstoši. Nometnēs bojā gāja mazākais 50~000~cilvēku.

Vairākumā poļu publikāciju par 1919--1920.~gada karu padomju karagūstekņu bojā eja izskaidrota ar objektīviem, no varas neatkarīgiem apstākļiem~--- daudzajos kara gados sagrauto Polijas tautsaimniecību, valstī valdošo vispārējo medicīniskā personāla un zāļu, pārtikas, apģērba, kurināmā, dzīvei piemērotu telpu utt. trūkumu.

Lūk, tikai viena liecinieka (A.~Mackeviča) stāstītais 1920.~gada 5.~martā: ``\dots{}ēdiens nometnēs tiek izsniegts tāds, ka neviens normāls cilvēks nespēs nodzīvot vairāk vai mazāk ilgstošu laiku. Sastāv tas no nelielas rupjas maizes porcijas, svarā līdz 1/2 mārciņai, un kausa viras dienā. Līdzīgas drīzāk samazgām, nekā virai. Vairumā gadījumu šīs samazgas, ko sauc par viru, tiek izdalītas nesālītas. Uz aukstuma un bada pamata slimības izplatās līdz neticamiem apjomiem. Medicīniskās palīdzības nekādas\dots{} Desmiti cilvēku mirst ik dienas.''

Protams, šie apstākļi veicināja dažādu epidēmiju izcelsmi un augsto mirstību karagūstekņu nometnēs. Taču diemžēl jārunā arī par poļu sabiedrībā uzkrājušos naidu pret ``krievu'' apspiedējiem, kuru nācās izjust ne tikai jaunajiem apspiedējiem~--- boļševikiem, bet arī viņu armijā mobilizētajiem vienkāršajiem krievu un arī citu tautu iedzīvotājiem, kuri personīgi nebija vainojami nedz par Krievijas caru realizēto politiku Polijā, nedz par boļševiku pasaules revolūcijas plāniem, kuru ceļā atradās arī poļu zemes. Krievu vēsturniece J.~Jakovļeva, atsaucoties uz krievu gūstekņu atmiņām, konstatējusi, ka vienāda poļu attieksme bija kā pret sarkanarmiešiem, tā baltgvardiem~--- poļiem tie pirmkārt bija krievi, tātad~--- ``mūžīgie ienaidnieki''. Neticis ņemts vērā arī tas, ka ģenerālis A.~Deņikins bija sekmējis ne tikai poļu bēgļu atgriešanos Polijā, bet arī palīdzējis nosūtīt uz Poliju no poļu bēgļiem Krievijas dienvidos saformētu brigādi. Poļos ieaudzinātais naids pret visu krievisko bijis visaptverošs. Poļu nometnēs nonākušie krievu baltgvardi liecināja, ka poļi, kuri ienīda veco, carisko Krieviju, šo naidu pārnesa arī uz baltgvardiem. (Pēc poļu zinātnieka Z.~Karpusa datiem 1921.~gada februārī Polijā atradās gandrīz 25~tūkstoši internēto krievu ``balto'' armiju karotāju). Par krievu ``balto'' virsnieku un karavīru nožēlojamo stāvokli poļu nometnēs 1920.~gada 21.~decembrī vēstulē J.~Pilsudskim rakstīja viens no ``balto'' līderiem, nerimtīgais boļševiku pretinieks B.~Savinkovs, kurš šai laikā dzīvoja Polijā un vadīja tur ``Krievu politisko komiteju'' (``\rutxti{Русский политический комитет}''). Tas vēlreiz apliecina poļu politikas antikrievisko, bet ne tika antiboļševistisko virzību. Tiesa, krievu vēsturnieks N.~Rajskis norādījis, ka Polijā gūstekņi tika dalīti 4 kategorijas, kur vissliktākos apstākļos atradās sarkanarmieši~--- komunistiskās partijas biedri. Taču gūstekņu subjektīvajās atmiņās šīs atšķirības varēja arī neizpausties, tik briesmīga dzīve nometnēs bija tiem visiem.

Ir zināmi fakti, ka jau ceļā uz nometnēm gūstekņu vidū mirstība sasniedza līdz 40\%. Ja nometnēs nonāca ap 150~tūkstoši, tad, pat ierēķinot apstiprinātas ziņas par to, ka ap 12~tūkstošu ukraiņu tika atlaisti mājup un vairākus tūkstošus sagūstīto poļiem atkaroja Sarkanā armija, tas nozīmē, ka ap 30~tūkstošu liktenis nav noskaidrots un vismaz daļa no tiem jāuzskata par gājušiem bojā.

Nenākas šaubīties, ka tie, kam izdevās izdzīvot poļu nometnēs, uz visu mūžu atcerējās piedzīvoto un nevarēja to piedot.

Ģenerālis, vēlāk vēstures zinātņu doktors N.~Antoņenko stāstīja par 1944.~gadā piedzīvotu gadījumu. Kad Sarkanā armija iegāja Polijas teritorijā, I.~Baltkrievijas frontē, kurā viņš dienēja, 18~poļu-padomju kara laikā poļu karagūstekņu nometnēs bijušo ieslodzīto sarkanarmiešu vēlējušies nokļūt vietās, kur viņi bija ieslodzīti un zināja, ka ir vēl dzīvi daži no nometņu administrācijas pārstāvjiem. Viņus nācies ``ar spēku nomierināt, lai atriebības jūtas neizpaustos par izrēķināšanos ar poļiem''. Poļu puse tad solījusi, ka tiks noskaidroti visi bojā gājušie sarkanarmieši, tiks uzcelti viņiem pieminekļi, uzcelts dievnams, kur izteikt nožēlu \dots{} Taču vēlāk tas tika aizmirsts.

Saskaņā ar J.~Staļina un toreizējā Polijas līdera B.~Beruta kopīgu lēmumu, pieņemtu 1951.~gadā, bija paredzēts ne vēlāk kā 1953.~gadā, Tuholā, kur 1920.~gadā atradās no lielākajām nometnēm padomju ieslodzītajiem, izveidot attiecīgu memoriālu. Taču celtniecības darbi Tuholā, kas bija sākti 1952.~gadā, tika apturēti 1953.~gada 7.~martā, divas dienas pēc ziņām par Staļina nāvi.

1998.~gadā Polijas ģenerālprokurore un justīcijas ministre H.~Suhocka, atbildot uz Krievijas Ģenerālprokuratūras vēstuli ar lūgumu izmeklēt 82,5 tūkstošu Sarkanās armijas karavīru nāves cēloņus, paziņoja, ka ``izmeklēšanas par it kā notikušo boļševiku gūstekņu iznīcināšanu, ko no Polijas pieprasa Krievijas ģenerālprokurors, nebūs.''. Ģenerālprokurore gan neapšaubīja karagūstekņu bojā ejas faktu. Savu atteikumu viņa pamatoja ar poļu vēsturnieku ``neapšaubāmu'' konstatējumu, ka 16--18~tūkstošu karagūstekņu nāves cēlonis bija ``vispārējie pēckara apstākļi'' un piebilda, ka par ``nāves nometnēm'' un gūstekņu ``iznīcināšanu'' Polijā nevarot būt ne runas un ``nekādas speciālas darbības, vērstas uz gūstekņu iznīcināšanu, netika veiktas''. No tā laika oficiālā Polijas valdības nostāja nav mainījusies.

Vairāki krievu publicisti šādu mūsdienu Polijas vadītāju reakciju uz vēsturnieku atrastajiem faktiem par padomju karagūstekņu likteni Polijā salīdzināja ar PSRS vadības nevēlēšanos gadiem atzīt valsts atbildību par Katiņas (skat. šī darba III un IV nodaļu) noziegumu. Šāda ignorance izsauca tādu pat sašutumu Krievijas sabiedrībā, kādu Polijas sabiedrībā izraisīja PSRS vadības nevēlēšanās paskatīties faktiem acīs Katiņas sakarā. Jāatzīst, ka ar cieņu uzlūkojamā Polijas puses vēlme noskaidrot līdz galam Katiņas noziegumu būtu izpelnījusies vēl lielāku krievu sabiedrības atbalstu, ja Polija ar tādu pat dedzību palīdzētu Krievijas pusei noskaidrot 1919.--1921.~gada karagūstekņu likteni. Ne velti krievu vēsturnieks N.~Rajskis raksta, ka, ja pēc vairākām desmitgadēm tika atzīta staļiniskā režīma atbildība par poļu virsnieku apslepkavošanu Katiņā, citiem poļu upuriem, tad krievu pusei ir tiesības gaidīt atklātību no poļu puses, vēl jo vairāk tāpēc, ka no Krievijas gūstekņu bojā ejas poļu nometnēs ir pagājis daudz lielāks laiks.

Jāsaka, ka mainās arī poļu sabiedriskā doma. 2000.~gadā poļu publicists B.~Skaradzinskis konstatēja: ``Diemžēl [gūstekņu dzīves] apstākļi tiešām bija drausmīgi, veda pie augstas mirstības, un es nedomāju, ka pie mums bija izdarīts viss, lai traģēdijas mērogus novestu līdz minimumam. Mēs vieglprātīgi lielījāmies ar to, cik tūkstošu krievu ir saņemts gūstā, bet ne statistiskajās atskaitēs, ne avīzēs neviens neieminējās par to, cik starp viņiem ir smagi slimu un smagi ievainotu, bet tādi taču iet bojā pirmām kārtām. Poļu apziņā nepastāv pati problēma par mūsu atbildību par šiem cilvēkiem. Mēs viegli atlaižam sev grēkus uz tā pamata, ka mūsu karagūstekņus Krievijā, tai skaitā slimos un ievainotos, 1920.~gadā arī ne pārāk ieredzēja''. Šis publicists savā laikā Varšavas apkārtnē ilgstoši meklēja, bet tā arī neatrada 1920.~gadā kritušo sarkanarmiešu kapus, jo par tiem neviens savlaicīgi neparūpējās. Starpkaru periodā valdošajam slānim Polijā pat prātā nevarēja ienākt rūpēties par kritušo pretinieku atdusas vietām. Pēc Otrā pasaules kara par apbedījumiem nerūpējās arī propadomiskās Polijas varas iestādes, acīmredzot nevēloties atcerēties 1920.~gada notikumus.

Diemžēl nākas konstatēt, ka daudziem mūsdienu Krievijas politiķiem jautājums par poļu-padomju karā kritušo piemiņu lielākoties ir kļuvis vienīgi par politiskās cīņas līdzekli, cilvēciski-morālo pusi atstājot novārtā. Nav dzirdēts, ka kāds no mūsdienu Krievijas vadītājiem līdz 2010.~gadam būtu izteicis vēlēšanos apmeklēt poļu nometnēs mirušo tautiešu apbedīšanas vietas, izbūvēt tur piemiņas zīmes. Tiesa, 2010.~gadā notika starpgadījums, kurš ir pelnījis, lai par to pastāstītu sīkāk.

Līdz 2008.~gadam Osuvas (neliela pilsētiņa Varšavas tuvumā) vietējie iedzīvotāji pie 1920.~gadā kritušu 22~sarkanarmiešu kapa, atzīmēta ar pareizticīgo krustu, mirušo piemiņas dienā nolika ziedus un aizdedza sveces. Tad ar Polijas piemiņas un mocību aizsardzības padomes (\pltxti{Rada Ochrony Pamięci Walk i Męczeństwa}) iniciatīvu tika nolemts uzcelt nelielu memoriālu. Ideju radīt pieminekli, kuram bija jākļūst par saprašanās un attiecību normalizācijas simbolu ar Krieviju, atbalstīja toreizējais Seima priekšsēdētājs B.~Komorovskis, vēlāk Polijas prezidents. Piemineklis tika izgatavots no nerūsējoša tērauda pareizticīgo krusta formā, kuru aptver 22~asi elementu, atgādinoši durkļus. Uzraksts uz memoriālās plāksnes vēstīja, ka šeit atdusas 22 Sarkanās armijas karavīri, krituši 1920.~gada 14.~un 15.~augustā. Tēlnieks M.~Moderau, vairāku sabiedrībā atzītu skulptūru, tai skaitā 2010.~gada 10.~aprīlī aviakatastrofā pie Smoļenskas bojā gājušā Polijas prezidenta L.~Kačinska un viņa pavadoņu piemiņai veltītā pieminekļa militārajā kapsētā Varšavā, autors, krusta kā simbola izvēli izskaidroja ar to, ka kritušie bija audzināti pirms revolūcijas, kaut krita, nesot jau citu simboliku. (Starp citu, brāļu kapu pirms pārapbedīšanas atrokot, tika atrasts uz krūtīm nēsājams krustiņš.

Memoriāls tika uzbūvēts, 2010.~gada 15.~augustā, Polijas armijas un Varšavas kaujas atceres dienā, bija jānotiek tā atklāšanai. Bija paredzēts, ka to apmeklēs arī Krievijas pārstāvji. Taču atklāšana nenotika poļu radikālnacionālistu vainas dēļ. Jau pirms tam izraisījās diskаusija. Pret pieminekļa atklāšanu protestēja radikālnacionālisti, norādot, ka kritušajiem boļševiku karavīriem neesot bijis nekā kopīga ar pareizticīgo krustu, pieminekļa mets neesot saskaņots ar sabiedrību, iebilda pret durkļu simboliem. Naktī pirms atklāšanas uz pieminekļa tika uzkrāsotas trīs sarkanas zvaigznes, dienā ieradās nacionāl-konservatīvās un katoliski-klerikālās partijas \pltxti{Prawo i Sprawiedliwość} (``Taisnība un taisnīgums'') delegācija ar protesta transparantiem. Pilsētas mērijas pārstāvis, ieraugot apgānīto pieminekli, sašutis izsaucās: ``Kāds kauns. Tā nebija [vietējo] iedzīvotāju akcija. Mēs cienām kapus.'' Piemiņas un mocību aizsardzības padomes ģenerālsekretārs vēsturnieks A.~Kunerts paguva vienīgi dažu simtu metru attālumā no notikuma vietas apstādināt oficiālo Krievijas sūtņa Polijā vadīto delegāciju, kura parādīja sapratni un atgriezās mājup, lai nesaasinātu incidentu. Nākamajā dienā pilsētiņā kāds bija izlīmējis ``Taisnības un taisnīguma'' konkurenta Polijas premjerministra D.~Tuska attēlus padomju mundierī. Presē un Internetā ``Taisnības un taisnīguma'' biedri protestēja pret to, ka par viņu naudu tiek izbūvēti boļševiku karavīru apbedījumi. Lai kādi attaisnotu vandālismu avīzē ``\pltxti{Rzeczpospolita}'' (``Žečpospolita'') vēsturnieks, profesors A.~Novaks pēc garāka apcerējuma par Polijas un Krievijas vēsturiskajām attiecībām iebilda pret pareizticīgo krusta attēlošanu piemineklī kritušajiem sarkanarmiešiem, jo nevarot taču zināt, ka tie visi bijuši pareizticīgie. Tā esot pēcnāves iespiešana vienā kategorijā. Lai nu kā, pilsētas burgomistrs nosūtīja prokuratūrai iesniegumu par izdarīto noziegumu~--- kapu apgānīšanu, bet Piemiņas un mocību aizsardzības padomes ģenerālsekretārs A.~Kunerts paziņoja, ka piemineklis karavīriem tiks notīrīts, restaurēts un notiks arī tā atklāšanas ceremonija. Intervijā avīzei ``\pltxti{Gazeta Wyborcza}'' viņš uzsvēra: ``Katram karavīram ir tiesības uz kapavietu''. Jāatzīst, ka reakcija krievu preses izdevumus bija bez eksaltācijas, iespējams gan tāpēc, ka 2010.~gada 10.~augustā bija gājis bojā Polijas prezidents Ļ.~Kačinskis. Vietējās varas iestādes pieņēma lēmumu, kuram piekrita arī Polijas piemiņas un mocību aizsardzības padome, noņemt durkļus atgādinošos elementus. Oktobrī piemineklis tika rekonstruēts. Atklāšanas ceremonija notika 2010.~gada 2.~novembrī, tajā nepiedalījās neviens Krievijas pārstāvis, bet piketēja neliela protestējošo grupa ar transparantiem. Tadejādi labi iecerētais pasākums tā arī nekļuva par poļu un krievu samierināšanās un savstarpējās izpratnes simbolu.

Kā ievērojams sasniegums saprašanās meklējumos vērtējams XXI gadsimta sākumā Krievijas Federālās arhīvu aģentūras un virknes Krievijas Federācijas arhīvu noslēgtais divpusējais līgums ar Polijas Valsts arhīvu Ģenerāldirekciju par problēmas izpēti. Tā rezultātā 2004.~gadā iznāca krājums ``Sarkanarmieši poļu gūstā 1919.--1922.gg.'' (\rutxti{Красноармейцы в польском плену в 1919--1922 гг. Сб. документов и материалов. Москва~--- Санкт-Петербург, Летний сад, 2004. 912~с.}). Abu~--- Krievijas un Polijas~--- pušu vēsturnieku darba rezultātā ir noskaidrots, ka dokumenti neapstiprina to, ka poļu varas iestādes mērķtiecīgi realizēja padomju gūstekņu iznīcināšanas politiku. Līdz tam krievu puses minētais mirušo sarkanarmiešu skaits bija pārspīlēts. Krājumā publicētie materiāli rāda, ka Polijas varas iestādes griezās dažādās Eiropas instancēs ar lūgumu palīdzēt karagūstekņu nometņu apgādē. Taču sagrautās karā Eiropas valstis varēja sniegt tikai ierobežotu palīdzību. Poļu iestādes piedāvāja daudzus gūstekņus padomju pusei, taču padomju iestādes atteicās tos pieņemt pirms būs parakstīts miera līgums.

Tomēr atšķirības vērtējumos paliek. Krājumā kopējo padomju karagūstekņu skaitu poļu vēsturnieki vērtē kā 80--85~tūkstošus, krievu~--- kā 157~tūkstošus lielu. Nometnēs mirušo skaits, pēc poļu puses datiem bijis 16--17 tūkstoši, pēc krievu~--- vismaz 18--20 tūkstošu. Krājuma izdošana neapklusināja dažādus viedokļus. Krievu politiķis un publicists V.~Šveds raksta, ka ir ``absolūti skaidrs, ka poļu gūstā bojā gājušo sarkanarmiešu skaits, ko krājumā \citespace{} piedāvā poļu un daži krievu vēsturnieki, ir acīmredzami samazināts''. Arī citi krievu vēsturnieki pēc krājumā ievietotajiem dokumentiem secina, ka tikai divās nometnēs Stšalkovā (\pltxti{Strzałow}) un Tuholā (\pltxti{Tuchola}), mira vairāk nekā 30~tūkstošu padomju karagūstekņu. Tā, pētniece T.~Simonova apstrīd poļu vēsturnieka Z.~Karpusa secinājumus par mirušo skaitu Tuholas nometnē, kuri veikti uz nometnes garīdznieka sastādīto nāves aktu un kapsētā apglabāto sarakstu pamata, norādot, ka garīdznieks nevarēja noturēt aizlūgumus par mirušajiem komunistiem un, pēc aculiecinieku atmiņām, mirušos apbedīja brāļu kapos.

Minētajā krājumā atrodams, ka abu pušu vēsturnieki, izmantojot dažādus avotus, pat mainās vietām iespējamo upuru un izdzīvojušo skaita noteikšanā. Tā, ja Z.~Karpuss uzskaitīja, ka Krievijā pēc 1921.~gada 24.~februāra vienošanās atgriezās 66~762 cilvēki, tad pēc Sarkanās armijas štāba Mobilizācijas pārvaldes sīkas izziņas no poļu gūsta atgriezās 75~699 sarkanarmieši. Tātad attiecīgi upuru īpatsvars samazinājās. Šī skaitļu atšķirība liecina, ka pat bez politiskiem uzstādījumiem pastāv pietiekami daudz objektīvu grūtību (avotu pretrunīgums u.~c.), traucējošu saskaņot viedokļus uz zinātniska pamata.

Krājuma priekšvārdā G.~Matvejevs raksta arī par tiem pāri par 60~tūkstošu sarkanarmiešiem, kuri nenomira gūstā, bet nerepatriējās arī uz Padomju Krievijā. Viņš uzskata par noskaidrotu ap 53~tūkstošu likteni. Daļa no tiem nonāca dažādās armijās, kuras karoja pret ``sarkanajiem''. (Poļu vēsturnieks Z.~Karpuss norāda, ka uzreiz pēc nokļūšanas gūstā dažādos ``balto'' krievu, ukraiņu u.~c. formējumos, kuri karoja pret Sarkano armiju, iestājušies ap 25~tūkstoši gūstekņu.) Daļa karoja poļu pusē, daļu atbrīvoja Sarkanās armijas pretuzbrukumi, daļa rietumukraiņu un rietumbaltkrievu tika atlaista vai paši aizbēga mājup, daļu poļi, labi paēdinot un apgādājot ar proklamācijām izdalīšanai sarkanarmiešu vidū, atbrīvoja propagandas nolūkos. Daļa nesen neatkarību ieguvušo valstu, tai skaitā Latvijas, iedzīvotāju, kuri pret viņu gribu bija mobilizēti Sarkanajā armijā, tiešā ceļā atgriezās dzimtenē. Tādejādi poļu nometnēs bojā gājušo skaits it kā sarūk.

Poļu vēsturnieki vēl un vēlreiz uzsver, ka atšķirībā no Katiņas gadījuma, kur par poļu gūstekņu iznīcināšanu 1940.~gadā ir pieejami dokumenti, nav nekāda dokumentāla pamata runāt par Polijas valdības un militārās pavēlniecības mērķtiecīgu krievu gūstekņu iznīcināšanas politiku. (Šeit gan jāsaka, ka teorētiski dokumentu trūkums var nozīmēt, ka tie tikuši mērķtiecīgi iznīcināti, vai arī vispār nav bijuši nepieciešami. Tā, līdz šim brīdim taču nav atrasti Ā.~Hitlera parakstīti dokumenti par miljonu ebreju iznīcināšanu, taču domu, ka holokausts varēja notikt bez fīrera norādēm, nopietni vēsturnieki neizsaka.)

Tomēr, ja arī valsts un armijas vadība netika devušas norādījumus īpaši nežēlot gūstā nonākušos pretinieka kareivjus, tad pierādījumi, ka attieksme pret tiem bija necilvēcīga, tomēr ir atrodami. Ka sadzīves stāvoklis un aprūpe nometnēs bija slikti, atzina arī speciāla Polijas Seima komisija 1920.~gada martā. Arī turpmāk situācija neuzlabojās. Brestļitovskas nometnes komandants 1920.~gada rudenī ievestajiem gūstekņiem atklāti paziņoja: ``Jūs, boļševiki, gribējāt mums atņemt mūsu zemi,~--- labi, es jums došu zemi. Nogalināt jūs man nav tiesību, bet es jūs tā barošu, ka jūs paši nosprāgsiet.'' Diemžēl tāda poļu amatpersonu izturēšanās bija ne tikai pret sarkanarmiešiem. Krievu publicists A.~Pamjatnihs ir minējis faktu, ka Stšalkovas nometnes, kura izcēlās ar ļaundarībām un nežēlīgu attieksmi pret gūstekņiem, (par ko padomju pārstāvji protestēja pat vēl 1922.~gada janvārī), komandants tika par to arestēts un nodots tiesai. Polijas armijas Ģenerālštāba pulkvedis I.~Matuševskis vēstulē ģenerālim K.~Sosnkovskim par gūstekņu bēgšanu no nometnēm rakstīja: ``Šo bēgšanu izsauc apstākļi, kuros atrodas komunisti un internētie (malkas, veļas un apģērba trūkums, sliktā ēdināšana, ilgā gaidīšana uz izbraukšanu uz Krieviju). Īpaši slavena ir kļuvusi nometne Tuholā, kuru internētie un karagūstekņi sauc par ``nāves nometni'', kur jau gājuši bojā ap 22~tūkstošiem sarkanarmiešu.'' Poļu vēsturnieks A.~Veļevejskis ar norādēm uz J.~Pilsudska laika dokumentiem, 1994.~gadā avīzē ``\pltxti{Gazeta Wyborcza}'' rakstīja par poļu militāristu pavēlēm nošaut simtus padomju karagūstekņu, neņemt tos gūstā, nogalināt smagi ievainotos.

Krievu publicists V.~Šveds uz minētajā krājumā ``Sarkanarmieši poļu gūstā 1919--1922 gg.'' publicēto dokumentu pamata secina: ``Dokumenti un liecības ar lielu pārliecības pakāpi ļauj apgalvot par plānveidīgu un burtisku sarkanarmiešu iznīcināšanu ar aukstumu un badu poļu karagūstekņu nometnēs. Tāpat var formulēt secinājumu, ka Polijā krievu gūstekņu bojā ejas jau iepriekšējo nolemtību noteica nevis augstāko varas iestāžu lēmumi, bet vispārējais pretkrieviskais poļu sabiedrības noskaņojums~--- jo vairāk boļševiku nosprāgs, jo labāk.'' Ja secinājuma pirmo pusi, neraugoties uz piebildi ``ar lielu pārliecības pakāpi'' var uzskatīt par tendenciozu, tad otrajai tā puses ir objektīvs pamats. Citviet V.~Šveds vēl saasinājis tikko citēto secinājumu, rakstot, ka ``\dots{}sarkanarmiešu masu bojā eju poļu gūstā var vērtēt kā ar iepriekšēju nodomu izdarītu viņu iznīcināšanu. Ņemot vērā šo noziegumu skaidri izteikto nacionālo virzību, ir jāstāda jautājums par genocīdu attiecībā pret krievu un ebreju tautības gūstekņiem''. Jāsaka, ka iepriekšēja nodoma konstatēšanai tieslietu praksē ir vajadzīgi ļoti nopietni pierādījumi, kurus šajā gadījumā acīmredzot neizdosies atrast. Ar genocīda tēzes izvirzīšanu tiek ignorēti fakti par poļu-padomju kara ne tikai nacionālo, bet arī šķirisko komponenti, par to, ka Sarkanās armijas rindās karoja un poļu nometnēs mira visdažādāko tautību cilvēki. Patiesos publicista un arī politiķa V.~Šveda nolūkus atklāj viņa ``tiltiņš'' starp šiem un vēlākajiem notikumiem: ``Padomju valdības skarbie paziņojumi 1920.--1922.~gadā sakarā ar necilvēcīgo attieksmi pret gūstā saņemtajiem sarkanarmiešiem ļauj tāpat izdarīt secinājumu par to, ka padomju vadībai 1939.--1940.~gadā bija jāatgriežas pie poļu sagūstīto sarkanarmiešu bojā ejas jautājuma. Tas bija viens no cēloņiem, kāpēc daļa poļu karagūstekņu, kam bija kāds sakars ar represijām pret sagūstītajiem sarkanarmiešiem, 1940.~gada pavasarī tika represēti''. Citētie apgalvojumi nav pierādīti, tajos it kā ``vēsturiski'' secinājumi tiek klaji izmantoti staļiniskā režīma attaisnošanai.

Der minēt arī datus par \strong{latviešu gūstekņu stāvokli Polijā} 1919.--1920.~gadā, kurus rūpīgi apkopojis Ē.~Jēkabsons.

Latvieši poļu gūstā sāka nonākt jau 1919.~gada agrā pavasarī, bet sevišķi daudz~--- pēc poļu karaspēka veiktās Viļņas ieņemšanas 1919.~gada aprīlī. Kopumā 1919.--1920.~gadā Polijas armijas gūstā un internējumā nonāca vairāki tūkstoši Latvijas iedzīvotāju, galvenokārt~--- kā karagūstekņi, turklāt viņu dzīves apstākļi gūstā bija ļoti smagi. Viņi tika nosūtīti uz dažādām gūstekņu nometnēm, arī Polijas iekšienē. Saskaņā ar ziņojumiem, vispārējā poļu attieksme pret gūstekņiem un to dzīves apstākļi bija slikti, ko acīmredzot noteica notiekošā poļu-padomju kara ārkārtīgi nežēlīgais raksturs. Tomēr bija vērojami arī sevišķas nepatikas izpausmes un pat nežēlīgas attieksmes gadījumi no poļu administrācijas puses tieši pret latviešiem viņu nacionālās piederības dēļ, ko savukārt acīmredzot noteica Latviešu padomju divīzijas kā Sarkanās armijas spējīgākā triecienspēka un latviešu ievērojamāko lielinieku darbības radītais viedoklis par latviešiem, kas tika kultivēts un pārspīlēts gan ar Poliju cieši saistīto krievu ''balto'', gan visas pasaules presē. Robeža starp gūstekņiem un internētajiem nebija stingri noteikta un savstarpēji maz ar ko atšķīrās, īpaši attiecībā uz Sarkanās armijas gūstekņiem un t.s. padomju internētajiem ``civilgūstekņiem'' (faktiski~--- sagūstītajiem padomju iestāžu darbiniekiem).

Vēl 1919.~gada oktobrī~--- decembrī Stšalkovas gūstekņu nometnes slimnīcā un Dombes nometnes slimnīcā no dizentērijas un citām slimībām nomira 15 latvieši.

Īpašs gadījums bija sešu latviešu gūstekņu grupa, kurai kopā ar latviešu tautības sargu izdevās izbēgt no gūstekņu nometnes Grodņā un pēc divu nedēļu gājiena cauri Lietuvai, kura laikā viņu vienīgā pārtika bija ``saknes un augļi'', 1919.~gada septembra beigās ieradās Jelgavā. Viņi izskatījās atbaidoši~--- bez apaviem un apģērba driskās, turklāt stāstīja, ka labprātīgi pārgājuši poļu pusē kauju laikā pie Viļņas un Minskas.

Kad Latvijas varas iestādes lūdza atbrīvot latviešu gūstekņus, poļu puse vilcinājās, jo tai svarīga bija viņu politiskās pagātnes pārbaude. Latvijas pilnvarotais gūstekņu lietās leitnants P.~Strods 1919.~gada novembrī kopā ar Polijas Austrumu zemju komisijas pārstāvi apmeklēja vairākas gūstekņu nometnes, konstatējot sevišķi smagus dzīves apstākļus Pikuļices nometnē pie Pšemišļas, kur atradās 50~latviešu. Viņiem trūka veļas, virsdrēbju, tiem bija no aukstuma piepampušas sejas, melnas rokas.

Vēlāk no Stšalkovas, kur atradās apmēram 500~latviešu, vienā dienā tika atbrīvoti gandrīz tie visi, pagaidām atstājot nometnē tikai 27~cilvēkus, kuri tika turēti aizdomās par piederību Komunistiskajai partijai. Šie simti tika atlaisti bez pienācīga apģērba un iztikas līdzekļiem, nepaziņojot par to Latvijas pārstāvjiem. Pirms 1919.~gada Ziemassvētkiem Varšavā cits pēc cita sāka ierasties vilcienu sastāvi, pilni ar ``puskailiem'' latviešu tautības bijušajiem poļu gūstekņiem. Latvijas puse pieņemšanai arī nebija sagatavota, trūka materiālo līdzekļu, tādēļ bijušie gūstekņi ieradās Latvijā noskranduši un pusbadā. Pēc neilga atvaļinājuma viņi tika mobilizēti Latvijas armijā.

Praktiski viņus izglāba samērā labās Latvijas un Polijas starpvalstu attiecības, kuru rezultātā Polijas varas iestāžu attieksme pret latviešiem, sākot ar 1919.~gada nogali bija salīdzinoši laba un viņiem tika dota iespēja atgriezties dzimtenē.

Jāņem vērā, ka aprakstītā attieksme bija vērojama pret latviešiem, jaundibinātas valsts, kura karoja pret boļševikiem, iedzīvotājiem, uz laiku nonākušiem Sarkanajā armijā, un kuri nebija apzināti kaitējuši Polijai. 1920.~gada otrajā pusē attieksme pret Sarkanās armijas karavīriem~--- Padomju Krievijas pavalstniekiem, kura apdraudēja Polijas eksistenci, acīmredzot bija ievērojami sliktāka. Materiāls netieši rāda, ka pret politiskā pretinieka~--- Padomju Krievijas karavīriem, pat ja arī tie karoja nevis boļševistiskās idejas vārdā, bet vienkārši aizstāvēja savu zemi pret svešzemju iekarotājiem, daudzu poļu nostāja bija gaužām nežēlīga.

Lai rastu pilnīgāku priekšstatu par gūstekņu stāvokli abās karojošajās pusēs jāpasaka, ka padomju puse arī bija saņēmusi vairāk nekā 40 tūkstošu \strong{poļu karagūstekņu}. (Pēc KK(b)P CK rīcībā esošiem datiem uz 1920.~gada rudeni 33~nometnēs Krievijas teritorijā atradās ap 30~tūkstošu poļu karagūstekņu.) XXI gadsimta sākumā sagatavotos arhīvu dokumentu krājumos maksimālais poļu gūstekņu skaits padomju republikās tiek lēsts kā 42~tūkstoši. (Ir izskanējis arī it kā precīzs skaitlis, ka no 1918~līdz 1920.~gadam padomju gūstā nonācis pavisam 51~351 poļu karavīrs.) No tiem pēc kara Polijā atgriezās ap 35~tūkstoši (Ir arī citi dati~--- 27,6 tūkstoši, bet pēc Z.~Karpusa ziņām~--- 26,5 tūkstoši cilvēku). Ap 2~000 poļu palika KSFPR. Pārējo tūkstošu liktenis nav skaidrs. Var pieņemt, ka vismaz daļa no viņiem gāja bojā gūstekņu koncentrācijas nometnēs.

Pēc ziņām krievu literatūrā padomju nometnēs poļu gūstekņu stāvoklis bijis labāks nekā padomju gūstekņu~--- poļu nometnēs. Padomju pusē vairākums poļu gūstekņu uzlūkoti kā ``šķiras brāļi'' un represijas pret tiem neesot bijušas iedomājamas. Ja arī notikuši kādi ekscesi, tad pavēlniecība centusies tos novērst un sodīt vainīgos. Tā, 1920.~gada jūlija beigās, kad Sarkanā armija vēl uzbruka, poļu komunistu darbinieks J.~Marhļevskis sastapa sagūstītu poļu grupu, kura bija labi ģērbta un apauta, bet kuru konvojēja basi un noskranduši sarkanarmieši. ``Šie varoņi,~--- konstatēja J.~Marhļevskis,~--- atteicās no tā, lai atņemtu gūstekņiem zābakus''. Šādas ainas novēroja ne tikai J.~Marhļevskis, kas gan neliecina par to, ka nevarēja būt arī pretēji gadījumi.

1920.~gada 17.~jūnijā Padomju Rietumu frontes pavēlniecības izdotā pavēle skanēja: ``Komandieriem un komisāriem veikt visus pasākumus, lai pilnīgi saprotamais Sarkanās armijas karavīru naids pret poļu baltgvardiem nekādā ziņā netiktu attiecināts pret gūstekņiem. Gūstekņus ir jānosūta uz aizmuguri tajā apģērbā, kurā tie tika sagūstīti. Sarkanarmiešiem ir stingri jāapzinās, ka sagūstīts poļu karavīrs vairs nav ienaidnieks. Tas vairākumā gadījumu ir mobilizēts Polijas strādnieks vai zemnieks. Viņš ir nevis jāapvaino, bet viņam ir jāatsedz acis uz poļu muižnieku un buržuāzijas uzbrukuma krievu tautai negodu. Būt nesaudzīgam kaujā~--- bruņiniekam attiecībā pret uzvarētajiem~--- lūk, revolūcijas cīnītāju devīze''. Sarkanajā armijā pavēles par cilvēcīgu attieksmi pret gūstekņiem tika izdotas daudzkārt, bet jau šo pavēļu daudzums liecina, ka tās ne vienmēr tika pienācīgi pildītas.

KSFPR Darba un Aizsardzības padome 1920.~gada 15.~septembrī izdeva lēmumu par poļu karagūstekņu uzturēšanu. Pārtikas un citu mantu apgāde, medicīniskā aprūpe tiem pienācās pēc aizmugurē atrodošos sarkanarmiešu normām. 1921.~gada 2.~februārī, papildinot šo lēmumu, tika nolemts poļu karagūstekņiem, kuri strādāja rūpnīcās, izmaksāt pilnu darba algu.

Ievērojami nopelni sagūstīto poļu stāvokļa uzlabošanā bija Sarkanās armijas Politiskās pārvaldes Poļu sekcijai, kura izvērsa propagandas darbu gūstekņu vidū un tāpēc arī bija ieinteresēta, lai tie būtu apmierināti ar uzturēšanas apstākļiem.

Krievu vēsturnieks S.~Poltoraks, atsaucoties uz arhīvu dokumentiem, rāda, ka padomju izpildu komitejas gādāja poļu gūstekņiem pārtiku. Viņiem rīkoja sapulces un mītiņus par tēmām: ``Poļi un Padomju Krievija'', ``Par ko karo poļu strādnieks un poļu algādzis'', ``Vai poļu strādniekam vajag karu ar Padomju Krieviju''. S.~Poltoraks raksta par poļu gūstekņu ``pāraudzināšanas'' padomju garā mēģinājumiem, kuri, kā par brīnumu, bieži izdevušies. Viņš to skaidroja ar pašu ``pāraudzinātāju'' lai arī naivo, bet patieso ticību savai vēsturiskajai misijai un to, ka gūstekņu vidū izsētā aģitācijas un propagandas sēkla nereti krita auglīgā augsnē. Daļa gūstekņu agrāk bija atradušies sociāldemokrātisku organizāciju ietekmē un proletariāta vienotības idejas viņiem nebija svešas. Poļu gūstekņu vidū izdalīja aktīvistus, kuri agrāk bija sastāvējuši Vācijas Sociāldemokrātiskajā partijā, PPS, Bundā, un ļāva viņiem patstāvīgi veikt propagandu citu gūstekņu vidū. Jau 1920.~gada jūlijā Jaroslavļā poļu gūstekņi radīja savu klubu un pašdarbības teātri. Kad 1920.~gada 4.~augustā J.~Ļeščinskis uzstājās savu sagūstīto tautiešu priekšā Kostromā, viņš vispirms tos apsveica ar nokļūšanu gūstā, jo tagad tie varēšot ņemt līdzdalību notikumos dzimtenē, kur drīzumā pār Varšavu plīvošot sarkanais karogs. Šajā gūstekņu nometnē tika izveidots politisko nodarbību pulciņš un boļševiku partijas komiteja. Jau 8.~augustā gūstekņu sapulcē tika paziņots, ka viņi atbalsta komunisma ideju un padomju varu. Līdzīga politiskā darbība tika izvērsta arī citās gūstekņu nometnēs. Tā 1920.~gada 9.~septembrī Kijevas apgabala kara komisariātā tika apspriests jautājums par 300~poļu gūstekņu grūto stāvokli nometnē pie Kijevas. Tika pieņemts lēmums noregulēt visas sadzīves problēmas un izvērst politisko darbu.

Taču, acīmredzot nevajadzētu arī pārspīlēt sagūstīto poļu dzīves ``jaukumus'' padomju nometnēs. Varētu domāt, ka nometņu administrācija saprata, ka visa aģitācija un propaganda būtu nolemta neveiksmei, ja pret gūstekņiem tiktu pausta attieksme kā pret nīstiem ienaidniekiem, netiktu izrādītas rūpes par ciešamu dzīves apstākļu sagādāšanu, tāpēc centās radīt tiem maksimāli labus apstākļus. Taču pret to netieši liecina KK(b)P CK 1921.~gada 1.februārī apstiprinātais partijas guberņu komitejām izsūtāmās vēstules teksts ar grifu ``nav domāts publicēšanai presē'' (``\rutxti{не для печати}''). Tajā bija teikts, ka daudzi atbildīgi padomju, partijas un militārie darbinieki poļu karagūstekņus uzlūko tāpat ``kā pārējās koncentrācijas nometņu iemītnieku kategorijas.'' Vēstule prasīja pastiprināt partijas un kultūrizglītojošo darbu gūstekņu vidū.

Kad kļuva skaudrs, ka miera līgums ar Poliju tiks drīzumā parakstīts, akcents no poļu gūstekņu iesaistīšanas bruņotā cīņā pret Poliju tika pārbīdīts uz viņu pasaules uzskata izkopšanu, lai viņi, atgriezušies Polijā, veiktu tajā sociālistiskās revolūcijas sagatavošanas darbu. 1921.~gada janvārī Kijevas nometnes gūstekņi pieņēma rezolūciju, kurā solīja Polijā cīnīties ``lai uz šlahtiskās Polijas drupām uzceltu pilsētu un lauku strādnieku padomju Sarkano Poliju.'' Cik šādi paziņojumi bija patiesi, atliek tikai minēt, varbūt gūstekņi vēlējās vienkārši ātrāk atgriezties ``šļahtiskajā Polijā'' un uzsākt tur mierīgu dzīvi.

Plašu poļu gūstekņu dzīves ainu padomju republikās sniedz 2004.~gadā Maskavā izdotais dokumentu krājums (\rutxti{Польские военнопленные в РСФСР, БССР и УССР (1919--1922 годы). Документы и материалы.~--- М.: Институт славяноведения РАН, 2004}). Tas rāda gan proletāriskās solidaritātes iedibināšanas mēģinājumus, to dažādos~--- gan pozitīvos, gan negatīvos rezultātus, gan padomju iestāžu pūles idejiski audzināt gūstekņus, gan sliktos sadzīves apstākļus un lielo mirstību to vidū. Tomēr kopumā, zinot par nepietiekamo Sarkanās armijas apgādi, vājo tās kaujinieku disciplīnu, vispārējo slikto stāvokli padomju aizmugurē, rodas aizdomas, ka vispārējais poļu gūstekņu stāvoklis padomju nometnēs tomēr nebija daudz labāks par padomju gūstekņu stāvokli poļu nometnēs. Lai kā arī būtu, ir jāturpina pētījumi gan par vienu, gan otru gūstekņu likteni. 2008.~gada janvārī poļu-krievu vēsturnieku grupa turpināja sarežģīto kopējo vēstures problēmu izpētes darbu. Krievijas toreizējais premjerministrs V.~Putins 2010.~gadā paziņoja, ka ``saskaņā ar jaunākajiem datiem no bada un slimībām Polijas gūstā nomira 32 tūkstoši [padomju cilvēku].''

1921.~gada 24.~februārī starp KSFPR un USPR, no vienas puses, un Polijas Republiku, no otras puses, tika parakstīta vienošanās par repatriāciju. Tā stājās spēkā no parakstīšanas brīža un nebija jāratificē. Abas puses apņēmās iespējami īsā laikā repatriēt visus to teritorijā esošos ķīlniekus, karagūstekņus, civilgūstekņus, internētos, bēgļus un emigrantus. 1922.~gada aprīlī Polijas nometnēs vairs nebija neviena padomju karagūstekņa.

Tiesa, abas puses vēlējās pirmkārt panākt tai visvairāk vēlamo repatriantu nosūtīšanu mājup no otrās puses teritorijas. Tas kavēja, ieviesa juceklī visa darba organizācijā un tikai palēnināja to. Daudzi poļu bēgļi dzimtenē atgriezās neorganizētas pašplūsmas ceļā. No 1918. 1īdz 1920.~gadam vairāk nekā miljons poļu atstāja Sarkanās armijas ieņemtos apgabalus un teritorijas, kas nonāca Padomju Krievijas sastāvā pēc Rīgas miera. Bez tam, pēc 1924.~gada 21.~maijā apkopotiem padomju datiem no 1921.~gada aprīļa līdz 1924.~gada aprīlim no padomju republikām Polijā atgriezās vēl ap 1~000~000 cilvēku.

Poļi sastādīja 15~līdz 25\% no repatriantiem, pārējie bija galvenokārt baltkrievu un ukraiņu zemnieki. Acīmredzot ziņas nebija visai precīzas, jo pēc līdz 1924.~gada vidum strādājošās jauktās krievu-ukraiņu-poļu repatriācijas komisijas datiem no PSRS uz Poliju aizbrauca apmēram 300~tūkstošu poļu. Daudzi poļi nevēlējās atgriezties Polijā, labāk palika Padomju Krievijā. Acīmredzot tikai mazākā daļa no tiem to darīja ideoloģisku apsvērumu rosināti. To var secināt no tā, ka 1922.~gadā tikko izveidotajā Padomju Savienībā bez Jakutijas un Tālajiem austrumiem dzīvoja ap 5~600 poļu komunistu. (Viņu organizācijas vadīja speciāli politbiroji: 4~apgabalu, 16~guberņu un 20~apriņķu.) 1926./1927.~gadā PSRS teritorijā bez Aizkaukāza un Uzbekijas dzīvoja gandrīz 790~500~poļu. Vairāk nekā 526~000 no viņiem bija lauku iedzīvotāji. Pēc poļu avotiem PSRS teritorijā palika ap 1,5~miljoni poļu, kuri neizmantoja tiesības repatriēties uz dzimteni. Domājams, ka padomju dati ir tuvāk īstenībai, jo tie balstījās arī uz 1926.~gada iedzīvotāju skaitīšanas rezultātiem. (Tiesa, pēc poļu-padomju kara daudzi poļi, apzinot sevi par tādiem, tomēr skaitīšanā varēja nosaukt citu tautību, lai izvairītos no ieskaitīšanas ``panu'' kategorijā.)

\strong{1921.~gada 18.~martā Rīgā KSFPR} (kura uzstājās arī Baltkrievijas PSR vārdā) un Ukrainas PSR delegācija, no vienas puses, \strong{un Polijas Republikas delegācija}, no otras, \strong{parakstīja miera līgumu}, kurā noteica robežu starp Poliju un padomju republikām. Padomju~--- Polijas robeža izveidojās uz kauju rezultātā radušās frontes līnijas pamata. Nekāda cita, teorētiska pamata šai robežai nebija un arī nevarēja būt. Polija saņēma teritorijas, kuras tā varēja sarunu ceļā iegūt vēl 1919.~gada beigās, 1920.~gada sākumā. Turpmāk Polija formāli robežojās ar Baltkrievijas un Ukrainas padomju republikām, bet faktiski ar to pašu vien Padomju Krieviju, kura 1922.~gadā kopā ar citām republikām pārtapa par PSRS. J.~Pilsudska plāns par bufervalstu radīšanu ar Krieviju netika īstenots. Poļu tautai pēc Pirmā pasaules kara beigām nācās vēl ap trijiem gadiem karot, samaksājot ar desmitiem tūkstošu poļu dzīvību, lai iegūtu vismaz daļu no poļu politiskās elites iecerētās ģeopolitiskās telpas.

Interesanti, ka Antantes valstis, kas Versaļas līgumā bija nostiprinājušas sev tiesības atzīt Polijas austrumu robežu, to izdarīja tikai 1923.~gada martā. Tieši tad tika pabeigts Polijas Republikas valsts teritorijas noformēšanas process, kas sākās 1918.~gada novembrī. Taču Rietumvalstis neatsaucas uz Rīgas miera līgumu, tādējādi demonstrējot savu nevēlēšanos nepārprotami nostāties Varšavas pusē šajā teritoriālajā strīdā un distancējoties no poļu-padomju kara rezultātā noslēgtā līguma.

Vācijā dzīvojošais bijušais padomju diplomāts, politiķis un vēsturnieks V.~Faļins Rīgas mieru salīdzinājis ar Brestas mieru, kuru Padomju Krievija noslēdza ar ķeizarisko Vāciju 1918.~gada pavasarī. Arī 1921.~gadā, tāpat kā 1918.~gadā, boļševikiem trūka spēka turpināt karu, turklāt Tālajos Austrumos Japānas kontrolē vēl atradās plaši Krievijas apgabali, kas lika baidīties no imperiālistisko valstu vienošanās un pat karagājiena pret Padomju Krieviju. Brestas miers sadalīja līdz tam Krievijā dzīvojošās tautas, to pašu darīja Rīgas miers. Principiāli tas maz atšķīrās no Žečpospolitas dalīšanām XVIII un XIX gadsimtā, tikai ja agrāk galvenais cietējs bija poļu tauta, tagad tās vietā bija ukraiņu un baltkrievu tautas.

Starp citu, mūsdienu baltkrievu vēsturnieks I.~Krens pat izteicis izbrīnu, ka boļševiki toreiz savos centienos panākt pasaules revolūcijas uzvaru neatdevuši Polijai visu Baltkrieviju. Pēc autora domām, Rīgas miers kā reiz pierādīja, ka minētās cerības, kaut ietekmēja boļševiku nostāju, taču nebūt nebija vienīgās, kas to noteica.

\strong{Kara iznākumu starp Padomju Krieviju un Poliju}, kā jau visos karos, noteica spēku samērs. Protams, liela bija ārvalstu atbalsta loma. Ja ne sākotnējais Vācijas, pēc tam Antantes valstu vispusīgais atbalsts Polijai, poļu-padomju kara iznākums varētu būt arī citāds. Padomju varas nodibināšanas iespēja Polijā gan bija visai nenozīmīga, taču valsts robežas droši vien būtu citādas. Attiecīgi citāda būtu arī politisko spēku ietekmes aina valstī.

Poļu-padomju kara rezultātā poļiem atkal izdevās šķērsot Kērzona līniju, gan nesasniedzot 1772.~gada robežas. Karā neviena puse neieguva visu vēlamo.

Krievu vēsturnieks S.~Poltoraks atzīmējis, ka neviena no karojošajām pusēm neatzina savu sakāvi. Abas puses pasludināja, ka tās ir uzvarējušas. 1921.~gada februārī V.~Ļeņins deklarēja ``\dots{}mēs no kara izgājām kā uzvarētāji'', jo 1920.~gadā Padomju Krievija bija piedāvājusi Polijai izdevīgākus miera noteikumus, nekā tie, kurus tā izcīnīja ar Rīgas miera līgumu 1921.~gada martā. (Starp citu, 1921.~gada oktobrī, spriežot pēc avīžu ziņām par Poliju, V.~Ļeņins rakstīja: ``Ja viss ir tā, tad (padomju) revolūcijai Polijā ir jāuzvar un drīzumā'' un pat brīdināja, ka nedrīkst padoties provokācijai un sākt priekšlaicīgu sacelšanos. Acīmredzot šeit jārunā par avīžu ziņu nepilnību un arī paša V.~Ļeņina pārlieko revolucionāro optimismu.)

Autoritatīvais poļu vēsturnieks P.~Lossovskis uzskata, ka karā uzvaru neguva neviens: ``\dots{}Šajā karā gala rezultātā nebija uzvarētāju un uzvarēto. Taču svarīgākais tika sasniegts~--- Polija aizstāvēja savu neatkarību, kaut kara gaita bija ļoti mainīga, un uzvaras kauss nosliecās te uz vienu, te otru pusi.'' Pie teiktā gan jāpiebilst, ka kara uzsācēja Polijā tajā zaudēja 184~tūkstošus karavīru, iegūstot sliktāku robežu nekā tai 1920.~gada sākumā piedāvātā, un vēl vairāk saasinot jau tā nedraudzīgās attiecības ar Padomju Krieviju/PSRS.

Arī latviešu vēsturnieks A.~Miņins uzskata, ka gan Padomju Krievijai, gan Polijai karš bija ``uzvaras zaudējums'', jo vienā gadījumā netika realizēta pasaules revolūcijas ideja, tomēr tika radīta iespēja Padomju Krievijas valstiskuma nostiprināšanai ierobežotā teritorijā, kas gan nesakrita ar Krievijas impērijas ārējām robežām, bet otrā gadījumā tika radīta iespēja Polijas valstiskuma nostiprināšanai, tiesa, ne ``Lielās Polijas'' robežās.

Var gan atzīmēt, ka mūsdienās Polijā atrodas politiķi un arī vēsturnieki, kas Rīgas līguma noslēgšanu traktē kā kļūdu. Attiecībā uz Rīgas līgumu tiek lietoti apzīmējumi ``Jalta 21'', ``valsts nodevība'', ``uzvarētāju padošanās'' vai “politiskās muļķības akts''. Ar to mirusi ideja par federālu, daudznacionālu Poliju. Var sastapt viedokli, ka poļiem līgums bijis ne tikai neveiksmīgs, bet vienkārši amorāls, jo atstājis aiz valsts robežām, “padomju mežoņu žēlastībai'', ap pusotra miljona poļu. Jānožēlojot, ka ārpus Polijas palika arī apgabali ap Daugavpili ar samērā lielu poļu īpatsvaru. Daudzi poļi Padomju Savienībā gāja bojā gan kolektivizācijas laikā notikušā bada dēļ, citi kļuva par “poļu operācijas'' upuriem 1937./1938.~gadu terora laikā. Šie poļi esot “pārdoti'' apmaiņā pret tukšiem padomju solījumiem. Tiesa, skan arī reālistiskākas balsis. Polijā augstu vērtētais Polijas un Krievijas attiecību eksperts profesors A.~Novaks gan norāda, ka galvenais Polijai bija pēc iespējas ātrāk izbeigt karu ar Padomju Krieviju, novērst tās sabiedrošanos ar Lietuvu, ko Polija uzskatīja par savu pretinieci, un panākt Polijas drošību. Tas arī ticis panākts.

Poļu vēsturnieks Z.~Zaluskis visai aptveroši raksturoja poļu-padomju kara vietu tālākajās Padomju Krievijas/Padomju Savienības un Polijas attiecībās: ``Ne no tās, ne citas puses netika realizēta neviena programma, neviena cerība, netika sasniegti vēsturiskie mērķi un politiskie nodomi~--- ne galvenie, ieplānotie, ne arī tekošie, notikumu attīstības gaitā improvizētie. Varētu teikt, ka sacensības rezultāts līdzinājās nullei. Varētu, ja starp Varšavu un Maskavu, Kijevu nebūtu parādījušies jauni kapu lauki, izauguši jauni, noturīgi rūgtuma un neuzticības, sāpju un naida šķēršļi. \citespace{} Kauju laukos \citespace{} galīgi izplēnēja ilūzija, ka Polija var uzdāvināt brīvību Ukrainai un Baltkrievijai, atnesot to uz savu pulku durkļiem un atvedot no sevis atkarīgas valdības un nacionālās komitejas. Izklīda ilūzija, ka Polija ir spējīga organizēt jaunu kārtību zemēs, kuras kādreiz~--- pirms simtiem gadu bija ar viņu saistītas, bet tai laikā izrādījās tai pilnīgi svešas. Sabruka federālistiskā koncepcija, Valsts priekšnieka un Pirmā Maršala politiskā pamatkoncepcija. Un \latxti{vice versa} [latīņu~--- otrādi]: kauju laukos \citespace{} izplēnēja ilūzija par sociālistisko revolūciju Polijā, par palīdzību Eiropas revolūcijai. Ar Tuhačevska armijas atkāpšanos mūžībā pazuda Polijas pagaidu revolucionārā komiteja, poļu Sarkanās armijas aizmetņi un visa Polijas Padomju Republikas ideja''. Citētā doma gan jāprecizē, konstatējot, ka kara rezultātā lieli apgabali ar pārsvarā lietuviešu, baltkrievu un ukraiņu iedzīvotājiem nonāca Polijas varā.

Polijas formālā uzvara karā pret padomju Krieviju ar laiku pārvērtās par faktisku poļu tautas zaudējumu, jo uzvaras iejūsmināta, Polijas vadība atkārtoja Žečpospolitas pieredzi līdz tās sagrūšanai, mēģināja radīt savu mazu impēriju, līdzspēlēt Eiropas lielvalstu politikā, kas galu galā noveda pie atkārtotas valsts faktiskas likvidācijas.

Jau pēc 1920.~gada 12.~oktobrī noslēgtā pamiera līguma Polijas Komunistiskās strādnieku partijas CK uzsaukumā poļu tautai norādīja: ``Miljoni baltkrievu, lietuviešu, ukraiņu zemnieku un strādnieku pēc Rīgas līguma tiek vardarbīgi pievienoti Polijai \citespace{} līdz ar to neapgāžami ir pierādīts tas, par ko Polijā nemitīgi runāja tikai mūsu partija: Polija veda imperiālistisku, iekarošanas karu\dots{}''.

Rīgas līgums paredzēja Padomju Krievijas un Padomju Ukrainas atteikšanos ``no jebkādām tiesībām un pretenzijām no zemēm'', kuras palika līgumā nospraustās robežas Polijas pusē. Taču miera sarunu laikā Krievijas un Ukrainas delegācija pauda viedokli, ka Austrumgalīcijas likteni jānosaka tās iedzīvotāju brīvas gribas izteikšanas rezultātā. Kad tā nenotika, padomju puse vairākkārt starptautiski aktualizēja jautājumu par Austrumgalīcijas piederību. Jāpasaka, ka Tautu Savienības padome tikai pēc diviem gadiem~--- 1923.~gada 15.~martā oficiāli atzina Polijas austrumu robežu, kas, iespējams, liecināja, ka pat antiboļševistiski noskaņoto Rietumu politiķu acīs tā bija apšaubāma.

Pēc 1921.~gada 18.~marta Rīgas miera līguma Polijas robeža gāja tālu uz austrumiem no Kērzona līnijas. Baltkrievu un ukraiņu apdzīvotās teritorijas sadalīja Polija un padomju republikas, kas 1922.~gadā iegāja PSRS sastāvā. Polijas rokās nonāca gandrīz puse Baltkrievijas un ceturtā daļa Ukrainas ar Ļvovu. Polija ieguva lielu daļu Volīnijas un visu Austrumgalīciju. Šīs teritorijas poļi uztvēra kā savas valsts ``mežonīgās nomales''.

Līguma rezultātā Polijas iegūtajās vairāk nekā 200~000~km$^{2}$ lielajās strīdīgajās teritorijā tika izveidotas 8~vojevodistes. No tām tikai Belostokas vojevodistē bija poļu iedzīvotāju vairākums. Viļņas vojevodistē poļi dzīvoja blakus lietuviešiem, ebrejiem un baltkrieviem, Poļesjes~--- baltkrieviem un ukraiņiem, Novogrudokas vojevodistē vairākums bija baltkrievi, bet Volīnijas, Ļvovas, Staņislavas un Tarnopoles vojevodistēs vairākumā bija ukraiņi.

Rietumvalstis Polijas tiesības uz rietumukraiņu zemēm atzina ar lielām atrunām. Izvirzījās prasība Polijai piešķirt tām autonomiju, tika spriests arī par Tautu Savienības mandāta piešķiršanu uz šo zemju pārvaldi Polijai uz 25~gadiem. Apmierinošs starptautisks risinājums šai problēmai tā arī netika atrasts. (Jāatzīmē, ka ar miera līguma noslēgšanu Polija pameta savus formālos sabiedrotos~--- ukraiņu nacionālistus ar S.~Petļuru priekšgalā, kuri līdz ar to vairs nevarēja būtiski apdraudēt Padomju Ukrainas pastāvēšanu.) Arī lietuviešu apdzīvotās teritorijas sadalīja Lietuva un Polija. Ar to Polija kļuva par valsti, kur paši poļi sastādīja tikai ap 64\% no visiem iedzīvotājiem.

Ukrainas, Baltkrievijas un Lietuvas zemju sadale nozīmēja arī to, ka bija cietusi neveiksmi J.~Pilsudska ``Starpjūru valstu savienības'', kurā ietilptu poļu, lietuviešu, ukraiņu, baltkrievu un varbūt arī vēl citas zemes, izveides koncepcija. Jau 1920.~gada beigās J.~Pilsudskis saprata, ka viņam neizdosies realizēt savu mūža sapni~--- XVIII gadsimta Žečpospolita netiks atjaunota pat jaunā, modernā federācijas variantā. 3.~decembrī, līdzgaitnieku pulkā viņš paziņoja: ``Mums ir atkal jāieņem Kijeva un Minska, lai apvienotu federācijā vai savienībā ar Poliju Ukrainu, Baltkrieviju un Lietuvu. Par nožēlu, izdarīt to tagad es nespēju. Polija to nevēlas un tai ir svešas Jagelloņu idejas, bez tam Polija ir ļoti karā novājināta. Mūsu armija ir ļoti nogurusi \citespace{} un nabaga. Nē, mēs nevaram šo uzdevumu izpildīt.'' 1921.~gada maijā, jau pēc Polijas un Padomju Krievijas miera līguma noslēgšanas J.~Pilsudskis apmeklēja vienu no internēto bijušo sabiedroto nometnēm un lūdza piedošanu viņiem, ka nespēja realizēt savu austrumu projektu un palīdzēt tiem radīt neatkarīgu valsti.

Kad 1921.~gadā galīgi izveidojās starpkaru Polija, tā bija kas vidējs starp abām sākotnējām~--- jagelloņu un pjastu koncepcijām, bet tālu pārsniedza robežas, kur poļi bija neapstrīdamā vairākumā. Pastāv diezgan pamatots viedoklis, ka praksē uzvarēja R.~Dmovska nacionālās valsts koncepcija. Paradoksālā veidā J.~Pilsudskis, atrodoties valsts priekšgalā, ar savām, uz federatīvu valstu savienības radīšanu virzītām metodēm, bija novedis to tuvāk sava politiskā pretinieka R.~Dmovska~--- mononacionālas etniskas poļu valsts ideālam, kura savā sastāvā neietvertu ``pārmērīgu'' skaitu nacionālo mazākumu pārstāvju un kuras daudznacionālais raksturs, kā poļu politiķi cerēja, saglabātos tikai uz nenoteiktu laiku. Taču Polijai draudzīgu, pateicīgu tai par palīdzību sava valstiskuma radīšanā Ukrainas, Lietuvas un Baltkrievijas vietā par Polijas kaimiņiem kļuva divas padomju republikas un ne mazāk kā tās naidīgā Polijai Lietuvas Republika. Polija bija sabojājusi attiecības ar visiem austrumu kaimiņiem, izņemot Rumāniju un Latviju.

Pēc citu tautu apdzīvotās teritorijas aneksijas Polijai turpmāk bija jārēķinās ar vairāk vai mazāk atklātām Padomju Krievijas/Padomju Savienības revizionistiskajām prasībām. Līgumā noteiktā robeža kļuva par nepārvaramu barjeru attiecībās starp Poliju un Padomju Krieviju. Robeža sadalīja ukraiņu un baltkrievu tautas, pie tam sadalīja, nevis ņemot vērā iedzīvotāju vēlmes, bet ar militāru spēku. Radās jauns nacionālo antagonismu uzkurināšanas avots, kura likvidēšana būtu iespējama tikai abām pusēm realizējot konsekventu sabiedriskās dzīves demokratizāciju, kas diemžēl nenotika.

J.~Pilsudskis uzskatīja, ka turpmāk Padomju valdība, sastopoties ar iekšējām grūtībām, centīsies tās pārvarēt ar sasprindzinājuma uzturēšanu uz robežas ar Poliju, tāpēc viņš centās stiprināt Polijas armiju, maksimāli mazināt Seima kontroli pār to, pakļaut savai pārraudzībai, kas izsauca neapmierinātību Seimā. Paredzamais naidīgums attiecībās ar Padomju Krieviju/Padomju Savienību traucēja uzticības pilnu attiecību izveidi starp J.~Pilsudski un franču militāristiem. Pēdējiem Polija bija nepieciešama galvenokārt kā sabiedrotais pret Vāciju, bet J.~Pilsudskis uz Franciju raudzījās kā uz sabiedroto pret Polijas austrumu kaimiņu.

Polija izveidojās par unitāru valsti, kurā nacionālajiem mazākumiem tika liegta teritoriālā un kulturālā autonomija. Versaļas sistēmas ietvaros kā pretsvars Padomju Krievijai uz spēka pamata radītā daudznacionālā Polija apvienoja pēc nacionālās pašapziņas, reliģijas, kultūras, civilizatoriskās pieredzes atšķirīgas tautas. Polijas robežās, kuras nebija nedz vēsturiski, nedz etnogrāfiski, nedz ekonomiski pamatotas, palika tik daudz etniski austrumslāviem un konfesionāli pareizticīgajiem piederīgo iedzīvotāju, ka bija utopiski domāt par viņu asimilāciju un integrāciju nacionālā poļu valstī. Izveidojās paradoksāla situācija: no vienas puses poļi aizstāvēja moderno tautu pašnoteikšanās principu pret ``tautu cietumiem'', ko lielākā vai mazākā mērā bija īstenojušas bijušās tās sadalītājas, bet no otras puses centās izveidot savu daudznacionālu valsti, kurā citu tautību iedzīvotāji atradās nevienlīdzīgā stāvoklī ar poļiem. Tas, ka ASV prezidenta V.~Vilsona saprātīgais un it kā pašsaprotamais priekšlikums, ka neatkarīgajā Polijā jāieiet teritorijām ar neapšaubāmi poļu iedzīvotājiem, tika ignorēts, bija viens no būtiskākajiem jaunveidojamās valsts vājuma cēloņiem. Uz Poliju pilnā mērā varēja attiecināt D.~Loid Džordža novērojumu par ``mazo atbrīvoto tautu'' uzvedību, kuras ``vēl nepaguvušas pašas baudīt brīvības augļus, jau cenšas sev pakļaut citus''. Lielbritānijas premjers pat izdarīja secinājumu: ``Šīs tautas~--- ticiet man~--- ir lielāki imperiālisti, nekā Anglija un Francija, un, protams, lielāki nekā [Amerikas] Savienotās Valstis''. Varbūt secinājums ir tendenciozs, bet rūgtu patiesības graudu tas satur.

Boļševiki Krievijā tai laikā propagandēja federāciju kā sociālistisku valstu savienību. Taču padomju projekts vismaz vārdos neparedzēja jel kādas nācijas pārākumu pār citām. Tas darīja padomju federāciju pieņemamu arī nacionāli noskaņotiem cilvēkiem. Atšķirībā no PSRS ar tās vismaz formāli brīvajām savienotajām republikām Polijas plānotā federācija tā arī nekad netika nodibināta. Patiesības labad gan noteikti jāpasaka, ka ne tikai poļu vadītāji nesa atbildību par robežas noteikšanu atkarībā no spēku samēra. Neviena no konfliktā tieši un netieši (Rietumu lielvalstis) iesaistītajām pusēm valstu robežu noteikšanas jautājumā nerēķinājās ar tikko visu pušu atzītajām nāciju pašnoteikšanās tiesībām.

Rīgas miera līguma 7.pants paredzēja, ka Polija krievu, ukraiņu un baltkrievu tautības iedzīvotājiem, dzīvojošiem tās teritorijā, nodrošinās brīvas tiesības attīstīt savu kultūru, valodu, veikt reliģiskos rituālus. To pašu solīja Padomju republikas poļiem, dzīvojošiem to teritorijā. Faktiski šīs prasības neievēroja neviena no abām pusēm.

Vairākas Polijas politiskās partijas izturējās kritiski pret miera līgumu ar Padomju Krieviju. PKSP orgāns ``\pltxti{Czerwony Sztandar}'' (``Sarkanais Karogs'') 1921.~gada 20.~aprīlī, to izanalizējis, nosauca to par ``daļēju'', bet ``nestipru un nodevīgu'', ``poļu imperiālisma triumfu''. Arī daļa poļu sociālistu no PPS ar ebreju izcelsmes darbinieku F.~Perlu priekšgalā Seimā kritizēja kā poļu, tā krievu tieksmes inkorporēt kaimiņtautu teritorijas. Viņu koncepcija paredzēja nomaļu demokratizāciju, piešķirt tām pašvaldības un plašas autonomijas tiesības kultūras, izglītības un valodas jomā. Arī ļudoviešu \pltxti{Wyzwolenie} iestājās par ukraiņu un baltkrievu autonomijas tiesībām Polijas sastāvā. Taču toreizējos apstākļos šo koncepciju nebija iespējams realizēt.

Miera līgumā abas puses arī atteicās no kara zaudējumu atlīdzības prasībām, Polija tika atbrīvota no atbildības par bijušās Krievijas impērijas parādiem un jebkāda veida saistībām.

Padomju Krievija apņēmās atgriezt Polijai 1915.~gadā izvestās kultūras un materiālās vērtības. Kaut tās tika vestas uz Poliju jau no 1921.~gada gan kastēm, gan veseliem vagoniem, ar pilnīgu saistību izpildi padomju pusei nevedās, tāpēc 1925.~gadā tika noslēgta papildus vienošanās, pēc kuras kā ķīla par noklīdušajām vērtībām uz 15 gadiem Polijai tika nodotas Krievijas imperatoru regālijas, zelts, pērles, dārgakmeņi 30~miljonu zelta rubļu vērtībā. (Tagad šīs regālijas atrodas Maskavā Dimantu fondā.)

1927.~gadā Polija un PSRS parakstīja Ģenerālo vienošanos, ar kuru tika rezumēti Speciālās komisijas poļu kultūras vērtību noskaidrošanā, kuras bija jāatgriež dzimtenē, darba rezultāti. Turpmāk šīs vērtības tika atgrieztas uz šīs vienošanās pamata. Ģenerālā vienošanās paredzēja visu to kultūras vērtību atdošanu Polijai, kuras tika izvestas no tās teritorijas Žečpospolitas sadalīšanas gados. Padomju valsts atdeva Polijai uz tajā ietilpstošajām teritorijām attiecošos arhīvu materiālus, no privātajām bibliotekām izvestās grāmatas, poļu karogus, lielgabalus, aukstos ieročus, zaudētus dažādos cīņas par brīvību periodos, gleznas no poļu muzejiem utt. Interesanti, ka starp citām vērtībām poļu puse pieprasīja tai atdot arī tādus karogus no padomju muzejiem, kuru heraldiskās zīmes izpētot, vēsturnieks P.~Belaveņecs secināja, ka tie ir nevis poļu, bet zviedru karogi un piestādīja poļu pusei tik pārliecinošus argumentus, ka tā no prasības atteicās. Tiesa, 1932.~gadā prasība tika atjaunota un, kā raksta A.~Širokorads, ``lai nebojātu attiecības'' padomju puse šos karogus atdeva Polijai. Poļi saņēma arī Varšavas atslēgas un sudraba timpānus, ko A.~Suvorovam Varšavas maģistrāts pasniedza 1794.~gadā. (Tiesa, visas šīs vēsturiskās un kultūras vērtības 1939.~gadā nonāca vācu iebrucēju rokās, kas tās lielākoties vienkārši ``privatizēja''.) 1932.~gadā darbs bija gandrīz pilnībā pabeigts. Ja kādas kultūras vērtības nevarēja tikt atgrieztas to piederības vietā (piemēram, 1831.~gadā konfiscētais Varšavas universitātes numismātiskais kabinets u.c.), Polija saņēma to ekvivalentu citos priekšmetos.

Ievērojot to, ka Padomju Krievija atgrieza Polijai no tās izvestās kultūras vērtības, tā savukārt vēlējās saņemt no Polijas materiālus, kas sastādīja vērtību boļševiku acīs~--- Polijā esošos ar V~Ļeņina darbību saistītos dokumentus. Daļu no V.~Ļeņina arhīva boļševikiem izdevās atgūt jau 1924.~gadā sarunu rezultātā, kuras Varšavā veda J.~Haņeckis. 1924.~gada 27.~aprīlī Polijas iestādes, iepriekš tos nokopējot, nodeva dokumentus un materiālus (rakstus, to fragmentus, vēstules, piezīmes u.c.) padomju pusei. Kā pateicību par to Polija saņēma itāļu gleznotāja Kanaletto gleznu ``Staņislava Augusta [Poņatovska] ievēlēšana [1764]''. Taču padomju pusei atdoti tika ne visi dokumenti. Daļa no tiem, saistīti ar V.~Ļeņina līdzgaitnieku G.~Zinovjeva un J.~Haņecka darbību palika Polijā un Padomju Savienības pārstāvji lūdza nodot viņiem arī tos. Polijas ārlietu ministrs M.~Zamoiskis ieteica premjerministram V.~Grabskim tos atdot pēc piederības ``ar noteikumu, ka iepriekš tiks noteikts attiecīgs ekvivalents''. Faktiski par materiāliem, kuriem no Polijas puses bija visai nosacīta vērtība, poļi prasīja dokumentus (šai gadījumā t.s. Volīnijas aktus), kuri interesēja poļu zinātniskās iestādes. Divi J.~Haņecka braucieni uz Poliju 1933.~gadā veda pie vēl dažu V.~Ļeņina dokumentu nodošanas PSRS.

Kaut abas līguma slēdzējas puses 1921.~gada 18.~martā oficiāli atteicās no savstarpējām pretenzijām, oficiāli pauda ``pilnīgu savstarpēju cieņu pret valstisko suverenitāti'' un apņēmās ``atturēties no jebkādas iejaukšanās otras puses iekšējās lietās, sevišķi no aģitācijas, propagandas un jebkādas intervences vai tās atbalstīšanas'', KSFPR/PSRS gan atzina Polijas neatkarību un J.~Pilsudska valdības leģitimitāti, taču tikai uz laiku atteicās no plāniem likvidēt Versaļas sistēmu un īstenot pasaules revolūciju. Boļševiku pārliecība ļāva viņiem atteikties no komunistiskās ``aģitācijas, propagandas'' citas valsts teritorijā tikai vārdos. Līgums netika ievērots. (Pēc poļu armijas iestāžu ziņojuma, jau 1921.~gada jūlijā Nesvižas (baltkrievu \betxti{Нясвіж}) apriņķa triju ciemu iedzīvotāji, kurus poļu varas iestādes sodīja par 1920.~gadā notikušu mežu izciršanu, siena izvešanu no svešām (acīmredzot~--- muižnieku) pļavām, it kā ``padomju aģentu aģitācijas'' ietekmē griezās pie Padomju valdības pievienot šos ciemus Krievijai. Ticamāk gan, ka nekāda aģitācija nebija vajadzīga, lai pamudināti no poļu varas daudzkārt cietušos zemniekus pacensties mainīt pavalstniecību.)

Vērtējot karu, jāievēro, ka no abām pusēm poļu-padomju frontē cīnījās vairāk nekā miljons cilvēku. To atmiņā palika neizdzēšamas pēdas~--- rūgtums un savstarpēja nepatika, ko politiķi varēja sekmīgi ideoloģiski ekspluatēt. Poļu vēsturnieks Z.~Zaluskis ir rakstījis par to, ka ``abās pusēs Rīgas līguma robežai, abās pusēs frontei, kura nepieciešamības pēc kļuva par robežu, nostiprinājās stereotipi, vairāk raksturīgi karastāvoklim, nekā draudzīgam mieram. Terminam ``baltpoļi'' Padomju Krievijā bija tāda pati emocionāla nokrāsa, kā terminiem ``boļševiki'' un ``padomes'' Polijā \citespace{} Poļu-padomju attiecību reālos apstākļus izveidoja 1920.~gada notikumi. \citespace{} Cilvēki, kuri izstrādāja un realizēja politiku visa starpkaru perioda garumā gan Varšavā, gan Maskavā,~--- tie bija cilvēki, izveidojušies 1920.~gadā''.

Polijas armija 1918--1920.~gados bija lielākais militārais spēks, ar kuru nācās sastapties boļševiku varai. Tas pārsniedza katru atsevišķi ņemtu ``balto'' armiju. Ja ``baltie'' bija iekšējie ienaidnieki, tad Polija~--- neapšaubāmi ārējais. Un šeit jau darbojās nevis tikai oficiāli kurinātais ``šķiru naids'', bet arī patriotisma jūtas.

Jāpasaka, ka citu problēmas aspektu izcēlis krievu vēsturnieks B.~Nosovs. Viņš, balstīdamies uz 20.~gadu padomju propagandas tēzēm par Antantes karagājieniem, 14~valstu, kuru vidū bija arī Polija, dalību karā pret Padomju Krieviju, uzskata, ka padomju cilvēki poļu-padomju karu uztvēra tikai kā vienu no kopējā Pilsoņu kara epizodēm. Pēc šī darba autora domām šī ``epizode'', kas, atšķirībā no citām, beidzās ar sakāvi, ieņēma īpašu vietu ar to saistīto padomju cilvēku apziņā. Pie tam Polijas uzbrukums 1920.~gada pavasarī sākās tad, kad A.~Kolčaka un A.~Deņikina vadītās lielākās ``balto'' armijas bija jau sakautas, cilvēki cerēja uz drīzām Pilsoņu kara beigām. Tāpēc rūgtums pret poļiem, kuri taču bija saņēmuši brīvību, kā daudziem šķita, no Padomju Krievijas rokām, bija jo lielāks. To padomju cilvēki atcerējās ilgi. Atcerējās arī sakāves rūgtumu, ukraiņu un baltkrievu tautu sadalīšanu. Padomju cilvēku atmiņā palika arī Rietumu proletariāta atbalsts Polijas uzbrukuma upurim~--- Padomju Krievijai, kura mērogus padomju oficiālā propaganda mākslīgi uzpūta. Cerības uz starptautiskā proletariāta atbalstu nodzīvoja līdz pat 1941.~gadam.

Ja pēc Versaļas līguma pretenzijas pret Poliju palika Vācijai, pēc Rīgas līguma~--- Krievijai. Padomju puse joprojām uzskatīja, ka Polija ukraiņu un baltkrievu apdzīvotās zemes Polija ir sagrābusi ar spēku, netaisnīgi, griežot ``pa dzīvu miesu'', ka šī robeža ir ļoti izdevīga uzbrukumam Padomju Krievijai/PSRS, jo uz austrumiem no tās līdz Dņeprai iebrucēju priekšā nebija nekādu dabīgu aizsardzības līniju. Tāpēc ``panu'' Polija tika uztverta kā potenciālais ienaidnieks Nr.~1. Sarkanās armijas vajadzībām tapušā Polijas militāri ekonomiskā un politiskā apskatā, kurš nāca klajā 1929.~gadā, bija teikts: ``No visām kapitālistiskajām valstīm, kuras aplenc PSRS, Polijai militārā ziņā ir tai visbīstamākā.\citespace{} Polijas ģeogrāfiskais stāvoklis padara to par ērtu placdarmu uzbrukumam PSRS tai visapdraudētākajā virzienā uz labā krasta Ukrainu un Donbasu un uz Maskavu caur Baltkrieviju.'' Padomju Savienībā visu starpkaru periodu tika kultivēts ``plēsīgās'', ``imperiālistiskās'' ``buržuāziski-muižnieciskās'' Polijas tēls. Tas kļuva par Padomju valsts ``kapitālistiskā ielenkuma'', kurš pastāvīgi ved cīņu par Padomju Savienības iznīcināšanu, simbolu. Rietumu robežu PSRS nostiprināja daudz vairāk nekā kādu citu. Padomju Savienības robežu divos t.s. sevišķajos kara apgabalos~--- Rietumu un Kijevas~--- piesedza septiņpadsmit nocietināti rajoni, izvietoti divās rindās. Polijas tendences dominēt Austrumeiropā tika uztvertas kā militāri politiskas savienības gatavošana pret PSRS.

Savukārt Polijā aizvadītais poļu-padomju karš izsauca ne tikai patriotisko, bet arī tradicionālo antikrievisko, kā arī sakarā ar ``revolūcijas eksporta'' mēģinājumu saistīto antipadomisko noskaņojumu pieaugumu. Starpkaru Polijā gan pret Padomju Krieviju karojušo, gan vienkārši kara briesmas piedzīvojušo poļu vidū valdīja nepatika pret kaimiņvalsti, kas daudzkārt pastiprināja jau agrāk valdījušo negatīvo attieksmi pret carisko Krieviju. Katrs attiecību ar PSRS aspekts tika skatīts caur karā piedzīvotā prizmu, jebkura padomju iniciatīva starptautiskajā arēnā radīja neuzticību, aiz kuras slēpās bailes no jauna ``pasaules revolūcijas'' spēku uzbrukuma. Vienkāršie cilvēki neapzinājās, ka faktiski pati Polija cīņā par labvēlīgākām robežām bija izsaukusi ``sarkano'' prettriecienu, ka atbildība par piedzīvoto bija jāprasa ne tikai no Padomju Krievijas, bet arī savas valsts vadītājiem.

Poļu-padomju kara cēloņi, norise līdz Otrajam pasaules karam abās pusēs tā arī netika pienācīgi izanalizēti, apauga ar mītiem. Padomju puse nevēlējās runāt par piedzīvotajām sakāvēm, noskaidrot to vaininiekus. Arī J.~Pilsudskis tūlīt pēc kara beigām aizliedza rakstīt un izdot analītiskus un sintezējošus materiālus par to. Atļauts bija apgaismot tikai atsevišķas tēmas, piemēram, kādas kaujas vai karaspēka daļas vēsturi. Ar to kara tēma Polijā tika mākslīgi izņemta no abu valstu ideoloģiskā pretstatījuma.

No otras puses, kauja par Varšavu bija pirmā divsimt gadu laikā poļu uzvarēta liela mēroga kauja, poļu-padomju karš~--- pirmais šai pat periodā poļu uzvarētais karš. Radās ``1920.~gada'', ``brīnuma pie Vislas'' leģenda, kuru enerģiski veidoja J.~Pilsudska piekritēji. Pilnīgā tika aizmirsts pirms tam notikušais ``brīnums pie Dņepras'', kad Sarkanā armija padzina no Ukrainas vidienes poļu iebrucējus. Arī S.~Petļuras, P.~Vrangeļa armiju dalība karā, pat visas ``balto'' kustības objektīvā palīdzība Polijai, neļaujot boļševikiem koncentrēt visus spēkus pret ārējo ienaidnieku~--- Poliju, neatstāja paliekošu iespaidu poļu sabiedrības apziņā. Tautā radās sajūta, ka ``mēs vieni paši sakāvām boļševikus''. Apziņa par to, ka 1920.~gadā Polijai pirmo reizi gandrīz 300~gadu laikā izdevās militāri sakaut Krieviju, radīja poļos pārlieku pašapziņu par saviem spēkiem, par savu valsti kā lielvalsti, kura spēja uzvarēt citu lielvalsti~--- Krieviju. Polijas vadītājus šī sajūta bieži mudināja uz ārpolitiskiem soļiem, kuru realizācijai valstij nepietika reālā spēka.

Blakus pastāvēja vēl viena ideja~--- ``Eiropa neļaus mums iet bojā''. Ģenerālis M.~Veigans Polijas armijas Ģenerālštābā, franču instruktori armijas daļās, amerikāņu un itāļu brīvprātīgie lidotāji Polijas debesīs radīja poļos iespaidu, ka Rietumi kritiskā brīdī Polijai atkal palīdzēs. Ar to tika aizmirsta visa XVIII un XIX gadsimtu pieredze, kad Rietumvalstis izrādīja poļu tautai līdzjūtību, dažkārt deva kādus solījumus, bet, pirmkārt sargājot savas intereses, nopietnu palīdzību tā arī nesniedza. Šīs abas idejas vairāk valdīja vienkāršo cilvēku, nevis labāk informēto politiķu vidū, taču tās atstāja iespaidu uz visu tālāko valsts politiku, ļāva politiķiem manipulēt ar masu apziņu.

Jāuzsver arī, ka Polija nebija apmierināta ne tikai ar to, ka nebija sasniegtas 1772.~gada robežas (Jaunā Polijas valsts aizņēma 52\% no Žečpospolitas teritorijas 1772.~gadā.), bet arī ar boļševiku režīmu kaimiņvalstī, uzskatīja sevi par apdraudētu. Attiecības starp abām valstīm palika saspringtas, kas pēc divdesmit gadiem noveda pie PSRS dalības Polijas sadalē. V.~Vitoss aprakstīja J.~Pilsudska noskaņojumu: 1921.~gada Rīgas līgums ir Polijas sakāve, ja ne pats V.~Ļeņins, tad viņa mantinieki valdīs Varšavā. Taču arī pēc šāda secinājuma atteikšanās no iepriekšējās, uz Krievijas vājināšanu virzītās politikas nenotika.

Latviešu vēsturnieks Ē.~Jēkabsons pamatoti konstatējis, ka starpkaru periodā Padomju Savienība ``sekmīgi centās nepieļaut un faktiski nepieļāva Baltijas valstu un Polijas tuvināšanos, izveidojot aizsardzības savienību'', taču pie teiktā vēl jāpiebilst, ka šādu iespēju Padomju Savienībai sagādāja Polija ar savu agresivitāti, Viļņas apgabala sagrābšanu un pastāvīgajām pretenzijām uz vadošo vietu Austrumeiropā. Viļņas jautājumā padomju republikas atzina, ka ``strīdus teritorijas starp Poliju un Lietuvu, šo teritoriju piederības jautājums vienai vai otrai no šīm abām valstīm ir vienīgi Polijas un Lietuvas ziņā'', taču turpmāk padomju vadība varēja izmantot šo konfliktu savās interesēs.

Poļu~--- padomju attiecības 30.~gadu nogalē bez 1919.--1920.~gada kara, tā abpusējās pieredzes varēja būt pavisam citādākas, varētu abām pusēm dot lielāku iespēju politiskiem manevriem, radīt pavisam citādus nacistiskās Vācijas un pārējās Eiropas konfrontācijas apstākļus.

Jāpievienojas Z.~Zaluska secinājumam: ``Tik tiešām, Polija sāka šo [poļu-padomju] karu un uzvarēja to poļu karavīrs \citespace{} Taču kara varēja nebūt, tam nevajadzēja būt. Polija par to samaksāja ne tikai ar desmitiem tūkstošu gājušo bojā 1920.~gadā, bet arī ar miljoniem bojā gājušo Otrā pasaules kara gados. Kara gājienā uz Kijevu un kara gājienā uz Varšavu nobrieda pati lielākā poļu un krievu traģēdija~--- izolācijas traģēdija.''

Arī mūsdienās~--- kā raksta poļu publicists B.~Skaradzinskis~--- poļi attieksmē pret sen pagājušo 1920.~gadu ``dod priekšroku idilliskai pieejai: kā-nekā tas vainagojās ar uzvarām pie Vislas un Nemunas un ``mūžīgo'' Rīgas mieru. Pie mums klusē par to, ka šis miers ir divdesmit gadus nenoturējās. Neviens neatceras par to, ka Baltkrievijai un Ukrainai ``Rīga'' nozīmēja viņu zemju aneksiju.'' To, ka Polijai pievienotās zemes nonāca Rietumu ietekmes sfērā, izbēga no boļševizācijas un kolektivizācijas, tikai niecīga šo tautu daļa tolaik novērtēja atzinīgi.

Kopumā sava valstiskuma atjaunošanas gaitā Polija bija ar spēku sagrābusi Vācijas, Krievijas un Lietuvas zemes, pārkāpusi vāciešu, lietuviešu, baltkrievu, ukraiņu tautu pašnoteikšanās tiesības, pēc būtības pārkāpusi Parīzes miera konferences lēmumus. ``Notikušu faktu'' politika, kura ļāva Polijai savā sastāvā iekļaut virkni teritoriju gan rietumos, gan austrumos, kaimiņvalstīs radīja pamatotu iespaidu par Polijas valdības, un īpaši J.~Pilsudska, neprognozējamību un viltību, kā arī maz slēptas kaimiņvalstu revanša alkas.

Jaunā Polijas valsts 1921.~gada robežās aptvēra 388~600~km$^{2}$ teritoriju~--- tā bija ievērojami lielāka nekā, piemēram, Lielbritānija vai Itālija. Polijas sastāvā nonāca ap 260~000~km$^{2}$ bijušās Krievijas impērijas zemju, kas sastādīja ap 76,5\% no jaunās Polijas valsts teritorijas, ap 80~000~km$^{2}$ no bijušās Austroungārijas impērijas (ap 20,8\% no jaunās valsts) un 48,6~tūkstošu km$^{2}$ no Vācijas impērijas (11,9\%). Uz 1923.~gadu no 5~548~km Polijas kopējās robežas (citi dati~--- 5~534~km) 1~412~km robežas bija ar PSRS, 1~912~km~--- ar Vāciju, 984~km~--- ar Čehoslovākiju. Nelieli robežas iecirkņi Polijai bija arī ar Rumāniju~--- 349~km, ar Latviju (109~km) un Dancigas brīvpilsētu ar tai pieguļošo teritoriju (221~km). Uz robežas ar Lietuvu (507~km) līdz 1938.~gadam saglabājās kara stāvoklis. Polijas jūras robeža bija 140~km gara. Faktiski iegūtajām Polijas robežām bija stratēģiski neizdevīga konfigurācija. Bez visiem zināmā ``poļu koridora'', starpkaru periodā pastāvēja vēl divi ``koridori'', kuru piederība Polijai bija apstrīdama: viens starp Lietuvu un PSRS, caur kuru Polijas teritorija sasniedza Latviju, un otrs starp Čehoslovākiju un PSRS, caur kuru Polija robežojās ar Rumāniju. Tikai robežu ar Rumāniju varēja nosaukt par tādu, kas atdalīja valstis, kurām nebija nopietnu savstarpēju pretenziju.

Mūsdienu vācu vēsturnieks H.~f.~Zitzevics saskaitījis laika posmā no 1918.~gada 1.~novembra līdz 1938.~gada 1.~oktobrim 17~poļu veiktus karagājienus, sacelšanās un citus ``karam līdzīgus sarežģījumus'' (\detxti{kriegsähnliche Vericklungen}). Tos visus varēja nosaukt par agresijas aktiem pret kaimiņos dzīvojošajām nācijām. Ja ne jaunradītās Polijas valsts, bet precīzāk~--- tās vadītāju nemitīgie agresīvie centieni pēc iespējas palielināt savas valsts teritoriju uz citu tautu apdzīvoto areālu rēķina, kas radīja pamatotu neuzticību Polijai, pēckara situācijā varbūt varētu rasties iespēja izveidot vairāku ap Baltijas jūru atrodošos jauno Austrumeiropas valstu militāri-politisku savienību, kura varētu garantēt to tautām drošību no revanšistisko valstu~--- Vācijas un PSRS~--- centieniem likvidēt pastāvošo Versaļas sistēmu. Tādā gadījumā šo savienību varētu atbalstīt arī karā uzvarējušās Rietumvalstis~--- Lielbritānija un Francija un pasaules vēsture ievirzīties pavisam citā gultnē. Protams, tā būtu tikai iespēja, kuru varētu izjaukt ikviens minēto valstu diktatoriņš, bet, kā zināms, tādu tajās netrūka. Taču pateicoties Polijas demonstrētajai agresivitātei pat tāda iespēja nerādās. Īpaši Lietuvai atņemtās Viļņas jautājums starpkaru periodā neļāva izveidot nekādu aptverošu Baltijas reģiona drošības sistēmu.

20.~gadu sākumā pēc vērtējumiem Polijā bija ap 27~miljoniem iedzīvotāju, no kuriem vairāk nekā 8~miljoni bija pieskaitāmi nacionālajiem mazākumiem. 1921.~gadā poļi bija ap 65\% no visiem valsts iedzīvotājiem, ukraiņi~--- 15\%, ebreji 10\%, baltkrievi~--- 5\%, vācieši~--- 4\%, atlikušo procentu veidoja lietuvieši, krievi un čehi. Vācu literatūrā ir dati, ka 1921.~gada 30.~septembrī Polijā dzīvoja 69,2\% poļu, 14,3\% rutēņu (ukraiņu), 7,8\% ebreju, 3,9\% baltkrievu, 3,9\% vāciešu, 0,3\% lietuviešu, 0,2\% krievu, 0,1\% čehu. Dažādi izdevumi dod dažādus skaitļus, jo nacionālās piederības vietā skaitīšanā fiksēja ticības piederību. Tautas skaitīšanā iegūtie skaitļi bija tendenciozi, poļiem pieskaitīja daļu ebreju, visus baltkrievus-katoļus un vairākumu lietuviešu. Trijās Austrumgalīcijas vojevodistēs: Tarnopoles (\pltxti{Województwo Tarnopolskie}), Staņislavas (\pltxti{Województwo Stanisławowskie}) un Volīnijas (\pltxti{Województwo Wołyńskie}) dzīvoja no 50~līdz 70\% ukraiņu.

Vēlāk, jau 30.~gadu otrajā pusē Polijas iedzīvotāju skaits sasniedza 35~miljonus cilvēku. Pēc iedzīvotāju skaita un blīvuma Polija ieņēma 10.--11.~vietu Eiropā. Pēc Polijas Republikas oficiālajiem datiem tajā dzīvoja 19~miljonu (65\%) poļu, 4~miljoni ukraiņu, 2,1~miljons ebreju, 1,1~miljons vāciešu un tikpat baltkrievu. (Literatūrā gan ir arī dati, ka ukraiņu bija 6,5~miljoni (18,6\%), ebreju~--- 3,4~miljonu (9,7\%) jeb ceturtā daļa no viņu skaita pasaulē, baltkrievu~--- 2,5~miljonu (7,1\%), vāciešu~--- 1~miljons (2,9\%), pārējo~--- 0,6~miljoni (1,7\%). Kā redzam, oficiālie dati arī šoreiz bija tendēti samazināt nacionālo mazākumu un palielināt poļu īpatsvaru valstī.)

65\% iedzīvotāju bija katoļi, 12\%~--- pareizticīgie, 10\% grieķu katoļi-uniāti, Mozus ticīgie~--- 10\%, protestanti~--- 2,5\%. Lielākā daļa iedzīvotāju (75\%) dzīvoja laukos un 55\% no visiem iedzīvotājiem bija zemnieki.

Raksturīga nacionālo mazākumu izvietojuma īpatnība Polijā bija tā, ka viņi visi, izņemot ebrejus, apdzīvoja kompaktas teritorijas, ukraiņi~--- dienvidaustrumos, baltkrievi~--- austrumos, lietuvieši~--- ziemeļaustrumos, vācieši~--- Silēzijā un t.s. Poļu koridorā. Šādā situācijā loģiski izvirzījās nacionālā jautājuma risinājuma variants, kas paredzētu jaunveidojamās valsts robežās ietvert tikai tās teritorijas, kurās bija poļu vairākums un realizēt Polijas nepoļu apmaiņu pret kaimiņvalstīs palikušajiem poļiem. Tā 1922.~gadā izdarīja Turcija un Grieķija, kad pusotra miljona grieķu devās uz Grieķiju un pusmiljons turku uz Turciju. Galīgi teritoriālos jautājumus šī apmaiņa neatrisināja, taču 30.--40.~gados karš starp šīm valstīm neizraisījās. Starp citu, Padomju valsts atļāva visiem poļiem pēc viņu vēlēšanās atstāt to un līdz 1924.~gada 31.~jūlijam uz Poliju pārcēlās 1,1~miljons cilvēku (padomju teritorijā palika dzīvot aptuveni 1,4~ miljoni poļu). Taču valdošajiem poļu politiķiem iedzīvotāju apmaiņas risinājums nešķita pieņemams. Arī Polijas teritorijā dzīvojošie nacionālie mazākumi nealka pārcelties uz neprognozējamo Padomju valsti.

Jaunās valdības priekšā stāvēja sarežģīts uzdevums saliedēt agrāk trijās dažādās valstīs dzīvojušos dažādām nācijām piederošos pavalstniekus vienā dzīvotspējīgā struktūrā. Taču Polija kļuva par unitāru valsti, kura konsekventi noraidīja mazākumtautību tiesības uz teritoriālo un kulturālo autonomiju, līdz ar to vēl vairāk apgrūtinot pavalstnieku saliedēšanas uzdevumu.

1919.--1921.~gada karu rezultātā J.~Pilsudskis vēl vairāk sabojāja attiecības ar tādu ievērojamu politisko spēku kā nacionāldemokrāti un viņu sabiedrotie. Savu attieksmi pret J.~Pilsudski uz ļauno pusi mainīja arī virkne PPS vadītāju, īpaši tie, kuri bija iestājušies partijā pēc tam, kad viņš sāka no tās attālināties. J.~Pilsudskim, kura rīcībā nebija savas politiskas partijas, tas mazināja iespēju ietekmēt Seimu brīdī, kad pirmajā vietā izvirzījās jautājums par nākamo valsts iekārtu un viņa paša vietu tajā.

Dienu pirms Rīgas miera noslēgšanas ar Padomju Krieviju 1921.~gada 17.~martā pēc ilgiem strīdiem Polijas Seims pieņēma konstitūciju (parasti tā arī tiek saukta par \strong{``Marta konstitūciju''}). Par paraugu bija ņemta t.~s. Trešās Francijas Republikas konstitūcija. Nacionāldemokrāti vēlējās piešķirt konstitūcijai vairāk konservatīvu raksturu: ieviest divpalātu parlamentu, likumdošanas tiesības piešķirt vienīgi augstākajām likumdevējām iestādēm, ierobežot vietējo pašvaldību pilnvaras. J.~Pilsudska atbalstītāji savukārt noraidoši izturējās pret parlamenta augšpalātas~--- Senāta izveidi, taču gala rezultātā piekrita tam, ka Senāts tika ievēlēts analoģiski Seimam, un Senātam nebija būtisku pilnvaru. PPS, ``\pltxti{Wyzwolenie}'' u.c. partijas centās piešķirt konstitūcijai vairāk demokrātisku raksturu~--- radīt vienpalātas parlamentu, paredzēt likumdošanas iniciatīvu arī tautai, referendumu organizēšanu, plašākas tiesības pašvaldībām. Rezultātā pieņemtajai konstitūcijai bija kompromisa raksturs.

Polijas parlaments apstiprināja principiālos lēmumus par valstiskuma atjaunošanu. Konstitūcija ieviesa parlamentāro sistēmu. Polijas parlaments sastāvēja no divām palātām~--- Seima no 444~deputātiem un Senāta no 111~deputātiem, kurus vēlēja uz pieciem gadiem.

Galveno varu valstī konstitūcija nodeva Seima rokās. Likumdošanas iniciatīva piederēja tikai Seimam, ar neuzticības balsojumu tas jebkurā brīdī varēja gāzt valdību, kura bija atbildīga parlamenta priekšā, bet otrā palāta~--- Senāts~--- līdzdarbojās likumdošanas laukā (varēja atlikt likumu stāšanos spēkā).

Parlamenta abas palātas kopā~--- kopējā Seima un Senāta sēdē~--- uz 7~gadiem ievēlēja valsts prezidentu. Viņš izvirzīja premjerministra un ministru kandidatūras, pārstāvēja valsti starptautiskajā arēnā. Taču prezidentam nebija tiesību uzlikt \latxti{veto} Seima pieņemtajiem likumiem un tiesību pirms termiņa atlaist parlamentu (to viņš varēja izdarīt tikai ar 3/5 Senāta locekļu piekrišanu). Faktiski prezidenta funkcijas bija vairāk reprezentatīvas, jo parlaments nevēlējās pieļaut J.~Pilsudska, kurš izrādīja parlamentam maz cieņas, kļūšanu par stipru valsts galvu. Prezidents bija atbildīgs Seimam un nevarēja vienlaicīgi būt par armijas virspavēlnieku. Šis pēdējais nosacījums bija vērtējams kā atklāta pļauka J.~Pilsudskim. Formāli prezidents atradās izpildvaras priekšgalā, bet valdība viņam nepakļāvās. Vēl vairāk, viņš nevarēja veikt savas konstitūcijā noteiktās pilnvaras bez valdības piekrišanas, jo pastāvēja t. s. kontrasignācijas tiesības.

J.~Pilsudski īpaši satrauca jautājums par armijas nākotni. Arī pēc karadarbības beigām austrumos un Austrumsilēzijā viņš neuzskatīja Polijas nākotni par stingri nodrošinātu. Spēcīgāko valsts drošības balstu maršals saskatīja tieši armijā. J.~Pilsudskis nevarēja samierināties ar faktu, ka konstitūcijas autori ieviešot tajā visu bez izņēmuma, valdības locekļu, tai skaitā arī kara ministra, atbildības principu parlamentam, automātiski padarīja šo amatu par politisko, atkarīgu no partijām, to cīņām, kompromisiem, konjunktūras. Maršals bija pārliecināts, ka, būdams politiska figūra, kara ministrs neizbēgami ienesīs politiku armijā, kamēr tai, pēc viņa domām, bija jāpaliek apolitiskai. Tāpēc armijas izvešana ārpus kārtējās politiskajā cīņas, bija galvenais viņa rūpju objekts laika periodā no 1921.~gada janvāra līdz diktatūras izveidošanai 1926.~gadā.

Prezidenta nozīmētā valdība bija atbildīga Seima priekšā, kas bija sašķelts daudzās sīkās, savstarpēji konkurējošās frakcijās, kuras nespēja radīt stabilu parlamentāru vairākumu. Konstitūcija paredzēja demokrātiskas vēlēšanas, personas neaizskaramību, tiesu neatkarību, vienlīdzīgas tiesības vīriešiem un sievietēm. Senāta pilnvaras bija visai šauras, tam nebija likumdošanas iniciatīvas tiesību, Seima pieņemtos likumus tas varēja apturēt tikai uz 60~dienām.

Valsts administratīvajā uzbūvē tika pasludināts decentralizācijas princips ar samērā plašām pašvaldību tiesībām. Vietējā pašvaldība bija paredzēta gminu (\pltxti{gmina}~--- vietējā kopiena, pagasts), povjatu (\pltxti{powiat}~---  apriņķis) un vojevodistu (\pltxti{województwo}~--- apgabals) līmenī. Pašvaldību izpildvaras un kontrolējošām iestādēm~--- padomēm un seimikiem~--- bija jāveidojas vēlēšanu rezultātā. Tās savukārt nozīmēja izpildvaras iestādes~--- valdes, pilsētās~--- maģistrātus. Taču darbojās šīs iestādes centrālās varas pārstāvju uz vietām vecāko (\pltxti{voight}), stārastu (\pltxti{wódz}) un vojevodu (\pltxti{wojewóda)} uzraudzībā. Tiesa, no jauna pašvaldības tika veidotas tikai bijušajā Polijas karalistē, kur agrāk tādu nebija. Citur to organizācija, tās iestāžu sastāvs un administratīvais iedalījums maz mainījās. Tikai Galīcijā, kur muižas agrāk veidoja patstāvīgas administratīvās vienības, tagad tās tika pievienotas lauku gminām.

Polijas valsts administratīvais iedalījums veidojās pakāpeniski. 1919.~gada 2.~augustā tika pieņemts pagaidu likums (faktiski darbojās līdz 1937.~gadam), uz kura pamata no Polijas karalistes guberņām tika izveidotas 5~provinces. No bijušajiem Prūsijas/Vācijas valdījumiem tika izveidotas divas provinces. 1920.--1921.~gadā tika veidotas provinces Silēzijā un Galīcijā. Tikai 1925./1926.~gadā tika galīgi noteikts Viļņas apgabala statuss.

Galu galā Visa Polijas \strong{teritorija tika sadalīta 17 vojevodistēs} (Skat. \ref{tab:table4}.~tabulu), pie tam Varšava bija izdalīta atsevišķā administratīvā vienībā ar vojevodistes tiesībām. Vairākās austrumu un dienvidu vojevodistēs~--- Staņislavas, Tarnovas, Volīnijas un Poļesjes pat pēc oficiāliem datiem iedzīvotāju vairākums bija ukraiņi un baltkrievi. Vojevodistes dalījās vidējās un zemākās administratīvās pārvaldes vienībās: povjatos un gminās.

\noindent
\begin{table}[h!]
\caption{Polijas administratīvi teritoriālais iedalījums starpkaru periodā} \label{tab:table4}
\begin{tabularx}{\linewidth}{|p{2.5cm}|p{2.5cm}|p{2.5cm}|p{2.5cm}|}
\hline
\strong{Vojevodiste} & \strong{Administratīvais centrs} & \strong{Teritorijas lielums tūkstošos km$^{2}$ (1930.~gadā)} & \strong{Iedzīvotāju skaits tūkstošos (1931.~gadā)} \\
\hline
Belostokas (\pltxti{Białostockie}) & Belostoka (\pltxti{Białostok}) & 26,0 & 1~263,3 \\
\hline
Kelces (\pltxti{Kieleckie}) & Kelce (\pltxti{Kielce}) & 22,2 & 2~671,0 \\
\hline
Krakovas (\pltxti{Krakowskie}) & Krakova (\pltxti{Krakow}) & 17,6 & 2~300,1 \\
\hline
Ļubļinas (\pltxti{Lubelskie}) & Ļubļina (\pltxti{Lublin}) & 26,6 & 2~116,2 \\
\hline
Ļvovas (\pltxti{Lwowskie}) & Ļvova (\pltxti{Lwow}) & 28,4 & 3~126,3 \\
\hline
Lodzas (\pltxti{Łόdzkie}) & Lodza (\pltxti{Łόdź}) & 20,4 & 2~650,1 \\
\hline
Novogrudokas (\pltxti{Nowogrόdzkie}) & Novogrudoka (\pltxti{Nowogrόdek}) & 23,0 & 1~057,2 \\
\hline
Poļesjes (\pltxti{Poleskie}) & Bresta (\pltxti{Brećś nad Bugiem}) & 36,7 & 1~132,2 \\
\hline
Pomožes (\pltxti{Pomorskie}) & Toruņa (\pltxti{Toruń}) & 25,7 & 1~884,4 \\
\hline
Poznaņas (\pltxti{Poznańskie}) & Poznaņa (\pltxti{Poznań}) & 28,1 & 2~339,6 \\
\hline
Staņislavovas (\pltxti{Stanisławowskie}) & Staņislavova (\pltxti{Stanisławόw}) & 16,9 & 1~480,3 \\
\hline
Varšavas pilsēta (\pltxti{Miasto Warszawa}) & Varšava (\pltxti{Warszawa}) & 0,14 & 1~179,5 \\
\hline
Varšavas (\pltxti{Warszwskie}) & Varšava (\pltxti{Warszawa}) & 31,7 & 2~460,9 \\
\hline
Silēzijas (\pltxti{Ślaskie}) & Katovice (\pltxti{Katowice}) & 5,1 & 1~533,5 \\
\hline
Tarnopoles (\pltxti{Tarnopolskie}) & Tarnopole (\pltxti{Tarnopol}) & 16,5 & 1~600,4 \\
\hline
Viļņas (\pltxti{Wileńskie}) & Viļņa (\pltxti{Wilno}) & 29,0 & 1~276,0 \\
\hline
Volīnijas (\pltxti{Wołỷńskie}) & Lucka (\pltxti{Łuck}) & 35,7 & 2~085,6 \\
\hline
\end{tabularx}
\end{table}

% TODO add coats of arms
% doc 2, page 462


Konstitūcija saglabāja tiesības uz privātīpašumu, taču atļāva valsts interesēs tā piespiedu atsavināšanu, zemes pirkšanas un pārdošanas regulēšanu. Par agrāro reformu nekas nebija teikts, vienīgi bija atbalstīta ideja par nepieciešamību mainīt agrāro iekārtu valstī.

Konstitūcija pasludināja apziņas brīvību, taču bija uzsvērta katoļu kā valsts iedzīvotāju vairākuma ticības īpašā vieta. (Jau 20.~gadu sākumā latviešu vēsturnieks un diplomāts Alfrēds Bīlmanis atzīmēja, ka konstitūcijas atsevišķā pantā rakstītais, ka ``katoļu ticība ir valsts ticība'' nozīmēja, ka baznīca nav šķirta no valsts.) Tiesa, tā kā nepastāvēja ierobežojumi citu reliģisko konfesiju darbībai, to piekritēji kādas negatīvas sekas neizjuta. Visas reģistrētās baznīcas saņēma valsts dotācijas. Apziņas brīvības principam gan neatbilda lēmums par obligātu reliģijas mācību skolās.

Konstitūcijā bija ietverti arī panti par vārda, preses (cenzūra tika aizliegta), sapulču un organizēšanās brīvību, darba aizsardzību, pilsoņu sociālo nodrošināšanu slimības un bezdarba gadījumā. (Jau 1918.--1921.~gados Polijā bija pieņemti vairāki sociālie likumi, starp citu, deklarēta 46~stundu darba nedēļu. Tiesa, ik pa laikam tika mēģināts šo normu caurskatīt.) Pēc Konstitūcijas pieņemšanas tika pārtrauktas jebkuras diskusijas par rūpniecības un citu nozaru uzņēmumu nodošanu valsts īpašumā.

Poļu tiesību vēstures speciālists V.~Komarņickis XX gadsimta 20.~gados, vērtējot parlamenta dominējošo lomu 1919.~gada 20.~februārī pieņemtajā t.s. Mazajā konstitūcijā un 1921.~gada Marta konstitūcijā, to uzskatīja par raksturīgu iezīmi tautām, kuras bija tikko atbrīvojušās no ārējiem vai iekšējiem politiskajiem spaidiem. Reizē pētnieks pauda viedokli, ka parlamenta pārsvaru valsts varas jomā Polijā lielā mērā veicināja arī dziļi vēsturiski un psiholoģiski cēloņi. Jau šļahtas visvarenības laikos Polijā izveidojās noraidoša attieksme pret stipru izpildvaru. Nacionālās nebrīves gados, kad poļi kļuva par nacionālās apspiešanas objektu, gadsimtu gaitā kultivētā negatīvā attieksme pret varu dziļi iesakņojās sabiedrības apziņā. Rezultātā poļu prātos valdīja negatīva attieksme pret jebkuru varu un pēc Polijas valsts radīšanas 1918.~gadā maz kas bija mainījies. Pēc autora domām, minētais secinājums gan nav neapstrīdams. To pieņemot pilnībā, būtu grūti izskaidrot lielas poļu tautas daļas entuziasmu cīņā par savas valsts izveidi, kā arī pēc dažiem gadiem notikušo valsts apvērsumu un faktisko parlamentārās varas nozīmes krasu mazināšanu.

Jāpiebilst, ka vēsturnieki bieži runā par Polijas valsts atjaunošanu, par otro Žečpospolitu, nevis jaunas valsts izveidi. Tas nav pareizi, jo pirmā Žečpospolita nebija poļu, bet gan divu tautu: poļu un lietuviešu valsts. 1918.--1921.~gadā Polijas valsts veidojās uz atšķirīga no pirmās Žežpospolitas nacionālā pamata.

Vēl pirms Konstitūcijas pieņemšanas 1921.~gada februārī valsts priekšnieka un armijas virspavēlnieka J.~Pilsudska oficiālas vizītes laikā Parīzē starp Franciju un Poliju tika parakstīts politiskas savienības līgums un slepena militāra konvencija. Vācijas agresijas gadījumā pret kādu no šīm valstīm abas puses apņēmās darboties kopīgi. ``Gadījumā, ja Polijai draudētu karš ar Padomju Republiku, vai pēdējās uzbrukuma gadījumā'', Francija apņēmās palīdzēt Polijai nodrošināties pret Vāciju un ``sniegt viņai [Polijai] palīdzību aizsardzībā pret Padomju armiju''. Polija apņēmās uzturēt miera laikā 9~kavalērijas brigādes, minimums 30~kājnieku divīziju, kā arī atbilstošu skaitu vispārējo dažādu ieroču šķiru rezerves formējumu (lauka un smago artilēriju, tankus, sapieru karaspēku, aviāciju u.c.), attīstīt savu bruņojuma rūpniecību. Jākonstatē, ka galvenā konvencijas virzība bija vērsta pret Vāciju, taču tas, ka nebija noteikts ko īsti nozīmē ``gadījums, ja Polijai draudētu karš ar Padomju Republiku'', atstāja diezgan lielas iespējas arī Polijas politiķu un militāristu avantūrām austrumu virzienā. 1921.~gada 3.~martā Polija arī ar Rumāniju noslēdza militāru konvenciju, paredzot kopīgu cīņu pret ``agresiju no austrumiem'', tātad pret Padomju Krieviju. Ar šiem līgumiem J.~Pilsudskis ievirzīja Polijas ārpolitiku konsekventā pretpadomju gultnē. Acīmredzot tādu politiku Polija ieturētu pat tad, ja Krievijā būtu uzvarējuši ``baltie'', pārāk lielas bija poļu bailes no varenas Krievijas izveides. J.~Pilsudskis un viņa līdzgaitnieki starp Vāciju un Krieviju, ``lai kāda arī tā nebūtu'' (J.~Pilsudskis), izvietotās Polijas ģeopolitisko stāvokli uztvēra kā pastāvīgu apdraudējumu.

Vēl līdz šai dienai vēstures literatūrā notiek domu apmaiņa par to, kādi faktori bija svarīgākie Polijas valsts izveidē 1918.--1921.~gadā. Jau tikko nosauktā posma laikā poļu vēsturnieki sadalījās tajos, kuri par galveno neatkarības izcīnīšanas priekšnoteikumu uzskatīja labvēlīgu starptautisko konjuktūru un tajos, kuri uzsvaru lika uz pašu poļu gadu desmitiem ilgušajiem atbrīvošanās pūliņiem. 20.--30.~gados, īpaši t.s. sanācijas režīma laikā liela vēsturnieku daļa atzina, ka J.~Pilsudskim pieder galvenie nopelni valsts atjaunošanā. Citi, ideoloģiski un politiski vairāk saistīti ar nacionāldemokrātu partiju un tās līderi R.~Dmovski, atzina, ka tieši viņš spēja panākt Rietumu valstu (Francijas, Lielbritānijas, Itālijas, un, pirmkārt, ASV) vadītāju atbalstu Polijas valsts atjaunošanai. Taču pats R.~Dmovskis 1920.~gada februārī rakstīja: ``Ja mēs [nacionāldemokrāti] radītu valdību, teiksim, ar mani priekšgalā, mēs Poliju varētu pazudināt.

\begin{enumerate}
\item Neaizmirstiet, ka mēs dzīvojam revolucionārā laikmetā, kurš kādu laiku vēl turpināsies. Bet mēs neesam revolucionāri.

\item Mūs neieredz visi revolucionārie strāvojumi, kuri pret mums tikai stiprāk mobilizētos un niknāk cīnītos nekā pret valdībām, kuras radīja Pilsudskis\dots{}'' Autoram šis spriedums šķiet pamatots.
\end{enumerate}

Mūsdienās poļu vēsturnieku vairākums uzskata, ka J.~Pilsudskis, R.~Dmovskis, V.~Vitoss, I.~Paderevskis ir Polijas neatkarīgās valsts 1918.~gadā radītāji, kaut līdz šai dienai vieni galvenos nopelnus pieraksta J.~Pilsudskim, otri~--- R.~Dmovskim (pēdējo gan kļūst arvien mazāk).

Starpkaru laikā tikai ļoti nedaudzi poļu pētnieki norādīja uz revolūcijām Krievijā un Vācijā kā poļu atbrīvošanās centieniem labvēlīgu starptautisku faktoru nozīmi. Turpretī šā laika vācu vēsturiskajā literatūra var atrast apgalvojumus, ka poļiem valsts iekritusi klēpī kā dāvana no augšas, sabrūkot agrāk pastāvējušajām to sadalījušajām impērijām, par ko poļu politiķi nevarējuši sapņot pat pašos drosmīgākajos sapņos.

PSRS oficiālā vēstures doktrīna uzskatīja, ka ``Polija par savu izveidi un eksistenci ir pateicību parādā PSRS: tā radās un tika uzturēta kā vienots vesels mākslīgi, kā pretsvars Padomju Krievijai/PSRS. Ja PSRS vietā būtu kāds politisks buržuāziska tipa veidojums, Lielbritānija un Francija, iespējams, ar tādu pat vieglumu piekristu jaunai Polijas sadalei starp kaimiņiem, atstājot varbūt kādu ``neatkarīgu'' etnogrāfisku Poliju ar 10--12~miljoniem iedzīvotāju, līdzīgi kā tās Pirmā pasaules kara laikā piekrita atdot visu Poliju cariskās Krievijas varā.''

Pēc Otrā pasaules kara, t.s. ``sociālisma'' laikā arī oficiālā Polijas vēstures versija noteica, ka Polija varēja iegūt neatkarību tikai pateicoties V.~Ļeņinam, krievu boļševikiem un Oktobra revolūcijai. Nopietni, patstāvīgi domājoši vēsturnieki šo faktoru neignorēja, bet arī neabsolutizēja. Izcilais XIX gadsimta un XX gadsimta sākuma perioda Polijas vēsturē pētnieks S.~Keņevičs rakstīja, ka vispirms Vācija uzvarēja karā Krieviju, kuru sagrāva revolūcija, pēc tam pati Vācija cieta sakāvi no Francijas, Lielbritānijas, ASV un citu sabiedroto armijām. Pie tam vēsturnieks uzsvēra, ka visas trīs valstis, kuras sadalīja Poliju, pazuda no politiskās skatuves tajās notikušo revolūciju rezultātā. Netika noliegta arī poļu tautas brīvības cīņu nozīme. Domājams, ka šādai izsvērtai pieejai mūsdienās ir vislielākās iespējas kļūt par dažādus viedokļus vienojošu pamatu.

\section{Iekšpolitiskā attīstība neatkarības periodā. 1921--1939}

\epigraph
{Zelta laiki allaž ir pagājušie laiki.}
{Teodors Fontāne (\entxti{Theodor Fontane})}

\epigraph
{Visvairāk nekārtības rada tie, kuri ievieš kārtību.}
{Tadeušs Kotarbiņskis (\pltxti{Tadeusz Marian Kotarbiński})}

\epigraph
{Jo sliktāk iet lietas, jo spilgtāka pasaka par tām.}
{Oļess Buzina (\uktxti{Олесь Олексійович Бузина})}

\epigraph
{Bez šaubām, varas dievišķošana ir viens no sliktākajiem cilvēku piekoptās elku dievību pielūgšanas veidiem, palieka no apspiestības un verdzības laikiem.}
{K.~Poppers (\detxti{Karl Raimund Popper})}

\epigraph
{Cilvēki, kurus dievina kā dievus, ar laiku tiešām zaudē cilvēciskos vaibstus.}
{Staņislavs Ježijs Lecs (\pltxti{Stanisław Jerzy Lec})}

\epigraph
{Tā valsts ir stipra, kurai par vairogu kalpo taisnīgums.}
{Platons (\eltxti{Πλάτων})}

\epigraph
{Valstis iet bojā tad, kad nespēj atšķirt labus cilvēkus no sliktiem.}
{Antisfēns no Atēnām (\eltxti{Ἀντισθένης})}

\epigraph
{Galvenais komunisma propagandists bija pats kapitālisms~--- ar tā īstenoto, bet ne komunistu izdomāto ekspluatāciju, ar tā krīzēm, bezdarbu, pērkamo politikānismu, kariem.}
{Jevgeņijs Jevtušenko (\rutxti{Евгений Евтушенко})}

\epigraph
{Ja valsts mašīnas daba ir tāda, ka tā spiež jūs piedalīties netaisnā lietā, es varu teikt tikai vienu: neievērojiet likumus.}
{Henrijs Deivids Toro (\entxti{Henry David Thoreau})}

\epigraph
{Cik gan valstis, kurās nav taisnīguma, ir līdzīgas laupītāju baram. Tātad, kas gan ir valstis bez taisnīguma, ja ne laupītāju bari, tāpat kas gan ir laupītāju bari, ja ne valstis miniatūrā?}
{Svētais Augustīns (\latxti{Aurelius Augustinus Hipponensis})}

\epigraph
{Labsirdība vēl nekad nav atņēmusi spēkus un stingrību brīviem cilvēkiem. Lai būtu stiprai, nācijai nav jākļūst nežēlīgai.}
{Franklins Delano Rūzvelts (\entxti{Franklin Delano Roosevelt})}

\epigraph
{Nežēlība raksturīga likumiem, kurus diktējušas bailes.}
{Kārlis Markss (\detxti{Karl Heinrich Marx})}

\epigraph
{Neviens cilvēks necīnās pret brīvību,~--- cilvēki cīnās, pats lielākais, pret brīvību citiem.}
{Kārlis Markss (\detxti{Karl Heinrich Marx})}

\epigraph
{Spēks vienmēr pievelk cilvēkus ar zemu morāli, ģeniālus diktatorus nenovēršami nomaina nelieši. Tajā ir sava likumsakarība.}
{Alberts Einšteins (\detxti{Albert Einstein})}

\epigraph
{Neuzticies pūlim, kas Tev gavilē, viņi bļautu tik pat skaļi arī tad, ja tevi pakārtu.}
{Olivers Kromvels (\entxti{Oliver Cromwell})}

\epigraph
{Karu ved valdošā grupa pret saviem padotajiem, un kara mērķis nav izvairīties no savas teritorijas zaudēšanas, bet saglabāt valdošo iekārtu.}
{Džordžs Orvels (\entxti{George Orwell})}

\epigraph
{Īstam patriotam vienmēr ir jābūt gatavam aizstāvēt savu zemi pret savu valdību.}
{Eduards Ebbi (\entxti{Edward Paul Abbey})}

\epigraph
{Veselas tautas aptvertu šausmas, ja tās uzzinātu, cik sīki cilvēciņi pār tām valda.}
{Šarls Moriss de Taleirāns-Perigors (\frtxti{Charles-Maurice de Talleyrand-Périgord})}

\epigraph
{Demokrātijā muļķim ir tiesības balsot, diktatūrā valdīt.}
{Bertrans~Rasels (\entxti{Bertrand Arthur William Russell})}

\epigraph
{Pats lētākais lepnums ir nacionālais lepnums. Ar to inficētajā subjektā tas atklāj individuālo īpašību trūkumu, ar kurām viņš varētu lepoties; jo pretējā gadījumā viņš nebūtu pievērsies tam, kas bez viņa ir kopīgs vēl daudziem miljoniem cilvēku. Tie, kuriem ir lielas personiskās priekšrocības, kuri pastāvīgi novēro savu tautu, pirmkārt pamanīs tās trūkumus. Bet nožēlojams cilvēks, kam nav ar ko lepoties, tver vienīgo iespējamo un lepojas ar nāciju, kurai viņš pieder; viņš ar aizkustinājuma jūtām ir gatavs aizstāvēt visus viņas trūkumus un muļķības.}
{Аrturs Šopenhauers (\detxti{Arthur Schopenhauer})}

\epigraph
{Nacionālisms var būt milzīgs. Bet dižens~--- nekad.}
{Staņislavs Ježijs Lecs (\pltxti{Stanisław Jerzy Lec)}}

\epigraph
{Nacionālisms~--- tas vienmēr ir provinciālisms.}
{Oļess Buzina (\rutxti{Олесь Алексеевич Бузина})}

\epigraph
{Auto firmu pārstāvji pārdod auto, apdrošināšanas firmu pārstāvji~--- apdrošināšanu, bet ko pārdod tautas pārstāvji?}
{Staņislavs Ježijs Lecs (\pltxti{Stanisław Jerzy Lec)}}

\subsection{Parlamentārais posms. 1921--1926}

% page 466


Pēc jaunās valsts izveides 20.~gados sākās dažādo poļu apdzīvoto apgabalu konsolidācija, tika veidota vienota tautsaimniecības, komunikāciju, izglītības sistēma. Ne visai ievērojami, tomēr pieauga migrācijas procesi zemes iekšienē. Piemēram, augstākajās mācību iestādēs mācījās jaunieši no dažādiem Polijas novadiem.

Poļu publicistikā bieži tēlots, ka pēc Polijas neatkarības atgūšanas kalpot tai ieradās poļu patrioti no citām valstīm, kuri tur bija jau nodrošinājuši sev karjeru un sabiedrisko stāvokli. Šādam priekšstatam atbilst tādas personības kā G.~Narutovičs, I.~Moscickis, kuri bija dzīvojuši un strādājuši Šveicē un par kuriem ies runa turpmāk, arī vēl daži citi. Taču lielākā daļa to poļu, kuri atgriezās dzimtenē, nāca no trijām bijušajām Poliju sadalījušajām valstīm. Tās visas bija smagi cietušas karā, un tur attieksme pret poļiem, kuriem tagad bija sava valsts, daudzviet bija pasliktinājusies. Bijusī plašā Austroungārija bija sarukusi līdz mazajai Austrijai, tai nebija vajadzīgs nedz tik liels ierēdņu aparāts, nedz armija kā agrāk, kad tur redzamu vietu ieņēma poļu ierēdņi un virsnieki. No Padomju Krievijas bēga pat krievu inteliģence, arī attieksme pret poļu pilsonību pēc poļu-padomju kara tur bija nelabvēlīga. Tai piederīgajiem visdrošāk likās atgriezties dzimtenē. Viņi bija tās patrioti, kaut arī to memuāros bieži ieskanējās nostaļģija pēc agrākajiem gadiem, kad Polijas valsts nepastāvēja, bet viņu personīgā dzīve bija labāk nodrošināta.

No Galīcijas (bijušajām Austroungārijas zemēm) poļu ierēdņu kadri, mācību iestāžu pasniedzēji lielā skaitā devās strādāt uz citiem Polijas novadiem. Kopumā Polijas jaunā ierēdniecība veidojās no dažādu valstu ierēdniecības un pašu poļu joks par to vēstīja, ka no Prūsijas tā mantojusi rupjību, no Krievijas korumpētību un no Austrijas birokrātismu. Pēdējais bija pati raksturīgākā jauno poļu ierēdņu iezīme, jo tieši Galīcijas poļu administrācijas ierēdņi savas labākās sagatavotības rezultātā parasti ieņēma vadošos amatus. Pie tam, kā norādījis poļu tiesību vēsturnieks K.~Gržibovskis, pēckara ekonomiskā sabrukuma apstākļos, kad valsts kase bieži bija tukša, uz amatiem, kuri prasīja lielus reprezentācijas izdevumus, varēja pretendēt tikai turīgi ļaudis. Tāpēc tādu rokās nonāca diplomātija, virsniecības un ierēdniecības augstākie posteņi u.c.

Par visas poļu vienotības apziņas formēšanas centru kļuva armija. Tiesa, šī apziņa vairāk tika veidota jauniesaucamajos, agrāko impēriju virsnieki saglabāja savas ambīcijas, pārmeta kolēģiem, kas bija dienējuši citu valstu armijās, ka tie turpinot sekot Vācijas, Austroungārijas vai Krievijas armijas nederīgajām tradīcijām, esot neprofesionāli. Tāpat kā starp ierēdņiem arī virsnieku vidū vadošo vietu ieņēma dienējušie Austroungārijas armijā, pirmkārt tās poļu leģionos. Vācu literatūrā atzīmēta augstāko poļu virsnieku: pulkvežu un ģenerāļu salīdzinoši augstā politizētības pakāpe, kad tos varēja saukt ne tikai par politizētiem militāristiem, bet jau par militarizētiem politiķiem. Sākotnēji taču J.~Pilsudskim sekojošie poļu leģionu virsnieki bija studenti vai jauni speciālisti: juristi, mediķi, inženieri, zinātnieki, mākslinieki. Kara gados viņi ātri sasniedza augstas pakāpes armijā. Pēc Polijas valsts atjaunošanas, vai nu paliekot armijā, vai arī atgriežoties civildzīvē, viņi parasti turpināja nodarboties ar politiku. J.~Pilsudskis arī pēc poļu atbrīvošanas kara beigām centās tā vai citādi saglabāt savu iespaidu armijā, jo tā viņam kalpoja ne tikai ārpolitisko, bet arī iekšpolitisko mērķu īstenošanai~--- nācijas vienības stiprināšanai, savas ietekmes valstī izplatīšanai. Ar to varēja apspiest strādnieku un zemnieku uzstāšanās. Tā notika 1923.gadā Krakovā, Borislavā (\pltxti{Borysław}) un Tarnovā (\pltxti{Tarnów}), 1930.~un 1932.~gadā ukraiņu apdzīvotajās teritorijās, 1937.~gadā zemnieku nemieru laikā, kad karavīri saņēma pavēlēs šaut uz savas valsts iedzīvotājiem.

Liela loma valstī bija Polijas armijas Ģenerālštāba Otrajai nodaļai (\pltxti{Oddział II Sztabu Generalnego Wojska Polskiego}), kura nodarbojās ar izlūkošanu ārzemēs un ar pretizlūkošanu, tāpēc bija cieši saistīta ar politisko darbību. Daudzi ievērojami politiskie darbinieki zināmu laiku bija strādājuši šai nodaļā.

Polijas zemnieki, īpaši nabadzīgie un algādži, ļoti labprāt devās dienēt armijā. Tajā viņiem bija nodrošināts regulārs ēdiens, apģērbs, pat neliels atalgojums. Polijas armijas rietumu garnizonos dienošie austrumprovinču zemnieki nereti pirmo reizi dzīvē redzēja ūdensvadu un kanalizāciju, elektrisko apgaismojumu, tramvaju un daudz ko citu. Salīdzinājumā ar iepriekšējām ikdienas grūtībām dienests armijā daudziem likās kā Dieva dāvana.

1921.~gadā demokrātiski vēlētais Polijas Republikas Seims atzina visu pilsoņu vienlīdzību likuma priekšā un ar to formāli apstiprināja dižciltīgo privilēģiju atcelšanu. Taču savas īpašās identitātes apzināšanās ļāva šļahtai jau iepriekš pārdzīvot ne vienu vien katastrofu. Arī pēc Polijas valsts atjaunošanas šļahtiči saglabāja savu zināmo nošķirtību, kārtas lepnumu.

1921.~gada ``Marta konstitūcija'' atcēla tās agrāk pastāvošās tiesiskās normas, kas bija pretrunā ar jauno iekārtu. Tomēr bijušo dažādo valstu teritoriju organisku saplūšanu traucēja dažādo tiesību sistēmu ilgstoša saglabāšanās. Dažādās jomās darbojās agrāko metropoļu likumi. Vēl ilgi atsevišķās vojevodistēs darbojās līdz tam pastāvošie austriešu, vācu un krievu civilie un krimināllikumi, bet 10~bijušajās Polijas karalistes guberņās krievu civillikumu vietā~--- \entxti{Code Napoleon}. Likumu ievērošanu uzraudzīja ap 50~tūkstošu dažāda veida policijas aģentu. Likumu unifikācija ievilkās uz 20~gadiem.

Vēl 1939.~gadā Polijas vācu nelabvēļi, norādot, ka pēc Austrijas ``anšlusa'' (vācu \detxti{Anschluss~--- pievienošana, šai gadījumā Austrijas pievienošana Vācijai}) 1938.~gadā tās likumdošanas pielīdzināšana Vācijas likumdošanai notika dažu mēnešu laikā, pretēji lēnajai likumu pielīdzināšanai Polijā, un no tā secināja, ka nevis agrākās Polijas dalīšanas bija radījušas atšķirības poļu vidū, bet tās pastāvēja jau ilgstoši, arī vēl pirms Polijas dalīšanas. Jāpiekrīt, ka starp dažādiem apgabaliem atšķirības tik tiešām pastāvēja jau pirms dalīšanām, taču tikpat nenoliedzami ir tas, ka dalīšanas tās pastiprināja.

Eiforijas apstākļos par apvienotas, lielas Polijas izveidi iekšējās pretrunas valsts attīstībā sākotnēji nelikās tik bīstamas kā ārējie draudi. Taču politiskās pretrunas, sākot no 1921.~gada, arvien vairāk saasinājās. Ja cīņas par neatkarību laikā strīdīgie jautājumi tika atvirzīti otrajā un trešajā plānā, pēc valsts nodibināšanās tie atkal nostājās priekšplānā. 1921.--1922.~gadā Polijā norisa 1~520~streiku, kuros piedalījās ap 1~miljons cilvēku. 1923.~gadā notika 1~263~streiki ar 849~tūkstošiem dalībnieku. Tie turpinājās arī nākamajos gados.

Atjaunotās Polijas Republikas (\pltxti{Rzeczpospolita Polska} vai arī \pltxti{II Rzeczpospolita}) politiķiem, neraugoties uz dažu poļu deputātu īslaicīgu darbību Vācijas un Austroungārijas parlamentos un vēl īsāko darbību Krievijas Valsts domē, nebija pamatīgas parlamentārisma pieredzes. Tas ietekmēja valsts stabilitāti.

Jau teikts, ka Polija nekļuva par etniski monolītu valsti. Nacionālie mazākumi~--- ukraiņi, baltkrievi, ebreji, vācieši un lietuvieši sajuta sevi kā otrās šķiras pilsoņus. (Par to sīkāk attiecīgajā apakšnodaļā.) Tas radīja spriedzi sabiedrībā. Taču arī pašu poļu dažādu sociālo slāņu un politisko spēku vidū pastāvošās pretrunas 1921.~gada 17.~marta demokrātiskās konstitūcijas darbības apstākļos šķēla sabiedrību. Parlamentārās sistēmas izveide Polijā radīja labvēlīgākus nekā iepriekš apstākļus plašu masu iesaistīšanai sabiedriskajā dzīvē. Pieauga politisko partiju kā nacionālo, šķirisko, grupu u.c. interešu izteicēju loma.

Sevišķi sabiedrības konsolidāciju traucēja neatrisinātais agrārais jautājums. Lielais daudzums pastāvošo partiju veicināja ātru savstarpējo koalīciju un attiecīgi arī valdību maiņu. Valdību politiskā kursa izmaiņas īpaši 20.~gadu pirmajā pusē pastāvīgi satricināja valsti. (Tomēr minētās pretrunas sabiedriskajā sistēmā 20.~gados vēl nekļuva par lielām briesmām valstij. Tas notika tikai 30.~gados.)

Kā raksta poļu vēsturnieki, paradoksālā veidā strīdi starp sociālistiem, endekiem un ļudoviešiem noveda pie situācijas, kuru nevarēja iedomāties citās valstīs, proti~--- parlamentārisma zināmas hipertrofijas. Proporcionalitātes princips arī bija novests līdz galējībai, jo mazās partijas baidījās, ka citādi pārāk lielu svaru iegūs pati ietekmīgākā~--- nacionāldemokrātu (endeku) partija. (No 1919.~gada tā oficiāli saucās par Nacionāldemokrātisko savienību (\pltxti{Związek Ludowo-Narodowy}), no 1928.~gada līdz pat 1947.~gadam par Nacionālo partiju (\pltxti{Stronnictwo Narodowe}), taču faktiski saglabājoties tam pašam politiskajam kursam, ideoloģijai un lielā mērā arī sastāvam, partiju tautā turpināja saukt tāpat kā iepriekš~--- \pltxti{Endecja}, tās biedrus~--- par endekiem.) Vairs nebaidoties no revolucionārās kustības, nacionāldemokrāti centās saņemt valsts varu. Īpaši viņu frakcijas vadītājs Seimā S.~Grabskis cerēja, ka jaunas Seima vēlēšanas un savienība ar citām partijām dos endekiem vairākumu parlamentā. Taču tas bija maz ticams valstī, kur katrs trešais pilsonis nebija polis un neatbalstīja endeku nacionālismu. Endeki arī aktīvi cīnījās pret J.~Pilsudski. Visas parlamentā pārstāvētās partijas, bet īpaši endeki baidījās no pārliekas prezidenta, uz kura amatu visreālākais kandidāts bija J.~Pilsudskis, varas. 1922.~gada vasarā konflikts starp parlamenta deputātiem un valsts priekšnieku sasniedza tādu asumu, ka endeki 25.~jūlijā iesniedza priekšlikumu par neuzticības izteikšanu J.~Pilsudskim. To atbalstīja 186~deputāti, pret to nobalsoja~206.

20.~gadu sākumā arī zemnieku politiskā kustība bija ieguvusi jau ievērojamu spēku. Zemniecības tiesības uz vadošo vietu sabiedrībā un valstī šo partiju ideologi pamatoja ne tikai ar tās lielo skaitlisko sastāvu un nopelniem nācijas priekšā, bet arī tās it kā unikālajām morālajām, ētiskajām, fiziskajām un sociālajām īpašībām. Tiesa, zemnieku politiskā kustības īpatnība starpkaru periodā palika organizatoriskā, politiskā un ideoloģiskā sašķeltība. To neizdevās pārvarēt arī pēc 1931.~gadā notikušās triju lielāko zemniecības partiju: Poļu Zemnieku partijas~--- \pltxti{Piast} (\pltxti{Polskie Stronictwo Ludowe~--- Piast)}, Poļu Zemnieku partijas~--- \pltxti{Wyzwolenje} (\pltxti{Polskie Stronnitctwo Ludowe~--- Wyzwolenje}) un Zemnieku partijas (\pltxti{Stronnictwo Chłopskie}) apvienošanās. Līdz tam un arī pēc tam poļu zemnieku kustībā pastāvēja vēl dažādi grupējumi ar atšķirīgu orientāciju. Ļudoviešu kustībā starpkaru periodā pastāvēja dažādi strāvojumi: no konservatīva līdz konsekventi demokrātiskam. Katras no ļudoviešu partijām mēģinājumi pārkāpt reģionālos ietvarus, izplatīt savu ietekmi visā valstī neizbēgami radīja asu konkurences cīņu par biedru bāzi un elektorātu. Nesakrita arī politiskās pozīcijas. Ja Poļu Zemnieku partija \pltxti{Piast} ieņēma centristiskas pozīcijas, tad Poļu Zemnieku partija \pltxti{Wyzwolenje} u.c. bija diezgan kreisi tendētas.

Visi ļudoviešu grupējumi plaši izmantoja terminu \pltxti{lud} (tauta, zemniecība), taču katrs tajā ielika savu, nebūt ne līdzīgu saturu. Tā, \pltxti{Wyzwolenje} sastāvā bez zemniecības iekļāva arī strādniecību un pilsētu sīkburžuāziju. Pēc partijas biedru domām, valstij bija jākalpo tieši šo slāņu interesēm. Toties muižniekiem un fabrikantiem gan tika atzītas tiesības dzīvot neatkarīgajā Polijā, taču ar noteikumu, ka viņi zaudēs visas privilēģijas, ko bija ieguvuši sakarā ar izcelsmi un mantisko stāvokli, un atzīs tautas tiesību prioritāti. Toties Poļu Zemnieku partijas \pltxti{Piast} traktējumā \pltxti{lud}-a sastāvā ietilpa tikai zemniecība. Zemniecības tiesības uz varu Polijā tika pamatotas ar to, ka tā ir šīs zemes varas nesēja, jo vara viņai pieder ``pēc taisnības'', ``zemniecība ir pamats, uz kuru Polija var stingri balstīties''. Partija \pltxti{Piast} gan bija gatava sadarboties arī ar strādniecību un citiem slāņiem, kuri atzina Polijas neatkarību.

Pēc konstitūcijas pieņemšanas 1921.~gada martā Seims nevis atlaida sevi, kā tas bija paredzēts, bet 1921.~gada 18.~maijā pagarināja savu darbības termiņu. Par pamatu šādai Seima rīcībai bija visu lielāko partiju vēlme ietekmēt vēlēšanu gaitu, izmantojot savu pārstāvniecību valdībā, kā arī pieņemt tādu vēlēšanu likumu, kurš traucētu mazāko partiju un nacionālo mazākumu ietekmes pieaugumu. Izvērsās arvien jaunas kombinācijas un intrigas, tāpēc 15~mēnešos no 1921.~gada septembra, kad atkāpās Zemnieku partijas \pltxti{Piast} līdera V.~Vitosa vadītā koalīcijas valdība, līdz 1922.~gada decembrim pie varas pabija četras valdības. Faktiski tās bija ierēdņu valdības ar dažiem partiju reprezentantiem, taču visas tās bija atkarīgas no kombinācijām Seimā.

Pēc poļu-padomju kara visai smagā situācijā atradās PKSP. Pretēji komunistu cerībām Rīgas miera, 1921.~gada konstitūcijas pieņemšanas rezultātā Polijas valsts bija stabilizējusies. Taču radikāli kreisa nostāja partijā dzīvoja vēl ilgi. Partijā izvirzījās jautājums par sabiedrotajiem tās cīņā. Taču to iegūšanu traucēja šauri sektantiskie uzskati par PPS kā buržuāzijas aģentūru strādnieku šķirā, kas atgrūda no PKSP daļu strādniecības. PSKP 20~cilvēku lielā delegācija Kominternes III kongresā 1921~gada jūnijā-jūlijā, kaut galvenajos jautājumos atbalstīja KK(b)P, kongresa lēmumus par masu atbalsta iegūšanas nepieciešamību, vienotu strādnieku šķiras fronti, tomēr reizē demonstrēja visai kreisus uzskatus. Tie tika pārvarēti tikai ar lielām grūtībām. V.~Ļeņinam 1921.~gada 19.~oktobrī nācās nosūtīt poļu komunistiem vēstuli ar brīdinājumu nerīkot priekšlaicīgu sacelšanos pret buržuāziju. Sektantiskie uzskati traucēja komunistiem arī Seima vēlēšanu cīņā.

Vēlēšanu likums tika pieņemts 1922.~gada vasarā un atbilda galvenokārt nacionāldemokrātu vēlmēm. Seimam bija jāsastāv no 444~deputātiem, no kuriem 372~tika vēlēti daudzmandātu apgabalos, bet 72~--- pēc visas valsts saraksta. Mandāti pēdējā tika sadalīti starp partijām un blokiem proporcionāli vēlēšanu apgabalos iegūtajam mandātu skaitam (ja tas bija ne mazāks par 6). Senātā tika vēlēti 111~deputāti. Kandidātu saraksta izvirzīšanai vēlēšanās apgabalos bija nepieciešami tikai 50~vēlētāju paraksti.

Seima vēlēšanās aktīvās tiesības baudīja visi Polijas pilsoņi no 21~gada vecuma, izņemot regulārās armijas karavīrus. Pasīvās vēlēšanu tiesības bija visiem, arī karavīriem, no 25~gadu vecuma. Senāta vēlēšanās varēja piedalīties visi, izņemot karavīrus, no 30~gadu vecuma, ievēlēti varēja tikt ne jaunāki par 40~gadiem. Ievēlētajiem karavīriem dienests bija jāatstāj. Vēlēšanu apgabalos ar poļu tautības vēlētāju vairākumu vajadzēja mazāk balsu parlamenta deputāta mandāta iegūšanai nekā apgabalos ar nepoļu vairākumu. Daudzviet izkaisītajiem vācu un ebreju vēlētājiem gandrīz nebija pat teorētisku iespēju izcīnīt deputāta vietu.

Šāds vēlēšanu likums piespieda nacionālos mazākumus nodibināt kopēju vēlēšanu bloku. Tajā ietilpa visu galveno nacionālo mazākumu: vācu, ebreju, ukraiņu un baltkrievu partijas un grupas. Visas tās prasīja nacionālo vienlīdzību un brīvību. Baltkrievi un ukraiņi bez tam prasīja pārtraukt austrumu rajonu kolonizāciju, poļu pārceļotāju kustības veicināšanu.

1922.~gadā divpalātu parlamenta (5.~novembrī tika vēlēts Seims un 12.~novembrī~--- Senāts) vēlēšanās piedalījās 19~kandidātu saraksti. Seima vēlēšanās balsoja ap 8~813~000~cilvēku jeb 67,7\% no visiem balsstiesīgajiem. Senāta vēlēšanās piedalījās ap 5~589~000~vēlētāju jeb 61,5\% pie vēlēšanām pielaisto. Gandrīz puse austrumu apgabalu iedzīvotāju vēlēšanās nepiedalījās. Tā Galīcijas ukraiņu iedzīvotāji boikotēja vēlēšanas, protestējot pret poļu īstenoto reģiona okupāciju.

Seima vēlēšanās vissekmīgāk uzstājās 1922.~gada 16.~augustā izveidotais \strong{labējo partiju izveidotais bloks}, kurš tika nosaukts par Nacionālās vienotības kristīgo savienību (\pltxti{Chrześcijański Związek Jedności Narodowej}). Blokā ietilpa endeki~--- Nacionāli Demokrātiskās savienība (\pltxti{Związek Ludowo-Narodowy}), Polijas Kristīgi demokrātiskā partija (\pltxti{Polskie Stronnictwo Chrześcijańskiej Demokracji} jeb saīsināti \pltxti{Chadecja}), Kristīgi nacionālā partija (\pltxti{Stronnictwo Chrześcijańsko-Narodowe}) un citas sīkpartijas. Poļu vidū šis bloks parasti ieguva saīsinātu nosaukumu Kristīgi Nacionālā savienība (\pltxti{Chsześcijańska Jedność Narodowa}) jeb vēl īsāk \pltxti{Сhjena}. Izveidotā abreviatūra skanēja līdzīgi kā ``hiēna'' (\pltxti{hiena}), ko tūlīt savā propagandā izmantoja nacionāldemokrātu politiskie pretinieki.

Vēlēšanu kampaņas laikā labējie aicināja izveidot ``nacionālu'' vairākumu, atklāti uzstājās pret ekonomiskās dzīves regulēšanu no valsts puses, nosaucot to par ``sociālismu'', kritizēja ``nesaprātīgo likumdošanu sociālisma garā'', toties aicināja atbalstīt ``kristīgo rūpniecību un tirdzniecību''. Labējo lozungus ``polis~--- nozīmē katolis'', ``valsti~--- poļiem'' liela poļu daļa uztvēra labvēlīgāk nekā kreiso aicinājumus izbūvēt daudznacionālu sociālā taisnīguma valsti.

Bloks nosodīja J.~Pilsudski par tuvību kreisajām partijām, par atbalstu armijā saviem piekritējiem, par ``pārsniedzošu iespējas'' ekspansionismu austrumu virzienā. Bloks vēlēšanās ieguva ap 1/3~mandātu, kas tomēr nedeva viņiem iespēju veidot valdību.

Var atzīmēt, ka arī J.~Pilsudska piekritēji (pilsudčiki) izmēģināja savus spēkus vēlēšanās. Viņi radīja divas mēreni konservatīvas partijas, kuras uzstājās kā valstisko vērtību aizstāves, taču tā arī neieguva vēlētāju atbalstu. Seimā iekļuva tikai tie maršala piekritēji, kuri balotējās citu partiju sarakstos.

Taču no 1922.~gada sākās pilsudčiku politiski-organizatoriska konsolidācija. Sāka iznākt viņu preses izdevumi, bijušie POW (\pltxti{Polska Оrganizacja Wojskowa}) biedri nodibināja Poļu brīvības organizāciju (\pltxti{Polska Organizacja Wolności}). Tās locekļu vairākums pēc Polijas neatkarības atgūšanas atradās diezgan smagā situācijā. Pēc gandrīz septiņiem gadiem ierakumos, ar nepabeigtu augstāko vai pat vidējo izglītību viņi tikai ar grūtībām pielāgojās jaunajiem civilās dzīves apstākļiem, meklēja autoritāti, uz kuru varētu paļauties, un atkal atrada to J.~Pilsudska personā, ar kuru arī saistīja savu dzīves mērķu sasniegšanu. Augustā ar J.~Pilsudska dalību notika I~leģionāru kongress, kurš nolēma pievienoties Poļu brīvības organizācijai. Vēlāk šīs struktūras J.~Pilsudskis varēja izmantot cīņā par varu.

Ļudoviešiem (partijām ``\pltxti{Piast}'', ``\pltxti{Wyzwolenie}'', dažām sīkākām grupām) neizdevās vēlēšanu cīņas gaitā vienoties, kaut viņi visi iestājās par demokrātisko iekarojumu aizsardzību. Ļudovieši ieguva ap 26\% vēlētāju balsu.

Arī PPS propagandēja pilsonisko tiesību aizstāvības nepieciešamību, demokrātisku parlamentārismu uzskatot par izejas punktu ceļā par sociālismu. Taču partijas lomu cīņā par demokrātiju un sociāli ekonomisko progresu vājināja partijas ticība J.~Pilsudskim, kurš, kļuvis par valsts priekšnieku, faktiski bloķēja pirmo (I.~Dašinska un E.~Moračevska vadīto) t.s. ``tautas'' valdību pasludinātos sociālos pārkārtojumus. Rezultātā tie kļuva par abstraktiem postulātiem, kuri vairs nebija PPS ikdienas uzdevumu sarakstā. Partija savā parlamentārajā darbībā aprobežojās ar valsts pastāvēšanas sākumperiodā iekarotās sociālās likumdošanas aizstāvēšanu. Strādniecība bija politiski sašķelta. 1922.~gadā PPS Seima vēlēšanās saņēma 10,3\% balsu.

PKSP uzstādīja tās izveidotas legālas Pilsētu un lauku proletariāta savienības sarakstus 54 no 64 vēlēšanu apgabaliem, aģitēja par muižnieku un baznīcas zemju nodošanu zemniekiem, strādnieku kontroles ieviešanu rūpniecības uzņēmumos, visu tautu līdztiesību. Taču varas iestādes nevēlējās pieļaut komunistu dalību demokrātiskās vēlēšanās, arestēja ap 600~partijas aģitatoru. Tomēr par partiju nobalsoja ap 132~tūkstoši vēlētāju, partija ieguva divas deputātu vietas. Kā atzīmējuši poļu pētnieki Tadeušs un Vitolds Ržepecki, uz vienu iegūto vietu Seimā komunisti bija saņēmuši 60~785 balsis, sociālisti~--- 26~327 balsis, bet ``\pltxti{Chjena}''-s bloks~--- nedaudz vairāk par 18~000 balsu. Nākas secināt, ka par demokrātisku vēlēšanu sistēmu Polijā varēja runāt visai nosacīti.

Kaut vairākas Rietumukrainas partijas vēlēšanas boikotēja, un dažas nacionālo mazākumu partijas uzstājās atsevišķi, nacionālo mazākumu bloks (\pltxti{Blok Mniejszości Narodowych}) saņēma līdz 20\% balsu, un izcīnīja 66~deputātu vietas Seimā. Ebreju deputāti vien ieguva 34~vietas Seimā un 12~Senātā (dažādās publikācijās skaitļi par dažādu partiju un arī nacionālo mazākumu deputātu skaitu ir atšķirīgi, tomēr vēlēšanu ainu tas būtiski nemaina). Tas nekādi neatbilda nacionālo mazākumu vietai sabiedrībā, taču tomēr centristiem un kreisajiem, ja tie vēlējās ko panākt Seimā, vajadzēja ar šo mazākumu rēķināties. Nacionālo mazākumu pārstāvju lielais skaits būtiski mazināja relatīvās labējo uzvaras vērtību.

Tāpat kā agrāk nevienai partijai nebija Seimā vairākuma, kas ļautu uz ilgu laiku noturēt varu un realizēt savu politisko programmu.

Paziņojot, ka prezidenta pilnvaras ir pārāk nenozīmīgas, J.~Pilsudskis atteicās balotēties, jo nevēlējās kļūt par valsts prezidentu, pakļautu Seimam. Tiesa, jau iepriekš minētais padomju vēsturnieks P.~Olšanskis uzskatīja, ka J.~Pilsudskis, apzinoties nacionālo mazākumu netīksmi pret viņa realizēto nacionālo politiku, nebija pārliecināts, par kuru no mazākajiem ļaunumiem: viņu vai nacionāldemokrātu kandidātu, izvēlēsies balsot mazākumu deputāti un tāpēc neriskēja ar sakāvi prezidenta vēlēšanās. Lai nu kā, 1922.~gada decembrī pēc ilgām cīņām parlamentā piektajā vēlēšanu kārtā ar neapstrīdamu vairākumu par \strong{prezidentu} tika \strong{ievēlēts} bezpartijisks zinātnieks ar pasaules atpazīstamību, līdz tam sabiedrisko darbu un ārlietu ministrs, kurš par tādu kļuva ar J.~Pilsudska atbalstu, \strong{Gabriels Narutovičs}. Viņš saņēma 289~balsis, bet viņa konkurents endeks grāfs M.~Zamoiskis~--- 227. G.~Narutoviču atbalstīja kreisās partijas, centristi, arī nacionālo mazākumu pārstāvji. G.~Narutovičs bija ļoti spējīgs cilvēks, pagātnē neatkarīgs uzņēmējs, kļuvis par Cīrihes Tehnoloģiskā institūta profesoru Šveicē. Viņš bija izaudzis demokrātiskās tradīcijās, labi orientējās Rietumeiropas politikā un kultūrā, piedalījies arī Lietuvas politiskajā dzīvē, tās neatkarības proklamēšanā 1918.~gadā. Vēlēšanās viņu atbalstīja arī J.~Pilsudskis.

Poļu nacionālisti, uzstājoties pret Polijas parlamenta lēmumu, tūlīt izraisīja pret prezidentu rīdīšanas kampaņu. Daudzos avīžu materiālos un speciāli organizētos mītiņos pār viņa galvu bira rupjas lamas un apvainojumi. Viņu esot uzspieduši cittautieši, poļu tautai viņa ievēlēšana nozīmējot ``smagu apvainojumu'', ``vardarbību pret poļu politisko domu''. \pltxti{Chjenas} bloka partijas paziņoja, ka ``neatbalstīs valdību, kuru veidos prezidents, kuru [poļiem] uzspieda citas nācijas: ebreji, vācieši, ukraiņi.'' Nacionāldemokrātu prese centās lasītājiem iegalvot, ka G.~Narutovičs pēc nacionālās piederības neesot polis, ``pat neprotot runāt poliski'', esot ``ebreju pasaules kapitāla ieliktenis''. Pilnīgi veltīgi bija skaidrot, ka pēc izcelsmes G.~Narutovičs ir poļu šļahtičs, nacionālistiski noskaņotajam pūlim nebija šaubu par pretējo, saprāta argumenti uz to neiedarbojās. Ksendzs, Satversmes seima un pirmā Seima deputāts K.~Lutoslavskis retoriski jautāja: ``Kā gan ebreji iedrošinājās uzspiest Polijai savu prezidentu?'' Tam sekoja citas publikācijas, kur G.~Narutovičs tika saukts par masonu un bezdievi. Endeki apgalvoja, ka parlamenta lēmumu var neievērot, G.~Narutovičs nedrīkst uzņemties prezidenta pienākumus, jo esot ievēlēts ar ``ebreju balsīm''. Tas nozīmēja, ka zināma daudznacionālās valsts tautas daļa nedrīkst piedalīties tikai poļiem piederošajā politikā.

Poļu labējo rīcību iedvesmoja 1922.~gada oktobrī itāļu fašistu realizētais maršs uz Romu (itāļu: \ittxti{Marcia su Roma}). Prezidenta zvēresta pieņemšanas dienā 11.~decembrī Seimu burtiski aplenca endeku piekritēju pūlis, kurš mēģināja uzcelt improvizētas barikādes, uzbruka deputātiem un senatoriem, kuri atbalstīja jaunievēlēto prezidentu, (daži no viņiem guva savainojumus), gānīja G.~Narutoviču, apmētāja viņu kortežu netīra sniega pikām. Viena no tām trāpīja viņam sejā. Divi huligāni ar nūjām rokās pat uzlēca uz viņa ekipāža pakāpieniem un tikai apsardzes priekšnieks tos nogrūda. Labējiem dauzoņām vēl brutālāk izrīkoties neļāva pēc PPS aicinājuma ieradušies Varšavas uzņēmumu strādnieki. Viens no viņiem, aizstāvot prezidentu, zaudēja dzīvību, daudzi tika ievainoti. Policija to visu vēroja ar apbrīnojamu vienaldzību. Pašā Seimā prezidenta zvēresta nodošanas laikā daudzi soli izrādījās tukši, jo labējie atteicās piedalīties ceremonijā. Bālais prezidents, stipri uztraucies, bet spēcīgā un pārliecinošā balsī nodeva zvērestu uz krucifiksa. Taču tūlīt pēc inaugurācijas ``īsto poļu'' nometne sāka prasīt prezidenta atkāpšanos. Iniciatori bija endeki Seimā.

Neraugoties uz to, ka prezidenta zvērests tika nodots, amatā G.~Narutovičam bija jāstājas pēc tam, kad 14.~decembrī J.~Pilsudskis nodotu viņam valsts galvas pilnvaras. Līdz tam 11.~decembra pēcpusdienā notika konfidenciāla valdības sēde ar J.~Pilsudska piedalīšanos. Viņš paziņoja: ``Es nevaru nodot varu brīdī, kad sūda brāļu banda grauj kārtību, apvaino prezidentu, bet valdība slīgst bezdarbībā; dodiet man pilnvaras un es nomierināšu ielu; ja nedosiet, es iešu viens un nomierināšu~--- šādos apstākļos es nevaru piekāpties.'' Taču valdība nevēlējās lai maršals atstātu rezidenci Belvedēra pilī, apvīts ar nemieru nomierinātāja oreolu. Pilnvaras viņš nesaņēma. Tiesa, viņš arī nedevās viens pats nomierināt nacionālistu pūli, ko neviens viņam neliedza. Viens no nacionāldemokrātu ideologiem S.~Stronskis rakstīja, ka J.~Pilsudskis pametis G.~Narutoviču kā traucēkli ceļā uz valsts atdzimšanu. Līdzīgā tonī tēmu uzturēja arī citas labējo avīzes.

14.~decembrī notika valsts galvas pilnvaru nodošanas ceremonija. Taču, kā raksta Seima priekšsēdētājs M.~Ratajs, naktī no 16. uz 17.~novembri, kad vēl jauna valdība nebija izveidota, virsnieki~--- J.~Pilsudska atbalstītāji bez pretestības sāka ieņemt valdības iestāžu ēkas un pat izpildīt dažas valdības funkcijas. Vienlaikus J.~Pilsudskis atkal deva saprast, ka viņš ir gatavs vadīt valdību un uzsvēra, ka notiks neizbēgamu sadursme starp strādniekiem un endeku kaujiniekiem. Taču PPS vadība, pirmkārt I.~Dašinskis, noraidīja šo plānu, nevēloties strādnieku organizāciju faktisku iesaistīšanu sazvērestībā. Bez strādnieku atbalsta pilsudčiki neuzdrošinājās izraisīt ``spēka'' cīņu.

1922.~gada 16.~decembrī, mākslas izstādes apmeklējuma laikā jaunievēlēto \strong{prezidentu nošāva} neveiksmīgs gleznotājs un mākslas kritiķis-modernists, Kultūras ministrijas ierēdnis, labēji noskaņots fanātiķis E.~Ņevjedomskis, kurš jaunībā bija Nacionālās līgas biedrs. E.~Ņevjedomski zinoši cilvēki uzskatīja par psihiski nelīdzsvarotu personu, kura viegli pakļāvās citu ietekmei. Taču ir dati, ka 1918.~gadā vēl Vācijas radītajā Polijas pārvaldes sistēmā viņš bija glezniecības un skulptūras nodaļas vadītājs, bet poļu-padomju kara laikā dienēja Polijas armijas Galvenā štāba pretizlūkošanas pārvaldē. Tas liek apšaubīt viņa kā nelīdzsvarotas personas lomu. Tiesas procesa gaitā 1922.~gada 30.~decembrī slepkava liecināja, ka pirms tam viņš plānojis nogalināt valsts priekšnieku, bet no jau detalizētā plāna atteicies, kad uzzinājis, ka J.~Pilsudskis neizvirza savu kandidatūru prezidenta amatam. Reportāžas no tiesas zāles labējo avīzes rakstīja par ``E.~Ņevjedomska varonību'', viņa ``svēto pārliecību'', ``cilvēka, kurš par visu vairāk mīlēja Poliju, traģēdiju''. Slepkava tika notiesāts un 1923.~gada 31.~janvārī nošauts. Pēc soda izpildes baznīcās tika noturētas daudzas mesas, aizlūdzot par slepkavu, kurš esot kritis ``par upuri pilsoņa pienākumam''. Slepkavas apbedīšanā piedalījās ap 10~000 cilvēku, nacionālisti viņu pasludināja par mocekli.

Kā raksta mūsdienu poļu vēsturnieki, morāla atbildība par prezidenta slepkavību bija jāuzņemas nacionālistiem, pirmkārt jau endekiem, kuri savā propagandā rupji apvainoja G.~Narutoviču, nevairījās no melīgiem apgalvojumiem, ka viņš esot ``žīdu prezidents'' utt. Taču cerēto viņi nesasniedza, labēji-centrisku vairākumu radīt Seimā neizdevās, jo centristi, kuri balsoja par G.~Narutoviču, norobežojās no endekiem. Pēc G.~Narutoviča nāves politiskajā dzīvē pieauga savstarpējs naids un neiecietība. Prezidenta vēlēšanu norise, un vēl jo vairāk~--- viņa nogalināšana spilgti demonstrēja Polijas parlamentārajai sistēmai piemītošos trūkumus~--- partiju sašķeltību, kreisā un labējā spārna, kā arī poļu vairākuma un nacionālo minoritāšu pretrunas. Ā.~Varskis rakstīja: ``Polijas lāsts ir tas, ka tajā demokrātija pastāv uz papīra, ka šī demokrātija dzima ne buržuāziskā revolūcijā, bet bailēs no proletāriskās revolūcijas \citespace{} ka, neraugoties uz papīra pastāvošo demokrātiju, patiesībā mums ir valsts administrācija, tiesa un armija, ko vada caru, Habsburgu, Hoencollernu juridiskais gars un juridiskais burts.'' Mūsdienu publicistikā ir izteikts uzskats, ka reizē ar pirmo Polijas prezidentu tika nogalināts arī demokrātiska dialoga gars. Grāmatas autoram gan šķiet, ka šis gars bija gana vājš jau valstij veidojoties 1918.~gadā. Uzvara karā pret kaimiņvalstīm ar spēka, bet ne kompromisa, savas valsts daudznacionālās sabiedrības vienošanas metodēm to vēl vairāk vājināja. Tomēr jāatzīst, kaut G.~Narutoviča nogalināšana iezīmēja demokrātijas gara vājināšanos, tā izpausmes sabiedriskajā dzīvē turpinājās. Opozīcija vardarbīgām metodēm politikā pastāvēja līdz pat valsts sakāvei 1939.~gadā.

Lielo satraukumu valstī pēc G.~Narutoviča noslepkavošanas mēģināja izmantot arī J.~Pilsudska piekritēji lai atkal panāktu varas pāreju viņa rokās, taču prezidenta pienākumus izpildošais Seima priekšsēdētājs M.~Ratajs jau prezidenta nogalināšanas dienā izveidot valdību uzdeva ģenerālim V.~Sikorskim. Viņa izveidotā Vispārnacionālā (varavīksnes) valdība ieveda Varšavā kara stāvokli, spēja atjaunotajā valstī radīt nosacītu stabilitāti. Reizē V.~Sikorskis realizēja zināmu kompromisu, par kara ministru nozīmējot K.~Sosnkovski, kurš savukārt par Ģenerālštāba priekšnieku apstiprināja J.~Pilsudski.

Par jauno prezidentu ievēlēja agrāko sociālistu, bijušo J.~Pilsudska līdzgaitnieku, pirms kara atgājušu no politikas un darbojošos kooperatīvajā kustībā, labējo ļudovieti no partijas ``\pltxti{Piast}'' S.~Vojcehovski, kurš kļuva par kompromisa kandidatūru un kuru atbalstīja J.~Pilsudskis. Par S.~Vojcehovski nobalsoja tās pašas frakcijas parlamentā, kuras balsoja arī par G.~Narutoviču, ieskaitot arī nacionālo mazākumu pārstāvjus. Šoreiz labējie parlamenta izvēli uztvēra mierīgi. Politiskā krīze it kā beidzās bez īpašām sekām, ja neskaita to, ka nacionālajiem mazākumiem bija nepārprotami parādīts, ka viņi ir otrās šķiras pilsoņi, kuriem nav tiesību vienlīdz ar poļiem izlemt valsts attīstības svarīgākos jautājumus. Pēc 1922.~gada vēlēšanām vairākums poļu partiju kā nerakstītu likumu pieņēma prasību izslēgt nacionālo mazākumu deputātus no parlamentārajām kombinācijām (valdību sastādīšanā, kompromisiem likumprojektu veidošanā u.c.) Ar to nacionālo mazākumu deputāti bija nolemti darboties tikai opozīcijā. Tas sašaurināja demokrātijas un parlamentārisma bāzi.

Notikušo J.~Pilsudskis uztvēra kā signālu sev pavisam atteikties no valstiskiem pienakumiem. J.~Pilsudskis, kurš pēc valsts galvas pilnvaru nolikšanas vēl palika Polijas armijas Ģenerālštāba priekšnieka un Aizsardzības padomes priekšsēdētāja amatos, 1923.~gada 30.~maijā vīlies atkāpās no pirmā, bet 2.~jūlijā, pēc tam, kad par kara ministru bija kļuvis S.~Šeptickis, ar kuru bijušais valsts priekšnieks atradās asā konfliktā, arī no otrā amata. 3.~jūlijā viņam par godu sarīkotā atvadu banketā maršals tieši apvainoja nacionāldemokrātus kā līdzvainīgus G.~Narutoviča nogalināšanā un paziņoja, ka viņš kā karavīrs nevar aizstāvēt šos kungus, tāpēc arī nolēmis aiziet no armijas un valsts dienesta. J.~Pilsudskis paziņoja, ka ir visu atdevis Polijai un tagad vēlas baudīt ģimenes dzīvi. Abas atkāpšanās prezidents S.~Vojcehovskis pieņēma. Faktiski šajā situācijā J.~Pilsudskis pielietoja jaunu cīņas taktiku~--- bez pretošanās uz laiku atstāja politiskās cīņas lauku. Pastāv divi galvenie viedokļi par bijušā valsts priekšnieka rīcības motīviem. Lūk, tie: J.~Pilsudska atturība politikā parlamentārisma posmā izskaidrojama: 1) kā spilgta demokrātisma izpausme; 2) izplānota taktika, lai atbrīvotos no visām partijiskām un ideoloģiskām saistībām, nogaidot, līdz viņa pretinieki netiks galā ar valstī pieaugošajām pretrunām un konfliktiem, un pēc tam kļūtu ar šķīrējtiesnesi politisko partiju strīdos. Autors uzskata, ka tālākie notikumi pierāda otrā skaidrojuma pareizību.

Maršals apmetās viņam uzdāvinātajā villā Varšavas Sulejuvkas (\pltxti{Sulejówek}) piepilsētā un it kā noslēdzās no politiskās dzīves. 1923.~gadā J.~Pilsudskis demonstratīvi atteicās no maršala pensijas, nododot to labdarības mērķiem, uzturēja ģimeni no ienākumiem par savu literāro darbību un publiskām lekcijām. Tie nebija augstāki par armijas kapteiņa algu. Publikas priekšā viņš parasti parādījās vecā poļu strēlnieka frencī, vilcienā brauca otrās klases vagonā. Tas deva vēlamo efektu, cēla viņa popularitāti. Taču reizē viņš vēroja politiskās norises, viņā auga kritisks noskaņojums par valdošo eliti, pie varas esošajiem politiskajiem spēkiem. 1923.~gadā kādam no saviem piekritējiem J.~Pilsudskis ``izkratīja sirdi'': ``Es, protams, zinu, ka ar šo sabiedrisko domu, ar šo ``inteliģences'' viedokli iet dažādi \dots{} tas ir pārāk jūtīgs pret varas iespaidu un noraida tos, kam nav varas. Es to novēroju arī attiecībā pret sevi. Manu ``lielo apbrīnotāju'' skaits no šī sabiedrības slāņa kļūst arvien mazāks. Šo gļēvo atkritēju skaits kļūs jo lielāks, jo ilgāk es stāvēšu malā. Tas ir tik cilvēcīgi un \dots{} ļoti poliski.''

J.~Pilsudskis un viņa līdzgaitnieki neradīja atsevišķu politisku partiju, paklausīgu savam vadonim. Taču viņš saglabāja ietekmi uz vairāku partiju darbiniekiem un daudziem ģenerāļiem un virsniekiem. Daži no viņa piekritējiem darbojās partijās, kas tiem deva zināmu, bet ne pārāk lielu iespēju ietekmēt Seimā izskatāmo lietu gaitu. Tomēr visas valdības bija spiestas rēķināties ar to, ka turpat blakus dzīvo un notikumus vēro cilvēks, kuru liela daļa Polijas pilsoņu uzskatīja par valsts dibinātāju, nacionālu varoni un cienīgu būt valsts vadībā. Taču visus valdību piedāvājumus sadarboties, ieņemt amatus armijā un valsts aparātā viņš noraidīja, izvirzot arvien jaunas valdībām un prezidentam nepieņemamas prasības. Vienlaikus viņa piekritēji izvērsa skaļu propagandas kampaņu, ka armija bez sava maršala nespēj sekmīgi attīstīties.

Bloks, kurš ievēlēja divus pirmos Polijas prezidentus, nebija stabils. Vairākums Seimā bija ar pagaidu raksturu, jo tajā ietilpa nacionālo mazākumu pārstāvji, drīz sākās jaunas koalīcijas meklējumi, kuru veidotu tikai poļu partijas. Spēku samērs Seimā noteica to, ka valdība varēja būt tikai kādas partiju koalīcijas radīta, vai arī ārpusparlamentāra~--- profesionāļu jeb ierēdņu veidota. Līdz 1926.~gada maijam nomainījās piecas tādas valdības. To straujās maiņas cēloņu vidū bija arī lielās domstarpības starp dažādām partijām par tādām problēmām kā: agrārā jautājuma risināšanas ceļi un tempi, krīzes parādību rūpniecībā pārvarēšanas metodes, attieksme pret mazākumtautībām, 1921.~gada konstitūcijas caurskatīšanu utt.

Ļudovieši no partijas \pltxti{Piast} 1923.~gada 17.~maijā noslēdza līgumu ar labējo bloku~--- Nacionālās vienotības kristīgo savienību. Šis tā saucamais ``Lanckaronas pakts'' (\pltxti{Pakt lanckoroński}~---  Pakta nosaukums cēlies no tā, ka tā noteikumi tika apspriesti kāda senatora īpašumā, kurš pārstāvēja Lanckaronas vēlēšanu apgabalu) paredzēja, ka:

\begin{enumerate}
\item Seima vairākumam jābūt poliskam, valdībai jāsastāv tikai no poļiem,

\item visās valsts iestādēs jālieto tikai poļu valoda,

\item visās profesionālajās, vidējās un augstākajās mācību iestādēs poļiem jāieņem priviliģēts stāvoklis, nacionālajiem mazākumiem tiek ieviesta procentu norma,

\item visās vietējās pašvaldībās tiek garantēts poļu ietekmes pārsvars,

\item agrārā reforma jāveic 10~gadu laikā, ik gadu 200~000 hektāru valsts un privāto zemju sadalot parcelēs (Valsts centrālajās vojevodistēs paliekošās saimniecības nedrīkstētu būt lielākas par 180~hektāriem un nomalēs~--- par 400~hektāriem. No reformas tiktu izslēgtas arī līdz 1~000~ha lielas muižas, ja tajās bija lauksaimnieciskās produkcijas pārstrādes uzņēmumi. 1925.~gadā šāds likums ar mazām izmaiņām arī tika pieņemts.) Tātad tika paredzēta nacionālo mazākumu tiesību sašaurināšana. Arī agrārā reforma tika ierobežota, jo faktiski tika dalīts ļoti maz zemes un īpaši šķēršļi tās sadalei likti nomalēs. Partija \pltxti{Piast} par tiesībām piedalīties valdības veidošanā samaksāja ar atteikšanos no radikālas agrārās reformas īstenošanas. Kā atzīmējis P.~Olšanskis, kaut agrārā reforma bija ar ļoti ierobežotu raksturu, tomēr muižnieku pamatmasa no šī laika atbalstīja J.~Pilsudski un viņa vadītos politiķus.
\end{enumerate}

Līguma īstenošanas rezultātā bija spiesta demisionēt V.~Sikorska valdība un 1923.~gada 28.~maijā tās vietā nāca vispirms V.~Vitosa valdība. Tā gan ilgi nenoturējās, daļa no partijas \pltxti{Piast} deputātiem uzstājās pret agrārās reformas ierobežošanu un izstājās no valdības, kura ar to zaudēja parlamenta vairākumu. Pēc tam no 1923.~gada decembra līdz 1925.~gada novembrim pastāvēja centriskā V.~Grabska valdība. Tai nācās nodarboties ar inflācijas apkarošanas, agrārās reformas, (1923.--1925.~gadā valsts pārdzīvoja ievērojamas ekonomiskas grūtības, par kurām sīkāk tiks runāts apakšnodaļā par tautsaimniecība attīstību) savas~--- Polijas ostas celtniecības, nacionālajām attiecībām. Šai valdībai bija arī jānoteic sava attieksme pret Vatikānu.

Bija nepieciešams noslēgt līgumu ar Vatikānu par agrāk pastāvējušo baznīcas diocēžu robežu izmaiņām, jo bija taču radušās jaunas valstis ar savām robežām un tagad iznāca, ka bīskapu un arhibīskapu draudžu vadības tiesības sniedzās pāri valstu robežām. 1925.~gada 10.~februārī \strong{Konkordāts} (speciāls līgums) \strong{ar Vatikānu} baznīcas diocēžu robežas pielīdzināja jaunajām valsts robežām. Bez bijušajām Gnēzenes (\pltxti{Gniezno})~--- Poznaņas (\pltxti{Poznań}), Ļvovas (\pltxti{Lwów}) un Varšavas (\pltxti{Warszawa}) arhibīskapijām (metropolijām) tika radītas vēl divas: Viļņas (\pltxti{Wilno}) un Krakovas (\pltxti{Kraków}). Valsts piekrita Vatikāna tiesībām nozīmēt bīskapus, neprasot varas iestāžu viedokli, uzņēmās sabiedriskajās skolās ieviest obligātu reliģisko apmācību. Katoļu baznīcas īpašumu nedrīkstēja atsavināt bez attiecīgo baznīcas instanču piekrišanas. Visa Polijas teritorija dalījās ap 6,5 tūkstošos katoļu draudžu, ap 63,5\% visu iedzīvotāju bija katoļi.

Jāpiebilst, ka no 1922.~gada valstī pastāvēja arī autokefālā (no grieķu \eltxti{αυτός}~— pats un \eltxti{κεφαλή}~— galva, patstāvīga pareizticībā baznīca, kanoniski un administratīvi pilnīgi neatkarīga no citām baznīcām) pareizticīgā baznīca, kuru 1925.~gadā oficiāli atzina arī Polijas valdība. Tās 2~000 garīdznieku un 1~500 draudžu ietekmē atradās galvenokārt ap 3~miljonu pareizticīgo valsts austrumu vojevodistēs.

J.~Pilsudskis sākotnēji ieņēma it kā nogaidošu pozīciju pret V.~Grabska valdību pat tad, kad par kara ministru kļuva viņa senais konkurents ģenerālis V.~Sikorskis, taču jau 1925.~gadā maršals izvērsa V.~Sikorska kritiku, panāca, ka par kara ministru kļuva ģenerālis L.~Žeļigovskis. Pret V.~Grabska valdību bija noskaņota armijas virsniecība, jo, pieņemot taupības programmu, tā uz nenoteiktu laiku atlika paredzēto militāro amatpersonu atalgojuma reformu. Ar to J.~Pilsudska piekritēju skaits armijā auga.

20.~gadu vidū Polijā saasinājās politiskā, ideoloģiskā un administratīvā cīņa pret Komunistisko partiju, kura savā II kongresā 1923.~gadā beidzot bija atteikusies no nihilistiskās pozīcijas nacionālajā jautājumā, kas novēloti, bet tomēr ļāva tai cīnīties par lielāku ietekmi masās.

1925.~gadā pēc Iekšlietu ministrijas datiem bija 3,5~tūkstoši politieslodzīto, apvainotu par piederību Komunistiskajai partijai. Lielākā politieslodzīto daļa, kā to atzina Seima komisija, daudzus mēnešus bija spiesta pavadīt iepriekšējā ieslodzījumā, gaidot izmeklēšanas rezultātus. Polijas cietumi (to bija 341) bija pārpildīti. 1925.~gadā Polijas valdība ieveda lauku kara tiesas. Komunistu organizācijās plaši tika iesūtīti policijas aģenti~--- provokatori. Tāpat kā carisma apstākļos partija pret tiem cīnījās ne tikai ar atmaskošanu, bet arī ar ieročiem. 1925.~gada 6.~martā divus jaunus komunistus: F.~Pilarčiku un M.~Gaičiku, ko apvainoja šāda provokatora, kura dēļ bija apcietināti desmiti revolucionāru, slepkavībā, aplenca policija. Abi nepadevās, septiņas stundas atšaudījās. Tikai ar smacējošas gāzes palīdzību izdevās pārvarēt šo divu revolucionāru pretošanos.

17.~jūlijā policija Varšavā mēģināja arestēt komunistus V.~Hibneru, H.~Rutkovski un V.~Kņevski, kuriem tāpat bija dots uzdevums nogalināt provokatoru. Provokators bēga uz Ļvovu, taču arī tur viņu atrada un nogalināja komunists N.~Botvins. Visi trīs izrādīja bruņotu pretošanos, bija upuri policistu un civiliedzīvotāju vidū. Tā kā bija cietuši arī civilisti, lauku kara tiesa visiem trijiem piesprieda pakāršanu, nevis nošaušanu kā parasti to darīja ar politiskajiem pretiniekiem. Spriedums tika izpildīts, kaut par arestētajiem iestājās ne tikai PSRS, bet arī Rietumu kultūras darbinieki~--- komunisti A.~Barbiss, Ž.~R.~Bloks, L.~Aragons. (Polijas Tautas republikas pastāvēšanas laikā starpkaru periodā bojā gājušie komunisti tika uzskatīti par varoņiem, viņu vārdos nosauktas ielas. 1950.~gada 21.augustā V.~Hibneru un viņa biedru bojā ejas 25~gadu atceres dienā pie Varšavas Citadeles sienas viņiem tika atklāts piemineklis.)

Tai pat laikā jāpasaka, ka komunisti meklēja sabiedrotos arī citās politiskajās organizācijās. 20.~gados viņiem izdevās iegūt zināmu ietekmi dažās kreisi noskaņotās ļudoviešu grupās, jo, kā jau teikts, ļudoviešu kustībā starpkaru periodā pastāvēja dažādi strāvojumi: no konservatīva līdz konsekventi demokrātiskam. Tas vēlāk~--- jau pēc Otrā pasaules kara~--- ļāva komunistiem atrast sabiedrotos arī daļā zemniecības.

Sakarā ar krīzi un augošo lielkapitāla neuzticību valdībai V.~Grabskis 1925.~gada 13.~novembrī iesniedza atlūgumu. Šai sakarā 18.~novembrī Rīgā iznākošā avīze ``\rutxti{Сегодня}'' (``Šodiena'') par stāvokli Polijā rakstīja: ``\dots{}kritisks stāvoklis šeit jau gandrīz kļuvis par normu. Un tāpēc asas krīzes šeit pārdzīvo kā krīzes kvadrātā. \citespace{} Publikā tūlīt izplatījās baumas par kaut kādu apvērsumu. Pagaidām nekāda ``apvērsuma'' nav, bet, kā saka, tas jūtams gaisā.'' Pavisam drīz izrādījās, ka avīzes priekšnojautas bija pareizas.

Jāsaka, ka V.~Grabska valdība bija visilgāk pastāvošā līdz 1926.~gada maijam, tās sasniegumi bija arī visbūtiskākie. (Blakus zlota stabilizācijai kā zināms sasniegums jāmin arī agrārais likums, kura galīgais variants tika pieņemts jau pēc V.~Grabska demisijas.) Viņa vietā valdības vadību uzņēmās bezpartijisks premjerministrs A.~Skrinskis. Sākotnēji likās, ka tāda valdība apmierina J.~Pilsudski, jo tajā ietilpa arī daži viņam tuvi cilvēki. Taču pakāpeniski pieauga \strong{J.~Pilsudska domstarpības ar pie varas esošajiem politiķiem}.

Trīs gadus, pavadītus Sulejuvkas villā, bijušais valsts priekšnieks galvenokārt izmantoja mīta par sevi nostiprināšanai. Viņa raksti, 1924.~gadā izdotā grāmata ``1920.~gads'', dažādas uzstāšanās un intervijas bija pilnas ar pašslavinājumiem. Piemēram, 1926.~gada 10.~februārī intervijā avīzei ``\pltxti{Kurier Poranny}'' (``Rīta Kurjers'') viņš paziņoja: ``Es esmu cilvēks, kurš ar savu darbu piespiež sevi cienīt, neatkarīgi no to vai citu apsvērumu noteiktas attieksmes pret mani. Es biju Valsts priekšnieks un pratu 1918.~gadā nodrošināt kārtību un valsts attīstību, neraugoties uz vispārēju klīšanu notikumu virpulī, kuru priekšā bailīgumā pārmežoniskojusies un nebrīvē pagrimusī sabiedrība izskatījās tā, ka sāka trīcēt tikai no domas par Polijas kā neatkarīgas valsts pastāvēšanu. Es pratu apvīt mūsu karogus ar tik dižu uzvaru slavu, kādas nepazina pat mūsu vecvectēvi, un panācu tās tad, kad gļēvulība mūsu balto ērgli padarīja par dzeltenu no bailēm \citespace{} Beidzot, pretēji poļu tradīcijai, es pratu, aizejot no valsts darbības, neko nepieprasīt no mūsu valsts, bet nodrošināt dzīvi sev un savai ģimenei ar paša darbu. Un labi saprotu, ka tādu cilvēku ir grūti neievērot, kamēr vien viņš elpo un dzīvo\dots{}''

Protams, J.~Pilsudska propagandiskais uzbrukums nevarēja palikt bez atskaņām. Piemēram, ģenerāļa S.~Šepticka 1925.~gadā izdotā grāmatā bija parādīts, ka tās pavēles poļu-padomju kara laikā, kuras J.~Pilsudskis pēc kara sevišķi kritizēja, tika izdotas pēc tiešiem viņa paša norādījumiem, kaut S.~Šeptickis pret tām iebilda. J.~Pilsudska radāmo mītu pamatoti apstrīdēja izcilais ģenerālis J.~Pilsudska bijušais līdzgaitnieks V.~Sikorskis grāmatā ``\pltxti{Józef Piłsudski jako wódz i dziejopis}'' (Juzefs Piklsudskis kā vadonis un vēsturnieks), kuru viņš izdeva ar pseidonīmu Karols Pomorskis (\pltxti{Karol Pomorski}). Par J.~Pilsudski V.~Sikorskis rakstīja: ``Dziedot sev kā virspavēlniekam, valstsvīram un cilvēkam slavas himnas, viņš apvainojumu dubļiem un nicinājuma vārdiem apvelta visus un visu, kas nav saistīti ar viņa svītu: ģenerāļus, armiju, tautu, bet pēdējā laikā nav žēlojis pat nacionālo simbolu~--- Balto Ērgli. \citespace{} Mīts, kuru pats sev vēlas radīt Maršals Pilsudskis, ir nolemts sakāvei pirmajā sadursmē ar vēsturisko patiesību, iegūtu kaut vai no viņa paša radīta avota kāds ir ``1920.~gads''.'' Savu tēzi V.~Sikorskis ilustrēja ar desmitiem piemēru, vispusēji kritizējot J.~Pilsudska secinājumus, konfrontējot tos ar dokumentiem. Arī V.~Sikorska grāmata nesastāvēja no ``tīras'' patiesības, tās mērķis bija nevis zinātnisks, bet propagandisks~--- radīt ``melnu'' mītu par konkurentu. Taču tajā nebija primitīvas demagoģijas, kas pildīja nacionāldemokrātu preses izdevumu lappuses.

Tolaik netika veiktas iedzīvotāju aptaujas, netika veidoti politiķu reitingi, taču domājams, ka poļu vairākumu stiprāk ietekmēja J.~Pilsudska pozitīvais devums valsts neatkarības izcīnīšanas periodā, sarūgtināja nebūšanas un neveiksmes parlamentārajā periodā, par kurām viņu nevarēja vainot, un tāpēc viņa popularitāte palika augsta. Uz to balstoties viņš varēja gatavot varas pārņemšanu savās rokās. Poļu vēsturnieki atzīst pazīmju esamību, ka J.~Pilsudskis 20.~gadu pirmajā pusē vairākus gadus gatavoja valsts apvērsumu. Redzami soļi šai virzienā bija sperti 1925.~gadā. Pēc apvērsuma gan tika apgalvots, ka maršala un viņam uzticīgo karaspēka daļu uzstāšanās bijusi spontāna reakcija uz negatīvajām parādībām, kuras arvien plašāk izplatījās Žečpospolitā. Protams, atzīties, ka apvērsums tika gatavots jau ilgstoši, neviens nevēlējās. J.~Pilsudska bijušie padotie vēl no strēlnieku savienību un leģionu laikiem pakļautību izprata kā uzticību saviem komandieriem. Viņi uzticējās tikai ``savējiem'', un uz virsniekiem, kuri bija dienējuši citās Austrungārijas karaspēka daļās vai pat Krievijas armijā, skatījās ``no augšas''. Šādu uzticību leģionāri ienesa arī Polijas armijā. Maršala piekritēji armijā pastāvīgi uzturēja savstarpējus sakarus, gatavoja augsni varas saņemšanai arī kreisajās partijās, kurām bija seni sakari ar bijušo valsts priekšnieku.

Arī nacionāldemokrāti nebija apmierināti ar politisko situāciju valstī. R.~Dmovskis un labējo nometne gatavoja t.s. ``Dižās Polijas nometnes'' (\pltxti{Obόz Wielkiej Polski}),~--- endeku vadītas organizācijas, kura atteicās no V.~Grabska aizstāvētā parlamentārisma, izveidi, lai ar ārpusparlamentāru spiedienu ieviestu autoritāru valdību. R.~Dmovskis 1925.~gadā Parīzē zondēja, kā Francija izturēsies labējo īstenota apvērsuma gadījumā. Tāpat arī ļudovietis no partijas \pltxti{Piast} V.~Vitoss un sociālists I.~~Dašinskis kritizēja hipertrofēto parlamentārismu un izpildvaras novājināšanos.

Šajā situācijā 1925.~gada beigās J.~Pilsudskis jau skaidri saskatīja rīcības plānu nākotnei. 15.~decembrī viņa līdzgaitnieks K.~Svitaļskis dienasgrāmatā ierakstīja: ``Komandanta plāns ir sekojošs: valdības krīzi censties pārvarēt bez Seima. Iekļūt armijā. Uzstāties, visticamāk, kara ministra lomā, asi un brutāli pret Seimu. Seimu neatlaist, bet retāk to sasaukt. Sēdot kabinetā, vērtēt tā locekļus, orientācijai, ar kuru var iet, ar kuru nē. Pie varas var nonākt 1926.~gada rudenī. Tad var rīkot vēlēšanu kampaņu.'' Pastāv arī citas liecības, ka 1926.~gada sākumā J.~Pilsudskis sarunās ar sev tuviem cilvēkiem sevi izpauda kā nākamais Polijas saimnieks. Tādejādi, jau no 1925.~gada beigām J.~Pilsudskis bija ieinteresēts izsaukt jaunu valdības krīzi un pat noteica tās termiņu. PPS politiķis I.~Dašinskis 1926.~gada janvārī brīdināja tautu: ``Morālais vadonis var sagraut Poliju''.

Apvērsuma nobriešanu veicināja smagais ekonomiskais stāvoklis (1926.~gada aprīlī bezdarbnieku skaits sasniedza 400~000), ko valdība nespēja pārvarēt. Arvien biežāk notika bezdarbnieku sadursmes ar policiju. Tā, 1926.~gada 9.~februārī Kališā bezdarbnieku demonstrācijas sadursmē ar policiju un karaspēku tika nogalināti 9 un vairāk nekā 60~strādnieki ievainoti. Atkal pieņēmās spēkā inflācija. Sarunas par ārzemju kredītu piešķiršanu vilkās garumā. Valstī auga neapmierinātība ar korupciju.

Sabiedrības sadalīšanās antagonistiskās nometnēs neatkarīgās Polijas pastāvēšanas pirmajos piecos gados ne tikai netika pārvarēta, bet arī ievērojami padziļinājās un ieguva stabilu raksturu.

Kopumā demokrātija izrādījās neērta poļu valdošajiem slāņiem. Poļu pilsonība bija neapmierināta ar nenoteikto spēku samēru politiskajā arēnā, kur tautas zemāko slāņu intereses pārstāvošo partiju klātbūtne parlamentā draudēja ar radikālu sociāli ekonomisko programmu aktualizēšanos, bet divu galveno pilsonības grupējumu (endeku un pilsudčiku) savstarpējā politiskā cīņa ilgstoši nedeva vienu noteiktu uzvarētāju. Pēc krievu vēsturnieces I.~Mihutinas vārdiem arī pilsētas sīkburžuāzija demokrātiskā parlamentārismā redzēja sev negatīvas puses. Strādnieku šķiras iekarojumos (tiesības uz streikiem, apdrošināšanas tiesības u.c.) tā saskatīja savu materiālo nedienu cēloni. Arī zemniecības, kuru pilsētu sīkpilsonība agrāk bija radusi redzēt uz zemākā sabiedriskās attīstības pakāpiena, iespaida pieaugums demokrātisku vēlēšanu rezultātā to neapmierināja. Ar laiku arvien skaidrāka kļuva redzama tendence valsts varas sistēmu vairāk pielāgot mantīgo augšslāņu vajadzībām.

Var arī piebilst, ka pēc 1922.~gadā notikušā itāļu fašistu apvērsuma, 1923.~gadā tāpat apvērsuma rezultātā Spānijā varu bija saņēmis ģenerālis M.~Primo de Rivera, arī Bulgārijā militāristi gāza likumīgo valdību un nodibināja labēju A.~Cankova valdību. Šādi precedenti arī Polijā stimulēja ekstrēmistiskās aprindas. Pāri Polijai vēlās neatklātu atentātu un sprādzienu vilnis. Valdība tajos apvainoja galvenokārt komunistus. Veltīgos vainīgo meklējumos policija arestēja desmitiem nevainīgu cilvēku.

Endeki vēlējās atveseļot ekonomiku, veicinot privātkapitāla darbību un pazeminot darba algu. Viņu finanšu ministrs J.~Zdehovskis budžeta deficīta pārvarēšanai paredzēja par 10\% paaugstināt uz laiku visus nodokļus (izņemot nodokli uz īpašumu), muitas nodevas; ieviest nodokļus graudu malšanai un uz elektrogaismekļiem; samazināt algas valsts kalpotājiem, pensijas un pabalstus pensionāriem un invalīdiem; atlaist no armijas daļu virsnieku un samazināt par 18~tūkstošiem dzelzceļnieku skaitu.

Sociālisti (PPS) turpretī piedāvāja segt izdevumus uz materiāli nodrošināto slāņu rēķina. Viņi bija spiesti ievērot komunistu propagandu, ka Polijas buržuāzija nedomā par tautu. Sociālisti piedāvāja palielināt nodokļus uz īpašumu; samazināt izdevumus armijai, saīsinot dienesta laiku tajā divas reizes; paplašināt sabiedriskos darbus bezdarbniekiem. Kad šo plānu partneri valdošajā koalīcijā atteicās realizēt, PPS atstāja valdību, aprīlī valdības koalīcija izjuka, kas 5.~maijā izsauca A.~Skšinska valdības demisiju. Izraisījās kārtējā valdības krīze. Tā notikumi attīstījās pilnā saskaņā ar J.~Pilsudska plānu.

Īpatnēju nostāju pauda kreisie ļudovieši. Viņi noraidīja jebkuras valdības pastāvēšanas likumību, kura savā darbībā neorientētos uz plašu lauku un pilsētu iedzīvotāju slāņu interešu apmierināšanu. Tāpēc viņi uzstājās pret labēji-centriskām valdībām, pat ja tās bija izveidotas atbilstoši demokrātiskajai parlamentārajai procedūrai. Nebūdami revolucionāri, viņi tomēr pieļāva iespēju gāzt tādu valdību vardarbīgā ceļā. To var vērtēt kā vienu no cēloņiem, kāpēc Poļu Zemnieku partijas~--- \pltxti{Wyzwolenje} un Zemnieku partija (\pltxti{Stronnictwo Chłopskie}) atbalstīja J.~Pilsudska apvērsumu. Šo partiju līderi cerēja, ka pēc V.~Vitosa valdības gāšanas viņi spēs īstenot savus tautas Polijas izbūves plānus. Tikai partija \pltxti{Piast} palika tīri parlamentāro metožu piekritēja sociāli-ekonomisko un politisko jautājumu risināšanā. Eksistējošo Polijā valsts iekārtu arī 1926.~gadā tā uzskatīja par pilnībā atbilstošu zemniecības interesēm. \pltxti{Piast} piedalījās visu pēc V.~Grabska valdības atkāpšanās 1925.~gada novembrī pastāvējušo parlamentāro valdību sastāvā. Šajā brīdī \pltxti{Piast} pēc saviem uzstādījumiem stāvēja tuvāk buržuāziskajām nekā kreisajām ļudoviešu partijām. Tā apvērsuma laikā ļudoviešu partijas izrādījās dažādās barikāžu pusēs.

Ar pilsudčikiem saistītā prese ikdienas rakstīja, ka labējie gatavo ``fašistiskas diktatūras'' nodibināšanu, skandēja par deputātu un valsts ierēdņu korupciju, valsts kases iztukšošanu, nekompetenci. Pastāv diezgan daudz liecību par to, ka J.~Pilsudskis pēdējos mēnešos pirms valsts apvērsuma veda sarunas ar ietekmīgiem rūpnieku un muižnieku pārstāvjiem, solot, nonākot pie varas, panākt likuma par agrārreformas pārskatīšanu, ierobežot sociālo likumdošanu u.c. Šie viņa soļi liecināja par maršala vēlmi piesaistīt sev tās iedzīvotāju grupas, kuras tradicionāli atbalstīja nacionāldemokrātus. Tai pat laikā viņš nepārtrauca attiecības ar saviem piekritējiem kreisajās partijās, arī ar kreisajiem arodbiedrību darbiniekiem, kuri varēja ar dzelzceļnieku streiku paralizēt visu dzīvi valstī, tai skaitā arī militāros pārvadājumus.

Otrs apvērsumu veicinošs faktors bija nebeidzamie strīdi par armijas vadību gan miera, gan, it īpaši, kara gadījumā, kuros pastāvīgi iejaucās J.~Pilsudskis. Atklātie politiķu mēģinājumi viņu~--- cilvēku ar augstāko militāro rangu, turēt tālāk no armijas tiešās vadības, izsauca biežus un atklātus maršala dusmu izvirdumus. Endeki turpretī kategoriski iebilda pret J.~Pilsudska pielaišanu armijas vadībai viņa neprofesionalitātes dēļ. Viņi uzskatīja, ka tā būtu valsts militāra katastrofa, līdzīga tai, kura norisa 1920.~gadā J.~Pilsudska uzbrukuma Kijevai rezultātā. 1926.~gada sākumā J.~Pilsudskis izvērsa arvien asāku valdības un partiju varas (\pltxti{partyjnictwo}) kritiku. 5.~maijā J.~Pilsudskis Francijas sūtnim Polijā J.~Larošam paziņoja, ka uzskata par nepieciešamu radīt tādu valdību, ``kura nostāsies pāri partijām un pirmkārt pāri Seimam'', ka ``parlamentārisms gandrīz visur pārdzīvo krīzi, kuru izārstēt var autoritāra vara''.

Labie un centristi gribēja atgriezties pie koalīcijas, kura 1923.~gadā radīja V.~Vitosa valdību. J.~Pilsudskis šo darbību uzskatīja par pret viņu vērstu provokāciju. (J.~Pilsudskis pret V.~Vitosu izjuta krasu antipātiju, kritizēja, un, teiksim tieši, rupji lamāja viņu kā nevienu citu politiķi. J.~Pilsudskis V.~Vitosu sauca par ``poļu zemnieku ļauno garu'', ``lielāko Polijas dumiķi'' un ``viskorumpētāko subjektu, kādu es esmu sastapis''. Acīmredzot maršals sajuta zemnieku vadonī cienīgu sāncensi.) Pats V.~Vitoss par atmosfēru pirms jaunas savas valdības radīšanas rakstīja; ``Skšinska valdības demisijas izraisītā krīze turpinājās ļoti ilgi, to pavadīja haoss, jezga, kas pastiprinājās ar katru dienu \citespace{} Kuluāros valdīja drūma un vienlaikus dīvaina noskaņa. Žurnālisti zobojās par Seimu, par deputātiem un pat par prezidentu \citespace{} Prezidents Vojcehovskis staigāja sadrūmis, darot zināmu nodomu nolikt savas pilnvaras''.

Šādā situācijā J.~Pilsudskis plānoja militāru demonstrāciju, kurai acīmredzot bija jāliek prezidentam nodot varu viņam. Vēl vecās valdības kara ministrs L.~Žeļigovskis 18.~aprīlī izdeva pavēli par speciāli atlasītu karaspēka daļu, kuras komandēja J.~Pilsudskim uzticami militāristi, manevriem maršala vadībā. Manevriem bija jāsākas 10.~maijā. Taču šajā dienā tika izveidota jauna V.~Vitosa valdība. Tas izsauca visu kreiso un arī J.~Pilsudska piekritēju sašutumu.

Kreisie apgalvoja, ka valdībai Seimā gan ir vairākums, taču tautā tai nav atbalsta. Parlamentārās iekārtas trūkumi, biežā valdību maiņa, (Kopš valsts izveides līdz 1926.~gadam Polijā bija darbojušās 15~valdības ar 11~dažādiem Ministru prezidentiem. 1925.~gadā Polijā pastāvēja 95~partijas, 1926.~gadā no tām 21 bija pārstāvēta Seimā.), korupcija izsauca tautas vilšanos par ``seimakrātiju''. Neviens faktiski neaizstāvēja 1921.~gada marta Konstitūciju, taču pieņemt jaunu konstitūciju arī nevarēja, jo domstarpības par tās principiem bija pārāk lielas.

V.~Vitosa valdībā kara ministra vietu ieņēma J.~Pilsudska atklāts pretinieks ģenerālis J.~Malčevskis, un maršala piekritēju pozīcijas armijā bija apdraudētas. Uztraucošs maršalam un viņa piekritējiem bija dienu pirms tam izskanējušais V.~Vitosa paziņojums: ``Lai beidzot Maršals Pilsudskis iznāk no aizsega, lai izveido valdību, lai uzņem savā koalīcijā visus radošos spēkus, kuri ar visu dvēseli jūt līdzi valsts interesēm. Ja tas netiks darīts, radīsies iespaids, ka viņš nav ieinteresēts ieviest valstī kārtību.'' Šāds aicinājums paredzēja valdības izveidi likumīgā, parlamentārā ceļā, kas, kā jau atzīmēts, J.~Pilsudskim, kurš nevēlējās būt parlamenta atkarībā, nebija pieņemami. Tagad maršals faktiski nevarēja atkāpties, nezaudējot reputāciju.

1926.~gada 11.~maijā intervijā avīzei ``\pltxti{Kurier Poranny''} J.~Pilsudskis publicēja skarbu uzbrukumu Seima politiķiem par viņu patmīlību un neizdarībām. Acīmredzot pēc V.~Vitosa norādījuma cenzūra avīzes numuru konfiscēja.

Faktiski J.~Pilsudskis vēlējās panākt Valsts prezidenta un politisko partiju līderu piekrišanu esošās politiskās sistēmas korekcijai, atzīstot viņam īpašu, Konstitūcijā neparedzētu valsts drošības garanta un vienīgā bruņoto spēku interešu izteicēja statusu. Zīmīgi, ka valdība un prezidents bija pārliecināti, ka J.~Pilsudskis arī šoreiz atturēsies no izšķirošas rīcības, aprobežojoties tikai ar skaļiem paziņojumiem presē. Tāpēc valdība 11.~maijā neveica nevienu soli, kas varētu kavēt apvērsumu.

12.~maija rītā maršals devās uz Belvedēru, lai pieprasītu prezidentam atlaist V.~Vitosa valdību. Taču prezidents S.~Vojcehovskis, neko nenojaušot, neilgi pirms tam bija devies uz savu ārpilsētas rezidenci.

Radās dīvaina situācija: ilgi lolotais plāns varēja izjukt tikai tāpēc, ka nebija kam izklāstīt savas prasības. Tad nu J.~Pilsudskis nolēma rīkoties noteikti, uzņēmās sev uzticamo karaspēka daļu komandēšanu un virzīja tās (15 pulkus) uz Varšavu. Pirmkārt vajadzēja ieņemt tiltus pār Vislu, lai pilsētas kreisajā krastā iegūtu kontroli pār valdības ēkām.

Iespējams, ka arī tagad bija plānota tikai spēka demonstrācija, lai piespiestu prezidentu S.~Vojcehovski atcelt valdības vadītāju V.~Vitosu. Tas sagrautu labējo koalīciju un pārvērstu J.~Pilsudski par arbitru politiskajos strīdos. No partijām PPS un \pltxti{Wyzwylenie} atbalstīja J.~Pilsudski. Dienas otrajā pusē maršalam uzticamās vienības Varšavā ieņēma Vislas tiltus.

Uz Poņatovska tilta Varšavā notika dramatiska prezidenta S.~Vojcehovska un maršala J.~Pilsudska tikšanās. Abus vienoja draudzīgas attiecības vēl no politiskās darbības sākuma. Maršals apgalvoja, ka viņa uzstāšanās nav vērsta pret valsti, bet tikai pret V.~Vitosa valdību, pieprasīja to atlaist. Taču prezidents tomēr vēlējās saglabāt likumību, valsts konstitucionālos pamatus, atteicās pakļauties dumpīgā maršala prasībai. Ar to J.~Pilsudska organizētā militāristu uzstāšanās vairs nebija kvalificējama kā stihiska akcija pret nepopulārām partijām, bet kļuva par atklātu dumpi, kurā militāristi, aizmirsuši savu zvērestu, centās diktēt savu gribu konstitucionālām valsts iestādēm. Maršalam atlika cerēt uz to, ka uzvarētājus netiesā, uz bruņotu spēku, arī uz Varšavas iedzīvotāju un kreiso partiju atbalstu.

Kreiso vidū bija darbinieki, kuri saskatīja faktisko bezdibeni starp šo partiju un bijušā sociālista J.~Pilsudska uzskatiem. Piemēram, PPS Izpildu komitejas priekšsēdētājs F.~Perls tās kongresā paziņoja: ``Pilsudskis nav mūsu cilvēks. Nevar no viņa taisīt kreisa parauga Musolini'' (Itālijas fašistu vadonis). Taču kopumā kreisajās partijās un to atbalstītāju rindās stiprāki izrādījās noskaņojumi par labu J.~Pilsudskim. Viņa akciju atbalstīja ievērojama daļa valsts iedzīvotāju. Varšavā manifestantu pūļi izteica savu solidaritāti ar J.~Pilsudska rīcību, jo uzskatīja, ka ar to tiek novērsta V.~Vitosa labējas diktatūras nodibināšanās. Tikai rietumu rajonos, kur stiprāka bija labējo un centrisko partiju ietekme, pārsvarā bija negatīva reakcija.

Sākās bruņotas sadursmes, kas veda pie upuriem no abām pusēm. (Kaujās krita 379~cilvēki, tai skaitā 164~civilisti, 920 tika ievainoti.) J.~Pilsudskis nebija rēķinājies, ka arī armijā bija notikusi šķelšanās~--- daļa no tās atbalstīja J.~Pilsudski, daļa~--- konstitucionālo labējo un centristu valdību, daļa~--- palika neitrāla. Piemēram, Krakovas kara apgabala pavēlnieks ģenerālis J.~Hallers un Poznaņas kara apgabala pavēlnieks J.~Dovbors-Musņickis mēģināja palīdzēt V.~Vitosa valdībai, taču tas neizdevās dzelzceļnieku streika dēļ. Ģenerālis V.~Sikorskis, kurš maija apvērsuma laikā komandēja Ļvovas kara apgabalu, deva pavēli karspēka daļām būt kaujas gatavībā, bet acīmredzot, novērtējis spēku samēru, tālāk nerīkojās, nesniedza atbalstu likumīgajai valdībai. J.~Pilsudskim tuvais ģenerālis K.~Sosnkovkis, kurš komandēja korpusu Poznaņā, par apvērsuma plāniem nebija informēts, acīmredzot nonāca morālas dilemmas priekšā~--- ko atbalstīt, likumīgo varu vai savu bijušo ``brigadieri'' J.~Pilsudski, un tāpēc neveiksmīgi mēģināja izdarīt pašnāvību, kā rezultātā mēnešiem ilgi viņa dzīvība bija apdraudēta.

Panākumus J.~Pilsudskim lielā mērā nodrošināja PPS, kura organizēja 14.~maijā sākušos dzelzceļnieku streiku, tāpēc valdība nespēja savlaicīgi pievilkt palīgspēkus, kamēr J.~Pilsudskis tos saņēma. PPS sludināja, ka apvērsums esot militārās demokrātijas revolūcija pret izkurtējušo labējo un centristu vadīto valsti. Kreisie cerēja, ka uzvarās gadījumā J.~Pilsudskis atlaidīs Seimu, noteiks jaunas tā vēlēšanas, kurās kreisie varētu uzvarēt un nonākt pie varas. Taču J.~Pilsudskis centās saglabāt apvērsuma militāro raksturu, lai arī turpmāk nenonāktu atkarībā no politiķiem.

Kad sazvērnieki sagrāba gandrīz visu Varšavu, ģenerāļi S.~Hallers (ģenerāļa J.~Hallera brālēns, tai laikā Ģenerālštāba priekšnieks) un T.~Razvadovskis ierosināja turpināt cīņu, evakuējot prezidentu un valdību uz Poznaņu. Valdības vadītājs V.~Vitoss iebilda, uzskatot, ka pilsoņu karam būs postošas sekas. Viņu atbalstīja citi ministri. Valdības rīcībā bija vēl daudz spēku, bet to pielietošana nozīmētu pilsoņu kara izraisīšanu, bet sliktākajā gadījumā~--- mazākumtautību sacelšanos vai, kā tika uzskatīts, pat Vācijas vai PSRS iebrukumu. To šie politiķi nevēlējās. (Šis arguments bez īpaša iemesla vēl joprojām tiek izmantots zinātniskajā literatūrā, kaut Polijas robežas garantēja Tautu savienība, un maz ticams, ka šai laikā militāri vājās Vācija un PSRS būtu nolēmušas uzsākt atklātu agresiju pret savu kaimiņu.) Uzklausījis visus, T.~Vojcehovskis paziņoja, ka noliek prezidenta pilnvaras, kuras pēc Konstitūcijas pārgāja Seima maršala M.~Rataja rokās. Atkāpās arī valdība.

M.~Ratajs paziņoja par to, ka viņš pieņēmis valdības atkāpšanos un drīzumā tiks nozīmēta jauna valdība un šai sakarā notiks tikšanās ar J.~Pilsudski. Tas liecināja, ka prezidenta vietas izpildītājs neuzdrošinājās pieņemt lēmumus, nenoskaidrojis maršala gribu. Šie notikumi nozīmēja likumīgās varas kapitulāciju bruņota spēka priekšā. Parlamentārisma posmu nomainīja J.~Pilsudska diktatūra.

Pats J.~Pilsudska sekojoši formulēja likumīgās varas likvidēšanu un apvērsumu attaisnojošos argumentus: ``Es izlēmu to [apvērsumu~--- V.Š.] sarīkot atbilstoši savai sirdsapziņai, un neuzskatu par vajadzīgu šai sakarā sniegt paskaidrojumus. Galvenie cēloņi patreizējam stāvoklim Polijā~--- nabadzībai, iekšējam un ārējam vājumam~--- bija zagšana, kura palika nesodīta. Polijā pār visu valdīja indivīda un partijas intereses, valdīja nesodāmība par visām ļaunprātībām un neģēlībām.'' Pieminēja viņš arī valsts arvien lielāko atkarību no ``jaunbagātniekiem''. Ar to diktators mēģināja iestāstīt savas valsts neatkarību atjaunojušajiem, bet no grūtībām nogurušajiem vidusmēra poļiem, ka viņš zina vainīgos un novērsīs nebūšanas. Visās jaunajās valstīs, arī starpkaru Polijā, uz kādu laiku šis arguments darbojās neatvairāmi. Masās valdīja dziļš sarūgtinājums par valstī valdošo korupciju. Plaši izplatījās uzskats, ka tikai ``stiprs'' vīrs pie varas spēs to apkarot.

PPS darbiniece, publiciste un vēsturniece L.~Celkoša ir atzīmējusi arī to, ka 1926.~gada maija apvērsums sakrita ar fašisma uzplaukumu Itālijā, bet poļu nacionāldemokrāti atbalstīja fašisma ideoloģiju. J.~Pilsudskis ar savu apvērsumu viņus apsteidza, tāpēc arī demokrāti un kreisie~--- ļudovieši, PPS atbalstīja apvērsumu.

Arī vēl mūsdienās vēsturnieku vidū turpinās strīdi par starpkaru parlamentārās iekārtas vērtējumu Polijā.

XX gadsimta 20.~gados poļu tiesību vēstures pētnieks V.~Komarņickis uzskatīja, ka Marta konstitūcijas radītā parlamentārā sistēma tā arī netika pilnībā realizēta, jo Seimā nespēja noformēties stabils vairākums, spējīgs pārvaldīt valsti. To, ka parlamentārisms Polijā nenostiprinājās, pēc viņa viedokļa, pierādīja ``speciālistu'' valdību veidošana, kas nozīmējis parlamentārās sistēmas noliegšanu. Bez stingra vairākuma parlamentā šīs valdības, ka arī koalīciju veidotie ministru kabineti nespējuši ieturēt konsekventu un ilglaicīgu politiku, kas arī novedis pie visas Marta konstitūcijas radītās politiskās sistēmas krīzes.

Par Polijas parlamentārismu rakstīja arī pazīstamais vēsturnieks V.~Konopčinskis, kurš bija Seima deputāts un varēja tieši novērot tā darbību. V.~Konopčinskis uzskatīja, ka parlamentārajai demokrātijai Polijā piemita būtiski trūkumi, kas traucēja tai efektīvi funkcionēt. Viņš neslēpa savu negatīvo attieksmi pret poļu kreisajiem, īpaši sociālistiem. Pēc viņa domām tie bija atbildīgi par parlamentārās demokrātijas likvidēšanu un turpmāko valsts attīstības ceļu.

Starp Polijas starpkaru politiskajai vēsturei veltītajiem pētījumiem jāatzīmē strādnieku un sociālistiskās kustības darbinieka, vēsturnieka un publicista Ā.~Pruhņika 30.~gados publicēts darbs, kurš vairākkārt tika izdots arī pēc kara un ietekmēja daudzu politisko procesu vērtējumus. Tajā bija pamatīgi raksturota Polijas partiju sistēma, partiju spēku samērs Seimos, taču klaji izpaudās arī autora simpātijas savai~--- sociālistiskajai partijai.

Pēc Otrā pasaules kara Polijas vēsturiski-juridiskajā literatūrā nostiprinājās vienkāršota, varētu teikt~--- staļiniska, pieeja valstiski politiskajām parādībām. Tiesību vēsturnieka A.~Gvidža grāmatā 1921.~gada Konstitūcija tika vērtēta kā reakcionāra, jo aiz tās ārēji demokrātiskajiem principiem slēpusies aizmaskota ekspluatatoru tirānija. Jau 60.~gados M.~Petšaks secināja, ka Polijas Seimā bija vieglāk pierunāt dažādas partijas ``gāzt'' valdību, nekā iesaistīt tās valdības koalīcijas veidošanā. To viņš skaidroja ar asu sociālu konfliktu esamību poļu sabiedrībā. No otras puses, zinātnieks arī redzēja Polijas parlamenta modeļa institucionālos trūkumus. Apmēram tais pat gados iznāca arī pazīstamā tiesību vēstures speciālista profesora A.~Ajnenkela darbs, kurā bija analizēti arī dažādu politisko spēku priekšlikumi veikt pārmaiņas valsts politiskajā sfērā un secināts, ka labējās un centriskās partijas attīstījās reakcijas virzienā, bet kreisās uzstājās kā parlamentārisma un demokrātijas aizstāves. Pētnieks arī norādīja, ka liels grūtību avots Polijas politiskajā dzīvē bija nacionālie mazākumi. Nebija radīta politiska sistēma, kura ļautu savienot valsts un nacionālo mazākumu intereses. A.~Ajnenkels uzsvēra, ka to izdarīt nepieļāva pirmkārt poļu labējie. Vēlāk A.~Ajnenkels arī rakstīja, ka ``Seims savu patstāvību saglabāja pārāk īsu laiku, zaudēja savas pilnvaras pārāk dramatiskos apstākļos, lai varētu tam dot maksimāli objektīvu vērtējumu. Tomēr, var riskēt apgalvot, ka tas, neraugoties uz dažādām kļūmēm, bija orgāns, kas funkcionēja samērā labi.'' Arī mūsdienu krievu vēsturnieks M.~Bulahtins, vērtējot Polijas valstiski politiskās attīstības problēmas 20.~gadu pirmajā pusē, uzskata, ka poļu parlamentārisma grūtības liecināja ne par parlamentārās demokrātijas kā tādas krīzi, bet tikai par 1921.~gada konstitūcijas radītā parlamentārisma modeļa krīzi. Tas Polijas politiskajos apstākļos parādīja zemu funkcionālo efektivitāti, kas veicināja sabiedrības vilšanos parlamentārismā, veda pie tās atsvešināšanās no parlamentārisma kopumā, radīja barojošu augsni antiparlamentārai propagandai un antidemokrātisko spēku pieaugumam. Parlamentārisma vājums XX gadsimta 20.~gadu pirmajā pusē būtiski ietekmēja arī valsts turpmāko attīstību.

Tādejādi var izdalīt trīs historiogrāfiskās tradīcijas Polijas parlamentārās demokrātijas vērtējumā. J.~Pilsudska piekritēji viennozīmīgi negatīvi vērtēja šādu valsts uzbūves modeli. Otro pieeju raksturo kritiska pieeja ar 1921.~gada Konstitūcijā iedibinātās parlamentārās demokrātijas modelim, saskatot tajā institucionālus trūkumus un funkcionālu vājumu. Par trešo tradīciju var uzskatīt pieeju, kura cīņā par konstitūcijas uzlabošanu saskata reakcijas un demokrātijas cīņas izpausmes. Tāda pieeja dominēja poļu historiogrāfijā ``sociālisma'' periodā. Mūsdienās, neraugoties uz augsto J.~Pilsudska vērtējumu poļu sabiedrībā, galvenokārt tiek caurskatītas un precizētas abas pēdējās pieejas.

Lai nu kā, apvērsuma veiksmīgā norise rādīja Polijas politiskās sistēmas nenobriedušo raksturu.

J.~Pilsudskis ar saviem piekritējiem spēja likvidēt demokrātisko iekārtu, taču reizē deva triecienu savai kā ``visas Polijas balss'' reputācijai. Notikumi parādīja, ka nebūt ne visa Polijas tauta atzina viņu par neapstrīdamu vadoni.

\subsection{``Sanācijas'' režīms. 1926--1939}

14.~maijā tika izsludināts pamiers. 15.~maijā bez konsultācijām ar prezidenta pilnvaru turētāju Seima maršalu M.~Rataju un politisko partiju līderiem J.~Pilsudskis uzdeva izveidot it kā pāri partijām stāvošu speciālistu valdību savam piekritējam matemātikas profesoram K.~Bartelam. Profesors 1926.~gadā bija piecu cilvēku lielas Seima deputātu grupas, atšķēlušās no ļudoviešu partijas \pltxti{Wyzwolenie}, līderis. Prezidenta vietas izpildītājs~--- Seima maršals M.~Ratajs jauno valdību tūlīt apstiprināja. Šāda prakse turpmāk kļuva tradicionāla. Valdību krīžu vairs nebija, bet tās mainījās pēc maršala gribas. (Līdz 1939.~gadam nomainījās 12~valdības.) Nebija partiju ķīviņu par portfeļiem un citu nīstās ``seimokrātijas'' atribūtu. Formāli visās valdībās līdz pat savai nāvei pats J.~Pilsudskis ieņēma kara ministra vietu, kā arī faktiski vadīja Polijas ārpolitiku. (Visiem premjeriem vajadzēja rēķināties, ka J.~Pilsudskis pats izlems visus ārpolitikas jautājumus. To zināja arī Polijas aizrobežu partneri.) Divreiz: 1926.--1928.~un 1930.~gadā J.~Pilsudskis ieņēma arī premjerministra vietu. Pilnīgi mainot valdības komandu, J.~Pilsudski parādīja, ka viņš nejūt kādu pateicību par savu uzvaru nevienai no partijām. Militārais diktators piesavinājās tiesības iecelt valdību bez apspriešanās ar parlamentu. Un politiskās partijas tam piekrita.

K.~Bartela valdībā nebija agrāko valdību locekļu, citu ievērojamu politiķu, arī no partijām, kuras atbalstīja apvērsumu. Visi ministri bija jauni izpildvaras koridoros. Viņa valdīšanai tika dots apzīmējums ``\pltxti{bartlowanie}'', ar ko saprata J.~Pilsudska izvēlētas un pēc viņa norādījumiem strādājošas valdības pastāvēšanu. Tādejādi J.~Pilsudskis, leģionāri un citi viņu atbalstošie spēki ne tikai atņēma tautas masām iespēju likumīgā, demokrātiskā ceļā ietekmēt savas valsts likteni, bet atbīdīja no varas arī buržuāziskās politiskās partijas, pārstāvošas vismaz daļas no valdošā slāņa intereses.

Kāpēc J.~Pilsudskis uzreiz nelikvidēja parlamentārismu un visu varu nekoncentrēja savās rokās, liecina viņa izteikums kādam vācu politiķim pēc tam, kad Vācijā 1933.~gadā varu bija pārņēmuši nacisti. Lūk, J.~Pilsudska teiktais: ``Es nesaprotu, kāpēc fīrers grib visus valsts vadības pavedienus apvienot savās rokās. Taktiski ir pareizāk, ja pašam nav jānes par visu pilna atbildība''. Tas viņam gan netraucēja visai bieži ar rupjām lamām vērsties pret parlamentu. 1928.~gadā kādā savā rakstā avīzē ``\pltxti{Gazeta Polska}'' (``Polijas Avīze''), kurā gāja runa par Seimu, diktators 37~reizes dažādos savienojumos izmantoja vārdu ``sūdi'' (poļu: \pltxti{gówno}), kuru ne pirms, ne pēc tam valdības oficiozs nelietoja. Arī sarunās ar saviem piekritējiem savu tautu viņš sauca par ``idiotu tautu'', kura ir jāaudzina ar nūju. Viņa atbalstītāji tāpat atzina: ``Tas ir vienīgais tonis, ko pie mums [Polijā] saprot''.

Pēc apvērsuma J.~Pilsudskis centās izvairīties no galējībām, pēc iespējas ātrāk atgriezties pie legālisma. To diktēja nevis cieņa pret likumību, bet politiskās intereses. Likumīgā prezidenta un premjera padošanās piespieda arī uzvarētājus vismaz skata pēc atteikties no atriebības. Tie virsnieki, kuri cīnījās pret apvērsuma dalībniekiem, to darīja nevis aizstāvot konkrēto V.~Vitosa valdību, bet gan likumīgo, legālo varu. Agrāko varas nesēju pilnvaru nolikšana un to nodošana jaunās rokās it kā atrisināja pretrunu. Apvērsuma pretinieki varēja atgriezties kazarmās, pakļauties J.~Pilsudskim kā jaunajam kara ministram, jo apvērsuma dalībnieku pusē tagad it kā nostājās konstitūcija un likums. Lai kaut kā nolīdzinātu tautā un arī armijā radušos plaisu, J.~Pilsudskis bija spiests izteikt savu cieņu arī cīņā pret viņu kritušajiem. Vienā no pirmajām pavēlēm pēc apvērsuma J.~Pilsudskis rakstīja: ``Vienā un tai pašā zemē iesūcās mūsu asinis, zemē, kura tik dārga gan vieniem, gan otriem, kuru mīl abas puses. Tad lai nu šīs karstās asinis, pašas vērtīgākās Polijā kareivju asinis kļūst zem mūsu kājām par jaunu brālības sēklu. \citespace{} Lai Dievs ir žēlīgs un atlaiž mūsu grēkus un novērš savu sodu, bet mēs ķersimies pie darba, kurš nostiprinās un atjaunos mūsu zemi''. Simbolisku nozīmi ieguva lēmums par kopēju svinīgu sēru ceremoniju, apbedījot abās pusēs kritušos apvērsuma upurus. 17.~maijā Varšavā tika izsludinātas sēras un notika savstarpējās cīņās kritušo militāristu apbedīšana ar valdības locekļu piedalīšanos. Kara ministrs J.~Pilsudskis tajā gan nepiedalījās. Samierināšanās ieceri diskreditēja arī aizlūguma laikā garnizona baznīcā notikušais skandāls, kurš ieguva plašu sabiedrisku rezonansi. Popularitāti baudošais bijušo Poļu leģionu kapelāns, 1920.~gadā iznākušo atmiņu, kurās tika slavināts J.~Pilsudskis, autors J.~Panašs norāva no krūtīm militāros apbalvojumus un nometa pie viena no apvērsuma vadītājiem ģenerāļa G.~Orliča-Drešera kājām. (Vēlāk J.~Panašs tuvinājās zemnieku kustībai, kļuva par vienu no pret J.~Pilsudski vērstās opozīcijas simboliem).

J.~Pilsudska apvērsumu nosodīja ģenerālis J.~Hallers, tāpēc jau 1926.~gada jūlijā tika atlaists no armijas. Atvaļināja ģenerāli S.~Šepticki. Vairāki ģenerāļi tika arestēti un ieslodzīti Viļņas cietumā, tai skaitā Polijas armijas Ģenerālštāba priekšnieks 1918.--1920.gadā T.~Rozvadovskis. Oficiāli viņiem tika izvirzītas apsūdzības kriminālnodarījumos, presē izvērstas melu kampaņas. Faktiski tā bija atriebība par pretošanos J.~Pilsudskim un izstumšana no politiskās cīņas lauka. Pēc nepierādītām hipotēzēm ģenerāļa T.~Rozvadovska ātru nāvi izsauca noindēšana. Ģenerāli J.~Malčevski tiesa attaisnoja, viņš tika atbrīvots, taču 1927.~gadā atvaļināts no dienesta. It kā par dalību krāpšanā, piegādājot armijai gāzmaskas, bet pēc būtības par pretošanos apvērsumam, tika degradēts par ierindnieku un uz pieciem gadiem cietumā notiesāts bijušais J.~Pilsudska līdzbiedrs Poļu leģionā ģenerālis M.~Žimerskis. Pēc atbrīvošanas viņš kļuva par Polijas Komunistiskās partijas biedru, vēlāk Armijas Ļudovas komandieri Otrā pasaules kara gados, pēc tam Polijas maršalu un Polijas Tautas armijas virspavēlnieku.

Par kauna traipu jaunajai varai kļuva ģenerāļa V.~Zagorska pazušana. Pēc atbrīvošanas no cietuma 1927.~gada augustā viņš ieradās Varšavā un mīklaini pazuda. Par to tika izvirzītas dažādas versijas. Oficiālā viņu apvainoja dezertēšanā. Cita versija skanēja, ka V.~Zagorskis kā bijušais Austroungārijas armijas ģenerālštāba izlūkdienesta darbinieks pratis inscenēt savu pazušanu. Taču J.~Pilsudska pretinieki apvainoja maršalu slepenā politiskā slepkavībā. J.~Pilsudska biogrāfi poļu vēsturnieki Nalenči uzskata, ka V.~Zagorski nogalināja maršala pakļautie, bet pats viņš neko nedarīja, lai noziedznieki tiktu sodīti. Gluži pretēji, viņi strauji kāpa pa karjeras kāpnēm. Arī mūsdienu Polijas presē atzīmēts, ka ģenerālis V.~Zagorskis ``pārāk daudz zināja'', tai skaitā arī par J.~Pilsudska ilggadīgo sadarbību ar Austroungārijas izlūkdienestu un par to, kā Vācijas izlūkdienests panāca J.~Pilsudska atbrīvošanu no ieslodzījuma 1918.~gada nogalē, kas arī ļāva viņam nonākt pie varas. Mīklainā, pareizāk~--- aizdomīgā, nāvē mira arī vēl ap pusotra desmita Polijas armijas virsnieku, īpaši to, kuri stāvēja J.~Pilsudska ceļā Poļu leģiona formēšanas gaitā Pirmā pasaules kara gados.

Tika atvaļināti no dienesta arī citi ģenerāļi, to vidū J.~Dovbor-Musņickis, kuru savulaik J.~Pilsudska vietā par savu vadoni gribēja redzēt Poznaņā sacēlušies poļi. Kopumā no armijas tika atlaisti 37 ģenerāli. Daži atkāpās pēc saviem ieskatiem. Tiesa, lielākā daļa likumīgās valdības aizstāvju turpināja kalpot Polijas armijā, un daži no viņiem pat ieguva ģenerāļa pakāpi.

Interesanta bija maršala rīcība pret ģenerāli V.~Sikorski, kam pirms apvērsuma viņš presē uzbruka visasāk. Acīmredzot ar nolūku konsolidēt armiju, J.~Pilsudskis atteicās no atriebības viņam. Līdz 1928.~gada martam V.~Sikorskis palika Ļvovas kara apgabala pavēlnieka amatā. Taču tā kā viņš grāmatā ``Pie Vislas un Vikras: poļu-krievu kara studija'' (\pltxti{``Nad Wisłą i Wkrą'': studium z polsko-rosyjskiej wojny, Lwow-Warszawa-Krakow, 1928}) izklāstīja savus uzskatus par kauju pie Varšavas, kuri nesakrita ar oficiālajiem, J.~Pilsudski slavinošajiem, viņa liktenis ar to bija izlemts. Maršals gan rīkojās tam laikam visai oriģināli, nevis atvaļināja V.~Sikorski no armijas, kas ļautu tam nodarboties ar politisku darbību, bet atstāja kara ministrijas rīcībā, nepiešķirot nekādu amatu. V.~Sikorskis tai pašā 1928.~gadā devās uz Franciju, kur bija opozīcijā Polijā nodibinātajam režīmam. Atbīdīts no armijas, viņš rakstīja grāmatas savas militārās specialitātes jomā, kuras tika tulkotas arī krievu valodā. V.~Sikorska grāmata ``Nākotnes karš'' (\pltxti{Przyszła Wojna}, 1934) tika kā mācību grāmata izmantota Sarkanās armijas Ģenerālštābā.

Armijā notika augstāko kadru nomaiņa. Virsnieki, nākuši no t.s. ``Vīnes karaskolas'' tika nomainīti ar ``leģionāriem''. Tā, no 1928.~līdz 1931.~gadam leģionāru skaits armijas ģenerāļu vidū pieauga no 44~cilvēkiem (54\% no visiem) līdz 54 (74\%), kara apgabalu pavēlnieku vidū no 5 (50\%) līdz 8 (80\%). J.~Pilsudska piekritēji taisīja karjeru, bet viņa pretiniekus vai nu atvaļināja, vai vienkārši nepiešķīra tiem augstākas pakāpes un amatus. Tomēr, vadoties no vēlmes piesaistīt sev visu armiju, J.~Pilsudskis neveica tik kardinālu virsnieku korpusa tīrīšanu, kādu daudzi gaidīja. Kadru pārbīdes, kuras viņš izdarīja, neatbilda uzvaras mērogiem. Viņš atbrīvojās tikai no tiem ģenerāļiem, kuru antagonisms pret viņu bija nepārvarams un kuri nekādi nevēlējās samierināties ar notikušo.

Poļu vēsturnieku vidū pastāv viedoklis, ka J.~Pilsudskis gandrīz pilnībā atteicās no represijām pret zemākajiem virsniekiem, īpaši nevajāja tos, kuri bija cīnījušies pret viņu. To virsnieku, kuri bija izcēlušies kaujās pret apvērsuma dalībniekiem, piemēram, pulkveža V.~Andersa, karjera gan palēninājās, bet neapstājās. Tomēr ar augstāk minēto viedokli pretrunā ir fakti, ka daudzi virsnieki bija spiesti atstāt dienestu. No 1926.~līdz 1934.~gadam virsnieku korpuss zaudēja 6~032 cilvēku, to skaitā arī daudz gados jaunu. Noplaka arī līdz apvērsumam maršala piekritēju nometnē tik izplatītā spriedelēšana par morāles lomu valstī un armijā. Konflikta politisko un ideoloģisko cēloņu uzsvēršanas vietā tagad tas tika tēlots kā brālīga ķilda, kura diemžēl beidzās ar kautiņu. Tas viss atļāva ātri nomierināt armiju.

Kopumā pēc apvērsuma militāristu loma valstī pieauga. Jau tas vien, ka, lai izvēlētos virsnieka karjeru, bija jābūt vidējai izglītībai, nodrošināja zināmu šķirisku atlasi. 1926.~gada 5.~augustā valdība paaugstināja virsnieku algas, kas bija samērojamas ar rietumvalstu virsnieku atalgojumu. Armija, īpaši tās virsnieku korpuss (ap 18~000 kadru virsnieku), bija un palika viens no Polijas valsts nozīmīgākajiem balstiem.

Ministrijās, citās centrālajās, kā arī vietējās iestādēs tāpat notika diezgan būtiska kadru nomaiņa, iepriekšējo vietā nozīmējot tādus ierēdņus, kas pirmajā vietā stādīja maršala autoritāti. Tā dažu mēnešu laikā tika nomainīti vienpadsmit no septiņpadsmit vojevodām. Kā savā dienasgrāmatā atzīmēja Seima priekšsēdētājs M.~Ratajs, nomaiņām netika pat meklēts pamatojums. ``Padotais, kuram priekšnieks bija nodarījis pāri, vienkārši rakstīja par viņu denunciāciju, ka tas ir endeks jeb pjastovietis, dažkārt piebilstot, ka arī ``zaglis'', kuru sen laiks padzīt no darba. Nesenējie nelabojamie endeki un pjastovieši demonstratīvi izstājās no savām organizācijām un līdzīgi dažādiem konvertītiem [cilvēkiem, kas pārgājuši citā ticībā] iestājās par ``maija kārtību'' vēl cītīgāk par pašu Pilsudski''.

Neraugoties uz plašu iedzīvotāju slāņu cerībām, politiskie ieslodzītie netika amnestēti, netika veiktas kādās izmaiņas valsts politikā ``kresos''. Bez sekām palika kreisi orientēto deputātu pieprasījums par amnestiju, par neatkarīgo preses izdevumu vajāšanu. Tādu sanācijas režīma nevērīgu attieksmi pret slāņiem, kuri zināmā mērā bija sekmējuši varās pāreju tā rokās, vismaz daļēji varēja skaidrot ar tai laikā sākušos rosību ekonomikā, par ko tiks runāts attiecīgajā apakšnodaļā.

1926.~gada maija apvērsuma rezultātā Polijā tika izveidots politisks režīms, kurš likvidēja tos principus, kuri atradās 1921.~gada Marta Konstitūcijas pamatā. Parlamentārais valsts funkcionēšanas mehānisms tika nomainīts ar autoritārām pārvaldes metodēm.

Taču visiem par brīnumu maršals atteicās atlaist Seimu un sarīkot jaunas vēlēšanas. Tam bija dažādi skaidrojumi.

\begin{enumerate}
\item J.~Pilsudskis nerīkoja parlamenta vēlēšanas, baidīdamies no kreiso uzvaras tajās (pats J.~Pilsudskis bija bezpartijisks, savus agrākos sabiedrotos~--- kreisos par līdzvērtīgiem partneriem vairs neatzina). Vēlēšanas maršals vēlējās sarīkot tad, kad valsts pēc apvērsuma būs nomierinājusies.

\item Sakompromitētā parlamenta pastāvēšana bija J.~Pilsudskim izdevīga. Jaunās parlamenta vēlēšanās uz panākumiem varētu cerēt tikai tās partijas, kuras atbalstīja maršalu~--- PPS, kreisie ļudovieši no partijas \pltxti{Wyzwylenie} un Zemnieku partijas (\pltxti{Stronnictwo Chłopske}). Taču tad J.~Pilsudskis pats būtu atkarīgs no šiem relatīvi kreisajiem spēkiem, bet viņš vēlējās paturēt lēmējvārdu visos strīdos. Viņam, kā rakstīja viens no ievērojamiem PPS darbiniekiem A.~Prohņiks, izdevīgāks bija iepriekšējais Seims, ``vājš, bez autoritātes \citespace{} ciestā zaudējuma un savas paša nespēcības pazemots \citespace{} nekā jauns, spēcīgs ar tikko tam izrādīto masu uzticību''. Maršals bija konservators un ignorēja parlamentārās darbības metodes. Diktatora attieksmi pret parlamentu uzskatāmi demonstrēja viņa vārdi, veltīti Seima deputātiem 29.~maijā, divas dienas līdz Nacionālās sapulces sanākšanai: ``Tagadējā situācijā es varētu jūs neielaist Nacionālās sapulces zālē, izsmiedams jūs visus, taču es pārbaudīšu, vai pār Poliju vēl var valdīt bez pletnes palīdzības''. Polijas parlamentu viņš vērtēja kā ``politisko prostitūtu Seimu''. Valdības veidošanā, daudzu svarīgu jautājumu izlemšanā J.~Pilsudskis faktiski nerēķinājās ar parlamentu.

\end{enumerate}

Taču J.~Pilsudska varas režīma izveides īpatnība bija tā, ka reizē ar nicinošās attieksmes pret Seimu demonstrēšanu maršals varēja balstīties uz tā deputātu vairākumu. Kā uzskata maršala biogrāfi poļu vēsturnieki D.~un T.~Nalenči, parlamenta padevību nevar izskaidrot tikai ar tā locekļu vēlmi izbēgt no iespējamām represijām. To noteikusi arī ticība, ka maršals nevēlas Polijā ievest diktatorisku režīmu, kurš neatstātu nekādu vietu parlamentārismam. Cits J.~Pilsudska biogrāfs krievu vēsturnieks G.~Matvejevs jau plašāk atklāja parlamentāriešu rīcību, atbalstot maršalu. Vieni izskaidroja balsošanu par J.~Pilsudska priekšlikumiem ar savu pārliecību par to nepieciešamību valsts attīstībai, citi akli ticēja maršala ģēnijam, trešie, kaut to slēpa, vadījās no savtīgām interesēm, baidījās zaudēt privilēģijas, ko deva deputāta vieta. Pret J.~Pilsudska režīmu, kaut dažādu motīvu vadīti, iestājās galvenokārt endeku, komunistu un slāvu nacionālo mazākumu pārstāvji. Tikai viņi zināmā mērā ļāva parlamentārismam Polijā ``saglabāt seju''. Ja tā rīkotos arī citi deputāti, iespējams, politiskā attīstība valstī norisinātos citā virzienā. Tad režīmam nāktos balstīties tikai uz spēku, galvenokārt armiju, kur tūlīt pēc apvērsuma J.~Pilsudska pozīcijas vēl nebūt nebija visspēcīgas.

31.~maijā parlamentā (kopīgā sēdē Seims un Senāts) pēc PPS priekšlikuma J.~Pilsudskis tika izvirzīts par prezidenta amata kandidātu. Nacionāldemokrāti par savu kandidātu izvirzīja Poznaņas vojevodu A.~Bņinski. Ar to endeki rādīja savu opozicionāro nostāju J.~Pilsudskim, nevēlēšanos samierināties ar jauno politisko realitāti. Zālē atradās 546~deputāti. Ar 292~balsīm parlaments ievēlēja J.~Pilsudski par valsts prezidentu. A.~Bņinskis saņēma 193~balsis. (61~biļetens izrādījās neaizpildīts).

Ziņa par J.~Pilsudska ievēlēšanu prezidenta amatā izsauca stihiskus viņa piekritēju mītiņus un demonstrācijas. Taču drīz sekoja sensacionāla ziņa~--- ievēlētais atteicās no amata. Vēstulē Seima maršalam M.~Ratajam J.~Pilsudskis skaidroja atteikšanos gan ar atmiņām par nogalinātā pirmā Polijas prezidenta G.~Narutoviča traģisko figūru, kuru viņš nespējis pasargāt no briesmīgā likteņa, gan ar vēlēšanos veikt tiešu darbu, no kura pastāvošā konstitūcija prezidentu turēja atstatus un tas neatbilda J.~Pilsudska raksturam. Būtībā pēc tam, kad parlaments faktiski bija izpļaukājis pats sevi un leģitimējis valsts apvērsumu, J.~Pilsudskis norobežojās, nevēlējās saņemt pilnvaras no tā. Acīmredzot maršalam vairāk rūpēja kontrole pār bruņotajiem spēkiem, bet prezidenta amats automātiski lika to nodot, lai varbūt arī personīgi uzticama, bet tomēr cita cilvēka rokās.

Piemērota prezidenta kandidāta meklējumos, premjers K.~Bartels ierosināja un J.~Pilsudskis pieņēma priekšlikumu prezidenta amatā ievēlēt savu piekritēju, ķīmijas profesoru I.~Moscicki, savā laikā studējušu Rīgas Politehniskajā institūtā (Savās atmiņās I.~Moscickis rakstīja, ka Rīga gadsimtu mijā bijusi ``sava veida galvaspilsēta austrumu pierobežas novadu poļiem), un bez zinātnieka karjeras pierādījušu sevi arī kā veiksmīgu menedžeru (viņš bija valsts slāpekļa savienojumu rūpnīcas Hožuvā (\pltxti{Chorzów}) direktors). I.~Moscickim nebija izteiktu politisku ambīciju, aiz viņa nestāvēja arī kādi noteikti politiski spēki. Par to, ka diktators grib viņu izvirzīt par prezidenta kandidātu, īsi pirms tam I.~Moscickim paziņoja nevis pats J.~Pilsudskis, bet premjers K.~Bartels. Parlamentā vēlēšanu otrajā kārtā par I.~Moscicki nobalsoja 291, par otru kandidātu A.~Bņinski 200 cilvēku. I.~Moscickis kļuva par prezidentu (palika amatā no 1926.~gada līdz pat 1939.~gada septembrim). J.~Pilsudskis tā vēlreiz parādīja visai valstij savu vārdu par parlamentāriešu vairākuma bezprincipialitāti pareizību.

6.~augustā ar prezidenta dekrētu tika izveidots bruņoto spēku ģenerālinspektora amats. Ģenerālinspektoram kara gadījumā bija jākļūst par augstāko virspavēlnieku, par savu darbību viņš nenesa atbildību nedz valdības, nedz Seima priekšā. Šo amatu, tāpat kā kara ministra portfeli, ieņēma un līdz dzīves beigām saglabāja J.~Pilsudskis. Viņam bija pakļauts arī armijas Ģenerālštābs. Tā maršals nodrošināja savās rokās armijas vadību un arī pilnīgu kontroli pār ārpolitiku.

1926.~gada 2.~augustā ar 246~balsīm ``par'' un 96 ``pret'' Seima pieņemtie Konstitūcijas labojumi deva prezidentam tiesības atlaist parlamentu; parlamenta sesija, veltīta budžeta pieņemšanai, tika ierobežota, bet valdība ieguva tiesības uz pastāvēšanu arī budžeta neapstiprināšanas gadījumā, rīkojoties iepriekšējā gada budžeta ietvaros. Prezidents ieguva arī tiesības izdot dekrētus ar likuma spēku, ar noteikumu, ka šos aktus pēc tam bija jāapstiprina Seimam. Tas viss nostiprināja izpildvaru, kaut arī valdība joprojām palika atbildīga Seima priekšā. Kā jau teikts, faktiski Seims un Senāts legalizēja J.~Pilsudska un citu militāristu apvērsumu.

Tas, ka apvērsums notika brīdī, kad Polijā sākās ekonomisks pacēlums un atslāba sociālā spriedze, veicināja militāristu režīma nostiprināšanos. Arī īpašus diplomātiskus sarežģījumus Polijai apvērsums it kā neradīja. Pat PSRS ārlietu ministrijas pārstāvis Polijas pilnvarotajam pārstāvim Maskavā paziņoja, ka ``mēs raugāmies uz Pilsudski un viņa izveidoto valdību ``bez aizspriedumiem'' un esam gatavi turpināt mūsu tuvināšanās politiku ar Poliju''. Taču pēc autoritārā apvērsuma tomēr saasinājās Polijas attiecības gan ar PSRS, gan Lietuvu, gan citām valstīm, par ko būs runa aiznākamajā apakšnodaļā.

1926.~gada 2.~oktobrī maršals pats uzņēmās valdības vadību. Tās sastāvā bija iekļauti viņam tuvi darbinieki no dažādām partijām, kā arī divi konservatīvo aprindu pārstāvji. Faktiski J.~Pilsudskis bija valsts vadītājs ar lielu varu. Tiesa, šajā laikā pasliktinājās viņa veselības stāvoklis. 1928.~gadā viņš cieta no insulta ar ķermeņa labās puses paralīzi, pēc kuras īpaši roka tā arī neatguva pilnīgu kustīgumu. Lika sevi manīt arī aknu slimība. J.~Pilsudska kā premjera pienākumus praktiski pildīja K.~Bartels u.~c. diktatora uzticības personas. Taču J.~Pilsudskim, arī tad, kad viņš pats nevadīja valdību, valstī bija lielāka loma nekā premjeram un prezidentam.

Kādreiz saka, ka līdz pat nāvei 1935.~gadā J.~Pilsudskis bija Polijas ``pelēkais kardināls'' [Tā no Luija XIII laika sauc ietekmīgus cilvēkus, kuri darbojas slepus un parasti neieņem formālus amatus], taču tas nav pareizi, jo viņš savu varu neslēpa un arī tauta to redzēja. Jāatzīmē, ka J.~Pilsudskim nevajadzēja ievērot īpašus drošības pasākumus. Labos laika apstākļos uz darbu viņš gāja kājam adjutanta pavadībā, sasveicinoties ar garāmgājējiem, īpašu uzmanību veltot bērniem. Tikai 1930.~gadā, kad iekšlietu ministrs F.~Slavojs Skladkovskis ziņoja par it kā pastāvošo sazvērestību PPS rindās (par ziņu patiesumu pastāv dažādas domas), diktators piekrita pastiprināt drošības pasākumus.

Būtiskas pārmaiņas sabiedrības dzīvē bieži ir saistītas ar varas pārņemšanu atsevišķu personu rokās, tāpēc Polijā bieži runā par ``Pilsudska ēru'' (vai vēlāk~--- par ``Gomulkas ēru''). Šāda pieeja noder, lai īsi un kodolīgi iezīmētu attiecīgo posmu valsts vēsturē, kaut to personifikācija arī noved pie daudzveidīgo attiecīgajā laikā notikušo procesu vienkāršošanas. Tāpēc autors lietojis nevis šo nosaukumu, bet zinātniskajā apritē iegājušo terminu ``sanācijas režīms''.

Polijā pēc apvērsuma tika pasludināts morālās atveseļošanās jeb \strong{sanācijas lozungs}. Pēc būtības tas bija tikai cits nosaukums autoritārismam.

\strong{Autoritārisms} (no latīņu \latxti{auctoritas}~--- vara, ietekme, kr. \rutxti{власть}, fr. \frtxti{autoritaire}~---  valdonīgs) ir politiskās varas sistēma, kam raksturīga varas koncentrācija valsts vadītāja rokās. Ar šo terminu apzīmē antidemokrātisku valsts pārvaldes sistēmu, kas izveidota, lai nodrošinātu bezierunu pakļaušanos. Demokrātija tiek likvidēta vai ierobežota. Valsts kontrolē saimniecisko, sabiedrisko un garīgo dzīvi. Starpkaru periodā daudzās valstīs autoritāro režīmu priekšgalā nostājās jau labi pazīstami politiķi, kuri bija populāri sabiedrībā kopš valstu dibināšanas un brīvības cīņu laika, piemēram, 1926.~gadā Lietuvā~--- A.~Smetona, 1934.~gadā Igaunijā~--- K.~Petss, 1934.~gadā Latvijā~--- K.~Ulmanis. Kā autoritārās valsts formas var izdalīt: a)~karaļa diktatūru (monarhiju); b)~prezidenta diktatūru; c)~militāru diktatūru. Taču nevar likt vienlīdzības zīmi starp autoritāru valsti un diktatūru.

Diktatūrām var būt arī paveidi. Daži pazīstamākie: nacistiskais režīms, ko nodibināja Ā.~Hitlers Vācijā; B.~Musolīni fašistiskais režīms Itālijā; J.~Staļina diktatūra PSRS, citu ``komunistisko'' līderu diktatūras Austrumeiropā un citur; dažādi autoritārie režīmi (Latvijā, Lietuvā, Polijā u.c.). Jāuzsver, ka 1926.~~gada apvērsumu nevar saukt par fašistisku, arī sanāciju Polijā nevar raksturot kā fašistisku režīmu, kā to darīja komunisti. Pat jau t.s. ``pārbūves'' PSRS laikā~--- 1986.~gadā padomju vēsturnieki I.~Jažborovska un N.~Buharins rakstīja, ka ``J.~Pilsudskis padarīja savu diktatūru par zemes pakāpeniskas fašizācijas ieroci''. PTR pastāvēšanas laikā atsevišķos izdevumos tika runāts pat par valdošo slāņu ``totalitāro praksi''. Taču sanācijas periodā nebija vienas partijas kundzības, saimniecību valsts necentās pilnībā kontrolēt, netika pilnīgi iznīcinātas politiskās brīvības. Pareizāk to saukt par autoritāru diktatūru.

Autoritārisma iezīmes:

\begin{enumerate}
\item Autoritāro režīmu priekšgalā atrodas vadoņi, politiķi vai ģenerāļi, kuri, lai gan bauda popularitāti daļā tautas, tomēr nelīdzinās tādiem harizmātiskiem līderiem kā Ā.~Hitlers un B.~Musolini. Valsts vadītājam pieder plašas pilnvaras un viņš realizē savu politiku jebkuriem līdzekļiem~--- gan humāniem, gan antihumāniem. Varas aparāts balstās uz vecajiem valdošajiem slāņiem, notiek tikai ierobežotas pārmaiņas elitē, kuras rokās atrodas reālā vara~--- atšķirībā no fašistiskajiem režīmiem, kuri rada jaunu valdošo slāni.

\item Atšķirībā no fašistiskajiem režīmiem autoritārajos nepastāv masu partija ar saviem bruņotiem spēkiem, nenotiek plašu masu mobilizācija varas atbalstam. Tomēr tiek ierobežota vai aizliegta opozīcijas partiju darbība, parlaments atlaists vai tā tiesības ierobežotas, ieviesta cenzūra.

\item Lielākā vai mazākā mērā noris politiskās represijas. Neliela cilvēku grupa valda ar vardarbības palīdzību, balstoties uz armiju un policiju, kas nepielaiž pie varas ne tikai kreisos, bet arī labējos (arī fašistiskos) radikāļus. Plaši tautas slāņi ir pasīvi vai arī pretojas.

\item Ideoloģijā valda nacionālisms, par augstāko vērtību tiek izsludināta valsts, kuras interesēs ir pieļaujama personības brīvības ierobežošana. Tiek izvērsts Vadoņa kults. Atšķirībā no fašistiskajiem režīmiem, kuri centās ieviest jaunu, ``laicīgu'' reliģiju, autoritārās diktatūras cenšas nostiprināt tradicionālo reliģiju.

\end{enumerate}

Dažus vārdus par konkrēto vārda \strong{sanācija} (no latīņu \latxti{sanati}: ārstēšana, atveseļošana, attīrīšana) nozīmi: Ekonomikā ar to saprot uzņēmuma, bankas vai nozares finansiālā stāvokļa uzlabošanas un konkurētspējas palielināšanas pasākumus (aizdevumus, pabalstus u.tml.), ko parasti veic valsts par nodokļu maksātāju līdzekļiem. Polijā ar šo vārdu apzīmēja J.~Pilsudska pēc valsts apvērsuma iedibināto režīmu, kas bija vērsts uz to, lai ``atveseļotu'' un attīrītu valsti no daudzo partiju korumpētības un nekārtībām, ieviestu tajā karavīru tikumus: patriotismu, disciplinētību, pienākuma apziņu un nesavtību. J.~Pilsudskis prasmīgi aizstāja tautas masu un to partijas uztraucošo sociālo jautājumu risināšanu ar solījumiem iznīdēt korupciju, valdošās ļaunprātības, atspirdzināt morālo klimatu valstī. Reāli šī ``sanācija'' drīzāk bija vispārīga nostāja bez skaidri izstrādātas rīcības programmas. Faktiski tā bija valsts aparāta attīrīšana no nekompetentiem vai nevēlamiem darbiniekiem, kuri J.~Pilsudskim nesimpatizēja. J.~Pilsudskis uzņēmās tādu kā ārsta lomu pār visu sabiedrību, vienpersoniski noteikdams valsts ``ārstēšanas'' jeb atveseļošanas no krīzēm kursu.

``Sanācija'' bija politiska sistēma, kura apvienoja sabiedrības daļā populārā J.~Pilsudska personīgo diktatūru, kurš kontrolēja armiju un izpildvaru, ar saglabātām parlamentārās iekārtas paliekām darbojošos politisko partiju un pārstāvniecības iestāžu veidā. Pēdējās, gan atdevušas daļu tiesību izpildvarai, tomēr varēja vēl izrādīt zināmu pretestību diktatūrai.

Formāli ``sanācija'' apvienoja sevī 3 procesus:

\begin{enumerate}
\item armijas pasargāšanu no politiskas ietekmes (praktiski armija tika pārvērsta par J.~Pilsudskās personiskās varas realizācijas ieroci);

\item valsts pārvaldes ``veselīgu'' attīrīšanu un profesionalizāciju~--- amatos iecēla zinātniekus (matemātiķus, ekonomistus, inženierus) nevis politiķus. Pastāvēja valsts tehnokrātija~--- zinātnieku dominēšana valsts pārvaldē;

\item Režīms centās izveidot visu slāņu un partiju apvienību, kura ``paceltu kopējās valsts intereses pāri atsevišķām partejiskām un sociālām interesēm''.
\end{enumerate}

% page 490


Pastāvēja Vadoņa kults~--- J.~Pilsudski dēvēja par ``nācijas varonīgo tēvu'', ``gudro vadoni'' un ``labvēlīgo aizstāvi''. 1928.~gada novembrī Varšavas pilsētas padome nolēma Sakšu laukumu pārdēvēt par J.~Pilsudska laukumu. Sabiedriskās iestādēs sāka uzstādīt viņa krūšutēlus. Sākotnēji atļauju vēl jautāja pašam maršalam, bet saņemot atbildes: ``Uzstādiet, ko vēlaties'', liekus jautājumus vairs neuzdeva. Polijas pilsētas piešķīra viņam Goda pilsoņa nosaukumu, skolnieki nosūtīja viņam fotoalbūmus par savu laimīgo dzīvi sanācijas režīma apstākļos. 1934.~gada sākumā Jagelloņu universitāte Krakovā izvirzīja J.~Pilsudska kandidatūru Nobeļa prēmijas saņemšanai, bet tā paša gada jūnijā Krakovas pilsētas padome nolēma viņam par godu uzbērt kurgānu tuvējā Soviņecas (\pltxti{Sovinec}) augstienē. To uzbūvēja 35 metru augstumā 1934.--1937.~gadā no zemes, atvestas no kaujas laukiem, kur bija cīnījušies poļi.

J.~Pilsudski valdīja ar ``stingru roku'', neatzina nekādu disidentismu. Reizē notika politisko pretinieku vajāšana. Piemēram, 1934.~gadā Berezas Kartuskas (\pltxti{Bereza Kartuska}) apkaimē tika radīta īpaša nometne, kurā tika turētas ``personas, kas rada draudus drošībai un sabiedriskajai kārtībai''. (Turpmāk par to tiks stāstīts sīkāk.) Galvenie ``sanācijas'' režīma pretinieki bija revolucionārie spēki, kā arī partijas, kuras līdz apvērsumam veidoja labēji-centrisko vairākumu Seimā, īpaši~--- nacionāldemokrāti.

Nosacīti autoritārā režīma 13~gadus var sadalīt trijos posmos:

\begin{enumerate}
\item Pārejas, šķietama parlamentārisma posms ilga līdz 1930.~gada novembra Seima vēlēšanām, kurām jau bija raksturīga plašu terora pielietošana;

\item Tālākas patvaļas izvēršanās, cīņas pret politiskajām partijām posms līdz J.~Pilsudska nāvei 1935.gadā, kas sakrita ar pastāvošās sistēmas leģitimāciju ar t.s. ``Aprīļa konstitūciju''.

\item Nākamo posmu raksturoja J.~Pilsudska epigonu mēģinājumi saglabāt radīto režīmu un pārvērst Poliju par lielvaru, kas beidzās ar kara izraisīšanos.
\end{enumerate}

Valdības rīcībai pirmajos gados pēc apvērsuma bija raksturīga mērena stingrība, kas Pasaules ekonomiskās krīzes jeb Lielās depresijas gados pārauga par stingru un agresīvu antiparlamentāru komandēšanas stilu. Rūpes par poļu tautas ``audzināšanu'' pārņēma armija un tāda iestāde kā ``Fiziskās un psihiskās audzināšanas un militārās izglītošanas valsts birojs'' \pltxti{(Państwowy Urząd Wychowania Fizycnego i Przysposobienia Wojskowego}). Pats maršals turpināja sava mīta spodrināšanu. 1927.gadā bijušo leģionāru kongresā viņš sacīja: ``Es esmu izdomājis daudz skaistu vārdu un apzīmējumu, kuri dzīvos arī pēc manas nāves un kuri poļu tautu noved idiotu līmenī''. Kā redzam, pašam viņam nebija ilūziju. Tika ierobežota vārda brīvība~--- masu informācijas līdzekļi tika pakļauti stingrai cenzūrai. Arodbiedrības un opozīcijas laikraksti tika paciesti tik ilgi, kamēr to aktivitātes netika uzskatītas par graujošām.

Taču blakus formāli likumīgiem līdzekļiem pret opozīciju ``sanācijas'' režīms lietoja arī brutālas metodes. Jau naktī no 1926.~gada 30.septembra uz 1.oktobri virsnieku grupa iebruka Seima deputāta ekonomista J.~Dzehovska, kurš aktīvi aizstāvēja valsts izdevumu, tai skaitā armijas vajadzībām, samazināšanu, dzīvoklī un viņu smagi piekāva. Kad. M.~Ratajs šo gadījumu pieminēja sarunā ar J.~Pilsudski, tas atbildēja, ka tā nav viņa darīšana, lai pats J.~Dzehovičs meklē savus pāridarītājus, bet viņš kaut kāda ``neģēļa'' dēļ neļaus turēt aizdomās visu virsniecību. Šādas metodes, kā piespiest ``cienīt'' ``vadoni'' un viņa realizēto politiku, viens no pirmajiem izjuta arī ievērojamais poļu dramaturgs, satīriķis un publicists A.~Novačinskis, kurš arī pēc maija apvērsuma nepārtrauca uzbrukumus maršalam. Par to viņu virsnieki zvērīgi piekāva, izsitot aci. Satīriķis pieklusa, bet J.~Pilsudska piekritēju rokās nonākusī Temīda bez panākumiem ``meklēja'' ``nezināmos noziedzniekus''. Citu opozicionāru žurnālistu un populāru rakstnieku T.~Dolenču-Mostoviču arī ``nezināmi noziedznieki'' izveda uz māla karjeru pie Varšavas, kur smagi piekāva. Viņa uzstāšanās presē pārtrūka, taču sanāciju viņš neiemīlēja, bet gan sāka pret to uzstāties smalkāk, publicējot garstāstus ar zemtekstu, izsmejošu varas iestādes. Par virsnieku vardarbīgu uzbrukumu objektiem kļuva ne mazums režīma kritiķu no politiķu un publicistu vidus, bet nevienā gadījumā noziedznieki tā arī netika atrasti.

Tajos atsevišķajos gadījumos, kad Seims un Senāts pieņēma lēmumus, neatbilstošus valdošo aprindu vēlmēm, izpildvara gandrīz vienmēr atrada veidu, kā rīkoties, ignorējot parlamentārās demokrātijas pamatnormas. Pēc būtības tā bija ``mīksta'' diktatūra pret gatavajiem pakļauties, bet vajadzības gadījumā tā varēja lietot arī represijas. Taču jāuzsver, ka pirmais veiksmīgais uzbrukums demokrātijai izraisīja arvien jaunus tās ierobežojumus. ``Mīkstā'' diktatūra ar laiku kļuva arvien ``cietāka''. Īsā laikā cietumos nonāca vairāk nekā 6 tūkstoši politieslodzīto.

Tomēr pirmajā acu uzmetumā politiskā sistēma pēc apvērsuma mainījās nenozīmīgi. Seima kritiķi J.~Pilsudskim varēja pārmest, ka viņš saglabāja parlamentārās sistēmas netraucētas darbības šķietamību, kad tā prestižs bija nokritis ļoti zemu. J.~Pilsudskis likuma burtu pielāgoja saviem mērķiem, saglabājot konstitucionālisma šķietamību, taču vajadzības gadījumā bija gatavs pielietot arī represijas. Konstitūcijas burts netika pārkāpts, bet tās gars zuda. Vairumā gadījumu likumi formāli tika ievēroti, taču to būtība bija pilnībā sagrozīta. Galvenā loma parlamenta pārvēršanā par ``piekto riteni'' valsts pārvaldē neapšaubāmi piederēja pašam J.~Pilsudskim. Asistēja viņam prezidents, ministri un vairākums deputātu. Jāpiebilst, ka parlamenta arvien pieaugošā diskriminācija nesastapa kaut cik ievērojamu masu, tai skaitā arī demokrātiskās inteliģences, kura pirms 1926.~gada apvērsuma bija diezgan aktīva, protestus.

J.~Pilsudskis, neapšaubāmi, bija pats populārākais cilvēks Polijā. Viņš bija daudz darījis Polijas valsts izveides labā, viens no Polijas valsts atjaunotājiem, dibinājis tās armiju. Krīzes un vilšanās gaisotne sabiedrībā radīja prasību pēc autoritātes. 20.gadu sākumā J.~Pilsudskis paspēja īstajā laikā noiet no politiskās skatuves Tādejādi, saglabājot savu autoritāti un cieņu, viņš formāli nebija ``sasmērējies'' parlamentārās krīzes laikā. Citi politiskie līderi~--- nedz sevi kā ievērojamu teorētiķi pierādījušais R.~Dmovskis, nedz veiklais taktiķis V.~Vitoss krīzes apstākļos nebija parādījuši pietiekamu izlēmību, gatavību uzņemties atbildību par valsts likteni.

J.~Pilsudskis savā labā izmantoja pretrunas poļu sabiedrībā, kuru bija sašķēluši seši ekonomiskās un politiskās krīzes gadi. No vienas puse~--- masas viņu joprojām uzskatīja par kreiso spēku pārstāvi un tāpēc saskatīja viņā glābiņu (kreisie taču iestājās par sociālo taisnīgumu). No otras puses~--- elitei viņš bija vajadzīgs, lai novērstu sociālo revolūciju pēc Krievijas parauga, lai neļautu Polijā nodibināties ``proletariāta diktatūrai''.

Sākotnēji pēc apvērsuma J.~Pilsudskis baudīja kreiso partiju atbalstu. Taču viņu pašu šī savienība arvien vairāk neapmierināja. Labējās un centriskās partijas pat apvainoja sociālistus un ļudoviešus par centieniem uzpūst sabiedrībā radikālus noskaņojumus, izmantot apvērsumu saviem mērķiem, piespiežot J.~Pilsudski pildīt kreiso programmatiskās nostādnes. Taču viņš netaisījās realizēt kreisu programmu, drīz uzsāka manevrus ar labējiem. Jau pirmie J.~Pilsudska diktatūras pastāvēšanas mēneši parādīja, cik iluzoras bija PPS un kreiso ļudoviešu cerības, saistītas ar valsts apvērsumu. Arvien atklātāka režīma sadarbība ar muižniecības un kapitālistiskajām aprindām, atteikšanās no parlamenta atlaišanas un ārkārtas vēlēšanām, zemnieku un strādnieku valdības izveides prasības ignorēšana, atteikšanās no radikālas agrārās reformas virzīja šīs partijas opozīcijā autoritārajam režīmam. Tomēr atšķirībā no komunistiem (PKP) kreisās nekomunistiskās partijas vienmēr uzsvēra savu apņēmību darboties likumības ietvaros.

\strong{Sociālisti (PPS)} aktīvi protestēja pret parlamenta tiesību ierobežošanu. 1926.~gada nogalē PPS pieņēma lēmumu par pāreju opozīcijā. Taču partija nevis cīnījās par pastāvošā režīma likvidāciju, bet tikai par valdības reorganizāciju, atsevišķu visvairāk sakompromitēto ministru nomaiņu. Viens no PPS līderiem T.~Arciševskis skaidrojot, kāpēc partija uzstājas pret diktatūru kā pārvaldes formu, paziņoja: ``Pastāvēja briesmas, ka strādnieku vidū iesakņosies centieni nodibināt diktatūru un viņi teiks: ja jau ir buržuāzijas diktatūra, lai tad ir arī proletariāta diktatūra. Pret to mēs arī uzstājāmies.''

Arī \strong{ļudovieši} vēlējās darboties tikai likumīgi. Viens no viņu līderiem M.~Ratajs teica: Mēs ilgu gadu garumā centāmies iekalt zemnieku galvās, ka tikai likumīgā ceļā viņi varēs kaut ko sasniegt''. Šai laikā aktīvākā ļudoviešu partija \pltxti{Wyzwolenie} bija samērā kreisa, PPS sabiedrotā, distancējās no ``sanācijas'' un pārgāja tai opozīcijā. Tāpat pārējie kreisi noskaņotie ļudovieši, kuri līdz 1926.~gadam centās panākt tādas pozīcijas valdībā, kuras ``atbilstu'' zemniecības īpatsvaram sabiedrībā un ļautu nodibināt ``zemnieku'' varu, jutās vīlušies. Apvērsums liedza viņiem tādu iespēju. Arī partijā ``\pltxti{Piast}'', kurai nepārtraukti uzbruka gan kreisās partijas, gan J.~Pilsudska piekritēji, parādījās tendence atiet no sadarbības ar endekiem un nodibināt ciešākus kontaktus ar centriskajām partijām un arī kreisajiem ļudoviešiem. Tiesa, ļudoviešiem vēl bija jāpārdzīvo diezgan sāpīgs savu rindu attīrīšanas process no cilvēkiem, kuri vēlējās sadarboties ar sanācijas spēku nometni. Ļudoviešu partijās notika šķelšanās, J.~Pilsudska piekritēji izstājās no to rindām.

Apvērsums ienesa sajukumu kreiso radikāļu~--- \strong{komunistu} rindās. Polijas Komunistiskās partijas vadība ar Ā.~Varski priekšgalā 1926.~gada pavasarī J.~Pilsudski uzskatīja par ``demokrātiskās un radikālās sīkburžuāzijas vēlmju izteicēju'', pretstatot viņu labējiem, kuri pārstāvot ``šļahtu un buržuāziju'' un gatavojuši labējas diktatūras nodibināšanu. Viņš kvalificēja apvērsumu kā sīkburžuāzijas dumpi, kas dodot triecienu lielkapitālam un muižniecībai, atbalstīja to, uzskatot par ``ceļu uz revolucionāri-demokrātisku diktatūru''. Varšavā daļa komunistu pat piedalījās ielu sadursmēs pret V.~Vitosa valdības piekritējiem. Viņu šie komunisti uzskatīja par fašistisku. Taču jau drīzumā Kominterne poļu Kompartijas vadības nostāju atzina par kļūdainu.

Jau 1926.~gada 22.~maijā Kominternes Izpildu Komitejas sēdē poļu komunistu lozungu par ``Pilsudska nometnes'', ``Pilsudska revolucionārā karaspēka'' atbalstīšanu raksturoja kā izrietošu no ``situācijas absolūti nepareiza politiskā vērtējuma''. 1926.~gada 27.~maijā VK(b)P CK Politbirojs pieņēma lēmumu uzskatīt balsošanu par J.~Pilsudska kandidatūru prezidenta amatam par ``nepieļaujamu nekādos apstākļos'', ieteikt poļu komunistiem kopā ar nacionālo mazākumu pārstāvjiem izvirzīt savu kandidātu. 31.~maijā Politbirojs apstiprināja vēstuli Polijas Komunistiskās partijas CK, kurā komunistu balsošanu parlamentā par labu J.~Pilsudskim atzina par ``noziegumu'', par ``rupju politisku kļūdu ar oportūnistisku raksturu'', un informēja par KIIK vienbalsīgi pieņemto lēmumu ``nebalsot par J.~Pilsudski''. PKP ceturtajā kongresā 1927.~gadā partijā notika šķelšanās ``labējos''~--- vairākumā ar Ā.~Varski un ``kreisajos''~--- mazākumā ar J.~Ļenski priekšgalā. ``Kreisie'' apgalvoja, ka J.~Pilsudska apvērsums bija fašistisks, ``labējie''~--- ka tas ievadīja militāras diktatūras sākumu, kurai ir tendence pāraugt fašistiskā. Abu frakciju cīņa norisa līdz 20.gadu beigām. Tā ieguva tik asu raksturu, ka KIIK atkal un atkal bija spiesta atgriezties pie šī jautājuma, mēģinot pārtraukt frakciju cīņu. Uz laiku darbā PKP CK tika nosūtīti VK(b)P Ck pārstāvji, tai skaitā arī Vilhelms~Knoriņš. 1930.~gadā PKP piektais kongress paziņoja par drīzas revolūcijas tuvošanos. Šai pat kongresā sociālisti tika raksturoti kā fašistiska partija, ko no mūsdienu viedokļa var uzskatīt par acīm redzamu aplamību. Arī citi legālie, parlamentā pārstāvētie opozicionārie sanācijai politiskie grupējumi, neraugoties uz to konfliktiem ar pastāvošo režīmu, tika vērtēti kā fašistiski vai profašistiski. Tas arī PTR pastāvēšanas laikā poļu vēsturniekam J.~Kovaļskim ļāva runāt par sektantisma tendencēm poļu komunistu vidū.

Partija tomēr spēja atzīt savu ``maija kļūdu'' un uzstāties pret sanācijas režīmu, organizēt streikus (Lodzā vispārēji streiki notika 1928., 1933., 1936.~gadā). Neraugoties uz represijām (1927.~gadā tika aizliegtas ar PKP saistītās legālās kreisās organizācijas, arī 1924.~gadā izveidojusies Neatkarīgā Zemnieku partija (\pltxti{Niezależna Partia Chłopska})), komunistiem izdevās saglabāt savu pārstāvniecību (pāris deputātu) Seimā līdz pat 1935.~gada vēlēšanām. Tika radīti daudzi legāli kreisi preses izdevumi. Partijas centrālais orgāns bija nelegālā avīze ``\pltxti{Czerwony Sztandar}'' (``Sarkanais Karogs'') un žurnāls ``\pltxti{Nowy przegląd}'' (``Jaunais Apskats''). 20.~un 30.~gadu mijā frakciju cīņa partijā sāka pierimt. 1929.~gadā par tās ģenerālsekretāru kļuva J.~Ļenskis. 30.~gados partijā bija ap 20~000 biedru. Tiesa, daudzi atradās ieslodzījumā. 1930.~gadā valstī arestēja 3~775~partijas biedrus, 1931.~gadā~--- 3~507, bet 1932.~gadā~--- 6~982. PKP turpināja kritizēt sociālistus un ļudoviešus. Uz 1934.~gada 1.~maiju publicētājā PKP CK uzsaukumā skanēja aicinājums uz masu revolucionāro cīņu par Polijas padomju republiku, kā arī teikts: ``PSP un \pltxti{Stronnictwo Ludowe} arī runā ``par strādnieku-zemnieku'' vai ``zemnieku-strādnieku'' valdību, taču cenšas mūs pārliecināt mīlēt mūsdienu buržuāzisko valsti, bet strādniekam un darba zemniekam mīlēt buržuāzisko valsti nozīmē mīlēt pātagu, ar kuru viņš tiek šaustīts.''

Tikai 30.gadu vidū PKP taktika mainījās. Atbilstoši Kominternes VII kongresa lēmumiem tā iestājās par komunistu, sociālistu un citu spēku apvienošanā cīņā pret fašismu (kā jau teikts, kaut valstī pastāvošo režīmu pēc būtības nevarēja saukt par fašistisku, visus 30.gadus komunistiskā propaganda J.~Pilsudska nodibināto režīmu sauca par fašistisku).

\strong{Endeki} neizrādīja īpašu neapmierinātību ar parlamentārisma ierobežošanu, (Pēc dažiem vērtējumiem, to viņi arī nevarēja izdarīt, jo balstījās nevis uz tautas masām, bet plānu poļu buržuāzijas slāni, kurš bruņotā sadursmē varēja vienīgi zaudēt), taču palika maršala pretinieku pozīcijās. Faktiski jaunā politiskā režīma un labējo sabiedrisko spēku mērķi sakrita. Taču to sasniegšanas līdzekļi atšķīrās. Pilsudskieši arī nevēlējās dalīties varā ar uzvarēto t.s. ``nacionālo nometni''. Endeki bija neapmierināti ar savu atstumšanu no varas. Par viņu galveno līdzekli cīņā par varu kļuva vardarbības atzīšana: apvērsums varas saņemšanai, terors pret politiskajiem pretiniekiem. Viņu līderis R.~Dmovskis atrada daudz pamācoša un viņa partijai derīga totalitārajos režīmos fašistiskajā Itālijā un (pēc 1933.~gada) arī nacistiskajā Vācijā. Viņa piekritēju vidē ``nacionālajā nometnē'' tika iedibināta hierarhiska vadība un stingra disciplīna, paplašinātas militarizētas vienības. 1926.~gada decembrī viņi, cenšoties apvienot visus labējos ``sanācijas'' pretiniekus, radīja organizāciju ``Dižās Polijas nometne'' (\pltxti{Obόz Wielkiej Polski}), kura pastāvēja līdz 1933.~gadam, kad ``sanācijas'' režīms to aizliedza. Kad tika redzamas J.~Pilsudska režīma atsevišķas neveiksmes, endeku balsīs bieži skanēja ļauns prieks: ``Vai tad mēs jums neteicām, ka ir jāklausa mūsu partijai, nevis J.~Pilsudskim?'', ``Mēs būtu visu izdarījuši citādi''.

Lai atņemtu saviem labējiem oponentiem sociālo bāzi, valdošais grupējums sāka radīt jaunu sabiedrisko spēku koalīciju, kura paustu mantīgo intereses. Līdz apvērsumam J.~Pilsudskim nebija plašas, labi organizētas aizmugures. Viņš varēja rēķināties tikai ar šauru savu veco līdzgaitnieku atbalstu no Leģiona un Poļu kara organizācijas, kura tagad ieņēma ievērojamu vietu armijā un pat kreisajās un centriskajās politiskajās partijās. Ne velti viens no populāriem poļu publicistiem, rakstnieks K.~Prušinskis 1941.~gadā polemiskā aizrautībā izsaucās: ``Vai tad Lodzas rūpnieki, kā uzskatīja vieni, vai Nesvižas (poļu \pltxti{Nieśwież}, baltkrievu \betxti{Нясвіж}, Niaswi; domātas toreizējās Polijas austrumzemes~--- V.Š.) zemes īpašnieki, kā skaidroja citi, pēdējā laikā noteica visu Polijā? Nē, visu vadīja militāristi.'' Rakstnieks gan daļēji pārspīlēja militāristu nozīmi. Ne ``visu'' viņi noteica. J.~Pilsudska atteikšanās no ``riskantiem sociālajiem un ekonomiskajiem eksperimentiem'' nodrošināja viņam sadarbības iespēju ar rūpnieku, finansistu un tirgotāju aprindām. To pārstāvji ieguva ministru portfeļus jau pirmajās pēc apvērsuma izveidotajās valdībās. Tautā, atceroties, ka J.~Pilsudska veidotajā leģionā bija tikai trīs brigādes, runāja, ka šie vīri nāk ``no 4.~un 5.~brigādes''.

Kontakti tika veidoti arī ar tiem kapitāla pārstāvjiem, kuri tradicionāli orientējās uz endekiem. ``Sanāciju'' sākotnēji atbalstīja arī ievērojama daļa ebreju, vācu un ukraiņu buržuāzijas grupējumu.

Maršals sabiedrotos meklēja arī konservatīvo aprindās, kuri gan kopš parlamentārisma ieviešanas Polijā 1919.~gadā politikā spēlēja visai maznozīmīgu lomu. Taču tā kā viņš naidojās ar endekiem, savu konservatīvo uzskatu pierādīšanai viņam nācās ``draudzēties'' ar šiem labēji-konservatīvajiem spēkiem. Neformāla savienība tika noslēgta bijušajam sociālistam J.~Pilsudskim, kurš gan nekad nenoliedza savu šļahtisko izcelsmi, 1926.~gada oktobrī kņazu Radzivillu (lietuviešu \lttxti{Radvilos}, poļu \pltxti{Radziwiłłowie}) muižā Ņesvižā (\pltxti{Nieswiez)} tiekoties ar vairāku poļu magnātu ģimeņu pārstāvjiem.

Sākotnēji ``sanācijas'' piekritējiem par labu nāca zināmi ekonomiskie panākumi, kurus izsauca starptautiskās konjunktūras uzlabošanās. Salīdzinoši labā 1926.--1928.~gadu ekonomiskā situācija noveda pie pagaidu iekšējas stabilizācijas. Izmantojot samērā labvēlīgo ekonomisko konjunktūru, kuru sanācijas darbinieki iztēloja par savu sasniegumu un minēto pagaidu politisko stabilizāciju, uz 1928.~gada martu tika nozīmētas nākamās parlamenta vēlēšanas.

Pirms tām pilsudskieši organizēja \strong{Bezpartijisko bloku sadarbībai ar valdību} (\pltxti{Bezpartyjny Blok Wspołpracy z Rządem}). Teorētiski bloka izveide tika pamatota kā J.~Pilsudska idejas par sabiedrības vienību ``augstāko, pāri šķirām stāvošo interešu vārdā'' īstenošanas sākums. Tā sanācijas lozungs tika papildināts ar valsts pārākuma ideju pār kolektīvajām un individuālajām interesēm. Kalpošana valstij tika deklarēta kā svarīgāka par partiju interesēm, bet faktiski šis bloks kalpoja kā atbalsts J.~Pilsudskim. Tā uzdevums bija nevis iesaistīt politikā tautu, bet gan pasargāt režīmu no tam antagonistiskas sociālas un ideoloģiskas ietekmes, nodrošināt režīma stabilitāti. Šīs idejas apvienoja tos konservatīvos, nacionālistus, ļudoviešus, sociālistus, sindikālistus, kuri sekoja J.~Pilsudska kursam. Pēc personīga J.~Pilsudska rīkojuma vēlēšanu vešanai no valsts kases blokam tika iedalīti 8~miljoni zlotu. Viss valsts administratīvais aparāts kalpoja bloka vēlēšanu cīņas vešanai.

Vieni (arī Latvijā 1929.~gadā iznākušas grāmatas par Poliju autori K.~Bürgers un J.~Ošs) bloku, neraugoties uz nosaukumu, uzskata par J.~Pilsudska piekritēju uz vēlēšanām organizētu partiju. Citi domā, ka tas drīzāk bija vēlēšanu saraksts nekā reāla politiska partija. Jāuzsver, ka bloks bija organizatoriski visai amorfs, principiāli atšķīrās no tādām organizācijām kā nacistu partija Vācijā (\detxti{Nationalsozialistische Deutsche Arbeiterpartei}) vai fašistiskā partija Itālijā (\ittxti{Partito Nazionale Fascista}). Bloka kodolu veidoja J.~Pilsudska piekritēji no leģionāru vidus un pārbēdzēji no bijušajiem kreisajiem, labējiem un centristiem. Vācu autors R.~Vingendorfs kā trīs bloka sastāvdaļas izdalījis:

\begin{enumerate}
\item t.s. ``pulkvežu grupu'' tai laikā ar pulkvedi V.~Slaveku priekšgalā;

\item bijušo sociālistu grupu ar agrāko premjeru J.~Moračevski vadībā;

\item konservatīvo grupu ar kņazu J.~Radzivillu un profesoru A.~Križanovski priekšgalā, kurai pieslējās galvenokārt lielie zemes īpašnieki.
\end{enumerate}

Protams, katra grupa sev izdevīgā garā interpretēja kopējo devīzi. Visvairāk t.s. Bezpartijiskajā blokā bija vienkāršu karjeristu un cilvēku, kuriem vispār nebija nekādu savu uzskatu, bet kuri bija gatavi sekot J.~Pilsudska harizmai. Vēlēšanu kampaņu bloka vārdā veda valdība, plaši izmantojot administratīvo aparātu uz vietām. Propaganda tika veikta ar lielu vērienu, ievērojamā mērā finansēta no valsts budžeta.

Pats J.~Pilsudskis, iespējams, cerot uz pārliecinošiem panākumiem vēlēšanās, piekrita balotēties Bezpartijiskā bloka sarakstā kā pirmais kandidāts. Faktiski ar to vēlēšanas tika pārvērstas par savdabīgu plebiscītu, kuram bija jāatspoguļo tautas attieksme pret pastāvošo režīmu un tā vadītāju. ``Sanācijas'' piekritēju uzstāšanās ar savu patsāvīgu Bezpartijiskā bloka programmu galīgi sagrāva kreiso partiju cerības, ka J.~Pilsudskis ar laiku būs spiests griezties pie tām pēc palīdzības un piekritīs pieņemt to noteikumus.

\strong{1928.~gada} vēlēšanas notika ar daudziem starpgadījumiem. Vēlēšanu kampaņas laikā tika pieņemts likums par kara stāvokli, kas deva valdībai tiesības apturēt Konstitūcijā garantētās pilsoņu politiskās tiesības, veikt masu arestus, pielietot bruņotu spēku. Virknē vēlēšanu iecirkņu tika pasludināti par nederīgiem endeku saraksti. Daudz likumu pārkāpumu, viltojumu bija austrumu rajonos, kur varas iestādes negribēja pieļaut nacionālo mazākumu un komunistu panākumus. Opozīcijas darbinieki pārliecinoši pierādīja patvaļu un ļaunprātības vēlēšanu komisiju darbībā par labu Bezpartijiskajam blokam.

Un arī tad Bezpartijiskais bloks neieguva vēlētāju vairākuma balsis, par to it kā nobalsoja tikai 1/4 vēlētāju (Faktiski pat 1/4 vēlētāju diezin vai balsoja par ``sanācijas'' nometnes sarakstiem), tas saņēma tikai 122 no 444~deputātu mandātiem, kļūstot gan par stiprāko no 20~frakcijām Seimā, taču pat bez vienkārša vairākuma (līdz tam trūka 93~mandātu). Salīdzinājumā ar 1922.~gada vēlēšanām tas bija liels J.~Pilsudska piekritēju panākums, taču no viņu tuvāko un perspektīvo uzdevumu viedokļa pareizāk runāt par varas nesēju sakāvi.

Interesanti, ka pirmo reizi ``valdības'' sarakstā bija iekļauti un tika ievēlēti arī daži ebreju deputāti un tādejādi parlamentā ebreji sadalījās valdības piekritējos un opozicionāros, kas vājināja to kopējās pozīcijas.

Smagu sakāvi piedzīvoja nacionāldemokrāti, iegūstot tikai 37~mandātus. Endeki, kuri 1928.~gadā mainīja partijas oficiālo nosaukumu uz ``Nacionālā partija'' (``\pltxti{Stronnictwo Narodowe}''), vairs nebija galvenās briesmas valdībai. Toties ievērojami nostiprinājās ļudovieši (``\pltxti{Piast}'' gan zaudēja mandātus, taču citas ļudoviešu partijas sasniedza salīdzinoši labus rezultātus) un sociālisti. Par vienu mandātu mazāk nekā 1922.~gadā ieguva nacionālie mazākumi. Pēc vēlēšanām PPS, ļudoviešu partijām ``\pltxti{Wyzwolenie}'' un no tās 1926.~gadā atšķēlušamies Zemnieku partija (\pltxti{Stronnictwo Chtopskie}) Seimā bija pat par vienu mandātu vairāk nekā Bezpartijiskajam blokam. Tomēr kreiso parlamentāro partiju bloks ne par soli netuvinājās savam sapnim iekarot Seimā vairākumu un izveidot savu koalīcijas valdību (zemnieku un strādnieku, kā vēlējās ļudovieši, vai strādnieku un zemnieku, kā postulēja sociālisti).

Par radikāli kreisajiem spēkiem balsoja ap 920~000 (citi dati~---vairāk nekā 850~000) vēlētāju (iepriekšējās vēlēšanās tādu bija ap 132~000), kas piespieda varas iestādes atzīt septiņu komunistu ievēlēšanu.

Vēlēšanas atspoguļoja jaunu politisko spēku izvietojumu parlamentā. Praktiski visu virzienu partijas, protams, izņemot Bezpartijisko bloku, atradās opozīcijā režīmam. Endeki gan bija noskaņoti kareivīgi, taču tā bija galvenokārt personāla sāncensība, antidemokrātiskā sistēma kā tāda neizraisīja īpašus iebildumus. Komunisti uzstājās par antidemokrātiskā režīma gāšanu un proletariāta diktatūras nodibināšanu, galvenokārt kritizējot sociālistus un ļudoviešus kā ``fašisma kreiso paveidu'' un tā neļaujot sev nodibināt ar tiem sadarbību. Sociālisti un ļudovieši cerēja parlamentārā ceļā pārveidot pastāvošo režīmu. Sākotnēji viņi pat cerēja izveidot Seimā vairākumu kopā ar Bezpartijisko bloku, taču rūgti vīlās.

Vēlēšanu rezultāti deva pamatīgu triecienu arī maršala autoritātei. Ja J.~Pilsudskis vēlējās kā iepriekš autoritāri valdīt Polijā, viņam vajadzēja piespiest Seimu kapitulēt vai to atlaist. Toreizējais iekšlietu ministrs F.~Slavojs Skladkovskis atmiņās rakstīja: ``Dažas dienas līdz Seima sēžu sākumam mani izsauca pans Maršals, kurš paziņoja sekojošo: ``Es personīgi prezidenta vārdā atklāšu Seimu. Viņiem tikai jāpieņem zvērests un jāievēl priekšsēdētājs. Un vairāk nekā. Ja man mēģinās traucēt, es ievedīšu Seimā policiju\dots{}''.

Konflikta ieganstus nenācās ilgi gaidīt. Jau atklājot Seima pirmo sēdi 1928.~gada 27.~martā komunistu deputāti Ā.~Varskis u.c. ieradušos tajā J.~Pilsudski sagaidīja ar sašutuma saucieniem, skandēja ``Nost fašistisko Pilsudska valdību''. Pie lietas ķērās F.~Slavojs Skladkovskis. Parādījās policija, komunistu deputātus vardarbīgi izveda no zāles. Tas izsauca sašutumu arī citos deputātos, kuri palīdzēja komunistiem aizsargāties. Kārtību ievest izdevās tikai ar nu jau bruņotas policistu papildvienības ierašanos. Vēlāk komunistu deputātiem vienam pēc otra atņēma deputāta imunitātes tiesības un nodeva tiesai. Par Seima priekšsēdētāju sacensībā ar J.~Pilsudska ieteikto K.~Bartelu tika ievēlēts sociālists I.~Dašinskis, kurš gan nedomāja par pretošanos režīmam, bet gan izteica pārliecību, ka ``valdība un Seims harmoniski sadarbosies''. Tomēr J.~Pilsudskis neuzdrošinājās uzreiz atlaist parlamentu, jo vēlēšanas bija demonstrējušas, ka sabiedrībā sanācijas režīmam ir visai vājš atbalsts, un nebija skaidrs vai jaunas vēlēšanas dotu viņa piekritējiem labākus rezultātus nekā iepriekš. Tāpēc maršals uzsāka Seima, tā deputātu apvainošanas kampaņu.

Iejaucās arī negadījums. 1928.~gadā naktī no 17. uz 18.~aprīli J.~Pilsudskis pārcieta insultu, uz divām dienām bija ievietots slimnīcā. Pastāvēja draudi, ka jebkuru mirkli insults var atkārtoties. Ārstēšanās kopumā bija sekmīga, taču pilnībā atjaunot labās rokas kustīgumu tā arī neizdevās. Maršals piemītošajā sev diktatora manierē 1928.~gada 25.~jūnijā atteicās no premjerministra pienākumiem. F.~Slavojs Skladkovskis atcerējās tās dienas valdības sēdi: ``Apsveicinājušies ar panu prezidentu mēs apsēdāmies ap galdu, un Kamandants ņēma vārdu. Kā kabineta vadītājs viņš, griežoties pie pana prezidenta, visu klātesošo ministru vārdā iesniedza demisiju. Taisnības labad jāsaka, ka neviens no mums, ministriem, pat nenojautām par pana Maršala nodomu iesniegt atlūgumu, taču mēs, protams, ar cieņu noklusējām. Pans Maršals atlūgumu pamatoja ar savas veselības stāvokli\dots{}''. (Par premjerministru J.~Pilsudska vietā atkal kļuva K.~Bartels.) Vēl pēc divām dienām J.~Pilsudskis sniedza interviju, kurā apgalvoja, ka atjaunotajai Polijai no parlamenta nekad neesot bijis nekāda labuma, pirmo Satversmes Seimu nosauca par ``politisko prostitūtu Seimu'', nākamos apsūdzēja par tādas atmosfēras radīšanu, kad pat mušas aizmigušas no deputātu runām. No šī laika militāristi~--- J.~Pilsudska piekritēji, ko tautā sauca par ``pulkvežiem'', (izņemot dažus mēnešus, kad valdības priekšā atradās jau agrāk minētais mērenais K.~Bartels) spēlēja galveno lomu valsts politiskā kursa noteikšanā. Ar to maršals uz politiskās skatuves uzveda savus uzticamākos cilvēkus no leģionāru aprindām.

1929.~gada februārī parlamentā izvirzījās jautājums par vēlēšanu mahinācijām. Atklājās, ka valdība vēlēšanu kampaņas laikā bija pārskaitījusi Bezpartijiskajam blokam 8~miljonus zlotu no bez Seima piekrišanas palielinātā valdības rīcības fonda. Seims pieņēma lēmumu finanšu ministru G.~Čehoviču saukt pie konstitucionālās atbildības, kaut ministrs bija rīkojies pēc premjera J.~Pilsudska norādījumiem. G.~Čehovičs 1929.~gada 8.~martā bija spiests atstāt savu posteni. Taču viņa notiesāšanu neizdevās panākt. Lietas izskatīšanas gaitā Valsts tribunāla priekšā kā liecinieks uzstājās arī J.~Pilsudskis, kurš pats uzņēmās visu vainu, reizē atļaujoties rupjus izteikumus pret tribunālu, Seimu un arī Konstitūciju. 1929.~gada jūnijā Valsts tribunāls atteicās tālāk skatīt G.~Čehoviča lietu. Tā tas izvairījās no lēmuma pieņemšanas, kuram, ievērojot, konstitūcijas pārkāpšanu, bija jābūt visai bargam. Taču valdības un opozīcijas attiecības pasliktinājās. Opozīcijai izdevās zināmā mērā atbildēt uz sanācijas piekritēju pārmetumiem par valsts līdzekļu izšķērdēšanu pirmsapvērsuma periodā, parādīt sabiedrībai, ka visi sanācijas apgalvojumi par savu nākšanu pie varas valsts atdzimšanas interesēs ir nekas vairāk kā parasta demagoģija.

Faktiski šie notikumi iezīmēja pāreju pie t.s. ``stingrās rokas'' politikas. Tas kļuva redzams jau 1929.~gada aprīlī ar K.~Bartela vadītās valdības atkāpšanos un tā saucamās \strong{``pulkvežu valdības''} (tajā darbojās pulkveži V.~Slaveks, A.~Pristors, K.~Svitaļskis) izveidi. 1929.~gada 31.~oktobrī visas Seima telpas, izņemot sēžu zāli, kur pulcējās deputāti, ieņēma bruņotu virsnieku grupas. To mērķis~--- deputātu piekaušana un Seima atlaišana vai tikai kārtējā tā noniecināšana un diskreditācija~--- tā arī palicis nezināms, jo, neraugoties uz J.~Pilsudska bezceremoniālo spiedienu, Seima priekšsēdētājs I.~Dašinskis atteicās ``durkļu, revolveru un zobenu draudu'' apstākļos atklāt sēdi, tā tika pārcelta un notika 5.~novembrī. Seims netika padzīts, acīmredzot ``sanācijas'' režīma vadītāji pēdējā brīdī neuzdrošinājās to darīt. Žurnāliste S.~Dubnova-Ērliha tā raksturoja situāciju valstī: ``Nospieda politiskā atmosfēra. Zeme raustījās neatrisinātu problēmu tīklos, bezatbildīgā vadība gāzelēdamās ``soļoja ar labo'', lauku un pilsētu nabadzīgo vidū auga un neatrada izeju bezspēcīgs sašutums''.

Centriskās un kreisās partijas 1929.~gada 14.septembrī noslēdza savienību, kuru nosauca par \strong{``Centroļevu''} (``\pltxti{Centrolew}'') un kura paziņoja, ka valdība tikai traucē sabiedrības ievelētajiem uzticības vīriem veikt soļus krīzes pārvarēšanas virzienā. Viens no ļudoviešu līderiem V.~Vitoss deklarēja: ``Mums ir jāaizstāv parlamentārisma principi''. Savienībai Seima diskusijās bija jādod pretspars ``sanācijas'' politikai, jāpiespiež J.~Pilsudski atlaist ``pulkvežu valdību'' un atgriezties pie mazāk autoritāras politikas. (Blokam Seimā piederēja 183 no 444~vietām, relatīvs vairākums salīdzinājumā ar citiem spēkiem.)

1929.~gada decembrī sakarā ar atklātajiem finanšu disciplīnas pārkāpumiem opozīcijas partijas panāca K.~Svitaļska valdības atkāpšanos. Tam bija ievērojama psiholoģiska nozīme, taču K.~~Svitaļska vietā par valdības vadītāju atkal kļuva K.~Bartels. Tālāks spiediens uz J.~Pilsudski neizdevās, jo Seima sēdes pastāvīgi tika atceltas. 1930.~gada martā J.~Pilsudskis atkal izveidoja jaunu V.~Slaveka valdību, kurai nebija vairākuma Seimā.

Šādā situācijā ``Centroļevs'' nolēma organizēt masu akcijas, demonstrācijas visā valstī. Visbūtiskākā bija t.s. Krakovas kongresa tautas tiesību un brīvību aizstāvībai sasaukšana 1930.~gada 29.~jūnijā. Kongress bija sasaukts likumīgi, to ievadīja plašs organizatorisks un propagandas darbs. Papildus tūkstotim delegātu no vietām, Centroļeva deputātiem un senatoriem Krakovā sapulcējās no 25~līdz 50~tūkstošiem to atbalstītāju. Kongresā J.~Pilsudskis tika raksturots kā diktators, pieprasīta viņa nozīmētās valdības atlaišana. Pieņemtā kongresa noslēguma rezolūcija bija ieturēta diezgan stingrā tonī. Tā pasludināja ``cīņu par tiesībām, brīvību un diktatūras likvidēšanu''. Režīms tika brīdināts pret mēģinājumiem veikt valsts apvērsumu un ar teroru nomākt demokrātijas atjaunošanas kustību, solot visnoteiktāko pretestību, līdz pat fiziskai. Par kritikas objekts kārtējo reizi kļuva prezidents. Opozīcija liedza tiesības uz morālo autoritāti visiem režīma pīlāriem~--- J.~Pilsudskim, valdībai un prezidentam. Pēc kongresa beigām notika masu mītiņš, kurā piedalījās aptuveni 30~000 cilvēku un tika pieņemta līdzīga rezolūcija. Pēc tam mītiņa dalībnieki devās pie Ā.~Mickēviča pieminekļa, ceļā izkliedzot saukļus ``Nost ar diktatūru!'', ``Pilsudski pie karātavām!'', ``Nost ar sulaini Moscicki!'' u.c. Krakovas kongresa panākumus apvienotā opozīcijas uzskatīja par sabiedrības, kuras liela daļa jau izjuta pasaules ekonomiskās krīzes sekas, apliecinājumu gatavībai cīnīties pret ``sanācijas'' režīmu. Taču šis kongress kļuva par ``Centroļeva'' pēdējo panākumu. Opozīcijai īsti nebija skaidrs, kā turpināt cīņu. Poļu komunisti saskatīja bloka darbībā tikai vēlmi ar citiem līdzekļiem nodrošināt ekspluatatoru šķiras varu, tāpēc ieņēma sektantisku pozīciju, kas tos izolēja no blokā apvienotajām masām. Kad pēc vasaras politiskās darbības pieklusuma ``Centoļevs'' atkal sāka gatavot masu akcijas (septembra vidū bija paredzētas līdzīgas manifestācijas vairāk nekā 20~pilsētās), valdība pārtvēra iniciatīvu. Jau 11.~augustā J.~Pilsudskis, atgriezies no atvaļinājuma, lika valdības vadītājam sastādīt ``lietas'' par visiem ``Centroļeva'' deputātiem. 25.~augustā J.~Pilsudskis atkal kļuva par valdības vadītāju, 29.~augustā pierunāja prezidentu I.~Moscicki atlaist Seimu un pasludināt jaunas vēlēšanas, ko nākamajā dienā prezidents arī izdarīja un noteica uz novembri jaunas vēlēšanas.

% doc2, page 498


Ar to tika ievadīts otrais posms ``sanācijas'' režīma vēsturē. 1930.~gada septembrī ar ieganstu, ka ``Centroļevs'' plānojot veikt valsts apvērsumu, ar tiesiskās procedūras pārkāpumiem tika arestēti 18~bijušie deputāti, kuri pēc parlamenta atlaišanas zaudēja imunitāti, to skaitā arī tādi populāri darbinieki kā V.~Vitoss un V.~Korfantu. Arestējamo sarakstu personīgi sastādīja J.~Pilsudskis. Kopā arestēja 88 (dati dažādos izdevumos nedaudz atšķiras) vadošos opozīcijas: sociālistu, ļudoviešu un nacionāldemokrātu politiķus. Viņus aizveda uz Brestas cietoksni, piekāva un sarīkoja šķietamu gatavošanos nāves soda izpildei.

Pēc dažiem datiem oktobrī bija arestēti ap 1~000~cilvēku, starp tiem 64~deputāti. (Pēc citām ziņām kopumā pirms vēlēšanām apcietināja ap 5~000 ``Centroļeva'' dalībnieku, nacionāldemokrātu un nacionālo mazākumu pārstāvju. Par pēdējo vajāšanām skat. nākamo apakšnodaļu.) Vēsturnieks, bijušais Seima deputāts Ā.~Prohņiks pastāstīja, kā vienu no arestētajiem vēl ceļā uz Brestu konvojs ieveda mežā un ar saucieniem ``Tu iedrošinies apvainot Čehoviču, tu iedrošinies pacelt balsi pret panu maršalu'' piekāva līdz bezsamaņai un tādā stāvoklī nogādāja Brestā.

Pats J.~Pilsudskis notikumus, saistītus ar politisko pretinieku ieslodzīšanu Brestā, attaisnoja kā nepieciešamus, citādi valstī ``notiktu briesmu lietas''. Viņš arī sūkstījās: ``Nezinu, kāpēc Dievs man licis dzīvot Polijā\dots{}'', ``Ir lietas, kuras jūs, poļi, neesat spējīgi saprast! \citespace{} Cik gan daudz man izdotos izdarīt, ja es valdītu pār citu tautu''. Pēc opozīcijas sagrāves J.~Pilsudskis atkal atteicās no valdības vadītāja amata, atsaucoties uz veselību. Demonstrācijas, kas atbalstīja opozīciju, tika izklīdinātas, daudzi to dalībnieki arestēti.

1930.~gada 16.~novembrī notika Seima vēlēšanas. Tajās piedalījās aptuveni 75\% vēlētāju. Vēlēšanas vairs ne pēc kādiem kritērijiem nevarēja nosaukt par ``brīvām''. Gandrīz pusmiljons vēlēšanu biļetenu, galvenokārt vojevodistēs ar poļu iedzīvotāju pārsvaru, tika pasludināti par spēkā neesošiem. Plaši izplatīta bija vēlēšanu rezultātu viltošanas prakse. Tautā šīs vēlēšanas pēc analoģijas ar uzvaru 1920.~gadā kaujā par Varšavu saņēma ironisko nosaukumu ``brīnums pie urnas''.

Bezpartijiskais bloks un ar to sadarbojušās mazās partijas ieguva 55,6\% (249 no 444) no vietām Seimā. ``Centroļevs'' ar 102~deputātu vietām tālu atpalika. Visvairāk mandātu zaudēja PPS. Toties pusotras reizes savu mandātu skaitu palielināja nacionāldemokrāti, kuriem ``sanācija'' šai laikā veltīja maz uzmanības. Nacionālie mazākumi spēja izcīnīt tikai 33~mandātus.

Režīms ieguva iespēju Seimā pieņemt tam tīkamus likumus, nerēķinoties ar opozīciju. Taču šādu iznākumu nodrošināja vēl atklātāka nekā 1928.~gadā vēlēšanu rezultātu viltošana, spiediens uz vēlētājiem. Pret to atskanēja daudzu poļu intelektuāļu protesti. Viens no visskaļāko rezonansi ieguvušajiem bija 45~Jagelloņu universitātes profesoru protests. Protestēja arī izcilie literāti M.~Dombrovska, T.~Žeļeņskis, A.~Sloņimskis, J.~Tuvims u.~c. J.~Pilsudska ``baltā'' leģenda sāka būtiski atšķirties no viņa vērtējuma sabiedrībā. Strauji kritās viņa popularitāte.

Partijas \pltxti{Piast} atteikšanās no bloka ar nacionālajiem demokrātiem un kreiso ļudoviešu partiju apzinātas sadarbības ar ``sanāciju'' neiespējamību, kopējā darbošanās ``Centroļeva'' sastāvā, pavēra perspektīvas ļudoviešu kustības apvienošanai. Lielā mērā iespēju apvienoties noteica arī zemnieku vēlme panākt vienotību, kura ļautu tiem aizstāvēties pret lielīpašnieku konkurenci. Ļudoviešiem zuda cerības atgriezties pie varas parlamentāras cīņas ceļā, viņi bija spiesti rēķināties, ka tiem tiks pilnībā atņemtas legālas darbības iespējas, pat viņu uzticība tīri konstitucionālām darbības metodēm tika pakļauta stingram pārbaudījumam.

Jau pēc Seima vēlēšanām 24.~novembrī pret naudas ķīlu tika atbrīvoti pirmie Brestas cietoksnī ieslodzītie opozīcijas politiķi. Pēdējo šīs grupas ieslodzīto atbrīvoja gan tikai 1930.~gada decembra beigās. Sabiedrībai kļuva zināma patiesība par viņu pieredzēto šais mēnešos. No dažādām pusēm izskanēja satrauktus balsis, protesti, prasības sodīt vainīgos. Nacionāldemokrāti ierosināja veikt parlamenta izmeklēšanu par ``Brestas lietu''. Ar interpelāciju pie valdības vērsās ``Centroļeva''’ partijas, sīki aprakstot cietumsargu attieksmi pret aizturētajiem. Vardarbību pret ieslodzītajiem Brestā aktīvi nosodīja universitāšu profesori, radošo profesiju cilvēki, ārsti un juristi, arodbiedrības un sabiedriskās organizācijas. Atkal dzīve bija pieradījusi oficiālās mitoloģijas par valdošo režīmu kā augstas morāles paraugu neatbilstību patiesībai.

1931.~gada rudenī sākās un 1932.~gada janvārī noslēdzās ``Centroļeva'' līderu tiesas process. Brestas cietoksnī ieslodzītie 11~politiķi tika apsūdzēti par nolūku vardarbīgi gāzt likumīgo valdību. Taču apsūdzētie pārvērtās par apsūdzētājiem, atgādinot poļiem par visiem Konstitūcijas un likumības pārkāpumiem, ko bija izdarījis valdošais režīms. Apsūdzības pārstāvji bija spiesti atzīt, ka apsūdzētie necentās mainīt valsts iekārtas pamatus. Opozīcija cīnījās par valdības maiņu legāliem līdzekļiem. Arī tiesa atzina, ka nav iegansta, lai apšaubītu tiesājamo lojalitāti tēvzemei, aprobežojās ar sīkiem, bieži nosacītiem sodiem. V.~Vitosam tika piespriesti 3~gadi cietumsoda. Process demonstrēja, ka Polijas tiesu sistēmu kontrolē politiskais režīms, deva papildus impulsu ``sanācijas'' pretinieku vēlmei mainīt lietu stāvokli valstī. Taču grupa opozīcijas līderu ar V.~Vitosu priekšgalā bija spiesti emigrēt uz ārzemēm (visbiežāk demokrātisko Čehoslovākiju) pēc tam, kad stājās spēkā tiesas lēmums par viņiem piemērojamo cietumsodu. Tā atkal radās poļu politiska emigrācija. Pats V.~Vitoss palika emigrācijā no 1933.~līdz 1939.~gadam.

Meklējot iespējas sekmīgāk cīnīties pret sanācijas režīmu, 1931.~gada 28.~februārī apvienojās trīs galvenās ļudoviešu partijas: \pltxti{Polskie Stronnictwo Ludowe~--- Piast} un \pltxti{Polskie Stronnictwo Ludowe~--- Nowe Wyzwolenie} un Zemnieku partija (\pltxti{Stronnictwo Chłopske}), izveidojot vienu Zemnieku partiju \pltxti{Stronnictwo Ludowe}, kura iestājās par parlamentārisma atjaunošanu, brīvām vēlēšanām, pilsoņu tiesību ievērošanu, tiesu sistēmas un pašvaldību neatkarību, atmeta iebildumus pret privātīpašuma tiesību aizskaršanu, prasīja nekavējoties realizēt radikālu agrāro reformu bez kādas kompensācijas muižniekiem. Partijas Galvenās padomes pieņemtā programma paredzēja valsts tautsaimniecības agrāru raksturu, rūpniecībā galvenokārt bija jāattīsta uzņēmumi, kuri pārstrādātu lauksaimniecības produkciju. Tika precizēts arī partijas raksturs. Tā pozicionējās nevis kā poļu tautas (\pltxti{ļud}-a), pirmkārt zemes īpašnieku, bet gan kā lauku iedzīvotāju, poļu zemnieku partija. Tika atcelts postulāts, ka zemniecība neatkarīgajā Polijā ir ieņēmusi valsts saimnieka pozīcijas.

Partijā gan saglabājās dažādas domstarpības, jo pastāvot nosacītam spēku līdzsvaram, neviena no ļudoviešu partijām nevarēja panākt, lai jaunā organizācija pārņemtu tās nostādnes. Partijā palika arī vēl cilvēki, kuri bija gatavi sadarboties gan ar lielburžuāziju, gan muižniecību, ieskaitot arī ``sanācijas'' režīma atbalstītājus. Taču pakāpeniski 30.~gados ļudoviešu vidē nostiprinājās konsekventa demokrātisma piekritēji, ``sanācijas'' pretinieki.

Partijas centrālais izdevums ``\pltxti{Zielony Sztandar}'' (``Zaļais karogs'') ar rūgtumu rakstīja, ka ``divpadsmit neatkarīgās Polijas pastāvēšanas gadi pilnīgi viennozīmīgi runā par to, ka neviens par zemniecību nerūpējas un uzlabot tās stāvokli negatavojas, tāpēc viņiem pašiem ir jāuzņemas cīņa par atbrīvošanos no nabadzības un tumsības. 1932.~gada beigās izdevums publicēja V.~Vitosa rakstu, kurš uzsvēra, ka ļudoviešu kustība negatavojas cīnīties pret pilsētu, darba inteliģenci un strādniecību, kuru stāvoklis maz atšķiras no zemniecības stāvokļa. Tā iestājas par demokrātijas atjaunošanu, pret sociālo netaisnību un sabiedrisko attiecību ``anarhizācijas'' briesmām. V.~Vitosa raksts pēc būtības saturēja nosodījumu zemniecības separātisma idejai, pretenzijām uz tās īpašo lomu sabiedrībā. Viņu atbalstīja arī daži citi ļudoviešu ideologi. No šī laika ļudoviešu kustībā arvien vairāk izpaudās kapitālisma nolieguma, idejas par sadarbību ar strādniecību un citām (nekomunistiskajām) darbaļaužu grupām, domas par to, ka radikālu agrāro reformu var īstenot tikai tautas pārvaldītā Polijā, kuras pamatu veidos zemniecība. Tagad ļudovieši teorētiski atzina arī neparlamentāro metožu pieļaujamību cīņā pret sanācijas režīmu. 1934.~gada septembra beigās \pltxti{Stronnictwo Ludowe} Galvenā padome izskatīja jautājumu par zemniecības diktatūras kā pārejas no ``sanācijas'' režīma uz parlamentāro demokrātiju sasniegšanas līdzekļa pieļaujamību. Galvenā padome gan šai jautājumā ``principiālu apsvērumu dēļ'' ieņēma negatīvu nostāju. Tomēr diskusija par zemniecības diktatūras kā varas pārejas formas nodibināšanas mērķtiecību, lai nodrošinātu zemniecības noteicošo ietekmi uz valsts politiku, ļudoviešu vidū turpinājās.

1932.--1934.~gadā Stronnictwo Ludowe sastāvēja no 100~līdz 120~tūkstošu (salīdzinājumam~--- 1932.~gada beigās Polijas Komunistiskās partijas rindās bija ap 12~000 biedru), bet 1939.~gadā~--- ap 200~tūkstošu biedru, tās jaunatnes organizācijā~--- ap 90~tūkstošu.

Tomēr jāsaka, ka legālā opozīcija kopumā, radusi cīnīties tikai ar tiesiskām metodēm, nebija gatava sanācijas realizētajam notikumu pavērsienam, vardarbības vilnim no varas iestāžu puses. Vairākums šīs opozīcijas pārstāvju, izjūtot atbildību par valsts likteni, neuzdrošinājās pat domāt par nelegālu cīņas formu izmantošanu pret sanāciju. Toties J.~Pilsudski un viņa sekotājus neatturēja arī neparlamentāras cīņas līdzekļu pielietošanas vajadzība, lai tikai izstumtu no politiskās cīņas arēnas partijas, kuras viņi uzskatīja par bīstamām saviem plāniem.

Tāpat jāatzīmē, ka režīma uzvaru vēlēšanās nodrošināja ne tikai represijas pret opozīciju, tās vājums, bet arī lielas vēlētāju daļas vienaldzīgā attieksme pret ``Centroļeva'' un sanācijas cīņu, kuras nozīmi viņi nesaprata un nesaistīja ar savu personīgo likteni. Sava ietekme bija vēlētāju nodarbinātībai ar ikdienas rūpēm, no agrākiem laikiem mantotajai un režīma propagandas uzturētajai J.~Pilsudska autoritātei.

Daļa sabiedrības simpatizēja ``stipras rokas'' politikai. Plaši iedzīvotāju slāņi vēl pietiekami nenovērtēja ekonomiskās krīzes grūtības, atcerējās 1926.--1929.~gada ekonomiskos panākumus, cerēja uz to turpināšanos.

30.~gadu sākumā no demokrātisma un liberālisma Otrajā Žečpospolitā maz kas bija palicis pāri. Polija ātri pārvērtās par autoritāru nacionālistisku valsti, kurā ērti dzīvot bija tikai turīgiem poļiem, un arī tad~--- tikai tai viņu daļai, kuri atbalstīja oficiālo varu vai arī politikas vietā priekšroku deva personīgajai dzīvei vai radošai darbībai.

Pavisam 1931.~gadā par politisku darbību arestēja 16~743, 1932.~gadā~--- jau 48~829 cilvēkus. Ja 1928.--1930.~gadā pēc nepilniem datiem sadursmēs ar policiju dzīvību zaudēja 187, tad 1932.~gadā jau 367~cilvēki. Šie notikumi radikāli mainīja politisko klimatu Polijā.

Seima tiesības tika arvien vairāk ierobežotas. 1933.~gada 23.~martā pats Seims piešķīra prezidentam tiesības izdot dekrētus ar likuma spēku. Bez sarežģījumiem, kaut opozīcijai demonstratīvi atturoties, I.~Moscickis tika otro reizi ievēlēts par prezidentu.

``Sanācijas'' režīma otrā fāze iezīmēja tā norieta sākumu. Tam nebija ko pretstatīt trūkuma, bezdarba un inflācijas pieaugumam. Auga streiku skaits, arī kreisā un labējā radikālisma piekritēju pulki. No šī laika, kā runāja laikabiedri, Polijas likteni arvien vairāk noteica trīs ``k'':
\begin{enumerate*}[label=\arabic*)]
\item J.~Pilsudska klans, ko pirmkārt veidoja militāristi;
\item J.~Pilsudska kliķe~--- politiķu grupa, kura rīkojās viņa uzdevumā;
\item J.~Pilsudska klaka (no franču \frtxti{claque}~--- plaukstu sasišana, plaukšķināšana~--- cilvēku grupa, salīgta izrādes vai atsevišķa aktiera šķietamu panākumu radīšanai vai otrādi~--- ``izgāšanai'')~--- pakalpiņi, kuri pēc norādījumiem ``no augšas'' tika sarīdīti pret politiskajiem pretiniekiem.
\end{enumerate*}

% page 502
% TODO Example inline enum


Kā atmiņās rakstīja jau minētā žurnāliste S.~Dubnova-Ērliha, valstī tika izvērsts J.~Pilsudska ``personības kults, kurš, protams, nekādi nebija salīdzināms ar divu kaimiņu despotu kultu, kuri izlēja asins jūras, taču tik un tā nevarēja neizsaukt protestus: īpašu sapīkumu izsauca sīka personīga naida izraisītas represijas. Protestu presē nekavējoties apspieda cenzūra''.

Atmiņu autore arī atstāsta gadījumu, kad ģimnāzijās klasē, kurā mācījās viņas dēls, ieradās tās direktors, teica patētisku runu par J.~Pilsudski un piedāvāja skolniekiem ``brīvprātīgi'' parakstīt vēstuli diktatoram, kurā bija izteiktas kvēlas uzticības jūtas. Skolniekus parakstīt vēstuli izsauca alfabēta kārtībā. Kad rinda nonāca līdz S.~Dubnovas-Ērlihas dēlam, viņš paziņoja, ka nevar to darīt, jo nepiekrīt J.~Pilsudska politikai. Pēc dažām dienām zēna vecāki saņēma ģimnāzijas pedagoģiskās padomes paziņojumu, ka viņu dēlam tiek liegta iespēja apmeklēt mācības līdz padome nebūs izskatījusi viņa ``lietu'', ka esot atklāta viņa nelabvēlīga ietekme klasē, kura acīmredzot esot rezultāts viņa audzināšanai nevēlamā garā. Tika izteikts pieņēmums, ka zēna vecāki labāk vēlēsies pārvest viņu uz citu skolu. Vecāki, aprunājušies ar dēlu, tā arī izdarīja.

Sarežģīti procesi norisa opozīcijas rindās.

Pieauga komunistiskās teorijas popularitāte. Komunisti nosodīja kapitālistiskās ražošanas haosu, uzskatīja parlamentāro demokrātiju par tautas krāpšanu. Viņu ietekme ievērojami pieauga, īpaši nacionālo mazākumu vidū. Tāpēc varas iestādes viņus nesaudzīgi vajāja.

Arī sociālisti vismaz vārdos uzskatīja par iespējamu nekavējoties nacionalizēt galvenās tautsaimniecības nozares, pat ieviest pagaidām ``proletariāta diktatūru''.

Pārmaiņas norisa ļudoviešu kustībā, kura izjuta pastāvīgu kā labējo, tā kreiso spēku spiedienu. Tajā notika šķelšanās, to atstāja daudzi pazīstami darbinieki, partijas vadībā izvirzījās jauni spēki, kuri cīnījās par to, lai zemniecība uzņemtos iniciatīvu tādas iekārtas nodibināšanā, kura atbilstu ne tikai tās, bet visu darba cilvēku, arī strādnieku interesēm. Īpaši ļudoviešu jaunatnes organizācija ``\pltxti{Wici}'' (\pltxti{Związek Młodzieży Wiejski}~--- Poļu lauku jaunatnes savienība) vairs nesaistīja zemniecības un valsts nākotni ar kapitālistisko iekārtu. Tās piedāvājamais sabiedrības attīstības modelis~--- t.s. ``trešais'' ceļš atšķīrās kā no pastāvošā kapitālisma, tā arī no PSRS realizētā sabiedriskās attīstības ceļa. Krievu vēsturnieks G.~Matvejevs uzskata, ka organizācijas ``\pltxti{Wici}'' izstrādātais sabiedriskās iekārtas modelis pēc pamatparametriem sakrita ar sabiedrības tautas-demokrātiskās attīstības ceļa koncepciju, kuru revolucionāri-demokrātiskie spēki mēģināja realizēt Vidus- un Dienvid-Austrumeiropā pēc Otrā pasaules kara. (Par to stāstīts grāmatas IV nodaļā). Ja \pltxti{Stronnictwo Ludowe} ideoloģiskās koncepcijas pēc minētā vēsturnieka domām var pieskaitīt buržuāziski reformistiskās ideoloģijai, tad tās jaunatnes organizācijas \pltxti{Wici} piedāvājamo sabiedrības attīstības modeli var uzskatīt par sabiedrības tautas demokrātiskas pārveides koncepciju.

Daļa kristīgo demokrātu simpatizēja korporatīvas valsts radīšanas idejai, kurā sabiedrība, organizēta pēc nozaru principa, pārņemtu valsts funkcijas.

Nacionāldemokrāti sadalījās dažādās grupās, kuras lielākā vai mazākā mērā nonāca pie tradicionālā parlamentārisma noliegšanas. Viņu vidū pārsvarā valdīja ideja par hierarhiskas ``valsts~--- tautas'' radīšanu. Pēc tam, kad valdība atlaida R.~Dmovska radīto ``Dižās Polijas nometni'' (\pltxti{Obóz Wielkiej Polski}, 1926--1933), tika radīta ``Nacionāli radikālā nometne'' (\pltxti{Obòz Narodowo-Radykalny}). Kustība orientējās uz itāļu fašisma paraugu. Tās ideoloģija jau balstījās uz brīvās konkurences, kuru agrāk endeki neapšaubīja, noliegšanu ekonomikā, propagandēja vairāku saimniecības nozaru pārvalstiskošanu. Kustība pieprasīja ārzemju, bet it īpaši ebreju kapitāla ekspropriāciju un nodošanu poļu rokās, arī kultūras un zinātnes pakļaušanu valstij. Taču tās biedru skaits bija tikai pāris tūkstošu. 1934.~gadā kustības kaujinieki, t.s. ``baltkrekli'' mēģināja ``ieviest kārtību'' Varšavas un dažu citu pilsētu strādnieku rajonos, taču dažādās strādnieku organizācijas šoreiz uzstājās vienoti un padzina labējos radikāļus. Par vardarbīgu rīcību ``Nacionāli radikālā nometne'' tika aizliegta un nogāja pagrīdē.

Ievērojot opozīcijas partiju aktivitāti, \strong{``sanācijas'' režīma pasivitāte} jaunu ideju izvirzīšanā bija ļoti uzkrītoša. Parlamentārisma noraidīšana neradīja pamatu kādai jaunai pozitīvai platformai. Režīms krīzi centās pārvarēt ar taupību ekonomikā un represiju pastiprināšanu politikā. Visas ``pulkvežu valdības'' turpināja šo kursu.

Valdības un valsts aparāta normāla darbība bieži cieta no J.~Pilsudska iejaukšanās valsts lietās, kuras viņš turpināja arī tad, kad atteicās no valdības vadītāja amata. Viņš novecoja, izrādīja maz iniciatīvas, taču tai pat laikā traucēja darboties citiem. ``Sanācijas'' politika tika turpināta vairāk pēc inerces. Režīma interesēs bija radīt priekšnoteikumus tā funkcionēšanai arī pēc tam, kad J.~Pilsudskis atstās politisko cīņu lauku. Tuvākie J.~Pilsudska līdzgaitnieki zināja, ka viņa veselība pasliktinās, tāpēc steidzami gatavoja papildinājumus Konstitūcijai, jo vēlējās lai režīms balstītos ne tikai uz maršala personīgo autoritāti, bet arī uz it kā konstitucionālām normām.

Virkne vēsturnieku ir izteikuši savas 1935.~gada Konstitūcijas pieņemšanas versijas. Viena no tām skan, ka, būdams jau gados un nevesels (slimības fakts tika rūpīgi slēpts no sabiedrības), J.~Pilsudskis 30.~gadu sākumā bija norūpējies par savas varas varbūtējo mantinieku, taču nevarējis atrast piemērotu kandidatūru. Daudzi spējīgi cilvēki no viņa apkārtnes (Piemēram, nākamais premjerministrs Otrā pasaules kara gados V.~Sikorskis u.c.) bija nonākuši opozīcijā, toties paklausīgie vadonim nelikās pietiekami spējīgi. Beidzot ``Vectēvs'' (\pltxti{Dziadek}), kā tuvākie sauca J.~Pilsudski, nolēmis, ka lai pilnībā ieviestu kārtību valstī, viņam jākoncentrē savās rokā ne tikai faktiskā, bet arī formālā vara un tikai tad jānodod tā pēctecim. Tāpēc arī 1935.~gada Konstitūcija pārvērta Poliju no parlamentāras par prezidentālu republiku. Izmaiņas Konstitūcijā bija paredzētas J.~Pilsudska varas stiprināšanai. Cita versija skan, ka pats J.~Pilsudkis nemaz necentās konstitucionāli nodrošināt varas pāreju cita līdera rokās pēc savas nāves, neinteresējās par konstitūcijas sagatavošanas gaitu, atdodot tās uzraudzību Bezpartejiskā bloka priekšsēdētājam V.~Slavekam.

Lai nu kā, 1931.~gada 29.~aprīlī sanāksmē Belvederā, kurā piedalījās I.~Moscickis, V.~Slaveks, K.~Svitaļskis, A.~Pristors un J.~Beks, J.~Pilsudskis paziņoja ka gadu gaitā pēc valsts apvērsuma viņš būtiski ir mainījis iekšējās attiecības valstī. Tagad, pēc diktatora domām, galvenais ir konstitucionālais jautājums, taču tas ir jārisina saskaņā ar likumu, tāpēc arī notika pirmstermiņa parlamenta vēlēšanas. Gadījumā, ja pat tagad to izdarīt neizdodas, viņš neapstāsies spēka izmantošanas priekšā, bet pabeigs darbu līdz galam.

Faktiski konstitūcijas projektu gatavoja viņa varas mantinieki savās interesēs. Jaunajā konstitūcijas projektā tika atmestas daudzas demokrātiskās normas, kuras saturēja 1921.~gada konstitūcija, tika nostiprināta sabiedrības virsslāņu un to ``vadoņa'' kundzība.

Bezpartejiskajam blokam nebija parlamentā nepieciešamā vairākuma jaunas konstitūcijas pieņemšanai. Taču 1934.~gada 26.~janvārī, izmantojot to, ka opozīcijas deputāti, izpildot agrāk pieņemto lēmumu par jaunas konstitūcijas sagatavošanas boikotēšanu, atstāja sēžu zāli, atstājot tikai divus novērotājus, bloks panāca projekta apstiprināšanu dažu minūšu laikā. Valdība faktiski nelikumīgi panāca konstitūcijas pieņemšanu pāris minūtēs, it kā tas būtu likums par paipalu aizsardzību jeb obligātu mājlopu potēšanu.

Pēc tam gan J.~Pilsudskis lika izvērst šķietamu diskusiju par Konstitūcijas labojumiem. Tie bija jāizskata Senātā, pēc tam Seimā. Senāts to izdarīja 1935.~gada 16.~janvārī, bet 24.~martā labojumi jaunās konstitūcijas projektā bija jānobalso Seimā. Neraugoties uz visu administratīvo resursu izmantošanu, tie nesaņēma 2/3 balsu (1921.~gada Konstitūcijas gars it kā to prasīja, kaut pamatlikumā nekas nebija teikts par konstitucionālo labojumu pieņemšanas kārtību). Tad Seima priekšsēdētājs paziņoja, ka ``praktisku apsvērumu dēļ'' konstitūcija skaitās pieņemta ar vienkāršu balsu vairākumu.

1935.~gada 23.~aprīlī, skanot 101~artilērijas zalvei, prezidents I.~Moscickis parakstīja jaunās, t.s. \strong{Aprīļa Konstitūcijas} tekstu. 1935.~gada Aprīļa konstitūcija būtiski atšķīrās no 1921.~gada Pamatlikuma. Priekšplānā tautas vietā tajā tika izvirzīta valsts, kurā vadoša vieta piederēja prezidentam. Viņš savās rokās apvienoja ``vienotu un nedalāmu valsts varu''. Konstitūcija saglabāja parlamentu, taču izpildvaru lielā mērā atbrīvoja no atbildības tā priekšā. Izpildvaru pirmkārt pārstāvēja prezidents, viņš arī formēja valdību, nozīmēja 1/3 senatoru. Prezidentu ievēlēja Seima un Senāta izvirzīta 80~cilvēku liela vēlētāju kolēģija, par valsts likteni viņš bija atbildīgs tikai ``Dieva un vēstures'' priekšā, varēja arī pats izvirzīt varas pārmantotāja kandidatūru. Ja vēlētāju kolēģija izvirzītu citu kandidātu, bija jānotiek vispārējai nobalsošanai tikai par šiem diviem kandidātiem. Kara gadījumā pat šāda procedūra nebija nepieciešama. Faktiski varas dalīšanas vietā tika ieviests prezidenta ietekmes pārsvars pār visām citām valsts iestādēm. PKP CK locekļa kandidāts J.~Bruns 1937.~gadā publicētā brošūrā rakstīja: ``Mūsdienu Polija nav republika ar plašu prezidenta varu, kā tas rakstīts sanācijas konstitūcijā, bet ir maršala mūža dzimtmuiža''. Tā bija izteikti prezidentāla pārvaldes sistēma, savienota ar cilvēku grupas diktatūru, kas bija saistīti ar kopīgu pagātni (PPS kaujas organizācija, leģions, 1919.--1920.~gada poļu-padomju karš) un arī kopīgu uzdevumu saglabāt varu kā tagadnē, tā nākotnē.

Pilsoņu tiesības, kuras bija ietvertas Aprīļa konstitūcijā, neietekmēja administratīvo un tiesu iestāžu darbību, praksē tās ieguva arvien vairāk fiktīvu raksturu. Visi vēlētāji tika iedalīti divās kategorijās: parastos un ``labākajos dēlos'', kuriem bija plašākas tiesības.

Ar 1935.~gada Aprīļa konstitūcijas ieviešanu tika pabeigta pāreja no ``sanācijas'' režīma agrīnā modeļa uz tā nobriedušu formu. Ja līdz tam iedibinātais varas modelis varēja darboties tikai balstoties uz J.~Pilsudski ar viņa autoritāti, leģendu un spēju ietekmēt sabiedrību un opozīciju, tad nobriedušai formai, pēc diktatora domām, bija jājūtas pietiekami komfortabli arī bez šāda ``sanācijai'' labvēlīga faktora.

Tiesa, kā vērtē mūsdienu poļu politologs V.~Kuleša, Aprīļa Konstitūcija bijusi ``oriģināls mēģinājums radīt jaunas tiesību attiecības, kuras vienādā mērā atmeta kā liberālismam un parlamentārajai demokrātijai raksturīgos principus un institūtus, tā arī toreizējo totalitāro valstu idejas''. Tomēr nākas secināt, ka 1926.~gada maija apvērsums, J.~Pilsudska autoritārā režīma nodibināšanās sākotnēji veda pie parlamentārās demokrātijas būtiskas ierobežošanas, bet pēc tam arī pie tās konstitucionālo pamatu likvidēšanas. Revolucionārie spēki Polijā neatzina 1935.~gada Konstitūciju ideoloģisku (konstitūcijas totalitārais raksturs), politisku (naidīgo centriski--labējo spēku varas apstiprinājums), tiesisku (konstitucionālo normu un Seima reglamenta pārkāpšana pieņemšanas gaitā) cēloņu dēļ.

Tomēr Aprīļa Konstitūcijai, neraugoties uz savu nedemokrātismu, no juridiskā viedokļa apstrīdamo, pēc būtības~--- nelikumīgo, pieņemšanu, bija sava nozīme Polijas valstiskuma saglabāšanā Otrā pasaules kara laikā. Šī Konstitūcija deva iespēju oficiāli saglabāt Polijas valstiskumu pēc valsts okupācijas, funkcionēt valdībai emigrācijā. Pēdējā, formāli pastāvot līdz pat 1990.~gadam, balstījās tieši uz Aprīļa Konstitūciju. Tas arī mūsdienu poļu politiķiem ļauj runāt par pārmantojamības (nepārtrauktības) principu starp Otro un Trešo Žečpospolitu.

Pēdējos dzīves gados J.~Pilsudskis bija norūpējies arī par sava tēla saglabāšanos tautas atmiņā. Viņš izsauca savu līdzgājēju vēl no Pirmā pasaules kara gadiem A.~Sļivinski, kurš profesionāli nodarbojās ar vēsturi, un palūdza uzrakstīt savu biogrāfiju, jo ar izbrīnu lasot, ko par viņu rakstot dažādi ļaudis. ``Tāda Pilsudska, kādu stāda priekšā kā mani pielūdzēji, tā arī mani pretinieki, nav un nekad arī nav bijis. Es gribētu, lai kaut kas no patiesības nonāktu līdz maniem pēctečiem''. Par vienu no biogrāfijas avotiem, bez citiem, bija jākļūst arī paša maršala atmiņu stāstījumiem. Diemžēl notika tikai trīs maršala sarunas ar biogrāfu. Pēc diktatora nāves A.~Sļivinskis publicēja tikai atsevišķus izvilkumus no sarunu pierakstiem, kuri dod priekšstatu par maršala vērtējumiem gan sev, gan poļu tautai, kuru viņam nācās vadīt starpkaru periodā, jo pēdējie, pēc publicētāja viedokļa bija tik ``briesmīgi'', ka viņš atteicās tos darīt zināmus atklātībai.

1935.~gada maijā J.~Pilsudska veselības stāvoklis ievērojami pasliktinājās. Viņu mocīja aknu vēzis. Pēdējās savas dienas viņš pavadīja invalīdu ratiņos. Tajos arī uzrakstīja savus pēdējos vēlējumus. Tie sākās vārdiem: ``Nezinu, vai gribēs mani apbedīt Vavelā. Lai jau!'' Tālāk viņš izteica vēlējumu, lai viņa sirds tiktu apglabāta Viļņā karavīru kapos, un turp tiktu pārvests viņa mātes zārks, kurš pirms tam bija apglabāts toreizējās Lietuvas robežās. Vēl 10.~maijā, uz brīdi atguvis samaņu, J.~Pilsudskis centās dot pēdējos norādījumus: ``Lavals [1935.~gada maijā Francijas ārlietu ministrs~--- V.Š.] \dots{} man vajag \dots{} Krievija''. 1935.~gada 12.~maijā, desmitajā valsts apvērsuma gadadienā 67~gadu vecumā, pirms vēl jaunā konstitūcija bija stājusies spēkā, \strong{J.~Pilsudskis nomira}. Viņa pēdējā griba tika izpildīta, neraugoties uz to, ka daļa katoļu garīdzniecības iebilda pret bijušā kalvinista apbedīšanu katoliskā svētvietā. Apglabāts viņš ir Krakovā Vavelas katedrālē blakus karaļiem, bet viņa sirds~--- kopā ar māti kapsētā Viļņā, kur apbedīti poļu leģionāri, kuri cīnījās par Viļņas polisku ``atbrīvošanu'', zem melna granīta kapakmens ar uzrakstu ``Māte un dēla sirds'', savas smadzenes viņš novēlēja izpētei Stefana Batorija universitātei (\pltxti{Uniwersytet Stefana Batorego}) Viļņā.

Viņa apbedīšanas ceremoniāla mērogus var salīdzināt ar V.~Leņina apbedīšanu PSRS. Pirms apbedīšanas trīs dienas J.~Pilsudska zārks bija novietots Belvederā, no viņa atvadījās tūkstošiem cilvēku. Godu maršalam atdeva arī ārzemēs. Francijā, Vācijā un PSRS par viņu tika izsludinātas oficiālas sēras. Ā.~Hitlers personīgi piedalījās aizlūgumā Sv.~Jadvigas baznīcā (\detxti{Sankt-Hedwigs-Kathedrale}) Berlīnē.

Pēc diktatora nāves viņa piekritēji turpināja izvērst viņa personības kultu. Sakarā ar to, ka 1938.~gada janvārī kādā Viļņas avīzē bija publicēts cenzūras nepamanīts raksts, kurā par mirušo diktatoru tika runāts kā par ākstu, bet Polijas armijas inspektors Viļņas apgabalā lika sev padotajiem virsniekiem piekaut žurnālistus, kas arī tika izdarīts, šī lieta nonāca Seimā. 1938.~gada 7.~aprīlī tas pieņēma speciālu likumu par J.~Pilsudska labā vārda aizstāvēšanu, kas paredzēja ieslodzījumu līdz 5~gadiem. Tas bija vienīgais normatīvais akts Polijas valsts vēsturē, sargājošs konkrētas personas reputāciju, un kļuva par vienu no jau mirušā diktatora kulta sastāvdaļām. Valdošais režīms gan paguva piemērot jauno likumu tikai reizi. Kad 1939.~gada jūnijā viena no lielākajām kristīgi-demokrātiskajām avīzēm ``\pltxti{Polonia}'' (``Polija'') Katovicē publicēja rakstu ``\latxti{Paderewski~--- Pater Patriae}'' (``Paderevskis~--- Tēvzemes tēvs'') tas izsauca valdošo slāņu sašutumu, jo neviens cits kā tikai J.~Pilsudskis varēja tikt tā godināts, un nekādā gadījumā I.~Paderevskis, kurš pēdējos gados bija nostājies pret maršalu. Nākamo likuma pielietošanu~--- jau par anekdotes stāstīšanu par maršalu~--- izjauca Otrais pasaules karš.

Taču pēc J.~Pilsudska nāves ar laiku pat daudzi viņa piekritēji sāka pārskatīt savus uzskatus. Rakstnieks J.~Tuvims 1944.~gadā kādā vēstulē rakstīja: ``Jo vairāk domāju par tā perioda, kurā tieši vai aizklāti dominēja Pilsudskis, uzskatiem un garu, jo vairāk mani pārņem izbrīns, ka mēs, laikabiedri, ļāvām sevi piemānīt un saindēties ar ``romantisko fluīdu'', ko izstaroja Pilsudskis. Godīgi atzīstos, ka arī es ilgi piederēju pie saindētajiem. Tam bija daudz cēloņu, un sevišķi, pārsteigums par neatkarību, emocionāla maldīšanās daudzsejainajā vēsturiskajā liekulībā, kuru mana paaudze pieņēma par tīru patiesību, briesmīgs poēzijas un ``misticisma'' balasts, kuru mēs gluži velti bijām paņēmuši no XIX gadsimta tā vietā, lai atbrīvotos no tā, tas ir~--- atrastu tam vietu tur, kur tam būtu jābūt: māksliniecisko, literāro vērtību sfērā, bet nevis, kā tas bija pie mums, pārliecinātu sevi un citus,~--- ka mēs vēsturē esam kas īpašs, ka mums ir kāda tur ``vēsturiska misija'', ``sūtība'' un tamlīdzīgas blēņas.'' Cits poļu publicists E.~Novakovskis darbā ``Maršala mīts'' rakstīja: ``Poļi, tik vētraini dzīvojoši no vēsturiskiem mītiem, nemīl, kad nomazgā zeltījumu pat tad, kad, notīrot to, stāv pilnīgi objektīvās pozīcijās.''

Pēc krievu emigrācijas rakstnieka un publicista M.~Aldanova vērtējuma J.~Pilsudskis dzīvoja ``Polijai, karam, slavai'', bet viņa dzīve rāda, cik ``maz vietas mūsdienu politikā ieņem principi un cik daudz vietas tajā ieņem kaislības''.

J.~Pilsudskis (un arī viņa sekotāji) par viņa radītā režīma uzdevumu uzskatīja poļu tautas audzināšanu atbilstoši saviem priekšstatiem par to, kādam jābūt savas valsts patriotam, taču pēc J.~Pilsudska biogrāfa krievu vēsturnieka G.~Matvejeva vērtējuma, viņš iztika``\dots{}bez atklātas vardarbības pār to. Pēckara Polijas valdītāji līdz 80.~gadu beigām tāpat nepieļāva sabiedrības reālu kontroli pār savu darbību, taču pie tam vēl centās to salauzt pār celi, izskaust no tās ``brīvības instinktu'', ko maršals neatļāvās darīt. Tas arī ir viens no cēloņiem, kāpēc Pilsudska piemiņa sabiedrībā dzīvo vēl līdz šai dienai''. Tiešām, mūsdienās daudzi poļi J.~Pilsudski patur piemiņā kā Polijas valsts atjaunotāju un vadītāju. Arī Latvijā Daugavpilī J.~Pilsudska vārdā nosaukta poļu vidusskola, un pie ēkas, kur viņš 1920.~gada janvārī tikās ar ģenerāli J.~Balodi, uzstādīta piemiņas plāksne.

No otras puses, 9~gadu ilgā J.~Pilsudska vienpersonīgās valdīšanas pieredze Polijā ir vēl viens apliecinājums tam, ka pat tās diktatūras, kuras tiek nodibinātas ne tāpēc, lai apmierinātu kāda līdera ambīcijas, bet ar augstiem valstiskiem un sabiedriskiem mērķiem, gala rezultātā pretēji to radītāju nodomiem, nes vairāk ļaunuma nekā labuma tautām. Tās izjauc neatkarīgi domājošiem, talantīgu, ietekmīgu politiķu un pārvaldnieku dabisko atlasi varas elitē. Tā rezultātā pēc diktatora nāves vara paliek to cilvēku rokās, kuri veidoja viņam tuvāko, t.s. iekšējo apli, bet nav raduši domāt un pieņemt atbildīgus lēmumus patstāvīgi. Šajā ziņā Polija nebija izņēmums. Ir zināms, ka visās valstīs ar nedemokrātiskām valdības formām diktatora nāve kļūst par liktenīgu notikumu. Arī Polijā ``sanācija'' ēra tiek dalīta divos periodos: J.~Pilsudska dzīves laikā un pēc viņa nāves.

J.~Pilsudskis nomira, atstājot pie varas politiķu kliķi, kas sastāvēja no vairākām grupām, kuras vienoja padevība maršalam, bet ne kāda aptveroša pamatideja. Pēc J.~Pilsudska nāves ``sanācijas'' režīms turpināja pastāvēt, taču valdības nometnei sava vadoņa nāve bija pamatīgs satricinājums, sākās trešais, faktiska norieta posms tā vēsture.

J.~Pilsudska piekritēji sašķēlās trīs grupās:

\begin{enumerate}
\item Kreisajā sanācijas piekritēju grupā (\pltxti{Lewica sanacyjna}), ar V.~Slaveku priekšgalā, kura centās panākt kompromisu ar opozīciju. To parasti sauca par ``pulkvežu grupu'', jo tajā ietilpa virkne jaunāko J.~Pilsudska sekotāju, kuri bija sasnieguši šo militāro pakāpi.

\item T.s. ``Pils'' (``\pltxti{Zamek}'') grupā, kura apvienojās ap prezidentu I.~Moscicki, kurš dzīvoja Varšavas pilī, no šejienes arī nosaukums.

\item Labējā sanācijas piekritēju grupā (\pltxti{Prawica sanacyjna}), jeb t.s. ``ģenerāļu grupa'', kurā ietilpa E.~Ridz-Smiglija piekritēji, armijas vecākie virsnieki.
\end{enumerate}

Pēc lielas politiķu daļas uzskatiem prezidents I.~Moscickis politiski neatbilda tiem pienākumiem, kuri viņam bija jāpilda pēc konstitūcijas. Tika gaidīts, ka viņš atkāpsies par labu premjerministram V.~Slavekam, vienam no J.~Pilsudska tuvākajiem līdzgaitniekiem, taču tas nenotika. Armijas ģenerālinspektora posteni no J.~Pilsudska mantoja ģenerālis E.~Ridzs-Smiglijs. Taču viņš neieņēma reizē arī kara ministra amatu, tā pozīcijas armijā bija vājākas nekā savā laikā J.~Pilsudskim.

Valdošo slāņu vidū pastāvēja domstarpības par tālāko valsts attīstības ceļu. Vieni uzskatīja par nepieciešamu ieturēt inflācijas politiku eksporta veicināšanai, pastiprināt valsts kapitālisma elementus tautsaimniecībā, meklēt vienošanos ar pilsonisko opozīciju, endekiem, ļudoviešiem u.c. Otri iestājās par deflācijas politiku, iebilda pret tuvināšanās mēģinājumiem ar opozīciju. Robeža starp šiem grupējumiem bija visai izplūdusi, tie šķēlās vairākās sīkākās grupiņās.

Nosacīti izveidojās divi spēcīgi varas centri~--- viens bija prezidents I.~Moscickis, otrs~--- armijas ģenerālinspektors, vēlāk~--- 1936.~gadā maršala pakāpi saņēmušais un jau miera laikā par armijas virspavēlnieku prezidenta ieceltais E.~Ridzs-Smiglijs.

Pēc J.~Pilsudska nāves jau 1935.~gadā pieauga streiku un strādnieku demonstrāciju skaits, kuras dažkārt beidzās ar asins izliešanu. Kreisajā spārnā komunisti aicināja aizstāvēt neatkarību un demokrātiju, radīt Tautas fronti cīņai pret fašismu. Tiem tuvinājās daļa sociālistu, ļudoviešu un arī daļa radikālās inteliģences, kura agrāk atradās J.~Pilsudska ietekmē. Arvien lielāku popularitāti ieguva lozungs ``Ļevicas vidū nav ienaidnieka!'' (\pltxti{Ļevica}~--- apkopojošs termins visiem poļu kreisajiem~--- komunistisko, sociāldemokrātisko, citu kreiso partiju un demokrātisko kustību~--- dalībniekiem starpkaru un kara periodā) 1935.~gada jūlijā~--- augustā PKP un PPS panāca vienošanos par to, ka abas partijas atturēsies no savstarpējiem uzbrukumiem un kopīgi cīnīsies par politieslodzīto atbrīvošanu. Kominternes VII kongresā viens no poļu komunistu vadītājiem J.~Ļenskis atzīmēja, ka Polijā sakarā ar agrārā jautājuma asumu un zemnieku kustības vērienu Tautas fronte var izveidoties pirms vienotas strādnieku frontes.

1935.~gada septembrī bija jānotiek jaunām parlamenta vēlēšanām. Pēc V.~Slaveka iniciatīvas tika izstrādāts jauns vēlēšanu likums. Pēc tā politiskās partijas vairs nevarēja izvirzīt savus kandidātus parlamenta vēlēšanās. To varēja darīt tikai vēlēšanu komisijas, kuras sastāvēja no administrācijas, saimniecisko un arodorganizāciju pārstāvjiem. Apgabalu vēlēšanu komisijas noteica divus kandidātus uz vienu deputāta vietu Seimā. Panākt opozīcijas kandidātu iekļaušanu šādu kandidātu sarakstā bija ļoti grūti. Tāpēc opozīcijai faktiski nebija citas izvēles kā boikotēt vēlēšanas. Šādā situācijā režīma un opozīcijas reālo ietekmi noteica nevis saņemto deputātu mandātu skaits, bet vēlētāju skaits, kas devās uz vēlēšanu urnām. Senāta vēlēšanās vēlētāju skaits tika ierobežots, ieviests izglītības cenzs (augstākā vai vidējā profesionālā izglītība), no 21 uz 24~gadiem paaugstināts vēlētājiem nepieciešamais vecums. Trešo daļu Senāta locekļu nozīmēja prezidents.

Šādos apstākļos tādas opozicionārās partijas kā PPS, ``\pltxti{Stronnitctwo Ludowe}'' u.c. nolēma vēlēšanas boikotēt. Boikota lozungs ieguva plašu popularitāti masās. Komunisti, vadoties no agrāk pieņemtajām nostādnēm, ka arī visgrūtākajos apstākļos jāizmanto visas legālā darba iespējas, gatavojās piedalīties vēlēšanās, taču, redzot tautas noskaņojumu, tomēr atteicās no dalības tajās un pievienojās boikotētājiem. Tā kā faktiski visa opozīcija vēlēšanas boikotēja, valdība pasludināja, ka aģitācija par vēlēšanu boikotu ir valsts noziegums. Tomēr ar nu jau mirušā diktatora sekotājiem~--- Bezpartijisko bloku sadarbojās tikai vācu nacistiskās ievirzes un dažas labējās rietumukraiņu un ebreju nacionālistiskās partijas. ``Sanācijas'' režīms sāka pieņemt totalitārismam raksturīgas iezīmes.

Seima vēlēšanas 8.~septembrī un nedēļu vēlāk notikušās Senāta vēlēšanas parādīja, ka režīmu atbalsta mazākā sabiedrības daļa. Saskaņā ar oficiālajiem datiem, vēlēšanās piedalījās tikai 45,6\% vēlētāju (ap 7,6~miljoni cilvēku, dažādos izdevumos skaitļi ir nedaudz atšķirīgi.) Taču arī šī statistika ir apšaubāma, jo paši augstākie vēlēšanu dalības rādītāji nāca no austrumu rajoniem. Ir aprēķini, ka vēlēšanās piedalījies tikai 37\% vēlētāju, citi dati~--- 48\% vēlētāju. Pēc dažiem datiem Varšavā pie vēlēšanu urnām devās tikai 23\% no vēlētājiem, Lodzā~--- tikai 13\%. Par derīgām tika atzītas tikai 35,5\% nodoto balsu. Ap 2~miljoni vēlētāju, kuri valdības draudu dēļ uz vēlēšanām ieradās, nodeva tukšus biļetenus. Tos pieskaitīja par Bezpartijisko bloku nodotajiem. Vēlēšanu rezultāti liecināja, ka par sanāciju balsoja tikai valsts kalpotāji, virsniecība, daļa pilsonības, valstij piederošajos uzņēmumos strādājošie un pastāvēja diezgan liela iespēja izveidot Tautas fronti pret sanācijas režīmu.

Varēja šķist, ka ir piepildījies J.~Pilsudska sapnis un varas rīcībā ir tai pilnīgi paklausīgs Seims, taču būtībā vēlēšanu rezultāti nozīmēja ne tikai Bezpartijiskā bloka, īpaši tā līdera V.~Slaveka sakāvi, bet arī visa ``sanācijas'' režīma sakāvi.

Iedzīvotāju mazākās daļas ievēlētais parlaments nebaudīja autoritāti. Samazinājās tā pieņemto likumu, iesniegto interpelāciju un priekšlikumu skaits.

V.~Slaveks plānoja radīt aktīvāku, disciplinētāku ``sanācijas'' organizāciju, taču izdarīt to nespēja. Atbilstoši I.~Pilsudska kuluāros izteiktajai gribai valsts prezidenta posteni bija jāieņem V.~Slavekam. Taču laikā, kad tika pieņemts jaunais pamatlikums, Pilsudskis jau bija smagi slims un nespēja (vai arī nevēlējās) organizēt valsts galvas nomaiņu.

Pēc J.~Pilsudska nāves viņa valdošajā grupā valdīja bažas, ka valstī varētu izraisīties nemieri un pat sacelšanās. Tādēļ I.~Moscicka nomaiņa šādos apstākļos bija nevēlama. Lēmums par valsts vadītāju maiņu tika atlikts līdz parlamenta vēlēšanām. Taču jau pirms vēlēšanām I.~Moscickis skaidri paziņoja, ka viņš nedomā atkāpties. Izrādījās, ka prezidentam, kurš deviņus gadus sabija diktatora ēnā, bija savas ambīcijas un uzskati, ka viņš varētu kalpot Polijas interesēm ne sliktāk kā V.~Slaveks. Pēc vēlēšanām ``pulkvežu'' grupai nebija citas izvēles kā atzīt savu sakāvi. I.~Moscicka atstādināšana būtu galīgi sakompromitējusi ``sanācijas'' nometni pēc tās realizētā trika, pieņemot Konstitūciju, un morālās sakāves vēlēšanās.

V.~Slaveks zudēja arī premjera posteni. Viņa vietu oktobrī ieņēma ģenerālis M.~Zindrams--Koscjalovskis, kuru atbalstīja E.~Ridzs-Smiglijs. Bezpartijiskais bloks sabruka. Var teikt, ka ``pulkvežu grupu'' pie varas nomainīja ``ģenerāļu grupa''.

Jaunu politisku ideju grupai nebija. Visi mēģinājumi radīt valsti stiprinošu, nacionāli ietonētu masu kustību, kura varētu aizvietot Bezpartijisko bloku, palika bez panākumiem. Valdība gan uzstājās ar solījumiem ``sadarboties'' ar sabiedrību, par cīņu pret ``karteļiem'' (monopolistu grupām), cenu samazināšanu uz rūpniecības precēm, taču bez sabiedrības pieprasītās politiskas amnestijas, nekas cits būtisks tā arī netika paveikts. Budžeta izdevumi galvenokārt (46\%) tika piešķirti aizsardzības vajadzībām, ap 10\% nācās tērēt dažādu ārzemju parādu atmaksai. Izdevumi veselības aizsardzībai, pabalstiem bezdarbniekiem arvien vairāk samazinājās. Līdzekļus budžetam valdība centās iegūt paaugstinot nodokļus, paplašinot ar nodokļiem apliekamo skaitu, pazeminot algas valsts un komunālajos uzņēmumos strādājošajiem strādniekiem un kalpotājiem. Pilnīgi saprotams, ka šie pasākumi nevairoja valdības popularitāti sabiedrībā.

Opozīcija pārgāja uzbrukumā. Pilsoniskās opozicionārās partijas mēģināja politisko cīņu pārcelt uz ielām. Diemžēl, bez pārmetumiem mazspējīgajai valdībai arvien asāk izskanēja uzbrukumi ebrejiem kā visu nelaimju vaininiekiem.

1936.~gada martā sākās plaši streiki Lodzā, pēc tam Krakovā. 23.~martā tur sadursmēs ar policiju tika nogalināti astoņi strādnieki. 25.~martā nogalināto apbedīšanā piedalījās desmiti tūkstošu cilvēku no visas apkārtnes. Aprīlī streika laikā Ļvovā atkal tika nogalināts tā dalībnieks. Notika plaša demonstrācija sakarā ar viņa apbedīšanu, atkal policija lietoja šaujamieročus, pēc preses datiem bija nogalināti 49~demonstrācijas dalībnieki, ap 300~cilvēku bija ievainoti. Kad strādnieki aplenca Brigidku (poļu \pltxti{Brygidki}, ukraiņu \uktxti{Бригідки}) cietumu, policija aicināja palīgā pilsētas garnizonu. Tā komandieris atteica palīdzību, atsaucoties uz to, ka viņš nav pārliecināts, vai karavīri šaus uz demonstrantiem. Pret strādniekiem vērstās represijas maz līdzēja. Ļvovas notikumu liecinieks bijušais premjers K.~Bartels, izsaukts uz Varšavu, lai sniegtu ``privātu'' ziņojumu valdībai, konstatēja: ``Tie vairs nav nemieri~--- tas ož pēc revolūcijas''. Taču iekšlietu ministrs F.~Slavojs Skladkovskis sakarā ar interpelāciju Seimā par šiem notikumiem paziņoja: ``Policija šāva un šaus''.

Zemnieku neapmierinātība sasniedza kulmināciju 10~dienas ilgajā zemnieku streikā 1937.~gada augustā, kas aptvēra desmit vojevodistes un kurā pēc dažādiem vērtējumiem piedalījās ap 8--10 miljoniem zemnieku. Tas beidzās ar bruņotām sadursmēm. Oficiāli tika paziņots par 42~nogalinātajiem un 34~ievainotajiem, faktiski cietušo skaits bija lielāks. Arestēja ap 4~000 zemnieku.

``Sanācijas'' režīmam pretojās ne tikai Tautas frontes radīšanas piekritēji ar komunistiem priekšgalā, bet arī komunistu pretinieki: daļa kristīgo demokrātu, ļudoviešu, sociālistu un endeku. Otrie 1936.~gada sākumā radīja t.s. \strong{``Moržas fronti''} (\entxti{Front Morges}~---  no \entxti{Morges} pilsētas nosaukuma Šveicē, kur I.~Paderevska villā notika šo organizāciju pārstāvju sanāksme). To vadīja emigranti: ģenerāļi V.~Sikorskis, J.~Hallers, komponists I.~Paderevskis, politiķi V.~Vitoss, V.~Korfantu. Šie politiķi nespēja vienoties ar PPS, kopdarbību nepieļāva nacionālistu un sociālistu nepārvaramās pretrunas. ``Moržas frontei'' nebija atbalsta plašākās masās.

Izrādījās neiespējama arī Moržas frontes sadarbība ar komunistiem. Poļu komunistu vairākums nespēja iedomāties, kā var sadarboties ar bijušajiem J.~Pilsudska līdzgājējiem. Komunisti bija aktīvi, spēja nodrošināt arī savu rindu pieaugumu. 1936.~gadā neraugoties, uz arestiem, partijā bija 33~tūkstoši biedru. Viņi aktīvi uzstājās pret sanācijas mēģinājumiem vienot sabiedrību ap J.~Pilsudska varas mantiniekiem. Tā, 1936.~gada novembrī PKP CK aicinājumā bija teikts: ``Runā, ka ģenerālis Ridzs-Smiglijs ir taisnīgs, ka viņš ierobežos neliešus un izbeigs nelikumības. To pašu kādreiz runāja par Pilsudski. Taču vai nelietība, noziegums un apspiešana ir kādreiz plosījusies vairāk nekā Pilsudska ``sanācijas'' laikā?''

Taču \strong{Polijas komunisti} saņēma negaidītu triecienu no ārienes.

1935.~gada rudenī Komunistiskās Internacionāles Izpildu komitejas Sekretariāts apsprieda kā Polijas Komunistiskā partija pilda Kominternes VII kongresa lēmumus. Tas tika darīts sekmīgi, taču pēkšņi apspriešana guva negaidītu ievirzi. Tajā izskanēja pārmetumi poļu komunistiem spiegošanā un provokācijā, tika ierosināts likvidēt Kominternes poļu sekciju, kura nespējusi aizkavēt ``spiegu un provokatoru iekļūšanu VK(b)P''. Tika nolemts nepieļaut poļu komunistu emigrāciju uz PSRS. Poļu komunistu Padomju Savienībā tai laikā bija ap 5~000. (Pavisam Padomju Savienībā 30.~gadu vidū dzīvoja ap 18~tūkstošu poļu politemigrantu: visvairāk, protams, komunistu, bet arī kreiso sociālistu un radikālo zemniecības politiķu grupas). 1935.~gada decembrī VK(b)P CK Politbirojs pieņēma lēmumu par t.s. ``logu'' uz PSRS~--- Polijas robežas, kurus IeTK bija ierīkojis Kominternes vajadzībām, lai tā varētu sazināties ar PKP, RBKP, RUKP. Šo partiju pārstāvji, kuri šai laikā atradās PSRS, vairs nevarēja to atstāt. 1935.~gada beigās sakarā ar KIIK Sekretariāta lēmumu par citu valstu komunistisko partiju biedru pārcelšanu uz VK(b)P ap 2~000 no tiem iekļāvās VK(b)P vietējās organizācijās, pārējie palika ārpus tām, jo tika atzīts, ka PKP ir ``piesārņota'' ar trockistiem un poļi vietējās organizācijās ir nevēlami.

1936.~gada 3.~janvārī KIIK sekretārs D.~Manuiļskis nosūtīja VK(b)P CK vadošo kadru nodaļas vadītājam N.~Ježovam vēstuli, izvirzot jautājumu par līdzekļiem PKP pārbaudei, ``\dots{}kura, kā Jums zināms, pēdējos gados ir galvenais spiegu un provokatoru elementu piegādātājs PSRS''. 28.~janvārī KIIK sekretariāts atzina ``stāvokli Polijas Kompartijā par ļoti nopietnu, kurš prasa radikālus līdzekļus partijas atveseļošanai un vadības nostiprināšanai centrā un uz vietām.'' KIIK piedāvāja izslēgt no PKP CK personas, ``kuras neizsauc politisku uzticību'', ``atlaist visvairāk ar provokāciju skartos partijas organizācijas posmus'' utt. Kominternes poļu sekcija tik tiešām tika likvidēta.

Polijas Komunistiskās partijas centieni izvērst legālu darbību PSRS izsauca personīgu J.~Staļina nepatiku. Viņš šai sakarā izteica viedokli, ka vajadzētu izsūtīt poļu komunistus no Padomju Savienības, lai tie varētu ierasties pie Polijas varas iestādēm, kas savukārt tos ieslodzītu cietumos vai Bereza Kartuskas nometnē. Kad to uzzināja poļu komunistu vadītāji, viņi lūdza J.~Staļinu tos pieņemt un mainīt nostāju. Tikšanās beidzās ar to, ka padomju diktators kā alternatīvu poļu komunistiem piedāvāja masveidā doties karot uz Spāniju 1936.~gada rudenī veidojamo Internacionālo brigāžu sastāvā.

Poļu kreisie aktīvi iesaistījās 1936.--1939.~gada Spānijas pilsoņu kara notikumos. Polijas Komunistiskās partijas Centrālā Komiteja izvirzīja lozungu ``Spānijas lieta~--- tā ir Polijas lieta''. Kaut Polijas valdošie slāņi dažādi centās nomelnot republikānisko Spāniju, faktiski atbalstīja ģenerāļa F.~Franko vadītos dumpiniekus, sūtot tiem apbruņojumu, ap 5~tūkstoši demokrātiski noskaņotu poļu, galvenokārt no zemnieku, sociālistu un komunistu partijām, devās kā brīvprātīgie uz Spāniju lai cīnītos pret F.~Franko falangistiem. Tas tika darīts pat neraugoties uz to, ka Polijas valdība viņiem atņēma pilsonību. No atsevišķiem poļu brīvprātīgo bataljoniem cīņās pie Madrides tika izveidota 13.~Internacionālā, poļu revolucionāra J.~Dombrovska vārdā nosauktā brigāde. Komandēja to poļu izcelsmes padomju ģenerālis K.~Sverčevskis, kuru tur sauca par ``ģenerāli Valteru''. Apmēram 10\% visu Internacionālo brigāžu cīnītāju bija poļi. Polijā tika vākti līdzekļi palīdzībai Spānijas bērniem un trūkumcietēju ģimenēm.

1956.~gadā Polijas Tautas republikas Valsts padome nodibināja medaļu ``Par jūsu un mūsu brīvību'', ar kuru tika apbalvoti Internacionālo brigāžu cīnītāji, kuri Pilsoņu kara (1936--1939) Spānijā laikā karoja republikāņu rindās. 1996.~gadā sakarā ar Pilsoņu kara Spānijā sākuma 60.~gadadienu visi Internacionālo brigāžu veterāni, tai skaitā arī poļi, ieguva Spānijas goda pilsonību.

1937.~gada 16.~augustā aptaujas kārtībā tika apstiprināts KIIK prezīdija lēmums par Polijas Komunistisko partiju.

28.~novembrī to noformēja ar KIIK protokolu un KIIK ģenerālsekretārs G.~Dimitrovs par to informēja J.~Staļinu, kurš lēmumam piekrita, bet iebilda pret tā publicēšanu. Lēmumā ``uz neapstrīdamu dokumentālu datu pamata'' bija konstatēts, ka ``poļu kompartijas vadošajās iestādēs vairākus gadus darbojās ienaidnieki, poļu fašisma aģenti \citespace{} pilsudčiki komunistiskajā partijā iesūtīja savus spiegus un provokatorus, kuri maskējās kā opozicionāri elementi''. Lēmumā bija teikts, ka, lai radītu tiem autoritāti strādnieku vidū ``fašisms savus aģentus ne reizi vien ieslodzīja cietumā, inscenējot viņu tiesāšanas komēdiju, lai, rodoties pirmajai izdevībai, atbrīvotu vai nu organizējot to bēgšanu vai apmainot pret spiegiem un diversantiem, notvertiem nozieguma vietā PSRS''. Tā PKP vadībā ``iesēdās spiegu un diversantu banda'', kas savukārt ``grāva vietējās partijas organizācijas, diskreditēja komunistisko partiju masu acīs, traucēja poļu un padomju tautu tuvināšanos''. Uz šādu secinājumu pamata arī tika nolemts ``Polijas komunistisko partiju atlaist'', bet ``visiem godīgajiem komunistiem'' ieteikt līdz partijas atjaunošanai darbības smaguma centru pārcelt uz masu organizācijām, ``cīnoties par vienotas strādnieku kustības izveidi un antifašistiskas frontes radīšanu Polijā''.

Traģisku likteni piedzīvoja Padomju Savienībā palikušie poļu komunisti. Staļinisko represiju gaitā viņi tika apcietināti. Aresti notika vairākos viļņos. Pirmais norisa jau 1933.--1934.~gadā, divi nākamie 1935.--1936.~un 1937.~gadā. Arestēti tika vairāk nekā 3~000 komunistu un daudzi viņu ģimenes locekļi. Vadošie komunisti no Polijas tika izsaukti uz PSRS un tur noslepkavoti. PSRS IeTK īstenoja t.s. ``poļu operāciju'' kā pirmo un mērogos apjomīgāko no ``nacionālajām operācijām'', kuras IeTK vērsa pret PSRS dzīvojošo nacionālo mazākumu pārstāvjiem. (To acīmredzot var izskaidrot ar to, ka Polija 20.--30.~gados tika uzskatīta par pašu bīstamāko no tiešajiem PSRS kaimiņiem, ka pārbēdzēju no Polijas bija vairāk nekā no citām kaimiņvalstīm, kas pēc 1937.--1938.~gadā izplatītās loģikas liecināja, ka arī bāze Polijas izlūkdienestam bija plašāka nekā citu valstu izlūkdienestiem.) PSRS iekšlietu tautas komisāra N.~Ježova pavēlē par poļu arestu bija iekļauti: ``\dots{}politemigranti un apmainītie ieslodzītie no Polijas, \citespace{} bijušie Polijas Sociālistiskās partijas un citu politisko partiju biedri, \citespace{} visi Padomju Savienībā palikušie poļu armijas karagūstekņi, \citespace{} visaktīvākie vietējie pretpadomju un nacionālistiskie elementi poļu rajonos''.

Pēc staļinisko represiju pētnieka A.~Gurjanova minētiem datiem 1938.~gada sākumā Polijas Komunistiskajā partijā bija ap 7~tūkstošiem biedru, no kuriem ap 5~tūkstošiem atradās PSRS. 1938.~gada jūnijā PKP CK gan vēl publiskoja aicinājumu, kurā norādīja uz Polijai draudošajām briesmām no Vācijas. Taču 1938.~gadā PSRS staļinisko represiju gaitā (pēc t.s. ``poļu sarakstiem'') bojā gāja poļu komunistu vadītāji Ā.~Varskis, J.~Umšlihts, J.~Ļenskis u.c. Pavisam šai laikā dzīvību zaudēja 69\% no PKP CK locekļiem un locekļu kandidātiem, kuri galvenokārt tika nošauti bez tiesas.

Kopā pēc PSRS IeTK 1938.~gada datiem poļu izcelsmes ieslodzīto skaits PSRS sastādīja 134~519 cilvēku, no kuriem no 50~000~līdz 67~000 tika nošauti. Pēc pētnieku N.~Petrova un A.~Roginska jau pēc PSRS sabrukuma savāktajiem datiem t.s. ``poļu operācijā'' 1937.--1938.~gadā tika notiesāti 139~815 cilvēki (ap 10\% no visiem 1937.--1938.~gadā notiesātajiem), no tiem uz nāvi~--- 111~071. Poļu vēsturnieks D.~Bočkovskis raksta, ka ``poļu lietā'' no 143~810 notiesātajiem nošauti tika 111~1091.

Kaut ne visi ``poļu operācijā'' skartie bija poļi pēc tautības, tā bija pati asiņainākā vēršanās pret kādas tautas pārstāvjiem staļiniskajā Padomju Savienībā. Galvenokārt poļiem no PSRS rietumu līdz pat austrumu robežām tika izvirzītas apsūdzības par spiegošanu Polijas labā.

Faktiski tik apjomīgu spiegošanu pret PSRS Polija nevarētu veikt, pat sasprindzinot visus savus spēkus. Armijas Ģenerālštāba Otrajā (Izlūkošanas) nodaļā (\pltxti{Oddział II Sztabu Generalnego Wojska Polskiego}) 1923.~gadā strādāja 58~virsnieki, 1928.~gadā 47~virsnieki, 64~zemāku pakāpju karavīri un 37~civilpersonas, kas nodarbojās ar izlūkošanu, pretizlūkošanu, diversiju gatavošanu pret Austriju, Čehoslovākiju, Lietuvu, Rumāniju, Ungāriju, Vāciju un PSRS. Protams šie darbinieki vadīja aģentu tīklus, taču to skaits nekādi nevarēja sniegties desmitos tūkstošu. 1939.~gadā, pat pēc 1936.--1939.~gadā veiktā izlūkošanas darbība pret visreālākajiem potenciālajiem ienaidniekiem: Vāciju un PSRS bija ievērojami paplašināta. Otrajā nodaļā strādāja no 100~līdz 250 virsnieku un ap 500 civilo darbinieku. Visā Polijas izlūkošanas tiklā ārzemēs darbojās ap 4~000 aģentu. Tādejādi PSRS pret poļiem izvirzītās lietas gandrīz visas bija vienkārši safabricētas.

Kopā ar PKP ``atlaida'' arī Rietumu Baltkrievijas un Rietumu Ukrainas komunistiskās partijas. Polijā dzīvojošie komunisti gan turpināja darboties sabiedriskajās organizācijās, bet kreisā kustība savu organizatorisko spēku zaudēja. Tiesa, pēc pusgada Parīzē pēc Kominternes Izpildu komitejas iniciatīvas tika sapulcināta poļu komunistu grupa, pamatā no bijušajiem Internacionālo brigāžu cīnītājiem Spānijā, lai sagatavotu partijas atjaunošanu. Vēl pēc dažiem mēnešiem~--- 1939.~gada 16.~maijā Kominternes Izpildu komiteja arī pieņēma attiecīgu lēmumu, taču drīzumā sācies Otrais pasaules karš nedeva iespēju to īstenot. (Tikai 1956.~gadā, kad Kominterne vairs nepastāvēja, Padomju Savienības Komunistiskā partija un dažas citas kompartijas atzina lēmumu par nepamatotu.) Dabīgi, ka notikumi PSRS sarežģīja komunistu darbību arī pašā Polijā, traucēja pretestību autoritārajam režīmam tajā. Poļu tautības pilsoņu, tai skaitā komunistu, noslepkavošana Padomju Savienībā kompromitēja ne tikai Komunistisko partiju, bet arī Nacionālās frontes ideju.

Polijas Komunistiskās partijas atlaišana 1938.~gadā strauji pasliktināja visu poļu kreiso partiju sadarbības iespējas. Tas izpaudās arī Otrā pasaules kara laikā.

Polijā pastāvēja vairākas strādnieku partijas. Nacionālā strādnieku partija (\pltxti{Narodowa Partia Robotnicza}~--- NPR), kura bija radīta 1920.~gadā, pretendēja uz strādnieku apvienošanu, propagandēja nacionālās solidaritātes ideju un sasniedza pat līdz 100~000 biedru, galu galā tomēr tā arī nespēja apvienot ievērojamu strādnieku daļu vairākumu un 1937.~gadā apvienojās ar kristīgajiem demokrātiem (\pltxti{Polskie Stronnictwo Chrześcijańskiej Demokracji, Chrześcijańska Demokracja}~---  PSChD), izveidojot Darba partiju (\pltxti{Stronnictwo Pracy}) ar kristīgi-demokrātisku raksturu, kurā strādnieku pārstāvniecība jau bija krietni mazāka.

Ja sākotnēji Polijas Sociālistiskās partijas biedru skaits nemitīgi samazinājās (1925.~gadā tajā bija 65~tūkstoši biedru, bet 1931.~gadā~--- 32~tūkstoši un 1933.~gadā~--- tikai 19~tūkstošu), tad 30.~gadu otrajā pusē ar partijas aktīvāku iesaistīšanos streiku cīņā tās biedru skaits sāka pieaugt. 1935.~gadā partijā to bija 25~tūkstoši, bet 1938.~gadā~--- jau 40~tūkstošu. Partijā iezīmējās kreisais spārns ar N.~Barlicki un S.~Dubuā priekšgalā, kurš bija gatavs uz sadarbību arī ar komunistiem. Tomēr Polijas strādnieku šķira starpkaru periodā palika sašķelta.

Neizdevās organizēt arī strādniecības un zemniecības sadarbību. Polijas Sociālistiskā partija un ļudovieši (\pltxti{Stronnictwo Ludowe}) ``sanācijas'' apstākļos bija pārdzīvojušas zināmu ideoloģisku evolūciju, taču ne tik radikālu, lai nopietni uztrauktu valdošo režīmu. Tāpat kā iepriekšējos gados viņu opozīciju ierobežoja bailes satricināt valstiskuma pamatus. Turklāt katra no šīm partijām pretendēja uz līdera vietu ``anti-sanācijas'' spēku priekšgalā, kā rezultātā starp tām nevarēja izveidoties cieša sadarbība. Abas partijas nākotni redzēja atšķirīgi. PPS cīnījās par sociālistisku republiku, kurā dominētu strādnieku intereses, SL~--- par zemnieku republiku, kurā strādnieku intereses būtu pakļautas zemnieku interesēm.

Bez rezultātiem palika arī demokrātiski noskaņotās poļu inteliģences mēģinājumi izveidot spēcīgu pretošanās kustību ``sanācijai''. 1937.~gada beigās sāka veidoties demokrātiskie klubi, kuri 1939.~gadā apvienojās Demokrātiskajā partijā (\pltxti{Stronnictwo Demokratyczne}), kuras līderi bija M.~Mihalovičs, M.~Kvasņevskis un citi.

Vēsturnieki ir nonākuši pie secinājuma, ka pret sanāciju vērstajiem sabiedriski politiskajiem spēkiem pietrūka gan spēku, gan vienotības, gan prasmes, lai gāztu sanācijas režīmu, pavērstu Polijas vēsturi demokrātiskas attīstības virzienā. Jāuzsver, ka sanācijas režīms lielā mērā balstījās uz represijām. Tā pēc oficiālajiem Polijas datiem 1937.gadā tiesas par noziegumiem pret valsti notiesāja 3~755 cilvēkus, tai skaitā par valsts nodevību 2~945~cilvēkus.

Kaut ``sanācijas'' režīma pretinieki nebija vienoti, tam nācās manevrēt. Par galvenajiem valdības nometnes pārstāvjiem kļuva I.~Moscickis un E.~Ridzs-Smiglijs (dažkārt tiem blakus min arī ārlietu ministru J.~Beku). Pēc zināmām svārstībām E.~Ridzs-Smiglijs nolēma nostiprināt savas pozīcijas, apvienojot ``stipras rokas'' piekritējus ``sanācijas'' nometnē, tikt vaļā no par viņu gandrīz 20~gadus vecākā prezidenta un saņemt varu savās rokās. Ja I.~Moscickis pārstāvēja vairāk liberālus uzskatus, tad zināmu kompromisu iezīmēja pēdējās starpkaru valdības ar premjeru ģenerāli F.~Slavoju-Skladkovski nozīmēšana 1936.~gada 15.~maijā. Valdībā ietilpa ``pulkvežu'' grupas pārstāvji un arī J.~Pilsudska kursa atbalstītāji no politiskā spektra kreisā spārna.

Bezpartijiskais bloks sadarbībai ar valdību (\pltxti{Bezpartijny Blok Wspołpracy z Rządem}) tika atlaists. Pēc tam ``stingrās rokas'' piekritēji ar E.~Ridzs-Smiglija atbalstu 1937.~gada februārī kā pakāpi uz vienotas partijas izveidi radīja Nacionālās vienības nometni (\pltxti{Obòz Zjdnoczenia Narodowego}). Tai izdevās piesaistīt pat daļu endeku. Formālais biedru skaits sasniedza divus miljonus. Tomēr lielākās partijas~--- PPS un nacionāldemokrātu vairākums atteicās no līdzdarbības. Tika izstrādāti plāni kā panākt vienošanos ar endekiem un pievarēt kreisi-centrisko opozīciju. Lai to izdarītu nepietika izvērst represijas, kas draudēja ar stāvokļa saasināšanos, pat pilsoņu karu. Taču 1938.~gada janvārī Nacionālās vienības nometne atteicās no tālejošiem politiskiem plāniem, no sadarbības ar endekiem, arvien vairāk zaudēja enerģiju un pārvērtās par ko līdzīgu agrāk pastāvējušajam Bezpartijiskajam blokam.

Gan valdošajās, gan opozīcijas aprindās auga noraidošs noskaņojums pret 1935.~gada Konstitūciju, tika cerēts, ka pēc jaunām vēlēšanām Konstitūcija varētu tikt demokratizēta.

Taču 1938.~gada 13.~septembrī ar prezidenta dekrētu abas parlamenta palātas tika atlaistas un jaunas parlamenta vēlēšanas notika novembrī, īsi pēc tam, kad Polija bija sagrābusi Čehoslovākijai piederošo Cešinas apgabalu (Skat. apakšnodaļu ``Polijas ārpolitika'') un valdošais režīms izvērsis šī ``triumfa'' propagandu. Opozīcija boikotēja vēlēšanas, taču tas nenotika tik asi, ar tik plašu vērienu kā trīs gadus iepriekš. Šai laikā ārējās briesmas nedaudz nolīdzināja iekšējo strīdu asumus. Vēlēšanās oficiāli it kā piedalījās 67\% vēlētāju (ir gan arī dati, ka vēlēšanās piedalījušies tikai 1/3~vēlētāju), taču tāpat kā iepriekš tika realizēts spiediens uz vēlētājiem, bija daudz falsifikācijas gadījumu.

Padomju izlūkdienests no sava poļu aģenta saņēma ziņas, ka vēlēšanās piedalījās 46\% vēlētāju, taču valdošos poļu politiķus uztraucis ne tik daudz šis mazais skaits, cik tas, ka no šo 46\% vēlētāju nodotajām balsošanas zīmēm 15\% nācās atzīt par nederīgām. Iekšlietu ministrija uzskatīja, ka nederīgās zīmes nodevuši jaunieši un tie J.~Pilsudska piekritēji, kuri ``nevienu neatzīst starp jaunajām autoritātēm un sāk atiet no valdošā bloka''. Šie 15\% vēlētāju tā vietā, lai atzīmētu, ko viņi vēlas ievēlēt, uz balsošanas zīmēm rakstīja ``pretvalstiskus un apvainojošus uzrakstus''.

Nacionālās vienības nometnes kandidāti ieguva ap 80\% deputātu mandātu. Parlamenta vēlēšanu rezultātu neatbilstību patiesajai tautas attieksmei pret ``sanācijas'' režīmu parādīja municipālās vēlēšanas, kuru 1.~kārta notika jau 18.~decembrī. Pēc oficiāliem datiem tajās piedalījās ne vairāk par pusi vēlētāju, pie tam 56\% mandātu ieguva opozicionārās partijas, kuras pirms tam boikotēja parlamenta vēlēšanas. Taču, kad 1938.--1939.~gadā notikušajās municipālajās vēlēšanās ievēlētās pilsētu un lauku padomēs sanācijas piekritēji izrādījās mazākumā, valdība atlaida tās, neapstiprināja to vadītājus. Tā bez jebkāda paskaidrojuma valdība neapstiprināja likumīgi ievēlētos Lodzas un Krakovas prezidentus.

Svarīgs poļu tautas demokrātisks iekarojums bija sabiedrības dalība vietējo tiesu darbībā. Līdz 1926.~gada apvērsumam zvērināto tiesā tiesneša funkcijas laiku pa laikam varēja pildīt katrs no iedzīvotāju ievēlētajiem tiesas locekļiem. Šādas tiesas ``austriešu'' Polijā eksistēja jau izsenis. Pēc Polijas neatkarības atjaunošanas I.~Dašinska, E.~Moračevska un V.~Vitosa valdības šo zvērināto tiesu sistēmu izplatīja uz visu Poliju. Ar sanācijas režīma izveidi sākās zvērināto tiesu kompetences ierobežošana līdz pat to likvidēšanai. 1939.~gadā tika noteikta kārtība, ka visus tiesnešus nozīmēja valdība.

Kara briesmu priekšā ``Moržas fronte'' mīkstināja savu opozicionāro nostāju. Daži pazīstami emigranti atgriezās Polijā. (Tiesa, ne visiem tas izdevās. Piemēram, 1938.~gadā Polijā vēlējās atgriezties ģenerālis V.~Sikorskis. Taču toreizējais premjerministrs ģenerālis F.~Skladkovskis atteicās to atļaut, paziņojot, ka ``viņam nepietiks policijas, lai apsargātu Sikorska personu''. Starp citu, vēlāk, jau sākoties karam, V.~Sikorskis vēlreiz vēlējās atgriezties dzimtenē un doties uz fronti. Šoreiz to nepieļāva maršals E.~Ridzы-Smiglijs.)

Saasinoties Vācijas pretpoliskajai politikai, ļudovieši, no vienas puses, pastiprināja kritiku pret sanācijas režīma sadarbību ar Vāciju, taču, no otras puses, sāka ``apslāpēt'' savu pret šo režīmu vērsto iekšpolitisko darbību. Polijai kritiskajās dienās 1938.--1939.~gadā ļudovieši gandrīz pilnībā pārtrauca sanācijas kritiku un atbalstīja tā ārpolitiku. Ļudovieši atbalstīja sanāciju konfliktā ar Lietuvu, apsveica poļu karaspēka ievešanu Cešinas Silēzijā. Partijas prese rakstīja: ``Tad, kad runa iet par valsts aizsardzību ārēju briesmu priekšā, izdziest atšķirības starp nometnēm''. Opozīcija piedāvāja radīt daudzpartiju valdību, lai nosargātu nacionālo neatkarību, taču sanācijas piekritēji noraidīja priekšlikumu. Faktiski sanācija, atsaucoties uz ārējām briesmām, varēja turpināt demokrātisko spēku apspiešanu. Tiesa, arī sanācijas darboņu realizētā iekšējā politika kļuva nosvērtāka, mazāk radikāla. 1935.--1939.~gadā Polijā tā arī neizveidojās totalitāra valsts pēc Vācijas vai Itālijas parauga, nenotika valsts fašizācija. Tam par šķērsli bija kā valdošo grupējumu neviendabīgums un sašķeltība, visas tautas akceptēta līdera trūkums, opozīcijas partiju, kam bija atbalsts sabiedrībā, pretestība, plašu masu antifašistisks noskaņojums.

1939.~gada sākumā un vasarā sabiedrībā pārsvarā valdīja patriotisma gars, gatavība nest upurus dzimtenes labā. Taču dziļi simptomātisks bija fakts, kas lika apšaubīt it kā gandrīz sasniegto galveno politisko spēku relatīvo samierināšanos ārējo briesmu priekšā. 1939.~gada 3.~aprīlī J.~Pilsudska līdzgaitnieks pulkvedis V.~Slaveks, palicis bez politiskiem amatiem un esot opozīcijā E.~Ridzs-Smiglija grupai, izdarīja pašnāvību. Maršals gan piedalījās V.~Slaveka bērēs, taču nelaiķa draugi neatļāva viņam nest zārku. V.~Slaveka apbedīšana faktiski kļuva par to agrāko J.~Pilsudska piekritēju manifestāciju, kuri atradās opozīcijā esošajiem valsts prezidentam un armijas vadītājam. Pastāv gan arī nepierādīta versija, ka V.~Slaveks bija viens no pret valdību vērstas sazvērestības organizatoriem un beidza dzīvi, kad viņam draudēja neizbēgams arests.

\strong{``Sanācijas'' režīma vērtējums kopumā ir pretrunīgs.} Kā \strong{pozitīvo} var atzīmēt, ka: nodibinājās \strong{nosacīta politiska stabilitāte}, tika pārtrauktas parlamentārās ķildas; sabiedrība netika kontrolēta pilnīgi visās darbības jomās, preses un apziņas brīvība bija daudz mazāk ierobežota nekā citās autoritārās valstīs, norisa valsts birokrātiskā aparāta profesionalizācija; budžeta sabalansēšana, ar daudzu uzņēmumu pārvalstiskošanas īstenošanu it kā tika rasta izeja no ekonomiskās krīzes. Polijā netika pielietots tik plaša apjoma terors, kā tas bija totalitārajos režīmos, teroram bija izlases raksturs, par tā upuriem kļuva nevis vienkārši citādi domājošie, bet tikai radikālākie režīma pretinieki. Pie tam revolucionāro un ``separātisko'' (nacionālo) organizāciju aizliegšanas prakse nebija sanācijas režīma izgudrojums, bet tika mantota no iepriekšējās, nosacīti demokrātiskās Polijas.

Galvenā sanācijas \strong{negatīvā iezīme} bija \strong{demokrātijas trūkums} (cenzūra, politisko pretinieku vajāšana, varas koncentrēšana vienas personas rokās). Lēmumi tika pieņemti bez saskaņošanas ar sabiedrības vairākumu. Režīms uzspieda sabiedrībai politisku pasivitāti, tauta tika vadīta kā aitu bars. Pēc autora domām, šis trūkums pārsvēra visus ``sanācijas'' panākumus. Taču par to var diskutēt bezgalīgi.

Mūsdienās \strong{par} t.s. \strong{Otro Republiku} (pareizāk~--- starpkaru Poliju) \strong{poļu tautā valda pozitīvs vērtējums}. Kaut svarīgākie jautājumi~--- agrārais, industrializācijas, nacionālo mazākumu~--- ne tuvu nebija atrisināti, tautas atmiņā šis laiks saglabājies kā neatkarības periods pēc nācijas, tās zemes sadalīšanas svešu varu starpā. Pie tam neatkarība saistījās ar lieliem teritoriāliem ieguvumiem, kuri ļāva poļos pastāvēt lepnumam par savu ``lielvalsti''. Neraugoties uz īso Polijas valsts pastāvēšanu starpkaru periodā, poļu tautas apziņā šis laiks iegāja ja ne kā vispārējas labklājības, tad tomēr kā relatīvi laimīgas dzīves periods.

Sociālajā jomā neatkarības gadi neienesa Polijas sabiedriskajā struktūrā būtiskas izmaiņas. Lielākais sasniegums bija tas, ka pakāpeniski izzuda atšķirības starp agrāk dažādām valstīm piederošajiem poļu novadiem. Svarīgi bija, ka vesela poļu paaudze dzīvoja apstākļos, kad valdīja uzskats, ka neatkarība ir normāls valsts stāvoklis. Tāpēc tās zaudējums II pasaules karā tika uzskatīts par pagaidu parādību.

Otro Žečpospolitas bojā eju pamatos noteica nevis sanācijas valdības kļūdas iekšpolitikā, bet gan ārpolitiskā situācija, taču sagaidīt smagos pārbaudījumus Otrā pasaules kara laikā poļu tauta varēja arī vienotāka, spēcīgāka, nekā tas bija realitātē. Vainojams par to, ka tas tā nebija, ir sanācijas režīms. Savas eksistences noslēguma posmā sanācijas režīms pastāvēja lielā mērā pateicoties t.s. spēka struktūru lojalitātei un masu politiskās inerces spēkam. Tas neatrada kopēju valodu nedz ar kreisajiem, nedz labējiem, bet centās prezentēt sevi kā virs partijām stāvošu vispārnacionālo interešu aizstāvi. Taču šīs it kā vispārnacionālās intereses galvenokārt atbilda pašu par politiķiem kļuvušo militāristu interesēm. Ļoti populāras personības vadītais režīms šo popularitāti izmantoja lai nodrošinātu valstī stabilitāti, kura līdzinājās stagnācijai. Izveduši valsti no politiskā haosa 20.~gadu pirmajā pusē, J.~Pilsudskis un tā sekotāji pasaules ekonomiskās krīzes laikā un pēc tās tā arī nespēja piedāvāt tautai perspektīvu tālākās attīstības programmu.

\subsection{Nacionālās attiecības starpkaru Polijā}

Visā starpkaru periodā nacionālās problēmas Polijā palika neatrisinātas. Jau ar pirmajiem atjaunotās nacionālās Polijas soļiem tīri objektīvi bija jāsarežģījas attiecībām starp jauno valsts iedzīvotāju vairākumu~--- poļiem un agrāk tās dažādos rajonos valdošo valstu pamatnāciju~--- vāciešu un krievu~--- nacionālajiem mazākumiem, kas tagad, zaudējuši savas privilēģijas, ar nostaļģiju atcerējās agrākos laikus. Vāciešiem un krieviem, degradētiem par nacionālajiem mazākumiem, bija grūti samierināties ar jauno stāvokli.

Bet tādi ievērojami nacionālie mazākumi kā ukraiņi, baltkrievi, arī lietuvieši blakus Polijai redzēja savas valstis. Tas viņos radīja dabīgas cerības, ka agri vai vēlu to iedzīvotāji nāks palīgā Polijā apspiestajiem tautiešiem.

Pārvarēt šo objektīvo atsvešinātību bija iespējams tikai ar izsvērtu valsts nacionālo politiku, kura neatstātu vietas iespējamiem pārmetumiem par jaunu nacionālo privilēģiju radīšanu. Taču sagaidīt šādu politiku nevarēja nedz no agrākās nacionālās apspiestības cietušajām poļu tautas masām, kuras alka atlīdzinājumu par agrāk nodarītajām pārestībām, nedz no Polijas valdošajiem slāņiem, kuri ar nacionālajiem lozungiem bija panākuši valsts atjaunošanu, un arī turpmāk centās ar tiem nostiprināt savu vadošo vietu tajā.

Skaistais lozungs ``Par jūsu un mūsu brīvību'', kuru poļu nacionālās kustības darbinieki plaši propagandēja cīņas pret carismu laikā, tagad neatbilda poļu elites interesēm. Tās realizētās nacionālās politikas būtību starpkaru periodā var izteikt kā ``Poliju poļiem!''

Poļu teritoriju sadalīšanas periodā nacionālo apspiestību tai vai citā mērā izjuta visi poļu iedzīvotāju slāņi, tai skaitā arī darbaļaudis. Uzticība nacionālajai idejai, gatavība cīnīties par savas zemes neatkarību bija raksturīga vairākumam poļu zemnieku, lielai daļai poļu strādnieku. Tāpat kā citi iedzīvotāju slāņi viņi, protams, neaizmirsa savas šķiriskās intereses, taču, kad vajadzēja cīnīties par savas valsts neatkarību, viņi bija gatavi nest lielākus upurus nekā to valstu iedzīvotāji, kuras nebija izjutušas svešzemju kundzību. Bez šīs starpkaru Polijas īpatnības ievērošanas grūti saprast tās iekšējās un ārējās politikas objektīvo nosacītību, tādas parādības kā J.~Pilsudska milzīgā popularitāte, PPS ietekmes pārsvars poļu strādnieku kustībā, stiprais patriotisms gan poļu sabiedrības augšslāņos, gan arī apakšslāņos. Starpkaru periodā poļu sabiedrības apziņa raksturojās ar saasinātām nacionālajām un arī nacionālistiskajām jūtām.

Daudzi vēsturnieki un etnopsihologi pamatus tam atrod poļu nacionālajā pašapziņā. XIX~gadsimtā, sargājot savu nacionālo identitāti, poļi piesavinājās uzskatu, ka polis~--- tas ir ``mūsu cilts pārstāvis'', runājošs mūsu valodā, ievērojošs mūsu tikumus un parašas, ir katolis. Poļi sīksti aizstāvēja savu valodu, ieradumus un ticību gan pret sadalītājvalstu antipolisko politiku, gan straujajām izmaiņām pasaulē. Pašus sevi poļi pārsvarā redzēja kā augstsirdības un nevainības iemiesojumus. Ja pagātnē viņi arī bija kļūdījušies, tad tikai paužot naivu ideālismu un labestību pret svešiniekiem, bet tie šīs poļu īpašības svešinieki izmantojuši lai viņus pakļautu. Pēc valstiskās neatkarības atgūšanas šādi priekšstati turpināja pastāvēt plašos poļu slāņos. Ārzemēs skanošā kritika par slikto attieksmi pret cittautiešiem, antisemītismu tika uzņemti kā apmelojumi, jo poļi jau ``pēc savas dabas'' neesot spējīgi uz ko tādu. Šādu uzskatu aizstāvjos vēl lielāku sašutumu izraisīja no pašu poļu vidus nākusī kritika. Poļu valodā runājošie rakstnieki, publicisti, citi inteliģences pārstāvji nevarēja būt ``īsteni poļi'', ja tie ļaunprātīgi ``apmeloja'' poļu varonību un augstsirdību. Un visbiežāk šie poliski runājošie kritiķi izrādījās ebreji!

Autoram gan jāatzīmē, lielākā daļa vainas par dažādu nacionālo kopienu saskaņas trūkumu, sadarbības traucējumiem bija jāuzņemas nevis poļu tautai kā tādai, bet gan tās virsslāņiem. Nacionālo pretrunu saasināšanos lielā mērā noteica Polijas valdošo slāņu kļūdainā nacionālā politika.

Ar Polijas valsts jeb Otrās Žečpospolitas izveidi nepiepildījās cerības, ka pēc Krievijas un Austroungārijas impērijas sabrukuma varētu izveidoties brālīga poļu, lietuviešu, ukraiņu un baltkrievu tautu savienība. Kaut Žečpospolita bija daudznacionāla valsts, faktiski tā apzināti realizēja vienas nācijas valsts modeli. Poļu politiskā elite izturējās pret pārējām slāvu tautām ar pārākuma apziņu, tajās radušās nacionālās kustības uzskatīja par mākslīgi radītām un dzīvotnespējīgām. Polijas valdošie slāņi, stādot plānus par agrākās Žečpospolitas atjaunošanu, neņēma vērā, ka Polijas sadalītas eksistences pusotrā gadsimtā poļi bija lielā mērā zaudējuši savu ietekmi uz citām tautām, kuras, savukārt, bija zaudējušas ieradumu uzskatīt poļu virsvaru par pašsaprotamu un pakļauties tai. To atmiņā bija saglabājušās tikai sliktās atmiņās par Žečpospolitu ar tās nacionālajiem spaidiem. Arī atjaunotās Polijas varas iestādes pašas daudz paveica, lai ukraiņos, baltkrievos, arī citos mazākumos radītu negatīvu attieksmi pret poļiem un Polijas valsti, lai tie sevi uzskatītu par svešķermeni Polijā, kuram nebija vēlmju sadarboties ar poļu sabiedrību. Ukraiņu vēsturnieks J.~Daškevičs pirms Otrā pasaules kara pastāvošo Polijas valsts doktrīnu raksturoja kā klaji imperiālistisku (pārmetot šādas izpausmes arī mūsdienu Polijai, kas gan izsauca poļu vēsturnieku asus protestus).

Lai panāktu konsensu ar nacionālajiem mazākumiem, Polijai bija jāatsakās no sava nacionālā rakstura un jākļūst par dažādu nacionālu kolektīvu apvienību. Tas būtu iespējams, ja Polijas valsts būtu radusies kā brīvprātīga tautu savienība. Taču 1918.--1921.gadā ar spēku iegūtās robežas nedeva šādu iespēju. Kā raksta krievu vēsturniece L.~Miļakova ``Par anahronismu izrādījās cerības, ka pēc Krievijas impērijas un Austroungārijas sabrukuma radīsies ``brālīga Žečpospolita'' kā poļu, ukraiņu un lietuviešu tautu savienība''. 1918.~gadā katra no šīm tautām centās radīt savu nacionālu valsti un, protams, apvienot tajā maksimāli daudz savas tautas piederīgo, ar to nonākot konfliktā ar citām tautām. Ja šķiriskajos un nacionālajos konfliktos iesaistītās puses pilnībā ievērotu visu nāciju pašnoteikšanās principu, arī V.~Vilsona 14~punktu prasību, ka Polijai ``jāietver visas teritorijas ar neapstrīdami poļu iedzīvotājiem'', tad varētu nospraust taisnīgākas valstu robežas, panākt sabalansētāku situāciju rajonos ar jauktu nacionālo sastāvu.

Taču ir jāapzinās, ka šāda prasība bija nereāla apstākļos, kad katras tautas nacionālisti redzēja tikai savas tiesības un nevēlējās tādas pat tiesības saskatīt kaimiņtautām. Vācietis R.~Vingendorfs poļu nacionālisma saknes starpkaru periodā saskatīja ``viņu zemapziņā palikušajās bezspēcīgajās dusmās par savu apspiestību nebrīves gados'', kas pēc savas valsts atjaunošanas izpaudās kā citu tautu nacionālo jūtu un interešu neievērošana. Citādi viņš nevarēja izskaidrot, kā tauta, kura pēc vairāk nekā simtgades apspiestības spēja saglabāt savu nacionālo apziņu, pēc atbrīvošanās kļuva savukārt par citu tautu apspiedēju. ``Poļi nacionālās jūtas 1918.gadā bija un līdz šodienai ir no verdzības atbrīvotā jūtas, kura pirmā vajadzība ir tagad paverdzināt citus.'' Polijas vēsturi skatījušais mūsdienu krievu publicists D.~Žukovs raksta: ``Poļu nāciju raksturo tādas īpašības kā nacionālais egoisms un tieksme uz nacionālo maksimālismu. Šīs īpašības palīdzēja poļiem panākt Polijas atdzimšanu (kaut arī ne pilnībā Žečpospolitas robežās), pie tam nerēķinoties ar virknes citu kaimiņu tautu interesēm''. Pēc šī darba autora domām, abi tikko minētie autori pārlieku uzsver tikai psiholoģiskos momentus, piedēvējot to noteicošu iedarbību uz visiem poļu tautas pārstāvjiem, un reizē neņem vērā dažādu slāņu kopumā un to atsevišķu pārstāvju garīgās un materiālās intereses. Tā, baltkrievu, ukraiņi zemnieku apspiešanā galvenokārt bija ieinteresēta poļu muižniecība, kaut neliela daļa poļu inteliģences, strādniecības spēja pacelties pāri nacionālajiem aizspriedumiem, tāpēc diezin vai ir pareizi runāt par visas poļu nācijas ``zemapziņu'', ``nacionālo egoismu'' u.tml. Jārunā arī par atsevišķu iedzīvotāju slāņu, grupu interesēm. Piemēram, uzņēmēji, izmantojot valsts austrumu rajonu iedzīvotāju nabadzību, nemaksāja viņiem pat to samērā niecīgo darba algu, kādu saņēma strādnieki Polijas vidienē. Tā, ja 30.~gadu vidū Varšavā par darba stundu strādnieki vidēji saņēma 0,96~zlotu, tad Volīnijā šis atalgojums par darba stundu bija 0,39~zloti.

Vēsturnieki ir izteikuši arī pavisam citus vērtējumus par nacionālo konfliktu cēloņiem, nebūt nemeklējot tikai vairākuma nācijas atbildības pierādījumus. Tā, vācu vēsturnieks V.~Johmans nacionālo konfliktu saasināšanās cēloņus saskatīja Pirmā pasaules kara sekās un īpaši jaunas politiskas iekārtas izveidošanā Centrāl- un Austrumeiropā. Vainu par konfliktu izraisīšanos vēsturnieks pirmkārt meklēja vācu mazākuma vadītāju un vācu garīdzniecības nepārprotami izteiktajā pārliecībā par vācu kultūras un vācu valstiskuma pārākumu, kas izsauca citu tautu, arī poļu, atbildes reakciju~--- saasinātu nacionālismu.

Lai kā arī nebūtu, sapņi par visu politiķu saprātīgumu, nacionālo ambīciju ierobežošanu bija utopiski. Polija, tāpat kā tās kaimiņvalstis, nevarēja izveidoties kā monoetniska valsts.

Lai neradītu atkal un atkal no jauna uzliesmojošus nacionālus konfliktus, bija tikai viens ceļš --nacionālo mazākumu brīvas attīstības garantiju starptautiskas sistēmas radīšana. Tas arī bija iezīmēts Parīzes miera konferencē.

Jau teikts, ka Versaļas konference bija noteikusi pēc kara jaunizveidoto saistības nacionālo mazākumu aizsardzības jomā. Arī Polijas t.s. Marta konstitūcija noteica visu nāciju pārstāvju vienlīdzīgās tiesības valsts un pašvaldību amatu sadalē. Rīgas miera līgums (1921) paredzēja, ka Polija krievu, ukraiņu un baltkrievu tautības iedzīvotājiem, dzīvojošiem tās teritorijā, nodrošinās brīvas tiesības attīstīt savu kultūru, valodu, veikt reliģiskos rituālus. Polijas valsts iestādes visus šos lēmumus pilnīgi ignorēja.

Jau iepriekš pieminētais t.s. Mazais Versaļas līgums jeb Nacionālo mazākumu līgums Tautu Savienības vārdā deva tiesiskas garantijas mazākumtautību eksistencei un attīstībai. To parakstījušās valstis apņēmās piešķirt saviem nacionālajiem mazākumiem pilsonisko un politisko vienlīdzību, nodrošināt noteiktas nacionālās tiesības (lietot dzimto valodu oficiālajās iestādēs, tai skaitā tiesās, līdzekļu piešķiršanu skolu radīšanai ar dzimto mācību valodu u.c.) Ebrejiem bija paredzētas pat īpašas tiesības, kā iespēja ievērot sabatu (ivritā: \hetxti{שַׁבָּת}, \emph{šabbat}~--- sestdiena kā atpūtas un dievkalpojumu diena). Šai dienā ebrejus nedrīkstēja izsaukt uz tiesu vai kādu citu iestādi. Nacionālie mazākumi gan nevarēja tieši kā juridiskas personas griezties pie Tautu Savienības, bet tikai ar Tautu Savienības locekļu starpniecību, kas uzņēmās to interešu aizstāvību. Vācu minoritātes aizstāvību varēja uzņemties Vācija, bet ebrejiem, lietuviešiem, ukraiņiem un baltkrieviem šādas aizstāves nebija (PSRS Tautu Savienībā līdz 1934.~gada septembrim nesastāvēja). Tomēr no Polijas uz Tautu Savienību tika sūtīti simti petīciju par vācu, ebreju, ukraiņu un baltkrievu stāvokli tajā. Līdz 1924.~gadam no Polijas tika saņemtas 218~petīcijas. Tiesa, no tām tikai 11~tika izskatītas Tautu Savienības Padomē, kas nerunāja par labu mazākumtautu tiesību ievērošanai. Tikai no Augšsilēzijas vien laikā no 1920.~līdz 1931.~gadam Tautu Savienībā tika iesniegtas 64~petīcijas par vāciešu stāvokli. Radītajai nacionālo mazākumu tiesību aizstāvības sistēmai bija daudz trūkumu, gala rezultātā tā darbojās ar ļoti būtiskiem ierobežojumiem. Faktiski minoritātes stāvoklis bija atkarīgs no tās mītnes valsts labās gribas.

Pēc poļu diplomāta A.~Visocka atmiņām, J.~Pilsudskis bija visai kritiski noskaņots pret Tautu Savienību, kura traucēja Polijai realizēt visas tās vēlmes. Par Tautu Savienības un Polijas attiecībām 1933.~gada 21.~jūlijā viņš izteicās: ``Mums ir trīs jautājumi, kuri mūs nevieno, jo ietilpst tās [Tautu Savienības] kompetencē, konkrēti~--- Danciga, Augšsilēzija un nacionālās minoritātes. Šo vairāku gadu laikā mēs esam tomēr vairākkārt sev uzdevuši jautājumu, vai vispār no Tautu Savienības neizstāties, ja mēs pārliecinātos, ka tā var kļūt vai vēlēsies kļūt par mums naidīgu vai bīstamu instrumentu.'' Arī poļu vairākums šīs mazākumu tiesību garantijas, nepieciešamību ietvert tās arī 1921.~un 1935.~gada konstitūcijās uzņēma kā Polijas diskrimināciju, tiesību ierobežošanu tai pašai lemt iekšējos jautājumus.

Diemžēl arī poļu zinātnieki Tautu Savienības radīto nacionālo mazākumu tiesību aizsardzības sistēmu aplūkoja tikai no poļu nacionālā viedokļa un saskatīja tajā gandrīz vai vienīgi trūkumus. Kā pirmā jau 1919.~gadā par šo jautājumu iznāca poļu vēsturnieka S.~Kutšebas monogrāfija ``\pltxti{Kongres, traktat i Polska''} (Kongress, līgums un Polija), kura arī noteica poļu historiogrāfijas attieksmi pret ``mazo Versaļas līgumu'' uz gadu desmitiem. Zinātnieks vērtēja līgumu kā Polijas suverenitāti pārkāpjošu, īpaši kritizēja Lielbritāniju par tās nostāju Polijas robežu jautājumā, pie tam pats autors aizvēra acis uz Polijas pārmērīgajām prasībām. S.~Kutšeba pamatoti kritizēja līgumu par to, ka tas attiecās tikai uz dažām Centrālās un Austrumeiropas valstīm, bet nelika ievērot nacionālo mazākumu tiesības arī t.s. ``vecajām'' demokrātijām, taču tas viņam un viņa sekotājiem kalpoja par argumentu nevis cīņai par līguma nosacījumu izplatīšanu arī uz lielvalstīm, bet gan prasībai atcelt tā uzliktos pienākumus Polijai.

Polijas valsts nepildīja arī tās un Padomju Krievijas 1921.~gada miera līguma 7.~panta prasību, ka krievu, ukraiņu un baltkrievu tautības iedzīvotājiem, dzīvojošiem Polijas teritorijā, būs nodrošinātas brīvas tiesības attīstīt savu kultūru, valodu, veikt reliģiskos rituālus. PSRS nelaida garām izdevību uzstāties ukraiņu un baltkrievu interešu aizstāvja lomā. 1924.~gada 10.~maijā PSRS ārlietu tautas komisārs G.~Čičērins notā Polijas vēstniekam Maskavā norādīja, ka nacionālo mazākumu vajāšanas Polijā ir sasniegušas tik masveidīgu un sistemātisku raksturu, ka to prese ir uzstādījusi it kā paradoksālu prasību atcelt tos Polijas Konstitūcijas pantus, kuri it kā garantē nacionālajiem mazākumiem vienādas tiesības ar poļiem, jo tie nekad nav ievēroti un tikai maldina gan šos nacionālos mazākumus, gan sabiedrisko domu ārzemēs. Tālāk notā tika minēti piemēri par Polijas varas iestāžu varmācību pret baltkrievu un ukraiņu tautības piederīgajiem. Pēc 1922.~gada Seima vēlēšanām, kurās liela daļa austrumu nomaļu iedzīvotāju nepiedalījās, pret viņiem tika sūtītas armijas vienību soda ekspedīcijas, kuru dalībnieki piekāva, arestēja gan tos, kuri nepiedalījās vēlēšanās, gan arī tos, kuri balsoja par nacionālo mazākumu bloka pārstāvjiem. Notā bija minēti arī piemēri par slēgtajām nacionālo mazākumu skolām un baznīcām (kas gan no pret reliģiju karojošās PSRS valdības puses bija jāvērtē kā klaja liekulība), par poļu preses rīdīšanas kampaņām pret nacionālajiem mazākumiem. Tika norādīts, ka Polija pilnībā ignorē tās un Padomju Krievijas 1921.~gada miera līguma 7.~panta prasības. Polijas ārlietu ministrs M.~Zamoiskis noraidīja PSRS apvainojumus kā iejaukšanos Polijas iekšējās lietās, norādot, ka PSRS iedzīvotāji, kā piederoši nacionālajiem mazākumiem, tā arī savienoto republiku vairākumiem, nebauda ne pilsoniskās tiesības, ne reliģisko brīvību. Ministrs paziņoja, ka Polijas valdība turpmāk atstās bez atbildes Padomju Savienības pretenzijas sakarā ar 1921.~gada miera līguma 7.~pantu. Protams, Polijas ministrs ar pilnām tiesībām runāja par nedemokrātisko iekārtu PSRS, taču viņš nevarēja apgalvot, ka Polijā nacionālo mazākumu tiesības tiek ievērotas. Pēc šīs notu apmaiņas 1924.~gada jūnijā Latvijas sūtnis Polijā Mārtiņš Nukša uz Rīgu ziņoja: ``Padomju Krievija, atsaukdamās uz Rīgas līguma 7.~paragrāfu (pareizi: pantu -V.~Š.), bija atradusi par iespējamu iejaukties un maija sākumā iesniegusi Polijai protesta notu par minoritāšu skolu un laikrakstu slēgšanu, pareizticības vajāšanu, pierobežas joslas spaidu kolonizāciju utt. Ar to iesākās vairākkārtīga, diezgan asa notu apmaiņa, pie kam Polija savās notās deklarē, ka S.S.S.R. (pareizi: PSRS~--- V.Š.) iejaukšanos minoritāšu jautājumā uzskata kā netiesisku un atsaukšanos uz Rīgas līguma 7.pantu kā aplamu. Tonis apmainītās notās poļiem asāks, nekā krieviem\dots{}''.

Polijas nacionālās politikas trūkumus saskatīja ne tikai PSRS, bet arī citu valstu pārstāvji. Latvijas pārstāvniecības Polijā pagaidu vadītājs Pēteris~Oliņš Latvijas Ministru Prezidentam un ārlietu ministram Z.~Meierovicam 1921.~gada 1.~augustā ziņoja: ``\dots{}jākonstatē, ka vissekmīgāk nemieru ap sevi rada paši poļi ar savu fatālo kolonizācijas un polonizēšanas politiku baltkrievu un ukraiņu apgabalos''. Tāpat stāvokli vērtēja arī pragmātiski noskaņotie poļu politiķi. Tā 1922.~gada 5.~maijā Seimā deputāts ļudovietis Rudzinskis paziņoja, ka valsts austrumu nomalēs netiek ievērotas nekādas konstitucionālās normas, visu laiku rodas iespaids, ka šīs teritorijas vēl atrodas kara stāvoklī, ik uz soļa var sastapt vārda, preses un citu pilsonisko brīvību ierobežojumus. Arī kultūr-izglītības, nacionālā un agrārā politika noskaņo šo novadu iedzīvotājus dumpīgi pret Polijas valsti.

Kopš Polijas robežu izkarošanas laika visu starpkaru periodu nacionālie antagonismi bija pārsvarā pār tautu mierīgas sadzīvošanas tendencēm. Demokrātiskais tautu pašnoteikšanās princips Polijā drīzumā pārtapa par autoritāri-diktatoriskas varas pretenzijām. Polijas partijas no politiskā spektra centra līdz galēji labējam flangam vadījās no sakāpināta \strong{nacionālisma ideoloģijas}, kuras dabīga sastāvdaļa bija nacionālo mazākumu ierobežošana un to vajadzību neievērošana. Īpaši nacionālistiski, pat šovinistiski noskaņoti bija nacionāldemokrāti. Kā uzsvēris vācu vēsturnieks E.~Meijers, pēc viņu uzskatiem pret nacionālajiem mazākumiem vērstajai politikai bija jābūt tikai ar vienu mērķi~--- veicināt to iznīkšanu. To varēja panākt polonizācijas ceļā, vai arī, kas galvenokārt attiecās uz vāciešiem un ebrejiem, panākot viņu aizbraukšanu no valsts.

Lielākais politiskais spēks, kurš stāvēja radikāla nacionālisma pozīcijā, bija nacionālie demokrāti, kuri centās darboties visās poļu sociālajās grupās. Viņi atzina tikai ``etniski tīru'' valstu pastāvēšanu, tāpēc cīnījās par ap trīs miljoniem lielā ebreju mazākuma eksistences ekonomisko pamatu iznīcināšanu, viņu uzņēmumu boikotu, izglītības iespēju mazināšanu, provocēja grautiņus. Agrāk Prūsijai piederējušajā Polijas daļā viņi aktīvi kultivēja naidīgumu pret vietējiem vāciešiem. Slāvu mazākumus kā it kā nespējīgus radīt savas valstis endeki vēlējās asimilēt. Šie uzskati nacionālajā jomā nereti ``pielipa'' pat nacionāldemokrātu pretiniekiem citos politiskajos jautājumos.

Tomēr patiesības labad jāsaka, ka, tāpat kā gandrīz katrā strīdā, arī poļu un mazākumnāciju konfliktos vainīgie, atbildīgie par tiem bija abās pusēs. Taču no vairākumā esošā spēka, kura rokās atradās arī valsts vara, šī atbildība prasāma daudz lielākā mērā. Tā, ebreji, kuri vienmēr Polijā bija dzīvojuši kā nacionālais mazākums, sākotnēji poļu valstiskumu vērtēja cerīgi un izturējās lojāli, taču gadu gaitā ar antisemītisma pieaugumu Polijas augšslāņu politikā arī viņi mainīja savu nostāju. Valstī dzīvojošie lietuvieši vēlējās panākt Polijas sagrābtās Viļņas un tās apgabala atdošanu Lietuvai. Vācu iedzīvotāji, vērtējot dzīvi Polijā un Vācijā, gandrīz visu laiku, ar nelieliem izņēmuma periodiem, atradās opozīcijā Polijas valstij. Ukraiņi, zaudējuši cerības uz patstāvīgas valsts izveidi, prasīja vai nu viņu apdzīvoto apgabalu pievienošanu Ukrainas Padomju Sociālistiskajai republikai vai vismaz teritoriālu autonomiju Polijas sastāvā. No līdzīgas perspektīvas neatteiktos arī baltkrievi. Taču poļi šos mazākumtautu nacionālos centienus uztvēra kā viņu valstiskuma apdraudējumu un centās apspiest.

Jau parlamentārās demokrātijas periodā no 1921.~līdz 1926.~gadam Polijas nacionālistiskās valdības nebūt necentās ievērot t.s. Marta konstitūcijas 95.~un 96.~pantā noteiktās visu nāciju pārstāvju vienlīdzīgās tiesības valsts un pašvaldību amatu sadalē. Kaut nacionālo mazākumu bloks 1922.~gada vēlēšanās bija guvis ievērojamus panākumus, Volīnijas vojevodistē ieguvis gandrīz 80\%, bet Poļesjes, Viļņas un Novogrudokas vojevodistēs vairāk nekā 75\% balsu, jebkura nacionālo mazākumu iesaistīšana valdību koalīcijās tika izsludināta par nacionālo nodevību. Nevienā no 31 valdībām, kuras pastāvēja no 1918.~gada 14.~novembra līdz 1939.~gada 1.~septembrim un kurās kopā darbojās ap 350~ministru, nebija neviena nacionālo mazākumu pārstāvja. Nebija neviena apgabala vadītāja (\pltxti{wojewoda}) vai kopienas (\pltxti{gmina}) vadītāja (\pltxti{starosta})~--- nepoļa. Labējo partiju pārstāvji pat izteicās par to, lai valdības veidošanā, neuzticības votuma pieņemšanā, budžeta apstiprināšanā nacionālo mazākumu deputātu balsis vispār neņemtu vērā. Polijas nacionālo mazākumu politikas ``atslēgas vārds'' bija ``\pltxti{polytika kresowa}'', kas nozīmēja nolīdzināšanu un piespiedu asimilāciju. Valsts austrumos Polijas valdība pret nacionālajiem mazākumiem atjaunoja tos paņēmienus, kādus bijusī Krievijas cariskā valdība savā laikā lietoja pret poļiem. Kā cariskie ierēdņi realizēja rusifikāciju, tā arī jaunās valsts ierēdņi realizēja polonizāciju. Kaut valdība nacionālās politikas mērķus neuzdrošinājās skaidri formulēt, var teikt, ka faktiski pret vāciešiem Poznaņā un Pomerānijā tika realizēts to padzīšanas kurss, pret ukraiņiem un baltkrieviem~--- vismaz daļējas to asimilācijas virziens, attiecībās pret ebrejiem vispār nevarēja saskatīt kādu noteiktu mērķi.

30.~gados Polijas politika kļuva arvien nacionālistiskāka. Valdības aprindas, lai gan apkaroja nacionāldemokrātus, nacionālajā jautājumā realizēja daudzas to nostādnes. Polijas dažādo valdību pozīcija 30.~gados pēc poļu pētnieka A.~Čojnovska vērtējuma veica radikālu evolūciju ``līdz zemes polonizācijas plānam, tuvam nacionāldemokrātu koncepcijām''

Pēc tam, kad 1933.~gadā Vācija bija izstājusies no Tautu Savienības un 1934.~gadā, sakarā ar parakstīto vācu-poļu deklarāciju (par to sīkāk tiks runāts apakšnodaļā par ārpolitiku) Ā.~Hitlers uz laiku atteicās no vācu mazākuma atbalsta Polijā, 1934.~gada 13.~septembrī Polijas ārlietu ministrs J.~Beks Tautu Savienībā Ženēvā paziņoja par Polijas atteikšanos ievērot līgumu par nacionālo mazākumu aizsardzību. Faktiski tā bija tikai jau līdz tam agrāk realizētā kursa oficiāla atzīšana.

Krievu vēsturniece I.~Mihutina norādījusi, ka Polijas valdībai pirmkārt bija svarīga prestiža puse: attiekšanos no Tautu savienības kontroles par nacionālo mazākumu aizsardzības līguma izpildi tā uzskatīja kā formālu nolīdzināšanu tiesībās ar lielvalstīm, pār kurām nebija tādas kontroles. Poļu prese J.~Beka paziņojumu traktēja kā Polijas starptautiskās autoritātes pieauguma liecību. Jāatzīst, sakarā ar to, ka līguma saistības bija uzspiestas tikai jaunajām Austrumeiropas valstīm, Polijai bija zināmas morālas tiesības atteikties no to izpildīšanas, pirms tās nav uzņēmušās arī citas valstis, taču pragmātiski domājošie politiķi Polijas aktu uztvēra kā soli miera līgumu revīzijas virzienā, ko centās panākt agresīvās valstis. Arī pēc poļu vēsturnieka V.~Mihoviča domām, J.~Beka deklarācija, kaut tieši vērsta tikai pret minoritāšu līgumu, faktiski grāva visu Versaļas sistēmu kopumā. Līdz ar šo soli Polija pievienojās tām valstīm, kuras nepieņēma Versaļas sistēmu, kuras eksistence taču nodrošināja arī pašas Polijas pastāvēšanu. Latviešu vēsturnieks Aivars~Stranga, norādot, ka J.~Beka paziņojuma iegansts bija Padomju Savienības uzņemšana Tautu Savienībā bez jebkādām saistībām nacionālo mazākumu aizsardzības jomā, kā dziļāko tā cēloni minējis ``augošo Polijas lielummāniju, iedomājoties sevi par piekto lielvalsti Eiropā, kura var nerēķināties ar pagrimušo Tautu Savienību''. Domājams, ka blakus minētajiem faktoriem, Polijas valdības rīcību ietekmēja arī citi: vēlme izpatikt ietekmīgajām antisemītiski noskaņotajām poļu aprindām, atbrīvot sev rokas nacionālo mazākumu politikas jomā.

Pret valsts realizēto nacionālo polonizācijas politiku dažādo mazākumtautu intereses sakrita. Taču sākotnēji kāda kopēja darbības koncepcija, nerunājot jau par konkrētu programmu, nepastāvēja. Izejas punkts pirmajiem nacionālo mazākumu kontaktiem bija izpratne, ka vairāk nekā 30\% Polijas iedzīvotāju ir nepieciešama parlamentāra politiska pārstāvniecība, lai pretotos nacionālos mazākumus diskriminējošajam poļu nacionālismam. Parlamenta nobalsošanās mazākumu deputātu balsīm nereti bija izšķiroša nozīme. Tas izsauca poļu vairākuma vēlmi mainīt vēlēšanu kārtību par ļaunu nacionālo mazākumu tiesībām. Neraugoties uz poļu nacionālistu naidīgumu, pēc Seima un Senāta vēlēšanām 1922.~gada novembrī vāciešu, ebreju, ukraiņu un baltkrievu pārstāvji izveidoja nacionālo mazākumu bloku, kurš bija pietiekami nozīmīgs, lai pēc piecām vēlēšanu kārtām palīdzētu ievēlēt par prezidentu G.~Narutoviču. Taču sekojošā viņa kā poļu nacionālistu acīs Polijas ``kauna simbola'' nogalināšana un rīdīšanas kampaņas izvēršana pret mazākumiem, īpaši ebrejiem, visai uzskatāmi izgaismoja nacionālo mazākumu apdraudēto stāvokli.

Pēc šī nacionālisma uzliesmojuma uz vairākiem gadiem nacionālo attiecību jomā iestājās relatīvs miers, kaut nedraudzīgā nacionālā politika pret mazākumiem, īpaši antisemītisms, turpināja pastāvēt.

Tā Polijas prezidents S.~Vojcehovskis, apmeklējot valsts austrumu provinces, vienmēr uzsvēra, ka Polijas saimnieks ir tikai poļu tauta, poļu tautības iedzīvotāji un neviens cits. Arī citi augstu stāvoši politiķi un ierēdņi apgalvoja to pašu, taču parasti piemetināja, ka visiem nepoļu iedzīvotājiem Polija ar starptautiskiem līgumiem un savā Konstitūcijā solījusi nodrošināt tikai drošību, mierīgu un brīvu attīstību. Tā vesela trešdaļa valsts likuma priekšā līdztiesīgu pilsoņu tika pasludināti par savās tiesībās ierobežotiem salīdzinājumā ar pilntiesīgajiem poļiem.

Mazākumu bloks parlamentā pārtrauca darboties, taču pastāvēja iespēja tā atjaunošanai. Labējās poļu partijas aktīvi apkaroja ebrejus. Pat centriskā partija ``\pltxti{Piast}'' pret ebrejiem izturējās noraidoši.

Neilgi pirms 1926.~gada maija apvērsuma nu jau pēc ukraiņu pārstāvju iniciatīvas tika nodibināts ``Otrais nacionālo mazākumu bloks Polijā''. Ar tā palīdzību vācu, ukraiņu, baltkrievu un lietuviešu politiķi cerēja gūt panākumus parlamenta vēlēšanās. Atzīmējama šo mazākumu, īpaši ukraiņu, nacionālistu atturīgā attieksme pret ebrejiem. Tikai pakāpeniski jau pēc t.s. ``sanācijas'' režīma nodibināšanās ebreju parlamentārieši varēja iesaistīties mazākuma bloka darbībā. 1927.~gada oktobrī notika oficiāla bloka konstituēšanās, tā preses izdevums bija avīze ``\latxti{Natio}'' (Nācija, no latīņu \latxti{natio}~--- cilts, tauta), taču tā iznāca tikai īsu laiku.

Tā kā demokrātijas apstākļos nacionālie mazākumi nesagaidīja tolerantāku nacionālo politiku, tie 1926.~gadā, cerot uz attieksmes maiņu, sākotnēji atbalstīja 1926.~gada apvērsumu, taču vīlās. J.~Pilsudskis gan pasludināja liberalizāciju un solīja novērst nacionālo mazākumu, pirmkārt jau ebreju, saimniecisko diskrimināciju, taču radikāla valsts un nacionālajiem mazākumiem piederošo pilsoņu attiecību maiņa nenotika. Pirms apvērsuma pieņemtie diskriminējošie likumi palika spēkā. ``Sanācijas'' režīms ar nacionālajiem mazākumiem tikpat kā nesadarbojās. Parlamenta vēlēšanās 1928.~gada 4.~un 11.~martā par nederīgām izsludināto balsu vidū vairākumu sastādīja tās, kuras bija nodotas par nacionālajiem mazākumiem, tomēr kopumā vēlēšanas noritēja salīdzinoši korekti. No tajās dalību ņēmušajiem 78\% (Seima) un 64\% (Senāta) vēlētājiem nacionālie mazākumi saņēma 19,3\% balsu. Taču tas nevarēja mīkstināt zaudējumus nacionālo tiesību laukā.

1930.~gadā vēlēšanas Seimā notika jau tādos politiska terora apstākļos, kad par kādu mazākumtautību bloku nevarēja vairs ir domāt, mazākumu pārstāvji zaudēja daudz balsu. Diemžēl 30.~gados par nacionālo mazākumu kopdarbību savu tiesību aizstāvēšanā bija palikušas vienīgi atmiņas. Pasliktinājās attiecības starp mazākumiem un poļiem. Piemēram, 1933.~gadā reaģējot uz to, ka Vācijā pie varas bija nācis Ā.~Hitlers, poļu un ebreju aprindas Lodzā kopīgi demolēja vācu ģimnāziju, izdevniecību un avīzes ``\detxti{Freie Presse}'' (``Brīvā Prese'') redakciju. Pēc J.~Pilsudska nāves viņa epigonu valdīšanas laikā nacionālo mazākumu, īpaši ebreju, stāvoklis vēl vairāk pasliktinājās.

Neraugoties uz kopīgajām nelīdztiesības iezīmēm, atsevišķu nacionālo mazākumu stāvoklī Polijā \strong{bija} arī \strong{ievērojamas atšķirības}.

Īpaši naidīga attieksme Polijā valdīja pret \strong{krievu} kultūru. Pirms Pirmā pasaules kara Krievijas sastāvā esošajās poļu zemēs bija izvietoti ap 100~000 krievu karavīru un virsnieku. Lielu grupu veidoja krievu ierēdņi, tirgotāji, arī strādnieki. 1905.~gadā Varšavā dzīvoja apmēram 10\% krievvalodīgo. Pēc neatkarības atgūšanas Polijā palikušie pareizticīgie, kopā ar vecticībnieku kopienu sastādīja ap 100~līdz 140~tūkstošu lielu krievu kopienu. Viņiem bija savi deputāti, sabiedriskās organizācijas (piemēram, Krievijas tautas apvienība (\rutxti{Российское народное объединение}), krievu labdarības biedrības, kuras atbalstīja krievu kultūru un izglītību., Viņi izdeva savus periodiskos izdevumus (piemēram, avīzes ``Mūsu Laiks'' (``\rutxti{Наше время}''), ``Par Brīvību'' (``\rutxti{За свободу}''), politisku nedēļas izdevumu ``Zobens'' (``\rutxti{Меч}''), sieviešu un bērnu žurnāli ``Pavards'' (``\rutxti{Очаг}'') un ``Dzimtais Vārds'' (``\rutxti{Родное слово}''). 1936.~gadā Polijā darbojās 10~krievu skolas, tai skaitā arī ģimnāzija un licejs. Pastāvēja savas apvienības arī kultūras darbiniekiem: Krievu literātu un žurnālistu savienība (\rutxti{Союз русских литераторов и журналистов}) un Krievu dramatiskā studija (\rutxti{Русская драматическая студия}). Bijušie Krievijas armijas ģenerāļi sastādīja ne mazāk kā 40\% no neatkarīgās Polijas bruņoto spēku augstākā komandējošā sastāva. Krievijas augstskolu beidzēji ieņēma redzamu vietu Polijas inteliģences rindās.

Taču blakus šīm pozitīvajām parādībām eksistēja arī negatīvas. Pēc 1921.~gada 17.~marta miera līguma noslēgšanas sākās krievu kultūras iestāžu, skolu slēgšana, pareizticīgo baznīcu iznīcināšana. Poļu sabiedrības apziņā bija izveidojies pareizticīgo baznīcas kā Polijas valstiskumam naidīga krievu iestādījuma tēls. Tāpēc valdība uzskatīja, ka valsts attiecību noregulēšanai ar šo baznīcu jānoris kā ``apspiestības, kuru realizēja poļiem naidīgā Krievijas impērija, likvidēšanai''.

Piemērs tika parādīts ar svētā Aleksandra Ņevska katedrāles (\pltxti{Sobór św. Aleksandra Newskiego}) iznīcināšanu Varšavā Sakšu laukumā (no 1928.~līdz 1939.~gadam un tagad~--- J.~Pilsudska laukums). Tās celtniecības iniciators savā laikā bija Varšavas ģenerālgubernators feldmaršals J.~Gurko, kurš uzrakstīja speciālu lūgumu imperatoram Aleksandram~III, pamatojot to ar argumentu, ka Varšavas pareizticīgo dievnami var uzņemt tikai desmito daļu pilsētas pareizticīgo iedzīvotāju, kuru toreiz bija ap 43~000. Imperators piekrita katedrāles celtniecībai, kuru vienlaikus varētu apmeklēt 2~500 cilvēku. Naudu tai vāca visā Krievijā. Ziedoja cara ģimene, tirgotāji, strādnieki un zemnieki. Pēc izcilā krievu arhitekta L.~Benua projekta būvētajai katedrālei tika likti pamati 1894.~gadā un durvis tā vēra 1912.~gadā. Altāra gleznojumu un pārējo glezniecības darbu vadību uzticēja V.~Vasņecovam. Dievnamā bija ap 10~000 mākslas priekšmetu. Katedrāles augstums sasniedza 70~m un tā bija tai laikā pati augstākā ēka Varšavā. Tas īpaši tracināja poļus. J.~Pilsudskis 1894.~gadā paziņoja: ``Šī krāšņā celtne viņiem [krieviem] būs triumfa simbols, mums tā ir godu laupošas verdzības zīme''. Kad 1915.~gadā vācieši ieņēma pilsētu, poļi pārvērta pareizticīgo katedrāli par katoļu dievnamu. Taču jau 1919.~gada sākumā, trīs mēnešus pēc Polijas neatkarības iegūšanas, Varšavas maģistrāts nolēma likvidēt lielāko daļu pareizticīgo baznīcu Varšavā, arī minēto katedrāli. Galvenais motīvs bija tas, ka katedrāle, izvietota vienā no galvenajiem laukumiem Varšavā, atgādina par krievu kundzību Polijā. Vietējie ``mākslas zinātnieki'' katedrāli atzina par ``mazvērtīgu'', pie tam tā aizņemot pārāk daudz vietas. Katedrāle tika grauta vairāku gadu garumā ar ap 15~000 ``mazu'' sprādzienu. Pilnībā to sagrāva 1926.~gadā. Vēlāk daži mozaīku fragmenti izdaiļoja katolisko Marijas Magdalēnas baznīcu priekšpilsētā Prāgā, bet katedrāli rotājušās jašmas kolonas galu galā uzstādīja uz J.~Pilsudska kapa Krakovā. Citas vērtības, kas atradās katedrālē, tika izlaupītas. Kā iznīcināja citas, mazākas pareizticīgo baznīcas, aprakstījis notikumu laikabiedrs pareizticīgo baznīcas arhibīskaps Serafīms, kurš gan nebija notikumu aculiecinieks, bet izmantoja citu ticīgo stāstījumus: ``Pati baznīcu iznīcināšana notika vandaliski. Negaidot nakti, ložmetēju apsardzībā, ieradās policistu vienība, dažkārt kopā ar arestantiem (acīmredzot~--- kā darbaspēku~--- V.Š.) un līdz rītam baznīca tiek nojaukta līdz pamatiem. Pie tam ir bijuši gadījumi, kad notiek svēto relikviju un citu mūsu cienīto reliģisko priekšmetu zaimošana. Garīdzniekus un draudzes locekļus, kuri mēģina aizstāvēt svētumus, policija piekauj, arestē un nodod tiesāšanai. Visas sūdzības augstākajām varas iestādēm paliek bez atbildes.''

Kad Polijas puse 1922.~gadā pārmeta Padomju Krievijai katoļu dievnamu darbības traucēšanu, padomju puse ne tikai varēja atsaukties uz savu likumdošanu, kas paredzēja visu īpašumu, tai skaitā baznīcas, nacionalizāciju, bet to, kuri tiek izmantoti kulta vajadzībām, nodošanu draudžu lietošanā, taču arī norādīt uz daudziem faktiem, kad Polijā tika vajāta pareizticīgo baznīca.

Tā, 1921.~gada beigās no 38~Holmas apriņķa pareizticīgo draudzēm tikai 11 varēja netraucēti noturēt dievkalpojumus, tai pat apriņķī ar dinamītu tika uzspridzināta sv. Kirila un Metodija baznīca, pareizticīgo garīdzniekiem atņemti to īpašumi.

Pat Romas pāvests Benedikts XV bija spiests iejaukties un savā vēstījumā Polijas katoļu bīskapiem 1921.~gada 16.~jūnijā aicināt ``lai katrs garīdznieks neskopojas ar brāļa mīlestības izrādīšanu pret citiem garīdzniekiem, kaut arī tie atšķirtos politiskā, nacionālā vai reliģisko ritu ziņā.'' Taču ``brāļa mīlestības izrādīšanas'' vietā poļu garīdzniecība izvērsa cīņu pret pareizticību. Jau 20.~gados Romas katoļu bīskapijas iesniedza Polijas valsts tiesām 724~prasības pret pareizticīgo konsistorijām, pieprasot, atgriezt katoļu baznīcas gandrīz 700~pareizticīgo baznīcas, draudzes ēkas, baznīcas zemes. Oficiālais iegansts šīm prasībām bija fakts, ka agrāk (pirms trešās Polijas dalīšanas 1795.~gadā) šis īpašums vairākumā gadījumu vai nu piederēja Romas katoļu vai grieķu-katoļu (uniātu) baznīcu kopienām vai arī atradās to lietošanā. Parasti tiesas šīs prasības gan noraidīja.

1929.~gadā Polijas varas iestādes mēģināja veikt ``lieko'' pareizticīgo baznīcu nojaukšanas akciju. Plānots bija iznīcināt 97~tempļus, taču iedzīvotāju, garīdzniecības un Ukrainas parlamentāriešu pretestības dēļ sagrauti tika tikai 23.

1935.~gada beigās Varšavā tika izveidota nacionālo jautājumu komiteja, kurā tika pasludināts virziens uz pareizticīgo baznīcas polonizāciju. Tās pirmais darbības rezultāts bija 1936.~gada 17.~oktobrī Polijas Pareizticīgās Baznīcas Sinodei uzspiestais lēmums par atsevišķu liturģiju tulkošanu poļu valodā. Nākamais solis bija Polijas iestāžu lēmums slēgt Valsts pareizticīgo garīgo semināru Kremeņecā~--- vienīgo iestādi Volīnijā, kur gatavoja pareizticīgos priesterus. Tajā pašā laikā tika plānots atvērt valsts teoloģisko liceju Varšavā. Šie soļi izraisīja sašutumu ukraiņos, taču uzsāktās politikas atbalstītāji uzmanību protestiem nepievērsa.

Neatgriezenisks kaitējums pareizticīgajiem iedzīvotājiem Polijā tika nodarīts 30.~gadu beigās. 1938.~gadā sākās t.s. ``\pltxti{revindikācija}'' (\pltxti{rewindykacja/akcja rewindykacji cerkwi prawosławnych w II Rzeczypospolitej})~--- Polijas Pareizticīgās baznīcas mantas iznīcināšana vai nodošana Romas katoļu baznīcas vai vietējo varas iestāžu īpašumā, bet plašākā nozīmē~--- Ukrainas pareizticīgo iedzīvotāju ``atgriešana'' katolicismā. Lai īstenotu akciju, ar rakstiem presē par ``ukraiņu briesmām'', kur pareizticīgo baznīca tika apsūdzēta gan par poļu ukrainizāciju, gan rusifikāciju, tika aktīvi sagatavota Polijas sabiedriskā doma.

Īsā laikā pierobežas teritorijās (Helmas, Poļesjes, Volīnijas novados) izvērsās oficiāli nemotivēta, izpildes metožu ziņā brutāla pareizticīgo iedzīvotāju ``pievēršana'' katoļticībai un pareizticīgo baznīcu masveida iznīcināšana. Akciju veica galvenokārt robežsargi ar armijas augstākās pavēlniecības piekrišanu. Varas iestāžu slēgtie dievnami tika aizzīmogoti, priesterus sodīja par dievkalpojumu noturēšanu tajos. No pareizticīgo garīdzniekiem pieprasīja, lai viņi ticības mācības stundās dievvārdus skaidrotu poļu valodā. Paralēli izdarīja spiedienu uz iedzīvotājiem, kas atteicās pāriet katolicismā, pat draudot ar izsūtīšanu.

Ukraiņiem, kuri bija sašutuši par viņu konstitucionālo tiesību pārkāpšanu, iestādes paskaidroja savu rīcību ar to, ka tiktu atjaunots vēsturiskais taisnīgums. Tika apgalvots, ka viņu senči bijuši poļu kolonisti~--- katoļi (mazūri), bet cara valdība piespiedusi tos pāriet pareizticībā.

``Revindikācijas'' kampaņa izraisīja sašutumu Ukrainas iedzīvotāju plašu aprindu vidū. Uz Varšavu tika sūtītas draudžu locekļu un garīdznieku apelācijas ar protestiem pret baznīcu iznīcināšanu. Pret kampaņu protestēja arī poļu inteliģences pārstāvji. Kaut Romas katoļu baznīca neiebilda pret represijām, kas vērsās pret pareizticīgo baznīcu, pat tās garīdznieki ne vienmēr vēlējās sadarboties ar varas iestādēm, saprotot, ka tas skars baznīcas reputāciju, un pēc akcijas beigām norobežojās no tās, paziņojot, ka valdība to netika saskaņojusi ar katoļu baznīcu.

Pareizticīgo baznīcu masveida slēgšana un iznīcināšana Helmas un Poļesjes novados norisa 1938.~gada pavasarī-vasarā. Jūnijā-jūlijā šajos divos novados iznīcināja 112~baznīcas (citi dati~--- 127), tostarp 98, kuras darbojās jau pirms Brestas baznīcas ūnijas (1596) noslēgšanas. Kopumā šajās zemēs 1937.--1938.~gadā likvidēja līdz 200~pareizticīgo baznīcu. Kaut baznīcu iznīcināšana tika pakāpeniski pārtraukta, uzbrukumi pareizticībai turpinājās. Helmas un Poļesjes novadus sagaidīja tālāka polonizācija, ne tikai asimilējot ukraiņus, bet arī palielinot tajos poļu īpatsvaru.

Tikai Ļubļinas vojevodistes 1938.~gada vasarā teritorijā tika iznīcināta 91~pareizticīgo baznīca, 10~kapelas un 26~lūgšanu nami. Līdz 30.~gadu beigām 144~pareizticīgo baznīcas kļuva par katoļu dievnamiem. No tām, pēc vojevodistes vadības statistikas, 122~Romas-Katoļu baznīcas pakļautībā nokļuva ``stihiskas \pltxti{revindikācija}s ceļā, ko veica Romas-katoļticīgie iedzīvotāji vai garīdzniecība''.

Šādu faktu esamību konstatēja arī poļu vēsturnieki. Piemēram, Z.~Zaluskis rakstīja, ka tikai Volīnijā par katoļu dievnamiem tika pārvērstas 139~kādreiz uniātu, bet pēc tam vairāk nekā 100 gadu bijušas pareizticīgo baznīcas, iznīcinātas 189~pareizticīgo baznīcas, bet darboties palika tikai 151 pareizticīgo dievnams. Gribot negribot prātā nāk paralēles ar tikpat vandālisku baznīcu iznīcināšanu Padomju Savienībā.

Sagrauto svētnīcu vidū bija īstas arhitektūras un vēstures pērles. 2005.~gada septembrī poļu avīze ``\pltxti{Pseglend}'' rakstīja: ``Neraugoties uz to, ka 1938.~gada notikumiem bija noziegumu pret cilvēcību pazīmes, neviens Polijas pārstāvis līdz šim laikam nav atvainojies par nelikumīgo pareizticīgo svētnīcu iznīcināšanu starpkaru periodā un par represijām pret pareizticīgo konfesijas poļu pilsoņiem.''

Jāpasaka, ka, pēc daudzu ukraiņu pētnieku domām, šī ukraiņus apvainojošā Polijas valdošo aprindu realizētā kampaņa kļuva par vienu no iemesliem asiņainai konfrontācijai starp ukraiņiem un poļiem Volīnijā Otrajā pasaules karā un tā sauktajam ``Volīnijas slaktiņam'' 1943.~gadā.

Smagi cieta arī pareizticīgo kapi. Tā 1841.~gadā Varšavā Voļas (\pltxti{Wola}) rajonā oficiāli tika atklāti kapi, kur jau bija apbedīti krievu karavīri, krituši 1831.~gadā. Līdz Pirmajam pasaules karam tur apbedīja Varšavā mirušos: kā krievu aristokrātus un augstākos ierēdņus, tā arī trūcīgos. 20.~gados daudzi ``iekarotāju'' pieminekļi tika sagrauti. Varšavas prese ziņoja par drausmīga vandālisma ainām kapos. Par laimi, bija arī tādi cilvēki, kuri apzinājās to vēsturisko un kultūras vērtību, tāpēc kapi tomēr saglabājās. Mūsdienās tie, kaut arī daudz cietuši, ir viens no mazskaitlīgiem krievu nacionālā mazākuma kultūras pieminekļiem, saglabātiem no XIX gadsimta.

Jāievēro, ka starpkaru Polijā apmetās samērā daudz krievu ``balto'' emigrantu. Pēc ziņām no statistiskajām gada grāmatām, datiem, kurus savos ziņojumos minēja poļu administratīvās iestādes un specdienesti, Polijā dzīvoja vairāki desmiti tūkstošu (1922.~gadā~--- ap 50~tūkstošu, bet 1924.~gadā~--- 60~800) krievu emigrantu un bēgļu no boļševiku varas. Starp viņiem bija jau minētie B.~Savinkovs un S.~Bulak-Balahovičs, arī pazīstami rakstnieki kā D.~Filosofovs un M.~Arcibaševs, kurš, starp citu, sevi pasludināja par T.~Kostjuško mazmazdēlu.

Pret šiem bēgļiem Polijā valdīja divēja attieksme: no vienas puses tie bija iespējamie sabiedrotie cīņā pret boļševikiem konflikta gadījumā ar PSRS; no otras puses poļu acīs tie bija agrāk Poliju apspiedušās cariskās Krievijas vainas mantinieki, pret kuriem attieksme varēja būt naidīga, nievājoša, bet ne draudzīga. Poļu atmiņas par agrākajiem pāridarījumiem neveicināja kontaktu nodibināšanu ar ``balto'' emigrāciju. Krievu bēgļu lielākais vairums par poļu viesmīlības trūkuma pamatu uzskatīja ``imanento poļu rusofobiju''. Rīgas miera līgumu šie emigranti vērtēja kā divkārt nelikumīgu, jo: pirmkārt, tas nostiprināja citas valsts (Polijas) rokās agrāk Krievijai piederošās zemes; otrkārt, Polijas valsts to bija parakstījusi ar ``uzurpatoriem''~--- boļševikiem. Pēdējo paraksts zem līguma tika uzskatīts par Krievijas ``valsts nodevību'', Polijas paraksts~--- tās kā boļševiku līdzdalībnieces vainas pierādījums. ``Balto'' vairākums apmešanos Polijā uztvēra kā pagaidu etapu savā dzīvē, centās pārcelties uz valstīm (Vāciju, Franciju), kur koncentrējās ``baltās'' emigrācijas politiskā un intelektuālā darbība, tāpēc, mēģinot turpināt cīņu pret boļševikiem, maz vērības pievērsa Polijas likumiem un interesēm. Baltemigrantu radikālisms attiecībās pret PSRS, revizionisms pret Rīgas miera līgumu lika poļu varas iestādēm izturēties pret šīm grupām bez īpašas labvēlības. Polijas varas iestādes tādēļ ne vienu reizi vien bija spiestas izsūtīt krievu emigrantu (galvenokārt monarhistu) grupas no valsts.

Vairākumam emigrantu nebija iztikas līdzekļu, viņi vienkārši centās izdzīvot. Ap 40~000 strādāja tikai pagaidu darbu vai bija bezdarbnieki. Polijas pilsonības trūkums daudziem bija nepārvarams šķērslis darba meklējumos. Viņi arī bija pirmie, kurus atlaida no darba dažādās vietējās varas iestādēs. Tā, 1923.~gadā Varšavā no pilsētas pašvaldības iestādēm pirmkārt atlaida krievus bez Polijas pavalstniecības. Bēgļu no Krievijas iekārtošana darbā bija tieši atkarīga no varas realizētās politikas. Pat nelielais krievu ārstu, inženieru, arhitektu un citu speciālistu skaits, kuram izdevās daudz maz ciešami iekārtoties Polijā, cieta no diskriminācijas. Viņi visi pelnīja mazāk nekā poļu kolēģi, galvenokārt pilsonības trūkuma dēļ.

Tomēr jārunā arī par citu tendenci. Tie krievi, kuriem bija poliski skanošs uzvārds, varēja pretendēt uz pavalstniecības maiņu. Krievi varēja viegli pārvērsties par poļiem vai ukraiņiem, diskriminācijas vietā iegūstot priekšrocības. Vēl lielākas priekšrocības deva piederība katoļu konfesijai. Starpkaru Polijā daudzus ieceļotājus no austrumiem, kuri bieži pat nespēja brīvi runāt poliski, tā arī sauca par ``katoļiem''.

Tā kā kurss uz ``vienotas poļu nācijas'' izveidi nozīmēja kresu polonizāciju un to derusifikāciju, krievu atrašanās Polijas teritorijā bija nevēlama. Viņus dažādā ceļā centās izspiest no valsts. To pārstāvjiem netika atļauts veidot savas politiskās partijas, kuru programmās bija vērojama tieksme mainīt pastāvošo stāvokli, tika ierobežota pārvietošanās brīvība valsts austrumu rajonos. Polijas varas iestādes nerādīja ieinteresētību emigrantu sadzīves apstākļu uzlabošanā, kāda bija vērojama citās Eiropas valstīs. Gandrīz nebija dotāciju labdarības organizācijām, stipendiju studentiem un skolniekiem, atvieglojumu izglītības iestāžu organizēšanā. Bēgļi cieta ne tikai no materiālās nenodrošinātības, tos izsekoja poļu specdienesti. Tas viss noteica to, ka attiecības starp poļiem un krievu emigrantiem pārsvarā veidojās negatīvas. ``Derusifikācija'', kura poļu acīs tikai atbilda vēsturiskajām taisnīgumam, krievu emigrantu vidū tika uztverta kā netaisnīga, pazemojoša un no Polijas atsvešinoša. Tā rezultātā katra puse otras puses uzvedību uztvēra ar aizdomām, ar īpašu uzmanību. Pilnīgi attaisnojamās bez dzimtenes palikušo krievu bēgļu kopienas aktivitātēs poļu vērtētāji saskatīja ``rusifikātoriskas tendences''. Un otrādāk, gandrīz katru poļu iestāžu soli, kaut arī tas bija pamatots ar valstiskām interesēm, emigranti uztvēra kā ``poliskās rusofobijas'' izpausmi.

Aptverošu poļu varas iestāžu rīcības pret krievu mazākumu Polijā vērtējumu ir devis ``balto'' ģenerālis A.~Deņikins: ``Ja pāršķirstām dažas vēstures lappuses, visas pirmsrevolūcijas rusifikācijas pērles pilnībā nobāl nežēlīgo un mežonīgo polonizācijas spaidu priekšā, kas nospieda krievu zemes, kuras pēc Rīgas miera nonāca Polijā \citespace{} Poļi mums atmaksāja, var teikt, ar pārpalikumu!''

Vācijas sakāve karā, Polijas valsts izveide, citas pārmaiņas pasaulē, arī sava stāvokļa maiņa pārsteidza pilnīgi nesagatavotus daudzus agrākos Vācijas pavalstniekus. Protams, \strong{Poznaņas vācieši} un arī novada ebreju lielākā daļa, kuri sevi pēc valstiskās piederības uzskatīja par vāciešiem, tagad jutās apdraudēti. Šie slāņi, kad bija redzams, ka Poznaņas provinces iekļaušana Polijas valstī ir neizbēgama, lielākoties bija gatavi lojāli sadarboties ar poļiem. Tā, piemēram, ebreji sākotnēji bija noskaņoti ticēt poļu Tautas padomes solījumiem par ebreju vienlīdzību jaunajā valstī, tāpēc, kad jaundibināmajā Poznaņas universitātē tika ieviesti ``\latxti{Numerus clausus}'' (no latīņu~--- ``slēgtais skaits'', šeit~--- ebreju uzņemšanas ierobežojumi mācību iestādēs), tāpat kā tas bija cariskās Krievijas augtākajās un vidējās mācību iestādēs, vietējie ebreji bija ļoti vīlušies.

Jāņem vērā, ka pastāvēja objektīvi konfliktu cēloņi. Poļi pamatoti baidījās no iespējas nonākt ekonomiski spēcīgākās Vācijas atkarībā, no jau iepriekš lielāku turību ieguvušo Polijas vāciešu vadošā stāvokļa saglabāšanās un nostiprināšanās nu jau Polijas valsts ekonomikā. Tāpēc Polijas varas iestādes ar plašu poļu slāņu atbalstu centās ne tikai ekonomiski, bet arī administratīvi izspiest vāciešus no viņu pozīcijām. Poļi cerēja nomainīt vācu kapitālu, īpaši Augšsilēzijā, ar franču investīcijām. Zināmā mērā šis process arī norisa, taču nedeva cerētos panākumus. Poļu acīs Poznaņas vāciešu palikšana savās dzīves vietās varēja tikai stiprināt Vācijas revizionistiskos noskaņojumus, tāpēc viņi tika lielākoties uzskatīti par nevēlamiem. Poznaņā politiski lielāko ietekmi paturēja tradicionāli pretvāciski noskaņotie nacionāldemokrāti un kristīgie demokrāti, bet uz samierināšanos gatavajiem sociālistiem šeit bija maza ietekme. Par poļu nacionālistu attieksmi pret vāciešiem runā redzamā Krakovas jezuītu darbinieka J.~Rostvorovska 1927.~gadā publicēts raksts, kurā autors nožēloja, ka vietējie vācieši nevarot saprast, ka poļi pēc kara neko nav Vācijai atņēmuši, bet ``tikai atkal atguvuši niecīgu daļu no tā, kas bija vai nu viņu valstiskais netaisnīgi atņemtais īpašums vai poļu gadsimtiem ilgi pret ģermanizāciju aizstāvētā teritorija''. Pēc viņa viedokļa ``vācieši aizmirst, ka vācu mazākums, ko viņi aizstāv, nav vietējais, bet ieceļojis elements, vēl vairāk~--- iekarotāju viļņa un sveša karojoša nacionālisma sastāvdaļa.'' Ar šādu ``iejūtīgu'' argumentu poļu autors diezin vai spēja iekarot kāda vācieša sapratni.

Līdz 1920.~gadā Poznaņas apgabals tika oficiāli nodots Polijai, to pēc dažām ziņām jau atstāja 2/3 tur dzīvojušo vāciešu. Par pamatu tam bija Polijas pasākumi, pirmkārt jau vācu tautības ierēdņu atlaišana. Kopā ar piederīgajiem tie sastādīja ap 100~000. Protams, viņi meklēja iztiku Vācijā. Polija nedeva uzturēšanās atļaujas tiem vāciešiem, kuri apgabalā bija iebraukuši pēc 1908.~gada, kad tika pieņemts zemes atsavināšanas likums. Viņi gan varēja lūgt Polijas pilsonību, bet lielākoties lūgumi tika noraidīti. Pilsonību saņēma tikai tie vācieši, kuri pēc Versaļas miera Polijai piešķirtajos apgabalos bija apmetušies pirms 1908.~gada un pēctecīgi nodzīvojuši tur vismaz 12~gadus. Visi tie, kuri atbilstoši miera līgumam izvēlējās sev Vācijas pilsonību, tika piespiesti Polijas teritoriju atstāt, pat ja tā bija viņu dzimtene. Pie tam vācu īpašumi daļēji tika konfiscēti un noslēgtie ilglaicīgie nomas līgumi anulēti. Pēc toreizējās Polijas tautas skaitīšanas datiem 1921.~gadā Poznaņas apgabalā bija palikuši tikai 327~846 vācieši. 1926.~gadā pēc vācu pašu aprēķiniem tur bija palikuši vēl tikai 224~254 vācieši. Salīdzinot ar 1910.~gadu, viņu skaits bija mazinājies par 454~785 cilvēkiem, ar to viņu īpatsvars iedzīvotāju vidū krities no 38,7\% līdz 14,3\%.

Kopā no 1919.~līdz 1926.~gadam no Polijas rokās nonākušajiem Poznaņas~--- Pomerānijas (t.s. ``Poļu koridora'') apgabaliem izceļoja no 600~000~līdz 800~000 vāciešu. (Vācu vēsturniece A.~Šmidte-Roslere raksta, ka līdz 1925.~gadam bijušās Prūsijas provinces atstāja apmēram 800~000 vāciešu, to vidū arī ap 200~000 bijušo ierēdņu un karavīru tika izraidīti. L.~Dralle min skaitļus, ka no 1919.~līdz 1923.~gadam, pirmkārt no bijušajām Pozenes un Rietumprūsijas provincēm, izceļojuši ap 700~000 vāciešu. Ukraiņu vēsturnieks un publicists S.~Lozuņko pat apgalvo, ka 1919.--1925.~gados Poliju atstājuši vairāk nekā 1,25~miljoni vāciešu un ``naidu pret Poliju izjuta gandrīz vai visi vācieši''. (Acīmredzot, atspoguļojot vāciešu izciesto pārestību, šis vērtējums tomēr ir pārāk emocionāls.) Skaitļi nedaudz atšķiras, taču ilustrē vienu tendenci. Kā raksta vācu vēsturnieks H.~f.~Zitzevics, jau 1919.~gadā Polijas kultūras ministrs S.~Grabskis uzstādījis t.s. ``Pozenes programmu'' novada ātrai ``atvācošanai'' (``\detxti{Entdeutschung}''), un bieži bijis iespaids, ka Polijas valdība apzināti vai neapzināti kopēja to, ko pirms tam poļi pārmeta O.~f.~Bismarka ģermanizācijas programmai. Starpkaru perioda aculiecinieks R.~Vingendorfs apgalvoja, ka no 1921.~līdz 1939.~gadam visā Polija puse tās pilsoņu~--- vāciešu tika izstumta, bija spiesta atstāt valsti. Citi vācu autori atzīst, ka pret Polijas vāciešiem bija izvērsts ``saimniecisks un administratīvs spiediens'', taču izceļošana notika galvenokārt ``brīvas'' izvēles ceļā. No poļu puses kā bijušie ģermanizatori uzlūkotie vācieši cerēja, ka dzimtenē atradīs labākus dzīves apstākļus. Īpaši vāciešu skaits samazinājās pilsētās: Poznaņā un Bidgoščā. Pilsētas, kuras vēl 1918./1919.~gadā bija vāciskas, dažos gados ieguva polisku iedzīvotāju vairākumu. Saprotams, ka no Polijas Vācijā atgriezušos vāciešu negatīvā nostāja pret Poliju un poļiem ietekmēja pārējo vācu sabiedrību, sekmēja saspringtas attiecības starp Poliju un Vāciju.

Uz laukiem 1914.~gadā Poznaņas apgabalā vācu īpašnieku rokās bija gandrīz 1~000~000~ha zemes, bet pēc Vācijas pavalstnieku zemes īpašumu likvidācijas, pārceļošanas līgumu anulēšanas, agrārās reformas to rokās palika vairs tikai ap 300~000~ha. Mainījās vācu iedzīvotāju sociālā struktūra. Ja pirms 1914.~gada viņu vidū dominēja ierēdņi, amatnieki un tirgotāji, zemnieku un īpaši strādnieku bija samērā nedaudz, tad XX gadu otrajā pusē pārsvarā bija lauku iedzīvotāji. Ārstu, skolotāju un garīdznieku gan bija pietiekoši, bet ierēdņu vietās vācieši netika pielaisti. Iespēja paturēt zemes īpašumus kļuva par vāciešu izdzīvošanas jautājumu. Ja poļi agrārās reformas īstenošanā saskatīja tikai Vācijas ķeizarvalsts laikā realizētās agrārās politikas likumīgu labošanu, tad gan pēc emigrējušo, gan pēc palikušo vāciešu viedokļa viņiem bija nodarīta netaisnība. Taču gan uz Polijas kā valsts pastāvēšanu, gan Vācijas realizēto revizionistisko politiku Poznaņas vāciešu vidū vienota viedokļa nebija.

Polijas valsts realizētā nacionālā politika Poznaņā izpaudās arī kā asimilatoriski-poloniska citu tautu (vācu un ebreju) skolu skaita samazināšana. Vācu autori~--- laikabiedri poļu izglītības politiku tieši raksturoja nevis kā virzītu uz tautas izglītības līmeņa paaugstināšanu, bet gan uz iespējami ātrāku un pilnīgāku izglītības iestāžu, kuras darbojās citās nacionālajās valodās, likvidāciju.

Piemēram, ja 1919.~gadā pastāvēja 1~600 vācu skolas, tad 1928.~gadā to bija jau tikai 300. Bijušajās Prūsijas zemēs darbojās 1920.~gadā Polijā pieņemts likums, ka vismaz 40~vācu bērniem kopienā ir jāierīko sava vācu skola vai vācu klase poļu skolā. Ja divus gadus šis minimums netika sasniegts, vācu skolu vai klasi varēja slēgt. Taču straujā vāciešu izceļošana un varas iestāžu ierobežojošā prakse, kad skolu iecirkņi bieži tika tā noteikti, lai minimums nebūtu sasniedzams, veda pie drastiska skolu skaita samazināšanās. 1933/34.~mācību gadā Poznaņā bija vēl 84~vācu skolas un 118~klases poļu skolās. 1938.~gadā skolu skaits bija samazinājies līdz 60 un klašu~--- līdz 92. Šai gadā tikai 20\% vācvalodīgo bērnu varēja apmeklēt skolas dzimtajā valodā.

Arī valsts centrā (t.s. ``kongresa Polijā''), kur nekāda apdraudējuma valsts valodai, ko nacionālisti vienmēr izvirza par ieganstu mazākumtautību izglītības iestāžu vajāšanai, nebija, vācu skolu skaits starp 1919.~un 1930.~gadu samazinājās no 564 uz 63. Daļēji tās tika pārvērstas par vācu-poļu skolām vai arvien vairāk mācību priekšmetu tika mācīti poļu valodā. 1935./36.~mācību gadā pastāvēja vēl 11~skolas ar vācu kā mācību valodu, 46 ar poļu un vācu valodu un 189~skolas ar poļu mācību valodu un 2~vācu valodas mācību stundām nedēļā.

Īpaša problēma bija nākamās vācu skolotāju paaudzes sagatavošana. Pirms Polijas valsts radīšanas pastāvējušie skolotāju sagatavošanas semināri pakāpeniski tika slēgti, daudziem vācu skolotājiem, tā kā tiem nebija Polijas pilsonības vai viņu poļu valodas zināšanas bija nepietiekamas, tika atņemtas atļaujas strādāt savā profesijā. Ārpus Polijas izturētie eksāmeni netika atzīti, vācu skolotājiem bija jāiztur ``lojalitātes pārbaudījumi''. Kā kādreiz poļu patrioti radīja nelegālas poļu skolas, tā tagad Poznaņas apgabalā ceļojoši vācu skolotāji, kuriem nebija valsts izdotas atļaujas strādāt izglītības laukā, iesaistījās dzimtās valodas apmācībā. 1936.~gadā Poznaņā darbojās 32~ceļojoši skolotāji un 87 viņu palīdzes. Poļu iestādes centās šo darbu kavēt ar augstiem naudas sodiem un pat brīvības atņemšanu, taču bez panākumiem.

Sevišķas grūtības dzimtās valodas mācību organizēšana radīja lauku bērniem. Apmācības lielākoties varēja notikt tikai ziemas mēnešos, jo pārējos gadalaikos bērni bija nodarbināti lauksaimniecībā. Poznaņas lauku rajonos 1933.~gadā gandrīz 2/3~vācu kolonistu bērnu vecumā no desmit līdz piecpadsmit gadiem bija analfabēti.

Mazākumtautības un Polijas valsts skolu jautājumā turējās pie pilnīgi pretējiem viedokļiem. Mazākumtautības, baidoties no skolu kļūšanas par asimilācijas, pārtautošanas līdzekli citas nācijas rokās, centās panākt to iespējami lielāku brīvību no valsts ietekmes. Turpretī Polijas valsts skolās saskatīja instrumentu savas nacionālās politikas realizēšanai, centās nostādīt nacionālo mazākumu skolas iespējami stingrākā savā kontrolē, pēc iespējas ierobežojot privātskolu darbību. Polijas privātskolu likums 1932.~gadā deva tiesības valsts iestrādēm slēgt jebkuru privātskolu, kurās pēc šo iestāžu viedokļa mācības vai audzināšanas darbs tika īstenots ``nelojālā garā''. Tomēr kādu 1933.~gadā izdotu nolikumu, pēc kura vēsture un ģeogrāfija bija jāmāca tikai poļu valodā, pēc vācu sūdzības Tautu Savienībā nācās atcelt.

Saruka arī nacionālo mazākumu preses, savienību, biedrību darbības apjoms. Tāpat kā izglītībai, Polijas valdība lika dažādus traucēkļus dažādu vācu biedrību darbībai. Daudzas no tām formālu iemeslu dēļ tika aizliegtas. Par nozīmīgāko Poznaņas vācu kultūras organizāciju attīstījās Poznaņas Vēsturiskā biedrība (\detxti{Posener Historische Gesellschaft}), kura attīstīja plašu publicistisku darbību. Ietekmīgākie vācu preses izdevumi atradās Lodzā. 1938.~gadā šeit iznāca divdesmit periodiski izdevumi. Pretestību asimilācijai veicināja dažādas vācu jaunatnes organizācijas. 1926.~gadā ``Ceļotājputnu'' \detxti{(Wanderfogeln}) un ``Taku meklētāju'' (\detxti{Pfadfindern}) grupas apvienojās organizācijā ``Vācu jaunatne Polijā'' (\detxti{Deutschen Jugendschaft in Polen}). Lodzas apgabalā izveidojās Vācu Sociālistiskās strādnieku jaunatnes grupas. Vācu jaunieši organizējās arī dažādās vingrošanas un citās sporta biedrībās.

Lielu atbalstu vācu mazākumam Polijā sniedza Vācijas valsts un sabiedrība. Spēcīgas vācu etniskās grupas eksistence Poznaņā bija svarīgs t.s. Veimāras Republikas valdošo spēku arguments par labu Versaļas miera līgumā noteiktās robežas ar Poliju revīzijai. Lai kavētu vāciešu izceļošanu no Polijas un stiprinātu viņu pozīcijas, tiem tika organizēts masīvs finansiāls atbalsts. Taču vācu mēģinājumi ar vācu mazākuma tiesisku atbalstu Tautu Savienībā un ar no 1925.~līdz 1934.~gadam vesto saimniecisko karu Poliju saimnieciski vājināt un ievadīt tās atteikšanos no strīdīgajiem apgabaliem palika bez sekmēm. Otrādi~--- Poznaņas saimnieciskā norobežošana no tradicionālajiem lauksaimnieciskās produkcijas tirgiem Vācijā veda pie papildus grūtībām vietējiem vācu zemniekiem, poļu rūpju pieaugums par savu rietumu robežu drošību negatīvi ietekmēja arī vācu mazākuma stāvokli.

Kā atzīst mūsdienu poļu vēsturnieks J.~Holcers, Polijas rokās nonākušajā Silēzijas daļā vietējos iedzīvotājos Vācijas sakāves karā rezultātā simpātiju pieaugums par labu Polijai bija vērojams tikai īsu brīdi. Pirmie pēckara gadi nedāvāja ekonomisku uzplaukumu nevienai no abām valstīm. Taču labklājības līmenis augstāks bija Vācijā. Pārgāja arī silēziešu psiholoģiskais šoks par Vācijas sakāvi. Pieauga simpātijas Vācijas labā. Polijas piekritēju skaits gāja mazumā. 1930.~un 1932.~gadā saasinājās Polijas valsts attiecības ar vācu mazākumu Augšsilēzijā. J.~Pilsudskis tāpēc lika noturēt militārus manevrus pie Vācijas robežas. Vācijā tas izraisīja bailes no iespējama Polijas uzbrukuma.

Atšķirīgs no bijušajām Prūsijas provincēm bija vācu stāvoklis t.s. Viduspolijā jeb agrākajā ``kongresa Polijā'' un Volīnijā. Kā uzsver vācu autori, šo apgabalu vācieši Polijas valsts atjaunošanu fuztvēra kā vēsturiskā taisnīguma triumfu, to apsveica un bija gatavi lojāli pildīt pilsoņu pienākumus, protams, saņemot arī pilsoņu tiesības. Padomju-poļu karā, kad Sarkanā armija nonāca pie Varšavas, daudzi Viduspolijas vācieši brīvprātīgi pieteicās Polijas armijā, lai valsti glābtu no boļševikiem. 1919.~gadā Satversmes sapulcē tika ievēlēti arī 2~vācieši. Taču drīz pēc vēlēšanām Polijas varas iestādes Vācu savienību (\detxti{Deutsche Verein}) atlaida, vācieši palika bez savas politiskas organizācijas. Tomēr abi deputāti piedalījās 1919.~gada februārī t.s. ``mazās konstitūcijas'' pieņemšanā. 1921.~gada septembrī tika nodibināta ``Polijas vāciešu savienība'' (\detxti{Bund der Deutshen Polens}), kura izveidoja diezgan plašu vietējo organizāciju tīklu. Šai pat gadā tika nodibināt arī Polijas vācu strādnieku partija (\detxti{Deutsche Arbeiterparetei Polens}). Savienību varas iestādes gan atteicās reģistrēt, taču tās preses orgāns ``Brīvā Prese'' (\detxti{Freie Presse}) turpināja darboties bez traucējumiem.

Vācu pilsoniskie slāņi, pirmkārt jau Lodzas novadā, pēc dažiem neveiksmīgiem mēģinājumiem 1924.~gada jūnijā nodibināja Vācu tautas savienību Polijā (\detxti{Deutsche Volksverband in Polen}), kas sevi pozicionēja kā politisku partiju, kuras darbības lauks ir visa Polija, taču faktiski tā darbojās tikai bijušajā t.s. kongresa Polijā. 1930.~gadā tajā bija 305~vietējās grupas ar 24~000~biedru.

Pēc hitleriskā režīma nodibināšanās Vācijā vācu mazākums Polijā kļuva arvien atkarīgāks no tā atbalsta. Pēc Vācijas izstāšanās no Tautu Savienības un Vācijas un Polijas neuzbrukšanas līguma noslēgšanas 1934.~gadā vācu mazākumam Polijā vairs nebija Vācijas atbalsta Tautu Savienībā un tā vadītāji bija spiesti būt atturīgāki sūdzībās par nacionālo nevienlīdzību Polijā. Vienlaikus kontroli pār ārzemēs dzīvojošajiem tautiešiem centās iegūt Vācijas Nacionālsociālistiskā strādnieku partija (\detxti{Die Nationalsozialistische Deutsche Arbeiterpartei}). Nacionālsociālistiskās kustības panākumi Vācijā ietekmēja arī vāciešus Polijā. 1931.~gadā tika nodibināta Jaunvācu partija (\detxti{Jungdeutsche Partei}), kuras darbinieki izteica neapmierinātību ar agrākajiem konservatīvajiem etniskās grupas vadītājiem. Mērenie elementi gan pretojās nacionālradikāļu ietekmes izplatībai, 1934.~gadā apvienojās ``Polijas vāciešu padomē'' (\detxti{Rat der Deutschen in Polen}), kura īstenoja pret nacionālsociālistiem virzītu kursu, taču pakāpeniski zaudēja atbalstu tautiešos. Tomēr līdz pat Otrajam pasaules karam vācu mazākums Polijā palika politiski sadrumstalots. Vairākums vāciešu piederēja t.s. klusējošajam mazākumam. Polijas valdība 30.~gadu otrajā pusē uzturēja kontaktus ar abām grupām, cenšoties tos ``izspēlēt'' vienu pret otru. Kaut 1937.~gada novembrī Vācija un Polija pat noslēdza savstarpēju nacionālo mazākumu aizsardzības līgumu, kurā abas puses apsolīja nodrošināt attiecīgajam mazākumam (Vācija~--- poļu, Polija~--- vācu) brīvu valodas, kultūras, saimniecisku un reliģisku attīstību, taču praksē tas tā arī netika īstenots. Vācu mazākuma Polijā pretrunas ar valsts realizēto nacionālo politiku palika spēkā. Kā atzīmējis R.~Vingendorfs, visas Polijas valdības no 1920.~līdz 1939.~gadam bijušas vienotas, darbojoties ar mērķi iznīcināt vācu kultūras elementus un tāpēc vācieši tiem poļu nodarītos sarūgtinājumus nekad neaizmirsīšot. To liecina arī vācu organizāciju veikto tautiešu skaitīšanas rezultāti. Ja 1926.~gadā Polijas teritorijā sevi pie vāciešiem pieskaitīja 1~280~000 cilvēku, tad 1938./39.~gadā šis skaitlis bija sarucis līdz 1~030~000.

Atzīmējams gan, ka vācu mazākuma pārstāvju atmiņās atrodams daudz piemēru par labām attiecībām ar kaimiņiem~--- poļiem. Turpretī historiogrāfijā spekulatīvi pārvērtētas konfliktu situācijas. Taču nav arī noliedzams, ka Poznaņas apgabalā, Pomerānijā, bijušajā vācu Augšsilēzijā vismaz daļa vāciešu līdz 1939.~gadam nebija samierinājušies ar nacionālā mazākuma stāvokli un cerēja, ka kādreiz viņu dzimtene atkal ietilps Vācijas sastāvā. Pēc Polijas Iekšlietu ministrijas vērtējumiem 1938.~gadā, kad saasinājās sasprindzinājums starp abām valstīm~--- Vāciju un Poliju~--- ap 20\% vācu iedzīvotāju Polijā izturējās lojāli pret valsti, kuras pilsoņi viņi bija. Tātad ap 80\% vācu iedzīvotāju Iekšlietu ministrija par lojāliem neuzskatīja.

To, ka krievu un vāciešu stāvoklis Polijas valstī pasliktinājās, poļu nacionālisti varēja izskaidrot ar vēsturisko pagātni, poļiem nodarītajām pārestībām, bet to, kāpēc nepiepildījās citu, tāpat carisma apspiesto nacionālo mazākumu cerības pēc carisma likvidācijas baudīt līdztiesību, poļu vēsturniekiem arī mūsdienās izskaidrot ir grūti.

Poļu vidū bija plaši izplatīts naidīgums pret \strong{ebrejiem}. Pēdējais parādījās gan kā prasība pēc Polijas ebreju asimilācijas vai emigrācijas, gan kā boikota akcijas, citi ekscesi pret ebrejiem, taču būtiski atšķīrās no ``rasiskā antisemītisma'' (\detxti{rassisches Antisemitismus}) kaimiņzemē Vācijā pēc tam, kad tur pie varas nāca nacisti. Neraugoties uz apspiešanu, ebreji Polijā varēja atrast iespējas politiskai darbībai, nerunājot jau par kulturālās un reliģiskās dzīves izkopšanu. Latviešu valodā vispamatīgāko ebreju dzīves ainu Polijā ir atklājis profesors Aivars~Stranga. Viņa dati un vērtējumi izmantoti arī šajā grāmatā.

Ebreji Polijas valstī bija lielākais nacionālais mazākums (vairāk kā 8--9\%, 1931.gadā~--- 10\%) pēc ukraiņiem (ap 18\%). Atšķirībā no pārējām mazākumtautībām, ebreji dzīvoja visās Polijas daļās, galvenokārt pilsētās un miestiņos, bija pilsētu iedzīvotāji (\pltxti{naród miejski}). 1931.gadā viņi sastādīja 26,4\% Polijas pilsētu iedzīvotāju. Lielpilsētās ebreji sastādīja līdz trešdaļai iedzīvotāju. Varšavā (pēc absolūtā ebreju skaita Varšava atpalika vienīgi no Ņujorkas) un Lodzā 30\% iedzīvotāju bija ebreji, Ļvovā un Viļņā~--- pat 40\%. (Citi dati~--- 1938.~gadā Polijā mita 3,5~miljoni ebreju, kas sastādīja 9,8\% no visiem iedzīvotājiem. Ap 2,5~miljoni ebreju dzīvoja pilsētās, kur veidoja 27,3\% no pilsētu iedzīvotājiem, ap 1 miljons~--- laukos, sastādot 3,2\% visu lauku iedzīvotāju. Kopējā ebreju iedzīvotāju skaita ziņā Polija šai laikā atpalika vienīgi no ASV, kur bija apmetušies ap 4 miljoni ebreju). Polijas austrumos 20.~gados ebreji dažviet veidoja pat pilsētu iedzīvotāju vairākumu.

Lielā mērā poļu naidīguma pret ebrejiem saknes bija rodamas apstāklī, ka ap 40\% no Polijas buržuāzijas bija pēc tautības ebreji. Tas kalpoja par auglīgu augsni antisemītisma izplatībai poļu sabiedrībā.

Dominējošu lomu ebreju saimnieciskajā darbībā ieņēma tirdzniecība, amatniecība un vieglā rūpniecība. Ebreju sastāvā, no vienas puses, bija neliela mantīgā buržuāzija, bet no otras puses, pilsētu sīktirgotāju un amatnieku masa. 30.~gadu sākumā 42\% no viņiem bija nodarbināti rūpniecībā un amatniecībā un 37\% tirdzniecībā. (Ir arī citi dati, ka ar tirdzniecību nodarbojās 34,5\% visu darbaspējīgo ebreju iedzīvotāju, kad no pārējiem Polijas iedzīvotājiem šai sfērā bija nodarbināti tikai 1,5~līdz 2\%.) 6\% ebreju strādāja t.s. brīvajās profesijās, 6\%~--- pārvaldē un apkalpojošajā sfērā, 5\% bija nodarbināti transportā un sakaru nodrošināšanā. Toties lauksaimniecībā viņi bija pārstāvēti tikai ar 4\% nodarbināto. Tas, ka tirdzniecībā 62\%, rūpniecībā un amatniecībā 26\%, veselības apkopē 17\%, bet izglītībā 20\% visu strādājošo bija ebreji, izsauca poļu lielu neapmierinātību.

Salīdzinot ar ukraiņiem un baltkrieviem, ebreji bija vairāk gatavi sadarboties ar Polijas valsti, kaut ebreju jaunatnes vidū bija plaši izplatītas arī komunistiskās idejas. Ievērojot agrāko gadu pieredzi un antisemītiskās noskaņas poļu sabiedrībā, ebreju aprindas bija tās, kuras jau no paša valsts pastāvēšanas sākuma centās panākt nacionālo mazākumu tiesību nodrošināšanu. Viduspolijā, Galīcijā un Volīnijā, kur ebreji dzīvoja vairāk norobežoti, viņi arī vairāk cieta no antisemītisma.

Polijai jau XVIII gadsimtā nebija svešs antisemītisms, taču Žečpospolitas laikā ebreju stāvoklis salīdzinājumā ar to citās valstīs bija pietiekamu labs un daudzi no viņiem līdzās poļiem cīnījās pret Austrijas, Krievijas un Prūsijas varām, piedalījās poļu sacelšanās 1830.~un 1863.~gadā. Taču jau tad daudzviet poļu zemākajos slāņos pastāvēja naids gan pret turīgajiem ebrejiem~--- muižnieku īpašumu pārvaldniekiem un nomniekiem, tirgotājiem, kuri tika uzskatīti par poļu ekspluatatoriem, gan pret ebreju amatniekiem un sīktirgotājiem~--- poļu konkurentiem. Šo naidu bieži savtīgos nolūkos izmantoja virsslāņi, lai novirzītu tautas neapmierinātību no sevis pret šiem mierīgos apstākļos viņiem tik derīgajiem ebrejiem. Īpaši antisemītisms pieauga XIX gadsimta otrajā pusē, kad augot poļu pilsonībai, tai nācās sastapties ar jau izveidojušos ebreju uzņēmēju spēcīgu ekonomisko konkurenci. Augošais poļu nacionālisms blakus Austrijas, Krievijas un Prūsijas apspiedējvalstīm kā ienaidnieku uzlūkoja arī blakus dzīvojošos politiskās tiesības nebaudošos, bet saimnieciskajā dzīvē par poļiem veiklākos un izdarīgākos ebrejus. Kā uzsvēris profesors A.~Stranga, ``poļu agrīnajam antisemītismam bija skaidri izteikta ekonomiskā daba''. To iespaidoja arī lielas Polijas daļas atrašanās Krievijas sastāvā, kur jau pēc imperatora Aleksandra II noslepkavošanas izraisījās grautiņu vilnis. Poļu neapmierinātībai ar ebreju kaimiņiem pamatu radīja arī tas, ka daļa Polijas ebreju neprata poļu valodu un daļai no viņiem Pirmā pasaules kara gados bija provāciska (Poznaņā) vai prokrieviska (Galīcijā) nostāja. Taču antisemītisms pastāvēja ne tikai nosauktajos novados. Tā, jau 1919.~gadā tai poļu apdzīvotajā teritorijā, kura līdz tam atradās Austroungārijas sastāvā, tika atlaisti no darba visi ebreji~--- valsts dzelzceļa kalpotāji.

Jau Parīzes miera konferences laikā (1919) notika vairāki antisemītiski incidenti. Pinskā, Viļņā tika nogalināti vairāki desmiti ebreji, arī poļu zemnieki slepkavoja ebrejus. Tas bija pirmais impulss, kas lika ebreju starptautiskajām organizācijām pieprasīt garantijas mazākumtautībām un gala rezultātā noveda pie īpaša Polijas līguma ar Sabiedrotajām valstīm parakstīšanas par nacionālo, valodas un reliģiozo mazākumu aizsardzību, par kuru jau runāts iepriekšējā apakšnodaļā, un kurš poļu valdošajiem poļu politiķiem tik ļoti nepatika. Neidializējot līgumu, jākonstatē, ka tieši pašu poļu antisemītisma izpausmes bija to izsaukušas. 1920.~gada augustā kaujas par Varšavu laikā, kad tūkstošiem Polijas ebreju pieteicās brīvprātīgajos, lai karotu par dzimteni pret Sarkano armiju, daudzi no viņiem piedzīvoja pazemojošu neuzticību~--- 17~000 ebreju brīvprātīgo tika nevis nosūtīti uz fronti, bet internēti kā ``neuzticami'' elementi. Nodarīto netaisnību gan cietušie, gan ebreju kopiena kopumā neaizmirsa vēl ilgi.

Ebreju vēstures pētnieks S.~Ettingers konstatējis, ka pēc kara beigām Polijā ebreju stāvoklis daudzās jomās pat pasliktinājās.

Ebreju stāvoklī bija vērojamas atšķirības pēc Polijas agrākā dalījuma. Poznaņas apgabalā pastāvēja daudz draudzīgākas ebreju un vāciešu attiecības kā citur. XIX gadsimtā ebreji šeit bija gandrīz asimilējušies ar vāciešiem. Poznaņas ebrejiem, tāpat kā vāciešiem, kuri nevēlējās pakļauties polonizācijai, palika divas galvenās alternatīvas: vai nu izceļot uz Vāciju, vai sakļauties savā diasporā un mēģināt nosargāt savu identitāti pret jaunās Polijas valsts spiedienu. Vairums ebreju, kuri bija gan saglabājuši ebreju reliģiju un kultūru, bet reizē arī kļuvuši par vāciešiem, apgūstot to kultūru, izvēlējās pirmo ceļu, taču, kā zināms, tas nedeva vēlamo, jo pavisam drīz Vācijā pie varas nāca Ā.~Hitlers. Poznaņā vietējo ebreju izceļotāju vietā gan nāca ebreju ieceļotāji no Polijas austrumu apgabaliem, taču tie nespēja kompensēt izbraucēju skaitu. Ja 1910.~gadā provincē dzīvoja 26,5~tūkstoši ebreju (1,3\% visu iedzīvotāju), tad 1934.~gadā to bija tikai 9~000 (0,4\%).

Daļa Polijas ebreju~--- ap 10\% XX gadsimta 20.~gados jau bija asimilējušies. Viņi deva lielu ieguldījumu poļu kultūrā. Poļu valodā rakstīja un Polijas literatūras vēsturē iegāja ebreji B.~Šulcs, J.~Tuvims, J.~Vittlins. Taču daudzi poļi pat uz asimilētajiem ebrejiem raudzījās ar aizdomām, lika tiem šķēršļus dzīvē tāpat kā ortodoksālajiem ebrejiem. Spēcīgais antisemītisms Polijā diskreditēja asimilāciju Polijas ebreju vairākuma acīs. Asimilētie ebreji parasti bija ebreju lielburžuāzijas pārstāvji un ebreju intelektuāļi, bet asimilācija pilnīgi neskāra Polijas ebreju proletāriskos un sīkpilsoniskos slāņus.

Polijas saistības dot atbalstu nacionālajām skolām palika neizpildītas, ebreji neieguva kulturāli-nacionālo autonomiju. Tika ievērota vienīgi viņu reliģiskā autonomija, par kuru bija norūpējušies ortodoksāli ticīgie judaisti. Pastāvēja ebreju diskriminācija, ebrejus gandrīz nepieņēma valsts dienestā, izņemot zemākajos amatos. Ebreju zinātniekus, piemēram, Šī darba pirmajā nodaļā minēto izcilo vēsturnieku Š.~Aškenazi, antisemītiski noskaņotie kolēģi atteicās ievēlēt par profesoru Varšavas universitātē. Kāda Polijas ebreja (M.~Mazora) atmiņās lasāms: ``Polijā nebija neviena ebreja~--- policijas aģenta vai tiesu izpildītāja. Un ne tikai~--- pat durvju sarga vai tramvaja vadītāja darbs bija ebrejiem nepieejams.'' Ebreji nedrīkstēja dienēt kara aviācijā, flotē, sakaru dienestā, iestāties kara skolās.

Taču tai pat laikā ebreju dzīve starpkaru Polijā bija daudzšķautnaina. Šai laikā Polijā pastāvēja ebreju skolas un kultūras iestādes, iznāca ap 30~ebreju avīžu, darbojās vairāk nekā 20~ebreju teātri (1936.~gadā no 103~teātriem Polijā 15~bija ebreju teātri). Varšavā no 1928.~gada darbojās Ebreju zinātņu institūts, kurš bija gan pētnieciska iestāde, gan pirmā ebreju augstskola valstī. Viļņā 1925.~gadā tika izveidots Jidiša kultūras un valodas pētniecības institūts.

Tūlīt pēc Pirmā pasaules kara beigām atjaunojās ebreju politiskā aktivitāte, kura nenorima līdz pat Otrajam pasaules karam. Jau 1918.~gada decembrī cionisti un reliģiskā partija ``\emph{Mizrahi}'' (saīsinājums no ivrita \emph{Merkaz ruhani}~--- garīgais centrs) organizēja Varšavā ``Pagaidu nacionālo centru'' un nacionālu ebreju frakciju, kura cīnījās pret valdības mēģinājumiem ierobežot nacionālo mazākumu pārstāvniecību ar nelabvēlīgu tiem vēlēšanu likuma pieņemšanu. Partija sadarbojās ar citiem nacionālajiem mazākumiem.

1925.~gadā, valdība, cenšoties sašķelt nacionālo mazākumu bloku Seimā, piedāvāja ebrejiem kompromisu~--- atzīt ebreju kopienu autonomiju, atcelt ierobežojumus ivrita un jidiša lietošanai, oficiāli atzīt un subsidēt ebreju skolas, garantēt sestdienas atpūtas ievērošanu utt. Taču, plaši izreklamējusi šos priekšlikumus ārzemēs, valdība neko nedarīja, lai tos īstenotu. Ebreju skolas nesaņēma finansiālu valsts atbalstu, to beigšanas apliecības netika valstiski atzītas.

Pieaugošās nacionālās nevienlīdzības un nodokļu sloga iespaidā daudzi ebreji emigrēja. Daļa pārcēlās uz Palestīnu, kur šis emigrācijas vilnis pēc izglītības un reliģijas lietu ministra (1923 un 1925--1926) S.~Grabska (viņa brālis~--- Vladislavs Grabskis tai laikā (1920, 1923--1925) bija premjerministrs) uzvārda tika saukts par ``Grabska emigrāciju''.

Pēc ``sanācijas'' režīma nodibināšanas, kurš vismaz vārdos nebija tik antisemītisks kā nacionāldemokrātu politiķi, ebreju vienotība mazinājās. Vismaz sākotnēji daudzi Polijas ebreji sāka cerīgāk raudzīties nākotnē, ticot, ka J.~Pilsudskis ierobežos viņa politiskā pretinieka R.~Dmovska kultivēto antisemītismu. Ebreju cerības daļēji attaisnojās. Vēl no cariskās Krievijas mantotās ebrejus diskriminējošās normas tika atceltas, 1928.~gada Seima vēlēšanās dažas ebreju partijas uzstājās ar savu kandidātu sarakstu, kurš gan guva ievērojami mazāk deputātu vietu nekā iepriekšējās vēlēšanās. Toties ebreju uzvārdi parādījās valdības sarakstā, kas bija bezprecedenta gadījums Polijas politikā. Polijas ebreju virsotne sadalījās valdības atbalstītājos vai vismaz neitrāli pret to noskaņotajos un opozicionāros. Daļa ebreju partiju sāka īstenot sadarbības taktiku ar ``sanācijas'' režīmu. Turpretī Bunda piekritēji un cionistu lielākā daļa noraidīja J.~Pilsudska autoritārismu. Ebreju līdzšinējās relatīvās vienotības izjukšana noveda pie viņu kopējās ietekmes mazināšanās.

Visa starpkaru periodā Polijas valsts realizētās, un pēc sanācijas režīma nodibināšanas pastiprinātās valsts intervences ekonomiskās politikas rezultāts bija sīko ebreju tirgotāju un amatnieku slīgšana nabadzībā. Tiem būtisku kaitējumu nodarīja valsts uzņēmējdarbības attīstība. Polijas valdība ne tikai ieviesa valsts monopolu uz tabaku, alkoholiskajiem dzērieniem, sāli un sērkociņiem, bet tās rokās atradās arī raktuves, mežrūpniecība, ķīmiskā ražošana, uzņēmumi, kuri nodarbojās ar armijas saimniecisko apgādi utt. 1925.~gadā valsts uzņēmumu strādnieki sastādīja 15--16\% visu pilsētu iedzīvotāju. 30.~gadu vidū valdības rokās koncentrējās ap 25\% rūpniecības un tirdzniecības. Vissmagāko triecienu ebreju uzņēmēji un tirgotāji saņēma spirta ražošanā un tirdzniecībā. Šai jomā 1933.~gadā viņu loma bija kritusi no agrākajiem 60\% uz niecīgajiem 2\%. Šī etatisma (no franču \frtxti{État}~— valsts~--- ideoloģija, kura absolutizē valsts lomu, tās iejaukšanos visās sabiedriskās dzīves jomās, arī ekonomikā) politika veda pie ebreju izspiešanas no viņu ekonomiskajām pozīcijām. Poļu vēsturnieks K.~Kavaļecs konstatē, ka poļu vidū izplatīto antiebrejisko noskaņojumu pieauguma pamatā bija ne tikai antisemītisms ar savu mitoloģiju, bet arī ekonomiska argumentācija.

Kā norādījis vācu vēsturnieks F.~Golževskis, par ilgstošajām saimnieciskajām grūtībām labējās, ieskaitot valdošās, poļu aprindas, kā arī daļa ļudoviešu labprāt vainoja mazākumtautības, īpaši ebrejus. Poļu vēsturnieks J.~Tomaševskis ir uzsvēris: ``Visas \citespace{} nacionālistisko aprindu koncepcijas, kuras vainagoja tēze, ka valsts ekonomiskā stāvokļa uzlabošana ir iespējama tikai pēc atbrīvošanās no ``svešiem elementiem'', bija absurds. Īstā problēma bija darba vietu trūkums, ko izsauca kapitāla trūkums un lauku vajadzības, un šiem faktoriem nebija nekā kopīga ar Polijas Republikas iedzīvotāju struktūru.''

Augot valsts uzņēmumu svaram rūpniecībā un tirdzniecībā, kas bija galvenais ebreju iztikas avots, viņu īpatsvars šais jomās ātri saruka. Taču ebreju ekonomiskā iespaida samazināšanās neradīja, kā to it kā varēja paredzēt, antisemītisma vājināšanos, bet, gluži otrādi, tā pieaugumu.

Kā uzsvēris A.~Stranga, jau tā saimnieciski vājajā Polijā ar atpalikušu sociālo struktūru, kur amatnieku un sīktirgotāju skaits bija nesamērīgi liels, vietējo ebreju sociālā struktūra bija vēl arhaiskāka. Ebreju amatnieku, bodnieku lielais skaits bija pretrunā ar ekonomikas modernizācijas prasībām. Tas arī bija galvenais cēlonis plašu ebreju masu nabadzībai. Taču paši viņi savās nedienās vainoja pret viņiem vērsto valsts ekonomisko politiku. Valsts intervenci ekonomikā un pret ebrejiem vērstos soļus viņi nosauca par ``balto boļševismu'' (\pltxti{biały bolszewizm}).

Blakus poļu antisemītiem pret ebrejiem ļoti aktīvi vērsās ukraiņu antisemīti Galīcijā. Pasaules ekonomiskajai krīzei sākoties saasinājās ukraiņu zemnieku nabadzība. Kā viens tās pārvarēšanas līdzeklis tika uzlūkota ukraiņu kooperatīvu radīšana, kuri centās izspiest no tirdzniecības ebrejus. Izvērsās atturības kustība ar nolūku izspiest ebrejus no alkohola tirdzniecības. Taču ar to cīņa pret ebrejiem neaprobežojās. Ebreji veikalnieku boikotu nomainīja fiziski uzbrukumi viņiem un viņu mantai, kurus provocēja Ukraiņu nacionālistu organizācija (\uktxti{Організація Українських Націоналістів}) jeb saīsināti OUN. Tā rezultātā daudzi ebreji bija spiesti pamest lauku rajonus un doties uz pilsētām. No antisemītu puses skanēja prasības realizēt ``pilsētu tīrīšanu'' no ebrejiem, to galīgu izspiešanu no saimnieciskās dzīves.

Jau pirms J.~Pilsudska nāves valdība pastiprināja finansiālo un ekonomisko spiedienu uz ebreju uzņēmējiem, samazinājās ebreju politiskais iespaids Polijas parlamentā, vājinājās ebreju starptautisko organizāciju iespējas ietekmēt Polijas valdības nostāju. Kaut 1935.~gadā pieņemtās jaunās t.s. Aprīļa Polijas Konstitūcijas 7.~pants aizliedza etnisko un reliģisko mazākumu diskrimināciju, tā izpilde netika nodrošināta. Sākot ar 1935.--1936.~gadu, jaunai ekonomiskai krīzei skarot Poliju, kas tikko bija atkopusies no 1929.--1932.~gada pasaules krīzes, pieauga ekonomiskais antisemītismas.

Saimniecības ``polonizācijas'' gaitā 30.~gadu otrajā pusē parādījās antireklāmas uzraksti ``Poļi, neiepērkaties pie ebrejiem''. 1936.~gada jūnijā premjerministrs F.~Slavojs Skladkovskis Seimā pasludināja ekonomisko cīņu (\pltxti{walka ekonomiczna}) pret ebrejiem. Augustā sekoja valdības rīkojums veikalu īpašniekiem izkārt pie veikala savu uzvārdu, rakstītu tā, kā tas bija darīts dzimšanas apliecībā. Rīkojuma mērķis bija darīt ebrejus atpazīstamus un veicināt viņu uzņēmumu boikotu. Arī katoļu baznīca faktiski atbalstīja ebreju ekonomisko boikotu. Labējā prese un sabiedriskā doma asi nosodīja tos poļus, kuri nepiedalījās boikotā. 1936.--1937.~gadā tā rezultātā strauji samazinājās ebrejiem piederošo veikalu skaits.

Pēc J.~Pilsudska nāves valdības piekritēji, kuri agrāk pretstatīja savu politiku no varas atbīdīto nacionāldemokrātu nostājai, faktiski tagad pārņēma no tiem daudzas idejas nacionālajā jautājumā, arī daļu no antisemītiskajiem lozungiem, tiesa~--- mērenākā formā. Tagad ebreju jautājums tika pilnīgi atdalīts no visu minoritāšu jautājuma; aizstāvot poļu nācijas konsolidācijas mērķi, sanācijas darboņi to definēja stingri etniski un konfesionāli, ebrejus atstājot ārpus tās. Lai citējam kāda ebreja (Z.~Zimaka) atmiņas par dzīvi Polijā sanācijas periodā: ``Ebreji bija nostādīti uz poļu sabiedrības zemākā pakāpiena. No rītiem atskanēja avīžu iznēsātāju balsis~--- ``Jauns līdzeklis pret utīm un žīdiem!''.

1937.~gada sākumā tika izveidota jauna masu organizācija ``Tautas Vienības Nometne'' (``\pltxti{Oboz Zjednoczenia Narodowego}'' jeb OZON), kura pasludināja ``jauno nacionālismu''~--- poļu interešu prioritāti savijumā ar antisemītismu. Šī ``Nometne'' bija pirmā valdības izveidota organizācija, kurā bija aizliegts uzņemt ebrejus, pat tos, kuri sen bija pieņēmuši kristīgo ticību un arī poļu kultūru. Strauji pieauga to profesionālo organizāciju skaits, kuras savos statūtos iekļāva paragrāfu ar aizliegumu uzņemt par to biedriem ebrejus. Arī profesionālās ārstu, juristu un inženieru savienības pārņēma \latxti{numerus clausus} (no latīņu~--- skaitlisks ierobežojums, procentu norma) un izslēdza no savām rindām ebrejus.

Jau 1935.~gadā valdības kultivētais poļu nacionālistiskais noskaņojums noveda pie tā, ka liela daļa studentu iestājās par atsevišķu t.s. ``geto solu'' (``\pltxti{geto lawkowe}'') ierīkošanu ebreju studentiem. (Jāatzīmē, ka tai pat 1935.~gadā Vācijā stājās spēkā bēdīgi slavenie ``Nirnbergas likumi'' (\detxti{Reichsbürgergesetz, Gesetz zum Schutze des deutschen Blutes und der deutschen Ehre}) pret ebrejiem.) Uzbrukumi ebreju studentiem ar lozungu ``Ebreji~--- uz Palestīnu'' aptvēra daudzas Polijas augstskolas. 1937.~gadā, kad universitātēs pieauga pret ebreju studentiem vērstu incidentu skaits, izglītības ministrs atļāva iekārtot šos t.s. ``geto solus''. Ebreji nedrīkstēja sēsties solos, kuri bija paredzēti ``āriešiem''.

Poļu studentu organizācijas iestājās par \latxti{numerus clausus} principa ieviešanu augstskolās, lai samazinātu ebreju studentu īpatsvaru līdz ebreju iedzīvotāju īpatsvaram valstī. Tam par iemeslu bija tas, ka ebreji sastādīja līdz 25\% no studentu skaita, kad viņu īpatsvars visu iedzīvotāju vidū bija tikai ap 8--9\%. Rezultātā 30.~gadu otrajā pusē Polijas augstskolās, tāpat kā agrāk cariskajā Krievijā pret ebrejiem tika realizēta \latxti{numerus clausus} sistēma. Ja 1928./1929.~gadā ebreju jauniešu universitātes vēl bija 20,4\%, tad 1937./1938.gadā to īpatsvars bija nospiests uz 10\% (citi dati: 1937.gadā~--- uz 7,5\%), kas~--- neievērojot šīs etniskās grupas īpašo struktūru~--- aptuveni atbilda ebreju īpatsvaram iedzīvotāju vidū. Atsevišķas fakultātes vispār atteicās uzņemt ebreju studentus.

Gods kam gods, daļa poļu inteliģences pārstāvju mēģināja aizstāvēt demokrātijas un tolerances principus. Vairāk nekā 50~poļu universitāšu mācību spēku publicēja deklarāciju ``\pltxti{Wobronie kultury polskiej}'' (``Aizsargājot poļu kultūru''), uzsverot Polijai vēsturiski raksturīgo toleranci pret ebrejiem. Taču šovinistiski noskaņotie poļu studenti pielietoja vardarbību, lai izspiestu ebrejus no ``āriešu'' soliem. Kur antisemītu izlēcieni saņēma pretsparu, notika asiņaini kautiņi. Kur aktīvai pretošanās kustībai trūka spēku, ebreju studenti atteicās sēsties šais viņiem ``atvēlētajos'' solos, ilgas lekciju stundas pavadīja stāvot kājās.

Diemžēl arī Romas katoļu baznīca Polijā pastiprināja savu tradicionālo reliģisko antisemītismu. Polijas baznīcas galva kardināls A.~Hlonds gan nosodīja fiziskus uzbrukumus ebrejiem, taču aicināja uz morālu cīņu pret tiem~--- jo ebreji esot ateisma un komunisma iemiesojums, turpinot no poļu-padomju kara laika pastāvošā stereotipa par īpaša ``ebreju komunisma'' (\pltxti{žydokomuna}) uzturēšanu. Viņš arī atbalstīja ebreju ekonomisko boikotu. 1937.~gadā Polijas bīskapu sinode pieprasīja ebreju bērnu segregāciju (no latīņu \latxti{segregatio}~--- atdalīšana, piespiedu kārtā īstenota dažādu rasu pārstāvju nodalīšana). Kā uzsvēris A.~Stranga, pēdējo pirmskara gadu laikā liela daļa poļu nācijas pieņēma pozīciju, ka Polija pieder poļiem un katoļiem. Lielas problēmas tas radīja arī asimilētajiem ebrejiem, pret viņiem kā ``neīstiem'' poļiem negatīvi izturējās šovinistiskie poļi, kā savas tautas ``nodevējiem''~--- ticību un ebreju pašapziņu saglabājušie tautieši. Piemēram, visasākie antisemītu uzbrukumi vērsās pret vienu no lielākajiem poļu dzejniekiem (ar ebreju izcelsmi) J.~Tuvimu par ``ebrejisku dzeju poļu valodā''.

Antisemītistiskos incidentos 1935.--1937.gadā tika nogalināti 79 un ievainoti 500~ebreji. 1938.~gada decembrī 117~Polijas Seima deputātu, tātad~--- vairāk nekā ceturtā daļa no visiem, parakstīja iesniegumu ar prasību pēc iespējas ātrāk atrisināt ``ebreju jautājumu''. Ārlietu ministrija gatavoja plašas ebreju emigrācijas programmu, kurai tikai formāli būtu jānes brīvprātīgas kustības vārds, bet pēc būtības tā būtu deportācija. 1936.~gadā ārlietu ministrs J.~Beks Tautu Savienībā publiskoja ideju par nepieciešamību nodrošināt katru gadu vismaz 80~000~līdz 100~000 ebreju izbraukšanu no Polijas, ko finansiāli būtu jānodrošina Tautu Savienībai vai arī Lielbritānijai (Tās rokās bija Palestīnas pārvalde). Lai atbrīvotos no, kā antisemītiem šķita, ``svešķermeņa'' Polijā, viņi atbalstīja arī cionistu centienus organizēt ebreju izbraukšanu uz Palestīnu. Radās arī cits~--- plāns pārcelt ebrejus uz Madagaskaru. Ar Francijas aizjūras teritoriju ministra M.~Motes atbalstu, kurš izteicās par salas nodošanu Polijas aizbildniecībā, Polijā tika radīta komisija, lai izpētītu iespēju piemērot salu ebreju no Polijas pārceļošanai. 1937.~gada maijā komisija devās uz Madagaskaru, sabija tur divarpus mēnešus un atzina salas ziemeļu daļu par piemērotu kolonizācijai. Viens no Otro pasaules karu pārdzīvojušajiem Polijas ebrejiem M.~Mazors 1952.--1954.gadā tapušās atmiņās rakstīja, ka Polijā 30.~gados bija ļoti populārs romāns par to, kā ebrejus izsūta uz Madagaskaru, un no poļu mutes ebreji bieži varēja dzirdēt ko līdzīgu: ``Kad tad jūs beidzot aizbrauksiet uz Madagaskaru?'' Ideju par Polijas ebreju pārvietošanu uz Madagaskaru Polijas ārlietu ministrs J.~Beks apsprieda ar savu Francijas kolēģi I.~Delbosu, kurš izrādīja atturīgu izpratni par Polijas valdības centieniem. Uz Francijas rēķina apmierināmas Polijas pretenzijas atzina Itālijas ārlietu ministrs G.~Čiano. Idejas par ebreju deportāciju uz Madagaskaru piekritējs bija arī Ā.~Hitlers. Ar viņu 1938.~gadā par šo tēmu sarunājās Polijas sūtnis J.~Lipskis, kurš, ziņojot par to ārlietu ministram J.~Bekam, rakstīja, ka par jautājuma atrisināšanu solījis ``fīreram'', ka ``mēs par to uzcelsim viņam greznu pieminekli Varšavā''.

Taču ar Otrā pasaules kara izraisīšanu saistītie notikumi neļāva projektam attīstīties tālāk. Nākt pretī ``Polijas interesēm'' bija gatava arī Lielbritānija, kura piekrita atļaut ik gadu pārcelties uz Palestīnu 50~tūkstošiem Polijas ebreju. Polijas valdība plānoja, ka viņi būtu spiesti atstāt lielāko daļu savas mantas Polijā, apņemtos apgādāt Poliju ar izejvielām un citām precēm no savām paredzamajām mītnes zemēm. Atzīmējams, ka pret ebreju izraidīšanu iestājās tikai PPS un nedaudz vēlāk arī nelielā 1939.~gada aprīlī dibinātā Demokrātiskā partija (\pltxti{Stronnitctwo Demokratyczne}). Aizsteidzoties notikumiem priekšā kā likteņa ironiju var pieminēt to, ka pēc Polijas sakāves 1939.~gada septembrī Vācija Polijas teritoriju pārvērta par ebreju koncentrēšanas un iznīcināšanas vietu.

Tika veidota pret ebrejiem vērstu likumu sistēma. 1939.~gada februārī Polijas Ārlietu ministrijas vadības uzdevumā Konsulārā departamenta Emigrācijas politikas nodaļa sagatavoja ebreju emigrācijas veicināšanas programmu, kurā bija apskatītas divas jautājuma atrisināšanas iespējas un to sekas. Viena iespēja bija vienpusīga akcija, Polijai ``uzsākot karu ar [domāts~--- pret~--- V.Š.] ebrejiem gan savā valstī, gan arī pasaulē''. Taču šis ceļš varēja izraisīt Polijas finansiālu un tirdzniecisku boikotu no Rietumu demokrātisko valstu puses. Otra iespēja tika nosaukta par ``konstruktīvu sadarbības ceļu ar ebrejiem'', kaut tās pamatus pēc poļu ierēdņu domām bija jāveido tādiem principiem, kurus par ``sadarbību'' varētu nosaukt tikai vislielākais ciniķis un demagogs. Lūk, šie ``principi'':

\begin{enumerate}
\item ebreji apņemas panākt pusotra miljona ebreju izbraukšanu no Polijas 30~gadu laikā;

\item ebreji piešķir Polijai ilgtermiņa kredītus 300--500~miljonu dolāru apjomā;

\item Polijas valsts garantē Polijas ebrejiem politisko un saimniecisko līdztiesību, t.i. faktisku privilēģiju piešķiršanu.
\end{enumerate}

Tādejādi ebrejiem par padzīšanu no savas dzimtās zemes būtu jāsamaksā miljoniem dolāru, bet līdztiesības nodrošināšana tika nosaukta par ``privilēģiju piešķiršanu''. Lai panāktu ebreju kopienas piekrišanu šiem principiem, tika ieteikts izveidot likumdošanu, kuras ceļā Polijas ebrejiem jebkurā brīdī varētu atņemt viņu ``priviliģēto stāvokli''. Šis ``likumdošanas instruments'' tika raksturots arī sīkāk: ``Pretebreju likumu mērķim jābūt ne tikai formālai, bet arī faktiskai ebreju izstumšanai no mūsu politiskās un saimnieciskās dzīves. Vienlaikus likumiem ir jāatvieglo mums emigrācijas politikas realizāciju''. Likumdošanas ceļā bija paredzēts revidēt Polijas pilsonības likumu, (ir aprēķini, ka ar likumu bija jāatņem pilsonība ap pusmiljonam Polijas ebreju) ieviest likumus par aizliegumu strādāt virknē profesiju, citās savukārt ieviest ``proporcionalitātes principu'', kā arī pieņemt likumu par ``ebreju emigrācijas nodokli''. Profesors A.~Stranga pamatoti atzinis, ka šis ``emigracionisms bija aroganta un nereālistiska pieeja ``ebreju jautājuma'' risinājumam: emigrācija no Polijas bija neliela, ebreju absolūtais vairākums vēlējās palikt un dzīvot savā dzimtenē.'' Šī darba autoram vēl jāpiezīmē, ka, kaut ebreju izsūtīšanās programma arī palika vienīgi projekta līmenī, nevar viennozīmīgi pateikt, ka tas notika pateicoties valdības saprātīgumam, kaut kādam humānisma līmenim, zem kura tā nevarēja nolaisties. Iespējams, ka tikai tāpēc, ka Polijas vadībai nācās risināt aktuālās attiecībās ar Vāciju, pret Polijas ebrejiem jau 1939.gadā netika vērsts klaji antisemītisks ``likumdošanas instruments'', kura darbības rezultātā to stāvoklis būtiski tuvinātos ebreju stāvoklim nacistiskajā Vācijā. Galējā ebreju izstumšanas programma gan tika noraidīta, taču valsts ``poliskā rakstura nostiprināšanas'' kampaņa turpinājās.

Visbīstamākā bija pieaugošā grautiņu kustība. 30.~gadu sākumā nacistu panākumu Vācijā ietekmē aktivizējās arī antisemītiskie elementi Polijā. Ebrejiem naidīgais noskaņojums, kuru uzkurināja ne tikai valdībai opozīcijā esošā Nacionāldemokrātiskā partija, bet arī daļa katoliskās baznīcas kalpotāju un arī daļa valdības aprindu, sasniedza vandalisku uzbrukumu, grautiņu un citu ekscesu pakāpi, kuru rezultātā bija arī slepkavības. Sākot ar 1933.~gadu kļuva arvien biežāki stihiski vai arī labējo ekstrēmistu organizēti fiziski uzbrukumi ebrejiem. No 1933.~gada maija līdz 1936.~gadam ap 1~300~ebreju bija guvuši savainojumus antisemītisko uzbrukumu laikā. No 1935.~gada vidus masveida uzbrukumi ebrejiem Polijā kļuva par izplatītu parādību.

Pēc poļu avīžu ziņām 1935.--1936.~gadā ebreju grautiņi bija atzīmēti 150 pilsētās, to upuru skaits sniedzās simtos. Tā, 1936.~gada martā nežēlīgs grautiņš norisa Kelces vojevodistes Pšitekas (\pltxti{Pshitek)} pilsētā. Pēc citām ziņām no 1935.~līdz 1937.~gadam Polijā notika 16~grautiņi, kuros dzīvību zaudēja 118 ebreju. Apkārtnes iedzīvotāji ne tikai izlaupīja ebreju bodītes, bet arī ielauzās to mājas un nogalināja saimniekus. 1937.--1938.~gadā grautiņu un uzbrukumu vilnis vēl pieauga.

Poļu nacionālo godu glāba Varšavas strādnieki, kuri 1938.~gadā sāka iekļauties daudznacionālās vienībās kopā ar ebrejiem, lai aizsargātu viņus no vardarbības.

Polijas valdošajos slāņos valdošo antisemītismu īpaši izgaismo pirmskara notikumi, kad ebreju jautājums 30.~gadu beigās apgrūtināja arī Polijas un Vācijas attiecības. Desmiti tūkstošu Polijas ebreju, kuri līdz tam mierīgi dzīvoja un strādāja Vācijā, tagad nacistu pieaugošo vajāšanu rezultātā centās atgriezties savā dzimtenē. Taču tai pašā laikā valdība Varšavā, daļēji kooperējoties ar Polijas cionistu organizācijām, centās intensificēt Polijas ebreju emigrāciju uz Palestīnu. Pēc Polijas valdības viedokļa, ebreji, kuri no Vācijas atgriezās dzimtenē, bieži bez pietiekamām poļu valodas zināšanām, tikai saasinātu etnisko sasprindzinājumu un palielinātu bezdarbnieku jau tā ievērojamo skaitu.

1938.~gadā Polijas valdība atteicās uzņemt valstī 15~000--20~000 Polijas izcelsmes ebrejus, kurus Vācija izraidīja no tai ar ``anšlusu'' (vācu \detxti{Anschluss}~--- pievienošana, aneksija) pievienotās (1938.~gada marts) Austrijas teritorijas. Lai pēc Austrijas ``anšlusa'' pieaugošo ebreju straumi ierobežotu, Polijas ārlietu ministrijas Konsulārā nodaļa nolēma ārzemēs dzīvojošajiem Polijas pilsoņiem iebraukšanas atļauju valstī piešķirt tikai izņēmuma gadījumos. 1938.~gada aprīļa sākumā tika pieņemts likums, pēc kura Polijas pilsoņiem, dzīvojošiem ārzemēs, divu nedēļu laikā bija jāpārreģistrē savas pases. Oktobra sākumā tika izdoti atbilstoši pasu noteikumi. Pēc tiem katram ārzemēs dzīvojošam Polijas pavalstniekam, kas vairāk par pieciem gadiem dzīvoja ārzemēs, varēja atņemt pilsonību vai vismaz liegt iebraukšanas atļauju Polijā. Likuma labojums galvenokārt bija vērsts pret Polijas pilsoņiem~--- ebrejiem, kuri, dzīvojot Vācijā un, izjūtot arvien lielākus pret sevi vērstus spaidus, varēja izlemt atgriezties Polijā. Tagad šāda iespēja tiem tika liegta. (Lai nerastos iespaids, ka Polija viena rīkojās nehumāni, jāpasaka, ka arī pārējās Eiropas valstis 30.~gados slēdza savas robežas no nacionālsociālistiskās Vācijas bēgošajiem ebrejiem. 1938.~gada vasarā notika 32~pasaules valstu konference, kurā apsprieda ebreju emigrācijas jautājumu. Pat tās valstis, kuras šai laikā uzņēma emigrantus, atteicās ielaist ebrejus. Notika tas bez kāda izteikta antisemītisma. Galvenais arguments atteikumam bija tas, ka valstis nevar uzņemties finanšu slogu par ebreju iekārtošanu jaunajās dzīves vietās, un jau tā tajās esot intelektuāļu un komersantu pārprodukcija.)

Taču Polijas rīcība īpaši uztrauca Vācijas valdību, jo tā bažījās, ka desmiti tūkstošu Polijas ebreju, kuri uzturējās Vācijā, tur vēl ilgstoši būs jāpiecieš, kad Vācijas ebrejiem jau bija atņemts īpašums, visas tiesības un pret viņiem tika izdarīts spiediens, lai panāktu to izceļošanu.

Pirms pārreģistrācijas termiņa beigām nacistiskās Vācijas varas iestādes aizturēja ap 17~tūkstošu Polijas ebreju, kuri lielākoties jau otrajā vai pat trešajā paaudzē dzīvoja Vācijā, un, ļaujot ņemt līdzi tikai pašas nepieciešamākās sadzīves lietas, steidzamā kārtā deportēja uz Polijas robežu un piespieda to šķērsot. Taču Polijas varas iestādes atteicās viņus ielaist valstī. Poļu robežsargi dzina šos pēc likuma Polijas pilsoņus atpakaļ uz Vāciju. Uz robežas notika no cilvēcības viedokļa pretīgas ainas. Trīs dienas cilvēki atradās zem klajas debess šaurajā neitrālajā joslā.

Galu galā Polija ar noteikumu, ka vairāk deportāciju nebūs, tomēr bija spiesta savus pilsoņus uzņemt, bet izvietoja tos nometnēs un daudzu mēnešu garumā viņi tur dzīvoja smagos apstākļos. Tūkstoši cilvēku 1938./1939.~gada ziemas mēnešus pavadīja, dzīvojot barakās necilvēcīgos apstākļos, pilnīgā neziņā par to, kas tos sagaida nākotnē. Pārtika viņi tikai no Polijas Sarkanā krusta un ebreju labdarības organizāciju ziedojumiem.

Starp deportētajiem atradās 17.~gadīga ebreju jaunekļa H.~Grinšpana, kurš nelegāli, bez attiecīgiem dokumentiem uzturējās Parīzē, vecāki, kā arī brālis un māsa. Uzzinājis par radinieku likteni, H.~Grinšpans 1938.~gada 7.~novembrī nopirka revolveri, ieradās Vācijas sūtniecībā un lūdza tikšanos ar tās darbinieku. Kad pie viņa iznāca sūtniecības trešais sekretārs E.~Fon~Rats, H.~Grinšpans viņam raidīja piecus šāvienus vēderā, no kuriem E.~f.~Rats pēc divām dienām nomira. Viņa bērēs Diseldorfā piedalījās Ā.~Hitlers un J.~f.Ribentrops. Atentāts kalpoja kā iegansts t.s. ``Kristāla naktij'' (vācu \detxti{Kristallnacht}~--- nosaukums cēlies no sadauzīto ebreju veikalu skatlogu stikliem), kad no 9.~uz 10.~novembri Vācijā tika sarīkoti ebreju grautiņi, kuros nogalināja ap 90~cilvēku, simtus sakropļoja, ap 30~000 ebreju arestēja un nosūtīja uz koncentrācijas nometnēm, nodarīja milzīgus materiālus zaudējumus. (H.~Grinšpanu gatavojās tiesāt, pēc Vācijas iebrukuma Francijā viņš nonāca vācu koncentrācijas nometnē un viņa liktenis pēc 1942.~gada nav zināms.)

Tā šajos notikumos savu attieksmi pret ebrejiem demonstrēja ne tikai nacistiskā Vācija, bet arī ``sanācijas'' režīms Polijā. Tā politika bija virzīta, kā raksta poļu vēsturnieks J.~Tomaševskis ``pret cilvēkiem \citespace{}, kuri vadoties no rasistiskiem motīviem tika vajāti Trešajā reihā [t.i.~--- nacistiskajā Vācijā]. Polijas pilsoņi, kuri atradās īpaši apdraudētā stāvoklī, zaudēja Polijas valsts atbalstu kā reiz brīdī, kad tas viņiem bija visvairāk vajadzīgs.''

Latvijas sūtnis Polijā Ludvigs Ēķis 1939.~gada janvārī Latvijas Ārlietu ministrijai ziņoja: ``Kādus 7~000 poļu žīdu Vācija esot uzgrūdusi ``poļiem uz kakla''. Vācijā vēl tomēr palikuši ap 150~000 tā saucamo Polijas žīdu, kuri tagad esot zaudējuši Polijas pilsoņu tiesības \citespace{} Tie, bez šaubām arī Vācijā nevarēšot palikt, bet Vācijai šie 150~tūkstoši poļu žīdu būšot jānotransportē kaut kur citur. Poļu-vācu sarunas žīdu jautājumā, kuras turpinoties ar pārtraukumiem jau ilgu laiku, būšot turpināties droši vien līdz tam laikam, kamēr šie 150~000 tagadējo žīdu--bezpavalstnieku (līdzšinējie poļu žīdi) būšot atrast sev citur mājas.'' Tagad mēs zinām, ka vairākums no viņiem ``mājas'' atrada viņsaulē. Aizraidīja viņus turp nacisti, bet arī Polijas ``sanācijas'' režīmam jānes daļa atbildības par šo Polijas pilsoņu traģisko likteni. Tiesa, tobrīd vēl neviens, arī Polijas valdība un poļu antisemīti, nespēja paredzēt Holokaustu, un nedrīkst pielīdzināt to vainu ebreju priekšā nacistu vēlākajiem noziegumiem.

Jāuzsver, ka Polijas ebrejiem valsts neatkarības gadi starp abiem pasaules kariem bija pēdējais relatīva miera periods pirms nākamajām kara un Holokausta šausmām. Ebreju skolu un kultūras iestāžu, ebreju avīžu un teātru skaits Polijā nekad ne pirms, ne pēc tam nebija tik liels kā šajā laikā. Paliek spēkā fakts, ka antisemītisms Polijā līdz Otrajam pasaules karam nesasniedza tik augstu pakāpi kā kaimiņzemē Vācijā. Polijā arī nekad neizveidojās tik masveidīgas rasistu un antisemītu kustības kā Ungārijā un Rumānijā.

\strong{Ukrainas un Baltkrievijas rietumu daļai}, vai pēc poļu terminoloģijas, Polijas ``austrumu kresiem'' (tagadējā Rietumu Ukraina un Rietumu Baltkrievija, kuras kādreiz ietilpa Polijas sastāvā) nācās izciest atpalikušas nomales, agrāra izejvielu rajona lomu. Strādnieki šeit saņēma 1,5~līdz 2~reizes zemāku algu kā Polijas centrālajos rajonos. Apmaksātu atvaļinājumu, sociālās nodrošināšanas darba cilvēkiem faktiski nebija. Faktiski valdība īstenoja pret mazākumtautībām vērstu politiku.

1924.~gada 1.~jūlijā valdība iesniedza Seimam likumprojektus, kuri paredzēja noteiktos apstākļos ukraiņu un baltkrievu valodu lietošanu administratīvajās iestādēs un tiesās, kā arī bilingvālas pasniegšanas sistēmas ieviešanu austrumu vojevodistu skolās. Kā to saprata nacionālie mazākumi, zem piekāpšanās valodu jautājumā izkārtnes faktiski skolu likumprojektam bija jākļūst par asimilatoriskās polonizācijas politikas ieroci, tā slēptais mērķis bija pakāpeniska visu vēl saglabājušos ukraiņu un baltkrievu skolu pārveidošana par poļu skolām, bet nekādā ziņā ukraiņu un baltkrievu valodu ieviešana poļu skolās. 9.~jūlijā, uzstājoties likumprojektu apspriešanas gaitā, viens no ukraiņu deputātiem uzsāka runu ukraiņu valodā. Sēdes priekšsēdētājs PPS pārstāvis E.~Moračevskis viņam atņēma vārdu. Kāda deputāta~--- baltkrieva paziņojums, ka Rietumukraina un Rietumbaltkrievija bija pievienotas Polijai ar ieročiem, izsauca labējo deputātu fizisku uzbrukumu nacionālo mazākumu pārstāvjiem. Protestējot pret vardarbību, daudzi mazākumu deputāti, kā arī daži partijas ``\pltxti{Wyzwolenie}'' pārstāvji atstāja zāli. Toties vairākums atbalstīja valdības likumprojektus.

Saglabājās faktiski savstarpēji nesaistītas divu tipu skolas. Turīgo iedzīvotāju bērni varēja apmeklēt ģimnāzijas, kas deva iespējas iestāties augstskolās. Toties ap 20\% visu bērnu beidza tikai viengadīgu skolu. Mūsdienu poļu vēsturnieks J.~Tomaševskis par starpkaru periodā realizēto poļu nacionālo politiku raksta: ``Poļu vēsturnieks, kurš pazīst cīņu par dzimto valodu, kura tika vesta visos [agrāk] sadalītajos [Polijas] apgabalos, var tikai ar kauna sajūtu studēt dokumentus par ukraiņu skolu ierīcību Polijas Republikā, kurai, pēc daudzu tās dibinātāju sapņiem bija jākļūst par taisnīguma valsti visiem pilsoņiem.''

Pret mazākumtautībām vērstas politikas tiešie realizētāji kļuva \strong{poļu kolonisti} (vietējais nosaukums ``osadņiki'' (poļu~--- \pltxti{osadnicy}, vienskaitlī \pltxti{osadnik}~---  ieceļotājs, kolonists)~--- agrākie poļu armijas karavīri, virsnieki, viņu ģimeņu locekļi, kā arī civilie pārceļotāji, kuri pēc poļu-padomju kara saņēma zemes gabalus Rietumu Ukrainā un Rietumu Baltkrievijā. Pēc A.~Gurjanova datiem 85\% šo kolonistu bija poļi, tomēr to vidū bija arī ukraiņi un baltkrievi. Ievērojot ierēdņu, dzelzceļa darbinieku, citu kalpotāju pieplūdumu Volīnijā 18 poļu valdīšanas gados poļu tautības iedzīvotāju skaits pieauga no 7,5 uz 15\%. Līdzīga situācija bija arī citās austrumu vojevodistēs. 1919.--1929.~gadā 77~tūkstoši poļu kolonistu saņēma ap 600~000~ha vietējiem iedzīvotājiem~--- ukraiņiem un baltkrieviem atņemtas labākās zemes. Kolonistiem piešķirtie zemes gabali (15~līdz 40~ha) pēc lieluma vairākkārt pārsniedza ne tikai nabadzīgāko, bet arī vidējo vietējo zemnieku zemes īpašumus. Pavisam kolonistu skaits sasniedza vairākus simtus tūkstošu. Pēc baltkrievu zinātnieku datiem 1921--1939.~gg. no etniskajām poļu zemēm uz Rietumu Baltkrieviju vien bija pārcelti ap 300~000 kolonistu. Austrumgalīcijā un Volīnijā viņu bija ap 200~000. Nereti viņi pildīja arī policijas funkcijas attiecībā pret vietējiem iedzīvotājiem. Tas viss kopā radīja plašu pretpolisku un pretkolonistu noskaņojumu vietējos ukraiņu un baltkrievu zemniekos. Poļu kolonistus vietējie iedzīvotāji ienīda kā ``kulakus'' (kr. \rutxti{кулак}~--- turīgs zemnieks, kurš pie bagātības ticis, ekspluatējot algotu darba spēku, nodarbojoties ar starpniecību un augļošanu), jo tie bija saņēmuši vietējo zemnieku apstrādātās zemes, kā karavīrus~--- iekarotājus, kuru rezultātā baltkrievi un ukraiņi nokļuva poļu okupācijā, kā polonizatorus, vietējās kultūras vajātājus. Poļu kolonizatoru darbība radīja ukraiņos un baltkrievos tādu naida uzplūdumu, ka viņi sāka masveidā dedzināt poļu muižnieku saimniecību ēkas.

Kad 1924.~gadā Polijas austrumos bija neraža un ne tikai zemnieki, bet arī vietējie poļu magnāti Radzivilli, Potocki u.c. nebija samaksājuši nodokļus, poļu varas iestādes tos mēģināja piedzīt no vietējiem zemniekiem. To ievākšanā piedalījās policija un pat karaspēks. Kā atbilde uz to iedegās partizānu cīņa. Rietumbaltkrievijā un Viļņas apgabalā vien to rindās cīnījās ap 10~000~cilvēku. (Pēc Polijas armijas Ģenerālštāba Izlūkošanas nodaļas analītiķu datiem Viļņas apgabalā, Poļesjē, Naļibokas, Belovežas un Grodņas gāršās darbojās 5-6 tūkstoši partizānu.) Kaut tikai par dalību partizānu akcijās draudēja 8~gadus ilgs cietumsods, pēc poļu vēsturnieka-emigranta V.~Lapiņska datiem partizānu karš turpinājās līdz 1926.~gadam. 1924.~gada septembrī izraisījās sadursmes Koveļas (poļu \pltxti{Kowel}, ukraiņu \uktxti{Ковель}), Luckas (poļu \pltxti{Łućk}, ukraiņu \uktxti{Луцьк}), Ternopoles (poļu \pltxti{Tarnopol}, ukraiņu \uktxti{Тернопіль},) Viļņas apkārtnē. Oktobrī zemnieku sacelšanās izplatījās uz Svencjanu (lietuviešu \lttxti{Švenčionys}, poļu \pltxti{Święciany}) un Ļvovas rajonos. 9.septembra naktī partizāni organizēja divu pret viņiem nosūtītu bruņuvilcienu sadursmi.

Runājot par pamatu baltkrievu un ukraiņu zemnieku neapmierinātībai un viņu radītai partizānu kustībai, jāmin arī Padomju Krievijas/PSRS iestāžu inspirēta darbība pret Polijas valsti. Jau pēc Rīgas miera noslēgšanas Padomju Krievija gatavoja un iesūtīja Rietumukrainā kaujiniekus, kas organizēja pretošanos Polijas varas iestādēm. Iespējams, ka tas notika sakarā ar iepriekšējā nodaļā pieminēto V.~Ļeņina vēlmi atriebties par Polijas atbalstu S.~Bulak-Balahoviča iebrukumiem padomju teritorijā. Arī Polijas armijas Ģenerālštāba Otrās (izlūkošanas) pārvalde veidoja un iesūtīja padomju teritorijā bruņotas pretpadomju vienības.

Abi izlūkošanas dienesti: Polijas un padomju veda slepenu karu pēc principa ``aci pret aci un zobu pret zobu''.

Sarkanās armijas Izlūkošanas pārvalde jau pēc poļu-padomju kara beigām uzsāka kaujas vienību radīšanu un iesūtīšanu Polijas austrumu vojevodistēs ar uzdevumu izveidot tur masveida pretošanās kustību poļu varas iestādēm. Bija cerēts, ka profesionālie kaujinieki iejūsminās apspiestās masas un to cīņas rezultātā boļševikiem izdosies atgriezt Rietumbaltkrieviju un Rietumukrainu padomju republiku sastāvā. Šo vienību darbība ieguva kopēju nosaukumu ``aktīvā izlūkošana''. Uz vēl vāji apsargātās robežas tika veidoti speciāli ``logi'' šo ``aktīvo izlūku'' pārejai uz kaimiņvalsts teritoriju. No 1921.~gada aprīļa līdz 1924.~gada aprīlim uz poļu-padomju robežas tika fiksēti 259 Sarkanās armijas Izlūkošanas pārvaldes kaujinieku robežas pāriešanas gadījumi. Par veikto operāciju slepenības pakāpi liecina tas, ka par to nezināja pat padomju drošības iestādes: Valsts Politiskā pārvalde un tās mantiniece Valsts Apvienotā Politiskā pārvalde. Kad 1923.~gada oktobrī pārvaldes priekšnieka vietnieks J.~Unšlihts pieprasīja no pārvaldes Kontrizlūkošanas nodaļas priekšnieka A.~Artuzovam informāciju par partizānu kustību Austrumpolijā, tā darbinieki tā arī nespēja sagatavot skaidru atskaiti, tajā kopā bija ``sadiegti'' fakti gan par padomju ``izlūku'', gan stihisko nemieru dalībnieku, gan pretpadomiski noskaņoto ukraiņu nacionālistu, gan vienkārši kriminālnoziedznieku darbību.

Poļesjes vojevodistes austrumu un Novogrudokas un Viļņas vojevodistu dienvidu daļās darbojās padomju kaujinieku vienības K.~Orlovska un S.~Vaupšasa vadībā. Jau 1922.~gada vasarā par viņu darbību runāja ``visa Polija''. Tika iznīcināti policijas iecirkņi, dedzinātas muižas, uzspridzināts dzelzceļa tilts, iznīcinātas lokomotīves. Visur, kur tas bija iespējams, tika sasaukti mītiņi un sapulces, kur oratori aicināja zemniekus atbalstīt sacēlušos, veidot pašaizsardzības vienības. Padomju kaujinieku atbalsts sekmēja arī vietējo nemiernieku darbību.

1923.~gadā šī darbība vēl pastiprinājās, bet tās kulminācija bija 1924.gadā. Šī gada augustā S.~Vaupšasa vadītā vienība uzbruka Stovbcu pilsētiņai (baltkrievu \betxti{Стоўбцы}), sagrāva tur vietējo garnizonu, povjata policijas pārvaldi, pilsētas policijas iecirkni, ieņēma cietumu un atbrīvoja tajā ieslodzītos ievērojamus poļu un baltkrievu komunistu vadītājus S.~Mertensu un P.~Korčiku. Kaujinieku vienība K.~Orlovska vadībā sagrāba vilcienu, ar kuru brauca Poļesjes vojevoda S.~Dovnarovičs un atbruņoja apsardzi. Kaut vilcienā atradās vairākas augstu stāvošas personas, nevienu no tām nenogalināja. Taču, kā ziņoja padomju avīze ``\rutxti{Ленинградская правда}'' (``Ļeņingradas Taisnība''), vojevoda saņēma pērienu~--- 40 cirtienu ar pletni. Pēc tam Varšavas vara viņu atlaida no amata par gļēvulību. Atņēmuši pasažieriem naudu, dārglietas un ieročus, kaujinieki nozuda. Pavisam no 1924.~gada aprīļa līdz novembrim ``aktīvās izlūkošanas'' dalībnieki veica vairāk nekā 80~liela mēroga kaujas operāciju, bet visa gada laikā ap 200~uzbrukumu un diversiju. Polijas varas iestādes labi apzinājās šādas ``aktīvās izlūkošanas'' bīstamību. 1922.~gadā ``kresos'' tika ieviesti nāves sodi. Par ierēdņu nogalināšanu, uzbrukumiem valsts iestādēm, diversijām uz dzelzceļa to izpildīja 24~stundu laikā. Kad 1924.~gada septembrī grupa ukraiņu nacionālistu Ļvovā sarīkoja neveiksmīgu uzbrukumu pret Polijas prezidentu S.~Vojcehovski, tas kalpoja kā iegansts represiju izvēršanai gan pret ukraiņu un baltkrievu nacionālistiem, gan padomju kaujiniekiem, gan pret nacionālajiem mazākumiem kopumā. Viļņas kara apgabala pavēlnieks ģenerālis E.~Ridzs-Smiglijs saņēma plašas pilnvaras. Tā kā Polijas pierobežas policija nespēja tikt galā ar ``aktīvās izlūkošanas'' grupu un vietējo partizānu darbību, 1924.~gada septembrī tika nodibināts Pierobežas aizsardzības korpuss (\pltxti{Korpus Ochrony Pogranicza}), kurš jau tai pašā gadā piedalījās 89~nopietnās sadursmēs. Korpusa spēki sastādīja 24~kājnieku bataljonus un 20~kavalērijas eskadronus, kopā ap 25~tūkstošu karavīru un virsnieku. Rietumbaltkrievijā un Rietumukrainā tika nostiprināti arī Polijas armijas garnizoni. Protams, notika padomju kaujinieku aresti, daži tika sodīti ar nāvi, citi saņēma pamatīgus cietumsodus. Faktiski austrumu vojevodistēs pastāvēja, kaut formāli nepasludināts, kara stāvoklis. Kaut sakarā ar ``daudz-maz normālu'' attiecību ar kaimiņvalstīm, arī Poliju, nodibināšanos Sarkanās armijas Izlūkošanas pārvalde vairākkārt deva instrukcijas pārtraukt ``aktīvās izlūkošanas'' veikšanu tajās, daudzi kaujinieki tās neizpildīja. Taču drīz kāds notikums padomju diversantu darbībai Polijā pārvilka svītru. 1925.~gada janvāra sākumā padomju kaujinieku vienību, tērpta Polijas armijas uniformās, tikta piespiesta pie padomju robežas un, lai izbēgtu iznīcināšanu, devās tai pāri un sagrāva padomju robežsargu posteni, kurš padomju kaujiniekus noturēja par poļu iebrucējiem un pretojās tiem cik spēja. Polijas un arī Rietumvalstu prese. Jebkura nākamā kādas padomju kaujinieku grupas kļūda varēja novest pie plaša starptautiska skandāla ar neprognozējamām sekām. Šai sakarā notika KK(B)P CK Politbiroja ārkārtas sēde. Tika radīta komisija incidenta apstākļu noskaidrošanai. 1925.~gada 25.~februārī KK(b)P CK Politbirojs apstiprināja lēmumu par ``aktīvās izlūkošanas'' grupu likvidēšanu Polijas teritorijā. Kaut iepriekšējo instrukciju nepildīšanu Politbirojs skaidroja gan ar iepriekšējos gados iedibinātājām tradīcijām, gan arī ar vietējo komunistisko partiju vājo vadību pār stihisko zemnieku kustību, tagad visas vietējo kaujinieku vienības bija jānodod tieši šo vietējo komunistisko partiju pārziņā, lai tās vadītos vienīgi no ``dotās valsts revolucionārajām interesēm'', atsakoties no to izmantošanas PSRS militāro iestāžu labā. PSRS pierobežas joslu bija ``jāattīra'' no kaujiniekiem, kuri nekontrolēti šķērsoja robežu. PSRS militāro mērķu sasniegšanai kaimiņvalstīs ``aktīvās izlūkošanas'' vienību vietā bija jārada īpaši konspiratīvi atbalsta punkti, kas gatavotos ``destruktīvai darbībai'' ienaidnieka aizmugurē iespējamā kara laikā. Tie savukārt nedrīkstēja būt saistīti ar komunistiskajām partijām. Acīmredzot šī lēmuma pieņemšanu lielā mērā izsauca Polijas protesti, jo Politbiroja lēmumā bija uzsvērts, ka Polijas valdībai nav nekādu pierādījumu par ``aktīvās izlūkošanas'' grupu saistību ar PSRS varas iestādēm, bet gan tikai minējumi. Tāpēc tika izteikts norādījums visiem pārmetumiem dot ``stingru pretsparu''. Tikai mūsdienās pētnieki varēja iegūt un publicēt materiālus par padomju ``aktīvās izlūkošanas'' darbību.

1925.~gada jūnijā tika izformēta S.~Vaupšasa vadītā vienība. Nedaudz vēlāk PSRS atgriezās arī K.~Orlovskis. (Starp citu, gan abi nosauktie, gan K.~Orlovska vienības kaujinieks V.~Koržs Otrā pasaules kara gados kļuva par Padomju Savienības varoņiem. Tā, V.~Korža vadībā 1941.~gada 28.~jūnijā notika pirmā padomju partizānu vienības kauja ar vācu iebrucējiem, kurā bez zaudējumiem no partizānu puses tika iznīcināts viens vācu tanks, saņemti divi gūstekņi. Jau pēckara gados K.~Orlovskis kļuva arī par Sociālistiskā darba varoni, bet viņa vadītais kolhozs par pirmo PSRS kolhozu~--- miljonāru. Viņa vārdā Baltkrievijā tika nosauktas ielas un iestādes.)

Pēc padomju specdienestu militārās operāciju pārtraukšanas visa pret Poliju vērstās, arī faktiski militārās, darbības vadība tika nodota Rietumukrainas un Rietumbaltkrievijas komunistiskajām partijām. Partizānu organizācijas, palikušas bez vadības un finanšu līdzekļiem, degradējās. To biedri pauda neapmierinātību, uzsāka darbību uz savu roku, kas bieži pārvērtās par vienkāršu laupīšanu. Izmantojot apstākļus, kad profesionāli sagatavotu kaujinieku uzbrukumi vairs nedraudēja, Polijas iestādes izvērsa cīņu pret vietējiem partizāniem. Tikai Novogrudokas vojevodistē vien tika arestēti ap 1~400 cilvēku.

\strong{Rietumu Baltkrievijai} (baltkrievu \betxti{Заходняя Беларусь}) raksturīga pastāvīga parādība bija bezdarbs. Bezdarbnieku pabalsti tika maksāti tikai daļai no tiem. Tā 1927.~gadā to saņēma tikai 13,6\% reģistrēto bezdarbnieku. Grūts bija arī zemnieku liktenis. Zemes īpašums bija sadalīts ļoti nevienmērīgi. Mazāk kā 1\% tās īpašnieku rokās atradās lielākā zemes daļa. Daudz bija arī bezzemnieku. Valdība centās šeit nostiprināt poļu tautības iedzīvotāju ekonomisko stāvokli, viņiem pārdeva valsts īpašumā esošo zemi. Bijušajiem poļu-padomju kara dalībniekiem līdz 45~ha zemes piešķīra vai nu bez maksas, vai ar atvieglotiem noteikumiem. Pavisam uz Rietumu Baltkrieviju no etniskās Polijas pārvietoja 300~000 poļu. 30.~gados laukos pastiprinājās sociālā diferenciācija. Baltkrievi uzmanīgi sekoja tam, kā blakus Padomju Baltkrievijā tika likvidēts muižnieku īpašums, kā tur dzīvo zemnieki. Nabadzība veicināja emigrācijas pieaugumu.

1931.~gadā pilsētās dzīvoja tikai ap 15\% novada iedzīvotāju. Ja Rietumbaltkrievijā dzīvoja ap 13\% Polijas iedzīvotāju, tad strādnieku šeit bija tikai 1,8\% no visiem Polijā strādājošajiem, saražotā produkcija sastādīja 3\% no valstī saražotās. Vairumā uzņēmumu strādāja 5~līdz 20~strādnieki. Lielākie uzņēmumi bija Pinskas sērkociņu fabrika, Grodņas tabakas fabrika, Lidas gumijas izstrādājumu fabrika. 1936.~gadā novadā bija 25 tūkstoši bezdarbnieku.

Smago sociāli ekonomisko stāvokli papildināja arī vietējo iedzīvotāju diskriminācija pēc nacionālās pazīmes. Kā aprēķinājuši baltkrievu zinātnieki, pēc Rīgas miera Polija anektēja 112~km$^{2}$ Baltkrievijas teritorijas ar 4~miljoniem iedzīvotāju, kuru vidū bija 67\% baltkrievu, 12--13\% poļu, 5--6\% ukraiņu, ap 9\% ebreju, 2\% krievu un 3\% lietuviešu, taču valsts iestādēs baltkrievu valodu nedrīkstēja lietot, valsts dienestā baltkrievus nepieņēma. Poļi veidoja 75\% no novadā dzīvojošās inteliģences, baltkrievi~--- 5\%. Baltkrievu nacionālā pašapziņa palika vēl nenoformēta. Tūlīt pēc Rīgas miera līguma noslēgšanas gan darbojās ap 400~baltkrievu sākumskolas, 7~pilnas skolas un 3~skolotāju semināri. Valdība pat daļēji finansēja baltkrievu kultūras darbību, subsidēja periodisko presi. Taču jau pēc pāris gadiem tā sāka mainīt politiku pret ``austrumu kresiem''. Vairs neļāva dibināt jaunas baltkrievu skolas. Sāka slēgt nepoļu skolas, nepoļu preses izdevumus.

Pēc vācu publicistu T.~Cilha un J.~Volmera datiem jau līdz 1923.~gadam vairāk nekā 350~baltkrievu skolas tika aizstātas ar poļu skolām. Rietumbaltkrievijā uz 2,5~miljoniem baltkrievu palika tikai 37~baltkrievu sākumskolas. (Līdzīgus datus dod arī mūsdienu baltkrievu vēsturnieki.) Jau 1925.--1926.~mācību gadā to skaits bija samazinājies līdz 25. Baltkrievu skolu vietā šeit tika atklātas 3~380 poļu skolas. Cīņas rezultātā 1927--1928.~gadā valdība tomēr bija spiesta atvērt 29~baltkrievu un 49~jauktas poļu--baltkrievu skolas. Eksistēja arī 4~baltkrievu ģimnāzijas. Pēc krievu vēsturnieces J.~Jakovļevas datiem 1938.--1939.~mācību gadā bija palikušas 5~kopīgas poļu-baltkrievu vispārizglītojošās skolas, 44~skolas un 1~ģimnāzija, kur kā vienu no disciplīnām mācīja arī baltkrievu valodu. (Jāsaka, ka dažādu autoru sniegtie dati par skolu skaitu atšķiras, taču visi tie demonstrē baltkrievu skolu izskaušanu, poļu izglītības iestāžu uzspiešanu.)

Lūk, citāts no Belostokas vojevodas H.~Ostaševska ziņojuma Polijas Iekšlietu ministrijai ``Poļu valdošā stāvokļa nostiprināšanas problēmas Belostokas vojevodistē'': ``Agri vai vēlu baltkrievu iedzīvotāji ir jāpolonizē. Viņi veido pasīvu masu, bez plašas tautas pašapziņas, bez savām valstiskuma tradīcijām. Vēloties šo procesu paātrināt, mums ir jāpārvar senā baltkrievu kultūra.''

Baltkrievu pārstāvju uzstāšanās Seimā saturēja sūdzības par apvainojumiem, spaidiem no dažāda līmeņa varas iestāžu puses. Cīņa par bērnu apmācību baltkrievu valodā faktiski pārvērtās cīņā par etnisko izdzīvošanu, jo bērnus spieda ne tikai mācīties, bet arī ikdienā runāt poļu valodā. Par Rietumu Baltkrievijas politisku un kultūras centru kļuva Viļņa. Tomēr izglītības baltkrievu valodā, tāpat kā izglītības vispār, problēma bija ļoti asa. 1931.~gadā baltkrievu apdzīvotajos novados analfabētu vecumā virs 10~gadiem bija 43\% visu iedzīvotāju, apmēram 13\% skolas vecuma bērnu to neapmeklēja. Arī 1939.~gadā 35\% visu iedzīvotāji neprata lasīt un rakstīt. Uz 1939.~gada 1.~septembri novadā 129,8~tūkstoši bērnu, galvenokārt no baltkrievu ģimenēm, palika ārpus izglītības sistēmas. 1927.~gadā tika izdoti 23~baltkrievu avīzes un žurnāli, 1930.~--- 12, 1932.~--- 8, bet 1939.~gadā palika tikai propoliski un klerikāli izdevumi. Nepastāvēja neviena baltkrievu teātra vai muzikāla iestādījuma.

Poļu varas iestādēm nepatika baltkrievu katoļu garīdznieku darbība baltkrievu nacionālās pašapziņas stiprināšanā. Tika sperti soļi viņu darbības ierobežošanai. Piemēram, ksendzu V.~Godļevski arestēja un viņam piesprieda divus gadus cietumsoda par pretvalstisku darbību. Pēc tam Polijas katoļu baznīcas vadība pakāpeniski baltkrievu garīdzniekus pārcēla darbā Polijas vidienē, bet to vietā nosūtīja etniskos poļus. Līdz 1936.~gada jūnijam vairāk nekā 1~300~pareizticīgo dievnamu tika pārvērsti par katoliskiem. Katoļu baznīca izplatīja uzskatu, ka būt katolim nozīmē būt polim. Siltas jūtas pret poļiem baltkrievos tas neradīja. Redzams baltkrievu nacionālās kustības darbinieks J.~Maļeckis rakstīja: ``Kā gan varēja baltkrievu vidū attīstīties polonofīlija, ja poļu valdība ciniski pareģoja, ka pēc 50~gadiem Žečpospoļitas teritorijā nebūs nekā baltkrieviska, un tāpēc Viļņas, Novogrudskas un Belostokas vojevodas slāpēja tautā pašas minimālākās baltkrievu kultūr-nacionālās darbības izpausmes''.

Savu stāvokli Polijas baltkrievi varēja salīdzināt ar poļu stāvokli Padomju Baltkrievijā, kur tie tāpat kā baltkrievi Polijā, sastādīja vienu no etniskajām pamatgrupām. T.s. ``iesakņošanas'' (krievu~--- \rutxti{корениза́ция }~--- padomju iestāžu politika nacionālajā jautājumā 20.gados un 30.gadu sākumā, vērsta uz padomju varas nostiprināšanu un nekrievu tautu nacionālās pašapziņas stiprināšanu, nacionālo kadru sagatavošanu) realizācijas laikā, Baltkrievijas Sociālistiskajā Padomju republikā blakus ``baltkrievizācijai'' tika ievērotas arī poļu etniskās vajadzības. 1921.~gadā BSPR Centrālā izpildu komiteja ar speciālu lēmumu atzina poļu valodas, tāpat kā baltkrievu, krievu un jidiša valodu, valsts valodas statusu. Arī Baltkrievijas Konstitūcijā poļu valoda tika fiksēta kā viena no valsts valodām. Eksistēja atsevišķas poļu un jauktas poļu-baltkrievu skolas. 1924.~gadā Baltkrievijā bija 136~poļu skolas. No 1922.~gada Minskā darbojās poļu pedagoģiskais tehnikums, Baltkrievijas Valsts pedagoģiskā institūta poļu nodaļa, 1920.--1930.~gadā tika izdotas 6~poļu avīzes. Kompaktas poļu apmešanās rajonos tika organizētas poļu ciemu padomes. No 1932.~līdz 1936.~gadam uz Koidanovas (\rutxti{Койдановский}) rajona pamata pastāvēja Poļu nacionālais rajons kā autonoma administratīvi-teritoriāla vienība. Nacionālajā rajonā visos administratīvajos līmeņos iestādēs tika ieviesta poļu valoda. Rajons gan tika izveidots vietā, kur poļu īpatsvars nebūt nebija vislielākais Baltkrievijā, toties atradās tas tuvu Polijas robežai, varēja kalpot par padomju realizētās nacionālās politikas skatlogu. Tā kā netālu bija dzimis jau mūžībā aizgājušais F.~Dzeržinskis, pilsēta ieguva un vēl tagad nes viņa vārdu. Tika gatavota arī otra poļu rajona dibināšana. 30.~gadu beigās staļiniskais režīms visu šo ``iesakņošanas'' politiku varmācīgi aprāva. Taču 20.~gados to nedz vienā, nedz otrā robežas pusē nevarēja paredzēt.

Saprotams, ka smagais sociāli ekonomiskais un nacionāli politiskais stāvoklis, valsts varas realizētais \strong{kurss uz polonizāciju} Rietumu Baltkrievijā \strong{izsauca pretošanos}. Tās gaitā dažādi politiski spēki stādīja dažādus mērķus. Ja nacionālistu kustība meklēja atbalstu Rietumos, tad daļa sociāli orientēto spēku skatījās austrumu virzienā.

Kā konstatējusi poļu vēsturniece D.~Mihaļuka, ``ierobežojumi kultūras un izglītības darbībā, katoļu iecelšanas atbalstīšana valsts un vietējās pārvaldes iestādēs, pareizticīgo iedzīvotāju tiesību ierobežošana vietējās varas struktūrās, politiskās vajāšanas, kā arī Maskavas iejaukšanās Polijas starpnacionālajās attiecībās radīja apstākļus baltkrievu komunistiskās kustības attīstībai Polijā.'' Rietumu Baltkrievijā komunisti uzstājās par visas baltkrievu tautas apvienošanu Baltkrievijas padomju sociālistiskās republikas (BPSR) (baltkrievu \betxti{Беларуская Савецкая Сацыялicтычная Рэспублiкa)} sastāvā. 1921.--1922.~gadā Brestā, Grodņā, Pinskā u.c. bija radītas komunistiskās partijas organizācijas, tika atjaunotas vai radītas arodbiedrības. Arī Polijas Komunistiskās partijas II kongress 1923.~gadā uzsvēra baltkrievu strādnieku un zemnieku tiesības atbrīvoties no ``muižnieciski-kapitālistiskās Polijas'' un pievienot savas zemes Padomju Baltkrievijai. 1923.~gada oktobra beigās Viļņā notika Rietumu Baltkrievijas komunistisko organizāciju konference, kurā tika pasludināta Rietumu Baltkrievijas Komunistiskās partijas (RBKP) kā Polijas komunistiskās partijas sastāvdaļas nodibināšana. No 1924.~gada RBKP uzņēma kursu uz bruņotas sacelšanās gatavošanu. Taču Polijas valdība rīkojās enerģiski, partizānu, bet pēc būtības~--- teroristu, cīņai trūka pēc miera un stabilitātes noilgojušās tautas atbalsta un RBKP drīz bija spiesta aicināt savus biedrus pārtraukt bruņotu cīņu.

20.~gadu sākumā baltkrievu vidū diezgan ievērojama ietekme bija Baltkrievu sociālistu-revolucionāru (eseru) partijai. Tajā bija vairāki tūkstoši biedru un tā balstījās galvenokārt uz zemniekiem. Partija prasīja veikt demokrātiskas reformas, muižnieku zemju sadali, nodrošināt nacionālu līdztiesību, neatkarīgas Baltkrievijas Republikas izveidi, kura būtu apvienota ar Lietuvu vienā valstī.

Visas baltkrievu neapmierinātības izpausmes Polijas varas iestādes apspieda ar dzelžainu roku. Tā, 1923.~gadā Polijā tika piespriesti 109~nāves sodi. Vairākums no notiesātajiem bija baltkrievu patrioti.

Parlamentāru opozīciju Polijas valdībai veidoja Baltkrievijas sociāldemokrātiskā partija, kura balstījās galvenokārt uz nacionālo baltkrievu inteliģenci. Partija uzstājās ar sīki zemnieku saimniecību apvienošanu kooperatīvos, lielo muižnieku zemes īpašumu atsavināšanu parlamentārā ceļā un sadali mazzemniekiem. Partijas programma saturēja arī demokrātisko brīvību prasības, tās mērķis bija demokrātiskas baltkrievu republikas izveide, kas apvienotu zemes, kurās vairākumā dzīvoja baltkrievi.

Ierobežotākas sociālās un politiskās prasības izvirzīja Baltkrievijas kristīgā demokrātija. Tai bija ietekme katoļticīgo zemnieku vidū. Partija noraidīja baltkrievu tautas dalījumu šķirās, uzstājās par evolucionāru attīstības ceļu.

Der atzīmēt, ka 1922.~gadā Polijas Seimā tika ievēlēti 11, bet Senātā~--- 3~baltkrievi no kopējā 444 deputātu un 111~senatoru skaita. 1928.~gadā Seimā ievēlēja 10, bet Senātā~--- 2~baltkrievus. 1935.~un 1938.~gadā augstākajās valsts pārstāvniecības iestādēs netika ievēlēts neviens baltkrievs. Līdz ar to parlamentāro cīņas ceļu par savām tiesībām daudzi baltkrievi pamatoti uzskatīja par mazperspektīvu.

Kopīga visām baltkrievu sīkpilsoniskajām partijām bija nacionālo tiesību, neatkarīgas demokrātiskas Baltkrievijas Republikas izveides prasība, uzstāšanās pret apvienošanos ar BPSR. Taču 20.~gados iezīmējušies Baltkrievijas padomju sociālistiskās republikas kā nacionālas valsts panākumi, fakts, ka 1924.~un 1926.~gadā Krievija tai nodeva no sava sastāva 17~apriņķus, kuros dzīvoja galvenokārt baltkrievi, kā arī atbrīvošanās kustības radikalizēšanās Rietumu Baltkrievijā, noveda pie domstarpību pieauguma šo partiju biedru vidū. Sakarā ar to 1924.~gadā pārtrauca savu darbību Baltkrievijas sociāldemokrātiskā partija, Baltkrievu sociālistu-revolucionāru partija pašlikvidējās un atzina BPSR. Daļa tās biedru iestājās RBKP.

1925.~gadā Rietumu Baltkrievijā tika radīta legāla demokrātiska masu organizācija~--- Baltkrievu zemnieku-strādnieku kopiena (baltkrievu \betxti{Беларуская Сялянска-Работніцкая Грамада}, poļu \pltxti{Białoruska Włościańsko-Robotnicza Hramada}), kura drīz vairāk nekā ar diviem tūkstošiem organizāciju aptvēra no 100~līdz 150~tūkstošu biedru (pati masveidīgākā politiskā partija Baltkrievijas vēsturē, protams, izņemot padomju laiku Baltkrievijas Komunistisko partiju) un 1925.--1927.~gg. skaitījās lielākā revolucionārā organizācija Eiropā. Tā cieši sadarbojās ar RBKP.

Baltkrievu zemnieku-strādnieku kopiena prasīja Rietumu Baltkrievijas pašnoteikšanos un visu baltkrievu zemju apvienošanu, demokrātiskās brīvības un muižnieku zemju konfiskāciju un sadalīšanu, kooperācijas veicināšanu, strādnieku kontroles ieviešanu pār ražošanu utt. Kaut organizācija par saviem mērķiem cīnījās tikai ar legāliem līdzekļiem, 1927.~gada martā J.~Pilsudska vadītā valdība organizāciju ne tikai aizliedza, bet brutāli sagrāva, arestējot simtus aktīvistu. Neievērojot parlamentāro imunitāti, apcietināja arī piecus Seima deputātus, kuru represēšana bija rupjš darbojošās Konstitūcijas pārkāpums. Par ``attaisnojumu'' ``sanācijas'' režīmam var pateikt, ka tāpat tas izrīkojās arī ar Neatkarīgās zemnieku partijas (\pltxti{Niezależna Partia Chłopska}) pārstāvjiem, kuri aizstāvēja poļu nabadzīgo zemnieku intereses.

30.~gados Polijas politika pret ``austrumu kresiem'' kļuva vēl autoritārāka un agresīvāka. Pret nacionālajiem mazākumiem tika īstenota to kultūras apspiešanas politika. Notika baltkrievu kultūras iestāžu likvidācija. Tika slēgti 140~pareizticīgo dievnami. 100~km pierobežas zonā ar PSRS iedzīvotājus ar izsūtīšanas draudiem spieda pieņemt katoļu ticību. Taču krīzes gados notika strādnieku streiki, aktivizējās zemnieki, kuri aizstāvoties pret valdības sūtītajām soda ekspedīcijām atsevišķas vietās organizēja partizānu grupas. Valdība ieveda ārkārtējās tiesas. 1931.~gada Rietumu Baltkrievijā tika piespriesti 44, bet 1932.~gadā~--- 76 nāves sodi.

1935.gadā baltkrievi zaudēja pēdējo deputāta vietu parlamentā. Ar ieganstu, ka baltkrievu kultūras biedrības nodarbojas ar nelegālu komunistisko propagandu, arvien vairāk to tika slēgtas, bet 1938.gadā pareizticīgajiem Polijas iedzīvotājiem aizliedza iepirkt zemi. Polijas politiku baltkrievu apdzīvotajos novados raksturoja to polonizācija ar speciālo robežapsardzes vienību, (\pltxti{Korpus Ochrany Pogranicja}), kolonistu, katoļu baznīcas un poļu skolas palīdzību. Rietumbaltkrievu dzejnieks M.~Tanks 1938.~gada aprīlī par poļiem rakstīja: ``Panu Polijas pastāvēšanas gados izaugusi vesela paaudze, saindēta ar lielvalstisku šovinismu, katolicisku un nacionālistisku garu \dots{} un bez poļu oficiālās politikas tā nepazīst neko \citespace{} Daži no politiskajiem procesiem un skopajiem ziņojumiem par pacifikācijām [latīņu~--- \latxti{pacificatio}~--- nomierināšana, par tām sīkāk~--- turpmāk. V.Š.] pirmo reizi uzzināja, ka zem viena jumta ar tiem, tikai aiz režģiem slēgtiem logiem, dzīvo miljoni citu tautību cilvēku, kuriem atņemtas visas cilvēku tiesības\dots{}''

Arī \strong{Rietumu Ukrainā} (ukraiņu~--- \uktxti{Західна Україна}), kura aptvēra Austrumu Galīciju un Volīniju, tika īstenota \strong{polonizācijas} politika. 1922.~gada 26.~septembrī Polijas Seims atzina Austrumgalīcijas ukraiņiem autonomijas tiesības. Rietumvalstu sabiedroto Sūtņu padome Londonā 1923.~gadā, lemjot jautājumu par Austrumgalīcijas nodošanu Polijai, arī paredzēja, ka Polijā ukraiņiem būs plašas pašpārvaldes tiesības. Polijas valsts iestādes šos starptautiskos lēmumus, tāpat kā citus iepriekš minētos, pilnīgi ignorēja.

Jaundibinātajā Polijas valstī 20.~gadsimta divdesmitajos un trīsdesmitajos gados ukraiņi dzīvoja divās dažādās vēsturiski un politiski izveidojušās daļās. Lielākā ukraiņu daļa dzīvoja bijušajā Austroungārijas Galicijā~--- Ļvovas, Staņislavas un Tarnopolas vojevodistēs. Šie ukraiņi pēc savas reliģiskas pārliecības pārsvarā piederēja pie grieķu—romiešu Uniātu Baznīcas. Agrākās cariskās Krievijas Volīnijas un Holmas guberņās~--- Polijas Volīnijas un Ļubļinas vojevodistēs dzīvoja galvenokārt pareizticīgie ukraiņi. Taču 20.--30.~gados visā Rietumukrainā bija ļoti nemierīgi. Poļu nacionālisti ukraiņus citādi nesauca kā par rusīniem, bet Rietumukrainu~--- par Mazpoliju.

Ukraiņi bija lielākā mazākumtautība Polijā un pret viņiem arī tika vērsts vairākums diskriminējošo likumu. Daļai konfliktu saknes bija atrodamas vēl poļu~--- ukraiņu cīņās 1918.--1920.~gadā. Galīcijā pēc asajām poļu un ukraiņu sadursmēm saglabājās starpnacionāls sasprindzinājums. Galīcijas likteni, kā jau iepriekš minēts, izlēma sūtņu konference Parīzē 1923.~gada 14./15.~martā. Taču tā paredzēja piešķirt novadam autonomiju (to jau pirms tam bija solījis arī Polijas Seims). Polijas varas iestādes šo prasību tā arī neizpildīja. Tas, ka Polijas valdība neievēroja tās 1919.~un 1923.~gadā Rietumu sabiedrotajiem dotos solījumus piešķirt autonomiju Austrumgalīcijai, radīja īpašu sarūgtinājumu ukraiņos. Jau 1921.~gada 25.~septembrī Austrumgalīcijas apmeklējuma laikā, kad J.~Pilsudskis, kurš līdz dzīves beigām necieta, ka viņam blakus atrodas miesassargi, kopā ar Ļvovas vojevodu pēc pilsētas rātsnama apmeklējuma sēdās automobilī, no pūļa izlauzās jauns cilvēks un caur auto aizmugures stiklu sāka šaut uz valsts galvu. Pēc trešā šāviena revolveris atteicās darboties un terorists tika sagūstīts. Tas izrādījās ukraiņu nacionālists, kaujinieks S.~Fedaks. Viņam piesprieda sešus gadus cietumsoda. Tas neapturēja ukraiņu nacionālistu cīņu pret poļu varu. Kad 1921.~gadā Polijā sākās pirmā tautas skaitīšana, liela daļa Austrumgalīcijas un Volīnijas iedzīvotāju~--- ukraiņu to boikotēja, tāpat kā 1922.~gada vēlēšanas Seimā.

Polonizācijas politikas rezultātā tika radīti šķēršļi ukraiņu skolu darbībai. Polijas valsts pakāpeniski, bet noteikti \strong{mazināja ukraiņu iespējas iegūt izglītību} dzimtajā valodā. 1924.~gadā Rietumukrainā bija palikušas tikai 2~646~ukraiņu skolas un 7~ģimnāzijas, kad poļu mazākumam bija 7~211~skolas un 57~ģimnāzijas. Administratīvais spiediens, naudas trūkums un 1924.~gada Polijā pieņemtais skolu likums noveda pie tā, ka 1934./1935.gadā eksistēja vairs tikai 487~šādas skolas. Tātad to skaits bija samazinājies par apmēram 80\%. (Literatūrā atrodami arī citi dati. 1922.~gadā Austrumgalīcijā pastāvēja 2~247 poļu skolas, 2~426 ukraiņu, bet pēc 3 gadiem poļu skolu bija 2~325, ukraiņu 745, divvalodīgu~--- 745. Kā redzams, skaitļi, lai gan atšķirīgi, atklāj vienu un to pašu tendenci.) Tikai līdz 1930.~gadam pastāvēja trīs privātas ukraiņu ģimnāzijas.

Poļu vēsturnieks un publicists Z.~Zaluskis rakstīja: ``Volīnijā, laukos, kur 80\% sastādīja ukraiņu zemnieki, no 2~tūkstošiem vispārizglītojošo skolu tikai 8~bija ukraiņu, bet 500~skolās ukraiņu valoda tika izmantota kā otrā''.

Nacionālās diskriminācijas politika turpinājās arī augstākās izglītības jomā, aizliedzot mācību priekšmetu pasniegšanu ukraiņu valodā Ļvovas universitātē. Pirms Pirmā pasaules kara Ļvovas universitātē bija astoņi ukraiņu valodas profesori un docētāji Ukraiņu valodas katedrā. Tikai 1928./1929.~gadā tā tika atjaunota, bet nosauca to par ``Rutēņu filoloģijas katedru''. Līdz 1933.~gadam universitātē nedrīkstēja strādāt ukraiņu tautības docētāji. Faktiski Rietumukrainas kultūrpolitiskajā centrā ar daudzām zinātniskajām un kultūras iestādēm teju pilnīgā atšķirtībā no citām etniskajām grupām dzīvoja 160~000 ebreju, 120~000 poļu un 80~000 ukraiņu. Varšavā ukraiņiem bija vieglāk kontaktēties ar poļiem nekā Ļvovā.

Turpretī kaimiņvalstī~--- Padomju Savienībā, kur dzīvoja ap 1,2~miljoniem poļu (1926.~gada tautas skaitīšanā PSRS teritorijā bija reģistrēti 782~300 poļu izcelsmes pilsoņi. Tiesa, tikai 46\% no viņiem sarunājās poļu valodā. Polijas Ārlietu ministrijas analītiķi lēsa, ka 1930.~gadā PSRS dzīvoja 955~000 poļu), tai skaitā 650~tūkstošu Ukrainā un ap 300~tūkstošu Baltkrievijā), 20.gados un 30.gadu sākumā tika radīts nacionālo ciemu padomju tīkls. Tika organizēti arī poļu nacionālie rajoni. Lielā mērā to noteica padomju-poļu sāncensības loģika. Poļu komunisti~--- internacionālisti Padomju Savienībā, kuri aktīvi piedalījās šai procesā, stādīja uzdevumu neitralizēt ukraiņu un baltkrievu ``nacionālismu'', kā arī radīt pretsvaru vietējo poļu iedzīvotāju vēlmei doties uz savu nacionālo valsti~--- Poliju. Baltkrievijā un Ukrainā jau no 1921.~gada sākās poļu nacionālo Padomju organizēšana. 1923.~gadā Ukrainā jau darbojās 26~poļu ciemu padomes. 1924.~gadā Ukrainas PSR Tautas komisāru padome noteica iedzīvotāju skaita normas nacionālu padomju izveidei: 500~poļu~--- ciema padomei un 10~000~--- nacionālam rajonam, kur poļi būtu iedzīvotāju vairākums. 1927.~gadā Ukrainā jau pastāvēja 139~poļu ciema padomes (Baltkrievijā~--- 24, bet visā PSRS~--- 164). 1929.~gadā Ukrainā jau bija 153~poļu ciema padomes, kuras aptvēra apmēram 1/3 Ukrainā dzīvojošo poļu.

1925.--1935.~gadā Padomju Ukrainā pastāvēja speciāls poļu J.~Marhļevska vārdā nosauktais administratīvais rajons ar 40,5~tūkstošiem iedzīvotāju, 70\% no tiem bija poļi. 1932.~gadā Baltkrievijā tika organizēts F.~Dzeržinska vārdā nosaukts poļu autonomais rajons. Rajons pastāvēja līdz 1938.~gadam.

Poļi bija pārstāvēti PSRS Padomju kongresos. IV Vissavienības padomju kongresā (1927) bija ievēlēti 7, V kongresā (1929)~--- 19, VI (1931)~--- 13 poļi. Tika radīts tiesu aparāts poļu valodā.

Ukrainā 1927.~gadā pastāvēja 337~poļu skolas, 1930.~gadā~--- jau 381. Kopā PSRS 1927.~gadā pastāvēja 550~poļu skolu (Ir arī citi dati, ka pavisam PSRS darbojās vairāk nekā 670 poļu sākum- un vidusskolu, 7 speciālās skolas ar apmācību poļu valodā.) Nacionālās skolas aptvēra apmēram 1/3 poļu bērnu. Poļi dzimtajā valodā mācījās 2~augstskolās (Minskas Pedagoģiskajā institūtā un Kijevas Sociālās audzināšanas institūtā.) Darbojās trīs poļu profesionālie teātri, viena centrālā vissavienības, 16~rajonu un 6 republikāniskās avīzes. Ik gadu tika izdoti simti grāmatu poļu valodā. Republiku Zinātņu akadēmiju sastāvā funkcionēja poļu zinātniskie institūti. Tas ļāva izpausties poļu nacionālajai pašapziņai.

Protams, arī tāda PSRS raksturīga parādība kā kolhozi nepagāja secen padomju poļiem. Jau 1927.~gadā Ukrainā pastāvēja 13~poļu kolhozi, 1928.~gadā organizēja vēl 8. 1929.~gadā Ukrainā to bija jau 62, bet visā PSRS~--- vairāk nekā 100. Poļu kolhozos, tāpat kā citos, stājās pārsvarā nabadzīgie zemnieki. Pēc masveida kolektivizācijas sākšanās poļu zemnieku liktenis PSRS bija tāds pat, kā visu citu tautību zemniekiem. Interesanti, ka J.~Marhļevska vārdā nosauktajā administratīvajā rajonā kolektivizācijas rādītāji bija sliktāki nekā vidējie valstī. Tā, 1932. gadā sabiedriskoto zemnieku saimniecību skaits kopā valstī sastādīja 61,5\%, Ukrainā~--- 72\%, bet Marhļevska rajonā tikai 16,9\%.

Ar J.~Staļina diktatūras izveidošanos 30.~gadu sākumā etnisko mazākumu statuss PSRS būtiski mainījās. Ja 20.~gados J~ Staļins atbalstīja poļu autonomijas~--- J.~Marhļevska vārdā nosauktā nacionālā rajona izveidi Ukrainā, tad 30.~gados~--- jau poļu deportāciju no tā. Nacionālie rajoni un ciemu padomes tika ātri likvidēti. Nedz PSRS 1936.~gada Konstitūcijā, nedz pēc tās pieņemtajās savienoto un autonomo republiku konstitūcijās nebija vairs atrodami kādu nacionālo mazākumu struktūru pieminējumi. Tika slēgtas poļu skolas, klubi un bibliotēkas, poļu teātris Minskā.

Poļi Padomju Savienībā deportācijai tika pakļauti jau no 1933.~gada sākuma, bet masu kārtībā~--- veseliem ciemiem no 1935.~gada, kad no toreiz pierobežā esošajiem Kijevas un Viņņicas apgabaliem uz republikas austrumu rajoniem pārvietoja 41~650 cilvēkus, no kuriem ap 60\% sastādīja poļi un vācieši. 1936.~gadā Ukrainā tika likvidēti poļu nacionālie rajoni, tika nolemts pārcelt poļu un vācu saimniecības uz Kazahstānu. Pēc P.~Poļana datiem 1936.~gadā no abiem jau minētajiem Ukrainas apgabaliem no 800 metru pierobežas zonas, poligonu un nocietināto rajonu teritorijām pārvietoja vēl apmēram 45~tūkstošu poļu un vāciešu uz Kazahiju, kur bija skarbāki klimatiskie apstākļi un mazas iespējas iekārtoties darbā.

No Baltkrievijas poļus pārvietoja uz Sibīriju un Komi autonomo republiku. V.~Zemskovs, atsaucoties uz arhīva dokumentiem, uzskata, ka 1936.gadā pavisam no PSRS-Polijas pierobežas tika izsūtīti 69~283 poļi. Kā pamats represiju piemērošanai kalpoja aizdomas par ``neuzticamību'' vai radu esamība Polijā. Tika arī izvirzīti standartveida apvainojumi par sadarbību ar Polijas izlūkdienestu.

Izsūtīto poļu dzīve lielā mērā bija atkarīga no viņus uzņemošo apgabalu iedzīvotājiem. Neraugoties uz padomju varas iestāžu nežēlību, vienkāršie ļaudis izsūtītos sagaidīja bez naida, bieži dalījās ar savu jau tā trūcīgo mantu.

Vēl nežēlīgāka izrēķināšanās ar daudziem poļu tautības iedzīvotājiem norisa 30.~gadu beigās staļinisko represiju kulminācijas laikā. Rietumukrainas un Rietumbaltkrievijas Komunistisko partiju mēģinājumi izvirzīt šo zemju pašnoteikšanās lozungus izsauca pret tām represijas un apvainojumus par mēģinājumiem saķildot PSRS ar Poliju. Starptautiskās vēsturiski izglītojošās, tiesību aizsardzības un labdarības biedrības ``Memoriāls'' (\rutxti{Международное историко-просветительское, правозащитное и благотворительное общество ``Мемориал'')} Poļu komisijas pētnieki apgalvo, ka līdz 1939.~gada 17.~septembrim par staļinisko represiju upuriem PSRS kļuva vairāk nekā 180 tūkstošu poļu un ieceļotāju no Polijas vai cilvēku, apvainotu par sakariem ar Poliju. (Pēc citiem datiem, trīsdesmitajos gados par masveida represiju upuriem kļuva vairāk nekā 240~tūkstošu poļu, no tiem 110~tūkstošu tika nošauti). Taču 20.~gados un 30.~gadu sākumā tas viss vēl bija priekšā, PSRS it kā pozitīvais piemērs ietekmēja arī nacionālo mazākumu noskaņojumu Polijā.

Jāatzīmē, ka Polijas valdošie slāņi ar aizdomām un nepatiku raudzījās uz tautiešiem Padomju Savienībā. Tikai divi piemēri. 1929.~gadā Polijā no 16 maija līdz 30.~septembrim Poznaņā tika organizēta Vispārējā nacionālā izstāde (poļu: \pltxti{Powszechna Wystawa Krajowa w Poznaniu}, ang.: \entxti{Polish General Exhibition}), pieskaņota Polijas neatkarības 10-gadei. To apmeklēja aptuveni 4,5~miljoni skatītāju, tostarp 200~tūkstoši ārzemju viesu un gandrīz 1~000 žurnālisti, kas pārstāvēja laikrakstus no 30~valstīm. Šim laikam bija pieskaņots arī to poļu kongress Varšavā, kuri dzīvoja aiz valsts robežām. Organizācijas komiteja ielūdza poļus no visas pasaules valstīm, arī PSRS. Padomju poļiem bija iedalītas 14~vietas. PSRS izvērsās aktīva kampaņa lai sagatavotos kongresam. Padomju poļu vidū tika vākti līdzekļi braucienam uz Varšavu. Protams, tika domāts gan par poļu dzīves PSRS, gan visas padomju iekārtas propagandu kongresā. Taču Organizācijas komiteja nobijās no šādas ideoloģiskas iedarbības un ar ieganstu, ka padomju delegācija ievēlēta nedemokrātiskā ceļā (2\% no visiem padomju poļiem~--- t.s. bijušajiem ``ekspluatatoriem, kulakiem un nepmaņiem'' bija atņemtas vēlēšanu tiesības), atteicās pieņemt delegāciju. Rezultātā ap 800~000 PSRS dzīvojošo poļu tika atņemta iespēja nosūtīt savus pārstāvjus uz tautiešu kongresu.

Tāpat tika mēģināts strikti ierobežot Polijas pilsoņu braucienus uz PSRS. 1927.~gadā pēc kreiso partiju iniciatīvas tika izveidota ``CK Polijas strādnieku un zemnieku delegācija nosūtīšanai uz PSRS'' (\pltxti{CK dla wysłania delegacji robotnikow i chłopów Polski do ZSPR}). Sakarā ar Oktobra apvērsuma 10.~gadadienu tika izveidota delegācija no 12~cilvēkiem. Taču tās locekļiem tika aizliegts izbraukt. Tomēr daļa delegācijas locekļu nelegāli devās uz PSRS, piedalījās Maskavā PSRS draugu kongresā, apmeklēja virkni Ukrainas un Baltkrievijas centru. Vēl viens mēģinājums nosūtīt delegāciju uz PSRS notika 1929.~gadā. Un atkal atļauja netika saņemta, delegācija atkal nelegāli, pa apkārtceļu caur Dancigu, devās uz PSRS. Ne visi no ievēlētajiem 100~delegātiem spēja šķērsot robežu. Tomēr daudziem tas izdevās, viņi tika sadalīti grupās, bez Maskavas un Ļeņingradas apmeklēja arī Pievolgas, Kaukāza, Baltkrievijas, Donbasa, Centrālās Āzijas pilsētas. Atgriezušies viņi stāstīja par redzēto līdzpilsoņiem. Valsts iestādes šo darbību uzskatīja par ``graujošu'', apsūdzēja un bargi notiesāja nepakļāvīgos.

Nacionālo mazākumu apspiešana, smagais ekonomiskais stāvoklis Rietumu Ukrainā radīja labvēlīgu augsni ukraiņu nacionālistu ietekmes pieaugumam. Vietējā ukraiņu nacionālajā kustībā Polijas teritorijā arvien stiprākas pozīcijas ieguva galējās nacionālistiskās grupas ar antipolisku un antisemītisku virzību. Kā atzīmējis nacionālo attiecību pētnieks Visvaldis~Lācis: ``Tā kā ukraiņu protesti toreizējā Tautu Savienībā Ženēvā un arī vecvecās ukraiņu patiešām nacionāl-liberālās partijas deputātu protesti Polijas Seimā nekādus rezultātus nedeva, nebija jābrīnās, ka poļu policijiskais terors un starptautisko līgumsaistību neievērošana sekmēja ukraiņu nacionāl-radikālas kustības nodibināšanos, kuras sauklis bija tāds, ka pret teroru jāatbild ar teroru.''

Ukraiņu nacionālisti izveidoja savas nelegālas cīņas organizācijas. 1920.~gadā nodibinājās ``Ukraiņu kara organizācija'' (ukr. \uktxti{Ukraïnska Viйskova Orгanizacija}), kuras vadība atradās Vīnē, bet pati tā darbojās Austrumgalīcijā ar teroristiskām metodēm. Tieši tās kaujinieki veica neveiksmīgus atentātus 1921.~gadā pret J.~Pilsudski un 1924.~gadā pret prezidentu S.~Vojcehovski.

Jau atzīmēts, ka sevišķi grūts nacionālo nomaļu zemnieku stāvoklis kļuva pasaules ekonomiskās krīzes gados. 1930.~gadā Rietumukrainā sākās stihiskas zemnieku sacelšanās. Izputējušie zemnieki, nonākuši uz izdzīvošanas robežas, uzbruka muižnieku īpašumiem, nodokļu ievākšanas iestādēm. Pēc poļu pilsoniskās preses atzinumiem tikai 1930.~gada pēdējās augusta dienās zemnieki Rietumukrainā nodedzināja 62~muižas. Dedzināšanas turpinājās arī turpmāk. Radikāli noskaņoto ukraiņu nacionālistu idejas rada atbalsi jaunatnē.

Mērenie ukraiņu nacionāl-demokrātiskie spēki, kuri iestājās par Austrumgalīcijas autonomiju, boikotēja vairākumu vēlēšanu, tomēr bija gatavi noslēgt ar valdību t.s. ``cietokšņa mieru'', kas ļautu viņiem līdzrunāt pārvaldes jautājumos. Mēģinājumi rast sapratni poļu un ukraiņu starpā, ko propagandēja atsevišķi Polijas ierēdņi un Ukraiņu grieķu-katoļu (uniātu) baznīcas galva 1900.--1944.~gadā Galīcijas metropolīts Andrejs, rada izpratni vienkāršo ļaužu vidū, taču izsauca niknu pretestību abās pusēs nacionālajām barikādēm esošajos radikāļos. Mērenība neapmierināja samērā lielu daļu iedzīvotāju, kuri vairāk sliecās atbalstīt aktīvu cīņu pret apspiedējiem.

1929.~gadā no ``Ukraiņu kara organizācija'' J.~Konovaļca vadībā tika radīta ``Ukraiņu nacionālistu organizācija''~--- OUN (\uktxti{Організація Українських Націоналістів)}, kas veda enerģisku nelegālu cīņu pret Polijas varas iestādēm, faktiski karu pret Polijas valsti. Drīz izveidojās arī radikāla jauniešu grupa, kuru vadīja S.~Bandera. OUN atbalstīja, arī finansiāli, daudzas ukraiņu emigrantu organizācijas ārzemēs. Savu finansiālo stāvokli uzlabot un darbību aktivizēt Ukraiņu nacionālistu organizācija centās arī ar Polijas valsts (pirmkārt banku un pasta iestāžu) un privāto īpašumu aplaupīšanu. ``Ukraiņu kara organizācija'' iekļāvās tajā kā kaujas organizācija, izvēršot aktīvus uzbrukumus poļu ierēdņiem, pasta pārvadājumiem. Piemēram, tikai 1930.~gada jūnijā tika sarīkoti 220~ugunsgrēki (dedzinātas tika galvenokārt poļu kolonistu sētas). 31.~jūlijā vien notika trīs uzbrukumi pasta transportiem.

Ukraiņu nacionālisti uzsver, ka šīs abas ukraiņu nosauktās organizācijas radās kā ukraiņu sabiedrības radikālo elementu atbilde uz poļu apspiedēju politiku pret ukraiņiem, atbilde uz teritoriju polonizāciju, kuras ukraiņi uzskatīja par savām. Tā kā viņu protesti parlamentā, presē un publiskos pasākumos nedeva rezultātu, viņi ķērās pie ārkārtējiem līdzekļiem. Tāpēc tika dedzinātas poļu kolonistu, poļu skolotāju mājas, saimniecības ēkas, spridzināti tilti, bojātas sakaru līnijas u.tml. Kopā 1930.~gadā Austrumgalīcijā tika nodedzinātas ap 2~200 poļu mājas.

Šeit gan jāsaka, ka Ukraiņu nacionālistu organizācija teroru vērsa ne tikai pret oficiālajām Polijas iestādēm un poļu iedzīvotājiem, bet arī pret ukraiņu darbiniekiem, kuri vēlējās atrast saprašanos ar Polijas valdību, tāpat arī pret komunistiem, pret ukraiņu kreisajām organizācijām, vietējās inteliģences pārstāvjiem, visiem, kuri nepiekrita radikāli nacionālistiskajiem uzskatiem. Organizācija centās darboties arī PSRS teritorijā.

Atbildot gan uz zemnieku kustību, gan uz ukraiņu nacionālistu teroru, Polijas valdība 1930.~gadā uzsāka t.s. \strong{``pacifikācijas politiku''}. Tā bija virzīta uz vardarbīgu pakļaušanās uzspiešanu, uz nacionālo mazākumu piespiedu nomierināšanu, vismaz daļēju nacionālo mazākumu politisko, saimniecisko un kulturālo iniciatīvu, interešu un tiesību apspiešanu.

1930.~gada septembrī J.~Pilsudskis nolēma izmantot praktiski valstī pastāvošo ārkārtējo stāvokli, kad parlaments bija atlaists, kas sabiedrību atstāja gandrīz pilnīgi vienaldzīgu, kad tika gatavoti daudzu režīma kritiķu aresti, lai ar spēku apspiestu ukraiņu nacionālās atbrīvošanās kustību austrumu Galīcijā. 1.~septembrī viņš pavēlēja pēc 15~dienām sākt provinces ``pacifikāciju''. Ieganstu tam deva 1930.~gada jūlijā Ukraiņu kara organizācijas ar terora līdzekļiem uzsāktā cīņa pret provinces polonizāciju.

Soda operāciju veica policijas un armijas vienības un tā ilga līdz 30.~oktobrim. Pēc oficiālajiem Polijas Iekšlietu ministrijas datiem pacifikācijas akcijas tika veiktas 16~Galīcijas apriņķu (povjatu) 450~ciemos, 53~ciemos~--- arī atkārtoti.

``Pacificētāji'' pirmkārt vērsās pret visu, kas atgādināja par vietas ukrainisko piederību, norāva izkārtnes ar ukraiņu sabiedrisko un kultūras organizāciju nosaukumiem, demolēja bibliotēkas, lasītavas, kooperatīvus. Tika slēgtas trīs ukraiņu ģimnāzijas, arestēti bijušie deputāti. Ukraiņu nacionālās kustības aktīvisti tika publiski pērti, kratīšanu gaitā zemnieku sētās izdauzīti logu stikli, izbērti graudi, sabojāti sadzīves priekšmeti. Bija sieviešu izvarošanas gadījumi. ``Pacifikācijas'' gaitā tika pielietots kolektīvās atbildības princips. Pavisam tika arestēti 1~739~ukraiņu (pēc citiem datiem~--- ap 5~000), 1~143 no tiem tika apsūdzēti terorismā un pretvalstiskās aktivitātēs, tika konfiscēti aptuveni 2,5~tūkstoš vienību šaujamieroču. Pēc dažādiem datiem Austrumu Galīcijā poļu policijas un armijas kavalērijas vienības veica reidus 800~ciemos, tos izlaupot, ap 500~ēku nodedzināja, ap 50~cilvēku nogalināja, ap 4~000~ievainoja vai sakropļoja. Faktiski tā bija otrā, pēc bruņotā konflikta 1918.--1919.~gados, poļu-ukraiņu kara stadija Galīcijā.

Taču ``pacifikācijas'' operācija nesasniedza mērķi. Province netika nomierināta. Ukraiņu politiķi griezās ar sūdzību Tautu Savienībā un kaut tā arī nesaskatīja, ka Polija būtu pārkāpusi savas saistības nacionālo mazākumu tiesību ievērošanas jomā, Polijas starptautiskais tēls cieta. Polijas realizētā ``pacifikācijas'' kampaņa un pasaules sabiedrības vienaldzība pret to tikai nostiprināja daļā radikālo Galīcijas ukraiņu politiķu ticību tam, ka kompromisu ar Polijas administrāciju panākt nav iespējams, ka ukraiņiem ir nepieciešams meklēt palīdzību un atbalstu pie Polijas ienaidniekiem.

Tur, kur nacionālā kustība tika apspiesta īpaši vardarbīgi, nacionālās piederības pazemošana veda pie atbalsta pieauguma Ukraiņu nacionālistu organizācijas kaujiniekiem. Kopumā kolektīvās atbildības principa pielietošana pret tūkstošiem nevainīgu zemnieku veda pie ukraiņu iedzīvotāju naida pieauguma pret Polijas valsti.

1932.~gada jūnijā un jūlijā Ļvovas vojevodistes 19~ciemos sacēlās ap 30~000 cilvēku. Ieganstu radīja valsts ierēdņu un vietējā magnāta no Potocku (\pltxti{Potocky}) dzimtas mēģinājums bez iztikas līdzekļiem palikušo trūcīgo zemnieku darbaspēku darba klausības veidā izmantot ceļu celtniecībā. Zemnieki, apbruņojušies izkaptīm, ķēdēm, daži arī ar no Pirmā pasaules kara laika noslēptajiem ieročiem, pretojās. Sacelšanās apspiešanai tika nosūtīti policija, 4~000 karavīru un 3~lidmašīnas, kas bombardēja sacēlušos ciemus. Cīņas ilga 20~dienas. Ap 800~cilvēku tika arestēti. Vēl viena ``pacifikācijas akcija'' norisa tai pat 1932.~gadā Volīnijā. Tur izvērsās tik apjomīga partizānu kustība, ka tās apspiešanai bija jānosūta soda ekspedīcija triju kājnieku pulku, četru kavalērijas eskadronu, četru bruņumašīnu, vairāku lidmašīnu sastāvā. Šādās akcijās varas iestādes nereti sadūrās ar karavīru nepakļaušanos, nevēlēšanos piedalīties tajās. Tāpēc armijā atsevišķos kara apgabalos nācās ieviest lauku kara tiesas. 1932.~gadā Tautu Savienība nosodīja Polijas valdības soļus pret ukraiņu iedzīvotājiem, taču tas neko nemainīja.

Ievērojamā poļu sabiedriski-politiskā darbinieka J.~Kuroņa dzīves biedre D.~Kuroņa ir atstāstījusi gadījumu, kad viņas vīrs atcerējās savus bērnības gadus. Viņa tēvs bija paņēmis četrgadīgo Jaceku līdzi uz laukiem~--- kādu ukraiņu ciemu, kur poļu karaspēks un policija bija veikuši soda ekspedīciju. Mazā Jaceka atmiņā neizdzēšami iegūla ainas ar pienu, izlietu grāvjos, mājām, kuras tika iznīcinātas ar visu mantību, par jaunietēm, pakārtām aiz kājām, par pazemojumiem, kurus izcieta vīrieši. Vēlāk J.~Kuroņs par to runāja ar aizžņaugtu rīkli, it kā visu būtu redzējis savām acīm. Bet visu to viņš nevarēja redzēt, jo visiespaidīgākās soda ekspedīcijas ukraiņu ciemos notika vēl pirms viņa piedzimšanas, bērnam par to darbību tika tikai stāstīts. Taču armijas, policijas, ierēdņu iznīcināmās brīvības, sabiedrības neatkarības idejas no tā laika bērna acīs bija ieguvušas svētas lietas nozīmi, par kuru viņš cīnījās visu savu dzīvi.

1932.~gadā Ukraiņu nacionālistu organizācijas priekšgalā nonāca S.~Bandera un terors pret Poliju kļuva labi sagatavots. 1933.~gada septembrī OUN īstenoja t.s. ``skolu akciju'', kad tās biedru ietekmēti, bērni skolās atteicās atbildēt uz poļu valodā uzdotiem jautājumiem, aicināja poļu skolotājus atgriezties Polijā, izmeta no skolām poļu valstiskuma simbolus utt. 28.~septembrī, protestējot pret ukraiņu izglītības polonizāciju, OUN nogalināja skolu kuratoru.

1934.~gada 15.~jūnijā vispārējā S.~Banderas vadībā OUN organizēja atentātu pret Polijas iekšlietu ministru B.~Peracki, stingrās līnijas piekritēju attiecībās pret ukraiņu nacionālistiem un ``pacifikācijas'' politikas īstenošanas vadītāju. Kādā Varšavas kafejnīcā ukraiņu nacionālists G.~Maceiko viņam iešāva trīs lodes pakausī. Slepkavam izdevās aizbēgt uz Argentīnu, taču atentāta organizatori tika arestēti. Polijas prezidents I.~Moscickis pēc nāves piešķīra ministram brigādes ģenerāļa pakāpi un apbalvoja ar Baltā Ērgļa ordeni. Varšavā un Krakovā, sērojot par terora upuri, ielās izgāja tūkstošiem cilvēku. Kaut vainīgs bija ukraiņu nacionālists, vaina tika uzkrauta visiem nacionālajiem mazākumiem, kā arī labējiem poļu radikālnacionālistiem. Polijas kreisās organizācijas turpretī uzskatīja, ka aiz B.~Peracka slepkavības stāv poļu specdienesti un atentātu salīdzināja ar reihstāga aizdedzināšanu Vācijā 1933.~gadā, kam sekoja plašas represijas. Arestēti tika vairāk nekā 1~000 banderoviešu, arī pats S.~Bandera. Taču viņa rīkojumi turpināja darboties. Jūlijā pēc viņa pavēles tika nogalināts Ļvovas ukraiņu ģimnāzijas direktors I.~Babijs. Tas izsauca lielu rezonansi ukraiņu sabiedrībā. Visas legālās partijas nosodīja slepkavību. Metropolīts A.~Šeptickis publicēja rakstu, kurā paziņoja: ``\dots{}nav neviena tēva vai mātes, kuri nenolādētu vadītājus, kuri ved jaunatni noziegumu neceļos'', ``ukraiņu teroristi, kuri sēd drošībā aiz novada robežas, izmanto mūsu bērnus vecāku slepkavošanai, bet paši slavas oreolā priecājas par tādu izdevīgu dzīvi''. Ar to OUN romantiski-ideālistiskajam oreolam, kuru centās radīt pagrīde, tika dots spēcīgs trieciens.

1936.~gada sākumā notika atbildīgo par B.~Peracka nogalināšanu tiesas process, kurā S.~Banderam un vēl diviem viņa līdzgaitniekiem piesprieda nāves sodu, bet deviņiem cilvēkiem piesprieda cietumsodu no 7~līdz 15~gadiem. Process izsauca plašu rezonansi pasaulē. Polijas valdība neuzdrošinājās izpildīt spriedumu, uzsāka sarunas ar legālajām ukraiņu partijām par ukraiņu-poļu attiecību ``normalizāciju''. S.~Banderam u.c. nāves sodu nomainīja ar mūža ieslodzījumu. Faktiski S.~Bandera cietumā pavadīja piecus gadus, līdz pēc Vācijas uzbrukuma Polijai 1939.~gada septembrī tika atbrīvots.

Izmantojot B.~Peracka nogalināšanu kā ieganstu, reizē ar citiem pasākumiem Polijas valdība pret saviem politiskajiem pretiniekiem spēra soli, kuru nekādi nevar uzskatīt par tiesisku. 1934.~gada 17.~jūnijā notika Polijas valdības sēde, kurā tika pieņemts lēmums par \strong{Berezas-Kartuzskas} (mūsdienās~--- Baltkrievijas Republikas Brestas apgabalā) \strong{nometnes} (\pltxti{Miejsce Odosobnienia w Berezie Kartuskiej}) ierīkošanu. Tai pašā dienā prezidents I.~Moscickis parakstīja rīkojumu ``jautājumā par personām, kuras apdraud drošību, mieru un sabiedrisko kārtību''. Uz dokumenta bija vēl 11~ministru, tai skaitā arī J.~Pilsudska, paraksti. Dekrētā dibināmais iestādījums skaidri dēvēts par koncentrācijas nometni, pēckara poļu, kā arī padomju un krievu literatūrā arī lietots šis nosaukums, kaut mūsdienās daudzi poļu autori no tā izvairās. Poļesjes vojevoda V.~Kosteks-Bernackis 2.~jūlijā izdeva rīkojumu, ar kuru aizliedza apkārtējiem iedzīvotājiem tuvoties, fotografēt koncentrācijas nometni, uzturēt kontaktus ar tajā ieslodzītajiem.

Nometnei apkārt bija vairākkārtējs dzeloņstiepļu žogs ar sargtorņiem, kuros atradās ložmetēji. Par uzraugiem dienēja speciāli atlasīti policisti, kuri izcēlās ar īpašu nežēlību pret arestētajiem. Tiesa, kāds bijušais ieslodzītais vēlāk atcerējās, ka kāds policists, vērojot necilvēcīgās mocības, kādas cieta nometnes iemītnieki, nevarēja novaldīt asaras. Taču tas bija izņēmuma gadījums.

1966.~gadā Minskā iznāca bijušo nometnē ieslodzīto atmiņu un dokumentu krājums ``\rutxti{Они не стали на колени}'' (``Viņi nekrita ceļos''), kurā detalizēti parādīta nometnes iemītnieku atkarība no administrācijas patvaļas, kuras mērķis bija sadragāt ieslodzīto cilvēcisko pašcieņu, salauzt viņus fiziski un morāli, padarīt par paklausīgiem režīma rīkojumu pildītājiem.

Attiecībā pret ieslodzītajiem tika pielietoti dažādi fiziskās un psiholoģiskās ietekmes līdzekļi. Piemēram, par tādu kļuva viņiem dabīgo vajadzību nokārtošanas organizācija. Vienas kameras iemītnieki (ap 30~cilvēku) visi kopā tika sadzīti nelielā tualetē, kur varēja uzturēties ne ilgāk par minūti. Katram nācās iekārtoties, kā nu viņš varēja. Vissmagākā situācijā nonāca pēdējā ieslodzīto grupa. Kamēr rinda nonāca līdz viņiem, tualetē vairs nebija kur kāju nolikt. Tos, kas kaut sekundi aizkavējās tualetē, uzraugi nežēlīgi piekāva un lika rokām savākt ekskrementus un sasviest klozetpodos. Pēc tualetes cilvēki metās pie kulbas ar ūdeni. Taču arī mazgāšanās laiks bija ierobežots. Knapi apskalojuši netīrās rokas un saslapinājuši seju, ieslodzītie bija spiesti stāties ierindā un skriešus mesties uz ēdnīcu. Tur tāpat triju minūšu laikā nācās steigā norīt karsto duļķaino ūdeni (kafiju) vai skābo strebjamo, piekožot 400~g. maizes porciju. Visi protesti tika apspiesti. Apsargi ņirgājās ``\pltxti{Ty kurwo inteligencka, możesz obiad z gównem jeść}'' (``Tu inteliģentā kuce, vari ēst pusdienas ar sūdiem).'' Pēc tam sekoja jaunas komandas.

Arī pirts apmeklējums pārvērtās par mocībām. Iemīļota uzraugu izprieca bija kailu, tikko nomazgājušos cilvēku dzenāšana pa nometnes laukumu, spiežot pildīt dažādus fiziskus vingrinājumus. Nomazgāt pēc līšanas pa laukumu pielipušos dubļus uzraugi neļāva. Tika dota komanda apģērbties un ar to arī ``mazgāšanās'' bija galā. Šāda ``fizkultūra'' notika regulāri arī citā laikā. Par katru kļūmi, vilcināšanos sekoja sods it kā par nepakļaušanos pavēlēm~--- skanēja komanda: ``saliekt muguru'', lai ar gumijas stekiem varētu ērtāk piekaut ``vainīgos''. Ja komanda netika izpildīta, nepaklausīgais saņēma papildus sitienus. Īpašas mocības bija jāizcieš slimajiem, fiziski vājajiem, invalīdiem. Medicīniskā palīdzība viņiem faktiski netika sniegta, bet fiziskais vājums kalpoja kā papildus iemesls lai par viņiem paņirgātos. Piemēram, cilvēkus, kuri knapi varēja staigāt, spieda skriet. Kad viņi krita, tos ar kāju spērieniem piespieda piecelties un turpināja ņirgāšanos.

Nelielās kamerās ar cementa grīdu mita līdz 40 cilvēku. Svētdienās arestētajiem bija tiesības sēdēt vai gulēt blakus savam matracim, taču bija aizliegts pāriet uz citu vietu vai staigāt pa kameru. Tajā bija jāvalda pilnīgam klusumam. Protams, kad uzraugu nebija tuvumā, risinājās klusas sarunas, par kurām atkal draudēja sods~--- komandu ``gulties'', ``celties'' izpilde u.c. Vēl viens soda veids bija ievietošana karcerī, kur ieslodzītajiem nācās gulēt uz kailas grīdas, bet kādreiz arī visu laiku atrasties stāvus, jo grīda tika laistīta ar ūdeni, lai tie nevarētu apsēsties.

Ieslodzītie bija spiesti samierināties ar trūcīgu ēdienu (radinieki kādreiz gan saņēma atļaujas pārtikas paku nosūtīšanai nometnes iemītniekiem, taču viņu pašu rokās tās nenonāca, uzraugi rīkojās ar tām pēc saviem ieskatiem) un smagi strādāt, veicot bezjēdzīgu darbu~--- piemēram, rokot nevienam nevajadzīgas bedres, pārvietojot no vienas vietas un atpakaļ zemes kalnus. Jebkurš darbs tika pārvērsts par ņirgāšanos pār cilvēkiem. Par mēģinājumiem kaut brīdi atvilkt elpu sekoja iekaustīšana ar steku. Ieslodzītie savukārt mērķtiecīgi sabotēja pavēļu pildīšanu. Tā kaut kādā mērā tika taupīti viņu fiziskie spēki, tika stiprināta cīņas griba. Visnepatīkamākā bija komposta gatavošana, kad ar kailām rokām bija jāsajauc nopļauta zāle ar ekskrementiem, šis maisījums jānes uz uzraugu dārziem un tur jāizkaisa. Pēc šāda darba nedrīkstēja ne nomazgāties, ne arī nomainīt apģērbu. Uz komposta gatavošanu policisti parasti nosūtīja rakstniekus, ārstus, inženierus, komunistiskās partijas un arodbiedrību aktīvistus. Par mazāko pārkāpumu šos ieslodzītos spieda līst uz vēdera pa zemi, pārklātu ar kompostu.

Valdība centās nepieļaut informācijas par zvērībām nometnē nonākšanu sabiedrībā, taču īlenu maisā bija grūti noslēpt. Jau 1934.~gada 18.~novembrī senatoru grupa iesniedza valdības vadītājam pieprasījumu, kurā atainoja stāvokli nometnē un norādīja, ka Polijas Konstitūcijā ``sodi saistīti ar fiziskām mocībām, ir aizliegti''. Polijas Seimā daži pilsonības deputāti nometni sauca par ``XX~gadsimta apkaunojumu''. Ievērojot protestus Polijā un ārzemēs pret zvērībām Berezas-Kartuzskas nometnē, 1936.~gadā režīms tajā nedaudz tika mīkstināts. Arī 1938.~gadā, kad ar Kominternes Izpildu komiteja lēmumu tika atlaista Polijas Komunistiskā partija, nometnes administrācija, acīmredzot cerot uz komunistu pretošanās izbeigšanu, uz dažām dienām pārtrauca represijas. Taču, kaut no brīvības tika atgādāts lēmums par nometnē pastāvošās partijas organizācijas atlaišanu, tās biedri nolēma to neatlaist, darīt visu, lai saglabātu revolucionāro garu savās rindās. Tad arī ieslodzīto mocības atjaunojās ar dubultu spēku.

Nometnē ieslodzītos atļāva trīs mēnešus turēt bez tiesas sprieduma, vienkārši pēc policijas vai vojevodistes vadītāja rīkojuma. Juridisku līdzekļu, kā izvairīties no ieslodzījuma nometnē, vienkārši nebija. Administrācijai bija tiesības ieslodzījumu pagarināt, atstājot cilvēkus tajā vēl uz trijiem mēnešiem. Tā faktiski cilvēkus nometnē varēja turēt nenoteiktu laiku. Trīs mēneši nereti pārvērtās par gadiem. Kaut pēc dekrēta ieslodzījums nometnē neskaitījās soda līdzeklis, faktiski tas bija daudz bargāks un mokošāks nekā parasts cietumsods. Ieslodzītie varēja panākt savu atbrīvošanu, ja viņi atteicās no ``pretvalstiskās'' darbības un parakstīja attiecīgu solījumu. Taču radikālo organizāciju biedru vidū tā skaitījās nodevība un parakstīt deklarāciju nozīmēja zaudēt pašcieņu, arī savu līdzbiedru uzticību. Tikai mazākā ieslodzīto daļa parakstīja attiecīgās deklarācijas. Turpretī citos naids pret ``sanācijas'' režīmu kļuva tikai dziļāks.

Ieslodzīto skaits svārstījās no 100~līdz 900. Tā, 1934.~gadā nometnē atradās ap 200~ieslodzīto, tai skaitā 120~ukraiņu nacionālistu, 40~galēji labējās poļu organizācijas Nacionālā radikālā nometne (\pltxti{Obóz Narodowo-Radykalny}) biedru un ap 40~RUKP biedru. Vēlāk ieslodzīto sastāvs mainījās un palielinājās komunistu īpatsvars.

Tika ieslodzīti arī cilvēki, kuri vienkārši kritizēja valdību. Tā, avīzes ``\pltxti{Słowo}'' (``Vārds'') redaktors, vēlākais Polijas valdības emigrācijā premjers S.~Cat-Mackevičs 1939.~gada 22.~martā tika arestēts un bez tiesas sprieduma 17~dienas turēts un spīdzināts Berezas-Kartuzskas koncentrācijas nometnē, jo viņš kritizēja Polijas ārpolitiku, pauda viedokli, ka Polijas armija ir vāja, un poļi varēs sevi aizstāvēt ne ilgāk kā divas nedēļas, kas drīzumā gandrīz burtiski piepildījās.

% page 554


1939.~gada 30.~augustā pirms Otrā pasaules kara sākās arī citu ``neuzticamo elementu'' aresti un Berezas-Kartuskas nometnē ievietoja no 1,5~līdz 2~tūkstošiem vāciešu, tai skaitā ap 500~sieviešu, no 2~līdz 3~tūkstošiem ukraiņu organizāciju biedru u.c. Arī kara laikā uzraugu izturēšanās pret ieslodzītajiem nemainījās. Piemēram, ieslodzītās baltkrievu sievietes 1939.~gada septembrī nometnē pavadītās nedēļas vēlāk atcerējās kā nelāgu murgu. Visā nometnē pavadītajā laikā viņas reizes trīs-četras saņēma virumu, toties maizi~--- ne reizes, smagi cieta no ūdens trūkuma. Tāpat kā vīriešus par mazāko it kā nepaklausību viņas tika piekautas. Otrās nedēļas beigās daudzas sievietes tiktāl zaudēja spēkus, ka nevarēja piecelties. Atbrīvošanas dienā vīrieši viņas no kamerām iznesas uz rokām. Vēl ilgi uz viņu ķermeņiem saglabājās no uzraugu stekiem nometnē iegūtie zilumu. Toties daži notriektie vācu lidotāji, kuri arī bija ievietoti Berezas-Kartuzskas nometnē, jutās tajā kā viesos. Viņi pilnā formā pastaigājās pa nometni, ar mutautiņiem māja pāri tai pāri pārlidojošām vācu lidmašīnām.

Nometne turpināja pastāvēt līdz pat brīdim, kad apsardze, baidoties no Sarkanās armijas, naktī no 17.~uz 18.~septembri to atstāja likteņa ziņā. To atklājuši, 18.~septembra rītā kreisi noskaņotie ieslodzītie izlauzās no kamerām un devās pretī Sarkanajai armijai. Vietējie iedzīvotāji viņus sirsnīgi uzņēma un pacienāja, kas liecināja par viņu simpātijām ``sanācijas'' režīma upuriem.

Pēc poļu vēsturnieka V.~Slešinska datiem piecos nometnes pastāvēšanas gados no nežēlīgās apiešanās mira 13~cilvēki. Pavisam visā pastāvēšanas laikā nometnē bija ieslodzīti līdz 10~000 cilvēku.

Pēc masu arestiem 1934.--1935.~gadā OUN vadība atradās ieslodzījumā (daļa patvērās ārzemēs), daudzas vietējās organizācijas pārtrauca pastāvēt. Tomēr vācu vēsturnieks M.~Aleksanders atzīst, ka 1936./1937.~gada ziemā attiecības ukraiņu un poļu vidū bija tik saspringtas kā pilsoņu kara priekšvakarā.

Jau teikts, ka Polijas valdība centās izskaust visu krievisko. Tas izpaudās arī reliģiskajā jomā. Valsts atbalstīja katolicismu pret pareizticīgo ticību. Taču cietēji bija galvenokārt ukraiņi un baltkrievi. 1933.~gada decembrī pareizticīgo Varšavas-Holmas eparhijas garīdznieku un laicīgo sanāksmē izskanēja skaitļi, ka Polijas valsts pastāvēšanas 12~gados pareizticīgie tās teritorijā zaudēja ne mazāk par 500~dievnamiem. Pavisam pēc ukraiņu un baltkrievu nacionālo organizāciju datiem 1~300 pareizticīgo baznīcas tika ``pārprofilētas'' par katoļu baznīcām, dažkārt pat ar asins izliešanu. Ticīgie bija spiesti noturēt dievkalpojumus privātmājās, kapličās un pat zem atklātas debess. Kad blakus PSRS 1929.~gadā boļševiki uzsāka jaunu karagājienu pret reliģiju, arī Polijā aktīvi izvērsās pareizticīgo baznīcas vajāšana. 1929.~gada vasarā katoļu baznīcas iestādes iesniedza Polijas tiesās 724~iesniegumus par pareizticīgo dievnamu un cita baznīcas īpašuma atsavināšanu. Tam pamatā bija it kā vēlme atgūt XIX gadsimtā netaisnīgi katoļu, bet īpaši uniātu baznīcai atsavināto un pareizticīgajām baznīcām nodoto īpašumu. Netika ņemts vērā, ka vēl senāk~--- XVI--XVII gadsimtā dievnami tika atsavināti pareizticīgajai un nodoti uniātu baznīcai. Tiesa gan, īpašu panākumu uniātiem nebija. Uniātu baznīcas kongresā 1927.~gadā tika minēts, ka visā Polijā to draudzēm bija klāt nākuši ap 20~000 jaunu locekļu, taču pēc diviem gadiem tādu bija tikai 16~500 cilvēku, tātad mazāk kā divus gadus iepriekš. Tai pat laikā notika arī pāreja pareizticībā. No 1923.~līdz 1933.~gadam pareizticīgo vidū bija ap 60~000 bijušo uniātu. Acīmredzot šādus rezultātus nepoļu vidū izsauca vardarbīgā cīņa pret pareizticīgo baznīcu.

Spilgti šīs cīņas uzplūdi norisa 1938.~gada jūnijā un jūlijā, kad Holmas novadā, kur gadsimtiem bija blakus dzīvojuši katoļi un pareizticīgie, ``pēc katoliskās sabiedrības'' pieprasījuma ar armijas un policijas spēkiem tika iznīcinātas ap 150~lauku pareizticīgās baznīcas. (V.~Lācis raksta, ka ``ar poļu Katoļu Baznīcas klusu piekrišanu sagrāva 200 ukraiņu dievnamu''.) Protams, par pareizticīgo dievnamu iznīcināšanu lasot, nav jāaizmirst arī tas, kā cariskajā Krievijā privilēģijas tika dotas pareizticīgajai baznīcai, katolicisms tikai paciests, pareizticības stiprināšanai tika mākslīgi radīta Holmas guberņa, taču katoļu dievnami tad netika noārdīti. Divdesmitajā gadsimtā civilizētā Eiropa (ārpus PSRS) vēl nepazina šādu vandālismu.

Ukraiņu nacionālisti cīņā pret Polijas valsti meklēja sabiedrotos ārpus tās robežām. Sadarboties ar Padomju Savienību, kura pauda atbalstu ukraiņiem Polijā, viņi nevarēja principiālu apsvērumu dēļ. Cerēt uz blakus esošo Čehoslovākiju arī nevarēja, jo tās sastāvā ietilpa Aizkarpatu Ukraina, kuras iedzīvotāji arī labprāt redzētu izveidojamies apvienotu visu ukraiņu valsti. Vienīgais iespējamais sabiedrotais varēja būt Vācija, kurai bija savi rēķini pret Poliju sakarā ar notikušajām robežu izmaiņām un vācu mazākuma stāvokli Polijā. Ukraiņu nacionālistu organizāciju centri atradās Berlīnē un saņēma nacistisko specdienestu palīdzību.

Blakus ukraiņu nacionālistiem cīņu pret Polijas varu izvērsa ukraiņu komunisti. 1924.~gadā Volīnijas komunisti pat uzsāka gatavot plašu bruņotu sacelšanos, kuru gan novērsa poļu policija, veicot plašus arestus. Tomēr, ievērojot Jaunās ekonomiskās un ``ukrainizācijas'' politikas panākumus blakus Polijai Padomju Ukrainā, pieauga arī Rietumu Ukrainas Komunistiskās partijas (RUKP) ietekme. 1923.~gadā tā iekļāvās Polijas komunistiskās partijas sastāvā. Partija gan darbojās pagrīdē, taču tās ietekmē atradās arī legālas organizācijas, 1923.--1924.gadā~--- Ukrainas sociāldemokrātiskā partija, bet 1925.--1932.gadā~--- Ukrainas zemnieku un strādnieku sociālistiskā apvienība (\uktxti{Українське Селянсько-Робітниче Соціалістичне Об'єднання}). Tiesa, padomju boļševiku vadoņi jaucās RUKP lietās, tajā bija domstarpības, šķelšanās, tā nespēja efektīvi darboties. Ar 20.~gadu beigām, ar kolektivizāciju PSRS un tās izraisīto briesmīgo badu Ukrainā, daļa Rietumukrainas komunistu atbalstītāju no tiem novērsās. 1933.gadā RUKP vadītājus izsauca uz PSRS un arestēja. Represēja daudzus rietumukraiņu komunistus, kuri dzīvoja PSRS. RUKP biedru skaits samazinājās no 4~300~līdz 2~600.

1926.~gadā Austrumgalīcijā un Volīnijā ukraiņu komunisti izveidoja legālu Ukraiņu strādnieku-zemnieku apvienību (ukraiņu \uktxti{Українське Селянсько-Робітниче Соціалістичне Об'єднання~--- ``Сель-Роб''}), taču 1932.~gadā ``sanācijas'' režīms to slēdza.

Sakarā ar Polijā pieņemto likumu par ``sabiedriski bīstamo elementu izolāciju'' un Berezas Kartuzskas koncentrācijas nometnes radīšanu RBKP un RUKP partija izvirzīja ideju radīt vienotu strādnieku un antifašistisku Tautas fronti. 1936.~gadā RUKP Centrālā komiteja izdeva manifestu ``Par brīvību, par laimīgu ukraiņu tautas nākotni'', kurā bija izklāstīta Rietumu Ukrainas antiokupācijas Tautas frontes platforma: komunisti atbalstīja Austrumgalīcijas autonomiju demokrātiskas Polijas sastāvā, uzskatot to kā etapu ceļā uz pašnoteikšanos. Taču radīt patiešām plašu tautas kustību komunistiem neizdevās nedz Baltkrievijā, nedz Ukrainā. Kā jau teikts, 30.~gadu otrajā pusē Polijas Komunistiskajai partijai briesmīgu triecienu deva staļiniskais režīms Padomju Savienībā. Daļa komunistu pēc partijas likvidācijas turpināja darboties arodbiedrībās, citās sabiedriskās organizācijās, taču iepriekšējās mērķtiecības viņu rīcībā vairs nebija.

No vienas puses, Polija tā arī palika unitāra valsts, kura liedza saviem nacionālajiem mazākumiem tiesības uz teritoriālu un kultūras autonomiju. (Atsevišķi vēsturnieki ir rakstījuši par Poliju arī kā koloniālu valsti ar iekšējām kolonijām valsts austrumos.) No otras puses, Polija starpkaru periodā tā arī nekļuva par integrētu valsti. Kara sākumā nacionālie mazākumi Polijā sastādīja 14~miljonu cilvēku jeb 40\% iedzīvotāju. 7~austrumu vojevodistēs (Ļvovas, Staņislavas, Ternopiļas, Volīnijas, Poļesjes, Novogrudskas un Viļņas), kas veidoja ap pusi no Polijas teritorijas, poļu bija krietni mazāk nekā puse no iedzīvotājiem.

Gan valsts ekonomiskā politika, gan galvenās poļu politiskās partijas savā darbībā ignorēja austrumu nomaļu iedzīvotāju vairākuma intereses. Jau mūsdienās franču-poļu pētnieks D.~Bovua konstatēja ``Polijas politika starpkaru periodā Podolijā un Volīnijā ievērojamā mērā turpināja to, kas [pirms tam] notika muižnieku saimniecībās. Valdība vadījās no pārliecības , ka ukraiņi~--- tā ir ``mūsu tauta''. Viņu labā daudz kas tika darīts [piebildīsim~--- no poļu viedokļa~--- V.Š.], taču tomēr par maz, lai apmierinātu viņu gaidas \citespace{} Daudzu gadu garumā neizdevās samierināt poļu cerības ar tautu cerībām, kuras apdzīvoja kresus.'' Tāpēc nebūt nenācās brīnīties, ka terminu ``tautu cietums'', kuru pirmo reizi pielietoja marķīzs A.~de Kustins grāmatā ``\frtxti{La Russie en 1839}'' (``Krievija 1839.~gadā'') attiecībā pret Krieviju (vēlāk tā sāka saukt arī Austroungāriju), čehu autors J.~Vozka 1932.~gadā attiecināja uz tā laika Poliju, tā arī nosaucot savu grāmatu ``\cstxti{Polsko žalář národů}'' (Polija~--- tautu cietums).

Zem poļu nacionālā jūga nacionālo mazākumu nacionālisms tikai auga, meklēja atbalstu ārzemēs (Vācijā, PSRS, arī citās valstīs), tas veidoja bīstamu materiālu jebkurām pret Poliju vērstām intrigām. Poļu valstiskuma pretinieki (vācietis R.~Vingendorfs) arī 1939.~gadā apgalvoja, ka arī tad tā esot bijusi tikai karikatūra par valsti, no dažādiem ielāpiem vienkārši kopā sadiegts, bet nevis līdz vienotībai un kopīgai valstiskai dzīvei izaudzis veidojums. Īpaši neveiksmīga kara gadījumā Polijai vienmēr vajadzēja rēķināties, ka tās armijas aizmugurē var uzliesmot nacionālo mazākumu, pirmkārt ukraiņu, sacelšanās. Iesaukti poļu armijā daudzi ukraiņi izdevīgā brīdī bija gatavi pavērst ieročus pret apspiedējiem.

Nobeigt stāvokļa analīzi var ar poļu vēsturnieka Z.~Zaluska konstatējumu: ``Pretēji Dmovska cerībām Polija izrādījās nespējīga sev piesaistīt un asimilēt nacionālos mazākumus Austrumos~--- ne nacionālā, ne kulturālā, ne ekonomiskā nozīmē. Tai pat laikā ar savu klātbūtni tā veicināja nacionālistisku izkropļojumu attīstību šo tautu atbrīvošanās kustībās, savstarpēju naidu.\citespace{} Baltkrievu tieksme Minskas, bet ukraiņu~--- Kijevas virzienā kļuva par neatgriezenisku procesu, kas pieauga ar šo abu tautu vispirms jau zemnieku, bet zināmā mērā arī strādnieku, nacionālās un šķiriskās apziņas nostiprināšanos.'' Nākotne pilnībā apstiprināja citēto spriedumu.

\section{Tautsaimniecības attīstība}

\epigraph
{Kapitālista un strādnieka vienlīdzība ir tāda pat, kā divu cīkstoņu vienlīdzība, no kuriem vienam ir sasietas rokas, bet otram iedots ierocis rokā.}
{Ļevs Tolstojs (\rutxti{Лев Николаевич Толстой})}

\epigraph
{Tautas galēja nabadzība gandrīz vienmēr ir tās vadoņu noziegums.}
{Pjērs Buasts (\frtxti{Pierre Boiste})}

\epigraph
{Likums ir kā zirnekļa tīkls: mazie kukainīši tajā iet bojā, lielie izķepurojas.}
{Tautas gudrība}

\epigraph
{Mēs visi zinām~--- gandrīz katra pārliecīga bagātība ir veikli maskētā veidā nozagta nabagajiem.}
{Aksels Munte (\detxti{Axel Martin Fredrik Munthe})}

\epigraph
{Bagātība tā ir daudzu uzkrājumi viena rokās.}
{Julians Tuvims (\pltxti{Julian Tuwim})}

\epigraph
{Taisnība ir izsalkušo pusē.}
{Viktors Igo (\frtxti{Victor-Marie Hugo})}

\epigraph
{Nacionālās krīzes īpatnības. Slīkst ne visi, bet glābti tiek tie, kuri neslīkst.}
{Staņislavs Jankovskis (\rutxti{Станислав Янковский})}

\epigraph
{Es ticu, ka bankas mūsu brīvībām ir bīstamākas par armijām.}
{Tomass Džefersons (\entxti{Thomas Jefferson})}

\epigraph
{Paši lielākie cūkas parasti prasa no cilvēkiem, lai viņi būtu eņģeļi.}
{Juliāns Tuvims (\pltxti{Julian Tuwim})}

\epigraph
{Cilvēku par cilvēku padara tikai sociālās nevienlīdzības apzināšanās.}
{Aleksandrs Bloks (\rutxti{Александр Александрович Блок})}

\epigraph
{Tiem, kuri uzlikuši acīm klapes, ir jāatceras, ka komplektā ietilpt arī iemaukti un pātaga.}
{Staņislavs Ježijs Lecs (\pltxti{Stanisław Jerzy Lec})}

\epigraph
{Katrs godīgs cilvēks izjūt kaunu par valdību, pie kuras dzīvo.}
{Henrijs Luis Menkens (\entxti{Henry Louis Mencken})}

% page 558


Polijas ekonomika starpkaru periodā sāka savu patstāvību zem hiperinflācijas, pēckara rūpniecības panīkuma un izpostītās lauksaimniecības sloga. Polijas oficiālā statistika ziņoja, ka kara laikā bija nodedzinātas, sagrautas 1~800~000 ēkas, palikuši 4~500~000~ha izpostītas, atmatā stāvošas zemes.

Pēc pirmā Pasaules kara Polija bija agrāri-industriāla valsts.

1921.~gadā lauksaimniecībā bija nodarbināti 72,3\% visu darba spējīgo iedzīvotāju. Pēc 1927.~gada datiem, kad valdīja salīdzinoši labvēlīga ekonomiskā konjunktūra, 64\% iedzīvotāju bija nodarbināti lauksaimniecībā un tikai 9,2\% rūpniecībā, 7,3\% amatniecībā un 6,4\% tirdzniecībā un transportā. No 1921.~līdz 1939.~gadam lauku iedzīvotāju īpatsvars samazinājās no 75~līdz 70\%. Tā kā poļu rūpniecības produkcijas agrākie noieta tirgi kļuva nepieejami, bet jaunus vājā Polijas industrija nevarēja iekarot, tai bija nepieciešams pietiekami plašs iekšējais tirgus. Zemniecība kā reiz arī varēja kļūt par rūpniecības preču galveno pieprasītāju. Agrārā jautājuma atrisināšanai bija gan politiska, gan ekonomiska nozīme. No muižniecības zemes īpašuma likvidācijas un zemnieku saimniecību nostiprināšanas bija atkarīga preču attiecību veidošanās laukos un arī tautsaimniecības attīstība kopumā.

Starp dažādiem valsts reģioniem lauksaimniecībā eksistēja milzīgas atšķirības.

Ja Polijas vidienē salīdzinoši sekmīgi attīstījās vidējo zemnieku saimniecības, tad Rietumgalīcijā knapi vilka dzīvību mazsaimniecības un pundursaimniecības, turpretī austrumu rajonos pastāvēja galvenokārt ekstensīvi apsaimniekojamas muižnieku lielsaimniecības.

Pa zemes ceļiem no rietumiem uz austrumiem Poliju šķērsojošie ceļinieki gandrīz ik pēc simt kilometriem varēja vērot jaunas, atšķirīgas zemnieku saimniecību iezīmes. Valsts rietumos varēja skatīt poļu zemniekus, kuri gan atšķīrās no šo apgabalu vācu zemniekiem, taču daudz mazāk uzkrītošā veidā nekā aiz bijušās agrākās Polijas karalistes robežas dzīvojošie. Aiz šīs bijušās robežas lauksaimniecības darba rīki bija citādi; neraugoties uz aramzemes līdzību, tās apstrāde atšķirīga; arī mājlopiem bija cits izskats. Viss bija mazāks, sliktāk kopts nekā rietumos. Visur vēl bija vērojami trīslauku sistēmas pielietošanas piemēri, lauksaimniecības mašīnas gandrīz netika izmantotas. Poļu zemnieki vēl joprojām audzēja tās pašas augumā mazās govis, kā Pirmā pasaules kara laikā. Turpinot ceļojumu, Vislas apgabalā pēkšņi parādījās zemnieku saimniecības ar citu raksturu. Tie bija bijušo vācu ieceļotāju ciemati, kuri paši gan bieži jau runāja poļu valodā. Tomēr viņus un viņu saimniecības ražības, labklājības ziņā varēja atšķirt no poļu kaimiņu sētām. Tālāk uz austrumiem vietējo iedzīvotāju dzīves veids bija vēl atšķirīgāks. Zemnieku vietā, kuri lepojās ar savu zemes īpašnieka stāvokli, parādījās Austrumpolijas pundursaimniecību īpašnieki, kuri bija vairāk mājamatnieki, nekā zemnieki. Viņiem zemes knapi pietika, lai nodrošinātu sev pārtiku.

Polijas lauksaimniecība, neraugoties uz labvēlīgākiem klimatiskajiem un citiem dabas apstākļiem nekā daudzās citās Eiropas valstīs, atradās vienā no pēdējām vietām Eiropā ražības ziņā.

Pēc 1921.~gada datiem, kuros gan vēl nebija iekļauti Augšsilēzijas un Viļņas apgabali, kopumā vidējo un mazsaimniecību īpašnieki sastādīja 87\% no visiem zemes īpašniekiem, taču viņu īpašumā bija tikai 31,9\% no visas privātīpašumā esošās zemes. Pēc padomju vēsturnieku datiem Polijā rentablas bija tikai tās saimniecības, kuru valdījumā atradās ne mazāk kā 5~ ha zemes un 64\% no visiem zemes īpašniekiem bija trūcīgie zemnieki. Turpretī 19~000 cilvēku (0,6\% no visiem), kuru īpašumā bija vairāk nekā 100~ha zemes, kopumā piederēja vairāk 13~miljonu ha (44,8\%) no visas privātīpašumā esošās zemes. (Pēc citiem datiem 4~255~muižu, kuras bija lielākas par 500~ha, īpašnieku rokās atradās 31\% lauksaimniecībā izmantotās zemes.) Skat. \ref{tab:table5}. tabulu.

\noindent
\begin{table}[h!]
\caption{Saimniecību skaits un lielums Polijas lauksaimniecībā XX~gs. 20.~gadu 2.~pusē} \label{tab:table5}
\begin{tabularx}{\linewidth}{|p{4cm}|p{4cm}|p{4cm}|}
\hline
\strong{Saimniecību lielums hektāros} & \strong{Saimniecību skaits absolūti} & \strong{Dotās kategorijas saimniecību skaits \% no kopējā skaita} \\
\hline
Lielākas par 100~ha & 18~916 & 0,58 \\
\hline
No 50~līdz 100~ha & 11~163 & 0,34 \\
\hline
No 20~līdz 50~ha & 76~436 & 2,34 \\
\hline
No 10~līdz 20~ha & 311~529 & 9,55 \\
\hline
No 5~līdz 10~ha & 733~526 & 22,49 \\
\hline
Zem 5~ha & 2~110~609 & 64,7 \\
\hline
Kopējais saimniecību skaits & 3~261~909 & 100 \\
\hline
\end{tabularx}
\end{table}

% page 559


Sadrumstalotās sīksaimniecības radīja agrāro pārapdzīvotību, slēptu bezdarbu, nedeva iespēju pienācīgi apgādāt ar pārtiku pilsētas un eksportēt lauksaimniecības produkciju. Seims pārskatīja jau 1920.~gada 15.~jūlijā pieņemto likumu par \strong{agrāro reformu}. 1920.~gadā to pieņēma tāpēc, ka uz Varšavu virzījās Sarkanā armija un vajadzēja kādu propagandas ieroci pret boļševikiem. Tā paredzēja pakāpenisku lielo zemes īpašnieku~--- muižnieku zemes gabalu sadali sīkās saimniecībās ar daļēju atlīdzību. Kad Sarkanā armija tika atsviesta, pirmkārt jau endeki agrāro reformu vērtēja par pārāk radikālu. Tā kā Konstitūcijā bija pasludināta privātīpašuma neaizskaramība, tas ļāva apstrīdēt agrārās reformas atbilstību Konstitūcijai.

Polijas pilsonība, kura agrāk bija nevis pašapliecinājusies cīņā pret feodālismu, bet pakāpeniski iespiedusies jau pastāvošajās ekonomiskajās un politiskajās struktūrās, arī tagad~--- valstiskās patstāvības apstākļos, neraugoties uz to, ka agrārās iekārtas modernizācija varētu sekmēt arī citus ekonomiskos procesus, palika nomaļus no cīņas par feodālisma palieku likvidēšanu laukos. Arī pēc atbrīvošanās no triju monarhiju uzspiestās kārtības saglabājās vēl pirms valstiskās neatkarības atjaunošanas realizētais kompromiss ar muižniecību.

Turpretī zemnieku partijas~--- ļudovieši centās izmantot savu pārstāvniecību parlamentā lai nostiprinātu zemniecības ekonomisko pamatu uz feodālo īpašumu rēķina. Tomēr ļudovieši~--- zemnieku partiju biedri tā arī nespēja īstenot savus priekšstatus pat taisnīgu agrāro iekārtojumu valstī. Poļu Zemnieku partija ``\pltxti{Wyzwolenje}'' bija par īpašumu formu daudzveidību lauksaimniecībā. Privātīpašumam bija jābūt valdošajam, taču vidēju un sīku zemnieku saimniecību formā. Lielos zemes īpašumus pēc šīs partijas biedru domām bija jāatsavina bijušajiem īpašniekiem un jāsadala starp bezzemniekiem un mazo saimniecību īpašniekiem par nenozīmīgu izpirkšanās maksu. Arī Poļu Zemnieku partija ``\pltxti{Piast}'' iestājās par piespiedu zemes atsavināšanu lielīpašniekiem un tās sadalīšanu bezzemniekiem un sīkzemniekiem. Perspektīvā pēc reformas par lauksaimniecības pamatu bija jākļūst 14~ha un lielākām zemnieku saimniecībām. Tātad bija jāsaglabā tikai tās saimniecības, kuras spētu ražot tirgum. Tajos gadījumos, kad to prasītu valsts un nācijas intereses, tika pieļauta arī valsts īpašuma uz zemi saglabāšana. ``\pltxti{Piast}'' programmā gan nebija atbilžu uz virkni jautājumu, piemēram, vai lielsaimniecības jālikvidē pilnībā, kas notiks ar vidēji lieliem muižnieku īpašumiem u.c.

Taču nepieciešamo radikālo agrāro reformu nebija iespējams realizēt, jo pret to kategoriski iebilda labējās partijas un muižniecība austrumos, kura tur pildīja polonisma balsta lomu. Tieši austrumu vojevodistēs atradās grāfu Zamoisku (\pltxti{Zamoyski}), kņazu Radzivillu (\pltxti{Radziwiłłowie}), grāfu Potocku (\pltxti{Potoccy}), kņazu Čartorijsku (\pltxti{Czartoryscy}) desmitiem līdz simtiem tūkstošu ha lielās muižas. 1920.~gada 15.~jūlijā Seima pieņemtais likums netika pildīts. 1925.~gada 28.~decembrī tas tika būtiski labots, mīkstinot muižniekiem nepieņemamo ``radikālismu''. Parcelācijai (zemes sadalīšanai sīkākos gabalos) ik gadu tika paredzēts nodot tikai 200~000~ha, jeb, citiem vārdiem sakot, piespiedu zemes atsavināšana par atlīdzību notika tikai tad, ja gada laikā zemes tirgū piedāvājums bija mazāks par 200~000~ha. Pēc ukraiņu vēsturnieka I.~Kundjubas datiem vidēji ik gadu parcelācijai tika nodoti ne vairāk kā 130~000~ha muižnieku zemju, bet starp 1932.~un 1938.~gadu zemniekiem tika sadalīti 85~tūkstoši ha zemes.

Ārpus rūpnieciskajiem apgabaliem nesadalītas varēja palikt 180~ha, bet austrumu rajonos pat 300~ha lielas saimniecības. Atvieglinājumus piešķīra saimniecībām, kuras nodarbojās ar kokmateriālu izstrādi, spirta dedzināšanu, cukura ražošanu utt. Tās varēja paturēt līdz 700~ha zemes. Poļu muižnieki Polijas austrumu rajonos, tāpat kā vācu muižnieki tās rietumu rajonos aktīvi cīnījās pret savas zemes atsavināšanu.

Pēc dažiem datiem starpkaru periodā parcelācijai tika pakļauti 654~000~ha zemes. Kopumā zemnieku rokās nonāca 2~600 tūkstošu ha, kas sastādīja mazāk par 7\% no visas zemes platības. Tomēr par 50~ha lielāku saimniecību, kuras sastādīja tikai 0,5\% (citi dati~--- 0,6\%) no visām saimniecībām, īpašumā palika 25,8\% apstrādājamo zemju un ap 45\% no visas zemes. Pie tam par iegūstamo zemi zemniekiem nācās ilgstoši maksāt. Pazīstamais ekonomists un politiķis V.~Grabskis atzina, ka ``procenti, kurus tagad maksā zemnieki, pērkot zemi parcelācijas rezultātā, bieži kļūst par slogu, smagāku nekā dzimtbūšana.''

Kā norādījis vācu vēsturnieks T.~Šīders, praktiski agrārreforma tika izmantota pret intensīvi saimniekojošajām vācu rokās esošajām lielsaimniecībām Rietumpolijā, kuras dažos gados zaudēja lielāko daļu (līdz 1928.~gadam~--- 93\%) no visā valstī parcelācijai nododamās teritorijas, kad poļu muižniecība, kuras rokās Polijas austrumos atradās veseli 4~miljoni~ha, zaudēja pavisam maz. Atsevišķi vācu autori īpaši sašuta par Rietumpolijā vācu lielīpašniekiem piederējušās zemes nodošanu ``tūkstošu'' poļu bezzemnieku no bijušās Polijas karalistes, Austrumpolijas un Galīcijas rokās. Ar to notikusi ``Āzijas~--- Eiropas robežas pārbīde uz rietumiem''. Arī Prūsijas poļi Poznaņas un Pomorjes novados ar naidu un nicinājumu skatījušies uz ienācējiem kā zemākas kultūras pārstāvjiem. Tāpat ar ``kongresa Polijas'' poļiem nevarējuši saprasties arī agrākajā Austroungārijas impērijā dzīvojušie goraļi jeb kalnieši (poļu \pltxti{górale}, slovāku \sktxti{gorali}, no poļu vārda ``kalns'' \pltxti{góra}, slovāku~--- \sktxti{hora)}, kuri mita Tatru kalnos un rutēņi no Galīcijas. Goraļi un rutēņi ar saviem lopiņiem gadsimtiem dzīvojuši cieši saistīti ar dabu, turpretī Viduspolijas zemniekiem ilgstošā dzimtbūšana esot laupījusi šo personīgo saikni ar zemi, uz kuras tie dzīvojuši, un dzīvniekiem, ar kuriem ik dienu tiem nācās saskarties. Protams, šie secinājumi nesa vācisko vērotāju personības zīmogu. Viņi bija ieinteresēti saskatīt jaunās Polijas valsts vājuma punktus, uzturēt cerības uz tās sabrukumu.

Tomēr agrārā reforma neapmierināja ne tikai vācu lielsaimniekus, bet arī lielu daļu poļu zemnieku. Polijā 47\% zemnieku saimniecību, kas aizņēma gandrīz 11~miljonus ha, cieta no skrejgabalu sistēmas paliekām. Pie tam 20\% no tām bija sadalītas 5~līdz 20~atsevišķos skrejgabalos. To likvidēšanas darbi norisa, taču diezgan lēni: no 1919.~līdz 1926.~gadam bija pārdalīti tikai ap 450~000~ha zemes. Zemes pārdale, jaunu zemnieku saimniecību ierīkošana palēnām notika galvenokārt centrālajā Polijā. Parcelācijas rezultātā radās arī daudz nerentablu zemnieku saimniecību. Augstā zemes cena, neraugoties uz pastāvošo kreditēšanas sistēmu, faktiski pārvērta par fikciju likumā pasludinātās bezzemnieku un mazzemnieku pirmās rokas tiesības, iepērkot jaunus zemes gabalus. No 1920~līdz 1938.~gadam tikai ap 735~000 zemnieku saimniecību saņēma zemes piedalījumu no agrārās reformas fonda. Tika radītas tikai 154~000 jaunas ap 9~ha lielas zemnieku saimniecības. 1931.~gadā pastāvēja ap 2~612~000~līdz 10~ha lielas saimniecības. Starp tām 750~000 bija pundursaimniecības, kuru lielums nesasniedza 2~ha. 1927.~gadā 43,8\% zemnieku saimniecību nebija zirgu, 30,7\%~--- cūku, 14,2\%~--- govju. Kā rādīja Finanšu ministrijas dati par 1927.~gadu, līdz 5~ha lielas zemnieku saimniecības vidēji nodokļos maksāja 2,38~zlotus zemes nodokļa no 1 hektāra, bet muižnieku latifundijas, lielākas par 2~000~ha,~--- 1,01~zlotu no ha.

Kaut pakāpeniski pieauga lauksaimniecības inventāra daudzums, minerālmēslojuma lietošana (Ja 1921.~gadā tos izmantoja 99~000~t, tad 1927.~~gada~--- jau 958~700~t, tātad~--- gandrīz 10~reizes vairāk), sīksaimniecības nevarēja dot labus rezultātus. Nav brīnums, ka Polija, neraugoties uz pietiekami labvēlīgajiem klimatiskajiem un augsnes apstākļiem, ražības ziņā atradās vienā no pēdējām vietām Eiropā.

Arī pēc Pirmā pasaules kara no Polijas turpinājās emigrācija. Ja pirms kara emigranti galvenokārt devās uz Vāciju un Amerikas valstīm, pēc tā sakarā ar emigrācijas ierobežošanu ASV un pēckara ekonomiskajām grūtībām pasaulē, galvenais emigrācijas virziens bija Francija. To regulēja starpvalstu vienošanās. 1924.~gadā no Polijas emigrēja 75~tūkstoši cilvēku, 1925.~gadā~--- 81~tūkstotis, 1926.~gadā~--- jau 168~tūkstoši. Pie tam no šiem 168~tūkstošiem 70~tūkstošu devās uz Franciju, 45~tūkstoši~--- uz Vāciju. Sākās arī sezonas poļu laukstrādnieku iebraukšana Latvijā.

Interesantas atmiņas par vietējo poļu un no Polijas iebraukušo algādžu attiecībām Latvijā (Liepājā) ir atstājis bijušais liepājnieks, aizsargs, vācu okupācijas laika Liepājas politiskās policijas priekšnieks, vēlāk emigrants Kārlis Siljakovs: ``Kad Latvijas pēdējos gados mūsu valdība kā darba spēku lauksaimniecībā sāka ievest poļu laukstrādniekus, likās, ka Liepājas poļu pašlepnums tā kā drusku sašķobījās un viņu saites ar Poliju palika vaļīgākas, bet tas viss bija tikai šķitums. Ievestos laukstrādniekus uzskatīdami kā savai kārtai nepiederīgus, viņi tos atstūma no savas vides un nereti lietas nogāja pat tik tālu, ka apkārtnes lauksaimniekam uz Liepāju līdz atbraukušais poļu laukstrādnieks bez sava saimnieka palīdzības, pat poļu veikalā nevarēja nopirkt kabatas nazi, jo veikala īpašnieks atsacījās ar viņu runāt poliski.''

Lielākais ārzemēs dzīvojošo poļu skaits tomēr joprojām palika Vācijā. Pēc tās varas iestāžu datiem Vācijā 1925.~gadā bija 803~000 poļu, pēc pašu poļu vērtējumiem~--- 1,5~miljoni. Tiesa, lauksaimniecības mehanizācijas rezultātā Vācijā ar 1918.~gadu gandrīz pārtrūka poļu sezonas strādnieku ieceļošana tajā. Tikai otrajā vietā pēc poļu skaita bija PSRS, pēc 1926.~gada tautas skaitīšanas datiem~--- ap 788~000, pēc 1939.~gada skaitīšanas~--- ap 630~000 cilvēku. Neraugoties uz Viļņas apgabala okupāciju un aneksiju, Lietuvā pēc oficiāliem datiem 1923.~gadā bija ap 66~000, pēc poļu vēsturnieka J.~Tomaševska vērtējuma~--- ``apmēram 200~000''. Ap 60~000 poļu dzīvoja Latvijā.

Tomēr, neraugoties uz emigrāciju, pēc dažiem aprēķiniem, ap 150~000~cilvēku ik gadu palielināja ``lieko darba roku'' skaitu. 1931.~gadā valstī uz 100~ha aramzemes vidēji bija 88,3~cilvēku, kas divas reizes pārsniedza lauku iedzīvotāju skaitu, kurš tajā laikā tika uzskatīts par maksimāli pieļaujamu. 1939.~gadā cilvēku, kuri dzīvoja laukos, pārsniedzot lauksaimniecības pieprasījumu pēc darbaspēka, bija jau 5,3~miljoni. Viņi pārtika no gadījuma darbiem, radinieku žēlastības. Auga to bezzemnieku slānis, kuri bija atkarīgi no muižnieku labvēlības, strādāja to labā par pārtiku, atļauju ganīt savu gotiņu mežmalā utt. Tāpēc nebija brīnums, ka jau 1929.~gadā ļudoviešu presē pirmoreiz izskanēja kritika pret kapitālismu kā ``ekspluatācijas un netaisnības'' iekārtu. Kā galvenais līdzeklis pret šiem ļaunumiem tika uzskatīta kooperācija.

Pirmā pasaules kara laikā kooperatīvā kustība, kura Polijā pirmos soļus spēra XIX gadsimta 60.--70.~gados, bija sagrauta. Pēckara periodā tā strauji atjaunojās. Tā, ja patērētāju kooperācijā 1914.~gadā bija ap 40~tūkstošu biedru, bet 1918.~gadā ap 61~tūkstoti, tad 1926.~gadā to bija jau ap 600~tūkstošiem. Lauksaimniecības kooperācija sāka izvērsties 1920.~gadā Galīcijā. 20.~gadu otrajā pusē Polijas Republikas lauksaimniecības kooperatīvu federācijā (\pltxti{Polska Federacja Spółdzielni Rolniczych}), kura apvienoja dažādus~--- kredīta, tirdzniecības un ražošanas kooperatīvus, sastāvēja jau 677~tūkstoši biedru. Taču būtiski uzlabot visu savu biedru stāvokli kooperatīvi nespēja.

Zemniecību nospieda pārlieku smags nodokļu slogs. 1924.~gadā austriešu rakstnieks J.~Rots apceļoja Galīciju un publicēja savus ceļojuma iespaidus. Viņš rakstīja: ``Ir grūti dzīvot. Galīcijai ir jāpabaro vairāk nekā astoņi miljoni iedzīvotāju. Zeme ir bagāta, iedzīvotāji ir nabagi. Viņi ir zemnieki, tirgotāji, sīkamatnieki, ierēdņi, karavīri, virsnieki, pārdevēji, banku kalpotāji, muižu īpašnieki. Ir par daudz tirgotāju, par daudz ierēdņu, par daudz karavīru, par daudz virsnieku. Visi galu galā dzīvi no vienīgās ražojošās šķiras~--- zemnieku rēķina''. Loti līdzīgu stāvokļa vērtējumu Polijas laukos 1927.~gadā deva Lielbritānijas parlamenta deputāts Dž.~Bekets, kurš bija apmeklējis Rietumukrainu: ``Mēs apmeklējām Volodimiru-Volinska (\uktxti{Володимир-Волинський}~--- apgabals tagadējās Ukrainas ziemeļrietumos). Es pazīstu Indiju, un jūs, protams, esat dzirdējuši par briesmīgu nabadzību Indijas laukos. Taču man nekad nav nācies redzēt tik nospiedošu un izmisīgu nabadzību \citespace{} Tagad mums kļūst skaidrs, kāpēc Polija uztur tik daudzskaitlīgu armiju''.

Valstī pastāvēja ap simts dažādu nodokļu veidu, bet zemnieki maksāja arī ap 60~dažādu nodevu un sodu maksājumu (par iebraukšanu pilsētā, par ratu stāvēšanu tirgū un uz ielas, par preces nosvēršanu, par lopu kaušanu, par sanitāri-veterināro apkalpošanu, par to, ka uz ratiem nebija uzkarināts reģistrācijas numurs utt.) Minskas Mazoveckas (\pltxti{Mińsk Mazowiecki}) apriņķī bija gadījums, kad zemnieks, lai atrastu pircēju, kam pārdot zirgu, bija spiests vairākkārt iebraukt pilsētā. Gala rezultātā viņš pārdeva zirgu par 12~zlotiem 50~grošiem, bet viņa nomaksāto nodevu summa sastādīja 10~zlotu. Fiskālā sistēma nesa viduslaiku iezīmes. Viduslaiku palieka bija atsevišķu valsts nodokļu ievākšanas tiesību pārdošana privātām personām. Tā notika ar ceļu nodokli Galīcijā un Cešinas Silēzijā.

Jau starpkaru periodā dažādi autori ar dažādiem politiskajiem uzskatiem kritizēja valsts politiku lauksaimniecības jomā.

Vācietis R.~Vingendorfs rakstīja: ``Saimniecisko politiku nenosaka poļu tauta. Poļu valsts kā tāda ir pilsētu inteliģences produkts. Tā nosaka arī valsts saimniecisko politiku. Taču inteliģencei ir tikai viens liels mērķis, padarīt savu valsti iespējami līdzīgu Rietumeiropai, arī ārēji radīt tādu pat izskatu, un beidzot pārvarēt mazvērtības sajūtu Rietumu priekšā. Tikai tā var izskaidrot, ka 20~pastāvēšanas gados poļu valstī, kurā ir daudz izmantojamas aramzemes un pārsvarā zemnieciski iedzīvotāji, pēdējie saņēmusi tik maz kā neviens cits slānis.'' Īpaši šis autors bija sašutis par to, ka gandrīz nekas netika darīts lai veicinātu lauksaimniecisko ražošanu Austrumpolijā, paceļot tās tehnisko līmeni no XVII gadsimta līdz XX gadsimta standartiem. Austrumpolijā turpināja pastāvēt feodālas muižas un zemnieku pundursaimniecības. Viduspolijā netika veikti nepieciešamie meliorācijas darbi, lai aizsargātos pret parastajiem pavasara plūdiem. Zemnieki ne Austrumpolijā, ne Viduspolijā netika agrotehniski izglītoti. Tāpēc arī esot radies stāvoklis, ka vēl 1938.~gadā 77,7\% no visa Polijas eksporta sviesta tika saražoti agrāk Prūsijai piederējušajos apgabalos, jo pārējos novados trūka tehnisko iekārtu lai saražotu eksportam derīgu sviestu.

Protams, ne viss tikko izskanējušajā kritikā atbilda patiesībai, ne visu pie labākās gribas varēja panākt divdesmit gados. Taču tam, ka Polijas valdošais slānis (R.~Vingendorfs to aplami identificējis ar inteliģenci) pirmkārt bija norūpējies par savas varas nostiprināšanu, nevis tautas dzīves līmeņa paaugstināšanu, diemžēl jāpiekrīt.

Par agrāro stāvokli Polijā starpkaru periodā interesantu vērtējumu ir sniedzis ASV vēsturnieks T.~Snaiders. Viņš norāda, ka trijām kaimiņvalstīm PSRS, Polijai un Vācijai bija jārisina kopēja politekonomiska problēma~--- ko darīt apstākļos, kad pastāvīgi draudēja bada iespēja, bija dārgi pārtikas produkti, valstu austrumos pastāvēja atpalikusi tautsaimniecība, orientēta uz laukiem. Padomju lēmums bija kolektivizācija jeb zemnieku pārvalstiskošana. Savukārt Ā.~Hitlers problēmas risinājumu saskatīja Vācijas impērijas radīšanā, kur ar pārtiku to nodrošinātu iekarotās zemes. Polijas pozīcija bijusi kaut kur pa vidu. Tā vēlējusies atņemt daļu zemju sekmīgi strādājošajām [muižnieku] saimniecībām un atdot to zemniekiem. T.~Snaiders skaidro: pēc Pirmā pasaules kara pati dzīve prasījusi lai valsts būtu arī zemnieku valsts, tas nozīmē~--- bija jāpārliecina zemnieki, ka dodot kaut ko valstij, identificējoties ar nāciju, kļūstot lojāliem nacionālajai valstij, viņi arī kaut ko saņems par to. Bet vienīgais ievērojamais žests, kas spējis viņus par to pārliecināt, bija agrārā reforma, muižnieku zemju sadale un tās nodošana zemniekiem. Nacionālās pašapziņas pieņemšana pašā elementārākajā līmenī nozīmēja zemes iegūšanu apmaiņā par lojalitāti valstij. Taču šo apmaiņu tā arī neizdevās izdarīt un tieši šeit Polijas valsts piedzīvoja neveiksmi. Agrārā reforma neizdevās pamatā politisku, bet daļēji arī etnisku cēloņu darbības rezultātā, tā kā kresos valsts nevēlējās piešķirt zemi nepoļu izcelsmes personām. Bet bez pamatīgas agrārās reformas nebija izredžu panākt, lai sabiedrības vairākums nopietni izturētos pret valsti. Tas tā arī neizturējās. Pēc Lielās depresijas Žečpospolitai vairs nebija izredžu īstenot lauksaimniecības reformu. Pieauga politiskā opozīcija. Toties nacionālisms bija bezauglīgs: poļu nacionālisti gandrīz visi apgalvoja, ka zemniecība ir nācijas sirds, taču nebija spējīgi neko darīt tās labā. Protams, amerikāņu vēsturnieka skatījums ir visai shematisks, taču tuvina objektīvā stāvokļa izpratnei.

Sekas Polijā pastāvošajai agrārajai iekārtai bija zemnieku radikalizēšanās, to kustības aktivizēšanās. Kā atzīmējis angļu vēsturnieks N.~Deiviss, zemnieku sabiedrība arvien vairāk polarizējās, jo padziļinājās konflikti starp konservatīvo zemes īpašnieku interesēm un radikālajām zemnieku partijām; zemniecībai uzgūlās jauns valdības birokrātijas un ārvalstu kompāniju prasību slogs, to plosīja šķiru un etniskie protesti. Komunisti, izanalizējuši agrārās attiecības, zemes sadalījumu dažādu īpašnieku kategoriju vidū, nonāca pie secinājuma, ka pastāvošo agrāro attiecību saglabāšanā objektīvi bija ieinteresēta tikai ļoti neliela grupa lauku iedzīvotāju. Secinājums atbilda patiesībai, bet valdošie slāņi to nevēlējās redzēt.

Kaut jaunā Polijas valsts saspringtu militāru un diplomātisku cīņu rezultātā ieguva vairākus rūpnieciski attīstītus apgabalus, tie paši par sevi vēl nepavēra sekmīgas rūpnieciskās attīstības iespējas. \strong{Rūpniecība}, neraugoties uz dažādu izejvielu atradņu lielo skaitu, salīdzinājumā ar Rietumeiropas valstīm Polijā bija attīstīta krietni vien vājāk. Ievērojama loma tautsaimniecībā bija akmeņogļu ieguvei, metalurģiskajai, tekstilrūpniecībai un pārtikas rūpniecībai. Šis nozares koncentrējās vairākos rajonos: smagā un kurināmā rūpniecība Dombrovas baseinā (\pltxti{Zagłębie Dąbrowskie}) un Augšsilēzijā, tekstilrūpniecība~--- Lodzas rajonā, mašīnbūvniecība~--- Varšavā, kara rūpniecība~--- Kelces-Radomas apgabalā. Citas pilsētas bija otršķirīgi rūpniecības centri.

Visā starpkaru periodā Polijas rūpniecība nesasniedza līmeni, kurš attiecīgajās poļu apdzīvotajās teritorijās bija pirms kara. Galvenie cēloņi tam bija iekšējā tirgus šaurība, ekonomisko sakaru vājināšanās ar bijušajām Polijas sadalītājvalstīm, bet īpaši ar PSRS, valsts mērogā vienota tautsaimniecības mehānisma izveides grūtības, poļu preču nespēja konkurēt pasaules tirgū ar citu, attīstītāku valstu ražojumiem.

1921.~gadā strādnieku skaits sasniedza ap 700~tūkstošu, pieaugot augstas konjunktūras gados. 70\% visu rūpniecības strādnieku bija koncentrēti trijos apgabalos: Dombrovas, Silēzijas un Lodzas, kuri sastādīja tikai ap 20\% no valsts teritorijas. Kā jau iepriekšējā nodaļā rādīts, dažādu apgabalu ekonomiskās attīstības līmenis bija dažāds.

Lai minam tikai pāris piemēru. Valsts austrumu rajonos, izņemot spirta un degvīna, miltu un stērķeles ražošanas pusmājamatnieciskos uzņēmumus, gandrīz nebija citu ražotņu. Pilnībā izpalika modernā rūpniecība.

No 50~dzelzceļa līnijām, kuras no agrākās Vācijas Pozenes provinces un agrākās Austroungārijas Galīcijas provinces veda uz bijušo Polijas karalisti, tikai 10~turpinājās tās teritorijā. Pēc neatkarības atjaunošanas visi šie apgabali bija vienā valstī.

1926.~gada sākumā Poliju PSRS delegācijas sastāvā (Starp Poliju un PSRS, neraugoties uz vēsajām attiecībām, notika apmaiņa ar oficiālām arodbiedrību, literātu, saimniecības darbinieku, arī politiķu delegācijām, pēc vienošanās abu valstu teritorijas pārlidoja otras valsts lidotāji, darbojās vienošanās par apmaiņu ar politieslodzītajiem) apmeklēja ievērojamais padomju saimnieciskais darbinieks Jūlijs Kārlis Daniševskis. Savā ziņojumā viņš rakstīja: ``Kopējais iespaids par Polijas rūpniecību ir pretrunīgs, no vienas puses, bez šaubām, Katovices un Sosnovices (\pltxti{Katowice} un \pltxti{Sosnowiec}~--- pilsētas Polijas dienvidos, Augšsilēzijā, akmeņogļu ieguves un melnās metalurģijas centri) rajonos ir lieliski iekārtotas mašīnu, lokomotīvju un vagonu būvniecības, ķīmiskie un citi uzņēmumi, kā arī lieliskas ogļu šahtas, kurās mēs arī nolaidāmies. To pašu var teikt par Lodzas tekstilrūpnīcām, kuras ir iekārtotas pēc tehnikas pēdējā vārda, pat ar jaunākajiem darbgaldiem, izgatavotiem 1925.~gadā. Taču, no otras puses, var redzēt, ka noieta šīm precēm nav. Noliktavas piepildītas, un ne tikai tekstiluzņēmumos, bet arī metalurģiskajos. Pat nelielā rūpniecības centrā~--- Ļubļinas~--- mēs varējām vērot to pašu parādību lauksaimniecības mašīnu un svaru rūpnīcās. Tekstilrūpnīcas strādā tikai daļēji, vienā, labākajā gadījumā divās maiņās, un arī tad 3--4 dienas nedēļā. Vienā maiņā strādā arī Sosnovices ogļu raktuvēs. Tai pašā laikā, neraugoties uz formāli pastāvošo astoņu stundu darba dienu, sarunās noskaidrojās, ka ar to vai citu ieganstu atsevišķas grupas, un diezgan plašas, tekstilrūpniecībā strādā 10 un pat 12~stundas.''

Rūpniecības attīstībai piemita samērā ierobežots raksturs. Ja akmeņogļu, cukura, elektroenerģijas ražošana dažus gadus pēc kara jau pārsniedza pirmskara līmeni, tad čuguna un tērauda kausēšana, naftas ieguve, kokvilnas audumu izlaide visu starpkaru periodu tā arī nesasniedza 1913.~gada apjomu. Tā, piemēram, lielākajā tekstilrūpniecības pilsētā Lodzā ražošana laikā no 1918.~līdz 1939.~gadam samazinājās par 75~procentiem tādēļ, ka tradicionālais noieta tirgus~--- Krievija~--- vairs nebija pieejams. 1924.~gadā rūpniecības kopapjoms sasniedza tikai 56\% no 1913.~gada produkcijas.

Lieli, samērā moderni rūpnieciskie uzņēmumi Polijā sadzīvoja ar daudziem sīkiem, tehniski atpalikušiem. Tai laikā, kad Polijas iedzīvotāji sastādīja 1,6\% no visiem zemeslodes iedzīvotājiem, Polijas rūpnieciskās produkcijas apjoms nepārsniedza 0,5--0,7\% no pasaulē saražotā apjoma.

Polijas rūpniecībā jau tradicionāli visai augsts bija ārzemju kapitāla ieguldījumu apjoms. Poļu vēsturnieks L.~Grosfelds raksta, ka 1929.~gadā visu Polijas akciju sabiedrību kapitālā ārzemju kapitāla daļa sastādīja 33,3\%, bet 1933.~gadā~--- jau 44,2\%. Bez tam Polijā darbojās arī ārvalstu uzņēmumu filiāles. Pēc minētā tautsaimniecības vēsturnieka aprēķiniem ārzemju kapitāla kontrolē esot bijuši ap 68\% visa Polijas akciju kapitāla.

Augšsilēzijā valdošais bija vācu kapitāls (Izņēmums bija amerikāņu-vācu apvienību ``\pltxti{Wspólnota Interesów Górniczo-Hutniczych}'' (Kalnrūpniecības un metalurģijas kopīgās intereses), kuru kontrolēja A.~Harimana un F.~Flika finanšu grupas.) Naftas rūpniecībā un Dombrovas baseina ogļrūpniecībā valdošais bija franču kapitāls. Metalurģijā un naftas rūpniecībā ievērojamas bija amerikāņu kapitāla pozīcijas. Ķīmiskajā un tekstilrūpniecībā darbojās britu kapitāls. Ārzemju kapitālam bija liela loma arī Polijas banku sistēmā. Jāatzīmē gan, ka ārvalstu kapitālisti gandrīz neradīja jaunus uzņēmumus, bet uzpirka jau esošos vai to akcijas. Pēc 1930.~gada statistikas datiem ārzemju ieguldījumi jaundibinātu kompāniju akcijās sastādīja 11~miljonus zlotu, bet jau agrāk pastāvējušu kompāniju akcijās~--- 94~miljoni zlotu. Daļu peļņas ārzemju investori izveda, daļu izmantoja citu uzņēmumu akciju uzpirkšanai, konkurences novēršanai. Ārzemju uzņēmēji un baņķieri bija arī Polijas uzņēmumu kreditori. 1928.~gadā oficiāli ārzemju koncerniem izmaksājamās un izvedamās no valsts dividendes tika vērtētas kā 53~081~000~zlotu, bet procentu maksājumi par piešķirtajiem kredītiem~--- 110~321~000~zlotu. 9~gados (1924--1932) Polijas maksājumi ārzemju kapitālistiem sastādīja 3,5~miljardus zlotu.

Faktiski ārzemju kapitāla pieplūdums veda nevis pie kapitālieguldījumu pieauguma, bet pie līdzekļu izvešanas dividenšu un procentu veidā pieauguma. Tā, pēc oficiāliem datiem, Polijā tika ieguldīti 164,8~miljoni zlotu, bet izvesti aiz robežas tika 326,3~miljoni zlotu. Gadījās, ka ārvalstis tieši iejaucās Polijas iekšējās lietās. Tā, bijušais Polijas premjerministrs K.~Svitaļskis savā dienasgrāmatā fiksējis interesantu faktu. 1929.~gada 26.~augustā pie viņa kā pie premjerministra ieradās ASV sūtnis un izsacīja pretenzijas sakarā ar streiku metalurģiskajā rūpnīcā, kura ietilpa viņa tautieša A.~Harimana vadītajā koncernā.

Pēc oficiāliem datiem vislielākais strādnieku skaits Polijā bija 1928.~gadā~--- gandrīz miljons (980~000) cilvēku. (Pēc citiem datiem 1928.gadā kalnrūpniecībā nodarbināto strādnieku bija 1~422~000, metalurģiskajā rūpniecībā~--- 616~000, metālapstrādē~--- 806~000, pārtikas rūpniecībā~--- 637~000, kokapstrādē~--- 411~000, ķīmiskajā rūpniecībā~--- 313~000.) Taču no tiem ne mazāk kā 100~000 (citi dati~--- 200~000) pastāvīgi bija bez darba. Viņu dzīves līmenis bija viens no viszemākajiem brīvā tirgus Eiropā. Piemēram, tekstilizstrādājumu patēriņš poļu iedzīvotāju vidū bija piecas reizes, sviesta patēriņš~--- trīs ar pusi reizes zemāks nekā Vācijā.

Polijā bija samērā attīstīta strādnieku likumdošana: pastāvēja kolektīvie darba līgumi, apdrošināšana pret slimībām, negadījumiem un bezdarbu, darbojās darba inspekcija. Taču likumi ne vienmēr tika pildīti pilnā mērā~--- organizācijas, līdzekļu trūkuma rezultātā. Sieviešu darbs faktiski tika apmaksāts par 30\% zemāk nekā vīriešu. Pēc gada darba tika piešķirts 8~dienu atvaļinājums. Pēc triju gadu nostrādāšanas tas tika pagarināts līdz 15~dienām. Garīgā darba strādniekiem ik gadu pienācās mēnesis atvaļinājuma. Bezdarbnieka pabalsts tika maksāts tiem, kuri pirms darba zaudēšanas gada laikā bija nostrādājuši ne mazāk nekā 20~nedēļas, 30--50\% apmērā no vidējās darba algas, taču ne ilgāk par 17 nedēļām. Apdrošināšana pret bezdarbu neattiecās uz laukstrādniekiem, jaunākiem par 18 gadiem, mājkalpotājiem, kalpotājiem, darbu zaudējušajiem slimības, darba nespējas gadījumā vai streika rezultātā.

Ar strādnieku sociālo aizsardzību nodarbojās arodbiedrības. Taču tās bija ļoti sadrumstalotas un ar dažādu politisku nokrāsu. Daļa arodbiedrību atradās PPS ietekmē. Taču, ja 1922.~gadā tajā bija ap 459~tūkstošu biedru, līdz 1926.~gadam to skaits samazinājās līdz 200~tūkstošiem. Tikai 1927.~gadā sakarā ar zināmu ekonomiskās aktivitātes un attiecīgi arī strādnieku kustības pieaugumu atkal bija vērojams biedru pieaugums. Blakus minētajām šķiriski orientētajām arodbiedrībām pastāvēja arī nacionālas poļu un ebreju arodbiedrības, kuras noliedza jebkādu politisku darbību, atzina tikai ekonomisko cīņu, arī kristīgas, baznīcas ietekmē esošas arodbiedrības, kuras apvienoja galvenokārt nerūpnieciskos strādniekus. 1924.~gadā visās arodbiedrībās kopā sastāvēja ap 1~300~000 biedru, no tiem ap 400~000 bija fabriku strādnieki.

Daudz smagu problēmu Polijai bija jārisina \strong{finanšu jomā}. Liela daļa budžeta joprojām tika tērēta militārajiem izdevumiem. Polija ne tikai baidījās no abiem kaimiņiem~--- Vācijas un Padomju Savienības, bet vēlējās arī Francijai pierādīt, ka ir vērtīgs, militāri stiprs tās sabiedrotais.

1923.~gadā jau V.~Sikorska valdības laikā, pirms jaunās valdības ar ļudovieti V.~Vitosu priekšgalā nākšanas pie varas, strauji \strong{pieauga inflācija}. Ja 1918.~gadā 1 ASV \$ maksāja 9,8~poļu markas, 1919.~gadā par to jau bija jāsamaksā 90~markas. Turpmāk inflācija turpināja augt. (Skat. \ref{tab:table6}.~tabulu) 1923.~gadā jau tika izdotas naudas zīmes 100~000, 250~000, 500~000 un 1 miljona marku vērtībā, bet gada beigās pat 5 un 10~ miljonu marku vērtībā.

\noindent
\begin{table}[h!]
\caption{ASV dolāra vērtība poļu markās attiecīgajā gadā} \label{tab:table6}
\begin{tabularx}{\linewidth}{|p{4cm}|p{4cm}|}
\hline
\strong{Gads} & \strong{1 ASV dolāra vērtība poļu markās} \\
\hline
1919.~g. & 90 \\
\hline
1921.~g. & 6~000 \\
\hline
1923.~g. maijs & 52~000 \\
\hline
1923.~g. jūlijs & 140~000 \\
\\\hline
1923.~g. novembra sākums. & 2~000~000 \\
\\\hline
1923.~g. novembra beigas & 5~000~000 \\
\hline
1924.~g. janvāris & 9~300~000 \\
\hline
\end{tabularx}
\end{table}

% page 566


Ja 1918.--1923.~gada periodā valsts no tiešajiem un netiešajiem nodokļiem bija saņēmusi \$214~miljonus~, no citiem budžeta ienākumiem~--- \$236~miljonus, tad papīra naudas emisija deva tai \$868~miljonus, tātad naudas izlaide divreiz pārsniedza visus budžeta ieņēmumus. Ievērojamais ekonomists finanšu ministrs V.~Grabskis ierosināja stabilizācijas pasākumus, kuriem bija jāskar faktiski visi iedzīvotāji, bet valdība atteicās pieņemt viņa programmu, kura labējiem nepatika, jo tika skartas to elektorāta intereses. V.~Grabskis tāpēc 1923.~gada jūnijā iesniedza atlūgumu. Inflācija pirmkārt ``sita'' pa darba ņēmēju kabatu. Pēc kāda poļu ekonomista aprēķiniem, uz strādnieku, kuri sastādīja ap 12\% Polijas iedzīvotāju, pleciem gūlās 24\% to zaudējumu, ko 1923.~gadā nesa inflācija. Mantīgie nemaksāja nodokļus, par eksportpreču pārdošanu ieņemto valūtu atstāja ārzemēs. Ienākumi no tiešajiem nodokļiem 1923./1924.~gadu mijā sastādīja tikai 8\% no valsts izdevumiem. Valdība kapitālu trūkumu kompensēja ar naudas emisiju, paaugstināja netiešos nodokļus (uz sērkociņiem, cukuru, degvīnu), ieviesa arī jaunu zemes un īpašuma nodokli. Notika masveida atlaišanas. 1923.~gada oktobrī sākās streiki. 29.~oktobrī, cerot salauzt strādnieku pretošanos, V.~Vitosa valdība nolēma ieviest uz dzelzceļa ārkārtējo stāvokli un lauku kara tiesas. Strādnieku spiediena rezultātā PPS un arodbiedrību vadība bija spiesta piekrist vispārēja streika izsludināšanai valstī. 5.~novembrī šis streiks sākās.

Vairākās pilsētās izvērsās strādnieku sadursmes ar policiju. Sevišķi asas tās bija Krakovā. Policija atklāja uguni uz streikojošo strādnieku demonstrāciju. Divi cilvēki tika nogalināti, taču demonstranti neapstājās, uzbruka policijai. Aiz policijas stāvēja armijas divas rotas, taču karavīri, kuru vidū, kā vēlāk noskaidrojās, bija daudz jauniešu no ukraiņu un baltkrievu ciemiem, nemetās policijai palīgā, bet bez pretestības atdeva savus ieročus. Tos ieguvuši, strādnieki spēja atvairīt ulānu uzbrukumu, pat sagrāba vienu no trijām pret tiem sūtītajām bruņumašīnām, vairākus ložmetējus, kopā ap 500~šauteņu. Streikojošo sadursmes ar armijas daļām notika arī Borislavā (\pltxti{Borysław}) un Tarnovā (\pltxti{Tarnów}). Novembra streika rezultātā ielu kaujās gāja bojā ap 40~cilvēku, vairāki simti tika ievainoti. Tiesa, visi politiskie spēki, izņemot komunistus, baidījās no pilsoņu kara un meklēja kompromisa lēmumu. Tika atcelts kā vispārējais streiks, tā valdības ievestais ārkārtējais stāvoklis un lauku kara tiesas.

Kad 1923.~gada 19.~decembrī tika izveidota ``ārpus partiju'' valdība ar formāli bezpartijisko V.~Grabski priekšgalā, tā par galveno uzdevumu stādīja tautsaimniecības atveseļošanu. V.~Grabskis bija pirmais Polijas premjerministrs, kurš centās īstenot politiku visas sabiedrības interesēs. Tādēļ reformu slogu, kas tika uztverts kā pirmais solis finanšu sistēmas stabilizēšanā, viņš paredzēja uzkraut vienlīdzīgi visām sabiedrības grupām. Viņa izstrādātā programma finanšu sistēmas reformai tika paredzēta uz trim gadiem. Daži nepieciešamie likumi Seimā bija pieņemti jau līdz tam, citu pieņemšanu V.~Grabskis bija iecerējis panākt veikt parlamentā pēc iespējas īsākā laikā. Ar uzticības balsojumu Seimā 1924.~gada 5.~janvārī viņa valdību atbalstīja labējie un centrs, PPS un kreisā zemnieku partijas atturējās, pret balsoja bija vienīgi komunisti un mazākumtautību pārstāvji.

Gandrīz divus gadus V.~Grabskim izdevās lavierēt starp labējiem un kreisajiem. Kad V.~Grabska valdība no Seima saņēma pilnvaras stabilizēt valūtu, tā pacēla nodokļus turīgākajiem iedzīvotāju slāņiem. Mantīgie slāņi, neraugoties uz to, ka šie nodokļi uzlika tiem krietni lielākas finansiālās saistības, tos tomēr atzina, jo inflācija, pāraugot hiperinflācijā (inflācijā ar augstu vai ārkārtīgi augstu tempu), bija arī tiem neizdevīga. Jau pēc dažiem mēnešiem~--- 1924.~gada 28.~aprīlī V.~Grabskim izdevās nomainīt poļu marku pret zlotu (attiecībā 1,8~miljonu marku: 1 zlotu.) Zlota vērtība tika pielīdzināta Šveices frankam, kas to padarīja par vienu no dārgākajām valūtām Eiropā. To jau varēja brīvi konvertēt. 1924.~gada budžetā naudas emisijas loma bija minimāla. Taču uz laiku maksājumus pārtrauca vairākas bankas, palielinājās nestrādājošo uzņēmumu skaits. Ar finanšu reformu bija saistīts tarifu pieaugums dzelzceļa pārvadājumiem. Tā rezultātā kļuva neiespējams kokmateriālu eksports. Piemēram, to pārvešana no Volīnijas uz Dancigu maksāja 60\% no tā cenas Dancigā. Vēl vairākus mēnešus pieauga oficiāli reģistrēto bezdarbnieku skaits (1924.~gada septembrī tas sasniedza 165~000), kritās darba alga.

1924.~gads izvērtās smags Polijas tautsaimniecībai. Lauksaimniecībā valdīja neraža, jo daudzi apgabali cieta plūdos. Polija graudu eksporta vietā bija spiesta tos importēt. Būtiski pazeminājās cenas Polijas eksportprecēm: kokmateriāliem, cukuram, oglēm. Tas viss negatīvi ietekmēja darba cilvēku materiālo stāvokli.

Valūtas stabilizācija nebalstījās uz visas saimnieciskās sistēmas atveseļošanos. Finanšu stāvokļa nokārtošana atkailināja citus tautsaimniecības trūkumus: rūpniecības tehnisko atpalicību un zemo ražīgumu, zaudējot inflācijas prēmijas, mazinājās tās konkurences spēja ārējos tirgos. Turpretī iekšējais patēriņš pieauga maz, tāpēc inflācijas beigas izsauca krīzi rūpniecībā, tirdzniecībā, transportā, amatniecībā. Kritās ražošana. 1924.~gadā tā sastādīja 84,8\%, bet 1925.~gadā~--- 86,5\% no 1923.~gada līmeņa. Tas izsauca nodarbinātības samazināšanos. 1924.~gadā lauksaimniecības produktu cenas sāka atpalikt no pieaugošajām rūpniecības preču cenām. Šai gadā ievērojamu valsts daļu skāra arī neraža. Kviešus 1924.~gadā novāca 61,6\%, bet rudzus~--- 78,6\% apmērā no 1923.~gada ražas. Lai mazinātu pārtikas krīzi pilsētās, valdība sāka veicināt labības un miltu ievešanu no ārzemēm. 1925.~gada janvārī~--- aprīlī labības ievedums sasniedza 10\% no visa importa vērtības. Neizdevās likvidēt valsts budžeta deficītu. No 1924.~gada augusta līdz 1925.~gada augustam ārējās valsts tirdzniecības deficīts sasniedza 646,7 miljonus zlotu.

Otrais inflācijas vilnis bija mazāk jūtams, taču krīzes pazīmes saglabājās. No 1925.~gada janvāra līdz jūlijam banknošu zelta saturs kritās no 65 uz 49\%. Valsts bankā samazinājās ārzemju valūtas daudzums. Ja oficiāli 1925.~gada augustā \$1 maksāja 5,18~zlotus, neoficiālais maiņas kurss bija 7,2~ zloti pret \$1. Polijas bankai nācās paaugstināt oficiālo kursu līdz 5,73~zlotiem. Decembra beigās \$1 maksāja jau 8,45~zlotus. Samazinājās ražošana. Piemēram, ja 1924.~gada pirmajā pusgadā ieguva 16~432~000~t ogļu, 1925.~gadā tai pašā laikā gada ieguve samazinājās uz 14~950~000 t. Pieauga bezdarbnieku skaits.

1925.~gada 14.~novembrī demisionēja V.~Grabska valdība. Viņa vietā par premjerministru kļuva A.~Skšinskis. Taču ko būtiski mainīt viņš nespēja. 1925.~gada decembrī oficiāli bija reģistrēti 311~000 bezdarbnieku. 1925.~gada 17.~novembrī visu fabrikantu, banku, tirdzniecības, amatu un zemkopības apvienību pārstāvji iesniedza prezidentam S.~Vojcehovskim memorandu, kurā nosauca valsts ekonomisko stāvokli par katastrofālu un raksturoja to sekojoši: ``Ekonomiskā dzīve Republikā pamirst ar katru dienu, katru stundu. Bezdarbs rūpniecībā nemitīgi aug, bezdarbnieku skaits jau ir pārsniedzis 213~tūkstošu cilvēku, kas sastāda 30~procentu no visiem rūpniecībā strādājošajiem. Līdzīgas situācijas nav nekur Rietumeiropā. Tie strādnieki, kuriem ir darbs, bieži ir nodarbināti tikai 5, 4, 3 dienas nedēļā. Rezultātā uz 6~strādniekiem, nodarbinātiem ražošanā, iznāk 4~bezdarbnieki \citespace{} šī attiecība ar katru dienu palielinās. Drīz izrādīsies, ka mazākums baros vairākumu. Ar tādu uzdevumu vairs neviens netiks galā.'' Finanšu ministrs izdeva rīkojumu samazināt algas kalpotājiem, pensijas invalīdiem, dzelzceļnieku štatus un ieviest nodokļus elektrības un petrolejas apgaismojuma lietotājiem.

Tika taupīts arī citādi. Tā, Polijā izbraukšana uz ārzemēm bija saistīta ar prāviem izdevumiem~--- ārzemju pase maksāja daudz vairāk nekā, piemēram, Latvijā. Ar to tika panākts, ka daudzi tūkstoši vasarnieku, kas citādi aizplūstu uz ārzemju kūrortiem, aizvezdami līdzi arī prāvas naudas summas, palika Polijā.

1925.~gada beigas un 1926.~gada sākums iezīmēja ekonomiskās krīzes virsotni. Valsti satricināja demonstrāciju un mītiņu vilnis.

Tomēr ar 1923.~gada finanšu katastrofu un to pavadošo revolucionāro krīzi beidzās ar Pirmā pasaules kara seku pārvarēšanu un valsts tapšanu saistītais nestabilitātes periods. 1924.~gadā ar naudas reformu un konsekventu saimniecisko politiku tika apturēta straujā inflācija un atjaunots ekonomiskais līdzsvars, kas kļuva par pamatu ekonomikas augšupejai 1926.--1929.~gadā. (1926.~gadā zlota kurss ar stabilizācijas aizņēmuma palīdzību tika noteikts: 8,91~zlots = \$1. 1~kg tīra zelta maksāja 5~924,24~zlotus.)

Īslaicīgā ekonomiskā stabilizācijas perioda laikā tika meklēti risinājumi pastāvošajām problēmām.

Viena no tādām bija \strong{parādu nomaksa}. Parādus sastādīja:

\begin{enumerate}
\item Poļu nacionālās komitejas Parīzē izdevumi J.~Hallera armijas finansēšanai. Tie bija \$25,8~miljoni un sastādīja ap 5,8~\% no visa parāda.

\item Versaļas un Sen-Žermenas līgumi paredzēja, ka Polijai jānomaksā sava daļa no Vācijas un Austroungārijas reparācijām, atbilstoši šo valstu teritoriju, kuras pēc kara iegāja Polijas sastāvā, ekonomiskajai vērtībai. Taču vēlāk šī summa daudzkārt mainījās un gala rezultātā kā tāda vairs nefigurēja.

\item Vislielākie parādi Polijai radās 1919.~un 1920.~gadā sakarā ar poļu-padomju karu. Aizņēmumos tika saņemti \$246,1~miljons jeb 55,9\% no visa parāda.

\item 1921.--1926.~gadā parāds pieauga vēl par \$91,8~miljoniem (21,8\%).
\end{enumerate}

1924.~gada 14.~novembrī Vašingtonā tika parakstīts līgums par Polijas parāda, kura lielāko daļu sastādīja samaksa par ASV piegādēm 1919.--1921.~gadā, nolīdzināšanu. Uz 1933.~gada 1.~janvāri šis parāds sastādīja 3~289,6 miljonus zlotu. Tas bija jānomaksā līdz 1984.~gadam. Bez tam Polija saņēma vēl virkni aizdevumu. No 1918.~līdz 1927.~gadam bija realizēti arī 22~iekšējie aizņēmumi, to kopējā summa sasniedza \$33~miljonus. 1927.~gadā sekoja jauns aizņēmums no ASV, kurš paredzēja, ka valdībai tā atmaksas nodrošināšanai jāpaaugstina nodokļi, Polijas bankas darbu un valsts budžeta izpildi (ik gada ceturksni) kontrolēja amerikāņu finansistu pārstāvis. Ekonomists un politiķis S.~Grabskis 1929.~gadā presē rakstīja: ``Starptautiskā ārvalstu kapitāla kontrole pār mūsu ekonomiku liedza valdībai iespēju sekmīgi cīnīties pret tirdzniecības bilances deficītu. Ārzemju kapitāls ir ieinteresēts, lai Polija būtu amerikāņu, angļu, vācu preču patērētāja, bet nevis amerikāņu, angļu, vācu uzņēmumu konkurente starptautiskajos tirgos. \citespace{} Mēs jau esam zaudējuši neatkarību kredītu politikas jomā''.

Kopumā Polijas parādi salīdzinājumā ar citām Eiropas valstīm nepavisam nebija lieli, taču to nomaksu traucēja tautsaimniecības attīstības grūtības. Ļoti lielu Polijas valsts budžeta daļu sastādīja izdevumi militārām vajadzībām. 1922.~gadā militārie izdevumi bija 39\%, 1923.~gadā~--- 38,7\%, 1924.~gadā~--- 42\%. Tai pat laikā tādās kaimiņvalstīs kā Rumānija, Somija, Latvija, Igaunija, nerunājot jau par Rietumu lielvalstīm, šie izdevumi bija krietni mazāki. Tāpēc padomju autors I.~Troickis, kurš Sarkanās armijas vajadzībām sagatavoja Polijas militāri ekonomisku un politisku apskatu, secināja, ka ``Polija ir pati militāristiskākā valsts pasaulē'' un to pēc militārajiem izdevumiem var salīdzināt vienīgi ar Japānu. Šo vērtējumu pēc viena kritērija~--- militāro izdevumu īpatsvara valsts budžetā gan nevar atzīt par korektu.

\strong{Problēmas} pastāvēja Polijas \strong{ārējā tirdzniecībā}.

Pēc pirmā acu uzmetuma varēja likties, ka Polijai, kuras atsevišķas sastāvdaļas agrāk bija vismaz daļēji integrētas triju kaimiņvalstu ekonomikā, jāturpina uzturēt ar šīm valstīm ciešas ekonomiskās saites, jāizmanto iepriekšējo sakaru pastāvēšana savas ekonomikas attīstībai. Taču tā tas nenotika.

Austroungārija bija sadalīta vairākās savstarpēji konkurējošās valstīs, kurām sakari ar jauno Polijas valsti bija vajadzīgāki ne vairāk kā ar citām kaimiņvalstīm.

Agrākās Krievijas vietā bija izveidojusies PSRS ar citādu ideoloģiju, citādu~--- plānveida saimniecību. Tas apgrūtināja agrāko tirgus attiecību atjaunošanu. Padomju valdība vēl 1921.~gadā piedāvāja Polijai noslēgt tirdzniecības līgumu, taču 1922.~gada martā sākušās sarunas beidzās neveiksmīgi, jo Polijas delegācija pieprasīja padomju valstī pastāvošā ārējās tirdzniecības monopola atcelšanu. Padomju puse tam nevarēja piekrist.

Ievērojama daļa poļu uzņēmēju, īpaši pirms kara saistītu ar Krievijas tirgu, labprāt tajā atgrieztos, neraugoties uz minētājām atšķirībām, taču tam par šķērsli bija Polijā valdošo politiķu vairākuma uzskati, kuri miera noslēgšanu ar Padomju Krieviju vērtēja ne kā jaunu attiecību veidošanas posma sākumu, bet gan tikai kā pagaidu pamieru. Polijas armijas Ģenerālštābs vērtēja padomju-poļu tirdznieciskās attiecības kā ``padomju propagandas'' avotu un tāpēc uzstājās par to ierobežošanu līdz minimumam. Tikai reta valdība, kā A.~Poņikovska vadītā (1922.~gadā), centās normalizēt ekonomiskās attiecības ar austrumu kaimiņu. Tikai pierobežas rajonos pēc kara beigām radās stihisks varas iestāžu nekontrolēts tirgus. Turpretī oficiāli no Polijas importēto preču apjoms visā padomju importā 1921.~gadā sastādīja mazāk nekā 0,1\%. Avīze ``\pltxti{Kurier Polski}'' 1921.~gada decembrī rakstīja: ``Ja kaut ceturto daļu no tā, ko mēs tērējam armijai, mēs uzdrošinātos ieguldīt mūsu dalības palielināšanai Austrumeiropas saimnieciskās dzīves atdzīvināšanā, tas nodrošinātu mūsu nākotni daudz lielākā mērā kā cietokšņi un lielgabali''. Polijas un PSRS valdību delegācijas vairākkārt ilgstoši veda sarunas par starpvalstu tirdzniecības līguma noslēgšanu, bet līdz 1939.~gadam tās visas bija neveiksmīgas.

Ekonomisko attiecību attīstība gan norisa neatkarīgi no sarunu norises, taču vispārējās sarežģītās abu valstu politiskās attiecības ietekmēja arī tirdzniecību. 1931.~gadā Polijas eksports uz PSRS sastādīja 6,7\% no visa Polijas eksporta, bet imports no PSRS~--- 2,5\% no visa Polijas importa. 1932.~gadā tirdzniecības apgrozījums samazinājās. Pēc abu valstu neuzbrukšanas līguma noslēgšanas 1932.~gada 25.~jūlijā apmēram gadu to ekonomiskā sadarbība uzlabojās, bet pēc Polijas un Vācijas tuvināšanās 1934.~gadā atkal pasliktinājās. 1934.~gadā Polijas eksports salīdzinājumā ar 1931.~gadu samazinājās par 80\%, bet imports~--- par 50\%. Polijas un PSRS tirdzniecības līgums tika parakstīts tikai 1939.~gada februārī, pusgadu līdz Otrā pasaules kara sākumam un praktiski tam vairs nebija nozīmes.

Arī ar Vāciju Polija bija konfliktējusi par robežām, politiskās pretrunas arī vēlāk traucēja saimniecisko attiecību attīstību.

Jau 1922.~gadā Polijā tika pieņemts lēmums izbūvēt jaunu ostu nelielā ciematā Gdiņā uz ziemeļiem no Dancigas brīvpilsētas. Nākamajā~--- 1923.~gadā V.~Sikorska valdība noslēdza līgumu ar Franciju par Gdiņas ostas un zemūdeņu bāzes tajā izbūvi. 1924.--1930.gg. tur galvenokārt par valsts līdzekļiem izvērsās plaša celtniecība. Tika uzsākta tieša dzelzceļa līnijas~--- t.s. ``ogļu maģistrāles'' būvniecība no Augšsilēzijas uz Gdiņu (pabeigta 1933.~gadā). Lieli nopelni Gdiņas attīstības koncepcijas izstrādē un īstenošanā bija rūpniecības un tirdzniecības ministram, vēlāk premjera vietniekam un finanšu ministram F.~Kviatkovskim. Ja 20.~gadu sakumā Gdiņa bija neliels zvejniekciems, 1935.~gadā tajā mita jau 70~000 iedzīvotāju, tās ostā varēja redzēt kuģus ar visdažādāko valstu karogiem. 40~regulāras kuģu satiksmes līnijas saistīja Gdiņu ar svarīgākajām ostām pasaulē.

Vācu laikabiedri gan uzskatīja, ka Gdiņas ostas celtniecību nosacīja nevis saimnieciskā nepieciešamība, bet ārpolitiskas ambīcijas. Vislas upes tirdzniecības rajonu gadsimtu gaitā sekmīgi bija apkalpojusi Dancigas (Gdaņskas) osta. Pie tam pirms Pirmā pasaules kara Vislas upi izmantoja triju valstu tirgotāji, tagad tā visa atradās Polijas rīcībā. Versaļas konferencē poļu prasība pēc Dancigas tika pamatota ar Vislas kuģniecības nodrošināšanu, un vismaz daļēji pamatojoties uz šo prasību arī tika radīta Dancigas brīvpilsēta. Taču 30.~gadu beigās preču kravas, kuras izgāja caur 1896.~gadā ierīkotajām Vislas slūžām 18~km no Dancigas, salīdzinājumā ar pirmskara periodu, bija būtiski samazinājušās (sasniedza tikai ap 75\% no agrākā apjoma). Tādejādi izrādījās, ka Versaļas konferencē Dancigas iegūšanai Polijas īpašumā izmantotais saimnieciskais pamatojums bija izrādījies nepareizs. Jau 1933.~gadā preču apgrozība caur Gdiņu pārsniedza to Dancigā. Gdiņa kalpoja Polijai ne tikai eksportam, bet arī ārzemju preču ievešanai. Ja 1931.~gadā caur Gdiņu tika importētas preces par 106~miljoniem zlotu, tad 1932.~gadā~--- jau par 133~miljoniem, bet 1933.~gadā~--- par 297~miljoniem. Tai pašā laikā importa vērtība caur Dancigu samazinājās no 263~līdz 124~miljoniem zlotu.

Gdiņas uzplaukums gan saasināja Polijas attiecības ar Vāciju, kura cīnījās par Dancigas atgūšanu. Taču vēstures gaita pierādīja, ka perspektīvā poļu pūliņi attaisnojās. Arī pēc tam, kad pēc Otrā pasaules kara Danciga (Gdaņska) nonāca poļu rokās, Gdiņa nezaudēja savu nozīmi. Mūsdienās trīs centri (``\pltxti{Trojmiasto}''): Gdaņska, Sopota un Gdiņa ir viena no Polijas aglomerācijām, kuras attīstās visstraujāk.

Jāņem vērā arī, ka poļu un vācu attiecības Dancigā bieži apgrūtināja abu valstu politikas realizāciju. 1925.~gada janvārī Dancigā izvērsās vācu un poļu iedzīvotāju sadursmes. (Par ieganstu tām kalpoja Polijas pasta kastīšu izkāršana pilsētas ielās, jo brīvpilsēta tika iekļauta Polijas pasta iestāžu darbības sfērā. Vācu nacionālisti šīs pastkastītes pārkrāsoja ķeizariskās Vācijas krāsās, ar to arī izprovocējot nemierus.)

Situācija arī turpmāk nebija labvēlīga ilgstošai attīstībai, ekonomiskās grūtības neveicināja valdības stabilitāti. Īpaši tās kļuva jūtamas, kad no 1925.~gada jūnija \strong{ar Vāciju} izvērsās atklāts \strong{``muitas karš''}. 1925.~gadā Polija ieviesa jaunus muitas tarifus, kuri ierobežoja Vācijas, Austrijas un Čehoslovākijas rūpniecības izstrādājumu importu. Vācija atbildēja ar simetriskiem pretpasākumiem, atbildei ierobežoja Polijas lauksaimniecības ražojumu importu. Pēc Vācijas statistikas datiem, kuri gan neņēma vērā tirdzniecību caur Dancigu, Polijas ievedums Vācijā 1923.~gadā sastādīja 357~miljonus vācu marku, bet Vācijas ievedums Polijā~--- 400~miljonus marku. 1933.~gadā šie skaitļi bija attiecīgi 55,9 un 55,5~miljoni marku. Tā kā Polija ieņēma otršķirīgu vietu Vācijas ārējā tirdzniecībā (tās eksports uz Poliju sastādīja tikai 6\% visa Vācijas eksporta, toties Vācija atradās pirmajā vietā Polijas eksportā un importā~--- ievedums no Vācijas sastādīja 32,4\% no visa Polijas importa, bet izvedums uz to~--- 42,4\% visa Polijas eksporta), ``muitas kara'' smagumu nācās iznest Polijas tautsaimniecībai. Valstī radās lieli ogļu uzkrājumi. Šajā laikā beidzās vienošanās par Augšsilēzijas ogļu un rūpniecības ražojumu piegādi Vācijai. Vācijas valdība atteicās garantēt līdzšinējo 6~miljonu tonnu poļu ogļu eksportu uz Vāciju, kas sastādīja 40\% no visa Polijas eksporta vērtības. (1923.~gadā uz Vāciju izveda 8~miljonus tonnu ogļu, bet 1924.~gadā~--- tikai 6,7~miljonus tonnu.) Tiesa, angļu ogļraču streiks 1926.~gadā ļāva izvērst poļu ogļu eksportu uz Skandināviju. Tomēr Vācija kavēja poļu preču virzību uz savu tirgu, bet tas bija poļu eksporta galvenais virziens. Polijas zaudējumi sasniedza līdz 1/4 no tās eksporta. Samazinājās cukura izvešana, kurš ieņēma otro vietu Polijas eksportā. Tā kā Polija bija spiesta samazināt nepieciešamo iekārtu un izejvielu importu, lai iegūtu valūtu, tā par lētu cenu eksportēja visu, ko vien varēja pārdot ārzemju tirgos.

Polija savukārt veica vairāku Vācijas preču ieveduma ierobežojumus. Par to bija sašutuši vācu nacionālisti: Polijas kā agrāras valsts tirdzniecības partnerim esot jābūt Viduseiropai, bet tās tirdzniecības ceļiem uz Balkāniem un Itāliju jāiet caur Vāciju, nevis pa Vācijas teritoriju uz Baltijas jūru. Taču Polija cenšoties vest galveno tirdzniecības apjomu pa jūru. Polijas muitas pasākumi esot vienīgi cīņas pasākumi pret Vāciju. Zlota stabilitāte balstīja V.~Grabska popularitāti, bet ``muitas karš'' veicināja zlota devalvāciju. Tas noveda pie tā, ka 1925.~gada novembrī viņa valdība krita.

Poļu~--- vācu muitas karš ilga vairākus gadus. Kā atzīmēja poļu vēsturnieks J.~Krasuskis, 1928.gadā notika galīga situācijas, kura pastāvēja no konflikta sākuma, noskaidrošanās. Šai laikā neviena no pusēm nebija ieinteresēta sasniegt vienošanos. Vācijā to nevēlējās agrārieši un ogļu ražotāji, Polijā~--- rūpnieciskās aprindas. Sarunas par tirdzniecību bija tīrā šķietamība, kaut arī neviena puse nevēlējās uzņemties atbildību par atklātu to neveiksmi.

Tikai 20.~gadu beigās Polijai beidzot pavērās iespēja uzlabot attiecības ar Vāciju. 1929.~gada pavasarī tika noslēgts Polijas un Vācijas līgums par atteikšanos no savstarpējām pretenzijām par kara zaudējumu atlīdzību. Polijas valdība cerēja, ka tas atvērs ceļu tirdzniecības līguma noslēgšanai, jo līdz šim Vācijas valdība galvenokārt savas Austrumelbas rajona muižniecības interesēs atteicās atvērt Vācijas tirgu visu Polijas lauksaimniecības ražojumu ievešanai. Šīs un vairāku citu starpvalstu vienošanos noslēgšana 1928.~un 1929. gadā gan vājināja agrāko ekonomiskā kara saspringumu, taču neatrisināja galvenās Polijas un Vācijas ekonomiskās pretrunas.

1930.~gada 17.~martā pēc ilgstošām sarunām beidzot tika parakstīts Vācijas un Polijas tirdzniecības līgums. Likās, ka beidzot atrasta izeja no nemitīgā vācu-poļu attiecību sasprindzinājuma, no piecus gadus ilgušā savstarpējā ``muitas kara''. Abas puses savas ekonomiskās attiecības turpmāk apņēmās būvēt uz vislielākās labvēlības principa pamata. Pēc līguma Polija varēja vairāk eksportēt lauksaimniecības, īpaši lopkopības, produkciju. Reizē gan līgums nostādīja grūtā situācijā vairākas Polijas rūpniecības nozares, kuras ``muitas kara'' laikā bija brīvas no Vācijas rūpniecības konkurences. Taču vienošanās tā arī nekad nestājās spēkā. Drīzumā pēc tā parakstīšanas Vācijas agrāriešu protekcionistiskā kursa piekritēju spiediena rezultātā valstī tika paaugstināti muitas tarifi, kas Polijai radīja vēl lielākas grūtības lauksaimniecības produkcijas eksportam nekā iepriekš. Pēc tam Reihstāgs atteicās ratificēt tirdzniecības līgumu ar Poliju. Tomēr ``muitas karš'' bija apturēts, ``uzskrūvēto'' muitas nodevu vietā stājās līdzsvaroti muitas tarifi. Jāsaka, ka ``muitas karš'' arī veicināja Polijas tautsaimniecības konkurences spējas nostiprināšanos, tā pārstāja būt atkarīga no Vācijas tirgus, taču šis rezultāts parādījās tikai pēc dažiem gadiem.

1926.--1929.~gadā Polijas tautsaimniecība, pateicoties starptautiskajai konjunktūrai, attīstījās samērā sekmīgi. 1926.~gada februārī krīzes virsotne bija jau aiz muguras un kopš 1926.~gada vidus Polijas tautsaimniecība sāka atkopties, taču tas nebija ``sanācijas'' režīma nopelns. Pēc pēckara stabilizācijas krīzes pārvarēšanas 1926.~gads bija pirmais normālas saimnieciskās attīstības gads. Kaut Polijas attīstības tempi vēl joprojām atpalika no attīstītajām valstīm (ja 1926.~gadā pasaules [bez PSRS] rūpnieciskās ražošanas indekss par 17\% pārsniedza 1913.~gada līmeni, tad Polijā tas bija zemāks nekā 1913.~gadā par 40\%). Sākotnēji K.~Bartela valdība apspieda inflāciju ar valdības izdevumu taupību, apcirpa valsts dzelzceļa izdevumus, pavairoja eksportu.

Pieauga poļu ogļu eksports, jo parādījās jauni to tirgi. Lielbritānijas ogļraču streika (1926.~gada maijs-septembris) rezultātā Polija varēja kāpināt savu akmeņogļu eksportu. Britu ogļraču streiks deva Polijas oglēm noieta tirgu Baltijā un Skandināvijā, taču galvenā Polijas ogļu noņēmēja izrādījās pati Lielbritānija. Ja 1924.~gadā tā Polijā iepirka 9~000 t ogļu, tad 1926.~gadā~--- jau 2,7~miljonus~t. Ja kopumā 1926.~gada pirmajā pusē Polija eksportēja 3,5~miljonus t ogļu, tad gada otrajā pusē~--- jau 10~miljonus~t. Strauji pieauga transporta pārvadājumi, jaunus pasūtījumus saņēma mašīnbūvniecības u.c. nozaru uzņēmumi. Ekonomisko rosību veicināja ārzemju kapitālu ieguldījumi Polijas rūpniecībā, galvenokārt iegūstošajā.

Sakarā ar ogļu eksporta pieaugumu aktīva kļuva Polijas ārējās tirdzniecības bilance. 1926.~gadā izvedums par 707~miljoniem zlotu pārsniedza ievedumu. No 313~tūkstošiem 1926.~gada janvārī līdz 442~tūkstošiem 1926.~gada novembrī pieauga rūpniecībā nodarbināto strādnieku skaits. 1926.~gada jūlijā Poznaņā tika izgatavota Polijā pirmā lokomotīve no poļu materiāliem. 1927.~gada janvārī darbu uzsāka pirmā Polijā automobiļu rūpnīca. (Tiesa, vēl 1935.~gada 1.~jūlijā Polijā bija tikai 25~212~kravas auto. Valsts ar 10,5~kravas automobiļiem uz katriem 10~000~iedzīvotājiem tālu atpalika no citām lielākajām Eiropas valstīm. Piemēram, Vācijā bija 59~kravas automobiļi uz katriem 10~000 iedzīvotājiem.) Arī 1926.~gada apmierinošā raža veicināja stabilizāciju sabiedrībā. Tas viss palīdzēja cīņā pret ekonomiskajām grūtībām.

1926.--1929.~gadu periods starpkaru laikā bija vienīgais, kurš Polijai it kā solīja tālāku ilgstošu uzplaukumu. No ekonomiskā viedokļa tas bija pats gaišākais posms valsts vēsturē. Bija stabilizējies arī starptautiskais stāvoklis, tāpēc auga investīcijas un cēlās zlota kurss. Mazinājās ``muitas kara'' ar Vāciju nelabvēlīgās sekas. Polijā ieradās ASV ekonomiskais eksperts E.~Kemmerers, kā rezultātā 1927.~gada oktobrī Polija saņēma finansiālu, tā saucamo ``stabilizācijas'' aizdevumu (\$62~miljoni un £2~miljoni), no starptautiska banku konsorcija, kurā galveno lomu spēlēja ASV un Lielbritānijas bankas. Taču aizņēmuma noteikumi Polijai bija mazāk labvēlīgi nekā citām tai laikā aizņēmumus saņēmušajām ekonomiski stabilākajām valstīm~--- Vācijai, Austrijai u.c. Polijas tautsaimniecības nesakārtotības pazīmes palika. Īpaši sāpīgi to izjuta nemantīgie.

Ekonomisku grūtību apstākļos, kad Polijā trūka kapitāla investīcijām, uzņēmēji centās apvienot spēkus. Pastāvīgi pieauga dažādu Polijas uzņēmumu noslēgto vienošanos skaits par tirgu sadali, par cenu, par ražošanas apjomu, importu un eksportu, peļņas sadali utt. Piemēram, 1934.~gadā 175~karteļu vienošanās, kuras darbojās valsts iekšienē, un 50~starptautiskās tieši noteica cenas, 147~--- maksājumu nosacījumus, 113~--- noieta kvotas, 138~--- cenas un maksājumu nosacījumus, 97~--- cenas, maksājumu nosacījumus un noieta kvotas. 58~poļu un 14~starptautiskie karteļi regulēja peļņas sadali, 87~starptautiskas vienošanās regulēja eksportu, 75~--- importu (47~gadījumos starptautiskās vienošanās paredzēja Polijas eksporta ierobežošanu, 4~vienošanās to vispār aizliedza. (10~gadījumos poļu uzņēmēji saņēma kompensācijas par eksporta ierobežošanu, bet 80~firmas uz valsts iekšienē noslēgto vienošanos pamata saņēma kompensāciju par ražošanas ierobežošanu.) Šīs dažādās apvienības vienkāršības dēļ turpmāk visas sauktas par karteļiem.

Karteļu lomu uzskatāmi parāda piemērs. Līdz cementa karteļa izveidei 1926.~gada 100~kg cementa maksāja 2,5~zlotus, tai pat gadā pēc karteļa izveides cēlās līdz 6,5~zlotiem, bet 1927.~gadā sasniedza 7,7~zlotus un tādā līmenī turējās līdz 1933.~gadam. Kad valdība bija spiesta 1933.--1934.~gadā likvidēt cementa karteli, dažos mēnešos cenas uz cementu kritās četras reizes.

Pēc 1937.~gadā izdotā oficiālā krājuma ``\pltxti{Mafy Rocznik Statystczny}'' (``Statistikas gadagrāmata'') datiem karteļu veidošanās aina bija sekojoša: Skat. \ref{tab:table7}.~tabulu.

\noindent
\begin{table}[h!]
\caption{Karteļu skaits Polijā 1919.--1936.~gadā} \label{tab:table7}
\begin{tabularx}{\linewidth}{|p{4cm}|p{4cm}|p{4cm}|}
\hline
\strong{Gads} & \strong{Nacionālo karteļu skaits} & \strong{Starptautisko karteļu skaits, kuros piedalījās Polija} \\
\hline
1919 & 9 & 3 \\
\hline
1920 & 11 & 4 \\
\hline
1921 & 13 & 4 \\
\hline
1922 & 15 & 4 \\
\hline
1923 & 23 & 6 \\
\hline
1924 & 31 & 6 \\
\hline
1925 & 40 & 6 \\
\hline
1926 & 53 & 8 \\
\hline
1927 & 64 & 14 \\
\hline
1928 & 77 & 23 \\
\hline
1929 & 100 & 40 \\
\hline
1930 & 133 & 48 \\
\hline
1931 & 168 & 69 \\
\hline
1932 & 187 & 75 \\
\hline
1933 & 215 & 89 \\
\hline
1934 & 233 & 94 \\
\hline
1935 & 268 & 104 \\
\hline
1936 & 274 & 105 \\
\hline
\end{tabularx}
\end{table}

% page 574


Viens no vadošajiem Polijas starpkaru perioda ekonomistiem, bijušais ministrs un vairāku lielu uzņēmumu vadītājs J.~Glivics 1930.~gadā rakstīja, ka ``nebija gandrīz nevienas [rūpniecības] nozares, kurā nenotiktu koncentrācijas un sindikātu, karteļu, konvenciju un tamlīdzīgu apvienību izveides process''. Valsts struktūras atbalstīja monopolizācijas procesu ekonomikā.

Tā, apgabala tiesa 1930.~gadā noraidīja kādas tapešu fabrikas, kura neietilpa kartelī, sūdzību pret tapešu ražojošu uzņēmumu karteli par ``darbību, kas virzīta pret iedibinātajām tradīcijām''. Tiesas lēmumā bija teikts: ``Kartelizācija ir tendence, kura atbilst mūsdienu saimnieciskās dzīves attīstības tendencei. \citespace{} Karteļu pastāvēšanu Polijā pārsvarā izsauc nepieciešamība, nevis savtība'' \citespace{} ``ja jau sindikātu, trestu un karteļu pastāvēšanu pieļauj valsts, tad ir arī jāpieļauj firmu, apvienotu karteļos, noteikti cīņas līdzekļi pret firmām, kuras neietilpst šajās apvienībās''. 1927.~gadā Polijas valdības radītas komisijas atskaitē bija norādīts, ka karteļi necenšas veikt pasākumus ražošana racionalizācijai, tai pat laikā neļauj sekmīgi darboties labi tehniski apgādātiem, bet finansiāli vājiem uzņēmumiem; apgrūtina produkcijas realizāciju ar lieliem komisijas izdevumiem; ierobežo ražošanu, lai ``uzskrūvētu'' cenas uz produkciju; nenodarbojas ar darba apstākļu uzlabošanu. Taču neraugoties uz to visu, minētā komisija iestājās par valsts atbalstu kartelizācijai. Ievērojams tā laika poļu ekonomists S.~Dzevuļskis 1933.~gadā kartelizācijai veltītā grāmatā rakstīja: ``Polijā valdība bez šaubām īsteno politiku, kura atbilst karteļu interesēm.'' To pierādīja arī augstu cenu pastāvēšanu uz karteļos apvienoto uzņēmumu produkciju arī krīzes gados, kas veicināja patēriņa krišanu. Tā, 1931.~gada beigās, kad lauksaimniecības preču cenu indekss sastādīja 53,3 (par 100 pieņemtas 1928.~gada cenas), kopējais rūpniecības cenu indekss~--- ap 88, karteļos ietilpstošo uzņēmumu preču cenu indekss bija 107. Monopoli visiem spēkiem aizstāvēja savus ienākumus.

Polijas banku un uzņēmumu uzraudzības padomju raksturīga īpatnība bija lielais to sastāvā ietverto t.s. ``veco'', ``vēsturisko'', aristokrātisko dzimtu pārstāvju skaits. Tā, kņazs S.~Ļubomirskis bija 13 (tai skaitā 2~banku, 3~kalnrūpniecības kombinātu, 4~mašīnbūves un cementa ražošanas, apdrošināšanas u.c. kompāniju) uzraudzības padomju loceklis. Septiņi citi kņazi Ļubomirski ietilpa vēl 15~dažādu uzraudzības padomju sastāvā. Grāfs A.~Potockis bija 11~uzraudzības padomju loceklis, bet vēl 12 citi grāfi Potocki piedalījās 26 uzraudzības padomju darbībā. Minēto padomju, kā arī uzņēmumu valžu un direkciju locekļi bija arī daudzi augstākie valsts ierēdņi.

Milzīgu ietekmi ieguva Polijas rūpniecības, kalnrūpniecības, tirdzniecības un finanšu centrālā asociācija ``Leviafāns'' (\pltxti{Centralny Związek Polskiego Przemysłu, Górnictwa, Handlu i Finansów ``Lewiatan''}), no 1932.~gada Polijas Rūpniecības centrālā asociācija \pltxti{(Centralny Związek Przemysłu Polskiego}). Asociācija lobēja nodokļu un sociālo iemaksu pazemināšanu ražotņu īpašniekiem, valsts atbalsta palielināšanu, darba ņēmēju tiesiskās aizsardzības samazināšanu. Asociācijas biedri, lieluzņēmēji starpkaru periodā bija Seima un Senāta deputāti, iekļauti valdību sastāvā, īpaši ``sanācijas'' gados ieņēma augstus valsts amatus. Tā, viens no ``Leviafāna'' dibinātājiem, tā valdes priekšsēdētājs, 16~lielu Polijas un starptautisku banku un rūpniecības kompāniju uzraudzības padomju loceklis un priekšsēdētājs A.~Vežbickis reizē bija Seima deputāts (1919--1927 un 1935--1937), kādu laiku bija rūpniecības un tirdzniecības ministrs (1918), valdības delegācijas loceklis Parīzes miera sarunās, no 1926.~gada Ministru padomes Ekonomiskās komitejas Konsultatīvās komitejas priekšsēdētājs. ``Leviafāna'' Centrālās padomes loceklis, 12~lielu kompāniju uzraudzības padomju loceklis un priekšsēdētājs Č.~Klarners 1924.--1925.~gadā pildīja finanšu ministra, 1925.~gadā~--- rūpniecības un tirdzniecības ministra pienākumus, daudzkārt bija senatora amatā. Šo uzskaitījumu var turpināt ilgi.

Uzņēmēji ietekmi valsts iestādēs izmantoja, lai gūtu maksimālu peļņu. Piemēram, 1948.~gadā Londonā iznākušajā bijušā Polijas aviācijas ģenerāļa L.~Raiska atmiņu grāmatā ``\pltxti{Słowa prawdy o lotnictwie polskim, 1919--1939}'' (Patiesības vārdi par Polijas Gaisa spēkiem 1919--1939.) cita starpā bija pastāstīts, kā firma ``\pltxti{Plage i Laśkiewicz}'' ieguva valsts pasūtījumu uz 300~lidmašīnu izgatavošanu triju gadu laikā. Avansu veidā firma saņēma gandrīz visu līgumā paredzēto summu. Taču izrādījās, ka šai laikā tika izgatavotas tikai 70~sliktas kvalitātes lidmašīnas, kuras bieži avarēja un lidotāji tās sauca par ``lidojošajiem zārkiem''. Kad firma bankrotēja, Gaisa spēku pavēlniecība noīrēja tās rūpnīcu par nesamērīgu augstu maksu. L.~Raiskis min arī citus piemērus, kā uzņēmēji iedzīvojās uz valsts rēķina. Poļu profesors F.~Cveigs 1944.~gadā Londonā izdotā darbā rakstīja, ka Polija palika ekonomiskā ziņā atpalikusi ``daļēji ārzemju kapitāla spiediena, daļēji pilnīgās valsts interešu atstāšanas novārtā no valdības puses rezultātā''.

Stāvokli Polijā būtiski ietekmēja \strong{pasaules ekonomiskā krīze}, kura sākās 1929.~gadā. Sabruka mīts par it kā ``sanācijas'' radīto ekonomisko brīnumu. Polijai krīze atnesa salīdzinoši lielākas grūtības nekā citām valstīm. Ogļrūpniecība 1929.~gadā bija noslogota par 74\%, tekstilrūpniecība~--- tikai par 50\%. Vairākums strādnieku~--- tekstilnieku uz darbu devās tikai divas līdz četras dienas nedēļā. 1929.~gadā visas pasaules (bez PSRS) rūpnieciskās ražošanas indekss pret pirmskara 1913.~gadu sastādīja 150, tad Polijā tas veidoja tikai 99,3. [Pēc citiem datiem 1929.~gadā pasaulē (bez PSRS) rūpnieciskās ražošanas līmenis bija par 46\% augstāks kā pirms kara, bet Polijā šai gadā rūpnieciskā ražošana sastādīja 86,56\% no pirmskara līmeņa.] Ja 1933.~gadā pasaulē šis indekss bija 79,8 (Polijas kaimiņos~--- Vācijā~--- 60,7, Čehoslovākijā~--- 60,2), tad Polijā tas samazinājās līdz 55,6. Tas bija zemāks nekā kaimiņvalstīs. Ja ogļu~--- galvenās eksportpreces~--- izvedums 1929.~gadā sastādīja 5,6~miljonus t, tad 1931.~gadā tas samazinājās līdz 2,3~miljoniem~t. Skandināvijā poļu ogļu vietā atkal atgriezās Anglijā iegūtās. Eksporta tirgus zaudēja poļu cukurs. Bankrotēja daudzas firmas un uzņēmumi. 1929.--1932.~gadā uzņēmumu skaits Polijā samazinājās par 23,7\%. Ekonomika smagi cieta. Piemēram, 1932.~gadā ražošanas jaudas ogļrūpniecībā bija noslogotas par 60\%, čuguna ražošanā~--- par 85\%, mašīnbūvniecībā~--- par 90\%, cukurrūpniecībā par 58\%.

1932.~gadā Polijā rūpniecības produkcijas tika saražots par 1/3 mazāk kā 1929.~gadā, tā sastādīja tikai 53\% no pirmskara apjoma. Salīdzinājumā ar 1929.~gadu 1933.~gadā rūpnieciskās ražošanas indekss Polijā bija 55. Poļu ekonomists B.~Civinskis 1932.~gadā konstatēja: ``Mūsu rūpniecības finansiālo stāvokli 1932.~gadā var raksturot kā progresējoša bankrota stāvokli; ar krīzes sākumu Polijas rūpniecības aprīkojuma attīstība pa augšupejošu līniju pārtrūkst. Sākas kustība atpakaļ~--- saražotā patēriņš, jaunu kapitālieguldījumu vairs nav.'' Par rūpniecības tehnisko stāvokli 1937.~gadā viņš rakstīja: ``Neatjaunināmie, novecojušie rūpniecības uzņēmumi ar katru gadu zaudē savu ražošanas vērtību, kļūst par vecu grabažu muzejiem, bet nevis dižas nācijas ekonomisko bruņojumu.''

Jau 1928.~gadā Polijā bija rudzu, kartupeļu neraža. Turpmāk kritās zemnieku saimniecību ienesība. Krīze rūpniecībā un lauksaimniecībā bija cieši saistīta ar maksājumu krīzi. Tirdzniecības bilances deficīts 1929.~gadā sastādīja 854 miljoni zlotu.

Specifiska finanšu krīzes iezīme Polijā bija poļu valūtas~--- zlota zelta paritātes fiktīva saglabāšana, kad daudzas kapitālistiskās valstis veica savu valūtu devalvāciju. Tai laikā, kad zelta krājumi pasaulē pieauga, Polijas bankas zelta un valūtas rezerves sāka kust. Ja 1929.~gadā Polijā bija zelts 476~miljonu zlotu vērtībā, tad 1933.~gadā tā krājumi bija 471~miljonu zlotu vērtībā. Vēl straujāk kritās ārvalstu valūtas krājumi. Ja 1928.~gadā Polijas bankā bija ārzemju valūta 714~miljonu zlotu vērtībā, tad 1933.~gadā tās krājumi sastādīja vairs tikai 88~miljonus zlotu. Rezultātā zelta un valūtas nodrošinājums zlotam no 83,5\% 1929.~gadā kritās līdz 58,2\% 1932.~gadā. Zlota devalvācija netika veikta tāpēc, ka varēja rasties tāds stāvoklis, kad valdība nevarētu noteikt devalvācijas robežu un kurss kristu arvien zemāk, kā tas notika 1923.~gadā. Arī iedzīvotāju vairākums vēl no pēckara inflācijas laika krasi negatīvi izturējās pret devalvāciju un inflāciju, kas vestu pie tā, ka plaši iedzīvotāju slāņi zaudētu savus kaut nelielos uzkrājumus. ``Sanācijas'' režīmam ar tā jau šauro sociālo bāzi tas draudētu ar lieliem nemieriem.

Pēc poļu ekonomistes un statistiķes T.~Maleckas datiem, 1929.~gadā Polijas ārējais parāds sastādīja 34\% nacionālā ienākuma. Galvenais kreditors bija ASV, kurai 1930.~gadā piederēja 38,2\% no kopējā Polijas parāda.

Jau 1930./1931.~gada budžets tika izpildīts ar 63 miljonu zlotu deficītu. 1931./1932.~gadā budžeta deficīts sasniedza 205~miljonu zlotu apjomu. Nākamā gada budžeta deficīts bija 242~miljoni zlotu. Četru gadu laikā valsts kases ienākumi saruka par 40\%. Valsts izdevumi, neraugoties uz visiem taupības pasākumiem, samazinājās tikai par 27\%.

Finanšu aktivitāti raksturoja arī privāto banku skaita samazināšanās krīzes periodā. Ja 1928.~gadā Polijā darbojās 52~privātās bankas, tad 1933.~gadā to bija palicis tikai 39.

Par pasaules ekonomisko krīzi nebija atbildīga ``sanācijas'' valdība. Jau minētais Č.~Klarners pat aicināja lepoties ar to, ka Polija nav veicinājusi pasaules ekonomiskās krīzes izraisīšanos, jo tās daļa pasaules tirdzniecības pagrozījumā sastādīja tikai 0,8\% (Citi dati~--- 0,7\% līdz krīzei, un 0,5\% pēc tās) un nevienai poļu precei nebija izšķirošas nozīmes pasaules tirgos. Taču krīzes laikā Polijas vadība parādīja bezspēcību, modernas ekonomikas likumu neizpratni. Valdība baidījās no inflācijas un ierobežoja valsts izdevumus. Tā arī bažījās, ka ar kapitālu aizplūšanas uz ārzemēm, ar zlota konvertēšanas ārzemju valūtā nepieļaušanu tiks grauta ticība Polijas stabilitātei. Tāpēc tie valsts ekonomikas regulēšanas līdzekļi, protekcionisms un intervencionisms, kurus plaši lietoja citas tirgus ekonomikas valstis, Polijā tika pieņemti pārāk vēlu un bez vajadzīgās noteiktības. Polijā nebija kā līdzīga ASV realizētajam F.~D.~Rūzvelta) ``Jaunajam kursam'' (\entxti{New Deal}), kad blakus citiem pretkrīzes pasākumiem norisa cīņa pret ``negodīgu konkurenci'', un ``neveselīgām attiecībām'' starp biznesu un uzņēmumu darbiniekiem. Pārprodukcijas krīzes apstākļos no diviem veidiem kā iegūt peļņu~--- uz ražošanas racionalizācijas pamata samazināt cenu uz preces vienību vai samazināt ražošanu un saglabāt augstas cenas~--- uzņēmēji-monopolisti izvēlējās otro. Tāda politika nebūtu iespējama pastāvot brīvai konkurencei, taču bija iespējama monopolu kundzības apstākļos.

Oficiāli tika uzskatīts, ka Polija realizē deflācijas politiku. Taču tā bija tikai daļa patiesības. Tā bija deflācijas politika pret strādājošajiem un izpaudās viņu darba algu samazināšanā, darba dienas pagarināšanā, viņu agrāko sociālo iekarojumu atņemšanā. Nevar saukt par deflācijas politikas realizēšanu augstu eksporta prēmiju, subsīdiju izsniegšanu, nodokļu parādu norakstīšanu uzņēmējiem, galvojumu piešķiršanu kredītu saņemšanai maksātnespējīgām kompānijām, augstu karteļu cenu atzīšanu. Lai veicinātu eksportu tika lietotas dempinga (mākslīgi pazeminātas) cenas. Ja 1932.~gadā Polijas iekšienē cukura mazumtirdzniecības cena bija 160~zlotu par 100~kg, tad eksportcena~--- 16~zlotu. 1933.~gadā valsts iekšienē 1~t ogļu (bez pievešanas maksas) bija 25~zloti, bet eksportcena~--- 10~zlotu. Faktiski starpību no savas kabatas sedza poļu patērētāji. Tāpēc nebija brīnums, ka cukura (un arī dažu citu preču) patēriņš Polijā bija pats zemākais Eiropā (ja 1929.~gadā viens Polijas iedzīvotājs patērēja 11,9~kg cukura, tad 1933.~gadā~--- 8,6~kg) un vairākkārt atpalika no patēriņa tādās valstīs kā Lielbritānija, Dānija, Čehoslovākija u.c.

Uzņēmēji centās krīzes smagumu pārlikt uz darba ņēmēju pleciem. Piemēram, lielu uzņēmumu vadītājs E.~Lansbergs uzstājās par apdrošināšanas bezdarba gadījumiem likvidēšanu, jo, pēc viņa domām, ``nodrošināšana bezdarba gadījumā ir pamatcēlonis, kas izsauc bezdarbu''. Šo kursu atbalstīja arī valdība. 1932.~gada decembrī premjerministrs A.~Pristors, uzstājoties Seimā, paziņoja: ``Atbilstoši saimniecisko sfēru prasībai, mums ir jāmazina izdevumi sociālajiem pasākumiem''. Tas tika teikts tad, kad poļu strādnieku dzīves līmenis un alga bija viena no zemākajām Eiropā. Jau līdz 1928.~gadam svarīgākajās rūpniecības nozarēs poļu strādnieku reālās darba algas indekss nesasniedza 1914.~gada līmeni. Pēc Starptautiskā darba biroja datiem piecu gadu laikā no 1924.~līdz 1928.~gadam Polijā, neraugoties uz svārstībām, darba alga samazinājās par 2\%. (Blakus Vācijā tai pat laikā darba alga pieauga par 55\%. Arī otrā Polijas kaimiņvalstī Rumānijā darba alga pieauga par 30\%). Pie tam no 1925.~gada vidus līdz 1926.~gada vidum Polijā darba algas strādniekiem sistemātiski tika aizturētas. Uzņēmēju parādi strādniekiem sastādīja miljonus zlotu. Tautā runāja, ka strādnieki uzstājas kā ``fabrikantu baņķieri''. Skat. \ref{tab:table8}. tabulu.

\noindent
\begin{table}[h!]
\caption{Polijas lielo un vidējo rūpniecības uzņēmumu strādnieku darba alga 1929--1935.~gadā} \label{tab:table8}
\begin{tabularx}{\linewidth}{|p{4cm}|p{4cm}|}
\hline
\strong{(1928.~gada līmenis pieņemts par 100) Gads} & \strong{Darba alga} \\
\hline
1929 & 109 \\
\hline
1930 & 89 \\
\hline
1931 & 70 \\
\hline
1932 & 51 \\
\hline
1933 & 46 \\
\hline
1934 & 49 \\
\hline
1935 & 51 \\
\hline
\end{tabularx}
\end{table}

Salīdzinājumu ar citām valstīm rāda strādnieku nožēlojamo stāvokli Polijā. Skat. \ref{tab:table9}.~tabulu.

\noindent
\begin{table}[h!]
\caption{Darba algas salīdzinājums Polijā un citās valstīs 1929.~gadā (darba alga Lielbritānijā pieņemta par 100)} \label{tab:table9}
\begin{tabularx}{\linewidth}{|p{4cm}|p{4cm}|}
\hline
\strong{Valsts} & \strong{Darba alga (salīdzinājumā ar Lielbritāniju \%)} \\
\hline
ASV & 259 \\
\hline
Zviedrija & 122 \\
\hline
Vācija & 82 \\
\hline
Čehoslovākija & 65 \\
\hline
Polija\footnote{Darba alga Polijā šeit pieņemta 14,5 zlotu dienā, kas pēc L.~Grosfelda vērtējuma krasi atšķīrās no faktiskās, daudz zemākās izpeļņas.} & 51 \\
\hline
\end{tabularx}
\end{table}

% page 577


Kā redzam, 1929.~gadā poļu strādnieka alga bija divas reizes zemāka nekā vidējā angļu strādnieka alga un par 40\% zemāka nekā vācu strādnieka alga. Krīzes laikā stāvoklis vēl vairāk pasliktinājās. Faktiski katrs trešais strādnieks zaudēja darbu, bet pārējiem samazināja algu.

Samazinājās arī garīgā darba darītāju apmaksa. Viņu mediālais (par mediālo šeit saukts tāds atalgojuma līmenis, par kuru augstāku saņēma puse no visiem garīgā darba darītājiem un kuru nesasniedza otra puse) atalgojums samazinājās no 260~zlotiem mēnesī 1928.~gadā līdz 189~zlotiem 1935.~gadā (par 30\%).

Krīzes apstākļos 1933.~gada 22.~martā tika pieņemts likums par izmaiņām 1919.~gada 18.~decembra likumā par darba laiku rūpniecībā un tirdzniecībā (``\pltxti{Ustawa z dnia 22~marca 1933~r. w sprawie zmian w ustawie z dnia 18~grudnia 1919~r. o czasie pracy w przemyśle i handlu.}''). Izmaiņas atcēla 46~stundu darba nedēļu. Tika mainīti arī citi likumi. Piemēram, strādniekiem vairs neapmaksāja svētdienas un svētku dienas, kuras iegadījās viņu atvaļinājuma laikā. Rūpniecības nozarēs ar sezonas darba raksturu tika ieviesta 10~stundu darba diena.

1932.~gadā tika pieņemts likums, saskaņā ar kuru sabiedriskās vajadzības nodrošinošo uzņēmumu strādniekus par streikošanu varēja sodīt ar cietumsodu līdz 5~gadiem. 1933.~gada janvārī stājās spēkā jauns Civilkodekss, kurš atzina tikai individuālos līgumus starp uzņēmējiem un strādniekiem, ar to dodot triecienu arodkustībai. Tomēr strādnieku kustība turpinājās. Piemēram, 1933.~gada martā norisa plašs Lodzas u.~c. pilsētu tekstilnieku streiks. Policijas sadursmēs ar streikojošajiem pieci strādnieki tika nogalināti.

1933.~gada 28.~marta likums atņēma tiesības saņemt slimības pabalstus visiem strādniekiem lauksaimniecībā, arī dažām citām strādnieku kategorijām, piemēram, praktikantiem. Citiem strādniekiem slimības pabalstu izmaksas termiņš samazinājās no 39~līdz 26~nedēļām. Slimības pabalstu lielums samazinājās no 60\% līdz 50\% no darba algas. Bērnu dzimšanas pabalsti samazinājās vairāk nekā divas reizes. Ja agrāk darba devējs apmaksāja 3/5 no apdrošināšanas pret slimību iemaksas, bet pats darba ņēmējs 2/5, tagad gan darba devējs, gan ņēmējs maksāja 1/2. Šādu taupības uz strādājošo rēķina piemēru uzskaitījumu var turpināt.

Valsts izdevumi sociālajai nodrošināšanai saruka no 155~miljoniem zlotu 1931./1932.~gadā līdz 65~miljoniem zlotu 1934./1935.~gadā. Tai pat laikā pensijas invalīdiem un algas valsts kalpotājiem samazinājās no 164~līdz 103~miljoniem zlotu. Premjerministrs F.~Slavojs Skladkovskis savā grāmatā ``\pltxti{Strzępy meldunków}'' (``Ziņojumu druskas'', 1936), ir pastāstījis kā tika sastādīts budžets krīzes laikā. Izrādās, ka tas notika uz īsu J.~Pilsudska pavēļu pamata. Diktators paziņoja, ka ``pirmā nelaime ir ierēdņi, pārāk uzpūstie štati un augstās algas''. Tāpēc divu gadu laikā viņu likmes tika samazinātas par 30\%. Var piebilst, ka valsts aparātā, kā arī valsts un municipālajos uzņēmumos bija nodarbināti simti tūkstoši strādnieku un kalpotāju (1/3~garīgā darba darītāju un 21\% strādnieku). Lieluzņēmējus šāds stāvoklis apmierināja. Viens no asociācijas ``Leviafāns'' pārstāvjiem R.~Bataļja 1933.~gadā ciniski rakstīja: ``Vai tad jautājumu par sociālās apgādes paplašināšanu un attīstību \citespace{} par pieprasījuma pēc grāmatām pieaugumu, par inteliģences sadzīves apstākļu uzlabošanu u.tml. var salīdzināt ar jautājumu par smagās rūpniecības pastāvēšanas nepieciešamību.''. Faktiski runa gāja nevis par smagās rūpniecības pastāvēšanu, bet tās nesto peļņu.

Šīs politikas rezultātā katastrofāli samazinājās plašu masu patēriņš. Patēriņa jomā uz vienu iedzīvotāju Polija ieņēma vienu no pēdējām vietām Eiropā, taču tā ieņēma pirmo vietu patēriņa krišanā krīzes laikā.

Jau minētais bijušais ministrs Č.~Klarners 1935.~gadā publicētā grāmatā ``\pltxti{U podstaw bezrobocia}'' (``Pie bezdarba saknēm'') ir salīdzinājis patēriņu krīzes laikā Polijā un citās valstīs. Skat. \ref{tab:table10}. tabulu.

\noindent
\begin{table}[h!]
\caption{Preču vidējais patēriņš (kg) gadā uz vienu iedzīvotāju Polijā, Lielbritānijā un Vācijā 1927.--1931.~gadā} \label{tab:table10}
\begin{tabularx}{\linewidth}{|p{3cm}|p{3cm}|p{3cm}|p{3cm}|}
\hline
\strong{Prece} & \strong{Polija} & \strong{Lielbritānija} & \strong{Vācija} \\
\hline
Kvieši & 49 & 159 & 81 \\
\hline
Rudzi & 164 & - & 118 \\
\hline
Kartupeļi & 762 & 109,9 & 655,8 \\
\hline
Cukurs & 11,3 & 46,6 & 28,1 \\
\hline
Ogles & 886 & 3~762 & 1152 \\
\hline
Kokvilna & 2,1 & 13,2 & 4,0 \\
\hline
Elektroenerģija (KWh) & 84,0 & 351,2 & 773,2 \\
\hline
\end{tabularx}
\end{table}

Cenu saglabāšana augstā līmenī vai pat to pieaugums veda pie plašu iedzīvotāju slāņu straujas pirktspējas krišanas. No 1928./1929.~līdz 1933./1934.~gadam gaļas, tauku, miltu, apģērba un apavu patēriņš samazinājās trīs-četras reizes. Zināmu priekšstatu par to sniedz \ref{tab:table11}.~tabula.

\noindent
\begin{table}[h!]
\caption{Pārtikas produktu personīgais patēriņš gadā uz vienu iedzīvotāju Polijā un Vācijā 1932.~gadā} \label{tab:table11}
\begin{tabularx}{\linewidth}{|p{3cm}|p{3cm}|p{3cm}|}
\hline
\strong{Produkts} & \strong{Polija} & \strong{Vācija} \\
\hline
Piens & 84 l & 150 l \\
\hline
Olas & 58 gab. & 152 gab. \\
\hline
Sviests & 2,4 kg & 45,9 kg \\
\hline
Maize & 179 kg & 112 kg \\
\hline
Kartupeļi & 203 kg & 154 kg \\
\hline
\end{tabularx}
\end{table}

% page 579


Laukos algādži un trūcīgie zemnieki pārtika galvenokārt no kartupeļiem, kuru vārīšanai sāls taupīšanas nolūkos vairākkārt tika izmantots sālīts ūdens, mazus bērnus ziemā ietina maisos ar salmiem, lai tie nenosaltu nekurinātā miteklī. 3/4 zemnieku bērnu slimoja ar tuberkulozi un rahītu. Jaundzimušo bērnu mirstības līmenis strādnieku ģimenēs vidēji sastādīja no 35~līdz 50\%. Tai pat laikā pat smagākajos krīzes gados netika samazināti izdevumi armijas vajadzībām, kura taču balstīja valdošo režīmu. Neskarts palika arī policijas budžets (1931.~gadā izdevumi policijas uzturēšanai sastādīja 135~miljonus zlotu).

Ievērojami cieta smagās rūpniecības strādnieki, īpaši metalurgi un ogļrači. 1933.~gadā bezdarbnieku skaits pārsniedza nodarbināto skaitu lielrūpniecībā un vidējos uzņēmumos. No darbu saglabājušajiem katrs trešais (tekstilrūpniecībā~--- katrs otrais) strādnieks bija nodarbināts nepilnu nedēļu. Reālā darba alga 1934.~gadā bija zemāka nekā puse algas 1928.~gadā.

Neraugoties uz atbalstu uzņēmējiem, mazinājās investīcijas, auga bezdarbs. Palīdzība bezdarbniekiem bija slikti organizēta. Interesanti atzīmēt, ka oficiālā statistika rādīja kā strādājošo, tā arī bezdarbnieku skaita samazināšanos. Ja 1931. gadā skaitījās 313~000~bezdarbnieku, tad 1932.~gadā bija reģistrēti tikai 220~000~bezdarbnieku. Pabalstus saņēma tikai 6\%~bezdarbnieku, to apjomi pakāpeniski saruka, saņemšanas laiks samazinājās līdz 13~nedēļām. Strādniekiem, nesaņemot darbu un arī bezdarbnieka pabalstu (1931.~gadā to saņēma 107~000~bezdarbnieku (1/3~oficiāli atzīto), 1933.~gadā~--- tikai 49~000), nebija jēgas stāties uzskaitē darba biržās. Iknedēļas bezdarbnieka pabalsts samazinājās no 16~zlotiem 1931.~gadā līdz 10~zlotiem 1933. gadā. Kā atzīst krievu vēsturnieki, atsevišķos gados bezdarbs sasniedza apmēram pusi no strādnieku kopējā skaita.

Oficiālajiem datiem par bezdarbnieku skaitu nevar uzticēties. Tā 1930.~gada martā oficiāli bija 659~000~strādājošo un 290~000~bezdarbnieku, kas kopā sastāda 849~tūkstošus strādnieku. Taču 1932.~gada jūlijā bija palikuši vairs tikai 471~tūkstotis strādājošo un arī bezdarbnieku skaits samazinājās līdz 218~000, kas kopā veido 689~000. Kur divu gadu laikā bija pazuduši 260~000~strādnieku, nemaz nerunājot par jauniešiem, kuriem šo gadu laikā bijā jāstājas ražotāju rindās, tā arī palika Polijas oficiālās statistikas noslēpums. Pēc oficiālās statistikas datiem Polija 1931.~gadā bija 313~000~bezdarbnieku. Taču pēc tautas skaitīšanas datiem 1931.~gada 9.~decembrī ārpus lauksaimniecības bija 681~tūkstotis bezdarbnieku. Lai arī daļa no viņiem varbūt strādāja gadījuma darbus, tomēr jākonstatē, ka oficiālie dati ir pazemināti (pēc būtības~--- viltoti).

Zināmu priekšstatu par reģistrēto bezdarbnieku skaita dinamiku sniedz \ref{tab:table12}.~tabula.

\noindent
\begin{table}[h!]
\caption{Oficiāli reģistrēto bezdarbnieku skaits Polijā (tūkstošos cilvēku)} \label{tab:table12}
\begin{tabularx}{\linewidth}{|p{3cm}|p{3cm}|p{3cm}|p{3cm}|}
\hline
1920.~gads & 52 & 1930.~gads & 299 \\
\hline
1921.~gads & 32 & 1931.~gads & 313 \\
\hline
1922.~gads & 75 & 1932.~gads & 220 \\
\hline
1923.~gads & 68 & 1933.~gads & 343 \\
\hline
1924.~gads & 162 & 1934.~gads & 414 \\
\hline
1925.~gads & 311 & 1935.~gads & 403 \\
\hline
1926.~gads & 190 & 1936.~gads & 466 \\
\hline
1927.~gads & 165 & 1937.~gads & 470 \\
\hline
1928.~gads & 126 & 1938.~gads & 456 \\
\hline
1929.~gads & 185 & ~ & ~ \\
\hline
\end{tabularx}
\end{table}

% page 580

Pieauga arī daļējais bezdarbs. Ja apstrādājošajā rūpniecībā 1928.~gadā nepilnu darba nedēļu bija nodarbināti 12,1\% visu strādnieku un kalpotāju, tad 1929.~gadā tādu bija 15\%, 1930.~gadā~--- 20,5\%, 1931.~gadā~--- 24,7\%, bet 1932.~gadā~--- 35,7\%. Pēc poļu tautsaimniecības vēsturnieka L.~Grosfelda vērtējuma grūtākajos krīzes gados rūpniecībā un sīkrūpniecībā, tirdzniecībā strādājošo strādnieku, mājkalpotāju un arī garīgā daba darītāju vidū bezdarbnieku skaits pārsniedza 1~miljonu cilvēku.

Krīzes apstākļos ļoti apgrūtināta bija strādnieku ekonomiskā cīņa, jo par streikošanu viņi tika nekavējoties atlaisti. (Tomēr 1929.--1933.~gadā notika 2~336~streiki, streikoja 20,4~tūkstoši uzņēmumu, streikotāju skaits sasniedza vairāk nekā 1~miljonu cilvēku.) Strādnieki meklēja jaunas cīņas formas. Izplatījās rūpnīcu ieņemšana, kas neļāva tur ievest streiklaužus. Starptautiskajā strādnieku kustībā šī cīņas forma ieguva ``poļu streika'' nosaukumu. Pirmoreiz to poļu strādnieki pielietoja 1931.~gada janvārī. Sākot no 1932.~gada vasaras ``poļu'' jeb ``okupācijas'' streiki kļuva masveidīgi. 1935.--1937.~gadā tie sastādīja apmēram 1/2 no visiem streikiem.

Uzņēmumu ieņemšana kļuva tik populāra cīņas forma, ka to pielietoja arī garīgā darba darītāji. Piemēram, vairākkārt neatstājot savas darba telpas, streikoja Varšavas maģistrāta kalpotāji. Komēdijas teātra aktieri nedēļu neatstāja teātra telpas, prasot izmaksāt aizturēto algu.

Netika atstātas novārtā arī citas streiku formas. Tā, 1932.~gada 16.~martā PPS līderi un Arodbiedrību Centrālā komisija sarīkoja vispārēju streiku. Kopā streikoja ap 300~000~strādnieku.

Aktīva streiku kustība turpinājās arī pēc 1933.~gada. Streiku statistiku rāda \ref{tab:table13}.~tabula.

\noindent
\begin{table}[h!]
\caption{Pārskats par streikiem Polijā 1930.--1936.~gadā} \label{tab:table13}
\begin{tabularx}{\linewidth}{|p{2cm}|p{2cm}|p{2cm}|p{2cm}|p{2cm}|}
\hline
\strong{Gads} & \strong{Streiku skaits} & \strong{Streikiem aptverto uzņēmumu skaits} & \strong{Streikotāju skaits} & \strong{Streiku rezultātā zaudēto darba dienu skaits} \\
\hline
1930 & 312 & 1~185 & 48~000 & 273~000 \\
\hline
1931 & 357 & 1~154 & 107~000 & 601~000 \\
\hline
1932 & 504 & 6~219 & 314~000 & 2~101~000 \\
\hline
1933 & 631 & 7~282 & 343~000 & 3~829~000 \\
\hline
1934 & 946 & 9~416 & 363~000 & 2~356~000 \\
\hline
1935 & 1~165 & 11~631 & 450~000 & 2~008~000 \\
\hline
1936 & 2~056 & 22~016 & 675~000 & 3~950~000 \\
\hline
\end{tabularx}
\end{table}

% page 581

Kā rāda tabula, streiku vēriens pieauga. 1934.~gadā Polijā streikotāju skaits bija lielāks nekā viņu skaits Austrijā, Beļģijā, Čehoslovākijā, Francijā, Igaunijā, Lielbritānijā, Norvēģijā, Rumānijā, Somijā, Šveicē un Zviedrijā kopā.

Vērsās plašumā bezdarbnieku nemieri. 1931.~gada janvārī Radomā (\pltxti{Radom}) viņi uz vairākām dienām sagrāba pārtikas noliktavas un organizēja produktu sadali. Februārī Varšavā demonstrācijā izgāja ap tūkstoti bezdarbnieku bērnu, nesot lozungus ``Darbu mūsu tēviem!'', ``Maizi mums un mūsu mātēm!''. Policija piekāva un izdzenāja demonstrantus. Pēc nepilnīgiem Komunistiskās Internacionāles Izpildu komitejas datiem no 1930.~gada oktobra līdz 1932.~gada martam Polijā notika 617~bezdarbnieku uzstāšanās, kurās piedalījās vairāk nekā 550~000~cilvēku.

Varšavā 1932.~gadā ik dienas līdz 20~cilvēku bada un bezdarba dēļ beidza dzīvi pašnāvībā. Pēc oficiāliem, nepilnīgiem datiem cietumos ieslodzīto skaits no 30 tūkstošiem 1928.--1930.~gadā pieauga līdz 56~tūkstošiem 1935.~gadā. 1931.~gadā lauku kara tiesas, kuras līdz tam darbojās tikai okupētajā Rietumukrainā un Rietumbaltkrievijā, tika ieviestas visas Polijas teritorijā. Pēc oficiāliem datiem lauku kara tiesas 1931.~gadā nāves sodu piesprieda 39~cilvēkiem, 1932.~gadā~--- 127~cilvēkiem. Pēc preses datiem tikai 1931.~gada četros mēnešos tika izpildīti 106~nāves spriedumi, tai skaitā 44~tikai vienā Viļņas vojevodistē.

Smagi cieta ne tikai strādnieki, bet arī sīkražotāji. 1931.~gadā no sīkražošanas nestajiem ienākumiem pilsētās dzīvoja 60\%, laukos 82\% iedzīvotāju. Pēc 1931.~gada no darba lauksaimniecībā vēl joprojām pārtika 60,6\% visu iedzīvotāju, un tikai atsevišķos novados šis skaitlis bija mazāks nekā 50\% (Augšsilēzijā un Cešinā ar lauksaimniecību nodarbojās 12,2\%, Poznaņā~--- 47,1\% un Lodzas novadā~--- 48,6\%iedzīvotāju).

Polijā kā pārsvarā agrārā valstī krīze izsauca īpaši smagas sekas. Lauksaimniecības produkcijas cenas kritās straujāk nekā rūpniecības izstrādājumu cenas. Lauksaimnieki smagi cieta no cenu krišanās uz savu produkciju. Turpretī cenas uz daļu rūpniecības ražojumu saglabājās augstas, jo tās sargāja ražotāju karteļu vienošanās. 1933.~gada decembrī lauksaimniecības produktu cenas bija kritušās līdz 34,4\% no 1928.~gada cenu līmeņa, bet sāls cena sastādīja 104,5\%, dzelzs~--- 84,5\%, minerālmēslojuma~--- 82,1\% no 1928.~gada līmeņa. Ja 1927./1928.~gadā zemniekam, lai iegādātos 10~l petrolejas, bija jāpārdod 13~kg rudzu vai 2,6~kg cūkgaļas vai 18~l piena, tad 1934./1935. gadā jau 33~kg rudzu vai 7,7~kg cūkgaļas vai 31~l piena.

Valsts Lauksaimniecības bankas piešķirtie aizdevumi no 88~miljoniem zlotu 1927.~gadā izauga līdz 167~miljoniem 1928.~gadā un 253~miljoniem 1929.~gadā, taču tikai 9\% no šīs summas saņēma mazzemnieki. Pēdējie bija spiesti meklēt aizņēmumus no privātiem avotiem un maksāt procentu likmi līdz 40--50\% (dažkārt arī augstāku).

Prēmiju sistēma par lauksaimniecības produkcijas eksportu bija labvēlīga tikai lielajiem eksportētājiem.

1933.~gadā, saņemot prēmijas par eksportu, no Polijas tika izvesti 500 tūkstoši t graudu un miltu. Tāpēc cena par 100 kg rudzu iekšējā tirgū pacēlās līdz 24 zlotiem. Tas bija izdevīgi lielražotājiem~--- muižniekiem un turīgajiem zemniekiem, kuri, pateicoties savam ekonomiskajam spēkam un piešķirtajiem kredītiem, varēja noturēt graudus līdz pavasarim nepārdotus. Turpretī sīkie un pat daļa vidējo zemnieki no tā cieta, tāpat kā pilsētu iedzīvotāji. Tūlīt pēc ražas novākšanas, kad vairākums zemnieku, kuri cieta no nodokļu un parādu sloga, bija spiesti pārdot ne tikai pārpalikumu, bet arī daļu no graudiem, nepieciešamiem pašu patēriņam, biržas cena graudiem kritās līdz 13--14~zlotiem, bet vietējā tirgū pat līdz 6--8~zlotiem par 100~kg. Zemnieki laukos, pārdodot graudus, saņēma 10--12~grošu par 1~kg kviešu, bet patērētāji Varšavā par 1~kg baltmaizes maksāja 1~zlotu. Zemnieki pārdeva 1~litru piena par 7~grošiem, bet strādnieki Varšavā to pirka par 35--40~grošiem. Ar lauksaimniecības produkcijas uzpirkšanu un tālāk pārdošanu iedzīvojās arī daļa pastāvošo kooperatīvu, kuros vadība bija muižnieku un turīgo zemnieku rokās. Toties nodokļu smagums gūlās uz sīkajām zemnieku saimniecībām. Tieši tās ieslīga arvien lielākos parādos. Ja 1929./1930.~gadā nodokļi, ko maksāja zemnieks no 1~ha, līdzinājās 280~kg rudzu cenai, tad 1932./1933.~gadā~--- 380~kg cenai. Lai samaksātu zemes nodokli par 5~ha lielu saimniecību 1927.~gadā zemniekam vidēji bija jārealizē 4,4\% ražas, tad 1934.~gadā~--- 13,9\%.

Jau 1930./1931.~gadā zemnieku saimniecību, it īpaši mazo, rentabilitāte būtiski samazinājās. Tikai turīgās 15--30~ha zemnieku saimniecības guva vidēji 8~zlotu ienākumus uz 1~ha. Vidējie un sīkie zemnieki nesaņēma pat atlīdzību par savu darbu. Paātrinājās iedzīvotāju pauperizācija (slīgšana nabadzībā). Norisa lauksaimniecību degradācija. Samazinājās tehnisko kultūru sējumi, mājlopu skaits, kritās vidējā ražība. Ja 1925.~gadā zirga nebija ap 45\% saimniecību, tad 1931.~gadā tādu saimniecību bija jau 75\%. Pēc Polijas ekonomiskās konjunktūras un cenu pētniecības institūta datiem lauksaimniecības mašīnu un inventāra iepirkumu indekss (1928.~gads =100) 1932.~gadā nokrita līdz 8,8. Tas nozīmē, ka Polijas lauki 1932.~gadā iepirka 11~reizes mazāk mašīnu un inventāra kā 1928.~gadā. Poļu sabiedriskais darbinieks, Seima deputāts V.~Zaklika šai sakarā runāja par ``poļu zemnieka neizmērojamo spēju atteikties ne tikai no kulturālajām, bet arī no elementārajām cilvēciskajām vajadzībām''.

Tiesa, krīzes atnestās neērtības izjuta arī lielie zemes īpašnieki. Piemēram, kņazam A.~Radzivillam krīze nozīmēja nepieciešamību no savas pils 140~ istabām apkurināt tikai 28, atlaist 40 no 300~viņa muižas saimniecībā strādājošajiem, tomēr paturot visus pils kalpotājus un savus personīgos sulaiņus.

Jāpasaka arī par bezdarbu lauksaimniecībā strādājošo algādžu vidū un t.s. slēpto bezdarbu uz laukiem. Četru krīzes gadu laikā darbu zaudēja ap 25\% laukstrādnieku. Emigrācijas samazināšanās un it īpaši sezonas darba emigrācijas ierobežošana uz Vāciju (1926.--1930.~gadā uz Vāciju vidēji ik gadus sezonas darbā devās ap 362,6~tūkstoši poļu laukstrādnieku, bet 1932.~gadā tādu bija tikai 400~cilvēku. Ja līdz krīzei no valsts izbrauca par ap 100~000 cilvēku vairāk nekā iebrauca tajā, tad 1931.--1933.~gadā Polijā iebrauca par 11~000 cilvēku vairāk, nekā izbrauca) radīja laukos līdz 1~miljonam bezdarbnieku. Jau līdz krīzei slēptais bezdarbs jeb ``lieko darba roku'' skaits pēc dažādiem vērtējumiem Polijā sastādīja no 5~līdz 9~miljonu cilvēku. Viņi nevarēja atrast pilna laika noslogotību nedz laukos, nedz pilsētā. Laukos viņi palika tāpēc, ka šeit vismaz daļa no viņiem apstrādāja savus niecīgos zemes gabaliņus, piepelnījās gadījuma darbos un varēja dzīvot pusbada dzīvi. Šo cilvēku ekonomiskais stāvoklis krīzes laikā drastiski pasliktinājās. Pēc Starptautiskā agrārā institūta aprēķiniem Polijā lauksaimniecībā nodarbināto strādnieku alga 1933.~gadā sastādīja 582~zlotus, kas līdzinājās 29\% no 1928.~gada līmeņa. Darba diena oficiāli ilga 12--13~stundas. Turīgie zemnieki algādžiem krīzes gados gandrīz pilnībā pārtrauca izmaksāt algu naudā, pārejot uz samaksu natūrā.

Īpaši smags strādnieku un zemnieku stāvoklis bija valsts austrumu vojevodistēs. Piemēram, strādniekiem šeit maksāja ievērojami mazāk nekā valsts centrā. Tā, Varšavā nedēļā strādnieks vidēji pelnīja 39, bet Volīnijā~--- 11~zlotus. Varšavā mazāk kā 10~zlotus pelnīja 4,7\% strādnieku, Volīnijā~--- 41,5\%.

Poļu prese ziņoja par badu t.~s. guculu apdzīvotajos rajonos, (poļu \pltxti{Huculszczyzna}, ukraiņu \uktxti{Гуцульщина}), kur 1932.~gadā to cieta 88,6\% saimniecību iedzīvotāji. Avīze ``\pltxti{Nowe godziny}'' (``Jauna stunda'') rakstīja: ``Visi cilvēki piepamst no bada un mirst. Kopā ar badu izplatījies vēdertīfs un tuberkuloze.'' Jau iepriekšējā apakšnodaļa runāts par ukraiņu zemnieku cīņām Ļvovas vojevodistē 1932.~gadā. Nākamajā, 1933.~gadā kā protests pret nežēlīgu nodokļu un parādu piedzīšanu notika zemnieku sacelšanās Galīcijā, kura aptvēra teritoriju no Gorlices (poļu \pltxti{Gorlice}, ukraiņu \uktxti{Горлиці}) pie Čehoslovākijas robežas līdz pat Krakovas apriņķim. Notika asiņainas sadursmes ar policiju. Ieradās policijas papildspēki, karaspēka daļas no citiem apvidiem, sākās masu aresti un sacēlušos ciemu ``nomierināšanu''.

Poļu publicisti nevarēja nekonstatēt, ka Polijas tautsaimniecībā pastāv briesmīgi paradoksi. Tā, viens no viņiem S.~Prostvils 1932.~gadā rakstīja: ``Polija ir ekonomisko pretrunu zeme. Melno dimantu novadā, ogļu pārprodukcijas zemē~--- ogles ir dārgas. Zemē, kurā cukurbiešu ir atliku likām, cukurs ir dārgs, dārgāks nekā zemēs, kuras importē cukuru. Zemē, kur ir cilvēku pārprodukcija, kur ir liels iedzīvotāju pieaugums, kur ir daudz cilvēku ar milzīgu apetīti, lauksaimniecības produkti neatrod noietu. Lauksaimniecības zeme ar piena upēm un ķīseļa krastiem ir pilna ar bada cietējiem''. Zināmu kopieskatu par Polijas iedzīvotāju stāvoklī sniedz nacionālā ienākuma (angļu \entxti{national income}) salīdzinājums ar citām valstīm. Skat. \ref{tab:table14}.~tabulu.

\noindent
\begin{table}
\caption{Nacionālas ienākums uz vienu iedzīvotāju dažādās valstīs (ievērojot valūtas devalvāciju atsevišķās valstīs) 1929.~un 1932.~gadā} \label{tab:table14}
\begin{tabularx}{\linewidth}{|p{2cm}|p{2cm}|p{2cm}|p{2cm}|p{2cm}|}
\cline{1-5}
\multirow{2}{*}{Valsts} & 1929 & 1932 & 1929 & 1932 \\
\cline{2-5}
& \multicolumn{2}{p{4cm}|}{absolūti zlotos} & \multicolumn{2}{p{4cm}|}{\% no Polijas nacionālā ienākuma uz 1 iedzīvotāju} \\
\cline{1-5}
ASV & 5~700 & 3~426 & 687 & 731 \\
\cline{1-5}
Lielbritānija & 4~577 & 2~876 & 551 & 613 \\
\cline{1-5}
Francija & 2~081 & 1~727 & 251 & 368 \\
\cline{1-5}
Vācija & 2~552 & 1~515 & 308 & 324 \\
\cline{1-5}
Latvija & 1~043 & 674 & 126 & 144 \\
\cline{1-5}
Polija & 830 & 472 & 100 & 100 \\
\cline{1-5}
\end{tabularx}
\end{table}

% page 583


Tiesa, tabulā izmantotie aprēķini vēl ir visai optimistiski. Citi aprēķinu dod bēdīgāku ainu. Pēc Polijas ekonomiskās konjunktūras un cenu pētniecības institūta aprēķiniem nacionālais ienākums Polijā uz vienu iedzīvotāju (pēc nodokļu atskaitījuma, kas vērtēti kā 15\%) 1929.~gadā bija 714~zloti, bet 1933.~gadā 280~zloti.

Zemāko ekonomiskās krīzes punktu Polija sasniedza 1933.~gada martā, pēc tam sāka iziet no tās. 1933.~gadā tika saražots 55,4\% rūpnieciskās produkcijas no 1928.~gada līmeņa. Tautsaimniecības atveseļošanās norisa ļoti lēni. (No 1933.~gada sāka atlabt rūpniecība, bet lauksaimniecība~--- tikai 1935.~gadā. Taču agrārais sektors Polijas ekonomikā joprojām bija galvenais, tajā bija nodarbināti vairāk nekā 60\% iedzīvotāju un tas ražoja lauvas daļu no nacionālā kopprodukta (1929.~gadā 71\%). Tiek uzskatīts, ka pilnībā Polijas saimniecība no krīzes atkopās tikai 1935.~gadā. Tomēr arī tad, ja rūpnieciskā ražošana citās brīvā tirgus valstīs 1937.~gada vidū sasniedza ap 96\%, tad Polijā~--- 85,3\% no 1929.~gada līmeņa. Polijā vēl joprojām nedarbojās 5 no 19~domnām, 9 no 55~martenkrāsnīm, bija samazinājies metalurģisko rūpnīcu skaits. Pēc poļu vēsturnieku Z.~Landau un J.~Tomaševska datiem 1938.~gadā no 90~ogļu šahtām darbojās 60, no 37~dzelzs rūdas ieguvēm~--- 26, no 31~cinka un alvas iegūšanas rūpnīcām~--- 10, no 86~cukura rūpnīcām~--- 61.

Tautā vēl ilgi saglabājās atmiņas par 20./30.~gadu mijas saimniecisko katastrofu. Tās bija dzīvas arī vēl pēc Otrā pasaules kara. 1929.--1933.~gada krīze gan sakompromitēja tirgus kapitālismu, noveda arī pie ``sanācijas'' režīma popularitātes krišanās, tā tomēr neveicināja pāreju atpakaļ pie parlamentārās demokrātijas. Tāpat kā citās valstīs (arī Latvijā) izplatījās ideja par plānveida ekonomikas, ko kontrolētu valsts, realizācijas nepieciešamību. Tā tika cerēts ieviest kārtību tautsaimniecībā. Šī ideja rada atbalsi visās politiskajās partijās.

Pēc ekonomiskās krīzes būtiski palielinājās valsts investīcijas tautsaimniecībā. No 1927.~līdz 1936.~gadam tās pieauga no 51~līdz 433~miljoniem zlotu. 1936.~gadā tika nacionalizēta (pret atlīdzību) virkne ārzemju uzņēmumu ar ārvalstu kapitālu. Tomēr daļa Polijas rūpniecības palika ārzemju kapitālistu kontrolē, kuri daļu savu dividenžu izveda. Pēc oficiāliem datiem 38,4\% Polijas akciju sabiedrību kapitāla piederēja ārzemju investoriem.

Šai laikā Polija bija guvusi zināmus panākumus transporta attīstībā. Tam pamats bija ielikts jau pirms valsts neatkarības atjaunošanas, par ko jau runāts. Vācu okupācijas laikā bijušajā Polijas karalistē esošais platsliežu dzelzceļu tīkls tika pārbūvēts atbilstoši Rietumeiropas standartam. 30.~gadu vidū Polijas rīcībā bija vairāk nekā 18~000~km dzelzceļa līniju (Vācija~--- vairāk nekā 58~000~km), taču apskatot dzelzceļa tīkla biezumu uz 100~km$^{2}$, aina bija cita. Polijā šis rādītājs bija 5,2~ km, (Vācijā 12,4~km, PSRS~--- 0,4~km). No 1920.~līdz 1935.~gadam dzelzceļa līniju garums Polijā pieauga par 1~400~km. Pasažieru vagonu skaits auga no 7~000~līdz 12~000, preču vagonu~--- no 58~000~līdz 160~000. Pirmās lokomotīves pēc kara nācās pasūtīt Filadelfijā (ASV), 30.~gadu vidū Polija jau pati apmierināja savas vajadzības un vēl eksportēja lokomotīves un vagonus uz Rumāniju, Beļģiju un PSRS.

Eksportēta poļu produkcija tika galvenokārt pa jūru. Piemēram, 1937.~gadā 66\% preču eksports norisa pa jūru. Tas ``sanācijas'' režīma propagandistiem radīja pamatu popularizēt ideju, ka ``Polija ir jūras valsts''. Taču tika noklusēts šādas ``jūras tirdzniecības'' patiesais raksturs. Gandrīz visas poļu preces tika izvestas ar ārzemju kuģiem. Ienākumi no preču pārvadāšanas gāja ārzemju kuģu īpašnieku kabatās. Tā, Gdiņas ostā starp ienākušajiem kuģiem vidū 1.~vietu skaita ziņā ieņēma Zviedrijai piederošie. 1930.~gadā Polijas tirdzniecības flote sastāvēja no 25~kuģiem, 1937.~gadā~--- no 59~vidējiem un maziem tvaikoņiem ar kopējo 97~440~t ietilpību (Ir arī citi dati, ka 1935.~gadā Polijai bija 96~kuģi.) Salīdzinājumam, Somijai piederēja 417~tvaikoņi ar 555~000~t tilpumu. Pat Latvijas Republikas tirdzniecības flotē īsi pirms Otrā pasaules jau bija 103~kuģi, no tiem 89~tvaikoņi un divi motorkuģi (pārējie bija motorburinieki un burinieki).

Iekšējie ūdens ceļi 30.~gadu vidū Polijā sastādīja 14~177~km, no tiem 4~900~km bija kuģojami, pārējos varēja izmantot tikai plostu pludināšanai. Pa ūdens ceļiem pārvadāja galvenokārt labību, cukuru un ogles. Taču pārvadājamais preču apjoms nesasniedza pat 1\% no pa dzelzceļu pārvadājamo preču apjoma, tātad~--- nebija īpaši nozīmīgs.

Aviācijas satiksme Polijā attīstījās kopš 1922.~gada. 1929.~gadā tika organizēts uzņēmums Polijas aviācijas līnijas (\pltxti{Polskie Linie Lotnicze}, LOT), kura uzturēja sakarus starp 7~Polijas pilsētām: Varšavu, Katovici, Poznaņu, Gdiņu, Viļņu, Krakovu un Ļvovu. (Maršrutu Varšava~--- Poznaņa~--- Berlīne tā apkalpoja kopā ar vācu ``\detxti{Lufthansu}''). Izmantotas tika ``\detxti{Junkers}'' un ``\detxti{Fokker}'' firmu ražotās lidmašīnas. Dzelzceļa, jūras un aviācijas transporta attīstības nepieciešamība stimulēja tehniskās domas attīstību.

Tomēr rūpnieciskās ražošanas attīstība norisa lēni. Pēc čuguna, tērauda un krāsaino metālu ieguves, elektroapgādes, rūpniecisko iekārtu skaita uz vienu iedzīvotāju Polijā palika vienā no pēdējām vietām Eiropas valstu vidū. 1936.~gadā ogles Polijā ieguva par 25\% mazāk nekā 1913.~gadā. Ja pasaulē naftas ieguve no 1913.~līdz 1936.~gadam pieauga apmēram 5~reizes, Polijā tā samazinājās no 1~100~tūkstošiem~t līdz 500~tūkstošiem~t., kaut arī naftas atradnes nebūt nebija izsmeltas. Priekšstatu par dažādu produkcijas veidu ražošanu Polijā dod \ref{tab:table15}.~tabula.

\noindent
\begin{table}
\caption{Dažādu produkcijas veidu ražošana Polijā 1913.~un 1936.~gadā} \label{tab:table15}
\begin{tabularx}{\linewidth}{|p{2cm}|p{2cm}|p{2cm}|p{2cm}|p{2cm}|}
\cline{1-5}
\multirow{2}{*}{Produkcijas veidi} & \multicolumn{2}{p{4cm}|}{Ražotās produkcijas apjoms tūkstošos t.} & \multicolumn{2}{p{4cm}|}{Produkcija uz 1 cilvēku kg.} \\
\cline{2-5}
& 1913.~gads & 1936.~gads & 1913.~gads & 1936.~gads \\
\cline{1-5}
Dzelzs rūdas ieguve & 500 & 500 & 16,6 & 14,9 \\
\cline{1-5}
Tēraudkausēšana & 1~700 & 1~100 & 65,7 & 32,8 \\
\cline{1-5}
Velmējumu izlaide & 1~200 & 800 & 40,0 & 23,9 \\
\cline{1-5}
Cinka ražošana & 192 & 93 & 6,4 & 3,1 \\
\cline{1-5}
Alvas ražošana & 45 & 15 & 1,5 & 0,4 \\
\cline{1-5}
Superfosfāta ražošana & 196 & 113 & 6,5 & 3,3 \\
\cline{1-5}
\end{tabularx}
\end{table}

% page 585



Polija turpināja eksportēt galvenokārt rūpniecības izejvielas un lauksaimniecības produktus. Piemēram, 30.~gadu vidū Polija izaudzēja 14\% no visas pasaules rudzu (esot šai rādītājā tikai aiz PSRS un Vācijas), 15\%~--- kartupeļu ražas. 1937.~gadā izejvielu izvedums sastādīja 42\%, bet pusfabrikātu un lauksaimniecības produktu izvedums~--- 43\% no visas eksportējamās produkcijas vērtības. Pie tam ap 25\% no visa eksporta vērtības deva kokmateriāli. Turpretī rūpniecisko ražojumu eksports sastādīja tikai 15\% no Polijas eksporta vērtības.

Vājā ekonomiskā attīstība tieši ietekmēja valsts ienākumus. Tiešie nodokļi 1934.--1935.~gada budžetā (525~miljoni zlotu) sastādīja ap 25\% visu valsts ienākumu, bet galvenos ienākumus nesa netiešie nodokļi, kuru smagumu visos laikos iznes vienkāršie patērētāji. Tikai ienākumi no valsts monopoliem (uz sāli, tabaku, spirtu, sērkociņiem) sastādīja 646~miljonus zlotu (ap 30\% visu ienākumu).

Strādnieku reālā alga 1935.~gadā izauga salīdzinājumā ar 1932.~gadu, taču sastādīja tikai 66\% no 1929.~gada algas. Nemazinājās bezdarbnieku skaits. Polija bija vienīgā valsts pasaulē, kur bezdarbs turpināja augt arī no 1933.~līdz 1938.~gadam. 1935.~gadā bija ap 400~000 bezdarbnieku, 1938.~--- ap 450~000.

Ļoti aktuāls pirmskara gados Polijā bija dzīvokļu jautājums un dzīvokļi pēc Eiropas standartiem bija pārapdzīvoti. Ja Lielbritānija bija 4,3\%, bet Vācijā 5,6\% pārapdzīvotu dzīvokļu, tad Polijā tādi bija 44,7\% visu dzīvokļu.

Seimā 1937.~gada sākumā pēc lielrūpniecības un vidējās rūpniecības uzņēmumu apsekošanas tika ziņots: ``Uz 10~tūkstošiem iedzīvotāju mums ir 3,1 ārsts \citespace{} Iedzīvotāju pārtika ir slikta \citespace{} Algu krišanās pēdējos 6~gados sastāda kopumā 22~procentus''.

Streiku skaits 30.~gadu otrajā pusē tikai pieauga. Ja 1934.~gadā valstī kopā notika 946~streiki, 1935.~gadā jau 1~165, 1936.~gadā~--- 2~056, bet 1937.~gadā~--- 2~078. Uzņēmumu skaits, kuros notika streiki, šai periodā pieauga no 9~tūkstošiem līdz 25~tūkstošiem.

Joprojām cenas uz lauksaimniecībā nepieciešamajiem ražojumiem palika nesamērīgi augstas salīdzinājumā ar cenām uz pašu zemnieku saražoto pārtiku. Skat. \ref{tab:table16}.~tabulu

\noindent
\begin{table}
\caption{Rūpniecības un lauksaimniecības preču cenu salīdzinājums zlotos 1927./28.~un 1937./38.~gadā} \label{tab:table16}
\begin{tabularx}{\linewidth}{|p{2cm}|p{2cm}|p{2cm}|p{2cm}|p{2cm}|}
\cline{1-5}
\multirow{2}{*}{Rūpniecības ražojumi} & \multicolumn{2}{p{4cm}|}{Rudzu graudi (kg)} & \multicolumn{2}{p{4cm}|}{Cūkgaļa (kg)} \\
\cline{2-5}
& 1927./28.~g. & 1937./38.~g. & 1927./28.~g. & 1937./38.~g. \\
\cline{1-5}
Arkls & 100 & 141 & 21 & 34 \\
\cline{1-5}
Sāls (10 kg) & 8 & 16 & 17 & 37 \\
\cline{1-5}
Petroleja (10 l) & 13 & 18 & 26 & 44 \\
\cline{1-5}
Ziepes (10 kg) & 52 & 59 & 11 & 17 \\
\cline{1-5}
Superfosfāts (100 kg) & 31 & 46 & 6 & 11 \\
\cline{1-5}
\end{tabularx}
\end{table}

% page 587


Pēc padomju vēsturnieku vērtējuma šīs t.s. ``cenu šķēres'' Polijā bija tik plaši atvērtas kā nekur citur Eiropā (Tiesa, netika ņemta vērā pati PSRS, kur nodibinātie kolhozi atradās pilnīgā atkarībā no valsts un tika izmantoti līdzekļu gādāšanai industrializācijas vajadzībām.) Prēmijas par labības un cukurbiešu eksportu uz ārzemēm saņēma galvenokārt akciju savienībās apvienojušies lielie zemes īpašnieki, kuri lielā daudzumā ražoja preču produkciju. Tāpēc nebija brīnums, ka ap 20~miljonu poļu zemnieku pirka rūpniecības izstrādājumus par trīs reizes zemāku summu nekā to darīja 5,5~miljoni Čehoslovākijas zemnieku.

Pat pēc agrārās reformas īstenošanas skaitliski nelielās muižniecības īpašumā atradās ap ceturtdaļa labākās aramās un citas lauksaimniecībā izmantojamās zemes. Visā starpkaru periodā muižnieki zaudēja tikai ap 20\% no saviem zemes īpašumiem. Tos iepirka galvenokārt turīgā zemniecība. 1935.~gadā 2~111 tūkstošu zemnieku saimniecību bija mazākas par 5~ha. (No 1925.~līdz 1935.~gadam no 5~līdz 10~ha lielo zemnieku saimniecību skaits samazinājās par 19\%) No tām 1~109~tūkstoši bija mazākas par 2~ha. Laukos bija nodarbināti 2~800~000 algādžu (17\% lauku iedzīvotāju), kam vispār nebija savas zemes. Vēl joprojām 40\% zemnieku saimniecību nebija zirga, bet 20\% saimniecību~--- tikai viens zirgs. Tikpat neiepriecinoša aina bija vērojama govju skaita ziņā. 14\% zemnieku saimniecību tādu nebija, bet 50\% saimniecību bija tikai viena govs.

Toties zemnieku parādi valstij un bankām, pēc oficiāliem datiem, 1936.~gadā sasniedza gandrīz 4~miljardus zlotu, bet parādi augļotājiem vairākkārt pārsniedza šo skaitli. (Zemnieku parādi par nenomaksātajiem nodokļiem vien 1934.~gadā sastādīja 140~ miljonus zlotu, bet 1938.~gadā šie parādi izauga līdz 608~miljoniem zlotu.) Sevišķi smagi lauksaimniecību skāra 1936.~un 1937.~gada neražas. Šajos gados Polijas lauki iepirka četras reizes mazāk mašīnu, iekārtu, minerālmēslojuma un divas reizes citu rūpniecības preču nekā 1928.~gadā. Piemēram, pēc 1936.~gada sausuma vairāk nekā 100~000~cilvēku badojās Viļņas vojevodistē, pārtika no sacirstiem salmiem, sajauktiem ar bojātiem kartupeļiem un nezālēm, kā arī koku mizas. Galveno graudkopības kultūru (rudzu un kviešu) ražošana uz vienu iedzīvotāju Polijā no 1911.~gada līdz 1936.~gadam samazinājās no 255 kg līdz 244~kg gadā. Var piebilst, ka, ja 1913.~gadā kviešu ražība vidēji Polijā bija 12,4 cnt/ha, tad 1931.--1938.~gg. periodā tā samazinājās līdz 11,9~cnt/ha. Cukurbiešu ražošana no 1911.~gada līdz 1936.~gadam samazinājās par 32\%.

Pēc ``Deklarācijas par spēka nepielietošanu starp Poliju un Vāciju'' parakstīšanas 1934.~gada 26. janvārī (par to tiks runāts turpmāk) 7.~martā tika parakstīts arī ``Protokols par poļu-vācu ekonomisko attiecību normalizēšanu'', ar kuru tika izbeigts abu valstu ``muitas karš'' un noslēgts ``saimniecisks miers''. Ar to tika likvidēta barjera Polijas lauksaimniecības produkcijas ceļā uz Vācijas tirgu. Taču attiecību noregulējumam bija arī otra puse~--- pieauga Polijas atkarība no Vācijas. 1938.~gadā tā atkal ieņēma pirmo vietu Polijas importā (23\%) un eksportā (24,1\%).

30.~gadu otrajā pusē Polijas tautsaimniecība varēja izmantot daļēju starptautiskās konjunktūras uzlabošanos, centās plānot ilgtermiņa valsts pasākumus. Kā vienam no sabiedrības konsolidācijas līdzekļiem vajadzēja kalpot valdības izstrādātajam Centrālā rūpniecības rajona (\pltxti{Centralny Okręg Przemysłowy}) uz dienvidiem no Vislas un Sanas saplūšanas vietas Piekarpatu zemienē ap Sandomiras (\pltxti{Sandomierz}) pilsētu četrgadu attīstības plānam. Tā izbūve sākās 1937.~gadā. To sauca arī par ``drošības trīsstūri'', jo rajons atšķirībā no Augšsilēzijas rūpniecības rajona atradās diezgan tālu no apdraudētajām valsts rietumu robežām. Šeit tika būvētas aviācijas un mašīnbūves rūpnīcas, ķīmiskās rūpniecības uzņēmumi. Plāna ietvaros bija jāuzceļ arī 17 uz militārās produkcijas ražošanu orientēti uzņēmumi, kam bija jāražo moderni ieroči, gan pēc licencēm, gan pašu poļu konstruēti. Vismaz daļēji šim plānam par paraugu bija ņemts no 1936.~gada Vācijā realizētais četrgadu plāns. Finansēts plāns tika ar valsts, daļēji arī no Francijas aizņemtiem līdzekļiem. Līdz 1939.~gadam Polijas Centrālajā rūpniecības rajonā bija jau uzcelti vai sākti būvēt vairāki uzņēmumi. Taču izejvielu atradnes atradās salīdzinoši tālu no šī apvidus. Daudzi kritizēja Polijas pūliņus, jo nebija pirmkārt izbūvēta infrastruktūra: pievadceļi u.c. 1937.~gada 29.~oktobrī Krakovas avīze ``\pltxti{Czas}'' (``Laiks'') rakstīja: ``Tagadējā stāvoklī jaunais industrijas centrs atrodas sava veida salas stāvoklī, bez saistības ar pārējo zemi. Bez saistības, burtiskā nozīmē, bez pārspīlējuma.'' Izņemot bijušās Vācijas teritorijas, Polijas lauku ceļi, skatot rietumeiropieša acīm, atradās autosatiksmei nepiemērotā stāvoklī, kur pat Amerikā ražotie auto lauza šasijas. Daudzi tūkstoši poļu izceļoja uz Centrālo rūpniecības rajonu darba meklējumos, taču tur sastapās ar situāciju, ka, viņiem ierodoties, strauji cēlās dzīvokļu un pārtikas cenas. Nespējot sevi nodrošināt, ieceļotāji atkal devās uz laukiem. Vācu autori šo situāciju raksturoja kā mājas celšanu nevis no pamatiem, bet gan no jumta. (Tiesa gan, viņi arī norādīja, ka Polijai vispār nav vajadzīga industrializācija, tā varētu attīstīt tikai lauksaimniecību, mainot tās produktus pret Vācijas rūpniecības izstrādājumiem, nevis veidot tādas ``mākslīgas konstrukcijas kā Gdiņa, jūras politika, industrializācija utt. '' Polijas saimnieciskajai politikai tika pārmests, ka tā veidota ne pēc saimnieciskajām, bet ārpolitikas prasībām. Jāatzīst, ka šādu ``draudzīgu'' ieteikumu realizēšana varēja vest tikai kā pie saimnieciskas, tā politiskas Polijas atkarības no Vācijas.)

Lai nu kā, 1936.--1938.~gadu posms bija pirmais reālais ekonomiskais uzplaukums starpkaru Polijas vēsturē, kas saistīts ne tikai ar pirmskara potenciāla atjaunošanu, bet arī plaša mēroga rūpnieciskās ražošanas modernizāciju.

Tiesa, kaut rūpniecība zināmā mērā tika modernizēta, tai pietrūka attīstības iespēju. Bez jau augstāk minētajiem cēloņiem to traucēja neattīstītā enerģētiskā bāze. Elektrības ražošana salīdzinājumā ar attīstītajām Eiropas valstīm Polijā bija zemā līmenī. Tai laikā kad Vācija 30.~gadu vidū uz vienu iedzīvotāju tika izmantots 441~kilovats elektroenerģijas, tad Polijā~--- tikai 84~kilovati. Aiz Polijas elektroenerģijas patēriņa ziņā palika tikai PSRS un Rumānija. Pie tam no 64~lielākajām elektrostacijām 31~atradās valsts rietumos, un 26~centrā, bet austrumos, kur vajadzība pēc enerģijas bija ne mazāka,~--- tikai viena vienīga elektrostacija.

Kopumā Polija starpkaru periodā bija valsts ar vidēju kapitālisma attīstības līmeni, taču Latvijas žurnāls ``Ekonomists'' 1938.~gada 2.~jūnijā konstatēja: ``Jau tagad Polijas fabriku rūpniecība kopā ar kalnrūpniecību ieņem ļoti redzamu vietu un nodarbina gandrīz 1/4 iedzīvotāju, bet šī rūpniecība koncentrēta samērā nelielā dienvidrietumu Polijas daļā, kur jau priekš kara [domāts Pirmais pasaules karš] izauga spēcīga rūpniecība: tā ir Poznaņas, Varšavas, Lodzas un Katovicas apkārtne.\citespace{} Turpretim gandrīz 2/3~visas valsts teritorijas ar 48\% iedzīvotāju gandrīz nemaz nav industrializēta \citespace{} Rietumu daļa ir bagātāka un kulturālāka, bet austrumu daļa cieš no tehnikas un civilizācijas trūkuma.''

Žurnāls arī atzīmēja, ka Polijas lielākā nelaime ir lielais skaits nenodarbinātu vai maznodarbinātu cilvēku. Polija arī pēc Pirmā pasaules kara palika lauksaimnieciska zeme. Taču, ja pirms kara liela lauku iedzīvotāju daļa izceļoja vai devās meklēt darbu kā laukstrādnieki uz Vāciju un citām zemēm, tāpat augošās pilsētas uzsūca ievērojamu daļu šo lieko laucinieku daļu, tad 30.~gados sakarā ar ekonomiskajām grūtībām ārvalstīs emigrācija uz tām mazinājās, arī Polijas pilsētas pēc pasaules ekonomiskās krīzes spēja uzņemt apmēram 2~reizes mazāk cilvēku nekā 20.~gados. Žurnāls rakstīja: ``Varētu domāt, ka lielais cilvēku daudzums lauksaimniecībā radīs ļoti intensīvu un ražīgu lauksaimniecību Polijā, tomēr dzīves tiešamība rāda pavisam citu ainu. \citespace{} Samērā zemā zemes ražība un samērā lielais cilvēku daudzums rada daudzos apgabalos ļoti trūcīgus dzīves apstākļus un ļoti zemas algas. Polijai nav pēckara gados izdevies laimīgi atrisināt lauksaimniecības jautājumu un jūtami celt lauku iedzīvotāju labklājību \citespace{} samērā lielā zemes platība nespēj pienācīgi uzturēt un paēdināt visus lauku iedzīvotājus.\citespace{} Agrārreforma Polijā \citespace{} bijusi ļoti mērena un devusi darba iespējas tikai niecīgai daļai jaunu zemnieku; tāpat pastāvošas lauku sīksaimniecības (parceles) nav ieguvušas jaunus nepieciešamus piegriezumus. \citespace{} Zemes gribētāju, kas spējīgi zemi apstrādāt, ir vismaz 1~milj. vīriešu cilvēku: bez tam pāri par 1,1~milj. sīksaimniecību pamatoti prasa pēc piegriezumiem.''

No 1921.~līdz 1939.~gadam Polijas iedzīvotāju skaits pieauga no 27~līdz 35~miljoniem, bet nacionālā ienākuma pieaugums būtiski atpalika no iedzīvotāju skaita pieauguma. Tiešas sekas tam bija darba cilvēku dzīves līmeņa stagnācija vai pat pazemināšanās. Pēdējā īpaši jūtama bija austrumu kresos. Zemais ienākumu līmenis atstāja iespaidu arī uz iedzīvotāju veselību. Pēc 1938.~gada Polijas Veselības ministrijas gadagrāmatas (\pltxti{Rocznik lekarski Ministerstwo Zdrowia}) datiem 1937.~gadā 87\% bērnu līdz 10~gadu vecumam slimoja ar tuberkulozi. Tautas veselību smagi apdraudēja venēriskās slimības. No tām cieta 27\% iedzīvotāju.

Arī 30.~gadu beigās nerima streiku kustība. 1934.--1938.~gadā Polijā notika 7~792~streiki, kas aptvēra 78,7~tūkstošus uzņēmumu. Streikos piedalījās vairāk nekā 2,3~miljonu cilvēku.

Starpkaru Polijas ekonomisko vājumu ne vienu reizi vien atzīmēja tā laika poļu zinātnieki. Piemēram, ievērojamais poļu ekonomists S.~Svaņēvičs, kurš Otrā Pasaules kara laikā bija represēts PSRS, bet pēc kara dzīvoja Rietumos, 30.~gadu otrajā pusē, pabijis Vācijā, vērtējot Austrumprūsijas, Lietuvas un Polijai piederošo Viļņas un Grodņas novadu attīstību, konstatēja, ka ``starpība starp reģiona atsevišķam daļām ir milzīga. Austrumprūsijā es vēroju rosīgu celtniecību, bet pie mums~--- stagnācija un marasms.''

Ir saprotams, ka izteikti negatīvs par 20.-30.~gadu Polijas ekonomisko attīstību bija poļu komunistu viedoklis. 1957.~gadā Polijas Apvienotās strādnieku partijas pirmais sekretārs V.~Gomulka norādīja: ``Masu bezdarba posts pilsētās un ``liekie'' cilvēki laukos, posts, kurš ne tikai nemazinājās, bet sistemātiski, kaut zigzagveidīgi, arvien vairāk pieauga un saasinājās, kā drūma ēna gulstas uz visiem 20~pirmskara Polijas pastāvēšanas gadiem \citespace{} Polija bija klasiska ekonomisko emigrantu zeme''. Poļu publicists A.~Vasiļevskis 1985.~gadā rakstīja, ka divdesmit gadu ilgušais neatkarīgās Polijas periods, kurai tad neviens netraucēja izmantot Rietumu valstu darbinieku~--- no F.~Foša līdz N.~Čemberlenam~--- labvēlību, beidzās ar to, ka neraugoties uz visas sabiedrības piepūli, nacionālais kopprodukts izrādījās mazāks nekā tajās pašās poļu apdzīvotajās teritorijās saražotais 1914.~gadā. ``Pati godkārīgākā mūsu vēsturē likme uz tuvināšanos Rietumiem nepadarīja Poliju par pēc Rietumu mērauklas civilizētu valsti. Tieši pretēji, nolēma to vēl lielākai atpalicībai''. Skarbs ir ebreju izcelsmes poļu komunista, žurnālista un diplomāta V.~Groša vērtējums: ``Līdz septembra [domāts Otrā pasaules kara sākums 1939.~gada septembrī] Polija bija vāja, atpalikusi, rūpnieciski un kulturāli neattīstīta valsts. Polija bija viena no pašām nabagākajām un vājākajām valstīm Eiropā un ieņēma starp tām vienu no pēdējām vietām''.

Šie konstatējumi pēc būtības sakrīt ar virknes ārzemju vērtētāju (arī Latvijas žurnālistu) un pētnieku viedokli. Piemēram, vācu vēsturnieks M.~Aleksanders kā būtisku rādītāju par valsts attīstību atzīmē to, ka līdz 1939.~gadam no Polijas emigrēja vairāk nekā 2~miljoni cilvēku. Mūsdienās vēsturnieki konstatē, ka rūpniecības pamatnozaru ražošanas kopapjoms un ražošana uz vienu iedzīvotāju Polijā pat pirmskara konjunktūras apstākļos 1938.~gadā nesasniedza 1913.~gada līmeni. Valstī esošais ap 5~300 lokomotīvju, ap 27~000~auto (no tiem 7~000~kravas), ap 500~traktoru parks lielā mērā bija novecojis. Armijas rīcībā bija ap 900~tanku, mazāk kā 1~000 lidmašīnu, ar ko nepietika, lai 1939.~gadā atvairītu ienaidnieku. Kā secina vācu vēsturnieks P.~Bonings, 1939.~gadā Polija ``pārkāpa industrializācijas slieksni, taču palika pārsvarā agrāri strukturēta.''

Augstāk izklāstītais materiāls rāda, ka Polijas tautsaimniecība par spīti grūtībām attīstījās, tomēr pirms Otrā pasaules kara Polija nebija kļuvusi par attīstītu industriālu valsti.

\section{Poļu kultūras izaugsme}

\epigraph
{Nevar būt ne patriotiskas mākslas, ne patriotiskas zinātnes. Abas, kā visi augstākie devumi, pieder visai pasaulei.}
{Johans Volfgangs fon Gēte (\detxti{Johann Wolfgang von Goethe})}

\epigraph
{Nav nacionālās zinātnes, tāpat kā nav nacionālās reizināšanas tabulas.}
{Antons Čehovs (\rutxti{Антон Павлович Чехов})}

\epigraph
{Katrs mākslas darbs pieder savam laikam, savai tautai, savai videi.}
{Georgs Vilhelms Frīdrihs Hēgelis (\detxti{Georg Wilhelm Friedrich Hegel})}

\epigraph
{Maldi ir nenovēršami, pārāk daudzi tur atrod savu izdevīgumu.}
{Pjērs Buasts (\frtxti{Pierre Boiste})}

\epigraph
{Trīskāršs slepkava ir tas, kurš nogalina domu.}
{Fjodors Burlackis (\rutxti{Фёдор Михайлович Бурлацкий})}

\epigraph
{Sabiedrība griežas ap sevi kā suns, kurš mēģina noķert savu asti, taču šī kustības šķietamība nav attīstība.}
{Antonio Gramši (\entxti{Antonio Gramsci})}

\epigraph
{Komercliteratūra ir narkotika tautai, līdzīga opijam.}
{Antonio Gramši (\entxti{Antonio Gramsci})}

\epigraph
{Patriotisms ir nepilnvērtīgs, ja tas noraida lielāko cilvēces daļu.}
{Jehudi Menuhins (\entxti{Yehudi Menuhin})}

\epigraph
{Katrs rakstnieks, kurš stājas zem partiju karogiem, agri vai vēlu nonāk dilemmas priekšā~--- vai nu pakļauties vai aizvērties.}
{Džordžs Orvels (\entxti{George Orwell})}

\epigraph
{Saprātīgus cilvēkus vareni valdnieki bieži ienīst.}
{Roterdamas Erasms (\latxti{Desiderius Erasmus Roterodamus})}


Starpkaru desmitgades daudz nozīmēja poļu kultūras un nacionālās pašapziņas attīstībai. Polijas valstiskuma atjaunošana radīja labvēlīgākus nekā līdz tam apstākļus poļu nacionālās kultūras un zinātnes attīstībai. To gan ierobežoja jaunās valsts materiālās iespējas. Tāpēc sekmīgāk attīstījās tās kultūras un zinātnes nozares, kurām nebija vajadzīgi lieli ieguldījumi, vai arī kurām bija prioritāra nozīme valdošā slāņa acīs: ar kurām varēja prezentēt nacionālos sasniegumus, kas pildīja ideoloģiskus, propagandas uzdevumus nacionālās pašapziņas stiprināšanā.

Literatūras un mākslas attīstībā bija vērojamas pretrunīgas tendences. Tieksme apgūt jaunākos Eiropas kultūras sasniegumus savijās ar mēģinājumiem radīt laika garam atbilstošu nacionālu stilu, reālisma jaunu formu meklējumi ar eksperimentiem bez dziļāka mērķa.

Attīstījās poļu \strong{literatūra}. 1920.~gadā radās Profesionālā poļu literātu savienība (\pltxti{Związek Zawodowy Literatów Polskich}), 1924.~gadā~--- PEN-kluba poļu sekcija. No 1933~līdz 1939.~gadam pastāvēja Polijas Literatūras akadēmija (\pltxti{Polska Akademia Literatury}), kuras uzdevums bija veicināt literāro darbību un popularizēt sasniegumus tās jomā. Pastāvēja vairāki literāri žurnāli. No tiem ietekmīgākais bija ``\pltxti{Wiadomości Literackie}'' (``Literārās Ziņas''), kas iznāc Varšavā 1924.--1939.~gadā.

Taču starpkaru perioda literatūra nebija vienota. Literārās parādības var iedalīt divos dažādos periodos. Pirmajam no tiem~--- divdesmitā gadsimta divdesmitajiem gadiem bija galvenokārt raksturīgs optimisms, dzīvesprieks, vitalitāte un eiforija par kara beigām, atgūto brīvību, cerības uz Polijas spožu nākotni.

Otrais posms~--- trīsdesmitie gadi pēc pasaules ekonomiskās krīzes atzīmējami ar pakāpeniski pieaugošu civilizācijas krīzes sajūtu un vispārēju pesimismu. Krīzes ietekmē cilvēkiem zuda psiholoģiskās un sociālās stabilitātes un drošības sajūta. Radās arī bažas par Polijas kā valsts drošību pasaulē. To atspoguļoja arī literatūra. Katastrofas priekšvakars vispilnīgāk bija izteikts S.~Vitkeviča un Krakovas dzejnieku-avangardistu darbos. To tumšie pravietojumi piepildījās 1939.~gadā, kad sākās Otrais pasaules karš.

Literatūras attīstībā pirmie starpkaru perioda gadi bija raksturīgi ar dažādu virzienu, skolu un grupu rašanos. Īpaši vētraini šis process risinājās \strong{dzejā}. Mazskaitlīgas grupiņas radās un drīz pazuda, to manifesti bija ļoti nenoteikti, meklējumus atspoguļojoši. To dalībnieki mainīja uzskatus.

Dzejā liela ietekme bija vecākās paaudzes dzejniekam \strong{L.~Staffam}. Viņa poētiskie simboli bija piepildīti reālām sajūtām, optimisma caurstrāvoti. L.~Staffa sabiedriskās lirikas krājumā ``\pltxti{Tęcza łez i krwi}'' (``Asaru un asins varavīksne'', 1918) atspoguļojās Pasaules kara pieredze. Starpkaru periodā iznākušajos dzeju krājumos ``\pltxti{Wydokie drzewa}'' (``Augstie koki'', 1932), ``\pltxti{Barwa miodu}'' (``Medus krāsas'', 1936) dzejnieka daiļradē iezīmējās pavērsiens uz ikdienas tematiku, parādījās skepticisms un arī humors.

1920.~gadā dzejnieki J.~Tuvims, J.~Lehoņs, A.~Sloņimskis, J.~Ivaškevičs un K.~~Veržinskis nodibināja \strong{grupu ``\pltxti{Skamander}''} (Mitoloģiskas upes vārds, kura, kā rādīts ``\emph{Iliādā}'', aptecēja Troju. Nosaukumu piedāvāja J.~Lehoņs). Vēlāk tai pievienojās vēl virkne dzejnieku. Tā publicēja savus periodiskos izdevumus, kuri ieguva plašu popularitāti inteliģencē. Dzejnieku, kas saistījās ar grupu ``Skamander'', darbi pauž prieku un aizraušanos ar ikdienas dzīvi. Grupa kā galveno lozungu izvirzīja ``būt šīsdienas dzejniekiem'', runāja par ``veco un jauno sakņu sintēzi''. Kā paraugs ``\emph{skamandritiem}'' pirmkārt kalpoja L.~Staffa dzeja. Grupas dalībnieki pasludināja, ka ``mākslas diženums izpaužas nevis tēmās, bet formās'', viņi vēlējās būt ``šīs dienas dzejnieki'', panākt veco un jauno tradīciju sintēzi. Reizē, sekojot principiālai nostājai neizvirzīt kādu noteiktu programmu, kas varētu ierobežot brīvību, grupas dalībniekus vienoja vēlme savienot dzeju ar ikdienu, dzejnieka kā vārda meistara izpratne, kurš nevis paceļas virs pūļa, bet piedalās politiskajās un sabiedriskajās norisēs.

``\emph{Skamandriti}'' arī noraidīja patriotiskas un pilsoniskas kalpošanas ideju sabiedrībai kā politisko apstākļu uzspiestu. (Var teikt, ka viena no visas starpkaru perioda poļu literatūras spilgtākajām iezīmēm bija ``dumpis'' pret ``pienākumu''.) Līdz ar to grupa ``\pltxti{Skamander}'' darbība dažkārt nonāca pretrunā ar valdošo interesēm. Piemēram, J.~Tuvima poēmu ``\pltxti{Bal w Operze}'' (``Balle operā''), kura satīriskā gaismā atainoja Polijas valdību, cenzūra aizliedza. J.~Tuvims~--- vārda brīnumdaris, prata piepildīt dzeju ar varenu lirisku spēku. Pirmskara gados viņa dzejā izskanēja asa fašisma kritika.

J.~Lehoņs interesanti pārveidoja romantisko tradīciju. A.~Sloņimska dzeja bija racionālistiska un retoriska. J.~Ivaškevičs pēc neilgas aizraušanās ar ekspresionismu nostiprinājās estētisma pozīcijās. K.~~Veržinska dzeja pārdzīvoja sarežģītu evolūciju: agrīnajās dzejās esošais dzīvi apstiprinošais lirisms vēlāk piekāpās traģiska skanējuma priekšā.

``\pltxti{Skamander}'' bija pati plašākā starpkara perioda poļu literārās inteliģences apvienība. Grupa darbojās līdz 1926.--1928.~gadam, taču arī pēc tās sairuma sabiedrība ``\emph{skamandritus}'' uztvēra kā zināmu kopumu.

1929.~gadā, paužot savu neapmierinātību, ka patriotisms tiek ekspluatēts valdošās elites interesēs, dzejnieks J.~Tuvims norobežojās no savas agrākās kreisi centriskās pārliecības un aicināja tautiešus:

\begin{verse}
Metiet nost savas šautenes uz ietves!\\
Nafta ir viņu, bet asinis~--- mūsu!\\
Un no galvaspilsētas uz galvaspilsētu\\
Lai skan sauciens, ka mēs vairs neliesim asinis.\\
Pani, jūs mūs nepiemuļķosiet!
\end{verse}

Šie vārdi nebūt nebija kvēls marksistisks aicinājums: J.~Tuvims savu poēmu bija iecerējis kā pacifisku deklarāciju. Tomēr tā palīdz saprast, kādēļ pēc kara J.~Tuvims bija gatavs sadarboties ar komunistu varu.

Kamēr diezgan tradicionālie dzejnieki-\emph{skamandriti} un to sekotāji iekaroja plašu lasītāju aprindu simpātijas, ne tik skaļu, taču ne mazāk svarīgu apvērsumu dzejā veica \strong{avangardistu} (no franču: \frtxti{avant-garde}~--- priekšējā vienība, ar šo vārdu parasti apzīmē cilvēku darbus vai idejas, kas ir eksperimentālas un inovatīvas un kas neievēro sabiedrībā pastāvošās normas, galvenokārt tieši mākslā, kultūrā) grupas, galvenokārt futūristi un t.s. ``Krakovas avangards'' (\pltxti{Awangarda Krakowska}).

Grupa darbojās 1922.--1927.~gadā un pulcējās ap Krakovā iznākošo avīzi ``\pltxti{Zwrotnica}`` (``Pārmija''). Sākotnēji tās pamatvirziens bija opozīcijas izveide pret agrāk valdošajiem romantisma un ``\pltxti{Młoda Polska}'' (``Jaunā Polija'') virziena izstrādātajiem paraugiem. Viens no XX~gadsimta avangarda kustības izvirzītajiem uzdevumiem bija jaunu redzējumu dzejā izveide par skaidru programmu. Avangarda dzejnieki centās atrast teorētisku pamatojumu savai jaunradei. Viņi vadījās no pieņēmuma, ka daiļrades process nav nejaušība, bet noteiktu uzskatu izpausme, kurus var priekšstatīt skaidri formulētās tēzēs. Tāpat kā citas XX gadsimta sākuma avangardistu grupas, arī ``Krakovas avangards'' izstrādāja un 1922.~gadā publicēja savu manifestu ``\pltxti{Miasto. Masa. Maszyna}'' (``Pilsēta. Masa. Mašīna''), kurā aprakstīja jaunās dzejas priekšnoteikumus. Grupa bija arī vienīgā no dažādajām poļu mākslinieku grupām, kas līdz pat detaļām izstrādāja savu dzejas sistēmu. Tās lozungs bija ``\pltxti{Najmniej słów}'' (``Mazāk vārdu'').

Viens no avangarda virziena priekšvēstnešiem Polijā un Krakovas avangardistu grupas \pltxti{Awangarda Krakowska} dibinātājiem, tās galvenais teorētiķis un programmas autors bija dzejnieks, mākslas kritiķis un literatūras teorētiķis T.~Peipers. Viņš publicēja trīs dzejoļu krājumus, kuri veidoja ievērojamu pienesumu poļu konstruktīviskajai dzejai.

Grupas sastāvā darbojās arī J.~Prubošs u.c. Savos agrīnajos darbos ``\pltxti{Śruby}'' (``Skrūves'',1925), ``\pltxti{Oburącz}'' (``Abām rokām'', 1926) J.~Prubošs prezentēja savu dzeju kā jaunu valodas sistēmu, kas raksturojās ar īsiem, bet reizē sarežģītām metaforām. 30.~gados dzejnieka jaunradē (krājumos ``\pltxti{W głąb las}'' (``Dziļi mežā'', 1932), ``\pltxti{Równanie serca}'' (``Sirds vienādojums'', 1938) parādījās sociālā problemātika, sabiedriskā protesta motīvi. Tomēr viņš palika konstruktīvisma pozīcijās, abstrakcionisma un sirreālisma ietekmēts.

1921.~gadā iznāca dzejnieka-futūrista S.~Mlodožeņeca pirmā grāmata~--- ``\pltxti{Kreski i Futureski}'' (``Rindas un futureskas''). 1926.~gadā viņš uzsāka ciešu sadarbību ar laikrakstu ``\pltxti{Kurier Warszawski}'', kur publicēja daudz aprakstu par lauku dzīvi. S.~Mlodožeņecs rakstīja arī lugas, darbus bērniem.

B.~Lesmians savā dzejā stingri pieturējās simbolisma likumiem un dzejā apcerēja filozofiski- metafiziskas tēmas. Pazīstamākās no viņa grāmatām bija ``\pltxti{Łąka}'' (Pļava, 1920) un ``\pltxti{Napój cienisty}'' (``Ēnu dzēriens'', 1936).

1920.~gadā kā satīrisku dziesmu un skiču autors savu literāro darbību uzsāka J.~Bžehva. Viņš sadarbojās ar pazīstamiem kabarē. 1926.~gadā iznāca viņa dzejoļu krājums, 1938.~un 1939.~gadā~--- ļoti populāri dzejas krājumi bērniem. J.~Bžehva arī tulkoja krievu dzejnieku darbus poļu valodā, šo darbību ar panākumiem turpināja arī pēc Otrā pasaules kara,

Bērnu literatūras jomā darbojās un panākumus guva, vairākas prēmijas saņēma K.~Makušinskis. Pēc Otrā pasaules kara daudzi viņa darbi tika filmēti, guva plašu popularitāti.

Dzejniece M.~Pavļikovska-Jasnoževska evolucionēja no vitāla optimisma līdz trauksmaini-katastrofāliem motīviem. Viņas izsmalcinātajām miniatūrām raksturīgs aforisms, ironija. Bez dzejas viņa rakstīja arī drāmas un komēdijas.

Ievērojamākais protesta dzejnieks pret sabiedrisko-politisko stāvokli Polijā 20.~gados bija V.~Broņevskis. Jau apzinīgās dzīves sākumā sociālistiski ievirzīts, pēc prezidenta G.~Narutoviča nogalināšanas viņš tuvinājās padomju paraugiem lirikā, darbojās kreisajos laikrakstos, atklāti izteica neapmierināšanu ar Polijā iedibināto politisko režīmu, tika varas iestāžu vajāts, 1931.~gadā arestēts, lai gan drīz atbrīvots.

Daudzi mākslinieki zaudēja ilūzijas par savas valsts politiku, nacionālajām tradīcijām un nacionālo literatūru. Amerikāņu vēsturniece M.~Šore ir izsekojusi vairāku poļu dzejnieku evolūciju no mākslinieciskā avangarda līdz kreisajiem uzskatiem vai, citiem vārdiem, no atziņām, ka ``Dievs ir miris'' un ``reālisms ir beidzies'' līdz pārliecībai, ka padomju komunisms aizpildīs radušos tukšumu.

Starpkaru otrajā desmitgadē literatūrā ienāca jauna lirisku dzejnieku paaudze, šā laika sociālās un ideoloģiskās krīzes liecinieku. Šai laikā savu literāro darbību uzsāka arī vēlākais Nobela prēmijas laureāts (1980) literatūrā Č.~Milošs, kurš debitēja 1930.~gadā Viļņā. Arī vēlākos darbos Č.~Milošs kavējas atmiņās par Viļņu un mācību laiku tajā. 1934.~gadā Poļu rakstnieku savienība Viļņā piešķīra viņam balvu par debiju dzejā.

\strong{Prozā} starpkaru periodā sadzīvoja tradicionālās formas un stilistiski, ka arī jaunas struktūras meklējumi.

Divdesmito gadu pirmajā pusē izcila loma bija \strong{S.~Žeromska} darbībai. Būdams ļoti jūtīgs pret sociālo netaisnību, viņš savos tā laika darbos reālistiski rādīja vienkāršo cilvēku ikdienas apstākļus, aicināja valdošos slāņu veikt reformas, uzlabot tautas dzīvi. 1924.~gadā viņš tika izvirzīts Nobela prēmijai, taču to par 1904.--1909.~gadā publicēto romānu ``\pltxti{Chłopi}'' (Zemnieki) saņēma V.~Reimonts. 1925.~gadā S.~Žeromskis saņēma valsts literāro prēmiju par romānu ``\pltxti{Wiatr od morza}`` (Vējš no jūras).

Pēckara periodā viņš arī uzrakstīja romānu ``\pltxti{Przedwiośnie}'', (Pavasara priekšvēstnesis, 1924) vienu no augstākajiem poļu kritiskā reālisma sacerējumiem, tuvu sociālistiskām idejām. Tajā varonis ir no revolucionārās Krievijas atgriezies jauns cilvēks, kurš pārdzīvo milzīgu vilšanos sapņos redzētās sociālā taisnīguma valsts vietā sastopoties ar Polijas reālo īstenību. Vēstījums noslēdzās ar komunistu vadītas poļu strādnieku demonstrācijas attēlojumu. Viens no ļudoviešu kustības vadītājiem T.~Nočņickis par romānu rakstīja, ka tas radīts ``satricinājumā par to, ka brīvajā neatkarīgajā Polijā atradās tik daudz šakāļu, hiēnu, plēsoņu, benžu un tik maz cilvēku''. Šis romāns kļuva par impulsu daudzu radošo inteliģentu~--- jau minētā V.~Broņevska, L.~Kručkovska, V.~Vasiļevskas u.~c. uzskatu radikalizēšanās virzienā. Toties nacionālradikāļu politiķi un publicisti no nacionāldemokrātu aprindām organizēja veselu dubļu liešanas kampaņu pret romānu un tā autoru. Ik dienas tika publicētas S.~Žeromska darbu nosodošas atsauksmes. Lūk, daži izteicieni no tām: ``Poļa sirdī nav vietas boļševistiskām simpātijām'', ``Polis, idejiski iegremdēts šai purvā, pārstāj būt polis'', ``Tas ir kaut kāds padomju kāpurs, kurš barojas poļu zemē un ar savu elpu saindē gaisu.'' Šie ``kritiķi'' paziņoja, ka rakstnieks slavinot ``morālu un fizisku izvirtību'', pamācīja viņu, ka tam jātiek vaļā no pārlieka humānisma, kurš aizēnojot horizontu un traucējot pareizi spriest par poļu tautas interesēm: ``Jūs varat pie tam palikt sociālists vai pat komunists, bet esat tik laipns, atmetiet dažādas ebrejiskas iniciatīvas un ebrejisku ideoloģiju! Kļūstiet kaut par velnu, bet lai tas būtu mūsu dzimtais velns, bet ne velns ar peisām [garas matu šķipsnas uz deniņiem, ortodoksālo un ultralabējo ebreju tradicionāls frizūras elements]. Un pie reizes atmetiet jebkuru līdzjūtību pret puspoļiem un pārtrauciet ar tiem jebkurus sakarus''. Kā mūsdienās rakstīja ievērojamais poļu sabiedriskais darbinieks Ā.~Mihņiks, par pamatu šādai vajāšanai kļuva vispārēja vilšanās pēc neatkarības atgūšanas. Tā tika gaidīta kā panaceja [mitoloģisks universāls līdzeklis pret visām slimībām] pret visām likstām, bet patiesībā nabadzība valstī neizzuda, nelieši nekļuva godīgi, kukuļņēmēji neatteicās no kukuļiem. Cīnītāji par brīvību~--- gan patiesie gan šķietamie~--- bija vīlušies. Savu neapmierinātību viņi bija gatavi izliet pār jebkuru tiem netīkamu personību. Šoreiz cieta arī poļu literatūras klasiķis S.~Žeromskis.

Tradicionālo stilu rakstniecībā pārstāvēja M.~Dombrovska un S.~Nalkovska.

\strong{M.~Dombrovskas} pirmajiem stāstu krājumiem ``\pltxti{Dzieci ojczyzny}'' (Tēvzemes bērni, 1921), ``\pltxti{Gałąź czereśni}'', (Kiršu zars, 1922) piemita audzinošs un patriotisks raksturs. Tiem sekoja poētiskas noveles grāmatā ``\pltxti{Uśmiech dzieciństwa}'', (Bērnības smaids, 1923), kurā attēloti bērni. Turpmāk sekoja vēl vairākas bērniem veltītas grāmatas. Labākie rakstnieces stāsti iznāca krājumos ``\pltxti{Ludzie stamtąd}'', (Cilvēki no turienes, 1925) un ``\pltxti{Znaki życia}'', (Dzīvības pazīmes, 1938). Apjomā lielākais rakstnieces sacerējums bija triloģija ``\pltxti{Noce i dnie}'' (Naktis un dienas, 1932--1934).

\strong{S.~Nalkovska} bija debitējusi 1906.~gadā ar romānu ``\pltxti{Kobiety}''(Sievietes) un arī vēlāk daudz tēloja sieviešu likteni. No 1918.~gada viņas sacerējumos bija plaši izvērsta patriotisma, Polijas valstiskuma atjaunošanas tēma. Rakstnieces darbos savijās reālistisks stāstījums ar psiholoģiskā romāna elementiem. Ar romānu ``\pltxti{Teresy Hennert}'' (1923) viņa vērsa asu kritiku pret valdošo slāni un tā iekšējām cīņām par varu. 30.~gados rakstniece aktīvi uzstājās pret ``sanācijas'' režīmu, piedalījās akciju organizēšanā pret politiskajām represijām.

Visas Eiropas literatūras attīstību ievērojami ietekmēja trīs starpkara periodā aktīvi strādājoši poļu rakstnieki: S.~Vitkevičs, V.~Gombrovičs un B.~Šulcs.

\strong{S.~Vitkevičs} bija rakstnieks, dramaturgs, gleznotājs, filozofs un mākslas fotogrāfs, pirmās Polijā avangardisko mākslinieku grupas ``\pltxti{Formiści}'' biedrs, ``tīrās formas'' teorijas radītājs. Zakopanē viņš nodibināja portretu darbnīcu ``\pltxti{Firma Portretowa S.~I.~Witkiewicz}'', daļa no tiem tika radīta narkotiku iespaidā. Viņa daiļradi ir grūti klasificēt, pieskaitīt kādam virzienam, taču vistuvāk stāvēja tā ekspresionismam. Nereti viņu arī pieskaita sirreālistiem. Viņa jaunrade ir tuva F.~Kafkas darbiem, pilniem absurda un baiļu no ārpasaules, taču bez atklāta pesimisma. S.~Vitkeviča romānos ``\pltxti{Pożegnanie jesien}'' (Atvadas no rudens, 1927), ``\pltxti{Nienasycenie}'' (Negausība, 1930) abstrakti-utopiskā formā tēlotas cilvēces drāmas un civilizācijas sabrukums. Rakstnieka darbi tulkoti daudzās valodās un ir ekranizēti.

\strong{V.~Gombrovičs} kļuva slavens tikai savas dzīves beigu posmā, tomēr mūsdienās tiek uzskatīts par vienu no ietekmīgākajiem autoriem poļu literatūrā. Dzīvojot Varšavā, V.~Gombrovičs iekļāvās vietējo intelektuāļu aprindās. 1937.~gadā tika publicēts romāns ``\pltxti{Ferdidurke}'', kurš ieguva atzinību poļu literārajās aprindās. Kopš Otrā pasaules kara sākuma viņš dzīvoja ārzemēs, kur arī tapa viņa nozīmīgākie darbi. Pirmais no tiem Polijā tika publicēts tikai 1986.~gadā. Atbrīvošanos, iziešanu ārpus ``poļu ietvariem'' rakstnieks pārvērta par savas novatoriskās jaunrades galveno tēmu. Starptautiska atpazīstamību rakstnieks sagaidīja tikai 60.~gados, pateicoties daudzajiem viņa darbu tulkojumiem un lugu uzvedumiem.

\strong{B.~Šulcs} bija mūsdienu Ukrainas teritorijā dzimis ebreju izcelsmes poļu rakstnieks, kritiķis un grafiķis. Pazīstamākie darbi~--- stāstu krājums ``\pltxti{Sklepy cynamonowe}'' (Kanēļa veikaliņi, 1934) un romāns ``\pltxti{Sanatorium ``Pod Klepsydrą''}'' (Sanatorija ``Zem smilšu pulksteņa'', 1937). Viņa sacerējumos mazas provinces pilsētiņas ikdienas dzīve kļūst par fantastisku līdzību par pasaules likteņiem. Viņa sacerējumi ir pārcelti uz teātra skatuvi, ekranizētas, kļuvušas par animācijas filmu, muzikālu skaņdarbu pamatiem. Kara laikā viņu ebreju geto teritorijā nošāva kāds esesietis. Virkne viņa literāro darbu un gandrīz visi grafiskie darbi ir pazuduši.

Izcila personība poļu literatūrā bija \strong{A.~Strugs}. Savos pirmajos pēckara darbos ``\pltxti{Odznaka za wierną służbę}'' (Apbalvojums par uzticīgu dienestu, 1921), ``\pltxti{Mogiła nieznanego żołnierza}'' (Nezināma karavīra kaps, 1922) viņš aprakstīja Pirmo pasaules karu un J.~Pilsudska veidoto leģionu kauju ceļu. Vēlāk rakstnieks piedzīvoja vilšanos par neatkarīgās Polijas realitāti, rakstīja darbus, kas parādīja buržuāziskās valsts ikdienu. Tam veltīti romāni ``\pltxti{Pokolenie Marka Świdy}'' (Marka Svidas paaudze, 1925), ``\pltxti{Fortuna kasjera Śpiewankiewicza}'' (Kasiera Spivankēviča karjera, 1928). 1933.~gadā A.~Strugs atteicās no dalības Polijas Literatūras akadēmijā. 1924.~gadā kļuva par Cilvēka un pilsoņa tiesību aizsardzības līgas (\pltxti{Liga Obrony Praw Człowieka i Obywatela}) priekšsēdētāju, no 1936.~gada~--- par Starptautiskās organizācijas palīdzībai revolūcijas cīnītājiem komitejas (\pltxti{Komitetu Centralnego Międzynarodowej Organizacji Pomocy Rewolucjonistom~--- MOPR}) locekli.

Par ievērojamu rakstnieku un žurnālistu kļuva ķīmijas, matemātikas un fizikas skolotājs \strong{L.~Kručkovskis}. Dzejoļus viņš sāka publicēt jau 1918.~gadā. 1932.~gadā iznāca viņa romāns ``\pltxti{Kordian i cham}'' (Kordians un bezkauņa), kas kļuva par literāru atklāsmi, revidējot pastāvošos priekšstatus par politisko un morālo stabilitāti valstī.

Pēc tam viņš radīja drāmu ``\pltxti{Bohater naszych czasów}'' (Mūsu laiku varonis, 1935), publicēja esejas kreisā spārna žurnālos un laikrakstos. Sekoja romāns ``\pltxti{Pawie pióra}'' (Pāva spalvas, 1935)~--- par Galīcijas laukiem pirms Pirmā pasaules kara, un romāns ``\pltxti{Sidła}'' (Lamatas, 1937) kuru pats autors nosauca par ``pēdējo gadu pirms Otrā pasaules kara~--- krīzes, bezdarba un fašisma dokumentu''. Jau tad paralēli ar literāro darbu L.~Kručkovskis nodarbojās arī ar publicistiku. Tā, 1938.~gadā iznāca viņa brošūra ``\pltxti{Dlaczego jestem socjalistą}'' (Kāpēc esmu sociālists) un rakstu krājums ``\pltxti{W klimacie dyktatury}'' (Diktatūras atmosfērā).

Redzama figūra poļu literatūrā bija \strong{J.~Kaden-Bandrovskis}~--- bijušais J.~Pilsudska adjutants un leģiona Pirmās brigādes hronists, senators, Polijas Literatūras akadēmijas ģenerālsekretārs. Viņš aktīvi piedalījās daudzās Polijas sabiedriskajās norisēs, bija konsekvents komunistu pretinieks, ilgstoši atbalstīja J.~Pilsudski, viņa iedibināto ``sanācijas'' režīmu, taču neslēpa arī savu vilšanos par to, savus rūgtos vērojumus valsts augstākajās aprindās attēloja romānos.

Kaut galvenās poļu literatūrā arī starpkaru periodā bija demokrātiskās tendences, Polijā bija arī literāti, kuri aizstāvēja radikālnacionālistiskus uzskatus, vērsās pret nacionālajiem mazākumiem. Diemžēl tieši šie rakstnieki saņēma valdošo slāņu atbalstu. Starpkaru periodā, galēji labējās politiskās nacionālistiskās partijas un strāvojumi, kas dažādos sabiedrības slāņos centās ieaudzināt šovinistisku noskaņojumus, īpaši rusofobiju, evolucionēja pēc Itālijas parauga veidojamas fašistiskās ideoloģijas virzienā. Te var runāt kā par ``sanācijas'' režīma, tā nacionāldemokrātu pārstāvjiem. Tos atbalstīja ne tikai muižniecības un lielburžuāzijas aprindas, bet arī ievērojama daļa pilsētu vidusslāņu, arī atsevišķi strādniecības un zemniecības pulki. Nebija izņēmums arī literāti.

``Sanācijas'' gados virknē periodisko izdevumu parādījās aicinājumi literāriem tieši propagandēt valdošo režīmu. Bija samērā daudz tādu, kas atsaucās aicinājumam. Polijas Literatūras akadēmija savā paspārnē kā reiz centās pulcināt ``sanācijas'' režīma atbalstītājus. Faktiski notika poļu literātu uzpirkšana, ko atviegloja viņu grūtais materiālais stāvoklis. Grāmatu tirāžas Polijā bija minimāls. 2~000~eksemplāru tirāža jau skaitījās liela.

Pēc 1934.~gada Polijas un Vācijas līguma daļā poļu inteliģences valdīja arī jūsma par Vācijas fīreru. Tā ievērojamais poļu dzejnieks un prozaiķis J.~Ivaškevičs pat uzrakstīja himnu, veltītu Ā.~Hitleram.

Taču Polijā fašisma kļūšanai par valsts ideoloģiju bija nopietni šķēršļi. Jau no nacionālās atbrīvošanās cīņas laika darbojās ne tikai nacionālistiskās, bet arī demokrātiskās, sociālistiskās tradīcijas. Pret fašismu cīnījās kreisie ļudoviešu darbinieki, sociālisti, komunisti. Opozīcijā fašistiskām pārvaldes metodēm atradās ne tikai liela daļa strādniecības un zemniecības, bet arī daļa pilsētu sīkīpašnieku, inteliģences un kalpotāju slāņu.

Nacionālistiskā un ``sanācijas'' ideoloģija saņēma daļas inteliģences pretsparu. Tā jau 20.~gados savu literāro darbību sāka J.~Pilsudska līdzstrādnieka L.~Vasiļevska meita, poļu valodas skolotāja, vēlāk vairāku izdevumu ''\pltxti{Robotnik}'' (Strādnieks), ``\pltxti{Naprzód}'' (Uz priekšu), ''\pltxti{Dziennik Popularny}'' (Populāra Dienas Lapa) redaktore \strong{V.~Vasiļevska}. 1934.~gadā iznāca viņas romāns ``\pltxti{Oblicze dnia}'' (Dienas seja), kurā stāstīta stihiskas revolūcijas problēma. 1936.~gada maijā viņa organizēja poļu inteliģences kongresu Ļvovā (kādreiz arī saukts par Antifašistisko kultūras darbinieku kongresu), kurš izteicās par cīņu pret ``sanācijas'' režīmu. Tā kā rakstnieces tēvam bija lieli nopelni valdošo slāņu priekšā (L.~Valsiļevskis bija viens no PPS līderiem un 1918.--1919.~gadā bijis Polijas ārlietu ministrs), meita par to netika arestēta, kaut pati ar tēvu atradās ``kara stāvoklī''.

Pret pastāvošo iekārtu un nacionālistisko virzienu literatūrā cīnījās arī citi, pārsvarā sociālistiski un komunistiski noskaņoti literāti. \strong{S.~Dombrovskis} darbojās kreisajā presē, 1929.~gadā debitēja ar dzejas krājumu ``\pltxti{Pożegnanie Termopil}'' (Atvadas no Termopilām). \strong{S.~Vigodskis} (savu literāro darbību 1928.~gadā uzsāka presē, bet viņa pirmais dzejas krājums ``\pltxti{Apel}'' (Aicinājums) 1933.~gadā iznāca Maskavā. Kara gados viņš bija ieslodzīts vairākās koncentrācijas nometnē, pēc tam atgriezās Polijā, bet 1968.~gadā antisemītiskās kampaņas rezultātā emigrēja uz Izraēlu. \strong{V.~Kovaļskis} bija aktīvs ļudoviešu kustības dalībnieks, darbojās publicistikā. Publicēja esejas, stāstus.

20.~gadu beigās un 30.~gadu sākumā vairāki poļu literāti, vajāti par savu pārliecību, pārcēlās uz dzīvi PSRS.

\strong{V.~Vandurskis} aktīvi darbojās Polijas komunistiskajā presē, bija tās izdevumu redaktors, vairāku lugu autors, 1929.~gadā emigrēja uz PSRS, bet jau 1933.~gadā tika kā ``spiegs'' un ``kontrrevolucionārs'' arestēts un nošauts.

Dzejnieks un rakstnieks \strong{B.~Jasenskis}, sākumā būdams futūristu grupas dalībnieks, 1923.~gada Krakovas sacelšanās ietekmē politiski radikalizējās, emigrēja uz Franciju, kur iestājās Francijas Komunistiskajā partijā. Tur kļuva pazīstams kā romāna ``\pltxti{Palę Paryż}'' (Es dedzinu Parīzi, 1929) autors, tika arestēts un izsūtīts no Francijas, devās uz Padomju Savienību, no 1934.~gada kļuva par Padomju rakstnieku savienības valdes locekli, aktīvi darbojās Proletāriskās literatūras starptautiskajā birojā. Ievērojams viņa panākums bija romāns ``\rutxti{Человек меняет кожу}'' (poļu \pltxti{Człowiek zmienia skórę}, latv. \lttxti{Cilvēks maina ādu}, 1933--1934) par pārveidojumiem padomju Vidusāzijas republikās.

Uz PSRS devās arī dzejnieks \strong{S.~Stande}. Tur iznāca S.~Standes dzejoļu krājumi, literatūrkritiski darbi. Arī B.~Jasenskis un S.~Stande gāja bojā staļinisko represiju gaitā.

Var teikt, ka literatūra Polijā attīstījās par spīti augstajām grāmatu cenām. Salīdzinājumā ar citām dzīves izmaksām grāmatas bija ļoti dārgas. 30.~gadu sākumā no grāmatas cenas 45--48\% sastādīja papīra un drukāšanas izdevumi, 8--10\%~--- autora atlīdzība, bet 45\%~--- tirgotāju pelņa.

Tika drukāti daudz un dažādi periodiski izdevumi. Tā, 1928.~gada sākumā Polijā iznāca jau apmēram 2~200~laikraksti un žurnāli. No tiem poļu valodā bija 1~750 (79,8\%), jidiša un ivrita valodā~--- 170 (7,7\%), vācu~--- 135 (6,1\%), ukraiņu~--- 85 (3,8\%), baltkrievu~--- 25 (1,1\%), krievu~--- 15 (0,7\%), lietuviešu~--- 5 (0,2\%), angļu~--- 5 (0,2\%), itāļu~--- 2 (0,8\%),esperanto~--- 5 (0,2\%), čehu~--- 1 (0,04\%). Apmēram 839~izdevumi bija veltīti dažādām speciālām nozarēm (lauksaimniecībai, dārzkopībai, tirdzniecībai, amatniecībai, rokdarbiem utt.). Bagātīga bija reliģiskā prese~--- vairāk nekā 100~izdevumu, no tiem 90~katoļu. Zinātniska satura izdevumu bija vairāk nekā 50, tie iznāca reizi mēnesī vai reizi trijos mēnešos.

Poļu \strong{teātra attīstība} starpkaru periodā norisinājās visai auglīgi, lai arī pretrunīgi. Kā sēnes pēc lietus izauga jauni teātri. Līdz Otrajam pasaules karam Polijā darbojās 70~stacionāru un pārvietojamu teātru, kuri uzveda izrādes poļu valodā. 1938.~gadā to skatītāju skaits sasniedza 5,5~miljonus. Ar citām valstīm, pat ar par nedraudzīgu uzskatīto PSRS, notika apmaiņa ar labāko teātru viesizrādēm.

Režijā jaunus ceļus meklēja izcilie skatuves mākslinieki L.~Šillers, J.~Osterva un jau minētais S.~Vitkevičs.

\strong{L.~Šillers} 20.~gados vadīja virkni teātru Varšavā, 30.~gados arī Poznaņā, Lodzā, Ļvovā. Augstu māksliniecisko līmeni L.~Šillera režijā sasniedza Varšavas Poļu teātra (\pltxti{Teatr Polski w Warszawie}) izrādes. Īpaši slavens kļuva viņa A.~Mickēciča ``\pltxti{Dzjady}'' uzvedums 1934.~gadā, kurš izraisīja plašu sabiedrisku atbalsi. Šo darbu viņš uzveda arī Ļvovā (1932), Viļņā (1933), Sofijā (1937). Režisors talantīgi izmantoja strauja kāpinājuma un groteskas līdzekļus. Piemēram, tēlojot ``augšslāņu'' pārstāvju pretstatu patriotu grupai Varšavas salona skatā, epizode beidzās ar to, ka pēkšņi aiz skatuves skanēja trauksmes zvans, revolucionārās dziesmas ``\pltxti{Warszawianka}'' melodija, nodzisa gaisma, bet aiz trim krustiem slīdēja mākoņi sarkanā atspīdumā, un tuvojošās revolūcijas nobiedētie Varšavas augšslāņu pārstāvji bēga. ``Sanācijas'' režīma apstākļos Ā.~Mickeviča darba dumpīgās, antidespotiskās virzības atklāsme izsauca valdošo slāņu klaju neapmierinātību. Režisoram tika pārmesta drāmas ``neizpratne un boļševizācija''.

\strong{J.~Osterva} Varšavā nodibināja eksperimentālo ceļojošo teātri ``\pltxti{Reduta}''. Viņa uzvedumos publika nebija nošķirta no aktieriem. Pēc tam J.~Osterva radoši strādāja arī citos teātros. J.~Ostervas vadītais kolektīvs ar panākumiem uzveda arī starpkaru periodā populārākā dramaturga J.~Šaņavska lugas, kurās autors sekoja poļu teātra tradīcijām. ``\pltxti{Reduta}'' bija viesizrādēs arī Rīgā.

Turpretī \strong{S.~Vitkevičs} savos dramatiskajos darbos centās iemiesot modernisma principus. Viņa spalvai pieder vairāk nekā 30~lugu, kurās groteski deformēta dzīves īstenība. S.~Vitkeviču var uzskatīt par ``absurda teātra'' priekšteci, taču viņa darbos izpaudās centieni atspoguļot dzīves jaunās realitātes, jaunos zinātnes un tehnikas sasniegumus, to ietekmi uz cilvēku.

Otrajam pasaules karam sākoties, S.~Vitkevičs beidza dzīvi pašnāvībā. Viņa kapa vieta tā arī palika nezināma. (1988.~gadā Zakopanē it kā tika apbedīti viņa pīšļi, taču vēlāk atklājās, ka tās ir kādas sievietes mirstīgās atliekas.) Tagad S.~Vitkevičs tiek uzskatīts par vienu no nozīmīgākajiem mūsdienu poļu rakstniekiem. Pēc Otrā pasaules kara viņa sekotājs bija ievērojamais poļu teātra darbinieks T.~Kantors (\pltxti{Tadeusz Kantor}, 1915--1990).

Starpkaru periodā Polijā darbojās arī vairāki ebreju teātri.

Polija piederēja tām zemēm, kur jau salīdzinoši agri sāka attīstīties \strong{kinomāksla}. Uz 1918.~gadu poļu kinoprodukcija jau sasniedza 70~dažādu sižetu lentas, neskaitot zinātniski-izglītojošās un dokumentālās. Taču šis mantojums bija ar zemu māksliniecisko vērtību, jo filmas tika uzņemta gandrīz vienīgi komerciālās interesēs. Arī pirmajos Polijas Republikas pastāvēšanas gados uzņemtās kinofilmas neatstāja kaut cik ievērojamas pēdas poļu kultūras vēsturē. Trūka finanšu līdzekļu, valsts iestādes arī neizrādīja interesi par kinomākslu. Tikai pēc 1926.~gada aktivitāte pieauga. Tomēr poļu kino neizvirzīja tādus talantīgus režisorus, kādi darbojās teātra jomā.

Nacionālo mērķu propagandu mēģināja veikt režisors \strong{R.~Boļeslavskis}. Jau 1920.~gadā viņš uzņēma pret boļševistisko Krieviju vērstas īsmetrāžas propagandas filmas, pēc tam divas pilnmetrāžas mākslas filmas ``\pltxti{Bohaterstwo Polskiego Skauta}'' (Poļu skauta varonība) un ``\pltxti{Cud nad Wisłą}'' (Brīnums pie Vislas). Taču nespējot iegūt Pilsētas teātra vadītāja vietu Lodzā, viņš pameta Poliju, 1923.~gadā apmetās Ņujorkā, nodibināja ``\entxti{American Laboratory [Stage] Theatre}'' (Amerikas eksperimentālais teātris), pēc tam strādāja Holivudā. Viņa ieguldījums kinomākslā ir atzīmēts ar zvaigzni Holivudas Slavas alejā.

Vērtīgu pienesumu poļu kinomākslai deva darbinieki, kuri centās pretstatīt savu mākslu komercializācijas tendencēm. Lai gan viņi bija saistīti ar dažādām organizācijām un centriem, grupa ieguva kopēju nosaukumu ``Avangards''. Piemēram, režisors \strong{A.~Fords} uzņēma dokumentālu filmu ``\pltxti{Droga młodych}'' (Jauno ceļš), kuras fabula stāstīja, kā kādā sanatorijā poļu bērni kalnraču streika laikā ielūdz ebreju bērnus un dalās ar viņiem visā, kas vien tiem pieder. Filma beidzās ar aicinājumu uz trūcīgo solidaritāti visā pasaulē. Šādu ievirzi cenzūra iztulkoja kā ``komunistisku tendenci'' un filmu aizliedza.

Redzamu vietu poļu kino ieņēma populāru literatūras darbu ekranizācijas. Ja 20.~gados no visām filmām ekranizējumi sastādīja ap 40\%, tad 30.~gados~--- jau ap 80\%. Tika filmēti S.~Žeromska, H.~Senkeviča, G.~Zapoļskas u.c. darbi. Taču, neraugoties uz iesaistītajiem lieliskajiem aktieriem, filmas mākslinieciskajā ziņā atpalika no to pamatā liktajiem oriģināldarbiem. Pozitīvi var vērtēt dažus J.~Lejtesa darbus. Viņš pirms kara uzņēma septiņas filmas, bet pēc tā turpināja karjeru kino un televīzijā Izraēlā un ASV. Pieminēšanas vērtas ir režisora J.~Gardana filmas.

Līdz Otrajam pasaules karam Polijā tika uzņemtas gandrīz 400~mākslas filmas. Tomēr tikai dažām no tām piemita augsta mākslinieciskā vērtība un nopietns sabiedrisks saknējums. Kopumā poļu kino līdz Otrajam pasaules karam augstas pasaules līmeņa virsotnes nesasniedza.

Taču atsevišķi izcili mākslinieki varēja sasniegt arī pasaules slavu. Tā poļu aktrise B.~Halupeca ar pseidonīmu \latxti{Pola Negri} kļuva slavena Holivudā atveidojot ``liktenīgo'' sieviešu (franču~--- \frtxti{femme fatale}) lomas.

\strong{Mūzikas} jomā poļu mākslinieki centās saglabāt un attīstīt nacionālās tradīcijas.

Visizcilākā poļu muzikālās kultūras figūra starpkaru periodā bija \strong{K.~ Šimanovskis}. Viņš radīja baletu ``\pltxti{Harnasie}'' (Atamani), kantāti ``Stabat Mater'', četras simfonijas, sonātes klavierēm, koncertus vijolei u.c. Kā izcils pianists viņš bieži sniedza koncertus ārzemēs. Vairākus gadus viņš vadīja Varšavas konservatoriju, centās reformēt muzikālās izglītības sistēmu, nonāca konfliktā ar citiem pedagogiem, arī finansiālās grūtībās un mira no tuberkulozes. Neatkārtojams K.~Šimanovska darbu izpildītājs bija vijolnieks P.~Kohanskis.

20.--30.~gadu poļu muzikālajai kultūrai bija raksturīga jaunās paaudzes, kas grupējās ap K.~Šimanovski un ``\pltxti{Jauno poļu komponistu asociāciju}'' (\pltxti{Stowarzyszenie Młodych Muzyków Polaków}), kura 1926.~gadā tika dibināta Parīzē un kurā darbojās P.~Perkovskis, K.~Sikorskis, J.~Maklakevicš, T.~Šeligovskis, B.~Voitovicš, S.~Vehovicš, diezgan asa sāncensība ar vecākās paaudzes komponistiem P.~Rītelu, E.~Moravski, E.~Mlinarski. Salīdzinoši neitrālās pozīcijās atradās B.~Sabeļskis un J.~Koflers.

Pēc ilgāka pārtraukuma pianista karjerā ar koncertu vienā no prestižākajām klasiskās mūzikas koncertzālēm Karnegi hallē (\entxti{Carnegie Hall}) muzikālo darbību atsāka \strong{J.~Paderevskis}. Panākumi bija milzīgi un viņš turpināja koncertdarbību. Tiesa, uzzinājis par 1926.~gada apvērsumu, viņš atkal pievērsās politiskajai darbībai, atklāti uzstājās pret ``sanāciju''.

F.~Šopēna skaņdarbus atskaņoja izcilie vecākās paaudzes pianisti A.~Mihalovskis, J.~Sļivinskis, J.~Turčinskis, kā arī vairāki jaunāki J.~Paderevska skolnieki. 30.~gados izvirzījās izpildītājmākslinieki: pianisti H.~Štompka, S.~Spinaļskis, Z.~Dževeckis, vijolnieces I.~Dubiska un J.~Uminska, operdziedātāji~--- liriski dramatiskais tenors J.~Kepura, kolaratūrsoprāns A.~Sari. Eiropas atpazīstamību ieguva dziedātāja (liriskais kolaratūrsoprāns) E.~Bandarovska-Turska, labākā poļu komponistu vokālo skaņdarbu izpildītāja.

Universitātēs strādāja muzikologi A.~Hibinskis, Z.~Jahimeckis u.c.

Paplašinājās muzikālo biedrību, orķestru loks, auga vidējo un augstāko muzikālo skolu, konservatoriju (Katovicē, Poznaņā) skaits. Tika iedibināta muzikālo konkursu rīkošanas tradīcija.

Kopš 1927.~gada Varšavā regulāri tika rīkots F.~Šopēna vārdā nosauktais pianistu konkurss. Pirmajā debitēja D.~Šostakovičs, vēlākajos gados tajos piedalījās daudzi izcili mūziķi, arī ievērojamākie poļu izpildītāji kā J.~Oleiničaks, K.~Cimermans u.c. Kopš 1935.~gada notika H.~Veņavska vārdā nosauktais vijolnieku konkurss.

1934.~gadā Varšavā tika nodibināts F.~Šopēna institūts, kurš pētīja dižā komponista daiļradi, plānoja izdot viņa sacerējumus, taču materiālu grūtību rezultātā līdz karam tā arī neviens sējums neiznāca. 1939.~gadā Varšavā pirmo reizi tika sarīkots Starptautiskais laikmetīgās mūzikas festivāls, kura tika izpildīti arī jauno poļu komponistu sacerējumi.

\strong{Glezniecībā} turpināja darboties kā jau iepriekšējā periodā sevi pierādījuši mākslinieki, tā arī jauni talanti. Dažādos posmos pastāvēja vairākas mākslinieku grupas ar savu īpatnēju pieeju mākslas problēmām. Izpaudās kubisma un ekspresionisma, arī citu virzienu ietekme. Tā, \strong{T.~Makovska} mākslinieciskais rokraksts veidojās veco Holandes meistaru, naivā reālisma, kubisma un poļu tautas glezniecības iespaidā, ir reizē metaforisks, fantastisks un lirisks. Ja agrīnajos T.~Makovska darbos var atpazīt citu mākslinieku ietekmes, tad vēlīnie ir raksturīgi tikai un vienīgi T.~Makovskim. Viņš izstrādāja savu māksliniecisko valodu, kurā izmantoja XX gadsimta sākumā populāros balagāna, ielu teātra, cirka tēlus. Viņa gleznas piepilda neveiklas būtnes ar ģeometriskām formām, raibās galvas segās, ar caurumiem acu vietā, līdzīgas koka lellēm.

Viens no galvenajiem virzieniem Polijas glezniecībā bija \strong{kolorisms}. 1924.~gadā grupa Krakovas Mākslas akadēmijas studentu, lielākoties bijušā impresionista un simbolista gleznotāja J.~Pankeviča audzēkņu, nolēma doties uz Parīzi, lai tur studētu labāko krāsas meistaru darbus un atjaunotu interesi par krāsu (latīņu \latxti{color}) kā glezniecības pamatu. Viņi arī rīkoja savu darbu izstādes Parīzē, Ženēvā.

Viens no grupas pārstāvjiem \strong{J.~Cibis} pavadīja Parīzē 10~gadus. Pēc atgriešanās dzimtenē viņa gleznas izstādēs ieguva izcilu vietu, jo tās pacēla reālistisko glezniecību jaunā līmenī. Viņš virtuozi pārvaldīja krāsu gammu, centās izteikt savu dabas izjūtu. J.~Cibisa stiprākā puse bija ainavas. Viņš variēja savus mīļākos motīvus, taču nekad neatkārtojās. J.~Cibis gleznoja arī klusās dabas, portretus, izmēģināja spēkus monumentālajā glezniecībā.

Līdz 1939.~gadam Parīzē uzturējās \strong{A.~Naht-Samborskis}, gleznoja ainavas, klusās dabas, žanra gleznas. Viņa daiļradi grūti definēt, tajā izpaudās gan vācu ekspresionisma, gan poļu kolorisma ietekme.

\strong{S.~Vališevskis}, neraugoties uz īso daiļrades posmu atstāja aiz sevis daudz spilgtu, pilnu dzīves prieka darbu.

Cits grupas dalībnieks \strong{P.~Potvorovskis} 1937.~gadā Vispasaules izstādē Parīzē ieguva sudraba medaļu. (Turpat Lielo balvu (\frtxti{Grand Prix}) saņēma viena no poļu labākajām portreta meistarēm O.~Bozņanska, kura gleznoja impresionistu manierē).

Turpināja darboties arī vecākās paaudzes mākslinieki. Lieliskus Polijas dabas skatus gleznoja J.~Falats, S.~Kamockis, S.~Zukovskis un daudzi citi.

Daudz un radoši strādāja \strong{L.~Vičulkovskis}. Starpkaru periodā viņš galvenokārt nodarbojās ar grafiku.

Galvenokārt gravīras gatavoja V.~Skočilass, kuru aizrāva Tatru kalniešu tematika. Viņa ``folkloras'' virziena gravīras piesaista uzmanību ar savu monumentalitāti, mākslinieciskās koncepcijas vienotību, bet reizē smalku dekoratīvismu. Tās atstāja ietekmi uz poļu grafikas skolas veidošanos.

30.~gados arī glezniecībā pieauga sociāli-kritiskās tendences. Liekulību un necilvēcību atmaskoja B.~Linkes satīriskie zīmējumi. 1938.~gadā viņa izstāde tika slēgta kā sociāli kaitīga.

Par brīvās domas un mākslinieciskās jaunrades apspiešanu PPS avīze ``\pltxti{Robotnik}'' (``Strādnieks'') 1932.~gada 22.~jūlijā rakstīja: ``Atklājot gadskārtējo Mākslas akadēmijas izstādi Krakovā, notika sekojošais incidents: izstādē bija eksponētas Akadēmijas studenta Ļevicka grafiskās etīdes, kuras bija atzīmējis profesors Voinarskis. Starp tām bija darbi par tādām tēmām kā policijas cīņa pret pūli Krakovas ielās šī gada martā, ksendzs biktskrēslā u.tml. Pēc akadēmijas vadības rīkojuma šie zīmējumi tika noņemti. Pulkstens 3 dienā akadēmijas ēkā ieradās politiskās policijas pārstāvji, kuri konfiscēja Ļevicka darbus un pēc tam veica viņa dzīvoklī kratīšanu. Pēc kratīšanas izmeklētājs deva pavēli par Ļevicka arestu.'' Tā politika iejaucās arī mākslā.

Reklāmas grafiku un plakātus darināja \strong{V.~Jastžebovskis} u.c.

\strong{Skulptūras} jomā starpkaru periodā Polijā diezgan izteiktas bija dažādas stilizācijas formas, taču reālistiskās tradīcijas arī netika zaudētas.

Ievērojams poļu skulptors šajā laikā bija \strong{K.~Duņikovskis}, kurš Varšavā bija Mākslas akadēmijas profesors. Līdz 1939.~gadam viņš radīja daudz monumentālu tēlu. Viņa daiļradē darbi, kuros stilizācija un deformācija bija novesti līdz galējībai, mijās ar reālistiskiem poļu vēstures veidotāju portretiem. (1940.~gadā viņš tika ieslodzīts vācu koncentrācijas nometnē, bet izdzīvoja, ražīgi strādāja arī Polijas Tautas republikā)

Cits izcils skulptors \strong{E.~Vittigs}, jaunībā izbaudījis impresionisma un simbolisma ietekmi, vēlāk tomēr nonāca pie neoklasiskajiem principiem. Viens no viņa veiksmīgākajiem darbiem ir piemineklis bojā gājušajiem lidotājiem Varšavā. Tā projekts tika radīts jau 1923.~gadā, izstādīts tā modelis 1929.~gadā Nacionālajā izstādē Poznaņā. Pats piemineklis tika atklāts 1932.~gadā. Tas bija pirmais skulpturālais darbs Eiropā, veltīts bojā gājušajiem lidotājiem. Pēc Varšavas sacelšanās sakāves 1944.~gadā nacisti pieminekli iznīcināja. 1967.~gadā tas tika no jauna atklāts, nu jau jaunā vietā.

Avangardistiskas konstruktīvisma tendences parādījās K.~Kobro u.c. mākslinieku skulptūrās.

Sarežģīta bija \strong{arhitektūras} attīstība. 20.--30.~gados neoklasicisma un funkcionālisma tendences (angļu apdrošināšanas biedrības ``\entxti{Prudential}'' debesskrāpis Varšavā, Jagelloņu bibliotēka (\pltxti{Biblioteka Jagiellońska}) Krakovā) sadzīvoja ar modernistisku stilizāciju tautas celtniecības garā (Polijas paviljons Starptautiskajā mūsdienu dekoratīvās un rūpnieciskās mākslas izstādē (franču \frtxti{Exposition Internationale des Arts Décoratifs et Industriels Modernes}) Parīzē 1925.~gadā. Celtniecības centrs 20.--30.~gados bija Varšava. Šeit tika uzceltas ievērojamas sabiedriskas ēkas. Tā, Č.~Przubiļskis projektēja Nacionālā teātra un Kara ministrijas, M.~Laļevičs~--- Ģeoloģiskā institūta un Lauksaimniecības bankas, R.~Šverčinskis~--- Sakaru ministrijas ēku. Jau 20.~gados Polijā plaši izplatījās konstruktīvisma stils. Tajā tika celti t.s. ``īres'' nami, arī savrupmājas. Daži arhitekti, kā J.~Košus-Vitkevičs u.c. centās izkopt nacionālo stilu.

Tika veikti mēģinājumi regulēt celtniecību pilsētā, radīt racionālu tās rajonu plānojumu, sākās masu celtniecība funkcionālisma garā. Taču šādu projektu realizācijas iespējas, galvenokārt kapitāla trūkuma dēļ, bija ierobežotas.

Starpkaru periodā Polijā salīdzinājumā ar tās pakļautības gadiem uzplauka \strong{zinātniskā darbība}, īpaši matemātikas un dabaszinātņu, pirmkārt ķīmijas, jomā. Galvenie zinātnes centri bija universitātes Varšavā un Krakovā. Dibinājās jauni centri~--- universitātes Poznaņā, Viļņā, Ļvovā kā arī Katoliskā universitāte Ļubļinā un Brīvā poļu universitāte (\pltxti{Wolna Wszechnica Polska}) Varšavā ar filiāli Lodzā, Kalnu-metalurģiskā akadēmija (\pltxti{Akademia Górniczo-Hutnicza}) Krakovā. Padziļinājās pētījumu specializācija. Polijas Zinību akadēmija (\pltxti{Polska Akademia Umiejętności}), kam bija zinātniskas biedrības statuss, pārstāvēja poļu zinātniekus aiz robežas.

Daļa poļu tautības zinātnes darbinieku atgriezās dzimtenē no agrāk to pakļāvušo valstu pilsētām. Īpaši daudz repatriantu~--- poļu inteliģentu atgriezās no boļševiku pārvaldītās Krievijas, uzsāka darboties atjaunotās Polijas zinātnē un izglītības sistēmā. Piemēram, universitātēs katedras vadīja filologi J.~Boduēns de Kurtenē, T.~Zeļinskis, filozofs un sociologs L.~Petržickis.

Stāvoklis zinātnē sarežģījās ekonomisko krīžu gados, īpaši grūti tai klājās vispasaules krīzes laikā. Nacionālās kultūras fonda (\pltxti{Fundusz Kultury Narodowej}), kura uzdevums bija veicināt zinātnes un mākslas attīstību, izdevumi samazinājās no 2,5~miljoniem zlotu 1928./1929.~gadā līdz 1~089~000~1932./1933.~gadā. Ap 600~pastāvošajām zinātniskajām laboratorijām un 27~augstākajām mācību iestādēm tad drastiski trūka finansējuma. Daudzi izcili zinātnieki meklēja pielietojumu saviem spēkiem ārzemēs.

Polijā pārsvarā attīstījās tās nozares, kurām nevajadzēja daudz līdzekļu.

\strong{Matemātika} starpkaru periodā bija tā zinātnes nozare, kurā poļu zinātnieki guva visizcilākos panākumus. Izveidojas plaša funkciju teorijas skola. Kopš 1920.~gada Varšavā sāka iznākt starptautisks žurnāls ``\latxti{Fundamenta Mathematicae}'' (Fundamentālā Matemātika). Līdz 1939.~gadam iznāca 32~tā sējumi. 1931.~gadā poļu matemātiķi izdeva veselu sēriju vispārinošu darbu ar kopējo nosaukumu ``\pltxti{Monografie matematyczne}'' (Monogrāfijas matemātikā). Poļu matemātiķi izdeva savus zinātniskos žurnālus (``\latxti{Studia Mathematika}'' u.c.), publicējās starptautiskos izdevumos, pasniedza ārvalstu universitātēs, piedalījās starptautiskos matemātiķu kongresos.

Pazīstams pasaulē kļuva viena no funkcionālās analīzes pamatlicējiem S.~Banaha vārds. Kopā ar viņu Ļvovā radoši strādāja H.~Šteinhauss, S.~Mazurs, J.~Šauders, V.~Orličs. Varšavā pētījumus kopu teorijas un topoloģijas laukā veica V.~Serpiņskis, S.~Janiševskis, S.~Mazurkevičs, K.~Kuratovskis, S.~Sakss, K.~Borsuks u.c.

Matemātikas panākumi veicināja loģikas attīstību. Radās t.s. Varšavas--Ļvovas matemātiskās loģikas skola. Izcilākie tās pārstāvji bija J.~Lukasevičs, S.~Lesņevskis, A.~Tarskis.

\strong{Fizikas} pētījumu galvenie centri, galvenokārt optikā un teorētiskajā fizikā atradās Krakovas un Ļvovas universitātēs. Drīz pazīstama kļuva eksperimentālās fizikas nodaļa Varšavas universitātē. Šeit notika pētījumi atomārajā un molekulārajā optikā, kodolfizikā u.c. M.~Volfke 1920.~gadā izteica hologrāfijas ideju. S.~Peņkovskis pētīja luminiscenci, A.~Jablonskis~--- molekulāro optiku, Č.~Belobžeskis~--- teorētisko fiziku un astrofiziku, V.~Rubinovičs~--- kvantu teoriju.

1920.~gadā izveidojās Polijas fizikas biedrība. Starpkaru periodā Varšavā notika divi starptautiski fiziķu kongresi.

Krakovas un Varšavas universitātēs, Varšavas Politehnikumā un Krakovas Kalnu akadēmijā sekmīgi attīstījās pētījumi \strong{ķīmijas} jomā. Ļvovas Politehniskajā institūtā panākumus guva vēlākais Polijas prezidents I.~Moscickis. Viņš, piemēram, izstrādāja metodes slāpekļskābes iegūšanai no gaisa, gazolīna iegūšanai no dabas gāzes u.c.

V.~Svetoslavskis pētīja termoķīmiju, M.~Centeršvērs~--- ķīmisko kinētiku, K.~Dzevonskis~--- krāsvielas, B.~Šiškovskis~--- šķīdumu ārējo spriedzi.

Tiesa, daudzi ievērojami poļu ķīmiķi turpināja strādāt ārzemēs. Pirmkārt jānorāda uz pasaulslaveno fiziķi un ķīmiķi divkārtējo Nobela prēmijas laureāti \strong{M.~Sklodovsku-Kirī}.

Polijā tika veikti pētījumi vispārējā ģeogrāfijā, kartogrāfijā, ģeomorfoloģijā un klimataloģijā. Ģeoloģiski tika pētīts Karpatu kalnu rajons, atklāts Piebugas ogļu baseins.

Bioloģijas jomā tika izstrādātas parazitoloģija, bioķīmija, botānika, zooloģija, hidrobioloģija, tika pētīti dabas aizsardzības jautājumi, attīstījās lauksaimniecības zinātne.

Mikrobioloģijas jomā darbojās viens no mācības par asins grupām izstrādātājs bakteriologs un imunologs L.~Hiršfelds un R.~Veigels, kurš atklāja vakcīnu pret tīfu.

Tika organizēti zinātniski centri terapijas, pediatrijas, ķirurģijas, neiroloģijas, psihiatrijas jomā.

\strong{Tehnisko zinātņu} attīstībai Polijā starpkaru periods nebija visai labvēlīgs, jo valsts rūpniecības attīstības līmeņa ziņā un tehniskās nodrošināšanas jomā neatradās sava laika pirmrindas sasniegumu līmenī. Tomēr zinātnieki centās palīdzēt valsts attīstībai. 1920.~gadā tika dibināta Tehnisko zinātņu akadēmija Varšavā, 1927.~gadā tika organizēts aerodinamikas institūts, 1928.~gadā~--- radiotehniskais un ķīmiskais, 1932.~gadā~--- radija institūts.

Darbojās vairāki ievērojami zinātnieki. S.~Brila pētīja metināšanas teoriju, uzbūvēja pirmo Polijā pilnībā no dzelzs konstrukcijām sametināto tiltu, piedalījās pirmo debesskrāpju celtniecībā. J.~Groškovskis darbojās radiotehnikas jomā un deva ieguldījumu radara darbības principu un uzbūves izstrādē.

Pozitīva tehnisko zinātņu attīstību ietekmēja valstij piederošo rūpniecības uzņēmumu darbība, īpaši nozarēs, saistītās ar ieroču ražošanu.

Ievērojami sasniegumi tika gūti aviācijas attīstībā. 1929.~gada 16.~oktobrī lidotājs F.~Žvirka ar lidmašīnu RWD-2 uzstādīja pasaules rekordu, sasniedzot 4~004~m augstumu. Šai gadā viņš un S.~Vigura Eiropā nolidoja 5~000~km, sasniedza daudzas sportiskas uzvaras, bet 1932.~gadā gāja bojā aviokatastrofā. 1930.~gada 22.~decembrī Starptautiskajā aviācijas salonā Parīzē poļu lidmašīna PZL P-6 tika atzīta par labāko iznīcinātāju pasaulē. 1934.~gadā B.~Orļinskis ar lidmašīnu PZL P.24 uzstādīja pasaules rekordu starp iznīcinātājiem, sasniedzot ātrumu 416~km/h. 1938.~gadā aviācijas izstādē Belgradā pirmo vietu, apsteidzot angļu, vācu un itāļu lidmašīnas, ieguva poļu bumbvedējs ``\pltxti{Łoś}'' (Alnis~--- pilns nosaukums ``\pltxti{PZL.37 Łoś}'')

20.--30.~gadi kļuva par \strong{filozofiskās domas} uzplaukuma laiku Polijā. Viens no poļu ietekmīgākajiem domātājiem XX gadsimtā bija \strong{T.~Kotarbiņskis}, daudzu darbu autors ētikas, izziņas teorijas, loģikas jomās, daudzus gadus pētīja ``prakseoloģiju''~--- darba organizācijas jautājumus, ir daudzu spārnotu aforismu autors. Viņš centās neopozitīvisma principus interpretēt materiālisma garā. Pie tam viņš uzskatīja materiālismu tikai par vienu no varbūtējākām mācībām, kura pielietojama zinātnieku vienošanās rezultātā. 30.~gados viņš aktīvi piedalījās poļu sabiedrības dzīvē, vēršoties pret nacionālismu, klerikālismu un antisemītismu. Viņam ir lieli nopelni antireliģiskās apgaismības ētikas propagandā, cīņā pret ``sanācijas'' ideoloģijas uzspiešanu augstākajā izglītībā. Jau pēc Otrā pasaules kara T.~Kotarbiņskis kļuva par Polijas Zinātņu akadēmijas prezidentu.

K.~Tvardovskis radīja t.s. \strong{Ļvovas-Varšavas skolu}~--- ievērojamāko poļu filozofijas vēsturē. Veica pētījumus metodoloģijā, eksperimentālajā psiholoģijā u.c. Pie minētās skolas piederēja arī V.~Tatarkevičs, kurš galvenokārt darbojās ētikas un estētikas jomās. Savos darbos izstrādāja mūzikas teorijas, dzejas un vizuālās mākslas jautājumus.

Otrs ievērojams virziens poļu filozofijā bija \strong{tomisms}, kas balstījās uz Akvīnas Toma mācību. Doktrīnas centrā ir ne tik daudz ticības dogmāti, cik mācība par to, kā tos saprast ar prātu. Tā centri bija Krakovā un Katoliskajā universitātē Ļubļinā. Šī virziena pārstāvji veica pētījumus viduslaiku filozofijas un reliģijas vēsturē, tika izstrādātas izziņas, ētikas un estētikas problēmas. Katoliskā filozofija tika izmantota cīņā pret materiālistisku un revolucionāru ideju izplatību.

Atzīts pasaulē zinātnieks bija sociologs F.~Znaņeckis, kurš 1920.~gadā atgriezās Polijā no ASV un vadīja socioloģijas fakultāti Poznaņas universitātē. Tur viņš arī organizēja Poļu socioloģijas institūtu (\pltxti{Polski Instytut Socjologii}), sāka izdot žurnālu ``\pltxti{Polski Przegląd Socjologiczny}'' (Poļu Socioloģijas Vēstnesis). F.~Znaņeckis uzturēja ciešus kontaktus ar amerikāņu kolēģiem. Ar kara sākumu pārcēlās uz ASV.

Nacionālas valsts izveide pavēra iespējas uzlabot pētījumu organizāciju \strong{valodniecības} jomā. Tika atvērtas lingvistiskas katedras universitātēs, pieauga lingvistisko izdevumu skaits. Pēc 1919.~gada valstī atgriezās virkne ievērojamu poļu lingvistu. Starp tiem ievērojamākais bija jau minētais J.~Boduēns de Kurtenē. Viņš strādāja Krakovā, kura kļuva par Polijas valodpētniecības centru. J.~Boduēns de Kurtenē lika zinātniskus pamatus poļu valodas fonētikai un fonoloģijai.

Polijas neatkarības atjaunošana radīja labvēlīgus apstākļus \strong{vēstures zinātnes} attīstībai. 1925.~gadā vispārnacionālu raksturu ieguva Poļu vēstures biedrības (\pltxti{Polskie Towarzystwo Historyczne}, radās jau 1886.~gadā Ļvovā, 1913.~gadā izveidoja filiāli Krakovā), kura koordinēja pētījumus vēstures jomā. 1925.--1933.~gadā tajā darbojās ap 1~300 biedru. Valstī atgriezās plašas arhīvu materiālu kolekcijas no PSRS, Vācijas, Austrijas, Šveices. Uzsāka darbību valsts arhīvi un bibliotēkas. Pieejami kļuva daudzi privāti arhīvi. Universitātēs koncentrējās lielākā daļa profesionālo vēsturnieku. Ar pirmajiem Polijas neatkarības gadiem Polijas Zinību akadēmija Krakovā uzsāka dokumentu krājumu izdošanu par viduslaiku vēsturi un XIX~gadsimta vēsturi.

Līdzās jau agrāk iznākošajiem, tika izdoti arī jauni vēsturei veltīti periodiski izdevumi. Vēsturiskie žurnāli ``\pltxti{Kwartalnik Historyczny}'' (Vēsturiskais [gada] Ceturksnis) un ``Przegląd Historyczny'' (Vēsturiskais Apskats) kļuva par Poļu vēstures biedrības centrālajiem izdevumiem. Sāka iznākt specializēti žurnāli un krājumi ``\pltxti{Niepodleglosc}'' (Neatkarība) ``\pltxti{Wojskowy Przeglcjd Historyczny}'' (Kara Vēstures Apskats), ``\pltxti{Kronika Ruchu Rewolucyjnego}'' (Revolucionārās Kustības Hronika), ``\pltxti{Roczniki Dziejow Spolecznych i Gospodarczych}'' (Sociālās un ekonomiskās vēstures gadagrāmata) u.c. Paplašinājās pētījumu tematika. Līdzās ar iepriekšējam periodam raksturīgajām tēmām plašāk tika pētīta jauno un jaunāko laiku vēsture.

Atjaunojās poļu vēsturnieku kongresu sasaukšana (1925.~gadā Poznaņā notika IV, 1930.~gadā Varšavā~--- V, 1935.~gadā Viļņā~--- VI~kongress.). Liels notikums poļu vēsturnieku dzīvē bija VII~Starptautiskā vēsturnieku kongresa sanākšana Varšavā 1933.~gadā. Ar ievērojamiem panākumiem varēja lepoties poļu arheologi un etnogrāfi. Tika pētītas gan ekonomikas, gan sociālā, gan politiskā un kultūras vēsture.

Līdz ar politiskajām pārmaiņām mainījās arī vēsturnieku uzskati. Tā J.~Kostrževskis, kurš kā vācu arheologa G.~Kosinnas skolnieks pirms Pirmā pasaules kara divsējumu darbā arheoloģiskos atradumus Austrumvācijā un Rietumpolijā I~gadsimtā pirms mūsu ēras pierakstīja austrumģermāņiem, pēc Polijas valsts nodibināšanas atzina, ka t.s. lužiciešu kultūra (vācu \detxti{lausitzer Kultur}, poļu \pltxti{kultura łużycka}) I~gadu tūkstotī pr. m. ēr.~--- I g. t. m. ēr.) esot bijusi poļu kultūras daļa un poļi tad apdzīvojuši plašus vēlākās Vācijas apgabalus.

Vēstures metodoloģiskajā bāzē starpkaru periodā pārsvarā bija neopozitīvisma un neoromantisma tendences. Aktīvi tika pētīta viduslaiku vēsture. Vēsturnieku vidū norisa diskusijas jautājumā vai Polijā pastāvēja feodālisms? Zinātnieki, kas to identificēja ar lēņu sistēmu, uz šo jautājumu atbildēja noraidoši. Piemēram, O.~Balcers apgalvoja, ka Polijā nebija lēņu sistēmas, kas neļāva nodibināties tautas saiknei ar monarhu, un tāpēc Polija bijusi vairāk saliedēta kā Rietumu valstis. Z.~Voicehovskis uzskatīja, ka tai laikā, kad Rietumos karaļa kalpotāji kļuva par viņa lēņa vīriem, Polijā feodālā sadrumstalotība attīstījusies, dalot valsti starp valdnieka dinastijas locekļiem. Tāpēc notikusi tieša pāreja no dzimtu iekārtas uz kārtu valstiskumu.

Turpretī K.~Tumarņickis uzskatīja, ka feodālisma ietekme Polijā galvenokārt izpaudusies uz zemākajiem slāņiem: pilsētniekiem un zemniekiem, bet šļahta neļāvusi pilnā mērā feodalizēt poļu sabiedrību.

Daudzi vēsturnieki savus spēkus veltīja ``Jagelloņu idejas'' propagandai, kas reizē bija mēģinājums izplatīt vienotu orientieri gan ārpolitikā, gan pagātnes pētniecībā. Vesela paaudze jau pirms kara savu darbību uzsākušo Krakovas skolas vēsturnieku, kā O.~Haļeckis, V.~Konopčinskis, S.~Kutšeba u.c. mēģināja pamatot Polijas attiecības ar kaimiņiem~--- Baltkrieviju, Lietuvu, Ukrainu ar vienotu pagātni un vēsturisko atmiņu. Tai pat laikā ``Jagelloņu ideja'' kalpoja Polijas ``vēsturisko tiesību'' uz rietumukraiņu un rietumbaltkrievu zemēm aizstāvībai, poļu tautas audzināšanai naidā pret citām, īpaši krievu, vācu un ebreju tautām.

S.~Kutšeba turpināja vēsturiski-tiesiskos pētījumus. Viņa spalvai piederēja viens no pirmajiem darbiem par Polijas valstiskuma atjaunošanu, kurā, neraugoties uz atsevišķu politiķu un partiju kritiku, tika idealizēta pilsonisko partiju, bet mazināta masu loma šai procesā.

Saglabājās tradicionālā poļu vēsturnieku interese par politisko vēsturi. Ievērojamāko skolu XIX~gadsimta Polijas vēstures pētniecībā radīja \strong{M.~Handelsmans}. Viņš savu pētniecības metodi sauca par ``sociopsiholoģisku''. Viņš bija pārliecināts, ka vēstures procesā nevar runāt par likumu darbību, bet gan tikai par attīstības tendencēm. Vēsturiskos vērtējumus viņš uzskatīja par relatīviem, to patiesuma pakāpi saistīja ar zinātnes attīstības līmeni un izziņas metožu pilnveidi. Kopumā būdams subjektīvs ideālists, dažos savos darbos viņš tiecās atzīt vēstures materiālistisku izpratni, noliedza neoromantisku, apoloģētisku nacionālās vēstures interpretāciju.

Dažādi tika vērtētas XIX~gadsimtā darbojošās vēsturiskās personības. Piemēram, Ā.~Skalkovkis kritizēja T.~Kostjuško darbību, toties atzinīgi izteicās par marķīzu A.~Veļepoļski. Turpretī M.~Handelsmans poļu tautas ``nacionālās gudrības'' iemiesojumu saskatīja kņazā Ā.~Čartorijskī.

Starpkaru periodā aktivizējās arī ekonomiskās vēstures problēmu izstrāde. Izveidojās divi to izpētes centri. Vienu~--- Ļvovā vadīja F.~Bujaks. Viņš centās radīt vienotu Polijas sociāli ekonomiskās attīstības ainu, pētīja Žečpospolitas ekonomiskā pagrimuma cēloņus, Galīcijai aktuālās emigrācijas, zemnieku stāvokļa, rūpalu un tirdzniecības attīstības problēmas. Otrs centrs izveidojās Poznaņa, kur ar saviem skolniekiem strādāja J.~Rutkovskis. Viņš pirmais no nemarksistiskajiem vēsturniekiem uzstājās ar feodālās iekārtas ekonomiskās vēstures koncepciju, pārejot uz sabiedriski-ekonomisko formāciju atzīšanas pozīcijām. Jau pirms Otrā pasaules kara J.~Rutkovskis sagatavoja plašu ``Polijas ekonomisko vēsturi'' (``\pltxti{Historia gospodarcza Polsci}''), kura tika izdota jau pēc kara.

% page 605


Jauns virziens bija kara vēstures pētījumi. No pozitīvā objektīvisma pozīcijām to aprakstīja V.~Tokaržs, īpašu uzmanību veltījot 1794.~gada un 1830.--1831.~gada sacelšanos izpētei. Viņš rūpīgi rekonstruēja faktus, nonākot pie secinājuma, ka sacelšanos noteica kā iekšējie, tā ārējie faktori, ka cīņa par neatkarību bija saistīta ar sabiedrisko pārmaiņu nepieciešamību.

Tika veikti pētījumi kultūras vēstures jomā. Pasaules atpazīstamību ieguva A.~Brukners. Viņš sagatavoja ievērojami pilnīgāku, nekā priekšgājēji ``Poļu kultūras vēsturi'' (``\pltxti{Dzieje kultury polskiej}'') parādot tās saistību ar lietuviešu, krievu, čehu u.c. kultūrām.

Ne visi vēsturnieku vērtējumi bija tikai pašu autoru pārdomu, simpātiju un antipātiju rezultāts. Lielā mērā tie bija atkarīgi no valstī valdošās vispārējās atmosfēras. Sevišķi tas bija redzams pētījumos par jaunāko laiku vēsturi. J.~Pilsudskis, viņa līdzgaitnieki~--- leģionāri tika atzīti par tiešiem agrāko poļu tautas nacionālās atbrīvošanās tradīciju mantiniekiem, reizē ``neredzot'' vai ``aizmirstot'' šo vēsturisko personāžu daudzos kļūdainos, pat noziedzīgos soļus.

Pēc 1926.~gada valsts apvērsuma pastiprinājās vēsturnieku atiešana no objektīvisma pamatiem. Radās grupa vēsturnieku, cieši saistītu ar ``sanācijas'' režīmu. Tās kredo formulēja pats J.~Pilsudskis rakstā ``\pltxti{Poprawki do historii}'' (``Labojumi vēsturei''). Vēsturiskā procesa centrā tika izvirzīta ``stipra personība'', faktiski saknēja aicinājums pārrakstīt vēsturi. Tāda pieeja bija tuva neoromantiķiem. Spilgts šādas pozīcijas pārstāvis bija S.~Zakševskis, kurš kā aktīvs J.~Pilsudska sekotājs slavināja stipru varu, vadoņu lomu vēsturē. Pieauga šovinistiskās tendences, īpaši atspoguļojot nepoļu tautu vēsturi šļahtas pārvaldītajā Žečpospolitā.

1930.~gadā tika radīts Austrumeiropas pētniecības institūts (\pltxti{Instytut Naukowo-Badawczy Europy Wschodniej}), kurš to pētīja caur poļu lielvalstisko interešu prizmu.

Īpaši ``sanācijas'' historiogrāfija koncentrēja uzmanību karu vēstures pētniecībai. Sevišķi aktīvi tika rakstīts par Pirmā pasaules kara un poļu-padomju kara notikumiem. Virknē darbu tika absolutizēta J.~Pilsudska un viņa radīto leģionu loma valsts neatkarības atjaunošanā. Zināmu vietu šai historiogrāfijā ieņēma XIX gadsimta poļu sacelšanos pētniecība. To sociālie aspekti gan tika apieti, uzsvars likts uz sacelšanos pretkrievisko virzību. Sacelšanos sakāvju cēloņi tika saskatīti ``vadoņu'', ``spēcīgu personību'' trūkumā. Pētot viduslaiku sižetus šādi vēsturnieki izskaistināja karaļu-iekarotāju lomu, slavināja Žečpospolitas agresīvu austrumu politiku, uzsvēra lielvalstiskuma tradīcijas. Tā Austrumeiropas pētniecības institūta ģenerālsekretārs O.~Gurka VI~poļu vēsturnieku kongresā (1935) slavināja Krakovas vēsturnieku skolas sasniegumus, izcēla to, ka tā par Žečpospolitas bojā ejas cēloni uzskatīja stipras varas trūkumu. Toties tos vēsturniekus, kuri neatzina ``sanācijai'' tīkamās koncepcijas, viņš apvainoja ``nacionālā bezatbildībā''. (Tiesa, 1945.~gadā O.~Gurka atgriezās Polijā, darbojās kā diplomāts un mācību spēks, acīmredzami atzīdams komunistu varu par pietiekami ``stipru'' un ``nacionāli atbildīgu''.)

Cits šīs grupas pārstāvis V.~Pobogs-Maļinovskis pētīja endeku darbību līdz 1918.~gadam. Parādot to politiskās programmas trūkumus, viņš centās izcelt J.~Pilsudska sekotāju nopelnus. Poļu leģionu vēsturi, īpašai akcentējot J.~Pilsudska lomu, pētīja V.~Lipinskis.

Oficiālajai pilsoniskajai vēstures zinātnei pretojās, savu vēstures redzējumu centās izplatīt marksistiski noskaņotie sabiedriskie un zinātnes darbinieki. Daudzu vēsturisku darbu autors bija PPS kreisā spārna darbinieks Ā.~Prohņiks. Ar pseidonīmu H.~Svoboda (\pltxti{Henryk Swoboda}) viņš publicēja grāmatu ``\pltxti{Pierwsze piętnastolecie Polski niepodległej}'' (Neatkarīgās Polijas pirmie piecpadsmit gadi), kurā parlamentāro periodu pretstatīja J.~Pilsudska diktatūrai, analizēja demokrātisko spēku kļūdas pirms 1926.~gada apvērsuma.

Sociālo kustību vēsturi pētīja sociologs, folklorists un kultūras vēsturnieks S.~Čarnovskis. Viņš pētīja reformācijas un kontrreformācijas sociālos aspektus, darbā ``\pltxti{Kultura}'' (1938) viņš mēģināja izdalīt sociālos faktorus Polijas kultūras attīstībā. Viņš arī izteica domas par pretrunu dialektiku, par sabiedrības polarizāciju kā sociālo attīstību virzošu spēku.

Materiālistiska vēstures izpratne bija likta pamatā poļu komunistu vēsturiskajiem darbiem. Vairākus darbus par Polijas vēsturi uzrakstīja J.~Marhļevskis, kurš šai laikā dzīvoja Padomju Krievijā/Savienībā. Viņš skatīja revolucionārās kustības attīstību un perspektīvas Polijā, cenšoties argumentēti pierādīt nepieciešamību revolucionārā ceļā atrisināt valstī pastāvošās sociālās pretrunas. Neraugoties uz zināmu procesu vienkāršošanu, viņš demonstrēja sabiedriskās attīstības galveno faktoru izpratni.

Arī citi poļu komunisti, dzīvojoši Padomju Savienībā, pētīja revolucionārās un strādnieku kustības vēsturi Polijā, vāca un publicēja dokumentus par to. 1926.~gadā Maskavā poļu valodā sāka iznākt žurnāls ``\pltxti{Z Pola Walki}'' (No Cīņas Lauka).

Radošu pieeju demonstrēja viens no Komunistiskās partijas vadītājiem J.~Bruns. Viņš izsekoja poļu nacionālās atbrīvošanās kustības attīstībai XIX un XX~gadsimtā, parādīja tā šķirisko saturu, kritizēja buržuāziskās spekulācijas nacionālajā jautājumā, rādīja strādnieku šķiras sociālo un nacionālo prasību attiecības, analizēja PKunLSD galveno ideologu uzskatus nacionālajā jautājumā. 1925.~gadā J.~Bruns Polijā tika ieslodzīts cietumā, apmaiņas kārtībā izsūtīts uz PSRS, dzīvoja Rietumos, 1941.~gadā atgriezās PSRS. Līdz šai dienai Varšavā viņa vārdā ir nosaukta iela.

Marksistu nopelns ir tādu jautājumu (pirmkārt sabiedrisko kustību un strādnieku cīņas vēsturē) izvirzīšana, kurus nevēlējās atspoguļot pilsoniskā historiogrāfija.

Abās starpkaru desmitgadēs sekmīgi attīstījās arī \strong{poļu inteliģences starptautiskie}, oficiālie un neoficiālie, \strong{kultūras sakari} ar demokrātisko valstu valstiskajām un sabiedriskajām organizācijām, zinātnes, literatūras, mākslas darbiniekiem. Tiesa, sarežģītāk veidojās poļu un tādu kaimiņvalstu kā Vācija, Lietuva, PSRS inteliģences sakari. Tiem traucēja valstu politiskās nesaskaņas. Īpaši nelabvēlīgi politiskie apstākļi ietekmēja Polijas un Padomju Savienības kultūras darbinieku attiecības Tā, tikai daži izcili zinātnieki un kultūras darbinieki: literāts un literatūrzinātnieks M.~Zdzehovskis, mākslas zinātnieks M.~Vaļickis saglabāja uzticību draugiem~--- ``moskaļiem'' un nekad arī nelika vienlīdzības zīmi starp pirmsrevolūcijas un staļinisko Krieviju.

Neseno poļu-padomju karu neaizmirsa neviena puse. Turklāt Polijas valdošās aprindas tēloja Padomju režīmu kā cariskās Krievijas impērisko tieksmju mantinieku, kurām pievienoti vēl pasaules revolūcijas realizācijas utopiskie centieni. Arī Padomju Savienība daudz ko mantoja no cariskās Krievijas, tradicionālajam negatīvajam poļa stereotipam pievienojās vēl apzīmējumi ``baltpolis'', ``pans''. Šiem buržujiem, kontrrevolucionāriem, muižniekiem tika pretstatīti darba tautas pārstāvji. Vajadzēja vairākus gadus, lai nodibinātos abu tikko savā starpā karojušo valstu kultūras darbinieku kontakti.

Tikai no 1927.~gada sākās padomju cilvēku: žurnālistu, izglītības darbinieku u.c. grupu braucieni uz Poliju. Pirmais padomju literāts, kurš apmeklēja Poliju, bija V.~Majakovskis. 1927.~gadā notika arī padomju rakstnieku L.~Seifuļinas un I.~Ērenburga vizītes Polijā. I.~Ērenburgs Polijā ieraudzīja vienīgi antisemītismu, militāristu kundzību, nacionālismu un provinciālismu literatūrā, krievu kultūras un valodas diskrimināciju. viņš secināja, ka ``lielummānija skārusi ne cilvēka~--- tautas smadzenes''. Arī citi padomju rakstnieki, XX gadu beigās~--- 30.~gadu sākumā pēc Polijas apmeklējuma to tēloja pēc viena šablona~--- Polija ir naidīga PSRS valsts, kurā aiz ārējā spožuma, militāristu zobenu un piešu žvadzoņas, modernām kafejnīcām, dāmu krāšņajiem tērpiem slēpjas darbaļaužu nabadzība.

Jau 20.~gados uz Polijas ekrāniem parādījās Padomju Savienībā ražotās filmas. Līdz Otrā pasaules kara sākumam poļu skatītāji varēja noskatīties apmēram 70~padomju kinolentes. Taču cenzūras (daudzas filmas tika demonstrētas stipri ``apgraizītas'') un policijas darbības (gadījās, ka padomju filmu izrādēs to skatītāji tika arī arestēti) rezultātā plašam skatītāju lokam to saturs tika atklāts nepilnīgi vai arī bija vispār nepieejams. 30.~gadu pirmajā pusē notika arī atsevišķu poļu mākslinieku, literātu, arhitektu un citu kultūras darbinieku grupu braucieni uz PSRS.

Smagas problēmas starpkaru perioda Polijā pastāvēja \strong{tautas izglītības} laukā. 1921.~gada 17.~marta Konstitūcija gan nostiprināja jau agrāk pasludināto vispārējo obligāto izglītības pienākumu, pirmajā pēckara desmitgadē notika pavērsiens, taču stāvoklis radikāli nemainījās. 30.~gadu vidū joprojām trūka skolotāju un skolu ēku. Skolotāju atalgojums samērā ar citām valstīm bija ļoti pieticīgs, tas teju pārsniedza iztikas minimumu.

Pēc oficiālajiem datiem 1921.~gadā 33\% Polijas iedzīvotāju vecumā no 10~gadiem bija analfabēti. 1931.~gadā 23\% no iedzīvotājiem, vecākiem par 10~gadiem, joprojām neprata lasīt un rakstīt. Faktiski analfabētisma līmenis bija vēl augstāks.

Īpaši smags analfabētisms bija nacionālo mazākumu iedzīvotāju~--- ukraiņu, baltkrievu, lietuviešu vidū. Četrās~--- galvenokārt baltkrievu un ukraiņu apdzīvotajās Viļņas, Novogrudokas, Poļesjes un Volīnijas vojevodistēs 1921.~gadā analfabētu bija 64\%, pārsvarā ukraiņu apdzīvotajā Staņislavas vojevodistē~--- 46\%, bet agrākās Prūsijas teritorijā esošajās vojevodistēs~--- tikai 4,2\%. Pēc 1926.~gada maija apvērsuma ``kresos'' sākās ukraiņu un baltkrievu skolu masu likvidācijas kampaņa. 1931.~gadā austrumu vojevodistēs analfabētu skaits vēl joprojām sasniedza 41\%, bet Staņislavas vojevodistē~--- 36,6\%. 30.~gadu vidū pēc Polijas skolu statistikas datiem mācības 21~806 tautskolās noritēja poļu valodā, 768 --vācu. (Pēc citiem datiem pastāvēja ap 800~sabiedrisku un 180~privātu vācu skolu). 183~skolās mācīja ivritā, 177~--- jidišā. Tikai 8~skolās bērni mācījās krievu valodā.

Ebreju iedzīvotāju vidū analfabētus bija grūti atrast, taču tas nav jāuzskata par Polijas izglītības sistēmas, bet gan par ebreju diasporas nopelnu.

Oficiāli bija ieviesta obligāta 7-gadīga izglītība poļu (nevis dzimtajā) valodā, taču faktiski tā netika īstenota. 1921./1922.~gadā Polijā viena sākumskola bija uz 986~iedzīvotājiem, bet 1929./1930.~gadā~--- jau uz 1~211~iedzīvotājiem. Sakarā ar iedzīvotāju pieaugumu bija jāceļ tūkstošiem jaunu skolu ēku, taču tā vietā ap 1~000~skolu tika slēgtas.

Ekonomiskās krīzes gados, to bērnu skaits, kuri neapmeklēja skolas, pastāvīgi pieauga. 1928./1929.~mācību gadā skolas neaptvēra ap 5\% bērnu, 1929./1930.~gadā 8\%, 1930./1931.~gadā 19\%, bet 1931./1932.~gadā jau 25\%. Pēc poļu vecāku savienības datiem 1932.~gadā ap 1~200~000~bērnu neapmeklēja skolas (valdība atzina, ka to nedarīja 600~000~bērnu). No 1929.~līdz 1932.~gadam tika slēgtas 37~ģimnāzijas, vidējo mācību iestāžu audzēkņu skaits samazinājās par 40~tūkstošiem. 60\% vidusskolu bija privātas ar augstu mācību maksu, kuras varēja apmeklēt tikai turīgu vecāku bērni. Bezdarbnieku ģimenēs tikai 46\% skolas vecuma bērnu apmeklēja to.

Viens no aktīvākajiem skolotāju organizāciju dalībniekiem Č.~Vicehs savās atmiņās norādījis, ka 20.~gadu beigās, kad bērnu pieplūdums Polijas skolās strauji pieauga un tās izrādījās pārpildītas, varas iestādes neatrada citu izeju kā atteikties no vispārējās izglītības principa un pazemināt apmācību līmeni sākumskolā. Simti tūkstošu bērnu (vairāk nekā 1~miljons) zaudēja iespēju mācīties, bet tūkstoši skolotāju~--- darbu. Ja 1928./29.~mācību gadā ārpus skolas palika 3,6\% bērnu, tad 1936./37.~mācību gadā~--- jau 11,7\%.

``Sanācijas'' režīma izglītības ministra J.~Jendžejeviča 1932./1933.~gadā īstenotā izglītības reforma padarīja zemnieku jauniešiem vidējo izglītību gandrīz pilnībā nepieejamu. Pēc šīs reformas skolas dalījās divās pakāpēs. Pirmajā pakāpē bija jāmācās septiņi gadi, taču zināšanas tā deva četrgadīgās sākumskolas līmenī. Laukos formāli pastāvošā septiņgadīgā vispārizglītojošā skola faktiski tika reducēta uz divgadīgu mācību kursu. 76\% skolu laukos bija šīs pirmās pakāpes skolas. Otrās pakāpes skolu, kuras deva izglītību 7-gadīgās skolas līmenī, bija ne vairāk kā 8\% no visām. 1937.~gadā 77\% poļu skolu, kuras galvenokārt atradās laukos, nesniedza izglītību augstāku par 4~klasēm.

Lielākajās pilsētās saglabājās četrgadīgas ģimnāzijas, kurās varēja iestāties pēc sešu klašu pabeigšanas vispārējā skolā. Taču 30.~gados pat šo vispārizglītojošo skolu vecākajās klasēs tikai 13\% bērnu nāca no strādnieku un zemnieku ģimenēm. Augstskolās no šādām ģimenēm studēja tikai 8\% jauniešu.

Avīze ``\pltxti{Głos Nauczycielski}'' (Skolotāja Balss) 1933.~gada 18.~decembrī rakstīja: ``Ap miljonu bērnu neapmeklē skolu. Skolotāju skaits samazinājies par 4~tūkstošiem. Šaurās, aukstās un mitrās skolu telpās risinās īsts karš. Par šī kara upuriem ik gadus kļūst simti jaunu cilvēku dzīvību. Plosās tuberkuloze, bet skolotāji rūpīgi slēpj savu slimību, lai nezaudētu maizes kumosu, vai, labākajā gadījumā, netiktu atvaļināti ar 40~zlotu lielu pensiju''.

Premjerministrs (1934--1935) L.~Kozlovskis vienkārši paziņoja, ka Polija neies bojā izglītības trūkuma dēļ. Lai uzlabotu stāvokli, varas iestādes aicināja upurēties \dots{} skolotājus. Visās skolās sāka izmantot bezalgas darbu~--- kandidātu uz skolotāja amatu neapmaksātu ``ražošanas'' praksi. Skolās, kur bija vairāki skolotāji, vienā klasē vidēji mācījās 60~līdz 70~skolēnu, bet skolās ar vienu skolotāju~--- 120 un vairāk bērnu. Pēc statistikas Polijā uz 1~skolotāju bija 63~skolnieki, kas divas reizes pārsniedza noteiktās normas. ``Sanācijas'' režīma laikā analfabētu skaits gandrīz nemazinājās. 1938.~gadā tādu bija 20\% visu iedzīvotāju. Vairāk nekā 5~miljoni cilvēku neprata ne lasīt, ne rakstīt, ap 1~miljonu cilvēku prata tikai lasīt, bet ne rakstīt.

Starpkaru perioda beigās gan visi poļu, bet ne visi ukraiņu, baltkrievu bērni mācījās skolās. Polijā skolu nacionālajiem mazākumiem ar mācībām dzimtajā valodā 1937.~gadā bija mazāk nekā tais pašās teritorijās 1913.~gadā.

Demokrātiski noskaņotie sabiedrības slāņi cīnījās par izglītības sistēmas demokratizāciju. 1932.~gadā kreisi noskaņotie izglītības darbinieki radīja savu biedrību ``\pltxti{Nowe drogi}'' (Jauni ceļi). 1937.~gada martā Poļu skolotāju savienības Ārkārtējā kongresā kreisie elementi guva pārsvaru savienības Galvenās valdes vēlēšanās. Valdība ievēlēto valdi atbīdīja no savienības vadības, ieceļot tajā savu komisāru. Atbildot uz to sākās skolotāju streiks. Tajā ievērojama loma bija jaunai skolotājai un iesācējai rakstniecei V.~Vasiļevskai.

Kopumā starpkaru periodā Polijā darbojās 16~valsts un 11~privātās augstākās mācību iestādes, kuras beidza 83~tūkstoši cilvēku. (30.~gadu vidū to skaits Polijā, pastāvēja 24~augstskolas: 13~valsts un 11~privātas. Valsts universitātes pastāvēja Krakovā, Viļņā, Varšavā, Ļvovā, Poznaņā; 2~tehniskas augstskolas~--- Varšavā un Ļvovā. Varšavā darbojās arī Lauksaimniecības augstskola, Mākslas akadēmija, Valsts zobārstniecības institūts, Politikas augstskola, Žurnālistikas un avīžu izdevniecības augstskola; Ļvovā~--- Veterinārais institūts. Krakovā~--- Kalnrūpniecības akadēmija. Privātas bija Varšavā: Brīvā poļu universitāte, Tirdzniecības augstskola; Ļubļinā Katoliskā universitāte, Ļvovā Arējās tirdzniecības akadēmija, Viļņā Austrumeiropas institūts.)

Pieauga augstāko mācību iestāžu studentu skaits. Tiesa, galvenokārt auga poļu skaits, lielā mērā uz citu nāciju pārstāvju rēķina. Samērā nelielais augstskolu beidzēju kontingents ar arvien lielākām grūtībām varēja atrast darbu specialitātē. Universitāšu diplomu īpašnieki, kuri to iegūšanai bija veltījuši lielas pūles, nereti bija studējuši, cīnoties ar trūkumu, ieguvuši tuberkulozi, gadiem ilgi nīka bezmaksas pedagoģiskajā praksē vai stažējoties pie advokātiem cerībā nākotnē iegūt īstu darbu, amatu, darbības lauku.

Nenākas pārspīlēt poļu, īpaši zemnieku, \strong{masu izglītības un kultūras} līmeni. Tā, francūzis grāfs O.~d’Ethegons, kurš dienesta uzdevumā divdesmitajos gados apmeklēja Poliju, uzrakstīja grāmatu ``\frtxti{Pologne, Pologne\dots{}}'', (Polija, Polija\dots{}, 1925). Līdz tam autors par poļiem sprieda pēc to emigrantiem Parīzē. Saskaroties ar dzīves īstenību pašā Polijā, viņš bija dziļi vīlies. Par poļu attiecībām ar baznīcu grāfs stāstīja: ``Kad vairākumam zemnieku pajautā, vai viņi pieder Romas katoļu vai grieķu katoļu baznīcai, parastā atbilde ir: ``Nē, mums ir poļu reliģija''. Vārds ``katolicisms'' zemniekiem neko neizsaka. Valoda un reliģija poļu vairākumam ir vienīgās nacionālās iezīmes. Joprojām darbojas pamatlikums: ``Kas ir katolisks, tas ir polisks''.

Tomēr valsts iedzīvotāju dzīvē ienāca arī jaunas vēsmas. 1925.~gadā Polijā Radiofonijas sabiedrība atklāja pirmo valstī radiostaciju. Nākamajā gadā oficiāli sāka strādāt Poļu Radio pirmā raidstacija. Programma sākās ar F.~Šopēna Polonēzi A-dur. 1931.~gadā Polijā sāka darbu tai laikā visjaudīgākā Eiropā raidstacija, kuras raidījumi aptvēra jau visu Poliju. 1939.~gadā radio abonementu skaits sniedzās jau nedaudz pāri miljonam.

Masu kultūrā arvien lielāku lomu cilvēku dzīvē ieņēma sports, pārstājot būt par bagātnieku privilēģiju. Plašu izplatību guva strādnieku sporta klubi. Veselīgu dzīves veidu un nodarbošanos ar sportu aktīvi propagandēja lauku jaunatnes savienības.

Poļu diska metēja H.~Konopacka 1926.~gada vasarā uzstādīja divus jaunus pasaules rekordus~--- 34,15 un 34,90~m~--- diska mēšanā. Nākamajā gadā atkal laboja rekordu, sasniedzot 39,18~m. 1926.~un 1930.~gadā viņa kļuva par Vispasaules Sieviešu spēļu (notika 1922--1934) uzvarētāju un visbeidzot kā pirmā poļu sportiste ieguva zelta medaļu Amsterdamas olimpiskajās spēlēs (1928). (1940.~gadā viņa emigrēja uz ASV, nodarbojās ar glezniecību). Plašu atpazīstamību pasaulē ieguva poļu skrējējs J.~Kusočinskis (gāja bojā Otrajā pasaules karā Pretošanās kustības cīnītāju rindās)~--- Losandželosas Olimpisko spēļu (1932) čempions 10~000 metru skrējienā.

Populārs bija futbols, airēšana, jāšanas sports.

\section{Polijas ārpolitika}

\epigraph
{Visi līgumi starp lielvalstīm pārstāj būt obligāti, kad tie nonāk pretrunā ar cīņu par izdzīvošanu.}
{Otto fon Bismarks (\detxti{Otto Eduard Leopold von Bismarck-Schönhausen})}

\epigraph
{Tā saucamajai miera politikai \citespace{} bija tikai tik daudz nozīmes, cik tā gatavoja totālam karam.}
{Erihs Ludendorfs (\detxti{Erich Friedrich Wilhelm Ludendorff})}

\epigraph
{Diplomātija un spēks ir vienas un tās pašas monētas divas puses.}
{Bils Klintons (\entxti{William Jefferson ``Bill'' Clinton})}

\epigraph
{Laupītājs paliek laupītājs, vai viņš darbojas viens vai barā, bet valsts, kura ved netaisnīgu karu, ir vienkārši liels laupītāju bars.}
{Bendžamins Franklins (\entxti{Benjamin Franklin})}

\epigraph
{Tieksmi pēc varas rada bailes. Tie, kuri nebaidās no saviem kaimiņiem, neredz nepieciešamību iegūt varu pār viņiem.}
{Bertrans Rasels (\entxti{Bertrand Arthur William Russell})}

\epigraph
{Ja cilvēki būtu patiešām diženi, pasaulē būtu daudz vietas.}
{Staņislavs Ježijs Lecs (\pltxti{Stanisław Jerzy Lec)}}

\epigraph
{Smaga kļūda bieži iegūst nozieguma svaru.}
{Lūcijs Annejs Seneka (\latxti{Lucius Annaeus Sĕnĕca})}

\epigraph
{Tie, kuri atsakās no brīvības pagaidu drošības labā, nav pelnījuši nedz brīvību, nedz drošību.}
{Franklins Delano Rūzvelts (\entxti{Franklin Delano Roosevelt})}

Polija starpkaru periodā bija t.~s. sanitārā kordona (fr. \frtxti{cordon sanitaire}) sastāvdaļa. Atrodoties starp Vāciju un Padomju Savienību, tā nevarēja bez bažām baudīt mieru. Divas lielās kaimiņvalstis: Vācija un Krievija/PSRS nespēja samierināties ar milzīgo teritoriju zaudēšanu par labu Polijai. Šo valstu politika lielā mērā bija vērsta uz Polijas novājināšanu un perspektīvā~--- sadalīšanu.

J.~Pilsudskis par Polijas galveno ienaidnieku uzskatīja Padomju Krieviju/PSRS, otrs ievērojams līderis~--- R.~Dmovskis un viņa vadītā nacionāldemokrātu partija~--- Vāciju. Šādos apstākļos milzu loma valstisko interešu aizstāvēšanā bija pašas Polijas realizētajai ārpolitikai. Kā atzīmējis poļu publicists A.~Vasiļevskis, Polijas sabiedriskajā domā starpkaru periodā valdīja uzskats, ka Polijai ir jābūt par Rietumu aizsargvalni pret boļševistiskajiem Austrumiem. Šī izpratne gan sakrita ar jau minēto \frtxti{cordon sanitaire} būtību, taču pēc autora domām, poļu valstsvīri šai laikā tomēr vienmēr pirmajā vietā stādīja nevis kādas abstraktas Rietumu, bet pašas Polijas intereses, kā nu viņi tās saprata.

Ņemot vērā starptautiskos draudus, daudz uzmanības tika veltīts armijas vajadzībām. Taču nepietika līdzekļu lai pienācīgi sagatavotos nākamajam karam. Tāpēc vēl jo vairāk svarīgi bija īstenot ārpolitiku, kas ļautu stiprināt valsts drošību.

20.~gadu sākumā Polijai bija samērā vēsas attiecības ar Lielbritāniju. 1921.~gada augustā Polijas ārlietu ministrs K.~Skirmunts uzskatīja, ka Lielbritānijai ``ebreju un vācu aprindu ietekmē'' ir pretpoliska nostāja, tās valdība arī baidās no franču hegemonijas kontinentā un tāpēc nostājusies pret Polijas interesēm. Arī vēlāk, kā atzīmējis Polijas vēstures pētnieks Ē.~Jēkabsons, Lielbritānija realizēja kopumā pret Poliju nelabvēlīgu politiku, kuras pamatā bija pieņēmums, ka, pirmkārt, Polija ir iekšēji nestabila, uz kaimiņvalstu rēķina paplašināta valsts, otrkārt, ka mieru iespējams saglabāt, apmierinot lielvalstu (šajā gadījumā~--- Vācijas) prasības.

20.~gados turpinājās Polijas sadarbība ar Franciju. Francija sniedza palīdzību Polijas kara flotes izveidē, izsniedza aizņēmumu militārām vajadzībām. Līdz 1932.~gadam Polijā darbojās Francijas militārā misija ar padomdevējas funkcijām. (Pēc 1926.~gada apvērsuma gan tās loma gandrīz pilnībā izsīka.) Tāpēc Polijas Komunistiskās partijas III konference 1923.~gadā pat atzina, ka Francijas sūtnis Varšavā spēlē tādu pat lomu kā Krievijas imperatores Katrīnas II sūtņi XVIII~gadsimtā, kas gan bija ievērojams pārspīlējums. Taču, kā atzīmējis latviešu vēsturnieks Kārlis Počs, kurš pētījis lielvalstu un Polijas attiecības starpkaru periodā, ``Francija palika uzticīga savam principam~--- par katru Polijas apmierināto prasību lika tai samaksāt ar ekonomisku piekāpšanos Francijai.'' (Pēdējā laikā gan franču historiogrāfijā ir parādījies arī viedoklis, ka 1921.~gada pakts nekad nav nodrošinājis reālu abu valstu savienību.)

Jāuzsver J.~Pilsudska apsvērums, ka Polija spēs izdzīvot blakus lielajiem, potenciāli bīstamajiem kaimiņiem tikai tad, ja tā ne vien balstīsies uz savienību ar Franciju, bet arī pati būs pietiekami militāri stipra, lai iespējamajam uzbrucējam radītu risku ar militāru spēku nesasniegt vēlamo.

Oficiāli Polija ar lielākiem vai mazākiem panākumiem realizēja \strong{``līdzsvara politiku''} (\pltxti{politika rόwnowagi}), cenšoties ievērot vienādu distanci no abiem kaimiņiem: Vācijas un PSRS.

Taču šīs politikas īstenošanu \strong{traucēja} pašas Polijas valdošo slāņu \strong{sapņi par Polijas vadītu valstu konfederāciju}, kura sniegtos no Bulgārijas līdz Skandināvijai jeb, citiem vārdiem, no Melnās līdz Baltijas jūrai. Tiesa, šie plāni paredzēja valstu savienības radīšanu Polijas vadībā, nevis citu valstu okupāciju, taču kaimiņvalstis, ievērojot savu vēsturisko pieredzi, kā reiz baidījās no Polijas varmācības. Viļņas piederības jautājums faktiski padarīja neiespējamu Polijas tuvināšanos ar Lietuvu, bet Cešinas [čehu \cstxti{Těšín}, vācu \detxti{Teschen}, poļu \pltxti{Cieszyn}] jautājums sarežģīja attiecības ar Čehoslovākiju. Latvija kopā ar Rumāniju bija vienīgās Polijas kaimiņvalstis, ar kurām tai nebija nopietnu teritoriālu un politisku konfliktu. Latvijai bija ar Poliju kopējs draudīgs ienaidnieks~--- Padomju Savienība, tādēļ arī Latvija attiecībās ar Poliju izrādīja zināmu sadarbības vēlmi, tomēr sadarbība nekad nekļuva sevišķi cieša. 1921.~gada otrajā pusē Latvijas un Polijas attiecības sarežģījās. Augusta beigās vizītes laikā Lietuvā Latvijas premjers un ārlietu ministrs Z.~Meierovics preses pārstāvjiem atklāti paziņoja: pēdējā Helsinku konferencē ir noskaidrojies, ka Polija ir drauds Latvijai un pašreiz nav iespējas izveidot plašo (piecu valstu: Somijas, Igaunijas, Latvijas, Lietuvas, Polijas) Baltijas Antanti. Starp citu, cēloņi šādai nostājai varēja būt arī tajā laikā Polijā klīstošā prasība pārskatīt Latvijas zemes reformu, kura skāra Polijas muižniecības intereses Latgalē, kā arī Polijas neatlaidīgā tieksme nodrošināt sev izeju pie Daugavas uz austrumiem no Daugavpils.

1922.~gada 17.~martā Varšavā Igaunijas, Somijas, Latvijas un Polijas ārlietu ministri parakstīja līgumu ar kuru viņu pārstāvētās valdības apņēmās ``neslēgt nevienu līgumu, kurš būtu tiešā vai netiešā veidā vērsts pret kādu no viņu valstīm. Valstis, kurās bija mazākumtautības, kuras piederēja skaitliski lielākajai tautībai vienā no pārējām valstīm, vienojās ``garantēt šīm minoritātēm visas tiesības un brīvības, kuras nodrošina tām nacionāli kulturālo organizāciju saglabāšanu un brīvību''. Pārstāvētās valstis arī pasludināja, ``ka gadījumā, ja vienai no tām uzbruks bez iemesla, pārējās saglabās labvēlīgu neitralitāti attiecībā pret valsti, kurai uzbrukts un nekavējoties savstarpēji sazināsies par soļiem, kurus vajadzēs spert.'' Līgums bija augstākais sasniegums centienos 20.~gadu sākumā izveidot Baltijas valstu savienību, kaut arī bez Lietuvas piedalīšanās. Seku gan līgumam nebija, jo tas savstarpējo pretrunu un atšķirīgo ārpolitisko nostādņu dēļ tā arī netika ratificēts. 1923.~gada martā Latvijas sūtnis Polijā Mārtiņš Nukša ziņoja Latvijas valdībai par jauniem Polijas teritoriāliem tīkojumiem: ``Polija centās nostiprināt savu iespaidu Klaipēdā un rauga arī dabūt savās rokās Grīvas dzelzceļa mezglu''. Taču kad Lietuva faktiski pievienoja Klaipēdu (1923--1924), Polija baidījās tieši iesaistoties konfliktā ar Lietuvu, jo Francija šajā laikā bija okupējusi Rūras apgabalu, PSRS koncentrēja karaspēku pie Polijas robežas un Polijai šai laikā varēja draudēt finanšu katastrofa. Pie tam Polijas valdībai gribot-negribot nācās saskatīt analoģiju starp ģenerāļa L.~Žeļegovska gājienu uz Viļņu un lietuviešu realizēto Klaipēdas ieņemšanu.

Vilinošā iespēja Polijai kļūt par ziemeļu un dienvidu kaimiņvalstu saistelementu un par vadošu spēku Austrumeiropā nepiepildījās pašas Polijas izraisīto konfliktsituāciju dēļ. Jāatzīst, ka poļu lolotie Austrumeiropas valstu konfederācijas plāni, ievērojot pašas Polijas saimnieciskās grūtības, augsto bezdarbu, iekšējo etnisko sasprindzinājumu, nespēju visās jomās sekmīgi veikt modernizāciju, kā arī kaimiņvalstu neieinteresētību iecerētās valstu sistēmas izveidē, jāvērtē kā klaji utopiski, balstīti vienīgi atmiņās par valsts dižo pagātni.

Nepiepildoties Polijas centieniem radīt ``Starpjūru valstu savienību'', tai nācās pastiprināti rūpēties par savu drošību. Pirmās draudu pazīmes jaunizveidotajai Polijas valstij parādījās jau ar 1922.~gada Padomju Krievijas un Vācijas vienošanās noslēgšanu, kad padomju un vācu pārstāvji \strong{Rapallo} pilsētiņā Dženovas tuvumā parakstīja \strong{līgumu} par diplomātisko un tirdzniecisko attiecību nodibināšanu, atteikšanos no savstarpējām pretenzijām par veciem parādiem. Ar šo soli abas valstis pārvarēja diplomātisko izolāciju un vienlaikus apliecināja savas revizionistiskās tieksmes attiecībā uz pēckara pasaules starptautisko kārtību. Rapallo līguma noslēgšana rādīja, ka arī XX gadsimtā atkal, tāpat kā XVIII gadsimtā, var rasties starptautisks grupējums, kas savulaik noveda pie Polijas valstiskuma likvidēšanas.

Mūsdienās atslepenotie padomju militārās izlūkošanas dokumenti rāda, ka Sarkanās armijas IV (Izlūkošanas) pārvaldes priekšnieks Jānis~Bērziņš saviem darbiniekiem deva norādījumus sadarboties ar Vācijas reihsvēra izlūkdienestu, lai apmainītos izlūkdatiem par Poliju, tās armiju. Šīs padomju--vācu attiecības turpinājās līdz Ā.~Hitlera nākšanai pie varas 1933.~gadā. Polija baidījās, ka šīs abas valstis varētu vēl vairāk tuvināties, izveidot pret Poliju vērstu savienību.

Apzinoties, ka agri vai vēlu Padomju Savienība vēlēsies atgūt poļu-padomju karā zaudētās teritorijas, poļu valstsvīri ilgstoši PSRS uzskatīja par ienaidnieku Nr.~1. J.~Pilsudskis ne vienu reizi vien uzsvēra: ``Krievijas sadalīšana ir Polijas valsts interešu pamatā Austrumos''. Tāpēc padomju vēsturniekiem bija pamats apgalvot, ka Polijas politika bijusi vērsta galvenokārt pret Padomju Savienību. Polijas valdība tikai gandrīz 3 gadus pēc Rīgas miera līguma noslēgšanu ar \strong{Padomju Krieviju}~--- 1924.~gada februārī nozīmēja savu sūtni Maskavā (līdz tam viņa funkcijas pildīja pilnvarots lietvedis~--- \frtxti{Charge d’affaires}), tai pat laikā turpinot sniegt atbalstu pretpadomju grupējumiem, kuri darbojās pierobežā. Tā, piemēram, KSFPR 1921.~gada aprīlī protestēja pret pretpadomju organizāciju saglabāšanu un darbību Polijā, jūlijā pieprasīja izsūtīt no valsts ievērojamākos krievu ``baltos'' emigrantus, taču protestus nācās atkārtot vēl un vēl. Polija iesniedza pretprasības. Tikai tad, kad KSFPR paziņoja, ka miera līgumā Polijai paredzēto reparāciju izmaksas sākums būs atkarīgs no pretpadomju organizāciju līderu izraidīšanas no Polijas, lieta izkustējās no sastinguma punkta.

Kompromisa rezultātā šie emigranti oktobra beigās (kaut Polijas valdība bija apņēmusies viņus izsūtīt jau 7./8.~oktobrī) atstāja Poliju. Tiesa, valsti atstāja nebūt ne visi pretpadomju darbinieki. Pēc Rīgas miera parakstīšanas pēc aptuveniem aprēķiniem Polijā palika arī no 35~līdz 40~tūkstošiem ukraiņu, kuri bija karojuši neatkarīgas Ukrainas valsts piekritēju rindās. Daļa no viņiem~--- ap 10~tūkstošu atradās poļu nometnēs. Kaut Ukrainas SPR valdība Polijas valdībai iesniedza protestus par petļuroviešu darbības atbalstīšanu, ar poļu specdienestu palīdzību no ukraiņiem~--- petļuroviešiem tika radītas vienības, kuras 1921.~gada oktobrī-novembrī no Polijas teritorijas iebruka Padomju Ukrainā ar mērķi izraisīt tur sacelšanos. Tikai novembra beigās Sarkanā armija spēja šīs ukraiņu nacionālistu vienības izspiest no padomju teritorijas un tās atgriezās Polijā, kur tika vēlreiz internētas.

T.s. Ukrainas Tautas Republikas (UTR) armijas karavīri internēto nometnēs, kuras pastāvēja līdz 1924.~gadam, saglabāja savu organizāciju, baudīja poļu varas iestāžu labvēlību. UTR politiķi pat būvēja plānus pārvērst Austrumgalīciju par sava veida ukraiņu Pjemontu, taču Polijas puse šādus plānus neatbalstīja, bet traktēja šo rajonu kā Polijas integrētu sastāvdaļu.

Padomju valdība turpināja iesniegt protesta notas, prasīja lai Polijas valdība godīgi pilda Polijas un Padomju Krievijas 1921.~gada miera līgumu, pārtrauc atbalstīt dažādas ``balto'' un nacionālistu vienības, kas ielauzās padomju teritorijās, terorizēja to iedzīvotājus. Poļu un padomju vēsturnieku 1966.~gadā izdotajā dokumentu krājumā par padomju-poļu attiecībām ievietotie dokumenti pārliecinoši pierāda, ka Polijas armijas Ģenerālštābs un arī valdība vadījās nevis no starptautisko tiesību normām, bet savām pragmātiskajām, uz izlūkdatu iegūšanu un Padomju varas vājināšanu vērstajām interesēm. Pretboļševistisko vienību gatavošana un darbība jau padomju teritorijā notika ar oficiālo Polijas iestāžu ziņu un tiešu atbalstu.

Spilgts piemērs bija atamana, bijušā UTR armijas ģenerāļa J.~Tjutjuņņika uzturēšanās Polijā, no kurienes viņa vadītie ukraiņu nacionālisti vairākkārt iebruka Padomju Ukrainas teritorijā. Padomju varas iestāžu rokās nonāca un tika publicēti dokumenti, kuri tieši apliecināja minētā ukraiņu nacionālistu vadītāja un Polijas armijas štāba sadarbību. Padomju drošības dienests notvēra un pārvervēja ģenerāļa aģentus, kuri viņam ziņoja, ka Ukrainā pastāv slepena pagrīdes organizācija, kas tikai gaida uz viņa ierašanos lai uzsāktu plašu pretpadomju sacelšanos. Tā pēdējais, jau galīgi bezcerīgais mēģinājums izraisīt ``Visukrainas sacelšanos'' notika 1923.~gada vasarā un atkal ne bez Polijas armijas Ģenerālštāba atbalsta. Kad 1923.~gada 26.~jūnijā J.~Tjutjuņņiks kārtējo reizi nelegāli pārgāja robežu, viņš tika notverts, piekrita sadarboties ar padomju iestādēm, nosūtīja ukraiņu nacionālistu vadītājiem vēstules, kurās izklāstīja, ka pretpadomju cīņa ir bez izredzēm uz panākumiem un tāpēc viņš pāriet padomju pusē. Vēlāk ģenerālis apmācīja padomju virsniekus, 1924.~gadā publicēja grāmatu ``\uktxti{З поляками про України}'' (Ar poļiem pret Ukrainu). Tas gan viņu nepaglāba no priekšlaicīgas nāves. 1929.~gadā čekisti viņu arestēja un 1930.~gadā nošāva.

Cīņā pret pār robežu nākošajām bandām Sarkanā armija no savas puses izvērsa īpašu pretdarbību. 1922.~gada 18.~martā Republikas Revolucionārā Kara padome apstiprināja armijas virspavēlnieka S.~Kameņeva priekšlikumu griezties pie Ārlietu Tautas komisariāta ar ierosinājumu brīdināt Rumānijas un Polijas valdības, ka atkārtoti S.~Petļuras un B.~Savinkova bandu uzbrukumi izsauks padomju pierobežas apgabalu iedzīvotāju tādu pat pretdarbību, bet padomju militārās iestādes varēs pilnībā garantēt, ka no padomju puses uz Poliju un Rumāniju nekādi bandu iebrukumi netiks pieļauti, tikai tad, ja padomju robežas paliks neaizskartas. Turpretī atkārtotu uzbrukumu gadījumā padomju teritorijai vietējās militārās iestādes paziņo, ka tās nekādi nespēs garantēt Rumānijas un Polijas robežu neaizskaramību, jo ``pārāk stingra politika šai ziņā'' sastapsies ar vietējo iedzīvotāju neizpratni.

Tiesa, padomju prasības izraidīt no Polijas teritorijas pretpadomju teroristisku formējumu dalībniekus turpinājās līdz pat 1939.~gadam, tomēr situācija uz robežas uzlabojās, bruņotas sadursmes kļuva retas.

Polija un PSRS savstarpēji centās ietekmēt iedzīvotājus otrpus robežas. Šajā cīņā veidojās visai neparastas kombinācijas. Par padomju ``aktīvo izlūkošanu'' jau runāts. Bez tam padomju izlūkdienests sadarbojās ar Lietuvas specdienestiem, atbalstīja ukraiņu nacionālos nemierniekus Galīcijā, kuri toreiz karoja pret poļiem. Tam visam vajadzēja destabilizēt stāvokli Polijā.

Padomju izlūkdienesta rezidents ar diplomātiskā darbinieka piesegu Polijā bija M.~Laganovskis. Viņš bija polis pēc tautības, bijušais sociālists, pārgājis komunistu rindās, par varonību Pilsoņu karā apbalvots ar Sarkanā karoga ordeni, viņu augstu vērtēja F.~Dzeržnskis. 1923.~gadā M.~Laganovskis saņēma uzdevumu radīt Polijā kaujas organizāciju. Tās vadībā nostājās divi poļu virsnieki, kuri pirms tam bija karojuši Polijas armijā: poručiks V.~Baginskis un podporučiks A.~Večorkevičs. V.~Baginskis bija pat apbalvots ar Kara nopelnu ordeņa Sudraba krustu (\pltxti{Krzyż Srebrny Orderu Wojennego Virtuti Militari}). Organizācijas uzdevums bija veikt diversijas un teroristiskus aktus, lai veicinātu revolucionārās cīņas pieaugumu Polijā. Politisko un sabiedrisko organizāciju telpās notika sprādzieni. Bumbas sprāga it kā bez politiskās orientācijas ievērošanas, pēc pulksteņa svārsta principa: te~--- labējā, tad atkal kreisā organizācijā. Radās iespaids, ka darbojas divas ekstrēmistiskas organizācijas, kuras mērķtiecīgi iznīcina savus politiskos pretiniekus. Daži poļu publicisti šo laiku Polijas vēsturē ir nosaukuši par ``bumbu periodu''. 1923.~gada 2.~augustā V.~Baginski un A.~Večorkeviču arestēja, apsūdzot par sprādzienu organizēšanu Varšavā un Krakovā aprīlī un maijā. Vēlāk prokuratūra viņus apsūdzēja arī dalībā sprādziena organizēšanā Varšavas citadelē 13.~oktobrī (tātad jau krietnu laiku pēc abu aresta), kur uzsprāga munīcijas noliktavas un bija 25~nogalinātie un 40~ievainotie. Vēlāk noskaidrojās, ka par traģēdijas cēloni kļuva kāda arsenāla strādnieka pieļauts drošības prasību pārkāpums, taču uz viena cilvēka (policijas aģenta-provokatora) liecību pamata tika tiesāti abi virsnieki. Valdība pēc šī sprādziena publicēja uzsaukumu iedzīvotājiem, kurā rakstīja, ka cietoksni uzspridzinājusi ``noziedzīga roka'', taču norādīt, kam tā pieder, pierādījumu trūkuma dēļ nespēja. Avīze ``\pltxti{Kurier Poranny}'' (``Rīta Kurjers'') rakstīja, ka galvenā liecinieka tiesā sniegtās liecības bija tik juceklīgas, ka tās nebija grūti atspēkot. Neraugoties uz to, abiem virsniekiem piesprieda nāves sodu. Avīze ziņoja, ka Polijā ir cilvēki, kuri ir ieinteresēti, lai lietas apstākļi nekļūtu plaši pazīstami. Vēlākos gados vēsturnieki secināja, ka kriminālvajāšana bija policijas provokācijas rezultāts, bet galvenais liecinieks bija pilna laika policijas informators.

Valsts prezidents S.~Vojcehovskis notiesātajiem nāves sodu nomainīja pret 15~gadu ieslodzījumu cietumā. Padomju Savienība ierosināja abu poļu komunistu apmaiņu pret diviem par pretpadomju darbību arestētiem Polijas pilsoņiem. (Jau 1922.~gada 3.~novembrī Polijas valdība bija pieņēmusi Padomju valdības priekšlikumu par personu apmaiņu, kurā bija ieinteresētas abas puses. Šai sakarā no Polijas cietumiem tika atbrīvoti un uz PSRS izsūtīti ap 340~poļu revolucionāru, notiesātu uz ilgstošiem termiņiem cietumā. Arī vēlākajos gados~--- 1928., 1932.~--- notika apmaiņa ar politieslodzītajiem.) Polijas valdība apmaiņai piekrita, tai bija jānotiek 1925.~gada 29.~martā pierobežas dzelzceļa stacijā. Poļu ierēdņu pavadībā ķēdēs ieslēgtie V.~Baginskis un A.~Večorkevičs tur tika nogādāti, bet tad viens no apsargiem policists J.~Muraška abus apsargājamos nošāva.

Uz slepkavību, izdodot uzsaukumu, nekavējoties atsaucās Polijas Komunistiskās partijas Centrālā komiteja. Valdība tika apvainota par to, ka ar vienu roku gatavojot ieslodzīto apmaiņu, ar otru tā sagatavoja viņu noslepkavošanu. ``Večorkeviču un Baginski nogalināja teroristiska organizācija. Klātesošie valdības pārstāvji bija mēmi nozieguma liecinieki \citespace{} Slepkavību veica teroristiska banda saskaņā ar PPS un buržuāzijas nodomiem. \citespace{} Buržuāzija nespēj valdīt citādi kā tikai ar terora palīdzību. Kur nepietiek ar ``likumīgām'' tiesām un slepkavībām, tur darbojas buržuāzijas fašistiskas organizācijas''. Kareivji tika aicināti atriebt V.~Baginska un A.~Večorkeviča nāvi.

Arī Padomju Savienība apvainoja Poliju norunas nepildīšanā un apmaiņai paredzētos poļus tai neatdeva. Padomju sūtniecības darbinieki Varšavā par sprādzieniem turēja aizdomās M.~Laganovska vadītos ļaudis, taču viņš to kategoriski noliedza.

Kad J.~Murašku tiesāja, apsūdzētais paziņoja: ``Es atzīstos divu traku suņu slepkavībā. Nogalināju viņus lielā uztraukumā.'' Polijas tiesa atzina, ka slepkava darbojies afekta stāvoklī un piesprieda viņam divus gadus viegla režīma cietumā, taču pēc gada par priekšzīmīgu uzvedību atbrīvoja. Pēc soda izciešanas, lai ar grūtībām, viņam izdevās atrast darbu robežapsardzē. Pēc dažām ziņām 1939.~gada beigās viņš kļuva par Vācijas policijas darbinieku un to nogalināja poļu pagrīdnieki.

Mūsdienās V.~Baginska un A.~Večorkeviča vārdā Baltkrievijas pilsētās ir nosauktas ielas.

Ar Polijas varas iestāžu atbalstu pret PSRS darbojās dažādas emigrantu organizācijas. Jau Pirmā pasaules kara laikā J.~Pilsudskis bija iznēsājis ģeopolitisku koncepciju, kura, kā arī uz tās pamata vēlāk izveidotā kustība, ieguva nosaukumu \strong{``prometeisms''} (``\pltxti{prometeizm}''). Poļu militārā vēsturnieka E.~Karaskeviča izklāstā ``prometeisma'' mērķis bija pretdarboties ``krievu draudiem'' Polijai, izmantojot dažādu etnisko grupu pretrunas Krievijas valsts sistēmā, gala rezultātā ar nacionālistisko kustību palīdzību panākot Krievijas, vēlāk Padomju Savienības novājināšanos un sadalīšanos. Polijai bija jāatbalsta nekrievu tautu, kuras dzīvoja Krievijā (Padomju Savienībā) neatkarības iegūšana. Nosaukums bija veiksmīgs. Mīts par Prometeju~--- cīnītāju par cilvēces laimi imponēja ne tikai revolucionāriem, bet arī liberāļiem, kuri kādreiz ienīda Krievijas impēriju, bet pēc tam arī PSRS. Plāns bija cieši saistīts ar J.~Pilsudska iecerēto ``\pltxti{Międzymorze}'' (``\latxti{Intermarium}'' jeb ``Starpjūru valstu savienības'') izveides ideju. ``Prometeisma'' kustībai bija jārada pretkrieviska josla no Baltijas līdz Melnajai jūrai. Mūsdienu krievu specdienestu vēstures pētnieki V.~Biļiņins, A.~Zdanovičs un V.~Korotajevs uzsver, ka ``prometeisma'' darbība izvērsās īpaši sakarā ar poļu-padomju karu, cenšoties vairāk vai mazāk no Polijas atkarīgās atsevišķas nacionālās kaimiņu valstiņas iesaistīt pretpadomju darbībā.

Pēc 1926.~gada apvērsuma graujošo darbību pret PSRS vadīja J.~Pilsudska tuvs līdzgaitnieks ģenerālis J.~Stahevičs. Poļu valdošās aprindas nodibināja kontaktus ar dažādiem PSRS dzīvojošo, bet boļševiku varu nepieņēmušo nekrievu tautu ``valdībām'' emigrācijā un nacionālajiem centriem. 1926.--1928.~gadā tika likti pamati sadarbībai un materiālai palīdzībai šīm emigrantu grupām. Ideoloģiskā darba PSRS teritorijā nostādņu izstrādei 1926.~gadā Varšavā tika radīts Austrumu institūts (\pltxti{Instytut Wschodni w Warszawie}), kurš ik gada ceturksni izdeva žurnālu ``\pltxti{Wschod}'' (``Austrumi'').

Emigrantu organizāciju aģentu saviem centriem sniegtās ziņas tika nodotas arī Polijas izlūkdienesta rīcībā. Polijas armijā uz līgumu pamata tika uzņemti dienestā nacionālās emigrācijas grupu pārstāvji. Tādejādi nacionālie centri būtu apgādāti ar pieredzējušiem militāriem kadriem gadījumam, ja padomju teritorijā notiktu pretboļševistiskas sacelšanās. Pirmie virsnieki-gruzīni ar īpašu līgumu tika uzņemti Polijas armijā 1921.--1923.~gadā. Dažādos laika posmos viņu skaits svārstījās no 35~līdz 80~cilvēkiem. 1926.~gadā tādā pat statusā tika uzņemti vairāk nekā 30~ukraiņi-petļurovieši. 1938.~gada beigās Polijas armijā uz līgumu pamata dienēja 40~ukraiņu virsnieku. Pastāvēja arī nelielas virsnieku~--- azerbaidžāņu un kalnu Kaukāza iedzīvotāju grupas.

20.~gadu beigās Varšavā nacionālistisko organizāciju pārstāvji ar Polijas iestāžu atbalstu radīja klubu ``\pltxti{Prometeusz}'', (Prometejs). Tā oficiālais nosaukums skanēja: ``Prometejs~--- Krievijas apspiesto tautu: Azerbaidžānas, Donas, Karēlijas, Urālu, Ingrijas, Krimas, Kubaņas, Kaukāza, Turkestānas un Ukrainas līga''. Tika reorganizētas militārās sadarbības formas, tika veidoti nacionālistu štābi. Piemēram, ukraiņu štāba darbības principi paredzēja, ka tas ir Polijas armijai draudzīga kara organizācija, kura darbojas Polijas teritorijā un izmanto Polijas armijas Ģenerālštāba tehnisku un materiālu palīdzību. Ukraiņu štāba uzdevumi bija mobilizācijas--operatīvā darbība, izlūkošana un propaganda. (Faktiski Polijas politikas vadītājiem maz rūpēja ukraiņu nacionālistu mērķi. Viņi stādīja savus mērķus. Tā, 1934.~gada jūnijā Polijas Ārlietu ministrijas Austrumu nodaļas vadītājs T.~Šetcels sarunā ar Bulgārijas pilnvaroto lietvedi paziņoja, ka Polija ``rēķinās ar to, ka ja Tālajos Austrumos sāksies [Japānas-PSRS] karš, PSRS tiks sakauta, un tad Polija iekļaus savās robežās Kijevu un daļu Ukrainas.'') Lūk, kas par ``Prometeju'' bija teikts Polijas armijas Ģenerālštāba Izlūkošanas nodaļas 1937.~gada augustā izstrādātā dokumentā: ``Prometejisms ir visu bez izņēmuma Krievijas apspiesto tautu kustība \citespace{} lai izsauktu nacionālu revolūciju PSRS teritorijā \citespace{} ``Prometejs'' mobilizē biedrus pēc pašu gribas un ar pašu atbildību, neuzņemoties nekādu politisku atbildību attiecībā pret nacionāliem centriem \citespace{} ``Prometejam'' jābūt tiesībām demonstrēt nacionālu radikālismu lai visefektīvākā veidā radītu revolucionāru dinamiku. To nedrīkst vainot par radikāli nacionālistiskām tendencēm un tās nedrīkst nepareizi vērtēt kā fašistiskas\dots{}''. Faktiski arī 30.~gados nepastāvēja robeža starp dažādu nacionālu kustību atbalstīšanu un pret PSRS vērstu graujošu darbību.

Aktīvi pret PSRS darbojās pašas Polijas izlūkošanas un pretizlūkošanas dienesti. Polijas izlūkdienests patstāvīgi spiegoja Padomju Savienībā. Piemēram, ir zināms gadījums, kad poļu izcelsmes Sarkanās armijas virsnieks brigādes komandiera pakāpē (\rutxti{комбриг}) ar Sarkanā karoga ordeni apbalvotais B.~Kontrums, nosūtīts mācīties Sarkanās armijas Kara akadēmijā, 1922.~gada februārī uzņēma sakarus ar Polijas militāro atašeju R.~Vaļikovski, apgādāja viņu ar informāciju par Sarkano armiju, bet pēc tam, inscinējis savu nāvi, 1922.~gada beigās nelegāli pārbēga uz Poliju.

Tur B.~Kontrums dienēja policijā. Otrā pasaules kara laikā karoja pret nacistiem Francijā, Norvēģijā un Polijā, bija arī Varšavas sacelšanās dalībnieks, nokļuva vācu gūstā, izbēga, atgriezās Polijā 1947.~gadā, dzīvoja Varšavā, strādāja rūpniecībā.

1948.~gada 13.~oktobrī B.~Kontrumu arestēja Polijas Valsts drošības dienests. 1952.~gada 9.~oktobrī PTR Augstākā tiesa apstiprināja Varšavas apgabaltiesas B.~Kontrumam piespriesto nāves spriedumu par ``noziegumiem'', izdarītiem policijas dienesta laikā pirmskara Polijā, un ``pretvalstisko darbību'' Otrā pasaules kara laikā. 1953.~gada 2.janvārī B.~Kontrums tika pakārts. Pēc nāves reabilitēts 1957.~gada 3.~decembrī.

Teiktais tomēr nenozīmē, ka divdesmitajos gados nebija nekādu iespēju kaut daļēji uzlabot abu valstu attiecības. Šādas iespējas vismaz teorētiski pastāvēja. Jau 1922.~gada decembrī Maskavā notikušajā konferencē par atbruņošanos Polija iesniedza padomju pusei priekšlikumu noslēgt līgumu par neuzbrukšanu. Polija centās noslēgt kolektīvu vienošanos ar Polijas, Baltijas valstu, Somijas un PSRS dalību un aizstāvēja šo priekšlikumu vairāku gadu garumā. PSRS saskatīja šai priekšlikumā Polijas vēlmi radīt pretpadomju bloku un iegūt vadošu lomu Polijas-Baltijas reģionā. Padomju vadība iebilda pret Polijas projektu, iestājoties par divpusēja līguma noslēgšanu.

1924.~gadā, kad sākās plaša starptautiska PSRS kā valsts atzīšanas kampaņa, arī Polija nedaudz mainīja savu attieksmi pret to. Pozitīvas pārmaiņas abu valstu attiecībās iezīmēja Polijas sūtņa L.~Darovska darbība Maskavā. Taču tās neieguva noturīgu raksturu, jo Polijas ārlietu ministra 1924.--1926.~gadā A.~Skšinskis realizētajā politikā savijās paternālisms pret mazākajām kaimiņvalstīm un pretpadomju tendences ar vēlmi uzturēt labākas kaimiņattiecības ar PSRS.

Faktiski Polijas~--- Padomju Savienības politiskās sarunas, kuras norisa 1924.--1928.~gadā no abām pusēm bija tikai diplomātiska spēle ar mērķi katrai pusei nostiprināt savas pozīcijas Eiropā, atrisināt savas drošības un pilnvērtīgas ekonomiskās attīstības problēmas. Tā, 1924.~gada 28.~augustā KK(b)P CK Politbirojs deva instrukciju, vedot sarunas ar Polijas pārstāvjiem, ``neatteikties no PSRS--Polijas robežu labošanas principa.''

1925.~gada janvārī PSRS pilnvarotais P.~Voikovs tikšanās laikā ar Polijas ārlietu ministru A.~Skšinski piedāvāja noslēgt PSRS--Polijas vienošanos par neuzbrukšanu un atteikšanos iesaistīties kādās kombinācijās, kuras būtu naidīgas otrai pusei. Janvāra beigās Polijas valdība jau atkārtoti ierosināja PSRS noslēgt vienošanos par neuzbrukšanu ne tikai ar Poliju, bet vairākām valstīm: Poliju, Baltijas valstīm, Somiju, un Rumāniju. PSRS to iztulkoja kā Polijas mēģinājumu radīt savu protektorātu pār minētājām valstīm ar nolūku organizēt karu pret PSRS. Pēdējās ārpolitikā tai laikā izpaudās vēlme noslēgt individuālus līgumus ar atsevišķām valstīm, kas neļautu tām ``veidot citu citai naidīgas grupas un kombinācijas'', ``likvidēt politisko savienību un grupu sistēmu, kura neizbēgami ved pie kara''. Var ilgstoši strīdēties, kurš plāns tolaik vairāk atbilda miera nostiprināšanas interesēm. Abas puses 20.~gados nespēja vienoties un noslēgt neuzbrukšanas līgumu, jo nebija ieinteresētas nekavējoties panākt savstarpēju tuvināšanos. Poļu-padomju sarunām šajā periodā nebija patstāvīgas nozīmes, tās bija atkarīgas no vispārējās starptautiskās situācijas attīstības. Bet tā nebija labvēlīga PSRS un Polijas attiecību normalizēšanai. Gluži pretēji, 1926.~gada martā Buharestē Polija un Rumānija noslēdza līgumu. To papildināja slepena militāra konvencija, kura paredzēja abu valstu savstarpēju militāru palīdzību iespējamā kara pret PSRS gadījumā.

Interesantu liecību par Polijas politiķu attieksmi pret PSRS sniedz gadījums 1925.~gada beigās. Kā savās atmiņās rakstīja toreiz zemnieku deputāts A.~Fiderkevičs, šai laikā Polijas Seima deputātu daļas vidū sāka nostiprināties doma par nepieciešamību Polijai izvērst sadarbību ar PSRS. Tika sastādīts 18~deputātu no dažādām frakcijām saraksts, kuriem bija jāapmeklē PSRS, lai nospraustu ekonomiskās sadarbības principus. Ideju atbalstīja toreizējais Polijas premjerministrs un ārlietu ministrs A.~Skšinskis, taču pēkšņi kļuva zināms, ka PPS deputāti atteikušies piedalīties delegācijā. Izplatījās ziņa, ka braucienam uz PSRS pretojies J.~Pilsudskis, kurā sadzīvoja gluži saprātīgs boļševiku vērtējums (1926.~gadā viņš atzina, ka komunistiskās iekārtas nomaiņā Krievijā Polija nav ieinteresēta, jo katra cita iekārta, acīmredzot būs mazāk labvēlīga Polijai) ar nesamierināmu naidu pret visu krievisko un padomisko. Viņš ietekmēja poļu sociālistu un savu piekritēju citās partijās pozīciju, nosaucot nodomu par braucienu uz PSRS pat par ``mežonīgu ieceri''.

Pēc 1926.~gada maija valsts apvērsuma Polijā tās politika pret PSRS būtiski nemainījās. Polijas attiecības ar PSRS palika vēsas. Pat tādas tās tika būtiski apdraudētas 1927.~gadā, kad Lielbritānija pārtrauca diplomātiskās attiecības ar PSRS, bet Varšavā tika nogalināts PSRS pilnvarotais pārstāvis P.~Voikovs. Slepkava uz padomju diplomātu izšāva 6~reizes, 2~šāvieni bija trāpīgi. Jau krītot, P.~Voikovs paguva vēl divas reizes izšaut uzbrucēja virzienā, taču netrāpīja. Slepkavam krievu emigrantam B.~Koverdam, kurš kā nozieguma motīvu minēja P.~Voikova dalību cara ģimenes nogalināšanā un sarkanā terora realizēšanā kopumā, poļu tiesa piesprieda mūža ieslodzījumu, taču 1937.~gadā viņš tika amnestēts un atbrīvots.

Daļā poļu tautas PSRS pilnvarotā pārstāvja nogalināšana izsauca sašutumu un līdzjūtību cietušajam. Uz vēstniecību atvadīties no mirušā stundām stāvēja cilvēku rindas. Policija bija apmulsusi un pēc dažiem neveiksmīgiem mēģinājumiem izklīdināt sanākušos, sāka pildīt savus tiešos pienākumus un uzturēt kārtību atnākušo vidū. Šķirsta ar PSRS pārstāvja mirstīgajām atliekām nosūtīšanas uz dzimteni laikā notika plašas demonstrācijas. To izdzenāšanai tika iesaistītās visas policijas rezerves. Daudzi demonstranti tika pamatīgi piekauti, taču arestu bija maz.

Sakarā ar PSRS pārstāvja nogalināšanu radās nopietns konflikts PSRS un Polijas attiecībās, pārtrūka sarunas par neuzbrukšanas pakta un tirdzniecības līguma noslēgšanu. Polijas valdība konfliktu noregulēja, izsūtot no valsts vairākus redzamākos ``balto'' emigrantu organizāciju pārstāvjus un brīdinot šīs organizācijas par nepieciešamību ievērot Polijas starptautiskās saistības un tiesiskās normas. Šī brīdinājuma neievērošanas gadījumā ``baltajiem'' draudēja viņu organizāciju un preses izdevumu slēgšana un vainīgo izsūtīšana no valsts. Tomēr dažādi konflikti turpinājās. Tā, 1927.~gada rudenī tika veikts atentāts pret PSRS pilnvaroto pārstāvi Polijā A.~Uļjanovu, 1928.~gada pavasarī tika gatavots atentāts pret jauno PSRS pilnvaroto D.~Bogomolovu. 1928.~gada maijā baltemigranti šāva uz PSRS tirdzniecības pārstāvi Polijā A.~Ļizarevu. 1930.~gada aprīlī bija mēģinājums Varšavā uzspridzināt PSRS diplomātisko pārstāvniecību. Atkārtoti uzbrukumi PSRS konsulātam Ļvovā 1933.~gadā prasīja cilvēku upurus.

Sakarā ar tās pilsoņu tiesāšanu Padomju Savienībā Polija pārtrauca ar to sarunas par neuzbrukšanas līguma noslēgšanu (tās netika atjaunotas līdz pat 1931.~gada augustam), tāpēc Maskavā, pēc dažu krievu vēsturnieku teiktā, pat esot radušās bailes par kara sākšanos. (Jau XX~gadsimta 70.~gados padomju vēsturnieki atzina, ka nekādu karu pret PSRS 20.~gados J.~Pilsudskis neplānoja.) 1927.~gada maijā PSRS ārlietu tautas komisārs G.~Čičerins, tiekoties ar Francijas ārlietu ministru A.~Briānu, norādīja uz Polijas īpašo nozīmi PSRS ārpolitikai un nepieciešamību to atturēt no avantūrām. A.~Briāns atbildēja, ka Francija jau agrāk apvaldīja un arī turpmāk apvaldīs Poliju.

PSRS vadība arī izrādīja interesi par Kelloga~--- Briāna paktu (līgumu par atteikšanos izmantot karu nacionālās politikas īstenošanai), kuru 1928.~gada augustā Parīzē parakstīja sākotnēji 15~valstis (vēlāk~--- vēl~48), tai skaitā arī Polija. PSRS nebija uzaicināta parakstīt līgumu, bet ``sacēla skandālu'' un septembrī tam pievienojās. Paktam bija jāstājas spēkā 1929.~gada jūlijā, bet PSRS valdība uzstājās kā iniciatore priekšlikumam, lai kopā ar to Polija un Lietuva paraksta protokolu par pakta priekšlaicīgu stāšanos spēkā. Polija toties vēlējās lai tādu protokolu parakstītu ne tikai nosauktās, bet arī citas kaimiņvalstis. PSRS nācās piekrist Polijas priekšlikumam. 1929.~gada februārī Maskavā PSRS, Polija, Rumānija, Igaunija, Lietuva, Irāna un Turcija parakstīja protokolu par pakta priekšlaicīgu stāšanos spēkā. Tas bija ievērojams Polijas solis uz samierināšanos ar PSRS.

Polijas diplomātiem darbu Padomju Savienībā sarežģīja poļu politiķu, ierēdņu, preses visai naidīgā izturēšanās pret PSRS.

Tā, 1927.~gadā Polijas vēstnieks PSRS S.~Pateks ziņojumā ārlietu ministram A.~Zaleskim sūdzējās par bijušā Polijas vēstnieka PSRS (1925--1927) S.~Ketšinska paziņojumiem presē, ka krievi ir krāsainie, tikai nokrāsoti balti, ka PSRS turpina carisma aizsākto politiku, ka PSRS ir neskaitāmu nejēdzību valsts, ka gan politiski, gan ekonomiski tā tuvojas katastrofai. Saprotams, lai cik lielā mērā šie spriedumi arī netuvinātos patiesībai, diplomāts par savas uzturēšanās valsti, lai arī jau bijušo, tā izteikties nedrīkstēja. Viņa pēctecim Maskavā tas traucēja sarunu vešanu ar padomju pārstāvjiem. 1928.~gadā notika diplomātisko notu apmaiņa starp PSRS un Polijas valdībām sakarā ar Volīnijas gubernatora (vēlāk arī Polijas iekšlietu ministra) H.~Jozefovska izteikumiem, kurus Polijas prese iztulkoja kā aicinājumu radīt Polijai draudzīgu Ukrainas valsti, bet PSRS~--- kā aicinājumu graut PSRS integritāti.

Šeit gan jāpiebilst, ka abu valstu labu kaimiņattiecību uzturēšanai tāpat traucēja daudzi Padomju vadības tieši veikti un tās inspirēti it kā spontāni padomju sabiedrisko organizāciju soļi. Piemēram, 1927.~gadā ar padomju Sarkanā Karoga ordeni par ``varonīgu cīņu zem Padomju Sociālistisko Republiku Savienības un Polijas strādnieku un zemnieku brālības karoga'' tika apbalvots Polijas Seima bijušais deputāts S.~Lancuckis, kuram 1924.~gadā Seims atņēma deputāta neaizskaramību un izdeva tiesai par aicinājumu gāzt valdību un nodibināt strādnieku un zemnieku varu. (Jau tolaik PSRS tika izvērsta plaša masu kampaņa viņa aizstāvībai.) Tiesa viņu attaisnoja, bet tad 1925.~gadā Seima vairākums atkal nobalsoja par viņa izdošanu tiesai par pret valdību vērstu uzstāšanos mītiņā. Šoreiz viņš tika notiesāts un par apbalvošanu ar padomju ordeni uzzināja cietumā. Bija arī citi līdzīgi fakti. Tā, Baltkrievijas PSR Centrālās Izpildkomitejas priekšsēdētājs A.~Červjakovs 1927.~gada Oktobra apvērsuma gadadienai veltītās demonstrācijas laikā poļu valodā izsaucās: ``Lai dzīvo Polijas Padomju republika!'' Šādu uzskaitījumu varētu turpināt ilgi. Saprotams, ka Polijas valdība šādus gadījumus vērtēja kā starptautiskajā praksē nepieņemamus soļus.

Lai nu kā, fakti rādīja, ka Polijas vadītāji pārvērtēja boļševiku pretinieku iespējas PSRS graušanā un nenovērtēja tās iespējamās briesmas, ko valstij nākotnē varēja radīt saspringtās attiecības ar PSRS. Visas starpkaru divdesmitgades pieredze rādīja, ka J.~Pilsudska koncepcija kā stiprināt Polijas drošību, kuras pamatideja bija uzdevums panākt maksimālu Krievijas/PSRS novājināšanu, bija nereāla. Poļu komunisti nepārtrauca pārmest savas valsts valdībai atkarību no starptautiskā imperiālisma, tādas politikas realizēšanu, kas atņem Polijai neatkarību, pārvērš to par imperiālistisko lielvalstu piedēkli, valsti, ar kuras interesēm var nerēķināties. Kā rādīja vēlākie notikumi, Polijas patstāvību ārpolitikā poļu komunisti vērtēja neadekvāti zemu.

Arī \strong{Polijas un Vācijas starpā valdīja dziļa neuzticība}. Polija kā viens no Pirmā pasaules kara uzvarētājiem (Pēc vācu viedokļa~--- par tādu Polija bija kļuvusi pilnīgi nepelnīti) dabīgi tapa par revizionistisku tīkojumu objektu. Kā rakstīja poļu vēsturnieks J.~Kozenskis, ``Versaļas miers ne tikai neapmierināja nevienu no pusēm, bet arī radīju bīstamu un neizdevīgu situāciju Polijai. Pastavīgi draudi Gdaņskai, Veimāras republikas atteikšanās no atzīšanas lika Polijai līdz Otrā pasaules kara sākumam pastāvīgi būt uzmanīgai.''

Pēc vācu politiķu viedokļa Polija kā valsts radās, izmantojot konjunktūras rakstura Rietumvalstu protekciju. Vācija atteicās \latxti{de~iure} atzīt jauno robežu ar Poliju un savu bijušo zemju atsavināšanu. Vācu vēsturnieks J.~K.~Henšs norādīja: ``Prasība pēc austrumu robežas revīzijas saņēma visu vācu partiju, izņemot VKP [Vācijas Komunistiskā partija] beznosacījumu atbalstu, izšķiroši ietekmēja visu Veimāras republikas valdību ārpolitiku un pēc sakāves karā nedrošajai un saplosītajai vācu sabiedrībai kļuva par vienu no nedaudzajiem nacionālās integrācijas faktoriem.''

Kā raksta amerikāņu vēsturnieks G.~Veinbergs: ``Nav nekādu pierādījumu tam, ka kāds no vadošus amatus Veimāras republikā ieņemošajiem stipras un neatkarīgas Polijas pastāvēšanu uzskatīja Vācijai par labu nākam.'' Poliju Vācijā sauca par ``sezonas'' valsti, uzskatot par rīcības nespējīgu tās ģeogrāfiskā stāvokļa un Rietumvalstu nenoteiktās pozīcijas dēļ.

Kā atzīmējis vācu vēsturnieks A.~Lavatu, par ``sezonas'' valsti vācu nacionālisti sauca ne tikai Poliju, bet arī paši savu t.s. Veimāras republiku, tikai ar to atšķirību, ka no pēdējās drupām bija jārodas jaunai Vācijas lielvalstij, bet Polijai bija atkal jānogrimst nebūtībā. Tikai pēc poļu-padomju kara un īpaši pēc J.~Pilsudska īstenotā valsts apvērsuma 1926.~gadā Vācijā sāka ataust atskārsme, ka Polija nav ``sezonas'' valsts, kas pastāv tikai no ``sabiedroto žēlastības''.

Faktiskais Vācijas armijas pavēlnieks (\detxti{Chef der Heeresleitung}), ģenerālis H.~f.~Sēkts jau 1922.~gadā paziņoja; ``Polijas eksistence ir nepanesama, nesavienojama ar Vācijas pastāvēšanas nosacījumiem. Tai ir jāpazūd un tā pazudīs. Pie tā to novedīs pašas iekšējais vājums un Krievija~--- ar mūsu palīdzību''. Arī Vācijas kanclers J.~Virts 1922.~gadā deklarēja: ``Poļus ir jālikvidē (\detxti{Polen muß erledigt werden}). Es neslēgšu nekādus līgumus, kuri varētu poļus stiprināt \citespace{} šajā ziņā es esmu pilnīgi vienisprātis ar militāristiem\dots{}''. J.~Virts par Poliju runāja kā par ``Versaļas bastardu''. Stājoties Vācijas prezidenta amatā, P.~f.~Hindenburgs savā programmas runā reihstāgā pieprasīja atdot Vācijai Dancigu, Poļu koridoru un Augšsilēziju. Poļu izlūkdienests arī zināja, ka Padomju Savienībā notiek Vācijas armijas jaunāko ieroču veidu izmēģinājumi, kuri varētu tikt lietoti arī pret Poliju. Šādu attīstības iespēju turēja prātā arī valdošie Polijas politiķi. Lai būtu gatavībā, viņi vēlējās turēt tautu kara gatavībā (\pltxti{naŕod pod bronią}).

Polijas atbalsta punkts bija Francija, kura Poliju uzlūkoja kā svarīgāko ``sanitārā kordona'' daļu. Taču abu valstu pūlēm stabilizēt stāvokli Austrumeiropā ļoti traucēja Polijas konflikti ar gandrīz visām kaimiņvalstīm un viņu izturēšanās pret nacionālajiem mazākumiem pašā Polijā, kas bieži tika apskatīta gan Eiropas valstu presē, gan Tautu Savienībā.

Šādos apstākļos Vācijas ārlietu ministram G.~Štrēzemanam ar gudru politiku pret Rietumvalstīm izdevās to valdībās mainīt neuzticību pret Vācijas revizionistiskajiem plāniem austrumos. Iezīmējusies Vācijas tuvināšanās ar Rietumu valstīm, kas Polijai draudēja ar starptautisku izolāciju, izsauca lielas rūpes Polijā.

1924.~gadā ar ASV starpniecību Vācija saskaņoja ar Lielbritāniju un Franciju tā saucamo ``Dauesa plānu'', kas paredzēja, ka Vācija maksās reparācijas, saņemot ASV kredītus. Tas veda pie Vācijas nostiprināšanās.

Otrs, vēl bīstamāks solis Polijai bija t.s. \strong{Lokarno vienošanos} noslēgšana 1925.~gada oktobrī. 1925.~gada 28.~februārī franču ziņu aģentūra ``\frtxti{Havas}'' pirmā paziņoja sensacionālo vēsti par to, ka Vācija ir slepeni piedāvājusi Francijai parakstīt drošības paktu, kurš apstiprinātu vācu-franču robežu neaizskaramību (Vācija ar to atzītu Elzasas-Lotringas piederību Francijai), abu valstu atteikšanos no kara kā nacionālas politikas līdzekļa un gatavību visus strīdus risināt mierīgā ceļā; paktu garantētu Lielbritānija un varbūt arī Itālija. Nekas šajā slepenajā piedāvājumā nebija teikts par garantijām Vācijas austrumu robežām ar Poliju un Čehoslovākiju, kuras abas bija Francijas partneres. Tas, par ko ziņoja ``\frtxti{Havas}'', bija izteikts G.~Štrezemana plānā, par kuru viņš vispirms bija slepeni informējis Londonu un pēc tam tikpat slepeni to piedāvājis Parīzei. Plāns iemiesoja G.~Štrezemana centienus pēc pakāpeniskas Versaļas līguma demontāžas (\detxti{Der Abbau}) sadarbošanās~--- nevis konfrontācijas~--- ar Lielbritāniju un Franciju ceļā, taču nekas nevarēja viņu piespiest piedāvāt līdzīgu paktu arī Polijai un Čehoslovākijai.

1925.~gada oktobrī Šveices kūrortā Lokarno tika noslēgts Reinas garantiju pakts (līgums starp Vāciju un Franciju, Lielbritāniju, Beļģiju un Itāliju), ar kuru Lielbritānija un Itālija garantēja Vācijas un Francijas un Vācijas un Beļģijas robežas, bet Vācijas un Polijas robežas~--- nē.

Ne velti līdz pat šim laikam daļa vēsturnieku uzskata, ka ar to Rietumvalstu vadītāji parādīja Vācijai iespējamo tās agresijas ceļu uz austrumiem. Labi zinādams, ka briti nav gatavi garantēt nekādu līdzīgu paktu austrumos (``Austrumu Lokarno''), Lielbritānijas ārlietu ministrs O.~Čemberlens tomēr liekulīgi paziņoja, ka jauni reģionāli pakti pārklās visu pasauli un galu galā tiks savienoti vispārējā globālā drošības sistēmā Tautu Savienības Padomes kontrolē. Taču nekas nemainījās britu stingrajos priekšstatos par to, ka viņu drošības robeža kontinentā ir Reina, nevis Visla. Neviena cita valsts līdz pat 1939.~gadam nesaņēma britu garantijas.

Lokarno vienošanās tikai apstiprināja Lielbritānijas un Francijas vēlmi risināt savas problēmas uz Eiropas ``mazo valstu'' rēķina.

Polija Lokarno līgumu noslēgšanu uztvēra kā savu interešu neievērošanu, ko asi pārmeta Rietumiem, sevišķi Francijai. Poļus biedēja tas, ka Francija, saņēmusi starptautiskas garantijas par savas robežas ar Vāciju neaizskaramību, koriģēs savas saistības pret Poliju un Čehoslovākiju. Pēc Francijas~--- Polijas 1921.~gada līguma Vācijas un Polijas konflikta gadījumā Francijai bija nekavējoties jāpalīdz Polijai, ja Vācija izraisītu neizprovocētu agresiju. Taču tas būtu pretrunā ar Reinas garantiju paktu, pēc tā noslēgšanas Polija zaudēja tiesības uz Francijas palīdzību Vācijas uzbrukuma gadījumā. Polijā to uztvēra kā Francijas atkāpšanos no 1921.~gadā sniegtā solījuma. Reinas pakts teorētiski pat ļāva tādām tā dalībniecēm kā Lielbritānija un Itālija nostāties Vācijas pusē.

Lokarno vienošanās būtiski izmainīja Polijas valsts statusu ārpolitikas jomā. No reģionālo starptautisko attiecību subjekta Polija kļuva par objektu Eiropas valstu politikā, tās drošību tagad garantēja nevis tiešs līgums ar Franciju, bet gan sarežģīta procedūra, ko paredzēja Tautu savienības statūti. Lokarno vienošanās simbolizēja Polijas ārpolitikas, tās politikas, kuras galvenais mērķis bija iegūt garantētas robežas ar Vāciju, neveiksmi.

Būtībā Lokarno pabeidza Francijas hegemonijas periodu kontinentā, aprēķini uz kuru bija Polijas ārpolitikas pamatā. Briesmas valsts robežai uztrauca visu poļu presi. Avīzē ``\pltxti{Słowo}'' (``Vārds'') 1925.~gada 4.~oktobrī rakstā ``Ceļs uz Lokarno'' tika norādīts, ka ``Žečpospolitai ir iespēja nenodzīvot līdz savai desmitgadei''.

Jāuzsver, ka PSRS uz konferenci nebija aicināta. Padomju diplomātija konferences rezultātus oficiāli gan novērtēja kā mēģinājumu radīt vienotu kapitālistisko valstu pretpadomju fronti. Taču būtībā PSRS Lokarno vienošanās sekas neuzskatīja par sev katastrofālām, turpināja sadarboties ar Vāciju, reizē šantažējot to ar iespēju tuvināties Polijai, kura bija uztraukusies par Rietumvalstu nelojalitāti.

Sākās Vācijas remilitarizācija, kam par ieganstu kalpoja it kā nepieciešamība būt gatavībā atvairīt iespējamos Polijas draudus. Ar katru gadu pēc Reinas garantiju pakta noslēgšanas pieauga nedrošības sajūta, it īpaši tādās valstīs kā Austrija, Čehoslovākija, Lietuva, daļēji arī citās Baltijas valstīs.

Polija par līgumu bija īpaši uztraukta, jo neilgi pirms tam Vācija centās uzsākt tiešas sarunas ar Poliju robežu jautājumā, prasot tai atdot Vācijai visu Pomerāniju (\pltxti{Pomorze}), kuras lielākā daļa jau no XVII gadsimta ietilpa Prūsijā, bet pārējā tajā iegāja pēc Napoleona sakāves 1815.~gadā. Tikai pēc Versaļas miera noteikumiem Polija ieguva dažus poļu apdzīvotus Pomerānijas rajonus~--- jau vairākkārt minēto t.s. ``Poļu koridoru''. Tā vietā Vācija tagad solīja nodrošināt Polijai brīvu tirdzniecību Baltijas jūras ostās, kā arī vienu dzelzceļa līniju uz šīm ostām. Pat Parīzē zināmās aprindās šis priekšlikums tika uzņemts labvēlīgi. (Faktiski šis pats projekts tikai citā variantā atdzima J.~f.~Ribentropa priekšlikumos Polijai 1938.~gada oktobrī).

Protams, šādi piedāvājumi izsauca satraukumu poļu sabiedrībā. Polijas Seima Kara lietu komisija pat pieprasīja, lai valdība nekavējoties izstrādātu likumprojektu par vispārēju iedzīvotāju sagatavošanu kara gadījumam. Polijas premjerministrs V.~Grabskis paziņoja, ka Polija Vācijas priekšlikumus var apspriest tikai kopā ar pārējām Versaļas līguma dalībvalstīm. Lielvalstis uz Poliju izdarīja spiedienu, lai panāktu tās piekrišanu garantiju paktam. Kaut poļu prese kopumā, izņemot atsevišķus nacionāldemokrātu izdevumus, izteica oficiālo, pozitīvo Lokarno līguma novērtējumu, J.~Pilsudskis un armijas Ģenerālštābs vērtēja to negatīvi. Šaubas par iespējamo Francijas palīdzību Polijai kara gadījumā izteica Seima maršals M.~Ratajs. Polija jutās spiesta meklēt jaunus sabiedrotos un citas drošības garantijas. 1926.~gada pavasarī Polijas ārlietu ministrija pat izsūtīja visiem Polijas sūtņiem ārzemēs speciālu cirkulāru, kurā bija tieši pateikts, ka ``var uzliesmot atklāts karš''. To, ka Polijas valstsvīru rūpes nebija bez pamata, pierāda mūsdienās zināmie fakti, ka pirmie Vācijas reihsvēra (bruņoto spēku) plāni par ``poļu koridora'' likvidēšanu, Augšsilēzijas atgriešanu Vācijas sastāvā tapa jau 1925.~gada nogalē, pāris mēnešu pēc Lokarno līguma noslēgšanas. Vācijas kanclers un ārlietu ministrs G.~Štrēzemans Vācijas sūtnim Londonā rakstīja, ka poļu jautājuma atrisināšana ir ``iespējams, pats svarīgākais Eiropas politikas jautājums vispār''.

Polija panāca vienīgi arbitrāžas līguma noslēgšanu ar Vāciju. Abas apņēmās risināt strīdīgos jautājumus bez militāra spēka pielietošanas, taču G.~Štrēzemans uzsvēra cerību uz savstarpējo robežu revīziju par labu Vācijai. 1927.~gadā Vācija tika uzņemta Tautu Savienībā, ieguva pastāvīgu vietu tās Padomē. Polijai turpretī bija tikai t.s. ``puspastāvīgā'' vieta, tā bija Padomē jāievēl. Tiesa, ievēlēšanu skaits nebija ierobežots. Vācijas panākumi starptautiskajā arēnā, jūtīgi aizskāra poļu nacionālo lepnumu un lāva apšaubīt valsts integritāti nākotnē.

1927.~gada septembrī Tautu Savienība, lai arī mīkstinātā formā, pieņēma Polijas iesniegto rezolūcijas projektu, kurš aizliedza izmantot karu kā starptautisku konfliktu risināšanas līdzekli. Polijas valdība to vērtēja kā ievērojamu savu panākumu, jo pieņemtajai deklarācijai bija jāaizsargā Polija no Vācijas uzbrukuma.

Reizē tai laikā pārspīlēts ``vācu briesmu'' tēlojums ļāva Polijas valdībai viegli panākt militāro izdevumu palielināšanu. Tā, 1929.~gadā tie sasniedza 38\% no valsts budžeta. Salīdzinājumam, Francijas budžetā šai laikā tie sastādīja 25\%. Tiesa, tai pat laikā poļu sabiedrībā bija arī bravūrīgi noskaņojumi. Piemēram, 1927. gadā Toruņā iznāca anonīma autora t. s. ``nākotnes'' [fantastisks] romāns ``1975'', kurā bija visos sīkumos aprakstīts poļu maršs uz Berlīni. 1929.~gadā vienas no poļu nacionālistiskajām organizācijām ``\pltxti{Liga Mocarstwowa}'' (Lielvalstiskās attīstības līga) izdevums ``\pltxti{Mocarstwowiec}'' (Lielvalstnieks) ``par iespējamo Polijas sadursmi ar Vāciju rakstīja: ``Karā pret vāciešiem nebūs gūstekņu un vietas cilvēciskām jūtām. Mēs pārsteigsim pasauli ar neparastiem, cilvēka mēru pārsniedzošiem asiņainiem upuriem, ar kuriem mēs nomazgāsim kaunu par apvainojumiem, ko mēs esam piedzīvojuši no vāciešiem. Mums ir nepieciešams Polijas armijā radīt nesamierināmu, līdz nežēlībai kāpinātu cīņas garu.''

Polijai draudošā izolācija spieda to meklēt attiecību uzlabošanos ar PSRS un Čehoslovākiju. Profesors K.~Počs atzīmējis, ka ``sanitārā kordona'' valstu centieni piemēroties Lokarno sistēmai neieguva garantijas līguma formu. Šķērslis tam bija Vācijas pretdarbība, Anglijas atturība un pašu ``sanitārā kordona'' valstu pretrunas. Šo pašu iemeslu dēļ arī franču un poļu diplomātijai neizdevās realizēt nevienu no Polijas-Baltijas valstu savienības projektu variantiem. Savukārt profesors A.~Stranga norādījis, ka garantiju jautājums ``kā reti kurš, izgaismoja Polijas vājumu''.

Polija traucēja reālas Baltijas valstu savienības izveidi. 1924.~gada jūnijā Latvijas sūtnis Polijā M.~Nukša uz Latviju ziņoja: ``\dots{}man liekas, ka pat tajā gadījumā, ja Polijai nebūtu bijušas tādas asas attiecības ar Lietuvu, kā tagad, tā būtu strādājusi pret triju mazo Baltijas valstu savienības ideju. Polija, domājams, arvien vēlēsies redzēt Lietuvu izolētā stāvoklī, uzskatīs to par savas nākamās ekspansijas objektu; to ietilpinot Lielajā Polijā, pēdējā taptu par īstu lielvalsti ar ostām, floti un noteikšanu pār Baltijas jūru. Ja tā nedomā visi poļi, tad tomēr liela daļa, klusībā~--- varbūt pat vairākums. \citespace{} Polija nevēlas tiltu starp Krieviju un Vāciju pār Baltijas valstīm, bet gan vēlas tādu pār Poliju''. Viņai labprāt patiktu kontrolēt sauszemes satiksmi starp kaimiņiem austrumos un rietumos.\citespace{} Visa Polijas rūpniecība ražo tādas kvalitātes preces, ka tā var cerēt uz noņēmējiem vienīgi Krievijā un Balkānos, mazā mērā arī Baltijas valstīs.''

Jau 1926.~gadā Polijas-Lietuvas konflikts tālāk saasinājās, jo jūnijā Lietuvas valdība paziņoja, ka starp Lietuvu un Poliju nebūs nekādu attiecību, iekams Viļņa netiks atdota Lietuvai. Lietuvai šāda pret tuvējo un galveno ienaidnieku~--- Poliju vērsta politika bija gluži dabīga. Rietumvalstis gan izdarīja spiedienu uz Lietuvu, lai panāktu tās piekāpšanos strīdā ar Poliju, taču 20.~gadu otrajā pusē tas nedeva rezultātus.

Īpaši Polijas~--- Lietuvas attiecības pasliktinājās pēc autoritārā apvērsuma Lietuvā 1926.~gada decembrī, pēc kura pie varas Lietuvā nāca antipoliskās un provāciskās politikas atbalstītāji. 1927.~gada otrajā pusē abu valstu attiecībās iestājās dziļa krīze. Cenšoties pacelt autoritārā režīma prestižu, Lietuvas valdība veica nacionālistiskus pasākumus pret poļu skolām, uz ko J.~Pilsudskis atbildēja ar vēl plašākām represijām pret lietuviešu skolām Viļņas apgabalā. Radās nopietnas poļu~--- lietuviešu kara briesmas. Lietuva Tautu Savienībā apvainoja Poliju agresīvos nodomos un Polijai bija starptautiski jātaisnojas. J.~Pilsudskis intervijā Polijas Telegrāfa aģentūrai rupji uzbruka Lietuvas Ministru prezidenta A.~Voldemāra personai. Poļu maršals pavisam nediplomātiski paziņoja: ``Lietuvas valdības priekšsēdētājs Voldemāra kungs modina manī bažas sava prāta stāvokļa dēļ: manuprāt, viņš ir nepieskaitāms. Uzspiestajās sarunās ar mums viņš atzina, ka atrodas kara stāvoklī ar mums un prasīja, lai Polija par karastāvokļa nomaiņu viņam ar kaut ko samaksātu, tā, it kā mēs šo karu būtu zaudējuši. Var teikt, ka viņš ir vienkārši bezkaunīgs. \citespace{} Kad es piebildīšu, ka šis lietuviešu superpatriots ilgāku laiku bija krievs, pēc tam vācietis un beidzot lietuvietis, katram ir jāatzīst, ka mums ir darīšana ar gandrīz klīnisku gadījumu''. Polija centās panākt Francijas un Lielbritānijas atbalstu. Taču PSRS noteikti iestājās par Lietuvas neatkarības saglabāšanu, reizē iesakot tai pārtraukt kara stāvokli ar Poliju, kādā tā formāli atradās kopš 1920.~gada. Pret Polijas nodomiem pakļaut Lietuvu uzstājās arī Vācija un poļu nacionālisti atlika savu agresīvo plānu īstenošanu. 1927.~gada 9.~decembrī starptautiska spiediena rezultātā J.~Pilsudskis sēdās pie sarunu galda ar it kā ``nepieskaitāmo'' Lietuvas valstsvīru. Ženēvā notika J.~Pilsudska un A.~Voldemāra tikšanās, kurā pēdējais bija spiests paziņot, ka starp abām valstīm nepastāv kara stāvoklis. Lietuvas~--- Polijas konflikta apspriešana Tautu Savienības decembra sesijā un abu diktatoru~--- J.~Pilsudska un A.~Voldemāra tikšanās mazināja saspīlējumu abu valstu starpā. Lietuva paziņoja par kara stāvokļa pārtraukšanu. Taču A.~Voldemārs nepārstāja uzstāt, ka Lietuvai jābūt neatkarīgai valstij ar galvaspilsētu Viļņā. Tiesa, Rietumvalstu spiediena rezultātā 1929.~gada septembrī A~Voldemārs bija spiests atstāt premjerministra posteni, taču poļu--lietuviešu konflikts palika neatrisināts. Lietuvā plaši izplatījās ``viltīgo poļu'' stereotips.

Vienlaicīgi Rietumvalstīs tika apspriesti arī dažādi kompensācijas varianti, ar kuru palīdzību varētu atrisināt Viļņas piederības jautājumu, kā arī jautājumu par papildus izeju Polijas pie jūras, kas savukārt ļautu nokārtot t.s. ``Poļu koridora'' problēmu un Polijas --Vācijas attiecības kopumā. Profesors K.~Počs parādījis, ka arī tādās Vācijas apdraudētās valstīs kā Polija un Čehoslovākija tika apspriesti dažādi apmaiņas un kompensāciju varianti, protams, uz citu valstu rēķina. Tā, Polija bija gatava atbalstīt t.s. ``anšlusu'' (vācu \detxti{Anschluβ}~--- pievienošana, savienība~--- šeit domāta Austrijas pievienošana Vācijai), ja Vācija atteiktos no Dancigas (Gdaņskas). Savukārt Čehoslovākija centās novirzīt Vācijas uzmanību uz Dancigu.

1928.~gada beigās J.~Pilsudskis un Čehoslovākijas ārlietu ministrs, vēlākais valsts prezidents E.~Benešs apmainījās ar ``komplimentiem'' abām valstīm. J.~Pilsudskis paziņoja: ``Čehoslovākija ir mākslīgi radīta, zemes gabals, kuru es nosauktu \detxti{Niemandsland}~--- nevienam nepiederoša zeme\dots{}'' (Šo atziņu pēc J.~Pilsudska nāves pārņēma arī Polijas ārlietu ministrs J.~Beks). E.~Benešs nepalika atbildi parādā, pasludinot: ``Ar Poliju Čehoslovākija nevar saistīties, jo tā ir Centrālās Eiropas Balkāni \citespace{} Galu galā, Polija ir valsts uz sezonu.''

Arī pret Latviju Polijai bija pretenzijas. 1926.~gada 20.~maijā Latvijas sūtnis Polijā M.~Nukša ziņoja: ``Protams, ka tie poļu muižnieki, kuru īpašums atrodas Ilūkstes apriņķa sešos pagastos \citespace{} ved, vismaz šeit, propagandu par šo pagastu pievienošanu Polijai, jo tas ir viņu interesēs.''

Tikai 1929.~gada 12.~februārī Rīgā tika parakstīta Latvijas un Polijas vienošanās par savstarpēju pastāvošās robežas atzīšanu, ar ko Polija oficiāli atteicās no teritoriālām pretenzijām uz Ilūkstes apriņķa Grīvas pilsētu un 6 pagastiem, kas bija iekļauti Latvijā 1920.~gada jūlijā, Polijas armijai poļu-padomju kara gaitā atkāpjoties. Taču, kā norādījis latviešu vēsturnieks Alberts~Varslavāns, ar savstarpēju robežu atzīšanu Latvija faktiski uzskatīja par likumīgu Viļņas apgabala aneksiju, ko bija veikusi Polija. Līgumam bija arī konfidenciāls (no sabiedrības slēpts) pielikums, kurā Latvija apņēmās izmaksāt agrārās reformas rezultātā īpašumus zaudējušajiem Polijas pilsoņiem finansiālu kompensāciju par 250~tūkstošiem ha poļu muižniekiem atsavināto zemi 5~miljonu zelta Ls apjomā jeb 30~Ls par hektāru. Polija toties atteicās izvirzīt vai atbalstīt jebkādas jaunas muižnieku prasības. Samaksa par vienu hektāru tiešām bija ļoti neliela, taču galvenais bija pats šis princips, kurš neatbilda Latvijā veiktās agrārās reformas saturam. Latviešu vēsturnieks A.~Stranga pamatoti raksta: ``Kompensācijas poļu muižniekiem (tiesa, ne visiem, bet tikai Polijas pavalstniekiem, kuri vienlaikus bija arī lielākie zemes īpašnieki) no tiesiskā viedokļa bija netaisnīgs izņēmums, radies no [Latvijas valdības] vēlēšanās uzturēt labas attiecības ar Poliju, kura~--- kā jau vienmēr tai bija raksturīgi~--- nereti pat aroganti iejaucās Latvijas iekšējās lietās vai pat apšaubīja Latvijas robežu integritāti.'' Latvija nokārtoja maksājumus līdz 1937.~gada 18.~novembrim.

Kā konstatējis Ē.~Jēkabsons, Latvijā dzīvojošie poļi nebija mazāk lojāli pret Latvijas valsti kā krievi, vācieši, ebreji, taču viņu īpašais pagātnes statuss Latgalē un Ilūkstes apriņķī joprojām uzturēja spriedzi starp nacionāli noskaņoto poļu sabiedrības daļu un varas iestādēm. 1938.~gada rudenī pēc Latvijas politpārvaldes ziņām vairāki vietējās Poļu jaunatnes savienības darbinieki pieļāva iespēju, ka kara gadījumā Polija atņems Latvijai Daugavpili un varbūt pat visu Ilūkstes apriņķi. Šādas runas pastiprinājās pēc tam, kad Polija ieņēma Čehoslovākijai piederošo Cešinas apgabalu. Ziņu patiesumu apstiprināja Polijas konsuls Daugavpilī. Tātad faktiski šādas baumas radīja Polijas agresīvā politika. Polijas sūtniecība Rīgā bija spiesta pielikt lielas pūles, lai pārliecinātu Latvijas valdību par savas valsts miermīlīgajiem nodomiem pret Latviju.

Arī \strong{30.~gados} Polijas ārpolitikā bija daudz sarežģījumu. Pasaules ekonomikas krīze papildus nopietnām ekonomiskām un sociālām grūtībām izraisīja Polijai bīstamas pārmaiņas spēku samērā Eiropā. Francijai, kas bija viens no galvenajiem Versaļas sistēmas balstiem un drošības garants saviem sabiedrotajiem Austrumeiropā, arvien grūtāk bija tikt galā ar Berlīnes uzvedības kontroli. Šajā sakarā Parīzē sāka apsvērt iespēju tuvināties Maskavai, kura sāka pārvērsties par jaunu spēka centru kontinenta austrumos. Tas grāva veco, gan pret Vāciju, gan PSRS vērsto Polijas ārpolitisko stratēģiju, būvētu uz aprēķiniem par Francijas hegemoniju Eiropā un Versaļas sistēmas stabilitāti. 1931.~gadā, sakarā ar Japānas agresiju Ķīnā, mainījās PSRS attieksme pret kaimiņvalstīm rietumos, arī pret Poliju. 1932.~gadā PSRS parakstīja virkni neuzbrukšanas līgumu ar Somiju, Latviju, Igauniju.

1931.~gada otrajā pusē arī J.~Pilsudskis sāka koriģēt valsts ārpolitiku, lai atbrīvotos no Francijas aizbildniecības un pārvērstu Poliju par vadošo spēku Austrumeiropā. (Padomju vēsturnieki uzsvēra, ka sakarā ar izmaiņām ārējā un iekšējā stāvoklī esot mainījusies tieši Polijas attieksme un tā bijusi spiesta veikt dažādus manevrus, lai uz laiku uzlabotu attiecības ar PSRS.) Sarunas par līgumu tika sāktas Maskavā 1931.~gadā. 1932.~gada 25.jūlijā tika parakstīts Polijas un PSRS līgums par neuzbrukšanu (poļu \pltxti{Pakt o nieagresji Polska~--- ZSRR}). Abas puses atzina otras puses suverenitāti, kopējās robežas un teritoriālo integritāti. PSRS un Polijas līguma pirmajā pantā tās paziņoja, ka ir atteikušās no kara kā nacionālās politikas instrumenta, apņēmās ``savstarpēji atturēties no jebkādām agresīvām darbībām vai uzbrukuma viena otrai, gan patstāvīgi, gan kopā ar citām valstīm''. Bija noteikts, ka par šajā pantā noteiktajām saistībām pretējām darbībām tiks atzīts ``jebkurš vardarbības akts, kas vērsts pret Otras Līgumu slēdzošās puses teritoriālo integritāti un neaizskaramību vai politisko neatkarību pat gadījumā, ja šīs darbības būtu veiktas bez kara pieteikšanas un izvairoties no visām tā izpausmēm''. Trešais līguma pants noteica, ka ``katra no Līgumu slēdzošajām pusēm apņemas nepiedalīties nekādos līgumos, kas no agresijas redzes viedokļa ir Otrai pusei acīmredzami naidīgi.'' Tādejādi Polija vēlreiz juridiski nostiprināja savus teritoriālos ieguvumus, fiksētus 1921.~gada Rīgas miera līgumā.

Līgums sākotnēji tika noslēgts uz trīs gadiem ar tā automātisku pagarināšanu uz nākamajiem diviem gadiem, ja neviena no pusēm nebrīdināja par tā denonsāciju sešus mēnešus pirms līguma termiņa beigām. Taču jau 1934.~gada 5.~maijā līgums tika pagarināts līdz 1945.~gada 31.~decembrim.

Ņemot vērā notiekošo muitas karu ar Vāciju, vispārēju pakāpenisku atteikšanos no Versaļas sistēmas un saišu ar Franciju vājināšanos, Polijā šis nolīgums tika uzskatīts par tās diplomātijas lielu panākumu. Līgums, kaut sabojāja Polijas draudzīgās attiecības ar Rumāniju, kopumā tomēr nostiprināja Polijas pozīcijas, ļāva tai asi reaģēt uz Ā.~Hitlera nākšanu pie varas 1933.~gadā. Tas arī sekmēja poļu un vācu sarunas, kuru rezultātā pēc 18~mēnešiem tika noslēgts Polijas un Vācijas līgums par neuzbrukšanu un par kurām turpmāk tiks runāts sīkāk. (Padomju vēsturniece V.~Parsadanova norādīja, ka, kā vēlāk izrādījies, Polijai šis līgums bija svarīgs kā trumpis tās tālākajās sarunās ar hitleriešiem.) Šī līguma saturs arī jāatceras, kad tiks runāts par Polijas agresiju pret Čehoslovākiju 1938.~gada rudenī un PSRS agresiju pret Poliju 1939.~gada septembrī.

Tikai poļu komunisti neuzbrukšanas līguma ar PSRS noslēgšanu vērtēja kā poļu imperiālisma dubultspēles izpausmi. Partijas ģenerālsekretārs J.~Ļenskis rakstīja: ``Ir jāatmasko fašistu un sociālfašistu [domāti sociālisti~--- V.Š.] attieksme pret paktu par neuzbrukšanu''. Tik tiešām, Varšava neapturēja savu atbalstu pretpadomju organizācijām, t.~s. kluba ``\pltxti{Prometeusz}'', kustībai, tikai vairāk centās to vairāk paturēt noslēpumā.

Zīmīga bija padomju puses attieksme pret līgumu. 1931.~gada beigās notikušajā J.~Staļina sarunā ar vācu rakstnieku E.~Ludvigu otrais izteica vācu politiķu bažas par to, ka līguma ar Poliju noslēgšana nozīmēs PSRS draudzīgo attiecību ar Vāciju pārtraukšanu, jo apstiprinās Versaļas līguma koncepciju, kurš sāpīgi aizskāra Vācijas intereses. J.~Staļins savukārt centās apstrīdēt šādu viedokli un paziņoja: ``Mēs vienmēr apgalvojām, ka esam gatavi noslēgt līgumu par neuzbrukšanu ar jebkuru valsti. \citespace{} Brīdī, kad poļi paziņoja par savu gatavību vest ar mums sarunas par neuzbrukšanas paktu, mēs, dabīgi, piekritām un uzsākām sarunas.'' Reizē padomju diktators centās nomierināt vācu politiķus: ``Vai tā ir Versaļas sistēmas atzīšana? Nē. Vai, varbūt, tā ir robežu garantija? Nē. Mēs nekad neesam bijuši Polijas garanti un nekad tādi nekļūsim, tāpat kā Polija nebija un nebūs mūsu robežu garants. Mūsu draudzīgās attiecības ar Vāciju paliks tādas pat, kādas bija līdz šim. Tāda ir mana stingra pārliecība.''

Kaut uz 1932.~gada neuzbrukšanas līguma pamata 1934.~gada jūlijā-augustā notika padomju lidotāju eskadriļas vizīte Polijā, Polijas karakuģi jūlijā apmeklēja Ļeņingradu, bet septembra sākumā padomju karakuģi~--- Poliju, arī PSRS kara un jūras lietu tautas komisārs K.~Vorošilovs sarunā ar Vācijas sūtni H.~f.~Dirksenu apstiprināja, ka līgums ar Poliju ``neizsauks izmaiņas Padomju Savienības draudzīgajās attiecībās ar Vāciju.''

Noteikti gan jāpasaka, ka PSRS vadītāji arī nebija atmetuši domas par citu zemju, arī Polijas, revolucionarizēšanu.

Pie Kominternes Izpildu komitejas 30.~gadu sākumā tika organizēti militāri kursi. Tajos ārzemju komunistisko partiju pārstāvji apguva bruņotas sacelšanās organizēšanas iemaņas, ielu un partizāņu cīņas taktiku. 1931.--1932.~gadā kursos mācījās 194~cilvēki, no tiem 56~poļi. Daļa no tiem tika nosūtīta pagrīdes darbā uz dzimteni. 1934.~gada februārī Polijas Komunistiskā partija saņēma Kominternes Izpildu komitejas Politsekretariāta direktīvas aktivizēt darbu Polijas bruņotajos spēkos, pārliecināt karavīrus solidarizēties ar strādnieku un zemnieku kustību, atteikties no soda funkciju izpildes pret tiem. Īpaši bija uzsvērta nepieciešamība karavīru vidū propagandēt Padomju varas nodibināšanas lozungu. Pēc būtības Kominternes Izpildu komiteja deva poļu komunistiem direktīvu gatavot armiju cīņai par pastāvošā režīma gāšanu. Iznāca, ka Kominterne, kura tikai formāli bija patstāvīga, bet faktiski pildīja VK(b)P vadības norādes, darbojās pretēji PSRS valdības ar Poliju noslēgtajiem oficiālajiem līgumiem. Var pieņemt, ka Polijas valdība, ja arī nebija pilnībā informēta par Kominternes aktivitātēm, tomēr tās darbību saistīja ar PSRS, tāpēc nevarēja tai uzticēties.

Tādejādi, neraugoties uz noslēgto neuzbrukšanas līgumu, arī 30.~gados turpinājās 20.~gadu savstarpējais naidīgums, tikai piesegts ar oficiālās diplomātijas plīvuru.

Lai minam tikai vienu piemēru, kurš uz citu, svarīgāku notikumu fona var likties kā gauži nenozīmīgs incidents, taču tas raksturoja poļu-padomju attiecību gaisotni. 1938.~gada vasarā Polijas-PSRS attiecību pasliktināšanās atmosfērā (iespējams, ka situāciju noteica sasprindzinājums Eiropā sakarā ar Vācijas pretenzijām Čehoslovākijai) pret vairākiem padomju diplomātiem un viņu ģimenes locekļiem Polijā tika vērstas policijiskas akcijas. Sūtniecības apmeklētājus, izejot no tās, aizturēja un veda uz policijas komisariātu. Arī diplomātus, viņu ģimenes locekļus aizturēja, pratināja, dienas gaitās pavadīja policija, sūtniecības preses birojam nepiegādāja poļu avīzes, nepielaida strādāt poļu apkalpojošo personālu utt. Padomju puse atrada savu atbildes veidu. PSRS sūtniecība iesniedza notu Polijas Ārlietu ministrijai, bet 28~poļu darbiniekus Padomju Savienībā, kuriem nebija diplomātiskā statusa, arestēja, Polijas sūtniecībai atņēma transportu, kas sūtnim pat neļāva apmeklēt PSRS Ārlietu ministriju, radīja virkni citu neērtību. Augusta vidū incidents sasniedza apogeju, pēc tam tika izsmelts un diplomātiskā dzīve ieguva ierasto veidu.

Joprojām sarežģītas palika Polijas attiecības ar Vāciju. 1930.~gada 10.~augustā tā saucamo okupēto Vācijas teritoriju ministrs G.~Treviranuss, uzstājoties pie Reihstāga sapulcējušos nacionālistisko demonstrantu priekšā, paziņoja, ka pastāvošā Vācijas un Polijas robeža dara neiespējamu mieru starp abām valstīm. Viņš ilgi runāja par ``nedziestošo Vācijas brūci austrumu flangā'', par ``zaudētajām, bet nākotnē Vācijai atgriežamajām zemēm Austrumos''. 14.~augustā Polijas ārlietu ministrs A.~Zaleskis iesniedza Vācijas kancleram H.~Bruningam oficiālu protestu sakarā ar G.~Treviranusa runu. Kaut tika publicēts oficiāls valdības komunikē, kurā bija apgalvots, ka valdība neko nezinājusi par gatavojamo sava locekļa uzstāšanos, jau 16.~augustā G.~Treviranuss, uzstājoties mītiņā Ķelnē paziņoja, ka viņam nav ``nekāda iemesla'' atsaukt kaut ko no iepriekšējā runā teiktā. Vācijas ministra antipoliskā akcija noveda gandrīz pie diplomātisko attiecību saraušanas starp abām valstīm.

Simpātijas Eiropā ne vienmēr atradās Polijas pusē. Itālijas diktators B.~Mussolini intervijā amerikāņu avīzes ``\entxti{Saturday Review}'' (``Sestdienas apskats'') korespondentam paziņoja; ``Polija labi darītu, ja mainītu savu pozu, lai neapdraudētu savu eksistenci''. B.~Mussolini vienā no savām publiskajām runām arī izteicās, ka Vācijai acīmredzot nebūs grūti sev labvēlīgi noregulēt ``poļu koridora'' un Augšsilēzijas jautājumu, un no savas puses pat ieteica Ā.~Hitleram to darīt, bet toties neko neuzsākt pret Austriju. Šis Itālijas diktatora paziņojums pēc Čehoslovākijas sūtņa Lielbritānijā J.~Masarika vārdiem pamatīgi sadusmoja Polijas diktatoru J.~Pilsudski, kurš drīzumā savukārt paziņoja, ka Poliju Austrija neinteresē un Vācija, ja tā vēlas, var to ``savākt'', kas savukārt izraisīja B.~Mussolini dusmas. D.~Loids Džordžs ekonomiskās krīzes apstākļos gan runāja par piecām Elzasām-Lotringām [apgabals, kura dēļ Francija iesaistījās Pirmajā pasaules karā]: ``poļu koridoru'', Augšsilēziju, Viļņu, Rietumukrainu un Rietumbaltkrieviju. Bet arī viņš ``poļu koridoru'' nosauca kā pirmo.

Starp citu, toreiz sabiedrība Vācijā uzskatīja par iespējamu Polijas preventīvu uzbrukumu. 1932.~gadā ar pseidonīmu Hans Nitrams iznāca H.~Martina grāmata ``\detxti{Achtung! Ostmarkenrundfunk! Polnische Truppen haben heute nacht die ostpreussische Grenze überschritten.}'' (Uzmanību! Austrumapgabala radio! Poļu karaspēks šodien naktī šķērsoja Austrumprūsijas robežu.) Vācijas Kara ministrijā (\detxti{Reichswehrministerium}), kura deva svētību grāmatas iznākšanai, uzskatīja to par pienesumu tautas ``garīgai mobilizācijai``. Kad grāmata parādījās Austrumprūsijas un Polijas robežapgabalu grāmatveikalu skatlogos, tā tika uztverta kā reportāža par patiesu notikumu, vietējo vācu iedzīvotāju vidū izplatījās panika, turpretī daļa poļu tautības iedzīvotāju piedzīvoja eiforiju. Toreiz Polijas iebrukums Vācijā likās daudz reālistiskāks nekā sešus gadus vēlāk ASV O.~Velsa pēc H.~Velsa romāna ``\entxti{The War of the Worlds}'', (Pasauļu karš) sagatavotā radiopārraide par marsiešu iebrukumu, kura tāpat izsauca klausītāju paniku.

Arī Polijas valdības izturēšanās Dancigas jautājumā deva Vācijai pamatu apsūdzēt Poliju agresijā. Tā Polijas ārlietu ministrs A.~Zaleskis Dancigas pilsētas Senāta priekšsēdētājam paziņoja: ``Dancigas jautājumu var izlemt tikai poļu armijas korpuss''. J.~Pilsudskis šai laikā īpaši centās nostiprināt attiecības ar Franciju, demonstrējot nepiekāpību teritoriālajos jautājumos, gatavību aizstāvēt poļu intereses Dancigā, ar to riskējot nonākt militārā konfliktā ar Vāciju. Kad Dancigas pilsētas Senāts 1932.~gadā atteicās pagarināt konvenciju, kas ļāva Polijas flotei izmantot Dancigas ostu kā savu, 14.~(citi dati~--- 15.) jūnijā ostas reidā, pārkāpjot starptautisko tiesību normas, ieradās Polijas mīnkuģis ``\pltxti{Wicher}'' (Viesulis) ar pavēli apšaudīt pilsētā tuvāko sabiedrisko ēku, ja tiks aizskarts Polijas karogs. Par laimi līdz konfrontācijai nenonāca, 15.~jūnijā Senāts atcēla savu lēmumu. Tomēr Dancigas vadība griezās ar sūdzību Tautu Savienībā, kura nosodīja Polijas rīcību.

1932.~gada 14.~septembrī Vācijas delegācija demonstratīvi atstāja Ženēvas atbruņošanās konferenci, jo tā atteicās atzīt Vācijas vienlīdzīgas tiesības uz bruņošanos. Citas Rietumu lielvalstis 5 valstu (Lielbritānijas, Francijas, Itālijas, ASV un Vācijas) sanāksmē 11.~decembrī bija spiestas piekrist Vācijas tiesībām uz ``vienlīdzību bruņošanās jomā'' un deklarēt drošību visām Eiropas valstīm. Tas vēl vairāk padarīja nedrošu Polijas stāvokli. Kā rakstīja poļu vēsturnieks J.~Krasuskis ``Ar šo līdzekli Veimāras republikas ārpolitika sasniedza gandrīz visus savus mērķus: Augšreinas atbrīvošanu, reparāciju likvidāciju un vienlīdzību. Lai pievienotu Austriju un atņemtu Polijai Pomorji un Silēziju republikai vienkārši nepietika laika.''

Draudīgs izaicinājums Varšavai kļuva B.~Mussolīni 1932.~gada oktobrī formulētā ideja par četru valstu Itālijas, Francijas, Lielbritānijas un Vācijas pakta noslēgšanu un to sava veida direktorāta izveidi Eiropā, kas uzņemtos pilnu atbildību par mieru saglabāšanu kontinentā, tostarp arī nevardarbīgām robežu maiņām. J.~Pilsudski sāpīgi aizskāra Polijas prasības pēc lielvalsts statusa, kam bija arī liela loma tās iekšpolitiskā stāvokļa stabilizēšanā, ignorēšana no nosaukto valstu puses.

Uz ``četru pakta'' priekšlikumu un 11.~decembra deklarācijas pieņemšanu J.~Pilsudskis reaģēja ar steidzamu valdības un armijas vadības sanāksmes sasaukšanu Varšavā. Pēc franču žurnālistes Ž.~Tabui liecības tika apspriests tikai viens jautājums~--- vai var un kādā mērā paļauties uz savienību ar Franciju un pieņemts secinājums: ``Nē, nepieciešams pacensties pēc iespējas ātrāk mainīt mūsu politiku tuvināšanās ar Vāciju virzienā''.

Faktiski šī politikas maiņa jau bija ievadīta. 1932.~gada 2.~novembrī Polijas ārlietu ministru A.~Zaleski, kura vadībā Polijas diplomātija bija pretdarbojusies Vācijas plāniem, J.~Pilsudskis nomainīja ar citu politiķi~--- pulkvedi J.~Beku, kam vajadzēja signalizēt Vācijai par Polijas ārpolitikas maiņu. Frankofīla A.~Zaleska, kuram starptautiskajā arēnā bija augsts personīgais prestižs, nomaiņa manāmi uztrauca politiski-diplomātiskās aprindas vairākumā Eiropas valstu. Tāpēc Polijas pārstāvjiem ārzemēs tika doti norādījumi skaidrot, ka iepriekšējais ministrs atstājis savu amatu veselības apsvērumu dēļ, sakarā ar sirds slimību. Tiesa, A.~Zaleskis šķīrās no dzīves pēc 44 gadiem 93~gadu vecumā\dots{}

Starp citu, vēlāk~--- 1939.~gadā padomju diktators J.~Staļins, tāpat nomainot ārlietu ministru M.~Ļitvinovu pret V.~Molotovu, tikai gāja J.~Pilsudska pēdās, signalizēdams Ā.~Hitleram par jaunu valsts kursu starptautiskajās attiecībās.

Ā.~Hitlers, 1933.~gada sākumā nākot pie varas (viņš kļuva par Vācijas kancleru 1933.~gada 30.~janvārī) jau 3.~februārī slepenas tikšanās laikā ar bruņoto spēku virspavēlniecības pārstāvjiem paziņoja, ka viņš izmantos saņemto varu, lai iegūtu Vācijai ``dzīves telpu'' (\detxti{Lebensraum}) austrumos. Taču reizē viņš bija gatavs iesaistīties sarunās ar Poliju, lai pacenstos to izmantot Vācijas ārpolitisko mērķu īstenošanas labā. Šai ziņā ir raksturīgs fakts, ko savās atmiņās min bijušais Vācijas sūtnis PSRS H.~f.~Dirksens. Sarunā ar viņu Ā.~Hitlers sapņaini ieminējies: ``Ja tikai mēs varētu vienoties ar Poliju! Pilsudskis ir vienīgais cilvēks, ar kuru to varētu īstenot!'' Diplomātiskajā darbībā Ā.~Hitlers pacentās īstenot savu nodomu. Vācija pagaidām nekādus nopietnus draudus Polijai neizrādīja, demonstrēja gatavību noslēgt ar to vienošanos. Tiesa, savienot prasību pēc ``poļu koridora'' un uzturēt labas attiecības ar Poliju bija sarežģīti.

Lai piesaistītu vācu nacionālistus, pret Poliju vērsti paziņojumi no Ā.~Hitlera un viņa līdzgaitnieku mutes nerimās. Britu avīzes ``\entxti{Sunday Express}'' korespondentam 12.~februārī Ā.~Hitlers paziņoja: ``Polijai ir jāatdod mums koridors un pēc iespējas ātrāk''. 1933.~gada marta vidū Ā.~Hitlers intervijā kādai angļu avīzei atkal izteica asu prasību paātrināt Vācijas-Polijas robežas revīziju. Tūlīt sekoja gan atsaukums, kur kanclera pretpoļu tonis bija ievērojami mīkstināts, tomēr par viņa agresivitāti šaubas neradās.

Uz Ā.~Hitlera runu oficiālu atbildi J.~Beks deva Seima ārlietu komisijas sēdē: ``Mūsu attieksme pret Vāciju būs tieši tāda pati kā Vācijas attieksme pret Poliju. Praksē šajā jomā viss ir atkarīgs vairāk no Berlīnes nekā no Varšavas.'' Šo paziņojumu jau iepriekš bija rūpīgi režisējis pats J.~Pilsudskis. J.~Beks vēlāk rakstīja: ``Maršals man ieteica neņemt verā Eiropas preses drudžaino uztraukumu un vācu jautājumu skart enerģiski, bet vienlaikus mierīgi, ievērojot mēra sajūtu vārdu izvēlē, lai neradītu iespaidu, ka nākotne sevī slēpj vienīgi negatīvas perspektīvas. Mūsu attieksmei pret Vāciju ir jābūt atkarīgai no Vācijas attieksmes pret mums.'' Kā secina krievu vēsturnieks G.~Matvejevs, J.~Pilsudskis uzskatīja, ka Ā.~Hitlers tuvākos 4--5 gadus būs spiests risināt valsts iekšējos jautājumus un viņa ārpolitika būs tāla no saukļa ``\detxti{Drang nach Osten}'' realizācijas. Šim viedoklim atbilda arī sākotnēji nacistiskā diktatora realizētā ārpolitika, sadarbības ar PSRS krasa ierobežošana, atteikšanās būvēt savu diplomātiju uz Lokarno līgumu un Vācijas 1926.~gada līguma ar PSRS pamata, kurus Varšavā vērtēja kā pretpoliskus. Šķita, ka Polijai ir radusies iespēja savu drošību garantēt nevis ar tāda nedrošā sabiedrotā kā Francija palīdzību, bet patstāvīgi realizējot politiku, kas ļautu Polijai balansēt starp Vāciju un PSRS, neļaujot tām tuvināties uz pretpoliska pamata.

Tomēr Ā.~Hitlera nākšana pie varas Polijā izsauca nopietnas bažas, nervozitāti, kuru tās valdītāji centās maskēt ar sava spēka demonstrāciju. 1933.~gada 1.~februārī J.~Beks Tautu savienības Padomes sesijā paziņoja, ka Polija turpmāk neatzīs ``Līgumu par nacionālo mazākumu aizsardzību'' (``\pltxti{Traktat Mniejszošciowy}''), nepiedalīsies sēdēs, kur tiks izskatītas vācu puses sūdzības par nacionālo mazākumu stāvokli Polijā. Kā atbilde uz to izskanēja Dancigas Senāta 16.~februāra lēmums uzteikt 1923.~gada līgumu ar Poliju par speciālo ostas policiju un pārņemt to savās rokās. (Pēc krievu vēsturnieka D.~Oficerova ziņām Vācijas armijas pavēlniecība februāra beigās~--- marta sākumā izstrādāja Dancigas sagrābšanas plānu.) Šajā situācijā J.~Pilsudskis atkal nolēma demonstrēt Polijas militāro spēku. Naktī no 5.~uz 6.~martu~--- nākamajā dienā pēc Reihstāga vēlēšanām Vācijā, kurās uzvarēja nacionālsociālisti, pretēji Versaļas miera noteikumiem poļu atbalsta punktā ``Vesterplate'' (\pltxti{Westerplatte}) pie Dancigas tika izcelts desants. J.~Pilsudskis bija licis uz laiku ar vairāk nekā 100~cilvēkiem pastiprināt tur esošo Polijas armijas kontingentu. Kā rakstīja poļu vēsturnieks M.~Vojcehovskis, ar to J.~Pilsudskis bija iecerējis paveikt vienlaikus divus darbus: ``izdarīt spiedienu uz Ā.~Hitleru un vienlaikus dot atbildi Rietumu valstīm, kuras bija labvēlīgas pret Vācijas revanšistiskajām prasībām’’. Pasaules presē izskanēja doma, ka Polija ar to esot pietuvinājusi degli pulvera mucai, kuras sprādziens draud ar pasaules karu. Asa Vācijas reakcija gan izpalika, taču Danciga griezās pie Tautu Savienības Augstā komisāra un pastiprināja savu t.s. palīgpoliciju. Tautu Savienības Padomē Ženēvā izvērsās asa diskusija. Tautu Savienības spiediena rezultātā Dancigas varas iestādes garantēja Polijas tiesības pilsētā, bet J.~Pilsudskis bija spiests izbeigt savu ``Vesterplates afēru'' un izvest papildus karaspēka kontingentu. Pēc tam maijā notikušajās Dancigas Senāta vēlēšanās pārliecinoši uzvarēja vācu nacionālsociālisti.

Martā J.~Pilsudskis ar uzticamu militāristu starpniecību piedāvāja Francijai kopīgu uzstāšanos pret Vāciju. Poļu politiķiem Vācijas-Polijas robežas revīzija šķita ļoti iespējama. Tāpēc J.~Pilsudskis demonstrēja gatavību kopdarbībā ar Franciju un ar Tautu Savienības piekrišanu okupēt Vācijai piederošos Augšsilēzijas un Austrumprūsijas apgabalus. Tiem vajadzētu kalpot par ķīlu, ka Ā.~Hitlers ievēros Versaļas miera līguma noteikumus. Tomēr J.~Pilsudska plāni sadūrās ar Francijas valdības neizlēmību, faktiski noraidošu attieksmi, tāpēc oficiālas saruna par šādu iespējamo rīcību pat netika uzsāktas. Arī Polijā konkrēti pasākumi, sagatavojot preventīvu karu pret Vāciju, netika veikti.

Pētnieku vidū izteikti dažādi viedokļi, vērtējot J.~Pilsudska nodomus. Piemēram, vācu vēsturnieks H.~Rooss uzskatīja, ka J.~Pilsudskis nedomāja vest īstu karu pret Vāciju, bet vēlējās veikt tikai ``policijisku akciju'', līdzīgu Francijas īstenotajai Rūras okupācijai 1923.~gadā.

Poļu vēsturnieku vidū pastāv vairāki viedokļi: ka J.~Pilsudska signāli Francijas virzienā 1) bija vērtējami vairāk kā brīdinājumi Vācijai ņemt vērā Poliju kā sarunu partneri, nostiprināt pēdējās pozīcijas paredzamajās abu pušu sarunās; 2) to uzdevums bija pārbaudīt Polijas savienības ar Franciju stingrību; 3) J.~Pilsudskim bija divi alternatīvi plāni attiecību attīstībai ar Vāciju, un viens no tiem arī vēlāk tika realizēts.

Padomju vēsturnieks D.~Kļimovskis atzina, ka Polijas diktatoram bija vairāki nodomi: pārbaudīt Polijas savienības ar Franciju stingrību, pārliecināt opozīciju Polijā, ka uz Franciju nevar droši paļauties, paradīt nacistiem, ka Polija var rīkoties enerģiski, ja tie ignorēs tās piedāvājumus sadarboties.

Mūsdienu krievu vēsturnieks D.~Oficerovs savukārt uzskata, ka Polijas diktatora piedāvājums Francijai nenozīmēja J.~Pilsudska gatavību preventīvam karam pret Vāciju, visticamāk viņš šādu iespēju pat nopietni neizskatīja, taču bija gatavs veidot plašu antihitlerisku fronti. Reizē minētais pētnieks Francijas atteikumu kopīgi vērsties pret Vāciju uzskata par 1933.~gada rudenī notikušās Polijas un Vācijas tuvināšanās cēloni.

Taču tā paša gada 7.~aprīlī Vācijas ārlietu ministrs barons K.~f.~Neirāts, uzstājoties valdības sēdē ar ziņojumu par ārpolitisko situāciju, paziņoja, ka valdības mērķis ir valsts austrumu robežas revīzija un saprašanās ar Poliju nav nedz iespējama, nedz arī vēlama. K.~f.~Neirāts Polijas darbības novērtēja tikai kā spēlēšanos ar domu par preventīvu karu pret Vāciju. Tāpēc grūti spriest, cik lielā mērā Vācija J.~Pilsudska draudus uztvēra nopietni.

Jau 1933.~gada 2.~maijā J.~Pilsudskis griezās pie Ā.~Hitlera ar ultimatīvu jautājumu, vai viņš atzīst Vācijas~--- Polijas robežu. Fīrers (vācu \detxti{Führer}~--- vadonis) atbildēja, ka Vācijas valdība negatavojas lauzt pastāvošos līgumus. Ā.~Hitlers vēlējās demonstrēt savu miermīlību, tai pat laikā izvēršot valsts iekšienē bruņošanos. Tas ļāva panākt Vācijas un Polijas taktisku tuvināšanos.

Bez Francijas noraidošās attieksmes pret preventīvu karu pret Vāciju J.~Pilsudska nostāju ietekmēja arī Rietumeiropas valstu attieksme pret B.~Mussolini jau 1932.~gadā izvirzīto un 1933.~gada 14.~martā apstiprināto t.s. ``četru pakta'' (Lielbritānijas, Francijas, Itālijas un Vācijas) projektu. 1933.~gada 15.~jūlijā Romā tika arī noslēgts Lielbritānijas, Francijas, Itālijas un Vācijas līgums jeb ``četru pakts'', kura oficiālais mērķis bija noregulēt dalībvalstu domstarpības, panākt to saskaņotu politiku visos strīdīgajos starptautiskajos jautājumos. Taču, apstiprinot valstu saistības pildīt Lokarno līgumus, ``četru pakts'' paredzēja Vācijas vienlīdzību bruņošanās jomā, jeb, citiem vārdiem, piešķīra Vācijai tiesības bruņoties. Pakts gan netika ratificēts, jo tā dalībvalstu vidū pretrunas tā arī netika pārvarētas, taču gan Polijai, gan PSRS, gan citām Austrumeiropas valstīm lika justies apdraudētām. Pakts pastiprināja Polijas vadības aizdomas, ka nosauktās ``lielās'' valstis krīzes gadījumā būs gatavas ziedot ``mazo'' valstu intereses. Pats aizvainojošākais poļiem bija tas, ka Polija kā reiz nebija pieskaitīta ``lielajām'' valstīm. Rezultātā Polija centās panākt drošību, noslēdzot ar Vāciju līgumu.

Tai pat 1933.~gada 14.~oktobrī Vācija galīgi atteicās turpināt dalību atbruņošanās konferences darbībā Ženēvā un pasludināja par savu izstāšanos no Tautu Savienības, taču Ā.~Hitlers deklarēja gatavību noslēgt līgumus par neuzbrukšanu ar citām valstīm.

Kā savā grāmatā ``\detxti{Gespräche mit Hitler}'' (Sarunas ar Hitleru, 1940) rakstīja toreizējais Dancigas Senāta priekšsēdētājs H.~Raušnings, kurš vēlāk izstājās no nacistu partijas un emigrēja, tad pat nacistiskās partijas virsotņu sanāksmē fīrers arī teicis, ka viņš parakstīs jebkurus līgumus par neuzbrukšanu un draudzību, taču nopietni par vienošanos ar Poliju nedomā līdz tam laikam, kamēr netiks likvidēti draudi Rietumos. Tiesa, H.~Raušnings tikās ar Ā.~Hitleru tikai 4~reizes un lielākoties citu cilvēku klātbūtnē, tāpēc virkne vēsturnieku uzskata, ka šajā grāmatā citētie Ā.~Hitlera izteikumi ir autora izdomājums. Lai kā arī nebūtu, Ā.~Hitlera tālākā rīcība liecināja, ka konkrētais izteikums bija iespējams. Polijas vadītājs J.~Pilsudskis, tāpat kā daudzu citu zemju līderi, acīmredzot pārāk zemu novērtēja fīrera liekulību un neprognozējamību.

Citu Eiropas valstu galvaspilsētās auga noskaņojums uzsākt Vācijas revanša tendenču ierobežošanas pasākumus. Radās labvēlīga situācija nacistiskās Vācijas starptautiskajai izolācijai. 1933.~gada oktobrī Francijas ārlietu ministrs Ž.~Pols-Bonkūrs griezās pie PSRS ārlietu tautas komisāra M.~Ļitvinova ar priekšlikumu padomāt par iespēju noslēgt vienošanos par savstarpēju palīdzību sakarā ar Vācijas atklāto gatavošanos karam. Atbilde bija pozitīva un padomju vadība 1933.~gada decembrī pieņēma lēmumu izvērst darbību lai izveidotu efektīvu kolektīvās drošības sistēmu Eiropā un novērstu kara draudus. Bija nodomāts iesaistīt sistēmā veselu virkni Eiropas valstu, noteikti arī Franciju un Poliju. Taču Vācijas diplomātija, apzinoties, ka šādas sistēmas izveide var apdraudēt visus turpmākos Vācijas vadības plānus, apsteidza Franciju un PSRS, izspēlējot ``Polijas kārti''.

Pēc paša J.~Pilsudska vārdiem, pēc Vācijas izstāšanās no Tautu Savienības un tam sekojošās Vācijas izolācijas viņš uzskatījis, ka ir radušies nosacījumi pārvarēt sasprindzinājumu abu valstu attiecībās. Poļu vēsturnieks K.~Lapters rakstīja, ka ``Varšava, vērtējot situāciju, nolēma izmantot Vācijas grūtības, lai paātrinātu jautājuma atrisināšanu par poļu-vācu vienošanās noslēgšanu''.

Kā, balstoties uz J.~Pilsudska izteikumiem, atzīmē vācu vēsturnieki, no Polijas viedokļa austrieša Ā.~Hitlera nākšana pie varas pavēra cerību, ka ``protestantiski-prūsiskā negatīvā politika savienībā ar Padomju Krieviju varētu beigties''. J.~Pilsudskis, kurš pats ilgstoši bija sadarbojies ar Austroungārijas varas iestādēm, Ā.~Hitlera ``austrieša mentalitāti'' atzina par mazāk bīstamu Polijai, salīdzinot ar ``prūsisko'' pasaules redzējumu. Tā kā nacistu nākšana pie varas pārtrauca Vācijas un PSRS sadarbību, J.~Pilsudskis nolēma to izmantot Polijas labā. Ā.~Hitlera antikomunismā viņš saskatīja garantiju Vācijas un PSRS vienošanās neiespējamībai un līdz ar to~--- iespēju Polijai saglabāt ``līdzsvaru'' starp abiem tai~--- un arī savstarpēji~--- naidīgajiem kaimiņiem.

Katrā ziņā J.~Pilsudska acīs padomju briesmas Polijas eksistencei bija pirmajā vietā. Arī Polijas Ārlietu ministrijā šai laikā nostiprinājās viedoklis, ka ar Ā.~Hitlera nākšanu pie varas tālāka vācu-padomju sadarbība vairs nebūs iespējama, to attiecību normalizācijas perspektīvu poļu politiķi uzskatīja par nereālu. Tieši abu Polijas lielo kaimiņvalstu sadarbība Varšavā izsauca ne mazākas bažas kā Padomju Savienībā arvien vēl lietotais pasaules revolūcijas lozungs un nacistu pasludinātais kurss uz vācu ``dzīves telpas'' paplašināšanu. Poļu politiķi uzskatīja, ka šādos apstākļos ``līdzsvara politikas'' starp Vāciju un PSRS realizācijai ir jānostiprina Polijas drošība. Tai bija jāved pie Polijas kā Francijas un Lielbritānijas partnera lomas pieaugšanas, kas savukārt stiprinātu Polijas pozīcijas attiecībās ar Vāciju un PSRS. Ja abas pēdējās nespētu atdzīvināt ar Rapallo līgumu uzsākto sadarbību, rastos iespēja panākt stabilas starptautiskās attiecības Centrālajā Eiropā, kas veicinātu arī Polijas robežu drošību.

15.~novembrī Polijas sūtnis Berlīnē J.~Lipskis nodeva Ā.~Hitleram J.~Pilsudska mutisku vēstījumu. Maršals paziņoja par vēlmi saglabāt Vācijas un Polijas starpā labas kaimiņu attiecības, taču viņu uztraucot Vācijas izstāšanās no Tautu Savienības, jo dalību tajā viņš uzskatot par vienu no Polijas drošības pamatakmeņiem. Kamēr Vācija bija Tautu Savienības locekle, Polija varēja justies zināmā drošībā atbilstoši Tautu Savienības statūtiem un tajos paredzētajām procedūrām konfliktu izraisīšanās gadījumos. Tāpēc viņš (J.~Pilsudskis), pirms spert soļus Polijas drošības pastiprināšanai, nolēmis griezties pie Ā.~Hitlera ar jautājumu, vai tas neredz iespēju ``kompensēt poļu-vācu attiecībām nodarīto zaudējumu''. Fīrers pozitīvi reaģēja uz J.~Pilsudska vēstījumu. Pēc 12~dienām viņš nosūtīja J.~Pilsudskim savu līguma projektu. 1934.~gada 9.~janvārī savukārt savu projektu Vācijai iesniedza Polijas valdība.

Jau minētais H.~Raušnings citā savā grāmatā ``\detxti{Die Revolution des Nihilismus. Kulisse und Wirklichkeit im Dritten Reich}'', (Nihilisma revolūcija. Trešā reiha aizkulises un īstenība, 1938) rakstīja, ka 1933.~gada 13.~decembrī viņam bijusi slepena, neprotokolēta saruna ar J.~Pilsudski, kurā tas izrādījis gatavību vienoties ar Vāciju un pat sliecies uz priekšlikumu veikt kopīgu karagājienu pret PSRS. 1951.~gadā H.~Raušnings ``precizēja'' savus izteikumus: J.~Pilsudskis teicis, ka viņš neizbēgami karošot pret PSRS un tāpat neizbēgams esot Vācijas konflikts ar PSRS. Tā bijusi ``atklāta izlūkošana, vai nav iespējams noslēgt militāru savienību un veikt kopīgu akciju''. Atkal jāsaka, ka pilnībā uzticēties šiem izteikumiem nevar. Turklāt, pēc Francijas sūtņa Polijā liecības, J.~Pilsudskis jau nāves priekšvakarā ne vienu reizi vien atkārtoja savu secinājumu, ka Vācija tuvākajos 12~gados nebūs spējīga radīt uzbrukumam spējīgu armiju. Tiesa, iespējams gan, ka šis secinājums bija domāts tieši un tikai franču sūtņa ausīm.

1934.~gada 26.~janvārī Polija un Vācija parakstīja uz desmit gadiem \strong{``Deklarāciju par spēka nepielietošanu starp Poliju un Vāciju''}, kurā abas valstis solīja izskatīt tās skarošos jautājumus tiešās sarunās, bet ``ne pie kādiem apstākļiem strīdīgo jautājumu atrisināšanai nepielietos spēku''. Faktiski visai pasaulei tika paziņots, ka ar līgumu ne tikai likvidēts pēc Lokarno konferences pastāvošais sasprindzinājums abu valstu attiecībās, bet arī darīts gals ``tūkstošgadīgajam'' vācu un poļu savstarpējam naidīgumam. Presē pēc tam tika sludināts ``mūžīgā miera'' un ``divu tautu draudzības'' sākums. Polija pat uzņēmās Vācijas interešu aizstāvību tās atstātajā Tautu Savienībā.

Jāatzīst, ka, kā tolaik šķita, Polija bija ieguvēja. Galvenais ieguvums bija pārliecība, ka Vācija, atsakoties no ekspansijas Polijas virzienā, pārorientēsies citos virzienos, kas skars Rietumu intereses, kā rezultātā Francijas un Lielbritānijas acīs pieaugs Polijas kā starptautisko attiecību subjekta loma. Tādejādi ``sanācijas'' režīms cerēja piespiest divas minētās lielvalstis tādu pat statusu atzīt arī Polijai.

Bez šī šķietamā, bija arī reāli ieguvumi, tika izbeigta no 1925.~gada vestais Polijas~--- Vācijas ``muitas karš'', uzlabojās Polijas attiecības ar Dancigas brīvpilsētu un poļu mazākuma stāvoklis Vācijā.

Taču ne mazāk svarīgs bija cits aspekts. Noslēdzot vācu-poļu līgumu, ``sanācijas'' Polija bija pirmā valsts, kura uzsāka sadarbību ar nacistisko Vāciju. Savukārt Vācijai tas bija pirmais ievērojamais ārpolitiskais panākums, deva tai iespēju iziet no ārpolitiskas izolācijas, kurā tā atradās pēc nacistu nākšanas pie varas.

Poļu zinātnieks S.~Svaņēvičs, kurš līguma parakstīšanas brīdī bija tā dedzīgs piekritējs, vēlāk konstatēja, ka tas tika parakstīts laikā, kad norisa sarunas par ``četru pakta'', kuram bija jāmazina Vācijas apetīte pret Poliju, ratifikāciju, un ``torpedēja'' to. Toties deklarācija par spēka nepielietošanu gan neatcēla Vācijas pretenzijas uz ``poļu koridoru'', bet, pēc pētnieka viedokļa, deva ``desmit gadus kompromisa meklēšanai un sabiedriskās domas sagatavošanai abās zemēs šī kompromisa pieņemšanai''. (Šeit gan jāpiezīmē, ka faktiski Polija neieguva desmit gadu miera, un abas puses arī nemeklēja un negatavoja sabiedrisko domu kompromisam.)

Ā.~Hitlers pamatoti rēķinājās, ka rūpnieciski attīstītā Vācija uz 1939.~gada beigām būs gatava uzbrukuma operācijām, bet agrārvalsts Polija pat pēc visoptimistiskākajiem vērtējumiem būs gatava karot ne agrāk kā 1940.~gadā. Vācu vēsturnieks G.~Šrams konstatējis: ``Hitlers ar miermīlīgām metodēm, ar vācu-poļu tuvināšanos sasniedza to, ko Veimāras republika ar vācu-padomju par miermīlīgu diezin vai nosaucamu kopdarbību tā arī nespēja sasniegt. Polija kļuva par galīgi nederīgu tai iecerētajai lomai uzstāties kā cariskās Krievijas pēctecim~--- Francijas sabiedrotajam austrumos. Hitlers \citespace{} kā pirmais nonāca pie drosmīgās domas, ka Versaļas sistēmai, kuras sabrukumu līdz šim veltīgi bija mēģināts sagraut ar pretpolisku politiku, var likt krist ar poļu palīdzību.'' Jau 1938.~gada 26.~septembrī, pāris dienu pirms bēdīgi slavenās Minhenes vienošanās noslēgšanas, par kuru tiks vēl runāts, Ā.~Hitlers teica: ``Visgrūtākā problēma bija mūsu attiecības ar Poliju \citespace{} Mums neizdotos panākt samierināšanos, ja Polijai būtu demokrātiska konstitūcija. Labi, ka Polijā nav demokrātijas\dots{}''. Neapšaubāmi, ka šai gadījumā fīreram bija taisnība.

Polijai neuzbrukšanas līgums it kā nozīmēja savas rietumu robežas lielāku drošību, laika iegūšanu tālākai ekonomiskai attīstībai, kā arī lielāku ārpolitisko patstāvību. Vācu vēsturnieks E.~Meijers uzskatīja, ka J.~Pilsudskis ar līguma noslēgšanu vēlējies traucēt Vācijas un citu Rietumvalstu vienošanos par robežu izmaiņām uz Polijas rēķina. Padomju vēsturnieks D.~Kļimovskis, atsaucoties uz dokumentiem, norādīja, ka J.~Pilsudska piekritēji ``26.~janvāra līniju'' uzskatīja arī par Polijas lielvalstisko tendenču intensifikācijas līdzekli.

Taču, vērtējot minēto deklarāciju, jārēķinās, ka ``miermīlības'' demonstrācija Ā.~Hiltleram faktiski neko nenozīmēja, jo tās tekstā taču ne ar vārdu netika pieminēts pēc Lokarno līguma palikušais atklātais svarīgais jautājums~--- vai Vācija atzīst savas austrumu (tai skaitā arī ar Poliju) robežas. Atbildes trūkumu uz šo jautājumu varēja iztulkot (kā to arī daudzi darīja) tā, ka pakts maskēja abu līgumslēdzējpušu teritoriālās pretenzijas pret trešajām valstīm~--- Austriju, Čehoslovākiju, Lietuvu, PSRS. Pakta esamība netraucēja Vācijai vēlāk izvirzīt Polijai teritoriālas pretenzijas. 1939.~gada 22.~augustā augstāko armijas komandieru sanāksmē Ā.~Hitlers starp citu paziņoja: ``Noslēdzot paktu ar Poliju, man bija nolūks tikai iegūt laiku\dots{}'' Faktiski to, ka Polijas robežas paliks apdraudētas, nacisti tās vadībai neslēpa. Kā savos pēc kara uzrakstītajos memuāros liecināja nacistiskās Vācijas ārlietu ministrs J.~f.~Ribentrops, ``\dots{}iepriekšējās sarunās par vācu~--- poļu deklarāciju vācu puse izdarīja kategorisku aizrādījumu: ieplānotā vienošanās nekādā gadījumā neiekļauj Vācijas atzīšanu savām pastāvošajām robežām, gluži pretēji, ar to ir jātiek radītam pamatam visu problēmu atrisināšanai, tātad, arī teritoriālo.''

Vācijas-Polijas pakts radīja pamatīgu robu pēc Pirmā pasaules kara pastāvošajā Eiropas kolektīvās drošības sistēmā, tas ļāva apšaubīt arī Polijas miermīlīgumu. Polijas ieguvumi no J.~Pilsudska 1933.--1934.~gadā īstenotā ārpolitiskā manevra bija īslaicīgi. Neapšaubāmi, ka vairāk no poļu-vācu pagaidu tuvināšanās ieguva Ā.~Hitlers.

Virkne vēsturnieku ir izteikuši pamatotu viedokli, ka šis pakts ne tikai nesaturēja Polijas robežu garantijas, bet radīja šķēršļus Polijas iekļaušanai kolektīvās drošības sistēmā un faktiski atvēra iespējas spiedienam uz to. Ā.~Hitlers tā izkļuva no starptautiskas izolācijas, bet nebūt negatavojās līgumu pildīt. Viņš varēja nedomāt par Vācijas austrumu robežu, bet koncentrēties uz Austrijas un Čehoslovākijas sagrābšanu. (Iespējams, novērst Ā.~Hitlera uzmanību tieši šai virzienā arī bija viens no J.~Pilsudska mērķiem, līgumu slēdzot. Lai nu kā, J.~Pilsudska atklāsmes liecina, ka viņš acīmredzami nepietiekami novērtēja draudus Polijai, ko radīja Hitlera nākšana pie varas, dodot fīreram iespēju izmantot līgumu ar Poliju, lai demonstrētu Vācijas miermīlību brīdī, kad tā izstājās no Tautu savienības.)

Padomju vēsturnieki uzsvēra, ka noslēgtais poļu-vācu līgums nebija gluži parasts. Atšķirībā no pastāvošās prakses, arī 1932.~gada PSRS un Polijas neuzbrukšanas līguma, Polijas un Vācijas deklarācija paredzēja pakta palikšanu spēkā arī tai gadījumā, ja viena no līgumslēdzējvalstīm pirmā iesaistītos karā ar kādu trešo pusi. Deklarācijas īpašo raksturu atzīmēja daudzi politiskie komentētāji un diplomāti. Padomju izlūkdienesta aģents, par šo deklarācijas punktu, gan tulkojot to tikai kā Polijas gatavību ievērot neitralitāti, ja Vācija izprovocēta būtu spiesta uzsākt karu, ziņoja savai vadībai Maskavā. Kā norādījis krievu vēsturnieks L.~Ņežinskis, šāda punkta esamība radīja jautājumu, vai tas nenozīmēja, ka abas puses bija gatavas arī palīdzēt viena otrai gadījumā, ja šī otrā puse nostātos uz agresijas ceļa. Visvairāk no šāda pavērsiena apdraudētas varēja justies blakus esošās Austrija, Čehoslovākija, Lietuva un arī Padomju Savienība. Minēto valstu presē un politiskajās aprindās tika diskutēts jautājums, kā līguma noslēgšana ietekmēs Vācijas un Polijas attieksmi pret robežojošajām ar tām valstīm.

Sakarā ar šo deklarācijas īpatnību diplomātu, politiķu un žurnālistu aprindās sāka kursēt baumas, ka vācu-poļu paktam ir pievienoti vai nu slepeni panti vai notikusi apmaiņa ar Ā.~Hitlera un J.~Pilsudska vēstījumiem par Vācijas un Polijas slepenajiem nodomiem uzbrukt PSRS. Piemēram, šādu iespēju pieļāva franču žurnāliste un vēsturniece Ž.~Tabui. Franču avīzes izteicās visai noteikti par ``slepenu poļu-vācu vienošanos''. Sociālistu avīzes ``\frtxti{Le Populaire}'' (``Populārā'') 1934.~gada 8.~februārī rakstā ``Pilsudskis un Hitlers'' lasāms: ``Pats būtiskākais ir jautājums: par kādu cenu Pilsudskis un viņa banda noslēdza vienošanos ar Hitleru? Vai Polija dos Vācijai rīcības brīvību Austrijas jautājumā? Vai tā pieņems ``tehnisko'' Vācijas sadarbību Ukrainā, par kuru tā jau sen sapņo?''

1935.~gadā profranciski noskaņotais J.~Pilsudska politiskais oponents ģenerālis J.~Hallers izteica pārliecību, ka ``tagad vairs nav nekādu šaubu, ka starp Vāciju un Poliju pastāv slepens militārs līgums, virzīts pret PSRS''. Pēc J.~Hallera vārdiem šis līgums kara gadījumā ar PSRS nodrošināja Vācijai ``Polijas piejūras apgabalā īpašu etapa punktu un īpašu spēku vācu militāro transportu apkalpošanai organizēšanu.'' Tādā kārtā iekārtojušies poļu teritorijā, ``vācieši pēc kara darbības beigām automātiski iegūs visu šo teritoriju. Pilsudskis nepiešķir vajadzīgo uzmanību Polijas rietumu robežām un ir gatavs atteikties no Piejūras, lai īstenotu savus fantastiskos plānus attiecībā pret Ukrainu un Lietuvu''. 1935.~gada 18.~aprīlī franču kreisā avīze ``\frtxti{Le Bourbonnais republicain}'' pat publicēja it kā slepenu pielikumu vācu-poļu deklarācijas tekstam, kuru tai bija nodevis parlamenta deputāts, bijušais ministrs L.~Lemurjē. 20.~aprīlī šo tekstu nodrukāja arī padomju avīzes ``\rutxti{Правда}'' и ``\rutxti{Известия}''. Polijas oficiālās personas noteikti noraidīja šāda dokumenta esamību. Tā pieminējumi historiogrāfijā ir ārkārtīgi reti. Cik autoram zināms, poļu vēsturnieki to apiet ar klusēšanu. Bez komentāriem to ir pieminējusi padomju vēsturniece I.~Mihutina. 2004.~gadā par to ir rakstījis krievu vēsturnieks S.~Morozovs. Viņš konstatē, ka oriģināli dokumenti arhīvos nav saglabājušies, taču pastāv daudz netiešu datu, ka slepena vienošanās tomēr pastāvēja. Kā vērtē ukraiņu žurnālists S.~Lozuņko, var pat izdarīt pieņēmumu, ka slepenā poļu-vācu vienošanās stiprināja A.~Hitlera apņēmību radikāli vērsties pret Versaļas sistēmas iedibināto kārtību Eiropā. Taču aizdomas, lai cik arī tās atbilstu laikmeta kontekstam, nevar uzskatīt par pierādījumu.

Pasaules sabiedrība bija pilnībā pārsteigta, uzzinot, ka Ā.~Hitlers un J.~Pilsudskis ir savstarpēji apņēmušies nākamajos desmit gados atteikties no spēka pielietošanas savu valstu attiecībās.

Īpaši Francijā, kuru ar Poliju saistīja sabiedroto līgumattiecības, poļu-vācu tuvināšanās izsauca daudzu politiķu attieksmes maiņu pret savu stratēģisko sabiedroto. Nostiprinot Vācijas starptautiskās pozīcijas, vācu-poļu līgums deva pamatīgu triecienu Francijas pēc Pirmā pasaules kara radītajai militāro savienību sistēmai. Neraugoties uz līgumā ietverto apgalvojumu, ka tas nav pretrunā ar abu līgumslēdzējvalstu iepriekšējām saistībām, līgums faktiski grāva uzticību 1921.~gada franču-poļu līgumam. Kā rakstīja franču žurnāliste un vēsturniece Ž.~Tabui, Berlīnē 1934.~gada 26.~janvārī vācu-poļu deklarācijas parakstīšanas laikā esošajam Francijas premjerministra kancelejas vadītājam no Polijas sūtņa J.~Lipska nācās uzklausīt vārdus: ``Turpmāk Polijai Francija ne tikai nebūs vajadzīga, bet viņa pat nožēlo, ka viņai bija dārgi jāmaksā par tās pakalpojumiem''. Pēc tās pašas autores liecības Polijas ārlietu ministrs J.~Beks, izsakoties par iespējamo Francijas un PSRS līguma (1935) noslēgšanu, norādīja: ``Jūs zināt, franču-poļu savienība vairs neinteresē Poliju. Galu galā tieši Francijai ir sevišķi vajadzīga Polija \citespace{} Bet kas attiecas uz Krieviju, man trūkst epitetu, lai raksturotu naidu, kādu pie mums jūt pret to!'' J.~Beks arī uzsvēra, ka Polija ar Vāciju vienojusies, lai atdarītu Francijai par tās izturēšanos. Lūk, daži viņa izteikumi: ``Francija tik daudz mūs interesējošu jautājumu izlēma bez mazākās vienošanās ar mums, pat bez paziņošanas''; ``Francijai [par Polijas ar Vācijas vienošanās gatavošanu] netika ziņots, jo tā nevienu reizi vien tā ir rīkojusies ar Poliju''; ``Francijai mazāk par visiem jābrīnās par Polijas rīcību~--- ir Versaļā, ir Lokarno, ir turpmāk tā visai maz rēķinājās ar Polijas interesēm. Tagadējais Polijas valdības solis nedaudz atskurbinās pašpārliecinātos francūžus.'' (Nevilšus prātā nāk maza bērna loģika, kurš iekāpj peļķē, atriebjoties mammai, ka tā viņam nenopirka konfektes.)

1934.~gada aprīlī, kad Varšavā bija ieradies Francijas ārlietu ministrs L.~Bartū, J.~Beks arī viņam paziņoja, ka franču-poļu savienība vairs Poliju neinteresē un ka tieši Francijai bija nepieciešama Polija. Atpakaļceļā L.~Bartū sarunā ar saviem līdzstrādniekiem secināja, ka ``ir nepieciešams radīt Austrumu Lokarno, kura balstītos uz Krieviju, un nav vērts tērēt laiku ar Poliju''.

Francijas ārlietu ministra secinājums bija loģisks. Neviena cita valsts, lai arī ne nesavtīgi, nebija tik daudz darījusi Polijas atdzimšanas labā. Polijai vajadzēja lielā mērā būt pateicīgai Francijai, par tās interešu lobēšanu Parīzes miera konferencē, par izejas pie jūras iegūšanu, par būtisko palīdzību sakāves novēršanā karā pret Padomju Krieviju 1920.~gadā, par ievērojamu bijušo Vācijas teritoriju, pat pretēji tautu pašnoteikšanās tiesību principam un Lielbritānijas pretdarbībai, piešķiršanu Polijai. Polija bija pateicību parādā Francijai par to, ka tā varēja nostiprināt savus iekarojumus Austrumgalīcijā, Rietumukrainā un Rietumukarainā un Viļņas apgabalā. Franči apbruņoja poļu armiju, apmācīja tās virsniekus.

Franču acīs Polijas --Vācijas līgums nozīmēja to, ka ar vienu triecienu Ā.~Hitlers bija izsitis no Francijas pret Vāciju būvētās militāro savienību sistēmas svarīgāko posmu~--- Poliju. To pašu Poliju, kuru Francija bija palīdzējusi radīt un ``stutēt'' kā potenciālu pretvācu otro fronti. Tagad Ā.~Hitlers no kara divās frontēs varēja nebaidīties, Polija tam kļuva par drošu aizmuguri. Turpmāk Polija, vēršoties pret Čehoslovākiju, palīdzēja Ā.~Hitleram graut arī ar Francijas atbalstu 1920.--1921.~gadā radīto t.~s. Mazo Antanti (Čehoslovākijas, Rumānijas un Dienvidslāvijas savienība, kuras mērķis bija pretoties Ungārijas paplašināšanās centieniem, bet kura deva Francijai iespēju izveidot otru fronti konflikta gadījumā ar Vāciju.) 1933.--1934.~gadā Polijas Kara ministrija pārtrauca sakarus ar franču bruņošanās rūpniecību, ar firmu ``\detxti{Schneider-Creusot}'', militārie pasūtījumi tika novirzīti uz Lielbritāniju un Zviedriju. Ja līdz 1933.~gadam poļu karakuģi tika būvēti tikai Francijas kuģu būvētavās, tad 1935.~gadā ievērojamu pasūtījumu saņēma britu kuģu būvētāji. Zaudētājos bija arī Čehoslovākijas rūpniecība, firma ``\cstxti{Škoda}'' Polijas kara ministrijas personā zaudēja senu klientu. Ja, piemēram, padomju historiogrāfijā bija ierasts tikai pārmest Francijai par palīdzības nesniegšanu Polijai 1939.~gada septembrī, tad autoram jāuzsver, ka pati Polija pēc tās līguma ar nacistisko Vāciju bija daudz darījusi, lai vājinātu Francijas aizsardzības spējas. Tāpēc nenākas brīnīties, ka franči 1939.~gadā nedega vēlmē karot par Poliju, kura no 1934.~gada faktiski bija darbojusies pret Francijas interesēm.

Sabojājās Polijas attiecības ar Čehoslovākiju. Toreizējais Polijas vēstnieks Prāgā V.~Gžibovskis ziņoja uz Varšavu, ka Čehoslovākijā Poliju uzlūko par valsti, kura ``salauza vienoto pretvācu fronti un palīdzēja Vācijai iziet no starptautiskās izolācijas''. Polija sāka realizēt stingrāku kursu pret kaimiņvalsti Čehoslovākiju. Pirms poļu un vācu deklarācijas parakstīšanas 1934.~gadā attiecības starp Poliju un Čehoslovākiju bija salīdzinoši korektas, 1933.~gadā Varšava pat piedāvāja Prāgai noslēgt vienošanos par militāru sadarbību. Taču kad mazinājās spriedze Polijas attiecībās ar Vāciju, poļu militārās un diplomātiskās aprindas arvien noteiktāk sāka uzskatīt Čehoslovākiju par valsti bez ilgas pastāvēšanas perspektīvas. Čehu avīzes rakstīja, ka vācu-poļu līgums ``atraisīja Polijai rokas''.

Poļu presē pret Čehoslovākiju tika vērsta propagandas kampaņa, apvainojot čehoslovākus par poļu iedzīvotāju diskrimināciju Cešinas apgabalā. Tas noveda pie konflikta situācijas divpusējās attiecībās un iedzīvotāju deportācijām no abām pusēm. Polija armijas Galvenā štāba (\pltxti{Stab Glowny}, tā no 1928.~gada saucās Ģenerālštābs) Otrā (Izlūkošanas) nodaļa ar Ārlietu ministrijas atbalstu 1935.~gadā uzsāka nelegālu poļu organizāciju radīšanu Čehoslovākijas teritorijā. Polijā tika izveidotas teroristu treniņu nometnes, pāri robežai uz kaimiņvalsti tika sūtīti ieroči un sprāgstvielas, radiostacija Katovicē veica aktīvu antičehu propagandu. Polijā apmācīti sabotētāji piedalījās situācijas destabilizācijā Čehoslovākijā 1935.~gadā notikušajās prezidenta vēlēšanās. Acīmredzot tikai saistībā ar J.~Pilsudska nāvi graujošās aktivitātes tika īslaicīgi iesaldētas, bet atkal atsāktas trīs gadus pēc Austrijas ``anšlusa'' Vācijai 1938.~gadā.

Liels nemiers bija jūtams Baltijas valstīs. ``Baltijas valstis, un visupirms Lietuva,~--- no Ženēvas uz Varšavu ziņoja Polijas delegāts Tautu savienībā E.~Račinskis~--- tagad īpaši jūt Vācijas ekspansijas draudus \citespace{} Lietuva, meklējot līdzekļus sava stāvokļa nostiprināšanai griezās pie Igaunijas un Latvijas ar politiskas savienības priekšlikumu. Šis priekšlikums satur sevī pretpolisku dzeloni''.

Var piebilst, ka līgums negatīvi ietekmēja arī vācu mazākuma stāvokli Polijā, jo Vācija vairs neuzstājās kā tā aizstāve.

Poļu opozicionāri negatīvi novērtēja J.~Pilsudska politisko soli. Poļu literāts S.~Voitkevičš publicēja 1969.~gadā Vroclavā atrastu ģenerāļa V.~Sikorska pirmskara laika dienasgrāmatu. 1934.~gadā ģenerālis tajā atzīmēja, ka ``no šīs minūtes Polija iet vācu diplomātijas pēdās''. Un tālāk: ``Mēs aizmirstam par to, ka tas, kas Vācijai šai gadījumā ir glābiņš, Polijai nozīmē pazudināšanu šā vārdā pilnā nozīmē''. Arī padomju izlūkdienests 1935.~gadā saņēma no saviem aģentiem Polijā ziņojumu, kurā bija norādīts, ka J.~Pilsudskim opozīcijā esošie ģenerāļi J.~Hallers un V.~Sikorskis uzskata, ka starp Vāciju un Poliju ir jau parakstīts arī slepens militārs līgums, kurš Vācijai kara gadījumā pret PSRS nodrošina atbalsta punktus Polijas piejūras apgabalos, bet Polija par to kā kompensāciju saņems Lietuvu un Klaipēdu. (Šie secinājumi gan neatbilda patiesībai.)

Vācu-poļu līgumu atklāti nosodīja Kristīgo demokrātu partija (\pltxti{Polskie Stronnictwo Chrześcijańskiej Demokracji}) Nacionālā strādnieku partija (\pltxti{Narodowa Partia Robotnicza}), ļudovieši un sociālisti. Aktīvi pret ``sanācijas'' režīma ārpolitiku uzstājās Zemnieku partija (\pltxti{Stronnictwo Ludowe}). 1935.~gada decembrī savā kongresā tā pieņēma speciālu rezolūciju par ārlietām. Kongress nosodīja valdības politiku, kura ``novedusi pie gandrīz pilnīgas Polijas izolācijas un milzu neuzticības tai'', atzīmēja, ka ``vispārējs miers ir iespējams tikai organizējot savstarpēju palīdzību pret uzbrukumu.'' 1936.~gada jūnijā vairāk nekā 100~tūkstošu zemnieku kongress pārvērtās par pret valdību vērstu demonstrāciju. Sapulcējušies prasīja demokrātiskas brīvības un sadarbības ar hitlerisko Vāciju pārtraukšanu. Zemnieku partija valsts ārējās drošības problēmu saistīja ar ``sanācijas'' režīma likvidēšanu, savienības ar Franciju un Čehoslovākiju nodibināšanu un draudzīgām attiecībās ar Padomju Savienību.

Arī daudzi endeki uzstājās pret poļu-vācu tuvināšanos. Partijas iekšienē pieauga domstarpības ārpolitikas jautājumos.

Īpaši asi aizdomas par Polijas valdības agresīvajiem nodomiem izteica poļu komunisti. Komunisti paziņoja, ka viņi ir pret to, lai Polija ``kļūtu par caurstaigājamu pagalmu hitleriešu ģenerāļiem, kuri grib pakļaut kā poļu tautu, tā arī Padomju Savienības tautas''. Partijas Centrālas komitejas aicinājumā 1934.~gada februārī bija teikts: ``\dots{}šis līgums nozīmē, ka poļu un vācu imperiālisms ir izdarījis ievērojamu soli uz priekšu pa vienošanās ceļu par kopīgu uzstāšanos pret PSRS \citespace{} Starp Poliju un Vāciju pastāv lielas pretrunas. Un ja tagad, neraugoties uz šīm pretrunām, neraugoties uz savstarpējiem draudiem un nežēlīgo naidu, fašistiskie ķēdes suņi atklāti un publiski ir saostījušies, tad ir skaidrs kā diena, ka Hitlers un Pilsudskis ir vienojušies par kopīgu uzbrukumu PSRS.'' Pie tam par galveno agresīvo spēku poļu komunisti uzskatīja Polijas valdību. Tamlīdzīgi vienpusēji spriedumi tikai traucēja komunistiem ietekmēt masas.

Daudz precīzāks bija PPS Galvenās padomes paziņojums 1935.~gada novembrī: ``Beka ārpolitika, saistot Poliju ar ``trešo reihu'', ne tikai novājina zemi, vājinot saites, kuras to savieno ar Franciju, Mazo Antanti, PSRS un pat ar Baltijas valstīm, bet arī pastiprina kara iespējas, ievelkot Poliju revizionistiskajā un fašistiskajā vācu-ungāru blokā [Tiesa, šāds bloks nodibinājās tikai 1938.~gadā, kad Vācija kopā ar Poliju un Ungāriju sagrāba daļu Čehoslovākijas teritorijas.~--- V.Š.] Šajā situācijā strādnieku šķirai ir visiem spēkiem ir jācenšas saglabāt miers, nesaudzīgi apkarojot fašismu kā galveno kara faktoru.'' Polijas Sociālistiskās partijas XXIV~kongress Radomā 1937.~gada sākumā pieņēma partijas programmu, kurā kā galvenais bija izvirzīts uzdevums gāzt ``sanācijas'' diktatūru, nomainīt valsts un sabiedrisko iekārtu, pavēršot to demokrātijas attīstības virzienā, mainīt valsts ārējo politiku.

Visai uzskatāmi poļu sabiedrības attieksmi pret 1934.~gada 26.~janvāra līgumu parādīja t.s. ``Ļubļinas anketa''. 1936.~gada augustā Miera aizstāvības komiteja Ļubļinā izplatīja anketu, kurā starp citiem bija arī jautājums kādu ceļu aptaujātie uzskata par efektīvāko valsts neatkarības aizstāvēšanai: līgumu par savstarpējo palīdzību ar Franciju, Čehoslovākiju un PSRS vai savienību ar Vāciju. 5~872~aptaujātie atbildēja, ka Polijas neatkarība ir atkarīga no tās attiecībām ar Franciju, Čehoslovākiju un PSRS, tikai 7~izteicās par savienību ar nacistisko Vāciju. Protams, nezinot cik pilnīgi ar aptauju bija aptverti dažādi iedzīvotāju slāņi, grūti spriest par visas sabiedrības noskaņojumu pilnībā, taču acīmredzot var secināt, ka valdības ārpolitisko kursu sabiedrības vairākums neatbalstīja.

Jau pēc Otrā pasaules kara īsu, bet saturīgu Vācijas-Polijas līguma vērtējumu deva V.~Čērčils: ``Vācu-poļu pakts ļāva nacistiem koncentrēt uzmanību vispirms uz Austriju, bet pēc tam uz Čehoslovākiju, kas nesa graujošas sekas abām šīm nelaimīgajām valstīm. Tas uz laiku vājināja saiknes starp Franciju un Poliju un traucēja solidaritātes veidošanos Austrumeiropas valstu vidū.'' Pēc autora domām izcilā britu valstsvīra teiktajam var piebilst, ka minētais līgums arī ievērojamā mērā ļāva nacistiem nostiprināties Vācijā, traucēja vispārējas drošības sistēmas izveidei Eiropā, un gala rezultātā bija solis uz Otrā pasaules kara izraisīšanu. Vācijas-Polijas ``draudzība'' izrādījās 1939.~gada septembra Polijas katastrofas prologs.

Tomēr jāuzsver, ka J.~Pilsudskis, kurš jau 1932.~gadā bija noslēdzis neuzbrukšanas līgumu ar PSRS, ar 1934.~gada līgumu ar Vāciju turpināja Polijas ``līdzsvara politiku'', panāca dažus miera gadus valsts iekšējai stiprināšanai. Tūlīt pēc neuzbrukšanas pakta ar Vāciju noslēgšanas, 1934.~gada februārī Polijas ārlietu ministrs J.~Beks devās uz Maskavu un, kā jau agrāk teikts, pagarināja Polijas neuzbrukšanas līguma ar PSRS vēl uz desmit gadiem. Ja PSRS vadība 20.~gados arī sevišķi necentās uzlabot attiecības ar Poliju, jo nevēlējās apdraudēt attiecības ar Vāciju, tad 30.~gados sakarā ar nacistu nākšanu pie varas Vācijā, situācija mainījās, PSRS sāka meklēt tuvināšanos ar Franciju, sāka baidīties no Vācijas, mainīja arī attieksmi pret Poliju. Tā ``līdzsvaru'' Polija uz laiku panāca. 1934.~gada 7.~martā sanāksmē Belvedērā J.~Pilsudskis konstatēja, ka, pateicoties minētajiem neuzbrukšanas līgumiem Polijai izdevies panākt tādu savas drošības līmeni, kāds līdz šim nebija sasniegts. Tai pat laikā Polijas diktators uzskatīja, ka labas attiecības ar Vāciju pastāvēs ne vairāk kā četrus gadus (Kā zināms, viņa paredzējums piepildījās) un, ja viņa vairs nebūs, Poliju gaida smagi laiki, jo tikai viņš viens spējot nepieciešamības gadījumā atlikt jautājuma izlemšanu vai uzstādīt to citādāk. J.~Pilsudskis, protams, saprata, ka viena lieta ir noslēgt līgumu, pavisam cita~--- uzticēties Ā.~Hitleram. Acīmredzot maršals gaidīja, ka noslēgtais līgums vismaz uz laiku novērsīs vācu teritoriālās pretenzijas un, iespējams, arī Vācijas iejaukšanās briesmas Polijas iekšpolitikā ar vācu mazākuma tiesību aizsardzības ieganstu. (Tik tiešām, kad 1934.~gada 13.~septembrī J.~Beks Tautu Savienībā Ženēvā paziņoja par Polijas atteikšanos ievērot līgumu par nacionālo mazākumu aizsardzību, viens no poļu diplomātiem~--- J.~Beka līdzgaitniekiem~--- J.~Lukaševičs to nosauca par pirmo derīgāko no vācu-poļu līguma rezultātiem.)

Interesantas ziņas par pēdējām J.~Pilsudska dzīves dienām sakarā ar viņa realizētās ``līdzsvara politikas'' rezultātu satur padomju aģenta ziņojums par stāvokli Polijā 1938.~gadā. Pēc tā bijušais Polijas militārais atašejs Parīzē pulkvedis J.~Blešinskis stāstījis, ka mirstot J.~Pilsudskis vairākkārt atkārtojis: ``Zjuk (pamazināmais vārds no ``Juzefs''~--- tā viņu bērnībā sauca ģimenē un arī līdzbiedri Polijas Sociālistiskajā partijā), tu zaudēji, Zjuk, tu zaudēji''. Lieta tā, ka pirms tam caur Varšavu uz Maskavu devās Francijas ārlietu ministrs P.~Lavals, kurš tur 2.~maijā arī noslēdza Francijas un PSRS līgumu par savstarpēju palīdzību. J.~Pilsudskis, būdams smagi slims, vēl neveiksmīgi mēģināja traucēt minētā līguma noslēgšanu, gribēja piecelties no gultas, kas izsauca iekšēju asiņošanu. Pirms nāves J.~Pilsudskis apzinājies, ka viņa vestā ārējā politika ar Francijas un PSRS līgumu par savstarpēju palīdzību noslēgšanu ir bankrotējusi. Par notikumu atstāsta precizitāti nevar galvot, taču pēc būtības Francijas--PSRS līgums tik tiešām daļēji bija atbilde uz Polijas --Vācijas līgumu un citiem Polijas soļiem.

Līgumu noslēgšana ar kaimiņvalstīm~--- Vāciju un PSRS varētu nest Polijai arī stabilu ieguvumu, ja laiks tiktu izmantots valsts aizsardzības spēju nostiprināšanai, taču tam nepietika līdzekļu un arī J.~Pilsudska novecojušie priekšstati traucēja modernizēt armiju. ``Līdzsvara'' politiku Polijai galu galā tā arī neizdevās līdz galam realizēt.

Abas kaimiņvalstis laiku izmantoja efektīvāk un Polija Otro pasaules karu sagaidīja sliktāk sagatavojusies nekā tās. 1934.~gadā J.~Pilsudskis, šajā laikā būdams jau smagi slims, apzinājās stāvokļa nopietnību. Viņš brīdināja: ``Polija sēž uz diviem krēsliem. Tas nevar turpināties ilgi. Mums jāzina, no kura nokritīsim vispirms un kad tas notiks.''

Poļu vēsturnieks S.~Gregorovičs, kurš ir pētījis Polijas attiecības ar PSRS un Vāciju 30.~gados, uzsver, ka Polijai attiecību uzlabošana ar PSRS varētu palīdzēt tās uzlabot arī ar Vāciju, taču tikai tai gadījumā, ja šī Polijas tuvināšanās PSRS nebūtu ar pretvācisku raksturu. PSRS turpretī vēlējās poļu-padomju attiecībām piešķirt tieši šo pretvācisko raksturu. Šai ziņā Polijas un PSRS intereses atšķīrās. Polija vēlējās uzlabot attiecības arī ar Vāciju, neļaut tai atjaunot Rapallo kursu, arī celt Polijas kā partnera nozīmību jebkurās starptautiskās sarunās. PSRS turpretī baidījās no Polijas nonākšanas Vācijas ietekmē. S.~Gregorovičs 1934.~gada 26.~janvārī Polijas un Vācijas parakstītajā neuzbrukšanas līgumā saskata tikai ``līdzsvara politikas'' pēc formulas ``Ne par soli tuvāk Maskavai kā Berlīnei'' konsekventu realizāciju.

Pret šo uzskatu ir arī iebildumi. Piemēram, krievu vēsturnieks V.~Sokolovs raksta: ``Polijas valdība 30.~gados nekad nepieturējās pie vienādas attāluma principa no Maskavas un Berlīnes, kā kādreiz mīl teikt daži vēsturnieki. Pietiek atcerēties, cik reizes J.~Beks bija Berlīnē un cik reizes viņš uzņēma augstus Vācijas viesus Varšavā. Tai pat laikā Maskavā viņš ar vizīti bija tikai vienu reizi~--- 1934.~gadā, par cik visa tālākā viņa politika bija vērsta uz PSRS izolēšanu un pretpadomju savienības izveidi. Tam ar nemieru sekoja Maskavā, fiksējot katru Varšavas soli.''

Tāpēc PSRS šādos apstākļos sāka realizēt pretvāciskas frontes un Vācijas starptautiskas izolācijas politiku. Jau VK(b)P XVII~kongresā 1934.~gada janvārī J.~Staļins paziņoja, ka PSRS ``orientējas uz PSRS un tikai uz PSRS \citespace{} un ja PSRS intereses prasa tuvināties ar tām vai citām valstīm, kuras nav ieinteresētas miera pārkāpšanā, mēs uz to ejam bez kādām šaubām''. (Tiesa, vēlākie notikumi parādīja, ka PSRS vadība ir gatava tuvināties arī tām valstīm, kuras kā reiz ir ieinteresētas miera pārkāpšanā, ja tikai tas atbilda dotajā brīdī pastāvošajām PSRS interesēm.) Galvenais faktors, kurš noteica PSRS politikas maiņu bija nacistu nākšana pie varas Vācijā, tās soļi uz savas ietekmes pastiprināšanu Eiropā. Metodes, ar kādām nacistiskā Vācija darbojās starptautiskajā arēnā, nevarēja neizsaukt bažas padomju vadībā. Tas noteica vēlmi tuvināties ar Franciju un tās sabiedrotajām, vēlmi radīt starptautiskās drošības sistēmu.

Viens no šīs politikas soļiem bija PSRS iestāšanās Tautu Savienībā. Kā uzņemšanas noteikumus tā izvirzīja kolektīvu Tautas savienības biedru uzaicinājumu un pastāvīgas vietas piešķiršana Savienības Padomē. Tā ko Polija pati pretendēja uz pastāvīgā Padomes locekļa vietu, tā pret izmaiņām savienībā izturējās nelabvēlīgi. Tikai Lielbritānijas un Francijas paziņojums, ka tās pilnībā atbalsta PSRS uzņemšanu, lika Polijai atteikties no aizkulišu opozīcijas. Taču apmaiņai pret PSRS kandidatūras atbalstu tā pieprasīja lai padomju diplomātija atzītu, ka abu valstu attiecību pamatu veido agrāk noslēgtās vienošanās. 1934.~gada 18.~septembrī PSRS tika uzņemta Tautu Savienībā.

Atšķirības Polijas un PSRS politikas mērķos galvenokārt izpaudās jautājumā par t.s. \strong{Austrumu paktu}. Tā tika saukts 1934.~gada sākumā Francijas ārlietu ministra L.~Bartū projekts par savstarpējas palīdzības līguma noslēgšanu starp PSRS, Čehoslovākiju, Latviju, Lietuvu, Poliju un Somiju. Lai izvairītos no Vācijas pārmetumiem par to, ka plānota tiek tās ``ielenkšana'', L.~Bartū piedāvāja arī tās iekļaušanu paktā. Bijušais Vācijas sūtnis Maskavā, bet 1935.~gadā Tokijā, H.~f.~Dirksens atzīmēja, ka ``austrumu pakts'' noslēdza durvis austrumu robežu caurskatīšanai''.

Projekta realizācijas mēģinājumi tomēr izsauca sarežģījumus. Ne Vācijas, ne Polijas valdībai tāda ``durvju noslēgšana'' nebija pa prātam. Polijas valdība ļoti labprāt caurskatītu robežas gan ar Čehoslovākiju, gan Lietuvu, gan PSRS. ASV sūtnis PSRS V.~Bullits 1934.~gada jūlijā ASV valsts sekretāram K.~Hellam rakstīja, ka J.~Pilsudska atteikšanās no dalības paktā ir saistāma ar padomju-japāņu kara gaidām, jo maršals vēlas sev saglabāt brīvas rokas austrumos, lai ``atjaunotu tur agrāko Polijas diženumu''. Cerības, ka notikumu attīstība novedīs pie PSRS-Japānas konflikta, aizmigloja poļiem acis. Polija enerģiski attīstīja militāri-tehnisko sadarbību ar Japānu, veicināja abu valstu izlūkdienestu darbības pret PSRS koordināciju. (Par Japānas lomu Vācijas un Polijas plānos pamatīgāk ir rakstījis vēsturnieks S.~Morozovs.)

Ja Francija un PSRS ar Austruma pakta palīdzību cerēja ievērojami uzlabot drošības līmeni Austrumeiropas reģionā, Polija pakta īstenošanas gadījumā zaudētu to ``neatkarību'' starptautiskajā arēnā, kuru, pēc J.~Pilsudski vārdiem, izdevās panākt 1934.~gada janvārī. Pamatojoties uz šiem apsvērumiem, maršals nolēma boikotēt Parīzes priekšlikumu.

Rezultātā Vācijas valdība, atsaucoties uz valsts ierobežojumiem bruņošanās jomā, 1934.~gada 8.~septembrī atteicās līgumā piedalīties. 27.~septembrī arī Polija atteicās no dalības tajā, norādot, ka tā varētu to darīt vienīgi tad, ja līgumam pievienotos Vācija. Galu galā Austrumu pakts tā arī netika noslēgts. Vēloties to vai nē, Polija otro reizi atbalstīja hitlerisko Vāciju, kas šajā projektā redzēja draudus tās ekspansīvajām centieniem.

Krievu vēsturnieks V.~Sokolovs uzskata, ka tieši J.~Beks ar to uz ilgu laiku apturēja kolektīvās drošības idejas īstenošanu. Netieši šo viedokli apstiprina virkne poļu vēsturnieku. Viens no viņiem~--- V.~Materskis norādījis, ka Parīzē J.~Bekam tika pārmests, ka neatbalstot Austrumu paktu, viņš grūž J.~Staļinu uz savienību ar Ā.~Hitleru. Kā raksta cits poļu pētnieks T.~Šarota, franču acīs J.~Beka politika arvien vairāk atspoguļoja nevis ``līdzsvara'' starp Vāciju un PSRS ievērošanu, bet iezīmēja Polijas pārvēršanos par Vācijas partneri. A.~Skšipeks norādījis, ka Polijas valdības negatīvu attieksmi pret Austrumu pakta projektu izsauca J.~Pilsudska nevēlēšanās tuvināties PSRS, politikas izmaiņa par labu Vācijai, Francijas un PSRS savienības skeptisks vērtējums, jo tā, kā viņam likās, mazināja Polijas kā Francijas sabiedrotā lomu, beidzot, kategoriska nevēlēšanās izlaist Sarkano armiju caur Poliju Francijas, Polijas un PSRS kara gadījumā pret Vāciju. Acīmredzot galvenais bija tas, ka Polija baidījās no Francijas un PSRS tuvināšanās, kas mazināja Polijas lomu Francijas-Polijas savienībā, kā arī sabojāt tikko uzlabotās vācu-poļu attiecības.

Austrumu pakta neveiksmi sabiedrība vērtēja kā vācu-poļu tuvināšanās apliecinājumu. Padomju vēsturnieks D.~Kļimovskis Polijas veikto Eiropas drošības sistēmas radīšanas iespējas zem ``Austrumu pakta'' nosaukuma ``torpedēšanu'' vērtējis kā pirmo etapu ceļā uz Polijas katastrofu 1939.~gadā. (Vai tas bija pirmais, var strīdēties.)

Polijas un PSRS uzskati atšķīrās arī jautājumā par franču-padomju mēģinājumiem veidot daudzpusīgu drošības sistēmu Eiropā, par abu valstu noslēgto savstarpējās palīdzības līgumu 1935.~gada 2.~maijā.

Polijas izturēšanās savukārt izsauca PSRS neapmierinātību un kārtējo abu valstu attiecību pasliktināšanos. Poļu vēsturnieks S.~Gregorovičs akcentē, ka Polijas sadarbība ar kādu no tās spēcīgākajiem kaimiņiem objektīvi nostādītu vājāko Poliju atkarībā no šī kaimiņa, bet tā laika ģeopolitiskajā realitātē neatkarīgas politikas zaudēšana vestu arī pie pašas neatkarības zaudēšanas. Turpretī Maskava Polijas ārpolitikā pirmkārt saskatīja pretpadomju kursu, vēlmi izolēt PSRS no Eiropas jautājumiem.

S.~Gregorovičs PSRS politiku starpkaru periodā raksturo kā stratēģiski virzītu uz Polijas valsts graušanu. Pēc lasītāja priekšā esošās grāmatas autora domām, šiem poļu vēsturnieka kopumā loģiskajiem konstatējumiem un secinājumiem tomēr ir vajadzīga piebilde. To, ko viņš 30.~gadu vidū sauc par padomju (un arī franču) pretvācisko kursu, citiem vārdiem un ar lielākām tiesībām var nosaukt par kolektīvās drošībās sistēmas veidošanas mēģinājumiem. Realizējot savās interesēs līdzsvara politiku, Polija neņēma vērā, ka 30.~gados starptautiskā situācija jau bija mainījusies, jo Vācijā pie varas bija nākuši nacisti, kuri apdraudēja visu pasauli, ka pret šīm briesmām varēja cīnīties tikai vairāku valstu apvienotiem spēkiem. Šādā situācijā PSRS nedraudzīgajai politikai pret Poliju, bez nosodāmajiem impēriskajiem mērķiem bija arī vismaz daļējs attaisnojums~--- Polija taču traucēja kolektīvās drošības sistēmas izveidi Eiropā. Ģenerālis V.~Sikorskis 1936.~gada pavasarī savā dienasgrāmatā rakstīja: ``Itālijas un Trešā reiha aliansei ir jāpaātrina Eiropas apvienošana koalīcijā pret Vāciju. Koalīcijai pēc iespējas ātrāk ir jāiekļauj sevī Padomju Krieviju. Bez balstīšanās uz Krieviju miermīlīgais bloks neuzvarēs''. Par nelaimi Polijai V.~Sikorskis bija atstumts no politiskas ietekmes. Protams, kādus īpašus pārmetumus Polijas valdībai par tradicionālās politikas realizēšanu arī 30.~gados neklājas izvirzīt, jo tikai atskatoties no tālākas gadu distances ir saprotama notikumu jēga. Diemžēl S.~Gregorovičs nav vēlējies saskatīt pretrunu starp Polijas īslaicīgajām un visas Eiropas ilglaicīgajām interesēm.

Arī turpmāk J.~Pilsudska un viņa epigonu realizētā politika palika pretrunīga. Blakus drošības apsvērumiem to lielā mērā ietekmēja arī lielvalstiskās ambīcijas. Par to, ka J.~Pilsudskis, un līdz ar viņu arī Polijas valdība, maz uztraucās par agresiju Eiropā, ja tā vērsās pret citām valstīm, uzskatāmi liecina poļu diplomāta, tās sūtņa Itālijā A.~Visocka atmiņas par sarunu ar J.~Pilsudski 1933.~gada 21.~jūlijā. J.~Pilsudskis sūtni pamācīja: ``Mussolini kungs ļoti interesējas par Austriju un ir ļoti pret anšlusu [Austrijas pievienošanu Vācijai]. Nu tad tā, Jūs nekad nejaucaties šajās lietās un ne par ko neuztraucaties, jo mēs esam gatavi pārdot visu šo anšlusu, tikai mums par to jāsaņem laba samaksa. Tāpat mēs pārāk nedusmosimies, ja~--- un te maršals sāka smieties kā uzjautrināts bērns~--- ja Dienvidtirole atgriezīsies pie savas vācu Tēvijas.'' Šos vārdus nevilšus nākas atcerēties, runājot par 1939.~gadu, kad PSRS par ``labu samaksu'' ļāva Vācijai sagraut Poliju.

% page 642


1935.~gada sākumā Polijā, kārtējo reizi ``uz medībām'' ieradās ``otrais vīrs'' Vācijā, Prūsijas ministru prezidents H.~Gērings. Sarunā ar J.~Pilsudski viņš, pēc Polijas ārlietu valsts sekretāra (ministra vietnieka) J.~Šembeka liecības, kurš pierakstīja pārrunas, ``izteica domu par kopēju vācu un poļu gājienu pret Krieviju, norādot uz priekšrocībām, kuras šī akcija dotu Polijai Ukrainā''. H.~Gērings ``deva saprast, ka šādā gadījumā Ukraina kļūs par Polijas ietekmes sfēru, kad Krievijas ziemeļrietumu daļa būs vācu ietekmes sfēra''. Protams, ir jāsaprot, ka nedz Vācija vēl bija gatava karam pret PSRS, nedz vecais Polijas diktators bija tik viegli noskaņojams uz starptautiskām avantūrām.

J.~Beka realizētā Polijas ārpolitika arī 30.~gadu otrajā pusē pilnībā sekoja J.~Pilsudska agrākajiem norādījumiem un maz mainījās. Valsts prezidents I.~Moscickis joprojām bija amatieris diplomātijā, armijas ģenerālinspektors E.~Ridzs-Smiglijs interesējās tikai par tiem ārpolitikas jautājumiem, kas bija tieši saistīti ar armijas modernizāciju un valsts aizsardzības spējām, bet premjerministrs nodarbojās ar iekšējām lietām. No parlamenta faktiski izraidītā opozīcija zaudēja iespēju ietekmēt Ārlietu ministrijas politiku.

Kaut arī Vācijas politiķi baidījās, ka Francija pēc maršala J.~Pilsudska nāves centīsies Poliju attālināt no Vācijas, tā nenotika un Polija turpināja maršala iedibināto ``attiecību līdzsvaru'' starp Vāciju un PSRS. J.~Beks mēģināja turpināt politiku ``vienlaicīgi sēdēt uz diviem krēsliem'': viņš vienlaikus demonstrēja apņēmību pildīt 1921.~gada Polijas un Francijas savienības līgumu un darīja visu iespējamo, lai nepasliktinātu attiecības ar Vāciju.

Maršala J.~Pilsudska nāve 1935.~gada maijā maz ko mainīja Polijas attieksmē pret PSRS. Palika spēkā ``naidīgās atturības politika pret PSRS'' (V.~Materskis). Līdz 1938.~gadam savstarpējā neuzticība tikai turpināja pieaugt. Kontakti faktiski pastāvēja tikai minimālā līmenī. Liela nozīme šai situācijā bija PSRS notiekošajām represijām. Kā secināja Polijas ārlietu ministrs pulkvedis J.~Beks, milzīgie zaudējumi, kurus represiju gaitā nesa PSRS politiskie, militārie un saimnieciskie kadri, uz daudziem gadiem atvirza PSRS iespējamā uzbrukuma draudus Polijai, tā savu austrumu flangu var uzskatīt par drošu un tāpēc var koncentrēties citiem uzdevumiem.

Kad 1935.~gada oktobrī Itālija uzsāka agresiju pret Etiopiju, pārējās Rietumvalstis nevēlējās aktīvi iesaistīties tās apturēšanā. Tautu Savienība pēc Lielbritānijas priekšlikuma vienīgi noteica pret Itāliju ekonomiskas sankcijas, kurām bija spiesta pievienoties arī Polija. Taču tās pilnībā pildīt Polija nevēlējās. Polija bija Itālijā pasūtījusi un gandrīz jau pilnībā apmaksājusi tvaikoni un automobiļu virsbūves savas armijas vajadzībām. Pie tam uz Itāliju bija novirzīts 3/4 visa Polijas ogļu eksporta, to pārtraucot, Polijai nāktos ciest lielus finansiālus zaudējumus. Tāpat kā Polijai, savi iegansti nebojāt attiecības ar Itāliju atradās arī citām valstīm, tādēļ sankcijas izrādījās neefektīvas. Rumānijas sūtnis Polijā K.~Visoianu 1936.~gada janvārī konstatēja: ``Ar savu tagadējo politiku Polija ved spēli, tā grib būt laba šodien pret vāciem, rīt pret frančiem un tā bez gala. Ar šādu politiku Polija panāk tikai to, ka valda neuzticība un ka [tagadējos] apstākļos, sevišķi Eiropas austrumos, nenonāk pie pilnīga miera un nostiprināšanās.'' 1936.~gada 26.~jūnijā J.~Beks paziņoja Tautu Savienībai, ka Polija atsakās no pret Itāliju vērstās sankciju politika. Polija bija pirmā valsts, kura starptautiski atzina Itālijas realizēto Etiopijas sagrābšanu. Polijas valstsvīri nevarēja iedomāties, ka drīzumā pati Polija var kļūt par agresijas upuri. Kad 1936.~gada jūlijā sākās dumpis pret likumīgo varu Spānijā, Polijas diplomātija faktiski atbalstīja Vācijas iejaukšanos Spānijas lietās.

Kaut Polijas Ārlietu ministrija ar J.~Beku priekšgalā realizēja salīdzinoši neatkarīgu ārējo politiku, kura arī pēc J.~Pilsudska nāves savos pamatprincipos bija balstīta uz viņa iedibināto ``attiecību līdzsvaru'' starp Vāciju un PSRS, tomēr līdz pat mūsdienām daudzi poļu pētnieki, (iespējams, pat vairākums) uzskata, ka J.~Beks ievērojami vairāk distancējās no Maskavas, nekā no Berlīnes, un sakarā ar to viņa realizēto politiku nevar saukt par ``līdzsvara'' politiku. Vairākums poļu vēsturnieku kā uz vienu no galvenajām J.~Beka kļūdām norāda uz to, ka viņš pēc Ā.~Hitlera nākšanas pie varas un idejiski-propagandiska kara izvēršanās starp nacistisko Vāciju un staļinisko PSRS nemaz neņēma vērā jauna Rapallo izlīguma vai vēl kādu citu abu valstu tuvināšanās iespēju. Tāpat daļa poļu vēsturnieku mūsdienās norāda uz to, ka ``līdzsvara'' politika bija visai vāja konstrukcija, kura nespēja kļūt par šķērsli Ā.~Hitlera un J.~Staļina imperiālistiskajiem centieniem. Arī kanādiešu vēsturnieks M.~Karleijs uzskata, ka 30.~gados J.~Beks arvien vairāk sliecās uz sadarbību ar Vāciju un konfrontāciju ar PSRS. Krievu vēsturnieks S.~Slučs domā, ka Polijas novirze ārpolitikā uz tuvināšanos Trešajam reiham ne tikai pārkāpa ``līdzsvara'' politikas principus, bet deva triecienu arī visai starptautisko attiecību sistēmai Eiropā, lielā mērā neitralizējot vienu no svarīgākajiem Francijas drošības sistēmas elementiem.

Vienlaikus Polija arī pēc J.~Pilsudska nāves centās spēlēt ievērojamu, dažkārt pārspīlēti, nepamatoti ievērojamu lomu reģionā un pat visā Eiropā. J.~Beks to formulēja kā priekšstatu par ``trešo Eiropu'' starp Vāciju un Padomju Savienību. Tāpat kā agrāk, arī pēc J.~Pilsudska nāves Polijas ārpolitika vērtēja Lietuvu un Čehoslovākiju kā šķēršļus Polijas ietekmes izplatībai visā reģionā starp Baltijas un Melno jūru.

Krievu vēsturnieks S.~Morozovs Polijas un Vācijas attiecības pēc 1934.~gada 26.~janvārī Polijas un Vācijas līguma dalīja divos periodos. Ja 1934.~gada ziemas~--- 1935.~gada pavasara periodā tās bijušas vienlīdzīgu partneru attiecības, tad pēc PSRS~--- Francijas savstarpējās palīdzības pakta noslēgšanas 2.~maijā līdz 1938.~gadam Vācija pret Poliju jau izturējās kā pret jaunāko partneri, izmantojot tās sliktās attiecības ar Čehoslovākiju.

Sākotnēji Polijas vadība nepietiekami novērtēja Ā.~Hitlera režīma dinamismu. Polija nepretojās vispārējās karaklausības ieviešanai Vācijā 1935.~gadā. 1936.~gada martā, kad Vācija ieveda savu karaspēku demilitarizētajā Reinas zonā, kā ieganstu izmantojot 1935.~gadā maijā parakstītā Francijas~--- PSRS savstarpējās palīdzības līguma ratifikāciju, Polija gan vārdos nosodīja Vācijas rīcību, bija it kā gatava palīdzēt Francijai, ja tā pretotos Vācijas solim, taču Francija piekāpās bez cīņas, noteiktas reakcijas nebija arī no citu Rietumu lielvalstu puses. Arī Polija kā īpaši nereaģēja, nevēloties pasliktināt attiecības ar Vāciju. Diplomātiski noformējot savu soli, kurš pārkāpa Versaļas līguma un Lokarno vienošanos nosacījumus, Vācijas ārlietu ministrs K.~f.~Neirāts uzaicināja Lielbritānijas, Francijas, Beļģijas un Itālijas sūtņus un paziņoja par Vācijas atteikšanos no Lokarno līgumiem. Tā vietā Vācijas valdība paziņoja par gatavību noslēgt neuzbrukšanas līgumus ar Franciju un Beļģiju. Pie tam Vācijas un Polijas 1934.~gada līgums tika nosaukts kā vienīgais konstruktīvais paraugs miera nodrošināšanai Eiropā. Turpretī D.~Kļimovskis rakstīja, ka 1934.~gada Vācijas-Polijas līgums, kurš formāli turpināja eksistēt līdz pat tā denonsācijai no Ā.~Hitlera puses 1939.~gada aprīlī, pēc Reinas zonas okupācijas faktiski pārvērtās par ``dzīvu līķi''. Ar šo soli nacistiskā Vācija bija izkļuvusi no visriskantākās zonas, kad kaimiņvalstis pret to varēja vērsties ar militāru spēku un piespiest ievērot starptautiskos līgumus. Vācija uzsāka (faktiski turpināja~--- V.Š.) strauju bruņošanos un Eiropa virzījās pretī jaunam karam.

Gluži mierīgi uz starptautiskajiem notikumiem Polijas valdošais slānis gan nevarēja raudzīties. 1936.~gada maijā Polijas vadītāju sanāksmē pie prezidenta tika nolemts pret PSRS saglabāt politisko kursu, balstītu līgumā par neuzbrukšanu, pēc iespējas attīstīt tirdzniecību, taču reizē ar to pārtraukt jebkādu sadarbību kultūras laukā kā ``kaitīgu un sekmējošu komunisma un padomju ideju propagandu''. Bailes no ``komunisma un padomju ideju propagandas'' liecināja par to, ka poļu vadošie politiķi apzinājās objektīvās augsnes esamību šīm idejām Polijā. Taču reizē 1936.~gada jūnijā Polijas armijas ģenerālinspektors E.~Ridzs-Smiglijs deva pavēli Galvenajam štābam izstrādāt karadarbības plānu pret Vāciju. Ārlietu ministrs J.~Beks noraidīja piedāvājumu Polijai iestāties pret PSRS vērstajā Antikominternes paktā (25.11.1936). Tādejādi Polija turpināja līdzsvara politiku.

Polijas plāni karam ar kaimiņvalstīm ``Rietumi'' (``\pltxti{Zachόd''}) un ``Austrumi'' (``\pltxti{Wschόd}'') bija orientēti uz aizsardzību. Arī nacionālās audzināšanas darbs, neraugoties uz lielas daļas politiķu lielvalstiskajām pretenzijām, Polijā pamatos bija vērsts uz valsts aizsardzības stiprināšanu (\pltxti{ku obronie Rzeczpospolitej}), bet nevis uz uzbrūkoša kara vešanu. To apzinājās arī vācu ģenerāļi. Vācijas kara ministrs V.~fon Blombergs 1937.~gadā izteicās, ka ``Vācijai nav jārēķinās ne ar kādu iespējamu triecienu \citespace{} Ka Polija varētu uzsākt karu pret mums, iespējams Krievijas pusē, pie tagadējā politiskā stāvokļa ir vairāk nekā neiespējami''. Tomēr valdība, arī ārlietu ministrs J.~Beks, pārvērtēja Polijas spēkus, uzsvēra tās lielvalstiskos centienus, kuri galvenokārt bija virzīti uz Austrumiem. Ar Berlīni Polijā rēķinājās vairāk nekā ar Maskavu.

Kad Vācija 1937.~gadā izvirzīja jautājumu par savu bijušo koloniju atgūšanu, koloniju jautājums kļuva aktuāls arī Polijā. Tās valdošie slāņi apgalvoja, ka Polijai, kā katrai sevi cienošai Eiropas valstij, ir jāiegūst kolonijas, kaut vai Āfrikā. Kā atzīmējis Ē.~Jēkabsons, tas tika darīts, jo J.~Beks uzskatīja, ka aktuālā politiskā situācija pasaulē prasa no Polijas aktīvāku ārējo politiku, bet daļēji arī iekšpolitisko prasību dēļ~--- uzstādīt prasību pēc kolonijām lika smagā ekonomiskā situācija, pārapdzīvotība laukos un iespaidīgā laukstrādnieku sezonas emigrācija. Taču jāpasaka arī, ka prasība pēc savām kolonijā Polijā izskanēja jau krietni agrāk.

Jau 1918.~gada 1.~oktobrī pēc atvaļināta Krievijas flotes kontradmirāļa K.~Porembska iniciatīvas tika izveidota ``Kuģniecības veicināšanas biedrība ``Poļu karogs'' (\pltxti{Stowarzyszenie na Polu Rozwoju Żeglugi ''Bandera Polska''}), kuras mērķis bija poļu flotes celtniecības propaganda. Tas var likties dīvaini, jo tai brīdī Polijai, kura faktiski vēl nemaz neeksistēja kā valsts, nebija arī izejas pie jūras (Tā radās tikai pēc t.s. ``Poļu koridora'' izveides 1919.~gadā). Taču ticība šādai notikumu attīstībai pastāvēja. 1919.~gadā K.~Porembski pieņēma dienestā jau Polijas Kara ministrijas jūras sekcijā (pēc tam~--- departamentā). Vēlāk biedrību pārdēvēja par ``Poļu navigācijas līgu'' (\pltxti{Liga Żeglugi Polskiej}). 1924.~gadā organizācija vēlreiz mainīja nosaukumu un sāka dēvēties par ``Jūras un upju līgu'' (Liga Morska i Rzeczna). Tā stādīja uzdevumu darboties jūras un upju flotes attīstības labā un popularizēt jūras problēmas sabiedrībā. Kā jau iepriekš minēts, 1921.~gadā tika uzsākta ostas paplašināšana un modernizācija Dancigai blakus esošajā Gdiņā. To prasīja ne tikai Polijas vēlme iekļūt lielvalstu (\pltxti{mocarstwo}) skaitā, bet arī nepieciešamība panākt ekonomisko neatkarību no citu valstu labvēlības ostu izmantošanā. 1928.~gada februārī tika izveidota ``Koloniālo pionieru savienība'' (\pltxti{Zwiazek Pionerow Kolonialnych}), kuras uzdevums bija propagandēt koloniju ieguvi Polijai. Šai pat gadā savienība kā atsevišķa sekcija iegāja ``Jūras un upju līgas'' sastāvā, kura 1930.~gadā tika pārdēvēta par ``Jūras un koloniju līgu'' (\pltxti{Liga Morska i Kolonialna}). To vadīja ģenerālis M.~Zaruskis. Organizācija centās palīdzēt nostiprināt Polijas floti, īpaši~--- Kara floti, kā arī cīnījās par koloniju ieguvi. 1933.~gadā valdība līgai piešķīra tiesības administrēt līdzekļu vākšanu ``Jūras aizsardzības fondam''. Divu gadu laikā ``Jūras un koloniju līga'' fondam savāca vairāk kā 8~miljonus zlotu. Tieši par fonda savāktajiem līdzekļiem notika Polijas zemūdenes ``\pltxti{Orzel}'' (Ērglis) izbūve Holandē.

1935.~gadā līgas izdotais žurnāls ``\pltxti{Morze}'' (``Jūra'') rakstīja: ``Mēs, poļi, tāpat kā itāļi, stāvam ātri pieaugošā iedzīvotāju skaita izvietošanas un izmantošanas problēmas priekšā. Mums, poļiem, tāpat kā itāļiem, ir tiesības prasīt, lai mums tiktu atvērti tirgi eksportam un reģioni pārceļošanai, lai mēs varētu iegūt izejvielas, nepieciešamās nacionālajai tautsaimniecībai, ar noteikumiem, ar kādiem tās saņem visas koloniālās lielvalstis.'' 1936.~gada septembrī Tautu Savienības sēdē Polijas ārlietu ministrs J.~Beks nolasīja aicinājumu paplašināt Savienības Mandātu komisijas, kas lēma jautājumus par Vācijai un Osmāņu impērijai atņemto koloniju izmantošanu, biedru skaitu, lai visas ieinteresētās valstis varētu saņemt savu daļu. Vācijas diplomāti ziņoja uz Berlīni, ka Polijas vadība izskata iespējamību sakarā ar to, ka Polija bija ieguvusi daļu Vācijas teritoriju, griezties pie Tautu Savienības ar priekšlikumu nodot Polijai daļu (līdz 9\%) bijušo Vācijas koloniju~--- Togo un Kamerūnu, kuras ``tāpat nevienam nav vajadzīgas''. Tai pat laikā ``Jūras un koloniju līga'' sāka organizēt ``jūras nedēļas''. Tās norisa ievērojamāko Polijas personību aizbildniecībā. 1937.~gadā tādas bija ģenerālis K.~Sosnkovskis, kurš kļuva par Līgas protektoru, Valsts prezidents I.~Moscickis (kļuva par Līgas goda biedru), maršals E.~Ridzs-Smiglijs, Polijas primass (1926--1948) A.~Hlonds. Nedēļa tika rīkota ar devīzi ``Mums ir vajadzīga stipra flote un kolonijas''.

1938.~gadā Polija blakus ``Jūras nedēļai'' plaši atzīmēja arī speciālu Koloniju dienu. 18.~aprīlī Varšavā un daudzās citās pilsētās notika demonstrācijas ar prasību ``piešķirt dižajai poļu nācijai kolonijas''. Dievnamos šai sakarā notika svinīgi aizlūgumi. Kinoteātros tika rādītas filmas par koloniju tēmu. 1939.~gada 10.~februārī, kad Gdiņā tika nolaista ūdenī zemūdene ``\pltxti{Oržel}'', ģenerālis K.~Sosnkovskis uzsvēra flotes nozīmi nākamo Polijas koloniālo valdījumu aizsardzībai. 1939.~gada 11.~martā Polijā tika publicēta programma koloniālajā jautājumā, kurā bija deklarēts, ka ``Polijai, tāpat kā citām Eiropas valstīm, jābūt pieejamām kolonijām''. Ar 1937.~gadu Polijas koloniālās pretenzijas plaši izskanēja arī aiz valsts robežām, taču Rietumu valstis nesteidzās ievērot un vēl jo vairāk atbalstīt Polijas pretenzijas. Galvenie Eiropas un ASV preses izdevumi veltīja maz uzmanības Polijas koloniju iegūšanas centieniem un Tautu Savienība nereaģēja uz Polijas tīkojumiem iegūt daļu agrāko Vācijas koloniju. Tomēr daļa Polijas iedzīvotāju ar entuziasmu uztvēra valsts vadības ambīcijas. Ja 1938.~gada sākumā ``Jūras un koloniju līgai'' bija ap 690~tūkstošu biedru, tad gada beigās pēc plašajiem propagandas pasākumiem~--- jau 841~tūkstotis, bet pirms 1939.~septembra~--- jau gandrīz miljons.

Var saskatīt zināmu saistību starp koloniālajiem centieniem ar Polijas nostāju pret Tautu Savienībā 1938.~gada janvārī diskutējamo Savienības statūtu 16.~punktu par kolektīvo drošību. Virkne valstu, kuras vēlējās iegūt brīvas rokas iespējamo starptautisko konfliktu un militāro sarežģījumu gadījumā prasīja tā atcelšanu. Starp šīm valstīm bija arī Polija.

Tai pat laikā Polijas vadība, izrādot lielvalstiskas pretenzijas, maz domāja par attiecību uzlabošanu ar savu austrumu kaimiņu~--- PSRS. Polijas ārlietu ministra vietnieka J.~Šembeka, kurš 1937.~gada 16.~februārī bija klāt Polijas armijas ģenerālinspektora maršala E.~Ridzs-Smiglija sarunā ar H.~Gēringu, piezīmēs lasām: ``Maršals \citespace{} nedomā, ka Padomes perinātu agresīvu plānu. Taču viņam ir iespaids, ka kādu sarežģījumu gadījumā Eiropā, Padomes izmantotu šo gadījumu, lai veicinātu sabiedrības nemierus vājākajā un pretoties mazāk spējīgākajā posmā.'' Tālāk sekoja ar E.~Ridzs-Smiglija roku pierakstīta piezīme: ``Nav iespējams iedomāties, ka jebkāda konflikta gadījumā Polija varētu nostāties boļševiku pusē''. Var secināt, ka Polijas valstsvīrs diezgan precīzi vērtēja esošo un prognozēja nākamo PSRS vadības nostāju, taču viņa piezīme liecina arī par Polijas vadītāju nespēju pacelties augstāk pār ideoloģisko nepatiku pret boļševismu un Krieviju kā tādu. (Minētajās sarunās H.~Gērings savam sarunu partnerim paziņoja, ka lielas briesmas ne tikai Polijai, bet arī Vācijai draud no austrumiem, šīs ``briesmas rada ne tikai boļševisms, bet arī Krievija kā tāda, neatkarīgi no tā vai pastāv tajā monarhiska, liberāla vai kāda cita iekārta.) Tas arī atšķīra Polijas vadības pozīciju, no tādu Rietumu demokrātiju vadītāju kā Lielbritānijas vēlākais premjers V.~Čērčils un ASV prezidents Fr.~D.~Rūzvelts nostājas, kuri Otrajā pasaules karā, domājot pirmkārt par savu tautu interesēm, bija gatavi un arī noslēdza savienību ar diktatoru J.~Staļinu. Par to liecināja arī 1937.~gada 31.~augustā Polijas Ģenerālštāba izstrādātā direktīva N~2304/2/37, kas atklāja Polijas politikas mērķus pret PSRS, kurus tā mēģināja realizēt ar PSRS dzīvojošo tautu radikālnacionālistu, apvienotu klubā ``Prometejs'', palīdzību. Direktīva runāja par ``Krievijas apspiestajām tautām'', paredzēja, ka ````Prometejs'' visiem spēkiem veicina centienus, virzītus uz to, lai izsauktu nacionālu revolūciju PSRS teritorijā''. Kā redzam, 20.~gados aizsāktais ``prometeisma'' kurss turpinājās arī tad, kad tiešas briesmas draudēja jau pašai Polijai. Vēl var piebilst, ka 1939.~gadā, kad Sarkanā armija iegāja Polijas teritorijā, padomju čekisti tur atrada poļu izlūkdienesta dokumentus, kuri saturēja datus par 186~tā darbiniekiem un aģentiem, kuri darbojās PSRS teritorijā.

Var šķist ka Vācijas mēģinājumiem realizēt Austrijas anšlusu bija jāsatrauc valdība Varšavā, bija jāvirza to atmest J.~Pilsudska formulu par ``līdzsvaru'', kura tolaik jau kļuva par acīmredzamu anahronismu. Taču pēc Polijas vadības uzskatiem visa līdzšinējā starptautisko attiecību attīstība Eiropā neatstāja vietu pieņemamai alternatīvai. Polijas līderi uzskatīja, ka PSRS nevar kļūt par iespējamu sabiedroto, Francijas kļūst arvien vājāka, ko parādīja arī Parīzes piekrišana Austrijas anšlusam. Lielbritānijai par Poliju bija ierobežota interese. Kā uzskatīja Varšavā, palika tikai viens ceļš: cik vien iespējams, uzturēt attiecības ar Vāciju iepriekšējā līmenī, un mēģināt iegūt maksimālu labumu no tām.

Polijas valdība bija gatava atbalstīt Vācijas pretenzijas pret Austriju, pretī saņemot Vācijas atbalstu poļu plāniem pret Lietuvu. (J.~Beks 1937.~gadā vairākkārt izteicās, ka ``normālu attiecību'' nodibināšanai starp Poliju un Lietuvu galvenokārt traucē ``padomju ietekme uz Kauņu'', kas vājinot Lietuvas pozīcijas pret Vāciju Mēmeles (Klaipēdas) jautājumā.) Saprašanās tika noformēta ar J.~Beka paziņojumu Vācijas reihstāga priekšsēdētājam un virknes citu amatu izpildītājam H.~Gēringam 1938.~gada janvārī un H.~Gēringa paziņojumu Polijas sūtnim Berlīnē J.~Lipskim martā. Kad Vācija 1938.~gada 14.~februārī Polijai bija paziņojusi par Austrijas anšlusa gatavošanu, Polija varēja plānot arī savus soļus pret Lietuvu. Kad 1938.~gada 12.--13.~martā Vācija veica Austrijas \emph{anšlusu}, arī Polija izmantoja situāciju un demonstrēja savu ``varenību''.

Kā ieganstu izmantojot līdz galam nenoskaidrotu incidentu uz Lietuvas robežas, kurā dzīvību zaudēja viens no poļu robežsargiem (viņš, pēc Polijas puses oficiālā ziņojuma, it kā sekojot robežpārkāpējam, bija 3 metrus iegājis Lietuvas teritorijā), uz robežas tika koncentrēts Polijas karaspēks.

Radās iespaids, ka Polija uzskata Lietuvu kā kompensāciju par Polijas piekrišanu Austrijas ``anšlusam''. Poļu žurnāls ``\pltxti{Przeglad Powszechny}'' (``Vispārējs Apskats'') rakstīja: ``Polijas sabiedriskā doma instinktīvi nojauta, ka sakarā ar ``anšlusu'' mums jāsaņem kaut kāda kompensācija. \citespace{} Pēc teritorijas un iedzīvotāju skaita Lietuva gan būtu Polijai visai nenozīmīga kompensācija, taču kvalitatīvi, ar savu ģeopolitisko stāvokli, Lietuva ir visai vērtīga.''

Pārējās Baltijas valstis: Latvija un Igaunija šajā konfliktā turējās sāņus, nevis atbalstīja Lietuvu, bet faktiski nostājās Polijas pusē. Tā 1938.~gada 12.~martā Latvijas ārlietu ministrs Vilhelms Munters izsauca Lietuvas sūtni Rīgā un paziņoja: ``Austrijas notikumi mani pamudina, un šoreiz jau oficiāli, valdības vārdā, vērst Lietuvas uzmanību uz draudošām briesmām, kas var ļoti drīz pienākt arī jau pie Lietuvas robežām, un ieteikt šai sakarībā nopietni apsvērt jautājumu par attiecību nokārtošanu ar Poliju.''

Lietuva varēja cerēt vienīgi uz Padomju Savienību. 16.~martā PSRS ārlietu tautas komisārs M.~Ļitvinovs izsauca Polijas sūtni Maskavā V.~Gžibovski un lūdza informēt Polijas valdību, ka, ievērojot stāvokļa nopietnību, Padomju valdība spiesta norādīt Polijas valdībai, ka Padomju Savienība ir ieinteresēta saglabāt Lietuvas valsts pilnīgu~--- ne tikai nominālu, bet arī faktisku~--- neatkarību un tāpēc nevarētu palikt vienaldzīga pret notikumiem uz Polijas un Lietuvas robežas, ja Lietuvas neatkarība tiktu apdraudēta.

Tomēr \strong{Polijas valdība} 1938.~gada naktī no 16.~uz 17.~martu \strong{iesniedza ultimātu Lietuvai} ar prasību nodibināt diplomātiskas attiecības, garantēt poļu mazākuma tiesības Lietuvā, atzīt \latxti{status quo}, atcelt Lietuvas konstitūcijas paragrāfu, kurš Viļņu pasludināja par Lietuvas galvaspilsētu. Polijai tas nozīmēja savu pozīciju nostiprināšanu un izredžu uz Polijas vadītas Austrumeiropas valstu koalīcijas palielināšanos. Ultimāta noraidīšanas gadījumā Polija draudēja Lietuvu okupēt.

Tas viss notika tikai pāris dienu pēc Vācijas realizētā Austrijas ``anšlusa'', vēl priekšā bija Minhenes vienošanās un Čehoslovākijas iznīcināšana, Itālijas realizētā Albānijas okupācija, bet Polija jau rīkojās kā valsts, gatava tūlīt uzsākt agresiju pret kaimiņvalsti~--- Lietuvu. Tomēr Polijas valdībai nācās rēķināties ar Vācijas ieinteresētību Mēmeles (Klaipēdas) pievienošanā. 17.~martā Polijas sūtnim Berlīnē J.~Lipskim tika dots norādījums informēt H.~Gēringu par Polijas valdības gatavību ievērot Vācijas intereses ``iespējamās akcijas pret Lietuvu'' kontekstā. Ā.~Hitlers bija gatavs atzīt Lietuvas sagrābšanu, izņemot Mēmeles pilsētu ar apkārtējo apgabalu, kuriem bija jāietilpst Vācijā. Tika plānots, ka poļu un vācu karaspēks attiecīgajos Lietuvas rajonos ieies vienlaicīgi.

Varšavā un dažās citās Polijas pilsētās ar valdības sankciju notika pret Lietuvu vērstas demonstrācijas, kuru dalībnieki izvirzīja prasību doties ``gājienā uz Kauņu''. Viņi aicināja maršalu E.~Ridzs-Smigliju: ``\pltxti{Wodzu, prowadź na Kowno!}''(Vadoni, ved [mūs] uz Kauņu!). Dažas labējās poļu avīzes izteica domu, ka Polijai jāiegūst Lietuva kā kompensācija par Vācijas veikto Austrijas \emph{anšlusu}. Vienīgi komunisti uzdrošinājās konsekventi uzstāties pret iespējamo Polijas agresiju pret Lietuvu. 1938.~gada martā PKP CK aicinājumā tautai bija teikts: ``Polijas aizstāvēšanas vietā mēs īstenībā redzam mēģinājumu terorizēt vājo lietuviešu tautu un iekalt viņu nebrīves ķēdēs.'' Tai pat laikā citā PKP CK aicinājumā izskanēja brīdinājums: ``\dots{}pašreiz Polijas iesaistīšanās jebkurā, pat uzvarošā, karā kopā ar hitlerisko Vāciju vai bez tās pret Lietuvu, Čehoslovākiju, Padomju Savienību vai citu mierīgu valsti nostādīs mūsu zemi pilnīgā atkarībā no Hitlera, kā arī var kalpot Hitleram kā izejas moments lai sagrābtu ``koridoru'' un [Polijas] rietumu zemes.'' Taču komunistu balsis nomāca valdības diriģētā propaganda. Darbojās senā imperiālistisko valstu loģika apmierināt savas intereses uz vājāko rēķina. Var tikai minēt, kā būtu attīstījušies notikumi, ja ultimāts tiktu noraidīts.

Lietuva bija militāri daudz vājāka par Poliju, tai nebija starptautiska atbalsta. Vienīgi PSRS ārlietu tautas komisārs M.~Ļitvinovs 1938.~gada 18.~martā Polijas sūtnim V.~Gžibovskim norādīja, ka ultimāts prasa Lietuvai ne tikai nodibināt diplomātiskas attiecības ar Poliju, bet arī tās atteikšanos no sava viedokļa Viļņas apgabala un citos jautājumos, bet tādas prasības ir līdzvērtīgas Lietuvas izvarošanai. PSRS valdības pārstāvis Polijas sūtnim paziņoja, ka tā iesaka Polijai neapdraudēt Lietuvas brīvību un neatkarību, pretējā gadījumā varētu tikt denonsēts Polijas un PSRS neuzbrukšanas līgums.

Desmitgadēm ilgi padomju vēsturnieki centās radīt iespaidu, ka tikai PSRS iejaukšanās nepieļāva to, ka Lietuvu sagrābtu Polija un Vācija. (1938.~gadā tā domāja arī daļa citu valstu politiķu. Piemēram, Čehoslovākijas ārlietu ministrs K.~Krofta uzskatīja, ka PSRS iejaukšanās izglāba Lietuvu no pievienošanas Polijai.) Turpretī mūsdienu krievu vēsturnieks S.~Slučs uzskata, ka PSRS vadība, apzinoties savu mazo starptautisko ietekmi, baidījās iesaistīties kādā konfliktā un demonstrēja savas ``līdzdalības šķietamību'' (``\rutxti{видимость участия}''), kura ne tikai neapturēja sekojošos notikumus, bet arī negatīvi ietekmēja PSRS iespēju vērtējumu pasaulē, kam bija nopietnas sekas tuvojošās Čehoslovāskijas krīzes laikā. Domājams, ka minētajam vēsturniekam ir liela daļa taisnības.

Lai kā arī nebūtu, Lietuva neuzskatīja PSRS atbalstu par pietiekamu un bija spiesta piekāpties, nodibināt ar Poliju diplomātiskas attiecības, apmainīties sūtņiem. Polija aprobežojās ar šīs prasības izpildi, atteicās no iebrukuma. Latvijas diplomāts Aleksandrs~Birznieks domāja: ``Tagad paredzams, poļiem atklāsies plašas iespējas ``legālā'' ceļā iespiesties Lietuvā, izplatīt tur savu iespaidu un ar savu nospiedošo fizisko un garīgo pārspēku galu galā panākt savu mērķi~--- Lietuvas pilnīgu pakļaušanu Polijai.'' Paredzējums gan nepiepildījās, jo Polija pati drīz kļuva par agresijas upuri un Eiropas mērogā būtiskas nozīmes tās gribas uzspiešanai Lietuvai pirmskara situācijā nebija.

Nākamais Polijas pretenziju objekts bija Čehoslovākija, ar kuru jau ilgstoši pastāvēja domstarpības par Cešinas apgabala piederību. Jau 1934.~gada rudenī pēc Polijas~--- Vācijas neuzbrukšanas līguma noslēgšanas Polija veica militārus manevrus, kuros gatavoja karaspēku Čehoslovākijas sabrukuma gadījumam. 1936.~gadā Polija atkal atdzīvināja savas pretenzijas Čehoslovākijai, paziņoja, ka vēlas iegūt Cešinu, panākt Slovākijas neatkarību. Tās teritorijā tika iesūtītas diversantu vienības, kas spridzināja dzelzceļus, nogalināja policistus. Diversijas notika arī Aizkarpatu Ukrainas teritorijā, kura piederēja Čehoslovākijai un nekad nebija bijusi etniska poļu teritorija. 1938.~gada 26.~martā Cešinā tika radīta ``Poļu savienība'', kuras mērķis bija atdalīt šo apgabalu no Čehoslovākijas un kura darbojās līdzīgi kā K.~Henleina vadītā nacistiskā Sudetu Vācu partija.

Bez tam, kā rakstīja krievu vēsturnieks V.~Volkovs, Polija Čehoslovākijā saskatīja galveno šķērsli plāniem pierādīt sevi kā lielvalsti un apvienot savā vadībā Donavas baseina valstis. No šīs pozīcijas arī izrietēja bezkompromisa nostāja pret Čehoslovākiju līdz pat tās valstiskuma iznīcināšanai. Uz pret Čehoslovākiju vērsta pamata Polija attīstīja attiecības ar Ungāriju, kura arī pretendēja uz slovāku zemēm.

Ievērojams kreisā spārna publicists O.~Dluskis jau 1937.~gada beigās Parīzē izdotā rakstā, brīdinot, ka Polijā valdošais slānis vienīgo izeju no krīzes redz militāru avantūru ceļā un atsaucoties uz it kā J.~Beka inspirētu labējā publicista A.~Bohenska grāmatu ``\pltxti{Między Niemcami a Rosją}'' (Starp Vāciju un Krieviju, 1937), izklāstīja Polijas valdības mērķus: atraut no Čehoslovākijas slovākus, viņu ``neatkarīgajai'' valstij balstoties uz Poliju, panākt Aizkarpatu Ukrainas piedalīšanu Ungārijai, bet Sudetu apgabala~--- Vācijai. O.~Dluskis tālredzīgi brīdināja: ``Šādas politikas vārdā poļu valdošie slāņi ir pilni apņēmības ziedot poļu tautu un tās neatkarību. Jo ir pilnīgi skaidrs, ka aiz šiem Čehoslovākijas sadalīšanas plāniem, kurus ``sanācija'' ir sagatavojusi saskaņā ar hitlerisko Vāciju un par kuriem tā jau ir samaksājusi ar faktisku atteikšanos no Gdaņskas [pēdējais nekorektais apgalvojums bija pieņemts kreiso aprindās~--- V.Š.], slēpjas Hitlera bruņotā dūre, kura iebrucēja apetīte sniedzas ne tikai uz Čehoslovākiju, bet arī uz Poliju''.

Arī Polijas Komunistiskās partijas 1938.~gada aprīlī izdotajā 1.~Maija uzsaukumā tautai bija uzsvērts: ``Tai laikā, kad Hitlers gatavojas lēcienam pret brālīgo Čehoslovākiju, lai, saņemot Poliju spīlēs, uzsāktu mūsu zemes sadali, sanācijas-endeku targoviešu [no Targovices konfederācijas nosaukuma, kura 1792.~gadā pasludināja Četrgadu seima pieņemto t.s. 3.~maija Konstitūciju par spēkā neesošu un ar krievu karaspēka palīdzību vērsās pret karali un Seimu, targoviešu nosaukums tika lietots kā dzimtenes nodevēju sinonīms~--- V.~Š.] banda izvērš niknu šovinistisku kampaņu pret Čehoslovākiju, nekaunīgi šantažē brālīgo čehu tautu.''

Latvijas diplomāts A.~Birznieks jau 1938.~gada 19.~martā Latvijas Ārlietu ministrijai ziņoja: ``Pirmais, ko man nācās dzirdēt sarunās ar poļiem pēc Vācijas iebrukuma Austrijā, bija tas, nu esot pienācis laiks pārmācīt čehus. Ja tikai Vācija kaut kādā veidā uzmākšoties čehiem, Polija nekavējoties brukšot čehiem virsū, lai atņemtu viņiem Cešinu un vēl kādas tur čehu pilsētiņas un apgabalus, kur dzīvojot poļi. Šis naids pret čehiem ir tik vispārīgs un dziļš visā poļu sabiedrībā, ka liekas pietiktu ar visniecīgāko iemeslu, lai izsauktu asiņainu izrēķināšanos ar savu brāļu tautu.''

Var jau apšaubīt A.~Birznieka sprieduma objektivitāti par poļu dziļo naidu pret čehiem, taču jau tāda vērtējuma rašanās vien uzskatāmi liecina, ka valsts izvērstā propaganda pret tām tautām, ar kurām Polijai bija teritoriāli strīdi, bija visai efektīva. Tāpēc arī Polijas vadība bez svārstībām varēja izmantot Čehoslovākijas grūto stāvokli, kāds tai radās sakarā ar Vācijas plašo propagandu par grūto Sudetu vāciešu likteni.

Pie tam vācu un poļu propaganda šai gadījumā neatšķīrās no abu valdību centieniem. Jau 21.~martā J.~Beks intervijā Lielbritānijas presei izteicās, ka ``250~tūkstošiem poļu'' Čehoslovākijā pienākas tāda pat autonomija, kādu prasa Sudetu vācieši.'' 22.~martā viņš apvainoja Čehoslovākijas valdību Komunistiskās Internacionāles biroja darbības pieļaušanā Prāgā. (Šai sakarā var tikai piebilst, ka kopš Oktobra apvērsuma Krievijā daudzu pilsonības politiķu briesmīgākais apvainojums pretiniekiem bija pārmetums sadarbībā ar komunistiem, neatkarīgi no tā, vai šajos pārmetumos bija kaut grans patiesības. Tāpat rīkojās arī Polijas ārlietu ministrs.)

Polijas diplomātija 30.~gados attiecībās pret Čehoslovākiju demonstrēja sīkumainību un nespēju saprast, ka draudzīgas attiecības, savienība ar to varētu kļūt arī par pašas Polijas drošību stiprinošu faktoru. Abu valstu savienība, pat ja tai nebūtu Francijas atbalsta, varētu nostādīt diezgan ievērojamu spēku vācu agresijas ceļā. Par sliktākajām savstarpējām attiecībām nevar pilnībā attaisnot arī Čehoslovākijas pusi, taču Polijas vaina ir lielāka kaut vai tāpēc, ka tā pretendēja uz lielvalsts statusu, tātad tai arī bij’ jārīkojas tālredzīgāk, jo ``kam vairāk dots, no tā arī vairāk tiek prasīts''. Kā norāda vācu vēsturnieks G.~Rode, J.~Beks kalkulēja izmantot gaidāmo Čehoslovākijas sabrukumu, lai iegūtu Polijai Slovākiju kā sabiedroto un pie reizes arī strīdīgo Cešinas (čehu \cstxti{Těšín}, vācu \detxti{Teschen}, poļu \pltxti{Cieszyn}) apgabalu.

Polijas propaganda apgalvoja, ka ar Cešinas atņemšanu Čehoslovākijai tiktu atjaunots vēsturiskais taisnīgums un apgabalā dzīvojošie poļu iedzīvotāji apvienotos ar tēvzemi. (Šai sakarā var piebilst, ka daudzus poļu vēsturniekus šokējošā sakarā ar Polijas sagrāvi 1939.~gadā izteiktā V.~Molotova frāze par to, ka Polija bija Versaļas sistēmas kropls bērns, faktiski tikai atkārtoja J.~Pilsudska frāzi par ``mākslīgi un kroplīgi radīto Čehoslovākijas Republiku''.)

Jāsaka, ka kaimiņvalstīs dzīvojošo poļu mazākumu nacionālo jūtu izmantošana pret ``apspiedējiem'', arī tad, ja tie realizēja daudz demokrātiskāku politiku nekā poļu valdītāji, kā tas bija Čehoslovākijas gadījumā, labi raksturo ``sanācijas'' vadītāju cinismu. ``Patriotisma'' (faktiski~--- šovinisma) uzplūdos arī daļa opozīcijā esošo Polijas politisko partiju biedru padevās ``panākumu reibonim'', neprotestēja pret savas valsts nostāšanos agresoru rindās. Turpretī demokrātiskās Čehoslovākijas autori poļu mazākuma problēmu apgaismoja nevis vietējo varas iestāžu politikas kontekstā, kā to darīja poļu propagandisti, bet gan salīdzināja to tiesības ar tām, kādas bija nacionālajiem mazākumiem, kuri dzīvoja citās Eiropas valstīs. No tā izrietēja, ka poļu mazākuma stāvoklis Čehoslovākijā bija daudz labvēlīgāks nekā, piemēram, poļu mazākumam Vācijā, ukraiņu Ungārijā vai tai pašā Polijā.

Kad Slovākijas Tautas partijas vadonis ksendzs A.~Hlinka 1937.~gadā viesojās Polijā, viņš tika apbalvots ar augstu Polijas ordeni. 1938.~gada maijā Polijā demonstratīvi labvēlīgi tika uzņemta Amerikas slovāku delegācija, kura iestājās par Slovākijas autonomiju Čehoslovākijas sastāvā. (Tikai A.~Hlinkas nāve 1938.~gada 16.~augustā un Ungārijas ambīcijas pēc Slovākijas neļāva šiem plāniem piepildīties un Polijai vajadzēja samierināties ar Cešinas sagrābšanu.)

1938.~gada septembrī, kad Ā.~Hitlers pieprasīja Čehoslovākijai atdot Vācijai Sudetu apgabalu, Čehoslovākija cerēja uz Francijas un PSRS atbalstu, ar kurām tai bija līgums (1935) par savstarpēju palīdzību. Savukārt Polija neizrādīja gatavību atbalstīt savu sabiedroto Franciju, kura bija arī Čehoslovākijas sabiedrotā.

Jau 1938.~gada 22.~maijā, kad Francijas ārlietu ministrs Ž.~Bonnē izsauca pie sevis Polijas sūtni Parīzē J.~Lukaševiču, lai pajautātu, kāda būs Polijas politika Vācijas un Čehoslovākijas sadursmes gadījumā, pēdējais atbildēja: ``Mēs nekustēsimies no vietas''. (Francijas un Polijas 1921.~gada līgums nesaturēja nekādas abu pušu saistības kara gadījumā pret Čehoslovākiju. Ja Francija uzbruktu Vācijai, lai palīdzētu Čehoslovākijai, no Polijas viedokļa tā kļūtu par agresoru.) Ž.~Bonne painteresējās par Polijas attieksmi pret PSRS, uzsverot padomju atbalsta svarīgumu Čehoslovākijai tādas Polijas ``pasivitātes'' gadījumā. (PSRS varēja iesaistīties konfliktā Čehoslovākijas labā vai nu pēc Francijas pieprasījuma, izpildot savas Francijas-PSRS 1935.~gada 2.~maija līgumā noteiktās saistības, vai arī uz pēc Čehoslovākijas prasības pieņemta Tautu savienības lēmuma pamata, atbilstoši 1935.~gada 16.~maija PSRS-Čehoslovākijas līgumam.) Polijas sūtnis atkal bija kategorisks: ``Poļi vienmēr ir uzskatījuši krievus par ienaidniekiem \citespace{} un nepieciešamības gadījumā mēs ar spēku pretosimies jebkuram krievu iebrukumam mūsu teritorijā, ieskaitot krievu aviācijas lidojumus''. Tādu pat pozīciju Francijas sūtnim Varšavā L.~Noelam izteica maršals E.~Ridzs-Smiglijs, kurš paziņoja, ka Polija uzskata Krieviju, neraugoties uz to, kas tajā valda, par savu ``ienaidnieku numur viens''.

Ja Polijai būtu vēlēšanās, tā varētu izjaukt gatavojamo savas kaimiņvalsts Čehoslovākijas sadalīšanu. Polijai tikai vajadzēja atklāti un stingri nostāties pret agresiju, vai pat tikai, vadoties no sava kā Francijas sabiedrotās pienākuma, izteikt labvēlīgu neitralitāti gatavojamās agresijas upurim, publiski neuzsvērt, ka tā caur savu teritoriju neizlaidīs Sarkano armiju palīgā Čehoslovākijai. Pat neielaižot Sarkano armijai savā teritorijā, bet atļaujot to izdarīt savai sabiedrotajai Rumānijai, Polija varētu sekmēt Ā.~Hitlera centienu apvaldīšanu. Lieta tā, ka Francijas sūtņa Buharestē spiediena rezultātā Rumānijas ārlietu ministrs N.~Petresku-Komnens piekrita variantam: ``Rumānija nevar izlaist Sarkano armiju, taču, ja padomju lidmašīnas lidos augstu virs Rumānijas, tās nevarēs redzēt''. Taču iejaucās Polija. Francijas ārlietu ministrs Ž.~Bonnē varēja tikai nopūsties: ``Rumāniju šai jautājumā saista tikai Polijas iebildumi.'' Kad viņš centās pārliecināt Polijas sūtni, ka ``ja Polija nevēlas nekādi palīdzēt Čehoslovākijai, tad lai vismaz netraucē Rumānijai'', Polijas sūtnis lika saprast, ka ``Polija arī tam nepiekritīs un Rumānija bez tās nevar pieņemt nekādu lēmumu''. No Polijas pozīcijas, no tās gatavības ļaut Sarkanajai armijai nonākt saskarē ar Vācijas armiju bija atkarīga arī Francijas gatavības sniegt palīdzību Čehoslovākijai. Sekojot savai sabiedrotajai Francijai, pret Vāciju būtu spiesta uzstāties arī Lielbritānija.

Kā raksta kanādiešu vēsturnieks M.~Dž.~Karlejs, Polijas valdības viedoklis bija tāds, ka jebkura Francijas agresīva darbība vai padomju karaspēka pārvietošanās uz Čehoslovākiju, teiksim, caur Rumānijas teritoriju, var likt Polijai uzstāties nacistiskās Vācijas pusē. Polijas valdītāji bija gatavi atteikties no savienības ar Franciju, lai tikai iegūtu laupījumu uz Čehoslovākijas rēķina un neļautu PSRS tai palīdzēt.

Francijā bija darbinieki, kuri labi apzinājās, ka uz Poliju nevar paļauties. ``Ja Polija nevēlēsies, sekojot Francijas paraugam, iesaistīties militārā savienībā ar Padomēm \citespace{} brīnišķīgi, mēs to izdarīsim arī bez [Polijas]''~--- izteicās ģenerālis M.~Veigāns 1933.~gadā. ``\frtxti{Tant pis pour la Pologne}'' (Neapskaužams būs Polijas liktenis), ja poļu valdība Čehoslovākijas konfliktā nostāsies nacistu pusē, 1938.~gadā sacīja Francijas armijas Ģenerālštāba priekšnieks M.~Gamelēns. Arī Francijas premjerministrs E.~Daladjē 1938.~gada 25.~maijā padomju sūtnim Francijā J.~Suricam atzinās, ka ``ne tikai nevar rēķināties ar poļu atbalstu, bet nav pārliecības vai viņi nedos triecienu no aizmugures'', ka, ``pretēji poļu solījumiem, viņš netic poļu lojalitātei pat tieša Vācijas uzbrukuma gadījumā Francijai''. Taču kad izraisījās krīze, Francijas valdība nespēja sekot šiem loģiskajiem secinājumiem. Notika tas tāpēc, ka, kā brīdināja sūtnis L.~Noels, ja Francija izdarītu spiedienu uz Poliju, prasot piešķirt Sarkanajai armijai tiesības iziet caur tās teritoriju, vai denonsētu Francijas-Polijas vienošanās, Polija noteikti nostātos Ā.~Hitlera pusē. Francijai izdevīgāk bija nelikt Polijai izvēlēties starp PSRS un Vāciju, jo pēc L.~Noela viedokļa izvēli nebija grūti paredzēt.

Neraugoties uz to, ka Francija bija noslēgusi savstarpējās palīdzības līgumus ar Čehoslovākiju un PSRS (1935), Francijas valdība nekad neizvirzīja jautājumu par aizsardzības problēmu kopīgu apspriešanu. Kad 12.~maijā PSRS ārlietu tautas komisārs M.~Ļitvinovs Francijas ārlietu ministram Ž.~Bonnē piedāvāja padomju un franču ģenerālštābiem sākt sarunas par tehniskiem jautājumiem, ieskaitot jautājumu par padomju karaspēka izlaišanu cauri Polijai un Rumānijai, Francijas pārstāvis par šo priekšlikumu solīja ziņot savai valdībai. Taču atbilde tā arī nesekoja. Kad PSRS ārlietu ministrs M.~Ļitvinovs uzdeva jautājumu Parīzei, kāda būs ar Poliju militārā savienībā saistītās Francijas pozīcija, ja Polijas uzbrukuma Čehoslovākijai rezultātā sāksies PSRS un Polijas karš, Francijas atbilde nesniedzās tālāk par paziņojumu, ka tādā gadījumā Francija neuzskatīs par savu pienākumu palīdzēt Polijai. Franču valstsvīri un diplomāti nevēlējās noslēgt savienību ar PSRS. Francijas premjerministrs E.~Daladjē sliecās uz Vācijas ``nomierināšanu'', uzskatot, ka pēc kara darbības beigām, ``neatkarīgi no tā, kas uzvarēs, Francijā, tāpat kā Itālijā un Vācijā, neizbēgami notiks revolūcija. Padomju Krievija nepalaidīs garām iespēju izplatīt pasaules revolūciju uz mūsu valstīm\dots{}'' Francijas sūtnis Maskavā R.~Kulondrs uzskatīja, ka britu-franču-padomju alianse varētu uzvarēt karu pret Vāciju, taču tad Sarkanā armija sagrautu Poliju un padomju ietekme izplatītos Centrālajā Eiropā, iespējams arī uz Vāciju un pat Franciju. Acīmredzot bailes no šādas uzvaras pār nacistisko Vāciju, tāpat kā bailes no zaudējuma tai, paralizēja Francijas vadītājus.

Lielā mērā konflikta atrisinājums bija atkarīgs no Francijas sabiedrotās Lielbritānijas, bez kuras atbalsta Francija neuzdrošinājās patstāvīgi uzstāties pret Vāciju. Taču jau 1938.~gada martā Lielbritānijas valdības ārlietu komiteja bija pieņēmusi lēmumu, kuru britu Ārlietu ministrijas (\entxti{Foreign and Commonwealth Office}) pastāvīgais sekretārs A.~Kadogans savā dienasgrāmatā formulēja ļoti skaidri: ``Čehoslovākija nav pat viena britu grenadiera piešu vērta.'' 21.~maijā turpat viņš rakstīja: ``Nolemts, mums nav jākaro''. Rietumu lielvalstis nevēlējās atbalstīt agresijas objektu~--- Čehoslovākiju.

Padomju un krievu vēsturnieks Ļ.~Bezimenskis aprakstījis, kā 1968.~gadā, kad Lielbritānijā pēc 30~gadu noilguma termiņa tika atvērti arhīvi, N.~Čemberlena arhīvā tika atrasts dokuments, kuru 1938.~gada 30.~augustā bija sastādījis viņa palīgs H.~Vilsons un kurā runa gāja par sagatavotu plānu ``Z''. Par to bija informēti tikai daži cilvēki Lielbritānijas vadībā. Plāns paredzēja, ka rodoties ``asai situācijai'', Lielbritānijas premjeram N.~Čemberlenam pēkšņi jādodas pie Ā.~Hitlera un sarunās ar viņu jānovērš visi iespējamie konflikta ar Vāciju iegansti. Faktiski tas bija vēlāk realizētās Minhenes vienošanās stūrakmens. Jau 15.~septembrī N.~Čemberlens devās pie Ā.~Hitlera. 16.~septembrī notika Lielbritānijas valdības sēde šaurā sastāvā. Viss tika saskaņots. 17.~septembrī valdība pilnā sastāvā atbalstīja ``pašnoteikšanās principu''~--- tā divkosīgi tika nosaukta Sudetu apgabala atraušana Čehoslovākijai.

Mūsdienu krievu publicists J.~Muhins ar izsmieklu raksta par to, ka 70~miljoni cilvēku lielā Vācija ar kara draudiem ``nobiedēja'' Britu impēriju, kurā dzīvoja 532 miljoni cilvēku un Francijas koloniālo impēriju, kurā dzīvoja 109 miljoni iedzīvotāji tik lielā mērā, ka šīs valstis atteicās no iespējamās 170 miljonu lielās Padomju Savienības palīdzības iespējamā karā pret agresoru. Taču šāds vienkāršots vērtējums un spēku salīdzinājums nav korekts. Pietiek norādīt tikai uz vienu apstākli. Koloniju esamība bija ne tikai metropoļu spēka apliecinājums, bet lielā mērā arī to vājuma izraisītājs. Var piebilst, ka tikpat vienkāršojošs situāciju bija 1939.~gada oktobrī izskanējušais PSRS ārlietu tautas komisāra V.~Molotova skaidrojums par Otrā pasaules kara sākumu, proti,~--- Lielbritānija ar saviem 47~miljoniem iedzīvotāju pārvalda 480 miljonu cilvēku kolonijās, tāpat Francija ar 42 miljoniem cilvēku metropolē pārvalda 70~miljonu kolonijās. Bailes zaudēt savu pasaules valdītāju stāvokli tad arī likušas Francijas un Anglijas valdošajām aprindām sākt pasaules karu. Pēc autora ieskata tuvāks patiesībai ir cits vērtējums: bailes par savu koloniālo impēriju nākotni bija viens no cēloņiem, kuri lika Rietumvalstīm meklēt iespēju izvairīties no kara ar Vāciju.

1938.~gada 11.~augustā Polijas sūtnis J.~Lipskis tikās ar H.~Gēringu. Pēc sarunas viņš ziņoja Polijas ārlietu ministrijai: ``Sakarā ar krievu problēmu viņš vispārējos vilcienos teica, ka pēc čehu jautājuma atrisināšanas tā kļūs aktuāla. Viņš atgriezās pie domas, ka padomju-poļu konflikta gadījumā Vācija nevarētu palikt neitrāla, nesniedzot Polijai palīdzību.'' Šo frāzi droši varēja interpretēt kā apstiprinājumu tam, ka, ja PSRS mēģinātu aizstāvēt Čehoslovākiju, Vācija kopā ar Poliju karos pret PSRS. Francijas sūtnis Vācijā A.~Fransuā-Ponsē 1938.~gada 22.~septembrī Francijas Ārlietu ministrijai adresētā ziņojumā rakstīja: ``Polija un Ungārija pievienojas Vācijai Čehoslovākijas vajāšanā. Francija un Anglija, kuras mēģināja piekāpjoties un visādi apmierinot Vācijas prasības, glābt čehu valsts pastāvēšanu, izrādījās triju valstu, kuras centās panākt Čehoslovākijas sadali, vienotas frontes priekšā.'' Tādejādi Francijas sūtnis skaidri iezīmēja agresorus. Tie bija Vācija, Polija un Ungārija.

J.~Beks, lai iegūtu Polijai jaunas teritorijas, bija gatavs neievērot to iegūšanas ceļā draudošās briesmas pašai Polijai un iztēlot Ā.~Hitleru gandrīz vai par miermīlīgu politiķi. Tā, 1938.~gada augustā uz Igaunijas sūtņa Polijā H.~Markusa aizrādījumu, ka Čehoslovākijai sabrūkot, Polijas robežas ar Vāciju kļūs garākas un Polijas stāvoklis grūtāks, Polijas ārlietu ministrs atrunājās ar to, ka ``Hitlers ir tagad atturīgāks, mazāk agresīvs nekā agrāk.'' Franču vēsturnieks M.~Murins rakstīja, ka ``Pulkveža Beka dīvainā politika bija jauns nacionālistiskās orģijas, lielvalstiskā kompleksa izpaudums''.

Valdību atbalstīja konservatīvā prese, norādot, ka Čehoslovākija esot Polijai daudz kaitējusi un ``iesaistīties militārā avantūrā'' lai to aizstāvētu, būtu ``smaga, ļoti smaga kļūda''. Vēl vairāk~--- 1938.~gada 8.--11.~septembrī Volīnijā notika lieli Polijas armijas manevri, ar kuriem Polija demonstrēja, ka tā ir gatava apstādināt Sarkano armiju, ja tā dosies palīgā Čehoslovākijai. Manevros piedalījās Polijas bruņoto spēku ģenerālinspektors maršals E.~Ridzs-Smiglijs un kara ministrs ģenerālis T.~Kaspšickis. Pēc manevru noslēguma 19.~septembrī Luckā (poļu \pltxti{Łuck}, ukraiņu \uktxti{Луцьк}) notika, pēc poļu preses rakstītā, ``neatkarīgās Polijas vēsturē lielāka militārā parāde'', kura ilga septiņas stundas.

J.~Beks 1938.~gada 19.~septembrī Polijas sūtnim Berlīnē J.~Lipskim lika paziņot Vācijas valdībai, ka Polijas valdība ar savu nostāju ``paralizējusi padomju intervences iespējas plašākā nozīmē Čehijas jautājumā. Mūsu spiediens Buharestē radīja vajadzīgo iedarbību. Mūsu veiktos manevrus Volīnijā Maskava uztvēra kā brīdinājumu.'' (Savs pamats runāt par Polijas lomu ``padomju intervences'' paralizēšanā J.~Bekam bija: baltkrievu publicists A.~Tarass norādījis, ka Sarkanās armijas spēki rietumu virzienā 1938.~gada septembrī skaitliski gan bija pārāki pār vērmahtu, taču vājāki par Vācijas un Polijas apvienotajiem spēkiem.) Tālāk J.~Beka norādījumos sūtnim sekoja pasāžas: ``Polija uzskata Padomju iejaukšanos Eiropas lietās pār nepieļaujamu''; ``Čehoslovākijas Republiku mēs uzskatām par mākslīgu veidojumu, kurš atbilst dažām doktrīnām un kombinācijām, bet neatbilst Centrālās Eiropas tautu patiesajām vajadzībām un veselajām tiesībām''; ``pagājušās vasaras laikā Polijas valdība četras reizes atteicās pievienoties starptautiskām intervencēm Čehoslovākijas aizsardzībai''.

Uz to Ā.~Hitlers ar tā paša J.~Lipska starpniecību 20.~septembrī pamudināja Polijas valdību uzstādīt Čehoslovākijai arī savas prasības. Jau nākamajā dienā~--- 21.~septembrī tika iesniegta Polijas nota Čehoslovākijai, kurā Polijas valdība atgādināja savas pretenzijas uz Čehoslovākijai piederošo Silēzijas Cešinu. Ā.~Hitlers to komentēja atzinīgi: ``Poļi ir īsti vīri. Pilsudskis būtu ar viņiem apmierināts''. Ā.~Hitlers, acīmredzot cerēja uz ilgstošu sadarbību ar Poliju. Padomju aģents Tokijā R.~Zorge ziņoja uz Maskavu par to, ka no Vācijas militārā atašeja viņš saņēmis ziņas, kas atspoguļoja Ā.~Hitlera tā laika aprēķinus, proti: Vācijai pēc Sudetu jautājuma atrisināšanas nākamā būs Polijas problēma, taču tā starp Vāciju un Poliju tiks izlemta draudzīgi sakarā ar to abu kopīgo karu pret PSRS.

Turpretī PSRS jau no krīzes sākuma demonstrēja atbalstu Čehoslovākijai. 1938.~gada maija vidū Čehoslovākijas Komunistiskās partijas ģenerālsekretārs K.~Gotvalds apmeklēja Čehoslovākijas prezidentu E.~Benešu un paziņoja viņam par J.~Staļina izteikto gatavību sniegt militāru palīdzību Čehoslovākijai pat tai gadījumā, ja pretēji līgumam to nedarīs Francija, bet Polija un Rumānija atteiksies izlaist caur savu teritoriju padomju karaspēku. J.~Staļins uzsvēra tikai vienu šādas palīdzības priekšnoteikumu~--- Čehoslovākijai pašai bija jāaizstāvas un jālūdz padomju palīdzība.

2.~septembrī Francijas pilnvarotais lietvedis Maskavā apjautājās PSRS ārlietu ministram M.~Ļitvinovam, kādu palīdzību Krievija ir gatava sniegt Čehoslovākijai Vācijas uzbrukuma gadījumā, ņemot vērā grūtības, kas var rasties sakarā ar Polijas un Rumānijas neitralitāti. Atbilde skanēja, ka PSRS ir gatava izpildīt savas saistības, bet grūtības, vismaz saistībā ar Rumāniju, ir pārvaramas. Tautu Savienības Padomes vairākumam tikai vajadzēja nolemt, ka Čehoslovākija ir Vācijas agresijas upuris un tas varētu noteikt Rumānijas pozīciju jautājumā par padomju karaspēka izlaišanu caur savu teritoriju palīgā Čehoslovākijai. (Turpretī uz Francijas ārlietu ministra Ž.~Bonnē jautājumu, kādu palīdzību Francijai konflikta gadījumā var sniegt Lielbritānija, tika saņemta atbilde, ka pirmajos sešos kara mēnešos Lielbritānija palīgā Francijai varētu nosūtīt divas kājnieku divīzijas un 150~lidmašīnu. V.~Čērčils šai sakarā atzīmējis, ja Ž.~Bonnē meklēja ieganstus, lai atstātu čehus likteņa varā un viņa meklējumi vainagojās panākumiem.) 20.~septembrī PSRS vadība deva savam sūtnim Čehoslovākijā direktīvu: ``Uz Beneša jautājumu, vai Padomju valdība tiešām nekavējoties palīdzēs Čehoslovākijai, ja Francija izrādīsies tai uzticama un arī sniegs palīdzību, Padomju Savienības valdības vārdā varat dot apstiprinošu atbildi.''

21.~septembrī PSRS aizsardzības tautas komisārs K.~Vorošilovs izdeva direktīvu Kijevas sevišķa militārā apgabala Kara padomei veikt Sarkanās armijas daļu manevrus pie Polijas robežas. Tas, ka šie manevri nenotika reizē ar Polijas armijas manevriem, liecināja par Kremļa nevēlēšanos pieļaut incidentus, kuri varētu rasties, abpus robežai notiekot lieliem karaspēka manevriem, un ar to riskēt ar militāru konfliktu. Izpildot 22.~septembrī Čehoslovākijas ārlietu ministra K.~Krofta izteikto vēlējumu, Padomju Savienības ārlietu ministra vietnieks V.~Potjomkins 23.~septembrī izsauca Polijas pilnvaroto lietvedi PSRS un paziņoja, ka agresijas gadījumā pret Čehoslovākiju PSRS valdība var bez brīdinājuma denonsēt 1932.~gada Polijas un PSRS neuzbrukšanas līgumu. Nākas atzīt šāda brīdinājuma neefektivitāti, tas nemazināja Polijas apetīti. Polijas valdība rēķinājās, ka Lielbritānija un Francija pametīs Čehoslovākiju savam liktenim un tad Padomju Savienība nespēs neko uzsākt. Šoreiz poļu aprēķins attaisnojās. Padomju avīze ``\rutxti{Известия}'' (``Ziņas'') šai sakarībā gan norādīja, ka ne tikai Čehoslovākijā ir poļu nacionālais mazākums, bet Polijā ir daudz ukraiņu. Un kas notiks, ja ukraiņi pieprasīs plebiscītu, lai noteiktu savu nākotni?'' Taču vārdiem darbi nesekoja.

Padomju autori uzsvēra, ka PSRS vairākkārt izteica gatavību palīdzēt Čehoslovākijai, ja tā pati pretosies. Piemēram, krievu vēsturnieks M.~Semirjaga norādīja, ka septiņos mēnešos pirms Čehoslovākija galu galā piekāpās Vācijai, Padomju Savienība divdesmit reizes pa dažādiem kanāliem izteica savu gatavību izpildīt savu sabiedrotā pienākumu Čehoslovākijas priekšā un sniegt tai nepieciešamo militāro palīdzību. Taču pastāv arī diezgan pamatots viedoklis, ka, vairākkārt demonstrējot gatavību uzstāties pret Vācijas agresiju, PSRS maz ar ko riskēja. Pirmkārt, pēc 1935.~gada 2.~un 16.~maija līgumiem starp PSRS un Franciju, PSRS un Čehoslovākiju jautājuma izšķiršana faktiski bija atkarīga no Francijas. Otrkārt, PSRS militāras palīdzības sniegšanu Čehoslovākijai apgrūtināja tas, ka starp abām valstīm atradās Polija un Rumānija, bet šīs valstis nepiekrita Sarkanās armijas izlaišanai caur savu teritoriju. Par nesankcionētu Sarkanās armijas iziešanu caur poļu un rumāņu teritoriju nevarēja būt ne runas, jo tad PSRS varētu tikt atzīta par agresoru. PSRS diplomātu pūliņi, lai konflikta risināšanā iesaistītos Tautu Savienība un pieņemtu lēmumu par kolektīvu vēršanos pret agresiju, nedeva rezultātus. Krievu vēsturnieks M.~Narinskis atzīmējis, ka lielākā daļa krievu pētnieku, tāpat kā viņš pats, uzskata PSRS iejaukšanos Vācijas-Čehoslovākijas konfliktā bez pašas Čehoslovākijas lūguma un bez labvēlīgas Francijas pozīcijas par nereālu.

Polija pēc PSRS brīdinājuma saņemšanas pieprasīja no Berlīnes atbildi uz jautājumu, vai tā var rēķināties ar Vācijas atbalstu konflikta gadījumā ar PSRS. Tika saņemta pozitīva atbilde. 22.septembrī Ā.~Hitlers paziņoja Lielbritānijas premjerministram N.~Čemberlenam, ka pieprasa no Čehoslovākijas teritoriālu piekāpšanos ne tikai Vācijas, bet arī Polijas un Ungārijas labā. Citiem vārdiem, Polija pieprasīja savu laupījuma daļu. Uzstādot blakus Vācijas ultimātam savas prasības, tā kļuva par Vācijas sabiedroto pret Čehoslovākijas Republiku. Faktiski notikušo nākas traktēt kā nacistiskās Vācijas un ``sanācijas'' Polijas slepenu vienošanos par Čehoslovākijas sadali.

Ja līdz tam Polijas valdība tikai deklarēja gatavību nostāties Vācijas pusē, tad tagad, saistoties ar Vāciju, tā reāli parādīja gatavību par nelielu citas valsts teritoriju nonākt konfliktā ar PSRS, kļūt ne tikai par politisku, bet arī militāru Vācijas partneri. 24.~septembrī J.~Beks Lielbritānijas sūtnim Varšavā H.~V.~Kennardam paziņoja, ka Polija negrasās ``saistīt sev rokas'' Cešinas jautājumā, ``tai nav nekādu kareivīgu nodomu, bet tā nevar samierināties ar to, ka Vācijas prasības tiks īstenotas, bet Polija nesaņems neko''.

Interesanti, ka Lielbritānijas sūtnis Polijā tās Ārlietu ministrijas Politiskā departamenta direktoram T.~Kobiļanskim jau 1938.~gada rudenī ļoti pravietiski pateica, ka pēc Čehoslovākijas sadalīšanas būs pienākusi kārta arī ceturtajai [netiek rēķinātas Polijas dalīšanas XIX gadsimtā~--- V.Š.] pašas Polijas sadalīšanai. Arī PSRS Ārlietu tautas komisāra vietnieks V.~Potjomkins Vācijas sūtnim Maskavā paziņoja: ``Čehoslovākijas iznīcināšanā piedalījušās valstis vēl ļoti nožēlos, ka piekāpās kareivīgajam nacionālismam. Pienākusi rinda Polijai, jo tajā dzīvo ``ļoti daudz'' vāciešu\dots{}''

Lai piespiestu Čehoslovākijas valdību atdot strīdīgo teritoriju, poļu puse nevairījās arī no provokācijām. Pirmkārt, tika formēts brīvprātīgo leģions ``\pltxti{Zaolzie}'' (poļu \pltxti{Zaolzie}, čehu \cstxti{Zaolží}, Zaolší, vācu \detxti{Olsa-Gebiet}~--- ``aiz Olsas upes''~--- Cešinas Silēzijas austrumu daļa), kurā iestājās volontieri kā no Silēzijas, kā arī citām Polijas daļām. Ar to Polija demonstrēja, cik ļoti tā ir norūpējusies par čehu ``apspiestajiem'' poļiem Cešinā. Otrkārt, tika izveidota 36~tūkstošu vīru liela armijas grupa ``\pltxti{Śląsk}'' (``Silēzija''), kuru komandēja ģenerālis V.~Bortnovskis. Šis nopietnais militārais spēks bija paredzēts iebrukumam suverēnas valsts teritorijā, ja tā atteiktos brīvprātīgi padoties Varšavas spiedienam. Papildus šiem diviem redzamajiem ``argumentiem'' bija arī trešais, kurš tika paturēts slepenībā. Tās bija poļu diversantu grupas, kurām 1938.~gada 19.~septembrī tika noteikti divi uzdevumi. Pirmkārt, iekļūt Čehoslovākijas teritorijā un izvērst tur teroristisku darbību, spridzināt tiltus un dzelzceļu stigas, uzbrukt žandarmērijas posteņiem un patruļām, lai destabilizētu situāciju u.c. Otrkārt, diversantiem bija jāprovocē vietējie poļu tautības iedzīvotājus uz sacelšanos un tādējādi jāparāda (tāpat kā to darīja vācieši Sudetos), ka poļi nevēlas dzīvot ar čehiem vienā valstī.

Reizē ar notas iesniegšanu Polija koncentrēja karaspēku pierobežā. Naktī uz 1938.~gada 25.~septembri Koņske miestā poļu teroristi apmētāja ar rokas granātām un apšaudīja ēkas, kurās atradās čehoslovāku robežsargi, divi nami nodega. Pēc divu stundu ilgas apšaudes uzbrucēji atkāpās uz Polijas teritoriju. Analoģiskas sadursmes tai naktī norisa arī citās Cešinas apgabala vietās. Taču, ja poļu kaujinieki vairāk vai mazāk veiksmīgi tika galā ar pirmo uzdevumu, tad otro tiem paveikt neizdevās. Tāpēc Polijas vadībai atlika vienīgi militāri diplomātiskais ceļš, kā atrisināt strīdīgās Cešinas Silēzijas problēmu.

29.--30.~septembrī sekoja Lielbritānijas, Francijas, Vācijas un Itālijas vadītāju tikšanās Minhenē (Šai sakarā jāpasaka, ka J.~Beks pirms starptautiskas konferences sasaukšanas Čehoslovākijas jautājumā cerēja, ka tajā varēs piedalīties arī Polijas delegācija, tā paužot ilūzijas par savas valsts politisko nozīmi lielvalstu vērtējumā.) un tika panākta bēdīgi slavenā \strong{Minhenes vienošanās}, kuras rezultātā notika lielas daļas Čehoslovākijas teritorijas sadale starp Vāciju, Poliju un Ungāriju. Kad konferences dalībnieki pieprasīja, lai Čehoslovākija piekāpjas Vācijas prasībām, agresijas upuris neuzdrošinājās pretoties, kaut arī spēku samērs starp to un Vāciju bija gandrīz līdzīgs un kalnu apvidus bija labvēlīgs čehoslovāku aizsardzībai.

Vēlākais Čehoslovākijas prezidents L.~Svoboda par to rakstīja: ``Spēku samērs starp Čehoslovākijas Republiku un Vāciju 1938.~gadā bija sekojošs: 45~čehoslovāku divīzijas pret 47~nacistu; \citespace{} 1582~ čehoslovāku lidmašīnas pret 2500~ienaidnieka, 469~mūsu tanki pret 720~fašistu, 2~miljoni apmācītu čehoslovāku kareivju pret 2~200 tūkstošiem vāciešu. No šiem skaitļiem pat nespeciālistam ir skaidrs, ka līdz divkāršam pārsvaram, kurš nepieciešams uzbrūkošai operācijai, bija tālu.''

Minhenes vienošanās, kurā Čehoslovākijas liktenis tika izlemts bez tās pārstāvju klātbūtnes, kļuva par vienu no kliedzošākajiem un ciniskākajiem trešās valsts tiesību neievērošanas gadījumiem jaunāko laiku vēsturē. Latvijas ārlietu ministrs V.~Munters jau pāris mēnešus pēc līguma noslēgšanas reālistiski konstatēja: ``Viena lieta ir skaidra: kara iznākumā radītā politiskā sistēma Eiropā ir sabrukusi. \citespace{} Minhenes nolīgums ne tikai sankcionēja teritoriālas cesijas par labu Vācijai, bet ievadīja vēl tālejošākas teritoriālas pārmaiņas, kuru apmērs pat vēl nav galīgi noteikts un kuru politiskās, stratēģiskās un saimnieciskās sekas vēl nemaz nav pārredzamas.''

Lielbritānija sasniedza kāroto mērķi~--- karš nesākās. Radās perspektīva panākt vienošanos citos starptautiskos jautājumos ar nacistisko Vāciju.

Francija gan atvirzīja kara briesmas, taču šī vienošanās faktiski izjauca pirms tam Francijas radīto drošības sistēmu Austrumeiropā un t.s. Mazo Antanti (Čehoslovākijas, Dienvidslāvijas un Rumānijas savienību). Pārstāja darboties Francijas līgumi ar PSRS un Čehoslovākiju. Faktiski bruka Francijas un Polijas savienība.

Var jau teikt, ka ar Minhenes vienošanos Lielbritānija un Francija tomēr ieguva gadu atelpas, kurā varēja stiprināt savus spēkus, labāk sagatavojoties atvairīt draudošo Vācijas agresiju. Taču V.~Čērčils jau pēc Otrā pasaules kara, ar tā nestajām zināšanām un pieredzi bruņojies, deva sekojošu vērtējumu: ``\dots{}atelpas gads, kurš it kā tika iegūts Minhenē, nostādīja Angliju un Franciju salīdzinājumā ar hitlerisko Vāciju daudz sliktākā stāvoklī, nekā tas, kurā viņas atradās Minhenes krīzes brīdī.'' Tikpat negatīvi kā V.~Čērčils pret Minhenes vienošanos attiecās Francijas politiķi Ž.~Mandels un P.~Reino.

Čehoslovākija bija spiesta piekāpties apvienotā Vācijas, Itālijas, Lielbritānijas un Francijas spiediena priekšā.

PSRS saņēma pamatīgu triecienu prestižam, nonāca pilnīgā izolācijā. Ar to nerēķinājās pasaules politikas jautājumu lēmēji. Tā nevarēja vairs uzticēties Francijai, kura tikko bija nodevusi citu savu sabiedroto~--- Čehoslovākiju. Tika zaudēts militārās sadarbības mehānisms ar Rietumiem, kā arī labs partneris~--- Čehoslovākija.

Vācija guva jaunas teritorijas, nostiprināja nacistisko režīmu ekonomiski un politiski.

Polija ieguva teritorijas, uz nacionālisma pacēluma viļņa uz brīdi nostiprināja savu iekšpolitisko stāvokli, taču nonāca lielās ārējās briesmās.

Padomju vadība novērtēja Minhenes vienošanos kā Rietumu valstu piekrišanu Vācijas agresijas novirzīšanai uz austrumiem, gala rezultātā~--- pret PSRS. Savu pilnīgo bezspēcību parādīja Tautu Savienība. (Pēc Tautu Savienības Statūtu 16.~panta, ja kāda no Savienības loceklēm uzsāka karu, pārējām bija pienākums saraut ar agresoru tirdzniecības un finansiālās attiecības, kā arī sagatavot karaspēku iekļaušanai Tautu Savienības bruņotajos spēkos, kuru uzdevums bija piespiest agresoru cienīt tās lēmumus.) Savienību kā tādu šai gadījumā faktiski aizstāja nosaukto četru Rietumvalstu vadītāju sanāksme. Parasti šādas izmaiņas pasaulē notika pēc ievērojamiem kontinentāla mēroga kariem. Tagad nacistiskā Vācija guva panākumus bez viena šāviena, iegūstot zemes, kuras nekad nebija piederējušas Vācijai.

Minhenē Ā.~Hitlers piekrita parakstīt arī divpusēju Lielbritānijas un Vācijas deklarāciju, kas pēc būtības pildīja līguma par neuzbrukšanu funkcijas. Lielbritānijas premjerministrs N.~Čemberlens, atgriezies pēc Minhenes tikšanās Londonā, paziņoja publikai, ka viņš atvedis ``mieru mūsu paaudzei'' (\entxti{peace in our time}). Lielbritānijas parlaments ar 366~balsīm pret 144 un Francijas parlaments ar 535 pret 75~balsīm atbalstīja Minhenes līguma ratificēšanu.

Tiesa, jau toreiz pastāvēja arī cits vērtējums. V.~Čērčils pēc Minhenes konferences norādīja, ka Anglijai bija jāizvēlas starp karu un negodu, tā izvēlējās negodu, bet iegūs arī karu (``\entxti{England has been offered a choice between war and shame. She has chosen shame, and will get war.}''). Viņš konferences rezultātus vērtēja kā ``pilnīgu un galīgu sakāvi'', jo līdzsvars Eiropā bija izjaukts.

Aprakstot Čehoslovākijas sadalīšanu 1938.~gada rudenī, V.~Čērčils deva arī pirmskara Polijas vadības izvērstu raksturojumu. ``Polijas tautas varonīgās rakstura īpašības,~--- rakstīja ievērojamais britu politiķis,~--- ``nedrīkst mūs spiest pievērt acis tās neprātīgumam un nepateicībai, kas gadsimtu gaitā nesuši tai neizmērojamas ciešanas \citespace{} Ir Eiropas vēstures noslēpums un traģēdija, ka tauta, kura ir spējīga uz jebkuru varonību, kuras atsevišķi pārstāvji ir talantīgi, drosmīgi, valdzinoši, pastāvīgi demonstrē tādus trūkumus gandrīz visās savas valsts dzīves jomās. Slava sacelšanos un posta periodos, neģēlība un negods triumfa periodos. Drosmīgākos no drosmīgākajiem pārāk bieži ir vadījuši neģēlīgākie no neģēlīgākajiem! Un tomēr vienmēr ir pastāvējušas divas Polijas: viena no tām cīnījās par taisnību, otra verdziski izdabāja nelietībai!'' Var nepiekrist dusmu brīdī izteiktam spriedumam, bet ieklausīties izcilā politiķa vārdos noteikti ir vērts.

Ja agrāk Rietumu autoru darbos tika uzsvērts, ka agresora ``nomierināšanas politikai'' (angļu \entxti{Appeasement-Politik}, vācu \detxti{Beschwichtigungspolitik}) nebija alternatīvas, tad jau no XX gadsimta 80., 90.~gadiem vairāk tiek runāts par tās kļūdainību.

Var stādīt jautājumu~--- vai 1938.~gada rudenī situācija bija tāda, ka Lielbritānijai un Francijai bija jāizvēlas: vai nu iestāties par Čehoslovākiju, faktiski nespējot tai palīdzēt, rezultātā pieredzot tās iznīcināšanu un vēl iekļūstot nesagatavotām karā ar Vāciju, vai arī zaudēt Čehoslovākiju, bet pavisam vai uz laiku novēršot savu iesaistīšanos karā? Jāpiekrīt krievu vēsturniekam A.~Djukovam, kurš uzskata, ka Lielbritānija un Francija 1938.~gadā varēja izglābt Čehoslovākiju. Francijai bija tikai pašai jāizpilda savs Čehoslovākijas sabiedrotā pienākums, un jāpieņem PSRS palīdzība, jo kā Čehoslovākijai, tā Francijai bija noslēgti ar PSRS savstarpējās palīdzības līgumi. Acīmredzot pietiktu ar stingras apņēmības demonstrēšanu, lai apturētu Vācijas un tās sabiedrotās pret Čehoslovākiju Polijas agresiju. Ja militārs konflikts arī sāktos (kas teorētiski nebija izslēgts) Vācija un Polija būtu krietni vājākas par apvienotajiem Čehoslovākijas, Francijas, PSRS spēkiem, kuriem droši vien vismaz ar Vācijas blokādi pievienotos arī Francijas sabiedrotā Lielbritānija. Lielbritānijai un Francijai risks šādos apstākļos būtu minimāls. Arī citi mūsdienās jautājumu pētījušie vēsturnieki atzīst, ka 1938.~gadā A.~Hitlera pretinieki: Čehoslovākija, Lielbritānija, Francija, PSRS (un, iespējams, citas valstis, kuras būtu gatavas uzstāties pret Vāciju) visas kopā gan ekonomiski, gan militāri bija daudz stiprākas par Vāciju, un, ja karš būtu sācies 1938.~gadā, tas prasītu daudz mazāk upuru un beigtos daudz ātrāk, nekā tas bija realitātē Otrā pasaules kara laikā.

Taču notika citādi~--- Rietumvalstis kapitulēja, par ko vēlāk dārgi samaksāja. Vismaz daļēji arī šīs Čehoslovākijas nodevības rezultātā Francija cieta smagu sakāvi 1940.~gadā, bet Lielbritānija gadu bija spiesta vienatnē karot pret Vāciju.

Tiesa, ASV bijušais valsts sekretārs H.~Kisindžers pauž viedokli, ja sāktos karš Čehoslovākijas dēļ, diezin vai demokrātiskās valstis varētu uzvarēt, jo karotu pret tautu pašnoteikšanās tiesībām [domātas Sudetu vāciešu pašnoteikšanās tiesības~--- V.Š.] un demokrātiskā sabiedrība neesot bijusi sagatavota sakāvēm, kuras neizbēgami būtu kara sākumā. Uz to gan jāsaka, ka nāciju pašnoteikšanās tiesību jautājums, pastāvot Versaļas sistēmai, bija tik lielā mērā ``samudžināts'', ka izšķirties, kad notiek uzstāšanās par kādas nācijas taisnīgo tiesību ievērošanu, kad~--- pret to, bija tiešām grūti, bet tomēr Rietumu sabiedrībā, kurā vēl no Pirmā pasaules kara laikiem bija dzīvas bailes no ``huņņu''~--- vāciešu barbarisma, antisemīta Ā.~Hitlera iztēlošanai par Sudetu vāciešu nacionālo pašnoteikšanās tiesību aizstāvi nebija lielu perspektīvu. Ja bija jāizvēlas, kam simpatizēt~--- mazajai miermīlīgajai Čehoslovākijai, kurai draudēja Vācijas agresija, vai Sudetu vāciešiem, kuri objektīvi uzstājās kā nemiera cēlāji Eiropā un Ā.~Hitlera plānu veicinātāji, eiropiešu vairākuma līdzjūtība bija čehoslovāku pusē. Cita lieta, ka karot par kādas citas nācijas interesēm Rietumu ``demokrātiskās sabiedrības'' nevēlējās. Protams, piemēram, Francijas ``demokrātiskā sabiedrība'' nebija gatava sakāvēm, taču tāda tā nebija arī ne 1914., ne 1940.~gadā. H.~Kisindžera retorika šeit vērsta tikai uz to, lai kādi attaisnotu Minhenes līguma slēdzējus.

Autors piekrīt tiem vēsturniekiem, kas uzsver, ka ar Minhenes vienošanos kolektīvās drošības ideja, kuru starpkaru periodā centās realizēt daļa Rietumu politiķu, tika atmesta un pēc tās starptautiskajās attiecībās sāka valdīt rupja spēka laikmets.

30.~septembrī Polijas valdība publiskoja savu ultimātu Čehoslovākijai, kuru tā nākamajā dienā pieņēma. Pirms atbildes saņemšanas 1.~oktobrī Polijas ārlietu ministrs J.~Beks sarunā ar Vācijas sūtni Varšavā H.~Moltki interesējās, vai Polija var cerēt uz Vācijas labvēlīgu pozīciju Polijas bruņota konflikta gadījumā ar Čehoslovākiju. Neraugoties uz atbildes izvairīgo raksturu, notikumi risinājās tālāk. 1939.~gada 1.oktobrī nacisti iesoļoja Sudetu apgabalā, 2.~oktobrī poļi~--- Cešinas apgabalā. Var atzīmēt, ka pirms tam~--- 1.~oktobrī Čehoslovākijas vēstnieks Padomju Savienībā Z.~Firlingers apmeklēja PSRS Ārlietu tautas komisariātu un uzdeva jautājumu, ko gatavojas darīt PSRS sakarā ar Polijas ultimātu Čehoslovākijai. Ārlietu tautas komisāra vietnieks V.~Potjomkins izvairīgi atbildēja, ka ir jānogaida, ko darīs lielvalstis, kuras ir parakstījušas Minhenes vienošanos. Līdz ar to tika it kā pilnībā ``aizmirsts'' brīdinājums, kuru PSRS bija izteikusi Polijai pirms 8~dienām. PSRS nevēlējās būt ierauta konfliktā vienatnē. PSRS attiecības ar Poliju pēc tās ekspansijas pret Čehoslovākiju gan palika rezervētas, taču, ja ievēro jau pirms tam esošo to slikto stāvokli, būtiski nesaasinājās. Pie tam pēc drīzumā pašai Polijai uzstādītajām Vācijas prasībām Polijas diplomātija centās uzsvērt savu vēlmi uzlabot attiecības ar PSRS.

Poļu karspēkam ieņemot Cešinu, norisinājās sadursmes, kurās krita viens poļu virsnieks un trīs karavīri. Polija ieguva strīdīgo teritoriju, kurā dzīvoja 80~000~poļu un 120~000~čehu (Literatūrā sastopami arī citi skaitļi. Piemēram, tur mita 156~tūkstoši čehu un 77~tūkstoši poļu.) Pēc tam Polijas valdība no Čehoslovākijas pieprasīja vēl jaunas teritoriālas piekāpšanas, nu jau Slovākijas teritorijā. Saskaņā ar abu valdību vienošanos novembrī Polija ieguva nelielu teritoriju (226~km$^{2}$) Slovākijas ziemeļos ar 4~280~iedzīvotājiem. Ir saprotams, ka slovākos pieauga pretpoļu noskaņojums.

Pasaules sabiedriskā domā Poliju pamatoti uzskatīja par vienu no Čehoslovākijas sadales līdzdalībniecēm. Tā, Latvijas sūtnis PSRS Fricis Kociņš 1938.~gada 14.~oktobrī uz Rīgu ziņoja, ka diplomātiskā korpusa aprindās Maskavā par Polijas rīcību neizskanēja neviens pozitīvs atzinums pat no tās draugiem. Taču pašā Polijā šī svešās valsts teritorijas sagrābšana tika vērtēta kā ievērojams politisks panākums. Lielākā poļu daļa gan pašā Polijā, gan ārzemēs atbalstīja Cešinas apgabala atraušanu Čehoslovākijai un tā aneksiju Polijas sastāvā. Līdzīga poļu nostāja1938.~gada nogalē bija vērojama arī Latvijā.

Kā liecinājis jau minētais K.~Siljakovs: ``Izteikta poļu lojalitāte pret latviešiem turpinājās līdz 1938.~gada 1.~oktobrim, tas ir, dienai, kad Hitlera armija iesoļoja Čehoslovākijai piederošā Sudētijā un Polija tanī pašā laikā sev pievienoja Čehoslovākijas Cešinas apgabalu.

No šī brīža Liepājas poļu kolonijas piederīgie vairs nejutās kā mūsu zemes pavalstnieki, bet lielvalsts Polijas piederīgie, un tā vietā un nevietā uzsvēra savu polisko izcelsmi, ko sevišķi ievēroja tie latvieši, kas dienu dienā bija ar tiem kopā,~--- poļu savstarpējā valoda latviešu vidē vairs nebija latviešu, bet poļu~--- viņu vīri sāka demonstratīvi lepoties ar Polijas ērgli svārku pogu caurumā un sieviešu kārtas piederīgie šo Polijas valsts simbolu sāka nēsāt pēc savas gaumes un visdažādākās variācijās,~--- gan iekārtu kakla ķēdītē blakus kristīgās ticības simbolam~--- krustam, gan izstrādātus kā saspraužamās adatas vai kā piekariņu pie rokas sprādzēm.

Sākot ar 1939.~gada sākumu, kad Vācijas prasības Dancigas koridora dēļ pret Poliju pieauga, tam proporcionāli līdz auga arī Liepājas poļu šovinisms, kas vācu iebrukumu laikā Polijā sasniedza tādu vārīšanās punktu, ka par poļu gaitām un Liepājas vācu tautības iedzīvotājiem un viņu īpašumu drošību nepārtraukti nomodā bija jābūt Liepājas policijai.''

Vairums poļu uzskatīja šo ``sanācijas'' režīma sasniegumu kā pierādījumu Polijas ieiešanai lielvalstu statusā. Ārlietu ministram J.~Bekam tika piešķirts augstākais valsts apbalvojums~--- ``Baltā ērgļa'' ordenis (\pltxti{Order ``Orła Białego''}), viņš kļuva par divu Polijas universitāšu goda (\latxti{honoris causa}) doktoru. Jāņem vērā, ka Polija šoreiz tik tiešām rīkojās patstāvīgi, neraugoties uz Rietumu lielvalstu Lielbritānijas un Francijas noraidošo nostāju. Kā norāda J.~Muhins~--- ja Vācija vēl varēja atsaukties uz Lielbritānijas un Francijas piekrišanu Sudetu okupācijai, tad Polijai Cešinas okupācijai tādas piekrišanas nebija. Polija ar to nostājās uz tiešas agresijas ceļa. Atbalstāms vēsturnieka Ē.~Jēkabsona atzinums, ka J.~Beka arguments par to, ka Čehoslovākijas valdība neilgi pirms Minhenes vienošanās bija principā piekritusi traktēt poļu minoritāti savā valstī līdzīgi kā vāciešus, un tāpēc Polijai esot bijušas arī oficiālas tiesības pieprasīt to pašu, ko saņēmusi Vācija, ``jāatzīst par formālu, un apgabala pārņemšanas laiks un apstākļi vēl šodien nostāda Poliju \citespace{} zināmā tālaika Vācijas sabiedrotā lomā\dots{}''.

Interesantu aspektu Polijas vadītāju politikā atklāj tās sūtnis Berlīnē J.~Lipskis ``pilnīgi slepenā'' ziņojumā J.~Bekam Nr. 1/182/38, kuru viņš uzrakstīja pēc tikšanās ar J.~f.~Ribentropu 1938.~gada 27.~septembrī. Tajā ir frāze: ``Gadījumā, ja Vācija ieņemtu visu Čehoslovākiju, Ribentrops uzskata par lietderīgu\dots{}''. Tātad J.~f.~Ribentrops jau 27.~septembrī runāja par iespējamo ``visas Čehoslovākijas ieņemšanu'' un gandrīz pusgadu pirms faktiskajiem notikumiem ``sanācijas'' režīma vadītāji zināja, kam gatavojās Ā.~Hitlers, un tomēr nevis izdarīja attiecīgus secinājumus par Polijas drošību, bet turpināja sadarboties ar Vācijas fīreru, kurš jau kopš Minhenes vienošanās brīža bija gatavs to lauzt. Te nu acīmredzot var izmantot franču teicienu ``\frtxti{c'est plus qu'un crime, c'est une faute}'' (Tas ir vairāk nekā noziegums, tā ir kļūda), kuru pieraksta Napoleona I ministriem Ž.~Fušē vai Š.~Taleirānam, bet patiesībā tas pieder citam šī laika darbiniekam B.~de~la~Mertam.

Pēc Otrā pasaules kara komunistiskie poļu autori principiāli nosodīja Polijas agresiju pret Čehoslovākiju. Tā, ebreju izcelsmes poļu žurnālists un diplomāts V.~Grošs rakstīja: ``Lomai, kura bija Ridzs-Smiglija un Beka sanācijas režīmam Čehoslovākijas sadalīšanā, nav cita piemēra mūsu tautas vēsturē. Mūsu zeme, kura pusotra gadsimta cieta nebrīvi, kura XIX gadsimtā iedvesmoja dzejniekus un bija cīnītāju par brīvību karogs, zeme, kuras liktenis vienmēr bija saistīts ar cīņu ``par jūsu un mūsu brīvību'',~--- šī zeme pēc savu valdītāju gribas tika pārvērsta par šakāli, kurš rāva asiņainas atliekas no lielāku plēsoņu galda. Nebija pagājuši ne divdesmit gadi kopš atbrīvošanās no iekarotāju jūga (vesela paaudze vēl atcerējās savus jaunības sapņus par neatkarību), kad mums nācās pārdzīvot šo drausmīgo apkaunojumu!'' Poļu komunisti šādu Polijas valdības nostāju saistīja ar ``sanācijas'' politiķu šķirisko naidu pret padomju kaimiņvalsti savienojumā ar bērnišķīgiem sapņiem par ``lielvalstisku'' Poliju ``no jūras līdz jūrai'' un vienkāršu politisku naivumu. Jāatzīst, ka skaidrojums bija vienkāršots, sarežģītajā starptautiskajā spēlē pirms Otrā pasaules kara nebija viegli atrast valstiski izdevīgāko ārpolitisko kursu. Pēc autora domām, tieši izdevīguma meklējumi, kuri pārauga agresivitātē, neļāva Polijas līderiem izvēlēties valstij visdrošāko ārpolitisko kursu~--- pievienošanos kolektīvās drošības politikas aizstāvjiem, Tautu Savienības pārvēršanu par miermīlīgo spēku sadarbības iestādi, nepiedalīšanos vardarbībā pret citām tautām, bet to aizstāvēšanu, jo Polijas brīvība bija tieši saistīta ar citu tautu brīvību un drošību. A.~Miņins atzīmējis, ka 30.~gados kolektīvās drošības ideja Tautu Savienības formā izgāzās, jo vairākas valstis izaicināja Tautu Savienības principus un palika nesodītas, attiecīgi Japāna 1931.~gadā, Itālija 1935.~gadā, Vācija 1938.~gadā un PSRS 1939.~gadā (veicot militāru iebrukumu Somijā). Diemžēl kolēģis nav ievērojis, ka ``Tautu Savienības principus'' 1938.~gadā kopā ar nacistisko Vāciju ``izaicināja'' arī ``sanācijas'' Polija. Sīks laupītājs arī ir laupītājs.

Mūsdienās Polijā nelabprāt atceras Minhenes vienošanos un Polijas dalību pret Čehoslovākiju vērstajā agresijā. Daudzi poļi gribētu aizmirst vēstures lappusi, kad Polija pret Čehoslovākiju faktiski uzstājās kā nacistiskās Vācijas sabiedrotā. Arī poļu vēsturnieki parasti nelabprāt piemin Polijas dalību Čehoslovākijas sadalē. Mūsdienu Polijā izdotā mācību grāmatā vidusskolām šis fakts noklusēts. Taču nevar izsvītrot notikušo kā nebijušu. Bez tam mūsdienās poļu vēsturnieku vidū Polijas politikas attiecībās ar Čehoslovākiju vērtējumi nav stabili.

Vēl nesen gandrīz visi poļu pētnieki Varšavas 30.~septembra ultimātu Prāgai un pēc tam Polijas armijas ieiešanu Cešinas apgabalā uzskatīja par lielu kļūdu. Ievērības vērta bija 2009.~gada 1.~septembrī Polijas prezidenta Ļ.~Kačinska publiskā atvainošanās par Polijas dalību Čehoslovākijas sadalē 1938.~gada oktobrī, atzīstot par ``kļūdu un grēku'', kā arī aicinot šim grēkam ``nemeklēt attaisnojumus, pat ja pastāv tāda iespēja''. (Tiesa, šis paziņojums acīmredzot tika veikts, lai sagaidītu no Polijā vizītē esošā Krievijas premjera V.~Putina līdzīgu paziņojumu par padomju veikto agresiju 1939.~gada septembrī, bet tas tā arī nesekoja.)

Taču arī šodien Polijā var dzirdēt balsis, kas cenšas ja ne attaisnot šo toreizējo Polijas valdītāju soli, tad vismaz daļēji izskaidrot to ar sarežģītajām poļu-čehoslovāku attiecībām. Daļa poļu vēsturnieku, tiesa~--- to mazākums, ir tendēti šīs Polijas ārpolitikas kā fatālas kļūdas vērtējumu mīkstināt, izskaidrot ar dažādiem apstākļiem. Tomēr vismaz daļa poļu vēsturnieku mūsdienās atzīst, ka Cešinas aneksija, ko Polijas valdība toreiz atspoguļoja kā spīdošu panākumu, patiesībā ievērojami pasliktināja Polijas situāciju, kad drīzumā Vācija izvirzīja prasības pret pašu Poliju. Piemēram, A.~Friške atzīst, ka Cešinas atgriešanas [Polijai] prasība, ``kura tika izvirzīta laikā, kad šī valsts [Čehoslovākijas Republika] tika Reiha visvairāk apdraudēta, bija katastrofāla kļūda. Tā atviegloja Rietumu lielvalstīm pieņemt Minhenes lēmumus un ietekmēja Čehoslovākijai pretošanās gribas vājināšanos''.

Gandrīz visi eiropieši mūsdienās zina, kas 1938.~gadā notika Sudetos, taču daudz mazāks skaits cilvēku apzinās, ka tas pats notika Cešinas apgabalā, tikai vācu karavīru vietā Cešinā, Karlovi Varos (Čehu \cstxti{Karlovy Var}, poļu \pltxti{Karlowe Wary}, vācu \detxti{Karlsbad}) un citās apgabala pilsētiņās iemaršēja poļu karavīri. Vācu publicists J.~Volmers konstatējis, ka Cešinas apgabalā Polijas valdība realizēja tādu pat ``tautiešu politiku'' (\detxti{Volksgruppenpolitik}) pret poļu mazākumu, kā Otrā pasaules kara gados nacionālsociālistu politiķi pret etniskajiem vāciešiem Polijā, radot tiem priviliģētu stāvokli salīdzinājumā ar pārējiem iedzīvotājiem. Pirmā poļu reforma Cešinas apgabalā bija čehu skolu aizliegšana. Notikumu aculiecinieks čehu ģenerālis A.~Vehireks vēlāk atcerējās: ``Poļi nesaudzīgi vajāja čehus, terorizēja, atlaižot no darba, izsvieda no mājām, konfiscēja mantu. Viss, kas bija čehisks, tika iznīcināts. Čehu valoda un pat sveicieni tika aizliegti. Par sveicienu ``\detxti{Nazdar}'' sodīja ar 4 zlotiem. Tad nu čehi sveicināja viens otru ar saucienu ``Četri zloti''. Čehu nosaukumus likvidēja pat uz kapiem\dots{}''. Rezultātā no Cešinas apgabala bēga 30~tūkstošu čehu un 5~tūkstoši vāciešu. Līdzīgi notikumi risinājās arī Tatru kalnos esošajos četros slovāku ciemos, kuru iedzīvotāji savā laikā bija nobalsojuši par pievienošanos Čehoslovākijai. Tiesa, arī jaunizveidoto poļu iestāžu darbiniekiem, policistiem, robežsargiem Cešinā nācās izjust vietējo iedzīvotāju naidu tik lielā mērā, ka Latvijas sūtnis Polijā Ludvigs Ēķis rakstīja par to ``terorizēšanu'', par ko Polijas valdība sniedzot protestus Čehoslovākijas valdībai Prāgā. Arī ASV militārais atašejs Polijā V.~Kolberns ziņojumā Vašingtonai vēstīja, ka ``čehu atbalstītāju darbība ir pieņēmusi organizēta terorisma formu un tās rezultātā gandrīz katru dienu Cešinas Silēzijā notiek incidenti.'' Tiesa, noteikti jāuzsver, ka pēc tam, kad Polija bija okupējusi Cešinu, neviens cilvēks netika nogalināts. Poļu ieviesto okupācijas režīmu nevar salīdzināt ar nacistu daudz briesmīgāko režīmu to sagrābtajās zemēs, kad bojā gāja miljoniem cilvēku.

Jāpiezīmē, ka sakarā ar Cešinas okupāciju saasinājās arī vācu un poļu nacionālistu attiecības Polijā. Līdz 1938.~gada rudenim Čehoslovākijā dzīvojošais vācu mazākums atbalstīja poļu mazākuma kultūru Cešinas apgabalā. Piemēram, Vācu teātra ēka Cešinā bieži tika piešķirta poļu izrāžu uzvešanai. Taču poļu ``pateicība'' bija visai savdabīga. Pēc tam, kad Polija okupēja Cešinu, sākās karagājiens arī pret visu vācisko šai apgabalā: skolām, kultūras biedrībām, grāmatveikaliem. 1939.~gada jūnijā poļu nacionālisti ielauzās jau slēgtajā Vācu teātrī un sadragāja piecas marmora piemiņas plāksnes. Uz vienas no tām bija iegravēti visu ziedotāju teātrim vārdi, uz citām~--- līdzšinējo vietējās teātra biedrības vadītāju vārdi.

Galvenais Polijas ieguvums bija sagrābtā apgabala rūpnieciskais potenciāls. 1938.~gada beigās šai teritorijā esošie uzņēmumi ražoja gandrīz 41\% no visa Polijā saražotā čuguna un gandrīz 47\% tērauda. Vērtējot notikušo, V.~Čērčils savos memuāros rakstīja, ka Polija ``ar hiēnas alkatību ņēma dalību čehoslovāku valsts aplaupīšanā un iznīcināšanā''.

Tomēr poļu radikālnacionālisti ar jau esošo ``panākumu'' vēl nebija apmierināti. Vācija taču bija saņēmusi vairāk teritoriju! Skanēja balsis, ka Polijai pat vajadzētu nodibināt savu protektorātu pār Čehiju un Slovākiju. Tika izdarīts spiediens uz valdību, lai tā panāktu ``vismaz'' kopīgu robežu ar Ungāriju, t.i.~--- no Čehoslovākijas Polijai par labu tiktu atsavināta daļa Aizkarpatu Ukrainas. Tas pēc Polijas imperiālistu viedokļa palielinātu Polijas starptautisko svaru. Kā parādījis poļu vēsturnieks H.~Batovskis, Polijas valdība tūlīt pēc Minhenes uzskatīja, ka pienācis laiks uzstāties kā pretendentei uz ``Čehoslovākijas mantojumu''. Tāpēc tā izvērsa diplomātisko aktivitāti, lai uz Čehoslovākijas teritorijas rēķina panāktu Polijai kopējas robežas ar Ungāriju izveidi, Polijas robežas ar Rumāniju paplašināšanu, reizē noregulējot arī Rumānijas un Ungārijas attiecības. Taču uz lielāko daļu no Aizkarpatu Ukrainas (čehu \cstxti{Země Zakarpatskoukrajinská}) un daļu Slovākijas pretendēja Ungārija, kura arī pieprasīja no Čehoslovākijas šīs teritorijas. Tā 2.~oktobrī izteica gatavību sākt sarunas, bet reizē griezās pie Dienvidslāvijas un Rumānijas ar lūgumu izdarīt spiedienu uz Ungāriju.

Tomēr bailes no ukraiņu irredentisma (mūsdienās~--- valsts, partijas vai politiskās kustības politika, kas vērsta uz nācijas apvienošanu vienā valstī) spieda Polijas vadītājus izteikties par Piekarpatu teritorijas pievienošanu Ungārijai. 1938.~gada oktobrī J.~Beks centās panākt, lai Vācija Ungārijai dotu zaļo gaismu Aizkapatu okupācijai un Polijas-Ungārijas robežas radīšanai. Ungārija lūdza Polijai palīdzību ar karaspēku. Izvērsās poļu un ungāru teroristu provokācijas pret Čehoslovākiju tai piederošajā Aizkarpatu Ukrainā, kura 11.~oktobrī bija saņēmusi autonomiju. 19.~oktobrī Polijas Galvenā štāba priekšnieks ģenerālis V.~Stahevičs deva pavēli uzsākt poļu diversantu darbību Čehoslovākijas teritorijā. 56~diversiju gaitā poļu kaujinieki uzspridzināja 13~tiltus, 27~telefona līnijas, nogalināja 23~cilvēkus, 15~ievainoja, aizveda 20~gūstekņus.

Dienvidslāvija un Rumānija, kuras jau sen uzstājās pret pārmērīgo Ungārijas apetīti, paziņoja Vācijai un Itālijai par savu vēlmi ierobežot Ungārijas pretenzijas ar teritorijām, kur bija pārsvarā ungāru iedzīvotāji. Šos argumentus labprāt uztvēra Vācijas valdība, kas uzskatīja, ka kopēja Ungārijas un Polijas robeža būs pretrunā ar Vācijas interesēm. Pie tam nacistu vadībā jau sāka nobriest plāni par pilnīgu Čehoslovākijas likvidāciju un tālāku Vācijas ietekmes izplatību Austrumeiropā. Tāpēc poļu valdītāju jaunajiem agresīvajiem nodomiem nebija lemts piepildīties.

1938.~gada 2.~novembrī Vācijas un Itālijas ārlietu ministri J.~f.~Ribentrops un G.~Čiano, uzstājoties kā arbitri Ungārijas un Čehoslovakijas strīdā par Aizkarpatu Ukrainas teritoriju tā saucamās Vīnes arbitrāžas (\detxti{Wiener Schiedsspruch}) gaitā, noteica, ka Čehoslovākijai jānodod Slovākijas dienvidu rajoni un Aizkarpatu Ukrainas dienvid-rietumu rajoni ar Užgorodas (ukraiņu \uktxti{Ужгород}, ungāru \hutxti{Ungvár}) un Mukačevas (ukraiņu \uktxti{Мукачеве}, ungāru \hutxti{Munkács}) pilsētām Ungārijai. Pēdējā ieguva ap 12~000~km$^{2}$ teritoriju ar 1~miljonu iedzīvotāju. Tomēr daļā Aizkarpatu Ukrainas ar galvaspilsētu Hustu (čehu, poļu, vācu \detxti{Chust}, ukraiņu \uktxti{Хуст}, ungāru \hutxti{Huszt}) saglabājās autonomija. Vācijas un Itālijas pārstāvji arī ieteica radīt ``Aizkarpatu Ukrainas'' valsti. Apgabalā izveidojās autoritārs režīms. Tika radīta vietējā armija~--- Karpatu Seča (\uktxti{Карпатська Січ}) ar 2~000 karavīriem.

Pirmā Vīnes arbitrāža pēc Minhenes vienošanās kļuva par vēl vienu pakāpienu Čehoslovākijas valsts iznīcināšanā. Diversantu darbība tika pārtraukta un arī poļu vienības 722 kaujinieku sastāvā atsauktas. Polija pagaidām nesaņēma jaunus ieguvumus. Šie notikumi tikai paātrināja visa poļu-vācu attiecību kompleksa caurskatīšanu. Notikumi ļāva noprast, ka Vācija šai Aizkarpatu Ukrainai gatavojas pievienot citus ukraiņu apdzīvotus apgabalus, tai skaitā arī Polijai piederošus, un veidot no Vācijas atkarīgu marionešu valsti. Tas dotu Vācijai spiediena līdzekli gan pret PSRS, gan Poliju. Nenoteiktības stāvoklis turpinājās vairākus mēnešus, līdz galīgai Čehoslovākijas likvidācijai.

Izmantojot kā ieganstu augstāk aprakstītās Polijas aktivitātes, savas prasības tai savukārt izvirzīja Vācija.

Izrādījās, ka kopējā vēršanās pret Čehoslovākiju, Sudetu un Cešinas apgabalu atņemšana tai bija pēdējais Polijas un Vācijas sadarbības piemērs starpkaru periodā. 1938.~gada oktobrī Vācijas un Polijas attiecībās notika pavērsiens. Ā.~Hitlers lika J.~f.~Ribentropam izteikt Polijas valdībai savu vēlmi. Jau 1938.~gada 24.~oktobrī, kad Polijas sūtnis Berlīnē J.~Lipskis Vācijas ārlietu ministram J.~f.~Ribentropam izteica Polijas vēlmi iegūt robežu ar Ungāriju, otrais, neko nesolot, savukārt deklarēja Vācijas vēlēšanos pievienot savai teritorijai Dancigu un izbūvēt dzelzceļu un autoceļu caur ``poļu koridoru'', kas savienotu Prūsiju ar pārējo Vācijas teritoriju. Kamēr Polija nebija izbūvējusi Gdiņas ostu un ap 2/3~sava jūras tirdzniecības apjoma veda caur Dancigu, vietējos vāciešus šāds stāvoklis apmierināja. Kad no 1928.~gada arvien lielāka Polijas tirdzniecības daļa norisa caur Gdiņu, Dancigas ekonomiskais stāvoklis būtiski pasliktinājās. Tautu Savienība piespieda Poliju izdalīt Dancigai kvotu kravu apgrozībā, taču jau tas vien, ka Polija jebkurā brīdī varēja ekonomiski ``nožņaugt'' Dancigu, izvirzīja jautājumu par tās atgriešanu Vācijai. Ja jau Polijai bija sava osta, kāpēc gan bija mākslīgi jāsaglabā Dancigas brīvpilsētas statuss? Argumentu par labu Dancigas piederībai Polijai tās vadītājiem nebija.

Trīs nedēļas pēc Minhenes vienošanās Ā.~Hitlers jau plānoja panākt plašāku vienošanos ar Poliju, tāpēc izsauca uz savu rezidenci Berhtesgādenē Polijas sūtni un īsi raksturoja iespējamo piedāvājumu. Ā.~Hitlers ierosināja Polijai atteikties no savām tiesībām uz Dancigu un atļaut pāri savai teritorijai izbūvēt Berlīnes--Kēnigsbergas automaģistrāli. Ja poļi tam piekristu, viņi varētu pievienoties sev izdevīgai politiskai un ekonomiskai savienībai, kas būtu vērsta pret Padomju Savienību. Pēc angļu vēsturnieka N.~Deivisa pieņēmuma Ā.~Hitlera rasu un ideoloģiskie aizspriedumi ļāva viņam cerēt uz Polijas drīzu piekrišanu, jo ``galu galā, poļu pulkvežiem taču bija jātiek galā ar lielāko ebreju kopienu visā Eiropā, turklāt Polija bija noskaņota krasi antikomunistiski, tāpēc viņš uzskatīja, ka Polijas un nacistiskās Vācijas partnerība ir ļoti dabiska''. Taču, kā uzsvēris N.~Deiviss, Ā.~Hitlers slikti pazina poļu dabu un nezināja, ka poļu nacionālisti pret Vāciju ir tikpat naidīgi noskaņoti kā pret Krieviju. Taču pats galvenais~--- Ā.~Hitlers nevarējis iedomāties, ka no maršala J.~Pilsudska mantiniekiem viņš sagaidīs pilnīgi citādu atbildi nekā no Čemberlena un Beneša. Pulkveži negrasījās locīties bijušā jefreitora priekšā. Viņi bijuši gatavi cīnīties un cīnoties iet bojā.'' Tiktāl N.~Deiviss. Šo vērtējumu nākas komentēt.

Pats nepatīkamais Polijas bija tas, ka tās situācija 1938.~gada beigās bija ievērojami sliktāka nekā pirms 1934.~gada janvāra: ``Starpjūru'' bloku izveidot nebija izdevies, Versaļas sistēma bija praktiski iznīcināta, tai skaitā ar pašas Polijas piedalīšanos. Francija un Lielbritānija, īstenojot Vācijas nomierināšanas politiku, kļuva par vēl mazāk uzticamām Polijas sabiedrotajām, nekā pirms tam, attiecības ar Padomju Savienību tā ir palika aukstas.

Faktiski 1934.~gada 26.~janvāra vācu-poļu līgums, kurš neapstiprināja pastāvošo Vācijas un Polijas robežu neaizskaramību, jau 1938.~gadā pavērsās pret Poliju. J.~Lipskis pēc sarunas ar J.~F.~Ribentropu gan ziņoja J.~Bekam, ka J.~f.~Ribentrops Polijai piedāvājis sadarbību ``koloniālajā jautājumā, ebreju pārvietošanā no Polijas un kopīgos pasākumos pret Krieviju Antikominternes pakta (Noslēgts 1936.~gada 25.~novembrī starp Vāciju un Japānu ar nolūku nepieļaut Komunistiskās Internacionāles ietekmes pieaugumu pasaulē. Vēlāk tam pievienojās arī citas valstis.~--- V.Š.) ietvaros.'' Uzaicinājums iesaistīties Antikominternes paktā bija līdzvērtīgs paziņojumam, ka Ā.~Hitlers vairs nav apmierināts ar Polijas līdzsvara politiku starp Vāciju un PSRS. Taču Polija, kura pati pretendēja uz lielvalsts statusu, nekādā gadījumā nevēlējās kļūt par Vācijas jaunāko partneri. Kad to tās sūtnis Berlīnē J.~Lipskis 1938.~gada 19.~novembrī arī oficiāli paziņoja, neizbēgama kļuva Vācijas un Polijas attiecību saasināšanās. Ā.~Hitlers jau 24.~novembrī deva pavēli sagatavot brīvpilsētas Dancigas ieņemšanas plānu, kuru varētu realizēt politiski izdevīgā brīdī.

Tiesa, vēl 1939.~gada 5.~janvārī Ā.~Hitlers sarunā ar J.~Beku pieļāva iespēju ļaut Polijai piedalīties Čehoslovākijas palieku sadalē, prasot par to Dancigas pievienošanu Vācijai un eksteritoriālu auto- un dzelzceļu caur ``poļu koridoru''. Pilnīgi iespējams, ka līdz 1939.~gada pirmajiem mēnešiem Ā.~Hitlers bija paredzējis gadījumā, ja Polija bez ierunām pakļausies viņa prasībām, ``jaunās kārtības'' iedibināšanas gaitā piešķirt tai tādu pat statusu, kāds vēlāk tika piešķirts Slovākijai, Ungārijai un Rumānijai. Tāpat iespējams, ka prasība Polijai atteikties no Dancigas, atļaut Vācijai izbūvēt eksteritoriālu autostrādi un dzelzceļu pa ``poļu koridoru'' bija līdzeklis kā izzondēt Polijas pozīciju~--- vai tā pakļausies Vācijas hegemonijai, īpaši pēc tikko īstenotā Vācijas un Polijas sadarbības piemēra Čehoslovākijas krīzes sakarā, par kuru Polija tika atalgota ar Cešinas apgabalu. Kā uzskatīja padomju vēsturniece T.~Grigorjanca, acīmredzot līdz 1938.~gadam Ā.~Hitleram bija nodoms vienoties ar Polijas valdošajām aprindām un izmantot Poliju kā savas ārpolitikas instrumentu, pārvēršot to par Vācijas vasali.

Taču 1939.~gada 8.~janvārī Karaļa pilī Varšavā Polijas vadības sanāksmē Ā.~Hitlera piedāvājums tika noraidīts. Ar to Polija kļuva par pirmo valsti, kura noraidīja Vācijas fīrera diktātu. Jāpiekrīt jau pieminētajam krievu vēsturniekam M.~Narinskim, ka tas dara godu toreizējai Polijas vadībai. 1939.~gadā Polijas valdība noraidīja Vācijas priekšlikumus ne tikai tāpēc, ka tā neticēja Ā.~Hitlera solījumiem un saprata, ka teritoriāla piekāpšanās tikai izraisītu jaunu prasību uzstādīšanu, kā tas jau bija Čehoslovākijas gadījumā. Arī Dancigas un satiksmes ceļu ar to kontroles jautājums Polijai nebija būtiski. Brīvpilsētu joprojām kontrolēja Tautu savienība, un poļiem bija sava osta Gdiņā. Varšavai atteikšanās no tiesībām uz Dancigu galvenokārt bija prestiža jautājums. Izšķiroši bija citi apsvērumi.

Ā.~Hitlera priekšlikuma pieņemšana nozīmētu Polijas atteikšanos no saistībām PSRS priekšā, atteikšanos no līdzsvara politikas, pakļaušanos Vācijas politikai. Vācijas priekšlikumi bija arī kopējās tās un Polijas pret Čehoslovākiju vērstās rīcības tiešs rezultāts. Ja jau Polija bija padevusi Vācijai ``mazo pirkstiņu'' un kopā ar to aplaupīja Čehoslovākiju, Vācija vēlējās paņemt ``visu roku'' un panākt globālu Polijas pakļaušanos. Iekļaušanās Antikominternes valstu blokā nozīmētu Polijas suverenitātes ievērojamu ierobežošanu un tās pārvēršanos par Vācijas satelītu.

Pārāk cieša tuvināšanās Vācijai un kopīga vēršanās pret Padomju Savienību būtu Poliju starptautiski izolējusi, nozīmētu atteikšanos no līdz tam Polijas realizētās ``līdzsvara politikas'' un nostādītu Poliju \latxti{de facto} Vācijas vasaļa lomā. Šāda politika Polijas valdībai, kura vadījās gan no reālās situācijas, gan no zināmām lielvalstiskām ilūzijām, bija nepieņemama. Ārlietu ministrs J.~Beks vēstulē kādam vecam draugam atklāti izklāstīja nepiekāpīgās nostājas pamatojumu: ``Vai nu mēs pacelsim mūsu valsti līdz lielvaras augstumiem, vai arī mēs nokļūsim~--- tā vai citādi, \latxti{nolens volens} (latīņu~--- gribot negribot)~---  Slovākijas [šai laikā Vācijas marionešvalsts~--- V.Š.] vai līdzīgā stāvoklī.''

Tiesa gan, pēc krievu vēsturnieka V.~Pihalova domām, tam, ka Polija atteicās oficiāli iestāties Antikominternes paktā, nebija būtiskas nozīmes, jo jau līdz tam tā bija atbalstījusi agresīvās valstis, gan Itālijas realizēto Etiopijas sagrābšanu, gan Vācijas un Itālijas palīdzību Franko vadītajiem dumpiniekiem Spānijā, gan Japānas iebrukumu Ķīnā, gan Austrijas pievienošanu Vācijai un Čehoslovākijas sadali. Jāsaka, ka šim secinājumam nevar piekrist, jo, izņemot Čehoslovākijas sadales gadījumu, Polijas nostāja nebūt nebija tik viennozīmīga. Arī Polijas atteikšanos no piekāpšanās Vācijai Dancigas un t.s. ``poļu koridora'' jautājumā, nenākas uzskatīt par pēkšņu Vācijai draudzīgās Polijas ārpolitiskā kursa maiņu. Polija tikai turpināja jau J.~Pilsudska laikā aizsākto ``līdzsvara'' politiku.

Zīmīgi, ka šādos apstākļos dažos Polijas slāņos valdīja eiforija par nākotnes perspektīvām. Polija armijas Galvenā štāba Otrā (Izlūkošanas) nodaļa 1938.~decembrī prezentēja Polijas ārpolitikas Austrumos izstrādni. Atsaucoties uz J.~Pilsudski, tajā bija deklarēts: ``Krievijas sadalīšana atrodas Polijas valstisko interešu austrumos pamatā.'' Tālāk uzdevums tiks konkretizēts: ``\dots{}mūsu iespējamā pozīcija tiks ietverta sekojošajā formulā: lai kas arī nepiedalītos [Krievijas] sadalē, Polija šajā svarīgajā vēsturiskajā brīdī nevar palikt pasīva. Viņai ir pietiekami savlaicīgi jāsagatavojas tam fiziskajā un garīgajā ziņā.\citespace{} Galvenais uzdevums~--- Krievijas novājināšana un sagrāve''. Protams, tas nebija oficiālais valdības kurss. Taču, šajā gadījumā droši var teikt: kas politiķiem bija prātā, militāristiem~--- uz papīra un jau darbībā.

Šai laikā Padomju Savienībā masveidā tika iesūtīti poļu izlūkaģenti. PSRS IeTK Robežsardzes karaspēka Kijevas apgabala sešos iecirkņos no 1939.~gada 1.~janvāra līdz 13.~jūnijam tika aizturēti 34~robežas pārkāpēji. Vairākums no viņiem izrādījās poļu izlūki. Tai pat gadā vairāki poļu aģenti tika atmaskoti Kijevā, Taškentā un Novosibirskā. 1939.~gada rudenī, kad IeTK rokās Polijas teritorijā nonāca vietējā izlūkdienesta dokumenti, tur tika atrasti dati par 186~poļu izlūkaģentiem, kuri darbojās PSRS teritorijā.

Arī nākamajā tikšanās reizē ar J.~f.~Ribentropu 26.~janvārī J.~Beks vēlreiz apstiprināja Polijas pretenzijas uz Padomju Ukrainu un izeju uz Melno jūru. Jādomā, ka Polijas ārlietu ministrs atspoguļoja visas Polijas vadības nostāju. Taču uz J.~f.~Ribentropa kārtējo piedāvājumu pievienoties Antikominternes paktam J.~Beks tikai diplomātiski solīja vēlreiz padomāt par to, pasakot arī, ka ``mēs [poļi] ļoti nopietni izturamies pret mūsu neuzbrukšanas līgumu ar Krieviju un uzskatām to kā ilglaicīgu lēmumu''. Par to uzzinājis, Ā.~Hitlers ar nožēlu noteica: ``Gudrais maršals Pilsudskis nomira pārāk agri\dots{}'' Ukraiņu vēsturnieks un publicists S.~Lozuņko uzsvēris, ka Polija vēl 1939.~gadā stādīja mērķi sagraut PSRS, sagrābt daļu tās teritorijas, šo mērķu īstenošanas gaitā netika izslēgta arī savienība ar nacistisko Vāciju, ar kuru Polija bija tik ``sekmīgi'' sadarbojusies Čehoslovākijas teritoriju aneksijā (zinot to, nenākas brīnīties par PSRS attieksmi pret Poliju). Autoram tomēr jāatzīmē, ka blakus šiem pastāvošajiem tālejošajiem mērķiem Polijas valdība konkrētajos apstākļos turpināja realizēt ``līdzsvara politiku'' starp Vāciju un PSRS. Interesants ir vēl viens šīs politikas realizācijas aspekts. Pēc J.~Beka negatīvās atbildes saņemšanas atpakaļceļā uz Berlīni J.~f.~Ribentrops secināja: ``Tagad, ja mēs [vācieši] negribam nokļūt pilnīgā aplenkumā, paliek pāri tikai viena izeja: apvienoties ar Krieviju.''

Arī pēc pirmās Vīnes arbitrāžas Polijā vēl pastāvēja cerības iegūt jaunas teritorijas. Kad 1939.~gada marta sākumā aktuālas kļuva slovāku pretenzijas uz atdalīšanos no Čehoslovākijas, J.~Beks paziņoja, ka Polijas valdība pilnībā jūt līdzi slovākiem. 11.~martā J.~Beks Senāta ārlietu komisijā teica runu, kurā uzsvēra Polijas koloniālās pretenzijas, pēc tam atkārtoja prasību, lai Čehoslovākijas valdība līdz galam apmierina Polijas un Ungārijas teritoriālās prasības. Viņš izteica arī cerību, ka ``Gdaņskas [Dancigas] problēma'' tiks atrisināta apmierinoši Polijai. Ar to tika pārvērtēti savi, kā arī nepietiekami novērtēti Vācijas spēki.

Taču līdz pretenziju realizācijas mēģinājumiem nenonāca, jo 1939.~gada martā Vācijas prasības, par kurām jau 1938.~gada oktobrī runāja J.~f.~Ribentrops, tika atkārtoti izvirzītas ultimatīvā veidā. Poļu vēsturnieks J.~Krasuskis secinājis, ka J.~Beka iluzorie plāni Poliju nostādīt kā ``trešo spēku'' starp Vāciju un Padomju Savienību nevarēja beigties citādi kā ar fiasko.

Runājot par Polijas ambīcijām, tomēr jāatzīst, ka blakus avantūristiskiem ārpolitiskiem plāniem Polijas valdošajās aprindās pastāvēja arī reālistiskāki centieni. Polijas diplomātija spēra soļus attiecību normalizēšanai ar PSRS.

Tā kā PSRS bija piedzīvojusi neveiksmi mēģinājumos apturēt Vācijas un Polijas agresiju pret Čehoslovākiju, tā bija spiesta piekrist Polijas priekšlikumiem. 1938.~gada 27.~novembrī tika publicēts komunikē par to, ka abu valstu attiecību pamatā paliek iepriekš noslēgtie līgumi, bet 1932.~gada neuzbrukšanas līgums tiek pagarināts līdz 1945.~gadam. Tiesa, Vācijas sūtnis Varšavā H.~Moltke 1939.~gada 13.~februārī intervijā vācu žurnālistam noteikti paziņoja: ``1938.~gada novembra deklarācijai nav nozīmes un tādā nozīmē tā tiek saprasta arī Berlīnē. Stāvoklis ir pilnībā skaidrs. Mēs zinām, ka vācu-padomju konflikta gadījumā Polija stāvēs mūsu pusē. Tas ir pilnīgi noteikti.'' Protams, tas bija tikai atsevišķa diplomāta vērtējums, taču, domājams, tas atspoguļoja Vācijas vadītāju tā brīža uzskatus.

1939.~gada 14.~martā, Ā.~Hitlera pamudināti, slovāku nacionālisti pasludināja Slovākijas ``neatkarību'', tās atdalīšanos no Čehijas. Faktiski Slovākija (\cstxti{Slovenská republika}) kļuva par Vācijas marionešvalsti. Rietumu demokrātisko valstu vadītāji sakarā ar Čehoslovākijas ``iekšējo sabrukumu'' atsauca savas garantijas Čehoslovākijas valsts pastāvēšanai. Vēl vairāk~--- N.~Čemberlens pārstāvju palātā paziņoja, ka Čehoslovākija beigusi savu pastāvēšanu ``iekšēja sabrukuma rezultātā'' un pasludināja savas valdības gatavību sekot iepriekšējai ārpolitiskajai līnijai, pat uzsvēra, ka ``nevienam neļaus novirzīt sevi no šī kursa''. (``Šo kursu'' raksturoja tas, ka Anglijas banka palīdzēja Ā.~Hitleram pārdot Čehoslovākijas aplaupīšanas rezultātā iegūto zeltu gandrīz 1~miljarda sterliņu mārciņu vērtībā (mūsdienu cenās), lai par iegūto valūtu pirktu stratēģiskos materiālus.)

Interesanti, ka tai pat dienā~--- 14.~martā padomju diplomāts G.~Astahovs savā dienasgrāmatā izdarīja ierakstu par sarunu ar Latvijas pilnvaroto lietvedi (franču \frtxti{Сhargé d’affaires}) Ēriku Īgenbergu. Latviešu diplomāts, norūpējies par to, vai PSRS ``atļaus vāciešiem sagrābt Rīgu'', atzīmējis, ka pēdējā laikā pasliktinājušās attiecības starp Poliju un Latviju. Viņa rīcībā bija ziņas, ka Varšavā saasinājusies divu politiķu grupu ar dažādu ārpolitisko orientāciju cīņa. Viena iestājoties par cīņas pret pieaugošajiem Vācijas draudiem pastiprināšanu, neapstājoties attiecību ar Vāciju pasliktināšanās perspektīvu priekšā, otra iestājoties par tālāku tuvināšanos ar Vāciju, atteikšanos no ``koridora'', rēķinoties ar kompensācijas saņemšanu citviet. Šai sakarā poļus neapmierinot tikai Lietuvas sagrābšana, viņi uzskatot par nepieciešamu ``pieķert'' līdzi arī daļu Latvijas, nodrošinot izeju pie jūras Liepājā. Līdz ar to latvieši gaidot savu kārtu, uzskatot, ka sākotnēji vācieši atrisināšot poļu-lietuviešu problēmas, pēc tam nākšot Latvijas kārta. Domājams, ka latviešu diplomāta domu izklāsts padomju diplomāta dienasgrāmatā atbilda pirmā teiktajam, kaut viņa vērtējums varēja balstīties arī uz nepilnīgām ziņām.

Šādos apstākļos Ā.~Hitlers piespieda Čehoslovākijas prezidentu E.~Hahu Berlīnē parakstīt vienošanos, kura pilnībā likvidēja valsts neatkarību. 15.~martā Vācijas karaspēks bez cīņas iegāja Čehijā, tika izveidots Bohēmijas un Morāvijas protektorāts (vācu \detxti{Protektorat Böhmen und Mähren}, čehu \cstxti{Protektorát Čechy a Morava}). Ar to bija pierādījusies agresora ``nomierināšanas politikas'' aplamība, sagrauta Lielbritānijas un Francijas valdību kā miera uzturētāju reputācija. Ar Čehoslovākijas valsts likvidēšanu, tās teritoriju sadali, Bohēmijas un Morāvijas protektorāta izveidi nacistiskā Vācija ieguva teritorijas, kuru iedzīvotāju pamatmasa nebija vācieši. Šie notikumi bija lūzuma moments situācijas attīstībā Eiropā. Tieši ar šiem notikumiem faktiski sākās nacistu iecerētā citu Eiropas tautu pakļaušana. Taču, kā norādījusi krievu vēsturniece N.~Naročņickaja, Rietumu historiogrāfijā un arī sabiedrības apziņā Čehoslovākijas sagrābšana un sadalīšana netiek traktēta kā Eiropas kara sākums, jo, to atzīstot, nāktos atzīt arī to valstu atbildību, kas sankcionēja robežu pārdali. Tāpēc arī visas pasaules acu priekšā Vācijas un tās sabiedroto sagrābtā un sadalītā Čehoslovākija netiek uzskatīta par nacistiskās Vācijas upuri tādā mērā, kādā par to tiek uzskatīta vēlāk agresijai pakļautā Polija.

Mūsdienās vairākums poļu vēsturnieku piekrīt uzskatam, ka Lielbritānijas un Francijas realizētā agresora nomierināšanas, piekāpšanās politika (\entxti{appeasement policy}), tā vietā lai mazinātu, tikai vairoja viņa ``apetīti''. Arī Polijas valdība neprotestēja pret Vācijas realizēto faktiski visas Čehoslovākijas sagrābšanu, kaut arī nožēloja, ka pati Polija no upura saņēma pārāk maz. Poļu prese pat apsveica Čehoslovakijas pazušanu no pasaules kartes. Lūk, ko rakstīja avīze ``\pltxti{Kurjer Poranny}'' (``Rīta Kurjers'') 1939.~gada 16.~martā: ``Čehoslovākijas republika, radīta 1918.~gādā, ir pārstājusi pastāvēt. Neizturēja pārbaudījumu arī kroplīgais veidojums, radīts četru lielvalstu Minhenes tikšanās laikā pagājušā gada septembrī \citespace{} Poļu tauta ar gandarījumu sagaida šo faktu.''

Līdz tam Čehoslovākijas sastāvā ietilpstošās Aizkarpatu Ukrainas Seims 15.~martā pasludināja savas valsts~--- ``Karpatu Ukrainas'' (\uktxti{Карпатська Україна}) izveidi. Tā varētu pastāvēt tikai ar Vācijas labvēlību. Tomēr Ā.~Hitlers izšķīrās par šī apgabala nodošanu Ungārijai, kura arī 14.~martā (citi dati~--- 17.~martā) to okupēja un anektēja, ukraiņu karavīri nespēja efektīvi pretoties. Ar to tika novērsta iespēja izveidoties neatkarīgai ukraiņu valstij, kas varētu arī Polijas ukraiņos izsaukt vēlēšanos tai pievienoties, bet Polija ieguva kopēju robežu ar Ungāriju. Amerikāņu pētnieks, Pulitcera prēmijas laureāts H.~Bolduins uzsvēra, ka ``Polija un Ungārija, kā maitasputni plēsa gabalos nospiesto, mirstošo [Čehoslovākijas] valsti''.

Izklāstīto notikumu rezultātā sabruka poļu valdošo slāņu cerības pievienot Polijai lielākus slovāku un rutēņu (Karpatu ukraiņu) apdzīvotos apgabalus. Taču galvenais~--- Polija tagad faktiski no trijām pusēm bija Vācijas un tās sabiedroto Slovākijas un Ungārijas ielenkumā. Vācijas armija vien varēja uzbrukt Polijai no Austrumprūsijas ziemeļos, Vācijas rietumos, Sudetu apgabala un Bohēmijas un Morāvijas protektorāta dienvidos. Tā varēja Poliju saņemt ciešās spīlēs.

Pēc Čehoslovākijas galīgas likvidācijas gluži dabīgi izvirzījās jautājums~--- kurš būs nākamais Vācijas agresijas upuris? Kurp virzīsies Ā.~Hitlers~--- uz austrumiem vai rietumiem? Atbildei ļoti svarīga bija Polijas pozīcija, kura veidotu otro fronti Vācijas uzbrukuma Francijai gadījumā, un kura nodalīja Vāciju no PSRS.

Jau 8.~martā Vācijas politisko, saimniecisko un nacistu partijas augstāko darbinieku sanāksmē Ā.~Hitlers paziņoja, ka pēc Čehoslovākijas okupācijas ``sekos Polijas kārta. Mums nenāksies gaidīt šeit stipru pretestību. Vācijas kundzība pār Poliju nepieciešama, lai apgādātu Vāciju ar lauksaimniecības produkciju un Polijas oglēm.'' 21.~martā Polijas sūtnim Berlīnē oficiāli tika iesniegtas Vācijas prasības par Dancigas iekļaušanu Vācijas sastāvā un eksteritoriāla koridora radīšanu Polijas piejūrā. Nedēļu pēc Čehoslovākijas sabrukuma, 1939.~gada 21.~martā Polijas sūtnis izsaukts otrreiz un viņam paziņots, ka fīrers esot saniknots par to, ka nekas šai jautājumā nevirzās uz priekšu.

Kas varētu draudēt par Vācijas prasību neņemšanu vērā, parādīja Vācijas rīcība pret Lietuvu. Jau pēc Minhenes Lietuva centās pēc iespējas uzlabot attiecības ar Poliju, jo pārāk nospiedoši tai draudēja Vācijas iebrukums, lai atņemtu Klaipēdu. Taču 22.~martā Ā.~Hitlers piespieda Lietuvu atteikties no Mēmeles (Klaipēdas), kur vācieši bija iedzīvotāju lielākā etniskā grupa. Mēmelē ieradās virkne nacistu vadoņu, drīzumā arī Ā.~Hitlers, kurš šeit parakstīja dekrētu par Mēmeles un tās novada iekļaušanu Vācijas valsts sastāvā.

No savas dibināšanas 1252.~gadā Mēmele līdz 1919.~gadam ietilpa vācu zemju sastāvā. Tikai XIX gadsimta beigās tā kļuva arī par lietuviešu kultūras centru. 1910.~gadā Mēmeles apgabalā 44\% iedzīvotāju bija lietuvieši, kaut pašā pilsētā dominēja vācieši. Pēc Pirmā pasaules kara Klaipēdas apgabals kļuva par Tautu Savienības mandātteritoriju. Teritoriju pārvaldīja autonoms landtāgs franču virsvadībā. 1923.~gada janvārī Lietuvas armija ieņēma Klaipēdu un pievienoja to Lietuvas valstij. 1924.~gada martā Tautu Savienība akceptēja Klaipēdas iekļaušanu Lietuvā. Var teikt, ka 1939.~gadā Mēmele (Klaipēda) tika atgriezta Vācijai. Toties 1945.~gadā pēc uzvaras pār Vāciju PSRS daļu Austrumprūsijas~--- Klaipēdu un tās apgabalu nodeva Lietuvas PSR.

Neviena valsts neatbalstīja Lietuvu tai grūtajā brīdī. Vienīgi atzīmējams, ka pēc tam, kad Klaipēdas aneksija jau kļuva par faktu, Polijas tirdzniecības ministrs piedāvāja Lietuvai eksteritoriālu zonu Gdiņas ostā, kā arī palīdzību ostas iekārtošanā Sventajā (lietuv. \lttxti{Šventoji}, mūsdienās Palangas sastāvā.)

Iespējams, ka tieši Mēmeles aneksija, Vācijas iespiešanās Baltijā un tuvošanās padomju robežām piespieda PSRS vadību pārvērtēt savu iepriekšējo kursu un, turpinot panākt kolektīvās drošības sistēmas radīšanu, reizē mēģināt realizēt tam alternatīvu variantu~--- vienoties nevis ar Lielbritāniju un Franciju, bet gan ar Vāciju.

Tam, ka tieši 1939.~gada martā Ā.~Hitlers iesniedza savas prasības Polijai, bija savs pamats. Polijas akcija pret Čehoslovākiju bija nostādījusi Poliju nosacītā starptautiskā izolācijā. Bez tam Vācijas četrgadu bruņošanās plāns jau nesa rezultātus. Pēc tā īstenošanas 1940.~gadā Ā.~Hitlers plānoja karu pret Franciju. Līdz tam vajadzēja novērst iespēju, ka Polija šai karā varētu nostāties Francijas pusē.

Rietumu lielvalstu, īpaši Lielbritānijas, vadītāji, kuri Minhenē bija piekāpušies nacistiskajai Vācijai, faktiski atdevuši tai Čehoslovākiju, 1939.~gada pavasarī pēc tās likvidācijas saņēma savu valstu sabiedrības nosodījumu. Rietumu lielvalstu valdības vairs nevarēja piekāpties Vācijai, agresora ``nomierināšanas'' politika Rietumu sabiedrības acīs jau bija pierādījusi savu kļūdainību. Tā 29.~martā Polijas sūtnis Lielbritānijā E.~Račinskis J.~Bekam ziņoja, ka notikumi Čehoslovākijā ietekmējuši arī britu ``samierinātāju'' [ar Vāciju] nostāju, kuri cerēja, ka ``lieta nonāks līdz krievu-vācu karam, kurš novājinās abas puses, ne bez netieša izdevīguma Rietumu valstīm. Notikumu ātrā norise, kas dod Vācijai vērtīgus iekarojumus bez asins izliešanas, ir atsegusi šādu spriedumu vājās vietas''. Rietumu lielvalstu valdības tagad bija spiestas demonstrēt nepiekāpību un atbalstu Polijai.

Lielbritānijas valdība 19.~martā ierosināja parakstīt četru pušu: Lielbritānijas, Francijas, Padomju Savienības un Polijas deklarāciju, kurā tās apņemtos gadījumā, ja rastos draudi kādai Eiropas valstij zaudēt politisko neatkarību, sarīkot savstarpējas konsultācijas par kopēju pretdarbību.

1939.~gada 15.~martā Vācijas veiktā Čehijas okupācija lika arī Polijas valdībai pārvērtēt savas prognozes par Ā.~Hitlera tālākajiem nodomiem. 23.~martā Polija veica daļēju mobilizāciju Dancigas apkārtnē. Taču pirms tam 22.~martā J.~Beks sarunā ar Lielbritānijas sūtni Polijā H.~Kennardu uzsvēra Polijas nevēlēšanos iekļauties kādās kombinācijās ar PSRS dalību un deva mājienu, ka, ja PSRS tiktu apieta, Polija varētu iesaistīties ciešākā savienībā ar Lielbritāniju un Franciju. 23.~martā, dienā, kad Vācijas karaspēks iegāja Mēmeles apgabalā, J.~Beks Polijai piedāvātās četru pušu deklarācijas vietā, kas no Polijas viedokļa saturēja tikai miglainus solījumus un tikai provocētu Ā.~Hitleru, ierosināja parakstīt divpusēju: britu-poļu paziņojumu.

Acīmredzot tas, ka Polijas valstsvīru skati, tāpat kā gadsimtus iepriekš, bija vērsti Ukrainas, Baltkrievijas un Lietuvas virzienā, viņi vēlējās pakļaut šīs teritorijas, nosacīja viņu vēlmi panākt izlīgumu ar Vāciju. Tas arī noteica Polijas valdības nespēju aiz nereāliem valsts paplašināšanas plāniem visā pilnībā saskatīt tai konkrēti draudošās briesmas. 25.~martā J.~Beks piedāvāja Ā.~Hitleram panākt kompromisu Dancigas jautājumā. Kā norāda poļu vēsturnieks un politologs S.~Debskis, Dancigas jautājumā Polija bija gatava piekrist Tautu Savienības aiziešanai no brīvpilsētas un divpusēja līguma ietvaros vienoties par pilsētas statusu ar Vāciju. Jautājumā par eksteritoriālo autostrādi Polija arī bija gatava diskutēt ar Ā.~Hitleru. Tomēr reizē Polijas ārlietu ministrs arī demonstrēja Vācijai savas valsts nepiekāpību teritoriālajos jautājumos. 26.~martā J.~Beks sarunā ar Vācijas sūtni Varšavā H.~Moltki paziņoja: ``Ja Hitlers skars Dancigu, tas būs karš''. Un piebilda: ``Mēs neesam čehi''. Kad Ā.~Hitleram par to tika paziņots, viņam uznāca dusmu lēkme. J.~Beka izteikumi, valdības lēmumi atspoguļoja lielas daļas poļu noskaņojumu. Mūsdienās Jagelloņu universitātes Krakovā profesors A.~Novaks uzsvēris, ka tikai divi sabiedriski darbinieki un publicisti aizstāvēja viedokli, ka par katru cenu nepieciešams vienoties ar Vāciju. Viens bija profesors V.~Studņickis, kurš arī Otrā pasaules kara laikā aizstāvēja sadarbības ideju ar nacistiem, taču vēlāk viņu par protestiem pret nacistu vardarbību Gestapo arestēja, pēc tam gan atbrīvojot. Otrs bija avīzes ``\pltxti{Słowo}'' (``Vārds'') redaktors S.~Mackevičs, kurš iestājās par piekāpšanos Vācijai ``koridora'' jautājumā, bija pat uz 17 dienām ieslodzīts Bereza Kartuskas nometnē. Abi viņi beidza savu dzīvi emigrācijā Rietumos.

Kā liecinājis poļu ekonomists S.~Svaņēvičs, viņam nācās dzirdēt arī tādus poļu izteikumus, ka ``Hitlers, lūk, nav ieinteresēts pilnīgā Polijas iznīcināšanā, jo tās administratīvais aparāts un armija var viņam noderēt plānu realizēšanai Austrumos. Šī viedokļa piekritēji pieļāva varbūtību, ka Hitlers sagrābs daļu Polijas, piemēram, Gdaņsku un Pomorji, lai piespiestu mūs uzsākt sarunas, kurās viņš ar mums runātu no spēka pozīcijām. \citespace{} Un tas izskatījās loģiski, ja Hitlers gatavojās sākt karu pret Krieviju. Tādā gadījumā mūsu austrumu robežas būtu lielisks placdarms vācu uzbrukumam kā Maskavas, tā Ukrainas, tā arī padomju rūpniecisko rajonu dienvidos virzienā. Bez tam, šķita, ka Hitleram tiešām ir kaut kādi plāni attiecībā pret Poliju. Pat viņa pretenzijas uz Gdaņsku un autostrādes celtniecību Pomorjē bija taču daudz mērenākas nekā tie paziņojumi par ``poļu jautājumu'', kuri skanēja Veimāras republikas laikā.''

Pats S.~Svaņēvičs gan bijis ``pārliecināts par to, ka poļu-vācu karš neizbēgami novedīs pie tā, ka mūsu austrumu teritoriju okupēs Padomju Savienība.'' Šī viņa pārliecība balstījās nevis padomju ārpolitikas ``zināma kareivīgumā'', bet gan vairāk ``veselajā saprātā: Staļins nekad nepieļaus, lai pie Minskas, ceļā uz Maskavu, slikti apbruņotās Polijas vietā, kura bija spiesta transportēt ložmetējus zemnieku ratos, pēkšņi atrastos varenas Hitlera tanku daļas.''

Par oficiālās Polijas militārās propagandas virzienu priekšstatu dod ``Polijas armijas [avīzes]'' (\pltxti{Polska Zbrojna}) ievadraksts ``Mēs esam gatavi'' 1939.~gada marta beigās. Tajā tika paziņots, ka poļi, atšķirībā no čehiem, nejūtas zemāki nekā zemeslodes varenākās tautas. Svešu divīziju skaits poļus nespēj iebiedēt, jo viņu pašu armija, tās apbruņojums un poļu tautas kareivīgais gars ir pietiekams, lai nodrošinātu Polijai uzvaru. Ka tā nebija tikai žurnālistu retorika, liecināja aizsardzības ministra vietnieka ģenerāļa J.~Gluhovska izteikums kādā svarīgā sarunā, ka Vācijas vērmahts ir liels blefs, jo Vācijai trūkt apmācītu rezervju, lai papildinātu savas vienības. Uz jautājumu, vai viņš tic, ka Polija ir militāri pārāka par Vāciju, J.~Gluhovskis atbildēja: ``Pats par sevi saprotams!'' Der atzīmēt, ka tai pat 1939.~~gada pavasarī arī Vācijas Ģenerālštāba priekšnieks F.~Halders, uzrunājot vācu virsniekus, sacīja: ``Skaita ziņā Polijas sauszemes spēki ir piektie lielākie Eiropā. Nebūs aplami, ja teiksim, ka tā kara gadījumā var uzstādīt aptuveni 70~divīziju, pie tam ļoti spēcīgu~--- otru lielāko Eiropā~--- kavalēriju. Bet, mani kungi, tas viss ir tikai šķietamība. Jo pirmkārt Polija nedrīkst nonākt pie visu savu rezervju izsmelšanas, un otrkārt iet runa par tikai lielu baru trūcīgi apmācītu ierindas sastāvu, kuram karā pret vācu reihsvēru nav nekādu izredžu. Šeit mēs runājam arī par vispārējo izglītību. Ja neskaita Rumāniju, poļu karavīrs ir pats dumjākais Eiropā.'' Kā pusgadu vēlāk izrādīsies, vairāk pamatots bija F.~Haldera vērtējums.

Tādejādi Varšava galīgi atteicās pieņemt Vācijas priekšlikumus par teritoriālo noregulējumu. 28.~martā tā paziņoja, ka Dancigas statusa maiņa tiks traktēta kā uzbrukums Polijai. Ar to atkal tika demonstrēta Polijas vēlme saglabāt ``līdzsvara politiku'', ko gan varēja arī traktēt kā nevēlēšanos pat neatkarības zaudēšanas draudu priekšā sadarboties ar Padomju Savienību.

Objektivitātes labad jāpasaka, ka mūsdienās Polijā daudzi neizjūt kaunu par tā laika ambīcijām, pastāv arī tādi uzskati, ka Polija izdarījusi kļūdu, nepiekrītot Ā.~Hitlera priekšlikumiem un kopā ar to nevēršoties pret PSRS. Tā, 2005.~gadā viens no ievērojamiem poļu vēsturniekiem P.~Večorkovičs 1938.~gada oktobrī Vācijas Polijai uzstādītās prasības nosauca par ``piedāvājumu paketi'', kuru pieņemot Polija neciestu nekādus ievērojamus zaudējumus, jo Gdaņska (Danciga) toreiz neesot bijusi poļu pilsēta, bet autostrādes ideju caur ``poļu koridoru'' 30.~gados izvirzījusi poļu diplomātija, mēģinot normalizēt poļu vācu attiecības. Viņš atklāti izteica nožēlu par to, ka 1939.~gadā Polija, kaut arī mēģināja, tā arī nesasniedza vienošanos ar Ā.~Hitleru un kopā ar to nesakāva tik nemīlamo Krieviju. P.~Večorkovičs rakstīja, ka, piekrītot Vācijas priekšlikumam, ``Mēs varējām atrast gandrīz tādu pat vietu Reiha [Vācijas valsts~--- V.Š.] pusē kā Itālija un droši vien labāku nekā Ungārija un Rumānija. Gala rezultātā mēs atrastos Maskavā, kur Ādolfs Hitlers kopā ar Ridzs-Smigliju pieņemtu uzvarējušā poļu-vācu karaspēka parādi. Protams, skumju asociāciju izsauc Holokausts. Taču, ja labi par to padomā, var nonākt pie secinājuma, ka ātra Vācijas uzvara varētu nozīmēt, ka tas vispār nenotiktu.''

Jākonstatē, ka šo poļu vēsturnieku nemulsināja perspektīva Polijai atrasties nacistu nometnē, ka viņš bija gatavs nolaisties līdz klajai holokausta falsifikācijai, tīšuprāt neredzot, ka vēl pirms Vācijas sakāvju sērijas sākšanās~--- jau no 1939.~gada nogales Polijas pilsētās tika ierīkoti geto, 1940.~gada maijā tika likti pamati Osvencimas nāves nometnei, ka 1941.~gada 31.~jūlijā, kad nacistiskais režīms vēl izbaudīja uzvaru eiforiju, reihsmaršals H.~Gērings jau nosūtīja Vācijas Galvenās drošības pārvaldes vadītājam R.~Heidriham direktīvu, ``veikt virspusēju organizatorisku un materiālu sagatavošanos galīgai ebreju jautājuma atrisināšanai vācu ietekmes zonā Eiropā''. Autors šeit gribētu vēl piebilst, ka, ja vēl labāk ``par to padomā'', pēc Maskavas iekarošanas varētu sekot ne tikai ebreju genocīds, bet arī tāds pat nacistiskās Vācijas genocīds pret poļu tautu.

Taču, kā norādījis poļu vēsturnieks M.~Zahariass, 1939.~gadā Ā.~Hitlers jau tiecās izraisīt karu un pretēji J.~Beka cerībām, nekādi racionāli apsvērumi, ieskaitot iespēju, ka Rietumvalstis iesaistīsies konfliktā, nespēja mainīt viņa politiku. Pie tam J.~Beks nerēķinājās ar iespēju, ka Rietumvalstis varētu uzskatīt Polijas fronti tikai kā aizsegu, kurš ļautu tām mierīgi mobilizēt savas armijas, nostiprināt pretgaisa aizsardzību, tā sagatavojoties vairākus gadus ilgstošam karam. Kad 1939.~gada 26.~martā J.~Lipskis nogādāja Berlīnē memorandu ar Polijas piekrišanu autostrādes caur ``poļu koridoru'' celtniecībai, tiesa, bez eksteritorialitātes tiesībām, J.~f.~Ribentrops piedraudēja Polijai ar ``zināmas valsts''~--- Čehoslovākijas likteni.

Var apgalvot, ka 1939.~gada pavasarī jautājumu par kara uzsākšanu pret Poliju fīrers jau bija izlēmis. Jau 25.~martā Ā.~Hitlers paziņoja Vācijas sauszemes spēku pavēlniekam V.~F.~Brauhičam, ka ``tuvākajā laikā sekojošajiem lēmumiem ir jārada īpaši labvēlīgi politiski priekšnosacījumi. Polijai jātiek tā sakautai, lai nākamajos gadu desmitos ar to nevajadzētu rēķināties kā politisku faktoru.'' Fīreram vienīgi vēl nebija stingras apņemšanās pilnībā likvidēt Polijas valstiskumu. Ā.~Hitlers nevienu reizi vien izteicās par t.s. ``Rest-Staat'' (pārpalikuma valsts)~--- poļu marionešu valsts saglabāšanu, ar kuru viņš varētu operēt sarunās ar Rietumiem. Nacistu propaganda plaši pievērsās Dancigas problēmai.

Tad sekoja notikums, kas radīja iespaidu, ka Polijai ir paveicies, jo negaidīti nāca tāds Lielbritānijas solis, uz kuru poļu politiķi nebija cerējuši~--- Polijai tika piešķirtas Lielbritānijas garantijas.

Pēc vienas versijas 29.~martā angļu avīzes ``\entxti{News Chronikle}'' (``Jaunumu Hronika'') Berlīnes korespondents J.~Kolvins nodeva Lielbritānijas Ārlietu ministrijai informāciju, ka ir neizbēgams Vācijas uzbrukums Polijai. Ārlietu ministra vietnieks A.~Kadogans ierosināja piešķirt Polijai garantijas, kuras piespiestu Vāciju pirms uzbrukuma stingri padomāt un, iespējams, no tā atteikties. Priekšlikums tika pieņemts. J.~Kolvina informācija gan vēlāk neguva apstiprinājumu, bet garantijas tika piešķirtas.

Pēc citas versijas ierosinājumu piešķirt Polijai un Rumānijai garantijas izsauca baumas par sagaidāmo Vācijas ultimātu Rumānijai.

31.~martā Lielbritānijas premjerministrs N.~Čemberlens parlamentā deklarēja tās gatavību aizstāvēt Polijas neatkarību. Premjers paziņoja: ``Ja notiks darbība, kas tieši apdraud Polijas neatkarību, un ja Polijas valdība uzskatīs par nepieciešamu izrādīt pretestību ar saviem nacionālajiem spēkiem, Anglijas valdība uzskatīs par savu pienākumu sniegt Polijas valdībai visu iespējamo palīdzību, kas vien būs tās spēkos.'' N.~Čemberlens atzīmēja, ka Čehoslovākijas resursus jau izmanto Vācija, ja arī Polijas resursi nonāks tās rokās, tad tas var radīt ļoti nopietnas sekas Lielbritānijai. Lūk, premjera vārdi: ``Mūsu politikas ģenerāllīniju attiecībā pret Vāciju nosaka ne kādas noteiktas atsevišķas valsts aizsardzība, kura varētu nonākt Vācijas apdraudētībā, bet vēlme novērst vācu kundzības izveidi pār kontinentu, kā rezultātā Vācija kļūtu tik spēcīga, ka varētu apdraudēt mūsu drošību. Vācijas kundzība pār Poliju un Rumāniju palielinātu tās militāro varenību un tāpēc mēs piešķiram garantijas šīm valstīm''. Vēl N.~Čemberlens piebilda, ka viņš runā ne tikai Lielbritānijas, bet arī Francijas vārdā. Tā Lielbritānijas un Francijas garantijas tika piešķirtas Polijai, pēc V.~Čērčila vārdiem ``\dots{}tai pašai Polijai, kura tikai pusgadu agrāk ar hiēnas alkatību piedalījās Čehoslovākijas valsts aplaupīšanā un iznīcināšanā''.

Vācu publicists un vēsturnieks J.~Šusters uzskata, ka garantijas Polijai (vēlāk arī Rumānijai, Grieķijai un Turcijai) N.~Čemberlens piekrita piešķirt tāpēc, ka ticēja tā pareizi iztulkot vēstures mācību. Pēc Pirmā pasaules kara Lielbritānijai taču tika pārmests, ka tā 1914.~gadā nepietiekami skaidri parādījusi ķeizariskajai Vācijai, ka, ja tā uzbruks Francijai, tai būs jākaro arī ar Lielbritāniju. Vadoties no šī apsvēruma, pēc divdesmit pieciem gadiem tad arī Polija saņēma Lielbritānijas garantijas. Ievērojamais britu vēsturnieks A.~Teilors rāda, ka Polija kā sabiedrotais bija īpaši nepieciešama Francijai, tieši tāpēc arī Polija Lielbritānijas acīs ieguva līdz tam nebijušu nozīmi.

Citu skaidrojumu, kāpēc mainījās N.~Čemberlena nostāja pēc Čehoslovākijas sagrābšanas, devis krievu vēsturnieks A.~Šubins. Pēc viņa domām, Ā.~Hitlera atteikšanās no ``ukraiņu jautājuma atrisināšanas'' bija daļēja piekāpšanās poļiem, kurus varēja uztraukt ukraiņu valsts radīšana, jo Polijā pašā ietilpa plaši ukraiņu apdzīvoti apgabali. Taču Rietumu politiķiem Berlīnes atteikšanās no ukraiņu tēmas (runa par Aizkarpatu Ukrainu~--- V. Š.) liecināja par to, ka Ā.~Hitlers ir atteicies no tiešas virzības uz PSRS. N.~Čemberlens uzskatījis, ka vācu fīrers, pakļāvis sev vairākas Austrumeiropas valstis, var tagad vērsties pret Rietumiem. Tieši tāpēc arī 17.~martā N.~Čemberlens paziņojis, ka atsakās no ``nomierināšanas'' politikas, bet 31.~martā devis Polijai Lielbritānijas garantijas, ja tā tiks pakļauta ``tiešai vai netiešai agresijai'', ar ``netiešu agresiju'' saprotot līdzīgu notikumu attīstību, kā ar Čehoslovākiju 1939.~gada martā.

Lielbritānijas štābu priekšnieki gan norādīja, ka vajadzības gadījumā Polijai nebūs iespējams palīdzēt, un par izdevīgāku sabiedroto uzskatīja Padomju Savienību, taču N.~Čemberlens viņu viedokli neņēma vērā.

Britu premjera lēmums zināmā mērā nostādīja Lielbritāniju atkarībā no Polijas realizētās politikas, jo jautājumu par to, vai kas apdraud Polijas neatkarību, lēma Polijas valdība, un Lielbritānija varēja uzsākt rīkoties tikai pēc tās lēmuma. Faktiski Lielbritānija bija piekritusi J.~Beka vēlējumam saglabāt Polijas ``līdzsvara politiku''. Tāpēc atskanēja kritika no dažādām pusēm, vēl jo vairāk tāpēc, ka īpašas simpātijas Rietumeiropas demokrātiskajās valstīs Polija vairs nevarēja izsaukt. Taču jau nākamajā dienā pēc N.~Čemberlena deklarācijas avīze ``\entxti{The Times}'' (``Laiks'') komentēja vienošanos: ``Jaunās saistības, kuras ir uzņēmusies Anglija, neliek Lielbritānijai aizstāvēt tagadējās Polijas robežas katru collu. Izšķirošais vārds šīs deklarācijas tekstā ir nevis ``integritāte'', bet ``neatkarība''''. Šādā izpratnē garantijas Polijai ļāva lielvalstīm pēc vajadzības mainīt tās teritorijas lielumu, kas arī vēlāk notika.

Interesanti, ka pirms Lielbritānijas premjerministra N.~Čemberlena deklarācija par garantijām Polijai izskanēja parlamentā, Lielbritānijas ārlietu ministrs E.~Halifakss iepazīstināja ar to saturu PSRS sūtni Londonā I.~Maiski un ieteica Padomju valdībai uzstāties ar analoģisku paziņojumu. Padomju diplomātija pēc sarunām ar poļu pārstāvjiem no tā izvairījās, norādot kā cēloni uz Polijas nevēlēšanos sadarboties. (4.~aprīlī tika publicēts Padomju Savienības Telegrāfa aģentūras paziņojums, ka pretēji franču avīžu apgalvojumiem, PSRS nav uzņēmusies ``kara gadījumā apgādāt Poliju ar kara materiāliem un slēgt savu izejvielu tirgu Vācijai.'')

3.~aprīlī Londonā ieradās J.~Beks. Briti vēlējās, lai Polija atbalstītu Rumāniju, ja tai draudētu nacisti. Taču Polijas ministrs to atteicās darīt. Britu garantijas J.~Bekam jau bija ``kabatā'', tāpēc viņš nejuta nekādu vajadzību ``jaukt vienu ar otru'' pat attiecībā ar Polijai draudzīgo Rumāniju. 6.~aprīlī Lielbritānija un Polija parakstīja pagaidu vienošanos par savstarpējo palīdzību. (Tā 1939.~gada 25.~augustā tika nomainīta ar pastāvīgu līgumu).

Taču atzīmējams, ka maija sākumā notika britu un franču militāro darbinieku sarunas. Abas puses atzīmēja Vācijas militāro pārākumu uz zemes un gaisā, un, pamatojoties uz šo faktu, izvēlējās stratēģiju, kas ļāva kāpināt to militāro spēku, paredzēja arī sabiedroto armiju iespējamu atkāpšanos ilgstoša, ekonomiski nogurdinoša kara pret Vāciju sākumposmā. Palīdzība Polijai nebija abu valstu prioritāšu skaitā. Vēl vairāk, 4.~maijā tika nolemts: ``Polijas liktenis būs atkarīgs no kara galīgā iznākuma. Tas, savukārt, būs atkarīgs no mūsu spējas panākt Vācijas iespējamo sakāvi, bet nevis no mūsu spējas samazināt spiedienu uz Poliju kara sākumā.''

Berlīnē Lielbritānijas garantijas Polijai tika interpretētas kā pret Vāciju vērsts Polijas līgums ar Apvienoto Karalisti un Polijas gatavība aktīvi piedalīties karā pret Vāciju, kas nozīmē atteikšanos no 1934.~gada Polijas līguma ar Vāciju par neuzbrukšanu un ļauj to denonsēt.

J.~f.~Ribentrops Lielbritānijas un Polijas vienošanos no Polijas puses nosauca par ``kliedzošu'' vācu-poļu 1934.~gada 26.~janvāra deklarācijas pārkāpumu, kura izslēdza spēka lietošanu starp abām valstīm, jo pēc Lielbritānijas un Polijas līguma par savstarpējo palīdzību Polijai Anglijas un Vācijas bruņotas sadursmes gadījumā ``bija pienākums uzbrukt'' Vācijai. Protams, nacistu diplomāts neņēma vērā to, ka pati Vācija ar savām teritoriālajām pretenzijām piespieda Poliju meklēt aizsardzību Lielbritānijā.

Lielbritānijas paraugam sekoja arī Francija, demonstratīvi apstiprinot 1921.~gada Francijas un Polijas savienības līguma spēkā esamību un paziņojot, ka ``Francija un Polija ir devušas savstarpējas tiešas garantijas to vitālo interešu tieša vai netieša apdraudējuma gadījumam''. Tā kā Rietumu lielvalstis bija uzņēmušās saistības pret Poliju, tām vajadzēja rīkoties arī militārajā jomā. Aprīlī un maijā Varšavu apmeklēja apvienota sabiedroto misija. Tā noslēdza vienošanos, ka Vācijas uzbrukuma gadījumā Polijas armijas uzdevums būs apturēt vērmahta uzbrukumu tik ilgi, līdz sabiedrotie Rietumos sagatavos uzbrukumam savas armijas.

Francijas un Polijas militārās sadarbības līgumā bija paredzēts, ka Vācijas agresijas gadījumā pret Poliju ``Francijas armija automātiski uzsāks savu dažādo [ieroču] šķiru bruņoto spēku darbību sekojošā veidā:

\begin{enumerate}
\item Francija uzreiz pēc iepriekš izstrādāta plāna uzsāks gaisa spēku darbību.

\item Tiklīdz daļa no Francijas spēkiem būs gatava (apmēram trešajā dienā), Francija pakāpeniski attīstīs uzbrukuma darbību ar ierobežotiem mērķiem.

\item Tiklīdz iezīmēsies pret Poliju vērstais Vācijas trieciena galvenais virziens, Francija ar saviem galvenajiem spēkiem sāks uzbrukumu Vācijai.'' Francijas Ģenerālštāba priekšnieks ģenerālis M.~Gamelēns solīja, ka piecpadsmitajā dienā pēc mobilizācijas ``lielākā daļa mūsu spēku'' (franču ``\frtxti{le gros de nos forces}) pāries Francijas un Vācijas robežu.
\end{enumerate}

% page 674


Kā parādīja vēlākie notikumi, Francijas apņemšanās palika tikai uz papīra. Francijas saistībām pret Poliju nebija lielāka svara kā savulaik tās saistībām pret Čehoslovākiju.

Jāpiebilst, ka gan Lielbritānijā, gan Francijā bija arī reāli situāciju vērtējoši politiķi, kuri norādīja, ka Polija nespēs izrādīt nopietnu pretestību Vācijai, Rietumvalstīm tāpēc ir jāorientējas uz nopietnu sadarbību ar PSRS. Lielbritānijas valdības piekāpību Polijai nosodīja tāds pieredzējis politiķis kā D.~Loids Džordžs.

Pēc paziņojuma par garantiju piešķiršanu Polijai notika šāda toreizējā premjerministra N.~Čemberlena un politikas veterāna, bijušā Lielbritānijas premjerministra D.~Loida Džordža saruna:

``Loid Džordžs pajautāja premjerministram par PSRS lomu antinacistiskā bloka veidošanā.

``Poļiem un rumāņiem nepatīk tāda ideja,'' atbildēja Čemberlens. ``Tāpēc Krievijas dalība vispār ir problemātiska.''

``Tad kāpēc jūs uzņematies risku ieraut Britāniju karā pret Vāciju?''

``Briesmas ir minimālas,'' atbildēja Čemberlens. ``Mums ir informācija, ka vācieši neriskēs karot divās frontēs.''

``Bet kur tad ir šī ``otrā fronte''?'' pajautāja Loids Džordžs.

``Polija,'' atbildēja Čemberlens. Uz to Loids Džordžs sāka smieties un izzobot Čemberlenu, norādot, ka Polijai nav pieklājīgas aviācijas, ka bruņojums poļu bruņotajos spēkos ir visai viduvējs, ka ekonomiski un iekšpolitiski Polija ir vāja.''

1939.~gada 3.~aprīlī, uzstājoties pārstāvju palātā, D.~Loids Džordžs paziņoja: ``Ja Krieviju neiesaistīja [līgumā] tikai noteiktu poļu jūtu dēļ, kuri negrib samierināties ar krievu klātbūtni savā valstī, mums ir jāuzstāda tāda klātbūtne kā noteikums, un ja poļi nav gatavi pieņemt šo vienīgo noteikumu, kuram pastāvot, mēs varam sniegt viņiem rezultatīvu palīdzību, tad viņiem pašiem jānes [par to] atbildība.'' Taču pie varas Rietumvalstīs atradās tādi cilvēki kā Lielbritānijas valdības vadītājs N.~Čemberlens, kurš nepareizi vērtēja Vācijas un Polijas spēku samēru. D.~Loida Džordža brīdinājums tika ignorēts. Ja tas tā nebūtu, notikumi pasaulē varētu ievirzīties citādāk.

Vēsturiskajā literatūrā ir izteikti ļoti dažādi viedokļi par Lielbritānijas garantijām Polijai. Ir izskanējis viedoklis, ka Lielbritānijas un Francijas garantijas Polijai lika tām iestāties karā pret Vāciju bez nepieciešamības, jo Ā.~Hitlers pēc Polijas sagrāves varētu veikt agresiju pret PSRS, un Lielbritānija un Francija varētu palikt malā. Angļu vēsturnieks Dž.~Grenvils gan norādījis, ka argumenti, ka abas Rietumu valstis tiešām vēlējušās panākt Vācijas uzbrukumu PSRS, ir visai apšaubāmi. Pie tam konkrētie fakti ļauj noraidīt pieņēmumu, ka Ā.~Hitleram tai laikā bija tāds nodoms.

Tuvs iepriekšminētajam ir viedoklis, ko izteica amerikāņu žurnālists V.~Šīrers, kurš bija gadu desmitiem sekojis notikumiem Polijā. 1939.~gada 6.~aprīlī viņš savā dienasgrāmatā ierakstīja: ``Šajā nedēļā es ēdu-dzēru ar daudziem poļiem~--- no Ārlietu ministrijas, militāristiem, Pilsudska vecajiem leģionāriem, kuri vada poļu radio. Un visi viņi nevēlas saprast, ka tā ir nepieļaujama greznība būt vienlaikus gan par Krievijas, gan Vācijas ienaidniekiem, ka viņiem ir jāizvēlas, ka, ja viņi iesaistīs palīdzības sniegšanā Krieviju kopā ar Franciju un Lielbritāniju, tad viņi ir glābti.'' Acīmredzot, ņemot vērā Polijas valdītāju rīcības stūrgalvību, amerikāņu žurnālists Anglijas garantiju piešķiršanu Polijai komentēja sekojoši: ``Var jau apdrošināt pulvera rūpnīcu, ja tajā tiek ievēroti drošības noteikumi, taču apdrošināt rūpnīcu, kura pilna ar trakajiem, ir nedaudz bīstami''.

Pastāv uzskats, ka Lielbritānijas galvenais uzdevums bijis novilcināt kara sākumu un tas arī izdevies, jo Rietumos (pret Franciju) īsts (nevis ``dīvainais'', ``sēdošais'') karš sākās tikai 1940.~gada 10.~maijā.

Īpašus pārmetumus Lielbritānijas valdībai gan politiķu, gan vēsturnieku vidū ir izsaucis apstāklis, ka reāli Lielbritānija kā jūras lielvalsts ar lielu floti, bet mazu sauszemes armiju nekādu būtisku palīdzību Polijai nespēja sniegt. Bez tam garantijas Polijai un Rumānijai deva zināmu drošības sajūtu arī PSRS, jo Vācijai, lai uzbruktu tai, palika tikai vienīga šaurs koridors caur Baltijas valstīm, kas šādu uzbrukumu pārvērta par praktiski neiespējamu. PSRS turpmāk vairs nebija ieinteresēta savu drošības interešu vārdā meklēt kompromisus ar Rietumu lielvalstīm, bet par atbalstu tām varēja prasīt adekvātu pretimnākšanu savām interesēm. Ar to Lielbritānija, bet pēc tam arī Francija, kas tai pievienojās, bija radījušas labvēlīgu situāciju PSRS valdībai, kura nu varēja ``izspēlēt'' dažādus politiskās spēles variantus.

Bez tam atzīmējams, ka pēc Polijas galīgas atteikšanās piekāpties Vācijai (26.~marts), pēc N.~Čemberlena uzstāšanās parlamentā (31.~marts), pēc militāra plāna uzbrukumam Polijai izstrādes sākuma (aprīlis) jau agrāk dažu vācu diplomātu lolotie plāni par iespēju kādi vienoties ar Padomju Savienību izplatījās arī Ā.~Hitlera tuvāko līdzgaitnieku vidū, bet vēlāk pārņēma arī pašu fīreru. No šī laika sākās vācu diplomātu zondāža, pētot padomju puses gatavību vest sarunas ar Vāciju. Aprīlī notika pirmais zondāžas mēģinājums: Vācijas Ārlietu ministrijas valsts sekretāra E.~Veiczekera, vēlākā Vācijas Federatīvās Republikas prezidenta R.~Veiczekera tēva, saruna ar PSRS sūtni Berlīnē A.~Merekalovu.

Uz šiem pārmetumiem vismaz daļēji ir atbildējis amerikāņu vēsturnieks G.~Veinbergs, kurš norādījis uz trim secinājumiem, no kuriem vadījās Lielbritānijas valdība, pasludinot garantijas.

\begin{enumerate}
\item Ja pirms tam Lielbritānija centās par katru cenu novērst pat maza kara izraisīšanu Eiropā, jo vienalga lai kur tas izraisītos~--- austrumos vai rietumos, tas varētu aptvert visu kontinentu un Lielbritānija tik un tā tajā tiktu ierauta, tad pēc Čehoslovākijas iznīcināšanas Lielbritānijai bija jāatsakās no vecās stratēģijas un jārod jauna. Jaunās~--- 1939.~gada stratēģijas uzdevumi bija: skaidri demonstrēt savu pozīciju, tā iebiedējot Vāciju, uzmundrinot potenciālo upuri un cenšoties piesaistīt tā pusē citas valstis;

\item Ja 1938.~gadā Vācijas propaganda par it kā apspiestajiem 3~miljoniem Sudetu vāciešu atstāja ietekmi uz Lielbritānijas sabiedrību un arī valdību, tad 1939.~gadā, zinot, ka vācu mazākumam Polijā tiešām ir par ko sūdzēties, briti apzinājās, ka šīs vācu mazākuma sūdzības ir Vācijas propagandas līdzeklis. Bija tikai viens solis kā nacistiskā Vācija varētu pierādīt pretējo~--- atjaunot Čehoslovākijas neatkarību, ko tā nekad nedarītu. Britu sabiedrība ietekmēja valdību stingra kursa pret Vāciju ieturēšanas virzienā;

\item Lielbritānijas demonstrētajam atbalstam Polijai bija nepieciešami sabiedrotie, jo pati Lielbritānija to nodrošināt pietiekamā apjomā lai aizsargātu Polijas neatkarību, nespēja. Tāpēc tika meklētas tuvināšanās iespējas ar PSRS.
\end{enumerate}

Savukārt profesors Ē.~Jēkabsons pamatoti uzskata, ka Lielbritānija un Francija bija ieinteresētas, lai Polija nenonāk atkarībā no Vācijas, jo tādā gadījumā Vācija iegūtu pilnīgu drošību austrumos un dienvidos, līdz ar to~--- pilnīgu rīcības brīvību savai pret britiem un frančiem vērstajai politikai rietumos. Polija Lielbritānijai un Francijai bija vajadzīga kā militārs spēks pret Vāciju. Tāpēc tai arī tika solīta militāra palīdzība.

Krievu vēsturniece V.~Beļajeva uzsvērusi, ka garantiju piešķiršana Polijai nenozīmēja to, ka N.~Čemberlens būtu atteicies no savas agresora ``nomierināšanas'' politikas. Apstākļi laikā pēc Minhenes vienošanās noslēgšanas bija mainījušies, ``nomierināt'' Ā.~Hitleru ar Polijas teritoriju bija neiespējami~--- sašutums pasaulē un arī Lielbritānijā par Čehoslovākijas iznīcināšanu bija pārāk liels. Tad arī Lielbritānijas valdība noslēdza līgumu ar Poliju. Tā mērķis bija iepriekšējais~--- nekādā gadījumā nepieļaut atklāta militāra konflikta izraisīšanos. Garantijas Polijai bija galvenokārt spiediena līdzeklis uz Vāciju, lai tās vadība apzinātos, kādas briesmas tai draud kara pret Poliju izraisīšanas gadījumā. Garantijas nesasniedza cerēto, jo Lielbritānijas politika palika iepriekšējā. Atstāstītais spriedums ir loģisks un atspoguļo lietas būtību. Iespējams, ka garantijas Polijai tiktu piešķirtas arī bez izteiktas Polijas pretestības Vācijas spiedienam, tik nokaitēta bija atmosfēra pēc Čehoslovākijas valstiskuma likvidācijas.

Ukraiņu publicists S.~Lozuņko par Lielbritānijas un Francijas katastrofālu kļūdu uzskata garantiju piešķiršanu Polijai, pirms tam nepiesaistot garantiju līgumam PSRS. Lielbritānijas valdībai tai brīdī rokās bijušas visas iespējas izdarīt diplomātisku spiedienu uz Poliju, lai tā piekristu sadarboties ar PSRS. Vajadzējis tikai kā garantiju piešķiršanas nosacījumu izvirzīt prasību piesaistīt savienībai PSRS, jo tikai tādā gadījumā Ā.~Hitlers tiktu iespiests spīlēs, kuras varētu apturēt viņa plānoto agresiju. Acīmredzot esot bijis sākotnēji jānoslēdz trīspusēja Lielbritānijas, Francijas un PSRS vienošanās par kopīgu pretošanos nacistiskās Vācijas agresijai, bet pēc tam jānostāda Polija izvēles priekšā, vai nu tā pieņems triju valstu garantijas, vai paliks viena pret vienu ar Vāciju. Jāatzīst, ka konkrētajos apstākļos šāda notikumu gaita nebija reāla. Tur jau tā lieta, ka Rietumvalstis nevēlējās noslēgt nedz tās saistošu līgumu ar PSRS, nedz izdarīt spiedienu pret Poliju, lai piespiestu to pieņemt padomju puses tai izvirzītās prasības.

Tāpat kā agrāk PSRS brīdinājums neatturēja Poliju no agresijas pret Čehoslovākiju, tā tagad Lielbritānijas garantijas nespēja atturēt Vācijas agresiju pret Poliju. Gluži pretēji, 31.~marta Lielbritānijas garantijām Ā.~Hitlers atbildēja ar naktī uz 4.~aprīli izdoto pavēli sagatavoties uzbrukumam Polijai. 1939.~gada 28.~aprīlī runā reihstāgā Ā.~Hitlers paziņoja par 1934.~gadā ar Poliju parakstītā neuzbrukšanas līguma denonsāciju. Savā laikā Vācijai noderējušais līgums tagad tai vairs nebija vajadzīgs. Kā denonsācijas ieganstu Ā.~Hitlers izmantoja tikko Polijai piešķirtās Lielbritānijas garantijas: ``\dots{}kuras, noteiktos apstākļos, piespiedīs Poliju uzsākt kara darbību pret Vāciju gadījumā, ja notiktu sadursme starp Vāciju un citu valsti, kurā savukārt piedalīsies Anglija. Šis apstāklis ir pretrunā ar vienošanos, kuru es kādu laiku atpakaļ noslēdzu ar maršalu Pilsudski \citespace{} Tāpēc es uzskatu, ka vienošanos vienpusīgi ir lauzusi Polija un tā, tādā kārtā, vairs nedarbojas.''

Interesanta ir Vācijā dzīvojošā bijušā padomju diplomāta, politiķa un vēsturnieka V.~Faļina versija par Lielbritānijas garantijām Polijai, kuru viņš pamato ar nacistu vadoņu izteikumiem. Proti: 1939.~gada jūnijā H.~Gērings sarunā ar Lielbritānijas sūtni Vācijā N.~Hendersonu izmeta frāzi: ja Londona vēl pagaidītu ``kaut 10~dienas'' ar garantiju piešķiršanu Polijai, situācija izveidotos pavisam citāda. To pašu atkārtoja Ā.~Hitlers Itālijas ārlietu ministram G.~Čano 12.~augustā un Ā.~Hitlera vietnieks nacistu partijā R.~Hess lordam-kancleram D.~Saimonam 1941.~gada maijā. Pēc R.~Hesa vārdiem, poļi sliecušies pieņemt vācu nosacījumus gan Dancigas, gan ``poļu koridora'' jautājumā, taču mainījuši savu pozīciju Lielbritānijas spiediena rezultātā. Nacistu vadoņu izteikumus apstiprinot arī franču avoti. Pēc V.~Faļina viedokļa, briti apsteidza J.~Beka sagatavoto soli~--- vizītes laikā Lielbritānijā paziņot, ka Polija ir zaudējusi uzticību Parīzei un Londonai un pārslēgsies uz sadarbību ar nacistisko Vāciju. Tāpēc jau nedēļu pirms J.~Beka ierašanās Londonā N.~Čemberlens nāca klajā ar savu paziņojumu par garantijām Polijai, ja tai tiks uzbrukts. V.~Faļins uzskata, ka tikai Lielbritānijas spiediens licis Polijas valdībai atteikties no piekāpšanās Vācijai. Arī poļu vēsturnieks P.~Večorkovičs uzskatīja, ka britu garantijas ``pārsvītrojušas'' iespējamo poļu-vācu vienošanos, jo līdz tam laikam poļi gan nevēlējās pieņemt vācu priekšlikumus, bet sarunas vēl bija iespējamas. Pēc būtības nosauktie vēsturnieki ar to vismaz daļēji pievienojās jau J.~f.~Ribentropa Nirnbergas tiesas procesā izteiktajam apgalvojumam, ka Vācijas~--- Polijas kara sākuma 1939.~gada septembrī ``izšķirošais faktors bija angļu garantijas Polijai \citespace{} Šīs garantijas apvienojumā ar poļu mentalitāti padarīja neiespējamas sarunas ar Poliju vai kaut kādas savstarpējas attiecības ar poļiem''.

Pēc šī darba autora domām, noteicošā notikumu attīstībā Polijas attiecībās ar Vāciju tomēr bija nevis Lielbritānijas, bet pašu Polijas politiķu izšķiršanās. Darbojās inerces spēks. Līdz pēdējam brīdim Polijas vadītāji nevēlējās iesaistīties nedz Vācijas vadītā pretpadomju savienībā, nedz Rietumvalstu un PSRS veidojamā drošības sistēmā pret Vāciju (izņēmums bija Polijas savienība tikai vienīgi ar Rietumvalstīm), tā turpinot iepriekšējo ``līdzsvara politiku'' arī tad, kad valstij vairs nebija resursu tās īstenošanai.

1934.~gada pakta denonsēšanu var vērtēt arī kā visas iepriekšējās J.~Pilsudska un J.~Beka diplomātiskās darbības rezultātu un Polijas ``līdzsvara politikas'' sabrukumu. 1939.~gada aprīlī-jūnijā tapa Ā.~Hitlera direktīvais plāns ``\detxti{Fall Weiß}'' (Plāns Baltais), kurš paredzēja dot triecienu Polijai ne vēlāk kā 1939.~gada septembrī. Ne pats Ā.~Hitlers, ne Vācijas prese neslēpa, ka Poliju gaida Čehoslovākijas liktenis. 2.~maijā Vācijas Ārlietu ministrijas Austrumu biroja padomnieks P.~Kleists sarunā ar vācu žurnālistu atklāja, ka Ā.~Hitlers jau februārī secinājis, ka sarunu ceļā Poliju nevar piesaistīt Vācijai un nolēmis ar spēku nostādīt to ``uz ceļiem''. Vācijas priekšlikumi Polijai izteikti jau ar pārliecību, ka tie tiks noraidīti. Aprēķins bijis pareizs un Vācija varēja tikt vaļā no līguma par neuzbrukšanu. 23.~maijā vērmahta vadītāju sanāksmē Ā.~Hitlers paziņoja, ka ``Polija vienmēr būs mūsu pretinieku pusē''. Kaut vēl 18.~maijā Polija bija paziņojusi Vācijai, ka tā ``var būt pārliecināta, ka Polija nekad neļaus savā teritorijā ieiet nevienam Padomju Krievijas karavīram'', ar to cenšoties pierādīt ka tā ir ``barjera pret boļševismu'', pēc fīrera domām Polija nevarēja pretendēt uz šādu lomu. ``Tātad, jautājums par to, lai pažēlotu Poliju, atkrīt,~--- rezumēja Ā.~Hitlers,~--- atliek lēmums: pirmajā izdevīgajā gadījumā uzbrukt Polijai''. Tādejādi Vācijas lēmums par uzbrukumu Polijai bija pieņemts jau vairākus mēnešus pirms 23.~augustā noslēgtā Molotova-Ribentropa pakta. Tieši 23.~maiju krievu vēsturnieks A.~Djukovs uzskata par datumu, kad tika pieņemts principiāls lēmums par karu pret Poliju. Padomju izlūks R.~Zorge jau 1939.~gada maijā ieguva ziņas par Vācijas lēmumu uzbrukt Polijai septembrī. Protams, to tūlīt uzzināja PSRS, bet ne Polijas vadība.

5.~maijā Polijas ārlietu ministrs J.~Beks Seimā paziņoja: ``Mieram \citespace{} ir sava cena, bet tā ir izmērāma. Mēs Polijā nepazīstam jēdzienu ``miers par katru cenu''''. Jāuzsver, ka Polija, piekrītot noslēgt līgumu par savstarpēju palīdzību ar Lielbritāniju, nostājās uz kolektīvās drošības ceļa, kaut turpināja atteikties no sadarbības šai jomā ar PSRS. Vācu vēsturnieks J.~Hakers 1983.~gadā izdotā grāmatā izteica viedokli, ka tieši 1939.~gada martā Polija no potenciāla Vācijas partnera pārvērtās par nopietnu šķērsli pēdējās ekspansijai uz austrumiem, kas arī bijis izšķirošs pagrieziena punkts uz Vācijas un PSRS tuvināšanos.

Maršals E.~Ridzs-Smiglijs tagad pret gaidāmo padomju uzbrukumu gatavoto plānu ``Austrumi'' nolika pie malas un pavēlēja sīkāk izstrādāt plānu ``Rietumi''. No 1939.~gada marta-aprīļa Polijā lielāka uzmanība tika veltīta sagatavotības palielināšanai karam ar Vāciju. Līdz tam galvenie līdzekļi tika ieguldīti nocietinājumu celšanai Austrumos un visi kara varianti tika izstrādāti paredzot pirmkārt sadursmi ar PSRS. Polijas militāristi uzskatīja, ka viņi ir gatavi stāties pretī Vācijas karaspēkam arī bez sabiedrotajiem, kas gan jāvērtē kā smaga savu spēku pārvērtēšana un pretinieka spēku nenovērtēšana. 1939.~gada aprīļa vidū nacionālistiskā organizācija ``\pltxti{Polski Związek Zachodni}'' (Poļu Rietumu savienība), kuru stingri kontrolēja Ārlietu ministrija, izdeva skrejlapu, kas aicināja tautiešus paaugstināt valsts aizsardzības spējas, domājot par Polijai atņemtajiem (un Vācijai piederošiem) apgabaliem: Marienburgu, Ermlandi (\pltxti{Warmia}), mazūru, kašubu zemēm, Pozenes robežapgabalu, Opoļjes Silēziju. To varēja saprast kā aicinājumu panākt minēto zemju pievienošanu Polijai.

1939.~gada aprīlī un maijā notikusī vācu un poļu diplomātu apmaiņa ar asām vācu prasībām un tikpat asām poļu atbildēm nedeva rezultātu. Kā raksta N.~Deiviss, Ā.~Hitlers pat nedomāja uzklausīt Rietumu diplomātus vai savus līdzgaitniekus, piemēram, H.~Gēringu vai B.~Musolīni, kuri vēlējās paildzināt mieru. 23.~maijā notikušajā militāristu sanāksmē viņš plaši izrunājās par Vācijai nepieciešamo ``dzīves telpu'' (\detxti{Lebensraum}) Austrumos un par to, ka drīz sāksies karš. 14.~jūnijā viņš saviem ģenerāļiem lika sagatavoties kara darbībai astoņu nedēļu laikā. 22.~augustā, kad šīs astoņas nedēļas bija aizritējušas, kādā citā sanāksmē Berghofā viņš paziņoja, ka ``karu labāk sākt tūlīt''.

Šādā atmosfērā Polijā pieauga naids pret vāciešiem. Kā ziņoja Latvijas sūtnis Polijā L.~Ēķis, Polijas pilsētā Bidgoščā (vācu \detxti{Bromberg}) un nesen Čehoslovākijai atņemtajā Boguminā (morāvu \cstxti{Bohumin}, vācu \detxti{Oderberg}) notika poļu un vācu iedzīvotāju sadursmes, kur Boguminā civilā ģērbusies poļu policija vācu demonstrācijā ``iejaukusies ar savām gumijas nūjām un piekāvusi vāciešus tādā mērā, ka kādi 25~vīri palikuši guļot uz pilsētas laukuma.''

Poļu politiķu vairākums pārvērtēja savas valsts militāros spēkus. Tā 1939.~gada 3.~jūlijā Cešinā notika Nacionālās vienības nometnes (\pltxti{Obòz Zjdnoczenia Narodowego})~--- organizācijas, kura tobrīd veidoja valdību un it kā pārstāvēja tautu, sapulce. Tajā minētās organizācijas Silēzijas iecirkņa pārstāvis klaji paziņoja: ``Nākamajā karā Polija nezaudēs neko no saviem ieguvumiem, gluži pretēji, šajā karā, kurš mums beigsies uzvaroši, mums jāatgūst Opoles (poļu \pltxti{Opole}, vācu \detxti{Oppeln}) Silēzija un Austrumprūsija.'' (Var tikai pabrīnīties, ka šādu pārliecību daudzi poļi saglabāja arī pēc 1939.~gada septembra sagrāves.)

Tikai ļoti nelabprāt, kad karš jau stāvēja durvju priekšā, Rietumvalstu valdības uzsāka sarunas ar PSRS vadību par kolektīvas drošības sistēmas radīšanu. Un arī tad lielā mērā tikai sabiedriskās domas spiediens lika Lielbritānijas un Francijas valdībām par to izšķirties. Francijas ārlietu ministrs Ž.~Bonne uzsvēra: ``Tagad Francijas un Britānijas sabiedriskajā domā veidojas tāda varena kustība par vienošanos ar PSRS, un visā pasaulē, ieskaitot Franciju, milzīgā cilvēku daudzumā, pat ar pašiem mērenākajiem uzskatiem, tādā mērā pieaug pārliecība, ka tieši no tā ir atkarīgs miers, ka sarunu neveiksmes gadījumā ir nepieciešams par katru cenu vainu par to uzkraut Padomju Savienībai''. Ar šī uzdevuma izpildi jau pirms Otrā pasaules kara nodarbojās Rietumvalstu politiķi, pēc kara~--- arī daudzi to angažētie vēsturnieki.

Kad 11.~aprīlī Lielbritānija zondēja PSRS, ar ko tā varētu nepieciešamības gadījumā palīdzēt Vācijas apdraudētajām valstīm, PSRS ārlietu tautas komisārs M.~Ļitvinovs, rakstot par to vēstulē padomju pilnvarotajam Parīzē, atzīmēja, ka Londona un Parīze acīmredzot uzskata, ka tas atbilstu padomju interesēm un turpināja: ``Taču savas intereses mēs paši apzināsimies un darīsim to, ko tās mums diktē. Kāpēc gan mums dot iepriekš solījumus, negūstot no šiem apstākļiem sev it nekādu labumu?'' Tā redzam, nostāja bija gauži pragmātiska. Lielbritānija 14.~aprīlī oficiāli ierosināja PSRS nākt klajā ar paziņojumu, ka ``agresijas gadījumā pret kādu Padomju Savienības kaimiņvalsti Eiropā, kura izrādītu pretestību [agresijai], tā varēs rēķināties ar Padomju valdības palīdzību, ja tā izrādīsies vēlama, pie tam šī palīdzība tiks sniegta tādā veidā, kurš tiks uzskatīts par visērtāko''. Arī Francija no savas puses iesniedza līdzīgu priekšlikumu. Faktiski tas nozīmēja, ka PSRS no savas puses piešķirs garantijas ne tikai Polijai un Rumānijai, bet arī Baltijas valstīm un Somijai, kurām nebija nekādu Lielbritānijas un Francijas garantiju. PSRS, saprotot, ka konflikta gadījumā Rietumvalstis atkal varētu tāpat kā Minhenē vienoties ar agresoru, bet PSRS tiktu ierauta karā, noraidīja vienpusīgu garantiju piešķiršanu, taču piedāvāja noslēgt visu triju lielvalstu kopīgu savstarpējās palīdzības līgumu gadījumam, ja pret kādu no tām tiks izdarīts uzbrukums; sniegt atbalstu visām valstīm, kuras robežojās ar PSRS; steidzīgi noslēgt vienošanos par konkrētām savstarpējās palīdzības formām un vienošanos par to, ka dalībvalstis neslēgs separātu mieru.

1939.~gada 16.~maijā Lielbritānijas valdība izskatīja triju karaspēka šķiru štāba priekšnieku memorandu, kurā tie Lielbritānijas-Francijas un PSRS savstarpējas palīdzības līguma iespējamo nenoslēgšanu nosauca par ``diplomātisku sakāvi, kurai būs nopietnas militāras sekas''. Lielbritānijas ārlietu ministrs E.~Halifakss savukārt uzskatīja, ka politiskie argumenti pret pakta noslēgšanu ar PSRS pārsver militāros argumentus par to. Premjerministra N.~Čemberlena pozīcija bija vēl kategoriskāka~--- 19.~maijā parlamentā viņš paziņoja, ka ``drīzāk atkāpsies no amata, nekā parakstīs līgumu ar Padomēm''.

Rietumvalstu un PSRS sarunu vešanu traucēja savstarpējā neuzticība, abu pušu vēlme panākt pēc iespējas vairāk, pretī piedāvājot pēc iespējas mazāk. E.~Halifakss 1939.~gada 4.~jūlijā valdības sēdē paziņoja: ``Mūsu galvenais mērķis sarunās ar PSRS ir panākt, lai Krievija nenodibinātu kādas attiecības ar Vāciju''.

1939.~gada jūlija sākumā N.~Čemberlens sarunājās ar savu draugu aviācijas ministru K.~Vudu. Kad pēdējais pajautāja: ``Kas jauns sarunās par paktu?'', N.~Čemberlens sapīcis atmeta ar roku un atbildēja: ``Es vēl neesmu atmetis cerību, ka izdosies izvairīties no šī nelaimīgā pakta parakstīšanas.''

V.~Čērčils ir publicējis 1939.~gada 26.~martā viņam nosūtīto N.~Čemberlena vēstuli, kura atklāja toreizējā premjerministra nostāju un nevēlēšanos noslēgt kolektīvās drošības līgumu ar PSRS. N.~Čemberlens rakstīja: ``Man jāatzīstas, ka Krievija manī izsauc dziļu neuzticību. Es nemaz neticu tās spējai realizēt iedarbīgu uzbrukumu, pat ja tā to gribētu. Un es neuzticos tās motīviem, kuriem, pēc manām domām, ir maz kas kopīgs ar mūsu brīvības idejām. Tā vēlas tikai saķildot visus pārējos. Bez tam, daudzas no mazajām valstīm, īpaši Polija, Rumānija un Somija, izturas pret to ar naidu un aizdomām.''

ASV iekšlietu ministrs H.~Ikess savā slepenajā dienasgrāmatā ierakstīja: ``Anglija turpināja sevi krāpt ar ilūziju, ka viņai izdosies izraisīt Krievijas un Vācijas sadursmi un tā pašai iziet sausai no ūdens''. Pat instruējot 2.~augustā Lielbritānijas militārās delegācijas vadītāju sarunās Maskavā admirāli P.~Draksu, E.~Halifakss uzdeva viņam ``vilcināt sarunas iespējami ilgāk''. ``Ilgāk''~--- nozīmēja līdz septembra beigām, kad, pēc Lielbritānijas vadītāju viedokļa, to nepieciešamība vispār atkristu klimatisko apstākļu dēļ, kas būtu nelabvēlīgi Vācijas iebrukumam Polijā, bet līdz nākamā gada pavasarim situācija varēja mainīties. P.~Draksam izsniegtās rakstiskās instrukcijas skanēja: ``Britu valdība nevēlas uzņemties kādas konkrētas saistības, kuras varētu sasaistīt mums rokas tādos vai citādos apstākļos. Tāpēc ir jācenšas militāro vienošanos novest līdz pašiem vispārīgākajiem formulējumiem''. (Padomju izlūkdienests saņēma informāciju par instrukciju saturu, un, sākoties sarunām, padomju delegācija uzreiz piedāvāja ļoti konkrētu militārās vienošanās projektu.) Faktiski ar to Rietumvalstu vadītāju uzskatos jau pašas sarunas, arī tās vainagojošs līgumu ar PSRS, kurš jau pats par sevi, ar savu eksistenci vien, neatkarīgi no abu pušu konkrētās gatavības piedalīties militārā konfliktā, bet nevis to rezultātā panākamā konkrētā vienošanās par kopdarbību, kļuva par līdzekli, ar ko apturēt Vācijas agresiju, piespiest Ā.~Hitleru pārtraukt agresīvos soļus, likt viņam meklēt miermīlīgus strīdīgo jautājumu risinājumus. (Jāsaka gan, ka šādi Rietumvalstu aprēķini bija utopiski, nedz nacistiskais režīms bija spējīgs būt miermīlīgs, nedz Vācijas tautsaimniecība, bruņošanās izdevumu novājināta, iztikt bez iekarojumiem. Arī Rietumu valdību cerības panākt PSRS piekrišanu karot pret Vāciju tai nelabvēlīgā situācijā, bez citu valstu atbalsta neattaisnojās.) Kā raksta amerikāņu vēsturnieks T.~Fārenbahs, N.~Čemberlens un E.~Daladjē ``nebija zaudējuši cerību, ka Hitlers un Staļins uzsāks cīniņu viens pret vienu'' un Rietumi varēs ``novērot abu pretinieku cīņu no sev droša attāluma''.

Ir saprotams, ka arī Rietumvalstu sabiedrība nealka iesaistīties plēsonīgās Polijas, kura tikko bija piedalījusies Čehoslovākijas saplosīšanā, aizstāvēšanā pret otru, lai arī lielāku plēsoņu~--- Vāciju. Tikai sapratne par to, ka ar Polijas sagrābšanu Ā.~Hitlers ``nenomierināsies'', lika sabiedrības lielākajai daļai iestāties par agresijas apturēšanu.

Neraugoties uz jau gūto negatīvo pieredzi, Rietumvalstīs bija arī politiķi, kuri, cenšoties nomierināt agresoru, prasīja, lai Polija piekāpjas atsevišķos jautājumos. Viens no izlīgšanas ar Vāciju sludinātājiem Francijā, vēlākais kolaborants M.~Dea 1939.~gada 4.~maijā avīzē ``\frtxti{L’Œuvre}'' (``Darbs'') publicēja rakstu ``\frtxti{Faut-il mourir pour Dantzig?}`` (``Vai jāmirst par Dancigu?'') Tajā viņš argumentēja, ka Francijai ir jāizvairās no kara ar Vāciju. Raksts un tā izraisītā diskusija nepalika bez ietekmes, taču tautas vairākums vairs nevēlējās atbalstīt agresora nomierināšanas politiku. Kā rādīja Franču sabiedriskās domas un tirgus pētniecības institūta (\frtxti{Institut français d'opinion publique}) 1939.~gada jūlijā veiktajā aptaujā uz jautājumu: ``Vai jūs domājat, ka gadījumā, ja Vācija mēģinātu ieņemt Dancigu, mums, ja vajadzēs, pat vardarbīgi, tā ir jāattur?'' 76\% atbildēja apstiprinoši, 17\% noraidoši, 7\% nebija sava viedokļa. Tad pat Lielbritānijā uz Britu sabiedriskās domas institūta (\entxti{British Institute of Public Opinion}) jautājumu ``Ja Vācija un Polija sāks karu par Dancigu, vai mums jāpilda savas saistības un jācīnās Polijas pusē?'' apstiprinoši atbildēja 76\% aptaujāto, noraidoši~--- 13\%, nebija sava viedokļa~--- 11\%. Tātad, kad jau pēc Vācijas iebrukuma Polijā septembra sākumā Francijas un Lielbritānijas valdības pieteica Vācijai karu, tās rīkojās atbilstoši savu tautu vairākuma gribai.

Amerikāņu vēsturnieks D.~Hogans 1961.~gadā izdeva darbu, kurā par kara izraisīšanu galvenokārt vainoja Lielbritānijas un Polijas ārlietu ministrus E.~Halifaksu un J.~Beku. Mūsdienās viedokļu plurālisma apstākļos krievu publicistikā ir izskanējis uzskats, ka Ā.~Hitlera prasības esot bijušas pamatotas, jo pat uz ``poļu koridoru'', kurš pēc Versaļas līguma sadalīja Vāciju divās daļās, Vācija nepretendējusi, bet Polija, atsakoties tās pildīt, bijusi ``galvenais konfliktu izraisītājs''. Publicists D.~Žukovs par Vācijas iebrukuma pamatu uzskata ``poļu nepārtraukto ņirgāšanos pār vācu mazākumu, kā arī poļu nevēlēšanos noregulēt jautājumu par izsenis vācu pilsētu Dancigu.'' Viņš arī raksta, ka J.~Beks 1939.~gada 6.~aprīlī, parakstot savstarpējās palīdzības līgumu ar Lielbritāniju, saistīja Poliju ar pēdējo un ``noraidīja pilnībā pieņemamos Vācijas priekšlikumus par Dancigas jautājuma noregulēšanu, reizē ar to atsakoties no draudzīgu attiecību nodibināšanas ar reihu''. Arī krievu vēsturnieks A.~Širokorads uzskata, ka Polijas valdības vaina Otrā pasaules kara izraisīšanā nav mazāka kā Vācijas, Itālijas un Japānas valdību vaina. Šim uzskatam var atrast to apstiprinošus argumentus (Polijas agresīvā politika no valsts tapšanas brīža, starptautiskas drošības sistēmas veidošanas traucēšana tiešā kara priekšvakarā), taču, pēc grāmatas autora domām, pilnībā atzīt tā pamatotību nav iespējams. Par ``galveno konfliktu izraisītāju'' neapšaubāmi jānosauc Ā.~Hitlers, citu valstu līderiem piederēja tikai izpalīgu loma. Polijas valdības rīcībā esošie resursi bija pārāk vāji, lai tā varētu realizēt tikpat agresīvu politiku kā nosauktās trīs valstis. Polija vienkārši nespēja līdzināties tām. Tās vainu pielīdzināt minētajām trijām agresorvalstīm nozīmētu grēkot pret patiesību.

Ja par Versaļas līguma taisnīgumu Vācijas-Polijas robežas nospraušanā (kur galvenā loma bija Rietumu lielvalstīm, nevis Polijai) var arī diskutēt, tad 1939.~gadā demonstrēto Polijas nepiekāpību, pilnībā atsakoties izpildīt Vācijas prasības, var loģiski pamatot.

Poļu vēsturnieki parasti uzsver poļu valstsvīru patriotismu, nacionālo lepnumu. Tā, Nacionālā piemiņas institūta prezidents J.~Kurtika ir citējis J.~Beka uzstāšanos Seimā 1939.~gada 5.~maijā: ``Cilvēku, tautu un valstu dzīvē ir tikai viena lieta, kas ne par kādu naudu nav nopērkama, un tas ir gods.''

Padomju vēsturnieki Polijas valdības nepiekāpību Dancigas un ``poļu koridora'' jautājumā skaidroja ar tās bailēm, ka piekāpšanās gadījumā varētu sekot tāds tautas sašutuma uzvirdums, kas aizmēztu ``sanācijas'' režīmu. Versija ir loģiska, iespējams, ka darbojās arī tādi apsvērumi, taču pierādīt to nevar, vēl jo vairāk tāpēc, ka Otrā pasaules kara priekšvakarā pašu poļu sabiedrībā pieauga patriotisks noskaņojums, iekšējās nesaskaņas kļuva it kā mazāk aktuālas. (Par to liecināja salīdzinoši augstā~--- 67,3\% vēlētāju dalība 1938.~gada Seima vēlēšanās, kuras aicināja boikotēt gan PPS, gan zemnieku partija (``\pltxti{Stronnictwo Ludowe}''), gan nacionāldemokrāti. 1939.~gadā ar ``sanācijas`` politiķiem sadarbojās gan ``ļudoviešu'' vadonis V.~Vitoss, gan opozīcijā esošais ģenerālis L.~Žeļigovskis.) Valdošais režīms to apzinājās un arī izmantoja. Bez tam padomju vēsturnieki rakstīja, ka Polijas valdošās aprindas no iespējamās tuvināšanās ar PSRS baidījās arī tāpēc, ka tā varētu pastiprināt Padomju Ukrainas un Padomju Baltkrievijas pievilcību un ietekmi Polijas austrumu nomaļu tautu acīs. Pēc autora viedokļa, ja šādu apsvērumu (attiecībā uz nemantīgajiem ukraiņiem un baltkrieviem tam bija savs pamats) poļu politiķi tiešām ņēma vērā, tad tas tomēr nebija viens no svarīgākajiem. Ar savā valstī dzīvojošajām apspiestajiem nacionālajiem mazākumiem valdošie poļu politiķi bija pieraduši nerēķināties.

Jau daudz pārliecinošāks ir amerikāņu vēsturnieka G.~Veinberga viedoklis, ka Polijas augšslāņi bija gatavi meklēt kompromisus pat teritoriālos jautājumos, taču nekādā gadījumā nevēlējās pieslieties Antikominternes paktam, ar to simboliski pakļauties Vācijai, jo tas nozīmētu zaudēt suverenitāti.

Pamatots ir arī jau minētā angļu vēsturnieka A.~Teilora vērtējums, ka Polijas politiķi bija mācījušies no ``Minhenes'' priekšvēstures, kad Ā.~Hitlers arvien ``uzskrūvēja'' savas prasības, līdz panāca sev vēlamo. Tāpēc arī poļi nevēlējās piekāpties Vācijai ne par collu, jo bija guvuši pieredzi no Čehoslovākijas likteņa: ir tikai viens veids kā nepiekāpties pārāk daudz~--- nepiekāpties nemaz.

Tādejādi Polijai, tās vadībai nekādi nevar pārmest uzstāšanos pret Vācijas ``pamatotajām'' prasībām un ``konflikta izraisīšanu''. Cits jautājums, cik efektīvi Polijas valdība aizstāvēja savas valsts intereses. Lielā mērā Polijas valdītāju nostāju noteica viņu vidū izplatītās lielvalstiskās ambīcijas, kaut par lielvalsti 1939.~gadā ārzemju valdības Poliju neuzskatīja.

Galvenā bija un palika problēma par Polijas gatavību sadarboties ne tikai ar Rietumiem, bet arī ar PSRS. Polijas valdība uzskatīja, ka PSRS vajadzētu palikt Polijas ``sabiedroto sabiedrotā'' pozīcijā. Tas parādījās jau aprīlī J.~Beka sarunās Londonā. No Polijas viedokļa bija pietiekami, lai kara gadījumā Polijai domāti kara materiāli plūstu caur PSRS, un ar to tiktu noslēgta vienošanās par ekonomisku sadarbību. Iespējams, Polijas valdība nevēlējās PSRS dalību piedāvātajā antihitleriskajā koalīcijā arī tāpēc, ka tad strauji samazinātos Polijas loma Eiropā. Zibenskara iespēju Polijas vadība pat neņēma vērā.

Jau 4.~aprīlī uz Polijas sūtņa Maskavā V.~Gžibovska apgalvojumu par to, ka tad, kad radīsies vajadzība, Polija griezīsies pēc palīdzības pie PSRS, PSRS ārlietu tautas komisārs M.~Ļitvinovs atbildēja, ka tā pēc palīdzības var griezties arī tad, kad būs jau par vēlu. 10.~maijā savas vizītes laikā Varšavā PSRS ārlietu tautas komisāra pirmais vietnieks V.~Potjomkins izteica Polijai PSRS valdības priekšlikumu sadarboties pret nacistu agresiju.

1939.~gada 11.~maijā Polijas valdība ar sava sūtņa Maskavā V.~Gžibovska starpniecību atbildēja, ka Polija iebilst pret trīspusējām garantijām tai un ``neuzskata par iespējamu noslēgt savstarpējās palīdzības paktu ar PSRS\dots{}'' Pēc pāris dienām Polijas sūtnim Londonā E.~Račinskim, kurš ziņoja, ka britu sabiedrība arvien vairāk izdara spiedienu uz N.~Čemberlena valdību, prasot savienību ar PSRS, J.~Beks atbildēja: ``Kas attiecas uz mūsu dalību šāda veida līgumos, tad mēs tāpat kā iepriekš pieturamies pie uzskata, ka līgums par savstarpēju poļu-padomju palīdzību paātrinātu konfliktu.'' Polijas ārlietu ministram varēja izrādīties taisnība, ja tiktu noslēgts izolēts poļu~--- padomju savstarpējās palīdzības līgums, taču tur jau tā lieta, ka sarunas notika par Lielbritāniju, Franciju, PSRS, Poliju, Rumāniju un Baltijas valstis aptverošu kolektīvās drošības sistēmas radīšanu, kura varētu apturēt agresoru. 25.~maijā padomju sūtnis N.~Šaronovs sarunā ar J.~Beku atkal paziņoja: ``Mēs, protams, esam gatavi palīdzēt, taču lai palīdzētu rīt, jābūt gataviem šodien, t.i. iepriekš jāzina par vajadzību palīdzēt.'' (Pēc būtības padomju sūtnis atkārtoja jau 4.~aprīlī izskanējušo M.~Ļitvinova brīdinājumu.) J.~Beks šo paziņojumu atstāja bez atbildes, bet 12.~augustā pat lielmanīgi paziņoja, ka ``Padomju Savienības palīdzība Vācijas uzbrukumu atvairīšanai Polijai nav vajadzīga''.

Rietumvalstu un PSRS sadarbības uzsākšanu 1939.~gadā vasarā ievērojami traucēja jautājums par netiešu agresiju, tās definīciju un tās pielietošanu praksē. Vārdi ``netieša agresija'' bija ņemti no Lielbritānijas garantijām Polijai. Ar tiem tika saprasts tas, ko Ā.~Hitlers bija paveicis ar Čehoslovākiju. Viņš tieši neuzbruka šai valstij, bet ar Lielbritānijas, Francijas un Itālijas palīdzību piespieda to kapitulēt un izprovocēja Slovākijas atdalīšanos. Taču 9.~jūlijā V.~Molotovs deva savu ``netiešas agresijas'' izpratni. Tā būtu situācija, kad upuris ``ar citas valsts spēka pielietošanas draudiem vai bez tiem piekrīt'' veikt darbību, ``kuras sekas ir šīs valsts teritorijas un spēku izmantošana pret to vai pret vienu no līgumslēdzējām pusēm''. Šāda ``netiešas agresijas'' izpratne izraisīja Rietumu partneros, kā arī tajās valstīs, kuru drošībai bija jātiek ``garantētai'', nepatīkamas pārdomas par PSRS nolūkiem. Iznāca, ka PSRS var iebrukt savu kaimiņvalstu teritorijā tiem brīvprātīgi (``bez draudiem'') tuvinoties Vācijai un pie tam vēl saņemt Lielbritānijas un Francijas palīdzību. Runa pirmkārt bija par Baltijas valstīm, kuras padomju vadība uzlūkoja (blakus Rumānijai) par vienu no iespējamā Vācijas uzbrukuma PSRS placdarmiem. Turpretī Rietumu valstis nemaz nevēlējās karot ar Ā.~Hitleru, ja tas sāktu cīņu ar J.~Staļinu par Baltiju.

Rietumvalstu viedokli izteica Lielbritānijas ārlietu ministrs E.~Halifakss: ``\dots{}veicinot Krievijas iejaukšanos citu valstu lietās, mēs varam nodarīt neaprēķināmus zaudējumus savām interesēm, kā mājās, tā arī visā pasaulē.'' Briti baidījās piešķirt PSRS licenci apdraudēt kaimiņvalstis un izplatīt komunismu pasaulē. Padomju vadība savukārt baidījās no Vācijas agresijas kaimiņvalstīs, vai ar šo valstu piekrišanu, vai arī bez tās, tāpēc vēlējās nodrošināties arī pret tādu agresiju, kuru agresors varēja iztēlot pat par palīdzību upurim.

Jūlija sākumā Francijas sūtnis Maskavā P.~Nadžiars ierosināja atrisināt jautājumu par garantijām Baltijas valstīm un to pieminēšanu slepenā protokolā pie Rietumvalstu un PSRS līguma. Kā atzīmējusi I.~Fleišhauere, ``tas bija pirmais nopietnais priekšlikums par slepena protokola iekļaušanu līgumā. Tas nāca no Francijas puses.'' Sākotnēji padomju puse nepieņēma priekšlikumu par slepenu protokolu, jo padomju varas sākotnē vēl V.~Leņins bija noraidījis slepeno diplomātiju kā padomju valstij nepieņemamu. Taču nedaudz vēlāk tas tika realizēts noslēdzot Vācijas-PSRS 1939.~gada 23.~augusta līgumu. Pragmatisms izrādījās svarīgāks par ideoloģiskajiem principiem.

Nespējot atrisināt jautājumu par ``netiešu agresiju'', abas puses piekrita koncentrēt uzmanību nevis uz politiskajiem, bet militārajiem vienošanās aspektiem. V.~Molotovs šai sakarā izteica viedokli, ka ``netiešas agresijas definējums atradīsies''.

12.~augustā Maskavā sākās triju pušu: Francijas, Lielbritānijas un PSRS militāro misiju sarunas (angļu: \entxti{triple alliance negotiations}), par kuru gaitu vēsturiskajā literatūrā ir bagātīgs materiāls. Tāpēc šeit tikai par Polijas lomu šajās sarunās. Tieši tās pozīcijas dēļ sarunas nonāca strupceļā. Sarežģīts bija jautājums par Sarkanās armijas izlaišanu caur Polijas u.c. kaimiņvalstu teritoriju, lai tā varētu pildīt iespējamās PSRS saistības un stāties pretī agresoram.

Padomju valdība šo jautājumu izvirzīja jau 1937., kā arī 1938.~gadā sakarā ar iespējamo PSRS palīdzību Čehoslovākijai. Tagad tas tika ierosināts vēlreiz. Jautājums Padomju Savienības vadībā bija rūpīgi izstrādāts. Visos piecos Vācijas iespējamās agresijas gadījumos: pret Franciju, Lielbritāniju, Poliju, Turciju, un PSRS (izmantojot Somijas un Baltijas valstu teritoriju), Sarkanās armijas izlaišana caur Poliju un Rumāniju tika uzskatīta par obligātu. Tāpat kā Lielbritānijas un ASV sauszemes spēki Pirmajā pasaules karā nevarētu sadarboties ar Francijas armiju, ja tie netiktu ielaisti Francijas teritorijā, tā arī Sarkanā armija nevarētu sadarboties ar Lielbritānijas, Francijas un Polijas bruņotajiem spēkiem, ja tā netiktu ielaista Polijas teritorijā. PSRS vadība deklarēja, ka Polijas aizliegums Sarkanajai armijai iziet caur savu teritoriju neļauj radīt efektīvu Austrumu fronti pret Vāciju. Ja Austrumu fronte tiktu sagrauta, tas ļautu Vācijai savus spēkus virzīt vai nu uz rietumiem, vai tālāk uz austrumiem. Tas bija ne tik daudz Polijas vai Rumānijas drošības jautājums, cik Francijas un Padomju Savienības drošības jautājums. Padomju valdība negatavojās uzstāties pret Vāciju bez precīzām militārām garantijām. Trīspusējās sarunās PSRS delegācija norādīja, ka Anglijai un Francijai ir jāiegūst Polijas apņemšanās kara gadījumā izlaist padomju bruņotos spēkus caur savu teritoriju.

Rietumvalstu nostāja bija citāda. Jau sākoties trīspusējām sarunām, Francijas delegācijas vadītājs Ž.~Dumenks norādīja, ka Polijai, Rumānijai un Turcijai pašām jāizstāv sava teritorija, bet trijām lielvalstīm jādodas tām palīgā, ``kad tās to lūgs''.

Padomju vēsturnieks A.~Manusevičs šai sakarā bez pamatojuma rakstīja, ka Rietumvalstis vēlējušās, lai Vācijas uzbrukuma gadījumā Polijai PSRS pieteiktu karu Vācijai, bet neveiktu nekādas militāras darbības līdz vācu karaspēka iziešanai pie padomju robežas. Ar to Rietumvalstis vēlējušās nostādīt PSRS Vācijas trieciena priekšā, ziedot Vācijai Poliju, lai tikai ierautu PSRS karā Eiropā ar Vāciju, bet Āzijā ar Japānu. Faktiski tas bija Krievijas Pilsoņu karā izplatīto uzskatu nekritisks atkārtojums, ka imperiālisti vēlas visus, ko vien izdotos, ``uzrīdīt'' padomju valstij.

Uz to padomju delegācijas vadītājs K.~Vorošilovs atzīmēja, ka Polija un Rumānija var nevis savlaicīgi lūgt palīdzību, bet kapitulēt, toties Lielbritānijas, Francijas un Padomju Savienības interesēs nav pieļaut ``lai papildus Polijas un Rumānijas bruņotie spēki tiktu iznīcināti'', tāpēc ir jau iepriekš jāvienojas par padomju karaspēka piedalīšanos šo valstu aizsardzībā no agresijas. K.~Vorošilovs precizēja, ka runa iet par padomju karaspēka iziešanu tikai caur noteiktiem apgabaliem: t.s. Viļņas koridoru (domāts Polijai toreiz piederošais Viļņas apgabals) ziemeļos un pēc iespējas, caur Lietuvu uz Prūsijas robežām, kā arī, ja radīsies nepieciešamība, arī caur Galīciju dienvidos. Bija paredzēts ar padomju karaspēku nobloķēt Vācijas iespēju iebrukt Baltijā un dienvidos ``nogriezt'' tai pieeju Ukrainai un Rumānijas naftai.

Francijas Kara ministrija uzziņā par sarunu gaitu atzina, ka padomju delegācija stingri norobežo Sarkanās armijas virzības zonas poļu teritorijā un vadās ``vienīgi no stratēģiskiem apsvērumiem,'' turpretī Polijas ārlietu ministrs J.~Beks un Polijas armijas Galvenā štāba priekšnieks ģenerālis V.~Stahevičs izrādot ``nesamierināmu naidīgumu''. Ģenerāļa Ž.~Dumenka rakstītajā ``Dienasgrāmatā par franču militārās misijas uzturēšanos Maskavā 1939.g.'' atreferēti padomju argumenti par to, ka Lielbritānijai un Francijai ir savienības līgumi ar Poliju, ka tieši pret Poliju, Rumāniju, Franciju un Lielbritāniju pastāv militāra uzbrukuma draudi un tāpēc ir pilnīgi dabīgi, ka Lielbritānijas un Francijas valdībām ir jāatrisina jautājums par Sarkanās armijas iziešanu caur Polijas un Rumānijas teritorijā ar šo valstu valdībām. Sekoja Ž.~Dumenka vērtējums situācijai: ``Padomju atbilde bija ārkārtīgi skaidra; un, mums par nelaimi, tās loģika bija nepielūdzama. Bija pašapmāns vest sarunas ar Padomju Savienību, neatrisinot iepriekš, vismaz stratēģiskajā plānā, jautājumu par krievu-poļu sadarbību''.

Polijas valdībai Sarkanās armijas ieiešana Polijas teritorijā, lai arī ierobežotās zonās, bija pilnīgi nepieņemama. Pēc Otrā pasaules kara poļu vēsturnieks Z.~Zaluskis rakstīja, ka E.~Ridzs-Smiglijs un J.~Beks bijuši pārliecināti, ka Polijas un Padomju Savienības cīņā pret nacismu sadarbības rezultātā Polijas tautas masās norisināsies sociālā radikalizācija un to politiskā aktivizācija, kas ar laiku novedīs pie buržuāzijas varas krišanas. Tā kā neviena šķira nevar tiekties uz pašnāvību, Polijas ``sanācijas'' režīma vadītāji arī noraidījuši PSRS palīdzību, taču ar to viņi arī veikuši Polijas neatkarības politisku slepkavību, nolemjot to militārai sakāvei un okupācijai. Minētais vēsturnieks Polijas vadītāju nostāju iztēlojis kā tīri šķirisku, kā bailes no savas tautas. Pēc autora viedokļa, vērtējums neatbilst patiesībai. Iepriekšējā poļu-padomju kara pieredze taču rādīja, ka poļu valdošajiem slāņiem bija vairāk jābaidās no ārējām, nevis iekšējām briesmām. Arī tagad, kad brieda karš ar Vāciju, nebija pamata domāt, ka poļu tautas nostāja būs citāda kā poļu-padomju kara laikā. Vairāk kā no savas tautas radikalizēšanās bija jābaidās no PSRS palīdzības, jārūpējas, lai no tās neciestu Polijas nacionālās intereses. Taču acīmredzot vēl vairāk bija jābaidās no tiešajiem Vācijas draudiem.

Polijas vadība bažījās, ka pēc kara PSRS varētu atteikties atstāt Poliju, miera konferencē pieprasīt ``savu daļu'' uz tās teritoriju rēķina. ``Sarkanie'' taču ieņemtu Lietuvai atņemto un tolaik Polijai piederošo Viļnas apgabalu, kas ļautu tiem iesaistīt miera sarunās Lietuvu, kā arī Ukrainai atņemto Galīciju, kas varētu atdzīvināt ``ukraiņu jautājumu''. Faktiski Polija, kura 1919.--1920.~gadā pati ar spēku bija sagrābusi daļu Ukrainas un Baltkrievijas, ne bez pamata baidījās, ka Padomes kādu dienu tāpat ar spēku vēlēsies tās atgūt. Polijas valdība, vadoties no 1920.~gada poļu-padomju kara pieredzes, pamatoti bažījās, ka reiz iegājusi Polijā (vēl jo vairāk teritorijās, kuras Polija savulaik bija pati sagrābusi ar spēku), Sarkanā armija to neatstās. Maršals E.~Ridzs-Smiglijs norādīja, ka 1920.~gadā ``Sarkanajai armijai līdzi nāca visa administrācija, politiskais, propagandiskais un cits aparāts \citespace{} ieiešana tūlīt novestu pie daļas valsts okupācijas un mūsu pilnīgas atkarības no Padomēm \citespace{} Krievija vēlētos iegūt kompensāciju par 1920.~gadu un Polija zaudētu austrumu kresus.'' Pēc maršala domām ``padomju vadība, izvirzot militāros jautājumus, grib kontrabandas ceļā ievilkt sarunās politiskos jautājumus un tā vest sarunas, lai tās ievilktu vai izjauktu. Padomēm nav nodoma karot ar Vāciju un iekļauties uzbrukumā tai.''

Maršalam E.~Ridzs-Smiglijam pieder arī vārdi: ``Ar vāciešiem mums draud briesmas zaudēt mūsu brīvību. Ar krieviem mēs zaudēsim savu dvēseli''. Kā uzskata krievu publicists S.~Drožžins, maršals kļūdījās. Ja ne padomju karavīri, kuri atbrīvoja Poliju no nacistiem, tad visticamāk poļu nācija pēc neilga laika vispār pārstātu pastāvēt. Autoram jāatzīst, ka minētā varbūtība bija pilnīgi reāla. Ja pēc ebreju iznīcināšanas nacisti būtu ķērušies arī pie citu tiem netīkamu tautu likvidācijas, tad nu poļi iznīcībai būtu nolemti vieni no pirmajiem.

Polijas nevēlēšanās ielaist Sarkano armiju savā teritorijā cēloņus īsi un izteiksmīgi ir raksturojis Polijas sūtnis Parīzē J.~Lukasevičs. Kad 1939.~gada 18.~augustā Francijas ārlietu ministrs Ž.~Bonne jautāja sūtnim, kādi Polijai ir šīs nevēlēšanās cēloņi, sūtnis atbildēja: ``Beks nekad neatļaus krievu karspēkam ieņemt tās teritorijas, kuras mēs viņiem atņēmām 1921.~gadā. Vai jūs, franči, ielaistu vāciešus Elzasā-Lotaringijā?'' Tādejādi faktiski runa gāja nevis par poļu teritorijas sargāšanu no padomju draudiem, bet par to, lai nekādā gadījumā neatdotu 1920.~gadā Polijas iegūto laupījumu. (Starp citu, bijušais Francijas sūtnis Varšavā L.~Noels 1984.~gadā publicētā rakstā norādīja, ka uz Ž.~Bonne vēlējumu, lai Polija tomēr atrastu sadarbības ceļus ar PSRS, J.~Lukasevičs arī atbildējis, ka tas neesot iespējams, pirms Krievija nav atdevusi teritorijas, kuras ietilpa Žečpospolitā pirms tās pirmās dalīšanas 1772.~gadā. Kad Ž.~Bonnē iebilda, ka sadursmes briesmas ar Vāciju izsauc ``padomju palīdzības nepieciešamību'', J.~Lukaševičs uz to paziņoja, ka ``kara pirmajās dienās nevis vācieši, bet poļi ielauzīsies Vācijas iekšienē.'') Polijas sūtnis Parīzē neslēpa savu apmierinājumu par trīspusējo sarunu Maskavā buksēšanu. Tai pat 18.~augustā telegrammā J.~Bekam viņš ziņoja, ka Polijas nepiekāpīgās pozīcijas dēļ Maskavā ``notiek formālas sēdes, bezsaturīgas un nebūtiskas''.

Kā par poļiem rakstīja amerikāņu publicists H.~Bolduins, kurš kara gados strādāja avīzē ``\entxti{The New York Times}'',: ``Viņi bija pārāk lepni un pārāk pašpārliecināti, dzīvojot pagātnē. Daudzi poļu karavīri, savas tautas kareivīgā gara un tradicionālā naida pret vāciešiem piesātināti, runāja un sapņoja par ``maršu uz Berlīni''. Šīs poļu cerības labi atspoguļo vienas no karavīru dziesmām vārdi ``\dots{}tērpti tēraudā un bruņās, Ridzs-Smiglija vadīti, mēs maršēsim uz Reinu\dots{}''. (Pats maršals E.~Ridzs- Smiglijs 1939.~gada vasaras vidū nonāca tik tālu kareivīguma demonstrēšanā, ka paziņoja: ``Polija vēlas karu ar Vāciju un Vācija nevarēs izvairīties no tā, pat ja tā to vēlētos.``) Tāpat 20.~augustā Polijas Galvenā štāba priekšnieks ģenerālis V.~Stahevičs paziņoja angļu militārajam atašejam, ka ``nevar būt ir runas par to, lai izlaistu padomju karaspēku pāri Polijas robežai''. Tai pat 20.~augustā J.~Beks informēja Polijas sūtni Francijā J.~Lukaševiču par savu atbildi Lielbritānijas un Francijas sūtņiem, kuri atbalstīja PSRS prasību par tās armijas izlaišanu caur Polijas teritoriju. J.~Beka atbilde skanēja, ka ir ``nepieļaujami, lai šīs valstis apspriestu jautājumu par citas suverēnas valsts teritorijas militāru izmantošanu. Poliju ar Padomēm nesaista nekādi militāri līgumi, un Polijas valdība slēgt tādu līgumu negatavojas.'' Lielbritānijas sūtnis Varšavā H.~Kennards ziņoja par J.~Beka teikto: ``Polijas politikas princips vienmēr ir bijis nepieļaut sveša karaspēka atrašanos poļu teritorijā, un Polijas valdība iebilst pret krievu, tāpat kā vācu, karaspēka aiziešanu caur poļu teritoriju. Viņš [J.~Beks] uzskata, ka maršals Vorošilovs tagad mierīgā ceļā grib iegūt to, ko viņš gribēja iegūt ar ieroču spēku 1920.~gadā''. Tas tātad bija viens Polijas nepiekāpības cēlonis. Otrs, kā to parādīja J.~Lukaševiča pašapzinīgais paziņojums, bija savas valsts spēku (un arī Rietumu sabiedroto palīdzības) pārvērtēšana, Vācijas sasniegumu nenovērtēšana.

Jau no mūsdienām raugoties uz minētajiem vēsturiskajiem notikumiem, jāsaka, ka Baltijas valstu pieredze 1939.--1940.~gadā, kad tās bija spiestas uzņemt savā teritorijā PSRS sniegto ``danajiešu dāvanu'' [no latīņu teiciena ``\latxti{Timeo danaos et dona ferentes}'' (Baidos grieķu, pat ja tie nes dāvanas)]~--- padomju karaspēka kontingentu~--- pilnā mērā apstiprināja poļu bažas. Taču ne velti arī 1939.~gadā Polijas atteikšanos sadarboties ar Padomju Savienību Francijas sūtnis Varšavā L.~Noels nosauca par ``politisku slimību''. Tai gan bija savs loģisks pamatojums, taču tobrīd tā veda pie maldīga secinājuma un attiecīgi kļūdainas rīcības. Poļu vēsturnieks V.~Materskis norāda, ka J.~Beks 1939.~gadā nonāca nacistiskās Vācijas vadītāju antikomunistiskās retorikas gūstā un pat domu nepieļāva par iespējamu vācu-padomju pret Poliju vērstu vienošanos. PSRS starptautisko stāvokli Polijas vadītāji vērtēja kā ļoti grūtu, tādu, kas nedod tai manevra iespējas. Perspektīvu, ka PSRS un Vācijas attiecības varētu normalizēties, viņi uzskatīja par nereālu, jo zināja Ā.~Hitlera negatīvo attieksmi pret Krieviju un vēl jo vairāk pret PSRS. Bez tam Berlīne taču bija piedāvājusi Polijai dalību pretpadomju blokā. Tam vajadzēja pārliecināt poļus par nacistu fīrera nemainīgo pretpadomju nostāju. Arī padomju puse centās radīt Polijas vadībai iespaidu, ka pārmaiņas padomju politikā nav gaidāmas un konflikta gadījumā ar Vāciju var rēķināties ar PSRS ``labvēlīgu neitralitāti'' un izejvielu piegādi. (Tādu nostāju taču demonstrēja PSRS ārlietu tautas komisāra pirmais vietnieks V.~Potjomkins jau minētās viņa vizītes 1939.~gada 10.~maijā laikā Varšavā.) Šie uzskati, savienojamā ar represiju nodarītā ļaunuma lomas pārvērtēšanu PSRS militāri-ekonomiskajam potenciālam, arī kalpoja par J.~Beka vadītās Polijas ārpolitikas ``slimības'' pamatu, izskaidro šīs politikas nespēju aiz viena ļaunuma~--- Sarkanās armijas ieiešanai Polijas teritorijā~--- saskatīt vēl lielāku ļaunumu~--- Polijas sakāvi karā un neatkarības zaudēšanu.

Kad Rietumvalstis piedāvāja noslēgt militāru konvenciju bez Polijas piekrišanas, to darīt atteicās jau padomju puse. Sarkanajai armijai bija jāsaņem tiesības ieiet Polijas teritorijā jau pirmajā kara dienā, nevis pēc Polijas armija sakāves. Padomju puse varēja paredzēt, ka pēc Vācijas uzbrukuma Polijai Lielbritānija un Francija gan pieteiks Vācijai karu, taču pa īstam karot nesteigsies, atstājot PSRS vienu pret vienu ar nacistisko Vāciju.

PSRS vadībai bija arī citi cēloņi būt neapmierinātai ar Rietumu lielvalstu rīcību. Piemēram, Francijas sūtnis Maskavā P.~Nadžiars 1939.~gada 2.~augustā ziņoja savas valsts ārlietu ministram Ž.~Bonnē gan par V.~Molotova neapmierinātību ar Londonas un Parīzes avīžu pārlieko informētību, kas nevar būt nejauša, gan atgādināja to, ka abu Rietumvalstu sūtņi Maskavā jau divus mēnešus iepriekš bija brīdinājuši savas valdības par sarunu publiskošanas briesmām, par abu valdību visslepenāko instrukciju saviem sūtņiem nonākšanu presē. Arī Lielbritānijas delegācijas loceklis, sūtnis Maskavā V.~Sidss ziņojumā savam ārlietu ministram rakstīja, ka viņš izrādījies ``neiespējamā'' stāvoklī V.~Molotova priekšā, kad Londonas avīzēs izrādījās publicēti Lielbritānijas priekšlikumi par slepena pielikuma pie paredzamā līguma noslēgšanu gandrīz vai tai pašā brīdī, kad viņš tos piedāvāja padomju valdībai.

Jau 7.~augustā PSRS vadība saņēma informāciju, ka starp 15.~un 20.~augustu ir jābeidzas vācu karaspēka koncentrācijai uzbrukumam Polijai. 13.~augustā padomu aģenti ziņoja uz Maskavu: ``Karagājiena pret Poliju sākumtermiņš vēl nav zināms. Tiek uzskatīts, ka Polijai iesniegs vēl priekšlikumu, kuram tā jāpārliecina par pretošanās veltīgumu''. Vajadzēja rēķināties, ka jau no 25.~augusta var sākties karš. Lielbritānija šādu informāciju saņēma dienu vai divas vēlāk. Taču radikālu izmaiņu trīspusējās sarunās nebija.

Līdz šim lasītājiem no pirmskara notikumu virknes mazāk ir zināms, kā 1939.~gada vasarā risinājās Polijas attiecības ar Dancigas brīvpilsētu. 29.~jūlijā pēdējā nosūtīja Polijas valdībai protesta notu pret poļu muitnieku rīcību Dancigā, kuri pārkāpjot savas pilnvaras, bija ``palaiduši rokas''. Polijas diplomātiskais pārstāvis Dancigā pēc saskaņošanas ar J.~Beku 4.~augustā iesniedza pilsētas senāta priekšsēdētājam un valdības vadītājam A.~Greizeram ultimātu. Tajā tika draudēts, ka, ja 5.~augustā pilsētas valdība nedos stingru solījumu, ka turpmāk tā vairs nekad neiejaukties poļu muitnieku darbībā (kuriem tika dots rīkojums, pildot dienesta pienākumus, turpmāk nēsāt līdzi ieročus), Polija pārtrauks visu ārzemju pārtikas produktu piegādes Dancigai, kura pati ražoja tikai nelielu daļu sev nepieciešamās pārtikas. Faktiski tie bija draudi, ka Polija piespiedīs badoties pilsētas iedzīvotājus. Pēc Ā.~Hitlera pieprasījuma, kurš vēl nebija gatavs karam, A.~Greizers piekāpās. 6.~augustā~--- Poļu leģionu dienā maršals E.~Ridzs-Smiglijs Krakovā teica lielu svētku runu. Viņš apliecināja, ka Polija ir gatava atbildēt par visām strīda ap Dancigu sekām. Tad nu sapulcējies pūlis kā pēc komandas sāka skandēt: ``Atdodiet mums Dancigu! Atdodiet mums Dancigu!'' (Turpretī, kā liecinājis amerikāņu žurnālists V.~Šīrers, Dancigas iedzīvotāju noskaņojums bija pavisam cits. 11.~augustā, pabijis Dancigā, viņš savā dienasgrāmatā ierakstīja: ``Cilvēki grib pievienošanu Vācijai. Taču ne par kara cenu un ne par sava galveno tirdzniecības vārtu ar Poliju cenu. Bez šīs funkcijas, kura pavājinājusies pēc tīri poliskās ostas Gdiņas celtniecības divpadsmit jūdžu attālumā, viņi ir nolemti badam, ja Vācija neiekaros Poliju. Tāpat kā visi vācieši viņi to vēlas vienalga kādiem līdzekļiem''.) Rezultātā, kad uz Dancigas robežas sadursmē tika nogalināts poļu karavīrs, naktī uz 16.~augustu Polijas karaspēks saņēma pavēli tūlīt un bez brīdinājuma šaut uz katru, kurš pāries robežu ar Dancigu.

Gadiem vestais psiholoģiskais karš pret Vāciju bija sasniedzis virsotni, liela daļa poļu tautas tāpat kā 1918.--1920.~gadā bija gatava atbalstīt arī avantūristiskus savas valdības soļus. 14.~augustā Polijas varas iestādes uzsāka vāciešu masu arestus Augšsilēzijā. Tūkstoši cilvēku spaidu kārtā tika nosūtīti uz valsts iekšieni. (Šai sakarā latviešu lasītājiem acīmredzot nevar nenākt prātā Latvijas un citu nesen anektēto padomju republiku daļas civiliedzīvotāju izsūtīšana 1941.~gada naktī no 13.~uz 14.~jūniju uz PSRS Sibīrijas vai Kazahijas ziemeļu rajoniem. Gan ``sanācijas'' režīms Polijā, gan staļiniskais režīms PSRS, gatavojoties karam, veica līdzīgus pasākumus.) Vācu iestādes un uzņēmumi Polijā pārtrauca darbību, daudzi vietējie vācieši, baidoties no vēl lielākām vajāšanām, mēģināja bēgt uz Vāciju.

Tiesa, Polijas vadība vēl joprojām cerēja vienoties ar Vāciju. Jau minētais Polijas sūtnis Berlīnē J.~Lipskis 18.~augustā izvirzīja ārlietu ministram J.~Bekam jautājumu par iespējamām sarunām ar nacistu vadītājiem. J.~Beks tās uzskatīja par lietderīgām un 20.~augustā J.~Lipskis steidzami izlidoja uz Varšavu, lai apspriestu ar J.~Beku iespējamo sarunu saturu. (To daži padomju vēsturnieki tendenciozi interpretēja kā Polijas valdītāju gatavību piekrist kam līdzīgam Vācijas sadalītās un daļēji sagrābtās Čehoslovākijas liktenim, lai kļūtu par Ā.~Hitlera nozīmētiem vietvalžiem. Visa nākamā kara gaita gan runā pretī šādam apgalvojumam.)

Polijas vadītāji gan pārvērtēja savas valsts spēkus, gan nesaprata, ka viņu Rietumu sabiedrotajiem nedrīkst uzticēties. Lielbritānijai bija ļoti maza sauszemes armija, bet Francijā ļoti izplatīts bija uzskats, ka francūžiem nav jēgas mirt par ``poļu koridoru''. Kā savā ziņojumā ārlietu ministram rakstīja Lielbritānijas sūtnis Maskavā V.~Sidss, Lielbritānijas, Francijas un PSRS sarunu laikā Polijas sūtnis Maskavā V.~Gžibovskis viņam paziņojis, ka ``to sekmīgs nobeigums praktiski nav nepieciešams Polijas drošībai''. Poļi cerēja, ka Rietumvalstis turēs savus solījumus, un tas nodrošinās uzvaru karā pret Vāciju. Lielbritānija un Francija uzturēja šo ilūziju poļu vadītājos, kaut saprata, ka nespēs palīdzēt Polijai. Tās stūrgalvība stādīja Rietumvalstis neērtā situācijā. Viņi nevarēja atteikties no saviem solījumiem, nezaudējot savu lielvalstisko prestižu. Ja viņi jau 1939.~gada vasarā pamestu Poliju likteņa ziņā, tā vai nu padotos Ā.~Hitlera prasībām vai arī neizbēgami ciestu sakāvi, kas arī vēlāk notika un ko Rietumvalstis paredzēja jau iepriekš. Acīmredzot Rietumvalstu vadītāji cerēja ar savu demonstrēto gatavību palīdzēt Polijai vienkārši iebiedēt Ā.~Hitleru un atturēt viņu no kara izraisīšanas. Taču viņi pievīlās savās cerībās un pievīla Poliju.

Padomju historiogrāfijā viennozīmīgi valdīja versija, ka vainojamas par kolektīvās aizsardzības līguma nenoslēgšanu bija Rietumvalstis, jo tās neko nedarīja, lai panāktu iespēju Sarkanajai armijai iziet caur Polijas teritoriju. Šo versiju intervijā avīzei ``\rutxti{Известия}'' (``Ziņas'') 27.~augustā deva PSRS delegācijas vadītājs K.~Vorošilovs, vēlāk, 1949.~gada februārī tā tika izvērsta Padomju Informācijas biroja ``vēsturiskajā uzziņā'' ``\rutxti{Фальсификаторы истории}'' (Vēstures falsifikatori) un atkārtota daudzos padomju vēsturnieku darbos. Tā vēl 1989.~gadā A.~Manusevičs izteica viedokli, ka, ja Lielbritānija un Francija, Polijai noraidot militāru vienošanas ar PSRS, būtu piedraudējušas Polijai ar atteikšanos no savām garantijām, tad tā būtu spiesta pieņemt Padomju Savienības palīdzību, kura sekotu automātiski, neraugoties uz to, ka Rietumvalstis nepildīja savas saistības. Šāds spriedums balstās vairāk uz pieņēmumiem, nekā reālās situācijas adekvātu vērtējumu. Nedz Rietumvalstis varēja bez prestiža un arī savas drošības vājināšanas atsaukt tikai nesen piešķirtās garantijas, nedz Polijas vadītāji bija gatavi padoties tādam spiedienam, nedz arī Padomju Savienības atbalstam Polijai bija jāseko automātiski pēc uzbrukuma tai. Par kopdarbības nosacījumiem kā reiz taču bija jāvienojas Maskavā notiekošajās sarunās un PSRS vadītāji nebija ieinteresēti palīdzēt Polijai vienatnē, bez Lielbritānijas un Francijas bruņoto spēku iesaistīšanās konfliktā. Arī latviešu izcelsmes krievu vēsturnieks V.~Sīpols norādījis, ka, Rietumu mērķis bija vienīgi ļaut turpināties sarunām. To sapratuši arī abu valstu pārstāvji sarunās. Kā 27.~augustā ministram Ž.~Bonnē atklāja sūtnis P.~Nadžiars, ģenerālis Ž.~Dumenks viņam bija skaidri pateicis, ka ``mūsu aprēķinos nekad neietilpa pozitīva [domāts~--- tieša militāra~--- V.Š.] Krievijas palīdzība Polijai.''

Mūsdienās krievu vēsturnieces I.~Jažborovska un V.~Parsadanova atzīst, ka Rietumvalstis, kaut arī izturējās pret PSRS visai atturīgi, tomēr būdamas ieinteresētas Padomju Savienības atbalstā pret iespējamo agresoru, ne tikai neatbalstīja Polijas apņemšanos neielaist savā teritorijā Sarkano armiju, bet centās pierunāt Polijas valdību piekrist, kaut ``klusējošā'' formā, tās ieiešanai valstī. Tāpat krievu vēsturnieks L.~Gibianskis norādījis, ka augusta vidū Lielbritānija un Francija beidzot parādīja lielu ieinteresētību noslēgt militāru vienošanos ar PSRS, jo neredzēja citu iespēju aizstāvēt Poliju no tuvā vācu uzbrukuma draudiem. Rietumvalstu vadītāji saprata Polijas vadītāju bažas par padomju militārās klātbūtnes Polijā sekām, taču vērtēja Vācijas draudus par bīstamākiem un uzskatīja par nepieciešamu riskēt ar padomju palīdzības pieņemšanu, cerot uz britu-franču kontroli pār Sarkanās armijas darbību Polijā. (Tiesa, padomju vēsturnieks I.~Ostoja-Ovsjanijs rakstīja, ka Francija vēlējās saņemt no Polijas vienīgi kādu paziņojumu par labo gribu, kas neuzliktu tai nekādas saistības, taču, domājams, šis apgalvojums nepietiekami novērtēja Francijas bailes no kara ar Vāciju vienai pret vienu.)

Ar Polijas vadības nostāju bija neapmierināti Rietumu diplomāti, kuri redzēja tuvojošās reālās kara briesmas. Francijas sūtnis Maskavā P.~Nadžiars 17.~augustā telegrafēja savam kolēģim Polijā L.~Noelam: ``Pats mazākais, ko mums vajadzētu panākt no poļiem, ir tas, lai viņi neieņemtu pozīcijas, kuras izprovocētu mūsu sarunu ar krieviem pārtraukšanu \citespace{}, kas uzreiz novestu pie tādām sekām, kuru nopietnību kā viņiem, tā arī mums, viņu garantiem, viņi var iedomāties''.

17.~augustā arī Lielbritānijas sūtnis Varšavā H.~Kennards saņēma savas valdības norādījumus atbalstīt sava franču kolēģa L.~Noela pūliņus panākt Polijas valdības piekrišanu padomju karaspēka iziešanai caur Polijas teritoriju. Taču 19.~augustā J.~Beks paziņoja L.~Noelam Polijas valdības galīgo negatīvo atbildi.

Londonā ar plāna ``\detxti{Weiß}'' detaļām iepazinās 18.~augustā, tūlīt arī informējot par to Polijas vadību. Taču arī tas nemainīja pēdējās nostāju. Lielbritānijas ārlietu ministrs E.~Halifakss 20.~augustā, telegrafējot savas valsts sūtnim Varšavā H.~V.~Kennardam par tuvākajā laikā iespējamo uzbrukumu Polijai, atzīmēja: ``Ja Beka k-s domā, ka viņš var novērst vai mazināt šāda uzbrukuma iespējamību vienkārši ar to, ka viņš atturēsies pieņemt palīdzību no Padomju Savienības, es domāju, ka viņš kļūdās. Bet kad viņš saka, ka, ja karš patiesi sāksies, tad stāvoklis, iespējams, būs cits, un poļu pozīcija, iespējams, mainīsies, es domāju, ka viņš nonāk pretrunā pats sev. Ja Polija spēs pieņemt padomju palīdzību, ja sāksies karš, tad nav skaidrs, kāpēc tā nevar piekrist padomāt par šādas palīdzības pieņemšanu, kad karš ir tuvu.'' E.~Halifakss arī norādīja, ka formāla angļu-poļu līguma noslēgšana, uz ko steidzināja Polijas puse, ``diezin vai spēs būt pretsvars negatīvajam efektam, patiesi katastrofālajam, kāds būs angļu-franču-padomju sarunu galīgai neveiksmei.'' Lielbritānijas ārlietu ministrs lūdza šos apsvērumus izklāstīt J.~Bekam. Tas tika izdarīts. Vēlu vakarā H.~V.~Kennards telegrafēja E.~Halifaksam J.~Beka negatīvo atbildi: ``Polijas valdība iebilst pret krievu karaspēka iziešanu caur tās teritoriju''.

Francijas premjerministrs E.~Daladjē 18.~augustā ASV sūtnim Parīzē V.~Bullitam paziņoja, ka uzskata ``atteikšanos no patiesas militāras krievu palīdzības par lielāko muļķību no poļu puses''. 21.~augustā E.~Daladjē pat nosūtīja Francijas delegācijas vadītājam ģenerālim Ž.~Dumenkam norādījumu parakstīt konvenciju, kas paredzētu Sarkanās armijas iziešanu caur t.s. Viļņas koridoru un, nepieciešamības gadījumā, arī Polijas un Rumānijas dienvidaustrumu teritorijām. 22.~augustā Ž.~Dumenks par to paziņoja K.~Vorošilovam. Taču nedz Lielbritānijas, nedz Polijas un Rumānijas piekrišanas nebija. Tādejādi Lielbritānijas un Francijas gatavība ``pēdējā brīdī'' parakstīt militāru konvenciju nedeva rezultātus.

Paralēli Rietumvalstu un PSRS sarunām diplomātiskās aktivitātes norisa arī starp Vāciju un PSRS. Jau 1939.~gada aprīļa vidū Vācijas diplomāti saņēma A.~Hitlera uzdevumu iniciēt vācu-padomju attiecību ``jaunu Rapallo etapu''. Tiesa, to toreiz vēl bija paredzēts uzsākt ``pēc poļu jautājuma atrisināšanas'', taču vēlāk plāni tika koriģēti. Mūsdienu historiogrāfijā turpinās strīds par to, kura puse Vācija vai PSRS pirmā parādīja iniciatīvu par iespējamās politiskās vienošanās noslēgšanu.

Bieži vēstures literatūrā uzsver, ka Vācijas un PSRS sarunām durvis pavērusi J.~Staļina runa VK(b)P XVIII kongresā 1939.~gada 10.~martā, kurā viņš norādīja, ka konfliktam starp Vāciju un PSRS nav nekādu redzamu cēloņu. Taču jāuzsver, ka, kaut J.~f.~Ribentrops apgalvo pretējo, pēc vairāk ticamas Vācijas sūtniecības PSRS padomnieka G.~Hilgera un Vācijas ārlietu ministrijas darbinieka K.~Šnurres liecības, Vācijas ārpolitikas vadītāji tai ilgstoši nepievērsa uzmanību, kaut Vācijas sūtniecība par šo padomju diktatora uzstāšanos nosūtīja uz Berlīni sīku ziņojumu. Acīmredzot Vācijas tai laikā gatavotās Čehoslovākijas aneksijas gaitā netika veltīta uzmanība tām iespējām, kuras politikā pavēra J.~Staļina runa. Tikai pēc M.~Ļitvinova nomaiņas 1939.~gada 3.~maijā PSRS ārlietu tautas komisāra amatā ar V.~Molotovu Ā.~Hitlers sāka interesēties par iespēju noregulēt attiecības ar PSRS. Vācu vēsturniece I.~Fleišhauere 1990.~gadā izdotā pamatīgā pētījumā ``\detxti{Der Pakt: Hitler, Stalin und die Initiative der deutschen Diplomatie 1938--1939}.'' (Pakts: Hitlers, Staļins un vācu diplomātijas iniciatīva 1938--1939.) ir izsekojusi veselai virknei pēc vācu diplomātu un vēlāk arī nacistu līderu iniciatīvas veiktu padomju diplomātu realizētās politiskās līnijas zondāžu, kuru uzdevums bija noskaidrot vai un ar kādiem nosacījumiem PSRS būtu gatava uzlabot attiecības un noslēgt vienošanos ar Vāciju. Tikai viens zondāžas mēģinājums (I.~Fleišhauere izdala 10 tādus kontaktus) notika pirms M.~Ļitvinova atcelšanas no amata, visi pārējie~--- jau pēc tās.

Padomju puse izturējās ļoti piesardzīgi, atsauca pat savu sūtni uz Maskavu, vācu diplomātiem nācās vest sarunas ar padomju pilnvaroto lietvedi Vācijā G.~Astahovu vai meklēt kontaktus ar valstsvīriem Maskavā. Tikai kad augustā Maskavā ieradās Lielbritānijas un Francijas militārā delegācija, Ā.~Hitlers nolēma ņemt savās rokās iniciatīvu lai piesegtu Vācijas aizmuguri pret PSRS.

Pēc krievu vēsturnieka Ļ.~Bezimenska domām Padomju Savienības staļiniskā ārējā politika jau, mazākais no 1939.~gada aprīļa, bet faktiski jau agrāk, bija pakļauta mērķim iegūt laiku. Un, iespējams, kā uzskata minētais vēsturnieks, J.~Staļins bija nolēmis mācīties no N.~Čemberlena. Kāpēc gan neiegūt laiku, lai labāk sagatavotos karam, ar Ā.~Hitlera ``nomierināšanu''? Ja N.~Čemberlens ieguva laiku, ziedojot britiem draudzīgo Čehoslovākiju, kāpēc gan J.~Staļinam to neizdarīt, ziedojot Hitleram Padomju Savienībai naidīgo Poliju?

30.~maijā Vācijas Ārlietu ministrija savam sūtnim Maskavā nosūtīja sekojošu instrukciju: ``Pretēji agrāk noteiktajai politikai mēs tagad esam nolēmuši uzsākt konkrētas sarunas ar Padomju Savienību.'' Kopš jūlija beigām Berlīne aktīvi zondēja iespējamās vienošanās ar PSRS noteikumus. Pakāpeniski iezīmējās abām pusēm pieņemams risinājums. 1939.~gada 29.~jūlijā Vācijas Ārlietu ministrijas valsts sekretārs E.~Veiczekers pilnvaroja Vācijas sūtni Maskavā F.~Šulenburgu nodot V.~Molotovam: ``Jebkurā Polijas jautājuma attīstības gadījumā, vai mierīgā ceļā, kā mēs to vēlamies, vai jebkurā citā ceļā, tas nozīmē ar spēka pielietošanu no mūsu puses, mēs būsim gatavi garantēt visas padomju intereses un sasniegt saprašanos ar valdību Maskavā.'' Augusta sākumā no vācu diplomātu puses izskanēja arī ideja par slepenu protokolu, pievienojamu abu valstu līgumam. 2.~augustā J.~f.~Ribentrops sarunā ar padomju pilnvaroto lietvedi G.~Astahovu ar šādu protokolu ierosināja nodalīt abu valstu intereses ``visā garumā no Melnās līdz Baltijas jūrai''. 4.~augustā (citi dati~--- 3.~augustā) jēdziens ``slepens protokols'', kurā tiktu fiksēta abu pušu vēlme uzlabot politiskās attiecības, izskanēja no Vācijas Ārlietu ministrijas Ekonomikas nodaļas Austrumeiropas referentūras vadītāja K.~Šnurres mutes. Vēlāk šo slepenā protokola ideju jau izvirzīja J.~Staļins un V.~Molotovs.

1990.~gadu sākumā bijušais PSRS ārlietu ministra V.~Molotova palīgs un tulks V.~Pavlovs, kurš piedalījās vācu-padomju sarunās, stāstīja, ka protokola iniciatīva nākusi no PSRS diktatora J.~Staļina puses. To vēlāk Nirnbergas procesā liecināja arī Vācijas ārlietu ministrs J.~f.~Ribentrops. Taču vairāku vācu diplomātu liecības pilnībā atspēko šo izteikumu. Kā uzsvērusi I.~Fleišhauere, rietumu vēsturnieku apgalvojums, ka padomju puse sarunās bijusi aktīva puse, radās ``aukstā kara'' gados.

Radušajos apstākļos Vācijas priekšlikumi padomju vadībai likās pievilcīgāki nekā Rietumvalstu nenoteiktie piedāvājumi. Arī Latvijā iznākošās avīzes ``Brīvā Zeme'' ievadrakstā bija konstatēts: ``Nav arī ne mazāko šaubu par to, ka Vācija bija gatava uz tālejošām koncesijām, lai panāktu izlīgumu ar Padomju Savienību. Nostādīta jautājuma priekšā, vai tikt ierautai karā ar Vāciju, vai palikt ārpus šī kara, Padomju Savienība, protams, izšķīrās par otru ceļu\dots{}'' Nešauboties var teikt, ka notikumu gaita bija pievedusi pie situācijas, kad sadarbības uzsākšanā bija ieinteresēti abu valstu vadītāji. Kaut padomju vadītāji bija pilni neuzticības pret nacistisko Vāciju un sākotnēji nesteidzās parādīt savu piekrišanu līgumam ar Vāciju, sarunu ar Lielbritānijas un Francijas delegācijām ievilkšanās garumā spieda pieņemt izšķirošu lēmumu.

VK(b)P CK Politbiroja lēmums par sarunu uzsākšanu ar Vāciju paralēli sarunām ar Lielbritāniju un Franciju tika pieņemts 11.~augustā. Politbirojs nolēma ``uzsākt oficiālu vāciešu izvirzīto jautājumu apspriešanu, par ko paziņot Berlīnei''.

Mūsdienās ir zināms, ka padomju izlūkdienesta darbinieks D.~Bistroļetovs atrada atslēgas Lielbritānijas, Francijas, Vācijas, Itālijas u.c. valstu ārlietu ministriju diplomātiskajiem kodiem un PSRS vadība regulāri lasīja to saraksti. Tas nenozīmē, ka J.~Staļins zināja visus Rietumvalstu noslēpumus, taču, pieņemot lēmumus, varēja balstīties uz ļoti daudziem viņam zināmiem faktiem. Piemēram, par 1939.~gada 3.~aprīļa Ā.~Hitlera lēmumu par uzbrukumu Polijai viņš uzzināja pēc 10 dienām. Īpaši vērtīgas ziņas padomju izlūkdienestam piegādāja t.s. ``Kembridžas piecīši'' K.~Filbijs, D.~Maklīns, A.~Blants, G.~Berdžess, D.~Kernkross. Viņi ziņoja arī par to, ka Lielbritānijas valdība nevēlējās panākt vienošanos par militāru savienību ar PSRS, par tiem kontaktiem, kuri pastāvēja starp Lielbritānijas un Vācijas ``pusoficiālajām'' aprindām.

Tādam slimīgi aizdomīgam politiķim kā J.~Staļins gan šie ziņojumi, gan pati angļu-franču-padomju sarunu gaita nevarēja neradīt pārliecību, ka Rietumvalstis tās ved tikai tāpēc, lai saistītu rokas padomju diplomātijai, atņemtu tai manevra iespējas starptautiskajā arēnā. No Kremļa viedokļa dabīgi bija jārodas jautājumam: bet ja nu Rietumvalstis grib izpildīt savas saistības Polijas priekšā ne ar saviem, bet Sarkanās armijas spēkiem? Kā raksta poļu pētnieks S.~Debskis: ``Daudz kas runā par labu tam, ka vismaz Francijā radās tāda veida projekti. \citespace{} iespējams, bija arī ziņas par to, ka Londonā tās [Lielbritānijas garantijas Polijai~--- V.Š.] tulko tikai kā ``demonstrāciju'' un negatavojas spert nekādus soļus, lai realizētu Polijai dotās saistības.'' Kā rakstīja V.~Čērčils: ``Minhene un daudz kas cits pārliecināja Padomju valdību, ka ne Anglija, ne Francija nekaros, pirms viņām nebūs uzbrukts, un pat tādā gadījumā no viņām būs maz jēgas. Tuvojošās vētra bija gatava teju, teju sākties. Krievijai bija jāparūpējas par sevi.'' Šim secinājumam var tikai pievienoties.

Arī savas valsts iespēju vērtējums pēc 1937.--1938.~gada represiju viļņa, iespējams, virzīja padomju diktatoru pie secinājuma, ka valsts nav gatava lielam karam, ka pat nepilnīga tuvināšanās Vācijai PSRS var dot daudz vairāk nekā sadarbība ar Rietumu valstīm. Tas viss veda pie alternatīva kursa. Jebkura līguma noslēgšana ar Vāciju uzreiz nostādītu PSRS neitralitātes pozīcijās Vācijas un Rietumvalstu sadursmes gadījumā. Pie tam tai laikā neviens nevarēja paredzēt Vācijas veiksmīgos ``zibenskarus'' pret Poliju, Norvēģiju, Franciju, Dienvidslāviju. Gluži pretēji, varēja domāt, ka Vācijas karš ar Rietumiem būs smags un ilgstošs. S.~Debskis uzskata, ka, ja Kremlī būtu informācija par Rietumvalstu iejaukšanās nenovēršamību Vācijas uzbrukuma gadījumā Polijai, līdz PSRS un Vācijas līgumam vai vismaz PSRS iebrukumam Polijā 1919.~gada 17.septembrī un kārtējai Polijas sadalei nenonāktu. Tā tas tiešām varētu būt, ja Lielbritānija un Francija būtu gatavas karot par Polijas interesēm (taču, kā mēs zinām, tā nebija). Tomēr apstākļos, kad draudēja grūti prognozējams militārs konflikts kaimiņteritorijās, par PSRS vadības prioritārajiem ārpolitiskajiem uzdevumiem kļuva 1) palikt ārpus karā ierauto valstu loka, 2) gūt maksimālu labumu no kara. Abi šie uzdevumi bija saistīti ar politiku pret Poliju. Padomju diktators bija ieinteresēts veicināt konfliktu starp Trešo reihu un Rietumvalstīm un nebija ieinteresēts saglabāt \latxti{status quo}. Versaļas sistēmas krahs deva viņam iespēju atgūt agrāk Krievijas impērijai piederējušās teritorijas, bet optimālā variantā revolucionarizēt visu Eiropu.

Augusta vidū notika vairākas sanāksmes starp V.~Molotovu un Vācijas sūtni Maskavā F.~f.~Šulenburgu. Tika apspriests ekonomisko attiecību stāvoklis, jautājums par politiskā pakta noslēgšanu un ideja par J.~Ribentropa ierašanos Maskavā.

Izšķirošā diena J.~Staļina lēmuma pieņemšanā bija 1939.~gada 19.~augusts, kad, pēc vairāku bijušo padomju darbinieku liecībām, viņš guvis galīgu pārliecību, ka Lielbritānija negatavojas noslēgt ar PSRS efektīvu savienību, bet sarunas par to britu valdībai nepieciešamas tikai lai novilcinātu laiku.

Šai sakarā nebūs lieki pieminēt, ka vēl nesen vēsturnieku vidē šad tad tika minēts arī tāds ``dokuments'' kā J.~Staļina runa 1939.~gada 19.~augustā VK(b)P CK Politbiroja sēdē, kurā viņš it kā pamatojis PSRS ieinteresētību kara izraisīšanā. Aģentūras ``\emph{Havas}'' paziņojumu par J.~Staļina runu tajā viņš pats kategoriski noliedza. 1994.~gadā tika publicēts it kā šīs runas teksts minētajā sēdē. Tas tika uztverts kā sensācija, kura gan ātri noplaka. Runā J.~Staļins jau 19.~augustā it kā sīki informēja klātesošos par ietekmes sfēru sadali, kura notika tikai dažas dienas vēlāk. Gan franču valoda, kurā dokuments sastādīts, gan tā atrašanās vieta (PSRS Valsts centrālā Sevišķā arhīva, vēlāk Krievijas valsts kara arhīva trofeju fondos) liecina par to, ka šo J.~Staļina runas tekstu ir sastādījuši aģentūras ``\emph{Havas}'' darbinieki vai viņu informatori. Pie tam viņi parādījuši 23.~augustā noslēgtā PSRS un Vācijas neuzbrukšanas līgumam pievienotā slepenā protokola satura nepārzināšanu. (Pēc šī viltojuma iznāk, ka Vācija piešķir PSRS pilnu rīcības brīvību Baltijas valstīs un neiebilst pret Besarābijas pievienošanu PSRS, gatava piešķirt PSRS kā ietekmes zonas arī Rumāniju, Bulgāriju un Ungāriju.) 2004.~gadā žurnālā ``\rutxti{Отечественная история}'' (``Tēvzemes vēsture'') №1 Krievijas ZA Slāvistikas institūta darbinieks S.~Slučs, pēc autora domām, visai pārliecinoši pierādīja, ka ``dokuments'' ir falsifikācija. Pēc šīs publikācijas atsaukšanās uz minēto runu ir mazinājušās. Piemēram, britu vēsturnieks Dž.~Robertss konstatēja, ka runas saturs acīmredzami liecina, ka tas ir viltojums. Ir gan vēl arī vēsturnieki, kuri apstrīd S.~Sluča secinājumu. Taču, kā konstatējis ungāru izcelsmes amerikāņu vēsturnieks Dž.~Lukāčs, pat par J.~f.~Ribentropa vizīti lielai daļai VK(b)P CK Politbiroja locekļu pirms 21.-22.augusta nekas netika teikts. Vienīgi iespējams, ka 19.~augustā J.~Staļins par starptautiskajiem jautājumiem neoficiāli apspriedās ar dažiem Politbiroja locekļiem.

Pēc pētnieces I.~Fleišhaueres domām, cēlonis tam, ka tieši 19.~augustā J.~Staļins piekrita pakta ar Vāciju parakstīšanai, bija kā Vācijas vadības arvien lielākā uzstājība, gan viņa informācija par britu un franču nespēju panākt Polijas piekrišanu padomju karaspēka ieiešanai tās teritorijā, gan sarežģītais stāvoklis Tālajos Austrumos, kur risinājās konflikts pie Halhingolas.

19.~augustā Vācijas un PSRS valdības apmainījās ar neuzbrukšanas līguma projektiem. Tai pat 19.~augustā J.~f.~Ribentrops Itālijas sūtnim Berlīnē B.~Attoliko arī nodeva Ā.~Hitlera atbildi Itālijas ``dučem'' (itāļu \ittxti{duce}~--- vadonis) B.~Mussolini uz iepriekš saņemto vēstījumu, ka Itālija nespēj vest Eiropas mēroga karu. Ā.~Hitlera vēstījums dod skaidras atbildes par viņa tuvākajiem nodomiem. Vācijas ``fīrers'' tajā izklāstīja savu kredo (latīņu \latxti{credo}~--- ticu, cilvēka personīgā pārliecība, pasaules uzskats):

\begin{enumerate}
\item lēmums uzbrukt Polijai ir pieņemts, tas netiks pārskatīts;

\item Polijas konflikts paliks lokāla konflikta līmenī, jo Lielbritānija un Francija neuzdrošināsies uzbrukt ``ass'' valstīm (Domātas Vācija un Itālija. Terminu ``ass'' pirmo reizi lietoja B.~Mussolini 1936.~gada 2.~novembrī, kad pēc Itālijas un Vācijas draudzības līguma parakstīšanas viņš runāja par Romas-Berlīnes asi ap kuru pārējās Eiropas valstis varēs riņķot.);

\item ja Rietumu lielvalstis tomēr sniegs militāru atbalstu Polijai, ``ass'' valstīm diezin vai radīsies vēl labāka izdevība, kā ar tām norēķināties;

\item karš, pat ja tā mērogi izplatīsies, būs īslaicīgs, jo ``ass'' valstīm ir pārsvars.
\end{enumerate}

% page 693


Acīmredzot Ā.~Hitlers blefoja, centās iedrošināt B.~Musolini, taču pēc teiktā var secināt, ka vēl pirms neuzbrukšanas pakta ar PSRS noslēgšanas Ā.~Hitlers bija pieņēmis viennozīmīgu lēmumu par uzbrukumu Polijai. Krievu vēsturnieki G.~Kumaņevs un E.~Škļars ir norādījuši, Ā.~Hitlera pozīcija 1939.~gada vasaras beigās bija jau stingri nosprausta un neatkarīgi no pakta uzbrukums Polijai nebija vairs novēršams, viņš ``ne tikai nevēlējās, bet arī nevarēja atteikties no uzbrukuma tai''. Protams, PSRS kā iespējama pretinieka izolēšana atviegloja Vācijai uzbrukumu Polijai. Taču apgalvojums, ka bez tās Vācija neuzdrošinātos uzbrukt Polijai, nav balstīts uz faktiem, bet tikai pieņēmumiem.

20.~augustā Berlīne pieņēma padomju līguma projektu ar papildu protokolu tam. Tas ietvēra faktiski visu, ko vien Berlīne varēja toreiz vēlēties: neierobežotu neitralitāti, kas pieļāva vācu uzbrukumu Polijai, kā arī tādu neuzbrukšanas saistību izpratni, kas garantēja savstarpējo tirdzniecību arī kara gadījumā un izslēdza jebkādu vienas no līgumslēdzējvalstīm kooperāciju ar citām lielvalstīm pret līguma partneri. Vācijai, tāpat kā Lielbritānijai, nebija vajadzīga PSRS kā sabiedrotais. Tās abas vēlējās, lai Padomes vienīgi nenostājas pretinieka pusē. Tiesa, par šo ``labdarību'' un vācu pusei tik nepieciešamā varianta ``atrašanu'' un piedāvājumu PSRS pieprasīja augstu cenu. Neuzbrukšanas līgums varēja stāties spēkā tikai tādā gadījumā, ja vienlaicīgi tiktu parakstīts speciāls slepens protokols par ārpolitikas jautājumiem. ``Protokols'', bija uzsvērts Maskavas neuzbrukšanas līguma projektā, ``ir pakta sastāvdaļa''. 21.~augustā J.~Staļins deva piekrišanu Vācijas ārlietu ministra J.~f.~Ribentropa vizītei Maskavā ne vēlāk ka 23.~augustā, ar to izdarot izvēli par labu vienošanās noslēgšanai ar Vāciju. Slepenā protokola teksts tika lielā steigā sastādīts un saskaņots J.~f.~Ribentropa vizītes laikā Maskavā.

21.~augustā Ā.~Hitlers piedāvāja Londonai pieņemt sarunu vešanai H.~Gēringu, bet Maskavai~--- J.~Ribentropu. Abas puses piekrita. Par Ā.~Hitlera priekšlikumu par H.~Gēringa vizīti uz Londonu Maskavā kļuva zināms dažu stundu laikā. Lielbritānijas sūtnis Vācijā N.~Hendersons 21.~augustā ziņoja Londonai: ``Veikti visi sagatavošanās darbi, lai Gērings slepeni ierastos ceturtdien 23-jā. \citespace{} Viss virzās uz to, ka norisināsies dramatisks notikums, un mēs tikai gaidām apstiprinājumu no vācu puses.'' 22.~augustā N.~Čemberlens nosūtīja Ā.~Hitleram steidzamu vēstījumu, kurā, gan paziņojot, ka viņa valdība izpildīs saistības Polijas priekšā, ja Vācija lietos pret to ieročus, norādīja, ka viņš ir gatavs vest sarunas par Vācijas un Polijas attiecību strīdīgajiem jautājumiem, kā arī plašākām starptautisko attiecību problēmām, kas skar Lielbritānijas un Vācijas intereses. Daudzi to tulkoja kā jaunu ``Minhenes'' piedāvājumu, mājienu Vācijai, ka tā varēs saņemt no Polijas visu, ko vien vēlēsies, ja atturēsies no vardarbības. Tā, ASV sūtnis Parīzē V.~Bullits uz Vašingtonu telegrafēja, ka britu diplomāti veic ``delikātu Polijas nodevības sagatavošanu, izmantojot līdzekļus, analoģiskus tiem, kādus viņi tik sekmīgi pielietoja attiecībā pret Čehoslovākiju''.

Taču, rēķinoties ar nepieciešamību pirmkārt parakstīt līgumu ar PSRS, 22.~augustā Ā.~Hitlers atcēla H.~Gēringa lidojumu, par to gan paziņojot Londonai tikai 24.~augustā. Ā.~Hitlers, jau praktiski panācis vienošanos ar J.~Staļinu, savus uzdevumus risināja ar lielāku vērienu. Viņam, pēc visa spriežot, pat nebija nekādu šaubu par kara pret Poliju nepieciešamību. Ā.~Hitlera izvēli noteica: pirmkārt, pārliecība, ka Vācijas bruņotie spēki ātri sakaus Poliju, pat ja to atbalstīs Lielbritānija un Francija (Turpretī PSRS nostāšanās Polijas pusē nozīmētu katastrofu Ā.~Hitlera plāniem); otrkārt, cerība, ka vienošanās ar PSRS ļaus lokalizēt Vācijas un Polijas karu, atturēs Lielbritāniju un Franciju no tiešas iejaukšanās konfliktā un ļaus Vācijai izturēt Rietumvalstu ekonomisko blokādi.

22.~augustā Ā.~Hitlers paziņoja bruņoto spēku komandieriem, ka Polija ir jāiznīcina, pat ja tas vedīs pie konflikta ar Lielbritāniju un Franciju. (Viņš gan atkal piebilda, ka angļi un franči, kā viņam šķiet, negatavojas karot.) Ā.~Hitlers deklarēja: ``Mans pakts [domāta 1934.~gada 26.~janvāra deklarācija~--- V.Š.] tika iecerēts tikai kā līdzeklis tam, lai iegūtu laiku''. Berlīnē uzskatīja, ka ātra uzvara pār Poliju tiks gūta arī tad, ja Lielbritānija un Francija nostāsies tās pusē. Šo valstu politika, to bruņoto spēku stāvoklis pārliecināja nacistus, ka tās neuzsāks aktīvu kara darbību, kā arī notika. Uzstājoties vērmahta vadītāju priekšā, Ā.~Hitlers paziņoja: ``Minhenē es iepazinos ar šīm nožēlojamajām būtnēm~--- Daladjē un Čemberlenu. Viņi ir pārāk lieli gļēvuļi, lai uzbruktu. Viņi neies tālāk par [Vācijas] blokādi.'' Tai pat laikā tika uzskatīts, ka Vācija nedrīkst uzbrukt Polijai, ja tai palīdzēs arī PSRS, jo tad spēku samērs izveidotos tāds, ka Vācijas karaspēks varētu ciest sakāvi. Kad 1939.~gada augustā Ā.~Hitlers redzēja, ka turpmāk gaidīt Lielbritānijas un Francijas faktisku atbalstu kā Čehoslovākijas likvidācijas gadījumā vairs nevar, viņš forsēja notikumus, vienojoties ar J.~Staļinu.

Tā kā Polijas nepiekrišana padomju karaspēka iziešanai caur tās teritoriju faktiski bloķēja militārās sarunas Maskavā un britu-franču-padomju savienība nekļuva par realitāti, J.~Staļins acīmredzot apzinājās, ka tai brīdī pastāv divas alternatīvas: vai nu padomju-vācu līgums, vai arī britu-vācu līgums. Protams, apgalvot, ka H.~Gēringam izdotos tādu noslēgt, nevarēja, taču riskēt ar šādu varbūtību J.~Staļins kā valsts vadītājs arī nedrīkstēja.

Kā kodolīgi formulējis krievu vēsturnieks A.~Šubins, 1939.~gada augustā briti savus aprēķinus pamatoja ar secinājumu, ka Ā.~Hitlers spēj vienoties ar Lielbritāniju, bet nespēj ar PSRS. Franči~--- ar secinājumu, ka J.~Staļins spēj vienoties ar Lielbritāniju un Franciju, bet ne ar Ā.~Hitleru. Ā.~Hitlera aprēķini balstījās uz secinājumu, ka Rietumi tik un tā neuzsāks karu, tāpēc viņam svarīgāk ir vienoties ar J.~Staļinu. J.~Staļina aprēķini balstījās uz pretrunām, pastāvošām starp divām ``imperiālistu'' grupām (patiesībā šai gadījumā no imperiālistiskām pozīcijām uzstājās tikai Vācija) un viņš secināja, ka vienošanās ir jānoslēdz ar to no tām, kura būs gatava vairāk piekāpties PSRS priekšā. Pēc šī vērtējuma, tūlīt gan seko vēsturnieka konstatējums, ka ``vismazāk tika ņemta vērā Polijas pozīcija. Tas aizvainoja poļu politiķus. Lielvalstis aizmirsa, ka mazā, bet lepnā nācija var paralizēt to plānus, lai ko tas arī tai maksātu''. (Tiesa, pēdējais teikums tēlo Polijas vadītājus par gauži nesaprātīgiem, kas neatbilst patiesībai. Viņi, gluži tāpat kā Rietumvalstu vadītāji, neticēja Ā.~Hitlera un J.~Staļina vienošanās iespējai.)

Tomēr, kad 22.~augustā tika paziņots par gatavojamo PSRS un Vācijas neuzbrukšanas pakta parakstīšanu, Francijas ārlietu ministrs Ž.~Bonnē uzdeva Francijas sūtnim Polijā L.~Noelam vēlreiz mēģināt panākt, lai Polija kaut klusējot atzītu tās sabiedrotās Francijas delegācijas vadītāja ģenerāļa Ž.~Dumenka tiesības parakstīt vienošanos, kura apmierinātu Padomju Savienību. Ministrs sūtnim uzdeva ``visenerģiskākā veidā uzsvērt, ka Polijai nav ne morālu, ne politisku tiesību atteikties izmēģināt šo pēdējo līdzekli glābt mieru''. Tikai 23.~augustā J.~Beks Lielbritānijas un Francijas sūtņiem arī pauda nosacītu gatavību izskatīt sadarbības iespējas ar Sarkano armiju kara gadījumā. Viņa formulējums skanēja: ``Franču un angļu štābi ir pārliecināti, ka kopējas darbības pret agresoru gadījumā sadarbība starp PSRS un Poliju noteiktos apstākļos nav izslēgta. Ņemot to vērā, štābi uzskata par nepieciešamu sastādīt kopā ar padomju štābu jebkurus plānus.'' Francijas sūtnis Varšavā L.~Noels 23.augustā telegrafēja Z.~Bonnē: ``Es nosūtu uz Maskavu sekojošu telegrammu: Polijas valdība piekrīt tam, lai ģenerālis Dumenks pateiktu sekojošo: ``Esam pārliecināti, ka kopējas darbības pret agresiju gadījumā Polijas sadarbība ar PSRS uz tehnisku noteikumu pamata, kuri ir jāsaskaņo, ir iespējama''. Francijas sūtnis Maskavā P.~Nadžiars savukārt telegrafēja ministram Ž.~Bonnē, ka ``šī piekāpšanās nāk pārāk vēlu. Bez tam tā ir nepietiekama, jo neļauj atsaukties uz pašas Polijas valdības lēmumu.'' P.~Nadžiaram bija taisnība. ``Mēs zinām,~--- raksta itāļu vēsturnieks L.~Kanfora,~--- ka krievi jutās apkrāpti tā priekšā, kā ar nolūku nekonsekventi veda sarunas angļi un francūži. Viņi atkārtoja Brest-Ļitovsku [tur 1918.~gada 3.martā Padomju Krievija noslēdza separātu mieru ar Centrālajām valstīm: Vāciju, Austroungāriju, Turciju un Bulgāriju] pavisam citos apstākļos, izvedot sevi no tuvojošamies kara.''

PSRS valdību Rietumvalstu vilcināšanās, Polijas pozīcijas ievērošana neapmierināja un tā noslēdza neuzbrukšanas līgumu ar hitlerisko Vāciju. Vēlāk, runājot 31.~augustā PSRS Augstākās padomes sesijā, PSRS Ministru padomes priekšsēdētājs un ārlietu ministrs V.~Molotovs izskaidroja Polijas pozīciju kā sekošanu Anglijas un Francijas norādījumiem, kuras nevēloties PSRS nostiprināšanos ar savstarpējās palīdzības pakta noslēgšanu. Apgalvojums nebija korekts. Polijas valdība, jau ieguvusi Lielbritānijas un Francijas garantijas, varēja realizēt savu relatīvi pastāvīgo politiku pret PSRS. Rietumvalstis to nenostādīja notikuša fakta priekšā kā savulaik Minhenē to izdarīja ar Čehoslovākiju. Lēmumu ar kategorisku atteikumu jebkādu Sarkanās armijas daļu ielaišanai Polijā pieņēma pati Polijas valdība pretēji skaidri izteiktām Rietumvalstu vēlmēm. Attiecīgi par trīspusējo sarunu izjukšanu atbildība tai arī jānes pilnā mērā. Lielbritānijas, Francijas un PSRS sarunas parādīja, ka kaut Rietumu lielvalstis nebija gatavas sadarboties ar PSRS uz līdztiesības pamatiem, pilnā mērā uzņemties saistības apturēt nacistu agresiju, tās tomēr nebija kā vienīgās vainojamas galīgā sarunu izjaukšanā. Šis ``gods'' pienākas arī Polijas diplomātijai ar J.~Beku priekšgalā. H.~Kennards pamatoti norādīja: ``Jautājums par krievu karaspēka iziešanu [caur Poliju] izrādījās tas akmens, pret kuru atduras visi priekšlikumi par kolektīvas alianses izveidi Austrumu Eiropā''. Var atzīmēt, ka to saskatīja arī politiskie novērotāji Latvijā. 1939.~gada 17.~oktobrī avīzes ``Brīvā Zeme'' oficiālajā ievadrakstā ``Vienīgi pareizais ceļš'', kura pamatmērķis bija attaisnot Latvijas realizēto ārpolitiku, bija norādīts, ka pirmskara vasaras notikumos ``\dots{}izšķiroša nebija Baltijas valstu, bet Polijas nostāja, kas noraidīja jebkādu padomju palīdzību, līdz ar to padarīja par neiespējamu angļu-franču vēlēšanos piesaistīt Padomju Savienību blokā pret Vāciju.''

Kaut Vācijas un PSRS neuzbrukšanas līguma (vācu \detxti{Deutsch-sowjetischer Nichtangriffspakt}, krievu \rutxti{Договор о ненападении между Германией и Советским Союзом)}, parasti saukta pēc tā parakstītāju vārdiem par \strong{Molotova-Ribentropa paktu}, un slepenā protokola pie tā parakstīšanas laiks bija pēc dažādām ziņām no pulksten 22-iem 23.~augustā līdz pulksten 2-iem 24.~augustā, vēsturē tas iegājis ar 23.~augusta datumu.

Daudzi no zinātnieku darbiem, kuri veltīti šim līgumam, vēl joprojām atrodas ideoloģisko pozīciju, tajā vai citā valstī vispārpieņemto uzskatu vai ``vēsturiskās modes'' ietekmē. Redzama autoru tendence notikumu gaitu ``iestrādāt'' savu politisko uzskatu sistēmā. Var atrast pilnīgi pretējus, klaji polarizētus uzskatus. PSRS un citu t.s. ``sociālistisko'' valstu vēsturnieki galvenokārt uzsvēra, ka pakta noslēgšana Padomju Savienībai bija nepieciešamība, ka to tai uzspieda Lielbritānijas, Francijas un Polijas realizētā politika. Arī mūsdienu Krievijas vēsturnieku vairākums uzskata, ka pakta noslēgšanu izsauca nepieciešamība nodrošināt valsti pret iespējamo Vācijas agresiju vai vismaz iegūt laiku aizsardzības stiprināšanai. Turpretī mūsdienu Polijas historiogrāfijā valda tendence pielīdzināt J.~Staļina un Ā.~Hitlera rīcību pret Poliju, iztēlot PSRS kā Polijai neapšaubāmi piederošas teritorijas okupētāju.

Par to, cik lielā mērā, no ilglaicīgas perspektīvas raugoties, pakta noslēgšana bija PSRS ārpolitikas un diplomātijas veiksme vai neveiksme, var strīdēties bezgalīgi. Taču ar pakta noslēgšanu PSRS pagaidām palika ārpus lielvalstu kara Eiropā, ieguva zināmu manevra brīvību starp karojošajām valstīm, kā arī iespēju vismaz daļēji apmierināt tīkojumus atjaunot cariskās Krievijas robežas.

Francijas sūtnis Maskavā P.~Nadžiars konstatēja: ``\frtxti{Apres Munich, c`est la rédu berger à la bergère}.'' (Pēc Minhenes krievi mums atmaksāja ar triecienu pret triecienu)~--- tātad pakts tika vērtēts kā Padomju Savienības atbilde Rietumvalstīm par Minheni. Jautājums nebija par to, ka J.~Staļins vairāk ticētu Ā.~Hitleram nekā Rietumvalstu līderiem. Viņš neticēja nevienam. Viņš, tāpat kā Rietumvalstis 1938.~gadā, vienkārši gribēja izvairīties no kara, upurējot citas valstis. 1938.~gadā tika upurēta Čehoslovākija, 1939.~gadā~--- Polija. Tikai Čehoslovākija bija miermīlīga valsts, pret kuru Rietumu lielvalstīm nebija nekādu pretenziju, tās pašas tai bija piešķīrušas Vācijas apstrīdētās robežas; turpretī Polija savu agresivitāti bija parādījusi kā 1919.--1920., tā arī 1938.--1939.~gadā, un PSRS to kopš nodibināšanas uzskatīja par nevēlamu kaimiņu, kurš sagrābis padomju tautu teritorijas. Krievu vēsturnieks A.~Djukovs, salīdzinot 1938.~gada un 1939.~gada situācijas, konstatējis: Čehoslovākija vienmēr bija lojāla pret Rietumvalstīm, tomēr tās to nodeva. Polija ilgus gadus realizēja atklāti nedraudzīgu politiku pret PSRS. PSRS nevēlējās riskēt, glābjot valsti, kura atbalstīja dažādus padomju nedraugus, arī pati gandrīz neslēpa vēlmi paplašināties uz PSRS rēķina. Krievu vēsturnieks un politologs tā I.~Šiškins atzīmējis, ka J.~Staļins, tāpat kā savā laikā Katrīna II, pilnībā izmantoja starptautiskajā politikā atvērušās iespējas. T.s. Molotova-Ribentropa pakts nodrošināja Vācijas neiejaukšanos notikumos bijušās Krievijas impērijas telpā un Rietumbaltkrievijas un Rietumukrainas praktiski mierīgu atkalapvienošanos ar Padomju Krieviju. Bet pēc Polijas agresijas 1918.--1920.~gadā J.~Staļinam bija vēl mazāk cēloņu stādīt Polijas intereses augstāk par baltkrievu un ukraiņu interesēm un padomju valsts drošību nekā Katrīnai Lielajai. Tajā pašā laikā tas pats pakts nodrošināja nacistiem rīcības brīvību pret Poliju un tās sadalīšanu.

No mūsdienu augstumiem skatoties uz pirmskara periodu, var jau teikt, ka J.~Staļinam vajadzēja turpināt pieturēties kolektīvās drošības politikai un dot vairāk laika Rietumu līderiem mainīt savu politiku. Taču neviens nevar apgalvot, ka pēdējie uz to bija spējīgi. Lielbritānijas un Francijas valdības jau ilgstoši tika brīdinātas par vācu-padomju vienošanās iespēju. ``Cik reizes es netiku runājis par to!~--- vēlāk atcerējās Francijas sūtnis Maskavā 1933.--1936.~gadā Š.~Alfāns,~--- Vienojieties ar PSRS ar savstarpējo palīdzību, citādi krievi vienosies ar vāciešiem.'' Tās nebija viena diplomāta domas. ``Nepieciešamība noslēgt vienošanos mums ir vairāk neatliekama nekā viņiem [krieviem],~--- 1939.~gada 20.~jūlijā rakstīja Lielbritānijas pārstāvis PSRS V.~Strengs. Pretēji PSRS mēs uzņēmāmies saistības [domāta Lielbritānijas valdības deklarācija aizstāvēt Polijas neatkarību~--- V.Š.], kuru izpildi no mums var pieprasīt jebkurā laikā.'' Šie saprātīgie brīdinājumi Rietumu lielvastīs netika ņemti vētrā. Īpaši Lielbritānijas valdība joprojām meklēja ceļus kā izvairīties gan no kara, gan līguma ar Padomju Savienību. Rezultātā 1939.~gada 25.~augustā Francijas sūtnis Maskavā (1939-1940) P.~Nadžiars savam kolēģim L.~Noelam ar rūgtumu atzina: ``Tiešām, grūti iedomāties, kā varēja cerēt panākt no PSRS, lai tā uzņemtos saistības pret Vāciju, kura tik laipni izturējās pret to, ja mūsu garantijas saņēmušie poļi un rumāņi neko negribēja dzirdēt par krievu palīdzību''.

Jau krietni vēlāk~--- 1943.~gada decembrī Teherānā V.~Čērčils paziņoja: ``Es ļoti labi sapratu Krievijas situāciju kara sākumā un, ņemot vērā mūsu vājumu kara sākumā, Minhenes politiku un to faktu, ka Francija nodeva Čehoslovākijai dotās garantijas, es sapratu, ka padomju valdība nevarēja šādos apstākļos riskēt ar savu pastāvēšanu''.

Jau memuāros analizējot 1939.~gada augusta notikumus V.~Čērčils pamatoti rakstīja: ``Nevar pateikt, kam tas [pakts] izraisīja lielāku nepatiku~--- Hitleram vai Staļinam. Abi apzinājās, ka tas varēja būt tikai pagaidu apstākļu diktēts līdzeklis. Antagonisms starp divām impērijām un sistēmām bija nāvīgs. Staļins, bez šaubām, domāja, ka Hitlers Krievijai būs mazāk bīstams ienaidnieks pēc kara gada pret rietumu valstīm. \citespace{} Par labu Padomēm var teikt, ka Padomju Savienībai bija vitāli svarīgi cik iespējams tālāk atvirzīt vācu armijas izejas pozīcijas, lai krievi iegūtu laiku un varētu savākt spēkus no visiem savas kolosālās impērijas nostūriem. Krievu prātos ar sakaitētu dzelzi iededzinātas katastrofas, kuras viņu armijas cieta 1914.~gadā, kad viņi metās uzbrukumā vāciešiem, vēl nepabeiguši mobilizāciju. Bet tagad viņu robežas bija krietni vien tālāk uz austrumiem nekā pirmajā karā. Viņiem vajadzēja ar spēku vai viltu okupēt Baltijas valstis un ievērojamu daļu Polijas, pirms viņiem uzbruks. Ja viņu politikai arī bija raksturīgs auksts aprēķins, tad tai paša laikā tā bija arī lielā mērā reālistiska.'' Tātad V.~Čērčils, kurš labi pārzināja lietas apstākļus, atzina J.~Staļina politikas apstākļu uzspiesto raksturu un pragmatismu. Kā konstatējis krievu publicists V.~Kožinovs, no tā var secināt, ka V.~Čērčils, būdams J.~Staļina vietā, rīkotos tāpat.

J.~Staļinam nācās izvēlēties vai riskēt ar ieraušanu karā ar Vāciju tūlīt~--- neizdevīgos apstākļos, bez sabiedrotajiem, vai atlikt to uz vēlāku laiku. Viņš izvēlējās atlikšanu. Tiesa, var piebilst, ka 1939.~gadā vēl nepastāvēja reālas briesmas, ka gadījumā, ja nebūtu parakstīts Vācijas un PSRS neuzbrukšanas līgums, Vācija pēc Polijas sakaušanas varētu turpināt virzību uz Austrumiem un uzbrukt PSRS. Ā.~Hitlers tad vēl nebija gatavs karam pret PSRS. Vēsturnieku pētījumi pierāda, ka ne 1939.~gada augustā, ne nākamajos mēnešos nekādu konkrētu uzbrukuma plānu pret PSRS Vācijai vēl nebija. Plānā ``\detxti{Weiß}'' nebija nekādu norādījumu par to, ko darīt vērmahtam, ieejot Polijas austrumu rajonos~--- Rietumukrainā un Rietumbaltkrievijā. Nebija nevienas norādes par iespēju karagājienu pret Poliju pārvērst par karagājienu pret PSRS. Tāpēc nevar apgalvot, ka tikai ar Molotova-Ribentropa pakta noslēgšanu PSRS izvairījās no tūlītēja kara. Domājams, ka arī J.~Staļins to apzinājās un rīkojās ne tikai centienu izvairīties no tūlītēja kara, bet arī lielvalstisku motīvu vadīts, vēloties paplašināt PSRS teritoriju un tās ietekmi. Taču iespējams arī, ka aizdomīgums par Rietumvalstu gatavību darīt visu, lai tikai iznīcinātu PSRS, bija tādā mērā pārņēmis padomju diktatoru, ka viņš pārspīlēja tūlītēja Vācijas uzbrukuma draudus Padomju Savienībai.

Jāatceras arī, ka 1939.~gada augustā pilnā sparā norisa Sarkanās armijas un Mongolijas Tautas Republikas karaspēka vienību konflikts ar Japānas armiju pie Halhingolas upes (mongoļu \emph{Chalchyn gol}), kurš bija sācies jau 1939.~gada maijā un kuru pēc apjoma varēja salīdzināt ar Vācijas-Polijas karu 1939.~gada septembrī. Sarkanā armija tajā izmantoja apmēram četras reizes vairāk tanku nekā bija visā Polijas armijā. Kopumā kaujas notika no 1939.~gada 8.~maija līdz 15.~septembrim~--- četras reizes ilgāk nekā Vācijas karš pret Poliju. Pēc padomju ``Vēsturiskās enciklopēdijas'' (``\rutxti{Советская историческая энциклопедия}'') datiem japāņi zaudēja ap 25~tūkstošus vīru, bet padomju un mongoļu armija ap 10~tūkstošiem (pēc japāņu datiem tie zaudējuši 67~tūkstošus, PSRS~--- 18,5~tūkstošus), vairāk nekā Vācija 1939.~gada septembrī zaudēja Polijā.

Minhenes vienošanās laikā nedz Lielbritānijai, nedz Francijai nedraudēja tūlītējs karš divās frontēs, taču tās bija gatavas ziedot Čehoslovākiju Ā.~Hitleram. PSRS vadībai 1939.~gada augustā vajadzēja rēķināties ar karu divās frontēs. PSRS vadītājiem 1939.~gadā bija vismaz tikpat liels, ja ne lielāks, pamats raizēties par savu drošību kā Lielbritānijas un Francijas vadītājiem 1938.~gadā.

Jāņem arī vērā, ka jau pati PSRS sarunu uzsākšana ar Vāciju nozīmēja, ka Padomju politikas iespējas tika ierobežotas. PSRS vadība no zināma brīža nostādīja sevi stāvoklī, kad atbildēt ar ``nē'' Vācijas priekšlikumiem normalizēt attiecības un noslēgt neuzbrukšanas līgumu nozīmētu parādīt naidīgumu un izaicināt Vāciju uz uzbrukumu. PSRS situācija šai aspektā bija līdzīga Polijas situācijai, kura noraidīja visus Vācijas priekšlikumus. No PSRS drošības viedokļa bija tikai divas izejas: vai nu vienošanās ar Rietumvalstīm par kolektīvo drošību, vai ar Vāciju par neuzbrukšanu. Trešā~--- palikt vienatnē~--- nedeva nekādu drošību. Reāli vērtējot Rietumvalstu nodomus un, kā vēlāk izrādījās, pārvērtējot Polijas un arī Francijas pretošanās spējas Vācijai, PSRS vadībai alternatīvas līguma noslēgšanai ar Vāciju bija samazinājušās līdz minimumam, varētu teikt~--- bija tikai teorētiskas.

Tiesa, attiecībā pret hitlerisko Vāciju J.~Staļins pieļāva to pašu kļūdu, ko Rietumvalstu līderi~--- cerēja, ka ir iespējams atlikt neizbēgamo sadursmi, līdz Padomju Savienība būs uzkrājusi pietiekami spēkus. Krievu vēsturnieks un diplomāts J.~Kvicinskis rakstīja: ``Ja Staļins arī kļūdījās, tad tikai mierīgās pauzes ilguma vērtējumā, kuru deva pakta ar Vāciju noslēgšana. Pati pakta noslēgšana bija, neapšaubāmi, tajos apstākļos pareizs un mērķtiecīgs politisks un militārs lēmums.''

Starp citu, J.~Muhins konstatējis, ka līdz 1941.~gada vasarai nacistiskās Vācijas karos ar Poliju, Franciju u.c. Rietumeiropas valstīm tik vēlamā, bet iztrūkstošā otrā (Austrumu) fronte bija nevis Ā.~Hitlera, bet ``Lielbritānijas un Polijas nopelns''. (Pēc J.~Muhina izteikuma ``trulie un nelietīgie Polijas politiķi savu muļķīgo ambīciju vārdā \citespace{} atteicās no vienīgā tai laikā drošā sabiedrotā~--- PSRS, nedeva iespēju radīt antihitlerisko savienību, izsauca Otro pasaules karu''. Šis izteikums gan jāvērtē kā, labākajā gadījumā, publicistisks pārspīlējums.) Turpretī tas, ka PSRS no kara sākuma līdz pat tā beigām karoja tikai vienā frontē, pēc J.~Muhina viedokļa, bijis J.~Staļina nopelns. Arī šis secinājums jāpapildina, uzsverot ne tikai PSRS un Rietumvalstu interešu sakritību, bet arī šo valstu tālaika vadītāju skaidro politisko gribu.

Par Molotova-Ribentropa pakta noslēgšanu var runāt arī kā par lielu Rietumvalstu ārpolitikas neveiksmi. Ne velti V.~Čērčils atzina: ``Tas fakts, ka tāda vienošanās izrādījās iespējama, liecina par angļu un franču politikas un diplomātijas vairāku gadu garumā izgāšanos.'' Jau 1941.~gadā, kad Lielbritānija reālo notikumu attīstības rezultātā atradās vienā koalīcijā ar PSRS, Lielbritānijas ārlietu ministrs 1938.--1940.~gadā E.~Halifakss par nevēlēšanos noslēgt konkrētu vienošanos ar PSRS savā dienasgrāmatā atzina, ka ``tā bija kļūda, mēs nedrīkstam atkārtot to''.

Šī darba autors piekrīt uzskatam, ka trīspusējo angļu-franču-padomju sarunu neveiksme nav meklējama kādas no pusēm ļaunprātībā vai vēlmē piemānīt sarunu partnerus. Gan Rietumvalstis, gan PSRS bija ieinteresētas nodrošināties pret agresijas briesmām. Taču abu pušu pieeja problēmas risinājumam nesakrita. Ja Rietumu valstis vēlējās apmierināties ar vispārēju, Ā.~Hitlera iebiedēšanai domāta līguma noslēgšanu, padomju vadība no savas puses uzskatīja par neattaisnojamu risku noslēgt jebkādu savienību bez stingri noteiktām visu pušu saistībām. Tā konsekventi pieturējās pie šīs prasības. Galu galā šī atšķirīgā pieeja nozīmēja, ka trīspusējās sarunas jau no paša sākuma bija nolemtas neveiksmei. Ā.~Hitleram bija ievērojamas manevra iespējas, kuras ļāva viņam izolēt Poliju un uzbrukt tai. Taču šīs iespējas viņam pavēra ne tikai viens 23.~augusta līgums, bet gan kā Rietumu līderu, gan padomju diktatora, gan arī pašas Polijas vadītāju politika.

Molotova-Ribentropa pakta noslēgšanas vēsture jau sen ir apaugusi mītiem. Šis process sākās jau 1939.~gada vasarā, kad britu un franču valdības pašas organizēja informācijas noplūdi presē, lai sagatavotu sabiedrisko domu iespējamajai sarunu ar PSRS izjukšanai un uzkrautu vainu par to PSRS.

Pirmais, kurš PSRS un Vācijas neuzbrukšanas līgumam un slepenajam protokolam pie tā piešķīra pārlieku milzīgu nozīmi bija jau minētais nacistu ārpolitikas realizētājs, Vācijas Ārlietu ministrijas valsts sekretārs E.~f.~Veiczekers, kurš 1939.~gada 14.~oktobrī rakstīja, ka bez šī līguma ``karš nebūtu izraisījies''. Vēlāk, lai noveltu visu atbildību uz PSRS, šo neapšaubāmi kļūdaino un arī melīgo secinājumu uztvēra gan par Ā.~Hitlera ``nomierināšanas politiku'', 1938.~gadā Austrijas un Čehoslovākijas, bet 1939.~gada septembrī faktiski arī Polijas tautas atstāšanu vienu agresora priekšā, savu starptautisko saistību nepildīšanu atbildīgie Rietumu politiķi, gan viņu realizēto kursu apkalpojošie vēsturnieki. Pēc britu vēsturnieka M.~Mazovera vārdiem, apgalvojums par to, ka bez pakta nebūtu kara, ir ``pārāk iekrāsots ar mūsdienu politiskajām tieksmēm, lai to uztvertu nopietni.''

Tik tiešām, 23.~augusta līgums Rietumeiropā izsauca sašutuma vētru. Kā atzīmējis krievu vēsturnieks A.~Proņins, ``katram politiski domājošam cilvēkam bija skaidrs, ka padomju-vācu pakts nozīmēja iebrukumu Polijā, t.i. Otrā pasaules kara sākumu. Tas nemazina Vācijas vainu Otrā pasaules kara izraisīšanā. Padomju Savienība pilnībā piekrita karam starp citām tautām, lai iegūtu daļu no tā augļiem''. Taču, kāpēc gan PSRS nedrīkstēja darīt to pašu, ko gadu agrāk bija darījušas Lielbritānija un Francija? Ja demokrātiskās Rietumvalstis, vairoties no ieraušanas karā, atdeva Ā.~Hitleram veselu valsti (Čehoslovākiju), kāpēc lai diktators J.~Staļins rīkotos citādāk Polijas gadījumā, ja par šo atdošanu viņš vēl ieguva pamatīgas teritorijas? Britu vēsturnieks A.~Teilors šai sakarā rakstīja: ``Rietumos tika sacelts troksnis par Padomju Krievijas noziegumu, noslēdzot vienošanos ar vadošo fašistisko lielvalsti. Grūti bija saprast britu un franču politiķus, kuri aktīvi atbalstīja Čehoslovākijas sadali un pat centās panākt jaunu vienošanos uz Polijas rēķina.''

Tiesa, jāatzīst, ka bija starpība abu pušu (pirmajā gadījumā~--- PSRS, otrajā gadījumā~--- Rietumvalstu) reakcijā pret vienu un otru Ā.~Hitlera ``nomierināšanas'' mēģinājuma izpildījumu. Pirmajā gadījumā J.~Staļins par Rietumvalstu gatavību aizstāvēt Čehoslovākiju ilūzijas neloloja, par ko liecināja solījums palīdzēt pēdējai, arī negaidot Tautu Savienības Padomes lēmumu par šādas palīdzības nepieciešamību. Otrajā gadījumā Rietumeiropas valstu politiķi cerēja, ka, lai kā arī nepavērstos notikumi, Lielbritānija varēs tos vērot no savām salām, Francija,~--- aizslēpusies aiz Mažino līnijas, bet reāli pret nacistisko Vāciju karot nāksies PSRS. Tagad šīs Rietumvalstu cerības sabruka. Tas nevarēja neizraisīt vilšanos piedzīvojušo Rietumu līderu sašutumu.

Jau minētais krievu vēsturnieks un diplomāts J.~Kvicinskis norādījis, ka ``padomju diplomātijas māksla izpaudās tajā, ka tā tajā brīdī izvairījās no konflikta ar Hitleru un ļāva karā ar to vispirms iesaistīties angļiem un frančiem (pretēji to plāniem un aprēķiniem). Tas bija izšķirošs priekšnoteikums vienotas antifašistiskas frontes izveidei pēc Vācijas uzbrukuma Padomju Savienībai.''

Krievu vēsturniece N.~Naročņickaja norādījusi, ka apzīmējumi ``kara pakts'', ``sadalīšanas pakts'' u.c., kuram neesot bijis precedenta vēsturē, nopietnam vēsturniekam var izsaukt tikai ironiju. Viņa min veselu gūzmu piemēru no Vestfāles miera (1648) līdz Deitonas līgumam (1995), kad līgumslēdzējas ar spēku dalīja teritorijas, neievērojot to tautu intereses. Pētniece konstatē, ka padomju-vācu līgums mainīja Ā.~Hitlera iecerēto uzbrukumu secību, padarīja to Rietumiem nepatīkamu, ar to ietekmējot arī pēckara valstu konfigurāciju, neļaujot anglo-sakšiem ieiet Austrumeiropā ne kara sākumā, ne pēc antihitleriskās koalīcijas uzvaras. Līdz ar to nepiepildījās Rietumu cerības neļaut PSRS ietekmei izplatīties Austrumeiropā. Tieši tāpēc Rietumu ideologi cenšas demonizēt Molotova-Ribentropa paktu.

Polijā 1939.~gada 23.~augusta PSRS un Vācijas neuzbrukšanas līgums un tā slepenais protokols ir kļuvuši par sava veida nacionālās traumas simbolu, ar to tiek saistīta Polijas ``ceturtā'' dalīšana (vēsturiski ar to tiek ``aizmirstas'' pēc XVIII gadsimtā notikušajām trim Žečpospolitas dalīšanām vēl arī Napoleona I un Vīnes kongresa veiktās Polijas teritorijas dalīšanas, kuras ``ceturto dalīšanu'' faktiski pārvērš par sesto) un vēlākā atrašanās nacistiskās Vācijas okupācijā un PSRS atkarībā. Šī trauma arī mūsdienās izpaužas arī kā sakāpināts aizdomīgums pret iespējamo kaimiņvalstu aizbildniecību pār Poliju, sava veida fobija pret tām.

Tā 2005.~gadā VFR kanclera G.~Šrēdera un Krievijas Federācijas prezidenta V.~Putina noslēgto līgumu par gāzes vada būvi zem Baltijas jūras, apejot Polijas teritoriju, tās ārlietu ministrs R.~Sikorskis 2006.~gadā salīdzināja ar ``Hitlera-Staļina paktu''. Lūk, viņa teiktais: ``Polija ir īpaši jūtīga pret koridoriem un līgumiem pāri tās galvai, \citespace{} Tā ir Molotova-Ribentropa tradīcija. Tas bija 20.gadsimtā. Mēs nevēlamies atkārtojumu.'' Pie tam diemžēl tiek aizmirsts par Polijas toreizējās valdības noraidošo attieksmi pret sadarbību ar PSRS cīņā pret agresiju, atkal parādās Polijas ambīcijas būt stāvokļa noteicējai Eiropas centrā. Arī daudzi poļu vēsturnieki nevēlas saskatīt vismaz ``skabargu'', bet patiesībā~--- ``baļķi'', savas valsts ``acī''.

Atbilstoši minētajiem mītiem, Padomju Savienības vadība ilgstoši meklēja iespēju noslēgt minēto paktu, slepus un nodevīgi vienojās ar nacistiem. Britu-franču-padomju sarunu laikā padomju puse uzstādījusi arvien jaunas prasības, lai dotu laiku Vācijai izšķirties. Otrā pasaules kara sākumu noteicis tieši Molotova~--- Ribentropa pakts. Šādā notikumu gaitas interpretācijā ēnas pusē paliek Rietumvalstu vilcināšanas, nenoteiktības taktika sarunās ar PSRS, Polijas atsacīšanās no Sarkanās armijas palīdzības, kad Polijas likteni uz lielvalstu cīņas par ietekmi pasaulē fona lielā mērā noteica poļu vairākuma mūžsenā nepatika un pat naids pret Krieviju.

Mūsdienās Rietumos 1938.~gada Minhenes vienošanās ir sākta vērtēt daudz objektīvāk nekā līdz PSRS sabrukumam. Tā, angļu vēsturnieka D.~Reinoldsa grāmatā ``\entxti{Summits: Six Meetings That Shaped the Twentieth Century}'' (Samiti: sešas tikšanās, kuras noteica XX gadsimta seju, 2007), kā šāda pasaulei svarīga tikšanās minēta tieši Minhene, bet nevis Molotova-Ribentropa pakta parakstīšana. Grāmatu, veltītu Minhenes līguma noslēgšanai arī uzrakstīts krietni vairāk, nekā veltītu vācu-padomju attiecībām, nerunājot jau par 1939.~gada 23.~augusta līgumu kā tādu. Vēstures zinātnē Minhenes vienošanās tiek viennozīmīgi vērtēta kā Rietumu politikas un diplomātijas izgāšanās, kuru tikai minimāli var attaisnot ar militāra rakstura apsvērumiem. Vēl vairāk~--- vārds ``Minhene'' ir kļuvis par sugas vārdu, kapitulācijas un politiskas sakāves sinonīmu. Minhenes vienošanās ne tikai apkaunoja Lielbritāniju un Franciju, kuras atteicās palīdzēt Čehoslovākijai, bet izrādījās arī, mazākais, stratēģiska kļūda, ja ar to tika cerēts novērst karu, bet varbūt arī tīša ļaunprātība, ja ar to bija plānots, atpērkoties no agresora~--- nacistiskās Vācijas, novirzīt tās agresiju uz austrumiem un palikt malā no konflikta. Lielbritānija un Francija tik un tā pēc gada tika ierautas karā, bet situācija tām bija jau daudz nelabvēlīgāka nekā 1938.~gadā.

Arī pēc Minhenes vienošanās Lielbritānijas un Francijas valdību ieturētā politika nebija bez vainas. Mūsdienu kanādiešu vēsturnieks M.~Dž.~Karleijs uzsvēris, ka tieši ``Nevila Čemberlena opozīcija aliansei ar PSRS bija izšķirošais faktors trīspusējo sarunu izjukšanai 1939.~gada vasarā un ienesa ievērojamu ieguldījumu Otrā pasaules kara tuvināšanā.'' M.~Dž.~Karleijs norādījis uz apgalvojuma, ka V.~Molotovs Lielbritānijas un Francijas sūtņiem izvirzījis arvien jaunas prasības, ko lieto abu minēto valstu politikas apoloģēti, neatbilstību patiesībai. Padomju politikas pamati bija stingri noteikti jau 1935.~gadā. Jautājumi par ``netiešu agresiju'', par Sarkanās armijas iziešanas tiesībām caur citu valstu teritoriju, arī par garantijām Baltijas valstīm nebija jauni un negaidīti izvirzīti.

Kanādiešu vēsturnieks arī atzīmējis, ka viens no svarīgākajiem Otrā pasaules kara cēloņiem blakus Versaļas līgumam, 1929.--1932.~gadu krīzei, Rietumvalstu domstarpībām un vājumam (šo uzskaitījumu varētu turpināt~--- V.Š.) bija Rietumeiropā valdošais antikomunisms. Tas bija pamatā Lielbritānijas, Francijas un PSRS kooperācijas neiespējamībai. Rietumu politiķi starp pasaules kariem redzēja vistiešāko kopsakaru starp karu un revolūciju. Par to viņi izteicās pirms Otrā pasaules kara, par to rakstīja pēc tā. Šī dominante gan sastapa pretestību, taču dažkārt izšķirošajos momentos uzvarēja antikomunisms. Ja Rietumvalstu politiķus~--- N.~Čemberlenu, E.~Daladjē u.c. neuzskata par muļķiem, bet tādi viņi nebija nekādā gadījumā, tad viņu politiku attiecībā pret PSRS 1939.~gadā jāvērtē nevis kā smagu kļūdu, bet gan kā smalku, aprēķinātu risku, kurš tomēr neattaisnojās. Tieši antikomunisms piespieda Rietumvalstis visu laiku meklēt kompromisus ar nacistiem, ar to mazinot savu drošību un gala rezultātā iestiegot karā. Jāpiebilst, ka antikomunisms būtiski ietekmēja arī ``sanācijas'' Polijas politiku.

Tomēr arī mūsdienās aukstā kara paliekas liek daudziem vēsturniekiem ``caur pirkstiem'' skatīties uz antikomunisma ietekmi pirmskara gados, uz to destruktīvo lomu, kura tam bija Lielbritānijas un Francijas ārpolitikā. Vēl daudzi Rietumu vēsturnieki un arī viņu mūsdienu sekotāji austrumos nekādi nevēlas atzīt, ka pie Molotova~--- Ribentropa pakta noslēgšanas noveda Minhenes vienošanās un angļu-franču-padomju sarunu nesekmīgums. Viņi aizstāv pretēju uzskatu~--- tieši padomju-vācu neuzbrukšanas pakts noteicis trīspusējo sarunu par antihitleriskās savienības izveidi neveiksmi, Ā.~Hitlers un J.~Staļins vienojušies, jo starp nacistisko un staļinisko režīmu neesot bijis starpības vai tā bijusi ļoti neliela.

Tādus uzskatus noteica pēc Otrā pasaules kara sācies aukstais karš, antikomunistiskā kursa turpinājums Rietumos. Šī aukstā kara ideoloģija radīja vēlmi noslēpt Lielbritānijas un Francijas valdību atbildību par kara tuvināšanu 1938.--1939.~gadā.

Arī mūsdienās jauno, no PSRS ietekmes atbrīvojušos Austrumeiropas valstu, tai skaitā Polijas, un dažu bijušo PSRS republiku historiogrāfijā ir vērojami centieni mazināt Minhenes vienošanās ietekmi uz tālākajiem notikumiem, Otrā pasaules kara izraisīšanu, priekšplānā izvirzot vienīgi Molotova~--- Ribentropa paktu kā visu nelaimju aizsākumu. Tā, piemēram, poļu vēsturnieks un žurnālists K.~Voucickis, rakstot par vēsturisko atmiņu Eiropā, uzsvēris, ka 1939.~gada 23.~augusta pakta vēsturiskā nozīme vēl joprojām tiekot nepietiekami novērtēta, kaut tas ``gandrīz visu Austrumeiropu ierāva Pasaules kara haosā''. Tai pašā laikā viņš ne ar vārdu nepiemin Minhenes vienošanos.

Arī Latvijas ZA loceklis, profesors Inesis~Feldmanis, sekojot daudzu Rietumu vēsturnieku koncepcijām, uzsver, ka Minhenes vienošanās un Molotova~--- Ribentropa pakts ``ir grūti salīdzināmi (lai neteiktu nesalīdzināmi) lielumi'', kur ``Minhenes vienošanās bija it kā mēģinājums modificēt pēc Pirmā pasaules kara noteikto starptautisko kārtību~--- tā saukto Versaļas sistēmu.'' Teiktajam varētu piekrist, taču tālāk profesors I.~Feldmanis raksta, ka Lielbritānija un Francija ``atšķirībā no Vācijas un Padomju Savienības netiecās pēc teritoriāliem ieguvumiem un citu tautu pakļaušanas. Šo valstu līderi pamatā nedomāja XVIII gadsimta, bet gan XX gadsimta diplomātijas kategorijās.'' I.~Feldmanis gan atzīst, ka britu un franču īstenotā diplomātija bija ``apšaubāma un neveiksmīga'', ka Lielbritānijas un Francijas vadītāji ``neriskēja atbildēt uz Vācijas izaicinājumu, bet piekrita kompromisam'', taču uzskata, ka ``strikti nosodīt viņus par to ir pagrūti, jo likme tomēr bija ļoti augsta, un iespējamā alternatīva viņiem pamatoti šķita pārāk bīstama un tobrīd neadekvāta'', ka ``bez atrunām ir pagrūti akceptēt historiogrāfijā visai izplatītu apgalvojumu, ka tieši Minhenes vienošanās bija Otrā pasaules kara prologs''.

Pēc autora domām, šie cienījamā vēsturnieka spriedumi prasa kritisku komentāru, jo līdzās neapstrīdamiem secinājumiem tie satur neattaisnojamu minēto Rietumu valstu pirmskara imperiālistiskās ārpolitikas apoloģētiku.

Pirmkārt, par diplomātijas kategorijām, kurās domāja Lielbritānijas un Francijas līderi. Minhenes tikšanās, kurā lielvalstu vadītāji mainīja Eiropas politisko karti, nejautājot ``mazo'' valstu viedokli, vairāk atbilda XVIII~gadsimta (arī citu iepriekšējo gadsimtu) politiskās domāšanas tradīcijām.

Daži piemēri. Lielbritānija kopš 1792.~gada atradās kara stāvoklī ar Franciju un britu flote sagrāba visu valstu kuģus, ja radās aizdomas, ka tās tirgojas ar Franciju. 1800.~gada beigās Dānija, Krievija, Prūsija un Zviedrija radīja savienību, kuras mērķis bija realizēt bruņotas neitralitātes politiku un aizstāvēt savas brīvas tirdzniecības tiesības. Lielbritānija uzskatīja, ka šī savienība apdraud tās intereses, iesniedza Dānijai ultimātu ar prasību no tās izstāties. Dānija to noraidīja. Tad 1801.~gada 2.~aprīlī britu eskadra admirāļu D.~Pārkera un H.~Nelsona vadībā ieradās Kopenhāgenas ostā un bez kara pasludināšanas uzbruka nesagatavotajai Dānijas flotei, pirms tā nebija apvienojusies ar Krievijas un Zviedrijas flotēm. Dāņu flote tika sakauta un Dānija bija spiesta izstāties no bruņotas neitralitātes savienības. 1807.~gadā, kad Dānija bija nonākusi Napoleona ietekmē, Lielbritānija baidījās, ka atjaunotā Dānijas flote varētu (!) apdraudēt Lielbritāniju. 1807.~gada septembrī britu flote šoreiz jau admirāļa D.~Gambjē vadībā trīs dienas apšaudīja Kopenhāgenu, ieskaitot arī dzīvojamos kvartālus, piespiežot Dāniju atdot savu floti britiem. Minētajās sadursmēs bija daudz cilvēku upuru.

Gluži tāpat Lielbritānija rīkojās 1940.~gada 8.~aprīlī, kad tās flote, lai traucētu dzelzs rūdas nogādi Vācijai, pārkāpjot starptautiskās tiesības un nerēķinoties ar Skandināvijas valstīm, mīnēja Norvēģijas teritoriālos ūdeņus. 1940.~gada jūlija sākumā, kad sakautās tās bijušās sabiedrotās Francijas flote draudēja nonākt Vācijas rokās, neraugoties uz to, ka Francijas valdība ar maršalu H.~Petēnu un Francijas flotes pavēlnieks admirālis F.~Darlāns vairakkārt bija apliecinājuši, ka neviens franču kuģis nenonāks vācu rokās, Lielbritānijas flote realizēja operāciju ``Katapulta'' (``\entxti{Catapult}''), uzbrūkot britu un franču ostās Francijā un Āfrikā stāvošajai franču flotei. Ja britu ostās operācija tika realizēta sekmīgi, upuru skaits bija mazs, tad citās ostās gāja bojā simti franču jūrnieku.

Tā, Lielbritānija rīkojās vienādi kā XIX gadsimta sākumā, tā XX gadsimta vidū: amorāli, bet pragmātiski.

Kā zināms, Lielbritānija un Francija, neraugoties uz daudzām līderu nomaiņām, vēl ilgi pēc Otrā pasaules kara principiāli neatteicās no saviem vardarbīgi sagrābtajiem teritoriālajiem, koloniālajiem ieguvumiem, pat vēl mūsdienās pēc pasaules koloniālās sistēmas neatgriezeniska sabrukuma cenšas paturēt atsevišķas agrāk sagrābtās teritorijas. Lai pieminam tikai XVIII gadsimta sākumā Spānijai atņemto un vēl tagad Lielbritānijai piederošo Gibraltāru un Francijas aizjūras teritoriju Jaunkaledoniju, kurā vēl XX~gadsimta 90.~gados uzliesmoja aborigēnu cīņas par neatkarību. Abas valstis arī šodien, izmantojot dažādas atrunas, liedzas atgriezt pēc piederības bijušajās kolonijās salaupītās kultūras vērtības. Abas valstis arī pēc Otrā pasaules kara un arī XXI~gadsimtā nereti ir ar militāru spēku iejaukušās citu, īpaši Āfrikas, valstu iekšējās lietās. Starptautiskā sabiedrība šie fakti lielākoties tiek interpretēti dažādi, kaut tie liecina, ka Lielbritānija un Francija vadītāji vēl mūsdienās nav pilnībā atmetuši XVIII un XIX gadsimta domāšanu.

Otrkārt, tieši uzvarētājvalstis Pirmajā pasaules karā, un pirmkārt jau Lielbritānija un Francija, nesa galveno atbildību par Versaļas sistēmas izveidi, kura sevī imanenti saturēja briesmas pasaules mieram. Nedz Vācija, nedz Padomju Krievija nepiedalījās šīs sistēmas radīšanā, vēl jo vairāk~--- protestēja pret to ar visiem to rīcībā esošajiem līdzekļiem.

Kā uzsvēris krievu vēsturnieks M.~Meļķuhovs, Lielbritānija un tās sabiedrotā Francija tik tiešām ar Minhenes vienošanos mēģināja, uz jauna, varētu teikt, Minhenē radīta balansa starp Eiropas lielvalstīm pamata modernizēt Versaļas sistēmu, taču darīja to ar nolūku saglabāt savu vadošo lomu pasaules politikā, bet nevis miera saglabāšanas interesēs.

Nevar arī iedomāties, ka abas lielvaras~--- Lielbritānija un Francija~--- mēģinātu 1938.~gadā ``modificēt pēc Pirmā pasaules kara noteikto starptautisko kārtību'', it kā labojot vācu nācijai nodarīto ``netaisnību'' Sudetu jautājumā, ja tās nejustos spiestas to darīt, piekāpjoties Ā.~Hitlera šantāžas priekšā. Jā, ``iespējamā alternatīva'' šo lielvalstu līderiem ``pamatoti šķita pārāk bīstama un tobrīd neadekvāta'', taču šāda situācija bija radusies, galvenokārt viņu pašu vai viņu priekšteču pie valsts stūres darbības rezultātā un to atrisināt bija viņu kā valstsvīru pienākums. Taču N.~Čemberlens un E.~Daladjē to risināja uz sava sabiedrotā~--- Čehoslovākijas suverenitātes rēķina, tā rezultātā tieši grūžot Ā.~Hitleru uz jaunām avantūrām. Kā konstatējis vācu vēsturnieks K.~Hildebrands, ar Minhenes vienošanos bija ielikti pamati Molotova-Ribentropa pakta noslēgšanai. Amerikāņu vēsturnieks un publicists V.~Šīrers pamatoti jautāja: ``Ja Čemberlens rīkojās godīgi un cildeni, nomierinot Hitleru un atdodot viņam Čehoslovākiju 1938.~gadā, tad kāpēc Staļins gan rīkojās negodīgi un zemiski, pēc gada nomierinot Hitleru ar Poliju, kura tāpat jau atteicās no padomju palīdzības?'' Citiem vārdiem: ja demokrātiskās valstis varēja rīkoties amorāli, bet pragmātiski, kāpēc lai tā nedarītu PSRS diktators J.~Staļins? Minhenes vienošanās bija plats solis uz jaunu karu. ``Ja nebūtu Minhenes, nebūtu arī pakta ar Hitleru''~--- vēlāk atzīmēja arī pats J.~Staļins sarunā ar V.~Čērčilu un tas ne vārda neiebilda. Jāpiekrīt krievu ekonomistam un politiķim I.~Jirgenam, ka ``Minhenes darījums starp Lielbritāniju, Franciju, nacistisko Vāciju un Itāliju \citespace{} bija patiesais lūzuma punkts Eiropas starptautiskās politikas sistēmā. Pēc Minhenes bija jau saprotams, ka mazai nācijai nav iespēju pretoties lielu valstu lēmumiem.''

Treškārt, Lielbritānijas un Francijas politikas aizstāvji uzsver, ka Molotova-Ribentropa pakta rezultātā PSRS guva laupījumu~--- Rietumukrainu, Rietumukrainu, vēlāk arī Baltijas valstis, bet Francija un Lielbritānija, slēdzot Minhenes vienošanos, rīkojušās nesavtīgi, vienīgi miera saglabāšanas interesēs. Mūsdienu krievu publicistikā ir izteikts viedoklis, ka parakstot Minhenes vienošanos, N.~Čemberlens un E.~Daladjē mierināja sevi ar ilūzijām, ka novērš lielu karu, bet J.~Staļins 1939.~gada augustā jau zināja, ka Vācija ir pieņēmusi lēmumu uzbrukt Polijai, un noslēdzot līgumu ar Vāciju, viņš šim uzbrukumam iededz zaļo gaismu, ar to kļūstot par tiešu kara izraisīšanas dalībnieku, solot neiejaukties notikumos laupītājam, kurš grib nogalināt un aplaupīt kaimiņu, par to pretī saņemot arī daļu laupījuma. Taču paskatoties dziļāk, analoģija par enerģiskā laupītāja Ā.~Hitlera piekrāptiem vientiešiem N.~Čemberlenu un E.~Daladjē un savtīgo laupītāja atbalstītāju un līdzlaupītāju J.~Staļinu sabrūk. Plaši pazīstamo N.~Čemberlena frāzi, atgriežoties Londonā no Minhenes, ka viņš atvedis ``mieru mūsu paaudzei'' varēja pateikt vai nu mazinformēts, vai ne visai gudrs, vai arī liekulīgs cilvēks. Jo ko tad nozīmēja atdot Ā.~Hitleram Sudetus? Gar to robežu ar Vāciju bija izbūvēta tam laikam modernu nocietinājumu līnija, kurai nonākot vācu rokās, pārējā Čehoslovākijas teritorija izrādījās neaizsargāta. Sudetos bija koncentrēta liela Čehoslovākijas aizsardzības rūpniecības daļa, kuras nonākšana Vācijas rokās ievērojami stiprināja tās militāro potenciālu. Visu to zinot, N.~Čemberlens nevarēja cerēt, ka viņa parakstītais līgums uz ilgu laiku nodrošinās mieru Eiropā, ka drīz netiks sagrābta arī pārējā Čehoslovākijas daļa. Viņš un E.~Daladjē vienkārši nevēlējās ļaut notikumiem ieraut savas valstis karā. Stāsta, ka E.~Daladjē, pēc Minhenes atgriežoties Parīzē un ieraugot lidostā milzu pūli, nodomājis, ka cilvēki atnākuši protestēt pret tikko noslēgto darījumu. Uzzinot, ka tie atnākuši viņu sveikt, premjeram nevilšus izspruka vārdi: ``Ir gan muļķi!'' Vismaz tie nebija liekulīgi.

Jāatgādina, ka Lielbritānija un Francija jau Pirmā pasaules kara rezultātā bija guvušas pamatīgus ``kumosus'': no Tautu Savienības ieguvušas savā pārvaldē t.s. mandātteritorijas (faktiski~--- kolonijas), kuras pirms tam bija piederējušas Vācijai un Turcijai, Francija atņēmusi Vācijai Lotringu un Elzasu, kur par pēdējās piederību Francijai vai Vācijai varētu strīdēties bezgalīgi. Turpretī Krievija, kura arī bija nesusi milzīgu Pirmā pasaules kara smagumu, pārtopot par PSRS, teritorijas zaudēja. Tāpēc 1938.~gadā Lielbritānijas un Francijas valdības bija norūpējušās par jau 1919.~gadā iegūtā saglabāšanu, bet PSRS vadība 1939.~gadā~--- par 1920.--1921.~gadā zaudētā atgūšanu. (Šai sakarā var piebilst, ka kaut arī visi reālie Polijas militārie plāni austrumu virzienā, ieskaitot arī pēdējo~--- 1939.~gada pavasarī atjaunoto, bija ar aizsardzības raksturu, taču, kā norādījis Z.~Zaluskis: ``\dots{}poļu valdošie slāņi vēlējās noturēt to, ko bija ieguvuši 1920.~gadā''.)

No morālā viedokļa īpašu starpību starp pēc Pirmā pasaules kara gūto Lielbritānijas un Francijas (tāpat arī Polijas) laupījumu un 1939.~gada PSRS laupījumu grūti saskatīt. Abos gadījumos lielvalstis nerēķinājās ar tautu pašnoteikšanās tiesībām, rūpējās tikai par savām interesēm.

Jāpiekrīt britu vēsturniekam N.~Deivisam, ka Minhene pārliecināja Ā.~Hitleru, ka turpmākā mierīgā agresija viegli atnesīs jaunas dividentes. Poļu vēsturnieks E.~Duračinskis uzskata, ka Minhenes vienošanās kā visspilgtākā piekāpšanās politikas izpausme ienesa jaunu strāvu starptautiskajās attiecībās. Destabilizācijas process pārgāja kritisko robežu. Jautājums ``kad būs karš?'' pārstāja būt tīri teorētisks. Vislielāko ietekmi tas atstāja uz Poliju: tā kļuva par to starptautiskās konjunktūras izmaiņu, kuru sākumu iezīmēja Minhene, upuri. Minhenē tika dots izšķirošs trieciens Versaļas sistēmai, kura, neraugoties uz visiem saviem trūkumiem tomēr stabilizēja situāciju Eiropā, tās centrā esošajām valstīm deva kaut nosacītu drošības sajūtu. Minhene pārliecināja Ā.~Hitleru, ka viņš arī turpmāk var īstenot savu prasību politiku, jo visticamāk tā nesastaps Rietumu pretošanos, bet vienīgo reālo draudu viņa plāniem var radīt tikai hipotētiskā Lielbritānijas un Francijas savienība ar PSRS. Ā.~Hitlers nonāca pie secinājuma, ka viņš šos draudus spēs novērst, reanimējot Rapallo politiku.

Sudetu sagrābšana būtiski stiprināja Vācijas militāro un ekonomisko varenību, padarīja Čehoslovākiju par praktiski aizstāvēties nespējīgu, tā faktiski iznīcinot tās lomu citu valstu (arī Francijas) aizsardzības sistēmā pret Vācijas agresiju. No šī viedokļa raugoties, ja Padomju Savienība, 1939.~gadā anektējot Rietumukrainu un Rietumbaltkrieviju, vismaz uz laiku atņēma šo kumosu Vācijai, kad Lielbritānija un Francija savām rokām stiprināja Vāciju, atdodot tai Sudetus. Pēc visas Čehijas (faktiski arī Slovākijas) sagrābšanas Ā.~Hitlera guvums vēl vairāk pieauga. (Starp citu, pirms Vācijas iebrukuma PSRS no 21~vērmahta tanku divīzijas 5 bija nokomplektētas ar Čehoslovākijā ražotiem tankiem.)

Bez tam Čehoslovākijas likvidācija radīja tādu stratēģisku konfigurāciju, kur Polija no trijām pusēm atradās Vācijas un Slovākijas ielenkumā. Arī par to nesa atbildību kā Minhenes līguma slēdzēji, tā Polijas valdība, kas kopā ar Vāciju sagrāba čehu teritorijas. PSRS uz Minheni nebija aicināta, nekādu atbildību par kaunpilno darījumu nenesa.

Kad 1938.~gadā Polija dabūja gabalu Čehoslovākijas teritorijas, tās valdītāju prieku aizēnoja tikai tas, ka ieguvums nebija lielāks un Polija netika uzaicināta Minhenes konferencē pievienoties četrām lielvalstīm, tātad netika atzīts tās lielvalsts statuss. Tāda bija tā reālā Polija, kuru, pēc dažu vēsturnieku uzskatiem, PSRS vajadzēja glābt par katru cenu, pat ja tas bija pretrunā ar pašas PSRS interesēm.

Kā jau augstāk klāstīts, tikai tad, kad Vācija sagrāba visu Čehiju un radīja marionešvalsti Slovākiju, tad arī Lielbritānijas un Francijas vadītāji vismaz daļēji saprata, kādā situācijā viņi bija nostādījuši savas valstis un mieru pasaulē. Tikai tad N.~Čemberlens 1939.~gada 31.~martā parlamentā paziņoja, ka Vācijas agresijas gadījumā pret Poliju Lielbritānijas valdība sniegs tai palīdzību ar visiem tās rīcībā esošajiem līdzekļiem. Lielbritānijai sekoja Francija, kuru ar Poliju arvien vēl saistīja 1921.~gada savstarpējās palīdzības līgums. Katrs nākamais Ā.~Hitlera solis automātiski ievilka šīs valstis karā. Šeit pilnībā jāpiekrīt profesoram I.~Feldmanim, ka šāda diplomātija bija ``apšaubāma un neveiksmīga''.

Ceturtkārt, nav šaubu, ka J.~Staļins un viņa režīms vēlējās izmantot radušos situāciju savās interesēs, ne tikai novilcinot PSRS iesaistīšanos karā, bet arī gūstot jaunas teritorijas, tāpat kā Polija izmantoja situāciju 1938.~gadā, gūstot Cešinas apgabalu. Var jau pieņemt, kā tas nereti tiek darīts īpaši publicistikā, ka J.~Staļins 1939.~gada vasaras situācijas izveidošanos bija jau iepriekš paredzējis un tagad viņam atlika tikai pieprasīt no abām ieinteresētajām pusēm samaksu par PSRS labvēlību, reāli to saņemot no Vācijas uz sīkāko Austrumeiropas valstu rēķina. Ja arī tiešām PSRS diktators redzēja tik tālu uz priekšu, tad jāsaka, ka to darīt viņam ļāva notikumos iesaistīto valstu realizējamā un paredzamā politika. Tomēr jākonstatē, ka vēl 1939.~gada vasarā situācija varēja pavērsties dažādi. (Protams, kad nākas runāt par vēstures gaitas alternatīvām, jāuzsver, ka visi secinājumi tiek būvēti ``uz smiltīm''. Taču šāda vēstures ``modelēšana'' ļauj arī skaidrāk izvērtēt reālo vēstures gaitu.)

Teorētiski, sākoties karam, notikumi varēja pavērsties visdažādākajos virzienos. J.~Muhins kā reālu iespējamību 1995.~gadā apskatīja pat variantu, kad pēc Vācijas uzbrukuma Polijai Lielbritānija, Francija un PSRS pieteiktu Vācijai karu, bet pēc dažām dienām Polija, Lielbritānija un Francija noslēgtu ar Vāciju separātu mieru, PSRS paliekot kara stāvoklī. Pēc tam Polija pieteiktu PSRS karu un kopā ar Vāciju sadalītu Ukrainu. PSRS šis variants esot bijis pat ļoti reāls, jo Padomju Krievija 1918.~gadā pati noslēgusi ar Vāciju Brestā separātu mieru, pametot savus sabiedrotos. (Tātad tā būtu sabiedroto atmaksa par boļševiku rīcību 1918.~gadā, noslēdzot Brestas mieru.) Arī Polijas pagrieziens par 180~grādiem bijis reāls, jo vēlāk Otrā pasaules kara beigu posmā, cietušas sakāvi, Itālija, Somija, Rumānija izdarījušas pagriezienu par 180 grādiem un karojušas ar Vāciju. Šāds spriedums gan kvalificējams kā tīras spekulācijas paraugs, jo neņem vērā nedz Lielbritānijas, Francijas un Polijas sabiedrības noskaņojumu, nedz politiskās inerces spēku, kas nekādi neļauj valstīm mainīt kursu par 180~grādiem dažu dienu laikā.

Tas pats publicists apskata arī variantu, kad Ā.~Hitlers nevis noslēgtu neuzbrukšanas līgumu ar PSRS, bet savienībā ar Poliju, Ungāriju, Rumāniju, Itāliju un Japānu (kā J.~Muhins neprecīzi raksta~--- ``ar visu Antikominternes paktu'') uzbruktu PSRS 1939.~gadā. Fīreram tad nevajadzētu rēķināties ar otras frontes iespējamību savā aizmugurē. Atkal jāsaka, ka šī teorētiskā iespēja bija nereāla, kaut vai tāpēc, ka Polija taču kā reiz atteicās pievienoties Antikominternes paktam. Alternatīvu vēstures variantu izklāstā minētais publicists šeit pametis zinātniskās analīzes ceļu, pat ignorējis vēstures faktus. Minētās alternatīvas šī darba autors uzskata par nereālām. Taču pastāvēja citas, jau krietni vien tuvākas konkrētās situācijas attīstības iespējām.

Pirmā iespēja: Lielbritānijas un Francijas valdības, pretēji dažu to līderu nespējai reāli izvērtēt situāciju, tomēr noslēgtu savienību ar PSRS, piekāpjoties tās Polijas un Baltijas valstu suverenitāti aizskarošo prasību priekšā. Tas būtu no minēto Rietumvalstu iepriekšējās politikas izrietošs jauns to pazemojums un sabiedrotās Polijas un arī Baltijas valstu nodevība, taču, iespējams, vismaz uz laiku novērstu Vācijas agresiju un radītu citām valstīm labvēlīgākus kā līdz tam priekšnoteikumus tās atvairīšanai nākotnē. Vācijas agresijas aizkavēšana dotu iespēju gan tās potenciālajiem upuriem, gan Rietumu lielvalstīm, gan PSRS labāk sagatavoties karam. Vispār novērst karu acīmredzot neizdotos. Polijas un arī Baltijas valstu liktenis tad būtu atkarīgs no konkrēto notikumu virzības, kara gadījumā~--- uzvarētājvalstu labās gribas un spēku samēra, tāpat kā tas notika pēc nacisma sagrāves 1945.~gadā.

Otrā, vairāk teorētiska nekā praktiska, iespēja: Lielbritānija un Francija, noslēgtu savienību ar PSRS, bet tas neatturētu Ā.~Hitleru no agresijas pret Poliju. Tad nu Lielbritānija un Francija varētu kopā ar PSRS iesaistīties karā pret agresoru~--- Vāciju (ar Rietumu frontes izveidi Francijā), kas nozīmētu arī ar kopīgu vienošanos sankcionētu Sarkanās armijas ieiešanu Polijas teritorijā ar acīmredzamu nolūku tajā palikt. Karš tad notiktu starp Vāciju un Lielbritānijas, Francijas, Polijas un PSRS bloku. Acīmredzot Vācija tiktu sakauta, bet tālāko notikumu gaitu, arī konkrēto Polijas pēckara statusu un robežu atkal noteiktu spēku samērs starp uzvarētājvalstīm.

Trešā iespēja, vairs ne tīri teorētiska, bet diezgan reāla. Lielbritānija un Francija varētu noslēgt savienību ar PSRS, nedomājot to nekavējoties pildīt. Ja Vācija, neraugoties uz Lielbritānijas un Francijas noslēgto savienību ar PSRS, tik un tā uzbruktu austrumu virzienā un Polijas teritorijā uzsāktu kaujas ne tikai pret Polijas bruņotajiem spēkiem, bet arī pret Sarkano armiju, ir pat ļoti ticama iespēja, ka Rietumvalstis varētu vest tādu pat ``sēdošu'' karu, kā to darīja tad, kad Vācija reāli uzbruka vienai pašai Polijai. Lielbritānija un Francija taču līdz tam nealka un arī pēc kara sākšanās nevēlētos pēc iespējas ātrāk sākt karot pret Vāciju, kad tās varētu faktiski no malas (no Britu salām un Mažino līnijas) noraudzīties Vācijas cīņā pret Poliju un PSRS, lai iejauktos tikai tad, kad tās iznākums jau būtu skaidri redzams. No Rietumvalstu interešu viedokļa sākotnēji tas būtu tām vislabvēlīgākais variants, cietējas būtu Vidus- un Austrumeiropas tautas. Vācija taču arī nesteigtos atklāt otro fronti pret Rietumvalstīm. Pavisam citādā situācijā atrastos PSRS, ja vērmahts pēc nenovēršamās Polijas sagrāves nonāktu pie robežas, kura atradās daudz tuvāk padomju centriem nekā tas bija 1941.~gadā. Vai ar Polijas atļauju tās teritorijā novēloti iegājušās Sarkanās armijas daļas varētu vērmahta tālāku virzību aizkavēt, grūti spriest, jo, kā zināms, lielam karam tolaik nebija gatava nedz Vācija, nedz PSRS. Otrajai tad tiktu uzveltas gan savas, gan Polijas un Rietumvalstu negatavības karam sekas. 1939.~gada septembrī PSRS vadība pilnībā varēja pārliecināties par savu aizdomu pret Rietumvalstīm pamatotību, savu sliktāko paredzējumu trīspusējo sarunu laikā piepildīšanos. Ja PSRS būtu noslēgusi savienības līgumu ar Rietumvalstīm, tai, tāpat kā Polijai, nāktos rūgti vilties, gaidot ātru un efektīvu to palīdzību. Par to liecina arī Rietumvalstu vilcināšanās ar Otrās frontes atklāšanu līdz 1944.~gadam.

Ja Ā.~Hitlers arī apturētu vērmahta virzību pie padomju robežas, nebija nekādas garantijas, ka Japānas valdošie slāņi necenstos par katru cenu realizēt ideju par vienlaicīgu uzbrukumu PSRS gan no austrumiem, gan rietumiem. Japāna varēja ne tikai pati doties uzbrukumā PSRS, bet arī pieprasīt lai to dara Vācija. Japāna mēģināja savu plānu realizācijā iesaistīt arī citas valstis. Tā 1939.~gada 30.~jūnijā ASV prezidents F.~D.~Rūzvelts brīdināja PSRS sūtni Vašingtonā K.~Umanski, ka Japānas puse viņam piedāvājusi nākotnē izvērst kopīgu japāņu-amerikāņu darbību Austrumsibīrijas dabas bagātību ekspluatācijā.

Ceturtā iespēja~--- PSRS nenoslēgtu līgumu nedz ar Lielbritāniju un Franciju, nedz Vāciju. Krievu vēsturnieks V.~Grosuls raksta par to, kas tad notiktu. Tādā gadījumā visu Poliju sagrābtu Vācija un Baltijas valstis atrastos tās kontrolē. Bijušās Polijas-PSRS robežas vietā būtu Vācijas--PSRS robeža. Šādā gadījumā Baltijas valstu armijas, kā arī Rietumukrainā un Rietumbaltkrievijā izveidoti formējumi Vācijas karā pret PSRS, kurš varēja sākties arī pirms 1941.~gada jūnija, piedalītos pirmās pusē. (Tēze par Baltijas valstu armiju iespējamo cīņu Vācijas pusē ir, mazākais, apšaubāma, bet kas līdzīgs baltiešu \detxti{Waffen SS} leģioniem tik tiešām varētu kampaņā piedalīties.~--- V.Š.). Ņemot vērā šādu iespēju~--- kā raksta V.~Grosuls~--- ``Pārmest PSRS starptautisko vienošanos pārkāpšanu periodā, kad no citu lielvalstu puses notika burtiski starptautisko aktu pārkāpumu kaskāde, nozīmētu nolemt valsti zaudējumam lielā karā''. Nākas piekrist, ka tāda iespēja bija tuva realitātei. PSRS sakāve karā ar Vāciju, supervarenas nacistiskās Vācijas izveidošanās, tās kundzības nostiprināšanās Eiropā nekādi neatbilda nedz Eiropas, nedz visas pasaules tautu interesēm.

Piektā iespēja faktiski tika realizēta dzīvē. PSRS noslēdza vienošanos ar Vāciju, reizē tiekot vaļā no nepatīkama kaimiņa~--- Polijas (tiesa, tā vietā saņemot vēl nepatīkamāku~--- Vāciju) un iegūstot kādreiz Krievijas impērijai piederējušās Polijas teritorijas un laiku līdz Vācijas agresijai pret PSRS.

Nemazinot Padomju Savienības vadības atbildību par to, kas notika starp 1939.~gada augustu un 1941.~gada jūniju, reizē jākonstatē, ka visi 1939.~gada un vēlākie notikumi nenorisa tukšā vietā, tie bija tās pretstāves daļa, kura starp Rietumiem un Padomēm sākās ar boļševiku nākšanu pie varas 1917.~gadā. (Ja skatās dziļāk, šo pretstāvi var atrast jau ar Maskavijas kā nopietna spēlētāja uznākšanu uz Eiropas vēstures skatuves.) Galu galā nedz PSRS, nedz Rietumiem nebija tiesību pārmest otrai pusei par to, ka starp tām valdīja naidīgums, ka katra no tām attiecīgajā periodā centās ``nomierināt'' Vāciju. Lielbritānija un Francija baidījās no komunisma izplatības, no komunistu vienošanās ar nacistiem, Padomju Savienība~--- no pret sevi vērsta kapitālistisko valstu bloka izveides.

Jāuzsver arī, ka 23.~augusta līguma noslēgšana lielā mērā veicināja to, ka tika izbeigts konflikts pie Halhingolas. PSRS un Vācijas neuzbrukšanas līguma parakstīšana Japānā izsauca dziļu sašutumu pret tās sabiedroto Antikominternes paktā (1936)~--- Vāciju. Vācijas iesaistīšanās iespējamā Japānas karā pret PSRS kļuva maz ticama. Japānas valdība demisionēja un pēc japāņu armijas graujošas sakāves augusta beigās jaunā valdība 1939.~gada 9.~septembrī piedāvāja PSRS uzsākt sarunas un noslēgt pamieru. 1940.~gada 9.~jūnijā tika panākta principiāla vienošanās par robežlīniju konflikta rajonā un 1941.~gada 13.~aprīlī parakstīts PSRS un Japānas līgums par neitralitāti. Faktiski Molotova-Ribentropa pakts pārsvītroja Antikominternes paktu. 1941.--1945.~gadā PSRS nenācās karot divās frontēs. Ja no 1937.~gada līdz 1939.~gada septembrim Tālajos austrumos bez plaši pazīstamajiem konfliktiem pie Hasana ezera un Halhgolas upes periodiski notika sīkākas sadursmes, tad pēc 1īguma noslēgšanas tās izbeidzās līdz pat 1945.~gada 8.~augustam, kad Sarkanā armija, pildot PSRS saistības Rietumu sabiedroto priekšā, devās uzbrukumā Japānai.

Krievijai vēstures gaitā pēc sakāvēm bija nācies slēgt miera līgumus, kuri nu nekādi neatbilda tās interesēm, bet kurus tai uzspieda stiprākās valstis. Lai minam Tilzītes (1807), Parīzes (1856), Portsmutas (1905), Brestas (1918), Rīgas (1921) līgumus. Visi šie līgumi bija uzspiesti Krievijai. Faktiski 1939.~gadā noslēgtais līgums ar Vāciju arī bija uzspiests PSRS. Uzspieda to gan Vācijas draudi, gan Rietumvalstu un Polijas nevēlēšanās ņemt vērā PSRS intereses. Tad nu rezultātā tapa līgums, ar kuru PSRS atguva divus gadu desmitus atpakaļ zaudētās teritorijas, bet zaudētājos bija gan visai agresīvā Polija, gan miermīlīgās Baltijas valstis. Ja Polijas vadītāji cīnījās par pusotru gadsimtu atpakaļ (1772. gadā) zaudētām robežām, tad padomju vadītāji ar vismaz tikpat lielām tiesībām varēja cīnīties par savas valsts 1914.~gada robežām. Diemžēl, ne vieni, ne otri ir nedomāja painteresēties par šais teritorijās dzīvojošo cilvēku gribu. Visi minētie līgumi tikai fiksēja spēku samēru un tāpēc bija ar lielākā vai mazākā mērā pagaidu raksturu. Tikai nacisma sagrāve ļāva nodibināt ja ne taisnīgas, tad vismaz stabilas robežas. Vainot par notikušo tikai PSRS vadību, neminot arī citus vaininiekus, ir vēsturiski nekorekti. Molotova-Ribentropa paktu zināmā mērā var salīdzināt ar 1934.~gada Polijas un Vācijas līgumu, kurš pašā sākumā pārtrauca pret hitleriskās Vācijas agresīvo kursu vērstas starptautiskas drošības sistēmas veidošanos. Tikai pret ``sanācijas'' Poliju starptautiskais sašutums nebija tik liels, jo vēl nebija skaidri redzams, kurp var novest, lai arī pagaidu, vienošanās ar Ā.~Hitleru.

Rietumu valstu apoloģēti arī izmanto faktu, ka PSRS un Vācijas neuzbrukšanas līgumā nebija norādes, ka tas zaudē spēku, ja kāda no līgumslēdzējpusēm izraisa karu pret trešo valsti. Protams, tas šajā gadījumā veicināja Vācijas agresiju pret Poliju. Taču jau tā laika PSRS valdības vadītājs V.~Molotovs norādīja, ka šādu pārmetumu izteicēji ``aizmirst'', ka tādas norādes nebija nedz Polijas un Vācijas neuzbrukšanas līgumā 1934.~gadā, kuru Vācija pret Polijas gribu anulēja 1939.~gadā, nedz Anglijas un Vācijas deklarācijā par neuzbrukšanu, parakstītā Minhenē 1938.~gada 30.~septembrī, kura faktiski deva Ā.~Hitleram pilnīgu rīcības brīvību attiecībā pret citām valstīm. V.~Molotovs uzdeva retorisku jautājumu: ``Kāpēc gan PSRS nevar darīt to, ko atļaujas darīt gan Polija, gan Lielbritānija?'' Pēc autora viedokļa, pārmetumus bija pelnījušas visas minētās valstis. Visas tās rūpējās tikai par sevi, nevis starptautisko drošību.

Poļu vēsturnieks J.~Krasuskis grāmatā ``\pltxti{Tragiczna niepodległość Polski 1918--1947}'' (Polijas traģiskā neatkarība 1918--1947, \pltxti{Toruń, Wydawnictwo Adam Marszałek}, 2007) devis šādu vērtējumu: ``PSRS noslēdza paktu ar Vāciju, lai novērstu no sevis tās trieciena galveno smagumu, bet Rietumu valstis piemēroja aizsardzības stratēģiju, lai šo smagumu pārsviestu uz PSRS.'' Līdzīgus secinājumus var atrast arī citās Rietumu autoru grāmatās.

Kaut formāli 23.~augusta līgums bija tikai par neuzbrukšanu, krievu vēsturnieks S.~Slučs uzskata, ka tas arī bez tam pievienotā slepenā protokola bija pretrunā ar starptautisko tiesību normām. 1932.~gada PSRS un Polijas līguma 3.~pantā abas puses apņēmās ``nepiedalīties nekādos līgumos, kuri no agresīvā skatu punkta būtu klaji naidīgi otrai pusei''. Turpretī PSRS un Vācijas 1939.~gada 23.~augusta līguma 2.~pants skanēja: ``Gadījumā, ja viena no Līgumslēdzējpusēm izrādīsies kara darbības objekts no kādas trešās valsts puses, otra Līgumslēdzējpuse nekādā veidā neatbalstīs šo valsti.''. Pēc S.~Sluča domām 1939.~gada 23.~augusta līgums konkrētajā situācijā, kad PSRS vadība bija informēta par Polijai draudošo Vācijas uzbrukumu, bija ``naidīgs'' Polijai. Šim secinājumam var piekrist, vienīgi piebilstot, ka abu valstu savstarpējais naidīgums, kaut vairāk vai mazāk aizmaskotā veidā, būtībā pastāvēja jau visu starpkaru periodu. Jākonstatē, ka 23.~augusta līgums akceptēja Vācijas agresiju, it kā izprovocētu ``no kādas trešās valsts [Polijas] puses''.

Krievu vēsturniece N.~Ļebedeva šai sakarā uzskata, ka līgumu var kvalificēt nevis kā līgumu par neitralitāti, bet gan par sadarbību. Daļēja taisnība tā ir, taču šeit gan jāuzsver, ka zināma PSRS un Vācijas politiska un militāra sadarbība sākās jau notikumu attīstības rezultātā, kad Polija tika sakauta un Rietumu valstis to pameta likteņa varā. Līdz tam pakts pildīja neuzbrukšanas līguma funkcijas. Poļu filozofs un publicists B.~Lagovskis pamatoti norādījis, ka Molotova~--- Ribbentropa pakts nenozīmēja savienību starp PSRS un Trešo reihu, starp Staļinu un Hitleru, bet no abām pusēm tā bija divpusēja militārā viltība, divpusēja krāpniecība, kuras mērķis bija panākt sev labākas sākuma pozīcijas priekšā stāvošajā savstarpējā karā.

Krievu žurnālists, starptautisko jautājumu speciālists P.~Fjodorovs uzskata, ka abi dokumenti~--- kā Minhenes vienošanās, tā Molotova-Ribentropa pakts ir amorāli, taču pirmais bija ļaunāks. Pēc viņa domām, Minhenes vienošanās tieši pagrūda agresora uz karu, bet Molotova-Ribentropa pakts tika noslēgts jau neizbēgama kara priekšvakarā. Bez tam pakts glāba daudzus Austrumeiropas ebrejus no iznīcināšanas nacistu nāves nometnēs. Jāatzīst sprieduma loģika.

Lai nu kā, skatīt Molotova-Ribentropa paktu kā pavērsiena punktu visā Eiropas vēsturē nav pamata. PSRS bija milzīga loma pasaules vēsturē, bet nav nekādas vajadzības to pārvērtēt. Pakts noteica vienīgi PSRS iesaistīšanās karā laiku un apstākļus. Citi notikumi pamatos attīstījās tā, kā tos noteica citu Eiropas lielvalstu valdību rīcība. Ā.~Hitlers veikt uzbrukumu Polijai bija stingri nolēmis jau 1939.~gada aprīlī. Viņš vienīgi gaidīja izdevīgu situāciju.

Kādā neoficiālā sarunā, kura notika 17.--19.~jūnijā, un kuras saturu padomju izlūki nodeva Maskavai, jau minētais J.~Ribentropa kancelejas Austrumu biroja padomnieks P.~Kleists stāstīja: ``Fīrers nepieļaus, lai angļu-franču-krievu sarunu iznākums par paktu [savstarpējās palīdzība līgumu~--- V.Š.] ietekmētu viņa gribu poļu jautājuma radikālas atrisināšanas gaitā. Vācu-poļu konfliktu Berlīne atrisinās sarunu par paktu, kā sekmīga, tā arī nesekmīga iznākuma gadījumā.'' Reizē P.~Kleists apgalvoja, ka Ā.~Hitlers, vērtējot PSRS vadības pēdējo mēnešu nostāju pret Vāciju, neticot, ka PSRS varētu ņemt dalību Lielbritānijas un Francijas militārā darbībā pret Vāciju. 19.~jūnijā tādā pat sarunā Vācijas Ārlietu ministrijas darbinieks paziņoja, ka Vācijas militārā akcija pret Poliju sāksies augusta beigās~--- septembrī ar masveida triecieniem no visām pusēm un konflikts tiks atrisināts pirms angļi un franči paspēs atgūties. Arī citi vācu izcelsmes dokumenti liecina, ka Ā.~Hitlers plānoja sakaut Poliju pirms Rietumvalstu iejaukšanās. J.~f.~Ribentrops atmiņās liecināja, ka Ā.~Hitlers augusta sākumā Itālijas ārlietu ministram G.~Čiano izteica vēlmi atrisināt Dancigas un ``poļu koridora'' jautājumu ``pie visiem apstākļiem''. 14.~augustā Ā.~Hitlers paziņoja H.~Gēringam, feldmaršalam V.~Brauhičam un admirālim E.~Rēderam par lēmumu uzbrukt Polijai ne vēlāk kā pēc divām nedēļām.

Tieši atbildīga par Otrā pasaules kara izraisīšanu bija Vācija, bet netieši arī citas lvalstis, kuras pret agresoru pielietoja ``nomierināšanas'' politiku.

Atsevišķi jārunā par \strong{slepeno protokolu} pie PSRS un Vācijas neuzbrukšanas līguma. Kaut pasaules valstīm jau bija nelāga pieredze ar dažādiem slepeniem līgumiem, kurus pēc Oktobra apvērsuma publicēja boļševiki, tie nebūt nebija izslēgti no diplomātiskās prakses. Minētais protokols pie Molotova-Ribentropa pakta bija tikai viens no tiem.

Kā savos memuāros rakstīja J.~f.~Ribentrops, vienošanās, kuras skar citas valstis, nevar tikt iekļautas sabiedrībai atklājamos līgumos, tāpēc arī slēdz slepenus līgumus. Konkrētais slepenais protokols citstarp ticis noslēgts arī tādēļ, ka vācu-krievu līgums pārkāpa vienošanos starp Krieviju [PSRS] un Poliju, kā arī 1936.~gada līgumu starp Franciju un Krieviju, kurš paredzēja savstarpējas konsultācijas, pirms puses slēdza līgumus ar citām valstīm. (Par PSRS līgumu ar Poliju pārkāpšanu J.~f.~Ribentropam var piekrist. Turpretī par PSRS 1935.~gada līgumu ar Franciju, kuru nacistu diplomāts datē ar 1936.~gadu, kad tas tika ratificēts, viņš situāciju vienkāršojis. Faktiski sarunas par šī līguma praktisko pielietošanu ar Lielbritānijas piedalīšanos taču notika 1939.~gada vasarā Maskavā, par ko jau augstāk rakstīts, un beidzās bez rezultātiem, tāpēc var uzskatīt, ka ``savstarpējās konsultācijas'' bija notikušas, taču beigušās negatīvi.

Par sarunu gaitu J.~f.~Ribentrops memuāros rakstīja: ``Jau sarunu pirmajā daļā Staļins paziņoja, ka vēlas noteikt zināmas interešu sfēras. Kā zināms, ar interešu sfēru tiek saprasts, ka ieinteresētā valsts ar pie šīs sfēras piederošo valstu valdībām ved sarunas, kuras skar tikai tās [pirmās valsts] intereses, bet otra valsts paziņo par savu kategorisko neieinteresētību. Pie tam Staļins paziņoja, ka šo valstu iekšējo struktūru viņš nevēlas skart. Uz Staļina interešu sfēras prasību, es, domājot par Poliju, atbildēju: poļi kļūst arvien agresīvāki, un būtu labi tam gadījumam, ja viņi novedīs lietu līdz karam, noteikt šķīrējlīniju, lai vācu un krievu intereses nesadurtos. Šī demarkācijas līnija tika noteikta pa Vislas, Sanas un Bugas upēm [J.~f.~Ribentrops kļūdās, Bugas upe netika pieminēta.~--- V.Š.]. Pie tam es apsolīju Staļinam, ka no Vācijas puses tiks darīts viss, lai jautājumu ar Poliju noregulētu diplomātiski-mierīgā ceļā.''

Slepenais protokols pie pakta paredzēja Austrumeiropu sadalīt padomju un vācu interešu sfērās.

Pie tam jāpasaka, ka ietekmes un interešu sfēru sadale tai laikā nebija nekas īpašs starptautiskajā praksē. Jau padomju varas pirmajos mēnešos V.~Leņina vadītā valdība publicēja Krievijas impērijas līgumus ar citām valstī, kuros pasaule tika dalīta interešu sfērās. Piemēram, tika publicēts Krievijas--Lielbritānijas 1907.~gada līgums un konvencija par ietekmes sfērām Irānā, Afganistānā, Tibetā u.c., Krievijas 1916.~gada slepenais līgums ar Lielbritāniju un Franciju par Turcijas sadali u.c.

Starptautisko tiesību speciālists G.~Tunkins pasvītrojis, ka pirms Otrā pasaules kara ``starptautiskās tiesības saturēja normas un institūtus, kuri bija tautu paverdzināšanas līdzekļi'', pie tādiem pieskaitot arī ``ietekmes sfēras.''

Arī starpkaru periodā bija starptautisku darījumu piemēri, kuriem bija sāpīgas sekas trešo valstu likteņos.

Šādos darījumos piedalījās arī t.s. demokrātiskās Rietumu valstis. Kā piemēru bijušais rīdzinieks, krievu vēsturnieks V.~Simindejs min tā saucamo ``Romas paktu'' (franču \frtxti{Accord franco-italien}, itāļu \ittxti{Accordi franco-italiani}), kurš tika noslēgts starp demokrātisko Franciju un fašistisko Itāliju 1935.~gada 7.~janvārī. Francijas valdība atdeva Itālijai daļu franču koloniju pie tolaik neatkarīgās Etiopijas robežām. Slepenajā protokolā pie šīs vienošanās bija fiksēta noruna par ``Itālijas rīcības brīvību attiecībā pret Etiopiju''. Savukārt Itālija solīja pārtraukt pret Franciju vērstu propagandu Tunisijā un saglabāt \latxti{status quo} Donavas baseinā, kā arī kopā ar Franciju garantēt Austrijas robežu neaizskaramību. Tālaika Francijas ārlietu ministrs P.~Lavals ciniski atzina: ``Labāk virzīt Itāliju Āfrikas tuksnešos, nekā ielaist to Balkānos.''

Kad 1938.~gada martā Polija draudēja Lietuvai ar karu, Vācija nodrošināja savas interešu sfēras neaizskaramību, paziņojot Polijas valdībai, ka tā kara gadījumā pretendē uz Mēmeles (Klaipēdas) apgabalu, bet ar pārējo Lietuvas teritoriju Polija var rīkoties pēc saviem ieskatiem.

Dažus mēnešus vēlāk, dalot Čehoslovākijas teritoriju, noskaidrojās, ka Vācija un Polija abas pretendē uz vienu un to pašu čehu teritoriju~--- Bohuminas pilsētu (vācu \detxti{Oderberg}, poļu un čehu \cstxti{Bogumin}). Vācija tolaik piekāpās savai sabiedrotajai pret Čehoslovākiju Polijai.

Polijas valdība emigrācijā 1941.~gada vasarā sarunās ar PSRS pārstāvjiem centās panākt slepena protokola parakstīšanu par PSRS-Polijas 1939.~gada robežu atjaunošanu.

Jau Otrā pasaules kara gaitā 1944.~gada oktobrī V.~Čērčils piedāvāja J.~Staļinam noteikt ietekmes sfēras pēckara Eiropā, tai skaitā Polijā.

Ja par pasaules sadali interešu sfērās, tai skaitā slepenu vienošanos ceļā, runā plašākā nozīmē, tad, kā norādījis krievu vēsturnieks un politologs M.~Demurins, jāmin gan Centrālās un Austrumeiropas kartes pārdale pirms 1918.--1919.~gada Parīzes konferences un tās laikā, gan virkne vienošanos, kas noveda pie Varšavas līguma, Savstarpējās Ekonomiskās Palīdzības Padomes, PSRS sabrukuma, Deitonas vienošanās par Dienvidslāvijas likteni, vienošanās par NATO izplatību Austrumeiropā. Tādejādi vēsture ir devusi neskaitāmus starptautisku līgumu piemērus, kuri smagi ietekmēja trešo valstu stāvokli. Šādus līgumus slēdza un pasauli ``ietekmes sfērās'' dalīja un dala arī valstis, kuras pieņemts uzskatīt par demokrātiskām.

Slepenais protokols pie Vācijas-PSRS neuzbrukšanas līguma paredzēja Polijas austrumu apgabalu, Igaunijas, Latvijas un Somijas iekļaušanu PSRS interešu sfērā, bet Lietuvas un Polijas rietumapgabalu~--- Vācijas interešu sfērā. Formāli vienas valsts ``interešu sfērā'' otra valsts apņēmās nekādos apstākļos neizrādīt kādu aktivitāti (piemēram, veikt kapitālieguldījumus, saņemt koncesijas, atbalstīt iespējamās sacelšanās utt.). PSRS arī neuzņēmās veikt pret Poliju militāru darbību un ar to pārkāpt ar Poliju noslēgtos un spēkā esošos līgumus, kā arī savas starptautiskās saistības. Krievu vēsturnieks Ļ.~Bezimenskis pat konstatējis, ka ``pēc satura neviens no [slepenā protokola] punktiem neiziet ārpus tai laikā plaši izplatītas prakses.''

Taču konkrētajos apstākļos 1939.~gada vasaras nogalē abas lielvalstis faktiski bija vienisprātis, ka ``interešu sfēra'' nozīmē ``teritoriāli politiskas pārkārtošanas gadījumā'' brīvību sagrābt un anektēt attiecīgo valstu teritorijas. 23.~augustā Padomju Savienība un Vācija dalīja interešu sfēras uz papīra, lai vēlāk norobežošanu pārvērstu arī realitātē. Polijā abu lielvalstu interešu sfēru robeža tika nosprausta pa upju līniju: Nareva, Visla un Sana (28.~augustā robežlīnija tika precizēta~--- pa upēm ``Pisa, Nareva, Visla un Sana''), vienlaikus fiksēta arī tās valstiskuma iespējamā izzušana. Neuzbrukšanas līgums kopumā ar slepeno protokolu ne tikai nodrošināja Vācijas aizmuguri, tai izvēršot karu Eiropā jau pēc Polijas sagrāves, bet arī PSRS interešu ievērošanu. Padomju Savienībai pirmo reizi savā vēsturē izdevās panākt, ka viena no pasaules lielvalstīm~--- Vācija~--- atzina padomju intereses Austrumeiropā. Ar to PSRS tuvinājās mērķim atgūt lielvalsts statusu, kuru tā bija zaudējusi pēc Krievijas sakāvēm Pirmajā pasaules karā un boļševiku režīma nodibināšanās. Pakts ar slepeno protokolu izšķirīgi ietekmēja arī Baltijas valstu likteni.

Ja par pašu neuzbrukšanas līgumu vēl var spriest dažādi, tā pielikums~--- slepenais protokols pēc būtības (tāpat kā daudzas citas līdz un pēc tam dažādu valstu slēgtas slepenas vienošanās) bija cinisks darījums uz vairāku valstu rēķina (nevis uz vienas, kā Minhenes vienošanās gadījumā). Īpaši negatīvi to vērtē to valstu vēsturnieki, kuru intereses tas pārkāpa.

Poļu vēsturnieki pauž vienotu un pamatotu viedokli, ka slepenā protokola saturs liecina par tā parakstītāju nodomiem neievērot pastāvošās starptautisko tiesību normas. Šie nodomi drīzumā pārvērtās konkrētā rīcībā.

Tāpat poļu vēsturnieki, izskatot Molotova-Ribentropa pakta un slepenā protokola pie tā sekas, to stāda notikumu virknē, kuri noveda pie Otrā pasaules kara. Tiesa, tālāk viņu viedokļi šķiras.

Vieni uzskata, ka parakstot līgumu un slepeno protokolu, Maskava uzņēmās tādu pat atbildību par kara izraisīšanu kā Berlīne. Pēc autora domām, ar to netiek ņemts vērā, ka karu uzsāka Vācija, bet tās pretinieki no 1939.~gada septembra~--- Polija, Lielbritānija un Francija~--- nebija paveikuši visu iespējamo tā novēršanai.

Otri apgalvo, ka, nebijis Vācijas un PSRS līguma, Ā.~Hitlers nedotu pavēli par uzbrukumu Polijai. Apgalvojumu nav iespējams nekādi pierādīt. (Autors uzskata, ka tā spriežot tiek nenovērtēts fīrera avantūrisms; viņa jau tad izdotās pavēles, kad par paktu pat tā slēdzēji vēl nevarēja iedomāties; kā arī saspringtais Vācijas iekšpolitiskais un ekonomiskais stāvoklis, jau paveiktā Vācijas tautsaimniecības pārorientācija uz kara saimniecību.)

Trešie uzskata, ka Ā.~Hitlers arī bez pakta uzbruktu Polijai, pamatoti rēķinoties ar Rietumu pasivitāti. Viņam tas bija tikai sava veida pārapdrošināšanās līgums. Šim viedoklim pievienojas arī autors, kaut jāatzīst, ka visi šie viedokļi ir tikai pieņēmumi. Diemžēl vēsturnieki nespēj droši apgalvot, kā attīstītos notikumu gaita, ja kāds to attīstības posms tiktu radikāli mainīts.

Līdzās teiktajam arī jāuzsver, ka slepenais protokols bija uzrakstīts visai amorfos izteicienos. Tajā nebija fiksētas kādas noteiktas Polijas stāvokļa izmaiņas. Lūk, teksts, kas attiecas uz Poliju: ``Teritoriāli politiskas pārkārtošanas gadījumā Polijas valstij piederīgos apvidos Vācijas un PSRS interešu sfēras aptuveni norobežo upju līnija: Nareva, Visla un Sana. Jautājumu, vai abu pušu interesēm atbilst neatkarīgas Polijas valsts saglabāšana un kādām jābūt tās robežām, galīgi varēs noskaidrot tikai tālākās politiskās attīstības gaitā. Jebkurā gadījumā abas valdības šo jautājumu risinās draudzīgas saprašanas garā''. Starp citu, ar to šis līgums atšķīrās no Minhenes vienošanās, kura tieši paredzēja atsevišķu Čehoslovākijas apgabalu nodošanu Vācijai.

No mūsdienu skata punkta raugoties, daži pētnieki ir tendēti uzskatīt, ka Ā.~Hitlers un J.~Staļins jau iepriekš labi zināja, kas notiks tuvākajā laikā. Taču tā tas nebija. 1939.~gada augusta situācija bija tik neskaidra, samezglota, ka visu valstu politiķi, tai skaitā arī abi augstāk minētie, centās parakstīt tādas vienošanās, kuras vēlāk varētu interpretēt dažādi, atbilstoši arī neparedzētiem apstākļiem. 24.~augustā neviens vēl nevarēja būt drošs, kā attīstīsies situācija. Pat pēc kara sākuma visu septembra pirmo pusi Vācija nevarēja būt droša par J.~Staļina tālāko rīcību, līdz viņš nelika Sarkanajai armijai iesoļot Polijā. Jau tas vien rāda, ka slepenajam protokolam varēja būt vairākas interpretācijas.

Objektīvi (ja izslēdz Ā.~Hitlera subjektīvo apsēstību ar kara nepieciešamību) vēl varēja rasties iespēja vismaz pagaidām atrisināt pastāvošās pretrunas vienošanās ceļā, starptautiskā, varbūt arī līdzīgā kā 1938.~gadā Minhenē, konferencē. Teorētiski arī Polija pēdējā brīdī varēja vēl noslēgt vienošanos vai nu ar Berlīni vai Maskavu. Katram gadījumam Vācijas un PSRS neuzbrukšanas līgums un slepenais protokols pie tā bija aprēķināts dažādām iespējamām situācijām.

Starptautisko tiesību speciālists, ukraiņu profesors V.~Makarčuks uzskata, ka ``Slepenais protokols par ietekmes sfēru sadali starp nacistisko Vāciju, no vienas puses, un PSR Savienību, no otras, bija sastādīts tā, ka formāli nepārkāpa tai laikā pieņemto starptautiski-tiesisko dokumentu sastādīšanas praksi.'' Minētais jurists secina, ka, ja arī slepenā protokola teksts gadījuma dēļ vai ar nodomu no Vācijas puses tiktu publiskots, tas varētu radīt PSRS drīzāk morālus, nevis starptautiski-tiesiskus zaudējumus. Pats par sevi protokols nedeva iespēju iztēlot PSRS par agresoru. Vēlme mainīt robežas savā labā un arī praktiski soļi uzstādītā mērķa īstenošanai pati par sevi vēl nebija ne toreizējo, ne mūsdienu starptautisko tiesību pārkāpums. Jautājums bija par pieļaujamu vai nepieļaujamu, ievērojot starptautiskās tiesības, mērķa sasniegšanas līdzekļu pielietošanu. Zem Vācijas, Itālijas, Polijas un Ungārijas spiediena Rietumu lielvalstis taču pieļāva Eiropā pastāvošo robežu izmaiņas (Minhene, pirmā Vīnes arbitrāža). Pie tam šīs izmaiņas notika, izmantojot draudus pielietot spēku (nemaz nerunājot par to, ka dažādos agresīvo valstu izprovocētajos incidentos tas arī tika lietots), neievērojot ieinteresētās puses pretestību. Polijas valdība atklāti darbojās pēc Ā.~Hitlera noteikumiem, pirmo reizi, kad kopā ar Vāciju pēc Minhenes sadalīja Čehoslovākiju, otro reizi, kad klaji atbalstīja Karpatu Ukrainas aneksiju, ko veica Ungārija. Ar to Polija nostādīja sevi blakus Vācijai (precīzāk~--- aiz tās, jo Polija tomēr iniciatīvu neuzņēmās un bija krietni vien vājāka), atbalstot starptautiski sankcionētu spēka izmantošanu lai mainītu pastāvošās robežas.

Arī pašas Polijas robežu jautājumā, ja tikai Ā.~Hitleram tas šķistu izdevīgi, bija iespējams sasaukt vēl vienu ``miera'' konferenci pēc Minhenes konferences parauga, lai īstenotu teritorijas pārdali ``miera ceļā.'' Tāda iespēja pastāvēja arī pēc Molotova-Ribentropa pakta noslēgšanas, tikai šoreiz tajā, ņemot vērā noslēgto paktu, varētu piedalīties ne tikai Rietumu valstis, bet arī PSRS. 26.~augustā PSRS sūtnis Londonā I.~Maiskis ierakstīja savā dienasgrāmatā: ``Gaisā ož pēc jaunas Minhenes \citespace{} Ja Hitlers parādīs kaut minimālu gatavību vienoties, var atkārtoties pagājušā gada notikumi''. Sūtnis neko nezināja par slepeno protokolu, sprieda pēc britu varas vīru izturēšanās un tāpēc kļūdījās. Ā.~Hitleram miera līgums ar Rietumvalstīm nebija vajadzīgs, viņš vēlējās izraisīt karu. ``Otra Minhene'' nenotika, taču arī ar šādu iespēju, pirms karš patiesi nebija sācies, PSRS vadībai bija jārēķinās.

To, ka nacistiskā Vācija 1939.~gada 1.~septembrī uzbruka Polijai, nevar izskaidrot tikai ar PSRS vadības piekrišanu šai agresijai. Nevienā Vācijas Ārlietu ministrijas dokumentā, kuri visi ir publicēti, nav pat mājiena par to, ka PSRS vadība tā vai citādi bija devusi savu piekrišanu karam pret Poliju. Var jau teikt, ka ar paktu bija dota šāda klusējoša piekrišana tālākai Ā.~Hitlera rīcībai. Taču viss varēja pavērsties citādi, ja citādi beigtos Vācijas-Polijas karš. Kaut Rietumukraina un Rietumbaltkrievija pēc pakta noslēgšanas 1939.~gada septembrī, bet Baltijas valstis un Besarābija 1940.~gadā nonāca PSRS sastāvā pēc Molotova-Ribentropa pakta ar tā slepeno protokolu tapšanas, ``pēc'' vēl nenozīmē ``tāpēc''. Šie notikumi tikai netieši bija saistīti ar Molotova-Ribentropa paktu un slepeno protokolu pie tā. Pakts tikai pavēra iespēju šādai notikumu attīstībai. ``Pakts nebija nedz savienība, nedz vienošanās par Polijas sadalīšanu'', rakstīja ievērojamais britu vēsturnieks A.~Teilors ``Minhene tiešām bija savienība sadalīšanai: briti un franči nodiktēja sadalīšanu čehiem. Padomju valdība neveica tādu akciju pret poļiem. Viņi apsolīja palikt neitrāli, ko vienmēr lūdza poļi, un ko pieprasīja Rietumvalstu politika''. Tikai tad, kad Polijas liktenis jau faktiski bija izšķirts, kad Rietumvalstis tai nesniedza solīto palīdzību, kad vērmahts tuvojās PSRS robežām, tā virzīja Sarkano armiju Polijas teritorijā. Ja Francija, Lielbritānija reāli ietu Polijai palīgā, PSRS varētu vismaz 1939.~gadā palikt arī bez Rietumukrainas un Rietumbaltkrievijas, vēstures gaita virzītos citādi. Teritoriju sadale atkarībā no spēku samēra notiktu kara beigās, kura gaitu nebija iespējams prognozēt. Citiem vārdiem, J.~Staļins 23.~augustā bija atstājis sev brīvas rokas. Domājams, ka, Francijai un Lielbritānijai reāli iesaistoties karā, PSRS nepārkāptu neitralitātes robežas, ļautu karojošajām valstīm noskaidrot spēku samēru, savstarpēji noasiņojot, un attiecīgi rīkotos tikai tad, kad kara rezultāts būtu skaidri redzams. Neiejaukšanās politika šai gadījumā atbilstu PSRS interesēm.

Arī krievu vēsturnieks Ā.~Djukovs norādījis, ka, kā padomju--vācu līgums par neuzbrukšanu, tā arī slepenais protokols pie tā, nereglamentēja kādu teritoriālu Eiropas pārdali, bet sadalīja tikai ``interešu sfēras''. To, ka padomju ``interešu sfērā'' nonākušās teritorijas, arī tai laikā Polijai piederošās Rietumukraina un Rietumbaltkrievija var tikt pievienotas PSRS, minētie dokumenti neparedzēja (iekļaušana ``ietekmes sfērā'' nebūt nenozīmē aneksiju). Viņš arī uzsver, ka līdz mūsdienām tiešas juridiskas pakta sekas nav saglabājušās. Mūsdienu Eiropas valstu robežas radās pēckara teritoriālā noregulējuma pasākumu, Eiropas drošības un sadarbības apspriedes (1975) atzītās pēckara robežu neaizskaramības un 1991.~gadā notikušā PSRS sabrukuma rezultātā. Ja runā par tādām ``pakta sekām'' kā teritoriju pievienošana PSRS 1939.~gadā, tad jākonstatē, ka uz Krieviju, Igauniju un Latviju, Centrālāzijas un Aizkaukāza valstīm tās mūsdienās neattiecas. PSRS 1939.~gada ieguvumus mūsdienās izmanto Lietuva (Viļņas apgabals), Baltkrievija (valsts rietumu apgabali), Ukraina (valsts rietumu apgabali). Pārējās valstis ir atzinušas to teritoriālo integritāti (Izņemot Krievijas Federācijas izdarīto Krimas pievienošanu 2014.~gadā). Tāpēc dīvains šķiet profesora I.~Feldmaņa apgalvojums: ``No starptautisko tiesību aspekta Molotova-Ribentropa pakts bija neskaitāmas reizes noziedzīgāks par nelikumīgo Minhenes vienošanos, kuras mērķis tomēr bija miera saglabāšana. Turpretī 23.~augusta pakts atļāva gan Ā.~Hitleram, gan arī Staļinam (šo faktu nedzkāpēc bieži ignorē) sākt karu, uzbrūkot kaimiņiem\dots{}''

Pirmkārt, Minhenes vienošanās tik tiešām ``atļāva'' Ā.~Hitleram sagrābt sākotnēji Sudetus, bet pēc tam ``bez atļaujas'' sadalīt visu Čehoslovākiju. Zinot, kas notika vienošanās rezultātā, apgalvot, ka tās mērķis ``bija miera saglabāšana'', var tikai ar neskaitāmām atrunām. Otrkārt, formāli Molotova-Ribentropa pakts nedeva kaut kādu ``atļauju'' nedz Ā.~Hitleram, nedz J.~Staļinam. Tāpat taču var rakstīt, ka Rietumvalstis, atsakoties no reālas palīdzības Polijai 1939.~gada septembrī, ``atļāva'' Vācijai to iznīcināt. Treškārt, jau iepriekš parādīts, ka Ā.~Hitlera uzbrukums bija ieplānots pusgadu iepriekš un to apturēt varēja tikai visu apdraudēto valstu: Polijas, Francijas, Lielbritānijas un PSRS kopēji pūliņi. Ceturtkārt, PSRS karaspēka ieiešanu Polijas teritorijā par ``kara sākšanu'' pret Poliju 1939.~gadā neatzina nedz pati Polija, nedz tās sabiedrotās Lielbritānija un Francija, nedz mūsdienu starptautiskās tiesības. Piektkārt, ja par PSRS ``uzbrukumu kaimiņiem'' tiešām var runāt, arī tad jāatceras konkrētais vēsturiskais konteksts. Ja nosodām 1939.~gada PSRS uzbrukumu Polijai, tāpat jānosoda kopš Polijas izveides tās veiktie daudzie agresīvie akti. Ne velti ``no starptautisko tiesību aspekta'' Lietuvas, Baltkrievijas, Ukrainas robežas mūsdienās ir nepārprotami atzītas, kaut daļa to teritorijas kādreiz piederēja Polijai.

Uz brīdi stādīsimies priekšā, ka Molotova-Ribentropa pakta nosodītāji par spīti veselajam saprātam varētu panākt starptautisku piekrišanu ``pakta seku'' likvidēšanai, pirms līguma eksistējušo valstu robežu atjaunošanai. Tas nozīmētu pieļaut ne tikai nozīmīgu Ukrainas, Baltkrievijas un Lietuvas teritoriju nodošanu Polijai, bet arī tagadējo Polijas rietumu teritoriju atdošanu Vācijai, tātad~--- tādas nacionālās netaisnības atjaunošanu, kādu iedibināja Versaļas, Rīgas līgumi un citi 20.~gadu sākuma regulējumi. Domājams, ka, šādu iespēju iedomājoties vien, pasaules sabiedrības vairākums tādu soļu piekritējus pasludinātu par tālu turamiem no reālās politikas. Ne velti PSRS sabrukuma un Baltijas valstu valstiskuma atjaunošanas periodā Lietuvas delegācija dažādos baltiešu saietos vienmēr panāca, ka netika nosodītas Molotova-Ribentropa pakta sekas, jo tad lietuviešiem draudēja šķiršanās no Viļņas novada. Lietuvas politiķu praktiskā rīcība atbilst krievu diplomāta A.~Duļjana secinājumam, ka paktam vēsturiskā perspektīvā bija liela praktiska nozīme XX gadsimta beigās, veicinot nacionālo valstu (Lietuvas, Ukrainas, Baltkrievijas) atjaunošanos un nodibināšanos to tagadējās robežās. Mūsdienu krievu vēsturnieks, politologs un politiķis V.~Ņikonovs ir pievērsis uzmanību faktam, ka 1999.~gadā zinātniskā konferencē Kijevā sakarā ar Molotova~--- Ribentropa pakta 60.~gadadienu tā tika uztverta kā Ukrainas atkalapvienošanās gadskārta. (Šeit gan jāpiebilst, ka mēneša laikā pēc pakta noslēgšanas notikumi varēja gūt dažādus pavērsienus un ``Ukrainas atkalapvienošanās'' varēja arī nenotikt.)

Jāsaka, ka Minhenes vienošanās un Molotova-Ribentropa pakta salīdzinājums nav par labu Rietumu valstīm.

Loti precīzi ir pateicis V.~Čērčils: ``Pakts bija tieši tik cinisks, cik neizbēgams.'' Pēc vienas no starptautiskajās attiecībās izplatītajām spēles noteikumu formulām (piedēvētas dažādiem autoriem~--- no O.~f.~Bismarka līdz S.~Kubrikam) uz pasaules skatuves lielās valstis rīkojas kā laupītāji, bet mazās kā prostitūtas. Neapšaubāmi, PSRS kā 1939.~gadā, tā arī citkārt, tāpat kā citas lielvalstis, uzvedās kā laupītāja. (Tālāko frāzes daļu nevajadzētu saprast pārāk burtiski, kaut pilnībā noliegt to arī nevar.) Nav pamata iztēlot J.~Staļinu kā nezinīti, kurš nevarēja paredzēt pakta iespējamās sekas. Neapšaubāmi, viņš rēķinājās ar realitātē notikušo variantu. Viņš apzinājās, ko dara un vēlējās gūt no tā maksimālu ``peļņu''. To pierāda daudzi dokumenti, tapuši starp 1939.~gada 23.~augustu un 1941.~gada 22.~jūniju. Saprotot savas pozīcijas amoralitāti, padomju vadība dažādi maskēja noslēgtā pakta un tam pievienotā slepenā protokola jēgu.

Tā, 22.~augustā, vēl pirms pakta un protokola parakstīšanas padomju preses dienests paziņoja, ka ``sarunas par neuzbrukšanas līgumu ar Vāciju nekādā mērā nespēj pārtraukt vai palēnināt angļu-franču-padomju sarunas. Runa iet par miera lietas veicināšanu: vienas virzītas uz starptautiskā saspīlējuma mazināšanu, otras~--- uz ceļu un līdzekļu sagatavošanu cīņai pret agresiju, ja tā notiks''. Paziņojums atbilstu patiesībai, ja nākamajā dienā tapušajā neuzbrukšanas līgumā kā vairākumā iepriekšējo padomju neuzbrukšanas līgumu būtu iekļauts pants, kas dotu iespēju to nekavējoties anulēt, ja kāda no parakstītājvalstīm veiktu agresiju pret trešo valsti, un, galvenais, tam netiktu pievienots slepenais protokols par Eiropas sadali ietekmes sfērās.

J.~f.~Ribentrops savos memuāros stāsta, ka sarunu gaitā viņš pajautājis J.~Staļinam par to, kas tad notiks ar Lielbritānijas un Francijas delegācijām, kuras vēl veda sarunas Maskavā. J.~Staļins atbildēja: ``No tām pieklājīgi atvadīsies''. Uz citu jautājumu, kā savienot Vācijas un PSRS neuzbrukšanas līgumu ar PSRS un Francijas līgumu par savstarpējo palīdzību, J.~Staļins atbildēja: ``Krievijas intereses ir augstākas par visām citām''. Līdz ar Molotova-Ribentropa pakta noslēgšanu Lielbritānijas un Francijas delegāciju Maskavā vestās sarunas pārtrūka, kaut formāli tās varēja turpināties. Sarunā ar Francijas sūtni Maskavā P.~Nadžiaru V.~Molotovs paziņoja, ka ``līgums par neuzbrukšanu ar Vāciju nav nesavienojams ar savienību par savstarpējo palīdzību starp Lielbritāniju, Franciju un Padomju Savienību''. 24.~augustā Lielbritānijas Ārlietu ministrija telegrammā britu sūtniecībai Vašingtonā atsaucās uz ārlietu tautas komisāra vārdiem, ka ``pēc kāda laika, piemēram, pēc nedēļas, sarunas ar Franciju un Angliju var tikt turpinātas.'' Arī padomju historiogrāfijā ilgstoši tika uzturēts apgalvojums, ka pēc Molotova-Ribentropa pakta noslēgšanas augusta beigās Padomju valdība tāpat kā iepriekš bija gatava turpināt sarunas ar Lielbritāniju un Franciju par savstarpējās palīdzības līguma noslēgšanu. Tas tika argumentēts ar V.~Molotova konfidenciālu paziņojumu Francijas sūtnim 23.~augustā un ārlietu ministra vietnieka S.~Lozovska paziņojumu Ķīnas sūtnim 26.~augustā. Taču, kad 25.~augustā admirālis P.~Drakss un ģenerālis Ž.~Dumenks ieradās pie PSRS aizsardzības tautas komisāra un padomju delegācijas sarunās vadītāja K.~Vorošilova un Sarkanās armijas Ģenerālštāba priekšnieka B.~Šapošņikova un admirālis P.~Drakss uzdeva jautājumu vai padomju militārā misija uzskata par iespējamu turpināt sarunas arī mainījušās politiskās situācijas apstākļos, K.~Vorošilovs viennozīmīgi atbildēja, ka ``tālākas sarunas zaudē jebkuru jēgu''. Tiesa, tai pat 25.~augustā K.~Vorošilovs intervijā presei it kā atspēkoja ziņu aģentūras \entxti{Reuters} izplatīto ziņojumu, ka viņš paziņojis angļu un franču militārajām misijām, ka sakarā ar padomju~--- vācu līguma parakstīšanu Padomju valdība uzskata tālākas sarunas ar Angliju un Franciju par bezmērķīgām. PSRS aizsardzības tautas komisārs uzsvēra: ``Nevis tāpēc pārtrūka militārās sarunas ar Angliju un Franciju, ka PSRS noslēdza neuzbrukšanas paktu ar Vāciju, bet otrādi, PSRS noslēdza paktu par neuzbrukšanu ar Vāciju, starp citu, tā apstākļa rezultātā, ka, militārās sarunas ar Franciju un Angliju nonāca strupceļā nepārvaramu domstarpību rezultātā.'' Šos vārdus 27.~augustā citēja laikraksts ``\rutxti{Правда}''.

Tādejādi var jau uzskatīt, ka V.~Molotova, S.~Lozovska un arī K.~Vorošilova paziņojumi bija tikai diplomātisks aizsegs sarunu ar Rietumvalstīm pārtraukšanai. Taču tāpat var arī izteikt pieņēmumu, ka, iespējams, PSRS vadība bija gatava noslēgt sava veida pārapdrošināšanās līgumu par sadarbību ar Rietumu valstīm, (varbūt arī vairs necīnoties par tiesībām ievest Sarkano armiju Polijā), jo starptautiska konflikta gadījumā tad PSRS varētu izvēlēties, kuru līgumu pildīt~--- ar Vāciju vai Lielbritāniju un Franciju noslēgto. Šeit tad varētu izpausties J.~Staļina atbildes J.~f.~Ribentropam jēga, ka ``Krievijas intereses ir augstākas par visām citām''.

Tiesa, no Lielbritānijas un Francijas viedokļa to savienība ar PSRS būtu pilnīgā pretrunā ar PSRS un Vācijas neuzbrukšanas līgumu, kurā abas puses apņēmās atturēties no savstarpējiem uzbrukumiem ``gan atsevišķi, gan kopā ar citām valstīm''; bet, ja vienai no pusēm uzbruktu trešā valsts, otra puse solīja nekādā mērā neatbalstīt šo trešo valsti, kā arī nepiedalīties nekādos grupējumos, vērstos pret partnervalsti. Sarunu turpināšana ar PSRS arī pēc Molotova-Ribentropa pakta noslēgšanas nekādi neatbilda Lielbritānijas un Francijas valdību plāniem, kad nevis tām, bet PSRS vadībai piederētu izšķiršanās, kā reaģēt uz gaidāmo Vācijas agresiju pret Poliju, un tās atsauca savus pārstāvjus no Maskavas. N.~Čemberlens, acīmredzot vēl cerot ``nomierināt'' Ā.~Hitleru, paziņoja: ``Galvenais ir tas, ka, ja Anglija liks mierā Hitleru viņa sfērā [Austrumeiropā], tad viņš liks mierā mūs.''

Maskavā augusta beigās norisa PSRS Augstākas Padomes sesija, kura ratificēja neuzbrukšanas līgumu ar Vāciju. Šai sakarā izskanējušajā V.~Molotova runā Polija tika nosaukta par vainīgu par to, ka netika panākta drošības līguma noslēgšana, kas tikai daļēji atbilda patiesībai. Entuziasmu zālē izsauca V.~Molotova sarkasms attiecībā pret Lielbritānijas un Francijas sociālistisko partiju līderu vēlmi, lai PSRS karotu sabiedroto pusē pret Vāciju: ``Lai pakaro paši, bez PSRS. Mēs paskatīsimies, kas viņi par karotājiem''.

Vai PSRS 1939.~gadā varēja izglābt Poliju?~--- Tikai teorētiski. Tad viņai vajadzētu vienai pašai (ar nelielu poļu atbalstu un, ja Polija pēc savas armijas sakāves beidzot piekristu ielaist Sarkano armiju savā teritorijā) karot pret Vāciju. Lielbritānija un Francija varētu tāpat kā Čehoslovākiju nodot arī Poliju un PSRS. Jāņem vērā, ka padomju vadītāji pēc Minhenes vienošanās baidījās no tā, ka Lielbritānija un Francija varētu atkal noslēgt vienošanos ar Vāciju. Risks, ko uzņemtos šādā gadījumā PSRS, būtu milzīgs. Bez sadarbības ar Rietumiem karot ar Vāciju par Polijas, ar visām tās sagrābtajām lietuviešu, baltkrievu un ukraiņu zemēm, brīvību un teritoriālo integritāti, no PSRS puses būtu neprāts.

Mūsdienu krievu vēsturnieks V.~Dorošenko uzskata, ka Polijas pastāvēšana starp Vāciju un PSRS otrajai bija pēdējais drošības faktors un tai [Polijas pastāvēšanai] vajadzētu būt par padomju ārpolitikas stūrakmeni, ja vien J.~Staļins rēķinājās ar aizsardzības, nevis uzbrukuma politiku. Vēsturnieks uzskata, ka J.~Staļins pieļāva kļūdu, izrietošu no vācu-padomju attiecību perspektīvas nepareiza novērtējuma. Toties cits mūsdienu krievu vēsturnieks M.~Narinskis secinājis: ``Vienkārši Staļins uzskatīja vienošanos ar Hitleru par izdevīgāku un īstenoja šo darījumu. Viņš ieguva nevis laiku, bet teritoriju''. Lūk, kā sprieda V.~Molotovs: ``Ja mēs 1939.~gadā neietu pretī vāciešiem, viņi ieņemtu visu Poliju līdz robežai. Tāpēc mēs ar viņiem vienojāmies. Viņi bija spiesti piekrist. Tā bija viņu iniciatīva~--- pakts par neuzbrukšanu. Mēs nevarējām aizstāvēt Poliju, jo tā nevēlējās ar mums vest lietas. Nu un par cik Polija nevēlas [sadarboties], bet karš ir deguna priekšā, dodiet mums kaut to daļu Polijas, kura, kā mēs uzskatām, bez ierunām pieder Padomju Savienībai''.

Taču šāds pragmatisms nevar attaisnot līgumu laušanu, agresiju pret nepatīkamu, bet starptautiski atzītu kaimiņvalsti. 1989.~gada decembrī jautājumu par PSRS un Vācijas neuzbrukšanas pakta politisko un juridisko vērtējumu izskatīja PSRS Tautas deputātu kongress. Tas atzina, ka slepenais protokols jau sākotnēji bija prettiesisks dokuments, to parakstošo personu slepena vienošanās (V.~Molotovam nebija oficiālu pilnvaru parakstīt slepenus protokolus nedz 23.~augustā, nedz arī pēc mēneša~--- 28.~septembrī. Šie protokoli netika apspriesti valdībā, tie netika iesniegti ratifikācijai PSRS Augstākajā Padomē.). Ne pēc sastādīšanas metodes, ne pēc satura slepenais protokols ``neatbilda padomju ārpolitikas ļeņiniskajiem principiem''. (Šis PSRS Tautas deputātu kongresa formulējums gan slēpj PSRS, tāpat kā citu impēriskas politikas realizētājvalstu, bezprincipialitāti.~--- V.Š.) Protokolā notikusī PSRS un Vācijas ``ietekmes sfēru'' nodalīšana bija pretrunā ar trešo valstu neatkarību un suverenitāti. Kongress atzina slepenos protokolus (23.~augusta un vēlākos) par spēkā neesošiem no to parakstīšanas brīža.

Tiesa, kā norādījis jau minētais krievu vēsturnieks un politologs M.~Demurins, PSRS Tautas deputātu kongresa lēmumā bez pierādījumiem bija kļūdaini apgalvots, ka protokolus ``Staļins un viņa apkārtējie izmantoja ultimātiem un spēka spiedienam pret citām valstīm un viņu priekšā uzņemto tiesisko saistību pārkāpšanai''. Speciālista kritikai piekrist var tikai daļēji~--- slepenos dokumentus tieši izmantot norādītajā veidā nebija iespējams tāpēc vien, ka tie bija slepeni, taču jau mājieni un starp valdībām cirkulējošās ziņas par to esamību, piemēram, spieda Baltijas valstu valdības 1939.~gada septembra beigās, oktobra sākumā parakstīt savstarpējās palīdzības paktus ar PSRS.

Ar neuzbrukšanas līguma noslēgšanu PSRS faktiski pievienojās Rietumu valstīm, kuras, realizēja agresora ``nomierināšanas politiku'' (jāatzīst, ka nosaukums nav precīzs, gan Rietumvalstis, gan PSRS agresoru ``nomierināja'' tikai vienā~--- savā virzienā, citos ļaujot tam plosīties), ļāva Ā.~Hitleram vienu pēc otra graut potenciālos Rietumu sabiedrotos (Lielbritānija un Francija~--- Austriju, Čehoslovākiju, PSRS un arī Rietumu valstis~--- Poliju), līdz pašas atradās aci pret aci ar agresoru. Vēl vairāk, PSRS rīcība 1939.~gadā līdzinājās Polijas rīcībai 1938.~gadā, izmantojot labvēlīgo situāciju un sagrābjot Cešinu. Tikai PSRS laupījums bija lielāks. Latvijas ZA loceklis I.~Ronis konstatējis: ``Ko patiesībā nozīmēja Ribentropa--Molotova pakts? To, ka tiek atjaunotas divas~--- Vācijas un Krievijas~--- impērijas to sākotnējās robežās ar nelielām niansēm. Nekā cita tur nav. Pārējais ir vēsturnieku fantāzijas.''

Var ilgstoši strīdēties, kurš diktators īsākā un ilgstošākā perspektīvā vairāk ieguva no Molotova-Ribentropa pakta. Sākotnēji ieguvēji bija abi. Tālākais bija atkarīgs no ļoti daudziem faktoriem, bet visvairāk no tā, cik lielā mērā abiem agresoriem izdosies nostiprināt savus militāros spēkus. Krievu pētnieks V.~Ņikonovs raksta: ``Kā vēsturnieks pēc izglītības, kas gatavojis izdevumu ``V.~Molotovs'' trijos sējumos un izpētījis tūkstošiem arhīvu dokumentu, es varu apliecināt: ja pakts netiktu noslēgts, Hitlera karš pret PSRS būtu sācies 1939., nevis 1941.~gadā. Jautājums bija tāds: kur Vācijas armija apstāsies pēc uzbrukuma Polijai? Ja šo līniju nebūtu izdevies nospraust, tā varētu atrasties Urālu rajonā. Laikā no 1939.~gada līdz 1941.~gadam PSRS spēja divkāršot savu militāro potenciālu.'' Tiesa, tas ir viena pētnieka spriedums un šeit nākas piebilst, ka Vācija šai laikā, pateicoties iekarojumiem, savu potenciālu kāpināja straujāk nekā PSRS. Kā parādīja veiksmīgais Vācijas iebrukums PSRS 1941.~gadā, līdz tam labāk sagatavoties savstarpējam karam bija izdevies Ā.~Hitleram. Taču 1939.~gadā to nevarēja droši paredzēt. Sava taisnība ir tiem, kuri uzsver~--- ja nebūtu pakta, Vācijas karaspēkam 1941.~gadā nenāktos ar kaujām noiet simtus km pa tai laikā jau Padomju Savienības rokās esošo Rietumbaltkrieviju un Rietumukrainu un tas varētu nonākt pie Maskavas daudz agrāk nekā tas notika realitātē, un tad 1941.~gada novembrī Sarkanajā laukumā varētu notikt nevis Sarkanās armijas, bet vērmahta parāde. Antihitlerisko koalīciju tas nostādītu ārkārtēju, ja ne pat katastrofālu grūtību priekšā. Pēc pakta noslēgšanas notikumu gala rezultāts ir zināms~--- nacisti nonāca līdz Maskavai un Staļingradai, taču karš beidzās Berlīnē. Neviens nevarēja garantēt, kāds būtu rezultāts, ja Padomju Savienībai, neparakstot paktu ar Vāciju, nāktos ar to karot jau pirms 1941.~gada jūnija un no Maskavai daudz tuvākas robežas~--- pārāk dažādi varēja attīstīties notikumi. Protams, tāpat nevar kalkulēt Polijas iespējamo likteni.

J.~Staļins, noslēdzot vienošanos ar Vāciju kopā ar slepeno protokolu, rīkojās ne sliktāk un ne labāk kā jebkurš valsts vadītājs, kurš pirmkārt ievēro savas valsts intereses, kā nu viņš tās saprot, tai pat laikā neievērojot citu valstu intereses un vajadzības gadījumā pārkāpjot agrāk noslēgtos starptautiskos līgumus. Savas valsts teritorijas paplašināšana jebkuriem līdzekļiem visos laikos taču tika uzskatīta par valsts vadītāju nopelnu. Izņēmumus var atrast, bet to ir daudz mazāk nekā pretēju piemēru. Jau 1937.~gada novembrī J.~Staļins saviem līdzgaitniekiem paziņoja: ``Krievu cari izdarīja daudz slikta. Viņi aplaupīja un paverdzināja tautas. Viņi veda karus un sagrāba teritorijas muižnieku interesēs. Bet viņi paveica vienu labu darbu~--- radīja milzīgu valsti~--- līdz Kamčatkai. Mēs saņēmām šo valsti mantojumā.'' Arī 1939.~gadā J.~Staļins darbojās atbilstoši šai Krievijas impērijas mantinieka, milzīgas valsts valdnieka loģikai. Taču tādām koloniālajām impērijām kā Lielbritānijai un Francijai nebija morālu tiesību pārmest padomju impērijas valdniekam savas impērijas paplašināšanu. (Jau Otrā pasaules kara noslēguma posmā Lielbritānijas ārlietu ministrs A.~Īdens, ziņojot parlamentam par Jaltas konferences (1945.~gada 4.~--- 11.~februārī) rezultātiem pārmeta saviem priekšgājējiem: ``Vai gan šodien kāds šaubīsies par to, ka, ja vienotība starp Krieviju, Britāniju un Savienotajām valstīm, nodibināta Jaltā, būtu pastāvējusi [jau] 1939.~gadā, tad šis karš nekad neizraisītos?'')

Mūsdienu literatūrā daļa vēsturnieku un publicistu, tai skaitā arī V.~Suvorovs (Rezuns), apgalvo, ka J.~Staļins jau ilgstoši gatavoja karu pret Vāciju un plānoja to uzsākt 1941.~gada vasarā. Nav autora uzdevums iedziļināties šai jautājumā. Taču fakts ir fakts, ka karš starp Vāciju un PSRS sākās 1941.~gada 22.~jūnijā ar Vācijas iebrukumu PSRS, neatkarīgi no tā, gatavoja vai negatavoja J.~Staļins uzbrukumu Vācijai. Ja karu dažas dienas vai nedēļas vēlāk sāktu PSRS, arī tad tas sāktos tai teritoriāli labvēlīgākās pozīcijās nekā tās bija 1939.~gadā pirms pakta.

Var piekrist V.~Suvorovam, kad viņš raksta: ``\dots{}savienība ar Hitleru bija taktisks paņēmiens, uzmanības novēršanas manevrs. Bet Staļina stratēģiskā iecere~--- Vācijas sagrāve un Eiropas atbrīvošana no brūnā mēra. Būt Hitlera draugam ir kauns un neiedomājama nelietība. Uzbrukt Hitleram~--- svēta lieta.'' Šim viedoklim būtu jāpiekrīt vismaz visām uz laiku Ā.~Hitlera, nacistu jūgā nonākušajām daudzcietušajām tautām.

1939.~gadā J.~Staļins, iespējams, cerēja uz to, ka karš izvērsīsies Rietumos, ilgs tur vairākus gadus, tāpat kā Pirmais pasaules karš, un novedīs pie jaunas revolūcijas, kur PSRS varēs izpildīt savu ``revolucionāro pienākumu'' palīdzēt tautām ``atbrīvoties no ekspluatācijas''. Teorētiski iespējams, ka viņš jau 1939.~gadā paredzēja, ka Ā.~Hitlers iekaros gandrīz visu Eiropu un jau tad iecerēja pats sākt uzbrukumu Vācijai. Taču vairāk iespējams, ka pie šāda soļa nepieciešamības viņš nonāca vēlāk, Ā.~Hitlera sekmīgo iekarojumu rezultātā un redzot, ka ar viņu vairs nevarēs vienoties, lai PSRS labā iegūtu vēl jaunas teritorijas. J.~Staļina cerības uz ilgu karu Rietumos var salīdzināt ar Lielbritānijas cerībām uz savu salu nesasniedzamību un Francijas cerībām uz Mažino līnijas drošību. Tā vietā lai kopīgi vērstos pret agresoru gan Rietumvalstu vadītāji, gan PSRS diktators cerēja, ka nelaime piemeklēs citus. Kurš gan toreiz varēja paredzēt, ka Polijas, Francijas, Lielbritānijas, Holandes, Beļģijas, Norvēģijas, Grieķijas, Dienvidslāvijas armijas tiks pa vienai metodiski sakautas, šo valstu (izņemot Lielbritāniju) rūpniecība sāks strādāt Vācijas labā un PSRS 1941.~gada vasarā atradīsies faktiski viena nacistisko uzbrucēju priekšā.

Pakts bija zināms kompromiss, kas Padomju Savienības pusei nesa gan zaudējumus, gan ieguvumus. Morāli zaudējumi tika ciesti gan padomju, gan starptautiskās sabiedrības acīs par strauju pagriezienu no antinacistiskās politikas, ko viennozīmīgi slavināja kreisā propaganda, uz izlīgumu ar Ā.~Hitleru. J.~Staļinam ar pilnām tiesībām revolucionāro ideālu nodevību varēja pārmest un arī pārmeta visi tie ideālistiski domājošie komunisti, kuri pēc 1937.--1938.~gadā PSRS notikušās plašās komunistu apslepkavošanas vēl uzskatīja viņu par komunistu. Ne velti PSRS pakts ar nacistisko Vāciju pasaules komunistiskajās partijās izsauca to biedru rindu sarukšanu. No šīm partijām, arī neko nezinot par slepeno protokolu, izstājās godīgi, principiāli un revolucionāri noskaņoti cilvēki, kuri nespēja samierināties ar cinisku, lai arī PSRS interešu diktētu vienošanos ar komunistu niknākajiem ienaidniekiem~--- nacistiem. Komunistiskajās partijās palika vai nu tie, kuri akli ticēja PSRS un Kominternes vadībai, bet faktiski~--- J.~Staļinam, tam, ka viss, ko viņš dara, nāk par labu komunisma uzvarai pasaulē, vai arī cilvēki, kuri vienkārši nevēlējās saskatīt, ka pats J.~Staļins un arī citi pēc staļiniskajām ``tīrīšanām'' dzīvi palikušie padomju komunistu vadītāji ārpolitikā rīkojās gluži tāpat kā viņu konkurenti~--- kapitālistisko valstu līderi. Trīsdesmito gadu beigās kā iekšpolitikā, tā arī ārpolitikā pastāvošos Padomju Savienības vadības orientierus nekādi nevar nosaukt par sociālistiskiem, atbilstošiem klasiskajiem marksistu priekšstatiem par sociālismu. Galvenais šajos orientieros bija jaunu teritoriju ieguve, lielvalstiskās ietekmes izplatība uz citām zemēm. Pragmatisms guva virsroku pār ideoloģiskajiem apsvērumiem.

Krievu vēsturnieks J.~Novopašins rakstīja, ka, ja sociālistiskais elements būtu noteicošais PSRS pirmskara valstisko interešu noteikšanā, tad padomju vadītājiem būtu jācenšas nevis palikt malā no konflikta, nevis kopā ar Vāciju jāpiedalās prettiesiskās agresīvās darbībās, bet šai konfliktā ``enerģiski jānostājas antifašistisko spēku pusē''. Tomēr, reālistiski vērtējot situāciju, jāsaka, ka 1939.~gadā principiāli marksistiska politika no PSRS interešu viedokļa diezin vai nestu tai panākumus. Kur gan bija tie ``antifašistiskie spēki'', kuru pusē PSRS varētu nostāties? Lielbritānijas, Francijas un Polijas valdībās tādu nebija. Arī Eiropas tautu vairākums, vēl neizjutis nacisma briesmu darbus uz savas ādas, diemžēl neieņēma skaidri antifašistisku pozīciju. Cīnīties vienatnē pret nacismu varētu uzdrošināties tikai avantūristiski noskaņoti komunistiski fanātiķi. Grūti spriest, kāda 1939.~gadā varētu būt tiešām sociālistiskas PSRS konkrētā ārpolitika un kādas tās sekas. Vēsturniekiem nākas vērtēt reālo PSRS vadības realizēto ārpolitiku, nevis tās pseidosociālistisko retoriku.

Lai nu kā, J.~Staļinam un viņa sekotājiem nācās aizstāvēt vienošanās ar Vāciju nepieciešamību gan demokrātiski, gan komunistiski noskaņoto cilvēku priekšā. Pavisam grūti to bija darīt pēc PSRS iebrukuma Polijā 1939.~gada 17.~septembrī un reālas iesaistīšanās pasaules pārdalē. Staļinisms diktēja arī padomju vēsturnieku uzskatus. Tikai ``pārbūves'' gados PSRS pirmo reizi radās iespēja izskatīt pirmskara un Otrā pasaules kara sākumposma notikumus un Polijas Republikas likvidāciju ne tikai caur tās valdošo aprindu kļūdainās un pat noziedzīgās politikas rezultātu prizmu, bet arī kā divu Polijas kaimiņvalstu vienošanos par tās sadali.

Taču visi staļiniskā režīma noziegumi, vispārliecinošākie to nosodījumi nespēj attaisnot Lielbritāniju un Franciju par to realizēto agresora ``nomierināšanas'' politiku, kā arī Poliju, (tāpat kā Baltijas valstis) par nevēlēšanos veidot efektīvu kolektīvās drošības sistēmu kopā, lai ar totalitāru, demokrātiskajai cilvēces daļai nepieņemamu, bet vienīgo valsti, kuras iesaistīšana šai sistēmā ļautu apturēt nacisma uzvaras gājienu jau 1939.~gadā. Morāli, staļinisko totalitārismu nosodoši argumenti nevar attaisnot citu valstu vadītāju politiskas kļūdas, politiku, kas neatbilda savas valsts un visas cilvēces interesēm. Ja runā par literatūrā izskanējušajiem pārmetumiem PSRS par to, ka tā ar Molotova-Ribentropa pakta noslēgšanu izraisīja Otro pasaules karu, tad jāprecizē, ka pirmos soļus šai virzienā ar agresora ``nomierināšanu'' veica Lielbritānija un Francija, PSRS tām tikai pievienojās. Pēc būtības, visas šīs valstis rūpējās, lai pašas izbēgtu no kara, nerēķinoties ar visas Eiropa drošību.

Tikai Lielbritānijas, Francijas un PSRS, (bet faktiski arī Polijas) kolektīvās drošības līgums varēja atturēt Ā.~Hitleru no Otrā pasaules kara uzsākšanas. 1989.~gadā tolaik vēl padomju vēsturniece V.~Parsadanova uz jautājumu, vai Polija bija vainīga par to, ka tāds līgums netika noslēgts, atbildēja, ka kopumā jau ir vainīga, taču ne lielākā mērā kā Minhenes līgumu parakstījušās Rietumu lielvalstis. Acīmredzot šim spriedumam ir jāpiekrīt.

Poļu vēsturnieks V.~Materskis gan raksta, ka Polijas realizētajai ``līdzsvara'', politikai, ``ar visu tās teorētisko, papīra raksturu, praksē nebija alternatīvas, vēl jo vairāk tāpēc, ka 30.~gados Polija bija tikai starptautiskās spēles priekšmets, bet nevis subjekts. Bez tās ne akmens nepaliktu no poļu ārpolitikas pamatpostulāta par neatkarīgas valsts eksistences saglabāšanu\dots{}'' Šādam secinājumam nevar piekrist. Pat 30.~gadu beigās Polijas vadītāji savu valsti nebūt neuzskatīja par ``starptautiskās spēles priekšmetu''. Ka tāda tā nebija, viņi arī pierādīja gan ar Cešinas sagrābšanu, gan ar ieguldījumu pretdarbībā Lielbritānijas, Francijas un PSRS savstarpējās palīdzības līguma noslēgšanai. Tāpēc arī 2019.~gadā baltkrievu vēsturniece L.~Ņikitenkova atzinusi, ka tieši Polijas vadības nepareizais starptautiskās situācijas vērtējums, Polijas lomas pārvērtēšana pirmskara Eiropas drošības starptautiskajā sistēmā, kā arī Lielbritānijas un Francijas realizētā piekāpšanās politika Vācijai, nostādīja pasauli kara priekšā.

Polijas realizētajai ``līdzsvara'', Rietumvalstu, bet pēc 1939.~gada 23.~augusta arī PSRS realizētajai agresora ``nomierināšanas'' politikai uz citu valstu rēķina, arī mazo valstu neitralitātes politikai vienmēr pastāvēja alternatīva~--- kolektīvās drošības politika. Cita lieta, ka Polijai, un ne tikai tai, šī politika nebija pieņemama, ja tā bija jāveido kopīgi ar PSRS.

Ja arī mēs uzskatītu, ka ``30.~gados Polija bija tikai starptautiskās spēles priekšmets, bet nevis subjekts'', tad uzreiz arī jāpasaka, ka tieši ``sanācijas'' režīms nes lielu atbildību par to, ka Polija no starptautisko attiecību subjekta kļuva par to objektu. Tās valdība nespēja pacelties no šauri šķiriskajām interesēm līdz patiesi nacionālajiem, pat vispārcilvēciskajiem uzdevumiem. Polija baidījās zaudēt suverenitāti, nevēlējās pieļaut padomju karaspēka ieiešanu tās teritorijā, tās rezultātā iespējamo ``poļu zemnieku komunizāciju''. Polija (tāpat kā PSRS) bija slima ar 1920.~gada sindromu. Šīs slimības rezultātā pat tad, kad sabruka ``līdzsvara'' politika un vairs nebija iespējams balansēt starp Vāciju un PSRS, Polija noraidīja iespēju sadarboties ar pēdējo. Tā deva savu ``ieguldījumu'' iespējamā Lielbritānijas, Francijas un PSRS līguma izjaukšanā. Liela daļa patiesības ir krievu žurnālista S.~Kņazeva secinājumā: ``Paradoksāli, bet fakts: ja neskaita pašu Vāciju, tad atrast citu valsti, kura vairāk izdarīja Otrā pasaules kara sākšanas labā nekā Polija, nav iespējams. Kopumā, kara sākumu Eiropā jau sen laiks skaitīt nevis no Hitlera uzbrukuma Polijai, bet kopš vācu-poļu agresijas pret mierīgo Čehoslovākiju 1938.~gadā.''

Vajadzētu paturēt prātā, ka pēc Molotova~--- Ribentropa pakta noslēgšanas viena agresīva valsts~--- nacistiskā Vācija uzbruka otrai agresīvai valstij~--- ``sanācijas'' Polijai. Poļu vēsturnieki gan to nevēlas atcerēties. Poļu vēsturnieks P.~Madajčiks raksta, ka poļu vēsturiskajā atmiņā ar 23.~augustu datētais un 24.~augustā Maskavā parakstītais Molotova-Ribentropa pakts ir vissvarīgākais nesenās vēstures atmiņu objekts līdzās 1945.~gada Jaltas konferencei, Katiņas slaktiņam, 1944.~gada Monte Kassino (\ittxti{Monte Cassino}) ieņemšanas kaujai Itālijā, kur krita daudzi poļu karavīri, un 1944.~gada Varšavas sacelšanās cīņām. Jāatzīst, ka poļu vēsturnieka minētā ``poļu vēsturiskā atmiņa'' ir visai selektīva, saturoša nezinātniskus atzinumus, tajā nav nedz J.~Pilsudska 1934.~gadā noslēgtā līguma ar nacistisko Vāciju, nedz Minhenes konferences un tās rezultātā notikušas Polijas agresija pret kaimiņvalsti Čehoslovākiju, nedz Polijas ātrās sagrāves 1939.~gadā, poļu armijas izvešanas no PSRS uz Irānu un divu gadu stāvēšanas angļu komunikāciju apsardzes dienestā, poļu īstenoto ebreju slaktiņu Kelcē, Jedvabnē u.c., AK (\pltxti{Armia Krajowa}) sadursmju ar padomju partizāniem, poļu vienprātīgi veiktās vācu un ukraiņu iedzīvotāju piespiedu pārvietošanas kara rezultātā. Pēc tā visa objektīvs notikumu vērotājs nevarēs piekrist ``poļu vēsturiskajā atmiņā'' palikušajam vērtējumam.

Lai nu kā, Polijas stāvoklis pēc Molotova-Ribentropa pakta parakstīšanas kļuva traģisks. Polijas Ārlietu ministrijas Preses nodaļas pirmā reakcija uz Molotova-Ribentropa pakta noslēgšanu bija paziņojums, ka valsts drošība ar to netiek aizskarta, jo Polijai ar PSRS ir 1932.~gada neuzbrukšanas līgums. ``Pēc būtības maz kas ir mainījies'',~--- Francijas sūtnim Varšavā L.~Noelam paziņoja J.~Beks. Kad ziņojumu par to redzēja Francijas sūtnis P.~Nadžiars, viņš secināja: ``Nevar iedomāties lielāku muļķību''. Arī Varšavas radio paziņojums par Molotova-Ribentropa pakta noslēgšanu skanēja: ``Dabīgi, ka līguma noslēgšana neatstās nekādu ietekmi uz Polijas lēmumiem Eiropas politikas jautājumos. Polijai nekad nav bijis vēlēšanās kaut kādi sadarboties ar Krieviju un tā atkārto savu nemainīgo ``nē''. Tā uzskata, ka Padomes nevar būt derīgs partneris Eiropas politikā. Pakta noslēgšana apstiprina Polijas ārpolitikas tēzi: Austrumos Eiropa beidzas ar Poliju. Tas viss nostiprina Polijas prestižu''. Turpmāk J.~Beks Polijas diplomātiskajām pārstāvniecībām citās valstīs lika Molotova-Ribentropa paktu komentēt sekojoši:

\begin{enumerate}
\item Tas nekādā ziņā nemaina Polijas faktisko pozīciju sakarā ar to, ka Polija nekad nerēķinājās ar Padomju palīdzību.

\item Tas nekādi nemaina Polijas politisko līniju, tāpat kā netraucē tās savstarpējās attiecības ar sabiedrotajiem.

\item Tas ir Padomju divkosīgās spēles pierādījums~--- tā acīmredzami izvairās no pilnas iesaistīšanās kādā no buržuāzisko valstu pusēm, vienlaikus labprāt redzot Eiropas kara iespēju.
\end{enumerate}

Pirmajā acu uzmetienā šķiet, ka J.~Beka instrukcijā situācija bija vērtēta it kā adekvāti. Turpmāk Polijai, kura nerēķinājās ar PSRS palīdzību, to arī nesaņēma; pakts tai netraucēja veidot savstarpējās attiecības ar sabiedrotajiem (tikai tās cerēto palīdzību nesniedza, kas ļāva arī PSRS sagrābt savu ``kumosu'' no Polijas teritorijas); PSRS tiešām izvairījās ``no pilnas iesaistīšanās kādā no buržuāzisko valstu pusēm'', domājot par savām interesēm. Taču to, kas notiks ar Poliju, tās ārlietu ministrs nespēja paredzēt. Mūsdienās analizējot J.~Beka vērtējumus, poļu vēsturnieki nešaubās par to, ka Polijas valdība nepietiekami novērtēja Molotova-Ribentropa līguma nozīmi, bet J.~Beks to atstāja gandrīz bez ievērības. Ministra secinājumu, ka ``tas nekādā ziņā nemaina Polijas faktisko pozīciju'' viņi vērtē kā principiālu kļūdu. Stiprais pārsteigums, ar kādu vēlāk~--- 17.~septembrī J.~Beks sagaidīja ziņu par Sarkanās armijas ieiešanu Polijas teritorijā, bija tiešs šī 23.augusta līguma kļūdainā vērtējuma rezultāts.

Turpretī Polijas sabiedrībā ziņas par paktu izsauca īstu sprādzienu. Prese sāka runāt par ceturto Polijas dalīšanu, izklāstot tās noteikumus visai tuvu slepenajā protokolā paredzētajam variantam. Tomēr presē skanēja arī V.~Sikorska izteiktās optimistiskās domas, ka pakts kā ``divpusēja politiska krāpniecība'' nebūs ilglaicīgs, tas ``nenoteiks PSRS un Vācijas izturēšanos nākamajos notikumos''. Taču labvēlīgam risinājumam, savas valdības cerībām gan uz Polijas armiju, gan sabiedrotajiem liela daļa tautas neticēja.

Uz poļu oficiālās preses saceltā uzbudinājuma par padomju-vācu pakta noslēgšanu fona naktī no 23.~uz 24.~augustu Polijā sākusies daļējā mobilizācija, baumas par Vācijas karaspēka ieiešanu Dancigā vismaz daļā poļu iedzīvotāju izraisīja paniku: nekavējoties veidojās rindas pie pārtikas veikaliem, pāris stundu laikā tika izpirkti milti, putraimi, cukurs un sāls, naudas noguldītāji piepildīja banku un krājkašu telpas. Daļa varšaviešu metās ārā no pilsētas, pārpildot vilcienus, jau tā pilnus ar mobilizētajiem. Pat uz ielām notika automobiļu, motociklu, velosipēdu un zirgu rekvizīcijas, kuru barus veda pa pilsētas ielām. Izrādījās, ka poļu militāristi mobilizācijai nebija gatavi. Karaspēka apgādē bija lieli trūkumi. Rekrūši, saņēmuši iesaukšanas pavēstes un ieradušies savu daļu izvietošanās vietās, tās vairs tur nesastapa, jo tās bija jau devušās uz citu vietu. Daudzviet valdīja haoss un nespēja orientēties. Tam visam klāt nāca vēl neziņa, kā Vācijas uzbrukuma gadījumā rīkosies PSRS.

Polijas armija rīkojās atbilstoši plānam ``Rietumi'', izvietojot lielāko karaspēka daļu valsts rietumos pie Vācijas robežas. Taču bija skaidrs, ka ar to nepietiek, jo saistībā ar Čehoslovākijas okupāciju dienvidos bija radusies vēl jauna 300~km gara fronte. Viss Polijas aizsardzības plāns būtībā balstījās uz aprēķinu, ka Rietumu sabiedrotie drīz pret Vāciju radīs otru fronti.

Vācijā uzbrukuma Polijai sagatavošanas darbi bija jau pabeigti. 1939.~gada 25.augustā Ā.~Hitlers Lielbritānijas sūtnim N.~Hendersonam ar uzspēlētu sašutumu paziņoja, ka visi starpgadījumi ir ``no poļu puses izsaukti'' un ``Vācija ir izšķīrusies pie jebkādiem apstākļiem novērst šo maķedonisko (acīmredzot bija domāti Eiropā plaši izplatītie priekšstati par maķedoniešiem kā laupītāju tautu~--- V.Š.) stāvokli pie savas austrumu robežas.'' Tai pat 25.augustā Ā.~Hitlers izdeva pavēli Vācijas armijai 26.~augusta rītā veikt iebrukumu Polijā, tomēr divi apstākļi neļāva plānam īstenoties kā bija iecerēts.

Lielbritānija arī vēl pēc Ribentropa-Molotova pakta noslēgšanas turpināja savu uz kara novēršanu virzīto politiku un 25.~augustā apstiprināja 6.~aprīlī Lielbritānijas un Polijas parakstīto savstarpējās palīdzības līgumu. Ar to tā vēlreiz brīdināja Vāciju, ka iesaistīsies konfliktā, ja Vācija uzbruks Polijai. Kā konstatējis profesors Ē.~Jēkabsons, ``Maskavā noslēgtais pakts vismaz paātrināja Polijas un Lielbritānijas militārās savienības līguma noslēgšanu 25.~augustā.'' Līgumā bija ierakstīts: ``Gadījumā, ja viena no Līgumslēdzējpusēm tiks iesaistīta kara darbībā citas Eiropas valsts agresijas rezultātā, otra Līgumslēdzējpuse nekavējoties sniegs kara darbībā iesaistītajai Līgumslēdzējpusei visu nepieciešamo atbalstu un palīdzību''. Taču kā izrietēja no līgumam pievienotā slepenā protokola, ar ``citu Eiropas valsti'' bija domāta tikai Vācija, bet ne kāda trešā valsts. (Iespējams, šāda traktējuma pieņemšanu izsauca jau 17.~aprīlī sakarā ar trīspusējo sarunu gatavošanu PSRS ārlietu tautas komisāra M.~Ļitvinova iesniegtais priekšlikums Lielbritānijas valdībai izskaidrot, ka solītā palīdzība Polijai nozīmē tās sniegšanu vienīgi agresijas gadījumā no Vācijas puses.) Vēlāk slepenā atruna ļāva Lielbritānijai nepieteikt karu Padomju Savienībai, kad Sarkanā armija 17.~septembrī pārkāpa Polijas robežu. Izrādījās, ka pēc būtības Lielbritānijas valdība ar Polijas un Lielbritānijas Savstarpējās palīdzības līguma slepenā protokola noslēgšanu bija ne tikai izvairījusies no nepieciešamības pasludināt kara stāvokli pret PSRS, bet arī sākusi būvēt pamatus Antihitleriskās koalīcijas izveidei nākamajos gados. Jāatzīmē, ka par šo slepeno protokolu neko nezināja PSRS vadība, kura uzskatīja, ka Lielbritānijas garantijas Polijai var iegūt arī pretpadomju raksturu. Padomju diktators J.~Staļins saglabāja savu neuzticības un naidīguma pilno izturēšanos pret Lielbritāniju.

Vācijas agresijas pret Poliju gadījumā Lielbritānijas valdība, saprotot, ka britu vēlētāji tai neatļaus īstenot ``jaunu Minheni'', bija gatava vismaz politiski nostāties Polijas pusē.

Tiesa, kā norādījis krievu vēsturnieks A.~Širokorads, 30.~gadu beigās un 40.~gadu sākumā Lielbritānijā tika izstrādāti dažādi politiskie varianti, iespējams, ļoti kompromitējoši tās vadību savas un ārvalstu sabiedrības acīs. Ne velti ievērojama daļa Lielbritānijas valdības dokumentu nav atslepenota vēl līdz šai dienai, kaut pēc likuma tas bija jāizdara pēc 30~gadiem. Taču vēsture rēķinās galvenokārt ar realitāti.

Tāpat kā 23.augustā cieta neveiksmi Lielbritānijas un Francijas mēģinājums ar savienības noslēgšanu ar PSRS iebiedēt Vāciju, 25.augustā cieta neveiksmi Ā.~Hitlera mēģinājums šokēt Lielbritānijas valdību ar Molotova-Ribentropa pakta noslēgšanu, cerības to atturēt no oficiālas savu garantiju Polijai apstiprināšanas.

Tai pat 25.~augustā Ā.~Hitlers saņēma ziņu no Itālijas diktatora B.~Musolini, ka Itālija nav gatava karam, tai trūkst tam nepieciešamo resursu. Abu ziņu rezultātā uzbrukums Polijai tika atlikts no 26.~augusta uz 1.~septembri.

Ar uzbrukuma termiņa pārcelšanu saistās vairākas interesantas epizodes. Vācijas uzbrukuma plāns paredzēja strauju virzīšanos Polijas iekšienē, lai neļautu pretinieka armijai atiet uz austrumiem. Uzdevuma izpilde bija atkarīga no divu atslēgas pozīciju robežjoslā savlaicīgas nonākšanas vācu rokās. Tās bija: dzelzceļa tunelis un stacija pie Jablunkovas pāreju (vācu~--- \detxti{Jablunka Pass bei Mosty}, čehu~--- \cstxti{Mosty u Jjablunkova}, poļu \pltxti{Jablunkovsky Przełęcz}) ar dzelzceļa tuneli Polijas dienvidrietumos 1938.~gadā sagrābtajā Cešinas apgabalā, kurai bija svarīga stratēģiska nozīme, tā kalpoja kā sava veida tilts no Čehijas ziemeļiem uz Polijas dienvidiem, un tilts pār Vislu pie Diršavas (\detxti{Die Weichselbrücke Dirschau}) ziemeļrietumos. Šo šauro vietu sagraušana varētu būtiski aizkavēt vācu karaspēka virzību. Tāpēc bija sagatavotas divas operācijas šo mezgla punktu ieņemšanai.

Atbilstoši uzbrukuma pavēlei, Ā.~Hitlers izdeva 25.~augustā, ap 30~cilvēku liela vācu izlūkdienesta vienība, daļēji tērpta poļu formās un daļēji civilapģērbā, kas bija Hāgas vienošanos par kara vešanu (1899--1907) pārkāpums, leitnanta H.~Hercnera vadībā saņēma pavēli naktī no 25.~uz 26.~augustu ieņemt Jablonkovas pāreju un nepieļaut tās uzspridzināšanu un atbrīvot no šķēršļiem dzelzceļu. Vienība pulksten 21-os atstāja savu bāzi un devās uz galamērķi. Pavēles atsaukums, kurš sekoja pēc stundas, to vairs nesasniedza, jo vienība darbojās kalnainā apvidū, kur raidstacijas nespēja uztvert pavēli par operācijas atsaukumu. Pēc pusnakts vienība šķērsoja Slovākijas~--- Polijas robežu, trīs stundas vēlāk ieņēma dzelzceļa staciju un tuneļa dienvidu galu, atbruņojot un ieslogot pārejas apsargus, pēc tam veltīgi gaidīja uz vācu pamatspēku ierašanos. Taču poļu sakarniekiem izdevās nosūtīt ziņu par vācu iebrukumu un poļu armijas vienības devās uz pāreju. Satraucies par vācu pamatspēku izpalikšanu, pēc pieciem rītā vienības komandieris ar lokomotīvi nosūtīja ziņnesi uz savu štābu. Bet bija jau par vēlu, pēc brīža šai apkārtnē izvietotie poļu spēki jau aplenca vācu vienību. Vācu diversantu grupa, nesagaidījusi agrāk apsolītos papildspēkus, bija spiesta atkāpties kalnos un ar diviem ievainotajiem atgriezās Slovākijā. (Citi dati~--- tikai komandierim ar dažiem saviem vīriem izdevās izbēgt, pārējie tika sagūstīti un atbrīvoti nākamajās dienās sekojušo poļu~--- vācu sarunu laikā.) Nepatīkamo gadījumu vācu puse izskaidroja ar ``pārskatīšanos'', kas notikusi mežainā apkārtnē, kur grūti izsekot robežai. Vācijas virsnieki atvainojās Polijas pusei par kādas ``neaprēķināmas [personas] izraisītu starpgadījumu'' (``\detxti{von einem Unzurechnungsfähigen verursachten Zwischenfall}''). Var piebilst, ka 1.~septembra rītā Jablonkovas tuneli poļi tomēr paguva uzspridzināt dažas minūtes pirms vācu armijas vienību ierašanās.

Neizdevās arī otrais pasākums pret tiltu Diršavā (vācu \detxti{Dirschau}, poļu \pltxti{Tczew}). Visla šeit veidoja Polijas austrumu robežu ar Dancigas brīvpilsētu. 2~km uz austrumiem jau uz Vācijas zemes atradās Marienburgas jeb Malborkas pils (vācu \detxti{Marienburg}, poļu \pltxti{Zámek w Malborku}). Tāpēc jau 1939.~gada maijā poļu karavīri tiltu katram gadījumam mīnēja. Vācieši plānoja pirms uzbrukuma Diršavai pieprasīt lokomotīvi, kurai Marienburgā kā parasti būtu jāpārņem tranzītā caur poļu koridoru ejošu vilcienu, tad poļu apkalpi nomainīt ar vācu, kura vācu kaujas grupu un sapierus ar vilcienu aizvestu līdz Diršavai, kur tie arī sagrābtu un atmīnētu tiltu, kamēr vācu triecienlidmašīnas veiktu uzlidojumu stacijai. Atšķirībā no Jablunkovas šeit naktī no 25.~uz 26.~augustu pēdējā brīdī izdevās atsaukt pāragro uzbrukuma pavēli. 31.~augusta pusnaktī poļu lokomotīve atstāja Diršavas staciju. Kā plānots, tās apkalpe Marienburgā tika arestēta un bija spiesta atdot savus formas tērpus. Vilciens atstāja Marienburgu ar vairāk kā stundas novēlošanos un devās uz Diršavu bruņuvilciena pavadībā, uz kura pievienošanu uzstāja H.~Himlers. Arestētas tika arī kāda pretī braucoša poļu vilciena apkalpe un poļu muitas ierēdņi, kuri pamanīja aizdomīgo parastā kravas vilciena konvoju. Taču kāds starp Marienburgu un Diršavu esošās Simonsdorfas (\pltxti{Szymankowo}) stacijas ierēdnis paguva kolēģus Diršavā brīdināt. Poļu garnizonam tika pavēlēts apturēt satiksmi pa tiltu un to uzspridzināt. Tas arī tika izdarīts. Tilta aizstāvju un ar vilcienu ieradušos vācu karavīru starpā sākās apšaude, kamēr vācu pikējošie bumbvedēji bombardēja staciju. Pilsēta kapitulēja 2.~septembra rītā. Atriebjoties par neizdevušos pasākumu, 13 dzelzceļa darbiniekus kopā ar diviem ģimenes locekļiem un pieciem vai sešiem muitas ierēdņiem jau 1.~septembra naktī nogalināja vācu nacistu trieciennieki.

1939.~gada 28.~augustā īsi pirms pusnakts Polijas vācu mazākuma pārstāvji Vācijas militārās izlūkošanas~--- abvēra uzdevumā organizēja bumbas sprādzienu Tarnovas (\pltxti{Tarnόw}) dzelzceļa stacijā. Tajā dzīvību zaudēja 22~cilvēki. Tā kā pilsēta atradās nevis pie robežas, bet dziļi valsts iekšienē, sprādziens izraisīja vietējo iedzīvotāju neizpratni un sašutumu. Policijai izdevās arestēt daļu no atentāta sarīkotājiem. Izvirzās jautājums, kādēļ sprādziens tika īstenots, jo tas nekalpoja nekādiem militāriem mērķiem. Viens no poļu vēsturniekiem ir izteicis pieņēmumu, ka atentāta organizētāji apzināti vēlējās izraisīt poļu iedzīvotāju naidu pret vācu mazākumu, tā provocējot uzbrukumus tam. Taču pastāv arī iespēja, ka tā bija vienkārši vācu izlūkdienesta kārtējā neveiksme kara sākumposmā.

Pēdējās pirmskara dienās Ā.~Hitlers turpināja pūliņus ``iedzīt ķīli'' starp sabiedrotajiem, panākt, lai Lielbritānija un Francija tomēr pamestu Poliju savam liktenim. Viņa mērķis bija, ātri sakaujot Poliju, izvairīties no kara divās frontēs. Bija jādod iegansts Lielbritānijas valdībai atteikties no garantijām Polijai, savstarpējās palīdzības līguma ar Poliju. Tāpēc Polijai tika piedāvāti jauni, it kā mīkstināti vienošanās noteikumi par Dancigu un ``poļu koridoru''. (Ā.~Hitlers prasīja vienīgi Dancigas atdošanu Vācijai un plebiscīta sarīkošanu ``koridora'' teritorijā pēc gada, kad kaislības būs norimušas.'') Taču tā bija tīra krāpniecība, domāta iespaida radīšanai pasaules un īpaši vācu tautas acīs, ka fīrers dara visu miera saglabāšanai. Ā.~Hitlers par to skaidri izteicās savu apkārtējo priekšā. Viņš sprieda, ka Polija, ievērojot agrāko pieredzi, nekādiem priekšlikumiem nepiekritīs. Vērmahta sauszemes spēku Ģenerālštāba priekšnieks ģenerālis F.~Halders 28.augustā savā dienasgrāmatā ierakstīja, ka, iespējams, angļi, rēķinoties ar savām saistībām Polijas priekšā, vedīs ``šķietamu karu'' un fīrers par to ``neapvainosies''. Ģenerālis pareizi prognozēja Lielbritānijas politiku.

Notika dzīva apmaiņa ar paziņojumiem starp Berlīni un Londonu. Lielbritānija, par spīti visam, vēlējās panākt Vācijas un Polijas sarunas domstarpību Dancigas un ``poļu koridora'' jautājumā apspriešanai. Kā, saskatot paralēles ar Čehoslovākijas krīzi 1938.~gadā, uzskata Dž.~Grenvils, Lielbritānija un Francija būtu piekritušas atzīt jebkuras Polijas teritoriālas piekāpšanās Vācijai, ja tā nelietotu varu to iegūšanai. Lielbritānijas valdības sēdē 26.~augustā tās sūtnis Vācijā N.~Hendersons aizstāvēja domu, ka ``Mūsu garantiju Polijai vērtība ir tajā apstāklī, lai dotu Polijai iespēju nonākt pie noregulējuma ar Vāciju''. Tātad Rietumvalstīs netika izslēgta miera saglabāšana uz Polijas rēķina.

Taču J.~Beks acīmredzot uzskatīja, ka Polijas un Francijas 1921.~gada līgums un Lielbritānijas 1939.~gada 30.~martā Polijai piešķirtās garantijas un Polijas--Lielbritānijas 25.~augusta savienības līgums atturēs Ā.~Hitleru no uzbrukuma Polijai, lai savukārt nesaņemtu triecienu no Lielbritānijas un Francijas, bet, ja šis uzbrukums tomēr notiktu, Polija nepaliks vienatnē, tai tiks sniegta ievērojama palīdzība no Rietumu sabiedroto puses, kopā ar kuriem tā uzvarēs karu pret nacistisko Vāciju.

29.~augustā kārtējā apspriedē I.~Moscickis, E.~Ridzs-Smiglijs un J.~Beks konstatēja, ka jārēķinās vai nu ar Vācijas uzbrukumu Polijai, vai starptautisku kompromisu uz Polijas rēķina. Sabiedroto pārstāvji tika informēti, ka Polijā paredzēts izsludināt vispārēju mobilizāciju, taču pēc to lūguma tā pagaidām tika atlikta, kā rezultātā kara sākuma brīdī Polijas bruņotie spēki nebija pilnā kaujas gatavībā. Tās pašas dienas vakarā Ā.~Hitlers Lielbritānijas valdībai paziņoja, ka viņš nākamajā dienā Berlīnē sagaida Polijas pārstāvi ar attiecīgām pilnvarām, ar kuru varētu noslēgt Vācijas~--- Polijas vienošanos. J.~Beks atteicās turp doties, pēc vācu publicista un vēsturnieka J.~Festa vārdiem, baidoties, ka viņu tur sagaidīs Austrijas federālā kanclera K.~Šušniga un Čehoslovākijas prezidenta E.~Hahas liktenis, kad attiecīgi 1938.~gada februārī un 1939.~gada martā Ā.~Hitlers piespieda viņus atdot Vācijai varu pār savām valstīm. Lielbritānijas sūtnis Varšavā H.~V.~Kennards ziņoja uz Londonu par J.~Beka paziņojumu, ka ``uz Berlīni \citespace{} viņš nebrauks, jo nevēlas pārdzīvot to, ka pārdzīvoja prezidents Haha.'' Sarunu vietā Polija izsludināja vispārēju mobilizāciju.

Poļu vēsturnieki, izteikdami kritiku par J.~Beka politiku, reizē atzīmē arī vienu viņa vēsturisku nopelnu~--- kopā ar pārējo pirmskara Polijas valdību viņš stingri pateica ``nē'' Ā.~Hitlera ultimātam, ar to nedodot Lielbritānijas un Francijas valdībām īstenot kādu jaunu Minhenes variantu, bet liekot tām pieteikt karu nacistiskajai Vācijai. Jau minētais Nacionālā piemiņas institūta prezidents J.~Kurtika uzsvēris, ka tieši J.~Beks ar savu diplomātiju galu galā iepiņķēja Hitleru karā, kas norisinājās divās frontēs uzreiz un galu galā iznīcināja nacistisko ``Tūkstošgadu reihu''. T.~Šarota J.~Beku pieskaita pie antihitleriskās koalīcijas dibinātājiem un runā par viņa pienesumu arī sabiedroto uzvarā 1945.~gadā. Pēc šis grāmatas autora uzskatiem poļu vēsturnieka apgalvojuma pirmajai pusei vismaz daļēji var piekrist, taču no uzvaras pār nacismu viedokļa viņa pienesums tai tomēr jāvērtē ar negatīvu zīmi.

Kad Polija atteicās no sarunām Berlīnē un izsludināja vispārēju mobilizāciju, Ā.~Hitlers varēja apvainot Poliju par atteikšanos no ``mierīga noregulējuma'' un turpināt pūliņus pierunāt Rietumu lielvalstis atteikties no palīdzības tai. (Tomēr kara pirmajās dienās atklājās, ka vismaz vārdos Rietumvalstis pamest Poliju vairs nevarēja.) 31.~augustā Ā.~Hitlers izdeva pavēli vērmahtam uzbrukt Polijai 1.~septembra rītā. Vācija, iebrūkot Polijā, pārkāpa 1919.~gada Versaļas miera noteikumus un veselu rindu citu starptautisku vienošanos.

Vācijas vadībai bija jāatrod formāls casus belli (latīņu~--- \latxti{iegansts karam}). Ā.~Hitleram bija gluži vienalga, vai viņa sagādātais \strong{uzbrukuma Polijai iegansts} kādu pasaulē pārliecinās. Vienīgais uz ko viņš varēja cerēt, bija sakarā ar it kā notikušo Polijas provokāciju Rietumu valstu atturēšana no savu garantiju pasludināšanas par spēkā esošām un kara pieteikšanas Vācijai. Uzbrukuma naktī Ā.~Hitlers bija pavēlējis inscenēt vairākus ``robežkonfliktus''. Vācijas drošības dienests sagatavoja veselu virkni provokāciju visā Vācijas~--- Polijas robežas garumā. Kara iegansta sagādāšanas plānošanu un sagatavošanu uzņēmās Vācijas drošības policijas (SD) vadītājs R.~Heidrihs.

Vienā no variantiem bija iecerēts, ka par uzbrukuma ieganstu kalpos poļu soļi pret nacistu vadītāju Dancigā A.~Forsteru, kuru 23.augustā ievēlēja par Dancigas brīvpilsētas ``valsts vadītāju'' (``\detxti{Staatsführer}''). Realizēts gan tika cits scenārijs.

Silēzijā tika izraudzīti trīs punkti par mērķi šķietamajiem poļu uzbrukumiem. Viens no tiem bija Gleivicas (vācu \detxti{Gleiwitz}, poļu \pltxti{Gliwice}) radiostacija. Tika nolemts, ka speciāli sagatavota vācu policijas darbinieku un SS vīru grupa sagrābs netālu no Polijas robežas esošo radiostaciju, noraidīs ziņu par Polijas armijas iebrukumu Vācijā un aicinājumu tajā dzīvojušajiem poļiem sacelties.

Pēc 25.~augusta uzbrukuma pavēles saņemšanas viņi jau šķērsoja Polijas robežu, no kuras bija jānotiek uzbrukumam, kad pēdējā brīdī tomēr tos paguva atsaukt atpakaļ. Arī naktī no 31.~augusta uz 1.~septembri ne viss norisa kā bija plānots. Daļa radiostacijai tuvāko māju iedzīvotāji daļēji bija evakuēti, pārējiem ar lielu troksni, šaudīšanos un komandsaucieniem poļu valodā ``uzbruka'' H.~Heidriha ļaudis. Taču pārējā pilsēta no tā visu gandrīz neko nemanīja. Neinformēto raidstacijas personālu uzbrucēji viegli pievarēja, taču izrādījās, ka no Gleivicas radioraidītāja jau no 1935.~gada netika pārraidīts vairs neviens šeit uz vietas sagatavots raidījums un ``iebrucēji'' tikai ar pūlēm sameklēja kādu mikrofonu, kurā varēja nolasīt savu iepriekš sagatavoto paziņojumu. Ir pat ziņas, ka ar lielām pūlēm izdevās noraidīt tikai vienu frāzi: ``Gleivicas raidstacija atrodas poļu rokās.'' Pie tam raidītāja jauda bija maza, to dzirdēja tikai pilsētas apkārtnē un H.~Heidrihs Berlīnē pie radioaparāta velti gaidīja ziņas no Gleivicas.

Pēc uzbrukuma raidstacijā bija jāpaliek ``pierādījumiem'', ka uzbrucēji tik tiešām bijuši poļi. No tuvējās koncentrācijas nometnes uz Gleivicu atveda uz nāvi notiesātus kriminālnoziedzniekus, kurus ietērpa poļu armijas formā un izdarīja injekcijas, pēc kurām tie zaudēja samaņu. Viens līķis tika atstāts radiostacijas telpās, pārējos nogalināja un pameta tuvējā apkārtnē it kā notikušajā kaujā kritušus.

Vācu diversantu grupas vadītājs A.~Neijoks pēc kara Nirnbergas procesā liecināja: ``Apmēram 1939.~gada 10.~augustā Heidrichs, \detxti{Sipo} (\detxti{Sicherheitspolizei}~--- drošības policija) un \detxti{SD} (\detxti{Sicherheitsdienst}~--- drošības dienests) šefs personīgi man pavēlēja sagatavot viltus \strong{uzbrukumu radiostacijai Gleivicā} Polijas robežas tuvumā un likt visam tā izskatīties, ka uzbrucēji bija poļi \citespace{} 31.~augustā dienas vidū es saņēmu kodētu pavēli veikt uzbrukumu tai pat vakarā pulksten astoņos. Bez tam, šī uzdevuma izpildei es saņēmu uzdevumu griezties pie Millera (\detxti{Geheime Staatspolizei}, arī \detxti{Gestapo}~--- slepenpolicijas vadītājs) pēc ``konserviem'' (Tā kodeti tika saukti it kā incidentā ``bojā gājušie'' vācu koncentrācijas nometņu ieslodzītie) \citespace{} Mēs ieņēmām raidstaciju kā bija pavēlēts, trīs vai četras minūtes bijām ēterā, dažas reizes izšāvām no pistolēm un aizgājām\dots{}''

1.~septembra rītā Vācijas ziņu birojs (\detxti{Deutches Nachrichten Büro}) izplatīja paziņojumu ``Poļi izdarījuši uzbrukumu radiostacijai Gleivicā''. Tajā bija teikts: ``Breslava. 31.~augusts. Šodien ap pulksten astoņiem vakarā poļi uzbruka un sagrāba radiostaciju Gleivicā. Ar spēku iebrukuši radiostacijas ēkā, viņi paguva griezties [pie klausītājiem] poļu un daļēji vācu valodā. Taču pēc dažām minūtēm tos sagrāva policija, kuru bija izsaukuši radioklausītāji. Uzbrukums radiostacijai acīmredzot bija signāls vispārējam poļu partizānu uzbrukumam Vācijas teritorijai. Kā izdevās noskaidrot, gandrīz vienlaikus poļu partizāni šķērsoja Vācijas robežu vēl divās vietās. Tās bija labi apbruņotas vienības, kuras acīmredzot atbalstīja poļu regulārās daļas. Drošības policijas apakšvienības, kas apsargā robežu, ir uzsākušas kauju ar iebrucējiem. Sīvas sadursmes turpinās.''

Ā.~Hitlers griezās pie vācu tautas, paziņojot, ka Polija veikusi uzbrukumu Vācijas teritorijai un no tā brīža Vācija atrodas ar Poliju kara stāvoklī. Savā runā reihstāgā Ā.~Hitlers (Skat. dokumentu Nr.~19.) runāja par pavisam 14 ``šausmu darbiem'', kurus poļu naktī paveikuši uz robežas, ``starp tiem trijiem visai smagiem''. Tikai pēc kara jau Vācijas Federatīvajā Republikā tika noskaidrota patiesībā par šiem notikumiem. Polijā tie ieguva nosaukumu ``Gleivicas provokācija'' (\pltxti{Prowokacja Gliwicka}), 70.~gadu beigās par to tika uzņemta filma ``\pltxti{Operacja Himmler}'' (Operācija Himlers).

Ir saprotams, ka jautājums par poļu ārpolitikas vērtējumu laika posmā no 1933.~līdz 1939.~gadam ir kļuvis par poļu vēsturnieku asu diskusiju objektu. Jau XX gadsimta 80.~gados daudzi poļu vēsturnieki sāka aizstāvēt Polijas valdības ``līdzsvara politiku'', velkot paralēles ar to līdzsvara politiku, kuru Eiropā tradicionāli ilgstoši realizēja Lielbritānija. Toties krievu pētnieks S.~Morozovs savā monogrāfijā par Polijas un Čehoslovākijas attiecībām (\rutxti{Польско-чехословацкие отношения. 1933--1939. Что скрывалось за политикой ``равноудаленности'' министра Ю.~Бека. Москва, 2004}) ir parādījis, ka daudzi mūsdienu Polijas historiogrāfijas pārstāvji (S.~Gregorovich, M.~Kamiński, M.~Zacharias) vēlas notušēt ``sanācijas'' režīma agresīvo pusi, noklusēt tos plānus, kurus tā laika poļu vadoņi vēlējās īstenot uz Čehoslovākijas un citu kaimiņzemju rēķina. S.~Morozovs rāda, ka J.~Pilsudskis un viņa sekotāji faktiski realizēja avantūristisku politiku, ar mērķi atjaunot Žečpospolitas kādreizējo varenību. Kā atzīmējis vācu profesors M.~Berndts: ``Ne gluži bez pamata poļi komunisti 1945.~gadā, lai distancētos no pirmskara Polijas, par ``Otrās republikas eksperimentu'' runāja kā par ``grandiozu kļūdu''.

No otras puses, jāatzīst, ka ``sanācijas'' vadoņus grūti vainot par to, ka viņi neizdarīja visu iespējamo, lai radītu Eiropā kolektīvās drošības sistēmu, kas varētu reāli savaldīt lielāko agresoru~--- nacistisko Vāciju. Polija bija iespiesta starp divām varenākām valstīm, uz kuru draudzību tā nevarēja cerēt, jo jau pats neatkarīgas Polijas valsts rašanās un pastāvēšanas fakts, vēl jo vairāk tās apstrīdamajās robežās, izraisīja abu kaimiņvalstu vadītāju vēlmi līdz ar visu pastāvošo Versaļas sistēmu likvidēt šo pēc Pirmā pasaules kara radušos valsti ar lielām pretenzijām. Pat ja Polijā pirms kara valdītu demokrātisks režīms, diezin vai tas palīdzētu tai izglābties no okupācijas.

% page 728


\section{Dokumenti 2.~nodaļai}

\subsection{Izraksti no Polijas delegācijas\footnote{Polijas delegācija Parīzes Miera konferencē darbojās kopš tās darba sākuma. Tā bija izveidota uz Poļu nacionālās komitejas bāzes. Šo 1917.~gadā Parīzē izveidoto komiteju Rietumu valstis bija atzinušas par atdzimstošā Polijas valstiskuma pārstāvi. Tā sastāvēja no Poļu nacionāldemokrātiskās partijas vadītājiem. Tādējādi konferencei iesniegtās notas Polijas robežu jautājumā reprezentēja nacionāldemokrātu uzskatus, tādēļ J.~Pilsudskis un valdība, dažas notā izteiktās prasības neatbalstīja. Tomēr prasība par Galīciju un vēl dažām teritorijām bija kopīga gan nacionāldemokrātiem, gan J.~Pilsudskim un viņa domubiedriem.} notas Miera konferencei Polijas rietumu robežas jautājumā (Parīzē, 1919.~gada 28.~februārī)}

Polijas valsts atjaunošana ir jāuzskata par:

\begin{enumerate}
\item taisnīguma atjaunošanas un atmaksas aktu par sadalītājvalstu noziegumiem (1772--1793--1795);
\item poļu nacionālo spēku, kuri attīstījušies par spīti sadalītājlielvalstu destruktīvajiem centieniem, attīstības rezultātu un Austroungārijas sabrukuma, Krievijas impērijas sabrukuma un Vācijas sakāves konsekvenci;
\item nepieciešamību izveidot spēcīgu un patiesi neatkarīgu valsti, kas atrastos starp Vāciju un Krieviju.
\end{enumerate}

No pirmā punkta izriet, ka par atjaunotās Polijas valsts robežu noteikšanas pamatu ir jākalpo Polijas teritorijai pirms pirmās sadalīšanas (1772.~g.).

Otrais punkts liek veikt zināmas izmaiņas vecajās robežās, iekļaujot Polijas teritorijā dažas zemes, kuras atrodas ārpus tās 1772.~gada robežām, taču kurās poļu iedzīvotāji ir apliecinājuši savu dzīvotgribu, vai atstājot ārpus Polijas valsts tās 1772.~gada teritorijas zemes, kurās poļu tautas spēki pēc dalīšanām ir ievērojami samazinājušies.

Trešā punkta konsekvence ir zināmas izmaiņas vecajās robežās, piemērojot tās jaunās valsts ekonomiskajām prasībām un tās aizsardzībai.

Polijas rietumu zemes, ko dalīšanu laikā sagrāba Austrija un Vācija, ir sekojošas:
\begin{enumerate*}[label=\arabic*)]
\item Galīcija,
\item Poznaņas apgabals,
\item Rietumprūsija, izņemot Štumas (\detxti{Stuhm}) apriņķi, Sušas (\detxti{Rosenberg}) apriņķi, kā arī Kvidzinas (\detxti{Marienwerder}) apriņķa austrumu daļu,
\item Austrumprūsijas Varmijas (\detxti{Ermland}) daļa.
\end{enumerate*}

Galīcija ir sens apgabals Polijas valstī, kurai tā piederējusi kopš vissenākajiem laikiem līdz pat 1772.~gadam. Izņēmums attiecībā uz tās austrumu daļu ir laika posms 1091.--1340.~gadā.

Galīcijā ir 8~025~675 iedzīvotāju (1910), kas dzīvo 78~490~km$^{2}$ teritorijā, un saskaņā ar pēdējās austriešu tautas skaitīšanas datiem 3~200~000 no viņiem lieto ukraiņu, bet 4~700~000~--- poļu valodu.

Ukraiņu valodā runājošie iedzīvotāji apdzīvo apgabala austrumu daļu un ir 58,6\% no turienes iedzīvotāju kopskaita. Pēdējos piecdesmit gados viņu vidū iezīmējušies divi nacionāli virzieni, kuri ir vērsti pret poļu varu: viens no tiem krievu propagandas iespaidā centās panākt ukraiņu valodas aizstāšanu ar krievu valodu, otrs, kuru sākumā atbalstīja Vīnes valdība, bet vēlāk~--- Berlīnes valdība, tiecās izveidot ukraiņu nāciju. Pēdējo divdesmit gadu laikā, pateicoties vācu iespaida pieaugumam Viduseiropā un Austrumeiropā, panākumus guvusi ukraiņu kustība.

Būdama ierocis vācu un poļu konfliktā, ukraiņu kustība ir ieguvusi izteikti negatīvu raksturu. Pretēji citām nacionālās atbrīvošanas kustībām, tā ļoti maz paveikusi intelektuālajā un iedzīvotāju ekonomiskās organizācijas jomā, bet savu enerģiju izlietojusi gandrīz vienīgi politiskajās cīņās.

Jāuzsver, ka 91\% ukraiņu iedzīvotāju ir sīkzemnieki vai laukstrādnieki. Tas izskaidro apstākli, kāpēc pat austriešu varas laikā visas Galīcijas administrācija atradās poļu rokās; zemē nebija nekādu citu uz pārvaldes funkcijām spējīgu faktoru.

Tāds ukraiņu iedzīvotāju stāvoklis Galīcijā ne ar ko nepamato ukraiņu kustības vadoņu pretenzijas uz Austrumgalīcijas atraušanu no Polijas valsts. Tajā pašā laikā ukraiņu tautības tiesības ir neapstrīdamas, un nav nekādu šaubu, ka Polija tās pilnībā ievēros.

Pilnīgi droši var paredzēt, ka brīdī, kad izbeigsies Polijai naidīgā austriešu--vācu ietekme, poļiem un ukraiņiem būs labas attiecības.

Galīcijai ir īpaši svarīga vieta pārējo Polijas reģionu vidū dēļ tās dabas bagātībām. Galīcijas rietumos ir svarīgi ogļu un sāls krājumi, bet austrumos~--- nafta un potašs. Tās dienvidaustrumos Polijai ir robeža ar Rumāniju, kuras tautu ar poļiem saista tradicionālas draudzības saites un kopīgas intereses. Atbildīgajā un grūtajā stāvoklī, kādā abas valstis atrodas, tās ir spiestas tiekties pēc ciešas saprašanās un harmoniskas savstarpējas sadarbības.

Sākot ar Galīcijas un reizē ar to arī Polijas dienvidu robežu, Polijas valsts robeža iet gar Karpatiem, kuri kopā ar Baltijas jūras piekrasti veido Polijas vislabākās dabiskās robežas. Šī robeža atbilst vecajai 1772.~g. robežai un tikai rietumos ir jāmaina, gan ģeogrāfisku, gan etnogrāfisku apsvērumu dēļ.

Spišas (\pltxti{Zips}), Oravas (\pltxti{Arwa}), Trenčinas (Čačas apkārtnē) pagastu daļas veido ķīli, kas iekļaujas starp Silēziju un Galīciju, un kādreiz ietilpa senajā ungāru valstī. Tās atrodas Polijas valsts dabiskajās robežās, tajās ir poļu iedzīvotāji un tādēļ Polija uz tām pretendē.

Uz rietumiem no Galīcijas Polija izsaka savas pretenzijas uz Cešinas / kņazisti (\detxti{Teschen}), kura ir Austrijas Silēzijas sastāvdaļa.

Cešinas (\detxti{Teschen}) kņaziste aptver 2~282~km$^{2}$ [lielu teritoriju] un to apdzīvo 435~000 iedzīvotāju (1910), no kuriem 54,8\% ir poļi, 27,1\% čehi, bet 18\% vācieši. Čehu iedzīvotāji ir koncentrēti Frīdas apriņķī (78\%), kurš atrodas zemes dienvidrietumos, kamēr poļi apdzīvo Cešinas (77\%), Frīštates (64\%) un Beļskas (78\%) apriņķi; vācieši ir izkaisīti pilsētās un zemes tirdzniecības centros.

Starp poļu zemēm, kuras nebija Polijas valsts sastāvā dalīšanu laikā, Cešinas zeme pirmā uzsāka nacionālo atdzimšanu. Jau 1848.~gadā Slāvu kongresā\footnote{Slāvu kongress notika 1848.~gada jūnijā Prāgā, tajā Viduseiropas slāvu tautu pārstāvji iestājās par nacionālu līdztiesību un vērsās pret vācu nacionālisma izplatīšanos Austroungārijas slāvu zemēs.} Cešinas zemes pārstāvis demonstrēja savu poļu tautību, bet pirmais šajā zemē izveidotais dienas laikraksts bija ``\pltxti{Dziennik Polski}'' (``Poļu avīze'') (1848.~g.). Kopš šī laika Cešinas poļi ir arvien vairāk nostiprinājuši saiknes ar pārējām poļu zemēm.

Polija pieprasa visu Cešinas kņazistes poļu daļu; atzīst čehu tiesības uz Frīdas apriņķi. Sakarā ar to, ka oficiālā statistika ir neprecīza, par ko viegli pārliecināties, salīdzinot 1900.~un 1910.~gada tautas skaitīšanas rezultātus, tikai uz vietas strādājoša komisija var noteikt precīzu robežlīniju. \citespace{}

Kopumā Polija Cešinas zemes teritorijā prasa robežu no bijušās Ungārijas robežas (\pltxti{Lomna Gurna}), gar Frīdas apriņķa robežu līdz pat Grušovai, kur ieiet Prūsijas Silēzijas teritorijā.

Augšsilēzija ir pati nozīmīgākā no zemēm, kuras neietilpa Polijas valstī 1772.~g., bet kuras pievienošanu Polija tagad pieprasa.

Silēzija piederēja Polijai kopš tās pastāvēšanas sākuma līdz pat 14.~gadsimtam, kad Silēzijas poļu Pjastu nama kņazi\footnote{Pjasti (\pltxti{Piasty})~--- Polijā līdz 1370.~g. valdošā karaļu dinastija.} pakļāvās Čehijas karalim, vēlāk~--- Vācu ķeizaristei\footnote{Domāta Vācu nācijas Svētā Romas impērija (\detxti{Heiliges Romisches Reich Deutscher Nation}).}. No 14.~gadsimta Silēzija tika pakļauta pakāpeniskai ģermanizācijai, kurai pretoties spēja tikai Cešinas kņaziste, Augšsilēzija (Opoles reģentija) un daži apriņķi Vidussilēzijā (Vroclavas reģentija).

Ģermanizācijas process Lejassilēzijā un Vidussilēzijā tika pabeigts Prūsijas varas laikā, kad šīs zemes ieņēma Fridrihs II\footnote{Fridrihs II Lielais (\detxti{Friedrich II der Große}, 1712--1786)~--- Prūsijas karalis (1740--1786).}.

Daļa no Prūsijas Silēzijas, uz kuru Polija tagad pretendē, ietver Opoles reģentiju, izņemot Nissas, Grotkovas apriņķi, daļas no Nemodlinas un Prudņikas apriņķa, tālāk~--- Namislovas un Sicas apriņķi, daļa no Milicas apriņķa Vroclavas reģentija. Šī teritorija aptver 12~000~km$^{2}$ [lielu teritoriju] un tajā dzīvo 2~100~000~iedzīvotāju (1910.~g.), no kuriem, saskaņā ar oficiālo Prūsijas statistiku, 67\% ir poļi. Patiesībā poļu īpatsvars ir daudz lielāks; šīs teritorijas vidusdaļā un austrumu daļā tas sasniedz 90\%.

Šīs aizmirstās teritorijas nacionālā atdzimšana ir sākusies pēdējos piecdesmit gados. Līdz ar šī gadsimta sākumu Augšsilēzija sāka ievēlēt deputātus Reihstagā, kuri kopā ar Poznaņas apgabala un Rietumprusijas pārstāvjiem veidoja Poļu deputātu grupu un cīnījās pret valdības pretpolisko politiku.

Bez nacionālas dabas motīviem par labu Silēzijas pievienošanai Polijas valstij runā arī tās ģeogrāfiskais stāvoklis. Silēzijas teritorija iestiepjas dziļi poļu zemēs, starp Poznaņas apgabalu, Kongresa karalisti\footnote{Kongresa karaliste (\pltxti{Krolestwo Kongresowe})~--- pēc Vīnes kongresa (1815) izveidotā Polijas Karaliste (\pltxti{Krolestwo Polskie}), līdz 1917.~gadam faktiski Krievijas impērijas sastāvdaļa.}, Galīciju un Cešinas kņazisti, turklāt tā ir saistošais posms starp Poliju un Čehoslovākijas Republiku.

Augšsilēzijai ir pirmšķirīga nozīme no ekonomiskā aspekta. Kā raktuvju rajons, kurš apņem Polijas karalisti un Galīciju, no visas vācu produkcijas tas dod 23\% ogļu, 49\% alvas un 67\% cinka. \citespace{}

Poznaņas apgabalu Prūsija sagrāba pirmajā (1772.~g.) un otrajā (1793.~g.) Polijas dalīšanā. Turpmāk tas ietilpa Varšavas kņazistē (1807.--1815.~g.), bet pēc Vīnes kongresa (1815.~g.) atkal tika iekļauts Prūsijā.

Poznaņas apgabalā ir 2~100~000 iedzīvotāju, kuri dzīvo 29~000~km$^{2}$ [lielā teritorijā]. Oficiālā vācu statistika vāciešu īpatsvaru lēš ap 38\%. Patiesībā vietējo vāciešu skaits nepārsniedz 20\% no iedzīvotāju kopskaita.

Poznaņas apgabals ir bijis Polijas valstiskuma šūpulis, vissenākās poļu civilizācijas zeme, tā rezultātā~--- savā sabiedriskajā dzīvē visattīstītākā zeme. Lielākā daļa iedzīvotāju ir rietumzemju civilizācijas līmenī, demonstrē vislielāko enerģiju un organizatoriskās spējas, kā arī ļoti karstu poļu patriotismu. Minēto iemeslu dēļ apgabalam ir pirmšķirīga nozīme Polijas valsts atjaunošanā, un Polijas valsts pieprasa tā pievienošanu pilnībā.

Šajā vietā robeža iet pa vēsturisko Polijas 1772.~gada robežu, sākot no punkta, kurā izbeidzas jaunā robeža Silēzijā līdz pat punktam, kur tā sasniedz Rietumprūsiju.

Rietumprūsija, izņemot zemāk minētos apriņķus, ietilpa Polijas valsts sastāvā no 1466.~g.~līdz pat pirmajai Polijas dalīšanai (1772.~g.), kad Prūsija to paņēma sev. Polijā līdz pat otrajai dalīšanai (1793.~g.) palika tikai Dancigas\footnote{Danciga (\detxti{Danzig})~--- Gdaņskas vāciskais nosaukums.} pilsēta.

Rietumprūsijas teritorija ir 25~550~km$^{2}$ un [tur dzīvo] 1~703~000 iedzīvotāju, no kuriem, saskaņā ar oficiālo vācu statistiku, poļi it kā veido tikai 35\%. Šis skaitlis ir daudz mazāks, nekā īstenībā.

Rietumprūsijas iekļaušana Polijas valstī, no vienas puses, izraisītu iebraukušo [vācu] ierēdņu un karavīru emigrāciju, no otras puses~--- celtu terorizēto un ģermanizēto iedzīvotāju pašapziņu, kam sekas, visticamāk, būtu ievērojamas poļu majoritātes atklāšanās šajā zemē.

Šīs provinces vācu iedzīvotāji ir galvenokārt pie Vislas grīvas un dienvidrietumu apriņķos (Valeskas un Čluhovas). Šie divi apriņķi var būt ievērojamā to daļā atdoti Vācijai, apmaiņā Polija pieprasīs Pomerānijas austrumu daļu~--- Lenborkas, Bitovas apriņķi un Slupskas apriņķa daļu; divus pirmos no šiem apriņķiem Polija zaudēja 1657.~g., bet tie visi no ģeogrāfiskā viedokļa tai ir nepieciešami. \citespace{}

Rietumprūsija veido ģeogrāfisku veselumu ar Austrumprūsiju, kuras lielākā dala ar Kenigsbergu kā centru ir pilnībā pārvācota. Šīs no Vācijas ar poļu teritoriju Rietumprūsija atdalītās vācu teritorijas jautājumam priekš poļu tautas ir pirmšķirīga nozīme; no šī jautājuma atrisināšanas ir atkarīgs Polijas liktenis. Gadījumā, ja Austrumprūsijas vācu daļa joprojām paliktu neatņemama Vācijas valsts sastāvdaļa, šāds lietu stāvoklis būtu līdzvērtīgs vācu varas saglabāšanai Rietumprusijas poļu teritorijās ar Dancigas ostu. No Polijas attiecību viedokļa ar citām valstīm, ja to no jūras atdalītu vācu zemes, [tā] paliktu Vācijas žēlastībā, kā tas izpaužas pašreizējos apstākļos.

Polijas patiesas neatkarības nodrošināšanas nosacījums, kas tai ļautu ieņemt nopietnu un atbildīgu stāvokli Austrumeiropā un izvairīties no vācu ietekmes, būtu vienīgi Austrumprūsijas atdalīšana no Vācijas valsts ar sekojošu Rietumprūsijas atdošanu atpakaļ Polijai (ievērojot augstāk minētos robežas izlabošanas principus), kā arī tai līdz pirmajai dalīšanai Austrumprūsijā piederošās Varmijas atdošanu Polijai.

Turklāt Polija pieprasa atdot tai visu Austrumprūsijas dienvidu joslu, kura nekad nav ietilpusi Polijas valsts sastāvā, taču kuras iedzīvotāji neapšaubāmi ir poļi, turklāt viņu īpatsvars sasniedz 71\% (saskaņā ar oficiālo skolu statistiku).

Šīs poļu teritoriju daļas iedzīvotāji visvēlāk sāka piedalīties tautas dzīvē. Nacionālās pašapziņas atmoda tur ir vērojama tikai kopš 20.~gadiem.

Šīs provinces lietuviešu iedzīvotāju ziemeļaustrumu daļai pie Nemunas upes grīvas ir jābūt iekļautai Lietuvā un kopā ar Lietuvu apvienotai Polijas valstī. Pārējā Austrumprūsija, kas veido vācu zemi ar galveno pilsētu Kēnigsbergu (platība 19~000~km$^{2}$ un 1~070~000 iedzīvotāji) var pastāvēt neatkarīgi kā republika Tautu Savienības protektorātā. \citespace{}

\emph{Romāns Dmovskis}\footnote{Romāns Dmovskis (\pltxti{Dmowski}, 1864--1939)~--- Polijas valstsvīrs, nacionālā labējā spārna ideologs. 1917.--1919.~gadā Poļu nacionālās komitejas priekšsēdētājs Parīzē, 1919.~gadā Polijas pārstāvis Miera konferencē, 1923.~gadā ārlietu ministrs. Nacionāldemokrātiskās partijas, vēlāk citu uz tās bāzes izveidoto labējo organizāciju vadītājs.}

\pltxti{Avots: Wybor tekstow žrodłowych z historii Polski i powszechnej 1918--1945. Białystok, 2004. S.~23--28.}

Publicēts no: Polijas ārējā politika 1918.--1937.~gadā. Vēstures avotu krājums / Sastādījis, tulkojis un komentējis Ēriks Jēkabsons. Rīga, 2009. 19--23.~lpp.

\subsection{Izraksts no KSFPR valdības pārstāvja J.~Marhļevska slepenā ziņojuma V.~Ļeņinam par sarunām ar J.~Pilsudska pilnvaroto 1919.~gada 18.~oktobrī}

Viņš bija pārsteigts, ka tas padomju valdībai nav skaidrs, jo poļi sniedza pierādījumus. Proti, kad boļševiku armija atkāpās no Kijevas, poļi varēja nogriezt tai ceļu uz Maziru. Tā būtu pilnīga šīs armijas bojā eja. Otrkārt, poļi nevirzījās uz Daugavpili, lai gan bija iespēja. Pēc viņa teiktā, Polijas pavēlniecība cerēja, ka tas tiks saprasts Maskavā kā zīme un pierādījums, ka poļi nevēlas iet tālāk par viņu ieņemtajām pozīcijām un nevēlas palīdzēt Deņikinam. Es piezīmēju, ka šādas ``simboliskas zīmes'' ne vienmēr ir skaidras, un norādīju, ka pēc Polijas laikrakstu ziņām Varšavā uz sarunām ieradies Deņikina pārstāvis. Viņš atbildēja, ka, protams, Polija dotajā brīdī (viņš to ļoti uzsvēra) nevar strauji saraut attiecības ar Deņikinu, tādēļ pārstāvi gan uzņēma, taču viņš neko nesasniedza. Es piezīmēja, ka pēc tām pašam ziņām, pastāv vervēšanas biroji Viļņā Judeničam, Varšavā Deņikinam. Viņš atbildēja, ka tas tā bija. Bija pat 1~krievu bataljons frontē (Luņiņecas virzienā), bet tas ir noņemts, karavīri atbrīvoti, virsnieki ievietoti koncentrācijas nometnē. Pēc tam norisa parasta saruna: Polijai Deņikins ir nāves briesmas, u.t.t. Tad viņš jautāja, vai es uzskatu, ka ir iespējams panākt klusu vienošanos par karadarbības pārtraukšanu? Es atbildēju, ka šāda vienošanās ir ļoti neapmierinoša: kādu iemeslu dēļ kaut kur var notikt sadursme, un tūlīt pat līgums var tikt uzspridzināts. Viņš piekrita, ka šāds līgums nevar turpināties ilgi, bet nedēļu vai pat mēnesi var noturēties; Es jautāju, kāpēc gan, ja Igaunija, Latvija un Lietuva var vest sarunas, kāpēc to nevar darīt Polija? Viņš atzīmēja, ka Polija pašlaik to nevar, bet drīzumā varēs. Es jautāju, vai Antante traucē. Viņš izvairījās no atbildes, norādot, ka situācija ar Antanti ir ļoti sarežģīta, tiek risināti ļoti svarīgi jautājumi, un šodien Polijai vēl nav pilnīgas rīcības brīvības, bet drīz tai būs ``brīvas rokas''. Es centos gūt precīzāku atbildi, vai šeit ir runa par Silēziju, vai čehiem, bet nesaņēmu atbildi. Viņš tikai atzīmēja, ka visa robeža ar bijušo Austriju Polijai vēl nav noteikta. Atgriežoties pie jautājuma par slepeno vienošanos, es paziņoju, ka uzskatu to par iespējamu. Viņš, šķiet nopriecājās un sāka iztaujāt par formu. Es teicu, ka par to ir iespējams vienoties. Piemēram, sarunas tad norisinātos starp militārajiem darbiniekiem, viņi spēs noteikt robežas, pāri kurām armija neies. Garantijām tomēr ir būt nopietnām, piemēram, mutvārdu solījumam, ko no vienas puses dotu Pilsudskis, no otras Ļeņins, ka noslēgtie nosacījumi tiks ievēroti. Viņš uzskatīja, ka šāda forma ir piemērota.

Tulkots no: \rutxti{История России. ХХ~век: 1894--1939.~--- Москва: Астрель: АСТ, 2009, C.~686.}

\subsection{Izraksti no žurnālā ``\rutxti{Новый мир}'' (``Jaunā pasaule'') PSRS 1931.~gadā (N~5--6) publicētajām un žurnālā ``\rutxti{Новая Польша}'' (``Jaunā Polija'') Polijā 2001~gadā (N~11) daļēji publicētajām bijušā Polijas karagūstekņu nometnēs ieslodzītā J.~Podoļska (N.~Valdena) atmiņām}

Naktī iziet nokārtot vajadzības mēs baidījāmies. Sargi kādreiz nošāva divus puišus, kuri pirms rītausmas izgāja ārā no barakas, apsūdzot viņus par mēģinājumu bēgt. Bēdīgi slavenie mēģinājumi bēgt un priekšniecības apvainošana maksāja dzīvību nevienam vien simtam mūsu gūstekņu. Aizdomīgos bieži pārcēla uz soda baraku~--- to dzīvs neatstāja gandrīz neviens. \citespace{}

Pirms aizvešanas mēs veda uz pirti. Ņirgpilnās higiēniskās peldes maksāja dzīvību ne vienam vien ieslodzītajam. Stundām drebi kails aukstā uzgaidāmajā telpā, tad~--- strūkla remdena vai pārmērīgi karsta ūdens~--- un jau dzen tālāk. Uz slapja ķermeņa uzvelc slapju, smirdīgu apģērbs, ko kā netīru rituli izmet no dezinfekcijas iekārtas.

Pēc pirts mēs nodalīja ar niknu konvoju no pārējiem ieslodzītajiem. Vairāki cilvēki tika nošauti par mēģinājumu nodot zīmīti aizbraucējiem.\citespace{}

Nevaru nosaukt precīzu mūsējo skaitli, kas pabija poļu gūstā, bet diez vai ir kļūdaini apgalvot, ka uz katru, kas atgriezās Krievijā, palika divi, aprakti Polijā.\footnote{Nākamie atmiņu izraksti, kuri rāda autora J.~Podoļska, pēc tautības ebreja, ciešanas no poļu antisemītisma, Polijā izdodamajā žurnālā ``\rutxti{Новая Польша}'' netika pārpublicēti.} \citespace{}

Durvis atvērtas. Ar kliedzieniem un lāstiem ienāca vairāki apakšvirsnieki. Es nosaucu savu uzvārdu un dienesta pakāpi amatu armijā, ko biju apdomājis savā ieslodzījumā.

``Žīds?'' ar satricinātu naidu iemeta man [sejā] viens polītis.

``Nē.''

``Bet kas esi, suņa dzimums!''

``Tatārs,'' es teicu, ātri ņemot vērā bioloģiskās funkcijas, kas tuvināja musulmaņus ebrejiem. Pēkšņā pāreja islāmā ne vienu reizi vien man vēlāk daudz palīdzēja. Kur polis ebreju piekāva līdz nāvei, kādu citas tautības cilvēku viņš varēja labā noskaņojumā piekaut tikai līdz pusnāvei. \citespace{}

``Atceros, lielajās stacijās pie mūsu vagona pienāca kungi ar nūjām, ``dāmas no sabiedrības''. Visvairāk ``piemērotos'' ieslodzītos vilka ārā no vagona, sita un skrāpēja. Īpaša piekrišana bija ebrejiem. Ar sliktu dūšu atceros, kā šie zvēri nāca pie manis. Sākās pastāvīgais dialogs.

``Žīds?''

``Nē.''

``Tiešām?''

``Slimoju ar tīfu,'' es teicu ar vājprātīgā izmisumu. Tas deva vēlamo efektu, publika ātri lika mani mierā, sakot: ``Nu, un nosprāgsti, viņu nošaut vajadzētu.'' Kāds šļahtiču jauneklis tik tiešām vēlējās izmēģināt ar mani savu revolveri. Kāds viņu apturēja.

Tulkots no: \rutxti{Станислав Куняев. Шляхта и мы //

% http://www.voskres.ru/bratstvo/kunyaev2.htm // 30.12.2009
}

\subsection{Izraksts no Varšavā 1920.~gadā iznākošās avīzes ``\pltxti{Kurjer Nowy}'' raksta ``\pltxti{Czy to prawda}?'' (``Vai tā ir taisnība?) par latviešu karagūstekņu dzīvi Polijā}

Pirms dažām dienām caur Varšavu izgāja latviešu vienība, kurā bija vairāki simti bijušo Sarkanās armijas karavīru, piespiedu kārtā mobilizētu laikā, kad Latviju bija ieņēmis padomju karaspēks. Šie karavīri ar virsniekiem priekšgalā brīvprātīgi pārgāja Polijas pusē, lai tādā veidā atgrieztos dzimtenē.

Poļu daļas viņus ļoti labvēlīgi uzņēma, bet pirms nosūtīšanas uz nometni viņiem izdeva liecības, ka viņi brīvprātīgi pārgāja poļu pusē.

Taču pa ceļam uz nometni sākās sistemātiska laupīšana, no ieslodzītajiem novilka visu, dažreiz tos atstājot tikai apakšveļā. Dažiem izdevās saglabāt dažas no savām mantām, bet viņiem visu atņēma Stšalkovas nometnē. Viņi palika gandrīz bez drēbēm un apaviem, ģērbušies lupatās un basām kājām.

Bet tas vēl nebija nekas, salīdzinot ar latviešu sistemātisko spīdzināšanu. Tā sākās ar 50 sitienu nozīmēšanu ar pletni, savītu no dzeloņstieplēm, pie tam viņiem tika paziņots, ka latvieši kā ``žīdu algotņi'' dzīvi nometni neatstās. Vairāk nekā desmit ieslodzīto nomira no asins saindēšanās. Pēc tam uz trijām dienām ieslodzītos atstāja bez pārtikas un, draudot ar nāvi, aizliedza doties pēc ūdens. Divi ieslodzītie: Lācis un Škurins tika nošauti bez kāda iemesla. Draudi, iespējams, tiktu izpildīti, un neviens no latviešiem nebūtu atstājis nometni dzīvs, ja kapteinis Vāgners un un poručiks Maļinovskis netiktu nodoti tiesai. Pēc izmeklēšanas komisijas ierašanās situācija nedaudz uzlabojās, bet daudzi nomira no slimībām, aukstuma un bada.

Tagad, visbeidzot, viņi tika nosūtīti uz dzimteni. Atlaida viņus ziemā tajās pašās lupatās, kurās viņi ieradās nometnē.

Avots; ``\pltxti{Kurjer Nowy}'' (Warszawa), ne vēlāk kā 1920.~gada 16.~janvārī.

Tulkots no: \rutxti{Прудникова~Е., Чигирин~И.~--- Катынь. Ложь, ставшая историей. 2011. // %http://www.litmir.me/br/?b=174257&amp;p // 23.12.2014
}

\subsection{Bijušo latviešu karagūstekņu liecība par uzturēšanos Polijas gūstā 1919.--1920.~gadā}

Mēs, apakšā parakstījušies, no Polijas gūsta Strelkovā\footnote{Domāta Stšalkova (Strelkova~--- šīs apdzīvotās vietas krieviskais nosaukums).} pārnākušie Latvijas pilsoņi apliecinām sekošo:

Pēc gūstā krišanas visiem tika atņemtas drēbes un zābaki. Poļu zaldāti nesaudzīgi sita visus gūstekņus, it īpaši latviešus.

Strelkovā [1] kara gūstekņu lēģerī poļu priekšniecība apgājās ar kara gūstekņiem sliktāki kā ar lopiem: visu pagājušo vasaru (maiju, jūniju un jūliju) gūstekņi tika sistemātiski pērti, it sevišķi latvieši. Sišanu izdarīja sekošā kārtā: bez jebkāda iemesla izdzen visus latviešus no baraka, sastāda rindā, pavēl iziet no rindas pa nodaļām dažus soļus uz priekšu, tad pavēl izģērbties, un poļu zaldāti ar no telefona drātīm savītiem sitamiem (nagaikām) apstrādā savus upurus, sizdami līdz asinīm. Vairāki dabūja līdz simtu cirtienu. Ja kāds kliedza, viņam virsnieks (poručiks Maļinovskis) uzspēra ar kāju uz kakla un ar revolveri draudēja nošaut.

Tādā kārtā tika izpērti visi latvieši 18.~un 19.jūnijā: viena daļa tika pērta 18.jūn[ijā] no pusdienas līdz vakaram un otra~--- 19.jūn[ijā] no rīta līdz pusdienai. Šī pēršana notika poručika Maļinovska vadībā. Reti kāds latvietis šādā vai tādā veidā iztika cauri bez pēriena: tā, piem[ēram], Pēteris Ciemiņš, uzdodamies par krievu, tika cauri ar veselu ādu. Bet turpretim Maļinovskis pavēlēja nošaut latvieti Arnoldu (jeb Arvīdu) Lāci par to, ka viņš, nevarēdams paciest sitienus, kliedza. Maļinovskis teica, ka tas būšot citiem par biedinājumu.

Strelkovā bija ap 500 latviešu.

Dzīve lēģerī ir šausmīga: baraki veci, sapuvuši~--- laiž cauri vēju un lietu. Vasarā vēl varēja kaut kādi nebūt iztikt, lai gan mitruma dēļ daudzi gūstekņi saslima un nomira, pēc lietus ūdens bija jānes ārā ar spaiņiem. Grīdu barakos nav un katrā baraka galā ir pa lodziņam. Ziemai iestājoties, baraki kurināti netika, lai gan krāsnis bija pilnīgi kārtībā. Ja kādu pieķēra, kas patvarīgi aizkurinājis krāsni, viņam neizdeva maizi vairākas dienas no vietas. Baraka pārzinis nereti dabūja nosēdēt karcerī tamdēļ, ka nebija raudzijies, lai krāsnis netiktu kurinātas. Bija gadījieni, kad baraka pārzini par to pat sita.

Šādos apstākļos bija jānodzīvo viss laiks un, salam iestājoties, bez bada jācieš arī aukstums, jo visas labākas drēbes jau sen bija atņemtas.

Vasarā maizes izdeva 1/2--3/4 M[ārciņas], dažreiz pat tikai vienreiz nedēļā. Pēc jaunās ražas sāka izdot pa 3/4 M[ārciņas] gandrīz katru dienu. Siltu ēdienu deva vienreiz pa dienu, pa litram putraimu putru vai no rāceņiem, bietēm un burkāniem vārītu zupu. Dažreiz ēdiens bija labāks, dažreiz sliktāks. Kad daudzi sāka slimot (ar tīfu), ēdienu drusku uzlaboja. Tomēr ēdiens bija tāds, ka nebijām ne paēduši, ne arī badā varējām nomirt. Sevišķi grūta tomēr bija sistemātiskā spīdzināšana: daudzreiz poļu zaldāti garāmejot bez kaut kāda iemesla uzbruka un sita mūsu gūstekņus.

Pagājušo gadu, jūlija beigās ieradās kāda Sabiedroto (amerikāņu) komisija, kura sūdzības izmeklēja un visiem poļu zaldātiem atņēma sitamos. Sistemātiskā sišana tika izbeigta, bet sists tomēr vēl tika.

Paraksti:

\begin{enumerate}
\item Pēteris Adama d. Ciemiņš, Palsmanes pag., Cēsu apr., Grušļu mājas.
\item Jānis Alberta d. Bemfelds, Sоpeles pag., Dobeles apr., Kalna Košķu mājas.
\item Otto Otto d. Dimiņš, Ilzenes pag., Valkas apr., Čonkas mājas.
\item Jānis Jēkaba d. Kolbergs, Gaiķu pag., Kuldīgas apr., Zubru muiža.
\item Andrievs Antona d. Vaivods, Vitebskas g[uberņa], Daugavpils apr., Vorkovas pag., Lielkursīšu sādža.
\end{enumerate}

\emph{P.~S.} Gūstekņu apģērbi ir galīgi noplīsuši, daudzi no viņiem ir tikai kreklos. Pie izsūtīšanas poļi viņiem nekādas drēbes nav izdevuši.

Avots: LVVA, 2574. f., 3.~apr., 13.~l., 43.~lp.

Pārpublicēts no: Ēriks Jēkabsons. Latviešu tautas karagūstekņi Polijā 1919.--1920. gadā // Latvijas Arhīvi. 2006. N~2.~130--131.~lpp.

\subsection{Izraksts no Polijas 4.armijas pavēlniecības II nodaļas 1920.~gada 28.~aprīļa ziņojuma armijas Virspavēlniecībai par baltkrievu attieksmi pret poļu okupācijas karaspēku.}

Slepeni. \citespace{}

No poļu karaspēka ieiešanas laika pareizticīgās Baltkrievijas (Minskas novada) teritorijā attiecības un politiskie noskaņojumi ir būtiski mainījušies~--- mainījušies mums par sliktu.

Sagaidīja mūs vispār laipni, kā armiju~--- atbrīvotāju, kas nes likumību, dzīves normalizēšanos, cenu pazemināšanos un lielāku drošību. Šie pozitīvie aspekti tomēr iedzīvotāju atmiņā ilgi nepalika. Viņi drīz pie tiem pierada, atšķirības starp Poliju un Boļševikiju zaudēja savu asumu un sāka izpausties mūsu kundzības negatīvā puse

Karaspēks tikai ieradās un ieradās, sablīvējot frontes līniju un ceļot etapus. Mūsu armija~--- sastāvā pēc nodienētajiem gadiem jauna,~--- pakāpeniski apgūst veca zaldāta psiholoģiju, ka, jau reiz viņš pilda militāro dienestu, tad viss ir viņam pakļauts, un viņa veic rekvizīcijas un vardarbību pret iedzīvotājiem, kurus nepazīst, kam nesimpatizē, kurus uzskata par svešiniekiem.

Civilā administrācija tiek izvēlēta tā, ka senais baltkrievu zemnieku priekšstats: ``polis~--- tas ir pans'' turpina valdīt visā Minskas novadā, tikai viens vecākais nav muižnieks, visus svarīgākos amatus apriņķu pārvaldē ieņem muižnieki. Visi Minskas rajona priekšnieka Račkeviča, mūsdienīga, vēl līdz šim iedzīvotāju simpātijas baudoša cilvēka, centieni mainīt šo situāciju ir neauglīgi, jo pēc sagatavotības administratoru amatiem Minskas novadā ir derīgi gandrīz vienīgi muižnieki. Zemākais personāls (uzraugi, policijas darbinieki) tiek izraudzīts pārsvarā no bijušajiem krievu ierēdņiem un sīkajiem šļahtičiem, kurus raksturo pērkamība, kontrabanda, administratīva neprasme un vienlaikus aizraušanās ar pātagu un uzrunu: ``tu, bauri''. Žandarmēriju visur ienīst. Visi pasākumi sakarā ar virspavēlnieka aicinājumu\footnote{Acīmredzot domāts 1919.~gada 28.~aprīļa J.~Pilsudska aicinājums bijušās Lietuvas lielkņazistes iedzīvotājiem, kurā viņš solīja piešķirt tiem tiesības pašiem lemt savas nacionālās problēmas.} drīzāk nes ļaunumu nedemokrātiskās un neapmierinošās izpildes dēļ; municipālā kampaņa iedzīvotājus noskaņoja pret jebkurām vēlēšanām. Igumenskas apriņķa [baltkrievu \betxti{Ігуменскі павет}] padomē no 25 locekļiem ir 1 zemnieks, baltkrievu skola ir vieta, kur notiek bezgalīgi skandāli un konflikti ar inspektoriem, kuri nesaprot tās nozīmi.

Poļu sabiedrība kopumā savos entuziasma pilnajos centienos savienoties ar Poliju neapzinās to, ka tai kā nacionālai minoritātei vajadzētu nokārtot attiecības ar vietējiem iedzīvotājiem, nevis padziļināt antagonismu, ka nākotnē par mūsu lielāko spēku var izrādīties morālā autoritāte, bet nevis nikns šovinisms un neiecietība.

Pēc atgriešanās savās muižās poļu muižnieki ar savu atriebību kalpiem un zemniekiem, izvairīšanos no piegādēm un šķūtīm, zemes iznomāšanu ar augļotāja ienākumiem, masveida mežu pārdošanu spekulantiem izsauc ne tikai baltkrievu zemnieku, bet arī sīko zemes īpašnieku~--- šļahtiču sašutumu, palielinot šķiru naidu, padziļinot tumšajās baltkrievu masās nepatiku pret Poliju, jo ``Polija~--- tie ir pani''.

Augstāk izklāstītās mūsu kļūdas izmanto savā aģitācijā krievu un baltkrievu eseri, bet praksē~--- boļševiki.

Tā rezultātā viss Minskas novads ir pārklāts ar kaujas grupu tīklu un mežos aug bruņotas partizānu vienības \citespace{}.

Ģenerālštābam piekomandēts rotm[istrs] un nodaļas priekšnieks Stamirovskis\footnote{K. Stamirovskis}

Avots: \pltxti{Centralne Archiwum Ministerstwa Spraw Wewnẹtrznych, R-20, t.~4a, d.~54.}

Tulkots no: \rutxti{Документы и материалы по истории советско-польских отношений. Т.~V.~Москва, 1963. C.~19--20.}

\subsection{
Izraksti no KSFPR Revolucionārās kara padomes priekšsēdētāja Ļ.~Trocka ziņojuma apvienotajā KSFPR VCIK, Maskavas Strādnieku un zemnieku un sarkanarmiešu deputātu padomes, arodbiedrību un fabriku-rūpnīcu komiteju sēdē 1920.~gada 5.~maijā, atkārtota 10.~maijā mītiņā Gomeļā}

Mēs ne tikai atzinām un sankcionējām Polijas neatkarību, bet mēs aizstāvējām šo neatkarību pret visvareno vācu militārismu. Mūsu delegācija Brest-Ļitovskas sarunās atteicās atzīt par neatkarīgas Polijas pārstāvi Kuharževska\footnote{J.~Kuharževskis (\pltxti{Jan Kucharzewski}, 1876--1952)~--- poļu politiķis, 1917.~gadā Polijas Reģentu padomes loceklis, novembrī kļuva par marionešu valdības premjerministru.} valdību~--- šo Berlīnes plēsoņu nabaga lietvedi. Vācu imperiālisms ļoti bija vajadzīga~--- pirmkārt lai ietekmētu savu darbaļaužu masu sabiedrisko viedokli~--- tieša vai netieša okupācijas, vardarbīgā režīma Polijā, kurš tika tēlots kā poļu tautas pašnoteikšanās, atzīšana no mūsu puses. Šajā laikmetā Hohencollernu aģenti jau mēģināja nozagt Krievijas revolūcijai šo formulu un padarīt to par savu sagrābšanu un vardarbību aizsegu. Mēs bija pārāk vāji, lai palīdzētu apspiestajai Polijai ar ieročiem rokās. Bet mēs bijām kopā ar Polijas tautai pret tās apspiedējiem,~--- mēs pretstatījām laupītāju meliem mūsu revolucionāro patiesību par Poliju. \citespace{}

Kas attiecas uz Poliju, tad, kad tā ir sagrāba Lietuvu, Baltkrieviju un ievērojamu daļu no Ukrainas un Lielkrievijas teritorijas, tas ir, uz pagājušā gada beigām, šķita, ka tā ir sasniegusi zināmu līdzsvaru starp savu apetīti un saviem spēkiem. Miera ar Poliju mums nebija, bet militārās operācijas netika izvērstas, bet tām bija vairāk vai mazāk nozīmīgu atsevišķu izlūkošanas vienību darbības raksturs. Nopietnas rīcība nebija. Abās pusēs arvien vairāk nostiprinājās pārliecību, ka kara nebūs, tas ir pabeigts, un drīzumā diplomātija rezumēs rezultātus un parakstīs sliktu vai labu, pilnīgu vai nepilnīgu, bet mieru. Tā mūsu Rietumu fronte stāvēja cerībā uz ātru mieru, un mūsu padomju diplomātija darīja visu iespējamo, lai paātrinātu šī miera iestāšanos. \citespace{}

Tagad, kad Pilsudskis, visvairāk agresīvo imperiālisma elementu bīdīts, metās paverdzināt Ukrainu, viņš nolēma atkal slēpties aiz muguras tā paša Petļuras\footnote{S.~Petļura (\uktxti{Симон Васильович Петлюра}; 1879--1926)~--- ukraiņu politiskais un militārais darbinieks. 1919.~gada 10.~februārī faktiski kļuva par vienpersonisku Ukrainas diktatoru.} muguras. Un nožēlojamais ``hetmanis'', kurš pārdeva sevi austriešu un vācu ģenerāļiem un arī angļu-franču imperiālistiem, protams, nepalaida garām gadījumu pārdot sevi Polijas paniem. Pēc tam, kad Pilsudskis martā sagrāba Moziru, Kaļinkovičus, Ovruču un Rečicu, viņš no 23.~aprīļa uzsāka plašu uzbrukumu Volīnijas~--- Kijevas frontē, sagrāba Žitomiru un Žmerinku, bet galvenos spēkus virzīja uz Kijevu. Pašlaik Polijas karaspēks tieši un nepastarpināti apdraud Kijevu un visu Ukrainu, un reizē arī cieši ar garīgās, materiālās un militārās vienotības saiknēm ar to saistīto Padomju Krieviju. \citespace{}

Poļu strādnieku šķira negribēja un negrib karu. Polijas zemniecība no jaunā režīma, Pilsudska režīma un viņa sabiedrotajiem šļahtičiem saņēma tikai drupatas, vai precīzāk, tikai solījumus sniegt drupatas. Tā nevar ilgi iestāties par šo režīmu, tā nevar tiekties pēc kara, kurš atnesīs tai valsts nodokļu pieaugumu, un, ja karš ievilksies, tad arī pilnīgu grimšanu nabadzībā, pilnīgu [spēku] izsīkumu.

Tas viss ir neapstrīdams, bet šodien tas vēl nav skaidrs visai Polijas zemniecībai. Tajā vēl ir spēcīgi nacionālie aizspriedumi. Polija, pēc ilgas apspiestības vēl pārāk nesen ieguva neatkarīgas republikas stāvokli. Nacionālais jūtīgums joprojām ir pārāk svaigs, valsts neatkarības medusmēnesis vēl nav pagājis, un uz šiem noskaņojumiem mēģina izveidot savu politiku ``poļu valsts priekšnieks'' Pilsudskis. Vēl saglabājas visvairāk atpalikušo zemnieku masu neuzticība un naids pret Krieviju un krievisko, jo viņu apziņā un atmiņās Krievija un krieviskais~--- ir cars un cariskais. Lūk no šāda vēsturiskā kapitāla cer gūt procentus ``valsts priekšnieks'' Pilsudskis.\citespace{}

Biedri, es gribētu, lai galvenais secinājums, ko jūs gūsiet no šīs sanāksmes, būtu tas, ka cīņa, kas mums stāv priekšā, būs smaga un sasprindzināta. Poļu buržuāzija zina, ka tā, uzbrūkot mums, likusi savu likteni uz spēles. Un tie, kas stāv aiz viņas, zina, ka baltgvardu Polija, kas apspiež poļu proletariātu, kurš gadu desmitiem ir saistīs ar Petrogradas un Maskavas proletariātu kopīgā revolucionārā cīņā~--- ka šī baltgvardu Polija cenšas starp mums un Eiropu uzcelt barjeru. Poļu šļahtiči saka, ka krievus, šos barbarus un skitus, ir nepieciešams atvirzīt tālāk uz austrumiem. Bet mēs tiecamies uz rietumiem, pretī Eiropas proletariātam, kurš zina, ka mēs varam ar to sastapties ne citādi kā tikai pār baltgvardu Polijas līķi, brīvā un neatkarīgā strādnieku un zemnieku Polijā.\citespace{}

[sekojošais izraksts no Ļ.~Trocka runas mītiņā Grodņā]

Biedri, es šodien biju pie Rečicas. Man tur, mūsu frontē, stāstīja par tām neaprakstāmajām zvērībām, kuras veic poļu baltgvardu virsniecība un poļu kulaki pret sagūstītajiem un ievainotajiem sarkanarmiešiem. Viņi vairs neatzīst gūstekņus vispār. Tie pakar ne tikai komunistus, tie pakar katru vidusmēra bezpartijisko sarkanarmieti, kurš nonācis viņu rokās, tie iznīcina pat ievainotos un slimos. Biedri, es jautāju, vai tas nav pārspīlēts, vai tie nav meli, jo arī ienaidnieku nedrīkst apmelot. Man teica: ``tad atnāca tāds un tāds, tur izlavījās kāds un tāds~--- viņi visi ir cilvēki, kuri pelna uzticību, viņi redzēja un visi apstiprina šīs šausminošās zvērības''.

Kā mēs atbildējām? Šodien pat mēs Republikas Revolucionārās kara padomes vārdā izdodam pavēli visam Rietumu frontes karaspēkam atteikties no atriebības par visiem šiem noziegumiem bezpalīdzīgajiem poļu gūstekņiem. Ja mēs ņemsim poļu strādnieku vai zemnieku gūstā, tad, biedri, lai tiek nocirsta tā sarkanarmieša roka, kurš pret nebrīvē esošu un neapbruņotu, ievainotu un slimu pacels nazi. Mēs cīnāmies tikai pret bruņotiem. Kāpēc poļu šļahta apšauj mūsu sagūstītos strādniekus un zemniekus? Tāpēc, ka viņa zina, ka godīgs sarkanarmietis vienmēr paliks magnātu un varmāku niknākais ienaidnieks. Bet ja mēs gūstā saņemto poļu strādnieku un poļu zemnieku nosēdināsim sev līdzās un pateiksim viņam mūsu patiesību par Pilsudska un viņa magnātu meliem, tad šis poļu strādnieks vai poļu zemnieku pēc dažām nedēļām vai pēc dažām dienām kļūs par niknāko Pilsudska ienaidnieku. Tā mēs pārvētām par revolucionāriem vācu karavīrus, kuri pēc tam sacēlās pret Vilhelmu, un austriešu un ungāru, un Kolčaka un Deņikina [karavīrus]~--- tie visi izgāja mūsu skolu. Mēs gūstekņus nenošāvām, bet pārvērtām par apzinīgiem cīnītājiem. Tāpēc arī poļu leģionāriem, poļu strādniekiem un zemniekiem, kuri nonāk mūsu gūstā, nav jābaidās no zvērībām un nāves sodiem~--- nē, mēs nesīsim viņiem komunisma gaismu, mācības gaismu par visu strādājošo brālību. Ja viņi pie mums nāca ar plēsonīgā imperiālisma dzelteno karogu, tad viņi dosies prom no mums zem revolūcijas un komunisma sarkanā karoga. Nepieciešama nežēlīga cīņa kaujās un augstsirdība attiecībā pret sagūstīto ienaidnieku. Nežēlīgs naids~--- pret magnātu, kapitālistu; draudzīgi pastiepta roku~--- poļu strādnieku masām. Mēs neļausim ielauzties pie mums, bet mēs arī nepacelsim rokas pret poļu tautas neatkarību. Un mēs ticam, mēs zinām, ka Polijas Republika no šī kara iznāks citāda.\citespace{}

Tulkots no: \rutxti{Л. Троцкий. Война с Польшей. Доклад на объединенном заседании ВЦИК, Московского Совета Раб., Кр. и Красноарм. Депутатов, правлений профессиональных союзов и фабрично-заводских комитетов 5 мая 1920 г. //

%http://www.magister.msk.ru/library/trotsky/trotl782.htm // 22.05.2015
}

\subsection{Izraksti no Strādnieku Zemnieku Sarkanās armijas štāba IV pārvaldes priekšnieka J.~Bērziņa\footnote{Jānis Bērziņš (krievu: \rutxti{Ян Карлович Берзин}, dzimis Pēteris Ķuzis; 1889--1938), latviešu komunists un profesionāls revolucionārs. PSRS militārā izlūkdienesta vadītājs (1924--1935, 1937).} ziņojuma PSRS aizsardzības tautas komisāram K.~Vorošilovam}

1927.~gada 27.~janvārī

Pilnīgi slepeni

\begin{enumerate}
\item \citespace{} Lai novilcinātu [gaidāmo] mūsu Savienības karu ar kapitālisma pasauli un uzlabotu militāro-politisko situāciju, ir mērķtiecīgi un nepieciešami \citespace{}

\item traucēt Polijas-Vācijas strīdīgo jautājumu (Danzigas koridors, Augšsilēzija u.c.) atrisināšanu;

\item traucēt Polijas un Baltijas valstu savienības noslēgšanu;

\item atturēt Vāciju no galīgas pārejas mums naidīgā nometnē;

\item veicināt Francijas un Lielbritānijas attiecību saasināšanos. \citespace{}

\end{enumerate}

Avots: Krievijas valsts kara arhīvs (\rutxti{Российский государственный военный архив~--- РГВА, Ф.~33987. Оп.~3. Д.~128. Л.~26. Подлинник}).

Tulkots no: \rutxti{Cоветско-польские отношения в политических условиях Европы 30-х голов XX~столетия. Сборник статей. Москва: ``НАУКА'', 2001. C.~60.}

\subsection{Izraksti no rabīna Jēkaba Beikera (\entxti{Jacob L.~Baker}) atmiņām par dzīvi Jedvabnē starpkaru periodā}

Es piedzimu Jedvabnē 1914.~gadā un pavadīju tur pirmos divdesmit savas dzīves gadus. \citespace{} Strādājām mēs smagi, strādājām laukā, un mūsu kaimiņi-poļi to redzēja. Par to viņi mūs cienīja. Mēs dzīvojām ar viņiem saskaņā un draudzībā, kā kaimiņi. \citespace{} Es varētu minēt daudzus piemērus, kādā saskaņā dzīvoja ebreji un poļi Jedvabnē. Mēs uzticējāmies viens otrām. Mūsu dzirnavas atradās ārpus pilsētas, lai to spārni labāk ķertu vēju. Kad bija ebreju svētki, un mēs ilgu laiku pavadījām sinagogā~--- kam mēs atstājām mūsu bērus? Kaimiņiem poļiem, un viņi darbojās ar tiem kā savējiem. Iespējams, tas ir labākais piemērs \citespace{}

Es tad jau biju pieaudzis un redzēju, kā no 30.gadu vidus auga poļu naidīgums pret ebrejiem. Tā sākās ebreju veikalu piketēšana, ierobežojumi mājlopu rituālai kaušanai, ``geto solu'' ierīkošana.

Arī Jedvabnē Nacionāli demokrātiskās jaunatnes grupas stāvēja ar metāla stieņiem pie ebreju veikaliem, lai poļi neko tur nepirktu. Sākās uzbrukumi ebrejiem, notika slepkavības. Es atceros vismaz divu ebreju bēres, kurus nogalināja poļu huligāni. Mēs dzīvojām arvien lielākās bailēs. Ir skumji to teikt, jo Polija~--- manas mājas, un mani bērni vēl tagad dzied poļu dziesmas. Bet 30.[gadu] beigās gandrīz visi ebreji gribēja atstāt Jedvabni, bēgt no Polijas. Aiz bailēm no vajāšanām un gaidāmā kara. Es 1938.~gada februārī aizbraucu uz Ameriku: tur jau bija mans brālis Jehuda. \citespace{}

Avots: \pltxti{``Rzeczpospolita'', 2001, 10--11 martā}

Tulkots no: % http://www.novpol.ru/index.php?id=110 // 30.12.2009.

\subsection{Izraksti no Vācijas sūtņa Varšavā Hansa Ādolfa von Moltkes politiskā ziņojuma Ārlietu ministrijai 1939.~gada 28.~martā}

Satrauktais noskaņojums, kuru Polijā izsauca Čehijas jautājuma atrisināšana, \citespace{} ir vēl ievērojami pieņēmies spēkā. Polijas politiskās un militārās aprindas acīmredzot uzskata, ka tagad \citespace{} vācu apvērsuma briesmas pret Dancigu ir jūtamas gaisā. Iedzīvotāju vidū izplatās mežonīgas baumas, piemēram, ka notiek cīņas starp vācu un poļu karaspēku Oderbergā\footnote{Oderberga~--- pilsēta Vācijā, netālu no Polijas robežas.}, ka ministrs Beks ir arestēts un līdzīgas fantastiskas ziņas. Nopietnāka nozīme ir kara noskaņojuma rašanās faktam, kas tiek veicināts ar antivāciskiem paziņojumiem presē, un daļēji ar griezīgi asu oficiālo propagandu, īpaši provincē, kas jau ir novedusi pie daudziem incidentiem. Plašās aprindās tiek uzskatīts, ka šodien karš ir kļuvis neizbēgams un tūliņ izcelsies.

Praktiskie pasākumi, ko veic valdība, palīdz paaugstināt pastāvošo kara psihozi. Pēdējo nedēļu laikā ir notikusi triju gadu 1912, 1913 un 1914 rezervistu iesaukšana. \citespace{}

Oficiālajai militārajai propagandai ir īpaši raksturīgs sauklis ``Mēs esam gatavi''. Ar to, cita starpā, tiek norādīts, ka poļi, atšķirībā no čehiem, nejūt mazvērtības sajūtu pret pasaules spēcīgākajām tautām.

Ārvalstu divīziju skaits nebaida poļus, jo viņu pašu armija, tās apbruņojums un poļu tautas kareivīgais gars ir pietiekams, lai, lai nodrošinātu poļiem uzvaru. Daudzas citas publikācijas, kas parādās katru dienu presē ieturētas tai pašā garā un toņkārtā. \citespace{}

Ka tā nav tikai preses propaganda, liecina oficiāls paziņojums, ko kara ministra vietnieks Gluhovskis [\pltxti{Gluchowski}] ir sniedzis nopietnās sarunās, kurās viņš izteicās, ka vācu armija ir liels blefs, jo Vācijai neesot apmācītu rezervju lai ar tām papildinātu vienības. Jautāts, vai viņš uzskata, ka Polija patiesi ir militāri pārāka par Vāciju, Gluhovskis atbildēja: ``Bet pats par sevi saprotams.``

Avots: \detxti{Akten zur deutschen Auswärtigen Politik 1918--1945. Aus dem Archiv des Deutschen Auswärtigen Amtes, Serie~D, Bd.~VI. Baden-Baden 1956, Dok.~116, S.~119--120}

Tulkots no: \detxti{Jochen Böhler. Der Überfall. Deutschlands Krieg gegen Polen. Frankfurt am Main, 2009, S.~35--36.}

\subsection{Polijas ārlietu ministra J.~Beka telegramma Polijas sūtnim Francijā J.~Lukasevičam 1939.~gada 20.~augustā}

Franču un angļu vēstnieki griezās pie manis sakarā ar Francijas un Lielbritānijas-padomju štābu sarunām, kuru laikā Padomes pieprasīja piešķirt tām iespēju nākt saskarsmē ar vācu armiju Pomorjē, Suvalku novadā un austrumu Mazpolijā. Šo nostāju atbalsta angļu un franču demaršs.

Es atbildēju, ka tas nav pieņemams, ka šīs valstis apspriež citas suverēnas valsts teritorijas militāru izmantošanu. Poliju ar Padomēm nesaista nekādi militāri līgumi, un Polijas valdība šādu līgumu noslēgt negatavojas.

Francijas sūtnis teica, ka tādā gadījumā viņi atbildēs Padomēm, ka Polijas valdība ir atteikusies no [jautājuma] apspriešanas vai, ka Francijas valdība neuzņēmās izdarīt formālu demaršu, jo bija pārliecināta par negatīvu atbildi. Es atstāju jautājumu par atbildi Padomēm Francijas un Anglijas ieskatiem, atrunājot, ka atbilde nedrīkst radīt ieganstu pārpratumiem.

Beks

Avots: \rutxti{Документы и материалы кануна Второй мировой войны 1937--1939~гг. в 2-х томах. Январь-август 1939~г. Москва. Политиздат. 1981, Т.~2. С.~314.}

Tulkots no: % http://www.doc20vek.ru/taxonomy/term/116 // 22.05.2015

\subsection{Izraksts no Lielbritānijas ārlietu ministra E.~Halifaksa telegrammas Lielbritānijas sūtnim Varšavā H.~V.~Kennardam 1939.~gada 20.~augustā}

1. Ja Polijas valdības galīgā atbilde [Lielbritānijas priekšlikumam izlaist Sarkano armiju caur Polijas teritoriju, lai uzsāktu kara darbību pret Vāciju] būs nelabvēlīga, radīsies ļoti nopietna situācija. Padomju valdība uzskata, ka ir bezjēdzīgi turpināt militārās sarunas, kamēr nav saņemta atbilde, un sanāksme Maskavā attiecīgi tiek atlikta uz vairākām dienām. Ja atbilde būs noraidoša, sarunas, visticamāk, tiks izjauktas, un Lielbritānijas un Francijas mēģinājumi panākt vienošanos ar Padomju Savienību būs neveiksmīgi. Esmu pārliecināts, ka šāda neveiksme iedvesmos Hitleru sākt karu, kurā Polija uzņemsies pirmā uzbrukuma smagāko nastu. No otras puses, es esmu pilnīgi pārliecināts, ka militāri politiska līguma noslēgšanai ar Padomju Savienību būs mērķis viņu atturēt no kara.

2. Beka k-s acīmredzot nepiekrīt šim viedoklim un uzskata, ka Polijas piekrišana Padomju priekšlikumam novestu pie tā, ka Vācija nekavējoties izsludinās karu. Es uz situāciju skatos nedaudz citā gaismā. Mēs, varbūt, esam uz kara robežas, kurā dažās nākamajās nedēļās vai pat dažās nākamajās dienās Polija var kļūt par iznīcinoša uzbrukuma upuri. Ja Beka k-s domā, ka viņš var novērst vai mazināt šāda uzbrukuma iespējamību vienkārši ar to, ka viņš atturēsies pieņemt palīdzību no Padomju Savienības, es domāju, ka viņš kļūdās. Bet kad viņš saka, ka, ja karš patiesi sāksies, tad stāvoklis, iespējams, būs cits, un poļu pozīcija, iespējams, mainīsies, es domāju, ka viņš nonāk pretrunā pats sev. Ja Polija spēs pieņemt padomju palīdzību, ja sāksies karš, tad nav skaidrs, kāpēc tā nevar piekrist padomāt par šādas palīdzības pieņemšanu, kad karš ir tuvu.

3. Es pilnībā saprotu visus negatīvos aspektus un risku, atļaujot padomju karaspēkam ienākt Polijas zemē; tomēr šķiet, ka šādam riskam var dot priekšroku salīdzinājumā ar Polijas neatkarības iznīcināšanas risku, ja tā nesaņems palīdzību, lai cik arī nevēlama tā būtu; un būtībā labākais veids, kā izvairīties no nepieciešamības pieņemt šādu palīdzību, ir atvieglot angļu-franču-padomju savienības noslēgšanu, piekrītot izstrādāt šādas palīdzības plānus, kas tiktu sniegta vajadzības gadījumā.

4. Pēdējo dažu dienu laikā Polijas valdība ir uzstājusi, ka ir vēlams pēc iespējas ātrāk noslēgt oficiālu Anglijas un Polijas līgumu, pamatojoties uz to, ka tas radīs uzticību Eiropā; taču Viņa Majestātes valdība uzskata, ka Anglijas un Polijas līguma noslēgšanas pozitīvais efekts, visticamāk, nespēs līdzsvarot patiesi katastrofālo angļu-franču-padomju sarunu Maskavā galīgo neveiksmi.

Viņa Majestātes valdība pēdējos kritiskajos mēnešos ir veltījusi visnopietnākās pūles, lai izveidotu tā saucamo miera fronti, kuras viens no mērķiem ir saglabāt Polijas neatkarību. Šie centieni tiks apdraudēti, un var tikt pilnībā iznīcināti, ja Polija šajā ļoti svarīgajā etapā, neraugoties uz visām grūtībām, kuras es pilnībā saprotu, nedos savu pienesumu; tai jāizdara viss iespējamais, lai saglabātu Polijas neatkarību, atvieglotu un pilnībā padarītu efektīvu tai piedāvāto palīdzību.

5. Es priecāšos, ja jūs ar pilnu nopietnību paziņosit šos apsvērumus Beka k-m. Mēs piedzīvojam ļoti kritisku brīdi, kad viņa lēmums var kļūt par izšķirošu faktoru.

Avots: \rutxti{Документы и материалы кануна Второй мировой войны 1937--1939 гг. в 2-х томах. Том~2. Январь-август 1939~г. Москва. Политиздат. 1981. С.~316--317.}

Tulkots no: % http://militera.lib.ru/docs/da/kanun/index.html // 8.11.2011.

\section{Ilustrācijas II nodaļai}

% TODO images

\chapter{Polija Otrajā pasaules karā}

\epigraph
{Katrā gadsimtā ir savi viduslaiki.}
{Staņislavs Ježijs Lecs (\pltxti{Stanisław Jerzy Lec})}

\epigraph
{Politika ir karš bez asinsizliešanas. Karš ir politika ar asinsizliešanu.}
{Mao Dzeduns (\zhtxti{毛澤東})}

\epigraph
{Dižs laikmets var ietvert pamatīgu skaitu mazu cilvēku.}
{Staņislavs Ježijs Lecs (\pltxti{Stanisław Jerzy Lec)}}

\epigraph
{Vāja paaudžu atmiņa stiprina leģendas.}
{Staņislavs Ježijs Lecs (\pltxti{Stanisław Jerzy Lec)}}

\epigraph
{Politikā noteikta mērķa labad var noslēgt savienību pat ar pašu velnu~--- vienkārši jābūt pārliecinātam, ka tu piemānīsi velnu, bet nevis viņš tevi.}
{Kārlis Markss (\detxti{Karl Marx})}

\epigraph
{Kad kādai valstij vēsturē nav panākumu, ar ko lepoties citu valstu un tautu priekšā, tiek runāts par šīs tautas upuriem, lai izsauktu citu tautu simpātijas un varbūt arī lai gūtu materiālus labumus.}
{Vitālijs Tretjakovs (\rutxti{Виталий Товиевич Третьяков})}

\epigraph
{Lai kā mums nepatiktu to atzīt, ir lielu un ir mazu valstu vēsture. To, ko mēs uzskatām par drausmīgu starptautisko tiesību pārkāpumu, izšķirīgos brīžos darījušas un joprojām dara daudzas citas lielvalstis ar daudzām citām mazākām valstīm. Tad tiek noslaucīts no galda jautājums par tiesībām un par veidu, kādā tas tiek darīts. Darbojas cita~--- spēka loģika.}
{Pēteris Krupņikovs}

\epigraph
{Nekad tik daudz nemelo, kā kara laikā, pēc medībām un pirms vēlēšanām.}
{Oto fon Bismarks (\detxti{Otto von Bismark})}

\epigraph
{Ne par vienu karu tik daudz nemelo un nav melojuši, kā par Otro pasaules un mūsu Lielo Tēvijas karu, jo to rezultāti neapmierināja ne uzvarētos vāciešus, ne rietumu sabiedrotos.}
{Jūlijs~Kvicinskis (\rutxti{Юлий Александрович Квицинский})}

\epigraph
{Pirmais kara upuris vienmēr ir patiesība.}
{H.~V.~Džonsons (\entxti{Hiram Warren Johnson})}

\epigraph
{Jo interesantāks laikmets vēsturniekam, jo laikabiedram bēdīgāks.}
{Nikolajs Glazkovs (\rutxti{Николай Иванович Глазков})}

\epigraph
{Lielā Kartāga veda trīs karus. Pēc pirmā tā joprojām bija varena. Pēc otrā tā joprojām bija apdzīvojama. Pēc trešā tā vairs nebija atrodama.}
{Bertolts Brehts (\detxti{Bertolt Brecht})}

\epigraph
{Katrs karš ir riebīgs un nejēdzīgs! Cēlums ir tikai atsvabināšanās kariem, kad tauta nomet kāda apspiedēja jūgu. Bet tas vairs nav karš, tā ir revolūcija.}
{Rainis}

\epigraph
{Godīgums politikā ir spēka rezultāts, liekulība~--- vājuma rezultāts.}
{Vladimirs Ļeņins (\rutxti{Владимир Ильич Ленин})}

\epigraph
{Kad visas citas tiesības tiek atņemtas, tiesības sacelties ir neapšaubāmas.}
{Tomass Peins (\entxti{Thomas Paine})}

\section{Polijas teritorijas kara sākumposmā 1939--1941}

\epigraph
{Nepastāv dabīgas valsts robežas.}
{Kārlis Poppers (\entxti{Karl Raimund Popper})}

\epigraph
{Baidies no tā, kurš no tevis baidās.}
{Persiešu sakāmvārds}

\epigraph
{Būtu bērnišķīgums nodoties ilūzijām: visus ierobežojumus, visus starptautiskos līgumus, kuri var tikt noslēgti miera laikā, kā sausas lapas noslaucīs kara vējš.}
{Džulio Duē (\frtxti{Giulio Douhet})}

\epigraph
{Tikai muļķi var būt nesatricināmi savā pārliecībā.}
{Mišels Montēns (\frtxti{Michel Eyquem de Montaigne})}

\epigraph
{Valdītājiem nevis jāvaino cilvēkus patriotisma trūkumā, bet jādara viss iespējamais lai viņi kļūtu patrioti.}
{Tomass Makolejs (\entxti{Thomas Babington Macaulay})}

\epigraph
{Ja mēs neatbildam par savu ciltsbrāļu darbību, tad arī jēdziens \emph{nācija} zaudē vispār jebkādu dzīves jēgu.}
{Aleksandrs Solžeņicins (\rutxti{Александр Исаевич Солженицын})}

\epigraph
{Karš nevar būt taisnīgs, jo taisnīgi karot nav iespējams, pat ja jūs karojat par taisnību.}
{Tadeušs Kotarbiņskis (\pltxti{Tadeusz Marian Kotarbiński})}

\epigraph
{Daudzskaitlīgākās armijas nebūt nav tās stiprākās. Gadās, ka uzvarētāja zābaks agrāk piederēja uzvarētajam.}
{Staņislavs Ježijs Lecs (\pltxti{Stanisław Jerzy Lec})}

\epigraph
{---} % TODO wat
{Emīls Žirardēns (\frtxti{Émile de Girardin})}

\epigraph
{Tumšos laikos ir grūti palikt ēnā.}
{Staņislavs Ježijs Lecs (\pltxti{Stanisław Jerzy Lec})}

\subsection{1939.~gada septembra kampaņa}

Kaut ir pagājuši jau gadu desmiti, taču, kā atzīmējis vācu vēsturnieks J.~Bolers, vēl joprojām ne viss ir skaidrs par Otrā pasaules kara notikumiem, pastāv daudzi vēsturiski mīti gan par 1939.~gada kampaņu (vācu: \detxti{Septemberfeldzug}, poļu: \pltxti{kampania wrześniowa}), gan vēlākajām norisēm.

Gatavošanās karam faktiski norisa visu starpkaru periodu. Pēc poļu-padomju kara Polijas armija tika samazināta no 1~miljona līdz 500--550~tūkstošiem cilvēku. Tika plānots to samazināt līdz 17~tūkstošiem virsnieku un 275~tūkstošiem karavīru, taču sakarā ar nestabilo situāciju Eiropā 1938.~gada vidū tajā bija 338~tūkstoši, bet 1939.~gada vidū~--- 434~tūkstoši virsnieku un karavīru. Dienestā ik gadus tika iesaukti ap 200~000~cilvēku. Tas nozīmēja, ka visi kara dienestam pakļautie tā arī nepabija armijā un daļa no tiem palika bez militārās apmācības.

Pēc vācu autora R.~Vingendorfa datiem Polijas armijā 1939.~gada sākumā bija 10~korpusu, kuri vienoja 90~kājnieku pulkus (273~bataljonus), 15~kavalērijas brigādes (40~pulkus ar 200~eskadroniem), 32~vieglās artilērijas pulkus (324~zirgu vilkmes, 4~motorizētās baterijas), 11~smagās artilērijas pulkus (98~zirgu vilkmes, 11~motorizētās baterijas), 6~zirgu vilkmes, 19~motorizētas zenītartilērijas baterijas. Ārpus korpusiem pastāvēja 5~bruņumašīnu un tanku pulki, 4~sapieru brigādes (14~bataljoni), dažādas citas palīgvienības. Aviācija ar aptuveni 1~000~lidmašīnām bija sadalīta trijās: Varšavas, Poznaņas un Krakovas grupās. Bez tam ar robežas apsardzi nodarbojās 42~bataljoni un 20~eskadroni. Pastāvēja arī 10~žandarmērijas pulki (70~eskadronu). Mūsdienās literatūrā var lasīt, ka pirms kara sākuma Polija mobilizēja 39~divīzijas, 16~atsevišķas brigādes, tai bija ap 870~tanku un tankešu, 4~300 lielgabalu un mīnmetēju, līdz 400~kaujas spējīgu lidmašīnu. Pastāv arī citi dati, kur daži rādītāji atšķiras, bet kopainu tas nemaina.

Armijas lielumu ierobežoja valsts budžeta iespējas. Polija miera laikā armijas vajadzībām tērēja vairāk nekā 30\% valsts budžeta, pēdējos četros gados pirms kara šī daļa sasniedza 50\%. Citas valstis, spēcīgākas un bagātākas par Poliju, militārām vajadzībām tērēja mazāk: Francija~--- 14\%, Lielbritānija~--- 16\%, Itālija~--- 22\%. (Pēc dažiem datiem, Ā.~Hitlers, forsējot gatavošanos karam, izmantoja 45~reizes vairāk līdzekļu nekā Polija.) Bez tam poļu iedzīvotāji, nereti uz savas labklājības rēķina, ziedoja aizsardzībai ievērojamus līdzekļus. 1939.~gadā Polijas militārie izdevumi tika palielināti, taču armijas mobilizācijas plāns līdzekļu trūkuma dēl saīsināts. Līdz 1939.~gada septembrim tas bija izpildīts tikai par 35\% un nespēja kvalitatīvi mainīt bruņoto spēku struktūru. Kara ministrijai armijas vajadzībām 1939./1940.~gadā bija iedalīti 800~miljonu zlotu. Naudas trūkums veda pie tā, ka likumā paredzētais 2~gadu dienesta laiks iesaucamajiem tika samazināts uz 18~mēnešiem.

Pēc poļu iestāžu aprēķiniem kara gadījumā valsts varēja rēķināties ar ap 4~miljoniem apmācītu karavīru. Taču visus tos nebija iespējams apgādāt pat ar vieglajiem ieročiem. 1939.~gadā Polijas kara rūpniecībā bija nodarbināti gan vairāk nekā 32~tūkstoši cilvēku, taču armijas vajadzības viņi nespēja apmierināt. Polijas kara rūpniecības mobilizācijas iespējas bija nepietiekamas. Pēc armijas štāba virsnieku ziņojumiem kara vajadzības valsts varēja apmierināt tikai daļēji: strēlnieku ieroču ziņā par 50\%, ložmetēju~--- par 45\%, artilērijas~--- 70\%, artilērijas munīcijas~--- 40\%, zenītartilērijas munīcijas~--- 30\%. Kaut viens no ``sanācijas'' režīma darbības virzieniem bija Polijas armijas pārapbruņošana, tas to tā arī nespēja paveikt.

Faktiski pat šauteņu un karabīņu pietika tikai 1,5~miljona karavīru apbruņošanai. Pat ložmetēju skaita dubultošana nebija iespējama. Poļu iestādes gan apgalvoja, ka ieroču pietiks, taču šī propaganda bija domāta galvenokārt tam, lai paaugstinātu Polijas kā sabiedrotā lomu ārvalstu acīs. Jau no 1939.~gada marta vojevodistu vadītāji armijas štābam lūdza piešķirt šautenes, kuru trūkums jau bija izjūtams. Pirmskara mēnešos galvenokārt pierobežā tika izveidoti daži Nacionālās aizsardzības bataljoni (\pltxti{Bataliony Obrony Narodowej}). To dalībnieki saņēma vienīgi šautenes, nekādus citus ieročus. Protams, kara laikā šādas vienības bija nolemtas milzīgiem neattaisnotiem zaudējumiem.

Moderna kara vešanai Polijas armijai trūka kustīguma. Strīdu ``zirgs vai motors'' bija izšķīris jau pats J.~Pilsudskis par labu zirgam. Par nākamā kara darbības vietu viņš uzskatīja PSRS teritoriju ar vāji attīstītu ceļu tīklu un lieliem līdzenumu attālumiem un viņa lēmumu ietekmēja ne tik daudz senā poļu tautai raksturīgā zirgu mīlestība, cik priekšstati, ka grūti izbraucamā apvidū motoram nav redzamu priekšrocību, zirgs ātrāk izvilks lielgabalu no dubļiem nekā traktors, kavalērijas vienība ātrāk atradīs ceļu purvā nekā bruņumašīna. Pie tam atkrita degvielas sagādes problēma. Nācās rēķināties arī ar poļu bruņojuma rūpniecības ierobežotajām iespējām.

1939.~gada septembrī poļiem bija ap 650~vieglo tanku. Polijas austrumu apgabalu apgāde ar degvielu ar valsts rīcībā esošajiem līdzekļiem un ceļiem bija pilnīgi neiespējama.

Polijas Republikas gaisa spēki (\pltxti{Siły Powietrzne Rzeczypospolitej Polskiej}) jau 1919.~gadā sāka ārzemēs iepirkt Pirmā pasaules kara laika lidmašīnas. Poļu lidotāji lidoja ar angļu, franču, vācu, austriešu un itāļu ražotām lidmašīnām. 1925.~gadā Polijā tika radīta aviācijas darbinieku skola. Jaunas lidmašīnas tika iepirktas galvenokārt Francijā, ražotas pēc licences arī Polijā. 1939.~gada sākumā poļu aviācija bija organizēta 6~pulkos, taču augustā šī organizācijas forma tika atcelta, atsevišķas eskadriļas tika nodotas sauszemes daļām. 1939.~gada februārī Polijas Seimā kara ministrs T.~Kaspžeckis paziņoja, ka Polijai ir lieliski lidotāji un labi konstruktori, izcilas lidotāju skolas un personāla izglītības iestādes, taču salīdzinājumā ar kaimiņvalstīm trūkst lidmašīnu.

Polijas aviācijas rūpniecība ražoja vairāku tipu lidmašīnas, taču tikai nedaudzi no tiem atbilda toreizējam pasaules līmenim.

1937.~gadā tika konstruēts un 1938.~gadā sākts ražot salīdzinoši labs poļu bumbvedējs \pltxti{PZL.~37 ``Los''} (Alnis). Pavisam līdz karam tika izgatavotas ap 120~šādas lidmašīnas. To maksimālais ātrums bija 440~km/h. (Taču jau 1938.~gadā Vācijā un Francijā uzbūvētie bumbvedēji sasniedza ātrumu lielāku kā 500~km/h.) Vēlāk kara laikā šo bumbvedēju pielietošana nenesa lielus panākumus. 24~no tiem tika notriekti lidojumā, 2~iznīcināti uz zemes. Tomēr divu nedēļu laikā ``Alņi'' veica ap 100~kaujas lidojumu, nometa 150~t bumbu, notrieca 6 vācu iznīcinātājus \detxti{Messerschmitt Bf.~109}.

Tāds pat trūkums bija iznīcinātājam ``\pltxti{Wilk}'' (Vilks), kurš ar ātrumu 460 km/h nevarēja sekmīgi cīnīties ar bumbvedējiem, kuru ātrums pārsniedza 500~km/h. Savam laikam atbilstošu iznīcinātāju ``\pltxti{Jastrząb}'' (``Vanags'') ražošanai bija jāsākas tikai 1940.~gadā. Īpaši smags stāvoklis bija motoru būves jomā. Poļu lidmašīnas, tāpat kā auto, tika apgādātas ar franču, angļu vai amerikāņu būvētiem motoriem, kurus lielāku remontu nepieciešamības gadījumā vietējie uzņēmumi nespēja savest kārtībā.

1939.~gada pavasarī Polijā izsludinātais valsts iekšējais aizņēmums aviācijas attīstības vajadzībām pildījās tikai ar lielām grūtībām. Vajadzēja pielietot piespiedu līdzekļus. Preču ievešanas no ārzemēm atļaujas saņēma tikai tās firmas, kuras varēja uzrādīt kvītis par iemaksām valsts aizņēmumam. Tomēr par spīti tam, aizņēmums ienesa tikai 400~miljonus zlotu. Poļu nedraugi vācieši to skaidroja ar to, ka poļu tautas finansiālās iespējas esot gaužām mazas, to liela daļa atrodoties ebreju rokās un uzticība Polijas valūtai esot tik vāja, ka kapitāli parasti tiekot ieguldīti nekustāmajā īpašumā, bet nevis valsts vērtspapīros. Acīmredzot tā bija daļa patiesības, bet tai jāpievieno arī neticība Polijas valsts stabilitātei sarežģītajos starptautiskajos apstākļos. Tiklīdz sāka izplatīties baumas par iespējamo karu, no apgrozības sāka pazust sudraba nauda. Tāpēc 1939.~gada jūlijā vajadzēja laist apgrozībā 50~miljonus tikko izkaltu sudraba zlotu. Sākoties karam, Polijas gaisa spēkos bija ap 750~lidmašīnu (literatūrā ir arī citi dati~--- 850 un pat vēl vairāk lidmašīnu), taču kaujas spējīgas bija tikai ap 400.

Polijas flotes lielākie kuģi bija iznīcinātāji (mīnkuģi), flotes kopējā tonnāža (18~000 t), ieskaitot arī zemūdenes, atbilda pusei viena liela līnijkuģa tonnāžas. Tā nespēja pilnvērtīgi pretoties Vācijas flotei. 1939.~gada 24.~augustā Lielbritānijas valdība nosūtīja maršalam E.~Ridzs-Smiglijam priekšlikumu evakuēt no Baltijas jūras uz Britu salām mūsdienīgākos Polijas kara flotes kuģus. Sākotnēji E.~Ridzs-Smiglijs atbildēja noraidoši, taču vēlāk mainīja savu lēmumu. Vēl līdz 1939.~gada 1.~septembrim Lielbritānijā ieradās trīs Polija mīnkuģi (``\pltxti{Błyskawica}'', ``\pltxti{Burza}'' un ``\pltxti{Grom}''), vēlāk tiem pievienojās zemūdenes ``\pltxti{Wilk}'' un ``\pltxti{Orzel}''. Pēdējā bija internēta Tallinā, taču spēja aizbēgt un kā \latxti{de~facto} jau nepastāvošas valsts flotes vienība 1939.~gada 14.~oktobrī ienāca Skotijas ostā Rosaitā (\pltxti{Rosyth}). Mīnkuģi ``\pltxti{Wicher}'' un mīnu izlicēju ``\pltxti{Gryf}'' nogremdēja vācu aviācija pirmajās kara dienās, divi tralētāji (``\pltxti{Mewa}'' un ``\pltxti{Rybitwa}'') karoja līdz septembra otrajai pusei. Trīs zemūdenes (``\pltxti{Ryś}'', ``\pltxti{Sęp''} un ``\pltxti{Żbik}'') pēc kauju beigām tika internētas Zviedrijā. Pašā Polijā pēdējā flotes vienība Helas pussalā padevās vāciešiem 2.~oktobrī. Vienīgais poļu flotes panākums bija vācu mīnu tralera M~85 uzspridzināšana, kurš uzdūrās zemūdenes ``\pltxti{Żbik}'' izliktai mīnai.

Austrumos diezgan iespaidīgu spēku sastādīja Pinskas upes flotile, kura 20.~gadu sākumā bija radīta kara gadījumam ar PSRS. Tās galvenais triecienspēks bija 6 monitori un 3 bruņotas lielgaballaivas. Vienā no flotiles baržām, kopā ar citiem kuģiem nogremdētā 1939.~gada septembrī un oktobrī izceltā, tika atrasta arī ķīmisko šāviņu noliktava. Grūti teikt, kā tie tiktu pielietoti citādas notikumu attīstības rezultātā.

Jau 1936.~gadā Varšavā Sakšu pilī, kur atradās Polijas Aizsardzības ministrija, tika debatēts par aizsardzības plānu ``\pltxti{Zahod}'' (Rietumi), gatavotu pret iespējamo Vācijas uzbrukumu. (Pastāvēja arī plāns ``\pltxti{Wschód}'' (Austrumi) karam pret PSRS). Taču tikai no 1939.~gada sākuma Polijas vadība rēķinājās ar Vācijas militāra uzbrukuma iespējām. Tikai no 1939.~gada marta Ģenerālštābs galveno uzmanību koncentrēja plāna ``\pltxti{Zahod}'' pilnveidei.

Sakarā ar Vācijas uzbrukuma draudiem poļu karaspēks slepenu mobilizāciju sāka 1939.~gada 23.~martā. Mobilizēja 4~divīzijas un 1~kavalērijas brigādi. Radīja vairāku armiju pavēlniecības. Tika pastiprināti garnizoni pierobežā. Pasākumu pamatā bija 1938.~gada aprīlī pieņemtais mobilizācijas plāns ``W'', kurš paredzēja daļēju mobilizāciju miera apstākļos. 23.~augustā sākās slepena pamatspēku mobilizācija, bet 26.~augustā~--- mobilizēto spēku pārgrupēšana uz apdraudētajiem rajoniem. 30.~augustā tika dota pavēle par izejas pozīciju ieņemšanu pirmajam ešelonam.

Mobilizācijas pasākumus Polijas vadība veica pat slepus no Rietumu sabiedrotajiem, kuri uzskatīja, ka šie pasākumi var pagrūst Vāciju uz karu. Kad 29.~augustā Polija gribēja pasludināt atklātu mobilizāciju, Anglija un Francija panāca, ka tā tika atlikta līdz 31.~augustam. Taču tā kā bija veikta slepena mobilizācija, mobilizācijas plāns bija izpildīts par 60\%. Kā rakstījis vācu vēsturnieks G.~Rode, vācu armijas uzbrukums Polijai nebija negaidīts. Izlūkdienests Polijas Ģenerālštābu par gaidāmo uzbrukumu bija informējis līdz pat atsevišķām detaļām, uzbrukums tika gaidīts jau 31.~augustā, tāpēc 1.~septembrī visas pie robežas izvietotās poļu daļas atradās kaujas gatavībā savās paredzētajās pozīcijās. Tomēr pēc citām ziņām poļu karaspēka izvēršana nebija pabeigta, 46,8\% poļu daļu atradās savos kara izejas rajonos, taču nebija paguvušas ieņemt tām paredzētās pozīcijas. 1939.~gada 1.~septembra rītā Polijai pierobežā bija 22~kājnieku, 3~kalnu strēlnieku, 10~kavalērijas divīzijas, 1~bruņu-motorizēta brigāde, kā arī vairāki Nacionālās aizsardzības bataljoni (Citi dati dod par dažām vienībām citādus skaitļus, bet kopainu starpība nemaina.) Uz 1939.~gada 1.~septembri Polijas bruņoto spēku skaitliskais sastāvs sasniedza 840~000 vīru, (Pēc citiem datiem~--- ap 900~000 karavīru), kad plānots mobilizēt bija 1~350~000.

Jau kara dienās Polijas armijas sastāvs sasniedza gandrīz pusotra miljona cilvēku. No mobilizācijas plānā paredzētajām 39~kājnieku divīzijām, 11~kavalērijas un 2~motorizētajām brigādēm vācu uzbrukuma rezultātā neizdevās mobilizēt tikai divas kājnieku (44. un 45.) divīzijas. Toties virs plāna tika saformētas mazākais 3~divīzijas, vairākas brigādes un kājnieku pulki un divas kavalērijas brigādes. Kopumā bija mobilizēti spēki, kas atbilda ap 55~divīzijām. Taču tām trūka bruņojuma, ne visas varēja iesaistīties kaujās.

Ja salīdzina pretinieku spēku samēru tikai pēc kaujās iesaistīto daļu skaita, tas nebūt nebija tik nelabvēlīgs Polijai. AP 60~vācu divīzijām un brigādēm pretī stāvēja ap 50~poļu divīzijas un brigādes, kā arī robežsardzes karaspēks (ap 3~divīzijas) un 50~Nacionālās aizsardzības (\pltxti{Oborona Narodova}) jeb zemessardzes bataljonu, tiesa, ar mazām kaujas spējām. Taču ļoti stiprs bija vācu pārspēks artilērijā, tankos un aviācijā.

Ikdienā noskaņojums Polijā bija pretrunīgs. 1939.~gada augustā daudzi Polijas iedzīvotāji bija nodarbināti ierakumu un nocietinājumu ierīkošanas darbos pie valsts rietumu robežas. Jau no 30.~gadu sākuma visā valstī parastas bija gaisa trauksmes un gāzu uzbrukuma mācības. Tagad tam pievienojās arī patversmju no gaisa uzbrukumiem, tanku šķēršļu joslu būvniecība. Taču līdz pat kara sākumam valdīja cerība, ka līdz atklātai kara darbībai tomēr nenonāks. Cilvēki nedz varēja, nedz arī gribēja stādīties priekšā, ko poļu tautai nozīmēs karš ar nacistisko Vāciju. Rietumvalstu garantijas, kaut arī tās pastāvēja tikai uz papīra, stiprināja optimistisku skatu uz nākotni.

Pirmajās septembra dienās Polijā izplatījās visdažādākās baumas, kuras neļāva iegūt daudz-maz skaidru priekšstatu par patieso stāvokli. Nebija skaidrs: vai Polija viena pretojas ienaidniekam vai Rietumvalstu karaspēks jau ir iemaršējis nacistiskajā Vācijā? Vai poļu daļas stāv jau pie Berlīnes vai arī vācu armija ir jau ielauzusies Polijas iekšienē? Vai vācu karavīri izturēsies vienaldzīgi pret vietējiem iedzīvotājiem, kā tas bija Pirmajā pasaules karā, vai arī no viņiem draud briesmas?

3.~septembris, kad Lielbritānija un Francija pieteica Vācijai karu, bija poļu vislielākā cerību diena. Ar Vācijas uzbrukumu Polijai karš sākās vispirms kā 2 valstu ierobežots karš . Tikai līdz ar Lielbritānijas un Francijas oficiālo iesaistīšanos karadarbībā lokālo Polijas~--- Vācijas karu no 3.~septembra pieņemts uzskatīt par Otro pasaules karu. Varšavā, kura jau bija piedzīvojusi pirmos gaisa uzbrukumus, cilvēku pūļi gavilēdami devās ielās, pulcējās pie Lielbritānijas un Francijas sūtniecībām. Cilvēkiem likās, ka glābiņš ir tuvu. Daudzās poļu dokumentālās filmās par Otro pasaules karu atkārtojas Polijas Telegrāfa aģentūras kinooperatora 1939.~gada 3.~septembrī nofilmētie kadri pie Francijas sūtniecības Varšavā. Milzīgs, entuziasma pārpildīts pūlis piepilda ielu: uz augšu paceltas dūres, prieka saucienos atvērtas lūpas. Caur pūli gandrīz nepamanīti virzās smagie auto, tie ved karaspēku uz fronti. Pāri visiem stāv neliela auguma Francijas sūtnis, kurš māj ar roku, pateicas par sveicieniem. Taču skatītājiem paliek nepamanīts, ka sūtnis klusē. Viņam nav ko teikt.

Polijā bija jūtams arī patriotisks pacēlums. Bijušie 1919.–1921.~gada Polijas karu veterāni, īpaši 1920.~gada brīvprātīgie, nu jau kā jaunākie virsnieki~--- poručiku (poļu \pltxti{porucznik}) un podporučiku (\pltxti{podporucznik}) pakāpēs atkal iestājās armijā 1939.~gada septembrī. Pastāvēja cerības karu noslēgt Berlīnē.

1939.~gada beigās krievu ``balto'' emigrantu avīze publicēja kāda līdz tam Polijā dzīvojoša krievu inteliģenta, bijuša ``balto'' virsnieka atmiņas par septembra notikumiem Polijā. Citēsim fragmentus no tām. ``Viss 1939.~gada augusts pagāja mobilizācijas skubā. Gan prese, gan vara, gan ierindas cilvēki pilnīgi nopietni apsprieda jautājumu par pilnīgu Vācijas sagrāvi. Lūk, izplatīts uzskats: ``Vācijā režīms brūk, revolūcija degungalā, bads, uz Poliju bēg tūkstoši (!!) vācu dezertieru; atradās ``aculiecinieki'', kuri ``savām acīm'' bija redzējuši šos tūkstošus vācu virsnieku un karavīru, pārejošus vācu-poļu robežu. Polijas armijai tikai ir jādod vienlaikus trieciens gan Austrumprūsijai un Berlīnei, un viss sabruks. Danciga tiks ieņemta dažu stundu laikā, pēc nedēļas mūsu kavalērija dzirdīs zirgus senajā poļu Kruļevcā (\pltxti{Królewiec}~--- Kēnigsberga), bet pēc divām nedēļām mēs būsim pie Berlīnes sienām. Protams, karš beigsies 2--3 nedēļās, ja nepiemānīs angļi un francūži, taču ja viņi arī tagad neiesaistīsies, tiksim galā arī bez viņiem. Ievērojot briesmīgus revolūcijas draudus, vācieši būs spiesti kapitulēt un Polija nospēlēs milzīgu vēsturisku lomu, atjaunojot to stāvokli, kāds pastāvēja līdz XVII gadsimtam, kad mūsu karaļi savām rokām piešķīra hercogu titulus teitoņu markgrāfiem.'' Iebilst tam visam nebija nekādas jēgas. Ja jūs visas šīs sarunas uztvērāt skeptiski, uz jums sāka skatīties ``šķībi'', ja jūs sākāt tās apstrīdēt, tad jūs varēja turēt aizdomās par nelojalitāti.''

Par Padomju Savienības iespējamo rīcību tika spriests tikpat ``pārliecinoši'': ``Boļševiki baidās no Polijas kā no uguns. Mums austrumos būtībā nav nepieciešama nekāda armija. Viens pats Robežapsardzības korpuss tiks galā ar robežas uzraudzīšanu. Vienīgais par ko bažīsies boļševiki, tad par to, lai Polija pēc Vācijas sakaušanas nevirzītos pret boļševikiem. Taču mēs viņus arī tāpat nostādīsim uz ceļiem. Kad Polija būs pilnīgi nodrošināta rietumus, viņai pēc Vācijas sagrāves būs jāpiebeidz boļševisms un Polijai būs jāizpilda otra vēsturiska misija~--- piebeigušai boļševismu, jāatjauno nacionāla Krievija, ar kuru mēs gribam dzīvot draudzībā un saskaņā''. Iespējams, ka pēdējā norāde par ``draudzību un saskaņu'' ar Krieviju bija speciāli izteikta sarunu biedra~--- krieva ausīm. Bet citādi jāatzīst, ka poļu vairākuma apziņu varas turētājiem bija izdevies novest līdz tādam pat ``patriotisma'' līmenim, kāds tas bija, piemēram, pirms Pirmā pasaules kara Vācijā. (Pirms iebrukuma Polijā 1939.~gadā Vācijā tautas sajūsmas uzplūdu nebija.) Diemžēl tautas masas vēl līdz šim laikam savtīgiem ``patriotisma'' sludinātājiem ļauj sevi itin viegli apkrāpt.''

Īpaši grūtas pārdomas mocīja Polijas ebrejus: satraucošas ziņas par ebreju vajāšanām Vācijā pēdējos gados izraisīja arvien lielākas bažas, bet vai tad arī Polijā etniskas nesaskaņas un konflikti nebija dienas kārtībā? Daudzi ebreji domāja, ka tiem vācu okupācijas gadījumā būs pat priekšrocības salīdzinājumā ar poļiem~--- viņi taču runāja vācu valodai tuvajā jidišā, kas lāva ar iekarotājiem saprasties. Pirmajā pasaules karā vācu armija viņus bieži izmantoja kā tulkus.

Pēc nacistiskās Vācijas uzbrukuma visā valstī sociālisti, komunisti, arodbiedrību darbinieki mītiņos un sapulcēs prasīja, lai valdība dod iespēju poļu strādniekiem ar ieročiem rokās stāties pretī iebrucējiem. Varšavā, Silēzijā, Piejūras novadā veidojās strādnieku pašaizsardzības vienības, kaut arī saņemt atļauju to organizēšanai nenācās viegli.

Vēlāk ``sociālistiskajā'' Polijā kā komunista un patriota simbols tika propagandēts bijušais leģionārs, cīnītājs par Polijas neatkarību, kurš agri vīlās izveidotajā valsts iekārtā un 1921.~gadā iesaistījās komunistiskajā kustībā, M.~Bučeks. Starpkaru periodā viņš tika trīs reizes arestēts, Polijas cietumos pavadīja 16~gadu, ar to kļūstot par politieslodzīto ar vislielāko stāžu. 1935.~gadā M.~Bučeks nepiekrita atteikties no Polijas pilsonības, lai izkļūtu no cietuma un dotos uz PSRS. Kara pirmajā dienā amnestēts, viņš pieteicās tuvākajā armijas daļā, taču tika atkārtoti arestēts un tikai pēc divām dienām atbrīvots. Viņš gāja bojā naktī no 9.~uz 10.~septembrī, kopā ar karavīru grupu piedaloties sadursmē pret vāciešiem Varšavas apkaimē. Tā, poļu komunists M.~Bučeks krita kaujā, aizstāvot to Poliju, pret kuras iekārtu pirms tam bija cīnījies. Domājams, viņš saprata, ka, cīnīties par brīvu Polijas attīstības ceļa izvēli var tikai pastāvot brīvai valstij.

Varšavā strādnieku partiju (kā PPS, tā arī bijušās PKP) aktīvisti organizēja strādnieku Brīvprātīgos bataljonus Varšavas aizstāvībai. Daļa no tiem tika reorganizēta Strādnieku brigādē, kurā bija ap 6~000 cilvēku. Līdzīgus brīvprātīgos formējumus radīja arī Gdiņas strādnieki. Tur viena no šāda bataljona rotām ieroču trūkuma dēļ apbruņojās tāpat kā XVIII un XIX gadsimta nacionālo sacelšanos dalībnieki~--- ar izkaptīm, vertikāli piestiprinātām kātos. Taču 12.~septembrī varas iestādes, baidoties no masu iniciatīvas, aizliedza brīvprātīgo pieteikšanos un strādnieku vienību tālāku veidošanu. Bija tūkstošiem brīvprātīgo, gatavu ar ieročiem rokās aizstāvēt Varšavu, taču virsnieki tikai plātīja rokas, atbildot, ka ieroču nav. Brīvprātīgie raka ierakumus un prettanku grāvjus Varšavas pievārtē, vēlāk tūkstošiem nonākot vācu gūstā.

1935.~gadā izstrādātais un 1938.–1939.~gadā papildinātais aizsardzības plāns paredzēja, ka Vācijas karaspēks virzīsies no Baltijas jūras piekrastes~--- Austrumprūsijas un Austrumpomerānijas virzienā un no Ziemeļsilēzijas Lodzas~--- Varšavas virzienā. Plāns paredzēja nodarīt maksimālus zaudējumus vērmahtam pierobežas kaujās, atkāpties uz Vartas (\pltxti{Warta})~--- Vidavkas (\pltxti{Widawka}) upju līniju, līdz tiktu sagaidīts Rietumu valstu uzbrukums Vācijai no aizmugures, kas to ierautu karā divās frontēs. Pēc tam Polijas armijai bija jādodas pretuzbrukumā uz Berlīni un Austrumprūsiju.

Šis plāns neievēroja ne iespējamo sabiedroto nevēlēšanos nekavējoties iesaistīties konfliktā, ne pēc Čehoslovākijas likvidēšanas radušos jauno iespēju vērmahtam uzbrukt Polijai arī no Slovākijas teritorijas, ne arī varbūtību, ka Polijai pašai var nākties karot divu fronšu karā pret Vāciju un PSRS.

Polijas teritorija nav piemērota plaša mēroga aizsardzības kara vešanai. Izņemot Poļesjes purvus austrumos un Karpatu augstieni dienvidos, tajā gandrīz nav dabīgu ūdens šķēršļu. No apmēram 5~700~km sauszemes robežas 2~700~km veidoja robeža ar Vāciju, 120~km~--- robeža ar Bohēmijas un Morāvijas protektorātu (vācu \detxti{Protektorat Böhmen und Mähren}; čehu \cstxti{Protektorát Čechy a Morava}), vairāk nekā 1~400~km~--- ar PSRS. Polijas–Vācijas robeža bija praktiski neaizsargāta. Polijai trūka līdzekļu lai tik plašā frontē celtu vērienīgas aizsardzības līnijas (esošās atradās galvenokārt Silēzijā un Cešinas apgabalā), tās militārā doktrīna paredzēja ātru karaspēka pārdislokāciju un pretuzbrukumus.

Ienaidnieks apdraudēja Poliju ne tikai no valsts rietumos esošās Vācijas pamatteritorijas, bet arī ziemeļos no Austrumprūsijas un dienvidos no Silēzijas, Morāvijas un Slovākijas. Tādejādi Polijas teritoriju faktiski no trijām pusēm aptvēra pretinieka karaspēks. Savā laikā jau franču ģenerālis M.~Veigans kara ar Vāciju gadījumā bija ieteicis poļiem atkāpties aiz Narevas~--- Vislas~--- Sanas upēm. Tādā gadījumā poļu karaspēkam būtu vieglāk aizsargāties aiz ūdens šķēršļiem, arī frontes līnijas garums samazinātos līdz apmēram 600~km, poļu armija neļautu sevi aplenkt vācu tanku ķīļiem, tiktu iegūts laiks sagaidīt Rietumvalstu palīdzību. Taču tas nozīmētu zaudēt Silēzijas ogļu ieguves rajonu, daudzas rūpniecības pilsētas (tekstilrūpniecības centru Lodzu u.c.), kā arī 30.~gados sākto attīstīt Centrālo rūpniecības rajonu (\pltxti{Centralny Okręg Przemysłowy}), kur bija salīdzinoši daudz uz militārās produkcijas ražošanu orientētu uzņēmumu. Kā vēlāk rakstīja vācu ģenerālis (pēc tam ģenerālfeldmaršals) E.~Manšteins, lēmuma par stratēģiski svarīgu rajonu faktiski brīvprātīgu atstāšanu varētu apdraudēt gan pašas Polijas valdības pastāvēšanu, gan arī nekādā gadījumā nestimulētu franču un angļu tautu vēlmi steigties palīgā poļiem. Pēc ģenerāļa kā militāra speciālista viedokļa Polija nedrīkstēja izvietot savus galvenos spēkus robežas tuvumā, kur tie tūlīt saņēma pamatīgus triecienus un tika ``samalti'', tos vajadzēja atvilkt tūlīt, kā bija redzami vācu armijas galvenie triecienu virzieni, lai koncentrētu spēkus tajos. Taču, kā pats ģenerālis atzīmējis, tas nespētu garantēt, ka Polija tā izbēgtu kaunpilno sakāvi, jo Rietumi taču pameta to savam liktenim un laika novilcināšana maz ko dotu.

1939.~gada 1.~septembrī pulksten 04:45, atbilstoši plānam ``\detxti{Weiß}'',Vācijas karaspēks uzsāka uzbrukumu visas vācu-poļu robežas garumā, kā arī no Sudetu, Morāvijas un Slovākijas teritorijas. Frontes līnija sastādīja ap 1~600 km. Dažas minūtes vēlāk Vācijas mācību līnijkuģis ``\detxti{Schleswig-Holstein}'' uzsāka poļu jūras kara bāzes Vesterplatē (vācu \detxti{Westerplatte}) pie Dancigas apšaudi. Šie lielgabalu šāvieni tiek uzskatīti par pirmajiem jaunajā karā.

Vācu karaspēks apbruņojuma un motorizācijas ziņā bija daudz pārāks pār Polijas armiju. Vācija raidīja pret Poliju 44 no 52~savām kadru divīzijām, tai skaitā 6~bruņutanku un 2~motorizētas; vairāk nekā 2~500 tanku un 1~390~kaujas lidmašīnu, kopā apmēram 1,5~miljonu vīru. T.s. ``Rietumu vaļņa'' aizsardzība pret Franciju bija uzticēta 25 rezerves divīzijām, kuras atbalstīja tikai 8 kadru divīzijas. Tur nebija neviena vācu tanka.

Vācija jau no paša sākuma veda totālu karu, neievērojot starptautiski atzītos un agrāk pierastos kara vešanas noteikumus. Jau 1.~septembrī tika bombardētas pilsētas un apdzīvotas vietas, neraugoties uz to, vai tajās atradās militāri objekti un ražotnes vai nē. Šis totālais karš tika vests ne tikai pret karaspēku, bet pret visiem pretinieka valsts iedzīvotājiem. Vācu aviācija Polijas kara laikā pilnīgi vai daļēji izpostīja vairāk nekā 150~poļu pilsētas un apdzīvotas vietas.

Kaut poļu vēsturnieki apgalvo, ka Polijas bruņotie spēki cīnījās ārkārtīgi sīvi, fakti tam neatbilst. Par 1939.~gada septembra cīņām baltkrievu vēsturnieks ar poļu saknēm A.~Usovskis raksta: ``Šis karš, neraugoties uz tā neparasti īso raksturu, ir pilns ar mītiem, ko ir ``uzcepuši'' poļu pēckara ``minhauzeni'' un Krievijā sakautie vācu ģenerāļi. Pirmie rakstīja par poļu armijas varonību, par varoņdarbiem, par gara cēlumu, otrie~--- par kliedzošām muļķībām, kuras šī kara dažu dienu laikā veica Polijas armija''. Tiesa, vācu ģenerāļu memuāros tiek augstu vērtētas Polijas armijas pretestība. Piemēram, E.~Manšteins raksta, ka ``poļi cīnījās ar ārkārtīgu drosmi pat bezcerīgās situācijās''. Taču vācu ģenerāļi bija ieinteresēti izcelt savus nopelnus uzvarā pār tik drosmīgu pretinieku, bet īsā kara gaita nebūt neliecināja, ka visi tik tiešām karoja ar degsmi un varonību. Tikai atsevišķās vietās bija vērojama sīva pretestība, tur ienaidniekam tika sagādāti arī smagi zaudējumi. Citur poļu armijas daļas atkāpās bez lielākas pretestības.

Pirms paša kara sākuma Polijas aviācija paguva pārbāzēties uz nelieliem nomaskētiem lidlaukiem. Vāciešiem no gaisa izdevās iznīcināt tikai treniņa un palīgdienesta lidmašīnas. 1.~septembrī jau ap pulksten 7.00~poļu lidotājs V.~Gniss notrieca (vai, kā uzskata daži pētnieki, vismaz sašāva) pirmo Vācijas lidmašīnu Otrā pasaules kara laikā.

Pēc otrās kara dienas Vācijas aviācijas štābs varēja ziņot, ka tai ir izdevies padzīt poļu aviāciju no tās miera laika bāzēm un sakarā ar izvietošanu nesagatavotos lidlaukos būtiski ierobežot tās darbību, taču iznīcināta kara pirmajās dienās poļu aviācija tomēr nebija. Tā turpināja karot, kaut bez izredzēm uz lielākiem panākumiem.

Vēsturē kā zināms varonības paraugs ir iegājušas cīņas pie \strong{Vesterplates} (\detxti{Westerplatte})~--- nelielas pussalas Baltijas jūrā Dancigas tuvumā.

% page 793



Tā kā padomju-poļu kara laikā Dancigas ostas dokeri atteicās izkraut Polijai atsūtītos ieročus un munīciju, to Polija izmantoja par ieganstu, lai 1924.~gadā ar Tautu Savienības atļauju blakus Dancigai Vislas ietekā jūrā Vesterplattes pussalā izbūvētu munīcijas noliktavas (\pltxti{Wojskowa Skłanica Tranzytowa}) un izvietotu garnizonu. No 1939.~gada maija līdz augustam garnizons, neievērojot Tautu savienības lēmumu, tika pastiprināts un sasniedza ap 210~vīru. No 1939.~gada pavasara vācieši gatavojās Vesterplates nocietinājumu ieņemšanai.

25.~augustā Dancigas līcī ar ieganstu atzīmēt šeit pirms 25~gadiem notikušo kāda vācu kuģa bojā eju iegāja vācu mācību līnijkuģis ``\detxti{Schleswig-Holstein}'', kurš bija bruņots ar smago artilēriju un uz kura atradās desantam sagatavota 225~cilvēku liela jūras kājnieku rota (\detxti{Marinesturmkompanie}). Pašā Dancigā tika gatavoti vietējie policijas (līdz 6~000 cilvēku) un SS (līdz 2~000 cilvēku) spēki. Šurp slepus nosūtīja ieročus. Vācu uzbrukuma plāns paredzēja poļu nocietinājumu apšaudi ar artilēriju, kam bija jāsalauž poļu karavīru pretošanās gars. Pēc tam pussalai uzbruktu Dancigas esesiešu bataljons, divas policistu rotas un jūras kājnieku rota. Šiem spēkiem bija jāieņem pussala. Taču poļu garnizons izrādīja sīkstu pretošanos. Vācu uzbrucēji cieta lielus zaudējumus. Vakarā bija jau 82~nogalināti un ievainoti (citi dati~--- 82~nogalināti un vairāk kā 100 ievainoti) vācu karavīri. Neraugoties uz Ā.~Hitlera prasību propagandas nolūkos pēc iespējas ātrāk ieņemt Vesterplati, uz triecienlidmašīnu Ju-87 uzlidojumiem, uz vēl vairāku vācu kuģu iesaistīšanu apšaudē, poļu garnizons noturējās vēl vairākas dienas. Tikai 7.~septembrī tika pacelts balts karogs.

Pēc kapitulācijas nocietinājumu komandants majors H.~Suharskis saņēma tiesības arī gūstā saglabāt virsnieka zobenu, kas liecināja, ka viņa virsnieka cieņu atzina arī pretinieki. (Var gan atrast arī materiālus, ka komandants, kurš bija saņēmis pavēli, ja iespējams, noturēties divpadsmit stundas un to izpildījis, otrajā aizstāvēšanās dienā bija gatavs padoties, taču pārējie virsnieki to nepieļāva, pat uz laiku atstādinot viņu no komandēšanas. Vācu vēsturnieks J.~Bolers šādu apgalvojumu parādīšanos saista ar pēc ``sociālistiskās'' sistēmas sabrukuma atklājušos dažu novēlotu ``ideoloģisko cīnītāju'' vēlmi ``atmaskot'' minēto sistēmu, kurai esot bijis izdevīgs no vienkāršas ģimenes cēlies varonis, pretstatā viņa vietniekam, kurš nācis no dižciltīgajiem un tiešām esot pelnījis apbrīnu.) Taču Vesterplates varoņi, kuru veikums ir iemūžināts daudzās filmās un grāmatās, neraugoties uz pēdējā laika debatēm, poļu tautā joprojām tiek pieminēti ar godu.

Poļu zaudējumi pie Vesterplates sastādīja 15 kritušos, 13 smagi un 25 līdz 40 viegli ievainotos. Vācieši turpretī zaudēja ap 400 nogalināto un ievainoto. Dažiem Vesterplates aizstāvjiem izdevās bēgt no vācu gūstekņu nometnēm un viņi turpināja karot Rietumos un Padomju Savienībā izveidotajos militārajos poļu formējumos, kā arī pašā Polijā AK (\pltxti{Armia Krajowa}) rindās, daži atkal nonāca kā vācu, tā padomju nometnēs. Daļu no vācu gūstekņu nometnēm atbrīvoja veselības stāvokļa dēļ. Galu galā lielākā daļa Vesterplates aizstāvju pārdzīvoja karu. Polijas Seims visus Vesterplates aizstāvēšanas dalībniekus par vīrišķību apbalvoja ar ordeni ``\latxti{Virtuti Militari}''. Mūsdienās 1939.~gada cīņu vietā ir ierīkots Vesterplates muzejs.

Pašā Dancigā tās gauleiters (\detxti{Gauleiter der NSDAP}~--- augstākais nacistu partijas vadītājs kādā apgabalā) A.~Forsters jau 1.~septembrī lika pilsētā izkārt plakātus, kuri vēstīja, ka brīvpilsētas statuss tajā vairs nav spēkā; izsludināja ``Likumu par Dancigas atkalapvienošanos ar Vāciju'' (\detxti{Gesetz der Wiedervereinigung Danzigs mit dem Großdeutschen Reich}) un paziņoja, ka visa vara atrodas viņa rokās. Tūlīt arī tika arestēti tūkstoši poļu un ebreju pilsoņu. (Daļu no ap 1~500 arestētajiem poļiem nošāva.) Nākamajā dienā nacisti pilsētas pievārtē sāka ierīkot Štuthofas koncentrācijas nometni (\detxti{Konzentrationslager Stutthof}).

Par poļu karavīru varonību it kā liecina fakti par poļu kavalērijas uzbrukumiem vācu bruņutanku daļām. Vienīgā joma, kur Polija varēja uzstādīt 11:1 lielu pārspēku pār vērmahtu, bija kavalērija. Gan Polijā, gan vērmahtā (un arī Sarkanajā armijā) zirgi tika galvenokārt izmantoti transportā, arī izlūkgājienos nereti tika sūtīti kavalēristi. Taču kaujās viņi parasti piedalījās kājnieku ierindā, kā to arī prasīja poļu armijas reglaments. Objektīvi kavalērijas daļas bija vienas no labākajām Polijas armijā, elites karspēks, labāk sagatavots, labāk bruņots (to pastiprināja ar artilēriju, ložmetējiem, bruņu tehniku), ar labākām manevra iespējām. 1939.~gada septembrī poļu kavalēristi veica, mazākais, sešus uzbrukumus jātnieku ierindā. Divos no tiem kaujas laukā atrādās arī vācu bruņu tehnika: tanki un bruņumašīnas.

Tā, 1.~septembrī kaujas pie Krojantiem (\pltxti{Bitwa pod Krojantami}) laikā 18.~Pomerānijas ulānu pulks pēkšņi kādā meža ielokā burtiski uzgrūdās vācu 76.~kājnieku divīzijas motorizētajam bataljonam un, savlaicīgi nespējot pagriezt zirgus atpakaļ, bija spiesti virzīties uz priekšu, uzbrūkot vācu kājniekiem, kuri bēga, bet tālāk poļi uzdūrās tankiem, kas nekavējoties atklāja uguni. Poļi bija spiesti atkāpties mežā, zaudējot apmēram 1/3 sastāva.

Mīts radās nākamajā dienā, kad itāļu kara korespondents I.~Montanelli iztaujāja vācu karavīrus par sadursmes norisi un aprakstīja poļu karavīru varonību, kuri ar zobeniem un pīķiem uzbrukuši vācu bruņu tehnikai. Daži ``aculiecinieki'' stāstīja, ka poļu kavalēristi ar zobeniem cirtuši vācu tankus, domājot, ka tie ir izgatavoti no papes. Arī vēlāk it kā fiksēti poļu kavalēristu uzbrukumi padomju tanku kolonām. Pēc saņemšanas gūstā poļu kavalēristi stāstījuši, ka virsnieki viņiem iegalvojuši, ka ``boļševikiem tanki ir izgatavoti no finiera''. (Galvot par šo ziņu patiesumu autors neuzņemas, bet poļu propagandas ietekmē šādu gadījumu iespējamība nav izslēgta.) Ziņas sāka strauji izplatīties. Nacionālsociālistiskā propaganda stāstu izmantoja lai tēlotu poļu karavīru pārlieko pašpārliecinātību un fanātismu. Vācu žurnāls ``\detxti{Die Wehrmacht}'' 13.~septembrī rakstīja, ka poļi nopietni nenovērtējuši vācu armijas spēku, bet poļu propaganda stāstījusi, ka vācu bruņu tehnika vispār esot bez bruņām, un to tad arī centušies izmantot poļu kavalēristi. 1941.~gadā vācieši par šo tēmu uzņēma propagandistisku filmu ``\detxti{Kampfgeschwader Lützow}'' (``Lutcova kaujas eskadrons'').

Par poļu kavalēristu uzbrukumu ar aukstajiem ieročiem vācu tankiem 3.~septembrī savos memuāros rakstīja arī ģenerālis H.~Guderiāns. Padomju Savienībā pēckara laikā šo gadījumu minēja kā poļu komandieru nesaprātīguma apliecinājumu. Kinorežisors A.~Vajda šos notikumus 1959.~gadā mākslinieciski atveidoja filmā ``\pltxti{Lotna}'' (zirga vārds), taču arī viņš nespēja izvairīties no propagandiskiem štampiem un šo filmu kritizēja kara veterāni. Pēc kara kritušajiem 18.~ulānu pulka kavalēristiem pie Krojantiem ir uzcelts piemineklis. Šeit ik gadus notiek jātnieku uzvedums, kurā tiek attēlota minētā kauja.

19.~septembrī pie Vulkas Venglovas (\pltxti{Wólka Węglowa}) 14.~ulānu pulka komandieris nolēma, izmantojot pārsteiguma momentu, jātnieku ierindā izrauties no vācu aplenkuma Varšavas virzienā. Taču tās izrādījās vācu motorizēto kājnieku pozīcijas, pie tam tanki atradās tuvumā. Kauja ilga 18~minūtes. Poļu kavalēristi izrāvās no aplenkuma, zaudējot ap 20\% sastāva~--- 105~kritušos un ap 100~ievainotos. Vācu zaudējumi sastādīja 52~kritušos un 70~ievainotos. Arī citos gadījumos uzbrukumi jātnieku ierindā parasti notika pēc pašu kavalērijas daļu komandieru iniciatīvas, tikai vienā gadījumā poļu kavalēristu eskadrons devās uzbrukumā kļūdas rezultātā, pārpratis vecākā komandiera pavēli.

Neraugoties uz atsevišķiem varonības pierādījumiem, jau pirmajās kara dienās Polijas karaspēks zaudēja vairākas pierobežas kaujas.

Kā uzskata baltkrievu publicists un vēsturnieks A.~Usovskis, daudzi Polijas armijas Galvenā štāba virsnieki jau no 2.~septembra raudzījās uz karu kā uz zaudētu. Šai dienā E.~Ridzs-Smiglijs izmetis frāzi, ka poļu armijas sagrāve ir nenovēršama. Tie, kas cerēja uz labāko, arī neredzēja kādu citu līdzekli, kā vērst notikumus Polijai pozitīvā virzienā, kā vienīgi efektīvu Rietumu palīdzību. E.~Ridzs-Smiglija direktīva Galvenajam štābam, izdota jau 3.~septembrī, noteica: ``Sakarā ar radušos stāvokli un problēmu kompleksu, kuras izvirzījusi notikumu gaita, mūsu bruņoto spēku atiešanas ass jāorientē ne vienkārši uz austrumiem, Krievijas virzienā, kuru pakts saista ar vāciešiem, bet uz dienvid-austrumiem, uz sabiedrotās Rumānijas un Ungārijas, kura izturas labvēlīgi pret Poliju, pusi\dots{}''. Ar to maršals orientējis savus virsniekus ne uz aizsardzības organizēšanu, bet uz bēgšanu Rumānijas un Ungārijas robežas virzienā. Baltkrievu autors šādu direktīvu nevis iznīcināt ienaidnieku, bet organizēt bēgšanu, uzskata par nodevīgu. Šeit gan jāsaka, ka Galvenajam štābam domātā direktīva karaspēka daļām netika sūtīta, tās saņēma konkrētas pavēles par pretestību ienaidniekam. Acīmredzot, vērtējot minēto direktīvu, publicists ņēmis virsroku pār vēsturnieku.

Kara sākumā nekārtības gadās visās armijās. (Ir zināms, ka ģenerālis H.~Guderiāns 1.~septembrī, vadot tanku korpusa uzbrukumu dažiem izolētiem poļu bataljoniem, bija spiests personīgi drudžaini apbraukāt savus pulkus, ar savu iejaukšanos glābjot situāciju, un tikai 2.~septembrī pirmās dienas panika viņa korpusā bija pārvarēta.) Tādu nekārtību daudz bija arī poļu armijā. Gadījās, ka poļu karavīri apšaudīja evakuējošos bijušos Silēzijas sacelšanās dalībniekus, kuri baidījās palikt vācu ieņemtajā teritorijā. Daudzviet skanēja saucieni; ``Nodevība! Spiegi! Diversanti raida signālus!''

Vācu motorizētās vienības sekmīgi grāva poļu armiju ``\pltxti{Łódź}'' un ``\pltxti{Kraków}'' aizsardzības līnijas. Jau 5.~septembrī vērmahta daļas pārrāva armijas ``\pltxti{Łódź}'' fronti. Tās brigādes un pulki pārvērtās bēgošu karavīru pūļos, starp kuriem uz austrumiem virzījās bēgošo vietējo iedzīvotāju bari. Pavēlniecība vairs nevarēja vadīt atkāpšanos, jo vairākums virsnieku bija pametuši savas daļas un arī bēga uz Lodzu un Varšavu. Armija ``\pltxti{Łódź}'' nespēja noturēt pozīcijas Vartas (\pltxti{Warta}) un Vidavkas (\pltxti{Widawka}) upju krastos.

Galvenokārt vācu aviācijas triecienu rezultātā poļu armija ``\pltxti{Modlin}'' bija dezorganizēta un sāka haotiski atkāpties Vislas upes virzienā. Kritisks kļuva arī armiju ``\pltxti{Kraków}'' un ``\pltxti{Prusy}'' (Prūsija) stāvoklis. 6.~septembrī tika ieņemta Krakova. 6.~septembrī Polijas armijas pavēlniecība pavēlēja armijām ``\pltxti{Pomorze}'' (Piejūra), ``\pltxti{Poznań}'', ``\pltxti{Kraków}'' un ``\pltxti{Prusy}'' atkāpties uz pozīcijām pie Vislas. Tai pat dienā valdība atstāja valsts galvaspilsētu un devās uz Ļubļinu, no turienes uz Kremeņecu, bet 13.~septembrī uz Zaļeščikiem (ukr. \uktxti{Заліщики}, poļu \pltxti{Zaleszczyki}, tagad Ukrainā)~--- pilsētu pie pašas Rumānijas robežas. Prezidents I.~Moscickis jau 1.~septembrī bija pārcēlies uz \pltxti{Blotas} ciemu pie Varšavas, 8.~septembrī~--- uz Olikas (poļu \pltxti{Ołyka}, ukr. \uktxti{Олика}) pilsētu Volīnijā, 14.~septembrī~--- uz Zaluču (\pltxti{Załucz}). 7.~septembrī arī Galvenais štābs ar virspavēlnieku priekšgalā pārvietojās no Varšavas uz Brestu.

Tāpēc būtiski pasliktinājās tā sakari ar armijām, pat ar Varšavā palikušo Galvenā štāba priekšnieku ģenerāli V.~Staheviču. No Brestas cietokšņa nebija telefona sakaru ar fronti, tikai pēc 12~stundām no Varšavas tika ierīkota telefona līnija ar vienu no armijām. No Varšavas uz Brestu pārvesto radiostaciju izveda no ierindas vācu bumbvedēju uzbrukumi, tāpēc ģenerālis V.~Stahevičs saņemtos no frontes ziņojumus uzdeva motociklistam pa ar bēgļiem pilnajiem ceļiem nogādāt Brestā, kur E.~Ridzs-Smiglijs un štābs pieņēma lēmumus, pēc tam tie tika sūtīti uz Bugas upes flotiles štābu, kuras rīcībā bija strādājoša radiostacija, kas tos nosūtīja uz Kara flotes štābu Varšavā. Tikai pēc tam tie nonāca pie V.~Stahēviča, kurš varēja nosūtīt virspavēlnieka rīkojumus karaspēkam, kad tie sen bija zaudējuši kādu jēgu.

Francijas pārstāvis pie Polijas armijas Galvenā štāba 10.~septembrī uz Parīzi ziņoja: ``Šeit valda pilnīgs haoss. Polijas virspavēlniecībai gandrīz nav sakaru ar karojošajām armijām un lielākajām daļām \citespace{} [tai] nav nekādas informācijas par ienaidnieka virzību un pat par savu pašu karaspēku tā informēta visai nepilnīgi vai vispār nav informēta. Ģenerālštābs [Tai laikā~--- Galvenais štābs~--- V.Š.] sadalījies divās daļas. Polijas armija faktiski ir sakauta jau pirmajās dienās.''

Poļu aviācijas iznīcinātāju brigāde, kura aizstāvēja Varšavu, 3.~septembrī notrieca 3~vācu lidmašīnas, 5.~septembrī~--- 9, 6.~septembrī~--- 15. Taču ar E.~Ridzs-Smiglija pārvietošanos uz Brestu turp tika pārcelta arī šī brigāde un arī daļa artilērijas. Tagad vācieši varēja bombardēt Varšavu daudz brīvāk, ko viņi arī darīja. Pavisam Varšavā, galvenokārt bombardēšanas rezultātā, gāja bojā ap 20~000 civiliedzīvotāju, taču acīmredzot maršala un ap viņu esošo ģenerāļu drošība tika vērtēta augstāk.

Padomju sūtnis Londonā I.~Maiskis dienasgrāmatā 7.~septembrī aprakstīja savus iespaidus par karu: ``Vācieši Polijā ir guvuši milzīgus panākumus. [Poļu] koridors ir ieņemts, Krakova ir kritusi. Lodza atrodas zem uguns, motorizētas vācu daļas jau pietuvojušās Varšavai. \citespace{} Poļu vājums ir pārsteidzošs. Viņi atkāpjas visur. Viņi nekur nevar palikt ilgāk par pāris stundām, ilgākais 1--2 dienas. Nav lielu cīņu. Tiesa, vācieši nomāc viņus ar motorizāciju un aviāciju, bet tomēr \dots{} Nē, ir skaidrs, ka lietas būtība nav tikai motorizācijā. Šķiet, ka Polijas valsts ir caurcaurēm sapuvusi. Tik lielā mērā sapuvusi, ka tās armija nespēj nodrošināt nekādu nopietnu pretošanos ienaidniekam pat ar saukli par valsts neatkarības aizsardzību.''

Taču arī šī situācija ilga tikai līdz 10.~septembrim, kad E.~Ridzs-Smiglijs devās Rumānijas virzienā, vispirms uz Volodimiru-Volinsku (\pltxti{Włodzimierz Wołyński}), 13.~septembrī uz Mlinovu (ukr. \uktxti{Млинів}) 15.~septembrī uz Kolomiju (ukr. \uktxti{Коломия}). No 12.~septembra vairs nepastāvēja vispārēja armijas vadība. Jau 9.~septembrī Polijas valdība ievadīja sarunas ar Bukaresti par iespēju atkāpties uz Rumāniju, 9.--11.~septembrī Polijas valdība uzsāka sarunas arī ar Francijas valdību par patvēruma piešķiršanu. Polijas zelta krājumi 13.--16.~septembrī tika pārsūtīti uz Rumāniju. Arī valdība, diplomātiskais korpuss, militārā pavēlniecība virzījās uz Rumānijas robežu. 16.~septembrī sākās poļu-rumāņu sarunas par Polijas valdības tranzīta iespējām uz Franciju. Tas viss tikai pastiprināja haosu valstī un grāva armijas cīņas garu.

Turpretī vērmahta vadītāji varēja lepoties ar panākumiem. Jau 5.~septembrī Vācijas sauszemes spēku Ģenerālštāba priekšnieks ģenerālis F.~Halders savā dienasgrāmatā ierakstīja: ``Ienaidnieks faktiski sakauts'' un 10.~septembrī atzīmēja, ka vācu karaspēks šķērso Bugas un Sanas upes, kas nākotnē kļuva par Vācijas austrumu robežu ar PSRS. Tādejādi desmit dienās vācieši teritoriāli gandrīz pilnībā bija sasnieguši savus mērķus.

8.~septembrī vācu 4.~tanku divīzijas pirmās daļas sasniedza Polijas galvaspilsētas Varšavas priekšpilsētas, bet maršals E.~Ridzs-Smiglijs pavēlēja visam poļu karaspēkam pēc iespējas atkāpties uz dienvidaustrumiem, lai izveidotu aizsardzības līniju netālu no Rumānijas robežas. Viņš esot cerējis, ka šajā mežainajā apkārtnē, iestājoties rudenim, vācu motorizēto daļu virzība kļūs lēnāka un caur Rumānijas teritoriju poļu armija varēs saņemt kara materiālus no Rietumu sabiedrotajām valstīm.

Starp citu, atkāpšanās steiga noveda pie tā, ka pie Jaslo (\pltxti{Jasło}-pilsēta Polijā, Piekarpatu vojevodistē), uzspidzinot tiltu un ierīkojot aizsprostu vācu daļām, tika izmantotas bumbas ar kara vešanas noteikumos aizliegto iprītu. Tā rezultātā 2~vācu karavīri tika nogalināti un 12~ievainoti. Tā kā tā bija atsevišķa poļu virsnieka kļūda, vācu armija atturējās veikt atriebības pasākumus, arī pielietojot ķīmiskos ieročus.

Pret Poliju karoja arī 1939.~gada martā pasludinātās Slovākijas Republikas armija. Slovākijas uzbrukums Polijai sākās reizē ar Vācijas karaspēka ieiešanu Polijā. Slovākijas armija ``\sktxti{Bernolak}'' (\sktxti{Slovenská Poľná Armáda skupina ``Bernolák''}) triju divīziju sastāvā ar ap 50~tūkstošu kopējo karavīru skaitu Slovākijas aizsardzības ministra ģenerāļa F.~Čatloša vadībā ieņēma tās zemes, kuras Čehoslovākija bija zaudējusi konfliktā ar Poliju 1920.--1924.~gadā, kā arī sakarā ar Minhenes vienošanos 1938.~gadā. 11~mēnešus pēc tam, kā Polija 1938.~gadā sagrāba Čehoslovākijas teritorijas, Slovākijas avīzēs parādījās fotogrāfijas: slovāku lidotājs uz notriekta poļu bumbvedēja fona, slovāku karavīrs ar šķiltavām aizdedzina Polijas karogu, apmierināti slovāku karavīri fotografējas uz poļu karagūstekņu fona. Faktiski slovāku sadursmes ar poļu armiju notika tikai pirmajā kara nedēļā, taču formāli konflikts beidzās 6.~oktobrī, dienā, kad Vācijas pavēlniecība oficiāli paziņoja par uzvaru karā. Kopā slovāku daļas karā zaudēja 18~kritušos, 46~ievainotos un 11~bez vēsts pazudušos.

Vācu militārā izlūkošana (\detxti{Abwehr}) cīņās pret Poliju iesaistīja arī speciāli sagatavotas neregulāras daļas, kuru vidū bija 4~000~pirms kara Slovākijā apmācītu ukraiņu nacionālistu no vienības ``Austrumgalīcija'' (\detxti{Ostgalizien}).

Kā norādījis krievu vēsturnieks A.~Širokorads, 1939.~gada septembrī vērsties pret Poliju ar bruņotu spēku bija gatava arī Lietuva, kura pie tās robežas uzstādīja visas trīs savas divīzijas. Polija tām pretī stādīja divas savas divīzijas. Taču PSRS valdība nevēlējusies, lai Lietuva draudzētos ar Vāciju pret Poliju un pēc attiecīga padomju diplomātiska demarša Lietuvas armija palika savās pozīcijās. Turpretī amerikāņu vēsturnieks G.~Veinbergs uzsvēris, ka Lietuva pati atteicās sekot Vācijas aicinājumam uzbrukt Polijai, jo tad tai būtu jānostājas nacistu pusē, kuri nesen tai bija atņēmuši Klaipēdu, pret Lielbritāniju un Franciju. Domājams, ka Lietuvas valdība nebija noskaņota riskēt ar iespējamo valsts ieraušanu karā un karaspēku pierobežā koncentrēja savas drošības interesēs, nevis lai līdzīgi Polijai 1938.~gada oktobrī mēģinātu gūt labumu no Eiropas pārdales.

10.~septembrī Vācijas armija bija ieņēmusi ap 40\% no poļu apdzīvotās teritorijas un gandrīz visus galvenos ekonomiskos centrus. Jau 1939.~gada septembra pirmajās dienās Ā.~Hitlers lika izsaukt no armijas pazīstamo juristu, kurš vēl pirms Ā.~Hitlera nākšanas pie varas bija uzstājies daudzos tiesas procesos, aizstāvot nacistus (vairāk nekā 40~procesos aizstāvēja arī personīgi Ā.~Hitleru), vēlāk pildījis dažādus, arī impērijas ministra bez portfeļa, pienākumus, H.~Franku un iecēla viņu par civilās administrācijas priekšnieku pie vērmahta virspavēlnieka Polijā.

Tomēr 10.~septembrī poļu armijas ``\pltxti{Poznań}'' un ``\pltxti{Pomorze}'' ģenerāļa T.~Kutšebas vadībā uzsāka kontruzbrukumu pie Bzuras upes. Tā sākums bija sekmīgs. Vāciešu pārsteigums un poļu skaitliskais pārsvars nesa poļiem panākumus pirmajā kaujas fāzē. Taču jau 12.~septembrī vācieši atkal pārgāja uzbrukumā, aplencot poļus. Tikai nelielai poļu armijas daļai izdevās izlauzties no aplenkuma, lai pievienotos Varšavas un Modlinas cietokšņa (\pltxti{Twierdza Modlin}) aizstāvjiem. Pēdējās poļu armijas paliekas pie Bzuras padevās vai tika sagūstītas līdz 19.~septembrim.

12.~septembrī vācu motorizētās daļas tuvojās Ļvovai. 14.~septembrī noslēdzās Varšavas aplenkums. 14.~septembrī arī sākās kaujas par Brestas cietoksni, kurš tika ieņemts 17.~septembrī. 16.~septembrī vācu armijas daļas, kuras uzbruka no ziemeļiem un dienvidiem, sastapās pie Bugas un tā ielenca lielāko daļu Polijas armijas. 18.~septembrī tika ieņemta Ļubļina.

Tur, kur poļu varas iestādes bija aizbēgušas un nebija pat kaut kādas pagaidu varas, izveidojās īslaicīgs varas vakuums, kura laikā notika laupīšanas, uzbrukumi mierīgajiem iedzīvotājiem. Pie tam ne vienmēr trūkumcietēji uzbruka turīgajiem. Šos notikumus pārdzīvojušie ebreji varēja stāstīt, kā viņu veikaliņus izlaupīja poļu kaimiņi, pirms vēl vācieši bija pārņēmuši kontroli pār konkrēto apdzīvoto vietu. Šīs laupīšanas faktiski visur turpināja vācu karavīri un policisti. Tā kā dažviet pirms vācu ienākšanas tika atvērtas cietumu durvis, blakus dezertējušajiem karavīriem pa zemi klīda arī kriminālas bandas, kuras tāpat savtīgos nolūkos izmantoja labvēlīgos apstākļus.

Jau 16.~septembrī Polija faktiski karu bija zaudējusi, lielākā daļa Rietumu un Centrālās Polijas atradās vācu rokās, ceļi bija pārpildīti ar bēgļu straumēm, gaisa telpu gandrīz pilnībā kontrolēja vācu aviācija, Polijas karaspēkā un sabiedrībā valdīja pilnīgs haoss un neizpratne. Kā secināja V.~Čērčils, divu nedēļu laikā Polijas armija kā organizēts spēks faktiski beidza pastāvēt. Tomēr Austrumpolijā un apgabalā starp Vislu un Bugu vēl cīnījās apmēram 25~poļu divīzijas~--- apmēram puse no sākotnējās Polijas armijas, pārpalikušie spēki, arī ielenktā un nežēlīgi bombardētā galvaspilsēta, turpināja izrādīt pretestību.

Pēc jau iepriekš citētā anonīmā krievu emigranta atmiņām par kara sākumu ``Polijas avīzēs parādījās bravurīgi ziņojumi ar virsrakstiem ``Poļu aviācija virs Berlīnes'', ``Britu flote pie Dancigas'', ``Angļu-franču spēku desants Austrumprūsijā'', ``Poļu kavalērija uzvaroši virzās uz Kēnigsbergu'', ``Zigfrīda līnija pārrauta, ceļš uz Berlīni atvērts''. Taču pēc dažām dienām nolija auksta ūdens šaltis. Darbojās taču radio, klausoties kuru bez abu pušu neizbēgamajiem pārspīlējumiem no neitrālajām raidstacijām varēja uzzināt patiesību, pienāca arī Rīgas avīzes. Dzīves īstenība lika sevi manīt. Taču līdz pat pēdējām dienām poļi neļāvās grūtsirdībai. [Avīzēs bija redzami virsraksti] ``Stratēģisks manevrs. Vācieši lamatās.'' ``Vācieši tiks iedzīti Vislā'', ``Vācu tanki iestrēgst poļu ceļos'', ``Visla~--- poļu Mažino līnija''.

Par kara norisi ierindas virsnieka acīm ļauj spriest Polijas armijas rotmistra (kavalērijas kapteiņa) J.~Kļimkovska atmiņas, uzrakstītas 1945.--1947.~gadā. Autors stāsta, kā viņš kara sākumā no savas brigādes tika nosūtīts uz Lodzu pēc pavēlēm, jo sakaru ar armijas štābu nebija. 4.~septembrī pulksten 3-jos dienā viņš ieradās Lodzā. Ap ģenerāļa J.~Rummela komandētās armijas štābu nebija nekādas sardzes un patruļu. Kāda ienaidnieka patruļa vai ukraiņu nacionālistu vai vietējo vāciešu vienība varēja saņemt gūstā armijas pavēlniecību, sagrābt operatīvos plānus un pavēles un neviens tam nevarētu pretoties. Nekādas pavēles J.~Kļimkovskis štābā nesaņēma, tikai solījumu tās nosūtīt brigādei tad, kad armijas štābs būs saņēmis instrukcijas no virspavēlniecības. Taču pavēles nepienāca arī turpmāk, brigāde paātrinātā tempā uzsāka atkāpšanos, neizrādot nekādu iniciatīvu ierīkot slēpņus vai kā citādi mēģināt aizkavēt ienaidnieka virzību Polijas teritorijā. Autors raksta: ``Nepieļaut sevi apsteigt, ielenkt~--- tā bija kaut kāda mānīja, kaut kāda psihoze, kura aptvēra mūsu komandieru prātus un dvēseles un aizsedza visu. Vārdu sakot, notika sacensība ar vācu bruņutanku daļām~--- kas ātrāk nonāks līdz Varšavai~--- viņi vai mēs. Neviens nedomāja par to, lai aizturētu ienaidnieku kaut uz dažām stundām, ja ne uz vairākām dienām vai ilgāk.'' Pēc divām ne visai smagu kauju dienām un pēc vairākām dienām marša bez kaujām un pat bez saskarsmes ar pretinieku komandiera neprasmes dēļ no brigādes kaujas sastāva, kā raksta autors, ``gandrīz nekas nepalika pāri''. 12.~septembrī pārpalikušie brigādes spēki saņēma pavēli nosegt t.s. ģenerāļa V.~Andersa grupas, kura līdz galam tā arī netika saformēta, aizmuguri, kamēr tā dos triecienu Mazoveckas Minskai (\pltxti{Mińsk Mazowiecki}). Trieciens neizdevās (sākotnēji vācu karaspēku izdevās atstumt no pilsētas, taču viņu pretuzbrukums iznīcināja visus poļu panākumus) un brigāde pa ceļiem, kuros valdīja pilnīgs haoss, atkāpās uz Ļubļinu. Pilsēta tika sasniegta 13.~septembra vakarā, taču 14.~septembrī izplatījās ziņa, ka karaspēkam tā jāatstāj, civilās un militārās varas iestādes to jau bija izdarījušas. Evakuācijas vispārējais virziens bija uz dienvidiem, Ungārijas un Rumānijas robežu, kur it kā notikšot karaspēka koncentrācija. 15.~septembra pēcpusdienā brigādes atliekas sasniedza Vlodavas (poļu \pltxti{Włodawa}, ukraiņu \uktxti{Володава}) pilsētu. ``Karavīru šeit bija kā skudru un nekārtība valdīja vēl lielāka, nekā kur citur. Neviens nekomandēja, nedeva pavēlēs, nepārzināja situāciju. Neviens nezināja, ko darīt tālāk, un, galvenais, kas bija pats briesmīgākais, neviens necentās pārvaldīt stāvokli uz vietas.'' Tālāk armijas daļas virzījās uz dienvidiem, līdz 21.~septembrī nonāca Zamostjes (poļu \pltxti{Zamość}, ukraiņu \uktxti{Замостя}) pilsētas apkaimē. ``Anderss pavēlēja pamest visus vezumus, visu, kas varēja kavēt mūsu daļas, pat lauku virtuves un vezumus ar munīciju. Kareivjiem bija jāņem līdzi tikai tas, ko varēja aiznest viņš pats vai viņa zirgs. Tādejādi tika iecerēts saīsināt kolonas garumu un atvieglot pāreju pār robežu. Kaut kāda trula sāpe mūs nomāca no šiem rīkojumiem. Bija pilnīgi skaidrs, ka nekas netika gatavots cīņai pret vāciešiem. Nekādas ziņas mēs nesaņēmām un nezinājām, kas notiek apkārt mums.'' Citētās rindas bija tapušas pēc kara vēl nenorimušo emociju iespaidā, taču neievērot notikumu dalībnieka liecību nevar.

Kaujas par Varšavu sākās 8.~septembrī, kad ap pulksten 17.~pirmās vācu daļas ielauzās pilsētā, bet tika atvairītas. Trieciena mēģinājumi tika atkārtoti 9.~un 10.~septembrī. 12.~septembrī, novērtējusi pretošanos, vācu pavēlniecība atteicās ieņemt pilsētu triecienā. Taču jau 7.~septembrī Polijas virspavēlniecība propagandas daļas vadītājs bija izdevis aicinājumu galvaspilsētas iedzīvotājiem, kuri vēlas ar ieročiem rokās karot pret iebrucējiem, pamest Varšavu, evakuēties uz Vislas austrumu krastu, tur iestāties Polijas armijas daļās. Rezultātā daudzi brīvprātīgie pameta Varšavu, tā novājinot tās spēkus. Ģenerālis V.~Anderss, kurš 11.~un 12.~septembrī bija Varšavā, savās atmiņās rakstīja: ``Milzīgs sajukums radās iedzīvotājiem doto pretrunīgo pavēļu dēļ. Dažas proklamācijas ieteica visiem, kas spējīgi nest ieročus, pilsētu atstāt, kamēr citas deva gluži pretējus rīkojumus.'' 14.~septembrī Polijas galvaspilsēta tika pilnībā ielenkta. 22.~septembrī sākās uzbrukums Varšavai, plaši izmantojot aviāciju, uzlidojumos piedalījās pāri par 1~100~lidmašīnu, nometot 5~818~tonnu bumbu. Kaujas par pilsētu norisa līdz 28.~septembrim, kad poļu pavēlniecība bija spiesta parakstīt kapitulācijas aktu.

Pēc angļu publicista L.~Deitona ziņām 30~vācu transportlidmašīnu Ju-52, lidojot pāri pilsētai, no kravas lūkām izmeta uz to degbumbas, radot lielus postījumus. Savukārt angļu vēsturnieks M.~Gilberts stāsta par pilsētas iedzīvotāju smago stāvokli ar pārtiku (kad bumbas sprādzienā uz pilsētas ielas tika nogalināts zirgs, jau pēc stundas no tā bija palikuši tikai kauli), par lielo ievainoto skaitu (no 140~000~poļu karavīru 36~000~bijuši ievainoti), par vācu karaspēka vilcināšanos ieiet pilsētā (tajā nebija elektrības, ūdens) pēc tās kapitulācijas, jo tas nevēlējās uzņemties atbildību par ievainotajiem un civiliedzīvotājiem. Šajā laikā daudzi nomira. 1.~oktobrī pirms ieiešanas pilsētā vācieši pieprasīja 12~ķīlniekus~--- 10~kristiešus un 2~jūdus, lai nodrošinātos, ka uz vācu karavīriem netiks šauts. Reizē viņi ieveda Varšavā arī kara lauka virtuves, no kurām dalīja viru iedzīvotājiem, taču, kad šis process tika nofilmēts, virtuves tika aizvestas.

Zīmīgi notikumi risinājās Polijas armijas aizmugurē. Jau no 1939.~gada sākumu, kad saspringa \strong{attiecības} starp Vācijas un Polijas valstīm, tās saasinājās arī \strong{starp poļu un vācu tautības iedzīvotājiem Polijā}. Vācu mazākuma stāvoklis Polijā 1939.gada septembra sākumā bija mokošs. Sakarā ar kara draudiem vācu mazākumu poļi uzlūkoja ar arvien lielākām aizdomām. Vietējos vāciešos viņi saskatīja potenciālus spiegus un diversantus. Tas, ka nacisti centās izmantot Polijas etniskos vāciešus saviem mērķiem, deva aizdomām zināmu pamatu. Taču, kā vērtējis mūsdienu vācu zinātnieks J.~Bolers, fanatizētu vācu nacionālistu~--- ``brīvības cīnītāju'' vietējo etnisko vāciešu (\detxti{Volksdeutshe}) vidū bija tikai apmēram 1\%. Taču 1939.~gada vasarā daudzi poļi boikotēja vāciešiem piederošos veikalus, traucēja viņu pasākumus, uzbruka viņu organizāciju mītnēm. Tāpēc strauji pieauga poļu un vācu sadursmes. Protams, poļu masu rīcība bija daudz iedarbīgāka, jo preses cenzūru, uzraudzību, arestus varas iestādes vērsa galvenokārt pret vāciešiem. Pirmoreiz pēc 20~gadiem daudzi Polijas vācieši atkal bija gatavi pagriezt savai mītnes zemei muguru un pāri robežai devās uz Vāciju. Taču katrs pieaudzis bēglis poļu acīs bija potenciāls vācu militāristu informants un tāpēc spiediens uz Polijā palikušajiem vāciešiem vēl vairāk pieauga. Polijas Ārlietu ministrijā 1939.~gadā vairākus mēnešus ar lielām bažām vēroja notikumu attīstību nacionālo attiecību jomā, jo dažu etnisko vāciešu (\detxti{volksdeutsche} jeb ``vāctautiešu'') dalība pret Poliju virzītos pasākumos varēja izraisīt poļu asu vēršanos pret visu vācu mazākumu. Septembra sākuma notikumi šīs bažas apstiprināja.

Polijai piederošajā Silēzijas daļā no 1,4~miljoniem iedzīvotāju ap 160~000~bija vāciešu. Apmēram 1/3~iedzīvotāju bija gan vācu, gan poļu saknes, viņi runāja abās valodās un vairāk vai mazāk jutās saistīti ar abām valstīm. Šie cilvēki izjuta lielu rūgtumu, noraugoties kā politiķi, ierēdņi, no kuriem daudzi bija no Polijas vidienes pārcēlušies uz robežapgabalu, lai šeit ``stiprinātu poliskos elementus'', karam sākoties, kā pirmie pakoja mantas un bēga no Silēzijas. Vietējie poļu iedzīvotāji caur asarām ironizēja: ``Skatieties, kā viņi mūs pamet nelaimē! Ar pāris koferiem viņi pie mums atbrauca, un tagad viņi bēg ar līdz malām piekrautām kravas automašīnām!''

Regulāro poļu karaspēku, apvienotu operāciju grupā ``Silēzija'' (\pltxti{Grupa Operacyjna ``Śląsk''}), atbalstīja neregulārie poļu ``Brīvprātīgo sacēlušos nodaļu'' (\pltxti{Ochotnicze Oddziały Powstańcze}) formējumi. Kad karaspēkam draudēja aplenkums, maršals E.~Ridzs-Smiglijs 2.~septembrī deva pavēli atkāpties no industriālā apgabala. Brīvprātīgo vienības atpalika, jo tās šī pavēle nesasniedza vai to dalībnieki nolēma pašu spēkiem aizstāvēt dzimteni. Iznāca tā, ka 3.~septembrī pie Katovices šīs poļu brīvprātīgo vienības sastapās ne tikai ar vācu regulāro karaspēku, bet arī ar vācu izlūkdienesta no Silēzijas vācu mazākuma veidotajām vienībām (\detxti{Freikorpskämpfern}). Ar to šeit var runāt par pilsoņu kara elementiem. Par pašu cīņu norisi ir saglabājušās pretrunīgas ziņas, tāpēc poļu pētnieki, cenšoties noskaidrot patiesību, ir izteikuši atšķirīgus viedokļus. Īpaši strīdīga ir epizode par Katovices izpletņu lekšanas torņa izmantošanu cīņai pret vācu iebrucējiem, kaut arī masu apziņā tā ir kļuvusi par leģendu, poļu brīvprātīgo varonības simbolu.

1939.~gada septembrī Polijā nebija partizānu kara pret vācu karaspēku, kaut vācu karaspēka daļu ziņojumos it kā šāda kara gadījumi minēti visai bieži. Ja šie ziņojumi atbilstu patiesībai, tūkstošiem poļu iedzīvotāju ar ieročiem rokās vajadzētu klīst pa vācu ieņemto teritoriju. Taču nedz karaspēka, nedz vēlāk ieradušās vācu policijas ziņojumos nav minēts neviens uz vietas ar ieroci rokās notverts civiliedzīvotājs. Minētā Augšsilēzijas un Katovices ieņemšana ir izņēmums, jo tur pret vērmahtu cīnījās arī poļu brīvprātīgo formējumi.

Kā ir izpētījuši vācu vēsturnieki, vācu karavīri pirms kara sākuma tika brīdināti, ka, ieejot Polijas teritorijā, uz viņiem jebkurā brīdī var tikt vērsti civilapģērbos tērptu teroristu uzbrukumi. Tāpēc nenākas brīnīties, ka nepieredzējušo karavīru vidū valdīja nervozitāte. Kad kaut kur pēkšņi atskanēja šāvieni, tie tikta uztverti kā uz konkrēto vienību raidīti, karavīri bija gatavi atbildēt draudiem, apšaudot jebko, kas likās aizdomīgs. Visu piecu Polijā iebrukušo vācu armiju darbības rajonos štābiem ienāca ziņojumi par šādiem ``uzbrukumiem''. Patiesībā iebrucēji apšaudīja nevis reālus partizānus, bet savu iedomu tēlus. Bija gan arī karavīri, kuri šādā šaudīšanā saskatīja savu biedru pieredzes trūkumu, nevis par to piedēvēja vainu vietējiem iedzīvotājiem. Piemēram, kāda vācu 30.~kājnieku divīzijas karavīra stāsts liecināja: ``Bieži apmešanās vietās īpaši ar tumsas iestāšanos sākās mežonīga šaudīšanās \citespace{} Tā lielākoties bija izskaidrojama ar savu vienību nervozitāti. Pēc tūlīt izskanējušām pavēlēm izlādēt ieročus šaušana ātri beidzās. Vai šāva vietējie iedzīvotāji, neviens pārliecinoši nevarēja konstatēt.'' ``Partizānu'' psihozes esamību labi apzinājās vācu ģenerāļi. Artilērijas ģenerālis V.~f.~Reihenau, kurš Polijas karagājiena laikā komandēja 10.~armiju, sev padoto korpusu komandieriem rakstīja: ``Armijas daļu nervozitāte par šaudīšanos aiz muguras palikušajā apkārtnē ir jālikvidē. To bieži izraisām mēs paši. Atsevišķi īstenu partizānu uzbrukumi nedrīkst vest pie tā, ka visa daļa pie vismazākā iegansta uzsāk nekārtīgu šaudīšanos. Māju dedzināšana kā atriebības pasākums ir aizliegta.''

Panikas izraisītie šāvieni nereti ķēra arī savējos, kas vēl vairāk pastiprināja bailes no neredzamajiem ``partizāniem'', kā arī radīja vēlmi atriebties vietējiem iedzīvotājiem par saviem ievainotajiem vai pat nogalinātajiem biedriem. Uz tikai pašu vācu karavīru fantāzijā eksistējošo ``partizānu karu'' tika atbildēts ar tūkstošu vietējo poļu un īpaši ebreju civiliedzīvotāju nošaušanu. Šādas izrēķināšanās ar nevainīgiem civiliedzīvotājiem notika paša vērmahta spēkiem, negaidot speciālo ``tīrīšanas'' vienību (\detxti{Einsatzgruppen)} ierašanos. Pat kādas šīs vienības vadība uz Berlīni ziņoja: ``Kaut ir fakts, ka vācu vienībām uzbruka civilisti, taču gala rezultātā vairs nevar noskaidrot, kurš uz kuru šāva.''

Vērmahta pavēlniecība arī ienesa savu devumu izrēķināšanās gaitā ar iedomātajiem civilteroristiem. 1939.~gada 12.~septembrī Vācijas Ģenerālštāba pulkvedis E.~Vāgners izdeva ``Rīkojumu par ieroču piederēšanu'' (\detxti{Verordnung über Waffenbesitz}), ar kuru okupētajos apgabalos uz rietumiem no Vislas vidusteces, un uz ziemeļiem no Narevas, kur vairs nenorisa cīņas, jebkuru ieroču atrašanās vietējo iedzīvotāju rokās tika pasludināta par sodāmu. Taču minētajā apkārtnē sakarā ar vācu straujo virzīšanos uz priekšu aiz frontes līnijas bija palikuši tūkstoši poļu karavīru. Ar ``Rīkojumu par ieroču piederēšanu'' viņi visi tika pasludināti par partizāniem. Pēc sadursmēm ar šādām poļu karavīru grupām, neraugoties uz to, ka viņi nēsāja uniformu, sagūstītie šo grupu karavīri tika nošauti kā partizāni.

Diemžēl jārunā arī par pretēju procesu: poļu represijām pret vācu tautības Polijas pilsoņiem. Vāciešu biezi apdzīvotajos Polijas apgabalos Pomerānijā un Poznaņā Polijas varas iestādes veica viņu masu arestus, priekšroku dodot izglītotās pilsonības pārstāvjiem.

Jau 1939.~gada 30.~augustā, vēl pirms Otrā pasaules kara izraisīšanās Polijā sākās ``neuzticamo elementu'' aresti. Berezas-Kartuskas nometnē un cietumos ievietoja no 1,5 līdz 2~tūkstošiem vāciešu, tai skaitā ap 500~sieviešu, no 2 līdz 3~tūkstošiem ukraiņu organizāciju biedru u.c. Aresti turpinājās arī karam sākoties. Transporta infrastruktūrai vācu uzbrukuma rezultātā sabrūkot, arestētie 10~000 līdz 15~000~cilvēku kājām tika dzīti austrumu virzienā. Apsargi par mazāko nepaklausību draudēja šaut. Ejot caur apdzīvotajām vietām deportējamajiem nācās dzirdēt sev domātas poļu lamas, draudus un dažkārt saņemt arī sitienus. Tādā veidā vietējie vācieši varēja konstatēt, ka poļu varas iestādes ir iecerējušas radīt viņiem tādas grūtības, lai pēc iespējas lielāks skaits neizdzīvotu. Domājams gan, ka poļu iestādes vispirms bija norūpējušās, lai pēc iespējas ātrāk un pēc iespējas lielāks skaits vācu izcelsmes pilsoņu, kuri tika uzskatīti par potenciāliem ienaidniekiem, tiktu nogādāti austrumos, tālāk no frontes līnijas, kur no viņiem nedraudētu briesmas. Taču par pienācīgu vācu tautības Polijas pilsoņu aizsardzību varas iestādes nebija parūpējušās. Deportējamie vienīgi cerēja, ka tos paglābs uzbrūkošais vācu karaspēks. Pie Lovičas (\detxti{Łowitz}) pilsētas viņi tik tiešām sagaidīja atbrīvošanu. Tie vācieši, kuri varēja palikt savās dzīves vietās, līdz vācu armijas daļu atnākšanai dzīvoja pastāvīgās bailēs no laupīšanām un vardarbības.

Pamatu poļu neuzticībai vietējiem vāciešiem radīja mazas to vāciešu daļas rīcība, kura gatavoja ieroču noliktavas, uzbrukumus Polijas iestādēm pirmskara dienās, kā arī iesaistījās paramilitārās formācijās, kuras piedalījās cīņās jau kara laikā. Mūsdienu poļu vēsturnieku uz drošības iestāžu dokumentu pamata veiktie pētījumi rāda, ka Otrā kara priekšvakarā, tā laikā un arī pēc tā bieži izvirzītie pārmetumi vācu mazākumam Polijā par spiegošanu Vācijas labā un sabotāžu lielākoties bija bez pamata. Daudzi vācieši sakarā ar 1939.~gada vasarā izraisīto kara psihozi īpaši centās sevi parādīt kā lojālus Polijas pilsoņus. Piemēram, viņu devums bija neproporcionāli liels ziedojumos Polijas Nacionālās aizsardzības fondam. Taču, neraugoties uz to, Polijā sākās vācu mazākumam piederīgo, īpaši vadošo personu, aresti, pret tiem vērsti ekscesi, kuri noveda pie vairāku tūkstošu vāciešu nāves. Polijas valdības pasākumi, vērsti uz valsts aizsardzību pret tās ienaidniekiem, lielākoties skāra pilnīgi nevainīgus cilvēkus un starp tiem arī tādus, no kuru puses nekādi draudi neizrietēja,~--- bērnus, sievietes, sirmgalvjus. Kara pirmo dienu haoss veicināja cietsirdību pret deportētajiem, tie kļuva par poļu nacionālistu dusmu izvirduma mērķi. Uzraudzības no augšas trūkums ļāva par dzīvības un nāves jautājuma izšķīrējiem kļūt tam nesagatavotiem apcietināto konvojētājiem, kuru acīs viņu apsardzībai uzticēto, nometnēs ieslodzīt paredzēto Polijas vāciešu dzīvībām nebija nekādas vērtības. Ar to ekscesi bija jau iepriekš programmēti, kaut arī no Polijas valdības puses tieši neizsaukti. Liela daļas līdzatbildības par vācu civiliedzīvotāju upuriem bija jāuzņemas arī tām Vācijas iestādēm, kuras, izmantojot dažus tūkstošus ``vāctautiešu'' saviem mērķiem, apdraudēja visu vācu nacionālo mazākumu Polijā.

Vispazīstamākais un vienlaikus poļu un vācu vēstures literatūrā visplašāk diskutējamais gadījums notika 1939.~gada 3.~septembrī Bidgoščā (\detxti{Bromberg}), ko vācieši vēlāk nosauca par \strong{``Brombergas asins svētdienu''} (\detxti{Bromberger Blutsonntag}). Pilsēta pēc Pirmā pasaules kara atradās Polijas teritorijā. Divos starpkaru gadu desmitos iedzīvotāju etniskais sastāvs izmainījās drastiski. 1919.~gadā vācieši sastādīja 80\% iedzīvotāju, bet 1939.~gadā tikai 6,4\%~--- nedaudz mazāk nekā 10~000, bija vācieši. Tomēr Bidgošča joprojām palika nozīmīgs vācu mazākuma kulturāls un politisks centrs, tur darbojās daudzas tā organizācijas. Vēl 30.~gados vietējo poļu un vāciešu vidū valdīja diezgan labas kaimiņu attiecības. Stāvoklis mainījās, nacionālais sasprindzinājums pieauga, tuvojoties karam. Jau iepriekšējā nodaļā rakstīts par vācu un poļu iedzīvotāju sadursmēm 1939.~gada martā. Karam sākoties, pilsētā tika apcietināti simti par aizdomīgiem uzskatīti vācu mazākuma pārstāvji un 2.~septembrī kājām viņi tika konvojēti uz valsts iekšieni. Šīs dienas vakarā vācu armijas daļas atradās jau 15~km attālumā no pilsētas. Savukārt tā kļuva par sava veida adatas aci, caur kuru pie robežas sakautās poļu daļas plūda uz aizmuguri. 3.~septembra rītā, ap pulksten 10-iem vairākās vietās pilsētā bija dzirdami šāvieni, kuri izraisīja paniku atejošajās poļu daļās. Tad ``atbildot'' uz vācu diversantu darbību, kuri it kā bija šāvuši uz poļu armijas daļām, karavīri ielauzās vācu mazākumam piederošajās ēkās, meklējot ieročus un uz vietas izrēķinoties ar to īpašniekiem vai vienkārši tur dzīvojošajiem vāciešiem. Līdz 4.~septembrim tika nogalināti, pēc dažu mūsdienu vācu vēsturnieku domām, mazākais, trīs simti vāciešu pilsētā un tikpat pilsētas apkārtnē. Citi vācu vēsturnieki vērtē, ka Brombergā vien vācu civiliedzīvotāju upuru skaits sniedzās no 1~200 līdz 1~500. Polijas vēsturi pētījušais latviešu vēsturnieks Ē.~Jēkabsons uzskata, ka kara pirmajās dienās, atriebjoties par šķietamu vai reālu dalību pret poļu karaspēku vērstā Vācijas diversantu darbībā, poļi nogalināja apmēram 300~vietējos vācu iedzīvotājus, savukārt vācu karaspēks uzreiz pēc ienākšanas~--- 300-400 vietējos poļu civiliedzīvotājus, bet turpmākajos mēnešos tika nogalināti vēl 5~000~pilsētas poļu iedzīvotāju.

No poļu ``atriebējiem'' cieta arī daži poļu tautības iedzīvotāji. Piemiņas vērts ir vietējais poļu iedzīvotājs B.~Pulkovskis, kurš centās apstādināt vardarbības plūdus, nostājās no savām mājām izrauto vācu kaimiņu priekšā un reizē ar viņiem tika nogalināts. Izdzīvojušie Bidgoščas vācu pārstāvji pēc kara apzvērēja savu nevainību.

Polijas Tautas republikā par notikumiem Bidgoščā runāja kā par ``vācu diversiju'' un uzsvēra pēc tam, jau no 5.~septembra, notikušo vācu izrēķināšanos ar pilsētas poļu iedzīvotājiem, kad arī dažu dienu laikā tika nogalināti vairāki simti cilvēku. Pēc oficiālām poļu ziņām iekšējās sadursmēs septembrī esot nogalināti 238 poļu karavīri un 223 vācu ``piektās kolonas'' piederīgie.

Poļu vēsturnieks K.~Pospešaļskis, pamatojoties uz dokumentālajiem materiāliem, ko pēc kara apkopoja Izmeklēšanas komisija par hitleriešu noziegumiem Polijā, nonāca pie secinājuma, ka 1939.~gadā kopumā par upuriem kļuva aptuveni 2~000 vācu minoritātes pārstāvji. Taču, vācu pētnieki, izpētot materiālus par minētajiem notikumiem, dod citus vērtējumus. Piemēram, vācu vēsturnieka K.~Jansena dati liecina, ka septembra kampaņas laikā Polijā gāja bojā ap 4~000 vietējo vācu tautības civiliedzīvotāju (daži izdevumi šo skaitli paaugstina līdz 5--7 tūkstošiem).

Ar nacisma attaisnošanas mēģinājumiem skandalozi slavenais vācu vēsturnieks E.~Nolte pat rakstīja, ka Otrais pasaules karš sācies ar poļu genocīdu pret vāciešiem, un, ja karš turpinātos ilgāk par trim nedēļām, tad diezin vai kāds vācietis Polijā izdzīvotu.

1989.~gadā vācu žurnālists G.~Šauberts pirmais no vācu autoriem atzina, ka, iespējams, tik tiešām visa pamatā ir bijis vācu diversijas akts. Mūsdienās poļu vēsturnieku vairākums uzskata, ka provokācija vai nu no vācu mazākuma vai no Vācijas aģentu puses ir notikusi.

Tomēr 2003.~gadā Bidgoščas akadēmijas Vēstures institūta direktors V.~Jastržembskis, kurš agrāk aizstāvēja tradicionālos poļu vēsturnieku uzskatus, tos mainīja un paziņoja, ka nekādas vācu diversijas nav bijis, bet poļu karavīriem vienkārši ``neizturēja nervi'' un viņi savas dusmas izgāza uz nevainīgiem civiliedzīvotājiem. Par to V.~Jastržembski asi kritizēja daži citi poļu vēsturnieki, kuri nevēlējās revidēt agrākos priekšstatus. V.~Jastržembskis tika atcelts no Vēstures institūta direktora amata. 2004.~gadā Polijas Nacionālās atmiņas institūts (\pltxti{Instytut Pamięci Narodowej~--- Komisja Ścigania Zbrodni przeciwko Narodowi Polskiemu}) vēlreiz veica starp 3.~un 5.~septembri Brombergā notikušā izmeklēšanu. Kā atzīmējis jau minētais vācu vēsturnieks J.~Bolers, tajā liels uzsvars tika likts uz poļu armijas ziņojumiem un karavīru izteicieniem, kas runāja par it kā Brombergā redzētajiem vācu partizāniem. Taču ar šādu dokumentu palīdzību vācu ``piektās kolonas'' darbību nevar nedz pierādīt, nedz arī noraidīt. Brombergā dzimušais vācu pētnieks H.~Rasmuss ir sastādījis pilsētā dzīvojušo un nogalināto vāciešu sarakstu. Tajā iekļauti 358 personu, tai skaitā 39~sieviešu un 55~sirmgalvju, vārdi. Kaut dažiem poļu vēsturniekiem vēl joprojām gribētos iztēlot, ka notikumi Bidgoščā bija vācu diversantu izprovocēti, H.~Rasmuss ir norādījis, ka šādi mēģinājumi neiztur nekādu kritiku, runāšana par vācu diversijām ir kas līdzīgs vācu leģendai par Vācijas sakāvi Pirmajā pasaules kara ``dunča dūriena [ko it kā veica revolucionāri] mugurā'' rezultātā. Kā konstatējis vācu vēsturnieks G.~Rode, nav pierādījumu vācu mazākuma naidīgai rīcībai pret Polijas valsti 1939.~gada septembrī. Pat tie poļu vēsturnieku, kuri, uzskata, ka vācu ``piektā kolona'' eksistēja, var atsaukties vienīgi uz Vācijas abvēra (bruņoto spēku izlūkdienesta) plāniem, bet nevis to īstenošanas rezultātiem.

Aculiecinieku liecības nedod skaidru atbildi uz jautājumu, bija vai nebija Bidgoščā jel kāds objektīvs pamats poļu vardarbībai. Nav arī pierādīts, ka tā bija plānota diversija vai vācu priekšējo izlūku izraisīta apšaude. Neapstrīdams ir tas, ka pilsētas tuvumā jau atradās vācu armija, šādos apstākļos dažu ``vāctautiešu'' šaušanai uz atejošajiem poļu karavīriem nebūtu nekādas militāras nozīmes un tā varētu tikai apdraudēt mierīgos vācu tautības civiliedzīvotājus. Ja šāva iesūtīti diversanti, tad pret tiem arī bija jācīnās. Un pat ja daži no vietējiem vāciešiem būtu piedalījušies provokācijā, tas poļu karavīriem nedeva nekādas tiesības izrēķināties ar simtiem nevainīgo. Jākonstatē, ka šai gadījumā Polijas puses karotāji izdarīja kara noziegumu. (Tomēr poļu mācību grāmatās Bidgoščas pilsētas ``aizsardzības grupas'' darbība tiek pasniegta kā ``lielisks'' Tēvzemes mīlestības piemērs.)

Ā.~Hitlers jau pirms uzbrukuma Polijai rēķinājās ar šādu notikumu attīstību. Jau 1939.~gada 19.~jūnijā Vācijas Ārlietu ministra biroja padomnieks P.~Kleists sarunā ar kādu vācu žurnālistu stāstīja: ``Diemžēl, mums jārēķinās ar lieliem vāciešu, kuri dzīvo Polijā, zaudējumiem. Nesen Hitlers paziņoja, ka par katru nogalinātu vācieti viņš pavēlēs nošaut simt poļus. Sekojoši, ja poļi sarīkos vāciešu izkaušanu, viņi tiks pakļauti pašām nesaudzīgākajām atbildes represijām''.

Polijas vāciešu upuriem bija jādod ``attaisnojums'' fīrera iecerētajām etniskajām tīrīšanām. Jau 1939.~gada 22.~augustā augstāko armijas komandieru sanāksmē Ā.~Hitlers starp citu paziņoja: ``Es izdevu pavēli un likšu nošaut katru, kurš kaut ar vārdu iebildīs pret principu, ka kara mērķis ir nevis kādas noteiktas līnijas sasniegšana, bet ienaidnieka fiziska iznīcināšana. Tāpēc esmu sagatavojis~--- pagaidām tikai Austrumiem~--- savas ``miroņgalvas'' vienības un devis tām pavēli bez līdzcietības un žēlsirdības sūtīt nāvē visus vīriešus, sievietes un bērnus, kam ir poļu izcelsme un kas runā poliski. Tikai tādā veidā mēs iegūsim mums vajadzīgo dzīves telpu…''

Mūsdienu poļu vēsturnieks un politologs J.~Krols ir norādījis, ka H.~Himlera pavēle par ``poļu sacēlušos'' un ``partizānu'' iznīcināšanu tika izdota jau 3.~septembrī, pirms izplatījās ziņas par izrēķināšanos ar vācu civiliedzīvotājiem Brombergā (3./4.~septembrī). Taču pavēle attiecās tikai uz vērmahta aizmugurē, nevis frontē esošajām vienībām. Tā rāda, ka nebūt ne ``Brombergas asinssvētdiena'' izraisīja vācu teroru pret poļiem.

1939.~gada novembrī sastādītā Vācijas Ārlietu ministrija dokumentā ``Poļu zvērības pret vāctautiešiem Polijā'' (``\detxti{Die polnischen Graueltaten an den Volksdeutschen in Poolen''}) rakstīja par 5~437 ``stingri pierādītām slepkavībām''. Taču šis vācu tautības upuru skaits Polijas iedzīvotāju vidū vērtējams kā pārspīlēts. To vidū bija ieskaitīti arī Polijas armijā mobilizētie un kritušie. Tai pašā gadā iznākušajā tā otrajā izdevumā upuru skaits jau bija pieaudzis uz ``vairāk nekā 58~tūkstošiem''. Pēc Ā.~Hitlera norādījuma 1940.~gada februārī tika runāts par 12~875 nogalinātiem un ap 45~000~bez vēsts pazudušiem Polijas vāciešiem. Tas bija tiešs, apzināts viltojums, lai attaisnotu nacistu noziegumus Polijā. Poļu nogalinātie vācieši tika pasludināti par nacionālsociālistiem un kritušiem kareivjiem par tautu (\detxti{Volkstum}). Poļu ``zvēriskuma'' ``pierādījums'' nacistiem bija vajadzīgs kā pamatojums, lai poļus iztēlotu par zemcilvēkiem (\detxti{Untermenschen}) un attaisnotu pašu vācu nacistu veiktās nelikumības~--- tās vietējo poļu un ebreju civiliedzīvotāju slepkavības, ko veica vācu karavīri, policisti un arī vietējie etniskie vācieši vācu okupācijas pirmajās nedēļās, jo, lūk, esot darbojies likums ``aci pret aci''. Jau 1939.~gada septembra otrajā pusē mazpilsētas Pšemislas (\detxti{Przemyśl}) apkārtnē pie Sanas upes vācu einzatcgrupas (\detxti{Einsatzgruppen der Sicherheitspolizei und des SD})~--- Vācijas Drošības policijas un Drošības dienesta speciālās vienības, kuras veica masu slepkavības iebrukuma Polijā laikā 1939.~gadā, Balkānu kampaņas 1941.~gadā un galvenokārt pret Padomju Savienību vestā kara (1941-1945) laikā) apšāva vairākus simtus ebreju. Šī akcija iezīmēja nacistu antisemītiskās politikas jaunu attīstības pakāpi.

Vērmahta Ģenerālštāba pulkvedis E.~Vagners 12.~septembrī deva rīkojumu deportēt ebrejus pāri Sanas upei uz zonu, kura ``Molotova-Ribentropa paktā'' bija paredzēta iekļaušanai padomju ``ietekmes sfērā''. Taču šai laikā vērmahts vēl bija iesaistīts nopietnās kaujās un rīkojums palika neizpildīts. Priekškarpatos esošā Peremišļa septembra vidū nonāca vācu rokās. No 16. līdz 19.~septembrim einzacgrupas pilsētā veica ebreju arestus. Viņus daļēji izveda ar automašīnām, daļēji kājām izdzina dažādās vietās pilsētas tuvumā un nošāva. Šaušana ilga vairākas dienas, precīzs upuru skaits nav zināms. Pilsētā izvietoto vācu armijas daļu komandieri zināja par notiekošo, taču neprotestēja. 19.~septembrī 14.~armijas pavēlnieks ģenerālis V.~Lists izdeva pavēli ``pārtraukt pasākumus pret ebrejiem'', kas acīmredzot liecināja, ka arī karavīri bija tādos iesaistīti. 24.~septembrī arī Vācijas sauszemes spēku virspavēlnieks V.~f.~Brauhičs izdeva pavēli, kurā viens no punktiem skanēja: ``armijai piederīgajiem ir aizliegta dalība policijas veicamajās eksekūcijās.''

Jāsaka, ka vācu ierindas karavīru vidū, jau sākoties kara gājienam uz Poliju, kopumā valdīja antisemītiski noskaņojumi, arī uzskats: ``Kur ir ebrejs, tas ir partizāns''. Viņi bija gatavi pielietot ieročus pret visiem aizdomīgajiem, bet ebreji tiem šķita īpaši aizdomīgi. Tiesa, kad pēc pāris nedēļām kļuva skaidrs, ka nekādas partizānu kustības nav, strauji samazinājās armijas izpildīto ``sodu'' skaits.

Pēc Otrā pasaules kara vērmahta ģenerāļi savos daudzskaitlīgajos memuāros apgalvoja, ka visas ``akcijas'' pret ebrejiem esot īstenojušas SS vienības, bet armijas ``karavīru mundieris palicis tīrs''. Bet vācu ģenerāļi arī apgalvoja, ka Vācijas uzbrukuma kari esot bijuši ``aizsardzības cīņas'' (\detxti{Abwehrkampf}) pret boļševismu un aziātismu (\detxti{Asiatentum}). Aukstā kara gados, kad Vācija bija sadalīta un Rietumvācijas apbruņošana bija Rietumu interesēs, šāda argumentācija Rietumu sabiedrotajos atrada dzirdīgas ausis. VFR veidoto bundesvēru balstīja bijušo vērmahta virsnieku pieredze un autoritāte. Lai to balstītu, bija nepieciešams uzturēt mītu par ``sliktajām'' SS un ``labajām'' vērmahta vienībām, rehabilitēt arī kara noziedzniekus. Taču no 20.~gadsimta 70.~gadiem Rietumu vēsturnieki vairs neatbalsta šo mītu. Ir izauguši pierādījumu kalni arī par vērmahta dziļo iesaisti holokaustā. Leģendu par ``tīro'' vērmahtu vairs neaizstāv neviens nopietns vēsturnieks.

Arī ``vāctautiešu pašaizsardzības'' dalībnieki dažkārt asistēja einzatcgrupu veiktajās civiliedzīvotāju nošaušanās, taču viņu gatavību līdzdarboties nevajadzētu pārspīlēt. Gadījumos, kad poļu uzbrukumi saviem kaimiņiem~--- vietējiem vāciešiem bija prasījuši upurus (kā Brombergā), tādai līdzdalībai bija emocionāla motivācija. Daļā vācu iedzīvotāju tiešām varēja pamosties atriebības jūtas pret ``apspiedējiem'' poļiem, ja arī tiešie kaimiņi tiem neko ļaunu nebija izdarījuši. Daļa izdzīvojušo vāciešu vienkārši centās izmantot okupācijas varas atbalstu, lai iedzīvotos uz poļu iedzīvotāju rēķina. Esošos vācu mazākuma izrēķināšanās aktus ar kaimiņiem~--- poļiem bez atriebības par bijušajiem vai iedomātajiem pāridarījumiem noteica arī atsevišķu indivīdu vēlme iegūt priviliģētu stāvokli ``jaunajā kārtībā''. Taču bija arī tādi vietējie vācieši, kuri izpelnījās pārmetumus par pārāk tuvām attiecībām ar poļu tautības kaimiņiem. Tie etniskie vācieši, kuri vairījās piedalīties vardarbībās, bija spiesti rēķināties ar agrāko biedru nicinājumu, izslēgšanu no partijas u.tml.

Tā, dažādos Poznaņas novados vietējie nacistiski noskaņotie vācieši pēc poļu armijas daļu atiešanas un pirms vērmahta daļu ierašanas pašaizsardzībai dibināja vietējas drošības vienības. Šī ``pašaizsardzība'' nereti pārauga vardarbībā pret poļu kaimiņiem. Jau pirmās kara nedēļas beigās pēc Ā.~Hitlera norādījuma vācu mazākuma vienības tika apvienotas centralizēti organizētā ``Vāctautiešu pašaizsardzībā'' (\detxti{Volksdeutsche Selbstschutz).} Tās uzdevums bija kopot visus nākotnē SS vienībām noderīgos, kā arī ``atriebība par ``Brombergas asinssvētdienu'' un citiem poļu izlēcieniem pirms vērmahta ierašanās.''

10.~septembrī Brombergas tirgus laukumā tika sarīkota poļu masu nošaušana, tādas pat notika arī citās pilsētās, bojā gāja tūkstoši poļu civilistu. Acīmredzot ``atriebības'' akcijas kļuva nekontrolējamas, tāpēc oficiālās vācu iestādes 1939.~gada novembrī pat paziņoja par šo vācu ``pašaizsardzības'' vienību atlaišanu, kas gan ievilkās līdz 1940.gada sākumam. Daļā to dalībnieku ieplūda SS un policijas sastāvā. Faktiski minētās etnisko vāciešu pašaizsardzības vienības darbojās kopā ar Vācijas Drošības dienesta (SD) vienībām.

Taču slepkavošanas pieaugums oktobrī un novembrī liecināja nevis par spontāniem atriebības aktiem, bet gan nacistu plānveidīgi īstenotu ``nepilnvērtīgo rasu'' iznīcināšanas programmu.

Turpretī, kā poļu un vācu vēsturnieki aprēķinājuši, no kara darbības sākuma tikai līdz militārās pārvaldes beigām (27.~oktobrī vara okupētajās poļu teritorijās pārgāja vācu civilās administrācijas rokās) vācu armijas, einzatcgrupu un t.s. vietējo vācu iedzīvotāju pašaizsardzības vienības Polijā nogalināja no 12~līdz 18~tūkstošu poļu civiliedzīvotāju. Amerikāņu vēsturnieks A.~Rossino rakstīja, ka vēlāk veiktā izmeklēšanā atklājās, ka vērmahta karavīri vien~--- nevis SS vai drošības policija~--- 64 gadījumos nošāva poļu gūstekņus, reizēm masveidā. Pēc vācu vēsturnieka G.~Juberšēra ziņām 55~dienu laikā tikai vērmahts vien veica 311~poļu militārpersonu un civilistu masu nāvessodu izpildi. Citas Vācijas struktūras šai laikā veica 764~soda akciju, kurās gāja bojā 24~tūkstoši poļu pilsoņu.

Vācu vēsturnieks H.~f.~Zitzevics stāstījis par gadījumu, kad kāds vācu okupācijas aparāta ierēdnis, bijušais nacistu partijas trieciennieks naktī no 1939.~gada 22. uz 23.~oktobri pilnīgi piedzēries lika poļu gūstekņus izvest no Inovroclavas (poļu \pltxti{Inowrocław}, vācu \detxti{Hohensalza}) tiesas cietuma, kur tie bija ieslodzīti, un pilnīgi patvaļīgi 55~no tiem nošaut. Pie kam viņš pats ar pistoli līdzdarbojās slepkavībā. Kad gadījums kļuva zināms un izsauca pārējo ierēdņu sašutumu, gauleiters A.~Geislers griezās pie Vācijas Tieslietu ministrijas ar lūgumu gadījumu noregulēt, lai nesabojātu vainīgā ierēdņa dzīvi, bet viņam ar goda vārdu lika apzvērēt, ka tas 10 gadus nelietos alkoholu.

PTR historiogrāfijā tika pausts, ka nacistu atriebības akcijās par vācu civiliedzīvotāju nāvi esot nogalināti 10~500~poļu un 13~000~ievietoti koncentrācijas nometnēs. Ir arī dati, ka līdz 1939.gada beigām nogalināti tika vairāk nekā 40~000~poļu un mazākais 7~000~ebreju. 1939.~gadā pasaules sabiedrība neko nezināja par masu nošaušanām Polijā.

Nebija vienkāršas arī \strong{poļu}, no vienas puses, \strong{un ukraiņu un baltkrievu}, no otras puses, \strong{attiecības}. Tā, ukraiņu tautības deputātiem Polijas Seimā 1939.gada rudenī, protams, nebija ne vismazākās nojautas par Molotova~--- Ribentropa pakta esamību. Kad Vācija 1939.gada septembrī uzbruka Polijai, ukraiņu deputāti, baidoties no poļu pūļu un poļu policistu iespējamajām atriebības akcijām un jauna naida uzliesmojumiem pret ukraiņiem, šajā poļu tautai smagajā situācijā Seimā pasteidzās paziņot par ukraiņu lojalitāti pret poļu valsti. Taču jau pirms šīs formālās lojalitātes izpausmes pēc poļu policijas iepriekš sastādītiem sarakstiem tūkstošiem ukraiņu inteliģences pārstāvju un ukraiņu organizāciju dalībnieku 1939.gada augustā tika ieslodzīti koncentrācijas nometnē Bereza-Kartuskā un Polijas cietumos. Arī 1.~septembrī Ļvovas un Pšemislas pilsētās poļu drošības dienests veica ukraiņu inteliģences arestus. Ļvovā un citviet Galīcijā poļu jauniešu bandas uzbruka ukraiņiem. Dažās vietās izcēlās stihiskas sadursmes starp no vācu armijas bēgošajām poļu armijas un poļu policijas atsevišķām vienībām un ukraiņu iedzīvotājiem. Parasti ukraiņi bija spiesti aizstāvēties pret poļu uzbrukumiem, kuri bez jebkāda pamatojuma par Polijas militāro sakāvi vainu uzvēla ``ukraiņu nodevējiem'', kaut gan ne ukraiņu tauta, ne arī ukraiņu nacionālisti nekādi nebija vainojami ne par ar Versaļas miera līgumu radītās daudznāciju Polijas izveidošanu, ne arī tās sagrāvi karā.

Polijas armijā mobilizētie ukraiņi un baltkrievi, sākoties karam ar Vāciju nebūt nealka mirt par ``neatkarīgo Poliju'', kur pret viņiem bieži izturējās kā pret darba lopiem (\betxti{быдло}). Par viņu noskaņojumu var spriest pēc kara dienās radušās tautas dziesmas~--- častuškas.

\vspace{1.5em}

\noindent
\begin{minipage}{0.45\textwidth}
\betxti{Вы ня думайце, палякi,\\
Вас ня будзем баранащь,\\
Мы засядзем у акопах\\
I гарэлку будзем пiць}
\end{minipage}
\hspace{2em}
\begin{minipage}{0.45\textwidth}
Jūs, nedomājiet, poļi,\\
Ka mēs par jums karosim,\\
Mēs ierakumos sēdēsim\\
Un šnabīti ierausim.
\end{minipage}

Polijas armijas kaujas spējas nevarēja neietekmēt arī virsnieku un Polijas iedzīvotāju atšķirīgais nacionālais sastāvs, ko rāda 17.tabula

\noindent
\begin{table}
\caption{PSRS Kozeļskas un Starobeļskas nometnēs ieslodzīto Polijas armijas 8~394 virsnieku nacionālā sastāva 1940.~gadā salīdzinājums ar pirmskara Polijas iedzīvotāju nacionālo sastāvu} \label{tab:table16}
\begin{tabularx}{\linewidth}{|p{3cm}|p{3cm}|p{3cm}|}
\cline{1-3}
\multirow{2}{*}{Tautība} & \multicolumn{2}{p{6cm}|}{Procentuālais sastāvs} \\
\cline{2-3}
& Visas Polijas iedzīvotāji & Virsnieki \\
\cline{1-3}
Poļi & 60,0 & 97,4 \\
\cline{1-3}
Ukraiņi & 21,0 & 0,1 \\
\cline{1-3}
Ebreji & 9,0 & 1,9 \\
\cline{1-3}
Baltkrievi & 6,0 & 0,3 \\
\cline{1-3}
Vācieši & 3,0 & 0,1 \\
\cline{1-3}
Pārējie & 1,0 & 0,2 \\
\cline{1-3}
\end{tabularx}
\end{table}

% page 807


Protams, uzskaitītie 8~394 cilvēki nebija visi Polijas virsnieki. Ja būtu pieejami dati par visa Polijas virsnieku korpusa nacionālo sastāvu, iespējams, atsevišķu mazākumtautību īpatsvars par vairākām procenta desmitdaļām vai pat procentu mainītos, taču neapstrīdams tik un tā paliktu poļu tautības virsnieku nospiedošais pārsvars pār citu tautību pārstāvjiem. Tātad Polijas armijas virsnieki pēc tautības bija gandrīz vai vienīgi poļi. Spēku Polijas armijai, tās autoritāti mazākumtautību iedzīvotāju acīs tas, domājams, nepalielināja. Tās pat nacionālās un šķiriskās atšķirības, kas radīja savstarpēju atsvešinātību pārējā Polija sabiedrībā, radīja to arī armijā. Poļu un Polijas pilsoņu~--- cittautiešu nacionālās pretišķības, sākoties karam, neizzuda. Tas bija viens no Polijas vājuma cēloņiem.

Vislielākās poļu cerības kara gaitā saistījās ar tās \strong{Rietumu sabiedroto} solīto militāro \strong{palīdzību}.

Lielbritānija un Francija iesniedza Vācijai ultimātu ar prasību pārtraukt kara darbību un izvest savu karaspēku no Polijas un brīvpilsētas Dancigas teritorijas. Kad tas netika izdarīts, kā jau agrāk teikts, 3.~septembrī tās pasludināja Vācijai karu, formāli izpildot savus sabiedroto pienākumus pret Poliju.

Dienu pirms tam, jau 2.~septembrī Francijas valdība izsludināja mobilizāciju un sāka koncentrēt savu karaspēku pie Vācijas robežas. Tomēr jāatzīmē, ka Lielbritānijas valdība, īpaši premjerministrs N.~Čemberlens, pēc Vācijas iebrukuma Polijā divas dienas šaubījās, līdz 3.~septembrī pārstāvju palātā paziņoja, ka Lielbritānija atrodas kara stāvoklī ar Vāciju. Var pateikt arī citādi, tā, kā to izdarījis angļu vēsturnieks A.~Teilors: ``Pārstāvju palāta ar spēku uzspieda karu šaubās esošajai angļu valdībai''. Lai paliktu pie varas, N.~Čemberlena valdība piekāpās parlamenta un tautas gribai.

Poļiem vissāpīgākais bija tas, ka nepiepildījās viņu cerības uz Rietumu palīdzību. Tā bija visai noturīgas. Vajadzīgas bija vairākas dienas lai poļi aptvertu, ka Rietumvalstu vārdiem nesekos darbi. Taču tikai tad, kad Polijā iegāja vēl arī Sarkanā armija, bet Polijas valdība un militārā pavēlniecība bēga uz ārzemēm, apkrāpto poļu cerības uz ātru un efektīvu Rietumu palīdzību sabruka galīgi.

1939.~gada maijā Francijas un Polijas ģenerālštābu pārstāvju noslēgtā vienošanās paredzēja, ka Vācijas uzbrukuma gadījumā Polijai, Francijai nekavējoties jāuzbrūk Vācijai no gaisa, trešajā kara dienā jāiesaista kaujā sauszemes spēki, bet piecpadsmitajā dienā uzbrukumā jādodas Francijas galvenajiem spēkiem. Tad pat Lielbritānijas valdības un Ģenerālštāba pārstāvji apņēmās šādā gadījumā pret Vāciju nosūtīt savu bumbvedēju aviāciju. Taču jāatgādina, ka vēl 1939.~gada 4.~maijā Lielbritānijas un Francijas militāristi savstarpēji bija vienojušies, ka, sākoties Vācijas karam pret Poliju, nesniegs pēdējai tās gaidīto palīdzību. 4.~septembrī Francijas un Lielbritānijas armijas komandieri apstiprināja šo lēmumu. Francijas Ģenerālštāba priekšnieks M.~Gamelēns 5.~septembrī deklarēja, ka Polijai nav palicis iespēju turpināt pretošanos.

1939.~gada 5.~septembrī, kad bijušais Lielbritānijas Admiralitātes pirmais lords L.~Emeri ieteica Lielbritānijas aviācijas ministram K.~Vudam organizēt Švarcvaldes mežu aizdedzināšanu Vācijā, tas sašutis atbildēja: ``Nevar būt ne runas pat par to, lai bombardētu militārās rūpnīcas Essenē, kuras ir privātīpašums, vai komunikāciju līnijas, jo tas atgrūdīs no mums amerikāņu sabiedrību.'' (``Amerikāņu sabiedrība'' pati gan nekaroja, bet juta līdzi Ā.~Hitlera pretiniekiem un no tās nostājas bija atkarīga iespēja saņemt dažāda veida palīdzību.) Francūži Saarbrikenes frontes iecirknī izkāra milzīgus plakātus: ``Mēs šai karā neizdarīsim pirmo šāvienu''. Notika daudzi vācu un franču karavīru brāļošanās gadījumi, viņi viens otru apciemoja, mainījās ar pārtiku un spirtotiem dzērieniem. Kad kāds iniciatīvas pārpilns franču artilērijas pulka komandieris sāka iespējamo mērķu piešaudi, viņš gandrīz tika nodots lauku kara tiesai. Korpusa komandieris viņu rāja: ``Tu saproti, ko izdarīji?~--- Tu gandrīz uzsāki karu!''

Jāņem vērā, ka jau 1939.~gada augustā Lielbritānijai un Francijai bija 57~divīzijas un 21~brigāde pret Vācijas 51~divīziju un 3~brigādēm. Tikai Francijas rīcībā 12.~septembrī bija 78~divīzijas un 18~atsevišķi tanku bataljoni. (Dažādos izdevumos skaitļi nedaudz atšķiras, taču tas nespēj būtiski mainīt spēku samēru). Taču lielākā, kaujas spējīgākā Vācijas karaspēka daļa vērsās pret Poliju.

Polijas ārlietu ministrs J.~Beks kopš kara sākuma pastāvīgi deva rīkojumus Polijas pārstāvjiem Londonā un Parīzē viskategoriskākajā formā pieprasīt sabiedroto aktīvu rīcību, kā arī pats centās pārliecināt angļu un franču vēstniekus Polijā par to, kādas sekas varētu būt kaut vai svarīgāko Vācijas dzelzceļa mezglu sabombardēšanai. Tomēr šie pūliņi nedeva nekādus rezultātus, izņemot deklaratīvus paziņojumus un kārtējo reizi neizpildītus solījumus. 6.~septembrī jau Polijas armijas virspavēlniecība lūdza sabiedrotos veikt aviācijas triecienu Vācijas rūpniecībai un armijai, 7.~septembrī tie atbildēja ar piekrišanu. 10.~septembrī tika paziņots, ka aviācijas triecieni ir sākušies, taču tas neatbilda patiesībai. 8.~septembrī Francijas armijas vadība nolēma nesūtīt uz Poliju nevienu lidmašīnu, nebombardēt arī objektus Vācijā. Francijas aviācija aprobežojās ar izlūklidojumiem, bet 4.~septembrī veiktais 10~britu lidmašīnu uzlidojums Ķīlei nedeva rezultātus, tika zaudēta puse lidmašīnu.

Francijas rīcībā bija ap 3~300~lidmašīnu, no kurām apmēram puse bija jaunākās paaudzes, pilnā mērā atbilstošas kaujas prasībām. Arī Lielbritānija jau ar pirmo kara dienu varēja uzsākt Vācijas rūpniecības centru graušanu. Karā sākumā tikai metropolē vien tās rīcībā bija 1~476~lidmašīnas. Vēl 435~atradās kolonijās.

Taču bumbu uzbrukumi, izņemot minēto mēģinājumu bombardēt Ķīli, Vācijai nenotika. Briti nesteidzās palīdzēt Polijai, kurai bija devuši ``garantijas''. Britu karaspēks divu korpusu sastāvā tikai oktobra vidū ieņēma pozīcijas uz Francijas un Beļģijas robežas, diezgan patālu no frontes līnijas.

Francijas Ģenerālštāba priekšnieks ģenerālis M.~Gamelēns jau 5.~septembrī paziņoja, ka Polijai nav izredžu tālāk pretoties un ``tas ir pamats mūsu spēku saglabāšanai'', tāpēc arī nav jāpiegriež vērība sabiedrības vispārējam sašutumam par palīdzības nesniegšanu Polijai un nav jāuzsāk kara darbība. Jau 7.~septembrī Lielbritānijas armijas štāba priekšnieks V.~Aironsaids informēja valdību, ka karš turpināsies vismaz trīs gadus, tādēļ nav lietderīga nekavējoša Lielbritānijas iesaistīšanās karadarbībā. 9.~septembrī valdība šo nostāju akceptēja.

Tomēr, acīmredzot lai mazinātu sabiedrības sašutumu, 9.~septembrī Reinas upes ielejā sākās franču spēku uzbrukums. Tas aprobežojās ar piecu divīziju priekšējo vienību darbību dažu km dziļumā, ieņemot 12~apdzīvotas vietas. Vācieši bez pretošanās atkāpās, tā izvairoties no zaudējumiem un iesaistīšanās nopietnās kaujās. Francijas armija pat nenokļuva līdz Vācijas nocietinājumiem~--- Rietumu valnim jeb Zigfrīda līnijai (\entxti{Westwall} jeb \detxti{Siegfriedstellung}).

12.~septembrī Lielbritānijas Augstākā Kara padome pieņēma lēmumu nesniegt palīdzību Polijai, jo ar to nevar Poliju glābt no sakāves. Briti baidījās, ka pēc uzvaras pār Poliju Vācija drīzumā virzīs uzbrukumu uz rietumiem un spēki būs vajadzīgi šī uzbrukuma atvairīšanai. Tai pat 12.~septembrī pirms Lielbritānijas un Francijas premjeru tikšanās Abvilā (franču \frtxti{Abbeville}, pilsēta Francijā) eksperti secināja, ka Polija karu ir jau zaudējusi. Kad Abvilā notika sabiedroto vadības apspriede ar N.~Čemberlena, E.~Daladjē un M.~Gamelēna dalību, tika nolemts pārtraukt uzbrukumu ar pamatojumu, ka ``notikumi Polijā neattaisno tālāku kara darbību Sārā''; tātad tika nospriests, ka palīdzēt Polijai ir jau par vēlu. Polija par šo lēmumu netika pat informēta. Rietumvalstu rīcība kara pirmajās nedēļas ļauj izdarīt secinājumu, ka to militārās saistības sakarā ar Polijai sniegtajām garantijām pēc būtības bija Polijas krāpšana. Tā gan nebija ļaunprātīga, jo ar garantijām, palīdzības solījumiem Polijai Rietumvalstis centās atturēt Vāciju no uzbrukuma tai, arī stiprināt pašas Polijas cīņas sparu, taču gala rezultātā tomēr sekmēja Polijas vadības kļūdas stāvokļa novērtējumā un veda pie Polijas armijas sagrāves nevienlīdzīgā cīņā.

16.~un 17.~septembrī vācu armija atguva gandrīz visas līdz tam zaudētās teritorijas, jo franču spēki atkāpās aiz Mažino līnijas. 17.~oktobrī pēdējie franču karavīri atstāja Vācijas teritoriju. Vācu pusē visās šais sadursmēs uz Francijas robežas bija 196~kritušie un 356~ievainotie, bija zaudētas 11~lidmašīnas. Pavisam četros kara mēnešos rietumos Vācija zaudēja 696~kritušos, ievainotos un pazudušos bez vēsts.

Sākušos karu Francijā sauca par \strong{``dīvaino''} (\frtxti{drote de guerre}), Vācijā~--- ``sēdošo'' (\detxti{Sitzkrieg}), ASV~--- ``šķietamo'' vai ``neīsto'' karu (\entxti{phoney war}). Tas faktiski turpinājās līdz 1940.~gada maijam, kad sākās Vācijas uzbrukums rietumos. Gandrīz 9~mēnešus 5~miljoni franču un britu karavīru nīka nocietinājumos un ciematos pie franču-vācu robežas. Neviens nedomāja ne tikai par kara darbību, bet pat par nopietnām militārām apmācībām. Komandieri vairāk domāja kā organizēt karavīru atvaļinājumus, viņu izklaidi.

Pasaules vēsturē ir ne mazums piemēru, kad uzticams sabiedrotais pilda savu pienākumu pat tad, ja tas pirmajā brīdī nenāk par labu viņam pašam. Tā 1914.~gadā I~pasaules kara laikā Krievijas karaspēks, cenšoties palīdzēt Francijai, pie kuras galvaspilsētas Parīzes stāvēja Vācijas karaspēks, uzsāka uzbrukumu Prūsijā, nepabeidzis mobilizāciju. Nesagatavotais uzbrukums beidzās ar divu Krievijas armiju sagrāvi, taču Vācija bija spiesta pārsviest daļu savu spēku uz Austrumu fronti, kā rezultātā zaudēja kauju pie Marnas, bija spiesta atkāpties no Parīzes. Francija izturēja, turpināja karu, kas atbilda arī Krievijas ilglaicīgajām interesēm. Šo Krievijas impērijas rīcību 1914.~gadā Lielbritānija un Francija 1939.~gadā nevēlējās ņemt par paraugu.

Kā norāda krievu vēsturnieki, prasīt, lai ``civilizētās nācijas''~--- angļi un franči rīkotos, sekojot morāliem apsvērumiem, būtu naivi, taču tām bija jārīkojas ``saprātīgā egoisma'' principa vadītām. Ja arī tām trūka spēku, lai tūlīt sakautu hitlerisko Vāciju, tās varēja uzsākt kara darbību ar esošajiem spēkiem, lai tā palīdzētu Polijai, saglabātu to ilgstošāk kā sabiedroto, vājinātu Vāciju un vienlaikus pašas varētu vākt spēkus ilgstošam karam. Taču tā vietā Polija burtiski tika atstāta viena pret daudz spēcīgāko ienaidnieku.

Padomju vēstures literatūrā, tāpat kā tās masu propagandā, Lielbritānijas un Francijas un dalība kara darbībā no 1939.~gada septembra līdz 1940.~gada maijam tika vērtēta līdzīgi, kaut tas neatbilda īstenībai. Ja Francijai galvenā bija sauszemes robeža ar Vāciju un tā atklāti izvairījās no aktīvas kara darbības, tad Lielbritānijai galvenā spēku pielikšanas joma bija karš uz jūras. Uz Lielbritānijas floti terminu ``dīvainais'' vai ``sēdošais'' karš nevar attiecināt. Tā karoja un darīja to pat visai intensīvi. Vācija bija aplenkusi Britu salas ar zemūdenēm, kuras ik dienas gremdēja 2--3 britu tirdzniecības kuģus. Lielbritānijas kara flotes uzdevums bija tirdzniecības ceļu aizsardzība. Tā cieta ievērojamus zaudējumus, taču bloķēja Vācijas kuģniecību un neļāva Vācijas zemūdens flotei ``nožņaugt'' savu valsti ekonomiski. Arī Vācijas flote cieta ievērojamus zaudējumus. Taču nevar arī apgalvot, ka Lielbritānija, līdz ar savu aizsardzību centās arī maksimāli palīdzēt Polijai.

Uz Lielbritānijas un Francijas kara pret nacistiem pirmo fāzi pilnībā var attiecināt slaveno K.~Klauzevica aforismu: ``Karš ir politikas turpinājums, tikai ar citādiem līdzekļiem''. Domājams, ka Lielbritānijas un Francijas valdošajos slāņos un lielā mērā arī visā sabiedrībā darbojās vēl ``agresora nomierināšanas'' politikas recidīvs. Briti un franči sprieda: ja jau karš tieši neskar mani un manu valsti, tad tajā var arī tūlīt neiejaukties. Vienmēr taču palika cerība uz kādu ``debesu dāvanu''~--- piemēram, Ā.~Hitlera nāvi, nacistiskās Vācijas konfliktu ar staļinisko Padomju Savienību vai vēl kādu neprognozējamu pavērsienu. Protams, par to nedrīkstēja atklāti runāt, nodevību pret sabiedroto neviens taču negrib labprāt atzīt, bet Lielbritānijas un Francijas tuvākās egoistiskās intereses pārsvēra visus citus apsvērumus.

Poļu vēsturnieks P.~Večorkovičs uzskatīja, ka Polijas vadība nerēķinājās ar to, ka Lielbritānija un Francija nebija gatavas karam. Viņš arī devis savu šo valstu rīcības iztulkojumu. Abām nosauktajām lielvalstīm esot bija vajadzīgs laiks, lai panāktu bruņošanās ziņā Vāciju un tās bija gatavas to iegūt par katru cenu. Situācija Eiropā izskatījusies tā, it kā nestos ragavas ar medniekiem (frančiem un britiem), kuriem vēl bija jāpielādē bises, bet tiem pakaļ dzenas vilku bars (Vācija). Lai glābtu savu ādu, medniekiem vajadzēja vienu pēc otra nomest vilkiem upurus. Pirmais upuris bija Austrija, otrais Čehoslovākija, par trešo kļuva Polija. Tikai šajā gadījumā britiem, kuri jau bija nolēmuši karot, optimāls risinājums bija scenārijs, pēc kura Polija pretotos. Tādejādi, nevis Lielbritānija iesaistījās karā, lai izglābtu Poliju, bet Polija glāba Lielbritāniju. Lielbritānija jau gandrīz no paša sākuma zināja par slepeno protokolu pie Molotova-Ribentropa pakta, taču nepaziņoja par to poļiem. Ja Polija zinātu par Vācijas un PSRS vienošanos, tai atliktu tikai kapitulēt Vācijas priekšā. Tādejādi Lielbritānija un Francija, zinot par to, kas notiks 17.~septembrī (par Sarkanās armijas iesoļošanu Polijā), vienkārši norakstīja Poliju zaudējumos. Protams, šī versija nav pierādīta ar dokumentiem un acīmredzot britu un franču vadītāji vairāk cerēja uz Lamanšu un Mažino līniju kā šķēršļiem, kas apturēs Vācijas agresiju rietumu virzienā, taču vielu pārdomām tā rada.

Poļu profesors A.~Novaks, sekojot Rietumu stereotipam un mēģinot kaut daļēji attaisnot Rietumvalstis, bet visu vainu par Polijas nelaimēm uzvelt PSRS, radiointervijā sakarā ar 70.~gadadienu kopš nacistiskās Vācijas sakāves paziņoja, ka ``iespējams, Franciju un Angliju ietekmēja fakts, ka Padomju Savienība vienojās ar Hitleru. Un, iespējams, franču un angļu militārie priekšnieki nolēma, ka nav iespējas palīdzēt Polijai. Ja draudzības [starp PSRS un Vāciju] nebūtu, varbūt tas ietekmētu Rietumu sabiedroto rīcību.'' Citiem vārdiem sakot, ja PSRS, neraugoties uz Polijas noraidošo attieksmi pret sadarbību ar to, tomēr būtu gatava karot par Polijas interesēm, ziedot tam savu pilsoņu dzīvības, ``varbūt'' arī Lielbritānija un Franciju nelauztu savus solījumus Polijai. Šāda savu sabiedroto nodevības attaisnošana tikai liecina par dubultstandartu pastāvēšanu poļu sabiedrībā.

Jāpasaka, ka karot par ``sanācijas'' Poliju nevēlējās arī rietumu komunisti. 1939.~gadā Francijas Komunistiskā partija kvalificēja karu kā imperiālistisku no abām pusēm, tās deputāti parlamentā balsoja pret kara kredītu piešķiršanu, partija uzsāka pretkara kampaņu, aicināja karavīrus dezertēt. Partijas ģenerālsekretārs M.~Torēzs, iesaukts armijā, dezertēja un bēga uz PSRS, franču kara tiesa viņam piesprieda nāves sodu. (1940.~gadā viņš tika apžēlots un atjaunoja politisku darbību).

Turpinot arī pēc 3.~septembra realizēt veco politiku, Rietumvalstis palaida garām vislabvēlīgāko izdevību kad Vācija rietumos pret Franciju turēja tikai ierobežotas pieseguma daļas. Vācijas armijas grupā rietumos kopā ar rezervēm nebija pat puses to spēku, kurus kara darbībā uzreiz varētu iesaistīt viena pati Francija. ``Dīvainais karš'', ko veda Lielbritānija un Francija, faktiski nozīmēja sekošanu tēzei, lai katra valsts aizstāvas pret agresiju, kā prot un var. To, ka nedrīkst ļaut nacistiskajai Vācijai sakaut savus pretiniekus pa vienam, britu un franču vadītāji tai brīdī nevēlējās redzēt. Lūk, ko par Lielbritānijas un Francijas izturēšanos šajās traģiskajās dienās rakstīja V.~Čērčils: ``Pēc mobilizācijas abu pušu armijas stāvēja gar robežu aci pret aci bezdarbībā. Nebija nekādu gaisa uzlidojumu nedz Anglijai, nedz Francijai, izņemot dažus izlūkošanas lidojumus. Francijas valdība lūdza mūs tāpat rīkoties attiecībā pret Vāciju, lai atbildes bombardēšana nevērstos pret Francijas rūpniecību. Mēs tāpēc samierinājāmies ar lapiņu nomešanu, kurām vajadzēja atmodināt vāciešu morālās jūtas \citespace{} Francija un Lielbritānija palika pasīvas, kaut tai pašā laikā vācu kara mašīna ar visu spēku grāva un paverdzināja Poliju. Hitleram nebija pamata žēloties.'' Rietumu sabiedrotajiem 1939.~gada septembrī bija unikāla iespēja, izmantojot savu īslaicīgo ievērojamo pārsvaru Rietumu frontē, dot Vācijai triecienu, kas varētu būtiski mainīt visu Eiropas tautu likteni. Taču šo iespēju Rietumi neizmantoja, kam bija katastrofālas sekas 1940.~gada maijā.

Var piebilst, ka 1939.~gada augustā Ā.~Hitlera pozīcijas Vācijas valsts priekšgalā vēl nebūt nebija tik stabilas kā pēc vēlākajām uzvarām. Ir zināms, ka vācu militāristi neveiksmīga kara gadījumā ar Rietumvalstīm gatavoja sazvērestību pret Ā.~Hitleru, kuras iznākumu nevarēja paredzēt. Ja ar kara sākumu pret Poliju izvērstos arī karš Rietumos, nav izslēgts, ka sazvērnieki no vācu ģenerāļu vidus varētu panākt fīrera atstādināšanu, lai noslēgtu mieru ar Rietumvalstīm.

Rietumvalstis centās vienīgi minimalizēt zaudējumus, ko tiem nesa gaidāmā Polijas sakāve. Francijas sūtnis Varšavā L.~Noels jau 6.~septembrī ieteica pārcelt Polijas valdību uz Franciju, 11.~novembrī viņš apsprieda šo jautājumu ar J.~Beku, vienlaikus vedot sarunas ar Rumāniju par Polijas valdības internēšanu. L.~Noels jau bija iecerējis savu kandidātu Polijas premjerministra amatam~--- ģenerāli V.~Sikorski. 9.~septembrī tika parakstīta Francijas un Polijas vienošanās par poļu armijas veidošanu Francijā.

Ja Lielbritānija un Francija pildītu savas saistības un uzsāktu nopietnu kara darbību pret Vāciju, tam būtu nozīmīgas sekas visā tālākajā vēsturē. Nav iespējams aprēķināt jau 1939.~gadā teorētiski iespējamā Vācijas un kopēja Polijas un Rietumvalstu īsta, nevis ``sēdoša'' kara ilgumu un sekas, taču viens ir skaidrs~--- tā sekas nevarēja būt vēl ļaunākas, kā tās bija, Vācijai vispirms 1939.~gadā sagraujot Poliju un pēc tam 1940.~gadā arī Franciju. Pat ja abas minētās valstis karā divās frontēs pret Vāciju būtu sakautas jau 1939.~gadā, ko gan nekādi nevar apgalvot, arī tad Lielbritānija paliktu nacistiskajai Vācijai neaizsniedzama.

Skaidrs ir arī tas, ka Rietumvalstu godīgas ``spēles'' pret Poliju gadījumā, Ribentropa-Molotova pakta slepenais protokols izrādītos par bezvērtīgu papīra lapeli, jo PSRS pie šāda notikumu pavērsiena nesūtītu Sarkano armiju Polijas teritorijā, kas tai nozīmētu karu pret Rietumvalstīm, bet gan gaidītu Vācijas un tās pretinieku kara iznākumu. Pat ja piepildītos J.~Staļina sapņi par to, ka abas kapitālistisko valstu grupas viena otru novājinās un Sarkanajai armijai pavērsies ceļš uz Rietumiem, PSRS karš ar Somiju (1939--1940) parādīja, ka Sarkano armiju arī novājinātajos Rietumos nekādi panākumu lauri šai laikā negaidītu.

Ar ``savādā'' jeb ``sēdošā kara'' īstenošanu Lielbritānija un Francija laika (kā vēlāk izrādījās~--- astoņu mēnešu) ieguvuma vārdā ne tikai radīja sev milzīgas grūtības nākotnē, bet veicināja Vācijas kundzības izplatību pār gandrīz visu Eiropu. Ja tiek runāts par to, ka Molotova-Ribentropa pakts atļāva nacistiem izraisīt pasaules karu un par to ir jānes atbildība ne tikai tiem, bet arī boļševikiem, tad krietni lielāka atbildība kā boļševiku vadībai ir jānes Rietumvalstu demokrātiskajām valdībām, kuras savu tuvāko interešu vārdā 1939.~gada septembrī nepildīja sabiedrotā pienākumus pret Poliju, no kā pašas vēlāk arī cieta. Vilcināšanās īpaši smagi izmaksāja Francijai, pēc dažiem mēnešiem tā kļuva par Vācijas kārtējo upuri.

Historiogrāfijā ir izskanējuši arī argumenti, kuri attaisno Rietumvalstu bezdarbību 1939.~gada septembrī. Kādreiz nākas lasīt, ka Rietumu sabiedrotie nesniedza militāru palīdzību Polijai, jo saprata, ka efektīvas palīdzības nosacījums ir sadarbība ar Padomju Savienību, ko pati Polija bija noraidījusi. Taču šeit nu jāpiebilst, ka arī paši Rietumu sabiedrotie nebija pietiekami enerģiski mēģinājuši Polijas vadītājus pievērst sadarbībai ar Padomju Savienību. Tātad, iztēlot, ka Polija pati vien cieta par savām iepriekšējām kļūdām, nebūtu pareizi. Tā cieta arī par savu Rietumu sabiedroto kļūdām. Polija, noraidot padomju palīdzību, bija darījusi gandrīz visu iespējamo, lai saņemtu Rietumu palīdzību. Tā arī saņēma konkrētu solījumu, ka Francijas uzbrukums Vācijai sāksies piecpadsmitajā dienā pēc mobilizācijas sākuma. Taču tas palika tikai solījums.

Ir apgalvots, ka Lielbritānijas un Francijas palīdzība vienkārši nokavēja, Polija sabruka ātrāk nekā angļi un franči bija sagatavojušies aktīvai kara darbībai. Arī šis apgalvojums neatbilst patiesībai. Jau 1.~septembrī Francijas rīcībā bija ap divi miljoni bruņotu vīru (550~000 liela kadru armija un 1,5 miljoni rezervistu, iesaukti uz mācībām). Šī armija varēja uzsākt kara darbību, kas ievērojami atvieglotu Polijas armijas stāvokli un varētu ievadīt Otrā pasaules kara notikumus pavisam citā gultnē.

Krievu publicists J ~Muhins attaisno Francijas valdību par atteikšanos doties palīgā Polijai, jo tā 6.~septembrī no sava sūtņa Varšavā uzzinājusi, ka poļu vadītāji jau gatavojas bēgšanai. Šajā situācijā 15.~septembrī uzsākt pilnvērtīgu uzbrukumu Vācijai, kā tas agrāk bija plānots, pēc J.~Muhina domām, būtu bijis neprāts. Šādai Rietumvalstu attaisnošanai ar Polijas valdības gatavošanos bēgt gan nevar piekrist. Valdības nostāja pēc sabiedroto uzbrukuma varēja mainīties, bet armija un tauta 6.~septembrī vēl negrasījās padoties.

Bez tam tiek arī izvirzīts arguments, ka Rietumu sabiedrotie, domājot pirmkārt par izredzēm uzvarēt karā, gribēja tam labāk sagatavoties, ievērojami nostiprināt spēkus, lai 1941. vai 1942.~gadā sāktu izšķirošu uzbrukumu. Taču arī šāds arguments nepārliecina. Divu gadu laikā sasniedzamais sabiedroto spēku pieaugums nevarēja kompensēt 40~poļu divīziju zaudējumu, tās priekšrocības, kuras ieguva Vācija, okupējot Polijas teritoriju, izmantojot tās resursus, kā arī likvidējot kara draudus divās frontēs un iegūstot rīcības brīvību. Acīmredzot nebūs pārspīlēts apgalvojums, ka, ja Francija vēlējās nopietni karot ar Vāciju, visizdevīgākais brīdis tam bija 1939.~gada septembrī.

Rietumu sabiedrība kopumā bija noskaņota citādi. Tā sajuta nacisma briesmas, nosodīja sabiedroto bezdarbību Vācijas~--- Polija kara laikā. Arī vairākas komunistiskās partijas nosodīja Ā.~Hitleru kā agresoru. 1939.~gada septembra karš pārtrauca Ā.~Hitlera ``mierīgos'' iekarojumus, ar to sākās Eiropas tautu pretestība nacistiskajai agresijai. Pirmā šādu pretestību izrādīja poļu tauta. Diemžēl iespējas, ko Polijas pretošanās pavēra tās sabiedrotajiem, tā arī palika gandrīz neizmantotas.

Poļu vairākums 1939.~gadā bija sašutis par Rietumu sabiedroto nodevību. Profesors V.~Studņickis ir liecinājis: ``Kad es 1939.~gada 3.~novembrī ierados Varšavā, daudzkārt sastapos ar antipātiju, pat naidu pret Angliju \citespace{} Anglija mūs ievilka karā un nenodrošināja nekādu atbalstu \citespace{} Varētu Londona tā izskatīties kā mūsu [izpostītā] Varšava!''

Polijai karā pret Vāciju bija svarīga otra lielā kaimiņa~--- \strong{PSRS nostāja}. Padomju Savienība no Polijas~--- Vācijas kara pirmās dienas formāli ievēroja neitralitāti. Padomju sūtnis Varšavā N.~Šaronovs tāpat kā iepriekš apgalvoja J.~Bekam, ka PSRS ir gatava turpināt sarunas ar Lielbritāniju un Franciju un parakstīt militāru konvenciju. Atsaucoties uz K.~Vorošilova interviju avīzei ``\rutxti{Известия}'', viņš apjautājās, kāpēc Polija negriežas pie PSRS pēc palīdzības. Taču, kad Polijas sūtnis PSRS V.~Gžibovskis 5.~septembrī, kad vācu karaspēks jau tuvojās Varšavai, vizītē pie V.~Molotova lūdza piegādāt tai kara materiālus un atļaut tai militāro kravu tranzītu caur padomju teritoriju, atbilde saturēja solījumu tikai stingri ievērot darbojošos tirdzniecības līgumu, bet noraidīja tranzīta un kara materiālu piegādes iespēju, jo pastāvošajā starptautiskajā situācijā Padomju Savienība negribot tikt ierauta karā un tai esot jāgarantē sava drošība. Atteikšanos no kara materiālu tranzīta gan varēja traktēt kā 1932.~gada 25.~jūlija Polijas un PSRS neuzbrukšanas līguma pārkāpumu, jo tas paredzēja, ka agresijas gadījumā pret vienu pusi otra līgumslēdzējpuse ne tieši, ne netieši nesniegs palīdzību uzbrūkošajai valstij. PSRS atteikšanās nodrošināt tranzītu Polijai bija netieša palīdzība Vācijai. Tātad būtībā PSRS vadības nostāja nebūt nebija džentlmeniska.

Sākoties karam, 1939.~gada 1.~septembrī Vācijas sūtniecības Maskavā padomnieks G.~Hilgers nodeva PSRS ārlietu tautas komisāram V.~Molotovam Vācijas gaisa kara spēku Ģenerālštāba priekšnieka lūgumu, lai padomju radioraidītājs Minskā brīvajā no pārraidēm laikā raidītu īpašu signālu vācu lidotāju virs Polijas orientācijai, kā arī savu pārraižu laikā pēc iespējas vairāk atkārtotu vārdu ``Minska''. Pirms vēl nebija kļuvusi skaidra notikumu virzība, starptautiskais spēku sadalījums, tas, vai Lielbritānija un Francija šoreiz iestāsies pret Vāciju, padomju puse ar J.~Staļina tiešu piekrišanu daļēji apmierināja vācu puses lūgumu. Pirmā lūguma daļa tika noraidīta, bet otrā izpildīta. Formāli par to neviens PSRS nevarēja neko pārmest, bet zināma palīdzība ar to Vācijai tika sniegta.

PSRS vadība tomēr nesteidzās pēc iespējas ātrāk ieņemt teritoriju, kura Molotova-Ribentropa paktā bija noteikta kā PSRS ``interešu sfēra''. Vēl bija jākonstatē, vai ilgi Polija spēs pretoties Vācijai, kāda būs Lielbritānijas un Francijas reakcija uz kara sākumu, vai netiks organizēta ``otra Minhene'', jo viens no padomju vadītāju mērķiem bija iespēju robežās izvairīties no tā, ka Lielbritānija un Francija varētu pieteikt karu PSRS. Reizē PSRS vadība enerģiski gatavojās tālākajiem notikumiem.

Nevar piekrist nereti izskanējušajam apgalvojumam, ka padomju karaspēka iebrukumu Polijā noteica jau Vācijas un PSRS 23.~augusta līguma slepenais protokols. Jāuzsver, ka ne Molotova-Ribentropa pakts, ne slepenais protokols pie tā neuzlika PSRS pienākumu nedz palīdzēt Vācijai sakaut Poliju, nedz ieņemt kādas Polijas teritorijas. Nebija paredzētas nedz saistības ar karaspēku ieņemt savu ``interešu sfēru'', nedz kādi termiņi tā izdarīšanai. Staļiniskā PSRS vadība bija saglabājusi brīvas rokas. Taču, pateicoties šim protokolam, PSRS izdevās ierobežot Vācijas virzību uz austrumiem ar protokolā noteiktajām ``interešu sfēru'' robežām, neļaut Vācijai vienai pašai izlemt Polijas valsts likteni, un pašai PSRS palielināt savu teritoriju. Jākonstatē, ka J.~Staļins, piedaloties kopā ar Ā.~Hitleru Austrumeiropas sadalē, rīkojās uzmanīgi, centās nedot ieganstu apvainojumiem agresijā pret Poliju, viņa salīdzinājumam ar Ā.~Hitleru.

PSRS teorētiski varēja realizēt trīs rīcības variantus.

Pirmais ir tīri teorētisks. Baidoties no Vācijas panākumiem, cenšoties glābt Poliju, vai, izmantojot situāciju, ka Lielbritānija un Francija ir pieteikušas karu Vācijai un ``vienota imperiālistu fronte pret PSRS'' vairs nedraud, PSRS varētu sākt karu pret Vāciju lai sodītu agresoru, iznīdētu nacismu. Šādas notikumu attīstības piekritēji uzsver, ka PSRS varētu izmantot brīdi, kad Vācijas armija ir cietusi zaudējumus un šādu PSRS soli apsveiktu gan Rietumvalstis, gan Polija. Pēc viņu viedokļa, rezultāts būtu Vācijas sakāve jau 1939.~gadā, Austrumvāciju okupētu PSRS, bet Rietumvāciju~--- Rietumvalstis.

Jāsaka, ka šāda varianta īstenošana reāli iespējama būtu tikai tad, ja Polijas vadība pēc pirmajām sakāvēm attaptos un aicinātu Sarkano armiju ieiet savā teritorijā. Taču nekas neliecināja par tādas iespējas pastāvēšanu. Pretboļševistiski noskaņotie poļu vadītāji pat Polijai tiekot nacistu okupētai, sava lepnā gara mudināti, nevēlējās karot pret Vāciju kopā ar viņu nīsto PSRS.

Ja tomēr Polija kaut novēloti piekristu sadarboties ar Sarkano armiju savā teritorijā, PSRS iegūstot jau faktiski sakautu sabiedroto Poliju, iedzīvotos ne tikai tādā ienaidniekā kā Vācija, bet acīmredzot klāt nāktu arī Japāna, ar kuru Halhingolas konflikts vēl nebija noregulēts. Tad būtu piepildītas Rietumvalstu cerības, pašām paliekot malā, noraudzīties kā karo nacistiskā Vācija un militāristiskā Japāna pret boļševistisko PSRS. Šādā gadījumā PSRS nevarēja būt pārliecināta pat par to, vai Rietumvalstis vismaz vārdos turpinās atbalstīt Poliju ar tās uzradušos jauno sabiedroto PSRS, vai arī centīsies ātrāk noslēgt mieru ar Vāciju. Kara sekas būtu neprognozējamas. Var jau uzskatīt, ka Vācija 1939.~gadā nebija gatava ilgstošam karam, bet, kā rādīja vēlākais PSRS-Somijas karš, tāpat tam nebija gatava arī PSRS. Visas spekulācijas par iespējamo Vācijas -PSRS kara iznākumu 1939.~gadā ir atmetamas kā nepierādāmas.

Otrais. PSRS varēja nedarīt neko, noraugoties kā Vācija sagrauj Poliju, kurai praktiski nekādu palīdzību nesniedza nedz Francija, nedz Lielbritānija, ļaujot Vācijai ieņemt visas Polijas teritoriju (ka Vācija sakaus Poliju, jau pēc pirmās kara darbības nedēļas gandrīz nevienam nekādu šaubu nebija) un nonākt krietni vien tuvāk Maskavai, nekā vērmahts realitātē bija 1941.~gada 22.~jūnijā. Tā būtu atteikšanās no 23.~augusta līguma slepenā protokola pildīšanas un arī ar to saistītajiem plāniem par savas ietekmes un teritorijas izplatīšanu uz rietumiem. Šāda nekonsekvence būtu jāvērtē kā nepiedodama kļūda no padomju vadītāju puses, savas valsts interešu neievērošana.

Trešais variants realizējās dzīvē. Paņemt savu kumosu no brūkošās Polijas J.~Staļins gatavojās jau savlaicīgi, taču galīgo lēmumu pieņēma tikai tad, kad kļuva skaidrs, ka otras ``Minhenes vienošanās'' nebūs, bet poļi nespēs ilgstoši pretoties Vācijai.

Jau 1939.~gada 1.~septembrī VK(b)P CK Politbirojs nolēma par 76~ divīzijām palielināt SZSA strēlnieku divīziju skaitu, lai tas sasniegtu 173. 5.~septembrī padomju prese paziņoja, ka kārtējā demobilizācija (pavisam 210~632 cilvēkiem) no Sarkanās armijas tiek atlikta uz mēnesi. Jau septembra sākumā sākās papildus mobilizācija Sarkanajā armijā. 10.~septembrī tika publicēts paziņojums, ka ``noticis daļējs rezervistu iesaukums, jo vācu-poļu karš iegūst draudošu raksturu un prasa spert soļus valsts aizsardzībai''. Līdz decembrim tika iesaukts vairāk nekā 1~miljons cilvēku. Uz gadu tika palielināts dienesta laiks. Rezultātā 20.~septembrī Sarkanajā armijā dienēja ap 5,3~miljoniem cilvēku, no kuriem 659~tūkstoši bija jauniesauktie.

Lielbritānijas un Francijas kara pieteikums Vācijai 3.~septembrī lika Vācijas hitleriskajai vadībai mēģināt ``pārapdrošināt'' savu vienošanos ar J.~Staļinu par Polijas sadali. Vācijai bija izdevīgi iesaistīt notikumos PSRS.

Pēc krievu vēsturnieka Ļ. Bezimenska un ukraiņu vēsturnieces O.~Bucko domām Ā.~Hitlers rēķinājies, ka PSRS iebrukums Rietumukrainas un Rietumbaltkrievijas teritorijā izsauks Rietumu sabiedrībā sašutuma vilni un Lielbritānija un Francija pieteiks PSRS karu. Tiesa, pierādījumus šim savam uzskatam abi autori nesniedz, atstājot to pieņēmuma līmenī.

Jau 1939.~gada 3.~septembra vakarā Vācijas ārlietu ministrs J.~f.~Ribentrops telegrafēja sūtnim Maskavā V.~f.~Šulenburgam norādījumu izteikt PSRS vadībai priekšlikumu ``piemērotā brīdī'' uzstāties pret poļu spēkiem un okupēt padomju pusei paredzēto ``interešu sfēru''. PSRS vadītāji ļoti uzmanīgi reaģēja uz šo priekšlikumu. Viņi piekrita, ka ``piemērotā brīdī'' Padomju savienībai vajadzēs konkrēti rīkoties, taču uzskatīja, ka šis brīdis vēl nav pienācis, jo pārlieka steiga var kaitēt un ``veicināt mūsu ienaidnieku apvienošanos''. Kā savā ziņojumā rakstīja Vācijas sūtnis Maskavā V.~f.~Šulenburgs, viņi centās pasaules sabiedriskās domas priekšā izvairīties no apvainojuma par kopā ar Ā.~Hitleru veiktu agresiju pret Poliju.

Taču Vācijas armijas panākumi situāciju mainīja. Kad naktī no 8. uz 9.septembri J.~f.~Ribentrops aicināja V.~f.~Šulenburgu pasteidzināt padomju vadību, V.~Molotovs jau atbildēja, ka padomju bruņotie spēki gatavojas darbības uzsākšanai.

Par J.~Staļina plāniem septembra sākumā liecina viņa tikšanās 1939.~gada 7.~un 8.~septembrī ar Kominternes Izpildu Komitejas ģenerālsekretāru G.~Dimitrovu, kurā piedalījās arī padomju līderi V.~Molotovs un A.~Ždanovs. Tā ka Kominternes vadība sakarā ar sākušos Pasaules karu bija apjukumā, J.~Staļins deva Kominternei jaunu uzstādījumu: karš noris starp divām kapitālistisko valstu grupām par pasaules pārdali un tāpēc iepriekšējais dalījums demokrātiskās un fašistiskās valstīs ir zaudējis jēgu, bet Tautas frontes lozungs jānoņem. Poliju viņš neuzlūkoja kā agresijas upuri. ``Šīs valsts iznīcināšana tagadējos apstākļos nozīmētu, ka par vienu buržuāzisku fašistisku valsti kļūtu mazāk. Kas tur būtu slikts, ja Polijas sagrāves rezultātā mēs sociālistisko sistēmu izplatītu uz jaunām teritorijām un iedzīvotājiem.'' Raksturojis sākušos karu kā sadursmi ``starp kapitālistisko valstu divām grupām'', J.~Staļins turpināja: ``Mums nav nekas pret to, ka tās pamatīgi izkaujas un novājina viena otru. Nebūtu slikti, ja ar Vācijas rokām tiktu sašķobīts bagātāko kapitālistisko valstu (sevišķi Anglijas) stāvoklis. Hitlers, pats to nesaprotot un nevēloties, grauj kapitālistisko sistēmu.'' Uz šo spriedumu pamata tika sagatavota un 9.~septembrī apstiprināta Kominternes sekretariāta direktīva, kur par vācu-poļu karu viennozīmīgi tika paziņots: ``Patreizējais karš ir imperiālistisks, kura izraisīšanā vienādi vainojama visu karojošo valstu buržuāzija. Karu nevienā valstī nevar atbalstīt ne strādnieku šķira, ne vēl jo vairāk komunistiskā partija. \citespace{} ``Starptautiskais proletariāts nekādā gadījumā nevar aizstāvēt fašistisko Poliju, kura noraidīja Padomju Savienības palīdzību [un] apspiež citas tautības''.

9.~septembrī PSRS vadība saņēma vācu izplatīto ziņu par Varšavas ieņemšanu un apsveica Vāciju ar to. Sakarā ar šo ziņu PSRS Aizsardzības tautas komisārs K.~Vorošilovs un Sarkanās armijas Ģenerālštāba priekšnieks B.~Šapošņikovs parakstīja direktīvu Baltkrievijas un Kijevas kara apgabalu armijas grupām naktī no 12. uz 13.~septembri pāriet padomju~--- poļu robežu un uzdevumu ``ar zibenīgu triecienu sagraut pretī stāvošos pretinieka spēkus''. 10.~septembrī V.~Molotovs paziņoja Vācijas sūtnim V.~f.~Šulenburgam, ka jau mobilizēti 3~miljoni cilvēku, bet tomēr aizbildinājās, ka vēl ir vajadzīgs laiks armijas sagatavošanai. Taču ziņa par Varšavas krišanu izrādījās nepatiesa un direktīva tā arī palika Sarkanās armijas Ģenerālštābā. Varšavas ieņemšanas aizkavēšanās (tā kapitulēja tikai 28.~septembrī) veda pie Polijas robežas pārejas atlikšanas.

Sakarā ar to, ka Sarkanās armijas uzbrukums Polijai kavējās, J.~f.~Ribentrops 1939.~gada 12.~septembrī pat apsprieda iespējamību izprovocēt ukraiņu sacelšanos Austrumgalīcijā.

Starp citu, Vācijas militārā izlūkdienesta vadītājs admirālis V.~Kanariss jau augustā tikās ar Ukraiņu Nacionālistu organizācijas vadītāju A.~Meļņiku un paziņoja par ``neatkarīgas Ukrainas'' izveides iespēju. A.~Meļņiks pat sāka formēt nākamo koalīcijas valdību. Vācu izlūkdienesta skubināti jau 10.~septembrī Ukraiņu nacionālistu organizācijas (\uktxti{Організація Українських Націоналістів}) kaujinieki uzbruka poļu apsardzes vienībai pie Razvadovas (ukr. \uktxti{Розвадів}) dzelzceļa stacijas. Šīs sacelšanās apspiešanā bija iesaistīti ap 600~poļu policisti no Židačovas (ukr. \uktxti{Жидачів}) pilsētas. Kauju laikā tika nodedzināti divi ciemati. 11.~septembrī Ukraiņu nacionālistu organizācijas kaujinieki atbruņoja ap 500~poļu karavīru pie Stavčanu (ukr. \uktxti{Ставчани}) un Obrošinas (ukr. \uktxti{Оброшинo}) ciemiem. Var nosaukt vēl dažas ukraiņu nacionālistu uzstāšanās gadījumus.

14.~septembrī V.~Molotovs Vācijas sūtnim V.~f.~Šulenburgam paziņoja, ka Sarkanā armija sagatavojusies aktīvai darbībai ātrāk nekā tas bijis paredzēts, taču, kamēr Varšava aizstāvas, PSRS nevar pamatot sava karaspēka ievešanu Polijā ``ar poļu valsts sairšanu''. V.~Molotovs lūdza paziņot, kad ir sagaidāma Polijas administratīvā centra~--- Varšavas krišana. Acīmredzot padomju vadību ietekmēja arī ziņas, ka Polijas valdība ir atstājusi Varšavu. (Tiesa, tikai 17.~septembrī padomju telegrāfa aģentūra TASS saņēma drošas ziņas, ka prezidents I.~Moscickis un Polijas valdības locekļi atrodas pie Rumānijas robežas un ir lūguši atļauju ierasties Buharestē.) Padomju prese saņēma signālu sabiedriskās domas sagatavošanai jaunai notikumu attīstībai.

Padomju avīze ``\rutxti{Известия}'' (``Ziņas'') uzsvēra, kas poļu lidmašīnas pārkāpj padomju gaisa robežu. Vēl skaidrāku valodu runāja avīze ``\rutxti{Правда}'' (``Taisnība''). 14.~septembrī tā publicēja VK(B)P CK Politbiroja locekļa A.~Ždanova rakstu, kuru bija labojis pats J.~Staļins. Tas izskaidroja ``Polijas sakāves karā iekšējos cēloņus'', uzsverot tās daudznacionālo sastāvu un ekspluatatorisko iekārtu: ``\dots{}Kur cēlonis stāvoklim, kurš Poliju novedis bankrota priekšā? \citespace{} Polijas valdošo aprindu nacionālā politika raksturojas ar nacionālo mazākumu, īpaši ukraiņu un baltkrievu, apspiešanu un spaidiem \citespace{} Polijas valdošās aprindas, kuras dižojas ar savu it kā brīvības mīlestību, ir izdarījušas visu, lai pārvērstu Rietumu Ukrainu un Baltkrieviju beztiesiskā kolonijā, kura atdota poļu paniem izlaupīšanai \citespace{} Polijas nacionālie mazākumi nekļuva un nevarēja kļūt par valstiskā režīma uzticamu balstu \citespace{} ``Poļu valsts, dibināta uz tās teritorijā dzīvojošo baltkrievu un ukraiņu ekspluatāciju, izrādījās rīcībnespējīga un pie pirmajām militārajām neveiksmēm sāka sairt.'' Lasītājam tika iedvesta doma, ka Polijā pastāv pretrunīga situācija: no vienas puses~--- pie varas vēl turas reakcionārā Polijas valdība, no otras~--- jūgu nesošie ``asinsbrāļi'' baltkrievi un ukraiņi ar nepacietību gaida atbrīvošanu no PSRS ar Sarkanās armijas spēkiem. Šodien mēs zinām, ka PSRS vadība šajā laikā jau gatavoja ideoloģisku aizsegu iebrukumam Polijā, ka pašā PSRS 30.~gados nebūt nebija nacionālas vienlīdzības, taču tāpēc, ka patiesību par kādu ļaundari izsaka vēl lielāks ļaundaris, tā nepārstāj būt patiesība.

14.~septembrī K.~Vorošilovs un B.~Šapošņikovs nosūtīja direktīvas Baltkrievijas un Ukrainas frontēm par \strong{padomju karaspēka uzbrukuma Polijai sākumu} 17.~septembrī pulksten 5.00. PSRS interesēs vairs nebija vilcināties. Sarkanajai armijai tika dota pavēle ``iznīcināt un sagūstīt Polijas bruņotos spēkus'', kuri darbojās uz austrumiem no pirms tam saskaņotās līnijas, kas dalīja Poliju Vācijas un PSRS ``interešu sfērās''.

15.~septembrī J.~f.~Ribentrops deva V.~f.~Šulenburgam instrukciju nodot PSRS vadībai informāciju, ka Polijas armijas iznīcināšana tuvojas noslēgumam, Varšava tiks ieņemta tuvākajās dienās. Tika apsveikta gatavotā PSRS militārā iejaukšanās, taču padomju vadība arī brīdināta, ja šī iejaukšanās nenotiks, aiz vācu ietekmes sfēras var izveidoties politisks vakuums, kur var rasties apstākļi jaunu valstu izveidei.

Kā norāda ukraiņu jurists V.~Makarčuks, 1939.~gadā starptautiskajās tiesībās pastāvēja norma, kuras nav mūsdienās~--- valstu ``tiesības uz pašpalīdzību''. Šo tiesību ietvaros ``valsts, kura uzskatīja, ka cita starptautisko tiesību subjekta rīcība rada draudus tās dzīves svarīgām interesēm, (kuras varēja tulkot visai plaši) varēja atbilstoši pastāvošajām starptautiskajām tiesībām lietot spēku, lai novērstu šos draudus.'' Polijas armijas sakāve radīja reālus draudus PSRS drošībai. Vācijas karaspēks jau iegāja Rietumbaltkrievijas un Rietumukrainas teritorijā, kuras atbilstoši slepenajam protokolam pie 23.~augusta līguma ietilpa PSRS ``interešu sfērā''. Vācu daļām virzoties uz priekšu ar ātrumu 25--30~km diennaktī un nesastopot sev priekšā vienotu poļu fronti, tās 4 līdz 8~diennakšu laikā (21.--25.~septembrī) varēja ieņemt visu Rietumbaltkrieviju un Rietumukrainu. Radās jautājums~--- vai Ā.~Hitlers ievēros vienošanos? Viņš jau bija pārkāpis vairākus līgumus ar citām valstīm. Vai tāpat nenotiks arī ar šo? (Protams, par to, cik lielā mērā PSRS vadība domāja par ``pašpalīdzību'', cik par savu lielvalstisko plānu īstenošanu, var strīdēties bezgalīgi). Lai nu kā, 15.~septembra Vācijas brīdinājums par ``politiskā vakuuma'' ar apstākļiem jaunu valstu izveidei rašanos Polijas austrumu teritorijās lika Padomju vadībai rīkoties.

16.~septembrī pēc sarunas ar J.~Staļinu V.~Molotovs informēja V.~f.~Šulenburgu, ka nākamajā dienā Sarkanā armija šķērsos Polijas robežu visā tās garumā. (Starp citu, tieši 1939.gada 16.~septembrī tika izbeigta kara darbība starp PSRS un Mongolijas karaspēku, no vienas puses, un Japānas armiju, no otras puses, pie Halhingolas upes. Sarkanajai armijai vairs nevajadzēja baidīties no karadarbības divās frontēs.)

Īstenojot 14.~septembra K.~Vorošilova un B.~Šapošņikova direktīvas par padomju karaspēka uzbrukuma sākumu, Baltkrievijas frontes Kara padome 15.~septembrī izdeva pavēli, kurā bija teikts, ka ``baltkrievu, ukraiņu un poļu tauta lej asinis karā, kuru ar Vāciju uzsākusi Polijas muižnieciski-kapitālistiskā kliķe. Baltkrievijas, Ukrainas un Polijas strādnieki un zemnieki ir cēlušies cīņai pret saviem mūžsenajiem ienaidniekiem muižniekiem un kapitālistiem. Vācu armija ir smagi sakāvusi Polijas armijas galvenos spēkus. Baltkrievijas frontes armijas ar 1939.~gada rītausmu pāriet uzbrukumā ar uzdevumu~--- veicināt sacēlušos Baltkrievijas un Polijas strādnieku un zemnieku cīņu, lai gāztu muižnieku un kapitālistu jūgu un nepieļautu, lai Rietumbaltkrievijas teritoriju ieņemtu Vācija.'' Pavēlē bija arī noteikts: ``\dots{}Frontes tuvākais uzdevums ir iznīcināt un sagūstīt Polijas bruņotos spēkus, kuri darbojas uz austrumiem no Lietuvas robežas un līnijas Grodņa~--- Kobrina.'' Līdzīgu pavēli saņēma Ukrainas frontes karaspēks.

Kopumā abu: Baltkrievijas un Ukrainas fronšu karaspēks, kurš 16.~septembra vakarā atradās gatavībā Polijas pierobežā un 17.~septembrī šķērsoja robežu, pēc dažiem datiem sastāvēja no 21~strēlnieku un 13~kavalērijas divīzijām, 16~tanku, 2~motorizētajām brigādēm un Dņepras upes flotīlijas. Pēc 1993.~gadā Krievijas Federācijā publiskotajiem datiem karagājienā uz Rietumukrainu un Rietumbaltkrieviju piedalījās 466~516 padomju militārpersonu, taču pēc citiem datiem vēlāk pienāca papildinājumi un kopā iebrukumā Polijā piedalījās 617~588 cilvēki, 4~736 tanki, 4~959 lielgabali un mīnmetēji. Abi fronšu aviācijā skaitījās 3~298 lidmašīnas (faktiski to bija mazāk un kara darbībā tās piedalījās, veicot galvenokārt izlūklidojumus).

Pulksten 3.00 1939.~gada 17.~septembrī PSRS ārlietu tautas komisāra vietnieks V.~Potjemkins Polijas sūtnim V.~Gžibovskim nolasīja PSRS valdības notu, kura neilgi pirms tam bija koriģēta, ņemot vērā Vācijas vēlmes.

Koriģēts tika informatīvais aizsegs, kādu padomju vadība izmantoja Sarkanās armijas karagājiena izskaidrošanai. 10.~septembrī V.~Molotovs informēja V.~f.~Šulenburgu, ka Padomju valdībai ir nolūks izmantot vācu daļu tālāko virzīšanos, lai paziņotu, ka Polija ir sabrukusi un Padomju Savienībai tāpēc ir jāiet palīgā ukraiņiem un baltkrieviem, kurus ``apdraud'' Vācija. Šis arguments ir nepieciešams, lai Padomju Savienības intervence masu acīs šķistu pieklājīga un lai izvairītos no tā, ka Padomju Savienība izskatītos kā agresors. 15.~septembrī J.~f.~Ribentrops uzdeva V.~f.~Šulenburgam paziņot V.~Molotovam, ka par šādu Sarkanās armijas darbības izskaidrojumu nevar būt ne runas, jo tas Vāciju un PSRS visas pasaules priekšā nostādīs kā ienaidniekus un ``būs pretrunā vēlmei veidot draudzīgas attiecības, ko izteikušas abas valstis''. Nākamajā dienā V.~Molotovs informēja V.~f.~Šulenburgu, ka Padomju valdība atradusi jaunu formulējumu, lai attaisnotu Sarkanās armijas ievešanu Polijā, kurā nekas nebija teikts par Vācijas draudiem ukraiņiem un baltkrieviem. Taču informāciju Vācijai par Sarkanās armijas ievešanu Polijā padomju vadība nodeva pēdējā brīdī~--- naktī uz 17.~septembri. Vācijas vadītājiem vairs nebija laika protestēt pret padomju akcijas motivāciju, kurai arī bez Vācijas nosaukšanas vārdā bija pretvācisks raksturs, jo baltkrievu un ukraiņu tautas aizsardzība nebija iedomājama no kādas citas valsts, bet vienīgi no Vācijas.

Notā bija apgalvots, ka ``\dots{}Poļu-vācu karš ir atsedzis Polijas valsts pastāvēšanas iekšējo nespēju \citespace{} Varšava kā Polijas galvaspilsēta vairs nepastāv. Polijas valdība ir sabrukusi un nerāda dzīvības pazīmes. Tas nozīmē, ka Polija un tās valdība faktiski ir beigušas pastāvēt. Līdz ar to ir beiguši darboties līgumi, noslēgtie starp PSRS un Poliju. Palikusi pamesta un bez vadības, Polija ir kļuvusi par izdevīgu vietu dažādām nejaušībām un negaidītiem notikumiem, kas var apdraudēt PSRS.'' (Notā bija virkne faktoloģisku neatbilstību īstenībai, Varšavu vācu karaspēks vēl nebija ieņēmis, Polijas valdība vēl atradās savas valsts teritorijā, Polijas valsts \latxti{de~facto} gan jau bija sagrauta, taču \latxti{de~iure} tā pastāvēja.) Padomju valdība paziņoja, ka ``būdama līdz šim neitrāla'', tagad tā vairs nevar neitrāli izturēties pret varbūtējo apdraudējumu valstij un ukraiņu un baltkrievu iedzīvotāju, dzīvojošu Polijas teritorijā, neaizsargātību un tāpēc devusi armijai pavēli pāriet Polijas robežu un ņemt savā aizsardzībā šo tautu dzīvību un mantu. Faktiski J.~Staļins 1939.~gadā Polijas ukraiņu un baltkrievu apdzīvoto teritoriju pievienošanu PSRS pamatoja tāpat kā Krievijas imperatore Katrīna II 1792.~gadā (par ko jau rakstīts 1.~nodaļā). (Atzīmējams, ka arī vēlāk J.~Staļina realizētā Bukovinas (1940), un Aizkarpatu (1945) aneksija (atņemot šīs zemes Polijai, Rumānijai, Čehoslovākijai,) nenorisa zem ``krievu zemju'' apvienošanas, kā tas bija I pasaules karā, bet ukraiņu un baltkrievu tautu apvienošanas lozunga.) Taču ir saprotams, ka ne XVIII, ne XX gadsimtā nedz Krievijas, nedz PSRS valdnieki ne tik daudz rūpējās par krievu, ukraiņu un baltkrievu tautu interesēm, cik savas impērijas varenības kāpināšanu.

Polijas sūtnis V.~Gžibovskis, pretrunā ar visiem diplomātijas noteikumiem, atteicās notu pieņemt, jo ``tas būtu nesavienojams ar Polijas valdības cieņu''. Tiesa gan, otrs notas eksemplārs bija jau nosūtīts uz Polijas sūtniecību, nodots tās sargam. Polijas diplomāti nevarēja apgalvot, ka notu nav saņēmuši. Tai pat 17.~septembrī citu valstu sūtņi un vēstnieki, kuri atradās Maskavā, saņēma identiskas Padomju valdības notas, kurās bija runāts par notas iesniegšanu Polijas sūtnim (tās teksts bija pievienots) un paziņojumu, ka PSRS ieturēs neitralitātes politiku pret ``jūsu valsti''. Tādejādi J.~Staļins sūtīja Lielbritānijai un Francijai brīdinājumu, ka viņš negatavojas pret tām karot. Nosaukto valstu valdības viņu saprata pareizi.

Par padomju notu Polijai jāsaka, ka tāpat kā 1939.~gada martā Lielbritānijas premjers N.~Čemberlens pārstāvju palātā paziņoja, ka Čehoslovākija beigusi savu pastāvēšanu ``iekšēja sabrukuma rezultātā'' (bet nevis nacistu darbības rezultātā, kā tas bija patiesībā) un Rietumu demokrātisko valstu vadītāji sakarā ar šo ``iekšējo sabrukumu'' atsauca savas garantijas Čehoslovākijas valsts pastāvēšanai, līdzīgi PSRS vadība savā notā paziņoja, ka ``Polija un tās valdība faktiski ir beigušas pastāvēt'' jau 17.~septembrī. No 17.~septembra padomju varas iestādes neatzina Polijas sūtniecības darbinieku Maskavā diplomātisko statusu, traucēja to evakuāciju no PSRS, kad diplomātisko tiesību normas kā reiz prasīja sniegt viņiem palīdzību izbraukšanā. Padomju Savienības Tautas komisāru padomes priekšsēdētājs un ārlietu komisārs V.~Molotovs Maskavā 1939.~gada 1.~novembrī vēlreiz paziņoja: ``Starp Padomju Savienību un Vāciju noslēgtais neuzbrukšanas līgums mums uzlika par pienākumu ieturēt neitralitāti gadījumā, ja Vācija piedalītos karā. Mēs šo līniju konsekventi esam ieturējuši, kam nebūt nerunā pretim mūsu karaspēka ieiešana bijušās Polijas teritorijā, kas sākās 17.~septembrī. Kā zināms, mūsu karaspēks Polijas teritorijā iegāja tikai pēc tam, kad poļu valsts bija sairusi un faktiski izbeidza pastāvēt.'' Faktiski N.~Čemberlens un V.~Molotovs, motivējot~--- viens savas valsts atteikšanos no garantiju Čehoslovākijai pildīšanai, otrs~--- neuzbrukšanas u.c. līgumu ar Poliju pārkāpšanu, lietoja vienu un to pašu argumentu. Gan Rietumvalstu vadītāji martā, gan PSRS vadītāji septembrī redzēja tikai to, ko viņi vēlējās saskatīt. Tomēr daudzi no tiem, kuri pagātnē un arī mūsdienās izsaka neapmierinātību ar PSRS rīcību 1939.~gada septembrī, ``aizmirst'' nosodīt Lielbritāniju un Franciju, kuras ``pievēra acis'' uz faktisko Vācijas agresiju pret Čehoslovākiju pusgadu pirms tam.

Vēl viens salīdzinājums. 1939.~gada septembrī PSRS attiecībā pret Poliju būtībā rīkojās tāpat kā Polija pret Lietuvu 1920.~gada oktobrī un pret Čehoslovākiju 1938.~gada septembrī/oktobrī. Abas valdības~--- gan Polijas, gan PSRS rīkojās ar spēku, bet aizbildinājās ar savu tautiešu aizsardzību, pēc tam noformēja citu valstu teritoriju it kā brīvprātīgu pievienošanos savai valstij.

PSRS rīcību 1939.~gada septembrī padomju propaganda tēloja un arī atsevišķi mūsdienu krievu vēsturnieki traktē vienīgi kā nacionālo interešu noteiktu, Polijas 1920.~gadā Padomju Krievijai uzspiestā kara noslēgumu un poļu okupēto zemju, kuras Polija 20~savas valdīšanas gados tā arī nespēja ``sagremot'', atbrīvošanu. Faktiski šāda pieeja neatšķiras no poļu nacionālistu pieejas tai pašā 1920.~gadā. Tie taču tāpat atsaucās uz ``nacionālajām'' interesēm. Jākonstatē vienīgi, ka dažādo~--- gan poļu, gan padomju ``nacionālo'' interešu paudēji nerēķinājās ar vietējo iedzīvotāju interesēm, neļāva tiem pašiem noteikt savu nākotni.

Tie vēsturnieki, kuri pārmet PSRS, ka tā 1939.~gada septembrī nepalīdzēja, bet uzbruka Polijai, faktiski pielieto dubultstandartus. Tas, kas bija atļauts Polijai 1938.~gadā (Cešinas apgabala atņemšana Čehoslovākijai), izrādās, nav atļauts PSRS 1939.~gadā (Rietumukrainas un Rietumbaltkrievijas atņemšana Polijai) pat tad, ja tas tieši atbilda tās valstiskajām un, kā vēlāk izrādījās, arī visas antihitleriskās koalīcijas interesēm, jo ļāva vienai tās dalībniecei~--- PSRS 1941.~gadā paglābties no sagrāves. Tā, ilgstoši Polijā strādājušais Krievijas diplomāts J.~Ivanovs uzskata, ka 1939.~gada septembrī ``faktiski tika ielikti pamati vācu blickrīga neveiksmei Krievijā'', bet ``padomju karaspēka ieiešana Polijā, kā apstiprināja tālāko notikumu gaita, objektīvi atbilda cīņas interesēm pret fašismu, ieskaitot arī pašas Polijas intereses gadus uz priekšu''. Padomju rakstnieks K.~Simonovs, kurš tūlīt pēc Sarkanās armijas ieiešanas Polijas teritorijā kā žurnālists apmeklēja to, rakstīja: ``… no galvas neizgāja doma, kura toreiz nebija sveša daudziem: nu, bet, ja nebūtu [padomju valdības 1939.~gada 17.~septembra] paziņojuma, ja nebūtu vienošanās ar vāciešiem par demarkācijas līniju, [mēs] nenonāku līdz tai, \citespace{} kas tad ieietu šajās pilsētās un laukos, kas ieņemtu visu Rietumu Baltkrieviju, kas nonāktu sešdesmit kilometru attālumā no Minskas, gandrīz līdz pašai Minskai? Vācieši.'' Protams, tādā gadījumā PSRS robeža 1941.~gada jūnijā atrastos 300--400~kilometrus tālāk uz rietumiem. Tas ļautu Ā.~Hitleram daudz ātrāk realizēt padomju teritoriju sagrābšanu un, iespējams, jau vasarā fašistu armija sasniegtu Kaukāza naftas atradnes. Vācijas rīcībā tad atrastos ievērojami lielāki resursi un iespējas aktivitātēm Rietumu frontē un karadarbībai Vidusjūras reģionā. Nav izslēgts, ka padomju spēkiem būtu jāatkāpjas no Ļeņingradas un Maskavas. Pie šāda pavērsiena karš noteikti ieilgtu un būtu neiespējami prognozēt, kad un kā tas beigtos.

Ko gan var pārmest J.~Staļinam un viņa līdzgaitniekiem, ir tas, ka viņi gan sauca sevi par komunistiem, bet, sagrābjot jaunas teritorijas, rīkojās gluži tāpat kā viņu nīstie ``kapitālisti''. Pēc būtības gan vieni, gan otri rīkojās imperiālistiski.

Ukraiņu publicists S.~Lozuņķo atzīmējis, ka, ja Polija starpkaru periodā stādīja mērķi panākt PSRS sadali, neizslēdzot pat savienību ar nacistisko Vāciju, tad kādai gan bija jābūt PSRS attieksmei pret Poliju? Vai ne tādai pašai, kā Polijas attieksmei pret PSRS? Autoram gan jāatgādina, ka ne tikai Polija starpkaru periodā bija naidīga pret PSRS, arī PSRS tāda bija pret Poliju. Tam pamatā bija gan vēsturiski un politiski-ideoloģiski cēloņi, arī poļu-padomju kara sekas. Taču S.~Lozuņķo ir taisnība tai ziņā, ka 1939.~gadā gaidīt kādu citu PSRS vadības attieksmi pret Polijas nelaimi, pret Vācijas iebrukumu tajā nebija nekāda pamata. Pasaules politikā valstu līmenī kā dzinējspēks darbojas to intereses, nevis taisnīguma, līdzcietības, agrāko nodarījumu piedošanas un vēl citi cēli morāles principi. 1939.~gadā vienīgā baltkrievu un ukraiņu atkalapvienošanās alternatīva bija to šo tautu daļu, kuras dzīvoja Polijas teritorijā, pāriešana no poļu okupācijas varas Vācijas okupācijas varā un labvēlīgu pozīciju nodrošināšana vērmahtam, lai tas varētu uzbrukt PSRS. Padomju Savienībai tas bija nepieņemami. Bez sadarbības ar Rietumiem karot pret Vāciju par Polijas, ar visām tās sagrābtajām baltkrievu un ukraiņu zemēm, brīvību un teritoriālo integritāti, no PSRS puses būtu neprāts.

Strīdīgs ir jautājums par Polijas valdības rīcībspēju padomju notas saņemšanas brīdī. Piemēram, Krievijas Zinātņu Akadēmijas Slāvu studiju institūta līdzstrādnieks O.~Ņemenskis uzsver, ka 1939.~gada 17.~septembrī Polijas valsts kā tāda vairs nepastāvēja. Taču 17.~septembrī, vēl formāli atrodoties Polijas teritorijā, kaut jau ceļa jūtīs uz Rumāniju, Polijas valdība sagatavoja protestu pret padomju notu: ``Polijas valdība protestē pret padomju valdības notā izklāstītajiem motīviem, jo Polijas valdība pilda savas normālās funkcijas, bet Polijas armija dod pretsparu ienaidniekam''. Tiesa, savas zemes atstāšanas organizāciju diezin vai varēja nosaukt par ``normālu funkciju'' pildīšanu. Protestu izdevās publicēt aiz Polijas robežas nedēļu pēc valdības bēgšanas uz Rumāniju un internēšanas, kas arī par tās rīcībspēju neliecina.

Par to, ka Polijas valdība 17.~septembrī jau nespēja efektīvi pildīt savus pienākumus, liecināja arī fakts, ka 1939.~gada 18.~septembrī Polijas vēstnieki Londonā un Parīzē, izpildot J.~Beka rīkojumu un iesniedzot Lielbritānijas un Francijas valdībām notas ar paziņojumu, ka tiek gaidīts šo valstu ``aktīvs protests pret [PSRS] agresiju'', nespēja paskaidrot, vai Polija ar PSRS atrodas karastāvoklī. Vēl šodien Polijā par šo neskaidrību tiek vainota uz Rumāniju bēgusī Polijas politiskā elite.

Taču vismaz juridiski Polijas valdība pirms internēšanas nepārtrauca savu eksistenci. No šī viedokļa raugoties, PSRS valdības arguments par visu padomju-poļu līgumu automātisku darbības izbeigšanos ir apstrīdams.

Ja par II Pasaules karu uzskata karu starp Poliju, Angliju, Franciju un Vāciju, tad šajā karā PSRS 1939.~gada septembrī neiesaistījās, tā speciāli pasludināja savu neitralitāti. Pie tam šīs valstis, katra savu apsvērumu dēļ, atzina PSRS par nekarojošu pusi. Pēc mūsdienu terminoloģijas Sarkanās armijas darbību Polijas austrumu daļā nosacīti varētu pat nosaukt par miera saglabāšanas operāciju.

Taču, ja II Pasaules karu uzlūko kā starptautisko attiecību sistēmas nomaiņas procesu, kurš ietver lielvalstu un citu valstu savstarpējo karu kopumu par savas ietekmes paplašināšanu, to robežu izmaiņām, kuras izveidojās 1919.--1922.~gadā, tad PSRS 1939.~gada 17.~septembrī, protams, iesaistījās II Pasaules karā, taču nevis Vācijas pusē, kā to mēģina tēlot daļa vēsturnieku, bet gan kā trešais, patstāvīgs spēks, kurš darbojās savās interesēs. Taču šo iesaistīšanos nevar uzskatīt par dalību karā, pat neraugoties uz to, ka Sarkanā armija lietoja ieročus, cieta arī zaudējumus. Pat Polijas valdība, tās militārā pavēlniecība to neatzina par kara darbību. Acīmredzot PSRS, Sarkanās armijas agresīvās rīcības nosaukšanai ir jāmeklē kāds cits termins. Pēc autora ieskatiem šai rīcībai atbilstošāks būtu nosaukums \strong{``intervence''} (varmācīga iejaukšanās citas valsts iekšējās lietās).

% page 821


Nevar piekrist apgalvojumam, ka Sarkanā armija palīdzēja vērmahtam sakaut Polijas armiju. Šīs tēzes pamatā ir gan toreizējie padomju propagandas apgalvojumi, gan vēlākās poļu ģenerāļu subjektīvās liecības.

Tā, 1939.~gada 31.~oktobrī PSRS Augstākās Padomes sēdē V.~Molotovs paziņoja, ka Polijas valdošie slāņi lepojās ar savas valsts un armijas stiprumu, taču izrādījās pietiekoši ``īsa trieciena no Vācijas armijas, bet pēc tam~--- Sarkanās armijas puses, lai nekas nepaliktu pāri no šī kroplīga Versaļas līguma lolojuma, kurš dzīvoja uz apspiesto nepoļu tautību rēķina.'' Arī K.~Vorošilova pavēlē PSRS valsts svētkos 1939.~gada 7.~novembrī bija norādīts, ka ``ar strauju uzbrukumu Sarkanās armijas daļas sakāva poļu karaspēku, īsā laikā izpildot savu pienākumu Padomju dzimtenes priekšā.'' V.~Molotovs un K.~Vorošilovs pārspīlēja Sarkanās armijas lomu Polijas sagrāvē, bet V.~Molotovs arī nekorekti lika vienlīdzības zīmi starp starpkaru Polijas valdības neveiksmīgo politiku un Polijas valstiskumu kā tādu.

Poļu ģenerāļi ir rakstījuši, ka 17.~septembrī Polijas armija pret vērmahtu vēl varējusi cīnīties pašā Polijā, ka tā vēl nebūt nebijusi sakauta utt. Tā Polijas armijas Galvenā štāba priekšnieks ģenerālis V.~Stahevičs izteicās, ka bez Sarkanās armijas ienākšanas Polijā tās karaspēks vēl ilgi varētu pretoties vāciešiem. Arī ģenerālis V.~Anderss apgalvoja, ka septembra vidū sākušas trūkt vācu komunikācijas, Polijas armija vēl varējusi zināmu laiku pretoties, lai ļautu Lielbritānijai un Francijai dot savu triecienu Vācijai. Šādi apgalvojumi balstās uz pārāk optimistisku stāvokļa vērtējumu. Arī poļu vēsturnieki ir tendēti pārspīlēt Sarkanās armijas ienākšanas Polijā militāro nozīmi. f(Tā, poļu vēsturnieks P.~Večorkovičs ir izteicies, ka ``padomju uzbrukuma dienai kara gaitā bija izšķiroša nozīme.)

Var jau piekrist apgalvojumiem, ka bez Sarkanās armijas iebrukuma poļu pretošanās būtu ilgāka, taču šī pretestība būtu bez izredzēm uz uzvaru, tai nebūtu stratēģiskas nozīmes. Jā, ja Rietumvalstis Polijai ietu palīgā, kas piespiestu Vāciju daļu spēku pārsviest uz rietumiem, tad pretestībai būtu nozīme, tā varētu nest zināmus panākumus. Taču, ja Francija un Lielbritānija to nedarīja tad, kad viņām vajadzēja karot tikai ar Vāciju, tad pēc 17.~septembra, kad notika PSRS iebrukums Polijā, kad vismaz teorētiski varēja gaidīt PSRS nostāšanos Vācijas pusē arī pret Rietumvalstīm, tās tāpat nemēģināja glābt nelaimīgo Poliju.

Patiesībā vērmahts jau bija sakāvis Polijas armiju, kad 17.~septembrī Sarkanā armija pārsvarā bez lielākām kaujām sāka virzīties Polijas teritorijā. Kā liecināja vācu ģenerālis Ē.~Manšteins, Sarkanās armijas iesoļošana Polijā notika periodā, kad jau norisa kara darbību noslēdzošas kaujas un kara iznākums bija skaidrs. Par kādu lielu Sarkanās armijas nozīmi Polijas armijas sakāvē nevar runāt arī, ievērojot tos salīdzinoši mazos poļu spēkus, kuri bija izvietoti gar padomju robežu. Sarkanās armijas milzu pārspēka rezultātā tie tika ātri pievarēti, bet galvenie Polijas armijas spēki, kuri cīnījās pret Vāciju, ir nemēģināja organizēt pretestību arī Sarkanajai armijai. Tāpat vērmahta un Sarkanās armijas ciesto zaudējumu salīdzinājums ļauj spriest, kas tad iznesa galveno kara darbības pret Poliju smagumu. Arī E.~Ridzs-Smiglija 20.~septembrī izteiktais skaidrojums par to, ka 17.~septembrī viņš rūpējies par poļu karavīru pārsviešanu uz Franciju lai turpinātu cīņu, nonāk pretrunā ar iepriekš citēto poļu ģenerāļu izteikumiem. Tāpēc pārmetums PSRS par Polijas pretošanās Vācijai traucēšanu ir jāuzskata par mazpārliecinošu. Ja arī Sarkanā armija 17.~septembrī nebūtu iegājusi Polijai piederošajās teritorijās austrumos, Polijas armijai, atkāpjoties turp no poļu apdzīvotajiem rajoniem, sastopoties ar nedraudzīgu lielas daļas vietējo iedzīvotāju attieksmi, pēc jau pārciestajām sakāvēm nebūtu ne mazāko cerību uz sekmīgu karadarbības turpināšanu.

Jau 1939.~gada oktobrī vācu rakstnieks un kara ziņotājs R.~Bathe savai Vācijas~--- Polijas karam veltītajai grāmatai izvēlējās virsrakstu ``\detxti{Der Feldzug der 18 Tage}'' (``18 dienu karagājiens''), ar to norādot, ka Vācija uzvaru pār Poliju faktiski izcīnīja triju nedēļu laikā. Rakstnieks uzsvēra, ka Vācijas armijai šai laikā izdevās virknē aplenkuma operāciju poļu spēkus ``sašķelt, sagūstīt un iznīcināt''. Protams, zināmā mērā pārējās cīņas un ``tīrīšanas akcijas'', kas norisa vēl oktobrī ar to tika atstātas novārtā. Kā uzskatīja daži poļu un padomju vēsturnieki, iemesls, kāpēc cīņu laiks tika saīsināts, ir viegli saprotams~--- ar to tika mazināta Sarkanās armijas loma Polijas sagrāvē. Vācijas propagandai bijis izdevīgi tēlot, ka Polija jau bija sabrukusi, kad tajā iegāja Sarkanās armijas daļas. Tikai otrā pretinieka uzbrukums no aizmugures, apstākļos, kad nebija nekādas Rietumu sabiedroto palīdzības, padarījis Polijas stāvokli bezcerīgu. To sapratuši kā poļu politiķi, tā karavīri un civiliedzīvotāji. Esošie fakti gan liecina, ka politiķi un vismaz ģeneralitāte Polijas stāvokļa bezcerību saprata jau septembra pirmajās dienās. Citādi nebūtu notikusi to sasteigtā bēgšana no Polijas.

Septembra otrajā pusē, kad izplatījās ziņas par Sarkanās armijas ieiešanu Polijas teritorijā un Polijas valdības bēgšanu, daudziem poļu karavīriem cīņas griba sabruka kā kāršu namiņš, viņi vairs neturēja prātā pienākuma izpildi. Desmiti tūkstošu domāja tikai par to, kā nomainīt formu pret civilapģērbu, lai izbēgtu draudošajam gūstam. Taču, ja tas arī izdevās, nodevīgs izrādījās poļu armijā ieviestais īsais matu griezums. Jaunus vīriešus ar tādām frizūrām, kuri nevarēja uzrādīt civilpersonas dokumentus, vācu karaspēks nekavējoties arestēja. Tā kā vācu okupētajos Polijas apgabalos par ieroču glabāšanu bija paredzēts nāves sods, daudzas karaspēka daļas vispirms nometa ieročus un tikai tad izklīda. Ja civiliedzīvotāju mājās vai to apkārtnē atrada militāro formu vai ieročus, viņi nonāca nāves briesmās, jo tos varēja turēt aizdomās par poļu karavīru atbalstīšanu vai pat par dalību cīņās.

PSRS vadība Polijas sagrāvē nebija vainīga, tā tikai baudīja tās augļus. Par sagrāvi atbildība galvenokārt ir jānes Polijas politiskajai un militārajai vadībai, kura noveda valsti pie tik bēdīga rezultāta, bet vēlāk centās sevi attaisnot, apvainojot ``viltīgos boļševikus'' un ``tūļīgos francūžus''.

Līdz pat 1939.~gada 1.~septembrim poļu izlūkdienests nebija ziņojis par kādu īpašu padomju bruņoto spēku pārvietošanos Polijas robežas virzienā. Tikai 1.--5.~septembrī sāka pienākt šādi ziņojumi, taču poļu pavēlniecība šādu virzību vērtēja kā parastu padomju reakciju uz kara sākumu kaimiņvalstīs. 12.~septembrī Polijas kara atašejs Parīzē atsūtīja uz Varšavu izvilkumu no Francijas izlūkdienesta pārskata, kurā bija aģentūras ziņas, ka Padomju Savienība uzsāks pret Poliju un Rumāniju vērstu darbību. Neraugoties uz to, padomju akcija 17.~septembrī bija pilnīgs pārsteigums Polijas militārajām un valdības aprindām, kuras tai brīdī atradās Rumānijas pierobežā un gatavojās šeit aizstāvēt t.s. ``rumāņu placdarmu''.

Polijas bruņotie spēki 17.~septembrī uzbrukumu no austrumiem negaidīja. Tiem ap 1~400~km garajā robežā ar PSRS ik uz km bija ap 10~robežsargu. Vēl valsts austrumu daļā atradās ap 350~tūkstošu daļēji apbruņotu, pat vēl neietērptu mundieros jauniesaucamo, kuri izgāja apmācību.

Pie tam padomju tanku straujo virzību uz rietumiem nodrošināja arī plašs padomju propagandas darbs: paziņojumi par slāvu tautu savstarpējo palīdzību, paglābšanu no vācu iebrukuma, pat poļu nacionālie karodziņi padomju tanku lūkās. Poļu historiogrāfijā uz plašu liecinieku aptauju pamata konstatēts, ka padomju propaganda bija pretrunīga: no vienas puses, tika apgalvots, ka Sarkanā armija ir atnākusi palīgā Polijai, no otras~--- ``PSRS un Vācija ir saistītas nesaraujamā savienībā'', no trešās~--- Padomju Savienībai poļu ``austrumu teritoriju'' atgūšana ir revanšs par 1920.~gada sakāvi, Polija ir pārstājusi eksistēt uz visiem laikiem, taču, kad tā kļūs par sociālistisku republiku, padomju karaspēks tai palīdzēs.

Toreizējais Baltkrievijas K(b)P CK 1.~sekretārs P.~Ponomarenko, kurš bija Baltkrievijas kara apgabala Kara padomes loceklis un piedalījās armijas daļu vadībā, kuras iegāja Rietumbaltkrievijā, jau pēc tās pievienošanas Baltkrievijas PSR 1939.~gada novembrī BK(b)P CK plēnumā izteicās: ``Kara darbību mēs gatavojām mazliet, protams, agrāk \citespace{} Par godu mūsu partijas biedriem un tiem, kuri zināja, jāsaka, ka konspirācija tika ieturēta pilnībā, līdz pat brīdim, kad mūsu robežsargi 17.~septembra rītausmā šķērsoja robežu. Katrs poļu robežsardzes postenis to uztvēra kā vietēju incidentu un mēģināja organizēt pretsparu. Tikai tad, kad ieraudzīja, ka pienāk [Sarkanās armijas] regulārās daļas, sāka mukt, bet, kā saka, bija jau par vēlu.''

Sarkanajai armijai pārejot Polijas robežu, gar to izvietoto Polijas armijas daļu komandieri nezināja ko darīt, prasīja instrukcijas no sava Galvenā štāba. Kaujas spējīgākās poļu daļas jau cīnījās pret vērmahtu rietumos un organizēt efektīvu pretošanos Sarkanajai armijai nebija iespējams. Mēģinājumi nosūtīt pretī Sarkanajai armijai parlamentāriešus arī neko nedeva~--- tie vienkārši tika aizturēti.

Kā rakstīja Galvenā štāba priekšnieks ģenerālis V.~Stahevičs, ``poļu daļas bija boļševiku dezorientētas, jo tie parasti izvairās atklāt uguni, bet viņu komandieri apgalvo, ka viņi ir nākuši Polijai palīgā pret vāciešiem. Padomju karavīru pamatmasa nešauj, pret mūsējiem izturas ar demonstratīvām simpātijām, dalās ar papirosiem utt., visur uzsver, ka iet palīgā Polijai''. (Pret šo liecību, kura acīmredzot tapa uz apakšnieku nepilnīgu ziņojumu pamata, ir jāizturas uzmanīgi. Aprakstītajā veidā padomju karavīri varēja izturēties pret ``asins'' un ``šķiras'' ``brāļiem''~--- ukraiņu un baltkrievu darbaļaudīm, bet ne visiem Polijas pavalstniekiem, padomju propagandas ietvaros karavīri varēja stāstīt, ka iet palīgā Polijas tautai, bet nekādi ne valstij.)

Sakarā ar padomju agresiju prezidents I.~Moscickis Kosovas pilsētiņā (ukr. \uktxti{Косів}, poļu \pltxti{Kosów}, tagad Ukrainā) griezās ar aicinājumu pie tautas, apvainojot PSRS visu juridisko un morāles normu pārkāpšanā, aicināja poļus ``saglabāt stingru garu un vīrišķību cīņā pret bezsirds barbariem''. Viņš arī pasludināja par Polijas prezidenta rezidences un citu augstāko varas iestāžu pārcelšanu uz ``vienas no mūsu sabiedroto valstu teritoriju''. Taču \strong{nedz prezidents, nedz Polijas valdība karu PSRS nepieteica}.

Mūsdienās poļu avīze ``\pltxti{Zycie Warszawy}'' (``Varšavas Dzīve''), 1993.~gada 17.~septembrī publicēja viedokli, ka Polijas valdība pieļāva lielu kļūdu, formāli nepasludinot karu PSRS, jo tas ļautu internacionalizēt konfliktu. Tiesa, ar to tiek aizmirsts, ka Polijas un Lielbritānijas savstarpējās palīdzības līgumam pievienotajā slepenajā protokolā bija skaidri norādīts, ka tas ir vērsts vienīgi pret Vāciju. Pat ja Polijas valdība arī deklarētu kara stāvokli ar PSRS, Lielbritānija un bez šaubām arī Francija, kura bez Lielbritānijas atbalsta baidījās ko uzsākt, nevēlētos jau esošajam vienam pretiniekam~--- Vācijai pievienot vēl otru~--- PSRS. Lielbritānijas un Francijas valdības ``ieteica'' Polijas valdībai nepieteikt karu PSRS. Kas notiktu, ja tā neievērotu šo ``ieteikumu'', grūti spriest, bet var iedomāties, ka Polija pat varētu zaudēt savus sabiedrotos. Tādā gadījumā būtu būtiski apdraudēta Polijas valdības sekmīga darbība emigrācijā.

Notikumi 1939.~gada septembrī bija ļoti pretrunīgi. Izvēloties dažādus faktus, tos var traktēt gan kā vieglu Sarkanās armijas ``pastaigu'', kad poļu karavīri labprāt padevās gūstā, bet sarkanarmieši galvenokārt bija nodarbināti ar to, lai demonstrētu viņiem proletārisko internacionālismu, gan kā (vismaz atsevišķas) sīvas kaujas, kurās viss notika tāpat kā jebkurā karā.

Padomju prese tolaik divas nedēļas ilgušo Sarkanās armijas ieiešanu Rietumukrainā, Rietumbaltkrievijā un Viļņas apgabalā sauca par ``atbrīvošanas karagājienu''. Vēlāk visplašāk izplatītais nosaukums bija ``septembra kampaņa'', kas gan bija pārņemts no vācu propagandas. Poļu vēsturniekiem tas izraisa nepatiku, jo viņi uzsver, ka pretošanās (kaut nenozīmīga) turpinājās arī oktobra sākumā. Polijas Tautas republikas laikā tika lietots nosaukums ``Polijas aizsardzības karš'', kurš gan neievēroja, ka karš ir politikas turpinājums, bet Polijas pirmskara politika nebūt nebija tikai uz aizsardzību vērsta.

1972.~gadā iznāca poļu disidenta Ļ.~Močuļska grāmata ``\pltxti{Wojna Polska 1939}'', (``Poļu karš 1939''), kurā viņš attaisnoja poļu politiķu kursu 1939.~gada septembrī un zem Poļu kara nosaukuma apvienoja poļu cīņu gan pret Vāciju, gan PSRS.

Mūsdienās liberālie poļu un krievu vēsturnieki runā par ``padomju-poļu karu''. Tiešām, vārda plašā nozīmē var runāt par karu~--- gan ar bruņotām sadursmēm, gan kritušajiem, gūstekņiem un materiālajiem zaudējumiem. Taču jāpiekrīt krievu vēsturniekam A.~Širokoradam, ka vārda šaurā nozīmē termins ``karš'' šeit tomēr nav iederīgs. Precīzāk ir runāt par konfliktu, karagājienu vai intervenci. Sarkanās armijas rindās Polijā krita 1~173~cilvēki, 302~pazuda bez vēsts. Salīdzinājumam, pie Halhingolas tā zaudēja 6~831 kritušos un 1~143 pazudušos bez vēsts. Iznāk, ka Sarkanās armijas zaudējumi ``lokālajā konfliktā'' pret Japānu 8~reizes pārsniedza zaudējumus ``karā'' pret Poliju.

Ja Polijas puse pieteiktu PSRS karu, aicinātu tautu karot kā pret Vāciju, tā pret PSRS, Sarkanās armijas zaudējumi Polijā varēja būt krietni vien lielāki. Ir dati, ka Polijas armijā 17.~septembrī vēl skaitījās ap 650~000 karavīru. No tiem ap 440~000 karoja pret vāciešiem. Ja pārējie 210~000, lai arī bez izredzēm uz panākumiem, (bet poļi taču pagātnē tā karoja bieži!) sāktu īstu kara darbību pret Sarkano armiju, cīņas varēja izvērsties smagas. No poļu ultranacionālistu viedokļa tā arī vajadzēja rīkoties. Taču realitāte bija tāda, ka valsts sabrukumu vēl vienas~--- pretpadomju frontes atklāšana nevarētu novērst, bet vestu tikai pie daudzskaitlīgiem upuriem.

Bez tam, aicinot tautu uz cīņu pret vēl vienu ienaidnieku~--- PSRS, kad jau vislielākajam diletantam militārajā jomā bija redzams, ka Polijai karš pret Vāciju ir bezcerīgi zaudēts, un pie tam paši aicinātāji uz vēl vienu karu, bija vainojami par septembra sakāvi un paši meklēja patvērumu ārzemēs, viņi acīmredzot tikai pārklātu sevi ar negodu, sabojātu attiecības ar pret Vāciju karojošajām Rietumvalstīm, būtiski apgrūtinātu savu atgriešanos pie varas Polijā. Bet bēgošie politiķi, militāristi uz to cerēja. Ir pamats uzskatīt, ka, gan no tautas interešu viedokļa, gan no valdošā slāņa viedokļa, Polijas valdības rīcība, nepiesakot karu PSRS, bija pareiza.

Maršals E.~Ridzs-Smiglijs 17.~septembra vakarā pa radio noraidīja Polijas armijas daļām pavēli: ``Ir iebrukušas Padomes. Pavēlu īstenot atkāpšanos uz Rumāniju un Ungāriju pa īsāko ceļu. Ar Padomēm kara darbību neuzsākt, [tā pieļaujama] tikai gadījumā, ja no to puses ir mēģinājumi atbruņot mūsu daļas. Uzdevumi Varšavai un Modlinai, kurām jāaizstāvas pret vāciešiem, paliek bez izmaiņām. Daļām, pie kuru atrašanās vietām pienākušas Padomes, jāved ar tām sarunas ar mērķi izvest garnizonus uz Rumāniju un Ungāriju \citespace{} Robežas aizsardzības korpusa daļām un daļām, kuras piesedz ``rumāņu priekštilta nocietinājumus''~--- turpināt pretošanos.'' Šādu pavēli varēja pamatot ar neskaidrību par padomju patiesajiem nodomiem. Bez skaidra politiskās varas lēmuma, kā izturēties pret PSRS, Polijas armijas virspavēlnieks arī nevarēja tai dot konkrētākas pavēles. Direktīva slikto sakaru rezultātā nesasniedza visas karaspēka daļas. Taču pat tur, kur tā tika saņemta, dažviet tā izsauca šaubas un tika pat uzlūkota kā padomju provokācija. Faktiski šāda pavēle poļu karaspēkā tikai radīja vēl lielāku nekārtību un dezorganizāciju.

20.~septembrī, jau būdams internēts Rumānijā, E.~Ridzs-Smiglijs mēģināja izskaidrot savu pavēli ar to, ka ``pirmajā dienā boļševiki uz mūsu vienībām nešāva un tās arī neatbruņoja'' un viņš vēlējies ``atvilkt samērā daudz karaspēka uz Ungārijas un Rumānijas teritoriju'', lai varētu to pārvest ``uz Franciju un tur organizēt poļu armiju'', lai ``poļu karavīrs turpinātu piedalīties karā un lai, karu nobeidzot ar uzvaru, pastāvētu Polijas armija, kura pārstāvētu Poliju un viņas intereses''. Jāpiekrīt profesoram Ē.~Jēkabsonam, ka ``E.~Ridzs-Smiglija rīcību 17.~septembrī noteica haoss, apjukums un panika, kas valdīja visā Polijā.''

Jau minētais poļu vēsturnieks P~.Večorkovičs, kaut atzīst poļu spēku nesagatavotību valsts austrumu rajonu aizstāvēšanai, tomēr galveno lomu to ātrajā zaudēšanā piešķir E.~Ridzs-Smiglija ``fatālajai pavēlei ``Pret Padomēm nekarot''''. Kaut arī Polijas austrumzemju pilnvērtīga aizsardzība neesot bijusi iespējama, tomēr vajadzējis tur izveidot ``simbolisku Vesterplati''~--- radīt aizstāvēšanās punktu, kurš ilgstoši pretotos. Šajās zemēs esot bijušas vajadzīgas poļu ``Termopilas'' (kaujā pie Termopilām 489.~gadā p.m.ēr. grieķu-persiešu kara laikā grieķi ilgstoši aizturēja uzbrūkošos persiešus), lai visai pasaulei nodemonstrētu poļu tiesības uz šīm teritorijām. Par šādu vietu varējusi kļūt, bet nekļuvusi Ļvova. Pie tam šis vēsturnieks tēlo Sarkanās armijas daļas, kura iegāja Polijā kā ``salašņas'' dažādos mundieros (viena divīzija it kā devusies uz fronte civilapģērbā), ar šautenēm, iesietām virvē. Pēc staļiniskajām armijas tīrīšanām nebijis komandieru. Ja poļi 1939.~gadā karotu tikai pret padomēm, viņi bez problēmām tiktu galā ar agresiju un 17.~septembrī atrastos 150, bet varbūt 300~km, uz austrumiem no Polijas austrumu robežām. Šī darbaautoram sakarā ar šo izteikumu gan jākonstatē, ka poļu sabiedrībā un arī vēsturniekos vēl arvien ir dzīvs šļahtas gars, kurš tik bieži vēsturē jau novedis nāciju pie utopiskām ilūzijām un tām sekojošām sakāvēm.

Atšķirībā no virspavēlnieka armijas ``Varšava'' pavēlnieks ģenerālis J.~Rummels deva pavēli Polijas robežu pārgājušās padomu daļas uzskatīt par ``sabiedrotām'', par ko arī nosūtīja paziņojumu padomju sūtnim un lūdza tā paskaidrojumus, kā pret viņa pavēli izturas PSRS armija. Nenoliedzami, ka šī pavēle bija pretrunā ar maršala E.~Ridzs-Smiglija direktīvu un radīja vēl vairāk neskaidrību.

17.~septembrī pulksten 16-os Polijas armijas pavēlniecība nonāca pie secinājuma, ka armijas stāvoklis ir bezcerīgs un pilnīga sakāve neizbēgama. Kutas (poļu \pltxti{Kuty}, ukr. \uktxti{Кути}, tagad Ukrainā) pilsētiņā notika sanāksme, kurā piedalījās prezidents I.~Moscickis, armijas virspavēlnieks E.~Ridzs-Smiglijs, premjerministrs F.~Slavojs Skladovskis un ārlietu ministrs J.~Beks. Tika pieņemts lēmums evakuēties uz Rumāniju. 17.~septembra vakarā prezidents, kā arī valdība ar premjerministru priekšgalā šķērsoja Rumānijas robežu. Pēc 17./18.septembra pusnakts tiem sekoja virspavēlnieks maršals E.~Ridzs-Smiglijs un Galvenais štābs, pametot valsti un karojošo armiju.

Francijas sūtnis L.~Noels, kurš kopā ar Polijas valdību atradās Zaļeščiku (ukr. \uktxti{Заліщики}, poļu \pltxti{Zaleszczyki}, tagad Ternopoles apgabals Ukrainā) rajonā, uz robežas ar Rumāniju, raksturoja tur pastāvošo situāciju: ``Vietējie iedzīvotāji gaida drīzu padomju karaspēku ierašanos. Klīda baumas par Berlīnes un Maskavas noslēgto vienošanos par Polijas sadalīšanu. Tika nosaukti pat dažas ģeogrāfiskās detaļas par jaunās robežas atrašanos. Daudzi ebreju iedzīvotāji gaidīja Sarkano armiju ar atklātu nepacietību. Vieni ukraiņi saistīja savas cerības ar Hitleru, citi~--- ar Staļinu. Gandrīz visi cerēja uz ātru un pilnīgu varas nomaiņu, kas paglābs viņus no Polijas žandarmērijas, muitniekiem, muižniekiem.''.

18.~septembrī prezidents, valdība un virspavēlnieks jau bija atstājuši Polijas teritoriju, kaut kaujas vēl turpinājās. Polijas tauta un tās armijas daļas, kuras vēl turpināja pretošanos, tika pamestas savam liktenim.

Toties Polijas armijas daļu pretošanās turpināšana vērmahtam arī pēc 17.~septembra jāuzlūko gan kā sava karavīra pienākuma izpilde, gan kā patriotisma izpausme, nevēloties atzīt savas valsts sagrāvi.

Polijas sūtniecība Londonā jau 17.~septembrī sakarā ar Sarkanās armijas ieiešanu Polijas teritorijā sniedza paziņojumu par šī agresijas akta nosodījumu. Tajā bija teikts, ka ``Polijas valdība funkcionē Polijas teritorijā'', taču, kā atzinis Ē.~Jēkabsons, šai paziņojumā bija ``kļūdaini norādīta Polijas valdības darbības vieta~--- Polijas teritorija''.

Jāuzsver, ka Polijas vadības pārejas laikam pāri valsts robežai nebija būtiskas nozīmes (no PSRS armijas ievešanas Polijas teritorijā viedokļa), jo pavēle par Sarkanās armijas virzību bija atdota jau 14.~septembrī, brīdī, kad Polijas vadītāji vēl neapšaubāmi atradās Polijā. No otras puses, kaut Polijas valdība oficiāli pārgāja Rumānijas robežu tikai 17.~septembra vakarā, no notikumu attīstības viedokļa daudz būtiskāk bija tas, ka valsts pārvaldes institūti jau līdz 17.~septembrim bija zaudējuši iespēju realizēt savu varu, arī armija bija kļuvusi faktiski nevadāma, tai savu gribu uzspieda vērmahts un Sarkanā armija.

Polijas politiķi un militāristi cerēja no Rumānijas doties uz Franciju un tur ar sabiedroto atbalstu turpināt cīņu. Taču no ārienes viņu rīcība izskatījās kā vulgāra savas ādas glābšana. Kā raksta poļu vēsturnieks un žurnālists D.~Bališevskis šais dienās Polijas~--- Rumānijas robežu šķērsoja, iespējams, ap 50~000 valsts politiskās, ekonomiskās, militārās elites pārstāvju. Pārsteigti un sarūgtināti par tik strauju vadības bēgšanu, vairāki virsnieki, kuri apsargāja tiltu uz robežas, nošāvās.

Par Polijas valdības locekļu, vadošo militāristu u.c. rīcību, pametot savu valsti, tautu, sagūstītos poļu karavīrus likteņa varā, var ilgstoši strīdēties gan no morālā, gan juridiskā viedokļa.

Vieni to attaisnos ar vēlmi turpināt cīņu pret nacismu, turpretī citi par piemēru, kā jāaizsargā savas tautas intereses, nosauks Francijas valdības rīcību 1940.~gadā, kura, kapitulējot Vācijai, pacentās kaut daļēji aizstāvēt franču tiesības un neokupētajā Francijas daļā izveidoja no nacistiem atkarīgo t.s. Višī režīmu (fr. \frtxti{le régime de Vichy}).

Krievu publicists J.~Muhins pat pārmet Polijas vadībai ne tikai savas tautas, bet arī savu sabiedroto~--- Francijas un Lielbritānijas nodevību. Polijas valdītāji ievilkuši Franciju un Lielbritāniju karā, bet paši bēguši. (Protams, karā lielvalstis ievilka nevis Polija, bet Ā.~Hitlers.) Viņš arī uzskata, ka Polijas valdība bēgot radīja ``šausmīgu juridisku kāzusu'': Polijas valsts vairs nav, bet nav arī miera, kapitulācija nav parakstīta, karš it kā turpinās, bet tas nozīmē, ka Lietuvai, Latvijai, Rumānijai, Ungārijai un PSRS jāgādā par internēto poļu apsardzību, ēdināšanu, apģērbšanu, taču~--- cik ilgi? Aizbēgot no valsts, prezidents un valdība zaudēja savu statusu, kļuva par ``bijušajiem'', kuri nevienu nepārstāv. PSRS vismaz pati bija iegājusi ``bijušās Polijas'' teritorijā, tāpēc tās pienākums bija rūpēties par internētajiem. Bet ``uz rumāņu, ungāru, lietuviešu un latviešu kakla poļi bija uzsēdušies paši un bez ielūguma.'' J.~Muhins uzskata, ka visi citi Polijas valdības rīcības varianti būtu pieņemamāki: ja tā atietu kopā ar karojošo armiju kā Beļģijas valdība Pirmajā pasaules karā, aizpeldētu uz Angliju zemūdenē, kuras ekipāžu varētu uzskatīt par karojošās Polijas armijas daļu, pat padodoties gūstā vāciešiem valdība tomēr paliktu valdība, ar to poļus pieņēmušās valstis varētu noslēgt līgumus par zaudējumu atmaksu. Minētie pārmetumi Polijas valdībai tikai daļēji ir pamatoti, aizpeldēšana zemūdenē uz Angliju no Polijas dienvidu pierobežas nebija iespējama, padošanās agresoriem: Vācijai vai PSRS~--- poļiem nepieņemama. Bēgšana bija kaunpilna, taču ja arī prezidents, valdība paliktu uz valsts kuģa komandtiltiņa, kuģis vairs nebija glābjams.

Jau minētais baltkrievu vēsturnieks un publicists A.~Usovskis raksta, ka Polijas vadība bēga uz Rumāniju ne tāpēc, ka baidītos no vāciešiem, bet gan tāpēc, lai Polija ``pārstātu pastāvēt kā poļu tautas nacionāla valsts'', bet tai pat laikā \latxti{de~iure} paliktu kara stāvoklī ar Vāciju. Citiem vārdiem~--- ja vadība paliktu Polijā, Vācija ar to būtu spiesta parakstīt kapitulāciju, nodiktējot Polijai savus noteikumus. Taču ar to karš, kuru Vācijai bija pasludinājušas Lielbritānija un Francija, ``karātos gaisā''~--- ja jau pati Polija savu pilnvaroto pārstāvju personā piekristu tai nodiktētajiem noteikumiem, tad zustu pamats britu un franču iesaistei jau atrisinātajā Polijas un Vācijas konfliktā. Tiem atliktu tikai izbeigt karu Rietumu frontē. Turpretī bēgot no valsts poļu vadītāji devuši Rietumu sabiedrotajiem likumīgu ieganstu turpināt karu, jo Polija \latxti{de~iure} nebija uzvarēta, turpināja, kaut niecīgiem spēkiem, vest karu. Ar to Polijas vadītāji demonstrējuši, ka starptautiskā kapitāla intereses tiem bijušas svarīgākas, nekā 35~miljonu cilvēku liktenis, kuri viņiem bija uzticējuši valsts vadību un kuri, savas valdības atstāti, nonāca pilnīgā Vācijas noteikšanā.

A.~Usovskis it kā spriež loģiski, taču, neievērojot dažus apstākļus, ir izveidojis patiesībai neatbilstošu ainu. Viņš neredz (vai izliekas neredzot), ka, paliekot Polijā, tās pirmskara vadība kādus okupācijas režīma atvieglojumus tautai diezin vai varētu panākt. Pārāk dziļš naids bija izaudzis gadsimtu gaitā starp vācu un poļu valdošajiem slāņiem un daļēji arī tautām, pārāk svarīga Ā.~Hitlera acīs bija poļu apdzīvotā teritorija, lai tur kādi ierobežotu Vācijas interešu īstenošanu. Jau vēlākais Višī režīma liktenis Francijā no 1940.~gada jūlija līdz 1944.~gada augustam rādīja, ka vācu okupanti faktiski nerēķinājās ar vietējo valdību, piemēram, 1942.~gada novembrī okupējot visu Francijas teritoriju.

Mūsdienās lielākā poļu vēsturnieku daļa jau nešaubās par to, ka labākas izejas kā atkāpties uz Rumāniju Polijas politiskajai un militārajai vadībai nebijis. Variantu par padošanos Vācijai neviens neizskatīja, Sarkanās armijas daļu ierašanās pie Rumānijas robežas tika gaidīta ļoti drīz, tās arestētu visus poļu varas nesējus, bet pāreja uz Rumāniju atvēra poļu politiķiem iespēju nonākt Rietumos.

Jāatzīst, ka nepiekrītot kapitulēt Vācijas priekšā, Polijas vadība rīkojās savas valsts interesēs. Cits jautājums~--- cik tām atbilda veco ``sanācijas'' vadītāju centieni palikt valsts struktūru vadībā arī pēc Polijas okupācijas, kādas viņiem uz to bija morālas un politiskas tiesības? Tā kā Polijas valdība tā arī neparakstīja aktu par kapitulāciju Vācijai, tad tiek uzskatīts, ka Polija juridiski nezaudēja savu suverenitāti.

17.~septembra rītā, atsaucoties Polijas valdības lūgumam, Rumānija bija piekritusi uzņemt Polijas valdības locekļus kā politiskus bēgļus savā teritorijā un atvieglot viņiem tranzītu uz citu neitrālu valsti, taču neļāva viņiem politiski darboties kā Polijas valdībai. Taču I.~Moscickis sagatavoja uzsaukumu, lai vērstos pie Polijas tautas un mēģināja to nosūtīt uz ārzemēm. Polijas valdības locekļi sniedza paziņojumus, kuros bija pausta griba ``turpināt karu''. Tādejādi Polijas prezidents un valdība Rumānijas teritorijā nevēlējās atteikties no politiskas darbības. Rumānijas valdība izmantoja šo formālo iemeslu, lai, ievērojot uz to izdarīto Vācijas spiedienu, Polijas prezidentu, valdības locekļus, tāpat kā armijas virspavēlnieku ar visu štābu, internētu.

Jāsaka, ka līdz ar starptautiskās situācijas maiņu, ar Polijas valdības bēgšanu no savas valsts krasi mainījās Latvijas valdības attieksme pret Poliju. Ja 1938.~gada oktobrī tā Poliju apsveica sakarā ar tās dalību Čehoslovākijas dalīšanā (kā norādījis V.~Semindejs, tad to, iespējams, noteica cerība, ka Polija neizvirzīs līdzīgas pretenzijas un Latvijas dienvidaustrumu rajoniem), tad 1939.~gada 21.~septembrī, pat neraugoties uz Lielbritānijas un Francijas protestiem, Latvijas valdība pārtrauca attiecības ar Polijas sūtniecību, aizbildinoties ar to, ka tai nav sakaru ar savu valdību. Profesors A.~Stranga Polijas sūtniecības slēgšanu vērtē kā neatbilstošu Latvijas apdraudētības pakāpei. Šī soļa vienīgais motīvs bija izpatikt PSRS un Vācijai. Tādējādi Latvija morāli nostājās agresoru pusē. Profesors Ē.~Jēkabsons detalizētāk pētījis Latvijas un Polijas attiecības pēc 1939.~gada 17.~septembra. Kā viņš uzskata, ``faktiski kopš šī brīža kara sākumā izsludināto Latvijas neitralitāti vairs nevar uzskatīt par pilnībā spēkā esošu, kaut arī attiecībā pret savā teritorijā internētajiem poļu karavīriem Latvijas valsts neitrālas zemes pienākumu pildīja līdz pat savai bojāejai 1940.~gada vasarā''. Autoram vienīgi jāpiebilst, ka Latvijai, tāpat kā citām Rietumu valstīm, arī agrāk bija pietrūcis spara atklāti vērsties pret agresoriem Eiropā, bet to, ka Latvijas valdībai bija pamats raizēties par valsts drošību, parādīja jau notikumi 1939.~gada oktobrī~--- PSRS kontingenta izvietošana Latvijā.

Par Latvijas pilsoņu~--- poļu attieksmi pret notiekošo kaimiņvalstī stāsta no Latvijas emigrējušā bijušā liepājnieka Kārļa Siljakova 1982.~gadā ASV publicētie memuāri.

``Kad divas nedēļas pēc vācu iebrukuma Polijā tur iebruka arī krievi un pēc kam arī drīzi sekoja Polijas valsts sadalīšana starp abiem iebrucējiem, Liepājas poļi, tāpat kā pārējie liepājnieki, jutās vīlušies, kad Anglija un Francija nepieteica karu arī otram iebrucējam~--- krievam. Bet poļu sāpju mērs ar to nebija pilns. Sašutuši par mūsu valdības rīcību, kas tūlīt pēc Polijas sakāves slēdza tās diplomātisko pārstāvniecību, Liepājas poļos varēja novērot arī zināmu naidu pret mums.

Viņi nerunāja par to atklāti un skaļi, bet viena otra netaktiska darba biedra zoboti par Polijas ātro sakāvi, viņi apvainoti un dusmās pateica arī savas domas par mums, un tās mums nekādi komplimenti nebija. Atgādinot, ka Polija ir bijusi vienīgā valsts, kas par savu palīdzību mūsu atbrīvošanas cīņās no mums nav prasījusi ne santīma. Viņi izteica apslēptus draudus un lāstus mūsu valdībai, kas labu tagad bija atmaksājusi ar ļaunu.

Saslēgušies savā kopībā vienotāki kā līdz tam, viņi sāka organizēt palīdzību pārbēgušiem un Kara ostā internētiem poļu karavīriem. Pateicoties viņu palīdzības gribai un mūsu uzraudzības iestāžu piemiegtām acīm, līdz laikam, kamēr Kara ostu iznomājām krievu bāzēm, tur poļu internēto gandrīz vairs nebija. Katrs, kas tikai bija gribējis, ar Liepājas poļu palīdzību bija atradis ceļu uz Zviedriju.

Pēc krievu iesēšanās Kara ostā un vāciešu repatriēšanās, poļu tērpus vairs nerotāja Polijas ērglis un viņu nacionālā aktivitāte Liepājā uz āru vairs nebija jūtama. Atbilde uz jautājumu, kur palikusi viņu lepnā pašapziņa, gandrīz vienmēr bija viena un tā pati:~--- Kas nu mēs vairs par poļiem, šeit dzimuši un auguši, bet tas nepārliecināja, jo tanī pašā laikā viņi centās uzņemt sakarus ar vēl pirms neilga laika no sava vidus atstumtiem no Polijas ievestiem laukstrādniekiem.

Es nezinu, pa kādiem ceļiem tanī laikā poļi bija ieguvuši ziņas par notikumiem Polijā pēc tās sadalīšanas, bet Liepājas poļi bija tie, kas pirmie zināja stāstīt par vācu mežonīgām izrīcībām, it sevišķi pret žīdiem un cilvēku masveidīgu bezvēsts pazušanu krievu ieņemtā teritorijā.''

Mūsdienu historiogrāfijā PSRS rīcība 1939.~gada 17.~septembrī attiecībā pret Poliju tiek vērtēta pretrunīgi. Gandrīz visi autori nosoda PSRS vadības paziņojumu, ka Polijas valsts 17.~septembrī jau praktiski neeksistēja. Tiesa, šim pārmetumam ir arī iebildes. Tā, krievu vēsturnieks A.~Orlovs uzskata, ka PSRS ieveda savu armiju Polijā brīdī, kad Polijas valdība jau zaudēja spēju pārvaldīt valsti, bet ``valdība emigrācijā'' vēl nebija izveidota, tāpēc starptautiskās tiesības bija ievērotas. Viņa oponenti savukārt uzdod jautājumus, vai tad valdības maiņas brīdī starptautiskās tiesības nedarbojas, un vai tad valsts pastāvēšanu nosaka valdības spēja to pārvaldīt?

Daudzi vēsturnieki uzskata, ka Sarkanās armijas ieiešanu Polijas teritorijā noteica jau slepenais protokols par ``interešu sfēru'' sadali pie 1939.~gada 23.~augusta Vācijas un PSRS līguma.

Taču citi pauž viedokli, ka, neraugoties uz to, ka Vācijas iebrukums Polijā bija gaidīts, vērmahta sekmīgā darbība Polijā un tās armijas ātrā sagrāve izrādījās pārsteigums padomju vadībai, kurai steidzīgi bija jāsper atbildes soļi. Pieauga padomju robežu apdraudējums, radās bažas, ka Vācija var neievērot neuzbrukšanas paktu, tāpēc arī PSRS ieveda Polijā savu karaspēku. Arī redzīgākie Rietumu politiķi prata pareizi novērtēt PSRS mērķus. Tā, ASV valsts sekretārs K.~Halls savos memuāros PSRS iebrukumu Polijā raksturoja kā ``mērķa noturēt Hitlera leģionus no pievirzīšanās pārāk tuvu Krievijai motivētu''.

Šim secinājumam pretī nerunā arī tas, ka Vācijas vadība jau septembra sākumā uzstājīgi aicināja PSRS ieņemt savu ``interešu sfēru''. Padomju vadībai vajadzēja rīkoties tā, lai iegūtu reālu varu savā ``interešu sfērā'' un neiesaistītos karā. Tēlaini var teikt, ka tai nācās izkuģot starp diviem tai reāliem briesmoņiem: Scillu un Haribdu (Vāciju un Rietumvalstu bloku), bet novirzīšanās uz jebkuru pusi draudēja ar ieraušanu karā. Bija jāņem vērā ļoti daudzi faktori. Tā, krievu vēsturnieks M.~Semirjaga uzskatīja, ka Sarkanās armijas ieiešanas Polijā novilcināšanu noteica nepieciešamība labāk psiholoģiski sagatavot šim solim PSRS iedzīvotājus, PSRS vadības bažas par iespējamo konfliktu ar Lielbritāniju un Franciju, vēlme nomierināt pasaules sabiedrisko domu un sagaidīt Varšavas krišanu. Tam var piekrist, vienīgi piebilstot, ka tad, ja Lielbritānija un Francija reāli iesaistītos karā Polijas pusē, Vācija būtu spiesta dalīt spēkus, Polijai nedraudētu tik ātrs sabrukums, nāktos runāt ne vairs par paredzētā Sarkanās armijas karagājiena Polijā ``novilcināšanu'', bet gan atcelšanu.

Diezgan bieži, īpaši publicistikā, tiek pausts viedoklis, ka kopš Molotova-Ribentropa pakta parakstīšanas PSRS pārstājusi būt neitrāla, esot kļuvusi par nacistiskās Vācijas sabiedroto, atbalstījusi Vācijas agresīvās tieksmes, sadarbojusies ar to, palīdzot sakaut Poliju. Šeit jāatgādina, ka ne neuzbrukšanas līgums, ne slepenais protokols pie tā neparedzēja abu valstu kopīgu rīcību pret Poliju. Vācija vienatnē jau bija sakāvusi Poliju, pirms sākās Sarkanās armijas kustība rietumu virzienā. PSRS palīdzību Vācijai, kuru nevar noliegt, šais nepilnās trijās nedēļās var vērtēt kā maznozīmīgu.

Norādot uz vērmahta un Sarkanās armijas sadarbības gadījumiem Polijā, virkne autoru apgalvo, ka 17.~septembrī faktiski PSRS iesaistījusies karā Vācijas pusē. Arī tas neatbilst patiesībai, jo pat ņemot vērā slepenos protokolus pie līgumiem, PSRS ar Vāciju militāri-politiski savienību (ne 23.~augustā, ne arī vēlāk~--- 28.~septembrī) tomēr neizveidoja, kaut virzība uz to šķita vērojama. 1939.~gadā Eiropa izrādījās sašķelta trijos militāri-politiskos blokos: angļu-franču, vācu-itāļu un padomju (ja vienu valsti var saukt par bloku). Katram no tiem bija savi mērķi, kas arī tos virzīja uz karu. Protams, PSRS pārkāpa starptautiskās tiesības, kad 17.~septembrī izvērsa agresiju pret Poliju. Taču arī tad tā nevis nostājās Vācijas pusē, bet cīnījās Polijā par savām interesēm. PSRS un Vācijas intereses tikai daļēji un uz laiku sakrita, abas valstis nekad otru neuzskatīja par savu stratēģisku sabiedroto, kaut propagandā tādi apgalvojumi arī pavīdēja. Tām katrai bija savi atšķirīgi mērķi, kas arī noveda tās pie konflikta 1941.~gadā.

Jāņem vērā, ka Vācijas armija, tās vadītājiem neko nezinot par 23.~augustā notikušo ``interešu sfēru'' sadali, kara darbības rezultātā 17.~septembrī bija jau šķērsojusi PSRS ``interešu sfēras'' robežas Polijā un turpināja virzīties uz austrumiem, piemēram, ieņēma Belostoku (\pltxti{Białystok}), Brestu, pat sasniedza Ļvovu. PSRS nu nekādi nebija ieinteresēta, lai Vācija ar šīm teritorijām izrīkotos pēc sava prāta (vai nu pati pārvaldītu, vai radītu no Vācijas atkarīgas poļu vai ukraiņu marionešu valstiņas, vai atrastu vēl kādu variantu). PSRS iesaistīšanās Polijas notikumos 17.~septembrī, vēlāk notikusī vērmahta atvilkšana no padomju ``interešu sfēras'' dažās Vācijas aprindās pat izsauca nepatiku. Arī tāpēc par abu valstu savienību nevar runāt.

Ja ne visi, tad neapšaubāms vairākums poļu vēsturnieku, arī krievu vēsturniece N.~Ļebedeva, ukraiņu vēsturnieks O.~Ovčarenko u.c. uzskata, ka, piedaloties kopējā ar Vāciju agresijā pret suverēno Poliju, lai arī bez kara pieteikšanas, bet, ievedot Sarkano armiju Polijas teritorijā, sankcionējot kaujas darbību pret tās armiju, PSRS staļiniskā vadība līdz ar to izvērsusi agresīvu karu un pārkāpusi veselu virkni starptautisku līgumu:

\begin{itemize}
\item PSRS miera līgumu ar Poliju, parakstītu 1921.~gada 18.~martā Rīgā;

\item 1929.~gada 9.~februāra protokolu par Briāna-Kelloga pakta (noslēgts 1928.~gada 27.~augustā Parīzē), kurš aizliedza izmantot karu kā nacionālās politikas instrumentu, pirmstermiņa stāšanos spēkā;

\item 1932.~gada 25.~jūlija PSRS un Polijas līgumu par neuzbrukšanu un 1934.~gada 5.~jūlija protokolu, ar kuru tas tika pagarināts līdz 1945.~gadam; (Pirms uzbrukuma Polijai Vācija 1939.~gada 28.~janvārī bija uzteikusi 1934.~gadā parakstīto neuzbrukšanas līgumu ar Poliju, bet PSRS 1932.~gada līgumu ar Poliju par neuzbrukšanu 1939.~gada septembrī nebija denonsējusi.)

\item 1933.~gada Londonas konvenciju par agresora noteikšanu (1933.~gada 3.~jūlijā Afganistānas, Igaunijas, Latvijas, PSRS, Polijas, Persijas, Rumānijas un Turcijas parakstītā konvencija par agresijas noteikšanu saturēja pantu, ka par agresoru tika uzskatīta arī tā valsts, kuras bruņotie spēki, kaut arī bez kara pieteikšanas, iebruka citas valsts teritorijā. Šo pantu varēja attiecināt arī uz PSRS rīcību 1939.~gada 17.~septembrī pret Poliju.);

\item 1938.~gada 26.~novembrī Maskavā publicēto PSRS un Polijas valdību kopīgo komunikē, kurā vēlreiz tika apstiprināts, ka abu valstu attiecību pamatā atrodas 1932.~gada līgums par neuzbrukšanu.
\end{itemize}

Visi šie dokumentu saturēja pantus, kas uzlika par pienākumu to dalībniekiem atturēties no vardarbības pret citiem vienlīdzīgiem starptautisko tiesību subjektiem.

Tomēr jāsaka, ka minētie uzskati ir tendenciozi un apstrīdami.

Vispirms par 1933.~gada konvenciju par agresora noteikšanu. No 1933.~gada, kad tā tika parakstīta, līdz 1939.~gadam pasaulē daudz kas bija mainījies. Minētajai konvencijai bija jākļūst tikai par citas konvencijas~--- par bruņojuma samazināšanu un ierobežošanu~--- sastāvdaļu. Taču tā nekad netika parakstīta. Kopš romiešu tiesību rašanās spēkā ir aksioma, ka, nedarbojoties pamatlīgumam, nestājas spēkā arī tam pievienotais. Konvencijas ievērošana no to parakstīšanas puses būtu labas gribas žests, bet, ievērojot to, ka galamērķis~--- bruņošanās ierobežošana~--- nebija sasniegts, savā laikā veiktā konvencijas parakstīšana netraucēja padomju vadītājiem vairāk rūpēties par savas valsts interesēm nekā labas gribas demonstrēšanu.

Tālāk~--- par kara uzsākšanu, pārkāpjot veselu virkni līgumu. PSRS karu Polijai nepieteica. Arī Polijas valdība, (ne 17.~septembrī vēl Polijas teritorijā esošā, ne vēlāk emigrācijā nodibinātā) vismaz \latxti{de~iure} nevērtēja savas attiecības ar PSRS kā kara stāvokli, neraugoties uz simtiem nogalināto un desmitiem tūkstošu sagūstīto Polijas pilsoņu. Tāpat to nedarīja arī Polijas sabiedrotās un citas valstis. Pēc 1939.~gadā spēkā esošajām starptautiskajām tiesību normām PSRS karaspēka ievešana Polijā nevarēja tikt iztulkota kā kara pieteikšana. Kā atzīmējis ukraiņu tiesību speciālists V.~Makarčuks ``no tiesību, \latxti{de lege lata} (lat.~--- no likuma viedokļa) normas viedokļa'' PSRS armijas ievešanu Polijā nevarēja traktēt kā kara sākšanu un tā tas arī netika traktēts. Tātad, ja nebija kara, tad netika pārkāpts arī Briāna-Kelloga pakts.

Jāņem vērā arī pirms Otrā pasaules kara pastāvošās Tautu Savienības, kuras viens no mērķiem bija iespējamo karu novēršana, rīcība. Savienības uzdevums bija organizēt savus locekļus agresijas atvairīšanai. Saskaņā ar Savienības statūtu 16.~punktu, visām tās dalībvalstīm bija pienākums pārtraukt tirdzniecības un finansiālās attiecības ar agresoru. Tautu Savienība PSRS par agresoru neatzina un neviena valsts attiecības ar to nepārtrauca. Turpretī, kad pēc pāris mēnešiem PSRS uzsāka karu pret Somiju, Tautu Savienība to atzina par agresoru un izslēdza no savām rindām.

Tiesa, 1921.~gada Rīgas miera līgumu un 1932.~gada PSRS un Polijas līgumu par neuzbrukšanu PSRS pārkāpa. Tomēr līdz t.s. Atlantijas hartas (1941.~gada 14.~augusts) pieņemšanai neviena no lielvalstīm neapšaubīja apdraudēto valstu ``tiesības'' ar ieroču palīdzību pārņemt kontroli pār teritorijām, kas ir svarīgas to interesēm. Kara draudu priekšā arī Rietumu lielvalstis rīkojās, rūpējoties par savu drošību, savām interesēm, nevis par to rīcības atbilstību morāles vai starptautisko tiesību normām.

Tā, 1940.~gada 5.~februārī Lielbritānijas un Francijas Augstākā padome, sakarā ar PSRS un Somijas karu apstiprināja plānu, kurš paredzēja, pārkāpjot Norvēģijas un Zviedrijas neitralitāti un suverenitāti, militāri okupēt virkni Norvēģijas pilsētu un dzelzceļa līniju līdz Luleo (zv. \svtxti{Luleå}) pilsētai Zviedrijā. Plāns gan netika īstenots, jo PSRS un Somijas karš beidzās, bet Norvēģiju pasteidzās okupēt Vācija. Taču bija arī reāli gadījumi, kad neatkarīgas valstis tika demokrātisko lielvalstu okupētas. Tā, pēc tam, kad Vācija 1940.~gada 9.~aprīlī okupēja Dāniju, maijā vispirms Lielbritānijas, bet pēc tam arī tolaik neitrālās ASV karaspēks izcēlās Islandē, kura atradās ūnijas attiecībās ar Dāniju. 1941.~gada augustā ar ieganstu, lai nepieļautu Vācijas ietekmes pieaugumu valstī, Lielbritānija un PSRS ieveda savu karaspēku Irānā. Tādejādi ieganstu, ka kādas teritorijas var sagrābt un izmantot nedraudzīga valsts, šīs teritorijas okupācijai Otrā pasaules kara laikā izmantoja ne tikai totalitārā PSRS, bet arī demokrātiskās Rietumvalstis. To rīcība šai ziņā neatšķīrās no PSRS rīcības Rietumukrainā un Rietumbaltkrievijā (un arī nedaudz vēlāk Baltijā). PSRS rīkojās atbilstoši tai laikā spēkā esošajām starptautiskajām tiesiskajām normām. Ja uzskata PSRS rīcību 1939.~gada septembrī par nelikumīgu, par tādu pat jāuzskata arī Lielbritānijas un ASV rīcība 1940--1941.~gadā.

Cita lieta, ka PSRS ``atbrīvoto'' teritoriju iedzīvotāji nonāca staļiniskā režīma varā, nesa neskaitāmus upurus, ko nevar teikt, piemēram, par Rietumu okupēto Islandi. Bet par to~--- turpmākajās lappusēs.

Protams, forma, kādā PSRS lauza neuzbrukšanas līgumu ar Poliju, neatbilda tādiem gadījumiem paredzētajai procedūrai. No juridiskās puses tā bija tieša PSRS pieņemto saistību neizpilde. Vēl vairāk, padomju-poļu līgumā bija paredzēts, ka ``par rīcību, kura ir pretrunā ar uzņemtajām saistībām, tiks uzskatīts katrs vardarbības akts, kurš pārkāpj otras valsts teritoriālo vienotību un neaizskaramību, pat ja šī rīcība tiktu realizēta bez kara pieteikuma un neradot visas iespējamās tam sekojošās sekas''. No juridiskās puses PSRS rīcību tiešām var kvalificēt kā agresiju. Taču krievu vēsturnieki norāda, ka nedrīkst aizmirst to, ka Rietumukraina un Rietumbaltkrievija nonāca Polijas rokās pašas Polijas agresijas pret saviem austrumu kaimiņiem rezultātā, un PSRS rīcība 1939.~gada septembrī bez visa cita bija arī revanšs par 1919.--1920.~gadā zaudēto karu un tā rezultātā pazaudēto teritoriju atgūšana. Piemēram, krievu publicists V.~Kožinovs rakstīja, ka tas, kas notika 1939.~gadā, pēc vēsturiskās būtības bija nevis PSRS agresija pret Poliju, bet Polijas agresijas seku likvidēšana.

Tiesa, šiem vēl mūsdienās lielvalstiskās un impēriskās (ne tikai Krievijas, bet arī citu imperiālistisko lielvalstu) domu gaitas izteicēju uzskatiem, tas, ka PSRS 1939.~un 1940.~gadā tikai atgrieza savas, t.i. kādreiz Krievijas impērijā ietilpušās zemes, ar to atšķīrās no Vācijas, kura 1938.~gadā anektēja Sudetus, kas nebija tai piederējuši, bet pēc tam sagrāba citu valstu teritorijas, ir savs vājais punkts. Šāda domāšana nav savienojama ar uz tautu pašnoteikšanās principu balstītām starptautiskām tiesībām. Tikai nosacīti var runāt par to, ka Padomju Savienībai pievienoto bijušo Polijas teritoriju iedzīvotāji ieguva savu valstiskumu un labākus dzīves apstākļus. Valstiskuma izveide, par ko vēl tiks runāts turpmāk, tā formas izvēle netika nodota pašu baltkrievu un ukraiņu tautu rokās. Padomju Savienībā jaunpievienotajos apgabalos realizētās represijas, lai gan tika izvērstas pēc citiem kritērijiem nekā represijas ``sanācijas'' Polijā, nebija labākas par tām. Vienīgais, kas pēc autora domām, daļēji mīkstina PSRS vēsturisko vainu par 1939.~un arī 1940.~gada aneksijām, bija šo teritoriju neatdošana agresīvajai nacistiskajai Vācijai, kas tikai vairotu tās ekonomisko spēku un apetīti, bet to iekļaušana PSRS aizsardzības zonā pret sagaidāmo nacistiskās Vācijas agresiju. Saprotot, ka 1939.~gadā nekāda pilnīgi brīva baltkrievu un ukraiņu tautas gribas izpausme nebija iespējama, mūsdienās, vērtējot notikumu attīstību 1939.~gadā no pasaules civilizācijas interešu viedokļa, PSRS veiktā Rietumukrainas, Rietumbaltkrievijas un Viļņas apgabala aneksija jāvērtē kā mazākais ļaunums salīdzinājumā ar iespējamo visas Polijas teritorijas nonākšanu vācu nacistu rokās.

Bez tam, ārkārtējos gadījumos diezgan bieži tiek piesaukta būtisko apstākļu maiņas doktrīna (lat. \latxti{rebus sic stantibus}), kas paredz, ka viena no līgumslēdzējpusēm var no tā atteikties, ja līguma slēgšanas laikā pastāvošie apstākļi ir būtiski mainījušies, to maiņu nevarēja iepriekš paredzēt, bet šo apstākļu pastāvēšana veidoja būtisku pamatu pušu piekrišanai līgumam. Šai gadījumā PSRS bija iespēja atsaukties uz to, ka, parakstot neuzbrukšanas līgumu, tā bija rēķinājusies ar to, ka Polija saglabās savu suverenitāti un kalpos PSRS kā sava veida vairogs starp to un iespējamo agresoru (Vāciju). Kad Polija stāvēja suverenitātes zaudēšanas priekšā un vairs nespēja pildīt šo vairoga lomu, PSRS bija tiesības rīkoties, lai aizstāvētu savas intereses. Doktrīnas pielietošanas tiesiskumu var apšaubīt un abu strīdā iesaistīto pušu juristu viedokļi parasti atšķirsies. Taču jāatzīst, ka dzīvē starpvalstu līgumi darbojas tikai tik ilgi, kamēr tas ir izdevīgi abām pusēm. Šajā gadījumā PSRS intereses tās valdībai likās svarīgākas nekā līguma ievērošana. Faktiski PSRS rīkojās tai pat virzienā kā Lielbritānija un Francija, nepildot savas sabiedroto saistības pret Poliju. Kā atzinis vācu vēsturnieks G.~Manns: ``Šodien parakstīt līgumus un rīt tos lauzt, maldināt, apkrāpt, atsevišķus indivīdus noslepkavot, veselas rases iznīcināt,~--- tā tas vienmēr ir bijis starp cilvēkiem; \citespace{} Aplūkojiet pasaules vēsturi, nē, nevis kādai tai ir jābūt saskaņā ar kristīgo morāli, bet kāda tā patiešām ir, tad nevar noliegt šīs teorijas patiesumu."

Kaut Vācijas līderi PSRS iebrukumu Polijā gaidīja jau agrāk, tomēr arī tagad tie bija apmierināti. 1939.~gada 19.~septembra runā Dancigā Ā.~Hitlers triumfēja: ``Polija nekad vairs nepiecelsies.''

\strong{Rietumvalstu reakcija} uz PSRS rīcību 1939.gada septembrī bija visai uzmanīga.

Protams, Rietumu presē par 17.~septembrī sākto Sarkanās armijas ievešanu Polijā skanēja sašutums. 18.~septembrī Londonas ``\entxti{The Times}'' raksturoja PSRS soli kā ``naža dūrienu Polijas mugurā''. Lielbritānijas un Francijas attiecības ar PSRS saasinājās. Var teikt, ka brīdī, kad starp Rietumvalstīm un Vāciju norisa ``karsts'' karš, starp Rietumvalstīm un PSRS notika sava veida ``psiholoģisks'' karš. Taču tika publicēti arī raksti, kas paskaidroja padomju rīcību, tās pretvācisko virzību.

Rietumos bija arī reālpolitiķi, kuri pozitīvi novērtēja PSRS rīcību. Tā, jau 17.~septembrī Francijas sūtnis Varšavā L.~Noels rakstīja: ``PSRS noteikti bija jāieved sava armija, kamēr vēl nav par vēlu.''

Pēc Sarkanās armijas iesoļošanas Polijā Lielbritānijas valdība 18.~septembrī saskaņā ar Polijas un Lielbritānijas Savstarpējās palīdzības līguma slepeno protokolu konstatēja, ka atbilstoši britu-poļu 1939.~gada 25.~augusta līgumam Lielbritānijai ir jāaizstāv Polija tikai Vācijas agresijas gadījumā un ``Lielbritānijai pēc padomju invāzijas Polijā nav nekādu saistību iesaistīties karā ar PSRS''. Tāpēc Lielbritānijas valdība pieņēma zināšanai PSRS valdības paziņojumu, kurā tā piedāvāja neitralitātes ievērošanu ``pret valdībām, ar kurām Padomju valdība uztur diplomātiskās attiecības''~--- tātad arī Lielbritānijas valdību. Tika nolemts ``nesūtīt Krievijai nekādu protestu''. Briti aprobežojās ar paziņojumu par padomju rīcības nelikumību un savu agrāko saistību Polijas priekšā apstiprināšanu.

Notika viedokļu apmaiņa starp Londonu un Parīzi, un abas valdības nonāca pie secinājuma, ka kara pieteikšana PSRS nekādā gadījumā neglābs Poliju, bet var ciešāk piesaistīt Maskavu Berlīnei un negatīvi ietekmēt Vācijas ekonomiskās blokādes efektivitāti. Francija gan pieprasīja no PSRS ``papildus paskaidrojumus''. 19~.~septembrī Maskavā tika saņemta nota, kurā bija norādīts, ka ar PSRS valdības izvirzītajiem argumentiem uzbrukumu Polijai attaisnot nevar, prasīts pārtraukt Sarkanās armijas virzīšanos un izvest to no Polijas. Taču, kad Padomju vadība notu ignorēja, tas nekādas sekas neradīja.

Tai pat 19.~septembrī visām britu diplomātiskajām pārstāvniecībām tika sniegti valdības izskaidrojumi, ka Lielbritānija negatavojas pieteikt karu PSRS, bet tai ir jāpaliek pēc iespējas labākās attiecībās ar to. Britu diplomātiem ārvalstīs tika dots rīkojums Lielbritānijas valdības atturīgo rīcību pret padomju agresiju skaidrot ar to, ka formāli starp PSRS un Poliju karastāvokļa nav, kā arī ar to, ka Polijas valdība atstājusi valsti, tātad~--- lietot vienu no tiem argumentiem, kuri bija izmantoti PSRS valdības notā Polijas valdībai 17.~septembrī. Ar to vismaz daļēji padomju argumentācija tomēr tika pieņemta. Britu presei tika dots norādījums pārtraukt jebkuru pretpadomju propagandu.

Kad Polijas sūtnis Londonā E.~Račinskis Lielbritānijas Ārlietu ministrijā ierosināja Lielbritānijai saskaņā ar tās un Polijas līgumu pieteikt PSRS karu (ko pati Polija gan nebija darījusi), Lielbritānijas ārlietu ministrs E.~Halifakss varēja sūtnim paziņot, ka tā ir Lielbritānijas valdības prerogatīva~--- pieteikt vai nepieteikt PSRS karu. Bijušais britu premjers D.~Loids Džordžs 23.~septembrī vēstulē E.~Račinskim norādīja: ``Būtu noziedzīgs neprāta akts krievu virzīšanos nostādīt līdzās ar vācu virzīšanos.''

Arī Polijas oficiālais sabiedrotais pret PSRS~--- Rumānija PSRS rīcībā nesaskatīja agresiju un karu agresoram nepieteica.

Arī ASV vadība nolēma, ka 1921.~gada Rīgas līgumā noteikto Polijas robežu pārkāpšanu no padomju karaspēka puses nevar kvalificēt kā kara aktu. Ņemot vērā ilglaicīgās ASV intereses, uz PSRS netika attiecināts embargo (daļējs vai pilnīgs aizliegums izvest no kādas valsts vai arī ievest tajā preces, valūtu u.~tml.), ko prasīja valstī spēkā esošais likums par neitralitāti. Faktiski Rietumvalstu rezervētās nostājas pret PSRS pragmātiskos cēloņus atklāja ASV valsts sekretārs K.~Hells, kurš rakstīja: ``Kaut krievu uzbrukumu Polijai varēja atzīt par militāru akciju, prezidents Rūzvelts un es \citespace{} nevēlējāmies uzskatīt Krieviju par valsti, kas karo vienādā mērā kā Vācija, jo tā rīkojoties, mēs vēl vairāk grūstu Krieviju Hitlera apskāvienos.''

Tomēr, neraugoties uz visām atrunām, Lielbritānijas un Francijas vadītāju atturību, pasaules sabiedrība pret PSRS, tāpat kā savulaik pret Poliju Cešinas gadījumā, izturējās kā pret nacistu sabiedroto.

Taču atgriezīsimies pie notikumiem Polijā 17.~septembrī.

Pie robežas ar PSRS esošais poļu karaspēks~--- 17~kājnieku pulki un 6~kavalērijas eskadroni (pavisam Polijas austrumu vojevodistēs atradies ap 340~000 poļu kareivju un virsnieku, 540~lielgabalu, vairāk nekā 70~tanku, 5~bruņuvilcieni, ap 160~lidmašīnu)~--- mēģināja atkāpties uz robežu ar Rumāniju, neiesaistoties kaujās ar Sarkano armiju, taču tā kā reiz centās atbruņot poļu daļas, tāpēc poļiem atlika vai nu mēģināt izlauzties ar kauju, vai padoties. Poļu daļas parasti padevās bez kaujas. Taču norisa arī sadursmes.

Pēc sadursmēm uz robežas un robežsargu daļu atkāpšanās bieži padomju karaspēka priekšā vairs nebija apakšvienību, kuras spētu turpināt cīņu. Vairākums pilsētu un ciematu tika atdoti bez kaujas.

Var teikt, ka nopietnākās cīņa Sarkanajai armijai ar poļu daļām notika frontes ziemeļos. No turienes tikt līdz Ungārijai vai Rumānijai nebija iespējams, tāpēc poļu karavīri centās izlauzties uz Lietuvu un Latviju, kas atsevišķos gadījumos tiem arī izdevās.

Organizēta pretošanās, ilgāka par diennakti, norisa tikai dažos gadījumos: Viļņas, Grodņas, Tarnopoles pilsētās u.c. Sarkanā armija pieļāva kļūdu~--- mēģināja pilsētas ieņemt triecienā, bet tur poļu karaspēku atbalstīja daļa poļu iedzīvotāju, dažas no tām bija nocietinātas. Pārsvarā pretošanos izrādīja robežsardzes korpusa (\pltxti{Korpus Ochrony Pogranicza}), žandarmērijas un poļu zemessargu daļas. 18.~septembrī Grodņā sākās apšaude starp poļu karaspēku un prokomunistiski noskaņotiem pilsētniekiem. Pēdējiem izdevās atbrīvot politieslodzītos no vietējā cietuma. Pretēji E.~Ridzs-Smiglija pavēlei Grodņā nopietnu pretestību Sarkanajai armijai izrādīja ap 3~000 mācību vienības virsnieku un karavīru, kā arī policistu un zemessargu. Tikai 22.~septembrī Sarkanās armijas daļām izdevās ieņemt pilsētu. Kaujās par pilsētu tā zaudēja 57~kritušos un 159~ievainotos, 19~tankus un 4~bruņumašīnas. Kaujas laukā palika arī 644~poļi, gūstā krita 1~543~poļu karavīri. Poļu literatūrā ir ziņas, ka šeit pēc pilsētas ieņemšanas tikuši nošauti ap 300~tās aizstāvji, tādi pat nošaušanas gadījumi bijuši arī citur, kur poļu daļas pretojušās.

Parasti, ja nebija cerību tikt līdz Rumānijai vai Ungārijai, poļu vienības nodeva ieročus Sarkanās armijas daļām, pretī saņemot solījumu, ka personālsastāvs tiks atbrīvots. Pēc nepilnīgām ziņām jau līdz 21.~septembrim Sarkanās armijas un PSRS IeTK darbinieku rokās bija nonākuši ap 120~tūkstoši poļu karavīru un virsnieku. Gan poļu karavīri, gan virsnieki padomju gūstā devās visai mierīgi, nemēģinot slapstīties, lai, piemēram, no Baltkrievijas pārietu Lietuvas robežu, kas sākotnēji bija samērā viegli izdarāms. Viņi neuzskatīja Padomju Savienību par pretinieku.

Kaut jau pirms ``atbrīvošanas karagājiena'' sākuma 11.~septembrī PSRS iekšlietu tautas komisārs L.~Berija parakstīja pavēli par karagūstekņu arestu noformēšanu, bet 15.~septembrī tika parakstīta direktīva par IeTK grupu uzdevumiem ``kārtības nodrošināšanā, kaitnieciskās un kontrrevolucionārās darbības nepieļaušanā'' (pēc dažiem datiem Sarkanās armijas daļām sekoja ap 200~000 cilvēku no PSRS IeTK un robežapsardzes karaspēka daļām, kas nodarbojās ar ``kontrrevolūcijas'' apspiešanu). Sarkanās armijas daļās pirmajās ``atbrīvošanas karagājiena'' dienās valdīja neskaidrība, ko darīt ar milzīgajām poļu gūstekņu masām. Sarkanarmieši neievēroja stingrus gūstekņu konvojēšanas noteikumus. Valdīja sajukums un nekārtība. Baltkrievijas PSR iekšlietu ministrs L.~Canava, kurš sekojot armijas daļām, ieradās bijušajās Polijas teritorijās, uz Maskavu ziņoja: ``Rietumu rajonos salasījies daudzu tūkstošu liels no frontes bēgošo [poļu] karavīru daudzums, kuri aizšķērso ielas, bet ar [IeTK] operatīvās grupas spēkiem izolēt viņus nav iespējams. SZSA daļas viņus gūstā neņem, rezultātā neviens neveic filtrāciju un poļu karavīru kustība notiek brīvi''. Tikai 20.~septembrī abu fronšu pavēlnieki izdeva pavēles par gūstekņu konvojēšanas kārtību, to virzības ceļiem un punktiem, kur tie jānodod IeTK iestādēm. Taču straujā Sarkanās armijas virzība radīja lielas grūtības tām daļām, kuras nodarbojās ar gūstekņu transportēšanu. Karagūstekņu nometnes tika vēl tikai veidotas un nespēja uzņemt visus turp nosūtītos. Aizturētos bieži nebija ar ko barot, kur izmitināt zem jumta, nereti tas izrādījās staļļa vai kūts jumts.

Tāpēc jau 21.~septembrī PSRS aizsardzības tautas komisāra vietnieks G.~Kuļiks rakstīja J.~Staļinam, V.~Molotovam un K.~Vorošilovam: ``Gūstā saņemts ļoti daudz ierindnieku un virsnieku. Uzskatu, ka valdībai jādod rīkojums pēc uzskaites sagūstītos baltkrievus un ukraiņus atlaist mājās, jo barot viņus nav ar ko, konvojēšana prasa ļoti daudz cilvēku.'' Līdzīgus lūgumus sūtīja arī armijas komandieri. Atbildot uz to, aizsardzības tautas komisārs K.~Vorošilovs un SZSA Ģenerālštāba priekšnieks B.~Šapošņikovs 23.~septembrī izdeva pavēli: ``Sagūstītos Rietumbaltkrievijas un Rietumukrainas zemniekus, ja viņi uzrāda dokumentus, kas apliecina, ka viņi tiešām bija poļu mobilizēti, tiek atļauts atbrīvot.''

Sākumā tā arī notika. Taču Sarkanās armijas Politpārvaldes priekšnieks Ļ.~Mehliss 24.~septembrī ziņoja J.~Staļinam, ka pēc pavēles par karagūstekņu atlaišanu poļu karavīru ir pilni ceļi, daudzi no viņiem cenšas nonākt vācu ieņemtajos rajonos. Viņš prasīja, lai tiktu izpildīta pagaidām uz papīra palikusī pavēle par visu militāristu un policistu aizturēšanu un nodošanu IeTK. Acīmredzot J.~Staļins atbalstīja Ļ.~Mehlisu, jo 25.~septembrī K.~Vorošilovs atcēla savu 23.~septembra pavēli par karagūstekņu-vietējo zemnieku atbrīvošanu. 26.~septembrī ukraiņu un baltkrievu atbrīvošana tika pārtraukta. Poļu vēsturnieki uzsver, ka pēc ieroču nodošanas parasti visus poļus, arī tos, kuri padevās, saņemot solījumu tikt atbrīvotiem, pretēji solītajam nosūtīja uz karagūstekņu nometnēm. Arī pēc kaujām turpinājās visu to, kuri bija dienējuši Polijas armijā, tvarstīšana, lai nosūtītu viņus uz nometnēm. To vajadzībām parasti pielāgoja klosteru telpas.

Uz mājām ejošie bijušie poļu armijas karavīri atkārtoti tika aizturēti. PSRS IeTK grupas arestēja ne tikai karavīrus, bet faktiski visus, kas nēsāja formu. Polijā tādu bija daudz: policisti, ugunsdzēsēji, dzelzceļnieki, studenti, jaunatnes organizāciju dalībnieki, pasta darbinieki.

Mūsdienu poļu vēsturnieki apgalvo, ka, tā kā sarkanarmiešu vidū bijis izplatīts uzskats, ka Polijas valsts ir beigusi pastāvēt, tiem radies priekštats, ka attiecībā pret tās pilsoņiem nedarbojas arī starptautiskās tiesības un gūstekņus var iznīcināt, pielietojot ``sarkano teroru''. Pilnīga nesodāmība savienojamā ar ``šķiru cīņas'' ideoloģiju veicinājusi gūstekņu, arī ievainoto un slimo, nogalināšanu. Tādā veidā bojā gājušo skaits sniedzies tūkstošos. Tiesa, fakti par ``gūstekņu, arī ievainoto un slimo, nogalināšanu'', kas ``sniegušies tūkstošos'', netiek minēti.

Notika arī dažas Sarkanās armijas un vērmahta sadursmes. Dažviet, neraugoties uz Vācijas pavēlniecības veiktajiem drošības pasākumiem, dažas vācu daļas nokļuva zem padomju aviācijas triecieniem. Tāpēc Vācijas spēki steigšus tika atvilkti aiz ``interešu sfēru'' līnijas, par kuru PSRS un Vācijas vadītāji bija vienojušies 23.~augustā un kuru vērmahts, vajājot poļu daļas, bija jau paguvis pārkāpt. Dažos apvidos 19.~septembrī šī atkāpšanās bija jau notikusi. Taču pie Ļvovas vācu spēki, aizbildinoties ar nepieciešamību vispirms sagraut pilsētā no 10.~septembra ielenkto poļu armijas grupējumu, vēl palika uz austrumiem no nospraustās ``interešu sfēru'' robežas.

19.~septembrī arī Sarkanās armijas daļas pietuvojās Ļvovai. Poļu artilērija atklāja pret tām uguni. Daļa padomju tanku tomēr nonāca pilsētas centrā. Sākās padomju un poļu daļu komandieru sarunas. Taču pēkšņi poļiem uzbruka arī vācu daļas, izraisījās apšaude. Padomju tanki un bruņumašīnas izrādījās starp divām ugunīm~--- poļu un vācu. Kad padomju tankistu mēģinājumi pārtraukt apšaudi, iezīmējot sevi ar sarkaniem un baltiem karodziņiem, neizdevās, tie atklāja atbildes uguni. Tā rezultātā vācieši zaudēja trīs prettanku lielgabalus, bija krituši trīs virsnieki un ievainoti deviņu kareivji. Padomju puse zaudēja divas bruņumašīnas un vienu tanku, bija nogalināti trīs, bet ievainoti četri cilvēki. Pārpratums gan drīz noskaidrojās, bet sarkanarmiešiem un vācu karavīriem nācās savākt savus kritušos un ievainotos.

Vācijas varas iestāžu mēģinājumi panākt, lai, neraugoties uz slepenajā protokolā pie 23.~augusta līguma nospraustajām ``interešu zonām'', Vācijas rokās paliktu arī Rietumukrainas naftu saturošie rajoni, sastapa padomju valdības pretsparu. Tā norādīja, ka incidentus var izbēgt, ja Vācija atvilks savu karaspēku. 20.~septembrī Ā.~Hitlers bija spiests izdot pavēli par vērmahta atkāpšanos no Ļvovas, bet 21.~septembrī starp PSRS un Vāciju tika parakstīta vienošanās par vācu armijas daļu atvilkšanas kārtību uz rietumiem no agrāk noteiktās ``interešu sfēru'' līnijas'' pa Narevas~--- Vislas~--- Sanas upēm. Vācijas vadība bija spiesta atteikties arī no Ļvovas-Drohobičas (ukr. \uktxti{Дрогобич}) naftas rajona, ko septembra sākumā bija ieņēmis vācu karaspēks un ko Vācija labprāt paturētu savas rokās. Vācu ģenerāļi, kuri nebija informēti par slepeno protokolu, bija neapmierināti, bet pavēli izpildīja.

Poļu daļu pavēlnieks pilsētā ģenerālis V.~Langners uzskatīja par labāku pārtraukt pretošanos un atdot Ļvovu Sarkanajai armijai, paziņojot: ``Viņi, ģermāņi, ir visa slāviskā ienaidnieki, bet jūs esat slāvi''. (Taču padomju puse neizpildīja kapitulācijas noteikumus~--- atļaut karavīriem brīvi aiziet Ungārijas un Rumānijas virzienā. Viņi tika ievietoti gūstekņu nometnēs. Pašam ģenerālim gan izdevās aizbēgt uz Rumāniju.)

Interesanti, ka Beskidu (ukr. \uktxti{Бескиди}, poļu \pltxti{Beskidy}) dzelzceļa stacijā Rietumukrainā padomju karavīri sastapās ar\dots{} ungāriem. Ungārija pretendēja uz daļu Polijas teritorijas, tāpēc arī acīmredzot kāda to karaspēka daļa bija devusies pāri robežai. Ieraugot sarkanarmiešus, ungāru karavīri atklāja uguni. Taču sadursmē iesaistījās padomju bruņumašīnas un ungāriem caur tuvumā esošo tuneli nācās atkāpties uz savas valsts teritoriju. Sarkanās armijas daļas viņiem nesekoja.

22.~septembrī padomju tanku brigāde iegāja Brestā, kuru jau bija ieņēmušas vācu daļas. Tā kā pēc 23.~augusta līguma slepenā protokola pilsēta atradās padomju ``interešu sfērā'', vācu karaspēkam tā bija jāatstāj.

No ģenerāļa H.~Guderiāna atmiņām vēstures literatūrā ir iegājis stāsts par it kā Brestā notikušo kopējo vērmahta un Sarkanās armijas daļu ``parādi'' (šādu terminu lietoja H.~Guderiāns), uz kuru atsaucas vēsturnieki, kuri vēlas pasvītrot abu valstu sadarbību. Lietas būtība: Vācijas 19.~motorizētais korpuss, kuru komandēja ģenerālis H.~Guderiāns, 14.~septembrī ieņēma Brestas pilsētu, bet 17.~septembrī arī cietoksni. Taču pēc Molotova- Ribentropa pakta slepenā protokola pilsētai bija jānonāk PSRS ``interešu sfērā''. Tāpēc tika saskaņota pilsētas nodošanas kārtība. Ģenerālis H.~Guderiāns ierosināja sarīkot pilnvērtīgu karaspēka parādi, ar iepriekšēju karaspēka daļu nostāšanos un sekojošu svinīgu soļošanu pavēlniecības priekšā. Padomju augstākais virsnieks brigādes komandieris (\rutxti{комбриг}) S.~Krivošeins savās atmiņās stāsta, ka viņš piedāvājis citu pilsētas nodošanas procedūru: ``Pulksten 16-os Jūsu korpusa daļas pārgājiena kolonā ar karogiem (\rutxti{штандартами}) priekšgalā atstāj pilsētu, manas daļas, tāpat pārgājiena kolonā, ieiet pilsētā, apstājoties ielās, pa kurām virzās vācu pulki, un ar saviem karogiem salutē garāmejošajām daļām. Orķestri izpilda militārus maršus.'' H.~Guderiāns, kurš sākotnēji uzstājās par pilnvērtīgu parādi, gala rezultātā S.~Krivošeina piedāvājumam piekrita, ``tiesa, vienojoties, ka viņš kopā ar mani stāvēs uz tribīnes un sveiks garāmejošās daļas''. Tā arī notika. Sākumā virzījās vācu kājnieki, tad motorizētā artilērija, pēc tam tanki. Virs tribīnēm pārlidoja arī pāris desmitu vācu lidmašīnu. Šo ceremoniju nevar nosaukt par militāru parādi šī vārda pilnā nozīmē, kā to parasti dara uz PSRS ``atmaskošanu'' centrētie vēsturnieki. Faktiski tā bija nevis kopīga parāde, bet pilsētas nodošana padomju pusei, kad gar tribīni soļoja vienīgi vācu daļas. Tas bija vienīgais kopīgais PSRS un Vācijas tāda veida pasākums.

Citur padomju karaspēks nesastapa kaut cik nozīmīgu pretestību. Ir ziņas, ka no Polijas rietumiem uz austrumiem evakuējās dažādas Polijas ministrijas un citas iestādes. Arī policija saņēma pavēli izbraukt uz austrumiem un pāriet padomju iestāžu rīcībā. Sarkanās armijas ieņemtajā teritorijā policistus no Silēzijas arestēja tieši ešelonos un nosūtīja ``gūstā''.

Ģenerālis V.~Anderss liecināja, ka viņa vadītajā karaspēka grupā, kura mēģināja izlauzties caur padomju armijas daļām uz Rumānijas un Ungārijas pusi, ``27.~septembrī artilērija izšāva savas pēdējās granātas, šautenes~--- savas pēdējās patronas.'' Ģenerālis uzskatīja, ka Sarkanā armija, neļaujot poļu karaspēkam atkāpties uz Ungāriju un Rumāniju, līdz ar to neļāva izglābties no vācu un padomju gūsta vismaz 200 līdz 300~tūkstošiem cilvēku. Tiesa, Itālijas sūtnis Rumānijā saskatīja gluži pretēju ainu~--- Sarkanā armija dienvidu virzienā tīši virzījusies lēni, lai ļautu Polijas valdībai, daļai armijas atkāpties uz Rumāniju.

Kā rakstīja jau agrāk pieminētais rotmistrs J.~Kļimkovskis, pēc 17.~septembra Rumānijas virzienā atejošo poļu karavīru vidū ``nonāca vienīgi neskaidras ziņas, ka Sarkanā armija iegājusi Polijas teritorijā. Neviens precīzi nezināja, vai tā ir patiesība, un ja ir, tad ar kādu mērķi. Izplatījās visfantastiskākās baumas.'' Šādā situācijā poļu grupējumam ceļu aizšķērsoja vācu daļas. Pie Krasnobrodas (poļu \pltxti{Krasnobród}) pilsētas notika kauja, (viena no pēdējām kaujām Eiropā, kurā no abām pusēm tika izmantota kavalērija). Poļu grupa tika izkliedēta, daļa karavīru mēģināja nokļūt Ungārijā vai Rumānijā, bet lielākā daļa nolēma palikt Polijā. J.~Kļimkovska vadītā karavīru grupa pēc sadursmes ar vāciešiem naktī no 29. uz 30.~septembrī izklīda, pats viņš devās uz Ļvovu, pa ceļam sastopot jau Sarkanās armijas daļas, kuras poļu karavīriem ļāva mierīgi virzīties tālāk. Poļiem nācās vienīgi vairīties no satikšanās ar vietējiem ukraiņu iedzīvotājiem, kuri pret viņiem izturējās naidīgi. ``Tikai Sarkanās armijas klātbūtnei jāpateicas par to, ka šai laikā nenonāca līdz lieliem grautiņiem vai poļu masveidīgai izkaušanai''. V.~Anderss ar vairākiem virsniekiem vēl centās nokļūt līdz Ungārijas robežai, taču sadursmēs ar padomju daļām divreiz ievainots nonāca padomju gūstā, ārstējās padomju kara hospitālī. (Pēc J.~Kļimkovska teiktā V.~Andersa grupu aplenca vietējo ukraiņu vienība un tieši sadursmē ar to viņš tika ievainots, pēc tam lūdza padomju varas iestāžu palīdzību un devās padomju gūstā.) Oktobrī J.~Kļimkovskis pat vairākkārt apmeklēja V.~Andersu hospitālī. Kad rotmistrs paziņoja ģenerālim par savu vēlmi doties uz Parīzi pie ģenerāļa V.~Sikorska, V.~Anderss lūdza paziņot V.~Sikorskim par savu gatavību dienēt viņa pavēlniecībā, lamāja ``sanāciju'' un noslēgumā lūdza, lai viņš diplomātiskā ceļa V.~Andersu izdabū no Ļvovas un palīdz nonākt Parīzē, kas gan nebija iespējams.

Sarkanās armijas daļas virzījās uz rietumiem vidēji 30 līdz 40~km dienā. Tikai pirmajā tās Ģenerālštāba operatīvajā ziņojumos, publicētā 17.~septembrī bija runa par to, ka padomju daļas ar sajūsmu sveic vietējie iedzīvotāji, taču turpmākajos ziņojumos līdz 23.~septembrim tika runāts tikai par ieņemtajām pilsētām un apdzīvotajiem punktiem, gūstekņu skaitu. Tikai vēlāk plašu izplatību ieguva termins ``atbrīvošanas karagājiens'' (\rutxti{освободительный поход}). Bija jāsaprot, ka vietējie baltkrievu un ukraiņu iedzīvotāji tiek atbrīvoti no ``poļu panu'' kundzības.

Jāsaka, ka Padomju Savienības iedzīvotāju vairākums pozitīvi vērtēja ``baltkrievu un ukraiņu atbrīvošanu''. Rakstnieks K.~Simonovs liecinājis, ka brīdī, kad Vācija uzbruka Polijai, visas viņa un viņa darba biedru militārās avīzes redakcijā simpātijas bija Polijas puse, jo ``stiprākais uzbruka vājākajam un tāpēc, kaut pakts par neuzbrukšanu palika pakts, bet kurš gan no mums vēlējās fašistiskās Vācijas uzvaru\dots{}''. Taču gan pats rakstnieks, gan viņa biedri, atceroties padomju-poļu karu 1920.~gadā, Ukrainas un Baltkrievijas polonizāciju, PSRS un Polijas saspringtās attiecības starpkaru periodā, priecājās par Sarkanās armijas ieiešanu Rietumukrainā un Rietumbaltkrievijā. Var droši apgalvot, ka tā domāja daudzi PSRS iedzīvotāji, arī nebūdami Komunistiskās partijas biedri. Kā padomju propagandas ``sasniegums'' un padomju cilvēku naivums jāvērtē fakts, ka Sarkanajai armijai ieejot Polijā, tās karavīri jautāja vietējiem iedzīvotājiem: ``vai šeit ir pani?'' ``Pani'' bija šķiras ienaidnieki, no kuriem vietējie ukraiņi un baltkrievi viņiem bija jāatbrīvo.

Interesanti, ka vēlāk uz PSRS deportētā poliete S.~Cecinska atmiņās stāstīja, ka padomju izplatītajās skrejlapās bijis rakstīts: ``\rutxti{Polakam, panam, sobakam~--- sobachaya smert}'' (Poļiem, paniem, suņiem~--- suņa nāvi). Jādomā, ka atmiņu autore kļūdās un citāts ir vismaz neprecīzs. Padomju politdarbinieki varbūt arī lika vienlīdzības zīmi starp poļu paniem un suņiem, taču aicinājums nogalināt tos, pie tam nenodalot no darba cilvēkiem, nevis kā vispārinātu tēlu, bet kā konkrētas personas, bija pretrunā ar visiem Sarkanās armijas Politpārvaldes uzstādījumiem.)

Sarkanās armijas ienākšanu Rietumukrainā un Rietumbaltkrievijā \strong{vietējie iedzīvotāji} uzņēma dažādi. Lūk, kā par to rakstīja Baltkrievijas Komunistiskās partijas CK 1.sekretārs P.~Ponomarenko: ``Baltkrievu zemnieku noskaņojums lielisks, viņi cik var atbalsta Sarkano armiju \citespace{} aiz Belostokas [jau tīri poļu zemēs~--- V.Š.] iedzīvotāji sveicinās ar mums atturīgāk, mazāk zin krievu valodu, toties biežāk aiz stūra un no meža atskan šāvieni uz Sarkanās armijas karavīriem un komandieriem''. Arī septembrī~--- oktobrī tapušie padomju robežsargu komandieru ziņojumi IeTK vadībai liecina, ka daudzi vietējie ukraiņu un baltkrievu iedzīvotāji tiešām priecājās par Sarkanās armijas ienākšanu, ka poļu karavīru vairākums nevēlējās karot, labāk bēga prom no valsts vai padevās gūstā sarkanarmiešiem, taču poļu labējo partiju aktīvisti, daļa virsniecības, turīgo vietējo iedzīvotāju (ziņojumos~--- ``žandarmi un kulaki'') uzsāka aktīvu darbību (ziņojumos~--- ``teroru'') pret ukraiņu un baltkrievu tautības padomju piekritējiem. Bija gadījumi, kad atbildot uz to, Sarkanās armijas komandieri pavēlēja uz vietas nošaut ar ieročiem rokās sagūstītos poļu virsniekus, žandarmus, kolonistus. Tiesa, minēto pavēļu izdevēji pēc augstāko komandieru rīkojumiem vēlāk paši tika tiesāti un sodīti. Pēc vēsturnieku M.~Meļķuhova un V.~Sīpola datiem par marodierismu, gūstekņu nošaušanu bez tiesas un citiem noziegumiem tika notiesāti 49~Sarkanās armijas karavīri. Kara prokuratūra centās izvērst cīņu pret vardarbības izpausmēm, slepkavībām, marodierismu. Desmiti vainīgo tika sodīti ar cietumsodu un pat nošaušanu. Taču šī cīņa bija mazefektīva, jo pavēlniecība nepiešķīra spriedumiem pienākošos rezonansi. Drīzumā PSRS aizsardzības tautas komisārs K.~Vorošilovs un Ģenerālštāba priekšnieks B.~Šapošņikovs ierobežoja kara prokuratūras darbību, uzskatot par pietiekamu vainīgajiem piemērot administratīvos soda līdzekļus.

Kaut Polijas armijai lielu kauju ar Sarkano armiju praktiski nebija, pats tās ieiešanas fakts Polijas teritorijā, PSRS sadarbība ar Vāciju, pēc tam Rietumukrainas un Rietumbaltkrievijas atdalīšana no Polijas un pievienošana PSRS jau izsauca negatīvu poļu tautas reakciju. To nevarēja neietekmēt arī pirms tam Polijā vestā ilgstošā pretkrieviskā un pretboļševistiskā propaganda.

Mūsdienās poļu valodā publicēts liels daudzuma poļu liecību par Sarkanās armijas, PSRS Iekšlietu Tautas Komisariāta darbinieku veiktajām represijām pret poļu karavīriem, valsts iestāžu darbiniekiem, citiem poļu iedzīvotājiem. Domājams, šie materiāli pārsvarā atspoguļo gan no nosaukto padomju institūciju, gan no vietējo darbaļaužu rokām cietušo mantīgo poļu ciešanas. Diemžēl, ``šķiru cīņai'' piemita gaužām nepievilcīgas iezīmes. Mūsdienu poļu vēsturnieki ir tendēti vardarbībās aktus, Rietumukrainas un Rietumbaltkrievijas vietējo iedzīvotāju uzstāšanās gadījumus pret Polijas valsts iestādēm, poļu ierēdņiem, muižniekiem u.~c. 1939.~gada septembrī traktēt kā Padomju varas, tās Iekšlietu Tautas komisariāta, vietējo komunistu inspirētus, nevēloties saskatīt, ka bez tāda fona kā nabadzība un sociālā un nacionālā nevienlīdzība iedzīvotāji nebūtu viegli pamudināmi uz vardarbību.

Arī latviešu politiķis un publicists V.~Lācis apgalvo, ka ukraiņus Polijā vairāk uztrauca nevis poļu represijas, vērstas tieši pret viņiem, bet gan staļiniskās represijas pret Padomju Ukrainas iedzīvotājiem: ``Pēc drausmīgajiem saimnieciskajiem un politiskajiem terora aktiem pret ukraiņu tautu Padomju Ukrainā milzīgs ukraiņu vairākums Polijā, neraugoties uz poļu izturēšanos pret viņiem, tik un tā bija antikrieviski un antikomunistiski noskaņots daudz vairāk nekā antipoliski. Tāpēc pēc padomju karaspēka iebrukuma Rietumukrainā, tad, kad no padomju lidmašīnām tika nomestas skrejlapas ar vārdiem ``Tveriet izkaptis un cirvjus un apkaujiet poļu kungus!'', ukraiņi uz šiem saukļiem nereaģēja.'' Fakti gan runā pretī šādai mākslīgi konstruētai versijai. Apspiestie uz radušos izdevību ``nokārtot rēķinus'' ar apspiedējiem ``reaģētu'' arī bez īpašām padomju lidmašīnu kaisītām skrejlapām.

Vēsturnieki, kuri pētījuši notikumus 1939.~gadā Rietumukrainā un Rietumbaltkrievijā, uz vēsturisko avotu analīzes pamata secinājuši, ka Sarkanās armijas ieiešanu pārsvarā pozitīvi uztvēra ukraiņi, baltkrievi un ebreji, pārsvarā negatīvi~--- poļi. Grūti teikt, vai visumā turīgākie kā citu nāciju pārstāvji poļi ciestu vairāk vai mazāk no hitleriskā, nevis staļiniskā režīma nodibināšanas jau 1939., bet ne 1941.~gadā. Vācu kundzībai pastāvot, pret viņiem nebūtu represiju pēc sociāliem kritērijiem un, iespējams, tās būtu mazākas pēc politiskiem kritērijiem, taču vācu nacistu acīs viņi bija un palika zemākas rases pārstāvji, kuri nebija pelnījuši nekādas tiesības. Ukraiņi, baltkrievi, nemaz nerunājot par ebrejiem, kuri 1939.--1941.~gadā, tāpat kā poļi, arī cieta no staļiniskajām represijām, šajos gados tomēr nesa mazākus upurus, nekā vēlāk vācu okupācijas periodā.

Daudzās vietās vietējie ukraiņu un baltkrievu zemnieki, kā arī ebreji uzsāka stihisku izrēķināšanos ar poļu tautības zemes īpašniekiem, kolonistiem, robežsargiem, policistiem un pat mazākām karaspēka vienībām. Ē.~Jēkabsons atzinis, ka daļēji šīs akcijas var uzskatīt par iepriekšējo gadu nepārdomātās Polijas valdības nacionālās politikas sekām, kad dažādiem līdzekļiem tika mēģināts pārpoļot nacionālo minoritāšu pārstāvjus. Šī darba autors gan uzskata, ka poļu nacionālistu realizētā citu nāciju pārstāvju pārpoļošanas politika bija gana pārdomāta un mērķtiecīga, cerot ar spēku salauzt nepoļu pretestību. Polijas varas iestādes pašas vien bija vainīgas, ka liela daļa tai laikā Polijā dzīvojošo baltkrievu, ukraiņu, lietuviešu un ebreju ar prieku un cerībām sagaidīja Sarkano armiju. Viņu cerības nepiepildījās, bet daudzu no viņiem izturēšanās liecināja, ka neko īpaši labu šo nāciju pārstāvji nebija piedzīvojuši arī agrākās autoritārās Polijas valsts pastāvēšanas laikā.

Poļu liecinieki gan apgalvo, ka padomju propaganda tieši aicināja ``norēķināties par pagātni'', iznīcināt ``šķiras ienaidniekus'', laupīt muižas, sarīdīja tautības. Taču šīs pašas atmiņas liecina, ka īpaši aktivizējās ukraiņu nacionālisti, kuri ķērās pie poļu iznīcināšanas, ko pavadīja viduslaiku cienīgas spīdzināšanas. Tātad ``šķiriskā'' padomju propaganda diezin vai varēja izsaukt nacionālistiskus ekscesus, ja jau pirms kara ar nacionālistisku propagandu tiem nebūtu radīts noturīgs pamats cilvēku apziņā. Ja ekscesi rastos padomju propagandas rezultātā, tad to pamatā būtu šķirisks, nevis nacionāls naids.

22.~septembrī Sarkanās armijas Politpārvaldes priekšnieks Ļ.~Mehliss uz Maskavu ziņoja; ``Naids starp ukraiņiem un poļiem pieaug. Ukraiņi tagad ir aktivizējušies un atsevišķos apgabalos terorizē poļu zemniekus. Bija gadījumi, kad savstarpēji tika aizdedzināti ciemi, slepkavību un laupīšanu gadījumi. Dots rīkojums plaši izvērst darbu pret ukraiņu un poļu nacionālo naidu, gluži pretēji, apvienojot spēkus pret paniem-muižniekiem.''

Mūsdienās grupa poļu un ukraiņu vēsturnieku, veicot kopēju problēmu apspriešanu, ir nonākuši pie secinājuma, ka Rietumukraina un Rietumbaltkrievija 1939.~gada septembrī bija kaimiņu lielvalstu manipulāciju objekti. Viens no to uzdevumiem bija radīt Polijā nestabilitāti un nodrošināt tās sadalīšanu. Jāpiekrīt, ka Polijas kaimiņvalstu vadība pirms Otrā pasaules kara tiešām bija ieinteresēta situācijas Polijā destabilizācijā. Vācija jau ilgstoši sadarbojās ar ukraiņu nacionālistu organizācijām. PSRS vadība gan bija pati devusi triecienus vietējām komunistiskajām partijām, tomēr pret Polijas valsti vērstā un PSRS atbalstītā poļu, ukraiņu un baltkrievu komunistu ietekme, kaut mazinātā apjomā, bija palikusi. Rietumukrainas un Rietumbaltkrievijas teritorijā plašu ideoloģisku darbu 1939.~gada septembrī veica Sarkanās armijas politdarbinieki. Taču nedrīkst arī noklusēt, ka nestabilas situācijas radīšanā bija daudz darījuši, kā tas parādīts iepriekšējā nodaļā, gan poļu, gan ukraiņu nacionālistu vadoņi. Nacistiskā un padomiskā politiski-ideoloģiskā un politiskā darbība, protams, traucēja uz demokrātijas un nāciju vienlīdzības principiem balstītu tautu attiecību veidošanu Polijas ``austrumu kresos'', taču demokrātiju un vienlīdzību šeit tāpat necentās iedibināt nedz ``sanācijas'' režīms un poļu nacionālisti, nedz ukraiņu nacionālisti, kuri uzturēja ciešus kontaktus ar nacistiskās Vācijas struktūrām.

Kā atzīmējuši paši poļu vēsturnieki, lielā Polijas ziemeļ-austrumu daļā nebija ne tikai lielu armijas vienību, bet arī policijas, kura evakuējās kopā ar valsts administrāciju un pašvaldības iestādēm, kad Polijas valdība atstāja Varšavu. Apstākļos, kad Polijas valstiskās struktūras sabruka, bet padomju varas iestādes vēl nebija radītas, minētajās teritorijās nevarēja runāt par kādu sabiedrisko drošību. Notika ne tikai uzbrukumi Polijas armijas daļām, bet arī civiliedzīvotāju aplaupīšana, kā arī slepkavības. Tas turpinājās līdz tam brīdim, kad padomju varas iestādes aizpildīja administratīvās varas vakuumu.

Tā, poļu-ebreju izcelsmes amerikāņu pētnieka J.~Grosa 1988.~gadā izdotā darbā ``\entxti{Revolution from Abroad: Soviet Conquest of Poland's Western Ukraine and Western Belorussia}'' (Revolūcija no aizrobežas. Padomes sagrābj poļu Rietumukrainu un Rietumbaltkrieviju) izmantoti ap 20~tūkstošu V.~Andersa armijas karavīru atmiņas par padomju gados (1939--1942) pārdzīvoto. Uz to pamata J.~Gross apgalvo, ka, sākoties Sarkanās armijas virzībai Polijas teritorijā, tur jau pirms tās nonākšanas sākās slepkavības, piekaušanas un mantas iznīcināšana, kas turpinājās vēl vairākas dienas pēc tam. ``Ukraiņu un baltkrievu lauku apvidos bija dziļi iesakņojies etniskais naids. Tas aizpildīja asinīm vakuumu, kurš bija izveidojies poļu administrācijas kraha rezultātā. Ar poļiem, kuri sastādīja iedzīvotāju mazākumu un parasti apmetās lielās pilsētās un tikai nelielā skaitā bija pārstāvēti laukos, izrīkojās nežēlīgi. Tūkstoši tika nogalināti, bieži ar pārdomātu nežēlību.'' Pēc autora domām, minētais pētnieks, iespējams, neievēro to, ka PSRS pametušo V.~Andersa armijas karavīru atmiņas satur arī pašattaisnošanās motīvu par savu rīcību, atstājot grūtā brīdī (1942) PSRS, ko padomju puse faktiski traktēja kā nodevību. Šī pašattaisnošanās varēja ietvert arī pret poļiem vērstās ``nežēlības'' pārspīlēšanu.) Tomēr jāuzsver, ka J.~Gross, rakstot par poļu armijas atkāpšanos caur ``austrumu kresiem'', bija spiests konstatēt, ka ``savas pēdējās kaujas poļu armija savā teritorijā veda pret ukraiņiem, baltkrieviem, ebrejiem''.

Var atzīt, ka padomju aģenti un vietējie komunisti bija ieinteresēti radīt labvēlīgu Sarkanajai armijai un Padomju Savienībai noskaņojumu, demonstrēt ``darbaļaužu'' neapmierinātību ar ``panu'' Poliju, taču veiksmīgi to izdarīt viņi nevarētu, ja ne divi faktori:
\begin{enumerate*}[label=\arabic*)]
\item saspringtās starpnacionālās attiecības starp valdošo un turīgāko poļu iedzīvotāju slāni un nacionāli diskriminēto un sociāli apspiesto citu tautību vairākumu;
\item vietējo iedzīvotāju bažas par savu likteni, Polijas valstij sabrūkot.
\end{enumerate*}

Sarkano armiju ar prieku sagaidīja tikai kreisi noskaņotā baltkrievu un ukraiņu iedzīvotāju daļa, taču ukraiņu un baltkrievu vairākums labprātāk nonāca padomju republiku ar ukraiņu un baltkrievu vairākumu sastāvā, nevis hitleriskās Vācijas varā. Protams, nacionāli noskaņotā baltkrievu un ukraiņu daļa, kura par savu ideālu uzskatīja nacionālās suverenitātes iegūšanu, ar bažām raudzījās nākotnē, baidījās no jaunām represijām. Šos cilvēkus nevarēja iepriecināt Rietumbaltkrievijas un Rietumukrainas pievienošana padomju republikām. Ebreji bija gatavi padoties jebkurai antinacistiskai varai.

To, ka Polijas austrumu apgabalu pamatiedzīvotāju attieksme pret poļu mantīgajiem slāņiem nebija tikai boļševiku aģitācija izsaukta, bet tai bija dziļāki cēloņi, apliecināja arī Latvijas konsula Viļņā Feliksa Donasa ziņojums. Sarkanajai armijai ieejot lietuviešu apdzīvotajās Polijas teritorijās, ``zemnieki, pret kuriem muižnieki bija izturējušies slikti, izlaupīja muižas un noslepkavoja īpašniekus. Vairākos gadījumos garāmejošais lielinieku karaspēks izglāba muižniekus no nāves, apcietinot viņus un aizvedot, bet vēlāk atbrīvojot. Bija pat tādi gadījumi, ka karaspēks šāva uz laupītājiem. Daži zemnieki izdzina muižniekus no muižām, noslepkavoja viņus, bet atlikušo mantu un zemi sadalīja savā starpā. Ja attiecības bija bijušas labas, tad muižniekam atstāja dažas istabas, 1~zirgu, pāris govju, dažus hektārus zemes, izsniedzot katru dienu attiecīgu pārtikas devu (``\pltxti{pajok}''). Nodibinājās lauku komitejas, kuras regulēja lauku dzīvi.''

Šādā situācijā nebija brīnums, ka arī citādi pret boļševikiem simpātijas nebaudoši cilvēki izmantoja gadījumu, lai gan atriebtos agrākajiem apspiedējiem, vecajai, aizejošajai varai, gan arī katram gadījumam nodemonstrētu savu pozitīvo attieksmi pret jauno varu.

Juku periodā nacionālā ziņā galvenie cietēji, protams, bija poļi, bet galvenokārt tāpēc, ka līdz tam viņi piederēja priviliģētajiem slāņiem, bija zemes īpašnieki, kolonisti, policisti un dažāda ranga valsts ierēdņi. Līdz tam apspiestais iedzīvotāju vairākums viņiem tagad atriebās par agrākajiem esošajiem vai šķietamajiem pāridarījumiem.

Arī poļu vēsturnieks M.~Veržbickis, vairāku pēdējā desmitgadē tapušu darbu autors par poļu, ebreju un baltkrievu dzīvi padomju varas apstākļos 1939.--1941.~gadā, ir spiests atzīt, ka tieši baltkrievu un ebreju trūcīgāko slāņu pārstāvji 1939.~gada septembrī bija iesaistījušies pret poļu varas iestādēm virzītās darbībās.

Sarkanās armijas Politpārvaldes priekšnieks Ļ.~Mehliss ziņoja uz Maskavu, ka poļu virsnieki padodas padomju karavīriem divu iemeslu dēļ~--- lai nenonāktu vācu gūstā un lai izvairītos no baltkrievu un ukraiņu zemnieku izrēķināšanās ar tiem. Viņš minēja arī konkrētu gadījumu, kad poļu virsnieku grupa lūdza pastiprināt viņu konvoju no sarkanarmiešiem, baidoties no vietējo zemnieku uzbrukuma. Acīmredzot šādas situācijas bija iespējamas apgabalos, kur ``sanācijas'' režīma apstākļos bija notikušas zemnieku ``pacifikācjas'' akcijas. To apstiprina poļu virsnieka H.~Goržehovska stāstījums par to, kā septembrī sagūstītu viņu kopā ar citiem gūstekņiem padomju karavīri konvojēja uz nometni: ``Pēc tam mūs kājām dzina uz Rovno (ukr. \uktxti{Рівне}, poļu \pltxti{Równe}~--- pilsēta Ukrainā). Kā tagad atceros: kad mēs gājām caur pilsētu, daudzviet, pamatā uz ebreju bodītēm, karājās šauri sarkani karogi. Bija skaidri saskatāms, ka tie ir Polijas karogi, kuriem norauta augšdaļa. Ebrejietes un ukrainietes apšlakstīja mūs ar samazgām, kliegdamas: ``Beigas jūsu poļu valstij''''.

Poļu vēsturnieks M.~Kornats ir uzsvēris, ka no poļu nacionālās atmiņas viedokļa PSRS 1939.~gada septembrī parādīja sevi kā agresoru, piedalījās kārtējā Polijas sadalē, bet Polija kļuva par divu totalitāru režīmu upuri. (M.~Kornats faktiski liek vienlīdzības zīmi starp abiem~--- nacistisko un staļinisko režīmiem, sakot: ``Jā, Polijas pilsētas bombardēja vācu luftvafe, fašisti veica mierīgo iedzīvotāju masveida slepkavības, taču bija arī staļiniskā režīma, kurš valdīja Padomju Krievijā, noziegumi. Nedrīkst aizmirst par iedzīvotāju deportācijām uz austrumiem, par Katiņas apšaušanām\dots{}'') Vēsturnieks gan atzīst arī citus viedokļus, to, ka Polijas ebrejiem Padomju Savienība bija vienīgā cerība uz glābiņu no Vācijas nestā genocīda, ka Krievijas (patiesībā visas PSRS) iedzīvotājiem Lielais Tēvijas karš sākās 1941.~gada 22.~jūnijā.

Taču arī tad poļu vēsturnieka viedoklī ir vājas vietas. Viņš gan saka patiesību, bet ne visu patiesību. Runājot par poļu nacionālo atmiņu, viņš faktiski runā par vienas Polijas tautas daļas, droši vien lielākās, bet tomēr~--- tikai daļas, atmiņu, kuru mūsdienās gan ar dažādu līdzekļu palīdzību ``no augšas'' tiek mēģināts padarīt par visas nācijas atmiņu. Pieminētie ebreji (un arī baltkrievi, lietuvieši, ukraiņi) taču bija Polijas pilsoņi, par kuriem diemžēl vēl līdz mūsdienām poļu vidū ja ne valdošais, tad redzamākais ir uzskats kā par nepoļiem, ar kuriem poļi varēja faktiski nerēķināties 1939.~gadā un kuru atceres arī mūsdienās nevar veidot ``nacionālo atmiņu''. Tiek ``aizmirsti'' arī tie poļi, kuri, neraugoties uz staļiniskajām represijām, izvēlējās cīņas pret nacismu ceļu Sarkanās armijas rindās vai palīdzēja tai kaldināt uzvaru padomju aizmugurē; tie, kuri cīnījās kā pret ``sanācijas'', tā arī pret nacistisko režīmu un kuru cerības uz labāku dzīvi saistījās ar PSRS. Tādejādi it kā nevilšus tiek noņemta atbildība no autoritārā ``sanācijas'' režīma par sagrāvi karā. It kā pavisam dabīgi no poļu vēsturnieka mutes izskan apgalvojums, ka PSRS 1939.~gadā piedalījās Polijas dalīšanā, kaut PSRS kā 1939.~gadā, tā arī vēlāk atteicās no poļu apdzīvotajām teritorijām. Varētu saprast, ja šādi ``nacionālo atmiņu'' traktētu cilvēki, tālu stāvoši no zinātnes, taču no vēsturnieka ir jāprasa precīzāks un daudzšķautnaināks vēstures skaidrojums.

Ja bija iespēja, karu pret Vāciju faktiski jau zaudējušo Polijas armijas daļu karavīri bargi izrēķinājās ar ukraiņu un baltkrievu iedzīvotājiem, kuri bija sacēlušies pret Polijas varu. To parāda poļu vēsturnieka T.~Strzemboša dati par poļu 110.~kavalērijas brigādes ulānu rīcību, kad tie ``neņēma gūstā saķertos ar ieročiem rokās. Nošāva tos uz vietas.'' To varbūt vēl varētu izskaidrot ar kara apstākļiem, kad bēgot no Sarkanās armijas, nebija pat laika sasaukt kara tiesu, lai izšķirtu bruņoto civiliedzīvotāju likteni. Taču bija arī vēl briesmīgāks gadījums, ko jau vajadzētu klasificēt kā kara noziegumu.

Jau minētais vēsturnieks M.~Veržbickis visai lakoniski stāsta par t.s. ``Skidelas (\betxti{Скідаль}~--- pilsēta Baltkrievijā Grodņas rajonā) sacelšanās'' apspiešanu. Taču avīzē ``\rutxti{Советская Беларусь}'' (Padomju Baltkrievija) 2006.~gada 5.~augustā baltkrievu publicisti šos notikumus attēlojuši krietni plašāk, kas ļauj izvērtēt gan pašu sacelšanos, gan tās apspiešanu.

Lūk, šīs publikācijas saīsināts pārstāsts. Kad Skidelas (pilsētā pirms kara dzīvojis ap 7~000 iedzīvotāju, galvenokārt ebreji. Mazāko daļu~--- ap 40\% sastādījuši baltkrievi, poļi, krievi) apkārtnes ciemos iedzīvotāji uzdzirda radio vēsti par Sarkanās armijas karagājienu uz Rietubaltkrieviju, bijušie RBKP biedri, komjaunieši, to līdzjutēji devās uz Skidelu. Bija nolemts sagatavot karogus, organizēt cilvēkus, lai svinīgi sagaidītu Sarkano armiju. Sākās sadursmes ar policiju, kuru revolucionāri noskaņotie jaunieši centās atbruņot. Ieradās arī ap 20~bruņotu zemnieku. Bijušie RBKP biedri ātri organizēja revolucionāru komiteju un no ieņemtā maģistrāta balkona, kur tika izkārts sarkans karogs, pasludināja tur notiekoša mītiņa dalībniekiem varas pāreju tās rokās. Bruņotu zemnieku vienības ieņēma pastu, elektrostaciju, policijas iecirkni, blakus ciemā atbruņoja karaspēka vienību un policiju. Sadursmes laikā tika nogalināti vairāki poļu virsnieki un arī sacelšanās dalībnieki. Kad no Mastu (baltkr. \betxti{Масты}, \betxti{Мосты}) pilsētas sacelšanās apspiešanai pa dzelzceļu tika izsūtīts karaspēks un policija, sacēlušies izjauca sliežu ceļu, apstādināja vilcienu un atbruņoja karavīrus, kuri gandrīz nepretojās~--- tūlīt sāka izmest šautenes pa vagonu durvīm un logiem. Ieroči tika sadalīti apkārtējo ciemu zemniekiem. Taču kā jau diletanti, sacēlušies nebija nodrošinājuši pilsētas aizsardzību. 19.~septembrī no Grodņas ieradās un Skidelā iebruka, kā avīze raksta, ``fašistiskie rīkļu rāvēji'' (Padomju prokuratūras 1940.~gadā tapušā ziņojumā rakstīts par ``soda vienībām, kuru sastāvā bija poļu armijas karavīri un amatpersonas un žandarmu daļas''.) Viņiem pievienojušies arī vietējie poļu kolonisti. Sacēlušos pretestība tika salauzta un sākta izrēķināšanās. Bojā gāja 29~sacelšanās dalībnieki, kādam nogrieza ausis un izdūra acis, dažiem uz krūtīm izgrieza zvaigznes. Upuru vidū pārsvarā bija cilvēki ar slāviskiem uzvārdiem, bet bija arī pāris ebreju. 200~cilvēku tika piespiesti nogulties uz zemes. Soda ekspedīcijas dalībnieki pavēlēja to skūpstīt, ņirgājoties: ``Skūpstiet, žīdi, komunisti, poļu zemi, tā nekad nepiederēs jums''. Ēkas, kurās dzīvoja sacelšanās dalībnieki, tika nodedzinātas. 70~cilvēku tika apcietināti. Nākamajā dienā~--- 20.~septembrī Skidelā ienāca Sarkanā armija. 1940.~gada vasarā notika lietas izmeklēšana, tika arestēti 15~cilvēki, Skidelas poļu iedzīvotāji. 11~no viņiem, kuri pēc tiesas ieskatiem bija ņēmuši dalību sacelšanās apspiešanā, tika nošauti.

Jāsecina, ka Skidelā aktīvu darbību pret Polijas varu izvērsa galvenokārt kreisi noskaņotie mazākumtautību iedzīvotāji. Nesastopot nemaz vai sastopot tikai vāju pretsparu, viņi darbojās visai aktīvi. Iedzīvotāju vairākums, cik var spriest, izturējās pasīvi. Soda ekspedīcijas no Grodņas organizētāji vadījās no saviem parastajiem priekšstatiem~--- nepieļaut nekādu iedzīvotāju pašdarbību, vērstu pret Polijas iekārtu, nerēķinoties pat ar valsts sabrukumu, Sarkanās armijas tuvošanos. Tā vietā lai kaut kā mēģinātu pretoties Sarkanajai armijai, notika beztiesiska, šķiriska un nacionāla izrēķināšanās ar vietējiem kreisajiem, ko var raksturot kā pilsoņu kara uzplaiksnījumu. Zvērību veicēji bija pelnījuši sodu, taču diemžēl par 1940.~gadā notikušo padomju tiesas procesu nevar pateikt, cik tas ievēroja likumību. Iespējams, ka ``šķiriskā taisnīguma'' vārdā cieta arī nevainīgi cilvēki.

Vēl jāpiebilst, ka minētajā publikācijā nav pieminēti notikumi 22.~un 23.~septembrī, kad tās pašas Skidelās apkārtnē pēc baltkrievu publicista V.~Loseva datiem vietējie zemnieki 4~miestos izrēķinājās ar bijušajiem poļu leģionāriem-kolonistiem. Bojā gāja 42~cilvēki. Nav zināms, vai vainīgie arī tika kādi sodīti.

Dažviet Baltkrievijā RBKP bijušie biedri radīja kara revolucionārās komitejas, kuras atbruņoja policistus, uzņēmās svarīgu objektu apsardzi. Tika radītas pagaidu pārvaldes iestādes, tai skaitā zemnieku komitejas, strādnieku gvardes un zemnieku milicijas vienības. Kalpi, bezzemnieki un trūcīgie zemnieki sāka sadalīt muižnieku kustāmo un nekustāmo mantu. Tika sadalīti ap miljonu~ha zemes, 14~tūkstošu zirgu, 33,4~tūkstoši govju, 15,7~tūkstoši cūku.

Literatūrā sastopami apraksti, kā Polijas austrumu apgabalu iedzīvotāji sagaidīja Sarkano armiju ne tikai ar sarkaniem karogiem, bet arī ar svētbildēm. Tātad daļa iedzīvotāju tās atnākšanā pirmkārt redzēja savu atgriešanos Krievijā, atbrīvošanos no Polijas nacionālā un arī reliģiskā sloga. Viņus gaidīja smaga vilšanās, taču tā nebija viņu vaina, bet tas nebija arī Polijas nopelns.

Lūk, kā situāciju Rietumbaltkrievijas Belostokas apvidū pēc ziņas saņemšanas par Sarkanās armijas virzīšanos Polijas teritorijā aprakstīja jau agrāk pieminētais anonīmais krievu emigrants. ``Tai pat naktī sākās zemnieku nemieri. Apkārtnes muižnieki bēga. Zemnieki~--- starp viņiem daudzi poļu formā~--- apbruņoti šautenēm, uzstādīja sardzi. No rīta izveidojās kaut kas līdzīgs revolucionārai komitejai. Taču nekādu laupīšanu nebija. To nebija vispār mūsu rajonā. Vai nu boļševiki baidījās no tuvumā stāvošajiem vāciešiem un nesūtīja savus aģitatorus vai pārsvaru ņēma veselais saprāts zemniekos, taču visi ``ekscesi'' aprobežojās ar dažu policistu un voita (pagasta vecākā) arestu, pie tam voitu atlaida pret galvojumu, bet policistus pažēloja un deva tiem drīzumā aizbēgt.

Nākamajā dienā pienāca ziņas. Sarkanā armija izturas pilnībā miermīlīgi pret vietējiem iedzīvotājiem. Nekādu ālēšanos, visur kārtība, pat attieksme pret sagūstītajiem poļiem pilnībā pieklājīga. Taču vara pilsētās tiek nodota vissliktāko elementu rokās: [tie ir] politiskie noziedznieki, politiski neuzticamie poļu varas laikā, nikni jauni ebrejiņi~--- amatniekzeļļi.

Nekādu ārēju ālēšanos es nenovēroju un objektivitātes labad vēlreiz uzsveru, ka sarkanarmiešu daļas izturējās disciplinēti, un, acīmredzot, vietējos iedzīvotājus noskaņoja sev labvēlīgi.

Neapšaubāms pacēlums bija laukos: boļševiku solījumi, zemes dalīšana, daudzi sākotnēji populāri soļi veicināja zemnieku siržu piesaisti boļševikiem. Privāttirdzniecība propagandas nolūkos arī pagaidām netika pārtraukta. Tika vesta plaša propaganda [sakarā ar] Rietumbaltkrievijas Satversmes sapulces vēlēšanām.''

Tādejādi, vietējo iedzīvotāju attieksmi pret Sarkano armiju nevar vērtēt viennozīmīgi. Diemžēl, nenākas brīnīties, ka, ja līdz 1990.~gadam PSRS tika rakstīts tikai par draudzīgu sarkanarmiešu sveikšanu, tad pēc ``sociālisma'' sabrukuma gan Polijā, gan bijušajās padomju teritorijās pirmkārt tiek runāts un rakstīts par pēc tam sekojošajām represijām pret vietējiem iedzīvotājiem.

Pēc valdības bēgšanas no valsts Polijas armijas pretestība vērmahtam vēl turpinājās. Vācijas aviācijas ģenerālštāba 10.~septembrī izdotajā pavēlē bija norādīts, ka ``uzbrukums [pilsētām] ir jāuzskata kā atriebība par pret vācu karavīriem izdarītajiem noziegumiem''. Pēc 10.~septembra bombardēšana ievērojami pieauga. Šajās dienās Varšava kā pirmā no Eiropas lielpilsētām piedzīvoja gaisa uzbrukumu briesmas, kas prasīja simtiem upuru. 15.~septembrī Varšavu sāka apšaudīt arī vācu artilērija, koncentrējot ap pilsētu vairāk nekā 1~000 lielgabalu. 19.~septembrī iedzīvotāji tāpēc, ka pilsētas kapsētas bija pārslogotas, bija spiesti bombardēšanu un artilērijas apšaužu laikā kritušos apbedīt apstādījumos. Pēc četrām dienām pilnībā pārstāja darboties elektrības un ūdens apgāde.

Taču pilsētas aizstāvji nevēlējās padoties. Patrioti vēl joprojām cerēja uz Rietumvalstu palīdzību, pildīja virspavēlnieka maršala E.~Ridzs-Smiglija pavēli noturēt pilsētu par katru cenu. 22.~septembrī 8.~vācu armijā ieņēma izejas pozīcijas ap Varšavu un gatavojās triecienam. Tās pavēlnieks ģenerālis J.~Blaskovitcs bija norūpējies par iespējamo lielo upuru skaitu kājnieku daļās. Aviācijai bija jāpiespiež pilsētas aizstāvjus ātrāk padoties. 23.~septembrī sākās intensīva pilsētas bombardēšana. 25.~septembrī Ā.~Hitlers deva rīkojumu, ka politisku iemeslu dēļ pilsētu jāieņem nedēļas laikā. Šī diena pilsētas vēsturē iegāja kā ``pekles svētdiena''~--- virs Varšavas tika nomests 560~tonnu sprāgstvielu un 72~tonnas degbumbu. Degbumbas tika izmestas no transportlidmašīnas ``Ju-52''. Tā rezultātā dažas no tām skāra pat savas armijas daļas. Fugasa un degbumbas iznīcināja ne tikai militārus objektus, bet arī biezi apdzīvotus kvartālus.

Pilsētas prezidents (mērs) S.~Staržinskis, kurš 8.~septembrī tika iecelts arī par ``Varšavas aizstāvības civilo komisāru'' (\pltxti{Komisarz Cywilny Obrony Warszawy}), darīja visu iespējamo, lai stiprinātu varšaviešu cīņas garu. Slavena kļuva viņa radiouzruna pilsētniekiem aplenkuma laikā: ``Es gribētu, lai Varšava būtu diža \citespace{} Un šodien es redzu to dižu'' (\pltxti{Chciałem, by Warszawa była wielka \citespace{} I dziś widzję ja wielką}). Pēc kapitulācijas 8.~oktobrī vācieši viņu arestēja. Pēc kara beigām Polijā klīda dažādas versijas par viņa likteni. Visticamāk, viņš tika nogalināts turpat Varšavā. 1978.~gadā poļu režisors A.~Tross-Rastaveckis viņam veltīja filmu ``\pltxti{Gdziekolwiek jesteś, Pane Prezydencie}'' (Lai kur tu nebūtu, prezidenta kungs). 2003.~gadā S.~Staržinskis ar lielu balsu pārsvaru tika ievēlēts par ``gadsimta varšavieti''.

Par civiliedzīvotāju izturēšanos šais briesmu pilnajās dienās ir palikušas dažādas liecības.

Lūk, ko vietējā avīzē ``\pltxti{Kurier Warszawski}'' (``Varšavas kurjers'') rakstīja kāds amerikāņu kinodokumentālists J.~Brauns: ``Mēs visi amerikāņi, kuri esam sapulcējušies savā sūtniecībā, un kuriem ir iespēja sekot aplenktās galvaspilsētas dzīvei, nespējam atrast vārdus apbrīnai par Polijas civiliedzīvotāju izturību un varonību. Mums patīk jūsu miers, jūsu humors, jūsu neatlaidība. Varētu domāt, ka bez pārtraukumiem notiekošie vācu gaisa uzbrukumi, kas domāti šī gara salaušanai, novedīs pie pilsētas atdošanas. Taču pat tajos cilvēkos, kuri pelnījuši vislielāko līdzcietību, man nav gadījies atrast šādu nodomu. Gluži pretēji. Jo stiprāk jūs ienaidnieks vārdzina, jo nelokāmāka ir jūsu stāja''. Taču, lūk, ko rakstīja kāds anonīms autors 25.~septembrī: ``Bumbvedēji drāžas garām mūsu logiem, trīs degbumbas trāpa [kaimiņēku] jumtos \citespace{} Visā namā izbirst stikli. Tā vai citādi kāda bumba ķers arī mūsu māju, kāda starpība, notiks tas piecas minūtes agrāk vai vēlāk. Kaut nu taču tas notiktu! Kaut reiz tam būtu beigas!''

1939.~gada septembra aplenkuma laikā Varšavas civiliedzīvotāju vidū blakus visuresošas vardarbības pieredzei valdīja nenoteiktas bailes. Piemēram, vācu lidotāji gan neveica gāzu uzbrukumus pilsētai, taču bumbu detonāciju saceltie putekļu mākoņi un sprādzienus pavadošās nepazīstamās smaržas izsauca paniku. Daudzos cilvēkos valdīja bailes no indīgām vielām, ko piloti, saldumos vai desās noslēptas, it kā nometot virs pilsētas, vai kuras pašā pilsētā tiekot slepus izdalītas. Iedzīvotāji viens otru brīdināja, ka ar izpletņiem virs Varšavas izmesti vācu spiegi, kuri, iespējams, ir saindējuši akas. Nepazīstami cilvēki bieži tika noturēti par vāciešiem, tiem pieprasīja uzrādīt dokumentus. Īpašās aizdomās tika turēti Varšavā dzīvojošie evanģēlisko draudžu piederīgie. Visām šīm bailēm, aizdomām nebija pamata. Nav nekādu pierādījumu tam, ka Varšavā tiktu iesūtīti vācu spiegi vai diversanti. Un kāda gan jēga būtu to darīt, ja pilsēta tika bombardēta, gaidot drīzu tās padošanos.

Aplenkuma laikā dzīvību zaudēja ap 6~000 poļu karavīru un ap 25~000 civilistu. (Ir arī dati, ka Varšavas aizstāvēšanās cīņās 1939.~gada septembrī gāja bojā 45~000 cilvēku un tika izpostīti 15\% ēku.)

Pilsētas aizstāvji cieta ne tikai no bombardēšanas, artilērijas apšaudes, bet arī no ūdens un pārtikas trūkuma, draudēja izcelties sērgas. 27.~septembra rītā Varšavas kara apgabala pavēlnieks ģenerālis J.~Rummels piekrita pilsētas kapitulācijai, akts tika parakstīts 28.~septembrī.

1939.~gada 1.~oktobrī Varšavā iemaršēja ģenerāļa V.~f.~Reihenau komandētās 10.~vācu armijas karavīri. Pilsētas J.~Pilsudska laukumā notika militāra parāde. Dažas dienas vēlāk pilsētu apmeklēja Ā.~Hitlers, kurš to šķērsoja atklātā auto. Iedzīvotājiem, piedraudot ar ieročiem, bija aizliegts tuvoties ielām, kur notika izrāde. Piesardzības pasākumi nebija lieki, jo jau organizētie poļu pagrīdes spēki plānoja pilsētas centrā uzspridzināt Ā.~Hitleru. Sprādziens neizdevās, jo ielu nosprostošanas dēļ viens no atentāta gatavotājiem laikus netika paredzētajā vietā.

Jau pirms Varšavas kapitulācijas 1930.~gada 27.~septembrī pēc vieniem datiem~--- ar virspavēlnieka maršala E.~Ridzs-Smiglija, pēc citiem, ticamākiem,~--- ar ģenerāļa J.~Rummeļa pavēli tika radīts nelegāls ``Polijas uzvaras dienests'' (\pltxti{Słužba Zwycięstwu Polski}) ar stingru organizatorisku tīklu un ar ģenerāli M.~Karaševiču-Tokaževski priekšgalā. Tā bija pirmā poļu pagrīdes organizācija, radīta Otrā pasaule kara laikā. Tās galvenais mērķis bija pretestības organizācija vācu un padomju spēkiem, Polijas atbrīvošana, tās pirmskara robežu atjaunošana.

Varšava pārtrauca pretestību 29.~septembrī, 30.~septembrī kapitulēja Modlinas cietoksnis, 2.~oktobrī pārtrūka pretestība Helas pussalā. 5.~oktobrī kapitulēja pēdējais organizētās pretošanās centrs Kockā (\pltxti{Kock}~--- pilsēta Ļubļinas vojevodistē).

ASV oficiāli ieturēja neitralitāti karā. Taču tas netraucēja ASV armijas štāba priekšniekam ģenerālim Dž.~Maršallam apsveikt Vācijas sauszemes karaspēka Ģenerālštāba priekšnieku ģenerāli F.~Halderu ar uzvaru pār Poliju.

Karā \strong{abas puses cieta zaudējumus}, taču uzvarētājiem tie bija krietni mazāki.

Jau pēc Otrā pasaules kara veiktie pētījumi rādīja, ka 1939.~gadā kaujās ar vērmahtu bojā gāja ap 66~tūkstoši (ir arī citi dati, pēc vācu vērtējumiem~--- \strong{60~000 līdz 100~000}) poļu karavīru, pāri par 133~000 tika ievainoti. \strong{Vairāk nekā 400~000 karavīru, to vidū ap 18~000 (citi dati~--- ap 19~000) virsnieku, nonāca vācu gūstā.} (Ap 1~300 karavīru nokļuva slovāku gūstā.) Tūkstoši poļu gūstekņu: karavīru un virsnieku tika izvesti uz Vāciju, ievietoti karagūstekņu nometnēs. Attieksme pret viņiem nometnēs kopumā bija korekta. Tiesa, poļu virsnieki, kuriem izdevās no nometnes Desselē (\detxti{Dößel}) bēgt, un kuri gandrīz visi tika notverti Vācijas teritorijā, tika nogalināti. Taču tas bija viens no nedaudziem izņēmuma gadījumiem.

\strong{Ap 80~000--90~000 poļu karavīru izdevās izrauties uz} neitrālajām valstīm. Ap 30~tūkstošu karavīru atkāpās uz Rumāniju un 40~tūkstošu~--- uz Ungāriju, 14~000 uz Lietuvu, 1~600 uz Latviju. Viņi šajās valstīs tika internēti.

Karadarbībā gāja bojā apmēram 100~000 Polijas civiliedzīvotāju, liela daļa no tiem aviācijas uzlidojumos. \strong{Bez tam kā ``aizdomīgi elementi'' tika sagūstīti arī ap 200~000 civilistu.} Ap \strong{61~tūkstoti ebreju} uzreiz nodalīja no pārējiem gūstekņiem, ar viņiem apgājās daudz sliktāk nekā ar pārējiem.

Poļu vēsturnieki uzsver, ka Vācijai dzīvā spēka zaudējumi Polijā sastādīja ap 45~000 nogalinātu un ievainotu karavīru. Vācijas armija Polijā zaudējusi arī ap 1~000 tanku un bruņumašīnu (citi dati~--- ap 830 bruņu tehnikas vienību~--- apmēram 30\% no sastāva), 370~lielgabalu, vairāk nekā 10~000 transporta līdzekļu (ap 6~000 auto un 5~500 motociklu), vairāk nekā 700~lidmašīnu. (Taču ir arī citi poļu dati~--- kara triju nedēļu laikā poļu lidotāji esot notriekuši vismaz 126~vācu lidmašīnas, iespējams vēl 10, un bojājuši 14~lidaparātus. Paši poļi zaudējuši 70\% lidmašīnu un daudzus pilotus.) Daudzi vācu ieroči cietuši tādus bojājumus, ka tiem bija nepieciešams kapitālais remonts.

Z.~Zaluskis pat rakstīja, ka vācu zaudējumi septembra kampaņas laikā bija ne mazāki kā viņu zaudējumi nākamajos 20~mēnešos Norvēģijā, Francijā un Balkānos, līdz pat iebrukumam PSRS. Tas gan neatbilst patiesībai, Francijā vien vērmahts zaudēja vairāk nekā 45~tūkstošus kritušo un bez vēsts pazudušo.

Savukārt, vācu vēsturnieku ziņas rāda, ka vācu neatgriezeniskos zaudējumus nenākas pārspīlēt.

Pirmējie 1939.~gadā Vācijas pavēlniecības publiskotie dati par 10~572 kritušajiem, 30~322 ievainotajiem un 3~409 bez vēsts pazudušajiem, kas balstījās uz sanitāro iestāžu sniegtajām ziņām (un kurus poļu vēsturnieki pamatoti apšaubīja kā samazinātus), atšķīrās no vēlāk armijas iestāžu veiktās uzskaites. Tā, XX gadsimta 50.~gados uz Vācijas Ģenerālštāba materiālu pamata tika noteikts, ka krituši bija 16~343 vācu karavīri. Atšķirība no iepriekšējiem datiem acīmredzot bija izskaidrojama ar to, ka bez vēsts pazudušie tika pieskaitīti kritušajiem, un bija mirusi daļa ievainoto. Pārējie lielākoties atgriezās ierindā. Atšķiras arī vācu un poļu dati par vācu tehnikas zaudējumiem. Tā, vācu vēsturnieki uzsver, ka transportlīdzekļu un bruņu tehnikas (674~tanku) lielos zaudējumus noteica ``nodilums un neizbraucamie poļu ceļi''. Pārspīlēti bija arī poļu dati par vācu aviācijas zaudējumiem. Pēc vieniem datiem Vācijas aviācija Polijā zaudēja 285 lidmašīnas (15\% no tās sastāva), tai skaitā 109~iznīcinātājus un pikējošos bumbvedējus, pēc citiem~--- ap 600~lidmašīnu. Vēl pēc citiem vācu datiem poļu aviācija un zenītartilērija kopā notrieca ap 250~vācu lidmašīnu, bet pie tam poļu gaisa spēki zaudēja gandrīz visus iznīcinātājus un ap 80\% bumbvedēju).

Padomju ``atbrīvošanas karagājiena'' gaitā krita ap 3,5~tūkstošu poļu militāristu un civilpersonu, ap 20~tūkstošu tika ievainoti vai pazuda bez vēsts. (Pēc poļu vēsturnieka A.~Friškes datiem Sarkanās armijas darbības rezultātā nogalināti tika ap 2~500 poļu karavīru un policistu, kā arī vairāki simti civilistu. Tiesa, cik no militāristiem un civilistiem zaudēja dzīvību sarkanarmiešu, cik~--- sacēlušos vietējo iedzīvotāju darbības rezultātā, vēsturnieks nedod. Pēc Ē.~Jēkabsona sniegtajām ziņām padomju-poļu sadursmēs krita apmēram 6~000--7~000 poļu karavīru un civiliedzīvotāju, bet 10~000 tika ievainoti.)

Padomju puse oficiāli pasludināja, ka krituši bija 737 un ievainoti 1~862~cilvēki. Mūsdienās tiek atzīts, ka krita 1~173~cilvēki, 2~002~tika ievainoti, 302~pazuda bez vēsts. Pie tam liela zaudējumu daļa tika ciesta nedisciplinētības un neorganizētības dēļ. (Poļu puse šos skaitļus uzskata par samazinātiem, kritušo esot bijis vairāk~--- ap 2,5~tūkstošiem, kā arī vismaz 4~tūkstoši, pat 8~000--10~000~ievainoto).

Dati par Sarkanās armijas saņemto poļu gūstekņu skaitu dažādos izdevumos ievērojami atšķiras. Tā, Ukrainas un Baltkrievijas frontes atskaitījās par it kā sagūstītiem ap 454~700 cilvēku (1939.~gada 17.~septembrī~--- 2.~oktobrī Ukrainas fronte gūstā saņēma 394~498 poļu karavīrus, tai skaitā 16~723~virsniekus, Baltkrievijas fronte~--- 17.--30.~septembrī 60~202 cilvēkus, tai skaitā 2~066 virsniekus), taču šai skaitā tika iekļauti visi, kas nēsāja formas tērpu~--- policisti, dzelzceļnieki, skauti, ugunsdzēsēji, pastnieki utt. Poļu vēsturnieks J.~Zavodnijs minējis skaitli 230~670, pārējie pievienoti arestu rezultātā PSRS sagrābtajās teritorijās. PSRS IeTK Konvoja karaspēka materiāli rāda 226~391 cilvēku.

Bez tam lielā mērā neskaidrības izsauc tas, ka dažas dienas septembrī Sarkanās armijas daļas vietējos Rietumukrainas un Rietumbaltkrievijas iedzīvotājus, iesauktus Polijas armijā, vienkārši atlaida mājup, bet vēlāk atkal ``gūstīja''.

Var teikt, ka gūstā Sarkanā armija saņēma ap 240--250~tūkstošus poļu karavīru. Arī V.~Molotovs PSRS Augstākās padomes sesijā 31.~oktobrī, runājot par saņemtajiem gūstekņiem, nosauca skaitli ``apmēram 250~tūkstošu''. No 1939.~gada 25.~septembra~--- līdz 7.~oktobrim IeTK nometnēs, kuras tika organizētas no 22.~septembra, tika ieslodzīti 125~803~cilvēki. (Karagūstekņu un internēto atrašanās Iekšlietu ministrijas sistēmas pārraudzībā bija tai laikā parasta lieta. Tā bija Lielbritānijā, Ungārijā, Itālijā.) Nespējot tādu skaitu cilvēku apgādāt ar pārtiku, telpām, pat dzeramo ūdeni, steidzami bija jālemj to liktenis.

VK(b)P CK Politbirojs jautājuma izskatīšanai radīja komisiju A.~Ždanova vadībā. Tajā ietilpa arī L.~Berija un Ļ.~Mehliss. 2.~oktobrī komisija atbalstīja VK(b)P CK Politbiroja lēmuma projektu, kurš bija sagatavots uz L.~Berijas un Ļ.~Mehlisa priekšlikumu pamata. Projekts paredzēja Rietumbaltkrievijā un Rietumukrainā dzimušos gūstekņus pēc filtrācijas atlaist mājup. Karavīri, dzimuši ``Polijas vācu daļā'' bija jātur IeTK nometnēs līdz sarunu ar Vāciju noslēgumam par to nosūtīšanu mājup. Virsniekus, policistus, citus spēka struktūru darbiniekus bija jāizvieto speciālās nometnēs. Uz VK(b)P CK un PSRS TKP 3.~oktobra lēmuma pamata atbilstošu pavēli 4.~oktobrī izdeva K.~Vorošilovs un B.~Šapošņikovs. Rezultātā no nometnēm un gūstekņu savākšanas punktiem no 7. līdz 18.~oktobrim tika atbrīvoti līdz 43~tūkstošiem karavīru un apakšvirsnieku. Daļa sagūstīto (25~tūkstoši) tika aizturēta ceļu celtniecībai un darbam šahtās. (Atbilstoši 1929.~gadā parakstītās Ženēvas karagūstekņu aizsardzības konvencijai un PSRS pieņemtajam ``Nolikumam par karagūstekņiem'' gūstekņus varēja nodarbināt.)

Rietumukrainas un Rietumbaltkrievijas iedzīvotājus varēja noformēt par brīva līguma strādniekiem. Taču tendence neguva attīstību. Noslēdzot līgumu, bija jāpiešķir apdzīvojamā platība, kuras uzņēmumiem praktiski nebija. Jāņem arī vērā, ka bijušajiem Polijas un reizē jaunajiem PSRS pilsoņiem nebija profesionālo darba iemaņu un arī gribas strādāt. Daļa no tiem tieši atteicās strādāt, masveidā (vairāk nekā 1,4~tūkstoši cilvēku) bēga. Padomju iestādes saprata, ka viņu vardarbīga aizturēšana darbos rada pamatu pretpadomju aģitācijai, centās tos nomainīt ar bijušajiem policistiem no tiem pašiem apgabaliem. 1939.~gada beigās L.~Berija izdeva pavēli noskaidrot darba neefektivitātes, masveida atteikšanās no darba cēloņus. Nemieru iniciatori bija jānodod tiesai. Bija jāsoda arī padomju ierēdņi, kuri nenodrošināja nodarbināto pienācīgu apgādi. Tā padomju iestādēm šie cilvēki ne tikai kalpoja kā papildus darba spēks, bet radīja arī papildus rūpes.

Oktobra vidū PSRS un Vācijas valdības nolēma veikt karagūstekņu apmaiņu. Tika izveidoti apmaiņas punkti: viens Baltkrievijā un divi Ukrainā. Pavisam uz etniskās Polijas teritoriju tika repatriēti 42~492 cilvēki, kuri bija izteikuši vēlmi turp doties. Taču daudzi atteicās braukt mājup, lai nenokļūtu nacistisko okupantu varā.

Daļa poļu teritorijās dzīvojušo Polijas armijas karavīru, visbiežāk komunisti un ebreji, baidoties no nacistu vajāšanām, griezās pie padomju nometņu administrācijas ar lūgumu atstāt viņus PSRS. Daži lūdza piešķirt viņiem PSRS pilsonību. Taču Maskavā šos lūgumus uzņēma nelabvēlīgi. Nometņu vadībai tika uzdots izskaidrot gūstekņiem, ka viņiem jāatgriežas dzīves vietās un padomju iestādēm tie jānosūta turp.

Vēl 562~cilvēki vācu iestādēm pēc to pieprasījuma tika nodoti 1940.--1941.~gadā. Savukārt Vācija no 1939.~gada oktobra līdz 1941.~gada pavasarim padomju pusei nodeva 13~757~poļu karavīru, kuri visi tika atlaisti mājup. Acīmredzot tie bija ierindnieki vai zemākie komandieri.

Ja 1939.~gada septembra beigās IetK rokās bija ap 130~tūkstošu poļu gūstekņu, tad 19.~novembrī IeTK nometnēs palika 39,5~tūkstoši cilvēku.

Poļu vēsturnieks L.~Jerževskis īstenībai neatbilstoši apgalvoja, ka 180~tūkstošu gūstekņu esot deportēti uz PSRS. Jāsaka, ka arī mūsdienās literatūrā atrodami dažādi skaitļi. Tā, pēc vēsturnieces I.~Rosovskas-Jakubčikas ziņām karagūstekņu nometnēs tika paturēti ap 59~tūkstošu poļu.

Neskaidrs ir jautājums par Sarkanās armijas sagūstīto poļu karavīru statusu, jo to noteica nevis likums, bet 1939.~gada 20.~septembrī pie PSRS Tautas komisāru padomes izveidotās Ekonomiskās padomes pieņemts ``Nolikums par karagūstekņiem'' (\rutxti{Положение о военнопленных}). Ar Iekšlietu tautas komisāra L.~Berijas rīkojumu 19.~septembrī bija radīta PSRS IeTK Pārvalde karagūstekņu lietās (\rutxti{Управление по делам военнопленных НКВД СССP}). Pastāvošā situācija pret Poliju tika noteikta kā ``kara stāvoklis'', tātad pēc šīs pieejas saņemtie gūstā Polijas pilsoņi bija jākvalificē kā karagūstekņi. Tāpēc krievu vēsturnieks A.~Proņins, uzskatot, ka 1939.~gada septembrī starp Poliju un PSRS pastāvēja kara stāvoklis, no tā secina, ka PSRS rokās nonākušie Polijas karavīri bija jāatzīst par karagūstekņiem, attiecībā pret kuriem bija jāievēro visi šim gadījumam paredzētie starptautiskie noteikumi, kas gan netika darīts. Taču jāuzsver, ka no PSRS, tāpat kā no Polijas, puses oficiāli kara stāvoklis tā arī netika pasludināts.

Pēc poļu vēsturnieka Č.~Madajčika domām karadarbības bez kara pieteikšanas rezultātā padomju gūstā saņemtie poļu virsnieki ieņēma vidus stāvokli starp gūstekņiem un ieslodzītajiem, kaut padomju dokumentos viņi tika kvalificēti kā karagūstekņi. Viņi uzreiz tika nodoti PSRS IeTK aizgādībā. Uz Starptautiskā Sarkanā Krusta organizāciju netika nosūtīti gūstekņu saraksti, noraidīta tika arī virsnieku prasība izturēties pret viņiem atbilstoši starptautiskajām konvencijām. Tikai kad sakarā ar atsevišķiem saslimšanas gadījumiem ar tīfu tika veikta profilaktiskā potēšana, dažiem slimiem virsniekiem no Kozeļskas nometnes atļāva ārstēties Smoļenskas slimnīcās.

Savukārt krievu publicists J.~Muhins uzskata, ka Sarkanas armijas aizturētie un atbruņotie poļu karavīri jākvalificē kā internētie. Pamatā šim secinājumam gluži vienkārši argumenti: savstarpēji karoja taču Vācija ar Poliju, bet PSRS nekaroja, kara stāvoklis ne no PSRS, ne Polijas puses nebija pasludināts. PSRS tikai atguva agrāk Krievijai atņemtās teritorijas un ieviesa tajās kārtību. Padomju gūstā nokļuvušajiem poļu karavīriem juridiski nedrīkstēja piemērot karagūstekņu statusu, PSRS kā formāli neitrāla valsts viņus kā karojošas valsts pilsoņus varēja vienīgi internēt. (no lat. \latxti{internus}~--- iekšējs, šeit~--- ierobežot karojošo pušu karavīru pilsoņu brīvības līdz kara beigām).

Tiesa, arī šim secinājumam var iebilst. Lai arī juridiski kara starp Poliju un PSRS nebija, taču kara darbība tika vesta. Ja poļu karavīri atstātu savas valsts teritoriju, viņus varētu internēt, bet viņi taču to nedarīja. Sarkanās armijas komandieri atskaitījās par saņemto gūstekņu, nevis internēto skaitu.

Diemžēl, pastāvošais PSRS režīms juridiskus ``smalkumus'' faktiski neievēroja. Sagūstīto poļu karavīru stāvoklis, apiešanās ar viņiem nemainītos no nosaukuma maiņas. Internēti poļu karavīri tika neitrālajās valstīs (Rumānijā u.c.), bet PSRS pret viņiem faktiski izturējās kā pret gūstekņiem. (Literatūrā sastopamas ziņas, ka Padomju iestādes tikai 1941.~gada jūlijā sagūstītos poļu karavīrus sāka saukt par internētajiem. Krievu pētniece V.~Parsadanova gan apgalvo, ka jēdziens ``internētie'' padomju dokumentos tika lietots no 1939.~gada decembra. Cits pētnieks~--- V.~Šveds raksta, ka internētie IeTK nometnēs parādījās tikai pēc Baltijas valstu aneksijas 1940.~gada vasarā).

Bija gadījumi, kad poļu karavīrus, kuri ilgstoši izrādīja bruņotu pretestību Sarkanajai armijai, tiesāja padomju kara tribunāli. Kādā no spriedumiem, ar kuru tika notiesāti uz nāvi 22 poļu karavīri, bija teikts: ``Bandas mērķis bija sagraut uzticību Sarkanās armijas spēkam un bijušās Rietumu Baltkrievijas padomju teritorijā nodibināt iepriekšējo fašistisko tiesību iekārtu''. Tātad par pretošanos draudēja nevis gūsts, bet nāves sods. Tā padomju iestādes aizturētos kara gūstekņus ``likumīgi'' varēja nodot kara tribunāliem. Apgalvojums, ka Polijas valsts vairs neeksistē, formāli ļāva gūstā saņemtos poļu karavīrus uzskatīt nevis par karagūstekņiem, bet par kriminālnoziedzniekiem, kuri pretojās Sarkanajai armijai, tai ieviešot kārtību Rietumukrainā un Rietumbaltkrievijā. (Jāatzīst, ka, pat ja Polijas valdība būtu līdz pilnīgai valsts okupācijai atradusies Polijas teritorijā, ja tā būtu pieteikusi PSRS karu, padomju gūstā saņemto poļu karavīru stāvoklī maz kas mainītos. PSRS staļiniskā vadība maz, brīžiem pat nemaz, nerēķinājās ar starptautiskajām tiesiskajām normām.)

1939.~gada septembra notikumos iesaistītās puses ir devušas dažādu \strong{kara kampaņas vērtējumu}. Pēc poļu vēsturnieku vērtējumiem Polijā vācu vērmahts esot sastapis stipru un tā virspavēlniecībai negaidītu pretsparu. Daļēji tā bija taisnība. Par to var spriest pēc tā, ka Vācijas vadītāju jau 1939.~gadā plānotais uzbrukums Francijai šai gadā tā arī nevarēja notikt. Taču Vācijas bruņoto spēku Virspavēlniecības (\detxti{Oberkommando der Wehrmacht~--- OKW}) operatīvās pārvaldes priekšnieks ģenerālis A.~Jodls \strong{par Polijas izrādīto pretestību} paziņoja: ``Pretinieks, kurš pats par sevi bija visai ievērojams, sabruka negaidīti ātri''.

Poļu autori arī uzsver, ka Polijas pretošanās izrādījusies tikai par nedēļu īsāka nekā Francijas cīņa pret Vāciju 1940.~gadā, kaut Francijas un Lielbritānijas spēki bija daudz spēcīgāki gan skaita, gan bruņojuma ziņā. Tā, Polijas Tautas Republikas laikā vēsturnieks Z.~Zaluskis norādīja, ka Polijā pretošanās ienaidniekam ilga 35~dienas (tiesa gan, cīņas iznākums bija izšķirts 2--3 nedēļās~--- V.Š.), bet 1940.~gadā Rietumu frontē tā formāli ilga 46~dienas, taču faktiski tika pārtraukta jau pēc 38~dienām~--- 17.~jūnijā pēc maršala A.~Petēna paziņojuma par kapitulāciju. Septembrī vācieši līdz Varšavai nonāca 7~dienās, taču Varšava aizstāvējās vēl 21~dienu, turpretī Parīze vispār nepretojās.

Patiešām, spēku salīdzinājums it kā ļauj izdarīt secinājumu par izcili sīkstu poļu pretošanos. (Pielīdzinot divas brigādes vienai divīzijai, var teikt, ka Polijas armijā bija 48~divīzijas. Rietumu frontē 1940.~gadā karoja 107~franču divīzijas (neierēķinot palikušās Āfrikā), 14~angļu, 22~beļģu, 12~holandiešu, 2~poļu~--- pavisam 157~divīzijas. Polijā karoja 64~vācu divīzijas, Rietumos~--- 143. 1940.~gadā Rietumu sabiedroto pusē bija pārspēks tankos, kad Polijai nācās karot pret apmēram 15~kārtīgu vācu tanku pārspēku. Pēc zaudējumiem pret Poliju Vācijas mehanizētie spēki būtiski nepieauga. Pret Poliju Vācija varēja izmantot 34~tanku bataljonus (3~195 tankus), pret Rietumiem~--- 35 (3~385 tankus). Vācijas rīcībā 1939.~gadā karā pret Poliju bija 2~433 kaujas spējīgas lidmašīnas, 1940.~gadā uzbrukuma Rietumos laikā~--- 2~509.) It kā iznāk, ka Polijas armija karoja sekmīgāk nekā Rietumu sabiedrotie. Taču uzreiz jāsaka, ka no teiktā neizriet, ka franči u.c. tautu pārstāvji būtu cīnījušies sliktāk nekā poļi. To, ka Francija u.c. valstis ātri kapitulēja, var izskaidrot ar dažādiem objektīviem un subjektīviem cēloņiem, bet tas vairs neietilpst šī darba uzdevumos.

% page 849


Poļu vēsturnieku uzsāktajam Polijas un Francijas pretošanās salīdzinājumam var atrast arī citas puses. Piemēram, var salīdzināt Francijas un Polijas virsnieku zaudējumus un izdarīt secinājumus, kā darījis baltkrievu publicists un vēsturnieks A.~Usovskis. Parasti virsnieku zaudējumi pret ierindniekiem kara sākumā ir diezgan lieli, taču pēc tam, kad armija gūst pieredzi, šī attiecība samazinās un veido no 1:10 līdz 1:12. Francijā 1940.~gada kaujās virsnieku krita ļoti daudz~--- 10~virsnieku pret 23~ierindniekiem~--- fronte bruka un virsnieki ar personīgo paraugu centās vest kareivjus kaujā. 1939.~gada septembra kampaņā Polijā pēc A.~Usovska datiem uz vienu kritušo poļu virsnieku iznāca 32~ierindnieki. Poļu virsnieku salīdzinājumā ar ierindniekiem krita 12~reizes mazāk kā franču virsnieku. Ja armija sīvi cīnās aizsardzībā vai dodas pretuzbrukumos, virsnieku zaudējumiem jābūt lieliem, jo tie uztur savu padoto kaujas spējas arī ar personīgo piemēru. Ja virsnieku zaudējumi ir nelieli, tad armija ir vai nu augsti profesionāla un tās komandējošajam sastāvam nav vajadzības tīši riskēt ar dzīvību, vai arī tas bēg, pirms vēl sākuši bēgt ierindnieki. Polijas armiju 1939.~gada septembrī vērtēt kā augstu profesionālu nekādi nevarēja, kas liek domāt, ka poļu virsnieki pirmie bēga no kaujas lauka un tāpēc to kritušo īpatsvars bija tik zems.

1939.~gada maijā Polijas armijā dienēja ap 18,5~tūkstoši virsnieku, tai skaitā flotē~--- 358, aviācijā~--- 1~164, robežsardzē~--- 846. Bez kadru virsniekiem bija arī ap 60~tūkstošu rezerves virsnieku vecumā līdz 50~gadiem, ap 12~tūkstošu zemessardzes virsnieku līdz 60~gadu vecumam. Pēc poļu vēsturnieku vērtējumiem kara laikā iesaukti bija ap 40~tūkstošu virsnieku. Pēc aptuvenām aplēsēm vācieši saņēma gūstā ap 18~tūkstošu virsnieku, ap 10~tūkstošiem nonāca PSRS, ārzemēs nokļuva arī ap 10~tūkstošu. Ungārijā vien nonāca 12~ģenerāļu, vairāk nekā 6~000 virsnieku un jaunāko komandieru. Rumānijā patvērumu atrada ap 5~000 poļu virsnieku. Pat vēl pēc gada tajā mita 11~poļu ģenerāļu. Arī Latvijā pēc 1939.~gada septembra izrādījās 176~poļu virsnieki, taču līdz 1940.~gada vidum šeit esošo poļu karavīru un policistu skaits ievērojami pieauga.

Daļa no virsniekiem no Rumānijas un Ungārijas turpināja ceļu uz Rietumiem caur Dienvidslāviju un Itāliju vai pa jūras ceļu. Ungārijā vēl 1943.~gada vidū palika ap 1~500 poļu virsnieku. Rumānijā līdz 1941.~gada sākumam palika tikai ap 1~500 poļu virsnieku, 1941.~gada februārī Rumānija tos izdeva Vācijai. 1940.~gada jūnijā Polijas armijā Francijā dienēja ap 6~000 virsnieku, ap 600 bija nonākuši Tuvējos Austrumos. Pēc Francijas sagrāves 1940.~gada vasarā no tās evakuējās ap 7~300 virsnieku, ap 2~000 poļu nokļuva gūstā vai palika Francijā, vairāki simti tika internēti Šveicē. Vairāk nekā 6~000 poļu virsnieku nonāca Lielbritānijā. Poļu vēsturnieks Č.~Madajčiks pieņēma, ka Otro pasaules karu pārdzīvoja 10 līdz 11 tūkstošu (tātad~--- lielākā daļa) pirmskara kadru virsnieku.

Pēc Polijas Kara zaudējumu biroja (\pltxti{Biuro Odszkodowań Wojennyc}, 1945--1947) pētījumiem cīņās pret vērmahtu 1939.~gada septembrī pavisam krita 1~967 poļu virsnieku. Pēc 2005.~gadā publicētiem poļu vēsturnieku Č.~Gržeļaka un H.~Stančika ziņām gan esot krituši 3~300 virsnieku, taču šie dati varētu būt arī pārspīlēti. Toties vācu un padomju gūstā viņu nonāca 28 tūkstoši. Kā raksta J.~Muhins: ``Vācieši nogalināja vienu poļu virsnieku un 14 uzreiz padevās gūstā.'' (Gūstā gan daļa no viņiem padevās ne ``uzreiz'', bet Polijai pārtraucot pretošanos.) Gūstā saņemto virsnieku skaits tiešām bija liels, īpaši ievērojot, ka Polijas armija visu starpkara periodu tika gatavota karam, baudīja sabiedrībā augstu prestižu. (Salīdzinājumam, pēc 1944.~gada datiem Vācijas armija Polijas kampaņas laikā zaudēja 819 kritušus virsniekus un 15~450 karavīrus; tātad to attiecība bija 1:20, kas jāvērtē kā augsti kvalificētas armijas cienīgs rādītājs.)

Atsevišķi var pateikt par ģenerāļiem. Viņu līdz karam Polijas armijā bija 97 (vēl arī 3 kontradmirāļi). Krita no viņiem 5. (Tiesa, ir arī dati, ka no 49~poļu ģenerāļiem, tai skaitā jau atvaļinātiem, 1939.--1940.~gadu pārdzīvoja tikai desmit. Acīmredzot, šī uzskaite ir nepilnīga.) Vācu gūstā nonāca 33 ģenerāļi un 1 kontradmirālis. 24 tika internēti Rumānijā. 18 ģenerāļu un 1 kontradmirālis nonāca padomju gūstā. Pēc kāda poļu vēsturnieku aprēķina PSRS nometņu, cietumu, izsūtījuma režīmu izjuta 62~poļu ģenerāļi, no viņiem 51~gāja bojā. Iespējams, ka runa bija arī par jau pensionētiem ģenerāļiem.

Par ģenerāli J.~Oļšinu-Viļčinski poļu literatūrā ir dati, ka viņu noziedzīgā kārtā nošāva padomju karavīri. Turpretī krievu literatūrā stāstīts, ka īsi pirms ģenerāļa sagūstīšanas notika viņa komandēto armijas daļu sadursme ar Sarkanās armijas karavīriem. Poļi nogalināja 11 un ievainoja 14 sarkanarmiešus, aizdedzināja 4~padomju vieglos tankus. No viena tanka izlēca 3 cilvēku ekipāža un padevās poļiem, kuri tos nošāva. Par šo faktu poļu vēsturnieki klusē. Tā kā komandieris nes atbildību par savu padoto rīcību, tad kopā ar sievu un adjutantu uz Lietuvu bēgošā ģenerāļa automašīnai nonākot sarkanarmiešu rokās, viņi ģenerāļa sievu atbrīvoja, bet pašu ģenerāli un tā adjutantu pēc vienības komisāra pavēles nošāva. Ja viss patiešām notika tā, kā aprakstīts, tad acīmredzot būtu jāizmeklē abi incidenti un jāsoda abu pušu vainīgie par gūstekņu nogalināšanu kā par kara noziegumiem.

Z.~Zaluskis rakstīja, ka kara gaitā bija gan ceļš no Varšavas uz Rumānijas robežu, pilns ar ģenerāļu un pulkvežu mašīnām, piekrautām ar čemodāniem, gan arī lielākie kara pirmajā posmā (līdz 1941.~gada 22.~jūnijam) salīdzinājumā ar citām karojošo valstu armijām zaudējumi augstāko virsnieku~--- divīziju un brigāžu komandieru vidū. Tie liecinājuši gan par vācu prasmi iznīcināt Polijas armijas vienību komandpunktus, minēto poļu virsnieku neprasmi vadīt pakļautās daļas ar jaunāko virsnieku starpniecību, gan šo virsnieku centieniem, neraugoties uz neveiksmēm, tomēr organizēt pretsparu ienaidniekam, nežēlojot tam arī savas dzīvības. Taču var arī salīdzināt poļu kritušo ģenerāļu skaitu ar vēlākajiem vācu un padomju ģenerāļu zaudējumiem Austrumu frontē, kur vērmahts cīnījās ar Sarkano armiju.

Austrumu frontē 1941.~gadā krita 10 vācu ģenerāļu, 1942.~gadā~--- 20, 1943.~--- 27. Tātad mēnesī krita vidēji 2 ģenerāļi. Līdz Staļingradas kaujai Sarkanajai armijai neizdevās sagūstīt nevienu vācu ģenerāli. Arī pēc tam gūstā saņemts vācu ģenerālis bija liels retums. (Tā, Baltkrievijas operācijā 1944.~gadā 21 vācu ģenerālis tika sagūstīts, bet 10 tika nogalināti vai arī paši nošāvās.)

No 1941.~gada 22.~jūnija līdz 1942.~gada 1.~janvārim kaujās krita (vai arī pazuda bez vēsts) 71~padomju ģenerālis. Tātad ik mēnesi dzīvību zaudēja ap 12 ģenerāļu (3~mira, krituši gūstā). (Pavisam Sarkanajā armijā kara gados krita, mira no ievainojumiem, pazuda bez vēsts, gāja bojā gūstā, katastrofās, tika represēti 458 ģenerāļi un admirāļi, jeb ap 10\% no visiem, kuri šai laikā dienēja armijā.)

Salīdzinoši mazais kritušo poļu ģenerāļu skaits 1939.~gada septembrī neliecina par sīvu Polijas armijas pretošanos.

Zaudētu karu dalībnieku memuāri parasti atspoguļo to autoru personīgo pieredzi, tikai dažkārt pretendējot uz sintēzi. Tāpat tas ir ar 1939.~gada septembra kara poļu dalībnieku memuāriem. Gandrīz katrs memuārists attaisnoja savu darbību, norādot uz augstāk stāvošo komandieru kļūdām. Taču tāpat darīja arī šie augstāk stāvošie. Galu galā izrādījās, ka visām sakāvēm ir tikai viens vaininieks~--- virspavēlnieks. Polijas gadījumā tas bija maršals E.~Ridzs-Smiglijs. (Var atzīmēt, ka arī nacistiskās Vācijas militāristi pēc savas sakāves karā atrada vienu vaininieku~--- virspavēlnieku Ā.~Hitleru.) Noraidot šādus vienpusējus vērtējumus, pētnieki ir meklējuši citus sakāves cēloņus.

1939.~gada septembra kara kampaņu kā pirmais darbā ``\pltxti{Komentarze do historii polskich działań obronnych 1939~r.}'' (Komentāri poļu aizsardzības darbībai 1939.~gada septembrī. 1--3.d., Varšava, 1969.--1978.) nopietni analizēja viens no Varšavas aizstāvēšanas vadītajiem pulkvedis M.~Porvits. Analīze rādīja, ka kara gaitā no poļu puses bija pieļautas daudzas kļūdas un pārrēķināšanās, kurām bija nemazsvarīga loma valsts militārā sakāvē. Kļūdas pieļāva kā virspavēlniecība kopumā, tā arī atsevišķi armijas vadītāji. Mūsdienās no dažādām pusēm ir analizētas Polijas gan militāri-politiskās, gan stratēģiskās un taktiskās kļūdas.

Tā, M.~Porvits par vienu no stratēģiskām kļūdām uzskatīja nepareizu militārās attīstības plānošanu pirms kara. Milzu nauda tika piešķirta piecu zemūdeņu, četru mīnkuģu, viena mīnu licēja būvei. Visi šie kuģi kara laikā izrādījās praktiski nevajadzīgi, nauda bija izšķiesta veltīgi. Ja līdzekļi tiktu piešķirti aviācijas, tanku, artilērijas attīstībai, armija kļūtu krietni vien kaujas spējīgāka.

Varšavā uzskatīja, ka galvenais vācu trieciens vērsīsies pret Dancigu un Poznaņu, lai atgūtu pēc Versaļas miera zaudētās vācu zemes. Polijas valdošais slānis nesaprata, ka Ā.~Hitlera mērķis ir pilnīga Polijas iznīcināšana, tāpēc valsts nebija pietiekami sagatavota atvairīt triecienus citos virzienos. Bez tam Polijas vadītāji nebija rēķinājušies, ka Vācija varēs uzbrukt uzreiz ar lieliem spēkiem, uzskatīja, ka tai vajadzēs laiku mobilizācijai. Tika palaists garām brīdis, kad Vācija veica armijas mobilizāciju. Tādejādi Poliju lielā mērā pazudināja tās valdītāju pārliekā pašapziņa un pašpaļāvība. Ja ne minētās kļūdas, vērmahtam nāktos ciest krietni lielākus zaudējumus, laužot poļu pretestību.

Jau tā trūkstot spēkiem, Polijas virspavēlniecība daļu no tiem turēja austrumos uz padomju robežas. Šāda rīcība izrādījās ``ne šis, ne tas''. Ar šiem nelielajiem spēkiem nevarēja nedz atturēt PSRS no iebrukuma Polijā, nedz, tam sākoties, dot pienācīgu pretsparu Sarkanajai armijai. Ja nebūtu ukraiņu un baltkrievu iedzīvotāju nemieru draudu (par ko atbildība bija jānes Polijas vadībai), šīs divīzijas daudz lietderīgāk varētu izmantot rietumos.

Netika apbruņoti tūkstoši brīvprātīgo, kuri vēlējās karot pret nacistiem.

Kara sākumposmā Polijas armijas pirmais ešelons bija vairākās grupās caurmērā vienmērīgi izvietots gar valsts robežu. Piecpadsmit divīziju un brigāžu bija novietotas virspavēlniecības rezervē Varšavas, Lodzas un Radomas rajonā. Faktiski pavēlniecība atteicās no trieciengrupu radīšanas uzbrukuma operāciju veikšanai, ar to atdodot iniciatīvu pretinieka rokās. Tam nebija pārāk lielu grūtību pārraut vājās poļu aizsardzības līnijas.

Jau mūsdienu poļu vēsturnieks un žurnālists D.~Bališevskis raksta, ka vairākuma kļūdu cēlonis bija tas, ka kara darbības plāns tika balstīts vienīgi uz optimistiskāko notikumu attīstības scenāriju. Armijas daļu uzdevums bija neļaut sevi sakaut līdz karā iesaistīsies Rietumu sabiedrotie, ko gaidīja dažu dienu laikā. Tika domāts, ka pret Polijas armiju karos tikai daļa vērmahta, daļa Vācijas aviācijas utt. Šis autors arī secina, ka 1939.~gada septembra vidū karš jau bija zaudēts, kaut kaujas vēl turpinājās.

\strong{Stratēģiskas kļūdas tika pieļautas arī paša kara laikā.}

Tā, neraugoties uz to, ka Kara ministrijas pagrabos bija labi ierīkots komandpunkts ar laika garam atbilstošu sakaru mezglu, pārāk ātri (6.~septembrī) Galvenais štābs atstāja Varšavu, kas militārās vadības maksimālas centralizācijas apstākļos noveda pie karaspēka dezorganizācijas.

Poļu komandējošais sastāvs nebija sagatavots cīnīties pilnīga vai daļēja ielenkuma apstākļos. Gandrīz visas poļu daļas, kuras nokļuva vācu aplenkumā, tika iznīcinātas vai nonāca gūstā. Izņēmums bija vienīgi pulkveža S.~Mačeka vadītā motorizētā brigāde, kura izrāvās līdz Ungārijai, bet pēc tam pārcēlās uz Franciju un Lielbritāniju. Tai laikā, kad daļa augstāko virsnieku kopā ar saviem karavīriem gāja bojā vācu ierīkotajos ``katlos'', daudzi citi vienkārši bēga.

Kā rakstīja poļu vēsturnieks J.~Maciševskis, vismaz trīs poļu ģenerāļus: K.~Fabriciju, J.~Rummelu un S.~Domb-Bernacki, kuri pameta komandētās armijas savam liktenim, bet paši bēga, par to vajadzēja nodot kara tribunālam. Citi vēsturnieki šim sarakstam pieskaita arī ģenerāli V.~Bonč-Uzdovski, pulkvedi E.~Dojanu-Surovku u.c., taču poļu pavēlniecība nekādus lēmumus šai sakarā nepieņēma.

Milzīga stresa apstākļos atsevišķi komandieri pieļāva kļūdas, kurām bija katastrofālas sekas. Tā, karaspēka grupas komandieris ģenerālis M.~Boruta-Spehovičs kampaņas beigās pasteidzās pāragri atlaist savu štābu, kas noveda pie grupas sagrāves.

Kad Polijas pavēlniecība redzēja, ka nav iespējams evakuēt Galvenā štāba Otrās (izlūkošanas) nodaļas arhīvu, tā nolēma to iznīcināt un šai nolūkā izvest uz vienu no fortiem pie Varšavas, taču to ātri uzzināja Vācijas armijas izlūkdienests~--- abvērs. Tā operatīvā grupa ieradās, kad poļu virsnieki tikai vēl gatavojās sākt arhīva iznīcināšanu. Materiālus knapi izvietoja sešās kravas mašīnās un izveda uz Berlīni. Vācieši ieguva visu plašo poļu aģentu sarakstu gan Rietumeiropas valstīs, gan PSRS. Sagrābtie dokumenti tika ātri izpētīti. Tādā ceļā vācu izlūkdienestam izdevās atklāt visu poļu izlūkošanas tīklu ārzemēs, kurā bija ap 4~000 aģentu. Daļa poļu aģentu, izņemot tos, kuri darbojās pret Vāciju, tika pārvervēti un sāka darboties hitleriešu labā. Piemēram, viens no poļu komunistu vadītājiem M.~Močars savās atmiņās stāstīja, kā 1943.~gada aprīlī vācu politiskā policija Gestapo deva jūtamu triecienu Polijas Strādnieku partijas Lodzas organizācijai, arestējot vairākus tās biedrus. Kā vēlāk izrādījās, Gestapo kā provokatoru izmantoja vēl ``sanācijas'' laiku policijas aģentu. Poļu arhīvs arī ļoti palīdzēja vāciešiem, plānojot jūras desanta operāciju pret Angliju. Tas deva pilnīgu priekšstatu par britu karaspēka izvietojumu, izdevīgākajām desantēšanās vietām.

``Sociālistiskās'' Polijas laikā poļu komunistu vadītāji uzsvēra, ka Polijas sakāve 1939.~gada septembrī bija poļu buržuāzijas un muižniecības šauri šķiriskās pretpadomju politikas rezultāts, kura ``naidu pret pirmo strādnieku un zemnieku valsti~--- Padomju Savienību stādīja augstāk par savas valsts interesēm. Buržuāziski-muižnieciskās Polijas valdošās šķiras nespēja garantēt poļu nācijas pastāvēšanas elementāros pamatus\dots{}'' Kaut doma izteikta ``augsti'' patētiskā stilā, tajā ir liela daļas patiesības. Šķiriskais un nacionālais naids pret austrumu kaimiņvalsti neļāva Polijas valdošajām aprindām adekvāti novērtēt starptautisko situāciju. Polijas militārā katastrofa bija arī sekas Polijas atteikumam sadarboties ar PSRS, kas tai tika piedāvāts Lielbritānijas-Francijas-PSRS sarunās 1939.~gada pavasarī un vasarā.

Vācijas (un Slovākijas) spēku pārsvars bija galvenais Polijas sakāves cēlonis. Taču sava loma bija arī iekšējiem cēloņiem~--- pretrunām Polijas daudznacionālajā, sociāli neviendabīgajā sabiedrībā, arī subjektīvajām valdošo slāņu kļūdām gan ārējās, gan iekšējās politikas laukā.

Jau minētais Polijas Tautas Republikas vēsturnieks Z.~Zaluskis, analizējot kara notikumus, uzsvēra, ka 1939.~gada septembrī karu zaudēja Polijas pastāvošās iekārtas nespēja nodrošināt savu un poļu nācijas aizsardzību. Pēc viņa sprieduma 1939.~gada 1.~septembrī Polija atradās: katastrofālā ārpolitiskā izolācijā, katastrofālā stratēģiskā stāvoklī, katastrofālā militāri-ekonomiskā vājumā salīdzinājumā ar pretinieku. Polija karu bija zaudējusi jau pirms tā sākumam. Karš bija zaudēts jau tad, kad Polija izveidojās kā buržuāzijas un muižnieku valsts ar tendenci paplašināt savas robežas austrumu virzienā, pat nāvīgu briesmu priekšā nespējīga pārvarēt naidīgo attieksmi pret PSRS. Galvenos savus spēkus Polija koncentrēja uz austrumu robežas, lai nodrošinātu kundzību pār iekšējām kolonijām~--- Rietumukrainu un Rietumbaltkrieviju. Vēsturnieks konstatēja, ka, ieguvusi neatkarību, Polija kapitālistiskās iekārtas ietvaros nespēja panākt tās kapitālistiskās valstis, kuras bija to apsteigušas attīstībā. Tai nācās izciest ne tikai savas, bet arī visas kapitālistiskās pasaules krīzes. Polijas ekonomika nespēja nodrošināt vajadzīgo militārā potenciāla pieaugumu. Šo Polijas ceļu uz militāro katastrofu vadīja noteiktas šķiras~--- buržuāzija un muižniecība, to realizēja noteiktas politiskās partijas, viss sanācijas režīms. Tas izrādījās pietiekami, lai tauta ``sanācijas'' režīmam dotu iznīcinošu spriedumu.

Z.~Zaluskis arī norādīja, ka tūlīt pēc sakāves poļu buržuāziskie politiķi priekšplānā izvirzīja nevis sakāves cēloņu, bet sakāves vaininieku meklējumus, kas bija otršķirīga problēma, jo pat visstingrākā viena nozieguma vaininieku sodīšana nevar pasargāt no citiem noziegumiem, kurus veic citas personas. Tieši partijām, kuras konkurēja ar ``sanācijas'' spēkiem, jautājums par vainīgajiem bija ļoti svarīgs. Norādot uz ``sanāciju'', kā vainīgu par septembra katastrofu, tika aizmaskoti dziļākie sakāves cēloņi, kas pavēra iespēju bijušā režīma konkurentiem atkal pretendēt uz varu. (Plašu masu apziņā tik tiešām sakāves cēloņu izpratne palika ``sanācijas'' nosodīšanas līmenī. Tās konkurenti, arī endeki, kuri nacionālajā jautājumā piekopa krietni šovinistiskāku politiku kā ``pilsudskieši'', savu vainu par pirmskara Polijas valsts kursa neveiksmīgo rezultātu neatzina. Bija vajadzīgi vēl pārbaudījumu gadi, lai labējie politiskie spēki sakompromitētu sevi tautas acīs.)

Poļu vēsturnieks mēģināja atbildēt arī uz jautājumu, vai pirmskara periodā citu pilsonisko spēku vadīta Polija būtu spēcīgāka par ``sanācijas'' režīma vadīto? Vai būtu labākas tās starptautiskās pozīcijas? Vai Polija būtu ekonomiski stiprāka? Vai tai būtu stiprāka armija? Viņa atbildes: Iespējams, ka Polijai būtu stiprāka militārā savienība ar Franciju. Taču tas nenozīmētu, ka franču armija būtu izkustējusies no Mažino līnijas, ka franči būtu vairāk gatavi ``mirt par Dancigu''. Iespējams, franču Minhenes politikas piekritēju ietekmē Polija 1939.~gadā pat piekāptos Vācijas ultimātam. Līdz 1926.~gada maija apvērsumam realizētās Polijas ekonomiskās politikas rezultāti neliecināja, ka pret J.~Pilsudski noskaņotās partijas prata saimniekot labāk nekā sanācijas politiķi. Gan demokrātijas, gan autoritārā režīma periodā Polijā tika izsaimniekoti militārajām vajadzībām izdalītie līdzekļi, bet pēc apvērsuma tomēr tika uzceltas vairākas bruņojuma rūpnīcas. Tā ir pamats domāt, ka arī bez ``sanācijas'' režīma Polijas liktenis 1939.~gadā būtu tikpat traģisks.

Jāsaka, ka šajā marksistiskajā analīzē ir daudz patiesa. Taču tā tomēr nav visa patiesība. Galvenais~--- kā rāda daudzu valstu vēsturiskā pieredze~--- īsā laikā (bet Polijas rīcībā tie bija divi nepilni gadu desmiti) panākt civilizācijas ziņā apsteigušās valstis praktiski nav iespējams, tādi vai citādi vājuma punkti, kas ietekmē turpmāko attīstību, tomēr saglabājas. No vēstures mantotā ekonomiskās, sociālās un politiskās attīstības atpalicība neļāva Polijai, piemēram, apsteigt Vāciju militārā ziņā. Cita lieta, ka, izvēloties citas prioritātes, citus attīstības līdzekļus, Polija, iespējams, 1939.~gadā varēja sagaidīt Otrā pasaules kara sākumu ekonomiski spēcīgāka; ar viendabīgāku sociālo struktūru; ar sabiedrotajiem, kuri būtu vairāk gatavi to atbalstīt pret agresiju; ar labāku stratēģisko robežu konfigurāciju, bez Slovākijas Republikas (slovāku \sktxti{Slovenská republika}) Vācijas pusē, pat bez PSRS parakstīta Molotova-Ribentropa pakta un slepenā protokola pie tā. Taču tikpat iespējams, ka citi~--- ne ``sanācijas''~--- spēki tāpat būtu cietuši tādu pat politisku un militāru sakāvi. Šeit ir jāpiekrīt Z.~Zaluskim. (Zīmīgi gan, ka vēsturnieks pat nepieminēja teorētisku variantu, ka Polijā pirms Otrā pasaules kara sākuma pie varas varētu nonākt kreisie spēki (PPS, PKP, citas organizācijas. Ar to netieši tika atzīts, ka pēc 1920.~gada kara pieredzes Polijā praktiski nepastāvēja kreisa alternatīva.) Ne viss (pēc viņa viedokļa) bija atkarīgs no pašas Polijas. Tā, Rietumvalstis savu egoistisko interešu vārdā nepildīja palīdzības solījumus Polijai. PSRS~--- valsts, kuras objektīvās intereses prasīja apstādināt Vācijas agresiju uz austrumiem, vadība, reaģējot uz Polijas un Rietumvalstu iepriekšējo politiku, bija norūpējusies nevis par palīdzību kaimiņvalstij Polijai, potenciālā sava pretinieka Vācijas sagraušanu, bet vēlējās izmantot situāciju PSRS teritorijas paplašināšanai. Tā ne tikai Polijas valdītāju, bet veselas virknes valstu vadītāju kļūdas, savas valsts interešu izvirzīšana realizējamās politikas centrā neļāva apstādināt nacistisko agresiju jau 1939.~gadā. Poļu tautai tas nozīmēja piecu gadu ilgu nacistiskās Vācijas jūgu, milzumu cilvēku dzīvību, grandiozus materiālos zaudējumus.

Uz izklāstītās analīzes pamata Z.~Zaluskis izdarīja secinājumu, ka līdzīgu sakāvju kā 1939.~gada septembrī novēršanai Polijas sabiedrībai bija radikāli jāpārveido sava valsts sociālisma virzienā. Taču vēlākā vēsturiskā pieredze, kuru vēsturnieks diemžēl līdz galam neizdzīvoja, rādīja, ka Polijas pārveide staļiniskā, pēc tam birokrātiskā ``sociālisma'' virzienā ar lieliem upuriem risināja tikai dažas problēmas, vēl vairāk atstājot nerisinātas. Polijas attīstības pēc Otrā pasaules kara demokrātisku alternatīvu savas dzīves laikā Z.~Zaluskis nesaskatīja, vai, iespējams, neuzdrošinājās saskatīt.

Kā jau mūsdienās avīzē ``\pltxti{Tribuna}'' rakstīja poļu žurnālists P.~Skura, ja ``sociālistiskās'' Polijas pastāvēšanas gados par galveno Polijas sakāves cēloni septembra kampaņā parasti nosauca solītās Rietumu sabiedroto palīdzības iztrūkumu, tad mūsdienās ``ir modē akcentus pārvietot uz notikumiem, saistītiem ar padomju agresiju.'' Taču ar tautas atmiņu manipulēt nav tik viegli. Ap 60\% poļu pat nezin, kad īsti PSRS ieveda savu karaspēku Polijā. Tikai katrs ceturtais varēja nosaukt konkrēto datumu~--- 17.~septembri. Galvenokārt tie bija cilvēki, kuri paši bija pārdzīvojuši Otro pasaules karu. Toties 1.~septembri, kad sākās Vācijas iebrukums Polijā, atceras praktiski visi.

Jāuzsver, ka neraugoties uz armijas sagrāvi, visas valsts teritorijas okupāciju, Polija oficiāli nekapitulēja Vācijas un tās sabiedrotās Slovākijas priekšā. Bez pretestības kustības iekšzemē cīņu pret Vāciju turpināja Polijas valdība emigrācijā (PVE) un poļu formējumi Rietumu sabiedroto un PSRS armijās.

Ne bez pamata poļu vēsturnieks P.~Večorkovičs atzina, ka ``mēs izglābām angļus un frančus, jo Hitlers vairs nespēja,~--- kā viņš plānoja,~--- uzbrukt Rietumiem 1939.~gadā.'' (Vērmahta Virspavēlniecības direktīvā N3, izdotā 3.~septembrī, jau bija norādīts, ka pēc Polijas iekaroto rajonu ``nomierināšanas'' būs jāuzsāk ``karaspēka pielietošana Rietumos''. To pašu apstiprināja direktīva N4 arī 25.~septembrī. Citiem vārdiem, pēc Polijas jauns trieciens tika plānots Rietumos.) Vēsturnieks arī piebilda: ``Ja tā notiktu, jau 1940.~gadā viņš [Hitlers] savu karaspēku nosūtītu pret Padomju Savienību, kas, savukārt, nozīmētu šīs valsts iznīcināšanu. Tādā veidā poļi \citespace{} izglāba ādu Staļinam''. Tas, ka 1940.~gadā Ā.~Hitlers nenosūtīja vērmahtu pret PSRS, nebija atkarīgs vienīgi no Polijas armijas pretošanās 1939.~gada septembrī, te ar vēsturnieka muti runā nacionālā tieksme pārspīlēt savas valsts nozīmi pasaules notikumos, taču jāatzīst, ka Polijas pretestībai agresoram 1939.~gadā bija pozitīva nozīme antihitleriskās koalīcijas izveidē un uzvarā.

Poļu tautas upuri 1939.~gada septembrī nebija velti. Uzņemot pirmo Vācijas triecienu, Polija deva iespēju saviem Rietumu sabiedrotajiem vēl astoņus mēnešus gatavoties nu jau īstam karam ar nacistisko Vāciju, ko tie varēja izmantot lai pabeigtu karavīru apmācīšanu, kāpinātu militārās produkcijas ražošanu. Tiesa, Rietumvalstis neizmantoja pilnībā tām atvēlēto laiku. 1939.~gada septembra karā tika nodemonstrētas jaunās ``zibenskara'' vešanas metodes: masīvi aviācijas un dziļi tanku-motorizētā karaspēka triecieni, kas sašķēla pretinieka aizsardzības līnijas, piecreiz augstāki nekā iepriekš uzbrukuma tempi un ievērojami lielāks veicamo militāro operāciju dziļums, ko nodrošināja jaunie kara vešanas līdzekļi. Nākotnē tas vairs nevarēja būt pārsteigums. Polijas armijas pieredze, ja tā tiktu pamatīgi apgūta, rādīja arī veidus (vismaz taktikas ziņā) kā pretoties jaunajām kara vešanas metodēm. Taču Rietumos Polijas kara pieredzi nopietni neanalizēja, var teikt~--- pat gandrīz vispār neņēma vērā. Ziņojums par Vācijas pielietotajām kara vešanas metodēm, ko sastādīja Polijas virsnieki pēc septembra kampaņas, un ko ģenerālis V.~Sikorskis iesniedza Francijas virspavēlniecībai, Francijas sakāves dienās tika atrasts starp neatvērtajiem un neizskatītajiem Francijas armijas virspavēlnieka ģenerāļa M.~Gamelēna papīriem. Rietumu politiķi un militāristi nenovērtēja tās briesmas, kas viņiem pašiem draudēja pēc Polijas sakāves. Viņiem Austrumeiropā nebija tik svarīgu interešu, kuras būtu vērts nopietni aizstāvēt.

Lai kā daži poļu politiķi un arī vēsturnieki visu vainu par 1939.~gada sāpīgo sakāvi mēģina uzkraut tikai Poliju okupējušajām valstīm~--- Vācijai un PSRS, daļēji arī savus solījumus nepildījušajiem Rietumu sabiedrotajiem, jāatzīst, ka liela daļa atbildības tomēr jānes pašas Polijas pārlieku augstprātīgajiem un nepietiekami tālredzīgajiem politiķiem.

Sakāve Polijai nenozīmēja karadarbības beigas, poļu politiķi centās atjaunot un legalizēt valsts augstāko varu. Pēc bēgšanas no valsts un internēšanas Rumānijā I.~Moscickis atbilstoši 1935.~gada konstitūcijas 13.~pantā paredzētajam gadījumam, ja valsts prezidents nespēj pildīt savas funkcijas, 25.~septembrī atkāpās no amata (pateicoties tam 1939.~gada decembrī I.~Moscickis varēja pamest Rumāniju un apmesties Šveicē, kur arī pavadīja visus Otrā pasaules kara gadus) un par savu pēcteci nozīmēja ģenerāli un diplomātu, Polijas sūtni Itālijā B.~Veņavu-Dlugošovski. Taču viņš prezidenta amatā sabija tikai pāris dienu. Pret šo poļu darbinieku protestēja opozīcija, viņu Polijas vadības emigrācijā priekšgalā nevēlējās redzēt arī Lielbritānija un sevišķi Francija, faktiski izvirzot savu \latxti{veto}, kas pēc būtības bija klaja iejaukšanās Polijas iekšējās lietās (tāpēc Rietumu angažētie vēsturnieki to parasti ``aizmirst''). Tā faktiski ārvalstu spiediena rezultātā B.~Veņava-Dlugošovskis no amata atteicās un tas tika piedāvāts bijušajam Senāta maršalam (priekšsēdētājam) V.~Račkevičam, kurš tolaik uzturējās Parīzē. 30.~septembrī viņš amatu pieņēma un nekavējoties Polijas sūtniecībā Parīzē nodeva prezidenta zvērestu. Tātad faktiski Polijas prezidents nespēja iecelt savu pēcteci pat tajā kārtībā, kādu paredzēja nedemokrātiskā 1935.~gada konstitūcija. Viņam ``palīdzēja'' sabiedrotie, kuri tikko bija pametuši Poliju nelaimē. Jau kā jaunieceltais prezidents V.~Račkevičs atbrīvoja ģenerāli F.~Slavoju-Skladovski no premjerministra amata un uzticēja Polijas valdības emigrācijā (\pltxti{Rząd Rzeczypospolitej Polskiej na uchodźstwie})~--- PVE~--- veidošanu ģenerālim V.~Sikorskim, kurš arī spēlēja galveno lomu emigrācijā. 2.~oktobrī prezidents izdeva dekrētu par Seima un Senāta atlaišanu.

Zīmīgi, ka valdības priekšgalā izvirzījās tieši V.~Sikorskis. V.~Sikorskis bija ļoti piemērota kandidatūra valdības vadītāja amatam. Viņa, bijušā leģionāra nopelnus Polijas valstiskuma izcīnīšanā un nostiprināšanā neviens nevarēja apstrīdēt, taču tai pat laikā viņš bija J.~Pilsudska un ``sanācijas'' pretinieks, atstājis Poliju jau 1928.~gadā, un nenesa nekādu atbildību par ``sanācijas'' laikā notikušo Polijā, ko gan ``sanācijas'' piekritēji viņam nekad nepiedeva. V.~Sikorskis baudīja uzticību dažādos poļu sabiedrības slāņos, ieskaitot arī sociālistus un ļudoviešus. Viņam bija arī Francijas un Lielbritānijas valdošo aprindu atbalsts. Franču pilsētā Anžerā (\pltxti{Anger}) tika radīta Polijas valdība emigrācijā.

Steigu ar jaunās valdības formēšanu lielā mērā izsauca arī poļu politiķu bažas, ka Vācija un Padomju Savienība kopīgi vai arī katra atsevišķi var radīt Polijā marionešu valdību un parakstīt ar to separātu miera līgumu. Kopīgā vācu-padomju komunikē 1939.~gada 18.~septembrī taču bija teikts, ka abas puses ir gatavas palīdzēt poļu iedzīvotājiem ``pārkārtot savas valstiskās pastāvēšanas apstākļus''.

Valdības veidošanai emigrācijā jau bija vēsturiski precedenti. Pirmā pasaules kara laikā, kad visa Beļģijas teritorija bija Vācijas okupēta, Beļģijas valdība turpināja pastāvēt Francijas teritorijā un to atzina sabiedrotie. Čehoslovākijas bijušais prezidents E.~Benešs jau pirms Otrā pasaules kara sākuma bija Londonā radījis Čehoslovākijas valdības emigrācijā tehnisko aparātu un sarunās ar dažādu valstu pārstāvjiem zondēja iespējas to pārvērst par citu valstu atzītu valdību emigrācijā. Tomēr, kaut Otrā pasaules kara laikā valdību emigrācijā likumības atzīšana bija diezgan plaši izplatīta (tā noderēja Vācijas okupēto zemju tautu un arī pasaules sabiedriskās domas mobilizācijai pret agresorvalsti), taču no mūsdienu starptautisko tiesību viedokļa tādas valdības neatbilst starptautiskajiem atzīšanas kritērijiem, galvenokārt tāpēc, ka tās nav saistītas ar savas valsts teritoriju un iedzīvotājiem, tās neveic efektīvu pār tiem kontroli un tām nav varas pilnvaru.

1939.~gada 4.~oktobrī, ziņojot par noskaņojumu Parīzē strādājošo diplomātu aprindās, padomju sūtnis Francijā J.~Surics, rakstīja, ka Polijas jaunizveidotās ``valdības pilnvaras atklāti tiek apšaubītas. Neviens nešaubās ka Moscicka konstitucionālais akts ir izdots ``\latxti{ad~hoc}'' (lat.~--- speciāli dotajam mērķim) un ne Polijas teritorijā.''

Valdību emigrācijā galvenokārt veidoja pilsoniskie opozīcijas spēki, kuri bija cīnījušies pret ``sanāciju'', pēdējās 1938.~gada novembrī notikušās parlamenta vēlēšanas Polijā bija boikotējuši, jo ``sanācijas'' režīmu poļu sabiedriskā doma 1939.~gada rudenī pamatoti apvainoja rīcības nespējā un tam uzlika atbildību par sakāvi.

Starp citu, franču vairākums 1939.--1940.~gadā ne visai laipni uzņēma poļus, jo uzskatīja, ka karš bija sācies Polijas dēļ. Kā rakstīja J.~Kļimkovskis, franči uzskatījuši karu par ``lielu lietu, ko uzsākusi maza nācija''. Katrs otrais francūzis atkārtojis, ka viņam jākaro par Poliju, lai tai piederētu Gdaņska, katrs trešais vispār nevēlējies karot. Iespējams, J.~Kļimovskis pārspīlēja Polijas lomu franču apziņā, bet vismaz daļas no viņiem nevēlēšanos karot apstiprina arī daudzas citas liecības.

V.~Sikorskis Francijā izveidoja valdību formāli no četru Otrās republikas opozīcijas partiju: Nacionālās partijas (\pltxti{Stronnictwo Narodowe}), Zemnieku partijas (\pltxti{Stronnictwo Ludowe}), Polijas sociālistiskās partijas (\pltxti{Polska Partia Socjalistyczna~--- PPS}) un Darba partijas (\pltxti{Stronnictwo Pracy}) pārstāvjiem. Citus kritērijus norādījusi padomju vēsturniece V.~Parsadanova: valdības pamatu veidoja t.s. ``Moržas fronte'' (\pltxti{Front Morges}) un Tirdzniecības bankas grupa, kas bija galvenokārt ar probritānisku orientāciju. Pēc 1939.~gada sakāves pirmskara t.s. Moržas frontē ietilpstošo spēku biedri pārmeta ``sanācijai'', ka tā novedusi valsti pie katastrofas. ``Sanācijas'' atbalstītāji nepalika atbildi parādā, savukārt apvainojot savus oponentus valsts vājināšanā. Pieaugošais pret ``sanāciju'' kritiskais, daļēji pat naidīgais noskaņojums ļāva bijušajai legālajai opozīcijai daļēji īstenot savas pretenzijas uz varu (precīzāk~--- uz starptautisko un iekšpolitisko Polijas interešu pārstāvniecību).

Taču reizē ar mērķi apvienot visus pilsoniskās Polijas atjaunošanas piekritēju spēkus V.~Sikorskis noslēdza kompromisu ar agrāk valdošajiem spēkiem. Tā it kā nācijas, bet faktiski valdošo slāņu apvienošanas interesēs ``nacionālās vienotības valdības emigrācijā'' sastāvā tika iekļauti arī bankrotējušā sanācijas režīma pārstāvji. PVE koalīcijas raksturs demonstrēja, ka ``sanācijas'' režīma spēki, būdami galvenie vainīgie nacionālajā katastrofā, nevarēja turpmāk vieni paši pārstāvēt tautu. Tomēr virkni amatu saņēma arī mazāk sakompromitētie ``sanācijas'' režīma darbinieki. To pārstāvji bija ģenerālis K.~Sosnkovskis, (viņu V.~Račkevičs nozīmēja par savas vietas izpildītāju gadījumam, ja viņš pats nevarēs pildīt prezidenta pienākumus līdz kara beigām, ģenerālis kļuva arī par ministru bez portfeļa, vienlaikus par Polijas iekšējo lietu ministru komitejas priekšsēdētāju (\pltxti{przewodniczącym Komitetu Ministrów do Spraw Kraju}) un par pilsoniskās pretošanās kustības pavēlnieku (\pltxti{komendant główny Związku Walki Zbrojnej}), pulkvedis A.~Kocs, kurš decembrī kļuva par iekšlietu ministru) u.c. Par ārlietu ministru savas profranciskās orientācijas dēļ tika iecelts J.~Pilsudska sekotājs A.~Zaleskis. Viņa izveidotais diplomātu korpuss sastāvēja gandrīz tikai no J.~Pilsudska piekritējiem. ``Sanācijas'' politiķu grupa ievērojami ietekmēja valdības emigrācijā darbību ar savu cilvēku izvirzīšanu dažādos amatos. ''Sanācijas'' un opozīcijas politiķus vienoja mērķis atgriezties Polijā un saņemt varu. Kā rakstīja poļu vēsturnieks J.~Šļusarčiks, ``Londonas valdības darbības galvenais virziens bija gatavošanās cīņai par varu.''

Jau tas vien, ka 1939.~gada septembrī bija ciesta militāra katastrofa, kas pierādīja ne tikai visas pirmskara ārējās politikas bankrotu, bet arī iekšējās, saimnieciskās un sociālās politikas vājās vietas, faktiski liedza ``sanācijas'' režīma politiķiem tiesības uz valsts tālāku pārvaldi. Par sakāvi bija jānes politiska un morāla atbildība. Arī austrumu teritoriju zaudēšana faktiski bija šīs sakāves sekas. Taču uzņemties par to vainu ``sanācijas'' politiķi nevēlējās. Viņiem bija izdevīgi nostāties patriotu stājā un faktiski turpināt pirmskara politiku, censties poļu tautas masas iesaistīt cīņā par citu tautu vairākuma apdzīvotu teritoriju paturēšanu Polijas sastāvā, lai arī tas solīja tautai jaunus upurus. Polijas valdībai emigrācijā, bet īpaši ``sanācijas'' režīma darboņiem bijušo valsts austrumu teritoriju zaudēšana bija nepieņemama ne tik daudz no Polijas ekonomisko, politisko un prestiža interešu, cik no savas varas saglabāšanas viedokļa.

Taču faktiski arī neviena no ``sanācijai'' opozīcijā bijušajām pilsoniskajām politiskajām partijām nevēlējās noskaidrot septembra katastrofas cēloņus, caurskatīt starpkaru perioda iekšējās un ārējās politikas principus. Daudziem patriotiski noskaņotiem zemnieku kustības pārstāvjiem un sociālistiem bija grūti iedomāties citu ārējo politiku, kā to realizēja ``sanācijas'' režīms, atklāti nosodīt to par valsts novešanu līdz septembra katastrofai, starptautiskas drošības sistēmas izveides traucēšanu.

Šeit ir vietā pateikt par mēģinājumu noskaidrot atbildīgos par 1939.~gada septembra sagrāvi. Jau 1939.~gada 10.~oktobrī emigrācijā tika izveidota speciāla komisija septembra katastrofas cēloņu noskaidrošanai. Tās priekšgalā nostājās autoritatīvais ģenerālis J.~Hallers. 1940.~gada sakumā tika atjaunināts komisijas sastāvs, taču tās darbība tik un tā nenesa rezultātus. Acīmredzot, par 1939.~gada katastrofu atbildīgajiem arī emigrācijā pietika spēku nepieļaut komisijas sekmīgu darbību.

Faktiski tikai viens no V.~Sikorska komandas dalībniekiem~--- sociālās nodrošināšanas ministrs sociālists J.~Stančiks stingri iestājās par to, lai bijušie premjerministrs, ārlietu ministrs, kara ministrs un virspavēlnieks stātos kara tribunāla priekšā par savu konstitūcijā noteikto pienākumu nodrošināt valsts aizsardzību neizpildīšanu. Šis jautājums tika divreiz skatīts valdības sēdē. V.~Sikorskis un ārlietu ministrs A.~Zaleskis uzstājās pret priekšlikumu. Ģenerālis paziņoja: ``Gružus no istabas ārā nenesīsim. Atgriezīsimies pie šī jautājuma pēc kara''. Droši vien V.~Sikorskis saprata, ka lēmumu par tiesas sarīkošanu pār nosauktajām ``sanācijas'' režīma amatpersonām, ja par tādu paziņotu radio, piemēram, BBC (\entxti{British Broadcasting Corporation}), Polijas iedzīvotāju vairākums uzņemtu ar apmierinājumu un sapratni. Taču V.~Sikorskim un viņa atbalstītājiem bija svarīgi demonstrēt poļu vienotību Francijai un Lielbritānijai, un viņi baidījās radīt nelabvēlīgu iespaidu sabiedrotajos. Ir jāatzīst, ka ārējo briesmu priekšā emigrācija bija spiesta ievērot zināmu vienotību, taču nevar noliegt arī to, ka gala rezultātā vispareizākā politika izrādās tā, kura balstās uz patiesības atzīšanu. (Tiesa, Polijas armijā Francijā bija arī cilvēki, kuri bija gatavi pretoties ``sanācijas'' režīma darbinieku ietekmes saglabāšanai, vēlējās atvērt ceļu jauniem politiķiem. Viens no tiem, kapteinis T.~Kerns atklāti paziņoja, ka poļiem atgriežoties dzimtenē, uz tās robežas būtu jānošauj visus virsniekus no majora un augstāk. Paziņojums ``sanācijas'' aprindās izsauca milzu uztraukumu.)

Diemžēl poļu ``sanācijas'' režīmam opozīcijā bijušās poļu liberālās, demokrātiskās aprindas bija pārāk vājas, lai panāktu kā emigrācijā, tā dzimtenē ``sanācijas'' darboņu padzīšanu un savu nostiprināšanos pēc sakāves jaunveidojamās varas struktūrās. Īpaši vāja bijusī opozīcija bija militārajā jomā, kuru gan J.~Pilsudskis, gan viņa mantinieki sevišķi stingri turēja savās rokās. Bijusī opozīcija samierinājās ar to, ka tā gan piekļuva varas iestādēm emigrācijā un dzimtenē, bet nespēja tās pilnībā kontrolēt. Tas nozīmēja arī to, ka tā nevarēja līdz galam atklāti kritizēt savus bijušos politiskos pretiniekus, bet pēc tam sabiedrotos valdošajā koalīcijā. Tāpat tika pārņemtas J.~Pilsudska un ``sanācijas'' režīma pretenzijas uz austrumu teritorijām, kas Otrā pasaules kara gaitā dārgi izmaksāja poļu tautai. Arī iepriekšējā režīma iekšpolitikas vērtējumā nenotika būtiskas izmaiņas.

Poļu ekonomists L.~Landau, kurš jau kara gados rakstīja kara un okupācijas gadu hroniku (``\pltxti{Kronikę z lat wojny i okupacji}''), publicētu 1962.--1963.~gadā, jau 1939.~gada 21.~novembrī tajā par PVE ierakstīja: ``Sabiedrība ar labvēlību, kaut arī ar nepacietību novēro šīs valdības darbību, taču visnoteiktākā veidā pretotos pret tās atgriešanos Polijā. Valdība ir cita, nekā tā, kura bija Polijā pirms kara, taču pārstāv to pašu sistēmu, to pašu ideoloģiju\dots{}'' Pēc minētā Polijas armijas rotmistra J.~Kļimkovska atmiņām V.~Sikorska valdībā vislielākā ietekme bijusi agrākajiem Ģenerālštāba Otrās (Izlūkošanas un pretizlūkošanas) nodaļas darbiniekiem, kā arī augstākajiem ``sanācijas'' laika virsniekiem. Domājams, ka atmiņu autors bija tendēts pārspīlēt virsnieku lomu poļu emigrācijas vadībā, taču to, ka bankrotējušā ``sanācijas'' režīma darbinieki ieguva lielu ietekmi V.~Sikorska valdībā, tai pakļautajās struktūrās, pierāda liels skaits avotu un par to ne vienu reizi vien nāksies runāt arī šai grāmatā Otrajam pasaules karam veltītajās lappusēs.

PVE struktūra gandrīz dublēja pirmskara valdības variantu, emigrācijas apstākļos tā bija visai smagnēja. Pastāvēja deviņas ministrijas (iekšlietu, tieslietu, finanšu, rūpniecības, tirdzniecības un kuģniecības, izglītības un skolu problēmu), vairāki biroji, komitejas un komisijas.

Kā PVE labvēlīgs faktors jāatzīmē tas, ka nacistu nagos nenonāca valsts zelta rezerves gandrīz 80~miljonu ASV dolāru vērtībā~--- tās tika novietotas vispirms Stambulā, bet pēc tam Dakārā, Franču Rietumāfrikā (tagadējās Senegālas teritorijā). Kad Ā.~Hitlers 1940.~gadā okupēja Franciju, zeltu pārveda uz Kanādu, bet vēlāk tas tika nogādāts Londonā. Ar šiem uzkrājumiem poļu valdība (un tās darbinieki) emigrācijā vismaz daļēji bija finansiāli neatkarīga, citādi tai būtu klājies krietni grūtāk. Tiesa, tas nenozīmē, ka naudas pietika arī savu bruņoto spēku uzturēšanai, Polijas pilsoņu materiālai atbalstīšanai. Līdzekļus PVE darbam deva britu kredīti par 5\% gadā. Pēc V.~Čērčila vārdiem, Lielbritānija poļu vajadzību apmierināšanai kara gados iztērēja 120~miljonus sterliņu mārciņu.

Kaut jaunizveidotās valdības likumību varēja apšaubīt, Francija un Lielbritānija tās pilnvaras nekavējoties atzina, tām drīz sekoja lielākā daļa neitrālo valstu. Toties Padomju Savienība neatzina Polijas valdību emigrācijā līdz brīdim, kad pati kļuva par nacistiskās Vācijas uzbrukuma objektu un kad svarīga kļuva sabiedroto iegūšana Rietumos.

Par vienīgo savu legālo darbības pamatu valdība atzina 1935.~gadā pieņemto nedemokrātisko Aprīļa konstitūciju. Tā kā Polijas valdības emigrācijā leģitimitāte balstījās uz Konstitūcijas pamata, kuru opozīcija pirms tam bija noraidījusi kā ``nedemokrātisku'' un ``autoritāru'', no tā izrietēja secinājums, ka arī pati valdība ir ``nedemokrātiska'' un ``autoritāra''. (Šo pretrunu 1943.~gadā un vēlāk veikli izmantoja poļu komunisti un to sabiedrotie.)

Visi politiķu apgalvojumi par Polijas valdību emigrācijā likumību balstījās uz viņu pagaidu vienošanos par sadarbību kara apstākļos. Faktiski to varēja uzskatīt, augstākais, par pagaidu valdību, kamēr nebija notikusi demokrātiska nācijas gribas noskaidrošana.

To, ka PVE tika Rietumu lielvalstu starptautiski atzīta, nevar uzskatīt par tās leģitimitātes apliecinājumu. Rietumvalstis tikpat viegli kā 1939.~gadā šai valdībai atzīšanu deva, 1945.~gadā, tāpat pragmatisku apsvērumu vadītas, to liedza. To, kā PVE vērtēja ASV, parādīja prezidents F.~D.~Rūzvelts Jaltas konferencē 1945.~gada 7.februārī, kur viņš negaidīti paziņoja, ka faktiski no 1939.~gada nekāda Polijas valdība neeksistē un sabiedroto pienākums ir palīdzēt poļiem radīt apstākļus tādas valdības izveidei.

Tiesa, vēlāk PVE līdzīgā stāvoklī atradās arī citu nacistu sagrābto zemju valdības emigrācijā, kuras gandrīz visas kara laikā uzturējās Londonā. Valstu~--- nacistu pretinieku interesēs bija tās uzskatīt par savu nāciju pārstāvēm, tāpēc tās tika starptautiski atzītas un ar tām uzturēti diplomātiski sakari. Šādas valdības emigrācijā, nevarot pierādīt savi leģitimitāti, varēja tikai lūgt, bet ne prasīt savu zemju pilsoņu atbalstu. Tāpēc PVE centās sevi apliecināt kā nacionālo interešu aizstāvi, pasludinot ``sanācijas'' laika pārvaldes metožu noraidīšanu un solot pēc kara atbrīvotajā Polijā veikt virkni sociāli-ekonomisku reformu.

Par to, ka no dažādiem spēkiem (opozīcijas un ``sanācijas'' piekritējiem) radīta valdība nespēja būt stabila, liecināja vēlāk sekojošās krīzes.

Kad 1939.~gada novembrī E.~Ridzs-Smiglijs atteicās no virspavēlnieka amata, 7.~novembrī V.~Račkevičs iecēla V.~Sikorski arī par armijas virspavēlnieku un Bruņoto spēku ģenerālinspektoru. Par viņa vietnieku tika nozīmēts ģenerālis K.~Sosnkovskis, kurš vēlāk kļuva arī par Polijā izveidotās militārās pretošanās organizācijas oficiālo vadītāju. Kara gados PVE labā gandrīz visās Eiropas valstīs darbojās vairāk nekā 1~700~aģentu.

Sākoties karam pret Poliju, \strong{jautājumu par Polijas valsts likteni} Ā.~Hitlers vēl nebija izlēmis. Septembra sākumā, kad Ā.~Hitlers vēl cerēja panākt Rietumvalstu piekāpšanos, viņš neizslēdza iespēju radīt kādu poļu ``pārpalikuma valsti'' (\detxti{Reststaat}). 12.~septembrī J.~f.~Ribentrops sarunā ar izlūkdienesta šefu V.~Kanarisu izteicās, ka tāds risinājums Ā.~Hitleram šķiet vissimpātiskākais, jo tad ar šīs valsts valdību varētu vest sarunas par ``austrumu miera'' noslēgšanu. Taču bez šīs J.~f.~Ribentrops minēja vēl divas alternatīvas: Polijas sadali starp Vāciju un PSRS un vairāku atsevišķu nacionālu valstu, tai skaitā arī Rietumukrainas, izveidi. Turpretī J.~Staļins, kā uzskatīja padomju vēsturniece T.~Grigorjanca, bažījās, ka Ā.~Hitlers varētu noslēgt militāru pamieru ar poļu ``pārpalikuma valsti''. Tieši šīs bažas likušas padomju diktatoram izšķirties par ātrāku, nekā tas bija paredzēts agrāk, padomju karaspēka ievešanu Polijā, kaut V.~Molotovs 10.~septembrī bija paziņojis V.~f.~Šulenburgam, ka ``Sarkanajai armijai lai sagatavotos gājienam uz Poliju, vajag no divām līdz trīs nedēļām''.

Sarkanās armijas iesoļošana Polijā būtiski samazināja iespēju Ā.~Hitleram radīt vairāk vai mazāk nozīmīgu poļu ``pārpalikuma valsti''. Savā runā Dancigā 1939.~gada 19.~septembrī Ā.~Hitlers deklarēja, ka ieņemto poļu teritoriju ``galīgā valstiskuma forma'' pirmkārt ir atkarīga no Vācijas un Krievijas. ``Pilnīgi skaidri redzams,~--- viņš paziņoja,~--- ka Polija tai veidā, kādā tā radās Versaļas līguma rezultātā, nekad vairs neatjaunosies. To garantē arī Krievija.''

Kā rakstīja krievu vēsturnieks Ļ.~Bezimenskis, t.s. ``atbrīvošanas karagājiena'' uz Rietumukrainu un Rietumbaltkrieviju pamatojums ar ukraiņu un baltkrievu aizsardzības interesēm bija loģisks, taču tam bija trūkums~--- ar to nevarēja pamatot Sarkanās armijas virzību Vislas un Narevas virzienā, apgabalos ar viendabīgu poļu iedzīvotāju sastāvu. Tāpēc padomju puse arī nolēma, ka Sarkanajai armijai ir jāapstājas nevis pie Vislas un Varšavas, bet pie Bugas un Brestas. Iespējams, J.~Staļins, citi politiskie un militārie padomju vadītāji, no agrākiem laikiem atceroties poļu naidīgumu pret krieviem, nevēlējās paredzamā kara ar Vāciju gadījumā Sarkanās armijas aizmugurē atstāt naidīgi noskaņotus poļus. Kā par J.~Staļina lēmumu atteikties no bijušajām Krievijai piederējušajām poļu teritorijām uzrakstīja amerikāņu žurnālists V.~Šīrers: ``Labi atceroties Krievijas daudzu gadsimtu vēsturi, viņš saprata, ka poļu tauta nekad nesamierināsies ar savas neatkarības zaudēšanu.'' Taču J.~Staļins nevēlējās arī vienkārši uzdāvināt Ā.~Hitleram poļu teritorijas.

Jau trīspusējās britu-franču-padomju sarunās 1939.~gada augustā PSRS bija piedāvājusi palīdzību iespējamajiem Rietumu sabiedrotajiem caur Polijai piederošo t.s. ``Viļņas koridoru'', kam bija stratēģiska nozīme. Tagad PSRS vadība uzzināja, ka Lietuvas sūtnis Vācijā K.~Škirpa Berlīnē ved sarunas par Lietuvas pāreju Vācijas protektorātā. Šāda Viļņas atdošana Lietuvai celtu tajā Vācijas ietekmi. Jau kara otrajā nedēļā, vēl saglabājoties PSRS nogaidošai nostājai, Vācijā sākās Polijas jautājuma ``galīga risinājuma'' variantu apspriešana. 12.~septembrī apspriedē pie Vācijas bruņoto spēku Galvenās virspavēlniecības štāba priekšnieka V.~Keitela tika izskatīti varianti ar Viļņas rajona nodošanu Lietuvai un Galīcijas un citu Polijas pakļautībā bijušo ukraiņu rajonu pārveidi par ``neatkarīgiem'' veidojumiem. Tas neietilpa J.~Staļina plānos.

Vēl nesen vēstures literatūrā tika izteikts uzskats, ka, tikai ierodoties 28.~septembrī Maskavā, J.~f.~Ribentrops uzzināja par PSRS vadības priekšlikumu mainīt ``interešu sfēru'' robežas. Kā rāda V.~f.~Šulenburga 20.~septembra telegramma Vācijas Ārlietu ministrijai, tad PSRS vadība par to bija izšķīrusies jau 19.~septembrī. V.~Molotovs lika saprast, ka padomju valdība vēlas nekavējoties sākt sarunas par Polijas sadali. Tika ierosināts sarunas noturēt Maskavā, jo no padomju puses tās jāved personām, kuras ir apveltītas ar augstāko varu un nevar atstāt Padomju Savienību. 20.~septembrī sūtnis V.~f.~Šulenburgs ziņoja no Maskavas, ka J.~Staļins ir noteikti pret poļu ``pārpalikuma valsts'' radīšanu. 23.~septembrī J.~f.~Ribentrops ierosināja V.~Molotovam uzsākt sarunas par Polijas teritorijas sadali un griezās pie PSRS valdības ar lūgumu atteikties par labu Lietuvai, kura bija iekļauta Vācijas ``interešu sfērā'', no agrāk Polijas sastāvā esošā Viļņas apgabala, kuram bija jāiekļaujas padomju ``interešu sfērā''. 25.~septembrī J.~Staļins atbildēja ar pretpriekšlikumu un izsauktajam Vācijas sūtnim paziņoja, ka Polijas jautājuma noregulēšanas gaitā ir jāizvairās no visa, kas nākotnē varētu izsaukt rīvēšanos Vācijas un PSRS attiecībās. No šī viedokļa viņš iebilda pret ``patstāvīgas Polijas palieku'' saglabāšanu un ierosināja mainīt agrāk slepenajā protokolā pie neuzbrukšanas līguma fiksētās ``interešu sfēras''. J.~Staļins piedāvāja Vācijai starp Vislu un Bugu esošo teritoriju, kura aptvēra Lubļinas vojevodisti un Varšavas vojevodistes austrumu daļu, kuras iedzīvotāju sastāvs bija izteikti polisks, apmainīt pret Lietuvu, kurai savukārt bija jānonāk padomju ``interešu sfērā''. Tas bija gudrs un tālredzīgs lēmums. Ar to J.~Staļins atbrīvojās no pārmetumiem par etnisko poļu teritoriju sagrābšanu un rūpēm par tur mītošajiem poļu iedzīvotājiem. Vācijas piekrišanas gadījumā, PSRS ar Vācijas atbalstu nekavējoties stātos pie ``Baltijas valstu jautājuma'' risināšanas.

1939.~gada 27./28.~septembra Maskavas sarunu laikā ar J.~f.~Ribentropu J.~Staļins paskaidroja tam, ka Polijas sadale pa Vislu var kļūt par cēloni iespējamām nesaskaņām starp PSRS un Vāciju: Vācija acīmredzot radītu savās poļu teritorijās protektorātu, bet PSRS~--- autonomu ``sociālistisku'' poļu padomju republiku, bet tas, pēc J.~Staļina domām, varētu dot poļiem ieganstu izvirzīt jautājumu par ``atkalapvienošanos'', kas būtu ``pastāvīgu nemieru perēklis Eiropā''. Tāpēc būtu lietderīgāk poļu teritorijas atstāt vienās~--- Vācijas rokās. Toties PSRS ``interešu sfērā'' nonāktu Lietuvas teritorijas. Pēc neilgas apspriešanas Vācijas puse pieņēma priekšlikumu. Rezultātā tika panākta vienošanās par abu valstu ``interešu sfēru'' nodalīšanu nevis pa Vislas, bet Bugas upi. Padomju diktators šoreiz apspēlēja Vācijas diktatoru. Tas nevarēja atteikties no teritoriju apmaiņas, jo Vācijai karam pret Rietumiem bija nepieciešamas izejvielas, kuras vajadzīgajā daudzumā varēja piegādāt tikai PSRS.

1939.~gada 28.~septembrī Vācijas un PSRS pārstāvji Maskavā precizēja robežas līniju, būtiski mainot ``interešu sfēru'' dalījumu. Šo vienošanos abas puses fiksēja slepenā protokolā pie 1939.~gada 28.~septembrī parakstītā Vācijas un PSRS ``Līguma par robežu un draudzību''. Līgumu par draudzību PSRS noslēdza ar valsti, kura veda agresīvu karu pret Poliju un pēc tās sakāves turpināja karot pret Lielbritāniju un Franciju.

Līgumam bija pievienoti trīs protokoli: viens konfidenciāls un divi slepeni.

Ar konfidenciālā protokola parakstīšanu padomju puse apņēmās netraucēt padomju ``interešu sfērā'' dzīvojošo Vācijas pilsoņu un vācu tautībai piederīgo izceļošanu uz Vāciju, kā arī viņu mantas izvešanu. Vācija savukārt tādas pat saistības uzņēmās pret baltkrieviem un ukraiņiem, kuri bija nonākuši Vācijas pārvaldes sfērā. Uz šīs vienošanās pamata it kā varēja notikt brīvprātīga iedzīvotāju izvēle kurā valstī dzīvot. Kā norādījis krievu vēsturnieks V.~Simindejs, situāciju traģisku padarīja tas, ka faktiski nekāda brīvprātīguma nebija. Par kādu brīvprātīgumu gan varēja runāt, ja Polijas patriotam~--- vācietim vai baltkrievam bija jāizvēlas kurā valstī dzīvot~--- hitleriskajā Vācijā vai staļiniskajā PSRS, bet citas alternatīvas vienkārši nebija.

Viens no slepenajiem protokoliem precizēja pašu ``interešu sfēru'' robežas, kuras tagad lielā mērā sakrita ar etniskajām; otrajā bija konstatēts, ka Vācija un PSRS katra savā robežas pusē nepieļaus poļu aģitāciju, vērstu pret otru pusi.

Abu pušu vienošanās apkarot poļu aģitāciju, vērstu pret kādu no pusēm, objektīvi bija vērsta pret veidojošos poļu tautas antinacistisko nacionālās atbrīvošanas kustību, saistīja rokas poļu kreisajiem spēkiem, dodot iespēju pretošanās kustības iniciatīvu pārņemt pilsoniskajām aprindām.

To var saukt arī par Polijas sesto dalīšanu, kaut visa poļu apdzīvotā teritorija faktiski palika Vācijas rokās. PSRS piedalījās Polijas valsts bijušo teritoriju, bet ne poļu zemju sadalē. Ar 28.~septembra līgumu tā ieguva gandrīz visu Rietumukrainu (izņemot ap 40~km platu un 19,5~tūkstošu km$^{2}$ lielu tās rietumu zonu, kura atradās aiz Kērzona līnijas (t.s. \pltxti{Zakerzonie}) un kuru 1939.~gadā ieguva Vācija, bet 1945.~gadā atguva Polija) un Rietumbaltkrieviju, kā arī Vilņas apgabalu, (par tā likteni tiks vēl runāts), bet atteicās par labu Vācijai no poļu apdzīvotajām teritorijām, izveidojot robežu faktiski pa Rietumvalstu iepriekš starptautiski atzīto Kērzona līniju. PSRS pievienoto Polijas austrumu apgabalu (\pltxti{Kresy}) teritorijā 201,015~km$^{2}$ platībā dzīvoja aptuveni 13,3~miljoni cilvēku.

1939.~gada 28.~septembrī noteiktā PSRS un Vācijas robeža atbilda kā jau Parīzes miera konferencē 1919.~gada decembrī un Antantes Augstākās padomes konferencē Spā 1920.~gada jūlijā atbalstītajai Padomju Krievijas un Polijas robežlīnijai kopumā (ar dažām novirzēm), tā arī jau mūsdienu Polijas starptautiski atzītajai austrumu robežai, tāpēc runas par Rietumbaltkrievijas un Rietumukrainas ``netaisnīgu'' atņemšanu Polijai 1939.~gadā ir bez pamata.

Līdz ar 28.~septembra līgumu jautājums par poļu ``pārpalikuma valsti'' pilnībā tika nodots Vācijas rokās. Faktiski līgums nostiprināja Polijas valsts likvidēšanu. Pēdējo tās pieminējumu saturēja Ā.~Hitlera runa reihstāgā 6.~oktobrī, kurā viņš paziņoja, ka Vācijas valdības uzdevums bez ``jaunās kārtības'' radīšanas ir arī ``poļu valsts atjaunošana, kura gan ar savu struktūru, gan vadību garantēs, ka tur nekad neradīsies nedz ugunsgrēka perēklis pret Vāciju, nedz intrigu centrs pret Krieviju un Vāciju''. Īstenībā jautājums jau bija izlemts un nekādi nebija saistāms ar jēdzienu ``valstiskums''.

Kā uzskatīja krievu vēsturnieks V.~Volkovs, aiz daļējās ``interešu sfēru'' izmaiņas 1939.~gada 28.~septembrī stāvēja jaunas PSRS ārpolitiskās koncepcijas izstrādes pirmā posma noslēgums. Šī koncepcija vadījās no mērķa likvidēt ``poļu jautājumu'' valsts iekšienē, kā arī nodibināt stratēģisku kontroli pār Baltiju. Tāda robežjautājuma izšķiršana sakrita ar pirms Pirmā pasaules kara pastāvošo Krievijas impērijas robežu atjaunošanu. (Šeit gan nepieciešams precizējums, Varšava un citas poļu teritorijas, kuras pirms Pirmā pasaules kara Polijas karalistes sastāvā ietilpa Krievijas impērijā, tagad nokļuva Vācijas varā, un PSRS nebūt neieguva visas Krievijas impērijā ietilpstošās zemes.) Arī britu vēsturnieks A.~Bulloks rakstīja, ka ``ceturtā [faktiski sestā~--- V.~Š.] Polijas dalīšana atļāva Krievijai atgūt bijušās krievu teritorijas, ko poļi anektēja 1920.~gadā''.

Pēdējais PSKP CK ģenerālsekretārs M.~Gorbačovs pēcvārdā brošūrai ar viņa un poļu inteliģences pārstāvju tikšanās materiāliem 1988.~gada 14.~jūlijā t.s. Molotova-Ribentropa paktu vērtēja kā ``neizbēgamu'', taču tai pat laikā asi nosodīja 28.~septembra līgumu. Pēdējais tika raksturots kā ``politiska kļūda ar smagām sekām'' kā Padomju Savienībai, tā arī citām valstīm, tāpat arī pasaules komunistiskajai kustībai, kā ``tieša un izaicinoša atkāpšanas no ļeņinisma, atteikšanās no tā principiem''. Tamlīdzīgus uzskatus pauduši arī vēsturnieki.

Jaunākajā krievu historiogrāfijā padomju~--- vācu 28.~septembra līgums parasti tiek vērtēts klaji negatīvi, kā politiska kļūda. Atsevišķi krievu vēsturnieki, uzsverot, ka PSRS ``Līgumu par robežu un draudzību'' noslēdza ar karojošu valsti, uzskata, ka ar to PSRS atgāja no neitralitātes un uzsāka sadarbību ar Vāciju. A.~Bogaturovs atzīst, ka ar to ``sadarbība starp Maskavu un Berlīni tika noformēta ar pilnvērtīgu starpvalstu līgumu.'' Arī V.~Semirjaga uzskatīja, ka PSRS ar to faktiski iesaistījās militāri politiskā savienībā ar Vāciju. Savukārt pēc A.~Ņekriča viedokļa padomju~--- vācu attiecības no 1939.~gada septembra līdz 1940.~gada novembrim (V.~Molotova vizītei Berlīnē) veidoja ``it kā nepabeigtu militāri politisku savienību''. (Tiesa, A.~Ņekričs atzina arī līguma pozitīvās sekas: to, ka robeža tika noteikta pa ``Kērzona līniju'', ka PSRS ieguva Rietumukrainu un Rietumbaltkrieviju, rīcības brīvību Baltijā, radīja barjeru Vācijas agresijas priekšā Austrumeiropā.) Kā rakstījuši krievu vēsturnieki G.~Kumaņevs un E.~Škļars, slepeno protokolu pie 23.~augusta un 28.~septembra līgumiem, parakstīšana ar hitlerisko Vāciju nevar tikt attaisnota. Šie dokumenti pārsvītrojuši morālos principus, uz kuriem pēc 1917.~gada Oktobra revolūcijas, lai arī formāli, balstījās PSRS ārējā politika.

Par 1939.~gada notikumiem, 17.~septembrī uzsākto Sarkanās armijas ``atbrīvošanas karagājienu'' un 28.~septembra Vācijas-PSRS līgumu vēsturnieku strīdi turpinās arī mūsdienās.

2009.~gada 23.~septembrī Polijas Seims pieņēma īpašu rezolūciju, nosodot Sarkanās armijas 1939.~gada 17.~septembra iebrukumu Rietumukrainā un Rietumbaltkrievijā: ``Tad sākās puses Polijas padomju okupācija un pret Polijas pilsoņiem vērsts genocīds \dots{}''~--- Tāds ir mūsdienu Polijas oficiālais redzējums. Taču var pastāvēt arī citi~--- oficiāli valstiski un neoficiāli redzējumi.

Piemēram, mūsdienu krievu vēsturnieks J.~Pavļuks savā 2006.~gadā aizstāvētā kandidāta disertācijā rakstīja, ka ``tas, kas notika 1939.~gadā, vēsturiskā nozīmē bija nevis PSRS agresija pret Poliju, bet poļu agresijas seku likvidēšana''. Turpat viņš arī atzinis, ka ``uz Ukrainas vēsturiskās teritorijas atjaunošanas pamata tika noņemts jautājums par ceturto Polijas dalīšanu, jo no Polijas tika atdalīta teritorija, tai nepiederoša un par okupāciju runa nevarēja būt.''

Pēc autora domām, sadarbība starp Maskavu un Berlīni gan tika noformēta ar starpvalstu līgumu, taču par PSRS militāri politisku ``savienību'' ar Vāciju nevar būt runas, jo:

\begin{enumerate}
\item abi partneri neuzticējās viens otram, rūpīgi raudzījās, lai viena puse nepārkāptu otrās valstiskās intereses;

\item Vācija oficiāli ievēroja neitralitāti PSRS karā pret Somiju, savukārt, PSRS~--- Vācijas karos pret Franciju, Lielbritāniju u.c. valstīm. Tiesa, šī PSRS ``neitralitāte'' daudzos gadījumos faktiski tika pārkāpta, PSRS sniedzot Vācijai militāri-politiska un militāri-ekonomiska rakstura pakalpojumus (tāpat kā, piemēram, ASV Pirmā pasaules kara sākumposmā formāli bija neitrālas, bet sniedza ekonomisku un arī militāru palīdzību Antantes valstīm). Tomēr J.~Staļinam izdevās līdz galam nesabojāt attiecības ar Rietumvalstīm.

\item To, ka šāda ``savienība'' nepastāvēja, liecina arī Lielbritānijas valdības mēģinājumi pēc Francijas sakāves stiprināt savas attiecības ar PSRS. Kad 1940.~gada vidū PSRS ieradās Lielbritānijas sūtnis S.~Kripss, J.~Staļins viņam paziņoja, ka ``mums nav bloka ar Vāciju, kas paredzētu karu pret Angliju. Mums ir tikai pakts par neuzbrukšanu''. Protams, J.~Staļins ``aizmirsa'' līgumu ar Vāciju par ``draudzību'', taču jāatzīst, ka ``draudzība'' var būt dažāda, tā nenozīmē militāri politisku savienību. Zināms taču, ka it kā pastāvošā Vācijas un PSRS ``savienība'' beidzās ar Vācijas uzbrukumu PSRS.
\end{enumerate}

Par ``morālajiem'', ``ļeņiniskajiem internacionālisma'' un citiem padomju laikā cēliem epitetiem apveltītiem principiem PSRS iekšējā un ārējā politikā, kurus piemin daži vēsturnieki, jāsaka, ka, tāpat kā Rietumvalstis, arī padomju valsts tos neievēroja jau V.~Ļeņina dzīves laikā. Uzskatāms piemērs tam bija Sarkanās armijas iebrukums poļu apdzīvotajās teritorijās 1920.~gadā.

Taču jāsaka arī kas cits. Tas, ka Vācija un PSRS kopīgi cīnījās pret poļu atbrīvošanās kustību, tikai atkārtoja agrākās imperiālistisko valstu vienošanās par koloniju sadali. Piemēram, kā 1886.~gadā Lielbritānija noslēdza līgumu ar Vāciju par ietekmes sfēru sadali Klusā okeāna rietumu daļā, tā 1912.~gadā Francija ar Spāniju dalīja Ziemeļmaroku utt. Neviena no nosauktajām valstīm neievēroja sadalāmo teritoriju iedzīvotāju intereses, tāpat kā to nedarīja Polijas dalītājvalstis XVIII un XIX gadsimtā, tāpat kā to Polija nedarīja ar iekarotajām teritorijām XX gadsimtā. Bija tikai dabīgi, ka kārtējo reizi sadalot Poliju, tās sadalītājvalstis vienojās vērsties pret poļu nacionālās atbrīvošanās kustību. Jā, PSRS līdz ar to traucēja poļu kreisajiem spēkiem pārtvert iniciatīvu cīņā pret vācu okupantiem, taču no PSRS vadītāju viedokļa tai brīdī labu attiecību uzturēšana ar Vāciju bija svarīgāka nekā nākotnē paredzamo sabiedroto poļu kreiso atbalstīšana. Vēsture, tai skaitā mūsdienu, dod daudz piemēru, kad tuvāko interešu vārdā tiek aizmirsts par morāles normām, principiem un pat savām interesēm tālākā perspektīvā.

Bez tam, ja raugās no ``ļeņinisma principu'' viedokļa, tad jānorobežo PSRS mēģinājums palikt ārpus kapitālistisko valstu konflikta, kuru nevar nosodīt, tūlīt blakus nenosodot Rietumvalstu piekopto agresora nomierināšanas politiku, no konkrētajām politiskās darbības metodēm, kurus pielietoja staļiniskā Padomju Savienības vadība. Lietas būtība nebija tik daudz slepenajos protokolos un ``interešu sfēru'' noteikšanā un pārdalīšanā (bez tām 23.~augusta neuzbrukšanas līgums Padomju Savienībai zaudētu nozīmi un nevarētu kalpot kā ārpolitisks instruments), cik apstāklī, ka J.~Staļina politika praksē pārkāpa līgumos noteiktos ietvarus. Netika ievērotas PSRS anektēto teritoriju iedzīvotāju cilvēktiesības. PSRS faktiski neitralizēja to ekonomisko blokādi, kuru Rietumvalstis veidoja ap Vāciju. Tika uzsākta abu valstu drošības dienestu sadarbība, Vācijai tika izdoti vācu komunisti-antifašisti, arestēti staļiniskā terora uzplūdu gados. Nosauktie soļi nu nekādi neatbilda ``ļeņinisma principiem'', toties līdzinājās citu valstu piekoptajai politiskajai praksei.

Bez tam izvirzās jautājums~--- cik ilgi PSRS varētu pastāvēt, ja tā balstītos tikai uz tās deklarētajiem marksisma un ``ļeņinisma principiem''? Tāpat kā citas, arī tā saucamās demokrātiskās, valstis vajadzības brīdī ievēroja tikai lietderības ``principu'', tā arī J.~Staļins atmeta ``šķiriskos'' postulātus (arī 1939.~gada septembrī), lai nodrošinātu savas valsts intereses. Rīcība morāli nosodāma, bet saprotama. Ar ko gan daudz atšķīrās J.~Staļina rīcība attiecībā pret Rietumukrainu un Rietumbaltkrieviju no J.~Pilsudska rīcības pret Viļņu 1920.~gadā? Gan vienā, gan otrā gadījumā bija līgumu laušana, starptautisko tiesību neievērošana, ar skaistiem, patriotiskiem saukļiem aizmaskota vardarbība. Abos gadījumos cietušie sauca pēc palīdzības, abos gadījumos lielvalstu dienas kārtībā bija citas, svarīgākas problēmas nekā morāles aizstāvība, taisnīguma atjaunošana.

PVE griezās pie valstīm, kuras atradās kara stāvoklī ar Vāciju, un arī neitrālajām valstīm ar protesta notu, pieprasot nosodīt Vācijas un Padomju Savienības 28.~septembra līgumu, pasludināja par kopējās cīņas ar sabiedrotajiem turpināšanu par valsts pilnīgas neatkarības un suverenitātes atjaunošanu, uzsāka poļu armijas izveidi ārzemēs. Taču Rietumvalstu politiķu attieksme pret poļu aktivitātēm bija faktiski noraidoša.

Britu politiķi PSRS vadības rīcību 1939.~gada septembrī novērtēja kā Lielbritānijas interesēm atbilstošu un izrietošu no Krievijas vēsturiskajām tiesībām. Lielbritānijas kara kabinetam 25.~septembrī sagatavotā memorandā V.~Čērčils rakstīja: ``Kaut krievi ir vainīgi rupjā vārda laušanā nesenajās sarunās [domātas sarunas 1939.~gada augustā Maskavā], taču maršala Vorošilova prasība, saskaņā ar kuru krievu armijām, ja tās būtu Polijas sabiedrotās, būtu jāieņem Viļņa un Ļvova, bija pilnībā mērķtiecīga militāra prasība. To noraidīja Polija, kuras argumentus, neraugoties uz visu to dabīgumu, patreizējo notikumu gaismā nevar uzskatīt par apmierinošiem. Rezultātā Krievija kā Polijas ienaidnieks ieņēma tās pat pozīcijas, kura tā varētu ieņemt kā apšaubāms un aizdomās turams draugs. Starpība nemaz nav tik liela, kā tas varētu likties. Krievi ir mobilizējuši ļoti lielus spēkus un ir parādījuši, ka viņi var ātri un tālu pavirzīties no savām pirmskara pozīcijām. Tagad viņi robežo ar Vāciju un pēdējai nav nekādas iespējas atkailināt Austrumu fronti. Tās novērošanai nāksies atstāt lielu vācu armiju. Cik man ir zināms, ģenerālis Gamelēns [Francijas armijas virspavēlnieks Otrā pasaules kara sākumā~--- V.Š.] tās skaitu nosaka kā mazākais 20~divīzijas, bet tās var būt arī 25 un pat vairāk. Tāpēc Austrumu fronte potenciāli pastāv.'' Protams, V.~Čērčils nevarēja paredzēt notikumus divus gadus uz priekšu, bet pastāvošo situāciju viņš vērtēja reālistiski.

To pašu domu Admiralitātes Pirmais lords atkārtoja arī savā runā radio 1.~oktobrī: ``Krievija realizē aukstu savu pašas interešu politiku. Mēs dotu priekšroku tam, lai krievu armijas stāvētu savās tagadējās pozīcijās kā Polijas draugi un sabiedrotie, bet ne kā iekarotāji. Taču Krievijas drošībai pret nacistu apdraudējumu bija nepieciešams, lai krievu armijas stāvētu šajā līnijā. Lai kā arī nebūtu, šī līnija pastāv, ir radīta Austrumu fronte, kurai nacistiskā Vācija neuzdrošināsies uzbrukt. Kad Ribentropa kungs pagājušajā nedēļa bija izsaukts uz Maskavu, viņam nācās uzzināt un pieņemt to faktu, ka nacistu plānu īstenošanu pret Baltijas valstīm un Ukrainu ir galīgi jāpārtrauc \citespace{} Es nevaru pareģot, kā Krievija rīkosies. Tā ir mīkla, kuru ir ļoti grūti atminēt, taču atslēga tai ir. Šī atslēga ir Krievijas nacionālās intereses. Ievērojot drošības apsvērumus, Krievija nevar būt ieinteresēta, lai Vācija nostiprinās Melnās jūras krastos vai lai tā okupē Balkānu valstis un pakļauj Dienvid-Austrumu Eiropas slāvu tautas. Tas būtu pretrunā ar vēsturiski izveidojušamies Krievijas dzīvotspējas interesēm.''

Šoreiz V.~Čērčilu atbalstīja arī tādi politiķi, kuri 1939.~gada augustā vēl nostājās pret viņu. Tā, 4.~oktobrī, uzstājoties Lordu palātā un komentējot Rietumbaltkrievijas un Rietumukrainas pievienošanu PSRS, lords E.~Halifakss paziņoja: ``Vismazāk es vēlētos aizstāvēt padomju valdības rīcību tai brīdī, kad tā to veic, taču būs tikai taisnīgi atgādināt divus apstākļus: pirmkārt, padomju valdība nekad neuzsāktu tādu rīcību, ja nesāktu un, bez kara pieteikuma ielaužoties Polijā, piemēru nerādītu vācu valdība; otrkārt, ir jāatgādina, ka padomju valdības rīcība veda pie robežas pārnešanas pēc būtības līdz tai līnijai, kuru Versaļas konferences laikā ieteica lords Kērzons. Es netaisos aizstāvēt Padomju vai kādas citas valdības, izņemot savējās, rīcību. Es tikai minu vēsturiskus faktus, un uzskatu, ka tie nav apstrīdami.'' Arī Francijas vadība deva V.~Sikorskim skaidri saprast, ka tā netaisās iesniegt PSRS protestu sakarā ar tās robežas izmaiņām rietumos.

Faktiski ar 1939.~gada 17.~septembri, kad Sarkanā armija iegāja Polijas teritorijā, Rietumbaltkrievijas un Rietumukrainas iedzīvotāju dzīve asociējās ar viņu apvienošanos vienotā valstī ar saviem tautiešiem Baltkrievijas un Ukrainas PSR. 1939.~gadā vienīgā baltkrievu un ukraiņu atkalapvienošanās alternatīva bija to šo tautu daļu, kuras dzīvoja Polijas teritorijā, pāriešana no poļu okupācijas varas Vācijas okupācijas varā un labvēlīgu pozīciju nodrošināšana vērmahtam, lai tas varētu uzbrukt PSRS. Gan šīm tautām, gan Padomju Savienībai tas bija nepieņemami.

Piemēram, Padomju Baltkrievijas iedzīvotāji kopš 1945.~gada atzīmēja ``Baltkrievijas tautas apvienošanās dienu vienā valstī.'' Pēc PSRS sabrukuma šie svētki vairs oficiāli netika svinēti, taču daudzas sabiedriskās organizācijas iestājās par to atjaunošanu. Rezultātā 1999.~gada septembrī Baltkrievija atzīmēja Rietumbaltkrievijas ``atkalapvienošanās'' ar BPSR 60.~gadadienu, bet 2009.~gadā~--- 70.~gadadienu. Daudzās Rietumbaltkrievijas apdzīvotās vietās arī šodien ir sava ``17.~septembra iela''. 2012.~gadā 17.~septembrī baltkrievu publicists J.~Semaško, atzīstot, ka pēc 1939.~gada 17.~septembra sekoja represijas, kurās cieta daudzi cilvēki, arī baltkrievi, un daudziem no tiem, kuri gadu desmitiem bija cīnījušies par nacionālās atbrīvošanas ideāliem, tas kļuva par īstenu traģēdiju, tomēr uzsvēra, ka, neraugoties uz to, nedrīkst aizmirst, ka tautas apvienošanās ir augstākais vēsturiskā taisnīguma akts. Tas darīja galu pretdabiskajam stāvoklim, kad viena nācija bija sadalīta divās daļās. Līdz ar apvienošanos tika ievērotas tautas pamatintereses, saglabāta tās etniskā un teritoriālā vienotība, noteikta tās tagadne un nākotne. Šim vērtējumam nākas piekrist. No 2021.~gada Baltkrievijā 17.septembris, kad Sarkanā armija iegāja Polijas austrumu provincēs, kā rezultātā Baltkrievijas austrumi un rietumi apvienojās Baltkrievijas PSR sastāvā, tiek svinēts kā Nacionālās vienotības diena.

Krievu vēsturniece A.~Noskova uzsvērusi, ka ar PSRS teritorijas ``pavirzīšanu'' uz rietumiem, padomju vadība risināja arī ļoti sāpīgo ``ukraiņu jautājumu''. 1939.--1940.~gadā notikusī Rietumukrainas, Ziemeļbukovīnas un Besarābijas pievienošana PSRS ļāva izveidot Ukrainas PSR, kurā bija apvienoti gandrīz visi ukraiņu apdzīvotie apgabali (ārpus tās pagaidām palika tikai Aizkarpatu Ukraina. J.~Staļins piespieda Čehoslovākiju to atdot Padomju Savienībai 1945.~gadā). Tā kopojot gandrīz visas ukraiņu (protams, kopā ar blakus tiem dzīvojošajiem citu tautību pārstāvjiem) apdzīvotās zemes vienas Ukrainas PSR sastāvā, padomju vadība atņēma aktīvi provāciski un pretpadomiski noskaņotajai ukraiņu emigrācijai tik pievilcīgu mērķi kā ``vienotas Ukrainas'' radīšana.

Saprotams, ka ukraiņu nacionālistiskie politiķi pauda savu neapmierinātību ar jauno situāciju, ar jauno PSRS robežu. Tā, Ukraiņu Nacionālistu organizācijas vadītāji izteica neapmierinātību ar to, ka pēc 1939.~gada 28.~septembra Vācijas un Padomju Savienības līguma Galīcija nonāca PSRS rokās. Arī par Ukrainas PSR būtisku palielināšanu nacionālisti nekādu prieku neizrādīja, jo padomju iekārtu uzskatīja par naidīgu ukraiņu mērķiem. Vācijas amatpersonas gan viņus mierināja, ka drīz stāvoklis var mainīties to labā.

Rietumvalstu politiķi vairāk ieklausījās poļu nekā Vāciju atbalstošo ukraiņu viedoklī. Taču jau 12.oktobrī, pēc N.~Čemberlena runas parlamentā, kurā viņš paziņoja, ka Polijas valstiskumam jātiek atjaunotam, ko deputāti sagaidīja ar ovācijām, daudzi no viņiem pēc tam kuluāros padomju sūtnim I.~Maiskim atklāti izteicās, ka neviens Anglijā nopietni nedomā par Padomju Savienības rokās nonākušo apgabalu iekļaušanu Polijā.

Polijas valdības emigrācijā ārlietu ministram A.~Zaleskim, kurš 11.--12.~oktobrī veda sarunas Londonā, Lielbritānijas ārlietu ministrs E.~Halifakss un premjerministrs N.~Čemberlens paziņoja, ka poļi nevar rēķināties ar Lielbritānijas atbalstu teritoriju atgūšanā, kuras tika pievienotas PSRS. A.~Zaleskim tika atgādināts, ka 1919.~gadā Lielbritānijas piedāvātā t.s. Kērzona līnija nodalīja no Polijas zemes, kuras apdzīvoja pārsvarā ukraiņi un baltkrievi. (Briti atzina poļu prasību pamatotību tikai uz nelielu agrākās Polijas daļu uz rietumiem no Kērzona līnijas~--- Suvalku rajonu, kurš atrodoties uz ziemeļaustrumiem no Varšavas, piekļāvās Austrumprūsijai. Sadalot Polijas teritoriju, Vācijas vadība, atsaucoties uz to, ka Suvalku rajons no 1283.~līdz 1422.~gadam ietilpa Teitoņu ordenī, panāca tā iedalīšanu Vācijai un ar Sudāvijas (\detxti{Sudauen}) nosaukumu pievienoja Austrumprūsijai. Pēc Lielbritānijas domām, šī teritorija varēja tikt atdota Polijai pēc sarunām pēc kara paredzamajā miera konferencē.)

Starptautiski Polija bija nonākusi sarežģītā, neapskaužamā situācijā, jo tās vadība 17.~septembrī oficiāli nebija paziņojusi par kara stāvokli ar RSRS, nebija griezusies Tautu Savienībā, bet bēgusi uz Rumāniju. Visas PVE sūdzības bija vērstas Rietumu sabiedrotajiem, kuriem militārās palīdzības problēma Polijai, vēl jo īpaši pēc tās teritorijas sadales, nebija aktuāla. Svarīgāk bija radīt labvēlīgu situāciju starptautiskajā arēnā. Tāpēc, kā bija teikts Polijas sūtņa Lielbritānijā E.~Račinska 12.~oktobra ziņojumā par PVE ārlietu ministra A.~Zaleska sarunām ar Lielbritānijas valstsvīriem, ``gan Halifakss, gan Čemberlens pastiprināti uzsvēra, ka tiešajam mērķim, uz kuru patreiz tiecas sabiedrotie,~--- uzvarai pār Vāciju~--- jābūt arī poļu tuvākajam mērķim. Dotajā brīdī ir jāizvairās no jebkādām darbībām pret Padomju Krieviju, jo Krievija vēl nav pārgājusi Vācijas nometnē un pastāv cerība noturēt to neitralitātes stāvoklī. Zaleskim tika paziņots, ka poļiem jāstāv vienotā ierindā kopā ar sabiedrotajiem \dots{}''. Pārējo problēmu risināšana tika atlikta līdz miera konferencei nākotnē. E.~Halifakss lordu palātā paziņoja, ka ``vajag redzēt atšķirību starp Vācijas un PSRS rīcību Polijā.'' Gan Lielbritānijas, gan Francijas valdības neieteica Polijai saasināt attiecības ar Padomju Savienību. Arī ASV Valsts sekretārs K.~Halls, paziņoja, ka ne prezidents F.~D.~Rūzvelts, ne viņš negribētu saskatīt Krievijā valsti, ``karojošu tādā pat mērā kā Vācija, jo, to darot, mēs grūstu Krieviju dziļāk Hitlera apkampienos''. Šāda nostāja pilnībā atbilda Lielbritānijas un Polijas līgumam par savstarpējo palīdzību, kura slepenais protokols izslēdza līguma darbību Padomju Savienības uzbrukuma Polijai gadījumā. Arī J.~Staļins, ieguvis to, ko viņš tajā brīdī uzskatīja par optimālo, realizēja tādu politiku, kura ļāva Rietumiem nesaskatīt viņā Ā.~Hitlera sabiedroto, ja arī konkrētajā brīdī~--- attiecībā pret Poliju~--- abi diktatori rīkojās saskaņoti.

Šai sakarā varētu kārtējo reizi runāt par Rietumu valstu divkosību, sabiedrotās Polijas nodevību, taču šīs konkrētās situācijas izvērtējumā autors negribētu nostāties moralizētāja pozā un nosodīt Rietumvalstu rīcību. Minhenes vienošanās un Polijas pamešana likteņa varā uzbrūkošās Vācijas priekšā bija Rietumu valstu tieša nodevība, ar ka tās zaudēja iespējamos reālos sabiedrotos (Čehoslovākiju un Poliju) cīņā pret nacistisko Vāciju. Turpretī 1939.~gada oktobrī Lielbritānija un Francija rīkojās tiešām savas (un arī citu valstu, tai skaitā Polijas) nākotnes vārdā, pragmātiski nesaasinot attiecības ar Padomju Savienību, neatmetot iespēju nākotnē ar to kopā karot pret nacistiem.

Acīmredzot ņemot vērā Lielbritānijas un Polijas līgumam par savstarpējo palīdzību un tā slepeno protokolu, Polijas premjerministrs V.~Sikorskis pieņēma sabiedroto piedāvāto pozīciju. Jāatzīst V.~Sikorska skaidrais prāts un tālredzība, kad viņš piekrita Rietumu sabiedroto argumentācijai. 18.~oktobrī, griežoties pie poļu tautas, V.~Sikorskis un PVE uzsvēra nepieciešamību piedalīties sabiedroto vestajā antihitleriskajā cīņā, veidot valstu, tai skaitā slāvu valstu, kopdarbību Vācijas uzbrukuma atvairīšanā Austrumos, ``nodalīt Vāciju no Krievijas'' un izveidot stingru līdzsvaru Eiropā.

17.~oktobrī PSRS sūtnis Londonā I.~Maiskis, atsaucoties uz Lielbritānijas ārlietu ministra vietnieku R.~Batleru, ziņoja uz Maskavu, ka Polijas valdības emigrācijā ārlietu ministrs A.~Zaleskis nepretendē uz to, lai Polijai tiktu atdota Rietumukraina un Pietumbaltkrievija, toties pieprasa Austrumprūsiju (Lūk, kad Otrā pasaules kara laikā pirmoreiz izskanēja ideja palielināt Poliju uz Vācijas rēķina!), apgalvojot, ka citādi Polijas stāvoklis būtu ``neiespējams''. Pēc R.~Batlera vārdiem, Anglijas valdības aprindās uzskata, ka ``nevar būt nekāda jautājuma par Rietumukrainas un Pietumbaltkrievijas atdošanu Polijai. Ja rastos iespēja radīt pieticīga lieluma etnogrāfisku Poliju (kā Napoleona laiku ``Varšavas hercogiste'') ne tikai ar PSRS un Vācijas, bet arī ar Anglijas un Francijas garantijām, tad britu valdība uzskatītu sevi par pilnā mērā apmierinātu.'' Līdzīgi 27.~oktobrī izteicās N.~Čemberlena padomnieks H.~Vilsons.

Arī ārlietu ministrs E.~Halifakss 26.~oktobrī savā runā Lordu palātā konstatēja, ka jaunā PSRS robeža rietumos kopumā atbilst Kērzona līnijai, un tādejādi netiešā formā pateica, ka Lielbritānijai nav iebildumu pret to, ka PSRS ir ieņēmusi Rietumukrainu un Rietumbaltkrieviju.

PSRS Augstākās Padomes sēdē 1939.~gada 31.~oktobrī, V.~Molotovs paziņoja: ``Par vecās Polijas atjaunošanu, kas ir visiem skaidrs, nevar būt ne runas. Tāpēc bezjēdzīga ir tagadējā kara turpināšana ar bijušās Polijas valsts atjaunošanas lozungu''. Tālāk sekoja secinājums, ka Vācija cenšas ātrāk beigt karu un noslēgt mieru, bet Anglija un Francija ``uzstājas par kara turpināšanu un pret miera noslēgšanu''. Padomju politika tika tēlota kā ``atbalsts Vācijai tās miera centienos'' un reizē neitralitātes ievērošana tās vestajā karā. Pēdējo secinājumu var droši vērtēt kā liekulību, bet uzstāšanās pret ``vecās Polijas atjaunošanu'' bija pietiekami skaidra.

Kad novembra beigās notika V~.Sikorska un A.~Zaleska vizīte Londonā ar uzdevumiem vienoties par apgādes jautājumu nokārtošanu Francijā veidojamajām Polijas armijas daļām, atvieglot poļu divīzijas formēšanos Kanādā un iegūt Lielbritānijas valdības solījumu par Polijas agrākajās robežās atjaunošanu, pēdējais tā arī palika nesasniegts. Atbildot uz poļu pārstāvju priekšlikumu, E.~Halifakss viņiem ``nolasīja lekciju'' par Kērzona līnijas priekšrocībām, uzsvēra jautājuma par Polijas robežām izvirzīšanas nesavlaicīgumu. Avīzes ``\entxti{The Sunday Times}'' (Svētdienas Laiks) slejās britu liberāļu līderis D.~Loid Dzordžs apsveica Sarkanās armijas rīcību, nosodīja Polijas valdības imperiālismu. Kad Polijas sūtniecība Londonā nolēma dot uz šo rakstu oficiālu atbildi, to atteicās publicēt pat pusoficiālā ``\entxti{The Times}''.

Arī mūsdienās britu pētniece K.~Svorda, analizējot Lielbritānijas reakciju uz Padomju Savienības rīcību Austrumpolijā 1939.~gadā, uzskata, ka 1939.--1941.~gadā noformējās Lielbritānijas politikas pamati kara periodam. Tika atrisināts jautājums, kas gan ir britu galvenais ienaidnieks. Tā izrādījās Vācija. Britu salām radušies draudi virzīja V.~Čērčilu u.c. vadīties no secinājumiem, kas bija aprēķina pilni. Vadoties no šādiem apsvērumiem tika arī būvētas attiecības ar Polijas valdību emigrācijā.

Nevar pierādīt, ka ar Sarkanās armijas ieiešanu Rietumukrainā un Rietumbaltkrievijā, ar PSRS naidīgās Polijas iznīcināšanu PSRS reizē bija nostājusies arī pret Rietumvalstīm, apdraudot visu pastāvošo spēku samēru pasaulē. Iepriekš citētie britu politiķi to tā nevērtēja. Gluži pretēji, Vācijas-PSRS robežu viņi vērtēja kā potenciālu Austrumu fronti. Kā vēlāk izrādījās, viņiem bija taisnība. Jāsecina, ka apstākļos, kad PSRS vēl nebija ierauta karā, tālredzīgākie Rietumu politiķi saprata Rietumukrainas un Rietumbaltkrievijas atgriešanas prasības Polijai bezperspektivitāti un, iespējams, arī šādas prasības neleģitimitāti.

Tikai tad, kad PSRS Vācijas iebrukuma rezultātā bija nonākusi milzu grūtībās, poļu labējiem politiķiem atkal radās cerības vēlreiz atkārtot 1920.--1921.~gada situāciju, ar spēku sašķeļot ukraiņu un baltkrievu tautas, daļu to apdzīvotās teritorijas atkal pakļaujot Polijai.

Kā 1975.~gadā konstatēja angļu vēsturnieks A.~Teilors, 1939.~gadā Lielbritānijas Ārlietu ministrija norādīja, ka britu valdība, 1920.~gadā nospraužot ``Kērzona līniju'', tolaik (1920) uzskatīja ``tās teritorijas, kuras tagad [1939] ieņem Sarkanā armija'', ir ``likumīgi piederošas krieviem'', bet ``turpmāk nebija izdevīga gadījuma, lai atzītu Padomju Krievijas (1939) uzbrukumu par likumīgu, un \citespace{} šis jautājums pastāvīgi apgrūtināja attiecības starp Rietumvalstīm un Padomju Krieviju''. Jāpiezīmē, ka, Rietumvalstis šādu ``izdevīgu gadījumu'', lai atzītu PSRS rīcību 1939.~gadā par likumīgu, nebija ieinteresētas ``atrast'', jo to attiecības ar PSRS gan Otrā pasaules kara gados, gan pēc tā nekad nebija vienkāršas un ne līdz galam atrisināti jautājumi vienmēr varēja noderēt politiskajos strīdos. Arī mūsdienās, kad vairs nepastāv PSRS, šie strīdi un neatrisinātie jautājumi, daudziem politiķiem un arī vēsturniekiem vēl joprojām ļauj runāt par agresīvo ``Polijas sadalīšanu'' 1939.~gadā.

Pie tam Rietumvalstis līdzīgos gadījumos nebūt nerīkojās citādāk kā Padomju Savienība. Piemēram, XIX gadsimtā (1846--1848) ASV varmācīgi atņēma Meksikai gandrīz pusi tās teritorijās (tagadējos ASV Kalifornijas, Teksasas, Ņūmeksikas, Arizonas, Nevadas, Jutas štatus un daļu Vaiomingas štata), taču amerikāņi nekautrējas runāt par krievu ``imperiālismu''. Kā atzinis A.~Teilors, ``Krievijas tiesības uz Baltijas valstīm un Polijas austrumu daļu bija daudz pamatotākas nekā [Amerikas] Savienotajām Valstīm uz Ņūmeksiku'', bet angļi un amerikāņi ``piemēroja pret krieviem normas, kuras viņi neattiecināja uz sevi.''

Kaut valdošā t.s. pasaules sabiedriskā doma ir nosodījusi Padomju Savienības vadības līgumu ar nacistiem, nevar tāpēc likt vienlīdzības zīmi starp abām līgumslēdzējvalstīm. Abas tās neņēma vērā citās zemēs dzīvojošo tautu intereses. Taču PSRS rīcība vismaz daļēji atbilda baltkrievu, ukraiņu un lietuviešu tautu interesēm. Kā rakstīja krievu diplomāts un vēsturnieks V.~Faļins, viņš poļiem un citiem, kuri apvaino Padomju Savienību bijušo Polijas teritoriju sagrābšanā, vienmēr uzdodeva jautājumu: ar ko Ukrainas un Baltkrievijas sašķelšana bija labāka vai sliktāka par Polijas sašķelšanu? Atbildes nekad nebija.

Nacisti kaldināja plānus tālākai pasaules iekarošanai, PSRS vadītāji, vīlušies kolektīvās drošības sistēmas izveides mēģinājumos, būvēja ap savu valsti drošības joslas no jaunām teritorijām. Mūsdienās gan diezgan plašu izplatību ir guvis uzskats, ka J.~Staļins 1941.~gadā gatavoja PSRS uzbrukumu Vācijai.(Par to ir rakstījuši vācu publicists G.~Gilesens, vācu vēsturnieki J.~Hofmans, un V.~Mazers, austriešu sociologs E.~Topičs. Pēdējos gadu desmitos ļoti populāras, bet arī daudz kritizētas ir kļuvušas bijušā padomju izlūka, pārbēdzēja uz Lielbritāniju V.~Suvorova grāmatas, veltītas šai tēmai.) Diemžēl mērķis atmaskot J.~Staļinu, pierakstot viņam visus iespējamos noziegumus, ir novedis pie tendences attaisnot Ā.~Hitleru par uzbrukumu PSRS. Jautājums par to, vai J.~Staļins 1939.~gadā plānoja ``pasaules revolūcijas'' uzvaru, bet 1941.~gadā jau gatavojās uzbrukt Vācijai, vēstures zinātnē nav atrisināts un autors nevēlas ar to spekulēt.

Autora viedoklis: ir jānosoda kā Polijas agresija pret Padomju Krieviju 1919--1920.~gadā, tā padomju Krievijas agresija pret Poliju 1920.~gadā un PSRS agresija pret to 1939.~gadā. Abas valstis ievēroja tikai savas intereses, nerēķinoties ar to rīcības skartajām mazākajām tautām. Taču nosodīt tikai vienu pusi nav nedz taisnīgi, nedz zinātniski.

Jāņem vērā vēl kāds aspekts. PSRS gan noslēdzot 23.~augusta līgumu, gan 17.~septembrī ievedot karaspēku līdz tam Polijai piederošajās teritorijās, gan pievienojot šīs teritorijas PSRS, rīkojās savu drošības interešu vadīta. To var salīdzināt ar ASV rīcību 1962.~gada Karību krīzes laikā, kad ASV vadība ar prezidentu D.~Kenediju priekšgalā uzskatīja, ka tiek apdraudēta valsts drošība, un rīkojās neievērojot nekādus starptautiskos noteikumus, pārkāpjot valstu suverenitāti. Kā liecina vēlākie notikumi, PSRS vadībai 1939.~gadā bija daudz vairāk pamata bažīties par savas valsts drošību. nekā ASV vadībai 1962.~gadā, kaut ASV rīcība 1962.~gadā nekad nav bijusi pakļauta tādai kritikai kā PSRS rīcība 1939.~gadā.

Pēc 28.~septembra līguma un slepeno protokolu pie tā parakstīšanas palika neatrisināts jautājums par līdz Otrajam pasaules karam Polijai piederošo \strong{Viļņas pilsētu un apgabalu}, uz kuru visu starpkaru periodu aktīvi pretendēja Lietuva.

Latvijas konsuls Viļņā F.~Donass ziņojumā ārlietu ministram stāstīja par savu sarunu 1939.~gada 17.~septembrī ar turīgiem poļu zemniekiem Viļņas apgabalā. Viņi lamāja Polijas valdību, kurai, ja jau valsts neesot bijusi gatava karam, vajadzējis atdot Dancigu Vācijai. Zemnieki baidījās no vācu ienākšanas, ``jo tie rekvizējot zemkopības ražojumu daudz pamatīgāk un nežēlīgāk nekā lielinieki''. Zemnieki uzskatīja, ka ``vislabāk jau būtu, ja ienāktu lietuvieši, jo zem viņiem zemnieki varētu kārtīgi un bez bailēm ražot savus produktus un pārdot pilsētas tirgos.'' Ziņojums dod ieskatu arī par notikumiem pašā Viļņā, Sarkanajai armijai to ieņemot. 19.~septembrī pilsētā noritēja ielu cīņas starp Sarkanās armijas daļām un poļu patriotu pulciņiem, ar degmaisījumu pildītām pudelēm tika sadedzināti vairāki padomju tanki. Toties citur, pēc Latvijas konsula ziņojuma, ``bija salasījušies vietējie komunisti, viņu piekritēji un nedaudz ziņkārīgo, galvenā kārtā žīdi un šaubīga izskata personas, kopā ap pāris simtu cilvēku. Viņi saucieniem apsveica ienācējus\dots{}''. Blakus Latvijas konsulātam, ``cilvēki [konsuls diemžēl šajā gadījumā neapraksta cilvēku tautību un to, vai viņi bijuši mazāk, tikpat vai vairāk ``šaubīga izskata''~--- V.Š.] ``sāka laupīt produktus no kara noliktavām \citespace{} 2~ dienas nesa maisus ar miltiem un citiem produktiem.'' 20.~septembrī Viļņas ielās bija ``daudz strādnieku un vienkāršu cilvēku, kuri priecīgām sejām sarunājas ar Padomju Krievijas kareivjiem un virsniekiem. Turīgās aprindas nospiestas un bailīgas.''

Baltkrievijas vēsturiskajā literatūrā atzīmēts, ka Viļņas apgabalā un Baltkrievijas PSR dzīvojušie baltkrievi pēc Polijas sadalīšanas cerēja uz apgabala pievienošanu Baltkrievijai. BPSR presē tika drukāti raksti par Viļņas vēsturi un baltkrievu tautas tiesībām uz to, rādot, ka nozīmīgākie baltkrievu kultūras, literatūras, vēstures pieminekļi saistīti ar šo ``Lietuvas-Baltkrievijas valsts'' vēsturisko galvaspilsētu. Viļņas apgabala pilsētās un miestiņos notika mītiņi, kuros pieņēma rezolūcijas ar prasību pilsētu un apgabalu pievienot Baltkrievijai, pašā Viļņā noticis ``grandiozs'' mītiņš sakarā ar pilsētas pievienošanu BPSR, tajā piedalījušies 70~000 cilvēku jeb 1/3 pilsētas iedzīvotāju. Sāka iznākt avīze \betxti{``Віленская праўда''} (Viļņas Taisnība) baltkrievu valodā.

Taču baltkrievu cerības nepiepildījās. Neraugoties uz to, ka lietuviešu īpatsvars Viļņā sastādīja tikai dažus procentus, tās liktenis tika izlemts citādi. 1939.~gada 31.~oktobrī, uzstājoties PSRS Augstākās Padomes sēdē, V.~Molotovs paziņoja: ``PSRS nolēma nodot Viļņas pilsētu Lietuvai ne tāpēc, ka tur būtu lietuviešu iedzīvotāju pārsvars. Padomju valdība rēķinājās ar to, ka Viļņas pilsēta, kuru Polija ar spēku atrāva Lietuvai, ir jāpieder Lietuvai kā pilsētai, ar kuru, no vienas puses, ir saistīta vēsturiskā pagātne, bet, no otras puses, lai īstenotu lietuviešu cerības\dots{}'' 1939.~gada 10.~oktobrī PSRS parakstīja savstarpējās palīdzības līgumu ar Lietuvu, kas paredzēja ne tikai padomju armijas bāzu izvietošanu Lietuvā, bet arī Viļņas apgabala rietumu daļas nodošanu Lietuvai. Šī teritorija bija 6~909 km$^{2}$ liela, to apdzīvoja apmēram pusmiljons iedzīvotāju.

Pārējā Viļņas apgabala daļa tika iekļauta Baltkrievijas PSR. Pēc Lietuvas pievienošanas PSRS un Lietuvas PSR izveidošanās jaunizveidotajai padomju republikai 1940.~gada oktobrī papildus tika nodota nu jau Baltkrievijas teritorija 2637~km$^{2}$ platībā. Pēc Otrā pasaules kara robežas starp Lietuvas PSR un Baltkrievijas PSR tika vēlreiz precizētas, palielinot Lietuvai piešķirto Viļņas apgabala daļu. Šīs robežas saglabājušās arī mūsdienās.

Viļņas apgabala atdošana Lietuvai, tāpat kā Rietumbaltkrievijas un Rietumukrainas pievienošana Baltkrievijas un Ukrainas PSR, vismaz daļēji novērsa pasaulē iespaidu, ka PSRS rīkojas tikai savu lielvalstisko interešu vadīta.

Kā norāda mūsdienu baltkrievu publicisti, aiz ``labas gribas žesta'', apmainot ``baltkrievu Viļņu'' pret tiesībām izvietot Lietuvā 20~tūkstošu vīru karaspēku, slēpās ``PSRS lielvalstiskās intereses''. Bez tam Viļņas nodošana Lietuvai bijusi vērsta arī uz baltkrievu nacionālās kustības, kuras centrs esot bijusi Viļņa, apspiešanu. Komentējot teikto, jāatzīst, ka, protams, PSRS staļiniskā vadība Viļņu nodeva Lietuvai nevis taisnīguma principa, bet gan ``lielvalstisko interešu'' vadīta, taču to, ka šis solis bija vērsts pret baltkrievu nacionālo kustību, būtu grūti pierādīt tās vājuma dēļ. Tā nevarēja nopietni apdraudēt PSRS intereses.

Var teikt, ka Lietuvas teritorijas izmaiņas notika pretēji tās valdības gribai. Lietuvas valdošais slānis, kā rakstīja PSRS pilnvarotais Lietuvā V.~Semjonovs, ``bija gatavs atteikties no Viļņas, lai tikai nepieļautu Lietuvā uzturēties padomju garnizoniem''. Lietuviešu vidū skanēja viedoklis, ka boļševiki ``atdeva mums Viļņu, bet atņēma visu Lietuvu''. Tiesa, kā ziņoja Latvijā iznākošā avīze ``Brīvā Zeme'': ``11.~oktobrī visa Kauņa svētku rotā sakarā ar Lietuvas~--- PSRS līguma parakstīšanu. Iedzīvotāji nakti pavada pie radioaparātiem, gaidot ziņu par līguma parakstīšanu. Ielās pulcējas studenti un organizācijas, lai dotos uz manifestāciju vietām.'' Lietuvas propaganda nopelnus par Viļņas pievienošanu pierakstīja savam vietējam diktatoram A.~Smetonam, taču faktiski tie pienācās citam diktatoram~--- J.~Staļinam. Kopumā Lietuvas tauta apsveica Viļņas pievienošanu Lietuvai, bet bija neapmierināta ar padomju militāro bāzu ierīkošanu savā valstī.

1939.~gada 27.~oktobrī, baznīcu zvaniem zvanot, Lietuvas karaspēks iegāja Viļņas apgabalā. Taču tā karavīri izbaudīja īstu šoku~--- viņu valodu šeit gandrīz neviens nesaprata, daudzi viņus šeit negaidīja, bet uztvēra kā okupantus. 72~gadus vēlāk viens no Lietuvas armijas veterāniem (\lttxti{Jonas Vainavskas}) stāstīja: ``Kad mēs iegājām Viļņā, bija skumji. Mēs redzējām, ka neviens par to nepriecājās, kaut oficiāli arī tika apgalvots, ka mūs sagaidīja ziediem. Kam mēs tur bijām vajadzīgi.'' 28.~oktobrī notika oficiāla Lietuvas armijas sagaidīšanas ceremonija. Viļņas lietuviešu iedzīvotāji Katedrāles laukumā, kur bija uzbūvēta arī tribīne, un Ā.~Mickēviča ielā gaviļu saucieniem apsveica tajā ienākošo Lietuvas karaspēku. Taču, kā ziņoja Latvijas konsuls Viļņā, ``Tomēr ļaužu pūlis nebija tik liels un saucieni ne tik vispārīgi, kā agrākās poļu militārajās parādēs. Sevišķi to varēja novērot citās ielās, pa kurām gāja karaspēks. Tur pietrūka lietuviešu un poļi tikai noskatījās garām ejošā karaspēkā, bet nedemonstrēja nekādu skaļu sajūsmu. Poļu nostāju raksturo sekojošie vārdi, kādus dzirdēju vairākkārt: ``No trīs velniem (vāciešiem, lieliniekiem un leišiem) leiši vēl ir tie vislabākie''. Pēc Lietuvas karoga pacelšanas ceremonijas Gedimina kalnā runu teica ģenerālis V.~Vitkauskas. Viņš runāja par lietuviešu tautas prieku sakarā ar ``Viļņas pievienošanu''. Viņu gan saprata tikai viņa padotie karavīri, pilsētā ievestie lietuviešu policisti un daži pilsētnieki. Citās Lietuvas pilsētās sakarā ar Viļņas pievienošanu notika plašas demonstrācijas ar V.~Ļeņina, J.~Staļina, V.~Molotova, G.~Dimitrova portretiem.

Kā raksta poļu vēsturnieks P.~Madajčiks, poļu izcelsmes iedzīvotāji Viļņas reģionā, kuri šeit bijuši etniskais vairākums, tā nodošanu Lietuvai 1939.~gadā uztvēra kā okupāciju. Tiesa, vēsturnieks lieto poļiem izdevīgus iedzīvotāju skaitīšanu datus (Polijā inkorporētajā Viļņas reģionā lietuviešu izcelsmes iedzīvotāju bijis tikai vairāk nekā 10\%, Viļņas pilsētā 1916.~gadā lietuviešu bijis mazāk par 3\%, vēlāk to īpatsvars krities zem 1\%), ignorējot citus, kuri rāda arī ievērojamu baltkrievu, lietuviešu un ebreju klātbūtni. (Tā, pēc 1897.~gada vēl Krievijas impērijā veiktās tautas skaitīšanas Viļņas guberņā dzīvoja 56,1\% baltkrievu, 17,8\% lietuviešu un tikai 8,2\% poļu. Šī tautas skaitīšana, tāpat kā 1931.~gadā Polijā inkorporētajā Viļņas reģionā veiktā, kura bija balstīta uz iedzīvotāju ``valodu un reliģiju'' un atzina, ka reģionā dzīvo 59,7\% poļu, neapšaubāmi jāvērtē kā tendencioza, tāpat kā daudzi citi ieinteresēto personu nosauktie skaitļi.) P.~Madajčiks uzsver: ``Tā gan bija īslaicīga okupācija (no 1939.~gada 26.~oktobra līdz 1940.~gada 15.~jūnijam), bet tās laikā poļu sabiedrība ļoti vienoti nostājās pret Lietuvas varas institūcijām. Lietuvas rīcībā bija pārāk maz laika, lai mēģinātu dabūt savā pusē atsevišķas poļu grupas, iesaistot tās savas valsts celtniecībā'' (P.~Madajčiks poļu atbalstu Lietuvai sauc par kolaborāciju). Tādējādi poļu vēsturnieka nostāja ir skaidra: Lietuva veica okupāciju, poļi tai aktīvi un pasīvi pretojās. Arī cits poļu vēsturnieks D.~Roguts atzīst, ka vēl mūsdienās daži [poļu] vēsturnieki runā par Viļņas okupāciju no lietuviešu puses, kura sekojusi tās padomju okupācijai. Daudzās publikācijās un sabiedriskās diskusijās poļi ir pauduši līdzīgu nostāju: tā bija okupācija, no Lietuvas puses tas bija nepareizs un netaisnīgs solis, taču tā bija labāka par vācu vai padomju okupāciju un nav ar tām salīdzināma. Lai nu kā, kāds referendums, vismaz Tautas sapulces vēlēšanu formā, kā kaimiņos esošajās Rietumbaltkrievijā un Rietumukrainā, Viļņas apgabalā netika rīkots, kas tikai liecina, ka Lietuvas iestādes nevarēja cerēt uz apgabala iedzīvotāju nepārprotamu atbalstu tā pievienošanai Lietuvai. Poļu vēsturniekiem šī atšķirība acīmredzot neliekas kaut cik būtiska, citādi taču nāktos atzīt, ka Rietumukrainā un Rietumbaltkrievijā iedzīvotāju masveida atbalsts pievienošanai Ukrainas un Baltkrievijas PSRS pastāvēja.

Pēc aprakstītajiem notikumiem Viļņā sāka strauji augt lietuviešu iedzīvotāju īpatsvars. Notika tas, ieplūstot pilsētas rūpniecībā strādniekiem no citiem valsts apgabaliem, kā arī pārceļoties turp ierēdņiem, policistiem, garnizona karavīriem. 1939.~gada beigās Lietuvā poļu militāros bēgļus sāka ievietot koncentrācijas nometnēs, kuras gan oficiāli saucās par ``piespiedu darba nometnēm''. Desmiti tūkstoši poļu tika atlaisti no darba. 1940.~gada martā Lietuvas varas iestādes atņēma pilsonību visiem poļiem, kuri bija apmetušies Viļņa un tās apgabalā pēc 1920.~gada. Rezultātā ap 150~000 cilvēku zaudēja svarīgas politiskās un ekonomiskās tiesības. Tomēr drošības sajūtas paplašinātajā Lietuvā nebija. Periodā no Polijas sagrāves līdz 1940.~gadam Viļņā esot bijusi populāra anekdote: Kādu valodu mācās optimists? Lietuviešu valodu. Kādu valodu mācās pesimists? Poļu valodu. Un kādu valodu mācās reālists? Krievu valodu. Kā zināms, Lietuva, tāpat kā Latvija un Igaunija, 1940.~gadā kļuva par padomju republikām.

1940.~gada vasarā, kad PSRS anektēja Lietuvu, Viļņas apgabals tā arī palika Lietuvas PSR sastāvā. (Var piebilst, ka 1945.~gada februārī Sarkanā armija ieņēma arī 1939.~gada 23.~martā Vācijas okupēto Klaipēdas apgabalu, ko PSRS 1946.~gadā nodeva Lietuvas PSR.)

Rezultātā ap 48,5\% bijušās Polijas (bez Cešinas apgabala) teritorijas ar ap 20,1 miljoniem cilvēku nonāca nacistiskās Vācijas pārvaldē, ap 50\% teritorijas ar ap 13~miljoniem iedzīvotāju ieguva PSRS, bet 1,5\% teritorijas ar ap 0,5~miljoniem iedzīvotāju~--- Lietuva. Mūsdienās Lietuvas publicistikā var atrast PSRS nosodījumu par ``kaunpilno dalību Polijas sadalē''. Taču tīšuprāt tiek ``aizmirsts'', ka Lietuva ar Viļņas apgabala pievienošanu salīdzinājumā ar pirmskara laiku ievērojami paplašināja savu teritoriju un iedzīvotāju skaitu. Minētie publicisti tāpat ``aizmirst'', ka konsekventi pieņemot viņu pozīciju, Lietuvas valstij mūsdienās vajadzētu atdot Polijai kādreiz lietuviešu politiķu tik kāroto un no PSRS rokām saņemto Viļņas apgabalu. Lai nu kā, Lietuva, kura pēc PSRS sabrukuma atguva neatkarību, šodien var būt apmierināta ar vēstures gaitu. Ja Viļņa tās rokās nebūtu nonākusi jau 1939.~gadā, diezin vai pēc Otrā pasaules kara tās pretenzijas uz savu seno galvaspilsētu tiktu ņemtas vērā.

Tā 1939.~gada oktobrī toreizējo \strong{Polijas valsts teritoriju sadalīja Vācija, PSRS, Slovākija un Lietuva}. Mūsdienu poļu politiķi gan divas pēdējās valstis piemin ļoti nelabprāt (bieži to pilnībā noklusē), galveno savu neapmierinātību virzot pret Vāciju un PSRS un tās tiesisko mantinieci Krieviju. Taču patiesību nedrīkst noklusēt.

Jau 1939.~gadā Viļņā poļu iedzīvotāji izveidoja virkni nelegālu un puslegālu poļu organizāciju (savstarpējās palīdzības biedrību, biedrību palīdzībai kara upuriem u. tml.) 1939.~gada 1.~un 2.~novembrī Viļņā notika plašas poļu demonstrācijas pret Lietuvas varu. Policija tās izklīdināja, arestējot daudzus dalībniekus. Notikumi kalpoja tikai par ieganstu tālākai lietuviskošanai. 1.--4.~decembrī, atsaucoties baumām par poļu skolu slēgšanu, notika ``skolu streiks''. Poļu ģimnāzisti neapmeklēja nodarbības, pie skolām tika uzstādīti posteņi, kuri sūtīja mājup ``streiklaužus''. Līdz 8.~decembrim policija arestēja 174~streika dalībniekus, 10.~decembrī streiks tika izbeigts. Augstākais Viļņas lietuviskošanas punkts bija Stefana Batorija universitātes slēgšana 1939.~gada 15.~decembrī. Lietuvas drošības policija veica diezgan plašus arestus. Līdz Sarkanās armijas otrreizējai ieiešanai Viļņā 1940.~gada jūnijā tā apcietināja vairāk nekā 1~500 poļu. Kā atzīmējis poļu vēsturnieks P.~Ņivinskis, pēc Viļņas apgabala pārņemšanas Lietuvas pasākumi pret poliskumu, uzstāšanās pret poļu valodas lietošanu, pat dievkalpojumiem poļu valodā, noveda pie tā, ka ievērojama daļa Viļņas poļu iedzīvotāju tolaik uzskatīja lietuviešu ``okupāciju'' par ``briesmīgāko'' visā pilsētas vēsturē. Tas bija zināms paradokss, jo par kaut kādiem cilvēku upuriem poļi nevarēja žēloties. Acīmredzot poļu nacionālistiem nu jau Lietuvas Republikā bija nepanesami iejusties nacionālā mazākuma ādā.

\subsection{Vācu okupācijas režīms Polijas teritorijā. Pretošanās kustības sākums}

1939.~gada 6.~un 8.~oktobrī ar Ā.~Hitlera dekrētiem Vācijas iekarotās Polijas teritorijas tika sadalītas divās daļās: \strong{``iekļautajos [Vācijā] austrumu apgabalos''} (\detxti{Eingegliederten Ostgebiete}) un ``okupēto poļu apgabalu Ģenerālgubernatūrā'' (\detxti{Genaralgouvernement für die besetzten polnischen Gebiete}), ko no 1941.~gada 31.~jūlija vienkārši sauca par Ģenerālgubernatūru.

Autors uzskata, ka, neraugoties uz to, ka pēc 1939.~gada 28.~septembra PSRS un Vācijas ``Līguma par robežu un draudzību'' visa poļu apdzīvotā teritorija nonāca Vācijas rokās, tomēr ir likumsakarīgi runāt par Polijas sesto dalīšanu, jo Vācija daļu poļu zemju piedalīja savai teritorijai, bet no atlikušajām izveidoja t.s. Ģenerālgubernatūru. Tā kā abu daļu statuss bija atšķirīgs, par dalīšanu var runāt ar pilnām tiesībām.

Mūsdienu poļu vēsturnieks un politologs J.~Krols ir izdalījis četrus komponentus, no kuriem 1939./1940.~gadā sastāvēja nacistu propagandā tēlotais bijušās Polijas valsts un poļu tautas tēls:
\begin{enumerate*}[label=\arabic*)]
\item Polijas valsts bija mākslīgs Versaļas sistēmas veidojums, nespējīgs pastāvēt un traucējošs kārtību Eiropā;
\item Poļi ir bezatbildīga, nespējīga neko radīt nācija, kura ar savām pretenzijām apdraud Vāciju un kurai ir vajadzīgs valdītājs, kas to vadītu;
\item etniski-rasistiski poļu iedzīvotāju ir ``rasistiski mazvērtīgi'', ``pāržīdiskoti'', un iemieso sevī ``aziātisku garu'';
\item bioloģiski poļu pārliekais skaits apdraud vācu kaimiņus un viņu mazvērtīgā rakstura dēļ ir drastiski jāsamazina, jādeportē un to elite jāiznīcina.
\end{enumerate*}
Šos atzinumus, starp citu, tika mēģināts pamatot arī ar vēsturisku personību citātiem. Piemēram, Fridrihs II (Lielais) 1773.~gadā kādam savam laikabiedram bija rakstījis: ``Jūs vērtējat poļus tā, kā es vienmēr esmu tos uzskatījis: kā augstprātīgu, uzpūtīgu, pagrimstošu un zemāku rasi".

Ar pēc Pirmā pasaules kara tai atsavināto teritoriju atgūšanas ieganstu Vācija tieši pievienoja sev ap 1/3 bijušās Polijas (jeb ap 1/2 no tieši Vācijas iekarotās Polijas teritorijas). Vācijas sastāvā atgriezās bijušās Prūsijas zemes Danciga, Pozene, Rietumprūsija un ievērojama Silēzijas daļa. Taču Vācijai tika pievienotas arī Pomorjes, Augšsilēzijas vojevodistes, liela daļa Lodzas vojevodistes, rietumu daļa no Varšavas un Krakovas vojevodistēm, Mazūrija, Suvalku apriņķis, daži Kelces vojevodistes apriņķi un viss Dombrovas baseins. Ar to Polijas ekonomiski attīstītākie apgabali kļuva par Vācijas reiha sastāvdaļu. Šī Vācijas tieši anektētā Polijas teritorija~--- ap 92~000~km$^{2}$~--- bija apmēram 2~reizes lielāka par to, ko Vācija zaudēja 1918.--1922.~gadā. Šais teritorijās dzīvoja ap 9,5~miljoniem iedzīvotāju, no tiem 87\% poļu, 6,5\% vāciešu un 6,4\% ebreju. (Citi dati: 10,1~miljons cilvēku, no tiem 8,9~miljoni poļu, 600~tūkstošu (5,9\%) vāciešu, 603~tūkstošu ebreju, 11~tūkstošu ukraiņu un 21~tūkstotis citu tautību pārstāvju). Pirmskara Polijā šeit tika saražoti 100\% akmeņogļu, 97,5\% čuguna, 90\% tērauda un 70\% cukura. Šeit bija koncentrēti 80\% poļu rūpniecības uzņēmumu, arī labības raža bija krietni augstāka nekā valstī vidējā.

1939.~gada 26.~oktobrī notika Vācijas okupācijas zonas administratīva sadale. To veicot, nacionālsociālistiskā Vācija neievēroja nekādas agrākās tradīcijas, piemēram, Poznaņas apgabala saistību ar Prūsiju pirms 1918.~gada.

Vācijā iekļautie Polijas apgabali tika sadalīti četrās daļās.

\begin{enumerate}
\item Suvalku apgabals un Mazovijas ziemeļi (no agrākās Varšavas vojevodistes daļas tika izveidota Cihanuvas (poļu \pltxti{Ciechanów}, vācu \detxti{Zichenau}) administratīvā vienība (\detxti{Regierungsbezirk}) tika piedalīti Austrumprūsijai (\detxti{Reichsgau Ostpreußen}). Kopā tajā bija apmēram 16~200~km$^{2}$ ar 1~miljonu iedzīvotāju;

\item Danciga, agrākais ``poļu koridors'', Pomožes vojevodiste tika pievienota Dancigas-Rietumprūsijas valstsnovadam (\detxti{Reichsgau Danzig-Westpreuβen}). Apmēram 21~200~km$^{2}$ platībā šeit dzīvoja ap 2,15~miljonu iedzīvotāju, to skaitā ap 800~000~vāciešu un 120~000~kašubu (kašubu \emph{Kaszëbi}; poļu \pltxti{Kaszubi}, vācu \detxti{Kaschuben}~--- rietumslāvu etniska grupa, dzīvojoša Polijas ziemeļos);

\item Augšsilēzijās rūpniecības rajonu ar Katovici pievienoja Silēzijai (\detxti{Reichsgau Ober-Schlesien}). Tajā bija apmēram 10~600~km$^{2}$ ar 2,6~miljoniem iedzīvotāju. Dažus bijušos poļu apriņķus tieši pievienoja Opoles iecirknim (vācu \detxti{Landkreis Oppeln});

\item Ja iepriekš nosauktajos pārveidojumos bija vērojama mazāku agrāko Polijas teritoriju pievienošana lielākiem Vecvācijas (\detxti{Altreich}) apgabaliem, tad no lielākās Poznaņas provinces daļas kopā ar rietumu Viduspoliju ar Lodzas pilsētu un tās rūpniecības apgabalu tika izveidots ``valsts novads Poznaņa'' (\detxti{Reihsgau Pozen}). Ar 44~000~km$^{2}$ lielu teritorija un ap 4,5~miljoniem iedzīvotāju tas tālu pārsniedza agrāko Vācijas Pozenes provinci 1914.~gadā. 1940.~gada februārī to pārdēvēja par Vartas novadu (\detxti{Reichsgau Wartheland} vai \detxti{Warthegau}). Iedzīvotāji novadā sākotnēji sastāvēja no 85\% (3,6~miljonu) poļu, 8\% (360~000) ebreju, un 7\% (327~000) vāciešu. Vartas novadam bija jākļūst par Trešā reiha paraugnovadu, sava veida nacionālsociālistisku paraugteritoriju, kur eksperimentāli jāizveido situācija, kuru varētu tālāk izplatīt uz visu Vāciju. Tāpēc šeit poļu iedzīvotāju padzīšanai un ģermanizācijai bija jānotiek ar sevišķu sparu.
\end{enumerate}

Vācijā ``iekļautie austrumu apgabali'' tika pievienoti valsts saimnieciskajai un finanšu telpai, bet pasu un policijas lietās joprojām palika ``ārzemes''. Šiem apgabaliem 10~gadu laikā bija jātiek ``atpoļotiem'' (\detxti{entpoloniziert}).

Ar Polijas likvidēšanu no dienas kārtības automātiski pazuda Cešinas problēma. 1939.~gada oktobra sākumā, kad pacēlās jautājums par Cešinas reģiona likteni, Bohēmijas un Morāvijas protektorāts (vācu \detxti{Protektorat Böhmen und Mähren}, čehu \cstxti{Protektorát Čechy a Morava}) ar tā protektora K.~f.~Neirāta atbalstu izvirzīja uz to savas pretenzijas. Taču Ā.~Hitlers pēc Ārlietu ministrijas un H.~Himlera ieteikuma daļu reģiona piešķīra Slovākijai (Ā.~Hitlers vēlējās Slovākijai piešķirt arī Polijas Zakopanes kūrortapgabalu, taču Slovākijas prezidents J.~Tiso bija tik saprātīgs, ka no šādas apstrīdamas ``dāvanas'' atteicās), daļu Cešinas atstājot okupēto poļu zemju Ģenerālgubernatūras sastāvā. Tomēr kara gados Ā.~Hitlers tā arī nepieņēma galīgo lēmumu par Cešinas likteni.

Nacionālsociālistu plāni paredzēja paaugstināt vāciešu skaitu reiham pievienotajos Polijas apgabalos, pārvietojot turp vācu iedzīvotājus no citām Austrumeiropas valstīm~--- no Baltijas līdz Besarābijai. 1939.~gada 6.~septembrī (tūlīt pēc tam, kā vērmahts bija okupējis Rietumpoliju) Ā.~Hitlers reihstāgā teica kaismīgu runu: starp Eiropas tautām jābūt stingrām un skaidrām robežām, tāpēc Vācijai ir jāpārceļ uz jaunām dzīves vietām simtus tūkstošus etnisko vāciešu. Taču reizē nācās konstatēt, ka iekarotā ``dzīves telpa'' nav tukša, bet ir daudzu cilvēku dzimtene. Šo telpu vajadzēja ``attīrīt'' vācu tautas vajadzībām. Poļus no viņu dzimtenes nacisti izraidīja galvenokārt rasistisku apsvērumu vadīti. Tika uzskatīts, ka ``sveštautiešu'' iedzīvotāji var atstāt kaitīgu ietekmi uz Vācijas valsti un vācu tautas kopību. Jau sākotnēji šī ``attīrīšana'' bija saistīta ar visbriesmīgākajiem noziegumiem. Tā, vācu vēsturnieks V.~Rīss konstatējis, ka no 1939.~gada septembra līdz 1940.~gada pavasarim Vartas novadā un Rietumprūsijā t.s. ``eitanāzijas'' programmas ietvaros nošāva vai ar gāzi nogalināja vairāk nekā 7~700 psihiatrisko slimnīcu iemītnieku, kurus okupanti uzskatīja par ``liekēžiem'' (\detxti{überflüssige Esser}).

No savām mājām padzītajiem, par vācu kalpiem padarītajiem vietējiem iedzīvotājiem nodarītā netaisnība, pret tiem vērstā vardarbība tika attaisnota kā nāciju cīņa (\detxti{Volkstumskampf}).

Pielietojot seno iekarotāju principu~--- ``skaldi un valdi'', okupanti vēlējās sadalīt to okupēto zemju iedzīvotājus nošķirtās, savstarpēji naidīgās nacionālās grupās. Ar 1939.~gada 7.~oktobra Ā.~Hitlera dekrētu SS reihsfīrers H.~Himlers tika iecelts par ``Valsts komisāru vācu tautības stiprināšanai'' (\__txti{Reichskomisar für die Festugung deutschen Volkstums}). Pēc fīrera rīkojuma viņa uzdevumi bija:

\begin{enumerate}
\item īstenot visu valstsvāciešu un etnisko vāciešu ātrāku repatriāciju uz reihu;

\item dažādu vācu nācijai bīstamu ietekmju izskaušana;

\item jaunu vācu tautas izmetināšanas areālu, pirmkārt jau dienvidaustrumu Eiropā, veidošana ar repatriēto vāciešu pārvietošanu turp.
\end{enumerate}

Tika radīts speciāls dienests vācu tautības stiprināšanai (``\detxti{Dienststelle des Reichskommissars für die Festigung deutschen Volkstums}'').

Tūkstošiem etnisko vāciešu bija paredzēts izmitināt ``jaunajās austrumu teritorijās''~--- jaunveidotajos Dancigas-Rietumprūsijas un Vartas valstsnovados, pamatos Polijas teritorijā. Ar tām Austrumeiropas valstīm, no kurām bija paredzēta etnisko vāciešu pārceļošana, līdz Vācijas uzbrukuma sākuma Padomju Savienībai tika noslēgti attiecīgi līgumi. Atbilstoši šiem dokumentiem ap 360 tūkstošu etnisko vāciešu no Austrumeiropas valstīm tika pārcelti jau līdz 1941.~gada pavasarim. Pavisam kara gados tika pārvietoti vismaz 650~000 etnisko vāciešu. (Pēc oficiāliem H.~Himlera vadītā dienesta vācu tautības stiprināšanai datiem līdz 1944.~gada janvārim bija pārcelti 770~577 etniskie vācieši, bet līdz gada beigām, sakarā ar vērmahta atkāpšanos Sarkanās armijas priekšā, to skaits pieauga līdz pat miljonam. Tomēr jāievēro, ka vairāki desmiti tūkstošu citu valstu pilsoņu~--- nevāciešu tika nepamatoti, kaut lielākoties pēc pašu vēlēšanās, pieskaitīti etniskajiem vāciešiem un iekļauti pārvietojamo sarakstos.)

No mūsdienu viedokļa raugoties, šī vācu mazākuma no Ģenerālgubernatūras (īpaši Helmas novada), no Baltijas, padomju rokās nonākušajiem bijušajiem Polijas austrumu apgabaliem un Rumānijas (Besarābijas un Bukovīnas) gatavība pārcelties no gadsimtiem viņu apdzīvotām teritorijām uz pilnīgi nepazīstamiem apgabaliem ir grūti izskaidrojama. Taču tikai formāli pārcelšanās notika brīvprātīgi. Faktiski daļa šo etnisko vāciešu tikai ar dažādu represiju draudiem tika piespiesta pārcelties uz Vāciju. Tā, H.~Himlera 1940.~gada 19.~augustā izdotajā direktīvā tika atzīmēts, ka tie, kuri nepiekritīs uzreiz pārcelties, vēlāk tiks uzskatīti par bēgļiem bez kādām vāciešu privilēģijām. Jāatceras arī, ka daļa etnisko vāciešu pēc 1939.~gada 23.~augustā notikušās Austrumeiropas sadalīšanas Vācijas un PSRS interešu sfērās guva pirmo negatīvo dzīves pieredzi staļiniskā režīma apstākļos, tāpēc pretošanās gadījumā nacistu nodomiem baidījās zaudēt Vācijas atbalstu, kā arī cerēja Vācijas teritorijā nodrošināt sev saimnieciskās un kulturālās attīstības perspektīvas.

Teritorijas, kurās bija paredzēts izmetināt vāciešu, nacistiem bija jāattīra no vietējiem iedzīvotājiem: daļēji, ja runa gāja par slāvu iedzīvotājiem, vai pilnībā, ja par ebrejiem. Otrajā gadījumā var runāt par genocīdu, tieši Polijas ebreji no Vartas apgabala bija vieni no pirmajiem, kurus nosūtīja uz nāves nometnēm.

Vācu politikas pamatus Polijā noteica rasu teorijas principi, kurus nacistu vadoņi tulkoja pēc saviem priekšstatiem. Polija bija vienīgā vieta Eiropā, kur nacisti savā cīņā par ``dzīves telpu'' (\detxti{Lebensraum}) paguva pilnā apjomā visiem iedzīvotājiem piemērot savu rasu politiku.

Poļus kā slāvus, tāpat kā ģermāņus, romāņus, baltus u.c. Eiropas tautas, varēja pieskaitīt āriešiem (XX~gadsimta sākumā literatūrā tā sauca indoeiropiešus~--- tautu kopumu, kas runā indoeiropiešu valodās). Taču hitlerieši to neatzina. 1939.~gada 28.~septembrī Ā.~Hitlers sarunā ar nacistu ideologu un ``speciālistu'' rases jautājumos A.~Rozenbergu izteicās, ka poļi, izņemot ``plānu ģermāņu slāni'', sastāv no ``briesmīga materiāla'' (\detxti{furchtbaren Material}), bet poļu ebreji ir ``tas baismīgākais, ko vispār var iedomāties'' (\detxti{das grauenhafteste, was man sich überhaupt vorstellen kann}).

Ā.~Hitlers poļus arī raksturoja ar vārdiem ``kaitīgums'' (\detxti{Schädlichkeit}) un ``nepilnvērtība'' (\detxti{Minderwertigkeit}). {Ā.~Hitlera pārdomas par poļiem pēc Polijas sakāves fiksēja viņa propagandas šefs J.~Gēbelss: ``Fīrera spriedums par poļiem ir iznīcinošs. Drīzāk dzīvnieki nekā cilvēciskas būtnes, pilnīgi primitīvi un bezveidīgi. Un valdošā šķira~--- nožēlojams rezultāts, kas tapis sajaucoties zemākajiem slāņiem ar ārisko kungu rasi. Poļi ir neiedomājami netīrīgi. Viņu spēja saprātīgi spriest līdzinās nullei. Fīrers negrasās asimilēt poļus. \citespace{} Ja Heinrihs Lauva (Velfu dinastijas Saksijas hercogs, kurš karoja par zemēm uz austrumiem no Elbas) būtu iekarojis Austrumus \citespace{} rezultātā būtu tapusi stipri slāviskota vācu jaukteņu rase. Labāk tā, kā ir tagad. Tagad mēs zinām rasu iedzimtības likumus un varam attiecīgi rīkoties.''

Iznīcinošie poļu vērtējumi tika vērsti pirmkārt pret poļu augstāko slāni, kurš bija pēc iespējas pilnīgi jālikvidē. Savā ģermanizācijas politikā okupanti asi vērsās pret poļu kultūru un valodu, centās fiziski likvidēt poļu inteliģenci, uzskatot, ka, kamēr pastāvēs poļu inteliģence, būs dzīva arī poļu tautas dvēsele. 1939.~gada 24.~oktobrī Vācijas Propagandas ministrija izdeva direktīvu: ``Pat pēdējai govju meitai Vācijā ir jāieskaidro, ka poļi ir zemcilvēki. Poļi, žīdi un čigāni stāv uz tās pašas nevērtīgo cilvēku pakāpes. \citespace{} Šai nostājai vienmēr ir jāskan kā galvenajai vadlīnijai \citespace{} līdz Vācijā katrs ikvienu poli, vienalga, vai tas ir laukstrādnieks vai intelektuālis, jau zemapziņā uzlūkos kā parazītu.'' Tiesa, sadzīvē vācieši ebrejus un čigānus vērtēja vēl zemāk par poļiem.

Attiecībā par poļu masām pamatnostādnes deva nacistu partijas Rasu politikas pārvalde (``\detxti{Rassenpolitischen Amt}'') Berlīnē dokumentā ar nosaukumu ``Par apiešanos ar bijušo Polijas teritoriju iedzīvotājiem no rasu politikas viedokļa'' (\detxti{Über den Umgang mit den Bewohnern der ehemals polnischen Gebiete aus einem rassenpolitischen Standpunkt}). 25.~novembrī pārvalde izdeva priekšrakstu, kurā postulēja, ka poļi būtiski atšķiras no vācu tautas. Viņi satur tikai nedaudz ``nordisko asiņu'', tāpēc asimilācijai un ``ievācošanai'' (\detxti{Eindeutschung}) pēc pamatīgas un stingras pārbaudes ir derīgi tikai nedaudzi. Poļus bija jāģermanizē, jāpakļauj verdzībai vai arī jāiznīcina, atkarībā no tā, vai viņi dzīvoja Vācijai tieši pievienotājos apgabalos vai ārpus tiem~--- Ģenerālgubernatūrā. Tiem bija jāizzūd kā nācijai. SS reihsfīrers H.~Himlers 1941.~gada decembrī noteica, ka turpmāk poļu tauta vairs nedrīkst tikt apzīmēta kā ``pēc dabas radniecīga'' (\detxti{Artverwandt}), bet pret to jālieto termini ``neģermāņu'' (\detxti{nichtgermanisch}) un ``sveštautiešu'' (\detxti{fremdvolkisch}). Kādā nacistu partijas Rases politikas departamenta publikācijā 1941.~gadā bija pateikts: ``Rasistiski orientētas politikas pret svešiniekiem tālākais mērķis var būt tikai poļu lielākās daļas ``aizgādāšana projām (\detxti{Entfernung}) no valsts''. ``Aizgādāšana projām'' nozīmēja piespiedu deportāciju.

Vismasveidīgākās vardarbīgās masu deportācijas Eiropā norisa Polijas teritorijā. Jau kara pret Poliju laikā hitlerieši izstrādāja poļu izsūtīšanas plānus no Polijas rietumu vojevodistēm, kur dzīvoja līdz 7~miljoniem poļu jeb līdz 70\% no to iedzīvotājiem. Bija paredzēts ap 1~miljonu cilvēku izsūtīt līdz 1940.~gada februārim. Tiesa, vadoties no mainīgās starptautiskās un iekšējās situācijas, šos plānus nācās koriģēt.

Jau 1939.~gada beigās no Vācijas anektētājām teritorijām sākās \strong{poļu izraidīšana uz Ģenerālgubernatūru}. Īpaši plašus mērogus tā sasniedza Vartas novadā, kur vietējais gauleitars A.~Greizers asi iestājās pret poļu ģermanizācijas mēģinājumiem, uzskatot poļus par tai absolūti nederīgiem. (Tomēr daļa poļu, kura atbilda nacistu rasistiskajiem standartiem, tika izvesti uz Vāciju ģermanizācijai. Vartas novadā tādu izrādījās 31~423~cilvēks un Silēzijā~--- 14~tūkstošu cilvēku.) No 1939.~gada 1.~līdz 17.~decembrim no Vartas apgabala 41~pilsētas un apriņķa uz Ģenerālgubernatūru tika pārvietoti 87~883~cilvēki. Otrais deportēšanas etaps ilga no 1940.~gada 15.~marta līdz 1941.~gada 20.~janvārim. Tika pārvietoti 133~506~cilvēki~--- daļēji uz Ģenerālgubernatūru, daļēji~--- piespiedu darbā uz Vāciju. Trešajā etapā: 1941.~gadā un 1942.~gada janvāra sakumā bija paredzēts deportēt 330~000 poļu. Taču izdevās pārvietot tikai 130~826 cilvēkus. Līdz Vācijas uzbrukumam PSRS deportācijas norisa pilnā sparā, pēc tam transportlīdzekļu trūkuma dēļ masu deportācijas tika apturētas. Bez tam no Vartas novada 23~573 poļu izveda uz Franciju, lai izmantotu tur lauksaimniecībā.

% page 877



% ----- INFO PAGE -----
\newpage\thispagestyle{plain}
{\centering

% ----- LICENSE SIGNS BLOCK -----
\includegraphics[width=2.5em]{cc.pdf}~\includegraphics[width=2.5em]{by.pdf}~\includegraphics[width=2.5em]{nd.pdf}
% ----- LICENSE SIGNS BLOCK -----

\vspace{0.2em}

% ----- LICENSE BLOCK -----
{Šī grāmata tiek izplatīta saskaņā ar licenci \textbf{CC~BY~ND~4.0}.\par}
{Sīkāka informācija: \href{https://creativecommons.org/licenses/by-nd/4.0/deed.lv}{creativecommons.org/licenses/by/4.0}\par}
% ----- LICENSE BLOCK -----

\vfill

% ----- AUTHOR AND TITLE -----
{\large{Vitālijs~Šalda}\par}
\vspace{0.5em}
{\Huge\textbf{POLIJA}\par}
{\Large\textbf{XIX un XX gadsimtā}\par}
% ----- AUTHOR AND TITLE -----

\vfill

% ----- BOOK INFO -----
{Pirmais izdevums\par}
{Grāmatas informācija: \href{https://github.com/regnveig/polija}{github.com/regnveig/polija}\par}
% ----- BOOK INFO -----

\vfill

% ----- DOCUMENT INFO -----
{Rediģēšana un korektūra: \textit{Ņikita Skabcovs}\par}
{Darbvirsmas publicēšana: \textit{Emils Viesnā}\par}
{Vāka dizains: \textit{Emils Viesnā}\par}
{Izveidots ar \XeLaTeX\par}
% ----- DOCUMENT INFO -----

\vfill

% ----- PUBLISHER INFO -----
{Neatkarīgs izdevējs \textsc{Tsunami}\par}
{Rīga 2023\par}
% ----- PUBLISHER INFO -----

\vfill

% ----- QR CODES -----
~
% ----- QR CODES -----

% }
% % ----- INFO PAGE -----
%
% % ----- EMPTY PAGE -----
% \thispagestyle{plain}~
% % ----- EMPTY PAGE -----
%
% % ----- BACK COVER -----
% \includepdf[pages={1}]{cover_back.pdf}
% % ----- BACK COVER -----

\end{document}
